% !TeX program = xelatex
% ================================================================
%  《收获与播种》(Récoltes et Semailles) 中文译本 · 全书合订
%  收原书 p.23-929 全部译文, 依次为:
%    ① 前奏 "主题的呈现——或四个乐章的前奏" (原书 p.23-122, 第 1-4 章)
%    ② 第一部分 "愚妄与更新"           (原书 p.123-254, 第 5-12 章)
%    ③ 第二部 "葬礼(I)或中国皇帝的新衣" (原书 p.255-412, 第 13-16 章)
%    ④ 第三部 "葬礼(II)或阴阳之钥"      (原书 p.413-632, 第 17-18 章 18.1-18.2)
%    ⑤ 第四部 "葬礼(III)或四种运算"     (原书 p.633 扉起, 18.3 至全书末行 p.929)
%  另有五卷分册: vol-prelude / vol-part1 / vol-part2 / vol-part3 / vol-part4。
%  --------------------------------------------------------------
%  基类: mathtranslation.cls v3.2 (不改); 样式与宏见 common.tex,
%  封面见 cover-art.tex(与五卷同一套画, 只换配色与卷次文字)。
%  构建: xelatex ×4 (710 → 720 页, 目录页码第 4 遍才收敛)
% ================================================================
\PassOptionsToPackage{AutoFallBack}{xeCJK}
\documentclass[
    PaperSize=a4paper,
    HeadingFont=hei,
    Biblatex=false,
    BookTitleCN={收获与播种},
    Translator={数译},
    LogoFile=logo.pdf
]{mathtranslation}

% ================================================================
%  《收获与播种》中文译本 —— 各卷共用导言区
%  由 vol-prelude / vol-part1 / vol-part2 / vol-part3 / vol-part4
%  以及全书合订 main-all 共同 \input。基类 mathtranslation.cls v3.2 不改。
%  --------------------------------------------------------------
%  样式均在此改写, 以贴近原书:
%    · 章标题仿原书 "5. Travail et découverte"(块状左对齐黑体)
%    · 节/小节: 黑体左对齐, 编号后带句点
%    · 脚注: 人工编号照录原书——前奏与第一部分页内重新编号,
%      第二至四部各自连续编号并带(*)角标; 第一部分另有 \fnoteP
%      处理 "(n)" 式连续注号
%    · 边码: 原打字稿页码 \origpage{N}(前奏作 A1、P1、L1、罗马数字等)
%    · ⋄ 记号: 段首补写标记 \pmark 与上标形式 \dmark
% ================================================================

% 题名含空格, 不能走类选项(会被 kvoptions 吃掉), 在此赋值
\renewcommand{\BookTitleEN}{Récoltes et Semailles}
\renewcommand{\OriginalAuthor}{Alexandre Grothendieck}
\renewcommand{\TranslationDate}{\today}

% --------------------------------------------------------------
%  罕用字后备(如"祐"): Fandol 不载时退到 SimSun
% --------------------------------------------------------------
\setCJKfallbackfamilyfont{\CJKrmdefault}{{SimSun}}

% --------------------------------------------------------------
%  hyperref 元信息(各卷在 % ================================================================
%  《收获与播种》中文译本 —— 各卷共用导言区
%  由 vol-prelude / vol-part1 / vol-part2 / vol-part3 / vol-part4
%  以及全书合订 main-all 共同 \input。基类 mathtranslation.cls v3.2 不改。
%  --------------------------------------------------------------
%  样式均在此改写, 以贴近原书:
%    · 章标题仿原书 "5. Travail et découverte"(块状左对齐黑体)
%    · 节/小节: 黑体左对齐, 编号后带句点
%    · 脚注: 人工编号照录原书——前奏与第一部分页内重新编号,
%      第二至四部各自连续编号并带(*)角标; 第一部分另有 \fnoteP
%      处理 "(n)" 式连续注号
%    · 边码: 原打字稿页码 \origpage{N}(前奏作 A1、P1、L1、罗马数字等)
%    · ⋄ 记号: 段首补写标记 \pmark 与上标形式 \dmark
% ================================================================

% 题名含空格, 不能走类选项(会被 kvoptions 吃掉), 在此赋值
\renewcommand{\BookTitleEN}{Récoltes et Semailles}
\renewcommand{\OriginalAuthor}{Alexandre Grothendieck}
\renewcommand{\TranslationDate}{\today}

% --------------------------------------------------------------
%  罕用字后备(如"祐"): Fandol 不载时退到 SimSun
% --------------------------------------------------------------
\setCJKfallbackfamilyfont{\CJKrmdefault}{{SimSun}}

% --------------------------------------------------------------
%  hyperref 元信息(各卷在 % ================================================================
%  《收获与播种》中文译本 —— 各卷共用导言区
%  由 vol-prelude / vol-part1 / vol-part2 / vol-part3 / vol-part4
%  以及全书合订 main-all 共同 \input。基类 mathtranslation.cls v3.2 不改。
%  --------------------------------------------------------------
%  样式均在此改写, 以贴近原书:
%    · 章标题仿原书 "5. Travail et découverte"(块状左对齐黑体)
%    · 节/小节: 黑体左对齐, 编号后带句点
%    · 脚注: 人工编号照录原书——前奏与第一部分页内重新编号,
%      第二至四部各自连续编号并带(*)角标; 第一部分另有 \fnoteP
%      处理 "(n)" 式连续注号
%    · 边码: 原打字稿页码 \origpage{N}(前奏作 A1、P1、L1、罗马数字等)
%    · ⋄ 记号: 段首补写标记 \pmark 与上标形式 \dmark
% ================================================================

% 题名含空格, 不能走类选项(会被 kvoptions 吃掉), 在此赋值
\renewcommand{\BookTitleEN}{Récoltes et Semailles}
\renewcommand{\OriginalAuthor}{Alexandre Grothendieck}
\renewcommand{\TranslationDate}{\today}

% --------------------------------------------------------------
%  罕用字后备(如"祐"): Fandol 不载时退到 SimSun
% --------------------------------------------------------------
\setCJKfallbackfamilyfont{\CJKrmdefault}{{SimSun}}

% --------------------------------------------------------------
%  hyperref 元信息(各卷在 \input{common} 之后可用 \hypersetup 覆写)
% --------------------------------------------------------------
\hypersetup{
  pdftitle={收获与播种},
  pdfauthor={Alexandre Grothendieck / 数译 译},
  pdfsubject={Récoltes et Semailles — Réflexions et témoignage sur un passé de mathématicien}
}

% --------------------------------------------------------------
%  标题样式: 仿原书 —— 黑体左对齐, 编号后带句点
% --------------------------------------------------------------
\titleformat{\section}[block]
  {\normalfont\Large\headfont}{\thesection.}{0.5em}{}
\titlespacing*{\section}{0pt}{3.2ex plus 1.2ex minus .2ex}{1.6ex plus .2ex}
\titleformat{\subsection}[block]
  {\normalfont\large\headfont}{\thesubsection.}{0.5em}{}
\titlespacing*{\subsection}{0pt}{2.6ex plus 1ex minus .2ex}{1.2ex plus .2ex}
\titleformat{\subsubsection}[block]
  {\normalfont\normalsize\headfont}{\thesubsubsection.}{0.5em}{}
\titlespacing*{\subsubsection}{0pt}{2.2ex plus .8ex minus .2ex}{1ex plus .2ex}
\titleformat{name=\subsubsection,numberless}[block]
  {\normalfont\normalsize\headfont}{}{0em}{}
\setcounter{secnumdepth}{3}
\setcounter{tocdepth}{3}
% 为含长外文书名与行内公式的段落提供必要的断行余量
\setlength{\emergencystretch}{3em}

% 章标题样式: 原书章标题与节同级排印("13. A) HERITAGE..."),
% 故章标题取块状左对齐黑体, 字号介于节与部之间
\titleformat{\chapter}[block]
  {\normalfont\LARGE\headfont}{\thechapter.}{0.5em}{}
\titlespacing*{\chapter}{0pt}{3.6ex plus 1.4ex minus .2ex}{2ex plus .2ex}
% 目录中章号照原书用阿拉伯数字("13. A) ..." 而非 "第十三章")
\renewcommand{\CTEXthechapter}{\thechapter\hspace{.3em}}

% --------------------------------------------------------------
%  目录观感: 条目与引导点之间、引导点与页码之间各留 0.5em
%  内核 \@dottedtocline 让 \leaders 从标题末字起排, 第一个点几乎
%  贴着末字, 读起来像多标了一个句点; 页码前也没有余量。两处各补
%  0.5em。
% --------------------------------------------------------------
\makeatletter
\patchcmd{\@dottedtocline}{\nobreak\leaders}{\nobreak\hspace{0.5em}\leaders}%
  {\typeout{MT: TOC 引导点前已留 0.5em}}%
  {\typeout{MT-WARN: TOC 引导点前间距补丁未生效}}
\patchcmd{\@dottedtocline}{\hfill\nobreak\hb@xt@\@pnumwidth}%
  {\hfill\nobreak\hspace{0.5em}\hb@xt@\@pnumwidth}%
  {\typeout{MT: TOC 页码前已留 0.5em}}%
  {\typeout{MT-WARN: TOC 页码前间距补丁未生效}}
\makeatother

% --------------------------------------------------------------
%  正文半角标点后的间距
%  本书体例用半角 , ; : ? ! 与 ASCII 括号, 由 xeCJK 按 HalfRight
%  (右半留白) 处理。放宽到 0.45em, 与全角标点右半留白大致齐平。
% --------------------------------------------------------------
\xeCJKsetup{CJKecglue={\hskip 0.45em plus 0.06em minus 0.09em}}

% --------------------------------------------------------------
%  页眉页脚: 偶数页眉右对齐斜体节题; 页码页脚外侧
% --------------------------------------------------------------
\fancyhf{}
\fancyhead[RE]{\small\itshape\rightmark}
\fancyfoot[LE,RO]{\small\thepage}
\renewcommand{\headrulewidth}{0pt}
\renewcommand{\sectionmark}[1]{\markright{\thesection.\ #1}}
\fancypagestyle{plain}{%
  \fancyhf{}%
  \fancyfoot[LE,RO]{\small \thepage}%
  \renewcommand{\headrulewidth}{0pt}%
}

% --------------------------------------------------------------
%  原书版式记号宏
% --------------------------------------------------------------
% 边码: 原打字稿页码(G. 手稿系统里的自引页码), 排在页边
\usepackage{marginnote}
\newcommand{\origpage}[1]{\marginnote{\normalfont\scriptsize p.~#1}}

% ⋄ 记号: 段首/行内正常大小; 上标形式为作者后补文字的位置标记
\DeclareRobustCommand{\pmark}{\ensuremath{\diamond}\hspace{0.15em}}
\DeclareRobustCommand{\dmark}{\textsuperscript{\ensuremath{\diamond}}}

% 脚注人工编号: \fnote{795}{(*)}{文本} —— 编号与(*)角标均照录原书
\newcommand{\fnote}[3]{%
  \begingroup
  \renewcommand{\thefootnote}{#1#2}%
  \footnote{#3}%
  \endgroup
}

% 第一部分的脚注: 页内序号作上标(#1), 注文行首印 "序号(连续注号)",
% 正文在序号之后另行括注同一连续注号 —— \fnoteP{1}{(3)}{注文}
% 原书有三条注无连续注号而改标日期(ch07 注3 = 30 Septembre,
% ch09 注11/12 = 8 août), 第二参数留空, 日期写进注文之首。
\newcommand{\fnoteP}[3]{%
  \begingroup
  \renewcommand{\thefootnote}{#1}\footnotemark
  \renewcommand{\thefootnote}{#1#2}\footnotetext{#3}%
  \endgroup
}

% 正文呼号与注文行首标号不同(原书照录)
\newcommand{\fnotevariant}[4]{%
  \begingroup
  \renewcommand{\thefootnote}{#1#2}\footnotemark
  \renewcommand{\thefootnote}{#1#3}\footnotetext{#4}%
  \endgroup
}

% 大注块标题: \notehead{171³}{斜体评论} —— 正文随后另起
\newcommand{\notehead}[2]{%
  \par\bigskip
  \noindent{\headfont 注~#1}\ifblank{#2}{}{\hspace{0.5em}{\itshape #2}}%
  \par\nobreak\smallskip\nobreak
}

% ∗ ∗ ∗ 分隔
\newcommand{\asterisms}{%
  \par\vspace{1.2ex}%
  \noindent\hspace*{2em}$\ast$\hspace{1.2em}$\ast$\hspace{1.2em}$\ast$%
  \par\vspace{1.2ex}%
}

% 无编号子目的目录条目
\newcommand{\tocsubsub}[1]{%
  \phantomsection
  \addcontentsline{toc}{subsubsection}{#1}%
}

% 部扉: 仿原书部题页
\newcommand{\partdivider}[3]{%
  \clearpage
  \thispagestyle{empty}%
  \phantomsection
  \addcontentsline{toc}{part}{#1\quad #3}%
  \vspace*{6.5cm}
  \begin{center}
    {\fontsize{16pt}{26pt}\selectfont\headfont #1\par}
    \vspace{0.9cm}
    {\fontsize{22pt}{34pt}\selectfont\headfont #2\par}
  \end{center}
  \clearpage
}

% 献辞页: 仿原书部题后的题献, 独占一页且不入目录
\newcommand{\dedication}[1]{%
  \thispagestyle{empty}%
  \vspace*{4cm}%
  \begin{center}#1\end{center}%
  \clearpage
}

% --------------------------------------------------------------
%  各卷"译者说明"的共用体例段(各卷的范围段自写, 后接本段)
% --------------------------------------------------------------
\newcommand{\CommonNoteBody}{%
为便于与原书及学界通行的打字稿页码对勘, 本册保留:

\begin{itemize}
  \item 原书全部章节编号(章~1--18, 节如~5.2、18.2, 子目如~18.1.1)与原书小节标题;
  \item 原书脚注的编号与角标——
        \begin{itemize}
          \item 前奏(第~1--4 章)与第一部的注在各页内重新编号, 正文中的上标即该
                页序号; 第一部紧随其后的括注(如“(3)”)是该部分的连续注号, 与第~12
                章“第一部分的注”里的“注~3”一一对应, 注文行首照原书印作“1(3)”式。
                其中三条注原书未标连续注号而改标日期(ch7 注~3 = 30 Septembre、
                ch9 注~11/12 = 8 août), 本册照录其日期;
          \item 第二部第~13--16~章各自重新编号, 分别为~1--41、1--86、1--132、1--29;
                第三部第~17~章自~1~起至~12, 第~18~章自~1~重新起算并延续至~794;
                末段自~\textbf{795(*)}~续至~1053(*)。各段与原书一一对应;
        \end{itemize}
  \item 页边码“p.~1”“p.~A3”“p.~1066”式小注——这是作者引用自己打字稿原稿时
        所用的页码, 与本排印本页码不同。前奏四章的边码各有前缀: 第~1 章作
        \texttt{A1…A6}, 第~2 章作 \texttt{P1…P65}, 第~3 章作 \texttt{L1…L56},
        第~4 章作罗马数字 \texttt{I…XXII}, 自第一部起作纯数字 \texttt{1…171};
  \item 作者用以标记段落的~$\diamond$~记号, 以及标示后来补写文字的上标
        ~\textsuperscript{$\diamond$}~记号;
  \item 原书大小注(“注~44'”“注~46$^1$”“注~171³”等)的编号与层级。
\end{itemize}

第~12 章“第一部分的注”把第~5 至第~11 章各页脚注汇为一编, 并给每条注配上节题
(如“12.19. 耶稣与十二使徒”), 因此部分注文与正文脚注重复——这是原书的体例, 非
重出之误。这一部分属随笔式的自省文字, 同一句法语在正文脚注与章末汇编注中
各有措辞, 本册有意保留各自的译笔, 不强求统一。

除上述标记外, 译文力求贴近作者驳杂、跳跃而辛辣的文风; 数学名词的译法见各卷
术语表(如~\emph{topos}~不译, \emph{faisceaux pervers}~作“反常层”)。文中的日期
括注(如“(5~月~30~日)”“(1984年3月补写)”)为作者历次补写时的纪日, 一律照录。文中
另有一些对注号、页码的引用(如“n°~1695”“打字稿第~227~页”)沿袭作者打字稿的编号
系统, 与本排印本不尽一致, 亦无法在本书中落实, 悉照原文录出。

全套五卷(另有全书合订一册)的分工是:

\begin{itemize}
  \item \textbf{卷一·前奏}——主题的呈现: 或四个乐章的前奏(原书 p.23--122);
  \item \textbf{卷二·第一部}——愚妄与更新(原书 p.123--254);
  \item \textbf{卷三·第二部}——葬礼(I): 或中国皇帝的新衣(原书 p.255--412);
  \item \textbf{卷四·第三部}——葬礼(II): 或阴阳之钥(原书 p.413--632);
  \item \textbf{卷五·第四部}——葬礼(III): 或四种运算(原书 p.633--929)。
\end{itemize}
}
 之后可用 \hypersetup 覆写)
% --------------------------------------------------------------
\hypersetup{
  pdftitle={收获与播种},
  pdfauthor={Alexandre Grothendieck / 数译 译},
  pdfsubject={Récoltes et Semailles — Réflexions et témoignage sur un passé de mathématicien}
}

% --------------------------------------------------------------
%  标题样式: 仿原书 —— 黑体左对齐, 编号后带句点
% --------------------------------------------------------------
\titleformat{\section}[block]
  {\normalfont\Large\headfont}{\thesection.}{0.5em}{}
\titlespacing*{\section}{0pt}{3.2ex plus 1.2ex minus .2ex}{1.6ex plus .2ex}
\titleformat{\subsection}[block]
  {\normalfont\large\headfont}{\thesubsection.}{0.5em}{}
\titlespacing*{\subsection}{0pt}{2.6ex plus 1ex minus .2ex}{1.2ex plus .2ex}
\titleformat{\subsubsection}[block]
  {\normalfont\normalsize\headfont}{\thesubsubsection.}{0.5em}{}
\titlespacing*{\subsubsection}{0pt}{2.2ex plus .8ex minus .2ex}{1ex plus .2ex}
\titleformat{name=\subsubsection,numberless}[block]
  {\normalfont\normalsize\headfont}{}{0em}{}
\setcounter{secnumdepth}{3}
\setcounter{tocdepth}{3}
% 为含长外文书名与行内公式的段落提供必要的断行余量
\setlength{\emergencystretch}{3em}

% 章标题样式: 原书章标题与节同级排印("13. A) HERITAGE..."),
% 故章标题取块状左对齐黑体, 字号介于节与部之间
\titleformat{\chapter}[block]
  {\normalfont\LARGE\headfont}{\thechapter.}{0.5em}{}
\titlespacing*{\chapter}{0pt}{3.6ex plus 1.4ex minus .2ex}{2ex plus .2ex}
% 目录中章号照原书用阿拉伯数字("13. A) ..." 而非 "第十三章")
\renewcommand{\CTEXthechapter}{\thechapter\hspace{.3em}}

% --------------------------------------------------------------
%  目录观感: 条目与引导点之间、引导点与页码之间各留 0.5em
%  内核 \@dottedtocline 让 \leaders 从标题末字起排, 第一个点几乎
%  贴着末字, 读起来像多标了一个句点; 页码前也没有余量。两处各补
%  0.5em。
% --------------------------------------------------------------
\makeatletter
\patchcmd{\@dottedtocline}{\nobreak\leaders}{\nobreak\hspace{0.5em}\leaders}%
  {\typeout{MT: TOC 引导点前已留 0.5em}}%
  {\typeout{MT-WARN: TOC 引导点前间距补丁未生效}}
\patchcmd{\@dottedtocline}{\hfill\nobreak\hb@xt@\@pnumwidth}%
  {\hfill\nobreak\hspace{0.5em}\hb@xt@\@pnumwidth}%
  {\typeout{MT: TOC 页码前已留 0.5em}}%
  {\typeout{MT-WARN: TOC 页码前间距补丁未生效}}
\makeatother

% --------------------------------------------------------------
%  正文半角标点后的间距
%  本书体例用半角 , ; : ? ! 与 ASCII 括号, 由 xeCJK 按 HalfRight
%  (右半留白) 处理。放宽到 0.45em, 与全角标点右半留白大致齐平。
% --------------------------------------------------------------
\xeCJKsetup{CJKecglue={\hskip 0.45em plus 0.06em minus 0.09em}}

% --------------------------------------------------------------
%  页眉页脚: 偶数页眉右对齐斜体节题; 页码页脚外侧
% --------------------------------------------------------------
\fancyhf{}
\fancyhead[RE]{\small\itshape\rightmark}
\fancyfoot[LE,RO]{\small\thepage}
\renewcommand{\headrulewidth}{0pt}
\renewcommand{\sectionmark}[1]{\markright{\thesection.\ #1}}
\fancypagestyle{plain}{%
  \fancyhf{}%
  \fancyfoot[LE,RO]{\small \thepage}%
  \renewcommand{\headrulewidth}{0pt}%
}

% --------------------------------------------------------------
%  原书版式记号宏
% --------------------------------------------------------------
% 边码: 原打字稿页码(G. 手稿系统里的自引页码), 排在页边
\usepackage{marginnote}
\newcommand{\origpage}[1]{\marginnote{\normalfont\scriptsize p.~#1}}

% ⋄ 记号: 段首/行内正常大小; 上标形式为作者后补文字的位置标记
\DeclareRobustCommand{\pmark}{\ensuremath{\diamond}\hspace{0.15em}}
\DeclareRobustCommand{\dmark}{\textsuperscript{\ensuremath{\diamond}}}

% 脚注人工编号: \fnote{795}{(*)}{文本} —— 编号与(*)角标均照录原书
\newcommand{\fnote}[3]{%
  \begingroup
  \renewcommand{\thefootnote}{#1#2}%
  \footnote{#3}%
  \endgroup
}

% 第一部分的脚注: 页内序号作上标(#1), 注文行首印 "序号(连续注号)",
% 正文在序号之后另行括注同一连续注号 —— \fnoteP{1}{(3)}{注文}
% 原书有三条注无连续注号而改标日期(ch07 注3 = 30 Septembre,
% ch09 注11/12 = 8 août), 第二参数留空, 日期写进注文之首。
\newcommand{\fnoteP}[3]{%
  \begingroup
  \renewcommand{\thefootnote}{#1}\footnotemark
  \renewcommand{\thefootnote}{#1#2}\footnotetext{#3}%
  \endgroup
}

% 正文呼号与注文行首标号不同(原书照录)
\newcommand{\fnotevariant}[4]{%
  \begingroup
  \renewcommand{\thefootnote}{#1#2}\footnotemark
  \renewcommand{\thefootnote}{#1#3}\footnotetext{#4}%
  \endgroup
}

% 大注块标题: \notehead{171³}{斜体评论} —— 正文随后另起
\newcommand{\notehead}[2]{%
  \par\bigskip
  \noindent{\headfont 注~#1}\ifblank{#2}{}{\hspace{0.5em}{\itshape #2}}%
  \par\nobreak\smallskip\nobreak
}

% ∗ ∗ ∗ 分隔
\newcommand{\asterisms}{%
  \par\vspace{1.2ex}%
  \noindent\hspace*{2em}$\ast$\hspace{1.2em}$\ast$\hspace{1.2em}$\ast$%
  \par\vspace{1.2ex}%
}

% 无编号子目的目录条目
\newcommand{\tocsubsub}[1]{%
  \phantomsection
  \addcontentsline{toc}{subsubsection}{#1}%
}

% 部扉: 仿原书部题页
\newcommand{\partdivider}[3]{%
  \clearpage
  \thispagestyle{empty}%
  \phantomsection
  \addcontentsline{toc}{part}{#1\quad #3}%
  \vspace*{6.5cm}
  \begin{center}
    {\fontsize{16pt}{26pt}\selectfont\headfont #1\par}
    \vspace{0.9cm}
    {\fontsize{22pt}{34pt}\selectfont\headfont #2\par}
  \end{center}
  \clearpage
}

% 献辞页: 仿原书部题后的题献, 独占一页且不入目录
\newcommand{\dedication}[1]{%
  \thispagestyle{empty}%
  \vspace*{4cm}%
  \begin{center}#1\end{center}%
  \clearpage
}

% --------------------------------------------------------------
%  各卷"译者说明"的共用体例段(各卷的范围段自写, 后接本段)
% --------------------------------------------------------------
\newcommand{\CommonNoteBody}{%
为便于与原书及学界通行的打字稿页码对勘, 本册保留:

\begin{itemize}
  \item 原书全部章节编号(章~1--18, 节如~5.2、18.2, 子目如~18.1.1)与原书小节标题;
  \item 原书脚注的编号与角标——
        \begin{itemize}
          \item 前奏(第~1--4 章)与第一部的注在各页内重新编号, 正文中的上标即该
                页序号; 第一部紧随其后的括注(如“(3)”)是该部分的连续注号, 与第~12
                章“第一部分的注”里的“注~3”一一对应, 注文行首照原书印作“1(3)”式。
                其中三条注原书未标连续注号而改标日期(ch7 注~3 = 30 Septembre、
                ch9 注~11/12 = 8 août), 本册照录其日期;
          \item 第二部第~13--16~章各自重新编号, 分别为~1--41、1--86、1--132、1--29;
                第三部第~17~章自~1~起至~12, 第~18~章自~1~重新起算并延续至~794;
                末段自~\textbf{795(*)}~续至~1053(*)。各段与原书一一对应;
        \end{itemize}
  \item 页边码“p.~1”“p.~A3”“p.~1066”式小注——这是作者引用自己打字稿原稿时
        所用的页码, 与本排印本页码不同。前奏四章的边码各有前缀: 第~1 章作
        \texttt{A1…A6}, 第~2 章作 \texttt{P1…P65}, 第~3 章作 \texttt{L1…L56},
        第~4 章作罗马数字 \texttt{I…XXII}, 自第一部起作纯数字 \texttt{1…171};
  \item 作者用以标记段落的~$\diamond$~记号, 以及标示后来补写文字的上标
        ~\textsuperscript{$\diamond$}~记号;
  \item 原书大小注(“注~44'”“注~46$^1$”“注~171³”等)的编号与层级。
\end{itemize}

第~12 章“第一部分的注”把第~5 至第~11 章各页脚注汇为一编, 并给每条注配上节题
(如“12.19. 耶稣与十二使徒”), 因此部分注文与正文脚注重复——这是原书的体例, 非
重出之误。这一部分属随笔式的自省文字, 同一句法语在正文脚注与章末汇编注中
各有措辞, 本册有意保留各自的译笔, 不强求统一。

除上述标记外, 译文力求贴近作者驳杂、跳跃而辛辣的文风; 数学名词的译法见各卷
术语表(如~\emph{topos}~不译, \emph{faisceaux pervers}~作“反常层”)。文中的日期
括注(如“(5~月~30~日)”“(1984年3月补写)”)为作者历次补写时的纪日, 一律照录。文中
另有一些对注号、页码的引用(如“n°~1695”“打字稿第~227~页”)沿袭作者打字稿的编号
系统, 与本排印本不尽一致, 亦无法在本书中落实, 悉照原文录出。

全套五卷(另有全书合订一册)的分工是:

\begin{itemize}
  \item \textbf{卷一·前奏}——主题的呈现: 或四个乐章的前奏(原书 p.23--122);
  \item \textbf{卷二·第一部}——愚妄与更新(原书 p.123--254);
  \item \textbf{卷三·第二部}——葬礼(I): 或中国皇帝的新衣(原书 p.255--412);
  \item \textbf{卷四·第三部}——葬礼(II): 或阴阳之钥(原书 p.413--632);
  \item \textbf{卷五·第四部}——葬礼(III): 或四种运算(原书 p.633--929)。
\end{itemize}
}
 之后可用 \hypersetup 覆写)
% --------------------------------------------------------------
\hypersetup{
  pdftitle={收获与播种},
  pdfauthor={Alexandre Grothendieck / 数译 译},
  pdfsubject={Récoltes et Semailles — Réflexions et témoignage sur un passé de mathématicien}
}

% --------------------------------------------------------------
%  标题样式: 仿原书 —— 黑体左对齐, 编号后带句点
% --------------------------------------------------------------
\titleformat{\section}[block]
  {\normalfont\Large\headfont}{\thesection.}{0.5em}{}
\titlespacing*{\section}{0pt}{3.2ex plus 1.2ex minus .2ex}{1.6ex plus .2ex}
\titleformat{\subsection}[block]
  {\normalfont\large\headfont}{\thesubsection.}{0.5em}{}
\titlespacing*{\subsection}{0pt}{2.6ex plus 1ex minus .2ex}{1.2ex plus .2ex}
\titleformat{\subsubsection}[block]
  {\normalfont\normalsize\headfont}{\thesubsubsection.}{0.5em}{}
\titlespacing*{\subsubsection}{0pt}{2.2ex plus .8ex minus .2ex}{1ex plus .2ex}
\titleformat{name=\subsubsection,numberless}[block]
  {\normalfont\normalsize\headfont}{}{0em}{}
\setcounter{secnumdepth}{3}
\setcounter{tocdepth}{3}
% 为含长外文书名与行内公式的段落提供必要的断行余量
\setlength{\emergencystretch}{3em}

% 章标题样式: 原书章标题与节同级排印("13. A) HERITAGE..."),
% 故章标题取块状左对齐黑体, 字号介于节与部之间
\titleformat{\chapter}[block]
  {\normalfont\LARGE\headfont}{\thechapter.}{0.5em}{}
\titlespacing*{\chapter}{0pt}{3.6ex plus 1.4ex minus .2ex}{2ex plus .2ex}
% 目录中章号照原书用阿拉伯数字("13. A) ..." 而非 "第十三章")
\renewcommand{\CTEXthechapter}{\thechapter\hspace{.3em}}

% --------------------------------------------------------------
%  目录观感: 条目与引导点之间、引导点与页码之间各留 0.5em
%  内核 \@dottedtocline 让 \leaders 从标题末字起排, 第一个点几乎
%  贴着末字, 读起来像多标了一个句点; 页码前也没有余量。两处各补
%  0.5em。
% --------------------------------------------------------------
\makeatletter
\patchcmd{\@dottedtocline}{\nobreak\leaders}{\nobreak\hspace{0.5em}\leaders}%
  {\typeout{MT: TOC 引导点前已留 0.5em}}%
  {\typeout{MT-WARN: TOC 引导点前间距补丁未生效}}
\patchcmd{\@dottedtocline}{\hfill\nobreak\hb@xt@\@pnumwidth}%
  {\hfill\nobreak\hspace{0.5em}\hb@xt@\@pnumwidth}%
  {\typeout{MT: TOC 页码前已留 0.5em}}%
  {\typeout{MT-WARN: TOC 页码前间距补丁未生效}}
\makeatother

% --------------------------------------------------------------
%  正文半角标点后的间距
%  本书体例用半角 , ; : ? ! 与 ASCII 括号, 由 xeCJK 按 HalfRight
%  (右半留白) 处理。放宽到 0.45em, 与全角标点右半留白大致齐平。
% --------------------------------------------------------------
\xeCJKsetup{CJKecglue={\hskip 0.45em plus 0.06em minus 0.09em}}

% --------------------------------------------------------------
%  页眉页脚: 偶数页眉右对齐斜体节题; 页码页脚外侧
% --------------------------------------------------------------
\fancyhf{}
\fancyhead[RE]{\small\itshape\rightmark}
\fancyfoot[LE,RO]{\small\thepage}
\renewcommand{\headrulewidth}{0pt}
\renewcommand{\sectionmark}[1]{\markright{\thesection.\ #1}}
\fancypagestyle{plain}{%
  \fancyhf{}%
  \fancyfoot[LE,RO]{\small \thepage}%
  \renewcommand{\headrulewidth}{0pt}%
}

% --------------------------------------------------------------
%  原书版式记号宏
% --------------------------------------------------------------
% 边码: 原打字稿页码(G. 手稿系统里的自引页码), 排在页边
\usepackage{marginnote}
\newcommand{\origpage}[1]{\marginnote{\normalfont\scriptsize p.~#1}}

% ⋄ 记号: 段首/行内正常大小; 上标形式为作者后补文字的位置标记
\DeclareRobustCommand{\pmark}{\ensuremath{\diamond}\hspace{0.15em}}
\DeclareRobustCommand{\dmark}{\textsuperscript{\ensuremath{\diamond}}}

% 脚注人工编号: \fnote{795}{(*)}{文本} —— 编号与(*)角标均照录原书
\newcommand{\fnote}[3]{%
  \begingroup
  \renewcommand{\thefootnote}{#1#2}%
  \footnote{#3}%
  \endgroup
}

% 第一部分的脚注: 页内序号作上标(#1), 注文行首印 "序号(连续注号)",
% 正文在序号之后另行括注同一连续注号 —— \fnoteP{1}{(3)}{注文}
% 原书有三条注无连续注号而改标日期(ch07 注3 = 30 Septembre,
% ch09 注11/12 = 8 août), 第二参数留空, 日期写进注文之首。
\newcommand{\fnoteP}[3]{%
  \begingroup
  \renewcommand{\thefootnote}{#1}\footnotemark
  \renewcommand{\thefootnote}{#1#2}\footnotetext{#3}%
  \endgroup
}

% 正文呼号与注文行首标号不同(原书照录)
\newcommand{\fnotevariant}[4]{%
  \begingroup
  \renewcommand{\thefootnote}{#1#2}\footnotemark
  \renewcommand{\thefootnote}{#1#3}\footnotetext{#4}%
  \endgroup
}

% 大注块标题: \notehead{171³}{斜体评论} —— 正文随后另起
\newcommand{\notehead}[2]{%
  \par\bigskip
  \noindent{\headfont 注~#1}\ifblank{#2}{}{\hspace{0.5em}{\itshape #2}}%
  \par\nobreak\smallskip\nobreak
}

% ∗ ∗ ∗ 分隔
\newcommand{\asterisms}{%
  \par\vspace{1.2ex}%
  \noindent\hspace*{2em}$\ast$\hspace{1.2em}$\ast$\hspace{1.2em}$\ast$%
  \par\vspace{1.2ex}%
}

% 无编号子目的目录条目
\newcommand{\tocsubsub}[1]{%
  \phantomsection
  \addcontentsline{toc}{subsubsection}{#1}%
}

% 部扉: 仿原书部题页
\newcommand{\partdivider}[3]{%
  \clearpage
  \thispagestyle{empty}%
  \phantomsection
  \addcontentsline{toc}{part}{#1\quad #3}%
  \vspace*{6.5cm}
  \begin{center}
    {\fontsize{16pt}{26pt}\selectfont\headfont #1\par}
    \vspace{0.9cm}
    {\fontsize{22pt}{34pt}\selectfont\headfont #2\par}
  \end{center}
  \clearpage
}

% 献辞页: 仿原书部题后的题献, 独占一页且不入目录
\newcommand{\dedication}[1]{%
  \thispagestyle{empty}%
  \vspace*{4cm}%
  \begin{center}#1\end{center}%
  \clearpage
}

% --------------------------------------------------------------
%  各卷"译者说明"的共用体例段(各卷的范围段自写, 后接本段)
% --------------------------------------------------------------
\newcommand{\CommonNoteBody}{%
为便于与原书及学界通行的打字稿页码对勘, 本册保留:

\begin{itemize}
  \item 原书全部章节编号(章~1--18, 节如~5.2、18.2, 子目如~18.1.1)与原书小节标题;
  \item 原书脚注的编号与角标——
        \begin{itemize}
          \item 前奏(第~1--4 章)与第一部的注在各页内重新编号, 正文中的上标即该
                页序号; 第一部紧随其后的括注(如“(3)”)是该部分的连续注号, 与第~12
                章“第一部分的注”里的“注~3”一一对应, 注文行首照原书印作“1(3)”式。
                其中三条注原书未标连续注号而改标日期(ch7 注~3 = 30 Septembre、
                ch9 注~11/12 = 8 août), 本册照录其日期;
          \item 第二部第~13--16~章各自重新编号, 分别为~1--41、1--86、1--132、1--29;
                第三部第~17~章自~1~起至~12, 第~18~章自~1~重新起算并延续至~794;
                末段自~\textbf{795(*)}~续至~1053(*)。各段与原书一一对应;
        \end{itemize}
  \item 页边码“p.~1”“p.~A3”“p.~1066”式小注——这是作者引用自己打字稿原稿时
        所用的页码, 与本排印本页码不同。前奏四章的边码各有前缀: 第~1 章作
        \texttt{A1…A6}, 第~2 章作 \texttt{P1…P65}, 第~3 章作 \texttt{L1…L56},
        第~4 章作罗马数字 \texttt{I…XXII}, 自第一部起作纯数字 \texttt{1…171};
  \item 作者用以标记段落的~$\diamond$~记号, 以及标示后来补写文字的上标
        ~\textsuperscript{$\diamond$}~记号;
  \item 原书大小注(“注~44'”“注~46$^1$”“注~171³”等)的编号与层级。
\end{itemize}

第~12 章“第一部分的注”把第~5 至第~11 章各页脚注汇为一编, 并给每条注配上节题
(如“12.19. 耶稣与十二使徒”), 因此部分注文与正文脚注重复——这是原书的体例, 非
重出之误。这一部分属随笔式的自省文字, 同一句法语在正文脚注与章末汇编注中
各有措辞, 本册有意保留各自的译笔, 不强求统一。

除上述标记外, 译文力求贴近作者驳杂、跳跃而辛辣的文风; 数学名词的译法见各卷
术语表(如~\emph{topos}~不译, \emph{faisceaux pervers}~作“反常层”)。文中的日期
括注(如“(5~月~30~日)”“(1984年3月补写)”)为作者历次补写时的纪日, 一律照录。文中
另有一些对注号、页码的引用(如“n°~1695”“打字稿第~227~页”)沿袭作者打字稿的编号
系统, 与本排印本不尽一致, 亦无法在本书中落实, 悉照原文录出。

全套五卷(另有全书合订一册)的分工是:

\begin{itemize}
  \item \textbf{卷一·前奏}——主题的呈现: 或四个乐章的前奏(原书 p.23--122);
  \item \textbf{卷二·第一部}——愚妄与更新(原书 p.123--254);
  \item \textbf{卷三·第二部}——葬礼(I): 或中国皇帝的新衣(原书 p.255--412);
  \item \textbf{卷四·第三部}——葬礼(II): 或阴阳之钥(原书 p.413--632);
  \item \textbf{卷五·第四部}——葬礼(III): 或四种运算(原书 p.633--929)。
\end{itemize}
}


% --------------------------------------------------------------
%  封面参数(全书合订)
% --------------------------------------------------------------
% ================================================================
%  《收获与播种》统一封面 —— TikZ 原生矢量绘制(不依赖任何外部图片)
%  母题: 暮色中的麦田 —— 犁沟向地平线的灭点收束, 前景麦穗剪影,
%        一位播种者把一把种子撒向天空, 种子沿弧线落入田里。
%  用法(在 \begin{document} 之前, 依次):
%      % ================================================================
%  《收获与播种》中文译本 —— 各卷共用导言区
%  由 vol-prelude / vol-part1 / vol-part2 / vol-part3 / vol-part4
%  以及全书合订 main-all 共同 \input。基类 mathtranslation.cls v3.2 不改。
%  --------------------------------------------------------------
%  样式均在此改写, 以贴近原书:
%    · 章标题仿原书 "5. Travail et découverte"(块状左对齐黑体)
%    · 节/小节: 黑体左对齐, 编号后带句点
%    · 脚注: 人工编号照录原书——前奏与第一部分页内重新编号,
%      第二至四部各自连续编号并带(*)角标; 第一部分另有 \fnoteP
%      处理 "(n)" 式连续注号
%    · 边码: 原打字稿页码 \origpage{N}(前奏作 A1、P1、L1、罗马数字等)
%    · ⋄ 记号: 段首补写标记 \pmark 与上标形式 \dmark
% ================================================================

% 题名含空格, 不能走类选项(会被 kvoptions 吃掉), 在此赋值
\renewcommand{\BookTitleEN}{Récoltes et Semailles}
\renewcommand{\OriginalAuthor}{Alexandre Grothendieck}
\renewcommand{\TranslationDate}{\today}

% --------------------------------------------------------------
%  罕用字后备(如"祐"): Fandol 不载时退到 SimSun
% --------------------------------------------------------------
\setCJKfallbackfamilyfont{\CJKrmdefault}{{SimSun}}

% --------------------------------------------------------------
%  hyperref 元信息(各卷在 % ================================================================
%  《收获与播种》中文译本 —— 各卷共用导言区
%  由 vol-prelude / vol-part1 / vol-part2 / vol-part3 / vol-part4
%  以及全书合订 main-all 共同 \input。基类 mathtranslation.cls v3.2 不改。
%  --------------------------------------------------------------
%  样式均在此改写, 以贴近原书:
%    · 章标题仿原书 "5. Travail et découverte"(块状左对齐黑体)
%    · 节/小节: 黑体左对齐, 编号后带句点
%    · 脚注: 人工编号照录原书——前奏与第一部分页内重新编号,
%      第二至四部各自连续编号并带(*)角标; 第一部分另有 \fnoteP
%      处理 "(n)" 式连续注号
%    · 边码: 原打字稿页码 \origpage{N}(前奏作 A1、P1、L1、罗马数字等)
%    · ⋄ 记号: 段首补写标记 \pmark 与上标形式 \dmark
% ================================================================

% 题名含空格, 不能走类选项(会被 kvoptions 吃掉), 在此赋值
\renewcommand{\BookTitleEN}{Récoltes et Semailles}
\renewcommand{\OriginalAuthor}{Alexandre Grothendieck}
\renewcommand{\TranslationDate}{\today}

% --------------------------------------------------------------
%  罕用字后备(如"祐"): Fandol 不载时退到 SimSun
% --------------------------------------------------------------
\setCJKfallbackfamilyfont{\CJKrmdefault}{{SimSun}}

% --------------------------------------------------------------
%  hyperref 元信息(各卷在 \input{common} 之后可用 \hypersetup 覆写)
% --------------------------------------------------------------
\hypersetup{
  pdftitle={收获与播种},
  pdfauthor={Alexandre Grothendieck / 数译 译},
  pdfsubject={Récoltes et Semailles — Réflexions et témoignage sur un passé de mathématicien}
}

% --------------------------------------------------------------
%  标题样式: 仿原书 —— 黑体左对齐, 编号后带句点
% --------------------------------------------------------------
\titleformat{\section}[block]
  {\normalfont\Large\headfont}{\thesection.}{0.5em}{}
\titlespacing*{\section}{0pt}{3.2ex plus 1.2ex minus .2ex}{1.6ex plus .2ex}
\titleformat{\subsection}[block]
  {\normalfont\large\headfont}{\thesubsection.}{0.5em}{}
\titlespacing*{\subsection}{0pt}{2.6ex plus 1ex minus .2ex}{1.2ex plus .2ex}
\titleformat{\subsubsection}[block]
  {\normalfont\normalsize\headfont}{\thesubsubsection.}{0.5em}{}
\titlespacing*{\subsubsection}{0pt}{2.2ex plus .8ex minus .2ex}{1ex plus .2ex}
\titleformat{name=\subsubsection,numberless}[block]
  {\normalfont\normalsize\headfont}{}{0em}{}
\setcounter{secnumdepth}{3}
\setcounter{tocdepth}{3}
% 为含长外文书名与行内公式的段落提供必要的断行余量
\setlength{\emergencystretch}{3em}

% 章标题样式: 原书章标题与节同级排印("13. A) HERITAGE..."),
% 故章标题取块状左对齐黑体, 字号介于节与部之间
\titleformat{\chapter}[block]
  {\normalfont\LARGE\headfont}{\thechapter.}{0.5em}{}
\titlespacing*{\chapter}{0pt}{3.6ex plus 1.4ex minus .2ex}{2ex plus .2ex}
% 目录中章号照原书用阿拉伯数字("13. A) ..." 而非 "第十三章")
\renewcommand{\CTEXthechapter}{\thechapter\hspace{.3em}}

% --------------------------------------------------------------
%  目录观感: 条目与引导点之间、引导点与页码之间各留 0.5em
%  内核 \@dottedtocline 让 \leaders 从标题末字起排, 第一个点几乎
%  贴着末字, 读起来像多标了一个句点; 页码前也没有余量。两处各补
%  0.5em。
% --------------------------------------------------------------
\makeatletter
\patchcmd{\@dottedtocline}{\nobreak\leaders}{\nobreak\hspace{0.5em}\leaders}%
  {\typeout{MT: TOC 引导点前已留 0.5em}}%
  {\typeout{MT-WARN: TOC 引导点前间距补丁未生效}}
\patchcmd{\@dottedtocline}{\hfill\nobreak\hb@xt@\@pnumwidth}%
  {\hfill\nobreak\hspace{0.5em}\hb@xt@\@pnumwidth}%
  {\typeout{MT: TOC 页码前已留 0.5em}}%
  {\typeout{MT-WARN: TOC 页码前间距补丁未生效}}
\makeatother

% --------------------------------------------------------------
%  正文半角标点后的间距
%  本书体例用半角 , ; : ? ! 与 ASCII 括号, 由 xeCJK 按 HalfRight
%  (右半留白) 处理。放宽到 0.45em, 与全角标点右半留白大致齐平。
% --------------------------------------------------------------
\xeCJKsetup{CJKecglue={\hskip 0.45em plus 0.06em minus 0.09em}}

% --------------------------------------------------------------
%  页眉页脚: 偶数页眉右对齐斜体节题; 页码页脚外侧
% --------------------------------------------------------------
\fancyhf{}
\fancyhead[RE]{\small\itshape\rightmark}
\fancyfoot[LE,RO]{\small\thepage}
\renewcommand{\headrulewidth}{0pt}
\renewcommand{\sectionmark}[1]{\markright{\thesection.\ #1}}
\fancypagestyle{plain}{%
  \fancyhf{}%
  \fancyfoot[LE,RO]{\small \thepage}%
  \renewcommand{\headrulewidth}{0pt}%
}

% --------------------------------------------------------------
%  原书版式记号宏
% --------------------------------------------------------------
% 边码: 原打字稿页码(G. 手稿系统里的自引页码), 排在页边
\usepackage{marginnote}
\newcommand{\origpage}[1]{\marginnote{\normalfont\scriptsize p.~#1}}

% ⋄ 记号: 段首/行内正常大小; 上标形式为作者后补文字的位置标记
\DeclareRobustCommand{\pmark}{\ensuremath{\diamond}\hspace{0.15em}}
\DeclareRobustCommand{\dmark}{\textsuperscript{\ensuremath{\diamond}}}

% 脚注人工编号: \fnote{795}{(*)}{文本} —— 编号与(*)角标均照录原书
\newcommand{\fnote}[3]{%
  \begingroup
  \renewcommand{\thefootnote}{#1#2}%
  \footnote{#3}%
  \endgroup
}

% 第一部分的脚注: 页内序号作上标(#1), 注文行首印 "序号(连续注号)",
% 正文在序号之后另行括注同一连续注号 —— \fnoteP{1}{(3)}{注文}
% 原书有三条注无连续注号而改标日期(ch07 注3 = 30 Septembre,
% ch09 注11/12 = 8 août), 第二参数留空, 日期写进注文之首。
\newcommand{\fnoteP}[3]{%
  \begingroup
  \renewcommand{\thefootnote}{#1}\footnotemark
  \renewcommand{\thefootnote}{#1#2}\footnotetext{#3}%
  \endgroup
}

% 正文呼号与注文行首标号不同(原书照录)
\newcommand{\fnotevariant}[4]{%
  \begingroup
  \renewcommand{\thefootnote}{#1#2}\footnotemark
  \renewcommand{\thefootnote}{#1#3}\footnotetext{#4}%
  \endgroup
}

% 大注块标题: \notehead{171³}{斜体评论} —— 正文随后另起
\newcommand{\notehead}[2]{%
  \par\bigskip
  \noindent{\headfont 注~#1}\ifblank{#2}{}{\hspace{0.5em}{\itshape #2}}%
  \par\nobreak\smallskip\nobreak
}

% ∗ ∗ ∗ 分隔
\newcommand{\asterisms}{%
  \par\vspace{1.2ex}%
  \noindent\hspace*{2em}$\ast$\hspace{1.2em}$\ast$\hspace{1.2em}$\ast$%
  \par\vspace{1.2ex}%
}

% 无编号子目的目录条目
\newcommand{\tocsubsub}[1]{%
  \phantomsection
  \addcontentsline{toc}{subsubsection}{#1}%
}

% 部扉: 仿原书部题页
\newcommand{\partdivider}[3]{%
  \clearpage
  \thispagestyle{empty}%
  \phantomsection
  \addcontentsline{toc}{part}{#1\quad #3}%
  \vspace*{6.5cm}
  \begin{center}
    {\fontsize{16pt}{26pt}\selectfont\headfont #1\par}
    \vspace{0.9cm}
    {\fontsize{22pt}{34pt}\selectfont\headfont #2\par}
  \end{center}
  \clearpage
}

% 献辞页: 仿原书部题后的题献, 独占一页且不入目录
\newcommand{\dedication}[1]{%
  \thispagestyle{empty}%
  \vspace*{4cm}%
  \begin{center}#1\end{center}%
  \clearpage
}

% --------------------------------------------------------------
%  各卷"译者说明"的共用体例段(各卷的范围段自写, 后接本段)
% --------------------------------------------------------------
\newcommand{\CommonNoteBody}{%
为便于与原书及学界通行的打字稿页码对勘, 本册保留:

\begin{itemize}
  \item 原书全部章节编号(章~1--18, 节如~5.2、18.2, 子目如~18.1.1)与原书小节标题;
  \item 原书脚注的编号与角标——
        \begin{itemize}
          \item 前奏(第~1--4 章)与第一部的注在各页内重新编号, 正文中的上标即该
                页序号; 第一部紧随其后的括注(如“(3)”)是该部分的连续注号, 与第~12
                章“第一部分的注”里的“注~3”一一对应, 注文行首照原书印作“1(3)”式。
                其中三条注原书未标连续注号而改标日期(ch7 注~3 = 30 Septembre、
                ch9 注~11/12 = 8 août), 本册照录其日期;
          \item 第二部第~13--16~章各自重新编号, 分别为~1--41、1--86、1--132、1--29;
                第三部第~17~章自~1~起至~12, 第~18~章自~1~重新起算并延续至~794;
                末段自~\textbf{795(*)}~续至~1053(*)。各段与原书一一对应;
        \end{itemize}
  \item 页边码“p.~1”“p.~A3”“p.~1066”式小注——这是作者引用自己打字稿原稿时
        所用的页码, 与本排印本页码不同。前奏四章的边码各有前缀: 第~1 章作
        \texttt{A1…A6}, 第~2 章作 \texttt{P1…P65}, 第~3 章作 \texttt{L1…L56},
        第~4 章作罗马数字 \texttt{I…XXII}, 自第一部起作纯数字 \texttt{1…171};
  \item 作者用以标记段落的~$\diamond$~记号, 以及标示后来补写文字的上标
        ~\textsuperscript{$\diamond$}~记号;
  \item 原书大小注(“注~44'”“注~46$^1$”“注~171³”等)的编号与层级。
\end{itemize}

第~12 章“第一部分的注”把第~5 至第~11 章各页脚注汇为一编, 并给每条注配上节题
(如“12.19. 耶稣与十二使徒”), 因此部分注文与正文脚注重复——这是原书的体例, 非
重出之误。这一部分属随笔式的自省文字, 同一句法语在正文脚注与章末汇编注中
各有措辞, 本册有意保留各自的译笔, 不强求统一。

除上述标记外, 译文力求贴近作者驳杂、跳跃而辛辣的文风; 数学名词的译法见各卷
术语表(如~\emph{topos}~不译, \emph{faisceaux pervers}~作“反常层”)。文中的日期
括注(如“(5~月~30~日)”“(1984年3月补写)”)为作者历次补写时的纪日, 一律照录。文中
另有一些对注号、页码的引用(如“n°~1695”“打字稿第~227~页”)沿袭作者打字稿的编号
系统, 与本排印本不尽一致, 亦无法在本书中落实, 悉照原文录出。

全套五卷(另有全书合订一册)的分工是:

\begin{itemize}
  \item \textbf{卷一·前奏}——主题的呈现: 或四个乐章的前奏(原书 p.23--122);
  \item \textbf{卷二·第一部}——愚妄与更新(原书 p.123--254);
  \item \textbf{卷三·第二部}——葬礼(I): 或中国皇帝的新衣(原书 p.255--412);
  \item \textbf{卷四·第三部}——葬礼(II): 或阴阳之钥(原书 p.413--632);
  \item \textbf{卷五·第四部}——葬礼(III): 或四种运算(原书 p.633--929)。
\end{itemize}
}
 之后可用 \hypersetup 覆写)
% --------------------------------------------------------------
\hypersetup{
  pdftitle={收获与播种},
  pdfauthor={Alexandre Grothendieck / 数译 译},
  pdfsubject={Récoltes et Semailles — Réflexions et témoignage sur un passé de mathématicien}
}

% --------------------------------------------------------------
%  标题样式: 仿原书 —— 黑体左对齐, 编号后带句点
% --------------------------------------------------------------
\titleformat{\section}[block]
  {\normalfont\Large\headfont}{\thesection.}{0.5em}{}
\titlespacing*{\section}{0pt}{3.2ex plus 1.2ex minus .2ex}{1.6ex plus .2ex}
\titleformat{\subsection}[block]
  {\normalfont\large\headfont}{\thesubsection.}{0.5em}{}
\titlespacing*{\subsection}{0pt}{2.6ex plus 1ex minus .2ex}{1.2ex plus .2ex}
\titleformat{\subsubsection}[block]
  {\normalfont\normalsize\headfont}{\thesubsubsection.}{0.5em}{}
\titlespacing*{\subsubsection}{0pt}{2.2ex plus .8ex minus .2ex}{1ex plus .2ex}
\titleformat{name=\subsubsection,numberless}[block]
  {\normalfont\normalsize\headfont}{}{0em}{}
\setcounter{secnumdepth}{3}
\setcounter{tocdepth}{3}
% 为含长外文书名与行内公式的段落提供必要的断行余量
\setlength{\emergencystretch}{3em}

% 章标题样式: 原书章标题与节同级排印("13. A) HERITAGE..."),
% 故章标题取块状左对齐黑体, 字号介于节与部之间
\titleformat{\chapter}[block]
  {\normalfont\LARGE\headfont}{\thechapter.}{0.5em}{}
\titlespacing*{\chapter}{0pt}{3.6ex plus 1.4ex minus .2ex}{2ex plus .2ex}
% 目录中章号照原书用阿拉伯数字("13. A) ..." 而非 "第十三章")
\renewcommand{\CTEXthechapter}{\thechapter\hspace{.3em}}

% --------------------------------------------------------------
%  目录观感: 条目与引导点之间、引导点与页码之间各留 0.5em
%  内核 \@dottedtocline 让 \leaders 从标题末字起排, 第一个点几乎
%  贴着末字, 读起来像多标了一个句点; 页码前也没有余量。两处各补
%  0.5em。
% --------------------------------------------------------------
\makeatletter
\patchcmd{\@dottedtocline}{\nobreak\leaders}{\nobreak\hspace{0.5em}\leaders}%
  {\typeout{MT: TOC 引导点前已留 0.5em}}%
  {\typeout{MT-WARN: TOC 引导点前间距补丁未生效}}
\patchcmd{\@dottedtocline}{\hfill\nobreak\hb@xt@\@pnumwidth}%
  {\hfill\nobreak\hspace{0.5em}\hb@xt@\@pnumwidth}%
  {\typeout{MT: TOC 页码前已留 0.5em}}%
  {\typeout{MT-WARN: TOC 页码前间距补丁未生效}}
\makeatother

% --------------------------------------------------------------
%  正文半角标点后的间距
%  本书体例用半角 , ; : ? ! 与 ASCII 括号, 由 xeCJK 按 HalfRight
%  (右半留白) 处理。放宽到 0.45em, 与全角标点右半留白大致齐平。
% --------------------------------------------------------------
\xeCJKsetup{CJKecglue={\hskip 0.45em plus 0.06em minus 0.09em}}

% --------------------------------------------------------------
%  页眉页脚: 偶数页眉右对齐斜体节题; 页码页脚外侧
% --------------------------------------------------------------
\fancyhf{}
\fancyhead[RE]{\small\itshape\rightmark}
\fancyfoot[LE,RO]{\small\thepage}
\renewcommand{\headrulewidth}{0pt}
\renewcommand{\sectionmark}[1]{\markright{\thesection.\ #1}}
\fancypagestyle{plain}{%
  \fancyhf{}%
  \fancyfoot[LE,RO]{\small \thepage}%
  \renewcommand{\headrulewidth}{0pt}%
}

% --------------------------------------------------------------
%  原书版式记号宏
% --------------------------------------------------------------
% 边码: 原打字稿页码(G. 手稿系统里的自引页码), 排在页边
\usepackage{marginnote}
\newcommand{\origpage}[1]{\marginnote{\normalfont\scriptsize p.~#1}}

% ⋄ 记号: 段首/行内正常大小; 上标形式为作者后补文字的位置标记
\DeclareRobustCommand{\pmark}{\ensuremath{\diamond}\hspace{0.15em}}
\DeclareRobustCommand{\dmark}{\textsuperscript{\ensuremath{\diamond}}}

% 脚注人工编号: \fnote{795}{(*)}{文本} —— 编号与(*)角标均照录原书
\newcommand{\fnote}[3]{%
  \begingroup
  \renewcommand{\thefootnote}{#1#2}%
  \footnote{#3}%
  \endgroup
}

% 第一部分的脚注: 页内序号作上标(#1), 注文行首印 "序号(连续注号)",
% 正文在序号之后另行括注同一连续注号 —— \fnoteP{1}{(3)}{注文}
% 原书有三条注无连续注号而改标日期(ch07 注3 = 30 Septembre,
% ch09 注11/12 = 8 août), 第二参数留空, 日期写进注文之首。
\newcommand{\fnoteP}[3]{%
  \begingroup
  \renewcommand{\thefootnote}{#1}\footnotemark
  \renewcommand{\thefootnote}{#1#2}\footnotetext{#3}%
  \endgroup
}

% 正文呼号与注文行首标号不同(原书照录)
\newcommand{\fnotevariant}[4]{%
  \begingroup
  \renewcommand{\thefootnote}{#1#2}\footnotemark
  \renewcommand{\thefootnote}{#1#3}\footnotetext{#4}%
  \endgroup
}

% 大注块标题: \notehead{171³}{斜体评论} —— 正文随后另起
\newcommand{\notehead}[2]{%
  \par\bigskip
  \noindent{\headfont 注~#1}\ifblank{#2}{}{\hspace{0.5em}{\itshape #2}}%
  \par\nobreak\smallskip\nobreak
}

% ∗ ∗ ∗ 分隔
\newcommand{\asterisms}{%
  \par\vspace{1.2ex}%
  \noindent\hspace*{2em}$\ast$\hspace{1.2em}$\ast$\hspace{1.2em}$\ast$%
  \par\vspace{1.2ex}%
}

% 无编号子目的目录条目
\newcommand{\tocsubsub}[1]{%
  \phantomsection
  \addcontentsline{toc}{subsubsection}{#1}%
}

% 部扉: 仿原书部题页
\newcommand{\partdivider}[3]{%
  \clearpage
  \thispagestyle{empty}%
  \phantomsection
  \addcontentsline{toc}{part}{#1\quad #3}%
  \vspace*{6.5cm}
  \begin{center}
    {\fontsize{16pt}{26pt}\selectfont\headfont #1\par}
    \vspace{0.9cm}
    {\fontsize{22pt}{34pt}\selectfont\headfont #2\par}
  \end{center}
  \clearpage
}

% 献辞页: 仿原书部题后的题献, 独占一页且不入目录
\newcommand{\dedication}[1]{%
  \thispagestyle{empty}%
  \vspace*{4cm}%
  \begin{center}#1\end{center}%
  \clearpage
}

% --------------------------------------------------------------
%  各卷"译者说明"的共用体例段(各卷的范围段自写, 后接本段)
% --------------------------------------------------------------
\newcommand{\CommonNoteBody}{%
为便于与原书及学界通行的打字稿页码对勘, 本册保留:

\begin{itemize}
  \item 原书全部章节编号(章~1--18, 节如~5.2、18.2, 子目如~18.1.1)与原书小节标题;
  \item 原书脚注的编号与角标——
        \begin{itemize}
          \item 前奏(第~1--4 章)与第一部的注在各页内重新编号, 正文中的上标即该
                页序号; 第一部紧随其后的括注(如“(3)”)是该部分的连续注号, 与第~12
                章“第一部分的注”里的“注~3”一一对应, 注文行首照原书印作“1(3)”式。
                其中三条注原书未标连续注号而改标日期(ch7 注~3 = 30 Septembre、
                ch9 注~11/12 = 8 août), 本册照录其日期;
          \item 第二部第~13--16~章各自重新编号, 分别为~1--41、1--86、1--132、1--29;
                第三部第~17~章自~1~起至~12, 第~18~章自~1~重新起算并延续至~794;
                末段自~\textbf{795(*)}~续至~1053(*)。各段与原书一一对应;
        \end{itemize}
  \item 页边码“p.~1”“p.~A3”“p.~1066”式小注——这是作者引用自己打字稿原稿时
        所用的页码, 与本排印本页码不同。前奏四章的边码各有前缀: 第~1 章作
        \texttt{A1…A6}, 第~2 章作 \texttt{P1…P65}, 第~3 章作 \texttt{L1…L56},
        第~4 章作罗马数字 \texttt{I…XXII}, 自第一部起作纯数字 \texttt{1…171};
  \item 作者用以标记段落的~$\diamond$~记号, 以及标示后来补写文字的上标
        ~\textsuperscript{$\diamond$}~记号;
  \item 原书大小注(“注~44'”“注~46$^1$”“注~171³”等)的编号与层级。
\end{itemize}

第~12 章“第一部分的注”把第~5 至第~11 章各页脚注汇为一编, 并给每条注配上节题
(如“12.19. 耶稣与十二使徒”), 因此部分注文与正文脚注重复——这是原书的体例, 非
重出之误。这一部分属随笔式的自省文字, 同一句法语在正文脚注与章末汇编注中
各有措辞, 本册有意保留各自的译笔, 不强求统一。

除上述标记外, 译文力求贴近作者驳杂、跳跃而辛辣的文风; 数学名词的译法见各卷
术语表(如~\emph{topos}~不译, \emph{faisceaux pervers}~作“反常层”)。文中的日期
括注(如“(5~月~30~日)”“(1984年3月补写)”)为作者历次补写时的纪日, 一律照录。文中
另有一些对注号、页码的引用(如“n°~1695”“打字稿第~227~页”)沿袭作者打字稿的编号
系统, 与本排印本不尽一致, 亦无法在本书中落实, 悉照原文录出。

全套五卷(另有全书合订一册)的分工是:

\begin{itemize}
  \item \textbf{卷一·前奏}——主题的呈现: 或四个乐章的前奏(原书 p.23--122);
  \item \textbf{卷二·第一部}——愚妄与更新(原书 p.123--254);
  \item \textbf{卷三·第二部}——葬礼(I): 或中国皇帝的新衣(原书 p.255--412);
  \item \textbf{卷四·第三部}——葬礼(II): 或阴阳之钥(原书 p.413--632);
  \item \textbf{卷五·第四部}——葬礼(III): 或四种运算(原书 p.633--929)。
\end{itemize}
}
                       % 共用导言区
%      % ================================================================
%  《收获与播种》统一封面 —— TikZ 原生矢量绘制(不依赖任何外部图片)
%  母题: 暮色中的麦田 —— 犁沟向地平线的灭点收束, 前景麦穗剪影,
%        一位播种者把一把种子撒向天空, 种子沿弧线落入田里。
%  用法(在 \begin{document} 之前, 依次):
%      % ================================================================
%  《收获与播种》中文译本 —— 各卷共用导言区
%  由 vol-prelude / vol-part1 / vol-part2 / vol-part3 / vol-part4
%  以及全书合订 main-all 共同 \input。基类 mathtranslation.cls v3.2 不改。
%  --------------------------------------------------------------
%  样式均在此改写, 以贴近原书:
%    · 章标题仿原书 "5. Travail et découverte"(块状左对齐黑体)
%    · 节/小节: 黑体左对齐, 编号后带句点
%    · 脚注: 人工编号照录原书——前奏与第一部分页内重新编号,
%      第二至四部各自连续编号并带(*)角标; 第一部分另有 \fnoteP
%      处理 "(n)" 式连续注号
%    · 边码: 原打字稿页码 \origpage{N}(前奏作 A1、P1、L1、罗马数字等)
%    · ⋄ 记号: 段首补写标记 \pmark 与上标形式 \dmark
% ================================================================

% 题名含空格, 不能走类选项(会被 kvoptions 吃掉), 在此赋值
\renewcommand{\BookTitleEN}{Récoltes et Semailles}
\renewcommand{\OriginalAuthor}{Alexandre Grothendieck}
\renewcommand{\TranslationDate}{\today}

% --------------------------------------------------------------
%  罕用字后备(如"祐"): Fandol 不载时退到 SimSun
% --------------------------------------------------------------
\setCJKfallbackfamilyfont{\CJKrmdefault}{{SimSun}}

% --------------------------------------------------------------
%  hyperref 元信息(各卷在 \input{common} 之后可用 \hypersetup 覆写)
% --------------------------------------------------------------
\hypersetup{
  pdftitle={收获与播种},
  pdfauthor={Alexandre Grothendieck / 数译 译},
  pdfsubject={Récoltes et Semailles — Réflexions et témoignage sur un passé de mathématicien}
}

% --------------------------------------------------------------
%  标题样式: 仿原书 —— 黑体左对齐, 编号后带句点
% --------------------------------------------------------------
\titleformat{\section}[block]
  {\normalfont\Large\headfont}{\thesection.}{0.5em}{}
\titlespacing*{\section}{0pt}{3.2ex plus 1.2ex minus .2ex}{1.6ex plus .2ex}
\titleformat{\subsection}[block]
  {\normalfont\large\headfont}{\thesubsection.}{0.5em}{}
\titlespacing*{\subsection}{0pt}{2.6ex plus 1ex minus .2ex}{1.2ex plus .2ex}
\titleformat{\subsubsection}[block]
  {\normalfont\normalsize\headfont}{\thesubsubsection.}{0.5em}{}
\titlespacing*{\subsubsection}{0pt}{2.2ex plus .8ex minus .2ex}{1ex plus .2ex}
\titleformat{name=\subsubsection,numberless}[block]
  {\normalfont\normalsize\headfont}{}{0em}{}
\setcounter{secnumdepth}{3}
\setcounter{tocdepth}{3}
% 为含长外文书名与行内公式的段落提供必要的断行余量
\setlength{\emergencystretch}{3em}

% 章标题样式: 原书章标题与节同级排印("13. A) HERITAGE..."),
% 故章标题取块状左对齐黑体, 字号介于节与部之间
\titleformat{\chapter}[block]
  {\normalfont\LARGE\headfont}{\thechapter.}{0.5em}{}
\titlespacing*{\chapter}{0pt}{3.6ex plus 1.4ex minus .2ex}{2ex plus .2ex}
% 目录中章号照原书用阿拉伯数字("13. A) ..." 而非 "第十三章")
\renewcommand{\CTEXthechapter}{\thechapter\hspace{.3em}}

% --------------------------------------------------------------
%  目录观感: 条目与引导点之间、引导点与页码之间各留 0.5em
%  内核 \@dottedtocline 让 \leaders 从标题末字起排, 第一个点几乎
%  贴着末字, 读起来像多标了一个句点; 页码前也没有余量。两处各补
%  0.5em。
% --------------------------------------------------------------
\makeatletter
\patchcmd{\@dottedtocline}{\nobreak\leaders}{\nobreak\hspace{0.5em}\leaders}%
  {\typeout{MT: TOC 引导点前已留 0.5em}}%
  {\typeout{MT-WARN: TOC 引导点前间距补丁未生效}}
\patchcmd{\@dottedtocline}{\hfill\nobreak\hb@xt@\@pnumwidth}%
  {\hfill\nobreak\hspace{0.5em}\hb@xt@\@pnumwidth}%
  {\typeout{MT: TOC 页码前已留 0.5em}}%
  {\typeout{MT-WARN: TOC 页码前间距补丁未生效}}
\makeatother

% --------------------------------------------------------------
%  正文半角标点后的间距
%  本书体例用半角 , ; : ? ! 与 ASCII 括号, 由 xeCJK 按 HalfRight
%  (右半留白) 处理。放宽到 0.45em, 与全角标点右半留白大致齐平。
% --------------------------------------------------------------
\xeCJKsetup{CJKecglue={\hskip 0.45em plus 0.06em minus 0.09em}}

% --------------------------------------------------------------
%  页眉页脚: 偶数页眉右对齐斜体节题; 页码页脚外侧
% --------------------------------------------------------------
\fancyhf{}
\fancyhead[RE]{\small\itshape\rightmark}
\fancyfoot[LE,RO]{\small\thepage}
\renewcommand{\headrulewidth}{0pt}
\renewcommand{\sectionmark}[1]{\markright{\thesection.\ #1}}
\fancypagestyle{plain}{%
  \fancyhf{}%
  \fancyfoot[LE,RO]{\small \thepage}%
  \renewcommand{\headrulewidth}{0pt}%
}

% --------------------------------------------------------------
%  原书版式记号宏
% --------------------------------------------------------------
% 边码: 原打字稿页码(G. 手稿系统里的自引页码), 排在页边
\usepackage{marginnote}
\newcommand{\origpage}[1]{\marginnote{\normalfont\scriptsize p.~#1}}

% ⋄ 记号: 段首/行内正常大小; 上标形式为作者后补文字的位置标记
\DeclareRobustCommand{\pmark}{\ensuremath{\diamond}\hspace{0.15em}}
\DeclareRobustCommand{\dmark}{\textsuperscript{\ensuremath{\diamond}}}

% 脚注人工编号: \fnote{795}{(*)}{文本} —— 编号与(*)角标均照录原书
\newcommand{\fnote}[3]{%
  \begingroup
  \renewcommand{\thefootnote}{#1#2}%
  \footnote{#3}%
  \endgroup
}

% 第一部分的脚注: 页内序号作上标(#1), 注文行首印 "序号(连续注号)",
% 正文在序号之后另行括注同一连续注号 —— \fnoteP{1}{(3)}{注文}
% 原书有三条注无连续注号而改标日期(ch07 注3 = 30 Septembre,
% ch09 注11/12 = 8 août), 第二参数留空, 日期写进注文之首。
\newcommand{\fnoteP}[3]{%
  \begingroup
  \renewcommand{\thefootnote}{#1}\footnotemark
  \renewcommand{\thefootnote}{#1#2}\footnotetext{#3}%
  \endgroup
}

% 正文呼号与注文行首标号不同(原书照录)
\newcommand{\fnotevariant}[4]{%
  \begingroup
  \renewcommand{\thefootnote}{#1#2}\footnotemark
  \renewcommand{\thefootnote}{#1#3}\footnotetext{#4}%
  \endgroup
}

% 大注块标题: \notehead{171³}{斜体评论} —— 正文随后另起
\newcommand{\notehead}[2]{%
  \par\bigskip
  \noindent{\headfont 注~#1}\ifblank{#2}{}{\hspace{0.5em}{\itshape #2}}%
  \par\nobreak\smallskip\nobreak
}

% ∗ ∗ ∗ 分隔
\newcommand{\asterisms}{%
  \par\vspace{1.2ex}%
  \noindent\hspace*{2em}$\ast$\hspace{1.2em}$\ast$\hspace{1.2em}$\ast$%
  \par\vspace{1.2ex}%
}

% 无编号子目的目录条目
\newcommand{\tocsubsub}[1]{%
  \phantomsection
  \addcontentsline{toc}{subsubsection}{#1}%
}

% 部扉: 仿原书部题页
\newcommand{\partdivider}[3]{%
  \clearpage
  \thispagestyle{empty}%
  \phantomsection
  \addcontentsline{toc}{part}{#1\quad #3}%
  \vspace*{6.5cm}
  \begin{center}
    {\fontsize{16pt}{26pt}\selectfont\headfont #1\par}
    \vspace{0.9cm}
    {\fontsize{22pt}{34pt}\selectfont\headfont #2\par}
  \end{center}
  \clearpage
}

% 献辞页: 仿原书部题后的题献, 独占一页且不入目录
\newcommand{\dedication}[1]{%
  \thispagestyle{empty}%
  \vspace*{4cm}%
  \begin{center}#1\end{center}%
  \clearpage
}

% --------------------------------------------------------------
%  各卷"译者说明"的共用体例段(各卷的范围段自写, 后接本段)
% --------------------------------------------------------------
\newcommand{\CommonNoteBody}{%
为便于与原书及学界通行的打字稿页码对勘, 本册保留:

\begin{itemize}
  \item 原书全部章节编号(章~1--18, 节如~5.2、18.2, 子目如~18.1.1)与原书小节标题;
  \item 原书脚注的编号与角标——
        \begin{itemize}
          \item 前奏(第~1--4 章)与第一部的注在各页内重新编号, 正文中的上标即该
                页序号; 第一部紧随其后的括注(如“(3)”)是该部分的连续注号, 与第~12
                章“第一部分的注”里的“注~3”一一对应, 注文行首照原书印作“1(3)”式。
                其中三条注原书未标连续注号而改标日期(ch7 注~3 = 30 Septembre、
                ch9 注~11/12 = 8 août), 本册照录其日期;
          \item 第二部第~13--16~章各自重新编号, 分别为~1--41、1--86、1--132、1--29;
                第三部第~17~章自~1~起至~12, 第~18~章自~1~重新起算并延续至~794;
                末段自~\textbf{795(*)}~续至~1053(*)。各段与原书一一对应;
        \end{itemize}
  \item 页边码“p.~1”“p.~A3”“p.~1066”式小注——这是作者引用自己打字稿原稿时
        所用的页码, 与本排印本页码不同。前奏四章的边码各有前缀: 第~1 章作
        \texttt{A1…A6}, 第~2 章作 \texttt{P1…P65}, 第~3 章作 \texttt{L1…L56},
        第~4 章作罗马数字 \texttt{I…XXII}, 自第一部起作纯数字 \texttt{1…171};
  \item 作者用以标记段落的~$\diamond$~记号, 以及标示后来补写文字的上标
        ~\textsuperscript{$\diamond$}~记号;
  \item 原书大小注(“注~44'”“注~46$^1$”“注~171³”等)的编号与层级。
\end{itemize}

第~12 章“第一部分的注”把第~5 至第~11 章各页脚注汇为一编, 并给每条注配上节题
(如“12.19. 耶稣与十二使徒”), 因此部分注文与正文脚注重复——这是原书的体例, 非
重出之误。这一部分属随笔式的自省文字, 同一句法语在正文脚注与章末汇编注中
各有措辞, 本册有意保留各自的译笔, 不强求统一。

除上述标记外, 译文力求贴近作者驳杂、跳跃而辛辣的文风; 数学名词的译法见各卷
术语表(如~\emph{topos}~不译, \emph{faisceaux pervers}~作“反常层”)。文中的日期
括注(如“(5~月~30~日)”“(1984年3月补写)”)为作者历次补写时的纪日, 一律照录。文中
另有一些对注号、页码的引用(如“n°~1695”“打字稿第~227~页”)沿袭作者打字稿的编号
系统, 与本排印本不尽一致, 亦无法在本书中落实, 悉照原文录出。

全套五卷(另有全书合订一册)的分工是:

\begin{itemize}
  \item \textbf{卷一·前奏}——主题的呈现: 或四个乐章的前奏(原书 p.23--122);
  \item \textbf{卷二·第一部}——愚妄与更新(原书 p.123--254);
  \item \textbf{卷三·第二部}——葬礼(I): 或中国皇帝的新衣(原书 p.255--412);
  \item \textbf{卷四·第三部}——葬礼(II): 或阴阳之钥(原书 p.413--632);
  \item \textbf{卷五·第四部}——葬礼(III): 或四种运算(原书 p.633--929)。
\end{itemize}
}
                       % 共用导言区
%      % ================================================================
%  《收获与播种》统一封面 —— TikZ 原生矢量绘制(不依赖任何外部图片)
%  母题: 暮色中的麦田 —— 犁沟向地平线的灭点收束, 前景麦穗剪影,
%        一位播种者把一把种子撒向天空, 种子沿弧线落入田里。
%  用法(在 \begin{document} 之前, 依次):
%      \input{common}                       % 共用导言区
%      \input{cover-art}                    % 本文件
%      \CoverSetup{prelude}{卷一}{前奏}%
%                 {主题的呈现——或四个乐章的前奏}{p.23–122}
%  正文中再 \makecover 即可。配色关键字:
%      prelude / part1 / part2 / part3 / part4 / all
%  五卷共用同一幅画, 只换配色与卷次文字, 以保证"一套书"的观感。
% ================================================================

\RequirePackage{tikz}

% --------------------------------------------------------------
%  配色: 六套(前奏 / 四部 / 全书), 各为同一暮色母题的变调
% --------------------------------------------------------------
\newcommand{\CoverPalette}[1]{%
  \ifstrequal{#1}{prelude}{%
    \definecolor{cvSkyTop}{RGB}{10,26,48}%
    \definecolor{cvSkyMid}{RGB}{46,84,116}%
    \definecolor{cvSkyBot}{RGB}{232,164,92}%
    \definecolor{cvSun}{RGB}{255,224,168}%
    \definecolor{cvFieldTop}{RGB}{150,110,52}%
    \definecolor{cvFieldBot}{RGB}{58,36,18}%
    \definecolor{cvFurrow}{RGB}{40,22,10}%
    \definecolor{cvSil}{RGB}{30,18,10}%
    \definecolor{cvCream}{RGB}{247,240,228}%
  }{%
  \ifstrequal{#1}{part1}{%
    \definecolor{cvSkyTop}{RGB}{22,70,100}%
    \definecolor{cvSkyMid}{RGB}{92,150,168}%
    \definecolor{cvSkyBot}{RGB}{240,200,104}%
    \definecolor{cvSun}{RGB}{255,244,200}%
    \definecolor{cvFieldTop}{RGB}{176,140,58}%
    \definecolor{cvFieldBot}{RGB}{84,54,22}%
    \definecolor{cvFurrow}{RGB}{45,28,12}%
    \definecolor{cvSil}{RGB}{34,22,10}%
    \definecolor{cvCream}{RGB}{250,245,232}%
  }{%
  \ifstrequal{#1}{part2}{%
    \definecolor{cvSkyTop}{RGB}{46,20,30}%
    \definecolor{cvSkyMid}{RGB}{128,52,54}%
    \definecolor{cvSkyBot}{RGB}{226,120,64}%
    \definecolor{cvSun}{RGB}{255,206,150}%
    \definecolor{cvFieldTop}{RGB}{132,84,44}%
    \definecolor{cvFieldBot}{RGB}{48,26,18}%
    \definecolor{cvFurrow}{RGB}{36,16,14}%
    \definecolor{cvSil}{RGB}{26,12,12}%
    \definecolor{cvCream}{RGB}{248,236,224}%
  }{%
  \ifstrequal{#1}{part3}{%
    \definecolor{cvSkyTop}{RGB}{30,20,54}%
    \definecolor{cvSkyMid}{RGB}{92,62,124}%
    \definecolor{cvSkyBot}{RGB}{216,132,140}%
    \definecolor{cvSun}{RGB}{255,214,196}%
    \definecolor{cvFieldTop}{RGB}{128,90,74}%
    \definecolor{cvFieldBot}{RGB}{44,26,30}%
    \definecolor{cvFurrow}{RGB}{34,20,30}%
    \definecolor{cvSil}{RGB}{22,14,22}%
    \definecolor{cvCream}{RGB}{246,238,238}%
  }{%
  \ifstrequal{#1}{part4}{%
    \definecolor{cvSkyTop}{RGB}{12,28,28}%
    \definecolor{cvSkyMid}{RGB}{44,86,72}%
    \definecolor{cvSkyBot}{RGB}{186,168,86}%
    \definecolor{cvSun}{RGB}{240,232,168}%
    \definecolor{cvFieldTop}{RGB}{110,104,50}%
    \definecolor{cvFieldBot}{RGB}{30,34,20}%
    \definecolor{cvFurrow}{RGB}{26,32,20}%
    \definecolor{cvSil}{RGB}{14,20,14}%
    \definecolor{cvCream}{RGB}{240,238,226}%
  }{%
    % 默认: 全书(麦浪金褐)
    \definecolor{cvSkyTop}{RGB}{28,40,32}%
    \definecolor{cvSkyMid}{RGB}{128,110,64}%
    \definecolor{cvSkyBot}{RGB}{232,178,88}%
    \definecolor{cvSun}{RGB}{255,236,180}%
    \definecolor{cvFieldTop}{RGB}{168,124,52}%
    \definecolor{cvFieldBot}{RGB}{62,40,20}%
    \definecolor{cvFurrow}{RGB}{42,26,12}%
    \definecolor{cvSil}{RGB}{28,18,10}%
    \definecolor{cvCream}{RGB}{250,242,226}%
  }}}}}%
}

% --------------------------------------------------------------
%  一键设定:  \CoverSetup{配色}{卷次}{卷名}{副题}{原书页码范围}
% --------------------------------------------------------------
\newcommand{\CoverSetup}[5]{%
  \CoverPalette{#1}%
  \newcommand{\VolNum}{#2}%
  \newcommand{\VolName}{#3}%
  \newcommand{\VolSub}{#4}%
  \newcommand{\VolPages}{#5}%
}

% --------------------------------------------------------------
%  一支麦穗:  \cvWheatEar{x}{底 y}{高}{倾角}
%  茎略弯, 麦粒成对左右斜出, 顶端一芒三针
% --------------------------------------------------------------
\newcommand{\cvWheatEar}[4]{%
  \begin{scope}[shift={(#1,#2)}, rotate=#4]
    \draw[cvSil, line width=0.85mm, line cap=round]
      (0,0) .. controls (0.5,{0.35*#3}) and (0.8,{0.72*#3}) .. ({0.55*#3},#3);
    \foreach \cvK in {0,...,6}{%
      \pgfmathsetmacro{\cvT}{0.50+0.50*\cvK/6}
      \pgfmathsetmacro{\cvGY}{\cvT*#3}
      \pgfmathsetmacro{\cvGX}{0.55*\cvT*\cvT*#3}
      \pgfmathsetmacro{\cvGS}{1.0-0.34*\cvT}
      \begin{scope}[shift={(\cvGX,\cvGY)}, rotate=28]
        \fill[cvSil] (0,{1.9*\cvGS}) ellipse [x radius={0.92*\cvGS}, y radius={2.3*\cvGS}];
      \end{scope}
      \begin{scope}[shift={(\cvGX,\cvGY)}, rotate=-28]
        \fill[cvSil] (0,{-1.9*\cvGS}) ellipse [x radius={0.92*\cvGS}, y radius={2.3*\cvGS}];
      \end{scope}
    }
    \begin{scope}[shift={({0.55*#3},#3)}]
      \fill[cvSil] (0,1.5) ellipse [x radius=0.78, y radius=1.95];
    \end{scope}
    \draw[cvSil, line width=0.26mm, line cap=round]
      ({0.55*#3},{#3+2.3}) -- ({0.55*#3-1.7},{#3+9.5});
    \draw[cvSil, line width=0.26mm, line cap=round]
      ({0.55*#3},{#3+2.3}) -- ({0.55*#3+2.1},{#3+10.5});
    \draw[cvSil, line width=0.26mm, line cap=round]
      ({0.55*#3},{#3+2.3}) -- ({0.55*#3},{#3+11.5});
  \end{scope}
}

% --------------------------------------------------------------
%  封面
% --------------------------------------------------------------
\makeatletter
\renewcommand{\makecover}{%
  \begin{titlepage}
    \thispagestyle{empty}
    \begin{tikzpicture}[remember picture, overlay, x=1mm, y=1mm,
                        shift={(current page.south west)}]
      % ===== 天空 =====
      \shade[top color=cvSkyTop, middle color=cvSkyMid, bottom color=cvSkyBot]
        (0,105) rectangle (210,297);
      % ===== 田野 =====
      \shade[top color=cvFieldTop, bottom color=cvFieldBot]
        (0,0) rectangle (210,105);
      % ===== 落日与光晕 =====
      \begin{scope}
        \clip (0,105) rectangle (210,297);
        \fill[cvSun, opacity=0.14] (152,132) circle (46);
        \fill[cvSun, opacity=0.24] (152,132) circle (30);
      \end{scope}
      \fill[cvSun, opacity=0.95] (152,132) circle (16.5);
      % ===== 犁沟: 向地平线灭点收束 =====
      \begin{scope}
        \clip (0,0) rectangle (210,105);
        \foreach \cvI in {1,...,20}{%
          \pgfmathsetmacro{\cvDX}{3.1*pow(\cvI,1.55)}
          \draw[cvFurrow, opacity=0.30, line width=0.32mm] ({105+\cvDX},0) -- (105,105);
          \draw[cvFurrow, opacity=0.30, line width=0.32mm] ({105-\cvDX},0) -- (105,105);
        }
        \foreach \cvY in {8,18,31,47,68,92}{%
          \draw[cvFurrow, opacity=0.15, line width=0.3mm]
            (0,\cvY) .. controls (60,{\cvY+3.5}) and (150,{\cvY-3}) .. (210,{\cvY+1.5});
        }
      \end{scope}
      % 地平线
      \draw[cvSil, opacity=0.45, line width=0.45mm] (0,105) -- (210,105);
      % ===== 播种者剪影(立于地平线) =====
      \begin{scope}[cvSil]
        \draw[line width=1.9mm, line cap=round, line join=round]
          (55.5,116) -- (50,105)  (55.5,116) -- (60.5,105.5);
        \draw[line width=2.3mm, line cap=round, line join=round]
          (55.5,116) -- (55,123.5);
        \draw[line width=1.6mm, line cap=round, line join=round]
          (55.2,123.2) -- (61.5,126.6) -- (68,129.4);
        \draw[line width=1.5mm, line cap=round, line join=round]
          (55.2,123.2) -- (51.4,119.8) -- (48.6,116.2);
        \fill (55.6,127.3) circle (2.35);
      \end{scope}
      % ===== 种子: 沿贝塞尔弧撒出 =====
      \draw[cvCream, opacity=0.13, line width=0.5mm]
        (68,129.4) .. controls (100,168) and (160,162) .. (196,104);
      % 掌心那一把(尚未撒出)
      \foreach \cvI in {1,...,5}{%
        \pgfmathsetmacro{\cvJX}{1.1*sin(\cvI*53)}
        \pgfmathsetmacro{\cvJY}{1.1*cos(\cvI*83)}
        \fill[cvCream, opacity=0.9] ({68+\cvJX},{129.4+\cvJY}) circle (1.0);
      }
      \foreach \cvI in {1,...,26}{%
        \pgfmathsetmacro{\cvT}{0.035+0.92*\cvI/26}
        \pgfmathsetmacro{\cvU}{1-\cvT}
        \pgfmathsetmacro{\cvSX}{\cvU*\cvU*\cvU*68 + 3*\cvU*\cvU*\cvT*100
                                + 3*\cvU*\cvT*\cvT*160 + \cvT*\cvT*\cvT*196}
        \pgfmathsetmacro{\cvSY}{\cvU*\cvU*\cvU*129.4 + 3*\cvU*\cvU*\cvT*168
                                + 3*\cvU*\cvT*\cvT*162 + \cvT*\cvT*\cvT*104}
        \pgfmathsetmacro{\cvR}{1.55-0.85*\cvT}
        \pgfmathsetmacro{\cvJX}{1.7*sin(\cvI*97)}
        \pgfmathsetmacro{\cvJY}{1.7*cos(\cvI*71)}
        \fill[cvCream, opacity={0.88-0.28*\cvT}] ({\cvSX+\cvJX},{\cvSY+\cvJY}) circle (\cvR);
      }
      % ===== 前景麦穗(两排) =====
      \foreach \cvI in {1,...,11}{%
        \pgfmathsetmacro{\cvWX}{4+19*(\cvI-1)+1.6*sin(\cvI*2.1)}
        \pgfmathsetmacro{\cvWH}{23+7*sin(\cvI*1.7)}
        \pgfmathsetmacro{\cvWL}{3.2*sin(\cvI*2.9)}
        \begin{scope}[opacity=0.45]
          \cvWheatEar{\cvWX}{30}{\cvWH}{\cvWL}
        \end{scope}
      }
      \foreach \cvI in {1,...,10}{%
        \pgfmathsetmacro{\cvWX}{9+20.4*(\cvI-1)+1.5*sin(\cvI*1.3)}
        \pgfmathsetmacro{\cvWH}{41+10*sin(\cvI*2.3)}
        \pgfmathsetmacro{\cvWL}{4.5*sin(\cvI*3.7)}
        \cvWheatEar{\cvWX}{30}{\cvWH}{\cvWL}
      }
      % ===== 压暗一层, 统一色调并保证文字可读 =====
      \fill[black, opacity=0.10] (0,105) rectangle (210,297);
      % ===== 底部信息带 =====
      \fill[cvSil, opacity=0.72] (0,0) rectangle (210,30);
      \draw[cvCream, opacity=0.40, line width=0.4mm] (0,30) -- (210,30);
      % ===== 文字 =====
      \node[anchor=north west, inner sep=0, text=cvCream] at (22,268)
        {\fontsize{38pt}{45pt}\selectfont\headfont \ziju{0.08}\BookTitleCN};
      \node[anchor=north west, inner sep=0, text=cvCream] at (22.6,242)
        {\fontsize{15pt}{20pt}\selectfont\itshape \BookTitleEN};
      \draw[cvCream, opacity=0.55, line width=0.5mm] (22,233) -- (120,233);
      \node[anchor=north west, inner sep=0, text=cvCream] at (22,224)
        {\fontsize{11pt}{15pt}\selectfont 一位数学家对往昔的反思与见证};
      \node[anchor=north west, inner sep=0, text=cvCream, opacity=0.78] at (22.4,208)
        {\fontsize{10pt}{14pt}\selectfont \ziju{0.14}\VolNum};
      \node[anchor=north west, inner sep=0, text=cvCream] at (22,199)
        {\fontsize{30pt}{37pt}\selectfont\headfont \ziju{0.1}\VolName};
      \node[anchor=north west, inner sep=0, text=cvCream, opacity=0.92] at (22,180)
        {\fontsize{13pt}{18pt}\selectfont \VolSub};
      \node[anchor=north west, inner sep=0, text=cvCream] at (22,19)
        {\fontsize{11.5pt}{15pt}\selectfont \OriginalAuthor\quad 著};
      \node[anchor=north west, inner sep=0, text=cvCream, opacity=0.85] at (22,7)
        {\fontsize{9.5pt}{13pt}\selectfont \Translator\ 译\quad ·\quad
         \ifdefempty{\TranslationDate}{\today}{\TranslationDate}};
      \node[anchor=north east, inner sep=0, text=cvCream, opacity=0.85] at (188,13)
        {\fontsize{10pt}{14pt}\selectfont 原书 \VolPages};
    \end{tikzpicture}
  \end{titlepage}%
}
\makeatother
                    % 本文件
%      \CoverSetup{prelude}{卷一}{前奏}%
%                 {主题的呈现——或四个乐章的前奏}{p.23–122}
%  正文中再 \makecover 即可。配色关键字:
%      prelude / part1 / part2 / part3 / part4 / all
%  五卷共用同一幅画, 只换配色与卷次文字, 以保证"一套书"的观感。
% ================================================================

\RequirePackage{tikz}

% --------------------------------------------------------------
%  配色: 六套(前奏 / 四部 / 全书), 各为同一暮色母题的变调
% --------------------------------------------------------------
\newcommand{\CoverPalette}[1]{%
  \ifstrequal{#1}{prelude}{%
    \definecolor{cvSkyTop}{RGB}{10,26,48}%
    \definecolor{cvSkyMid}{RGB}{46,84,116}%
    \definecolor{cvSkyBot}{RGB}{232,164,92}%
    \definecolor{cvSun}{RGB}{255,224,168}%
    \definecolor{cvFieldTop}{RGB}{150,110,52}%
    \definecolor{cvFieldBot}{RGB}{58,36,18}%
    \definecolor{cvFurrow}{RGB}{40,22,10}%
    \definecolor{cvSil}{RGB}{30,18,10}%
    \definecolor{cvCream}{RGB}{247,240,228}%
  }{%
  \ifstrequal{#1}{part1}{%
    \definecolor{cvSkyTop}{RGB}{22,70,100}%
    \definecolor{cvSkyMid}{RGB}{92,150,168}%
    \definecolor{cvSkyBot}{RGB}{240,200,104}%
    \definecolor{cvSun}{RGB}{255,244,200}%
    \definecolor{cvFieldTop}{RGB}{176,140,58}%
    \definecolor{cvFieldBot}{RGB}{84,54,22}%
    \definecolor{cvFurrow}{RGB}{45,28,12}%
    \definecolor{cvSil}{RGB}{34,22,10}%
    \definecolor{cvCream}{RGB}{250,245,232}%
  }{%
  \ifstrequal{#1}{part2}{%
    \definecolor{cvSkyTop}{RGB}{46,20,30}%
    \definecolor{cvSkyMid}{RGB}{128,52,54}%
    \definecolor{cvSkyBot}{RGB}{226,120,64}%
    \definecolor{cvSun}{RGB}{255,206,150}%
    \definecolor{cvFieldTop}{RGB}{132,84,44}%
    \definecolor{cvFieldBot}{RGB}{48,26,18}%
    \definecolor{cvFurrow}{RGB}{36,16,14}%
    \definecolor{cvSil}{RGB}{26,12,12}%
    \definecolor{cvCream}{RGB}{248,236,224}%
  }{%
  \ifstrequal{#1}{part3}{%
    \definecolor{cvSkyTop}{RGB}{30,20,54}%
    \definecolor{cvSkyMid}{RGB}{92,62,124}%
    \definecolor{cvSkyBot}{RGB}{216,132,140}%
    \definecolor{cvSun}{RGB}{255,214,196}%
    \definecolor{cvFieldTop}{RGB}{128,90,74}%
    \definecolor{cvFieldBot}{RGB}{44,26,30}%
    \definecolor{cvFurrow}{RGB}{34,20,30}%
    \definecolor{cvSil}{RGB}{22,14,22}%
    \definecolor{cvCream}{RGB}{246,238,238}%
  }{%
  \ifstrequal{#1}{part4}{%
    \definecolor{cvSkyTop}{RGB}{12,28,28}%
    \definecolor{cvSkyMid}{RGB}{44,86,72}%
    \definecolor{cvSkyBot}{RGB}{186,168,86}%
    \definecolor{cvSun}{RGB}{240,232,168}%
    \definecolor{cvFieldTop}{RGB}{110,104,50}%
    \definecolor{cvFieldBot}{RGB}{30,34,20}%
    \definecolor{cvFurrow}{RGB}{26,32,20}%
    \definecolor{cvSil}{RGB}{14,20,14}%
    \definecolor{cvCream}{RGB}{240,238,226}%
  }{%
    % 默认: 全书(麦浪金褐)
    \definecolor{cvSkyTop}{RGB}{28,40,32}%
    \definecolor{cvSkyMid}{RGB}{128,110,64}%
    \definecolor{cvSkyBot}{RGB}{232,178,88}%
    \definecolor{cvSun}{RGB}{255,236,180}%
    \definecolor{cvFieldTop}{RGB}{168,124,52}%
    \definecolor{cvFieldBot}{RGB}{62,40,20}%
    \definecolor{cvFurrow}{RGB}{42,26,12}%
    \definecolor{cvSil}{RGB}{28,18,10}%
    \definecolor{cvCream}{RGB}{250,242,226}%
  }}}}}%
}

% --------------------------------------------------------------
%  一键设定:  \CoverSetup{配色}{卷次}{卷名}{副题}{原书页码范围}
% --------------------------------------------------------------
\newcommand{\CoverSetup}[5]{%
  \CoverPalette{#1}%
  \newcommand{\VolNum}{#2}%
  \newcommand{\VolName}{#3}%
  \newcommand{\VolSub}{#4}%
  \newcommand{\VolPages}{#5}%
}

% --------------------------------------------------------------
%  一支麦穗:  \cvWheatEar{x}{底 y}{高}{倾角}
%  茎略弯, 麦粒成对左右斜出, 顶端一芒三针
% --------------------------------------------------------------
\newcommand{\cvWheatEar}[4]{%
  \begin{scope}[shift={(#1,#2)}, rotate=#4]
    \draw[cvSil, line width=0.85mm, line cap=round]
      (0,0) .. controls (0.5,{0.35*#3}) and (0.8,{0.72*#3}) .. ({0.55*#3},#3);
    \foreach \cvK in {0,...,6}{%
      \pgfmathsetmacro{\cvT}{0.50+0.50*\cvK/6}
      \pgfmathsetmacro{\cvGY}{\cvT*#3}
      \pgfmathsetmacro{\cvGX}{0.55*\cvT*\cvT*#3}
      \pgfmathsetmacro{\cvGS}{1.0-0.34*\cvT}
      \begin{scope}[shift={(\cvGX,\cvGY)}, rotate=28]
        \fill[cvSil] (0,{1.9*\cvGS}) ellipse [x radius={0.92*\cvGS}, y radius={2.3*\cvGS}];
      \end{scope}
      \begin{scope}[shift={(\cvGX,\cvGY)}, rotate=-28]
        \fill[cvSil] (0,{-1.9*\cvGS}) ellipse [x radius={0.92*\cvGS}, y radius={2.3*\cvGS}];
      \end{scope}
    }
    \begin{scope}[shift={({0.55*#3},#3)}]
      \fill[cvSil] (0,1.5) ellipse [x radius=0.78, y radius=1.95];
    \end{scope}
    \draw[cvSil, line width=0.26mm, line cap=round]
      ({0.55*#3},{#3+2.3}) -- ({0.55*#3-1.7},{#3+9.5});
    \draw[cvSil, line width=0.26mm, line cap=round]
      ({0.55*#3},{#3+2.3}) -- ({0.55*#3+2.1},{#3+10.5});
    \draw[cvSil, line width=0.26mm, line cap=round]
      ({0.55*#3},{#3+2.3}) -- ({0.55*#3},{#3+11.5});
  \end{scope}
}

% --------------------------------------------------------------
%  封面
% --------------------------------------------------------------
\makeatletter
\renewcommand{\makecover}{%
  \begin{titlepage}
    \thispagestyle{empty}
    \begin{tikzpicture}[remember picture, overlay, x=1mm, y=1mm,
                        shift={(current page.south west)}]
      % ===== 天空 =====
      \shade[top color=cvSkyTop, middle color=cvSkyMid, bottom color=cvSkyBot]
        (0,105) rectangle (210,297);
      % ===== 田野 =====
      \shade[top color=cvFieldTop, bottom color=cvFieldBot]
        (0,0) rectangle (210,105);
      % ===== 落日与光晕 =====
      \begin{scope}
        \clip (0,105) rectangle (210,297);
        \fill[cvSun, opacity=0.14] (152,132) circle (46);
        \fill[cvSun, opacity=0.24] (152,132) circle (30);
      \end{scope}
      \fill[cvSun, opacity=0.95] (152,132) circle (16.5);
      % ===== 犁沟: 向地平线灭点收束 =====
      \begin{scope}
        \clip (0,0) rectangle (210,105);
        \foreach \cvI in {1,...,20}{%
          \pgfmathsetmacro{\cvDX}{3.1*pow(\cvI,1.55)}
          \draw[cvFurrow, opacity=0.30, line width=0.32mm] ({105+\cvDX},0) -- (105,105);
          \draw[cvFurrow, opacity=0.30, line width=0.32mm] ({105-\cvDX},0) -- (105,105);
        }
        \foreach \cvY in {8,18,31,47,68,92}{%
          \draw[cvFurrow, opacity=0.15, line width=0.3mm]
            (0,\cvY) .. controls (60,{\cvY+3.5}) and (150,{\cvY-3}) .. (210,{\cvY+1.5});
        }
      \end{scope}
      % 地平线
      \draw[cvSil, opacity=0.45, line width=0.45mm] (0,105) -- (210,105);
      % ===== 播种者剪影(立于地平线) =====
      \begin{scope}[cvSil]
        \draw[line width=1.9mm, line cap=round, line join=round]
          (55.5,116) -- (50,105)  (55.5,116) -- (60.5,105.5);
        \draw[line width=2.3mm, line cap=round, line join=round]
          (55.5,116) -- (55,123.5);
        \draw[line width=1.6mm, line cap=round, line join=round]
          (55.2,123.2) -- (61.5,126.6) -- (68,129.4);
        \draw[line width=1.5mm, line cap=round, line join=round]
          (55.2,123.2) -- (51.4,119.8) -- (48.6,116.2);
        \fill (55.6,127.3) circle (2.35);
      \end{scope}
      % ===== 种子: 沿贝塞尔弧撒出 =====
      \draw[cvCream, opacity=0.13, line width=0.5mm]
        (68,129.4) .. controls (100,168) and (160,162) .. (196,104);
      % 掌心那一把(尚未撒出)
      \foreach \cvI in {1,...,5}{%
        \pgfmathsetmacro{\cvJX}{1.1*sin(\cvI*53)}
        \pgfmathsetmacro{\cvJY}{1.1*cos(\cvI*83)}
        \fill[cvCream, opacity=0.9] ({68+\cvJX},{129.4+\cvJY}) circle (1.0);
      }
      \foreach \cvI in {1,...,26}{%
        \pgfmathsetmacro{\cvT}{0.035+0.92*\cvI/26}
        \pgfmathsetmacro{\cvU}{1-\cvT}
        \pgfmathsetmacro{\cvSX}{\cvU*\cvU*\cvU*68 + 3*\cvU*\cvU*\cvT*100
                                + 3*\cvU*\cvT*\cvT*160 + \cvT*\cvT*\cvT*196}
        \pgfmathsetmacro{\cvSY}{\cvU*\cvU*\cvU*129.4 + 3*\cvU*\cvU*\cvT*168
                                + 3*\cvU*\cvT*\cvT*162 + \cvT*\cvT*\cvT*104}
        \pgfmathsetmacro{\cvR}{1.55-0.85*\cvT}
        \pgfmathsetmacro{\cvJX}{1.7*sin(\cvI*97)}
        \pgfmathsetmacro{\cvJY}{1.7*cos(\cvI*71)}
        \fill[cvCream, opacity={0.88-0.28*\cvT}] ({\cvSX+\cvJX},{\cvSY+\cvJY}) circle (\cvR);
      }
      % ===== 前景麦穗(两排) =====
      \foreach \cvI in {1,...,11}{%
        \pgfmathsetmacro{\cvWX}{4+19*(\cvI-1)+1.6*sin(\cvI*2.1)}
        \pgfmathsetmacro{\cvWH}{23+7*sin(\cvI*1.7)}
        \pgfmathsetmacro{\cvWL}{3.2*sin(\cvI*2.9)}
        \begin{scope}[opacity=0.45]
          \cvWheatEar{\cvWX}{30}{\cvWH}{\cvWL}
        \end{scope}
      }
      \foreach \cvI in {1,...,10}{%
        \pgfmathsetmacro{\cvWX}{9+20.4*(\cvI-1)+1.5*sin(\cvI*1.3)}
        \pgfmathsetmacro{\cvWH}{41+10*sin(\cvI*2.3)}
        \pgfmathsetmacro{\cvWL}{4.5*sin(\cvI*3.7)}
        \cvWheatEar{\cvWX}{30}{\cvWH}{\cvWL}
      }
      % ===== 压暗一层, 统一色调并保证文字可读 =====
      \fill[black, opacity=0.10] (0,105) rectangle (210,297);
      % ===== 底部信息带 =====
      \fill[cvSil, opacity=0.72] (0,0) rectangle (210,30);
      \draw[cvCream, opacity=0.40, line width=0.4mm] (0,30) -- (210,30);
      % ===== 文字 =====
      \node[anchor=north west, inner sep=0, text=cvCream] at (22,268)
        {\fontsize{38pt}{45pt}\selectfont\headfont \ziju{0.08}\BookTitleCN};
      \node[anchor=north west, inner sep=0, text=cvCream] at (22.6,242)
        {\fontsize{15pt}{20pt}\selectfont\itshape \BookTitleEN};
      \draw[cvCream, opacity=0.55, line width=0.5mm] (22,233) -- (120,233);
      \node[anchor=north west, inner sep=0, text=cvCream] at (22,224)
        {\fontsize{11pt}{15pt}\selectfont 一位数学家对往昔的反思与见证};
      \node[anchor=north west, inner sep=0, text=cvCream, opacity=0.78] at (22.4,208)
        {\fontsize{10pt}{14pt}\selectfont \ziju{0.14}\VolNum};
      \node[anchor=north west, inner sep=0, text=cvCream] at (22,199)
        {\fontsize{30pt}{37pt}\selectfont\headfont \ziju{0.1}\VolName};
      \node[anchor=north west, inner sep=0, text=cvCream, opacity=0.92] at (22,180)
        {\fontsize{13pt}{18pt}\selectfont \VolSub};
      \node[anchor=north west, inner sep=0, text=cvCream] at (22,19)
        {\fontsize{11.5pt}{15pt}\selectfont \OriginalAuthor\quad 著};
      \node[anchor=north west, inner sep=0, text=cvCream, opacity=0.85] at (22,7)
        {\fontsize{9.5pt}{13pt}\selectfont \Translator\ 译\quad ·\quad
         \ifdefempty{\TranslationDate}{\today}{\TranslationDate}};
      \node[anchor=north east, inner sep=0, text=cvCream, opacity=0.85] at (188,13)
        {\fontsize{10pt}{14pt}\selectfont 原书 \VolPages};
    \end{tikzpicture}
  \end{titlepage}%
}
\makeatother
                    % 本文件
%      \CoverSetup{prelude}{卷一}{前奏}%
%                 {主题的呈现——或四个乐章的前奏}{p.23–122}
%  正文中再 \makecover 即可。配色关键字:
%      prelude / part1 / part2 / part3 / part4 / all
%  五卷共用同一幅画, 只换配色与卷次文字, 以保证"一套书"的观感。
% ================================================================

\RequirePackage{tikz}

% --------------------------------------------------------------
%  配色: 六套(前奏 / 四部 / 全书), 各为同一暮色母题的变调
% --------------------------------------------------------------
\newcommand{\CoverPalette}[1]{%
  \ifstrequal{#1}{prelude}{%
    \definecolor{cvSkyTop}{RGB}{10,26,48}%
    \definecolor{cvSkyMid}{RGB}{46,84,116}%
    \definecolor{cvSkyBot}{RGB}{232,164,92}%
    \definecolor{cvSun}{RGB}{255,224,168}%
    \definecolor{cvFieldTop}{RGB}{150,110,52}%
    \definecolor{cvFieldBot}{RGB}{58,36,18}%
    \definecolor{cvFurrow}{RGB}{40,22,10}%
    \definecolor{cvSil}{RGB}{30,18,10}%
    \definecolor{cvCream}{RGB}{247,240,228}%
  }{%
  \ifstrequal{#1}{part1}{%
    \definecolor{cvSkyTop}{RGB}{22,70,100}%
    \definecolor{cvSkyMid}{RGB}{92,150,168}%
    \definecolor{cvSkyBot}{RGB}{240,200,104}%
    \definecolor{cvSun}{RGB}{255,244,200}%
    \definecolor{cvFieldTop}{RGB}{176,140,58}%
    \definecolor{cvFieldBot}{RGB}{84,54,22}%
    \definecolor{cvFurrow}{RGB}{45,28,12}%
    \definecolor{cvSil}{RGB}{34,22,10}%
    \definecolor{cvCream}{RGB}{250,245,232}%
  }{%
  \ifstrequal{#1}{part2}{%
    \definecolor{cvSkyTop}{RGB}{46,20,30}%
    \definecolor{cvSkyMid}{RGB}{128,52,54}%
    \definecolor{cvSkyBot}{RGB}{226,120,64}%
    \definecolor{cvSun}{RGB}{255,206,150}%
    \definecolor{cvFieldTop}{RGB}{132,84,44}%
    \definecolor{cvFieldBot}{RGB}{48,26,18}%
    \definecolor{cvFurrow}{RGB}{36,16,14}%
    \definecolor{cvSil}{RGB}{26,12,12}%
    \definecolor{cvCream}{RGB}{248,236,224}%
  }{%
  \ifstrequal{#1}{part3}{%
    \definecolor{cvSkyTop}{RGB}{30,20,54}%
    \definecolor{cvSkyMid}{RGB}{92,62,124}%
    \definecolor{cvSkyBot}{RGB}{216,132,140}%
    \definecolor{cvSun}{RGB}{255,214,196}%
    \definecolor{cvFieldTop}{RGB}{128,90,74}%
    \definecolor{cvFieldBot}{RGB}{44,26,30}%
    \definecolor{cvFurrow}{RGB}{34,20,30}%
    \definecolor{cvSil}{RGB}{22,14,22}%
    \definecolor{cvCream}{RGB}{246,238,238}%
  }{%
  \ifstrequal{#1}{part4}{%
    \definecolor{cvSkyTop}{RGB}{12,28,28}%
    \definecolor{cvSkyMid}{RGB}{44,86,72}%
    \definecolor{cvSkyBot}{RGB}{186,168,86}%
    \definecolor{cvSun}{RGB}{240,232,168}%
    \definecolor{cvFieldTop}{RGB}{110,104,50}%
    \definecolor{cvFieldBot}{RGB}{30,34,20}%
    \definecolor{cvFurrow}{RGB}{26,32,20}%
    \definecolor{cvSil}{RGB}{14,20,14}%
    \definecolor{cvCream}{RGB}{240,238,226}%
  }{%
    % 默认: 全书(麦浪金褐)
    \definecolor{cvSkyTop}{RGB}{28,40,32}%
    \definecolor{cvSkyMid}{RGB}{128,110,64}%
    \definecolor{cvSkyBot}{RGB}{232,178,88}%
    \definecolor{cvSun}{RGB}{255,236,180}%
    \definecolor{cvFieldTop}{RGB}{168,124,52}%
    \definecolor{cvFieldBot}{RGB}{62,40,20}%
    \definecolor{cvFurrow}{RGB}{42,26,12}%
    \definecolor{cvSil}{RGB}{28,18,10}%
    \definecolor{cvCream}{RGB}{250,242,226}%
  }}}}}%
}

% --------------------------------------------------------------
%  一键设定:  \CoverSetup{配色}{卷次}{卷名}{副题}{原书页码范围}
% --------------------------------------------------------------
\newcommand{\CoverSetup}[5]{%
  \CoverPalette{#1}%
  \newcommand{\VolNum}{#2}%
  \newcommand{\VolName}{#3}%
  \newcommand{\VolSub}{#4}%
  \newcommand{\VolPages}{#5}%
}

% --------------------------------------------------------------
%  一支麦穗:  \cvWheatEar{x}{底 y}{高}{倾角}
%  茎略弯, 麦粒成对左右斜出, 顶端一芒三针
% --------------------------------------------------------------
\newcommand{\cvWheatEar}[4]{%
  \begin{scope}[shift={(#1,#2)}, rotate=#4]
    \draw[cvSil, line width=0.85mm, line cap=round]
      (0,0) .. controls (0.5,{0.35*#3}) and (0.8,{0.72*#3}) .. ({0.55*#3},#3);
    \foreach \cvK in {0,...,6}{%
      \pgfmathsetmacro{\cvT}{0.50+0.50*\cvK/6}
      \pgfmathsetmacro{\cvGY}{\cvT*#3}
      \pgfmathsetmacro{\cvGX}{0.55*\cvT*\cvT*#3}
      \pgfmathsetmacro{\cvGS}{1.0-0.34*\cvT}
      \begin{scope}[shift={(\cvGX,\cvGY)}, rotate=28]
        \fill[cvSil] (0,{1.9*\cvGS}) ellipse [x radius={0.92*\cvGS}, y radius={2.3*\cvGS}];
      \end{scope}
      \begin{scope}[shift={(\cvGX,\cvGY)}, rotate=-28]
        \fill[cvSil] (0,{-1.9*\cvGS}) ellipse [x radius={0.92*\cvGS}, y radius={2.3*\cvGS}];
      \end{scope}
    }
    \begin{scope}[shift={({0.55*#3},#3)}]
      \fill[cvSil] (0,1.5) ellipse [x radius=0.78, y radius=1.95];
    \end{scope}
    \draw[cvSil, line width=0.26mm, line cap=round]
      ({0.55*#3},{#3+2.3}) -- ({0.55*#3-1.7},{#3+9.5});
    \draw[cvSil, line width=0.26mm, line cap=round]
      ({0.55*#3},{#3+2.3}) -- ({0.55*#3+2.1},{#3+10.5});
    \draw[cvSil, line width=0.26mm, line cap=round]
      ({0.55*#3},{#3+2.3}) -- ({0.55*#3},{#3+11.5});
  \end{scope}
}

% --------------------------------------------------------------
%  封面
% --------------------------------------------------------------
\makeatletter
\renewcommand{\makecover}{%
  \begin{titlepage}
    \thispagestyle{empty}
    \begin{tikzpicture}[remember picture, overlay, x=1mm, y=1mm,
                        shift={(current page.south west)}]
      % ===== 天空 =====
      \shade[top color=cvSkyTop, middle color=cvSkyMid, bottom color=cvSkyBot]
        (0,105) rectangle (210,297);
      % ===== 田野 =====
      \shade[top color=cvFieldTop, bottom color=cvFieldBot]
        (0,0) rectangle (210,105);
      % ===== 落日与光晕 =====
      \begin{scope}
        \clip (0,105) rectangle (210,297);
        \fill[cvSun, opacity=0.14] (152,132) circle (46);
        \fill[cvSun, opacity=0.24] (152,132) circle (30);
      \end{scope}
      \fill[cvSun, opacity=0.95] (152,132) circle (16.5);
      % ===== 犁沟: 向地平线灭点收束 =====
      \begin{scope}
        \clip (0,0) rectangle (210,105);
        \foreach \cvI in {1,...,20}{%
          \pgfmathsetmacro{\cvDX}{3.1*pow(\cvI,1.55)}
          \draw[cvFurrow, opacity=0.30, line width=0.32mm] ({105+\cvDX},0) -- (105,105);
          \draw[cvFurrow, opacity=0.30, line width=0.32mm] ({105-\cvDX},0) -- (105,105);
        }
        \foreach \cvY in {8,18,31,47,68,92}{%
          \draw[cvFurrow, opacity=0.15, line width=0.3mm]
            (0,\cvY) .. controls (60,{\cvY+3.5}) and (150,{\cvY-3}) .. (210,{\cvY+1.5});
        }
      \end{scope}
      % 地平线
      \draw[cvSil, opacity=0.45, line width=0.45mm] (0,105) -- (210,105);
      % ===== 播种者剪影(立于地平线) =====
      \begin{scope}[cvSil]
        \draw[line width=1.9mm, line cap=round, line join=round]
          (55.5,116) -- (50,105)  (55.5,116) -- (60.5,105.5);
        \draw[line width=2.3mm, line cap=round, line join=round]
          (55.5,116) -- (55,123.5);
        \draw[line width=1.6mm, line cap=round, line join=round]
          (55.2,123.2) -- (61.5,126.6) -- (68,129.4);
        \draw[line width=1.5mm, line cap=round, line join=round]
          (55.2,123.2) -- (51.4,119.8) -- (48.6,116.2);
        \fill (55.6,127.3) circle (2.35);
      \end{scope}
      % ===== 种子: 沿贝塞尔弧撒出 =====
      \draw[cvCream, opacity=0.13, line width=0.5mm]
        (68,129.4) .. controls (100,168) and (160,162) .. (196,104);
      % 掌心那一把(尚未撒出)
      \foreach \cvI in {1,...,5}{%
        \pgfmathsetmacro{\cvJX}{1.1*sin(\cvI*53)}
        \pgfmathsetmacro{\cvJY}{1.1*cos(\cvI*83)}
        \fill[cvCream, opacity=0.9] ({68+\cvJX},{129.4+\cvJY}) circle (1.0);
      }
      \foreach \cvI in {1,...,26}{%
        \pgfmathsetmacro{\cvT}{0.035+0.92*\cvI/26}
        \pgfmathsetmacro{\cvU}{1-\cvT}
        \pgfmathsetmacro{\cvSX}{\cvU*\cvU*\cvU*68 + 3*\cvU*\cvU*\cvT*100
                                + 3*\cvU*\cvT*\cvT*160 + \cvT*\cvT*\cvT*196}
        \pgfmathsetmacro{\cvSY}{\cvU*\cvU*\cvU*129.4 + 3*\cvU*\cvU*\cvT*168
                                + 3*\cvU*\cvT*\cvT*162 + \cvT*\cvT*\cvT*104}
        \pgfmathsetmacro{\cvR}{1.55-0.85*\cvT}
        \pgfmathsetmacro{\cvJX}{1.7*sin(\cvI*97)}
        \pgfmathsetmacro{\cvJY}{1.7*cos(\cvI*71)}
        \fill[cvCream, opacity={0.88-0.28*\cvT}] ({\cvSX+\cvJX},{\cvSY+\cvJY}) circle (\cvR);
      }
      % ===== 前景麦穗(两排) =====
      \foreach \cvI in {1,...,11}{%
        \pgfmathsetmacro{\cvWX}{4+19*(\cvI-1)+1.6*sin(\cvI*2.1)}
        \pgfmathsetmacro{\cvWH}{23+7*sin(\cvI*1.7)}
        \pgfmathsetmacro{\cvWL}{3.2*sin(\cvI*2.9)}
        \begin{scope}[opacity=0.45]
          \cvWheatEar{\cvWX}{30}{\cvWH}{\cvWL}
        \end{scope}
      }
      \foreach \cvI in {1,...,10}{%
        \pgfmathsetmacro{\cvWX}{9+20.4*(\cvI-1)+1.5*sin(\cvI*1.3)}
        \pgfmathsetmacro{\cvWH}{41+10*sin(\cvI*2.3)}
        \pgfmathsetmacro{\cvWL}{4.5*sin(\cvI*3.7)}
        \cvWheatEar{\cvWX}{30}{\cvWH}{\cvWL}
      }
      % ===== 压暗一层, 统一色调并保证文字可读 =====
      \fill[black, opacity=0.10] (0,105) rectangle (210,297);
      % ===== 底部信息带 =====
      \fill[cvSil, opacity=0.72] (0,0) rectangle (210,30);
      \draw[cvCream, opacity=0.40, line width=0.4mm] (0,30) -- (210,30);
      % ===== 文字 =====
      \node[anchor=north west, inner sep=0, text=cvCream] at (22,268)
        {\fontsize{38pt}{45pt}\selectfont\headfont \ziju{0.08}\BookTitleCN};
      \node[anchor=north west, inner sep=0, text=cvCream] at (22.6,242)
        {\fontsize{15pt}{20pt}\selectfont\itshape \BookTitleEN};
      \draw[cvCream, opacity=0.55, line width=0.5mm] (22,233) -- (120,233);
      \node[anchor=north west, inner sep=0, text=cvCream] at (22,224)
        {\fontsize{11pt}{15pt}\selectfont 一位数学家对往昔的反思与见证};
      \node[anchor=north west, inner sep=0, text=cvCream, opacity=0.78] at (22.4,208)
        {\fontsize{10pt}{14pt}\selectfont \ziju{0.14}\VolNum};
      \node[anchor=north west, inner sep=0, text=cvCream] at (22,199)
        {\fontsize{30pt}{37pt}\selectfont\headfont \ziju{0.1}\VolName};
      \node[anchor=north west, inner sep=0, text=cvCream, opacity=0.92] at (22,180)
        {\fontsize{13pt}{18pt}\selectfont \VolSub};
      \node[anchor=north west, inner sep=0, text=cvCream] at (22,19)
        {\fontsize{11.5pt}{15pt}\selectfont \OriginalAuthor\quad 著};
      \node[anchor=north west, inner sep=0, text=cvCream, opacity=0.85] at (22,7)
        {\fontsize{9.5pt}{13pt}\selectfont \Translator\ 译\quad ·\quad
         \ifdefempty{\TranslationDate}{\today}{\TranslationDate}};
      \node[anchor=north east, inner sep=0, text=cvCream, opacity=0.85] at (188,13)
        {\fontsize{10pt}{14pt}\selectfont 原书 \VolPages};
    \end{tikzpicture}
  \end{titlepage}%
}
\makeatother

\CoverSetup{all}{全书}{合订本}{前奏 · 第一部 · 第二部 · 第三部 · 第四部}{p.23–929}

\hypersetup{
  pdftitle={收获与播种(全书)},
  pdfauthor={Alexandre Grothendieck / 数译 译},
  pdfsubject={Récoltes et Semailles — Réflexions et témoignage sur un passé de mathématicien}
}

% --------------------------------------------------------------
\begin{document}

\frontmatter
\makecover

% ---------------------------------------------------------------
%  译者说明
% ---------------------------------------------------------------
\chapter*{译者说明}
\markboth{译者说明}{译者说明}
\addcontentsline{toc}{chapter}{译者说明}

本册是《收获与播种》中文译本的全书合订, 连续收录原书第~23 至第~929 页
(至全书末行)的全部译文:

\begin{itemize}
  \item \textbf{前奏}——“主题的呈现: 或四个乐章的前奏”(Présentation des Thèmes
        ou Prélude en quatre Mouvements), 覆盖原书第~23 至第~122 页, 含第~1 至
        第~4 章: 权作序言、巡游一部作品(或孩子与母亲)、一封信、引言;
  \item \textbf{第一部}——“愚妄与更新”(Première partie: Fatuité et
        Renouvellement), 覆盖原书第~123 至第~254 页, 含第~5 至第~12 章: 劳动与
        发现、梦与做梦者、双重面孔、师与生、收获、孩子嬉戏、孤独的历险, 以及
        “第一部分的注”;
  \item \textbf{第二部}——“葬礼(I):或中国皇帝的新衣”(Deuxième partie:
        L'Enterrement (I) ou la robe de l'Empereur de Chine), 覆盖原书第~255 至
        第~412 页(第~13 至第~16 章);
  \item \textbf{第三部}——“葬礼(II):或阴阳之钥”(Troisième partie:
        L'Enterrement (II) ou la Clef du Yin et du Yang), 覆盖原书第~413 页起,
        收第~17~章“亡故者(仍未死去……)”及第~18~章的小节~18.1--18.2;
  \item \textbf{第四部}——“葬礼(III):或四种运算”(Quatrième partie:
        L'Enterrement (III) ou les Quatre Opérations), 自原书第~633~页部扉起,
        收小节~18.3--18.5.4~注~(c1)“记忆的失灵——或新编历史”第四段, 并续至
        全书末行(原书第~929 页)。
\end{itemize}

\CommonNoteBody

\makecontents

\mainmatter

% ================================================================
%  一 · 前奏 (第 1-4 章, 原书 p.23-122)
% ================================================================
\partdivider{前奏}{主题的呈现\\[0.35em]{\normalsize 或四个乐章的前奏}}{主题的呈现——或四个乐章的前奏}

\setcounter{chapter}{0}
\markboth{收获与播种·前奏}{}

% ================================================================
%  第 1 章 权作序言 (原书 p.27-30, 边码 A1-A6)
%  源文: _work/src1/ch01.md
% ================================================================
\chapter{权作序言}

% ---------------- p.27 ----------------

\noindent\pmark 1986年1月30日\par

\origpage{A1} 只差一篇序言没写,就可以把《收获与播种》交给印刷厂了。我发誓,我是怀着世上最好的意愿,想写出点像样的东西来的。这一回,要写点合乎情理的东西。三四页,不能再多,但要写得有分量,用来引荐这头一千多页的庞然大物。要写点能“抓住”厌倦了的读者的东西,让他隐约看到:在这令人生畏的“一千多页”里,也许有些叫他感兴趣的东西(甚至,谁知道呢,与他有关的东西?)。“抓人”不是我的风格,不是。可这一回,我偏要破一次例!那位“疯到敢冒这个险的出版商”(敢出版这头显然无法出版的巨兽),总得勉强收回成本吧。

结果,没成。我尽力了。而且不止一个下午,不像我原先打算的那样三下两下就完事。到明天,整整三个星期了,我一直在写它,纸页越堆越高。写出来的东西,肯定不是人们体面地称之为“序言”的那种。又一次搞砸了,没治!我这把年纪,改不了了——我也不是那块料,不是叫卖的料,也不是让人叫卖的料。即便只为讨人喜欢(讨自己喜欢,也讨朋友喜欢……)。

写出来的,是一种长长的、带评注的“巡游”,穿过我作为数学家的作品。这巡游首先是给“门外汉”的——给那个“对数学一窍不通”的人。也是为我自己,我此前从未有过这样一趟巡游的闲情。一针引一线,我发现自己被引到一些向来留在未言说之中的东西面前,并且把它们说了出来。说来也巧,这些也正是我觉得最本质的东西——在我的工作中,在我的作品里。这些东西没有一点技术性。我这份天真的“传达”是否成功,由你评判——这尝试本身大概也有点疯。我的满足与快乐,将是让你感受到了它们。这些东西,我那些博学的同行里,许多人已经不再感受得到了。也许他们变得太博学、太有名望了。这往往让人失去与单纯而本质之物的接触。

在这趟“巡游一部作品”中,我也稍稍谈到我的一生。也在这里那里,稍稍谈到《收获与播种》讲的是什么。在“巡游”之后的那封“信”(注明去年五月)里,我还会更详细地再谈。这封信原本是写给我从前的学生们,以及数学界里我“昔日的朋友”的。但它同样没有一点技术性。任何读者,只要想通过一份“现场”的叙述,了解究竟哪些前因后果最终促使我写下《收获与播种》,都可以毫无困难地读它。比起“巡游”,它更能让你预先尝到“大世界”——数学的“大世界”——里某种气氛的味道。也(和“巡游”一样)尝到我的表达风格的味道,据说有点特别。还有借这风格表达出来的那种精神——这种精神同样不是人人都欣赏的。

\origpage{A2} 在“巡游”里,以及在《收获与播种》各处,我谈数学工作。这工作我熟悉,是第一手的。我所说的多数东西,对一切创造性工作、一切发现的工作,肯定都是真的。至少对所谓“智力”工作是真的——那种主要“靠脑子”和靠写来进行的工作。这样的工作,其标志是理解的绽放与成熟,即对我们正在探究之物的理解。% ---------------- p.28 ----------------
但举一个处于另一极的例子:爱的激情,同样也是发现的冲动。它为我们打开一种所谓“肉体的”认识,这种认识同样在更新、在舒展、在加深。这两种冲动——不妨说,一种驱动着工作中的数学家,另一种驱动着情人——彼此之近,远过于人们通常的猜度,也远过于人们愿意向自己承认的。我希望《收获与播种》的这些页面,能在你的工作中、在你每日的生活里,帮你感受到这一点。

在巡游中,主要谈的是数学工作本身。相反,对于这工作所处的境况,对于在真正的工作时间之外起作用的那些动机,我几乎一言不发。这有可能给人的形象——给我的形象,或给数学家、给一般“科学家”的形象——添上一层固然讨喜、却已走样的光晕。像是“伟大而高尚的激情”,没有任何修正。总之,合于那位大大的“科学神话”(S 请大写!)的路数。那英雄的、“普罗米修斯式”的神话,文人与学者们争先恐后地掉了进去(如今还在掉)。也许只有史家有时能抵挡这诱人的神话。实情是:在“科学家”的种种动机里——那些有时驱使他毫不吝惜地投入工作的动机——野心与虚荣所起的作用,与在任何别的行当里一样重大,而且几乎同样普遍。其形式或粗或细,因人而异。我绝不认为自己是个例外。读我这份见证,我希望,在这一点上不会留下任何疑问。

\origpage{A3} 也确实,最贪婪的野心也无力发现哪怕一条数学陈述,更无力证明它——正如它(譬如)无力“让人勃起”(就字面意义而言)。无论男女,使人“勃起”的绝不是野心,不是想要出风头、想要展示一种力量的欲望——那力量在此是性方面的——恰恰相反!而是对某种既强有力、又极真实、同时又极细微之物的敏锐知觉。人们可以称之为“美”,而美只是那东西的千张面孔之一。有野心,并不必然妨碍人偶尔感到一个存在、一件事物的美,这我同意。但可以肯定的是:让我们感到这美的,不是野心……

第一个发现并驾驭了火的人,正是和你我一样的人。绝不是人们以“英雄”“半神”之类名目设想的那种人。肯定,像你像我一样,他尝过焦虑的咬啮,也尝过那种虚荣的药膏——它让人忘掉咬啮。但在他“认识”火的那一刻,既没有恐惧,也没有虚荣。英雄神话里的真相就在于此。当神话被用来替我们遮掩事物的另一面——同样真实、同样本质的另一面——它就变得寡淡无味,变成了药膏。

我在《收获与播种》里的用意,是谈这两个方面——谈求知的冲动,也谈恐惧及其虚荣的解药。我相信自己“理解”、至少是认识这冲动及其本性。(也许有一天我会惊奇地发现,自己原来错得多么厉害……)但说到恐惧与虚荣,以及由此派生的对创造力的暗中阻塞,我很清楚自己并未深入这大谜的底。我也不知道,在我余下的岁月里,是否会看到这奥秘的底……

在写作《收获与播种》的过程中,浮现出两个形象,分别代表人类这场冒险的两个方面。那就是孩子(亦名“工匠”),以及老板。在即将开始的巡游里,要谈的几乎只是“孩子”。副题“孩子与母亲”里的,也是他。这名字,我希望,会在巡游途中渐渐明朗起来。

\origpage{A4} 而在反思其余的全部篇幅里,相反,占据前台的主要是老板。他可不是白当老板的! % ---------------- p.29 ----------------
更确切地说,不是一个老板,而是些相互竞争的企业的老板们。不过也确实:所有老板在本质上都彼此相像。而一旦开始谈老板,也就意味着要出现“坏蛋”了。在反思的第(I)部分(“愚妄与更新”,紧接本篇导论性的部分——即“四个乐章的前奏”——之后),那个“坏蛋”主要是我。在其后三个部分里,则主要是“别人”。轮着来!

这就是说,除了深刻的哲学思考与(毫无悔意的)“忏悔”,还会有“硫酸肖像”(借用我一位同行兼朋友的说法——他自己也挨了几下……)。更不必说那些大手笔、毫不含糊的“手术”。罗贝尔·若兰(Robert Jaulin)\fnote{1}{}{罗贝尔·若兰是老朋友了。我似乎明白,面对民族学界的建制,他处在一个“白狼”般的地位,与我在数学“上流社会”面前的处境略有相似。}向我保证(半开玩笑地):我在《收获与播种》里做的是“数学界的民族志”(也许是社会学,我也说不准了)。当然,听说自己(连自己都不知道)做着高深学问,总是受用的!事实是:在反思的“调查”部分里(而且违背我的意愿……),我看见我正在写下的那些纸页上,数学界建制的一大半人列队走过,更不必说许多地位更卑微的同行与朋友。而在最近这几个月里,自从去年十月我把《收获与播种》的试印本寄出去之后,这事又“来了一遍”。我这份见证,真是像一块石头扔进了死水塘。回响各色各样,几乎什么调子都有(只除了无聊那一种……)。几乎每一次,都完全出乎我的预料。此外还有许多沉默,意味深长。显然,关于我那些从前的学生、以及处境或高或低的同行们脑子里在转什么,我过去(如今也还)有许多要学,而且是形形色色的功课——请原谅,我是说“数学界的社会学”!对所有已经前来为我晚年这部社会学巨著添砖加瓦的人,我要在此表达我的感激之情。

当然,对那些暖色调的回响,我格外敏感。也有少数几位同行,向我道出了他们的感动,或一种(此前未表达过的)危机感、衰败感——对他们自认所属的这个数学界内部的危机与衰败。

\origpage{A5} 在这个圈子之外,最早对我的见证报以热烈、甚至激动回应的人里,我想在此举出西尔维·谢瓦莱与卡特琳·谢瓦莱(Sylvie et Catherine Chevalley)\fnote{2}{}{西尔维与卡特琳·谢瓦莱是克劳德·谢瓦莱的遗孀与女儿;这部《收获与播种》的核心部分(ReS III,“阴阳之钥”)正是献给他的。反思中多处谈到他,以及他在我的历程中所扮演的角色。}、罗贝尔·若兰、斯泰凡·德利若尔日(Stéphane Deligeorge)、克里斯蒂安·布尔茹瓦(Christian Bourgois)。如果说《收获与播种》将获得比最初那份(只给极狭小圈子的)试印本更广的流传,那主要归功于他们。尤其归功于他们那种有感染力的确信:我竭力去把握、去说出的东西,是应当被说出的;而且它能够被听到的范围,比我那些同行(往往阴沉,甚至乖戾,毫不愿意自我质疑……)的圈子更大。于是克里斯蒂安·布尔茹瓦毫不犹豫地承担了出版这“无法出版之物”的风险,而斯泰凡·德利若尔日则给我这份难以下咽的见证以殊荣,把它收进“Epistémé”丛书,(眼下)与牛顿、居维叶、阿拉戈并列。(我想不出更好的同伴了!)对每一位在我格外“敏感”的时刻一再表示同情与信任的人,我很高兴在此说出我全部的谢意。

于是我们这就出发,开始一趟巡游一部作品的旅程,作为穿越一生的旅行的开场。是的,一趟漫长的旅行,一千多页,每一页都塞得满满。这旅行我花了一辈子去做,也没做完,又花了一年多一页一页地重新发现它。有时词语迟迟不肯来,来表达一份仍在躲闪、尚未被把握的经验的全部汁液——就像压榨机里堆得满满的、熟透而饱满的葡萄,有时似乎想躲开紧攥着它的那股力量……但即便在词语仿佛互相推挤、汩汩涌出的时刻,它们的推挤与涌出也绝不是信手拈来。每一个词都在途中被称过,或者事后被称过,若发现太轻或太重,就仔细调准。所以这份反思—见证—旅行,不是为着让人一天读一点、一个月读完的,也不是为着那个急于读到“最后一句话”的读者。

% ---------------- p.30 ----------------
《收获与播种》里没有“最后一句话”,没有“结论”,正如我的一生里、你的一生里也没有。有的是一坛酒,在我这个人的木桶里陈了一辈子。你喝下的最后一杯不会比第一杯或第一百杯更好。它们全是“同一杯”,又各不相同。而倘若第一杯坏了,整桶都坏了;那还不如喝好水(若找得到),别喝坏酒。

\origpage{A6} 但好酒不宜匆匆喝,也不宜站着喝。

% ================================================================
%  第 2 章 巡游一部作品——或孩子与母亲 (原书 p.31-53, 边码 P1-P37)
%  源文: _work/src1/task-ch02a.md
%  说明: 章内 Contents 目录不排印
% ================================================================
\chapter{巡游一部作品——或孩子与母亲}

\noindent 1986年1月\par

\section{(1) 事物的魔法}

% ---------------- p.31 ----------------
\origpage{P1} 小时候,我很喜欢上学。读和写、算术、唱歌(他拉一把小提琴给我们伴奏),还有史前人类和火的发现,都是同一位老师教我们。我不记得那时候我们在学校里曾经觉得无聊。他掌握着数字的魔法,也掌握着词语、符号和声音的魔法。还有韵脚的魔法,在歌谣里,在小小的诗里。韵脚里似乎藏着某种超出词语之外的奥秘。这情形一直持续到有一天
% ---------------- p.32 ----------------
有人向我解释,那不过是个再简单不过的“窍门”:所谓押韵,无非是让两个前后相接的说话片段以同一个音节收尾,于是它们就像中了魔法一样,变成了诗句。这可是一场启示!在家里,我身边有人接应,一连几个星期、几个月,我都玩着作诗。有一阵子,我说话只押韵了。幸好这毛病过去了。可即便今天,偶尔我还会写写诗——只是不大再去找韵脚了,除非它自己跑来。

另有一回,一个比我大的伙伴,已经在读中学了,教了我负数。那是另一种很好玩的游戏,但很快玩腻了。还有填字游戏——我花了好几天、好几个星期来编,越编越交错复杂。这游戏里结合着形式的魔法,以及符号和词语的魔法。但这份热情也离开了我,似乎没留下什么痕迹。

在中学里,先是在德国读第一年,后来在法国,我是个不错的学生,但算不上“出色的学生”。凡是我最感兴趣的,我不计工本地投入;不太感兴趣的,就往往被撂在一边,也不太在乎那位“老师”怎么评价。在法国上中学的第一年,1940年,我和母亲一起被关在芒德附近的里厄克罗集中营。那是战争年代,我们是外国人——用当时的话说,是“不受欢迎的人”。但集中营的管理当局对营里的孩子睁一只眼闭一只眼,尽管他们同样是“不受欢迎的人”。我们进出多少随自己的便。我年纪最大,也是唯一去上中学的,学校在四五公里外,不管下雪还是刮风,脚上穿着总是进水的破鞋。\origpage{P2} 我还记得第一次“数学考试”,老师给了我一个坏分数,为的是证明“三角形全等的三种情形”之一。我的证明不是书上那个,而他是虔诚地照书办的。可我心里很清楚,我的证明并不比书上那个更有说服力,也不更没说服力——我遵循的正是书里的精神,用的还是那套老掉牙的“把某个图形这样挪到另一个图形上去”。显然,教我的这个人觉得自己没有能力凭自己的眼光来判断(在这里,就是判断一个推理是否成立)。他必须求助于一种权威,此处即一本书的权威。这些做派一定深深触动了我,我才会记得这桩小事。此后直到今天,我有很多机会看到:这样的做派绝不是例外,而是几乎普遍的规则。这个话题有很多可说的——在《收获与播种》里,我不止一次以这种或那种形式触及它。可直到今天,不管我愿不愿意,每次再撞上它,我都感到手足无措……

战争最后几年,母亲仍被关在集中营里,我则待在“瑞士救援会”办的一所儿童之家,收容难民儿童,在利尼翁河畔勒尚邦。我们大多是犹太人;每当接到通知(由当地警察来通知)说盖世太保要来搜捕,我们就到树林里躲上一两夜,三三两两一小群,并不太意识到这真的关系到我们的性命。那一带塞文山区藏满了犹太人,许多人靠当地居民的互助活了下来。

在“塞文中学”(我在那里受教育)最让我吃惊的,是我的同学们对所学的东西兴趣之少。至于我,学年一开始就把课本狼吞虎咽地读完,心想这一回总算要学到真正有趣的东西了;而这一年余下的时间,我就尽量自己打发,与此同时,预定的教学大纲一个学期接一个学期,无情地照本宣科。不过我们的老师都很好。博物学老师弗里德尔先生,人品与学识都极出色。可他不会“镇压”,被人起哄到死,以致到了学年末,课已经没法再上下去了,他那无力的声音
% ---------------- p.33 ----------------
被一片喧闹淹没。说不定正因为这个,我才没成为生物学家!

\origpage{P3} 我相当多的时间——甚至上课的时候(嘘……)——都在做数学题。很快,书上的题就不够我做了。也许是因为它们久而久之彼此太相像了;但我相信,尤其是因为它们来得太像从天而降,就这么一个接一个排着队,不说自己从哪儿来、往哪儿去。那是书上的问题,不是我的问题。其实真正自然的问题并不缺。比如,已知三角形三边之长 $a$、$b$、$c$,这个三角形就确定了(不计其位置),那么就该有一个显式的“公式”,例如把三角形的面积表示成 $a$、$b$、$c$ 的函数。四面体也一样:已知六条棱的长,体积是多少?这一回我想我大概费了不少劲,但硬啃下去,终究还是做出来了。反正,一件事“抓住”了我,我就不数在上面花了多少小时、多少天,哪怕把别的全都忘了!(如今也还是这样……)

我们的数学书里最让我不满意的,是对长度(一条曲线的)、面积(一个曲面的)、体积(一个立体的)这些概念,完全没有认真的定义。我向自己保证,一旦有闲,就要补上这个空缺。1945到1948年间,我在蒙彼利埃大学念书,把最精壮的精力都花在了这上面。学院的课不合我的意。虽然从未对自己明说,我大概觉得老师们不过是在重复他们的书,就像芒德中学那位第一位数学老师一样。所以我也只是隔很久才去一趟学院,为的是了解那个老掉牙的“教学大纲”讲到了哪里。就那大纲而言,书是够用的;但同样清楚的是,它们根本不回答我给自己提的问题。说实话,它们连看都没看见那些问题,就像我的中学课本也没看见一样。只要它们随手给出计算长度、面积、体积的配方,靠单重、双重、三重积分(超过三维的情形则谨慎地绕开……),那么为它们给出一个内蕴的定义这回事,似乎就根本不成为问题——对我的老师们如此,对教科书的作者们也如此。

\origpage{P4} 以我当时有限的经验看,完全可能显得我是世界上唯一对数学问题怀有好奇心的人。无论如何,在那些年里,这就是我未曾说出口的信念;那些年我处于完全的智识孤独之中,而这份孤独并不使我难受\fnote{1}{}{1945到1948年间,我和母亲住在蒙彼利埃十来公里外的一个小村子里,迈拉尔格(属旺达尔格乡),四野是葡萄园中的一处僻静地。(父亲1942年在奥斯维辛失踪。)我们靠我那份微薄的学生助学金过活,过得很拮据。为了收支相抵,我每年去摘葡萄,摘完葡萄再做些“拾穗酒”,勉强脱手(据说这违反了当时的法律……)。此外还有一处园子,从不用侍弄,却大量供给我们无花果、菠菜,甚至(到最后)还有番茄——那是一位好心的邻居种下的,种在一片绚丽的罂粟正中间。日子过得美——只是有时紧得刚好差那么一点,比如要换一副眼镜架,或者一双穿到露出麻线的鞋。幸而母亲在集中营里待了那么久,身体虚弱多病,可以享受免费医疗。否则我们永远付不起看医生的钱……}。说实话,我想那时候我从未深究过:自己是否真是世上唯一可能对我所做的事感兴趣的人。我的精力已充分投入那场自己给自己出的挑战:发展一套能让我完全满意的理论。

我丝毫没有怀疑自己终能做到、终能找到事物的最终答案——只要我肯花力气去细究它们,把它们对我说的随说随记到纸上。比如说,体积的直觉是无可置疑的。它只能是一个现实的映照,这现实眼下还难以捉摸,但完全可靠。要把握的就是这个现实,如此而已——也许有点像有一天“韵脚”那个神奇的现实被把握住、“被理解”那样。

% ---------------- p.34 ----------------
动手的时候,我十七岁,刚出中学,我以为这不过是几个星期的事。结果我在上面待了三年。我还想方设法,硬是在大学二年级期末考砸了一门——球面三角(在“深入天文学”选修方向里,原文如此),因为一个愚蠢的数值计算错误。(得承认,一离开中学,我的计算就一直不怎么行……)正因如此,我才不得不在蒙彼利埃再待第三年,在那里读完学士学位,而不能马上去巴黎——据人向我保证,那是唯一有机会遇到懂行之人的地方,懂数学里什么才算重要。给我通风报信的苏拉先生还向我保证:数学里最后剩下的一些问题,二三十年前已由一位名叫勒贝格的人解决了。他恰好发展了一套测度与积分理论,(说真的,巧得可笑!)给数学画上了句号。

\origpage{P5} 苏拉先生教我们“微积分”,是个和善的人,对我一片好意。我并不认为他因此说服了我。我心里大概早已预感到:数学在广度与深度上都是无边无际的。大海有“句号”吗?总之,我从未动过念头去把苏拉先生提到的勒贝格那本书找来,而他自己大概也从未把这本书拿在手里过。在我心里,一本书里可能装的东西,和我为了满足自己对那些令我不解之事的求知欲、照自己的方式所做的工作,毫无共同之处。

\section{(2) 独处的重要}

一两年后,我终于在巴黎与数学界有了接触,最终在那里得知(在许许多多别的之外):我在自己角落里用凑合的手段做出来的工作,(几乎)就是“人人”都熟知的东西,叫做勒贝格测度与积分理论。在我向其谈起这项工作(甚至给他们看一份手稿)的两三位年长者眼里,这有点像我白白浪费了时间,重做了一遍“已经知道的东西”。我也不记得自己有过失望。那时候,为自己所做的工作去获取什么“信誉”、哪怕只是一句赞许或别人一点兴趣,这个念头大概还离我的头脑很远。何况我的精力已经够忙的了:要熟悉一个完全不同的圈子,尤其是要学会巴黎人眼中数学家的 ABC\fnote{2}{}{我在《收获与播种》第一部分(ReS I)“受欢迎的异乡人”一节(第9号)里,对这段有点粗粝的过渡时期作了简短的叙述。}。

然而,现在回想这三年,我意识到它们绝非白费。在不知不觉中,我在孤独里学到了数学家这一行当最本质的东西——那是任何老师都无法真正教给人的。我从未需要对自己说明,也从未遇到一个可以分享我求知之渴的人,可我却知道,用俗话说,“打骨子里”知道:我是一个数学家——一个在完全意义上“做”数学的人,就像人们“做”爱那样。数学对我来说,成了一位永远迎候我欲望的情人。这些孤独的岁月奠定了一份从未动摇过的信赖的基础——无论是二十岁到巴黎时发现自己的无知有多么深、要学的东西有多么浩瀚;还是二十多年后,我永远离开数学界时那些动荡的插曲;还是最近这些年,由我昔日最亲近的伙伴们操办的那场针对我本人和我的作品的、相当疯狂的“葬礼”
% ---------------- p.35 ----------------
(提前举行,而且办得干净利落)……\origpage{P6}

换一种说法:在这些关键的年月里,我学会了独处\fnote{3}{}{这个说法有些不确切。我从来不需要“学会独处”,原因很简单:在我的童年里,我从未丢掉出生时就在我身上、也就在每个人身上的那种天生能力。但这三年孤独的工作——在其中我得以按自己那套自发的严格要求的标准,给自己量了尺寸——却在我身上确认并重新奠下了一份信赖与安然的确据,这一次是在我与数学工作的关系上;这份确据与那些立法的共识和风尚毫无关系。我在《根与孤独》一注(ReS IV,第1713号,尤其第1080页)里有机会再次提到它。}。我的意思是:凭自己的眼光去接近我想认识的事物,而不是信赖那些明说的或默认的观念与共识——它们来自我自认为属于的某个或大或小的群体,或由于任何别的理由而被赋予对我的权威。中学也好,大学也好,那些无声的共识曾告诉我:对“体积”这个概念本身,没有必要提出问题,它被说成“人所共知”“显而易见”“没有问题”。我越过了它,当作理所当然——正如几十年前勒贝格想必也曾越过它。正是在这种“越过”的行动里,在总之做自己、而不只是做那些立法共识的表达式里,在不被它们为我们划定的那个强制性圈子关在里面——首先是在这种孤独的行动里,才有“创造”。其余的一切都是额外来的。

后来,在接纳我的这个数学界里,我有机会遇到许多人,无论是年长者,还是与我年纪相仿的年轻人,他们显然比我聪明得多、也“有天赋”得多。我欣赏他们学新概念时那种轻松,像玩一样,拿来耍弄,仿佛从摇篮里就认识它们——而我却觉得自己笨重迟钝,像一只鼹鼠,在由那些(人们向我保证)我理当学会的东西堆成的一座不成形的山里面艰难地开路,而那些东西的来龙去脉我根本抓不住。事实上,我一点也没有那种出色学生的样子:轻松通过名校会考,三下两下就消化掉那些吓人的教学大纲。

\origpage{P7} 何况,我那些更出色的同学大多成了有本事、有名望的数学家。然而,隔了三十或三十五年回头看,我发现他们并没有在我们这个时代的数学上留下真正深刻的印记。他们在已经现成的环境里做了些事情,有时是漂亮的事情,而那环境他们从没想过要去触动。他们不知不觉地成了那些看不见而专横的圈子的囚徒——那些圈子在某个圈子、某个时代里界定出一个“宇宙”。要跨越它们,他们本应在自己身上重新找到出生时即有的那种能力,那种我也有的能力:独处的能力。

小孩子呢,独处毫无困难。他天生就是孤独的,虽然偶尔有人做伴他也不讨厌,而且知道到了吃奶的钟点就吵着要妈妈的“咪咪”。他也清楚地知道,不用谁来告诉他,那“咪咪”是给他的,而他会吃。可我们常常已经和心里的那个孩子失去了联系。我们不断地从最好的东西旁边走过,不屑于看它一眼……

如果说在《收获与播种》里我除了对自己说话,还对别的什么人说,那绝不是对一个“公众”说。我在那里是对你、对读我的你说话,把你当作一个人,一个孤独的人。我想说话的,是你身上那个懂得独处的人,那个孩子,而不是别人。那孩子常常很远,我很清楚。他什么没见过,而且早就见惯了。他不知躲到哪里去了,要走到他跟前,常常不容易。人家会赌咒说,他早就死了,或者干脆从未存在过——可我确信他就在某处,活得好好的。

我也知道我被听懂的标志是什么。那就是:超越一切文化与命运的差异
% ---------------- p.36 ----------------
,我所说的关于我本人、我的生活的话,在你那里得到回声与共鸣;当你在其中也重新找到你自己的生活、你自己对自己的经验,也许还是你此前未曾留意的一面。这不是什么“认同”,不是去认同与你相距遥远的什么东西或什么人。而也许是,你稍微地,通过我对自身生活的重新发现,也重新发现了你自己的生活,那离你最近的东西——顺着《收获与播种》一页一页,直到我今天此刻正在写下的这些页面。\origpage{P8}

\section{(3) 内在的冒险——或神话与见证}

首先,《收获与播种》是一场对我自己、对我的生活的反思。因此它也是一份见证,而且是从两方面来说的见证。它是对我过去的见证,反思的主要重量正落在这过去上。但同时它也是对最切近的当下的见证——对我此刻正在写作的见证,对《收获与播种》这些页面在一个个钟点、一个个夜晚和白天里诞生的此刻的见证。这些页面是一场关于我的生活的长久沉思的忠实见证,那沉思如它实际所经历的那样延续着(并且此刻仍在延续……)。

这些页面没有文学上的野心。它们是一份关于我自己的文献。我只在极狭窄的限度内才允许自己动它(比如偶尔作些文体上的修饰)\fnote{4}{}{因此,对错误(事实上的,或视角上的,等等)的可能的更正,并不是对初稿加以修补的机会,而是放在页底注里,或在日后对所述情境的“回访”中进行。}。若说有什么野心,那也只是:是真的。而这已经很多了。

此外,这份文献一点也不是“自传”。你在里面不会知道我的出生日期(它除了用来画一张星象图,恐怕没什么意思),也不会知道我母亲和父亲的名字、他们以什么为生,不会知道那位曾是我的妻子的人和其他在我生命中有过重要位置的女人的名字,也不会知道这些爱情中生下的孩子们的名字,以及他们各自把自己的生活过成了什么样。这并不是因为这些事情在我生命里不重要,也不是因为它们今天已不再重要。而是因为,这场对我的反思既已如此开始并延续下去,我从未在某一刻感到被促使去稍稍着手描述这些我在这里那里掠过的事情,更不会去认真地罗列名字和数字。从未有一刻,我会觉得这能给当时我所追索的“话题”增添什么。(而在前面那几页里,我却像不由自主似的,写进了也许比接下来一千页里还要多的、关于我生活的物质细节……)

如果你问我,这千页一路追索的“话题”究竟是什么,我会回答:是讲述、并由此发现我一生所是、且现在仍是的那场内在冒险。这份关于一场冒险的叙述—见证,同时在我刚才说的两个层面上展开。一个是对过去那场冒险的勘探,从它的根与起源一直追溯到我的童年。另一个是这“同一场”冒险的延续与更新,在我写《收获与播种》的一个个瞬间和一天天里,作为对来自外部世界的一声猛烈呼唤\fnote{5}{}{关于这声“猛烈呼唤”的细节,见《信》,特别是第3至8节。}的自发回应。\origpage{P9}

外部的事实为反思提供养料,但只在这样的限度内:它们激起并引发内在冒险的又一次跃起,或有助于照亮它。我的数学作品被埋葬、被洗劫这件事(后面会详谈),就是这样一次挑衅。它在我心里激起了强大的自我反应的全员出动,同时又向我揭示了那些深的、此前未被察觉的纽带,它们仍把我与从我这里生出的作品连在一起。

% ---------------- p.37 ----------------
\origpage{P10} 确实,我属于“数学尖子”这件事,并不必然成为你对我这段特殊冒险感兴趣的理由(更不是好理由)——我改了圈子、改了生活方式之后跟同行们闹出的麻烦,也不是理由。而且并不缺少同行,甚至朋友,他们觉得当众摊开自己的“心境”(用他们的话说)是最可笑不过的。要紧的是“成果”。“灵魂”呢,也就是说我们身上那个经历着这些“成果”的“生产”、以及它的各种余波(既在“生产者”的生活里,也在他的同类的生活里)的东西,却遭人轻视,甚至被公然嘲弄。这种态度自命为“谦虚”的表现;我在其中看到的却是一种逃避的征兆,一种古怪的失序,而且是我们呼吸的空气本身所助长的。我当然不是为那种被对自己的暗暗蔑视所击中的人写作的——那蔑视使他看不上我能给他的最好的东西。一种对真正构成他自己生活之物的蔑视,对构成我生活之物的蔑视:那些驱动着心灵的浅层与深层运动,粗糙的或细微的,那正是经历着经验、并对之作出反应的“灵魂”,它僵住或舒展,它退缩或在学习……

一场内在冒险的叙述,只能由经历它的人来做,别人不行。但即便这叙述只给自己看,它也很少不滑进建造神话的车辙里,而那神话的主角就是叙述者本人。这样的神话,并非诞生于一个民族、一种文化的创造想象,而是诞生于那个不敢承担一份谦卑现实的人的虚荣,他乐于用一件构造物、一件他心智的作品去替换它。但一份真实的叙述(如果有的话),关于一场真正被那样经历过的冒险,是珍贵之物。而这并不因为围绕叙述者的某种威望(不论有理没理),而只因为它存在,带着它的真之品质。这样的见证是珍贵的,无论它来自一位有名望甚至显赫的人物,还是来自一个没有前途、拖家带口的小职员,还是来自一个普通刑事犯。

这样一份叙述若对别人有什么功效,那首先是通过对另一个人经验的不加粉饰的见证,让他重新面对自己。或者(换个说法)也许是抹去他心中(哪怕只在读一次书那么长的时间里)那份蔑视——他借它来看待自己的冒险,以及那个既是乘客又是船长的“灵魂”……

\section{(4) 风俗画}

说到我作为数学家的过去,后来又(像违心地)发现我那场巨大“葬礼”的种种波折与内幕,我不由自主地——并非刻意——画出某个圈子、某个时代的画卷:一个某些价值正在瓦解的时代,那些价值曾赋予人的劳动以意义。这就是“风俗画”的一面,围绕一桩在“科学”编年史上大概是独一无二的“社会新闻”画成。我前面说的话,我想已经足够清楚:你在《收获与播种》里找不到一份关于某桩不寻常“事件”的“卷宗”,好让你很快了解情况。可是我有一位朋友,一直在找那份卷宗,他闭着眼睛走过,什么也没看见,几乎错过了构成《收获与播种》的实质与血肉的一切。

正如我在《信》里详细得多地解释的,“调查”(或“风俗画”)主要在第二部和第四部展开:《葬礼(一)——或中国皇帝的新衣》和《葬礼(三)——或四种运算》。一页一页地,我在那里执拗地把一大批(至少可以说)有汁水的事实一件一件挖出来,又尽量随挖随“安顿”。渐渐地,这些事实聚成一幅总体画面,它逐步从雾中显出
% ---------------- p.38 ----------------
,颜色越来越鲜,轮廓越来越清晰。在这些逐日写下的注里,刚出现的“原始事实”与个人的回忆、与心理的、哲学的、乃至(偶尔)数学的评论和思考,不可分割地混在一起。就是这样,我也无能为力!

以我做过的那项工作(它让我绷紧了一年多)为基础,用“调查结论”的体例编一份卷宗,应该是几个钟头或几天的额外工作,视有兴趣的读者的好奇与要求而定。我确实试过一回,要编那份著名的卷宗。那是我开始写一个本该叫《四种运算》的注的时候\fnote{6}{}{那条原计划的注最终炸开成了《收获与播种》第四部(同名《四种运算》)里的大约70条注,足足铺开四百来页。}。可后来,不行,一点办法也没有。我做不到!这确实不是我的表达方式,到了晚年更是如此。而且我现在认为,有《收获与播种》,我为“数学共同体”做的已经够多了,可以毫无悔意地把编那份该编的“卷宗”的事,留给我以外的人(如果我的同行中有谁觉得自己与此有关)。\origpage{P11}

\section{(5) 继承人与建造者}

现在该在这里说几句我的数学作品了——它在我生命里占据了一个重要的位置,而且(令我自己惊讶)至今仍占着。《收获与播种》里我不止一次回到这部作品——有时用每个人都能清楚明白的方式,有时则用多少带点技术性的词语\fnote{7}{}{除了关于我过去作品的数学概览,这里那里还有一些含新数学发展的段落。最长的一篇是 ReS IV 第171(ix)号注里的《五张照片(晶体与 D-模)》。}。后一类段落,在很大程度上会“越过”不止“门外汉”的头顶,连对其中所谈的数学已经不怎么“在行”的同行数学家也一样。你当然可以干脆跳过那些你觉得太“专业”的段落。你也可以把它们浏览一遍,也许顺带捕捉到数学事物世界那“神秘的美”(这是一位非数学家朋友来信中的说法)的一线反光——它们像一座座“不可企及的奇异小岛”,在反思的浩瀚流动的水面上浮现……

大多数数学家,我刚才说过,倾向于把自己圈在一个概念框架里、一个一劳永逸地定下来的“宇宙”里——基本上就是他们求学时发现已经“现成”的那个。他们像一座宽敞漂亮的、一切都安置好了的房子的继承人:有客厅、厨房、作坊,有成套的锅碗瓢盆和随手可用的工具,凭这些,天哪,足够做饭和做点手工了。这房子在几代人手里是怎样一点点盖起来的,某些工具(而不是别的……)是怎样、为什么被设计和打造出来的,为什么这间屋子这样布置、那间那样安排——这些问题是这些继承人从不会想到去问的。这就是“宇宙”,就是必须生活在其中的“给定之物”,如此而已!它显得很大(而且多数时候,离把它的所有房间都走遍还远着呢),同时又很熟悉,尤其是:不可变易。他们忙活,是为了维护和装点一份家产:修一张摇晃的柜子,给外墙抹灰,磨快一件工具,甚至有时候,对比较有事业心的来说,在作坊里从头做一件新家具。而他们全心投入时,这件家具也可能美极了,整座房子也显得更美了。\origpage{P12}

更罕见的是,他们中有人会想到去改动一件备用的工具,或者甚至在需求的反复催促下,去设想并造出一件新的。这样做的时候
% ---------------- p.39 ----------------
,他几乎要连声道歉,因为他觉得自己这样有点像冒犯了家族传统应有的虔敬,觉得自己用一桩古怪的创新把它冲撞了。

房子里大多数房间的窗和百叶窗都仔细关着——大概是怕有一股别处来的风灌进来。而当漂亮的新家具这里一件、那里一件,还不算子嗣,开始把变得狭窄的房间塞满,一直侵占到走廊时,这些继承人里没有一个愿意意识到:他那熟悉而舒适的宇宙开始有点紧得转不开身了。与其承认这一点,他们各人宁愿勉勉强强地侧身挤过去、卡在那里:有人挤在路易十五式碗柜和藤编摇椅之间,有人挤在流鼻涕的娃娃和埃及石棺之间,最后还有一位,在无计可施之下,尽其所能地爬上一堆杂七杂八、摇摇欲坠的椅子和长凳……

我刚画的这幅小画,并非数学界所特有。它说明的是一些根深蒂固、由来已久的制约,在所有圈子和人类活动的一切领域里都能遇到,而且(据我所知)在所有社会、所有时代都是如此。我此前已有机会提及,我也绝不声称自己可以免于它。我的见证将表明,实情恰恰相反。只不过,在智力创造活动这个相对有限的层面上,我受这种制约的触动相当少\fnote{8}{}{我认为主要原因在于某种有利的氛围,它环绕着我的童年,直到我五岁。参见《无辜》一注(ReS III,第107号)。}——它可以叫做“文化盲”:看不见(也无法活动于)周围文化所固定的那个“宇宙”之外。

至于我,我觉得自己属于那样一脉数学家:他们自发的天职与快乐,是不断地建造新房子\fnote{9}{}{这个“待建之屋”的原型意象,第一次浮现并被表述出来,是在《阴,仆役,与新主子们》一注中(ReS III,第135号)。}。一路之上\origpage{P13},他们也忍不住随时发明并打造一切必需的工具、器具、家具和仪器,既为了从地基到屋脊把房子建起来,也为了给未来的厨房和作坊备足东西,把房子布置得能住、能舒服。然而,一旦连最后一条檐沟、最后一只凳子都安放好了,这位工匠却很少在这些地方久留——这里的每一块石头、每一根椽子,都带着加工并安放它的那只手的痕迹。他的位置不在那些现成宇宙的安宁里,无论它们多么宜人、多么和谐——无论是他亲手布置的,还是他的前人们布置的。别的任务已经在召唤他去新的工地:在某种需要的专横推动下,那需要也许只有他一个人清楚地感到,或者(更常见的是)他预感到某种只有他一个人预感到的需要。他的位置在露天里。他是风的朋友,不怕独自干活,一干几个月、几年,必要时干一辈子,除非有可喜的接替者来搭把手。他当然和所有人一样只有两只手——但那是两只随时能猜到自己该做什么的手,既不嫌最粗笨的活,也不嫌最精细的活,而且从不厌倦去认识、反复认识那些无数次召唤它们去认识的事物。两只手也许太少,因为世界是无穷的。它们永远也穷尽不了它!然而,两只手又很多……

我这个人历史不行,但若要我举出这一脉数学家的名字,我自然而然想到伽罗瓦和黎曼(上世纪),以及希尔伯特(本世纪初)。若要在当初接纳我进入数学界的年长者中找一位代表\fnote{10}{}{我在《受欢迎的异乡人》一节(ReS I,第9号)里谈到这些开端。},我最先想到的是让·勒雷(Jean Leray)的名字,尽管我和他的接触
% ---------------- p.40 ----------------
始终极为零星\fnote{11}{}{不过这并不妨碍我(继 H. 嘉当与 J.-P. 塞尔之后)成为勒雷引入的某个重大革新概念——层——的主要使用者和推广者之一,而层贯穿我整个几何学家的工作,是其中一件本质的工具。也正是它给了我钥匙,把(拓扑)空间概念扩展为 topos 概念,后面会谈到。\par 此外,在我看来,勒雷与我勾画的“建造者”肖像有所不同:他似乎并不倾向于“从地基到屋脊建造房子”。不如说,他忍不住在无人想得到的地方开出了广阔的地基,却把完成它们、在上面建造、以及房子建成后住进去(哪怕只住一阵)的事留给别人……}。

我刚才大致勾画了两幅肖像:一位是“恋家”的数学家,满足于维护和装点一份遗产;一位是建造者—拓荒者\fnote{12}{}{我刚刚偷偷地、“从侧面”把两个带有雄性回响的修饰语(即“建造者”和“拓荒者”)并置在一起,而它们表达的却是发现冲动相当不同的方面,其本性比这些名字所能唤起的更精细。这将在这次巡游—反思的后文、“发现母亲——或两个山坡”(第17号)那一步里显现。},他忍不住不断越过那些界定一个宇宙的“看不见而专横的圈子”\fnote{13}{}{而且,他也就这样在无意中给这个旧宇宙(不是为他自己,至少是为那些不如他好动的同类)指定了新的界限,划在新的圈子里——这些圈子当然更广,但和它们所取代的那些一样看不见、一样专横。}。也可以用两个有点一刀切、但颇有意味的名字称他们:“守成者”与“革新者”。两者在同一个延续于世代、世纪、千年的集体冒险中,都有存在的理由和要扮演的角色。在一门科学或一门艺术的繁盛期,这两种气质之间没有对立,也没有敌意\fnote{14}{}{尤其在数学界,在我亲身见证的时期(1948–1969),情况就是如此,当时我自己也是这个世界的一部分。1970年我离开之后,似乎出现了一种大范围的反应,一种对“想法”总体的“蔑视的共识”,尤其针对我引入的那些重大革新想法。}。它们不同,而且互相补充,就像面团与酵母互相补充。\origpage{P14}

在这两个极端类型之间(它们在本性上丝毫不相反对),当然有一整套中间气质。某位“恋家者”从不想离开一处熟悉的住所,更不会去揽下到天知道什么地方另建一座的活计;然而当住处真的开始变窄时,他会毫不犹豫地拿起抹刀,收拾一间地窖或阁楼,加高一层,必要时还在墙边接出一间规模不大的新附属房\fnote{15}{}{我的大多数“前辈”(例如《一笔情愿的债》导论第10节里谈到的那些)都属于这种中间气质。我特别想到亨利·嘉当、克劳德·谢瓦莱、安德烈·韦伊、让–皮埃尔·塞尔、洛朗·施瓦茨。除也许韦伊之外,他们对我眼看要登上的那些孤独冒险,都投以“善意的目光”,没有“暗暗的担忧或责备”。}。他骨子里不是建造者,但常常会带着善意的眼光——至少是不带暗暗的担忧与责备——看着另一位曾与他同住一座房子的人,如今
% ---------------- p.41 ----------------
在某个穷乡僻壤辛苦地聚拢梁木和石头,一副仿佛已经看见宫殿的神气……\origpage{P15}

\section{(6) 视角与眼光}

但还是回到我自己和我的作品上来。

如果说我在数学家的技艺上有所出众,那与其说靠解题的灵巧与坚持(解前人传下的问题),不如说靠我身上这种天然的倾向:它驱使我看见一些显然关键、却无人看见的问题,或者提取出那些缺失的“好概念”(往往在新概念出现之前,谁也没意识到它们缺失),以及那些谁也没想到的“好陈述”。常常,概念与陈述配合得如此完美,使我心里不可能怀疑它们是对的(顶多需要些修补)——于是往往,当这不是要发表的“部件活”时,我就省得再往前走,省得花时间去把证明弄妥;而那证明,一旦陈述和它的语境看清了,多半只是个“手艺”问题,甚至可以说是例行公事。值得注意的事情数不清,不可能把每一个的召唤都跟到底!尽管如此,在我已写出并发表的作品里,正式证明过的命题和定理数以千计,我想我可以说,除了极少的例外,它们全都进入了共同的遗产,被普遍当作“已知”,在数学各处被日常使用。

但我的独特天赋把我带向的,比发现新问题、新概念、新陈述更进一步的,是发现富有成果的视角——它不断引导我引入并或多或少地发展一些全新的主题。我觉得,这就是我对我那个时代的数学所贡献的最本质的东西。说实话,我刚才说的那无数问题、概念、陈述,在我看来只有在这样一种“视角”的光照下才有意义——更确切地说,它们是从中自发地、带着显然性之力生出来的;就像黑夜里亮起的一道光(哪怕是散漫的光),仿佛从虚无中生出它突然向我们揭示的那些或模糊或清晰的轮廓。没有这道把它们束成一束的光,那十个、一百个、一千个问题、概念、陈述,就会显得像一堆杂七杂八、没有形状的“心智小玩意儿”,彼此孤立——而不是一个整体的各个部分;那整体也许仍不可见,仍躲在夜的褶皱里,却已被清楚地预感到。

富有成果的视角,就是那样一种视角:它把那些无人感到的灼热问题,揭示为同一个包容它们、赋予它们意义的整体的活的各部分;并且(也许是对这些问题的回答)把那些如此自然、却无人想到去提取的概念,以及那些仿佛顺理成章、却在激发它们的问题和使之得以表述的概念尚未出现之前谁也不敢提出的陈述,一并这样揭示出来。在我们这门技艺里\fnote{16}{}{毫无疑问,情况不只在我们这门“技艺”里如此,而且(在我看来)在一切发现的工作里都是如此,至少当这工作处于智识认识的层面时。},比所谓“关键定理”更有力的发现工具,是富有成果的视角——更确切地说,它们不是工具,而是研究者本人的眼睛,那双热切地想要认识数学事物之本性的眼睛。

这样,富有成果的视角不外是那只“眼睛”:它同时让我们发现,又让我们在被发现者的多样性中认出统一。而这统一,真正是把这些多样之物联系起来并赋予生命的生活本身与气息。

但正如其名所示,一个“视角”本身仍是局部的。它向我们揭示一处风景或全景的一个侧面,而此外还有许多同样有效、同样“真实”的侧面。只有当对同一现实的互补视角彼此结合起来,当我们的“眼睛”增多,目光才更深入地进入对事物的认识。我们想认识的现实越丰富、越复杂,拥有好几只“眼睛”\fnote{17}{}{每个视角都会带来发展出一套表达它、并专属于它的语言。拥有好几只“眼睛”或好几个“视角”来把握一个情境,也等于(至少在数学里)拥有好几种不同的语言来圈定它。}以便在它的全部广度与全部精微中把握它,也就越重要。

而有时候,一束汇聚在同一片广阔风景上的视角,凭着
% ---------------- p.42 ----------------
在我们身上那种能在多中把握一的能力,使一件新事物有了身体;这件新事物超出每一个局部透视,正如一个活物超出它的各个肢体和器官。这件新事物,可以叫做一种眼光。眼光把体现它的那些已知视角统一起来,并向我们揭示另一些此前被忽略的视角,正如富有成果的视角让人发现并把握众多新问题、新概念、新陈述,把它们当作同一整体的部分。\origpage{P16}

换个说法:眼光对于它仿佛从中生出、并由它统一起来的那些视角,就像白昼明亮温暖的光对于太阳光谱的不同成分。一种广阔而深沉的眼光,像一处取之不竭的源泉,注定要去启发和照亮的工作,不仅是在某一天生出它、并使自己成为它的仆人的那个人的工作,也是世世代代人的工作——他们也许(像他本人一样)被它让我们隐约看见的那些遥远边界所迷住……\origpage{P17}

\section{(7) “大想法”——或树木与森林}

我数学活动所谓“多产”的时期,也就是有正式发表物作证的时期,在1950到1969年之间,共二十年。而在二十五年里,从1945年(我十七岁)到1969年(我将满四十二岁),我几乎把全部精力投进了数学研究。这投入无疑过分。我为此付出的是长期的精神停滞、逐步的“变厚”,在《收获与播种》的篇幅里我不止一次有机会谈到。然而,在一项纯智力活动的有限领域内,凭一种局限于数学事物世界的眼光的绽放与成熟,那些年仍是创造力旺盛的年份。

在我生命的这段漫长时期,几乎全部时间和精力都用在所谓“部件活”上:那种细致的打造、装配和磨合工作,是从头建造一座座房子所必需的,而一个内心的声音(或一个魔鬼……)命令我建造它们,并按工作进展随时向我口授一位总施工的方案。我被“手艺”的活计占住:时而当石匠,时而当砌工、木匠,甚至当铅管工、细木工和家具匠——我很少偷闲把那份除我之外无人可见的总图(后来才显出它的样子……)用白纸黑字记下来,那总图在日、月、年之中以梦游者般的稳当引导着我的手\fnote{18}{}{“梦游者”这个意象,是柯斯特勒(Koestler)那本出色的书《梦游者》的标题给我的(Calmann-Lévy 版),那本书是一篇“关于宇宙观念史的研究”,从科学思想的开端直到牛顿。这段历史中令柯斯特勒深受触动、并由他特别指出的一点是:我们的世界认识常常从某一点走到另一点,那另一点(从逻辑上、事后看来)显得近在咫尺,而道路却要经过最离奇的迂回,仿佛在挑战健全的理性;然而,尽管有这千百道似乎注定要让人永远迷路的迂回,出发去寻找宇宙之“钥匙”的人们,却带着“梦游者般的稳当”,不由自主地、往往甚至连自己也没有察觉,碰上了另一些他们远未预料到的“钥匙”,而那些竟然正是“对的”。\par 据我在周围所能观察到的,在数学发现的层面上,这些发现之路上的惊人迂回,是某些大格局研究者的事,但绝不是所有研究者的事。这也许是因为,两三个世纪以来,自然科学的研究,尤其是数学研究,已从与某一文化、某一时代相联系的宗教或形而上学强制前提中解放出来——那些前提曾是对宇宙的“科学”理解之展开(无论好坏)特别强大的制动。然而确实,数学中某些最根本、最显然的观念与概念(如位移、群、零这个数、文字运算、空间中一点的坐标、集合的概念,或拓扑“形式”的概念,更不必说负数与复数),都经过几千年才出现。这些都是那种根深蒂固、深深植入心灵之中的“阻塞”的明白迹象——它抵制全新的观念,即便这些观念天真地简单,似乎以显然之力自行强加于人,而且几代、乃至几千年都是如此……\par{}回到我自己的工作,我觉得其中的“失手”(也许比多数同行更多)只限于细节,通常很快就由我自己发现。那不过是纯属“局部”的“途中意外”,并不严重影响关于所考察情境的基本直觉的可靠性。相反,在观念与重大指导性直觉的层面上,我的作品似乎没有任何“失误”,尽管这听起来难以置信。每一刻,即便不能把握一条进路的最终归宿(它多半仍藏在目光之外),也至少能把握那些最富成果、引我直达本质事物的方向——正是这种从未失灵的可靠性,使柯斯特勒的“梦游者”形象重新浮上我的心头。}。应该说,我乐于在“部件活”里倾注一份充满爱意的细致,这活并不令我讨厌。此外,前辈们所宣讲和实行的数学表达方式,(至少可以说)给了工作的技术面以优先,而不大鼓励那些停留在“动机”上的“离题话”;甚至对那些仿佛要从雾里唤出某个也许富有启发性的形象或眼光的离题话也不鼓励——因为它们还没有化身为木头的、石头的或纯硬水泥的、可触摸的建筑,更像是梦的碎片,而不是工匠专注而尽责的活计。\origpage{P18}

从量上说,我在这些高强度多产的年份里的工作,具体表现为大约一万二千页的发表物——论文、专著或讨论班讲义\fnote{19}{}{从1960年代起,其中一部分发表物是与同行(尤其是 J. 迪厄多内)和学生合作写成的。}——以及数以百计(如果不是数以千计)的新概念,它们连同我提取它们时给的名字一起,进入了共同的遗产\fnote{20}{}{这些概念中最重要的那些,《主题概述》以及随附的《历史评注》都作了回顾,两者将收入《反思》第4卷。有些名字是朋友或学生向我建议的,比如“光滑态射”这个说法(J. 迪厄多内),或“site、champ、gerbe、lien”这一整套,那是让·吉罗(Jean Giraud)的博士论文里发展出来的。}。在数学史上,我想我大概是在我们这门科学里引入新概念最多的人,同时也正因如此,不得不发明最多新名字的人——为了把这些概念表达得细致,并尽我所能地富有提示性。\origpage{P19}

这些纯“量”的指标,当然只给出对我作品的一种极其粗略的把握,错过了真正构成其灵魂、生命与力量的东西。正如我刚才写的,我在数学里贡献的最好的东西,是我先得以窥见、随后耐心提取并或多或少加以发展的新“视角”。像我刚谈到的那些概念一样,这些新视角进入极其多样的情境,它们本身也几乎数不清。

然而有些视角比另一些更广阔,单凭它们就引发并包揽大量局部视角,涉及大量各不相同的具体情境。这样一个视角,也可以恰当地称为一个“大想法”。凭它自身的丰产性,这样一个想法生出成群的后代,这些想法都继承它的丰产性,但大多数(如果不是全部)比母想法范围更小。

至于表达一个大想法,也就是“说出”它,这往往是几乎和它被构想出来、和在构想者心中缓慢孕育一样精细的事——或者更确切地说,这费力的孕育与成形工作,恰恰就是“表达”该想法的那项工作:耐心地、日复一日地把它从诞生时环绕它的雾幕中提取出来,逐渐给它一个可触的形态,成为一幅在星期、月份和年份中不断丰富、稳固、精细起来的图画\origpage{P20}。仅仅给这想法命名,用某个醒目的说法,或用多少带点技术性的关键词,也许只要几行、甚至几页——但很少有人在尚未相当熟悉它的情况下,能听懂这个“名字”,并在其中认出一张面孔。而当这想法到了完全成熟的时候,也许一百页就足以表达它,让孕育它的那位工匠完全满意——也可能一万页经过长期打磨和掂量的文字仍不够\fnote{21}{}{1970年离开数学舞台时,我关于概形这一中心主题的全部发表物(其中不少是合作),加起来该有一万来页。然而这在我眼前看到的那个关于概形的宏大纲领里,只占不大的一部分。这纲领在我离开后便被无限期搁置了,尽管此前已经发展并发表、可供所有人使用的东西,除极少数例外,一开头就进入了共同的遗产,成为大家日常使用、视为“人所共知”的概念与结果。\par 我离开时已完成的那部分纲领,关于概形主题及其延伸与分枝,单是它本身就代表了数学史上所完成的最庞大的基础工作,而且肯定也是科学史上最庞大的之一。}。

而无论在哪种情形下,在那些为了把想法据为己有而了解了这项终于呈现出想法的蓬勃生长的工作的人当中——它像一片在林间荒地上长起来的宽阔乔木林——很有理由打赌:许多人会清楚看见所有这些挺拔有力的树,并且各有所用(有人为了爬上去,有人为了取梁木和木板,还有人为了在自己壁炉里烧旺火……),但很少有人懂得看见森林……

\section{(8) 眼光——或十二个主题的和谐}

也许可以说,“大想法”就是那样一种视角:它不仅显得新而丰产,而且把一个新的、广阔的主题引入科学,而这主题正是它的化身。而任何一门科学,当我们不把它当作权力与支配的工具,而当作我们这个物种跨越时代的认识冒险来听时,它不外就是这种和谐——在不同时代或广或狭、或贫或富——它在世代与世纪中展开,靠的是所有先后出现的主题之间那种精细的对位,那些主题像被从虚无中唤来,加入其中并在其中交织。\origpage{P21}

在我于数学中提取出的众多新视角里,回头看去,有十二个我会称之为“大想法”\fnote{22}{}{以下是这份十二个主导想法、即我作品的“主主题”的清单(按出现的时间顺序),供好奇的数学读者一阅。\par 1. 拓扑张量积与核空间。\par 2. “连续”与“离散”的对偶(导出范畴、“六种运算”)。\par 3. 黎曼–罗赫–格罗滕迪克瑜伽($K$-理论,与相交理论的关系)。\par 4. 概形。\par 5. Topos。\par 6. 平展上同调与 $\ell$-进上同调。\par 7. Motif 与 Motivic 伽罗瓦群(Grothendieck 的 $\otimes$-范畴)。\par 8. 晶体与晶体上同调,“德拉姆系数”“霍奇系数”的瑜伽……\par 9. “拓扑代数”:$\infty$-栈、derivator;topos 的上同调形式主义,作为新的同伦代数的灵感。\par 10. 温和拓扑学。\par 11. 远阿贝尔代数几何的瑜伽,伽罗瓦–泰希米勒理论。\par 12. 正多面体与各种正则构型的“概形”或“算术”视角。\par 除这些主题中的第一个(其中重要的一支属于我1953年的博士论文,并在1950至1955年间我的泛函分析时期得到发展)之外,其余十一个都是在我1955年以后的几何学家时期提取出来的。}。看我的数学家作品,“感受”它,至少是稍稍看见并“感受”其中某些想法,以及它们引入的那些构成作品的经纬与灵魂的大主题。

% ---------------- p.45 ----------------
事所必然,这些想法中有些比另一些“更大”(而那些也就因此“更小”!)。换句话说,在这些新主题里,有些比另一些更广,有些更深地潜入数学事物之谜的核心\fnote{23}{}{在这些主题中,就意义之广而言,我觉得最广的是 topos,它提供了把代数几何、拓扑学与算术综合起来的想法。就已产生的发展之规模而言,最广的是概形主题。(关于这一点,见第20页页底注(*)。)正是它为所论主题中的另外八个(即除第1、5、10之外的所有主题)提供了“典范的”框架,同时为代数几何以及代数—几何语言的自下而上的更新提供了中心概念。\par 在另一端,十二个主题中的第一个和最后一个,在我看来比其余的规模更有限。不过,就最后一个而言,它把一种新视野引入正多面体与正则构型这个极古老的主题,我怀疑一个全身心投入其中的数学家,一生也不足以把它穷尽。至于所有这些主题中的第一个,即拓扑张量积,它扮演的更多是一件现成新工具的角色,而不是后人发展的灵感来源。不过,直到最近这些年,我仍偶尔收到多少较近的工作的零星回响,那些工作(二三十年后)解决了我留下的某些问题。\par 在我看来,这十二个主题中最深的是 Motif,以及与之紧密相连的远阿贝尔代数几何和伽罗瓦–泰希米勒瑜伽。\par 从工具经过我亲手调校、臻于完善、并在近二十年研究的各个“前沿地带”日常使用的力量看,我觉得最值得注目的是“概形”与“平展上同调、$\ell$-进上同调”这两块。对一个消息灵通的数学家来说,我想从现在起几乎不可能怀疑:概形工具,以及由它产生的 $\ell$-进上同调工具,属于本世纪少数几项重大成就之列,它们在最近几代人之中滋养并更新了我们的科学。}。有三个(在我看来并非最不重要的)只是在我离开数学界之后才出现的,至今仍处胚胎状态;“正式地说”,它们甚至并不存在,因为没有任何正式发表物可以充当它们的出生证明\fnote{24}{}{这三者得到哪怕最简略勾勒的唯一“半正式”文本,是1984年1月为申请 CNRS 脱产而写的《纲领概要》。这一文本(导论第3节“罗盘与行囊”里也谈到它)原则上将收入《反思》第4卷。}。在我离开之前出现的九个主题里,最后三个我留在蓬勃生长之中,今天仍处在童年状态,因为(我走后)没有充满爱的手去照料这些“孤儿”的所需,它们在一个敌对的世界里被弃之不顾\fnote{25}{}{在我不久离开之后,这三个孤儿便被无声无息地埋葬了;其中两个后来又被大吹大擂地掘出,却不提工匠的名字——一个在1981年,另一个(鉴于此举毫无瑕疵的成功)就在次年。}。至于另外六个主题,在我离开前的二十年里已完全成熟,可以说(有一两点保留\fnote{26}{}{这“除了极少例外”,主要指格罗滕迪克式的对偶瑜伽(导出范畴与六种运算),以及 topos 的瑜伽。这一点将在《收获与播种》第二部和第四部(《葬礼(一)》与《葬礼(三)》)里详加讨论。})它们那时已经进入共同的遗产:尤其在几何人中间,今天“人人”都在整天、随时吟诵它们,甚至不再知道自己是在吟诵(就像茹尔丹先生说着散文)。它们是我们呼吸的空气的一部分,当我们“做几何”,或当我们做多少带点“几何性”的算术、代数或分析时。\origpage{P22}

我作品这十二个大主题绝不是彼此孤立的。在我眼里,它们属于一种精神与话题的统一,这统一像一个共同而持久的低音,贯穿我全部“写出的”和“未写出的”作品。写这些行的时候,我仿佛又听见了同一个低音——像一个召唤!——贯穿那三年“无偿的”、顽强的、孤独的工作,那时候我还不曾操心世上除我之外是否还有别的数学家,因为我当时完全被那召唤我之物的魅力攫住了……\origpage{P23}

这种统一,并不只是同一个工匠在他手下出来的作品上留下的印记。这些主题由无数既精细又显然的纽带彼此相连,就像那些各各清晰可辨的不同主题之间的相连:它们
% ---------------- p.46 ----------------
在同一片广阔的对位中展开、交缠——在一个把它们聚合、推向前去、并赋予每一个以意义的和谐之中,每一个的意义、运动与圆满都有所有其他者参与。每一个局部主题似乎都从这更广的和谐中生出,又随时从中重生,远不是这和谐显得像那些先于它存在的构成性主题的“总和”或“结果”。说实话,我无法摆脱这种感觉(大概很可笑……):在某种意义上,正是这尚未出现、却无疑早已确确实实“存在”于尚未诞生之物的幽暗怀抱中某处的和谐——正是它先后唤起了这些只有通过它才获得全部意义的主题,也正是它在我刚出少年、那些炽烈孤独的年月里,已经用低而急迫的声音呼唤我……

无论如何,我作品的这十二个主主题确实都像出于一种秘密的预定,共同汇成同一部交响曲——或者换一个比喻,它们体现为同样多的不同“视角”,全都汇入同一个广阔的眼光。

这眼光只是到1957、58年前后才开始从雾中浮现、显出可辨认的轮廓——那是紧张孕育的年份\fnote{27}{}{1957年正是我被引向提取“黎曼–罗赫”主题(格罗滕迪克版)的那一年——这个主题一夜之间把我捧成了“大明星”。这一年也是我母亲去世之年,因此也是我生命中一道重要分界线之年。那是我一生中创造力最旺盛的年份之一,而且不限于数学层面。那时,全部精力投入数学工作已有十二年。就在那一年,一种感觉浮现出来:关于什么是数学工作,我大致已经“转完了一圈”,也许现在是时候投入别的事情了。那显然是一种内在更新的需要,一生中第一次浮出水面。那时我想过要当作家,并且有好几个月停止了全部数学活动。最后我决定,至少还要把手上已有的数学工作写成白纸黑字,也就是几个月的事,顶多一年……\par 无疑,那一跃的时机还没有成熟。总之,一旦重新拾起数学工作,便是它重新攫住了我。此后又十二年,它再没有放开我!\par 这段间歇之后的一年(1958),也许是我数学生涯中最丰产的一年。新几何的两个中心主题正是在这一年绽放:概形理论的强力启动(这是我当年夏天在爱丁堡国际数学家大会上报告的主题),以及“site”概念的出现——它是关键的 topos 概念的技术性临时版本。隔了将近三十年回头看,我现在可以说:新几何的眼光真正诞生,就在这一年;它紧随这门几何的两件主要工具之后:概形(它代表旧“代数簇”概念的一次变形),以及 topos(它代表空间概念的一次更深层的变形)。}。也许有点奇怪:这眼光对我如此切近,如此“显然”,以至于直到一年前\fnote{28}{}{我头一次想到给这眼光一个名字,是在1984年12月4日的反思里,在《阴,仆役(二)——或慷慨》一注(ReS III,第637页)的第1361号子注中。},我都没想过给它一个名字。(而我的一大癖好,恰恰是不停地给向我显现的事物命名,作为把握它们的第一步……)确实,我说不出某一个特定时刻,可以被经历为这眼光出现的时刻,或事后可以认作那样的时刻。一种新眼光是那样广阔,它的出现无疑不能定在某一个特定时刻,而必须在漫长的年月里(如果不是在世代之中)逐步渗入并占有那个或那些细看与凝视它的人;仿佛在那些注定要被逐步取代的熟悉的眼睛后面,必须辛苦地长出新眼睛来。而这眼光也太大,谈不上像把握路上拐角处出现的第一个概念那样去“把握”它。所以,最终不必惊讶:给这样广阔、又如此切近而弥散的东西命名的念头,只是事后、只是当这东西达到完全成熟之后才出现。\origpage{P24}

说实话,直到两年前,我与数学的关系(除了教书的任务之外)仅限于做数学——跟从一股不断把我向前拉的冲动,进入一个不断吸引我的“未知”。我从未想过在这股冲劲中停下来,哪怕只停一瞬\origpage{P25}
% ---------------- p.47 ----------------
,回头看看也许已画出的走过之路,甚至为一部已过去的作品定位。(无论是把它定位在我的生命里,当作一件仍有深的、长期未被察觉的纽带把我与它相连的东西;还是把它定位在“数学”这场集体冒险里。)

还有一件怪事:要让我终于“放下”并重新认识这部半被遗忘的作品,或者仅仅想到给那曾是它灵魂的眼光一个名字,竟需要我突然面对一场规模巨大的“葬礼”的现实:那场以沉默和嘲弄进行的葬礼,埋葬的是那眼光,以及生出那眼光的那位工匠……

\section{(9) 形式与结构——或事物的道路}

不知不觉,这篇“前言”一针引一线,最后成了对我作品的一次正式介绍,主要是给非数学读者看的。已经陷得太深,退不回去了,我只剩下把“介绍”讲完!我想尽力至少说几句,谈谈那些我在前面篇幅里晃动过的奇妙“大想法”(或“主主题”)的实质,以及那道著名的“眼光”的本性——那些主导想法据说要汇入其中。由于无法求助于多少带点技术性的语言,我大概只能传递一个极其模糊的形象(如果真有什么东西能“传过去”的话……)\fnote{29}{}{这个形象不得不“模糊”,丝毫也不妨碍它是忠实的,也不妨碍它确实再现了所看之物的本质的某种东西(在此即我的作品)。反过来,一个形象无论多么清晰,也可能完全走样,而且只包含次要的东西,完全错过本质。所以,如果你“接住”了我关于自己作品所要说的东西(那肯定我心中形象的一部分会真的“传过去”),你就可以得意地说,你对我的作品中最本质的东西,把握得比任何一位博学的同行也许都更好!}。

传统上,人们区分宇宙之物的三类“性质”或“方面”作为数学思考的对象:数\fnote{30}{}{这里指的是所谓“自然数”0、1、2、3等等,或者(至多)指那些借它们、通过初等性质的操作表达出来的数(如分数)。这些数不像“实数”那样,可以用来度量一个可连续变化的量,比如直线上、平面上或空间中两个变动点之间的距离。}、量、形。也可以把它们叫做事物的“算术”方面、“度量”(或“分析”)方面、“几何”方面。在数学所研究的大多数情境里,这三个方面同时在场并紧密互动。然而,最常见的是其中一个明显占主导。在我看来,对大多数数学家来说(对认识他们或了解他们作品的人来说),他们的基本气质相当清楚:是“算术家”、“分析家”还是“几何家”——即便他们本领很多,在各种音区和调门里都工作过。\origpage{P26}

我最初那些孤独的思考,关于测度与积分理论,毫无歧义地属于“量”或“分析”一类。我引入数学的第一个新主题也是如此(它在我看来比另外十一个范围要小)。我之从“分析”这个“侧门”进入数学,在我看来并非由于我的特殊气质,而是由于可以叫做“偶然境遇”的东西:对于我这颗爱好普遍性与严格性的头脑来说,中学和大学向我提供的教学里最大的空缺,恰好在事物的“度量”或“分析”方面。

1955年是我数学工作的一个关键转折:从“分析”
% ---------------- p.48 ----------------
转向“几何”。我还记得那种震撼的印象(当然完全是主观的):仿佛我离开了干瘦而硌人的草原,忽然来到某种“应许之地”,那里丰饶的财富无穷无尽地增殖,凡是手愿意落下去采摘或挖掘的地方都是如此……而这种压倒性的丰饶印象,超出了任何尺度\fnote{31}{}{我用“压倒性的,超出任何尺度”这个词语组合,勉强译出德文的“überwältigend”及其英文对应词“overwhelming”。上一句里那个(并不恰当的)“震撼的印象”,也应当带着这层意味来理解:当面对一种非凡的壮丽、伟大或美时,我们心中激起的印象与情感突然把我们淹没,以致任何表达我们所感的意愿都仿佛预先被摧毁了。},在此后这些年里只得到确认和加深,直到今天。\origpage{P27}

这就是说,如果在数学里有一件东西(大概从来如此)比任何别的更令我着迷,那不是“数”,也不是“量”,而始终是形。而在形为了向我们显现而选取的千零一张面孔中,比任何别的更令我着迷、并继续令我着迷的,是数学事物中隐藏的结构。

一件事物的结构绝不是我们可以“发明”的东西。我们只能耐心地把它揭示出来,谦卑地去认识它,“发现”它。这项工作里若有创造性,我们若有时做着铁匠或不倦的建造者的活,那绝不是为了“塑造”或“建造”什么“结构”。那些结构根本不曾等我们才存在,才恰好是它们所是的样子!而是为了尽可能忠实地表达我们正在发现和探测的这些事物,以及这个不肯轻易交付的结构——我们摸索着,用也许还结结巴巴的语言,试着把它圈定。于是我们不断地“发明”适合的语言,以越来越精细地表达数学事物的内在结构,并借助这语言,“建造”那些理应说明已被把握和看见之物的“理论”,一点一点、从头到尾。这里有一个持续不断的往返运动:在把握事物与表达被把握者之间,那语言在工作中、在当下需要的不断压力下变得精细并重新被创造。

读者大概已经猜到,这些“从头建造”的“理论”,也正是前面谈到的那些“漂亮的房子”:我们从先人那里继承的房子,以及我们在事物的召唤与倾听中亲手建造的房子。而我刚才说到建造者或铁匠的“创造性”(或想象力),还得补充一句:构成其灵魂与秘密神经的,绝不是那个说“我要这个,不要那个!”、并乐于凭自己心意决断的人的傲气;那就像一个糟糕的建筑师,还没看见、感受过一块地,还没探明它的可能与要求,脑子里就已经准备好了图纸。研究者的创造性与想象力的品质,在于他注意力的品质,在于他倾听事物之声。因为宇宙之物从不厌倦谈论自己、向那个有心去听的人显现。而最美的房子,显出工匠之爱的那一座,并不是比别人更大或更高的那座。美的房子,是忠实映照事物隐藏的结构与美的那座。\origpage{P28}

\section{(10) 新几何——或数与量的联姻}

可我又跑题了——我本打算谈主主题,它们在同一位母眼光中结合起来,像许多条河流回归它们为其子嗣的大海……

% ---------------- p.49 ----------------
这广阔的统一性眼光可以描述为一种新几何。据说,这就是上个世纪克罗内克所梦想的那一种\fnote{32}{}{关于这个“克罗内克的梦”,我只是听人说起:有人(也许正是约翰·泰特)对我说,我正在实现那个梦。在我从前辈那里接受的教育里,历史性的参考极少见;哺育我的并不是阅读多少古代的、甚至同时代的作者,而主要是与其他数学家——首先是我的前辈们——的口头或通信交流。1958年概形理论突然而强劲的启动,其主要的、也许是唯一的外来灵感,是塞尔那篇以 FAC(《代数凝聚层》)这个缩写广为人知的文章,它发表于此前几年。除此之外,理论后来发展中的主要灵感都是从自身产生的,并在年复一年中自我更新,靠的只是内部简洁性与一致性的要求,以及一种努力:在新语境中说明代数几何里“人所共知”的东西(那些东西在我手中变形的同时被我逐步消化),以及这个“已知”让我预感到的东西。}。但现实(一个大胆的梦有时让人预感到或窥见它,并鼓励我们去发现它……)每次都在丰富与共鸣上超过最大胆、最深的梦。可以肯定,对这新几何的不止一个部分(如果不是全部),在它出现的前一天,没有人会想到——工匠本人也不比别人更早想到。

可以说,“数”适合把握“不连续”或“离散”的聚集体的结构:那些往往有限的系统,由可以说是彼此孤立的“元素”或“对象”组成,没有某种从一个到另一个的“连续过渡”原则。“量”则相反,是典型地能够“连续变化”的性质;因此它适合把握连续的结构与现象:运动、空间、各种“簇”、力场等等。于是,算术显得(大致上)是离散结构的科学,分析则是连续结构的科学。

至于几何,可以说,它以现代意义上的科学的形式存在了两千多年,一直“骑”在这两类结构上,即“离散的”与“连续的”\fnote{33}{}{说实话,传统上处在几何家注意力中心的是“连续”方面,而“离散”性质,尤其是数值的和组合的性质,则被略而不谈,或草草打发。十年前我惊叹地发现二十面体组合理论的丰富,而这个主题在克莱因论二十面体那本经典著作里连碰都没有碰到(大概根本没看到)。我看到几何家对自发进入几何的离散结构这种(长达两千年的)忽视的另一个显著迹象:群(尤其对称群)的概念直到上世纪才出现,而且它最初还是(由埃瓦里斯特·伽罗瓦)在一个当时并不被视为属于“几何”的语境中引入的。确实,直到今天,仍有许多代数学家始终不明白:伽罗瓦理论本质上正是一种“几何”眼光,它更新了我们对所谓“算术”现象的理解……}。而且长期以来,两种本属不同种类的几何——一种离散、一种连续——之间并没有真正“离婚”。不如说,那是对同样的几何图形作研究时的两种不同视角:一个强调“离散”性质(特别是数值的和组合的性质),另一个强调“连续”性质(比如在周围空间中的位置,或以点的相互距离来度量的“量”等等)。\origpage{P29}

直到上世纪末,才出现了一场离婚,随着人们有时称为“抽象(代数)几何”的东西的出现与发展。大致说来,它在于为每个素数 $p$ 引入一种“特征 $p$ 的(代数)几何”,照着前几个世纪传下来的(代数)几何的(连续)模型来仿制,然而所处的语境却显得不可还原地“不连续”“离散”。这些新的几何对象自本世纪初以来地位日益重要,尤其是由于它们与算术——离散结构的典型科学——的紧密关系。这似乎是安德烈·韦伊(André Weil)\fnote{34}{}{安德烈·韦伊,旅居美国的法国数学家,是“布尔巴基小组”的“创始成员”之一;《收获与播种》第一部里会不少谈到这个小组(偶尔也会谈到韦伊本人)。}作品中的一条指导性想法,也许还是主要的力量线(在他的书面作品中,它或多或少停留在默认状态,这也是应该的):即“代数”几何,尤其是与不同素数相联系的“离散”几何,应当提供钥匙,让算术得到大范围更新。他就是本着这种精神,在1949年提取出著名的“韦伊猜想”的。

% ---------------- p.50 ----------------
说实话,这些猜想令人目瞪口呆:它们让人们对这些离散本性的新“簇”(或“空间”)隐约看见某些类型的构造与论证\fnote{35}{}{(供数学读者参考。)这里指的是与微分簇或复簇的上同调理论相联系的“构造与论证”,尤其是涉及莱夫谢茨不动点公式与霍奇理论的那些。}的可能性——而在此前,这些似乎只在分析家认为配得上这个名号的“空间”,即所谓“拓扑”空间(连续变化的概念在其中通行)的框架里才可设想。\origpage{P30}

可以认为,新几何首先是这两个世界之间的综合——它们此前相邻、紧密相依,却又彼此分隔:“算术”世界,其中住着那些没有连续性原则的(所谓)“空间”;以及连续量的世界,其中住着本来意义上的“空间”,它们对分析家的手段敞开,(正因如此)也被他接受为配得上在数学之城里有住所。在新眼光里,这两个昔日分隔的世界只成了一个。

这种“算术几何”(我提议这样称呼这种新几何)眼光的第一个胚胎,就在韦伊猜想里。在我某些主要主题\fnote{36}{}{指四个“中间”主题(第5至8号),即 topos、平展上同调与 $\ell$-进上同调、Motif,以及(较小程度上)晶体。这些主题是我在1958到1966年间先后提取出来的。}的发展中,从1958到1969这些年,这些猜想始终是我主要的灵感来源。而且,在我之前,一边是奥斯卡·扎里斯基(Oscar Zariski),另一边是让–皮埃尔·塞尔,已经为“抽象”代数几何那些无法无天的空间发展出某些“拓扑”方法,其灵感来自此前用于大家那些“正经”空间的方法\fnote{37}{}{(供数学读者参考。)在我看来,扎里斯基在这方面的主要贡献是引入“扎里斯基拓扑”(它后来成了塞尔在 FAC 中的一件本质工具),以及他的“连通性原则”和他所谓“全纯函数论”——后者在他手中变成了形式概形理论,以及形式与代数之间的“比较定理”(第二个灵感来源是塞尔那篇奠基性的 GAGA)。至于我在正文中提到的塞尔的贡献,当然首先是他把层(由让·勒雷在十几年前、在一个完全不同的语境中引入)的视角引入抽象代数几何,即在那篇已经提到的另一篇奠基性文章 FAC(《代数凝聚层》)中。\par 借这些“回顾”之光,如果我要举出新几何眼光最近的“祖先”,立刻想到的就是奥斯卡·扎里斯基、安德烈·韦伊、让·勒雷和让–皮埃尔·塞尔这几个名字。其中塞尔起了独特的作用,因为主要正是通过他,我才不仅知道了他自己的想法,也知道了扎里斯基、韦伊和勒雷那些在新几何的绽放与发展中起过作用的想法的。}。

当然,他们的想法在我构筑算术几何的最初几步中起了重要作用;不过说实话,更多是作为出发点和工具(为了一个广阔得多的语境的需要,我不得不或多或少把它们从头重做),而不是作为一路继续滋养我的梦与计划的灵感来源。无论如何,一开始就很清楚:即便重做过,这些工具也远远达不到所需——连朝着那些奇妙猜想走出最初几步都不够。\origpage{P31}

\section{(11) 魔法扇——或天真}

在新几何的启动与发展中,两个关键的力量线是概形与 topos。它们几乎同时出现,紧密共生
% ---------------- p.51 ----------------
\fnote{38}{}{这次启动发生在1958年,第23页的页底注里谈到它。“site”或“格罗滕迪克拓扑”的概念(topos 概念的临时版本)紧随着概形概念出现。反过来,正是它提供了“局部化”或“下降”的新语言,在概形主题与工具的发展中每一步都要用到。更为内蕴、更具几何性的 topos 概念,起初在随后几年里仍是隐性的,主要从1963年起随着平展上同调的发展而浮现,并逐步在我这里确立为最根本的概念。},在新几何引人注目的兴起中像是同一根驱动神经,而且从它们出现的那一年起就如此。为了结束这趟对我作品的巡览,我至少还得就这两个想法说几句。

概形这个概念,是把此前人们摆弄的那无穷系列的(代数)“簇”概念(每个素数一个这样的概念\fnote{39}{}{这一系列里还应当包括 $p=\infty$ 的情形,它对应“特征零”的代数簇。}……)统摄进唯一概念的想得到的最自然、最“显然”的办法。此外,同一个“概形”(或新版“簇”)对每个素数 $p$ 都产生一个完全确定的“特征 $p$ 的(代数)簇”。这些不同特征的簇的集合,于是可以被看成某种“簇的(无穷)扇”(每个特征一把)。而“概形”就是这把魔法扇,它把自己的各种可能特征下的“化身”或“转世”彼此连起来,像许多不同的“扇骨”。由此,它提供了一个有效的“过渡原则”,把那些此前显得多少彼此孤立、互不相通的几何所属的“簇”彼此连起来。如今它们被统摄在一门共同的“几何”里,并由它相连。可以把它叫做概形几何,它是那种“算术几何”的最初草图,而后者将在随后几年里在其中充分展开。\origpage{P32}

概形这个想法本身天真地简单——如此简单、如此卑微,以致在我之前没人想到要弯得这么低。说实话,甚至“蠢”得可以:以致好些年里,尽管事情明摆着,对我许多博学的同行来说,这实在“不严肃”!而我也确实需要好几个月紧凑而孤独的工作,才能在自己的角落里说服自己:“它行得通”——我那股不可救药的幼稚使我执意要试验的新语言,那个蠢得可以的语言,确实适合以新的光亮和新的精细,并且从此在一个共同的框架里,把握那些与前此“特征 $p$ 的几何”相连的最早的几何直觉中的某些。这类练习,任何“消息灵通”的人都会预先判定为愚蠢而无望;而在我的所有同行和朋友中,大概只有我一个人会想到把它放进脑子,甚至(被一个秘密的魔鬼驱使……)不顾所有人反对把它做到底!

我并没有让自己被周围那些关于什么“严肃”、什么不严肃的立法共识分散注意力,而是像过去一样,单纯地信任事物卑微的声音,信任我身上那个懂得倾听的东西。报偿是即时的,而且超出一切期待。在这几个月里,我甚至没有“特意”,就摸到了强大而出人意料的工具。它们让我不仅能以更透彻的光重新找到(像玩一样)那些公认难懂的旧结果并超越它们,而且终于能够着手并解决一些“特征 $p$ 的几何”问题——此前用当时一切已知手段都显得无法企及\fnote{40}{}{概形理论这次“强力启动”的报道,就是我1958年在爱丁堡国际数学家大会上的报告。在我看来,这篇报告是概形视角最好的导论之一,(也许)能让一位几何读者有动力去勉强啃一啃那部(后来的)厚重论著《代数几何原理》——它详详细细地(一个技术细节都不放过)阐述了代数几何的新基础与新技巧。}。\origpage{P33}

% ---------------- p.52 ----------------
在我们对宇宙之物的认识里(无论数学的或别的),我们身上那革新的力量不外是天真。那是我们出生时人人分得的原初天真,它安息在每个人心里,常常成为我们蔑视的对象,也成为我们最隐秘的恐惧的对象。唯有它把谦卑与大胆结合起来,让我们深入事物之心,也让我们允许事物深入我们、浸透我们。

这种力量绝不是“天赋”异禀的特权——比如说有一种超常的脑力,能灵巧轻便地消化和摆弄大量已知的事实、想法和技术。这样的天赋当然可贵,对于像我这样出生时没有被“超乎一切尺度地”赋予它的人,当然值得羡慕。

然而,不是这些天赋,也不是最炽烈的野心配上无懈可击的意志,能让人越过那些把我们的宇宙关起来的“看不见而专横的圈子”。唯有天真能越过它们,不知不觉也不在意,在那些我们独自倾听事物、全神贯注于一场儿童游戏的时刻……

\section{(12) 拓扑学——或雾中的丈量}

我们刚看到,“概形”这个革新性的想法,让人能把与不同素数(或不同“特征”)相联系的不同“几何”彼此连起来。然而这些几何各自本质上仍是“离散”或“不连续”的,与过去几个世纪传下来的(可上溯到欧几里得的)传统几何相对照。扎里斯基和塞尔引入的新想法,在一定程度上为这些几何恢复了某种连续性的“维度”,它随即被刚出现、旨在把它们统一起来的“概形几何”继承。但就那些(韦伊的)“奇妙猜想”而言,还差得远。从这个角度看,这些“扎里斯基拓扑”粗疏到这种地步,几乎就像还停留在“离散聚集体”的阶段。显然缺少的是某种新原则,能把那些几何对象(或“簇”,或“概形”)与通常的、或“正经”的(拓扑)“空间”连起来;所谓正经的空间,姑且说,是那些“点”显得彼此清楚分开的空间,而在扎里斯基引入的那些无法无天的空间里,点却有一种讨厌的倾向,彼此黏成一团……\origpage{P34}

确实,只有这样一个“新原则”的出现,而且非它不可,才能让那场“数与量的联姻”,或者说“不连续几何”与“连续几何”的联姻得以完成——韦伊猜想里已经有了对它的第一次预感。

“空间”概念大概是数学里最古老的概念之一。它在我们对世界的“几何”把握中如此基本,以致两千多年来它或多或少一直是默认的。只是在上一个世纪里,这个概念才逐步摆脱直接知觉(对我们周围那唯一同一个“空间”的知觉)及其传统理论化(“欧几里得式”)的专制控制,获得自身的自主与动力。今天,它属于数学中最普遍、最常用的少数概念之列,无疑为无一例外的所有数学家所熟悉。它也是再千变万化不过的概念,有千百张面孔,取决于人们给这些空间加上什么类型的结构:从最丰富的(如那些可敬的“欧几里得”结构,或“仿射”与“射影”结构,或同名“簇”的“代数”结构,后者推广并柔化了前者),直到最贫乏的:那些任何“定量”信息要素似乎都一去不返地消失了
% ---------------- p.53 ----------------
,只剩“邻近”或“极限”概念之定性精髓\fnote{41}{}{说到“极限”概念,我在这里想的尤其是“取极限”那个概念,而不是(对非数学人士更熟悉的)“边界”那个概念。},以及形式直觉最难以捉摸的版本(所谓“拓扑的”)。在这些概念中,最贫乏的一个,也是过去半个世纪里一直充当某种广阔的公共概念怀抱、把其余一切统摄于其中的那个,就是拓扑空间。对这些空间的研究构成几何中最迷人、最有生命力的分支之一:拓扑学。\origpage{P35}

由一个(所谓“拓扑”)“空间”所体现的这种“纯品质”结构,乍看之下无论多么难以捉摸——因为一切定量性质的数据(特别是两点之间的距离)都缺席,使我们无法挂靠任何关于“大”“小”的熟悉直觉——上个世纪里人们毕竟还是把这些空间精细地圈在了一套仔细“量身剪裁”的语言的紧密而柔韧的网眼里。更妙的是,人们还从头发明并制造出各种“米尺”或“量绳”,无论如何也要给那些触手般伸展的“空间”挂上各种“度量”(叫做“拓扑不变量”)——那些空间像抓不住的雾,似乎躲开一切测量的尝试。确实,这些不变量中大多数、而且是最本质的那些,其本性比单纯的“数”或“量”更精细——它们本身就是多少有些精巧的数学结构,以多少复杂些的构造挂到所考虑的空间上。这些不变量中最古老、最关键之一,上个世纪已经(由意大利数学家贝蒂)引入,是由与该空间相联系的各个所谓“上同调”“群”(或“空间”)构成的\fnote{42}{}{说实话,贝蒂引入的不变量是同调不变量。上同调是它的一种大致等价、“对偶”的版本,出现得晚得多。这个方面压过了最初的“同调”方面,主要(无疑)是由于让·勒雷引入了层的视角,下面会谈到。从技术角度看,可以说我大部分几何学家的工作,都在于提取并或多或少地发展那些缺失的上同调理论,为各种空间与簇、尤其是为“代数簇”和概形而设。一路之上,我还被引向用上同调的语言重新解释传统的同调不变量,并由此让它们在一种全新的光下显现。\par 拓扑学家还引入了许多别的“拓扑不变量”,用来把握拓扑空间的这类或那类性质。除了空间的“维数”和(上)同调不变量,最早的其他不变量是“同伦群”。1957年我又引入了另一个,即所谓“格罗滕迪克”群 $K(X)$,它随即大行其道,其重要性(无论在拓扑学还是在算术中)至今不断得到确认。\par 一大批新不变量,本性比目前已知和使用的更精细,但我觉得是根本性的,都写进了我的“温和拓扑学”纲领(其极简的概述见《纲领概要》,将刊于《反思》第4卷)。这一纲领基于“温和理论”或“温和空间”的概念,它有点像 topos 那样,构成空间概念的(第二次)“变形”。它比后者显然得多(在我看来),也没有那么深。我预计它对“严格意义上的”拓扑学的直接影响会明显更有冲击力,并且会彻底改变几何拓扑学家的“手艺”,办法是深刻改造他所工作的概念语境。(正如代数几何里引入概形视角时所发生的那样。)此外,我把自己的《纲领概要》寄给了好几位旧友和著名的拓扑学家,但看来它并没有引起他们中任何一位的兴趣……}。\origpage{P36}

正是它们(确实,主要是“在字里行间”)出现在韦伊猜想里,构成其深刻的“存在理由”,并使它们(至少对我来说,被塞尔的解释“带进了门”)具有全部意义。但要把这样的不变量挂到这些猜想里出现的“抽象”代数簇上,以回应为那一事业的需要而提出的非常精确的要求——这还只是一个希望。我怀疑,除塞尔和我自己之外,没有别人(甚至,尤其是,安德烈·韦伊本人!\fnote{43}{}{说来矛盾,韦伊对上同调形式主义有一种顽固的、看来是发自内脏的“阻塞”——而这恰恰主要是他那些著名的猜想,从1955年起启发了代数几何中那些宏大上同调理论的发展(塞尔以其已在前面的页底注中提到过的奠基性文章 FAC 开了第一炮)。\par 在我看来,这种“阻塞”在韦伊那里属于一种普遍的厌恶:厌恶一切“大堆头”,厌恶一切近似于形式主义(当它无法用几页纸概括时)或多少有点复杂交错的“构造”的东西。他确实一点也不像“建造者”;显然,他在三十年代是被迫地动手发展代数几何最初的“抽象”基础,而那些基础(鉴于他这种性情)对使用者来说实在成了一张普罗克鲁斯忒斯之床。\par 我不知道他是否怪我得太过分,怪我投入去建造那些宽敞的住宅,使克罗内克的梦和他的梦得以化身为一套精细而有效的语言与工具。总之,他从未就他看见我从事的工作、或已经完成的工作,对我说过一句话的评论。我把《收获与播种》寄给他三个多月,还附上我亲手写的热情题献,也没有得到任何回音。})真的相信……\origpage{P37}

% ================================================================
%  第 2 章 巡游一部作品——或孩子与母亲 (后半, 原书 p.54-71, 边码 P38-P65)
%  源文: _work/src1/task-ch02b.md
%  说明: 原书注 49、注 60 的注文在源文中缺失, 正文保留裸上标
% ================================================================

% ---------------- p.54 ----------------
\origpage{P38} 还在此前不久,我们关于这些上同调不变量的构想,已因让·勒雷的工作而大大丰富、深刻更新——那工作是战争期间他在德国当俘虏时继续做的,时在四十年代前半。其决定性的创新想法,是空间上的(阿贝尔)层,勒雷给每个层配上一列相应的“上同调群”(所谓“以该层为系数的”上同调群)。这就好像那根一直用来“丈量”空间的、可靠的老式“上同调标尺”,忽然被增殖为多得难以想象的众多新“标尺”,各种尺寸、形状、质地都有,每一根都贴合于所论的那个空间,每一根都就它向我们交出精确无比的信息,而且是唯有它才能给出的信息。这正是我们对各类空间的研究方式发生深刻转变的主导想法,也肯定是本世纪出现过的最关键的想法之一。多亏让–皮埃尔·塞尔后来的工作,勒雷的想法很快结出最初的果实——在其出现后的头一个十年里,拓扑空间(尤其是与上同调紧密相关的所谓“同伦”不变量)理论有了一次令人瞩目的重新启动;而所谓“抽象”代数几何也有另一次同样根本的重新启动(靠塞尔那篇奠基性文章 FAC,1955年发表)。我自1955年起在几何方面的工作,正与塞尔这些工作相接续,因而也与勒雷那些创新想法相接续。

\section{(13) topos——或双人床}

勒雷引入的层的视角与语言,让我们以一种新的光照去看各类“空间”与“簇”。然而它们并未触及空间概念本身,只是让我们以新的眼睛更精细地把握这些人人早已熟悉的传统“空间”。可后来表明:这个空间概念,并不足以承载那些最本质的“拓扑不变量”——正是它们表达着“抽象”代数簇(例如韦伊猜想所适用的那些)的“形状”,乃至一般“概形”(旧日簇的推广)的形状。对于那场期待中的“数与量的联姻”来说,这就像一张分明过于狭窄的床:充其量只有未来夫妇中的一个(即新娘)能勉强挤进去安身,却从来容不下两个!而那尚待发现的“新原则”——好让仁慈的仙女们许下的婚约得以圆房——不是别的,也正是未来夫妇所缺的那张宽敞的“床”;而在此之前,竟没有人哪怕察觉到这一点。

这张“双人床”,随着 topos(广义空间)这个想法出现了(像被魔杖一点……)。这个想法在同一种拓扑直觉里,既容纳了体现着连续量世界的传统(拓扑)空间,也容纳了那些不知悔改的抽象代数几何学家所谓的“空间”(或“簇”),还容纳了无数其他类型的结构——它们此前似乎无可挽回地被铆在“不连续”或“离散”的聚集体的“算术世界”上。

正是层的视角做了那沉默而可靠的向导,那把有效(且绝非秘密)的钥匙,
% ---------------- p.55 ----------------
\origpage{P39} 毫不迟疑、毫不绕弯地把我领进洞房,领到那张宽大的婚床前。这床确实宽得(像一条宽广、宁静、深不见底的河……),宽到——
\begin{quote}
“国王的全部马匹\\
都能一齐来饮……”
\end{quote}
——正如一支老歌里唱的;那歌你肯定也唱过,至少听人唱过。而头一个唱这歌的人,对 topos 那种隐秘的美、那种宁静的力量,感受得比我那些博学的学生、比我旧日的任何朋友都要真切……

钥匙是同一把,无论在最初那套临时性的进路里(经由“site”这个十分方便、却并非内蕴的概念),还是在 topos 的进路里。现在我想试着描述一下 topos 这个想法。

试想一个给定(拓扑)空间上所有层所构成的“集合”,或者,愿意的话,由所有这些用来“丈量”它的“标尺”所构成的那套了不起的装备\fnote{44}{}{(给数学家看的)严格说来,这里说的是集合的层,而不是勒雷为了构造“上同调群”而引入、作为最一般系数的阿贝尔层。顺便说,我相信自己是头一个系统地使用集合的层的人(从1955年起,在堪萨斯大学写的那篇《具有结构层的纤维空间的一般理论》里)。}。我们把这一“集合”或这套“装备”看作带有它最显然的那种结构——如果可以说的话,那结构“一眼就看得见”;也就是所谓“范畴”结构。(非数学的读者不必因为不知道这个术语的技术含义而发慌。下文他完全用不着它。)正是这种可称为“(所论空间上的)层范畴”的“丈量上层建筑”,从此将被看作“体现”着空间最本质的东西。这在(“数学上的健全常识”看来)完全正当,因为事实是:人们可以凭一个拓扑空间所关联的这个“层范畴”(或这套“丈量装备”),把它整个地“重新构造”出来\fnote{45}{}{(给数学家看的)严格说来,这只对所谓“索伯(sobre)空间”成立。不过这类空间几乎包括了我们平常遇到的全部空间,尤其包括分析学家们所钟爱的一切“分离”空间。}。要核实这一点只是个简单的练习——当然,得先把问题提出来……有了这一条,就足以确信:只要我们出于某种理由觉得合适,从此便可以“忘掉”最初那个空间,只保留、只使用与之关联的那个“范畴”(或那套“装备”);它将被看作所要表达的“拓扑(或空间)结构”最恰当的化身。

正如数学中常有的事,我们在这里(靠“层”或“上同调标尺”这个关键想法)成功地把某个概念(这里是“空间”)用另一个概念(这里是“范畴”)表达了出来。每一次,当人们发现可以把一个概念(表达某一类情境)译成另一个概念(对应另一类情境)时,我们对这两个概念的理解都会同时加深——因为分别附着于这一个或那一个的特定直觉,会出人意料地汇流到一起。于是,一个“拓扑”性质的情境(由某个给定空间体现),在这里被译成了一个“代数”性质的情境(由某个“范畴”体现);或者,愿意的话也可以说:由空间体现的“连续”,被这个“代数”性质的范畴结构“译”出、“表达”出来了——甚至直到那“不连续”或“离散”性质的外包装。

但这里还有更多的东西。这些概念中的头一个,即空间概念,在我们看来多少是个“极大”的概念——它已经如此一般,以至很难想象还能找到一个仍然“合理”的扩张。然而,事实是:在镜子的另一面\fnote{46}{}{这里说的“镜子”,像《爱丽丝镜中奇遇》里那面一样,是把一个空间放在它前面、便给出该空间之“像”的那面镜子;像就是与之关联的“范畴”,被看作空间的一种“替身”,“在镜子的另一面”……},我们从拓扑空间出发所遇到的这些“范畴”(或“装备”),性质非常特别。它们确实享有一组特征极强的性质\fnote{47}{}{(给数学家看的)这里说的主要是我在范畴论中以“正合性”为名引入的那些性质(同时还引入了现代范畴论中一般的归纳极限与投射极限概念)。见《论同调代数中的若干要点》,《东北数学杂志》,1957年(第119–221页)。},使它们看上去像是其中最简单那一个的各式“仿制品”——那最简单的,就是从缩成单独一点的空间出发所得到的那个。话虽如此,一个推广了传统拓扑空间的“新型空间(或 topos)”,将被径直描述为一个“范畴”:它未必来自某个寻常空间,却仍然具备这样一个“层范畴”所应有的全部好性质(当然,这些性质是一次性地明文规定好的)。

% ---------------- p.56 ----------------
\origpage{P41} 这就是那个新想法。它的出现,可以看作下面这个其实近乎孩子气的观察的结果:在一个拓扑空间里,真正要紧的绝不是它的“点”或点的子集\fnote{48}{}{例如,人们可以造出非常“大”的 topos,它只有一个“点”,甚至根本没有“点”!},也不是点与点之间的邻近关系之类,而是这个空间上的层,以及它们所构成的范畴。总而言之,我不过是把勒雷最初的想法推到了它最终的结论——推完之后,再跨出那一步。

如同层这个想法本身(源自勒雷),如同概形这个想法,如同一切要来撞翻某种积习已深的看法的“大想法”,topos 这个想法也自有一种令人不知所措之处:它那么自然,那么“显然”,那么单纯(可以说,到了天真乃至简陋的极限,甚至到了“犯傻”的地步)——就是那种特别的性质,常让我们脱口而出:“哦,原来就这么点东西!”语气里一半失望,一半嫉妒;也许还外加一层“古怪”“不正经”的弦外之音,凡是因为出人意料的过度简单而叫人摸不着头脑的东西,往往就落得这评语。这也许是在提醒我们,想起童年那些久已埋没、久已否认的日子……

\section{(14) 空间概念的变异——或气息与信仰}

概形概念构成了“代数簇”概念的一次大扩张,因此它把我前人所传的代数几何从里到外更新了一遍。topos 概念则构成了空间概念的一次意想不到的外延,更确切地说,一次蜕变。由此,它承载着拓扑学、并在此之外还有几何学获得类似更新的许诺。而且,从现在起它已经在新几何的兴起中起了关键作用(主要通过由它而生的 $\ell$-进上同调与晶体上同调这两个主题,并通过它们,在韦伊猜想的证明中起作用)。像它的姐姐(而且几乎是孪生姐姐)一样,它具备任何一次丰产的推广所必备的两个互补特征,如下。

% ---------------- p.57 ----------------
\origpage{P42} 第一,新概念不太宽:也就是说,在新的“空间”里(人们更愿意叫它们“topos”,免得冒犯一些敏感的耳朵\fnote{49}{}{“topos”这个名字是(连同“topologie”或“topologique”之名一起)选定的,意在暗示:它正是拓扑直觉所适用的“那个对象”。凭这个名字唤起的丰富心象之云,应当把它看作大致相当于“(拓扑)空间”一词,只是更强调这个概念“拓扑”上的特殊性。(于是,有“向量空间”,但在另行通知之前,并没有“向量 topos”!)看来应当把两个说法一并保留,各存其特殊性。}),那些对旧日好空间来说司空见惯的、最本质的“几何”直觉与构造\fnote{50}{}{这些“构造”里,尤其包括一切熟悉的“拓扑不变量”的构造,上同调不变量也在其中。至于后者,我在已经提到的那篇文章(“东北”1955)里把该做的都做了,好让它们对任何“topos”都有意义。},大体上都能搬过来。换句话说,对这些新对象,我们手头有整套丰富的意象与联想、概念以及至少某些技术,而这些东西以前只限于老式对象。

第二,新概念同时又足够宽,能把许多情境包进来——这些情境在此之前并不被认为会产生“拓扑–几何”性质的直觉,而那正是过去只留给寻常拓扑空间的直觉(而且自有其道理……)。

从韦伊猜想的眼光看,这里的关键在于:新概念确实足够宽,使我们能给每个“概形”配上这样一个“广义空间”或“topos”(称为所论概形的“平展 topos”)。这个 topos 的某些“上同调不变量”(全是些“傻乎乎”的东西!)当时看来颇有希望提供“所需要的东西”,使这些猜想获得全部意义,而且(谁知道呢!)也许还能提供证明它们的办法。

正是在我此刻正在写下的这些纸页里,我有生以来第一次(作为数学家)有闲暇去述说(哪怕只是对自己述说)我数学作品中全部的主旋律与重大指导性想法。这使我更好地领会这些主题、以及它们所体现的种种“视角”,在统一它们、并生出它们的那幅宏大几何图景中各自的位置与分量。正是通过这项工作,新几何最初那强劲兴起中两个要害的创新想法才在充分的光照下显现出来:概形的想法,与 topos 的想法。

这两个想法里,如今在我看来更深的是第二个,即 topos 的想法。假如五十年代末我没有挽起袖子,日复一日固执地干下去,一干十二个长年,做出一件精细与力量都臻完美的“概形工具”——那么我几乎无法想象,在此后的一二十年里,别人最终会一直忍住不去引入那个明明非引入不可的概念(哪怕是被迫的),也不会好歹搭起几间“预制件”的破棚屋,聊胜于我这番有心一石一石拼合、亲手竖起的那一座座宽敞舒适的宅邸。然而反过来,在近三十年的数学舞台上,我看不出还有谁会有这份天真,或这份单纯,去(替我)跨出那另一步——那一步比什么都关键,它引入了 topos 这个如此孩子气的想法(哪怕只是“site”那个想法)。就算这个概念已经有人殷勤地奉上,连同它似乎暗藏的那点怯生生的许诺——我也看不出还有谁,无论旧友还是学生,会有那份气息,尤其是那份信仰,把这卑微的想法\fnote{51}{}{(给数学家读者)我说“把这卑微的想法做到底”,指的是把平展上同调当作通往韦伊猜想的一条进路这个想法。正是在这句话的启发下,我于1958年发现了 site 概念;而这个概念(或与之极相近的 topos 概念)以及平展上同调的形式体系,则是在我的推动下于1962至1966年间发展起来的(有几位合作者协助,他们将在适当的地方谈到)。\par{}我说“气息”和“信仰”,指的是“非技术”性质的品质,而在这里,它们在我看来正是本质的品质。在另一个层面上,我还可以加上我所谓的“上同调的嗅觉”,也就是当初在我身上发展起来的、用于建造上同调理论的那种嗅觉。我原以为把它传给了我的上同调学生们。离开数学界十七年之后回头一看,我发现他们当中没有一个人把它保留下来。}一直做到底(它表面上如此微不足道,而目标看来无限遥远……):从最初结结巴巴的起步,直到“平展上同调之掌握”完全成熟——在随后那些年里,它最终在我手中化成了这个。

% ---------------- p.58 ----------------
\section{(15) 国王的全部马匹}

\origpage{P43}\origpage{P44} 是的,河很深;我童年的水又宽又静,在那片我以为早已离开的国土上。国王的全部马匹都能在那里从容地、喝个痛快地一齐来饮,也饮不尽它!那水来自冰川,炽热如那些遥远的积雪,又有着平原黏土的柔和。我刚说到其中一匹马,一个孩子把它牵来饮水,它久久地喝了个够。我还看见另一匹来饮了一会儿,说不定是循着同一个孩子的足迹来的——可那一回没耽搁多久。一定是有人把它赶走了。也就这样了,可以说,再没有别的。然而我看见平原上流浪着无数干渴的马群——就在今天早上,它们的嘶鸣还把我从床上拽起来,时辰还早得不像话,而我已是快六十的人,喜欢清静。没办法,我只好起床。看着它们瘦成一副副饿马的骨架,我心里难受,而好水明明不缺,青草地也不缺。可是看来,有人对这片我记忆中殷勤好客的国土施了恶意的魔法,封住了通往那些丰沛之水的路。又或许这是本地的马贩子设的局,好把价钱压下来,谁知道呢?又或许这是个再也没有孩子牵着马去饮水的国度,马儿们只好渴着,因为再没有一个孩子找得到通往那条河的路……

\section{(16) Motif——或心中的心}

topos 这个主题出自概形主题,就在概形出现的那一年——但在广度上它远远超出了母题。正是 topos 主题(而非概形主题)才是那张“床”,或那条“深河”:几何与代数、拓扑学与算术、数理逻辑与范畴论、连续的世界与“不连续”或“离散”结构的世界,都在那里结亲。如果说概形主题好比新几何的心脏,那么 topos 主题就是它的外壳,或它的居所。它是我所构想的最宽广的东西,为的是用同一种充满几何回响的语言,精细地把握那些彼此相距最远的情境所共有的某种“本质”——那些情境来自数学事物这个辽阔宇宙的这一片或那一片。

然而 topos 这个主题远没有像概形主题那样走运。我在《收获与播种》里多处谈到此事,这里不是细说这个概念所遭遇的种种离奇变故的地方。不过新几何的两个主旋律主题都出自 topos 主题,那是两个互补的“上同调理论”,它们各自都是为着给韦伊猜想提供一条进路而构想出来的:平展(或“$\ell$-进”)主题,与晶体主题。前者在我手中化成了 $\ell$-进上同调工具,如今它已显出是本世纪最强有力的数学工具之一。至于晶体主题,在我离开之后被压缩到一种近乎隐秘的存在,最终(在需要的压力下)于1981年6月被发掘出来,在聚光灯下,用一个化名,而情形比围绕 topos 的那些还要离奇。

正如所预料的,$\ell$-进上同调工具成了建立韦伊猜想的基本工具。我自己证明了其中一大摞,而最后一步是在我离开三年后由皮埃尔·德利涅出色地完成的,他是我那些“上同调”学生中最出色的一位。

另外,约在1968年,我曾提炼出韦伊猜想的一个更强的版本,而且尤其是一个更“几何”的版本。这些猜想一直“带着”(如果可以说的话!)一层看来不可约化的“算术”色彩,而这些猜想的本意恰恰是要通过“几何的”(或“连续的”)中介,来表达和把握“算术”(或“离散的”)\fnote{52}{}{(给数学家看的)韦伊猜想从属于一些“算术”性质的假设,尤其因为所论的簇必须定义在有限域上。从上同调形式主义的观点看,这导致要给与这种情形相关联的 Frobenius 自同态一个特殊位置。在我的进路里,那些关键性质(“广义指标定理”一类)涉及任意的代数对应,并且对事先给定的基域不作任何算术性质的假设。}。
% ---------------- p.59 ----------------
\origpage{P45}\origpage{P46} 在这个意义上,我所提炼的那个版本,比韦伊本人的版本更“忠于”“韦伊的哲学”——那种从未写出、也很少被说出的哲学;在过去四十年几何的非凡兴起中,它也许是主要的、未明言的动机\fnote{53}{}{不过,在我1970年离开之后,出现了一股十分明显的反挫,具体表现为一种相对停滞的局面;我在《收获与播种》的字里行间不止一次有机会谈到它。}。我的改写,要点在于提炼出某种“精华”:在被称作“抽象”代数簇的框架内,那仍然应当有效的东西——出自对“寻常”代数簇有效的经典“霍奇理论”\fnote{54}{}{这里“寻常”的意思是:“定义在复数域上”。霍奇理论(所谓“调和积分”理论)是复数代数簇背景下已知上同调理论中最强的一个。}。我把这些著名猜想的这个全新版本(就代数闭链而言)称为“标准猜想”。

在我心里,这是在发展出 $\ell$-进上同调工具之后,朝这些猜想迈出的新的一步。但同时,而且尤其是:它也是通往那个在我看来仍是我在数学中引入的最深主题的可能进路之一\fnote{55}{}{这是最深的主题,至少在我数学家活动“公开”的时期(1950至1969年)里是这样,也就是说直到我离开数学舞台为止。我认为自1977年起发展起来的远阿贝尔代数几何与伽罗瓦–泰希米勒理论这一主题,深度与之相当。}:Motif 的主题(它本身也生于“$\ell$-进上同调主题”)。这个主题好比概形主题的心脏或灵魂,是它最隐蔽、最躲开目光的部分;而概形主题本身又处于新图景的核心。而标准猜想\fnote{56}{}{(给代数几何读者)必要时,这些猜想需要重新表述。更详细的评论见《工地巡览》(ReS IV 第178号注,第1215–1216页)以及《信念与知识》第769页的页底注(ReS III 第162号注)。}里提炼出的少数关键现象,可以看作构成了 Motif 主题的终极精华,看作这个在一切主题中最精微者的“气息”,看作新几何这颗“心中的心”的“气息”。

大致是这么回事。我们已经看到,对给定的素数 $p$,(尤其为着韦伊猜想)知道如何为“特征 $p$ 的(代数)簇”构造“上同调理论”,这有多重要。而那条著名的“$\ell$-进上同调工具”恰恰提供了这样一种理论,甚至提供了无穷多个不同的上同调理论——对每个异于特征 $p$ 的素数 $\ell$,各有一个。这里显然又有一件“缺失的理论”,它对应的正是 $\ell$ 等于 $p$ 的情形。为补这个缺,我专门设想了另一种上同调理论(前面已经提到过),叫作“晶体上同调”。另外,在 $p$ 为无穷这个重要的情形里,还另有三种上同调理论\fnote{57}{}{(给数学家读者)这些理论分别对应于 Betti 上同调(用超越方法定义,借助把基域嵌入复数域)、霍奇上同调(塞尔定义)和 de Rham 上同调(我定义);后两者已可上溯到五十年代(而 Betti 上同调则上溯到上个世纪)。}——而且谁也不能保证人们迟早不会再去引入新的上同调理论,其形式性质全都类似。于是,与寻常拓扑学里的情形相反,这里摆在我们面前的是一大堆多得令人惶惑的不同上同调理论。人们有一种很清楚的印象:在某种起初还相当模糊的意义上,所有这些理论都该“归于同一”,它们“给出同样的结果”\fnote{58}{}{(给数学家读者)例如,若 $f$ 是代数簇 $X$ 的自同态,它诱导出上同调空间 $H^{i}(X)$ 的一个自同态,那么后者的“特征多项式”应当整系数,并且不依赖于所选取的那个特定上同调理论(例如 $\ell$-进理论,$\ell$ 可变)。对于一般的代数对应也一样,只要 $X$ 假定是真且光滑的。可悲的真相(它也让人看出,在我离开之后,特征 $p>0$ 的代数簇的上同调理论处于多么可悲的荒废状态)是:这件事至今仍未得到证明,甚至在 $X$ 是射影光滑曲面、$i=2$ 这个特殊情形里也没有。事实上,据我所知,在我离开之后还没有人肯费心关注这个关键问题——它正是那种看来从属于标准猜想的典型问题。时尚的法令是:唯一值得注意的自同态是 Frobenius 自同态(德利涅用现成的家伙什把它单独处理掉了……)。}。

% ---------------- p.60 ----------------
\origpage{P47} 正是为了表达不同上同调理论之间这种“亲缘”直觉,我提炼出了与一个代数簇相关联的“Motif”概念。用这个词,我想暗示:在借助先天地可能的全部上同调理论、与簇相关联的这一大堆不同上同调不变量之下,有一个“共同的动机”(或“共同的理由”)。这些不同的上同调理论,就好比同一个“基本动机”(称为“Motif 上同调理论”)的若干不同的主题展开,各自带着自己的“速度”、“调性”和“调式”(“大调”或“小调”);而这个基本动机,同时又是所有这些不同的主题化身(也就是说,所有这些可能的上同调理论)中最根本、最“精细”的一个。这样,与一个代数簇相关联的 Motif,便构成了“终极的”、“头等的”上同调不变量;其他一切(与各种可能的上同调理论相关联者)都从它推出,如同种种音乐化身,或种种“不同的实现”。簇的“上同调”的一切本质性质,都早已能在相应的 Motif 上“读出”(或“听出”),于是,那些特殊化的上同调不变量(例如 $\ell$-进或晶体不变量)上我们所熟悉的性质与结构,便只是 Motif 内部性质与结构的忠实映照\fnote{59}{}{(给数学家读者)另一种看域 $k$ 上 Motif 范畴的方式,是把它想象成 $k$ 上有限型分离概形范畴的一种“阿贝尔包络范畴”。于是,与这样一个概形 $X$ 相关联的 Motif(或“$X$ 的 Motif 上同调”,我记作 $H_{mot}^{*}(X)$)便显出是 $X$ 的一种阿贝尔化的“化身”。这里的关键在于:正如一个代数簇 $X$ 可以有“连续变化”(它的同构类因而依赖于连续的“参数”,或“模”),与 $X$ 相关联的 Motif、或更一般地一个“变化的”Motif,也同样可以有连续变化。这是 Motif 上同调的一个侧面,它与一切经典上同调不变量的情形形成鲜明对照,包括 $\ell$-进不变量——唯一的例外是复数代数簇的霍奇上同调。\par{}这就让人想见,“Motif 上同调”是一个精细得多的不变量,比传统上纯粹拓扑的不变量更紧地圈定了 $X$ 的“算术形状”(如果我敢冒险用这个说法)。在我的 Motif 图景里,这些 Motif 构成一条极其隐蔽、极其精致的“绳索”,把一个代数簇的代数–几何性质,与由它的 Motif 所体现的“算术”性质联结起来。后者可以被看作一个就其精神本身而言“几何”性质的对象,但在其中,那些从属于几何的“算术”性质可以说被“裸露”出来了。\par{}这样,在我眼中,Motif 便是人们至今能够与一个代数簇相关联的最深刻的“形状不变量”,除了它的“Motif 基本群”之外。这两个不变量在我看来,都像是一种尚待描述的“Motif 同伦型”的“影子”(我在《工地巡览——或工具与眼光》那条注里顺便说了几句(ReS IV 第178号,见工地5(Motif),尤其是第1214页))。我觉得,正是后面这个对象,才该是任意代数簇那种难以捉摸的“算术”(或“Motif”)“形状”直觉最完美的化身。}。

以上用音乐比喻的非技术语言说出的,是一个仍然具有孩子般单纯的想法之精华,它既精细又大胆。我是在那些我认为更紧迫的基础工作之余发展这个想法的,名之曰“Motif 理论”或“Motif 的哲学(或‘瑜伽’)”,时为1963–1969年。这是一个结构丰富得令人着迷的理论,其中很大一部分至今仍是猜想性的\fnote{60}{}{这些年里,但凡有人愿意听,我就把自己的 Motif 图景讲给他听,却始终不肯费事把这件事白纸黑字地发表出来(服务于众人的其他任务,我并不缺)。这就使后来我的某些学生能更从容地劫掠,而在我那些对情形一清二楚的旧友们的温情注视之下。(参见下页页底的注。)}。

% ---------------- p.61 ----------------
\origpage{P48}\origpage{P49} 我在《收获与播种》里多次谈到这套“Motif 的瑜伽”,它格外让我牵挂。这里不是回头重述我在别处所说的地方。只需说:那些“标准猜想”再自然不过地从这套 Motif 瑜伽中流出。同时,它们也为 Motif 概念的一种可能的成形构造提供了一个进路。

这些猜想,在我看来——今天在我看来仍然如此——是代数几何中最基本的两个问题之一。这个问题,以及另一个同样关键的问题(所谓“奇点消解”),至今都还没有解决。然而,第二个问题在今天仍像一百年前一样,显得是一个显赫而令人生畏的问题;而我曾有幸提炼出的那一个,却被时尚的专断法令(就在我离开数学舞台之后的那些年里,而且正如 Motif 主题本身一样\fnote{61}{}{事实上,这个主题在1982年(晶体主题之后一年)被发掘出来,这一回用的是它原来的名字(而且形式狭窄,只在基域特征为零这一种情形里),却没有说出工匠的名字。这只是一个例子,还有若干个:某个概念或主题在我离开之后随即被当作格罗滕迪克式的幻影埋掉,随后十年或十五年里又被我的某些学生一个个发掘出来,带着谦逊的骄傲,而且(还用得着再说明吗)不提工匠……})归为格罗滕迪克式的可爱胡闹。不过我又在预先说后面的事了……

\section{(17) 发现母亲——或两个坡面}

说实话,我为了建立韦伊猜想而对它们本身所作的思考,一直是零星的。那片已开始在我面前展开、我努力去细看、去把握的全景,在广度和深度上都远远超出一次证明的假定所需,甚至远远超出这些著名猜想最初所能让人瞥见的一切。随着概形主题与 topos 主题的出现,一个意想不到的新世界骤然打开了。“那些猜想”在其中确实占据中心位置,有点像一片幅员辽阔、行省无数的帝国或大陆的首都;而那些行省大多数与这个辉煌显赫的地方只有最遥远的关联。我从未需要对自己明说,但我知道自己从此是一项大任务的仆役:去探索这个巨大而未知的世界,把握它的轮廓直到最远的边界;也去四处走遍,以顽强而有条理的心思清点那些最近、最易到达的行省,为它们绘制忠实而精确得一丝不苟的地图,连最小的村落、最小的茅屋也要有它的位置……

主要是后面这项工作吸走了我大部分精力——那是一项耐心而庞大的基础工作,只有我一人看得清楚,尤其是只有我一人“用五脏六腑感觉得到”。从1958年(概形主题与 topos 主题接连出现的那一年)到1970年(我离开数学舞台的那一年),占去我时间远为最多的正是它。

而且我常常咬着嚼子忍耐:被这样拖住,像被一块黏人的重物坠着,没完没了地做那些(一旦看清了要害之后)在我看来更近乎“后勤”而非向未知冲刺的活儿。我得不断地按住那股向前冲的冲动——那是拓荒者或探险者的冲动,他出发去发现、去探索一个个未知而无名的世界,那些世界不停地呼唤我,要我去认识它们、为它们命名。那股冲动,以及我投在它上面的精力(几乎像是偷来的!),始终少得可怜。

然而我心底很清楚:正是这份从“任务”那里(可以说)偷来的精力,才是最稀有、最精微的本质——我作为数学家的“创造”,首先正落在那里:落在那份专注的凝神上,去把握那尚未诞生之物
% ---------------- p.62 ----------------
\origpage{P50}\origpage{P51} 在一个温暖而不竭的孕育母体中幽暗、无形、潮湿的褶皱里初现的形状与轮廓的痕迹;它似乎在呼唤我,要成形、要化身、要诞生……在发现的工作里,这份专注的凝神、这份炽热的关切是一种本质的力量,正如太阳的温暖之于埋在沃土里的种子的幽暗孕育,之于它们卑微而奇迹般地迎光绽放。

在我作为数学家的工作中,我看得最清楚的是这两股力量或冲动,它们同样深,而(我觉得)性质不同。为了说它们,我用了“建造者”的形象,以及“拓荒者”或“探险者”的形象。把两者并排放着,我忽然觉得它们都太“阳”,太“男性”,甚至“大男子气”!它们带着神话般高傲的回响,或“重大时刻”的回响。它们肯定是被我身上旧日那种“英雄式”的创造工作观(那超级阳刚的观点)的残迹所激发的。就这样,它们给出一种浓重着色的、甚至僵硬的、“立正”式的看法,去看一个远比这流动、远比这谦卑、远比这“简单”的现实——一个活生生的现实。

然而在这股“建造者”的男性冲动里(它似乎不断把我推向新工地),我同时也清楚辨认出恋家者的冲动:深深依恋“那个”家的人的冲动。首先,那是“他的”家,是“亲近者”的家——是一个亲密的活体之所在,他觉得自己是它的一部分。只是到了后来,随着被感受为“亲近”的圈子逐渐扩大,它才也成为一座“众人的家”。而在这股“造房子”的冲动里(就像人们“做”爱那样……),也有,而且首先有,一份柔情。还有与那些材料接触的冲动:一件一件地塑造它们,带着爱怜的用心,而只有通过这种充满爱的接触才能真正认识它们。而当墙立起来、梁和屋顶架好之后,还有一种深深的满足:一间一间地安置,看着在这些厅堂、房间和小角落里,渐渐建立起这座活着的房子的和谐秩序——美、好客、宜于居住。因为房子,首先并且隐秘地在每个人心里,也是母亲——是环绕我们、庇护我们的东西,既是庇护所又是安慰;而且也许(更深一层,甚至就在我们正把它整个造出来的时候)我们自身也正是从那里来的,是它在我们的出生之前那些永远被遗忘的时光里庇护了我们、哺育了我们。那也是“怀抱”。

而刚才为了越过“拓荒者”这个显赫名号、为了把握它所遮盖的那个更隐蔽的现实,自然而然浮现出来的形象,也同样剥去了一切“英雄”腔调。那里出现的,仍然是母性的原型意象——哺育的“母体”及其无形而幽暗的劳作……

这两股在我看来“性质不同”的冲动,最终比我以为的更接近。两者都属于一种“接触的冲动”,把我们带向“母亲”的相遇:带向那位既体现着切近的、“已知”之物,又体现着“未知”之物的她。委身于这一股或那一股冲动,都是“重新找到母亲”。都是重新接上与切近者、与“或多或少已知者”的接触,同时也接上与“遥远者”、与那“未知”却同时被预感着、即将被认识者的接触。

这里的差别是色调上的、配比上的,不是性质上的。当我“造房子”时,占上风的是“已知”;当我“探索”时,是未知。这两种发现的“方式”,或者更确切地说,同一个过程或同一项工作的这两个方面,不可分割地连在一起。两者都本质,而且互补。在我的数学工作里,我辨认出在这两种进路之间、更确切地说在其中一个占上风的时刻(或时期)与另一个占上风的时刻(或时期)之间,有一种不断的往返运动\fnote{62}{}{我在这里就数学工作所说的话,对“静观”的工作也同样成立(《收获与播种》里到处都会谈到它)。我几乎不怀疑:这是出现在一切发现工作中的东西,包括艺术家(譬如说,作家或诗人)的工作。我在这里描述的两个“坡面”,也可以看作:一个是表达及其“技术”要求的坡面,另一个是接受(种种知觉与印象)的坡面,后者靠一份专注的凝神而变成灵感。两者在工作的任何时刻都在场,而且在其中一个占上风的“时节”与另一个占上风的“时节”之间,总有这种不断的“往返”运动。}。但同样清楚的是:在任何时刻,两种方式都在场。

% ---------------- p.63 ----------------
\origpage{P52} 当我建造、布置,或者清理、打扫、归置时,定调的是工作中“阳”的或“男性”的“方式”或“坡面”。当我摸索着探索那不可把握者、那无形者、那无名者时,我便是自己存在中“阴”的或“女性”的那一面。

对我而言,问题绝不是要贬低或否弃我天性中的这一面或那一面,两者都本质——那建造并生育的“男性”,和那孕育、并庇护缓慢幽暗的怀胎的“女性”。我“是”两者——“阳”与“阴”,“男人”与“女人”。但我也知道,创造过程里最精细、最精微的本质,在“阴”的、“女性”的那一坡面一边——那谦卑、幽暗、往往其貌不扬的坡面。

工作中始终(我相信)对我施以最强魅惑的,正是这一坡面。可当时的共识却鼓励我把最明澈的那部分精力投到另一坡面,投到那在可触的“产品”中化身并自我确证的坡面——且不说那些产品还是完成、收尾了的,轮廓分明,以凿过的石头般的明证来担保自己的实在……

拉开距离看,我很清楚这些共识是怎样压在我身上的,也很清楚我是怎样“承受这重量”的——柔顺地承受!我工作中“构思”或“探索”的那一部分,直到我离开的时候都还被压缩在最低限度,好吧。然而,在这样回望我作为数学家的作品时,有一点以惊人的明显浮现出来:使这部作品有其本质与力量的,恰恰是这个如今被轻视、甚至被嘲笑或居高临下地鄙夷的坡面:那是“想法”的坡面,甚至是“梦”的坡面,绝不是“结果”的坡面。在这些纸页里,我试着界定我对我那个时代的数学所贡献的最本质的东西,用的是一种拥抱森林、而不在几棵树上流连的眼光——我看到的不是一串“大定理”的荣誉榜,而是一把活生生的、丰产想法的扇子\fnote{63}{}{并非我的作品里缺少人们所谓的“大定理”,其中还包括解决了别人提出、而我之前无人能解的问题的定理。(我在“上涨的海……”那条注(ReS III 第122号)第554页的页底注(***)里回顾了其中几个。)但正如我在这趟“巡游”开头(“视角与眼光”那一站,第6节)已经强调过的,这些定理对我来说,只有在一个大主题的滋养性语境里才获得全部意义,而那主题是由某个“丰产想法”开启的。于是它们的证明就像从源头、毫不费力地从承载它们的那个主题的本性、“深度”中流出——正如河水之浪仿佛从它水域的深处温柔地生出,没有断裂,也不费力。我在已经引过的那条注“上涨的海……”里,用别的意象表达了完全类似的意思。},它们全都汇聚到同一个广阔的图景里。

\section{(18) 孩子与母亲}

当这篇“前言”开始变成一趟穿越我数学作品的巡游,还加上我那段关于“(血统纯正的)继承人”和“(无可救药的)建造者”的小小品时,这个没写成的前言也开始有了一个名字:它该叫“孩子与建造者”。接下来几天里越来越清楚:“孩子”和“建造者”是同一个人物。于是这名字变得更简单了:“孩子—建造者”。天哪,这名字倒不缺少气派,而且完全对我的胃口!

可只见反思表明:这位高傲的“建造者”,或者(更谦逊地说)那个玩造房子游戏的孩子,只是那个著名的“玩耍的孩子”的一副面孔,而他有两副。还有另一个孩子:他爱探索东西,爱到处乱掏,爱钻进沙堆,钻进泥泞无名的淤滩,
% ---------------- p.64 ----------------
\origpage{P53}\origpage{P54} 钻进最不可能、最荒唐的角落……大概是为了掩人耳目(哪怕只是瞒过我自己……),我起初给他起了个火亮的名字“拓荒者”,后来又换成“探险者”——这名字更接地气,却仍罩着一圈威望的光晕。简直要问:在“建造者”和“拓荒者–探险者”之间,哪一个更男性、更诱人!正面还是反面?

可再凑近一看,我们这位无畏的“拓荒者”原来是个女孩(我却乐意把她打扮成男孩)——她是水洼、雨、细雨和夜的姐妹,沉默,几乎看不见,因为她总退到阴影里——那个总被遗忘的人(人们不对她做出一副嘲弄的样子就已经不错了……)。而我也真有办法,一天又一天地把她忘掉——可以说,是双倍地忘掉:我起初只愿意看见那个男孩(那个玩造房子的孩子……)——即便后来我终于忍不住,到底还是看见了另一个,我看到的她,也仍然是个男孩……

至于我那趟巡游的好名字,这下就完全站不住了。那是个全“阳”的、全“大男子”的名字,一个瘸腿的名字。要让它站得不歪,就得把另一个也写进去。可怪的是,“另一个”并没有真正的名字。勉强贴得上的只有“探险者”,但那还是个男孩的名字,没办法。这里的语言是个婊子,它甚至不让你察觉就把你套住,显然跟祖传的偏见串通一气。

也许可以用“造房子的孩子与探索的孩子”敷衍过去。把这一点留在不言之中:一个是“男孩”,另一个是“女孩”,而这是同一个男孩–女孩的孩子,他在建造时探索,在探索时建造……可是昨天,除了这观照与探索、命名与建造的阴阳两坡面之外,事情的又一面还显现了出来。

宇宙、世界,乃至万有,归根到底是陌生而遥远的东西。它们并不真正与我们相干。在我们最深处,求知的冲动并不是把我们带向它们。吸引我们的,是它们可触而切近的化身,最近、最“肉感”的那个,充满深沉的回响与奥秘——那化身与我们血肉之躯的起源相混同,也与我们族类的起源相混同——也是那自古以来就等待着我们、沉默着准备接纳我们的她,“在路的另一端”。冲动从她那里涌出,欲望的道路从她那里伸展——从她,母亲,从那位生我们如同生这世界的她;而那些道路把我们带向与她的相遇,朝她奔去,为的是一次次折返,消融在她之中。

% ---------------- p.65 ----------------
\origpage{P55} 就这样,在一趟不期而至的“巡游”的路上拐个弯,我意外地重又遇到一个曾让我觉得亲切、又被我有些遗忘了的寓言——孩子与母亲的寓言。可以把它看作一个关于“生命寻求自身”的寓言。或者在更谦卑的个体存在层面上,看作一个关于“存在者寻求事物”的寓言。

这是一个寓言,也是一种祖辈经验的表现,那经验深植于心灵之中——是滋养深层创造层的原初象征中最有力的一个。我相信在其中认出了以原型意象的远古语言表达出来的东西:人身上创造力的气息本身,它鼓动着人的血肉与精神,无论在它最卑微、最短暂的表现里,还是在最辉煌、最持久的表现里。

这“气息”,正如使它化身的血肉形象一样,是世上最卑微的东西。它也是最脆弱的、最被人忽视、最遭人轻贱的东西……

而这气息在你一生中的种种遭际,不是别的,正是你的冒险,你一生中“认识”的冒险。表达它的那个无言的寓言,就是孩子与母亲的寓言。

你是孩子,出自母亲,被庇护在她之中,以她的力量为食。孩子从母亲那里奔出——那至近的、熟识的母亲——去与母亲相遇:那无边的、永远未知而满是奥秘的母亲……

\begin{center}
“巡游一部作品”终
\end{center}

\subsection*{尾声:看不见的圈子}
\tocsubsub{尾声:看不见的圈子}

\section{(19) 死亡是我的摇篮(或三个娃娃给一个垂死者)}

直到五十年代末 topos 视角出现为止,空间概念的演变在我看来都是一种本质上“连续”的演变。它似乎平顺地、没有跳跃地延续着,从欧几里得对我们周围空间的理论化开始,从希腊人传下的几何学开始——那几何学致力于研究生活在这个空间里的某些“图形”(直线、平面、圆、三角形等)。诚然,数学家或“自然哲学家”构想“空间”\fnote{64}{}{我写尾声时最初的打算,是附一份极简的梗概,略述其中某些“深刻变化”,并显出我在其中看到的这种“本质连续性”。我放弃了,免得把这趟巡游拉得太长——它已经比预定的长多了!我想在《反思》第四卷拟议的“历史评注”里再谈,那一回是为数学家读者写的(这完全改变了阐述的任务)。}的方式发生过深刻的变化。但这些变化在我看来全都属于一种本质的“连续性”——它们从未把(像所有人一样)依恋着熟悉意象的数学家,突然抛进一片陌生的地方。它们就像一个人身上那些也许深刻、却是渐进的变化:我们从这人的童年就认识他,从他会走路一直看到他成年、完全成熟。有些漫长平淡的时期里,变化觉察不到;另一些时期里,也许又很剧烈。但即便在生长或成熟最强烈的时期,即便我们有几个月、甚至几年没再见到他,任何时刻都不可能有丝毫怀疑、丝毫犹豫:我们重又见到的,仍然是他,一个熟识的、亲近的存在,哪怕容貌已经变了。

另外我想可以说:到本世纪中叶,这个亲近的存在已经大大地老了——像一个终于筋疲力尽、磨损殆尽的人,被一大堆他毫无准备的新任务压垮。也许他甚至已经安然死去,却没人费心去记下一笔、验明此事。“大家”还都做出一副在活人家里忙活的样子,弄得几乎就像他确实还活得好好的。

那么,请想想这对家里的常客是多么难堪:那位可敬的老人原本僵直地端坐在扶手椅里,忽然却看见一个身量不过三个苹果高的壮实小崽子在那儿撒欢,还顺口宣称——毫不带笑,仿佛理所当然——空间先生(从现在起,您尽可随意把那“先生”二字去掉……)就是他!要是他至少有点像家里的相貌也好啊,说不定是个私生子,谁知道……可一点也没有!一眼看去,没有任何地方让人想起那位大家那么熟悉(或者以为熟悉……)的老空间父亲,那位大家无论如何都确信(这总该是最起码的……)是永恒的父亲……

% ---------------- p.66 ----------------
\origpage{P56}\origpage{P57} 这就是那著名的“空间概念的变异”。这就是我至少从六十年代初起就不得不“看见”、把它当作明摆着的事的东西,却一直没机会把它对自己明确说出来,直到我写下这些行字的此刻。而忽然,只凭这个形象的说法、只凭它立刻唤起的联想之云,我以新的清晰看见:传统的“空间”概念,以及与之紧密相关的“簇”的概念(各类簇,尤其是“代数簇”),在我来到这一带的时候,已经老得那么厉害,简直就像已经死了……\fnote{65}{}{这个断言(有人会觉得它武断)得“加点盐”来听。它的有效程度,不多不少,正如下面那个断言(我稍后也把它算在自己账上):力学的“牛顿模型”(地上的或天上的)在本世纪初已是“奄奄一息”,那时爱因斯坦赶来救驾。事实是,即便在今天,物理学里大多数“寻常”情形下,牛顿模型都完全够用,而去找相对论模型才是发疯(鉴于测量所允许的误差幅度)。同样,数学里许多情形下,旧的熟悉的“空间”与“簇”概念仍然完全够用,不必去找幂零元、topos 或“温和结构”。但在两种情形里,对于尖端研究中越来越多出现的语境,旧的概念框架都已变得无法表达哪怕最“寻常”的情境。}。我还可以说:正是随着概形视角(连同它的子嗣\fnote{66}{}{(给数学家看的)在这个“子嗣”里,我尤其算上形式概形、各类“重数”(特别是概形重数或形式重数),最后还有所谓“刚性解析”空间(由泰特引入,依循我提供的一个“总施工”,灵感来自 topos 这个新概念,同时也来自形式概形的概念)。这份清单当然远不完备……},外加一万页基础工作为其代价)、随后是 topos 视角的接连出现,一个不声不响的危机局面终于被解开了。

在刚才那幅图景里,其实不该说是一个小崽子,作为一次突然变异的产品,而该说是两个。而且这两个孩子之间有一种无可否认的“家族相似”,尽管他们一点也不像那位故去的老人。再者,凑近看还可以说:概形这小家伙,好比是已故的空间父亲(别名“各类簇”)与 topos 这小家伙之间的一个“亲缘环节”\fnote{67}{}{另外,对这两个娃娃,还该再加上第三个更小的,它出现在不那么宽厚的年代:就是“温和空间”这小崽子。正如我在别处指出过的,它没拿到出生证明,而我还是在完全不合法的状态下把它算进了我有幸引入数学的十二个“主旋律主题”之中。}。

\section{(20) 瞥一眼对门的邻居}

这局面,我觉得与本世纪初爱因斯坦的相对论理论出现时的局面非常接近。那时有一个更触目的概念死胡同,表现为一个突然的、看来无法解决的矛盾。理所当然,那个要来重整混沌的新想法,是一个孩子般单纯的想法。值得注意的是(而且符合一个再重复不过的剧本……):在所有那些忽然紧张起来、拼命想“保住家什”的聪明人、显要人、大人物当中,没有一个人想到这个想法。非得是一个无名的年轻人,刚从(说不定)学生的阶梯教室板凳上下来,跑来(也许还为自己那份大胆有点不好意思……)向他那些显赫的前辈解释该怎么做才能“拯救现象”:只要不再把空间与时间分开就行了\fnote{68}{}{这样描述爱因斯坦的想法,当然有点简略。在技术层面上,需要弄清的是在新时空上放什么结构(不过这一点其实已经“飘在空中”了,有麦克斯韦的理论和洛伦兹的想法)。这里本质性的一步不是技术性质的,而是“哲学”性质的:意识到对于相距遥远的事件而言,同时性概念没有任何实验上的实在性。就是这个“孩子气的观察”,这句“可皇帝没穿衣服!”,让人跨过了那道著名的“限定一个宇宙的、专横而看不见的圈子”……}!从技术上说,当时一切都已具备,足以让这个想法绽开并被接受。而爱因斯坦那些前辈的体面也正在于此:他们确实懂得接纳新想法,没有过多训斥。这是一个迹象,说明那还是个伟大的时代……

% ---------------- p.67 ----------------
\origpage{P58} 从数学角度看,爱因斯坦的新想法平淡无奇。但从我们对物理空间的构想来看,那却是一次深刻的变异,一次突然的“背井离乡”。这是自2400年前欧几里得提炼出物理空间的数学模型以来的第一次此类变异;那模型自古至今被所有物理学家和天文学家(包括牛顿)原封不动地拿来供力学之用,用以描述地上与天上的力学现象。

爱因斯坦这个最初的想法后来大大深化了,化成一个更精微、更丰富、更柔韧的数学模型,借助了已有数学概念这座丰富的武库\fnote{69}{}{这里主要是指“黎曼簇”的概念,以及这种簇上的张量计算。}。随着“广义相对论”,这个想法扩展为一幅物理世界的宏大图景,以同一道目光拥抱亚原子的极小微观世界、太阳系、银河与遥远的星系,以及电磁波在时空中行进——那时空每一点都被其中的物质所弯曲\fnote{70}{}{这个模型区别于空间与时间的欧几里得(或牛顿)模型、也区别于爱因斯坦最初那个模型(“狭义相对论”)的最触目的特征之一,是时空的整体拓扑形状始终不确定,而不是由模型本身的性质强制规定。这种整体形状究竟是什么,在我看来(作为数学家)是宇宙学中最迷人的问题之一。}。在宇宙学与物理学的历史上,这是第二次、也是最后一次(继三个世纪前牛顿那第一次大综合之后),出现一幅用数学模型的语言表达的、统一的宏大图景,涵盖宇宙中全部物理现象。

不过,爱因斯坦这幅物理宇宙的图景,后来也被事态淹没了。需要解释的“全部物理现象”,从本世纪初起就有时间不断扩充!出现了大量物理理论,各自以或多或少的成功,去解释那一大堆“被观察的事实”这座巨大杂物堆里有限的一捆。而人们至今还在等那个大胆的孩子,他会在玩耍中找到那把新钥匙(如果真有那么一把的话……),那个梦想中的“蛋糕模型”,肯于一下子“走通”,把所有现象一齐拯救……\fnote{71}{}{人们把这样一种假想的理论称为“统一理论”,它终将把前面谈到的众多局部理论“统一”起来、调和起来。我觉得,那项尚待着手的基本思考,将要在两个不同的层面上展开。\par 1°)一个“哲学”性质的思考,关于“数学模型”这个概念本身——为一个现实片段而设的数学模型。自牛顿理论成功以来,物理学家的一条默许公理便是:存在一个数学模型(甚至是唯一一个模型,或“那个”模型),能完美地表达物理实在,没有“脱节”,没有含糊。这个统治了两个多世纪的共识,像是毕达哥拉斯“万物皆数”那活生生的洞见留下的一块化石残迹。也许这就是新的“看不见的圈子”,取代了从前的形而上学圈子,用来限定物理学家的宇宙(而“自然哲学家”这一族类看来已彻底灭绝,被计算机那一族高奏凯歌地取代了……)。只要肯停下来哪怕一瞬,就很清楚:这个共识的有效性绝非显而易见。甚至有非常严肃的哲学理由,导致先天地怀疑它,或至少为它的有效性划定极严格的界限。现在正是时候,把这个公理交给一番严格的批判,甚至“证明”——超出一切可能的怀疑——它没有根据:不存在唯一严格的数学模型,能说明至今编目的全部所谓“物理”现象。\par 一旦“数学模型”这个概念本身、以及这样一个模型的“有效性”概念(在测量所允许的某些“误差幅度”之内)被令人满意地界定清楚,“统一理论”的问题、或至少“最优模型”(意义尚待明确)的问题,才终于会被清楚地提出来。同时,对于选择这样一个模型所带有的(也许是必然带有的)任意性程度,人们大概也会有更清楚的想法。\par 2°)在我看来,只有在这样的思考之后,提炼一个比前人的模型更令人满意的明确模型这一“技术”问题,才获得全部意义。那时也许正是时候,摆脱物理学家的第二条默许公理;这条公理可上溯到古代,深植于我们感知空间的方式本身:那就是空间与时间(或时空)的连续本性,也就是“物理现象”在其中展开的那个“场所”的连续本性。\par 想必已是十五、二十年前,我翻着构成黎曼全集的那本薄薄的书,被他“顺带”的一句评论打动了。他在那里指出:空间的终极结构完全可能是“离散的”,而我们关于它所形成的“连续”表象,也许只是对一个更复杂的实在的一种(久而久之也许过度的)简化;对人类精神来说,“连续”比“不连续”更容易把握,因此它便充当了我们把握不连续的一种“近似”。这话出自一位数学家之口,洞察力令人吃惊,而且当时物理空间的欧几里得模型还从未受到质疑;从严格的逻辑意义上说,传统上反倒是“不连续”被用作通向连续的技术进路。\par 何况近几十年来数学的发展表明,连续结构与不连续结构之间的共生,远比本世纪前半人们所想象的更为紧密。无论如何,要找到一个“令人满意的”模型(或者必要时,一组这样的模型,尽可能令人满意地“彼此衔接”),无论它是“连续”的、“离散”的还是“混合”性质的——这样的工作肯定要动用极大的概念想象力,以及一种炉火纯青的嗅觉,去把握并揭示新型的数学结构。这种想象力或“嗅觉”,我觉得是稀罕之物,不仅在物理学家中如此(那里爱因斯坦和薛定谔似乎是难得的例外),甚至在数学家中也如此(而这话我是完全有把握说的)。\par 总而言之,我预计,那期待中的更新(如果它还会来的话……)更可能出自一位骨子里的数学家、同时又熟知物理学的重大问题,而不是出自一位物理学家。但尤其是:那需要一个人具备“哲学上的开放”,才能抓住问题的症结。这个症结绝不是技术性质的,而是一个“自然哲学”的根本问题。})。

% ---------------- p.68 ----------------
\origpage{P59}\origpage{P60} 把我对我那个时代数学的贡献,与爱因斯坦对物理学的贡献相比较,这一念头是出于两个理由逼到我面前的:这两份作品都得益于我们对“空间”的构想发生的一次变异(一者取数学意义,另一者取物理意义);而且两者都采取一种统一性图景的形式,拥抱极大一批此前看来彼此分离的现象与情境。我在此看到他的作品\fnote{72}{}{我绝不声称自己熟悉爱因斯坦的作品。事实上,他的著作我一本也没读过,他的想法我只是通过道听途说、而且极为粗略地知道。不过我觉得自己还是看得出“森林”,尽管我从未费劲去细看他的任何一棵树……}与我的作品之间一种明显的精神亲缘。

这种亲缘,我觉得丝毫不被一种明显的“实质”差异所否定。正如我前面已经暗示的,爱因斯坦式的变异涉及物理空间概念,而爱因斯坦是从已知数学概念的武库里取用,从不需要扩充它,更不必颠覆它。他的贡献在于:在当时的已知数学结构中,挑出那些最适合\fnote{73}{}{关于“奄奄一息”这个形容词的评论,见前面的一条页底注(第55页的注)。}充当物理现象世界之“模型”的,用来取代前人传下的那个垂死的模型。在这个意义上,他的作品确实是一位物理学家的作品,并且在此之外,还是一位“自然哲学家”的作品——取牛顿及其同代人所理解的意义。这个“哲学”维度在我的数学作品里是没有的:我从未被引去追问那些在数学事物的宇宙中进行的“理想”概念构造,与物理宇宙中发生的现象(甚至心灵中经历的事件)之间可能有什么关系。我的作品是一位数学家的作品,他有意避开了“应用”(于其他科学)的问题,也避开了我工作的“动机”与心理根源的问题。而且是一位被其十分独特的天赋所驱使、不断扩充其技艺根基处概念武库的数学家的作品。就这样,我甚至没有察觉、仿佛在玩耍之间,被引去颠覆了对几何学家来说最根本的那个概念:空间概念(以及“簇”概念),也就是我们对几何对象所居住的那个“场所”本身的构想。

新的空间概念(像一种“广义空间”,但那些本该构成“空间”的点已经或多或少消失了)在实质上一点也不像爱因斯坦引入物理学的那概念(后者对数学家来说毫不令人迷失)。相反,比较倒该与薛定谔发现的量子力学去做\fnote{74}{}{据我理解(通过从各方传回的回声),人们一般认为本世纪物理学有过三次“革命”或大动荡:爱因斯坦的理论、居里夫妇对放射性的发现、以及薛定谔引入量子力学。}。在这种新力学里,传统的“质点”消失了,代之以一种“概率云”,它在此处与彼处的周围空间中疏密不同,取决于质点
% ---------------- p.69 ----------------
\origpage{P61}\origpage{P62} 位于该区域的概率。在这种新眼光里,人们清楚感到我们构想力学现象的方式发生了比爱因斯坦模型所体现的更深的一次“变异”——这次变异并不是简单地把一个在肩宽上略嫌狭窄的数学模型,换成一个类似但裁得更宽、更合身的模型。这一回,新模型与那些好的旧传统模型像得如此之少,以至连力学的大专家数学家也必定突然感到背井离乡,甚至迷失(或者愤慨……)。从牛顿力学走到爱因斯坦力学,对数学家来说,想必有点像从地道的普罗旺斯土话走到最时新的巴黎黑话。而走到量子力学,我想,那是从法语走到汉语。

而这些“概率云”——取代了昔日那些令人安心的物质粒子——奇异地让我想起充满 topos 的那些飘忽不定的“开邻域”:它们像一个个消散的幽灵,环绕着一些想象出来的“点”;而不肯驯服的想象力仍旧不顾一切地攀附着那些点……

\section{(21) “唯一者”——或孤独的馈赠}

这趟到“对门邻居”物理学家家里的小小远足,对一个(像大多数人那样)对数学家的世界一无所知、却肯定听人说起过爱因斯坦和他那著名的“第四维”、甚至听说过量子力学的读者,也许可以充当一个参照点。说到底,即便发明者当初并未料到他们的发现会具体化为广岛,后来又化为军备与(所谓)“和平”用途的原子竞赛,事实是:物理学上的发现有形地、几乎立刻地影响着一般人的世界。数学发现的影响,尤其在所谓“纯”数学(也就是说,不以“应用”为动机)里,则不那么直接,肯定也更难界定。譬如我就不知道,我对数学的贡献曾“服务”于什么,比如说造出个什么机器。事情如此,我当然无功可居,但这不妨碍它让我安心。一旦有了应用,就可以肯定:头一个扑上去的是军方(其后是警察)——至于工业(哪怕是所谓“和平”工业),也不见得总好多少……

当然,无论为我自己导航,还是为一位数学家读者导航,更该做的其实是在数学史本身之中找“参照点”来定位我的作品,而不是到别处去找类比。这些天我想过这事,以我对那段历史相当模糊的了解为限\fnote{75}{}{我还是个小鬼的时候,就对历史不太上心(地理也一样)。(《收获与播种》第五部分(只写了一部分)里,我曾“顺带”探到我觉得这种对历史的局部“阻塞”的深层原因——我相信近些年来这阻塞正在消解。)何况我的前辈们在“布尔巴基圈子”里所接受的数学教育,也没让事情好转——那里偶尔出现的历史引用少而又少。}。在“巡游”途中,我已经有机会提到一条数学家的“谱系”,他们的性情是我所认同的:埃瓦里斯特·伽罗瓦(Evariste Galois)、黎曼、希尔伯特。如果我对本门技艺的历史知道得更多,说不定我能把这条谱系往过去延伸得更远,或者在其中插入几个我只有耳闻的名字。让我吃惊的是:我不记得曾听说过(哪怕只是通过比我更熟悉历史的朋友或同行的暗示)除我之外还有哪个数学家,带来过一大串创新想法,而这些想法并不是彼此多少脱节的,而是像一个宏大统一图景的各个部分(物理与宇宙学中的牛顿和爱因斯坦、生物学中的达尔文和巴斯德就是这样的情形)。在数学史上,我只知道两个“时刻”诞生过一幅宏大格局的新图景。其一是2500年前古希腊数学作为我们今天所理解的科学诞生之时。另一个首先是十七世纪微积分诞生之时,那个时代以牛顿、莱布尼茨、笛卡尔等人的名字为标志。据我所知,在这两个时刻之一诞生的图景,都不是一个人的作品,而是一个时代的集体作品。

% ---------------- p.70 ----------------
\origpage{P63} 当然,从毕达哥拉斯和欧几里得的时代到十七世纪初,数学有时间改变面貌;同样,从十七世纪数学家创造的“无穷小计算”到本世纪中叶,也是如此。但据我所知,在这两个时期(一个两千多年,一个三百年)里发生的深刻变化,从未具体化、凝聚为一幅表达在某一部给定作品里的新图景\fnote{76}{}{写完这些行字几个小时后,我忽然想到:我在这里竟没有想到布尔巴基先生那部(集体)论著所力图呈现的当代数学的大综合。(在《收获与播种》第一部分里还会大量谈到布尔巴基小组。)我觉得这出于两个理由。\par 一方面,这一综合仅限于把一大批已经知道的想法与结果来一番“整理”,并不从中贡献它自己独创的新想法。若说其中有新想法,那大概是“结构”概念的一个精确数学定义;它确实成了贯穿整部论著的一条宝贵线索。但这个想法在我看来,更像一位聪明而富有想象力的词典编纂者的想法,而不像一种语言获得更新的要素——那种更新会带来对实在(这里是数学事物的实在)的新把握。\par 另一方面,从五十年代起,结构这一想法就被事态超越了:随着“范畴”方法骤然涌入数学中某些最有活力的部分,例如拓扑学或代数几何。(这样,“topos”概念就拒绝进入布尔巴基那只结构的“口袋”,那口袋确实窄得在肩宽上不够用!)布尔巴基固然是在完全知情的情况下决定不去趟这“浑水”,但这样一来,它也就放弃了最初的雄心:为整个当代数学提供基础与基本语言。\par 不过,它倒是固定下了一套语言,同时也固定下某种写作风格与处理数学的方式。这种风格起初是某种精神的(很不完整的)反映,那精神是希尔伯特活泼而直接的遗产。五十、六十年代,这种风格终于占了上风——好的方面如此,(更主要地)坏的方面也如此。近二十年来,它终于变成一种僵硬的“教规”,其所标榜的“严格”纯是门面,而当年鼓舞它的那种精神,看来已一去不返。},不像物理学和宇宙学里随大综合而发生的情形:牛顿的,然后是爱因斯坦的,在它们历史中的两个关键时刻。

看来,作为我心中诞生的那幅统一性宏大图景的仆役,我在数学史上从开端到如今都算“独一无二”。抱歉,我这样子像是要把自己抬举得超过了允许的分寸!不过,让我自己松一口气的是,我似乎辨认出一个潜在的(而且是天赐的!)兄弟。

刚才我已经有机会提到他,作为我那些“性情上的兄弟”中的第一位:埃瓦里斯特·伽罗瓦。在他短暂而闪电般的一生中\fnote{77}{}{埃瓦里斯特·伽罗瓦(1811–1832)死于一场决斗,年二十一岁。我记得他有好几部传记。我年轻时读过物理学家英费尔德写的一部传记体小说,当时给我很深的印象。},我相信辨认出了一幅宏大图景的开端——正是那“数与量的联姻”的图景,在一幅新的几何图景之中。我在《收获与播种》别处提到\fnote{78}{}{见《伽罗瓦的遗产》(ReS I,第7节)。},两年前我心中如何突然出现这个直觉:在当时最强烈地魅惑我的那项数学工作里,我正在“接续伽罗瓦的遗产”。这个直觉此后很少被提起,却有时间在沉默中成熟。我三周来对作品的回顾性思考,肯定又对此有所助益。如今我相信自己与一位过去数学家认出的最直接的谱系关系,正是把我与埃瓦里斯特·伽罗瓦连起来的那一条。不管对错,我觉得我一生中十五年发展起来、并在我离开数学舞台之后十六年里仍在我心中继续成熟、继续丰富的那幅图景——这图景也是伽罗瓦不可能不去发展的那一幅\fnote{79}{}{另外我深信,一个伽罗瓦会走得比我远得多。一方面因为他那份完全卓异的天赋(就我而言,并未分得这份天赋)。另一方面因为他大概不会像我这样,让大部分精力被那些没完没了的活儿分去,去把已经多少到手的东西逐一细致地整理成文……},
% ---------------- p.71 ----------------
\origpage{P64}\origpage{P65} 假如他处在我这个位置待在这一带,而没有早逝粗暴地截断一段壮丽的冲势。

肯定还有一个理由,促成我这种“本质亲缘”的感觉——这种亲缘不能归结为单单“数学性情”,也不能归结为一部作品的显著方面。在他的生活与我的生活之间,我还感到一种命运上的亲缘。诚然,伽罗瓦死得愚蠢,才二十一岁;而我呢,快六十了,还决心活得长长久久。可这并不妨碍:埃瓦里斯特·伽罗瓦生前,正如一个半世纪之后的我,在官方数学界里都是一个“边缘人”。就伽罗瓦的情形,浅看之下会觉得这种边缘性是“偶然的”,他不过是还没来得及靠自己的创新想法和作品“站住脚”。就我的情形,在我数学家的最初三年里,我的边缘性是由于我对数学界本身的存在一无所知(也许是故意的……)——那个我本该去面对的圈子;而自我离开数学舞台以来的十六年里,它是深思熟虑的选择的结果。肯定正是这个选择,招来了报复性的“无懈可击的集体意志”:要把我的名字的一切痕迹从数学中抹去,连同我所侍奉的那幅图景一起抹去。

但越过这些偶然的差异,我相信在这种“边缘性”里辨认出一个共同的原因,我觉得它是本质的。这个原因,我不把它看作历史境况,也不看作“性情”或“性格”上的特点(这些在我与他之间,大概就像人与人之间可能有的那么不同),当然更不在“天赋”的层面(伽罗瓦显然天赋惊人,而我比较起来只是平平)。如果真有一种“本质亲缘”,我看见它在一个远远更谦卑、远远更基本的层面上。

在我一生中,我曾在少数难得的场合感到这样的亲缘。也正因它,我觉得自己与另一位数学家“相近”,他是我的前辈:克劳德·谢瓦莱\fnote{80}{}{我在《收获与播种》里多处谈到克劳德·谢瓦莱,尤其是“与克劳德·谢瓦莱相遇——或自由与善意”那一节(ReS I 第11节),以及“向克劳德·谢瓦莱告别”那条注(ReS III 第100号注)。}。我要说的联系,是某种“天真”或“单纯”的联系,我前面有机会谈到过它。它表现为一种倾向(常不为周围人所赏识):用自己的眼睛去看事物,而不是透过某个人类群体殷勤奉送的、有专利的眼镜去看——那群体或大或小,出于这样那样的理由被赋予了权威。

这种“倾向”,或这种内心的态度,不是成熟的特权,而是童年的特权。它与生命一同,是出生时得到的馈赠——一份谦卑而可畏的馈赠。一份常常被深埋的馈赠,有些人还能多少保住它,或者也许能重新找回它……

也可以把它叫作孤独的馈赠。

% ================================================================
%  第 3 章 一封信 (前半: 原书 p.73-91, 边码 L1-L30)
%  源文: _work/src1/task-ch03a.md
% ================================================================
\chapter{一封信}

\noindent 1985年5月\par

% ---------------- p.73 ----------------
\section{(1) 千页之信}

\origpage{L1}我在这里寄给你的这份文本,是我所在大学张罗着打字、并只印了有限几册的,可它既不是抽印本,也不是预印本。它的名字——《收获与播种》——已经够清楚地预告了这一点。我把它当作一封长信寄给你——而且是一封再私人不过的信。我之所以寄给你,而不是满足于让你哪天(如果有这份好奇心)在书店里卖的某本书里看到它(倘若有哪个出版商疯到肯冒这个险……),是因为在它里面我更多是对你说话,而不是对别人。写它的时候,我不止一次想到你——得说,这封信我写了一年多,是把自己整个儿放进去写的。这是我给你的一份馈赠;写的时候我尽心竭力,把我(每一刻)能拿得出的最好的东西给出来。这份馈赠会不会被接受,我不知道——你的回答(或者你的不回答……)会让我知道……

在寄给你的同时,我也把《收获与播种》寄给数学界里所有那些同事、朋友或(旧日的)学生:我曾在某个时候与他们密切相连,或者他们以某种方式出现在我的反思里,点了名也好,没点名也好。你很有可能就在其中;而只要你用心去读,不只用眼睛和脑子去读,你必定会在那些没有点你名字的地方认出自己。我也把《收获与播种》寄给另外几位朋友,是科学家也好,不是也好。

% ---------------- p.74 ----------------
\origpage{L2}这封你正在读的“导引之信”,它向你预告并呈上一封“千页之信”(这还只是开头……),也将权作前言。写这几行的时候,那篇前言还没有写出来。此外,《收获与播种》由五个部分组成(还不算一篇“抽屉式”的导论)。我这里寄给你的是第一部分(《愚妄与更新》)、第二部分(《葬礼(1)——或中国皇帝的袍子》)和第四部分(《葬礼(3)——或四种运算》)\fnote{1}{}{我把那些以某种名义出现在我的反思里、但我并不亲自认识的同行另作别论。我只给他们寄去《四种运算》(那一部分与他们关系更近),同时寄去由这封信构成的“O 分册”,以及《收获与播种》的导论(外加前四部分全体的详细目录)。}。我觉得这几部分与你关系更近。第三部分(《葬礼(2)——或阴阳之钥》)无疑是我这份见证里最私人的部分,同时也是在我看来——比别的部分更甚——具有“普遍”价值的部分,超越它诞生时周遭的那些特殊境况。第四部分(《四种运算》)里,我这儿那儿提到这一部分;不过第四部分可以独立来读,甚至在很大程度上可以独立于前面三个部分来读\fnote{2}{(*)}{总的说来你会看到,每一个“节”(在《愚妄与更新》里)或每一则“按语”(在《收获与播种》其后三个部分的任何一个里),都有它自己的统一与自治。它可以脱离其余部分单独读,正如人看着一只手、一只脚、一根手指或一根脚趾、或整个身体上任何一大块一小块,也能看出兴味、得到乐趣,同时却不会忘记:那只是一个整体的一部分,而唯有这个整体(它始终留在未言说之中)才赋予它全部的意义。}。如果你读了这里寄去的东西,想给我回信(这也是我的愿望),或者它让你想读读缺的那一部分,那就告诉我。只要你回信让我感到你的兴趣超出了纯粹表面的好奇,我很乐意把它寄给你。

\section{(2) 《收获与播种》的诞生(一份闪电式的回顾)}

\origpage{L3}在这封前信里,我现在想用几页篇幅(如果办得到的话)告诉你,《收获与播种》里讲的是什么——说得比那个副标题本身更详细些:“一位数学家对往昔的反思与见证”(我的往昔,你猜得到的……)。《收获与播种》里东西很多,各色人等在其中看到的想必也很不同:一趟发现一段往昔的旅行;一场关于存在的冥想;一幅某个圈子、某个时代的风俗画(或者说,是从一个时代悄悄而不可抗拒地滑向另一个时代的图画……);一场调查(有时近乎侦探小说,另一些时候则近乎在数学这座大都市的底层上演的斗篷与长剑小说……);一次大规模的数学漫游(它会播下不止一颗……);一部应用精神分析的实用论著(或者,随你挑,一本“精神分析小说”);一篇关于自我认识的颂词;“我的忏悔”;一本私人日记;一种关于发现与创造的心理学;一份起诉书(照例是毫不留情的……),甚至是在“数学的美好世界”里算的一笔总账(而且不给人留半点情面……)。可以肯定的是,写它的时候我一刻也没有感到无聊,而与此同时我见识了、也看尽了各色各样的东西。要是你那些要紧的事务让你有空读它,那我倒奇怪你会读得无聊。除非是逼着你读,谁知道呢……

显然,它并不只是对数学家说话。不过也确有一些时候,它更多是对数学家说的,而不是对别人。在这封写给“《收获与播种》之信”的前信里,我却正想把那些可能与你、作为数学家,关系较近的东西概括出来、并着重凸显出来。
% ---------------- p.75 ----------------
\origpage{L4}要这么做,最自然的方式就是照直对你讲讲:我怎样一针引一线,接二连三写出了刚才提到的那四五块“砖头”。

如你所知,1970年我离开了数学的“大世界”,起因是我所属的机构(IHÉS)里一桩军费的事。此后几年,我搞反军备、搞生态,那种“文化大革命”式的行动,你大概在这儿那儿也听到过一些回声;然后我几乎从流通中消失了,天知道丢到了外省的哪所大学里。有传闻说,我把时间都花在放羊和打井上。实情是:除了许多别的营生,我还像所有人一样,硬着头皮去系里上课(这就是我那点毫无新意的糊口活儿,今天还是)。偶尔,甚至能连着几天、几星期、几个月,断断续续地重做一些数学——我手头满是装满我鬼画符的纸盒,能看懂的想必只有我自己。但那都是些与我从前做的事,至少乍看之下,极为不同的东西。1955年到1970年之间,我心爱的题目是上同调,尤其是各色簇的上同调(特别是代数簇)。我判断自己在这个方向上已经做得够多,尽可以让别人不用我也能应付;既然要做数学,那也该换一张唱片了……

1976年,一种新的激情进入我的生命,其强烈不亚于当年的数学激情,而且说来还是这激情的近亲。这就是对被我称作“冥想”的东西的激情(事物总得有个名字)。这个名字——正如任何别的名字在这里也免不了——必然会招来无数误解。与数学中一样,这是一件发现的工作。关于它,我在《收获与播种》里这儿那儿有所表述。总之,显然这里有的是活儿,够我干到死。而且,不止一次,我当真以为数学已成往事,今后我只会去管那些更严肃的事——以为我要去“冥想”了。

然而我最终还是不得不承认(那是四年前),数学的激情并未就此熄灭。甚至,连我自己也不太明白是怎么回事,并且暗自吃惊:我这个人(近十五年来)早已不再想在有生之年发表一行数学,却忽然发现自己被卷进了写一部数学著作的活儿里,那部书显然没完没了,还要一卷一卷地出;既然已经上了路,我索性要把自己以为在数学上有话要说的东西,统统倒进一套叫《数学反思》的(无穷?)系列书里,从此再没人提它。

那是两年前,1983年春。当时我已经太忙于写《追寻栈》(第一卷)了——那本该也构成《数学反思》(数学的)第一卷——顾不上追问自己身上发生了什么事。九个月后,理所当然地,这第一卷算是写完了,差不多也就是说,只剩导论要写、全稿要重读、注要加——然后就可以付印……

% ---------------- p.76 ----------------
\origpage{L5}直到此刻,那部书仍未完成——一年半以来它一动没动。那篇还欠着的导论,已经越过了(打字稿)一千二百页的大关;等它真正干完,少说也有一千四百页。你已经猜到了,这篇所谓“导论”,不是别的,正是《收获与播种》。据最新的消息,它该构成那套预定“系列”的第一、二卷,再加上第三卷的一部分。于是那套系列就此改名,干脆叫《反思》(不一定非得是数学的)。第三卷剩下的部分,主要由数学文本构成,如今这些对我比对《追寻栈》还要烫手。《追寻栈》就先等着吧,等明年再弄它的注、索引,当然还要一篇导论……

第一幕终!

\section{(3) 老板之死——弃置的工地}

\origpage{L6}我觉得,现在是时候做几点说明了:为什么我那么突然地离开了一个我在其中显然过得挺自在、过了二十多年的世界;为什么我起了那个古怪的念头,要“回来”(像个还魂的人……),而在我缺席的这十五年里,人家本来过得挺好;最后,为什么一部六七百页的数学著作的导论,竟写到一千二(或者一千四)百页。而在这里,一进入正题,我大概就要让你难受了(抱歉!),甚至惹你恼火。因为毫无疑问,你跟我从前一样,喜欢把你所属的那个圈子看得很“玫瑰色”——那个圈子里有你的位置、你的名头,诸如此类。我知道那是什么滋味……而这一回,恐怕要有点刺耳了……

我在《收获与播种》里这儿那儿谈到我离去的那段插曲,没有太在其中停留。那次“离去”在那里更像是我作为数学家的生命里一个重要的断口——我作为数学家的种种事件,都不断地相对于这个“点”来定位,作为“之前”与“之后”。要把我从一个我深深扎根的圈子、从一条早已深深划定的“轨迹”上拽开,需要一记极猛的冲击。这冲击来自碰撞:在一个我深深与之认同的圈子里,撞上了某种形式的腐败\fnote{3}{}{这里说的是各国科学家整体——以“建制”为首——与军事机器之间毫无保留的合作,把它当作经费、声望与权力的方便来源。这个问题在《收获与播种》里只顺带掠过一两次,例如去年4月2日那则《敬重》(第179号,第1221—1223页)。};对这腐败,我直到那时都选择闭上眼睛(仅仅做到自己不参与)。事后回看,我明白:在那事件之外,还有一股更深的力量在我身上起作用。那是一种强烈的、内在更新的需要。这样的更新,在一个地位显赫的科研机构那温吞的学术氛围里,是无法完成、也无法持续的。我身后是二十年紧张的数学创造、是二十年过度的数学投入——同时也是二十年漫长的、在“密闭罐”里的精神停滞……不知不觉间,我窒息了——我需要的是开阔处的空气!我那“天赐的”离去,标志着一场漫长停滞的骤然终结,也是我身上那些深刻力量走向平衡的第一步——那些力量曾被弯折、被拧紧在一种极度失衡、僵固的状态里……那次离去,真真正正是重新出发——是新旅程的第一步……

如我所说,我的数学激情并未就此熄灭。它在一阵阵零星的思考中找到了表达,走的路与我“从前”为自己划定的全然不同。至于我留在身后的那部作品,那部“从前”的作品——无论是白纸黑字印出来的,还是那部也许更为本质、却还没能找到书写之途、未能成书的作品——它大可以显得,而且在我眼里也确实显得,已经与我脱了钩。在去年之前(那是在《收获与播种》里),我从没起过念头,要在那些零零星星传回来的回声上稍作停留。我很清楚:我在数学里所做的一切,尤其在我1955到1970年那个“几何的”时期所做的,都是些本该被做出来的事;而我看见或瞥见的那些东西,都是本该出现、本该被拉到光天化日之下的东西。我也清楚:我做的工作,以及我让人做的工作,都是做得好的工作,是我整个人都投进去的工作。我把全部力气与全部爱都放了进去;而(在我看来)它从此是自治的了——一个活生生的、强健的东西——不再需要我来当它的母亲。在这方面,我离去时心里完全平静。我毫不怀疑:我留下的那些写出来的与没写出来的东西,我是把它们留在了可靠的手里;那些手会照看着它们展开、生长、繁衍,依照它们作为活物、作为强健之物的本性。

在这十五年紧张的数学工作中,在我身上绽放、成熟、长大的,是一种宏大的、统一的眼光,
% ---------------- p.77 ----------------
\origpage{L7}它化入少数几个极简单的有力观念。那眼光是一种“算术几何”的眼光,是拓扑学、几何学(代数的与解析的)与算术的综合,我在韦伊猜想里找到了它的第一个胚胎。正是它,成了那些年里我主要的灵感之源——在我看来,那些年尤其是这样一段时期:我从中提炼出这门新几何的主导观念,并打制出它几件主要的工具。这眼光、这些有力观念,对我来说已成了某种第二本性。(而在与它们断绝一切接触近十五年之后,我如今发现,这“第二本性”在我身上依然活着!)它们对我来说如此简单、如此自明,以至于理所当然地以为“所有人”都随时间与我一道把它们吸收、据为己有了。只是最近,就在这几个月里,我才明白:无论是那眼光,还是那几条一直是我恒常向导的“有力观念”,在任何一个已出版的文本里都找不到逐字的书写,至多不过是字里行间。而尤其要紧的是:我自以为已经传达出去的那眼光,以及承载它的那些有力观念,在达到完全成熟二十年后的今天,仍然不为任何人所知。我这个工匠,这些我有幸发现之物的仆役,同时也是它们在其中仍然活着的唯一一人。

\origpage{L8}我打制的这样或那样的工具,被这儿那儿用来“撬开”一个公认困难的问题,就像撬开一只保险箱。那工具显然结实。可我知道它还有另一种“力”,不只是万能钳子的力。它是一个整体的一部分,正如肢体是身体的一部分——那整体是它由之而来的,赋予它意义,它从那里取得力与生命。你尽可以用一根骨头(如果够粗)去砸开一个脑壳,这是明摆着的事。但那不是它真正的功能,不是它存在的理由。而我看见这些工具散落各处,被这个人那个人抓在手里,有点像被仔细剔净、剥光的骨头,是他们从一具身体上扯下来的——一具活的身体,而他们装作不认识它……

我在这里用字斟句酌的话说出的东西,是经过长久反思之后的结论;但它想必早已被我一点一点、含含糊糊地感知到了,年复一年,在那个尚未成形、还在寻求进入自觉的思想与形象、进入清晰分节之言说的未言说层面。我曾判定:这段过去,说到底与我无关了。那些远远近近传到我这里来的回声,尽管都被过滤过,却仍然意味深长,只消我稍作停留。我曾以为自己不过是众多工匠中的一个,在五六个正开足马力干着的“工地”上忙活\fnote{4}{}{关于这些被弃置的“工地”,我最后在三个月前那组《荒废的工地》(第176$'$至178号)里做了梳理。一年前,在发现“葬礼”之前,这事已经谈过一次,就在我重新与自己的作品接上联系、谈到它后来的遭遇的第一则按语《我的孤儿们》(第46号)里。}——也许是个更有经验的工匠,是那个长辈:当年他曾独自在这些地方干了许多年,直到一拨称心的接班人来了;是长辈,不错,但说到底与别人并无不同。可是你看,这一位一走,那就像一家石作行在老板突然去世之后宣告破产:几乎可以说一夜之间,工地就空了。“工匠们”走了,每人腋下夹走一点小零碎,觉得自家用得着。钱箱也没了,于是从此再没有任何理由让他继续费力干活了……

这同样是一种说法,是从一年多的反思与调查里澄清出来的。但毫无疑问,早在我离去后的最初几年,这已经是“在什么地方”被感知到的东西了。除了德利涅关于弗罗贝尼乌斯特征值绝对值的那些工作(就是那个“威望问题”,我最近才弄明白……)——每当我在远远近近碰上一位旧日的亲近同伴,曾与我同在工地上干过的,我问他“那么……呢?”,得到的总是同一个意味深长的手势:双臂往上一抬,像是在讨饶……显然,所有人都在忙着比那些曾放在我心头的事更重要的事;也显然,就在所有人摆出一副忙碌而要紧的神气忙活的时候,其实没做出多少东西。最要紧的
% ---------------- p.78 ----------------
\origpage{L9}那一块已经不见了——那种赋予局部任务以意义的统一,还有,我想,一种温度。剩下的是一堆从整体上脱落下来的零散任务,各人在自己的角落里孵着自己的小宝贝,或者好歹让它生点利息。

即便我想抗拒,眼看一切都戛然而止,我当然还是难受:再也听不到有人谈起 Motif,谈起 topos,谈起六种运算,谈起德拉姆系数,也听不到有人谈起霍奇系数,谈起那个“神秘函子”——它本该把各个素数对应的 $\ell$-进系数,围绕德拉姆系数,在同一把扇面上彼此联系起来;听不到有人谈起晶体(除非是得知它们始终原地不动),也听不到有人谈起“标准猜想”以及其他那些我提炼出来、显然代表着关键问题的东西。就连随《代数几何原理》(靠着迪厄多内不倦的协助)开始的那项宏大的基础工作——那几乎只消顺着已经取得的势头继续下去——也被撇在一边:所有人都满足于搬进别人耐心拼装、架设、擦亮了的墙壁与家具里。工匠一走,竟没有一个人想到轮到自己挽起袖子、动手拿起抹刀,去盖那些还等着盖的许多房子——那些好房子,能住人的,为自己,也为所有人……

我又一次忍不住,用完全自觉的形象把话接了下去,这些形象是从反思的工作里析出、浮上来的。但我毫不怀疑:这些形象早已以某种形式存在于我存在的深处。我一定早已感到,对我作品、同时也对我本人的那场“葬礼”,有着某种悄悄逼近的真实性;而这真实在去年4月19日突然强加于我,带着不可辩驳的力量,并且就叫那个名字:“葬礼”。相反,在自觉的层面上,我几乎不会想到要为此愤慨,甚至不会为此悲伤。说到底,不管是不是旧日的“亲近者”,那只是当事人自己的事:他选择拿自己的时间干什么。如果当年似乎推动过他、启发过他的东西不再启发他了,那是他的事,不是我的事。如果同样的事,以完美的一致,无一例外地落在我所有旧日的学生头上,那仍然是他们各自的事;我还有别的猫要抽,没工夫去追究这能有什么意思,就这样完事!至于我留下的那些东西,那些与我仍有某种深刻而不被察觉的联系的东西——即便它们显然被弃置在那片荒废的工地上,我自己却很清楚:它们不是那种会怕“时间的凌辱”、会怕风尚起伏的东西。如果它们还没有进入公共的遗产(当年我本以为已经进入了),它们迟早会在那里扎根,十年或一百年,说到底无关紧要……

\section{(4) 一阵葬礼的风}

可是,尽管这些年来我一直乐于回避那场大规模“葬礼”的含混感知,它却不依不饶地提醒我记起它,以别的面貌,而且不那么无害——不只是对一部作品的冷淡而已。我渐渐得知——我也说不清是怎么得知的——那被遗忘的眼光里的好几个概念,不仅已经废弃不用,而且在一个体面的圈子里,成了居高临下鄙夷的对象。情形尤其如此的是那个关键的统一性概念:topos(广义空间)——它正处在新几何的核心,正是它提供了拓扑学、代数几何与算术共有的几何直觉,也正是它让我既能提炼出平展的、$\ell$-进的层论工具,又能提炼出晶体上同调的那些主导观念(诚然,此后它们多多少少被遗忘了……)。说实话,
% ---------------- p.79 ----------------
\origpage{L10}连我自己的名字,也年复一年地、悄悄地、莫名其妙地成了取笑的对象——成了没完没了、稀里糊涂的“轰炸”的同义词(正是针对那些著名的 topos,或者那些“Motif”——他拿这些玩意儿把你耳朵都磨出茧子来了,可谁也没见过……),成了长达千页地把一根头发劈成四瓣的同义词,成了关于那些人人都一向知道、根本不用等他说的事滔滔不绝、臃肿巨大的唠叨的同义词……大致就是这类调子,不过是压低了声音,靠言外之意,带着“在高蹈之士与高雅之交之间”应有的那份体贴。

在《收获与播种》里进行的那场反思中,我想我已摸到了在各人身上起作用的那股深刻力量——它就藏在那副取笑与俯就的神气背后,面对一部其分量、其生命、其气息都非他们所能企及的作品。我还发现(撇开我这个人身上那些给我的作品与命运打下印记的特殊之处)那个秘密的“催化剂”,它促使这些力量以这种形式显现出来:面对一个完好无损的创造力所给出的意味深长的迹象,露出那种漫不经心的轻蔑。总之,在这场被取笑和轻蔑捂得悄无声息的葬礼上,他就是那位大主祭。说来奇怪,他同时也是所有人当中离我最近的那一个——而且是唯一一个曾经领会过某种眼光、并把它据为己有的人,那眼光里满是生命与强烈的力量。不过我又抢先了……

说实话,那些这儿那儿传回我这里的“隐约取笑的小阵风”,并没有特别伤到我。直到三四年前,它们在某种意义上还是无名的。我当然从中看到一个不大令人愉快的时代征兆,但它们并没有真正质疑我,也没有在我心里引起焦虑或不安。相反,更直接触动我的,是那些与我本人拉开距离的迹象,它们这儿那儿来自我昔日数学界里的不少朋友——尽管我离开了那个曾与我们共有的世界,我仍觉得自己与他们由同情的纽带相连——除了共同的激情与一段共同的过去所造就的那些纽带之外。在这方面,每一次我固然都难受,可也并没有太往心里去,而且(就我记忆所及)从没想过要把这三串迹象联系起来:被弃置的工地(以及被遗忘的眼光)、“取笑之风”,以及不少旧日朋友拉开距离。我给他们每个人都写过信,没有一个人回信。再说,从那以后,我写给旧友或学生、谈那些搁在我心头的事的信,常常没有回音,这已经不稀奇了。新时代,新风尚——我又能怎样?我只好不再给他们写信。可是(如果你是其中一位)我现在正在写的这封信,将是那个例外——一番再次向你递上的话——这一次接不接受,你自己看吧……

\origpage{L11}某些旧友与我本人拉开距离的最初迹象,如果我没有弄错,要追溯到1976年。也正是那一年,另一“串”迹象开始出现,那是我在回到《收获与播种》之前还要谈的。说得更准确些,这后两串迹象是同时出现的。就在我写下这些的此刻,我觉得它们其实不可分割,说到底,它们是同一现实的两个不同侧面或“面孔”,那一年闯进了我亲身经历的场域。至于我方才正要谈的那个侧面,那是一种系统性的“拒不受理”,\dmark 不声不响、无可辩驳,由一个“毫无裂缝的共识”\fnote{5}{}{这个“毫无裂缝的共识”,在《愚妄与更新》里这儿那儿偶尔被提到,到下一部分《葬礼(1)》里才终于成为一份详尽的见证与一番反思的对象,其中有“第X支队列”即“灵车”,由那些“棺材按语”(第93—96号)以及按语《掘墓人——或整个教团》组成。后者为《收获与播种》的这一部分收尾,同时也构成反思这“第二次吐息”的第一个成果。}专门留给1970年以后那少数几个学生及其同类——他们凭着各自的工作、工作方式与灵感,清清楚楚带着我的影响的印记。
% ---------------- p.80 ----------------
\origpage{L12}也许正是借这个机缘,我第一次察觉到那阵“隐约取笑的气息”,它通过他们,针对的是某种数学风格与某种数学进路——一种(按照当时在数学建制里显然已成普遍的共识)本不该存在的风格与眼光。

这里同样,是某种在无意识层面上被清楚感知到的东西。就在同一年,同一个荒诞的剧本(演示的是:一篇显然出色的博士论文竟无法发表)连续重复了五次,带着马戏团噱头般执拗的滑稽,它最终还是硬闯进了我自觉的注意。如今回想,我明白:当时某种现实正以一种善意的执拗在“向我示意”,而我在装聋:“喂,睁眼看看吧,大个子,留神点你眼皮底下正发生的事,这事跟你有关,真的……!!”我稍稍振作了一下,看了一眼(就那么一小会儿),半是发懵,半是心不在焉:“哦,是啊,好吧,有点怪,看来这儿是有人跟谁过不去,有什么东西横竖过不去,而且还这么齐整,说实话简直难以置信!”

这事简直难以置信到了这个地步,以至于我赶紧把那个噱头、把那场马戏都忘了。确实,我并不缺别的有趣的事可做。可这并没挡住马戏在随后几年里继续提醒我记起它——不再是噱头的调子,而是暗中以羞辱为乐的调子,或者照准脸上抡一拳的调子;只不过这里的人物都是些高雅人士,那一拳在这里也就必然采取更高雅的形式,但一样有效,全凭那些高雅人士的创意……

我感受到如同“照准脸上抡来的一拳”(是别人挨的)的那个事件,发生在1981年10月\fnote{6}{}{这段插曲见按语《棺材之三——或有点过于相对的雅可比簇》(第95号),尤其第404—406页。}。那一次,而且是从种种新精神的迹象不断向我袭来以来的第一次,我被打中了——大概比拳头落在我自己身上还要更重;那一拳是别人挨的,而那个人是我所疼爱的。他多少有点像个学生,而且是一位天分非凡的数学家,刚做出了一些漂亮的东西——不过这是个细节,说到底。不是细节的,是当时我有三个“从前的”学生,直接与那件事站在一边,而当事人(并非没有道理地)把那件事当作一种羞辱与侮辱。另外两个旧日的学生,早已有机会居高临下地对待过他,像有钱人打发一个穷酸书生那样\fnote{7}{}{在上一则脚注所引的那则按语里顺带谈到过。}。还有一个学生,三年后也会跟上同样的步子(而且依旧是“照脸一拳”的风格)——不过这一点我当时当然还不知道。当时叫住我的东西,已经绰绰有余。仿佛我那段从未被审视过的数学家的过去,忽然借着我五个旧日学生中某些人的嘴脸,朝我做了一个狰狞的鬼脸——这些人如今都成了要人,有权有势,而且倨傲……

那本该是再也合适不过的时机,去停下来,去探问那突然以如此暴力向我招呼的东西究竟是什么意思。可在我身上的某个地方早已决定(而这事从来不必说出口……):那段“从前”的过去,说到底与我无关了,我没有必要在那里停留;如果它现在似乎用一种我认得再清楚不过的声音——轻蔑时代的声音——在招呼我,那一定是认错了人。然而,我连着几天、也许几周,被焦虑捆得紧紧的,却连承认它都不肯。(只是去年,借着写《收获与播种》使我重回到这段插曲,我才终于得知有这份焦虑,它一出现就被控制住了。)我没有去确认它、探问它的意义,而是四处活动,我写信给
% ---------------- p.81 ----------------
\origpage{L13}左右各方,写“那些非写不可的信”。当事人甚至还费心给我回了信,自然是些含糊其辞的信,根本不触及任何实质。浪头终于平息,一切恢复秩序。直到去年,我几乎没再想过这事。可这一次,它留下了一道伤,或者更像一根扎人的刺,让人避着不去碰;一根养着那道伤口的刺,而那伤口只等着愈合……

这大概是我作为数学家的生涯里经历过的最痛苦、最难熬的经验——当我有幸看到(却又不肯真正去认清我眼睛所看见的):“某个我疼爱过的旧日学生或同伴,以悄悄压碎另一个我疼爱的人为乐,并且在那人身上认出了我”。它当时给我的打击,肯定比去年我那些相当疯狂的发见还要重——那些发见(在浮浅的目光看来)也许显得更加不可思议……确实,这经验在同样的调子上引起了另外几次共鸣,只是没那么暴烈,当时也就有点“漏过去了”。

\origpage{L14}这也让我想起:同样是1981年,也是我与旧日学生中唯一一个在我离去后仍保持经常联系的人的关系发生急剧转折的一年;而这个人,十五年来在数学层面上一直是我“优先的对话者”。正是那一年,已经出现几年的“那些表示轻蔑的迹象”\fnote{8}{}{这段插曲见按语《两个转折》(第66号)。}“突然变得如此粗暴”,于是我中断了与他的一切数学交流。那是在前面那件“照脸一拳”的事之前几个月。事后回看,这巧合让我觉得触目;但我不认为当时自己做过任何联系。那些事被归进一个个“分开的抽屉”;何况有人已经宣布过,这些抽屉其实无关紧要——案子已经了结!

这也让我想起:还是这1981年6月,已经开过一场颇为精彩的研讨会,在不止一个意义上值得记住——那场研讨会着实配得上以“反常讨论会”这个抹不掉的名字进入历史(或者说进入历史剩下的那点东西……)。我是在去年5月2日认识它的(更确切地说,是它朝我砸下来的!),就在发现那有血有肉的“葬礼”(4月19日)两周之后——我立刻明白,我撞上了“封神”。当然,是一场葬礼的封神,但同时也是一种轻蔑的封神——轻蔑的对象,是我们这门科学存在两千多年来数学家伦理那条心照不宣、不可动摇的根基:那条最基本的规矩,即不得把从别人那里拿来的想法与结果当作自己的东西拿出来。而我此刻注意到这两个事件在时间上那引人注目的巧合——它们看上去性质与分量都大不相同——不禁愕然:这里显露出的,正是对一个人的敬重与对一门技艺、一门科学的基本伦理规则的敬重之间那种深刻而明显的联系;正是这些规则,使这门技艺的从事不成为一场“混战”,使那些以擅长此道著称、并在此定调的人整体上不成为一个毫无顾忌的“黑手党”。不过我又抢先了……

\section{(5) 旅行}

我想,关于我的“重返数学”、并一针引一线地关于《收获与播种》的写作所处的来龙去脉,我大致已经绕了一圈。是去年3月底,在《愚妄与更新》的最后一节(《一段过去的重量》,第50号)里,我才终于想到要追问这次意外重返的缘由与意义。说到“缘由”,最强的一条想必是那种既含混又迫人的感觉:那些强健有力的东西,我曾
% ---------------- p.82 ----------------
\origpage{L15}当年托付给一些充满爱意的手——“它们却在坟墓里,与风、雨、阳光的恩泽隔绝,在我看不见它们的这十五年里萎坏着”\fnote{9}{}{引自按语《坟前的旋律——或自得》(第167号),第826页。}。我一定是渐渐地明白了,而且此前从未想到要对自己说出来:最终撬开那些虫蛀的木板、放出那些被囚禁的活物的,不是别人,正是我自己;那些活物生来不是为了让它们在封闭的棺材里烂掉,而是为了在开阔的空气中绽放。而那些围绕这些包了衬垫、臃肿不堪的棺材(无疑是照着那位可敬的死者的模样……)的假作哀戚与暗含取笑的神气,想必也“终于在我心里唤醒了一根近十年来有点睡过去的斗志之弦”,以及投身混战的欲望……\fnote{10}{}{见《一段过去的重量》(第50节),尤其第137页。(**)。}。

就这样,两年前,一件本来只打算作几天、顶多几周的快速踏勘——对某一处被撇下的“工地”的探看——变成了一部N卷本的数学长篇连载,插进了那套著名的新丛书《反思》(“数学的”,等哪天把那个多余的限定词删掉再说)。而且,从我得知自己正在写的是一部要出版的数学著作的那一刻起,我也就知道:除了或多或少合乎惯例的“数学”导论之外,我还要给它加另一篇“导论”,性质更私人。我觉得重要的是,我要解释我的“重返”——那根本不是重返一个圈子,而仅仅是“重返”紧张的数学投入,重返发表我笔下写出的数学文本,为期不定。同样,我也想解释清楚我现在写数学时所怀的精神,它在某些方面与我离去前的写作精神大为不同——那是一种发现之旅的“航海日志”精神。更何况,还有些别的事搁在我心上,它们无疑与这些相关,但我觉得更为本质。我早已认定,我要慢慢来,把该说的话说出来。那些事还含混不清,对我来说,却与我将要写的这些卷书、以及它们将插入其中的《反思》的意义不可分割。绝不该把它们草草塞进去,像为占用了匆忙读者宝贵的时间而道歉似的。如果说《追寻栈》里有些东西值得他、也值得所有人知道,那恰恰是我留待在这篇导论里说的。要是二三十页不够把它们说完,我就写上四十页,甚至五十页,那又何妨——何况我并没有强迫谁来读我……

\origpage{L16}《收获与播种》就是这样诞生的。我1983年6月写下了预定那篇导论的最初几页,那是在写《追寻栈》第一卷中间的一段空档。然后去年2月我又接着干,那时我那卷书已经基本完成好几个月了\fnote{11}{}{其间我花了足足一个月,思考一个伪直线系统的“结构曲面”,那是用一条伪直线相对于这样一个系统的所有可能“相对位置”的全体来表示的。我还写了《纲领概要》,它将收入《反思》第三卷。}。我本以为这篇导论会是个机会,让我把心里还有两三处略微模糊的东西弄清楚。可我丝毫没料到,它会像我刚写完的那卷书一样,成为一趟发现之旅;而且这趟旅行所进入的世界,比我准备在那卷已成的书和随后几卷里踏勘的世界要丰富得多、广阔得多。日子一天天、一周周、一月月过去,我也不太清楚发生了什么,这趟新的旅行就这样继续下去:去发现一段过去(三十多年来被执意回避的……),发现我自己,发现我与这段过去的种种联系;也去发现数学界里某些曾与我亲近、而我了解得如此之差的人;最后甚至,在
% ---------------- p.83 ----------------
\origpage{L17}这股势头之中,还成了一趟数学发现的旅行——因为十五二十年来,我第一次有闲情回到某些我离去时留得滚烫的问题上去\fnote{12}{}{五六十年代,我常常压抑自己想投身这类油水十足、滚烫的问题的欲望,因为被没完没了的基础工作占住了;那些工作,没有人能、也没有人愿意替我接着做,而我离去之后,也没有人有心把它继续下去……}。总之我可以说,我在《收获与播种》的篇幅里进行的,是三趟彼此紧紧交织的发现之旅。而这三趟,没有一趟随着第一千二百来页那个句号而告终。我这份见证将要收拢的回声(甚至包括那由沉默构成的回声……)都将成为这趟旅行的“续篇”。至于它的终点,这趟旅行肯定属于那种从未被走到头的旅行——也许,甚至到我们死的那天也走不到头……

于是我终于又回到了出发点:如果办得到,事先告诉你,《收获与播种》里“讲的是什么”。可也真是,前头那些篇幅,即使我没去找,也已经多多少少告诉你了。也许更有意思的做法是顺着势头往下讲,而不是去“预告”。

\noindent 1985年6月\par

\section{(6) 阴影的一面——或创造与轻蔑}

前面那些篇幅,是上个月借着一段短促的“空档”写下的。其间,我终于把《四种运算》(《收获与播种》的第四部分)最后收拾停当——现在只差把眼前这封信、这封“前信”也写完(它同样摆出一副要长到无法无天的架势……),一切就终于可以付排、可以复印了。我都不再相信这事了,快一年半了,我一直“就要写完”那批出了名的按语!去年2月(还有前一年的6月),当我动手写这篇对一部数学著作来说性质有点不寻常的“导论”时,(我想)主要有三类东西是我当时想说的。首先,我想说明我重返数学活动的意图,说明我写《追寻栈》这第一卷(我刚宣布它写完)时所怀的精神,也说明我打算带着《反思》去进行一趟更广阔的数学踏勘与发现之旅时所怀的精神。从此对我来说,不再是为某个正在孕育中的数学新世界呈上一套细致入微、一丝不苟的基础。那更该是些“航海日志”,工作在其中日复一日地继续,什么也不隐瞒,就跟它真实进行的样子一样,带着它的卡壳与失手,带着它执拗的回头,也带着它突然的跃进——一件日复一日被不可抗拒地拉着往前的工作(尽管有无数意外与不测),仿佛被一根看不见的线牵着——被某种飘忽、执着而可靠的眼光牵着。一件常常在摸索中前行的工作,尤其在那些“敏感的时刻”:某种还没有名字、还没有面孔的直觉,几乎无法察觉地浮上来;或者在某个新旅程的起点,应着最初几个想法与直觉的召唤、追着它们走——它们往往飘忽,不肯被语言的网眼捉住,而恰好是那种能够轻柔地捉住它们的语言,常常还欠着。首先要做的,正是让这样一种语言从一片显然空无的、摸不着的雾霭中凝结出来。那还只是被预感、尚未被瞥见、更谈不上被“看见”和被手指触到的东西,一点点地从不可秤量之物中澄清出来,从它的阴影与雾霭的外套里脱出来,取得形状、血肉与重量……
% ---------------- p.84 ----------------
\origpage{L18}工作的这一部分,样子不起眼,甚至(许多时候)可以说干得糟糕,却同时也是最微妙、最本质的部分——真正说来,正是在那里,某种新东西出现了,靠着一种紧张的注意、一种关切、一种对这脆弱得无限微妙、即将诞生之物的敬重。这是创造的部分,是受孕与在养育之母体温暖的幽暗中缓慢孕育的部分:从那看不见的双重始基细胞起,变成无形的胚胎,又日复一日、月复一月地,在一种幽暗而紧张、看不见也无形状的劳动中,变成一个血肉俱全的新生命。

这也正是发现工作中“幽暗的”部分,那“阴”的或“女性”的部分。与之互补的一面,那“明亮”的、“阳”的或“男性”的部分,则更像抡起锤子或大槌干活,干在磨得锋利的凿子上,或干在钢火淬好的楔子上。(那是些现成可用的工具,效力早已得到验证……)两个方面各有其存在的理由与功能,彼此不可分割地共生——或者说得更好些,它们正是那对不可分离的原始宇宙力量夫妇中的妻与夫;它们不断更新的拥抱,不断重新唤起那幽暗的创造劳作:受孕、孕育与诞生——孩子、新生之物的诞生。

\origpage{L19}在那篇出了名的、私人的、“哲学的”导论里,我觉出需要说的第二件事,恰恰是关于创造工作的本性。多年来我已明白:这本性一般都被忽视,被随手拈来的套话、被祖祖辈辈的压抑与恐惧掩盖了。它究竟在多大程度上确实如此,我此后才渐渐发现,日复一日、月复一月,在《收获与播种》里进行的那场反思与“调查”的全过程中。正是在这场反思的“发令”之初,在1983年6月写下的那几页里,我第一次被一个外表平淡、实则令人目瞪口呆的事实攫住——只消稍微停留一下就会看到:我刚才说到的发现工作中那“最富创造性”的部分,在那些本该呈现这样一种工作(至少呈现它最摸得着的果实)的文本或言谈里,几乎哪儿都不露面;无论是教科书和其他教学文本,还是原创论文与论著,还是口头课程与讨论班报告等等。看来几千年来,自从数学与其他技艺、科学诞生之初起,就存在着某种“沉默的合谋”,围绕着这些“不可告人的劳作”——它们是一切新观念(或大或小)绽放的前奏,那些观念更新着我们对这个世界某一部分的认识,而我们就在这个不断创造着的世界里活着。

说白了,看来对那个侧面、那个阶段——一切发现工作(以及一般创造工作)中最关键的阶段——的认识所受到的压抑,竟如此有效,如此被那些明明亲身了解这种工作的人内化,以至于人们常常会发誓说:即便这些人,也已经把关于它的一切自觉记忆的痕迹都铲除了。有点像在一个清教到极点的社会里,一个女人会从自己的记忆里铲除掉:与每一个她自认为有责任去擤鼻涕、擦屁股的孩子相关的,那使她受孕的一刻(她勉强承受的拥抱),那漫长的几个月妊娠(被当作一件不成体统的事),以及那漫长的几个钟点分娩(被当作一场有点恶心的苦难,最后总算解脱)。

这个比喻也许显得夸张,而且,如果我把它用到我今天还记得的、我二十年前自己所属的那个数学圈子的精神上,也许确实夸张。可在我于《收获与播种》里的反思过程中,我能够认识到——尤其在这最后几个月(随着《四种运算》的写作)以触目的方式认识到——自我离开数学舞台以来,在今天那些我曾熟悉的圈子里、并且(在我看来,至少在很大程度上)在整个数学界里成为律令的那种精神,发生了令人目瞪口呆的败坏\fnote{13}{}{何况这种败坏绝不限于“数学世界”。在整个科学生活里都能看到它,而在科学生活之外,在当代世界的全球尺度上也能看到。朝这个方向所作的初步确认与反思,见开启阴阳之思的那则按语《肌肉与内脏》(第106号)。}。甚至可能,无论凭我极为特殊的数学个性,还是凭围绕我离去时的那些境况,我那次离去都在一场本已在进行的演变中起了催化剂的作用\fnote{14}{}{这就是上页脚注所引那则按语里考察的演变。它与(我本人和我作品的)“葬礼”之间的联系,在《阴的葬仪(阳葬阴之四)》《天赐的时机——或封神》《不认(一)——或召回》《不认(二)——或变形》(第124、151、152、153号)诸按语里出现并得到考察。另见较晚的(在 RS IV 里的)按语《无用的细节》(第171(v)号,(c)部分“那些什么也不像的东西——或干枯”)与《家庭相册》(第173号,c 部分“那唯一者——或默许”)。}——那场演变,我当时丝毫未察觉(我那些同事与朋友里也没有人察觉,也许只有克劳德·谢瓦莱(Claude Chevalley)是唯一的例外)。我在这里尤其想到的这种败坏的那个侧面(它只是众多侧面中的一个\fnote{15}{}{在《收获与播种》里最常处在注意中心的侧面,尤其是在两个“调查”部分(RS II,即“中国皇帝的袍子”;RS IV,即“四种运算”)里,而且也许正是最让我“目瞪口呆”的那个侧面,是这一行当伦理的败坏:它表现为肆无忌惮的掠夺、贬损与舞弊,行之于当下某些最负盛名、最出色的数学家之间,而且(在很大程度上)是在众目睽睽之下。至于另外一些更微妙的侧面——它们其实与这一侧面直接相关——我请读者去看已经引过的那则按语(第173号 c 部分)《那些什么也不像的东西——或干枯》。})是对以下这一切心照不宣的轻蔑,乃至毫不含糊的取笑:在数学里,凡是不像锤子砸在铁砧上或凿子上那种纯粹劳作的——那些最微妙的创造过程(而且往往样子最不起眼)所受的轻蔑;一切灵感、梦、眼光(无论它们多么有力、多么丰饶)所受的轻蔑;乃至(推到极处)一切观念所受的轻蔑,无论它被构想、被表述得多么清楚:凡是没有白纸黑字写出来、没有以纯粹而硬邦邦的命题形式被写出、被登记、被编目、可以进入那些灌满我们巨型计算机无穷记忆的“数据库”的东西。
% ---------------- p.85 ----------------
\origpage{L20}(借用 C. L. 西格尔(C. L. Siegel)的一个说法\fnote{16}{}{这个说法在上一页脚注所引的那则按语里被引用并加以评论。})曾经有过一次非同寻常的“压平”,数学思想的一次“收缩”:它被剥去一个本质的维度,剥去它全部的“阴影的一面”、那“女性”的一面。诚然,由于一种祖传的传统,发现工作的这一面在很大程度上一直隐而不显,可以说从没有人谈过它——但与梦的深处源头(那些滋养伟大眼光与伟大构想的源头)的活的接触,此前(据我所知)从未丧失过。看来我们从今起已经进入了一个干枯的时代:在那里,这一源头不是被从中心处切断,而是通往它的路被堵死了,凭的是普遍轻蔑那不容上诉的判决,以及嘲弄的那些代表们。

我们似乎正在接近这样一个时刻:在每个人身上,不仅将被铲除一切与源头相近的工作、那“女性气质”的工作(被讥为“稀里糊涂”、“软塌塌”、“不牢靠”——或者在另一极端,被讥为“琐屑”、“小儿科”、“瞎轰炸”……)的记忆,而且这工作本身连同它的果实也将被连根拔除:那构想、酝酿并生出新概念与新眼光的所在。那也将是我们的技艺之实践被缩减为干巴巴、空洞无物的脑力“举重与哑铃”表演的时代,缩减为在竞赛中“攻破”(“难度有口皆碑”的)问题的种种壮举不断加码的时代——那是一个“超猛男”式的、亢奋而贫瘠的肥大的时代,接在三百多年创造性的更新之后。

\section{(7) 敬重与坚毅}

可是我又离题了,预先说出了反思教给我的东西。我原本是从双重的意图出发的,这双重的意图在反思开始之前就已清清楚楚在我心里:一是作一份“意图声
% ---------------- p.86 ----------------
\origpage{L21}明”,二是(与前者紧密相连,正如刚才显出的那样)就创造工作的本性说点什么。然而还有第三个意图,它在自觉层面上想必没那么清楚,却回应着一种更深、更本质的需要。它是由那些有时令人不知所措的“招呼”引起的,那些招呼从我的数学家的过去向我传来,借着那些曾是我的学生或朋友的人的声音(或者至少是其中不少人的声音)。在表皮层面上,这需要表现为一种想“把肚子里的东西倒空”、想说几句“不中听的实话”的欲望。但更深处,肯定是一种终于要去认识一段过去的需要,而那段过去我此前一直选择回避。首先正是从这种需要里,生出了《收获与播种》。这场漫长的反思,就是我对自身这种求知冲动、对那从外部世界、从我头也不回地离开的“数学世界”不断重新传来的招呼,日复一日作出的“回答”。除了《愚妄与更新》最开头那几页——构成其最初两章(《劳动与发现》与《梦与做梦者》)的那几页——之外,而且从接下《恐惧的诞生》(第18页)的那一章起,伴着一份根本不在计划之内的“见证”,我认为,正是这种认识我的过去并完全承担它的需要,成了写作《收获与播种》时起作用的主要力量。

\origpage{L22}从数学家的世界传到我这里、并在《收获与播种》全过程中(尤其是在第二、第四部分所进行的“调查”中)以新的力量反复落回我身上的那种招呼,一上来就戴上了自得的面具,要不就是(经过“精细调配”的)鄙夷、取笑或轻蔑的面具:有时是针对我,而(更主要地)是针对那些竟敢受我启发的人(他们当然没有料到等着他们的是什么),那些被某道心照不宣、不可更改的判决“归了类”,算作与我有一手的人。于是我又一次看见那条“明显”而深刻的联系在此显现:对他人这个人的敬重(或缺乏敬重);对创造之举、对其某些最微妙最本质的果实的敬重;最后还有对科学伦理那些最明显的规则的敬重——那些规则扎根于对自己和他人最基本的敬重,我忍不住要称之为我们这一行当里的“体面规则”。这些无疑都是某种基本的、本质的“自我敬重”的各个侧面。如果我试着用一句斩钉截铁的话,来总结《收获与播种》就某个曾是我的世界、我曾与之认同二十多年的圈子所教给我的一切,我会说:那是一个失去了敬重的世界\fnote{17}{}{这里同样,这个说法不只适用于某个有限的圈子——在那里我有充分机会就近看到这事——我觉得它概括了当今整个世界某种普遍的败坏。(可与第 L19 页脚注比较。)在《收获与播种》里所进行的“调查”之总结这一更有限的框架内,这个说法出现在去年4月2日那则《敬重》(第179号)里。}。

这件事在此前那些年里已经强烈地被感觉到,即便没有被说出来。它只是不断得到确认、变得明确,总是出人意料,有时令人目瞪口呆,贯穿《收获与播种》的全过程。它清楚显现在这样一个时刻:一种“哲学的”、一般的反思,忽然变成了个人的见证(就在开启前面提到的《恐惧的诞生》那一章的《受欢迎的异乡人》一节(第9号,第18页)里)。

然而这种感知并不是以尖锐或苦涩的怨责的调子出现的,而是(凭写作的内在逻辑、凭写作所唤起的另一种态度)以一种追问的调子出现的:在这场败坏里,在今天我所确认的这种敬重的丧失里,我自己那一份是什么?正是这个追问,贯穿并支撑着《收获与播种》这第一部分,直到它最终化解为一个清楚而毫不含糊的确认\fnote{18}{}{见《竞技数学》与《旋转木马收摊了》两节(第40、41号)。}。在此以前,这场败坏对我来
% ---------------- p.87 ----------------
\origpage{L23}说,像是忽然“从天而降”,无法解释,因而更加令人愤慨、难以容忍。而在反思的过程中,我发现它一直悄悄地进行着,想必没有人能在自己周围或自己身上察觉它,贯穿整个五十与六十年代,也包括我本人。

对这个卑微事实的确认——它想必明显得很,平淡无奇——标志着见证中第一个关键转折,以及一种立竿见影的质变\fnote{19}{}{次日,见证便深化为一场对我的沉思,并在随后的几周里保持这种特殊的品质,直到《收获与播种》这“第一次吐息”结束(以《一段过去的重量》一节,第50号)。}。这是我关于自己数学家的过去、关于我自己所要学的第一件本质的事。这种对我在普遍败坏中所负责任的认知(其敏锐程度随反思的时刻而或深或浅),贯穿《收获与播种》的全过程,成了一条底色、一个提醒。尤其当我的反思显出是在调查一个时代的失宠与不公时,情形更是如此。与理解的欲望、与那种驱动并推动一切真正的发现工作的好奇心一道,正是这种卑微的认知(它一再在途中被忘掉,又不顾一切地浮上来,而且总在最料不到的地方……)使我的见证免于(我相信)沦为对世界忘恩负义的干瘪怨责,甚至免于沦为与某些曾是我的学生或朋友(或两者兼是)的人“算总账”。这种对自己的不徇情,也给了我那种内心的平静,或者说那种坚毅,使我免于徇情于人的陷阱,哪怕只是免于一种虚假“谨慎”的陷阱。凡是我认为在反思的某一刻该说的话——无论关于我自己,还是关于我某位同事、某个旧日学生或朋友,还是关于一个圈子、一个时代——我都说了,从不需要去硬压自己的迟疑。至于那些迟疑,每次我只消仔细审视它们,它们便消失得无影无踪。

\section{(8) “我的亲近者”——或默契}

\origpage{L24}这封信里,我不打算把《收获与播种》写作中(或它某一阶段中)的所有“要紧时刻”(或所有“敏感时刻”)逐一巡览\fnote{20}{}{关于《收获与播种》前三部分全体的一个简短回顾与总结,你会在这两组按语里找到:《暮熟的果实》(第179—182号)与《发现一段过去》(第183—186号)。}。我只想说:这件工作里有四个界限分明的大阶段,或者说四口“气”——像呼吸的几口气,又像一道涌起的浪涛里接连而来的几道浪;那道浪涛我也不知是怎么从一片广袤、无声、静止而又流动、无界无名的大海里涌出来的,那海便是“我”;或者更确切地说,是一片无限辽阔、无限深邃的大海,远过于它所承载、所养育的那个“我”。这几口“气”、这几道“浪”,化成了《收获与播种》目前已经写成的四个部分。每一道浪来了,都不是我呼唤来的,也丝毫未曾预料;而任何时候我都说不清它要把我带到哪里,又何时结束。而当它结束了、新的一道浪已经接了上来,我还总以为自己在某一段冲势的收尾(那也终将,在最终的最终,是《收获与播种》的收尾!),其实我早已被同一个宏大运动的另一口气抬起来、托着往前走了。只有事后回看,这运动才清楚地显出,才毫不含糊地显出:在那被当作行动、被当作涌动而经历的东西里,有一个结构。

而且,那运动当然没有随着我在《收获与播种》上打的那个句号(完全是暂时的!)而结束,
% ---------------- p.88 ----------------
\origpage{L25}也不会随着我写给你的这封信的句号而结束;这封信也是那运动的一个“时段”。它也不是在某一个1983年6月或1984年2月的某一天诞生的——那天我坐到打字机前,要写(或接着写)某部数学著作的某篇导论。它诞生了(或者更确切地说,它重又诞生了……),就在冥想进入我生命的那一天……

可是我又离题了,任由即时的形象与联想把我带走(甚至卷走……),而不是乖乖守着一条“话头”、一条预定的线索。我今天本来是要接着讲,尽可能简短地讲去年4月那场“发现葬礼”的故事:当时我已有两个星期相信《收获与播种》已经写完——而就在那三四个星期的工夫里,种种发现怎样一个接一个、一个比一个更大更不可思议地朝我劈头砸下来——大得、疯得,以至于此后好几个月,我都极难“相信我健全官能的证词”,极难摆脱面对明证时那种暗暗的难以置信\fnote{21}{}{我试着用那个故事《中国皇帝的袍子》来表达这种困难,见同名按语(第77号);后来又在那则《尽到的义务——或真相的时刻》(第163号)里回到这一点。}。这种隐秘而顽固的不信,直到去年10月(在发现“光彩夺目的葬礼”六个月之后)才终于消散,那是在我的朋友、旧日学生(不错,是隐形的)皮埃尔·德利涅(Pierre Deligne)来我家拜访之后\fnote{22}{}{这次拜访的叙述,见我刚引的那则按语(在上页脚注中)。}。我这才第一次面对那场葬礼,不再是借文本的中介——那些文本固然用意味深长的言辞对我说着贬损、掠夺与对一部作品的屠杀,说着某种数学风格与某种数学进路(以缺席的大师为其人身)的埋葬——而是这一次,以一种直接而可触的方式,带着熟悉的面貌,借着一个熟悉的声音,音调殷勤而天真。葬礼终于就在我眼前,“有血有肉”,带着那副我如今已认得的殷勤而平淡的面貌;但这是我第一次用新的眼睛、新的注意去看它。于是,那个在我前几个月的反思中已显出是我隆重葬仪上的大主祭、是那位“披着祭袍的祭司”、同时又是一场史无前例的“手术”的主要匠人与主要“受益者”、一部被交付给取笑与掠夺的作品的隐形继承人的人,就这样在我面前展开……

这次会面正落在《收获与播种》里“第三道浪”的开端,当时我刚投身那场关于阴与阳的漫长冥想,追着一串飘忽而执着的联想。当时,这段短短的插曲只留下几行回声的痕迹,一笔带过。然而它标志着一个重要的时刻,它的果实要几个月后才清楚地显现。

还有第二次这样面对“有血有肉的葬礼”的时刻。那不过是十天前的事,它又一次“在最后一刻”重新启动了一场总也不断重新开张的调查。这一次,只是一个打给让–皮埃尔·塞尔(Jean-Pierre Serre)的电话\fnote{23}{}{这几乎是逐字引自按语《掘墓人——或整个教团》(第97号,第417页)。}。这场“漫无边际”的谈话,以触目的方式、甚至超出一切预料地证实了我(不过几天前)刚刚长篇大论地对自己讲明白的东西\fnote{24}{}{在同一则按语(第173号)的 c 部分(“那唯一者——或默许”)。},而且几乎是违我本意地讲明白的:关于塞尔在我的葬礼里所扮演的角色,关于他内心里对那件“就在他眼皮底下”发生的事的一种“秘密默许”,而他从不装作看见或感觉到什么。

这里同样,理所当然,谈话是再“酷”不过、再友好不过的;而且塞尔对我这种友好的态度显然也再真诚不过、再真实不过。这并不妨碍这一次我真正看见,或者——我真想写——“摸到”了我终于不得不承认的那种“默许”;它无疑是“秘密的”(如我先前所写),但尤其
% ---------------- p.89 ----------------
\origpage{L26}是殷勤的,正如我当时所能毫无疑义地看到的那样。一种殷勤而毫无保留的默许:让该被埋掉的被埋掉;并且,凡在看来合宜之处,无论用什么手段,都要让一种真实的(塞尔对此有第一手的了解)、不受欢迎的生父,被一个伪造而受欢迎的生父替换掉……\fnote{25}{}{这几乎是逐字引自按语《掘墓人——或整个教团》(第97号,第417页)。}。

\origpage{L27}\pmark 这触目地证实了一年前已经出现的一个直觉,当时我写道\fnote{26}{}{这段引文出自同一则按语(见上页脚注),同在第417页。}:

\begin{quote}
“从这一光照下看\fnote{27}{}{所谓“这一光照”,指的是刚才谈到的那种蓄意,即不惜一切代价除掉“不受欢迎的生父”(甚至“不可容忍的”,这里沿用所引按语中的说法)。},主祭德利涅便不再显得是那个按照决定他自身生命与行动的那些深刻力量、塑造出一种风尚的人,而更像是那“合法的继承人”\fnote{28}{}{德利涅这个“继承人”的角色,是一个既隐形(德利涅已发表的文字里,没有一行能让人猜到他可能从我的嘴里学到过什么),同时又清楚地被所有人感觉到、承认的角色。这正是德利涅两面手法及其“风格”的典型侧面之一:他精熟地在这含混之上做戏,收取这种心照不宣的继承人角色所带来的好处,同时又不认那位已故的师父,并主持着规模宏大的埋葬作业。}的指定工具,服务于一种具有毫无裂缝之一致性的集体意志——那意志执着于一项不可能的任务:把我和我个人的风格从当代数学里抹掉。”
\end{quote}

如果说德利涅当时在我眼里是那种“具有毫无裂缝之一致性的集体意志”的“指定工具”(同时也是头一个、主要的“受益者”),那么塞尔现在在我眼里,就是这同一种集体意志的化身,是它毫无保留之默许的担保人;那默许给予一切数不清的舞弊与欺诈,直至那些大规模的集体蒙骗与肆无忌惮的侵占“作业”,只要这些都有助于那项“不可能的任务”——针对我这点微不足道的、已故的人身,或者针对另一位\fnote{29}{}{我在这里想到的是左格曼·梅布胡特(Zoghman Mebkhout),第一次谈到他是在导论第6节(“葬礼”),然后是在按语《我的孤儿们》(第46号)里,以及在(较晚写成的、在发现葬礼之后的)按语《一场教学的失败(二)——或创造与自得》和《一种不义与无力的感觉》(第44$'$、44$''$号)里。我是在构成葬礼“第七支队列”的那十一则按语《讨论会——或梅布胡特的层与反常》(第75—80号)中,一点一点发现那桩对梅布胡特先驱性工作的不义抹除与侵占的。关于这(第四次、也是最后一次)“运算”,一份更详尽的调查与叙述,构成了“四种运算”这场调查里最充实的部分,用的正是那个理所当然的名字:《封神》(第171(i)至171号诸按语)。}竟敢自称出于我、并且不顾所有人反对而摆出“格罗滕迪克的继承者”架势的人。

\pmark 葬礼之中有许多悖谬而令人不知所措的侧面,其中之一便是:它首先是——甚至可以说是唯独是——由那些曾是我的朋友或学生的人做成的,而在那个世界里,我从不曾知道自己有敌人。我想,首先正是由于这一点,《收获与播种》才与你关系格外近,而我这封正在写的信也才要反过来成为一次招呼。因为如果你是数学家,如果你是他们中的一个——曾是我的学生,或曾是我的朋友——那么你与这场葬礼大概并非无关:或者以行动,或者以默契,哪怕只是以你对我的沉默——对一件就在你家门口台阶前进行的事保持沉默。而如果(说来稀奇)你接受我这些卑微的话、接受它们带给你的见证,而不是闭门不出、把送信人当不受欢迎的人打发走,那么你或许就会明白:被所有人、并且有你参与(积极的参与,或心照不宣的默许)埋掉的,不只是另一个人的作品——那是他与数学相爱的果实与活生生的见证;而且,在这个埋葬(它从不叫出自己的名字……)之下更深、更隐秘的层面上,被埋掉的是你自己存在中活着而本质的一部分,是你原初的认识、爱与创造的力量;你乐于亲手借着另一个人的身体把它埋掉。

% ---------------- p.90 ----------------
\origpage{L28}在我所有的学生里,德利涅占着一个很特殊的位置,关于这一点我在反思中谈得很多\fnote{30}{}{关于这一点,尤其见 RS II 里那十七则一组的按语《吾友皮埃尔》(第60—71号)。}。他曾经是——而且远远是——最“亲近”的一个,也是唯一一个(不论学生与否)把某种宏大的眼光亲密地吸收过来、并据为己有\fnote{31}{}{这种“宏大的眼光”,德利涅确确实实“亲密地吸收过来、并据为己有”了;它对他也曾有过强大的魅力,而且至今仍不由自主地魅惑着他,尽管同时有一股强横的力量推着他去摧毁它,去炸开它根本的统一,去把散落的碎片据为己有。于是,他对一位被否认的、“已故的”师父那种隐形的敌意,正是他自身内部分裂的表现;这分裂深深地印在他离去之后的作品上——那部作品远远停在我所了解的他那相当惊人的才能之下。}的人;那眼光在我身上诞生、长大,远在我俩相遇之前。而在我所有与我共有数学激情的朋友里,最亲近的是塞尔——他同时也多少算个长辈——(同样远远地)最亲近,尤其是因为他在十年间在我的工作里扮演了一个独一无二的角色:某些重大投入的“引爆器”,以及五十、六十年代启发我数学思考的多数有力观念(直到我离去之时)的“引爆器”。他们两人与我这个人的那种很特别的关系,当然与他们各自非凡的才能不无关联;这才能使他们对同代以及其后几代的数学家拥有同样非凡的威望。撇开这些共同点,塞尔与德利涅的性情与作风,在我看来却再相异不过,在许多方面彼此正相对立。

不管怎样,如果说有哪位数学家曾以某种名义“亲近”我这个人、我的作品(而且,更有甚者,是出了名地如此),那正是塞尔与德利涅:一位是长辈,是在某种眼光孕育的关键时期我作品中的灵感之源;另一位是我最有天分的学生,而他反过来把我当作(并且一直当作,不管有没有那场葬礼……)他主要的(也是秘密的……)灵感之源\fnote{32}{}{关于这一点,见上页脚注。}。如果在我离去之后不久,一场葬礼就开动起来(并且正儿八经地变成了“死亡”),并且落实为一支没完没了的、大大小小“作业”的队列,服务于同一个目的,那么这事只可能是在这两位——前长辈与前学生(甚至前“门徒”)——协同而紧密团结的配合下完成的:一位悄悄而有效地主持着作业,同时召集我某些学生归队\fnote{33}{}{这里说的,确确实实是另外那五个以簇的上同调为主题(与德利涅一样)的学生。},他们正巴不得屠杀父亲(以那个肥硕臃肿、嗡嗡饶舌的超级娘们的怪诞可笑形象为靶子);另一位则毫无保留、无条件、无限制地对(四项)作业的继续开“绿灯”(那些作业是:贬损、屠杀、宰割、分赃,分一具取之不尽的尸首……)。

\section{(9) 清空}

\origpage{L29}正如我刚才已经暗示过的,我必须克服相当大的内心抵抗,或者更确切地说,要靠一种耐心、细致而顽强的工作让它们自行消解,才能与某些熟悉的形象分开——那些形象牢牢盘踞、惰性极大,几十年来在我身上(正如在所有人身上,想必也在你身上)取代了对现实直接而细腻的感知,具体说来,是对某个数学世界的感知,而一段过去与一部作品仍然把我与那个世界连在一起。这些形象、这些现成观念里最根深蒂固的一条是:一个国际知名的学者,甚至一个看上去是大数学家的人,竟然会去干大大小小的欺诈,这从一开始就被排除在外(哪怕只是作为例外,更不必说作为一种亲爱的
% ---------------- p.91 ----------------
\origpage{L30}习惯……);或者,即便他(又出于老习惯)自己不伸手去沾,竟会张开双臂欢迎别人搞起来的那些“(有时简直不顾一切体面感)”的作业,而在其中,由于这样或那样的原因,他还能得利。

我身上这种精神的惰性竟到了这个地步:仅仅不到两个月前,在整整一年之久的漫长反思之后,我才终于怯生生地瞥见:塞尔在这葬礼里也许也有份——而这件事如今在我看来已是明证,甚至与最近那次意味深长的谈话无关。如同对那个曾在我起步时善意接纳我的“布尔巴基圈子”的所有成员一样,在他身上尤其如此,对我来说存在着某种心照不宣的“禁忌”,围绕着他这个人。他代表着某种“优雅”的化身——这种优雅绝不限于形式,它还包括一种严格、一种一丝不苟的正直。

在我去年4月19日发现那场葬礼之前,我做梦也不会想到:我某个旧日的学生竟会在执业中做出不诚实的事,无论是对我,还是对任何人;而如果这假设落在我最有天分的那一个、也最亲近的那一个身上,在我看来就最为荒诞不经!然而,早从我离去之时起,直到其后这些年、直到今天,我一直有充分机会看清他与我的关系是何等分裂。也不止一次,我见他(看来只为取乐)运用那种使人气馁并羞辱人的权力,只要时机合适。每一次我都深受触动(大概比我愿意向自己承认的更深……)。这些都是某种深刻失序的相当意味深长的迹象;而这失序(我有充分机会确认)绝不限于他一个人,即使在我旧日学生中最狭小的圈子里也是如此。这种失序,由于丧失对他人的敬重,其触目与深刻程度,并不亚于那种表现为所谓“职业上的不诚实”的失序。尽管如此,\dmark 发现这样一种不诚实,对我仍是一个十足的意外、一记重击。

在那场令人目瞪口呆的揭示之后——随后又是一整串同一路数的“瀑布”般的发现——的几周里,我渐渐明白:在我某些学生中间\fnote{34}{}{见上页脚注。},某种舞弊早在我离去之前的那些年里就已经开始了。恰恰在我最有天分的那一个身上,这一点格外触目——就是他,在我离去之后定了调子,并且(如我刚才所写)“悄悄而有效地主持着作业”。隔了将近二十年回看,这舞弊如今在我看来已是明摆着的事,简直“扎眼”。如果说当时我选择闭上眼不看正在发生的事,一心在那个“一切都只是秩序与美”的世界里(如我乐于想象的那样)追逐“白鲸”,那么我今天确认:我当时未能承担起我应负的责任——对那些在我身边学一门我所爱的行当的学生应负的责任;那行当还不只是一门手艺,也不只是养成某种“嗅觉”。由于对出色的学生徇情,由于我乐于(凭一道心照不宣的判决)把他们当作“特殊的存在”、当作无可怀疑的人,我当时就以我那一份\fnote{35}{}{那份“贡献”,尤其见于按语《特殊的存在》(第67$'$号),以及《攀升》与《含混》两则按语(第63$'$、63$''$号),后来又(在一种略有不同的光照下)见于按语《逐出》(第169号)末尾。另一类“贡献”见于《愚妄与更新》,那是对天分不那么出色的年轻数学家的自得态度。这种对普遍败坏中所负一份责任的认识,在《竞技数学》一节(第40号)里达到顶点。}促成了我今天看到一个世界、以及一些曾为我所爱的人中间蔓延开来的那种(在我看来史无前例的)腐败。

当然,鉴于那些东西惰性极大,要与之分离——与人们习惯称之为“幻觉”的东西分离(那语调里不无几分惋惜……),而我宁愿称之为
% 下接 part1-ch03b.tex

% ================================================================
%  第 3 章 一封信 (后半: 原书 p.92-108, 边码 L31-L56)
%  源文: _work/src1/task-ch03b.md
% ================================================================
%
% ---------------- p.92 ----------------
\origpage{L31}现成的观念;关于我自己,关于一个我从前与之认同的圈子,关于一些我爱过、也许如今依然爱着的人——与这些观念“分离”,或者更确切地说,让它们从我身上自行脱落。这是一件工作,确确实实是工作,但从来不是一场搏斗——这工作带给我的,除了许多别的可贵之物,还有时而出现的悲伤时刻,却从来没有一刻悔恨,也没有一刻苦涩。苦涩是回避一种认识的手段之一,是回避一段经历所传消息的手段之一;是以关于世界、关于他人的另一种“幻觉”(不妨说是负片式的)为代价,把自己维持在对自身的一种顽固幻觉之中。

\origpage{L32}我是没有苦涩、也没有悔恨地看着这些现成的观念一件一件从我身上脱落,像卸下一件件累赘、甚至压人的重负——它们曾为我所“珍爱”,出于积习,也因为它们“从来”就在那里。它们确已变成了某种第二本性。但这“第二本性”并不是“我”。把它们一块一块地剥离,并不是撕裂,甚至也不是那种被剥夺了珍视之物的人的失落。“剥离”一词,在我这里是作为工作的报偿与果实而来的。它的标志,是一种当下即至的、有益身心的轻松,一种受欢迎的解放。

\section{(10) 运动中的四道浪}

果不其然,这封信与我动笔时所设想的全然不同。我本来主要是想就“葬礼”做一份小小的“概述”:大致情形如此如此,你信也罢不信也罢(我自己当初也几乎难以相信……),可事情就是这样,无可置疑,不管你乐意不乐意,某年某月某日某页,白纸黑字印在某期刊某本书上,去看就是了——再说,《收获与播种》里一切都已逐条拧开来细讲;见“四种运算”诸注——爱信不信!你要是宁可不去读我,自有别人替你去读……

结果这些全都没有。可这封信已经写到三十页的坎上了,而我原计划统共五六页。也不是我存心如此:一页一页写下来,我被引着对你讲的,竟是那些最本质的东西;而我当初那么急不可耐要倒空的“口袋”(这一回它可真是明明白白摆在头几页上!),到现在还没解开呢!它甚至不再挠我的手指了,那点欲望在路上就消散了。我明白了:这里不是地方……

说实话,《收获与播种》的第四部分(也是最长的一部分),名为“葬礼(3)”或“四种运算”,原本出自一则“按语”,当初正是打算做成一份小小的“概述”,好把去年那场突袭式(而且风风火火)的调查所揭示的东西大致勾勒一下——那场调查续在第二部分里(“葬礼(1)”,或“中国皇帝的袍子”)。我以为那不过是一则五页或十页的“按语”,不会更多。结果一环扣一环,调查又被它带着重新跑了起来,写成了将近四百页——几乎是我本该为之做摘要、做总结的那部分的两倍!于是那份小小的概述至今仍然缺着,而《收获与播种》里已有六百页是专门调查“葬礼”的。这确实有点蠢。不过,在把稿子交给印刷厂之前,总还来得及把它补进《导论》的第三部分(那部分早就不是在乎多十页二十页的时候了)。

《收获与播种》的五个部分(最后一部分尚未写完,恐怕几个月之内也写不完),是(三道)“冥想”之浪与(两道)“调查”之浪的交替。这里仿佛以一种缩影的方式,映照出我最近九年的生命——那九年同样也是“浪”的交替,那些浪涌自如今支配我生命的两种激情:冥想的激情与数学的激情。而说实话,\origpage{L33}%
% ---------------- p.93 ----------------
我方才用“调查”这个信手拈来的名字所称的那两个部分(或“两道浪”),恰恰是直接从我作为数学家的过往的扎根处涌出来的,被我对数学的激情、以及扎根于这激情之中的种种自我执念所驱动。

第一道浪,“愚妄与更新”,是与我的数学家过往的第一次相遇,由此引出对我当下的一次冥想——我刚刚发现,这当下正扎根于那段过往之中。诚然,这一点绝无半点预谋;可这一部分为《收获与播种》此后全部篇幅定下了“基调”,它像是一种内在的、天赐而不可或缺的准备,好让我承受紧随其后的“光华四射的葬礼”之发现——那是在第二道浪里,“葬礼(1)——或中国皇帝的袍子”。

其实,与其说这是一场“调查”,不如说这就是那场发现逐日展开的故事:它对我是怎样的冲击,我如何竭力面对这样不打招呼就劈头砸下来的东西,如何设法用我的亲身经历去安放那难以置信之物,直到它终于变得为我所熟悉,好歹让它变得可以理解。这一运动引出一个最初的、暂定的终点,在按语《掘墓人——或整个教团》(第97号)里:这是第一次尝试,要在一个多年来、如今更比任何时候都更尖锐地显出“对常理的可怕挑战”之面貌的东西里,辨认出一种解释、一种意义!

这同一场第二运动,也引出一次“病中插曲”\fnote{36}{}{这次插曲是两则按语的主题:《事件——或身体与精神》和《陷阱——或轻易与耗竭》(第98、99号),它们开启了名为《亡者(始终没有故去)》的“第十一支队列”。},逼我绝对静养,三个多月里中止了一切智力活动。那正是我自以为又一次即将把《收获与播种》收尾的时候(只差最后几件“后勤”活儿……)。去年九月底前后恢复常态活动,我正准备终于给那些还陷在困境里的按语做最后的收尾,仍以为只消再添两三则收尾的按语,其中一则讲我刚经历的那次“健康事件”。结果,一周又一周、一月又一月,又来了上千页——比已经写好的还多出一倍——而这一回,显然我还是没写完\fnote{37}{}{“始终没有写完”——哪怕只因为后面还应有一道第五部分,而在写这几行的时候,它还没有完成。}!

事实上,在这段漫长的中断期间,我几乎与一种(在搁下它时)再热不过(甚至滚烫!)的质料失去了接触;这几乎强迫我以新的眼睛回到那质料上去——如果我不想愚蠢地“了结”一个“程序”的最后收尾,而那个程序我早已与它断了活的联系。

第三道浪就是这样诞生的,诞生在《收获与播种》这一片浩瀚的运动之中——那是一道漫长的“冥想之浪”,主题是阴与阳,是事物的动态与人存在中“暗”与“明”的两面。这道冥想出自一种愿望:更深入地理解在“葬礼”中起作用的那些深层力量;然而它从一开始就获得了自身的自治与统一,并且一上来就朝向最普遍的东西,同时也朝向最私密个人的东西。正是在这道冥想之中,我发现了这件事(说实话它显而易见,只要人肯把问题提出来):在我自发地发现事物的进程里,无论在数学还是在别处,“基调”都是“阴”的、“女性”的;而且尤其重要的是,与最常见的情形相反,我始终忠实于我身上这种原初本性\fnote{38}{}{不过,这种“对我原初本性的忠实”绝非彻底的。长久以来,它仅限于我的数学工作;而在别的一切地方,尤其在我与他人的关系里,我随大流,抬高并优先看待我身上那些被感受为“男性”的特征,压抑那些“女性”的特征。关于这一点,在《一个生命的故事:三乐章的循环》(第107—110号)那一组按语里有相当详细的讨论,那组按语差不多是《阴阳之钥》的开篇。},从未让它弯折或修正,以
%
% ---------------- p.94 ----------------
\origpage{L34}迁就周遭圈子里被尊奉的主导价值。这一发现起初在我眼里只是一种有趣的好奇。只是渐渐地,它才显出是理解“葬礼”的一把关键钥匙。而且——这一点在我看来意义还要大——如今我极其清楚、不留丝毫疑问地看到:如果说,以绝非出众的智力天赋,我竟然能始终在数学工作中使出自己的全部本事,创造出一部作品,孕育出一种广大、有力而丰饶的视野,那么这完全要归功于这种忠实,归功于这种全不理会去迁就规范的心态——靠了它,我才怀着完全的信任委身于原初的求知冲动,不去削它、砍它,不减去任何造就其力量、其精细、其不可分割之本性的东西。

\origpage{L35}然而,在这道名为“葬礼(2)——或阴阳之钥”的冥想里,处在注意力中心的并不是创造力及其源头,而毋宁是“冲突”,是创造力的阻塞状态,是创造的能量因心灵中敌对力量(多半是隐秘的力量)的对抗而涣散。暴力的那些面貌——那种(表面上)“无端的”、“为了好玩”的暴力——曾不止一次让我在“葬礼”中不知所措,并让一大堆相似的亲身处境重新浮现。对这种暴力的经验,在我一生中就像是“冲突经验那坚硬的、不可化约的核”。我此前从未直面过这样一个可怕的谜:这种暴力在一般人类存在中、尤其在我自己的存在中,其存在本身与普遍性究竟是怎么回事。整个下半部(即“阴”的一面,或“衰”的一面),注意力都集中在这个谜上。正是在冥想的这一部分里,对“葬礼”之意义、对其中表达出的那些力量,一种更深的看法逐渐浮现出来。这大概也是《收获与播种》在自我认识层面上最丰饶的一部分:它让我接触到一些要害的问题与处境,并且恰恰让我感觉到那种“要害”的性质——而直到去年,这一点还一直被回避着。

在这段关于阴阳的没完没了的“离题”走完之后,我发觉自己大体上仍旧剩着那“两三则按语”要写(顶多加一两则,其中一则早已有了现成的名字:“四种运算”……),写完了就与《收获与播种》两清。

后来发生的事大家都知道了:这“最后几则按语”最终成了《收获与播种》最长的一部分,将近五百页。这就是这场运动的“第四道浪”。它也是“葬礼”的第三即最后一部分,我给它取名“四种运算”——这也是那一组按语(“对一具尸体做的四种运算”)的名字,那组按语构成了这第四次思考吐息的核心。在《收获与播种》里,这一部分才是严格意义上的“调查”部分——不过还带点调味:这场调查并不局限于纯粹的“技术”方面、即“侦探”方面;这里的思考,如同《收获与播种》别处一样,首先是由求知与求理解的欲望驱动的。这里的调子当然比“葬礼”第一部分要“带劲”得多——那会儿我还有点揉着眼睛、自问是不是在做梦呢!话虽如此,一页页披露出来的事实却常常来得正是时候,让许多此前在这里那里只是一带而过的东西当场得到说明,落实到具体而触目的例子上。也正是在这一部分里,数学上的离题占了重要位置,那是被一次重新接上的接触激发起来的(出于调查的需要)——与一种我十五年未见的质料重新接触。在光谱的另一端,还有对吾友左格曼·梅布胡特(Zoghman Mebkhout)遭遇的现场叙述(这一部分即献给他):他落到一伙手段高强、毫无顾忌的“黑手党”手里,而当初他投身于那个(固然引人入胜、表面看来又无伤大雅的)课题——各类簇的上同调——时,根本没有梦到过这些。为了给
%
% ---------------- p.95 ----------------
\origpage{L36}这整个“调查”部分那纠缠错综的按语、副按语、副副按语……的迷宫提供一条简明的线索,我请你去看目录(第167$'$号至1767号),以及这一叠里的第一则:《侦探——或玫瑰色的人生》(第167$'$号)。不过我要提一句:这则按语注明4月22日,后来有点被“事态甩在了后面”——因为一波三折,这场我当时以为(基本上)已经查到底的调查,又拖拖拉拉地继续了两个月。

\origpage{L37}这第四次吐息一口气延续了四个多月,从二月中旬直到六月底。主要是在思考的这一部分里,靠一种细致而顽强的“档案工作”,日子一天天过去、页面一页页累积,“葬礼”的现实才逐渐建立起一种具体的、可触的接触;总之,我才好歹与它“熟识”起来——尽管它在我身上引起的(并且继续引起的)那种发自脏腑的拒绝反应,一直阻碍着真正的认识。这段漫长的思考,以对德利涅来访的回顾开篇(本信中已经提到过),而以关于我与塞尔的关系、关于塞尔在“葬礼”中的角色的一番“最后一刻”的思考收尾\fnote{39}{}{见按语《家庭相册》(第173号)的 c、d、e 三部分,其中最后一部分注明的日期是6月18日(正好是十天前)。只有一则按语或按语片段日期更晚(即《一场屠杀中的五篇论文——或孝道》,第1767号,注明次日,6月19日)。你会注意到,在《收获与播种》这第四部分、即“调查部分”里,与别处相反,按语往往按逻辑顺序而非时间顺序相随。例如“葬礼”最后两则按语(构成结尾的《自深处》)注明4月7日,比我方才引的那则按语早两个半月。我还要指出:除了《葬礼(3)》真正的“调查”部分(第167$'$—1767号按语)——那部分构成殡葬大典的“第五个时辰”(《阴阳之钥》是第二个时辰)——按语都按写作顺序相随,极少例外。}。此前那个空缺,就在于我暗中把塞尔“排除在外”,借着我已说过的那个“禁忌”——如今我觉得,这或许正是直到上个月为止我对“葬礼”的理解中留下的最严重的空缺;于是这番“最后一刻”的思考,反倒成了《收获与播种》这“第四次吐息”带给我的最重要的东西:让我对“葬礼”及其间表达出的力量,有一种不那么单薄、更加充实的领会。

\section{(11) 运动与结构}

我想,关于《收获与播种》我有什么最要紧的话想对你说,都已兜了一圈说完了;为了让你先知道“这是怎么回事”,我说的肯定已经足够,足够让你判断:接下去那一千多页的信,是否“与你有关”——因而,你是否还要继续读下去。若答案是“是”,那么我觉得,再就《收获与播种》的形式补几句说明(尤其是实用方面的)是有用的。

这形式是某种精神的映照与表达,我竭力想把那种精神在前面的篇幅里“传递”出来。与我过去的出版物相比,如果说《收获与播种》里(以及它所出自的《追寻栈》里)出现了一种新的品质,那大概就是自发性。当然,这里有贯穿的线索,有大的追问,它们赋予整个思考以连贯与统一。然而思考是逐日进行的,没有预先立定的“纲领”或“计划”,从来没有什么“应该证明的东西”事先定给我。我的意旨不是证明,而是发现,是往一种未知的质料里更深地走进去,是让那些还只是被预感、被猜疑、被瞥见的东西凝结起来。我可以毫不夸张地说,在这件工作里,没有一天、也没有一夜的思考是在“预料之中”的范围内展开的——无论就观念、意象、联想而言,那些东西都是我坐到
%
% ---------------- p.96 ----------------
\origpage{L38}白纸面前、要在上面顽强地追一条线索、或者接上另一条刚冒出来的线索时,已经在那里的。每一次,思考中出现的东西都与我事先能够预言的不同——假使我竟敢勉力去预先描述我以为在眼前看见的东西的话。最常见的是,思考走上了开头全然未料的道路,通向同样未料的新风景。然而即便它沿着一条大致预定的路线走,几小时里旅程向我显示的东西,也与我上路时对它抱有的形象相去之远,正如一处真实的风景——带着它清凉的阴影与温暖的光的嬉戏,随着散步者脚步而微妙变换的透视,还有被一阵吹得草舞、吹得芦苇歌唱的微风送来的无数声响与无名芬芳……——这样一处活的、不可捉摸的风景,与一张明信片相去之远,无论那明信片多么美、多么成功、多么“准确”。

是那种一口气追下去的思考——一天之内、或一夜之内——构成不可分割的单位,可以说是活的、个体的细胞;在整片思考(《收获与播种》在此即是)之中。整场反思之于每一个这样的单位(或那些共同构成旋律的“按语”\fnote{40}{}{最初写《愚妄与更新》时,“按语”这个名字对我就是“注解”的同义词,起的是脚注的作用。出于排版上的方便,我宁愿把这些注解放到正文末尾(第1至44号按语,第141页与第171页)。这样做的原因之一是:这些“按语”或“注解”中有一些要占上一页乃至好几页,甚至比它们本该评论的正文还要长。至于思考“初稿”那些不可分割的“单位”,我当时因为找不到更好的名字,就把它们叫作“节”(比“段落”不那么令人生畏!)。\par 这一情形和文本的结构,随着下一部分而改变;那一部分最初叫“葬礼”,后来成了“葬礼(1)”(或“中国皇帝的袍子”)。这段思考接在双重按语《我的孤儿们》与《拒受一份遗产——或一个矛盾的代价》(第46、47号,第177、192页)之后,那两则按语是给《收获与播种》最后那一“节”(更确切地说,是给它本来要有的第一部分、即《愚妄与更新》)《一段过去的重量》(第50号,第131页)作注解的。后来,又给同一节添了别的注解(第44$'$、50号按语),还有别的按语来给《我的孤儿们》作注解,这些又依次生出新的作注解的按语;这一回还不算上真正的脚注——当预定的注解(写成白纸黑字之后仍然是)篇幅不大的时候。于是,从理论上说,《收获与播种》那一部分(当时它本该构成第二即最后一部分)就表现为一组给《一段过去的重量》那一“节”作的“按语”。由于惯性,这种“按语”(而非“节”)的细分法在随后三个部分里仍旧保留着;在那里,我给思考“初稿”作注解的手段兼用两种:脚注(当篇幅允许时),以及正文中引到的、写在后面的按语。\par 从排版上说,“按语”与“节”(在 RS I 里用作思考“初稿”的基本单位)的区别在于一个记号,如(1)、(2)等(即把按语编号放在括号里、排成上标,这是引注的通行做法),它或者放在该按语的开头,或者作为一种引注放在正文中提及它的适当位置。而“节”则用1至50的阿拉伯数字标出(不用那些令人生畏的下标与上标——出于实际需要,我在按语里不得不使用它们)。话虽如此,可以说《收获与播种》第一部分里“节”的功能与后面各部分里“按语”的功能并无本质区别。我在本信这一部分(《自发与结构》)中关于这一功能所作的评论,同样适用于 RS I 的“节”,尽管我一律用“按语”这个通名。\par 至于别的细则与约定,尤其是关于如何读《葬礼(1)》的目录,我请参见《导论》第7节(《葬礼的编排》),特别是第 xiv–xv 页。}……),正如一个活机体的身体之于它每一个个体细胞——那些细胞千差万别,每一个都占着一个只属于它的位置、担着一个只属于它的功能。不过有时候,在同一段一口气追下来的思考里,事后会觉察到一些重要的断口,把它分成好几个这样的单位或消息,于是每一个都得到自己的名字,由此获得自己的身份与自治。另一些时候则相反,一段由于某种原因(多半是偶然原因)被打断的思考,第二天或第三天又自发地接了下去;或者一段横跨两三天写成的思考,事后看来却像是从头到尾一口气写成的;仿佛只有睡眠的需要,才迫使我们违心地在那里插入某种断口(可以说是“生理性的”),它只由一行干巴巴的日期
%
% ---------------- p.97 ----------------
\origpage{L39}(有时还由好几行)标出来,标在所谈那则“按语”的某些相邻段落之间——这时那则按语便以一个独一无二的名字把自己区分出来。

\origpage{L40}就这样,《收获与播种》的每一则按语都有它自己的个体性,有使它区别于任何别的按语的面孔与功能。对每一则,我都设法用一个名字表达它独有的特性;那名字应当还原或唤起它“要说”的东西之本质,至少是某种本质。每一则,我真正认它,首先就是凭它的名字;此后每当需要它来帮忙,我也是凭这个名字唤它。

名字往往在我还没想到它之前就自己来了。正是它不期而至的出现,向我示意:我还在写的这则按语已近完成——它已经把要说的话说完,只等我写完眼下这一段……同样常见的是,在继续思考之前重读昨天或前天的按语时,名字也照样自发地出现。有时,在刚来的那则按语出现之后的几天或几周里,名字会有所变动,或者添上第二个我起初没想到的名字。许多按语有两个名字,表达对它所传消息的两种不同的照明,有时是互为补充的照明。这些双名中第一个向我呈现的,是在“愚妄与更新”开头不久,即《与克劳德·谢瓦莱相遇——或自由与善意》(第11号)。

只有两次,我在开始写一则按语之前就先有了名字——而且两次都后来被事态的发展搅乱了!

只有拉开距离——几周、甚至几月的距离——之后,在逐日相继的按语整体里,才会显出某种整体的运动与结构。我设法靠各种分组与再分组来把握这两者,每一组都有它自己的名字,那名字赋予它独有的存在、它的功能或消息;有点像方才那个比喻里同一个身体的器官与四肢,以及四肢的某些部位。于是,在《收获与播种》这“整体”里,有我已说过的五个“部分”,每一部分各有自己的结构:“愚妄与更新”分组成八“章”,I 至 VIII\fnote{41}{}{在《愚妄与更新》里,我有时把这些章称作《收获与播种》的“部分”;当然不要把它们与已说过的那五个部分相混,那五个部分是后来才出现的。};而构成“葬礼”的那三个部分合起来(它们也是几个月里逐渐显出来的……),则由一支漫长而庄严的、十二支“队列”I 至 XII 的行列组成。其中最后一支,更确切地说是那场“殡葬大典”(这就是它的名字)——前面十一支“队列”都朝它走去(肯定没怎么料到什么……)——其规模真是巨大,与它所隆重安葬的那部作品的规模相称:它囊括了 RS III(《葬礼(2)》)的几乎全部和 RS IV(《葬礼(3)》)的全部,连同它近八百页、一百五十来则按语(而当初,这场大名鼎鼎的典礼只预定包含两则!)。这场典礼由大主祭亲自(以他众所周知的谦逊……)手法娴熟地主持,分为九个彼此分开的“时辰”或礼仪环节,以《悼词》开场(这还用说),并以《自深处》收尾(理当如此)。这些“时辰”之中另有两个,一名《阴阳之钥》,一名《四种运算》,每一个都(远远)构成它所嵌入的《收获与播种》那部分(III 或 IV)的最大份额,并且顺带把自己的名字给了那一部分。

在《收获与播种》全书中,我像护着眼珠一样照料目录,不断改动它,以顾及源源不断涌来的未预料的按语\fnote{42}{}{这些未预料的按语里,尤其有那些“出自一则尺寸变得无法接受的脚注”的。最常见的是,我把它紧放在它所隶属的那则按语之后,给它同一个编号,必要时加上一个 $'$ 或 $''$ 号,甚至加上 $'''$——这就免去了把此后所有已写好的按语一并重新编号这桩干不下来的活儿!这些由一则脚注生出另一则的按语,在目录里前面带一个 $!$ 号(至少在《葬礼(1)》里如此)。},并
%
% ---------------- p.98 ----------------
\origpage{L41}且尽我所能让它精细地映照出思考的整体运动与其中显现的微妙结构。正是在第三、尤其是第四部分(刚才说到的“钥匙”与“四种运算”)里,这结构最为复杂、最为交错。

\origpage{L42}为了给文本保住自发性的品格、保住思考如其实际进行与实际被经历时那种未预料的面貌,我不愿意把名字放在每则按语前面——何况名字每次都是事后才出现的。所以我建议你:每读完一则按语,就翻到目录去看看这则按语叫什么;顺便也可以一眼看出它如何嵌入已经进行的思考,甚至嵌入尚未来到的思考。否则,你就有望无可救药地迷失在一堆看来难以消化、五花八门的按语里,其编号有时古怪,甚至可说令人生畏\fnote{43}{}{至于这些有时看来也许古怪的编号何以必要,我请你参见这封滔滔不绝的信里前一条脚注。};就像一个旅人在一座陌生的城市里迷了路(那城市是按世代与世纪的任性古怪地长出来的……),既没有向导,连一张地图也没有,帮他辨认方向。在准备付印的稿子里,我打算随文加入“章”以及别的按语组与节组的名字,只把按语(或节)本身除外。但即便那样,偶尔查阅目录在我看来仍不可少,免得迷失在几百则按语的一团乱麻里——它们一串接一串,绵延一千多页……

\section{(12) 自发与严格}

自发与严格,是同一品质那“暗”与“明”的两面。唯有二者的婚配,才生出文本或存在者的那种特殊品质——人们也许可以用“真之品质”这样的说法去唤它。如果说在我过去的出版物里,自发性(即便不是缺席,也是)被压缩到极小的份额,我并不认为,随着它在晚年于我身上绽放,严格性因此就变小了。毋宁说,它那位阴性的伴侣的完全在场,给了严格性一种新的维度、新的丰饶。

这种严格也针对它自己施行:它留意使自己必须做的那桩微妙的“筛选”——在意识领域里流过的大量东西中,不断地把有意义的、本质的东西从偶然的、次要的东西里澄清出来——不至于变厚、凝固成审查与自满的自动机制。唯有我们心中求知的渴望,才能唤醒并激发这样一种不笨重的警觉、这样一种灵敏,去对抗那巨大而无处不在的惰性,即那些由现成观念凿成的“(所谓)天然的坡道”——它们是我们的恐惧与我们所受的塑造的表达。

而这同一种严格、同一种警觉的注意,也朝向自发性,也朝向那些带自发性面貌的东西,为的是在那里同样分辨出成分:分辨出那些固然再自然不过的“坡道”,把它们与真正从存在深处、从原初的求知与行动冲动中喷涌出来的东西区别开来——那冲动把我们带向与世界的相遇。

在写作层面上,严格表现为一种持久的用心:用语言尽可能精细、尽可能忠实地圈定要表达的思想、情感、知觉、意象、直觉……;在所要表达之物轮廓分明之处,不满足于一个含糊或近似的词;而在所要表达之物仍被“还只是被预感”的迷雾围绕之处,也不满足于一个虚假精确(因而同样走形)的词。当我们试图就地捕捉它如其所是的样子时——也只有在那时——未知之物才向我们显露其真正的本性,也许甚至直到大白天光之中,如果它生来属于白昼、而我们的欲望又促它卸下阴影与雾的帷幕。我们的角色不是妄称描述并固定我们所不知、从我们手中逃逸的东西,而是谦卑地、炽烈地去认识那从四面八方围绕着我们的未知与神秘。

% ---------------- p.99 ----------------
\origpage{L43}这就是说,写作的角色不是记录一项研究的结果,而正是记录研究的过程本身——记录爱的工作,记录我们与我们的母亲、那“世界”、那“未知者”的爱的作品;她不知疲倦地召唤我们进入她之中,去认识她那取之不尽的躯体里的她,去到她之中那神秘的欲望之路把我们带去的任何地方。

为了传达这一过程,那些回头的修正——它们给写作的“初稿”、甚至第二稿第三稿添上细微差别、加以精确化、深化,有时加以改正——本身就是发现进程的一部分。它们构成文本本质的一部分,赋予文本全部意义。所以,放在《愚妄与更新》末尾的那些“按语”(或“注解”)——构成文本“初稿”的那五十“节”里处处引到它们——是文本不可分割的、本质的一部分。我竭力劝你边读边翻过去看,至少每读完一节的末尾(那节里出现了一处或几处对这类“按语”的指引)就去看。对《收获与播种》其余部分里的脚注也是如此;或者,当某则“按语”(在此充当“正文”)里引到后面某一则按语,那后一则便起“回看”它或注解它的作用。这一点,连同我请你不要在阅读中途离开目录的劝告,就是我要给你的阅读建议中最主要的一条。

最后一个问题,是实用性的,它将(有点平淡地)结束这封该收尾的信。有些时候有点“慌乱”,为了把《收获与播种》的各个分册准备好,交给系里的复印部去印,好让印刷(如果可能)赶在大假之前。匆忙之中,有一整页最后一刻的脚注要补进第2分册(《葬礼(1)——或中国皇帝的袍子》),却“跳”掉了。那些主要是改正某些直到最近、在写《四种运算》的过程中才出现的材料性错误。这些脚注里有一条比其他几条更重要,我想在此指出。那是对按语《受害者——或两种沉默》(第78$'$号,第304页)的一则注解。在那则按语里,我竭力——除了别的——去圈定我的一些印象(当然全是主观的):关于吾友左格曼·梅布胡特当时如何“内化”他所遭受的那种令人不安的掠夺;而这则按语被他觉得对他不公,因为我似乎几乎把他与那些掠夺者“装进了同一个袋子”。可以肯定的是,在那则按语里——它并不自称给出别的东西,只给出与某个特定“时刻”相关的印象——我只敲响了一口钟,而把另一些同样真实(也许还不那么可争议)的钟声留在未言之中(而且大概还当作不言自明的事)。不管怎样,对这一个微妙题目的思考,在一年之后,在按语《根与孤独》(第171号)里大大深化了。这一则没有引起左格曼的保留。关于同一题目的另一些思考材料,也见于《三块界石——或无辜》与《死页》(第171(x)与(xii)号)这两则按语。这三则按语都属于《封神》,那是《四种运算》里专讲对左格曼·梅布胡特作品的据为己有与窃取那场运算的部分。

我只剩下祝你阅读愉快了——也盼着轮到我读你的信!

亚历山大·格罗滕迪克

\origpage{L44}
\subsection*{后记:或一场辩论的语境与前情}
\tocsubsub{后记:或一场辩论的语境与前情}

% ---------------- p.100 ----------------
\origpage{L45}
\section{(13) 瓶装光谱仪}

这封信写完之后,整整七个月过去了;寄出去也已近四个月,连同随信而去的那个“大砖头”。而且每一本上都有我手写的题赠\fnote{44}{}{只有少数例外,主要是那些我本人不认识的同行,他们只收到试印本的第0和第4分册,作为他们积极参与我的“葬礼”的奖赏。}。像一只“漂流瓶”,更确切地说,像一整群这样的漂流瓶子,我的消息一路漂去、流转,直到这个我曾熟悉的数学微观宇宙最偏远的角落。随着日子、星期、月份一天天过去,直接与间接的回声传回来,于我不期然地,面前仿佛出现了一张数学界的巨大射线照片——那是由一架触手四伸的“瓶装光谱仪”拍下的,而我那些天真的“瓶子”就是它一根根游走的天线。于是(责任随身份而来!),我这个本来并不缺事可做的人,如今面前又多了一件新任务:破译这张射线照片,并且尽我所能,把我从中读出的东西报告出来。这将作为《收获与播种》的第六部分(也是最后一部分,我保证!)。倘若上帝假我以岁月,这一部分便将为我“晚年那部社会学巨著”加冕。眼下,先说几句初步的评论。

用来迎接我这支可怜的手工小舰队的,占压倒优势的,是一种甜腻又带三分凶气的调子,唱的是“瞧,格罗滕迪克老了,犯起偏执狂来啦”,或者“这位可真把自己当回事”——如此这般,戏就演完了!然而这种调子的信我只收到一封\fnote{45}{}{这封信出自一位曾是我的学生的人,而且他还是我的一位“同葬者”。},另有两封属于那种含蓄的、自得其乐的嘲弄\fnote{46}{}{出自两位当年在布尔巴基共事的同行,其中一位还是我初出道时善意接待过我的年长者之一。}。我的数学家收信人,包括那些曾是我的学生的人,大多数以沉默作答\fnote{47}{}{寄给数学家的共一百三十一份,至今有五十三个收信人有了动静,哪怕只是告知收到。这其中有六个是我过去的学生——另外八个人里,我一个也没得到音信。}——这沉默对我说了很多。

不过我已经收到了大量通信。绝大多数信是客气而尴尬的调子,那尴尬常常还想显得友好,仿佛是出于礼数。有两三次,我在这尴尬背后、仿佛被它滤过之后,感到一种依然活着的情感的温度。最常见的情形是:尴尬若不是以对好情感的郑重声明来表达(为他自己,或为别人),就是以恭维来表达——我这一辈子也没收到过这么多恭维!唱的是“伟大的数学家”、“了不起的篇章”(关于创造力“以及诸如此类”……)、“无可争议的作家”,等等等等。为了凑足分量,我甚至还有幸得到一句情真意切(绝非反讽)的恭维,夸我内心生活之丰富。不用说,在所有这些信里,我的通信人都小心翼翼地不进入任何问题的实处,更不会亲自卷入;调子倒更像是那种“被请来发表意见”的人(借用其中一封信的说法)面对一桩有点不干不净的事——何况还是假设的、甚至想象的——总之尤其要紧的是:一桩与他本人无关的事。可是当他装出要碰一碰其中一个问题的样子时,那是指尖轻触,并且尽可能把它推得离自己远远的——或者借给我一番好劝告之机,或者用谨慎的条件式,或者在不知\origpage{L46}%
% ---------------- p.101 ----------------
该说什么时用那套现成的套话,或者用别的什么法子。\origpage{L47}不过还是有些人透露说,也许发生过一些不太正常的事——同时小心地把是什么事、是什么人留在最大的含糊之中……

我也收到过坦率热烈的回声,来自十五六位新老朋友。有些人表达了感动,并不想掩饰它、也不想让它噤声。这些回声,以及另一些来自数学界之外、同样热烈的回声,将是我对一项漫长而孤独的工作的报偿——那工作不仅为我自己,也是为所有人做的。

在收到我这封信的一百三十来位同行中,有三个人真正在完整意义上作了回应:他们把自己放了进去,而不是限于对时下事件作一番远距离的评论。我还从一位非数学界的通信者那里收到一个这样的回声。这些才是对我的消息的真正答复。这也正是我最好的报偿。

\section{(14) 三只脚踩一个盘子}

我的好几位同行和数学家朋友曾表示希望,《收获与播种》能在数学界引发一场广泛的辩论:关于这个圈子里的风气状况,关于数学家的伦理,关于他工作的意义与目的。目前,至少可以说,事情并没有往那个方向走。从现在起(说一句应景的双关语),关于一场“葬礼”的辩论,看样子已经径自被一场辩论的“葬礼”取代了!

不管人们愿不愿意,尽管大多数人的沉默与冷漠,一场辩论却确确实实已经开了场。它恐怕永远不会具有真正公开辩论那样的规模,更不会(上帝保佑别这样!)具有“官方”辩论那种排场与僵硬。无论如何,已经有很多人抢先一步、三下两下就把它了结了——在还没弄清它是什么之前,就在自己心里把它关掉了;他们仗着那条万古不易、不可动摇的共识:“在这最好的世界里,一切都再好不过”(在此即数学的世界)。不过,也许终究会有一次质疑从外面来,渐渐地从一些“证人”那里来——他们不属于同一个圈子,因而不是这个圈子的群体共识的囚徒,因而(即便在自己心里)也不觉得自己被牵连在内。

在收到的几乎全部回声里,我都看到同一种混淆,涉及两个前提性的问题:《收获与播种》所提出(至少是默契地提出)的“辩论”究竟针对什么;以及谁有资格了解它并就它表态,或者说:谁有资格在充分知情的基础上形成意见。关于这一点,我想在此标明三个“坐标点”。当然,这并不能阻止那些执意要混淆的人继续待在混淆里。至少,对于那些想知道究竟是怎么回事的人,这也许能帮他们不被四面八方来的响动(甚至包括那些最善意的响动……)分了神。

a)一些真诚的朋友向我保证“一切最终都会摆平”(而“一切”,我猜,指的是那些不幸被弄坏了的“东西”……);我只要重返舞台,“用新工作让自己站住脚”,做做演讲等等——其余的别人会做。人们会宽宏地说“对这个该死的格罗滕迪克,咱们也的确有点不公”,然后不声不响地、带着或多或少的诚意去拨正航向\fnote{48}{(*)}{我已经不止一次注意到这类微妙的迹象,表明人们已经好好记下:狮子醒了……};甚至慈父般地拍拍他的肩膀,说他是“伟大的数学家”,好让这么个毕竟还算可敬的家伙平静下来——可惜他偏偏作势要发火,要掀起不受欢迎的风浪。

正如这些朋友所暗示的,这里根本不是要“松松绳子”、也不是要让人松松绳子。就我而言,我全然不需要恭维,不需要真诚的仰慕者,也不需要“盟友”——为“我的”事业,或为任何别的事业。问题不在我身上——我过得好得很;也不在我的作品——它自己会说话,哪怕是对聋子说。如果说这场辩论也(在别的之外)涉及我这个人、我的作品,那只是因为它们是对别的东西的揭示——通过一场葬礼的现实(的确是最具揭示性的现实之一)。

如果真有“某个人”在我看来该唤起警觉、不安与紧迫之感,那绝不是我这个人物,甚至也不是我的任何一位“同葬者”。那是一个集体的存在者,既不可捉摸又十分具体,人们常常谈起它,却从来小心翼翼地不去审视它——它的名字叫“数学共同体”。

最近这几周,我终于把它看成一个有血有肉的人,而它的身体正遭受着深处的坏疽。最好的食物、最精美的菜肴,在它体内都变成了毒,使病害传播得更广、扎得更深。然而它又有一种不可遏制的贪食,要永远吃得更多,仿佛这肯定是它自欺的一种方式——关于一种它无论如何也不愿去认识的病。对它说什么都是白费——最最简单的词也失去了意义。它们不再是消息的载体,只剩下引发恐惧与拒绝的喀嗒声……

b)我的大多数同行,甚至那些心肠好的老朋友,在贸然发表意见时,都要围上一圈谨慎的条件式,诸如“如果事情真是……那确实是不能容许的”——好让他们心满意足地再躺回去睡大觉。可我原以为我说得够清楚了……

\origpage{L48}隔着七个月的距离,如今我可以说明:对于《收获与播种》里报告并评论的几乎全部事实,它们的真实性不成为任何争论的对象。我会在后面回到那少数几个例外(它们也会各自在相应的地方被标明是这样的);至于其余全部事实,在《收获与播种》初版写成之后,与几位主要当事人的仔细对证(即皮埃尔·德利涅、让–皮埃尔·塞尔与吕克·吕西耶)消除了细节上的错误,并就材料性事实本身达成了毫不含糊的一致\fnote{49}{(*)}{我很高兴向这三位表达我的感谢,感谢他们在此事上表现出的善意;并确认他们在一切涉及材料性事实的问题上完全诚实。}。

所以,辩论根本不涉及事实的真实性——那不成问题——而涉及这样一个问题:这些事实所描述的作派与态度,是否应当被视为被认可的、“正常的”,还是不是。

这里说的是这样一些作派:在我的见证里,我(也许错了……)称之为丑闻式的;如背信弃义或滥用权力,如明目张胆的不诚实,不止一次达到了不义与无耻的程度。在了解了这些事实(十五年前还不可想象的事实)之后,我还得学的一件事——相当难以想象的一件事——就是:我的同行中的大多数,甚至包括那些曾是我的学生或朋友的人,今天竟把这些作派视为正常,视为完全体面。

c)我的许多同行与老朋友维持混淆还有第二种办法。调子是这样的:“抱歉,可我们不是这一行的专家——别要求我们去了解那些(天佑如此……)从我们头顶上飞过去的事实……”

我则断言:要了解那些主要事实,根本不必是“专家”(这回轮到我抱歉了!),甚至不必会背乘法口诀表或知道毕达哥拉斯定理。

% ---------------- p.103 ----------------
\origpage{L49}连《熙德》或拉封丹《寓言》都没读过也不要紧。一个发育正常的十岁孩子,在这件事上跟最负盛名的专家一样有能力(甚至比他还强……)\fnote{50}{(**)}{当然,我写《收获与播种》并不是给十岁孩子看的;若要对他说话,我会选一种他熟悉的语言。}。

\origpage{L50}让我举一个例子来说明这一点,就从“葬礼”里随手拈来一个\fnote{51}{(***)}{这是我在1984年4月19日发现的第一场“葬礼大手术”,也正是那一天,“葬礼”这个名字强加给了我。关于此事,见同一天写的那两则按语:《一个梦的回忆——或母题的诞生》与《葬礼——或新父亲》(RS III,第51、52号)。那里也有即将谈到的这本书的完整出处。}。要了解下面这几件事并就它们作出判断,根本不必知道“Motif”这个多面而极微妙的数学概念的来龙去脉,也不必小学毕业。

1°) 1963年至1969年间,我引入了“Motif”这一概念;并围绕这一概念发展出一套“哲学”与一套“理论”,它们部分地仍停留在猜想阶段。无论对错(此处无关紧要),我把 Motif 理论看作我为我那个时代的数学带来的最深的东西。再说,如今“Motif 的瑜伽”的重要性与深度已无人质疑(在它被近乎完全沉默了十年之后——那沉默从我离开数学舞台之后随即开始)。

2°)在唯一一本(1981年出版)、基本上以 Motif 理论为专题的书里(这个由我引入的名字就出现在书名中),唯一一处能让读者疑心我这点微末之人与书中长篇大论发展的那种理论有远近关系的段落,在第261页。那段话(两行半)是在向读者解释:那里发展的理论,与一个名叫格罗滕迪克的人的理论毫无关系(该理论在那里被提到第一次也是最后一次,没有任何别的出处或说明)。

3°)有一个著名的猜想,人称“霍奇猜想”(它究竟讲什么都无所谓),它的成立将意味着那本光彩夺目的书里所谓“另一种”Motif 理论,与我将近二十年前在众目睽睽之下发展出的那种理论相同(是它的一个极特殊的特例)。

我还可以加上 4°):那本书四位署名者中最负盛名的一位曾是我的学生,他多年来那些精彩的想法,除了从我这里,还能从哪里学来?——而他在书里把它们呈现得像是刚刚才想出来似的\fnote{52}{(*)}{我并不是说这本书里没有一些想法,甚至一些漂亮的想法,出自这位作者或别的合著者。但整本书的问题框架、赋予它意义的概念语境,乃至那套微妙的 X-范畴理论(被误称为“Tannaka 式”)——那套技术构成了全书的核心——这些都是我的成果。};5°)这两件事在消息灵通的人中间是众所周知的,可是在文献里白费力气也找不到一处书面痕迹,能证明那位出色的作者也许从我口中听到过什么\fnote{53}{(*)}{不过有一个例外:1977年塞尔写的一份报告里有一行,将在相应之处谈到。};以及 6°)那个微妙的算术问题——据主要作者本人向我解释,它构成全书的核心问题(而我的名字始终未被提起)——是我在六十年代沿着“Motif 的瑜伽”的路子提出来并澄清的,而作者是通过我才知道它的;我还可以再堆上 7°、8° 等等(在相应之处我当然不会不这么做)。

前面这些,对我的用意来说已经够了,我的用意就是:要了解这类事实并就它们作出判断,根本不需要什么专门的“能力”——“事情不在那个层面上”。这里起作用的能力,除了健全的理性(原则上人人都有)之外,就是我愿意称之为体面感的东西。

这里说的那本书,如今已是数学文献中被引用最多的著作之一,而它的“主要作者”也是当代最负盛名的数学家之一。话虽如此,而且看得清楚,在这桩事情里,如今在我看来最最值得注意的是:在这本书数不清的读者当中,包括那些第一手知道究竟是怎么回事的人,包括那些曾是我的学生、我的朋友的人——竟没有一个人看出其中有什么不正常。无论如何,直到我写这几行的今天,还没有一个人来向我表示,对这本光彩夺目的书有丝毫保留\fnote{54}{(**)}{总共只有两位同行(包括左格曼·梅布胡特)向我表示过这样的“保留”。这两位都不能算是这本书的“读者”。他们是出于好奇翻了翻,想弄弄清楚……}。

% ---------------- p.104 ----------------
\origpage{L51}至于我的同行和老朋友中那些从未把这书拿在手里、却以此为由申辩自己无能的人,我要对他们说:要当“专家”,根本不必——只要到随便哪一家数学图书馆去把那卷书借来,翻一翻,亲眼看看那无人争议的东西就行了……

\section{(15) 坏疽——或时代精神(1)}

\origpage{L52}这场“Motif 运算”只是同一水准的四场“大运算”之一,此外还有一大批规模较小、精神相同的。它远不是用来充实我那份“一个时代的风俗画”的集体蒙骗中“最大”的一桩,更不是最独特的一桩。它所做的,不过是趁富人不在(或已故……)之机,只掠夺了他的羊群;而不是(在普遍的漠然中)为了取乐,当着他的面去掐死穷人的母羊。甚至直到数学语言里如今已进入日常用法的层面,那些看来无害的书名、概念名、命题名,时时被人引用,它们本身就已经是一种蒙骗或欺诈\fnote{55}{}{我这里想到的尤其是那个古怪的缩写“SGA 4 ½”(分数还是有用处的!),它一个就抵得上双重欺诈(而且是当代数学文献中被引用最多的缩写之一);还有“维尔迪耶对偶”或“维尔迪耶对偶对象”、“德利涅–格罗滕迪克猜想”之类的名称,最后还有“Tannaka 式范畴”(这一回 Tannaka 本人并无过错,因为从来没人问过他……)。这些将在相应之处更详细地谈到。},并以各自的方式见证一个时代的不体面。

如果我自认为曾为“数学共同体”做过什么有用的事,那就是把若干不光彩的事实带到了大白天光之下——那些事实原本在暗处发酵。这肯定是人人都天天、或差不多天天、远远近近接触到的那类事实。他们之中有多少人肯费点闲工夫停下来——哪怕只是一瞬——去闻一闻空气,看一看?

那个自己曾撞上一些人的傲慢、另一些人(或同一批人)的不诚实的人,也许还自以为那是一种专门分配给他的特别霉运。把自己的经验与我的见证相对照,也许他会感到,这“霉运”也是他给一种时代精神起的名字,那种时代精神压着他,一如压着所有人。而且(谁知道呢!)也许这会促使他投入一场辩论——那场辩论与他有关,正如与我有关。

可是,如果我“拿到公共广场上摊开”的这堆“脏衣服”,除了引来一些人没有笑意的冷笑、另一些人客气的尴尬,在众人的漠然中别无反响,那么一个本来浑浊的局面就会变得非常清楚。(至少对那个还愿意用自己眼睛的人是这样。)真诚与体面的传统共识\fnote{56}{}{我说这些“出于真诚与体面的共识”,并不是说它们从未被违反过。但即便在被违反的时候,那也确实是“违反”,而共识本身仍然是被接受的。},在数学家之间的关系里、在数学家与其技艺的关系里,从此都将成为过去之物,“过时”之物。无需有哪个国际数学家组织再隆重宣告,这从此也将是既成之事、几乎是官方之事:如今,一切手段都允许,不再有任何保留与限制——对于数学界里掌权的那伙“靠内部推举结成的兄弟会”而言。一切观念上的捣鬼
%
% ---------------- p.105 ----------------
\origpage{L53}把戏,只为把那个只等着相信的漠然读者牵着鼻子走;一切父权名分的黑市交易;同伙之间的假引注,以及对注定沉默者的沉默;各种拉帮结派与形形色色的伪造,乃至在众目睽睽之下最粗劣的剽窃——是的,对这一切都说“阿门”,并且带着数学公共广场上所有“大名头”、所有大大小小老板的祝福,无论是用言语还是用沉默(若不用积极而殷勤的参与的话)。是的,对那里风靡一时的“新风格”说“阿门”!那曾经是一门艺术的东西,如今经(几乎)一致的赞同,变成了一个搅浑水、抢地盘的集市,在老板们慈父般的目光下。

曾经有一个时候,在数学家的世界里,权力的行使受到一致的、不可触犯的共识的限制,那是集体体面感的表达。这些共识与这种感情,如今都成了过时之物、被超越之物,肯定配不上计算机、太空舱与中子弹的光荣时代。

从此这将是既成而封定之事:权力,对于那伙掌握它的人而言,是一种任意专断的权力。

\section{(16) 公开认错——或时代精神(2)}

在《信》里,我想我已充分说明了我写《收获与播种》时所怀的精神,足以让人清楚:我绝不声称在那里做史家的工作。那是一份真诚的见证,关乎第一手的亲身经历,以及对此经历的思考。见证与思考都供所有人取用,包括史家——他可以把它当作众多材料之一来用。那时,把这份材料交付批判性分析、依其技艺的严格准则来做,便是他的事。

当然,应当区分狭义上的事实(“生事实”或“材料性事实”)与对这些事实的“评估”或“解释”——后者赋予它们意义,而这意义对一位观察者(或共同行动者)和对另一个人并不相同。大体上可以说,《收获与播种》的“见证”方面关涉事实,而它的“思考”方面关涉对事实的解释,也就是我为赋予它们意义所做的工作。在构成见证的“事实”里,我也归入“心理事实”,尤其是我的见证所映照的各种情感、联想与意象,无论它们发生在或远或近的过去,还是就在写作的当下。

对于我在《收获与播种》里描述或提及的事实,我区分三类来源。有些事实是记忆还给我的,其清晰或模糊程度因场合而异,有时还被变形。关于这些,我可以为我在写作之时的真诚立意担保,却绝不为全无错误担保。相反,我曾有机会指出其中若干细节错误,并在相应之处用后来补写的脚注标明。另一方面,有书面文献,尤其是书信,特别是正式发表的科学出版物,我有时以所期望的全部精确性引到它们。最后,还有第三方的证词。有时它补充我自己的记忆,让我重新唤起、加以精确化,有时加以改正。在某些少见的场合(我马上会回过来谈),这证词给我带来了与我已知信息全然不同的新信息。当我转述这样一份证词时,这并不意味着我有能力自始至终核查它的准确与可靠,而只是说:它相当可信地嵌入了我第一手所知的那张丰富的事实之网,足以让我确信(无论对错……)这份证词在本质上确实合乎真相。

% ---------------- p.106 ----------------
\origpage{L54}对于一个细心的读者,我想在任何时候都不会有困难去“分清”事实的记述与对事实的解释,并且(在前一种情况下)辨明我方才描述的三类来源中,是哪一类在起作用。

\origpage{L55}\asterisms

方才我提到一份第三方证词,我转述了它,却未能“自始至终核查其可靠性”——那就是左格曼·梅布胡特的证词,关于围绕他著作的那场大规模掉包手术。在《收获与播种》里提到的“材料性事实”中,如今尚可争议、或按我现在的判断需要改正的,只有某些仅凭梅布胡特一人证词而立的事实。为结束这篇后记,我要在此对试印本《收获与播种》中呈现的“梅布胡特事件版本”提出一些批评性的评论。更详细的评论与更正将各自在相应之处收入印刷版(即《收获与播种》的定本)。

我曾想充当其转述者的那个“梅布胡特版本”,在我看来本质上由以下两个论点构成:

1. 1972年至1979年间,梅布胡特是孤身一人\fnote{57}{}{1975年柏原的可构造性定理是例外,它在理论中的重要性毫无争议。但按梅布胡特的说法,那该是柏原对这一正在诞生的理论的唯一贡献。这一(不准确的)说法,曾由柏原其他论文的缺席所佐证——在那些论文里,他本该至少暗示某些主要想法。},在普遍的漠然中,从我的著作汲取灵感,发展出“$D$-模哲学”,把它作为我意义上的“上同调系数”的新理论。

2. 无论在法国还是在国际上,都有过一种一致的共识,要把他的名字与角色从这一新理论中抹掉——一旦它的意义开始被承认。

这个版本有很强的文献支持:一方面是梅布胡特那些完全令人信服的出版物;另一方面是其他作者许多出版物的署名(尤其是1981年6月吕米尼讨论会论文集),其中掉包的蓄意用心不容任何怀疑。最后,梅布胡特另外提供给我的那些更详细的细节(我在“葬礼(3)——或四种运算”那部分里转述了它们),虽然无法直接核实,却与某种整体的氛围完全吻合,而那氛围的真实性在我看来已不容任何怀疑。

我刚得知若干新的事实\fnote{58}{}{我感谢皮埃尔·沙皮拉与克里斯蒂安·乌泽尔,他们愿意提醒我注意这些事实,以及我对梅布胡特–柏原之争的陈述中那种偏颇之处。},它们表明对上面第 1° 点须作大大的修正。梅布胡特所处的孤立\fnote{59}{}{这种孤立首先来自我过去的学生对梅布胡特的想法与工作的漠然——他固执地装作在追随一位被一致共识判给遗忘的“先祖”……}确实是真的,但那是一种相对的孤立。在法国,有 J.-P. 拉米在同一个题目上的工作(梅布胡特对我只字未提),而尤其是:看来梅布胡特所发展并完成到底、且被他归为自己名下的某些重要想法,可能应归于柏原\fnote{60}{}{这些想法中最重要的是所谓“黎曼–希尔伯特”对应(用新风格的术语来说)在 $D$-模上的版本。相关猜想已由梅布胡特证明,并且(据沙皮拉向我肯定)也由柏原证明(而梅布胡特向我保证他的证明是唯一发表的)。关于证明优先权的问题,对我仍是一片迷雾,我放弃把余生用来把它弄清楚……\par 至于以 $\nabla^{\infty}$-模表述的姊妹命题,看来毫无疑问:想法的归属与证明都属于梅布胡特。}。这一来,就使我(过于)忠实地转述的那个梅布胡特版本里关于柏原–梅布胡特之争的某些情节,显得不可信或可疑了。

% ---------------- p.107 ----------------
\origpage{L56}毫无疑问,无论就“档案工作”的层面而言,还是就他设法完成到底的某些想法的构想而言,梅布胡特都是 $D$-模新理论的主要拓荒者之一,也许甚至是首要的拓荒者;无论如何,他是唯一一个全身心投入这项任务的人,而这项任务的真正意义当时还没向他显现,正如它也没向所有人显现。也确实,围绕这部著作发生的那场掉包手术——以吕米尼讨论会为顶点——在我看来仍是数学界这个世纪的大不体面之一。但若声称(像我出于真诚所做的那样)梅布胡特是独自担此任务的人,那就错了。相反,他是唯一一个有诚实与勇气明确指出我的想法与著作在他的工作中、在新理论的萌发中所占分量的人。

这篇后记不是细谈此事的地方——我会在相应之处细谈,包括一些有助于说明“梅布胡特版本”心理语境的评论。如果说“梅布胡特–柏原的纠纷”对我有某种意义,那仅仅在于它照亮了一个时代的整体氛围。而对我来说,甚至它那些变形、甚至通过那些促其浮现而起作用的力量,“梅布胡特版本”也显得是我为“一个时代的卷宗”提供的众多不那么可争议的材料之中一份有力的“时代征兆”。

我还须为我的轻率作公开认错:我呈现梅布胡特–柏原之争时,只用了一幅仅仅依据梅布胡特所提供的证词与文献的图画,而且仿佛那个版本不容任何怀疑。那个版本把一位第三方置于可笑乃至可憎的光线下,这就更有理由谨慎行事。为我的轻率、为这种缺乏健全的谨慎,我很愿意在此向柏原先生表示我最诚挚的歉意。

% ================================================================
%  第 4 章 引言 (原书 p.109-122, 边码 P.I-P.XXII)
%  源文: _work/src1/task-ch04.md
% ================================================================
\chapter{引言}

% ---------------- p.109 ----------------
\section{(I) 五叶草}

\subsection{1. 梦与成就}

\origpage{I} 到今年七月,就是三年前了,我做了一个不大寻常的梦。我说“不大寻常”,这印象只是事后才有的,是醒来回想时才有的。梦本身来到我这里,却像世上最自然、最显然不过的事,不敲锣不打鼓——以至于醒来时我差点没留意它,随手把它扔进遗忘的角落,好去办“当天的正事”。从前一天起,我就已经动身去做一场关于我与数学之关系的反思了。这是我有生以来第一次肯费这个力气去看一看——而且,我之所以在那个时刻动手,实在是因为我几乎是被逼无奈、硬被赶上路的!在此前的那些月和那些年里,发生过一些如此奇特——若不是说暴烈——的事情,那是些数学激情的爆发,不打招呼就闯进我的生活,以至于实在再也没法继续不去看看到底发生了什么。

我说的这个梦没有任何情节,也没有任何动作。它只有一幅图像,静止不动,同时又极其鲜活。那是一个人的头,侧面像。看得出它正从右向左看。那是一个成年男子,无须,一头乱发在头的周围形成一圈力量的光轮。从这颗头上透出来的,尤其是那种年轻、欢快的力量感,仿佛从那柔韧有力的颈项弧线里涌出(那弧线是猜出来的,而不是看出来的)。脸上的表情,与其说是个成年人,或者一个已经沉稳下来的人(不管成熟与否),不如说是个顽皮的小淘气,正为某件他刚刚得手、或正在盘算着要干的勾当而洋洋得意。从那里透出来的,尤其是那种强烈、含而不露、迸发为嬉戏的生的欢乐……

% ---------------- p.110 ----------------
\origpage{II} 并没有第二个人在场,没有一个“我”在观看或凝视这另一个——我们只看得见他的头。但有的是对这头颅、对它透出的东西的一种强烈感知。也没有任何人去感受印象、去评论它、给它命名,或者给被感知的那个人贴上一个名字,称之为“某某”。有的只是这极其鲜活的东西,这颗男人的头,以及对这东西同样鲜活而强烈的感知。

至于醒来时,我并不是有意地,就记起了过去这一夜的梦;这个男人的头的影像,并没有以特别的强度从那堆梦里凸显出来,它并没有挤到前面来,冲我喊、向我低语:你该看的正是我!当这个梦出现在我对一夜之梦的匆匆一瞥里,在被窝里温暖的静谧中,我自然有了醒着的心智那个习惯性的反应:给看见的东西安上一个名字。何况我根本不必去找:只要我一问,我立刻就知道了,梦里那个在那儿出现的男人的头,不是别的,正是我自己的头。

这一手不赖,我当时想,到底还是得服气:在梦里这样看见自己,就像那是另一个人!这个梦的到来,有点像我在散步时,纯属最偶然的巧遇,撞上了一株四叶草,甚至五叶草,于是我按规矩惊异上片刻,然后继续赶路,就像什么也没发生过。

至少,事情差点就这样过去了。幸而——在这种情形下我碰过许多次——我还是为了问心无愧,把这个“不赖”的小插曲白纸黑字记了下来,并开始了一场反思,它本该顺着前一天的思路继续下去。可是,一针引一线,那一天的反思最终只是让我沉浸于这个不起眼的梦、这幅唯一的图像的含义,以及它带给我的关于我自己的那条消息。

这里不是细说我这一日的冥想教了我什么、带给我什么的地方。或者更确切地说,不是细说这个梦教了我什么、带给我什么的地方——那是在我已经进入一种专注与倾听的状态、从而能够接纳它要对我说的话之后。梦与这份倾听的第一枚直接的果实,是一股骤然涌来的新能量。正是这股能量支撑了随后几个月里那场长时间的冥想,与之相抗的是内心里顽强的阻力,我得靠耐心而执拗的功夫把它们一块一块拆掉。

我开始留意自己某些梦,已有五年;而这是第一个不以“信使之梦”如今已可辨认的那副面目出现的信使之梦——那种梦有惊人的场景手段,有异常强烈、有时令人翻肠倒肚的视象。这一个却“酷”到了家,没有任何强迫你注意的东西,简直是悄无声息——取亦可,舍亦可,毫无纠纷。

\origpage{III} 几周之前,倒来过一个老派风格的信使之梦,调门是戏剧性的,甚至是野性的;它突然而即刻地终结了一段漫长的数学狂热。两个梦之间唯一表面的亲缘关系是:两者之中都没有观察者。这个梦以一则简短有力的寓言,指出当时我生活里正在发生的一件事,而我却不肯费心去留意它——说实话,我还小心翼翼地不去理会它。正是这个梦让我当时明白了有一件反思的工作刻不容缓;几周后我便投入其中,它随后持续了将近六个月。\origpage{IV}在这部反思—见证《收获与播种》的最后一部分里,我多少有机会谈到它;这部作品开启本卷,并给了本卷它的名字\fnote{1}{}{尤其参见第43节《扫兴的老板——或压力锅》。}。

% ---------------- p.111 ----------------
\origpage{V} 我之所以用这另一个梦、用这幅我自己的影像—视象(我在德语笔记里称之为“Traumgesicht meiner selbst”)来开篇这则引言,是因为最近这几周里,这个梦的念头不止一次回到我这里,那时,那场“关于一位数学家的往昔”的冥想正走向它的终点。说实话,回头看,这个梦以来的三年在我眼里像是澄清与成熟的年月,走向对它那条单纯而清澈的消息的实现。梦把“我本来的样子”显示给我。同样清楚的是,在我醒着的生活里,我并不完全是梦所显示的那个人——来自久远的重负与僵硬(过去和现在都)常常妨碍我完全而单纯地做我自己。这些年里,尽管这个梦的念头只是偶尔回到我这里,它一定以某种方式起了作用。它绝不是某种范本或理想,让我努力去与之相像,而是对一种欢乐的单纯的悄声提醒——那单纯“就是我”,它以许多方式显现,并被召唤着去摆脱那些继续压着它的东西,去充分地绽放。这个梦是一条既纤细又有力的纽带,一头连着仍被来自过去的许多重负压着的现在,一头连着这现在已孕育在胚芽里的、近在眼前的“明天”——一个从此刻起就是我、而且肯定一直就在我之中的“明天”……

肯定,如果说最近这几周里这个难得提起的梦又一次鲜明地在场,那是因为在某个层面上——那不是探测与分析的思想的层面——我一定已经“知道”:我正在做、并正引向终点的这件工作,这件接续并深化了三年前那件工作的工作,是向着实现它带给我的、关于我自己的那条消息迈出的又一步。

如今在我看来,这就是《收获与播种》的主要意义所在,这件将近两个月的高强度工作的主要意义所在。直到它完成了,我才意识到我做这件事有多么重要。在这件工作的过程中,我经历过许多欢乐的时刻,那往往是一种促狭的、爱开玩笑的、兴高采烈的欢乐。也有过悲伤的时刻,有过一些时刻,我重新经历这些年里痛苦地击中过我的挫折与伤心——但没有过哪怕一刻的苦涩。我离开这件工作时,带着那种知道自己把一件工作做到了底的人的完全的满足。没有任何一件事,无论它多么“小”,是我在那里回避掉的,或者是我心里看重、本该说出而竟没有说出的,以至于此刻会在我心里留下一点不满、一点遗憾的残渣——无论那多么“小”。

在写这份见证时,我很清楚它不会讨所有人喜欢。甚至很可能,我找到了办法把所有人,一个不落,都得罪了。然而这根本不是我的本意,连得罪任何人也不是。我的本意不过是去看那些简单而重要的事,那些日常的事——我作为数学家的过去(有时也包括现在)里的事——好最终(晚知道总比不知道好!)毫无疑义、毫无保留地发现它们过去是什么、现在是什么;并且一路走一路用简单的话说出我所看见的。

\subsection{2. 一场旅行的精神}

这场反思最终成了《收获与播种》,它起初是作为《追寻栈》(À la Poursuite des Champs)第一卷(正在完成中)的“引言”开始的——《追寻栈》是我自1970年以来就打算发表的第一件数学工作。头几页是我去年六月在一个空闲的空当里写的;不到两个月前,我又从搁下的地方把这反思接了起来。我意识到有不少东西值得看、值得说,所以指望的是一篇比较充实的引言,三四十页。可是此后将近两个月里——直到此刻,我正在为这最初本是引言的文字写一篇新引言——我每天都以为今天就
% ---------------- p.112 ----------------
会把它做完,至不济明天或后天也就完了。等到几周之后我开始逼近一百页这个关口,这引言就被擢升为“导论性的一章”。又过了几周,当所谓“这一章”的规模远远超出正在准备的这卷里其他各章(写这几行时,除最后一章外都已写完),我终于明白了:它的位置不该在一本数学书里,这场反思与见证在那里实在太挤。它真正的位子是在单独的一卷里,那将是“数学反思”系列的第1卷——这个系列我打算今后几年顺着《追寻栈》的势头继续下去\fnote{2}{}{(5月28日)这里指的是《收获与播种》第一部分《愚妄与更新》的文本。写这几行时,第二部分尚未写出。}。

我不会说《收获与播种》——这部《数学反思》系列的第一卷(接下来先是《追寻栈》的两三卷)——是《反思》的“导论”卷。不如说,我把这第一卷看作将来之物的地基,或者更确切地说,看作定下基调的那一卷,给出我踏上这趟新旅程时所怀的那种精神的那一卷;这旅程我打算在今后几年里继续走下去,而它把我带向哪里,我说不上来。

关于本卷主干部分的这些说明,最后再交代几点实用性质的情况。读者在《收获与播种》的正文里偶见“本卷”的提法,不必奇怪——那暗指的是《追寻栈》的第一卷(《模型史》),我当时还以为自己仍在为它写引言。我不愿去“改正”这些地方,因为我最看重的是保住文本的自发性,保住它作为见证的真实性——不仅是对一个遥远过去的见证,也是对正在写作的当下这一刻的见证。

同样出于这个理由,我对文本初稿所作的修改,只限于纠正文笔上的笨拙,或某些有时含混、妨碍理解我要表达之意的说法。这些修改有时把我引向比写初稿时更清晰、更精细的领会。至于对初稿稍具实质性的改动——为给它添些层次、说得更精确、作些补充,或者(有时)加以更正——则放在大约五十条编号的“注”里,集中排在反思的末尾,它们占了全文四分之一还多。我用(1)之类的符号去指这些注……在这些注里,我挑出二十来条,觉得它们(以篇幅或实质而论)重要得可与这篇反思自发组织成的五十个“节”或“段”中的任何一节相比。这些较长的注也收进了目录,排在五十节的清单之后。不出所料,对某些长注来说,又需要给注再加一条或几条注。这些注便接在该注之后,用的是同一类指号;只有相当短的注例外,它们就排在同一页上作“脚注”,用的是诸如(46\textsubscript{1})、(46\textsubscript{2})那样的指号。

给正文的每一节、以及每一条最充实的注取一个名字,我从中得到很大的乐趣——何况后来这甚至成了我理清头绪的必需品。不用说,这些名字都是事后才找到的;当初开始写某一节或某条稍长的注时,我一条也说不出它的实质是什么。我用来称呼八个部分(I至VIII)的那些名字(如“劳动与发现”等等)就更是如此——这八部分是我事后把构成全文的五十节归并而成的。

至于这八个部分的内容,我只作极简短的评论。头两个部分——I(《劳动与发现》)与II(《梦与做梦者》)——含有一些关于数学工作、以及一般而论关于发现工作的反思要素。
% ---------------- p.113 ----------------
\origpage{VII}在其中,我本人牵涉得零散得多,也间接得多,远不如后面几个部分。后面这几个部分才尤其具有见证与冥想的品格。\origpage{VIII}III至VI部分,主要是关于我1948至1970年间“在数学世界里”那段数学家往昔的反思与见证。推动这场冥想的动机,首先是想理解这段过去,以此努力去理解并承担一个在某些方面有时令人失望或令人茫然的现在。VII(《孩子嬉戏》)与VIII(《孤独的历险》)两部分,更多涉及我从1970年至今、也就是从我离开“数学家的世界”不再回头以来,我与数学之关系的演变。那些动机、那些力量与境况,它们使我(连我自己都吃惊)在中断十三年多之后,重新开始一种“公开的”数学活动(即写作并发表《数学反思》)。

% ================================================================
%  第 4 章 引言(后半) — 原书 p.114-122, 边码 P.IX-P.XXII
%  源文: _work/src1/task-ch04b.md
%  承接 chapters/part1-ch04.tex(该文件止于 4.1.2 之末)
% ================================================================

% ---- [补译] 4.1.3 节余下部分(p.113; part1-ch04.tex 中缺此节正文) ----
\subsection{3. 罗盘与行囊}

我还该就与《收获与播种》一起构成本卷、并与它同名的另外两篇文字说几句话。

《纲领概要》勾画出我近十年来所追寻的数学思考的主要主题。我打算今后几年里,至少把其中若干稍加展开,写成一连串不拘形式的反思——我已有机会提起过它们,即“数学反思”。这份概要是去年一月我写的一份报告的逐字复制,当时是为了支持我申请国家科研中心(CNRS)的一个研究员职位。我把它收进本卷,是因为这个纲领显然远远超出我卑微个人之力所能及——哪怕上天让我再活一百年,并且我选择把这一百年用来尽可能远地推进这些主题。

《主题概要》写于1972年,是为另一次申请(法兰西学院的一个教授职位)。它按主题勾画了我当时视为自己主要数学贡献的东西。这篇文字带着它写作时心境的印记——那阵子我对数学的兴趣,说得客气些,也只能算是极边缘的。所以这份概要顶多不过是一份干巴巴、有条有理的罗列(幸好它并不求详尽……)。它看不出受到任何眼光的照耀,也看不出被任何欲望的气息吹动——仿佛我在其中像交差一样逐一过目的那些东西(而这确实就是我当时的心境),从来不曾被一种活的眼光触碰过,也从来不曾有过那种激情:要把它们带到光天化日之下,而它们当时还只是躲在雾与影的帷幕背后被隐约预感着……

然而,我之所以还是决定把这份恐怕没什么灵感的报告收在这里,主要是想堵住(假如这事办得到的话)某些高来高去的同行和某种风气的嘴:自从我离开那个我们曾共同所属的世界以来,他们总爱居高临下地看他们亲切地称之为“格罗滕迪克货色”(grothendieckeries)的东西。据说,这词的意思是死抠一些太琐碎的小事,以致一位严肃而有品味的数学家不肯为它们浪费自己当然宝贵的时间。也许这份“难以消化的文摘”在他们眼里会显得更“严肃”一些!至于我笔下那些由一种眼光、一种激情所驱动的文字,它们并不是写给那些人的——某种风气把他们维持、并且为他们开脱在一种自得自满之中,使他们对那些令我着迷的东西毫无感觉。如果说我写东西不只是为自己,那是为这样的人:他们不觉得自己的时间和自己本人珍贵到了不能无休无止地去追寻那些谁都不屑一顾的显而易见之物,不能为每一件被发现之物的内在之美而欢欣,并且把这美与我们所知的任何别的、各有其自身之美的东西区分开来。

% ---------------- p.114 ----------------
\origpage{IX} 如果我愿意把构成本卷的三篇文字彼此定位,并定出各自在这趟旅行中——我随《数学反思》一同登上了这趟旅行——所扮演的角色,我就可以说:反思—见证《收获与播种》反映并描述了我踏上这趟旅行时所怀的精神,正是这精神给了它意义。《纲领概要》描述了我的灵感之源,它们为这趟驶入未知的旅行定出一个方向——当然不是说定下一个目的地——有点像一只罗盘,或者一根结实的阿里阿德涅线。最后,《主题概要》匆匆过了一遍一份行囊,那是我1970年之前作为数学家的往昔里积攒下来的;其中至少有一部分,在旅行的这个或那个阶段会有用、会受欢迎(就像我的上同调代数与 topos 代数那套本能反应,如今在《追寻栈》里已是我不可或缺的东西)。而这三篇文字相随的次序,连同它们各自的篇幅,恰好(并非我刻意安排)反映出我在这一旅程中给予它们的分量与重视——这趟旅程的第一程,已近尾声。

\subsection{4. 一场追寻显而易见之物的旅行}

我还该就这趟一年多一点以前启程的旅行——即“数学反思”——再说几句更周详的话。关于我怀着怎样的精神踏上这趟旅行,我在《收获与播种》开头八节(即这场反思的第I、II部分)里已说得相当详细;我想,这种精神从现在起就在本卷里,也在接续它的那一卷(即《模型史》,它是《追寻栈》的第1卷)里显露出来了——后者正在完成中。因此,我觉得在这篇引言里再就此发挥,已是多余。

我当然无法预言这趟已经启程的旅行会是什么样子;这事我会随着它一路走下去而一点点发现。眼下我并没有一个哪怕只及大要的预定路线,而且我怀疑近期内会冒出这么一条来。正如我前面说过的,那些大概会启发我思考的主要主题,都多少勾画在《纲领概要》里了——那份“罗盘文本”。在这些主题当中,还有《追寻栈》的那个主要主题,即“栈”;我希望就在今年之内把它(并且就此打住)走完一遍,写成两卷,也许三卷。关于这个主题,我在《纲领概要》里写道:“……这有点像一笔债,我要向一段科学的过去偿还:在那十五年里(1955至1970年间),上同调工具的发展始终是我在代数几何基础方面工作中不变的主旋律”。因此,在预定要写的主题当中,这一个最深地扎根于我的科学“过去”。也正是这一个,在逝去的这十五年里始终作为一种遗憾在场——它也许是我离开数学舞台时留下的未竟工作中最显眼的一处空缺,而我旧日的学生或朋友里,没有一个人费心去填补它。关于这项进行中的工作,感兴趣的读者可以去查阅《纲领概要》中相应的那一节,或者去查阅《追寻栈》第一卷(正在完成中)的引言(这一回是真正的引言!)

作为我的科学过去留给我的另一份遗产,尤其还有那个我格外放在心上的 Motif 观念:它至今仍等着走出那一直被囚于其中的黑夜——尽管它现身已有足足十五年了。如果没有任何比我位置更好的人(更年轻,而且拥有他所掌握的工具与知识)决定在最近这几年里动手,那么最终由我自己来干这里所必需的基础工作,也并非不可能。

我趁此机会说明:Motif 观念的命运(或者不如说,厄运……),以及我发掘出来的另外几个观念——它们在我看来(潜在地)是其中最多产的——的命运,构成了一篇将近二十页的回顾性反思,是《收获与播种》诸“注”中
%
% ---------------- p.115 ----------------
\origpage{X} 最长的一篇(也是最后几篇之一)\fnote{3}{}{这条双注(第46、47号)及其子注已收入《收获与播种》第二部分“葬礼”,后者是它的直接续篇。}。事后我把这条注分成两部分(《我的孤儿们》与《拒绝一份遗产——或一个矛盾的代价》),此外还有紧随其后的三条“子注”\fnote{4}{}{这里指的是第48、49、50号子注(第48$'$号注是后来补加的)。}。这连续五条注合起来,是《收获与播种》里唯一不作泛泛一提、而是真正谈论数学概念的部分。这些概念在这里成了机会,用来例示数学家世界内部的某些矛盾——而这些矛盾本身又映照出人本身之中的矛盾。有一个时候我曾想把这株触手蔓生的注从它所属的文本里分出去,归入《主题概要》。那样做的好处,是能给后者一个透视,并给一篇实在太像目录的文字注入一点生气。然而我还是没有这样做,为的是保住一份见证的真实性——而这条巨型注,不管我乐意不乐意,确实就是这份见证的一部分。

\origpage{XI} 关于《收获与播种》里所说我以何种心境着手“反思”,我想在这里只补一句话——对此我已在其中一条注(《年轻人的势利——或纯洁的捍卫者们》)里表过态,那里我写道:“我毕生作为数学家的抱负,或者不如说我的欢乐与激情,始终是去发现那些显而易见之物;在眼下这部著作里,这也依然是我唯一的抱负”(《追寻栈》)。对于我一年来随“反思”进行的这趟新旅行,这也同样是我唯一的抱负。开启这趟旅行的这部《收获与播种》也不例外(至少对我的读者而言是如此,如果真有读者的话)。

\subsection{5. 一笔情愿的债}

我想以几句关于本卷《收获与播种》两则题献的话,来结束这篇引言。

那则题献——“献给曾是我的学生们的人,我曾把自己最好的给了他们——也给了最坏的”——至少去年夏天就已在我心里,尤其是在我写开头四节的时候——那四节当时还该是一本数学书的引言。这说明我很清楚(其实已经清楚好几年了):有一个“最坏的”东西需要审视——而现在正是时候,再没有比现在更好的时候了!(不过我万万没想到,这个“最坏的”最终会领着我走完一场将近两百页的冥想。)

相反,“献给曾是我的长辈们的人”这则题献,是半路上才出现的,就像这场反思的名字本身(它后来也成了一卷的名字)一样。这则题献向我揭示出他们在我作为数学家的生命中所扮演的重要角色——这角色的效力至今仍然活着。这一点,在后面的篇幅里大概会相当清楚地显现出来——所以在这里不必就此多说。这些“长辈”,按他们在我二十岁那年(大致)进入我生命的次序,是亨利·嘉当、克劳德·谢瓦莱、安德烈·韦伊、让-皮埃尔·塞尔、洛朗·施瓦茨、让·迪厄多内、罗歇·戈德芒、让·德尔萨特。当时那个初来乍到的我,受到了他们每一个人的善意接待;此后他们中许多人给了我持久的友谊与情谊。我还须在此提到让·勒雷:在我与“数学家世界”的第一次接触时(1948/49年),他的善意接待同样是一种珍贵的鼓励。我的反思显现出一笔感恩之债,欠他们每一个人——这些“来自另一个世界、另有另一份命运”的人。这笔债绝不是负担。发现它时,它像一份欢乐来临,让我变得更轻。

\noindent 1984年3月底\par

% ---------------- p.116 ----------------
\section{(II) 一桩敬重之举}

(5月4日—……6月)

\subsection{6. 葬礼}

\origpage{XII} 一件意料之外的事,把一场本已走到头的反思又重新发动起来。它开启了一连串大大小小的发现,在过去这几周里逐步揭开了一个一直模糊不清的局面,并使它的轮廓愈发分明。这尤其使我周详而深入地走进了一些事件与处境——此前它们只是被顺带一提,或隐约暗示过。于是,前面提到过的那篇关于一部作品之沉浮的“十五页回顾性反思”(引言,第4节),就获得了意想不到的规模,多出了大约两百页。

由于事势所迫,也由于一场反思的内在逻辑,我在路上被引到把别人牵涉进来,其程度不亚于牵涉我自己。被牵涉得比谁都深的(除了我自己),是一个与我有着将近二十年友谊的人。关于他,我曾写道(那是一种委婉说法\fnote{5}{}{关于这个“委婉说法”的含义,见注《与众不同者》,第57$'$号。}):在这份扎根于共同激情的亲切友谊的最初几年里,他“多少有点像个学生”;而在很长一段时间里,在我内心深处,我把他看作某种“合法继承人”——所继承的,是我相信自己还能在数学上带来的东西,是在那部始终零碎的已发表作品之外的东西。很多人早已认出他是谁了:他就是皮埃尔·德利涅。

我并不为随这些注一起公开一份对一段私人关系的个人反思、并且就这样未经征求他同意而把他牵涉进来而道歉。一个长期隐而不显、混乱不清的局面,终于被摊到光天化日之下加以审视,在我看来是重要的,对所有人都是健康的。这样做,我是在提供一份见证——无疑是主观的,既不自称穷尽了一个微妙而复杂的局面,也不自称没有错误。它的首要长处(正如我过去那些出版物、或眼下正在写的那些的首要长处一样)是它存在着,供可能对它感兴趣的人取用。我所操心的,既不是要说服人,也不是要躲在那些所谓“确凿”的东西背后,以免于错误或怀疑。我所操心的是要真实——把我在每一刻所看见、所感到的东西就那样说出来——以此作为深入它们、理解它们的手段。

“葬礼”这个名字——用于所有那些与“一段过去的重量”相关的注的总称——在反思过程中以一种越来越强的力量自显出来\fnote{6}{}{在这场反思接近尾声时,另一个名字浮现出来,表达了某幅图景的另一个触目惊心的方面——那幅图景在过去五周里逐步在我眼前展开。那是一个童话的名字,我将在相应之处回到它:《中国皇帝的袍子》……}。在其中,我扮演那位预先的死者的角色,与几位(年轻得多的)数学界同行组成凄凉的送葬队伍;他们的作品落在我1970年“离去”之后,并且凭着某种风格、某种对数学的进路,带着我影响的印记。这些人里排在头一个的,是我的朋友左格曼·梅布胡特;他有幸——这运气沉重得很——要承受一个被当作“1970年之后的格罗滕迪克的学生”的人所遭遇的一切不利,却又没享受到与我接触、得到我鼓励与指点的那种好处,因为他只是通过我的文字、从我的作品那里“受教”的。那阵子(在他所出没的那个世界里),我已经俨然一副“死者”的模样,以致长期以来,连见上一面的念头显然都不曾出现过;一段持续的往来(无论私人的还是数学的),直到去年才终于结上。

% ---------------- p.117 ----------------
\origpage{XIII} 这并没有妨碍梅布胡特逆着一种专横的风气、逆着他那些(曾是我的学生的)长辈们的轻蔑,在近乎完全的孤立中,做出独创而深刻的工作——靠的是一种意想不到的综合:把佐藤学派的想法与我的想法合到一起。这项工作提供了一种新的入手之处,可以把握解析簇与代数簇的上同调,并预示着我们对这种上同调的理解将有一场大规模的更新。毫无疑问,假如梅布胡特在那些最适合做这件事的人那里,得到了他们当年在我这里得到过的热情接待与毫无保留的支持,那么这场更新早已是既成之事,而且已经完成多年了。至少,自1980年10月起,他的想法与工作为代数簇上同调理论的一场引人注目的重启提供了灵感与技术手段,使这门理论终于(除了德利涅围绕韦伊猜想的那些结果之外)走出了一段漫长的停滞期。

\origpage{XIV} 难以置信却又千真万确:近四年来,他的想法与结果被“所有人”使用着(与我的想法同等待遇),而他的名字仍被小心地无视、被默不作声地略过——恰恰是被那些第一手了解他的工作、并在自己的著作里根本性地使用它的人。我不知道数学界在别的什么时候可曾见过这样的失格:其中某些最有影响、最有声望的成员,在普遍的漠然中,带头蔑视数学家这行伦理中最被普遍承认的规矩。

我看见四个男人,才能出众的数学家,他们和我一起享有——也有资格享有——这场以沉默与轻蔑举行的葬礼的哀荣。而我在每个人身上都看见:轻蔑如何啃咬着那曾激励过他的美好激情。

除了这几个人,我尤其看见两个人:他们一个接一个站在数学公共广场的聚光灯下,在众人的簇拥中主持葬礼,而同时(在一种更隐蔽的意义上)他们也被埋葬——被自己的双手埋葬,与他们刻意埋葬的那些人一同被埋葬。其中一个我已经点过名了。另一个同样是从前的学生、从前的朋友:让-路易·维尔迪耶。1970年我“离去”之后,他与我之间的接触没有维持下来,除了职业层面上几次匆促的会面。所以大概正因如此,他在这次反思里只通过他职业生涯中的某些举动出现;至于这些举动可能的动机——就他与我的关系这一层面而言——则未经审视,而且我对此也完全摸不着头脑。

如果说有一个追问在逝去的这些年里始终逼着我,它曾是《收获与播种》的一个深刻动机,也在这整场反思中一路跟着我,那就是:在某种精神、某种风气的到来之中,我占了怎样一份——正是这种精神与风气,使得像我说的那样的失格成为可能,而这一切发生在一个曾属于我的世界里,在我作为数学家的二十多年里,我曾与它认同为一。这场反思让我发现:由于我身上某些虚浮的态度——它们表现为对才具平平的同事一种无言的轻蔑,表现为对我自己、对某些才具出众的数学家的自我迁就——我对这种精神并不陌生;而今天,我看见这精神在我曾爱过的那些人当中肆意铺展,也在我曾把自己热爱的一门手艺教给的那些人当中铺展;那些我没有好好爱、也没有好好教的人,如今在这个曾为我所珍爱、而我已经离开的世界里定调子(当他们不是在立法的时候)。

我感到一阵风刮起来,那是自得、玩世不恭与轻蔑的风。“它刮着,不管什么‘功绩’什么‘过失’,用它的气息烧灼卑微的志向,也烧灼最美的激情……”。我明白了:这阵风,是我参与播下的盲目而不经心的种子那丰饶的收成。而如果它的气息回吹到我身上,回吹到我曾交托给别人手里的东西上,回吹到我今天所爱的人身上,而且
%
% ---------------- p.118 ----------------
\origpage{XV} 回吹到那些敢于自称属于我、或仅仅是受我启发的人身上,那么这是事物的回还,我没有理由抱怨,而且它有许多东西要教我。

\subsection{7. 葬仪的序次}

于是,在“葬礼”这个名目下,我在目录里把那些与那个表面平淡无奇的节“一段过去的重量”(第50节)相关的主要“注”的浩浩荡荡的队列归拢在一起,从而让这个名字充分显出了它的意义——当初写《收获与播种》“初稿”的这最后一节时,这名字一上来就落到了我头上。

\origpage{XVI} 在这支亲缘关系错综的长长的注的队列里,过去四周里才加入的那些(注(51)至注(97))\fnote{8}{}{还须加上1984年5月12日的第104号注。第98号及其后的各注(上一句提到的第104号注除外)构成反思的“第三口气”,自1984年9月22日起;它们也标了日期。}显得与众不同:它们是唯一标了日期的(4月19日至5月24日)\fnote{9}{}{在同一天写的一连串连续的注里,只有第一条标日期。其余未标日期的注是第44$'$号至第50号(构成第I、II、III支队伍)。第46、47、50号注写于3月30日或31日,第44$'$、48、48$'$、49号注写于四月上半月,最后,第44$''$号注是标了日期的(5月10日)。}。我觉得最自然的做法,是按它们在反思中相继出现的时序\fnote{10}{}{有时我在这个时间顺序里作过幅度不大的调换,换成一个“所谓逻辑的”次序,只要我觉得反思进程的整体印象不会因此失真。不过,作为仅有的例外,我要指出十一条注(其编号前有一个感叹号!)——它们源自某条注之后的页底脚注,篇幅涨到了让人却步的规模,于是我把它们各自放在与之相关的那条注之后(除了第98号注,它关联的是第47号)。}把它们排出来,而不是按某种所谓的“逻辑”次序,也不是按先前各注中提到它们的先后次序。为了能找回后一种(绝非线性的)承继关系——即各条参与其中的注之间的谱系——我在目录里让每条的编号后面跟着(在它之前的那些注中)第一次提到它的那一条的编号\fnote{11}{}{当对某条注的引用(如(45))就出现在“一段过去的重量”这一节本身之中时,则把这一节的编号(50)加括号,置于该注编号之后,如“46(50)”。},若没有这种情况,则跟着它所直接续写的那一条的编号\fnote{12}{}{在目录里,若某条注是前一条注的直接续写(此时两号相连),其编号前加一个 * 号。例如“*47, 46”表示第47号注是第46号注的直接续写(而这里第46号并非紧邻在它前面的那一条,紧邻的是第46\textsubscript{9}号注)。\par 最后,我在目录里给那些后面不跟别的编号的注加了着重号,也就是说,给那些代表反思“新起点”的注——它们并不嵌在已经做过的反思的某个确定位置上。}。(后一种关系在正文里另有标示:在前一条注的末尾放一个指号,如注(46)末行末尾的($\Rightarrow$ 47),指向续写它的注(47)。)最后,某一注中某些稍具技术性的细节,则集中在该注末尾,编成子注,用接在该注原编号之后的连续下标标示——如注(46)《我的孤儿们》的子注(46\textsubscript{1})至(46\textsubscript{9})。

为了多少给“葬礼”的整体安排立起结构,也为了让人在蜂拥而至的众多注里能认得出路,我觉得在此场合宜于在队列中插入几个庄重而富于暗示的小标题,每个小标题都走在并引领一支或长或短的送葬队——由主题相同、彼此相连的连续若干注组成。

于是我有幸看见,在一支前来为我送葬的庄严长队里,十支\fnote{13}{}{(9月29日)其实最终有十二支队伍,若把灵车(X)和“死者(仍未过世)”(XI)算进去的话——后者在最后关头又溜进了队列……}队伍一支接一支地聚拢起来——有的卑微,有的气派,有的悔恨,有的暗暗欢腾,在这样场合下本也免不了如此。于是走上前来的有:身后的学生(人人都把无视他当作本分)、孤儿们(为这场合才刚被掘出)、时风及其名流(这是我活该)、诸Motif(我所有孤儿中最晚出生、也最晚被掘出的),
%
% ---------------- p.119 ----------------
\origpage{XVII} ……我的朋友皮埃尔谦逊地领着最要紧的那一支队伍,紧随其后的是“一致同意”——那些(默默地)协同一致的注——以及全体出动的(所谓)“乖戾”研讨会(它与身后的学生、即“无名学生”划清界限的方式,是夹在中间的几支捧着鲜花与花环的送葬队);最后,为了给这气派的队列一个体面的收场,走上前来的还有“学生”(绝非身后的,更绝非无名的),又名“老板”,身后跟着我那一帮忙忙叨叨的学生(带着大量铁锹与绳索),最后是灵车(车上赫然四口漂亮的橡木棺材,用螺丝牢牢拧紧,还不算那个掘墓人)……十支队伍终于齐了(总算齐了),缓缓走向葬礼仪式。

\origpage{XVIII} 仪式的高潮是悼词,由我的朋友皮埃尔亲自出马,以完美的分寸献上——他应众人之请主持葬礼,人人满意。仪式在一曲终局而决定性的“自深处”(De Profundis)歌声中结束(至少人们这样希望)——这歌由可敬的死者本人唱出,作为真诚的感恩;他瞒着所有人,从自己那场声势骇人的葬礼上活了下来,并且按说从中得了教训,心满意足——这份满足,便构成了那场值得纪念的葬礼的最后一个音符、最后一个和弦。

\subsection{8. 一个秘密的终结}

在这场(但愿是)最后的反思阶段中,我看出有必要在《数学反思》第1卷的“附录”里,除前面谈到的三篇之外,再附上另外两篇性质属于数学的文字\fnote{14}{}{此外,我想给《主题概要》(见引言第3节“罗盘与行囊”)加一份“评注”,就我对其中粗略过目的那些“主题”的贡献作些说明,也就那些在我的数学作品中主要思想力的生成中起过作用的影响作些说明。最近六周的回顾已经显出(令我自己也吃惊):在大多数这些思想的启动上,以及在我1955至1970年间给自己定下的某些“大任务”上,塞尔都起了“雷管”的作用。\par 最后,还有一篇性质属于(通常意义上的)数学的文字,也是唯一(偶然)出现在非技术性文本《收获与播种》里的——我要指出注《屠杀》(第87号)下的子注第87\textsubscript{1}号:我在那里以它应有的仔细,阐明凝聚层语境中为人熟悉的黎曼–罗赫–格罗滕迪克定理的一个“谨慎的”(猜想性)变体。这个猜想(与另外若干一起)出现在1965/66年度 SGA~5 讨论班的收尾报告里;那份报告(如同许多别的报告一样)在十一年后以 SGA~5 之名出版的文集中没有留下任何痕迹。这个关键讨论班落在我某些学生手里之后的种种遭际,以及这些遭际与某种“SGA~4½ 行动”的关联,在注第63$''$、67、67$'$、68、68$'$、84、85、85$'$、86、87、88号所进行的反思过程中逐步显露出来。\par 又及:另一条给出相当充实的数学评注的注,讨论的是为那些已知的、拥有所谓“六种运算”对偶形式的情形(在可能的范围内)抽出一个共同“topos 式”框架是否合宜——我还要指出注《赊账的论文与全险》(第81号)下的子注第81\textsubscript{2}号。}。

第一篇是一份带评注、分两部分的报告的复制,是我1968与1969年就 P.\ 德利涅的工作(其中有些至今仍未发表)写的,对应于1965/67/68三年间在法国高等科学研究所(IHÉS)的数学活动。

另一篇是一份“六种变率公式表”的草稿,汇集了我在1956至1963年间抽出的某种对偶形式(灵感来自庞加莱对偶与塞尔对偶)的共同特征;这份公式表后来表明,对于迄今遇到的一切上同调对偶情形,它都具有“普适”的品格。随着我离开数学舞台,这套形式似乎就废弃不用了,以致据我所知,还没有人(除我之外)肯费神哪怕只写出那些基本运算的清单、它们所给出的基本典范同构的清单,以及这些同构之间那些本质的相容关系的清单。

这份融贯的公式表草稿,对我将是显然的第一步,通向那幅“Motif 之梦的宏伟全景图”——十五年来,它一直“等着那位敢于挥笔把它画出来的数学家”。

% ---------------- p.120 ----------------
\origpage{XIX} 看样子,这位数学家不是别人,正是我自己。确实,该是时候了:将近二十年前诞生、在私密中交托出去的东西——它本不是为了成为一个人的特权,而是为了供所有人取用——终于要走出秘密的黑夜,再一次诞生于白昼的充足光线下。

确实,除我之外,只有一个人对这套“Motif 的瑜伽”有贴身的了解——在我离去之前的那些日子里、那些年里,他一天天、一年年从我口中把它学了去。在我有幸发现并带到光天化日之下的一切数学事物当中,这个 Motif 的实在在我眼里仍是最迷人的、最充满奥秘的——它就在“几何”与“算术”之间那深层同一性的核心处。而这套“Motif 的瑜伽”——那长久被忽视的实在把我引向的瑜伽——也许是我在作为数学家的生命头一段里抽出的最有力的发现工具。

但同样真实的是:这个实在,以及这套努力把它圈得尽可能近的“瑜伽”,从来不曾被我守为秘密。我被撰写基础这一必须完成的任务占住了(此后人人都很乐意把它们原封不动地用在自己日常的工作里),没有拿出那必不可少的几个月,为这套 Motif 的瑜伽写一份宏阔的整体草稿,从而把它交到所有人手里。然而在我出乎意料地离去之前的那些年里,我并没有放过机会——每逢相遇,就向愿意听的人讲它,首先是我的学生们;而他们(其中一人除外)都把它忘了,就像所有人都把它忘了一样。我讲它,不是为了推销什么会挂上我名字的“发明”,而是为了引人注意一个每走一步都显现出来的实在——只要你关心代数簇的上同调,尤其是关心它们那些“算术”性质,以及迄今已知的各种上同调理论彼此之间的关系。这个实在,与当年“无穷小”的实在一样可触——那东西被感知到,远早于严格语言的出现,而正是那种语言使人能完满地把握它并把它“确立”下来。而为了把握 Motif 的实在,我们今天绝不缺一套灵活而合用的语言,也不缺在数学理论建构方面成熟老练的经验——这些,我们的前人是缺的。

如果我从前在屋顶上大声喊过的东西落进了聋耳朵,如果某一个人的沉默不语所收到的回声,是所有那些装作对上同调感兴趣的人(而他们明明和我一样有眼睛有手……)的沉默与麻木,那么我也不能把责任只归到那一个选择把“好处”留给自己的人身上——那本是我为所有人而交托给他的。我们不得不承认:我们这个时代,其科学产出的狂飙足以与投入军备或消费品的劲头相比,却离我们十七世纪前人那种“大胆的活力”相去甚远——他们为了发展无穷小的计算,“不走那么多弯路”,不会因为这东西是否“猜想性”的顾虑而停住,也不等他们中某位显赫人物肯来给他们开绿灯,才去抓住每个人自己亲眼看见、第一手感到的东西。

\subsection{9. 场景与角色}

凭它自身的内部结构,也凭它那个特殊的主题,“葬礼”(如今已占《收获与播种》全文一半以上)在很大程度上、并且从逻辑观点看,独立于它前面那场漫长的反思。然而这只是一种纯粹表面的独立。对我来说,这场围绕着一场“葬礼”的反思——那葬礼渐渐从“未言”与“预感”的雾中走出来——与在它之前、它由之而生并且由之获得全部意义的那场反思,是不可分的。它起头时只是“顺路”匆匆一瞥一部作品的沉浮——那部作品我有点(很大程度上)已经看不见
% ---------------- p.121 ----------------
\origpage{XX} 了;后来,既未预谋也未寻求,它变成了一场对一段重要关系的冥想,而这又把我引向对那部作品在“曾是我的学生们”手中的命运的反思。把这场反思从它自发从中生出的那场反思里割出去,在我看来就是把它降格为一幅简单的“风俗画”(甚至是在数学“上流社会”里算的一笔账)。

\origpage{XXI} 确实,若有人偏要如此,整个《收获与播种》也能作同样的降格,变成一幅“风俗画”。当然,某个时代、某个圈子里通行的风俗,那些参与塑造该圈中人生活的风俗,自有其重要性,也值得描写。然而《收获与播种》的细心读者会明白:我的打算不是描写风俗,也就是说,不是描写某一个场景——它随时代而变、因地而变,我们的行动就在它上面展开。这场景在很大程度上界定并划出了我们身上各种力量可用的手段,使它们得以表达。场景以及它提供的手段(还有它强加于人的“游戏规则”)千变万化,而我们身上那些深层力量的本性——它们在集体层面上塑造场景,在个人层面上又在场景上表达自己——看来从一个圈子、一种文化到另一个圈子、另一种文化,从一个时代到另一个时代,都是同一个。如果说我这一生里,在数学之外、在爱女人之外,还有什么东西让我感到奥秘与吸引(不错,是晚近的事),那正是这少数几种力量的隐秘本性——它们有能力驱使我们行动,为“最好”也为“最坏”,为埋藏也为创造。

\subsection{10. 一桩敬重之举}

这场最终取名为“葬礼”的反思,起初是一桩敬重之举。敬重那些我发现了的东西——我看着它们在一片虚无中凝聚、成形,我是第一个尝到它们的滋味与劲道的人,我给它们起了名字,既为表达我对它们的认识,也为表达我的敬重。对这些东西,我给出了自己身上最好的部分。它们汲取着我身上的力量,生长、绽放,像许多粗壮的枝条从同一根活着的树干上迸发出来,那树干有着粗壮而众多的根。这些是活着的、在场的东西,不是可有可无的发明——它们紧密相依,构成一个活生生的统一体,这统一体由它们每一个构成,又给每一个以位置与意义、一个起源与一个终点。我很久以前就把它们留下了,没有丝毫不安,也没有丝毫遗憾,因为我知道我留下的东西是健全而强壮的,它根本不需要我,就能依着它自己的本性继续生长、绽放、繁衍。我留下的不是一袋金币——那是可以被人偷走的;也不是一堆工具——那是会生锈、会腐烂的。

然而,年复一年,当我自以为离那个我离开的世界很远的时候,却有一阵阵阴险的轻蔑、隐隐的嘲弄,一直飘到我隐居之处,在这里那里回旋,指着这些东西中的某些——我明知它们强壮而美,有它们的位置,有它们独一无二、任何别的东西都永远无法代替的功能。我感到它们像是敌对世界里的孤儿——那个世界患着轻蔑的病,专门纠缠没有铠甲护身的东西。这场反思就是在这样的心境中开始的,作为对这些东西的一桩敬重之举,并因此也是对我自己的一桩敬重之举——作为对它们与我之间那条深沉纽带的提醒:谁若乐于对这样一件由我的爱滋养出来的东西摆出轻蔑的样子,他所乐于轻蔑的正是我,以及一切从我这里生出的东西。

对于另一种人,情形也一样:他第一手知道那条把我和某件东西连起来的纽带——那东西他不是从我这里、又是从哪里学来的呢——却摆出一副把这条纽带视为无足轻重、或者干脆无视它的样子,或者(哪怕只是默默地、以省略的方式)为自己、或为别人僭取一种虚假的“父权”。我在其中
% ---------------- p.122 ----------------
\origpage{XXII} 清清楚楚看到一桩轻蔑之举:轻蔑那件从工匠手里生出的东西,轻蔑那默默而精细、使这件东西得以生出并由此诞生的工作,轻蔑那工匠,而首先(以一种更隐蔽也更本质的方式)轻蔑他自己。

如果我的“回归数学”只消让我重新想起这条纽带,只消在我心里唤起这桩当着所有人——当着那些摆出轻蔑样子的人,也当着无动于衷的见证者——的敬重之举,那么这次回归也就不算白费。

确实,我与我留下的那部写出来的和没写出来的(至少是未发表的)作品,真的失去了接触。开始这场反思时——我看那些枝条还相当分明,却不怎么记得它们原是同一棵树的一部分。说来奇怪:非得那幅我留下之物遭劫掠的图景在我眼前一点点展开,我才在自己心里重新找回那种意义——那种关于被如此劫掠、如此拆散之物的活生生的统一性的意义。一个人拿走了一些金币,另一个拿走了一两件工具,用来炫耀,甚至用来使用——但是那赋予我所留之物以生命和真正力量的那种统一性,却从每一个人、也从所有人手里漏掉了。不过我清楚知道有一个人,深深感到过这种统一性和这种力量,在他内心深处今天仍然感到它;而他却乐于耗费自己身上的力量,要去摧毁他在别人身上、通过别人的作品感到过的那种统一性。作品的美与创造之德,恰恰就在这种活生生的统一性之中。尽管遭了劫掠,我仍把它们找回来,完好如初,仿佛我刚离开它们——只是我成熟了,今天用一双新的眼睛看它们。

然而,如果说有什么东西真被劫掠、被损毁、被卸掉了它原本的力量,那是在那些人身上:他们忘了自己身上安歇着的力量,自以为在劫掠一件任由他们摆布的东西,其实他们只是把自己与那创造之德割断——那东西就在他们手边,一如它在所有人手边,却绝不任他们摆布,也不在任何人的权力之下。

这样,这场反思,以及通过它而来的那次不期而至的“回归”,也让我重新接触到一种被遗忘的美。正是完完全全感受到了这种美,才使那桩敬重之举有了全部意义——它在注《我的孤儿们》\fnote{15}{}{这条注(第46号)在收入“葬礼”的所有注中,按时间是最早的一条。}里表达得笨拙,而我刚才在这里,在完全知情的情况下,又把它重述了一遍。


% ================================================================
%  二 · 第一部 (第 5-12 章, 原书 p.123-254)
% ================================================================
\partdivider{第一部}{愚妄与更新}{愚妄与更新}

% 原书第125页献辞
\dedication{%
献给曾是我的长辈的人,\\[0.5em]
他们手足般地接纳我\\[0.5em]
进入那个曾属于他们、\\[0.5em]
而后也成为我的世界。\\[2em]
献给曾是我的学生的人,\\[0.5em]
我给了他们我身上最好的\\[0.5em]
——也给了最坏的……%
}

\setcounter{chapter}{4}
\markboth{收获与播种·愚妄与更新}{}

% ================================================================
%  第 5 章 劳动与发现 (原书 p.127-132, 边码 p.1-p.8)
%  源文: _work/src1/task-ch05.md
% ================================================================
% ---------------- p.127 ----------------
\chapter{劳动与发现}

\noindent 1983年6月\par

\section{(1) 孩子与好上帝}

\origpage{1} 我眼下正在做的这些数学笔记,是十三年来头一批我打算拿去发表的文字。读者不会奇怪:沉默了这么久之后,我的表达风格变了。然而这种表达上的变化,并不是工作风格或工作方法有什么变化的标志\fnoteP{1}{(1)}{(1984年3月补写)说我的工作“风格”与“方法”没有变,恐怕失于武断——毕竟我在数学上的表达风格已经深刻改变了。这一年来用于《追寻栈》的大部分时间,是坐在打字机前敲下一些思考;它们就是要几乎原样发表的(只除了后来补加的若干相当简短的注,借互见、勘误之类方便阅读)。没有剪刀,没有浆糊,不必费劲去炮制一份“定稿”(那份定稿尤其不该让任何通向它的过程透出影来)——这毕竟也是“风格”与“方法”的变化!除非把数学工作本身与写作、与陈述结果的工作割裂开来,而那是人为的,因为那不合事情的实情:数学工作与写作是拆不开的。} (1),更不是说我的数学工作本身的性质发生了转变。不但这性质一如既往——而且我确信:发现的工作,其性质在一个人与另一个人身上是相同的,它超越种种制约与气质所造成的差异,那些差异千变万化,无穷无尽。

发现是孩子的特权。我要说的是小孩子——那个还不怕弄错的孩子,不怕显得像个傻瓜,不怕不正经,不怕不跟大家一样。他也不怕他瞧着的东西不识趣:与他对它们的期待不一样、与它们本该是的模样不一样,或者更确切地说:与人们理所当然认定它们是的模样不一样。他不懂那些无声而毫无缺口的共识——它们是我们呼吸的空气的一部分,是所有明白人、所有公认的明白人的空气。天知道,自时间的黑夜深处以来,出过多少这样的人——明白人,公认的明白人!

我们的头脑里塞满了一种杂七杂八的“知识”——恐惧与懒惰、渴念与禁忌纠缠成一团;还有随处涌来的信息和按一按就出来的解释——一个封闭的空间,信息、渴念与恐惧在里面堆积,而旷野的风从来吹不进去。除了那点例行公事的“会做”,这“知识”的主要作用,似乎就是排掉活的知觉,排掉对世间事物的认识。它的效果,主要是一种巨大的惰性,一种往往压得人喘不过气来的重负。

% ---------------- p.128 ----------------
小孩子发现世界,就像他呼吸——呼吸的涨落让他把世界接进自己娇嫩的存在里,也让他投身于接纳他的世界。成年人也在发现,在那些稀少的时刻:他忘了自己的恐惧和知识,用睁得大大的、渴于认知的眼睛看事物、看自己,一双崭新的眼睛——一双孩子的眼睛。

\asterisms

\origpage{2} 上帝一边发现世界,一边创造世界;更确切地说,他永恒地\textbf{创造}着世界,一边创造一边发现——也一边发现一边创造。他创造了世界,又日复一日地创造着它,反反复复,千百万个千百万次,不歇息,摸索着,错了千百万个千百万次又拨正航向,不知疲倦……每一次,在这场用探针敲击万物、听万物回应的游戏里(它们回答:“这一下敲得不赖”,或者:“这一下你可是彻底犯浑”,或者:“顺顺当当,就这么来”),在又一次探针中——它修正或重提前一次探针,回应前一个回答……在造物主与万物之间这场无限对话的每一次往复里(这对话在创世的每一刻、每一处发生),上帝学习着、发现着,他越来越亲切地认识万物,而万物也在他手中渐渐有了生命、有了形状、不断变形。

这就是发现与创造的路数;亘古以来,似乎一直如此(就我们所能认识的而言)。一直如此——并不需要人类姗姗来迟地登场(不过是一两百万年前的事),撸起袖子插手进来——再加上最近那些众所周知的糟糕后果。

我们当中,不时有人发现这样或那样的事物。有时他会在自己的生命里重新发现——怀着惊叹——\textbf{发现}究竟是怎么一回事。每个人身上都有他所需要的一切,去发现这广袤世界上吸引他的一切,包括他身上那份奇妙的能力——世上最简单、最显然不过的东西!(然而许多人已经忘了的东西,就像我们忘了唱歌,忘了像孩子那样呼吸……)

每个人都能重新发现发现与创造是怎么一回事,而没有人能把它发明出来。它们先于我们就在那里,它们就是它们之所是。

\section{(2) 错误与发现}

回到我真正意义上的数学工作的“风格”,或它的“性质”、它的“路数”:它们如今与从前一样,都是好上帝亲自教给我们的——不用言语,教给每一个人,天知道在什么时候,也许远在我们出生之前很久。我照他的样子做。这也是每个人出于本能所做的,一旦好奇心驱使他去认识某一样东西——从那一刻起,这东西就被这份欲望、这份渴念占据了……

\origpage{3} 当我对一样东西感到好奇,数学的也好,别的也好,我就去问它。我去问它,不在意我的问题也许很蠢、或者会显得很蠢,也不要求它非得先经过反复掂量。问题常常采取断言的形式——一个其实只是探针的断言。我对自己的断言信几分?这当然取决于我在理解眼前之物上走到了哪一步。往往,尤其在一项研究的开头,这断言干脆就是错的——可还是得把它做出来,才能让自己信服它是错的。往往,只消把它写下来,
% ---------------- p.129 ----------------
错就一目了然;而在写下来之前,只有一片模糊,一种不适,而不是这份明白。这样就能少一分无知,再次进攻,带着一个也许不那么“离谱”的问题—断言。更常见的是:断言照字面理解是错的,但那个仍旧笨拙地试图借它表达出来的直觉是对的,虽然还模糊。这直觉会渐渐从一层同样不成形的脉石中析出——起初是些错误或不恰当的观念;它会渐渐走出“未被理解”的幽冥(那只是等着被理解),走出“未知”的幽冥(那只是等着被认识),取得一个只属于它自己的形式,让轮廓变得精细、鲜明——随着我向眼前这些事物提出的问题越来越精确、越来越切题,把它圈得越来越近。

但也会这样:靠这条路子,反复的探针会聚成局势的某一幅图像,从雾中显出相当清晰的线条,足以让人开始相信这幅图像很好地表达了实情——其实满不是那么回事,只要这图像沾上一个分量十足的错,一个足以叫它深深走样的错。从最初发现的“对不上号”出发——所得图像与某些确凿事实对不上,或与另一些同样为我们所信任的图像对不上——逐步查出这样一个错误观念的工作,有时很艰苦,而且往往伴随着越来越紧的张力,随着人逼近矛盾的结,那矛盾先还模糊,渐渐越来越刺耳——直到它终于爆开:错误的发现,某一种世界观的坍塌,像一次解放。\origpage{4}\textbf{错误的发现,是一切发现工作中最关键的时刻之一,是其中最富创造性的时刻之一},无论那是数学工作,还是发现自我的工作。正是在那一刻,我们对所探究之物的认识突然更新了。

\textbf{怕错与怕真,是同一回事。}害怕弄错的人,无力发现。正是在我们害怕弄错的时候,我们身上的错误才变得岩石一般不可动摇。因为在恐惧中,我们死死抓住某一天被我们判定为“真”的东西,或者抓住那些历来就被当作“真”呈现给我们的东西。当我们不是被害怕失去一种虚幻安全的恐惧所驱动,而是被求知的渴念所驱动,那时错误就像痛苦或悲伤一样,从我们身上穿过,永不定形,而它经过的痕迹,是一种更新了的认识。

\section{(3) 不可告人的劳作}

一切真正研究的自发路数,可以说从来不出现在那些本应传达、传递“被找到”之物的实质的文本或话语里——这绝不是偶然。文本与话语,多半只满足于记下一些“结果”,用的形式,在凡夫俗子眼里,必须显得像一条条严峻而不变的定律,亘古以来就刻在某种巨型图书馆的花岗岩法版上,由某位全知的神口授给那些入门者—书记—学者之流;口授给那些写学术书和同样学术的论文的人,那些站在讲坛高处、或在讨论班更小的圈子里传递知识的人。有没有哪怕一本教科书、一本供小学生、中学生、大学生乃至“我们的研究者”用的手册,能给不幸的读者一点点关于研究究竟是什么的概念——除了那个普遍接受的概念:研究嘛,就是当你很有学问、考过一大堆考试甚至竞考的时候,就是那些大脑袋,巴斯德和居里,还有诺贝尔奖得主什么的……而我们这些读者或听众,勉勉强强把大人物们为了人类的利益而惠予记下的“知识”吞下去,我们顶多(如果肯下苦功)够格通过年终考试,而且
% ---------------- p.130 ----------------
还不一定……

\origpage{5} 这样的人有多少呢——包括那些不幸的“研究者”自己,那些为论文或文章发愁的人,甚至包括我们当中最有“学问”、最有声望的人——有谁天真到看得出:“寻找”不是别的,正是怀着激情\textbf{去问}万物——就像一个孩子\textbf{想知道}他自己或他的小妹妹是怎么来到这世上的。寻找与找到,也就是:提问与倾听,是世上最简单、最自然不过的事,世上没有人拥有它的特权。这是我们一出生就全都收到的一份“馈赠”——它注定要以无穷无尽的面孔表达自己、舒展自己,此一时与彼一时不同,一个人与另一个人也不同。

当你冒险把这类话说出口,你在这一位和那一位那里收到的,从最确信自己是笨学生的笨学生,到最确信自己有学问、远在凡夫之上的大学者,都是一样的笑:一半尴尬,一半心照不宣,好像你刚讲了一个有点粗俗的笑话,好像你在显摆一种一眼就能看穿的天真;“这话都不错,谁也不该看不起,这我同意——可也别太过分啊——笨学生就是笨学生,他不是爱因斯坦,也不是毕加索!”

面对这样一致的赞同,我再坚持就不识趣了。真是不可救药,我又一次失去了闭嘴的机会……

不,这绝不是偶然:各类教育书、劝世书和各种手册,异口同声地把“知识”呈现得像是一出娘胎就穿得整整齐齐,从那些天才的头脑里出来——那些头脑为了我们的益处把它记了下来。也不能说这是恶意,即便在少数场合,作者相当“门儿清”,知道那幅图像(他的文本免不了要暗示的图像)与实情毫不相符。在这种场合,陈述有时不止是一堆结果和诀窍:有一股气息穿过它,有一种活的视野赋予它生命,而这视野有时便从作者传到专心的读者那里。但一种默契——力量看来相当可观——使得文本不留下半点它由之产生的那些工作的痕迹,即使它以镌刻般的笔力表达出对事物有时相当深刻的视野,而那视野正是这工作真正的果实之一。

\origpage{6} 说实话,有些时候,我自己也隐隐感到这股力量、这份无声共识的分量——就在我筹划写作、出版《数学反思》的时候。如果我试着去探一探这份共识所采取的默不作声的形式,或者更确切地说,去探一探我内心对我这计划的抵抗(它由这份共识引发)所采取的形式,立刻冒出来的词是“不体面”。这份共识,不知从何时起已内化在我心里,它对我说(而这是我头一回费心把它这几周、也许是这几个月来带几分固执地对我嘟囔的东西,拉到日光下,拉进我的视野里):“把一项发现工作中的起起落落、边角上那些弄砸了的摸索、总之那些‘脏衣服’,摊在别人面前,甚至当众摊开——这不体面。这只是浪费读者的时间,而读者的时间是宝贵的。再说,这还要多出好多页,得排版、得印刷——多浪费啊,科学印刷品价钱这么贵!非得虚荣到什么地步,才会这样摊开那些对谁都没意思的东西,好像连我那些胡乱的折腾也成了了不起的事——总而言之,一个摆谱的机会。”还有更隐秘的一层:“发表这样一份反思的笔记——照它真正进行的样子——是不体面的,就像在公共广场上做爱是不体面的,或者把分娩的劳作留下的血污床单拿出来展示、哪怕只是让它晾在那里,也是不体面的……”

这里的禁忌,采取的是性禁忌那副既阴险又专横的面孔。正是在写这篇导论的时候,我才开始隐约看见它非凡的力量,以及这本身即非凡的事实所含的意义——这事实正是那股力量的证明:发现的真正路数,如此
% ---------------- p.131 ----------------
令人不知所措地简单,一种孩子气的简单,却几乎无处透出;它被无声地抹去、忽视、否认。甚至在科学发现这个相对无害的领域里也是如此——谢天谢地,那不是发现自己的“小鸡鸡”之类的事——总之是一项可以放心交到所有人手里的“发现”,而且(人们会以为)没什么可隐瞒的……

倘若我愿意顺着这里出现的这根“线”走下去——这根线一点也不细,反而粗壮结实得很——它肯定会把我带到远得多的地方,远过那几百页同调—同伦代数;那些页面我终究会写完,交给印刷厂的。

\section{(4) 他人之无谬与对己之轻蔑}

\origpage{7} 刚才我谨慎地说“我的表达风格”变了,还让人以为这没什么可奇怪的——说到底,那是个委婉说法。您明白:十三年没写过东西,当然跟从前不一样了,“表达风格”总得变,必然的……区别在于,从前我“表达”(原文如此)得像所有人一样:我先做工作,然后倒过来重做一遍,把所有的涂改都仔细擦掉。擦着擦着又出现新的涂改,把整个工作搅得有时比初稿还糟。于是又得重做——有时三遍,甚至四遍,直到一切都无可挑剔。不但没有可疑的角落,也没有悄悄扫到合宜的家具底下的垃圾(我从来就不喜欢角落里的垃圾,既然都费心扫地了);尤其是,读定稿的时候,它给人的印象固然体面(像任何别的科学文本一样),那就是:作者(此处即区区在下)是化为肉身的无谬。他万无一失地正好落在“那些”正确的概念上,然后是“那些”正确的陈述上,一环扣一环,像上了油的马达一样低吟着,那些证明“咚”的一声落下来,每一个都恰好在它该来的时候!

\origpage{8} 效果如何,请设想一个毫无戒心的读者:比如说一个正在学毕达哥拉斯定理或二次方程的中学生,甚至我的一位在所谓“高等”研究或教学机构里的同行(明白人一听就懂,此致!),(不妨说)正费劲地读某位显赫同行的某篇文章!这类经验在一个小学生——甚至大学生或研究者——的一生中重复成百上千次,再加上家庭里以及全世界所有国家的所有媒体里那场合宜的合唱,其效果可想而知。只要肯留心,人在自己身上、在别人身上都能看到:那是一种对自身毫无价值的内心确信,与那些“懂行的人”、“干事的人”的能力和重要性恰成对照。

这份内心的确信,有时会被一种能力所补偿——记住不理解之物的能力,甚至某种操作上的灵巧:矩阵相乘,用“正题”和“反题”拼出一篇法文作文……总之是鹦鹉或耍把戏的猴子的本领,如今比以往任何时候都更受看重,由人人艳羡的文凭认可,由安稳的职业生涯报偿。但即便那位文凭缝满一身、位置坐得安稳、也许还满身荣誉的人,在心底里也并不上当,并不把这些重要、这些“价值”的假标志当真。甚至那位更少见的人——他把全部赌注押在某样真正天赋的发展上,在职业生涯中懂得施展身手、做出创造性作品——在心底里也并不信服自己名声的光彩,尽管他常常想借这光彩让自己、也让别人看走眼。一种从未被审视过的同样的怀疑,住在这一位和那一位身上,也住在随便哪个笨学生身上;一种同样的确信,也许他们永远不敢去认识它。

正是这份怀疑、这份未曾说出的内心确信,推动着这一位和那一位不断地超越自己——
% ---------------- p.132 ----------------
在荣誉或作品的积累中;并且把他们私下里啃噬着自己的那份对己之轻蔑,投射到别人身上(尤其是那些他们对其多少有些权力的人身上……)——这是一场注定不可能的逃脱尝试,靠的是积累他们优于别人的“证据”\fnoteP{2}{(2)}{(1984年3月补写)重读最后这两段,我有些不适之感,因为写它们的时候,我牵连的是别人,而不是我自己。显然,写的时候,“我自己也可能与此有关”这个念头并没有碰到我。我肯定什么也没学到:我就这样满足于把多年来在别人身上察知、并以许多方式看着被证实的东西,白纸黑字写下来(无疑还带着某种满意)。在这份反思的后续部分,我不得不回想起:我的一生中,对别人的轻蔑姿态并不缺少。我所把握到的他人之轻蔑与对己之轻蔑之间的关联,若在我这个人身上竟然缺席,那才奇怪;健全的理性(以及类似情形下我对自身视而不见的经验——我终究还是察觉到了)告诉我:肯定不该是这样!然而眼下这只是一条纯粹的推论,它唯一可能的用处,是促使我亲眼看一看究竟发生了什么,去看、去审视(如果它确实存在过、或还存在着)这份仍属假设的对己之轻蔑——它埋得如此之深,至今完全逃过了我的目光。真的,可看的东西并不缺!这一件忽然在我看来是最要紧的几件事之一,恰恰因为它藏得如此之深。……\emph{[(1984年8月)不过关于这一点,参见“屠杀”一注(n\textsuperscript{o}~87)末两段的反思。]}} (2)。

% ================================================================
%  第 6 章 梦与做梦者 (原书 p.133-146, 前半: 6.1-6.8)
%  源文: _work/src1/task-ch06a.md
%  说明: 本章后半 (p.147-152, 6.9-6.11) 见 part1-ch06b.tex
% ================================================================
\chapter{梦与做梦者}

\noindent 1984年2月\par

% ---------------- p.133 ----------------
\section{(5) 被禁的梦}

\origpage{9} 借着《追寻栈》(À la Poursuite des Champs)写作中断三个月的机会,我要把《导论》接着去年六月停下的地方写下去。相隔半年多,我刚把它仔细重读了一遍,并添了几个小标题。

写这篇《导论》时,我很清楚:这类思考免不了招来一大堆“误解”——而想事先堵住它们也是白费力气,那无非是在原有的误解之上再堆上一些!关于这一点,我只想补充一句:我丝毫无意向那种由千年惯例确立起来的科学写作风格开战——我自己就曾勤勤恳恳地奉行它二十多年,还把它当作数学家这一行当的本质部分教给我的学生们。不管对错,今天我也仍然这样看待它,并且继续这样教。甚至可以说,我倒更显得老派:我坚持工作要做到底,从头到尾一针一线地缝,不放过任何一个稍有阴影的角落。近十年来我不得不往酒里掺水,那实在是形势所迫!“规范的写作”在我看来仍然是数学工作的重要一环,既是一种发现的工具——用来检验并深化对事物的理解,没有它,这理解始终是粗略而零碎的——也是一种传达这种理解的手段。从教学的角度看,严格的表述方式,也就是演绎的方式(它绝不排除铺开宏阔画卷的可能),有明显的好处:简洁,查考方便。这些好处是实实在在、有分量的——当讲稿面向
% ---------------- p.134 ----------------
\origpage{10} 数学家(比如说)、而且是已经对所述主题的某些来龙去脉、或对与之十分相近的主题足够熟悉的数学家时。

然而,若一篇讲稿的对象是孩子、是年轻人,或是那些事先完全不“在行”的成年人——他们的兴趣已经觉醒,而且多半(并将继续如此,而且事出有因……)对一项发现性工作的真正进程全然无知——那么这些好处就完全成了幻影。更确切地说,是这样的读者:他们连这种工作的存在都不知道——而这工作其实人人做得到,只要他有好奇心和常识;我们对宇宙万物的智性认识,正是从这工作中不断生出、又不断再生,包括那些以宏伟的编排体系表达出来的认识,比如欧几里得的《几何原本》,或达尔文的《物种起源》。对这种工作的存在与性质全然无知,几乎是普遍现象,各级教师也不例外,从小学教师到大学教授。这是一个非同寻常的事实;去年写这篇《导论》的第一部分、开始这番反思时,它才清清楚楚地显现在我眼前——同时也让我隐约看到了这件令人困惑之事深深的根……

即便讲稿面对的读者在各方面都确实“在行”,仍有一件重要的东西,是“严格的”表述方式不许自己传达的。而这东西,在严肃人士的圈子里——尤其在我们这些科学家当中——也完全不受待见!我说的就是梦。梦,以及梦向我们耳语的那些景象——起初像梦一样不可触摸,而且往往迟迟不肯成形。长年累月、甚至整整一生的紧张工作,也许都不足以让某个梦的景象完全显现出来,不足以看它凝聚、打磨,直到获得钻石的硬度与光华。这就是我们的工作,无论我们是用手还是用精神做工的工匠。当工作完成、或工作的某一部分完成时,我们把那有形的成果呈现在所能找到的最亮的光下,我们为它欢喜,还常常以它为荣。然而,我们长久打磨的这颗钻石里,并没有那个在打磨中激励我们的东西。也许我们造出了一件极精确的工具、一件好用的工具——但工具本身就是有限的,就像人手所做的一切,哪怕它在我们看来很了不起。一个景象,起初没有名字、没有轮廓,淡薄得像一片雾霭,它引导着我们的手,让我们伏在活儿上,感觉不到时辰流逝,也许也感觉不到年月。一片从雾与幽暗的无底之海中无声地剥落下来的雾霭……我们身上那无界限的东西,就是她,就是这片海——只要我们的渴望使她受孕,她便随时准备孕育、随时准备分娩。梦正是从这场婚配中涌出,像蜷伏在养育它的子宫里的胚胎,等待着幽暗的劳作把它带向第二次诞生,带向天光。

一个蔑视梦的世界有祸了——那也是我们身上深刻的东西遭蔑视的世界。我不知道在我们这种文化(电视、电脑和洲际火箭的文化)之前,是否还有别的文化宣扬过这种蔑视。这大概正是我们区别于前人的许多地方之一——我们把前人彻底取代、几乎可以说从地球表面清除掉了。我不曾知道有哪一种别的文化是不尊重梦的,不曾知道有哪一种文化里,梦的深根不被所有人感受到、不被承认。在一个人或一个民族的一生中,可曾有哪件宏大的作品,不是生于梦、不是先由梦哺育,而后才在大白天里绽放的?然而在我们这里(是不是已经该说:到处都一样?),对梦的敬重被叫作“迷信”;而众所周知,我们的心理学家和精神病学家早已把梦翻来覆去、量了个遍——那点分量,想必勉强够塞满一台小型电脑的存储器。也确实,“在我们这里”,再没有人会生火,也没有人敢在自己家里看着孩子出生,或看着母亲、父亲死去——
% ---------------- p.135 ----------------
\origpage{11}有诊所和医院就是干这个的。谢天谢地……我们这个世界,如此自豪于它的力量——以原子百万吨级计,以图书馆和电脑里储存的信息量计——恐怕也正是这样一个世界:在这里,每个人的无能,面对生活中那些简单而本质之物时的恐惧与蔑视,达到了顶点。

\origpage{12} 幸好,梦就像最压抑的社会里那原初的性冲动一样,命硬得很!不管是不是迷信,它仍旧偷偷地、执拗地向我们耳语一种认识——那是我们醒着的精神太笨重、或太怯懦,把握不住的认识——并且给它启发我们形成的种种计划注入生命、插上翅膀。

我刚才似乎暗示过,梦往往不肯成形;那只是表面现象,并未真正触及事情的底。那“不肯”多半来自我们醒着的精神,来自它平常的“心态”——而且“不肯”这个词还是轻描淡写!那更像是一种深深的戒心,它掩盖着一种古老的恐惧——认识的恐惧。就本来意义上的梦而言,梦的消息越是贴近我们,越是因为一旦被听见就会带来我们人格深处蜕变的威胁而显得沉重,这恐惧就越是起作用,它立起的屏障就越是有效。但看来,即便在数学的“梦”这种相对无害的场合,这份戒心也依然在场、依然有效;以至于一切梦似乎不仅被逐出了文本(反正我连一部记下梦的文本也没见过),而且也被逐出了同事间的谈话——不论是小圈子里,甚至一对一的时候。

事情若果然如此,那当然不是因为数学的梦不存在、或不再存在了——若真如此,我们的科学早就贫瘠了,而绝非如此;这种表面上的缺席、这种沉默的合谋,其原因想必与另一种共识紧密相连——那就是要小心抹去一切痕迹、一切提及:抹去那借以完成发现、借以更新我们对世界之认识的工作。或者不如说,围绕着梦、也围绕着梦所唤起、所启发、所滋养的工作的,是同一个沉默。以至于“数学的梦”这个说法本身,在许多人看来都会是无意义的胡话——我们如此经常地被一些按电钮式的陈词滥调驱动,而不是被我们对某种极其简单、日常而又重要的现实的直接经验驱动。

\section{(6) 做梦者}

事实上,经验告诉我:当精神渴望认识梦,而不是逃避它(或手里拿着一副现成的框框去套它——那是一样的),梦就丝毫不会“不肯成形”——它肯让人细致地描述,肯交出它的消息:那消息永远简单,从不愚蠢,有时还叫人震动。恰恰相反,我们心中的“做梦者”是一位无与伦比的大师:他每一次都能找到、或凭空造出最合用的语言,来绕过我们的恐惧、摇醒我们的麻木;他的舞台手段变化无穷,从完全没有一丝视觉或感官要素,到最令人瞠目的场面调度。当他显现时,绝不是为了躲开,而是为了鼓励我们(几乎总是白费力气,而他的好意从不疲倦……)走出我们自己,走出他看见我们深陷其中的那种沉重——有时他还会不动声色地拿这种沉重寻开心,用滑稽的色调把它戏仿一番。倾听我们心中的“做梦者”,就是与我们自己交流——逆着那些强大的堤坝,它们非要禁止我们这样做不可。

但能做大的人,也能做小。既然我们能借梦与自己交流、向自己显现自己,那么,把数学的梦——比如说——那绝不私密的消息传达给别人,想必也同样简单易行;这种梦并不牵动
% ---------------- p.136 ----------------
\origpage{13} 力量相当的抵抗。

而说实话,在我做数学家的过去里,除了追随一片片从浓重致密的雾的织物上一一剥落下来的梦,把它们“梦”到底——直到它们最显明、最坚实的显现,无可辩驳的显现——我还做过别的什么?又有多少次,我为自己那份固执而急得跺脚:我总要把每块凝结着我的梦的宝石、或半宝石,一丝不苟地打磨到最后一个刻面——而不是听从一股更深的冲动:去追随那母体织物千变万化的幽径,直到梦与它显在的、可以“发表”的化身(总之,按通行的规范)之间那片朦胧的边界!况且,我当时正要听从那股冲动,投入一项“数学科幻”的工作——一部关于“Motif”理论的“某种白日梦”,那理论当时还纯属假设,而且直到今天仍然如此,而且理所当然:因为再没有另一个“白日做梦者”投身这场冒险。那是六十年代末;那时我的生活(我丝毫没有想到)正要拐一个大弯,此后十来年间,这转折把我的数学热情挤到了边缘的位置,甚至遭到否认。

但细究起来,《追寻栈》——沉默十四年之后的第一部出版物——恰恰属于那场从未写出的“白日梦”的精神,而它似乎成了那场梦临时的后续。当然,至少初看之下,这两个梦的主题之间的差异,是两个数学主题之间所能有的最大差异;何况头一个(即 Motif 那个)似乎位于“用现有手段可做之事”的地平线之外,而第二个——即著名的“栈”及其同类——看起来完全伸手可及。这些差异,或许可以称为偶然的、意外的差异,它们也许会比人们预料的早得多就烟消云散\fnoteP{1}{(3)}{我在这里特别想到已故的莫德尔(Mordell)猜想、泰特猜想、沙法列维奇(Chafarévitch)猜想——去年,法尔廷斯(Faltings)一篇四十页的手稿把这三个猜想一并证明了;而就在此前,“在行”的人们早已确立的共识是:这些猜想“够不着”!恰好,我那心爱的“远阿贝尔代数几何”纲领当作拱心石的那条基本猜想,正与莫德尔猜想相近。(甚至据说后者可以由前者推出——这就更说明,这个纲领不是给严肃人士看的故事……)} (3)。在我看来,这些差异对于两个主题各自可能引发的这类工作,影响相对很小——因为这里说的正是“白日梦”,或者用不那么挑衅的说法:把概念上的粗加工一直做到一个足够融贯、足够精确的整体景象,足以或多或少完整地使人确信:这景象在根本上确实符合事物的实际。就本书所展开的这个主题而言,这大致意味着:对这一景象之有效性的逐项验证,变成了一件纯属本行手艺的事。这当然可能要求巨大的工作量,其中也需要机巧和想象,而且想必还有意外的转折与前景,幸而会使它不至于成为纯粹的例行公事(用安德烈·韦伊的话说,一场“长长的练习”)。

总之,这正是我过去做腻了又重做的那类工作,它就在我指尖上,因此在我余下的岁月里不必再做。只要我再度投入数学工作,我的精力肯定最该用在“白日梦”的边缘地带。在这个选择里,启发我的并不是什么效益考虑(假如这种考虑能启发得了谁的话),而恰恰是一个梦,或一些梦。我身上这股新的冲劲若能显出力量,那力量必定是从梦里汲取的!

% ---------------- p.137 ----------------
\section{(7) 伽罗瓦的遗产}

\origpage{14} 看起来,在所有的自然科学中,只有数学把我所谓的“梦”、或“白日梦”判以一道看来绝对的禁令,而且这禁令已施行两千多年。在别的科学里,包括物理学这样公认“精确”的科学,梦至少是被容忍的,甚至(视时代而定)是被鼓励的,只不过用的是些更“拿得出手”的名字,比如:“思辨”、“假设”(如著名的“原子假设”,它出自一个梦——对不起,出自德谟克利特的一种思辨)、“理论”……从那个“不敢道出自己名字的梦”的身份,过渡到“科学真理”的身份,是经由几乎察觉不到的层层台阶、经由一个逐渐扩大的共识完成的。相反,在数学里,这几乎总是(至少如今)一次突变,靠的是一篇证明的魔杖一挥\fnoteP{2}{(4)}{况且就在今天,我们也还遇得到身份可疑的“证明”。格劳尔特(Grauert)那个以他命名的有限性定理,多年来就是这样:他那个证明谁也不曾读完(而愿意读的人并不缺!)。这一困窘后来由另一些更透明的证明解决了,其中有些还走得更远,它们取代了最初那个证明。类似而更极端的情形,是所谓的“四色问题”的“解决”:其中计算的部分是用电脑(以及几百万美元)了结的。所以那是一个“证明”,它的依据不再是理解一个数学情境而来的内心确信,而是对一台没有理解能力的机器的信任——而使用它的数学家既不知道这台机器的结构,也不知道它如何运作。即便计算再由别的电脑、按别的计算程序核实一遍,我也不认为四色问题就此了结。它只是换了一副面孔:如今要做的几乎不再是找一个反例,而只是要一个证明(当然,得是读得懂的证明!)。} (4)。在那样一些年代——数学定义与数学证明的概念还不像今天这样清楚,还不成其为(或多或少)普遍的共识——却也有些显然重要的概念,身份暧昧不明,比如“负”数(帕斯卡拒斥过它),或“虚”数。这种暧昧,至今仍反映在通行的语言里。

\origpage{15} 定义、陈述、证明、数学理论这些概念的逐步澄清,在这方面极为有益。它让我们意识到我们手头的工具——尽管简单得像小孩玩意——有多么强大:靠它们,我们能把那看似无法表述的东西表述得完美精确——几乎只要足够严格地使用日常语言就行。如果说有什么东西从我童年起就让我着迷于数学,那恰恰是这种力量:用词语去圈住、并完美地表达某些数学事物的本质——这些事物乍看之下形态如此飘忽、如此神秘,仿佛在言语之外……

然而这种力量、这种完美精确与证明所提供的资源,有一个糟糕的心理上的反作用:它们更加强化了对“数学的梦”的传统禁忌;也就是说,强化了对一切不肯以常规的精确面貌出现的东西的禁忌(哪怕牺牲的是更宏阔的景象)——那种精确要有形式完整的证明来担保它“货真价实”,否则(而且如今越来越如此……)就得有据信可以整理成形的证明概要来担保。偶尔的猜想,勉强可以容忍,条件是它得符合问卷式的精确要求,而问卷上只允许“是”或“否”两种回答。(而且,还用说吗,还得是那个敢提猜想的人在数学界里有门有面。)据我所知,还没有过一个例子:把一门在其本质部分明明白白是猜想性的数学理论,当作“实验”来发展。确实,按现代规范,十七世纪以来发展起来的整套“无穷小”演算——后来成了微分学与积分学——会显得像一场白日梦,而它最终要等到两个世纪之后,才靠柯西的魔杖一挥,变成
% ---------------- p.138 ----------------
严肃的数学\origpage{16}。这必然让我想起埃瓦里斯特·伽罗瓦(Evariste Galois)的白日梦——他不巧也碰上了同一位柯西,运气不佳;但这回不到一百年就够了:又一下魔杖(这回是若尔当,如果我没记错),就让这场梦获得了公民权,并为此改名为“伽罗瓦理论”。

从这一切得出的结论——这结论对“1984年的数学”并不有利——是:幸好牛顿、莱布尼茨、伽罗瓦这样的人(肯定还漏掉很多,我对历史并不内行……)当年没有被我们今天的规范绊住;那时候,他们只管发现,没工夫立规范!

伽罗瓦这个例子不请自来,触动了我的一根敏感的弦。我似乎记得,第一次听人说起他和他的奇崛命运时(那时我大约还在念中学或大学),一种兄弟般的亲切感就在我心里醒来了。和他一样,我感到自己心中有一份对数学的激情——也和他一样,我觉得自己是个边缘人,是那个(在我看来)把他抛弃了的“上流社会”里的异乡人。可我自己最终也成了这上流社会的一员,又在某一天离开它,没有遗憾……这份有点被遗忘的亲近感,最近又以全新的面貌重新出现在我身上——当时我正在写《纲领概要》(Esquisse d'un Programme,为着我申请作为研究员进入国家科学研究中心[CNRS]的事)。这份报告主要概述我近十年来一些主要的思考主题。在所有这些主题中,最让我着迷、也是我打算在今后几年着重发展的那个,正是数学的梦的典型——它还接上了“Motif 之梦”,为后者提供了一条新进路。写这份《草稿》时,我想起了这十四年里我连续做过的、为期最长的一次数学思考。它从1981年1月一直持续到6月,我给它取名叫《穿越伽罗瓦理论的长征》。一针引一线,我意识到:我这几年来时断时续地追着的那场白日梦,后来得了“远阿贝尔代数几何”之名,其实无非是伽罗瓦理论的一个延续,“伽罗瓦理论的终极归宿,而且想必是合乎伽罗瓦本人的精神的”。

当我写到被引的那句话所在的那一段、这种连续性在我眼前出现时,一阵喜悦穿过我,至今未散。那是在完全的孤独中做下来的工作所得到的报偿之一。它的出现,与不久前我从两三位同事兼旧友那里受到的、颇为冷淡的接待同样出乎意料——他们明明是“在行”的人,其中一位还曾是我的学生;我曾有机会跟他们谈起我正在发现的东西,那时我还“热着”,心里满是欢喜……

这让我想到:今天接续伽罗瓦的遗产,必定也要接受他当年那份孤独的风险。也许时代的变化比我们想的要小;不过,这份“风险”在我眼里往往并不成为威胁。我若有时因所爱之人的故作冷淡或轻蔑而难过、失落,那么相反,多年来孤独——数学的或别的孤独——从不曾让我觉得沉重。若有一位忠实的女友,我每次离开她后都不断渴望再回到她身边,那就是她!

\section{(8) 梦与证明}

但让我们回到梦,回到几千年来压在它身上的那道数学禁令。这也许是所有先验之见中最顽固的一条——它往往不明说,扎根在习惯里,裁决着什么东西“算数学”、什么不算。要过几千年,那些像孩童玩意一样简单
% ---------------- p.139 ----------------
而又无所不在的东西——某些几何图形的对称群、另一些图形的拓扑形态、数字零、集合——才获准进入圣殿\origpage{17}!

\origpage{18} 当我跟学生们谈球面的拓扑,谈由球面添上环柄所得的那些形态——这些东西并不叫小孩子吃惊,却叫他们不知所措,因为他们自以为知道“数学”是什么——我得到的第一个自发反应总是:这哪算数学!数学嘛,当然是勾股定理、三角形的高、二次多项式……这些学生并不比你我更蠢;他们的反应,和古往今来、直到今天的世上所有数学家的反应一样,除了毕达哥拉斯或黎曼这样的人,也许还有五六位别的。就连庞加莱——那可不是等闲之辈——也能用一套深有感触的“甲加乙”式哲学,“证明”无穷集合算不上数学!想必也曾有过一个时候,三角形和正方形算不上数学——那只是孩子或陶匠在沙上、在瓶罐的泥坯上画的图,别搞混了……

精神被自己的“知识”憋住,这种根子上的惰性当然不是数学家独有的。我有点偏离正题了:我要谈的是压在数学的梦身上的禁令,以及借由它、压在一切不以制成品的惯常面貌出现(不是开箱即用)的东西身上的禁令。我对别的自然科学所知的那一点,已足以让我估量出:同样严厉的禁令早就会把它们判给贫瘠,或判给乌龟般的进展——有点像中世纪,那时容不得对《圣经》的字句刨根问底。但我也很清楚:发现的深源,以及发现之进程在其一切本质方面,在数学里与在宇宙中我们的身体和精神所能认识的任何别的领域、别的事物里都是一样的。放逐梦,就是放逐源头——判它只能在暗处存在。

我还很清楚(这经验自我少年时代与数学初恋爱慕以来从未被推翻):在展开一种广阔或深刻的数学事物之景象时,正是这景象与理解的展开、这渐进的深入,始终先于证明,使证明成为可能,并赋予它意义。当一个情境——从最卑微到最宏阔——在其本质方面被理解了,对所理解之物的证明(以及对其他东西的证明)就像熟透的果子一样落下来。而像青果一样从认识之树上硬拽下来的证明,会留下一丝不满足的余味,一种渴望未得平息的失落。在我做数学家的一生中,有两三次,我不得不在别无办法时下定决心去拽果子,而不是摘果子。我不是说我做得不好,也不是说我后悔。但我所做得最好的、我所最爱的东西,我都是心甘情愿取来的,不是强夺来的。如果说数学给了我丰盈的快乐,在我成熟之年仍让我着迷,那靠的不是我从它那里硬拽来的证明,而是我在它身上感到的无穷奥秘与完美和谐——它随时准备向一只爱它的手、一道爱它的目光显露自身。

\section{(9) 受欢迎的异乡人}

我觉得时候到了,该谈谈我与数学家世界的关系。这与我与数学的关系完全是两回事。后者从我年少时起就存在,而且很强,远在我想到世上竟有一个数学家的世界、一个数学家的圈子之前。那是一整个复杂的世界,有它的学会、它的期刊、它的聚会、讨论会、大会,有它的头牌女伶和苦力,有它的权力结构,有它的灰衣主教,还有同样灰扑扑的一大群任人使唤、任人摊派的人,为论文或文章发愁;也有另一些更少见的人,他们有的是手段和想法,却
% ---------------- p.140 ----------------
处处碰壁\origpage{19},绝望地找不到一位有权有势者撑腰——那些忙忙碌碌、叫人畏惧的强人握有一种魔法:让一篇文章得以发表……

\origpage{20} 我是在1948年、二十岁上到巴黎时,才发现有一个数学世界的;我那只瘪瘪的手提箱里,装着蒙彼利埃大学的理学学士学位,还有一部行距排得极密、双面书写、不留页边的手稿(纸很贵!),那是我三年孤独思考的记录,思考的东西(后来我才知道)当时早已以“测度论”或“勒贝格积分”之名广为人知。因为从没遇见过别人,我直到进京城那天都还以为,世上只有我一个人在“搞数学”,因此也就是唯一的数学家。(在我眼里这两者是一回事,而且直到今天也还多少是。)我曾玩弄那些我称之为可测的集合(其实也没遇到过一个不可测的集合……),玩弄几乎处处收敛,却不知道拓扑空间是什么。有十来个互不等价的“抽象空间”与紧性的概念,我一直没搞清,那是从一本小册子里捞来的(作者好像叫阿佩尔,收在《科学与工业时事》丛书里),这本小册子是我碰上的——天知道怎么碰上的。那时我还没听人说过——至少没在数学的语境里听过——群、域、环、模、复形、同调(还有别的,就不数了)这类古怪或野蛮的词;忽然之间,它们不打招呼,一齐朝我压下来。这一击不轻!

我之所以在这一击之下“活了下来”,还继续搞数学、甚至以此为业,是因为在那遥远的年头,数学世界还不太像它后来变成的样子。也可能是我运气好,落进了这个未曾料想的世界里一个格外好客的角落。我带着蒙彼利埃学院一位教授的一封含糊的推荐信,那是苏拉先生(跟他的同事们一样,他也没怎么在课堂上见过我!),他当年是嘉当的学生(父亲还是儿子,我已说不太准了)。因为埃利·嘉当(Elie Cartan)当时已经“退出赛场”,他儿子亨利·嘉当(Henri Cartan)便成了我有幸遇见的头一个“同类”。我当时并不知道这是多么吉利的兆头!他接待我时的那份礼貌中透着善意,那是他的标志,几代高师学生都熟悉——他们有幸在他手下初试身手。从他当时给我指点学业方向的那些建议看,他大概没有意识到我无知到了何等地步。无论如何,他的善意显然是冲着人来的,不是冲着行李或可能有的天分,也不是(后来)冲着名声或名气。

接下来那一年,我成了“高师”嘉当一门课上的常客(讲流形上的微分形式体系),我拼命跟;也常去“嘉当讨论班”,目瞪口呆地看着他和塞尔辩论,满嘴“谱序列”(呜!),满黑板的图(叫做“图表”),画满了箭头。那是“层”、“甲壳”以及一整套军火库的英勇时代,那些东西的意义我完全摸不着边,可我还是勉勉强强逼着自己把定义和陈述吞下去,把证明一一验过。在嘉当讨论班上,谢瓦莱、韦伊也定期露面;而到了布尔巴基讨论班的日子(参加者和听众加起来,撑死也就二三十人),就看见这个有名的布尔巴基帮的其余成员——迪厄多内、施瓦茨、戈德芒、德尔萨特——像一伙有点吵闹的哥们儿那样闯进来。他们彼此以“你”相称,说着一种我几乎完全听不懂的共同语言,烟抽得凶,也爱笑;只差几箱啤酒,气氛就齐全了——而啤酒是由粉笔和海绵代替的。那气氛与法兰西公学院勒雷的课截然不同(讲肖德尔关于无限维空间中拓扑度的理论,可怜的我!),我是听嘉当的建议去听的。我曾到法兰西公学院找过勒雷先生,问他(如果我没记错)这门课要讲什么。他可能给我的解释,我一点不记得了,
% ---------------- p.141 ----------------
\origpage{21} 也不记得自己有没有听懂什么——只记得:在那里我也感到一种善意的接待,那是给随便哪个初来乍到的外乡人的。肯定正是这一点,而非别的什么,让我去上那门课,并且像在嘉当讨论班那样硬着头皮跟下去——尽管勒雷在那里讲的东西,我当时几乎完全不懂。

奇怪的是:在这个世界里我是新来的,它的语言我几乎不懂,说得更差,可我并不觉得自己是个外乡人。虽然我几乎没机会(而且当然没本事)跟韦伊、迪厄多内这样快活的家伙说话,或跟嘉当、勒雷、谢瓦莱这样气度更雍容的先生们说话,我却觉得自己被接纳了,几乎可以说:是他们中的一个。我不记得有哪一次被这些人中的某一位居高临下地对待,也不记得有哪一次,我的求知欲——以及后来,又一次,我发现的喜悦——被一份自得或一份轻蔑挡了回去\fnoteP{3}{(5)}{这个事实尤其值得注意:直到1957年前后,布尔巴基小组的成员中不止一位对我还颇有保留,而他们最终把我增选进去时,我相信也带着几分不情愿。一句善意的玩笑话把我归入了“危险专家”(指泛函分析)之列。我有时感到嘉当身上有一种更深的、未说出的保留——有那么几年,我一定给了他这样的印象:一个热衷于空洞浮泛的推广的人。他曾惊讶地发现,我为布尔巴基写的那一篇(也是唯一一篇)稍长的稿子(论流形上的微分形式体系)里,竟有一点相当实在的想法——而我提议由我来写时,他并不怎么热心。(多年以后,在从凝聚对偶的观点发展留数形式主义时,这一想法又派上了用场。)况且在布尔巴基的会议上,我多半是跟不上的,尤其是逐篇共同宣读稿子的时候,我实在没法跟上宣读与讨论的速度。可能我确实不太适合集体工作。不管怎么说,我在融入共同工作上的这份困难,或我因别的缘故在嘉当和别一些人那里可能引起的保留,都不曾在任何时候给我招来讥讽或碰钉子,或哪怕一丝居高临下——至多在韦伊那里有过一两次(此人确实是个特例!)。嘉当对我的态度,从未有一刻离开过一贯的和善,其中带着亲切,也带着他特有的那点幽默——在我,那点幽默始终与他这个人分不开。} (5)。若不是这样,我也不会像人们说的那样“成为数学家”——我会另选一行,在那里我能施展自己,而不必去面对轻蔑……

尽管“客观上”我是这个世界的异乡人,一如我在法国是个异乡人,却有一条纽带把我和这些出身别的阶层、别的文化、别的命运的人连在一起:一份共同的激情。我怀疑,在我发现数学家世界的那个关键年头,他们中没有一个人——连嘉当也没有,我勉强算他的学生,可他还有许许多多别的学生(而且是不那么跟不上的学生!)——在我身上察觉出那份同样盘踞在他们心中的激情。在他们眼里,我大概只是课堂和讨论班一大群听众中的一个,记着笔记,显然不太跟得上。如果说我多少有点与众不同,那就是我不怕提问——而那些问题多半只是暴露出我在语言和数学两方面都无知得惊人。回答可能简短,甚至带着惊讶,但当时那个目瞪口呆的冒失鬼,从未碰过钉子,从未被“放回原位”——无论在布尔巴基那伙人不拘礼数的圈子里,还是在法兰西公学院勒雷课堂上更严峻的氛围里。那些年里,自从我口袋里揣着一封给埃利·嘉当的信来到巴黎,我从未觉得自己面对着一个宗派、一个封闭乃至敌意的世界。如果说我领教过、深谙面对轻蔑时那种内心的紧缩,那不是在那个世界里;至少不是在那时。对人的尊重,是我在那里呼吸的空气的一部分。不必去挣得尊重,不必先证明自己才能被接受、被和气地对待。说来也许奇怪:只要是一个人,只要有一张人的面孔,就够了。

% ---------------- p.142 ----------------
\section{(10) “数学共同体”:虚构与现实}

\origpage{22} 因此毫不奇怪:也许就在那一年,在我内心深处,我已感到自己是这个世界的一员;而在随后的年月里,这种感觉越来越清楚。我喜欢用“数学共同体”这个名称来指它,这名字对我满载着意义。在写下这些字句之前,我从未有机会细想过自己给这名字赋予了什么意思,尽管我在很大程度上认同于这个“共同体”。现在清楚了:它对我而言,不外是那个善意接待我、把我当作他们中一员的世界的理想延伸——在空间上和时间上的延伸;而且是一个与我生命中最主要的激情之一相连的世界。

\origpage{23} 我逐步认同的这个“共同体”,并不是对最初接待我的那个数学圈子所作的完全虚构的外推。最初的圈子一点点扩大了;我的意思是:在最初接触之后的十年、二十年里,由于共同感兴趣的主题和人与人之间的投缘,我经常来往的数学家圈子不断扩大。具体说来,就是同事与朋友的圈子,或者不如说那个同心结构:从与我关系最近的同事(先是迪厄多内、施瓦茨、戈德芒,后来尤其是塞尔,再后来是安德烈奥蒂、朗、泰特、扎里斯基、广中平祐、芒福德、博特、迈克·阿廷这样的人,还不算布尔巴基小组的人——那个小组也在一点点扩大——以及六十年代起投到我这里来的学生们……),直到我在此处彼处有机会遇到、由或强或弱的投缘而或紧或松地联系着的其他同行——总之,正是这个由相遇与投缘偶然凑成的微观世界,构成了那个于我充满温度与回响的名字的具体内容:数学共同体。当我把它认同为一个活生生的、温暖的实体时,实际上我认同的就是这个微观世界。

直到1970年那个“大转折”之后——我该说,那第一次觉醒之后——我才意识到:这个舒适可亲的微观世界只占“数学世界”极小的一部分;而我乐于赋予那个世界、那个我一直不去了解、从没想到要去关心的世界的那些特征,都是虚构的特征。

况且在这二十二年间,这个微观世界本身也换了面孔,而它周围的世界也在变。我自己当然也变了——年复一年,不知不觉,我和我周围的世界一样变了。我不知道我的朋友和同事们是否比我更觉察到这变化——周遭世界里的、他们自己的微观世界里的、以及他们自身之中的变化。我也说不清这奇怪的变化是何时、如何发生的——它想必是悄悄来的,蹑手蹑脚:有名望的人是叫人怕的。我自己就叫人生畏——不是叫我的学生或朋友怕,也不是叫那些亲自认识我的人怕,至少是叫那些只通过名声认识我、而自己并未感到有相当的名声庇护的人怕。

我是在将近十五年前那次“觉醒”之后,才意识到数学世界里(在其他科学圈子里同样如此,如果不是更甚)盛行的恐惧的。在此前的十五年里,我逐步地、不知不觉地进入了“大老板”的角色,进入了数学界《名人录》的世界。也同样是不知不觉,我成了这个角色的囚徒,它把我与所有人隔开,只除了少数几位“同侪”和少数几个学生(而且还……)——那些确实“有心气”的人。只有当我走出这个角色,环绕着它的恐惧至少有一部分才落了下来。那些在我面前多年沉默的舌头,松开了。

它们带给我的见证,不只是恐惧的见证,也是轻
% ---------------- p.143 ----------------
\origpage{24} 蔑的见证。尤其是身居其位的人对别人的轻蔑——一种唤起恐惧、也滋养恐惧的轻蔑。

我几乎不曾体验过恐惧,却实实在在体验过轻蔑——那是在一个人的存在与生活分量很轻的年月。我本来乐于忘掉轻蔑的时代,可它偏偏来向我问好了!也许它从未停止过,只是我满足于换了个世界(在我看来如此),满足于把目光移开,或者干脆:除了那些激动人心、没完没了的数学讨论之外,装作什么也看不见、什么也听不见?那些日子里,我终于肯去了解:轻蔑就在我周围到处盛行,在我选为自己的世界、我所认同、我所背书、并且曾娇养我的那个世界里盛行。

上面这些话也许给人一个印象:我被那些几乎一夜之间涌向我的见证搅得心神不宁。其实不然。这些见证是在一个仍然浮浅的层面上被接收的。它们只是加在我刚得知的另一些事实上,或加在我早已知道、却一直避免去注意的那些事实上。今天,我会这样表述当时学到的教训:“科学家”们,从最显赫的到最无名的,都是些和所有别人一模一样的人!我曾乐于想象“我们”是更好的一类,我们身上多点什么——我足足花了一两年,才摆脱这个顽固透顶的幻觉!

\section{(11) 与克劳德·谢瓦莱相遇——或自由与善意}

在帮助我做到这一点的朋友中,只有一位属于那个我刚离开、不打算再回去的圈子\fnoteP{4}{(6)}{\emph{我在“生存与生活”的朋友们}\par 在这些朋友中,我大概还该算上皮埃尔·萨米埃尔(Pierre Samuel)——我先前主要是在布尔巴基里认识他的,与谢瓦莱一样;而他(和谢瓦莱一样)在“生存与生活”小组里起了重要作用。我不觉得萨米埃尔有那么迷恋那种关于科学家优越性的幻觉。我感到,他带来最多的,是他在共同工作、讨论和待人接物中投入的常识和那份笑呵呵的好脾气,还有他优雅地在一个倾向于激进分析与激进抉择的团体里扮演“讨厌的改良派”的角色。在我退出“生存与生活”之后,他又在那里待了一阵,充当同名公报的主编;当他感到自己的存在对这个团体已不再有用时,便体面地离开了(去了“地球之友”)。\par 萨米埃尔和我出身于同一个狭小的圈子,这并不妨碍他成为那些沸腾年月里的朋友之一——我相信我从他们那里学到了些东西(尽管我是个糟糕的学生……)。他的为人方式,和谢瓦莱的为人方式一样——尽管二人毫不相像——对我那些“功绩主义”的倾向来说,是比最锋利的分析更好的解药!\par 现在我觉得,对那个时期所有让我学到点东西的朋友,我学到的主要不是通过解释、讨论之类,而是通过他们各自的为人方式、他们那与我不同的敏感——而“某种东西”最终就这样传达了过来……说到这里,我特别记得的,除了谢瓦莱和萨米埃尔,还有德尼·盖吉(Denis Guedj,他对“生存与生活”小组有很大的影响力)、达尼埃尔·西博尼(Daniel Sibony,他一直与这个小组保持距离,同时从角落里用一半鄙夷、一半打趣的眼光注视着自己的演变)、戈登·爱德华兹(Gordon Edwards,他是1970年6月在蒙特利尔那场“运动”诞生的共同行动者,并且多年间付出了惊人的精力,维持着《生存与生活》公报的英文“美国版”)、让·德洛尔(Jean Delord,一位与我年纪相仿的物理学家,精细而热情,他喜欢我,也喜欢“生存”这个小圈子)、弗雷德·斯内尔(Fred Snell,另一位在美国布法罗的物理学家, 1972年我在那里逗留了几个月,曾在他乡间的家中做客)。\par 在所有这些朋友中,五位是数学家,两位是物理学家,而全都是科学家——这似乎表明,那些年里与我最接近的圈子仍然是科学家的圈子,尤其是数学家的圈子。} (6)。他就是克劳德·谢瓦莱。他不发议论,对我的议论也不感兴趣,然而我想我可以说:我从他那里学到的东西,比我刚才说的那些更重要、也更隐蔽。

在我相当经常地与他来往的那段日子(“生存”小组的日子,他带着打了折扣的信念加入了这个小组),他常让我不知所措。我说不清是怎么回事,但我感到他握有某种我抓不到的
% ---------------- p.144 ----------------
\origpage{25} 认识——对某些想必极其简单、极其本质之物的理解;这些东西当然可以用简单的话说出来,但理解并不会因此就从一个人“传”到另一个人那里。我现在明白,他和我之间有一种成熟程度上的差别,这使我常常在他面前觉得自己站得不稳,像在进行一场聋子的对话——而这并不是因为彼此缺少好感或敬意。

虽然他没有用这样的词说过(就我记忆所及),但在他看来想必很清楚:我当时——或独自、或凭“生存”小组(后来成了“生存与生活”)内部共同的思考与活动的逻辑——达到的那些“质疑”(关于“科学家的社会角色”、关于科学等等),骨子里仍然是浮浅的。它们当然涉及我生活的世界,甚至涉及我在其中扮演的角色——但并没有真正深刻地牵动我。在那沸腾的年月里,我看待自己的眼光一丝未变。我并不是在那时才开始认识自己的。只是六年之后,我才生平第一次摆脱一个顽固的幻觉——不是关于别人或周围世界的幻觉,而是关于我自己的。那是另一次觉醒,分量比预备了它的第一次更大。它是一连串相继“觉醒”的“瀑布”中最初几次之一;我希望,这瀑布在我余年里还会继续流下去。

我不记得谢瓦莱曾在什么场合提到过认识自己,或者更确切地说,“发现自己”。不过回过头看,很清楚:他准是很早以前就已开始认识自己了。他有时也谈起自己,只在谈这谈那时说上几个字,那份朴素叫人不知所措。他是我没听见说过一句套话的两三个人之一。他话不多,而他说的话所表达的,不是他采纳并据为己有的观念,而是对事物的亲身感知与理解。所以他常叫我不知所措——早在我们还在布尔巴基小组里碰面的时候就是如此。他所说的,往往冲撞了我所珍视的、并因此被我当作“真”的那些看法。他身上有一种我所缺乏的内在自主,我在“生存与生活”的日子里开始隐约察觉它。这种自主不属于智性、不属于话语的层面。它不是一件可以像观念、观点之类那样“采纳”的东西……幸好,我从未想过要把在别人身上察觉到的这份自主“据为己有”。我得找到我自己的自主。这同时意味着:我要学会(或重新学会)做我自己。但在那些年里,我对自己的不成熟、对自己缺乏内在自主,毫无觉察。我若最终发现了这一点,那么与谢瓦莱的相遇,必定是那些在我身上默默发酵的酵母之一——那时我正投身于一些大计划。播下这酵母的,不是议论,也不是言语。要播下它,只需我在路上偶然遇见的某个人不讲议论,只满足于做他自己。

我觉得,在七十年代初,当我们因《生存与生活》公报的出版而定期见面时,谢瓦莱曾不加坚持地试图向我传达一个消息,而我当时太笨手笨脚,抓不住它,或太埋头于我的活动任务。我隐约意识到,他有关于自由——关于内在自由——的东西要教我。我那时倾向于靠一些宏大的道德原则运转,而且从《生存》最初几期起就开始吹那把小号,觉得理所当然;而他对说教式的话语有一种特别的反感。我想,在《生存》创办之初,这是我最摸不透他的地方。在他看来,这种话语只是一种强制的尝试,叠加在窒息人的无数别的外在强制之上。当然,一个人尽可以花一辈子去讨论这样
% ---------------- p.145 ----------------
一种看法,讨论正反两面\origpage{26}。它完全冲撞了我的看法——那看法(可想而知)是由最高尚、最普遍的感情驱动的。我很难过:我无法理解,谢瓦莱——我最敬重的人,我多少把他当作战友的人——竟然会带着一种狡黠的乐趣,不去分享这些感情!

\origpage{27} 我不明白:真相、事物的实际,既不是善良感情的问题,也不是观点或偏好的问题。谢瓦莱看见了一件东西——那东西再简单、再真实不过——而我没有看见。并不是他在什么地方读到过它;看见一件东西,和读到关于它的东西,二者毫无共同之处。人勉强可以用手(盲文)或耳朵(若有人读给你听)去“读”一段文字,但要看见那件东西本身,只能用自己的一双眼睛。我不认为谢瓦莱的眼睛比我的好。但他用它们,我不用。我被我的善良感情之类的东西占得太满,没有闲工夫去看我的善良感情和原则对我自己、对别人——首先是对我自己的孩子们——起了什么效果。

他一定看得很清楚:我常常不用自己的眼睛,甚至一点也不想用。奇怪的是,他从未对我透露过这一点。还是他说了,只是我没听见?还是他忍住了,认为那是白费力气?又或者他压根儿没想到要说——说到底,我用不用自己的眼睛,是我的事,不是他的事!

\section{(12) 功绩与轻蔑}

我想借我自己有限的经验细看一下:轻蔑是何时、如何进入数学家世界的,尤其是在那个由同事、朋友和学生构成的“微观世界”里——它已成了我的第二祖国。同时我也想看看,在这场转变中,我这一份占了多少。

我想我可以毫无保留地说: 1948–49年间,在我前面谈过的那个数学家圈子里(对我来说,它的中心是最初的布尔巴基小组),我没有遇到一丝一毫的轻蔑,甚至没有一丝一毫的鄙夷或居高临下——无论对我自己,还是对任何一个来此学习数学家行当的年轻人,无论法国人还是外国人。在那里起旗舰作用的人物——凭地位或声望,如勒雷、嘉当、韦伊——我并不怕他们,我想我的同学里也没有人怕他们。勒雷和嘉当很有“体面先生”的派头,另当别论;至于其余那些随随便便闯进来、像哥们儿一样跟嘉当“你你我我”、显然又“在行”的家伙,我甚至花了相当一段时间才明白:他们每一个都和嘉当本人一样是大学教授,绝不像我那样挣一天吃一天,而是领着在我看来属于天文数字的薪水,而且还是个有国际影响、国际听众的数学家。

听从韦伊的建议,我在南锡度过了随后的三年。那里当时算是布尔巴基的一个总部:德尔萨特、迪厄多内、施瓦茨、戈德芒(稍后还有塞尔)都在大学里任教。跟我一起在那里的只有四五个年轻人(我记得有利翁斯、马尔格朗热、布吕阿、贝热,若有记错请原谅),所以远不像在巴黎那样“淹在人堆里”。气氛也因此更亲切,大家彼此都认识,我想彼此都以“你”相称。可我翻检记忆时,我亲眼看见一个数学家毫不掩饰地轻蔑对待一个学生的头一遭、也是唯一一遭,恰恰就在那里。那可怜人是当天从另一座城市赶来的,要跟他的老板一起工作。(他当时应该在准备博士论文,后来还体面地通过了,而且此后也
% ---------------- p.146 ----------------
有了相当的名声,我想。)\origpage{28} 那场面让我相当震惊。要是有人敢用这种腔调跟我说哪怕一秒钟的话,我立马就把他摔门轰出去!可这一回,我熟的只是那位“老板”,甚至跟他称兄道弟,而不是那个学生——我只是见过他。我这位年长的同行,除了广博的学养(不只是数学的)和锋利的头脑,还有一种斩钉截铁式的权威感,在当时(而且在之后相当长一段时间里,直到七十年代初)都令我敬畏。他对我有某种影响力。我不记得自己有没有就他的态度问过他什么;我只记得我从那场面里得出的结论:这个可怜学生一定是差劲透了,才该受这样的对待——诸如此类。我当时没有对自己说:这学生若是真的差劲,那该劝他去做别的,该停止跟他合作,但无论如何不该轻蔑地对待他。我认同了像这位显赫前辈那样的“数学强者”,而以那些“废物”为代价——似乎轻蔑他们是正当的。于是我走上了那条早已画好的路:与轻蔑默契,因为这对我方便,它还突出了这一点——我,被接纳入有功者、“数学强者”的兄弟会了!\fnoteP{5}{(7)}{上面这一段,是整篇导论里在我的原稿上涂改得最厉害、增补最多的一段。对那件事的描述、用词的选择,最初都是顶着来的,逆着潮流——显然有一股力量在推着人快点把那件事带过去,好像只为对良心有个交代,好“转到正经事上去”。这些都是一份抵抗的熟悉征兆——这里抵抗的是对这段插曲的澄清,以及对它作为某种内心态度之揭示者的意义的澄清。情形与这篇导论开头(第2段)所描述的十分相似:在数学工作中发现一处矛盾及其意义的那个“关键”时刻——那时起作用的是精神的惰性,是它不肯与一种错误或不充分的理解(但其中根本没有我们这个人牵涉在内)分离的厌恶。后者却是主动的,必要时还能别出心裁,做到无水也能把鱼淹死;而我刚才说的惰性只是一种被动的力量。在现在这个情形里,比在数学工作的情形里更甚:刚刚以其全部的简单、全部的自明显现出来的发现,紧接着就是一阵卸下重负的轻松感,一种解放感。这不止是一种感觉——不如说是对刚刚发生之事的一种敏锐而感激的知觉,那是一次解放。} (7)

当然,我不会比任何人更清楚地对自己说:那些试着搞数学却搞不成的人,活该被轻蔑!如果我在那个时候或任何时候听见有人说这种话,我会漂亮地驳他一顿,真心为如此惊人的精神无知而难过。事实是:我泡在暧昧里,同时在两张互不相通的牌桌上玩:一边是漂亮的原则与感情,另一边是:可怜的家伙,得差劲成什么样才落得这般对待(言下之意:这种事肯定不会落到我头上,那当然!)。

最后我觉得,我讲的那桩事件,尤其是我在其中扮演的那个(看来无伤大雅的)角色,其实是我身上一种暧昧的典型;这暧昧在此后二十年里跟了我整个数学家的生涯,直到1970年“觉醒”之后才消散\fnoteP{6}{(8)}{正如下文将变得清楚的,这种暧昧丝毫没有“在1970年觉醒之后消散”。这里有一个“自我”典型的战略撤退:它把“觉醒之前”那段时期丢给盈亏自负,而那段时期随即成了划分出一个无可指摘的“之后”的分界线!} (8)——而直到今天,写下这些字句的时候,我才清楚地察觉它。也真可惜,我当时竟没有察觉。也许对我来说时机尚未成熟。总之,那些当时传到我耳中的、关于轻蔑之统治的见证——对那统治,我选择了闭上眼睛——并没有把我个人牵涉进去,也没有把我那亲爱的微观世界里离我最近的那部分同行和朋友中的任何一位牵涉进去\fnoteP{7}{(9)}{这并不完全准确:在我最亲近的同行中至少有一个例外,后文会看到。这里有一种记忆特有的“懒”,它往往倾向于把那些与由来已久的、根深蒂固的看法“不合”的事实“随手放过”。} (9)。它们更像是这个调子:唉!竟然不得不%
% 原句在此页末尾中断, 下文接 part1-ch06b.tex (p.147)

% ================================================================
%  第 6 章 梦与做梦者 后半 (原书 p.147-152, 6.9-6.11)
%  源文: _work/src1/task-ch06b.md
%  注意:首段紧接 part1-ch06a.tex 末句 “唉!竟然不得不”, 两者之间不可插入空行
% ================================================================
% ---------------- p.147 ----------------
学到(或者:让人领教)这种事——谁会相信呢——真得是个混账(我差点要说:废物,对不起!),才会那样对待活生生的人!说到底,这跟另一副调子也差不了多少:只消把“废物”换成“混账”,把“让人那样对待”换成“那样对待”,戏法就变完了!而荣誉嘛,当然,毫发无损——对善事业的旗手来说!

\origpage{29} 从这里清楚显出的,是我与种种轻蔑态度的默契。它至少上溯到五十年代初,也就是我在嘉当和他朋友们那里受到善意接待之后的那些年份。后来,当轻蔑在四面八方变得司空见惯时,我之所以“什么也看不见”,是因为我不想看见——并不比那桩孤立而格外扎眼的事里更想看:那一回,真得下足功夫,才能装作既看不见、也觉不出!

这份默契与我新的身份紧密共生——一个受尊敬的群体成员的身份,有功者、数学强者的群体。我记得我特别满意,甚至骄傲:在我为自己选定、又把我增选进去的那个世界里,起作用的不是社会地位,甚至(才不是呢!)也不单是名声——名声还得是挣来的;你尽可以是大学教授、院士或随便什么,可你要只是个平庸的数学家(可怜的家伙!),那你就什么也不是;唯一算数的是功绩,是深刻而独到的想法,是技术上的娴熟,是广阔的视野,诸如此类!

这套功绩主义,我曾毫无保留地认同于它(而它当然始终隐而不宣、不曾明言),可它在我身上到底还是挨了狠狠一击——就在我说过的1970年那场著名觉醒之后的日子里。而且我也不能肯定,它从那一刻起就消失了、没留下痕迹。要做到那一点,大概得靠我清清楚楚地在自身中把它识别出来;可我当时似乎主要是在别人身上谴责它。说起来,最早提醒我注意这套意识形态的,正是谢瓦莱,还有德尼·盖吉(我也是通过“生存”认识他的)——他们管它叫“功绩主义”(méritocratie),或者类似的名字——以及它里面那种暴力、那种轻蔑。谢瓦莱对我说(那大概是我们头一次在他家里见面、谈“生存”的时候),正是因为这个,他才再也受不了布尔巴基的气氛,索性不再踏进去。回想起来,我确信他当时一定早已看出:我实实在在是这套意识形态的一分子,也许甚至看出它在我的某些角落里还留着痕迹。但我不记得他曾有过任何暗示。也许这一回,他仍然宁愿让我自己去把那些他替我勾出的 i 点上点——而我一直等到今天才把它们点上。晚做总比不做强!

\section{(13) 力与厚度}

\origpage{30} 我讲的那桩事,很可能也标志着我内心一次倒向的时刻:转向对功绩兄弟会或多或少无条件的认同,而以那些被视为废物的人为代价——或者,按早几代人的说法,只消说“没有天分”(这个词到我那时候已经不流行了):那些平庸、暗淡的人——顶多算是“共鸣箱”(韦伊在什么地方这样写过),供真正要紧的人的大想法在其中回响……我的记忆常常像个掘墓人,连当时耗掉大量心理能量的片断也一并埋掉;可它偏偏留下了这一段,而这一段跟任何别的记忆都没有直接关联,面目又那么平淡无奇——单凭这一点,那个“倒向”当时确曾发生的感受,就变得可信了。

不到五年前的一次沉思中,我终于意识到:那种“我们,伟大而高贵的精神……”的意识形态,以一种格外极端、格外凶猛的形式,从我母亲童年起就在她身上肆虐,并且主宰着她与他人的关系——她乐于从自己的高处俯视那些人,带着一种往往是鄙夷的、甚至轻蔑的怜悯。
% ---------------- p.148 ----------------
\origpage{31} 再说,我对我的父母怀着毫无保留的敬佩。在那个著名的“数学共同体”之前,我认同的第一个、也是唯一一个群体,就是缩小到母亲、父亲和我三个人的家庭——而我有幸被母亲认可为配得上他们做我的父母。这就是说,轻蔑的种子想必从我的童年起就已播进我这个人里了。也许时机已经成熟,可以追一追这些种子在我的童年和成年生活中经历的种种变故,以及其中一些种子长成的收成——幻觉的收成、孤立的收成、冲突的收成。但这里不是做这件事的地方,此处我追着一个更有限的目标。我想我可以说:这种轻蔑态度在我一生中,从未达到我在母亲一生中见到的那种激烈、那种毁灭性的力量(那是我肯费心去看我父母一生的时候看到的——在母亲去世二十二年之后,在父亲去世三十七年之后)。但现在正是时候,过了这村就没这店:在这里至少仔细考察一下,这种态度在我的数学生涯中占了什么位置。

在这之前,为了把上一段讲到的那桩事放进它的一般背景里,我想强调这样一个事实:在我的五十年代记忆里——甚至更晚的记忆里——它完全是孤立的。即便在今天,虽然我发觉在我曾经所属的那个圈子里,某些最基本的礼貌和对他人的尊重,有时正以一种叫人错愕的方式被侵蚀\fnoteP{8}{(10)}{比方说,我已数不清有多少封信——谈的既有数学问题,也有实际问题或私事——寄给我当作朋友的同事或旧学生,却从未收到回音。看来这并非只是专门留给我一个人的特殊待遇,而确实是一种风气变化的迹象——据一些同方向的反响来看。(诚然,那些反响涉及的情形是:寄数学信件的人,收信的那位知名数学家并不认识他……)} (10),老板对学生那种直接而不加掩饰的轻蔑,也还应当是相当罕见的事。至于五十年代,我很少有记忆指向那种情形:一个有名望的人物周围笼罩着恐惧,或者有谁摆出轻蔑、哪怕只是鄙夷的态度。要往这个方向翻找,我可以说:头一回被请到南锡迪厄多内家做客时——他对我一向有着那份体贴入微的和气——这个文雅亲切的人谈起他的学生来,那副样子真叫我有点发懵:简直全是些蠢货!给他们上课是件苦差,而他们显然什么也听不懂……1970年以后,我听到从阶梯教室那一侧传来的回声,才知道迪厄多内确实让学生们害怕。可是,尽管他以见解果决、并且常常以雷鸣般的直率把它们讲出来而闻名,我却从未见他有过伤人、羞辱人的举动——即使当着他不怎么看重的同事的面,即使在他那些传奇式的雷霆大怒之际;而那怒气,来得快,消得也一样快、一样轻易。

我不认同迪厄多内关于他学生的那些说法,可也没有同他的态度拉开距离——那态度被呈现为世上最理所当然的事,对一个热爱数学的人来说几乎是不言自明的。有这位年长同行满怀善意的权威作后盾,那种态度当时在我眼里,至少是人们在面对学生和教学任务时可以有理由采取的态度之一。

看来,迪厄多内和我一样,都浸透了这套功绩主义;但当我们面对一个血肉之躯时,这种主义的隔离效果在很大程度上就被抵消了:单是那个人的在场,就默默地向我们提醒一些比所谓“功绩”更本质的现实,并重新接上一根被遗忘的纽带。我们的大多数同事或朋友想必也一样——他们浸染那种如此普遍的优越感,并不比迪厄多内或我少。今天,他们中的许多人肯定还是这样。

% ---------------- p.149 ----------------
\origpage{32} 韦伊同样有让学生害怕的名声;而在五十年代我的微观世界里,他是唯一一个让我觉得连那些地位(或者干脆说性情)比较谦卑的同事也怕他的人。他时不时会摆出那种不容回嘴的居高临下,足以叫最牢靠的自信也发慌。我的敏感也在其中起了作用,于是有那么一两次,我们闹了阵短暂的别扭。我在他的做派里没有察觉出一丝轻蔑的意味,也没有察觉出存心伤人、压人的意图;那更像是被惯坏的孩子的作派,以叫人难堪为乐(有时还带点促狭),仿佛以此使自己确信他握有某种权力。再说,他对布尔巴基小组有一种真正惊人的支配力;有时他给我的印象是,他拿指挥棒指挥着它,有点像幼儿园女教师指挥一群乖孩子。

在五十年代,我只记得另外一次,感到轻蔑以一种粗暴、不加掩饰的方式表达出来。它来自一位外国同事兼朋友,年纪跟我差不多。他的数学力量非同寻常。几年前,这股力量已经相当明显了,而我当时却被他对那位大教授的顺从(在我看来近乎谄媚)所触动——那时他还只是那人的一名谦卑助手。他出众的才能很快为他赢得了国际声誉,以及一所格外显赫的大学里的一个关键职位。他在那里统治着一小支助手兼学生的队伍,那副绝对的派头,显然跟他当年的老板统治他和他的同伴们时一样。我问他(如果我记得不错)有没有几个学生(言下之意:跟着他做出好东西的学生),他带着一副假作随意的神气(我把它译成法语)答道:“十二件!”——所以“件”就是他用来指他的学生和助手的词。一个数学家同时有这么多学生在他指导下做研究,确实少见——而我的这位对话者想必暗自引以为傲,又想用那副漫不经心的神气把它藏起来,仿佛在说:“哦,不过十二件,连提都不值一提!”那大概是1959年前后。我当时想必已经有了一层不错的甲壳,可我还是感到一阵恶心!我准是当场就用某种方式对他说了,而且我不认为他为此记恨我。也许他与学生的关系并不像他那说法所能让人设想的那样阴森(我没有得到过他任何一位学生的证词),而他只是落进了自己那孩子气的欲望的陷阱:想在我面前炫耀他全部的荣光。回头看,我看得出这桩事必定是我们关系的一个转折点,而先前的那些关系是友谊的关系——我在他身上感到某种脆弱,还有一种细腻,它们吸引着我亲切的同情。这些品质被磨损了、被腐蚀了,被他那位重要人物的地位——受人敬佩、叫人害怕的地位。这桩事之后,我心里对他留下了一种不适——说到底,我觉得自己不属于他那个世界……

\origpage{33} 然而我们确实属于同一个世界——而且,我并不比他更觉察到:我也在变厚。关于这一点,我留下了一段鲜明的记忆,那是在1958年的爱丁堡国际数学家大会上。从上一年的黎曼–罗赫定理工作起,我就被晋升为大明星了;而且(当时我并不曾用明确的话对自己说过)我也是大会的明星之一。(我在会上作了一个报告,讲同年概形理论的强劲起步。)希尔策布鲁赫(当日的另一位明星,带着他自己的黎曼–罗赫定理)致开幕辞,为的是向那年即将退休的霍奇致敬。说到某处,希尔策布鲁赫暗示:数学主要靠年轻人的工作,而不是靠年长数学家的。这在年轻人占多数的大会厅里激起了一片轰然的赞同。我满心欢喜,当然也完全同意:我刚满三十岁,还满可以算年轻,而世界是属于我的!一兴奋,我准是大声叫喊,还在桌子上猛捶了好几下。碰巧我坐在霍奇夫人(Lady Hodge)旁边——她就是那位
% ---------------- p.150 ----------------
\origpage{34} 本该在这次场合受到致敬的卓越数学家的妻子,而他今年就要退休了。她睁大眼睛转向我,对我说了几句话,我已不记得了——但我准是从她惊讶的眼睛里看到了映出的那份没有分寸的厚度:它刚刚就在这位已近暮年的女士面前毫无节制地发作了一通。那时我感到了一点什么,而“羞耻”这个词给它的形象也许是走了样的——那更像是一条朴素的真相,关于当时的我这个人。那天,我大概再没有在桌子上猛捶过。

\section{(14) 恐惧的诞生}

我猜想,正是在那个时候——当我(并非刻意去求)开始被看作数学界的一个明星的时候——某种恐惧也必定开始围绕着我这个人,对许多我不认识或不太认识的同事来说。我是猜的:我无法用一个确切的记忆把它定位,无法用一幅曾击中我、并留在我记忆里的画面来定位,就像前面讲的那桩事那样(那想必标志着我头一回在接纳我的圈子里遇上轻蔑)。

这事想必是不知不觉地发生的,没有引起我的注意,也没有表现为某个特别而典型的、会被记忆保留下来的事件;而记忆打在它上面的那道光,也许跟打在那另一桩事上的一样,是刻意平淡的。我对那些过渡年月的记忆“整块地”还给我的,是这样一件事:前来接近我的人——无论是在我的讨论班之后,还是在布尔巴基讨论班之类的聚会上、在某个学术会议或大会上——往往得克服一种怯场,而在我们交谈(如果有交谈的话)的过程中,这种怯场多多少少还看得出来。当交谈超过几分钟,这份局促往往就在我们谈话、谈话热起来的过程中渐渐消失。有时——很少——也会发生这种情形:局促一直保持着,以至于成了交流的真正障碍,甚至在最不涉及个人的层面、在一场数学讨论的层面上;而我当时就模糊地感到对面有一种无力的痛苦,一种对自己恼火的痛苦。我说这一切时,并不是真的“记得”,而像隔着一层雾;可这层雾仍然把一些想必被记录过、又多半随即被排遣掉的印象还给了我。要我说出这份局促——这种恐惧的表现——何时出现,除了靠猜,我实在办不到。

我不认为这种恐惧源自我这个人,也不认为它只限于某种态度、某些使我区别于同事的举止。若真是那样,我想我在七十年代初就该听到些回声了——那时我走出一个自己一直扮演的角色,正是明星、“大老板”的角色。我相信,受人畏惧的是这个角色,而不是我这个人。而这个角色,连同它那圈与尊重毫无共同之处的恐惧光晕,在五十年代初并不存在,还没有存在——至少在我1948年一遇上就接纳了我的那个数学圈子里并不存在。

再说,在1970年那场“觉醒”之前,我根本不会想到把这种怯场、这种我有时在不那么亲近的同事身上遇到的局促叫作“恐惧”。它一显露,我自己就感到不自在,于是尽力去消除它。有一件事很值得注意,也典型地说明我亲爱的微观世界对这类事多么不关心:在我属于这个圈子的二十年里,我不记得有过一次,这个问题在我和某位同事之间被谈起,或者被别的人当着我的面谈起!\fnoteP{9}{(11)}{\emph{阿尔多·安德烈奥蒂(Aldo Andreotti)、约内尔·布库尔(Ionel Bucur)}\par{}当然,我这边也未必没有遗漏——何况我那时候那副十足“书呆子”的脾气,大概也不会鼓励别人跟我谈这类事,更不会促使我去记住一场很可能确实有过的、朝这个方向的谈话。可以肯定的是,至少可以说,这个问题被谈起的场合想必极其罕见(甚至不用“恐惧”这个名字去称呼它……),今天想必也一如其旧,尤其是在“上流社会”里。\par 在那个世界里的众多朋友中,除了谢瓦莱(他至少在六十年代期间必定意识到了这种恐惧的气氛),我觉得确实清楚察觉过它的,只有另外一个人,就是阿尔多·安德烈奥蒂。我认识他,以及他的妻子芭芭拉(Barbara)和他们的一对双胞胎孩子(当时还很小),是在1955年(我想是在芝加哥韦伊家的一次晚会上)。我们一直关系很密,直到1970年“大转折”的时候——那时我离开了曾属于我们的那个圈子,同他们也稍稍失去了联系。阿尔多极其敏感,而这份敏感在与数学、与像我这样的“书呆子”的往来中丝毫没有磨钝。他身上有一种对接近他的人自然生发的同情之才。这使他跟我所认识的数学圈内、乃至圈外的所有别的朋友都不一样。在他那里,友谊总是压过共同的数学兴趣(那些兴趣并不缺),而他也是极少数与我稍稍谈过我的一生、他也谈过他自己的一生的数学家之一。他的父亲和我的父亲一样是犹太人,他也因此在墨索里尼的意大利吃过苦,正如我在希特勒的德国吃过苦。我见他总是乐于鼓励和支持年轻研究者——那是在一个越来越难以被当权圈子接纳的气候里。他那种自然的关切,总先落到人身上,而不是落在什么数学“潜力”或名声上。他是我有幸遇到的最叫人亲近的人之一。\par 说到阿尔多,又勾起了对约内尔·布库尔的记忆——他也是出人意料地、未到年纪就被带走了;而且和阿尔多一样,我想,人们更多是把他当作一位乐于重聚的朋友来怀念,而不是当作数学讨论的伙伴。人能从他身上感到一种善良,旁边还有一份不多见的谦逊,一种不断退到后面的倾向。这样一个从不看重自己、也从不给人压力的人,最后竟当上了布加勒斯特大学理学院院长,真是个谜;也许是因为他压根儿没想过要推辞那些他远不渴求、却由同事们或政治当局放到他肩上的职务——那肩膀,说起来,是结实的。他是农民的儿子(在一个“阶级标准”很重要的国家里,这想必起了作用),身上带着农民的通达与朴素。他必定察觉到围绕着名人的那种恐惧,但想必也觉得这事理所当然,是权力地位的自然属性。不过,我不认为他本人曾让任何人畏惧——当然也不曾让他的妻子弗洛丽卡(Florica)、他们的女儿亚历山德拉(Alexandra),或他的同事、学生畏惧;而我能听到的反响,也正是这个方向。} (11)那层替我充作记忆的“雾”,也没有还给我
% ---------------- p.151 ----------------
\origpage{35} 任何印象,像是这类情形曾在我身上引起过有意识或无意识的满足。我不认为在意识层面有过;但我也不敢断言,在最初那些年里,我在无意识层面不曾偶尔被它拂过。若有,那也必定是短暂的,不曾反映到行为里、起一种把局促固定下来的作用。当然,这并不是说我演的那个角色里没有掺进我的自负!但即使我在这角色上不计本钱地投入,当时驱动我这份自我的,并不是想给“普通同事”留下深刻印象的野心,而是不断超越自己,以逼出我的“同侪”不断更新的敬意——而且,也许首先是那些在我还无从施展之前就信任我、把我当作他们中一员的年长者。看来,我面对自己成为对象的恐惧时所持的内心态度——尽力不去理会它,同时在它显露之处勉力把它消除——这态度可以看作整个六十年代我所属那个圈子(“微观世界”)里的典型态度。

在此后过去的十年到十五年里,局面又大大恶化了——至少从那个世界时不时传到我这里的迹象看是如此,从我得以近距离目击、甚至有时参与其中的情形看也是如此。不止一次,在我旧日的朋友或学生中那些与我曾最亲近的人那里,我撞见了轻蔑那种熟悉而无可辩驳的迹象;撞见了那种(表面上“无缘无故”的)使人泄气、羞辱人、压垮人的意志。一阵轻蔑之风起来了——我说不清是什么时候——吹拂着那个曾于我亲切的世界。它吹着,不管什么“功绩”或“过”,用它的气息烧焦卑微的天职,也烧焦最美好的激情。我旧日的伙伴们,各自带着“自己人”,受着坚固城墙的保护,安稳地(像我从前那样)坐在那层裹着各自的、毡子般悄然无声的恐惧里——他们中可有一个人感到这阵风?我确切地认识一个,只有一个,在我的旧友当中:他感到了,还跟我谈起过,只是没有叫出它的名字。还有另外一位,有一天也像是身不由己地察觉到了它,随即赶忙在翌日\fnoteP{10}{(12)}{“翌日”一词在此应按字面理解,不是隐喻。} (12)就把它忘掉。因为感到这阵风并承担它,对我的一位旧友来说,
% ---------------- p.152 ----------------
\origpage{36} 正如对我自己一样,也就是接受把目光落到自己身上。

\section{(15) 收获与播种}

我不打算——我再也不会打算——为一阵吹着的风而愤慨了,既然我已看清:我与这阵风并非无关,尽管我身上的一份自负一心要让我相信相反的事。而即便我与它无关,我的愤慨对受辱的人和施辱的人也都只会是一份太可怜的献礼——而这两种人,我一样爱过。

我与这阵风并非无关:在我为自己选定的那个世界里,我与轻蔑、与恐惧有过默契。对这些污迹,就像对许多别的污迹一样,闭上眼睛对我正方便——在我的职业生活中如此,在我的家庭生活中也一样。在这两者里,我收获了我所播种的——也收获了别人在我之前或与我一同播种的,既包括我的父母(以及我父母的父母……),也包括我那些“旧日的新朋友”。而今天,收获这些发了芽的播种的,还有我之外的别人:既有我的孩子们(以及我孩子的孩子),也有我今天的某个学生——他被我旧日的某个学生轻蔑地对待。

谈到播种与收获,我心里既没有苦涩,也没有认命,更没有自怜。因为我已懂得:即便在苦涩的收成里,也有一种实体的肉,只消我们肯吃,就能拿它来养自己。当这份实体被吃下,成了我们血肉的一部分,苦涩就消失了——它原来只不过是我们面对一份为我们准备的食物时的抵抗的征兆。

我也知道:没有哪一次收获不同时是对另一些收获的播种——而那些收获往往比前一次更苦。当那条看来没有尽头的链条——漫不经心的播种与苦涩的收获,代代相传、代代重演——出现在眼前时,我身上仍有某处会揪紧。但我不再像面对一种残酷而无从逃避的宿命那样被压垮、起而反抗,更不再像从前那样做它心甘情愿而盲目的囚徒。因为我知道,在我遇到的一切里都有一种滋养的实体,无论播种的是我的手还是别人的手——只消我去吃,去看它变成认识。对我的孩子们,对所有我爱过的人、以及此刻我爱着的人,情形也一样:当他们收获我在自负而漫不经心的年月里播下的东西,或者收获我直到今天偶尔还在播下的东西时,也是如此。

% ================================================================
%  第 7 章 双重面孔 (原书 p.153-166, 边码 p.37-p.56)
%  源文: _work/src1/task-ch07.md (= ch07.md)
% ================================================================
\chapter{双重面孔}

\section{(16) 沼泽与头几排}

% ---------------- p.153 ----------------
\origpage{37}可我还没有走到这场反思的尽头:反思的是,在轻蔑的出现与蔓延之中,我占了怎样一份——在这样一个世界里,我至今仍兴高采烈地以“数学共同体”之名来指称它。我现在感觉到,正是这场反思,才是我所能献给那个世界里我所爱过的人的最好的东西;此刻我正准备着——当然不是回到那个世界去,而是要在那里重新开口说话。

我想,我尤其还剩下这样一件事要做:审视一下,在我同他们一样仍是那个世界中一员的那些年月,我与这个那个成员之间,维持的是怎样一种关系。

如今回想起来,有一点让我吃惊:那个世界里有一整块人,我虽说常与他们擦肩而过,却始终没进入我的注意,仿佛他们并不存在。当年我大概把他们感知为某种“沼泽”——在我心里并没有明确的功能,我想,连“共鸣箱”的功能都谈不上;像一团灰色的、无名的人群:在讨论班和学术会议上,他们一成不变地坐在最后几排,仿佛生来就被分派在那儿;报告进行中,他们从不开口冒险提一个问题,因为他们事先就确信无疑:自己的问题只可能问得不着边际。倘若他们向像我这种被公认为“在行”的人提问,那准是在走廊里,在“权威们”显然无意彼此交谈的时候——这时他们便飞快地、像踮着脚尖似的把问题提出来,仿佛为占用我们这种大人物的宝贵时间而羞愧。有时那问题确实显得不着边际,我便(我想)试着用几句话说明为什么;常常它也是切中要害的,我同样尽力作答,我相信。两种情形下都一样:在这样一种心态里(或者我该说,在这样一种气氛里)提出的问题,很少会有第二个问题跟上,把它说得更确切、更深入。也许我们这些坐在头几排的人,在那种场合确实太匆忙了(尽管我们有时无疑也在努力不显出匆忙的样子),以至对面那份畏惧无法消散,交谈也就无从产生。我当然感觉得到,正如我的对话者在他那一侧也感觉得到:我们卷入的这局面有多么虚假、多么造作——只是我那时从未对自己把它说出来,而他也一样,想必从未说出来过。我们两个都像古怪的自动机那样运转,一种古怪的默契把我们连在一起:假装不知道那攫住我们之中一个人的焦虑,而另一个人则隐隐约约地觉察到了它——那是充满焦虑的空气里的一小份焦虑,这空气浸透了那些场所,所有人想必都像我们一样觉察得到,又都选择一致同意地视而不见\fnoteP{1}{(13)}{前面这段描述,显然只求勉强用具体的字眼,把记忆的那团“雾”交给我的东西还原出来;那团雾从未凝结成任何一个稍微确切的具体情形,好让我在这里给出稍微“写实”或“客观”一点的描述。若让这段文字说出这样的意思——那些不愿坐在头几排、或者没有明星与“权威”身份的同行,在同这样一位人物说话时必定满心焦虑——那就是曲解我的话了。在这个圈子里我认识的大多数朋友,显然都不是这样,甚至那些偶尔出入学术会议或讨论班的人也不是。毫无保留地成立的是:“权威”身份造成一道屏障、一道沟壑,把它与没有类似身份的人隔开;而这道沟壑极少消失,哪怕只是一场谈话的工夫。我还要补充:“头几排”与“沼泽”之间那种主观的区分(在我看来却实实在在),绝不能化归为社会学标准(社会地位、职位、头衔等等),甚至也不能化归为“身份”、名声的标准;它还映照出气质上、或者更难把握的禀性上的一些心理特点。我二十岁到巴黎时,就知道自己是个数学家,我做过数学;尽管有我曾有机会谈过的那些水土不服,我心底里觉得自己是“他们中的一员”——只是这一点只有我自己知道,而且我起初甚至并不确信自己还会继续做数学。如今我倒更倾向于坐到最后一排去(在难得出现这个问题的场合)。} (13)。

% ---------------- p.154 ----------------
\origpage{38}这种对焦虑的模糊感知,在我这里要等到1970年第一次“觉醒”之后才变成自觉的;那时,这片“沼泽”从我乐意让它一直待着的那片阴影里走了出来。这件事并不是出于什么有意的决定,我也没有立刻意识到,我就这样离开了一个圈子,进入另一个圈子——从“头几排”人们的圈子进入“沼泽”:突然间,我的新朋友大多正是那些一年前我还会不声不响地归入那片无名无状的疆域的人。所谓沼泽,忽然间动了起来,借着一张张朋友的面孔活了过来——这些朋友与我由一场共同的历险连在一起,另一场历险!

\section{(17) 特里·米尔基尔(Terry Mirkil)}

说真的,早在这场关键的转折之前,我就与几位同学(后来成了“同行”)结下了友谊,假如当年有人问我,我无疑会把他们归入“沼泽”(而且假如他们不是我的朋友……)。非得有这场反思、非得我去翻检记忆,我才想起他们来,那些零散的记忆才聚拢到一处。我认识这三位朋友,是在最初的那些日子里;那时我和他们一样,在南锡学着这门手艺——也就是说,那时我们还在同一条船上,还没有任何东西把我标举为一位“权威”。这多半不是偶然;而此后二十年里再没有过别的这样的友谊,多半也不是偶然。我们四个都是外国人,这肯定是一条不可忽视的纽带——我与那些和我一样被“空投”到南锡的年轻“高师人”的关系,就远没有那么私人化;我们几乎只在系里见面。三位朋友中有一位,一两年后移居南美。他和我一样是法国国家科学研究中心(CNRS)的研究专员,而我总有这么个印象:他自己也不太清楚他在“找”什么;久而久之,他在 CNRS 的处境变得有点尴尬。我们继续偶尔见面或通信,最后失去了联系。我与另外两位朋友的关系持续得更久,也更强烈,远不那么浮浅。何况数学上的共同兴趣在其中只扮演了极黯淡的角色,甚至等于零。

% ---------------- p.155 ----------------
\origpage{39}\origpage{40}与特里·米尔基尔(Terry Mirkil)和他的妻子普雷索西亚(Presocia)——她娇小纤弱,正如他粗壮结实;两个人都带着一股温柔的神气——我们在南锡常常一晚上、有时一整夜地唱歌、弹钢琴(那时是特里弹),谈音乐(那是他们的挚爱),也谈我们各自生活里这样那样重要的事情。老实说,不是最重要的那些——不是那些总被如此小心地闭口不谈的事。然而这份友谊给了我很多。特里有一份我所缺乏的细腻与辨别力,而我大部分精力早已倾注在数学上了。远比我更多地,他保留着对单纯而本质之物的感觉——太阳、雨、土地、风、歌唱、友谊……

特里在达特茅斯学院(Dartmouth College)找到一个合他心意的职位之后——那里离哈佛不算太远,从五十年代末起我常去那儿小住——我们仍继续见面、通信。其间我得知他容易陷入抑郁,这让他一次次长时间住进“疯人院”——在他唯一一次、也是极简略地向我提到此事的信里,他这样称呼那些地方;那是在又一次“可怕的长住”之后。我们见面时,从不谈这件事——只有一两次极偶然地,为了回答我的惊讶:他和普雷索西亚为什么不收养一个孩子。我不认为我脑子里曾冒出过这样的念头:他和我,我们可以谈谈问题的底;或者哪怕只是碰一碰——恐怕连“也许有些问题值得看看”这个念头也不曾有过,无论在我朋友的生活里,还是在我自己的生活里……在这些事情上,有一个不曾说出口、也无法逾越的禁忌。

渐渐地,见面和通信稀疏了。确实,我越来越成了工作和一个角色的囚徒,尤其是成了那份意志的囚徒——它变成了一种固执的念头,也许是对别的东西的一种逃避:在作品的累积中不断地超越自己——而与此同时,我的家庭生活正神秘地、不可挽回地败坏下去……

有一天,我从特里在达特茅斯的一位同行来信中得知,我的朋友自杀了(那已是他死去并被埋葬很久之后……);这个消息像穿过一层雾传到我这里,像来自一个极其遥远、而我不知何时已经离开了的世界的回声。也许那是我心里的一个世界,它早在特里结束自己生命之前就已死去——他的生命被一种焦虑的暴力摧毁,那焦虑他没能、或者不愿去化解,而我也没有能够、或者不愿意去猜出……

\section{(18) 二十年的自命不凡,或:不知疲倦的朋友}

% ---------------- p.156 ----------------
我与特里的关系,我想,在任何时刻都不曾因我们在数学界地位的差异、或因我由此而生的优越感而变味。这份友谊,以及那些年月里生活馈赠给我的另外一两位朋友(全然不管我“配不配”),肯定是当时少有的解药之一,用来对付一种隐秘的自命不凡——它由社会地位滋养,更由我对自己数学力量的自觉以及我自己给予它的价值滋养。我与第三位朋友的关系就不同了。这位朋友,以及后来他的妻子(他是在我们于南锡相识前后认识她的),在这些年里,每一次我们相遇,无论在他们家还是在我家,都对我表现出一种温暖的友谊,满含体贴与质朴。这份友谊里显然没有任何与地位或脑力相关的潜台词。然而,我与他们的关系,二十多年里始终带着我内心那种深刻的暧昧,带着我说过的那道分裂——它给我的数学家生活打上了印记。在他们面前,每一次我都忍不住感到他们亲切的友谊并作出回应,几乎是违心的!与此同时,二十多年里我居然做到了这样一手绝活:从我的伟大高处,轻蔑地俯视我的朋友。这想必从在南锡的头几年就这样开了头;而且很长时间里,我的成见也扩展到了他妻子身上,仿佛\origpage{41}\origpage{42}事先就可以认定:他的妻子只可能同他一样“无足轻重”。母亲和我之间,我们喜欢只用一个嘲弄的绰号来指他;这绰号想必在我母亲1957年去世之后很久还刻在我心里。如今我看出,至少在我这态度背后的诸般力量之一,是我母亲那坚强的个性在她一生中、乃至她死后近二十年里对我的支配——那些年里,我继续浸染在支配了她一生的那些价值之中。我朋友那温顺、和蔼、毫无争强好胜之心的天性,被不声不响地归入“无足轻重”一类,成了嘲讽式轻蔑的对象。直到现在,我第一次费心去审视这段关系究竟是什么样子,才发现这种面对他人温暖善意而执意自我隔绝的状态,范围有多大——它曾如此长久地给这段关系打上印记。我的朋友特里,并不比另一位朋友更好斗、更犀利,却幸好得到了我母亲的认可,不曾成为她嘲笑的对象——我猜,正因如此,我与特里的关系才得以在我心里毫无阻力地绽放。他对数学的投入并不更热忱,“天赋”也并不更出众,可我并没有因此就找到借口,用那层轻蔑与自得的硬壳把自己同他和他的妻子隔开!

在这另一段关系里,我至今仍无法理解的,是我朋友那份亲切的友谊,何以从未因他每次见面都必定感觉到的我的迟疑而灰心。然而今天我很清楚:除了那层硬壳和那份轻蔑,我还是别的东西;除了一块脑肌肉和由它生出虚荣的自命不凡,我还是别的东西。和他们一样,我心里也有那个孩子——那个我假装不知道的孩子,轻蔑的对象。我把自己与他隔断了,可他仍在我心里的某处活着,健康而茁壮,像我出生那天一样。我朋友们的感情,肯定是给那个孩子的——他们不像我那样与自己的根隔断。而肯定也正是他,在大老板背过身去的时候,偷偷地、匆匆地回应了那份感情。

\section{(19) 没有爱的世界}

% ---------------- p.157 ----------------
大老板老了,幸好,他也稍稍风化了一些,那小子从此可以自在些了。至于同这些真正经得起考验的朋友的关系,我想我在这里确实指到了我一生中最触目、最怪诞的一例:某种自命不凡(除了别的之外)在一段私人关系里所起的效用。也许我还在自欺,但我相信这也是唯一的一例:我与数学界(甚至别处)某位同行或朋友的关系,被自命不凡持久地占据,而不是让它满足于偶尔地、隐约而短暂地露一露头。此外我似乎觉得,我当时在数学界乐于来往的众多朋友当中,没有一个我能想象他也经历过类似的事——在一段与同行(朋友也罢,不是也罢)的关系里。在我所有的朋友当中,我也许是最不“酷”的,最“死用功”的,最不肯让一丝幽默透出来的(这本事我晚年才终于有了),最容易把自己看得无比严肃的。肯定也正是因此,我才不会太去寻求与我同类的人作伴(假定真有这样的人)!

奇怪的是,我的朋友们,不管属于“沼泽”与否,都容忍我,甚至喜欢我。这是一件好事,也是一件要紧的事,值得在这里说出来——尽管我们常常一连几小时、几天只见面讨论数学:感情在流动,就像今天它仍在流动一样,在当时的友人(凭一些有时纯属偶然的意气相投)与我之间流动;从1949年在南锡我被亲切接纳的那一刻起——在洛朗·施瓦茨(Laurent Schwartz)与伊莱娜·施瓦茨(Hélène Schwartz)家里(我在那儿有点像家里的一口人),在让·迪厄多内(Jean Dieudonné)家里,在罗杰·戈德芒(Roger Godement)家里(有一阵我也常去\origpage{43})。

这种包裹着我在数学界最初几步的温暖感情,我有点倾向于把它忘掉,它对我整个数学生涯都很重要。肯定是它,让我同那个以长辈们为化身的圈子的关系,也带上了同样温暖的色调;它给了我对这个圈子的认同以全部力量,也给了“数学共同体”这个名号以全部意义。

显然,对今天许多年轻数学家来说,事情是这样:在他们学习的年月、往往还远不止这些年月里,他们被切断了与一切感情、温暖之流的联系;他们的工作只映在一位疏远的老板眼中,映在他吝啬的评语里,有点像在读科研与工业部的一份通令——那通令剪掉了工作的翅膀,剥夺了它比一份阴郁而不牢靠的糊口营生更深的意义。

可我说着说着就提前了——说到七十、八十年代数学界那种失宠,也许是所有失宠中最深的一种;那个数学界里,发号施令的是我昔日的学生,以及我旧日朋友的学生。在那个世界里,老板常常把工作题目分派给学生,像扔给狗一根骨头——就这个,别的没有!像给囚犯分派一间牢房:你就在那儿服你的孤独刑!在那个世界里,某项细致而扎实的工作、多年耐心努力的结果,被那位无所不知、手握权力的人微笑着轻蔑地否掉:“这工作不合我胃口!”事情就这样定了。扔进垃圾桶吧,别再提了……

这样的失宠,我很清楚,在我所认识的那个圈子里、在我五十、六十年代常来往的朋友中间,并不存在。确实, 1970年我得知:在数学之外的科学界,这倒更像是家常便饭——甚至在数学里,显然也不那么罕见:公开露骨的轻蔑,明目张胆的滥用权力(而且无处申诉),甚至在某些有名望、我也曾有机会遇见的同行身上。但在我天真地当作“那个”数学界、或者至少当作这个世界忠实缩影的那个朋友圈子里,我没有见过这样的事。

然而,轻蔑的种子想必已经在那里了,由我的朋友和我播下,在我们的学生身上发了芽。而且不仅在我们的学生身上,也在我某些旧日的伙伴和朋友身上。但我的角色不是揭发,甚至不是斗争:腐败是不能靠斗争解决的。在我爱过的某个学生身上、或在某个旧日伙伴身上看到它,我心里有什么东西会揪紧——而往往,我不是去接受痛苦带给我的认识,而是拒绝这痛苦,挣扎,躲进拒绝和一种斗争姿态里:这种事不该存在!可它偏偏存在——而且我心底知道它的意义何在。从不止一个意义上说,我与此并非无关:我所爱的某个学生或旧日伙伴,乐于不声不响地踩一踩另一个我所爱、并且他在其中认出我的人。

我又跑题了,可以说跑了两重题——仿佛轻蔑的风只绕着我的住处吹!然而,正是它吹在我身上、吹在我亲近和珍爱的人身上的那口气,使我被它触及,认识了它。但现在还不是谈这个的时候,除非只对我自己,在沉默里谈。现在该是我重新拾起这场反思—见证的线索的时候了;它不妨就叫“追踪轻蔑”——追踪我心中和我周围的轻蔑,在我曾属于的那个数学界里,在五十、六十年代。

\section{(20) 一个没有冲突的世界?}

% ---------------- p.158 ----------------
\origpage{44}\origpage{45}我原打算用几行字谈谈“沼泽”,算是尽一份心,只说它在那儿,而我不常去——就像沉思中(数学工作里也如此)常有的事:你注视的那个“什么也没有”,结果却满是生命、奥秘,以及一向被忽略的认识。就像另一个“什么也没有”,它也像碰巧似的就在南锡(看来南锡真是我这个新身份的摇篮!):那个多半有点不中用的学生的“什么也没有”——人家是怎么对待他的,那可真该瞧瞧……刚才我还闪过一下这个念头,就在我写下(也许写得太快了点?)“这些失宠”在我们这儿“还不存在”的时候。这么说吧:我能报告的这类事,只有这独一无二的一桩;它像(这总得承认)我暗指的那种“失宠”,而我不愿过多逗留在细节描述上。受过这种对待的人,用不着我画图,就很清楚我说的是什么。那些没受过、但每次遇上都不急于闭上眼睛的人,也一样清楚。至于其余的人,那些穷开心地轻蔑别人的人,以及那些满足于闭上眼睛的人(就像我自己二十年间成功地做到的那样),哪怕给他们画一本画册也是白费……

我还剩下要做的是:审视这两个十年里我与同行、与学生之间的私人与职业关系,顺便也审视一下我所知道的、我最亲近的同行彼此之间以及他们与学生之间的关系。今天最让我吃惊的是:在所有这些关系里,冲突似乎根本不存在。我得马上补充:这在当时我觉得再自然不过——差不多是最起码的事。冲突,在善良意愿的人之间、在智性与精神上都已成年的人之间、以及诸如此类的人之间(再说一遍:最起码的事!),本就不该有。哪里真有了冲突,我便把它看作一种令人遗憾的误会:只要有应有的善意,把话说清楚,那准能在最短时间内了结,不留痕迹!我年轻时之所以选择数学作我最喜爱的活动,肯定是因为我感到:在这条路上,这种世界观最有可能不至于每每一头撞上令人不安的反证。毕竟,当人们证明了一件事,所有人就一致了——也就是说,善良意愿的人以及诸如此类的人,这自然不在话下。

碰巧,我确实感觉对了。这两个十年,是在我亲爱的“数学共同体”那个“没有冲突”(?)的世界的安宁里度过的;这段历史,也是我内心长期停滞的历史:眼睛和耳朵都塞住了,除了数学几乎什么也没学到——而在我的私人生活里(先是我和母亲之间的关系,然后是我在她死后随即建立的家庭),一场无声的毁灭正猖獗横行,那些年里我任何时候都不敢去看它一眼。可那是另一个故事了……我在这几行里常提到的1970年“觉醒”,不仅是我数学生涯的转折、一次圈子的彻底更换,也是(差一年)我家庭生活的转折。也正是那一年,我第一次在与新朋友的接触中,冒险偶尔——还非常偷偷地——瞥一眼我生活里的冲突。正是那个时候,一个疑问开始在我心里冒出,并在随后几年里成熟:我生活里的冲突,以及我有时在别人生活里感到的那种冲突,并不只是一场误会、一道用海绵一抹就掉的“污痕”。

这种(至少相对的)冲突之缺席,在我选为己有的那个圈子里,事后看来是相当了不起的事——而我后来终于知道:凡有人类生活的地方,冲突都肆虐横行,在家庭里如此,在工作场所也如此,无论是工厂、实验室,还是教授或助教的办公室。简直可以说, 1948年9月或10月,我懵懵懂懂地在巴黎下船,正好一头栽进了宇宙间唯一那座天堂般的小岛,那里的人们彼此相处没有冲突!

% ---------------- p.159 ----------------
\origpage{46}一下子,我觉得这事真了不起——在1970年以来我所知道的一切之后。它当然值得仔细看看——这是神话,还是现实?我看得见流动在我众多朋友与我之间、后来又流动在学生与我之间的那份感情,我用不着去编造它——可在这个仿佛把冲突放逐了的天堂世界里,我几乎不得不去编造冲突!

确实,在这场反思里,我毕竟有机会碰到过两种冲突的情境,每一次都像是我内心某种态度的显影剂:一个是南锡“不中用的学生”那桩事,其中当事人之间的是非曲直我并不清楚。另一个是我自己内心的冲突情境,一道分裂,在我与“不知疲倦的朋友”的关系里——但这一桩从未以人与人之间冲突的形式表现出来,而那本是冲突唯一被普遍承认的形式。值得注意的是,按这个词的通常意义,这些朋友与我之间的关系完全没有任何冲突——它任何时候都不曾有过一丝阴云。分裂在我心里,不在他们那里。

我继续清点。最先想到的之一:布尔巴基团体!在我或多或少定期参加它的那些年里——也就是直到五十年代末——这个团体对我而言体现了一种集体工作的理想:既尊重这项工作中看似极细微的细节,也尊重每个成员的自由。在布尔巴基团体的朋友们中间,我任何时候都不曾感到一丝强制的气味,无论是对我,还是对任何别人,无论资深成员还是受邀而来、试试看他和团体之间是否“来电”的客人。任何时候都不曾有过一丝影响力之争,无论是围绕议程上某个问题上的观点分歧,还是围绕对团体的霸权。这个团体没有首领而运转,而且据我所能察觉,显然没有人内心深处渴望扮演这个角色。当然,如同一切团体,某个成员对团体、或对其他某些成员的影响力大于别人。韦伊在这件事上扮演了一个特殊的角色,我谈过。他在场时,他有点像“主持人”\fnoteP{2}{(14)}{人们会以为这与“没有首领”的说法相矛盾,其实不然。对布尔巴基的老成员来说,我想,韦伊被看作这个团体的灵魂,但从来不是“首领”。他在场而且他乐意的时候,就如我说过的,他成了“主持人”,但他并不发号施令。他心情不好的时候,可以把他厌恶的某个题目上的讨论卡住;大不了等下一次大会韦伊不在时,再平心静气地把那题目捡起来,甚至第二天他不再作梗时就捡起来。决定由在场的成员一致作出,同时考虑到:一个人与所有其他人意见相左而站在正确一边,绝非不可能(甚至也不罕见)。这条原则对集体工作来说,似乎荒谬。了不起的是:它居然行得通!} (14)。有两次,我相信,我的敏感为此不快,我就走掉了——这是我所知道的、我自己造成的唯一“冲突”迹象。渐渐地,塞尔对团体的影响力变得与韦伊相当。在我还是布尔巴基成员的时候,这并没有造成两人之间的对抗局面,我也没有听说他们后来结下了仇怨。又过了二十五年再回头看,我所认识的五十年代的布尔巴基,在我看来始终是一个范例:在一个围绕共同计划形成的团体里,关系质量上取得了不起的成功。这个团体的品质,在我看来,比它产出的那些书的品质还要罕见。我这一生满是特权,其中之一便是遇到了布尔巴基,并在里面待了几年。我之所以没有留下来,绝不是因为冲突,也不是因为我所说的那种品质败坏了,而是因为一些更私人的工作更强烈地吸引着我,我把全部精力都献给了它们。何况这次离开,既没有给我与团体的关系投下阴影,也没有给我与其中任何一位成员的关系投下阴影。

我该逐一回顾一下我所卷入的冲突情境,那些使我与某位同行或某个学生对立起来的情境——1948到1970年之间。
% ---------------- p.160 ----------------
\origpage{47}\origpage{48}唯一稍稍值得一提的,是与韦伊那两次短暂的闹翻,前面已经说过。还有几片短暂的、极其短暂的阴云,落在我与塞尔的关系上,起因是我的敏感:有时一场谈话不再叫他感兴趣,他便有一种令人有时不知所措的洒脱,一下子把它掐断;或者对我投身其中的某项工作、我坚持的某种看法直说他没兴趣,甚至厌恶——也许我坚持得太多、太频繁了!这从未闹到翻脸的地步。除了气质上的差异,我们的数学意气相投格外强,他必定和我一样感到:我们彼此互补。

与我有着相当甚至更强意气相投的另一位数学家,只有德利涅。说到这儿,我想起1969年德利涅被任命到法国高等科学研究所(IHÉS)一事曾引起紧张,当时我并不觉得那是“冲突”(按说冲突会表现为翻脸,或同行关系中的一个转折)。

我想我已经清点完了——就人与人之间的冲突而言,就可由可见迹象察觉的冲突而言,在我常出入的那个圈子里,同行之间、或同行与学生之间的关系中,二十二年间就只有这些,信不信由你。也就是说,在我选中的那个天堂里没有冲突——那么,想必也就没有轻蔑?数学里又一个矛盾?

看来,我非得再仔细看看不可!

\section{(21) 一个守得很好的公开秘密}

昨天我肯定漏掉了一些小事,比如我与某位同行的关系里那些短暂的“冷淡”,尤其因我的敏感而起。我还该补上三四次我的自尊受挫的场合:有时同行和朋友在他们某篇文章里没有提到,我向他们说过的某个想法或结果想必在他们的工作中起过作用(在我看来如此)。我至今还记得,这说明那是个敏感点,也许并没有随年龄完全消失!除了一次,我都没有向当事人提起,他们的诚意当然不容怀疑。相反的情形肯定也发生过,只是我没听到回声。在我的“小宇宙”里,我不曾得知有哪一桩优先权问题成了翻脸或仇怨的由头,甚至不曾有过当事人之间尖酸刻薄的话。不过,唯一一次我进行了这样的讨论(在我看来那是个明摆着的例子),有过一场口角,倒把气氛澄清了,没有留下怨恨的残渣。那是一位格外出色的同行,他的本领之一,是能以惊人的速度吸收他所听到的一切;我似乎觉得,他常常有一种糟糕的倾向:把刚从别人嘴里听来的想法当成自己的。

这里有一个难题,它在所有数学家身上(而且不止在他们身上)都会以或强或弱的形式出现;它也不只是由于那种自我中心的习气——它促使我们大多数人(我也不例外)去归给自己一些“功劳”,无论真的还是想当然的。对一种情境(数学的或别的)的理解,无论我们怎样达到它,有别人帮助与否,本身都是一件本质上属个人的事,一种个人的经验,其果实是一种眼光,那也必然是个人的。一种眼光有时可以传达,但被传达的眼光不同于原初的眼光。话虽如此,要想辨明他人在自己眼光的形成中占的那一份,却需要极大的警觉。我自己肯定并不总有这份警觉——那是我最不关心的事,然而我却期待别人对我有这份警觉\fnoteP{3}{}{(9月30日)不过,事情的另一个侧面,见6月1日的注(晚于本文三个月),“暧昧”(n° 63),那里考察了对自己和他人有所宽纵的种种陷阱。}!
% ---------------- p.161 ----------------
\origpage{49}迈克·阿廷(Mike Artin)是第一个、也是唯一一个让我听到这件事的人。那天他带着一副泄露公开秘密的玩笑神气对我说:当你终于一把抓住一种实质、并在其中弄懂了点什么的时候,还费劲去分辨哪一份“属于自己”、哪一份“属于别人”,这既不可能,也完全徒劳。这话让我有点不知所措,因为这完全不合乎嘉当、迪厄多内、施瓦茨等人以身教给我的职业伦理。然而我隐隐感到,他的话里、以及他那笑盈盈的目光里,有一个我直到那个年纪都没能领会的真理:我的数学(尤其是数学产出)被自我强烈地占据着,而在迈克身上并非如此。他真给人一种印象:做数学就像一个玩得开心的孩子,同时又不忘吃喝。

\section{(22) 布尔巴基,或我的大运气——及其反面}

在往可见表面之下再潜一点之前,有一个判断现在就已摆在我面前:我五十、六十年代常出入的那个数学圈子,确实是一个“没有冲突的世界”,一点不假!这本身就是相当了不起的事,值得我稍作停留。

我得马上说明:这是一个非常狭小的圈子,是我数学小宇宙的核心部分,局限于我近身的环境中——我定期见面、关系最紧密的那二十来位同行和朋友。逐一检视他们时,我吃惊地发现:其中半数以上是布尔巴基的活跃成员。显然,这个小宇宙的核心与灵魂是布尔巴基。它差不多就是布尔巴基,加上与布尔巴基最亲近的那些数学家。六十年代我自己已不再是这个团体的成员,但我与某些成员的关系仍一如既往地紧密,尤其是与迪厄多内、塞尔、泰特(John Tate)、朗(Serge Lang)、卡蒂耶(Pierre Cartier)。此外我仍常去布尔巴基讨论班,或者不如说,正是那时我才成了常客;也正是在那时,我在那里做了我的大部分报告(关于概形理论)。

大概正是在六十年代,布尔巴基团体里的“调子”滑向了越来越明显的精英主义;我当时肯定是其中一分子,因此也就不可能察觉。我还记得1970年的惊讶:我发现布尔巴基这个名字本身,在数学界广大的阶层(此前为我所不知的阶层)里变得多么不受欢迎——它或多或少成了这些词的同义语:精英主义、狭隘的教条主义、对“典范”形式的崇拜而牺牲活生生的理解、晦涩、阉割式的反自发性,等等等等!而且布尔巴基不只是在“沼泽”里名声不好:六十年代,也许更早,我就偶尔从另一些心性不同、对“布尔巴基风格”过敏的数学家那里听到过回声\fnoteP{4}{(15)}{我并没有得到这样的印象:这种对布尔巴基“风格”的“过敏”,曾给这些数学家与我、或与布尔巴基其他成员及同情者之间的交流造成困难——假如这个团体的精神是宗派精神、是精英中的精英的精神,那倒会如此。在风格与时尚之外,这个团体的所有成员对数学实质都有敏锐的感觉,无论那实质来自何处。只是到了六十年代,我才记得我某位朋友把某些数学家称为“讨厌鬼”,因为他对他们的工作不感兴趣。这类事我几乎一无所知,便倾向于把这样的评价当作真话照收,被那份洒脱的自信镇住了——直到有一天我发现,那位“讨厌鬼”是个独到而深刻的心灵,只是没有福气讨我那位才华横溢的朋友喜欢。我似乎觉得,在布尔巴基某些成员身上,面对别人的工作(当自己不懂或懂得不透时)所持的那种谦逊(至少是保留)态度,已经磨损了。起初,那时那种“数学本能”还残留着——它使人感到一种丰富的实质、一项扎实的工作,而无须参照什么名声或名望。从各处传到我这里的回声看,我似乎觉得,谦逊也好,本能也好,今天在我曾属于的那个数学圈子里都已成了稀罕之物。} (15)。作为一个无条件的信徒,我对此感到惊讶, 
% ---------------- p.162 ----------------
\origpage{50}\origpage{51}也有点难过——我原以为数学能使人心一致!然而我本该记得:我起步的时候,硬吞一篇布尔巴基的文字并不总是轻松、也不总是启发人的,尽管它很干脆。那典范的文字,至少可以这么说,几乎不让人感到它写作时的气氛。如今我似乎觉得,这恰恰是布尔巴基文字的主要缺陷——连偶尔的一丝微笑也不能让人猜到:这些文字是人写的,而且是被远比某种对无情严格典范的无条件效忠更重要的东西连在一起的人写的。

不过,滑向精英主义的问题,以及布尔巴基写作风格的问题,在这里都是题外话。在这里让我吃惊的是:我选作职业圈子的这个“布尔巴基小宇宙”,是一个没有冲突的世界。这事在我看来更了不起,因为这个圈子里的主角们各有强大的数学个性,其中不少被视为“大数学家”,他们每一个都当然有分量形成自己的小宇宙,并且在其中当无可争议的中心和首领!\fnoteP{5}{(16)}{说实话,布尔巴基的成员里有好几位,肯定各有自己那个或大或小的“小世界”,在布尔巴基的小世界之外或之上。但在我自己的情形里,这样一个世界只是在我不再是布尔巴基的一员之后,在我全部精力都投入到属于我个人的工作之后,才在我周围形成——这也许并非偶然。} (16)这些强个性的人,两个十年里在同一个微小宇宙、同一个工作团体里和睦甚至亲切地共处——这才是让我觉得如此了不起、也许是独一无二的事。这与昨天谈布尔巴基时已经浮现的“非凡成功”的印象相合。

看来归根到底,我是有了这份非凡的运气:初次接触数学界时,我恰好落在时空里那个优越的地点——那里几年前刚形成了一个品质非凡的数学圈子,也许就品质而言独一无二。这个圈子成了我的圈子,并始终是我心目中理想“数学共同体”的化身;而这样一个共同体,当时(在我这个化身之外)大概并不比数学史上任何别的时刻更存在,也许除了某些同样狭小的团体(比如也许那个围绕毕达哥拉斯形成、精神却完全不同的团体)。

我对这个圈子的认同非常强,与我四十年代末诞生的数学家新身份不可分割。这是我超出家庭之外第一个被温暖接纳、被当作自己人接受的团体。还有另一条纽带,性质不同:我自己对数学的进路,在这个团体的进路里、在我新圈子里成员的进路里得到了确认。它并不等同于“布尔巴基式”的进路,但显然两者是兄弟。

此外,这个圈子对我而言还该代表那个理想的地点(或差不多是!),那个没有冲突的地点——对它的追寻无疑曾把我引向数学,引向一切科学中最不让我感到有一丝冲突意图的科学!而我刚才说到我的“非凡运气”,心里也明白这份运气有它的反面。如果说它让我得以发展出种种手段,在长辈已成同辈的圈子里以数学家的身份充分施展,那么它同样是一种受欢迎的手段,用来逃避我自己生活里的冲突,并造成一场长期的精神停滞。

\section{(23) 自深处}

% ---------------- p.163 ----------------
\origpage{52}这个“布尔巴基式”的圈子,肯定对我的为人、我对世界以及我在世界中位置的看法产生了强烈影响。这里不是试着界定这种影响、以及它如何在我生活中表现的地方。我只想说:我丝毫也不觉得我那些自命不凡的倾向、以及它们那种唯才是举式的合理化,是由我与布尔巴基的接触、由我进入“布尔巴基圈子”所激发的——至少四十年代末和五十年代不是。那些种子早就在我心里播下,在任何别的圈子里都会找到发展的机会。我报告过的“不中用的学生”那桩事,绝不是这个圈子里主导气氛的典型,恰恰相反,我再说一遍;它只是我自己身上一种暧昧态度的典型。布尔巴基的气氛是一种尊重人的气氛,一种自由的气氛——至少我是这样感受的;这种气氛本性上会打消并减弱一切趋向支配或自命不凡态度的倾向,无论是个人的还是集体的。

这个品质非凡的圈子不复存在了。它什么时候死的,我说不上来;想必没有一个人察觉,也没有一个人在内心深处为它敲响丧钟。我猜想,人身上一定发生了不知不觉的败坏——我们想必都“添了些年头”,变得干硬了。我们成了重要人物,有人听,有权势,有人怕,有人求。火花也许还在,但天真在路上丢失了。我们中也许有人会在死前重新找回它,像一次新生——但那个曾接纳我的圈子已不复存在,我若指望它复活,那是徒劳。一切都回到了秩序之中。

而尊重,也许也在路上丢失了。轮到我们有了学生时,也许已经太晚,最好的东西传不下去了——还有一星火花,但没有了天真,也没有了尊重,除了对“同辈”和“自己人”。

风可以刮起来,吹,烧——我们躲在厚厚的城墙后面,各人带着“自己人”。

一切都回到了秩序之中……

\section{(24) 我的告别,或:外国人}

% ---------------- p.164 ----------------
这场对我数学生涯的回顾,走上了与我预想完全不同的路。说实话,我本来根本没想到要回顾,只想用几行、顶多一两页说说:今天我与那个我已离开的世界是什么关系;也许还反过来说说:从我偶尔收到的回声看,我的旧友与我是什么关系。相反,我原本打算稍仔细地考察一下,我在那些紧张的数学工作年月里引入的某些想法和概念的遭际——它们有时颇为奇怪;或者我该说:我有幸瞥见、并把它们从全然未知的黑夜中拉进半明半暗、有时甚至拉到最明亮的日光之下的那些新类型的对象与结构!如今这番话,在这场已经变成对一段过去的沉思、变成一种更好理解并承担某个有时令人不知所措的当下的努力里,似乎显得走调。看来,原计划中对某个几何“学派”的反思,只好等一个更合宜的时机——那个学派在我的推动下形成,又几乎不留痕迹地挥发掉了\fnoteP{6}{}{这个“更合宜的时机”来得比预想的早,那场反思构成了《收获与播种》第二部分“葬礼”的内容。}。所以眼下,我关心的是把这场回顾进行到底:回顾\origpage{53}\origpage{54}我在数学界里的数学生涯,而不是对一部作品及其命运发议论。

在刚过去的五天里,我被反思笔记之外的别的任务占着,一个回忆却相当执拗地回来了。它将作为我停下的那篇《自深处》的尾声。

事情发生在1977年底左右。几周前,我被传唤到蒙彼利埃轻罪法庭,罪名是“无偿收留并供养一名身份不合法的外国人”(也就是说,一名在法国的居留证件不合规定的外国人)。正是这次传唤,让我知道了1945年那部规定外国人在法地位的法令中令人难以置信的一款;这一款禁止任何法国人以任何形式帮助“身份不合法”的外国人。这条法律,甚至在希特勒德国对犹太人也没有类似的条文,显然从未按字面执行过。由于一个非常奇怪的“巧合”,我有幸被选作第一个试验品,用来首次实施这条独一无二的条款。

有几天我一直发懵,像瘫痪了一样,深深地灰心。突然间,我仿佛被抛回三十五年前——那时人命不值几个钱,尤其是外国人的命……后来我反应过来了,我振作起来。几个月里,我把全部精力投进去,试图动员舆论,先是在我所在的大学和蒙彼利埃,然后在全法范围。正是在这段紧张活动的时期——为了一件后来证明注定失败的事业——发生了我今天可以称之为“我的告别”的那个插曲。

为了在全国范围行动,我给科学界五位特别知名的“人物”(其中一位是数学家)写了信,把这条法律告诉他们——它今天在我看来仍像我被告发那天一样难以置信。我在信里提议共同行动,表明我们反对一条邪恶的法律;这条法律等于把几十万住在法国的外国人置于法律之外,并把几百万其他外国人当作麻风病人一样,让民众对他们心存戒备——这些人一下子成了嫌疑犯,可能给那些不提高警惕的法国人招来最糟糕的麻烦。

奇怪的是,完全出乎我意料:这五位“人物”中没有一位给我回音。看来,我是有些东西要学的……

于是我决定去巴黎,趁布尔巴基讨论班的机会——在那里我准能遇到许多旧友——首先在我最熟悉的数学界里动员舆论。我原以为,这个圈子对外国人的事业会格外敏感,因为我所有的数学同行,和我一样,每天都要与外国同行、学生和留学生打交道;他们大多数、或者全部,都曾有过居留证件上的难处,都曾在警察总局的走廊和办公室领教过专横、而且往往是轻蔑。我曾把计划告诉洛朗·施瓦茨,他说讨论班第一天报告结束后会让我发言,把情况提交给在场的同行。

那天我就这样去了,公文包里一大叠传单,是给同行们的。阿兰·拉斯库(Alain Lascoux)帮我分发,在亨利·庞加莱研究所的走廊里,在第一场之前,以及两场报告之间的“幕间”。如果我没记错,他自己还做了一份小传单——他属于那两三位同行之列:听说了这件事之
% ---------------- p.165 ----------------
\origpage{55}\origpage{56}后,也动了心,在我去巴黎之前就联系了我,表示愿意帮忙\fnoteP{7}{(17)}{我投身的那场行动,得到的热烈回响和切实帮助,主要来自科学界之外。除了阿兰·拉斯库与罗杰·戈德芒的友好支持,我在这里还要特别提到让·迪厄多内的支持:他专程到蒙彼利埃出庭,为这件注定失败的事业在别的证词之外添上他热忱的证词。} (17)。罗杰·戈德芒也在其中,他甚至做了一份传单,标题是“诺贝尔奖得主进监狱?”。这在他是一番好意,可我们显然不在同一个频道上:好像丑闻在于动了一位“诺贝尔奖得主”,而不是随便哪个小人物!

布尔巴基讨论班第一天,人确实很多,极多或多或少与我相识的人,包括布尔巴基的旧友和旧伙伴;我想,他们大多数人想必都在场。还有好几位我过去的学生。我差不多有十年没见过这些人了,来的时候我很高兴有这个机会重见他们,哪怕一下子见这么多!不过,总会有机会再小聚的……

然而这次重逢“不是那么回事”,从一开始就相当清楚。许多手伸过来、握上了,这是肯定的;许多“哟,你也在这儿,什么风把你吹来了?”的问题,也是有的——但在快活的语调背后,有一种说不清道不明的尴尬气息。是因为我带来的事业他们并不真正感兴趣,而他们来这里是为了一场一年三次的数学典礼,需要全神贯注?还是与我来做什么无关,而是我这个人本身引起了那种尴尬——有点像一位还俗的神父出现在一群纯正的神学院学生中间所引起的那种尴尬?我说不好——也许两者兼有。在我这一边,我忍不住注意到,某些曾经熟悉、甚至曾是我的朋友的面孔上发生的变化。他们仿佛僵住了,或者塌陷了。我曾在他们脸上认识的那种灵动似乎消失了,仿佛从不曾有过。我觉得自己像是面对一些陌生人,仿佛从来没有什么把我与他们连在一起。我隐约感到,我们并不生活在同一个世界里。我原以为借这个不寻常的机会能找回一些兄弟,结果面对的却是陌生人。教养很好,这得承认:我不记得有过尖酸的话,也不记得有传单被丢在地上。事实上,好奇心帮了忙,分发的传单(几乎)都该被读过了。

然而那条邪恶的法律并没有因此受到威胁!我得到了我的五分钟,也许还占了十分钟,来谈那些在我看来是兄弟、却被叫作“外国人”的人的处境。阶梯教室里坐满了同行,比我做数学报告时还安静。也许对他们讲话的那份确信,已经不在我身上了。不再有从前那种同情与兴趣的流动。人群中想必有些赶时间的人,我一定是这样对自己说的;我缩短了发言,提议立刻与感到关切的同行们留下来,更详细地商议可以做些什么。

宣布散会时,是全体向出口的冲锋——显然,人人都有一班火车或地铁马上要开,绝不能错过!一两分钟之内,埃尔米特阶梯教室就空了,简直不可思议!空荡荡的大阶梯教室里,刺眼的灯光下,只剩下三个人,包括阿兰和我。第三个人我不认识,我敢打赌又是那种见不得人的外国人,交往可疑,而且身份还不合法!我们没有花时间对刚才在眼前上演的那一幕说得太多——它本身已经够雄辩了。也许我还是唯一一个不敢相信自己眼睛的人,而我的两位朋友当时很得体地对此不加评论。显然,我是初来乍到……

晚上在阿兰和他前妻雅克琳家里结束,盘点一下形势,把可以做的事过一遍;也彼此多认识一点。无论是这一天还是后来,我都没有费心把刚经历的这个插曲放进过去的脉络里去看。然而正是那一天,我一定是不言不语地明白了:某个圈子、某个我曾认识并爱过的世界已经不存在了;我原以为能重新找回的一份活生生的温暖,早已消散——多半消散很久了。

% ---------------- p.166 ----------------
这并没有阻止那个世界一年一年传到我这里的回声——那个温暖已逃逸的世界——许多次让我不知所措,让我痛苦地被打动。我怀疑这场反思对将来能否对此有所改变——除非也许是:我再被这样触动时,会少梗着脖子抗拒一点……

% ================================================================
%  第 8 章 师与生 (原书 p.167-174, 边码 p.57-p.66)
%  源文: _work/src1/task-ch08a.md
% ================================================================
\chapter{师与生}

\section{(25) 学生与研究计划}

% ---------------- p.167 ----------------
\origpage{57}我还没有把当年与其他数学家的关系巡视完——那时我觉得自己与他们同属一个世界,同一个“数学共同体”。剩下要看的,主要是我与我那些学生的关系——按我所经历的样子——以及我与另一些人的关系,那些人把我看作前辈。

总的来说,我相信可以毫不保留地说:我与学生的关系,是尊重的关系。至少在这一件事上,我相信,我当年做学生时从前辈们那里得到的东西,并没有随着岁月而变坏。由于我素有做“难”数学的名声(说实话,这是个再主观不过的概念!),而且被人认为比别家老板更苛刻(这一点就主观得少些了),来找我的学生从一开始就相当有劲头:“他们是想干的”!只有过一个学生,开头有点“提不起劲”,说不准他能不能起步——后来还是起了,而且没等我推一把,他自己就发动起来了。

就我所能记起的而言,凡是要求跟我工作的学生,我都收下了。其中有两个,几周或几个月之后发现我的工作方式不合他们的脾胃。说实话,现在我觉得那两次都是“卡住”的局面,而我当时却草率地把它读成了不适于数学工作的迹象。今天,我要作这类预言会谨慎得多。我毫不犹豫地把自己的印象告诉了这两位当事人,劝他们不要在一个——据我看——不合他们天性的行当里继续下去。事实上我后来得知,至少这两个学生中的一个,我是弄错了:这位年轻研究者后来在代数几何与数论交界处的一些难题目上赢得了名声。另一个学生是位年轻女子,她跟我这段不如意之后有没有继续,我就不清楚了。不能排除:我对她能力的看法说得太斩钉截铁,使她泄了气——而她也许和别人一样有能力做出好工作。我觉得我对这两个学生和对别人一样,是给了信任的。我缺的倒是%
% ---------------- p.168 ----------------
辨别力:面对那些肯定是“卡住”、而非“无能”的迹象时,没能分清两者\fnoteP{1}{(18)}{我想,这份缺乏辨别力,在这两回并不是出于我的疏忽,而更是出于不成熟、出于无知。直到大约十年之后,我才开始留意“卡住”的机制——在我自己身上,也在我亲近的人身上,在学生身上——并估量出它们在每个人的生活里扮演的巨大角色,而且不只是在中学或大学里。当然,我后悔这两回没有更成熟一些的辨别力;但我不后悔把自己那些印象——无论有没有根据——明明白白地说出来。当我在某一场合看出工作是敷衍着做的,把这些事照它们本来的样子点明,在我看来是必要而有益的。而如果在另一场合,我由此得出的结论是草率的、没有根据的,那么负有责任的也不只我一个。被这样摇了一下的学生,仍然有得选:要么从中吸取教训(也许头一回就是这样),要么就此泄气,也许还改了行(这也未必是坏事!)} (18)。

从六十年代初起,也就是在大约十年里,有十一个学生在我这里做了国家博士论文\fnoteP{2}{(19)}{\emph{耶稣与十二使徒}\par 从1970年至今,还有一个学生伊夫·拉德加耶里(Yves Ladegaillerie)在我这里准备并通过了论文。前一时期的学生是:P.贝尔特洛(P. Berthelot)、M.德马祖尔(M. Demazure)、J.吉罗(J. Giraud)、M.阿基姆夫人(Mme M. Hakim)、黄春生夫人(Mme Hoang Xuan Sinh)、L.吕西耶(L. Illusie)、P.茹阿诺卢(P. Jouanolou)、M.雷诺(M. Raynaud)、M.雷诺夫人(Mme M. Raynaud)、N.萨维德拉(N. Saavedra)、J.L.维尔迪耶(J.L. Verdier)。(其中六位还是在1970年之后才完成论文工作的,也就是在我很少能抽出时间指导他们的时期。)在这些学生当中,米歇尔·雷诺(Michel Raynaud)占一个特殊的位置:他论文工作所处理的那些本质性的问题与概念,是他自己找到的,而且他完全是独立地把它们发展出来的;所以我这个严格意义上的“论文导师”的角色,只限于读已完成的论文、组建评审委员会,并做其中一员。\par 在我自己出题目的时候,我很当心地只限于那些我足够熟悉、与我有足够深联系的题目,以便在需要时觉得自己能够支撑学生的工作。一个显著的例外是米谢勒·雷诺夫人(Michèle Raynaud)关于基本群的局部与整体Lefschetz定理的工作,这些定理是以适当的平展site上的1-栈的形式表述的。这个问题在我看来是困难的(后来也确实证明如此),而对于我提出的那些猜想(它们其实几乎无可置疑),我并没有证明的思路。这项工作一直延续到七十年代初;雷诺夫人(如同此前她丈夫的情形)发展出一套精细而独到的方法,没有得到我任何帮助,也没有得到别处的帮助。这项出色的工作还引出了一个问题:把雷诺夫人的结果推广到n-栈的情形——在我看来,在概形的框架里,这应当代表“弱Lefschetz定理”一类定理的自然归宿。不过,这里相应的猜想(它同样几乎无可置疑)的表述,本质上要用到n-栈的概念;而n-栈的追寻,据说正是本书的主要对象[其实那是《数学反思》第3卷,不是眼前这部《收获与播种》第1卷——见导论,p.(v)],正如“追寻栈”这个书名所表明的。我们大概会在适当的地方再回到这一点上,\par 另一个相当特殊的情形是黄春生夫人:我最初是1967年12月在河内遇见她的,那是我在被疏散的河内大学讲一个月“课程—讨论班”的时候。第二年,我向她提出了她的论文题目。她是在战争时期格外艰难的条件下工作的,与我的接触只限于零星的通信。1974/75年她得以到法国来(借温哥华国际数学家大会之机),并在巴黎通过了论文(评审委员会由嘉当(Henri Cartan)主持,成员还有施瓦茨(Laurent Schwartz)、德尼(J. Deny)、齐斯曼(M. Zisman)和我)。\par 最后我还得提到皮埃尔·德利涅和卡洛斯·孔图-卡雷尔(Carlos Contou-Carrère):这两位都多少算是学生,前者在1965–68年前后,后者在1974–76年前后。两人显然都有(而且一直有)不寻常的才能,他们运用这些才能的方式大不相同,运气也大不相同。来比雷之前,德利涅多少做过蒂茨(J. Tits)的学生(在比利时)——不过我怀疑,按通常意义上的“学生”来说,他在数学上不曾是谁的学生。孔图-卡雷尔做过圣塔洛(L. Santaló)的学生(在阿根廷),还曾有一段时间——多少算是——做托姆(René Thom)的学生!联系建立起来的时候,两人都已经有一个数学家的身量了;只是孔图-卡雷尔还缺方法、缺手艺。\par 我在数学上对德利涅的作用,只限于零零碎碎地让他知道我关于代数几何懂得的那一点点;而他学它,就像听人讲一个故事——仿佛他向来就知道。此外,还在于沿途提出一些问题,而这些问题他多半当场或在随后几天就找到了答案。这就是我所知道的德利涅最早的工作。1970年以后的(他的,以及我那些“正式学生”的),我仅凭极为零散、遥远的回声才有所知[其实那是《数学反思》第3卷,不是眼前这部《收获与播种》第1卷——见导论,p.(v)]。\par 按孔图-卡雷尔自己在论文开头所说的,我在他身边的作用,只限于把他引入概形的语言。无论如何,他近些年作为国家博士论文所准备的工作,我只是从很远的地方跟着;那题目是最时髦的那一类,超出了我的能力。在广阔世界里经历了几番波折之后,孔图-卡雷尔最近终于在最后一刻——而且(现在在我看来)是违着他的心意的——被逼得来求助于我,请我充当论文导师并组建评审委员会。(这就使他冒着一种风险:显得像是“1970年之后”的格罗滕迪克的学生——在这种情势下,这可能带来严重的不便……)我尽我所能履行了这项任务;这大概是我最后一次行使这个职能(在国家博士论文这一层面上)。在这个有点特殊的情形里,我格外高兴有让·吉罗(Jean Giraud)的友情协助:他也从自己的时间里抽出一两个月,细读那部篇幅庞大的手稿,并写了一份详实而热忱的报告。} (19)。各自选定了合意的题目之后,他们都干得起劲,而且(我是这样感觉的)都与自己选定的题目深深认同。然而有一个例外:有个学生选了一个“该做的”题目——也许并无真正的信念——而这题目也有它不讨人喜欢的一面:它要做的是把一批已经到手的思想作技术上的定稿整理,有时艰难,甚至干枯,再也没有多少意外、多少悬念可言,%
% ---------------- p.169 ----------------
\origpage{58}只是把它们摆出一个全景来\fnoteP{3}{(20)}{这让我想到莫妮克·阿基姆(Monique Hakim)选的题目:说实话,那也不更招人喜欢;我真不知道她是怎么保持住精神的!她固然有时也吃力,但无论如何没到让她变得忧郁或阴沉的地步,而我们之间的工作是在一种亲切、放松的气氛里做的。} (20)。

我当时被一项庞大计划的需要卷着走——我需要人手——所以在出这个题目时,必然是缺乏心理上的辨别力:它肯定不合这位学生特有的性情。他那一头呢,大概也没太看清自己上的是一条什么船!总之,他和我都没能及时看出:这事一开始就起步错了,不如换个题目重新来过。

显然,他干活没有真正的信念,而且始终带着一副多少有些忧郁、阴沉的神气,不曾放下。我想我当时已经到了这样一个地步:不太留意这类事情——而这类事情(我本该记得)在任何研究工作中,而且不只在研究工作中,都能造成天壤之别!我那时的角色,仅限于:工作显出拖沓的样子,我就心烦;工作又动了,我就宽慰地“呼!”一声;等到最后,预定的计划总算“收了尾”。

直到1970年我那次“醒悟”之后好几年,因为要和这位从前的学生通信(他当了教授——再说,在那宽厚的年月里,人人都当上了教授!),我才想到:那一回终究是有什么地方不对头,那也许算不上一次彻底的成功。今天,它在我看来是一次失败——尽管“计划收了尾”(绝非草草了事!),尽管有文凭,还有随之到手的位置。而且我要负很大一部分责任:我把一个计划的需要,放在了一个人的需要之前——而这个人,是怀着信任把自己交到我手上的。我刚才还拿来标榜的那种“尊重”(“毫无保留”!),我说我对学生们做到了——在这里始终是表面的,与尊重的真正灵魂相隔绝:那灵魂就是对一个人的需要怀着情意的关注,至少当这些需要的满足取决于我的时候。这里的需要,是工作中的一份快乐;没有它,工作就失了意义,变成强制。

在这次反思中,我有机会谈到过一个“没有爱的世界”,并且在自己身上寻找那个我拒斥的世界的胚芽。这里就有一粒,而且个头不小——至于它在别人身上怎样长起来,我今天就说不上来了。这种表面的、缺乏关注、缺乏真正爱的尊重,也是我给予自己孩子们的“尊重”。和他们在一起,我有这份特权:看着这粒种子长起来,看着它蔓延。我也多少明白了:面对收成,皱眉头是没有用的……

\section{(26) 严格与严格}

如果把这位学生除外——他肯定并不比别人“天分”差——我可以说,我与学生之间的关系是融洽的,常常甚至是亲热的。出于无奈,他们%
% ---------------- p.170 ----------------
\origpage{60}都学会了耐心对待我作为“老板”的两个主要毛病:一是字迹无法辨认(不过我相信,最后人人都学会了辨认我的字);二是——这一点当然更严重,而我很晚才发现——我从根子上难以跟随别人的思路,除非先把它译成我自己的意象,用我自己的风格重想一遍。我远更倾向于把自己深受浸润的某种对事物的眼光传给学生,而不是鼓励他们心中绽出一种个人的眼光——那眼光也许与我的颇为不同。与学生的关系中这层困难,到今天也还没有消失;但我觉得它的影响减弱了,因为我意识到了自己身上的这种倾向。也许我的气质——无论是先天的还是后天的——使我更适于孤独的工作(我起初十五年的数学活动,约1945至1960年,也正是这样过来的),而不那么适于做“师傅”,去和那些数学上的志业与个性尚未完全成形\fnoteP{4}{(21)}{也许更准确的说法是:就我这种气质而言,我至今缺的是必要的成熟,还不能完全承担起一个教育者的角色。我后天形成的气质,长久以来被一种过分占优势的“男性”(或“阳”)特征所标记;而成熟的一个方面,恰恰是以“女性”(或“阴”)为主的“阴阳”平衡。\par(后来补写。)比我缺一份成熟更甚,我看在我作为教育者的生活里,直到今天,我缺的是一种慷慨——一种比在时间和精力上的投入更为精细地表达出来的慷慨,也是更本质的慷慨。在我教学的第一时期,这份欠缺并没有以可见的方式表现出来(比如说,以一连串失败局面的积累),多半尤其是因为,那些选择来跟我工作的学生身上的强烈动机把它补偿了。相反,在第二个时期,从1970年至今,在我看来,对于我在研究层面(也就是从DEA这一层往上)的教学中所看到的整体失败,这份欠缺至少是原因之一,而且无论如何是最直接牵涉到我本人的那一项原因。关于这一点,见《纲领概要》第8段,以及第9段《一项教学活动的总结》,那里透出这项活动七八年来留在我身上的受挫感[并请对照后来补写的注(23iv)。]。} (21)的学生相处。不过也确实:我从小就喜欢教书;而且从六十年代到今天,我能有的那些学生在我生活里占了重要的位置。这也就是说,我的教学活动、我作为教师的作用,在我的生活里曾经有过、并且仍然保有一个很大的位置\fnoteP{5}{(22)}{不过也许不会太久了,因为我已经决定申请进入法国国家科学研究中心(CNRS),从而结束在大学环境里的教学活动——这几年来,那活动变得越来越成问题。} (22)。

在我教学活动的前一时期,我和任何一个学生之间都没有过表面上的冲突,哪怕只是关系里一阵短暂的“冷淡”也没有过。只有一次,我不得不对一个学生说:他工作不认真;要是照这样下去,我就没兴趣再和他一起干了。那事情究竟怎样,他当然和我一样清楚;他改了过来,这件事就了结了,没有留下一片乌云。另一次,已经是七十年代初,那时我最主要的精力投在“生存与生活”(Survivre et Vivre)小组的活动里;有个学生看了我刚写好的、关于他工作的论文评审报告(我习惯如此),发起火来,认为报告里某些议论是在质疑他工作的质量(这绝不是我的本意)。这一回是我把话收回、把方向拨正了,没有为难。当时我并不觉得这段小插曲会在我们的关系里留下阴影;不过也可能是我弄错了。这位学生和我的关系,比我和别的学生更少个人色彩(我说过的那个“忧伤的学生”除外);那只是一份良好的工作关系,仅此而已,我们之间并没有真正的暖意流过。不过,我并不认为是我身上一种无意识的不厚道,使我把那些他认为于己不利的议论写进了报告——他还说“这事他不会放过”,不像他的一个同伴,那人已经在我这里通过了论文。与另一位学生——他天性敏感而多情——我则有着一份%
% ---------------- p.171 ----------------
\origpage{61}格外亲热的关系;如果说我在关于他论文的报告里也写进了那类曾让他同伴大为不快的议论,那肯定不是出于不厚道!再说,对他们两位,也对我所有的学生,只要我对他们提交的工作不是完全满意,我就不会给答辩开绿灯。而且,那个时期我的学生,论文一通过,没有哪一个不是很快就找到了一个与他相称的位置。

直到1970年:我几乎随时都能抽出时间给学生\fnoteP{6}{(22')}{即便在1970年之后,当我对数学的兴趣在我的生活里变得零散而边缘时,我也不认为有过学生来找我、要我跟他一起工作而我推辞掉的场合。我甚至可以说,除了两三个例外,1970年之后我的学生对他们所做工作的兴趣,远低于我自己对他们题目的兴趣——即便在我除了去学院的日子几乎不操心数学的那些时期。因此,我在1970年之前随时能抽出时间给学生的做法,以及作为它主要标志的那种在工作上的极端严格,对于我后来大多数学生都会毫无意义——他们做数学没有信念,仿佛是靠着一种他们不得不对自己不断施加的努力……} (22')。时机成熟的时候,每逢用得着,我就和这一个或那一个待上一整天、必要时好几天,去啃那些还没弄妥的问题,或者一起把他们文稿的一稿又一稿过一遍。就我对这些工作时段所体验到的而言,我不觉得自己在里面扮演过拍板定案的“导师”角色;每一次都是共同的研究,讨论平起平坐地进行,直到双方都完全满意。学生投入的精力是巨大的,当然无法与我自己应投入的相比——不过我这边有更多的经验,有时还有训练得更敏锐的嗅觉。

然而,在我看来,对任何研究——无论智性的还是别的——的质量最要紧的东西,绝不是经验的问题。那是对自己的严格要求。我要说的这种严格,本质上是精细的;它不属于一丝不苟地遵从无论什么样的规范——严格的规范也好,别的规范也好。它在于对我们自身内部某种细微之物的极度关注,而那东西逃开一切规范、一切尺度。那细微之物,就是:对所考察之物的理解,有,还是没有。更确切地说,我要说的这种关注,是对每一时刻当下那种理解之品质的关注:从一堆杂乱堆叠的概念与陈述(假设的或已知的)发出的那片噪音,直到完全的满足、完美理解那种完成了的和谐。一项研究的深度——无论它的结果是片断的理解还是总体的理解——就在这份关注的品质里。这样一种关注,并不像遵守某条训诫、凭着“当心点”的刻意意图去留神的结果;在我看来,它是从求知的激情里自发地生出来的,是那些把求知的冲动与它那些自我中心的赝品区分开来的标志之一。这种关注有时也被叫作“严格”。这是一种内在的严格,独立于某一时刻在某门(不妨说)特定学科里可能通行的那些严格规范。如果说在这本书里我容许自己对那些严格规范(我曾教过它们,它们自有其存在的理由和用处)放肆一些,我并不认为这本书里这种更本质的严格,比我过去那些规范文体的著作里更少。而如果说我毕竟也许还能传给学生某种比一套语言和一种手艺更贵重的东西,那无疑就是这份严格、这份关注、这份严谨——如果说不是在与他人的关系和对自己的关系里(在那一层,它在我身上和在任何别人身上一样缺),那至少是在数学工作里\fnoteP{7}{(23)}{\emph{孩子与老师}\par“传递”这个词在这里并不真正符合事情的实际,它提醒我要采取一种更谦逊的态度。这种严格不是可以传递的东西,至多只能被唤醒或鼓励——而它从很小的年纪起就被忽视了、被挫伤了,被家庭环境,也被学校和大学。就我所能记起的,这种严格一直在我那些追寻里——至少是智性性质的追寻里;而我不认为它是我的父母传给我的,更不认为是学校的老师们、或我那些数学前辈们传给我的。在我看来,它属于“无辜”的禀性,因而也属于每个人生来就被分派到的东西。这份无辜很早就“青的、没熟的都尝遍了”,于是它不得不或深或浅地沉下去,以致在余生里往往再也见不到它多少痕迹。在我身上,由于一些我还没想到去探究的原因,某种无辜在智性好奇这个相对无害的层面上存活了下来;而在别处一切地方、在所有人身上,它都深深地沉了下去,神不知鬼不觉!也许“教育”——就其完整意义而言——的秘密,或者说谜,就在于重新与这份表面上消失了的无辜取得接触。但是,如果在教育者本人身上这份接触不是已经先在、或已被重新找到,那就谈不上在学生身上重新找到它。而这时教育者“传”给学生的,绝不是这份严格或这份无辜(它们在这两人身上都是天生的),而是一种尊重,一种对那件通常被弃置之物的默默的重新肯定。} (23)。

% ---------------- p.172 ----------------
\origpage{62}这确实是一件很微薄的东西,但也许毕竟聊胜于无。

\section{(27) 失误——或二十年后}

也许除了我谈过的那两个学生——与他们终究没能建立起工作关系——之外,我不记得其他那些来找我、要求跟我工作的学生,是带着“怯场”或畏惧来的。想必他们多少已经认得我,哪怕只在我于法国高等科学研究所(IHÉS)的讨论班上待过一阵子。倘若在我们关系开始时有什么拘束,它最后总会在工作过程中消散,不再留下痕迹。不过这里我要举两个例外。一个就是那位始终没能真正对工作生出兴味的学生,甚至在我们共同工作时,他也一直只是单音节地应声。也许他来的时候,我能留给学生的时间正要减少;和他也没有过那种对着稿子、一坐就是一个下午、一整天的时段。不,确实,我不记得有过那样的时段;我倒觉得我们多半是匆匆见一面,一两个钟头,看看他做到哪儿了。说到底,撞上我,他大概是最倒霉的一个!

我本来要说的另一个学生,则是在我还随时能抽出时间给学生的那个时期跟我工作的。我们的关系从一开始就融洽。他甚至属于那少数几个与我建立起友情的学生之列——我有时去他们家,他们也来我家,有点像两家之间的往来。确实,即便在这些情形里,关系也始终停留在相对表面的层次——至少就我这一方而言。在意识层面上,那时我对自己家里、自己屋顶下发生的事就已经所知无几;至于我那些数学家朋友——学生也好,不是学生也好——的生活,除了妻子和孩子的名字,我最终几乎一无所知(就连这些名字,我有时也会忘掉,而从没有人为此怪我!)。也许我曾是“书呆子”的一个极端例子;但我认为,在我所认识的那个数学圈子里,大多数关系——即便是友好的、亲切的——都停留在这种表面的层次:彼此最终所知极少,除了在“未言明”的层面上被感知到的那一点。这无疑是人与人之间的冲突在这个圈子里如此罕见的原因之一——尽管在我看来很清楚:分裂存在于我的大多数同事和朋友内部,也存在于他们的家庭内部,和我身上、和别处一样多。

我不认为我与这个学生的关系有别于我与别人的关系;当时我也不觉得反过来,他与我的关系与其他学生的关系有什么显著的差别,尤其是与那些同我结下友情的学生相比。只是不久以来,我才得以明白:那一定是一份比我对大多数别的学生更强的关系。一个未曾说出的冲突,其可见的表现像一次意外的启示般到来——那已是他做我学生将近二十年之后。直到那时,我才把它和一件早已被遗忘的“小”事联系起来。在很长一段时间里,也许甚至在整个那一段时期(也就是几年)里——%
% ---------------- p.173 ----------------
\origpage{63}\origpage{64}我们或多或少定期一起工作的那些日子——这个学生一直保留着某种“怯场”。它每次见面都表现出来,带着不会骗人的迹象。随后,在共同工作的过程中,这些迹象又相当快地消失了。这些不安的迹象当然让我不自在,而我觉得他更不自在。我们两个都装作不知道这事,这也理所当然。想必我们谁也不会想到去谈它,甚至不会想到在心里给这样一个古怪、显然值得一顾的处境一点注意!在他那边,也在我这边,这份“怯场”想必被感受为一处单纯的“失误”,一处本不该有的东西。那“失误”定期提醒我们想起它,但每一次,它都识趣地消失,让我们有时间安安稳稳地去忙正经事——数学——同时忘掉“那本不该有的东西”。我不记得自己哪怕一次停下来,对那失误的意义提出什么问题;我也确信,我的学生兼朋友那一头也是如此。想必在我们各自从幼年起在身边所见的一切里,没有任何东西能在他或我心里提示出另一种对待一件令人不快之事的态度,除了尽可能把它挪开、让它不再碍事。在那件事上,这完全做得到,甚至很容易;而我们完全一致地什么都没看见、什么都没感觉、什么都没听见。

然而,借这两三年来回到我这里的许多回声与相互印证,我明白:那被我们当作“本不该有”而挪开的东西,并不因此就不再存在、不再表现出来。有时回到我这里的东西,也同样“不该有”——可它“在”,而且现在已经不能一手背就把它拂开了……

\section{(28) 未完成的收获}

直到1970年第一次“醒悟”之前,我与学生的关系,正如我与自己工作的关系,是满足与快乐的源泉,是我生活中那份和谐感的可触、无可否认的基础之一——那份和谐继续赋予我的生活以意义,而与此同时,一种捉摸不定的毁坏正在我的家庭生活里肆虐。那时候,在我眼里这些关系里没有任何表面冲突的成分;当时它们当中没有哪一个,哪怕在一瞬之间,成为挫折或痛苦的起因。有一件事也许显得悖理:我与某个学生的关系里的冲突,偏偏要等到那次著名的醒悟之后才显出来——也就是在一个转折之后,那转折给了我的生活一种它从前不曾有过的开阔,也许还给了我这个人一点点柔韧的开端——人们会想,这些品质本该有助于化解或避免冲突,而不是挑起或激化它。

可是再仔细看,我就看出这悖理只是表面的;无论从哪个角度看,它都消散了。第一点:一个冲突要有机会化解,首先得它已经表现出来。冲突表现出来的阶段,相对于隐藏着、被忽视的冲突而言,是一种成熟;而后者其实也有表现,而且那些表现恰恰因为借它们表达出来的冲突仍被忽视而更为“有效”。第二点:一个冲突要能以可辨认的方式表现出来,首先得有一段距离缩小了,或者消失了。近十五年来我生活里发生的种种变化——尤其是在一次次“醒悟”过程中——在我看来,都是那种能缩小距离、抹去隔绝的变化。一个冲突面对一位有声望、受人仰慕的老板时难以表达,面对一个被剥去了权力位置的人(在这里是自愿剥去的)、一个从某个握有权威与声望的圈子里自我放逐的人,就自在多了——这个人越来越不被看作某种实体(比如数学)的化身或特权代表,而越来越被看作一个和别人一样的人:%
% ---------------- p.174 ----------------
\origpage{65}\origpage{66}一个不仅可能被伤到,而且越来越不愿把自己的伤或痛藏起来的人。第三点,也是最要紧的:我第一次醒悟以来的演变——尤其是在那个时期以及随后的几年里——本就容易在我那些旧学生井井有条的世界里激起(或者也许是唤醒)疑问、不安、一种“重新审视”。我有充分的机会看出:不仅在他们身上如此,在数学界我昔日的朋友与同伴中间也如此,有时甚至在一些只凭传闻认识我的科学同行中间也如此。

还得说:一个稍有深度的冲突能得到化解,是最罕见不过的事。最常见的情形是:尽管有一切休战与表面的和解,我们那些冲突越排越长的队伍仍一路跟着我们,几乎寸步不离地跟一辈子,最后才在殡葬工那双阴沉的手里把我们松开。我有时有幸看到一个冲突多少解开了,有时甚至看到它在认识中化解——但直到如今,这样的事还不曾在我与某个学生、或与数学界某位昔日朋友的关系中、并借这关系而发生。我也很清楚:这样的事是否终会发生,毫无把握,哪怕我还能再活一百年。

有一件事很值得注意:我与某段过去决裂的那个时刻——我是说离开IHÉS那个插曲(这个机构多少像是我周围形成的那个数学小宇宙的“母模”)——这个决定性的插曲,同时竟是我某个学生对我的深刻敌意第一次表现出来的场合。肯定是这一情形,使那个插曲格外难堪、格外痛苦,像一次分娩或一次出生发生在格外艰难的条件下。当然,那时我还不能在我后来学会的光亮里去看这个插曲——它的意义当时逃开了我。此后很久,这份痛苦的意外仍留在那里。然而,就在那一年的夏天,这场带着苦涩的离开已经显为一次解放——就像一扇门忽然大大敞开(只消我推它一下!),通向一个不曾料想的世界,唤我去发现它。此后每一次新的醒悟,也都是一次新的解放:发现一种从属、一种内在的束缚,并且重新发现一个巨大未知者的在场——它就藏在那本该“已知”之物的熟悉外表后面。但同样在这整整十五年里,直到今天,这份固执、沉默而无懈可击的敌意一直跟着我,成为我在数学生活里所知道的唯一而持久的挫折之源\fnoteP{8}{(23')}{不过七八年来,在我的数学生活里还有另一个慢性的“受挫之源”,只是它在这些年里表达得隐蔽得多。靠着同一种“令人受挫”的局面在我的教学活动中顽固地重复、积累,它最终还是显了出来,并且最后爆发成一种“受够了!”,使我几乎终止了一切所谓“研究指导”的活动。在这次反思的过程中,我有一两次触及这个问题,最后至少在末尾稍作考察。我在那里至少描述了这份受挫感,并考察了它在我“回到数学”这件事上扮演的角色(参看第50段“一段过去的重量”)。} (23')。也许我可以说,它是我为那第一次解放、以及随后的几次解放付出的代价。但我很清楚:解放与内在的成熟,是与“要付的代价”不相干的事;它们不是“赚”与“亏”的问题。或者换个说法:当收获被带到头、当它完成的时候,就没有损失——那看似“损失”的,已经变成了“收获”。而事情就变得清楚:我至今还没能把那场收获带到头;就在我写下这几行的此刻,它仍是未完成的。

% ================================================================
%  第 8 章 师与生 (原书 p.175-186, 后半: 8.5-8.8)
%  源文: _work/src1/task-ch08b.md
%  说明: 本章前半 (8.1-8.4) 见 part1-ch08a.tex
% ================================================================

% ---------------- p.175 ----------------
\section{(29) 敌父(1)}

1970年那个转折之后,在完全不同的环境里——一所外省大学——开始跟我做工作的那类学生,与从前的学生也很不相同。此外,在我这里做到国家博士论文水平的,只有两个人。其余人的工作,都在DEA或第三阶段博士论文的水平上。我还该算上一大批紧紧跟住某些“入门研究”课程的学生:这些课对他们来说,是提出往往出人意料的数学问题的机会,有时也是想象出别出心裁的方法去解决它们的机会。我遇到的最积极的参与,是在给一年级学生开的某些“选修课”上。相反,在那些已经在大学氛围里泡了几年的学生身上,某种清新、某种兴趣的能力、某种个人的眼光,已经或多或少褪了色。选修课的学生里,有几个显然具备成为出色数学家的素质。看当时的情形,我克制住自己,没有鼓励任何一个人走上那条路——尽管那条路本可以吸引他们,他们本可以在那里出类拔萃。

我跟那些为准备硕士文凭而听我这类“课程”的学生,关系大多没有延续到一年以上。每一次,我的印象都是:这些关系很快变得亲切而松弛,总的说来如此。只有一个被无孔不入的“怯场”缠住的学生是例外\fnoteP{9}{(23$''$)}{\emph{不敢玩}\par 这个学生曾和我一起做了一整年的DEA“实习工作”,而在他与我的工作关系里,直到最后都保持着“紧绷”状态。那是一种坦率友好的关系,贯穿着一种无可怀疑的相互好感。然而可就是有这份“怯场”;这份恐惧——它真正的起因肯定不是对我这个人的畏惧,尽管看上去像是如此。要不是这个学生自己跟我谈起,我也许根本不会察觉这件事;他谈,无疑是为了或多或少“解释”一下他在这一年里工作几乎完全卡住的原因。\par 正如在其他学生身上发生过的那样——他们和他一样,起初对某种几何实质抓得很牢——卡住恰恰出现在要“做零件活儿”的时候:也就是把成形的陈述白纸黑字写下来,或者仅仅是把握我提供的那些陈述的意义与含义;我提议把它们当作一种语言的基石、当作“游戏规则”接受下来。“学校式的”反应几乎总是把面对“做研究”局面的学生推向:把导师传下来的那些心照不宣的“游戏规则”当作既定的、又必须服从的东西接受下来,而这些规则恰恰不是要试着说明白的,更不是要弄懂的。这些心照不宣的规则所采取的具体形式,就是语义学或计算上的“配方”,不妨说,像预科班的题解书那样(或者任何一本普通教科书)。此外,学生还期待导师给出“证明……”这种形式的任务——那是他以往经验里遇见的唯一一种数学“思考”形式。(而且我不认为大多数职业数学家、以及其他科学家的禀性有本质不同——只不过“导师”换成了“共识”:共识规定当下的游戏规则,并把它当作不可更动的既定之物。这共识还规定哪些“问题”是要去解决的;在这些问题之间,每个人都觉得自己有随意挑选的自由,甚至允许自己在工作过程中修改它们,乃至发明另一些……)。我注意到,我面对有待探究的数学实质时那种完全不同的态度——因而也是面对学生时那种完全不同的态度——几乎必定会引发一种失措,而焦虑就是它的征兆之一。像一切焦虑一样,它倾向于给自己找一张面孔,投射到某个外部的“理由”上,不管那理由说得通说不通。焦虑最常见的一张面孔,恰恰就是恐惧。\par 这类困难在我教学生涯的第一个时期几乎没有出现,也许只有两例例外:那两次“师生”关系没超过几个星期;也许还有(我说不准)“忧伤的学生”那一例——他也许感到自己被“铆”在一个完全不启发他的题目上,而他本来完全有更换的自由。在那个(我也谈到过的)长期被某种怯场困扰的学生身上,原因显然在别处。他的工作丝毫没有卡住,恰恰相反,他对选定的主题极为自在,并在其上做出了一项有规模的基础性工作。那个时期我的学生大多是巴黎高师的旧生,他们与亨利·嘉当(Henri Cartan)的接触,已经让他们见识了另一种“接近数学”的方式。与之(可以说)相反的另一极,是我在蒙彼利埃大学任教的第二个时期:正是在一年级学生那里,我所说的焦虑最少干扰思考性的工作。在这些学生中的许多人身上,面对一种不同进路时的惊讶既不引起焦虑,也不引起封闭,相反带来的是敞开和兴致——终于可以做点有趣的事了!据我的观察,几年大学对一个学生创造性禀性的作用是彻底而毁灭性的。奇怪的是,在这方面,长年中学的作用似乎相对无害。原因也许是:\par 大学那几年正落在这样一个年纪:我们身上与生俱来的创造力,最迟必须通过个人的工作表达出来,否则就要永远沉船——至少就智力上的创造性工作而言是如此。肯定是出于一种健康的本能,我在做学生的那些年(也是在蒙彼利埃大学),几乎不踏进课堂,把几乎全部精力都投给了个人的数学思考。} (23'),
% ---------------- p.176 ----------------
\origpage{67} 其余则都是如此——那些本应正式在我指导下准备研究工作的学生,无论哪一个层次,也一样。与从前的学生相比,一个不同之处(还有许多别的!)是:我们的关系不那么局限于共同的数学工作。学生与我之间的交流,常常以不那么表面的方式牵涉到我们两个人本身\fnoteP{10}{(23$^{v}$)}{这种差别有一个格外醒目的征兆,出现在“外国人事件”之际——我曾有机会谈到过(第24节)。当时,许多与我完全陌生的人向我表达了同情,而我记不起1970年之前的哪个学生想到过以这种方式表示什么,更不用说在我投身的那场行动中向我提供任何帮助。相反,在我看来,第二个时期的学生或旧学生里,没有一个不曾向我表达同情与声援,还有几个积极参加了我在当地开展的运动。在这个狭小圈子之外,1945年法令那件事也在许多仅仅知道我名字的大学生中间激起了一定的波澜;在我被传唤的那一天,有不少人来到法院,以示声援。这后一个情形还提示:我所察觉到的“之前”与“1970年之后”学生态度上的差别,也许较少表达他们与我之间关系的差别,而更多表达的是一种心态的差别。显然,我“之前”的那些学生已经成了重要人物,而要让重要人物动心,是需要很多条件的……但我1970年离开IHÉS、投身战斗行动那一幕,似乎表明并不只是这一点。那是个他们中间还没有谁太像重要人物的时刻,然而我记不起他们中有谁对我所投身的活动表现出丝毫兴趣。我倒认为,那一定让他们所有人,无一例外地感到不自在。这仍然合乎“心态差别”的方向,但这不能只归因于社会地位的差别。} (23'')。因此,在我教学生涯的第二个时期,与某些学生的关系中的冲突成分表现得更清楚、更直接,甚至更激烈,这就不奇怪了。在第一个时期的旧学生里,有两个人后来表现出系统而毫不含糊的对抗态度(我曾顺带提到过),不过始终停留在未言说的层面,也许甚至停留在无意识的层面。在更长一些的第二个时期,有三个学生让我遭遇了对抗。其中两个,表现得相当尖锐。

这些学生中的一个,对抗是在一段最友好的关系里一夜之间冒出来的——那是在这位朋友不再是我的学生许多年之后。我猜测,冲突的原因与其说是我那不堪形容的举止和人格,不如说是一份压抑已久的不满:他的工作(曾经很出色)没有得到他本有权期待的那份接待。这就是“1970年之后”以我为老板的那份可疑特权的背面;他必定为此怨我,尽管即便在内心深处也不大肯向自己承认。

另一个学生,尖锐的对抗在工作一年半之后就已经出现,而此前的氛围看上去非常亲切。这是第一次也是唯一一次,学生与我之间的关系困难出现在他还是学生身份的时候。它使共同工作的继续成为不可能——而那份工作本来开局吉利,带着最好的兆头般的热情,而且应当说,题目是一个绝好的思考主题。我的感觉是:这位年轻研究者身上有一种暗中作祟的自信不足——不相信自己有能力做出好工作(这能力在我看来毫无疑问);而对抗以尖锐的音调爆发出来,是某种“向前逃窜”,为的是抢在一种他惧怕的失败之前动手,并把责任预先推到一位可憎的老板头上\fnoteP{11}{(23$'''$)}{\emph{两兄弟}\par 这个学生身上的对抗,一上来就采取了“阶级对抗”的形式:我是那个对他的数学前途拥有“生杀之权”的“老板”,可以凭自己的好恶决定……当然,事情本身只能证实这种看法,因为我很快就不再承担指导这个学生的责任(那已变得令人难受)。这把他置于一个微妙的处境——在眼下的世道里,找一个“老板”可不是那么显而易见的事,尤其是题目已经选定的时候。在另一个学生身上,合法期待落了空,对抗采取了类似的形式。我被感受为专横的“学阀”,不能容忍那些被他视为下属的人(学生或级别较低的同事)有任何反驳。\par 这样一种“阶级态度”,在我与第一个时期学生的关系里,从未有过一丝一毫的表现。显而易见的原因是:在1970年以前的形势下,学生通过论文之后毫无疑问会得到一个副教授(maître de conférences)的职位,因而享有与我相同的社会地位,即“大学教授”。数字最能说明问题:1970年以前开始跟我做工作的十一个学生,一完成工作就得到了副教授职位;而在我的指导下或多或少工作过的那二十来个学生里,没有一个获得这样的职位。诚然,他们之中只有两个人有足够的动力去做国家博士论文(而且两人的论文都很出色)。\par 因此,在这第二个时期,某些暧昧(其深层根源始终被遮蔽着)采取了阶级对抗的形式,采取了对“老板”的不信任(这种不信任被呈现、也被感受为“本能的”“发自肺腑的”)——这就不足为奇了。在那些或多或少当过学生样子的人里,有一个与我的友好关系持续了十来年,没有出现过表面上的对抗插曲,却仍带着同一种暧昧,表现为一种不信任的态度,被“储备”在明显的善意之后。说实话,我从未被这种例行公事式的“不信任”骗过;它在我看来,尤其是这位朋友认为应当给自己找的一个理由,好让自己不去冒险越出他选定的、界限分明的领地——无论在职业生活里,还是在整个生活里都是如此。可这件事他本来可以自由地去做,没有人(顶多只有他自己!)会要他交代。\par 这三个例子,也是我整个教书经验中仅有的、学生(或多少像学生的人)与我之间的关系中某种暧昧以“阶级态度”表现出来的情形。当这种态度出现在大学“团体”内部的同事之间时,显得格外暧昧:在这样一个团体里,双方都享有与常人处境相比极为悬殊的特权,而这些特权使级别(以及薪金)的差别显得相对微不足道。我还注意到,一旦当事人自己也被提升到前一天他还在指责别人的那种处境,这些态度就像中了魔法一样消失了(而且理所当然!)。\par 此外,在我于数学界内部(常常也在外部)目睹的多数——如果不是全部——冲突局面里,我都辨认出一种类似的暧昧。那些“有位置的人”,无论其级别是否合乎他们的期待(合理与否),都享有相当闻所未闻的特权,那是任何别的职业或生涯都无法提供的。那些没有位置的人,则向往同样的安全感和同样的特权(这并不必然妨碍他们对数学本身感兴趣,有时还做出漂亮的工作)。在眼下的世道里,抢位置竞争激烈,没有位置的人常被当作流浪汉对待:我不止一次感到那种默契——在乐于羞辱人的人与被羞辱的人之间;后者咽下去,缩起来。他苦涩与敌意的真正对象,并不是那个行使了权力的人,而恰恰是他自己:他自己缩了起来,并且把对方奉为那种他乐意行使的权力的持有者。乐于羞辱人的人,也是那个报复并补偿(却永远抹不掉……)一场长期承受、早已深埋并被遗忘的羞辱的人。而默许自己被羞辱的人,是他的兄弟和效仿者:他暗中羡慕对方,并在苦涩里埋下那场羞辱,以及羞辱带给他的、关于他自己的那条卑微的消息。} (23''')。

% ---------------- p.177 ----------------
\origpage{68} 在学生与我之间所有这些冲突的显现里,有一个共同点——我教人做数学家这一行快二十五年了——就是强烈的暧昧。无一例外,对抗都是在事后才表现出来的,往往来得很阴险,出现在一段就其本身而言无可怀疑的好感关系里。我甚至可以说,在所有这些情形里,正如在许多别的、并未表现出明确对抗成分的情形里一样,我这个人在他们身上曾产生、并且仍在产生强烈的吸引力。肯定正是这吸引力本身的力量,也滋养了对抗的力量,并保证了它的延续。在对抗采取强烈反感、愤然排斥形式的情形里,肯定也仍然如此;同样,在另一个相反的极端情形里也是如此:在友谊般的敬重这件必不可少的外衣之下,(机会合适时)表达的是一种潇洒而精心调配过的轻蔑的做作……

说实话,这样的暧昧局面并不限于我与某些学生或旧学生的关系。事实上,在我整个成年生活中,从三十岁起(也就是说,从我母亲去世起),它们一直大量存在。在我的情感生活或夫妻生活里是这样,在我与男人的关系里也是这样,更确切地说,主要是与那些明显比我年轻的男人。我终于明白:我身上有某种东西——是天生的还是后天得来的,我说不清——似乎使我倾向于充当父亲的形象。看来我必定有理想的体格和相宜的振动频率,正好造就一个完美的养父!应当说,父亲这个角色我穿起来就像手套一样合手——仿佛它生来就是
%
% ---------------- p.178 ----------------
\origpage{69}\origpage{70} 我的。我不会去数自己有多少次对另一个人进入这样的角色,双方默契得完美无间。大多数时候,这种父—子或父—女的角色分配始终留在未言说之中,甚至留在无意识之中;但也有些时候,它被或多或少清楚地说了出来。在某些情形里,我甚至并没有进入什么游戏,却已经成了父亲的样子——我确信如此——对自己意识层面与无意识层面正在酝酿的事情都一无所知。

我第一次察觉自己扮演养父的角色,是在1972年,“生存与生活”时期,当时我突然面对一位年轻朋友的强烈排斥态度。(有趣的巧合:那是个脱离了常规的数学系学生!)我对第三者的某种举止让他失望了。我想,我本来会毫不困难地承认他的失望是有道理的,承认我在那件事上缺少慷慨——但那份反应的激烈,当时简直把我打懵了。那就像一阵突然腾起的猛烈仇恨之火,而它几乎马上又落了下去,因为已经清楚:他并没有真正把我掀下马。(只差一点点,但这一点我没有让他看出来……)。我不知道当时怎么会直觉到:他把我这个人——被他恰如其分地理想化了的——当成了他与父亲之间未解决冲突的投射对象。这个突然的直觉后来被忘掉了,却没能阻止我在此后许多年里仍然带着同样的确信继续进入父亲的角色,丝毫没有提防。当然,后来当我面对冲突的征兆——阴险的或激烈的——时,也总是同样痛苦的惊讶,不敢相信自己的眼睛,也不敢相信别的。

正是经过一场持续六七个月的、关于我父母生平的紧张独处工作——它让我在一种意想不到的光照下看见他们这个人——我才明白:这种养父的角色是虚幻的;它要替代一个确实存在、却被宣布为(哪怕只是通过默契)“失职”的真正父亲(而且要替代得更好,这是预先就说定了的!)。那是在帮别人回避冲突,把冲突从它所在的地方——比如说,在他与父亲的关系里——转移到第三者(此处即我本人)身上,而第三者与它全然无关。自从那场沉思(它发生在1979年8月到1980年3月)之后,我对自己保持警觉,不再闭着眼睛听任自己那份不合时宜的父性天职发作。这并没有阻止那种虚假局面重演(比如我与那个不得不停止合作的学生之间的关系)——但我相信,如今已没有我这方面的默契。

如果把那个合法期待落了空的学生放在一边,那么对我来说毫无疑问:在我遭遇学生或旧学生对抗的所有其他情形里,那都是同一种“与父亲的冲突”原型的重演:那个既受仰慕又受畏惧、既被爱又被恨的父亲——那个要去对抗、去战胜、去取代、也许还要去羞辱的“他”……但也是那个你暗中想成为的人——夺走他的一股力量,据为己有——另一个自我,被畏惧、被憎恨、被逃避……

\section{(30) 敌父(2)}

在某些旧学生与我之间制造出对抗的,并不是1970年那个大转折——当时的背景是一片牧歌式、无云的过去。它只是让一些对抗变得可见,而那些对抗本来很难在一种典型的、更为循规蹈矩的老板—学生(或前老板—前学生)关系框架里表达出来。我猜想,这类冲突在科学界想必并不罕见,只是它们通常表现得比我所卷入的那些关系更迂回、更难辨认。

回想起来,我终究不觉得在与学生的那些关系里,我曾那么倾向于进入父亲的角色——甚至,我抓不住一个哪怕多少朝着那个方向的记忆。

% ---------------- p.179 ----------------
\origpage{71} 就我这个人而言,我似乎觉得,我投入一段师生关系的几乎全部精力,就是我也投入数学、投入实现一个宏大纲领的那份精力。在第一个时期,我只看到一个例子:我身上对某个学生这个人的兴趣——一种亲近或好感的性质——其力量堪与(如果不说是等同)数学兴趣相比。但即便在那个例子里,我也不觉得我对他进入了父亲的角色。至于我对他、或对其他学生,在某个层次上可能施加的影响,那类事情在我与学生的关系里是我从不留意的。(即便今天,我也倾向于不留心,无论对近几年跟我做过工作的学生,还是对别的人。)当然,在所有这些情形里,学生与我之间的关系绝不是“对称”的:至少在师生关系持续的那段时间里(而且多半还要更久),一个学生在我生命里的分量,与我必定在他生命里占据的分量无法相比;关系在我身上和他身上所调动的心理力量也无法相比。除了那五六个以清晰可辨的对抗征兆表现出这些力量的情形,我意识到:在二十多年教书生涯里,我那些不同的学生、后来是旧学生与我之间关系的性质,对我来说始终是个完全的谜!再说,探究这些谜也不是我分内的事,那更多是各人自己那一份的事。不过,既然关心起自己这个人来,也许还有些比“他与旧老板关系的来龙去脉”更灼人的东西可以看。无论如何,尽管我在学生面前没有表现出任何进入父亲角色的倾向,我多半还是常常在他们眼里多少充当了养父——鉴于我前面谈过的自己那份特殊的心理“剖面”,也鉴于一种局面的内在动力:在那种局面里,我起码不能不显出一个长辈的样子。

无论如何,在我提到的几个例子里,学生与我之间关系的这种特殊色调,在我看来毫无疑问。在职业生活之外,还有许多别的例子:不管有没有我的默契,我显然在比我年轻的男人或女人面前充当了养父——他们被这个人吸引,首先与我以相互的好感相连,而绝无亲缘关系。至于我自己的孩子们,我对他们的父性之弦是强烈的;从很小的时候起,他们在我生命里就占据着重要的位置。奇怪的讽刺是:我的五个孩子,没有一个接受“有我做父亲”这件事。在我能够近距离了解的四个孩子的生活里——尤其是近几年——他们与我关系中的这道分裂,是他们自身内在一道深刻分裂的反映;尤其是一种拒绝:拒绝他们身上一切与我这个父亲相像的东西……但这里不是探究这道分裂之根的地方:那些根既扎在一段被撕裂的童年里,也扎在我的童年、我父母的童年里;同样也扎在母亲的童年、她父母的童年里。这里也不是衡量其后果的地方——在他们自己的生活里,或在他们孩子的生命里……

\section{(31) 使人泄气的权力}

为了结束这趟粗略的巡游——穿过1948到1970年间我在数学界里的种种关系——我还得谈谈我与更年轻的数学家的关系:他们或多或少是新手,因而严格说来没有“同事”身份,而我也不曾对他们扮演“老板”的角色。所以,这里说的是些年轻研究者:我在IHÉS的讨论班上与他们相遇一两年,或在哈佛及别处的这类课程、讨论班上,有时也通过
%
% ---------------- p.180 ----------------
\origpage{72}\origpage{73} 通信——比如我收到一位年轻作者寄来的工作,他期待我给出评论,肯定也期待鼓励。

与初入门研究者的关系,属于一个不如“学生们的老板”那么显眼的角色,但同样重要——这是我后来才明白的。那个年代,我并不像这六七年以来那样意识到:对一位有名望的数学家来说,这个角色代表着一种相当大的权力。首先是鼓励、激励的权力:它既存在于明显出色的工作那里(尽管那工作也许因表述笨拙或“行内功夫”不足而吃亏),也存在于仅仅扎实的工作那里;它甚至存在于那种只代表一份极微薄贡献的工作那里——按一位正当壮盛、手段强劲、对主题经验老到、见闻广博的前辈的标准,那贡献甚或微不足道,乃至等于零。只要送到我们手里的工作是认真写成的——这通常从头几页就看得出来——鼓励的权力就在那里。

而使人泄气的权力同样存在,而且无论面对什么样的工作,都可以随意行使。这就是柯西对伽罗瓦、高斯对雅可比行使过的那种权力——它可不是从昨天才有的,杰出而令人畏惧的人物一直都在使用它!历史之所以把这两桩事传给我们,是因为付出了代价的这两个人有足够的信念和自信继续走自己的路,尽管那些当时在数学界呼风唤雨的人怀着的权威毫无善意。雅可比找到了一份杂志发表他的想法,伽罗瓦找到的则是他最后一封信的纸页,权当“杂志”。

今天,对一个无名或不太有名的数学家来说,要让人知道自己,肯定比上个世纪更难。而有名望的数学家的权力,不只在心理层面,也在实践层面。他有权接受或拒绝一篇工作,也就是说:给不给它发表的支持。对也好错也好,我觉得“在我那个时代”,五六十年代,拒绝并非不可上诉——如果工作呈现出“值得注意”的结果,它还有机会得到另一位权威的支持。今天肯定不再如此:如今甚至很难找到一位有影响的数学家,肯(以他乐意采取的心情)去通读一篇属于他专业领域的工作——如果作者尚未获得名声,或者没有一位他认识的同行推荐。

近几年,我曾见有影响、出色的数学家行使他们使人泄气的权力,拒绝这个、拒绝那个:既针对某些显然应当做出来的扎实工作,也针对某些规模可观、清楚显示出作者力量与原创性的工作。有好几次,这样行使自己任意权力的人,竟是我的一个旧学生。这无疑是我作为数学家的生涯里经历过的最苦涩的经验。

但我又跑题了!我本来要考察的是:在我诚心诚意地扮演“知名数学家”角色的那些年月里,我是怎样使用自己手中那份鼓励人和使人泄气的权力的。我还该补充:在1970年以后我的科学活动所处的那个更卑微的层次上——在外省大学里作为普通教师中的一员——这种权力并没有因此就不存在,无论是对我的学生,还是(说实话,很少)对偶尔通信的人。但就我眼前的题目而言,只有我数学生涯的第一个时期才是要紧的。

至于我与学生的关系,从我带过的第一个直到今天,我相信可以毫无保留地说:我做了力所能及的一切去鼓励他们
%
% ---------------- p.181 ----------------
\origpage{74} 在自己选定的工作上往前走\fnoteP{12}{(23$^{iv}$)}{\emph{一次教学的失败(1)}\par 这些字句写下之后,我有机会与1970年之后的两位旧学生谈过,想同他们一起探究我在蒙彼利埃大学研究层次上教学失败的原因。他们对我说:我那种倾向——低估同化那些对我熟悉、对他们却不熟悉的技术对他们可能意味着的困难——在他们身上产生了令人泄气的效果,因为他们感到自己始终达不到我对他们的期望。此外(这件事在我看来意义还要大),他们有时感到受了挫伤:当我把袖子里藏着的一个成形陈述“卖给他们”、向他们泄底的时候,而不是留给他们乐趣,让他们在已经十分接近的时刻靠自己的办法把它发现出来。那之后,他们剩下的就只有做“练习”了(而那并不特别让他们着迷)——也就是证明那个陈述。我先前在一则注(注21)里指出过的自己身上那份“缺乏慷慨”,正落在这一处,当时没有就此多说。主要正是这类扫兴之事,构成了我在他们两人身上——在那么出色的开端之后——研究兴趣消失这件事上的个人贡献。\par 我意识到,我在1970年以前并不比之后更慷慨。如果说那时我没有遇到同样的困难,那无疑是因为:那个时期来找我的那类学生,有足够的动力,甚至能从一场“长长的练习”里品出魅力——那是学习本行、并一路学到一大堆东西的机会;同样,对于一个我“泄了底”的起步陈述,他们也能靠自己的办法从中抽出一大群远远超出第一个的陈述。当我换了教学地点,我在向新学生提议的思考主题上做了必要的调整,办法是选择那些可以被直接直觉把握的数学对象,不依赖任何技术行囊。但这必要的调整本身并不够,因为禀性上的差别(我的新学生与从前的学生相比)比单纯的行囊差别还要重要。这还印证了前面(第25段开头)说过的一点:我在“导师”这个角色上有某种不足;这不足在我教学生涯的第二个时期,比在第一个时期显露得强烈得多。} (23'''')。即便在今天,老板与学生的关系里恐怕也很少是别的样子,尤其当这位老板拥有手段,能够培养出色的学生,并借他们的协助开垦大片待耕的荒地。几乎难以置信、然而千真万确的是:甚至存在这样一种极端情形——一位声望卓著的老板,乐于把曾经激励过他自己年轻时代的数学激情,在才华出众的学生身上熄灭掉。

但我又跑题了!现在要考察的,是我与那些不是学生的年轻研究者的关系。在那样一些关系里,名人身上的自我中心力量,不那么容易把他推向鼓励的方向,因为那个找上他的无名年轻人的成功,对他的荣耀毫无增益,或增益甚微。恰恰相反,我认为:在缺乏真正善意的情况下,仅凭自我中心力量的运作,几乎无一例外地会把人推向相反的方向——去行使使人泄气、拒绝的权力。在我看来,这恰恰就是那条在社会一切部门都能看到的普遍法则:证明自己重要的自我中心欲望,以及满足它时伴随的秘密快乐,通常在这样一种情形下更强、更受享用——手里的权力找到机会去造成别人的失望,甚至羞辱,而不是相反。这条法则在某些特殊处境里表现得格外粗暴,比如战争,比如集中营的世界,监狱或精神病院的世界,或者干脆就是我们这样一个国家里任何人都可能进的医院……但即便在最日常的处境里,我们每个人都遇到过印证这条法则的态度与行为。对这些态度的矫正,首先来自文化的矫正:在某个圈子里,关于什么算“正常”或“可接受”的行为有一种共识;其次则是非自我中心性质的力量,比如对某个特定的人的好感,或者有时,一种自发的善意,甚至与它所指向的人无关。这样的善意无疑是稀罕之物,无论你去哪个圈子里找。至于数学圈里的文化矫正,在我看来,在过去的二十年里已经被大大侵蚀了。至少在我所熟悉的
% ---------------- p.182 ----------------
\origpage{75}\origpage{76} 那些圈子里,情况肯定如此。

我执意跑题,真是没办法!我本来要说的不是一篇关于这个世纪的议论,而是一场对我自己、对我与那些或多或少刚刚入门、又不是我学生的研究者之间关系的沉思。我不认为我提到的那条“法则”在那些关系里找到过表现的机会。由于一些不必在此考察的原因,看来自我中心的力量,虽然在我身上和在任何人身上一样强,在我一生中却没有走那条路去损害别人(除了几桩追溯到童年的情形)。我甚至敢说——既然有机会考察这件事——我面对他人的基本调性是一种善意的调性:能帮的时候就帮,能缓解的时候就缓解,有能力鼓励的时候就鼓励。即便在与那位“不知疲倦的朋友”如此深刻分裂的关系里,我身上的自负也从未把我误导到这样的地步:让我想到(哪怕是无意地)去伤害他。(我本来是有可能那样做的,而且当然会“怀着世上最安宁的良心”。)我相信,在大多数情形里,这种普遍的善意(哪怕只是浮在皮肤表面的一点点)也标记了我在数学界里的关系,包括与那些虽不在学生之列、却可能需要我支持或鼓励的初入门数学家。

我相信,至少在五十年代、直到六十年代初,无一例外都是如此。我觉得至少在那个时期,这份善意并不限于像广中平祐(Heisuke Hironaka)或迈克·阿廷(Mike Artin)那样明显出色的年轻人(当时还没有任何名声证明他们的本领)。但它在六十年代也许或多或少地淡去了,在自我中心力量的作用下。关于这一点,若有任何见证传到我这里,我会格外感激。

我的记忆只交给我一个确切的例子,我这就说;而在这个例子之外,就是那团著名的“雾”——它不凝结成任何别的确切事例或事实,却交给我某种内心的态度。当另一位数学家“踩到我的地盘上”、却装作不问我一声,好像他在自己家里似的——这个毛头小子!——我会感到某种恼火。那多半是年轻人的情形:不大在行,却碰巧重新找到一些我多年来、而且是从高处就知道的东西,有时还偏偏在一些很特别的情形里。我想,这种事不至于常常发生,但也许还是有两三次,也许是四次,我说不准。如我刚才所说,我只记得一个具体的例子,也许是因为那局面与同一位年轻数学家反复出现过好几次,形式各不相同。我可以说:这位年轻研究者——他所属的大学在国外——在各方面都极其得体:他把自己刚做好的工作寄给我,而我理应是消息最灵通的人。每一次,我都反应得极为冷淡,理由就是我说的那个。我甚至说不准自己有没有坦率地告诉他:他做的东西我早就知道了,因此他不至少在前言里向我行个小小的屈膝礼就发表,让我很不痛快。当然,假如他是我的学生,这种作者的虚荣就不会起那么大作用:一方面因为与学生之间已经建立起好感,另一方面因为无论如何不言而喻,学生的工作里也有老板的想法,除非另有说明!我相信这局面与同一位研究者发生过两次,也许甚至三次;每一次我的态度都一样冷淡,一样令人泄气。如果我记得不错,我从未答应过推荐这位研究者的工作到某家杂志发表,也从未答应参加他的论文答辩委员会(我记得这个问题曾被提起过)。几乎可以说,我像是打定主意挑他当出气筒。最妙的是,
%
% ---------------- p.183 ----------------
\origpage{77} 他的工作每一次都完全站得住脚——我相信写得用心,我也没有任何理由假定那些想法不是他自己找到的:那些想法当时还不怎么满街乱跑,只是(或多或少)“众所周知”于一小撮在行的人之间,比如塞尔、卡蒂耶、我,还有一两个人。我不明白的是:这位年轻同事(他当然最终有了论文,也有了一个受之无愧的职位)怎么就没厌烦一次次找上我——那个每次“给他冷脸”的人,而且他显然从未为此怨过我。我还记得有一次,他对我的迟疑表示惊讶;显然他不明白发生了什么。要是他等着我的解释,那可有他难受的!他有一张好看的脸,有点像古典希腊的,非常年轻——线条相当柔和、安详,透出一种内在的平静……现在,我第一次试着去捕捉他这个人和他的相貌所散发的印象,忽然意识到:他真的很像我说过的那位“不知疲倦的朋友”;他们本可以是兄弟——这位与我同龄的朋友,调子是含笑的那种;而这位研究者,比他年轻二十岁,调子略略低沉,却绝不忧伤。这种相似未必没有起过作用:我把一份在另一个人身上找不到机会表达的轻蔑,投射到了这一位身上——因为另一个人以那样忠实的友谊征兆把我解除了武装!而我也确实得长出一副多么厚的甲壳,才没有被这个必定招人喜爱的年轻人的明显诚意和把事情做好的意愿所打动:他不厌其烦地再来,而我连一个微笑都不肯赏他!

\section{(32) 数学家的伦理}

我昨天讲的那个例子,现在我终于肯费心把它白纸黑字记下来,它在我看来具有相当大的意义,在某些方面比我先前讲到的另外三个例子(无疑同样典型)更大——在那三例里,自负的力量深深地扰乱了我身上一种自然的善意与尊重的态度。这一次,我动用了一个实实在在的权力位置(同时像所有人一样,装作不知道这份权力),用它去使一位善意的研究者泄气,拒绝一篇本应发表的工作。这就是人们所说的滥用权力。它并不因为不落入刑法某一条的制裁范围就不那么昭然。幸而当时的形势不如今天严酷,所以这位研究者能够——我相信,没费太大力气——在几位比我更善意的同行的支持下把工作发表出来;他作为数学家的生涯,也没有因为我那滥用权力的行为而受到严重扰乱,更不用说被毁掉。事后我为此高兴,却并不想把它当作“减轻情节”。在更严酷的形势里,我也许会更当心——但这只是一个推测,在这里没有多少关系。不过我仍相信自己可以说:我心里并没有暗藏的恶意,没有由我说的那种恼火引起的害人之心。我对那恼火的反应是“发自肺腑的”,对自己没有丝毫批评的意愿,更没有丝毫意愿去看一眼自己心里发生了什么,哪怕只是看一眼我的反应可能在另一个人的生命里有什么分量。我没有衡量自己手中的权力;与这权力相随的责任这个念头,在那段关系里从未掠过我的心头。这是一桩不负责任行为的典型,就像街头巷尾随处可遇的那样,在科学界和在别处一样。

我记住的这个同类情形中唯一的例子,也许只是几个类似情形中的一个极端例子。引发一种毫无善意的态度的,是虚荣心受到冒犯后的恼火:它
%
% ---------------- p.184 ----------------
\origpage{78}\origpage{79} 容不得“随便哪个初来乍到的人”僭取权利,闯进禁猎区,打走一点只属于这些地盘之主的小猎物……这种恼火有的是现成的合理化,而且派头更高贵,这可想而知。被摆出来的不是我这卑微的个人——不!——而是对艺术、对数学的爱;这个年轻人连“是个天才”这条借口都没有,倒更像笨手笨脚的那一类,他会把一切都弄糟,我们可倒霉了;他要是把事情做得比我好倒也罢了,可我筹划的那些漂亮安排全泡汤了,说真的,总得有点不知羞才行……!而始终如水印般透出的,是那支功绩主义的主导动机:只有最好的那些人(比如我)在我的地盘上才有居留权,或者那些把自己置于这类人中某一位保护之下的人!(至于那种不那么常见的情形——确实是另一位大人物踩到我的地盘上——那又是另一回事了:一天自有一天的烦恼!)在这个具体例子里,还有(对此我如今已不大有疑问)另一股朝同一方向起作用的力量,而它完全是无意识的;这股力量在我与起步时那位不知疲倦的朋友的关系里,已经起过很强的作用:对某一类人的排斥自动机制——那类人不合我从母亲那里接过来的“男子气”标准。但这个情形,虽然对理解我自己有它的意义和趣味,与我眼前的题目却关系不大:我的题目是要在我自己身上,在我还属于某个圈子时曾有过的态度和行为里,找出我今天在那个圈子里所看到的深刻败坏的典型征兆。

我刚考察的这个例子,之所以在我看来比我其他那些缺乏善意与尊重的例子意义更大,是因为正是在它这里,数学家这一行某种基本的伦理遭到了侵犯\fnoteP{13}{(24)}{我想说的这种伦理,同样适用于任何别的围绕研究活动形成的圈子;在那里,把自己的结果公之于众、并从中获得承认的可能性,对每个成员的社会地位、甚至对作为该圈子成员的“存活”,都是“生死攸关”的问题,连带对他和他的家庭意味着一切后果。} (24)。在我起步时接纳我的那个圈子里——也就是布尔巴基及其亲近者的圈子——我想说的这种伦理大体上始终是隐含的,但它毕竟在场、有生命,而且(在我看来)是不可触动的共识的对象。据我记忆所及,唯一把它用清楚明确的话说出来的人是迪厄多内,那无疑是我在南锡做客的最初几次之一。他也许在别的场合又回到这一点。显然他觉得这是一件重要的事;而我当时必定感受到了他对此的看重,因为三十五年后的今天我还能记得。仅仅凭借我那些年长的前辈这个群体的道德权威,以及显然表达了群体共识的迪厄多内,我必定已经默默地把这种伦理认作自己的,却从未给它一刻思考,也从未明白它的重要性何在。说实话,我甚至不会想到:对它加以思考也许有用——因为我早已确信,我的父母和我自己这个人,各自都是伦理态度的完美化身(或者相差无几),负责、全面,而且经得起一切考验\fnoteP{14}{(25)}{\emph{职业伦理的共识——兼及信息控制}\par 除了与迪厄多内的那次谈话,在我作为数学家的生涯里,我记不起有哪次谈话——我参与的或旁听的——谈到这一行的伦理、谈到行业成员之间关系中的“游戏规则”。(这里我排除那些关于科学家与军事机构合作的讨论,它们发生在七十年代初“生存与生活”运动前后。那些讨论其实并不涉及数学家之间的关系。我在“生存与生活”里的许多朋友,包括谢瓦莱和盖吉,都觉得:我当时——尤其起初——把重点放在这个我格外敏感的问题上,使我远离了更本质的日常现实,正是我在眼下这份反思里考察的那类现实。)在学生与我之间,从未谈过这些事。那种心照不宣的共识,我想只限于这一条规则:不把别人的、自己可能已经知道的念头说成是自己的。这共识,在我看来自古就存在,而且至今没有任何科学圈子里的人质疑过。但若缺少另一条补充规则——保障每个研究者有可能让人知道自己的想法和结果——第一条规则就只是一纸空文。在今天科学界,处于声望与权力地位的人,对科学信息握有一种任意处置的控制权。在我所熟悉的那个圈子里,这种控制不再被迪厄多内所说的那种共识所缓和;而那共识,也许从未存在于他充当代言人的那个狭小团体之外。处于权力地位的科学家,几乎能收到他认为值得收到的一切信息(往往还超出),而且对其中很大一部分,他有权阻止发表,同时坐享其利,把那些信息斥为“没有意思”、“多少已知”、“trivial”等等……我在注(27)里还要回到这个局面。} (25)。

% ---------------- p.185 ----------------
\origpage{80} 再说,迪厄多内并没有对我长篇大论——那既不是他的风格,也不是他布尔巴基的任何朋友们的风格。他多半是顺口提起,像提起一件理应不言自明的事。他只是强调一条再简单不过、表面上全然无害的规则,就是这样:任何人发现了一个值得注意的结果,就应当有权利、也有可能把它发表,唯一的条件是这结果尚未成为某份出版物的对象。所以,即便这个结果为一人或数人所知,只要他们不曾费心把它白纸黑字写下来并发表,从而使它可供(哼!)“数学共同体”使用,那么任何别的人(言下之意:包括那位著名的“初来乍到者”)用自己的办法找到这个结果(言下之意:不管他用的是什么办法、什么观点和光照,也不管这些在那些理应比他更在行的人看来是否“褊狭”……),就应当有可能按自己的办法和光照把它发表。我记得迪厄多内还补充说:这条规则若不被遵守,就会给最恶劣的滥用敞开大门——很可能正是在那个场合、经他的口,我才知道了高斯拒绝雅可比工作那桩历史案例——借口是雅可比的想法他早已知道。

这条如此简单的规则,是对迪厄多内身上(以及布尔巴基其他成员身上)那种“功绩主义”态度——正如在我身上一样——的根本矫正。遵守这条规则,是一种廉洁的保证。我很高兴能够说,从至今传到我这里的一切来看,这种根本的廉洁在布尔巴基最初群体的每一位成员身上都保持完好\fnoteP{15}{(26)}{布尔巴基的“创始成员”是亨利·嘉当、克劳德·谢瓦莱、让·德尔萨特(Jean Delsarte)、让·迪厄多内、安德烈·韦伊。他们全都健在,只有德尔萨特例外——他在五十年代英年早逝,而那个时期,这一行的伦理总体上还受到尊重。} (26)。我注意到,对于曾经属于布尔巴基群体或圈子的别的数学家,情况并不都是这样。而在我自己这个人身上,它并没有保持完好。\par{}重读本文时,我曾想删掉这一段:在里头我可能给人这样的印象,像在颁发“廉洁”(或不廉洁)证书,而当事人对此毫不在乎,也不是我该做的事。这段文字可能引起的保留,肯定是有道理的。但我还是把它留着,为的是见证的真实性,也因为这段文字确实再现了我的感受,即便这些感受并不合宜。

迪厄多内用再平实不过的话向我讲的那种伦理,作为一种特定圈子的伦理,已经死了。或者说,那个圈子本身也死了,与那构成它灵魂的廉洁一同死去。这份廉洁在某些孤立的人身上保存了下来;在另一些它曾经败坏的人身上,它重现了,或者将会重现。它在某个人身上出现或消失,是各人精神历险的关键插曲之一。但这历险得以展开的舞台,已经深刻地改变了。一个曾经接纳我、被我认作自己的、我曾暗自引以为傲的圈子,不复存在了。使它有价值的东西,在我自己身上死去了,或者至少被另一种性质的力量侵入并取代——而那远在规范它的默会伦理于使用惯例与信仰表白中被公开否弃之前。如果我此后能够惊讶、能够愤慨,那是出于故意的无知。从这个曾属于我的圈子里回到我身上的东西,本来有一个关于我自己的消息要带给我,而我乐意把它一直回避到今天。

% ---------------- p.186 ----------------
% (本页仅存书眉“8. Maître et élèves”,无正文)


% ================================================================
%  第 9 章 收获 (原书 p.187-196, 边码 p.81-p.94) —— 前半
%  源文: _work/src1/task-ch09a.md
% ================================================================
\chapter{收获}

% ---------------- p.187 ----------------
\section{(33) 笔记——或新伦理}

\origpage{81}当然,一条职业伦理的规则,只有凭一种内在的态度才有意义,那态度就是它的灵魂。它无力创造它所力求表达的那种尊重与公正的态度,最多只能在一个普遍认可这条规则的圈子里,帮助这样的态度维持下去。一旦缺少了内在态度,哪怕这条规则还挂在人们嘴边,它也已失去全部意义、全部价值。

任何注释,无论多么一丝不苟、多么细致入微,也改变不了这一点。

不久前,我昔日的朋友与伙伴中有一位好心地向我解释:如今这年头,唉,数学产出涨得没边儿(这是谁都知道的),“人们”实在不得不——愿意也好,不愿意也罢——在写出来、交上来准备发表的文章里做一番严厉的筛选,只发表其中很小的一部分。他说这话时一脸真诚的遗憾,仿佛他自己也多少是这不可避免的宿命的受害者——同时,他那神情里还带点这样的意思:他自己就属于“法国那六七个人”(唉,是不幸,可事情就是这样!),由他们决定哪些文章能发表、哪些不能。我上了年纪,话少了,便只是默默听着。这个题目上有许多话可说,但我知道那会是白费口舌。一两个月后我还得知,这位同行几年前曾拒绝推荐一篇投给《法国科学院院报》(CR)的笔记,那篇笔记的作者和题目(是我七八年前向他提议的)都让我放在心上。作者为展开这个题目花了两年工夫,这题目确实不时髦(尽管我始终觉得它同样切题)。我认为他做出了出色的工作(作为第三阶段博士论文提交)。我不是这位年轻研究者的“老板”——他碰巧天分极高(我不知道他是否还会在数学里施展这些天分,看看他所受到的接待……),他的工作是在与我毫无接触的情况下做出来的。但也确实,这个题目的来源不可能有任何疑问;这可怜人一开始就处境不妙,而且他肯定还蒙在鼓里!这位同行倒还做得有礼有节,至少还有这点,我也没指望他更多,“真诚地抱歉,可您也明白……”。两%
% ---------------- p.188 ----------------
\origpage{82}年工夫——一个干劲十足的新手研究者的两年工夫,换来的却是一篇三页的《院报》笔记——那篇笔记又能花掉多少公款呢?这里有一种一眼就能看出的荒谬:两者之间那种巨大的不相称。只要肯费心去查一查深层的动机,这种荒谬肯定就消失了。只有这位同行兼旧友才能探测他自己的动机,正如只有我才能探测我的。但不必走很远我就知道,需要照顾的并不是数学产出涨得没边儿,您懂的,也不是公款(或《院报》那位想象出来的“无名读者”的耐心)……

这同一篇预备投给《院报》的笔记,先前还有幸被送到“法国那六七个人……”中的另一位手上,那位把它退给了作者的“老板”,因为这些数学“不叫他觉得好玩”(原话如此!)。(老板虽然恶心,却很谨慎,自己地位本来就不太稳,两次都宁可咽下去,也不愿得罪人……)后来我有机会跟这位同行兼旧日学生谈起此事,才知道他当初还是费心仔细读了送来的笔记,也想过一想(那笔记想必勾起他不少回忆……),并且觉得其中有些陈述本可以用对使用者更友好的方式写出来。然而他并不屑于把宝贵的时间浪费在把意见转达给当事人上:大人物的一刻钟,对一个无名年轻研究者的两年工夫!数学确实叫他“觉得好玩”到足以让他抓住这个机会,重新接触笔记里研究的那种情形(那情形一定在他心里、也在我心里,唤起了各式各样丰富的几何联想),把给出的描述吸收进去,然后——以他的家底与本事,毫不费力地——挑出其中的笨拙与疏漏。他没有浪费时间:他对某个数学情形的知识,由于一个初出茅庐的研究者两年的认真工作,变得更精确、更丰富了;那工作,大师当然几天之内就能做完(只做大略,不做证明)。这些一到手,人就想起来自己是谁了——案子已经判了,无名先生两年的工夫,正好进垃圾桶……

有些人,这种风一吹过来,什么也感觉不到——可我今天仍然被它吹得透不过气。在那件事里,这肯定正是所追求的效果之一(瞧瞧那拒绝所披的漂亮外衣),但肯定不是唯一的效果。就在同一次谈话里,这位旧友带着一种谦逊的得意向我吐露:他只有“所陈述的结果叫他吃惊,或者他不知道该怎么证明”时,才肯把一篇笔记推荐给《院报》\fnoteP{1}{(27)}{\emph{年轻人的“势利眼”,或纯洁的捍卫者}\par Ronnie Brown 向我转述了 J.H.C.怀特黑德(J.H.C. Whitehead)的一句话(布朗是他的学生),说的是“年轻人的势利眼:他们以为,一个定理的证明若是 trivial,定理便是 trivial 的”。我昔日的许多朋友,不妨好好琢磨琢磨这话。这种“势利眼”在今天绝不限于年轻人,我认识不止一位有声望的数学家,就常常如此。我对此格外敏感,因为我在数学里(别处也一样……)做得最好的东西,那些我认为最富成果的概念与结构,以及我凭耐心而顽强的工作从中抽出的本质性质,统统落进了“trivial”这个形容词之下。(这些东西,若作者不是早已成了名人,如今哪一样也难被接受为一篇《院报》笔记!)我做数学家的毕生野心,或者不如说我的激情与快乐,一直是找到那些显而易见的东西;在本书里(包括在眼下这个导论性的章节里……)我唯一的野心也是如此。决定性的东西,往往首先在于看见那个尚未被看见的问题(无论答案如何,也无论答案是否已经找到),或者抽出一个陈述(哪怕是猜想)——它概括并包含一个尚未被看见、或未被理解的局面;至于它是否被证明,证明是否 trivial,完全无关紧要,是彻底的枝节,哪怕一个仓促的、暂定的证明后来被发现是错的。怀特黑德所说的势利眼,是那种饱食终日的浪荡子的势利眼:他非要先弄清一瓶酒贵得离谱,才肯承认它好。这些年里,不止一次,我被旧日的激情一阵阵攫住,捧出我最好的东西,却看着它被那种自命不凡打回来。我感到一种至今犹新的痛,一份快乐落了空——但我还不至于流落街头,而且幸而我并不想把自己写的一篇文章安插出去。\par 怀特黑德所说的势利眼,是一种滥用权力,一种不诚实,而不只是对事物之美的麻木或封闭——当它由一个有权势的人施加于一个任其摆布的研究者身上,那人完全有条件把后者的想法吸收过来、据为己有,同时却以“显而易见”或“trivial”、因而“没有意思”为借口,堵住它们的发表。我这里甚至没有想到通常意义上的抄袭那种极端情形,那在数学界想必还很罕见。然而从实际的角度看,对付出代价的研究者来说,情形是一样的,而使它成为可能的那种内在态度,在我看来也相去不远。只不过它更舒服些:它伴随着一种对他人无限优越的感觉,伴随着那种以数学不可触碰的纯洁之毫不妥协的捍卫者自居者的问心无愧与内心满足。} (27)。这大概正是他发表得很少的一个原因。假如他把自己的标准用在自己身上,%
% ---------------- p.189 ----------------
\origpage{83}他索性就什么也不发表了。(确实,以他眼下的处境,他根本不需要发表。)他什么都晓得,要叫他吃惊,大概就跟找到一条虽可证明、他却不知如何证明的命题一样难。(二十年里,这两样我大约只碰上过两三回,而且最近十年、十五年里一次也没有!)他显然为自己的“质量”标准自豪,这标准使他在数学家的行当里成了把要求推到极致的冠军。我在其中看到的是一种百毒不侵的自我陶醉,而且不止一次看到,在那种笑眯眯、随和厚道的外表后面,是对他人毫无节制的轻蔑。我也看得出,他从中得到了极大的满足。

这位同行的情况,是我在“新伦理”的代表人物中遇到的最极端的一例。但它仍然很典型。在这里,无论是我讲的那件事,还是使之合理化的那篇信念宣言,用常识来衡量都透着一种乌布式的荒诞——其尺度如此之大,以至于这位脑子如此出众的旧友,以及肯定还有许多地位不那么显赫的同行(他们只消不找他推荐笔记就行了),都已经看不见它了。因为要看,至少得去看。只要肯费心去看动机(首先是他自己的动机),荒诞就暴露在光天化日之下,而它们同时也不再是荒诞的——它们交出了自己那份朴素的、一目了然的含义。

如果说这些年里,面对某些态度、尤其是某些行为,我常常痛苦到这般地步,那肯定是因为我隐约在其中辨认出——像是被推到极致、直到可笑或可憎地步的漫画——一些曾经属于我自己的态度与行为,它们又通过我的某位旧学生或旧友回到我身上。不止一次,我心里那套老反射又启动了:去揭发,去同那被清清楚楚指认出来的“恶”作斗争——但即便我偶尔在这里那里顺从了它,那也是带着分裂的确信。说到底,我很清楚,去斗争,仍旧是在事物的表面打滑,仍旧是回避。我的角色不是揭发,甚至也不是“改善”我所身处其中的这个世界,或“改善”我自己这个人。我的天职是学习,是通过我自己去认识这个世界,并通过这个世界认识我自己。如果我的生命能给我自己或他人带来什么好处,那也只是在我懂得忠于这份天职、懂得与自己和谐的限度之内。是时候提醒自己这一点了,好切断我心里那些老机制——它们在这里想推着我去为某个事业辩护(就说是某种已死的伦理吧),或者去说服人(去说服人相信取代它的那种伦理大概“荒诞”吧),而不是去探测、去发现、去认识,或者把描述当作探测的手段。写前面那两三页时,我除了就今日取代昨日之种种的流行态度说几句话之外,并无更确切的目的,却始终感到自己在提防着自己,那种心境,就像一个人随时准备大笔一挥,把刚写完的东西统统划掉、扔进纸篓!不过我还是要把写下的保留下来:它并不假,可它终究造出了一种虚假的局面,因为我在其中牵连别人多于牵连自己。我心底感到,我写着的时候什么也没学到——肯定正是这一点造成了我的不适。总之,是时候回到一种更实在的反思了,一种教我自己的反思,而不是妄图去教别人或说服别人的反思\fnoteP{2}{(28)}{写前面那几页时,起先我夹在“把肚子里的东西倒空”的欲望与一种保留、一种谨慎之间。所以我一直停在差不多上,这肯定是我那份不适、那种“我什么也没学到”之感的主要原因。自从记下这份不适的那几行写出之后,我把这几页改写了两遍——它们此前让我内心不满——写得更清楚地把自己摆进去,也更往事情深处走。一路下来,我当真“学到了点什么”,而且我相信,同时我还成功地用手指触到了一件重要的东西,它既超出这一具体事例,也超出我自己这个人。} (28)。

% ---------------- p.190 ----------------
\section{(34) 淤泥与源泉}

\origpage{85}我觉得,大体上,我已经把自己同其他数学家的关系——各种年纪、各种等级的数学家,那是我还属于他们那个世界、属于数学家的世界的时候——都过了一遍;同时,而且尤其是,也把我以自身态度与行为在其中所占的那一份过了一遍:那是我在今天所看到的某种精神里所占的一份,而这精神肯定不是昨天才有的。在这番反思里,或者更确切地说,在这趟旅行里,我四次遇到一些情形,它们在我看来是我身上某些态度与暧昧的典型:在那里,对他人那种自发的善意与尊重,被一些自我性的力量搅乱了,甚至一扫而光,尤其是(至少在这四例中的三例里)被一种自得搅乱了。这自得主要依仗的是这样一种所谓优越:某种脑力,加上我在数学活动上那种不成比例的投入,据说赋予了我这种优越。它从一种普遍的共识中得到确认与支撑——那共识几乎毫无保留地推崇这种脑力和这种不成比例的投入。

在这几例情形中,最后一例——那个“踩进我的地盘的没规矩的年轻人”——我觉得对眼下的论述最为重要。前三例是我这个人的典型,或我这个人的某些侧面的典型,是在某个时期(也确实是在某种环境里)——但是,如我已经有机会一再说过的,我根本不认为它们是我当时所属那个圈子的典型。我也不认为它们是今天法国数学界的典型,这么说吧——比如说,我与“不知疲倦的朋友”之间那种关系所特有的那种慢性迷误,如今大概很少见,就像当年想必也很少见一样。相反,我在“没规矩的年轻人”这件事上的态度与行为,却是今天数学界里每天都发生的事的典型,无论往哪里看。有影响力的数学家对无名青年怀抱善意、给予尊重,如今倒成了极罕见的例外——如果这青年没福气做他的学生(即便如此也未必……),也不是一位地位相当的同事向他推荐的学生。这大概正是我1970年“觉醒”之后不久就已传到我这里的东西——那场觉醒解开了许多沉默者的舌头——但那时我听到的第一手证词对我仍然是遥远的,因为它们直接涉及的,既不是我本人,也不是我圈子里最亲近的朋友。从某个时候起(大约1976年),我开始受到不止是表面的触动:那时传到我这里的回声,或我亲眼目睹的事实,其主角是这些朋友中的某几位,甚至是一些变得显赫的旧学生;而当那些遭到恶意中伤的人是我很熟悉的人、不止一次是我的学生(1970年以后的学生,那是不用说的!),因而他们的命运牵动我的时候,就更是如此。在某些情况下,而且毫无疑问:那种缺乏善意,甚至那种炫耀式的轻蔑态度,至少因为这样一个事实而变本加厉——某位年轻研究者是我的学生,或者他冒着风险(不一定非是我的学生)去做我昔日的朋友和其他同行也乐于称之为“格罗滕迪克式的玩意儿”的事情……

70年代初,“那个没规矩的年轻人”又给我写过信,非常客气地问我(他本来完全没有义务问我任何事!)是否见怪:他想发表一个定理的证明,那定理人家告诉他是我的,却从未发表过。我记得我回信时还是像过去那样一肚子不快,大概既没说行也没说不行,还让人听出这样的意思——我没看他的证明(他本来%
% ---------------- p.191 ----------------
\origpage{86}\origpage{87}当然愿意给我看,可我全不在乎,我当时正忙着我的那些活动!):那证明对我的证明肯定毫无补益(可它至少有一点好处:白纸黑字写下来,连同陈述本身,可供数学界使用!)。这恰好说明,那场著名的“觉醒”还多么浮浅,对某些根植于自得、根植于“精英至上”态度的行为毫无影响——而同一时期,我肯定正在《生存与生活》(Survivre et Vivre)上那些情真意切的文章里、在公开辩论的发言里,揭发着这些行为……

这就很具体地回答了我先前悬置的一个问题。不如就在这里承认这条朴素的真理:这类自得的态度在我身上绝没有“一劳永逸地”被克服,我怀疑有朝一日会,除非到我死。如果说确有转变,那也不是因为虚荣消失了,而是因为(重新)出现了一种好奇——对我自己这个人、对我身上某些态度与行为的真实本性的好奇。正是凭这种好奇,我才变得多少能察觉自己身上的种种虚荣表现。这深刻地改变了某种内在的动力,也因此改变了“虚荣”的效果;所谓虚荣,就是那股常常推着我去掩盖或伪造我对现实那种健全而细腻的感知的力量,其目的是壮大我自己、把自己置于他人之上,同时却宣称恰恰相反。

也许某位读者会感到困惑,就像我有朝一日也曾困惑:在我作为数学家的生活里,虚荣那种阴险而顽强的存在(他或许有时也在自己的生活里瞥见过它),与我所谓的对数学的爱、或者说激情(它也许同样在他自己对数学、或对某个别人、别的什么事的体验里唤起回声),这两者之间有一种表面的矛盾。如果他确实困惑,那么他自身就具备所需要的一切,可以(像我从前那样)重新接触到事物本身的现实——那是他可以第一手认识的东西——而不是像一只被关在词语与概念的无尽笼子里的松鼠那样打转。

看见一汪浑水的人,难道会说水与泥是一回事?要认识那并非泥的水,只消上到源头,看一看,喝一口。要认识那并非水的泥,只消登上被太阳和风晒干的河岸,在手心里掰开、捻碎一团带沙粒的黏土。野心、虚荣,或多或少能决定一个人在生活中给某种激情——比如数学的激情——留多少地盘,能在回报丰厚时把它变得贪得无厌。但最贪得无厌的野心,本身却无力发现或认识哪怕最微小的一件事,恰恰相反!在工作的时候,当理解一点点开动、成形、加深;当混乱中一点点显出秩序,或者当原本熟悉的东西忽然露出陌生的面貌,继而令人不安,直到终于爆出一个矛盾,把一种看似不可动摇的对事物的看法掀翻——在这样的工作里,没有一丝野心或虚荣的痕迹。那时领着跳舞的,是某种来自比“自我”及其不断自我膨胀的饥饿(哪怕是“知晓”与“知识”的饥饿)远为深远之处的东西——肯定比我们这个人、甚至我们这个物种都远为深远。

那就是源泉,它在我们每个人心里。

\section{(35) 我的激情}

三大激情支配了我成年后的生活,此外还有一些性质不同的力量。我终于在这些激情中认出同一种深层冲动的三种表现;那是我心中认知冲动所走的条条道路中的三条,而在我们这个无限的世界上,呈现在它面前的路有无数条。

% ---------------- p.192 ----------------
\origpage{88}在我生命里最先出现的,是我对数学的激情。十七岁那年,刚出中学,我松开缰绳,任由一种简单的癖好去跑,它就展开成了激情,在随后的二十五年里支配了我生命的进程。我“认识”数学,远在我认识第一个女人之前(除了那个我生下来就认识的女人);如今到了成熟的年纪,我发觉它始终没有燃尽。它不再支配我的生命,我也不再声称要支配它。有时它睡去,有时我甚至以为它熄灭了,它却又悄没声息地回来,一如既往地火热。它不再像从前那样吞噬我的生命——从前我把生命交给它去吞噬。它继续在我生命里留下深深的印痕,就像他爱的女人在一个情人身上留下的印痕。

我生命里的第二种激情,是对女人的追寻。这激情常以追寻伴侣的面貌出现在我面前。直到这追寻快要结束的时候,我才懂得把两者分开:那时我知道,我所追的东西哪里也找不到,或者也知道:我就把它带在自己身上。我对女人的激情,要到我母亲去世之后才真正展开(那是我第一次恋爱五年之后,那次恋爱生了一个儿子)。就在那时,二十九岁,我成了家,又有了三个孩子。起初,对孩子的依恋是对母亲依恋不可分割的一部分,是那股从女人身上发出、把我吸引到她那里去的力量的一部分。那是这份爱的激情的果实之一。

我并没有把这两种激情在我身上的并存体验为冲突,无论起初还是后来。我必定隐隐感到了两者的深层同一,而它明确显现给我,要晚得多,是在第三种激情进入我生命之后。然而,这两种激情对我生命的影响,只能是大不相同。对数学的爱把我吸引到某个世界里,即数学对象的世界;那个世界当然有它自己的“现实”,可它不是人们的生活在其中展开的那个世界。对数学事物的切身认识,可以说没有教给我任何关于我自己的东西,更不用说关于别人的东西——朝向数学的发现之冲动,只能使我远离自己和他人。有时两三个人会在这同一种冲动中相通,但那是一种表面的相通,实际上使每个人既远离自己也远离别人。所以,对数学的激情在我生命里从来不是一种成熟的力量;我也怀疑这样的激情能促进任何人的成熟\fnoteP{3}{(29)}{我这里说的是对数学、或对另一种纯智力活动的紧张而长期的投入。反过来,这样一种激情的展开,可以是与我们身上一种被遗忘的力量重新相识的方式,是同一份不肯就范的质料较量一番的机会,也一路让我们借某种真正属于自己的东西,更新并充实我们的身份感——这样的展开,完全可能是内在旅程中、成熟过程中重要的一步。} (29)。如果说我长期以来给这激情在我的生命里留了如此不成比例的位置,那肯定也正因为:它让我得以逃避对冲突的认识,也逃避对自己的认识。

相反,性的冲动,不管我们愿不愿意,都把我们径直抛向与他人的相遇,径直抛进我们自身、也抛进对方身上那个冲突的结!而我生命里对“伴侣”的追寻,却是对无冲突的幸福的追寻——那不是认知的冲动、性的冲动,并不像我喜欢相信的那样,而是对认识他人身上和我自身之中的冲突的无尽逃避。(为了结束这虚幻的追寻,以及如影随形般伴着它的不安,这是我必须学会的两件事之一。……)幸好,人尽管逃,性总会很快把人带回来!

有一天,我不再拒不接受冲突通过我爱过的、或正爱着的女人,以及通过这些爱情所生的孩子,顽强地带给我的教训。当我终于开始%
% ---------------- p.193 ----------------
\origpage{89}\origpage{90}倾听、开始学习的时候,而且在那之后许多年里,事情一直是这样:凡我学到的东西,都是通过我爱过的或正爱着的女人学到的\fnoteP{4}{(30)}{近几年,接力的是我的孩子们,由他们来给一个有时不大情愿的学生讲授人类存在的奥秘……} (30)。直到1976年,四十八岁,对女人的追寻一直是我生命里唯一巨大的成熟之力。这成熟之所以只是在此后几年、也就是最近七年才完成,是因为我用尽手边的一切手段来保护自己、不去成熟(这手段我是从父母那里、从我认识的那些圈子里学会的)。这些手段中最有效的,就是我在数学激情上的投入。

第三种大激情进入我生命的那一天——1976年10月某个夜晚——学习的大恐惧消失了。那也是对再寻常不过的现实的恐惧,首先是对关乎我这个人的那些朴素真理的恐惧,或对关乎我所珍爱的人的真理的恐惧。说来奇怪,四十八岁之前,我从未在自己身上察觉这恐惧。我是在新激情、认知激情的新表现出现的那一夜发现它的。可以说,后者取代了那终于被认出的恐惧的位置。多年来,我在别人身上把这恐惧看得清清楚楚,却有一种奇怪的盲目,使我对自己身上的它视而不见。对看的恐惧,使我看不见这同一种对看的恐惧!我和所有人一样,牢牢依恋着某种关于自己的形象,那形象大体上从童年起就没有动过。我说的那一夜,也是这旧形象第一次塌下去的一夜。另一些与它相像的形象接替了它,凭着顽固的惯性维持几天、几个月,甚至一两年,然后在审视的目光下也塌下去。懒得去看,常常推迟这样一次新的觉醒——但看的恐惧再也没有回来。凡有好奇之处,恐惧便再无容身之地。当我对自己怀着好奇的时候,我并不更害怕自己会发现什么,就像我想知道一个数学情形的底细时一样:那时有的是一种快活的期待,有时急切,却又执拗,准备好迎接一切愿意来到面前的东西,无论预料的还是意外的——一种满怀激情的专注,守候着那些毫不含糊的迹象,凭它们在假、半真与也许的初始混乱中认出真来。

在对自己好奇之中,有爱;这爱不为任何恐惧所扰——不担心我们所看的东西不合我们所愿。说实话,对自我的爱早已在这一夜之前的几个月里静静地绽开;这一夜也是这爱取得行动形式的一夜,可以说,是变得敢作敢为的一夜,毫不客气地把行头与布景撞翻!如我所说,不久又有别的行头和布景像变戏法一样冒出来,然后轮到它们被撞翻,没有咒骂,也没有咬牙切齿……

这七年里,新激情在我生命中的种种表现,终于在我看来像是一浪接一浪起伏的波涛,像一次辽阔而安详的呼吸的一进一出。这里不是试着去描出它那条蜿蜒变化之线的地方,也不是去描出作为对位的数学激情表现之线的地方。我已放弃去安排这一种或那一种的进程——如今倒不如说是两者的这双重运动在安排我生命的进程——或者更确切地说,它就是那进程本身。

在新激情出现之前的那些月份里——孕育与丰满的月份——对女人的追寻已开始改换面貌。它开始摆脱曾浸润它的那份不安,就像一口气息终于从压迫它的重负下解脱,重新找回它自己的幅度与节律。又像一团火,因为无处出气而半明半灭地闷烧着,一阵凉风吹来,忽然舒展成噼啪作响、轻盈%
% ---------------- p.194 ----------------
\origpage{91}而活泼的火焰!

火已烧足。一种看似无法止息的饥饿得到了饱足。两三年来,那场追寻似乎确已烧尽,不留一丝灰烬,把场地让给两种激情的歌唱与对唱。其一,我少年时代的激情,三十年间一直用来让我与一个被否弃的童年分离。其二,是我成熟年纪的激情,它让我重新找回了孩子,也找回了我的童年。

\section{(36) 欲望与冥想}

我说过的那一夜,一种新激情取代了一种永远消逝的旧恐惧,也是我发现冥想的那一夜。那是我第一次“冥想”之夜,它是在一种专横而急迫的需要压迫下出现的——此前几天,我仿佛被一阵阵焦虑的浪头淹没了。也许像一切焦虑一样,那是一种“起飞”式的焦虑:它执拗地向我指出:关乎我这个人的那份朴素而显眼的现实,已经同那个维持了四十年、我却从未怀疑过的自我形象脱节。肯定有一份巨大的认知渴望,与那些可观的逃跑力量并排着,也与那种逃避焦虑、想跟从前一样安宁的欲望并排着。于是有了一场紧张的工作,持续了几个小时,直到它有了结局,而我当时还不知道正在发生的事意味着什么,更不知道我要往哪里去。在这场工作中,一个个遁词被认了出来;或者更确切地说,正是这场工作让这些遁词一个个现形,每一个都表现为一种内心确信——我终于肯费心把它白纸黑字记下来,仿佛这样才更能浸透它,而它此前一直待在一种有利的模糊里。我满心欢喜地记下它,毫无提防:它必定有本事叫我着迷——那时的心境,就是一个什么都不怀疑的人的心境,对他来说,把一条未说出口的确信白纸黑字写下来,这本身就不可辩驳地证明它真实,证明它有根有据。假如我心里没有那份不知分寸、甚至该说不知羞耻的欲望——我是说认知的欲望——我每次都会停在这个“happy end”上,而每一步确实就是在这happy end 的心境里结束的。然后,我倒霉了!我会心血来潮——天知道怎么回事、为什么——要把我刚满心满意写下的东西凑近再看一看:那不是白纸黑字写着嘛,只要重读一遍就行了!可一认真重读,傻乎乎地重读,我就觉得有点不对劲,觉得不那么清楚,咦咦!然后再费心凑近一点看,就清楚了:根本不是那么回事,可以说全是瞎编,我刚才是拿尿脬当灯笼了!每一次这局部的发现都像个大惊喜,“好家伙!这一条还真不含糊!”——一个快活的惊喜,它带着一股新的劲头,把反思重新发动起来。往前走吧,总会弄到底的,肯定马上就到,只要顺着这股劲儿接着走!小结一下,理一理……于是另一种内心确信已经升上来了,带着“事情底细”的全部外表,我们只愿意相信:这一回总该对了吧;好歹记下来,图个心安,再说记下这么中肯、这么有体会的东西本身就是乐事,不赞成的人准是心术不正,这么明显的诚意,没法做得更好了,就这样,完美!

这就是新的阶段终点,新的happy end,我本来会高高兴兴停在那里的,要不是那个调皮到极点的坏小子又闹了起来——他真是不可救药,竟又把鼻子伸进这最新的“底细”和happy end里。拦不住他,于是又开始了新的一步!

% ---------------- p.195 ----------------
\origpage{92}\origpage{93}就这样,四个小时里,阶段一个接一个过去,像一颗洋葱,我把它的层一层层剥掉(这是那一夜末尾来到我脑中的形象),终于剥到了核心——到达那条再简单不过、再显眼不过的真理;说实话,那真理本就明摆在那里,可我却好几天、好几个星期(说到底,是我一辈子)都成功地把它藏匿在这一堆互相遮掩的“洋葱层”底下。

那条朴素真理终于出现,是巨大的宽慰,是意料之外、彻彻底底的解脱。那一刻我知道,我触到了焦虑的结。这五天的焦虑确确实实被化解了、消融了,转变成刚刚在我心里成形的那种认识。焦虑不只是从我眼前消失——就像整个冥想过程中那样,也像此前五天里好几次那样;而它转变成的认识,根本不属于观念的性质,也不属于那种为了了结、为了安宁而作出的让步(就像同一夜里有几回我做过的那样);那不是一件我从外面拿过来、或取得之后添到自己身上的东西。那是十足意义上的认识,第一手的,朴素而显眼,从此成了我的一部分,就像我的血肉是我的一部分。而且,它是用清楚、毫不含糊的话表达出来的——不是一大段话,而是一句三四个字的、再傻不过的短句。这表达是刚刚进行的那场工作的最后一步;只要这最后一步没跨过去,它就一直是短暂的、可逆的。在这整场工作中,把正在成形的思想、正在出现的念头仔细地、甚至一丝不苟地表达出来,是这工作本质的一部分;而每一次重新出发,都是对刚刚走过那一步的反思——那一步是我凭刚刚写下的见证而知道的(没法把它藏进一个靠不住的记忆的迷雾里!)。

在发现与解脱那一刻之后的几分钟里,我也知道了刚刚发生的事的全部分量。我刚发现的东西,比这几天那条朴素真理还要珍贵。这东西,就是我身上的一种能力——只要我有兴趣——去认识我内心发生之事的底细,认识任何分裂、冲突处境的底细;也因此,就是完全凭自己的手段去解决我心中任何我懂得去觉察的冲突的能力。

这解决不是靠什么恩典的功效完成的——前些年我倾向于那样相信——而是靠一场紧张、执拗、细致的工作,用的是我平常的官能。如果说真有“恩典”,它不在于我们心中某个冲突突然而彻底地消失,也不在于对冲突的一种理解现成地送到我们面前(就像懒人国里烤好的鸡!),而在于这认知欲望的在场或出现\fnoteP{5}{(31)}{我这里想到的是认知欲望的“阳”的形式——那种探测、发现、为出现之物命名的欲望……正是被命名,才使出现的认识变得不可逆、不可磨灭(哪怕它后来被埋没、被遗忘,哪怕它不再起作用……)。认知欲望的“阴”的、“女性”的形式,则在于一种敞开、一种接纳,在于默默迎候一种在我们存在的更深层、思想所不能抵达之处出现的认识。这样一种敞开、以及一种一时抹去一切冲突痕迹的突然认识的出现,仍像一种恩典,它深深地打动我们,尽管其可见的效果也许短暂。不过我猜想,这种如此来到我们身上的、无字的认识,在生命某些罕见的时刻,同样不可磨灭,它的作用甚至延续到我们对它的记忆之外。} (31)。正是这欲望引导我、在几个小时里把我带到冲突的核心——正如爱的欲望让我们万无一失地找到通往所爱女人最深处的那条路。

无论是自我发现还是数学,没有欲望,一切所谓的“工作”都只是装样子,哪里也去不了。最好的情况下,它让乐在其中的人无休无止地“围着罐子转”——罐里的东西是留给饿了要吃的人的!如同世上的%
% ---------------- p.196 ----------------
\origpage{94}所有人一样,我也有欲望与饥饿都不在场的时候。当涉及对自我的认知欲望时,我对自身以及我所卷入的处境的认识便停滞不动,我行事不是凭着对事情的认识,而是听任一些根深蒂固的简单机制摆布,随之而来的种种后果也一并承担——有点像一辆由电脑、而不是由人驾驶的车。但无论是冥想还是数学,没有欲望的时候、没有这份饥饿的时候,我不会想去装模作样地“工作”。所以我不曾有过冥想哪怕几个小时,或者做数学哪怕几个小时\fnoteP{6}{(32)}{\emph{炉中百铁,或:干熬无济于事!}\par 在我还做泛函分析的时候,也就是直到1954年,我有时会没完没了地死抠一个解不出的问题,明明已经没了主意,只是在旧念头的圈子里打转,而那些念头显然已经“咬”不动了。至少有一整年就是这样,特别是关于拓扑向量空间里的“逼近问题”,那问题要到二十年后才被完全不同路数的方法解决——以我当时的地步,那些方法只能与我无缘。那时驱动我的不是欲望,而是一种固执,以及对自己内心所发生之事的无知。那是难受的一年——我一生中唯一一段做数学变得难受的时期!我需要这番经历才明白:“干熬”没有用——一旦工作到了停住的地方,一察觉停住,就得转去做别的——不妨等到更合宜的时候再回到悬置的问题上。那个时刻几乎总是不久就出现——问题自己成熟起来,我根本不必去碰它,只凭在那些看似与此毫无关系的问题上兴致勃勃地工作,就足够了。我确信,当时若一味死抠,哪怕十年我也什么也弄不出来!正是从1954年起,我在数学里养成了同时往炉子里放很多块铁的习惯。我一次只打一块,但靠着一种不断重现的奇迹,我在一块上的工作也惠及所有别的铁,它们各自等着自己的时辰。后来与冥想初次接触时,情形也一样,而且完全不是我事先有意安排的——要审视的灼热问题一天天多起来,随着反思的进行而不断增加。} (32),而没学到一点东西的;而且多半(若非说总是)是某种意外、无法预见的东西。这跟我有、而别人没有的什么官能毫无关系,只在于:我真心不想做的时候,不去装模作样地做。(正是这份“想”的力量,单独就造出了我别处谈到的那种要求:在工作中不满足于差不多,非得把一种理解——无论多么朴素——走到尽头才满意。)在要发现什么的地方,没有欲望的工作就是无意义、就是装样子,跟没有欲望地做爱一样。说实话,我从未受到过这样的诱惑:把精力浪费在假装做一件我根本不想做的事上——何况有那么多让人热血沸腾的事可做,哪怕是在该睡觉的时候睡觉(还有做梦……)。

我相信,也是在同一夜,我明白了:认知的欲望与认知、发现的能力是一回事。只要我们信任它、跟随它,就是欲望把我们一直带到我们想认识的事物的核心。也还是它让我们找到——甚至不必去找——认识这些事物最有效、也最适合我们这个人的方法。至于数学,看来自古以来写作就是一种不可或缺的手段,无论“做数学”的是谁:做数学,首先就是写\fnoteP{7}{(33)}{这并不意味着,工作中没有纸(或黑板——那不过是它的一个变体!)的那些时刻,在数学工作里就不重要。尤其在“敏感时刻”更是如此:那时一种新的直觉刚刚出现,要做的不是“零敲碎打的工作”,而是以一种更整体、更直觉的方式与它“相识”——这种相识由反思的那个非正式阶段预备。在我这里,这类思考多半在床上或散步中进行,我觉得它在整个工作时间里所占的份额相对不大。同样的观察也适用于我迄今所做的冥想工作。} (33)。在一切以智力为主的发现工作里,大概也是如此。但“冥想”——我指的是自我发现的工作——肯定不必如此。然而在我这里,直到目前,写作在冥想中一直是一种有效而不可或缺的手段。就像在数学工作中一样,它是物质载体,把
% (切片止于此句: 源文 p.196 末行 "elle est le support matériel qui fixe le", 下接 p.197)

% ================================================================
%  第 9 章 收获 (后半) 原书 p.197-208, 边码 95-110
%  源文: _work/src1/task-ch09b.md
%  接 part1-ch09a.tex (p.187-196)
% ================================================================
%
% ---------------- p.197 ----------------
\origpage{95} 反思的节奏固定下来,并充当注意力的坐标与集合点——否则我的注意力往往倾向于四散随风;而且,书写还给我们留下一道刚刚做完的工作的有形痕迹,我们可以随时求助于它。在一场旷日持久的冥想里,能够回头翻看那些书面痕迹常常是有用的:它们见证着前几天、甚至几年前某个时刻的冥想。

\origpage{96} 思想,以及对思想的细致表述,因此在我迄今所实践的冥想中扮演着重要角色。但冥想并不因此就限于一种单纯的思想工作。单凭思想,无力把握生命。思想特别擅长的,是侦测矛盾——在我们对自身的看法里、在我们与他人关系的看法里,那些往往大得滑稽的矛盾;但常常,它不足以把握这些矛盾的意义。对于被求知之欲所驱动的人,思想是一件往往有用而有效的工具,甚至不可或缺——只要人始终意识到它的限度,这限度在冥想中显而易见(在数学工作中则更为隐蔽)。要紧的是:在另一样东西出现的那些敏感时刻,思想要懂得退开,踮着脚尖消失——那另一样东西也许以一阵突然而深沉的情绪的形式出现,而与此同时,手也许还在纸上继续奔跑,在同一刻给它一个笨拙而结结巴巴的表达……

\section{(37) 惊叹}

这段关于发现冥想的回顾,完全是意料之外地来到这里的,几乎违背我自己的意志——我起初打算考察的完全不是这个。我本来想谈惊叹。那个夜晚如此丰富,满载着如此多的东西,也满载着面对这些东西的惊叹。其实在工作进行当中,就已经有一种惊叹:面对每一条被揭出的新遁词,一种难以置信的惊叹——那就像一件用粗白线草草缝成的蹩脚戏装,而我竟,简直难以相信!还一本正经地把它当成真得不能再真的东西!此后许多次,在随后的年月里,我又重新遇到这同一种惊叹,如同在冥想的第一夜,面对我所发现的事实的巨大,以及那些让我一直对之视而不见的托辞的粗劣。正是从它那些滑稽可笑的侧面,我开始发现我身上那个未曾被猜想到的世界——一个随着日子、月份与年岁的推移而显露出惊人丰富的世界。然而就在那第一夜,值得我惊叹的就已经不只是些闹剧式的插曲。正是在那一夜,我头一次重新接触到一个被遗忘的、沉睡在我身上的能力;它的性质我还不清楚,只知道它是一种能力,而且随时听我支配。

而此前那几个月,也已经充满了一种无声的惊叹——惊叹于我身上一直带着的某样东西,那东西想必从来就在那里,而我刚刚才重新与它接触。我感受到这东西,不是作为一种能力,而更像是一种隐秘的温柔,一种既极其安宁又令人不安的美。后来,在我发现自己那份被忽视了如此之久的能力的狂喜中,我忘掉了这几个默默孕育的月份——它们只留下几首散落的诗作证:一些情诗,它们夹在我那些冥想笔记中间,大概常常显得格格不入……

只是许多年之后,我才回想起这些惊叹于世界之美、惊叹于我感觉到安歇在自己身上的那份美的时光。那时我才知道:我在自己身上感觉到的这份温柔与美,以及我不久之后发现、并深刻改变了我一生的那份能力,乃是同一件事的两个不可分割的方面。

% ---------------- p.198 ----------------
而我现在也看出:在我们身上那多重的东西——创造性——之中,那温柔的、内敛的、沉默的一面,自发地通过惊叹表达出来。而人认识所爱的女人、她也认识他,同样是在对由所爱者揭示出来的、自身之中那份不可言说的美的惊叹之中。当对所探究之物、对所爱之人的惊叹缺席时,我们与世界的拥抱就被削去了其中最好的部分——被削去了那使它成为对自己、对世界的祝福的东西。不是惊叹的拥抱,是一种无力的拥抱,只是占有姿态的简单复制。它无力生出别的东西,只会生出更多的复制品——也许更大些、更粗些、更厚些,无所谓,永远不是一次更新\fnoteP{8}{(34)}{\emph{无力的拥抱}\par “拥抱”一词在我这里绝不是一个简单的比喻,日常语言在此成了一个深层同一性的反映。人们可以说——而且不无道理——那么“没有惊叹的拥抱是无力的”这话就不对:若按字面理解,大地早该人烟断绝、一片荒芜了。极端的例子是强奸,其中确实没有惊叹,然而有时一个存在就在被强奸的女人身上被孕育出来。由这样的拥抱生出的孩子,想必不能不带这印记,它属于他所分得的那份“包裹”,属于该由他去承担的东西;这并不妨碍一个新存在确实被孕育、并且出生了,确实有过创造,那是力量的标志。也确实,我曾见过某个充满自得的数学家发现并证明出漂亮的定理,那是一次并不缺乏力量的拥抱的标志!但同样真实的是:如果这位数学家的生命被他的自得闷死(在某个时期,我自己的生活在一定程度上就是如此),那么这些与数学的拥抱之果实,对他、对任何人都不是福泽。对一个强奸所生的孩子的父亲和母亲,也可以这么说。我若说“无力的拥抱”,首先指的是无力在那个自以为在创造的人身上生出更新:他其实只造出一个产品,一件外在于他的东西,在他自身之中没有深层的回响;这产品远不能解放他、不能在他身上造出和谐,反而把他更紧地捆在他所囚于其中的那份自得上,那自得不断推着他去生产、再生产。这是深层意义上的一种无力,藏在一种“创造性”的外表背后——那“创造性”说到底只是一种没有刹车的生产力。\par 我也有充分的机会认识到:自得、无力惊叹,其性质是一种真正的盲目,是天然的敏感与嗅觉的阻塞;这阻塞即使不是完全的、永久的,至少在某些特定情境中显而易见。这是一种状态,其中某位声名显赫的数学家有时会在恰恰是他最擅长的东西上,显得比最迟钝的小学生还要愚笨!在另一些场合,他又会做出技术精湛的奇迹。然而我怀疑他是否还有能力去发现那些能够更新一门学科或一门科学的简单而显然的东西。它们在他下面太远,他已不屑于再看它们一眼!要看见谁也不屑于看见的东西,需要一种他已经失去、或被他放逐了的无辜……这肯定不是偶然:近二十年来数学生产的惊人增长,新结果纷至沓来,令那位只想稍稍“跟上形势”的数学家不知所措;然而(就我从各处得到的回声所能判断的),我曾多少熟悉的那些重大思考主题,几乎没有一个是真正更新了的,没有一个是大幅度的转化(而不只是积累)。更新不是量的事,它与投入的量无关——那量可以用某些“水平”的某些数学家花在某个题目上的“数学家人日”数目来衡量。一百万个“数学家人日”也无力生出一件像零那样孩子气的东西,而零更新了我们对数的感知。只有无辜有这份力量,它可见的标志之一,就是惊叹……} (34)。正是在我们还是孩子、准备为世上事物之美和为我们自身之美而惊叹的时候,我们才也准备好了去更新,并准备好作为柔软而顺从的工具,握在工匠手中,好让经祂的手、并借我们,也许有些存在与事物得以更新。

我清楚地记得,在四十年代末及随后的年月里,在这个对我而言代表着数学界的、不拘礼数的小圈子里——它有时喧闹而自信,那种略显断然的口气并不罕见(不过其中从不掺进一丝自得)——在这个圈子里,任何时候都为惊叹留有一席之地。惊叹在谁身上最显眼?是迪厄多内。不论是他作报告,还是他只在下面听,每当那个关键时刻到来、一道豁口突然敞开,就看见迪厄多内喜出望外、满面放光。那是纯粹的惊叹,有传染性、不可抗拒——其中“我”的一切痕迹都消失了。此刻我提起他,才意识到:这惊叹本身就是一种力量,它在他整个人周围立即发生作用,像一种以他为源头的辐射。要说我见过哪位数学家动用一种强大而基本的“鼓励之力”,那就是他!在此之前我从未再想到这一点,但现在我记起来了:也正是带着这样的心情,他当初接纳了
% ---------------- p.199 ----------------
\origpage{97} 我在南锡的最初几项结果,解决了他和施瓦茨一起提出的问题(关于 $(F)$ 空间与 $(LF)$ 空间)。那些结果都很平常,既不天才也不出众,可以说并没有什么值得惊叹的。后来我见过气派完全不同的一些东西,被那些自以为是大人物的同行一句不容分说的轻蔑打掉。迪厄多内身上全然没有这类负担——无论那自命是否站得住。没有任何这类东西妨碍他连小事也为之欣喜。

\origpage{98} 这种欣喜的能力里有种慷慨,它既造福于愿意让它在自己身上舒展的人,也造福于他周围的人。这恩惠的施行,并不怀着讨谁欢心的意图。它单纯得像花香,像阳光的温暖。

在我认识的所有数学家中,这份“馈赠”在我眼里以最耀眼、最有感染力、也许也最起作用的方式出现的,就是迪厄多内;我说不准\fnoteP{9}{(35)}{这份“馈赠”不是任何人的特权,我们生来都有。当它在我身上似乎缺席时,那是我自己把它赶走了,而重新接纳它只取决于我。在我身上或在某某人身上,这份“馈赠”的表达方式与另一个人不同,也许不那么有感染力、不那么不可抗拒,但它并不因此就不在,我也说不出它是否就不那么起作用。} (35)。但在我乐于交往的数学界朋友中间,没有哪一个身上缺这份馈赠。它总能找到机会显露,也许方式更含蓄,但随时都在。每当我走向他们中的一位,想与他分享我刚发现的、让我着迷的东西,它就显露出来。

如果说我作为数学家的生活里也有过挫折与痛苦,那首先就是:在我爱过的某些人身上,再也找不回我曾在他们身上认识的那份慷慨、那份对事物之美——无论“小”或“大”——的敏感;仿佛那曾使他们整个存在颤动不已的东西,已经无声无息地熄灭,被一种自得闷死了——那种自得属于这样一个人:世界对他而言已不够美,不值得他屈尊为之一喜。

当然,也还有另一种痛苦:看见我昔日某位朋友以俯就或轻蔑的态度对待我今天的某位朋友。但说到底,这痛苦也是同一种封闭造成的。对一件事物之美——无论它多么卑微——敞开着的人,一旦感受到这美,就不由自主地也对构想或做成它的人怀有敬意。在一件出于人手的事物之美里,我们感受到做它的人身上那份美的反光,感受到他制作它时所倾注的爱。当我们感受到这美、这爱,我们心里就不可能有俯就或轻蔑——正如在一个我们感受到某个女人的美、以及这美所标志的她身上那份力量的时刻,我们不可能对她有俯就或轻蔑。

\section{(38) 回归的冲动与更新}

迪厄多内身上时或闪耀的那种欣喜,想必触动了我身上某种深沉而有力的东西,这记忆此刻才如此强烈、如此鲜活地回到我心头,仿佛我刚刚还亲眼目睹。(其实我已经将近十五年没怎么有机会见到迪厄多内了,只有一两次匆匆照面。)当然,在意识层面,我并未对此给予特别注意——那不过是我这位年长的同行兼朋友那外放的性格中一个有点动人的特点,有时几乎滑稽。相反,对我重要的是:我在他身上找到了那个完美的、我甚至可以说是梦想中的合作者,来把那些将要为我眼前展开的广阔远景充当基础的东西,以一丝不苟的细心、一种充满爱意的细心,白纸黑字地写下来。直到此刻,当我把这两件事一并提起,
% ---------------- p.200 ----------------
\origpage{99} 那条联系才突然在我眼前显现:使迪厄多内成为一项伟大任务之理想仆人的——无论是在布尔巴基内部,还是在我们为另一项奠基性大工程而进行的合作中——是慷慨,是他的工作中、在他那些重大投入的抉择中,全无一丝虚荣的痕迹。我经常看见他隐没在他甘为仆役的那些任务后面,毫不吝惜地把无穷的精力献给它们,不求任何回报。毫无疑问,他什么也不求,却在工作本身、在他投入其中的那份慷慨本身里,得到一种充实与舒展——所有认识他的人都必定感受过。

我常常感到从他身上闪耀出来的那种发现之欣喜,在我心里立刻与一种相似的欣喜联系在一起——我曾在一个很小的孩子身上目睹过那种欣喜。两个记忆争相涌来——两个都让我重新看见我那时还很小的女儿。第一幅画面里,她大概才几个月,想必刚刚开始会爬。她大概是从人家把她放坐的那块草坪,爬到了旁边一条砾石小径上。她在发现那些小石子,一种无声的——却是起作用的——狂喜:她双手大把抓起石子,往嘴里送!——另一幅画面里,她大概一两岁,有人刚往金鱼缸里撒了鱼食。鱼儿们争先恐后地游过去,大张着嘴,吞下那些悬浮着、慢慢沉向缸底的小小黄色碎屑。小家伙此前从没意识到鱼也像我们一样吃东西。她身上仿佛突然一亮,迸出一声纯粹欣喜的叫喊:“妈妈看,它们在吃饭!”这确实值得惊叹——她在一瞬间发现了一个大奥秘:我们与一切其他生物同属一族的奥秘……

\origpage{100} 一个小孩子的欣喜里有一种无法用言语传达的、有传染性的力量,一种从他身上辐射出来、作用于我们的力量——而我们大多数时候总是竭力躲开它。在内心里安静的时刻,人能感觉到这力量随时都在孩子身上。只是有些时候它的作用更强些。这种围绕孩子的“力场”,在新生儿身上、在生命最初的日子和月份里最强。大多数情况下,它在整个童年期间都还感觉得到,随年岁渐长而一丝丝磨损,直到青春期——那时往往已经看不出痕迹了。然而,人们也能发现它辐射在各个年龄的人周围:在某些人身上,在得天独厚的时刻;在少数另一些人身上,则像一种气息或光晕,时时刻刻环绕着他们这个人。我有幸在童年认识这样一个人,一位男子,如今已经去世。

我也想到另一种力量,或能力:人们有时感到它从某个女人身上辐射出来,尤其是在她舒展于自己的身体、与身体交融的时刻。我常想到的词是“美”,它道出了其中一个侧面。这种美与任何美的规范、与所谓“完美”毫无关系;它不是青春的特权,也不是成熟的特权。它更像是一个人身上某种深层和谐的标志。这种和谐常常仍是残缺的,然而它正是通过这种辐射显现出来——那是力量的标志。这是一种把我们引向它所自出的中心的力量——或者不如说,它在我们身上唤起一种深沉的冲动:回到那位“女人—母亲”的身体里去——我们在生命破晓时,正是从那里出来的。它的作用有时强不可挡,当它来自所爱的女人时,令人翻覆。但对于不刻意对它封闭的人,它在每一个让这种美、这种深层和谐在自己身上舒展的女人身上都感觉得到。

从孩子身上辐射出来的力量,与从那个爱着自己身体的女人身上散发的力量是近亲。后者不断生出前者,正如孩子不断从母亲生出。但童年的力量,其本性不是吸引,也不是排斥。这种力量对不躲开它的人所施的谦卑而谨慎的作用\origpage{101},
% ---------------- p.201 ----------------
是一种更新的作用。

\section{(39) 夜之美、日之美(或:奥吉亚斯的牛圈)}

对我某个孩子的惊叹的记忆,是在五十年代末、六十年代初。如果说后来出生的其他孩子没有给我留下类似的记忆,那也许是因为我自己惊叹的能力已经钝了,我变得太疏远,无法与某个孩子的欣喜相通,甚至无法仅仅做一个见证者。

我还从未想过要追索这种能力在我一生中——从童年到今天——的变迁。那里想必有一条线索,一个极其灵敏的“探测器”。我从未想过追随这条线索,想必是因为这种能力的性质如此谦卑、外观几乎如此微不足道,以至于我几乎不会想到去特别注意它——我当时正全神贯注于发现和探测我所谓的我一生的“伟大力量”(它们今天仍在我身上显现)。然而,这种外观如此谦卑的能力,却提供了一个胜过一切的标志,指示着我们身上那份最稀有、最珍贵的力量之在与不在……

在我整个成年生活中,我从未与这种力量完全隔绝。无论我的生活在别的方面变得多么荒芜,我总在爱中重新找到孩子的惊叹、发现的欣喜。穿过许多片沙漠,爱的激情始终是与某个我已离开的东西之间那条活生生而茁壮的纽带,一根脐带,默默继续以温热而慷慨的血液滋养我。而在很长一段时间里,对所爱女人的惊叹,也与对她所生育的新生存在的惊叹不可分割——那些崭新的存在,无限娇嫩而生机勃勃,见证着她的力量,也承继着它。

但我这里的用意,主要是想稍稍追索这种“无辜之力”在我作为数学家的生活中的变迁,即我属于“数学家的世界”那个时期, 1948到1970年。当然,惊叹从未像在爱的激情中那样渗透我的数学激情。奇怪的是,若要我回想数学工作中某个欣喜或惊叹的特定时刻,我一个也想不起来!我接近数学的方式,从十七岁全力投入时起,就是给自己提出大任务。这些任务从一开始就总是“整理”、大扫除的任务。我看见一片表面的混乱,一堆杂乱异质之物,有时是不可称量的雾气,它们显然应当有一个共同的本质,藏着一个秩序、一种尚被遮蔽的和谐,要靠耐心、细致、往往旷日持久的工作把它揭示出来。这工作常常先是拖把和大刷子的活计——那粗活已经耗去大量精力——然后才轮到用掸子做收尾;收尾我不那么着迷,但也自有其魅力,而且无论如何显然有用。日复一日的工作中有一种强烈的满足:看见那个被猜到的秩序一点点显露出来,它总是比先前瞥见和猜到的更为精致,质地更为丰富。工作不断充满意外插曲,它们多半从考察某个看似极微小、此前一直被忽视的细节中冒出来。常常,对某个“细节”的打磨,给多年前做过的工作投下意想不到的光。有时它也带来新的直觉,而对它的深化又成了另一项“大任务”的对象。

这样,在我的数学工作中(除了1954年前后那个我有机会谈过的“艰难的一年”),有一种持续的悬念,注意力始终被吊着。何况,对我那些“任务”的忠诚不许我跑得太远,我咬着嚼子,急不可耐地想
% ---------------- p.202 ----------------
\origpage{102} 把它们全都做完,终于纵身跃入未知、真正的未知——而这些任务的规模此时已经变得如此之大:要把它们做完,即便有那些终于赶来救援的善意相助,我的余生也不够用!

我的主要向导,是不断寻求一种完美的融贯、一种完整的和谐——我猜到它就在事物动荡的表面之下,我耐心地努力把它揭示出来,从不厌倦。那确实是一种敏锐的“美”感,它是我的嗅觉,也是我唯一的罗盘。我最大的快乐,与其说是在它出现在光天化日之下时凝视它,不如说是看着它一点点从它喜欢不断躲藏其中的阴影与雾气的衣袍下显露出来。当然,不把它带到最明亮的日光下,我是不肯罢休的。那时我有时也体会到凝视的圆满——当一切可闻的声音汇成同一片浩瀚的和谐。但更常见的是:被带到日光下的东西,立刻又成了再一次潜入雾中的动机与手段,去追寻那位永远神秘、永远未知的“她”的又一次化身——她不断呼唤我,为了再认识她……

\origpage{103} 迪厄多内的快乐与欣喜,我觉得,主要在于看见事物之美在光天化日之下显现;而我的快乐,首先在于在雾与夜的皱褶里追寻它。这也许就是迪厄多内与我接近数学的方式之间的深刻差异。对事物之美的感觉,至少在很长一段时间里,在我身上想必不比他弱——尽管它也许在六十年代里、在一种自得的作用下变钝了。但看来,在迪厄多内身上通过惊叹显现出来的对美的感知,在我身上采取了不同的形式:不那么静观,更爱动手,在所感受和表达的情绪层面上也不那么外显。若果真如此,我要追索的,就应当是我身上对数学事物之美的那份敞开之变迁,而不是那份神秘的“惊叹之馈赠”。

\section{(40) 竞技数学}

相当清楚的是,对数学事物之美的敞开,在我身上从未完全消失,即使在六十年代直到1970年——那时自得在我与数学、与其他数学家的关系中逐步占据了越来越大的位置。若没有最低限度对事物之美的敞开,我根本不可能作为数学家“运转”起来,哪怕只开到最低档——而且我怀疑,任何人若身上不多少还活着这种美感,能否在数学里做出有用的工作。依我看,造成这位数学家与那位数学家之间、或同一位数学家的这项工作与那项工作之间差别的,与其说是某种所谓的“脑力”,不如说是这种敞开或敏感的精细度、细腻度——研究者彼此之间不同,同一个研究者不同时刻也不同。最深刻、最丰产的工作,也是那种见证着最敏锐的敏感性的工作——那敏感性用于把握事物隐藏的美\fnoteP{10}{(36)}{在我看来,这种对美的精细敏感,与我曾有机会以“严苛”(对自身)或“严格”(就其完整意义而言)之名谈过的东西密切相连;我曾把它描述为“对我们自身之中某种精细之物的注意”,一种对所受探究之物的理解品质的注意。对一件数学之物的这种理解品质,不能与对该物特有之“美”或深或浅、或完美或粗略的感知分开。} (36)。

若果真如此,那就得相信,这种敏感在我身上一直鲜活到最后,至少某些时刻如此——因为正是在六十年代末\fnoteP{11}{}{(8月8日)核实之后看来,我关于 Motif 的思考之开端是在六十年代初,而不是六十年代末。},我开始瞥见并稍稍揭示出我所发现的最隐蔽、最神秘的数学对象:我称之为“Motif”的东西。
% ---------------- p.203 ----------------
\origpage{104} 它也正是在我作为数学家的生活中对我施以最大魅惑的东西(若不算最近几年的一些思考——它们其实与 Motif 的实在密切相关)。毫无疑问,若不是我的生活突然转向一条完全意外的道路,把我远远带离数学事物那宁静的世界,我最终会追随这强大魅惑的召唤,把那些一直把我囚禁着的“任务”撂在一边!

也许我可以说:在我工作室的孤独里,美感一直保持不变,直到1970年我第一次“觉醒”的时刻,并没有真正受到那种常常标示着我与同类关系的自得的影响?某种“嗅觉”甚至想必随年岁而愈发精细,在与数学事物每日而亲密的接触中。我们对事物所能拥有的亲密认识,有时使我们能把握超出当下所知的东西,并更深入地进入认识——这种认识或这种成熟,以及作为其最显眼标志的这种“嗅觉”,是向事物之美与真之敞开的近亲。它助长、刺激这样的敞开,而且它是此前一切敞开时刻、一切“真理时刻”的总和与果实。

那么,我接下来要考察的,就是:这种对美的自发敏感,在我与这位或那位同行的关系中、在它有机会显现的那些时刻,在多大程度上受到了或深或浅的扰乱。

记忆在这方面交给我的东西,并不能凝成一个有形的、确切的事实,让我在这里多少详尽地摆出来。这里的记忆仍然只是一团雾,但它还是交给我一个总体印象,我得设法把它界定出来。那是某种内心态度留在我身上的印象,它想必最终变成了第二天性,每当我接到一条数学信息、而所涉之物多少“在我的射程内”时,它就显现出来。说实话,从某个相对无害的侧面看,这种态度想必一直属于我,它属于某种气质,我也曾有机会顺带触及。说的就是这种反射:首先只肯了解一个陈述,绝不看它的证明,先试着把它安放到我所已知的东西里,看看用这些已知的东西来说,这陈述是否变得透明、显然。这常常使我以或深或浅的方式重新表述那个陈述,使之更一般或更精确,往往两者兼有。只有当我无法用我的经验和我的意象去“安放”这个陈述时,我才准备好(有时几乎违背我自己的意愿!)去听(或读……)那些来龙去脉——它们有时给出事情的“那个”道理,或者至少给出一个证明,懂不懂另说。

这是我接近数学的一个特点,我觉得它使我有别于当时布尔巴基的所有其他成员,也使我实际上不可能像他们那样插入集体工作。这特点想必也在我教书时构成一种障碍,我的所有学生想必都感受过——直到今天(多亏上了年纪)它才终于松动了一些。

我身上这个特点,想必已经是一种敞开上的欠缺。它只意味着一种部分的敞开,只肯接纳“来得正好”的东西,至少对别的一切都极不情愿接纳。在我的数学投入之抉择中,在我同意把时间花在这样或那样不期而至的信息上的抉择中,这种有意“部分封闭”的打算,今天比任何时候都更强。它甚至是一种必要,如果我想能追随那最令我着迷之物的召唤,而又不再把“我的命拿去吃”献给数学夫人!

% ---------------- p.204 ----------------
\origpage{105} 然而这团“雾”交给我的,不止这个特点——我几年前就已经察觉它了(晚知道总比不知道好!)。在某个时刻,这种反射变成了某种关乎荣誉的事:我要是在比说这句话更短的时间里还“拿不下”这个陈述(假定它对我还不算太熟悉),那才叫见鬼!如果陈述的作者是个有名有姓的陌生人,还多一层微妙:那可不行,我(毕竟人家以为我是在行的人!)岂能手里没有这一套!而常常我确实有,而且绰绰有余——那时我的态度就倾向于这个方向:“行了,您可以回去穿好衣服了——等您做得好一点再来吧!”

这恰恰就是我在“那个在我园子里乱走的毛头小子”一事上的态度。我甚至不敢发誓:他做的东西里没有一些有趣的细节,是我那些“秘密笔记”里所做的东西不曾覆盖的——不过这终归是无关紧要的事\fnoteP{12}{}{(8月8日)后来我发觉,这件事并不那么“无关紧要”:它构成了从“竞技态度”通向某种不诚实之开端的界线,而这条界线,我也许有几次越过了……}。说到底,这段插曲也照亮了我在这里考察的问题:对数学事物之美的那份敞开之深层扰乱。仿佛从我把某样东西“做出来”的那一刻起,它的美对我而言就消失了,剩下的只有一种虚荣,要求为它记功受利。(而我却又不屑于花时间去发表它——诚然,要发表的也太多了。)那是一种典型的占有态度,类似一个男人,认识了一个女人之后,再也感觉不到她的美,又去追逐另外一百个,却仍容不得别人认识她。这种态度我在爱情生活中是斥责的,自以为远在这样一份虚荣之上,却又小心翼翼地不去确认这个显而易见的事实:我对数学的态度恰恰正是如此!

\origpage{106} 我有个印象:我刚才在自己身上指出的这些粗鄙的竞争心性、不妨说“竞技”心性,想必大约就在它们于我自己身上变得普遍的那个时候,开始在我的数学“圈子”里变得普遍起来。要我在时间上定位它们出现的时刻,或它们变成这个圈子里所呼吸的空气之一部分的时刻,或我的学生们在与我接触中所呼吸的空气的时刻,我却无能为力。我唯一敢说的,是这大概在六十年代,也许早在六十年代初,或五十年代末。(若果真如此,我的学生全都摊上了——他们只能要么接受,要么拉倒!)要定位它,我还需要别的确切例子,而此刻它们完全逃出了我的记忆。

当然,这个卑微的实情,与我为自己同数学、同一般年轻研究者的关系所描绘的高贵形象完全相反。用来蒙骗我自己的那个粗劣托辞,灵感来自功绩主义:对于这个形象,我所留住的只是我与学生们的关系(他们为我的声望添砖加瓦,是我声望中最高贵的花朵!),以及我与那些特别出色的年轻数学家的关系——我懂得承认他们的功绩,并且像对待我的学生一样平等待他们,不等他们的头上戴上桂冠(当然,桂冠也没等多久——嗅觉这东西,有就是有,没有就是没有!)。至于那些既没福气做我学生、也没福气做我某位朋友的学生、又不是年轻天才的年轻人,我丝毫不在意我与他们是什么关系。他们不算数。

我相信,当我与年轻研究者有了个人的接触——无论是在讨论班上遇见他,还是他写信给我——上述实情多半会松动、缓和一些。从这一点看,“毛头小子”那件事也许是个稍显特殊、例外的情形。我觉得,对我刚才谈到的那些研究者,我大概会多少把他们看作“受我庇护”的人,这想必在我身上唤醒了一种更和善的态度。在这种情形里,我想出头的欲望也能找到出口:对当事人发表我的评论,建议他也许从更广阔的视角、或者更深地进入事情本身去重做他的工作。在这种情形下,那位在有限时间里多少扮演了学生角色的年轻研究者,很可能也从中有所得,并且对我与他的关系留下好印象。(这方面任何或正或反的回声,我都欢迎。)

% ---------------- p.205 ----------------
\origpage{107} 我在这里主要想到的是较年轻的研究者,而“竞技”态度绝不限于我与他们的关系,这不用说。但确实是在与年轻研究者的关系中,一位有名望的数学家的影响——无论心理上还是实际上的——往往最强,对他们未来的职业生涯也最有后果。

\section{(41) 克里希那穆提,或:成为羁绊的解放}

\origpage{108} 这一夜,我停在一份宽慰之情、一份极大的满足上——那是一个没有白费时间的人的满足!我突然觉得轻松而快乐——一种有时带点促狭的快乐,迸成顽皮的笑——一个捣蛋鬼恶作剧的笑。其实我并没做什么大事,我不过是从一个略微不同的角度,重看了一段已经“知道”的插曲,即那个著名的“毛头小子……”。这个角度显示出的,是我与数学本身的关系——在某些情形下,不只是我与数学家的关系。就这一点,已经足够让一个我曾珍爱的神话化为青烟。

说实话,这并非我第一次审视我与数学的关系。两年半以前,我就曾被引向为此花上几个星期或几个月。那时我(在许多别的事情之外)意识到:在我过去对数学的投入中,自我之力、自我膨胀之力有多么重要。但昨夜我刚刚点到了一个当时逃过我注意的侧面。现在我回过头来谈它,发现这个侧面——即我与数学的关系中那种嫉妒态度——与第一夜“冥想”结束时到来的那个“再平常不过”的发现接上了(当时我并不知道自己在冥想,就像茹尔丹先生不知道自己是在说散文……)。这很可能在我随后那份欢欣的狂喜中占了一份。即便没有被自觉地察觉,那也多少像是对我从前发现的某样东西在新光照下的再次确认——而这时心里的快乐,与数学里的一样:当你并没有去找,却从一条完全不同的路径撞上你认识的、也许多年前就发现的东西。每一次,这都伴随着一种深切的满足感:事物的和谐再一次显露,而我们对它的认识也或多或少被更新。

而且我相信,这一次我确实“转完了”!好几天来我都清楚地感到还有什么东西没被挖出来,却又说不清到底是什么。我没有试着硬来,我觉得只需让它自己来,让所追随的那条线自由展开,穿过一片片既熟悉又意外的风景。说意外,是因为我此前从未费心去看它们。我是以散步的步子走近那个剩下的“热点”的。而我相信,这是我刚刚走完、且已接近尾声的这趟旅程中的最后一个。

而刚一到达这一点,我就有了一种感觉,像一个登上观景台的人:从那里,他看见\origpage{109}自己刚刚走过的那片风景铺展开来,而在途中每一刻,他只能看见其中一角。
% ---------------- p.206 ----------------
\origpage{110} 而现在有了这种辽阔与空间的感觉,那是一种解放。

如果我试着用文字表述眼前这片风景交给我的东西,出来的就是这些:这些年来,在我作为数学家的生活中到来的一切——常常来得不是时候,也不受欢迎——都是我在还属于数学家的世界那些年月里所播种之物的收获与讯息。

当然,这件事,这些年来我已经对自己说过又说,在我刚写下的这些笔记里也说过。我对自己说过,多少类比于另一些执拗地到来的收获——我长久拒斥,最终接纳并化为己有。从我这样接纳的第一份收获起,甚至在我懂得冥想之前,我就明白了:每一份收获都必有其意义,而赌气只是回避一个意义、推迟一个了结的期限。这一认识对我很宝贵,因为它常使我免于自怜,也免于那种义愤——义愤往往是自怜的一种伪装形式。这认识在我身上像是一种半成熟,它还远没有终结那种积习难改的反射:当收获显得苦涩时加以拒绝。当我说“赌气没有用”,收获并没有因此被接纳。我也许既不自怜也不义愤,可我仍在“赌气”!只要那盘菜没吃下去,它就还没被接纳——而不吃,就是赌气。

接纳并吃下去是一件工作:某种能量在“工作”,一种工作在光天化日下或在暗处进行,有东西在转变……而赌气则是一种能量的浪费,那能量散掉了——散在“赌气”上!而吃、消化、吸收这件工作,是省不掉的。仅仅经历一些事件、仅仅“取得”或“获得”一种经验,与一件工作毫无共同之处。它只是一份可供工作的材料,做不做由你。从我遇到数学家的世界起,三十六年了,我一直在使用这份自由:回避一件工作,而可供吃、可供消化的材料却年复一年地增多。我从昨天起感到的这份欢欣的解放感,确实是一个标志:我面前那件工作——我不断地以别的工作或任务为借口推开的——终于做完了。也确实到时候了!

现在还不能确定事情果真如此,还不能确定没有剩下某个阴暗而顽固的角落逃过了我的注意、我还得回到那里去。但同样真实的是:这种解放感不会骗人——我一生中每一次感到它,事后都能证实它确实是一次解放的标志;是某种持久的、已经获得的东西,是一种理解、一种认识的果实,而那认识已成为我的一部分。若我愿意,我有自由无视这认识,把它埋在我想要的地方、按我想要的方式埋掉。但要毁掉它,却不在我、也不在任何人的能力之内——正如人不能毁掉一只果子的成熟,不能把它弄回一种已不属于它的青涩状态。

再一次证实我并不比别人“更好”,这是一大宽慰。当然,这也是我相当经常地对自己重复的话——可是重复和看见,毕竟不是一回事!没有孩子的无辜与灵动——他看见,就像他呼吸——要看见明显之事,往往需要一件工作;好了,做完了,我终于看见了这一件:我并不比几天前还让我“喘不过气”的那些同行或旧学生“更好”!想想我如今卸下的这副担子有多重吧!自认为比别人强,从某种意义说也许挺受用,但那也非常累人。那甚至是精力的极大浪费——正如每一次要维持一个虚构时那样。人很少意识到这一点,但光是顶着风浪维持那个虚构,就已经需要精力了——而明显之事在每一场合都冲着我那仔细塞住的耳朵喊:这是假的,看看吧,傻瓜!看见,有时也许是一件工作,
% ---------------- p.207 ----------------
但一旦做完就做完了。这就一劳永逸地省掉了两件事:一是像那样动不动就堵上眼睛和耳朵地四处走动——这本身也是桩苦差!二是每当有件我自己不小心搁在那里的东西落到我头上时,就把它当作不可容忍的侮辱而自苦。

受够了这套把戏!看清了把戏,就已经出来了。付了钱,没错,我有权一直在里面转到老死,甚至有义务——那也无妨,所有人都会这么对我说:权利、义务——看人下菜碟。所有这些是权利又是义务、是义务又是权利的东西,也累人得很,当我自认为比别人强的时候,它们就黏上我。说到底这很正常:既然是强者,就悄悄收着(这就是“权利”),并且“付账”,为了人类精神与数学的荣誉尽全部义务——真的很美,没错,荣誉、精神、数学,谁还能说得更好,好极了!再来一个!是很美,没错,可也非常累人,到头来会让人落枕。我落过枕了,现在够了——我把位置让给别人去挺直腰板吧。

这也正常(既然我谈到学生):学生超过老师。我曾为此恼怒,我真有精力可浪费!这些都结束了!

多么轻松!

% ---------------- p.208 ----------------

% ================================================================
%  第 10 章 孩子嬉戏 (原书 p.209-218, 边码 111-121)
%  源文: _work/src1/task-ch10.md
% ================================================================
% ---------------- p.209 ----------------
\chapter{孩子嬉戏}

\section{(42) 孩子}

肯定还有些角落,扫帚没扫到。没关系,它们自会引起我的注意,到那时再收拾也不迟。至于我那大名鼎鼎的“数学家过往”,大扫除已经做完,这一点毫无疑问。

既然我又一次看清我并不比旁人更好,那就别再掉进那个老掉牙的陷阱:把自己看成\textbf{比我自己更好}!——把如今的我,退出了那台旋转木马、诸如此类的我,看成比十五年前、或者十五天前的我更好。这十五年里我学到了些东西,这是肯定的;这十五天里也是,甚至从昨天起也是。我学到东西的时候,就成熟了,就不再完全是同一个人了。学到某样东西之后的我,并不比那样东西还摆在我面前时的我“更好”。一个更熟的果子并不比一个不太熟的、或者青的果子“更好”。一个季节并不比它前一个季节“更好”。最熟的果子味道可能更可口,也可能更不可口,这要看各人的口味。一年一年过去,我在自己的皮囊里觉得更自在了,想必我身上发生的变化“合我的口味”——可这些变化并不合我所有朋友或亲近者的口味。每次我重新做起数学,四面八方都来恭维,腔调是这样的:“他当初干吗要去弄别的东西呢!现在一切都归了位,早该如此!”看着一个人变,总叫人不安……

\origpage{111} 我学习,我成熟,我变化——以至于有时我很难在从前那个人身上认出自己,得靠一段回忆、或者旁人一次意想不到的见证,才重新发现他。我在变;同时也有某种东西始终“是同一个”。它一直都在,从永远起就在,肯定从我出生起,也许还更早。近几年,我觉得自己渐渐能把它认准了。我把它叫作“孩子”。凭这东西,此刻的我不比一生中任何别的时刻更好;它就在那里,尽管往往很难猜出它的在场。也凭这东西,\textbf{我不比任何人更好},也没有\textbf{任何人比我更好}。在某些时刻、在某些人身上,孩子更在场一些。而这是件很让人受用的事。这并不意味着某个人比另一个人“更好”,或者比他自己的另一个时刻“更好”。

常常,我做数学的时候,或者做爱的时候,或者冥想的时候,是孩子在玩。并不总是只有他一个在“玩”。可他不在的时候,就没有数学,没有爱,也没有冥想。装样子是没用的——这出戏我很少演。

% ---------------- p.210 ----------------
\origpage{112} 当然不只有孩子。还有“我”,还有“老板”,或者说“大老板”,爱怎么叫就怎么叫。对企业运转来说,老板肯定不可或缺。既然有个老板,总该是为着什么。他照管庶务,而且像所有老板一样,有个讨厌的倾向:越占越宽。他把自己看得无比要紧,非要处处压过对面的老板。占得宽也好,不宽也好,他不过是老板,他不是工人。他组织,他指挥,他也收账——当然!收益他当作自己应得之物收下,亏损他当作侮辱来承受。但他什么也不创造。只有工人有创造的力量,而工人不是别人,正是孩子。

老板和工人合得来的企业,很少见。多数时候,根本看不到工人的影子,他被关在天知道什么地方。是老板做出顶替他站在车间里的样子,结果可想而知。也常常是这样:工人明明就在,老板却对他开战,或是猛烈的战争,或是小股交火——这样的车间里出不了什么像样的东西!有时老板对工人也有一种带戒心的宽容,他一边嘀咕一边由着他干,眼睛却一刻也不离开他。那就像一场从未停歇的战争里,不断续签的停战。工人也就趁着这停战,勉强干上一点活。

我刚才那场冥想的功效,绝不意味着我对数学的那份占有态度就此如中了魔法一般消失了!至少,我得近得多地去看占有欲的种种表现——刚才我不过是叫出了其中一种的名字,轻轻碰到了它一下。这里不是地方:这篇“导论”已经变成了“导论性的一章”,而这一章自己也已经显得长了!不过昨夜有一件事“咔嗒”一声对上了,此刻我想稍稍回过头去说它一说;那是两三年前我带着几分惊讶记下的一件事。

当时我正钻在一个数学问题上,已说不上是什么问题了;某一刻(不知由于什么情况),我发现我盯着的问题也许早有人盯过,它很可能在某本书里白纸黑字地处理过,只消我去图书馆查一查。单是想到这种可能性,效果就如雷击,把我惊得发呆:转眼之间,欲望不见了。忽然之间,那个我也许已经花了好几星期、还打算再花上几星期的问题,对我完全失去了吸引力!那不是赌气,那是兴趣突然而彻底的丧失。就算那本书就在我手里,我也懒得翻开。

\origpage{113} 事实上,那种可能性并未得到证实,于是欲望又回来了,我照旧一路做下去,像什么都没发生过。我却仍然愣着。当然,如果我正在做的东西真真是我另做一件事所\textbf{必需}的,那就不会有那么惊人的兴趣跌落。我常重做已知的东西,明知或猜到它们已知,却一点也不在乎。那时我正处在一种冲劲里:照我自己的方式、按它们向我呈现的样子去做,比到书里、文章里去翻找更省事,尤其有意思得多。那时我是“顺势”奔向别的东西,奔向欲望把我带往的地方。当然,我够“在行”,知道自己最终要到的那个地方,任何书、任何文章里都没有。

这让我想起来:数学工作即使孤零零地做了许多年,也不是纯粹个人的、个体性的工作,不像冥想那样——至少在我身上不是。我在数学里追的那个“未知”,要能那样强有力地吸引我,就必须不只对我未知,而是对所有人都未知。数学书里写着的东西不是未知,即便我自己从未听说过它。读书或读文章从来不吸引我,只要\textbf{能躲我就躲}。它能告诉我的,从来都不是未知;
% ---------------- p.211 ----------------
我对它的兴趣,也没有欲望的那种质地。那是一种“权宜的”兴趣,是对可能对我有用的一条\textbf{信息}的兴趣,是某个欲望的工具,而这条信息根本不是那个欲望的对象。

细想之后,我不觉得我讲的那件事是什么嫉妒、占有的心性的迹象,也不是什么受了挫的虚荣的迹象。我心里没有半点赌气,没有半点失望,只是一个前一瞬还强烈的欲望突然消失了。那是在一个我压根没想过要发表任何东西的时期,也没想过有一天我会忽然又想发表点什么。这欲望不是虚荣的表现,不是急于积累知识、头衔和信誉的饕餮——它确确实实是一个真实的欲望,是沉迷于游戏的孩子的那种欲望。而忽然之间——什么都没有了!谁懂就懂吧,我是不懂……抱歉!

\section{(43) 扫兴的老板——或高压锅}

\origpage{114} \noindent\pmark 我觉得,这场对我数学家一生的回顾,终究是做完了。当然,我没有把题目做尽——那得要好几卷书,假如这样的题目真能“做尽”的话。那不是我的意图。我的意图是弄清楚:在某种我如今只能一阵一阵感觉到的“气息”的出现里,我到底是不是参与者、是不是共同行动者;如果是,又是以什么方式。现在我弄清楚了,这让人舒坦。再往前走,去深挖那些只是瞥见或触及的东西,可能会很来劲。可看、可做、可发现的有意思的东西太多了!至于我作为数学家的过去,我觉得,为了承担这段过去而\textbf{必须}看的,已经看到了。

肯定,若把这番冥想往深里做,我准会学到许多关于我当下的有意思的东西。这项工作几乎每走一步都让我感到的一件事是:我对这段过去还眷恋得多么深,它直到今天在我的自我形象里、以及在我与他人的关系里,分量有多重;尤其是在我与那些——从某种意义上说——被我抛下的人的关系里。这番工作过程中,我与这段过去的关系肯定变了,变得疏离了些,或者说轻了些。往后自有分晓。但很可能还会留有一份眷恋,只要我的数学激情还没烧尽、还没餍足——只要我还“做数学”。至于它会不会先我而熄灭,我一点也不想去猜、去预言……

十多年里,我一直以为这激情已经熄灭。更确切地说,是我径自宣布了它已熄灭。那正是我暂时停下做数学、并且重新发现世界的那一天!此后三四年,我全力投入一项如此紧张的活动,旧日的激情想必找不到一丝缝隙钻出来现身。那是紧张学习的几年,不过在某个层面上还相当浮浅。在那之后的几年里,数学激情以突如其来的、完全无法预料的发作显现出来。每次发作持续几个星期或几个月,而我执意不去理会它们那相当清楚的含义。我早就一次性决定了:做数学的那股馋劲,既然明摆着百无一用,从此就是过时的玩意儿,句号!可那“百无一用”的东西并不这么听——而我这边,照旧充耳不闻。

\origpage{115} \noindent\pmark 有件事也许显得悖谬:恰恰是在发现冥想之后(1976年)、在我的生命里进来一种新的激情之后,旧激情才重新露面,而且格外猛烈,几乎算得上暴烈——仿佛每次都是一只锅盖在过强的压力下崩开。又过了五年,在形势的推动之下——这话说得不假——我才费心去察看究竟发生了什么。那是我做过的最长的一次冥想,针对一个看来范围分明的问题:我花了六个月顽强而紧张的工作,才绕着一座冰山走了一圈;它那可见的尖顶终于碍事到逼着我——几乎是违心——去看个究竟。不得不确认的是一种冲突局面,而从表面看,那是两股力量、两种欲望的冲突:冥想的欲望,和做数学的欲望。

% ---------------- p.212 ----------------
在这场漫长的冥想中,我一步一步学到:做数学的欲望——我一向鄙弃它——和冥想的欲望一样,都是孩子的欲望,而后者我捧到了天上。孩子才不管大老板、大掌柜那份鄙弃,也不管他那份谦逊的傲气!孩子的欲望一个接一个,随着时辰和日子,像一支舞的动作,一个从另一个里生出来。它们的本性就是如此。它们并不彼此对立,正如一首歌的各个诗节并不彼此对立,或者一首康塔塔、一首赋格的各个乐章并不彼此对立。是老板这个蹩脚的乐队指挥,宣布这个乐章“好”、那个乐章“坏”,在本来和谐的地方制造出冲突。

这场冥想之后,老板安分了些,不太做出把鼻子伸到不关他事之处的样子。这一回工作拖得很长,而我原以为几天就能完。工作一做完,“结果”就显得明摆着,几句话就能说出来\fnoteP{1}{(37)}{我想,几乎不必补充:这项旷日持久的工作,一天一天地还显露出别的东西,远不止我刚才以碑铭式的简短形式交出的那个“结果”。冥想的工作如此,为一个打算考察的特定问题所推动的数学工作也莫不如此。所走的那条路(它通向、或不通向对原问题或多或少的澄清)上的种种曲折,往往比原问题本身、或比“最终结果”更有意思。} (37)。可要是有个眼尖的人,在工作之前或进行中就把这几句话告诉我,恐怕也一点都帮不上我。工作之所以拖这么久,是因为阻力又强又深。再说,老板这回被骂得狗血淋头,却一声没吭,因为当时的气氛根本不容他发火。可以肯定的是,这六个月用得其所,而且是省不掉的;正如一个女人省不掉九个月的妊娠,到头来生下一个像娃娃一样“明摆着”的东西。

\section{(44) 再次倒转蒸汽}

到这里,我已经一年半没冥想了,只在十二月里有那么几个小时,为的是看清一个紧迫的问题。而这一年来,我把大部分精力投在数学上。这一“潮”来得和别的潮一样,数学潮也好,冥想潮也好:它们来时不打招呼。就算打了招呼,我也从来听不见!老板对冥想还留着一点小小的偏好,看来不得不信:每一次,我原以为会永远持续的冥想潮,转眼已被数学潮接替;而那数学潮——在我看来不过是几天、最多几星期的事——却拖宕下去,铺展成几个月,也许甚至,谁知道呢,成几年。不过老板终于明白:这些节律不是他造出来的,想去调节它们,他什么也捞不着。

但也许老板那点“小小的偏好”终究翻转过来了,因为近一年来事情已经说定、已经决定:我至少要用几年工夫去“重做数学”,可以说是正式地——我甚至申请了国家科研中心(CNRS)的一个职位!更重要、而且一年前还完全想不到的是:我又开始发表东西了。即便在刚才提到的1981年那场冥想之后,在做数学的欲望不再被当作穷亲戚对待之后,我也不会想到自己还能重新发表数学。顶多是别的东西:一本谈冥想的书,或者谈梦与做梦者的书——再说,我当时忙得根本没闲心去写一本关于它的书!
% ---------------- p.213 ----------------
写它干吗?!

于是就有了一个相当重要的决定,它牵动我今后若干年的生命进程,而它是在一种捎带着的情形下作出的,我甚至说不清是何时、怎样作出的。\origpage{116}\origpage{117} 有一天,当打字的笔记已经攒成厚厚一沓时(瞧瞧!直到那时,我那些数学上的苦思,我都还只是手写了事……\fnoteP{2}{(38)}{这些笔记其实是那封写给……的长信的继续,那封信成了它的第一章。它们被打字机打出来,是为了让这位旧日的朋友、以及另外两三位(尤其是龙尼·布朗(Ronnie Brown))读得清楚,我本以为他们可能会感兴趣。再说,那封信从未收到回音,收信人根本没有读它;将近一年之后(我问他是否收到了),他竟真心感到惊讶:我居然会有一刻以为他会读它,鉴于人们该期待我做的那种数学……} (38),关于栈、同伦模型等等),事情就成了定局:这个,发表!既然要做,索性做足,开一个小系列的数学反思,名字现成,只要把字母大写:“数学反思”!大致就是这样——此刻那团大名鼎鼎的“迷雾”向我复述的,就是这些;这迷雾常常替我充当记忆。这一回的回忆,肯定大大删削过了。无论如何,值得注意的是:这件事做成的时候,我连停下来\textbf{看看}自己往哪儿走、是什么在推我、带我的工夫都没有……这倒是我还想再做一次的:顺着这场不期而至的冥想的势头,好让我能觉得它真正做完了。

立刻浮上心头的问题是:我刚指出的这件“值得注意的事”,是老板那种(所谓的?)“谨慎”的迹象——他无论如何都不肯插手(哪怕只是投去一瞥不检点的目光……)一个如此美妙、根本不需要他的自发运动,如此等等;还是\textbf{相反},这迹象表明他已经断然站了边,而那所谓的“小小的偏好”正把他朝数学的方向全力推?

只消把问题白纸黑字写下来,答案就自己出现了!并不是那小子——他这一回玩起的把戏也许比别的把戏更长久——就此宣布自己要连着 X 年不停手地干下去,老老实实地涂黑纸张,涂够所需的那么多页,凑成漂亮一套大写字头丛书里数目可观的若干卷!是老板把一切都预见到、安排好了,小子只消照办。也许小子巴不得如此,这没法预先知道——但那是次要问题。再说,小子的兴致至少在一定程度上取决于\textbf{境况},而境况主要取决于老板。

老板选了边,再清楚不过。再说他刚显出过某种灵活:一个多月来,一场冥想就在他和蔼的目光下进行。诚然,他的和蔼绝非无私,因为冥想的有形成果,我正在写的这些笔记,将成为他已经看见自己建起的那座塔最美的奠基石,而砌塔的石头,是那位看来很听话的工人-孩子乖乖凿好的。说真的,现在夸他“灵活”还早了点儿!三个月前那几小时冥想,一年半里头总共就这么点,未免太寒碜了!

\origpage{118} \noindent\pmark 然而,我并不觉得这一整段时间里,有过什么被压抑、被挫败的冥想欲望。十二月的那几个小时里,我把事情理清了,看见了该看见的;这就足以改变一个本来不清楚的局面。我重新接上中断的数学工作,而不必把别的什么一刀斩断。我不觉得有冲突蹑手蹑脚地又回来了——我是说,那个两年多前已经化解、这一次会以颠倒的形式重新出现的冲突。老板有偏好,这是他本性如此,也完全是他权利——他要是装作禁止自己有偏好,那才蠢呢(虽然比这更蠢的事也有……)。这不是冲突的迹象,尽管冲突常由它而起。事情到了这一步,实在看不出有什么该怪他不灵活的!

% ---------------- p.214 ----------------
这一点看清了,剩下的是试着圈定老板的“动机”——为了这次蒸汽倒转;它做得再低调不过,可仔细看来,其实相当惊人。

\section{(45) 不是上师的上师——或三条腿的马}

这立刻把我带回1981年7月到12月的那场冥想;在那之前,我刚在一种数学狂热里过了四个月。这段有点疯的日子(从数学上看还颇有收获\fnoteP{3}{(39)}{这尤其是《穿越伽罗瓦理论的长征》(La Longue Marche à travers la théorie de Galois)的时期,《纲领概要》(Esquisse d'un Programme)里谈到过它(第3节:“与儿童画相关联的数域”)。} (39))随着一个梦,一夜之间结束了。那梦以一个具有野性而不可抗拒之力的比喻,描绘了我生命里正在发生的事——一个关于这场狂热的比喻。信息明如闪电,可我还是花了两天紧张的工作,才接受它那明摆着的含义\fnoteP{4}{(40)}{\emph{来访}\par 关于这个梦的工作,是一封用英文写的长信,写给一位朋友兼同行,他前一天曾在我家匆匆过了一趟。做梦者用来从一片表面的虚无里唤起这个逼真得惊人的梦的某些材料,显然取自这位亲爱的朋友那次短暂来访——我已经快十年没见过他了。因此,工作的第一天,与我以往的经验相反,我竟以为可以断定:我做的这个梦关系到我这位朋友,多于关系到我自己——该做这个梦的是他,不是我!这是一种回避梦境信息的方式,而那信息(我本该凭过去的经验一下子就明白)只关乎我一个人。我最终在工作这第一阶段——浮浅的阶段——之后的夜里明白了这一点;第二天我在同一封信里接着写下去。自那封值得记住的信之后,我再没有收到过这位朋友的任何音讯,他是我有过的朋友中最亲近者之一。\par 这次工作也是唯一一次以信件形式(而且还用的是英文)进行的冥想,因此我手上再没有它的书面痕迹。这个插曲让我格外震动——在许多别的插曲之中,它们都表明:一切超出某种门面的工作的迹象,一切把极简单、却通常被大家当作义务去忽视的事实摆到光天化日之下的工作,会在旁人心里引起多大的不适与恐惧。这一点我后面还要回来谈(见第47节,“孤独的历险”)。} (40)。这之后,我知道了自己该做什么。此后六个月的工作里,我再没回到那个梦上去,可我所做的,无非是更深入地进入它的含义、完整地消化它的信息。梦后第三天,这信息只在相当表面、相当粗糙的层面上被理解了。我尤其需要深挖的,是“我的”关系——我是说老板的——与当下这两种欲望各自的关系,而这两种欲望在我看来相互敌对。

\origpage{119} 自那场冥想以来,我生命里发生了这么多事,以至于它显得像是属于一个很遥远的过去。如果我试着把从它那里学到的、关于“老板”动机的东西说出来,就是这样:从“第一次觉醒”(1970年)算起、到那时已经过去的十二年里,老板押的注,明摆着是“\textbf{那匹劣马}”:\textbf{在数学与冥想之间}(他喜欢把两者对立起来)\textbf{他选了冥想}。

 这是一种说法,因为“冥想”这东西连同这个名字,只是五年前的1976年10月才进入我的生命。但在1970年那副刷新过的、我珍爱的自画像里,冥想六年之后来得正是时候,用它的光彩烘托出某种态度或姿态——这姿态早已被察觉,却直到1981年这场冥想才被审视。我把它叫作“大师综合征”,也有人(有道理地)称之为我的“上师做派”。我采用前一个称呼而不是后一个,多半是因为它便于把事情的性质搅混,而我乐意让自己耽在这种混淆里。我身上确实从小就有一个教人的自发乐趣,它与学的自发乐趣毫不相悖;
% ---------------- p.215 ----------------
也毫无做作的成分。在我与学生的关系里起作用的,主要就是这股力量;这关系是浮浅的,却强而纯正——我所谓“纯正”,是指没有做作。是在我称为1970年“觉醒”之后——一个我曾熟悉的天地退缩到几乎消失,连同学生、连同那些让我“教”点什么、让我把我知道的、对我有意义有价值的东西讲出来的机会一起退去——正是在那时,“老板”尽其所能地报了仇:他不再去教数学(那玩意儿只配用来糊口,除此之外配不上我新的伟大),我仿佛看见自己在用我的生活与榜样,教某种“智慧”。当然,我很小心,无论对自己还是对别人都不把这话说出来;而当我听到这个方向的回声时,我当然要矢口否认,并为这些朋友或亲近者如此不理解我而难过。我怎么解释也没用,他们就是不肯懂,真是些再叫人遗憾不过的学生!

\origpage{120} 我读过克里希那穆提(Krishnamurti)的一两本书,印象很深,脑子转眼就把某种信息和某些价值吸收了下来\fnoteP{5}{(41)}{\emph{克里希那穆提,或变成了羁绊的解放}\par 要说我从这次阅读里得到的只是一套词汇,以及一种把它据为己有、最终又理所当然地用它替换现实的倾向,那并不准确。我拿到手的第一本克里希那穆提的书之所以让我如此震动(何况我不过有闲读了几章),是因为他所说的,彻底掀翻了许多对我而言不言自明的东西;而我立刻发觉,那些都是\textbf{老生常谈},从来就是我呼吸的空气的一部分。同时,这次阅读第一次让我注意到一些影响深远的事实,尤其是“对\textbf{现实}的逃避”这一事实——它是心智最强大、最普遍的制约之一。这给了我一把钥匙,去理解此前一直无法理解、并因此(在五六年之后发现冥想之前,我并未察觉这一点)滋生焦虑的种种处境。我立刻在身边各处看到了这种逃避的实情。这解开了一些焦虑,却没有改变任何本质的东西,因为我只在旁人身上看见那实情,同时(理所当然地)设想它在我身上并不存在,总之我是那验证规则的例外(而且对这个确实了不起的例外,我再没多问一句)。事实上,我对别人、对自己都毫无好奇。这把“钥匙”\textbf{只}在怀着深入之欲的人手里才开锁。在我手里,它变成了驱邪术和做派。\par 是在1974年初,我第一次面对这样一个明证:我生命里的破坏一步步跟着我,不可能只来自别人;我身上有某种东西在吸引它、滋养它、让它延续下去。那是一个谦卑而敞开的时刻,宜于更新。可那次更新当时还停在边缘,而且短暂,因为缺少一种\textbf{深入的工作}。那“\textbf{我身上的某种东西}”还很模糊。我分明看出那是爱的匮乏,但连“做一番工作、更切近地查明我身上的爱在何处、如何匮乏,它如何表现出来,造成了哪些具体后果”等等的念头,都不可能来自我至此认识的任何圈子、任何人,也不可能来自克里希那穆提。(恰恰相反,克氏喜欢强调一切工作的徒劳,他自动把它等同于自我的“急于成为”。)于是,我拿一份借来的“智慧”当全部罗盘,除了耐心等待“爱”像圣灵的恩典一样降临到我身上,我看不出还能做什么。\par 然而,我刚刚在一个浪头的最深处学到的那点朴素的真理,却激起了一个强劲的新能量之浪,其势头可比两年半后推动我第一次投入冥想的那个浪。这股能量当时并没有完全闲置。几个月后,就在我因一场及时降临的事故动弹不得的那时,它推动了一次(写下来的)反思:一生中第一次,我审视了那个曾是我与他人关系之未言明基础的世界观,它来自我的父母,尤其来自我母亲。那时我十分清楚地意识到,这个世界观已经破产,它无法说明人与人之间关系的实情,也无助于我这个人以及我与他人的关系之舒展。这次反思仍带着“克里希那穆提风格”的印记,也带着克氏那条禁忌:任何朝向理解的\textbf{真正的工作}都不许做。可它还是使几个月前诞生、起初模糊而难以捉摸的一种认识变得切实而不可逆转。这种认识,当时世上任何书、任何人都无法给我。\par 这次反思要够得上冥想的品质,最缺的是对我的这个人、对我在我自己眼中的\textbf{形象}投去的目光,而不只是对世界观的目光——也就是对一个公理系统、而我并未真正“有血有肉”地出现在其中的目光。它还缺的,是在\textbf{当下}、在反思进行的那一刻对我自己的目光(这次反思仍停留在真正的工作之下);这目光本会让我同时看出一种借来的文风,以及对这些笔记文学外表某种沾沾自喜,也就是缺乏自发性、缺乏真实。无论它多么不足,对我与他人的关系当下影响也相对有限,这次反思在我看来仍是一个阶段——鉴于起点,这大概是必经的阶段——通向两年后应当发生的更深的更新。正是在那时,我终于发现了冥想——在发现这第一个意想不到的事实时:\textbf{关于我自己这个人,竟还有东西可发现}——那些东西几乎完全决定了我一生的走向,以及我与他人的关系的性质……} (41)。无需再多,就足以让我相信一切都已
% ---------------- p.216 ----------------
到手(当然,嘴上还否认)。我不需要再读;无论口头还是笔下,我都能即兴来一套最纯正的克里希那穆提,一套滴水不漏、环环相扣的话语。可话语再漂亮、再严密,也从来一刻都不像管用,无论对我还是对别人。这样过了好几年,我一点也没做出要从中吸取教训的样子。

随着发现冥想,那套行话一夜之间从我身上脱落,不留痕迹。那时我才懂得话语与认识之间的全部差别。

大老板马上校正了准星:克里希那穆提打入冷宫,冥想捧上台面!自然是悄悄地;如今他得换一套全然不同的手法。时势变了:这小子如今在他腿脚间乱窜,而且有时候眼神挺利。想必小子当时正忙着别处。总之,直到五年之后,某口锅炸开了、小子跑来看出了什么事,大老板的把戏才被识破。

说到底这事过去不算久,才两年多一点:那位“不露相的上师”终于被揭穿——又一副伪装进了冷宫!可怜的老板,他几乎要赤条条地见人了。换个说法:那匹“冥想”马顶替了那匹无名的马(那匹马可千万不能叫“克里希那穆提式”!),如今它带来的押注回报实在少得可笑,尤其比起“数学”马当年那些可观的回报——那时老板还押在它身上。他之所以这么久守着那笔烂注,纯粹是惯性——他已经换过一次注,这本来就不多见,而且当时全靠一记猛击的事件才办到\fnoteP{6}{(42)}{\emph{有益的拔除}\par 所说的那桩“猛击”事件,是1969年底发现:我觉得自己属于的那个机构,有一部分经费来自军事部的款项——这与我基本的公理不相容(而且今天仍然不相容)。这一事件是一整串事件(一个比一个更能说明问题!)中的第一桩,其后果是我离开了高等科学研究所(IHES),并一针引一线地彻底换了圈子、换了投入。\par 在IHES的英雄年代,迪厄多内(Dieudonné)和我是它仅有的成员,也是仅有的为它赢得\textbf{信誉}和听众的人:迪厄多内靠编辑《数学出版物》(Publications Mathématiques)(第一卷早在1959年就出了,即莱昂·莫尚(Léon Motchane)创办IHES的次年);我靠“代数几何讨论班”。最初那些年,IHES的存在极其朝不保夕,经费没有着落(靠几家充当赞助人的公司的慷慨),唯一的场所是巴黎梯也尔基金会(Fondation Thiers)借给的一间屋子(借得明显不情愿),只在我开讨论班的日子用〔IHES最近为建所二十五周年印发的一本小册子(尼科·凯珀(Nico Kuiper)好意寄给我一册)对这些\textbf{艰难}的开端只字不提,也许觉得它们不配去年那场大张旗鼓的庆典的庄严〕。我有点觉得自己和迪厄多内一起,是我所属这个机构的“科学”\textbf{共同创办人},而且我打算在此终老!我最终对IHES产生了强烈的认同,而我的离开(由于同事们的漠然)被体验成一种从另一个“家”被拽离,后来才显出是一场解放。\par 回过头看,我意识到我身上想必早就有一种更新的需要,我说不清从何时起。我离开IHES的前一年,我的精力投入忽然翻了个方向:把前一天还烫手的工作、把我最着迷的问题都撂在一边,转而去搞生物学(是在一位生物学家朋友米尔恰·杜米特雷斯库(Mircea Dumitrescu)的影响下)——这肯定不是单纯的\textbf{巧合}。我当时是以在IHES内部长期投入的架势扑进去的(这与该机构多学科的宗旨相符)。那肯定只是对一种深得多的更新之需的宣泄,而这种更新在IHES那“科学蒸笼”的气氛里无法完成,它是在我前面提到过的那场“觉醒的瀑布”中发生的。前后共有七次,最后一次在1982年。“军费”那段插曲是天意,它触发了这些“觉醒”中的第一次。军事部,和我IHES的老同事们一样,最终都得到了我全部的感激!} (42)。老板们不太喜欢换注——而这一次甚至算是往回走,回到上一笔注。

从1973年我退居乡间起,这匹新马带来的回报,比起从前那匹,就开始少得可怜。三年后冥想的意外降临,又让它们略略活络。甚至还有1979年3月到7月那阵令人头晕的高峰——这一段我在这里不多说——那时我又一次显出使徒的模样,这一回是一种既古老又新鲜的智慧的使徒;这智慧唱在我自己作的一部诗作里,而我最终没有把它交到出版商手上\fnoteP{7}{(43)}{\emph{“我自己作的那部诗作”} 里面有许多我第一手就知道的东西,它们今天在我看来,在我的生命里、以及在一般“生命”里,和它写成时一样重要——当时我是打算发表的。我之所以没发表,主要是因为后来我发觉,它的形式患有一种刻意“作诗”的毛病,以致它的整体构思过于搭建,许多段落缺少自然,有些地方甚至僵硬或浮肿得叫人难受。这种时而浮夸的形式,是我当时心性的映照;在那里,说到底常常是“老板”在领着跳舞——当然是笨重地领……} (43)。可是两年之后,上师彻底歇了业,就好像“冥想”这匹马也断了一条腿(就它带给老板的回报而言)——不管有没有手法,都没法再扮上师了!

% ---------------- p.217 ----------------
\origpage{121} 这以后就没怎么再拖——三条腿的马进了冷宫,连同使徒诗人、“不是上师的上师”,以及“不敢说出自己名字的克里希那穆提”。数学万岁!

下文如何,我们饶有兴致地等着……

% ---------------- p.218 ----------------
% (原书 p.218 为章末空白页)


% ================================================================
%  第 11 章 孤独的历险 (原书 p.219-230, 边码 p.122-p.139)
%  源文: _work/src1/task-ch11.md
% ================================================================
% ---------------- p.219 ----------------
\chapter{孤独的历险}

\section{(46) 禁果}

\origpage{122} 记笔记的事,我不得不停了两天。仔细重读之后,我看先前那套戏路子大致上确实是实情的写照——只是这写照现在还得再往深里挖一挖。我尤其该把冥想与数学这两匹“马”各自的长处圈得更准些;也该试着弄明白:究竟是哪些事件、哪些机缘,最终扣动了老板押注的“翻转”——明明有惰性之力拽着他,要他没完没了地守着一注,哪怕那注已经输了。

也许还该探一探孩子自己的偏爱。如今这点已经说定:他时不时就想换一换局,而老板显然也留着最起码的一点伸缩,不至于硬逼着他总玩这个、从不玩那个。这几年,老板学会了把孩子算进去,学着跟他将就,不必等到锅子炸开。这还不是完全的和谐,但已不再是战争,更像是某种诚挚的协约——偶发的紧张往往让它松动,而不是变硬。

只要不被硬顶得太狠,孩子的偏爱天性上是相当柔软的。(这跟老板可不一样——他到头来学到的那一丁点柔软,全是违着自己的意愿、到了晚年才学会的……)不过孩子柔软,并不等于他没有偏爱,不等于他不会对这一样东西的兴致比对那一样更浓。

要把这里头看清,常常一点也不容易:孩子的欲望、老板的偏爱,乃至老板一劳永逸定下的东西,三者之间很难分清。从前我对自己说:冥想比数学更好、更重要、更严肃,如此等等,理由如此如此(不用问,都是最站得住的理由)——那是老板事后给自己找的好理由,好让自己相信他下的那一注确实“押对了”。孩子从不说这样东西比那样“更好”“更重要”。他不爱讲话。他想做什么,只要没人拦,他就去了,从不问这事“重要”还是“更好”。他的兴致,对这件事对那件事、这一刻下一刻,有强有弱。要察觉他的偏爱,听老板那些解释性的话毫无用处——他自称代表孩子说话,其实只能代表他自己。只有观察孩子在玩什么,或许才勘得破他的偏爱。而且就算这样也不那么容易:他兴冲冲地玩着这一样,并不总意味着,要是老板不伸他那根手指头去拨一拨,他就不会欢天喜地地去玩另一样。

% ---------------- p.220 ----------------
\origpage{123} 显然,比其他一切更吸引他的,是未知——是追进夜里那些含混不清的褶子中去,把未知的东西带到光天化日之下,他自己的未知,也是所有人的未知。我感觉到,我说“也是所有人的”的时候,这确实是孩子的欲望,而不是老板的虚荣——那种想叫满堂看客、也叫自己吃一惊的虚荣。还有一点也是说定了的:孩子每回从阁楼、从取之不竭的地窖的幽暗中带回来的,都是些“显然”的、孩子气的东西。它们越是显得显然,他就越高兴。要是它们不显然,那就是他没把活干到底,在幽暗与天光之间半路上停下了。

\origpage{124} 在数学里,“显然”的东西,也是那种迟早总有人会撞上的东西。它们不是“发明”,可以有人做、也可以没人做。它们本来就在那里,从来就在;人人都在它们边上走来走去而不留意,顶多绕着它们兜个大圈子,或者每次都绊一跤、从它们身上跨过去。过了一年,或者过了一千年,难免有人终于留意到那东西,围着它挖,把它刨出来,翻来覆去地看,擦干净,末了给它起个名字。这类工作——我偏爱的工作——每一回换个人也能做;而且,换个人早晚也非做不可\fnoteP{1}{(44)}{不用说,我在这里把一个假设搁在一边了——至少可以说,这个假设绝非不可能:一场原子战争,或诸如此类的另一种“乐子”,猝然闯了进来,一举而永久地终结了那场叫做“数学”的集体游戏,顺带也终结了许多别的东西……} (44)。

在我自己的发现这件事上就完全不一样了——在这场绝非集体的游戏“冥想”里。我所发现的东西,世上没有任何别的人,今天也好、任何别的时刻也好,能替我发现。它只属于我一个人去发现,也就是说,也属于我一个人去承担。这个未知,并不因为事理之必然、几乎不管我肯不肯费心去理会它,就注定要被知道。它在沉默里等着被知道的那一刻;有时时候到了,我听见它在呼唤——而蒙召去认识它的,只有我一个人,只有我里面的那个孩子。这不是一个“缓期执行”的未知。当然,我尽可以听从它的呼唤,也可以躲开,说“明天”,或者说“总有一天”。但那呼唤是冲着我、也只冲着我来的;除我之外没人听得见,除我之外也没人跟得上。

每当我听从了这呼唤,“企业”里就有什么东西变了,多多少少。效果是即时的,当下就被感到是一种恩惠——有时是一种突如其来的解放,一种巨大的宽释:卸下了一副我常常连自己都没察觉在扛着的重担,而它的真实,恰恰由这宽释、这解放显出来。在振幅小一些的音阶上,这类体验在一切发现的工作里都很常见,我也曾有机会谈到过。然而有一件事,使自我发现的工作(不论它做在光天化日之下,还是留在地下)区别于任何别的发现工作:它真的改变了“企业”本身里的什么东西。这不是量上的改变,不是产量的增加,不是产品在尺寸上、甚至在品质上的差别。这是老板与工人—孩子之间关系的改变。也许连老板自己也变了——如果这话除了就他与工人、与孩子的关系而言,还能有别的意思。比如他也许不那么盯着产量了——但这也仍是他与工人关系的一面,因为一种关切、也许一种敬意出现了,而那从前对他是陌生的。凡我冥想过的场合,变化都朝着老板与工人之间关系变得清朗、变得平静的方向。除了某些冥想停留在表面的情形——那是在某种当下而有限的需要压迫之下做的“应景”冥想——那清朗一直延续到今天,平静也是。

% ---------------- p.221 ----------------
\origpage{125} 这就使自我发现的工作带上了别的发现工作所没有的一层意义,尽管许多本质的方面是共通的。在自我的认识里,在自我发现的工作里,有一个维度,使它们区别于一切别的认识、一切别的工作。也许这就是知识之树上那颗“禁果”。也许冥想在我身上施展的那种魔力——更确切地说,是冥想向我揭示其存在的那些奥秘的魔力——正是禁果的魔力。我跨过了一道门槛,恐惧在那里消失了。认识的唯一障碍是一种惰性,一种有时大得可观、但终归有限的惰性,绝非不可逾越。这惰性,我几乎每走一步都感到它,阴魂不散,无处不在。它有时把我惹得冒火,却从没让我泄气。(数学工作里也一样——那里主要障碍也是它,只是分量轻得没法比。)这惰性成了游戏的基本配料之一;更确切地说,成了这游戏里的一位主角——这游戏精巧,而且全不对称,其中有两位——更确切地说三位主角:一边是纵身一跃的孩子,和那(化作惰性的)老板,他竭尽全力地刹车(一面还声称自己不在场);另一边是那美丽未知者隐约可见的身影,满是神秘,既近又远,一面躲闪,一面召唤……

\section{(47) 孤独的历险}

“冥想”在我身上曾经有极大的魔力——跟从前“女人”的吸引力一样大,而它似乎已经取代了后者。我刚写下“曾经”,这并不表示这魔力今天已经熄灭。这一年我把自己投进数学里,它只是退到了背景中。经验告诉我,这局面一天之内就可能翻过来,正如这局面本身也正是一次全然未料的翻转的结果。其实,在我经历过的四个漫长的冥想时期里(其中一个延续了将近一年半),每一次我都理所当然地以为,我会顺着这股势头一直走到最后一口气,尽可能深入地去探生命之谜、人之存在之谜。当笔记堆成吓人的一摞摞、眼看要把我的工作室淹掉时,我甚至终于请人量身定做了一个柜子来安放它们,而且留得很宽(凭一个快速的算术级数推算),好把往后几年里很快会添进去的也一并安放;要是我没记错,我留了大约十五年的余量(这已经相当可观了!)。这一回老板事情办得漂亮,论管家,这是漂亮的管家功夫!这一手,加上把所有与冥想工作或近或远沾边的私人文件来一次大刀阔斧的整理,是他最后着手、并且(几乎)圆满完成的一件差事——就在偏爱与押注翻转之前。这不禁让人琢磨:他是不是心里早有另一层打算,是不是已经看见一册册《数学反思》塞满了那些空着的格子——那格子据说本是留给日后的“笔记”的。

当然,冥想之爱、自我发现之爱,辽阔得足以填满我这辈子直到最后一天。数学之爱也确实没有被耗尽,但也许这份饥渴在往后几年里终会吃饱。我里面有什么在这样盼着,并且觉得数学是一道羁绊,拦着我走那场只有我一人能走的孤独历险。而且我觉得,我里面这个“什么”不是老板,也不是老板的哪一桩心血来潮(按本性,他是分裂的)。我觉得数学之爱还带着老板的印记;无论如何,跟着它走,只是让我的生命在一个封闭的圈子里打转;一个轻省的圈子,一种惰性的运动,绝不是更新。

% ---------------- p.222 ----------------
\origpage{126} 我曾追问:数学之爱在我生命里这样执拗地持续,究竟是什么意思。跟着它走的时候,它并不真的填满我的生命。它给人快乐,也给人满足,但它本身并不足以给人真正的舒展、一种饱满。像一切纯粹智力的活动一样,强烈而长期的数学活动,效果毋宁说是使人变钝。我在别人身上看到这一点,尤其在我自己身上——每当我重新扑进去的时候。这活动太零碎了,它只动用我们直觉、感受诸能力中极小极小的一部分,那些能力由于久不使用而变钝。长期以来我并不觉察,而显然我的大多数同行,当年也不比我更觉察。我想,正是从我单单做冥想以来,我才留意到这件事。稍一留意,它刺眼得很——大剂量的数学使人愚钝。即便在两年半前的那场冥想之后——那场冥想承认了数学之爱确实是一种爱,是我生命里重要的一件事——如今我献身于这爱时,仍留着一份保留,一份迟疑,不是完全的献出。我知道,所谓“完全的献出”其实会是一种退位,会是顺着一种惰性走,会是一种逃避,而不是献出。

对冥想,我身上没有丝毫这样的保留。我投身其中时,是全身心投入,这献出里没有一点分裂的痕迹。我知道,在献出的时候,我与我自身、与世界完全一致——我忠于我的本性,“我就是道”。这份献出对己对人都是好的。它把我向自己敞开,也向他人敞开,用爱去解开我里面还结着的结。

\origpage{127} 冥想把我向他人敞开,它有力量解开我与他的关系,哪怕对方仍旧结着。可是,要找到机会与人多少谈一谈冥想的工作、谈一谈这工作让我认识的这样那样的事,是极难得的。这绝不是因为那些事“太私人”。打个不贴切的比方:对某个时候让我着迷的数学,我只能跟一位具备必备行头的数学家谈,而且他还得同时乐意对之发生兴趣。有时一连几年,我被某些数学东西迷住,却碰不上(甚至不去找)另一个能与之谈论它们的数学家。但我很清楚:只要我去找,就会找到;即便找不到,那也只是运气或时机的问题;让我感兴趣的东西,不可能不使某个人、甚至某些人感兴趣,是十年后还是一百年后,说到底无关紧要。正是这一点,赋予我的工作以意义,哪怕这工作是在孤独中做的。假如世上再没有别的数学家,而且将来也不会再有,我不相信做数学对我还会有意义——我猜,对任何别的数学家、任何别的“研究者”,情形也一样。这就接上了前面那个观察:对我来说,“数学的未知”是还没有任何人知道的东西——它不取决于我一个人,而取决于一个集体的实在。数学是一场集体的历险,已经延续了几千年。

% ---------------- p.223 ----------------
\origpage{128} 至于冥想,要谈论它,“行头”的问题根本不存在;至少到我这个地步不存在,我怀疑它将来也不会存在。唯一的问题,是别人身上有没有一份兴趣,能与我身上这份兴趣相应。所以那是一种好奇心,对着自己和别人身上真正在发生的事——越过那些必不可少的门面;只要人真的想看看门面底下盖着什么,那些门面就藏不住多少东西。但我已经知道,一个人身上出现这种兴趣的时刻,“真相的时刻”,是稀少而短暂的。当然,常能碰到所谓“对心理学感兴趣”的人,读过弗洛伊德、荣格,还读过许多别人,巴不得来几场“有意思的讨论”。他们随身带着那份行头,或轻或重,就是人们所谓的“文化”。它是他们对自己那幅形象的一部分,又加强那幅形象,而那幅形象他们是从来不肯去查看的——跟另一位对数学、对飞碟或者对钓鱼感兴趣的人一模一样。我刚才想说的,不是这种“行头”,也不是这种“兴趣”——尽管同样的字眼在这里指的是性质不同的东西。

\origpage{129} 换一种说法:冥想是一场孤独的历险。它的本性就是孤独的。不仅冥想的工作是一种孤独的工作——我想,一切发现的工作都是如此,哪怕它嵌在一项集体工作里。但由冥想的工作诞生的认识,是一种“孤独”的认识,一种无法分享、更无法“传达”的认识;或者说,它若能被分享,也只在稀少的瞬间。这是一种工作、一种认识,它与最根深蒂固的共识逆向而行,它让每一个人、让所有人不安。这认识当然表达得简单,用些简单清亮的词。当我向自己表达它时,我在表达中学习,因为表达本身就是工作的一部分,由一股强烈的兴趣托着。可是这些同样简单清亮的词,一旦撞上冷漠或恐惧的紧闭的门,就无力向别人传达任何意思。就连梦的语言——气力完全不同,资源无穷,由一位不知疲倦、心怀善意的“做梦者”不断更新——也穿不过那些门……

没有哪一种冥想不是孤独的。只要有一丝想得到谁认可的操心,想得到证实、得到鼓励,就没有冥想的工作,也没有自我的发现。有人会说,任何真正的发现工作,在工作进行的那一刻,也是如此。不错。但在工作本身之外,他人的认可——无论来自亲近的人、同行,还是自己所属的整个圈子——这认可对于这工作在一个献身者生命中的意义,是重要的。这种认可、这种鼓励,是最有力的激励之一,正是它们使“老板”(还是用这个形象)毫无保留地开绿灯,让孩子尽情尽兴。主要正是它们决定了老板的投入。我自己在数学上的投入,情形也一样——那是被嘉当、施瓦茨、迪厄多内、戈德芒这些人,以及他们之后的另一些人的善意、温暖与信任鼓励起来的。相反,对冥想的工作,丝毫没有这样的激励。这是工人—孩子的一种爱,老板心底里还算厚道,多多少少容着它,因为它“不带来”任何东西。它当然结果,但那不是老板巴望的果子。只要他不在这件事上自己骗自己,事情就很清楚:他不会往冥想里投钱——老板天性合群!

只有孩子,天性上是孤独的。

\section{(48) 馈赠与接纳}

昨天谈到冥想的孤独本性时,有个念头掠过:我写了将近六个星期的这些笔记,末了竟成了某种冥想,可它们却是要拿去发表的。这当然必定以许多方式影响了这冥想的形式,尤其是那份求简洁的用心,还有那份求谨慎的用心。冥想的一个本质方面,即对工作当下我身上发生之事保持持续的留意,只偶尔显现,而且浮在表面。这一切肯定影响了工作的进程和品质。但我仍感到它有冥想的品质,首先凭它果实的性质,凭一种对我自己的认识浮现了出来(在此情形下,尤其是对某一段过去的认识),而那是我此前一直回避的。另一个方面是自发性:近五十个“节”或“段”,反思是自发地聚成这些的,而每一个,我在开头时都说不出它的实质会是什么;每一次,实质都只在途中才显露;每一次,工作都带出一些新的事实,或者以新的光照明了此前被忽略的事实。

% ---------------- p.224 ----------------
\origpage{130} 这工作最直接的意义,是与我自己的对话,因而是一场冥想。然而,这场冥想是要发表的,而且还要充当随后那部《数学反思》的“开场”——这一点绝不是可有可无的处境,绝不是在写作过程中不起作用的东西。对我而言,它是这工作意义的一个本质部分。昨天我暗示老板肯定能在其中“捞到他的好处”(他是“捞好处”的老手,什么都少不了他那一份!),这丝毫不意味着它的意义就缩到那一点上——缩成那匹著名的三脚马迟来的、几乎算死后才发生的“回归”!也不止一次,我感到一个行为深层的意思,有时会超出促成它的那些动机(显的或隐的)。而在这“回归数学”里,我猜度着还有另一层意思,不只是某些心理力量相加的结果——那些力量在某个时刻、出于某些缘由,在我这个人身上碰到了一起。

这场我正在进行的“冥想”,是要献给我在数学世界里认识并爱过的人——如果我感到它是那份隐约见到的意义的重要部分,那并不是因为我在盼着这馈赠被人接纳。它被不被接纳,不取决于我,只取决于它所寄予的那人。它被接纳,我当然绝非无所谓。但那不是我的责任。我唯一的责任,是在我所作的馈赠里是真诚的——也就是说,也是做我自己。

冥想让我认识的,是些卑微而显然的事,一些不登大雅之堂的事。也是那些我在任何书里、任何论著里都找不到的事——不论那书多么博学、多么深刻、多么天才——那些没有人能替我找到的事。我盘问过一团“雾”,我费心去听它,我学到了一点卑微的真相,关于一种“竞技态度”及其显然的意义——在我与数学的关系里,也在我与他人的关系里。就算我把《圣经》、古兰经、奥义书都“照原文”读过,再加上柏拉图、尼采、弗洛伊德和荣格,我是个学问广博精深的天才——这一切也只会让我离那真相更远,那是孩子气的、显然的真相。就算我把基督的话重复一百遍——“像小孩子那样的人有福了,因为天国是他们的”——再细细给它作注,那也不过仍旧使我离我里面的孩子更远,离那些让我不舒服、只有孩子才看得见的卑微真相更远。就是这些东西,是我所能献出的最好的。

而我很清楚:这样的话被说出来、被献出去,用的是简单清亮的词,却并不因此就被人接纳。接纳,并不只是收到一条信息,不管是尴尬地收到,还是甚至感兴趣地收到:“嘿,这谁想得到……!”或者:“说到底也没那么稀奇……”。接纳,常常是在献出者身上认出自己。是借着别人的这个人,与自己相识。

\section{(49) 一条裂缝的确认}

这段关于当前工作意义、关于馈赠与接纳的短小思考,像是反思脉络里的一段离题;或者更像是例证,说明某些方面——它们使“冥想”区别于一切别的发现工作,尤其是数学工作。昨天我明白了:这些方面有双重效果,也就是两种方向相反的效果:对“孩子”是一种独一无二的魔力,对“老板”则是全然的无趣。看来这双重效果就在事物的本性里,绝不可能靠什么妥协、什么安排来缓和。不管怎么做,孩子一跟从他真正的偏爱,老板就在里头捞不到他的好处——一点也捞不到!

% ---------------- p.225 ----------------
\origpage{131} 毫无疑问,已经发生的那次翻转,意义就在这里;它很可能在往后几年里把冥想从我生命中一扫而空(除了像三个月前那样的“应景冥想”)。我并不认为,因此那几年就会完全荒芜,正如过去这一年也并不荒芜。但同样确实的是:我在这一年里学到的东西(数学以外),跟此前四年中任何一年学到的东西相比,都微不足道。怪的是,我经历过的四个漫长冥想时期,每一个都是极其饱满的时光,没有任何迹象让人怀疑我心里还有什么落空的东西。然而,锅子既然炸了,就说明某处有压力,而这压力不会是当天才有的;它一定已经存在,在我看不见的某处,存在了几个星期或几个月,而那时我正专注而全然地沉在冥想里。

\origpage{132} 不过这里我是被笔头的冲劲带着走了(更确切地说,是打字机的冲劲)。实情是:(除了最后一个冥想时期——它正冲到半途,就被一连串事件与情势截断了)冥想的强度从某一刻起渐渐低落,就像一道浪,而后面另一道浪正等着接它的位置……说实话,饱满之感也随着这同一运动,只是有个分别:它只在“冥想浪”的时候在场,不在“数学浪”的时候在场。

我试图圈定的这个局面,在我看来已经不再是冲突的局面,但看来它里面仍旧含着冲突的胚芽、冲突的可能。如今它对我来说,也许是凭它对我生命进程的冲击,我内部分裂的最显眼的迹象。这分裂不是别的,就是老板与孩子的分裂。

我无法把它了结。既然它在这一处表现里已经被勘破,我能做的,只是留意它,在往后的月份和年份里追索它的迹象与演变。也许这份对数学的爱——不得不说,有点不合时宜——会在燃烧中烧尽(像我身上另一份爱已经烧尽那样……),让位给那唯一的爱:对自我的发现,对我命运的发现。

这份爱足够辽阔,我说过,足以填满我的生命——而我的整个生命肯定也不够把它耗尽。

\section{(50) 一段过去的重量}

我把《收获与播种》最后收尾完工,已经是几天前的事了——此前一个多月,我一直以为再过几天就要完工。这一回也是,在“收尾”之后,我仍不完全确定自己是否真的完工了——因为有一个问题我确实悬着没解决。那就是“弄明白:究竟是哪些事件、哪些机缘,最终扣动了‘老板’押注的‘翻转’”,使数学取代了冥想的位子,而且是与巨大的惰性之力相抗。这几天,我的思绪不请自来地、相当执拗地回到这个问题上——而这几天我其实已经开始接上另一些全然不同的事,包括一些数学问题(共形几何方面的)。不如再借一借这股冥想的“末势”,多少往深里挖一挖,把地方腾干净。

% ---------------- p.226 ----------------
当我试着“凭感觉”回答为什么“我又弄起数学来了”(指的是那种大规模的、打算长期干的投入,至少以几年计),好几个联想就冒出来。其中最有力的,也许要数我近六七年来在教学活动中终于感到的那种慢性落空之感。就是那种越来越强的感觉:我“大材小用”;而且常常是,我把力气、把我最好的东西投给一些闷闷不乐的学生,而他们对我能给的东西毫不在意。

\origpage{133} 我到处看见值得做的妙事,而它们只等着人去做。往往,一点可怜的行头就足够上手;倒是这些东西自己告诉我们:该发展出什么样的语言来圈住它们,该置办什么样的工具去挖它们。我没法不看见它们——只要由于教学活动而与数学保持着经常的接触(不论那层次多么低),哪怕在我对数学兴趣最边缘的那些时期,也是这样。每一件瞥见的东西背后,只要稍一发掘,又有别的好东西,遮住一些、又露出一些……不论是在数学里还是在别处,只要带着真正的兴趣把目光放上去,就看见一种富饶显露出来,一种深度展开,叫人猜出它们取之不竭。我说的那份落空,就是怎么也没法把这种富饶之感——深之感——哪怕一星半点传给学生:哪怕只是一点想至少把他们手边的东西摸个遍的兴致,在他们无论如何都打算投进所谓“研究”活动的那几个月或几年里,尽情尽兴地干一场——反正也是为了拿这个那个文凭。十年来我教过的学生里,除了两三个,看起来“尽情尽兴”这个念头本身就让他们害怕;他们宁愿几个月、几年地垂着手原地蹭,或者苦哈哈地干一份鼹鼠式的活,那活的来龙去脉他们一概不知,只要末尾有个文凭就行。关于这种创造力的瘫痪,有很多可说的,它与“天赋”或“才能”的有无毫无关系——这又接上了我反思的最开头,那里我曾顺带触及这类阻塞的深层原因。但这不是我在这里要说的话;我这里要说的,毋宁是确认一种慢性的落空状态——这些情形,在最近七年教学活动中不断重演,终于在我身上造成了它。

\origpage{134} “解决”这样一份落空,最现成的办法——至少就它是我身上那位“数学家”的落空、而不是那位教师的落空而言——就是我自己去动手做那些事,至少做其中一部分;那些事我曾绝望地盼着哪个学生终于有一天会接过去做。这也正是我多少在这里那里做过的事:或者借教学活动之余、借教学活动的机会,花几个小时、甚至几天作一次偶然的思考;或者在数学上大发作的时期(那有时来得像真正的爆炸……),一发作就是几个星期、几个月。这样偶发的、一阵一阵的工作,往往只能对一个问题作最初的一遍粗加工,只能给出最零碎的视野——不如说,那是对尚待进行的工作看得更清楚一些;而那工作本身始终搁在那里等着做,而且正因为看得更清,就显得更烫手。两个月前,我给出了一份总体草图,涉及那些我才刚刚开始约略摸到门径的主要题目。那就是《纲领概要》(Esquisse d'un Programme)——我此前已有机会提到过它,它最终将附在这份反思之后,二者合起来构成《数学反思》的第一卷。

相当清楚:单靠这份(不妨说“私下的”)勘察工作,不足以化解我的落空。“大材小用”这个感觉,无疑翻译出一种想有所行动的欲望(我想,它源于自我,也就是说,“老板”的欲望)。这里说的行动,与其说是对他人的行动(比如对我的学生:推动他们,向他们“传达点什么”,或者帮他们拿到某个文凭,好去申请某个职位,等等),不如说是“数学家的”行动:为发现某些无人料到的事实、为某个理论的绽出出一份力,等等……这立刻接上了前面那个观察:数学是一场“集体的历险”。

% ---------------- p.227 ----------------
\origpage{135} 如果我追问自己近十年来做数学时的心境——那是我生命中一段根本不会想到有朝一日还能重新发表东西的时期,而且当时也多少清楚:我眼下或未来的学生,没有一个会用到我的勘察工作——我立刻就看到:那绝非某个只为个人乐趣做事、或者被一种只关乎自己而与他人无关的内在需要所驱赶的人的心境。我做数学的时候,我相信我心里某处是明白的:这些数学是做来传达给别人的,是要成为一件更广大之物的一部分的,我为它出力,而那东西绝不是个体性的。这“东西”,我可以称之为“数学”,或者更好的说法,“我们对数学事物的认识”。这里“我们”一词,首先无疑具体地指我所认识、与我共有某些关切的那些数学家构成的团体;但毫无疑问,它也超出这个狭小的团体,正如它超出我个人。这个“我们”指的是我们的族类——就这一族类通过它在各个时代的某些成员,曾经并且正在关切数学对象世界的诸般实在而言。直到此刻我写下这几行之前,我从未想到过我生命中存在着这“东西”,更没想到要追问它的本性,追问它在我作为数学家与教师的生活里起什么作用。

我提到的那个想有所行动的欲望,在我身上、在我作为数学家的生活里,似乎采取这样的形式:把所有人都不认识的东西从暗处取出来——不只是我不认识的东西(如前所见)——而且,还要把它交给所有人支配,也就是去丰富一份共同的“家产”。换句话说,这是一种欲望:为那超出我个人的“东西”、那份“家产”的扩大与丰厚出一份力。

这份欲望里,当然不缺借自己的作品来扩大自己的欲望。从这一面看,我又碰上了那种对“成长”、对扩大的饥饿——那是自我、是“老板”的特征之一;那是他侵占性的一面,推到极处便是毁灭性的一面(参见注44,第13.1.1节,原书p.260)。不过我同时也明白:想增加那些(或久或暂)多少会挂上我名字的东西的数目,这欲望远不足以穷尽、远不足以覆盖那股更广大的欲望或力量——它推着我去为扩大一份共同家产出力。我觉得这样一份欲望本来可以得到满足(如果不是“在我的企业里”——那里老板还相当有侵占性——至少在一个更成熟的数学家身上),而他自己这个人的角色仍是无名的。这也许就是自我扩大倾向的一种“升华”形式:通过与一件超出自身之物的认同。除非这类力量本身并非源于自我,而是属于一种更细腻、更深的本性;除非它表达的是一种深刻的需要,独立于一切制约,见证着一个人生命与整个族类生命之间那深深的联结——那联结是我们个体存在意义的一部分。我不知道,而追究这样一些如此宽泛的问题,也不是我在这里要说的话。

我要说的,毋宁是(以一种更朴素的角度)考察一个与我这个人有关的具体局面:一个落空的局面,靠零星断续的数学活动得到局部的、暂时的排遣。这样,这局面的逻辑迟早该把我引向把我找到的东西传达出去。既然直到去年,我根本不肯为我的数学之爱付出那种大规模的、长期的投入——那本是必要的:得靠一番因事制宜的“逐件加工”,才能把我掘出的矿藏“开采”出来拿去发表——那么剩下的另一条路,就是把我至少最放在心上的那些东西,传达给几位足够“在行”的数学家朋友。

我想,要是在最近十年里我碰上一位数学家朋友,对我扮演对话者和消息来源的角色(就像五六十年代那许多年里,塞尔在很大程度所扮演的那样),同时又充当一道中转,把我能传给他的“消息”传出去(这角色从前塞尔不必扮演,因为我自己就担了!)——那么我想“在数学上有所行动”的欲望,就会得到足够的满足,足以化解我的落空,而我只消在数学上断断续续、有节制地投些精力,把最大的份额留给我新的爱。我第一次带着这样的期待(至少心里暗含这期待)去找一位数学家朋友,是1975年,最后一次是1982年,一年半以前。有趣的巧合:两次都是为了想办法“安置”(好让它传开去,谁知道呢,说不定最终还会有人发展它!)同一个“纲领”,同调与同伦代数方面的,其最初的胚芽要上溯到五十年代,而早在六十年代末以前就已经完全“成熟”(依照我内心的确信);这纲领的一个初步的、粗线条的展开,恰恰就是《追寻栈》的主题——而我此刻按说正该给它写导论!不管怎样,由于大概各不相同的缘由,我那些想重新找回一种“特权对话者”关系的尝试,都夭折了——那种关系从前(1970年以前)曾与塞尔有过,后来又与德利涅有过。不过有一个共同的情形:我准备给数学的时间相对有限。在我说的那两次(1975年和1982年),这肯定也促成了交流的瘸腿。其实,我主要想的是把什么东西“安置”出去,并不太费心去做那件必要的事——“(重新)跟上行情”,以便我这方面成为一个让通信者满意的对话者;而对方在行内那些通行技术上比我“在行”得多(至少可以这么说!)。

% ---------------- p.228 ----------------
\origpage{136} 我大可以把《致……的信》——它充当《追寻栈》的第一章,去年二月写的信,才一年多以前——看作我最后一次尝试:在我一位旧日朋友那里,为我如今某些想法与关切找到一点回响。从这封信开始(更确切地说,重新开始)的反思,后来竟要成为1970年以来第一篇注定要发表的数学文本——而这一点,我当时好几个星期都没有察觉。差不多一年之后,我才收到对这封内容充实的信的间接反应(参见注\fnoteP{2}{(38)}{这些笔记其实是那封写给……的长信的续篇,后来那封信成了第一章。它们打字打出来,是为了让这位旧日的朋友读得下去,另外还为了两三个人(尤其是龙尼·布朗[Ronnie Brown]),我以为他们会有兴趣。这封信其实从未得到回音,收信人也没读过它:将近一年之后(我问起他是否收到了),他真心感到惊讶——我居然会有一刻以为他可能读它,看看人们该指望从我这儿得到的是什么样的数学……} (38))。这反应比迄今任何一位数学同行写给我的信都更能说明问题,让我感到某些针对我这点微末之身的姿态——自我离开那个曾与他们共处的圈子以来,这些姿态在我的数学家朋友中间已经成了常事。这封信出自一位我曾当作朋友、怀着热忱的同情去信的人,而信里有一种蓄意的嘲弄,格外猛烈地让我想起一件事——那是我最近几年终于越来越清楚地意识到的事。此前,我主要有机会注意到的是对我这个人本身的疏远,在数学“大世界”里,首先是在那些曾与我或远或近算朋友的人中间(45)。那里说的不再是人际层面上的疏远,而是一种共识,其性质如同一种时尚,并且像时尚那样把自己摆成理所当然的东西,流行于一些多少“在行”的人之间:那种论千页一包一包的数学,还有那些我拿它们在一二十年里磨破了人耳朵的概念(46,47),全面权衡起来并不很严肃;那里头有大量装腔作势的轰响,值得的东西却没多少;而除开围绕概形与平展上同调这两个概念的那几大篇“普遍胡话”——它们有时候倒确实自有其用处;唉,这一点人还是乐意承认的——至少剩下的部分,忘掉更厚道;至于那些不顾雅趣、不顾显而易见的严肃准则,仍旧摆出要吹这种格罗滕迪克式小号的人,就该跟他们那位师父(不管认不认)扔进同一条口袋;他们要是受到应得的对待,也只能怪自己……

\origpage{137} 无疑,自1976年以来(50)、尤其是最近两三年来传到我这里的许多这一路的声音(我刚把它们“照原样”抄了下来),终于唤醒了我身上那根好斗的筋——最近十年里它多少睡过去了。它们像反射一样挑起了一个念头:投入混战,让那些什么也没弄懂的毛头小子闭上嘴——总之是个十足愚蠢的反射,就像公牛:只要给它看一块红布,在它鼻子前晃一晃,它立刻就抖擞起来、冲上去,忘了自己原本正安安静静走着的路,那可是它自己的路!不过我仍觉得这反射相当表皮,单靠它不足以让我动摇。再说,幸好,做数学显然比朝一块红布冲过去、四面挨刀要有滋味得多。可是做数学,而且不顾一切地坚持一种工作风格、一种看事情的路数——那些是我的——这也有点“投入混战”的意思;这是在那些轻蔑与拒斥的迹象面前挺立自己——那些迹象落到我身上,毫无疑问是回应我的旧友们在我身上感到、或者以为感到的轻蔑,那轻蔑若不是针对他们,至少也是针对他们毫无保留地认同的那个圈子。所以这多多少少也是在跟着那块红布跑,而不是走自己的路。

% ---------------- p.229 ----------------
\origpage{138} 这个念头,最近几周里已多次在我这里出现过;而今天这场反思,也许主要就是朝考察这一面去的。走着走着,又显出另一面,那里自我的力量无疑也占了大份,但它不像是一种单纯的好斗反射。不如说是一种我身上的欲望——此刻我还看不清它的本性——要给我最近十年或十二年所做的数学工作一个意义,或者要看到它取得它全部的意义;而这意义(我内心确信)不可能归结为私下里的乐趣或个人的历险。可是,即便这欲望的本性仍未被理解,既然我还没闲下来细看它,这场反思已足以向我表明:真正压在我身上、可以说硬掰着我的手让我往数学投入的那股力量,恰恰就在这欲望里——“翻转”之力。红布有没有,它照样起作用。如果它是依恋一段过去的迹象,那也是最近这十年的过去,也就是“1970年以后”的过去,而不是那些已经白纸黑字写下的东西、已经做成的东西的过去,不是1970年以前的。

\origpage{139} 骨子里,对这些东西、对将来“后世”会怎样处置它们(何况是否真有后世还很可疑……),我心里毫无不安。我在这段过去里感兴趣的,绝不是我在其中做过的事(以及它已有或将会有的运道),而恰恰是没有做成的事——在那份当年摆在我眼前的大纲领里,只有极小一部分,靠我的努力,以及靠那些有时愿意加入我的友人和学生的努力,得以实现。这份纲领本身,在我对数学事物的眼光与路数更新的同时,也未经预谋、未经寻求地更新了。年复一年,重心移动了,题目如此,我要说的东西本身也如此:如今我最首要的意图,不再是完成那些精细基础的大工程,而是去探那些最让我着迷的奥秘,比如“Motif”的奥秘,或者 $\mathbb{Q}$ 在 $\mathbb{Q}$ 上的伽罗瓦群的“几何”描述的奥秘。当然,一路走来,我至少免不了在这里那里勾画些基础,就像我已经开始做的那样(除了别的)在《穿越伽罗瓦理论的长征》里,或者像我现在正在《追寻栈》里做的那样。然而要说的东西变了,表达它的风格也变了。

% ---------------- p.230 ----------------
换个说法:最近十年里,在数学事物的世界里,我瞥见了一些神秘的、极美的东西。这些东西不是我私人的,它们生来就是要被传达的——瞥见它们这件事本身的意义,我这样感到,就在于把它们传达出去,好让人接过去、理解、吸收……然而传达它们,哪怕是向自己传达,也就是深入它们,多少发展它们——那是一件工作。我当然很清楚:我绝不会把这工作做到头,即便再给我一百年去投入。但这不该是我今天操心的事——在我余下活着、余下发现世界的时间里,我要拿多少年、多少个月投到那件工作上;而另一件工作正等着我,那是我唯一能做的。我生命的季节,不在我的权柄之内,也不是我的角色去安排的。

% ================================================================
%  第 12 章 第一部分的注 (原书 p.231-241, 边码 p.141-p.152)
%  源文: _work/src1/task-ch12a.md (12.1–12.22, 注 1–22)
% ================================================================
% ---------------- p.231 ----------------
% 原书 p.231-233 的章内目录(Contents)按体例不排印
\chapter{第一部分的注}

% ---------------- p.232 ----------------
% (章内目录续, 不排印)

% ---------------- p.233 ----------------
\section{$\varnothing$}

\notehead{1}{}
\origpage{141}(1984年3月补写)说我的工作“风格”与“方法”没有变,恐怕失于武断——毕竟我在数学上的表达风格已经深刻改变了。这一年来用于《追寻栈》的大部分时间,是坐在打字机前敲下一些思考;它们就是要几乎原样发表的(只除了后来补加的若干相当简短的注,借互见、勘误之类方便阅读)。没有剪刀,没有浆糊,不必费劲去炮制一份“定稿”(那份定稿尤其不该让任何通向它的过程透出影来)——这毕竟也是“风格”与“方法”的变化!除非把数学工作本身与写作、与陈述结果的工作割裂开来,而那是人为的,因为那不合事情的实情:数学工作与写作是拆不开的。

% ---------------- p.234 ----------------
\section{$\varnothing$}

\notehead{2}{}
\origpage{142}(1984年3月补写)重读最后这两段,我有些不适之感,因为写它们的时候,我牵连的是别人,而不是我自己。显然,写的时候,“我自己也可能与此有关”这个念头并没有碰到我。我肯定什么也没学到:我就这样满足于把多年来在别人身上察知、并以许多方式看着被证实的东西,白纸黑字写下来(无疑还带着某种满意)。在这份反思的后续部分,我不得不回想起:我的一生中,对别人的轻蔑姿态并不缺少。我所把握到的他人之轻蔑与对己之轻蔑之间的关联,若在我这个人身上竟然缺席,那才奇怪;健全的理性(以及类似情形下我对自身视而不见的经验——我终究还是察觉到了)告诉我:肯定不该是这样!然而眼下这只是一条纯粹的推论,它唯一可能的用处,是促使我亲眼看一看究竟发生了什么,去看、去审视(如果它确实存在过、或还存在着)这份仍属假设的对己之轻蔑——它埋得如此之深,至今完全逃过了我的目光。真的,可看的东西并不缺!

\pmark 这一件忽然在我看来是最要紧的几件事之一,恰恰因为它藏得如此之深……\fnote{1}{}{(1984年8月)不过关于这一点,参见“屠杀”一注(n\textsuperscript{o}~87)末两段的反思。}

\section{$\varnothing$}

\notehead{3}{}
我这里想到的尤其是那几条“已故”的猜想——莫德尔(Mordell)的、泰特(Tate)的、沙法列维奇(Chafarévitch)的:去年,这三条猜想在法尔廷斯(Faltings)一篇四十页的手稿里同时得到了证明,而就在那个时候,“圈内人”那套早已稳固的共识还裁定这些猜想“够不着”呢!恰好,我所珍爱的“远阿贝尔代数几何”(géométrie algébrique anabélienne)纲领的那条拱顶石式的基本猜想,就与莫德尔猜想相近。(据说莫德尔猜想甚至是这条基本猜想的推论——这就满可以说明:这个纲领并不是给严肃人预备的故事……)

\section{$\varnothing$}

\notehead{4}{}
再说,即便在今天,也还能碰到地位不明的“证明”。格劳尔特(Grauert)那条以他命名的有限性定理的证明,多年来就是如此:谁(而热心人并不缺!)也读不下去。这种困惑,后来由另一些更透亮的证明解决了;其中有些走得更远,接替了最初那个证明。类似而更极端的情形,是所谓“四色问题”的“解答”:其计算部分是用计算机(外加几百万美元)打发掉的。于是这里所说的“证明”,其根基已不再是理解一种数学处境所生的内心确信,而是对一台没有理解能力的机器的信任——而使用它的数学家,对这台机器的结构与运作一无所知。即便假定这计算由别的计算机、按别的计算程序复核过,我也并不认为四色问题就此了结。它只是换了一副面孔:人们不再去寻一个反例,而只去寻一个证明(当然,得是读得懂的证明!)。

% ---------------- p.235 ----------------
\section{$\varnothing$}

\notehead{5}{}
\origpage{143}\origpage{144}这件事之所以格外值得注意,是因为直到1957年,布尔巴基(Bourbaki)小组里还不止一位成员对我有所保留——而我相信,他们终究把我收编进来时,多少是有些勉强的。一句善意的玩笑话,把我归入“危险的专门家”(指泛函分析方面的)。有时我在嘉当(Cartan)身上感到一种未说出口的、更认真的保留——有那么几年,我一定给了他这样的印象:一个倾向于无端而肤浅地推广的人。我见过他大为惊讶:在我为布尔巴基写的头一篇(也是唯一一篇)稍长的稿子(论流形上的微分体系)里,他发现了一点还算有点内容的思考——当初我提议由我来写时,他并不怎么热心。(这份思考几年后对我又有用:我从凝聚对偶的观点发展留数形式体系时,用到了它。)再说,在布尔巴基大会上,我多半是跟不上的,尤其是在共同宣读稿子的时候:那些宣读与讨论的节奏,我实在跟不下来。也许我真的不适合集体工作。不管怎样,我这种难以融入共同工作的困难,或者我因别的缘故在嘉当等人那里引起的保留,从未在任何时刻给我招来讥讽或冷遇,哪怕一丝居高临下的影子也没有——顶多是韦伊(Weil)那里有过一两次(此人确乎另当别论!)。嘉当对我的善意始终如一,不曾稍减,那里面既有热诚,也有他那特有的、在我看来与他这个人分不开的一点幽默。

\section{《生存与生活》的朋友们}

\notehead{6}{}
在这些朋友里,我大概还该算上皮埃尔·萨米埃尔(Pierre Samuel)——从前我主要是在布尔巴基里认识他的,就像谢瓦莱(Chevalley)一样;他(和谢瓦莱一样)在“生存与生活”(Survivre et Vivre)小组里起了重要作用。我看萨米埃尔并不怎么热衷于那种科学家的优越感的幻觉。我感到,他带来的东西很多,主要靠的是他在共同工作、讨论、与人交往中那份健全的常识和笑呵呵的好脾气,也靠他能在这样一个倾向于激进分析与抉择的群体里,优雅地扮演“讨厌的改良派”的角色。我退出“生存与生活”之后,他还在那里待了一阵,担任同名刊物的主编;当他感到自己在这个群体里的存在已不再有用时,他就体面地离开了(去加入“地球之友”)。

萨米埃尔和我同出那个狭小的圈子;这并不妨碍他属于那些沸腾年月里让我觉得学到了一点东西的朋友之列(尽管我是个多坏的学生……)。他那种为人处世的方式,以及谢瓦莱的方式——两者其实毫不相似——对我那些“精英至上”的倾向来说,是比最犀利的分析更好的解药!

如今在我看来,对那个时期所有让我学到了一点东西的朋友来说,我学到的东西,更多是靠他们各自的为人、靠他们与我不同的感受力(其中“某种东西”终于传达了过来),而不是靠解释、讨论之类。说到这一点,我尤其记得——除了谢瓦莱和萨米埃尔——德尼·盖吉(Denis Guedj)(他在“生存与生活”小组里很有威望)、达尼埃尔·西博尼(Daniel Sibony)(他始终与这个小组保持距离,一边从角落里用半是轻蔑、半是打趣的眼光看着它演变)、戈登·爱德华兹(Gordon Edwards)(他是1970年6月蒙特利尔那场“运动”诞生的共同推动者之一,多年来施展出惊人的精力,维持着《生存与生活》刊物的英文“美洲版”)、让·德洛尔(Jean Delord)(一位年纪与我相仿的物理学家,细致而热忱,他对我连同“生存派”这个小圈子都怀有偏爱)、弗雷德·斯内尔(Fred Snell)(另一位在美国布法罗的物理学家,1972年我在那里住了几个月,做客他的乡间住所)。

% ---------------- p.236 ----------------
\origpage{145}这些朋友当中,五位是数学家,两位是物理学家,而全都是科学家——这似乎表明:那些年里离我最近的环境,仍然是科学家的环境,尤其是数学家的环境。

\section{$\varnothing$}

\notehead{7}{}
上一段是整个导论里第一段在我原稿上被狠狠划掉、并添了许多涂改的。对那件事的描述、对词句的选择,起初都是逆着毛捋、顶着潮流来的——显然有一股力量在推着,要把那件事快快带过,仿佛是了却一桩心事,好“去谈正经事”。这些都是一道抵抗的熟面孔;在这里,它抵抗的是把这段插曲弄清楚,以及弄清它作为内心态度之显影的分量。情形与这篇导论开头(第2段)所描述的完全相似:在数学工作中发现一个矛盾及其意义的那个“关键”时刻——那时充当“抵抗”的,是心智的惰性,是它不肯离开一个错误或不充分的视野(而在那里面,我们这个人丝毫未被牵连)。这种抵抗是主动的,必要时还会发明创造,好把一条鱼在没有水的地方也淹死;而我方才说的惰性,只是一种被动的力。在眼前这桩事上,比在数学工作中尤甚:刚刚以其全部单纯、全部显明浮现出来的那份发现,随即带来的是一种卸下重负的轻松感,一种解放感。这不只是一种感觉——不如说,是对刚刚发生之事的敏锐而心怀感激的知觉,而那件事,就是一次解放。

\section{$\varnothing$}

\notehead{8}{}
后文会看清:这份模棱两可,丝毫没有“在1970年觉醒的次日消散”。这里有一招“自我”典型的战略性撤退:把“觉醒之前”那段时间丢给盈亏自负,而“觉醒”立刻成了一道分界线,划出一个无可指摘的“之后”!

\section{$\varnothing$}

\notehead{9}{}
这不完全确切:在我最亲近的同行里至少有一个例外,后面会看出来。这里有一种记忆特有的“懒惰”:它往往倾向于把那些与一套熟悉且根深蒂固的看法“对不上”的事实“蒙混过去”。

\section{$\varnothing$}

\notehead{10}{}
比如,我写给同事或旧日的学生、那些我当作朋友的人的信——数学问题也好,实际事务或私事也好——多得我已数不清,却从未收到过回音。看来这并不只是专门留给我一个人的“优待”,而是一种风气变化的征兆——我听到的回声也是这个方向。(诚然,那些回声涉及的情形是:寄数学信的人,收信的那位知名数学家并不认识……)

% ---------------- p.237 ----------------
\section{阿尔多·安德烈奥蒂、约内尔·布库尔}

\notehead{11}{}
\origpage{146}\origpage{147}当然。我这边也可能有遗忘——何况我那时格外“硬派”(polar)的性情,恐怕也不会让人乐意跟我谈这类事情,也不会让我记住一场很可能发生过的、这个方向的谈话。可以肯定的是:关于恐惧的问题被提起(甚至不用这个名字来称呼它……),至少该是极其罕见的;今天大概同样罕见,尤其是在“上流社会”里。

在那个圈子里我的众多朋友中,除了谢瓦莱(他至少在六十年代该已察觉到这种恐惧的气氛),我觉得真正清楚地感受到了它的,只有阿尔多·安德烈奥蒂(Aldo Andreotti)。我认识他,认识他妻子芭芭拉(Barbara)和他们那对双胞胎孩子(那时还很小),是在1955年(我想是在芝加哥韦伊家的一次晚会上)。我们一直十分亲近,直到1970年那个“大转折”:我离开了从前属于我们的那个圈子,与他们也就有些失去了联系。阿尔多感受力极敏锐,而这份敏锐丝毫没有因为与数学打交道、与我这样的“硬派”打交道而变钝。他身上有一种对接近他的人的、自然而然的同感之才。这使他与我在这数学圈里、乃至圈外认识的所有别的朋友都不同。在他那里,友情总是压过共同的数学利益(那些利益并不缺),他是极少数与我多少谈过我的一生、他也谈过他一生的数学家之一。他父亲和我父亲一样是犹太人,他因此在墨索里尼的意大利吃过苦头,正如我在希特勒的德国。我见他总是愿意鼓励和支持年轻研究者,而那时的气候,已经越来越难让这些人被建制接纳。他那种自发的关切,总先落在人身上,而不是落在什么数学“潜力”或名声上。他是我有幸遇见的最令人眷爱的人之一。

这番对阿尔多的追忆,又勾起了约内尔·布库尔(Ionel Bucur)的记忆——他也是骤然辞世,未及天年;和阿尔多一样,人们(我想)更怀念的是他这个人、这位乐于重逢的朋友,而不是作为数学讨论伙伴的他。在他身上,人感到一种善良,旁边还有一份少见的谦逊,一种不断让自己退到后面的倾向。这样一个人——如此不把自己当回事、如此不想让谁敬畏——最后竟当上了布加勒斯特大学理学院院长,真是个谜;多半是因为他压根没想到要推辞那些他远不觊觎的职务,而他的同事或政治当局偏要把它们放到他肩上——顺带说,那肩膀是结实的。他是农民的儿子(在一个“阶级标准”要紧的国家里,这一定起了作用),也带着农民那份健全的常识与朴素。他肯定察觉得到围绕着名望之人的那份恐惧,但肯定也觉得这事理所当然,是权力位置的天然属性。不过我并不认为他本人曾让谁感到恐惧——无论是他妻子弗洛丽卡(Florica)、他们的女儿亚历山德拉(Alexandra),还是他的同事或学生;我能听到的回声也确实如此。

\section{$\varnothing$}

\notehead{12}{}
这里的“次日”(lendemain)一词要按字面理解,不是隐喻。

% ---------------- p.238 ----------------
\section{$\varnothing$}

\notehead{13}{}
\origpage{148}很清楚,前面的描述没有别的野心,只是想用具体的词语尽力把记忆的那片“迷雾”所交给我的东西还原出来;这片迷雾从未凝结成任何一个稍为具体的个例,好让我在这里给出某种稍微“写实”或“客观”的描述。若要从这段话里读出:那些不愿坐在前排的同行,或者那些没有明星或泰斗身份的同行,在跟后一类人说话时必定被焦虑缠住——那就是曲解我的意思了。在我于这个圈子里认识的多数朋友身上,显然不是这样,即便是那些偶尔出入讨论会或讨论班的人。毫无保留地成立的是:“泰斗”身份造成一道屏障、一道沟,把它与没有同等身份的人隔开;而这道沟极少消失,哪怕只有一场讨论的工夫。我还要补充:“前排”与“沼泽”之间那条主观的界线(我却觉得它十分真实),绝不能归结为社会学标准(社会地位、职位、头衔等等),甚至也不能归结为“身份”或名声的标准;它还反映着气质上、或者某些更难把握的禀性上的心理特点。我二十岁初到巴黎时,我知道自己是个数学家,我做过数学;尽管有人地生疏之感(我有机会谈到过),我在心底里觉得自己“是他们中的一个”——而只有我知道这一点,甚至一开始我都不确定自己会不会继续做数学。今天,我倒更愿意坐到最后一排去(在难得碰上这个问题的场合)。

\section{$\varnothing$}

\notehead{14}{}
人们可能以为这与“没有首领”的说法相矛盾,其实不然。对布尔巴基的老人们来说,我觉得韦伊被看作小组的灵魂,却从来不是“首领”。他在场、并且他乐意的时候,就成了我所说的“串场人”,但他并不发号施令。他心情不好的时候,会把自己厌恶的某个题目卡住,让别人到下一次大会、趁韦伊不在时再安安静静地重拾这个题目,甚至第二天等他不捣乱了再谈。决议由在场的成员一致通过,因为他们认为:一个人逆着所有其他人的一致意见而站在对的一边,这绝非不可能(甚至并不少见)。这条原则对集体工作来说也许显得荒谬。了不起的是:它居然行得通!

\section{$\varnothing$}

\notehead{15}{}
我并不觉得这种对布尔巴基风格的“过敏”,曾在这些数学家与我、或与布尔巴基其他成员和同情者之间造成交流上的困难——倘若小组的精神真是小圈子精神、是精英中的精英的精神,那倒会如此。在风格与时尚之上,小组的所有成员都有一种对数学实质的敏锐感觉,不论那实质来自何处。只是到了六十年代,我才记得我某位朋友把某些他不感兴趣的数学家称作“讨厌鬼”。这类事我另外几乎一无所知,于是往往把这类评价照单全收,被那份满不在乎的自信震住了——直到有一天我发现,某位“讨厌鬼”是个独到而深刻的心灵,只是没有讨得我那位出色的朋友欢心罢了。我觉得在布尔巴基的某些成员身上,面对别人的工作、在自己不懂或懂得不透的时候所持的谦逊态度(或至少是保留态度),是先被磨掉的;而那种“数学本能”却还留着——那种不必参照名声与声望、就能嗅出丰厚实质或扎实工作的本能。

% ---------------- p.239 ----------------
\origpage{149}从各处传来的回声看,我觉得这两样——谦逊与本能——在从前的那个数学圈子里,今天都已成了稀罕之物。

\section{$\varnothing$}

\notehead{16}{}
说实话,布尔巴基的成员里有好几位肯定有“自己的”小宇宙,大小不一,在布尔巴基小宇宙之外或之上。但就我自己的情形而言,这样一个“自己的宇宙”之所以要等到我不再是布尔巴基的一员、全部精力都投到属于我个人的工作上之后,才在我周围形成——这也许并非偶然。

\section{$\varnothing$}

\notehead{17}{}
主要是在科学圈之外,我才遇到对我所投入的那场行动的热烈回响,以及切实的帮助。除了阿兰·拉斯库(Alain Lascoux)与罗歇·戈德芒(Roger Godement)的友情支持,这里我还尤其要提到让·迪厄多内(Jean Dieudonné)的支持:他特地跑到蒙彼利埃的轻罪法庭去出庭,为一件注定失败的事业,把他那份热忱的证词加在别的证词之上。

\section{$\varnothing$}

\notehead{18}{}
我想,这份缺乏辨别力,在这两个场合并非出于我的疏忽,而是出于不成熟、出于无知。直到十来年后,我才开始注意那些阻滞机制——在我自己身上,也在我的亲人或学生身上——并掂量它们在每个人生活里所起的巨大作用,而且不只是在中学或大学里。当然,我后悔这两个场合没有更成熟一些的辨别力;但我并不后悔把自己的印象清楚地说出来,不论那些印象有没有根据。当我在某个场合看到一件做得不认真的工作,把这事照它的本来面目说出来,我觉得是必要而有益的。若在另一个场合,我得出的结论是仓促的、没有根据的,那么负有责任的也并非只有我一个。被这样摇醒的学生,当时还有得选:要么从中吸取教训(第一次也许正是这样),要么就此泄气,也许还改了行(那也未必是坏事!)。

\section{耶稣与十二使徒}

\notehead{19}{}
从1970年至今,还有一个学生伊夫·拉德加耶里(Yves Ladegaillerie)在我这里准备并通过了论文。第一时期的学生是:P.贝尔特洛(P.~Berthelot)、M.德马祖尔(M.~Demazure)、J.吉罗(J.~Giraud)、M.阿基姆夫人(Mme~M.~Hakim)、黄春生夫人(Mme~Hoang Xuan Sinh)、L.伊利西(L.~Illusie)、P.茹阿诺卢(P.~Jouanolou)、M.雷诺(M.~Raynaud)、M.雷诺夫人(Mme~M.~Raynaud)、N.萨阿韦德拉(N.~Saavedra)、J.-L.维尔迪耶(J.-L.~Verdier)。(其中六位是在1970年之后才完成论文工作的,也就是说,在我数学上的可用时间极为有限的时期。)在这些学生当中,米歇尔·雷诺(Michel Raynaud)占一个特殊位置:他论文工作所处理的那些本质性问题与概念,是他自己找到的,而且他是完全独立地发展它们的;我所谓“导师”的角色,因此仅限于读已完成的论文、组成评审委员会,并成为其中一员。

轮到我出题目的时候,我十分小心,只限于那些我有足够强的关系、因而觉得自己在需要时能撑住学生工作的题目。一个显著的例外,是米歇尔·雷诺夫人(Mme Michèle Raynaud)关于基本群局部与整体莱夫谢茨定理的工作,那些定理是在适当的平展\emph{site}上以 1-栈(1-champs)的措辞表述的。这个问题在我看来是困难的(后来也确实证明如此),而对我所提出的那些猜想,我并没有证明的想法(再说那些猜想也几乎不可能有什么疑问)。这项工作一直做到七十年代初;雷诺夫人(如同此前她丈夫的情形)发展出一套精细而独到的方法,没有得到我或任何别人的帮助。这项出色的工作还引出一个问题:把雷诺夫人的结果推广到 $n$-栈($n$-champs)的情形——在我看来,在概形的语境里,这应当就是“弱莱夫谢茨定理”一类定理的自然归宿。不过,这里相关猜想的表述(它同样几乎不可能有什么疑问)本质上用到了 $n$-栈的概念;而对它的追寻,按说正是本书的主要对象\fnote{2}{}{(其实指的是《数学反思》第3卷,而不是眼前这第1卷《收获与播种》——见导论,第(v)页。)},正如书名《追寻栈》所示。我们大概会在相应的地方再回到这一点。

% ---------------- p.240 ----------------
\origpage{151}\origpage{150}另一个相当特殊的情形是黄春生夫人。我最初是在1967年12月在河内遇见她的,那时我在疏散后的河内大学讲了一个月的课程—讨论班。第二年,我向她提出了她的论文题目。她是在战争时期格外艰难的条件下工作的,与我的接触仅限于断续的通信。1974/75年她得以来到法国(趁温哥华国际数学家大会之便),于是在巴黎通过了论文(评审委员会由嘉当主持,另有施瓦茨(Schwartz)、德尼(Deny)、齐斯曼(Zisman)和我)。

最后我还该提到皮埃尔·德利涅(Pierre Deligne)和卡洛斯·孔图-卡雷尔(Carlos Contou-Carrère):这两人都多少算是学生,前者在1965—68年前后,后者在1974—76年前后。两人显然都有(并且一直有)非同寻常的才能,而他们运用这才能的方式极为不同,际遇也极为不同。到比雷(Bures)来之前,德利涅在比利时多少算是蒂茨(Tits)的学生——不过我怀疑,按通常意义说,他在数学上做过谁的学生。孔图-卡雷尔做过桑塔洛(Santalo)的学生(在阿根廷),有一阵还做过托姆(Thom)的学生(多少算吧!)。接触建立的时候,两人都已具备一个数学家的身量,只差孔图-卡雷尔缺少方法与行当上的训练。

我在数学上对德利涅的作用,仅限于零零碎碎地把我所懂的那点代数几何讲给他听;他学起来就像听故事——仿佛他向来就知道。顺带我也提出一些问题,而他多半当场或在随后几天里就找到了答案。这就是我所知道的德利涅最早的工作。1970年以后的工作(他的,以及我那些“正式学生”的),我只从极为零散遥远的回声里略知一二\fnote{3}{}{(我尤其有机会翻阅过贝尔特洛和德利涅寄给我的几份抽印本,承他们的好意。)}。

至于我对孔图-卡雷尔的作用,按他自己在论文开头所说,仅限于把他引进概形的语言。无论如何,他这几年作为国家博士论文所准备的工作,我只是远远地跟着;那是一个再时新不过的题目,超出了我的能力。

% ---------------- p.241 ----------------
\origpage{152}孔图-卡雷尔在大世界里经历了几番波折之后,最近终于在最后关头——而且(现在在我看来)是违着他的心意——不得不求我帮忙,充当论文导师并组成评审委员会。(这有让他显得像“1970年以后”的格罗滕迪克学生的风险,而在一种处境里,这可能会带来严重的不便……)我尽力履行了这项任务;这大概是我最后一次行使这个职能(在国家博士论文这一级上)。在这个有点特殊的场合,我格外高兴有让·吉罗(Jean Giraud)的友情协助:他也拨出一两个月的时间,仔细读了那部厚厚的稿子,并写了一份详尽而热忱的报告。

\section{$\varnothing$}

\notehead{20}{}
这让我想起莫妮克·阿基姆(Monique Hakim)所选的题目,说实话那也不见得更有吸引力;我真不知道她是怎么保持住精神的!如果说她有时也吃力,那至少没有到让她郁郁寡欢、闷闷不乐的地步;我们之间的工作是在亲切而放松的气氛里进行的。

\section{$\varnothing$}

\notehead{21}{}
也许更确切的说法是:对我这样的气质来说,要完全担起一个教师角色,我缺的仍是那份必要的成熟。我后天形成的气质,长期带着“男性”(或“阳”)特征过度占优的印记;而成熟的一个方面,恰恰是一种以“女性”(或“阴”)为主的“阴阳”平衡。

(后来补写)我看出,比成熟更缺的,是一种慷慨——在我至今的教学生涯里,我缺的正是某种慷慨;这种慷慨的表现方式,比在时间与精力上随时可用更为细微,也更为本质。在我第一个教学时期,这份欠缺并没有明显地表现出来(比如说,通过一连串失败的情形积累起来),多半尤其是因为它被选择来跟我工作的学生身上强烈的动机补偿了。相反,在第二个时期,从1970年至今,我觉得这份欠缺至少是原因之一,而且无论如何是最直接牵涉到我自己的那一条原因——我在研究层面(也就是从DEA这一级起)的教学中所看到的整体失败,原因就在这里。关于这一点,见《纲领概要》(Esquisse d'un programme)第8段,以及第9段“一项教学活动的总结”;在那里面透出的是这七八年来这项活动留给我的挫败感\fnote{4}{}{(另可参见后来补写的注(23iv)。)}。

\section{$\varnothing$}

\notehead{22}{}
也许不会再有多久了,因为我已经决定申请进入国家科学研究中心(CNRS),从而结束大学里的教学活动——这几年来,这活动变得越来越成问题。

% ---------------- p.242 ----------------
\origpage{153}\origpage{154}
\section{$\varnothing$}

\notehead{22$'$}{}
即便在1970年以后,当我对数学的兴趣在我生活里变得时有时无、退到了边缘,我也不认为有过哪一回,学生来叫我同他一起工作、而我推辞掉的。我甚至可以说:除了两三个例外,我1970年以后的学生对自己所做工作的兴趣,远远低于我本人对他们题目的兴趣——即便在我除了去系里露面的日子几乎不碰数学的那些时期,也是如此。因此,我1970年以前对学生的那份随时可供支配,以及作为其首要标志的、工作上的极端严格,对我后来多数学生来说恐怕毫无意义:他们做数学并没有什么信念,仿佛那是一种他们不得不对自己持续施加的努力……

\section{孩子与老师}

\notehead{23}{}
这里“传递”一词,与事情的实情并不真正相符——实情把我拉回到一种更谦逊的态度上来。这种严谨并不是一件可以传递的东西,最多只能被唤醒、被鼓励,而它从小就被忽视、被打消,被家庭环境,也被学校和大学。就我所能回想起来的而言,这种严谨一直存在于我的种种追寻之中——至少是那些智力性质的追寻——而我不认为它是由我父母传给我的,更不认为它传自哪位师长,无论是在学校里,还是在数学界的前辈当中。我觉得它似乎属于“天真”的诸属性之一,因而也属于每个人生来便被分得的东西。这种天真很早就“吃够了苦头”,于是不得不或深或浅地潜下去,以致在往后的一生里往往再也见不到它多少痕迹了。在我身上,由于一些我还不曾想去深究的原因,某种天真在智力好奇心这个相对无害的层面上存活了下来;而在其余一切地方,它都深深地潜了下去,神不知、鬼不觉!就像在每个人身上那样。也许,完全意义上的“教学”的秘密——或者不如说它的奥秘——就在于重新与这种表面看来已经消失的天真接上联系。但是,如果这联系尚未首先存在于、或尚未被重新找回到教师本人身上,那就根本谈不上在学生身上把它重新找到。而这时教师“传给”学生的,绝不是这种严谨或这种天真(两者在师生身上都是天生的),而是一种尊重,一种对这件通常遭人弃置的东西的、无声的重新肯定。

\section{$\varnothing$}

\notehead{23$'$}{}
不过,七八年来,在我的数学生涯里还有另一个慢性的“挫败之源”,只是这些年来它表现得隐晦得多。靠着反复出现,靠着同一种“令人挫败”的情形在我的教学活动中顽固地累积,它终究变得显眼起来,末了还爆发为一声“受够了!”,使我实际上终止了一切所谓的“指导研究工作”。在我的这番反思途中,我有一两次顺带触及这个问题,最后总算在结尾处多少把它审视了一番。我在那里至少描述了这份挫败感,并考察了它在我“重返数学”一事中所起的作用(参见第50段“一段过去的重量”)。

% ================================================================
%  第 12 章 第一部分的注 (原书 p.243-254, 边码 p.155-p.171)
%  源文: _work/src1/task-ch12b.md (12.26–12.50, 注 23″–44)
% ================================================================
% ---------------- p.243 ----------------
\section{怕玩}

\notehead{23$''$}{}
这个学生曾和我一起做了一整年的DEA“实习工作”,而在他与我的工作关系里,直到最后都保持着“紧绷”状态。那是一种坦率友好的关系,贯穿着一种无可怀疑的相互好感。然而可就是有这份“怯场”;这份恐惧——它真正的起因肯定不是对我这个人的畏惧,尽管看上去像是如此。要不是这个学生自己跟我谈起,我也许根本不会察觉这件事;他谈,无疑是为了或多或少“解释”一下他在这一年里工作几乎完全卡住的原因。

正如在其他学生身上发生过的那样——他们和他一样,起初对某种几何实质抓得很牢——卡住恰恰出现在要“做零件活儿”的时候:也就是把成形的陈述白纸黑字写下来,或者仅仅是把握我提供的那些陈述的意义与含义;我提议把它们当作一种语言的基石、当作“游戏规则”接受下来。“学校式的”反应几乎总是把面对“做研究”局面的学生推向:把导师传下来的那些心照不宣的“游戏规则”当作既定的、又必须服从的东西接受下来,而这些规则恰恰不是要试着说明白的,更不是要弄懂的。这些心照不宣的规则所采取的具体形式,就是语义学或计算上的“配方”,不妨说,像预科班的题解书那样(或者任何一本普通教科书)。此外,学生还期待导师给出“证明……”这种形式的任务——那是他以往经验里遇见的唯一一种数学“思考”形式。(而且我不认为大多数职业数学家、以及其他科学家的禀性有本质不同——只不过“导师”换成了“共识”:共识规定当下的游戏规则,并把它当作不可更动的既定之物。这共识还规定哪些“问题”是要去解决的;在这些问题之间,每个人都觉得自己有随意挑选的自由,甚至允许自己在工作过程中修改它们,乃至发明另一些……)。\origpage{155}我注意到,我面对有待探究的数学实质时那种完全不同的态度——因而也是面对学生时那种完全不同的态度——几乎必定会引发一种失措,而焦虑就是它的征兆之一。像一切焦虑一样,它倾向于给自己找一张面孔,投射到某个外部的“理由”上,不管那理由说得通说不通。焦虑最常见的一张面孔,恰恰就是恐惧。

这类困难在我教学生涯的第一个时期几乎没有出现,也许只有两例例外:那两次“师生”关系没超过几个星期;也许还有(我说不准)“忧伤的学生”那一例——他也许感到自己被“铆”在一个完全不启发他的题目上,而他本来完全有更换的自由。在那个(我也谈到过的)长期被某种怯场困扰的学生身上,原因显然在别处。他的工作丝毫没有卡住,恰恰相反,他对选定的主题极为自在,并在其上做出了一项有规模的基础性工作。那个时期我的学生大多是巴黎高师的旧生,他们与亨利·嘉当(Henri Cartan)的接触,已经让他们见识了另一种“接近数学”的方式。与之(可以说)相反的另一极,是我在蒙彼利埃大学任教的第二个时期:正是在一年级学生那里,我所说的焦虑最少干扰思考性的工作。在这些学生中的许多人身上,面对一种不同进路时的惊讶既不引起焦虑,也不引起封闭,相反带来的是敞开和兴致——终于可以做点有趣的事了!据我的观察,几年大学对一个学生创造性禀性的作用是彻底而毁灭性的。奇怪的是,在这方面,长年中学的作用似乎相对无害。原因也许是:大学那几年正落在这样一个年纪——我们身上与生俱来的创造力,最迟必须通过个人的工作表达出来,否则就要永远沉船,至少就智力上的创造性工作而言是如此。肯定是出于一种健康的本能,我在做学生的那些年(也是在蒙彼利埃大学),几乎不踏进课堂,把几乎全部精力都投给了个人的数学思考\origpage{156}。

% ---------------- p.244 ----------------
\section{两兄弟}

\notehead{23$'''$}{}
这个学生身上的对抗,一上来就采取了“阶级对抗”的形式:我是那个对他的数学前途拥有“生杀之权”的“老板”,可以凭自己的好恶决定……当然,事情本身只能证实这种看法,因为我很快就终止了对这个学生(已变得令人难受的)责任。这把他置于一个微妙的处境——在眼下的世道里,找一个“老板”可不是那么显而易见的事,尤其是题目已经选定的时候。在另一个学生身上,合法期待落了空,对抗采取了类似的形式。我被感受为专横的“学阀”,不能容忍那些被他视为下属的人(学生或级别较低的同事)有任何反驳。

这样一种“阶级态度”,在我与第一个时期学生的关系里,从未有过一丝一毫的表现。显而易见的原因是:在1970年以前的形势下,学生通过论文之后毫无疑问会得到一个副教授(maître de conférences)的职位,因而享有与我相同的社会地位,即“大学教授”。数字最能说明问题:1970年以前开始跟我做工作的十一个学生,一完成工作就得到了副教授职位;而在我的指导下或多或少工作过的那二十来个学生里,没有一个获得这样的职位。诚然,他们之中只有两个人有足够的动力去做国家博士论文(而且两人的论文都很出色)。

因此,在这第二个时期,某些暧昧(其深层根源始终被遮蔽着)采取了阶级对抗的形式,采取了对“老板”的不信任(这种不信任被呈现、也被感受为“本能的”“发自肺腑的”)——这就不足为奇了。在那些或多或少当过学生样子的人里,有一个与我的友好关系持续了十来年,没有出现过表面上的对抗插曲,却仍带着同一种暧昧,表现为一种不信任的态度,被“储备”在明显的善意之后。说实话,我从未被这种例行公事式的“不信任”骗过;它在我看来,尤其是这位朋友认为应当给自己找的一个理由,好让自己不去冒险越出他选定的、界限分明的领地——无论在职业生活里,还是在整个生活里都是如此。可这件事他本来可以自由地去做,没有人(顶多只有他自己!)会要他交代……

这三个例子,也是我整个教书经验中仅有的、学生(或多少像学生的人)与我之间的关系中某种暧昧以“阶级态度”表现出来的情形。当这种态度出现在大学“团体”内部的同事之间时,显得格外暧昧:在这样一个团体里,双方都享有与常人处境相比极为悬殊的特权,而这些特权使级别(以及薪金)的差别显得相对微不足道。我还注意到,一旦当事人自己也被提升到前一天他还在指责别人的那种处境,这些态度就像中了魔法一样消失了(而且理所当然!)。

此外,在我于数学界内部(常常也在外部)目睹的多数——如果不是全部——冲突局面里,我都辨认出一种类似的暧昧\origpage{157}。那些“有位置的人”,无论其级别是否合乎他们的期待(合理与否),都享有相当闻所未闻的特权,那是任何别的职业或生涯都无法提供的。那些没有位置的人,则向往同样的安全感和同样的特权(这并不必然妨碍他们对数学本身感兴趣,有时还做出漂亮的工作)。在眼下的世道里,抢位置竞争激烈,没有位置的人常被当作流浪汉对待:我不止一次感到那种默契——在乐于羞辱人的人与被羞辱的人之间;后者咽下去,缩起来。他苦涩与敌意的真正对象,并不是那个行使了权力的人,而恰恰是他自己:他自己缩了起来,并且把对方奉为那种他乐意行使的权力的持有者。乐于羞辱人的人,也是那个报复并补偿(却永远抹不掉……)一场长期承受、早已深埋并被遗忘的羞辱的人。而默许自己被羞辱的人,是他的兄弟和效仿者:他暗中羡慕对方,并在苦涩里埋下那场羞辱,以及羞辱带给他的、关于他自己的那条卑微的消息。

% ---------------- p.245 ----------------
\section{一次教学的失败(1)}

\notehead{23$^{iv}$}{}
\origpage{158}这些字句写下之后,我有机会与1970年之后的两位旧学生谈过,想同他们一起探究我在蒙彼利埃大学研究层次上教学失败的原因。他们对我说:我那种倾向——低估同化那些对我熟悉、对他们却不熟悉的技术对他们可能意味着的困难——在他们身上产生了令人泄气的效果,因为他们感到自己始终达不到我对他们的期望。此外(这件事在我看来意义还要大),他们有时感到受了挫伤:当他们已经十分接近的时候,我不是把发现的乐趣留给他们、让他们靠自己的办法去发现,而是把袖子里藏着的一个成形陈述“卖给他们”、向他们泄底。那之后,他们剩下的就只有做“练习”了(而那并不特别让他们着迷)——也就是证明那个陈述。我先前在一则注(注 21)里指出过的自己身上那份“缺乏慷慨”,正落在这一处,当时没有就此多说。主要正是这类扫兴之事,构成了我在他们两人身上——在那么出色的开端之后——研究兴趣消失这件事上的个人贡献。

我意识到,我在1970年以前并不比之后更慷慨。如果说那时我没有遇到同样的困难,那无疑是因为:那个时期来找我的那类学生,有足够的动力,甚至能从一场“长长的练习”里品出魅力——那是学习本行、并一路学到一大堆东西的机会;同样,对于一个我“泄了底”的起步陈述,他们也能靠自己的办法从中抽出一大群远远超出第一个的陈述。当我换了教学地点,我在向新学生提议的思考主题上做了必要的调整,办法是选择那些可以被直接直觉把握的数学对象,不依赖任何技术行囊。但这必要的调整本身并不够,因为禀性上的差别(我的新学生与从前的学生相比)比单纯的行囊差别还要重要。这还印证了前面(第25段开头)说过的一点:我在“导师”这个角色上有某种不足;这不足在我教学生涯的第二个时期,比在第一个时期显露得强烈得多。

\section{$\varnothing$}

% ---------------- p.246 ----------------
\notehead{23$^{v}$}{}
\origpage{159}这种差别有一个格外醒目的征兆,出现在“外国人事件”之际——我曾有机会谈到过(第24节)。当时,许多与我完全陌生的人向我表达了同情,而我记不起1970年之前的哪个学生想到过以这种方式表示什么,更不用说在我投身的那场行动中向我提供任何帮助。相反,在我看来,第二个时期的学生或旧学生里,没有一个不曾向我表达同情与声援,还有几个积极参加了我在当地开展的运动。在这个狭小圈子之外,1945年法令那件事也在许多仅仅知道我名字的大学生中间激起了一定的波澜;在我被传唤的那一天,有不少人来到法院,以示声援。这后一个情形还提示:我所察觉到的“之前”与“1970年之后”学生态度上的差别,也许较少表达他们与我之间关系的差别,而更多表达的是一种心态的差别。显然,我“之前”的那些学生已经成了重要人物,而要让重要人物动心,是需要很多条件的……但我1970年离开IHÉS、投身战斗行动那一幕,似乎表明并不只是这一点。那是个他们中间还没有谁太像重要人物的时刻,然而我记不起他们中有谁对我所投身的活动表现出丝毫兴趣。我倒认为,那一定让他们所有人,无一例外地感到不自在。这仍然合乎“心态差别”的方向,但这不能只归因于社会地位的差别。

\section{$\varnothing$}

\notehead{24}{}
我要谈的这套伦理,同样适用于其他一切围绕研究活动形成的圈子:在那里面,能不能让人知道自己的结果、并从中得到承认,对每个成员的社会地位都是“生死”问题,甚至是他作为该圈子成员的“存亡”问题,而这对本人及其家庭意味着什么,也就可想而知。

\section{职业道德的共识——以及信息的控制}

\notehead{25}{}
除了与迪厄多内的那次谈话,在我作为数学家的生涯里,我想不起还有哪一次我参与或旁听的谈话\origpage{160},谈到了这一行的伦理、谈到同行之间关系里的“游戏规则”。(这里不算七十年代初围绕“生存与生活”运动发生的那些关于科学家与军事机构合作的讨论。那些讨论并不真正涉及数学家彼此之间的关系。再说,我在“生存与生活”里的许多朋友,包括谢瓦莱和盖吉,都觉得我当时——尤其在初期——如此强调这个我格外敏感的问题,使我偏离了更紧要的日常现实,恰恰就是眼下这份反思所考察的那类现实。)在我与学生之间,从未谈到过这些事。我相信,那份默契只限于一条规则:不把别人的想法——只要自己有机会知道——说成自己的。我觉得这条共识自古就有,直到今天在任何科学圈子里都没有被质疑过。可是,若缺少另一条补充规则——保证每个研究者都有可能让人知道自己的想法与结果——第一条规则就只是一纸空文。在今天科学界,处在声望与权力位置上的人,把持着对科学信息的任意支配权。
% ---------------- p.247 ----------------
在我从前熟悉的那个圈子里,这种把持已不再被迪厄多内所说的那种共识所节制——而那种共识,也许从未存在于他充当代言人的那个小圈子之外。处在权力位置上的科学家,实际上能收到他认为值得收到的一切信息(往往还超出);而对于其中很大一部分,他有权阻止其发表,同时自己留着这份信息的好处,把它打发为“没意思”“多少已经知道了”“平凡”之类……这种局面我在注(27)里还要再谈。

\section{$\varnothing$}

\notehead{26}{}
布尔巴基的“创始成员”是亨利·嘉当、克劳德·谢瓦莱、让·德尔萨特(Jean Delsarte)、让·迪厄多内、安德烈·韦伊。他们全都健在,只有德尔萨特例外——五十年代,他未及天年便被带走了;那个年头,这一行的伦理还普遍受到尊重。

重读本文时,我曾想删掉这一段:在这里,我可能给人一种印象,仿佛在颁发“正直”(或不正直)证书,而当事人对此毫不在乎,这事也不该由我来做。这一段可能引起的保留,肯定是有道理的。\origpage{161}可我仍旧留着它,为的是证词的真实,也因为这一段确实如实复现了我的感受,哪怕这些感受并不合宜。

\section{“年轻人的势利”,或纯粹性的捍卫者}

\notehead{27}{}
罗尼·布朗(Ronnie Brown)向我转述了J.H.C.怀特黑德的一句话(他是怀特黑德的学生),说的是“年轻人的势利:他们以为一条定理平凡,只因为它的证明平凡”。我从前的许多朋友,不妨把这句话好好想一想。这种“势利”如今绝不限于年轻人;我认识不止一位有声望的数学家,常这么做。我对此格外敏感,因为我在数学上(别处也一样……)做得最好的那些东西——我引入的、在我看来最丰饶的概念与结构,以及我靠耐心而执拗的工作从中拔出的那些本质性质——全都落进了“平凡”这个评语里。(如今这些东西,若作者还不是个名人,哪一件都很难被《科学院通报》(CR)当作一篇短讯接受!)我这一辈子作为数学家的抱负——不如说是我的激情与快乐——始终是找出那些显然的东西;在这本书里(包括眼前这一章引论性的文字里……),这也同样是我唯一的抱负。决定性的东西,常常首先就是看见那个尚未被看见的问题(不管答案是什么,也不管答案是否已经找到),或者提出一条陈述(哪怕是猜想性的),把一种尚未被看见、或尚未被理解的局面概括并容纳进来;这条陈述若被证明了,那么证明平凡与否无关紧要——那完全是附带的事——甚至一条仓促的、暂用的证明后来被证明是错的,也无关紧要。怀特黑德所说的势利,是那种厌倦了的享乐者的势利:他要先确认一瓶酒贵得惊人,才肯去欣赏它。这些年里,不止一次,旧日的激情一阵阵回来抓住我,我把自己最好的东西捧出去,却看它被那份自得打了回来。我感到一种至今仍活的痛;一份快乐落了空——可我并没有因此流落街头,而且幸而我并不需要替自己写的一篇文章找地方发表。

% ---------------- p.248 ----------------
\origpage{162}怀特黑德所说的势利,当它由一个有权者施于一个任其摆布的研究者身上时,就不只是一种对事物之美的麻木或封闭,而是一种滥用权力、一种不诚实:那位研究者完全在他的掌握之中,他尽可随意吸收、使用其想法,同时以“显然”“平凡”因而“没意思”为借口,堵住它们的发表。我这里甚至没有想到通常意义上的剽窃那种极端情形——在数学界,那想必还极罕见。可是从实际角度看,对为此付出代价的研究者来说,局面是一样的;而使这局面成为可能的内心态度,在我看来也差不太多。它只是更舒服些:因为它伴着一种对他人无限优越的感觉,伴着良心安然与内心满足——那是一个自命为数学之不可触犯的纯粹性的不妥协捍卫者的人。

\section{$\varnothing$}

\notehead{28}{}
写前面那几页时,我起初在“把心里的东西全倒出来”的欲望与一份保留、谨慎的顾虑之间分裂。于是我就停在差不多、大约莫的地方——这肯定是我那份不适、“我什么也没学到”之感的主要原因。自从写下那几行承认不适的文字之后,我把这几页重写了两遍——它们让我心里不满意;我把自己更清楚地放了进去,也把事往更深处推。一路下来,我确实终于“学到了点东西”;我也相信,与此同时我摸到了某件要紧的东西,它超出了这个具体个例,也超出了我这个人。

\section{$\varnothing$}

\notehead{29}{}
我这里说的,是在数学里、或在另一项纯智力活动里那种强烈而长久的投入。反过来,这样一份激情的展开——它可以是重新认识我们身上一种被遗忘的力量的方式,也是与一种不肯就范的实质较量的机会,并且一路下来,用某种真正属于我们自己的东西去更新、丰富我们的身份感——这样一种展开,很可能是内心路程上的一个重要阶段,是成熟过程中的一个重要阶段。

\section{$\varnothing$}

\notehead{30}{}
这几年,接力的是我的孩子们:由他们去教一个有时不大情愿的学生,人生的种种奥秘……

\section{$\varnothing$}

\notehead{31}{}
\origpage{163}我这里想到的是求知欲的“阳”的形式——那探问、发现、为显现者命名的形式……正是由于被命名,出现的知识才变得不可逆转、不可抹去(哪怕它后来被埋掉、被遗忘,不再活跃……)。求知欲的“阴”的、“女性”的形式,在于一种敞开、一种接纳,在于无声地迎候一种出现在我们存在更深层——思想到不了的地方——的知识。
% ---------------- p.249 ----------------
这样一种敞开、一种突然到来、一时间抹去一切冲突痕迹的知识的出现,仍像一份恩典:它触到深处,而它可见的效果也许转瞬即逝。不过我猜,这种在我们一生中某些稀有时刻如此到来的、无言语的知识,同样不可抹去;它的作用甚至延续到我们对它的记忆之外。

\section{炉中百铁,或:干熬无用!}

\notehead{32}{}
我还在做泛函分析的时候,也就是直到1954年,我有时会没完没了地死咬一个问题,解不出来;那时我已经没有新想法,只是绕着旧想法的圈子打转,而那些想法显然已经“咬不住”了。整整一年都是这样,尤其是拓扑向量空间里的那个“逼近问题”;它要再过二十年才被解决,用的完全是另一路的方法,以我当时所处的境地,那些方法只能与我擦肩而过。当时驱动我的不是欲望,而是一种固执,以及对自己内心所发生之事的不了解。那是难受的一年——我一生中唯一一段做数学变得难受的时期!我需要这段经历才明白:“干熬”没有用——一件工作一旦到了停顿点,而且只要停顿被察觉,就该转到别的事上去;至于悬置的那个问题,等到更合宜的时候再回来。这个时刻几乎总是很快就出现——问题自己成熟起来,我并没有做出碰它的样子;靠的只是兴致勃勃地做那些看来与它毫无关系的问题的功效。我确信,当时若我硬撑下去,哪怕十年也一事无成!从1954年起,我在数学上养成了一个习惯:炉子里总同时烧着许多块铁。我一次只打一块,可由于一种不断更新的奇迹\origpage{164},我在一块上做的工作也惠及所有其余的,它们各自等着自己的时辰。与我第一次接触冥想时也一样,而且完全不是我刻意为之:随着反思继续,有待审视的灼热问题,数目一天天增多。

\section{$\varnothing$}

\notehead{33}{}
这并不意味着:工作中纸(或黑板,那不过是它的一种变体!)不在场的那些时刻,在数学工作里就不重要。尤其如此的是那些“敏感时刻”:一种新的直觉刚刚出现,这时要做的是以一种更整体、更直觉的方式去“认识”它,而不是做“零件活儿”——而反思的这个不定形阶段,正是为那件事作准备。在我身上,这类思考主要在床上或散步时进行;我觉得它只占用于工作的全部时间中一个相对不大的部分。同样的观察也适用于我迄今所从事的冥想工作。

\section{无力的拥抱}

\notehead{34}{}
“拥抱”这个词在我这里绝不是单纯的比喻;日常语言在此映出的是一种深层的同一。人们不无道理地会说:这么说,“没有惊叹的拥抱无能为力”就不对了——若按字面如此,大地就该人口凋零,或者干脆荒无人烟。极端的例子是强奸:惊叹当然不在场,可有时被强奸的女人身上仍孕育出一个生命。由这样的拥抱生下的孩子,肯定免不了带着那个印记;那印记会是他分得的那一份“包裹”的一部分,该由他自己去承担。可这并不妨碍:一个新生命确实被孕育、被生了下来;确实有过创造,那是力量的标志。也确实如此:我见过某个满身自得的数学家找到并证明了漂亮的定理,那正是拥抱不乏力气的标志!但同时也不假:若某位数学家的生活被他的自得闷死(在某种程度上,我自己生活中某个时期也曾如此),那么他与数学的这些拥抱所结的果子,对他、对别人都不是什么福分。从强奸中生下的孩子的父亲与母亲,也可以这样说。我说“无力的拥抱”,首先指的是这样一件事:在那个自以为在创造的人身上,无力生出更新——他只造出一个产品、一件他之外的东西,在他自身里没有深的回响;
% ---------------- p.250 ----------------
\origpage{165}这产品不但不解放他、不在他里面造出和谐,反而把他更紧地捆在他身内那份他作囚徒的自得上,那自得不断催他去生产、再生产。这是深层上的一种无力,藏在“创造力”的表象背后——而那归根到底不过是无缰的生产力。

我也有充分机会看出:自得、无惊叹之能,其性质是一种真正的盲,是自然感受力与嗅觉的堵塞;这堵塞即使不是全面而永久的,至少在某些具体情形里是明显的。在这种状态里,某位有声望的数学家有时会在恰恰他擅长的那些事情上,显得比最死心眼的小学生还蠢!在另一些场合,他会做出技术精湛的奇迹。可我怀疑,他是否还有能力去发现那些单纯而显然、有力量更新一门学科或一门科学的东西。那些东西离他太低、太远,他已经不屑于去看!要看别人不屑于看的东西,需要一份他早已失去、或被他放逐的天真……近二十年间数学产出惊人地增长,新结果的纷繁令人无所适从,连只想稍微“跟上”的数学家都被淹没;而与此同时(就我从四处传来的回声所能判断的而言),我所多少熟悉的那些大思考主题,却几乎没有真正的更新、没有大幅度的转化(不只是积累式的转化)——这肯定不是偶然。更新不是量的事情;它与投入的量无关——那种量可以用若干“数学家人日”来计量,即某些“水平”的某些数学家花在某个题目上的日数。一百万个数学家人日,也无力生出像零那样孩子气的东西——而零更新了我们对数的感知。只有天真拥有这份力量;它的一个可见标志,就是惊叹……

\section{$\varnothing$}

\notehead{35}{}
这份“天赋”不是谁的特权,我们全都生来带着它。当它在我身上显得不在时,那是我自己把它赶走的;要不要重新接纳它,只在于我。在我身上,或在某某身上,这份“天赋”的表达方式与另一个人不同,不那么有感染力,也许不那么不可抗拒,但它照样在场;我也说不出它是否不那么有力。

\section{$\varnothing$}

\notehead{36}{}
\origpage{166}这种对美的细腻感受力,在我看来与一样东西紧密相连;
% ---------------- p.251 ----------------
我曾以“严要求”(对自己的)或“严格”(取其完整意义)之名谈到它,把它描述为“对我们自身中某种细腻之物的注意”,对所要探问之物的理解品质的注意。对一件数学东西的理解品质,离不开对该物特有之“美”的一种或深或浅、或全或欠的感知。

\section{$\varnothing$}

\notehead{37}{}
我想,几乎不必补充:这项长久的工作日复一日地显出许多别的东西,远不止我刚才以碑铭般的简省形式交出的那个“结果”。冥想的工作,与由某个打算考察的具体问题所驱动的数学工作,并无不同。所走那条路(它通向、或不通向对最初问题或多或少的澄清)上的种种曲折,往往比最初的问题或“最终结果”更有意思。

\section{$\varnothing$}

\notehead{38}{}
这些注其实是写给……的那封长信的继续,那封信成了本书第一章。它们用打字机打出来,是为了让这位旧日的朋友能读,也为了另外两三位(尤其是罗尼·布朗),我以为他们可能会感兴趣。再说,那封信从未得到回音,收信人也没有读它:将近一年之后(我问过他是否确实收到),他由衷地惊讶——我怎么会有一刻以为他会读它,鉴于人们该期待我做的那类数学……

\section{$\varnothing$}

\notehead{39}{}
这一时期,除别的之外,也正是《穿越伽罗瓦理论的长征》的时期;在《纲领概要》里(第3段:“与童画相联系的数域”)谈到过它。

\section{来访}

\notehead{40}{}
对这个梦所做的工作,是一封英文长信的内容;收信人是一位朋友兼同行,他前一天在我家匆匆过了一趟。做梦者用来从一片表面的虚无中唤出这个逼真得惊人的梦的那些材料,有些显然取自那位亲爱的朋友来访的短暂片段——我已有近十年没再见过他。\origpage{167}于是,在工作头一天,而且与我过去的经验相违,我竟以为可以由此断定:我做的这个梦关系到我的这位朋友,多过关系到我——该做这个梦的是他,不是我!这是一种回避梦境消息的办法;而那消息(凭我过去的经验,我本该一开始就知道)只关系到我,不关别人。在这第一阶段——浮浅的——工作之后的夜里,我终于意识到了这一点;第二天,我在同一封信里接着做了下去。自那封值得记住的信之后,我再没有收到这位朋友的任何音讯——他是我有过的最亲近的朋友之一。

% ---------------- p.252 ----------------
这项工作,是唯一一次以信件形式(而且用的还是英文)成形的冥想,因此我手里不再有它的书面痕迹。这段插曲让我格外震动——在众多别的例子当中,它们都表明:任何越出某层门面的工作迹象,任何把极简单的事实——而人们通常认为有义务对之视而不见——带到光天化日之下的工作,是多么在别人身上引起不适与恐惧。这一点我在后面还要谈(见第47段《孤独的历险》)。

\section{克里希那穆提,或变成羁绊的解放}

\notehead{41}{}
若说我从这次阅读中得到的只有一套词汇,以及一种把它据为己有、最终又顺理成章地用它取代现实的倾向,那是不准确的。我拿到手的第一本克里希那穆提(Krishnamurti)的书之所以那么震动我(而且我还只来得及读其中几章),是因为他说的东西彻底掀翻了我许多自以为当然的事;我立刻看出,那些都是老生常谈,一直是我所呼吸的空气的一部分。同时,这次阅读第一次把我的注意力引向一些分量很重的事实,尤其是对现实的逃避——它是心智最强大、最普遍的制约之一。这给了我一把关键的钥匙,去理解那些此前一直不可理解、并因此(在我五六年之后发现冥想之前,我并未察觉这一点)滋生焦虑的处境。我立刻就能在我周围各处看到这种逃避的现实。这解开了某些焦虑,却没有改变任何本质的东西,\origpage{168}因为我只在他人身上看到那份现实,同时(理所当然地)以为它在我身上并不存在,以为我总之是那条证实规则的例外(而且对这么一个确实了不起的例外,我没有再问任何别的问题)。事实上,我对别人、对自己都毫无好奇。这把“钥匙”只有在被深入之欲驱动的人手里才能开锁。在我手里,它成了驱邪和做姿态。

直到1974年初,我才第一次认清:我生活里的那种毁灭一步步跟着我,不可能只来自别人;我身上有某种东西在吸引它、喂养它、让它延续。那是一个谦卑与敞开的时刻,适宜于一次更新。可那次更新当时仍是边缘的、短暂的,因为缺少深入的工作。那个“我身上的某种东西”还很模糊。我看得出那是爱的匮乏;可是,要有一项工作去更贴近地界定:爱的匮乏在我身上何处、如何发生,它如何显现,它具体造成了什么后果,等等——这样的念头,既不可能来自我至今认识的任何圈子或任何人,也不可能来自克里希那穆提。(恰恰相反,K.喜欢强调一切工作的徒劳,他自动把工作等同于“我”的“成为之饥”。)就这样,我拿一份借来的“智慧”当唯一的罗盘,看不见还有别的事可做,只好耐心等着“爱”像圣灵的恩典一样降到我身上。

然而,我刚在一个浪的深处学到的那点卑微真相,已经激起了一股强劲的新能量之浪,可与两年半之后承载我第一次投入冥想的那股相比。这份能量当时并没有完全闲置。几个月后,我因一场凑巧的事故动弹不得,它承载了一份(写下来的)反思:我一生中第一次审视那套世界观——它一直是我与他人关系未曾言说的基础,来自我的父母,尤其来自我母亲。当时我非常清楚地看出:这套看法已经破产,它无法说明人与人之间关系的实际,也无法促成我这个人以及我与他人关系的舒展。这份反思仍带着
% ---------------- p.253 ----------------
“克里希那穆提风格”的印记,也带着克里希那穆提式的禁忌——禁止一切朝向理解的真正工作。可它毕竟把几个月前生出的一份认识变得可触而不可逆转——\origpage{169}那认识起初一直模糊、游移。那份认识,当时世上任何一本书、任何别的人都不可能带给我。

要够得上冥想的品质,这份反思最缺的,是对我自己这个人、对我关于自己的看法的注视——而不只是对世界观的注视,对一个公理系统的注视,而在那系统里我并不真的“有血有肉”地出场。它还缺另一样:在当下、在反思进行的那一刻对我自己的注视(而那份反思还停留在真正工作的此岸);那样一种注视,本来会让我既察觉一种借来的风格,也察觉我对这些笔记的文学面貌有某种自得,也就是说,缺少自发与真实。尽管它如此不足,尽管它对我与他人关系的直接影响相对有限,这份反思在我看来仍是一个阶段——考虑到起点,这大概是必需的——通向两年后该发生的那次更深的更新。正是在那时,我终于发现了冥想——同时发现这第一个未曾料到的事实:关于我自己这个人,是有东西可发现的;而那些东西几乎完全决定了我一生的走向,以及我与他人关系的性质。

\section{有益的剥离}

\notehead{42}{}
那个“有冲击力”的事件,是1969年底的发现:我觉得自己属于的那个机构,部分经费来自国防部的款项——这与我基本的公理不相容(今天也仍然不相容)。这个事件是一整串别的(一个比一个更有揭示力!)事件中的第一个;它们的结果是我离开高等科学研究所(IHÉS),并由此一环扣一环地彻底换了圈子与投入。

在IHÉS的英雄年代,迪厄多内和我是它仅有的成员,也是仅有的为它在科学界赢得信誉与听众的人:迪厄多内靠主编《数学出版物》(头一卷早在1959年就出了,即莱昂·莫尚(Léon Motchane)创办IHÉS的次年),我靠《代数几何讨论班》。最初那些年,IHÉS的存在极为飘摇:经费没有着落(靠几家\origpage{170}充当赞助人的公司的慷慨),唯一的场所是巴黎蒂耶尔基金会(Fondation Thiers)借出的一间厅(它显然很不情愿),专供我讨论班开讲的日子\fnote{5}{}{IHÉS最近为纪念其成立二十五周年出的一本小册子(承尼科·凯珀(Nico Kuiper)好意寄我一本),对这样艰难的开端只字不提——也许认为它们配不上这个场合的庄重吧;去年,这个周年被大张旗鼓地庆祝了一番。}。我多少觉得自己,连同迪厄多内,是我所属这个机构的“科学”共同创办人;我还真打算在那儿过完这辈子!我终于对IHÉS有了很强的认同;而我的离开(由于同事们的漠然),当时被体验为从另一个“家”被剥离,后来才显出是一次解放。

回过头看,我意识到自己身上想必早有更新的需要,我说不清是从什么时候起。我离开IHÉS的前一年,我的精力投入突然翻了个面:丢下那些昨天还烫手的工作、那些最让我着迷的问题,转而去投身生物学(是在一位生物学家朋友米尔恰·杜米特雷斯库(Mircea Dumitrescu)的影响下)——这肯定不是单纯的巧合。我投身其中,是打算在IHÉS内部做
% ---------------- p.254 ----------------
\origpage{171}长期投入的(这也符合该机构多学科的宗旨)。这肯定只是给一种深得多的更新需要找了个出口;那种更新在IHÉS那种“科学蒸笼”的气氛里无从完成,而是在我已提及的那场“觉醒的瀑布”中发生的。共有七次,最后一次在1982年。“军费”那段插曲是天赐的,因为它触发了这些“觉醒”中的第一次。国防部,以及我在IHÉS的旧同事们,最终都配得上我全部的感激!

\section{$\varnothing$}

\notehead{43}{}
“我写的那部诗作”里包含许多我第一手知道的东西;它们今天在我生活中、在一般的“生活”中,显得与它写成之时同样重要——那时我是打算把它发表的。我之所以没有发表,主要是因为后来意识到:那形式被一种刻意“作诗”的意图所累,以至于它整体构思过于搭建,许多段落缺少自发,有些地方甚至僵硬或浮肿得难受。这种有时浮夸的形式,是我当时禀性的映照:在那里面,常常确实是“老板”在领舞——笨重得很,那是不用说的……

\section{$\varnothing$}

\notehead{44}{}
不用说,我在这里是把这样一个假设撇开的——至少可以说,它并非不可能:一场原子战争,或另一场同类的“乐事”,突然闯来,把那个叫做“数学”的集体游戏一举粗暴地终结,顺带还终结别的好些东西……


% ================================================================
%  三 · 第二部 (第 13-16 章, 原书 p.255-412)
% ================================================================
\partdivider{第二部}{葬礼(I)\\[0.35em]{\normalsize 或中国皇帝的新衣}}{葬礼(I)或中国皇帝的新衣}

% 原书第257页献辞
\dedication{%
献给曾是我朋友的人,\\[0.5em]
既献给少数仍是朋友的人,\\[0.5em]
也献给众多前来为我的葬礼齐声唱和的人。\\[2em]
纪念那场难忘的讨论会……\\[0.5em]
以及全体教众……%
}

\setcounter{chapter}{12}
\markboth{收获与播种·第二部}{}

% ================================================================
%  第 13 章 A) 遗产与继承人 (原书 p.259-280)
%  13.1 身后的学生 | 13.2 孤儿们 | 13.3 时尚——或名人传
%  源文: _work/src2/source.md
% ================================================================
\providecommand{\fnoteplainbody}[3]{%
  \begingroup\renewcommand{\thefootnote}{#1#2}\footnotemark
  \renewcommand{\thefootnote}{#1}\footnotetext{#3}\endgroup
}

\chapter{A)~遗产与继承人}

\section{I.~身后的学生}

\subsection{1.~教学的失败(2)——或创造与自负}

\notehead{44$'$}{[该注由第(I)部分《愚妄与更新》第八章“孤独的历险”第~50~节所召,原书第~227~页]}

\pmark\origpage{173} 这段文字,在读了“一段过去的重量”这最后一节的那位朋友那里“叮”地一声响了\fnote{1}{(*)}{(5月10日)那位朋友不是别人,正是Zoghman~Mebkhout;他慨然允准我揭开我在本注所引那封信(1984年4月2日)的出处上本来觉得应当维持的匿名。}。他写道:“在你许多昔日学生身上,如你所说,那个咄咄逼人、甚至近乎破坏性的‘老板’的形象,依然强烈。于是就有了你的那种印象。”(我想,这里说的“印象”,就是这一节以及补足它的注~n\textsuperscript{o}s~46、47、50 的某些段落里所表达的东西。)他在信的更前面写道:“首先,我认为你暂时离开数学[\,!],离开得好,因为你和学生之间当时有某种隔膜(当然,Deligne除外)。他们都有点被惊呆了……”。

这是第一次,关于我1970年以前“老板”角色的这种论调传到我耳朵里——超出了客套的恭维!还在同一封信的更前面:“……我明白了,你的昔日学生[按:指那些‘1970年以前的’]不太明白数学的\textbf{创造}是怎么回事,而你或许有一份责任……的确,在他们那个年代,问题都已经摆好了……”\fnote{2}{(**)}{(5月10日)前面这段引文被删削得很厉害,是出于对通信人匿名的尊重。完整的引文以及对其真意的评说,见下一条注;起初因缺乏更详尽的信息,我没有领会这段话的真正含义。}。

我的通信人无疑是想说:是我在提“问题”,连同一并要发展的那些观念,而不让我的学生们自己去找到这些问题与观念;也正是在这一点上,我或许遮蔽了他们,使他们不识数学创造工作中那要紧的部分。这恰与一次谈话留下的印象相合\dmark{}\origpage{174}——那是与两位1970年\textbf{之后}的前学生的谈话,前面一条注(注(23iv))里谈到过。的确,我在向我走来的学生身上,首先寻找的是\textbf{合作者},来发展那些在我之内已然成形的直觉与观念,去给一辆已经在那里的大车“推轮子”;他们因而不必从某种虚无中(如我的通信人当年不得不做的那样)去拉动它。然而——让一种可触的、柔软而致密的东西,从不可触之物的迷雾中取得形体——这恰恰从来就是数学工作于我最迷人的方面,也是工作里我尤其感到“创造”在进行的那一部分:某种比一个单纯的“结果”更精微、更要紧的东西的“诞生”。

有时,当我看到昔日学生中的某一位,以轻蔑的态度对待这件弥足珍贵的东西,在他身上抖开 J.H.C.~Whitehead所说的那种“势利”(就是瞧不起自己“本来也会证明”的东西)\fnote{3}{(*)}{见注“年轻人的势利——或纯洁的捍卫者”(n\textsuperscript{o}~27),原书第247页。},我无疑并非全然无关——这样那样地,我脱不了干系。我教学的失败,在1970年之后的一段时期里触目经心,如今又以一种不同而更隐蔽的形状,显现在我第一时期的教学里——尽管按通常的说法,那一期呈现为完完全全的成功!这一点,近年来我不时有所瞥见,也在给几位前学生的信里提到过,而直到如今,没有从他们任何一位那里得到真正的回音。

不过,倘若说我交给学生们做的、他们同我一起做的工作,是纯技术性的、纯粹例行公事式的,无力调动他们的创造才能,这样说怕也不确切。我交到他们手里的,是一些切实而可靠的出发点,他们在其中享有全部的选择自由,可以由此跃起,正如我自己在他们之前曾经做过的那样。我不记得曾给哪位学生出过一个连我自己也不会乐于去做的题目;他们中任何一位与我同行时走过的艰难路段,也都不比我自己在数学家生涯中独自走过的那些更艰难;只要明白工作必须做、也没有别的路,我就没有因此气馁或尥蹶子。

\pmark\origpage{175} 同时我也觉得,我今天所确认的这场失败,其原因比我所出题目的类型、比这些题目在多大程度上含混或分明,更为精微。我在这场失败中的那一份,似乎更应归咎于我对待数学时自身的自负态度;这些态度,我在这场思考中已有机会检视过。即便不曾浸透与某位学生相伴的那项工作本身,它们也必定或多或少浸染了我本人周围的氛围或空气。自负,哪怕以世界上最“不起眼”的方式表达,总是指向一种封闭,一种对事物精微的本质及其美的麻木——无论那是“数学之物”,还是有生命的、我们本可以接待、鼓励、或者也可以从我们高处的优越俯视的人;对我们身边吹拂的气息及其毁伤性的效力——毁伤他人也毁伤我们自己——麻木不仁。

\subsection{2.~一种不义与无力的感觉}

\notehead{44$''$}{[这条注的出现,不遵守写作的时序]}

(5月10日)借着朋友所给的自由引证权——他的信中凡我认为有用的段落都可以引用——我在这里给出一段完整得多的引文\fnote{4}{(*)}{见前一部分“教学的失败(2)——或创造与自负”的第二条脚注,n\textsuperscript{o}~44$'$。},把那段被删削的引文放回它真正的语境:

\begin{quote}
“的确,75到80年间我非常孤立,除了向Verdier提过寥寥几个问题。但那段时间我并不怨你的昔日
学生们,因为没有人真正懂得这条纽带[按:即离散系数与连续系数之间的纽带]的重要性。一切都在
1980年10月改变了:人们发现了这条纽带对半单群的第一个头等重要的应用,即 Kazhdan--Lusztig
重数公式的证明——其中关键地用到了我们所论的范畴等价。这个等价在得到“Riemann--Hilbert对
应”的名字时,没有一句别的评说——毕竟,这再自然不过了!正是在那里,我明白了你的昔日学生
们不很明白数学的\textbf{创造}是什么,而你或许有一份责任。我至今仍有一种不义与无力的感觉。的确,
在他们那个年代,问题都已摆好。这个定理的应用数目,无论在平展拓扑的框架还是超越的框架里,
都很可观,却始终顶着“Riemann--Hilbert对应”的名字!我有种感觉,我的名字配不上这个结果——对\origpage{176}
许多人,尤其对你的昔日学生而言。但正如你从我的工作的引言里能清楚看到的,正是你的“对偶”
形式主义自然地引出了这个结果。不过同你一样,我并不为“可构造离散系数”与晶体系数(或完整
$\mathcal{D}$-模)之间这条纽带的将来担忧。显然,它无论在空间的上同调还是分析中,都适用于许许多
多的领域。”
\end{quote}

正是朋友的信中这一段,启发了我(除本注之外)后来的那道注“当值无名之辈与好上帝定理”。照这封
信的说法,我全然没有料到(如我在适当的地方所解释的),朋友身上那种“不义与无力的感觉”,所针
对的不仅是那种盲目的轻蔑态度——它\textbf{系统性地贬低}他的贡献(这种态度,在我某些昔日学生身
上,后来我也变得颇为熟悉了)——而且是一场真正的诈骗行动:干脆\textbf{变走}一条关键定理的父亲身
份。这一处境,仅仅八天前才向我揭示——关于这一点,见“不义——或归来的意义”及其后各注
(n\textsuperscript{o}s~75至80),它们集结在“讨论会——或Mebkhout层与反常”的标题之下。

\notehead{45}{}

由于我改变了圈子与生活方式,同旧日朋友们会面或其他接触的机会变得稀少。但这并没有妨碍
“拉开距离”的迹象以种种方式表现出来,彼此强弱不一。而在另一些人那里——如Dieudonné、Cartan、
Schwartz,实际上还有所有在我起步时如此善待我的“长者”——我没有感到丝毫这类东西。除他们之
外,我有种感觉:除极少数人之外,我旧日的朋友或学生在数学世界里,其与我之间的关系(无论有无机会表达),在我
退出那个曾经共同的圈子、共同的世界之后,都变得分裂了,带上了“矛盾的双重态度”。

\section{II.~孤儿们}

\subsection{1.~我的孤儿们}

\notehead{46}{[该注由第(I)部分《愚妄与更新》第八章“孤独的历险”第~50~节所召]}

\pmark\origpage{177} 我想借这个机会,在这里说几句关于数学观念与思想的话——在我带到光天化日之下的所有
这些东西当中,它们在我眼里(而且远远领先于其他)具有最广大的分量(注46$_{1}$)\fnote{5}{(*)}{读者可在注~n\textsuperscript{o}~46$_{1}$至46$_{9}$中找到对本注所巡视诸观念的若干更技术性的评说。此外,在我引入的\textbf{诸}观念之外,读者还可在注 n\textsuperscript{o}~88“尸身”里找到关于我著作中(在我“完全做到底”的那一部分著作之内)被其作者视为“执牛耳者”的那一部分的思考。}。首先
是五个紧密相联的关键观念,我将按其特殊性、丰富性(与深度)递增的次序,快速巡视一遍。

首先是同调代数中\textbf{导出范畴}的观念(见注48,原书第274页),以及把它用于一个“通用”的
形式主义——所谓“\textbf{六种运算的形式主义}”(即诸运算 $\overset{L}{\otimes}$、$Lf^*$、$Rf_!$、$R\mathcal{H}om$、
$Rf_*$、$Lf^!$)(注46$_{2}$)——
用于几何学中迄今引入的最重要的几类“空间”的上同调:“代数”空间(如概形、概形叠等)、解析空间
(复解析的,以及刚性的解析的及同类者)、拓扑空间(当然,还在等待各种“温和空间”的语境,肯定还
有其他许多,比如以小范畴的范畴(Cat)作为同伦模型的那一类……)。这一形式主义所涵盖的,既有
离散性质的系数,也有“连续”系数。

这一对偶形式主义及其无所不在性的逐步发现,是通过一场孤独、执着而严格的思考完成的,它从
1956年延续到1963年。正是在这场思考中,导出范畴的观念逐步成形,以及它在同调代数中所担当
的角色的理解也逐步成形。
这套形式主义所缺的,是我的上同调形式主义视野里对“离散系数”与“连续系数”之间那条可
以想见的纽带的理悟——超越局部系统及其以带可积联络的模或晶体模来解释的那个熟知情形之外
的纽带\origpage{178}。这条深刻的纽带,最初在复解析空间的框架里得到表述,是由Zoghman~Mebkhout
(将近二十年后)发现并建立的:以一边由“可构造”离散系数、另一边由“D-模”或“微分算子复形”
观念所构成的导出范畴来表述(见注46$_{3}$)。

将近十年间,由于得不到那些最有资格鼓励他的昔日学生的鼓励——用他们的关切和他们在与我
接触中积累的经验来支持他——Zoghman~Mebkhout在几乎完全的孤立中,坚持他那出色的研究。但这
并没有妨碍他揭示并证明一门正在普遍的冷漠中蹒跚诞生的新晶体理论的\textbf{两把钥匙定理}
\fnote{6}{(*)}{(6月7日)Mebkhout提醒我,在这两个定理之外,还应添上第三个,同样以导出范畴表述,即他(或许
不太恰当地)所谓的 D-模的“\textbf{双对偶定理}”——三者中最难的一个。关于Mebkhout的观念与结果及其
应用的概观,见 Le~Dung~Trang 与 Zoghman~Mebkhout, Introduction to linear differential Systems, Proc. of Symposia in Pure
Mathematics, vol.~40 (1983), part.~2, p.~31--63。}
——两者(这实在是个不好的兆头!)都以导出范畴表述:其一如方才所指,给出“可构造离散系数”与
晶体系数(满足某些“完整性”与“正则性”条件)之间的范畴等价(注48$'$);其二是“\textbf{那}条晶体
的整体对偶定理”,关于从一个光滑复解析空间(不必紧,这意味着额外的可观的技术困难)到一个点
的常值应用。这些是深刻的定理\fnote{7}{(**)}{(5月30日)第二个定理的证明,遇到超越语境惯有的技术困难,
需要求助于所谓“évétesques”技术(原书用语,此处暂保留法文);我猜它可以归入“困难”证明之列。第一个定理的证明是“显然”的——而且深
刻,动用了Heisuke~Hironaka奇点消解的全部力量。如我在“团结”一注(n\textsuperscript{o}~85)倒数第二段所提示的,定理
一旦提炼出来,“随便哪一个”知情的来者都能证明它。可对照“年轻人的势利——或纯洁的捍卫者”
(n\textsuperscript{o}~27)一注中Whitehead的评语。我写那道注的时候,仿佛在某种隐秘的先知的不声不响的口授之下,
还不曾料到现实会在多大的程度上超过我那些怯生生的、摸索式的提示!},为新照亮解析的与概形的
空间(目前为特征零)\origpage{179}的上同调,并许诺这些空间的上同调理论一场规模宏大的更新。它们最终为它们
的作者——在两次进入法国国家科学研究中心的申请被拒之后——换来一个研究员职位(相当于大
学的助理或助教职位)。

这十年间,没有人想到跟Mebkhout说一声“六种运算的形式主义”——我那些学生人人烂熟
\fnote{8}{(*)}{他们是在 SGA~4 与 SGA~5 的讨论班里第一手学到的,或经由 R.~Hartshorne的《留数与对
偶》(Residues and Duality)转手。}、却在任何地方都找不到整理成文文本的形式主义——而他正苦苦对付着超
越语境所带来的巨大的技术困难。他最终是去年从我口中得知它的存在的(以一份看来只有我
一人掌握的公式总表的形式),其时他正不辞辛劳、好意向我讲解他所做的一切;而我,早已不那么在
行上同调了……也没有人想到向他提示:也许先“转向”特征零的概形语境会更“合算”——在那里,
超越语境固有的困难消失,而理论的基本概念问题反而显得愈发清晰。也没有人想到向他指出(或者
哪怕注意到一件我自引入晶体之时就已知道的事
\fnote{9}{(**)}{(5月30日)不过我来得及把它忘了——直到去年与Mebkhout的第二次会面,才凭那次重逢之
力重新想起。(见“冥世重逢”一注,n\textsuperscript{o}~78。)}):光滑空间(解析的或概形的)上的“D-模”,不多
不少,正是“模的晶体”(只要对双方都不问“凝聚”与否的问题)——而后者是一个万能的观念,对带任
何奇点的“空间”与对光滑空间同样行之有效(注46$_{4}$)。

以Mebkhout所展现的才能(与那份少见的勇气),有一点对我说来再清楚不过:倘若处身于一种友善
的氛围,他建立特征零概形的晶体上同调语境中完整的“六种运算”形式主义,将毫无困难,而且会乐
在其中——这样一项宏大计划所需的全部本质观念(包括他自己的,Sato学派与我自己的),在我
看来本已集结在手。以他的才具,这不过是几年功夫的问题,正如平展上同调的一个通用形式主义的
发展是几年功夫的问题(1962--1965)——只要六种运算这条主线(外加换基的两把钥匙定理)已经在
手\origpage{180}。诚然,那是被同侪与见证者的热情与善意的潮流所托举的年月;而不是逆着那些大权
在握者的傲慢自足、逆流而上的劳作……

我来到我想谈的第二对观念:\textbf{概形},以及与它血缘最近的\textbf{topos}。后者是\textbf{site}观念的
更内蕴的版本——site 是我为了使“局部化”的拓扑直觉形式化而最先引入的。(“site”这个词,后来
是Jean~Giraud引入的;他也为使 site 与 topos 的观念具备全部必要的灵活性做了大量工作。)正是代数几
何的种种触目的需要,使我接二连三地引入了概形与 topos。这对观念,蕴含着对代数几何、算术,以及
拓扑的同样宏大的更新之潜能:通过一种\textbf{综合},把这些太久彼此隔绝的“世界”统一在一个共同
的几何直觉里。

以概形的观点、site(或“下降”)的语言,外加十二年关键的奠基工作(还不算我的学生们和其他
志愿者的劳作)对代数几何与算术的更新,二十年来已是既成之事:概形观念、概形的平展上同调观
念(即便还不是平展 topos 与平展叠的观念),终于进入了习惯,进入了公共遗产。

相反,那场同样会把拓扑囊括进来的宏大综合——尽管二十年來,它所需的本质观念与主要
技术工具在我看来已然集结待命\fnote{10}{(*)}{(5月15日)这些“本质观念与主要技术工具”已在 SGA~4
与 SGA~5 两场讨论班的宏大画卷中集结,时在1963至1965年间。这幅画卷的 SGA~5 部分,其成文与出
版所遭遇的离奇变故——十一年之后(1977年)以一种面目全非、满目疮痍的形状问世——为这幅宏
大视野落入“某种时尚”之手的命运,提供了一幅惊心的图像——或者不如说,落入某些我的学生之手:
正是他们最先确立了这种时尚(见下一条脚注)。这些变故及其意义,在最近四个星期的思考中逐步显
露,延续于“伙伴”、“清场”、“格格不入者”、“信号”、“翻转”、“沉默”、“团结”、“蒙骗”、“死者”、“屠杀”
诸注之中(n\textsuperscript{o}s~63$'$'$'$、67、67$'$、68、68$'$与84--88)。}——却还在等待它的时刻。\dmark{}\origpage{181}十五年来(自我告
别数学舞台以来),topos 这个丰饶的统一观念、这个强大的发现工具,被某种时尚\fnoteplainbody{11}{(*)}{(5月13
日)写完这几行之后的六个星期里,思考的继续显明:这种“时尚”最先是由我的某些学生确立的——恰
恰是那些最有条件把某种视野、一些观念与技术手段据为己有的人;他们选择了占有这些劳作的工
具,同时对催生这一切的视野、以及对这视野在其之内诞生的那个人,一并拒认。}打入所谓严肃
观念的冷宫。直到今天,对自己这门科学这种潜在的可观扩展、以及它所提供的新资源,有丝毫觉察的
拓扑学家仍然寥寥。

在这幅更新的视野里,拓扑学家日常摆弄的拓扑空间、微分流形等,连同概形(他听说过)与拓扑、微分或概形的叠(这些却无人谈论),不过是同一类非凡几何对象——\textbf{赋环 topos}(注46$_{5}$)——的诸多化身:拓扑、代数几何与算术
的种种直觉,汇聚于其中的“空间”,共同构成一幅几何视野。各种“模”叠——只要你睁着眼睛,就
会处处遇到(注46$_{6}$)——为它提供了数不清的醒目例子。对它们的深入研究,是深入这些几何对象(或
别的对象,如果有不是几何的对象……)的本质性质的一条头等主线:这些模叠所描述的,正是那些对
象的变动、退化与泛化的样态。然而这份丰富至今无人理会,因为使我们得以精细描述它的观念,不
在通行认可的范畴之列。

这场被拒认的综合\fnote{12}{(**)}{(5月13日)这场综合之遭“拒认”——无论其精神,还是使之成为可能的那
个关键观念——首先出自的,不是别人,恰恰是那项技术手段(概形语言与平展上同调理论的构造)的主
要使用者和受益者:终其全部著作,他受惠于这些手段。这就是Pierre~Deligne。凭他非凡的(源于其非凡
才能的)影响力,以及他相对于我的著作所占据的十分特殊的位置——他像是我的著作的一位隐含
遗产继承人——对我所引入的主要观念(概形观念与平展上同调除外)所设的那道不动声色的系统的
路障,收效极大,无疑在那场\textbf{埋葬}了这些观念的“时尚”的确立中扮演了头等角色;这些观念,已被
迫过了将近十五年的萎顿生活。他的著作,被这层暧昧深深地打上了烙印;这层暧昧,我在接续本注的思
考中第一次瞥见。(见“拒受遗产——或矛盾的代价”,注 n\textsuperscript{o}~47,原书第271页)我离去之后Deligne
著作中这道持续的障碍,这一鲜活而仍然含混的第一知觉,在整个关于\emph{葬礼}的思考中,以惊人的方
式得到了精确与确证;在那场葬礼中,我的朋友扮演着主祭的角色。}还带来一个意外的一面:\dmark{}\origpage{182}某些最
常见空间(注46$_{7}$)的熟知的同伦不变量(或更确切地说,它们的投射有限紧化),被发现配有无人料到的
算术结构,尤其是某些投射有限Galois群的作用……

然而,将近十五年来,在“大世界”里,谁要是贸然说出“topos”这个词——除非是开玩笑,或者他有
逻辑学家的托词——以居高临下的姿态鄙视他,竟成了合乎体面的做法。(逻辑学家是出了名的“与常人不同”,他们的怪癖人们
自会原谅……)以导出范畴的瑜伽来表达拓扑空间的同调与上同调,同样不曾进入拓扑学家中间:对他
们来说,Künneth 公式(对不是域的系数环)至今仍是一个由两条谱序列(或充其量一大串短正合列)
构成的系统,而不是某个合宜范畴中唯一的一个典范同构;换基定理(对真态射,或对光滑态射)他们同
样始终无视——而这些定理(在其近邻平展上同调的框架里)曾构成这门上同调“强力启动”的关键转
折(见注46$_{8}$,原书第270页)。当曾经为发展这门瑜伽出过力的人自己早就把它忘得一干二净,还对那
个想要使用它的可怜人冷眼相待的时候,我就不必惊讶了\fnote{13}{(*)}{(5月13日)后来的思考显明:处境开
始随1981年6月的吕米尼讨论会而改变:在那里,人们看到那些曾经“忘记”(或者不如说,埋葬了……)
这些观念的人,带着它们招摇过市——却依然对那位没有其人便绝不会有那场光辉讨论会的同一“可
怜人”冷眼相待。(关于这场难忘的讨论会,见注 n\textsuperscript{o}s~75与81。)}

第五个观念——也许比其他任何观念都更贴近我的心——是“\textbf{Motif}”。它与前四者不同之处在
于:“\textbf{那}个”好的 Motif 观念(哪怕只在基域上,更不必说任意基底概形),至今尚未成为某个令人
满意的定义的对象,即便为此承认人们会需要的一切“合理的”猜想也无济于事。或者不如说,分明
\dmark{}\textbf{那}条应当首先做出的“合理猜想”,就是\textbf{一条理论之\textbf{存在}}的猜想——它满足这样这样的
数据与性质,而对一个\origpage{183}“在局中”的人\fnote{14}{(*)}{(5月13日)我终于明白:直到今天,唯一(除我
之外)符合这个相当特殊的“多少在局中”的意义的人,是Pierre~Deligne。他有四年之久,一面听“我在
代数几何里所知道的那一点点”,一面日常地做我的 Motif 思考的心腹之交。的确,这些事我跟周围许多
别的同事也谈过,但看来没有一个足够“在行”,能消化那在我之内发展了数年的总体视野,或把
我的提示当作出发点,自己去发展一套视野与纲领(正如我自己当年从Serre的某些思想所给的两三
份“强印象”出发所做的那样)。也许我错了,但我觉得,关心代数簇上同调的人们,只要Deligne——他在
上同调中一言九鼎,而同时又是唯一被认为了解这些 Motif 底细的人——自己还在对 Motif 沉默不表,
就并不处于“把 Motif 当回事”的心理状态。(6月8日)核实之后发现,我最初的 Motif 思考可追溯到六
十年代初——因此它们持续了将近十年。},把它完全显式地陈述出来,本不是什么难事(而且会非常迷
人!)。其实在“离开数学”前不久,我离做这件事只差一步。

从某些方面看,处境与微分积分学英雄时代的“无穷小”颇为相似——不过有两点不同。首先,今天
的我们拥有建造精微数学理论的经验与一套有效的概念装备,那是我们的先辈所缺的。其次,尽管有这
些手段,而这一分明本质的观念出现已二十多年,却没有人屈尊(或敢于,尽管有人不屑于……)
撸起袖子,提炼出 Motif 理论的大致轮廓,正如我们的先辈直截了当地为无穷小计算所做的那样。然而
如今,对 Motif 就像当年对“无穷小”一样清楚:这些生灵是存在的;只要你关心代数簇以及这种簇的族
的上同调,尤其关心它们的“算术”性质,它们就在代数几何中随时显形。也许比对我谈到的那四种观念
更甚,“Motif”——所有观念中最特殊、最丰富的一个——系于一大堆形形色色的直觉:绝非含混,而往
往可以\dmark{}以完美的方式\origpage{184}表述出来(必要时,偶或承认一些 Motif 式的前提)。这些“Motif 式”
直觉中最迷人的一个,在我这里,是“Motif Galois群”:它在某种意义上,使人得以在(绝对意义的)域
与有限型概形的投射有限Galois群上“配一个 Motif 结构”。(为给这一观念以确切意义所需的技术工
作——以那些为 Motif 观念提供暂时基础的“前提”来表述——是在Neantro~Saavedra论“塔纳卡
范畴”的论文里完成的。)

时下的共识,对 Motif 观念比对它的三个受难兄弟(姐妹)(导出范畴、所谓“六种运算”的对偶形式
主义、topos)要稍许客气些:它没有被干脆斥为“涂鸦”\fnote{15}{(*)}{如我在前面一条脚注里提过的,导出范
畴三年前曾享受一场锣鼓喧天的掘墓礼(只是没有提我的名字)。topos 与六种运算还在等待它们的时
刻,Motif 也是——除了两年前被部分开棺出土、换上另一位父亲的那一小块(见注n\textsuperscript{o}s~51、52、59)。(5月13日)},但实际上这毫无区别:只要没办法“定义”一个 Motif、没法“证明”点什么,严肃的人
们就只能绝口不谈(怀着最大的遗憾——这话是听说了的,可人要么严肃要么不严肃……)。当然,只要
人们宣称谈它本身都不严肃,那就永远别指望建造 Motif 理论、就它“证明”点什么!

可那几个在局中的人(以及制造时尚的人)自己心里非常清楚:借助那些保密的前提,可以证明许多
东西。也就是说,从今天起——实际上,自从这一观念在Weil猜想的余波中出现起(它们倒是被Deligne
证明了,这总算是个加分项!)——\textbf{Motif 的瑜伽}就实实在在地存在。但它有着一门\textbf{秘密科学}的
地位:入门者肯定寥寥\fnote{16}{(**)}{(5月13日)我现在似乎终于明白,截至1982年,所谓“极少数入门者”,其实仅有Deligne一个人。的确,他披露了这门“秘密科学”中透过某些重要结果显现出来的部分;这些结果包含在这套瑜伽之中,他随着自己能够证明它们而逐一公布,借此获取功劳,同时隐瞒仍然保密的灵感来源。可是,十五年来竟无人行动起来,最终投身于一门规模宏大的Motif理论,这分明说明,我们的时代远远没有无穷小计算英雄时代那种大胆的活力!}。它尽管“不严肃”,却仍然\dmark{}\origpage{185}使这些罕见的入门者能够在许
许多多上同调的处境中说出“人们有权从中期待什么”。由此衍生出许许多多的直觉与局部猜想;借助“瑜伽”提供的理解之光,人们有时事后便能用手头现成的手段处理它们。Deligne的若干项工作受到这套瑜伽启发\fnote{17}{(*)}{(5月13日)我终于对文献稍有了解。如今我看到,Deligne的全部著作都扎根于这套瑜伽。我的文献抽样以及其他交叉核对使我推测:在他的全部著作中,唯一提及这一来源的地方,是1970年《霍奇理论I》中的一句简短文字,将我与Serre一并提到。(见注n\textsuperscript{o}s~$78'_{1}$、$78'_{2}$。)},其中一项(如果我没有弄错)是他最早发表的工作:证明代数簇的射影光滑态射的Leray谱序列退化(为证明所需,取特征零)。这个结果是由有关“权”的考虑所
提示的;这类考虑具有算术性质——这些正是典型的“Motif 式”考虑,只不过被表达在 Motif 的“几何”的术语里。Deligne证明
这个命题靠的是Lefschetz--Hodge理论,而对那动机(49)(如果我记得不错)只字不提——没有它,谁也不会想
到去怀疑如此难以置信的东西!

Motif 的瑜伽,恰恰最初正是从我得自Serre的这套“权的瑜伽”之中诞生的\fnote{18}{(**)}{我从Serre那里(六十
年代初?)得到的,是一个观念或直觉:有一个重要的东西要理解!它起了一次初始推动的作用;此后几
年的思考紧随其后,先围绕一个“权的瑜伽”,不久又围绕一个 Motif 瑜伽。}。是他让我领略了“Weil猜想”的全部魅力(后来有了“Deligne定理”)。他向我解释了(在假定所论特征下存在奇点消解的
前提下)人们如何能凭权的瑜伽,给任意域上每个概形(不必光滑也不必真)配上“虚拟 Betti 数”
——这一点当时给我极深的震动(注46$_{9}$)。我相信,正是这个观念,是那场围绕权的瑜伽的思考的起
点;那场思考在随后几年里,一直在我奠基成文的任务的边缘上继续。(这也是我在七十年代重新拾起
的思路,借助任意基底概形上的“虚Motif”观念,旨在至少为虚Motif建立“六种运算”的一个形式主义
\origpage{186}。)如果说这些年里我一直对Deligne(我的得天独厚的对话者)以及凡愿听的人\fnote{19}{(*)}{(4月10日)我觉得,Deligne是唯一“听进去”的人——而且他小心地把听进去的东西的专属特权留给了自己。诚然,另一方面,写下这些收尾文字时,我“落在”了事件后面:两年前,Motif瑜伽已有一场部分的开棺出土,却对我的角色不着一字!对此可参见由一项意外发现所引发的注n\textsuperscript{o}~50、51、59——那项发现,对这场持续了十二年的下葬的意味,投下一束(至少对我而言)意外的光。直到那时,我还只是相当混乱地觉察到某种下葬的存在,却不曾费心凑近去看看……}讲这个
Motif 瑜伽,那当然不是为了叫他们把它当作一门只属于他们自己的秘密科学。($\Rightarrow$
注47,原书第271页)

\notehead{46$_{1}$}{}

除了以当年我对Riemann--Roch定理所给的表述(以及我据此给它的两种证明),还有它的各种变体而
引入的观念与观点之外,我几乎不作例外。如果我记得不错,这类变体是在 SGA~5 讨论班1965/66年
的最后一次报告里展开的;而那次报告连同同一讨论班的其他若干报告,都遗失得一干二净。最有意思的一种,我看是可构造离散系数的一条变形;我不知道它此后在文献里是否得到过显豁\fnote{20}{(**)}{(6月6日)我重新找到了它——以一个近旁的形态,顶着“Deligne--Grothendieck猜想”这个动听的名字——在MacPherson1974年发表的一篇文章里。详情见注n\textsuperscript{o}~87$_{1}$。}。值得注意:这一条同样容许一个“Motif”变形,它实质上等于断言:对互异的素数$\chi$(与剩余特征互素),与$\chi$-adic可构造层相配的“示性类”(在正则概形$Y$的Chow环里),当这些层出自同一个“词”——同一个motif(比方说,都是对一个给定的$f:X\to Y$的$R^if_!(\mathbb{Z}_\chi)$)——时,全都相等。

\notehead{46$_{2}$}{}

这个形式主义,可以看作同调中一种“整体对偶”形式主义的某种精粹;以其最“有效”的形
状——撇清了所论“空间”与应用的任何光滑性前提,或态射的真性——它还需要用一个局部对偶形
式主义来补足:在其中人们区分所谓“\textbf{对偶化}”的对象或“复形”\dmark{}\origpage{187}(在 $Lf^{!}$ 运算下稳定
的观念)——并由此有一个“\textbf{双对偶定理}”(以 $R\mathcal{H}om$ 表述),适用于满足合宜的有限
性条件(关于次数,以及局部层面的凝聚性或“可构造性”)的系数。当我说“六种运算的形式主义”时,
我通过后者所暗指的,是对偶性——其“局部”与“整体”两个方面。

深入理解上同调中对偶性的第一步,是六种运算的形式主义在第一个重要情形——诺特概形及上同调凝聚的模复
形——中的逐步展开。第二步是发现(在概形的平展上同调中)这个形式主义对离散系数同样适用。这
两个极端情形,足以确立这种对偶性之\textbf{无处不在}的信念——对Poincaré型“对偶”的一切几何处
境的信念;这一信念,此后又得到Verdier、拉米、吕盖等人的工作的印证。当那股十五年
来压在这套形式主义上、阻碍其大规模使用的\textbf{阻碍}崩解之后,它必将在任意类型的系数上再
度得到确认。

这种无处不在,在我看来是一个分量极重的\textbf{事实}。它使我们不能不感到Poincaré对偶与Serre对偶之间有一种深刻的统一性——Mebkhout最终以所需的一般性确立了它。这种无处不在,使“六种运算的形式主义”成为
同调代数中理解“全方位”上同调对偶现象的基础结构之一\fnote{21}{(*)}{有兴趣的读者,可在本卷的附录里找到这一形式主义的一个概要。}。至于这套颇为精微的结构竟未被行业内
的人显式陈述出来——连“三角范畴”的“好”观念(Verdier的版本仍是一种十分暂时且不充分的形状)
也不曾有人提出——这一事实毫不改变什么;拓扑学家乃至偶或关心上同调的几何学家,仍然竞相无视
这套对偶形式主义的存在,以及奠基它的导出范畴的语言,也同样无济于事。

\notehead{46$_{3}$}{}

D-模与微分算子复形的观点,是由Sato修引入的,最初由他本人及其学派发展,所循的\dmark{}\origpage{188}观点(据我的理解)与Mebkhout所遵循的相当不同;而后者的进路更贴近我的进路。

“可构造性”的诸种观念——对(复解析的、实解析的、分片线性的)“离散”系数——据我记忆,是由我在五十年代末前后首次提炼出来的(几年后我在平展上同调的框架里又重提了它们)。当时我提过一个关于这种观念之稳定性的问题:它对实或复解析空间之间的一个真态射的高阶直像是否稳定——并且不知道这种稳定性在复解析情形是否已被确立\fnote{22}{(*)}{(5月25日)它已由 J.L.~Verdier确立,见“好的参考文献”一注(n\textsuperscript{o}~82)。}。在实解析情形,我当年设想的观念本来就不是对的,因为没有Hironaka的实亚解析集观念——它带有经由直像的本质先行稳定性。至于 $R\mathcal{H}om$ 一类的局部运算,当时清楚的是:在特征零优秀概形的框架里(利用Hironaka奇点消解)确立可构造系数之稳定性的论证,在复解析情形原封不动地成立,双对偶定理亦然(见 SGA~5~I)。分片线性的框架里,自然的稳定性与双对偶定理是“容易的练习”——我当年很高兴做它们,以此检验对偶形式主义的“无处不在”,那正是在平展上同调起步之时(其一大意外恰恰是发现了这种无处不在)。

回到半解析的情形:这一方向上稳定性定理(经由六种运算的可构造系数)的“好”框架,分明是“温和空间”的框架(见《纲领概要》第5、6节)。

\notehead{46$_{4}$}{}

当然,D-模的观点,连同 $\mathcal{D}$ 是一个凝聚环层的事实,为模的晶体揭示了一种比我所惯用的更隐蔽的“凝聚性”观念;它在不必光滑的空间(解析的或概形的)上仍有意义。称之为“\textbf{M-凝聚性}”($M$ 即Mebkhout)才算公道。于是,对一个多少“在局中”(而且健全的数学直觉丝毫不爽)的人来说,以下应当相当显然:那条推广了光滑情形“微分算子复形”的\dmark{}“好的系数范畴”\origpage{189},不应是别的,正应是模的晶体的“M-凝聚”导出范畴(一个晶体复形称为 M-凝聚的,若其各上同调对象如此)。它在没有光滑性假设下仍有合理的意义,并且应当同时涵盖寻常的“连续”(凝聚)系数理论与“可构造”离散系数理论(对后者引入合宜的完整性与正则性假设)。倘若我的视野是对的,Sato--Mebkhout理论相对于先前已知的晶体语境而言,其两个新的概念配料,就是模的晶体的这种 M-凝聚性观念,以及关于 M-凝聚晶体复形的(更深刻的)完整性与正则性条件。这些观念一旦到手,第一桩本质的工作,就是在晶体语境中发展六种运算的形式主义,以便涵盖二十多年前我已发展的两个特殊情形(寻常凝聚与离散)——而我某些前学生上同调学家,早就为了想必更重要的任务把它们忘得一干二净……

Mebkhout其实早已通过研读我的著作,最终得知“晶体”观念的存在;他也感到他的观点(至少在特征零)应当为这个观念提供一条好的进路——但这点建议落入了聋子的耳朵。从心理上说,很难设想他会投身于那桩势在必行的浩大奠基工程:他所处的,恰是那些以上同调权威自居的人的高傲冷漠的气候之中,而他们本是最有资格鼓励——或者说打消……——的人。

\notehead{46$_{5}$}{}

(5月13日)这里说的,主要是以一个\textbf{交换局部环}赋环的 topos。以拓扑空间上这样一个环层的数据来描述“簇”结构的想法,最先由 H.~Cartan引入,又被Serre在其经典工作 FAC(凝聚代数层)中采用。正是这项工作,为引我走向“概形”观念的思考提供了初始推动。Cartan的进路(经Serre采用)所缺的,为涵盖迄今出现的所有类型的“空间”或“簇”所缺的,正是 topos 的观念(即恰好是这样的“东西”:在其上“集合层”的观念有意义,并且具备那些熟悉的性质)。

\notehead{46$_{6}$}{}

\dmark{}作为 topos 的其他出色的例子——它们不是寻常的空间,\origpage{190}而且以“通行”观念似乎也没有令人满意的替代品——我要指出:一个拓扑空间对一局部等价关系所得的商 topos(比如流形的叶状结构,此时商 topos 甚至是一个“多簇”(multiplicité),即局部是一个簇);对几乎任何一种数学结构(至少那些“以有限投射极限与任意归纳极限来表述的”)的“分类”topos。当人们取一种“簇”结构(拓扑的、微分的、实或复解析的、Nash 的,等等,甚至取一个给定基底上光滑的概形结构)时,每种情形都会出现一个特别诱人的 topos,配得上“普遍簇”(所论那一种的)之名。它的同伦不变量(尤其它的上同调——配得上所论簇类的“分类上同调”之名)早就应当被研究和熟知;但到现在,这条路还根本没有被走过……

\notehead{46$_{7}$}{}

这说的是这样一些空间 $X$:其同伦型被“以自然的方式”描述为一个复代数簇的同伦型。后者于是可定义在复数域的一个子域 $K$ 上,$K$ 是素域 $\mathbb{Q}$ 的有限型扩张。投射有限Galois群 $\mathrm{Gal}(\overline{K}/K)$ 于是自然地作用在 $X$ 的投射有限同伦不变量上。往往(比如 $X$ 是一个奇数维同伦球面时)可取 $K$ 为素域 $\mathbb{Q}$。

\notehead{46$_{8}$}{}

(5月13日)当我在Serre的论文 FAC 里学到代数几何的最初几课(正是它把我“触发”向概形)时,换基这个观念本身在代数几何里还几乎无人知晓,除非在换基域这个特殊情形。随着概形语言的引入,这一运算变成了代数几何中最常用的运算,它在每个时刻都会出现。这一运算在拓扑学里至今仍几乎无人知晓(除一些很特殊的情形),在我看来,是拓扑学相对于来自代数几何的思想与技术的孤立的一个典型征兆(在许多其他征兆之中),也是“几何”拓扑学那套不充分基础的一个顽固遗产。

\notehead{46$_{9}$}{}

\dmark{}(6月5日)Serre的想法是:\origpage{191}对一个域 $K$ 上任意有限型概形 $X$,应能配上整数
$$h^i(X)\quad(i\in\mathbb{N})$$
他称之为“虚拟 Betti 数”,使得:
\begingroup\renewcommand{\labelenumi}{\textbf{\alph{enumi})}}
\begin{enumerate}
\item 对 $Y$ 一个闭子概形、$U$ 其补开集,
\[
h^i(X) = h^i(Y) + h^j(U).
\]
{\small [译按:原书此式右端印作 $h^j(U)$,疑应为 $h^i(U)$。]}
\item 对 $X$ 射影光滑,有
\[
h^i(X) = X\text{ 的第 }i\text{ 个 Betti 数}\text{。}
\]
{\small [译按:原书左端印作 $h(X)$;此处按前后定义保留译稿的 $h^i(X)$。]}
\end{enumerate}\endgroup
(比如经由与 $k$ 的特征互素的 $\chi$-进上同调定义)。如果承认 $k$ 上代数概形的奇点消解,那么立刻可见诸 $h^i(X)$ 由这些性质唯一确定。这样一个映射 $X\to(h^i(X))_{i\in\mathbb{N}}$($k$ 固定)的\textbf{存在},利用具真支撑的上同调形式主义,可本质上归约为基域有限的情形。在 $\mathbb{Q}_\chi$ 上有限维的、$\mathrm{Gal}(\overline{k}/k)$ 连续作用的向量空间的“Grothendieck 群”里工作,取 $X$ 的 $\chi$-adic(具真支撑的)Euler--Poincaré特征在该群中的像,$h^i(X)$ 即指 $EP(X,\mathbb{Q}_\chi)$ 的“权 $i$ 分量”的虚拟秩,其中权的观念由Weil猜想导出,外加奇点消解的一个弱形式。即使没有消解,Serre的想法也因Weil猜想的强形式(Deligne在“Conjectures de Weil II”中建立)而得以实现。

我沿着这条路做过启发式的思考,通向“虚拟相对概形”的六种运算形式主义——基域 $k$ 被换成一个多少任意的基底概形 $S$——以及这类($S$ 上有限表现的)虚拟概形的各种“示性类”观念。于是,我(为简化回到基域的情形)被引向考虑比Serre的更细的整数不变量,记作 $h^{p,q}(X)$,满足与上面 a)、b) 类似的性质,并且由通常的公式
$$h^i(X)=\sum_{p+q=i}h^{p,q}(X)$$
给出Serre的虚拟 Betti 数。

\subsection{2.~拒受遗产——或矛盾的代价}

\notehead{47}{[该注直接续注46,见 13.2.1]}


\pmark 值得注意的是,我方才巡视的五种观念之中,有四种(恰恰是那些被当作“不严肃之事”的观念)关系到上同调,\origpage{192}而且首先是\textbf{概形与代数簇的上同调}。不管怎样,这四种观念全都是代数簇的上同调理论的需要向我提示的——先是连续系数,后是离散系数。这就是说,在1955至1970这十五年间,我诸项工作的一个主要动机与一条不变的主导旋律,正是代数簇的上同调。

引人注目的是,这也正是Deligne至今仍视为其灵感主要来源的主题——如果我相信去年IHES小册子就此所说的话\fnote{23}{(*)}{(5月12日)然而,我刚发现,那本小册子里没有任何东西会让读者疑心:我的著作竟与代数簇的上同调、或与任何别的东西的上同调有什么相干!对此可参见即日所写的“葬礼悼词(1)——或恭维”一注(n\textsuperscript{o}~98)。所论小册子,即“有益的挣脱”一注(n\textsuperscript{o}~42)脚注里提到、并在刚提到的“葬礼悼词”注中得到稍细考察的那一本。}。我得知此事时颇为惊讶。诚然,当Deligne(在他关于Ramanujan猜想的漂亮工作(译按:原书人名拼作Ramanuyam,此处按通常拼写订为Ramanujan.)之后)发展了他对Hodge理论的卓越推广时,我还“在场上”,而且完全“在行”。那首先——姑且从这里开始!——对他和对我都一样,是朝着在复数域上以构造性方式建立Motif观念迈出的第一步。在我1970年“转向”之后的最初几年,我自然也听到了Deligne证明Weil猜想的消息(这同时证明了Ramanujan猜想),还有紧随其后的正特征下的“蛮横的”Lefschetz定理。我对他的期待不比这更少!我甚至确信,他应当同时证明了那些\textbf{标准猜想}——那是我在六十年代末提出的,作为(至少)在一域上奠基“半单”Motif观念的第一步,并把这些Motif的若干预期性质翻译成$\chi$-进上同调性质与代数闭链群性质的语言。Deligne后来告诉我,他对Weil猜想的证明肯定无法用来证明那些(更强的)标准猜想,而且他对如何着手也毫无想法。这应当已有十年左右了。从那时起,我再没有听说,在对代数簇上同调之“Motif”(或“算术”)方面\dmark{}\origpage{193}的理解上,发生过其他真正决定性的进展。了解Deligne的手段,我曾默然断定他的主要兴趣必然已转向别处——故而我读到情况并非如此时,深感惊讶。

有一点在我看来是无可怀疑的:至少二十年以来,若不同时在某种程度上扮演“Grothendieck的继承人”,人们已几乎不可能在代数簇上同调的理解上做出大面积革新的工作。Zoghman~Mebkhout已经切身领教了这一点,而(在某种意义上)Carlos~Contou-Carrère亦然——他很快明白,换题目才是上策(注47$_{1}$)。在那些人们无法回避的事情之中,恰恰首推在尽可能贴近Motif的各种系数语境里(这些Motif目前扮演着一种理想“地平线”的角色)发展那著名的“六种运算形式主义”:特征零的晶体系数(沿着Sato学派与Mebkhout的路线,做成“Grothendieck风味”),或特征$p$的晶体系数(主要由Berthelot、Katz、Messing以及一群显然干劲十足的更年轻的研究者研究);Deligne式的“分层pro-模”(promodules stratifiés)(它们显现为凝聚D-模或D-凝聚晶体这一“ind”观念的对偶化变体,即“pro”变体);最后是“Hodge--Deligne”系数(它们看来与Motif一样好,只除了一点:其定义是超越的,且局限于复数域上有限型的基底概形)……在另一端,则摆着这样一个任务:把Motif观念本身从笼罩着它的迷雾中(且不无原因……)提取出来,并且如果可能,去攻克像“标准猜想”这样精确的问题。(关于后者,我曾设想过,除其他途径外,发展一域上射影光滑簇的“中间Jacobi簇”理论,作为或许能得到迹的正性公式的一个办法——那公式是标准猜想的一个本质配料。)

直到我“退出数学”的那一刻,这些任务与问题还在我手里灼烧——滚烫而多汁的东西,其中任何一件、在任何时刻都从不曾在我面前显形为一堵“墙”、一个终止点\fnoteplainbody{24}{(*)}{(5月25日)然而,那本大名鼎鼎的庆典小册子里,一管我自信认得出的匿名手笔,却正是在友好地暗示这一点。参见“葬礼悼词(2)”一注,它接着前一条脚注所引的“葬礼悼词(1)”。}。它们是一个灵感与实质\dmark{}\origpage{194}永不枯竭的源泉——在那里只需朝涌出的地方一拉(而且它“四处涌出”!),如期而至的与不期而至的都会到来。以我这般有限的手段,但只要在工作中不分心,我清楚地知道,只要肯下功夫,一天、一年或十年,人能做成多少事。我也知道——因为我曾在一个他不分心的年代亲眼见过——Deligne的手段是什么,只要他愿意投入,他一天、一周或一个月能做成什么。可是没有人,连Deligne也不能,长此以往地一面从高处俯视那归根结底需要探查的对象本身,连同其语言与整套由某位先行者为这一目的而发展的工具库(何况还是在后者的协助之下,更不必说众多同样参与了劳作的人……),一面做出丰饶的、深刻革新的工作(注59)。

我也想到$M_{g,\nu}$(于$\mathrm{Spec}\,\mathbb{Z}$上)的“Deligne--Mumford”紧化——它参数化亏格$g$、带$\nu$个标记点的连通光滑代数曲线。引入这些对象\fnote{25}{(*)}{载 Pub.~Math.~36, 1969, 第75--110页。评论见注n\textsuperscript{o}~63$_{1}$。},是为了解决在全特征下证明模空间$M_{g,\nu}$的连通性这一问题,办法是从特征零做特化论证。在我看来,这些对象$M_{g,\nu}$(连同群$\mathrm{Sl}(2)$)是我在数学中遇到过的最美、最迷人的对象(注47$_{2}$)。单单它们的存在,带着如此完美的性质,在我眼中就是一桩奇迹(而且是得到了透彻理解的奇迹),其分量远非那桩本待证明的连通性事实可比。对我而言,它们以精髓的方式蕴含了代数几何中最本质的东西,即(几乎所有)一切代数曲线(在一切可以想象的基底域上)的全体——而这些曲线恰恰是一切其他代数簇的终极建筑材料。但所论这类对象——“$\mathrm{Spec}\,\mathbb{Z}$上真且光滑的多重簇”——仍逸出“通行”范畴之外,也就是说,逸出那些人们(出于从不检视的原因)肯于“承认”的范畴之外。常人谈起它们至多闪烁其词,而且一脸抱歉,仿佛还在做什么“泛泛的无稽之谈”,尽管人们小心地会说“stack”或“叠”,以避讳“topos”或“多重簇”这个禁忌的字眼。毫无疑问,这正是这些独一无二的珍宝自十多年前被引入以来,除了我本人在一些未发表的研讨班笔记\dmark{}里之外,再无人(就我所知)研究或利用的原因。取而代之的是,\origpage{195}人们继续要么用“粗糙的”模簇空间,要么用模多重簇的那些有希望是真概形的有限覆盖——然而无论前者还是后者,都不过是它们所自出的那些完美珍宝的某种苍白而跛脚的影子,而后者实际上仍遭放逐……

Deligne的四项工作——关于Ramanujan猜想、关于混合Hodge结构、关于模多重簇的紧化(与Mumford合作)、关于Weil猜想——每一项都是我们对代数簇之认识的一次革新,因而同时也是一个新的起点。这些奠基性的工作相继出现于几年之间(1968--73)。然而将近十年过去,这些巨大的界标并未成为向着预感与未知发起新一轮冲击的跳板,成为更广泛革新的手段。它们通向的是一种阴郁停滞的局面(注47$_{3}$)。这肯定不是因为十年前各家各人手头的那些“手段”像被施了魔法般消失了;也不是因为我们触手可及的那些美好事物突然化为乌有。而是:世界美是不够的——还得有人肯为之欢喜……

\notehead{47$_{1}$}{}

我在这里想到Contou-Carrère在五六年前一个颇有希望的开端:相对局部Jacobi簇的理论,它们与整体Jacobi簇(所谓“广义Jacobi簇”)的联系——对任意概形上的光滑且不必真的曲线概形——以及与Cartier的交换形式群及典型曲线理论的联系。除去Cartier的一次鼓励性回应之外,最有资格赏识Contou-Carrère第一篇注记的人们,对它的接待如此冷淡,以致作者从此绝口不提他储备着的那第二篇,并赶紧换了题目(尽管如此仍未能躲过其他的噩运)\fnote{26}{(*)}{(6月8日)见“棺材3——或有点太相对的Jacobi簇”一注(n\textsuperscript{o}~95)的子注(951)。}。

是我向他建议了局部与整体Jacobi簇这个题目的,作为迈向一个纲领的第一步——那纲领可上溯至五十年代末,尤其指向一个任意维的“阿代尔”对偶复形理论,由局部Jacobi簇构成(对应局部环\dmark{}\origpage{196}任意维的情形),以类比于诺特概形的剩余复形(由其一切局部环的对偶化模构成)。我这一上同调对偶纲领的这一部分(连同其他部分),由于六十年代涌入的种种显得更为紧迫的任务,多少被搁置进了被遗忘的角落。

\notehead{47$_{2}$}{}

确切说来,是所有这些多重簇的族所嵌入的那座“Teichmüller塔”,以及这座塔以基本群胚表述的离散或投射有限范式,构成我在数学中遇到过的最丰富、最迷人的单一对象。群$\mathrm{Sl}(2,\mathbb{Z})$,连同其投射有限紧化的“算术”结构(即$\mathrm{Gal}(\overline{\mathbb{Q}}/\mathbb{Q})$在其上的作用),可视为这座塔的“投射有限版本”的主要建筑材料。译按:原书此处群名的第二个字形似 $\chi$,这里依前文 $\mathrm{Sl}(2)$ 的记号书写。关于此可参见《纲领概要》中的指示(且待《数学反思》诸卷中专论这一主题的卷帙)。

\notehead{47$_{3}$}{}

这一“阴郁停滞”的观察,并不是一个对近十年间围绕概形与代数簇上同调的主要事件了如指掌的人深思熟虑的意见。这是一个“局外人”的单纯的整体\textbf{印象},它是我1982与1983年间同Illusie、Verdier、Mebkhout的谈话与通信中(除其他来源外)得出的。这一印象当然在许许多多方面有待细化。比如,Deligne1980年发表的“Conjectures de Weil II”,就代表一个实质性的新进展,虽说就其主要结果而言算不上一桩意外。看来特征$p>0$的晶体上同调也有进展,更不必说围绕相交上同调的那阵“抢滩”——它终于迫使某些人(尽管不情愿)回到导出范畴的语言,甚至使他们想起那些久遭否认的父系渊源……

\section{III.~时尚——或名人的生活}

\subsection{1.~本能与时尚——或强者的法则}

\notehead{48}{[该注由注46所召,原书第265页]}

众所周知,导出范畴的理论归于 J.L.~Verdier。在他着手\origpage{197}我提议的那项奠基工作之前,我只是以启发的方式、以这些范畴的一个临时定义(后来证明是对的)同它们打交道,对其内部本质结构也只有一个同样临时的直觉(这一直觉在预期的语境里被证明技术上是错的:“映射锥”并不\textbf{函子地}依赖于导出范畴中定义它的那个箭头——它只在相差一个不唯一同构的意义上定义它)。我在五十年代末发展的凝聚层的对偶理论(即凝聚框架里的“六种运算”形式主义)\fnote{27}{(*)}{其中还缺一个对非真态射的$Rf_!$运算(具真支撑的上同调);它是六七年后由Deligne引入的,有赖于他引入的凝聚pro-模语境——在我看来这是一个重要的新观念(后来被他成功地纳入其分层pro-模理论)。},其全部意义唯有在关于导出范畴观念的一项奠基工作之上才得以显豁,而那项工作是Verdier后来完成的。

Verdier学位论文的文本(迟至1967年才答辩),约二十页,在我看来是迄今所写的对导出范畴语言最好的导引,它把这一语言安置在其本质使用的语境之中(其中好几项使用出自Verdier本人)。它只是一项正在撰写的工作的导言,而那项工作最终也是后来才写成的。我可以自诩——即便不是唯一的一个——也是极少数能够作证曾亲手捧读过这部著作的人之一:仅凭这篇导言,其作者被授予国家博士头衔的正当性就得到了确立!这部著作(或者曾经是——我不知道是否还有一份副本存于某处……)是迄今唯一一部按导出范畴的观点给出同调代数的系统奠基的文本。

\pmark 也许只有我一个人在惋惜:无论是那篇导言还是那奠基的正文本身,都未曾\origpage{198}发表\fnote{28}{(*)}{(5月25日)在这些文字写成之后,我发现Verdier学位论文的最初胚芽——写于1963年,即答辩前四年——最终于1967年发表。对此可参见“同伙”与“赊账的学位论文与全险”两注,n\textsuperscript{o}~63$'''$与81。}——以致使用导出范畴语言所需的本质技术行装,散见于文献中三个不同的地方\fnote{29}{(**)}{这三处是:Hartshorne那场著名的凝聚对偶研讨班,其中有我在五十年代下半叶发展的对偶理论至今唯一发表的部分;SGA~4里Deligne的一两个报告;Illusie那部浩繁的学位论文中的一两章。}。这样一个系统参考文本的阙如——其分量本可比肩Cartan--Eilenberg的经典著作——在我眼中,既是1970年我退出数学舞台之后导出范畴形式主义所遭冷落的一个典型\textbf{原因},也是一个典型的\textbf{征兆}。

诚然,早在1968年就已显露(在SGA~5中发展的一种上同调迹理论的需要之际):原初形态的导出范畴观念,以及相应的三角范畴观念,对某些需要而言是不够的,更深彻的奠基工作尚待去做。Illusie在其学位论文中引入“滤导出范畴”,沿这一方向迈出了有用但仍属谨慎的一步(主要为了迹的问题)。1970年我的离去,似乎成了同调代数基础、乃至与之血肉相连的Motif理论基础的全部反思戛然而止、永不再续的信号(注48$_{1}$)。然而就前者而言,为大面积奠基所需的全部本质观念,似乎在我离去之前的年头里就已到手(注48$_{2}$)。(包括“求导者”(dérivateur)这个关键观念,或曰“导出范畴制造机”——它看来是至今所遇到的诸三角范畴之下共同的那个更丰富的对象;这一观念将近二十年后才在一个非加性框架里得到少许发展,载《追寻栈》第二卷的一章。)此外,尚待完成的奠基工作有相当一部分已由Verdier、Hartshorne、Deligne、Illusie做成,这些工作可以原封不动地用于一项综合,把已获的观念纳入“求导者”这个更开阔的视野。

\pmark 诚然,这十五年来\fnote{30}{(*)}{(5月24日)这“逝去的十五年”应当细化——参见注n\textsuperscript{o}~47$_{3}$,以及内容更详的“赊账的学位论文与全险”一注(n\textsuperscript{o}~81)。}对导出范畴观念本身的冷落——在某些人那里,它近于对一段过往的否认——正顺应了某种时尚:\origpage{199}这种时尚对一切奠基的反思都投以轻蔑的目光,无论其多么紧迫\fnote{31}{(**)}{(5月25日)关于这一时尚之出现与延续背后的诸种力量,可参见“掘墓人——或全体教众”一注,n\textsuperscript{o}~97。}。但从另一方面看,我也再清楚不过:平展上同调的发展——今天“人人”都在不加思索地使用它(哪怕只是通过已故Weil猜想的途径……)——若没有导出范畴、六种运算、site与topos的语言(它们最初正是为这一目的而发展)所构成的概念行装,更不必说SGA~1与SGA~2,是绝无可能的。同样再清楚不过的是:今天代数簇上同调理论中可以观察到的那种停滞,若非我昔日某些学生在那些年月里宁肯追随一场他们自己正是始作俑者的时尚、而不肯听从其作为数学家的健全本能——那场时尚久已借着他们的支持获得了法律的效力——原本既不可能出现,更不可能安营扎寨。

\notehead{48$_{1}$}{}

同样的话(带有某些保留)也可用于我那整个代数几何奠基纲领:它只有很小一部分得到实现,而且随我的离去戛然而止。这一停顿给我的震动,尤在于对偶纲领——我视之为特别多汁的部分。Zoghman~Mebkhout的著作,顶风冒雨地继续着,却仍处在这条纲领的脉络里(并因料想不到的观念的汇入而焕然一新)。1976年Carlos~Contou-Carrère的工作(注47$_{1}$,原书第273页已提及)亦然——那是他谨慎地无限期搁置了的工作。此外还有一项关于曲面fppf上同调对偶的工作(Milne)。我所知者,尽于此矣。

诚然,我从未想到要为1955至1970年间为我自行凸显的那项长期工作纲领写一篇概要,一如我为最近十二年写下《纲领概要》。我想原因很简单:从未出现过一个特别的机缘(如同如今我进入CNRS的申请),足以\dmark{}\origpage{200}促成这样一项陈述性的工作。在收入《模型的历史》(《数学反思》II)第I章附录的致Larry~Breen的信(1975年)里,可以找到一些指示,涉及我1970年前日程表上的某些理论(尤其是对偶理论)——这些理论至今仍在等待有人伸出双臂,把它们接引进公共的遗产之中。

\notehead{48$_{2}$}{}

同样的话对Motif理论也成立,只是有一点不同:它多半注定要在一段时期内保持其猜想形态。

\subsection{2.~当值无名之辈与好上帝定理}

\notehead{48$'$}{[该注由注46所召]}

按惯例,人们以完成了提炼并确立一项理论之关键定理的人的名字来命名这些定理;而Zoghman~Mebkhout的名字,看来却被判定配不上这条基本定理——那是四年顽强而孤独的工作(1975--79)的结晶,逆着当时的时尚与其长辈们的轻蔑而行。当这定理的分量再也无法被忽视的那一天,这些人兴致勃勃地把它称作“Riemann--Hilbert定理”,我完全相信(尽管Riemann与Hilbert肯定未曾要求过这份荣耀……)他们这样做自有极好的理由。毕竟(一旦需要之感——即理解一般离散系数与连续系数之间精确关系的需要——顶着普遍的冷漠显现出来,一旦它经由细致而耐心的工作得到提炼与明确,一旦经过相继的几个阶段那正确的陈述终于被提炼出来,一旦它白纸黑字地写下并被证明,而最终,当这件孤独的果实之定理在最不被人期待之处证明了自身——经过了这一切之后)这定理显得如此显然(对那些“本会证明它”的人来说,简直可以径直说“平凡”……),实在犯不着让一个模糊的当值无名氏的名字来占据记忆!

受了这一先例的鼓舞,我提议从今往后把任何一项理论中真正自然、真正基本的定理称作“亚当与夏娃定理”,或者干脆追溯到更远,把荣耀归于当得荣耀之处,径直称之为“\textbf{好上帝定理}”\fnote{32}{(*)}{在我作为数学家的一生中,我从未有过这份乐趣:在一名学生那里激发、哪怕仅仅是鼓励出一篇含有“好上帝定理”的学位论文——至少不曾有过深度与分量可比的。}。

\pmark 据我所知,除我之外,Deligne是Mebkhout之前唯一一个感到如下问题之价值所在的人:\origpage{201}在一个比分层模更宽广的框架里理解离散系数与连续系数的关系,以便能以“连续”的术语诠释任意的“可构造”系数。沿这一方向的第一番尝试,是Deligne1968或1969年在IHES主持的一场(未发表的)研讨班的内容:他在其中引入“分层pro-模”的观点,并给出一个比较定理(在复数域上),把超越的离散上同调与相伴的De~Rham型上同调联系起来——后者对任意特征零基底域上有限型的概形都有意义。(显然,他那时还不知道他那两位遥远的先行者Riemann与Hilbert的卓越结果……)因此,比起Verdier\fnote{33}{(*)}{Verdier作为Zoghman~Mebkhout学位论文的名义导师(并以此身份“恩准过几次讨论”),看来正是(除Mebkhout本人之外)围绕这条基本定理之父系渊源所发生的隐匿——以及在由Mebkhout所发展的D-模观点为代数簇上同调理论开启的革新中其“学生”应得的信誉——的主要当事人。然而我不知道他是否曾为此比Deligne更多一点不安。}或Berthelot\fnote{34}{(**)}{(5月25日)写这些文字时,我(略有踌躇地)没有把我的朋友Luc~Illusie的名字放进这份“本最有资格”向Zoghman~Mebkhout给予那些本属当然的鼓励的我方学生名单里。我当时未曾留意我内心某种不安——它本可以教会我:我这是在暗暗为一位我心怀感情的人开脱,装作卸下他那与我的其他“上同调学家学生”同样背负的责任。},Deligne都处在一个格外有利的位置,足以赏识Mebkhout1975年所投入的研究方向的全部价值,以及此后Mebkhout诸结果——尤其是“好上帝定理”——的价值;这条定理以连续系数的语言,给出了比他本人所提炼的更精妙、更深刻的离散系数领会。但这并未妨碍Mebkhout不得不在一种令人难堪的精神孤立中继续其工作,也未妨碍属于他的信誉(我更愿意说:尤其属于他的信誉)至今——五年之后\fnote{35}{(***)}{(5月25日)其实,这一隐匿首先是Deligne与Verdier本人的作品。对此可参见“不义——或归来的意义”一注,n\textsuperscript{o}~75。}——仍遭藏掖。

\subsection{3.~罐装的权与十二年的秘密}

\notehead{49}{[该注由注46所召]}

\pmark 核对之后(载 Publications~Mathématiques~35, 1968),我发现:在《Lefschetz定理与谱序列的退化判据》一文临近结尾处\origpage{202},有三行文字影射了某些“权方面的考虑”,正是它们使我(以一个略欠一般性的形式)猜想出了该文的主要结果。我怀疑这句谜语式的影射能否对任何人有用,或者在当时能被Serre和我之外的第二个人读懂——反正我们俩本来就已知情\fnote{36}{(*)}{(4月29日)对这篇文章的更细考察——它从不止一个角度看都富有教益——见“驱逐”一注(n\textsuperscript{o}~63)。}。

借机指出:早在六十年代末,在Weil猜想的余波里,一部非常精确的“权的瑜伽”——包括对$R^if_*$与$R^if_!$这类运算下权的行为的刻画——对我说来(因而也对Deligne)已是熟知的了。这部瑜伽的一部分最终在Deligne的“Conjectures de Weil II”(《Publications Mathématiques》1980)中得到确立(在$\ell$-进系数层的语境里;更自然的Motif框架尚待其确立)。如无差错,在前后相隔的这十二年间\fnote{37}{(**)}{(4月19日)我刚收到并饶有兴味地读了Deligne的论著目录,从中发现,早在1974年,他在温哥华大会上的一次报告中就谈到了“权”——如此,“权周围的秘密”是六年,而非十二年。不过在我看来,这一秘密与围绕Motif的类似秘密(1970--1982的十二年间)密不可分。这一秘密的意味,在今天的反思中、在n\textsuperscript{o}~51--52之后那则长篇双注里,刚刚被一束新的光照亮。},文献中没有留下一处踪迹——哪怕是最简短、最不完全的陈述——述说这部(当时还完全是猜想性的)权的瑜伽;在这整段时间里,它始终是寥寥(两三个?)知情人的专属特权\fnote{38}{(***)}{(5月25日)从反思中陆续出现的全部信息来看,这几个“两三个知情人”似乎到头来只有Deligne一人:看来他煞费苦心,把这部得自我的瑜伽的专属收益留给自己,直到1974年(见上一条脚注)——那时机缘已熟,他尽可以把它当作自己地里长出的观念来呈献,既不提我,也不提Serre(见注n\textsuperscript{o}~$78'_{1}$、$78'_{2}$)。

(1985年4月18日)在这些文字写成之后,我又有机会得知Deligne在尼斯国际数学家大会(1970)上的报告“Théorie de Hodge I”(《会议文集》t.1,第425--430页)。与我据手中残缺信息而有理由相信的相反,这篇文章早在1970年就陈述了权的瑜伽的相当一部分。关于这些观念的来源,他只满足于对Serre的一篇文章(它与本题其实无关)以及“Grothendieck的猜想性Motif理论”做了一句谜语式的、纯礼节性的提及。(可与注n\textsuperscript{o}~$78'_{1}$、$78'_{2}$对读。)至于权的观念在$R^if_!$与$R^if_*$这类运算下的行为这一关键问题,则连提都没有提,而且在1980年前引“Conjectures de Weil II”一文之前也不会提;在该文主要定理的关系中我的名字未被提及,同样,在前一条脚注(一年前的今天)所引的那次报告“代数簇上同调中的权”里,Serre的名字和我的名字也都未被提及。}。然而这部瑜伽构成理解代数簇上同调“算术”性质的第一把本质钥匙:它既是在一给定情形中辨认方向的\textbf{手段},也是做出种种预测的手段——其可靠性从未落空——\origpage{203}而同时,正因如此,它也代表着代数簇上同调理论中摆在我们面前最紧迫、最迷人的\textbf{任务}之一。这部瑜伽直到最终被确立(至少在某些重要方面)之前几乎无人知晓,在我看来,正是“\textbf{信息封锁}”的一个格外触目的例证:扮演这一角色的,往往恰恰是那些身居其位、肩负其职、理应确保信息广为传播的人\fnote{39}{(*)}{对此另可参见第32、33节“数学家的伦理”与“注记——或新伦理(1)”,以及与之相关的两注:“共识性的操守与信息控制”与“年轻人的势利,或纯洁的捍卫者”,n\textsuperscript{o}~25、27。}。

\subsection{4.~进步是挡不住的!}

\notehead{50}{[该注由第(I)部分《愚妄与更新》第八章“孤独的历险”第50节所召]}

我在这方面的最初经验,是我为设法让Yves~Ladegaillerie关于曲面同痕定理的学位论文获得出版而付出的那些徒劳努力所结出的意外果实——那项工作,说实话,丝毫不逊于我曾充当“导师”的那十一篇国家博士论文(“1970年以前”的,当然!)。如果我没记错,这些努力持续了整整一年多,主角中有相当一批我昔日的朋友(当然,还少不了我昔日的一个学生)\fnote{40}{(**)}{对此可参见“棺材2——或被切成一段段的剖视”一注,n\textsuperscript{o}~94。}。其主要情节,今天在我眼前仍如一场场通俗笑剧!

这也是我头一回遭遇到某种新精神与新习俗(它们在我昔日朋友圈里已成风气)——在这次反思中我已在这里或那里有机会影射过它们。正是在那一年(即1976年),我头一次(但不是最后一次)获悉:如今,认认真真去证明那些人人都在用、而先行者们从来满足于默认的精微事项(在这里,即曲面拓扑中不存在野性现象)\fnote{41}{(***)}{对此另可参见“注记——或新伦理”那一情节(第33节)。那则大名鼎鼎的“注记”的过错,恰恰在于把一些向来含糊其辞的观念与陈述挑明了;而我却曾一直含蓄地使用它们,来建立那些以我的名字命名的结果,而且将近二十五年来,人人都在毫不脸红地使用这些结果(这一点,那两位鼎鼎大名的同行其实心知肚明)。(6月8日)详情另见“棺材4——或没有鲜花也没有花冠的topos”一注(n\textsuperscript{o}~96)。所谓“以我的名字命名的结果”,是关于某些整体与局部投射有限基本群的生成与有限表现的结果,它们曾在SGA~1等处得到“证明”,所用的下降技术在没有理论上的、细致的论证的情况下仍属启发性质;那样的论证,在Olivier~Leroy那项(显然“无法发表的”)关于topos基本群的Van~Kampen型定理的工作中完成了。},被视为不够严肃(至少就一个无名之辈而言……)。或者,去证明一个把若干已知的深刻定理\dmark{}\origpage{204}作为特例或推论涵盖于其内的结果(这显然表明,那个所谓的新结果只能是已知结果的一个特例或一个容易的推论)。或者,仅仅是在一个结果的陈述里,或是在用一种语言描述另一种情形时,肯费心把自然的假设仔细表述出来(这是可叹的“无谓折腾”的标志),而不肯把自己局限于投某位高格调意见人士之所好的那种特例。(去年我还见到有人责备Contou-Carrère,说他在其学位论文里不该把一般概形、而只把基底域当作立脚点——同时又赏脸承认,他想必是在其临时导师的坚持下才不得不如此。不过,这样说话的人对门道也足够“在行”,本应知道:即便把自身限定在复数域上,证明的需要也会逼着人引入一般的基底概形……)

当今某种时尚的迷途,竟至于不仅鄙薄细致的证明(乃至干脆鄙薄证明本身),甚至常常鄙薄成形的陈述与定义。照纸张与被填鸭的读者的耐心目前的行情,不久之后,人们大概就不再肯为如此昂贵的奢侈品劳神了!顺着当前的趋势外推,应当能够预言这样一个时刻:届时出版物中将不再有人费神去明确陈述定义或定理,人们从此满足于用一些代码词来指称它们,把填补空白的差事留给自己那不知疲倦、天才横溢的读者,随他按自己的悟性行事。审稿人的差事也将随之大为轻松:他只需到“名人录”(Who is Who)年鉴里查一查,作者是否以可信著称(反正谁也没法反驳一个由空白与省略号组成的精彩论文),抑或相反,是一个不便启齿的无名氏——那他早就(如同今天已然如此、且由来已久)被当然地扫地出门……

% Original notes 27 and 28 have plain callouts but starred body labels.
\providecommand{\fnoteplainmark}[3]{%
  \begingroup\renewcommand{\thefootnote}{#1}\footnotemark
  \renewcommand{\thefootnote}{#1#2}\footnotetext{#3}\endgroup
}
\chapter{B)~皮埃尔与Motif}

\section{IV.~Motif——或一次诞生的下葬}

\subsection{一个梦的回忆——或Motif的诞生……}

\notehead{51}{[该注由注46所召]}

\pmark (4月19日)写下这些文字(它们是“我的孤儿们”一注,即n\textsuperscript{o}~46的收尾)还不到一个月,我已发现它们稍稍落在了事件的后面!我刚收到“Hodge Cycles, Motives and Shimura Varieties”(LN~900),作者为Pierre~Deligne、James~S.~Milne、Arthur~Ogus与Kuang-Yen~Shih——那是Deligne惠寄的,还附了一份他的论著目录。\origpage{205}这部汇集六个文本的文集1982年出版,是1970年以来一个有趣的新事实:Motif观念出现在标题里,并在正文中占有了一席之地——虽然还极其谦逊,主要是经由“MotifGalois群”的观念。当然,距离一部Motif理论的全景图还很远;十五二十年来,这幅全景图一直在等待那位勇敢的数学家,肯把它“描绘出来”——足够宽广,以充当一代或几代算术几何学家的灵感之源、引导线索与地平线,他们将有幸确立其有效性(或无论如何,发现Motif之实在的底细……)(注53)。

好像也是从1982年起\fnote{1}{(*)}{(5月25日)我又落后了——这一回是一年:转折发生在1981年6月的吕米尼会议;参见“不义——或归来的意义”一注,n\textsuperscript{o}~75。},时尚的风向对着导出范畴多少开始转了;Zoghman~Mebkhout(在一阵或许有点欣快的高谈里)已经看到它们即将“席卷数学的所有领域”。它们的用处——任何一个消息灵通的人,单凭数学本能,早在六十年代初就已看得很清楚——如今才开始被承认;这(在我看来)主要得归功于Mebkhout孤身的努力:整整七年,他顶着一场专横的时尚,以只信自己本能者的勇气,扛下了为人开路、承受初试代价这桩吃力不讨好的差事……

\pmark 引人注目的是,读这部第一次让(在我退出数学舞台十二年之后)Motif观念低调重返“通行”数学观念之列的出版物,任何不知情的读者都无从疑心:\origpage{206}我这个人,竟与这个长期遭禁的观念的诞生、与围绕它发展起来的一部丰富而精确的“瑜伽”,有过任何关联——在那里,它(以极其残缺的形态)宛如凭空出现,不提任何父系渊源(注51$_{1}$)。

仅仅三周之前,我花了一两页篇幅大谈Motif瑜伽——就连我的一个“孤儿”,而且是我比任何一个都更挂心的那个,也表明我当时显然完全没有说到点子上!毫无疑问,我是在做梦:当我自觉得记起多年孕育的岁月——一种先是微弱而难以捉摸、随后逐月逐年丰富与明确的视野,在那顽强努力的尝试中,试图把握诸种已知的上同调理论(注54)所各自体现的那个“motif”、那共同的本质——它们每一种都以自己的语言,谈论着作为其直接可触的表现之一的那个“motif”的本性。毫无疑问,我还在做梦:当我记起Serre的某个直觉给我留下的强烈印象——Serre被引向把一个投射有限Galois群(一个本质上属于离散性质的——或至少依定义可归结为一个\textbf{有限}群投射系统的——对象)看作这样一个系统的源头:一个庞大的$\ell$-进\textbf{解析}群、乃至$\mathbb{Q}_\ell$上\textbf{代数}群(经由合宜的代数包络)构成的投射系统;这些群甚至往往是约化群——随之,解析群与代数群的(Lie式)直觉与方法的整套武库也被引入。这一构造对每个素数$\ell$都有意义;而我感到(或者我梦见我感到……)这里有一个待探测的奥秘:对不同素数,这些代数群之间的关系如何;它们必定全都出自同一个代数群投射系统,位于所有这些基底域唯一自然的公共子域之上,即特征零的“绝对”之域$\mathbb{Q}$。既然我喜爱做梦,我就继续梦见:我记起自己经由一项肯定只是梦的工作(因为我什么也没有“证明”)进入了这个隐约瞥见的奥秘;我记起自己终于懂得,Motif观念如何提供了理解这一奥秘的钥匙——如何仅仅凭着一个范畴的在场(这里是给定基底概形\dmark{}\origpage{207}上的“光滑”Motif的范畴,比如给定基底域上的Motif),它带有与域$k$上一个代数pro-群的线性表示范畴上所见相似的内蕴结构(代数pro-群观念的魅力,先前也是Serre向我揭示的),人们就真能重建出这样一个pro-群(一旦有了一个合宜的“纤维函子”),并把“抽象”范畴诠释为它的线性表示范畴。

通往一部“MotifGalois理论”的这一进路,是我多年前为描述一个拓扑空间或概形(乃至任意一个topos——不过在这里我感到,自己要刺痛那些“topos并不能使他们开心”的娇贵的耳朵……)的基本群所找到的那条进路向我提示的:即经由所论“空间”的平展覆盖的范畴,以及其上的纤维函子。而“MotifGalois群”的\textbf{语言}本身(我也满可以叫它们Motif“基本群”——两类直觉,自五十年代末以来,对我是同一回事……),连同“纤维函子”的语言(它们与上文所说的那些“显明的化身”——即适用于一给定Motif范畴的各个“上同调理论”——恰相对应),正是为表达这些群的深层本性而造的语言,并且显然提示了它们与Galois群及寻常基本群的直接联系。

我还记得在纤维函子这场游戏里的乐趣与惊叹,还有在Galois群作用下、使不同纤维函子通过“扭转”彼此过渡的torseur——在一场格外具体而迷人的境况里,重新找回Giraud著作中所发展的非交换上同调的整套观念:由纤维函子构成的gerbe(这里的底是平展topos,或者更好,$\mathbb{Q}$的fpqc topos——非平凡而有趣的topos,若曾有过的话!),以及束住这条gerbe的“带”(lien,以群或代数pro-群为值),连同这条带的诸种化身——经由各样代数群或代数pro-群实现,对应于gerbe的各个“截面”,也就是各个上同调函子。比如,特征零概形的各个复点,经由相应的Hodge函子,生成了gerbe的那么多截面,以及从一个截面到另一个截面的torseur;这些torseur连同其上运作的pro-群,带有卓越的代数-几何结构,表达着Hodge上同调的特有结构——不过在这里\dmark{}\origpage{208}我抢在了Motif之梦的另一卷前面……那是这样的年代:今天引领时尚的人们尚未宣布,topos、gerbe之类不逗他们乐,谈这些就是胡扯(当然,就算他们这么说,也丝毫不会妨碍我在它们所在之处认出topos与gerbe……)。而如今十二年又过去了,同一些人却装出发现并教导的样子说:gerbe(即便还不是topos),确实与代数簇的上同调、甚至与阿贝尔积分的周期,大有关涉……

在这里我本可以唤起另一段记忆的梦(或另一个梦的记忆……),它也生于一记“强印象”(我分明深陷于主观性之中!):Serre就Weil猜想背后某种“哲学”所作的若干评论。把这些猜想翻译成$\ell$-进系数($\ell$变动)的上同调语言,令人隐约感到,相应的上同调带有卓越的结构——“权滤过”的结构\fnote{2}{(*)}{(1985年1月24日)对这段走了样的记忆的更正,见注n\textsuperscript{o}~164(I4)及子注n\textsuperscript{o}~164$_{1}$,其中对“权的瑜伽”的渊源有更准确的说明。}。不同$\ell$-进上同调所共有的那个“motif”,想必正是这一本质算术结构的终极载体;由此,这一结构取得了几何的面貌——即几何对象“motif”上的一种卓越结构。把它称作“工作”(尽管这分明仍旧不折不扣是猜谜的片段)肯定又是我在自欺——当需要“猜出”(唯一的向导,是一个正在成形的视野的内在自洽,借助于此处或彼处已知的或猜出的零星要素……)一个Motif的各个“化身”的上同调所特有的结构:权滤过\dmark{}\fnote{3}{(**)}{(1985年2月28日)这里我的头脑有一点点混乱。实际上,说的是与之紧密相连的按“层”(niveaux)的过滤。}在其中如何体现——从Hodge化身开始(在Hodge--Deligne理论尚未问世的时代,而且不无原因……\fnote{4}{(***)}{那是一个年轻的Deligne恐怕还没有在数学语境里听说过“概形”一词、也没听说过“上同调”一词的时候。(这些观念,是他自1965年起在与我的接触中认识的。)})。这使我在(梦中)得以看到,Tate关于代数闭链的猜想(这已是第三记“强印象”,它在Motif之梦里激励着做梦者!)与Hodge猜想(注55)汇入同一幅宏阔的图景,并从中提炼出两三条同一成色的猜想;我对某些人谈起过它们,他们想必已把它们忘了,因为我从此再没听到有人提起,正如再没人提起“标准猜想”一样\dmark{}。无论如何,\origpage{209}它们只是猜想(而且并未发表……)。其中一条并不涉及某个特定的上同调理论,而是给出:一域上非奇异射影簇的Motif上同调上的权滤过,经由该簇自身按给定余维数的闭子集所作的几何过滤的直接诠释(余维数扮演“重量”的角色)\fnote{5}{(*)}{(1985年2月28日)实际上,这里是按“层”(niveaux)的过滤(参见前一条脚注)。}。还有那项“猜出”权在六种运算下的行为的工作(我本该给“工作”打上引号,可又下不了决心!)(那些工作早已失传……)。在那里,我同样从不曾有发明之感,而始终是发现——不如说,是听事物向我言说,当我肯手握笔杆费心去听的时候。它们所说的,有一种不容错辨的精确。

然后有了第三个“Motif之梦”,它像前两个梦的联姻——这一回,是要诠释Motif范畴所带的各种附加结构(其中最首要的,恰恰就是权滤过):用MotifGalois群上的结构,以及以这些群为结构群、用来“扭转”一个纤维函子以(典范地)得到任一其他纤维函子的torseur\dmark{}\fnote{6}{(**)}{正如某个“空间”$X$在两个“点”$x$与$y$处的基本群$\pi_1(x)$、$\pi_1(y)$,经由“扭转”——即经由从$x$到$y$的道路类的群$\pi_1(x,y)$这个torseur——即可彼此约化……}的语言。我记得,在这里更谈不上什么猜谜,而是规规矩矩的数学翻译。它们是一道道未发表的关于代数群线性表示的“练习题”,我兴味盎然地做了几天、几星期,分明感到自己正一步步逼近一个多年来令我心醉神迷的奥秘!必须以表示的语言去把握与表述的观念里最精微的,或许是Motif的“极化”观念——我由Hodge理论取得启发,力图从中澄清在Motif语境中仍有意义的东西。这一反思,应当就发生在我思考“标准猜想”之表述的前后——两种思考\dmark{}\origpage{210}皆受启于Serre(永远是他!)关于Weil猜想的“kählérien”类似物的想法。在这样的境况里,当事物本身向我们低语:它们隐藏的本性为何,我们经由何种手段能最精妙、最忠实地表达它——尽管许多本质事实似乎还远在证明的射程之外——单纯的直觉告诉我们:把事物执拗地向我们低语的一切,径直白纸黑字写下来;我们越肯听写,写下的就越清晰!无须为完备的证明或构造操心——在这一阶段背负这类要求,无异于自绝于一项大面积发现工作中最精微、最本质的阶段:视野的诞生——它从一派表面的虚无中取得形态与实体。单单\textbf{书写}、\textbf{命名}、\textbf{描述}这一行为——哪怕起初只是描述那些难以捉摸的直觉、或那些迟迟不肯成形的“疑点”——就具有一种\textbf{创造性的力量}。当认识的热情投入理智所能把握的事物之中,这正是它手中的头号工具。在这一类事情的发现之路上,这一工作是所有环节中最具创造性的环节:它永远先于证明,并为我们提供证明的手段——或者不如说,没有它,“证明”某事的问题根本无从提出,因为凡是触及本质的东西,还没有一样得到表述与看见。

单凭表述努力的力量,无形的开始成形,变得可以审察;明摆着是错的与可能的得以澄清;而其中那些与已知的或猜出的事物之整体如此完美地协调的,又转而成为正在诞生的视野的坚实可靠的要素。这一视野在表述的工作中日益丰富、日益明确。十样仅仅被猜到的东西——没有一样(比如Hodge猜想)单独令人信服——却相互映照、相互补足,似乎汇入同一种尚带神秘色彩的和谐,便在这和谐中获得视野的力量。即便这十样到头来全被证明是错的,那通向这一临时视野的工作也没有白做;它使我们隐约瞥见、并容我们稍稍深入的那种和谐,不是幻象,而是实在,召唤着我们去认识它。正是经由这一工作,唯有经由这一工作,我们才得以与这一实在、这一隐藏而完美的和谐亲密接触。当我们懂得:事物之为此物自有其理\dmark{}\origpage{211},我们的天职是认识它们而非支配它们,那么,一个错误曝光的日子,就是欢腾的日子(注56)——正如某项证明使我们毫无疑义地确知:我们曾如此想象的某事,当真是实在本身的忠实而真正的表达——那一天同样是欢腾的日子。

在这两种情形里,这样的发现都是一项工作的报偿,没有那工作,发现无从发生。但它或许要待经年累月的努力之后才姗姗而来,或许我们永远听不到那最后的底细——那是留给后来者的——而工作本身就是它自己的报偿,在每一瞬间都因这一瞬间向我们揭示的东西而丰盈。

\notehead{51$_{1}$}{}

(6月5日)然而Zoghman~Mebkhout刚刚提请我注意:前引书第261页,在Deligne一篇“重写并补充一封致Langlands的信”的文章里,提到了“Grothendieck的Motif”。那里写道:“这里要说的不是Grothendieck的Motif——他那样用代数闭链来定义的Motif——而是\textbf{绝对Hodge Motif},同样用绝对Hodge闭链来定义。”这里点名“Grothendieck的Motif”(不加着重),不是作为灵感之源,而是为了与之划清界限,并强调说的是\textbf{另一回事}(人们特意加了着重)。这一撇清格外引人注目,因为Hodge猜想的成立(我假定Deligne,一如其信函式文章的任何读者——从最初的收信人Langlands起——都知道这条猜想)将蕴含这两个观念\textbf{同一}!!

当然,早在1964年——我发展MotifGalois群观念的那一年——我就清楚地知道,可以在同一个模型上发展“Hodge Motif”观念,以及相应的“Galois-Hodge Motif群”观念;后者后来由Tate独立引入(我说不好在他之前还是之后),并得到了“Hodge-Tate群”(联系于一个Hodge结构)之名。这桩粗劣的欺骗(但出自这样一位显赫人物,它似乎并不叫任何人感到不适)在于:以一个可笑的借口——即通往该观念的技术进路(经由绝对Hodge闭链,而非代数闭链)与我所(极临时地)采取的进路不同(倘若Hodge猜想是错的)——把一个新颖而深刻观念的父系渊源干脆整个地藏掖起来,连同我围绕Motif发展起来的一整片丰富的直觉织体。我发展了近十年之久的这部瑜伽\dmark{}\origpage{212},自1968年肇始,就是Deligne著作中灵感的主要来源。它的多产与作为发现工具的力量,早在我1970年离去之前就昭然若揭;而它的同一性,不依赖于为确立它某一部分有限内容之有效而采取的任何技术进路。Deligne有幸提炼出这样两条互不依赖任何猜想的进路。但他却没有那份诚实点明其灵感之源;相反,从1968年起他就极力对所有人遮掩它,以便把它的收益独占为己有,只等1982年(以默许的方式)认领这份信誉。

\subsection{下葬——或新父亲}

\notehead{52}{}

回到Motif之梦,我记得自己还曾把它大声梦见。诚然,梦的工作按其本性是孤独的劳作——但这项顽强地持续了多年的工作(在一项吞噬了我大部分时间的浩繁奠基写作的边缘上进行)的种种波折,有一位日复一日的见证者,比Serre亲近得多——Serre只满足于远远地看这些事……\fnote{7}{(*)}{(5月25日)不过,我对Motif的反思,其实早于Deligne的出现。我关于MotifGalois理论的手写笔记,标注的日期是1964年。}。关于这位日复一日的心腹之交,我在回顾中写道:他在六十年代中期“多少扮演过学生的角色”,而我“给他讲过我在代数几何里知道的那一点点”。我本可以补上一句:我连按这个词的通常意义我所不曾“知道”的东西也讲给他听——那些数学的“梦”(以Motif为主题,也以别的为主题),它们总能在他那里找到一只专注的耳朵和一颗警醒的心灵,和我一样渴求理解。

诚然,当我写道Pierre~Deligne曾“多少扮演过学生的角色”时,这仍旧全然是一种主观的印象(注57),据我所知,没有任何书面的、至少是印刷的痕迹可资佐证——那种痕迹本可以让随便哪个人疑心,Deligne从我嘴里学到过什么——而在这里,我却乐于记起:我同他谈数学,从来没有一次不从中学到点什么。(甚至在我同他不再谈数学之后,我仍继续经由他学习那些更困难、或许也更重要的事情,包括在我写下这些文字的今天……)。

\pmark 新近,一位第三者(他不知怎么猜到了这事或许会让我感兴趣!)告知我,Deligne和别人有一部文本,其中谈到Motif,或者至少谈到“塔纳卡范畴”;我向Deligne提了一句,他由衷地惊讶,我居然会对这类东西感兴趣。\origpage{213}翻阅他惠寄的那一册,我确信他的惊讶完全有理。显然,我这个人,同书中所论的主题全然无涉。至多在导言里一句话带过:某些“标准猜想”(当年是我提出的,天知道为什么)对一域上Motif范畴的结构会有一个推论……好奇的读者想多知道些就为难了:他在这整部书里,找不到关于这些猜想的任何说明或出处——它们此后就再没被提起;也找不到对唯一一篇我解释如何以标准猜想在域上构造Motif范畴的已刊文本的提及;也找不到1970年以前唯一另一篇谈到Motif的已刊文本——出自Demazure之手(如果我没记错,是在一场Bourbaki研讨班上),它遵循我那条特设的构造原则,只是视角稍有不同……\fnote{8}{(*)}{核实之后,我发现,除了关于标准猜想的几页(Algebraic Geometry, Bombay, 1968, Oxford Univ.~Press (1969), 第193--199页),我没有任何已刊的数学文本谈到Motif。在Demazure的报告(Bourbaki研讨班n\textsuperscript{o}~365,1969/70,继Manin的俄文报告之后)里,提到我1967年在IHES所作的几次报告,它们(我猜想)应当构成Motif 视野的第一幅总草图。关于标准猜想及其与Weil猜想的关系,有一篇比孟买会议通告更详细的报告,出自Kleiman之手(Algebraic Cycles and the Weil conjectures, 载 Dix exposés sur la cohomologie des schémas, Masson-North Holland, 1968, 第359--386页)。在1970年以前,除我之外,我没有见到任何关于标准猜想的反思——尤其朝着对它们的证明。存心无视这些关键猜想(我在孟买的草图里说过,我认为它们连同优秀概形的奇点消解,是代数几何里最重要的未决问题)的做法,在我看来,很大程度上造成了代数簇上同调理论给我的那种停滞印象——就传到我耳中的种种回响而言。}

\pmark 连Neantro~Saavedra——他有幸位列我“1970年以前的学生”——都得到了应有的引用\origpage{214}。他跟我做过一篇学位论文,论的是我记得我称之为“刚性张量范畴”、而他称之为“塔纳卡范畴”的东西。人们至今仍在纳闷:Saavedra是靠了怎样奇迹般的运气,恰好预见了十年后才绽开的DeligneMotif理论的需要!实际上,在他的论文里,他做的恰恰就是\textbf{那}项在技术上构成MotifGalois理论之钥匙的工作,一如 J.L.~Verdier的论文原则上就是\textbf{那}项在技术上构成上同调六运算形式主义之钥匙的工作。有一个(此外还有别的)对Saavedra有利的不同:他费心发表了自己的工作;他确实不曾有Hartshorne、Deligne与Illusie联袂的笔杆,来替他免除这道手续。然而十年之后,Saavedra的论文被从头到尾、几乎一字不易地复制进那部卓越的文集里,这一次执笔的是Deligne与Milne。倘若只是为了更正Saavedra工作中的两个具体之点(注58),这或许并非必要。但万事必有其所以然,我想我辨出了Deligne亲自不辞这份辛劳的缘由\fnote{9}{(*)}{对此可参见“白板”一注中的反思,n\textsuperscript{o}~67。}\fnote{10}{(**)}{(6月8日)而当事情涉及带有我的影响之痕迹的工作时,就更是如此——对此可参见“注记——或新伦理”那一情节,第33节。}:那分明同他自己那条在出版事宜上把严格推到极点的标准背道而驰;而众所周知,一涉及别人,他把这条标准执行得一丝不苟……

至于诸观念与Motif瑜伽本身的父系渊源,对不知情的读者(而知情的读者已开始变得稀少,终将一个个寿终正寝……)来说,这一父系渊源无可置疑——无须在这里去惊动遥远的Hilbert与Riemann,更不必去惊动上帝。倘若那位显赫的作者——他关于阿贝尔簇上绝对Hodge闭链的漂亮结果,俨然是Motif理论的起点,干脆说,是Motif理论的诞生——对其父系渊源只字不提,那是一种令他生辉的谦虚\dmark{}\origpage{215},与行业的惯例与伦理完美合拍:按后者的要求,应当把荣耀归于它分明所当得之处的差事(如有必要的话),留给别人:留给合法的父亲……

\notehead{53}{}

被这个孤儿的坎坷际遇所触动,并疑心不会再有别人去做这项工作——这项工作,至今显然只有我一个人感到其需要与规模——我料想,那位“勇敢的数学家”不是别人,正是我自己,待我走完《追寻栈》之路之后(我预计它还要占去我大约一年)。

\notehead{54}{}

从那以后,代数簇又出现了两种新的上同调理论(除Hodge--Deligne理论之外——后者是Hodge上同调在“Motif”精神里的自然延伸),即Deligne的“分层pro-模”理论,以及尤其是晶体理论——Sato--Mebkhout式的“D-模”版本,配上前文说过的“好上帝定理”(又名Mebkhout定理)所带来的新的照亮。这条通往可构造离散系数的进路,多半注定要取代Deligne先前的版本,因为它无疑更适合表达与De~Rham上同调的关系。况且,这些新理论并不在给定概形上光滑Motif的范畴上提供新的纤维函子,而是(尚需一项比迄今所做更深彻的奠基工作)提供一种方式,以精确把握复数域上有限型概形上一个Motif(不必光滑)的“Hodge”化身,或特征零域上有限型概形上的“De~Rham”化身。此外很有可能:$\mathbb{C}$上有限型概形的Hodge--Deligne系数理论(显然始终未见成文),终将显现为被包含在Sato--Mebkhout式晶体系数理论(同样未见成文;并另附一项关键的滤过数据)之中,或者更确切地说,显现为后者与$\mathbb{Q}$-向量空间的可构造离散系数理论的某种交……至于澄清Mebkhout式晶体理论与Berthelot及其他人在正特征下发展的理论之间的关系,那是Mebkhout早在1978年之前就已感到的任务,处在一种普遍冷漠的气候里;而在我看来,它是近期摆在我们面前、关乎对代数簇“那一个”上同调(唯一而不可分,即Motif的上同调!)之理解的最迷人的任务之一。

\notehead{55}{}

\dmark{}我尽管做梦,但我关于Motif与Hodge结构之关系的梦,却让我在不经意间把手指按在了Hodge“广义”猜想——按Hodge最初表述的样子——的一处不自洽上,并把它代之以一个更正的版本;这一回(我敢打赌),它应当既不会比代数闭链上那个“寻常的”Hodge猜想更错,也不会比它少错。\origpage{216}

\subsection{一场大屠杀的前奏}

\notehead{56}{}

我尤其想到——恰恰在代数簇上同调的语境里——Griffiths对一种诱人的想法之错误性的发现:长期以来人们以为,一个同调等价于零的闭链,会有一个代数等价于零的倍数。对一个全新现象的这一发现,当时给我的触动之深,竟使我花了整整一星期的功夫去力求吃透Griffiths的例子,把他的构造(那是超越的,在$\mathbb{C}$上)搬移成一种“尽可能一般”的构造,并且在任意特征的域上成立。这一推广并非一目了然,靠的是(如果我没记错)Leray谱序列与Lefschetz定理。

(6月16日)这一反思,曾是我有机会在平展语境里发展“Lefschetz束”的上同调理论的机缘。我这方面的笔记,已由 N.~Katz在SGA~7~II研讨班(由 P.~Deligne与 N.~Katz主持)的报告XVII、XVIII、XX里详细阐发(他费心指涉了这些笔记,而且紧随其步伐)。而在 P.~Deligne所写的该卷导言里——其中说,该卷的关键结果是报告XV(平展上同调中的Picard--Lefschetz公式)与XVIII(Lefschetz束理论)——作者却小心地不提:在这套Lefschetz束的“关键理论”里,我也有我的一份。读这篇导言,得到的印象是:对该卷所发挥的主题,我毫无贡献。

SGA~7那场漫长的研讨班,1967--69年间接续了1960至1967年间在我推动下发展的SGA~1至SGA~6诸研讨班,曾由Deligne与我共同主持,是我以消没闭链群的一个系统理论发令起跑的。由于各报告由不同的志愿者执笔、拖沓经年,研讨班的两个分卷(SGA~7~I与SGA~7~II)迟至1973年才出版,经手的是\dmark{}\origpage{217}Deligne。尽管在研讨班当时,大家已商定它将以一场共同主持的研讨班的面目出现;我离去之后,Deligne却向我表示了他那个(在我看来很奇怪的)愿望:把研讨班\textbf{一分为二},一部分I以我主持的名义出现,另一部分以他和Katz的名义。我如今在其中辨出一场“战役”,它预演了“SGA~4½之役”——那场战役的目标(除别的以外)是要让SGA~1至SGA~7整个奠基丛书显出这样一副面孔:仿佛它是一套随手采撷的文本汇编,我这个人只在其中扮演一个偶发的、乃至多余的角色——尽管就其精神与构思而言,这套丛书一如EGA《代数几何原理》丛书,是与我个人密不可分的。这一倾向在SGA~4½分卷里表现得非常清楚,甚至粗暴;尤其表现在对SGA~5研讨班的屠杀里——该分卷与那场屠杀密不可分。对此可参见(除别的以外)“白板”与“屠杀”两注,n\textsuperscript{o}~67与87,尤其是“遗骸……”(n\textsuperscript{o}~88)。

(6月17日)SGA~7研讨班的总体构思(我在其中绝不曾区分什么“I”“II”两部分,至今也仍不区分)出自我这里;另一方面,Deligne作出了重要贡献(在我1969年写的Deligne工作评述报告里有记载,见该报告n\textsuperscript{o}~13、14条),其中就研讨班的需要而言最关键的,是Picard--Lefschetz公式——经由从一个已知的超越情形做特化论证而证明。把研讨班切成两部分,无论在数学上还是就双方的贡献而言,都毫无道理——SGA~7的两个“块块”里,Deligne和我都有实质性的贡献。

当然,倘若Deligne继续我开创的SGA奠基丛书——它远未跑完自己的赛程!——我会喜出望外。这场“SGA~7之役”根本不是什么延续;我感到它是一记粗暴的“锯子”(或电锯……),\textbf{终结}了SGA丛书——用的正是一个摆明要与我个人划清界限的分卷,尽管它维系于我的著作,并且像其余各卷一样带着我的印记。当我的个人痕迹在其中被尽量藏掖的同时,对我著作的口吻,尚未达到“SGA~4½之役”那种几乎不加掩饰的轻蔑——后者是对SGA~4与SGA~5研讨班之统一体更深重的一记锯痕,也是对其中未刊的SGA~5部分大加劫掠的手段与借口;劫下的碎块,在Deligne与Verdier之间均分……

\notehead{57}{}

\dmark{}请允许我赶紧补充:同样的评语,也适用于另一位大手笔的数学家——十年之后,我斗胆说过(在注n\textsuperscript{o}~19里)他也“多少扮演过学生的角色”。 \origpage{218}

\notehead{58}{}

这让我想起:Springer的《Lecture Notes》(它曾发表过六七篇在我指导下完成的“1970年以前”的博士论文)从未肯发表Yves~Ladegaillerie“1970年以后”的那篇论文(理由:他们不发表博士论文!)。人们倒可以说,它把Saavedra的论文发表了第二遍……我还曾向Deligne提起Ladegaillerie那条四处碰壁的漂亮同位结果(心里暗暗指望他会出手相助发表)——却没能引起他的兴趣(理由:他对曲面拓扑不在行……)。

幕落……

\subsection{新伦理(2)——或抢夺的集市}

\notehead{59}{}

(4月20日)写下这些文字——它们记录了一桩矛盾及其代价——已有几个星期;我惊讶地发现,当事人早在两年前就已找到一条最简单的路子来“化解”那桩矛盾——难就难在想得到!人们可以称之为“预先下葬法”(读者可在昨天写成的双注(注50、51)里了解它,那是发现之后情绪未凉之时写的)。我感到遗憾的是:这位被预先下葬的死者,冷不防又重新出现在那有名的“数学舞台”上(它有时候,分明更像一个抢夺的集市……),恐怕会给这项光辉方法的干净利落的应用,带来技术上的麻烦!在前面的注里(“共识性的操守——与信息控制”,n\textsuperscript{o}~6),我(还有些含混地)感到:科学行业里那条最获普遍认同的操守规则,在掌管科学信息控制的人们不尊重“每个科学工作者有权让自己的观念与结果为人所知”这一权利时,“就形同虚设”。正是在反思的这个节点上,我还费心相当详细地描述了一个个案:对这一权利的藐视,在我眼里昭然若揭;而且我分明感到,这藐视已逼近对那条首要规则——获得普遍共识的那一条——的藐视。(见“注记——或新伦理”,第30节。)

\pmark \origpage{219}这已不是我头一回感到这种特殊的痛楚:当我看到这条首要规则的\textbf{精神}遭人藐视,而藐视者却处于“暂停”的豁免状态——无论靠的是其地位(高到不容置疑!)、其才具,还是形式上的漫不经心。我试着在相关于所引第30节的那则注里(“年轻人的势利——或纯洁的捍卫者”)勾勒这种痛楚。当人们可以放胆藐视我在那里谈到的那些“显而易见”的事,并且(如今我可以补上一句)以同样的精神,也藐视那些既未证明、也未以已刊并为众人所知的“猜想”的形式注册备案的事情(或许是深刻的事情),那么人们同样可以(鉴于其“分量之轻”\fnote{11}{(*)}{“好上帝定理”(又名Mebkhout定理)的遭遇,正是如此。(6月8日)而且,一如对Motif瑜伽那样,还费心巧妙地制造出对此物拥有父权的表象,却从不明言!就个案而言,可参见“变戏法的人”一注,n\textsuperscript{o}~75$''$;就这条光辉方法的一般路数或风格而言,可参见“暂停!”一注,n\textsuperscript{o}~77,以及紧随其后的“据为己有与轻蔑”一注,n\textsuperscript{o}~59$'$。})把它们视作公共财产(不言而喻,是平凡的),因而,到时候,以最大的漫不经心与最坦然的良心,把它们当作“自己的”——当然,前提是没有人想到要去据为己有一篇十页或一百页(或哪怕十行)的硬核证明——它确立一个“自己本来不会证明”的结果(注59$'$)。我不曾想到,自己对“形同虚设”这一点,竟感触得这样准、说得这样中,因为我有幸目睹:上面那个案子里那条模糊的“界限”,被人轻松越过——而且可以肯定,越过时同样怀着最坦然的良心,\textbf{鉴于其分量之轻:那是一个梦},而且甚至未经证明(尤其,未经\textbf{发表}……)\fnote{12}{(**)}{人们何必有所顾忌,既然事情本身分明表明:如今普遍的共识认为这完全正常——至少出自一位格调如此之高的人物!所谓“坦然的良心”,不多不少,正是这样一种感觉:觉得自己与自己所处行业里通行的共识相合。}。

幸好我还有点招架之力——必要的时候,我多少能表达出我的感受与想说的话;我(不论对错)挣得了一份信誉,因而,当我有话要说的时候,有一线希望被倾听,或者,若我觉得有必要,把它发表出来。相反,我此刻更痛切地体会到那种“不义与无力之感”——当一个人\dmark{}\origpage{220}受了损害却无处申诉,在“一切尽在掌握”的人们的专断面前,自觉手脚都被捆住——而后者还随心所欲地支配一切。

诚然,在我做数学家的生涯里,我也曾以同样坦然的良心有过可鄙的行为;在反思中我已有机会谈到那些个案——它们是这场反思从遗忘与从未检视过的暧昧的迷雾里重新唤起的。探查它们,我终于懂得:今天(而且由来已久)学生轻松愉快地超过老师,我无须惊讶;凡是与我有感情或情谊相连的人,我也无须否认。但对我、对所有的人来说,把猫就叫作猫,是健康的——无论这只猫出自我家,还是出自别人家。

\subsection{据为己有与轻蔑}

\notehead{59$'$}{}

(6月8日)对这件事,就我的朋友Pierre~Deligne而言,我已完全不这么确信了,因为我有机会确认:他终于滑进了对$\ell$-进上同调工具的“默认父权”游戏——也就是我所说的对平展上同调的“掌控”。从“SGA~4½之役”(其中我的名字还被提起,只是带着一种做作的、漫不经心的轻蔑,轻蔑于我著作中的这一核心部分——他自己的那一部分正源出于此),到“葬礼悼词”(其中凡“上同调”一词本身,凡与我的名字相关处,一概禁用),有一段令人瞩目的演变。(初期阶段见“白板”与“格格不入者”两注,终局阶段见“葬礼悼词(1)、(2)”两注。)

作为这场步步升级中的几个中间阶段:1981年,有那篇“值得纪念的文章”,论的是所谓“反常”层(对此可参见“不义——或归来的意义”与“暂停!”两注,n\textsuperscript{o}~75与77);次年,有LN~900里Motif的开棺出土(而“葬礼悼词”则再下一年,即1983年登场)。在我得以观察到的所有这些情形以及一些规模较小的情形里,那种使Deligne能够以完全坦然的良心据有他人观念之信誉的内在态度与“方法”,就是\textbf{轻蔑}(它部分地保持缄默,却又被巧妙地暗示出来)——对着人们即将据为己有的那个“轻”东西:它“轻”到甚至不值一提,尽管人们转过头去就要用它做出真正有力的东西——Weil猜想,所谓“反常”层的理论……一旦战役完成,据为己有成为既成事实、并为众人所接受,总有的是时间来修正口径,带着已据为己有的东西,谦虚地招摇过市。同一项贡献,只要似乎还沾着一个有待下葬者的名字,就遭到漫不经心的轻蔑;\dmark{}\origpage{221}一旦被他本人据为己有,就被捧上云端($\ell$-进上同调,Motif,Mebkhout的瑜伽尚待其时),或者被某位好哥们据为己有(导出范畴的瑜伽,对偶的瑜伽——由Verdier据有,Deligne积极怂恿)。

\section{V.~我的朋友皮埃尔}

\subsection{孩子}

\notehead{60}{}

\dmark{}(4月21日)再拾起这个关于回忆的梦——它不只是对一个视野之诞生的回忆……\origpage{223}我清楚地记得(尽管我忘掉了那么多事情!)那种每一次都焕然一新的乐趣:同那个很快就不止是心腹之交的人交谈——对于一切令我困惑的东西,对于日复一日在我与数学的恋情中豁亮起来、令我心醉的东西,他比任何“学生”都更是知交。他那份始终警醒的兴趣,他了解任何事情的从容(“仿佛他一直就知道……”),对我而言,是一泓永不枯竭的欢喜之泉。他的倾听是完美的,由那种也在我身上跳动着的、求理解的渴望所推动——一种高度醒着的倾听,是心心相印的标志。他的评点,总是迎着我自己的直觉或保留而来;即便并非如此,它们也往往在我力图穿透重重迷雾去把握的那个实在上,投下一道意想不到的光。如我在别处说过的,我提出的问题,他往往应声就有答案,或者在其后的几天、几星期里把答案发展出来。这就是说,倾听是双向的:当他反过来向我解释他找到的答案时——也就是解释事情的所以然,而那所以然,总是带着那种完美的自然、那种从容显现;同样的从容,曾常令我心醉地见于我的某些长辈,如Schwartz与Serre(同样也见于Cartier)。正是这同一种单纯,这同一种“自明性”,是我在理解数学事物的过程中始终追求的。无须明言:凭着这条进路与这份要求,他和我显然“同属一家”。

从我们相识之初,我就分明感到:他那(人们所谓的)“天资”,有一种极其少见的品质,远在我所拥有的那份寒微的手段之上;而就求理解的热情、就对待数学事物之理解的那份严格而言,我们却在同一个音高上。我还隐约感到——当时我尚无法把它表述出来——我在他身上觉察到的这股“力量”(我在自己身上也感到,只是较弱):那种“看见”人人视而不见的自明之事的力量,乃是童年的力量,是孩童眼里的天真。他身上有某种属于孩子的东西,比我认识的任何别的数学家都更外露;这绝非偶然。他给我讲过:有一天,我记得还在中学,他曾兴致勃勃地用定义去验算乘法表(顺水推舟,连加法表也一并验了)\dmark{}\origpage{224},对1到9这些数。他当然不指望有什么意外——如果说真有意外(愉快的,一如既往……),那就是:这证明居然可以在短短几页里、漂漂亮亮又完完整整地做成,前后不过半个钟头的事。当他笑着给我讲这件事时,我分明感到,那半个钟头用得其所——这件事,我今天比当时懂得更透。这段小故事当时打动了我,甚至震动了我(我想我未曾形于颜色)——我从中感到一种\textbf{内在的自主}、一种相对于既有知识的自由的征兆;这种自由,在我童年与数学的关系中,自最初接触起就已在场(注69)\fnote{13}{(*)}{而且我觉得,这份自由在我做数学家的一生中从未全然隐没,如今又像童年时那样在场了。两三年前,我向我的朋友重提乘法表那桩童年小事。我感到他为这一回忆的唤起而局促——它分明已不再符合他心中的自我形象。这份局促我并不真正惊讶,却仍然难过:又一件我明知道、却仍难于承认的事,再次得到了确认……}。这种彼此互为特权对话者的关系——我想,那会儿我们几乎天天见面\fnote{14}{(**)}{至少在我住在比雷时是如此,他当时住在IHES的一间单间。自1967年(我搬到马西)之后,我想我们仍一周会面一两次,至少直到我不再投身数学为止。}——持续了五年,从1965年(如果我没记错)到1969年(含)。我还记得那一年写一份详尽的评述报告时的乐趣:当时我提议把他增选为我自建校起(1958年)就一直工作的那个机构——我的大部分数学著作就是在那里完成的——的教授。那份报告我没有存本了(注64):我在其中评阅了我想有十几篇我朋友的著作,当时几乎都未发表(有许多后来也没发表),而在我看,其中大部分——如果不是全部——都当得起一篇优秀国家博士论文的主体分量。呈递这份雄辩的报告,我感到比呈递一份关于我自己工作的报告更骄傲、更欢喜(那种事,我一生只做过两次,而且每次都是勉强自己……)。其中许多著作,是对我提出的问题的回答(其中唯一已刊的,是前面提到的那篇关于概形真光滑态射之Leray谱序列退化的工作(注63))。而最重要的两篇\dmark{}\origpage{225}却是对Deligne自己提出的问题的回答;显然,其分量与一篇“优秀的国家博士论文”不可同日而语。那就是他关于Ramanujan猜想的工作(原书拼作Ramanuyam)(发表于Bourbaki研讨班),以及关于混合Hodge结构的工作——亦称“Hodge--Deligne理论”。

有一件事很奇怪,而且是我写这份流光溢彩的报告时远远没有料到的:不到一年之后,我就要离开那个机构——我正准备在那里增选我年轻而令人惊叹的朋友,并且打算在那里终老。还有(如今我把这两桩双生的情节放在一起想)另一件怪事,也绝不像是单纯的“偶然”:一两个月前,正是这位(如今已不年轻!)朋友,向我宣布了他自己要从同一机构离去——而恰好一年之前,我才重新拾起一种规律的数学活动,以一种不期然的“重返”走上数学舞台(虽还不是“上流社会”……)

不止一次,我在《收获与播种》里有机会谈到我的离去——谈那次“有益的挣脱”——以及紧随其后的“觉醒”;正是后者使那一情节成为我生命中的一个关键转折。在其后那些紧张的年月里,数学家的世界,连同我在其中爱过的人们,以及数学本身最令我心醉神迷的那个东西,都变得极其遥远——仿佛淹没在另一个“自我”的回忆的迷雾里,那个自我仿佛死去千年……

但无论在这桩情节之前,还是在这第一个大转折之后的年月里,我都知道:那个曾是(一点点\fnote{15}{(*)}{至于我为何有这份顾忌,不肯把(太!)耀眼的Deligne当作我的一个学生,见“格格不入者”一注(n\textsuperscript{o}~67)。})我的学生、(在很大程度上的)心腹之交与朋友的人,只须顺着体内一个游戏着的、渴求认识的孩子那股自发的冲劲,就能发现并唤起崭新而出人意表的世界,就能探测它们、认识它们内在的本性——并由此把它们向他自己的同类、也向他自己揭示。所以,在我离去之后(绝无回归之念!),倘若我看见“一位勇敢的数学家”、一位富有灵感的人,开始大笔勾勒(权作开端……)我隐约瞥见的那幅宏阔图景——而我尚且只画下了一组局部的、临时的草图——那个人只能是他:他把做成这件事的一切都握在手里!勾勒这幅大面积的第一图景——一位“总设计师”把关于\dmark{}\origpage{226}代数簇上同调的已知与猜出的本质事物,纳入一个共同的视野——对于那个心中视野已然呼之欲出、只待破除“尚未写下”之迷雾的人,是几个月的工作,甚至用不上几年。(不妨在此后经年、甚至经代——倘若需要几代的话——将它重做、深化,直到Motif之实在的底细被充分理解与确立。)而且我毫不怀疑,这项不久前还“在我手里灼烧”的工作,随时都会有人去做,至少在其后两三年内,趁着它还热乎。我离去之后,的确只剩一个人,凭着他对认识的冲劲本身,注定要做这件灼热而迷人的工作。不妨,在“总设计师”写就并经检验、大厦的营造或深或浅地推进之后,把继续这项工程的差事——无论它多么迷人——留给别人,自己投身别的冒险,在那个数学事物的世界里:那里的每一个转弯,都显露着一个崭新而无边世界的许诺,只要我们有睁开的新眼睛去看……

当我的生活还在那座温暖的科学暖房里流转——它把尘世的喧嚣隔绝在外——而当Deligne发展他对Hodge理论的推广时(那应当是1968或1969年),在我们之间,有件事不言而喻:那项工作,是把“Motif图景”的某个\textbf{部分}加以实现、检验与明确的一个头等步骤;而那幅图景,作为一个整体,从未被白纸黑字写下\fnote{16}{(*)}{至于那套Hodge--Deligne理论后来(据我所知)从未超出这第一次尝试的阶段,从未扩展成一套“Hodge--Deligne系数”(及其上的“六种运算”)理论、架在复数域上有限型概形之上——这与另一件怪事密不可分:那幅宏阔的“Motif图景”从未被勾勒出来,而它的存在本身,直到今天还被小心地讳莫如深……}。在我告别暖房之后的年月里,当数学对我已相当遥远时,我得知Weil猜想最终获得证明,毫不意外。(若说有意外,那是“标准猜想”没有在同一步里得到证明——尽管人们提出它们,恰恰就是为了逼近Weil猜想,同时也是作为在一域上至少建立半单Motif理论的一个手段\fnote{17}{(**)}{只是在这最后几年里,我才模糊地(但近来更精确地!)意识到:“标准猜想”,连同它们为之提供了第一条“构造性”进路的Motif观念本身,都被\textbf{埋葬}了——缘由如今在我眼前显得格外清楚。(亦可与前面一条脚注对读。)}。)我很清楚:无论是这第一击\dmark{}\origpage{227}——那初试身手的、Hodge式的一般系数理论——还是对若干关键猜想(在若干其他多少为人所知的猜想之中)的这项证明,都还没有让他使出全副本领——差得还远。而我并不着急,因为我的注意力的绝大部分,已倾注在别处。(→注61)

\subsection{下葬}

\notehead{61}{}

我有幸亲眼见到一股童年的冲劲初次绽放,带着大面积舒展的许诺。在其后的十五年里,我渐渐确认:这一许诺始终被不断推迟。他身上有那样一种精微的东西,我曾懂得去感觉、去辨认(尽管那会儿我对那么多事情都无感!)——一种与脑力全然异质的东西(脑力既碾压人也穿透人……),一件对一切真正的创造性的工作而言头等本质的东西。这种东西,我在别人身上有时也感到过;但在我认识的任何数学家身上,它都未曾以可与之相比的力量显现。而我(想当然地)指望,这东西会继续在他身上绽放、转化,并毫不费力地经由一部独一无二的作品表达出来——在那部作品里,我不过是一个寒微的先行者。然而又一桩怪事(而在这么多“怪事”之间,必定有一条简单而深刻的纽带):我眼看这“精微的东西”、这既非膂力亦非脑力的“力量”,随年月流逝而渐渐消退,仿佛被\textbf{埋葬}在一层层越来越厚的沉积之下——那是一层层\textbf{别的东西}的沉积,我对它们知之甚稔——世上最寻常的东西!它未必与脑力、与老到的经验或在某一学科里磨炼出的嗅觉不能相安;这些东西,凭其作品的积累——也许光灿夺目,肯定自有其力量与美——可以博得一些人的叹服、另一些人的畏惧,或兼而有之。但这并不是我当初说“舒展”或“绽放”时所想的。\textbf{那个}绽放,我所想的那一个,是天真的果实——渴望认识,并且随时准备为这个取之不竭的世界里或大或小的事物的美而欢喜,或为它的某个部分(比如数学事物的那个宏阔世界……)的美而欢喜。唯有它,才有深刻革新的力量:无论是自我的革新,还是对世上事物之认识的革新。在我看来,把这一点实现得圆满的,正是谦逊的Riemann\fnote{18}{(*)}{Riemann(1826--1866)的著作,装满区区一册薄薄的文集,约十篇论文(诚然,他四十来岁就去世了),其中大部分包含一些简单而本质的观念,深刻革新了他那个时代的数学。}。这份真正的绽放,与轻蔑无缘\dmark{}\origpage{228}:对他人(对那些自感远在自己之下的人……)的轻蔑,对那些太“小”、太自明以致不屑一顾之事的轻蔑,对那些被认定低于自己正当期望之事的轻蔑;或者还有对某个\textbf{梦}的轻蔑——那梦执拗地向我们诉说着我们口口声声说爱的事物……它与轻蔑无缘,正如它与滋养轻蔑的自负无缘。

诚然,凭着那令人叹服的“天资”,但更凭着那件不打动任何人、却\textbf{创造}着的精微之物,“学生”注定要远远超过“老师”。我毫不怀疑,在我离开那个地方——我在那里曾见证一次如此美丽的起飞——之后的年月里,Deligne将使一部宏阔而深刻的作品得以全副展开,而我曾是它的先行者之一。这样一部作品的回响,岁月流转,不会不传到我这里;而我自己,在远离数学的其他追寻中跋涉,对他将要发现的那些崭新世界的全部分量与全部美,只能有不太完整的领会。

但学生若在内心深处\textbf{否认}老师,努力在人前人后抹去老师的一切痕迹(无论这份馈赠是为了更好还是为了更坏……),学生就无法真正超过老师——正如儿子若否认父亲,就无法真正超过父亲。这一课,我主要是经由我与我的孩子们的关系学到的,但后来也经由我与昔日某些学生的关系;尤其是与那一个人——在所有人之中,我始终有所顾忌、不肯以“学生”之名相称的那一个;因为从相识那一刻起,我就分明感到,我有东西要向他学习——一如他也有东西要向我学\fnote{19}{(*)}{(6月14日)关于我身上这份执拗的蓄意——贬低我所能给的东西、否认一段师生关系的实在——见“格格不入者”一注,n\textsuperscript{o}~67$'$。明摆着:我的朋友从我的接触中所学到的东西(“仿佛他一直就知道”,当然!),与我经由他所学到的,没有共同的尺度。倘若我继续高强度的数学投资直到今天、我们之间保持着规律的数学接触,情形想必会是两样。}

但只是在相识将近十年之后——在1975年之后,尤其自从我开始沉思我所经历的、我所目睹之事的意味——我才开始感到那个继续为我所珍的人身上的这份\textbf{枷锁}。我也隐约感到:这份对我这个人、对我在其人生关键年月里所扮演角色的秘密否认,更深地,也是对他\textbf{自己}的否认。(想必每当我们否认、并想抹去\dmark{}\origpage{229}某件分明发生过、而其果实有待我们去摘取的事情之时,都是如此。)

然而,由于对“数学界时下在做什么”、对他自己在其中所做的事情\fnote{20}{(*)}{自1970年以来,我收到过Deligne的四份抽印本,都当即匆匆翻过(一如我至今仍不时收到的大多数抽印本)。凭这点东西,难以对一部数学著作——哪怕只就其总轮廓或主要主题——形成概念。}多少不“在局中”,在几个星期前的反思之前,我从不曾掂量:这份枷锁,对他倾其所有投注的事情——他的数学著作——也压得有多重。诚然,八九年来,我不止一次看到,一份蓄意的轻蔑口吻(对我)或鄙夷(对一些他有能力使之气馁的别人)仿佛把单纯的常识、把数学家的健全本能一笔勾销(注66)。而且他并非我旧学生中——带或不带引号——唯一一个让我目睹这类态度对着我所珍视的人(或对着别人)的。但在别人身上,它从不曾这样痛切地触动我。在刚过去的两个月的反思里,我不止一次影射过这份经验——“它是我做数学家的一生中所曾经历的最苦的经验”;我也说过,这场《收获与播种》的反思到头来就它教给了我什么。那份创痛如此剧烈,它就一个始终为我所珍的人教给我分量如此之重的东西(而我继续回避它同时就我自己、就我的过去所教给我的……),以至于这种态度会对他、乃至对那些被弄得气馁或受了羞辱者的数学“创造力”产生多大影响,反而成了全然次要、甚至不足挂齿的问题。

“拒受遗产——或矛盾的代价”那条注,是我头一份写下的反思:盘点经年累月、此处彼处以断片形式回到我这里的东西——既关于“行内现状”,也关于那位我认识得那么好又那么少的人的著作。那也是头一回,我在一眼之间终于看清:那份他携行了大概十五年以上的拒绝,在其数学著作本身中付出的全部“\textbf{代价}”、造成的全部重负。写那条注时,我其实还“落在后面”:因为两年以来(而且没有“人”认为值得通知我),Motif早已走出它们被囚禁十二年的秘密……而今天,在我写下这场对我数学家过去的反思的这一(我相信)最终阶段之时——在我\dmark{}\origpage{230}大致了解那部标志着这次悄然“重返”的值得纪念的文集之后两天——对这份沉重负担的知觉,变得触目惊心:那正是生来该飞翔的人所乐于日复一日、千回百转拖着的重量——一次轻柔轻盈的飞翔,欢快而无畏,迎向未知,为了他的喜悦,也为了托起他的风的喜悦……\fnote{21}{(*)}{我绝无意暗示:“飞翔”、发现世界,是少数特殊者的特权;肯定,我们人人天生受此召唤!只是这份能力,很少有机会哪怕在极有限的方向上(比如数学工作)绽放。但在某个人身上,我曾有幸看到这样一份能力(在“数学”的方向上)格外夺目,如奇迹般被保存,其后却随岁月而倒退。}

倘若他不飞翔,倘若他满足于做一个受人仰慕与畏惧的人,靠不断积累自己优于他人的证据,我无须为他担忧。倘若他拖着他乐于拖的重量,他肯定从中找到满足——一如我自己曾乐于拖重量,并且今天仍在拖那些我在路上尚未学会割舍的重量。我所能给他的,无论好的一面还是坏的一面,他取其所愿。他的选择只属于他自己,我无须为其担忧;甚至无须在此裁决它们是最好的还是最坏的(注62)。对一个人是“最好”的,对另一个人是“最坏”——有时对同一个人也是(只要他肯变,这诚然少见……)。

但我们所做的选择、表达这些选择的行为(即便我们的言语常常否认它们),我们都须自行承担风险与后果。它们常带给我们所期待的报偿(我们把它当作“最好”来领受);而这些报偿本身,有时终于显出其反面(我们把它驳回为“最坏”,甚至当作一桩冒犯)。当人们终于懂得:反面并非冒犯,往往便把它看作一笔该付的价钱,付得老大不情愿。但也会有这样的领悟:某些反面,并不是无情的收银员——为了快活时光,不情愿也罢,总得付账;它们是耐心而执拗的信使,不倦地回来,永远带着同一个讯息;一个不合时宜、不断被拒的讯息——因为比反面本身更甚的,是它那份谦卑的讯息总遭驳回,那份讯息在我们眼前才显为“最坏”:比一千个反面更坏,常常比一千次死亡、比整个宇宙的毁灭更坏——对整个宇宙,我们已经毫不在乎……

\dmark{}\origpage{231}终于有一天,当我们乐于接纳那份讯息,眼睛霍然睁开,看见:曾被畏为“最坏”的,是一场\textbf{解脱},一场宏大的释负——而那份压得我们喘不过气、我们蓦然卸下的重量,恰恰是昨天我们攀住不放、当作“最好”的东西。


\subsection{事件}

\notehead{62}{}

(4月21日)人们会对我说:既然我无须为他担忧,那为什么还要为一个只关乎当事人与我俩的个人关系,铺陈上这么多页!

我之所以感到需要这场对一段关系的某些重要侧面的回顾性反思,是一个具体的、与我切身相关的事件使然(尽管我是带着两年的时差才得知它的)。另一方面,这一事件处在公共领域之中,比那些成名的数学家(如Deligne,或我自己)对待名气较小者或初学者的日常行为举止,还要显豁得多(虽然后者对他人生活的影响,往往与眼下这个个案完全不在一个量级)。所说的事件(即Lecture~Notes的LN~900那部“值得纪念的文集”的出版,别名“下葬文集”),连同围绕着它的一切,在我看来——无论对错——是\textbf{不健全的}。在我看来,给某些来龙去脉作一份周详的见证,直探我今日所感知的事情的根底,对所有人都是有益的,首先要从“当事人”本人做起。经由这一见证与这场反思,我并不想说服任何人相信任何事情(那实在太累人,何况毫无希望!\fnote{22}{(*)}{(5月25日)如果我在这里觉得有必要重申,想说服人是“太累人”而且“毫无希望”的,那恐怕是因为:在我身上某个地方,说服的意图毕竟实实在在地在场,而且也被觉察到了。从4月19日(我得知“值得纪念的文集”LN~900那一天)到4月30日的整个反思,都带着一种内在紧张、乃至分裂的状态——面对一桩全然出乎意料的“事件”的冲击,我尽力消化它带来的讯息。这一紧张,最终随着4月30日的“事物的回归”一注(n\textsuperscript{o}~73)而解开:反思终于回到我自己的身上,随即为我提供了这则讯息的明摆着的钥匙。}),而只是想去理解那些我身陷其中的事件与境况。倘若它们能促使别人在老生常谈之外作一番真正的反思,这份见证就不算白写。

\subsection{驱逐}

\notehead{63}{}

\pmark\origpage{232} (4月22日)这篇文章\fnote{23}{(*)}{即Deligne论谱序列退化与Lefschetz定理的文章(Publications~Mathématiques~35, 1968),“罐装的权与十二年的秘密”一注(n\textsuperscript{o}~49)中已提及。}1968年发表于《Publications Mathématiques》,即我离开数学家世界的前两年。它的出发点曾是一条猜想——我同Deligne谈起过:关于谱序列的一个退化性质,在当时看上去相当不可思议,却经由“算术”的途径、作为Weil猜想的一个推论,变得可信起来。这一动机本身就有极大的趣味,因为它表明:从隐含在Weil猜想里的那部“权的瑜伽”(这部瑜伽最早由Serre在某些重要方面隐约瞥见)里,人们能榨出多少东西。从那时起,我就经常把它应用于各种类似的情形,从“算术的”论证榨取“几何的”(对代数簇的上同调而言)结论。只要Weil猜想尚未确立,这些论证就始终是启发性的,但它们仍然具有巨大的证明力,代表着一种头等的\textbf{发现手段}。Deligne对这条特殊猜想的“几何”证明——借助Lefschetz定理(当时只在特征零确立)——在一个全然不同的方向上另有其趣味,且不提它的头一项功劳:不依赖任何猜想。两条进路所指示的那种联系——一边是Weil猜想(以及当时在我看来其最迷人方面的权的瑜伽),另一边是Lefschetz定理——联系着两件看似互不相干的事情;这联系本身就极富教益。

对我眼下的处境而言,这里有趣的一点——直到今天它才向我显出全部意味——是:这篇文章的读者,几乎无从猜到我参与了其主结果的最初动机,更绝无可能从这篇文章里得知那动机是什么。(另可参见注49开头。)为这样一个结果作陈述,\textbf{自然而然}的路数(我确信,连作者本人最初也是这样)本应是:从那条(的确引人注目的)猜想出发,指出找到的第一条理由——它同样引人注目,而这正是终于把那部著名的权的瑜伽“推销”出去的大好机会,它本身的分量远比该文的主结果更重\dmark{}\origpage{233}\fnote{24}{(*)}{恰恰是这部瑜伽,此后六年里(据我看来)一直保持着秘密!(6月7日)而且(如后来所显现的),它当时被Deligne“记在自己名下”地呈献出来,对Serre与我都不置一词。(见注n\textsuperscript{o}~78\textsubscript{1}、78\textsubscript{2}。)};然后接着谈“Lefschetz定理”的观点\fnote{25}{(**)}{(6月17日)用(“奶牛”)Lefschetz定理来证明谱序列的退化,这一想法出于Blanchard;不过他只在一条苛刻得罕见地得到满足的假设之下——即纤维的有理上同调构成的局部系统是平凡的——才得到退化定理。我熟悉Blanchard的工作,自然向Deligne谈起过;他因此从Blanchard的想法取得启发来完成他的证明,尽管他没有读过那篇文章。Serre对Blanchard证明的记忆比我好,他向Deligne指出:他的证明实际上是Blanchard证明的一个不费力的改写。Deligne在其评注2.10里承认了这一点。但这条引用了Serre的评注,写法却给人这样的印象:仿佛他是事后才得知Blanchard的想法的——事实绝非如此。于是,这篇文章的两条主要\textbf{源头}都被藏掖了:一边是\textbf{算术的动机},它使人得以预见Blanchard结果的一个重大强化;另一边是Blanchard的\textbf{证明想法}——他把后者优雅地加以改写,得到一个Blanchard恐怕不敢指望的结果,也因此从未试图用自己的方法去“得到”它。},它使人得以在略更一般的条件下(基底概形任意,不必真且光滑于一个域)证明最初的猜想,但只限于特征零。而文章实际采取的陈述,却从同调代数的一般的漂亮道理开始(想必十分漂亮,而且以作者惯有的优雅陈述出来)——那些道理,他想必早已同所有人一样忘光了,俨然是Lefschetz定理公理化的路数。主结果(当然,也是唯一人人记得的结果)以系X的面目出现在文章中部;而快到结尾的某处,在“评注2.9”里,读者不知为何,“权”这个词与我的名字才被提起……

文章发表时给我的印象,我已毫无记忆——我在局中,想必只是匆匆浏览了一眼。我肯定感到了一种“保持距离”的意图,但也感到:我的朋友不愿冒被人看作某位“大师”的信徒(或“驹子”)的风险,这原是再自然不过的事\fnote{26}{(***)}{(5月26日)关于我身上的这种态度,可参见紧随本注之后的“高升”一注(n\textsuperscript{o}~63$'$)。(6月8日)把它与我朋友那套据有他人观念的独门手法放在一起看——这里就是第一个典型的例子——我进而明白:我朋友的动机,根本不是要在一个显赫的“大师”面前保持什么“自主”,而恰恰是要藏掖他人的观念在其观念之生成中的作用,以待来日把这些他人的观念也一并据为己有(在第二步里)。(对此可参见“变戏法的人”与“据为己有与轻蔑”两注,n\textsuperscript{o}~75$''$与59$'$。)至于在我的朋友身上,这一癖好得以畅行无阻地发展,我自身应负的那一份责任,可参见“高升”“暧昧”与“格格不入者”三注(n\textsuperscript{o}~63$'$、63$''$、67$'$),其中显出我对这位光彩夺目的年轻人Deligne所曾表现出的某种纵容。}。诚然,倘若他心中有那份对自身力量的安然自信\dmark{}\origpage{234},他原会毫不犹豫地写出一部分量更重、对人人更有用(肯定也对他自己更有用)的著作,而不必担心别人看不出他的本色……(注65)

次年,他第一部大面积的著作——论混合Hodge理论——发表时,情形也多少类似。(我当时认为,这项工作的分量,可与Hodge理论本身相比,把它看作一套“Hodge--Deligne系数”理论的起点——那套理论可惜始终不曾问世……)如我所说,对他和对我都再明显不过:这项工作的“动机”在于我此前数年抵达的Motif瑜伽——它是使这部瑜伽获得有形实现的第一番尝试。在他的著作里点明这层联系,依我看(当时想必也这样觉得),本会立刻使这项工作获得比它凭自身优点已具有的更宽广的分量;同时,这又是一次机会,把读者的注意力引向Motif的实在——它在Hodge结构的实在之后,随处可以感觉到(注63$_{1}$)。

只有带着事后的认识,这些省略才在六年的权的瑜伽沉默\fnoteplainmark{27}{(*)}{(1985年4月19日)关于“六年”与“十二年”的更正,对权而言见第302页脚注(***)(1985年4月18日所写的部分)。}、十二年的Motif沉默(不妨说,是禁令)\fnoteplainmark{28}{(*)}{(1985年4月19日)关于“六年”与“十二年”的更正,对Motif而言见子注“预先出土”(n\textsuperscript{o}~168$_{1}$)。}之后,在这场以LN~900下葬文集为形式的罕见的“重返”,在Hodge--Deligne理论惊艳的开局之后的停滞之中,显出全部意味……但没有人能怀着殡仪人的心肠做出大事!

无论如何,倘若1970年我离开IHES时心智更成熟一些,那么从那一刻起,我就会一清二楚:在那个五年间一直扮演着我最亲近之友的角色的人身上,存在着一种对我而言深刻的暧昧。况且,在一个平静安适的机构里,在同侪间彬彬有礼的情面外观之下,我的离去,最终对人人都是方便的——缘由我自信事后已能辨出,而且各人并不相同。\dmark{}\origpage{235}显然,这一离去对我的年轻朋友再合适不过:他新近就了位,本只需同我站在一边(面对另外三位终身同事迟疑的冷淡),就足以翻转一个未定的局面。当时我不明白这一切意味着什么,是因为我分明不愿去理解那些其实足够清楚、甚至足以自明的事情!一如我生命中常常发生的那样,那时我心中有一种焦虑(我从不曾以这个名字称呼它!),它向我提示:在一个有形而简单的实在,与一幅我不肯放手的实在图景之间,出现了“脱节”:即我在即将告别的机构里所扮演角色的图景,而更甚的,或许是我与朋友之间那层关系的图景。正是这种拒绝面对一项无可置疑的实在,以及作为这桩矛盾之征兆、我却死死攀住的焦虑,使得那次“有益的挣脱”的情节,在当下显得如此痛苦\fnote{29}{(*)}{关于这一情节,见注n\textsuperscript{o}~42。}。

老实说,由于从未曾为这段关系专门写过一篇反思(除了一些散见于给我朋友的书信里的反思的开端,而且没有一封得到回音……),我此前不曾注意到:我的朋友对我之关系中矛盾态度的最初征兆(诚然隐约,却绝不会错)至少可以上溯到1968年,即“大转折”前两年。那是一个这段关系显得完美无缺的时刻:在数学层面上是无云的契合,处身于一段单纯而亲切的友情之中。人们此刻尽可以讥讽那些关于天真、创造的孩子的美好“长篇大论”!

然而我深知,那种契合是一个\textbf{实在},绝非幻象;一如那“精微的东西”是实在——那股创造的力,其后的作品只给了它一个苍白的反照。“天真”与“冲突”是两个可以触摸的实在,略一警醒的知觉就能辨认,绝不是什么概念;而且在我看来,它们按其本性相互排斥,一个排斥另一个。然而毫无疑问,这两个实在,在不同的层面上\fnote{30}{(**)}{在另外两三次机缘里,我得以确认:在同一个人身上、在某个时刻,竟有这等共存——包括我自己的身上,在有些时刻。},共存于我朋友对我的关系之中。在我所说的那个年代,“冲突”似乎并未干扰数学的创造性——至少在\dmark{}\origpage{236}孤独中做的工作里、或在面对面的交谈中进行的工作里没有。诚然,在我方才谈到的那两篇文章里——它们毕竟是这项工作最切实的果实之列——“冲突”的印记已经分明可见。而带着十五年的事后的认识,经由这些天与这些星期的反思,我看清了:这道印记(无论多么隐约),以一种触目惊心的方式,预演了冲突将逐步压倒那股最初的冲劲、并在岁月里剥除其最稀有的本质——那造就大命运的东西——时所采取的特殊形式\fnote{31}{(*)}{这一节如此高贵的抒情,使我有点脱离了尘世的实在。倘若我在这里称这道“印记”为“隐约的”,那是因为我自己还裹在一层厚茧里,还难舍那副我视若珍宝的眼罩!一旦终于摆脱了它,我便看清:所说的“印记”,是一场粗劣的藏掖;我不肯看见它,是出于我身上的某种纵容——6月1日的“暧昧”一注(n\textsuperscript{o}~63$''$)里,我对此有清楚的认识。至于“冲突对我年轻而光彩的朋友那股最初冲劲的压制”,我几乎把它说成一场令人惋惜的宿命,仿佛这位可怜人是全然不自愿的受害者,从而——唉!——也一并取消了“大命运”的收益。然而,他对自己的命运,负有与我对我自己的命运同样的责任。倘若早在我离去之前,他就选择了做他导师的掘墓人(权且从这里开始)的角色,倘若种种情势(其中包括时代精神)有利于这一选择,把“打谁都行的大老板”的角色让他尽情享受,那么,他也同样选择了把威望与权力所能给人的特权——包括(悄悄地)碾压与巧取的特权——品尝到最后一滴。鱼与熊掌不可兼得;按照事情的本性,经由这一选择(在这选择里,他自有一些好同伴),他失去的是那些更精微、更不趋时的东西的收益……(6月初的一则未注日期的脚注)。}

\notehead{63$_{1}$}{}

(5月26日)亦可与注60末尾脚注(*)处的评语对读:那里确认了Hodge--Deligne理论的自然的发育所遭到的“封锁”——由于对我引入的某些观念核(这里是六种运算——Motif与它们密不可分)的拒斥态度;那态度与这里考察的同属一类,因此早在“Hodge理论I与II”发表时就已显形。

同样这种态度——竭尽所能(乃至过之!)抹去我的影响的任何痕迹——也见于与Mumford合作的那项工作(注n\textsuperscript{o}~47已提及):论模多重簇的Mumford--Deligne紧化。(这项工作同样先于我的离去。)它利用一条从$\mathbb{C}$上的拓扑结果(经由超越途径得知)过渡到特征$p>0$的结果的原理——那条原理,是我在五十年代末为基本群理论引入的\dmark{}\origpage{237}。早在六十年代初,我就建议利用这一方法,在全特征下证明模簇的连通性\fnote{32}{(*)}{(1984年9月)核实之后,这一情节确确实实在前引工作的导言里(第75页)得到了记载。}。然而这一想法碰到了一些技术困难,它们曾使Mumford却步,而在他们那项工作里被优雅地克服:办法是引入\textbf{模多重簇},以及一种具有完美性质的“紧化”。模多重簇这个想法本身,至少“在字里行间”,可以在我在Cartan研讨班上的“Teichmüller”报告里找到——那还是site与topos的语言尚不存在的年代。Deligne使用的语言本身——在他本可使用一整套为表达这类情形量身定做的site、topos、多重簇的语言的地方,却用了“algebraic stack”——相当清楚地显示了(带着事后的认识,并且借着后来规模大得多的几场“战役”之光照亮)抹去这项光辉工作中所运用的某些主要观念之来源的意图。肯定正是这种态度(我在“拒受遗产——或矛盾的代价”一注,即n\textsuperscript{o}~47里,头一回预感到它)产生了那记“电锯效应”,斩断了其后对模多重簇的反思——尽管在我看来,模多重簇属于迄今所提炼出的一切“具体的”数学对象之中最美、最基本之列。

顺带指出:我在五十年代末引入的那些论证,使人得以(有赖于Mumford--Deligne紧化)不仅在全特征下证明模多重簇的连通性,而且确定它们的“素于$p$的基本群”:即寻常Teichmüller群的“素于$p$的投射有限紧化”。

\subsection{高升}

\notehead{63$'$}{}

(5月10日)带着再多不到三个星期的事后认识,我如今看清:我当时那种自以为“体谅”的态度——对着那“再自然不过”的保持距离的意图——实际上是一种洞察力的欠缺,是对我年轻而光彩的朋友的一份纵容。倘若那时我信任自己健全的知觉,而不肯被那些含糊的老生常谈弄得眼花缭乱、哄骗自己——它们摆出“体谅”的姿态,甚至摆出“宽宏大量”的姿态(“总不能因为他不把我的名字捧出来,我就说他几句……”)——我那时就会看清我如今(十六年后)看清的东西。我可以称之为:对读者、对我自己、也对他自己,缺乏一份正直。简简单单地看事情,毫不害怕把它们\dmark{}\origpage{238}叫出它们的名字,我本会有能力像现在这样坦然谈论它们;而我的朋友当时本可从中取法——或至少他会懂得:即便凭他那般手段,他的长辈们(或至少其中一位)也期待他在工作中具有与他们自己相同的那份正直。于是我看清:在那次机缘里——它在我退出数学舞台之前,彼时我绝非“局外”,而且对年轻朋友无疑还有某种道德上的分量——我经由当时表现出的这份\textbf{松懈}\fnote{33}{(*)}{(5月28日)这里用“纵容”一词,比那个有点躲闪的“松懈”,更能表达我那份态度的性质。在我与我年轻而光彩的朋友的关系中的这份纵容,在昨天的反思里显现得更清楚;见“格格不入者”一注,n\textsuperscript{o}~67$'$。},没有尽到我对他的责任。这一松懈,在“Hodge理论II”发表时得到了印证:那是Deligne的学位论文工作,其中对Motif、对我,都不置一词。诚然,到那时,数学和我朋友其人,对我都已十分遥远,仿佛隔着一层浓雾!

从我得以观察到的我朋友的演变——精神上的与数学上的(这两个侧面血肉相连)——我如今看清:在我与他相识、并被他的智力手段、他的目光的敏锐、他在数学上理解的敏捷所打动的那个时刻,我丝毫没有在他身上辨出成熟的欠缺;也没有(后来)辨出,他社会地位的扶摇直上——在不过四年的工夫里,从一个无名的学生,一跃而成为数学世界的明星、终身教授,在一个已然声名显赫的机构里被授予可观的特权与权力——将对他产生怎样的影响。我不后悔为他铺平了这条路,还让它走得更快——但我确认:那是出于我自身的眼力与成熟之不足。我为他尽的这份“效劳”,算不得效劳。只要我的朋友自己尚未把那场收成——他在我漫不经心的协助下为自己预备的收成——进行到底,这就不曾是一场“效劳”。

\subsection{暧昧}

\notehead{63$''$}{}

(6月1日)自从在我与我的朋友皮埃尔的关系中确认了那份“松懈”(或按其间出现的更贴切的说法,那份“纵容”)以来的三个星期里,我已有机会在反思中更清楚地认识到我身上的某种不够严格、某种纵容。它们最初显现于我与那个人的关系——在我那里,我比对任何人都更把他当作“格格不入者”来对待——但也显现于我与另一些数学家的关系,在后者眼里我扮演着长辈的角色。迄今我在这个方向上侦察到的\dmark{}\origpage{239},表现为我身上的一种暧昧——恐怕也在那位扮演学生角色的人身上:在那些场合,他把得自我的观念与方法、乃至我交给他的一项完整工作的详细施工图据为己有,而不明确指出来源,有时甚至只字不提。这类情形相当常见,既见于六十年代,也见于我离去之后直到最近这几年的年月。我觉得,在所有这些情形里,在某个层面上,我感到了那份暧昧;它表现为一丝不安的阴影,直到最近这些天才头一回被检视。使我入局、参与某种心照不宣,并使我掠过那丝不安、从不予理会的动机,在于\textbf{顺应}我对自己、对所谓“宽宏大量”应有的样子所抱的某种形象。真正的宽宏,并非生于一种顺应成见的心态,也不是出于想要成为“宽宏的人”、并在自己与他人面前显得宽宏的心思。被压抑的不安,每一次都是明明白白的征兆:这份“宽宏”是做作的,是一种\textbf{姿态},而不是宽宏之真正的、毫无保留的自发赠予。

在这份不安里,我辨出两种来源不同的成分。一种来自“老板”、来自那个始终不得满足的“我”:他没能两头都赢——既分享一项他知道其中有自己(或大或小)一份的工作的信誉,同时又配得上某种品牌形象——其中赫然印着“宽宏大量”这条老生常谈的标签(此外还有许多别的标签)。另一种成分来自“孩子”、来自我身上那个不为姿态与门面所欺的人:他朴素地感到这境况的虚伪\fnote{34}{(*)}{(6月5日)当我在这里说,不安(部分地)来自“孩子”时,这种说法给人的是一幅不真实的图景。并不是对一个虚伪境况的天真知觉造成了什么不安。不安是一个\textbf{抵抗}的征兆——抵抗那份知觉,是(在某些层面上分明被知觉到的)实在(这里是:一个虚伪的境况)与我所攀住的实在\textbf{图景}之间的“脱节”(在这个案子里,图景即:我正在“宽宏大量”,而且不能不这样!)——我为了图景的利益,把那份不合时宜的知觉\textbf{推开}、压抑。在这个案子里,一旦我放弃抵抗、容许知觉进入意识之眼的视野,“不安”便随之止息,连同那个虚伪的境况。我本想说:“这是就涉及我当下的虚伪境况而言,而不是就处在过去的境况而言。”但转念一想,我发现:我方才谈到的那些“过去的”虚伪境况,直到今天仍以本来的面目在身,或者至少直到三天前的反思——仅仅因为它们从未被检视、从而从未被解决。我成了它们的囚徒,以至于一有机会,就机械地重演同样的境况。我对自己的“观”的力量的知识(我在“欲与观”一节,n\textsuperscript{o}~36里谈过它)当时对我毫无用处,因为我缺乏对日复一日身陷其中的境况、对知觉与其“拣择”的不停息的游戏的留意——那孩子与老板的游戏使知觉沉默……}。不但对我自己虚伪,而且\dmark{}\origpage{240}对他人也虚伪。

总之,我的“宽宏”在于入局:入一场游戏,游戏中对方把他得自他人的观念呈献为自己的,也就是呈献一幅他自己与某种实在的图景——而他和我都心知肚明,那是假的。于是我们成了同谋,共谋一桩可称之为“作弊”的事,双方各得其所。至少按“我那个年代”通行的共识,这是一桩“作弊”——而那些共识,依我看,今天人们仍在嘴边念叨。倘若被当作“我庇护的人”的发现来使用的是别人的观念,我肯定不会入这样的局\fnote{35}{(*)}{“我庇护的人”这个说法,曾被我昔日的一个学生用来称呼我当时的、刚在数学里做出漂亮东西的一个学生;它曾使我咬牙切齿。然而,细想之下,我此刻正在考察的这份暧昧的境况,确立的恰恰是一种虚伪的关系:其中一方俨然是另一方的“受庇护者”。}。然而,我给予默许,让诞生于我的观念被呈献为他人之物,并不改变事情的本质——唯一的区别是:这回我们是两个人在作弊,而不是一个人。撇开涉及我个人的这一面(我自己参与一场作弊,一种违背我自己也自认服膺的共识的行为)不谈,有一点再清楚不过:怂恿他人作弊(即便这作弊看似只由我们自掏腰包——而事实绝非如此),或至少怂恿他人对一个他自己也假装服膺的共识持一种暧昧的态度、实则阳奉阴违,这里没有任何宽宏可言。真正的宽宏,本性上造福所有人,首先要造福它于其中显现的那个人、它所施与的那个人。我的这份暧昧态度——在人身上诱发或怂恿暧昧,并使我得以摆出“宽宏”的架势,而按正常的道理,对方这时应当显出几分作弊者的模样(而事实上我们俩都在作弊)——这份态度,对我既非善举,对他亦然。

只需把事情检视一番,明摆着的道理就会显现,甚至无须求助于经验、求助于“事件的教训”。然而,终究是事件把我引到了这番检视,使我终于发现一个明摆着的道理——三十\dmark{}\origpage{241}年前我同样有能力发现它,那时地平线上还没有出现一个学生,要跟我学一门手艺,并在与我的接触中浸染于这门手艺练习中的某种精神。我已有机会谈到工作之中的那种“严格”——我相信自己做出过表率(见“严与苛”一节,n\textsuperscript{o}~26)。但今天我还察觉到,在“工作”本身之外,也有一份严格之不足,表现为暧昧,表现为我所说的纵容。我觉得,我的这份暧昧,不是从我的任何一位长辈那里传来的——他们(我相信)对我都有与自己对自己相当的要求。在这一特定态度的暧昧之外,我还在我这个人本身之中辨出一种暧昧——在《收获与播种》第一部分,我已有不止一次机会谈到它。这份暧昧,随着1976年禅观的发现开始消解;而它的某些征兆——表现为已然成习的态度与行为(尤其在我与学生的关系里)——想必一直延续到今天。

显然,我身上这份暧昧,在我的某些学生那里找到了肥沃的土壤。那种经由默契做成的事,如今似乎已成了“上流”数学界习俗中的一个基调:在那里,浑水摸鱼(有或没有“当事人”的同意)、乃至有组织的劫掠(当放胆为之的人属于那不可撼动的精英时),似乎已成了寻常得无人再表惊讶的惯例——而同时,人人也都绝口不提。我身上的“老板”很想划清界限、揭发、愤慨——可这样一来,我不过是在延续我身上那份同样的暧昧,它的丰饶收成,我今天有目共睹。

\subsection{同伙}

\notehead{63$'''$}{}

(4月24日)\fnote{36}{(*)}{本注出自“本能与时尚——或强者的法则”一注(n\textsuperscript{o}~48)的一条脚注——在那则注里,我曾断言Verdier关于导出范畴的工作从未发表,却不知道其学位论文的一个“第0状态”已于1977年发表。至于Verdier围绕着那本应构成其学位论文工作的理论所玩出的那些古怪的腾挪,其全貌可参见“赊账的学位论文与全险”一注,n\textsuperscript{o}~81。}两天前翻阅刚收到的Mebkhout的一份抽印本,我撞见一处对 J.L.~Verdier一项工作的引证,题为“导出范畴,第0状态”,发表于SGA~4½(Lecture~Notes n\textsuperscript{o}~569,第262--311页)。我到今天才得知这一发表,情有可原:此前我从未有过捧起这一卷的荣幸——无论Verdier,还是其作者Deligne\dmark{}\origpage{242},都不曾费心在出版时或后来寄我一本。我不知道,当年与我共同组成评审委员会、仅凭一篇17页的导言(至今未发表)就把“理学博士”头衔授予 J.L.~Verdier的 C.~Chevalley与 R.~Godement,十年之后,有没有份收到这部不同寻常的“学位论文”的“第0状态”(这一回有50页)!我记得有一天曾亲手捧读过一部几百页的严肃奠基工作,它满可以理直气壮地算一篇优秀的博士论文,并且大体对应于我1960年前后向Verdier提议的那项奠基工作——只除了一点:彼时“三角范畴”框架(他用来表达导出范畴的内蕴结构)业已显出不够用。

几乎无须说:我的名字,在这份学位论文的“第0状态”里无处可见。人们纳闷:它到那里去做什么。众所周知,导出范畴是Verdier引入的,为的是让他得以发展拓扑空间的所谓“Poincaré--Verdier对偶”与解析空间的所谓“Serre--Verdier对偶”;只待某位当值无名之辈\fnote{37}{(*)}{关于这位可疑人物的几点情况,可参见“当值无名之辈与好上帝定理”一注(n\textsuperscript{o}~48$'$)。}出头,为他自己发展一个把二者综合起来的理论,并且理所应当地(无名的学生至少也得做到这一步!)被称作“Poincaré--Serre--Verdier对偶”。经过这一切,我剩下的只有随大流,做几处势在必行的调整,以便在概形的平展上同调或凝聚上同调那个再特殊不过的框架里,发展Poincaré--Verdier对偶与Serre--Verdier对偶……

我刚刚才得知(图书馆到底有用!)SGA~4½的内容\fnote{38}{(**)}{关于这一卷,可参见“白板”一注,n\textsuperscript{o}~67。}——在那里,人们还不吝让我以合著者、或不如说Deligne的“合作者”(原文如此)的身份露面(既不费心告知我,更不费心征求我意见)。显然,这是五年后问世的“下葬文集”的先声——那部文集,几天前我有幸拜读(见由这一事件激发的注n\textsuperscript{o}~50、51及以下)。但我无需亲手捧起那部“预葬”文集,连同那篇不肯自报家门的幽灵论文的铁证,便早在去年就已懂得:这份“学位论文”的下一个状态,除了我,不会有别人来写。于是我就着手这项工程,即\textbf{《追寻栈》}\dmark{}\origpage{243},从我那位显赫的前学生十七年前停下来的地方接着做。

\subsection{授职}

\notehead{64}{}

(4月25日)然而昨天,我在系里的办公室里找回了一份副本。其实是前后相隔一年、写于1968年4月(?)与1969年4月的两份报告。我用十七页的篇幅,评阅了他在IHES三年科学活动中进行的十五项工作。其中有关于Ramanujan猜想的工作(原书拼作Ramanuyam),关于模site的紧化的工作,以及Hodge理论的推广。报告里评阅的这些著作(单是我刚点名的几项就足够了)见证了一种惊人的创造性,以完美的从容、仿佛游戏般地铺展。把Weil猜想的证明——那还是紧随这第一次向未知的冲刺之后——放到一边,我觉得,他后来的著作,只给了那次独一无二的起飞一副苍白的图像:一个手段出众的年轻心灵,享受着使其绽放的种种同样出众的条件。然而,不得不相信,这些“出众的条件”中的某种东西,必定喂养了另一股力量——一股与求知的冲动异质的力量;它最终反客为主、取而代之,使最初的冲劲改道并被吞并。而显然,这个“某种东西”与我这个人有关……\fnote{39}{(*)}{(5月26日)至于我身上的某种纵容——它喂养了这个“某种东西”——可参见比本注晚两周的“高升”一注(n\textsuperscript{o}~63$'$)。}

这篇简短的评述报告(我打算作为附录收入本卷)从不止一个角度看都饶有趣味,包括数学的角度(评阅的有些著作,至今仍未发表)。在报告的几处,我预言:Deligne只勾勒了轮廓、处理了关键之点的某些工作,将由未来的学生们去发展。这些学生从未出现,缘由是他后来对常人世界的关系起了变化\fnote{40}{(**)}{在我还常在IHES与他为伴的年月里(尤其在我的研讨班上),Deligne对其他数学家、尤其对来参加研讨班的青年研究者(往往是初学者)的关系,带着亲切。我在那里看到的对他人思想的开放——哪怕它表达得笨拙甚至混乱——同我们数学的两两对谈里一样。他有一种能力:循着他人的思想,进入他人的图像与语言;这能力我一直欠缺,而且(依我看)使他比我远更适合“老师”的角色——善于激发他人身上天职与创造性的绽放。}。在我评阅的那些想法之中,据我所知\dmark{}\origpage{244}唯一曾被别人发展出来的(那人由此俨然成了Deligne的学生),是上同调下降理论——由圣多纳在SGA~4里发展(因此还处在最初冲劲的时期),此后成了上同调武库里最常用的工具之一。有一处既有趣又有代表性的细节:后来成为Deligne文章的四项工作里有三项\fnote{41}{(*)}{四项工作中唯一不受我直接影响的一项,是从Weil猜想推出Ramanujan猜想的工作;它处在一个研究方向(模形式的方向)上,那是我数学修养里最严重的“空洞”之一。其余三项是:Leray谱序列的退化,Hodge--Deligne理论,模多重簇(与Mumford合作)——见“驱逐”一注(n\textsuperscript{o}~63)与子注n\textsuperscript{o}~63$_{1}$。},我在行文之间,以动人的细心,让读者感觉到这些工作与我引入的想法、我提出的问题的关系——仿佛要抢在前面,预先对付作者日后在他的文章里对这些内容的沉默(这几篇文章,在我写报告的时候,没有一篇已发表,而且我相信,甚至还没有成文)。

\subsection{死结}

\notehead{65}{}

(4月26日)同样清楚的是:把一部大面积的“瑜伽”(权的瑜伽,进而Motif的瑜伽)留在自己怀里——关于它,我对除他之外的人也不止一次谈起过,但唯有他亲密地消化了它、把握了它的全部分量——这本身授予他一种额外的“优势”:作为一件无可比拟的发现工具的独家持有者,用以理解代数簇的上同调。然而我不认为,在我还完完全全在场、活跃于数学世界、毫无迹象预示我无限期离去的那个时刻,这一诱惑起过什么决定性作用。它必定是随着、或在我的离去之后才出现;我的离去,是攫取一份遗产的“天赐良机”——那份遗产本是完全理当归他所有!——只要藏起遗产,连同它的来源。

正是在这里,我又一次看见那个深藏的矛盾的死结,在一个极端而格外触目的个案里,远远超出一切个案:我想谈的是无知、不屑、深深埋藏的怀疑——它们围绕着我们自己身上所怀的那股创造之力:那笔独一无二的遗产,比一个人所能传递的一切都更珍贵。正是这种无知、这种对我们身上最珍贵、最稀有者的阴险的异化\dmark{}\origpage{245},使我们得以嫉妒在他者身上知觉到的力量,使我们得以贪求那力量在他人身上的果实与外部印记——那股我们已在自己身上遗忘的力量。一旦这份嫉妒、这份\textbf{取而代之}的欲望扎下根来、找到滋蔓的机会,把本可用于创造性绽放的能量引向这一欲望,我们身上的异化就变得更深,安营扎寨。我们越接近“取代、排挤、炫目”这个觊觎的“目标”,就越远离、越割断我们身上那股精微的力量,也就剪断我们自己创造冲劲的翅膀。在我们顽强地抬高自己的持久努力中,我们久已忘记飞翔,忘记我们生来就是为飞翔的。

在与我的关系里,从我们相识那天起,我就感到我的朋友完全自在,没有任何征兆使我疑心:他会为我的名声或我这个人感到丝毫的震撼或目眩,或者他心中有任何未曾说出的怀疑——无论是关于他在数学领域的天赋,还是关于任何别的。同样,依我看,他也确实从我和我所处的环境——包括我的家庭——得到了一种友好的、亲切的接待,足以使他自在。但那份吸引我、也吸引别人的单纯而似乎毫无疙瘩的自然,肯定不是等到这次相识才出现并绽放的。他这个人散发的、使他如此可亲的印象,是一种和谐的平衡:他对数学的爱好,丝毫不带吞噬人心的女神的形象。在他旁边,我多少显得像个不知悔改的“书呆子”,甚至“粗鲁迟钝的家伙”——我还记得,他含蓄地惊讶于我与身边的大自然、与四季的节律之间缺乏深入的接触;而我穿行其间,几乎视而不见……

然而,那份深藏的“怀疑”——当时我绝无能力知觉它(甚至今天,处在类似的境况里,恐怕也未必)——必定在我的朋友身上,早在我们相识之前就已存在。带着事后的认识,我毫无歧义地看到它的第一个征兆早在1968年;其后这些年里,又有更清晰的征兆\fnote{42}{(*)}{(5月10日)实际上,还有一个“十分清晰”的征兆,早在1966年;见注n\textsuperscript{o}~82的脚注(*)(原书第329页)。}。不过这些都是“间接的”征兆——从我第一手观察到的东西里,没有一个以怀疑、缺乏自信的面目出现;恰恰相反,而且年复一年越来越如此:一种自满,一种\dmark{}\origpage{246}蓄意的轻蔑口吻,乃至鄙夷。但这样的“反面”恰恰揭示其对面——它们成双,反面是正面的影子。

我还经由第三者得知:对某位他从未有机会亲近相识的、声名显赫(而且以难相处著称)的数学家,他曾处于巨大的紧张期待之中,生怕那位大人物不把他看作与自己的伟大相称。这一证词,与我在年轻朋友身上亲眼所见的截然相反,以至于我很难相信它(那是1973年)。带着事后的认识,它却与我另行得知的那些分裂的征兆相互印证,而且全都指向同一个方向。这份分裂,以及我在其中扮演的角色——作为一桩在相识之前想必还弥散无形的冲突的某种凝固剂——在通常情形下,在与一个曾(以这样或那样的意义)是“老师”、或至少是传递者或托付者的人的关系演变里,恐怕会一直隐而不显。这样,我的离去,就成了一个所有人都未曾察觉、而我或许是唯一知情的冲突的\textbf{显影剂}。

而我今天的“归来”,则是第二个显影剂,无疑更加不合时宜。除了它迄今已就我自己的过去、我的现在、以及我曾爱过而至今仍与之相连的人们所教给我的,我全然无法想象它还将向我揭示什么。也不知道它将向那个人揭示什么——最近一个星期以来,他正处在我反思的这一最终阶段的中心;上个月我曾给这一阶段命名(而且我没想到说得这样准……),叫作“\textbf{一段过去的重量}”。

\subsection{两个转折}

\notehead{66}{}

(4月25日)我的朋友皮埃尔对我的关系中那种蓄意的轻蔑与对抗,只限于数学与职业的层面。个人的关系直到今天仍是一种情谊与敬重的关系,不止一次表现为打动了我的细致体贴,分明是真实而无保留的感情的征兆。

在我离开IHES之后那些紧张的年月里,IHES渐渐沉入遗忘,一如这一情节长期不为我所理解的教诲。于是又有十多年,我的朋友在我这里(想当然地)依然是我数学上的特权对话者;或者更确切地说,1970到1981年间,他是我唯一的对话者(除一桩情节外):在我数学活动断断续续的时期,每当需要对话者的时候,我想到的就是他\dmark{}\origpage{247}。

也是把他当作离我最近的数学家,在最初几次机缘里(1975至1978年间),我同样自然而然地向他求助,请他为与我一起工作的学生提供帮助、担保或支持。其中第一次,是1975年Sinh~夫人的学位论文答辩——她在越南,在异常艰难的条件下写成了论文。他是我为论文委员会联络的第一个人。他推却了,话里带着意思:这只能是一篇冒牌的论文,他绝无意为之作保。(不过我到底有办法说动Cartan、Schwartz、德尼与齐斯曼的好心,请他们为这出“把戏”助一臂之力——答辩在关切与温暖的同情的气氛中举行了。)在其后三年里,又经过三四次同类经验,我才终于明白:在我那位显赫而有影响的朋友身上,有一种蓄意的对抗姿态,对着我的“1970年以后”的学生,也对着只带有我的影响之印记的工作(至少“1970年以后”着手的那些)。我不知道,我在好几次这样的场合所觉察到的那种公然的轻蔑,是否也多多少少见于他对另一些数学家的关系——那些他视为远在自己之下的人。他引以为荣的那套无所节制的精英主义的精神本身,会使我作肯定的猜想。无论如何,自1978年起,我不再为任何事情向他开口。这并未妨碍他那份令人气馁的力量,又找到机会有效地施展。

大约也是在那一年,头一批征兆——起初很隐约——显现出一种对我自己的数学活动的轻蔑态度。第一次机缘,是我关于胞腔地图的反思,起于一个关于它们的发现,它曾使我如遭雷击(参见《纲领概要》第3节:“与儿童涂鸦相联的数域”)。这项发现(诚然“平凡”,本来毫无使我的显赫朋友动容乃至感兴趣之处)是另一个数学之\textbf{梦}的起点与第一份材料——那个梦,其规模可与Motif之梦相比,只是到三年之后(1981年1--6月)才随《穿越伽罗瓦理论的长征》开始成形。这些笔记与同期其他的笔记(手稿凡两千页)构成对那片“新大陆”的头一次巡视——一桩关于儿童涂鸦的平凡议论,使我隐约瞥见了它。

\dmark{}\origpage{248}在这项紧张的工作过程中,有两三次我给我的朋友写信,把我的一些\textbf{想法}告诉他,有时也提出一些技术性的问题。当他肯对我的问题开口时,他的评点始终同样清晰、同样中肯,见证着他年轻时便已打动过我的那同一副“手段”。但一份自满已经磨钝了那份曾经使我陶醉的求理解的渴望,也磨钝了那种经由“小”事把握大事的能力,那种倾听大小事物而把握或构想宏伟蓝图的能力。那种能力不属于理智、不属于单纯的“效率”、不属于对一门已成学科的“掌控”或对已知技术的“掌控”;它在理智的层面上,映照一件本质全然不同的事——孩子的那份\textbf{惊叹的天赋}。这份天赋在他身上仿佛已经熄灭,仿佛从来不曾有过。至少就他对我的关系而言是如此——而在此之前,对我的“1970年以后”的学生,已然如此。他已成了一个大人物;他对数学的进路,变成了不折不扣的那种“竞技”态度——一两个月前我才头一回考察过它,而且我自己对它也绝非陌路……

或许我本可以给自己找到一个理由,解释那在共同热情中的契合、那条曾经连接我们的深深纽带的显然缺席。我本会很满足——多半如此——在机会来临时,向朋友的巧思、向他关于数学事物世界的广博知识,提出或多或少技术性的问题,或者仅仅请教信息。但就在那一年(1981年),那份轻蔑的征兆突然变得如此粗暴\fnote{43}{(*)}{(5月28日)关于这第二个转折的新的照亮,另可参见“反常”一注,n\textsuperscript{o}~76。},以至于我对再同他交流数学问题——哪怕偶尔——全然失去了兴趣。($\Rightarrow$注67)

\subsection{白板}

\notehead{67}{}

(4月26日)正是在昨天写前面那些文字时,我把我们关系中的这个新转折,与1982年(也就是几乎恰好在这个苛刻的转折关头)Lecture~Notes那部“卓越文集”的出版联系了起来——它为我的数学下葬揭幕,没有鲜花,也没有花冠!既然我已被判定数学地“死去”,我的朋友仍在这里或那里继续回答一些归根结底已无从必要的数学问题,这算得上某种恩典……

\pmark \origpage{249}试着去倾听事件的意味,我感到,同样不是偶然:轻蔑、数学上的漠然(对着那些——他的“健全的数学本能”本该告诉他——灼热而多汁的东西!)的初次显现,至少就他对我的关系而言,差不多就落在“预葬文集”SGA~4½出版前后的时刻,即五年之前\fnote{44}{(*)}{关于此事,可参见前天“同伙”一注(n\textsuperscript{o}~63$'''$)。(6月5日)本注的反思,在本注及随后三则注(“白板”“格格不入者”“绿灯”“颠倒”)里得到继续;它们使人隐约看到“SGA~4½之役”的意味,及其与母研讨班SGA~5之“拆解”的联系。这一反思又在“我的学生们”的行列里、尤其在“我的学生们(1)--(7)”的续篇里重新拾起;在那里,一场对那个我的上同调学家学生们学艺的研讨班的真正屠杀的图景,渐渐显露。此外,整场战役里摊开着一种漫不经心的轻蔑——我对它在同一时期初次显现的那种“隐约的轻蔑”的知觉,在我朋友对我的关系里,只不过是一个非常苍白的反照。还有一桩联想,是一两个星期前来到的,关于我朋友对我的关系中那个“第一个转折”的时刻——1977年底或1978年间。正是在1978年,我的朋友得到了他那枚实至名归的“奖章”(为了Weil猜想的证明)。这一新头衔(系于一桩“困难尽人皆知的”猜想的证明)被我的朋友以何种方式内化,在“葬礼悼词”(关于我这亡故者的文章)及其对偶文章(关于他自己的)里,触目惊心;那两篇文章,诚然是整整五年之后,才在一个“重大场合”问世。对此可参见“葬礼悼词(1)——或恭维”一注,n\textsuperscript{o}~104。}。单是这一卷出版时的种种情势,就已足以见证一种蓄意的轻蔑姿态——既是隐约的,又是炫耀的。不劳烦征求我意见,甚至不告知我一声,就把我列为Deligne的“合作者”,而且小心提防不给我寄一册——单是这一点,在我看来,就比一篇演说更有说服力。更不必说,Deligne这部书,就其主体而言,本是要把我在十五年多以前发展的那些工作——那时我还没有听说过我这位光彩夺目的朋友的名字!——变得让更广大的读者更易亲近。一份轻蔑,以及继之而来的傲慢,必定由两头喂养:一头是我的缺席——它使我对一切毫无察觉,浑然不觉地“照单全收”;另一头是某种气候——它使这类张冠李戴得以“通行”,似乎引不起最轻微的评注。总之,关于这一卷,以及关于他筹备的那部下葬文集,我没有收到任何人(尤其在我自以为在数学家世界里还拥有的一众朋友中间)的只字回音。

\pmark \origpage{250}在导言里,作者更是开门见山地亮明底色。该卷的目的,是让“非专家”免于“求诸SGA~4与SGA~5那些芜杂的报告”,“剪除多余的细节”,“让使用者得以忘掉SGA~5——人们可以把它看作一系列枝蔓,其中一些还非常有趣”(对这些“枝蔓”倒还算客气!)。“SGA~4½的存在,不久将使人得以照原样出版SGA~5”——一句神秘的话,因为人们纳闷:这一出版(出的是一部他们劝人忘掉的东西)拖拖拉拉已有十二个年头,而且拥有一组自洽完美的结果(它们不曾等Deligne来提炼与证明),凭什么要以SGA~4½的存在为前提。在提出这个问题的同时,我也隐约看到一个简单的回答,以及那个可怜的SGA~5研讨班坎坷命运的也许一种解释(注68)(那是我于1965/66年间纵横发展起来的,比DeligneSGA~4½分卷的出版早了十一年\fnote{45}{(*)}{对这一“奥秘”的澄清,可见“绿灯”一注(n\textsuperscript{o}~68)的一条脚注(4月28日所写)。})。当人们被告知(第2页):SGA~5的原始版本里,“Lefschetz--Verdier公式只是猜想式地确立的”(这对Verdier未免刻薄——他本该会证明自己的定理,而那定理比SGA~5更早\fnote{46}{(**)}{(6月10日)对此的更精确说明,见“屠杀”一注(n\textsuperscript{o}~87)的子注n\textsuperscript{o}~87(原书如此,未印附号)。}),而且“此外,诸局部项在那里也未加计算”。这对非专家读者(本卷首先面向的读者)可能显得是一处令人遗憾的缺漏。略知门道的读者则很清楚:这些“局部项”,到今天也仍然没有“计算”出来;而且,倘若有人问那位光彩夺目、武断专行的作者,他在这句话里(就一般情形而言)所说的“计算”究竟何指,他自己也会窘迫万分\fnote{47}{(***)}{(6月10日)在一般的Lefschetz--Verdier公式里,对一个系数层与其自身的上同调对应,那些“局部项”(对应于不动点集的连通分支)在写出公式的一刻便毫无歧义地得到定义。这些局部项的“计算”问题,只在具体的个案里才取得精确的意义;其中最简单的之一,是Frobenius态射的情形——在那里,它们简单地由这些点上的纤维所诱导的自同态的寻常迹给出。这条公式,在口头研讨班里,已作为一个远更一般的公式的特例得到完整的证明。}(不过,显然没有人想到去问这句不客气的问题)。

一句含混的话——“这个研讨班(? )含有\dmark{}\origpage{251}另一个证明,这一份在Frobenius态射的特殊情形下是完整的”——似乎暗示:SGA~5终究没有给出(谁又能料到呢,在一部“枝蔓的文集”里!)它所宣布的那个主要“结果”——一条迹公式,从而蕴含Weil式的L函数有理性——的完整证明;幸好“这个研讨班”前来救场,亡羊补牢,犹未为晚……

在第4页,我们得知,“Arcata”诸报告的目的,是“给出平展上同调诸基本定理的证明,剥去包裹着它们的、SGA~4里肆虐的non-sense”\fnote{48}{(*)}{在“我那个年代”,英文里那个沿用已久的词“général non-sense”(泛泛之谈,意为:有时烦人却往往必要的通论)并没有贬义,而更多带着一点玩笑与憨直的意味。那个沿用已久的定语“général”在这里被“忘掉”,肯定不是偶然:这样一来就成了“non sensé”,用漂亮的法语说,不多不少正是“无意义”,还让人联想到“嗡嗡的空谈”(bombinage)和“胡扯”。}。他还有这份慈悲,不去细说SGA~4里肆虐的那份令人惋惜的non-sense(比如topos以及诸如此类的恐怖之物——读者可以庆幸自己福星高照:靠了这部光辉文集的天意降临,终于对那可惜的“矿渣”来了一场大扫除……)(注67$'$)(注67$_{1}$)。

我此刻把这一卷的导言与各章导言通读了一遍,把在我看来最明白地亮明底色的那些评语与自白抄录于此,另有两三句(路数是:枝蔓嘛,当然,不过“非常有趣”)在我看来主要是用来“把苦药裹上糖衣”的(而且糖衣果然顺利咽下)。比如,作者倒有这份诚实,开宗明义地说清楚:“若求完整的结果与详细的证明,SGA~4仍然不可或缺”。这一卷,无论其精神与动机多么暧昧,并不属于一场诈骗\fnote{49}{(**)}{(5月26日)不过,请看后天的“颠倒”一注(n\textsuperscript{o}~68$'$):在那里,我回到了这一印象——它被证明是仓促的。在反思的进程里,一场大面积的“SGA~4½--SGA~5”战役渐渐显露,它主要是为了Deligne的“收益”而做成,有我全部“上同调学家”学生的协助或默许。我所相信可以确认的那份“诚实”(依据导言第7行刚引用的那句话),在这里扮演的,是“对照线”的角色,用以障人眼目——是地道的“暂停!”风格。我的朋友从1968年起就在用这种风格(见“罐装的权与十二年的秘密”与“驱逐”两注,n\textsuperscript{o}~49与63)。另可参见“暂停!”与“中国皇帝的新衣”两注,n\textsuperscript{o}~77与77$'$。}式的战役。它的角色,在我看来,更近于一次试探——显然是结论性的:原来根本用不着这么多顾忌!

\dmark{}\origpage{252}有一种\textbf{荒谬性上的步步升级}(显然无人察觉!):从这一卷,到他筹备的那一卷(SGA~4½,以及LN~900)。在两处,人们都看到一个手段出众的人——一个生来就该去发现、去巡游、去探测浩瀚世界的人——执意去“重做”一位先行者的工作:先是我,然后是我的一位旧学生(Saavedra);而与此同时,他对这些先行者的工作——它们都是细心做出、并且深探到底的——拿不出任何本质的新东西。(他所能加进去的东西,总共恐怕二十页、三十页就写得下,依我看。)在第一个个案里,给出的理由还说得过去:让非专家的使用者得以“无泪地”亲近平展上同调\fnote{50}{(*)}{(6月10日)写这条注时,我才刚刚“登陆”,还没有感觉到“SGA~4½之役”的真正意味(及其与SGA~5坎坷命运的联系——对后者,我恰在那时才有了一次突然的预感)。后来我懂得:以那个骗人的名义SGA~4½发表的这堆五花八门的文本(见“颠倒”一注,n\textsuperscript{o}~68$'$),根本不是要充当SGA~4与SGA~5研讨班(它们构成我已刊数学著作的心脏)的一部“无泪”的普及读物,而是一次用它取代我已刊著作的行动:把我的著作呈现为一个多少有些含混的先声,把它自己呈现为关于平展上同调的真正杰作——仿佛这杰作当归功于Deligne。对这样一种冒名顶替的惊人表述(出自一支匿名的手笔),可在“SGA~4½”这次“试探”六年之后,“葬礼悼词(1)——或恭维”一注(n\textsuperscript{o}~104)里看到。},不必倚仗卷帙浩繁的SGA~4与SGA~5。(不过,人们头一回看到这位作者对常人世界如此体贴,竟压过了做数学的乐趣……)第二个个案里,这项工作实际上等于\textbf{照抄}Saavedra跟我做的学位论文!那篇论文本是一个完美的参考;其中一个陈述的证明有误、另一个陈述带有一条多余假设,这肯定不是重写整篇文章的理由。当然,对这样一件怪事,没有给出任何“理由”。

然而,我无需亲手捧起SGA~4½,就感到了这件表面上荒谬的事情的意味:Deligne在十年之后“重做”Saavedra的论文!它无疑同那件只稍逊荒谬的事的意味相同,而后者为它做了铺垫:Deligne(在十二年之后)对Grothendieck已刊著作的某一部分做一个“摘要”(边缘上还带几分居高临下)。恰恰是这一部分,倘若他还继续对代数簇的上同调(他始终无法摆脱它)感兴趣,他在任何情形下都不能装作可以没有。而Saavedra的论文,在所有已刊的、带着我的影响之印记的工作里,恰恰是他无论如何都绕不开的那一部——倘若他想把我发展的\dmark{}\origpage{253}MotifGalois群观念“记在自己名下”,并终于(十五年后!)开发利用这个显然至关重要的观念。经由先执笔SGA~4½,五年后又经由LN~900里那篇江河浩荡的Milne--Deligne文章(别名Saavedra),我的朋友沾沾自喜地给了自己一种虚幻的解脱感——解脱于他分明感到是痛苦负担的东西:不得不时时引证那个他正要取代、否认的人,哪怕仅仅引证某个引证他的人。

要抵达这一亲身的信念——关于这两个“荒谬”举动所共有的意味——我根本无需通读我这位多产朋友那(五十一项)出版物的全单;那份清单,十天前我才(头一回)收到。老实说,我甚至没有想到把我手里的四份抽印本再翻一遍\fnote{51}{(*)}{此外,还不算载于IHES《Publications Mathématiques》的那些工作——所长Nico~Kuiper,将近十五年来,一直好意地把它们寄给我。},为我相信知道的事找佐证。倘若今后我还会查阅我朋友的著作,那将是为了在里面找到别的东西——那些我在别处已经知道得足够多的东西之外的东西。那时我肯定会有学习美好数学的乐趣——那些数学,当年我曾有幸当面聆听、亲口从他那里学习,那乐趣还要更大!

\notehead{67$_{1}$}{}

\begingroup
\renewcommand{\labelenumi}{(\arabic{enumi})}
\begin{enumerate}
\item (6月14日)在SGA~4½里,我又记下了两桩微型的(细节上的)巧取。其一在“SGA~4,SGA~4½,SGA~5的引路线”(请欣赏这串富有暗示的排列!)里:作者写道(第2页),为了在平展上同调里确立“一个类似于凝聚对偶的对偶形式主义……Grothendieck利用了奇点消解与纯粹性猜想”,从而给人这样的印象:这一形式主义,最终是靠他Deligne本人才确立起来的——在(对许多应用已然足够)0维或1维正则概形上有限型概形的情形(见同一段)。他心里非常清楚:六种运算的形式主义(即整体对偶理论)是我不依赖任何“猜想”确立的;他所加的这条限制,只对双对偶定理(或“局部对偶”定理)才有依据——而那条定理,到头来,在SGA~5里(在Illusie的笔下)还改了姓,叫“Deligne定理”!

另一方面,第100页上有一节,题为“Nielsen--Wecken方法”——那正是我在代数几何里引入的方法,用来证明一条Nielsen--Wecken型的公式;这两位作者(在超越语境里)是用一种三角剖分技术证明它的,那技术在\dmark{}\origpage{254}代数语境里毫无用处。Deligne知道这一方法(连同Nielsen与Wecken二位先生的大名——那篇漂亮的德文文章,他根本无需去读!),是从我嘴里、在SGA~5那场“技术枝蔓”的研讨班上听来的——而SGA~4½正是要让人忘掉那个研讨班!在这一节里,既没有对SGA~5的暗示,也没有对我的暗示;关于这一方法的父权,读者可以自行选择:要么归于Nielsen--Wecken(倘若他消息极其不灵),要么归于这位光彩夺目又谦虚过人的作者。

有意思的是:整部书里,Verdier那“Woodshoie”式的证明(照录原书拼写)——关于一条迹公式,包括我所需要(为了Frobenius态射)的那个情形——只字未提。这一证明(显然已被人遗忘,让位给SGA~5里发展的那个更一般的方法)本是把我的L函数的上同调诠释完全立稳的缺失环节。显然,Deligne与Verdier之间有过(想必是心照不宣的)交易——Verdier把Weil猜想迹公式的信誉让给Deligne,作为交换,Verdier保有他在前一年(1976年)据为己有的SGA~5那一部分。(对此可参见“好的参考文献”一注,n\textsuperscript{o}~82.)另一项补偿:导出范畴与三角范畴的“第0状态”在SGA~4½里的发表——我的名字在那里同样缺席。况且四年之后,在Deligne笔下,代数几何里的平展对偶又得了“Verdier对偶”之名——Verdier这笔买卖做得不坏!(见“不义——或归来的意义”一注(n\textsuperscript{o}~75)末尾。)

\end{enumerate}
\endgroup

\subsection{格格不入者}

\notehead{67$'$}{}

(5月27日)\fnote{52}{(*)}{本注出自前一条注“白板”的一条脚注,是对它的补充,恰好一个月后的同一天写成。}所引的那些段落,连同围绕这部名叫SGA~4½的卓越文集之出版的全部情势,见证了我朋友的一种蓄意的嘲弄与轻蔑姿态——对着我的著作的核心部分:由SGA~4与SGA~5这两个血肉相连的研讨班构成的全体。从4月24日(见“同伙”一注,n\textsuperscript{o}~63$'''$)到5月18日(见“遗骸……”“……与身体”两注,n\textsuperscript{o}~88、89)的反思过程中显露出来的这些“情势”里,原初的SGA~5研讨班遭到劫掠——落实为1977年的那场屠杀式出版——绝非最轻的一项。(尤其见“屠杀”一注,n\textsuperscript{o}~87.)

我朋友这份蓄意的嘲弄,若想起以下事实,便取得全部意味:口头的SGA~5研讨班,正是年轻的Deligne\dmark{}\origpage{255}与概形、与上同调技术、尤其与对偶形式主义、与$\chi$-进上同调的初次接触——1965年,他21岁,在IHES登岸,目标十分明确:跟我学“代数几何”。正是在这场口头研讨班里,以及在两年前SGA~4的笔记里,他得以有幸第一时间学习那些直到今天仍主宰其著作的想法与技术\fnote{53}{(*)}{差不多同样的话,其实也适用于我的其他每一位上同调学家学生:Verdier、Illusie、Berthelot、Jouanolou——对此可参见“团结”一注及其后四注(n\textsuperscript{o}~85至89)。}。

“SGA~4½之役--SGA~5”、乃至我朋友皮埃尔对我个人之关系的这一本质的背景,显然在我写前一条注(“白板(1)”,n\textsuperscript{o}~67)时并不在场,在我之前关于葬礼的反思里也不在场。那个“年轻的Deligne”的记忆——他登上SGA~5研讨班,那时他还一切都等着学,而且确实(而且很快)学到了许多——只是在反思的最后阶段,仿佛违背我的本意,才浮上来。从我那个小小的数学“微观宇宙”里出现年轻的Deligne那一年起,我心中便立下一条蓄意的主意:不把他算进我的学生之列(仿佛把他算进去,就会对一位天资如此夺目的人有失谦逊);这一主意使我还一并低估了——或者不如说,直到最近这几星期才完全忽视——一个分明而具体的实在:人们通常用“师-生”这双重的称呼(而这称呼正是我一向拒斥的)所表达的那个实在\fnote{54}{(**)}{(6月14日)这一蓄意的主意,分明见于我四个月前在“耶稣与十二使徒”一注(n\textsuperscript{o}~19)里终于决定如何谈论他的那种方式——仿佛这样一来,我就在对那个乐于与我划清界限的人,违背了一项缄默或谦虚的义务……}。我曾乐于忘掉、忽视:确实有过某种从我到他身上的“传递”——某种对我、对他都具有巨大\textbf{价值}的东西,尽管对其价值所在,他与我的领会想必大不相同。在他与我之间那四年密切的数学接触中,我所传递的,是某种倾注了我自己最好部分的东西,一件由我的力量与我的爱滋养的东西——一件(我相信)我毫无保留地赠出、却不曾掂量、甚至或许不曾真正感到其分量的东西。

可以肯定,我所给予的,喂养了他身上那种求知的热情\dmark{}\origpage{256}——它与我身上跳动的那份热情同调——但也喂养了\textbf{别的东西};那东西我很久之后才感觉到,而且当时不曾把它与那桩已经发生、却乐于被我忽视的“传递”联系起来。换个说法:我所给予的,在另一个对我隐蔽的层面上,\textbf{也}被接受了——不是作为探测一个迷人的、永不枯竭的未知的工具,而是作为(起初)取而代之、(后来)确立支配、确立对他人毫不容情的“优势”的\textbf{工具}。

不去区分我朋友身上归于那个渴望发现的“孩子”的份额、与归于那个渴望取代、渴望支配(乃至碾压)的“老板”的份额,而只从表面看某些观念、技术、工具在一部著作中所占的份额——即便如此,最近六个星期里的一项意外发现仍然触目惊心:我朋友的著作,从我们相识那年起就振翅起飞,而它直到今天,仍在由我传递给他的东西滋养着。十五年前退出数学舞台时,我曾想象:我传给我这位“非学生”的朋友的那“一点点”(尽管我分明看到它在一次惊人的最初冲刺中的作用),将是第一个跳板,一次起飞将把他带得远远越过起点,使他\textbf{远离}我的著作、远离我这个人。而实际发生的却是:我的朋友直到今天,仍系在那起点上,\textbf{系}在那部他既要否认、要么交付嘲弄或遗忘、又要“利用”的著作上。这是对一个与父亲或母亲的冲突性纽带的典型个案:那根纽带,无限期地把那个乐于在自身培养这份冲突、而不肯冲向世界的人,拴在他注定要离开、要超越的人的轨道上……

今天我看清:经由那份蓄意把我的年轻朋友当作“格格不入者”、而不简单当作一个有幸比其他学生更有才具的学生来对待的主意,经由那份在他面前低估或忘记我所传递之物的分量(以及由此放进他年轻双手的\textbf{权力}……)的主意——经由我的这些态度,我在不知不觉中喂养了他身上的自负与冲突;而这两者,都对我隐蔽。与此同时,我入了一场双人游戏——或者不如说,是一场配合默契的双人游戏,我实在说不好是谁“开的头”(假如这问题有意义的话):我以“谦虚”自命,说我的年轻朋友实在太耀眼,不该做任何人的学生,说我所能给他的一点点实在不值一提——而他,则(早在我离去之前)与我这个人、与我的著作划清界限,在我的\dmark{}\origpage{257}纵容目光之下,否认那分明滋养过他的土壤。

只有在写本注的这一刻,我才终于看清这场游戏;对它的散漫知觉,不过是一两个星期前才在场的。我也看清:我身上的这份“谦虚”或“谦卑”,是假谦虚、假谦卑——是一种不单纯:不肯简简单单地把事情看成它们本来的样子。在这场游戏里,有对我的年轻朋友的纵容——播下的种子滋蔓了百倍!——还有更隐微的,对我自己的纵容:为一桩“特权的关系”、一桩非凡之至的关系垒起一个高台\fnote{55}{(*)}{可与5月10日的“高升”一注(n\textsuperscript{o}~63$'$)对读:在那里,我头一次知觉到,在我与我的朋友皮埃尔的那段关系里,有纵容这一配料。这份知觉始终零散而残缺,直到今天,才在本注“格格不入者”所进行的反思中明确起来。}。(正如一切不单纯,说到底,多半是一种对自己的纵容……)

\subsection{绿灯}

\notehead{68}{}

(4月27日)老实说,我从不曾反思过SGA~5研讨班那离奇坎坷的命运背后的意味。它1965/66年的口头进行,并未引起什么特别的困难;而由一批批(常常掉链子的)志愿者所做的成文工作,却拖了\textbf{十一年}\fnote{56}{(**)}{以我为口头报告所作的详细笔记为底本,把整个研讨班写定,对我而言本不过是几个月的工作。}!直到1976年,Illusie才终于把事情接过来,着手写定尚付阙如的部分,并出版全帙。今天——在这个研讨班之后将近二十年——我才头一回明白“这里面有东西要理解”。也许只有我一个人……

我冒出的第一个念头是:在研讨班那些多少积极的听众中间——他们也多少熟悉此前的SGA~1至SGA~4——必定发生过一种\textbf{饱和}现象:对着那滚滚而来、如一场无可抗辩的怒潮般的“Grothendieck学问”\fnote{57}{(***)}{这与我那班“多少有点被震懵了”的学生的印象相合;那一印象,见“教学的失败(2)——或创造与自负”一注(n\textsuperscript{o}~44$'$)所引的信。}。显然,某些执笔者信仰不足,他们想必不太明白这一切要走向哪里,我为什么偏偏要那么固执,整整一年,要把平展上同调的那些\dmark{}\origpage{258}本质形式性质翻来覆去、直到完全驾驭,连同与之相系的那一整套新观念的武库。尤其能说明那种普遍冷淡的是:研讨班的收官报告——提出未决问题与猜想(据我所知从未发表)的那一场——和梳理各语境中Euler--Poincaré型公式与Lefschetz公式的开篇报告,竟都没有留下痕迹。我不记得当时曾知觉到这种冷淡(甚至后来也没有,直到今天\fnote{58}{(*)}{(5月26日)正是在我多少重新“泡进”SGA~5研讨班的语境之后,我才想起一种不安的印象:那是我翻阅刚收到的、已出版的研讨班文本时产生的(那应当是1977年,出版之年)。那份“残缺”的印象(当时只停留在弥散、未成形的形态)主要地——也许甚至完全地(我当时没有花多少工夫细看,尽管那很值得!)——是由于开篇与收官两份报告的缺席,而尤其(我相信)是由于宣告这一缺席时的那份漫不经心,仿佛这是理所当然——何必费那份力气把它们收进去呢!我在某个层面上必定“感到了什么”,但直到这个月(将近七年之后!)我才费心让它浮上来、加以检视——在“屠杀”一注,以及紧随其后的“遗骸……”“……与身体”两注里。}),因为我一头扎在当时的任务里。

SGA~5的命运——它原初拥有与我的任何其他研讨班同样强的\textbf{统一性},却在随后的十一年不成文的岁月里渐渐被\textbf{拆解}(注68$'$)——本可以让我看到:我以那样的执拗所推进的那些大计划,我曾有几年的时间找到人手辅佐,却根本没有成为一项共同的事业,而始终只属于我个人。我的纲领在这里或那里唤起了些应景的合作,却从不曾在我当时任何一个学生身上变成一种观念之力——一种会促使他去做比他跟我做学位论文时更长程、视野更开阔的工作的力量;那篇学位论文在他生命中的主要作用,是让他学会他选定的数学家这门手艺。

依我看,唯一一个从整体上把握了(且不说据为己有)某种总体视野的人——超越就某类问题、或为发展某类工具而进行的特定“合作”的框架——是Deligne。这想必就是为什么,我(无须把事情明言)在他身上看到的,与其说是一个“学生”,不如说是一位\textbf{天造地设的继承人}。“继承人”这个词,比先浮现在我眼前的“接续者”更好地表达我想说的——后者会暗示一部受制于所受遗产的作品。相反,我把这份“遗产”感受为一个单纯的\textbf{馈赠}:我有能力做出它\dmark{}\origpage{259},以助一种个人视野的铺展;那视野还将从许多别的馈赠中汲取养分(事实上,早在我离去之前就已经如此),并且注定将毫不费力地超越先行与滋养它的一切。

回到SGA~5的悲惨命运:昨天掠过我心头的念头是,这一命运与Deligne对我个人、对我的著作之关系的暧昧,恐怕不无联系——尤其考虑到他那强大的数学人格,必定对我的全体学生施加过一种支配力\fnote{59}{(*)}{(4月28日)这一支配力的一个具体而雄辩的征兆是:SGA~5的出版,直到Deligne认为合适、向Illusie示意让他积极张罗的那一刻,才终于完成——也就是说,恰好在他自己需要把它当作其SGA~4½“摘要”的底本、以便取而代之的那个\textbf{时刻}。(对此可参见Illusie执笔的SGA~5导言的结尾。)这一下照亮了、并使前天我在“白板”一注(n\textsuperscript{o}~67)里还称之为“神秘的”那句声明取得全部意味:“SGA~4½的存在,不久将使人得以照原样出版SGA~5。”“照原样”在这里是一记尖刻,我恐怕是唯一(从前天起)感觉到并赏识其分量的人!(相对于原初研讨班,已刊版本代表的那场“拆解”,人们有目共睹。)}。他必定在这份研讨班笔记的坎坷命运中,在其被剥除那构成口头研讨班之统一与冲劲的一切之后,在内心深处得到了满足。然而反思之后,有一点很清楚:这些坎坷的第一性的、本质的原因,\textbf{不在}与会者中任何一个人的性情里。这一原因虽然我尚未清晰地辨明,但毫无疑问,它首先关乎我自己的\textbf{这个人},\textbf{也}关乎那些在65/66年间假装承担研讨班成文工作的人们。它肯定就在他们对我的关系之中,或者也许,在他们对我为他们所化身的某种做数学的方式(或某个纲领,或某种看待事情的视野)的关系之中。SGA~5的命运,如今在我眼前,是一种雄辩而执拗的\textbf{显影剂},显影某种我迄今不曾费心检视的东西——只因为我甚至不曾察觉它;直到此刻,我也只是隐约瞥见\fnote{60}{(**)}{(5月26日)这正是前面倒数第二条脚注里所说的那个“某种东西”;它在过去这几个星期的反思中、尤其从那个时刻(5月12日)起,终于浮出了水面——那天,自它1977年问世以来,我头一回费心看一看,“一场辉煌的研讨班”落到我的上同调学家学生们手里、经他们十一年后做成屠杀版之后,变成了什么。}。也许,这几行文字会促使这场集体厄运的某位当事人,把他对这件事的切身印象告诉我。

\pmark \origpage{260}然而,也许有一课——至少暂时地——我现在就能从SGA~5这一情节里吸取;这一情节先预演、随后例证了我所投身的那个著名“纲领”在我离去之后几乎全线的那种壮观的\textbf{停摆}。与我在六十年代那些兴高采烈的年月里或多或少相信的相反(那时我正为自己终于找到好心人辅佐而喜不自胜!),今天在我看来:经由顽强而缜密的工作使一个宏阔的个人视野取得肉身,本性上不可能是一场\textbf{集体}的冒险或事业。或者不如说:倘若真有“集体事业”,那也不是那种围绕同一个人进行十年、二十年(甚至三十年)的工作。只要这一视野注定要成为人人共享的遗产,它就会仅仅在需要的压力下、在此处或彼处,经由这个或那个人日复一日的工作而取得肉身——那个人对这位先行者,也许只知道名字(充其量如此!);而那位先行者的视野太过宏阔,凭他一己之臂,无法使之成形\fnote{61}{(*)}{(4月28日)也许,“我的一己之臂”本足以实现我在六十年代末所设想的那项宏阔的工作纲领——条件是,在随后的二三十年里,我让自己成为这一纲领的独家仆人。我今天感到庆幸的是,我没有走那条路——它本可能是我的路,而我如今清楚地看到它的陷阱与危险。}

\subsection{颠倒}

\notehead{68$'$}{}

(4月28日)作为这场拆解的一个例子(此外还有许多\fnote{62}{(**)}{(5月28日)我只是到5月12日的反思里,才下决心盘点这场“拆解”——见那则名字更贴切的注,“屠杀”(n\textsuperscript{o}~87)。}),我重新想起了SGA~5一场关键报告的命运:它最终由Deligne本人(我相信他从1965年起就揽下了这件事,为的是十一年后才“履行”自己的承诺……)依据我的口头报告写成,然后不经任何手续,直接收进了SGA~4½!说的是:与正则概形上的代数闭链相联系的上同调类的形式主义——它经由过渡到所论闭链之支撑上的“具支撑”上同调,而顺畅地展开。几乎与平展上同调中的一切构造一样(它们在许多别的语境里同样有用,并且在那里已成惯例),这一构造是我在五十年代末在凝聚上同调的框架里发展的(这里是Hodge上同调与De~Rham上同调;在“抽象”代数几何的框架里,它们头一回得到研究,是在我最早的几场Bourbaki报告之一里)。它如此\dmark{}\origpage{261}自然,以至于显然蕴含与上积的惯常相容性\fnote{63}{(*)}{(5月28日)凝聚框架里,见我的Bourbaki报告n\textsuperscript{o}~49(1957年5月),§40(原书如此)。在5月8日的“好的参考文献”一注(n\textsuperscript{o}~82)里,我发现:这些想法,连同我在同一场SGA~5研讨班里为闭链所配的同调类而发展的那些想法(以及许多别的),已被 J.L.~Verdier据为己有,对SGA~5研讨班的存在、对我的这个人,都不置一词。这场战役发生在1976年,即“SGA~4½之役”前一年(在我看来,它与后者紧密相连),而且是在1965/66年母研讨班SGA~5的全体往届听众与与会者的眼皮底下、在他们全都知道的情况下进行的。}。

写这些文字时,我发觉这场戏法(把这场关键报告偷运进SGA~4½)竟成就了这样一个光辉结果:Deligne——他明明参加了65/66年的SGA~5研讨班\fnote{64}{(**)}{(5月28日)而且,正是在那里,他头一回听说了他在那部海盗文集SGA~4½里陈述得如此光辉的那些事情!对此可参见昨天的“格格不入者”一注(n\textsuperscript{o}~67$'$)。相对于他朋友Verdier前一年的手法、以及他自己在别的场合使用过的手法,我的朋友在这里倒是守在公然劫掠的界限之内:因为他把我列为那场闭链报告的作者(诚然,附带那项光辉的成果:把我列为他的“合作者”),而且尚未装作全然不知道:平展上同调理论、迹公式等等里面,有我的一份。至于在这一方向上的决定性进展,可参见“葬礼悼词(1)——或恭维”一注(n\textsuperscript{o}~104)。}——\textbf{不在}封面我的“合作者”之列(这件事昨天翻阅已出版的Lecture~Notes n\textsuperscript{o}~589时就已打动了我);而我,反倒有资格(在研讨班十一年之后)以“Deligne的合作者”的身份露面。不得不说,这是一次相当绝妙的\textbf{颠倒}!在SGA~4½出版之时——我就这样在不知情中与之“合作”——我停止一切公开的数学活动已经七年;以至于我甚至从不曾过问那个可怜的SGA~5的出版,它对我来说,属于一段已被我抛在身后的过去……

(4月30日)至于SGA~5,它如今显出这样一副面孔:一堆有点五花八门、没头没尾的文本(头尾都在半路丢失了!),只有参照SGA~4½的文本才“站得住”。有一件事引人注目,而我恰在此刻才注意到:SGA~4½这个名字本身,就分明暗示这份文本\textbf{先于SGA~5},而后者只有参照前者才存在\fnote{65}{(***)}{(5月28日)关于这场把SGA~4½“强行塞进”SGA~4与SGA~5这两个不可分割的部分——它们构成我已刊著作的心脏——之间的“暴行”的更深意味,见“遗骸……”一注(n\textsuperscript{o}~88)。}。倘若这份文本的作者\dmark{}\origpage{262}的心思不那么暧昧\fnote{66}{(*)}{(5月28日)“暧昧的心思”这个说法,在这里分明是一个委婉语!},而他出于情感的缘由,执意要把他的“摘要”(“外加若干新结果”)塞进他曾扮演过角色的SGA丛书,那么顺理成章的名字,当然是SGA~5~1/5。

我在这里看到第二场戏法;它使我掂量出:Deligne在SGA~5的命运中所占的份额,比我三天前所想的还要重。这也使我收回前一天表达的感想——SGA~4½并不属于一场诈骗式的战役。显然没有人(从Illusie开始——他的诚意当然无可怀疑\fnote{67}{(**)}{顺便也该趁此机会,感谢Luc~Illusie:他以那样的细心与忘我,把某些陷于困境的报告的成文工作、以及“全包”的出版,进行到底——而且是在肯定谈不上鼓舞人心的条件下;其中,我的全然缺席,肯定不是最轻的一项!(5月26日)在n\textsuperscript{o}~84至89各注、尤其“屠杀”一注所继续的后续反思之光里,我对Illusie的这些感谢,取得了巨大的、未曾预料的喜剧维度——写这些文字时,我远没有预感到!诚然,我是顶着我身上的一种抗拒写下它们的;那抗拒尤其表现为(已在计划中的)致谢在注的“正文”里被“遗忘”,以致我不得不靠一条脚注来“补上”。这份抗拒,无疑出于那种不安:从我头一回亲手捧起这部名叫SGA~5的文集起,我就感到了它(而且我相信,直到最近几个星期,我再也没有捧起过它);那份不安,我在前一条注“信号”(原书此处题名如此)的一条脚注(今日5月26日所写)里谈过。这份疏忽,恰好例证了在禅观中、对当下正发生于自身之事保持警醒的留意的重要性。缺乏这份警醒,这里的反思就停留在禅观之下,停留在一个表浅的层面——而留意到那份抗拒,本会引导我去探它的根,并由此也更近地看一看那场美好的研讨班变成了什么(这件事,我两周之后才做)。})——察觉这场“战役”,这无疑要归功于我已领教过的那份“支配力”,我想,还有我朋友这个人身上的魅力:二者都使他高悬于一切怀疑之上!

\subsection{化圆为方}

\notehead{69}{}

(4月27日)十一二岁的时候,我正被拘留在里欧克罗集中营(芒德附近);我发现了用圆规作图的游戏,尤其为六瓣的玫瑰花饰心醉神迷——把圆周用圆规的开度在圆周上截取六次,分成六等份,恰好落回起点。这一实验的发现使我确信:\dmark{}\origpage{263}圆周的长,恰好等于半径的\textbf{六}倍。后来(我想是在芒德的中学,我最终还是上了学),我在教科书里看到,这个关系据说要复杂得多:人们有 $\ell=2\pi R$,其中 $\pi=3.14\ldots$;我坚信书错了,书的作者们(以及自古代以来想必走在他们前面的所有人!)肯定从来不曾做过这个极简单的图——它明明白白显示,人们干脆就有 $\pi=3$。典型的是,我发觉自己的错误(把弧的长与连接其两端点的弦的长混为一谈)是在我把我的惊讶——惊讶于先行者们的无知——向另一个人倾诉的时候(一位被拘留者,玛丽亚,她志愿给我上过几节数学和法语的私人课),恰在我准备向她演示为什么应当有 $\ell=6R$ 的一刻。

一个孩子对自身理解力所怀有的这份信心——宁可信任自己的能力,也不把在学校学到的、书里读到的东西照单全收——是一件珍贵的东西。然而它不断遭到周围人的挫抑。许多人会在我这里讲述的经验里,看到孩童僭妄的例子:它终究得在习得的知识面前低头——事实终于戳穿了某种可笑。但一如我亲身经历的这一情节,其中丝毫没有挫折感、没有可笑之感,而分明是又一次发现(在我草率地把 $\pi=3$ 这一错误公式当作解释的那次发现之后):一次错误的发现,以及同时,人们应当有 $\pi>3$ 的发现——因为显然,弧的长\textbf{大于}连接其两端点的弦的长。这条不等式,恰好又向着被驳回的 $\pi=3.14\ldots$ 那个公式的方向;于是它显得讲得通了,同时我那时想必隐约察觉:也许有一些并非那么愚蠢的人,曾经钻研过这个问题。此刻我的好奇心已经满足;我不记得当时还想对这个数——它想必重要到人们为它单配了一个字母——的来龙去脉知道得更多\fnote{68}{(*)}{(4月28日)前面的追述,唤起了别的回忆;它们表明,那个大名鼎鼎的 $\pi$,比我起初记得的更让我着迷。在一本书里(也许就是同一本)找到的近似值$344/133$(原书如此)打动了我——它那么漂亮,我简直难以相信它只是近似的!那时我除了分数之外不知道别的数,我为表达 $\pi$ 的那个既约分数的分子与分母会是什么模样而着迷——它们必定是些很不寻常的数!不必说,在这些关于化圆为方的孩童反思里,我没有走多远。}。

\dmark{}\origpage{264}这一经验,无疑是头几课之一:当我自己的理解似乎与公认的知识相抵触时,要有某种审慎——这样的局面可能值得仔细检视。审慎是经验的果实;它贴合并补足(而不改变)对自身认识与发现能力的自发信赖,以及认识到我们本就具有这种能力所带来的笃定。

\subsection{葬礼}

\notehead{70}{}

(4月28日)昨晚重新想到SGA~4½封面上那一幕——我在不知情中,以我显赫前学生的“合作者”的身份露面——这事在我看来难以置信到这种地步,以致我疑心自己的记忆背叛了我,怀疑我其实分明被征求过意见、并且没多想就点了头。但这一假定,与我直到去年还抱持的态度背道而驰到如此程度——即我根本无意再发表数学(更不用说以某人的“合作者”的身份,而且是某位与我的关系当时在我看来已带着深深的暧昧的人)——以致它比它本该“解释”的那件事还要“难以置信”得多;而那件事,归根结底,对我没有任何神秘或不可解释之处!为了求个心安,我还是查了我朋友1976年至今的信件(为数不多,很快查完),当然没有找到任何关于SGA~4½出版的暗示。我还是给当事人本人写了几行,请他对这出我不太欣赏的“闹剧”给我些解释……\fnote{69}{(*)}{(5月26日)我的朋友到底赏了我一份回答,它驱散了最后一丝疑云。他让我以“合作者”的身份露面,确确实实是因为那场SGA~5报告——他把它写成并收进了SGA~4½;而他不认为有必要为这次移交、或为以“合作者”的身份露面征得我的同意,也不认为有必要给我寄一册这本我曾“鼎力合作”的文集,鉴于“我已经七年不做数学了”。(6月5日)我刚收到(迟做总比不做好!)Contou-Carrère的一封信(日期为5月30日),回复我4月14日的一封信——我在信中问他(为求心安):在我的书里,他是否曾见过一份SGA~4½的副本。看来确有这样一份副本,Contou-Carrère留在了自己手里(除非是他买过而忘了?)。另一方面,Deligne的回答似乎确实证实:他不认为有必要寄我一份:“给你寄一份4½的副本,本确实可能是个好主意;我无疑料想你那时不会看出它的趣味”(5月15日的信)。}

三天前,当我在反思里提到三年前我与我朋友皮埃尔的关系中发生的那个转折——我失去了\textbf{兴趣}\dmark{}\origpage{265},不愿再同他交流数学问题(见“两个转折”一注,即注66)——我想起了一份当时强烈在场的印象。要把它安放好,我得先说明:在流逝的那十年里,当我的朋友事实上扮演着我唯一的数学对话者的角色时,我曾指望(一如我指望他扮演那个角色一样想当然):他会成为我向他吐露的数学反思与\textbf{想法}的中继站,把它们转达给可能感兴趣的数学家们。如我在别处解释过的(见第50节“一段过去的重量”),正是拥有这样一位对话者-中继站的感觉,赋予我那些断断续续的数学活动时期以比解馋更深的意味:把它们联系到一项超越我个人的共同冒险。无疑也是这份感觉,使得在那么长的时间里,我从未感到丝毫发表我所发现的东西的欲望,更从未感到丝毫对自己退出数学舞台的懊悔。(这样的懊悔,说到底,从未出现过;我是毫无蓄意地、“重新出现”在那个“舞台”上的,甚至在我自己察觉之前!)

至于我的朋友在多大程度上回应了这一期待,我说不好——很有可能,只要他对我还保持着那份数学上的热忱——它由好奇与亲切的同情共同推动,正是它使他在我与数学家世界的关系里(也在某种意义上,在我与数学本身的关系里)所扮演的那个特殊角色成为可能、并显得自然而然——他就一直扮演着那个角色。一两天前,当我向自己提出前面那个问题时,我收到了(仿佛作为部分的即答!)Larry~Breen的一封信,寄来1974与1975年间几宗通信的抄件,其中有Deligne1974年写给他的一张两行短笺,附着我刚写给Deligne的一封关于Picard叠形式主义的信的抄件,征询Breen对我那封信的意见。他在其中以“大师”一词指涉我;我从中听出一种半开玩笑、半亲切的声调。我不记得还有别的场合,有人把我1970年离去之后向我朋友吐露的事情,经他人之口回传给我。很可能有过,而我忘了;更不必说,即便在我数学活动的那些片段里,需要请教朋友的时候也相对少见;而且直到1977或1978年,我向他吐露的那些反思,分量都有限。所以,\dmark{}\origpage{266}严格说来,在那一刻之前,并没有多少东西可供“中继”\fnote{70}{(*)}{我可以算作例外的,是我关于分层结构的剖视理论的最初反思——七十年代初,我想必向Deligne提起过。他以一种纵容的同情接待了我在这方面的期望,有点像人们对一个什么都不怀疑的大孩子的态度。(他在与我的关系里常常是这种心绪,而且肯定常常有理!)我朋友的怀疑——它由他所熟知的某些野性现象的知识所推动,而那些现象是我所不知道的——当时并没有说服我;毋宁说,他向我指出的那些事实,使我从那时起就疑心:人们为“做拓扑”而通行的“拓扑空间”语境,不足以柔顺地表达某些我觉得本质的拓扑直觉,比如“管状邻域”。在其后十年里,我几乎没有机会回到这些反思,我的那些“疑点”被我淡忘了些许;它们经由我1981年12月--1982年1月的反思重新变得现实(并成为亲身的信念)——那场反思,是被为“Teichmüller塔”建立一套“剖视”理论的需要所激发的。(对此可参照《纲领概要》第5、6节。)(6月5日)作为另一个例外,我可以算上我关于虚拟相对概形与虚拟Motif(在一般基底概形之上)的反思——我记得曾向Deligne提起过。由于这些东西与他已决意埋葬的一部瑜伽(直到1982年的开棺出土)紧密相连,他没有表现出对我向他解释的那些想法有任何共鸣,这并不奇怪——那些想法当然使我心醉神迷;关于它们的几点指示,见注n\textsuperscript{o}~46$_{9}$。}。

事情在1977年起变化:那一年,自六十年代以来头一回,我强烈地“咬住”了一种异常丰饶的实体。那是我关于地图的反思的开端,以及(大约同时)顺藤摸瓜地,关于正多面体的一条新进路的开端(见《纲领概要》第3、4节)。也是从那一刻起,对我很清楚:我刚刚把手指按上去的那些事实,开启了无法预料的视野,其广度与深度,可与Motif观念诞生时所隐约瞥见的(以及后来远不止隐约瞥见的)相比。

奇怪的是,那一回,我竟然还是带着他将成为回声的期待去找我的朋友——尽管“Motif”这个名字周围已沉默了七八年,那沉默本已足够雄辩,足以教会我:我的期待是幻象!这份惊人的不察,恰好例证了我心中那条蓄意的主意(甚至在一两年前发现禅观之后):对我与数学、与数学家们的关系,一概不予留意——它们本该属于一段遥远而彻底翻篇的过去!我头一条向着这个\dmark{}\origpage{267}方向的反思\fnote{71}{(*)}{关于这场反思,见“扫兴的老板——或高压锅”(第43节)。}恰恰发生在1981年——那是我与我朋友关系中的第二个“转折”之年,我已有机会谈到它。但即便在那场延续了几个月的禅观里,与其他数学家的关系也只被轻轻拂过;而与其中那个无疑最亲近的人的关系(至少在我们共同热情的层面上)——据我所记,连拂都没有拂过。然而那本会非常有用!

但带着事后的认识、经由此刻的反思,有一点变得清楚:那一刻所发生的事——它曾那样使我惊讶、使我失望(一次隐约的轻蔑突然现身,而我本以为可以分享一次给了我深刻印象的发现那犹带体温的喜悦)——正是注定要发生的事。恰恰是我要传达的东西的\textbf{分量}——它曾激起我对一种与我的兴趣合调的兴趣的期待——应当在我朋友身上,头一回激发出\textbf{使我气馁}的反射。这一反射理应更加强烈,因为从SGA~4½出版那一刻起,我就已经被“预葬”了。三年后我再次进攻时——那时我的朋友(手持他关于绝对Hodge闭链的漂亮定理)正准备把下葬办得像模像样,随即有次年问世的“值得纪念的文集”\fnote{72}{(**)}{即Lecture~Notes~900那一卷;见“一个梦的回忆——或Motif的诞生”一注(n\textsuperscript{o}~51)。}——同一个反射又起了作用,只是粗暴得多。(这一情节当时断绝了我们数学层面的交流,但并没有“使我气馁”……)

在两种情形里,那份漠然显然是真诚的,一如它在别的情形里、对着我之外的人表达时那样。这已不是我头一回在他身上(或在别人身上)看到:与求知的渴望异质的力量,把那渴望中和掉,并取代数学家的嗅觉。

正是在这两次机缘——1978年,然后1981年——我头一回如电光石火般瞥见了我朋友身上那桩矛盾的“\textbf{代价}”;那桩矛盾,多年以来我就知道,但作为他著作中、他对数学事物理解中的一种阻碍与限制的分量,此前从不曾向我清楚显形。但只是在这一我已持续一个月的禅观——关于自我离去以来阴险地进行着的某场\textbf{下葬}——之中,这一\dmark{}\origpage{268}分量才终于渐渐显于全光之下。

在显见的层面上,这些天与这些星期里我所发现的下葬——几年前我已有预感,却从未想到把特定角色派给任何人——首先是对\textbf{我的数学著作}的下葬,经由它、并且首先,是对\textbf{我的这个人}的下葬。当然,在所有人中最有资格对这场下葬伸出援手(许多人在内心深处都在呼唤它)、并主持这场无名葬礼的,是那位不久前还在众人眼中俨然合法继承人的朋友。倘若他主持了葬礼,他肯定不是独自参加了葬礼!但更深一层:我的朋友在这些漫长的十二年间所悄悄下葬的,不是别人,正是\textbf{他自己};或者说,是他身上那样一种东西——那件不打动任何人、精微而难以捉摸、如一朵花或一枚果实的香气、无价的东西。($\Rightarrow$注71)

\subsection{墓}

\notehead{71}{}

但顺着联想的线索,我已离开本题——本题是要唤起某一记“强印象”;三天来,它的回忆执拗地回到我这里。这一印象,定位于我与朋友关系中那个“转折”的时刻:当我面对一种蓄意鄙夷姿态的征兆(既含蓄又粗暴地显明)——那些征兆使我在数学层面上结束了我们的关系。我当时懂得:时刻已到,延续这样一种关系,我再也无所期待;而“决定”自己做出了,没有分裂,没有懊悔,作为这份迟到的(而且非常不完全的)理解的第一个果实。

我心中没有愤怒,更没有苦涩。(我记得在我们交往期间,从未对我的朋友动过怒,也从未有过苦涩——除了我离开IHES的那桩情节的时刻;而且当时被包括在内的,也不止他一个。)但有一种悲伤:在同一个继续为我所珍的人的关系里,翻过这一页——尽管曾经把我和他系在一起的最强的纽带,已经干枯、已经死去。而作为此后几年里仍在的一根刺,还留着那份未被消化的挫折:我曾带来想要分享的喜悦,带给那个在我看来最亲近、最有资格分享它的人,却撞在一份自满的紧闭之门上。这份挫折,我看,最终是被我此刻继续着的禅观消解的。就在今天,它还在向我显示:发生在我身上的,正是注定要发生的;\dmark{}\origpage{269}而这份挫折的第一责任人,不是别人,正是我自己——是我自作主张,耽溺于某一实在的幻象图景,而不肯使用我健全的能力、用醒着的眼睛去看这实在。

正是在这份悲伤的底色上,也在这份期待的挫折的底色上,出现了那奇怪的印象——它当时不是作为一场反思(当时并没有发生)的果实或终点而来,而是作为一份直接的、无可置疑的直觉。那就是:我能在数学层面上对我朋友说的一切,以及多年来我已对他说过的一切,我是托付给了一座\textbf{墓}——或者已经托付给了一座墓。虽然我从不曾对任何人谈起这份印象,也从未在后来的反思里把它白纸黑字记下,我却清楚地记得:当时在场的,正是那座\textbf{墓}的意象,以及表达它的那个词(用法语),也就是我刚刚写下的那个词。这份“印象”或意象,想必在那一刻,作为某种理解的(不妨这样说的)视觉表达而涌现——那种理解,在某个层面上,想必早已成形、早已在场,是长年累月间必定发生过的一整组知觉的果实,尽管注意力不曾留住它们,记忆不曾登记它们;一些想必十分简单、十分分明的知觉,但我没有“留住”它们,因为它们对于我身上那位常常有权随意拣择的人……显得不合意。无论在当时还是后来,这份断然的意象,都未曾与任何具体的、可触摸的“事件”的记忆相连——一件朝着这意象方向去的事件,本可能在我心中唤起它;对这一骤然浮现的意象的记忆,此后想必也很少掠过我;今天是我头一回稍稍停留在它上面。

倘若当时没有浮现任何记忆或联想,那肯定是因为我连接待它们的最低限度的余裕都没有。奇怪的是,那时(如果我把时刻定位得准\fnote{73}{(*)}{(6月11日)一些旁证向我确认:确实如此。这个“第二个转折”,位于1981年的下半年。})我正投入一场关于我与数学之关系的禅观,而这一情节——它毕竟相当有力地,经由一个现在,向我诉说某段过去——竟没有使我想起中断反思的“线头”,好把方才发生之事的来龙去脉纳入反思;那件事,在我的一生里并非无关紧要。

\dmark{}\origpage{270}此刻(在唤起这一意象、并说出它一出现就与一切记忆或联想脱节……之后)浮现的第一个(老实说,也是唯一一个)联想,是我的Motif之“梦”所遭的命运——在数学诸视野之中最为我所珍的那一个。倘若那段过去,也许仍对我保有一种隐秘的支配,那正是经由这个梦——而这份隐秘的支配(我相信我在写下这些文字的时刻隐约瞥见它)本身,有着超越言词的、梦的力量。倘若在过往投资的遗产里——在一场对数学的满怀热情的投资的遗产里——曾在流逝的十年间出现一种未说出口而深刻的挫折,那正是看到:一片死一般的沉默,围绕着我视为活的东西,而我曾把它们当作活生生、精力充沛、只待冲向天光的东西,托付给我的朋友!我一走,能够而且有使命看护这次绽放的,是他,唯有他;把唯有他(与我一起)才亲密感到的东西,交到人人手中。而我虽然从不曾以这样或那样的字眼向自己说清这件事——甚至从不曾(据我所记)停下哪怕一念的工夫,想一想我留下之物的命运——但随着岁月流逝,在我身上的某个地方,我必定已经懂得:这个始终为我所珍的梦,我是把它托付给了一座“墓”。

于是,随着这一唤起、以及它在我身上激起的第一个联想,我看见一串别的联想循着它的尾流涌来,向我揭示:我分明刚刚触到一个命门——那也许正是头号要紧之处,我的数学家过去的(久被忽视的)重量由此施加。

但此刻,我想,还不是循着这些联想走的时候——反思的这一“最终”阶段已经拖得很长了。关于我的朋友皮埃尔,关于Motif,我想在这场反思里已经说得够多——肯定,对许多人的口味来说,已经太多!我相信,就这些注而言,是时候以一场对这场双重下葬之反思眼下教给我的东西的\textbf{盘点},来收束它们了。

\section{VI.~事物的回归——或全体一致的赞同}

\subsection{一脚踏进旋转木马}

\notehead{72}{}

(4月29日)\dotfill\fnote{74}{(*)}{我认为最好在这里免掉读者整整一页关于禅观的一般考虑——那是一种顾左右而言他的方式,是不肯进入正题的抵抗的征兆。}

\pmark \origpage{271}我觉得,就那桩占据着我的事情而言,应当做的描述与澄清的工作,就其“局部的图景”而言,主体已经完成。(况且显而易见:眼下这些注定要发表的注,只是实际工作的一幅缩影——在这里把汇成这一幅或那一幅“局部图景”的所有要素一一摆明,根本无从谈起。)当然,经由这同一项工作,某种总体图景也必定已经成形——还模糊,它等待着被表述,以便取得形态与生命,并向我说出它要对我说的话。自我昨天的反思以来,我感到它已蓄势待绽,推着我借给它一个声音。

老实说,昨天的反思(我刚刚重读)所主要教给我的,\textbf{不关乎任何别人,只关乎我自己}。我带着某种如释重负,看到反思回到对我自己的反思这一坚实土地——因为一个星期以来,它屡屡给我这样的感觉:牵连别人的这个人,多于牵连我自己。昨天的反思终于向我揭示了一件分明显而易见的事:即我对一段过去——我的“数学家的过去”——的执着之力,以及那部大名鼎鼎的Motif之“梦”在其中所扮演的特殊角色。

话一旦终于说出口,它的显而易见便跃然眼前——最新近、也许最清晰的征兆,便是那阵情绪:两年后才发现某桩“事件”之时所引爆的情绪,即Motif在我的前“学生”兼朋友的引导下,那场“悄悄的”(而且迟到的)“重返”数学兽栏!这阵情绪立刻表现为一场自以为已告终结的反思的重新开工——重新开工立刻落实为五十页回顾反思的洪流!于是(而这一确认,在这场不合时宜的重启中已数次来到我面前)似乎我还没有像一两个月前、在阶段性终点与一阵(绝非幻象的)解脱感的高涨中所以为的那样,“下了旋转木马”\dmark{}\origpage{272}——带着这样的教诲:“我并不比别人高明”,“学生超过老师,我无须惊讶”\fnote{75}{(*)}{见“打住这木马!”一节,n\textsuperscript{o}~41。}。然而这教诲并未阻止我惊讶——只要“学生”在一个我全然未曾预料的方向上超过我,就足够了!不过,倘若这教诲不曾阻止“我惊讶”,它在流逝的反思中却不止一次对我弥足珍贵,使我避开那些惯常的陷阱(或至少,避开这些陷阱中的\textbf{某一些})。

回到这份“支配”之力、这份我对Motif之梦的执着之力:这份执着在本书的许多别处已经显现——既在《收获与播种》里(其中Motif被不止一次谈起,而且用词已然雄辩之至),也在《纲领概要》里(在那里,“客观地”说本无须谈到Motif),或在《主题概要》里(在那里,Motif在一群活力充沛的雏鸡中间,有点像没有孵过的蛋)。在后一部文本里——它上溯到十二年之前,显然是怀着疏离的心境写就的——关于Motif的最后一段,据我看,是唯一能让人忽然感到一股暖意流过的段落……

引人注目的是:自我离去以来的这十四年里,这份执着从不曾向我显现;直到昨天,我才终于隐约瞥见这一显然的事实,直到今天才终于把它表述出来。在将近三年前的那场禅观里(1981年7月至12月),我终于确认了第一个显然的事实:即我身上对数学的那份热情的持存——它在流逝的岁月里,曾以十分雄辩的方式表达过自己。但我对一段过去的执着,据我所记,当时未被察觉,直到今天仍是如此。

然而,我必定是随着“一段过去的重量”那场反思才开始隐约瞥见它的——那场反思,是在对数学家过去的禅观似乎已经收官之时、出于求个心安而到来的(只除了一点:我还不曾知觉到这段过去的\textbf{重量}!)。而且我写它的时候分明感到,自己还停留在事情的表面,并未真正穿透。我随后不得不补加的那些注(先是注46、47)在当时把我引向一个方向——有好一阵,它使我离开我自己:先是把注意力系于一部数学著作(以及在我看来它最“重要”的那些方面),继而是那部著作的坎坷、他人在其中的角色,而不是我自己。

\dmark{}\origpage{273}我刚刚重读了“一段过去的重量”那场反思(第50节)。临近其结尾,我确实开始隐约瞥见:那股“翻转之力”(向着一种不止是偶尔的数学投入)可能是“对一段过去的执着”——对数学家的过去——所使;但更是对“最近十年的过去”、对“1970年以后”的过去,而不是对那些已白纸黑字写下的事物的过去、做成了的事物的过去、1970年以前的过去。几行之后,我却回忆起——但只是“顺带地”——在“我那时眼前的宏阔纲领……只有一小部分得到实现”。写下这几行时,我所想的,主要是“宏阔纲领”中那些直接可实现的部分;而它们的动机之力(!)却远不及“Motif之梦”所代表的那种。(它的论证——但绝不是它的表述——在当时,是“地平线上”的大任务之一……)

显然,我对“Motif之梦”的执着(一如想必一切执着),首先(即便不是完全)具有自我本位的性质。那是不止\textbf{贡献}于一项共同事业的欲望,也是看到这一贡献\textbf{得到承认}的欲望。假定“Motif的宏阔图景”确实已按我自六十年代末所见的那般幅度被勾勒出来,但我在这一视野的绽放中的那一份却被隐瞒——我的不快,想必不会小于(甚至可能大于?)我得知那部“值得纪念的文集”时所感到的(在其中,我分明看到被我提炼并带到光天化日之下的某些观念与想法被重新拾起,却(至少我是这样感到的)被剥去了那曾令我深深心醉的呼吸与炽烈的生命)\fnote{76}{(*)}{(6月14日)这份“不快”,在我看来,首先出于那份无耻之感:一种对某条纽带的蓄意蔑视,装作无视它、把它看作无足轻重。当一个人发现的观念或结果被他人重新发现时——这是常有的事——情形全然不同。}。只要那要看“得到承认”的自我本位欲望——它对着我的数学过去的这些或那些或远或近的事物——尚未烧尽,我自称“下了旋转木马”就还为时过早。数学的“木马”已不再\textbf{包容}我,一如它当年包容我、一如它包容我的某些朋友。但我肯定还在上面留着一只脚;而且我疑心,只要我还掺和做数学,这只脚就会留在上面!

\subsection{事物的回归(或一脚踩进菜盘)}

\notehead{73}{}

(4月30日)我方才重新想起了SGA~5研讨班的命运,以及这一命运与SGA~4½的出版相系的方式。一个曾经混乱的局面——我只是最近这些天、并以\dmark{}\origpage{274}顺带的一瞥检视过它——如今在我眼前显得非常清楚。我刚就此给三天前的反思补加了一条脚注\fnote{77}{(*)}{这条长得犯规的脚注,变成了一则独立的注:“颠倒”(n\textsuperscript{o}~68$'$)。}(见“信号”一注,即注68(原书此处题名如此)),而且我觉得,加上我前天已有的评注(同样以脚注的形式)与此前一天的反思(“白板”,注67),我已表述得足够清楚,无须再为如今已显出足够雄辩的局面再做一幅总结的全景\fnote{78}{(**)}{然而5月9日及其后几天,我还要回到它;见注n\textsuperscript{o}~84--89。}。

行文至此,有一点很重要,必须确认:SGA~5所遭“悲惨命运”的第一责任人、那份“被弃局面”之利用的第一责任人,不是别人,正是我自己。倘若各位“志愿者”(他们揽下了自己并不真想做的成文工作)显然对自己心里没数,那我也同样没有看清自己:我执意不听从一个已然雄辩的局面给出的教训,把事情托付给三心二意的“合作者”,而不肯亲自接手、亲自去做那从那时起就落在我头上的成文工作。毕竟,从口头研讨班结束,到我离开数学世界的那一刻,整整三年过去了(而后者立刻表现为:在其后的十四年里,我对我已刊著作几乎全然的漠然)。诚然,那三年里我完全被别的任务占据:SGA研讨班的延续(SGA~6与SGA~7),EGA的写作,对日复一日涌现的、常常多汁的问题的反思——而在这些问题之中,Motif总体视野的渐渐成熟……被这些任务缠住,我选择了对一场已成过去的研讨班的命运闭上眼睛——那场研讨班(连同前一年的SGA~4),是我就已完成的工作而言所能做出的最深刻的数学贡献,而且多半也是影响最广的一份。

在我一去不返地离开之后,局面只可能更加恶化,使我那些前学生中最显赫的一位得逞:他把那部大名鼎鼎的SGA~4½,插进SGA~4与SGA~5那non-sense与多余细节的矿渣之间,上演了那出绝妙的战役;并让我荣任那份以中心关键文本自居之物的“合作者”——那文本的使命(如他以那种使他生辉的天真所说)是慈悲地让人“忘掉”包裹着它的沉重矿渣……

\dmark{}\origpage{275}总之,我早在离去之前、并经由离去本身所做的种种选择,蕴含着对我已刊著作(或就SGA~5而言,待刊著作)命运的影响,同样也影响我“著作”中仍停留在梦的形态的那一部分——而且是\textbf{未发表}的梦。我不后悔我的选择;当我今天确认这些选择的某些不合我意的后果时,也轮不到我抱怨!倒是轮到我去检视这些后果(而且越不合我意,越要检视!),对事实\fnote{79}{(*)}{(5月28日)这里应读作“我所知道的那些事实”。从后天起,一些全然意外的新事实,将使关于葬礼的反思重新启动,并使我把它相关的注的篇幅增至三倍。}做一幅总体图景(这件事已经做成),并从中汲取它们能给我的教益。这是留给我要做的;今天的反思,至少,也许是在这个方向上的第一步。最近这几天里,我身上已做成了某些比照,我想首先把它们白纸黑字写下来。

在六十年代初的那段时期,推动我一般地在我的学生们身上投资的主动力,是找到“\textbf{人手}”来实现我的本能指认为紧迫而重要的“\textbf{任务}”(至少在我所特有的数学眼光里)的欲望。这份“重要性”肯定不是纯主观的,不是一个“趣味与颜色”的问题;而且常常(我相信),把我的某项任务当作己任的学生,分明感到它“够分量”,也许还感到它在更宏阔的蓝图中的位置。

然而,就这份推动力、这份催促我实现任务的内在动力而言,所关的并不是某种“客观的”重要性——Fermat猜想、Riemann假设或Poincaré猜想的“重要性”,曾使我全然无感,我并不真正“感到”它们。把这些任务与一切其他任务区分开来的,在我与它们的关系里,是它们是\textbf{我的}任务;是我感到了它们,并把它们变成我自己的。我很清楚:感到它们,曾是一项精微而深刻的工作、一项创造性的工作的终点,它曾使人得以把握作为某项任务之对象的关键观念与问题。它们曾经是——而且无疑(在很大程度上)至今仍是我这个人的一部分。把我和它们系在一起(或今天仍系在一起)的那条纽带,在我把某项任务托付给一个学生时,绝没有被剪断——恰恰相反,这条纽带获得了新的生命、新的活力!这条纽带无须说出\dmark{}\origpage{276}(而我在这里把它“说出”,哪怕只对自己,也是头一回)。这条纽带,对那个选择了与我一道工作、并就其自选的某项任务而工作的学生来说,是明摆着的;对我,以及(我确信)对任何别人,也一样。这是构想了一件事的人与这件事之间深深的纽带——它不被后来那些也把这件事变成“他们的”、并向它贡献出自己最好部分的人所削弱,而(依我看)恰恰被加强。这条纽带,我从不曾细心检视过。它在我看来,深深扎根于“自我”的本性之中,具有普遍的性质。人们有时装作无视这条纽带,仿佛自己高踞于这类小家子气之上——甚至可能,我也曾一度入过这种装模作样\fnote{80}{(*)}{可以肯定的是:我遵循着那种“得体的做派”——即对这类事情视而不见,因为它们有悖于时兴的形象!(5月30日)关于这条纽带,见“……与身体”一注,n\textsuperscript{o}~89。}。但最近这些年里(或最近这些天与这些星期里),那寥寥几次:当他人摆出无视这条纽带——他明明知道这条纽带——的态度,而这条纽带把我系于某项已由他人(或由我自己)完成或仅仅指定的任务时,我被触到了痛处。人们可以把这个痛处叫作“虚荣”或“自负”,再给它扣上别的字眼——我并不声称这些字眼在这里不合用;但无论人们给它什么名字,我谈论它毫无羞耻,也不因自己如今的样子而羞耻;而且我知道,我所谈的这件事,是天底下最普遍的事!当然,一个人对其“作品”的这份执着,其力量并不千人一律。在我的生活里——“做”自童年以来就是我几番重大精力投资的恒常焦点——这条纽带曾是强的,今天依然是强的。

于是我可以这样说:推动我与学生之间关系的主动力,是我把他们看作可欢迎的“人手”,用以实现“我的”任务。这一表述可以显得犬儒,而它只是表达一个明摆着的实在——我的学生们想必同我一样分明感到了它。这些任务是“我的”,丝毫不妨碍他们也把它们变成“他们的”——正是他们在自己身上对任务的这种认同,动员了完成任务所需的能量;一如对这个任务的认同,曾在我身上动员了使它诞生并成形的能量,并且继续动员我持续投注于这一主题的能量。这份能量是不可或缺的,否则我甚至无法作为“老师”而“运转”——也就是作为把一门手艺(它同时也是一门艺术)传授下去的长者;而这件事,离开可观的能量的动员,无从进行。在我做教师的全部过去里,我从不曾感到矛盾\dmark{}\origpage{277}:同一个任务,对我这里工作的学生而言深刻地是“他的”,同时又深刻地保持为“我的”。我不认为这一局面带丝毫冲突的性质,也不认为它曾给冲突的意图提供过攀附的机会\fnote{81}{(*)}{倘若,受到某种气候的怂恿,我的一个学生曾想在某项与我合作的工作里藏掖我曾有过的角色,那么事情发生在一个他早已不在学生之位的时候。}。在这种对同一任务同时投资、并认同于它的局面里,学生和我自己(依我看)都各得其所,处于一种全然清晰的工作关系之中;而就其自身而言(依我看依然如此),它不含有任何冲突性的成分。但在纯属个人的层面上,这份关系却停留在表面——这丝毫不妨碍它亲切,甚至友好,有时甚至充满感情。

对我的任务的投资,并\textbf{经由它们}对我的学生们-合作者们的投资,是(我说过)自我本位性质的(想必一如一切投资)。这些任务的实现,对“自我”来说,当然首先是一种自我扩张的手段:经由实现一项宏阔规模的总工程——“我的一己之臂”本无法完成到底的工程。从我数学生涯的某个时刻起,就有这份恒常的暧昧:一种同居、一种紧密的相互渗透,一边是“\textbf{孩子}”——他的求知与发现的渴望,他对隐约瞥见之事、对细察之事的惊叹;另一边是\textbf{自我}、那个“\textbf{老板}”——为自己的作品而欣喜,渴望扩张,渴望经由作品的倍增、或经由一项宏大规模的总营造的执拗而不懈的追求,来增添自己的荣耀!在这份暧昧里,我看到一种分裂,它继续压在我的生活上,给它打上深深的印记——一种也许将伴我终生的分裂。这样的分裂,当然不为我这个人所独有;但也许,在我那被“最好”与“最坏”一并充满的生活里,这份分裂采取了比在别人身上更极端的形式。

于是我可以这样说:对那个贪得无厌、渴望扩张的“自我”而言(它在场上的并非孤身,但确实在!),我的学生们首先是可欢迎的“合作者”,甚至不妨说“工具”——可欢迎的“人手”,用以营造一座会诉说“我的”荣耀的巍峨工程!\fnote{82}{(**)}{写下这句话时,我略有踌躇,字斟句酌,明知有人会把它夺走,当作那个可怕的学阀终于撕下面具的犬儒自白!但我也知道:我拦不住那个一心想淹死一条碍事的鱼的人,他自会从容行事。这不会阻止我继续我的话题:发现并说出明摆着的事情——包括上面写下的那份寒微的真相;它只会让那种从不曾费心朝自己内部看一眼的人感到意外。}

在我看来,这一点已经相当清楚地\dmark{}\origpage{278}在三年前那场关于我与数学之关系的禅观里(进而,关于一般的“做”)显形过,尽管后来我也许把它淡忘了些许。正是这件事,在这最后几天里在我的思想中在场,使我得以同另一件值得注意的事实比照:恰恰是我那时的一位学生(带引号的学生,姑且如此!);而且在所有人中与我最亲近的、也是唯一能不费力气地、就其整体地“感到”我身上那些宏阔蓝图的人——那些蓝图仿佛不容分说地催促我去实现它们——恰恰是他,在我离去之后(而且在内心深处,想必更早……)经年累月地实施了那场\textbf{葬礼}(Enterrement)——规模与著作(Œuvre)本身相当的葬礼(这两个法文词在这里用大写并不过分!),并且最终“主持了葬仪”(Obsèques,再多一个大写,好压得住秤!)。

在这份局面里,打动人的,是那种乌布王式的、巨大的、不可抗拒的\textbf{喜剧性}!最近几天里,我想必隐约感到了这份喜剧;但它只是在刚刚,在我为我自己那场庄严的葬礼放上最后一个大写的瞬间,才向我显出它的真面目——在一阵突然而不可抗拒的大笑里!恰恰是这份\textbf{笑},在反思的这个号称“最终”的阶段里,直到此刻一直缺席;在这里,主导的音色,宁可说是那位“体面先生”的愁容——他的正当期望落了空(甚至可以说被卑劣地欺骗),当愁容没有让位于那些抖得干净利落的嘲讽评点的时候(会说话的人自然会说话!)。我分明感到,经过这场长时间的、忧伤音调的枝蔓(这个词让我想起了什么……),我重新走上了正路。

而此刻,一个顺理成章的名字浮现出来,用以命名这条“注”(人们已不再清楚一条注是关于什么的注,但无所谓),是收束它的时候了。它将是:“\textbf{事物的回归}”。($\Rightarrow$注74)

\subsection{全体一致的赞同}

\notehead{74}{}

我终于感到——呼!——我触到了这一“最终阶段”的终点;它绵延了十二天,其中每一天都(一如从前)以“最后一天”的面目出现。也许,终场的话,就在几分钟前已经说出。我的(象征性的)下葬,是一场\textbf{事物的回归},一场由我亲手播种的收获。(而我那有血有肉\dmark{}\origpage{279}的下葬——倘若我有这份幸福死去,身后留下能为我送葬的活着的男男女女——也将是一次回归:回到我出生时就离开的某物之中……\fnote{83}{(*)}{(5月28日)这一与我自己的死亡的突然联想,带着力量来到了。我一度想把它推开,随后又想删去这突如其来的括号——它仿佛饭里的头发,来得不是地方。我克制住了,出于某种敬意。奇怪的是,第二天我得知:就在我继续反思的那个4月30日晚上,在我居住的市镇,一位朋友的重病缠身的姊妹去世了。我头一回见到德尼丝,就在当天,在她的灵床上。次日5月2日,我与我的朋友、与许多别的活着的男男女女一起,在一个明媚的春日里,把她送入土中……})。我觉得,还可能要补加的,多半只会是\textbf{尾声}方面的了。

那位大名鼎鼎的“最亲爱的学生”,并不是我亲爱的学生们中间唯一一个兴冲冲为我送葬的人;而那些分明伸过手的人们,在他们中间,恐怕也不是唯一出席葬礼而不以为苦的!但说到底,谁是张三谁是李四,于我何干!(关于这件事知道得更多——如果只是为此——不会教给我任何新东西。)我终于懂得了这场“事物的回归”;懂了,我也就领受了它的益处。

然而,我还没有取尽这份益处为我储备的全部实质。我还不清楚地辨明:究竟是我身上的\textbf{哪样东西},使得某些前学生觉得,下葬与葬礼于他们有利可图。难道仅仅是我谈过的那份“贪求”?它(依我看)与别的“老板”们的贪求相去不远;而当年他们与我初试身手时,曾不费力地与它相安(而且无疑甚至不曾察觉,至少在意识的层面上)?那么,是“机缘”(我的离去等等)“造就了贼”,并且\textbf{显影}了他们身上——一如在“众学生之首”身上——的一种\textbf{普遍的倾向}:一旦情势有利,就埋葬自己的“老师”或“父亲”?也许,我的“老师”(或“父亲”……)形象还超出了通常的尺度,而这一情势,曾有助于以完美的齐奏引爆这套“下葬综合征”?此刻我看不清!也许,我将收集到的回声(但愿如此)会使我看得更清,更好地消化此刻摆在我面前桌上的这份意外的食粮。

参加下葬与葬礼的,并非只有学生们,尽管据我所知,没有一个非前学生在其中扮演过突出的角色。显然,我的许多旧朋友从中各得其所。这回,这件事在我看来并不太神秘。

\dmark{}\origpage{280}如我已有机会顺带说过的,我不止一次确认:我不合时宜地退出数学舞台,在我昔日的朋友们身上造成了深深的困扰。那是这样一类东西所激起的困扰:人们在其中隐约感到一种对深刻的重新审视、对革新的\textbf{挑衅}。在这个案子里,自然,数学家中间的这份困扰在我的朋友们身上最烈——在那些认识我、并能感到我曾对那些至今仍是他们的价值投入过多大力量的人中间;更不必说,这些朋友中的每一位,自己也曾在这些价值中、并在这些价值所提供的丰厚“回报”中,做出并且继续做出一份力量相当的投资。早在“生存与生活”(Survivre et Vivre)时期一开始,我就已有充分的机会在别的科学工作者中间观察到这样的困扰。但这从未妨碍:每一次,当我在某位昔日朋友的身上——同情的纽带依然把我和他系在一起——确认出疏远的、有时是敌意的毫不含糊的征兆时,我都感到意外。使我的“弃绝”对某些人变得格外难以容忍的,恰恰是:我本该是他们中间的“最好者”之一,是他们最不可能疑心会给他们来这一手的人!(而且,我确实好几次以为感到,数学界的某位旧朋友身上有一丝\textbf{怨恨}的音调。)于是,很自然,他们在一场时尚里得到了满足:那场时尚宣告,所有这些“Grothendieck学问”,说到底不过是一大摞纸换不来多少东西,等等等等。一个人,无论多么显赫,不足以成一场时尚——要发起的时尚,还得回应许多别人心里的期待、暗藏的欲望,然后才成为共识、才取得支配地位\fnote{84}{(*)}{(5月28日)同一方向上,可参见5月14日的“掘墓人——或全体教众”一注,n\textsuperscript{o}~97。}。

在我离去之后的这十四年里,我也许一直倾向于低估这一离去在“上流社会”里造成的困扰——而对我而言,1970年6月的这次离去进行得如此自然,以至连一项“决定”都无须做出:新的任务一夜之间接了旧任务的班;旧任务忽然退后,仿佛被一段遥远的过去吸收殆尽!(诚然,在蒙彼利埃大学的同事们中间,我也不曾面对这样的困扰;他们构成一个与我离开的那个环境全然不同的世界。)也许,我同样低估了这份困扰在我的“1970年以前”的前学生们中间所能起的作用——他们中的许多属于同一个环境,并在其数学投资中“倾其所有”。很有可能\dmark{}\origpage{281},这份困扰在他们身上起的作用,不亚于在我以为还拥有于同一环境中的别的朋友们身上。无论如何,(我的某位旧朋友或学生与我之间的)每一种情形,都是独一无二的、不同于所有其他情形的个案;我能做的一般推测,只有非常有限而临时的分量。

回到更坚实的个案的地面,有一件事打动了我:我得以确认曾积极参与亲爱的老师下葬的那两位前学生,恰恰也正是最初以轻蔑的姿态、以使人气馁的意愿引起我注意的那两位——对着更年轻的数学家们,他们是“1970年以后的学生”,或者我的观念与数学进路的影响在他们身上清晰可见。这一巧合当然毫不奇怪(这当然不妨碍事件每一次都使我惊讶!)。另一桩有趣的巧合是:两人都是个人关系最友好、甚至最充满感情的那一类(对其中一位,这份关系以同样的音调一直延续到今天)。这与那条一般的确认相合:恰恰是最亲近的关系,最有力吸引并固定冲突的力量。

还有一桩巧合打动了我。在我将近二十五年来有过的学生中间,有两位,在我看来,既以其出众的“手段”、又以其与这手段相称的数学投资,区别于所有其他学生。(一项与我在我生命的二十五年里所做的那种力量相当的投资。)而且,对两位,我都曾有所顾忌,不肯把他们算进我的学生之列——尽管确实,他们都在与我的接触中学到了对他们有用的东西\fnote{85}{(*)}{(5月28日)这是一个委婉语,正如我后来身不由己地确认的那样!对此可参见昨天的“格格不入者”一注,n\textsuperscript{o}~67$'$。}。事情的本性是:两人都各自发现了他们自己的任务,无须我从我曾有(或仍有)的储备中向他们提出任务;两人的学位论文工作,都是在与我这个人无涉的情况下完成的\fnote{86}{(**)}{(5月28日)这并不完全准确。两人都在其工作中,本质上使用了我锻造、并且是他们在与我的接触中学会使用的工具。在这一角色之外,构成其学位论文的工作里的Hodge--Deligne理论(“Hodge理论II”,《Publications Mathématiques》n\textsuperscript{o}~40, 1972, 第5--57页),直接出自他得自我的Motif瑜伽——“混合Hodge结构”是对(在Motif的眼光里同样“明摆着的”)问题——“把复数域上不必半单的Motif观念‘翻译’成(适当意义上的)Hodge‘结构’”的“明摆着的”回答。在一次出色完成的“翻译练习”之外,这项工作中当然有“与我这个人无关”的原创而深刻的想法。但同样清楚的是:倘若没有我引入数学、并且他们在与我的接触中率先获得的那些想法与工具,那套Hodge--Deligne理论今天就无从存在(而且恐怕Deligne或我的另一个学生的几乎全部著作,也无从存在)。}。如此多的共同点\dmark{}\origpage{282}!作为相异点,我要说:两人中较年轻的(如果我没弄错)今天“荣誉等身”(其详尽清单,我替读者、也替当事人那份众所周知的谦虚省去了),而且是影响力最大的数学家之一,也就是说,也是最有权势者之一;另一位,眼下是代理助理,暂代一个职位,而该职位的正式任职者明年就要回来。还有别的相异点,它们在某种程度上解释了这一命运的悬殊——一如也有别的相似点,在这里不值得一一铺陈。倘还要举一条,那就是:在我有过的全部学生中,个人关系最亲近、最友好的,恰恰也是与这两人;而共同的热情,当时一下子就在他们每一位与我之间建立了强的纽带。我此刻想说的这份\textbf{巧合}是:据我所知,他们也是仅有的两位学生(带引号,自不待言!),对着那个“上流社会”,曾尽其可能地淡化和抹去这条通向我的、如此简单而分明的纽带。

这是一份真正触目的巧合;写下这些文字的时刻,它的意味仍从我这里滑脱。对两人,我都可以援引各自的时势缘由,彼此不同。而且很有可能,甚至可以说大概:在两人身上,在某个大概已不再是充分自觉的意图的层面上,这样一种缘由(一方的自负,另一方的谨慎)起过作用。但我怀疑,这个现成的解释,无论在这个个案还是那个个案里,能提供对事情的理解。肯定,在更深处,别的力量——那些真正的力量——曾在自负或怯懦的熟悉表象背后起过作用。肯定,表达这些力量的行为,有某种重要的东西要向他们每一位诉说。但肯定同样:在两个如此不同的人身上出现同样的行为——仿佛他们串通好了(鉴于二人命运的悬殊,这当然不可思议)——也有某种重要的东西要向我诉说,而且不关乎别人,恰恰关乎我自己。难道又恰恰是那没完没了的“\textbf{对父亲的拒斥}”的再现?然而,这种拒斥\dmark{}\origpage{283}本有许多可以选择的表达途径!抑或,正因为无意识那如此可靠的本能——它(当涉及“触碰”时)总能不偏不倚地直取最敏感、最不设防之处——竟使得两人不约而同地落在了\textbf{同一个}地方?实际上,我倾向于这样想。但这是一件推论出来的事,而不是一件看见的事;而由于缺乏那双有幸看得清楚而深的眼睛,我感到自己有点像个盲人,在黑暗中尽力摸索,尽力用他的手、他的耳朵、他的皮肤去“看”——可它们生来不是用来看的。

然而,为了不以这条\textbf{困惑}的音调收束(它有损于我的名声),而以一条令一位善意的、或许会有的读者欢欣的音调收束,我将只说出方才浮现的那个一锤定音的名字;在我看来,它分明表达了这场(关于一场下葬的反思之)\textbf{尾声}的各条考虑所共有的内容,即:

\begin{center}\textbf{全体一致的赞同!}\end{center}

\chapter{C)~上流社会}

\section{VII.~研讨会——或Mebkhout的层与反常性}

\subsection{不义——或归来的意义}

\notehead{75}{}

\dmark{}\origpage{285}(5月2日)看来我学无止境!我刚得知两份文本,它们(至少对我而言)以料想不到的光,照亮了前文已谈过的那场“藏掖”(对Mebkhout著作的藏掖)(“当值无名之辈与好上帝定理”一注,即注48$'$)。文本所揭示的,正是那两位显赫的同行兼前学生所扮演的角色——我曾确认他们心存轻蔑的漠然,却不曾怀疑其职业诚意。两份文本都属于\textbf{吕米尼研讨会}(1981年7月6日至11日)的会议录;研讨会题为\textbf{《奇异空间上的分析与拓扑》},刊于Astérisque~n\textsuperscript{o}~100(1982)。

第一份文本是研讨会的导言,署名\textbf{B.~Teissier}与\textbf{J.L.~Verdier}(即那位曾充当 Z.~Mebkhout正式的论文导师的同一位)。这份一页半的文本,开篇先解释某条“所谓Riemann--Hilbert对应”,它分明注定要在研讨会中扮演头等重要的角色(而它不是别的,正是“好上帝定理”,又名Mebkhout定理)。在这一对应里(这正是它的魅力与深度所在,也是必须引入导出范畴的原因):一个正则完整\textbf{模}(即集中在零次的一个正则完整复形),对应着一个由$\mathbb{C}$-向量空间的层构成的可构造复形;据称,它可由纯拓扑的性质来刻画,而这些性质,对一个任意域上定义的、不必光滑的簇上的平展层的可构造复形,同样保有意义。据解释,这就是研讨会“主旋律”——“\textbf{反常性、相交复形、纯粹性}”——的出发点:所谓“\textbf{反常}”的(\dmark{}\origpage{286}层)复形\fnote{1}{(*)}{(5月4日)关于这一奇怪的应用,见注n\textsuperscript{o}~76“反常性”。}不是别的,正是那些在“实质上”对应于(“Mebkhout式地”)最简单的正则完整微分算子复形——后者借一个单独的$\mathcal{D}$-模来表达——的层复形。

第二份文本,是\textbf{A.A.~Beilinson、J.~Bernstein与P.~Deligne}关于反常层的长文的一部分\fnote{2}{(**)}{(5月4日)我此后收到了全文;它确认了我手头这一部分已经向我表明的东西。};导言里称之为研讨会的中心工作。如同目录与我手头的其他页码所见证的,这篇文章宣告了导出范畴与三角范畴以骤然的强力重返公共广场——追随着Mebkhout那些鲜为人知的著作,与那道著名的“所谓Riemann--Hilbert”定理。

难以置信,却千真万确:两份文本里,Z.~Mebkhout的名字都缺席,正如它在参考文献里也缺席。我要说明:不仅 J.L.~Verdier对Mebkhout的工作了如指掌(这是理所当然的!),Deligne也同样了如指掌;而且,对这样一个对数学时事如此灵通的人——况且这还是与他切身相关的主题——甚至难以设想会是别样\fnote{3}{(***)}{我尤其提请注意:Mebkhout的著作与其“好上帝定理”,相对于Deligne(1969年的)较早的工作,是一个决定性的进展——而Deligne没有发表那些工作。对此可参见已引用的注n\textsuperscript{o}~48$'$。}。

B.~Teissier\fnote{4}{(****)}{(6月12日)B.~Teissier久已关注Mebkhout的工作,并因此是极少数对他持鼓励态度的人之一。因此,他对这桩巧取了如指掌,而且是明知故犯地给予了配合。他曾向Mebkhout辩白说,反正他“什么也改变不了”。}以及吕米尼研讨会的其他与会者情况如何——尤其是与Deligne共同署名前引文章的那两位\fnotevariant{5}{(****)}{(*****)}{(5月28日)我此后得知:A.A.~Beilinson与 J.~Bernstein是通过 P.~Deligne(1980年10月)与Mebkhout本人(1980年11月,在莫斯科的一场报告里,内容极为详尽)得知Mebkhout的结果的。早在1981年6月(原书此处如此)的吕米尼研讨会之前,这两位作者就在他们关于著名的Kazhdan--Lusztig猜想的证明里,本质上使用了好上帝定理。可与“一种不义与无力的感觉”一注(注n\textsuperscript{o}~44$''$)里所引的Zoghman~Mebkhout的信对读。(6月3日)关于研讨会全体与会者之团结的更多细节,见下一条注“研讨会”,n\textsuperscript{o}~75$'$。}——我不知道。似乎与会者中没有一个人曾如此好奇,想去了解把他们动员起来的那些观念与关键定理的父系渊源。

我料想,事情不言自明——有点(很有点)像次年将为Motif重返同一“公共广场”揭幕的Lecture~Notes的LN~900那样\fnote{6}{(******)}{对此可参见注n\textsuperscript{o}~51、52、59。}:父权属于发起并主持研讨会的那些光彩夺目的数学家中最光彩的那一位。对所有人来说,可以肯定的至少是:那绝不是Riemann,也不是Hilbert——否则,那场光辉的研讨会\dmark{}\origpage{287}就该在1900年召开,而不是1981年,即在无名的学生于Jean-Louis~Verdier门下答辩论文的两年之后。

我在这里得以确认的这类战役,如今也许已是家常便饭\fnote{7}{(*)}{我还想到另外两场同方向的“战役”,它们落实为LN~900(参见前一条脚注)的出版,以及五年前的SGA~4½(对此可参见注n\textsuperscript{o}~67、67$'$、68、68$'$)。(5月9日)关于与前面几场紧密相连的第三场这样的战役,见“好的参考文献”一注(n\textsuperscript{o}~82)——这次是 J.L.~Verdier手笔的另一篇“值得纪念的文章”。},并且完全得到认可,只要实践它的是占据高位的数学家,而付出代价的那位只是个无名之辈(尽管人们好心地邀他来,让他高兴)。这些实践者中的一位,倘若凭其手段也凭其作品,俨然一位大数学家(这使他立刻高悬于一切怀疑之上),也丝毫不改变事情的本性。肯定,我太守旧了——在我那个年代,这类战役叫作\textbf{巧取};而这一桩,在我看来,是容忍它的那一代数学家的一场\textbf{耻辱}。

天才的光辉并不减轻这样的耻辱;它给它添上一个新的维度——在我们的科学史上或许是独一无二的维度\fnote{8}{(**)}{在数学之外的任何一门科学或艺术的历史里,我也从未听说过类似的事情。}。它\dmark{}\origpage{288}可以使人隐约看到:在这一举动的荒谬与无端的表象背后(它出自一个命运已对他奢侈到无边、却仍然耽于巧取的人……),或许有别种力量在起作用,而不只是出风头的欲望,或羞辱、或使那个自觉无援无声的人绝望的无谓欲望。

既然我分明已深陷“风俗画”之中,我顺带指出(几乎当作理所当然):我的名字,在前引文本里同样缺席。然而,我能欣然确认:前引文章里没有一页(在我手头的那些页码里\fnote{9}{(*)}{(5月4日)还有别的页码,我此后也知道了。})不是深深植根于我的著作、带着它的印记的——直到我引入的那些记号,直到那些随处可见的观念所用的名字——那些名字,是我在它们尚未得名时就认识它们的时候,给它们取的。当然,合乎规矩的调整是有的:比如我在五十年代提炼的双对偶定理\fnote{10}{(**)}{平展对偶理论同样如此:在其慷慨的朋友Deligne的笔下,它也改姓“Verdier对偶”!},就为这场合改名为“Verdier对偶”——还是那个Verdier,没有搞错……\fnote{11}{(***)}{(5月5日)可与注n\textsuperscript{o}~48$'$、63$''$对读。在这场绵延将近十五年的漫长的葬礼中,也在这场刚刚于上个月、由主要的“被预葬者”所发现的葬礼中, J.L.~Verdier分明始终与他显赫的朋友不可分割:在这可悲的场合,他不计其数地为他撒下那些必不可少的万千花束。}

不过,我的名字不可能连间接出现都没有:经由对几部至今无可替代的文本(尽管有SGA~4½——它还不太胜任其使命)的偶尔引证,即EGA与SGA。(在“SGA~=~Bois~Marie代数几何研讨班”这个缩写的解释里,我的名字当然不出现;但在EGA里,人诚实不诚实总藏不住:全称是给出的,作者名单里包括我……)另一个打动我的细节,它见证了下葬综合征的强迫之力——在一位分明毫无“强迫症倾向”的人身上:我看到的仅有的两处对SGA的引证,每次都恪尽职守地特别写明“SGA~4里M.~Artin的定理……”,生怕哪位多心的读者会想:该定理说不定出自那位被小心隐去的人——尽管天知道,那场报告分明是做过了的,而且出自一位叫得出名字的作者!(注77)

\dmark{}\origpage{289}这一切,想必人们会认为,在今天的“上流社会”里,本是兵家常事。这场小规模交战虽然不能让我高兴(它也不是为此而设的……),却并不真正损害那位被预葬者:他那象征性的遗骸,就这样被交付给这场抢夺集市的种种偶然——而我不过两个星期前,才惊叹地发现这场集市。它并不以无力之中遭受的“不义”之感啃噬我的生活。它没有折断把我推向数学事物、推向周遭世界之事物的喜悦与冲劲;它没有在我身上烧掉这些事物的精微的美。我可以自认为幸运——而我也\textbf{确实}幸运……

我也为我那意味曾向我滑脱的意外的“归来”而喜悦。倘若这归来只需教给我这些天里已经学到的东西,它就不曾虚度——它已经使我满盈。($\Rightarrow$注76)

\subsection{研讨会}

\notehead{75$'$}{}

(6月3日)关于研讨会的其他与会者,我得知了一些细节,它们驱散了一切疑问。尽管官方日程里没有安排Mebkhout的任何报告,Verdier却不得不当场、在最后一刻请他做一个报告,以补某场官方报告的缺漏(那场报告曾托付给布日林斯基——他对D-模理论不甚了解)。Mebkhout由此得以陈述其观念与结果,尤其是好上帝定理,其陈述方式使人们对这条定理——以及与之相随的那套哲学——的父系渊源不留任何疑问;而正是它们,使得代数簇上同调壮观地重新启动,并落实为这场研讨会。于是,\textbf{研讨会的全体与会者,经由这场报告,都被告知了这一父系渊源}。我料想,所有人无一例外,此后也都读到了研讨会的会议录,尤其是导言与前引的Beilinson--Bernstein-Deligne文章。显然,没有一个人觉得有什么不对——或者即便觉得,也没有流露半点。Zoghman~Mebkhout在这方面的回音,一无所获。于是,研讨会的全体与会者,都满可以被视为这场在研讨会期间做成的欺瞒的共谋。

这场集体的欺瞒,早在研讨会当时就已分明,因为当Deligne在其关于所谓“反常”层的口头报告里不提Mebkhout的名字时,没有任何人觉得有什么不对。报告人只限于陈述好上帝定理,并说\dmark{}\origpage{290}他不会在报告中证明它。另一方面,他倒也(带着他那惯常的谦虚)把话说得很足:猜出他所称“反常”的层那些卓越的、事先难以预测的性质,“没有任何功劳可言”——那些性质,由他方才谈到的“Riemann--Hilbert对应”明摆着提示出来\fnote{12}{(*)}{可与所引文章的第10、11页对读。(6月7日)关于藏掖之技艺的细节,见下一条注“变戏法的人”,n\textsuperscript{o}~75$''$。}。人人都觉得正常:他避免点名那位有幸发现这条天赐对应的“有功之人”,并摆出一副父权在己的模样——而恰恰是他们刚被告知、或将在随后几天里被告知:事情绝非如此。人们想必认定:一个研讨会上的无名龙套,竟是一条如此卓越定理的作者,这必是某种不可容忍的错配;于是人人尽了一份力,拨正船头,确立一个共识,把父权归于那位分明是天造地设的人——那位\textbf{本应}是作者的人\fnote{13}{(**)}{(6月5日)况且一切都环环相扣!“学生”行列(继“研讨会”行列之后)中继续进行的反思,以及某种声调(尤其还见于最近与Deligne的一次简短的通信往来——见“葬礼”一注(n\textsuperscript{o}~70)的第一条脚注),向我表明:对Deligne、对我的其他上同调学家学生们来说,久已分明:平展上同调的发现与其掌控,那位本该的作者,同样是Deligne;而在某个层面上(即支配行为与态度的层面),他们深深相信:\textbf{归根结底}那是他,而我在旁边,不过扮演一个毛糙而笨拙的助手的角色——与其说有用,不如说妨碍一套理论(通向“Deligne定理——原Weil猜想”)的和谐展开,妨碍一种令所有相关者满意的角色的分配……}。

有一处很有代表性的细节:\textbf{Mebkhout的报告没有出现在研讨会的会议录里}。Verdier曾请Mebkhout不要成文他的报告,理由是:研讨会旨在陈述新结果,而Mebkhout的结果早已发表两年有余。

当人们不肯把自己关在技术话语里,而去看这场光辉研讨会里真正发生了什么——看推动各方的力量与胃口——人们会以为在看一部电影:某座遥远\dmark{}\origpage{291}的巨型都市的暗巢里,黑手党的天下。然而这是一幅地道我们自家的风俗画,演员们是法国与国际科学最高贵的翘楚。那位对一切发号施令的大主厨,不是别人,正是不久前还之于我扮演着谦和微笑的精神之子——或至少合法继承人(同样谦和,同样微笑)的那一位。至于那可差遣、可宰割的“软柿子”——在满是寸步不让的“硬汉”的世界里——由于一个我尚未完全参透其意味的奇怪“偶然”,也与我这个人紧密相连。他是我的“学生”,一如大主厨(而且同他一样,是带引号的“学生”……)——那是在我已被宣告死亡并下葬多年之后,进到我的门下来学艺的……

\subsection{变戏法的人}

\notehead{75$''$}{}

(6月7日)人们尽可以在那篇“值得纪念的文章”(前两条注里谈到的)里,欣赏那套漫不经心的藏掖的炉火纯青。作为全部工作之本质动机的那个范畴等价,头一回被引进,是在导言第4页(第10页,第9--15行)的一句话的转角处,不给它任何名字,随即转入“反常”层观念的一连串推论(第10、11页)。此后直到第16页末,它再未被提起;在那里我们读到\fnote{14}{(*)}{下面的引文里,着重号是我加的。}:

\begin{quote}
“应当指出,在以下几点上——\textbf{它们本应在这份笔记中占有一席之地}——我们未能完成任务。

——反常层与完整模的关系。如本导言所述,它曾起过重要的启发作用。关键陈述是4.1.9(\textbf{此处未证明})……”
\end{quote}

(随即转入别的“本应占有一席之地”的要点。)

我赶紧去看:作者们没肯劳神纳入其工作(或至少没肯纳入其证明)的这个“关键陈述”究竟是什么。找n\textsuperscript{o}~4.1.9……我碰到一条“评注4.1.9”——想必不是它;我找的是一个“关键陈述”,一个正式表述的定理或附注,并带有一个出处,表明作者们\textbf{在}哪里证明过或将证明它——既然他们不\textbf{在}这里证……可我怎么找,也没有“定理4.1.9”的踪迹——只有一个段落编着4.1.9的号。我姑且把这条“评注”读下去(不抱什么信念\dmark{}\origpage{292}——想必是编号出了错……):里面说“4.1.1在复上同调中的类似是真的……”糟糕,难道我还得回到4.1.1去看这说的是什么?我跳过去,接着往下读——瞧,我本已不再指望了:十一行之后,有一句话以“人们知道……”开头,以“诱导范畴……与反常层的范畴之间的一个等价”收尾。

呼——原来就是它!可我再怎么往下找,也找不到只字片语,来澄清这句谜语般的“人们知道……”。那位本不曾“知道”的读者,只觉得自己全然配不上这局面。对他来说,反正有一点是清楚的(且不论他配不配):这个“本应在那份笔记中占有一席之地”的结果——人们在此借一条技术评注的转角处“提醒”读者应当知道的东西——分明归于那“笔记”的作者们,或其中一位;也许是其中最显赫、并执笔了这篇文章的那位(有一种“家传的笔法”是不会错的……),也就是做了那场口头报告、而其人人皆知的谦虚当然不许他说“是我!”的那位——但人人都懂了,无须他说……

这立刻让我想起最近几星期反思里的几段记忆。头一段,是Deligne1968年的第一篇文章:在4月22日的“驱逐”一注(n\textsuperscript{o}~63)里——那是在LN~900谜底揭穿之后三天——我(十六年后)终于费心把它细看了一遍。我在这里又见到同一套笔法,想必经过十三年的“磨合”而略有变体。在那篇1968年的文章里——其主要灵感来自我——他在文章结尾处顺带、谜语式地提了我的名字,权作“手续完备”。在这里,他连这份小心都不再费——经验久已教他:完全用不着了!相反,在他年轻时那篇文章里,既然他感到非提我不可,他就以完全藏掖该文的最初动机作为补偿(连权的瑜伽一起,以便六年后以一个替补的父权把它端出来,再过八年,又等着Motif的开棺出土……)。无论如何,即便藏起(并留作独家收益……)这篇文章本质的算术动机,这篇文章也“站得住”,完全可读,配得上其作者把事情做得尽善尽美\dmark{}\origpage{293}的名声。而在这里,他所发展的理论,若没有那条启发性动机,便不可理解。于是他给出了动机,却以“关键陈述”这个定语来指涉它,同时把它从腿底下打发掉——不赐给它一个名字,不赐给它一个成形而冠以定理或命题之名的陈述;连“对应”(所谓Riemann--Hilbert的)都没有——他把这份差事留给了他的朋友Verdier与Teissier。他无须给它命名(鉴于其分量之轻\fnote{15}{(*)}{(6月14日)为安放这个“轻”,我提请注意:Deligne曾在IHES主持一场研讨班,试图发展可构造离散系数向连续系数的翻译,而未获令人满意的结果。对此可参见“当值无名之辈与好上帝定理”一注,n\textsuperscript{o}~48$'$。}——他肯定五分钟就能证明它!),也无须提名任何人——别人自会代劳,并且完全合他的心意。分明有一部瑜伽、一套哲学,作者以完美的娴熟与权威挥洒自如,而无须命名任何东西——那个他装作不屑的“轻”(“本应在这份笔记中占有一席之地”),他心里清楚:只要他知道在该沉默处沉默、肯等待,它自会外加送到。他头一回成功地玩这一手,那个“轻”是“权方面的若干考虑”,借一条谜语式评注的转角处影射了一下(且等六年后,再把权的哲学大张旗鼓地端出来)。据我所知,第二回是在我1970年离去之时——那个“轻”是“Motif之梦”,它在十二年间都不配得到一言的礼遇(想想吧——不过是个梦,而且还是个亡故者的梦,外加未经发表!),且待这一回发现真正的motifs(以及拿它们能做什么),再把那无可争议的父权——一如既往地谦虚——戴在自己头上\fnotevariant{16}{(***)}{(**)}{关于这套“以轻蔑来据有”的技术的其他评注,见次日之注,n\textsuperscript{o}~59$'$。}。

\subsection{反常性}

\notehead{76}{}

(5月4日)我清楚地记得,头一回听到“反常层”这个名字(译按:法语pervers还有“邪恶”的意思)——那应当是两三年前——它曾令我不适地一震;它在我身上激起一种不安之感。这份感觉,在此后两三次再听到这个不寻常的名字时,又出现了。那是一种内心的“却步”,停留在意识的表层;倘若当时我停下来检视它,它想必会表达为类似这样的话:给一件数学的东西起这样的名字,是什么主意!甚至给\dmark{}\origpage{294}别的任何东西或生灵起这种名字——严格说,唯有对人可以破例——因为明摆着,在宇宙间的一切“物”之中,唯我们人类,是这个字眼有时可以施用的对象……

我相当有把握(不敢说全然有把握):头一回对我谈起所谓“反常”层的,不是别人,正是Deligne本人——那是在吕米尼研讨会之后,他顺路到我家来的时候\fnote{17}{(*)}{倘若确实如此(如今我已深信不疑),那就得归功于我朋友的谦虚:因为(至少在意识的层面上)我并不曾疑心,引入并命名它们的人不是别人,正是他。是读了那篇“值得纪念的文章”,我才明白过来。(5月28日)老实说,文章本身对此也同样不置一词——正如它不置一词说Deligne是Riemann--Hilbert对应的父亲。然而,对“反常层”这一命名的父权,我从未有过怀疑;而且这一父权,此后确实得到了证实。}。那甚至想必是我们之间最后几场数学谈话之一——他到我家来之后,再没有过别的。恰恰是在那次来访时,显现了那个“征兆”;几个星期或几个月之后(当这一征兆在那次会面之后的数学通信往来中得到再次确认时),正是它使我断绝了数学层面上的交流\fnote{18}{(**)}{在纯属个人的层面上,这份关系仍以过去那种亲切友好的音调延续着,表面毫无变化。我的朋友习惯大约每两年,趁某次远足之便,来看我一次。去年夏天他又来过,那是一个可喜的机缘,使我也得以认识他的妻子列娜与他们尚且幼小的女儿娜塔莎。我记得那是从又一场吕米尼研讨会回来的路上;关于那场研讨会,我几乎没有听到什么回音(只有Mebkhout的几句阴郁而含糊的暗示——人们又一次好心地邀请了他,而他也想不出比再度入局更好的事……)。他们在我家住了两三天,接触自始至终都极好。}。(关于这一情节,见“两个转折”一注,n\textsuperscript{o}~66.)

回到那些所谓(我以为是错称的)“反常”层:显然,“正常地”说,这些层本该叫作“Mebkhout的层”,那才算公道。(不止一次,我曾把一些我提炼并研究过的数学观念,冠以先行者或同行的名字——他们与观念的联系,远不及Mebkhout与这个美好观念的联系来得近——何况这观念,依我看,其音色与其说“反常”,不如说“崇高”!)Deligne在发现并命名这个出自Mebkhout著作的观念之际的心境——他自己分明早已“奢侈到无边”,却仍准备巧取他人——这份心境,满可以正当地称之为“反常的”。肯定,我的朋友自己,在其内心深处,在某个不会为所炫示的门面所欺的层面上,也曾感受到这一点。在这个名字的授予里(乍一看,它显得悖谬),我感到一记\textbf{挑衅}之举,一种在如此完全的权力之中的醉意——它甚至可以炫耀(象征性地,借一个挑衅性的名字的陈列——其分明之极的真意,\textbf{没有一个人}肯去读!)它那巧取他人的“反常”本性。

\dmark{}\origpage{295}在我看来,绝非不可能:在某个深层,我曾知觉到我朋友身上这份心境的音色,而这曾促成我方才谈到的那份不安\fnote{19}{(*)}{我甚至倾向于认为确实如此。我不止一次在我身上确认:对事情的深层知觉,其精微与敏锐,与掠过意识层面、浮于意识表层的东西,全然不可同日而语。那个完全“醒着”的人,想必正是这些知觉被持续地整合进有意识的视野与有意识的经历之中的人——也就是那个完全按其真实的手段生活、而不只是靠这些手段中微不足道的一小部分生活的人。}。这份不安,尤其表现为:对他想必曾向我作的种种解释,我一概心不在焉;而我不相信,在这次会面之前,曾有哪次我没有持续专注地听他说话——尤其当事关数学的时候。我心中对这个被(天知道为什么)称作“反常”的观念,有某种阻塞——我并不真想听人谈它,尽管它与我一直(且在某种程度上依然)十分贴近的一些问题紧密相连。

干脆说吧:Deligne等人的这整篇文章,就是典型的、如假包换的“\textbf{Grothendieck学问}”,它们满可以同样出自我手(唯独主要观念的名字除外)!这大体上就是我在前一条注(n\textsuperscript{o}~75)第二部分里已表达过的;也是我通读所引文章的那一刻就已感到的——只是那份弥散的感觉,还不曾落实为这一触目的确认——它是我刚刚才做的。这份确认,以一种惊人的方式,使我重新对这桩深层的矛盾敏感起来:那个(在某种意义上)不能自已地复制着他所否认的人、并同化于他的人——那个他要埋葬的人,同时\dmark{}\origpage{296}也正是他\textbf{想成为}的人,而且(在某种意义上)他\textbf{就是}的人。

前天,在写前一条注(“不义——或归来的意义”)时,我已被这一巧合打动:我与朋友之间关系的这一转折——它骤然失去了一段共同热情中的契合,而那契合曾是这段关系之所以存在、并最强韧的内驱——恰逢我的朋友从那场值得纪念的研讨会归来之时;而那场研讨会的意味,刚刚向我显露。1981年7月我们相会时——在某个层面上,它与我们会面的其他场合同样友好、同样亲切——令我瞠目的,是那个“征兆”:论声调与神色,隐约,却又粗暴地分明,一份蓄意的轻蔑。它像是我的朋友预先支取的一笔\textbf{定金}——这一回,是在个人关系的层面上——对着那份隐含的、同样“隐约”(而且同样“粗暴地分明”)的轻蔑:那是他刚刚在吕米尼研讨会上、在当时一场当红名家间的技术精湛的盛大演出之中,公开对作为公众人物的我表达的。也是同一份“轻蔑”,刚刚对着那个胆敢(哪怕只有一点点)自称源出我门下的人表达出来(只是这一回,带着一种全然不同、“反常”得多的粗暴);而那人因此注定在我朋友皮埃尔眼里——至少在某个层面上——只成了“另一个Grothendieck”\fnote{20}{(*)}{在我们的个人关系里,我的朋友以我名字亚历山大的爱称(源于俄语)称呼我;这也是自我童年起,我的亲人、最亲近的朋友们对我的称呼。}——而如今,当务之急是不惜一切代价碾碎他……

\subsection{暂停!(Pouce!)}

\notehead{77}{}

(5月5日)翻阅这篇值得纪念的文章\fnote{21}{(**)}{关于“值得纪念的文章”,见注n\textsuperscript{o}~75。}——它(据说)主宰了1981年6月那场同样值得纪念的吕米尼研讨会——另一个细节打动了我。最后一章,顶着“从$\mathbb{F}$到$\mathbb{C}$”这个意味深长的名目,连篇累牍地描述一条卓越的原理——我把它引入代数几何,恐怕已有二十年了;那应当还在Motif观念诞生之前(Motif经由原Weil猜想,给出它最\dmark{}\origpage{297}深刻的例证)。这条原理断言:对涉及一域上有限型概形的某些类型的命题,只需在有限基底域上(亦即在一种“算术性质”的局面里)证明它们,即可推出其在任意域上、尤其在复数域上的有效性——在后一情形,所论的代数-几何结果有时可以经由超越途径改述(比如以整系数或有理系数上同调的语言,或以Hodge结构的语言,等等)\fnote{22}{(*)}{(5月6日)我记得,运用这样一条原理的第一个例子,见于Lazard关于仿射空间$\underline{E}$上代数群律的幂零性的定理(在任意域上)。它的证明曾深深打动我;我由它取得启发,处理了许多别的命题,并把它做成一套“哲学”,它主宰了我关于Motif理论的反思。}。我的朋友知道它,不是从别人、恰恰是从我嘴里,经年累月,就着许多例子\fnote{23}{(**)}{其中一个例子,见“驱逐”一注(n\textsuperscript{o}~63)。}。这条原理的父权(其初等形式甚至在EGA~IV里有明文——请别问我哪一节哪一条……)更是众所周知的\fnote{24}{(***)}{(6月5日)我自称这条原理的“父亲”,或许言过其实:据我所知,它的第一次应用当归Lazard(见前一条脚注)。一如在别的场合,我的角色,是感到他人一个想法的一般性,并把它系统化到使之成为一种“反射”或“第二天性”的地步。在权与Motif瑜伽的框架里,头一个使用这条原理的,很可能是Serre(而非我),凭他虚拟Betti数的那个想法——正是它,把我引上一部一般的权与Motif瑜伽之路。(Serre的那个想法,见注n\textsuperscript{o}~46\textsuperscript{9}。)诚然,流行的惯例是:一条已成通例的推理“原理”,其父权不归于可找到其最初踪迹的作者,而归于头一个知觉到其一般分量、把它系统化并推广的人。在这个意义上,可以说 N.~Katz的更正(下面这句话里谈到)把这条原理的父权归于我,是有道理的。}。以至于1978年赫尔辛基大会上为我这位光彩夺目的朋友授予菲尔兹奖章时, N.~Katz禁不住在颂扬 P.~Deligne的演说里顺带提了它一句,从而(不动声色地)更正了其显赫的受奖人那有点难堪的系统性“遗忘”。我是几天前才读到这篇演说的,与那篇“值得纪念的文章”同时。

无论如何,在这篇文章里,“算术”向“几何”过渡的哲学,被以这样的措辞呈献:任何一个不知情的读者都不可能怀疑,那位光彩夺目的主笔作者(请原谅\dmark{}\origpage{298}我的失言……)刚刚才发现这条分量如此之大的奇妙原理。

诚然,我没有为这条方法申请专利,而我光彩夺目的朋友也从未在任何地方说他是它那神奇的发明者;正如他并未明言他是那条著名的“对应”的父亲(请欣赏这个词,它散发着好一股十九世纪的馨香!)——那对应被谦虚地记在Riemann与Hilbert名下(这两位,倒也配得上给一位如此显赫的后继者的孩子们当教父);正如他在那部“值得纪念的文集”(LN~900,原书此处作LH~900)里也未指明:发明Motif、MotifGalois群以及与之相随的一整套哲学(他至今才掏出其中一段)的,分明是他;SGA~4½那部名著也同样无话可说——在那里,人们甚至再度不吝赐我以该卷“合作者”的荣衔:该卷从零开始如此光辉地发展平展上同调,还肯赏脸(尽管那两卷带着可惜的“多余细节的矿渣”等等)求助于SGA~4与SGA~5这两卷注定被遗忘的卫星卷,并且宽宏大量地承认后者的功劳:提供若干补充与技术枝蔓(其中有些甚至“非常有趣”)\fnote{25}{(*)}{关于“SGA~AT之役”(原书如此)的细节,见“白板”“格格不入者”“绿灯”“颠倒”四注(n\textsuperscript{o}~67、67$'$、68、68$'$)。}。

在所有这些情形里——也在最近五六年里我得以确认的许多别的微型个案里,尽管我从没想到过要\textbf{辨明我的不安}、给我所目睹或参与的事情起个名字\fnote{26}{(**)}{在一个个案里“辨明我的不安”的第一步,是不到三个月前在《收获与播种》里迈出的:那场(事后证明十分费力的——而且不无原因!)反思“注记,或新伦理”(第33节)。这场反思在那则注“年轻人的势利,或纯洁的捍卫者”(注n\textsuperscript{o}~27)里得到继续;然后,不到两周之前(在前一天发现“值得纪念的文集”(LN~900)所带来的冲击之下),又随注n\textsuperscript{o}~59“新伦理(2)——或抢夺的集市”重新开始。写那一注时,我心中对动用“抢夺的集市”这个相当粗鲁的字眼,还留有一丝踌躇。此后接连而来的发现向我表明:任何踌躇都不合时宜。}——在所有这些情形里,我认出同一套\textbf{做派}。我的朋友永远、彻头彻尾地处于“\textbf{暂停}”的豁免状态——他可以从容取用,带着那份由同侪与下侪的仰慕(无论如何是有根有据的)所给予的坦然良心;那仰慕,担保着完全的豁免。

\subsection{中国皇帝的新衣}

\notehead{77$'$}{}

\dmark{}\origpage{299}(5月7日)当然,那些看着我朋友Deligne行事、并且多少“在局中”的人——我是说那些并非刚下船、并非刚刚从当事者本人或其他当红名家(即便不都金光闪闪)的出版物里才学到“时下数学”的人,亦即那些了解来龙去脉的人——这些同行(毕竟,他们也还没有那么稀少!),在\textbf{某个层面上},是明白发生了什么的。在那几桩“有点扎眼”的个案里,他们必定感到了那份特殊的小小不安——那是我自己在这些“微型个案”面前、比起他们的个案微小百倍的情形面前,也不止一次感到过的。但他们所感到的是如此\textbf{庞大}、如此\textbf{不可思议},以至于它想必从未浮出水面——一如它最终开始在我身上浮出水面那样,经由一项\textbf{工作}:它表达为围绕着一个微型个案的两份文本(前一条脚注里谈到的那桩个案)。我确实没有听说,在我们的科学或任何别的科学的历史里,有过堪与比肩的事。对某些人来说,它不曾“浮出水面”,反而想必\textbf{成了风气},或至少被视为\textbf{正常}——只要一位分明是天才的、为万人所仰慕的人物,以天底下最自然的神色实践它,在众目睽睽之下,而这事从来不曾(据我所知)引起过只言片语的评注。

最近这几天里,我禁不住多次重新想起那个故事:“中国皇帝的新衣”。那位皇帝,受了无耻骗子的欺、也受了自己虚荣的欺,传谕说他将举行盛大的出巡,身穿举世未见的最华美的衣裳——那衣裳,是那些号称裁缝艺术家的骗子们斥巨资为他赶制的。而当他出巡,满朝臣子盛装华彩、前呼后拥,那些“艺术家”们打躬作揖,皇室倾巢而出——无论在仪仗里,还是在聚拢来瞻仰这第七大奇迹的人群里,竟没有一个人敢相信自己眼睛的见证;人人恪尽职守地赞叹,并为这身衣裳那无与伦比的华彩锦上添花——皇帝就披着它!直到一个混在人群里走迷了路的小孩失声喊道:“可是皇帝什么也没穿呀!”——刹那间,所有的人仿佛同一个声音,跟着这孩子喊道:“皇帝是赤身裸体的!”

我觉得自己就像那个孩子:他相信自己眼睛的见证,哪怕他所看见的,是如此闻所未闻、前所未见、为众人忽视并否认的东西。

\dmark{}\origpage{300}至于孩子的声音是否足以使某些人回到他们健全官能的谦卑见证,那是另一回事。故事就是故事:它向我们诉说着实在的某样东西——但它并不是实在\fnote{27}{(*)}{(6月14日)写完这条注之后,“中国皇帝的新衣”这个名字,向我显出是这场葬礼的一个天然的副题,表达它一个格外触目的方面。随后,反思转而面向我的全体学生,乃至数学界建制的“全体教众”,这一副题显得不再那么切合。然而,我终于明白:那则首先在我想到我的朋友Deligne时来到心头的寓言,同样适用于葬礼的各个侧面与各个插曲——它们每走一步,都在不可思议之中抵达乌布王式的荒诞(那不可思议,人人都恪尽职守地装作不见),而它却是真的。这个方向上的反思,尤其可参见“进步是挡不住的!”“研讨会”“受害者——或两种沉默”“玩笑——或权复形”“欺瞒”“掘墓人——或全体教众”各注(n\textsuperscript{o}~50、75$'$、83、85$'$、97)——没有一条特别关涉我的朋友皮埃尔。}。

\subsection{墓外相会}

\notehead{78}{}

(5月6日)才过去五天,我终于得到了那份慷慨的一包文献——来自我的朋友Zoghman~Mebkhout;其中尤其是那两份已被检视过的“值得纪念的研讨会”文本——一场围绕着一桩纪念碑式的\textbf{欺瞒}搭起来的研讨会!“不义——或归来的意义”那条注——我在其中尽力消化这桩新“事件”那令人难以置信的意味——就写在我收到这包文献的当天(5月1日的次日),犹带着发现的余温\fnote{28}{(**)}{连同“注记——或新伦理(1)”一节,这条注是仅有的不得不重写数遍的两篇注或节之一:因为头一稿(乃至第二稿)所“流出”的,还带着我一贯看待事物的方式所带有的全部惰性,远在有待检视的实在之下。}。

自4月19日——我终于得知Lecture~Notes那部“值得纪念的文集”(LN~900——见注51、52)——以来,这是关于宏大葬礼之种种盛典的第三项大发现;它在我看来分量也最大:既因它照亮了一些与我关系极近之人的作为,也因其作为一个时代(显然独一无二的时代——诚然,我对历史所知寥寥……)的“风俗画”所具有的含义。

第二项发现紧随第一项而来:即“Motif”的开棺\dmark{}\origpage{301}出土——它们被埋了整整十二年。继“值得纪念的文集”之后,轮到“值得纪念的研讨班”——那场从未举行过的“研讨班”,顶着假名字(无论SGA还是4½这个号),还附赠一篇幽灵论文的“第0状态”,更不必说(真正的)SGA~5研讨班的一场中心报告(它以晚出的面目出现,尽管比它早十二年);一场为战役的需要而被“借走”、不经任何手续的报告。这场光辉的战役,以及它在降临到可怜的SGA~5研讨班(从头、从尾、从中间被拆散!)身上的那些离奇坎坷中所扮演的角色,是在4月24日至30日间持续的一场反思中渐渐显露的。(对此可参见“同伙”“白板”“格格不入者”“信号”“颠倒”五注,n\textsuperscript{o}~63$'''$、67、67$'$、68、68$'$。)

那项发现刚消化完——与此同时,我“我的朋友皮埃尔”的回顾反思也接近尾声——而4月30日,我刚以“尾声——或全体一致的赞同”这条名字双重欢腾的“终注”,在这场无休无止的葬礼下面,骄傲地画上最终的句点(这回是肯定的——我总算走到了头!)——就在这时,这包倒霉的文献到了,把句点、尾声、版式、编号……统统推翻。对文献以及随附的评注与书信匆匆一瞥,便明明白白:我的句点完蛋了,那场头等葬礼的堂皇部署——我正准备修饰其最后细节——也完蛋了;我只好重新套上司仪的挽具……

不过天知道,我的朋友Zoghman有的是时间把局势告诉我!局势以潜伏的形态,想必已持续十年;以“急性形态”(这还是委婉语)至少三年——自那场研讨会起。在会上,他分明感到了风向,而无须等到次年、以其显赫的前老板兼庇护人为教父的那份高度官方的会议录问世。

他的论文答辩(1979年2月)之后几个月,他曾到过我住了六年的那个村子,想送我一份论文。不巧,几天前我刚离开(从此没有回去过,除路过之外……),退居孤独。他只遇到我的女儿,论文是她后来转交给我的。是第二年,我想\dmark{}\origpage{302},我们才终于正式相识,在蒙彼利埃大学;我们聊了一两个钟头。那时我对数学并不怎么上心,想必再不会记得一篇我只花了几分钟翻过的论文,也记不得其作者的名字。但这并没有妨碍接触的热忱。我清楚地记得一股即刻的相互好感。数学谈得不多(据我所记),谈的更多是或多或少的个人事情。Zoghman后来告诉我(此事我已忘却):他还是向我大致解释了一点D-模的“哲学”,而且他对这次会面很满意,感到我在从他那里获悉新事物时哪怕微微“震动”了一下,虽然这些新事物(在某种意义上)又是“意料之中”的。我印象最深的,是他的为人给我的印象——一种执拗而平静的力量之感,一个“闯将”的力量之感。在那个时刻,比起去年我们相会时、或其后的通信中,我更强烈地感到性情的深切投合——尤其在“闯将”这一点上。但两次会面之间的两三年,似乎已颇损耗了他……

我不记得在我们第一次短暂的会面时,Zoghman曾向我谈起他所工作的孤立,谈起那些曾是我的学生的“权威们”没有给他任何鼓励。倘若他有所暗示,想必不曾坚持。在当时,这件事本没有什么可让我惊讶的\fnote{29}{(*)}{(5月30日)这并不完全真——我把后来产生的失望情绪投射到了过去。我还记得,去年夏天与Zoghman相会时,我曾惊讶:我的上同调学家学生们(尤其Deligne、Verdier、Berthelot、Illusie)竟没有一个在Zoghman的工作中扶持他。这份惊讶,十来天后Deligne来我家时又一度重演(我想必向他提了Zoghman一句,没有遇到回音),后来又经与Illusie的一次电话再次重演。(对此可参见“欺瞒”一注,n\textsuperscript{o}~85$'$。)}。我说不好那是在1981年6月吕米尼研讨会之前还是之后\fnote{30}{(**)}{(6月3日)那是在之前——1980年2月,他论文答辩后一年。}。倘若在之后,他心里想必还压着滚烫的事情——而他全然没有那副模样。他更像一个知道自己想做什么、知道自己要什么的人,安静地走自己的路,不寻衅,也无人寻他的衅。

\dmark{}\origpage{303}后来我们没有继续通信。但我清楚地记得他;去年年初,我给他去了个短信,姑且一问:他是否有空着手去做一项为“温和拓扑”而设的宏伟奠基工作——它(依我看)只等着一个有他这般才具的人来做。Zoghman起初没有向我明说:他对这一前景并不真感兴趣;不过他似乎很高兴借这个机会再见一面。那时我离局太深,没有看清局势:我以为D-模理论已然功成、就此收官,一如凝聚对偶理论那样(注78$_{1}$),而Mebkhout也许正缺“大任务”。直到去年夏天我们会面,我才明白:就在他亲手开创的理论里,“大任务”并不缺——有些甚至还未动工,因为连被人看见都还没有!

无论如何,那是一个现成的机会,来第二次会面——这一回不是像第一次那样旋风式的。去年夏天,6月里,Zoghman在我家大概住了一个星期。在数学层面上,这次会面主要使我得以多少弄懂D-模的瑜伽。我“解冻”得很慢:与我旧日的上同调之恋已有些失联,而且尤其一头扎在《追寻栈》的写作里——它处在相当不同的音区。Zoghman看到我一面心不在焉地听,并不气馁;他不倦地卷土重来,带着动人的耐心。我想,当我终于明白那些大名鼎鼎的D-模不是别的,正是我很久以前所称的\textbf{模的晶体}、并因此于奇异空间上仍保有意义的那一刻,我被触发了。于是,我看见我那晶体-微分过去的整片直觉之网从被遗忘的深处浮起,我那“六种运算”过去有点生锈的反射,重新咬合……

于是轮到Zoghman有点跟不上;或者不如说,是事后他才决定,不肯把手指伸进那副齿轮(一如我的朋友皮埃尔不肯——尽管只要我还在场,他曾满腔热火……)。($\Rightarrow$注78$'$)

\notehead{78$_{1}$}{}

\dmark{}\origpage{304}然而,凝聚对偶确有一定数量的“精细”结果——尤其关于“对偶化微分模”的结构、它们与“朴素”微分模的关系、以及非光滑平坦情形的迹映射与留数映射——是我在五十年代末发展的,而且据我所知从未发表。尽管如此,就其实质而言,凝聚对偶理论(至少在概形的框架里),一如平展对偶理论(及其对局部紧空间的离散系数上同调的变体——由Verdier依平展的模型发展)、又如线性代数或一般拓扑,都显现为实质上\textbf{已完成}\fnote{31}{(*)}{(6月12日)就平展对偶而言,只要纯粹性猜想与“双对偶定理”尚未在全部一般性下得到证明,这就不完全真。}的理论,因而其本性是打磨停当、随时可用的\textbf{工具},而不是多少不为人知、有待探究与消化的\textbf{实体}。

\subsection{受害者——或两种沉默}

\notehead{78$'$}{}

我们的会面,是在一种友好信赖与亲切的气氛里进行的。然而这份气氛并未兑现其许诺。我如今明白:从那一刻起,我朋友心中的信赖就远非完全。那是那场大名鼎鼎的研讨会之后两年、会议录在Astérisque里刊出之后一年\fnote{32}{(**)}{(10月9日)Zoghman提醒我:这些“会议录”实际上迟至1984年初才问世。}——也就是他正付出一场骇人听闻的巧取的代价的时刻。可他直到仅仅四天前,才肯把这件事告诉我!去年他来的时候,是从又一场吕米尼研讨会\fnote{33}{(***)}{(5月7日)这里的记忆有一点点混乱——我倒觉得,他是正准备去参加那场研讨会。当然,从那时起,我所忆起的那些“辛酸而含糊的字眼”,缘由并不缺。但这辛酸,又因他在我家小住之后途经吕米尼而再度加剧。他从吕米尼回来后给我打了个电话,我由此得了些回声。从那一刻起,我就有一种十分清晰的直觉:他赶到吕米尼去,是为了享受被“那些人”(我没有细问是哪些人)怠慢的滋味——那些人慷慨地邀请他,为的是享受把他当无足轻重之辈来对待的乐趣。我想必把这话对他说了,或至少暗示了;这想必无助于改善当时我的朋友对我的心绪。}\dmark{}\origpage{305}回来的(这回干脆就是以D-模为主题的)——人们又一次慷慨地邀请了他,而他也赶紧应召而去。他谈起它,措辞既辛酸又含糊,话里带着意思:如今他“火中取栗”已毕,“一切都是别人做的”。我完全可以想象那幅图景——尤其Verdier,忽然想起三角范畴(连导出范畴也算上,反正一不做二不休!)的父权——那父权他搁置了十来十五年,勉强容忍他的“学生”Mebkhout在工作中使用它们……(注81)。

尽管他当时不肯把事情说清,Zoghman心里对Verdier想必积怨已深——这完全可以理解,看看他前老板那副不甚鼓励的样子。然而,我的其他上同调学家学生——Deligne、Berthelot、Illusie——也同样不曾屈尊对他所做的事情发生兴趣、多少扶持他一把。但人们几乎要说:对Zoghman而言,这本是理所应当——(人们会说)他从不曾在长辈们中间体验过别的态度。倘若他当时怨谁,那唯一而且不折不扣的,是Verdier。

从Zoghman的暗示(他分明无意细说)里,我懂了:“人们”在系统地把他的工作轻描淡写——一笔带过而已。这毕竟是人世间最寻常的事。对一件事的重要性的评判,在很大程度上是主观的;给自己的工作、给自己伙伴与盟友的工作,比对别人的工作——尤其对那些由于这种或那种缘故而人们有意贬低者的工作——记上更多的功劳与重要性,是寻常而几乎普遍的事。(而在这里,那个“缘故”对我而言实在算不得什么谜!)没有什么能使我疑心:远在这些寻常态度之外,这里有一场纯粹的巧取战役——其中根本谈不上“轻描淡写”,而分明是要把Mebkhout对那些使停滞之地重获生命的观念与结果的父系渊源\textbf{整个藏掖}起来……

然而,倘若天底下有一个人,我的朋友向他倾诉是再自然不过的,那正是我:在那几年执拗的工作里——有时辛酸,逆着时下的风尚——是我的著作给了他灵感;是我,在家里亲切地接待他,并且在尽力学习他乐于教给我的东西之时,轮到自己多少\dmark{}\origpage{306}做了他的学生\fnote{34}{(*)}{而且,Zoghman不但不对我谈他自己的下葬,也不对我谈我的下葬——尽管他置身头排,注视那场下葬的进行,已然将近十年!老实说,他的那些“庇护人”(边缘上还带着几分勉强)甚至曾让他亲手抬过装着我的遗骸的棺材的一角;但他们没有原谅他:他是满座宾客中唯一敢偶尔说出那个众人皆讳莫如深的名字的人!于是,我的朋友想必在与我的关系里自觉处境尴尬,而他未能在自己身上找到那份单纯,去承担一段(像我的一样)满载暧昧的过去,并不绕弯、清清楚楚地说话。谈他的下葬,也就谈我的下葬、谈他本人在其中扮演的角色……无论如何,我最终发现这场大名鼎鼎的葬礼的全部盛况,是顶着一种“沉默的共谋”才做到的;这一共谋既包括我的朋友Zoghman,也包括我的朋友皮埃尔,想必还包括我在数学“上流社会”里的大多数朋友。(6月3日)关于更多的细节,见随后的注n\textsuperscript{o}~78$''$。}。

我朋友怀着亲切的热忱来访之后,紧接着就有了一次“反弹”。我有这样的印象:他决意把八到十年间——在他那些曾是我学生的某些人所给的冷漠与轻蔑的刺痛下——积攒起来的疑忌与辛酸,移栽到我的身上。在其后的几个月里,我们之间的通信始终没有离开酸中带甜的音域——它最终停在一张贺年卡上,没有得到回音。

直到3月底,我才重新联络Zoghman,给他寄去“一段过去的重量”以及我当时为该节补加的各注(n\textsuperscript{o}~45、46、47、50)。我是想问他:是否同意我像已经做的那样,让他出现在关于我著作的那段简短反思里(在“我的孤儿们”一注,即n\textsuperscript{o}~46中)——尽管对人人都会一清二楚:我使用的是他给我的信息,而他或许认为那些信息是保密的。我全然没有把握:我的朋友会不会(如在他之前的别人那样)宁可“把自己压扁,也不愿招人嫌”。倘若真是那样,我会难过。

等他的回信,我觉得日子漫长——十天后才收到。我多少料到它还会是模棱两可的——但\dmark{}\origpage{307}这一回,它分明是热忱的。他毫无保留地表示同意,甚至为我在其中谈到他的那些措辞而动容。

在他那封长信(八页)的第6页,他才顺带地、就其定理那“令人印象深刻的”应用之众多(“无论在平展拓扑的框架还是超越的框架里”)提了一句:这条定理在文献里始终顶着“Riemann--Hilbert对应”之名\fnote{35}{(*)}{其信件的引文,见“一种不义与无力的感觉”一注,n\textsuperscript{o}~44$''$。}。他说得那么附带,而且那一手字迹潦草得仿佛故意,险些整个从我眼皮底下溜走!但我到底记住了;那真是一件怪事。怪到几乎难以置信;再说,也许我的朋友言过其实,分明他对人人都怀恨在心,包括我在内——尽管我心无恶意,这一点本是够清楚的。于是我补写了一条注(好个Zoghman,我还以为已经收尾了呢!),名之为“当值无名之辈与好上帝定理”,外加另外两条:“本能与时尚——或强者的法则”(写它时,在别人之中我也想到了他)与“罐装的权与十二年的秘密”。这条“当值无名之辈”的注,我起初写得并无完全的信念;Zoghman在我眼里纠结之至、充满矛盾,以致我自问:我这是要驶向哪里——我只做他的回声,对事实本身并没有多少第一手了解。我压根没有想到可能有巧取,更没有想到Verdier或Deligne本人牵连在内。Zoghman对我说过的话里,没有任何东西能提示这个……

然而,这两人都与那条好上帝定理系得太紧:不经过他们至少心照不宣的同意,其父系渊源几乎无从藏掖。这事想必在其后的几天里在我身上发酵。我想起Deligne曾对那个(十年后)由Zoghman解决的问题苦苦思索过;至于Verdier,他毕竟担任过研究导师;即便他没为学生费多少力气、宁可用冷淡与气馁来打发他,他至少应当知道这项工作里两条主要定理是什么——Zoghman肯定向他解释过,在那几场大名鼎鼎的、Verdier\dmark{}\origpage{308}“恩准过”的“会谈”里!于是我给这条注增补了一则评语,论Mebkhout的工作与Deligne一次较早尝试的关系,外加一条论Verdier角色的脚注。这同时也是对我朋友Zoghman的一次试探……

人们也许会想:这样一来,Zoghman必会抓住时机,亮出他藏了三年的炮群,让分明的真相轰然炸响,让受压迫者的大业凯旋!全然不是!十五天的沉默之后,来了一封信,谈数学里的一切,唯独不谈好上帝定理——或者不如说,关于它,他只限于按我的请求,给我其论文中的精确出处。(我到底想知道,这条我如此斩钉截铁地担保过的大名鼎鼎的定理,在哪里得到证明!)

我不得不在回信里,就“对我的著作的宏大巧取”——我刚发现的那桩(随“值得纪念的文集”LN~900,外加“预支了许多乐趣”:不日将在系图书馆里结识SGA~4½)——说上几句,才在又一个十天的沉默之后,使我的朋友终于引爆!

这一回他总算“倾囊而出”——而且是一\textbf{大}包:挑选得当的文献,使我(一个不常光顾图书馆、甚至连我在大学办公室里越堆越高的抽印本也不常翻的人……)得以对一种“气候”形成均衡的印象;在那气候里,许许多多的人,至今仍未参与我那漫长而庄严的葬礼\fnote{36}{(*)}{(6月12日)不过,Katz、Manin、Langlands似乎不在其列……(1985年3月)关于Katz的另一种说法,然而可参见“把一个个 i 点清”一注(n\textsuperscript{o}~164(II5))与“调动”(n\textsuperscript{o}~169)“插曲2”。(1985年4月)Langlands亦然,见“预先出土(2)”一注,n\textsuperscript{o}~175$'_{1}$。}。除了主要的“铁证”(那场大名鼎鼎研讨会的两篇文章,它们使那桩不可思议的欺瞒原形毕露)与另一篇“值得纪念的文章”(这一回出自Verdier的手笔\fnote{37}{(**)}{关于这篇文章,见“好的参考文献”一注,n\textsuperscript{o}~82。})之外,还有 N.~Katz论“菲尔兹奖得主”Deligne的演说,外加Langlands与Manin在同一届1978年赫尔辛基大会上的报告;然后是Deligne1970年尼斯大会上的“Hodge理论I”(其中第3行又一次影射“Grothendieck的猜想性Motif理论”(注78$'_{1}$))\dmark{}\origpage{309},以及同一位Deligne1974年温哥华大会上的“代数簇上同调中的权”(其中我的名字未被提及(注78$'_{2}$));最后还有与 A.~Borel(又一位老哥们——我同时得知他已回苏黎世……)的通信,以及Mebkhout在法国科学院报告上的两条短文,其中1980年那篇是其学位论文第五章(前一年答辩)的摘要,把好上帝定理更突出了一点\fnote{38}{(*)}{这条注的精确出处——Mebkhout的论文与好上帝定理——见“大部头与上流社会(或:鱼目与珍珠……)”一注,n\textsuperscript{o}~80。}。还不算一份文件——嘘!——以保密之名传达,我在这里一个字也不多说……

两封信随这份厚实的包裹而来(4月27日与29日),一封很长,两封都很有分量。如今他终于把“底”(这一回是真的!)亮出来了,Zoghman却仍旧像自我重新联络他以来那样,力劝我极端谨慎。倘若听他的,我会小心不把我的反思笔记公开,它们将只在我们二人之间保持绝对的秘密——至少牵连任何人的那一部分,因为“他们”“大权在握”,“人人都站在他们一边”\fnote{39}{(**)}{(5月30日)我兴之所至,在这里说得有点过头。Zoghman从未在任何时刻建议我不要发表我笔记的这一部分或那一部分。近来他甚至坚持:这些笔记理应以书的形式问世,以造福“后世”——一份限量印制、类似预印本的东西,在他看来有点“竹篮打水”。}!然而,我分明告诫过Zoghman:这些笔记——我寄给他的只是与他相关的摘录——注定要公之于众,而且越快越好。

一切要素似乎终于齐备,足以使受压迫者的公义大业凯旋;但那位“受害者”似乎仍在竭尽全力,存心把水搅浑——仿佛带着一种(人们会说的)事后的懊悔,后悔亮出了那道大名鼎鼎的“底”——直到命中注定的5月2日为止,Zoghman想必是它唯一的持有人。这份暧昧,几乎每一行都透出来(我几乎没有夸张),直到我刚收到的最后几封信——包括最后那封:他以一种阴郁的凯旋神色,把“值得纪念的文章”全份寄来(而先前那“大包”里,他到底只舍得与这份关键铁证的头二十页分手\fnote{40}{(***)}{(10月9日)Zoghman向我说明:其实他手头起初并没有全文的复印件,是后来才复制的。}。)

\dmark{}\origpage{310}至于皮埃尔朋友——我是说Deligne(对世人而言,他既非皮埃尔,也非“朋友”……)——他简直恨不能含着热泪唱起赞歌——这样一来,倒仿佛不再是Zoghman是“受害者”,而是可怜的Deligne:他被周围的人带坏了,毒害至深;唯一的恶人、把他带坏的人,是Verdier(而且……不妨顺着我的目光看……):分明是我“一定得罪过Verdier”,以致他恶意到只为害人而害人;更不必说,做他导师的也是我,授予他博士头衔、荣耀以及其他一切的也是我——总而言之,“绝对权力”的手段全在我!\fnote{41}{(*)}{这已不是我头一回听到“绝对权力”这样的调门,人们想借此让自己相信自己的无能,并为之辩护。倘若有人曾把对自己本人的“绝对权力”托付给什么人,那么对Zoghman而言,那不是别人,正是Zoghman自己!}

显然,倘若我的朋友怨谁,他真正怨的,并不是他那位显赫的前老板——那个人,他十年里有幸见面“会谈”总共不过三次(照他最近写信对我说的来理解)——一个遥远得令人眩晕、全然够不着的人;他怨的,是那个他想来就来、可以同吃同住的人……\fnote{42}{(**)}{(5月8日)我朋友对我的关系中,冲突的毫不含糊的征兆,恰恰出现在那次小住之后——那次他“分享了我的面包与住处”,气氛毫无保留地亲切,消除了我们第一次短暂会面想必不曾完全抹去的“距离”感——这肯定不是偶然。我在这里遇到一个我由来已熟悉的局面;我曾在“敌父(1)、(2)”两注(第29、30节)里(以相对一般的措辞)表达过它。当年我写它们,作为对先前反思的评语时,我不曾料到:我在其中描绘的那个原型局面,竟会始终处于一场尚在到来的漫长反思的中心——而那时,我还自以为已近旅程的尽头!}

Zoghman每前进一步,披露一点新东西,使我更多地了解他身为受害者的那场巧取局面(并能多少有助于解开它),我都感到那像一次\textbf{挣脱}——一场筋疲力尽的内心斗争的终点。有一个\textbf{角色},他似乎已以身家性命相认同,死死攀住,如同攀住他最珍贵的财产——那\textbf{受害者}的角色\dmark{}\origpage{311};而要保持这个角色,他就必须围绕这个角色以及为之辩护的局面,保持最绝对的秘密\fnote{43}{(*)}{(5月30日)自这些文字写成(5月6日)以来,我朋友的态度已发生急遽的演变;近来我不再觉察到他对受害者角色的攀附。当然,下面的文字(一如上面的文字)所关涉的,是我朋友生活中的某些情节,绝不声称刻画一种性情、或描绘一种恒常的成见。}。而此刻,在他不情愿的配合之下——那配合仿佛是被一个局面的逻辑硬挤出来的;而那个局面,不是别人,正是我带着这些多事的、关于一场风平浪静的葬礼的反思造成的——这个秘密行将终结,连同也许还有那个他不知从何时起便乐于固守的角色;在这样的时刻,他确实可能痛苦欲裂,比以往任何时候都更怨我。

我朋友Zoghman的这场“下葬”,是由\textbf{两种沉默}的合力做成的:每一种回应着另一种、又反过来挑衅它,在一轮毫无罅隙的回旋里,一方的角色紧紧贴合一方的角色——巧取者们与被巧取者。倘若我不止一次被这样一个景象攫住:“下葬者”同时、而且更深地,是他自己的“被葬者”,那我同样被另一幅景象攫住:在另一位朋友的身上,一个“被葬者”同时、而且更深地,是他自己的“下葬者”——与那些他乐于为之充当自愿受害者的人,紧密共谋。而且我看得很清楚:他自己那场巧取的第一责任人,不是别人,正是我的朋友Zoghman本人——三年来,他经由沉默,默许那些对他恣意妄为的人对他的羞辱。他手里握有一切,可以一战;他却三年来选择忘记自己还有手,未曾抵抗便告败北\fnote{44}{(**)}{(5月30日)这当然是一个有斗争家气质的人的主观看法;那种气质,在这位身上似乎阙如。不过,看来这些文字写成之后,战斗的心弦在我朋友身上已经苏醒,他已决意与那桩让他付出代价的不义一战。(1985年4月18日)关于对我朋友心境的一种不同、也不那么“苛刻”的照亮,另可参见“根”一注(n\textsuperscript{o}~171$_{3}$)。}

\notehead{78$'_{1}$}{}

这份简短的预注,我从不曾亲手捧读过;我读到的只是更详尽的“Hodge理论II、III”——发表于《Publications Mathématiques》。正因如此,我一直有这样的印象:Deligne不曾认为有必要哪怕暗示一下,他关于Hodge理论的那些想法之生成中,\dmark{}\origpage{312}Motif理论所起过的角色。我对自己说:倘若他有意愿提及我可能在他身边扮演过的角色\fnote{45}{(*)}{(5月30日)而且直到几星期前,我还在系统地淡化这一角色。对此可参见5月27日的“格格不入者”一注(n\textsuperscript{o}~67$'$),我在那里头一回觉察到我身上的这份态度,并隐约瞥见其意味。},他无疑会借着构成其学位论文工作的“Hodge理论II”来做——那正是提及这类事情的天赐良机\fnote{46}{(**)}{(5月30日)我也不记得曾被邀参加其论文答辩委员会。葬礼那时已然进行得热火朝天……}。我刚刚看到:他经由这句碑文式的话\fnote{47}{(***)}{Serre也经由指涉记号[3],在同一行里含蓄地露面——好奇的读者,在“Hodge理论I”的参考文献里能找到他的名字。这条草草的对照线,想必是从1968年到今天之间唯一的一处:对它一口气提到的那些“源头”(Serre(别名[3])、motifs、Grothendieck……)——无论多么谜语式——的一点影射。(5月28日)然而此后我又撞见另一处这样的影射;鉴于那个场合之特殊,它格外有趣。对此可参见“葬礼悼词(1)——或恭维”一注(n\textsuperscript{o}~104),以及紧随其前的“掘墓人——或全体教众”一注(n\textsuperscript{o}~97)的结尾——后者安放了那个“特殊的场合”。}——影射“Grothendieck的猜想性Motif理论”——一劳永逸地尽了提及我的手续,外带一个出处(Demazure在Bourbaki研讨班上的报告)。

又是无话可说!他从未想到要指明:他学得这套理论(且别忘了,整个是猜想性的!)是通过\textbf{别的来源},而不是Demazure那份寒碜的文本——后者绝无可能给出一套如此丰富的理论(整个是猜想性的!)的图像;那套理论,以底纹的形式贯穿Deligne其后围绕权的瑜伽的全部著作——且待那部海盗文集LN~900的步步升级, MotifGalois群(这回连一行含着亡故者名字的简短出处都没有……)才终于(十五年后)在那里开棺出土。

反思之后,在这句简短的引文里,我认出同一套“暂停!”做派——一句纯礼节性的引用,为的是把手续办清;附带的那个出处,绝无照亮读者之效(在这个案子里,即照出与那些恰好需要藏掖的想法之间那些分明而深刻的关系\dmark{}\origpage{313}——而它们在随后的十二年里也确实保持着藏掖),却\textbf{有欺瞒之效}\fnote{48}{(*)}{写这些文字时,一个联想不由分说地来到:早两年发生的、围绕“权”的第一起显影事件;它在“罐装的权与十二年的秘密”一注(n\textsuperscript{o}~49)的开头被谈到,并在“驱逐”一注(n\textsuperscript{o}~63)开头处得到更详细的论述。至于一般的“暂停!”做派,见“暂停!”一注(n\textsuperscript{o}~76,原书如此,题名对应注77)的反思。这套做派,开始对我变得非常熟悉了!}。

\notehead{78$'_{2}$}{}

这份文本\fnote{49}{(**)}{“代数簇上同调中的权”,P.~Deligne,1974年温哥华大会,《会议文集》,第78--85页。}我无需亲手捧读(我是仅仅几个星期前才知道它存在的),就知道我的名字不在其中。Serre的名字同样不在——而最早隐约瞥见“权的哲学”的人,正是他;那套哲学,是我随后以极大的细节提炼出来的。

\subsection{老板}

\notehead{78$''$}{}

(6月3日)Zoghman向我解释:他只是逐渐地、起初还模模糊糊地,才意识到围绕我的著作所做的那场“巧取”。1975年Verdier给他的手稿(见“好的参考文献”一注,n\textsuperscript{o}~82)对他曾如同天赐,尤其向他引入了可构造性观念及其本质性质,以及双对偶定理——他从后者取得启发,做出了D-模语境里的双对偶定理(或“局部对偶”定理)。只是多年之后,在读SGA~5时(屠杀版诚然,但还没有被屠杀到足以瞒过一个像他那样留心的读者),他才开始觉察到什么。长年以来,他对那位遥远的长辈满怀钦佩与感激,深信他大量汲取的那些观念属于后者。甚至似乎有多年之久,他确实深信:那套所谓“Verdier对偶”的对偶理论确实归Verdier——或至少归“Serre--Verdier”;同样,他所称“Poincaré--Verdier对偶”的那个对偶想法,也确实归Verdier。直到1979年前后(他答辩之年),他才开始觉察有什么不对——但我料想,他必定小心地不对其显赫的“老板”流露半点,一如在我们1980年2月与1983年6月会面时,不对我也一样。直到“反常”研讨会,1981年6月,当他开始感觉到\dmark{}\origpage{314}自己本人的著作正在被人藏掖,他才同时开始更清楚地认识到:自己闯进的是怎样一个世界\fnote{50}{(*)}{Zoghman那时对他前老板的评价,已低落到这种地步:他甚至深信,Verdier六十年代所做的一切(我在注n\textsuperscript{o}~81“赊账的学位论文与全险”的一条脚注里评阅过),或多或少是我在口授、或至少在提示。}!肯定,对他而言,我必属于这个世界:在那里,我的旧学生们(或至少其中某些)占据高位,以对待亡师的同样漫不经心,劫掠着遗孤学生。倘若真有区别,那不过是:我是亡者,而他们活得好好的,并且以确凿的方式证明着这一点……

我可以想象,甚至在“反常”研讨会之后,Zoghman仍难以相信自身健全认识能力的见证——尽管那见证分明地教给了他发生了什么。那篇大名鼎鼎的、由B.~Teissier与其“老板(原文如此)”Verdier署名的研讨会会议录导言,他直到1984年1月才亲手拿到。在驳回那份明摆着的事实近三年之后,震动想必愈发剧烈——我自以为是这样理解的。两个月之后我才重新联络他,3月底给他寄去“我的孤儿们”与“拒受遗产——或矛盾的代价”各注;又过了一个月,他才终于决定向我“亮底”,把“反常层会议之欺瞒”告诉我。

\subsection{我的朋友们}

\notehead{79}{}

瞧,我正准备把这场反思收尾并公之于众——它将终结Zoghman自己在那场他付出代价的巧取周围保持的沉默;而那场巧取的阴暗收益,他也照单全收\fnote{51}{(**)}{(5月30日)我提请注意:这场反思所由生的,是我朋友身上的一些心境;它们如今看来已被超越。(可与注n\textsuperscript{o}~78$'$的5月30日两条脚注对读。)}。也许它会令他不快,一如也许会令我的朋友皮埃尔不快——一旦它完成、文本誊清并印出,我就会把它亲手送到他手里\fnote{52}{(***)}{我原不曾想到,在我前面尚余的年月里,还有机会回到首都住上几天。但十多年来,我的朋友皮埃尔频频移步,到我蛰居的乡野深处来看我;于是这一回,趁着这个例外场合,轮到我移步了——那同时也是对一个屡次发出、而我从未应邀的邀请的回应。}。我能献给我朋友Zoghman与朋友皮埃尔的\dmark{}\origpage{315}最好的东西,他们两位多半会把它当作最坏的东西来接受:当作一场灾祸,或一桩冒犯。更何况,我的见证是公开的——一如他们两人的沉默,都是公开的仪式性行为,使这一位和另一位都须承担责任。

拒绝或接纳我的见证,是他们的选择;对Jean-Louis亦然——我曾把他算作我的朋友,一如今天的Zoghman与皮埃尔。这些选择与我切身相关,却不是我的选择。我毫无预言它们的兴致。我不久就会知道;我怀着强烈的兴趣、一份悬念——却没有一丝焦虑的阴影——期待着到来的星期与月份将带给我什么。我唯一的关切、我唯一的责任,是:我所献出的,确实是我所能献出的最好的东西——也就是说:真。

也许有人会惊讶:我竟毫不容情地谈论我以“朋友”之名相称的人;并在这名字里看出一句套话,甚至看出其中并不存在的讥讽声调。当我把Zoghman~Mebkhout或Pierre~Deligne称作“朋友”时,那是对同情的、亲切的、敬重的感情的重申——它们在我写下这些文字的时刻,在我身上在场。敬重告诉我:我无须“容情”对待一位朋友,一如我无须“容情”对待自己——一如我,他有资格与那份寒微的真相相遇;一如我,他无须容情。

倘若我不以“朋友”之名指涉Jean-Louis~Verdier,那绝不是因为我认为他比我的朋友Zoghman与皮埃尔、或比我自己更不“好”、更不“配”,而是因为生活恰巧把我们分开了。十五年前、乃至更久以前把我和他系在一起的同情的、亲切的感情,已被时间或多或少磨蚀,不曾有机会经由哪怕一点个人的接触而复活。我为重建这种接触做过的寥寥几次尝试,没有得到回音;我不知道,这些反思的阅读,会不会使一段已然僵死的关系复活。但即便此刻他于我不再是“朋友”,我也不认为:我像对待自己、对待我的朋友们那样毫不容情地对待他,就是不敬;而且我深知,反其道而行,于他、于任何人,都无益。更不必说:无论他还是我的朋友皮埃尔,倘若他们真想“自卫”(或进攻)而不肯冒险看自己一眼,并不缺手段,也不缺后援。更不必说:在他们有条件\dmark{}\origpage{316}令人气馁、或碾压人的地方,这一位与那一位,都不止一次这样做过,毫不容情,也毫无怜悯。

\subsection{大部头与上流社会(或:鱼目与珍珠……)}

\notehead{80}{}

(5月9日)况且,也该是时候了:我到底给出那条大名鼎鼎的Riemann--Hilbert——(不具名的Deligne)——亚当与夏娃——上帝——(尤其不是Mebkhout)定理的出处:人人都在大量引用它(包括我自己),而显然还没有人想到要问它在哪里得到证明。我自以为从我的朋友Zoghman那里懂得了:那条“值得纪念的定理”就在他的论文里;我果然在其目录里找到了它,顶着(确实土气、确实配得上莽汉的)名字“一个范畴等价”,第三章,第3节,第75页。更倒霉的是,它甚至没资格叫“定理”,而是叫“命题3.3”(而且更糟的是,我的名字出现在同一页上,还带着着重号!)。我承认,由于没有读过前面那75页来弄清脉络,我甚至不完全有把握是不是它——Zoghman向我确认是的;我信他\fnote{53}{(*)}{(1985年4月17日)最终显明:人们通常使用的“好上帝定理”的形式,并非此处所引定理的形式,而是一个相近的形式,可以用同样的方法证明。见“一个视野的绽放——或闯入者”一注(n\textsuperscript{o}~171$_{1}$),尤其其中今日所写的那条脚注。}。证明(据说)是同一篇论文第五章的对象——论文于1979年2月15日在巴黎第七大学答辩,答辩委员会由 D.~Bertrand、R.~Godement、G.~Houzel、Le~Dung~Tráng、J.L.~Verdier组成。有兴趣的相关人士,倘若尚未收到作者寄赠的副本(凡他无论对错地疑心可能感兴趣的人,他都寄了论文),只消向他开口,他将乐于从命……他当然给每一位我的前上同调学家学生都寄了副本;没有一位给过活着的音讯。想必他们其间都改了行,真不巧……

得承认:Zoghman做买卖的本事,确实不怎么样,不会以清楚而诱人的方式呈献他的货色——这是学得来的手艺;而他不像我的前学生们有那份运气,得以跟一位行业高手、且从不吝惜时间的高手学徒。可他也无从抱怨:他那“三次会谈”是得过了;说不定哪天,哪位“权威”会想起来,甚至为他那份难啃的大部头回一句“已收到”。他自己想必也发觉了:这份大部头很难啃(即便它对Riemann与Hilbert都并非白费……):他在法国科学院报告上发了一条短文——那毕竟短些——以引起人们对其大名鼎鼎定理的注意;题目嘛,您猜一百次:《论Hilbert--Riemann问题》!我早就知道,我的朋友Pierre~Deligne的历史修养并不比我强:他只需把年代顺序拨正\dmark{}\origpage{317},再贡献出那个漂亮的民间传说式的称呼“对应”,戏法就完成了——Zoghman,这可是他自找的……这条短文是1980年3月3日的,Série~A,第415--417页。

Verdier想必是在他“恩准”其学生的那“三次会谈”之一里(或在答辩时)得知这条定理的;但倘若真要问,他想必什么也没察觉。Deligne呢,他到底何时察觉到了什么,我说不好;但可以肯定:1980年10月他已知情,而据他自己说,Bernstein与Beilinson也是。Mebkhout还亲自去了莫斯科,向Beilinson与Bernstein(唯恐他们读起来费劲)详详细细解释其结果。我不知道是他们还是Deligne读过那篇论文或其后那篇科学院短文;但想必他们到底弄懂了里面的东西,因为次年那场大名鼎鼎的吕米尼“研讨会”恰好围着它转——出于天底下最大的偶然。

总结一下,并把我情报部门最新惠予的讯息算在内:对局势完全知情、并参与了那场所谓“反常层会议”之欺瞒的,至少有五人,即(按演员姓名的字母顺序)A.A.~Beilinson、J.~Bernstein、P.~Deligne、J.L.~Verdier与 Z.~Mebkhout——外加整整一场研讨会的与会者(原文称他们为acultes,疑为adultes“成年人”):肯定还都是光彩夺目的数学家,他们分明不求别的,只求被欺瞒、把鱼目当珍珠\fnote{54}{(*)}{(6月3日)实际上,看来研讨会的全体与会者,无一例外,都已在现场被告知了局势。对此可参见今天写成的“研讨会”一注,n\textsuperscript{o}~75$'$。}。这再次证明:我们数学家——从显赫的奖章得主到无名的学生——并不比张三李四多一丁点聪明或明智。

\section{VIII.~学生——别名老板}

\subsection{赊账的学位论文与全险}

\notehead{81}{}

\dmark{}\origpage{319}(5月8日)我觉得是时候了,该对“幽灵论文”这桩事作更周详的表述了;此前我只在两条注里(注48与63$'''$)“顺带”谈过。一位不细心的或不怀好意的读者,也许会说,我同时责备我的前学生 J.L.~Verdier两件自相矛盾的事:既“埋葬”了导出范畴,又把它们“发表”了(在SGA~4½里)并自居其父;正如同一读者会说,我责备 P.~Deligne,既“埋葬”了Motif,又把它们开棺出土(在LN~900里)。因此,或许并非多余:对1960年至今的局势作一番回顾。

1960或1961年前后,我向Verdier提议一项可能的论文工作:发展同调代数的新奠基,基于导出范畴的形式主义——那是我在此前几年里,为概形语境中的凝聚对偶形式主义之需而提炼并使用的。当时已经说清楚:在我向他提出的纲领里,前景中并没有什么严肃的技术困难;主要是概念性的工作,其出发点已然在手,或许需要相当规模的展开——规模堪比Cartan--Eilenberg那部奠基之作。Verdier接受了这个题目。他的奠基工作进展顺利,于1963年落实为一份关于导出范畴与三角范畴的“第0状态”,由IHES油印。那是一份50页的文本,1977年作为附录复制进SGA~4½(如注63$'''$所说)\fnote{55}{(*)}{单是这份文本,对一位天资出众的青年研究者两三年的工作而言,似乎嫌瘦。但Verdier当时的大部分精力,用于修习同调代数与代数几何不可或缺的基础——尤其通过听我的研讨班,以及两两对谈的工作。他对对偶形式主义的贡献(见下文)来得较晚:在我与Artin于SGA(1963/64)里详细发展了平展对偶的形式主义之后,当时我建议他(在其导出范畴的奠基工作之旁),在“寻常”拓扑空间、及这类空间的可光滑化态射的框架里,发展同一形式主义。大约在我以SGA~1开创“代数几何研讨班”丛书之际(1960年),Verdier联络了我,与Jean~Giraud、Michel~Demazure同时——问我有没有工作给他们;而他们敲的正是对的门!一件曾打动我的巧合:打从我写“我的孤儿们”一注(n\textsuperscript{o}~46)起,我记得他们三位联络我的时候,恰好刚组成一个小研讨班,名为“孤儿研讨班”(主题:自守函数,计算进路(原文approche calculs à brin de zinc,其中à brin de zinc措辞待考)),因为他们的老板(或CNRS的教父?)不打招呼就走了一年,让他们意犹未尽,又有些茫然。这片空,很快就被填上了……}

\dmark{}\origpage{320}答辩之所以没有在1963年、而是1967年举行,是因为难以想象:这份50页、尚待来临的奠基工作的胚芽,能构成一篇国家博士论文——这个问题当然根本不会被提出。也是出于同一理由,1967年6月14日答辩时(评审委员会由 C.~Chevalley、R.~Godement与我——主席——组成),根本不打算把这项工作作为论文呈上。提交委员会的文本,17页(加参考文献),以一项正在撰写的大工程的\textbf{导言}的面目出现。它勾勒这项工作的主要想法,并把这些想法安置在其众多使用的语境之中。第10、11页,对这项奠基工作所计划的章与节作了周详的描述。

倘若 J.L.~Verdier的理学博士头衔,是凭这份17页、勾勒着连他自己都说并非出于己手的想法的文本授予的\fnote{56}{(*)}{论文开头写道:“本论文是在 A.~Grothendieck的指导下写成的。它所包含的本质观念出自他。没有他最初的启发、他始终如一的帮助、他富有成果的批评,我不可能把这项工作进行到底。愿他在这里找到我深深的感激的表达。\par 我感谢Claude~Chevalley肯担任我论文答辩委员会主席,并有耐心读完这份文本。\par 我感谢 R.~Godement与 N.~Bourbaki,把我领进数学之门。”——按:“本论文”一词,恐怕只能指那项着手了的奠基工作的全体——所提交的文本构成其导言;而严格说来,那项工作在答辩之时并未“进行到底”。(5月30日)这一不自洽,恰好映出那份局势的暧昧;而那份局势,我作为论文导师——并且(照我手头那份论文封面上的说法)作为委员会主席——是第一责任人。我对一位光彩夺目的学生不够“严格”,有所纵容;这一纵容的方向,与我对Deligne所表现的(见“格格不入者”一注,n\textsuperscript{o}~67$'$)如出一辙,并且出了一份力,使同样的果实结了出来。},那显然是委员会与他之间的一纸本着相互信任订立的契约\dmark{}\origpage{321}:他承诺把这项工作做到底、并公之于众——对这项工作,他呈献了一篇光辉的导言。这份契约,候选人没有履行\fnote{57}{(*)}{越发引人注目的是:1983年12月, J.L.~Verdier竟拒绝了我请他加入Contou-Carrère论文答辩委员会的提议——委员会里还有 J.~Giraud,而我任研究导师;他的理由是:那篇论文(分明已全部写成,并经 J.~Giraud细心读过)与那个委员会,不经\textbf{巴黎}各大学论文委员会(原文如此)的核准,不具足够的严肃性的担保。}:他宣布的那份文本——一份按照一个已被证明有效的新观点写成的同调代数奠基文本——从未发表。

显然,倘若Verdier1961至1967年间的工作,仅限于写下1963年那份骨架式的“第0状态”,委员会绝不会想到接受这份“赊账的论文”。当时,其工作的成文想必已有相当的进展,可望一两年内完成;而出于实际的缘由,让Verdier不必等到那项本该为之奠基的工作完成就取得头衔,显得合宜。

还得补上一句:1964至1967年间,Verdier对对偶形式主义做出过几项有趣的贡献(注81$_{1}$);这些贡献,连同他本该继续的那项奠基工作,可以佐证人们给他赊的这笔账。单是他的那些对偶贡献,严格说,本就足以构成一篇像样的博士论文。然而这样一篇论文,绝不合我素来命题的风格:那些题目,都是对一个我感到其需要与紧迫的理论,做系统的、贯彻到底的发展(注81$_{2}$,原书误作82$_{2}$)。我不记得Verdier曾想过提出这样一篇“凭既有成果立论”的论文;我怀疑我是否会接受——因为这样一篇论文,与他我之间立下的“契约”毫无对应:我把导出范畴这个美好的题目托付给他,责成他发展大面积的奠基。

我承认我全部的责任,作为 J.L.~Verdier的论文导师与委员会主席:我轻率地(与 C.~Chevalley、R.~Godement一同——他们信任我所给的担保)把博士头衔授予一项尚未做成的工作\fnote{58}{(**)}{在这份责任之外,我还该补上另一份:在其后的两年里(在我退出数学舞台之前),我未曾督察Verdier是否确实履行了他所立的契约。得承认,我的精力当时全部投在我自己承担的奠基工作上——更不必说Motif的与其他的种种反思——以致那桩提醒他人尽其义务的不快差事,我顾不上多想。Verdier放弃出版那项预定工作的决定,我是七十年代初才知道的——那时我已全然不再关心数学,“反应”的念头根本无从起。}。

\dmark{}\origpage{322}倘若我今天确认这份轻率的某些果实,我原无权抱怨。但这并不妨碍我公开地做这份确认,并且声明:我的前学生 J.L.~Verdier的所作所为,只关他自己的责任,不关任何别人。

不履行对我、对那位曾信任他的委员会所立的契约,是一种埋葬的方式:埋葬我引入、并由他承揽以一项大工程去奠基的导出范畴的观点。这项工作也许做过,却从未交到使用者手里。这是对一整组他自己也曾协助发展的想法“画一道叉”。

Mebkhout的著作重新拾起导出范畴的观念,没有从Verdier那里得到任何鼓励(同样也没有从我其他任何一位身为“上同调权威”的学生那里)。对导出范畴的事实抵制,在我看来,直到1981年前后都是彻底的\fnote{59}{(*)}{(5月30日)这些带几分存疑语气的措辞,其实并不合宜。Zoghman~Mebkhout向我确认了(他是付出代价才知道的):我对“Grothendieck风格”的同调代数所受待遇的存疑表述,与实情完全相符。}——就在那年,它们随着需求的突然推挤,在那场值得纪念的吕米尼“研讨会”里强力重返(见注75)。

然而,Verdier那份“论文”的“第0状态”,早在此四年之前,即1977年,就已作为附录刊于SGA~4½卷(见注n\textsuperscript{o}~63$'''$)——也就是其论文答辩十年之后,而在(据我所知)Mebkhout是唯一在其工作中使用导出范畴的人的时刻,逆着此前七年的风尚。如无差错,他一直是唯一的一个,直到那场已提到的研讨会上围绕那条大名鼎鼎的“Riemann--Hilbert对应”的大“抢滩”——在那里,Deligne,别名Riemann--Hilbert,以这条“对应”之父的面目出现——原文如此;而Verdier(连同他那被其慷慨朋友大量引用的天赐“第0状态”),以导出范畴与两千年风格同调代数之父的面目出现——我这个人则绝口不提,更不必提Mebkhout\fnote{60}{(**)}{可与“同伙”与“不义——或归来的意义”两注(n\textsuperscript{o}~63$'''$与75)中的评语对读。}。

\dmark{}\origpage{323}在这些事件的光照下,我想我懂得了这份“第0状态”出人意表的发表缘由——SGA~4½的导言(还是那位一以贯之的朋友执笔的)说它“已经无处可寻”,而当时除了(也许)Zoghman~Mebkhout,谁也不关心去“找”它\fnote{61}{(*)}{不过,恰恰是4月底翻阅我刚收到的一份 Z.~Mebkhout著作的参考文献时,我才知道这份“第0状态”业已发表——那份上了年纪的文本的存在,我甚至已经忘了……}。于是恰恰有这么一个倒霉人,在他自己的角落里,逆着所有人;执拗地使用着这些过了时的观念,却没人弄得清他想达到什么——他终究那么执拗,以至于一个疑心开始冒头:万一哪天这个家伙端出几样够分量的东西呢,谁知道……毕竟,他斗胆引为灵感之源(在大师的著作旁边)的那个人,当年可是用这一切证明过、发现过东西的;那些东西,人们无法统统装作忘记,即便忘了它们的作者——而且大师本人,Jean-Louis~Verdier亲自,不正是靠着那条“Lefschetz--Verdier公式”起家奔向荣耀的吗?而没有这些该进纸篓的观念,那条公式他连写下来都费劲,更遑论证明……

我那影响巨大的前学生在这将近十年里(自从他摆脱了那道烦人的手续……)\textbf{押注导出范畴没有前途}、并且还将这样押注到关键时刻(那场大名鼎鼎的研讨会)的钟点,他想必认定谨慎为上(谁知道呢……),抢先于可能到来的事件——总之,来一份“全险”:发表的(当然不是那项本该终有一日构成一篇论文的大工程,而是)一份“证人文本”、一件“以备不时之需……”的证据;一份能证明他对一个\textbf{孤儿}之父权的文本——那孤儿,他曾乐于厌弃,并在等待事件来临之际,继续否认\fnote{62}{(**)}{倘若 J.L.~Verdier真有让埋葬了七年之久的导出范畴瑜伽为人所知的愿望,他本该选择发表构成其论文的那篇导言文本,而不是一份无人问津、只有放在那部瑜伽及其众多使用的底色上才见趣味的技术文本。但人们懂得:他绝无意愿在那份50页的证人文本旁边,附上其论文的那17页——里面有如今叫人难堪的断语,论的是那个绝不可提名的人的角色……}。

\notehead{81$_{1}$}{}

\dmark{}\origpage{324}所说的贡献是:
\begingroup
\renewcommand{\labelenumi}{\arabic{enumi})}
\begin{enumerate}
\item 局部紧空间语境中对偶形式主义的奠基;
\item Galois模语境中的相应形式主义(与J.~Tate合作);
\item 那条所谓Lefschetz--Verdier的\textbf{不动点公式};
\item 局部紧空间中的对偶。
\end{enumerate}
\endgroup

贡献2)与3)相对已知的东西,构成一个“意外”。最重要的贡献,依我看是3)。它的证明,从对偶形式主义(无论“离散”还是“连续”系数)轻易得出;这并不妨碍它构成我们上同调武库中“万用”公式的重要一味。这条公式的存在,是Verdier发现的;对我而言,那是一次(愉快的!)意外\fnoteplainmark{63}{(*)}{(1985年4月19日)关于这条漂亮的公式、它的角色、以及它在葬礼中的离奇坎坷,我在《收获与播种》第四部分的“真正的数学……”“……与‘无意义’”“把戏与创造”三注(n\textsuperscript{o}~169$_{5}$、169$_{6}$、169$'_{6}$)里还要谈到。}。

局部紧空间语境中的对偶形式主义,就其实质而言,是对我在概形平展上同调语境里所做之事的“势在必行的”移植(而且没有那种一切尚待去做的局面所固有的困难)。不过它带来一个有趣的新想法:把函子$f^!$(不先对$f$作可光滑化处理)直接构造为$Rf_!$的右伴随,外带一个存在性定理。这一办法后来被Deligne在平展上同调里采用,使他不借可光滑化假设,在那个框架里定义$f^!$。

这些评语,我想,足以说明:到1967年,Verdier已证明了他做原创数学工作的能力;这当然是为他赊账的决定性因素。

\notehead{81$_{2}$}{}

作为另一个例子,我指出:在局部紧空间的语境里,依“六种运算”与导出范畴那套“万用”形式主义的精神,对偶形式主义的详尽发展——Verdier在Bourbaki研讨班上的报告,可算它的胚芽。即便只就拓扑\textbf{流形}的语境而言,据我所知,Poincaré对偶形式主义至今也没有一部令人满意的参考文本。

\dmark{}\origpage{325}(6月5日)还有两个方向,我遗憾地看到,Verdier不曾认为值得把工作做到底——那工作他起手的势头已足够强,本可\textbf{收取其信誉}(我是指:在离散系数与局部紧拓扑空间的语境里,开启一套对偶形式主义),尽管那些本质观念并非出自他,而且他(一如对导出范畴)全无心肝去做一项任务的\textbf{仆人}、把一套完整的形式主义交到使用者手里(一如我在SGA~4、SGA~5、SGA~7三场研讨班里竭力所做的)。我所预见并建议他发展的对偶纲领,置于一般拓扑空间(不必局部紧)的框架、以及其间“分离的”、局部“可光滑化”的态射(即局部地,源嵌入$Y\times\mathbb{R}^n$,$Y$为靶空间)之中。这显然是同\textbf{任意}概形的平展上同调框架的类比所明摆着提示的。Verdier看得出:在局部紧空间的框架里,态射局部可光滑化的假设是多余的(这曾是一件意外)。但这并不妨碍:局部紧空间的语境(从而排除并非局部紧的“参数空间”)分明碍手碍脚。一个更令人满意的语境,应当同时涵盖Verdier所选的情形和我预见的情形:拓扑空间(乃至topos?)多少任意,而态射$f:X\to Y$须满足两条限制:1)分离;2)“局部可紧化”,即$X$局部嵌入$Y\times K$,$K$紧。

在这个语境里,一个“被允许的”态射的纤维,是任意的局部紧空间。下一步,则是允许$X$与$Y$不是拓扑空间,而是“拓扑多重簇”(即“局部如同一个拓扑空间”的topos),乃至任意的topos,同时对态射作合宜的限制(有待明确),使纤维成为\textbf{局部紧的多重簇},必要时附加补充条件(也许近于佐竹的$G$-流形的观点),比如(而且在万不得已时!)局部地取$(X,G)$之形,$X$为带一个\textbf{有限}算子群$G$的紧空间。

据我所知,连“寻常的”\dmark{}\origpage{326}Poincaré对偶,也还不曾在紧光滑拓扑多重簇(光滑:即局部如同一个拓扑流形)的情形里得到发展。有限群的分类空间的情形似乎表明:人们几乎只能指望在模去挠性的意义下得到一条(绝对整体的)对偶定理;更确切地说,须取一个$\mathbb{Q}$-代数作为系数环。除这一限制外,倘若Poincaré对偶(“六种运算”风格)在这个语境里照原样成立,我不会惊讶。从来无人看过它一眼(除了那些不可救药的微分几何学家——他们装模作样地看叶状结构的“叶空间”的上同调),这并不奇怪:看看我的上同调学家学生们——Deligne与Verdier领衔——对多重簇观念本身所设置的普遍抵制。

老实说,缺的是这样一场奠基的反思:在任意topos及其上的“离散”系数层的语境里,描述(倘能做到)topos态射的“真性”“光滑性”“局部真性”“分离性”诸观念,从而提炼出一个“被允许的”topos态射$f:X\to Y$的观念——对其,$Rf_!$与$Lf^!$两种运算(前者为后者的左伴随)各有意义,从而取得六种运算形式主义的惯常性质。这里,topos被视为不带环层的,或也许带着环层(必要时假定为常值或局部常值);先假定(至少起初)带环层的topos态射$f:(X,\mathcal{A})\to(Y,\mathcal{B})$满足$f^{-1}(\mathcal{B})\to\mathcal{A}$\textbf{是一个同构}(注81$_{3}$)。上面的反思提示:当我们把自己限于特征零的系数环层(即$\mathbb{Q}$-代数)时(译按:特征零本身不足以保证环为$\mathbb{Q}$-代数,例如$\mathbb{Z}$),“被允许的态射”的观念可以明显放宽,以涵盖这样的“纤维”:比如多重簇(拓扑的或概形的),而非寻常的“空间”(拓扑的或概形的)。

这个方向的最初尝试(除我处理过、随后Verdier依同一模型处理过的情形之外),归Tate与Verdier:在离散群或投射有限群的语境里。这一初步尝试的记忆,曾鼓励我去年继续这个方向上的反思:在小范畴(它们推广离散群、充当同伦模型)的语境里。虽然没有走远,这场反思仍足以使我相信:必须存在一套完整的六种运算形式主义,在小范畴的范畴(Cat)的语境里。(对此可参见\dmark{}\origpage{327}《追寻栈》,第VII章,第136、137节。)这样一种理论在(Cat)、乃至在Pro(Cat)里的发展,一如这类理论在拓扑或概形的空间与多重簇的语境里的发展,对我而言,其主要趣味在于:它是朝着在一般topos的语境里更好地理解“离散对偶”迈出的一步。

去年Illusie向我透露:他在半单纯空间(或概形)的情形里,曾同对偶理论的种种疑难苦斗。那听上去分明还是同一回事——在一个个案里察觉到一套六种运算形式主义的存在,并理解它。但似乎,一场奠基反思的前景本身,就有本事使我的前学生们——至少我的上同调学家学生们——人人浑身冰凉。我当年同他们费过力气,是怀着这样的信念:他们不会(在概念工作的层面上)恰好在与我同行的那个地点戛然而止;每当一个新的局面显出他们与其伙伴同我做的工作不够用时,不会只是搓手发愁。人们所做的概念工作,从长远看\textbf{总是}不够用的;数学的进步,恰恰在于重新拾起它、超越它——除此之外别无他途。1955到1970年间,年复一年,我发现前些年所做的已不敷需要,便立刻重新伏案——至少当别人(比如Mike~Artin,以他自己意义下的“代数簇”观点)还没有抢先把这件事做掉的时候。但似乎,我的学生们把我给他们做的榜样也一并埋葬了,连同我这个人、我的著作。

\notehead{81$_{3}$}{}

我记得,在(比如说)平展上同调的六种运算形式主义里,充当系数的环层局部常值的假设是多余的——本质的假设是:它们是素于各剩余特征的挠层,\textbf{并且}$f^{-1}(\mathcal{B})\to\mathcal{A}$是一个同构。一旦放弃后一假设,人们便进入一门(据我所知从未明言的)理论:它把“空间离散”对偶与“凝聚”对偶(关涉系数环层及其同态)混合起来。于是,人们设想:在概形(或更一般的topos)$X$、$Y$上,把系数环层$\mathcal{A}$、$\mathcal{B}$换成$X$、$Y$上的相对概形(不必仿射)$X'$、$Y'$,把带环层的topos态射\dmark{}\origpage{328}$(X,\mathcal{A})\to(Y,\mathcal{B})$换成如下类型的交换图表
\[
\begin{tikzcd}[row sep=2em, column sep=3em]
X' \arrow[r] \arrow[d] & X \arrow[d] \\
Y' \arrow[r] & Y
\end{tikzcd}
\]
并带一套这一类型语境中的“六种运算”形式主义。当$X$、$Y$等为点的topos时,人们应当找回惯常的凝聚对偶。

\subsection{好的参考文献}

\notehead{82}{}

(5月8日)说的是 J.L.~Verdier的文章“与一个闭链相配的同调类”,1976年刊于Astérisque~n\textsuperscript{o}~36(SMF),第101--151页。从某种意义上说,这篇相当不可思议的文章(尽管如今已没有什么应当再使我惊讶……)与Deligne等人的“反常文章”互为表里。除一点保留外,它实际上等于花五十页,在一个略有不同的语境里,\textbf{照抄}我在十年或十五年前纵横发展过的观念、构造与论证——术语、记号,应有尽有,一字不差!我几乎以为自己回到了1965/66年间SGA~5研讨班的一场:那些东西在那里被讲明过(显然,详尽到与会者吃不消\fnote{64}{(*)}{这一方向上的评语,见注n\textsuperscript{o}~68、68$'$“信号”与“颠倒”,我在其中检视这场研讨班成文工作的离奇坎坷,以及这些坎坷与Deligne“SGA~4½之役”的关系。随后的反思,向我显出这些坎坷、以及母研讨班经Verdier与Deligne合谋肢解的又一个未曾预料的侧面。两人的出版物为这场肢解揭幕,分别刊于1976与1977年——它们构成给Illusie的“绿灯”,让他去筹备(十一年之后……)SGA~5的出版(那一卷,如Deligne在SGA~4½里所说,“可以看作一系列枝蔓,其中一些还非常有趣”)。})整整一年。至少在那场研讨班之后,对所有多少“在局中”的人而言\fnote{65}{(**)}{关于我回到这份“仓促”印象的一场反思,见“沉默”一注(n\textsuperscript{o}~84)。},这一切都属于“众所周知”的领域。Verdier当然到场了,一如Deligne(唯一一个从不掉队的人——尽管那是他头一回把脚迈\dmark{}\origpage{329}进我的研讨班\fnote{66}{(*)}{这场研讨班的那一年(我想),正是我结识Deligne的一年;他当时该有十九岁。他很快“入了局”,甚至揽下成文我前一年平展对偶报告的差事(那些报告,他想必经由我的讲解与我的笔记而知),以及关于与一个闭链相配的上同调类的那场报告——后文在引过的“颠倒”一注(n\textsuperscript{o}~68$'$)里谈到,在本注里也还要谈到一点。以他那般手段、对主题的那份完全驾驭,他竟等了十一年才做成文,然后不经告知,把它收进其SGA~4½——如今回顾,这向我表明:早在1966年(而不仅是我曾以为的1968年——见“驱逐”一注,n\textsuperscript{o}~63)——即我们相识的头一年,我的朋友对我的关系中,便已存在一种深层的暧昧,而且从那一刻起就以全然清晰的方式表达着;而我一直回避去面对它,直到今天!}——总得有人这样……)。倒是有一件事:到1976年,那场大名鼎鼎研讨班的“成文(原文如此)”,由那些早已厌烦的“志愿者(原文如此)”拖着,已经拖了十年——我现在明白:其中一位“志愿者”到底还是按他自己的方式揽下了“成文”的差事,早在SGA~5于1977年出版之前!看来这位倒霉研讨班的坎坷,方便的还不止Deligne一人——他以自己的方式,从一盘散沙的局面里渔利。不过在那个时刻,Deligne倒还小心:一边把SGA~5的一场关键报告拆下来、名正言顺地并入其SGA~4½,一边到底还在其成文(关于与一个闭链相配的上同调类)里提一句“依Grothendieck的一次报告”。(诚然,他从中得到一份补偿:得以据此把我列为他的“合作者”!——见“颠倒”一注,n\textsuperscript{o}~68$'$。)

回到与一个闭链相配的\textbf{同调}类(别混淆!)——按标题,它是Verdier文章的对象——这个形式主义,我在口头研讨班上,分几次报告,以穷尽的细节发展过;而且面对的听众分明在求饶(照例只有Deligne一人,始终神采奕奕、毫无倦容……)。那是那一年我在平展框架的对偶形式主义上发展过的无数“长练习”之一;我感到需要完全驾驭一切在我看来必须彻底领会的要点。这里的趣味在于:拥有一个在不必正则的环绕概形上成立的形式主义——到正则情形的\textbf{上同调}类的过渡,以及与我那套使用具支撑上同调的老构造的联系(它立即给出与上积的相容性),都是直接的。我还发现:\dmark{}\origpage{330}研讨班的这一部分,属于未收入已刊版本的那一堆——想必Illusie(出一版拿得出手的(哼)本子的全部准备工作,到头来都压在他身上)该庆幸Verdier揽下了这件事,依样画葫芦(此处即:一字不改!)。

按照如今已经功德圆满的套话,“几乎无须说”:我的名字既不在正文里,也不在参考文献里(除了那条永恒的SGA~4引证——它含蓄地带着我,尽管人们早晚得找到替代……)。没有任何暗示提到一场与SGA~5这个缩写相应的“代数几何研讨班”——作者想必听说过它——尽管我分明记得见过他,殷勤地、懂事地做着笔记(像所有人一样;当然,除了Deligne……)。

况且,我说我的名字不见于正文,只夸张了一丁点:它有一次露面,神秘而碑文式,在第38页,第3.5节“基本的上同调类,相交”(人们终于抵达问题的死结了!)。那个出处是一句谜语式的话,其意我承认参不透:“系统地使用权复形(???又是这些该死的权!)的想法归Grothendieck,并经Deligne整理成型”——对这些神秘的“权复形”——据说我出过想法、而我此刻是头一回听说——再无任何解释。此后全书再未提及(此前37页里也从未提及)。谁懂谁懂!至于该节的内容,它一字不差地抄自十年前举行的那场SGA~5研讨班(而那时这一构造已有五六年的历史,见注n\textsuperscript{o}~68$'$);那场研讨班,他自然绝口不引。对Deligne的引证(他竟“整理成型”一条想法——当我的朋友还是中学生时,那条想法就已然现成!)是一朵“花”:这一想法想必之所以来到作者心头,是因为年轻的新来者Deligne确实揽下了成文我那场报告的差事(而且一搁十一年,为了人们知道的那些收益,见前引注)。这朵“花”,属于这对形影不离的朋友之间礼尚往来的交换。

不过,文章里有一个(想必)新颖而极有趣的结果(定理3.3.1,第9页):解析可构造离散层经由一个解析真态射的高阶直像的稳定性。Verdier约十五年前从我嘴里学得各种语境中的可构造性观念,以及\dmark{}\origpage{331}那条稳定性猜想——我在五十年代末提出过它(并对凡愿听的人谈起),那时我还没有结识他的荣幸。读这篇文章,一个不知情的读者(但这样的人已开始稀少……我恐怕又在重复自己)绝不会想到:作者并不是正在趁热端上他刚刚才发现的观念与陈述。他无须说那是他——不言自明。这就是那套大名鼎鼎的“暂停”做派,分明已经蔚然成风。

除了这个细节(依我之见,它合乎行业的新规范)之外,这篇文章想必总有十页上下(五十页里),围绕这个有趣的结果,呈献作者本人的工作。比例放到相应的大小,在Verdier身上、一如在Deligne身上,最打动我的,是他完全有能力做出漂亮的数学。即使在这篇令人痛心的文章里,前引定理也透出这一点。但他(一如其朋友)把自己保持在掘墓人的心境里,便与他显赫的朋友一样,只靠其手段中微不足道的一部分运转。一个(曾令我瞠目的)表面平庸的征兆,见于一位分明给出过机巧与嗅觉之证明的数学家:对他的“学生(原文如此)”Mebkhout的工作的分量,全无直觉——他乐于以自己的伟大俯视他,却从不曾自己做出深度与原创性堪比的工作\fnote{67}{(*)}{同一个令人瞠目的嗅觉迟钝,在这同一场合也见于Deligne:他仿佛到1980年才“感到风向”(Mebkhout那些想法的重要性),而Mebkhout在这一方向上自1974年起就工作着。我的朋友身上,自然嗅觉被自满所堵塞,我已有不止一次机会观察到;尤其自1977年(或1978年)起——那一年似乎构成第一个“转折”(对此可参见“两个转折”与“葬礼”两注,n\textsuperscript{o}~66、70)。}。这不是说他也许没有能力做到Mebkhout或我的地步。而是他从不曾给自己留一点机会去做大事——也就是说,放开缰绳,让一股热情牵引,而不是把数学及其天赋当作\textbf{工具},用以炫目、支配或碾压。直到此刻,他一直满足于照原样拾取那些现成而丰产的观念与观点。确实,他似乎已全然丧失了何为\textbf{数学创造}的感觉。

然而我记得,他与我一起工作的时候,这份感觉还在。没有任何外在于他的东西,妨碍这份感觉\dmark{}\origpage{332}重新浮出。一如其朋友——在那人身上,我常常感到那件精微而活泼之物的同样的隐没,被同样的自负所堵塞。

这篇50页的不可思议的文章,刊于一份有分量的杂志,为我照亮了“注记——或新伦理”那桩事故(第33节)新的一面:一篇\textbf{几页}的科学院短文——概括一项扎实而\textbf{原创}的工作,题材重要(依我浅见),是一位天资极高的青年数学家\textbf{两年工作}的果实——竟被两位权威以“毫无趣味”为由拒稿\fnote{68}{(*)}{细节见“棺材4——或没有鲜花也没有花冠的topos”一注,n\textsuperscript{o}~96。}。其中一位权威,不是别人,正是Pierre~Deligne——同一个Deligne,却肯一字不差地、亲自照抄我的一个学生那份不起眼的博士论文(他还恪尽职守地引用它)。(这份复制件,经一项显赫的署名拔高,成了那部同样显赫的丛书里“值得纪念的文集”LN~900中最厚重的一篇!对此可参见注52、67的末尾。)

分明,这幅“风俗画”日甚一日地丰满起来;而我为此甚至无须走出隐退、走街串巷去混迹“上流社会”。在几篇挑选得当的“大文本”里翻阅的零零星星几个钟头,已足以使我大开眼界……

\subsection{玩笑——或“权复形”}

\notehead{83}{}

(5月8--9日)我重新想到了那个“权复形”——它出现于Verdier值得纪念的文章里那条“出处——暂停”\fnote{69}{(**)}{见前一条注“好的参考文献”。}之中——一条俨然滑稽戏、纯粹无意义的出处。这条离奇的出处刚一入眼,一个联想便来到心头,并且一直在脑际萦回。我面对看似离奇、似乎挑战一切合理解释的东西——而其意味却又清清楚楚、并且被清晰地知觉到,只是在一个不同于常规逻辑的层面上——这已远不是头一回。常规逻辑,是我一生中几乎全部时间在意识层面运转的唯一逻辑——结果是:我不断地被“离奇”的、不可理解的事件甩在身后——它们那不可还原的离奇,令人焦虑!我的生活变化很大,自那一刻起(不到十年之前):我开始在我官能的一个更宽广的音区上生活。我懂得了:一切离奇、一切所谓的“无意义”都有\textbf{意味}——单单\textbf{知道}这一点\dmark{}\origpage{333},从此对无意义背后的意味保持好奇,就常常向我敞开其明摆着的意义。

在这“权复形”的无意义里,我感到一记与“反常层”的命名同类的\textbf{挑衅}\fnote{70}{(*)}{见“反常性”一注,n\textsuperscript{o}~76。}——此处的乐趣在于向自己证明:在一份有分量的杂志里、在一个自命为标准参考文献的文本里\fnote{71}{(**)}{而且,这个文本看来今天确实是一份标准参考文献——至少多年来,它一直是Zoghman的枕边书之一(他最近寄给了我)。正是在那里,他学得了可构造性观念(它在其定理里起本质作用);而且长年以来,他深信Verdier是这条对他至关重要的观念的神奇发明者。},人们\textbf{可以放胆}说一句与此相类的离奇话,而\textbf{没有一个人}会想到哪怕提出一个问题!而且我确信:这一挑衅里所含的赌注——文章发表八年以来的赌注——\textbf{赢了},直到今天:今天我是头一个向作者提出那个天真的问题的人。

当然,一桩离奇出现的时间(或地点)——在这里,恰是全文唯一一次指涉我个人的当口——绝非偶然;其形态也绝非偶然:这里借对一类观念——“权”——的影射(它与整篇文章的主题全然无关),并借对一个从未存在过的复合观念“权复形”的即兴,来实现。立时来到我心头的那个联想,很可能提供这桩离奇——在挑衅、权力示威之外——的更精确意味的钥匙。那是与另一处谜语式、同样纯礼节性的影射(但尚不带离奇的附加维度!)的联想——在注49开头所引的Deligne文章里\fnote{72}{(***)}{即“罐装的权与十二年的秘密”一注。从我们此处关心的角度对这篇Deligne文章作更周详的检视,见后文所引的“驱逐”一注(n\textsuperscript{o}~63)。}。那正是一处晦暗的影射,在一篇“权”一词严格缺席的文章里——除Serre与我,谁也无从看出——影射某些“权方面的考虑”,它们使我(文中特意写明:以一个欠一般的形式)猜出了该文的主结果。如我在更详尽的“驱逐”一注(n\textsuperscript{o}~63)里所解释的:在这份纯礼节性的影射背后,透出\textbf{藏掖}的意图——既藏我的角色,也藏那些他打算独占其收益的观念(关于“权”及其与一般上同调、尤其与Hodge上同调的关系)。这一意图,Verdier想必领会得格外透,因为他自己\dmark{}\origpage{334}就在同一个音高上“运转”(至少就他对我的关系而言;而依我看,这正是那对形影不离的朋友之间主要的黏合剂)。在两种情形里,诚实的做法本应是:开篇就明白指出主要观念、或促成文章的那个或那些问题的源头。

把这一点说清之后,这桩表面无意义的象征语言背后,我所知觉的意味便是:我可以毫无顾忌地当众陈列一桩明摆着的\textbf{无意义},同时借这个无意义、借对“权复形”那句荒谬的影射-出处,表达我真正的意图:我对Grothendieck在这项工作中的角色,无意透露分毫,一如Deligne当年对“权方面的\textbf{考虑}”的凭空影射——那条影射,对当时的读者,其意味不多一分,恰如我此刻发明的“\textbf{权}复形”,为着事情的需要、也随我的高兴!

我刚刚把昨天写的这条注誊清——其间我被Verdier的一个电话打断:那天我本想找他,正是为了当面问这个问题。我向他解释:我正在晚年勉力学一点上同调——这门学问我从来没弄懂过,他知道得很清楚——Mebkhout为供我学习,给了我一篇Verdier的旧文章;Mebkhout曾长期拿它当枕边书。我如今正尽力读它,但那里有这条谜语式的出处——他当然好意引我——可我全然不懂他指的是什么。

他非常高兴,甚至有点受宠若惊——可不是嘛,宽大的笑容从一副慈父般的欢容后面透出来——我竟然在我这把年纪,靠他那页旧纸学起上同调来了。我没有料到:当我说他分明知道我对上同调从未入门时,他会想到反驳我——那分明是久已有之的定论……至于那大名鼎鼎的“权复形”,我在听筒里再次感到他的宽大笑容(人们会说我编故事!),为终于有人(而且正是收信人本人)注意到那件被放过了这么久的事而欣喜。同时,又带着一丝窘迫——我想,那不是不知如何作答的窘迫,而是没能把一份乐趣\dmark{}\origpage{335}藏住的窘迫(一如人们对一桩略有荤味的故事所感到的乐趣……)。掉队如我,他实在无须在这一头费心!他不假思索地把话岔到Deligne身上(我没有提这个名字):Deligne在一篇文章里做过一个证明,其中还引了我——他记不太清在哪里;总之那里谈的是权,可不是嘛——当然他有点忘了——但那确实不是算术的权;这一点我全对,不是一回事……

声调欢容而不容置辩,他示意已经为我花了不少时间——略显匆忙,却不失那副从容、略带庇护的口吻。我为用这样一个有点愚蠢的问题打扰他表示歉意,并为他的解释道谢。我的歉意是真诚的,我的道谢也是——他确实把我想知道的一切都教给了我\fnote{73}{(*)}{即便摆出那副掉队的神气,我也并不真觉得自己在演戏(我没有这份天分);那全然自然——说实话,对这套我已近十五年不曾操持的玩意,我确实有点掉队!但我相信,即便老糊涂到该上灵车,我仍闻得出空核桃与满核桃的区别……}。

\section{IX.~我的学生们}

\subsection{沉默}

\notehead{84}{}

\dmark{}(5月9日)昨天我也许写得急了点:我在“好的参考文献”(见注82)里写道,作者兼前学生毫无顾忌地照抄的内容,“属于多少‘在局中’的人‘众所周知’的领域”。我曾试着为自己弄清,这些“多少在局中的人”究竟是谁——结论是:\textbf{不多不少,正是1965/66年间SGA~5那场研讨班的亲爱的听众们}——而且,如我已有机会说过的,这些听众多半或多或少掉队;而从这场研讨班的成文工作落在一群志愿者手里、我不肯觉察其信念不足的坎坷来判断,常常是“多”于“少”(当然,照例是Deligne例外)。事实上,只要SGA~5尚未成文出版、以便人们经由阅读而“入局”,就绝不会有别的“局中人”!这场研讨班实际上(天意弄人!)是在我最亲爱的两位学生与战友的两大“值得纪念的出版物”之后才出版的:其一,Verdier1976年的那篇文章——其中对由他执笔首次发表的观念之来源只字不提;那些观念,已被大量谈过\fnote{74}{(*)}{尤其见注n\textsuperscript{o}~67、67$'$、58、68$'$。};其二,Deligne随SGA~4½\origpage{338}。此后,人们才好心地邀请Illusie,来张罗其余部分的出版!

我已记不太清那场研讨班的与会者——比如Artin在不在场。我想,无论如何,我第一期的学生多半都在——当然Sinh~夫人与Saavedra除外(当时我尚未结识他们),或许还有Hakim夫人。此外还有Bucur(已故)、Houzel、Ferrand——我不算Serre——他对上同调的重型装备素来没有胃口,只是远远地、谨慎地偶尔露一脚。虽然除Deligne外,也许没有人十分清楚这一切将通向哪里,但我想,总该有十来位听众(不算很投入)至少跟得足够紧,可以被算作“在局中”\dmark{}。从昨天起在我脑际萦回的念头是:在所有这些“在局中”的人中间——他们因此俨然以上同调行家自居(即便不都像Illusie与Berthelot那样以“上同调”论文成名——那些论文分明够分量)——即便把Verdier与Deligne除外,想必仍有相当一批人,亲手捧读过Verdier这篇文章!Verdier身上某种神气,使我深信:从没有人向他暗示过,或许有什么地方不对劲。我也清楚:同样从没有人向我提过这件事——我是5月2日,即恰好一周前的今天,经由Mebkhout才知道这篇文章的存在的;而他对那场巧取,当然已知情多年。

这给十天前“全体一致的赞同”(为埋葬我寒微的身量而达成的)那句兴高采烈的确认,赋予一个十分具体的意味(注74)!这份赞同,把我“1970年以前”的学生中的许多人(乃至全体)——也就是今天在数学界引领风骚的许多人——都包括在内;它还把(或曾把)我的朋友Zoghman本人包括在内:他被当作上流社会的灰姑娘,逆着所有人死死攀住某种“对我的著作的忠诚”(借用他自己的说法\fnote{75}{(*)}{(6月7日)最近来访时通读关于葬礼的各注,Zoghman向我指出:他那个“对我的著作的忠诚”的说法,并不真正表达他的想法。他心里想的,毋宁是对自己的判断力与数学本能的信任——它们向他表明,我的著作为他带来了他所需要的某些想法。因此,那是一种对\textbf{他自己}的忠诚;而要做真正革新的工作,这确实是本质的东西。}),并有胆量与执拗,时不时以此自命,后果则有目共睹。谁能懂,就去懂吧!

总之,我暗示那份有分量的杂志刊登了一篇照抄“众所周知”的空壳文章,是我错了。作者当着(虽非所有、却)众多见证者的面照抄的东西,既未发表、也非“众所周知”(唯有一个闭链在凝聚框架里的上同调类除外——那个我久已发表);而那些观念,我再没脸去贬低:我并不曾认为,花一年工夫、在众多听众面前在一个研讨班里发展这些观念与其他观念,是虚掷光阴。想必Verdier的文章,是我所发展的观念与技术的一小部分的“摘要”,有用而做得好:恰好使它们得以进入“众所周知”的领域,成为把上同调(或同调)用于那些多少配得上“簇”之名的对象者的日课。\dmark{}\origpage{339}从这个观点看,Verdier做了有用的事\fnote{76}{(*)}{他做了,诚然,以原初SGA~5研讨班的“拆解”为代价——那场拆解,他与Deligne是主要的操办者与“受益者”。\par (6月7日)三天后(5月12日)的反思(见“屠杀”一注,n\textsuperscript{o}~87)显明:与这场拆解相比,所发生的更是一场“屠杀”,而Illusie的关联比Verdier还要直接——即便他并非“受益者”,而是代人行事。},我最终并无不满之处。然而,根据我从我的前学生、今天仍是朋友的人那里——在电话里——以及从他这个人身上我得以感到的许多别的事情(其中最“重”的、或至少最“轰动”的,是反常层会议之欺瞒)所感到的——我分明感到\textbf{有什么地方不对劲}。那场值得纪念的研讨会,从数学上说,想必在许多方面都极其辉煌。而“不对劲”之处,处在全然另一个层面。我可以试着以言词勾勒它,但我分明感到,那没有多大意义。在那个不感到这场研讨会不对劲的人那里——也在许多别的研讨会里,毫无欺瞒——当我做了这场“勾勒”的努力、哪怕完全成功地做到令我满意——他也不会多感到一丁点……

对我仍然悬而未决的问题是:这个“征兆”——它表现为一件如今想必相对寻常的琐事(一位作者把他人的未发表观念呈献为自己的)——这个征兆,是世风普遍败坏的征兆,因而只是今天数学界某种“时代精神”的典型征兆呢,还是它更想就我个人、就那个曾经是我、如今又回到我身上的人(经由昔日学生们对我的态度)教给我点什么。

两种可能的意义绝不互相排斥。倘若没有某种世风的状况怂恿,我的前学生们对我的关系,本无从循这条路径表达。况且,早在这个“征兆”之前,我已看到别的许多征兆;在我看来,它们在“风俗画”的层面上更雄辩。而在这个征兆里打动我的,是它区别于所有其他征兆的特殊之处:它似乎\textbf{同时牵连我的大部分昔日学生}。

这样的情势绝非偶然。把它干脆记在“世风败坏”(千真万确的败坏)的账上,便是回避其更个人的意味——那意味牵连我,一如牵连我的每一位前学生。倘若我说“每一位”,似乎超出了这个征兆的实际幅度,那是字斟句酌。因为这个征兆适时地提醒我:我的某一位昔日学生,竟不曾至少面对过这类局面,几乎是不可设想的。多年来,我感到一种关于我的某种“风向”,吹拂在我离开的数学家世界里(如今我分明看清这风的来源与缘由,我想)。不可能他们中的任何一位,从不曾感到这阵风的吹拂——无论是经由像这篇掘墓文章的发表这样的“事故”,还是经由任何别的场合。无论当事人愿意与否,这样一场遭逢,必然向他提出(或重新提出)他与我的关系的问题——我曾教他手艺。而我确认的征兆,在把我引到这件事的那一个之外,是:\textbf{那些曾是我的学生的人里,没有一个就此事给过我任何回音}\fnote{77}{(*)}{(5月31日)有意思的是,唯一一个曾向我暗示存在一场下葬的人,是一位非洲朋友——十年前,他跟我做过一篇第三阶段博士论文(即“1970年以后的学生”,身份寒微);我一直与他保持友好的交往。他有所暗示的那封信,应当是两三年前;在那个时刻,这毫无使我惊讶之处。我当时没有就他的印象问细节;他只是最近,才重新谈起它们。}。这是一桩“巧合”,其意味仍向我滑脱——但它不可能没有意味(注84$_{1}$)。

天开始破晓——我感到是停笔的时候了。我拿不准:在《收获与播种》里,此刻此地继续深挖这桩触目巧合的意味,是否合宜。这也许是留给别的明天的收成——倘若我今夜的反思,在我昔日学生中的某一位那里遇到回音。($\Rightarrow$注85)

\notehead{84$_{1}$}{}

(5月16日)我的昔日学生们之间的这份完美的协调——在对我的全然沉默之中——与其他征兆指向同一方向。其一是迎接“外来者”情节(见第24节)的全然同样的沉默——关于那份沉默,我在注n\textsuperscript{o}~23v里已多少追问过。另一方面,除Berthelot(他给我寄过许多抽印本)、Deligne(寄过四份,在其约五十项出版物里)与Illusie的一份之外,我的任何一位前学生的抽印本,我都没有收到过。关于他们对我之关系的两价性,这说明了许多。寄送抽印本——即便人们怀疑我是否会在自己的工作中使用它们\fnote{78}{(*)}{(5月31日)直到1976年,这甚至想必可以算是被排除的:七十年代初,我曾相当明确地说过,我不打算再重操数学。1976年在IHES的那场关于带除幂的De~Rham复形的报告,则相当清楚地表明我仍在关注数学。}——本会是最明显不过的方式,向那位教了他们手艺的人表明:这门手艺在他们手里并未僵死,而是活生生、在运转。但同样真确的是:至少对其中某些人,他们的出版物也见证了他们参与一场心照不宣的下葬;对那场下葬,最好还是别告知被预葬者——手艺也罢,非手艺也罢……相反,我却收到过好几位从事晶体上同调的作者的许多抽印本\fnote{79}{(**)}{(5月31日)说的是我不认识的年轻作者;我料想他们效法了Berthelot——在他们眼里,他想必是长辈。这里稍显奇怪的是:至少两年来(自1982年9月6--10日的吕米尼研讨会起),Berthelot正亲力亲为地埋葬我(对此可参见“共同继承人……”一注(n\textsuperscript{o}~91)5月22日的脚注)——莫非他对我的关系近来起了转折?我不记得收到过那篇晶体上同调及其盟友的综述文章的抽印本——文中对我的名字讳莫如深;他肯定小心没寄给我!},甚至收到过许多我只闻其名的分析家同行们的抽印本——当他们的工作重提(有时解决)我在三十年前或更早提出的问题之时;尽管明摆着我不会回到那个我已告别的题目,而且从“实用”的角度看,那些抽印本是白寄了。但这些同行想必感到了某种我的学生们不屑于感到的东西。——当然,在六十年代,我的学生们在我的一切出版物里都是头等的受赠者:无论我的文章,还是EGA与SGA两大系列;他们每一位(除Sinh~夫人,或许还有Saavedra)都应当拥有我1955至1970年间出版的全部著作(想必近一万页)。

诚然,我的前学生们有很好的同伴:我在“上流”数学界的任何一位旧友——包括那些著作与我的著作血肉相连的人、或在六十年代我的工作纲领的展开中扮演过角色的人——都不曾认为值得在我离开共同环境之后继续给我寄抽印本\fnote{80}{(***)}{(5月31日)当然,促使他们给我寄的心理缘由,远不及对我的学生们那样强;但,人们也许天真地想,总该比我的分析家同行们、甚至比我收到过抽印本的众多代数几何学家——我并不认识或只略识面——强些。分明,自我离开共同环境之后,“曾是朋友”这件事,在我数学界的旧友们身上造成或强化了我已有机会观察到的那种拒斥的自动机制。(关于这些态度——《收获与播种》里此处彼处顺带影射过——见5月24日的“掘墓人——或全体教众”一注,n\textsuperscript{o}~97.)}。最近,在我寄出《纲领概要》的十五二十位旧友(包括几位学生)之中——那份概要(除别的以外)向他们宣告:中断十四年之后,紧张的研究活动重新开始,且研究主题与我们当年共同推进的血脉相连——只有两位(Malgrange与Demazure)费心寄来几行致谢。我所收到的较周详(而且热忱)的些许回音,来自我近来结识的青年数学家,以及我的老朋友Nico~Kuiper——尽管他对我所做的这一类事情全然外行。他是经第三者得知这份文本的,并为我那不期然的“重返”由衷高兴\fnote{81}{(*)}{(5月31日)这几乎是唯一的、来自我的一位旧友(或前学生)的、对“重返”表示首肯的回音。这当然毫无奇怪之处:亡者现身,总归不合时宜地打断葬礼的顺畅进行……\par (6月17日)然而最近,我荣幸地收到Mumford一封热忱的信:他自称对《概要》里勾勒的想法感到“激动不已”(thrilled)且“非常兴奋”(very excited),并向我确认:我为其Teichmüller塔的组合描述所需要的那个关键技术结果,确实已得到证明。这是自1978年以来,头一回有我的一位旧朋友对我的“阿那贝尔式”想法产生共鸣——那想法的非凡分量(堪比Motif瑜伽),在我看来,自始便是明摆着的。\par (1985年3月28日)在这些文字写成之后,我还收到 I.M.~Gelfand一封非常热忱的信(注1984年9月3日),回复《概要》。}。

\subsection{团结}

\notehead{85}{}

(5月11日)那场倒霉的SGA~5研讨班的事,继续在我脑际萦回。“好的参考文献”\fnote{82}{(**)}{见注n\textsuperscript{o}~82。}分明给这段往事投下一束新的光,并同时赋予那场光辉的“SGA~4½之役”一个新的意味。

越想,SGA~5的事在我眼里分量越\textbf{重}。我的第一印象——那是仅仅几个星期前我刚“下船”时(见注68、68$'$)——是:65/66年那场研讨班的穷听众中间的一盘散沙,被我的朋友皮埃尔以其大名鼎鼎的战役加以利用;而在那场战役里,别人都无关。至于SGA~5的诸种不幸,既不怪他也不怪任何人,而宁可怪我自己:我没能——唉!——点燃我的志愿者-执笔者听众的热情,也没能替他们做那项他们执拗着不做、嘴上却说着“马上就动手”\dmark{}\origpage{343}的工作。然而近几天显明:毕竟有一位,其热情十年后才苏醒,把他在那里乐于取用的东西发表出来(对研讨班不置一词),从而在别的“志愿者”们还没下决心动手的时刻,为自己立了一份“好的参考文献”。

自昨天起,对我越来越清楚的,是:直接卷入这场研讨班之藏掖的,不止两个“恶人”,而是\textbf{我的每一位“上同调学家”学生}。如无差错,他们每一位都听过这场研讨班——即(按我的“上同调学家”学生们出现的先后为序):Verdier、Berthelot、Illusie、Deligne、Jouanolou。(Jean~Giraud不算;他在与SGA~5或其前驱SGA~4的主题相当不同的音区上工作。)

这场研讨班,是我首先\textbf{为了我的学生们的收益}而做的——即便有时他们哀求饶命——\textbf{我认为它不是一堆狗屎}。那一年里,他们每一位都学到了一大包“使用上同调的数学家”这门手艺!我在平展框架里、以周详得多的方式,重做我最初在凝聚框架里发展的那些观念——这些东西,除了这场专为他们的收益而办的研讨班,他们在别处无处可寻:在我之前,没有人费心做过;除我之外,没有人甚至感到有什么要做、为什么做。(照例,Deligne除外:他在这个研讨班本身里经月习得,比别人悟性快。)正是听过这场研讨班(及其前一场)并在自己身上多少消化了它,而没有别的,才使他们从此对对偶形式主义“在局中”——而且他们是\textbf{仅有的}在局者。这份\textbf{特权},依我看,为他们造成一份\textbf{义务}:即留意使这份特权不致滞留在他们独家手里;使那些从我嘴里学到、并在他们直到今天的全部后续工作中构成不可或缺行装的东西,在合理而惯常的期限内——至多一年,宽限两年——交到人人手中。

\dmark{}\origpage{344}人们会不无道理地说:首先轮到我来督察这件事。但我当年是真心接受学生及其他听众自荐承担成文工作的(对肯认真投入的人,成文只会带给他们最大的益处)——并非为了享受一项本属于我的工作由别人代劳、而我袖手旁观的好处。我继续(有Dieudonné的协助与其他人——其中1966/67年间还有Berthelot与Illusie)发展另一些在我看来同样紧迫的奠基文本,而那些,当年没有别人能代我做、或在没有我的协助下做\fnote{83}{(*)}{1960到1970年间,我必须以平均每年约一千页文本(EGA、SGA、文章)的节奏工作,其中几乎全部都成了通行参考文献(这一点,我在写它们时、或鼓励某位合作者在我的协助下写它们时,就已经很清楚)。}。这些文本本身,也成了不可或缺的参考文献——包括对我的“上同调学家学生们”:他们同所有人一样,很高兴在需要时找到现成的它们。

凭着经由与我的接触、经由听我的或参与我的研讨班而获得的对上同调观念与技术的驾驭,合力成文这场研讨班,与对“大名鼎鼎的数学共同体”所提供的服务相比——也许还有,稍后,对我的一份忠诚义务相比——是一件微不足道的小事。我说过,对我这个有手感的人,写定整场研讨班是几个月的活。五人分工,加上他们在那些年里各自积累的成文经验,手头又有我的详细手写笔记,每人的投入至多一两个月。他们比别的执笔者(比如Bucur——他分明巴不得把一件明显超出他能力的差事,交给更年轻、动机更直接的手)装备好得多。

只要我还在场(也就是其后的三年里),我理解“靠我”的反射可能起作用:理应由我统筹一切、跟“志愿者们”周旋。很可能,倘若我当时要求他们每人限期做两三场报告,我同样承担,以便了结此事,他们不至于推却。\dmark{}\origpage{345}自我退出数学世界的那一刻起,局面全然改观。他们发现自己成了\textbf{某份遗产的独家保管人}:既是默示的(没有遗嘱),又非常具体。诚然,从实际看,我的离去等于一场\textbf{消失}——我确实“亡故”了,意思是:除他们之外,无人知道这份遗产,无人能用它,无人(无论为善还是为恶……)关心它的命运。

倘若在我离去后的七年里,这份遗产始终隐没(除1976年的“好的参考文献”!),那是因为\textbf{我的学生们不曾在整段时光里让它变成公共的}。按比例说,这局面在我看来,与“Motif瑜伽”的局面相去不远:那部瑜伽,除我之外,唯有Deligne洞悉入微,而他判定把它留在自己怀里、独享其益。若说乍看有区别,那是:此处的“受益人”是五个而非一个;而且一个人所藏之物的深度,与五人共同藏起之物的深度,不可同日而语。

每个人的深层动机,我当然不知道——即便是Deligne的情形,我的领会仍模糊,而且想必仍将模糊。但在“实际”的层面上,Deligne的棋(连同“SGA~4½之役”以及其余)再清楚不过。同样清楚的,是这些战役\textbf{没有全体的团结便无以进行}。依我看,Jouanolou却不太在局中——他不俨然“权威”,我的印象是他久已离开上同调的泥沼(注85$_{1}$)。但我难以想象,Illusie与Berthelot不曾亲手捧读过SGA~4½与“好的参考文献”两者;他们跟我一样识字,也不比我笨。

倘若Illusie突然张罗起SGA~5的出版——恰在Verdier取用、Deligne取用、Deligne需要一座后勤基地以成其大名鼎鼎的SGA~4½(从中顺手牵走孕育了这份文本与其全部著作的那两场研讨班)之时——而Illusie手里已有十年功夫可以做这件事,这肯定不是偶然。倘若我1966年所做的收官报告——论未决问题与猜想——“不幸没有成文,其非常漂亮的开篇报告也没有[sic]:后者巡视了各语境(拓扑的、复解析的、代数的)中的Euler--Poincaré公式与Lefschetz公式”,这同样肯定不是偶然——而那是一场我不在行的下葬。同样不是偶然:把研讨班的一场关键报告截肢、不经任何手续送进SGA~4½,在Illusie与Deligne看来竟如此自然(而且只配在“细节变更”里顺带一提)\fnote{84}{(*)}{(5月16日)实际上,正如次日我就发现的(见注n\textsuperscript{o}~87):对母研讨班(或父研讨班!)SGA~5,在Verdier、Deligne与Illusie的手里,发生了一场真正的“屠杀”。}。

Luc~Illusie——我对他怀有如对Pierre~Deligne一样的感情,而他(一如Deligne)在我面前始终极其亲切\fnote{85}{(**)}{甚至在我1970年离去之后,Illusie对我仍有过细致的体贴——比如,长年以来,岁末节庆之际,他都寄给我非常漂亮的贺年卡。恐怕我不常回信致谢、报个平安——这些忠实的友情的信使,来自一个显得无限遥远的过去,而我已经与它失联。\par (5月16日)相反,Illusie从无继续或重建数学层面接触的念头;去年(我为数学问题联络他时),我感到他的迟疑。自我离去以来的十四年里,我只收到过他一份、唯一的一份抽印本,日期1979年。}——他的(有意识的与无意识的)意图,我不知道。但我确认:他在Deligne身边,成了一场\textbf{毫不脸红的欺瞒}的共同演员:那场欺瞒,把1965/66年的母研讨班SGA~5——恰恰是Deligne头一回听说概形、平展上同调、对偶及其他“枝蔓”的地方——说成一份顶着欺眼之名SGA~4½的文本集(晚八年写成,却佯装在前:既凭标题里的编号,也凭Lecture~Notes里的出版号,更凭作者那句不寻常的评语“它的存在(SGA~4½的)不久将使人得以\textbf{照原样}出版SGA~5”——着重号是我加的)的一种无形而近乎滑稽的附录——并且还摆出对那份寒碜的合集所自出的全部工作的毫不掩饰的轻蔑。

没有这些被如此潇洒地轻慢的工作,奠定其应得名望的Deligne的大部头著作,眼下没有一部写得出来——倘若真要问,一百年后也写不出来(Illusie与我的其他上同调学家学生们想必同样)。在这场“SGA~4½之役”的精神里,有一种\textbf{无耻}\dmark{}\origpage{347};Illusie(想必连自己都不觉察)为它作保;而它之所以能如此招摇,只因有一个\textbf{共识}的心照不宣的首肯。除Deligne本人之外,卷入这一共识的头一批人,恰恰曾是我的学生、某份遗产的主要受益人:那份遗产,在他们眼前,被交付给抢夺集市的风险与轻蔑。

还有那份武断的自满神气,那份慈父般的庇护神气——前天在我们通电话时,我还在我的前学生身上领教过\fnote{86}{(*)}{关于这次通话,见“玩笑——或‘权复形’”一注(n\textsuperscript{o}~83)。};还有我的朋友皮埃尔身上那份较隐约的居高临下——自“SGA~4½--SGA~5”双重大捷的翌日起我就领教过(而当时,乃至其后七年,我对它连一丝疑心都没有)——这些神气,绝不是孤独的产物,而分明仍是一个\textbf{从未被质疑过的共识}的征兆。这些神气,告诉我的不仅是Verdier与Deligne,还有我曾经的学生们全体;而首当其冲,是那些(就其工作主题与日常操持的工具而言)最直接相关的人。

“欺瞒”这个词——它不召自来——适时地使我想起另一场欺瞒,其中摊开着同样的犬儒:即所谓“反常”研讨会的那一场。如今,这两场在我眼前\textbf{血肉相连、不可分解——使两者皆成为可能的是同一个精神}。也许除Jouanolou(他已不太混迹“上流社会”)之外,我认为这些前学生是这场耻辱的共同责任人、团结一致的上同调学家。至于Berthelot与Illusie,没有任何东西允许我预设恶意或不诚(在Verdier与Deligne的个案里,那是不容置疑的)。但我至少确认一种盲目、一种健全官能之使用的堵塞;其深层缘由当然向我滑脱。倘若他们心中没有那份冷漠与轻蔑的蓄意,那么Zoghman~Mebkhout——七十年代唯一公开自称出自我的著作的人,而且题材恰与两人切身相关(尽管他们不屑于觉察)——本会享有最低限度的“善意推定”,使他们至少稍微看一看他所做的事情,从而觉察他自1974年起所投入的方向的趣味;那趣味是\textbf{明摆着的}\dmark{}\origpage{348}!两人都不屑于从一个装作还出自Grothendieck的无名小卒那里觉察任何东西。无名小卒的论文,是他亲手寄给他们的;我不知道他们开没开封,或读没读那些更短、更易消化的说明文字——总之,他们不屑于哪怕寄一张收条(Deligne亦然——分明是他定调子的\fnote{87}{(*)}{(6月12日)其间我得知:两人都不曾参加那场研讨会(1981年6月的吕米尼)。不过,可参见“欺瞒”一注,n\textsuperscript{o}~85$'$。})。

这当然不曾妨碍:与那场值得纪念研讨会的其他与会者一起,他们带着兴趣了解了那条卓越的“Riemann--Hilbert对应”,却想不到就其来源或父权提出丝毫问题,或至少(作为扎实的数学家)问一问它在哪里得到证明(注85$'$)。不过,在这一点上,我信得过Deligne:他自会乐于优雅地向他们讲解那份证明——对像他们那样的人,想必明白之至——恰恰是那种借Hironaka式奇点消解的证明,他们久已学得,而且不是从别人、正是从我这里(注85$_{2}$)。Riemann--Hilbert,Hironaka,阿布拉卡达布拉——戏法告成!

显然,一如Verdier、一如Deligne,他们已全然忘了何为\textbf{数学创造}:一个经月历年缓缓澄清的视野,揭示那件“明摆着”却无人看见的东西,在一个无人想到过的“明摆着”的陈述中成形(而在本案里,一位Deligne曾徒劳地试过整整一年……)——随后随便什么人都可以用现成的技术在五分钟内证明它,靠着他在某场遥远研讨班的长凳上有幸学得、如今却不屑(或不记得)想起的技术……

倘若我谈论Berthelot与Illusie时毫不容情,那并不是我尤其想给他们抹黑(在与他们的两位朋友清账之后)。我知道:他们既不比他们亲爱的同行们“更坏”、更蠢,也不比我更坏更蠢;而我在他们身上、在这一场合所察觉的那份嗅觉与健全判断的欠缺(有时,还有那份对他人必要的尊重的欠缺……)绝非痼疾,而是一个\textbf{选择}的效果。无疑,这一选择曾给他们带来中意的\textbf{收益}\dmark{}\origpage{349};而随着我的反思而来的另一份“收益”,对这一位或那一位,也许不合时宜。倘若如此,那只说明:它再度重演同一个选择——即在自身官能的微不足道的一小部分上运转的选择,不惜把鱼目当珍珠、把珍珠当鱼目,并无望地把空核桃(好哥们的)与满核桃(无名外人的)混为一谈。各人自便!($\Rightarrow$注86、87)

\notehead{85$_{1}$}{}

Jouanolou是我学生中除Verdier外唯一一个不曾把发表其论文放在心上的人。在我看来,这是对他所发展的奠基工作——即导出范畴观点下的$\chi$-进上同调——冷淡的征兆。由于他这一主题的工作,大部分是在我离去\textbf{之后}进行的——即在同学们以Deligne与Verdier为首、对我引入同调代数的观念(尤其导出范畴观念)之普遍冷淡已然亮出信号的时刻——气候并不鼓励Jouanolou认同于自己的工作、给它(应得的)发表之荣。而当这同一些Deligne与Verdier,追随着Zoghman~Mebkhout(别名Verdier的无名学生、别名Grothendieck的遗孤学生)的著作,终于(在大张旗鼓、互相吹捧之中)发现导出范畴的重要性时(见注n\textsuperscript{o}~75、77、81),那份遭冷落的Jouanolou论文,自“反常”研讨会起,重新恢复了全部的现实性;倘若1970年我离去之后,概形上同调理论的发展照常进行,那份现实性本不曾中断。一个触目的细节,例证Deligne在我离去之后其取舍的一场急遽“转向”:恰恰是Deligne本人(他分明深知在三角范畴框架里发展$\chi$-进上同调形式主义的重要性)向Jouanolou提供了一条技术关键想法,供一个有待研究的、成形的$\chi$-adic三角范畴的定义之用;那条想法在论文中得到发展。(对此可参见我1969年论Deligne工作的“报告”,第8节。)

(5月30日)关于Jouanolou的工作,另见“共同继承人……”一注,n\textsuperscript{o}~91。

\notehead{85$_{2}$}{}

一份有意味的“巧合”:恰恰是在同一场SGA~5研讨班里,这一整批人学得了那条证明原理——它既用于证明平展上同调里的双对偶定理(在拥有奇点消解的情形),也用于诸$R^if_*$无真性假设的\textbf{有限性}\dmark{}\origpage{350}定理,对$R\mathcal{H}om$、$Lf^!$亦然。(这些有限性定理,同样被从SGA~5的已刊版本里藏掖走、并进了SGA~4½;Illusie甚至不认为有必要在其导言里点明——我是写这些文字时才发觉的!)Zoghman没有听过那场研讨班的福气(他得到的是“好的参考文献”的替代品),是在另一处我使用过该办法的地方(为$\mathbb{C}$上光滑概形的De~Rham定理)学得它的。他其实也能在“好的参考文献”里学到:在那里,我的证明在解析框架里被照抄,以确立我的SGA~5学生们从此乐于称作“Verdier对偶”的东西(那条对偶,在我结识他的荣幸之前,我就已经知道了)。分明,一切都环环相扣!\textbf{同一份证明}(连同陈述,一并抄自我)被Verdier用作一种对偶的父权凭证——那种对偶,他除了在这场被肢解、被交付轻蔑的SGA~5研讨班里,无处学得;而它又被用来\textbf{对付}Mebkhout,成为(正因其“明摆着”)心照不宣的借口与手段,毫不脸红地巧取一项重要发现的信誉。

(5月30日)我记得,我头一回使用Hironaka式奇点消解、并领会消解作为证明工具的非凡力量,是为了对Grauert--Remmert定理做“三勺汤”式的证明——该定理描述某类复解析空间的有限覆盖上的复解析结构;以及在$\mathbb{C}$上有限型概形的情形里的类似陈述。(并不排除:恰恰在那次,原理是Serre向我提示的。)后一结果,是证明平展上同调与寻常上同调之比较定理的主要配料(其余,仰仗$Rf_!$形式主义归约;再稍加消解,从$Rf_!$过渡到$Rf_*$……)。

\subsection{欺瞒}

\notehead{85$'$}{}

(6月3日)实际上,我得知:他们本无须为这份父权提出问题,因为Berthelot与Illusie都是从Mebkhout嘴里学得好上帝定理的——前者在1982年2月,后者早在1979年(Mebkhout论文答辩之年)。虽然两人都不曾参加那场研讨会,他们却是那场研讨会之欺瞒的同谋,因为\dmark{}\origpage{351}他们不可能不知道那场围绕Mebkhout对好上帝定理尤其之父权的藏掖。我还可以想象:他们与研讨会的全体与会者一起,争先恐后地做了那场集体欺瞒的头一批受骗者——那场欺瞒,由他们的朋友Verdier与Deligne悉心操办(我的五位上同调学家学生中,有四位俨然是它的同谋)。至少就Illusie而言:去年夏天Mebkhout来我家之后,与Illusie通电话时,他对Mebkhout显然满不在乎,令我吃惊——他十分惊讶(几乎为他老师难过:他本指望老师有更好的判断力……)看到我在代数簇上同调理论的重启中,把头等角色派给Mebkhout。一些力量可观的共识,早已决定把Mebkhout归入无名小卒之列;我的朋友Illusie与这份三重矛盾活得逍遥自在,不问一个为什么:好上帝定理及其哲学的头等角色;围绕这些事情之父系的藏掖(他本人也在大帮人之中参与其中);以及他对Mebkhout的器量与角色的轻视——尽管他分明知道:Mebkhout是这些东西从未被点名的作者,而这些东西革新了数学的一个领域——在那个领域里,Illusie本人俨然权威。

我在这里重新找到那种健全意识与健全判断的完全堵塞——甚至在一桩表面上如对科学问题的评判般不涉私人的事情里;这种堵塞我已有不止一次机会影射,而每一次都重新令我错愕。而我在Illusie(想必还有许多别人)对Mebkhout——我的“遗孤学生”——的关系里确认的这份矛盾,肯定不是别的,正是他与我的关系中那桩更关键的矛盾的众多效果之一。正是这份矛盾——在他身上尤其,也在我的其他学生身上——在本书这支由我昔日学生们组成的葬礼行列的各注所进行的反思中,越来越清楚地显形……

\subsection{亡故者}

\notehead{86}{}

(5月11日)一如常有的情形,我是带着几分迟疑开始这场新反思的,主题“SGA~5--SGA~4½--反常性”——它似乎已被检视再检视到腻:“这会给一位早已听腻的读者留下糟糕的印象;再进这样的细节毫无雅致可言,SGA~5呀SGA~4½呀,这些都是过去了的事,不值得再饶舌……”

\dmark{}幸好我没有被这类熟悉的调门吓退——它想以“分明不值得”、“由它去吧”为口实,阻止我探一件事的底(至少就我当下所能探到的深度)。\origpage{352}倘若我曾发现过我认为有用而重要的东西,那从来是在这样的时刻:我懂得不听那份冒充“理性”、甚至“体面”的声音,而顺从我身上那份不体面的欲望——去看看甚至被认为“没意思”、寒碜,乃至粗劣、不体面的东西。我一生中不记得有哪一次,我曾后悔把某样东西看得更近一点——对着那些恨不得阻止我的积习反射。在《收获与播种》里,这些抑制反射比别的场合更强,因为这场反思注定要公之于众,这就立刻带来分寸(当我牵连第三方)与简洁(顾念读者)的约束。但我最终并不觉得,这些约束曾有任何一刻阻止我触及我想触及的东西、或把它深入到我感到欲望的地步。在曾一度显得棘手的临界情形里,我径直前行,带着这样的把握:必要时,我总还留着不把那场不客气的反思所“流出”的东西收进《收获与播种》的余裕。这些“边缘个案”,只在我犹豫是否牵连他人时出现;牵连我自己,则从未犹豫。但即使在第一种情形,事情竟是(这对我是一桩意外):我从不曾动用过那份“余裕”:《收获与播种》的文本,是我反思的全本——至少,是这场反思中找到通往书写之路的那一部分。

我感到:随着前一条注\fnote{88}{(*)}{即同日的“团结”一注,n\textsuperscript{o}~85。}那场简短的反思,局面大大澄清了。我是说:一个曾经故意弄得混乱的局面的某个本质侧面——我刚以“主题”的三重名目(SGA~5--SGA~4½--反常性)唤起过它——以全光显现在我眼前:即一种“团结”、一种“默契”;它们此前只被模糊地知觉过。这绝不意味着我自以为探明并理解了一个复杂局面的全部机关与来龙去脉——它至少直接而分明显然地牵连七个人:Zoghman~Mebkhout(在某种意义上充当某一局面的“显影剂”)、我的五位前上同调学家学生\dmark{}\origpage{353},以及我自己。我甚至不敢自称:对“SGA~5等等”这一局面里、在我自己身上起过作用的全部机关与动机——那场“倒霉的研讨班”距今已近二十年!——我都有所知觉。但我感到,比起昨天(或甚至今天早晨),我的状态好得多,足以理解并安放我希望将从这个或那个主要当事人那里到来的回声。

摆在我面前的主要问题(它想必在反思的另一个阶段已在场,如今带着新的力量重现)是(依我看)这个:经由我的学生们(差不多全员)进行的这场葬礼所发生的一切,是一件全然\textbf{非典型}的事,系于我的个人与我的独特命运的某些特点(如近十五年前我退出数学舞台,及围绕它的种种情势,等等)?抑或相反,是一件“再自然不过”的事,出于机缘的简单巧合——按“机会造就窃贼”的原理?我犹豫相信后者,尽管此刻辨不出、甚至隐约瞥不见:我这个人身上的哪个特殊方面,竟有这种本领,在我昔日学生们中间造就一份如此完美、如此一致的\textbf{赞同},去埋葬“老师”,并埋葬那些自称出自他、或著作分明带着他的印记(却并不因而是“他们的”)的人。是我这个人周围那圈“父亲”的光环吗——我已有机会谈过?抑或,我的离去这一件事本身,对他们的每一位置疑?此刻,我全然说不出来,因为缺乏那双会看的眼睛……也许到来的月份会就此教给我点什么\fnote{89}{(*)}{(5月30日)这个方向上的反思,见“掘墓人——或全体教众”一注,n\textsuperscript{o}~97。}。

最近三个星期里,我不止一次想到另一桩奇怪的“巧合”:葬礼“在其全部光辉中”的发现(经历LN~900--SGA~4½--SGA~5--反常层会议四幕,然后回到SGA~5与SGA~4½)——这项发现恰在此时完成:我刚结束一场对我的数学家过去、我与学生们之关系的深入反思。也就是说:恰在我刚刚——尽我官能之所能、并在当时已知事实的限度内,照它们经由常常雾蒙蒙的回忆所还原的样子——就这段过去“与自己算清账”的时刻\dmark{}\origpage{354}。换个说法:那恰是我终于\textbf{准备好}得知这件事、并从中获益的时刻。

“偶然”把事情安排得如此之好,以至禅观中连一个中断都没有。那场反思——它以对那些(依我所感)最重要的观念之命运的一场简短回顾开始\fnote{90}{(*)}{见3月31日的“我的孤儿们”与“拒受遗产——或矛盾的代价”两注(n\textsuperscript{o}~46、47)。}(那场反思还停留在某种模糊之中,只有某种基调执拗地浮出……)——这场反思在4月19日那个星期四,全然自然地继续下去。诚然,那还是带着“无耻”——借用方才的词,它也恰当地描摹了我当时感到的一桩事——的印象所激起的情绪:读那部“值得纪念的文集”LN~900之时。

在这“同一场”反思的新起点上,主发动机是“老板”——我的自尊、我的体面之感受了伤;写下我的情绪,我便在某种程度上从中解脱。分明是“自我”、“老板”在随后十天里领舞——那些日子,没有笑容也没有笑声,只有密不透风的严肃。无疑我必须经过那里、经过那十天的弯路,反思才回到它所离开的中心——我自己的身上。我仍记得那次回归的如释重负——像出了隧道重见天光!就在那时,我找回了笑与微笑,仿佛从未分离。那是4月29日。次日30日,当月最后一天,我欢欢喜喜地为反思的这一最终阶段画上句点。

那肯定也是我完全准备好接收下一个“包裹”的时刻——这一回由我的朋友Zoghman惠寄——“研讨会”包裹,两天后到达。到今天,我消化那份包裹的实质已进入第十天。而在这一阶段,尽管我为这场没完没了反弹的反弹而咬着嚼子,笑容却一天也不曾离我而去。而今天,我真的相信(诚然,已是第无数回!),是终于画句点的日子。

五天前,我就有过这同样的抵达终点之感,觉得剩下的只是总务活:这里那里补几条脚注,把涂改过多的页誊清\dmark{}\origpage{355}(那每一回都是思路尚欠清晰的征兆,它要求经由这看似机械的功夫各就其位;而文本经过誊清,总是带着一副新面目出来……)。那是在我刚写完如今那则“我的朋友们”注(n\textsuperscript{o}~79)之时;它自然而然地续成了几段“终曲”。但我终于把这些终曲与注的开头分开了。因为事实证明:那桩大名鼎鼎的总务活爆了:那些不打行距敲出的“脚注”,变成了真正的(而\textbf{非}脚注的)漂亮条注,颇具规模,不得不隔行重新打出,然后想法安置在此处或彼处。又过了几天,我才向明摆着的事实低头:继名为“研讨会”的行列之后,另一个行列正在成形,要加入送葬的队伍——而压轴的行列,将不是我心中已定的那场研讨会,而将由\textbf{学生}领队。而且就在今天,当第一个行列——已缩成一条注——刚添上第二条(“一种不义与无力的感觉”)时,我也知道了将由谁领队:是“遗孤学生”。于是,这支由一个小写的“学生”(遗孤的,合乎其寒微的身份)开道、又由一个大写的“学生”(这回毫不寒微)压阵的队伍,我看,终于齐备了!

现在,我想,是时候了——在第一次“假到站”之后——回到最后一曲《自深处》(De Profundis)的和声;它们今天比五天前更贴切。以下是我当时记下的;它们同样表达我此刻的心情。

(5月31日)最终,这又是又一次“假到站”——“终曲”这一回还是提前了!二十天过去,其间“总务活”不断爆成对被忽视的这一侧面或那一侧面的反思的重启。又有六条注加入本该压阵的“学生”行列。灵车在“学生”身后现身,载着四口棺材,掘墓人随行。分明还缺些东西,好给一支似乎谁也没在送葬的队伍以实体与意味。

经验使我变得谨慎;我等着事件到来,此刻不冒险预言:队伍是否终于齐备,或者是否还有一个被遗忘的行列,会在最后\dmark{}\origpage{356}一分钟钻进来,不肯错过那终极的仪式\fnote{91}{(*)}{(6月12日)谨慎确有其必要:一个新的行列“我的学生们”已从最初称为“学生”的行列里分离出来——后者改称“学生——别名老板”。}。

\subsection{屠杀}

\notehead{87}{}

(5月12日)\fnote{92}{(**)}{本注接续前一天的“团结”一注(n\textsuperscript{o}~85)的反思。}为了多少算上同调学家的读者的教益,也尤其为了我自己,我想把这场对一个辉煌研讨班的按部就班的劫掠——它落在我的两位前上同调学家学生手里,而在其余人的善意的目光之下\fnote{93}{(***)}{反思的后续将显明:这些“其余人”中的一位,曾为这场战役有效地出手,代人行事。}——的细节巡视一遍:正是在那场研讨班里,他们比全世界早十二年、从工匠本人手里,学得了使他们成名的手艺的基础与精微之处。

我的两场口头报告,从未以任何形态交到公众手里。其一是收官报告,论未决问题与猜想——它“不幸没有成文”,鉴于其分量之轻;而屠杀版的导言作者也确实认为,连提一提所论的是\textbf{哪些}未决问题与猜想,都是多余的。他又何必费那份心呢——那不过是些问题(人人尽可随意自设!)与猜想(连证明都没有!)(注87$_{1}$)。其二是开篇报告,它把研讨班一举安放进一个更宽广的语境(拓扑的、复解析的、代数的),并巡视了Euler--Poincaré、Lefschetz、Nielsen--Wecken各型公式——其中一些,构成研讨班的主要应用之一。导言作者用“……也如同……”一笔带过这场报告的消失,那种措辞对一份\textbf{漫不经心}的心境说得很多;在当时——研讨班的作者已从流通中消失七年之际——这种心境分明是理所当然。

有整整一组报告,是我做的:关于与一个闭链相配的同调类与上同调类的形式主义(上同调情形取正则环绕概形)\fnote{94}{(****)}{细节见注n\textsuperscript{o}~82“好的参考文献”。}。它们被公平地瓜分了:上同调归Deligne,同调归Verdier——后者还是稍微越界到上同调;作为回报,则不吝以那大名鼎鼎的“权复形”向Deligne道个小小礼\dmark{}\origpage{357}\fnote{95}{(*)}{见注83“玩笑——或权复形”。}。(更不必说他还把$R\mathcal{H}om$的有限性定理与双对偶定理据为己有——一字不差地抄自研讨班;反正,狮子的份额归Deligne,那是理所当然……)导言作者认为连提一提同调的报告也无须。确实无须:前一年,他的朋友Verdier已揽下补足那份缺失的“好的参考文献”的差事(对一场研讨班、对我,均不置一词)。

有关于$R^if_*$(f非真)运算之有限性定理的口头报告,以及作为推论,关于$R\mathcal{H}om_*$与$Lf^!$运算的报告。关键定理以Hironaka式奇点消解的技术证明(因而只在拥有消解的情形有效)。我所使用的这些论证,自研讨班以降已成通例(见注85$_{2}$)。Deligne得以在另一些更便利于应用、如今在大多数应用里得到验证的假设下,证明这些有限性定理,以及双对偶定理。人们本可以指望:他会要求把这些改进纳入那场他有幸学得平展上同调、学得其后来全部著作之基础的观念与技术的研讨班。但这一情势,却被拿来当作“理由”,把研讨班的这一部分截去。至于双对偶定理,从此在Illusie的笔下(在概形框架里)成了“Deligne双对偶定理”(报告I的导言)。这也算公道:解析情形里,Verdier前一年就早已把父权判给了自己(甚至无须费心另找一份证明)。

有一场报告,发展一条“泛Künneth公式”,它曾由Illusie成文。此前没有人想到过提炼这一类陈述;其灵感在于这样的直觉:“一般地”——即在基的一般点附近——一个相对概形表现得如一个“局部平凡纤维”(在拓扑语境里)。经由一条与前面提到的证明相仿的优雅证明,Deligne得以消除我所做的奇点消解假设。判给了——报告删除,以对同一位Illusie在那场号称“在前”的研讨班SGA~4½里的报告的一条引证“取而代之”。

\dmark{}\origpage{358}有一组报告,论非交换迹的形式主义;它是作为手段发展的,用以在从未被处理过的情形里显豁Lefschetz--Verdier公式的局部项。据说,这些报告最终由Bucur成文,其手稿“在一次搬家中天赐般地丢失了”——这简直成了通俗笑剧!\fnote{96}{(*)}{想必正是这一情势,当场启发了Deligne对SGA~5的那条光辉批评:说Lefschetz--Verdier公式(别忘了——它“还停留在猜想阶段”!!!)的局部项在那里竟连计算都没有!(关于这条批评的离奇——在知情的读者看来,它与Verdier前一年那条大名鼎鼎的“权复形”难分伯仲——见“白板”一注,n\textsuperscript{o}~67。如此一来,成风的是Verdier!)}而在Illusie执笔的SGA~5导言里,这些报告居然变成了“Grothendieck的\textbf{交换}迹理论,以光辉方式推广了Stallings的理论”(而它们分明是\textbf{非}交换的!)。这个口误\fnote{97}{(**)}{即把“非交换”迹理论之父权派给我的那个口误——那顶王冠,倒也没人等我来戴。它竟然一直保留到已刊版本里,越发引人注目:Illusie是我学生中工作上最一丝不苟、直到最后细节的一位,想必。}只能出自一位秘书欠佳(或太好……)的灵感;她想必与我朋友Ionel~Bucur的搬家公司过从甚密。(“光辉地”一词是我的笔插进去的,为的是更好地还原这个同样天赐的口误所万无一失地暗示的思想。)

我无从抱怨,既然Illusie把重做的活揽在了自己身上(而且,照他的说法,还是个“更精致的”版本,因为这按层论的方式调了味——不过我记得,Illusie,在我那个年代,你做过的“精致”创新,可比这一手多……)。想必他费了不少时日——如果我没记错,当年我把这台机器调试好,花过几个星期;说不定我的手稿也在同一场天赐的搬家里丢了,天知道哪位被我的滔滔不绝淹没了的亲爱听众,至少记下过能读懂的笔记……

有一件事值得注意,我此前不曾留意:他没有把这场报告插回它原定的报告XI的位置(那想必也是它在口头研讨班里的位置),而宁肯在那个位置留下一个大窟窿,把他的报告做成一场伪托的报告,称为“局部项的计算”。这标题倒分明对应我记得自己在口头研讨班里做过的东西——奇怪。但就在他为这场报告所写导言的第1行,作者赶紧为我们\dmark{}\origpage{359}澄清误会:“这份报告,成文于1977年1月,\textbf{不对应于研讨班的任何一场口头报告}”。随后接以Lefschetz--Verdier公式(这名字分明对我说过什么;而我记得自己分明纵横发展过一套非交换迹理论,正是为了在某些情形里计算“局部项”……),继以一条Langlands公式,以及一份1967年的Artin--Verdier证明(那毕竟是口头研讨班“终曲”之后的一年;它不可能不影响到这两位作者——其中至少一位,或许两位都听过研讨班)。到这一页的末尾,我们才顺带得知——与开头的宣告相反——还有“这份报告的第二部分,性质技术得多”(这种语言我在哪里读到过……),它(请欣赏这个措辞)“\textbf{受Grothendieck为建立曲线上某些上同调对应的Lefschetz公式所用的方法的启发}”,附一处对同一场研讨班报告XII的引证,尤其对不可或缺的SGA~4½的引证;分明,为了这么点东西,没有理由把它插回大窟窿的位置——方才那份“更精致的版本”想必已尽善尽美。Illusie与Deligne还肯以“启发”之源的礼遇引我,可谓客气;而他们朋友Verdier前一年的榜样,分明早已表明:这份顾忌,完全用不着了。

我回到Illusie为顶着SGA~5之名的那一卷所写的导言。我们在那里再度得知——一如Deligne在其SGA~4½导言里已宣告的——研讨班之所以终于出版,全靠\textbf{了他的朋友}:

\begin{quote}
“我感谢 P.~Deligne说服我,在报告III的新版里,写出Lefschetz--Verdier公式的一份证明,\textbf{从而排除了这份研讨班出版的障碍之一}”。
\end{quote}

我们又一次深陷闹剧——那位顺从的Illusie,在SGA~4½的导言里照单全收!倘若研讨班十多年不曾出版,那是因为(难就难在想得到)没有一个人(在1977年救场的Deligne之前)想到过:为那条(不无道理地)称为“Lefschetz--Verdier”的公式写一份证明,也许是个好主意——而那份公式的父权,除了他那位\dmark{}\origpage{360}形影不离的朋友、我的前学生Verdier亲自,\textbf{至少从}1964年起(注87$_{2}$)就骄傲地佩戴着;也就是说,至少早在我那场研讨班结束之前两年——它当时只待好心的志愿者,便能交到人人手中!

最后,作为研讨班的又一场(最后的?)残损:Serre就“(Serre--)Swan模”所做的漂亮报告消失了——那份报告题为“Brauer理论导引”。所幸Serre,眼见事态的走向,明智地把那份报告收进了他的书《有限群的线性表示》(Hermann,1971),把它交到数学公众手里。(注87$_{3}$)

这一回,我想,这幅画我巡视完了。这幅画,关乎一场研讨班的命运——我曾把自己的最好部分倾注其中(注88)\fnote{98}{(*)}{关于“自己的最好部分”这一说法的意味,见随后的“遗骸……”“……与身体”两注,n\textsuperscript{o}~88、89。其中第一条,把SGA~5研讨班——连同与之密不可分的SGA~4——安放为我“彻底完成”的著作部分的主干。};二十年后重逢,它已面目全非,被那些曾是它独家受益人的人屠杀——或者至少,被其中三人,并获得所有其他与会者的首肯。

我并不后悔再一次费心把渐渐强加于我注意力的东西追到底。我曾在对我的一位前学生的关系的长篇回顾之末,确认这场“事物的回归”\fnote{99}{(**)}{见同名注(4月30日,n\textsuperscript{o}~73)。},当时已隐约预感到:他不是唯一一个“兴冲冲为我送葬”的人——而只是到了此刻,我才得知它的气息、它的“气味”(借用当时我的一场梦里出现的说法)——一场\textbf{暴力}的气息。这份气息,既被那种呈献高度技术性实体的(表面超然、不动声色的)话语所掩藏,又被它所泄露\fnote{100}{(***)}{尤其指伴随SGA~5(Illusie执笔)与SGA~4½(Deligne执笔)的那些导言性文字里的话语。}。这场暴力所瞄准的,经由一副任人宰割的“遗骸”,恰恰是那位曾是“老师”、“父亲”的人本身——而与此同时,“学生们”久已占据了他那令人艳羡的位置,不曾遇到任何抵抗;并且\dmark{}\origpage{361}久已在他们中间选出新的“父亲”,注定取代旧者、统治他们。

我感到这份气息,然而它于我仍是异己的、未被理解的。要“理解”它,想必这份气息得在我身上活着,或曾在我身上活过。但四年前,我头一回感到并掂量了我生活中一桩我从未想过、总以为理所当然的事:即我童年时代对父亲的认同,\textbf{不曾}带有冲突的印记——在我童年的任何时刻,\textbf{我既不怕我的父亲,也不嫉妒他},同时对他怀着毫无保留的爱。这层关系——或许是给我的一生打下最深刻印记的关系(在这场四年前的禅观之前,我甚至不曾觉察它)——在我的童年,如同对另一个既是强力的又是慈爱的自我的关系——这层关系,不曾打上分裂与冲突的印记。倘若在我那常常撕裂的一生里,对我身上所怀之力量的认识始终活着;倘若在我那绝不免于恐惧的一生里,我既不惧怕任何一个人,也不惧怕任何一件事——我把它归功于这份寒微的情势;直到我年过五十,它仍不为人知。这份情势是一件无价的特权,因为对自身之内创造之力的亲密认识——它\textbf{就是}那股力量——正使力量得按其本性自由表达:经由创造,经由一种创造的生活。

而这份特权——它使我免于冲突最深的那几道印记之一——此刻却像一副枷锁,像我生活经验里的一个“\textbf{空缺}”。一个难以填补的空缺,在那里,许多别人拥有情感、意象、联想的丰富织体,向他们提供道路(只要他们好奇肯走),既深刻理解他人、也深刻理解自己;面对那些境况——我靠反复与旁证才能勉强把握,却仍像一个异乡人——我身上求知的欲望,忍饥挨饿。

\notehead{87$_{1}$}{}

(5月31日)这场收官报告——与开篇报告一起,想必是最有趣、最实的一份——显然并非对所有人都白费,如我读到MacPherson的文章“Chern classes for singular algebraic varieties”(Annals of Math. (2) 100, 1974, 第423--432页;1973年4月收到)时所看到的。我在其中,以“Deligne--Grothendieck猜想”之名,重新找到我在这场报告里、于概形框架中引入的主要猜想之一\dmark{}\origpage{362}。MacPherson在复数域上代数簇的超越框架里重述它,Chow环换作同调群。Deligne1966年在我那场报告里学得这条猜想\fnote{101}{(*)}{(6月6日)诚然,形态略有不同;见本注今日所写的后续。\par (1985年3月)Deligne本人提供的细节,见“把一个个 i 点清”一注,n\textsuperscript{o}~164(II~1)。}——那恰是他在那场研讨班里初次露面、开始熟悉概形语言与上同调技术的一年(见“格格不入者”一注,n\textsuperscript{o}~67$'$)——几年之后,把我算进猜想的命名里,想必就不再合宜了……

(6月6日)我借此机会在此写明:我在研讨班里、于概形框架中所陈述的猜想是什么;当时我肯定还指出了它在复解析框架(乃至刚解析框架)中的明摆着变形。我把它构想为一个“Riemann--Roch”型的定理,只是取离散系数而非凝聚系数。(Zoghman~Mebkhout还对我说过:他的D-模观点应当使人可以把这两个Riemann--Roch定理都看作包含在同一条晶状Riemann--Roch定理之中;后者因此在特征零,代表我在数学里引入的两条Riemann--Roch定理——一条在1957年,另一条在1966年——的自然综合。)取定一个系数环$\Lambda$(不必交换,但为简单设为诺特的;且为平展上同调之需,其挠性与所论概形的特征互素……)。对一个概形$X$,以
$$K.(X,\Lambda)$$
记由$\Lambda$-模的平展可构造层构成的Grothendieck群。利用函子$Rf_!$,该群函子地依赖于$X$(对诺特概形$X$与分离有限型的概形态射)。对正则$X$,我假设存在一个典范的群同态,扮演凝聚RR定理中“Chern特征”的角色\dmark{}\origpage{363}:
\begin{equation}\tag{15.1}\label{eq:discrete-chern-map}
ch_X: K.(X,\Lambda)\to A(X)\otimes_{\mathbb{Z}} K.(\Lambda),
\end{equation}
其中$A(X)$是$X$的Chow环,$K.(\Lambda)$是由有限型$\Lambda$-模构成的Grothendieck群。这个同态,应当由“离散Riemann--Roch公式”对正则概形间一个\textbf{真}态射$f:X\to Y$的成立而唯一确定;该公式写作与凝聚Riemann--Roch公式同样的样子,只是把Todd“乘子”换成相对全Chern类:
\begin{equation}\tag{15.2}\label{eq:discrete-rr-formula}
ch_Y(f_!(x)) = f_*(ch_X(x)\, c(f)),
\end{equation}
其中$c(f)\in A(X)$是$f$的全Chern类。不难看出:在一个拥有Hironaka强形式奇点消解的语境里,RR公式确实唯一确定诸$ch_X$。

当然,这预设我们处在Chow环有定义的语境里。(我不知道有谁哪怕尝试过,为并非有限型于一个域上的正则概形写一套Chow环理论。)否则,也可以在凝聚语境里惯常的“Grothendieck”环$K^{\circ}(X)$所配的分次环里工作,按惯常方式过滤(见SGA~6)。还可以把$A(X)$换成偶数次的$\ell$-进上同调环,即诸$H^{2i}(X,\underline{Z}_\ell(i))$的直和。这有一个缺点:引入一个人为参数$\ell$,并给出较“纯数字”的公式;而Chow环的魅力恰在其连续的结构——过渡到上同调时被摧毁。

早在$X$是代数闭域上光滑代数曲线的情形,$ch_X$的计算就牵涉Artin--Serre--Swan型精致的局部不变量。这就是说,那条一般的猜想是一条深刻的猜想;对其的追求,系于对这些不变量的高维类似物的理解。

评注。同样以$K^.(X,\Lambda)$记由Tor维数有限的$\Lambda$-平展层可构造复形构成的“Grothendieck环”(当$\Lambda$交换时,该环作用于$K.(X,\Lambda)$……),人们同样应有一个同态
$$ch_X: K^.(X,\Lambda)\to A(X)\otimes_{\mathbb{Z}} K^.(\Lambda),$$
它同样(mutatis mutandis)引出同一条Riemann--Roch(RR)公式。

\dmark{}现在设$\mathrm{cons}(X)$为$X$上可构造整值函数的环。人们以或多或少许套套逻辑的方式定义典范同态
$$K.(X,\Lambda)\to \mathrm{Cons}(X)\otimes_{\mathbb{Z}} K.(\Lambda),\qquad K^.(X,\Lambda)\to\mathrm{Cons}(X)\otimes_{\mathbb{Z}} K^.(\Lambda).$$
倘若现在我们只限于\textbf{特征零}的概形,那么(利用具真支撑的Euler--Poincaré特征)可以看出:群$\mathrm{Cons}(X)$对诺特概形的有限型态射是共变函子(此外,作为环函子,它是反变的——这与特征无关),而前面的套套逻辑态射是函子的。(这对应于一个“众所周知”的事实——我想,在SGA~5口头研讨班里并不曾证明——即\textbf{特征零}时,对代数概形$X$上一个局部常值的$\Lambda$-模层$\mathcal{F}$,其经由$f_!:K.(X,\Lambda)\to K.(e,\Lambda)\simeq K.(\Lambda)$的像等于$d\chi(X)$,其中$d$是$\mathcal{F}$的秩,$e=\mathrm{Spec}(k)$,$k$为假定代数闭的基域……)这立刻提示:Chern同态(1)与(1$'$)应当可以从套套逻辑同态(2)与(2$'$)出发、复合一个“普适”Chern同态(不依赖任何系数环$\Lambda$)
$$ch_X: \mathrm{Cons}(X)\to A(X)$$
而导出;于是,RR公式的两个“带$\Lambda$系数”的版本,便在形式上显现为包含于一条在可构造函数层面上的RR公式之中——后者总以同样的形态书写。

当我们对固定基域(重新允许任意特征)上的概形工作,或更一般地对固定的\textbf{正则}基底概形$S$(比如$S=\mathrm{Spec}(\mathbb{Z})$)上的概形工作时,与惯常写法(自1957年以来熟悉的凝聚框架)最贴合的RR公式形态,经由引入乘积
$$ch_X(x)\, c(X/S) = c_{X/S}(x)$$
而得到(其中$x$在$K.(X,\Lambda)$或$K^.(X,\Lambda)$中皆可);人们可以称之为$x$相对于基$S$的\dmark{}\origpage{365}\textbf{Chern类}。当$x$是$K^.(X,\Lambda)$的单位元、即取值$\Lambda$的常值层的类时,我们得到$X$相对于$S$的相对全Chern类在典范同态$A(X)\to A(X)\otimes K^.(\Lambda)$下的像。这么规定之后,RR公式等价于:对$S$上(有限型的)变动的正则概形$X$、固定$S$,这些相对Chern类的作法
$$c_{X/S}: K.(X,\Lambda)\to A(X)\otimes K.(\Lambda)$$
对真态射是函子的,(5$'$)的变形亦然。特征零时,这化归为相应映射
$$c_{X/S}: \mathrm{Cons}(X)\to A(X)$$
对真态射的函子性。正是在“绝对Chern类”映射(6)的存在与唯一这一形态下——在$S=\mathrm{Spec}(\mathbb{C})$的情形——该猜想呈现于MacPherson的工作之中;相关条件(这里与特征零的一般情形一样)是:a)(6)对真态射的函子性;b)$c_{X/S}(1)=c(X/S)$(在这里,即“绝对”全Chern类)。相对于我最初的猜想,MacPherson所呈献并证明的形态,有两点不同。其一是“减”:他不置于Chow环、而置于整系数上同调环——更确切地说,经由超越途径定义的整系数同调群。其二是“加”——也许正是在这里,Deligne对我最初的猜想做出了一份贡献(除非这份贡献出于MacPherson本人\fnote{102}{(*)}{(1985年3月)确实如此;见前一条脚注所引的注n\textsuperscript{o}~164。}):即对映射(6)的存在与唯一,无须把自己限于正则概形$X$,只要把$A(X)$换成整系数同调群。由此很可能:在一般情形亦然,只要以诺特概形$X$的\textbf{Chow群}(一般不再是环)充当$A(X)$(或更好:$A.(X)$)。或者换个说法:虽然$ch_X(x)$的启发式定义($x$在$K.(X,\Lambda)$或$K^.(X,\Lambda)$中)本质上地利用了环绕概形正则的假设,可一旦把它乘上“乘子”$c(X/S)$(当概形$X$有限型于固定正则概形$S$之上),所得乘积(4)似乎无须对$X$的任何正则性假设仍保有意义\dmark{}\origpage{366}——作为张量积
$$A.(X)\otimes K.(\Lambda) \quad\text{或}\quad A.(X)\otimes K^.(\Lambda)$$
中的元素,其中$A.(X)$记$X$的Chow群。MacPherson的证明精神(它不使用奇点消解)提示:(5$'$)的同态有可能有一个显式的“演算式”构造:“将就着”$X$的奇点本身,连同系数层$\mathcal{F}$(其类为$x$)的奇点,在$X$上“收取”一个取值于$K^.(\Lambda)$的闭链。这也合乎我1957年随凝聚Riemann--Roch定理引入的观念的精神:我在其中尤其做自交计算,小心不让所论的闭链“动”。一个显明的初步归约(经由把$X$嵌入一个$S$-概形)是$X$为正则概形$S$的闭子概形的情形……

此外,应当可能发展一条\textbf{奇异的}(凝聚)Riemann--Roch定理——这个想法我早已熟悉,说不出从何时起;只是我从不曾认真检验过它。某种程度上,正是这个想法(除去与“上同调、同调、卡积”形式主义的类比)曾在SGA~6里(1966/67)使我系统地引入$K.(X)$、$K^.(X)$与$A.(X)$、$A^.(X)$,而不满足于只用$K^.(X)$工作。我不记得1966年在SGA~5研讨班里是否也想过这类东西、是否在口头报告里暗示过。由于我的手写笔记已散失(也许在一次搬家里?),我恐怕永远无从知道了……

(6月7日)翻阅MacPherson的文章,有一件事打动了我:“Riemann--Roch”一词在其中未曾出现——这也正是我没有立刻认出我1966年在SGA~5研讨班里所做那条猜想的缘由;那条猜想,对我来说(现在仍是),是一个“Riemann--Roch”型的定理。写文章之际,MacPherson似乎甚至不曾觉察这层显然的亲缘。我料想缘由在于:Deligne——我离去之后,他以自己乐意的形态使这条猜想流通——曾尽力“抹去”它与Riemann--Roch--Grothendieck定理的显然亲缘。我自以为感到他如此行事的动机。一边,这削弱了这条猜想与我个人的联系,并使它目前流通的那个名称“Deligne--Grothendieck猜想”更可信。(附记:我不知道它在概形情形是否也在流通;倘若在,我很想知道顶着什么名称。)但我看更深的缘由,在于他那个执拗的念头:尽一切可能否认并摧毁我的著作、我的数学视野的根基处的统一性\fnote{103}{(*)}{可与“遗骸”一注(n\textsuperscript{o}~88)里关于“SGA~4½之役”之深层意味的评语对读:那场战役同样意在,经由外来文本SGA~4½对SGA~4与SGA~5这两个不可分割的部分——它们正发展着我的著作——的“强行塞入”,把我的著作围绕平展上同调的深层统一,炸成一堆无定形的“技术枝蔓”。}。这是一个触目的例子,说明在一个手段出众的数学家身上,一个与一切数学动机全然无关的执念,竟能使那种我称之为数学的“健全本能”的东西昏暗(乃至全然堵塞)。这份本能不可能不察觉一个“同一条”Riemann--Roch定理的“连续”与“离散”两个陈述之间的类比——我当年在口头报告里当然也强调过。如我昨天所写,这份亲缘想必不久将由一个成形的陈述(Zoghman~Mebkhout所猜)确认——至少在复解析情形——使人得以从一条共同的陈述导出两者。显然,在Deligne对着Riemann--Roch定理的那份“掘墓”心境里\fnote{104}{(**)}{这些心境——恰对着Riemann--Roch--Grothendieck定理——在“葬礼悼词”里表现得格外清楚;见“葬礼悼词(1)——或恭维”一注,n\textsuperscript{o}~104。},他根本不可能发现那个在解析框架里联结两者的唯一陈述,更无从在一般概形框架里提出类似陈述的问题。一如他在这样的心境里,无从提炼D-模这个丰产的观点——它以过于自然的方式,出自他正要埋葬的观念——甚至在多年之间,无从承认Mebkhout那项在他失败之处成功的丰产工作。

\notehead{87$_{2}$}{}

(5月31日)那是1957年——我的Bourbaki报告论L函数有理性之年;我在其中启发式地使用Verdier的结果(???)(尤其该个案里局部项的预期形态),而不等Illusie十三年后——应Deligne之邀——才肯去证明它。而且我觉得,当年Verdier把他那条超\dmark{}\origpage{368}一般的公式拿给我看——它来得如一场意外——时,他是用“六种运算”形式主义在几行里证明它的:这类公式,(几乎)写出来就是证明了它!倘若有什么“困难”,至多在核验一两处相容性的层面上\fnote{105}{(*)}{(6月6日)此外似乎:经由双对偶定理(其间已晋升为“Deligne定理”),Lefschetz--Verdier公式的初始证明曾依赖一条奇点消解假设;Deligne得以在一域上有限型概形的情形摆脱它。这正是浑水摸鱼的大好机会,给人这样的印象:SGA~5竟从属于那个“在先(原文如此)”的“研讨班SGA~4½”(而且它确实先出版了!)。}。况且,无论Illusie还是Deligne,都一清二楚:我在研讨班里为各种显式迹公式所给的证明\textbf{是完整的};它们绝不依赖Verdier那条一般公式——后者只扮演了“扳机”的角色,怂恿人们在尽可能一般的情形里显豁并证明迹公式。两人在这里的不诚,昭然若揭。对Deligne,我写“白板”一注(n\textsuperscript{o}~67)时就已清楚——但对不知情的读者,当然不清楚;同样不清楚的,还有那种放弃使用自己健全官能的知情读者。

(6月6日)至于Illusie,他完全入了其朋友的戏,试图搅浑水,摆出一副超技术的口头研讨班的样子——仿佛它连全部结果的完整证明都给不出,尤其迹公式。然而那些公式在65/66年间\textbf{确实}得到证明(而且是头一回),恰恰是在那里,他与Deligne有幸学得它们,以及与之相随的一整套精微技术\fnote{106}{(**)}{在顶着SGA~5之名出版的那一卷的导言第二段里,Illusie把围绕平展上同调Lefschetz公式的三场报告III、III~B、XII当作“研讨班的心脏”;而人们看到,在报告III~B的导言里,他小心注明(与事实相反)“这份报告不对应于研讨班的任何一场口头报告”;在报告III与III~B的导言里,他尽力给人这样的印象:它们从属于SGA~4½,而报告III被呈献为“猜想性的”!!实际上,SGA~5研讨班的全体,在技术上独立于报告III(Lefschetz--Verdier公式)——后者扮演启发性动机的角色;而报告III~B不是别的,正是Bucur那次搬家所造成的“窟窿”(报告XI);那窟窿,成了这场进一步肢解的可喜口实。为使人相信那场“技术枝蔓”研讨班的版本(其朋友Deligne的授意),Illusie小心砍掉了开篇报告——我在其中为将在这场研讨班里发展的各大主题画了一幅初步的全景;迹公式在其中只占一小角(因其朝Weil猜想方向的算术含义而取得特殊重要性)。关于这些“大主题”的概观,见后面的子注n\textsuperscript{o}~87$_{4}$(原书误作87$_{5}$)。}。

\dmark{}这使我想起:当然,我曾费心在研讨班里证明Lefschetz--Verdier公式\origpage{369}——那是最起码的,也是我打算发展的局部与整体对偶形式主义的一次格外触目的应用。这几天问题来到我心头:既然分明有十来份报告的成文工作,在我亲爱的学生们手里陷于困境,Deligne与Illusie要为其SGA~5出版的“障碍(原文如此)”命名,本有挑选的余地——为什么偏偏选中他们的好哥们Verdier的定理——他当时正把父权当作应得之物佩戴着,一如他佩戴导出范畴与三角范畴的父权,而那些他也从不曾费心成文(或至少交到公众手里)。这里有一种\textbf{荒谬之中的挑衅}(或我的前上同调学家学生群体里的某种集体犬儒——我认为他们在那场屠杀战役里全体同谋),它使我想起Verdier前一年光辉地发明的“权复形”(见同名注,n\textsuperscript{o}~83),或(在不义的音区里)Deligne给本该叫“Mebkhout的层”的那些层所取的名字“反常”(见“反常性”一注,n\textsuperscript{o}~76)。我在这样的发明里感到,是对整个数学共同体的支配与轻蔑之举——同时是一份\textbf{赌注};它分明一直赢到亡者不期现身之时——他俨然像一群沉睡者中间唯一醒着的人……

\notehead{87$_{3}$}{}

(6月5日)一场屠杀盘点之后,人们自会掂量Illusie在其顶着SGA~5之名的卷子的导言第2行里的这份自白:

\begin{quote}
“相对于最初的版本,唯一重要的变更涉及报告II[泛Künneth公式],未予复制;以及报告III[Lefschetz--Verdier公式],已全部重写,并增补一个编号III~B的附录\fnote{107}{(*)}{它被呈献为“研讨班的心脏”的一部分!(见前一条脚注。)}。除若干\dmark{}\origpage{370}细节修改与脚注增补之外,其余报告一仍\textbf{其旧}”(着重号是我加的)。
\end{quote}

这里,Illusie又一次甘当其那位罄竹难书的朋友另一句俏皮话的应声虫,即:SGA~4½的存在“不久将使人得以\textbf{照原样}出版SGA~5”(见“白板”一注,n\textsuperscript{o}~67)——而Illusie在其报告与导言里,处处尽力为这一冒名张目——即:SGA~5(他与朋友学艺之地)竟依赖于那部海盗文集SGA~4½,由其后十二年间拾得或劫掠的破烂与断枝拼成——每页转弯处都不吝对SGA~4½的大量引证……

最后的字,照例归Deligne:一个月前(5月3日)他给我写信,回复一份简短的信息之请(见“葬礼”一注,n\textsuperscript{o}~70的开头):

\begin{quote}
“总之,倘若SGA~4½这份文本问世时,你已七年不做数学[?!],那只是相应于SGA~5出版[?]的漫长迟延,\textbf{它太不完整,无法照原样有益地出版}。

希望这些解释令你满意。”
\end{quote}

它们即便不曾“令我满意”,至少令我大开眼界……

\notehead{87$_{4}$}{}

(6月6日)或许是时候指出:口头研讨班里发展过的主要主题有哪些;已刊文本,只有靠旁证才能对其形成概念。

I)对偶理论的局部侧面;其本质技术配料是(一如凝聚情形)双对偶定理(补以一条“上同调纯粹性”定理)。我的印象是:这后一条定理的几何意味——作为一条局部Poincaré对偶定理——我当年在口头研讨班里分明解释过,却被我的学生们从此全然忘却\fnote{108}{(*)}{核实之后:这一几何诠释,至少在Illusie的成文里得到了保留。}。

II)迹公式,包括比惯常迹公式(其两边为整数,或一般为某系数环——如$\underline{Z}/n\underline{Z}$或$\ell$-进环$\mathbb{Z}_\ell$、乃至$\mathbb{Q}_\ell$——的元素)更精微的“非交换”迹公式;它们置于一作用在所论概形上的\textbf{有限}群的代数之中,系数取于一个合宜的环(如前一括号里所设想者)。这一推广来得十分自然:因为即便在惯常型的Lefschetz公式里,对“扭转的”系数层,人们也被引至以一个Galois覆盖(一般分歧)替换原概形,以“解扭”系数,而Galois群作用于其上。如此,“Nielsen--Wecken”型的公式,便自然地进入概形语境。

III)Euler--Poincaré公式。一边,借助Serre--Swan模,对代数曲线的一条“绝对”公式作了周详研究(推广了驯分歧系数的情形;该情形给出较朴素的Ogg--Chafarévitch--Grothendieck公式)。另一边,有未发表的、深刻的“离散”Riemann--Roch型猜想;其中一条,七年后以一种杂糅的版本重现,顶着“Deligne--Grothendieck猜想”之名,由MacPherson经超越途径证明(见注n\textsuperscript{o}~87$_{1}$)。

我对这两大主题(Lefschetz公式、Euler--Poincaré公式)之间深刻关系所能做的评语,也一并散失,未留痕迹。(一如我的习惯,我把全部手写笔记留给了“志愿者们-执笔者(原文如此)”;口头研讨班的手写笔记,我本来有一整套——尽管有些简略——如今片纸无存。)

IV)与一个闭链相配的同调类与上同调类形式主义的详论,自然出自一般对偶形式主义与那条关键想法:在所论闭链上做“具支撑”的上同调,利用上同调纯粹性定理。

V)有限性定理(包括泛型有限性定理)与任意支撑上同调的泛Künneth定理。

研讨班还发展了一套从挠系数过渡到$\ell$-进系数的技术(报告V与VI)。那是研讨班最技术的部分;研讨班一般用挠系数工作,事后再“取极限”以导出相应的$\ell$-进结果。这一观点是一个临时的\dmark{}\origpage{372}权宜之计,且待Jouanolou的论文(至今未出版!)直接在$\ell$-进框架里给出所需的形式主义。

我不把对几个经典概形的一些计算、以及Chern类的上同调理论,算进“主要主题”;Illusie在其导言里把它们捧为研讨班“最有趣的”之一。由于日程已满,我在口头研讨班里不曾认为有必要在这些计算与这一构造上流连:只需把十年前我为Riemann--Roch定理、在Chow环的语境里给出的论证,几乎一字不差地搬来即可。另一方面,明摆着应当把它们纳入成文的研讨班,以便给平展上同调的使用者一份便利的参考。Jouanolou揽下了这项工作(报告VIII);他本该把它看作一项美差——既服务于数学共同体,又为自己学到一些日后必用的基本技术——却把它当作苦役:其成文一拖数年\fnote{109}{(*)}{(6月12日)翻阅那份报告,我进而确信:Jouanolou与我的其他上同调学家学生之间,有完美的默契。}。想必他的论文也如出一辙:至今仍是一份幽灵参考,一如Verdier的那份……“极限过渡”那一部分也不该算作研讨班的“主要主题”之一,因为它不与某个特定的几何想法相连。毋宁说,它反映平展上同调语境(区别于超越语境)特有的一项技术复杂:平展上同调的主要定理首先关涉挠系数(与剩余特征互素);要有对应于特征零系数环的理论——Weil猜想所必需——必须对$\mathbb{Z}/\ell^n\mathbb{Z}$型系数环取极限,以获得“$\ell$-进”结果。

把这一切说清之后,口头研讨班五个主要主题里,唯一一个似乎以完整形态出现在已刊文本里的,是主题I。主题IV与V干脆消失,被SGA~4½吸收,附带的收益是可以大量引证它、给人这样的印象:SGA~5依赖一份自居在先的Deligne文本。主题II与III以残缺形态出现在已刊卷里,并始终维持对SGA~4½文本的同一冒名(那份文本,实际上整个出自母研讨班SGA~4、SGA~5)。

\subsection{遗骸……}

\notehead{88}{}

\dmark{}(5月16日)SGA~4与SGA~5这两场前后相续的研讨班——对我而言,它们如\textbf{同一场}“研讨班”——从虚无中同时发展出:topos的\textbf{语言}——强大的综合与发现的工具;以及平展上同调——一件打磨停当、效率完美的\textbf{工具};其本质形式性质,彼时甚至比寻常空间的上同调理论理解得还透\fnote{110}{(*)}{即便把自己限于最贴近“簇”的空间,如可三角剖分的空间。}\origpage{373}。这一整体,代表我在数学里、就彻底完成的工作而言所能做出的最深刻、最革新的贡献。同时——它并无意于此,虽然一切始终以事情之明摆着的自然进行——这项工作代表我作为数学家在其著作中所完成的规模最大的技术壮举\fnote{111}{(**)}{若干困难或意外的结果,是由别人(Artin、Verdier、Giraud、Deligne)得到的;工作的某些部分是与别人合作完成的。这丝毫不减损(至少在我心里)我对这项工作在我全部著作中之地位的评判的分量。况且,我想在《主题概要》的一篇附录里,更周详地回到这一点,在分明已属必要之处把点补在i上。}。这两场研讨班,对我是密不可分的。它们在其统一之中,同时代表\textbf{视野}与\textbf{工具}——topos,与一套完整的平展上同调形式主义。

视野至今仍遭拒斥,而工具近二十年来深刻革新了代数几何其最令我心醉的一面——“算术”的一面:经由一种直觉、一套“几何”性质的概念与技术行装所把握的算术。

肯定不只是为了暗示其上同调“摘要”先于SGA~5部分,Deligne才给它戴上SGA~4½这个欺眼之名——反正要戴,大可以叫SGA~3½!在“SGA~4½之役”里,我感到一种意图:把那部他的全部著作皆出于斯的著作(那部他始终无法摆脱的著作!)——其分明而深层的统一,在SGA~4与(真正的)SGA~5两场研讨班的整体里昭然可见——呈献为一副\textbf{分裂}的面目(一如他自己也是分裂的……),被一份外来的、轻蔑的文本的强行塞入\textbf{一分为二};那份文本,妄想显出是活的心脏、思想与\dmark{}\origpage{374}视野的精髓——而它在那视野里不曾有过任何份额\fnote{112}{(*)}{这份思想——就其主导观念与本质结果而言——早在年轻的Deligne登台、来与我接触学代数几何与上同调技术(1965至1969年)之前,就已完全成熟。(5月30日)对此可参见“格格不入者”一注,n\textsuperscript{o}~67$'$。};而环绕它的两个“断片”,则成了近乎怪诞的附录,一堆对那部自居为中心与本质的著作——出自Deligne的笔下,而我寒微的身量则被客气地接纳(在彻底下葬之前)跻身“合作者”之列——的“枝蔓”与“技术补充”\fnote{113}{(**)}{见“绿灯”与“颠倒”两注,n\textsuperscript{o}~68、68$'$。}。

“偶然”把事情安排得妙极。这副“任人宰割的遗骸”——那场被“执笔者们”一直弃置不顾的“倒霉研讨班”,我离去时落在我的上同调学家学生们手里、听凭其处置的——那可不是老师著作中\textbf{随便}哪一部分!既不是SGA~1与SGA~2(我在其中独自发展、尚不自知那些即将成为未来主作“起飞”的两件不可或缺的技术辅助工具的东西),也不是SGA~3(我的贡献主要是为使概形的“四面八方”技术纯熟而做的无穷音阶与琶音——有时颇费力气),也不是SGA~6(系统发展我围绕Riemann--Roch定理与相交形式主义的十年旧观念),乃至SGA~7(按一场反思的内在逻辑,它从对中心工具的持有、对上同调的掌控中流出)。分明是我著作的\textbf{主干部分}——其成文工作尚未完成(而且是经他们的手……)——我把它,至少部分地,留在了我的上同调学家学生们手里。他们选择屠杀的,正是著作的这一主干;他们据为己有的碎块,忘了那使其获得意味、美与创造之德的那份统一(注90)。

同样不是偶然:他们手执七拼八凑的工具、否认使其从虚无中诞生的精神与视野,竟无一人能在革新之著作重生之处——逆着他们的冷漠与轻蔑——辨认出它。也无从解释:六年之后,当那件新工具终于被Deligne把握,他们竟全体一致地埋葬了那位在孤独中创造它的人——Zoghman~Mebkhout\dmark{}\origpage{375},被否认的老师之遗孤学生!同样不是偶然:在Deligne最初的冲劲回落之后(那冲劲,在几年间把他带向一套新Hodge理论的强力启动,与Weil猜想的证明),尽管他手段惊人、我的上同调学家学生们手段出众,我今天仍确认这场“阴郁停滞”——在一个无比丰饶、一切似乎尚待去做的领域里。这无须惊讶:近十五年来,灵感的主要来源与某些“大问题”\fnote{114}{(*)}{这份“灵感的主要来源”,当然是“Motif瑜伽”。它只在Deligne一人身上起过作用——他把它留在怀里独享“收益”,且以一种被剥除其大半力量的狭隘形态,拒斥这部瑜伽的某些本质方面。在由它激发、却遭忽视或被悄悄抹黑的“大问题”之中,我如今(虽不折不扣是局外人)看到:标准猜想;以及面向一切惯常类型的系数——或多或少贴近“Motif”本身(它们对后者扮演“普适”系数的角色:生出所有其他系数的那种)——的“六种运算”形式主义的发展。可与“我的孤儿们”一注(n\textsuperscript{o}~46)里的评语对读。},即便它们在场、人们每走一步都与它们遭逢,却被小心绕开、藏掖起来,如那位十五年来人们不断要埋葬的人的信使一般。

\subsection{……与身体}

\notehead{89}{}

(5月17日)那份曾在我身上活着的念头、那种事物的视野——我曾以为已把它传递出去——我看作一个活的身体:健全而和谐,带着活物更新的力量、构想与生育的力量。瞧,这个活的身体变成了一副\textbf{遗骸},在人们中间被瓜分——这条腿或这半边身子,被妥善制成标本,在某一家充当战利品;另一条,被剁碎,在另一家充当击头的棍棒或回旋镖;还有一条,天知道,原样下到自家的汤锅里(我们已经不在乎这个了!)——其余的,统统可以烂在垃圾堆里……

这就是——形象化的措辞诚然,但在我看来,恰当地表达事情的某种实在——终于向我显现的图景。那根棍棒,至多会在这里或那里敲裂一两颗头颅\fnote{115}{(**)}{(5月31日)而且它还满可以用来证明某条“困难尽人皆知的”定理!};但这些四散的碎块——战利品也罢、棍棒也罢、家常汤也罢——将永远不具那个活的身体那份如此简单、如此分明的力量:那创造一个新生命的、爱的拥抱的力量……

(5月18日)这个活的身体、这副碎块随风四散的“遗骸”的意象,想必是在刚过去的一星期里在我身上渐渐成形的。\dmark{}它在我的打字机笔下呈现的那副滑稽形态,丝毫不意味着这意象是半点\textbf{虚构}、一场顺着话头即兴而来的略带阴森的滑稽表演。这意象表达一份\textbf{实在};当它经书面表述取得物质形态的那一刻,我被深深地触动。\origpage{376}这份实在,想必早在我“离去”以来的十四年里、甚至更早,已在此处彼处以断片的方式被我获知。一些信息的断片,起初由一份心不在焉、倾注别处的注意力,登记在一个表浅的层面——但它们全都指向同一方向,想必已在更深的层面,聚成某种意象——一幅未曾表述的意象,我不曾费心去认识它,那时我有更要紧的事。这份意象,在3月底以来持续进行的反思中、即这六七个星期里,大大丰富并明确了。更确切地说:零散的信息要素,终于经一份全然在场的、有意识的留心的检视,渐渐聚成\textbf{另一幅}意象——在检视与探查的思维这个较表浅的层面上;经由一项工作,它看似独立于第一幅意象在更深层的在场。这项有意识的工作,六天前在“屠杀”之发生的骤然视野中达到顶点——那时我感到了一份\textbf{暴力}的“气息”、“气味”;我想,这是整场反思里的头一回\fnote{116}{(*)}{(6月12日)近些年里,我曾有过这样的时刻:感到我的某些前学生对另一些与我“同葬者”怀有暴力的意图;但从不曾感到一种出自集体意志(在这里结成五人)、指向我这个人、经由我的著作而行的暴力。}。也是在这样的时刻,一个健全和谐之躯分明遭“屠杀”的感觉,想必已在临近表面的层次出现;而那幅更深的弥散意象,想必也已开始浮出水面,为那幅成形中的意象带来一层血肉——一种单凭思维无力给出的“气味”。

这份“血肉”的侧面,昨夜又在一场梦里显现——正是受这场梦的推动,我此刻回到昨天写下的文字。梦里,我身体的几处地方被割了相当深的口子。起初是嘴唇与口腔里的割口,大量出血;我对着镜子用大量清水漱口,漱口水被血染得通红。然后是腹部的伤口,也大量出血,尤其其中一处,血一鼓一鼓地涌出,仿佛是动脉(做梦者无心顾及解剖学的写实)。我甚至想到:倘若照这样流下去,我兴许会躺倒在地的;我用手按住伤口,蜷起身子止血——血确实止住了奔流,渐渐结成一个血块与一大片痂。后来,我小心地揭开那片痂:细微的愈合已经开始。我的一个手指也被割了,上面缠着厚厚一团绷带……

我无意对这个梦做更精微、更周详的描述,也不打算在这里(或别处)深探它。这个梦“照原样”已经以惊人的力量向我显示:昨天我谈的那个“身体”——我在写它时,把它看作与我分离的东西,像一个我或许孕育、生出、如今已走向世界走自己路的——这个身体,至今仍是我这个人亲切的一部分:那是\textbf{我的}身体,由血肉与一种生命的力量构成,这力量使它经得起深的创伤、并自我再生。而我的身体,想必也是世上与我系得最深、最不可解的东西……

做梦者没有让我停留于“屠杀”与瓜分遗骸的意象。那份意象要还原的,是我曾强烈知觉到的、\textbf{他人}身上的一份意图与心境的实在,而不是我自己经历这场侵凌、这场残损——经由一件我仍与之紧紧相连的东西——的方式。我与之相系到何等地步,做梦者刚刚让我隐约瞥见。这与我在“事物的回归(或一脚踩进菜盘)”一注(n\textsuperscript{o}~73)的反思里(以较弱的力度)所知觉的东西相合:在那里,我试着稍稍勾勒“构想了一件事的人与这件事之间的深深纽带”的那份感觉——它出现在那天的反思里。在那场4月30日的反思(仅仅三周之前)之前,在我的一生里,我曾装作无视这条纽带,或至少淡化它,顺着时兴老生常谈那条划好的坡道。操心一件已脱手的作品的命运——尤其操心我们的名字是否还多少附着于它——被感受为一种小家子气、一种悭吝;然而,人人却觉得自然:一个我们养大的(并且自以为爱过的)亲生孩子,倘若选择否认他出生时所得的名字,我们会深受触动。

\subsection{继承人}

\notehead{90}{}

\dmark{}(5月18日)我不知道六十年代里,可曾有哪个学生(Deligne除外)感到了这份本质的统一——超越他与我一起进行的有限工作。也许某些人模糊地感到过;而这份知觉,在我离去之后的年月一过,便永无挽回地失落了。可以肯定的却是:从1965年我们初次接触起,Deligne就预感到了这份活的统一。\origpage{378}正是这份对一个宏阔蓝图之统一旨趣的精微知觉,无疑是他身上那份强烈兴趣的主要兴奋剂——对着我有可传授、可传递的一切。这份兴趣,在1965到1969年这四年持续的数学接触里,从不曾衰减\fnote{117}{(*)}{这一时期包括五年,其中一年(1966),我的朋友在比利时服兵役。}。它赋予我们之间的数学交流以我谈过的那份卓越的品质;与别的数学家朋友,我只是在寥寥的瞬间知晓过它。正是这份对本质之物的知觉,以及它在他身上激起的那份炽热的兴趣,使他得以仿佛游戏般学得我能教他的一切:既有技术手段(锌丝概形的技术,Riemann--Roch与相交的瑜伽,上同调形式主义,平展上同调,topos的语言),也有使其统一的总体\textbf{视野},最后是\textbf{Motif瑜伽}——它当时是这份视野的主要果实,是我到那时为止所发现的最强的灵感之源。

清楚的是:Deligne是我的学生中直到今天仍唯一的一个,曾在某个时刻(我想自1968年起)把我可传递的一切,就其本质的统一与其手段的多样性,完全消化并据为己有\fnote{118}{(**)}{当我谈“全部”时,应理解为:就视野与手段中的本质事物而言。这当然不意味着,不曾有一些未发表的想法与结果,我从未想过向他谈起。相反,我不认为1965--69年间的任何一场数学反思,是我没有“趁热”向我的朋友谈过的——始终怀着乐趣,也始终有所获益。}。正是这份情势——想必人人都感到了——使他俨然我著作天造地设的“合法继承人”。显然,这份遗产既不使他累赘,也不限制他——它不是一副重担,而给他添翼;我是说:它以其活力滋养着他生而有之的那些“翅膀”\dmark{}\origpage{379}——一如别的视野、别的遗产(诚然没那么个人……)也将滋养它们……

这份他曾在成长与起飞的关键年月里赖以为养的遗产,以及使其具有美与创造之德的那份统一——他曾那么懂得去感觉、它已成为他自身一部分的东西——我的朋友此后\fnote{119}{(*)}{奇怪的是,这份分裂,想必在我们相识的头一年就已在场(它已经表现为对SGA~5研讨班的一份暧昧态度——那是他与概形、与Grothendieck风格的上同调技术、与平展上同调的初次接触),最迟在1968年、以毫不含糊的形态(见“驱逐”一注,n\textsuperscript{o}~63)——也就是数学交流完美无缺、他的数学思想的起飞在我看来尚未打上冲突印记的时刻。他当时(“顺带地”)做出过许多有趣的贡献(我在SGA~4的导言里极乐于把它们捧出来),题材却是他在我离去之后竭力要埋葬的。}予以否认,不懈地藏掖遗产,否认并摧毁那份作为其灵魂的创造性的统一。在我的学生中间,他是第一个做出榜样的人:一面据有工具、据有“碎块”,一面拼命肢解它们所自出的统一、那个活的身体。他自己的创造冲劲,由此被他身上这份深层的分裂所阻滞、吸收、最终肢解;这分裂驱使他否认并摧毁恰恰构成其力量、滋养其冲劲的东西。

我看这份分裂经由三个互相支援、密不可分的效果表达。其一是能量涣散的效果:在否认、肢解、取代、藏掖的努力中四散。其二在于对某些观念与手段的拒斥——它们对本是他选作中心主题的那个主题的“自然”发展本质不可少\fnote{120}{(**)}{这份拒斥,尤其表现为:导出范畴与三角范畴被埋葬(直到1981年),六种变换的形式主义被埋葬(直到今天),topos的语言被埋葬(亦然);以及对同调代数与同伦代数那套宏阔奠基纲领的一种“以轻蔑来封锁”——那纲领,我现在(二十年后)试图以《追寻栈》给出一个草图,而他当然同样不曾不感到其需要。最后,即便他从Motif瑜伽(被埋葬直到1982年)汲取灵感,这部瑜伽仍被削去其一部分力量:它与六种变换的形式主义——其一个本质的形式侧面——脱了钩。依我看,这一侧面,也被严格地逐出Hodge--Deligne理论之外。}。其三,是执著于那个首当其冲的主题:在其中,要取代、排挤一位每走一步都在场的老师、并无休止地抹去他——恰恰是这个主题,被注入最强烈的那份矛盾,它主宰其数学家的一生\dmark{}\origpage{380}。

我第一手所知的东西,加上一份从不曾骗过我的素朴本能或嗅觉,使我清楚:倘若Deligne不曾被这份其工作内部的深层矛盾撕裂,今天的数学便不会是如今这副样子\fnote{121}{(*)}{写“今天的数学”这些文字时,我想的不只是我们今天对数学事物或多或少的认识。背景里还有对一个特定\textbf{精神}的念头——在数学家的世界里,尤其在被人们(不带讥讽或嘲弄的口吻)称为“上流”数学界的东西里:那个“定调子”、裁定何为“重要”乃至“正当”、何为不然的精神;那个控制信息手段、并在很大程度上控制前程的精神。也许我夸大了身处旗手之位的一个人对给定环境在给定时代的“时代精神”所能有的影响。Deligne的影响,依我看,堪比(无论好坏)二十年前Weil在接待我的那个环境里所显得具有的影响——那个环境,我曾二十年认同于它。(5月31日)可与“掘墓人——或全体教众”一注(n\textsuperscript{o}~97)的(互补)反思对读。}:它本会在其若干本质的部分,经历宏阔的革新——一如我曾亲身充当其主要工具的那一场革新——恰恰是同一位Deligne拼命对抗、使之改道的那一场!\fnote{122}{(**)}{(6月16日)我深信:单凭我引入数学的主导观念正常发展这一点——沿着六十年代已获得的冲劲(那冲劲,被“电锯效应”猛然切断,下两条注将谈到它……)——今天的数学,在我离去十五年后,本会在其若干本质的部分,不同于如今的样子……}

同样毫无疑问:他天造地设,本该成为一所强力几何学派的灵魂——延续那所曾围绕我形成的学派:一所汲取其母体活力、以及接我班者之创造力的学派。但这所曾围绕我形成的学派、这座曾环绕那些紧张的受业年月的滋养母体,在我离去后的翌日便已解体。所以如此,恰恰是因为在分明接我班的那个人身上\fnote{123}{(***)}{这份\textbf{既成事实的接班},以毫不含糊的具体征兆表达:他在IHES接了我的班(我是在他入所的次年离开的——见“驱逐”一注,n\textsuperscript{o}~63);并且他接手了代数簇上同调这一中心主题,用的正是我为这一目的、在十五年间(1955至1970年)发展的手段。},找不到那个也将是一场共同冒险所结合的群体的灵魂的人——为一项其规模超出每个人的手段的任务。

我的印象是:我离去之后,我的每一位学生都退回自己的角落——工作自然有的是,数学里从不缺工作——但这个“角落”不再嵌入一个整体,这份“工作”不再被一股潮流、一个更宏阔的旨趣所承载。肯定,自我离去起——甚至更早——我的大部分学生或前学生的目光,都转向那位天造地设的“继承者”:他们中间最光彩、也离我最近的一位。在这个敏感的时刻,我的朋友想必平生头一回感到那种对他人的权力突然握在自己手里:那种对一个学派的命运生杀予夺的权力——那学派,他出身于斯;与他同学四年的朋友们,想必期待他保证其延续。局面整个在他手里,是他来定调子……他也确实定了——经由摧毁遗产;而首先摧毁的,是那份信任与那份期待\fnote{124}{(*)}{(5月26日)在反思的进程中,我还察觉对我的默示继承人的另一份全然不同的“期待”——这一回不出自我的学生们,而出自“全体教众”——对此可参见“掘墓人——或全体教众”一注(n\textsuperscript{o}~97)的结尾。我几乎不怀疑:这两份方向相反的期待——一份系于一个非常特定的时刻,另一份贯穿葬礼的十四年——都是实在的。更有甚者,我倾向于认为:在我的不止一位昔日学生身上,两份期待想必同时在场:一面是期待在他们的最光彩者身上,也找到那保证一所学派、一部著作——他们在其中有位置、有份额——之延续的人;一面是期待看到(倘若可能)那个人的任何痕迹被抹去——他的离去,曾在划好的坡道的安稳之中,以如此的力量,突然向他们发问……}——那些曾与他同学于同一位老师的人,本不会不带给他的……

当然,许多人为Deligne的著作所折服,并非没有理由。但我也清楚:这部著作,在其令人惊叹的最初冲劲(以Weil猜想的证明告终)之外,离使出“全副本领”还很远。诚然,它见证一种少有的技术驾驭与从容,把他置于“最好者”之列。但它没有我年少时在他身上知觉到的那份寒微的德性——革新的德性。他身上曾怀有的这份德性、那份幼童的鲜活或天真,久已深深埋没、遭否认。我本要写:凭着这份“德性”与其不凡的天赋\dmark{}\origpage{382},也凭着他的天赋之舒展所曾受惠的那些特殊情势,Deligne本被召唤去“支配”我们时代的数学,一如Riemann、一如Hilbert,各支配其时代的数学。一些扎根于日常语言的积习思路,向我提示了“支配”这一意象;但它给出的,是对实在的一幅失真的把握。这些伟人当然曾完完全全“把握”“消化”“据为己有”其时代的已知数学;这无疑也给了他们一份对技术手段的出众驾驭。但倘若他们名副其实地显得“伟大”,那不是凭其技术的壮举——从一块顽石上“硬凿”艰难的证明。而是凭各自在数学的几个重要部分所带来的革新,经由简单而丰产的“想法”——也就是说:因为曾把目光投向那些简单而本质的事物——那些在他们之前无人肯垂顾的事物。那种\textbf{看见}简单而本质之事物的孩童般的能力——无论它们多么寒微、多么为众人所不屑——革新的力量、每个人身上创造的力量,正寓于其\textbf{中}。这份力量,在我认识的那个年轻人身上——一个无人识得、谦卑而热忱地爱着数学的人——以一种罕见的程度在场。岁月流转,这份寒微的“力量”似乎已从那位受人仰慕与畏惧、无碍地享有其威望、以及威望赋予他的(有时是独断的)对他人之权力的数学家身上,消失了。

这份在我朋友身上对一件极精微、极活泼、为众人所忽视而有创造之力的东西的\textbf{窒息},自我离去以来我已感到过多次;近年来尤甚。但需要这些星期的发现,以及我自3月底以来持续的反思(循着《收获与播种》的冲劲),才开始在其全部幅度上感到这份窒息在我朋友生活中的毁灭性效果——也在许多我曾亲近认识的其他人身上。不只对我“1970以后”的学生(及类似者)——他们领教过他的恶意(某些情形里,也许是无意识的),它对着每一个人施展,沉重地压在三人身上;而且——我现在隐约看到——也在我“1970以前”的学生中间:经由一道\textbf{旨趣上的连续性}的毁灭,以及整体感、统一感的毁灭——那份本赋予其工作以比一堆署着他们名字的抽印本的堆积更深、更广的意味的东西(注91)\fnote{125}{(*)}{(6月16日)这第二个侧面,只是在葬礼的反思中才向我显形。倘若我有幸见到一位显赫的数学家动用“令人气馁的权力”,那恰恰是在不久前还俨然我天造地设继承人的那位身上。写“令人气馁的权力”一节时,我想到他有过许多(在反思转回我自己之前),却还不曾有一丝疑心(至少在意识的层面上):这份权力,竟在那些人中间找到施展的机会——对他们,他本该(一如当年对我)扮演完美数学家的典范……}。

最近七年里,不止一次;最近这些星期、这些天里,又是不止一次,我在一种于某个层面上被感受为巨大的\textbf{糟蹋}的东西面前,感到一阵悲伤——当人身上与别人身上最珍贵的东西,仿佛刻意地被挥霍或窒息。然而,我终究也懂得:这样的“糟蹋”是人的境况的一个基调,它以这样或那样的形态,无处不在——在最寒微与最显赫的人们的生活里,一如在各族各邦的生活里。这份“糟蹋”——它不是别的,正是冲突、分裂在每个人的生活中的行动——是一种我刚刚才开始探其边的丰富而深刻的实体,一份属于我去“吃”去消化的食粮。由此,这份糟蹋,以及我在每一步遇到的别的糟蹋、以及每在路转弯处来到我面前、常常不合时宜的任何事物——这份糟蹋与这些不合时宜之物,身上都带着一份\textbf{益处}。倘若禅观有意义、有革新之力,那正是在这样的限度内:它使我得以领受那份(经我积习的折光)显为“为害”之物的益处,使我得以\textbf{以}那看似为毁灭而造的东西为食。

以自己的经历为食,任它更新自己而不总是绕开它——这就是充分承担自己的生活。我身上有这份力量;每一刻,用它或弃之于闲置,都由我。对我的朋友皮埃尔,对曾是我学生的每一个人,亦然——同我一样自由,以这份“糟蹋”为食;这份糟蹋,我在这场漫长禅观的最后几天里,正巡视完毕。对读到这些为读者而写的文字的人,亦然。

\subsection{共同继承人}

\notehead{91}{}

\dmark{}(5月19日)从四面八方传到我这里的、关于我昔日学生们的回音,稀薄之至。\origpage{384}我离去之后,几乎没有人肯给我一点活着的音讯——哪怕寄一份抽印本\fnote{126}{(*)}{(5月31日)对此可参见“沉默”一注之后的注n\textsuperscript{o}~84\textsuperscript{1}。}。然而,把我得到的这点东西聚拢起来,我可以形成一个概念——诚然极其粗略。也许在随后的几个月里它会明确起来,倘若这场反思促使他们中某些人现身。

我已有机会确认Deligne著作在我离去之后的深刻断裂,尽管在若干方面,他违背本意地显出接班人的样子,因而显出某种连续性。我曾感到这道断裂,必定深深波及我其余所有学生的工作。这正是我想稍稍接近地勾勒的印象。

这些学生中唯一一个其工作显然(至少乍看)延续了他与我一起所做的工作的,似乎是Berthelot\fnote{127}{(**)}{从Verdier在我离去后几年间所追随的对偶主题——在一个与我发展它的语境相近的解析语境——看,有一种如Berthelot情形的连续性印象。但我觉得,那多少是一种“例行公事的连续性”;而我主要寻找其征兆(或征兆之缺席)的,是一种创造性的连续——在未知中延续一股最初的冲劲……}。他也是长年来唯一不断给我寄许多抽印本的——甚至也许全部抽印本。它们都处在晶体上同调这个艰难的题目里;其系统的启动,是他论文的对象。然而我觉得,一如我其他“上同调学家”(交换的)学生,其著作同样带着对我引入的某些主要观念的冷淡:导出范畴(及Verdier提炼的三角范畴),六种运算的形式主义,topos(注91\textsuperscript{1})。如Zoghman~Mebkhout本人所说,他自己的著作——题材与Berthelot的如此贴近(注91\textsuperscript{2})——正处在这条线上:这些观念,连同Sato学派的观念。倘若它们不曾被我的上同调学家学生们——Deligne与Verdier领衔——否认,那么很可能早在七十年代初,Mebkhout的晶体理论(他1975年才着手发展,而且顶着这同一批学生的漠然)就已抵达六种\dmark{}\origpage{385}运算形式主义的完全成熟——而它至今仍未达到\fnote{128}{(*)}{(6月7日)贸然下这一评判,我曾有一番踌躇:它可能被解读为贬低Mebkhout理论的原创性。那与我的想法绝不相符;更何况,我对我的每一位上同调学家学生的手段都有极佳的评价(当他们不被外在于健全数学意识的成见所堵塞时)。我的朋友Zoghman本人,消除了我可能有的这份顾忌:他自己深信,“正常地”说,早在七十年代初,本该是我的学生们来发展他的理论。在某个层面上,他们自己想必是最先深信这一点的:那理论的本该的作者,是他们,或Deligne——而世风的普遍败坏再一助阵,便足够他们表现得俨然自己(或Deligne)确实就是作者!对此可参见“研讨会”与“欺瞒”两注,n\textsuperscript{o}~75$'$与85$'$。}。

我倒是记得曾向Verdier谈起那个使我着迷的问题:离散可构造系数与连续系数的关系;而毫无引起他兴趣的样子。这想必随后使Deligne着了迷,因为他用一整年(1969年)主持一场研讨班,以建立一部词典;那词典想必不曾使他满意,因为他随后把它弃之不顾。(见“当值无名之辈与好上帝定理”一注,n\textsuperscript{o}~48$'$。)况且此后,他被自己的下葬综合征“堵”到这种地步:直到1980年10月,他都未曾知觉Mebkhout工作的重要性——而当他终于觉察时,是以人们所知道的那副掘墓心境(见注n\textsuperscript{o}~75至76)。

据我所知,Verdier自答辩以来的著作,其实质仅限于在解析语境(它有时呈现额外的技术困难)里重做我在概形-凝聚框架里所做的事,不引入任何新想法。这甚至相当离奇:凭他分明发展出的反射、凭他消息如此灵通,他竟没有靠自己摇动手柄而撞上Mebkhout的理论——而且竟至少不曾认出:他的“学生”正在做些依我看有趣的事情,而且是他所错过的(一如Deligne所错过的)。

老实说,尽管我对离散系数与连续系数之关系的问题着迷,我却当真不曾疑心Mebkhout的晶体理论——它将在我离去后的十年间绽放。相反,有一个宏阔的主题,出自我五十年代(1955--1960)交换与非交换\dmark{}\origpage{386}上同调的反思;它仅仅被着手过(在“交换”的侧面,即以加性范畴的语言):在Verdier的工作里——六十年代初启动,答辩之后被弃置(见注n\textsuperscript{o}~81)。非交换的侧面,更晚在Giraud的论文里着手:他发展一种几何语言,以topos上的1-叠,表述维数$\leq 2$的非交换上同调。自六十年代下半叶起,这两道起跑线的不足就已显而易见:一边,“三角范畴”观念(Verdier提炼的)不足以表达与一个导出范畴相系的丰富的结构(该观念注定被\textbf{求导者}这个丰富得多的观念取代);另一边,需要为任意维数的非交换上同调发展一套几何语言,以topos上的n-叠与$\infty$-\textbf{叠}表述。人们(或我)感到需要这两条进路的综合;它将为同调代数与同伦代数充当共同的概念基础。这样的工作,也与Illusie的论文工作直接连续——两个侧面在其中都有代表。

经由求导者观念(非交换与交换框架皆有效),Bousfield--Kan论同伦极限的奠基工作(Lecture~Notes n\textsuperscript{o}~304,1972年出版),也处在这条弥散的纲领的脉络里——至少自1967年起,它只待人手来发展。去年1月——我尚不知道一个月后我将投身《追寻栈》——我把关于同伦型的“积分”的反思(同伦论者熟悉其名:“同伦(归纳)极限”)呈给Illusie;那时我还全然不知道Bousfield与Kan的工作的存在,也不知道这类运算已被别人考察过。显明的是,Illusie同样不知——尽管他被认为自1970年我“逝世”以来,始终泡在同调-同伦的水域里!由此可见,他与某些自然地属于奠基反思、属于他六十年代亲身推进的那条路线的实在,似乎失联到了何等地步\fnote{129}{(*)}{同伦型的“积分”这一观念,是在分层结构的剖视的语境里、于1981年底重新强加于我的。}。他想必给自己挖了个小窝,不再常出来……

\dmark{}\origpage{387}随着对topos观念本身、对全部“范畴的无意义”的轻蔑,毫不奇怪:Giraud如今对他头一项大工作主题已全然冷淡。诚然,Deligne两年前开棺出土Motif之际,曾装作忽然发现非交换上同调武库——gerbe、lien等——的趣味,仿佛是他刚刚引入它们,连同Motif与MotifGalois群\fnote{130}{(*)}{见“一个梦的回忆……——或Motif的诞生”一注,n\textsuperscript{o}~51。}。这种马戏,多半无力重燃他自己拼命扑灭的火焰……去年2月,我曾把一封二十页的信的抄件寄给Giraud;它成了《追寻栈》的开卷第一章。那是一场绝无技术性的反思;在其中,我得以“一跃而过”当年曾使Giraud(以及许多别人)却步的“炼狱”——即“非严格”n-范畴(我如今称之为n-叠)的观念:它始终停留在启发的层面,却分明是基本的。那就是《追寻栈》的起点。去年12月,在Contou-Carrère答辩之际,我们相遇(双方心境全然友好);我从Giraud那里得知:他连读这封信的好奇心都不曾有!我得到的印象是:他对这类东西已经一笔勾销。在那里可能有一片丰饶的实体——在他久已放弃的方向上——这个念头,似乎连掠过他心头都不曾。我试图让他明白——恐怕未能成功——那里有一项多汁而宏阔的工作,近二十年来等着有人做;而我到了暮年才终于着手,以至少在总体上勾出一个草图——听事物本身的口授——那片丰饶的实体;我这个“亡者”仍强烈地感到它,而我的学生们早已把它忘了。

Jouanolou同样放弃了一条研究方向——他刚以论文起手的方向。这一方向,成了由恰恰向他提供其选题之关键技术钥匙的那个人所倡导的时尚的轻蔑对象。三年前“反常”研讨会上围绕三角范畴的“抢滩”之际,同一位Deligne忽然(面不改色地)装作发现眼前那项宏大的奠基工作——其缺失忽然处处可感,而他正是十年来头一个劝阻它的人\dmark{}\origpage{388}——这样一项工作的需要,对我来说,早在1963/64年平展上同调起步时就已分明;对Deligne,也丝毫不差,自他开始听说$\ell$-进上同调与三角范畴之时——即次年他登上我的研讨班之时。问题在于:在$\mathbb{Z}_\ell$上(比如说,一个基底概形之上)构造“可构造三角范畴”、并在这一框架里发展“六种运算”的形式主义(这一步,我记得在Jouanolou的论文里已告完成)之外,还要做一项类似的工作:把基底环$\mathbb{Z}_\ell$换成一个多少任意的诺特$\mathbb{Z}_\ell$-代数,比如$\mathbb{Q}_\ell$或$\mathbb{Q}_\ell$的一个(代数?)扩张。这属于那些二十年来时机已熟、却始终等着有人去做的事情——且待吹在它们身上的轻蔑之风平息……

Raynaud夫人的工作的自然延续(平展上同调中的若干弱Lefschetz定理,以1-叠的语言),本会置于一个严格禁忌的$\infty$-叠语境里——不必提它!Sinh~夫人的工作亦然——1968年开始,1975年才完成——其自然延续,本会是一个所谓“monomiale”范畴的包裹$\infty$-范畴观念,或这类范畴的三角变形\fnote{131}{(*)}{见“一个梦的回忆……——或Motif的诞生”一注,n\textsuperscript{o}~51。}——不必想它!另一个是把她的工作转译成topos上叠的语言——多可怕!至于Monique~Hakim,她也不幸把论文做在一个题目上——自我不合时宜的离去以来,这题目如今在边缘上显得有点可笑:局部带环的topos上的相对概形,您说说看!她关于这一主题的小书,刊于《Grundlehren》(Springer),想必每年卖三四册——难怪我在那家名声不好,他们对接受我可能推荐的文本不再热心。对我而言,那是为一切“绝对”的“簇”观念(代数的、解析的等等)向一般“基底”做“相对化”的第一步试验——其需要,在我看来是明摆着的(注91\textsuperscript{3})。人们会说:直至今日,没有它也过得很好。但同样真确的是:人在这世上两百多万年,没有做数学也过得很好。总之,Monique~Hakim——她做论文的动机与我为她提议该题的动机并不相同——肯定绝无念头与那样一个主题保持哪怕一点接触;那主题,一旦脱离有利共识的语境、或脱离一股顶风冒雨坚持一种执拗而可靠视野的思想,便不再能对她有任何意味。

至于Neantro~Saavedra~Rivano,他似乎已全然从流通中消失——在世界数学家人名册(以及一切官方的名册)上,我甚至找不到他名字的踪迹。可以肯定的是:他那过于范畴家的论文题目,在那些裁定何为严肃、何为不然的先生们那里,想必不会有好名声。依我看,这篇论文最自然的延续,不多不少,正是那幅“Motif的宏阔图景”——对这个学生较为谦逊的抱负而言,那主题分明嫌大。然而他到底得了那份意外的荣幸:不过两年前,亲眼见自己的论文被其中一位大人物从头到尾重做。(对此可参见“下葬——或新父亲”与“白板”两注,n\textsuperscript{o}~52与67.)

最终,在我的十二位“1970年以前”的学生中,我唯一不太清楚其工作中是否有一道或多或少的苛烈或深刻的\textbf{断裂}——相对于他们在我的接触中所追随的工作——的,是Michel~Demazure与Michel~Raynaud(注91\textsuperscript{4})。我所知道的只是:他们继续做数学,而且(就其出众的手段而言,想必在意料之中)属于我方才所称的“上流”数学界。

前面这段从有时极其单薄的材料出发的简短反思,当然在很大程度上是假设性的、极其粗略的。我希望被提到的人肯原谅我或许粗疏的判断之误;倘若他们肯就此向我示意,我将乐于更正。在这里,我再次意识到:每个人的情形,肯定不同于所有其他人,代表一个远比一个如我般遥远的人所能合理把握的实在复杂得多的实在,更遑论以几行文字表达。所有这些保留说尽,我却仍感到这场反思并非无用——至少对我而言:以若干具体事实,稍稍安放昨天显现的那份仍属弥散的印象(它想必在未成形的层面,多年以来已在场):即一道\textbf{断裂}——它在我离去后的翌日降临于我的许多学生身上,并且在个人的层面上,映照一所“学派”的突然消失——一夜之间;在那道学派受业岁月的关键年月里,他们必定感到自己是它的一部分\dmark{}\origpage{390}。

\notehead{91\textsuperscript{1}}{}

(5月22日)我刚读到 P.~Berthelot的一篇综述文章,论CIRM吕米尼研讨会(1982年9月6--10日)“p-adic分析及其应用”;题为“刚几何与特征p代数簇的上同调”(24页),勾勒综合Dwork--Monsky--Washnitzer上同调与晶体上同调的主要想法。晶体上同调(受Monsky--Washnitzer启发)的出发点观念(连同其名字本身),以及经由引入由刚解析空间构成的site使之完备的观念——这些我在六十年代引入的想法,已成了所有从事这一主题者的日课,以Berthelot为首;其论文正在于发展并充实这些出发观念中的某些。这并不妨碍:我的名字,无论正文还是参考文献,都严格缺席。这样,第四位殡仪人学生清晰在案。下一个轮到谁?

(6月7日)有一件事值得注意:自我引入晶体上同调的起步观念十五年有余、自Berthelot确立该理论对真光滑概形确是“对的”理论的论文十有余年之后,人们仍未抵达我所称的晶体上同调的“掌控”局面——堪比为平展上同调在SGA~4与SGA~5研讨班里发展的那一种。对一个含对偶现象的上同调形式主义的(第一阶)“掌控”,我指的不多不少,正是对一套六种运算形式主义的完全占有。虽然我还不够“在局中”、无以赏识晶体语境特有的困难,但我不会惊讶:这一相对停滞的主因,在于Berthelot及其他人对这套形式主义观念本身的冷淡;这使他们(一如Deligne对其仍处童年状态的Hodge理论)忽视那头一个本质的“台阶”——通往一个全然“成年”的上同调形式主义所必达的台阶\dmark{}\origpage{391}。想必正是同一类心境,也使他不曾察觉Mebkhout的观点对他自己研究的趣味。

附记:当我在这里谈“晶体上同调”,而语境是放弃真性的假设时(对一个“全然成年”的形式主义,这是必需的),不言而喻:人们所用的晶状site,其对象是并非纯无穷小、而是“合宜的”(带除幂的)拓扑代数型的“增厚”(带除幂)。对晶状site的这一原初形态之推广的需要(对我来说,原初形态只是“对的”晶体理论的第一次逼近),我一开始就很清楚;而Berthelot是(连同出发观念一起)从我这里学得它的。这一关联的书面影射,见《主题概要》第5条。

\notehead{91\textsuperscript{2}}{}

相当离奇的是:除我之外,似乎没有人觉察,未具名的Mebkhout的理论是晶体理论的一个本质的新侧面。我——一个近十五年来已全然“脱钩”于上同调的人——却觉察到了,一旦Mebkhout去年费心向我多少解释了他所做的事情。而当我在Illusie面前(当作理所当然地)提起这一点时,他的神气,仿佛把这看作把两件分明毫无关系的东西(D-模与晶体)做了一点“边缘离奇”的拉扯。然而我第一手知道:他有数学家的嗅觉;我的其他学生(在这个案里即上同调学家们,从Deligne起)也有——但我确认:在某些局面里,这嗅觉对他们不再有任何用处……我越想,越觉得离奇:在这样的气候里,Mebkhout到底还是做成了他的工作,没让他的数学嗅觉被其长辈们的全然不理解——他们高踞其上……——缴了械。

\notehead{91\textsuperscript{3}}{}

主要是从五十年代末我在Cartan研讨班上关于复解析空间理论之基础、以及“Teichmüller式水平模簇”之精确几何诠释的报告起,我才懂得了对迄今所用的“簇”的通行观念(代数的、实或复解析的、微分的——或此后,“温和拓扑”中的变形)做双重推广的重要性。其一是放宽定义,以容许任意的“奇异性”与“标量函数”结构层中的幂零元——以我的概形奠基工作为模型。\dmark{}\origpage{392}另一项推广,是向合宜的局部带环topos之上的“相对化”(诸“绝对”观念,经由取一个点的topos为基而得到)。这一概念工作,成熟已逾二十五年,在Monique~Hakim的论文里着手,至今仍等着有人接手。一个格外有趣的个案,是相对刚解析空间的观念:它使人得以把寻常的复解析空间与变剩余特征的局部域上的刚解析空间,看作同一个相对刚解析空间的“纤维”;一如相对概形观念(它终于入了俗)使人得以把定义在不同特征域上的代数簇相互联结。

\notehead{91\textsuperscript{4}}{}

Demazure的论文工作,一如Raynaud的,本质上使用他们在与我的接触中学得的一套纯熟的概形技术;但各自工作的本质观念,不属于“Grothendieck式”的行装——这使他们的工作有别于我第一期其他学生的工作。很可能这一情势的后果,是其著作中的一份连续性,免于“老师下葬综合征”效果造成的断裂。这未必意味着这一综合征不曾以别的方式触及这一位或那一位。三年前,Raynaud对Contou-Carrère相对局部Jacobi簇工作的态度,曾使我吃惊。所宣布的结果深刻、艰难,而且美不胜收,远超对“众所周知”之事的简单推广;与Cartan典型曲线理论有出乎意料的联系;有漂亮的显式公式——全然在Raynaud的射程之内(也在我的射程之内)。他接待的那份冷淡,想必对Contou-Carrère的战略撤退起了决定性的分量:他放弃了那个他曾毫无保留地投入、而看来只会给他带来麻烦的题目……\fnote{132}{(*)}{细节见“棺材3——或有点太相对的Jacobi簇”一注(n\textsuperscript{o}~95)的子注n\textsuperscript{o}~95\textsuperscript{1}。}。我那封向他表示我(痛心地)惊讶于他对这些结果之美无动于衷的信,没有得到回音。

\subsection{……与电锯}

\notehead{92}{}

\dmark{}将近四年前,当我来到这一带安家时,离我家不远有一片漂亮的樱桃园。\origpage{393}我散步时常绕过去看看。我喜欢看那些茁壮的樱桃树——正当盛年,树干粗壮,仿佛自古就与这片土地长成一体;野草在树下自由滋蔓。它们想必没见过化肥,也没见过农药;樱桃时节,我就到那里去摘有味道的樱桃。树,想必有二三十棵。

有一天我再去,看见所有的树干都在人高的地方被锯断,树冠瘫倒在树干旁的地上,茬口朝天——一幅大屠杀的景象。用一把好电锯,想必顷刻而成,至多个把钟头。这样的场面我从未见过——人们锯树,一般总肯弯下腰,贴着地面锯。樱桃滞销,这我同意;这片园子想必也产不出几吨,这没问题——但这些树茬诉说的,可不是滞销与收成……

昨天,我又有了这同样的感觉:一根强劲的树干,根系深厚,树汁丰沛,枝条繁壮、延续着它的冲劲——被齐人高、干净利落地锯断,仿佛纯粹为了取乐。正是费心把主枝一根一根看过、看清每一根都被锯断,才终于使我看清发生了什么。那本该循着一股冲劲、循着一种根系深处内在的必然而舒展的东西,被一刀切断,不加手续,好让它在众人眼前被指定为笑柄。

这使我想起Zoghman谈起的那个“误会”:它想必发生在我的学生们(Deligne除外)与我之间。分明的是:冲劲与视野,都不曾从我传到我的任何一位学生身上(Deligne除外——他分明“除外”,名副其实!)。每一位都消化了一套技术行装——对他所选的题目做出色的工作有用(甚至不可或缺),日后兴许还用得上。是否有别的什么东西的起头——超越于此——我说不好。即便有起头,它面对电锯也绝无机会:电锯顷刻把它刈平……

我知道,只要还有人做数学\dmark{}\origpage{394}——而且除非彻底放弃人们做了两个多千年的那一类数学——他们早晚会不由自主地让我眼前这些僵卧的树枝一一复活。其中一些,已经被我的电锯朋友据为己有;倘若上天借他余年,很可能他还会照样处置另一些,甚至全部。但大多数,已全然不合他的风格。不过,也许他终究会厌倦于不断取代别人——那肯定极其累人,而且无利可图——而满足于做他自己(这已经不错了)。

% 第16章 D) 被埋葬者们;原书p397-412

\chapter{D)~被埋葬者们}

\section{X.~灵车}

\subsection{棺材1——或知恩的D-模}

\notehead{93}{}

\dmark{}(5月21日)我的反思在那些“根正苗红”的学生们——“1970以前”的——身上,已滞留两个星期。每一天,这场反思都摆出一副“最后一项补充”的面目,出于求个心安,补向一场(几乎)已告终结的反思。不止一次,是一条不起眼的脚注——贸然接上前一天或前两天的反思——越拖越长,直到一份独立“注”的规模。每一回,它都很快得名,把自己区别于所有其他的注,并恰好插进其葬礼行列的合适位置,仿佛它从来就在那里!每隔两天,我都得(每次都怀着乐趣)至少重做目录的结尾:目录仿佛已收尾,却随之添上两三位新的入列者,有时甚至一整个新的行列……

这支队伍的规模,渐渐变得令人不安:谁也不会肯把这全部读完!但它这样加长,老实说,并不是为了一个假想读者的可疑收益,而首先是为我自己的收益——一如我做数学之时。这些“最后一项补充”——我每一回都仿佛违背本意地登船——从不曾让我后悔。借着一次次的最后一项补充,我学到了许多别处无从学到的东西,倘若不这样“从材料出发”反思,我便无从学到这些。而这些东西,一件一件聚成一幅色彩鲜明的图景,宏阔而有多个侧面。此刻我仍看到它尚未全部完成——还有两处,似乎在讨最后一笔。

我觉得,谈过那些“根正苗红”的学生之后,如今也该稍稍谈谈\textbf{被埋葬者们}了——那些“与我一同,有权领受这场以沉默与轻蔑进行的下葬之殊荣”的人。一如我、一如那些兴冲冲下葬的人,这些被埋葬者也不是圣人,也没有殉道的使命。我相信,他们中没有一个不曾暗暗怪我\dmark{}\origpage{396}给他招来麻烦——那些我全然无意给他招来的麻烦(只因为他曾失于轻率,把赌注押在我身上,押在某种做数学的进路与某种做派上……)——或者,至少不曾试图与我划清界限,一旦认清赌注分明已输\fnote{1}{(*)}{(1985年2月)自我到蒙彼利埃以来,在校外,我总共得知七八份(简短的)出版物,呈献(以摘要的形态)一项与我一起做、受我启发的工作。我的名字,在所有这些里缺席。}\origpage{397}。何况我可以确认:那是白费力气——一旦被盯上,就完了;划清界限,只会喂养一份轻蔑、给它一份心照不宣的辩解,而不能解除它。同样不止一次、以多种方式,我看到下葬者与被葬者的角色相互毗邻、相互混淆\fnote{2}{(**)}{(9月2日)以各不相同的方式,他们每一位都曾在某个时刻,终于把对其工作的轻蔑内化并据为己有,默许那份藏掖其工作或把它归为“无趣”的共识。}。想必正是这些暧昧的侧面,造成我身上一份长久的迟疑:不肯比此处彼处顺带的影射更周详地谈“被埋葬者们”。很可能,除Zoghman之外,另外三位我认识的人里,没有一个会感激我在这里替他做“广告”——仿佛我给他招的麻烦还不够多似的。

一如《收获与播种》里的许多时刻,我终于越过我身上这份迟疑。我对自己说:即便对着那些曾因我而受苦的人——经由他们在某个时刻做出的选择;在那选择里,他们出于这样或那样的缘由,分明各得其所,尽管他们同我一样,不曾料到其选择所附带的种种不利——即便对着他们,我的角色,也不是帮他们回避一个实实在在的局面;无论愿意与否,他们都身陷其中,而这局面分明有其意味,哪怕它带着严重的麻烦。

在接上我那些亲爱的共同亡故者、共同被埋葬者的四口棺材的黑色系列之前,我或许该以一条不那么丧气的注,给读者换换心情。首先,在与我们大学数学研究所的“地方”层面的关系里,我完全没有体验到:我为一个职位的候选人所说的话、或候选人是我的学生(不言而喻,是“1970以后”的)、或其著作受我著作的影响,就必然于他不利。那种系统性抵制的态度,仅仅是“上流”数学界对我这个人的关系——推而广之,对那些显得与我“1970以后”相连的人——的特征。据我所知,这份抵制,自我离去以来的十四年里几乎毫无破绽,只除两桩寒微的例外。其一涉及一个学生:他开局颇有希望,本该跟我做一篇题材极为诱人的国家博士论文;他申请我们大学的助理教授职位,却被专家委员会驳回。他在全国一级被“捞了回来”,有Demazure的帮助——我曾就这位学生的工作给Demazure写过信\fnote{3}{(*)}{就晋升或取得职位与身份的“实际”层面而言,我1970年以来的教学活动的全部账目,归结为两桩取得职位并附带身份的事:一回是助理教授,另一回是助教。说来奇怪的讽刺:两回,这晋升都是当事人一切研究活动戛然而止、彻底停摆的信号。}。另一方面,有两次,《Topology》杂志接受了出自我的学生的文章:一篇“Stein因子分解与剖分”,出自Jean~Malgoire与Christine~Voisin;以及Yves~Ladegaillerie一篇将刊的文章,包含其1976年论文的中心结果(见注n\textsuperscript{o}~94)。

我尤其已有机会谈到过Zoghman~Mebkhout,这里只“存照”再谈\fnote{4}{(**)}{除《导言(6)》(葬礼)之外,谈到Mebkhout的还有:“我的孤儿们”“当值无名之辈与好上帝定理”“不义——或归来的意义”“反常性”“墓外相会”“受害者——或两种沉默”“大部头与上流社会”“赊账的学位论文与全险”各注(n\textsuperscript{o}~46、48$'$、75、76、78、78$'$、80、81)。}。Mebkhout想必自1974年起开始从我的著作汲取灵感,并且顶风冒雨、一直汲取到今天。我不知道我的哪一位“官方”学生做出了分量可比的著作——尽管Mebkhout的著作,必然带着其被迫在逆境中推进的印记。如我在《导言(6)》里所说,四年来,Mebkhout的观念与结果,人人都在用,而他的名字却被小心地藏掖\fnote{5}{(***)}{在那场下葬里充当掘墓人的,大有人在;1981年6月的吕米尼研讨会,几乎全员参与。除我的上同调学家学生们(对此可参见“我的学生们(2):团结”一注,n\textsuperscript{o}~85)之外,其职业诚意在此直接且严重地成为问题、而我有知情的,是 J.L.~Verdier、B.~Teissier、P.~Deligne、A.A.~Beilinson、J.~Bernstein。}。我的朋友何以能在承受轻蔑、继而承受不义——如一场不可抗拒的宿命——的同时继续做数学,这对我是谜——那份宿命,经由他想必(并且想必仍然)感到高悬其上、令人眩晕的人而来\fnote{6}{(*)}{当然,Zoghman~Mebkhout并不比我傻,而且他足够“在局中”,对我每一位上同调学家学生的工作都有精确的把握,能掂量其分量与限度,毫无理想化的倾向。但这并不妨碍:一些力量可观的抑制,使他连那样一个念头都不曾有——即公开质疑他们中的任何一个,哪怕在恶意分明之处。}\dmark{}\origpage{398}:那些人,他想必头一回听说时,俨然“赛场诸神”;而那时,他(一如当年的我)是一个生计拮据的寒微留学生。1979年答辩之时,他在奥尔良有一个助教职位。当时他竭力设法进CNRS,三次进攻——到第三回(1982年10月),人们终于肯给他一个研究员职位(相当于大学的助理教授或副教授)。这给了他,即便不是身份的担保,至少一份相对的安全。

在我所知的四位“共同被埋葬”的数学家之中,Mebkhout是唯一一个逆着所有人继续其工作的:他信任自己的数学本能,不被那份无情时尚所能提示的谨慎与机会主义的考虑绊住。他身上——他并非好斗的天性——有一种对自己判断的素朴的\textbf{信念},它同时也是一份\textbf{宽宏};而它(远甚于脑力“手段”)是做出革新而深刻的工作的首要条件。我对他的著作的看法,肯定仍不完全。就我所知其著作的主干部分而言,我觉得:以他那出众的手段,倘若处身于一份热忱而起作用的同情的气候里,他本可以在三年或四年里(而非十年)完成它,并把它带向更大的成熟——在喜悦而非辛酸之中。但三年也罢十年也罢,“成熟”与否,值得注意的,是革新的著作出现了,而且竟能在这样的条件下出现。

\subsection{棺材2——或电锯剖切}

\notehead{94}{}

Yves~Ladegaillerie1974年开始与我工作。那是“姑且一试”,在他生涯的一个低谷——我当时向他提出了几条关于1维拓扑复形在曲面中的嵌入的天真反思;那时我对曲面一无所知(除亏格观念),他懂得更少。这有点像Grothendieck学问(反正在我这里,一切总是这样开始……);它在他那里或多或少咬住了钩,直到有一天终于“翻了盘”——我已记不清何时何故。也许是在一个分明多汁的问题正在成形之时:关于紧致带边定向曲面中1维紧复形的同痕类的判定的一条关键猜想。真——假?悬念持续了六个月、一年;其间伊夫熟悉了(也顺带使我熟悉了)曲面理论的关键定理,同时推进其工作的“奠基”部分。已知的结果使猜想相当可信,却分明远远不够——因为那猜想蕴含巴厄与Epstein的那些难啃的结果,以及另一些有着离奇甚至可疑面目的东西。1975年夏,他终于证明了那条关键猜想。它实质上等价于:以基本群的语言,对一个可三角剖分的紧空间嵌入紧致带边定向曲面的全部同痕类,给出完整的代数描述\fnote{7}{(*)}{非定向情形的“类似”陈述是假的——分明这是一个精微的结果,被小心地“剖”在一组同样“可信”却仍为假的假设-结论之中。关于Ladegaillerie工作的其他评语,见《纲领概要》,尤其第3节开头。}。

\dmark{}从伊夫“咬钩”那一刻起,他一年、至多一年半做完了论文:结果、成文,一应俱全,而且几乎苛刻地完美。那是一篇光辉的论文,比我多数学生的论文薄,却与其余十一篇中的任何一篇同样实。答辩在1976年5月举行。

论文至今未出版。它虽不厚,据说还是太厚,不宜出版——这是人们给我的许多条极好的理由之一。我在“进步是挡不住的!”一注(n\textsuperscript{o}~50)里列举了几条。我为“安放”这篇倒霉论文所作努力的经过——它是我有幸激发的最好论文之一——可以写一本小书;它肯定富有教益,而我放弃写它。在那些有极好的理由忘掉去了解这些结果、并闭着眼睛把这一切埋葬的昔日好友之中,有(按登场的顺序)Norbert~A.~Campo、Barry~Mazur、Valentin~Poenaru、Pierre~Deligne——还没算经Springer公司转手的 B.~Eckmann\fnote{8}{(*)}{我与Eckmann并不相识;为使伊夫的论文在Lecture~Notes出版而做的通信,是与Springer主管LN的Peters博士进行的。我想,经由六十年代由我(SGA尤其是)或由学生(论文)出版的十五卷LN,我曾是凭我的担保为这套丛书尚在草创之际的前所未有的信誉与成功出过力的人之一。拒绝我所推荐的工作所给的理由(他们不出版论文),是句玩笑话。通信事务上的“新面貌”,我的头一份经验也出自这一情节:怀着一套真正令人叹为观止的默契, A.~Campo、B.~Mazur、 V.~Poenaru与Peters博士,对我第二封信,都竟免于赏我一复——当时我天真(我的悟性慢……)地卷土重来,而他们那封迟疑的回信已表明:他们没有费心了解Ladegaillerie工作导言里所陈述的结果。}\dmark{}\origpage{400}。中心结果终于要发表了——九或十年之后,削减到只剩骨头,以《Topology》一篇短文的形式(嘘——我在那份可敬杂志的编委会里有一位同谋……)。工作的其余部分,一边证明了人人久已不证而用的事情(当然,没有证明也过得很好!),一边发展典型的Grothendieck学问,全然背逆时下的惯例与良俗。我很清楚:倘若我的朋友Deligne不在未来十年里大张旗鼓地负责“发现”它们,别人在三五十年内也忍不住会重做它们——因为我健全的本能告诉我:那些是基本的东西。它们曾是我阿那贝尔式沉思中宝贵的引线;倘若上天借我余年,在《数学反思》发展阿那贝尔式代数几何瑜伽的那一部分里,我会有充分的机会指涉它们。

这段经历对我是一次显影,同类中的头一次——对某样东西的显影;那东西,直到葬礼的反思,我才充分认识。况且,自从我的心神另有所系,我倾向于把它忘了。Yves~Ladegaillerie——我有过的最光彩的学生之一——则从那一刻起就懂得了:要在今天的数学界被接纳,仅仅全力投入、并做出合乎一切卓越要求的工作,是不够的。他另有几手,七年来他投身于更接地气的任务与较少问题的回报。他运气好:早在与我那场不幸相遇之前,他就有助理教授的终身职位\dmark{}\origpage{401},给了他一份其厄运不曾危及的保障。去年,一点数学的火星似乎重新燃起,在一个与我近年来关心的问题十分贴近的主题上——瑟斯顿式的双曲几何及其与Teichmüller群的关系。甚至可能我们还会一起走一小段路;或者他走他个人的散步,纯粹为了乐趣,不期待数学本身所能给的回报之外的任何回报。他清楚:倘若他期待别的,他最好换一个对话者或同路人(甚至,换一段过去……)。

\subsection{棺材3——或有点太相对的Jacobi簇}

\notehead{95}{}

我与Carlos~Contou-Carrère的最初相遇,是在1973年我到蒙彼利埃之后不几天,数学所的走廊里。他会把我堵在某个昏暗的角落,在我来得及客气告退之前,把一股数学解释的洪流倾倒在我头上。他连珠炮般倾倒给我的东西,全从我头顶掠过;而他丝毫不曾察觉,或在我怯怯暗示之时有丝毫窘迫。他迫切需要对话者,而我并非他唯一的“违心对话者”。那又是一个我全然不沾数学的时期。有一两年,我一看见他(一眼就能认出)的身影在走廊尽头出现,就躲。直到林登——他曾在蒙彼利埃做过一年客座教授——向我透风:Contou-Carrère手段不凡,却眼看要触礁,因为不知道怎样使用它们。在那之前,Contou-Carrère倾倒给我的东西站不站得住、他有没有手段,连掠过我心头都不曾——一切那么遥远。也许林登的提示来得恰逢其时:我开始对数学问题重新发生一点兴趣。总之,我横下心来,问Contou-Carrère可肯给我解释一件他做过的事,好让我能懂。我疑心:我是头一个向他提这种要求的人——至少在他来法国已有好几年之后。要他显豁一件事不容易,却绝非不可能,而且值得。我很快发觉:林登没有看错——Contou-Carrère满脑子观念,只等着被提炼、被细心发展;而且在我们可能向他提交的几乎一切数学局面里,他都有一种即时的、非常可靠的直觉。凭这份直觉的敏捷与可靠——哪怕在他毫不熟悉的事情上——他超过我、令我惊叹;在别的学生身上,我以可比的程度体会到它的,只有Deligne\fnote{9}{(*)}{我不确信在其他数学家身上遇到过它——除Pierre~Cartier(他年轻时以这份卓越的能力深深打动过我)与Olivier~Leroy(下一条注将谈到他)。}。相反,他对书写有一道几乎完全的阻塞!不可思议:他做数学\textbf{不写}——天知道他怎么能哪怕做出那么一点,更不必说与他人的交流:在那方面,“触礁”是彻底的(见上文)。

倘若我有什么紧急而有用的事情要教给Contou-Carrère,那是书写的艺术——或更直截了当地,哪怕只是让他明白:数学是\textbf{写着}做的。我费了两年、也许三年,直到1976或1977年\fnote{10}{(**)}{(6月7日)核实之后:直到1978年2月。},不敢说是否全然成功。他第一部完全白纸黑字写成的大工作,是他的Schubert闭链论文——迟至去年12月(1983年)才答辩\fnote{11}{(***)}{那是一项长工作(我没有读过):他细心发展一些与我毫无关系的想法,其中给出一切“Schubert”型闭链的显式奇点消解——此前无人能做到。难得他有一次成形的写作,人们却责备它太详(更不必说他的陈述太一般……)!至于我,倘若我有批评,那将朝着相反的方向:虽然Contou-Carrère声称其方法适用于一切类型的半单群与Schubert闭链,他却只在一般线性群的情形做了工作——因此他在这件明确的事情上没有做到底:普适Schubert闭链的等变奇点消解\textbf{以及}这些Schubert闭链的奇点轨迹的描述。这份欠缺,在我看来,是那道对“从材料出发”的扎实工作、对书写的“阻塞”的一份遗产;那阻塞,长久以来是他的主要缺陷。}。1978年至今,我们的关系极为零星;我的角色,实际仅限于在他职业生涯中屡屡以这样或那样方式陷入困境之时——那职业,始终悬在最不牢靠的代理助教职位上——尽力扶持他。

有两三年,我曾试图向Contou-Carrère提供一套精确而柔顺的数学语言的基础,外加几条系统化的原则。带着这份行装,连同他的手段与观念的丰富,接续什么,他有真正的挑选余地。他没有从他自己的出发观念入手,而是接上了相对局部与整体Jacobi簇的理论——我曾把它作为可能的论文题目向他提过。一旦我放手让他自己干,他在仅仅一年之间做出了一项很漂亮的工作,其中一部分在一篇科学院报告短文里已有预告(注95\textsuperscript{1})。把这个矿脉做到底,本会代表几年令人着迷、强烈激励他的工作——足以同时学得概形技术的全部精微。那时我毫无疑心——明摆着,Cartier、Deligne、Raynaud三人都会热烈接待这项已完成的工作:深刻、艰难,且在好几个方面出人意料。Cartier确实满心欢喜地看到他的某些旧观念取得新的现实性。相反,是Raynaud的冷漠,一如Deligne的——后者把完整手稿在抽屉里压了六个月,不屑给一点活着的音讯\fnote{12}{(*)}{Contou-Carrère倒也先防了一手:在其短文里,对我的这个人——曾给了他出发纲领的人——只字不提。那也是白费:任他往里加多少自己的东西,有一种“做派”不会错,它系于某些主题;谁愿在今天的数学里混出前程,最好避开这些主题。(6月7日)向当事人查询之后,我发现:我在这里把围绕Contou-CarrèreJacobi工作的两个不同情节弄混了。细节与精确出处,见下一条注(n\textsuperscript{o}~95\textsuperscript{1})。}。这是二对一——足以感到风向。那有点太相对的Jacobi簇,被无限期弃置,任其自生自灭。电锯干得很出色……

这并没有为Contou-Carrère挡住厄运;为其写一份周详的报告,本可以再写一本小书,我由衷放弃。我想正是在这时——自我1970年离开那个机构以来(在那机构尚无自己的屋顶的年月里,我曾是四年间1958--62唯一在“实地”代表它、使它取信于人的人\dmark{}\origpage{404}),唯一的一次:我担当起为某人推荐一份邀请(为期一年)的责任,正当Contou-Carrère眼看要没有职位、流落街头。我知道:我推荐的人——一如当年的Hironaka、Artin或Deligne被我热情接待时那样无名——会像他们一样,为接待他的机构增光。当然,我不曾忘了这么说。对Contou-Carrère幸运的是:他的代理助教职位(诚然配不上一份如此挑剔机构的邀请之荣)最终得以续聘\fnote{13}{(*)}{我无从抱怨:五六年之后,借去年IHES建所二十五周年庆典之机,人们倒真赏了我这个人一份邀请,甚至让我在部长致辞的庄严招待会与随后一周的IHES居留之间选择——还包食宿(人家郑重向我担保)。我对老朋友Nico~Kuiper说:想到我,非常客气;但我这把年纪,不再出门旅行了……}。

对这一情节我并不太惊讶:当时我已了解Deligne的心境;况且Nico~Kuiper曾提醒我,在这个案子里一切取决于Deligne。(我甚至没有想到向他暗示:这事兴许也关涉科学委员会的其他成员——恰恰因为这是个案……)然而,在Contou-Carrère(我的“门生”——Verdier曾在一封信里想当然地这样称呼他……)的一切厄运之中,最触动我的情节,发生在1981年10月:事涉他申请佩皮尼昂的一个教授职位。佩皮尼昂的同行们(他在那里任代理助教)想必乐于身边有一个在任何数学分支里都从容、可咨询的人。一次教授职位出缺,他们把他列为唯一候选人——罕见之至的事,分明表明:他们想看的正是他、而非别人。Contou-Carrère在其阿根廷Santalo门下的博士论文之外,出版物相对少:主要是科学院报告短文,预告结果(其中一些深刻),而无证明。还没有人向他暗示过:在如今的时势下,在安身之前,最好有附完整证明的文章作“证据”——这一点我这边当然向他念叨过不少,只是从较不功利的角度\fnote{14}{(**)}{前一年,Contou-Carrère申请雷恩的一个教授职位;他认识Berthelot与Larry~Breen。当时大学咨询委员会认为他的申请可受理;但职位给了另一位候选人。没有人费心提醒当事人:倘若他想有职位的希望,就得发表他所预告结果的详细证明。次年委员会的驳回,对Contou-Carrère、对他的佩皮尼昂同行们、对我,都来得完全意外。带着事后的认识、在这场反思的光照下,我进而怀疑:随着其论文的成文(如今已被宣告“照原样不可出版”)与答辩,局面是否真的变了;他是否有希望在法国找到一个教授职位。}。总之,Contou-Carrère的申请\dmark{}\origpage{405}被大学咨询委员会判定不予受理,材料退回。当时使我倒吸一口凉气的,是:委员会主席(做出决定的全国机构)不以委员会名义、任何成员也不以个人名义,具有这点最起码的尊重:给当事人Contou-Carrère本人、或至少给佩皮尼昂数学所所长写几行,解释这次表决的意味——在毫无解释的情况下,这次表决只能被领受为对佩皮尼昂同行们选择的响亮否定,对其唯一候选人能否体面胜任其被提名之职的否定。委员会里有我的三位旧学生,其中两位亲自认识Contou-Carrère。他们当然知道他同他们一样是我的学生——何况材料里还包括我一份对候选人工作格外称道的报告。他们中没有一个、委员会其他成员中也没有一个,想到过这记不经任何手续的铡刀式表决所代表的羞辱,想到过对一位不比他们中任何人逊色的数学家的按部就班的鱼雷攻击。

正是这一事故,在我做数学家的一生中,头一回使我感到那种我已在反思中不止一次谈过的“气息”。四年前,经由“外来者”的情节,我已感到过它\fnote{15}{(*)}{见“我的告别——或外来者”一节,第24节。}。但那不是在我曾经所属的世界内部,吹在\textbf{他们自己人}身上——吹在一个毫无保留地认同于那个世界的人身上。我好几个星期、也许好几个月为此如病如恙。为摆脱当时攫住我的一份焦虑——我当时无心去认识它\fnote{16}{(*)}{我是在次年一场漫长的禅观期间,才意识到这份焦虑的;在那场禅观里,我发现了焦虑在我生活中的角色——它的在场(1976年之前是慢性的,1976年之后是偶发的),曾是我一生的“世上守得最好的秘密”。有一套效率极高的机制,藏掖焦虑的一切公认的征兆;那焦虑,对我自己、对我的亲人,都始终不为人知。},我四处奔走,东一封信西一封信,还写了一篇三十来页的文本《大脑与轻蔑》,带着黑色幽默的笔调;我终于放弃了发表它\fnote{17}{(**)}{劝我别发表它的,恰恰是那些我正准备为之开战的人——幸亏我有先见之明,在试图让它面世之前,把我的文本给他们看过。}。带着事后的认识,我看清:那本是就所遇之事的意义展开\textbf{禅观}的绝好时机。最滑稽的是:当时“妨碍”我觉察需要一场深入禅观的,恰是一场我那时正投入的长禅观\fnote{18}{(***)}{见“扫兴的老板——或高压锅”一节,第43节。}——而且是一场关于我与数学之关系的禅观(即便还不是关于我数学家过去的)!它被一个情节搅乱:生活以强力向我发问——而我在奔忙中回避那发问,随后重新扎进“禅观”。带着事后的认识,我看清:那场“禅观”当时配不上这个名字;它缺少真正禅观的一个本质维度:对\textbf{当下时刻}的我自身的留意。我当时“禅观”某些较为遥远的事件的意味,却忽视一份被压抑的焦虑(诚然,由于长期练习这类控制,它被完美地控制着)——那是我拒认那股被驳回的“气息”所带来的讯息的征兆。

但我正在离开本题。那场鱼雷攻击,当然有其必有的效果。佩皮尼昂的同行们被申饬了一次,就够了。显然,他们那里连代理助教的职位也没有了——至少对Contou-Carrère没有。他在蒙彼利埃临时找到一个替补职位,就本学年,编制者明年回来。

不过我并不太为他的前程担忧:一段时间以来,Contou-Carrère已有智慧抢在命运的打击之前,转而投向计算机。以他那出众的手段,想必久已高屋建瓴地驾驭这门学问,同时在闲余做他喜欢的数学。他是两个孩子的父亲;在如今的时势下,带着紧追不放的过去,数学分明是碰运气的事,甚至凶险。他大可以做一个辉煌的计算机学家:没人会因他多少当过我的学生而对他严苛。

\notehead{95\textsuperscript{1}}{}

(6月7日)大约77年底,我向Contou-Carrère提交了一份为相对局部与整体Jacobi簇理论的周详工作计划;其中包括,在局部情形,建议“重新拧紧”Jacobi簇与Cartier的ind-群,以找到一个具有更美泛性质的“完备”Jacobi簇,它应当是“自对偶”的。我没有任何可提出的证明,1978年2月之后也不再过问其工作:我发觉我的在场抑制而非刺激其才能。他在随后的一年里到底“启动”了;其第一篇注“相对曲线的广义Jacobi簇:构造与分解的泛性质”(论整体情形)发表于1979年7月16日(CRAS t.289, Série A - 203)。

次月,他为局部Jacobi簇找到了决定性的结果,却一年半不发表;然后在1981年3月2日的一篇科学院报告里发表“一半”(寻常的、未与Cartier群重新拧紧的相对局部Jacobi簇的泛性质),顶着(乍看不太有说服力的)名字“相对几何局部类域论”(CRAS t.292, Série I - 481)。至于完备局部Jacobi簇的理论——依我看更有趣味——存在一份科学院报告的底稿,从未发表,题为:“局部Jacobi簇,普适Witt双向量群与tame符号”。当然,1979年起我就知道他的结果:那是对我向他提议的临时纲领的一个完整实现;为它,须克服可观的技术困难,要求大量的想象力与技术力量。(如无差错)我只读过第一篇注;我奇怪他不发表后续,即局部部分——他从不曾就此清楚说明——但他分明对第一篇注所受的接待感到失望。1980年雷恩申请失败之后,鉴于附在他申请材料里的我的支持信\dmark{}\origpage{408}论及整体与局部相对Jacobi簇上的卓越结果,他想必还是认定谨慎为上(为准备次年佩皮尼昂的申请),至少再发表一篇关于局部Jacobi簇的注,如非和盘托出。又是两个月之后,1981年5月,他把第三篇注的底稿寄给Deligne与Raynaud(Cartier想必久已知情),我想是先探探路。(我不相信:自1979年8月他手握结果以来的任何时刻,他找Cartan呈送这第三篇注会有丝毫困难。)Raynaud与Deligne都不给他一点音讯——但1982年3月,Deligne给他寄来一篇文章的手稿“tame符号评注”,作者是加藤和也,献给Deligne;它在基底域的情形做出Contou-Carrère的理论,并猜想其对任意基底环有效。Contou-Carrère当时跟我谈起,自认深信Deligne把他的结果(不提其名,也不给证明的任何指示)通报了加藤。当时这事在我看来如此不可思议,以至于我没有把Contou-Carrère当真——如今我才明白:那完全是我光彩夺目的朋友Deligne惯常的“暂停!”做派。Contou-Carrère一副真正被冒犯的神气:有人竟“斗胆猜想”他俨然视为某种私产的东西。然而他自己的猜想得自我的传授,却同样不认为有必要在三条注的任何一条里影射我这个人\fnote{19}{(*)}{关于我在这类局面里、与我的某些学生之间常常扮演的那份共谋的角色,见“暧昧”一注,n\textsuperscript{o}~63$''$。}!他对我说来,想必理所当然;而Deligne对他做的单纯的动念却冒犯了他——但他不敢向当事人吐一个字。(我曾力劝他向Deligne说清;他却小心没有做……)

这些年来,我料想,他必定以某种方式对自己施加了暴力,不发表那些极美的结果——做它们时,他想必全力投入。倘若他这样对自己施暴,那是出于时势的考虑:分明不利于这一类Grothendieck学问。近几天他十分惊讶地收到同一位Deligne的一封信:他(不动声色地!)惊讶他没有发表关于“总体”Jacobi簇的那篇注,并向他索要他关于这一主题、甚至别的主题的一切。Zoghman~Mebkhout几天前已告诉我:Deligne正在使用这些东西,甚至在这一语境里点了Contou-Carrère的名。分明,时候快到了:Contou-Carrère终于将认下一个属于他的孩子——一个他谨慎地埋了近五年的孩子。也许,谁知道呢,两位“学生-仇敌”和解的钟点也已敲响:我的学生中最光彩的两位,一位是受勋的院士,一位是代理助教;然而(无论和解与否),久已是一对\textbf{兄弟}。

\subsection{棺材4——或没有鲜花也没有花冠的topos}

\notehead{96}{}

(5月22日)倘若我说我从未见过Olivier~Leroy,我几乎没有夸张。可以肯定的是:从他听说我那一刻起,他就决定像躲瘟疫一样躲我。他的缘由,我承认,向我滑脱。也许一个本能告诉他:我只会给他招麻烦;也许Contou-Carrère(他长久以来与他过从甚密)向他提过醒——我恐怕永远无从知道。但我到底有幸与他有过两次有实质内容的交谈;我记得很清楚。

第一次应当在1976、77年:我和Contou-Carrère不打招呼地登门看他,想聊聊数学——是否心里另有盘算,我说不好。不过,想必有言在先:奥利维耶考虑读第三阶段博士,而我袖子里有的是题目。因为在Contou-Carrère家见过他一两次,加上Contou-Carrère自己向我透露的,我有一个印象:奥利维耶悟性快,而且不只在数学上。三人那晚令人难忘。我很快向奥利维耶提起一个纲领:一个topos基本群的理论、以及topos框架里的Van~Kampen型定理;他一副感兴趣的样子。经由Contou-Carrère的代数几何研讨班,他想必有一点topos的常识;他似乎乐于有机会在一个具体理论的例子上“练练手”,用topos的语言。足足两三个钟头,我把我所见当发展的理论的周详总施工图倾倒给他;它随我谈论而渐渐丰满,同时,代数几何与拓扑的一大群具体局面在我身上浮起——那些局面,要在topos框架里表述;而每一回,我都得先向一个头一回听说的人“复述”\dmark{}\origpage{410}。那晚不止一次,Contou-Carrère(他几乎无所不读、胃口极好)眼神茫然、掉了队——连他也觉得一次太多;不止一次,我谨慎地问奥利维耶,是否今天到此为止、改日再续。我原可省了这份心——分明奥利维耶神清气爽,眼神明亮,完全自如;我甚至笑起来,因为难以置信他竟不垮——可他全然不垮!那是个也许二十岁的年轻人,只有一点概形的常识、一点拓扑与topos;我相信他倒也摆弄过不少无限离散群……老实说,那是三脚猫的底子;而凭这个,他竟能填上所有的空白,毫不费力地“感到”我这个老资格、凭十五年与这一主题的熟悉、在两三个钟头里全速讲述的东西。我从未见过这样的事——至多在Deligne身上,也许还有Cartier;Cartier年轻时在这一条线上也相当出众。

总之,分明是拍板成交:奥利维耶将就上述题目做他的第三阶段博士论文。然而他想必不曾料到路的尽头有什么。总之,在其后他埋头苦干的两年里、乃至之后,我再也没有见过他。他的名义导师是Contou-Carrère,当然;但有机会同这样一个机灵的人聊聊,本会使我高兴。实际上,连答辩的通知都不曾发我,而且我相信从未收到过那份论文的副本——但我记得亲手捧读过一份,是在一个有资格得到它的人手里\fnote{20}{(*)}{这场捉迷藏,愈发离奇,因为毫无疑问:在整个朗格多克,能与Contou-Carrère一道、对Olivier~Leroy的工作多少懂点什么的人,只有我。不必说,Leroy那份科学院报告的底稿,我也同样不曾亲手捧读。也许我在自欺;但我觉得:倘若我不曾被如此苛刻地隔离、以致实际上无从干预,我总归能找到办法,让那篇倒霉的短文发表——必要时经由Cartan或Serre;他们虽不在局中,却肯信任我,只要我担保工作的严肃性。(6月7日)想必很久之后我才知道Leroy已答辩,而且我太忙,无暇自问为何连通知都没人给我。直到Contou-Carrère本人的答辩之后,我才“翻然醒悟”——那场答辩,我被认定是其论文导师(x)。他想方设法,使评审委员会成员中唯独我没有资格得到其论文的正式定本!今天我终于收到一份副本——他以为(他写道)得到它“对我没有兴趣”……(x)更确切地说:有一两年,Contou-Carrère谨慎地同时耍着两位“导师”(谁也说不准……),两人彼此不知有一位“平行”导师。Verdier的导师角色,我是在最后关头才知道的:1983年春,当Verdier分明终究要他的命变得清楚时,Contou-Carrère才最终转而投奔我。}。答辩是在那篇科学院报告“沉没”之前还是之后,我说不好。我说的是那次沉没:在第33节“注记——或新伦理(1)”里,我相当周详但不点名地谈到它。经手这次沉没的两位数学家,是Pierre~Cartier(正是那位:其惊人的直觉之敏捷,在我谈起他年轻的同行的敏捷时重回我心头;而他却如此客气、怀着全部的遗憾,沉没了那份敏捷)与Pierre~Deligne——带着他那句历史性的话:这些数学“不逗他乐”。(它们在他年轻时,可是“逗过乐”的……)我还该加上Contou-Carrère本人:他没有为自己的学生抬一根手指——那会冒触怒权势人物的危险。他想必向Olivier~Leroy暗示:最好忘掉那篇不合时宜的论文。无论如何,清楚的是:Leroy分明对这一情节画了粗粗的一道——即便日后有机会发表的不只是一篇科学院报告、而是其工作的全部,我深信他也不会用\fnote{21}{(*)}{这道“粗线”的一个雄辩征兆:Olivier~Leroy两年前蒙彼利埃一个助教职位出缺时提交的申请材料里,既不提其第三阶段论文的题目,也不提曾是其导师的Contou-Carrère的名字。他也不提任何一项个人的工作,无论什么。显然,他当时还没决定要不要那份职位——于是,尽管其天赋惊人,职位给了另一位候选人:材料扎实、意图毫无疑义的一位。}。这一回,电锯又干得漂亮\fnote{22}{(**)}{一个有趣的巧合:最近我得到回音,Cartier有意把其一场Bourbaki报告献给我(我相信这是头一遭),而且那份报告恰恰专论topos理论——正是这些topos,被同一位Cartier判定不配进入一篇科学院报告。这莫非是最近这几年时尚风向的转变?肯定不是;而且一切都环环相扣:那份报告谈的是topos在逻辑中的使用!我的朋友Cartier这份感人的题献,在我看来,与去年在一个重大场合所致的葬礼悼词同一血脉(见“葬礼悼词——或恭维”一注,n\textsuperscript{o}~104):在那里,“topos”一词被提到(在一串抖得利落的恭维之中),紧接着急忙补上(作为唯一而雄辩的评语)它们“如今用于逻辑”——别处没有,自不待言,至少只要我那些乐善好施恭维的朋友们、凭其手中的权力,还能阻止的话……(Cartier报告的出处:Catégories, logiques et faisceaux, modèles de la théorie des ensembles, Bourbaki研讨班n\textsuperscript{o}~513,1978年2月。)我在某些人(那些定调子的人,如Deligne、Cartier、Quillen……)对topos在几何中这类革新而深刻的观念的那份居高临下的(与抵制的……)态度里,感到一种惊人的狂妄。且假定他们中哪怕一位有那份材料(或那份天真……),能像我经由引入平展topos与晶状topos、从虚无中引出代数簇的一种新的拓扑视野(并由此,在拓扑到来之前,引出代数几何与算术深刻革新的手段)——毫无疑问,他乐于在自己身上培养、在他人身上煽动的这份轻蔑的态度本身,就会解除那视野与革新的力量,只为一份自负效劳。}。

\dmark{}\origpage{412}尽管有这场厄运,1981年初的那几个月里,我到底有幸常见Leroy。那是一场我当时主持的微研讨班,论Teichmüller塔的代数-算术理论(《纲领概要》里略有谈及)。真正意义的听众,只有Contou-Carrère与Leroy。即便对一个超选的巴黎听众(我懂我在说什么),一整间屋子也未必有三四个不掉队的。老实说,我办那场研讨班——当时Contou-Carrère正全力打磨其Schubert闭链的想法——是为了Leroy,想着他或许会对如此辉煌的题目咬钩。分明他“感觉得到”我在做什么;但他(我相信)事先已决定不“咬钩”。奇怪的是他甚至肯费心来——某种东西必定使他着迷,一如我着迷;而他自己对究竟要什么,也不甚了然。当我明白他不会咬钩,我便止损了。面对两位观众(无论多么光彩)继续独白,我无意。我与Leroy的第二场也是最后一场交谈,就在那时。我相信,此后我再没有见过他。

Leroy与我之间,除七年前的那个晚上,没有过真正的数学讨论——这解释了为何对他的工作我几乎一无所知,除那份不合时宜的topos工作。他的厄运想必不曾增进他对像我、甚至Contou-Carrère或“上流”数学界其他人等的信任。我听说他在文学院办研讨班——那里有一群和气的、彼此相处融洽的数学人。他们讲授组合拓扑的想法——一个近十年来全在我射程之内的题目。既然我天性谨慎(可不是嘛!),我没有问他讲些什么,也不知道他是否打算发表。至于境况,他过着一种最不合法的生活(却既非外国人、亦无身份问题):到处带习题课,报酬(嘘……)不知出自哪些小金库,背着司库与审计法院。我想,他是否最终走数学这一行,尚未下定决心;长远看,无论有没有审计法院,这想必不是个舒坦的处境。倘若我那幅富于教益的葬礼图景——他在其中充当第四口编外棺材——能帮他驱散困惑,这一回是在知根知底的情况下,我会高兴。

\subsection{掘墓人——或全体教众}

\notehead{97}{}

(5月24日)我是顶着我身上某种迟疑,才终于决定点名提及我的某些亲近的旧友与同行——在数学界;我得以看到他们充当“掘墓人”(或“电锯”)的角色:从一开始就斩断某些身份寒微或无保障的数学家的尝试——他们试图重提我的某些观念、循其固有的逻辑加以发展,或者(如Yves~Ladegaillerie的个案)仅仅追随一种带着我的影响之印记的进路与做派。如我一说再说,这类牵连他人的迟疑——或仅仅是在未曾\dmark{}\origpage{414}征询其意见的情况下点其名的迟疑\fnote{23}{(*)}{比如,为收进一条注(注n\textsuperscript{o}~19)——其中将点名提及所有跟我做完国家博士论文的学生——我就有过这样的迟疑。我身上这份犹豫,想必来自我的许多学生对与我这个人的名字相连的抗拒;那份抗拒,想必几年前我已在未成形的层面知觉到了。我的旧学生们(带或不带引号)中,当时与我划清界限的意愿已被我清晰知觉的,只有Contou-Carrère(我恰在那时才在他身上发现它)与Deligne(在那里,此事自1968年起已相当分明,尽管我不曾疑心这份意愿将把他带向何方)。在Deligne的个案里,我迟迟不肯以“多少算个学生”的名义点名他的迟疑,格外强烈:我不愿显得在拿这样光彩的学生夸耀,而他自己则不肯让那把他与我、与我的著作连在一起的纽带显露分毫。我的反思,还使我进而明白:这条纽带,在我年轻朋友的生命与著作里所占的分量,远比我曾经猜想的要大。(6月1日)关于我身上这些蓄意的心境,见“三天之后”的3月27日(原书日期如此,与5月24日不合)的“格格不入者”一注(n\textsuperscript{o}~67$'$)。}——在《收获与播种》里并不少见。在每一个个案里,我最终检视了那份迟疑,并懂得它并无根据;其来源不是体贴,而是混乱,甚至怯懦。在我点名陈述他人的行为或态度的所有情形里(依我看),这些行为或态度绝非“保密”性质。它们关涉当事人的职业生活,连同其必然的种种后果——牵连其他同行(包括我自己)的职业生活(并由此,牵连他们的全部生活)。我点名的每一个人,对其行为与态度、及其含义的全部谱系(无论他乐于忽视与否),都与我对我自己的负有同样的责任。倘若其行为的某些后果以这样或那样的形态回到他身上——比如经由我的这个人在此充当媒介的一份公开的“问责”——他毫无根据感到被冒犯。倘若我的语言有时形象而辛辣,我的意图绝非论战,绝非冒犯或羞辱任何人,而宁可是描述事实、描述我感受它们的方式——作为一种敦促,促请每一个人(首先是我点名的那几位)从其自身一侧检视它们,而不是以这样或那样的方式把它们打发掉(一如《收获与播种》的反思之前我常常做的那样)。倘若被如此唤起的人选择感到被冒犯,那是关涉他自己的选择。这份选择可能使我难过——出自我所敬重、甚至怀有感情的人;但它不压垮我。我谈过的那份迟疑——我的事事观感中某种混乱的征兆——一经理解、并由此超越,便无踪无影。

在葬礼的反思里,没有任何时刻,我曾感到有一场针对我的著作、针对那些胆敢从中汲取灵感(而非仅仅借用工具、却讳言锻造并交到他们手里的工匠之名)的人而精心策划的宏大“阴谋”。没有阴谋;但有一个\textbf{共识}——在我所称的“上流”数学界里,它至今在我眼前毫无破绽。这个共识,除至多极罕见的例外,绝非由一份对我这个人、对我的著作的有意识的恶意喂养\fnote{24}{(*)}{我所认为的毫不含糊的恶意行为,我只在Deligne与Verdier的个案里得知。}\dmark{}\origpage{415}。仅在若干例外的个案里,它才表达为对前面各注所谈的四位“共同被埋葬者”中这一位或那一位的、毫不含糊的恶意。但肯定,这样的恶意之所以能在我昔日的某位学生身上滋蔓,之所以能不受阻拦地表达,全赖普遍共识的怂恿。

这份共识,在我的几乎全部旧友或旧学生身上,显现的不是“恶意”的态度,而是(我相信)全然无意识的机制:一种令人错愕的整齐划一、一种毫无破绽的效率,像扫麦秸一样扫走健全意识与数学家的健全本能,让位于纯然自动的\textbf{拒斥态度}\fnote{25}{(**)}{这些拒斥态度,当然从不以这样的面目出现——即便在我朋友Deligne或Verdier那样极端的个案里。在对我有意识的层面,它们几乎不可见;那个层面(如我已有机会说过的)在我昔日的朋友与学生身上,几乎总是(也许甚至总是)同情与敬重的心境(其中某位,还在尽力设法自我开脱)。这样的同情与敬重的心境,不只在有意识的“意见”的表浅层面在场,甚至在实际的好恶、以及对他人实际认识的更深层面上也在场(独立于人们力图把他人关进其中的那些形象)。我们在这里处在一个典型的\textbf{两价性}局面里(我几乎想说:集体性的两价性),粗看之下,什么也“看不见”!(可与“敌父(1)、(2)”(第29、30节)里的反思对读;在《收获与播种》里,我头一回在那里触及这个标记过我生命中许多关系的两价性侧面——不只在数学界。)然而,在具体显现的层面上(葬礼里已被充分检视),这些两价力量的“合力”,在我看来,不再有任何两价性;它分明带着“令人错愕的整齐划一与毫无破绽的效率”,以我正要进一步检视的“自动拒斥态度”的面目呈现。}。这样的自动态度,我疑心,不只由我这个人、由那些其数学“气味”多少让人想起我的人所触发——也对着任何不像已经持有某个“建制”的\textbf{心照不宣的担保}的数学家:无论他已身在其中的行列,还是俨然其中某一位的“门生”(借用Verdier笔下的说法)。我觉得\dmark{}\origpage{416}:在几乎全体数学家身上,最低限度的“数学开放”的心境——它是健全意识与数学健全本能得以启动的前提——\textbf{只对着已经持有这样一份担保的人才启动}。

这类机制,想必几乎是普遍的:不只在数学界,也在社会的所有部门,无一例外。它远远超出一切个案。倘若(依我看来)在我的个案里、以及在建制眼中俨然“我的门生”者的个案里,有某种例外的局面,那是因为:过去,我曾持有“他们自己人”之一的身份,连同对我、对“我的人”的那份惯常的“最低限度的开放”。这份身份,因我1970年的离去而被收回。或者更确切地说:经由我自己的选择——它在我离去后的年月里不止一次清楚地表达,并经由我直至今日的生存方式——我确实不再是一个“他们自己人”。事实上,我自己也不再感到自己是“他们自己人”;我离开了一个曾经共同的世界,绝无回归之念。今天,我的“回归数学”绝不是回归“他们中间”、回归建制,而是回归数学本身;更确切地说,回归一种持续的数学投资,回归对我数学反思的发表活动。

我才开始明白:我的离去,曾被我的旧友们、我的学生们感受为一种“开小差”,甚至一种“冒犯”\fnote{26}{(*)}{这样一种看待与感受事情的方式,在我的朋友Zoghman~Mebkhout身上表达得格外雄辩。正是由于这次开小差,我才要为他与上流数学界的种种挫折负责:唯有他,失去了那份曾环绕着今日乐于把他当拖鞋踢的人的“庇护”与支持。}。想必那是最简单的办法,用以打发我离去的意味、它可能在人们身上唤起的追问:借这样一份弥散的\textbf{受害}之感、借怨恨的自动反应,表达为一场\textbf{报复}之举(在意识的层面,它想必很少被当作报复、甚至当作行为来知觉):既然他与我们切割,我们就与他切割——我们不再给他与“他的人”以保留给“自己人”的“注意的自动机制”的收益——他与他们,将如初来乍到者一般,领受自动拒斥的严苛!

\dmark{}\origpage{417}局面(对我的旧友们与学生们)还更复杂:因为我不但曾是建制的一员,而且他们中任何一个人,要做一个数学家,不可能不随时使用我是作者的观念、想法、工具与结果。我不知道在我们的科学或任何别的科学的历史里,可曾有过如此难堪的悖论的例子!在这一光照下,那些电锯效应(绝不限于我的朋友Deligne)——为斩断任何发展带有我的印记的观念的意图(尽管这样的发展只会加深那份困惑)——如今在我眼前,显出被一种无情的内在逻辑所推动:自某个已然做出的选择——拒斥的选择——起的一份\textbf{必然}。同样,我看到人们四处做出的努力——把那些已进入公共遗产、人们无论愿意与否都离不开的观念、想法、工具与结果的来源,置于彻底的沉默之下——也属同一逻辑。我在Deligne那些非常扎眼的“战役”面前——他装作一项一项地把我的若干主要数学贡献的父权据为己有(或把残渣慷慨地记在某个形影不离的好哥们名下)——曾以为察觉的那份“漠然”——那绝不是漠然,而是一份\textbf{心照不宣的首肯}。Deligne不过是做着建制的集体无意识期待他做的事:\textbf{抹去}那个与所有人切割的人的名字,从而化解那个无法容忍的悖论——\textbf{以一份可容忍的假父权,替换一份真实却不可接受的父权}。

在这一光照下,主祭Deligne,显出的不再是一个按其生命的深层力量与行为塑造一场时尚的人,而分明是\textbf{工具}——凭借其“合法继承人”的角色天造地设的工具:服务于一份毫无破绽的\textbf{集体意志};那意志,执意于一项不可能的任务:把我的名字、我的个人做派,从当代数学中抹去。

我几乎不怀疑:这幅图景,就其实质而言,表达了事情的实在——至少在集体的层面上。肯定,我的“归来”——它出人意料地为一场曾进行得让所有人称心的下葬画上句点,或(即便不画句点)至少不合时宜、不可容忍地扰乱一场仿佛早已排定的仪式的进行——这份归来,将不只使主祭者中的某一位不适与不满,还将使为这场可悲场合齐聚的\origpage{418}全体教众为难!当然,我毫无概念:那鼎鼎大名的集体无意识,将发明怎样的“解数”,来打发这位可敬的亡者不期归来所造成的烂摊子——他从为其量身准备的舒适的棺材里突然爬出(不可容忍的丑闻),装模作样按他自己的方式主持自己的葬礼。但我信得过全体教众:他们会找出办法,打发掉数学大厦里这桩多出来的小小矛盾——它已经万事俱备,只差这个!

我如今,在每个人的意象与态度的层面上,已能相当清楚地知觉集体共识的反映及其总形态,知觉那份抹去、埋葬的集体意志。那是普遍使用的“两幅图景”体系——两幅相互矛盾的图景,人们同时依凭它们行事;我在《收获与播种》里,就我自己的个案,头一回谈到它。(见“功劳与轻蔑”一节,第12节。)我怀疑可有人会直言不讳、明明白白地说:“Grothendieck做的尽是冒牌数学,不提也罢,说正事吧。”照直说,那会太过分明地违背建制的公理——至少眼下如此。按事情预定的演变,二三十年后,这问题反正无从提出:届时连这个名字都无从提起——它早已被所有人遗忘。共同的战术——无论个人还是集体——是沉默:人们不想亡者,至少不作为数学家想;人们不谈他,不提他(除非不得已,借天赐的缩写SGA或EGA;且待这些引证被另一些没有亡者任何踪迹的引证替换……)。

然而,也有——想必例外——的场合,完全的沉默变得行不通。其中一个场合,我料想,曾是我的CNRS任职申请;它想必使不止一人为难\fnote{27}{(*)}{(5月26日)今天我恰好从Zoghman~Mebkhout的电话里得知:我在CNRS全国委员会的同事们为我出了一把力,为我安排了一个两年期的“接待职位”。我不知道他们是否怀着热情做的;总之,委员会里我的朋友中,没有一位把这份出力做到底——打来电话、捎来一句好话,把5月15日就已成定局的好消息告诉我。(9月)我到底经由CNRS一封8月16日的信得到通知:那是一项attaché de recherches职位,为期一年(而非两年)。}。另一场,将是《收获与播种》的预分发\fnotevariant{28}{(*)}{}{说的是由我的大学经手的一份限量印制(150份),在最近的同事与朋友中间分发。},且待其在《数学反思》第一卷里出版(倘若我的出版商不至于胆怯退缩,肯冒得罪整个科学建制的风险)。这些都是亡者本人不可容忍的出格所造成的场合:他不知好歹地走出了派给他的角色。另一个场合(对理解这场尚未被一个不守纪律的亡者搅乱的葬礼,也许更有教益),是去年“隆重”庆祝的IHES二十五周年庆典。作为“IHES的四枚菲尔兹奖章之首”,在这个庄严的场合,要把我全然略过,是困难的——尽管我在IHES存在的四个英雄年月里、为使它取得真实存在所扮演的角色,被略过了。为向我致敬而精心炮制的《葬礼悼词》——载于为此庆典印行的小册子(我已有机会两次指涉它)——依我看是这一类文体的典范:一种优雅而含蓄的方式,使人人满意地化解当代数学里的这份“小小矛盾”……

瞧,我忽然精神焕发——像开始嗅到马厩的马!将近两个星期前,我开始就这一富有教益的情节写一场反思,在一条当即得名“葬礼悼词——或恭维”的注里。经过几番该把它放在哪里的踌躇(它出自为葬礼所写的第一条注的一条迟到的脚注),看来最自然的安放之处,不是“编年”的位置,而是那场应为葬礼压轴的“葬礼仪式”。于是,不经我寻求,我三周来沿着最后三个行列——“研讨会”“学生”,以及刚刚入列的“灵车”(随葬礼的最终部分,即葬礼仪式而来)——所追的那根“线”,接上了头;这场仪式的首要印记,恰恰是那份我5月12日开始检视的《葬礼悼词》杰作\dmark{}\origpage{420},它如今是自然而然紧随本注之后的那条注\fnote{29}{(*)}{(1984年11月)由于一场意外的疾病情节,该注(n\textsuperscript{o}~104)与“本注”被一个新的行列隔开——“亡故者——仍未死去”(n\textsuperscript{o}~98--103)。}。

我终于(又一次?)触到目标了!同时,这场对一篇葬礼悼词的反思的开端,忽然取得一个新的维度。它不再只是一个强力的头脑为一桩执念效劳的精巧发明——在一场官方“重大场合”的贵宾们面前、对着他们的漠然或奉命的关注挥霍——它尤其是一份在所有场合中最微妙的场合、对着一份集体期待所做的完美而体贴周到的回答:那份期待,问的是对我的这个人应当采取的态度。倘若他那一代人中有人真正配得上全体教众毫无保留的感戴,那正是我的朋友Pierre~Deligne——以他那份绝无瑕疵的完美,履行了人们期待于他的角色。


% ================================================================
%  四 · 第三部 (第 17 章起, 原书 p.413 起)
% ================================================================
\partdivider{第三部}{葬礼(II)\\[0.35em]{\normalsize 或阴阳之钥}}{葬礼(II)或阴阳之钥}

\setcounter{chapter}{16}
\markboth{收获与播种·第三部}{}

\chapter{亡故者(仍未死去……)}

\section{事件——或身体与精神}

\notehead{98}{}

\dmark{}(9月22日)葬礼各注中日期最新的一条(几条脚注不算),是5月24日——距今已有四个月。其后的两个星期,直到6月10日,主要用于重读并补充、修改已写成的各注;其间还有Zoghman~Mebkhout一两天的来访:他来通读葬礼的全部各注,在我把它们交付打印之前,并给我他的评点。\origpage{421}我原以为,定稿会在6月初备齐,并(这还是乐观的……)在大学暑假大忙之前打好印出。我满心想在人人打点行装之前,把我的“五百页的信”寄给各位!

实际上,在我写这些文字的时刻,葬礼的文本仍未完成:一如四个月前,它还缺最后的两三条注——外加其间新添的一条\fnote{1}{}{(9月23日)实际上,看来这条预定要写的“注”,已爆成三条不同的注(n\textsuperscript{o}~99--101)。}:即我方才以下列文字着手的那一条,算是这段时间里所发生之事的一份简报。6月10日,一桩新的意外闯入《收获与播种》——这部从不缺意外的写作:我病倒了!一阵胁痛突然袭来(前一分钟我毫无预感),以不容分诉的力量把我按倒在床上。站立乃至坐姿,顿时变得十分痛苦;唯有躺卧似还可忍。真是蠢透了——尤其在此刻:一项十万火急的工作眼看就要收尾,更不必说别的!躺卧着打字,没门;就连躺着写字,也不是件轻松事……

我又花了将近两个星期——其间我竭力顶住一切、照旧工作——才向这明摆着的事实低头:我的身体已经衰竭,它执意要求——而我不肯听闻——完全的休息。

我之所以这么难以听话入耳,是因为我的精神依然清新、敏捷,跃跃欲试要继续冲刺,仿佛它有一条独立于身体的自主生命。它甚至清新敏捷到这种地步:连身体的睡眠需要都难以兼顾,一再把睡眠——这个搅局的煞星!——的期限,推到衰竭的极限。

我一生中,直到三四年前,经由深长睡眠而有的无限的恢复力,曾是那些有时失度的精力投资的坚实而救命的配重:只要睡眠有保证,人就无所畏惧;尽可以(而不失为清醒)全身心投入工作的狂欢、直到衰竭——大不了再用修复性睡眠的狂欢补回来!这份能力——我一生把它与工作能力、与发现的能力(两者想必血肉相连……)一并视为理所当然——近年终于磨损,有时甚至消失;缘由我此刻仍辨不真切,也还不曾真正下力气去探它。越来越常见的是:在一整天的打字机(或手写笔记)之后,顺从我那拒绝继续的身体之命,我决意去睡;可一躺下(连同躺卧对坐姿紧张的局部缓解),立刻重新点燃反思。它愈发精神,一连几个小时,乃至整夜(或者不如说,夜里所剩的那点时间……)。我尽管明白这套系统无利可图(姑且不说它长远是否“活得下去”),——因为(至少在我这里)没有书写支撑的长久反思,终归原地打转,常沦为一种反刍——恶习已然养成,且有加剧之势。我觉得,它成了我近些年生活中能量涣散的大本营;而别的涣散机制,却经年累月一一被清除。

倘若这一机制在我生活中扎下如此顽固的根,倘若这些年我一直肯付这样的代价,那肯定是我身上的某种东西从中得到了好处,时候一到还会再从中获益。凑近些检视这一局面,并非奢侈——四个月来,我不止一次险些动手。

\dmark{}这想必是一件紧迫的事。然而我终于懂得:还有更紧迫的。\origpage{423}我得先应付燃眉之急:重新接上与我身体断掉的接触,帮它从我终于感到并承认的衰竭状态里爬出来,找回失掉的元气。我那时懂得:为此,我必须在一段不确定的时间里,放弃一切智力活动——哪怕是沉思临我之事的意味。今天重新提笔的这些注,就是为它画上句号:我几番大投资中这段漫长而有益的“括号”——那些投资,曾有一段(自今年2月起)汇入“收获与播种”的写作。本注是对这四个月“括号”的最初一点思考,或无论如何,一份简要的汇报。

等到终于懂得需要完全的休息,大的疲惫已成了深深的衰竭。由于不曾肯听我身体那分明不容分诉的语言,最初两个星期里,我在体力疲惫的状态下硬挤出来的那几页可怜的评注与修改,是以一份回头看来近乎疯狂的精力投入换来的!总之,在这些壮举之后,我只得长卧数星期,每天只起身几小时,做些必不可少的实际事务。

引人注目的是,一旦终于\textbf{懂得}了完全休息的需要,我便毫不费力地彻底脱离了一切智力活动,没有丝毫“作弊”的念头。我甚至无须真正做什么决定——单是懂得,我就已然脱开了。前一天还吊着我胃口的那些任务,忽然显得十分遥远,仿佛属于一段非常遥远的过去……

然而当下并非空无。几个星期乃至几个月里,睡眠迟迟不肯来;我长卧数小时,表面全然无所事事,却不记得有哪一次觉得时间难熬。我与我的身体重新相识,也与最切近的环境重新相识——我的房间;有时是眼前那片晒着太阳的草皮或干草:我在屋旁、或在一次短暂(而谨慎……)的散步中随意躺下的地方。我久久地追随一只苍蝇在一束阳光里的舞蹈,或一只蚂蚁、或一些绿色或粉红色的半透明小虫在无穷草茎上的长征——那些草茎在我眼前缠成难解的丛林\dmark{}\origpage{424}。也正是这样的心境:借着寂静与大疲乏之便,人以关切之心,追随最轻微的风在自家肠肚里的踌躇巡游——总而言之,是这样的心境:人们借此与最基本、最本质的事物重新接触;在这样的心境里,人们充分掂量一场修复性睡眠的全部益处,甚至掂量“能够无碍地撒一泡尿”竟是怎样的奇迹!

身体的寒微运转是一件非凡的奇迹;只有当这运转以这样或那样的方式受到扰乱时,人们(有时违背本心)才会多少意识到它。很清楚,“技术上”说,我“健康问题”的底子,是睡眠之扰。这一扰动的深层缘由,当时向我滑脱,如今依然。我是摸索着,首先试图找回睡眠——那种我熟识的酣畅大觉;而它偏偏在我最需要它的时刻神秘地躲避我!直到不久前我才找回它。想必不必说:求助于药片的念头,根本不会来到我心头;即便我试过一些汤剂或橙花水(借这次机缘认识的),我心底也明白那只是权宜之计。更正经的一件事:我借此机会对饮食做了重要的调整:减少淀粉类,增加绿叶蔬菜与水果(生熟皆可);(适度地)重新把肉作为日常饮食的常客;尤其是,极其严格地削减脂肪与糖的消费——在这方面,我(一如富裕国家里的许多人)至少自战后以来,一直有一种系统性的失衡。帮助我认识到这种饮食调整对找回失衡生活之平衡的重要性——尤其还有每日几小时量级的体力活动的必要,以抗衡紧张的智力活动;否则后者动辄耗竭身体,把可用的生命力往上抽,造成强的阳失衡——的,是我的女婿艾哈迈德;他行中医,对这类事情有很好的“感觉”。

况且,艾哈迈德并不满足于给我一些好建议,外加一套阴阳辩证——四五年来,我既有充分的机会熟悉事物这份精微的动力学,对它相当敏感。一旦我身体好转到可以打理园子,\dmark{}\origpage{425}而他看到我亲力亲为、想把一处寒碜的小园子多少重新开动起来,他便抢先发起几项大面积的工程:开出新的地块,运土,移栽与播种,做梯田、垒挡土矮墙,重整堆肥……随着日子与星期流逝,我眼看这些整治任务在我这位不知疲倦的朋友的推动下于眼前铺展——它们足以让我忙上几年,乃至忙完余生!

这恰是我所需要的,也是长远而言我所需要的,用以配衡一种过于炽烈的智力活动。在这方面,每天强迫自己散步——如人们早就这样建议过我——帮助不大:散步时脑袋照旧研磨,一如在床上,几乎视而不见地穿过景色的美!相反,浇园——我得当心园子的长势;更好是锄一畦菜——我就不由自主地留心、并多少上心:体会土的质地,锄地如何改变它,菜蔬与“莠”草如何影响它,堆肥与铺盖如何影响它;久而久之,也体会我理应照料的那些植物的境况——那境况,在很大程度上,映照着我肯给它们多少留意。园艺这项活动,以及环绕它的一切,回应我身上两股强烈的渴望或心性:其一推我去从事一种每天能看见\textbf{手底下出东西}的行动(散步绝不是这样;某位同行兼朋友建议的举重哑铃,更不是……);其二推我去从事一种每时每刻都有机会\textbf{学习}的行动——在与事物的接触中学习。我似乎恰恰在“做”某件事情的境况里最善于学习——某个在我手底下成形、变形的“某件事情”……

一旦跨过严格意义的衰竭状态,我的康复——依我看——得益于两类活动;或不如说,得益于日常活动(屋里的与园里的)中两类重要而有益的因素。其一是\textbf{体力的付出}:哪怕动工之前我常觉疲惫、没劲——活越是“重”,叫我抡起\dmark{}\origpage{426}重镐或搬动大石头的——过后我越觉精神,带着一份沉沉的好疲惫。其二是与\textbf{活物}的接触:需要照料的植物;需要为它们备好、随后铺盖或锄松的土地;需要烹备、而我吃着与备餐时同样有滋味的食物;讨食、也讨它那份抚爱的猫;还有各样器皿与工具,乃至那些磕磕碰碰、常不甚规整的石块——得把它们翻来转去,垒成肯站得住的矮墙……

体力的付出与活物的接触——这恰是智力工作所缺的两个方面;正因如此,智力工作按其本性是不完整的、残缺的,走极端时——若不以别的东西补足与抵偿——甚至危险乃至有害。这是短短三年多来我第三回有机会认识到这一点。如今甚至十分清楚:我面对一道苛刻的期限:改变某种生存方式,找回一种平衡——使我存在的阴极、我的身体,不再被持续怠慢,而一味让给阳极、精神(或更直白地说,脑袋)——否则,就在最近几年里把命搭上。这是我的身体对我说的,说得再清楚不过!我如今到了生命的一个节点:某种基本“智慧”的需要,已成了一个\textbf{活命}的问题——按这个词的本来与字面的意义。这肯定是好事——否则,所谓“智慧”将永远被推给“下回”,让位给那种智力活动上的贪食症;那种贪食,是我整个成年生活中占支配地位的力量之一。

面对这样一道分明的期限——“改变,否则就死!”——我无须自问选择为何。正因如此,将近四个月里,我得以在从不觉得勉强自己的情况下,戒绝一切智力活动——数学也罢,非数学也罢。我没有对自己说,却知道:说到底,一个活着的园丁,总比一个死去的数学家(或一个死去的“哲学家”或“作家”,反正一样!)值钱。带点促狭,还可以补一句:甚至比一个活着的数学家还值钱!(不过,那就是另一段故事了……)

况且,我不认为哪天自己会被逼进这样一种“极限”境地:长远放弃一切智力活动——数学的也罢,沉思的也罢。毋宁说,未来这些年里最直接、最紧迫的实际任务,依我看恰恰是达到一种生活的平衡:两类活动——身体的与精神的——日复一日共存,谁也不变得吞噬一切、排挤对方。\origpage{427}我不自欺:自我童年起,我最强的投资正在“精神”这一侧;今天仍支配我生活的两大热情,也仍把我引向那一侧。这两大热情——数学的热情与禅观的热情——依我看,尤其是前者,若非唯一地,充当了我生活中的失衡因子:像某种至今仍带着那份恼人的“吞噬其余、独占一切”的倾向的东西。肯定并非偶然:1981年6月以来,我生活中标志着失衡局面的三桩“患病情节”,恰恰都落在数学热情当台主演的时期。

人们会说,最后这桩情节——它发生于《收获与播种》的写作之中——并不完全如此:那部书构成一场对我自己的反思,甚至不失为一场名副其实的禅观。但同样真确的是:这场对我数学家过去的反思,一直由我的数学热情喂养;尤其在那第二部分——葬礼——依我看,这份热情的自我本位成分,以格外强而持续的方式被牵扯在内。然而,即便回顾,我并不觉得这场反思曾在任何时刻采取一种吞噬性的、乃至疯狂的节奏与音高,一如前两次——那时我的身体终究被迫喊出一声不容分诉的“受够了!”。与我的一生分割开来看,我一年半以来的智力活动(自“重启”以来——先《追寻栈》的写作,继以《收获与播种》)显得以相当合理的节奏进行,没有忘了喝水,也没有忘了吃饭(不过有时,稍稍,忘了睡……)。倘若它终于酿成第三桩“健康情节”(用个委婉语),那无疑是植根于整整一生:那永恒的失衡——一颗过强的脑袋,把它的节奏与法则强加给一个壮健的、长期默默承受的身体\fnote{2}{(*)}{这里我应当把1974至1978这五年算作例外:那几年没有被什么大任务主宰,手工的活计占去了我时间与精力的不小一部分。}。

\dmark{}\origpage{428}过去两个月里,我充分体会了身体劳动那无可替代的益处:与朴素的活物亲密接触,听它们默默讲述书本或单凭反思无力教给我的简单而根本的事。借着这样的劳动,我重新找回了睡眠——这位比吃喝还要珍贵的伙伴;也随之恢复了元气,重获那一度仿佛突然消失的健壮。我还发觉,在我目前的人生阶段,若想把去年开始的这场新的数学冒险再继续几年,就非得把双脚稳稳扎在园子的沃土里,否则便会危及我的健康与生命。

未来几个月,我得建立一种新的生活方式,使身体的劳动与精神的劳动每天都有一席之地,彼此调和。要做的事多着呢!

\section{陷阱——或省力与衰竭}

\notehead{99}{}

(9月23日)昨晚我不得不强行收住,免得顺着势头写到凌晨两三点,重新落回那副我熟得不能再熟的齿轮。我神清气爽;倘若顺着自然的坡道,我甚至会写到拂晓!智力工作的陷阱——至少那种怀着热情、在一门久经熏陶终于如鱼得水的材料里从事的智力工作——就在于它\textbf{容易}得令人难以置信。拉吧,拉吧,东西总会来;只消拉就是;有时才勉强感到一点吃力、一点摩擦,表明它多少有点抵抗……

然而我记得,在我做数学家的年轻岁月,有一种挥之不去的沉重感、滞重感,须以顽强的努力去克服,努力过后留下一阵疲惫。那主要对应于我一生中工具不足、乃至不顺手的时期;或对应于稍后的另一段:当时我不得不多少痛苦地、在一种(主要是Bourbaki小组的)环境的压力下,习得一套多少“四面八方”的工具——那环境熟练地使用它们,而它们的所以然却不曾随时显现给我,有时甚至多年如此。我已有机会在《收获与播种》第一部分里,谈到那些有时颇为费力的年头(见“受欢迎的外来者”第9节,与“炉中百铁,或:干熬无济于事!”一注,n\textsuperscript{o}~10)。那主要是1945至1955年,与我的泛函分析时期重合\dmark{}\origpage{429}。(我觉得,在我1960至1970年间的学生们身上,这种对缺乏足够动机的学习——凭长辈的权威硬吞观念与技术——的抵抗,比我当年弱得多;干脆说,我全然没有觉察到。)

回到本题:尤其自1955年起,我常常有“飞翔”之感——做着数学,毫不费力,一如我昔日曾那样艳羡过的某些长辈,艳羡他们那份近乎奇迹的从容——那份从容,对我这个寒微而滞重的身量来说,显得全然高不可攀!今天我明白:这样的“从容”并非某种特殊天赋的特权(这样的“天赋”,我曾在某些人身上、在我自己毫无此种“天赋”的时刻见过);它宁可说是自行显现为一份果实:一份对某门材料(比如说数学)的炽热兴趣,与对其或多或少的长期熏陶之结合的果实。倘若“天赋”确在这份从容的出现中插了一手,那想必是经由时间这个因子——或长或短,因人而异(有时在同一个人身上,也因时机而异,这是真的……)——才得以在某一题目的工作中抵达完美的从容\fnote{3}{(*)}{然而,我认识好几位数学家,各自做出过深刻的工作,却从不曾给我这里所说的那份从容、“轻便”的印象——他们仿佛时时与一种无所不在的滞重搏斗,每走一步都得费力克服。由于这种或那种缘由,方才所说的那份“自然的果实”,在这些杰出的人身上不曾“自行显现”,尽管它本该如此。可见并非所有的结合都结果实……}。

总之,年岁越往前,我做数学时这份“轻便”之感越强——仿佛事物只消我们肯费心看一看、稍稍端详,便只等着向我们显露自己。这不是技术精湛与否的问题——分明,从这个角度看,我的状态比1970年“退出数学”时差得多:此后我有机会做的,主要是“忘却”当年所学;只在私下里偶一为之“做做数学”,而且怀着一种与往昔(至少乍看)十分不同的精神、就着十分不同的主题。我也不是说:只消我肯去啃某道名题(比如Fermat的、Riemann的或Poincaré的),就能在一两年\dmark{}\origpage{430}乃至三年之内\fnote{4}{(*)}{[原书此注文字与注3完全相同,照录。]然而,我认识好几位数学家,各自做出过深刻的工作,却从不曾给我这里所说的那份从容、“轻便”的印象——他们仿佛时时与一种无所不在的滞重搏斗,每走一步都得费力克服。由于这种或那种缘由,方才所说的那份“自然的果实”,在这些杰出的人身上不曾“自行显现”,尽管它本该如此。可见并非所有的结合都结果实……},笔直地开辟一条通往其解答的路!我所说的那份从容,不是那种通往一个预先定好的\textbf{目标}——证明某条猜想或给它一个反例……——的从容;它宁可说是那种使人得以跃入未知的从容:朝着某个方向——一个幽暗的本能告诉我们它是丰产的;带着一份亲身的把握——它从不会被辜负:旅途中的每一天每一时,都不会不带来一份新的认识的收成。\textbf{怎样的}认识恰好等着明天,乃至今天的下一个钟头,我们固然有所预感;而正是这份不断被赶到前头的“预感”、与它浑然一体的悬念,不断把我们抛向前方——此时,我们所开掘的事物本身,仿佛在把我们吸进去。已然知道的,总在精确、滋味与丰富上超出所预感的;而这份已知,又旋即成为起点与材料,供一份焕新的预感腾身而起,追逐一个渴望被认识的新未知。在这场发现事物的游戏里,我们每一刻所循的\textbf{方向}是已知的,而\textbf{目标}却被忘记——假如我们确实曾从一个目标出发、曾打算抵达它的话。这个“目标”,其实当时是一个\textbf{出发点},一份野心或无知的产物;它尽过自己的职责:推动“老板”,定一个初始方向,并引发这场游戏——在这场游戏里,目标其实没有份。只要踏上旅途不是为了一两天,而是长程的,那么时日与月份的流逝将向我们揭示什么、在一长串未知的波折尽头将把我们引向何处,对旅人是一个全然的谜;一个如此遥远、老实说如此可望而不可即的谜,以致他不大操心!倘若他有时凝望地平线,不是为了那件不可能的差事——预测一个终点,更不是为了随心所欲地裁定它——而是为了定位他此刻正到哪里,并在供他继续旅程的诸方向之中,选出从今往后他感到最灼人的那一个……

这就是我方才谈及的“难以置信的轻便”——就一门全然智性的工作(如数学)中的发现而言。它既不被\textbf{内在的抵抗}\fnote{5}{(*)}{然而,我认识一位天资出众的数学家,他与数学的关系是典型的冲突型的:每走一步都被强大的抵抗所阻滞,比如害怕某一期望(比如以一条猜想的形态)到头来被证明是假的。这样的抵抗,有时可以酿成一种真正的智力瘫痪。可与前一条脚注对读。}所\textbf{阻滞}(一如它在我所实践的禅观工作中如此常见),也不要求付出一种\textbf{体力}——那体力会产生疲劳,而疲劳终于发出毫不含糊的停止信号\dmark{}\origpage{431}。至于\textbf{智力}的付出(假定到了只剩“时间因子”这一层“抵抗”的地步,还能谈“付出”的话……),它似乎既不产生智力的疲劳,也不产生身体的疲劳。更确切地说:即便有身体的“疲劳”,人们也并不真把它当作疲劳来感受——除非偶尔的腰酸背痛(因长时间保持一个固定坐姿),以及诸如此类的附带麻烦。这些,换个姿势便轻易消除。而躺卧的姿势有一种不合时宜的功效:叫它们烟消云散,从而促成智力工作的重新开工——而不是那份亟需的睡眠!

然而,我终于明白:有一种比肌肉或神经的疲劳更精微、更阴险的身体“疲劳”——后者以无可辩驳的休息与睡眠之需,显出自身。这里的字眼,与其用“疲惫”,不如用“衰竭”更切;当然,要说定:这一状态并不按这个词的通常意义被知觉——通常意义指极端的疲惫,表现为光是起身、走几步路都须巨大的努力。这里说的是身体的能量为大脑所付出的一种“衰竭”;它表现为身体总的“张力”、其生命力水平的逐步低落。看来,由过度的智力活动(我是指:不以足够的身体活动——那种产生身体疲劳与休息之需的活动——来补偿的智力活动)造成的这种衰竭,是渐进而\textbf{累积}的。这些效果,同时取决于给定时期内智力活动的\textbf{强度}与\textbf{时长}。以我从事智力活动的强度,以我的年纪与体质,这一累积衰竭,似乎在一两年不间断的活动之后——没有规律的身体活动作补偿——就抵达一个危险的临界阈值。

在某种意义上,我所说的这份“轻便”只是表象。紧张的智力活动动用可观的能量,这很清楚:能量在某处被取用、“花费”在一项工作上。这个“某处”,似乎就在身体的层面:身体竭其所能“照单收下”(或不如说,\textbf{垫付})脑袋那不计成本地自开的(有时令人眩晕的)开销。身体所供能量正常的回收途径,是睡眠。脑袋一旦变得贪食,终于就要侵占睡眠——这就等于吞噬一份不再更新的能量资本。\dmark{}\origpage{432}智力工作之“轻便”的陷阱与危险在于:它孜孜不倦地怂恿我们越过这道阈值,或越过之后停留在外;而且这一越界,并不以疲劳、乃至衰竭的那些惯常而确凿的征兆,向我们报警。我明白:要察觉这道阈值的临近与越过,需要极大的警醒——而那时,人正全身心投入一场迷人的冒险。在身体的层面知觉这份能量的亏空,要求一种对身体的倾听状态;我常常欠缺它,也很少有人具备它。况且我怀疑:这样一种有意识的注意与身体之间的契合状态,能在任何人身上、在其生命中一段纯然智性活动独占、排斥一切身体活动的时期里绽放。

况且,许多智力工作者凭本能感到这种身体活动的需要,并相应安排生活:园子、木工、山、船、运动……像我这样,为了一个过于侵占的热情(或一种过强的倦怠)而忽视这份健全本能的人,迟早要付账。三年里,这是我第三回结账;我必须说,我付得并无怨言——或者更确切地说,心怀感激:每一桩新的患病情节都让我明白,我不过是在收割自己疏忽的果实;而且,它还带给我一份教益——那份教益,恐怕唯有它能给我。这最后一桩——它刚刚结束——所给我的主要教诲,也许是:是时候抢在前面,叫这样的“提点”从此不再必要;或者更具体地说:是时候照料我的园子了!

\section{向Claude~Chevalley告别}

\notehead{100}{[原书此处误印作“注~97”]}

昨天与今天的反思里,我起初有意搁下了一桩事件:它正落在患病情节的正中,在7月的最初几天——也就是说,当时我还卧床。那就是Claude~Chevalley的逝世。

我是经由《解放报》上一篇多少以此为题的含糊文章得知的;一位女友随手递给我,想着也许我会感兴趣。关于Chevalley,文中几乎什么都没有;倒有关于Bourbaki的几大段——他是其创始成员之一。得知消息时,我感到自己全然发懵。几个月来,我总看见自己眼看就要做完《收获与播种》——打好、印出、装订,一应俱全——然后火速上\dmark{}\origpage{433}巴黎,把一份还热乎的副本送到他手里!倘若天底下有一个人,我确信会怀着真正的兴趣、且常常怀着乐趣读我那块铺路石,那就是他;而且我全然没有把握,除他之外还有没有第二个!

自反思之初,我就发觉:Chevalley在我行程的一个关键时刻,给过我某种东西——某种播于一场沸腾之中、却在寂静里发芽的东西。当时我感到把我与他系在一起的那种东西,与其说是一份\textbf{感情}——比如感激、同情或爱——不如说是别样。这些感情当然在场,一如它们对那些早在二十多年前把我当作“自己人”接待我的“长辈们”中的这一位或那一位,也在场。使我与Chevalley的关系有别于我与他们中任何一位、与我的大部分朋友(不妨说与所有朋友)的关系的,是另一样东西。那是一种感觉,或者更恰切地说,一份知觉:一种\textbf{本质的亲缘}——超越文化的差异、超越自我们少年时代起便打在我们身上的种种烙印。我说不好,在我反思中谈到他的那几行文字里\fnote{6}{(*)}{见“与Claude~Chevalley相遇——或:自由与善良的情感”(第11节),及下一节“功劳与轻蔑”的最末一段。},是否透出这份“亲缘”的些许。在我生平那段——即那几行所指的——时期里,Chevalley显出的,也许更是一位“长辈”——这一次,是在对生活中某些基本事物的领会层面上——而不太是一位“亲人”。然而这份距离,我后来的成熟想必把它缩小、乃至取消了;一如在数学层面上,这早已是旧事——在我与他的关系中,一如在我与其他长辈的关系中。倘若我此刻试图以言词勾勒这份亲缘的意味,或至少它的一个征兆,来到心头的是这个:他与我,都是“独行的骑士”——各自走在自己的“孤独的历险”里的旅人。关于我自己的那一场,我在《愚妄与更新》最后那一章(同名)里有所陈述\fnote{7}{(**)}{这个方向上,尤其见“禁果”与“孤独的历险”两节,n\textsuperscript{o}~46、47.}。也许,对那些深知Chevalley的人(乃至对别人),反思的这一部分,比专门谈他的那一部分,更能暗示我想表达的东西。

与他相见、并与他稍稍谈谈,本必定使我比过去更好地把握这位朋友,更好地安放这份本质的亲缘,以及我们的差异。倘若除了Pierre~Deligne之外,还有一个人让我急于把《收获与播种》的文本亲手交到他手里,那正是Claude~Chevalley。\dmark{}\origpage{434}倘若有一个人的评点——无论是顽皮的还是辛辣的——会对我有特殊的分量,那还是他。7月的头一个星期的那一天,我知道:我不会有这份乐趣了——再不能把我最好的东西带给他,也再听不到他的声音了。

那件怪事——它想必助长了我得知消息时的那份\textbf{发懵}——是:在此前的几个月里,不止一次,在想着与Chevalley的一场临近的会面时,我记起他正受着健康的困扰;我心中有一种被不断推开的不安:那场会面也许不会有了,我的朋友也许会在我去看他之前消失。给他写信或打电话的念头当然掠过心头——哪怕只为问问他的健康、他的近况,对他说几句我正投入的那项工作、以及我打算为此去看他的事。而我把这个念头当作愚蠢、当作打扰推开了(反正真没有什么理由……等等),一如人们在诸如此类的情形里常常做的那样。这恰好例证:我这个人,一如许多别人,仍未能充分运用自己的能力——推开那股对事物的幽暗预知;它向我低语一份认识,而我太忙、也太懒,不肯去听……

\section{表面与深处}

\notehead{101}{}

(9月24日)前两天围着“患病情节”兜了两天圈子之后,该是接上6月间中断的线头的时候了——就从我搁下的地方接起。我那时预计,还有最后两条注要写:一条“葬礼悼词(2)”(它将接续并补全5月12日的“葬礼悼词(1)——或恭维”),以及一条终曲“De~Profundis”,我打算在其中为围绕葬礼的整场反思勾勒一份盘点。

这两条注预定要写的内容,在我病倒的那一刻还热乎着——我眼看就要把它们统统倒到纸上,只差把最后一手加到前面所有各注上,好有这种在“稳固而齐整的后方”上工作的感觉……在其后整整三个月里(自6月23日确切算起),我实际上停止了葬礼上的一切工作,只除偶尔的几处排印更正;它(唉)有点从我脑子里溜走了。我甚至觉得自己有点傻——总之有点难为情——竟一本正经地坐下来,去填那些在“罚抄标题”后面等着的一页页白纸,只凭这些标题出现在一份临时的目录里,而我曾轻率地在某份注定要发表的文本里此处彼处提到过它们。\dmark{}\origpage{435}尤其是“葬礼悼词(2)”的情形;就连方才重读初稿“葬礼悼词(1)”(别名“恭维”),也不曾足以焐热那份在它的角落里从容冷却了几个月的实体!

然而,就在写下那条注的5月12日的次日,及其后整整一个月里,我的手一直发痒,想更深地开掘那座我刚刚不知不觉按上手的新矿。去年,当Nico~Kuiper好意给我寄来IHES建所二十五周年的小册子时,我花了大约半个钟头翻了一遍(包括关于Deligne与我各半页的两篇短文),没发现什么特别的。唯一打动我的,是通篇没有一句影射IHES那艰难的最初年月——那时它的名声是在一间凑合的屋子里立起来的,而我(连同最早的几届代数几何研讨班)是唯一在“实地”代表它的人。几个月后,1984年3月,我写“有益的挣脱”一注(n\textsuperscript{o}~14)时,又想起了它。因为不敢相信自己的记忆,我出于求个心安,请尼科再寄我一册那本小册子(第一册已找不着了)。这是第二回翻那两篇短文,眼睛也许不那么匆忙了。然而,这一回我又没“入局”,实在是。顺带我多少惊讶地注意到:关于Deligne的那篇里说,“其工作的主导轴心是‘理解代数簇的上同调’”——谁信呢!而这事,我一忘就是一两个月(直到写“拒受遗产——或矛盾的代价”一注,即n\textsuperscript{o}~47时,我才被引得想起)。相反,我却没有觉察:在那篇谈我的短文里,“上同调”一词不曾出现,“概形”一词也一样。在我当时那种心不在焉的状态里,没有任何东西使我疑心:这篇平淡无奇、堆满夸张形容词的文本,竟充当着一份葬礼悼词的角色——而且是(外加)“以完美的手法呈献的”!手法完美到这种地步,以致我怀疑:这本小册子的读者们(它读来有点无聊,因为满纸都是场合所要求的、四面八方的蓄意涂脂抹粉,想必如此……)是否有人比我更早看出这一点——在我第一遍、第二遍阅读的时候。

这立刻汇入一条我不断重温的确认:每当我由于这样或那样的缘由,得以以一种\dmark{}\origpage{436}多少强烈而持续的注意,去端详某样我先前只满足于以“惯常的”、例行的注意“顺带”看过的事物——那些日复一日在我生活中掠过的大事小事——就会出现这样的局面。这种局面在禅观时期频频出现;禅观屡屡(而且多半“顺藤摸瓜”、并无蓄意)使我得以以前所未有的仔细,去检视某些白日或黑夜的事件(梦也在内):它们在我惯常的注意状态里或多或少未被察觉,或者其意味(往往清楚而分明)起初全然逃脱了我有意识的注意。

当我在这里说“多少强烈而持续的注意”,我心底所指的,其实是一双\textbf{醒着的眼睛},一双新的眼睛——不被思维的积习压重、也不被充当门面的“所知”压重的眼睛。只要由于这样或那样的缘由,我们得以用醒着的、专注的眼睛去看事物,它们便仿佛在我们眼前变了。在我们日常“注意”所呈递的事物那呆板光滑的表面的貌似平庸背后,我们骤然看见一道未曾料到的\textbf{深处}敞开、活跃起来。事物的这份深处的生活,不曾等我们费心去认识它才在——它无时不在;它是它们内在本性的一部分,无论所论的是数学对象、园子里的一片草坪,还是某个时刻在某人身上起作用的全部心灵力量。

\textbf{思维}是揭破并容我们探测表面之下的深处、事物的隐秘生活的一件工具——说它“隐秘”,只因为我们太懒,不肯看;太受抑制,不肯看见。它是一件有其长处、也有其短处与限度的工具。但无论如何,思维很少被当作发现的工具来使用。它最常见的功能,不是发现在我们身上、在事物之中的隐秘生活,而恰恰是遮掩它、把它凝固化。它是一件多用途工具,同时供“劳作的孩童”与“老板”使用。在这一位手里,它化作一面帆,善于捕捉我们欲望的风,把我们送往未知深处;在另一位手里,它凝成一只纹丝不动的锚,任漩涡与风暴都撼它不动……

反思本来有点走岔了;瞧,它回到了一个出发点——也就是昨天我停留过的那条确认:积习与烙印何等根深蒂固,使我远未发挥\dmark{}\origpage{437}自己的能力!(在这件事上,我还发现自己同伴有的是……)正是借着对葬礼的逐步发现——从LN~900那部文集这样扎眼的事实出发\fnote{8}{(*)}{见“一个梦的回忆——或Motif的诞生”一注,n\textsuperscript{o}~51,及随后的“下葬——或新父亲”一注。}——一份懒惰的注意终于醒了过来。读“拒受遗产——或矛盾的代价”一注(n\textsuperscript{o}~47),使我在5月12日第三遍(!)重读那两篇大名鼎鼎的“短文”。这一回,我总算注意到一个有点出格的细节:在那本庆典小册子里献给我的那篇酒神颂风格的小文里,通篇绝口不提“上同调”(也不提代数簇或概形)!这事足够滑稽,值得一脚注;我当即写了起来。写着写着,我又发觉一两处别的“滑稽”细节,它们还不曾咬住我的注意:那虽是第三遍阅读,它照旧停留于表面、机械——差不了多少,我不过是\textbf{重复}、\textbf{复制}此前的阅读而已。只是在写那本该是一条脚注、后来变成“葬礼悼词(1)”一注的文字时,我才渐渐被咬住了钩;一份\textbf{好奇}醒了,使我再一次回到那两篇文本上,这一回看得稍近了些。直到那一刻,方才所说的那种转变才发生:一道“深处”敞开了——在一场重大场合的锦衣华盖之下端出的酒神颂的平板门面背后,有一种炽烈的生命!正是这份好奇,把一道机械的、重复的、心不在焉的目光,变成一道“醒着”的目光……

而且,所说的“觉醒”也非一蹴而就;它是渐进的,随着那条(真不愧是)脚注里继续的反思的进程而成。老实说,直到那条注画上句点,它还不曾完成——尽管(我记得)当时夜已深,催我“收尾”\fnote{9}{(**)}{这一点想必尤其如此:就在当天,我已先经历了那场漫长而充实的“屠杀”反思(n\textsuperscript{o}~87);而且紧随其后的“葬礼悼词——或恭维”一注,在其结尾处还指涉了它。}。可我前脚刚画上那个点——或者至少到次日——我就发觉:葬礼悼词这个题目,我还远没有穷尽。\dmark{}\origpage{438}直到那一刻,我才\textbf{充分}感到:这两篇看似如此短小、如此平淡的文本,含义何等丰饶;干脆说,是两座真正的矿!而且,只要我肯倾听,它们要说的话,我还远没有巡完……

(9月25日)昨夜,反思又不得不中途收住——虽说它在我感觉里才刚刚起步。然而我已连续在打字机前坐了三个半钟头;一些隐约的小征兆开始向我示意:该起身动动了。

我清楚地记得头一回:我得以把一份“强烈而持续的注意”投向成文的文本,并且日复一日、一连数月,亲历一场“表面”的惊人蜕变——呆板而平板的表面活了起来,显出丰富而精确的意味,一道未曾料到的“深处”。那同时也是我头一场长程的禅观:怀着一场未知之旅的精神,旅程要多久就多久……出发的材料,是我父亲(流亡巴黎)与我母亲(其时还在柏林,带着当时五岁的我)之间1933/34年间的浩繁通信。我的目的,是“认识”我的父母。前一年我发现:我一生对他们怀有的那份敬佩——它终于凝成一种孝顺的虔敬——掩盖并维持着对他们的一份极大的无知。这份惊人的无知——我一生乐于把自己保持在其中——还是在那之后一年、1979年8月至1980年3月的那场长程禅观里,才向我显出其全部维度。

我于1979年7月间开始“整地”:其中主要是对这宗通信做第一遍通读,在我当时为一部“我自撰的诗作”\fnote{10}{(*)}{这部著作、以及它所代表的我生平的那段情节,见“非古鲁的古鲁,或三足马”一节(n\textsuperscript{o}~45)的结尾,及它所指涉的注n\textsuperscript{o}~43.}收尾之余进行。每天晚上,我花几个小时读三四封回信;兴趣当然是有的,而且当时我会毫不犹豫地说:读得很专注。然而我隐约发觉:我依旧陌生、外在于我所读的东西——真正的意味从我这里滑脱。我所读的常常相当离奇,\dmark{}\origpage{439}仿佛我眼前起居、招摇的这一男一女,与我自以为认识的那些人毫无瓜葛——我的记忆向我返还一幅清晰、分明、不可动摇的图像的那些人。由于缺乏一项耐心的、一丝不苟的、苛求的工作——随着前进逐步展开的——我只有震惊而已,为这些书信里(相对)不多的几处“够分量”咬住我表面注意的东西。如此登记下来的东西,只是叠加在“众所周知”之上;而那个“众所周知”,自我幼年起、直到那些天,始终是(当然,我从不曾察觉)我生活的、我的身份感的无形而不可动摇的基石。假定我当时就止步于这第一遍阅读,那么薄薄一层未消化的新“事实”——它就这样叠加在主宰层之上——想必很快会被磨蚀、冲走,在其后的岁月里不留多少痕迹。

在那项预备性工作之际,我的主要投资在别处——在一部吞噬了我大部分精力的著作的写作里。我分明意识到,在一项别的工作的边缘上做的工作有其限度;我迟早得从头到尾重新来过,做一项全力投入、逐件核查原始材料的工作。我预计那将是几个星期的事——实际上,我在其中连着泡了七个月:细细检读我父母留下的书信与文字;其中最“灼热”的部分,无疑是1933/34年的通信。而且,七个月之末,我终究还是收住了,发觉这个题目(“认识我的父母”)几乎取之不尽。如今更紧迫的,是\textbf{认识我自己}——借着我刚刚就我父母学到的这一切,并由此,至少间接地……认识我自己那段被遗忘的童年……

我刚花了近两个钟头,翻阅这场关于我父母的禅观——它始于1979年8月3日——的最初的笔记。与我匆匆记起的相反,我当时还没有认识到——除非极其模糊地——有必要“从头到尾”(如我方才所写)重看我此前一个月里读过的父母的书信与其他笔迹。至少,我在笔记里没有作这方面的任何暗示。在一两天的一场总结性反思——盘点这次阅读在我身上激起的多重印象,有些还相当混乱——之后,我全然没有摆出要以逐件核查原始材料的细致工作重做它的架势。\origpage{440}我宁可(仿佛理所当然地)接着以同样快的速度读\textbf{另一些}信(尤其是我父母1937/39年间的一宗浩繁通信),并继续一场由阅读印象喂养的平行反思。正是在8月与次月的进程中,顺藤摸瓜地,我开始懂得:对一封信(或一份生活的其他笔迹见证)下一番\textbf{功夫}是什么意思——那功夫使人得以把握其真正的、有时昭然若揭的意味;而那意味,写信的人却往往乐于无视、瞒过自己如同瞒过别人,神不知鬼不觉!——同时却又把它“摊在字里行间”,有时张扬、有时尖刻。而且,想必很少有这样的情形:那份影射或挑衅(有时凶狠……)竟不曾抵达收信人,不曾被他在某个层面上知觉并“照单收下”——尽管他自己也小心不让这份知觉、这份认识进入他目光的视野,而且他同样扬帆满舵,入的正是这同一出“神不知,鬼不觉!”的游戏。恰恰是那些最晦暗的段落——无一例外地,那些看似近乎低能(或疯狂……)、挑战一切合理解释的段落——在一双好奇的眼睛面前,显出意味最为丰饶:真正的矿,供给无可替代的钥匙,使人得以更深地穿入那一大堆表面的无意义背后那简单而分明的意味。这样的段落,在我父母的通信里、尤其在领舞的我母亲的信里,比比皆是;头一遍阅读——7月间——它们自然全从我头顶掠过。次月里,我开始此处彼处咬住它们。直到9月,几路旁证才使我明白:分明,在我本该从1933/34年的信里学到的东西中,我也许漏掉了某些本质的东西;它们把我带回那些信,催我对其中某些做第一遍“深读”。这一读,立刻把我自童年以来关于我父母其人、关于他们与我和我姐姐的关系曾经如何的图像,翻了个底朝天。

\section{书写的颂歌}

\notehead{102}{}

(9月26日)两天了,我一头扎进“自传式的追忆”——而我本来是要(趁热情已经冷却)续写某条关于某篇葬礼悼词的注的。我不知道这场枝蔓是否多少焐热了我的兴致!至少,该是到达我前天动笔时所瞄准的那个点的时候了——大方向是:“论阅读一份装作不说它要说之话的讯息的艺术”。\dmark{}\origpage{441}这类文本-讯息,远比我当年料想的常见……

不用说,只要人们还甘愿(一如我一生的大部分时光)把你听到或读到的一切当现款收下、按字面领受,并且无论在什么事、什么人身上,都不去寻找当事人明白说出的意图之外的别的意图,那么这门“艺术”的“如何”就无从谈起。相反,当人们有那种难以名状的感觉——某项声明、滔滔一番或一段叙述里,有什么“不对劲”,里面有猫腻,有什么东西“溜过去了”,而它照理是不该被说出的(您还想在这儿猜什么!)——问题就来了。有时,那是一种基本而令人错愕的知觉:一处不自洽、一件荒谬,有时如此巨大、又同时看似无从把握,以致仿佛挑战一切表述,濒临低能或谵妄的边缘。这类局面常常被焦虑压得满满——而每逢这样的局面,荒谬骤然闯入我的生活:一桩不可容忍、不可理解的、暗藏威胁的荒谬,每一次都撼动我那关于世界与我自己的宁静视野的根基——我对它的反应,千篇一律,是一股即刻涌来的焦虑(从不被承认为焦虑,却立刻被搅浑、藏掖在一场狂暴而绝望的愤怒的浪涛之下)。至少,直到我发现“禅观”之时是如此:自那以后,一份无畏而进取的好奇,拆除了这些愤怒与焦虑之浪的引信,并接替了它们……

正是好奇——即求知的欲望——使我在需要的压力下,自发地找到了这门破译一份混沌的见证文本的“艺术”——或者说得更谦逊些:找到了一套适合我那有限手段与我那份滞重的方法。尽管我又好奇又努力,这些含义深重的信,第一遍(甚至第二遍)读下来,本质的东西全从我头顶掠过——“我什么也没看出来”。有时,就某几处格外晦暗而令人错愕的段落,对着或许纷乱的印象写几句评语,笔锋所至,竟得以更深地进入一篇看似密不透风的文本的意味。写着写着,我有时不得不为了引用,抄录一些或长或短的段落:它们或以晦暗见称,或因为粗看之下出于这种或那种缘由给我“重要”的印象。\origpage{442}日子与星期流逝,我发觉:单是\textbf{全文照录}我所端详的文本的某一段,就以惊人的方式改变了我与该段的关系——朝着领会其真正意味的方向敞开。

这完全出乎意料,因为我最初的动机(至少在有意识的层面)纯粹是为了省事。我甚至记得,长年以来,我心里有一种压抑的不耐:把宝贵的光阴花在仅仅充当一名抄写员上!我按捺着焦躁,尽可能快地写,只盼早点抄完……可是,眼睛扫读成行的文字之快,与手逐字逐句转录它们之慢,没有共同的尺度。写得再快,“时间因子”也全然不同。而我疑心,这个“时间因子”并非以纯然机械的、定量的方式起作用——或者不如说,它只是一份更精微、更丰富的实在的一个侧面。至少在我这里,眼睛扫过别人构思并写下的文字的行动,与手一个字母一个字母、一个词一个词重写这同样一些文字的行为,也没有共同的尺度。可以肯定:手,与精神或思维之间,有一种深层的共生;随着书写之手的节奏,而且不带任何蓄意,精神不由自主地重新构成、重新思考同样的词——它们组接成承载着意味的句子,而这些句子又连缀成话语。只要一份求知的欲望鼓舞着这只誊抄字母、词语与句子的手,并同时鼓舞着那个以同一音高、在另一个层面上也“誊抄”着它们的精神——可以肯定,这双重行动,便在我这个甘当抄写者-执笔人的人与这份讯息之间,造成一种远比那件主要消极、不留任何有形踪迹的阅读之眼的行为亲密得多的接触。

这份摸索而来的直觉,与一条由来已久的确认同向——那就是:在我身上,工作者的思维的节奏(无论数学工作还是任何别的工作,包括我称为“禅观”的工作),最常见(倘若不是永远)是书写之手的节奏,而绝不是阅读之眼的节奏\fnote{11}{(*)}{这一情势在我身上起作用,似比在我的大多数数学家同行身上更强;它当年使我难以融入Bourbaki小组的集体工作:我发现自己全然无力跟上那些以它们进行的速度宣读的报告。况且,我从不曾真正喜爱过\textbf{读}数学文本,哪怕再漂亮。我理解数学的自发方式,始终是\textbf{做}它们,或\textbf{重做}它们(必要时,此处彼处借同事们供给的想法与提示,退而求其次,借书……)。}。而我的手(或有时,由我的双手操纵的打字机……)\dmark{}\origpage{443}留下的\textbf{笔迹}——它以那份不慌不忙、却也从不磨蹭的思维前进的节奏写出——是这份思维不可或缺的物质支撑:既是它的“声音”,也是它的“记忆”。况且我疑心,在大多数(倘若不是全体)“智力工作者”身上,情形或多或少也是如此(也许程度稍轻)。

\section{孩子与海——或信与疑}

\notehead{103}{}

(9月27日)不管怎样,事实摆在这里:一如我唯有经由书写才能“进入”一门数学理论,我也几乎唯有经由\textbf{重写},才开始进入一份文本-讯息、一份讯息的“字里行间”。我的第一场“就文本而做”的禅观工作发生转变——一道表面的平庸开始向一道活的深处敞开,荒谬开始取得意味——恰恰\textbf{从那一刻起}:我开始全文重写那份讯息,或(当其篇幅大到无法处理时)重写那些嗅觉使我感到关键的段落。

人们会对我说:既然没有可靠的“客观”判据来担保一份“诠释”的有效性——它被呈献为对(比如说)一份文本的一项(所谓的?)“工作”的结果或终点——人们尽可以让任何文本、任何话语说出自己想说的任何东西,给它编造自己喜欢的随便什么“讯息”。这当然千真万确,例子肯定比比皆是!况且我怀疑(或许除了一门划界分明如史学的学科——即便如此……),这样的判据是否能提炼出来。反正这也无济于事:既不能阻止任何人随意编造异想天开的诠释,也不能使任何人得以探测并发现一份讯息、一个局面、一桩事件的真正意味。规则与判据是一门\textbf{方法}的配料;方法自有其用处与重要性(不过常被高估,以牺牲一些性质全然不同的其他因子与力量为代价)——作为科学或技术知识之发展中的发现与巩固的工具,也作为任何一门手艺(开车或修车,等等)中的工具。相反,在认识自己与他人的层面上,方法的角色变得全然次要:“总管事务”自会跟上来,当本质的东西在场时。而取法于一项方法、或从它出发,乃至死死抱住它,丝毫无助于那份更本质的东西的显现——恰恰相反!

换个说法:那个出发去寻找某样\textbf{预先决定}\dmark{}\origpage{444}的东西的人(他将称之为“真的”,或“真理”),将毫不费力地找到它,甚至证明它,直到自己全然满意;而且他肯定会在路上遇到这一个或那一个同道,甚至一整群,欢欢喜喜与他结盟、分享信念与满足。他像那位捕蝶人:出发时网里就装着一只漂亮的蝴蝶(说不定还是制成标本的),打猎归来,喜滋滋地(而且全然满意地)把它从网里掏出来。

还有另一种人:面对一个未知,如一个赤身的孩子面对大海。当孩子想认识它,他便走进去,便认识了它——无论它温凉,平静或起伏。被某个未知之物吸引、出发去认识它的人,肯定会多少认识它。带网也罢,不带网也罢,他会找到真,或无论如何,找到\textbf{一点}真。他的错误一如他的发现,都是他行程中的台阶;或者不如说,是他与所欲认识者的\textbf{恋情}中的台阶。

我深知自己在说什么,因为我这一生,轮流地,既当过那位捕蝶人,也当过那个赤身的孩子。把二者区分开来,毫无困难。我怀疑“客观判据”在这里能帮上什么忙——事情比那简单得多!只消用自己的眼睛……

同样毫无困难的,是分辨我刚谈到的这段行程中相继的台阶、相继的澄清阶段:从那一“死”的阶段——没有任何预感浮上意识,使人无从疑心“某种东西”,超出困倦的眼睛呈递给我们的那层平面而无定形的表面——直到经由相继的“觉醒”,抵达对这份“某种东西”越来越精微、越来越贴身、越来越完整的把握。无论所论的是在数学事物之发现中的行程,还是在对自己与他人的发现中的行程,其本性并无本质不同。在一份\textbf{认识}中的\textbf{前进}之感——它渐渐加深(哪怕穿过一份错误之积累,耐心地、不倦地不断更正)——在这后一种情形,一如在前一种情形,同样无可辩驳。

这份\textbf{安然的确信}——这是一份内在心性的一面;其另一面,是一份\textbf{向怀疑敞开}:一种好奇的态度,排除对自己的错误的任何恐惧,使人得以不断侦察并更正它们。这一双重立足——这份接纳怀疑与从事发现同样不可或缺的\textbf{信}——的本质条件,是对这项探索将“产出”什么毫无恐惧(无论明显的\dmark{}\origpage{445}还是暗藏的):尤其是,毫无这样的恐惧——怕我们即将发现的实在动摇我们的确信或信念,怕它使我们的希望幻灭。这样的恐惧,如同一道深深的瘫痪,瘫痪我们的创造官能、我们更新的力量。我们可以在艰辛与痛苦中发现并更新自己;却不能在恐惧中发现——恐惧那即将被认识者、即将诞生者。(正如一个人若在某个时刻惧怕一个女人、或惧怕把他带入她体内的那件行为,便不能认识她、使她受孕。)这样的恐惧,在一项科学研究、或任何其主题并不多少深刻牵涉我们自身的研究的语境里,想必相对罕见。相反,当事关对自己或对他人的发现时,它却是最大的绊脚石。

然而,伴随着一次发现——无论大小——的那份感觉,在对自己的发现或对他人的发现的情形里,一如在一场非人称的研究(比如数学)的语境里,同样无可辩驳。我已有机会影射过这份感觉。它是刚刚发生的某事的知觉在情绪层面的反照——某样\textbf{新}东西的显现——而这份“某样东西”,显得与比如说一个数学陈述、或一个此前从未有人想过的观念或证明的显现,同样具体可触、同样无可辩驳(请恕我重复!)。况且我觉得,难以把伴随某次特定发现的那份感觉,与我方才谈到的伴随整场研究的前进之感,区分或分割开来。大的与小的发现,如同把一份前进落实下来的相继\textbf{台阶},如同我们必须相继越过的\textbf{门槛}。\origpage{446}前进不是别的,正是这一连串门槛的跨越、每一台阶向下一台阶的升进。

那份“感觉”——或更好地说,那份映照、还原这一过程的知觉——是一个可靠、不容置疑的“判据”:我不记得它曾把我引入歧途,无论在数学还是在禅观;我不记得曾带着事后的认识确认:这份感觉曾是幻象。它常常使人毫无残余地分辨真与假,或辨出假中之真、与那份自命为真者之中的假。但它尤其是一切真正探索中不可替代的\textbf{向导}——一位随时准备(只要我们肯费心求教)告知我们是否走错了路、抑或正走在好路上的向导。

对这份可靠向导的倾听状态,依我看,不是别的,正是我在反思的另一处\fnote{12}{(*)}{见“严与苛”一节(n\textsuperscript{o}~26),我在其中把“严格”谈作一场探索中“对每一时刻当下\textbf{理解之品质}的精微留意”。}所命名的“严格”。这份严格,依我看,无论所论的是一场数学探索中的苛求,还是自我认识中的苛求,并无本质不同——没有后者,那样的认识便无从谈起。但不用说,这绝不意味着:这份严格在某一智力工作层面的在场,就担保或标志着它在自我与他人的认识中的在场。实际上,相反的情形才是真的;我从无数机缘里——从我自身开始——确认了这一点。在那个领域里,我这里所说的“严格”,是与禅观同时来到我的生活中的。或者不如说,我实在无从区分二者。我生活中的禅观时刻,不是别的,正是我以这样一份对自己极端苛求的心境检视自己的时刻——多半经由我与他人的关系。

\chapter{XII 葬礼仪式}
% 第18章章内目录，按法文原书第433—436页条目顺序生成。
\par\bigskip\noindent{\Large\headfont 目录}\par\medskip
\begingroup\small\setlength{\parskip}{0pt}
\par\noindent\hspace*{0.00em}\hyperref[ch18:local:h1]{18.1.~葬礼悼词}\nobreak\dotfill\nobreak\hyperref[ch18:local:h1]{\pageref*{ch18:local:h1}}\par
\par\noindent\hspace*{1.35em}\hyperref[ch18:local:h2]{18.1.1.~(1) 恭维}\nobreak\dotfill\nobreak\hyperref[ch18:local:h2]{\pageref*{ch18:local:h2}}\par
\par\noindent\hspace*{4.05em}\hyperref[ch18:local:n1]{注~104}\nobreak\dotfill\nobreak\hyperref[ch18:local:n1]{\pageref*{ch18:local:n1}}\par
\par\noindent\hspace*{1.35em}\hyperref[ch18:local:h3]{18.1.2.~(2) 力量与光环}\nobreak\dotfill\nobreak\hyperref[ch18:local:h3]{\pageref*{ch18:local:h3}}\par
\par\noindent\hspace*{4.05em}\hyperref[ch18:local:n2]{注~105}\nobreak\dotfill\nobreak\hyperref[ch18:local:n2]{\pageref*{ch18:local:n2}}\par
\par\noindent\hspace*{0.00em}\hyperref[ch18:local:h4]{18.2.~阴阳之钥}\nobreak\dotfill\nobreak\hyperref[ch18:local:h4]{\pageref*{ch18:local:h4}}\par
\par\noindent\hspace*{1.35em}\hyperref[ch18:local:h5]{18.2.1.~(1) 肌肉与肚肠(阳葬阴(1))}\nobreak\dotfill\nobreak\hyperref[ch18:local:h5]{\pageref*{ch18:local:h5}}\par
\par\noindent\hspace*{4.05em}\hyperref[ch18:local:n3]{注~106}\nobreak\dotfill\nobreak\hyperref[ch18:local:n3]{\pageref*{ch18:local:n3}}\par
\par\noindent\hspace*{4.05em}\hyperref[ch18:local:n4]{注~106\textsuperscript{1}}\nobreak\dotfill\nobreak\hyperref[ch18:local:n4]{\pageref*{ch18:local:n4}}\par
\par\noindent\hspace*{1.35em}\hyperref[ch18:local:h6]{18.2.2.~(2) 一段生命的故事:一个三乐章的循环}\nobreak\dotfill\nobreak\hyperref[ch18:local:h6]{\pageref*{ch18:local:h6}}\par
\par\noindent\hspace*{2.70em}\hyperref[ch18:local:h7]{18.2.2.1.~(a) 天真(阴与阳的联姻)}\nobreak\dotfill\nobreak\hyperref[ch18:local:h7]{\pageref*{ch18:local:h7}}\par
\par\noindent\hspace*{4.05em}\hyperref[ch18:local:n5]{注~107}\nobreak\dotfill\nobreak\hyperref[ch18:local:n5]{\pageref*{ch18:local:n5}}\par
\par\noindent\hspace*{2.70em}\hyperref[ch18:local:h8]{18.2.2.2.~(b) 超级父亲(阳葬阴(2))}\nobreak\dotfill\nobreak\hyperref[ch18:local:h8]{\pageref*{ch18:local:h8}}\par
\par\noindent\hspace*{4.05em}\hyperref[ch18:local:n6]{注~108}\nobreak\dotfill\nobreak\hyperref[ch18:local:n6]{\pageref*{ch18:local:n6}}\par
\par\noindent\hspace*{4.05em}\hyperref[ch18:local:n7]{注~108\textsuperscript{1}}\nobreak\dotfill\nobreak\hyperref[ch18:local:n7]{\pageref*{ch18:local:n7}}\par
\par\noindent\hspace*{4.05em}\hyperref[ch18:local:n8]{注~108\textsuperscript{2}}\nobreak\dotfill\nobreak\hyperref[ch18:local:n8]{\pageref*{ch18:local:n8}}\par
\par\noindent\hspace*{2.70em}\hyperref[ch18:local:h9]{18.2.2.3.~(c) 重逢(阴的苏醒(1))}\nobreak\dotfill\nobreak\hyperref[ch18:local:h9]{\pageref*{ch18:local:h9}}\par
\par\noindent\hspace*{4.05em}\hyperref[ch18:local:n9]{注~109}\nobreak\dotfill\nobreak\hyperref[ch18:local:n9]{\pageref*{ch18:local:n9}}\par
\par\noindent\hspace*{2.70em}\hyperref[ch18:local:h10]{18.2.2.4.~(d) 接纳(阴的苏醒(2))}\nobreak\dotfill\nobreak\hyperref[ch18:local:h10]{\pageref*{ch18:local:h10}}\par
\par\noindent\hspace*{4.05em}\hyperref[ch18:local:n10]{注~110}\nobreak\dotfill\nobreak\hyperref[ch18:local:n10]{\pageref*{ch18:local:n10}}\par
\par\noindent\hspace*{1.35em}\hyperref[ch18:local:h11]{18.2.3.~(3) 伴侣}\nobreak\dotfill\nobreak\hyperref[ch18:local:h11]{\pageref*{ch18:local:h11}}\par
\par\noindent\hspace*{2.70em}\hyperref[ch18:local:h12]{18.2.3.1.~(a) 事物的动力学(阴阳的和谐)}\nobreak\dotfill\nobreak\hyperref[ch18:local:h12]{\pageref*{ch18:local:h12}}\par
\par\noindent\hspace*{4.05em}\hyperref[ch18:local:n11]{注~111}\nobreak\dotfill\nobreak\hyperref[ch18:local:n11]{\pageref*{ch18:local:n11}}\par
\par\noindent\hspace*{2.70em}\hyperref[ch18:local:h13]{18.2.3.2.~(b) 敌对的配偶(阳葬阴(3))}\nobreak\dotfill\nobreak\hyperref[ch18:local:h13]{\pageref*{ch18:local:h13}}\par
\par\noindent\hspace*{4.05em}\hyperref[ch18:local:n12]{注~111$'$}\nobreak\dotfill\nobreak\hyperref[ch18:local:n12]{\pageref*{ch18:local:n12}}\par
\par\noindent\hspace*{2.70em}\hyperref[ch18:local:h14]{18.2.3.3.~(c) 一半与全——或裂缝}\nobreak\dotfill\nobreak\hyperref[ch18:local:h14]{\pageref*{ch18:local:h14}}\par
\par\noindent\hspace*{4.05em}\hyperref[ch18:local:n13]{注~112}\nobreak\dotfill\nobreak\hyperref[ch18:local:n13]{\pageref*{ch18:local:n13}}\par
\par\noindent\hspace*{2.70em}\hyperref[ch18:local:h15]{18.2.3.4.~(d) 原型认识与制约}\nobreak\dotfill\nobreak\hyperref[ch18:local:h15]{\pageref*{ch18:local:h15}}\par
\par\noindent\hspace*{4.05em}\hyperref[ch18:local:n14]{注~112$'$}\nobreak\dotfill\nobreak\hyperref[ch18:local:n14]{\pageref*{ch18:local:n14}}\par
\par\noindent\hspace*{1.35em}\hyperref[ch18:local:h16]{18.2.4.~(4) 我们的母亲,死亡}\nobreak\dotfill\nobreak\hyperref[ch18:local:h16]{\pageref*{ch18:local:h16}}\par
\par\noindent\hspace*{2.70em}\hyperref[ch18:local:h17]{18.2.4.1.~(a) 行为}\nobreak\dotfill\nobreak\hyperref[ch18:local:h17]{\pageref*{ch18:local:h17}}\par
\par\noindent\hspace*{4.05em}\hyperref[ch18:local:n15]{注~113}\nobreak\dotfill\nobreak\hyperref[ch18:local:n15]{\pageref*{ch18:local:n15}}\par
\par\noindent\hspace*{2.70em}\hyperref[ch18:local:h18]{18.2.4.2.~(b) 挚爱者}\nobreak\dotfill\nobreak\hyperref[ch18:local:h18]{\pageref*{ch18:local:h18}}\par
\par\noindent\hspace*{4.05em}\hyperref[ch18:local:n16]{注~114}\nobreak\dotfill\nobreak\hyperref[ch18:local:n16]{\pageref*{ch18:local:n16}}\par
\par\noindent\hspace*{2.70em}\hyperref[ch18:local:h19]{18.2.4.3.~(c) 信使}\nobreak\dotfill\nobreak\hyperref[ch18:local:h19]{\pageref*{ch18:local:h19}}\par
\par\noindent\hspace*{4.05em}\hyperref[ch18:local:n17]{注~114$'$}\nobreak\dotfill\nobreak\hyperref[ch18:local:n17]{\pageref*{ch18:local:n17}}\par
\par\noindent\hspace*{2.70em}\hyperref[ch18:local:h20]{18.2.4.4.~(d) 安格拉——或告别与再会}\nobreak\dotfill\nobreak\hyperref[ch18:local:h20]{\pageref*{ch18:local:h20}}\par
\par\noindent\hspace*{4.05em}\hyperref[ch18:local:n18]{注~115}\nobreak\dotfill\nobreak\hyperref[ch18:local:n18]{\pageref*{ch18:local:n18}}\par
\par\noindent\hspace*{1.35em}\hyperref[ch18:local:h21]{18.2.5.~(5) 拒受与接纳}\nobreak\dotfill\nobreak\hyperref[ch18:local:h21]{\pageref*{ch18:local:h21}}\par
\par\noindent\hspace*{2.70em}\hyperref[ch18:local:h22]{18.2.5.1.~(a) 失乐园}\nobreak\dotfill\nobreak\hyperref[ch18:local:h22]{\pageref*{ch18:local:h22}}\par
\par\noindent\hspace*{4.05em}\hyperref[ch18:local:n19]{注~116}\nobreak\dotfill\nobreak\hyperref[ch18:local:n19]{\pageref*{ch18:local:n19}}\par
\par\noindent\hspace*{2.70em}\hyperref[ch18:local:h23]{18.2.5.2.~(b) 循环}\nobreak\dotfill\nobreak\hyperref[ch18:local:h23]{\pageref*{ch18:local:h23}}\par
\par\noindent\hspace*{4.05em}\hyperref[ch18:local:n20]{注~116$'$}\nobreak\dotfill\nobreak\hyperref[ch18:local:n20]{\pageref*{ch18:local:n20}}\par
\par\noindent\hspace*{2.70em}\hyperref[ch18:local:h24]{18.2.5.3.~(c) 伴侣——或“恶”之谜}\nobreak\dotfill\nobreak\hyperref[ch18:local:h24]{\pageref*{ch18:local:h24}}\par
\par\noindent\hspace*{4.05em}\hyperref[ch18:local:n21]{注~117}\nobreak\dotfill\nobreak\hyperref[ch18:local:n21]{\pageref*{ch18:local:n21}}\par
\par\noindent\hspace*{2.70em}\hyperref[ch18:local:h25]{18.2.5.4.~(d) 阳扮阴——或主人的角色}\nobreak\dotfill\nobreak\hyperref[ch18:local:h25]{\pageref*{ch18:local:h25}}\par
\par\noindent\hspace*{4.05em}\hyperref[ch18:local:n22]{注~118}\nobreak\dotfill\nobreak\hyperref[ch18:local:n22]{\pageref*{ch18:local:n22}}\par
\par\noindent\hspace*{1.35em}\hyperref[ch18:local:h26]{18.2.6.~数学的阴与阳}\nobreak\dotfill\nobreak\hyperref[ch18:local:h26]{\pageref*{ch18:local:h26}}\par
\par\noindent\hspace*{2.70em}\hyperref[ch18:local:h27]{18.2.6.1.~(a) 最“大男子气”的艺术}\nobreak\dotfill\nobreak\hyperref[ch18:local:h27]{\pageref*{ch18:local:h27}}\par
\par\noindent\hspace*{4.05em}\hyperref[ch18:local:n23]{注~119}\nobreak\dotfill\nobreak\hyperref[ch18:local:n23]{\pageref*{ch18:local:n23}}\par
\par\noindent\hspace*{2.70em}\hyperref[ch18:local:h28]{18.2.6.2.~(b) 陌生的美人}\nobreak\dotfill\nobreak\hyperref[ch18:local:h28]{\pageref*{ch18:local:h28}}\par
\par\noindent\hspace*{4.05em}\hyperref[ch18:local:n24]{注~120}\nobreak\dotfill\nobreak\hyperref[ch18:local:n24]{\pageref*{ch18:local:n24}}\par
\par\noindent\hspace*{2.70em}\hyperref[ch18:local:h29]{18.2.6.3.~(c) 欲与严格}\nobreak\dotfill\nobreak\hyperref[ch18:local:h29]{\pageref*{ch18:local:h29}}\par
\par\noindent\hspace*{4.05em}\hyperref[ch18:local:n25]{注~121}\nobreak\dotfill\nobreak\hyperref[ch18:local:n25]{\pageref*{ch18:local:n25}}\par
\par\noindent\hspace*{2.70em}\hyperref[ch18:local:h30]{18.2.6.4.~(d) 上涨的海……}\nobreak\dotfill\nobreak\hyperref[ch18:local:h30]{\pageref*{ch18:local:h30}}\par
\par\noindent\hspace*{4.05em}\hyperref[ch18:local:n26]{注~122}\nobreak\dotfill\nobreak\hyperref[ch18:local:n26]{\pageref*{ch18:local:n26}}\par
\par\noindent\hspace*{2.70em}\hyperref[ch18:local:h31]{18.2.6.5.~(e) 九个月与上涨的海}\nobreak\dotfill\nobreak\hyperref[ch18:local:h31]{\pageref*{ch18:local:h31}}\par
\par\noindent\hspace*{4.05em}\hyperref[ch18:local:n27]{注~123}\nobreak\dotfill\nobreak\hyperref[ch18:local:n27]{\pageref*{ch18:local:n27}}\par
\par\noindent\hspace*{2.70em}\hyperref[ch18:local:h32]{18.2.6.6.~(f) 阴的葬礼(阳葬阴(4))}\nobreak\dotfill\nobreak\hyperref[ch18:local:h32]{\pageref*{ch18:local:h32}}\par
\par\noindent\hspace*{4.05em}\hyperref[ch18:local:n28]{注~124}\nobreak\dotfill\nobreak\hyperref[ch18:local:n28]{\pageref*{ch18:local:n28}}\par
\par\noindent\hspace*{2.70em}\hyperref[ch18:local:h33]{18.2.6.7.~(g) 超级妈妈,还是超级爸爸}\nobreak\dotfill\nobreak\hyperref[ch18:local:h33]{\pageref*{ch18:local:h33}}\par
\par\noindent\hspace*{4.05em}\hyperref[ch18:local:n29]{注~125}\nobreak\dotfill\nobreak\hyperref[ch18:local:n29]{\pageref*{ch18:local:n29}}\par
\par\noindent\hspace*{1.35em}\hyperref[ch18:local:h34]{18.2.7.~阴阳的颠倒}\nobreak\dotfill\nobreak\hyperref[ch18:local:h34]{\pageref*{ch18:local:h34}}\par
\par\noindent\hspace*{2.70em}\hyperref[ch18:local:h35]{18.2.7.1.~(a) 颠倒(1)——或激烈泼辣的妻子}\nobreak\dotfill\nobreak\hyperref[ch18:local:h35]{\pageref*{ch18:local:h35}}\par
\par\noindent\hspace*{4.05em}\hyperref[ch18:local:n30]{注~126}\nobreak\dotfill\nobreak\hyperref[ch18:local:n30]{\pageref*{ch18:local:n30}}\par
\par\noindent\hspace*{2.70em}\hyperref[ch18:local:h36]{18.2.7.2.~回顾(1)——或一幅画的三联}\nobreak\dotfill\nobreak\hyperref[ch18:local:h36]{\pageref*{ch18:local:h36}}\par
\par\noindent\hspace*{4.05em}\hyperref[ch18:local:n31]{注~127}\nobreak\dotfill\nobreak\hyperref[ch18:local:n31]{\pageref*{ch18:local:n31}}\par
\par\noindent\hspace*{2.70em}\hyperref[ch18:local:h37]{18.2.7.3.~回顾(2)——或死结}\nobreak\dotfill\nobreak\hyperref[ch18:local:h37]{\pageref*{ch18:local:h37}}\par
\par\noindent\hspace*{4.05em}\hyperref[ch18:local:n32]{注~127$'$}\nobreak\dotfill\nobreak\hyperref[ch18:local:n32]{\pageref*{ch18:local:n32}}\par
\par\noindent\hspace*{2.70em}\hyperref[ch18:local:h38]{18.2.7.4.~(d) 父母——或冲突的心脏}\nobreak\dotfill\nobreak\hyperref[ch18:local:h38]{\pageref*{ch18:local:h38}}\par
\par\noindent\hspace*{4.05em}\hyperref[ch18:local:n33]{注~128}\nobreak\dotfill\nobreak\hyperref[ch18:local:n33]{\pageref*{ch18:local:n33}}\par
\par\noindent\hspace*{4.05em}\hyperref[ch18:local:n34]{注~128$^1$}\nobreak\dotfill\nobreak\hyperref[ch18:local:n34]{\pageref*{ch18:local:n34}}\par
\par\noindent\hspace*{2.70em}\hyperref[ch18:local:h39]{18.2.7.5.~(e) 敌对的父亲(3)——或阳葬阳}\nobreak\dotfill\nobreak\hyperref[ch18:local:h39]{\pageref*{ch18:local:h39}}\par
\par\noindent\hspace*{4.05em}\hyperref[ch18:local:n35]{注~129}\nobreak\dotfill\nobreak\hyperref[ch18:local:n35]{\pageref*{ch18:local:n35}}\par
\par\noindent\hspace*{2.70em}\hyperref[ch18:local:h40]{18.2.7.6.~(f) 箭与浪}\nobreak\dotfill\nobreak\hyperref[ch18:local:h40]{\pageref*{ch18:local:h40}}\par
\par\noindent\hspace*{4.05em}\hyperref[ch18:local:n36]{注~130}\nobreak\dotfill\nobreak\hyperref[ch18:local:n36]{\pageref*{ch18:local:n36}}\par
\par\noindent\hspace*{2.70em}\hyperref[ch18:local:h41]{18.2.7.7.~(g) 冲突之谜}\nobreak\dotfill\nobreak\hyperref[ch18:local:h41]{\pageref*{ch18:local:h41}}\par
\par\noindent\hspace*{4.05em}\hyperref[ch18:local:n37]{注~131}\nobreak\dotfill\nobreak\hyperref[ch18:local:n37]{\pageref*{ch18:local:n37}}\par
\par\noindent\hspace*{2.70em}\hyperref[ch18:local:h42]{18.2.7.8.~(h) 颠倒(2)——或暧昧的反叛}\nobreak\dotfill\nobreak\hyperref[ch18:local:h42]{\pageref*{ch18:local:h42}}\par
\par\noindent\hspace*{4.05em}\hyperref[ch18:local:n38]{注~132}\nobreak\dotfill\nobreak\hyperref[ch18:local:n38]{\pageref*{ch18:local:n38}}\par
\par\noindent\hspace*{1.35em}\hyperref[ch18:local:h43]{18.2.8.~主人们与仆人}\nobreak\dotfill\nobreak\hyperref[ch18:local:h43]{\pageref*{ch18:local:h43}}\par
\par\noindent\hspace*{2.70em}\hyperref[ch18:local:h44]{18.2.8.1.~(a) 天鹅绒的脚掌——或微笑}\nobreak\dotfill\nobreak\hyperref[ch18:local:h44]{\pageref*{ch18:local:h44}}\par
\par\noindent\hspace*{4.05em}\hyperref[ch18:local:n39]{注~133}\nobreak\dotfill\nobreak\hyperref[ch18:local:n39]{\pageref*{ch18:local:n39}}\par
\par\noindent\hspace*{2.70em}\hyperref[ch18:local:h45]{18.2.8.2.~(b) 兄弟与配偶——或双重签名}\nobreak\dotfill\nobreak\hyperref[ch18:local:h45]{\pageref*{ch18:local:h45}}\par
\par\noindent\hspace*{4.05em}\hyperref[ch18:local:n40]{注~134}\nobreak\dotfill\nobreak\hyperref[ch18:local:n40]{\pageref*{ch18:local:n40}}\par
\par\noindent\hspace*{4.05em}\hyperref[ch18:local:n41]{注~134$^1$}\nobreak\dotfill\nobreak\hyperref[ch18:local:n41]{\pageref*{ch18:local:n41}}\par
\par\noindent\hspace*{2.70em}\hyperref[ch18:local:h46]{18.2.8.3.~(c) 阴,作为仆人——与新主人}\nobreak\dotfill\nobreak\hyperref[ch18:local:h46]{\pageref*{ch18:local:h46}}\par
\par\noindent\hspace*{4.05em}\hyperref[ch18:local:n42]{注~135}\nobreak\dotfill\nobreak\hyperref[ch18:local:n42]{\pageref*{ch18:local:n42}}\par
\par\noindent\hspace*{2.70em}\hyperref[ch18:local:h47]{18.2.8.4.~(d) 阴,作为仆人(2)——与慷慨}\nobreak\dotfill\nobreak\hyperref[ch18:local:h47]{\pageref*{ch18:local:h47}}\par
\par\noindent\hspace*{4.05em}\hyperref[ch18:local:n43]{注~136}\nobreak\dotfill\nobreak\hyperref[ch18:local:n43]{\pageref*{ch18:local:n43}}\par
\par\noindent\hspace*{4.05em}\hyperref[ch18:local:n44]{注~136$^1$}\nobreak\dotfill\nobreak\hyperref[ch18:local:n44]{\pageref*{ch18:local:n44}}\par
\par\noindent\hspace*{1.35em}\hyperref[ch18:local:h48]{18.2.9.~天鹅绒中的利爪}\nobreak\dotfill\nobreak\hyperref[ch18:local:h48]{\pageref*{ch18:local:h48}}\par
\par\noindent\hspace*{2.70em}\hyperref[ch18:local:h49]{18.2.9.1.~(a) 天鹅绒的脚掌——或微笑}\nobreak\dotfill\nobreak\hyperref[ch18:local:h49]{\pageref*{ch18:local:h49}}\par
\par\noindent\hspace*{4.05em}\hyperref[ch18:local:n45]{注~137}\nobreak\dotfill\nobreak\hyperref[ch18:local:n45]{\pageref*{ch18:local:n45}}\par
\par\noindent\hspace*{2.70em}\hyperref[ch18:local:h50]{18.2.9.2.~(b) 颠倒(4)——或夫妻的马戏}\nobreak\dotfill\nobreak\hyperref[ch18:local:h50]{\pageref*{ch18:local:h50}}\par
\par\noindent\hspace*{4.05em}\hyperref[ch18:local:n46]{注~138}\nobreak\dotfill\nobreak\hyperref[ch18:local:n46]{\pageref*{ch18:local:n46}}\par
\par\noindent\hspace*{2.70em}\hyperref[ch18:local:h51]{18.2.9.3.~(c) 天真的暴力——或传递}\nobreak\dotfill\nobreak\hyperref[ch18:local:h51]{\pageref*{ch18:local:h51}}\par
\par\noindent\hspace*{4.05em}\hyperref[ch18:local:n47]{注~139}\nobreak\dotfill\nobreak\hyperref[ch18:local:n47]{\pageref*{ch18:local:n47}}\par
\par\noindent\hspace*{2.70em}\hyperref[ch18:local:h52]{18.2.9.4.~(d) 奴隶与傀儡——或闸门}\nobreak\dotfill\nobreak\hyperref[ch18:local:h52]{\pageref*{ch18:local:h52}}\par
\par\noindent\hspace*{4.05em}\hyperref[ch18:local:n48]{注~140}\nobreak\dotfill\nobreak\hyperref[ch18:local:n48]{\pageref*{ch18:local:n48}}\par
\par\noindent\hspace*{1.35em}\hyperref[ch18:local:h53]{18.2.10.~暴力——或游戏与刺棒}\nobreak\dotfill\nobreak\hyperref[ch18:local:h53]{\pageref*{ch18:local:h53}}\par
\par\noindent\hspace*{2.70em}\hyperref[ch18:local:h54]{18.2.10.1.~(a) 正义者的暴力}\nobreak\dotfill\nobreak\hyperref[ch18:local:h54]{\pageref*{ch18:local:h54}}\par
\par\noindent\hspace*{4.05em}\hyperref[ch18:local:n49]{注~141}\nobreak\dotfill\nobreak\hyperref[ch18:local:n49]{\pageref*{ch18:local:n49}}\par
\par\noindent\hspace*{2.70em}\hyperref[ch18:local:h55]{18.2.10.2.~(b) 机械与自由}\nobreak\dotfill\nobreak\hyperref[ch18:local:h55]{\pageref*{ch18:local:h55}}\par
\par\noindent\hspace*{4.05em}\hyperref[ch18:local:n50]{注~142}\nobreak\dotfill\nobreak\hyperref[ch18:local:n50]{\pageref*{ch18:local:n50}}\par
\par\noindent\hspace*{2.70em}\hyperref[ch18:local:h56]{18.2.10.3.~(c) 贪婪——或蚀本买卖}\nobreak\dotfill\nobreak\hyperref[ch18:local:h56]{\pageref*{ch18:local:h56}}\par
\par\noindent\hspace*{4.05em}\hyperref[ch18:local:n51]{注~143}\nobreak\dotfill\nobreak\hyperref[ch18:local:n51]{\pageref*{ch18:local:n51}}\par
\par\noindent\hspace*{2.70em}\hyperref[ch18:local:h57]{18.2.10.4.~(d) 两种认识——或对认识的恐惧}\nobreak\dotfill\nobreak\hyperref[ch18:local:h57]{\pageref*{ch18:local:h57}}\par
\par\noindent\hspace*{4.05em}\hyperref[ch18:local:n52]{注~144}\nobreak\dotfill\nobreak\hyperref[ch18:local:n52]{\pageref*{ch18:local:n52}}\par
\par\noindent\hspace*{2.70em}\hyperref[ch18:local:h58]{18.2.10.5.~(e) 秘密的神经}\nobreak\dotfill\nobreak\hyperref[ch18:local:h58]{\pageref*{ch18:local:h58}}\par
\par\noindent\hspace*{4.05em}\hyperref[ch18:local:n53]{注~145}\nobreak\dotfill\nobreak\hyperref[ch18:local:n53]{\pageref*{ch18:local:n53}}\par
\par\noindent\hspace*{2.70em}\hyperref[ch18:local:h59]{18.2.10.6.~(f) 激情与饥渴——或逐步加码}\nobreak\dotfill\nobreak\hyperref[ch18:local:h59]{\pageref*{ch18:local:h59}}\par
\par\noindent\hspace*{4.05em}\hyperref[ch18:local:n54]{注~146}\nobreak\dotfill\nobreak\hyperref[ch18:local:n54]{\pageref*{ch18:local:n54}}\par
\par\noindent\hspace*{2.70em}\hyperref[ch18:local:h60]{18.2.10.7.~(g) 心软的爸爸}\nobreak\dotfill\nobreak\hyperref[ch18:local:h60]{\pageref*{ch18:local:h60}}\par
\par\noindent\hspace*{4.05em}\hyperref[ch18:local:n55]{注~147}\nobreak\dotfill\nobreak\hyperref[ch18:local:n55]{\pageref*{ch18:local:n55}}\par
\par\noindent\hspace*{2.70em}\hyperref[ch18:local:h61]{18.2.10.8.~神经中的神经——或侏儒与巨人}\nobreak\dotfill\nobreak\hyperref[ch18:local:h61]{\pageref*{ch18:local:h61}}\par
\par\noindent\hspace*{4.05em}\hyperref[ch18:local:n56]{注~148}\nobreak\dotfill\nobreak\hyperref[ch18:local:n56]{\pageref*{ch18:local:n56}}\par
\par\noindent\hspace*{1.35em}\hyperref[ch18:local:h62]{18.2.11.~另一个自己}\nobreak\dotfill\nobreak\hyperref[ch18:local:h62]{\pageref*{ch18:local:h62}}\par
\par\noindent\hspace*{2.70em}\hyperref[ch18:local:h63]{18.2.11.1.~(a) 缓期中的宿怨——或事物的回转(2)}\nobreak\dotfill\nobreak\hyperref[ch18:local:h63]{\pageref*{ch18:local:h63}}\par
\par\noindent\hspace*{4.05em}\hyperref[ch18:local:n57]{注~149}\nobreak\dotfill\nobreak\hyperref[ch18:local:n57]{\pageref*{ch18:local:n57}}\par
\par\noindent\hspace*{2.70em}\hyperref[ch18:local:h64]{18.2.11.2.~(b) 天真与冲突——或绊脚石}\nobreak\dotfill\nobreak\hyperref[ch18:local:h64]{\pageref*{ch18:local:h64}}\par
\par\noindent\hspace*{4.05em}\hyperref[ch18:local:n58]{注~150}\nobreak\dotfill\nobreak\hyperref[ch18:local:n58]{\pageref*{ch18:local:n58}}\par
\par\noindent\hspace*{2.70em}\hyperref[ch18:local:h65]{18.2.11.3.~(c) 天赐的情势——或封神}\nobreak\dotfill\nobreak\hyperref[ch18:local:h65]{\pageref*{ch18:local:h65}}\par
\par\noindent\hspace*{4.05em}\hyperref[ch18:local:n59]{注~151}\nobreak\dotfill\nobreak\hyperref[ch18:local:n59]{\pageref*{ch18:local:n59}}\par
\par\noindent\hspace*{2.70em}\hyperref[ch18:local:h66]{18.2.11.4.~(d) 否认(1)——或提醒}\nobreak\dotfill\nobreak\hyperref[ch18:local:h66]{\pageref*{ch18:local:h66}}\par
\par\noindent\hspace*{4.05em}\hyperref[ch18:local:n60]{注~152}\nobreak\dotfill\nobreak\hyperref[ch18:local:n60]{\pageref*{ch18:local:n60}}\par
\par\noindent\hspace*{4.05em}\hyperref[ch18:local:n61]{注~153}\nobreak\dotfill\nobreak\hyperref[ch18:local:n61]{\pageref*{ch18:local:n61}}\par
\par\noindent\hspace*{2.70em}\hyperref[ch18:local:h67]{18.2.11.5.~(f) 布景——或“第二天性”}\nobreak\dotfill\nobreak\hyperref[ch18:local:h67]{\pageref*{ch18:local:h67}}\par
\par\noindent\hspace*{4.05em}\hyperref[ch18:local:n62]{注~154}\nobreak\dotfill\nobreak\hyperref[ch18:local:n62]{\pageref*{ch18:local:n62}}\par
\par\noindent\hspace*{2.70em}\hyperref[ch18:local:h68]{18.2.11.6.~(g) 另一个自己——或认同与冲突}\nobreak\dotfill\nobreak\hyperref[ch18:local:h68]{\pageref*{ch18:local:h68}}\par
\par\noindent\hspace*{4.05em}\hyperref[ch18:local:n63]{注~155}\nobreak\dotfill\nobreak\hyperref[ch18:local:n63]{\pageref*{ch18:local:n63}}\par
\par\noindent\hspace*{2.70em}\hyperref[ch18:local:h69]{18.2.11.7.~(h) 敌对的兄弟——或传递(2)}\nobreak\dotfill\nobreak\hyperref[ch18:local:h69]{\pageref*{ch18:local:h69}}\par
\par\noindent\hspace*{4.05em}\hyperref[ch18:local:n64]{注~156}\nobreak\dotfill\nobreak\hyperref[ch18:local:n64]{\pageref*{ch18:local:n64}}\par
\par\noindent\hspace*{4.05em}\hyperref[ch18:local:n65]{注~156$^1$}\nobreak\dotfill\nobreak\hyperref[ch18:local:n65]{\pageref*{ch18:local:n65}}\par
\par\noindent\hspace*{1.35em}\hyperref[ch18:local:h70]{18.2.12.~冲突与发现——或恶之谜}\nobreak\dotfill\nobreak\hyperref[ch18:local:h70]{\pageref*{ch18:local:h70}}\par
\par\noindent\hspace*{2.70em}\hyperref[ch18:local:h71]{18.2.12.1.~(a) 无恨亦无容赦}\nobreak\dotfill\nobreak\hyperref[ch18:local:h71]{\pageref*{ch18:local:h71}}\par
\par\noindent\hspace*{4.05em}\hyperref[ch18:local:n66]{注~157}\nobreak\dotfill\nobreak\hyperref[ch18:local:n66]{\pageref*{ch18:local:n66}}\par
\par\noindent\hspace*{2.70em}\hyperref[ch18:local:h72]{18.2.12.2.~(b) 领会与更新}\nobreak\dotfill\nobreak\hyperref[ch18:local:h72]{\pageref*{ch18:local:h72}}\par
\par\noindent\hspace*{4.05em}\hyperref[ch18:local:n67]{注~158}\nobreak\dotfill\nobreak\hyperref[ch18:local:n67]{\pageref*{ch18:local:n67}}\par
\par\noindent\hspace*{2.70em}\hyperref[ch18:local:h73]{18.2.12.3.~(c) 无缘之暴力的缘由}\nobreak\dotfill\nobreak\hyperref[ch18:local:h73]{\pageref*{ch18:local:h73}}\par
\par\noindent\hspace*{4.05em}\hyperref[ch18:local:n68]{注~159}\nobreak\dotfill\nobreak\hyperref[ch18:local:n68]{\pageref*{ch18:local:n68}}\par
\par\noindent\hspace*{2.70em}\hyperref[ch18:local:h74]{18.2.12.4.~(d) 藤井日达古鲁吉——或太阳与其行星}\nobreak\dotfill\nobreak\hyperref[ch18:local:h74]{\pageref*{ch18:local:h74}}\par
\par\noindent\hspace*{4.05em}\hyperref[ch18:local:n69]{注~160}\nobreak\dotfill\nobreak\hyperref[ch18:local:n69]{\pageref*{ch18:local:n69}}\par
\par\noindent\hspace*{2.70em}\hyperref[ch18:local:h75]{18.2.12.5.~(e) 祈祷与冲突}\nobreak\dotfill\nobreak\hyperref[ch18:local:h75]{\pageref*{ch18:local:h75}}\par
\par\noindent\hspace*{4.05em}\hyperref[ch18:local:n70]{注~161}\nobreak\dotfill\nobreak\hyperref[ch18:local:n70]{\pageref*{ch18:local:n70}}\par
\par\noindent\hspace*{2.70em}\hyperref[ch18:local:h76]{18.2.12.6.~(f) 信念与认识}\nobreak\dotfill\nobreak\hyperref[ch18:local:h76]{\pageref*{ch18:local:h76}}\par
\par\noindent\hspace*{4.05em}\hyperref[ch18:local:n71]{注~162}\nobreak\dotfill\nobreak\hyperref[ch18:local:n71]{\pageref*{ch18:local:n71}}\par
\par\noindent\hspace*{2.70em}\hyperref[ch18:local:h77]{18.2.12.7.~(g) 最烫的铁——或转折}\nobreak\dotfill\nobreak\hyperref[ch18:local:h77]{\pageref*{ch18:local:h77}}\par
\par\noindent\hspace*{4.05em}\hyperref[ch18:local:n72]{注~162$'$}\nobreak\dotfill\nobreak\hyperref[ch18:local:n72]{\pageref*{ch18:local:n72}}\par
\par\noindent\hspace*{2.70em}\hyperref[ch18:local:h78]{18.2.12.8.~(h) 无穷的链——或传递(3)}\nobreak\dotfill\nobreak\hyperref[ch18:local:h78]{\pageref*{ch18:local:h78}}\par
\par\noindent\hspace*{4.05em}\hyperref[ch18:local:n73]{注~162$''$}\nobreak\dotfill\nobreak\hyperref[ch18:local:n73]{\pageref*{ch18:local:n73}}\par
\endgroup\medskip



\section{葬礼悼词}
\label{ch18:local:h1}

\subsection{(1) 恭维}
\label{ch18:local:h2}

\notehead{104}{}
\phantomsection\label{ch18:local:n1}

\dmark{}(5月12日)\fnote{1}{(*)}{(5月18日)下面这条注,“出自一条(注n\textsuperscript{o}~47的)脚注——它膨胀到了犯规的规模”。我把它插在这里,心想这一次这个次序比编年的次序更自然。自这条注写成的时刻起,我就感到需要再把它发展些许——这将由紧随本注之后的一条注来完成;写下这些文字时,它还没有写成。这两条注的全套,从现在起就顶上了那个顺理成章的名字:“葬礼悼词”!}在这本\origpage{447}\fnote{2}{(**)}{(5月18日)说的是IHES(高等科学研究所)为庆祝建所二十五周年于1983年出版的小册子。注“有益的挣脱”(n\textsuperscript{o}~42)的脚注里已提到它;“拒受遗产——或矛盾的代价”一注(n\textsuperscript{o}~47)的开头再次提到;本注(葬礼悼词(1))即关涉后者(见前一条脚注)。}小册子里就我的著作所写的那篇小传里,“上同调”或“同调”一词竟不曾出现!“概形”一词也没有。文中当然谈到(情势所需——我当时毕竟充当着“IHES的第一枚菲尔兹奖章”)我著作的“巨人般面貌”、出版的卷数、提炼出的本质问题、“最大的自然一般性”(奇怪的法语!),极为讲究的术语,影射“Grothendieck群”(我敢打赌,又是一项“最大的自然一般性”!),甚至topos及其在逻辑中的用处(但尤其不在别处!)……却没有任何一处影射一项\textbf{结果},或一项我本可能发展、且或许派得上用场的\textbf{理论}——看来只好相信:那二十卷巨人般的著作,严格说是空的,或至多是一些(从未解决的)问题与观念的汇编,当然都带着最大的自然一般性:Grothendieck群,判给了(既然我的名字已经粘在后面),呈献为代数K理论的“始祖”(!)(而它当然与拓扑K理论毫无干系——对后者则讳莫如深\fnote{3}{(***)}{我关于Riemann--Roch定理的工作,是代数\textbf{K}理论的第一记强力启动,而全然不是什么“始祖”。拓扑\textbf{K}理论,恰恰诞生于我证明Riemann--Roch--Grothendieck定理的那一年(1957),继我在赫策布鲁赫研讨班上的报告之后。这位被讳莫如深的“后代”的“始祖”,当时还不满一岁!代数K理论(加上巴斯引入的函子$K_i$[原书此处仅印$K$],以补我引入的函子$K^0$)在其后的年月里发展起来,受“始祖”与其第一个“后代”的双重影响。况且,早在六十年代下半叶,我就有了一条通向描述高阶$K^i$(对一个原文称为“monomiale”的范畴,比如加性范畴)的进路,循着Sinh~夫人的论文的路线。它停留在启发的层面,因为它奠基于包裹Picard的$\infty$-\textbf{范畴}的直觉;而彼时(乃至此后)还没有人费心发展(非严格的)$\infty$-范畴的观念——即我如今以topos上的$\infty$-叠(点的topos)之名所称的观念。随着我即将在《追寻栈》里发展的、一套叠的上同调-同伦形式主义的奠基草图(正处在我1955至1965年间所发展观念的直线上),这条通往高阶K-不变量理论的“几何”进路,终将唾手可得。};至于Riemann--Roch定理,想必是“始祖”的后代们去打理的——那些做真正定理、做正经事情的人!

在一个时尚以轻蔑为荣、对一般性(被那句有点可笑的“最大的自然一般性”的措辞不动声色地讥讽着……)不屑一顾的时代,\dmark{}那支费心为我撰写葬礼悼词的匿名手笔,把如今正被交付轻蔑的东西\dmark{}\origpage{448},毫不吝惜地赏给了我\fnote{4}{(*)}{(5月18日)而且这还是我删削过的!我的葬礼悼词的全文引证,见“葬礼悼词(2)”一注。}。我也赏识了(也许我是第一个……)同一支匿名手笔在悼词的这一段里的全部幽默:

“他在IHES创建了一个代数几何学派,围绕他主持的研讨班聚集起来,\textbf{并以他传播自己观念时的那份慷慨为滋养}”(着重号是我加的)。遗憾的是,一如我的“巨人般著作”,这个被我滋养得如此之好的“代数几何学派”,严格说也是空的——一个名字也不曾提起;而且没有人来抱怨自己被遗忘了——至少没有向我抱怨。

然而我似乎记得,年轻的Deligne曾在1965年(他当时该有十九岁)至1969年间,忠实地出没于那场(想必空空如也的)研讨班,并在这一研讨班里、在我们的两两对谈里,既学得了概形的技术,也学得了上同调技术与平展上同调——也就是说,恰恰是那些在其著作(至少在我见过的那些部分里)的每一页上都被使用的工具。在同一本小册子里献给Deligne的那篇小传里,同样没有任何影射,能使读者疑心\dmark{}\origpage{449}他可能从我这里学到了什么。然而,引人注目的是,在这篇(这一回可毫不“葬礼”的)Deligne颂——“IHES的第三枚菲尔兹奖章”——里,我的名字被提到了三次!甚至有一处迂回的措辞,带着照例必须环绕我寒微身量的那份含混,影射我曾“在任意域上的几何中构造了上同调理论”——而且肯定还是“带着最大的自然一般性”:一股Grothendieck学问的浓重气味,劈面而来\fnote{5}{(*)}{(5月18日)在葬礼悼词里,谈到我曾对术语怀有的那份“高度留意”。在“最大的自然一般性”或“在任意域上的几何中的上同调理论”这类离奇的措辞的使用里,我分明察觉到把这份留意变成笑柄的意图。我对给予事物之名字的那份极度的用心,自然出于我对这些事物的敬重——事物的名字本应表达其本质,或至少其某个本质的侧面。从传到我耳中的回声看,这份敬重的态度——它如今俨然合乎时宜的轻蔑所针对的对象——已不止一次令我震惊;那轻蔑,有时借给一些重要观念取些阿布拉卡达布拉式的名字来表达。对此另可参见“反常性”一注(n\textsuperscript{o}~76)。}。把上下文完整引出来是值得的;这是此类体裁的一件小小的杰作:

\begin{quote}
“从那里出发[经典Hodge理论],并从Grothendieck所提示的$\ell$-adic类比出发[人们纳闷:Grothendieck哪里找来的时间,一边写着他那二十卷最大的自然一般性,一边学得起这样正经的东西],他[Deligne]提炼出混合Hodge结构的观念,并使任何复代数簇的上同调带有这一结构。于是在$\ell$-adic上同调里[?],对有限域上的簇,他证明了困难尽人皆知的Weil猜想。这一结果显得愈发令人惊讶[!!],因为Grothendieck在任意域上的几何中构造了上同调理论之后[人们实在纳闷他还到那里去找什么],曾把余下的猜想[???]归约为一系列猜想——它们今天一如既往地无从下手。”
\end{quote}

说白了:这些Grothendieck学问(其名字足以吓跑最顽强的博物学通才)非但对证明这个困难如此“尽人皆知”的惊人结果毫无贡献,而且至多恰到好处,又用\textbf{猜想}(他从来不做别的!)把我们缠住一层,而且(谁又能料到呢)无从下手,一如他突发奇想提出它们的那时,也一如今天。

\dmark{}\origpage{450}然而我记得自己曾着手对付过它们——这些无从下手的猜想;但那想必是因为我消息不灵。那是在我出发——抱歉,我是想说“亡故”——的时分;而我那些比我消息灵通的后裔们,则小心从不把鼻子伸进这类玩意,因为Deligne言之凿凿:那是无从下手的!

我分明认得这套做派:人们尽了本分,大量引用Grothendieck(无论他还是谁,都不能声称人们在这庄严的日子埋葬他);人们甚至来了一个“暂停”式的影射,提到某些“$\ell$-adic类比”,它们曾在混合Hodge理论的启动中起过作用。这想必是十三年前那句大名鼎鼎的碑文式半行之后的第二遭\fnote{6}{(*)}{这句“碑文式的半行”,见于Deligne1970年尼斯国际数学家大会的报告“Hodge理论I”。评语见注n\textsuperscript{o}~78\textsuperscript{2}。};两处影射,与某篇1968年文章里那些“重量方面的考虑”出奇地相似\fnote{7}{(**)}{对此可参见“罐装的权与十二年的秘密”一注(n\textsuperscript{o}~49)的开头,及“驱逐”一注(n\textsuperscript{o}~63)中更周详的检视。}:人们处于“暂停”的豁免状态,同时牵着读者的鼻子走!在这里,托庄严场合之福,那记“暂停”式的引证干得比消灭罪证更妙——这段文字想就那位大名鼎鼎的Grothendieck给出的印象,恰恰是这几年我感到的那股时尚之“风”所承载的;那股风,我\textbf{今天}就有机会感受到了\fnote{8}{(***)}{见同日的“屠杀”一注,n\textsuperscript{o}~87.},不再是在当众致辞、宾客盈场的葬礼悼词的音调里,而是在屠杀的音调里……

我继续引用——它值得:

\begin{quote}
“这条定理(前Weil猜想)有助于使$\ell$-adic上同调成为一件强大的工具,无须提名这件如此强大的工具的发明者——那位光彩夺目而谦虚的人……;它适用于一些表面上远离代数几何的问题,比如Ramanujam猜想。

最近,他研究了阿贝尔簇上的Hodge循环,朝着Grothendieck曾梦见的“Motif”理论迈出了第一步。他还证明了“相交上同调”的代数机制——Mac~Pherson与Goresky的拓扑理论。这使得把它移植到$\ell$-adic理论成为可能;在那里,它显出惊人的用处。”
\end{quote}

就这样,一支匿名的手笔(我猜是同一支),在那部“值得纪念的文集”问世一年之后\fnote{9}{(****)}{即1982年出版的Lecture~Notes~n\textsuperscript{o}~900那一卷;“一个梦的回忆——或Motif的诞生”与尤其“下葬——或新父亲”两注(n\textsuperscript{o}~51、52)谈到了它。正是在那一卷里,Motif被“开棺出土”(在对它们十二年的死一般的沉默之后),顶着一份(心照不宣的)替补父权。},终于弥补了该卷里的一处小“疏漏”。想必毕竟有人提了一个问题;Deligne在这里以他自己的方式尽了他弥补疏漏的手续(毕竟,当终于轮到正经数学的时候,还肯引这位做梦人\dmark{}Grothendieck一把,也算客气!)\origpage{451}而且照旧欺瞒读者:所谓“第一步”,早在1968年就迈出了——随着Deligne启动Hodge--Deligne理论;而那理论,扎的根恰是Motif瑜伽:在先前的四年里,他的确是经由与我的接触“吸收”了它。他那部著作所自出、他始终无法摆脱却偏要否认的这部瑜伽,实际上在第一段引文的迂回措辞里,被以“$\ell$-adic类比”之名打发了。一位既不极灵通也不极留心的读者,断然不会疑心这些“$\ell$-adic类比”(它们对Hodge--Deligne理论扮演过一个起点的角色——但至多如此……\fnote{10}{(*)}{这套Hodge--Deligne理论至今仍停留在童年状态:因为“Hodge--Deligne复形”的观念——架在$\mathbb{C}$上任意有限型概形之上——与这些“系数”的六种运算形式主义,始终不曾得到发展。对这样一套理论的需要,对Deligne一如对我,早在其混合Hodge结构的最初工作之前就已分明;它以分明的方式出自Motif瑜伽。但自我退出数学舞台起,Deligne身上便发展出一道对我引入同调代数的关键观念(导出范畴、六种运算——更不必说topos)的“阻塞”;这使得一套开局如此惊艳的理论的自然发育归于停顿。})与一门我确实梦到过的“Motif理论”(而且是一个精确得吓人的梦!)之间的联系——除非经由这一条:还是那位做梦人Grothendieck,(凭“最大的自然一般性”的本事!)硬要向真正的数学家们提示一些类比,而做正经工作,则是他们的事了。

至于那条大名鼎鼎的“‘相交上同调’的代数机制”,我们这就到了“反常层会议”的正中心\fnote{11}{(**)}{关于这场研讨会,见第VII行列“研讨会——或Mebkhout的层与反常性”。}(尽管“反常”一词并未出口)。对“IHES的四枚菲尔兹奖章”中的一位\dmark{},人们毕竟还戴上手套——考虑到场合的庄严;但对这同一位Grothendieck的遗孤学生,就无须顾忌了。在这个特殊的场合——聚光灯下、部长致辞、等等等等——我自己的下葬,不是经由沉默的下葬,而是经由\textbf{恭维}的下葬:调配得当,按方服用。\origpage{452}但明摆着:在 Mac~Pherson 与 Goresky 被点名之处,对遗孤学生Zoghman~Mebkhout,沉默才是规矩——一如两年前的反常层会议,一如今天依然如此。

\subsection{(2) 力量与光环}
\label{ch18:local:h3}

\notehead{105}{}
\phantomsection\label{ch18:local:n2}

(9月29日)那条“前面”的注——“葬礼悼词(1)——或恭维”(n\textsuperscript{o}~104)——写于5月12日,距今已四个多月。它当初是作为“拒受遗产,或矛盾的代价”一注(注n\textsuperscript{o}~47,3月末)的一条脚注起手的,为的是顺带记下我刚察觉的一桩小“怪事”。可是写着写着,我顺着行与页发觉:我正在评点的这两篇看似平淡的短文,竟是一座真正的“矿”\fnote{12}{(*)}{关于此事的回顾性评语,见9月24日“表面与深处”一注(n\textsuperscript{o}~101)的开头。}——我事先既没有料到,也没有刻意求之。那一天,也正是我刚画完一幅屠杀图景之日(注n\textsuperscript{o}~87);那幅图景,在此前的几个星期里,已从迷雾中渐渐显形。在这里,它骤然落实,单凭一段铺陈的描述之力取得肉身;此刻它正在用力唤我。屠杀,与“恭维”——对可敬的亡者的葬礼悼词——宛如同一幅触目图景的相辅相成的两联,在同一天显现!

这当然足以使我满盈!次日起,“我的手发痒”要乘胜追击,尤其要更深地探测那座我意外按上手的小小矿中珍宝。变得清楚的是:头一件该做的事,是把小册子里那两段文字全文引出——同时也是更好地接触这些文本、更好地浸透其真正讯息——那份“字里行间”的讯息——的最好办法\fnote{13}{(**)}{对此可参见“论破译一份讯息的艺术——或书写的颂歌”一注(n\textsuperscript{o}~102),它紧随前一条脚注所引的那条注之后。}。\origpage{453}甚至\dmark{}还没来得及把那两篇文本誊清,前一天的接触就足以在我身上唤起或唤醒好几个联想;我感到它们多汁。我急于追随它们,虽然还不大知道它们会把我引向何方……

最终,在其后的日子与星期里,我并没有沿着这条势头接下去——尽管整整这段时间里,我一直许诺自己近日就回到它。一桩意外的“健康事故”,使《收获与播种》的一切反思工作、乃至一切智力工作,中断了三个多月\fnote{14}{(*)}{对此可参见“事件——或身体与精神”与“陷阱——或省力与衰竭”两注,n\textsuperscript{o}~98、99.}。那几天刚刚敞开的、有利于沿这个方向继续反思的“热的时刻”,从此过去了。它是否会回来,甚至我是否肯下力气去“吹”(把热的吹回来!)以便不惜一切代价唤回它,都不确定。干脆说:我此刻真正的愿望,是走向那条终注,为整场名为葬礼的反思做一份临时\textbf{盘点}——并画上一道\textbf{最后的线}!至于本注,我至少要兑现我许下的(而且外加许给读者的)全文引证;或许还会就这两篇文本(或许还有把它们白纸黑字重抄这件事本身)在我身上激起的某些联想,给几句简要的指示。

所说的两篇文本(分别载于1983年题名《高等科学研究所》的庆典小册子第13与15页),属于一组“速写肖像”:自1958年建所以来在IHES任过职的“终身教授”与“长期访问者”,按到所的编年为序。那是相当简短的文本,各约半页,各带在所年月与身份(教授或长期访问者)、主要荣誉、主要兴趣领域与最重要的贡献,并(视情形)附某些合作者的名字。然而对我寒微的身量而言,关于这后三个“客观”侧面——兴趣领域、主要贡献、主要合作者或学生——的显著空白,却由那些酒神颂风格的“恭维”填满;其中一些,前一条注里已经检出并引证……\dmark{}\origpage{454}

我有幸开篇的这一系列,由以下数学家与物理学家构成:A.~Grothendieck、L.~Michel、R.~Thom、D.~Ruelle、P.~Deligne、N.H.~Kuiper、D.~Sullivan、P.~Cartier、H.~Epstein、J.~Fröhlich、A.~Connes、K.~Gawedzki、M.~Gromov、O.~Lanford。

我记得Dieudonné与我同时在IHES任教授;而在这份名单上我发现并非如此——他想必只满足于主持《数学出版物》。然而我现在发现,在小册子第3页,IHES的“履历”里,事实并非如此:Dieudonné分明同我一样,自1958年(直到1964年)是“终身教授”——至少名义上如此。一桩有点奇怪的矛盾!我把“履历”的头两“条”抄在这里,1958与1961:

\begin{quote}
\textbf{1958}高等科学研究所协会由Léon~Motchane在巴黎创建,辅以声名世界级的科学顾问团……与一批欧洲工业家。科学活动由两位数学家启动:Jean~Dieudonné(→1964)与Alexandre~Grothendieck(→1970),聘为终身教授。《IHES数学出版物》第1期出版。

\textbf{1961}承认为公益机构。

……
\end{quote}

顺带我注意到:在这份简短的履历里,似乎有必要提到(略带象征意味的)《数学出版物》第1期的出版(它由一篇24页的文章构成,作者 G.E.~Wall与这个新生机构并无特殊瓜葛),却不提代数几何研讨班——它们以家喻户晓的缩写SGA~1与SGA~2为人所知:我正是经由它们,在那个机构几乎还只“纸上存在”的年月里,开始独自支撑起它的科学声誉。况且,直到《数学出版物》大约第24卷为止,这些出版物的主体,是由《代数几何原理》(EGA)的相继各卷(1至4)构成的\fnote{15}{(*)}{我是其作者,与 J.~Dieudonné合作。};其余各卷,每卷围绕五十页上下打转(其科学水准之高,自不待言\dmark{}\origpage{455})。另一方面,在第19页(在“速写肖像”系列——Dieudonné缺席其中,天知道为什么\fnote{16}{(*)}{(9月30日)一个念头来到心头:缘由恐怕是这个:为了不必说,在所论的年月里(1958--1964),Dieudonné的时间,主要用于写《代数几何原理》(在那里,我不合时宜地俨然第一作者)与Bourbaki的成文工作——更不必说钢琴与烹饪(Dieudonné既是行家音乐家又是行家厨师);这些,在那本小册子里,当然无从谈起:它太挑剔了,连一个顺带的微笑都挤不进去……}——之后),人们读到一段排版极“广告招贴”的文字(配一幅诱人的照片:那套显赫出版物全帙堆成的巍峨书山):

\begin{quote}
数学出版物

正是Jean~Dieudonné,独自[!],自1959年起把《数学出版物》带上了世界卓越的峰巅。

自1979年起,它们以每年400页的定期刊物形式出版,由一个编辑部主持,主编是Jacques~Tits。

发行由……承担(等等)
\end{quote}

倘若《数学出版物》被以\textbf{这种}方式——就在那里,在这份显赫机构的庆典呈献里(它的主要天职,从来不是当一份期刊的出版者)——捧上云端,那么毫无疑问,这是为了叫某些人忘掉一桩不太愉快的事实\fnote{17}{(**)}{这可委屈了我的朋友尼科(他当时已是那所庆祝周年的机构的十二年所长);他想必(在这个场合一如在别的场合)对此什么也没看出来……}:倘若不是在某几个关键的年头里,某一位无名氏——他在自己的角落里执拗地推进自己的一些想法(它们有幸咬住了某些人,包括“上流社会”里的人)——顶风冒雨\fnote{18}{(***)}{顶风冒雨:在这四年里,我全然没有被那些持续的警告与流言吓住——说一场(如消息灵通的朋友们所暗示的……)全然不现实的冒险、甚或带点骗局味道的冒险,眼看就要破产!事实是:IHES当时没有任何财务或地产的根基;它的生命,始终悬在几位心境或好或坏的工业家的短期捐赠上。我对此并不太操心,只是信任创办人兼所长Léon~Motchane:他年复一年,凭财务戏法与“公共关系”的奇技,“保住了局面”。毕竟,在那个宽厚的年代,倘若它塌了,我很有希望很快找到一个不那么成问题的落脚点!相反,倘若我赢了我押在IHES上的这一注(有Dieudonné的鼓励——他认识Motchane,而我对他全然信任),我在IHES的职位,便比我所知道的任何别的位置都合我的心意。}\origpage{456}带给它一份担保与信誉——那是天底下最漂亮的协会章程、乃至最漂亮的“声名世界级的科学顾问”(原文如此),\dmark{}都无力给予的——那么,这家机构恐怕早已自生自灭,被人遗忘了很久。

(9月30日)这份庆典小册子(我会越来越熟识它的!)的“炫示”与“四面八方涂脂抹粉”的风格——抱歉,我是想说(极)高规格的“公共关系”——当然不是我朋友皮埃尔的风格,也不是尼科的:他们两位肯定有更要紧的事,犯不着撰写这类应景文字。相反,显而易见:占据着我的这两篇速写肖像——一篇写我,另一篇写Deligne——并非不经后者至少供给关键词就写成;单凭他是IHES唯一在行的人,就足够了。而同样对我分明的:这两篇文本,至少,未经这同一位Deligne先读过并点头,是不会付印的。于是,一开始对我就很清楚:这两篇文本,无论如何,首先映出我朋友的心境与意图——他力图就我的这个人、就他自己的这个人,向他自己、也向数学公众给出的形象。当然,正是以这个身份,这两段文字使我感兴趣。这份兴趣,并不取决于Deligne是否这些泄露心迹的文字的作者,或作者是别人(想必是把整本小册子“构思”成篇的那一位):后者出于这种或那种缘由,大概也认领了这个我的朋友执意要传递的“讯息”。

到底了,这两篇速写肖像如下,摘自题为“终身教授与长期访问教授的活动”的肖像画廊(第13--19页)。

\begin{quote}
Alexandre~Grothendieck,数学家,1958至1970年任IHES教授,菲尔兹奖章得主。

在他留在研究所的12年里,A.~Grothendieck革新了代数几何的基础与方法,并为它开启了新的应用,尤其是算术方面的。他在IHES创建了一个代数几何学派,围绕他主持的研讨班聚集起来,并以他传播自己观念时的那份慷慨为滋养。\dmark{}\origpage{457}其著作的巨人般面貌,映照在其出版物中:与Jean~Dieudonné合作的专著“代数几何原理”(8分册),以及与众多学生合作的“布瓦玛丽代数几何研讨班”12卷。

在代数几何中,他提炼出本质的问题,并给每一个概念以最大的自然一般性。所引入的观念,显出远超代数几何的本质性。它们往往显得如此自然,以致我们难以想象它们曾费过多大力气。倘若它们今天显得理所当然,这想必得益于他对术语的高度留意。

还应当记得:“Grothendieck群”——在代数几何中与相交理论相连,并在拓扑中被使用——是代数K理论的始祖。为把先前在$\mathbb{C}$上经拓扑途径证明的结果移植过来而在一般基底域上的代数几何中引入的topos,如今用于逻辑。

他在1970年离开IHES,那时他对数学的热情正在消退。难道应当认为,他沿着自己划定的道路提出的问题,已经变得过于困难了吗?

\begin{center}$\cdots$\end{center}

Pierre~Deligne,数学家,自1970年起任IHES教授,菲尔兹奖章、Henri~Poincaré金质奖章得主,法国科学院外籍院士。

他的研究以“理解代数簇的上同调”为主轴。若复代数簇$X$是非奇异射影簇,调和积分理论便在$H^*(X)$上给出一个Hodge结构。他由此出发,并借助Grothendieck所提示的$\ell$进类比,提炼出混合Hodge结构的概念,并使每个复代数簇的上同调都具有这种结构。在$\ell$进上同调中,因而对于有限域上的簇,他证明了难度之高众所周知的Weil猜想。这一结果显得更加令人惊讶,因为Grothendieck在构造了任意域上的上同调理论之后,曾把尚余的猜想归结为一系列猜想,而这些猜想今天仍与当时一样难以攻克。

这一定理促使$\ell$进上同调成为一种强有力的工具,可以用于看似远离代数几何的问题,例如Ramanujam猜想。

\origpage{458}最近,他还研究了阿贝尔簇上的Hodge循环,朝Grothendieck梦寐以求的“动机理论”迈出了第一步。他也证明了“交叉上同调”——Mac Pherson与Goresky的拓扑理论——的代数机制。由此得以将它移植到$\ell$进理论中,而它在那里显得出人意料地有用。

他目前对非交换调和分析感兴趣(实或$p$进Lie群——或有限经典群——以及某些齐性空间上的函数理论);这延续了他关于自守形式(Ramanujam猜想),以及他与G.~Lusztig合作开展的关于有限群表示的研究。

他吸收、洞察整个数学领域的速度极快;因此,面对人们向他提出的每个问题,他都能作出富于启发且有建设性的回应。
\end{quote}

这两篇文本,还当补上第三篇:在其中,Deligne与我一口气并列。我在小册子里散页的一张上找到它;标题与肖像画廊所在的那章相同——“IHES的研究取向”,副标题:“关于‘科学活动之前景’的概要”。它实质上是那幅肖像画廊的一个苛刻的“缩编”——这一回只收“终身教授”(在任的与往任的)\fnote{19}{(*)}{(10月1日)为凑“分量”,Connes也被塞了进去(尽管他不过是“访问者”);这总归给收藏家又添一枚“菲尔兹奖章”。相反,我的朋友Nico~Kuiper却被弃置不顾。为这场合隐身让位而出难题的,想必不会是他……},每人两三行。他们依次是(按被提及的顺序):我自己、Deligne、Michel、Thom、Ruelle、Sullivan、Connes、Lanford三世、Gromov。这与那幅更详尽的肖像画廊的次序相同;只除了一点:这一回,Deligne“升了位”,以便有幸与我一口气并列。有趣的细节:在这份文本里,被巡视的显赫人物的名字,全都带着着重号出现;唯独我寒微的身量例外\fnote{20}{(**)}{(10月1日)这一手光辉的排印花招(其意图或许并非有意)取得的排版效果是:将要引出的这一段,俨然是献给Pierre~Deligne的——他的名字在字形上俨然“终身教授”世系之首,把我排除在外;而我在其中俨然一位“合作者”、机构的局外人!编年的次序倒是遵守了,无可挑剔——然而产生的效果(而且肯定是有意求得的),是一场角色的\textbf{颠倒},在我身上唤起一些熟悉的联想(“颠倒”“驱逐”“暂停”等注里谈过,n\textsuperscript{o}~68$'$、63、77)。于是我同时也认出那套据有的\textbf{做派}——“暂停!”的做派——它向我明白无误地指出这份讯息的\textbf{真正}作者。}!以下就是关于我的朋友与我的那一段:

\origpage{459}\begin{quote}
Alexandre~Grothendieck深邃得近乎传奇的理论与Pierre~Deligne的光辉发现(两位都是菲尔兹奖章得主),经由“跨学科”的手段(上同调),把拓扑、代数几何与数论连在了一起。就在最近,这使得联邦德国的 G.~Faltings(他已在IHES工作过)得以证明一条艰深的定理:它在数论中立标,并照亮了那条大名鼎鼎的“Fermat定理”。
\end{quote}

顺带我注意到:“菲尔兹奖章”在这幅微型画廊里配得上一个大写的M;而“跨学科性”,自IHES建所之初,就是其创办人的一大偏爱的主题。或许正是托这份情势的福,在这份摘要里,人们似乎终究暗示:我的这个人,兴许与某一种称为“上同调”的“跨学科手段”有点关系(而它碰巧也是Deligne著作的“主导轴心”——不知是哪门子的偶然)。

\dmark{}可瞧,我又开始从细处抠这份文本了!那句应景的Faltings引证——他一夜之间以其轰动性的结果(这里称为“艰深”,仿佛问题就在这里;不过对我的话题无关紧要……)蹿升到科学时事的头排——也属于文本的“细处”:说到底,是那位抄写者的“落款”,不值得我停留。分明是对Deligne与我的第一句话,才载有该段的“讯息”之本质。

它就我朋友兼前学生身上的某些心境,向我透了不少——首先是一份深层的“Unsicherheit”(不安全感,缺乏底气,缺乏深层的立足)\fnote{21}{(*)}{这个来到我心的德文词“Unsicherheit”,在法语里没有对应词,(我相信)英语里也没有。其直译“不安全”几乎不能用来指一种心灵特质。那个否定的词“缺乏底气”,是另一种将就的近似。应当理解:这里的“底气”是在一个深层的层面上;其缺乏,可以在某些时刻被知觉,而表面上占上风的,却是一种有底气、完美从容的印象;它们如同一副保护性的甲壳,带着往往可观的惯性与“坚固性”,无懈可击……}。这里,一如在他署名的已刊文本里的任何一处\fnote{22}{(**)}{至少在我迄今过目的那些里。}、一如在前两篇速写肖像里,没有任何东西\origpage{460}能使人设想:我的朋友可能在某个时刻从我这里学到过什么。可瞧,他在这里,以清清楚楚的措辞,把自己呈现为一份宏阔的统一视野的\textbf{另一位父亲}——那视野是“取”自他人的\fnote{23}{(*)}{这里还有一层特别的讽刺:这份视野——在这里被当作他自己的“光环”取自他人——实际上,自“师父亡故”以来,恰恰被那位俨然继承人、一边划清界限一边否认遗产的人,交付给轻蔑,并被系统地阻挠。对此可参见“继承人”“共同继承人……”“……与电锯”三注(n\textsuperscript{o}~90、91、92);其他例证,见第X行列(灵车),由四口“棺材”1至4与掘墓人构成(注n\textsuperscript{o}~93至97)。},仿佛被一种亲身的信念所慑服:相信自己深深无力构想并让那些\textbf{他自己的}视野——同样宏大乃至更宏大——在他身上发芽绽放;仿佛为了“是”并“显得”伟大,他从此只剩下那份寒微的资源:把那道光环\textbf{记在自己的账上}——那道光环,他年轻时曾乐于环绕在一位显赫、如今已亡故(或无论如何,被一份天赐共识宣告如此……)的长辈身边。去攫取一道\textbf{光环},而不是让自己身上那些尚且无形无名的、等待出生并被命名的东西发芽绽放——而不是活出栖于他身上的、同样在等待的\textbf{他自己的力量}……

(10月1日)昨夜,我觉得自己又一次触到了那桩冲突的心脏——那桩冲突,早在《收获与播种》的最开头(在“他人之无谬与对己之轻蔑”一节,n\textsuperscript{o}~4),即将八个月前,我已以一般的措辞谈到它;而临近葬礼的开端(在“死结”一注,4月26日的n\textsuperscript{o}~65),我又“在一个极端而格外触目的个案里”重新找到了它。这又是一场不期而遇:在一处引文的转角——我求个心安,终究把它与前两处一道收了进来!几天前,重翻那本大名鼎鼎的小册子时,我就已注意到了这一段;当时它确实打动了我,但我没有停留。可昨天,一经白纸黑字写下,它立刻显得比那两段我刚抄的、本该构成我正在写的这条注之主题的应景文字,意味更重、更加触目。然而,在那两段文字里,并不缺使人“翻然醒悟”的地方;它们本会在四个月前还使我毫不迟疑地\dmark{}紧接着再写十页——倘若不是二十页——的联想。但忽然我觉得:我本可以这样展开的东西,说到底——至多除一处例外——\origpage{461}是已知的东西的再度确认,只是角度或许稍有不同;尤其:它们终究是些附带的侧面——那一类我已在5月“恭维”那条注里、乃至在整个葬礼反思中充分铺陈过的侧面。而第三段引文,却把我带回某种本质的东西——在那场漫长的“侦查”(葬礼上的工作是其中一部分)中,我一直有点失去视线的东西。

当时我一度想就此打住,至少不试图以言词把握那短短四行的碑文式句子对我说的话——它在某个层面上,分明已被“听见”。我终究越过了这份犹豫。词句上升得又慢又迟疑,而那起初弥散的印象,却随着书写渐渐澄清。一经白纸黑字写下、删去看似多余的部分,我便知道:我已把我“听见”的东西,把握到我所能把握的最好程度。

夜已深得犯规,真的该停在那里了。我心满意足地睡下,却还不确定:是否会把刚写下的东西,收进我这份注定要发表的见证。毕竟,我也满可以让读者——倘若他有兴趣穿透一份讯息的表面——自己去打捞他在其中听到的东西!只是到今天,我才知道我会把这段文字收进去:它分明表达着我(或自以为)对某件事的某种知觉或领会;那件事,在我看来是重要的,甚至是这场葬礼的一个深层内驱力的关键。

\section{阴阳之钥}
\label{ch18:local:h4}

\subsection{(1) 肌肉与肚肠(阳葬阴(1))}
\label{ch18:local:h5}

\notehead{106}{}
\phantomsection\label{ch18:local:n3}

(10月2日)我想再至少追一笔由三联式葬礼悼词(昨天我终于给出了全文引证)激起的联想。这一联想,早在5月13日——我刚写完“葬礼悼词(1)——或恭维”一注(n\textsuperscript{o}~104)的次日——就已强加于我。它触及事物的一个常常在不知不觉中错失的侧面;我只是近五六年来,才开始\dmark{}\origpage{462}真正意识到它。

在所检视的文本的字里行间,人们看到某些\textbf{价值}的崇拜在宣示自己。比如,关于Deligne所证明的Weil猜想,被突出渲染的,是它们的“\textbf{困难}”——而非其\textbf{美}、其单纯、其自Weil陈述它们的那一刻起便已敞开的宏阔视野\fnote{24}{(*)}{(10月3日)而且还被称作“尽人皆知的”困难!这几乎毫无意义,除非意在镇唬那些不在局中的人!一条猜想的“困难”,只有在其得到证明之后才可能真正掂量;而它的\textbf{多产},却可以预先感到,并且往往在证明之前,就已经由它所启发的种种工作,客观地显形。那些“大”猜想之有别于其他猜想,不在于其“困难”(那是未知的——姑且假定这个词还有点意义……),而分明在于其\textbf{多产}。顺带指出:这是事物的一个典型的“阴”的、女性的侧面,而“困难”是一个典型的“阳”的、“男性”的价值。}。我也想到这些隐约瞥见的视野在其被证明之前就已结出的果实,想到一旦越过通往其证明的漫长旅程的最后一步,便应时成熟的其他隐约瞥见的果实。正是这些猜想的美、它们内在的惊人自洽、以及它们使人隐约瞥见的此前无人料到的联系,使它们成为两代几何学家与算术学家如此强而多产的灵感之源。我著作的最深部分(无论“彻底做到底”的那部分,还是“Motif之梦”)都直接受启于此(经由Serre的中介——他善于领会并传递其猜想中所表达的那份视野的全部力量)。没有它们,无论$\ell$-进上同调,乃至topos的语言,恐怕都无从问世。或者不如说:我在自己生命中大约十五年间执意发展的那幅关于(代数)几何、拓扑与算术的“宏阔统一视野”,正是在这些“Weil猜想”里,找到了它的一幅最初的、触目惊心的草图。而随着这视野在幅度与成熟上步步增长,恰是这视野本身、连同它使人得以一一把握的那些此前隐藏的事物,一步一步向我低语:该做什么,该从哪个头“下手”眼前伸手可及的东西。Weil猜想证明中的最后一步,不多不少,正是一场漫长而迷人的旅程中的一步——这场旅程始于何时,我说不上;肯定远在我的诞生之前,而且在我的死亡之后,也远不会终结!

\dmark{}\origpage{463}但照所引文本里可以辨出的那股精神,人们会以为“Weil猜想”是一桩举重的事:瞧,要“生拉硬拽”举起的重量!两百公斤,可不是闹着玩的;困难尽人皆知;很多人试过,还没有一个人成功——直到“H日”(“H”如“赫拉克勒斯”)!结果是惊人的(注106\textsuperscript{1}):您掂掂,整整两公担——本来谁也不信有人举得起来……在论Faltings所证“艰深定理”的那句简短评语里,人们知觉到的,也是同一股精神:在那里,同样,在这件我们对事物之认识的新台阶的命名本身之中,被突出渲染的又是\textbf{困难},以博取芸芸众生的叹服——而非那自一座新攀上的峰巅敞开的视野\fnote{25}{(*)}{而最打动我的,从我亲手拿到Faltings证明三条关键猜想——包括此处所说的Mordell猜想——的预印本那一刻起,恰恰是其进路的惊人的\textbf{单纯}:他凭四十来页,证明了这些本被当作“高不可攀”的结果!(可与注n\textsuperscript{o}~3对读。)}。连“Mordell猜想”这个名字(诚然,数学之外的公众不认识它)似乎都不值得提——仿佛对猜想的领会与陈述(在这里,是Mordell的)是件附带事,因为“容易”。取而代之的,是对“Fermat定理”(据说它被“照亮”了)的一记凭空的展望。诚然,后者举世皆知(甚至在数学界之外),作为一块足足三百公斤的重量(它顶住了三个世纪的攻势)。

我头一个想说的要点是:这些文本里(带着合乎场合的审慎)被高扬的价值,是那些可以称作\textbf{肌肉}的价值——在这里,即“脑力肌肉”的价值:那种使人得以凭手腕之力,超越那些尽人皆知的“困难”纪录的价值。

这些价值,不只是这里被捧上云端的这位英雄的价值,也不只是某本庆典小册子的作者的价值(作者匿名;我自信认得出)。它们也是(似乎)越来越支配数学界、进而更普遍地支配科学界的价值。甚至越出这个还相对有限的世界的边界\dmark{}\origpage{464},可以说:它们也是、而且越来越是一种被称作“西方的”特定“文化”的价值\fnote{26}{(*)}{当我在这里谈到今日所见的我们“文化”的“价值”时,我指的当然是那些“官方的”价值——由学校、媒体、家庭所传递,并在各职业圈子里获得普遍共识的价值。这并不意味着这些价值被人人毫无保留地接受,也不意味着它们构成所有人的态度与行为中的基调。而且,正是怀着痛心,正派的人们、媒体与称职的专业文献(出于教育家、社会学家、精神科医生等的笔下),谈到某“一部分青年”:他们分明不“合调”,有损一幅体面的图景!}。久已以来,这种“文化”及其价值已征服了我们星球的表面,消灭了所有其他文化——其优越性无可辩驳的明证。这些价值的行星象征、其英勇的化身,就是密封甲胄里的宇航员:在千百万屏息凝神、瘫坐在屏幕前的电视观众的注视下,第一个踏上某颗无法想象的遥远而荒凉的行星。

这些价值——由于无从更贴切地把握,我姑且以一个带象征意味的粗略字眼“肌肉”称之——并非始于昨日。用人类学的行话说,人们也可以称之为“父权制的”。依我看,最早以不容置辩之力(!)宣示其至尊地位的成文文本之一,是《旧约》(尤其摩西诸书)。然而,只需读一读这份远古时代的迷人文献,便会发觉:“父权”价值的至尊——男人之于女人,或“精神”之于“身体”或“物质”——远未走到否定或轻蔑那些补充价值的地步(那些价值,当时或许还不曾被知觉为“对立的”或“敌对的”\fnote{27}{(**)}{比如,对母亲的崇敬,是犹太文化里一种扎根极深的传统;它想必对“圣典”里所高扬的那些(不妨这样说的)“官方”价值起着补偿的作用。这一传统,以一种改造而更炽烈的形态,重现于天主教传统之中,即对(童贞女!)马利亚的崇敬。})。我不知道这两组互补价值之坎坷浮沉的历史是否已被写出——追索这段历史,穿越一个个世纪与千年,从摩西的时代直到今天,想必是一件迷人的事。它想必也是一部历史:一种“价值”平衡——一边是“父权的”或“男性的”,另一边是“母权的”或“女性的”;一边是“肌肉”与“肚肠”,“精神”与“物质”——的逐渐败坏;这败坏分明朝着“阳刚”(或在东方传统辩证法里的“阳”)价值的方向,以牺牲“阴柔”(或“阴”)价值为代价。

\dmark{}\origpage{465}我觉得,我们这个时代的特点,正是这种文化败坏的变本加厉。这段历史的最近的几幕,是两大对抗的超级大国(浸透着本质上相同的价值)之间的“太空竞赛”与(尤其是核)军备竞赛。作为这类“力量”或“权力”的加码竞赛中一场狂飙演进的终极一幕与可能的结局,人们此刻就可以预告:某场行星规模的核浩劫(或别的什么——可挑选的余地有的是……)。它也许有这份功劳:一举而永远地解决所有问题……

不过,我在这里的主旨,并不是要描画一幅诱人的“世界末日”图景(这事不必等我),更不是要向“肌肉”、或向“大脑”(别名“精神”)宣战。我很清楚:连我的“肚肠”也从中捞不到什么!我珍爱我的肌肉与我的大脑——它们对我有用,自不待言;一如我也珍爱我的“肚肠”,它同样少不了。毋宁说,我觉得有用的,是在这里以几句话(倘能做到)说明:这场由周围文化所传递的、在这两类价值之间的深层冲突,是如何在我自己的身上上演的。用更接地气的话说,这也是我对自己这个人身上两个同样真实、同样具体、按本性不可分而互补、本身绝非敌对的\textbf{侧面}或\textbf{面孔}的态度(接受乃至推崇,或拒斥)的沿革。我可以称它们为我身上的“\textbf{男人}”与“\textbf{女人}”,或者(取两个不那么“有负担”、因而不太容易误导的称呼):“\textbf{阳}”与“\textbf{阴}”。

在大多数人身上,“骰子似乎早在幼年初期就已掷定”:那些将在一生中默默主宰态度与行为、带着一台调校完美的自动机之效率的本质机制,就在那时各就其位。这些机制的核心,是对我们身上这样或那样特质、或这样或那样深层冲动(带着阳的或阴的“署名”)、或这样那样一组给定署名的特质与冲动的“包裹”、乃至对“阳”包裹或“阴”包裹之全体的肯定或拒斥的机制。正是这些机制,在很大程度上,决定着我们“自我”之结构的其余一切选择(肯定或拒斥)机制。

\dmark{}\origpage{466}由于一些至今对我仍是谜的缘由,在我自己的个案里,“自我”(“老板”)与我这个人身上的“男性”与“女性”之间的关系史(无论有意识的还是无意识的;无论在“老板”本身之中,还是在“工人”之中——这两者都系于万物的阴阳双面)——这段历史,比寻常的要动荡得多。我在其中辨出三个时期。最后一个时期,在某种意义上,又接上了第一个时期:它铺展在我童年最初的五年。这第三个时期——我可以称之为\textbf{成熟}的时期——可以被看作向那个童年的某种“回归”,或看作与“\textbf{童年状态}”的逐渐重逢,与我存在之中“阴”与“阳”毫无风波的联姻之和谐的逐渐重逢。这些重逢,始于1976年7月,时年四十八岁——也正是在那一年,我(三个月后)发现了此前一直不为我所知的一种力量:禅观的力量\fnote{28}{(*)}{见“欲与禅观”与“惊叹”两节,n\textsuperscript{o}~36、37.}。

我的父母——母亲与父亲——各自身上占支配地位的价值,都是阳的价值:意志、智力(即:智力上的力量)、自制、对他人的支配、不妥协、“Konsequenz”(德语,意为对——尤其意识形态的——取舍的极端一贯)、“理想主义”——无论在政治层面还是实践层面……在我母亲身上,这份价值化自少女时代起就取得了一种加剧的力量;那是她对“她自己身上的女人”(并由此,对一般的女性)所怀的一份真正的憎恨的背面。她身上的这份憎恨,最终取得了一种愈发放肆、愈发摧毁性的力量——恰恰因为它终生全然被遮蔽。(我自己直到五年前——禅观出现在我生命里三年之后——才发现这些事。)在这样一个双亲的语境里,我竟能在我童年最初的五年里全然绽放——直到那场从父母身边的剥离、直到我的原生家庭(由我父母、比我年长四岁的姐姐与我构成)的毁灭——这是一个谜(然而是一个对我毫不含糊的事实);那毁灭,出于我母亲的意志,并借1933年政治事件的东风(倘可这样说)。

\notehead{106\textsuperscript{1}}{}
\phantomsection\label{ch18:local:n4}

\dmark{}(10月3日)Weil猜想可能不成立——对此,我与Deligne从不曾有丝毫怀疑;我也不记得听过谁表达过这样的怀疑。把那个“结果”(即这些猜想的证明)称作“令人惊讶”,再度暴露了那种蓄意哗众的用心。况且,自“拓扑”与平展上同调引入以来的任何时刻,我从不曾感到这些猜想高不可攀;毋宁说(自1963年起),它们不出近几年就必被证明。1970年我离去之时,我几乎毫不怀疑:Deligne——他是所有人中最有资格的——不会迟迟不去证明它们(他也确实没有迟疑),同时还有“代数循环的标准猜想”——那些更强的猜想(他却转而去败坏它们的名声)。

\origpage{467}Deligne对这后一些猜想的效验有所保留,他做得有理;我同他一样不确信。但一条猜想的分量,不取决于它到头来是真还是假,也不取决于其所谓的“困难”品格——那品格会使它“高不可攀”;那品格,本是纯主观的。它唯独取决于:该猜想把手指按上去的那个\textbf{问题}(在它被提出之前,无人知觉过这个问题)——这个问题,是否触及对我们的认识而言真正本质的东西。而这一点是明摆着的(至少对我而言!):只要这些猜想之效验的问题未获解决,人们就不可能有对代数循环的良好领会,也不可能有对代数簇上同调的那些所谓“算术”性质(或再或者说,“Motif几何”)的良好领会。今天,一如1968年的孟买大会之时,我把这个问题——连同奇点消解问题——视为代数几何里摆在我们面前的最根本的两个问题之一。我分明感到两者的分量!这份潜在的多产,必定会显现出来,只等人们不再满足于跌跌撞撞地绕开一条被判定“过于困难”的猜想,只等有人终于肯挽起袖子、跟它真刀真枪地干一场!

\subsection{(2) 一段生命的故事:一个三乐章的循环}
\label{ch18:local:h6}

\subsubsection*{(a) 天真(阴与阳的联姻)}
\tocsubsub{(a) 天真(阴与阳的联姻)}
\label{ch18:local:h7}

\notehead{107}{}
\phantomsection\label{ch18:local:n5}

\dmark{}(10月4日)我已有机会提到我一生最初五年的一个重要侧面,把它称作一项头奖的“特权”\fnote{29}{(*)}{见“屠杀”一注,n\textsuperscript{o}~87.}:一种与我的父亲之间深的、毫无问题的认同;它从不曾被恐惧或嫉妒触及。这一情势——连同这种对我父亲的认同的存在本身、及其无声的力量——我是仅仅四年前(在那场1979年8月至1980年3月关于我父母的禅观之后、对童年与我此生的禅观之中)才意识到的。这份认同,宛如另一种认同的宁静而有力的心脏:对我父母、我姐姐(她比我年长四岁)与我组成的那个家庭的认同。我对父亲一如对母亲,怀着无涯的敬佩与爱。他们的身量,对我而言,是万物的尺度。

这绝不意味着我对他们的态度是一种自动的赞许、一种景仰的憨态。我当然不知道他们是我万物的尺度;但我十分清楚,他们同我一样会犯错误;我心中没有任何恐惧,会妨碍我察觉一个分歧并把它明白表示出来。在我周围的种种冲突里,我不怕以我自己的方式表态。这丝毫触及不着那份信、那份安然的确信——它们构成我存在深处不可动摇的立足;毋宁说,这恰恰自发地出自那份信、那份确信本身。有过这样的时候:我父亲在无可奈何的愤怒发作之中——而我姐姐(不动声色地)以撩拨他为乐——动手粗暴地打她;而每一次,我都义愤填膺,怀着对我姐姐毫无保留的团结。我想,这些是我与“我的父亲”的关系里仅有的几片过路的浓云(与母亲之间则没有)。这并不是我赞成我姐姐那些有时挺可恶的把戏;我想它们也并不真正困扰我——对我而言,万物的尺度不是\textbf{她}。她的把戏(其缘由,想必同对我父亲一样,向我滑脱;而他回回“中计”,或我母亲则无论事先事后都小心不介入)——这些把戏,在某种意义上,对我并不真正产生后果。她是我的姐姐,她就是她那样,如此而已。可是\textbf{我的父亲}竟放任自己那种盲目的粗暴……

最亲近的三个存在——它们一起构成我最初年月的母体——被冲突撕裂着:冲突使每一位既\dmark{}\origpage{469}反对自己、又反对另外两位:我母亲与我姐姐之间,是一场阴险的、不动声色的冲突;我父亲与另一边的我母亲及我姐姐之间,则是一场激烈爆发的冲突;她们各用自己的方式(而且在我父母生前,谁也不曾装作察觉……)牵着他的鼻子走。神秘而非凡的是:在人生这些最敏感、最关键的年月里,被如此包围在冲突之中,这场冲突竟对我保持\textbf{外在};它在那几年的光阴里,不曾真正“咬进”我的存在,不曾在我身上安营扎寨。

我存在之中的那份分裂——它给我的生命打下的印记,一如给任何别人的生命——并不是在那几年里进驻我的;而是在其后的两三年里,从我六岁到八岁上下。在某个时刻(我自信可以定位到误差几个月;它当落在我的第八个年头),在同父母(他们不大肯给我一点活着的音讯)与姐姐分离了两年多之后,发生了一次\textbf{翻转}。那首先是一场与我\textbf{童年}的决裂——从那一刻起,它被一套有效的遗忘机制“埋葬”(那套机制,差不了多少,一直延续到今天)。在某个深的层面(还不是最深的层面……),我那时把我的父母判成了“外人”,一如我的童年从此被判为“外来”。在某种意义上,我\textbf{退位}了:为了被接纳进从此环绕我的那个世界,我决定变得像“他们”、像在那里定规矩的成年人——去取得并发展在那里逼人敬重的武器,去在一个唯有某种“力量”才被接受、被看重的世界里,以对等的武器作战……

况且,那也正是我父母偏爱的那种力量——环绕我最初年月的,正是他们。而在这里,我回到了那件“神秘的事”(我刚循着另一条由它激起的联想的线索,离开了它):我生命最初年月里\textbf{我身上分裂的缺席}。

也许,谜已不在于这份缺席,而在于这一点:我的父母——父亲与母亲——那时竟各自把我\textbf{整个地}接受、彻底地接受:接受我身上“男性的”、“男人的”部分,也接受“女性的”部分。或者换个说法:我的父母,各自被冲突撕裂着,各自否认着自己存在的一个本质部分——各自无力对自己、对对方、对我姐姐,怀有一份慈爱的敞开……然而他们竟对我——他们的儿子\dmark{}\origpage{470}——找到了这样一份敞开,一份毫无保留的接纳。

再换个说法:在我一生最初的五年里,任何时刻,我都不曾体验过\textbf{羞于我之所是}的感觉——无论在我的身体及其官能上,还是在我的冲动、我的癖好、我的行动上。任何时刻,我都不曾须否认我身上的什么,以换取周围环境的接纳、与之和平共处。

当然,我也做过不“过关”的事:一如所有的孩子,当我来劲时,我肯定也惹人烦、乃至让人受不了;而有时明摆着,我得拨正船头。我不立规矩,也无意立——我无须补偿什么隐秘的残缺。在我父母对我的爱里,本不会有谄媚的位置、对任性耍脾气的纵容——不会有无条件的首肯。但是,倘若我免不了挨我父亲或我母亲的申饬(一如反过来有时也发生),在那几年里,他们谁也不曾使我因某个不中其意的行动或举止而蒙羞。

在一道对父亲毫无歧义的深深认同的底色上,作为孩子的我,今天在我眼前,显出同时打着男性气概与女性气质的印记,两者同样强。

我觉得,在每一个存在、每一件事物之中——在其阴性与阳性品质那些不可分解而流变的联姻之中,那使他成为他的联姻;其精微的平衡,正是活在这存在或事物之中的美、那和谐——在这种阴阳的亲密结合里,往往(或许永远)有一个基调、一个“主音”:或阴,或阳。这个基调,在一个人身上未必容易辨出:因为那些或多或少的有效而彻底的压抑机制,会歪曲牌局,以一副借来的形象替换原初的和谐。比如,我四十年来的“品牌形象”,是一副几乎纯然男性气概的形象——而且从来不曾被质疑、甚至不曾被(无论我自己还是他人,依我看)辨认为如此;直到我四十八岁。然而我倾向于相信:与生俱来的那个基调,终生都在场——至少在某些或许永远无缘公开亮相的深层。在我自己的个案里,奇怪的是,直到今天我都说不好这主音是什么——也就是\dmark{}\origpage{471}那个浸透了我最初的童年、在我出生时就已是“我的”的主音。不止一次,各种征兆使我疑心:这个基调是“阴”的;当我的存在得以自发显现之时——在那些它摆脱了自我童年起积攒在我身上的种种制约的时刻——主宰它的,是“女性”的品质。换个说法:也许,在我身体与精神之中充当创造之力的那个东西——我有时称之为我身上的“孩子”或“工人”(以别于代表自我结构的“老板”,即我身上受制约的部分、我个人累积的制约之总和或结果)——这份力量,或许更“女性”而非“男性”(尽管按其本性与必然,它两者都是)。

这里不是逐一巡视这些“征兆”的地方。况且,要紧的并不是:我身上那个深层的主音是“女性”还是“男性”。要紧的毋宁是:在每一刻都知道\textbf{做我自己}——毫无迟疑地接纳我身上使我成为“女人”的那些特质与冲动,一如接纳使我成为“男人”的那些——并容它们自由地表达自己。

在孩子时代的最初几年,常有不熟的人把我当作女孩——而这从来不在我的心头造成丝毫的不自在、丝毫的不安全感。我想,主要是我的声音造成这种效果:一副非常清亮、尖细的嗓子——更不必说我有一头长发(多半蓬乱着),也许仅仅因为我母亲(她另有操不完的心)不常费心替我剪一剪。此外我壮得像头小牛,略带粗暴或冒险的游戏并不使我生厌;这丝毫不妨碍一份对寂静乃至孤独的偏好,以及一份对抱娃娃游戏的偏好\fnote{30}{(*)}{这一偏好在小男孩身上显得少见,我相信主要是因为它被周围环境系统地劝阻。}。我不记得有人为此嘲笑过我;但这类事想必难免此处彼处地发生过。倘若这类小插曲过去而不留创伤或羞辱的痕迹,那肯定是因为:它们在我身上找不到任何回音或放大器——找不到任何不安全感可资凭依;而“我之所是”被那些唯一真正于我有分量的人所接纳,是超出一切疑问的。嘲笑本伤不到我;它只会反噬那位发难者——在我眼里,那人只会显得十分愚蠢:竟装模作样地对天底下最自然的事情挑三拣四。

\dmark{}\origpage{472}况且我深知:这类有点古怪的蠢事绝非稀罕;单是裸体,就可以成为惊世骇俗的缘由!然而,就我记忆所及,我早已有机会见过我的母亲、父亲与姐姐的裸体;也有的是机会满足我那份正当的好奇:看看他们每一位连同我自己,究竟是怎生造出来的。分明,男人的或女人的身形构造里,没有任何惊世骇俗之处;它们在我看来,就那样,好得很——尤其(我对此毫不讳言)女人的身形。

\subsubsection*{(b) 超级父亲(阳葬阴(2))}
\tocsubsub{(b) 超级父亲(阳葬阴(2))}
\label{ch18:local:h8}

\notehead{108}{}
\phantomsection\label{ch18:local:n6}

(10月5日)1933年——我生命的第六个年头——安放了我一生中第一个关键的转折;它同时也是我母亲、我父亲的生命中,在他们彼此的关系、他们与孩子的关系中,一个关键的转折。那就是我们四人组成的家庭遭受暴力的、决定性毁灭的情节;四十六年后,头一个也是唯一一个对它作出盘点、并追索其来龙去脉的人,是我——经由我父母的通信、以及一两份苍白、费解而顽强的记忆残片,耐心地探掘并破译——那是在我父亲去世、我母亲去世很久之后\fnote{31}{(*)}{我父亲1942年死于奥斯维辛;我母亲卒于1957年。我在这里谈的那项工作,是在1979年8月至1980年10月之间进行的。}

在这里铺陈我在这项漫长的工作中就这一情节的分量与意味所学到的、所懂得的,并非我的题旨。三天前,我已影射过这一转折\fnote{32}{(**)}{见“肌肉与肚肠(阳葬阴(1))”一注(n\textsuperscript{o}~106)的末尾。},把它作为我身上阴阳联姻的历史中三大时期的第一个的粗暴终结。1933年12月,我被人急匆匆地打发到一个陌生的家庭——无论我还是把我从柏林送来的母亲,都不曾见过他们。实际上,她把我送去的这些素不相识的人,不过是最先肯收我做“寄宿生”的人——收费微薄之至,而且没有任何形式的担保说这笔钱哪一天真能付;而母亲则准备尽快去巴黎,与正在那里望眼欲穿等她的父亲会合。\dmark{}\origpage{473}父母之间是说定了的:一切都将会很好——对在布兰肯塞(汉堡附近)的我如此,对我姐姐亦然:几个月前,她终究也被打发到柏林的一家残疾儿童机构(人家倒肯收她,尽管她并不比我或我们的父母更残疾)。

作为奇异六个月的结局——那六个月,满载着隐隐的威胁与焦虑——我一夜之间发现自己置身于一个全然不同的世界:与我的生活里唯一认识过的那个世界——由我父母、我姐姐与我构成的那个——全然不同。我置身其中,像一群寄宿生中的一员:我们与那家人分桌吃饭,在那家的孩子们眼里,俨然二等的孩子;他们自成一个世界,居高临下地看我们。在我待在那里的五年里(直到1939年——大战前夜,迫于时势,我才终于去投奔父母),母亲偶尔寄来一封匆忙而拘谨的信;父亲则从未有过一行亲笔。

收留我的那对夫妇很快喜欢上了我。丈夫是位辞去圣职的前牧师,靠一份微薄的年金与拉丁文、希腊文、数学的私人授课度日;妻子生气勃勃,有时透着狡黠——两位都不是寻常人物,在许多方面都令人依恋。丈夫是一位文化视野开阔的人文学者,在政治上有点迷失,与纳粹政权有过节;后者到头来放过了他。战后我重新与他们联络,并保持着持续的往来,直到两位先后谢世\fnote{33}{(*)}{她两年前去世,享年九十九岁;我还有机会看到她的遗容——在葬礼的前一天,与她单独相对。}。

从他、尤其从她那里,一如从我的父母那里,我领受了最好的,也领受了最坏的。今天,隔着长长的回顾,我为那份“最好”感谢他们(一如我感谢我的父母),也为那份“最坏”。正是这份最先由我父母、继由他们二位给予的最好与最坏,构成了我在童年分得的那个浩繁“包裹”的绝大部分(一如人人都分得自己的一份),那是一个属于我的包裹……——它等着我去开箱、去检视。它们属于我过去的实体、我过去的财富;滋养我的当下,全在我自己。

我的新环境,方方面面都“得体”之至、因循之至;尤其对一切涉及身体、涉及性的东西,都带着惯常的压抑态度。\dmark{}\origpage{474}然而,我想,过了好几年,我才把那些态度内化、据为己有——比如以裸身为耻,连同一种对身体的暧昧关系。这份自幼灌输的羞耻,是一道深层分裂的一个侧面:身体遭一种心照不宣的轻蔑,而所谓的“文化”价值(与记诵之类的智力能力混为一谈)却被捧上云端。我身上的这道分裂,一直无人知晓,直到我四十八岁,它才开始消解。那是我一生的第二个大转折:它标志着我与自己之关系的历史中“第三个时期”的降临——也就是说,我与我的身体、与我身上的“男人”与“女人”之关系的历史。但在此之前,我已有充分的机会把这道分裂传给我的孩子\fnote{34}{(*)}{至少,传给了五个孩子中由我参与抚养的四个。第五个也是最后一个,由他母亲抚养;迄今,连相识的机会都还不曾有过一个。},又眼看他们把它传下去……

昨天我已影射过\fnote{35}{(**)}{见前一条注“力量的绽放——或联姻”(注n\textsuperscript{o}~107)的开头。}我身上终于发生的那次“翻转”。在脱离最初的家庭环境(或不如说,在那个环境被\textbf{毁灭})之后两年多——这次翻转,标志着惯常的压抑机制的到位;我的童年,曾有那份罕见的运气,一直免于此。迄今我侦察到两大具有压抑本性的力量,它们主宰了我的成年生活、我童年的大部分(注108\textsuperscript{1})。我自信可以说:它们的出现并非渐进;在我的个案里,这些机制多少是一夜之间、以其全部力量出现的——作为无意识层面上一个蓄意\textbf{选择}的后果。我先前把这一选择称作“退位”;但与此同时,它也是一个强劲的行动原则:“我要像‘他们’那样”(而不是“像我那样”),也意味着:我要押“脑袋”的宝——我的脑袋反正不比任何人的差——用他们自己的武器作战,并打赢“他们”!

这些机制之一——也是我在这里最关心的——是最寻常不过的机制之一:那就是对我身上“女性”特质的\textbf{压制}(或被时下共识感知为女性特质的那些),以让位于“男性”价值。这枚硬币的正面,当然是倾其所有投注于那些被感知为“男性”的特质与才能,并把它们发展到失度的地步:它们占去了一片大得离谱的天地。

\dmark{}\origpage{475}倘若这里有什么超出常轨,那当然不是这一双重机制的单纯\textbf{在场},也不是(依我看)严格意义的那份“压抑”成分的力量——即对“阴”的特质、态度、冲动的压制之力。在这件事上,与我母亲身上的情形没有共同的尺度:她的一生(以及她亲近者的生活)被她对那使她成为女人的东西的憎恨(终生不为人知)摧残着。我相信,我的为人处事,从来都不曾全然没有某种温和,乃至温情;它们执拗地磨圆我自童年起为自己裁定的那个人物的棱角,常为我赢得同情与好感。出格之处,毋宁在于我投资的失度、在于我倾注于任务的能量的\textbf{失度}——目不斜视,不分心!在正经工作之外,我的精神不断地被抛向完成、抛向工作的这一站或那一站的抵达。这份态度(德语作“Zielgerichtetheit”,英语作“aimdirectedness”),是卓越地一种阳的态度:一种\textbf{绷紧}的态度,对一切不显得与任务直接相关者\textbf{封闭}的态度。这份失度,足以在他人身上唤起某种“超人”或“超级男性”的形象\origpage{476}——诚然可叹地令人叹服!(照时行的价值);却立刻(多半在无意识的层面上)激起本能的防御反应、乃至对抗反应:面对这样一份力量的铺展——它被感受为有威胁乃至有侵略性,或无论如何危险(注108\textsuperscript{2})。尤其,这形象无可抗拒地唤起“\textbf{超级父亲}”的形象,随即发动那环绕着永恒的父亲的冲突而纠结起来的、暧昧的吸引与排斥的万千反应……这正是在那些\textbf{暧昧}关系之中\textbf{我的}贡献;那些关系,在我的生活里如此常见,而我在《收获与播种》里又有那么多次与它们照面。这份暧昧,因我身上阴特质的持存而加深,而非减弱:那些特质喂养着一份同情——倘若只有阳特质增生成一个巨型“超人”,它是无从唤起这份同情的。

于是我再一次确认:在这些永恒的“暧昧关系”里,我仍旧不过是在收割我自己播下的种子——哪怕每一次收成都显得意外(而且不合时宜……)!因为推动我身上的“老板”在作品的积累中不断自我超越的动机(或至少\textbf{一项}动机),难道不恰恰是迫使并不断刷新我的同侪(首先)与非同侪(此外)的敬意?听我最优秀的几位哀叹:他们跟不上我向前奔跑的那个节奏?!是的,我身上确有这份隐秘的欲望:在他人身上(一如在我自己身上)唤起这幅“超乎自然”的、失度的形象——一如它所映照的那个人本身;而它执拗地经由他人回返:或以明白响亮的言词,经由那份预期的赞颂(并被当作应得之物照单全收)——\textbf{而且},经由阴沉的敌意与冲突的幽暗深层之路……\fnote{36}{(*)}{(10月6日)干脆说:我方才把手指按上去的那份“隐秘的欲望”,直到今天仍未烧尽——尽管它终于(不过几年之前……)被察觉,尽管它今天已不如当年那样吞噬一切。}

\notehead{108\textsuperscript{1}}{}
\phantomsection\label{ch18:local:n7}

\dmark{}(10月6日)我是说:在我的生活中起过作用的那些具有压抑本性的力量,似乎主要地(倘若不是排他地)采取这两种特定形态之一:埋葬过去,以及抬高我身上“男性”的特质、贬抑“女性”的特质。我绝无意说,这两股“力量”——它们各自都具有压抑的本性(即:指向一种“压抑”、一份对某一实在的藏掖)——就是“主宰了我的生活”的仅有力量!那会忘掉我存在中全部非自我本位的侧面:那股求知的冲动——它在身体的层面与精神的层面同样表达着自己。(对此尤其可见“我的诸热情”第35节。)

即便在结构着自我的那些力量之中——即“老板”的流出物之中——至少有一份,其本性并非压抑,而且远早于那些压抑的力量、在我的生活中的作用还更要本质:那就是对我父亲的认同;它宛如我自身力量之感“宁静而有力的心脏”。这份认同,丝毫不带有抬高某些价值或品质(比如男性的)、贬抑另一些(“女性的”)的意味。不问父亲所信奉的价值如何,他这个人(直到1933年——那年他身上发生了一次翻转\fnote{37}{(*)}{引人注目的是:我父亲身上这一“翻转”(他时年43岁)的方向,竟是朝着一个\textbf{超级-阴}的状态,朝着一种帕夏式的被动;与我母亲——扮演着超级阳的角色——密切共谋。她转而照管他,撇下自己的孩子们。(他们被弃置不顾、自生自灭,至少直到1939年——那一年,迫于时势,违背本意,她终于把我接到身边……)我父亲的这种依附关系、我父母之间阴阳角色的这场颠倒,一直持续到我父亲1942年消失。})始终带着一份强的阴阳平衡:直觉与自发,在其中占有的份额,不亚于智力与意志。

最后,作为另一股自我本性的重要“力量”——它与那些压抑机制血肉相连(或者不如说,它本身就具有“压抑”的本性)——还应当算上那份永恒的\textbf{虚荣};它的作用,在我的生活里,一如在任何人那里,同样沉重。但这份“力量”的本性是如此普遍——一如它在每个人的生活里(以或粗或精的形态)所扮演的支配角色——以致在一份关于自我结构之力与机制在个人身上所采取的特定形态、并赋予自我其特殊面貌与立足的清单里,几乎不值得专列一条。\origpage{477}

\notehead{108\textsuperscript{2}}{}
\phantomsection\label{ch18:local:n8}

\dmark{}(10月6日)在这份“力量的铺展”里,没有任何现行意义的“侵略”意图——无论有意识还是无意识;只有一份无意识的欲望:要令人印象深刻,要博得敬意。诚然,“博得敬意”这个自发来到我心的措辞,本身已带着一份\textbf{强制}的意味,近于“侵略”。这份无意识的强制意图——它同样在无意识的层面上被知觉——想必常常被经验为一种侵略(即便这份经验始终被遮蔽,一如它所引发的对抗反应)。与此同时,这份经验想必常与一些类似的、上溯到童年的经验混同——在那里,主角是父亲;而他显身为压抑性权威的主要持有者,甚至俨然一个压倒性的对手,被艳羡也被憎恨。

即便没有这样的混同,即便撇开他人对我身上任何“强制”意图的知觉,也想必常常有一份知觉:在这份(至少在其精神与意图上)纯然阳的“力量铺展”里,有一种强的\textbf{失衡}、一种根本的不和谐。这份失度,对当事者本人——即我自己——是有害的;而且说到底,对他的肉体生存都分明“危险”(近年来的几桩健康事故,已充分证明了这一点!)。想必正是这一点,潜在我思想的底层,当我写道:这样一份“力量的铺展”被感受为“无论如何危险”——“按本性”危险,是一个尤其不足为法的例子……!单是这样的感受,肯定就足以引发“防御反应”——即便没有任何侵略或侵略的意图。

诚然,这样的暧昧关系,1976年之后又重现过,与我的一些学生之间,在那些毫无数学投资的时刻——那时我的生活中没有任何明显的“力量铺展”。同样真确的是:那些\textbf{过去}的“铺展”造就了一份\textbf{名声},它至今贴着我的皮肤,尤其在我的职业生活里;而且在某种程度上,它取代了对\textbf{当下}的我的知觉。更不必说:我在操持某些数学主题时习得的从容如此之甚,以致即便在我的数学时期之外、加上名声的推波助澜,这份从容或天然的驾驭,对一些动力不足的学生,也可以已经有“力量铺展”的效果,使他们尽管看见我身上也有一些亲切乃至令人安心的特质,仍把我感受为某种超人(边缘上还有点超级父亲!)\dmark{}\origpage{478}。

况且,作为我所说那份从容的背面,我常常倾向于低估:对某个学生而言,取得某份行装、发展某件工具,会有怎样的困难——这往往把他摆进一个对我的期望的尴尬境地。(对此可参见“教学的失败(1)”一注,n\textsuperscript{o}~23\textsuperscript{iv}。)这样的局面,想必常常是一项与父亲之间的虚伪关系的重要配料之一……

\subsubsection*{(c) 重逢(阴的苏醒(1))}
\tocsubsub{(c) 重逢(阴的苏醒(1))}
\label{ch18:local:h9}

\notehead{109}{}
\phantomsection\label{ch18:local:n9}

(10月9日)四天前,在收束前一条注\fnote{38}{(*)}{见“阳葬阴——或超级父亲”一注,n\textsuperscript{o}~108.}之时,我感到满心欢喜。我不期然地重新接上了一份直觉——它曾于某个星期日,1976年10月17日来到我心头(距今将近八年):即我身上某股“特定的力量”在我的生活中、一如在我母亲的生活中的毁灭性效果的直觉。那是我一生中头一回,为我的平生——尤其我的童年——作一场反思,哪怕极其简略。那也是我发现禅观力量两天之后的日子\fnote{39}{(**)}{见“欲与禅观”一节,n\textsuperscript{o}~39.};而且,自那一刻以来,这是我头一回动用这份久被忽视的力量。那一天,反思径自转向我的童年,并无蓄意,而是出于一股深层的冲动,仿佛被一个非常可靠的本能所驱动。只是带着事后的认识,我才掂量出:当时把我推向的,恰是我真正的力量之源,同时也是我身上冲突与分裂之源;一份深层的知的渴望推着我走向它。将近三年,我没有再回到那里——那些年,我被一些“日程事务”的问题独占了心思,竟不觉得自己一直逗留在我生活中那个冲突的周边,却执拗地远离其心脏:那段淹没在迷雾里的童年;它显得无限遥远……

我刚刚“斜着”重新翻阅了那场我生命中至关重要的禅观的那十八页——密度超常的十八页。正是在那场禅观之后的当夜,或不如说,在那禅观之夜过后的凌晨,我做了一个震撼人心的梦——那也是我一生中头一个被我从头到尾热烈破译其讯息的梦。当时我对自己的去向、对正在发生的事,同两天前“发现禅观”之时一样,毫无所知。\dmark{}\origpage{479}整整四个小时,我沉入那份经验、那个寓言式之梦的意味之中,穿过一层层越来越灼热的含义,才抵达讯息的心脏、其简单而分明的意味。

那不是“智力”之领悟的骤然咔嗒一声,甚至不是暗处或半明处的一线突然之光。那宁可像一股在我体内诞生的深深的浪,蓦地贯穿我全身;在它浩瀚的水面上,给我送来那份此前一直躲避着我的意味:我在这一刻重新找回了一个至亲至贵的存在——一个自我童年起便失落了的存在……

这一刻被经验为一次\textbf{诞生},一场深刻的更新。这份感觉,那一整天、其后的几天里,都十分强烈。带着八年的回顾,这一刻在我今天眼里,依然是我一生中头等的创造时刻,是我的精神历程中一个本质转折的时刻。它当然是由此前几天、几个月里的许多别的“时刻”所预备的。也许头一个先兆,是十多年之前的那次“有益的挣脱”——脱离那个我本打算终老其中的机构\fnote{40}{(*)}{见同名注,n\textsuperscript{o}~42.}。这些在先的时刻,在我眼前宛如那些配料,或不如说,供我调遣的\textbf{手段}:凭它们,我才得以越过眼前那道我未曾察觉的门槛——它安放在比我已经越过的其他门槛更深、更隐蔽的层面。几天或几个小时以来,一切都已齐聚,只待我越过它;而我尽可以越过它,也尽可以不越过——日复一日,一生如此……

而一旦这道门槛真的越过,路便向进一步的跨越敞开,向别的“觉醒”或“苏醒”敞开——其中每一个,按其本性,都是一次更新,而且多少是一次“新生”、一次再生。有过这样的情形:我把其中某一道回避了几个月乃至几年,才终于跨出那一步,一路卸下某种顽强的幻象——它一生拦在我与我生活、与我周围世界的全部滋味之间。肯定,也有一些,直到我写下这些文字的时刻,我仍在回避……

\dmark{}\origpage{480}在这几天的反思的眼光里,正是这个与我的童年重逢的时刻——人们曾以为它失落、死灭了漫长的一生——标志着我精神行程“第二个时期”的终结:那是我个人生活中\textbf{自我本位机制}占支配地位的时期,对着创造的力量、认识与更新的力量;那些力量,曾经历近四十年的近乎完全的停滞。那也是一个“特定的力量”占上风的时代:一股几乎纯然“男性”性格的力量,映照着周围世界里所尊崇的价值;以牺牲我存在中“女性”的深层侧面与力量为代价——它们遭忽视、遭压抑(谢天谢地,从未全然得逞!)。

对这份力量之毁灭本性的头一份直觉——它曾主宰我的生活,一如主宰我母亲的生活、也主宰在我生命中扮演过重要角色的其他女人的生活——这份直觉,在这些强烈成熟的日子里短暂现身;那无疑是借着“阴”的、“女性”的能量在我对事物的意识把握中重新涌出之便。与我匆匆记起的相反,这份现身,不发生在重逢前夜的那场禅观里,而在重逢之后几个小时——在一场对刚刚发生之事之意味的简短禅观里。这份直觉,在这场禅观那几页笔记的末尾诞生并成形。我知觉到这份“力量”(今天我会称之为“超级阳的力量”,即阳占压倒主导的力量)的毁灭本性——先在我母亲身上,继而在其他女人身上;然后接上这几行收尾的文字:

\begin{quote}
“至于我自己身上的‘力量’,正是它,使我成了一个久待的靶子与对象——在我年轻的一生里,等着M.、然后J.、然后S.~的隐秘的恨与怨——那份恨,早在她们认识我的存在之前许久,就已存进她们心里,在那失去慈爱的童年的张皇岁月里。”
\end{quote}

“童年”一词——在见证我一生中一个头等重要的日子的这最后一行里——也在此后近三年间,最后一次出现!至于那份关于我身上超级阳的力量——作为对抗反应乃至恨与怨的挑动者——之本性的直觉,它(依我看)往往沉入遗忘,直到最近这些天。更确切地说:它只在我对生活中某些重要关系的知觉里(尤其与我爱过的女人们的关系)保持在场。相反,它却不曾真正深入那些多少“不速而至”\fnote{41}{(*)}{或被如此对待……}\origpage{481}的冲突局面——尤其与某些学生的冲突;我在《收获与播种》里已有多次机会检视或提及。在整个这段反思里,这一个事实——即经由一种并非自愿的“挑衅”,我自己也给我此处彼处提及或检视的那些冲突局面,贡献了我自己的一份——常常全然被遮蔽;而对方的贡献,却显得一清二楚。这当然是一个最普遍的反射,不妨说是一个普遍的反射!最近这些天的反思,终于把它拆了引信,同时又使我重新在我自己身上侦察到它——使我突然,在道路的转弯处(在一场关于阴与阳的反思里……)与我自己的鼻子撞了个正着——与一个\textbf{特定的}我自己,至少。

四天前那场简短的反思,也只不过勉强触及我这个人身上的诸多侧面——那“人物”(自我童年起我扮演的角色)中阳失衡借以显形的诸多侧面;以及这份失衡有时对他人产生的碾压效果。尤其对那些阳型力量尚未扎根的人——首先对我自己的孩子。我在这里尤其想到某种“武断自信”的做派:在一切(而且为数众多)我或对或错地有自己一套看法、感受或定见的事情上,我都是照着它运转的。诚然,我从未想过把我的看法强加给任何人,对我的孩子更甚于对任何别人——而且,仗着我身上(至少在有意识的层面)没有任何强制意图,我大半生都无力觉察:然而我的这些为人处事之道(它们在我看来是自发的、自然的;我远没有辨明其复杂的本性……)——它们对我的孩子、对别人,竟有着与强制同样的效果;或者不如说,一个更阴险的效果:在对方心中激起或滋养一份\textbf{不安}——对其自身的感受、看法、意见之价值的不安——仿佛它们(面对我那份毫无破绽的自信、乃至我那份痛心的惊讶)竟连\textbf{存在}的资格都没有。

况且我预感到:我身上这一癖好的发展——尤其在与我的孩子们的关系里——可能相当复杂,与我婚姻生活的坎坷血肉交织。这里不是追索其奥秘的地方;也不是对我这个人身上显出这份失衡的其余侧面试作一份或多或少\dmark{}\origpage{482}完整清单的地方——前一条注里,我试着勾勒了其一个格外显眼的侧面:那份“力量的铺展”。

不应当以为:这份被培养了终生的失衡,连同它借以显现的众多心理机制,会像魔杖一挥,一夜之间烟消云散。我本来也不作此想——无论在重逢的那一天,还是在其后的日子与星期里。

(10月10日)那是冰河解冻的日子,由一股新能量的强劲涌入所托举——是内心工作与惊叹的日子:面对这些崭新的世界,它们日复一日在我眼前半启,诞生于日常琐事的寒微织体之中,在一双渴望看见的、强烈的眼睛的作用下铺展。那也是我开始第一次预感到一个未知存在之丰富的日子:它忽然向我发问,而就在前一天我还对它浑然不觉。我从刚刚向我显露的那些“片段”去把握它——它们在重逢的当下、在随之而来的那段无法预料的旅程中现身。我分明感到:我刚经历的这场“诞生”,只是某件全然未知之事的开端;或不如说,是某件曾被中断、被割断或窒息、如今神秘地重新发动之事的重启。说来,这场强烈的“生成”,在此前几个月便已重新运转,只是还处在内省的思考几乎未曾参与的层面……

这场重新活过来的生成、这项重新开始的工作,其深层的一面,是我身上“女人”与“男人”、阴与阳之间原初平衡的逐步恢复——贯穿其后的日日、周周、年年。在某种意义上,我可以说:自重逢之时起,“童年”或童年状态便一直潜在地在场;我对自己的本性、对自己本质上不可摧毁的统一,有一份深刻而不可磨灭的认识,超越了某种“分裂”的效果——它还常常扰动我存在的表面。用“孩子”或“童年”这个词来称呼这种存在的统一,\dmark{}\origpage{483}倒是好几年后才有的事,约在我开始于有意识的思考层面认识万物的阴阳双面之时。也正是在那时,我开始知道(或至少预感到):童年状态、创造状态,乃是阴阳力量与能量完美平衡的状态,是阴阳“联姻”的状态,表现为一种创造的和谐。

我觉得,在某个层面上,这种对我根本统一的认识,每时每刻都在场,也每时每刻都在起作用。诚然,它的作用随时刻而或显或隐、或强或弱;而它绝不意味着比较持久地消除、甚至整块摧毁那些自我本位的力量,即“老板”——甚至不意味着消除那些压抑的力量(它们构成“自我”的很大一部分,倘若还不是其全部……)。这些力量,暗中偷换环绕我的现实与在我身上展开的现实;它们默默而执拗地运作,顶风冒雨地维持种种顽固的幻象——没有它们支撑,那些幻象便会立即因自身的重量而坍塌……其中一些压抑机制被我逐一认出,随后消失。我摆脱了一些沉重压在我身上的幻象,也澄清了几个顽固的疑点:它们曾终生被扔进地下的垃圾窖(受“老板”照料……),从未受过检视。它们的讯息终于被听见之后,这些疑点便消失,留下一份宁静而喜悦的认识。我也认出一些威力极大的压抑机制,深植于自我之中;近年来我意识到:它们在我生活中的分量,今天一如往昔,依然可观。它们朝着阳失衡的方向,朝着某些阴的力量与官能被遮蔽的方向。我不知道这些机制哪一天会被拆除引信——而我知道,这全在我自己。无疑,它们将在我对生活中那场冲突的根源,进入得比我迄今所做的更深、更彻底的那一天——也唯独那一天——烟消云散。

眼下,鉴于我的生活正朝着一项重要的数学投资转向,我尽可以说:这事毫无进展的迹象!

\subsubsection*{(d) 接纳(阴的苏醒(2))}
\tocsubsub{(d) 接纳(阴的苏醒(2))}
\label{ch18:local:h10}

\notehead{110}{}
\phantomsection\label{ch18:local:n10}

\dmark{}(10月11日)一两天以来,我一直想用几句话\origpage{484}盘点一下:八年之后,我身上这场“阴阳平衡的逐步修复”进展到了哪里。

也许最重要的变化,在于对我这个人——如其每时每刻真实所是——的一份比过去\textbf{大得多}的\textbf{接纳}。换个说法:我身上的压抑机制已大为松弛。如我昨天所说,一些机制在被发现并理解之后消失了;另一些——我一生不曾察觉的——则在其日常显现中变得为我所熟悉。我看着它们运作,不是当作必须不惜代价拔除的敌人,而是当作我受制约存在的诸多侧面之一部分,由此,也是当下“给定”之丰富的一部分——那份“给定”,忠实地映照着我的过往:既有我的种种制约、我存在中分裂之根源的那段“古老的”历史,也有我成熟的那段较近的历史——即那项工作:经由它,我终于开箱取出、并“吃下”消化了父母及其继任者传给我的最初包裹。因此,我身上的这份“接纳”,不仅包括“孩子”的那些我长期忽视并压抑的冲动与特质(尤其那些映出我身上女性侧面的),也包括“老板”特有的那些压抑机制——也就是说,恰恰是那些积习已深的“不接纳”机制!接纳后者,与“栽培”它们或强化它们,毫无共同之处。相反,它是一步不可或缺的第一步:经由一份好奇而慈爱的注意的效果,把它们多少解开或拆除引信。这八年的经验使我深信:只要这份注意沉得够深、直达压抑的根,压抑便会消解、消失,并释放出可观的能量——那份直到那时还被冻结起来、用以顶风冒雨维持某一套压抑机制、以及用以维持它们的思维习惯与其他习惯的能量。

不过,这份对我的新接纳,首先并不显现于对我这个人那些本性上“纠结”的侧面。它来时不声不响,早在禅观发现之前——因此也早在紧随其后的那场“重逢”之前。那是1976年7月,在与一位名叫G.~的年轻女子的一段短暂恋情中——她为人处事,或许比我先前爱过的女人们略多几分“男子气”。偶然(? )安排的情势是:环绕这段恋爱的物质条件,竟把我摆进一个典型的“女性”角色。我料理家务、\dmark{}\origpage{485}准备晚饭,等那位伴侣结束一天漫长而劳累的工作归来:她得在山里看管一百五十只山羊,晚上还得挤奶。这个不寻常的持家人妻的角色,竟像手套一样合我。事情看似微不足道——然而当时就“咔嗒”了一声。这一联,在我心中系向了我恋爱生活中的某些冲动与欲望——它们当时并头一次在几首情诗里表达出来;诗中,恋爱的经验毫不含糊地显为“女性”的。我那时懂得了——不经反思、不经“努力”,没有丝毫迟疑或扭捏的念头:在我的身体里一如在我的欲望里,在我的感情里一如在我的精神里,我是女人,同时我也是男人——而我存在中这两个深的实在之间,没有任何冲突。在那几天里,主音是女性的;我怀着感激、在默默的惊讶中,接纳了这件事。想到它时,我心中有一种宁静的、极温柔的喜悦。

这份喜悦自足自立,无须形诸言词——无论对我自己,还是对他人。我不知道是否对那位以我为男情人(或许是女情人……)的女子谈起过。肯定,在某个层面上,她知道,无须我说。

这份喜悦没有走味;它活着,直到今天。它出自一份活的认识,如香气伴随着花。在我生活的某些时刻或时期,这份认识——以及作为其标志的这份喜悦——比别的时刻更在场,作用更强。但我不相信它有一天会弃我而去。

在其后的星期与年月里,当我此处彼处有机会谈到这一经验与这份认识时,每一次都像是把一件贵重的珍宝传给别人——在我感到他敞开着、哪怕只消几分钟肯领受我身上这份喜悦的些许之时。我从不曾感到一份使我难于启齿的扭捏——仿佛这是一桩多少有点有伤风化的事。(倘若我身上“男人”的实在与力量不是高悬于一切怀疑之上,或许有时倒会有这样的扭捏!)我还记得有一回,我分明是在招摇过市,极力炫耀自己一箭双雕——只差像人人一样来例假、并且立马生个娃了。

我这份新的女性身份,叠加在我的男性身份之上,给我的恋爱生活带来一种立竿见影的更新。它在我此后做过情人的那些女人身上,激起一份\dmark{}\origpage{486}极强的回响:唤醒了恋人身上的男性冲动——那些冲动终生被小心压抑,此前只得以“偷偷摸摸”的方式表达,如某种毛边,不配进入有意识的恋爱经验。

无意识的恋爱经验富含原型冲动,其中最强有力的一种,是对母亲的回归、对原始怀抱的回归。这样的原型,在恋爱经验的深层,男人身上有,女人身上也有。在女人身上,对在两人的恋爱经验中满足这样一份冲动的抵抗,比男人更强:在男人那里,它撞上一道关键的禁忌;在她那里,撞上两道。无论对他还是对她,这些冲动在共同经验中的满足,往往多半停留在象征层面,尤其,避着意识。当这样的原型与这份经验从深层浮上意识之光的表层、进入有意识之眼的视野,这份经验旋即转变,获得一个新的维度。同时,可观能量的解放随之而来——它们先前被压抑机制压缩着,或被压抑的任务束缚着。其效果,是\textbf{性爱}冲动的一份即刻的\textbf{解放}:表现为焕然一新的强度,与恋爱经验中一种新的丰满。

由以上所述,想必已经显明:对我自身的这份新接纳,是与对他人的接纳并行的。两者密不可分。当然,这里说的是完整意义上的“接纳”:它绝不意味着对这样那样被感受为(唉)不可避免的“毛病”或“缺点”的一份(往往是酸溜溜甜丝丝的)\textbf{容忍}——对那种人们只好“凑合着过”的不幸。在那样一种态度里,我首先感到的是一份认命,甚至一份退位;肯定不是喜悦之源,也不是一种认知的冲动——认知那值得认识的东西:在人们肯于容忍的那些“缺点”或“毛病”的平板表面背后,那份隐约感到的、未知的深处……

这里说的是一份喜悦的、创造性的接纳,但这绝不意味着这份接纳是全然的——昨天我已确认,事情并非如此。留心的读者,在《收获与播种》里想必已不止一次亲自确认过这一点;一如我此处彼处顺带地有所察觉:当我发现自己又一次面对我身上那个永恒的\textbf{拒斥}机制——对着任何以不讨人喜欢的外观现身者,无论在他人身上还是在我自己身上\dmark{}\origpage{487}。(不过当事关自己时,这机制最常有的效果,是压根不去知觉那件不讨人喜欢的事……)

我所说的接纳,扎根于一份对被“接纳”之物的\textbf{兴趣}——无论在自己身上还是在他人身上。接纳本身,是一种典型“阴”性的内在心性;而它在我身上带的这层“兴趣”的意味,则是“阳”性的——在阴阳无限交织的那份精微的中国辩证法里,这是“阴中之阳”……我差点顺势说:接纳(真正的接纳!)与这份兴趣、这份好奇,是纯粹的同一。然而稍微端详一下这件事,我发现确也还有另一种接纳方式,其本性比我所惯常的更全然是“阴”的。那像是对被接纳之物的\textbf{迎纳},而非一种向它扑去加以探测的冲动。(这份迎纳的细微差别,顿时间在我眼前显为“\textbf{阴}中之阴”——瞧,来了!)兴趣的冲动与迎纳的态度,都可以构成对他人或对自己的接纳的基调。两者共有的是\textbf{同情}。这也是爱的一种形态。倘若这里有什么深层的东西可供把握,那就是这份确认:\textbf{接纳是爱的一种形态}。对自己的爱,对他人的爱;两者密不可分……

除寥寥的瞬间之外,当事关我自己的时候,我的兴趣比事关他人时更为强烈地投入。正是这份对我自己的炽热兴趣,鼓舞着这八年来的那些长程禅观。诚然,自我认识正是认识他人与世界的心脏,而不是相反;我感到,我的新热情——禅观——从前到现在都把我引向事物的心脏、最本质之处。这些年,我对他人的兴趣,连同由此而来的接纳,显现得更零碎,也更迟疑。它的一种具体表现是:与人相处时,我不再那么爱说话,而愿意倾听。我一生的大部分时间几乎全无这种倾听的能力。即使在重逢的那个大转折之后,我也屡屡发现自己因为没有倾听、缺乏辨别,\dmark{}\origpage{488}在不合时宜的时候说了话;过了好一阵,这积习才开始退去。

后来它不再那样泛滥,甚至几乎消失,丝毫不是我给自己定下了什么规矩的结果(譬如:“只有在……时你才许开口”)。只是因为,在我明知说了也无益、对别人对自己都带不来什么——至少带不来我眼里有价值的东西——的时刻,我失去了开口的兴致。倘若如今我常能感觉到这些,想必是因为我变得更留心了。这也不是纪律的成果(“你可得竖起耳朵仔细听……”);它究竟怎样来的,我说不上。无论如何,我因此感觉更好,生活也更有趣(尤其更安静了!)。别人也因此感觉更好……

我想,我开始真正少说话,是在我身上那股总想纠正别人身上所谓“错误”(不论我判断得对错)的力量几乎消失之时——仿佛察觉并更正我自己的错误还不够似的!这也是那份推着我(有时至今仍推着我)想要不惜一切代价说服他人相信这个或那个的力量,而不是简单看一看:为什么某人执拗地宁肯信这个而不信那个(那个在我看来是“对的”,我很想说服他!)——或者,为什么我如此看重他信那个而不是这个。这份我们身上近乎普遍的力量——它不断推着我们到他人的赞许中(哪怕只是一个人的……)寻求我们所信为真者之正当性的确认——这份深深扎根于自我的力量,我相信,终于在我身上松手了。那是一大解脱,是一场惊人的能量涣散的终结。正是两年前,当我终于意识到这份力量在我生活中的分量、其本性、以及它所代表的那份惊人的能量涣散时,它被拆除了引信——我顿时如释“千钧重负”。毫无迟疑地面对他人向我们的人格回掷的回声,而不被一份对赞许或确认的欲望或“需要”(无论藏得多深)所缚——这才是真正地“摆脱了他”。正是这样一种需要或欲望,构成了那个真正的“挂钩”——不起眼却坚固无比——冲突由此得以在我们身上“挂住”;我们由此(无论我们愿意与否、承认与否)处于他人的依附之下、处于其\dmark{}\origpage{489}\textbf{随心所欲}之下——总之,由此他“拿住”我们,并且(不动声色地)随意摆布我们……

照正常的道理,对他人的接纳,理当也蕴含着对其看待事物方式的接纳——无论这些方式在我们看来对错,甚至当事关他看待我们这个宝贵自身的方式时(包括我们自己的看法……)。然而恰恰在这里最疼——那正是接纳他人的命门,而绝不在于接纳那些不直接关涉我们自身的、多少恼人的寻常“缺点”。况且常常,我们拒斥他人的某些“缺点”,主要是因为经由它们,我们直接感到自己被问责——单是面对一些(又是对错难辨)似乎与我们相反的为人之道。换言之,是我们身上的一份\textbf{不安}——它以虚荣的(或多或少显性或隐性的)反应显形——才是妨碍我们接纳他人的大障碍。而这份扎根极深、由虚荣的动作所补偿的不安,在我看来,与对我们自己的不接纳密不可分;它宛如其形影不离的影子。

于是,对自己的完整接纳,在此显为把我们向对他人的接纳敞开的那把钥匙。而这道刚刚向我显现的联系,又接上另一道我久已认识、或许始终认识的深刻联系:对自己的爱,是对他人之爱的宁静而有力的心脏。

\subsection{(3) 伴侣}
\label{ch18:local:h11}

\subsubsection*{(a) 事物的动力学(阴阳的和谐)}
\tocsubsub{(a) 事物的动力学(阴阳的和谐)}
\label{ch18:local:h12}

\notehead{111}{}
\phantomsection\label{ch18:local:n11}

(10月13日)昨天我没有继续写注。取而代之,我兴致勃勃地把一批“阴阳成对”的东西巡视了一遍。从脑子里随手冒出来的那些开始,我渐渐被咬住了钩,最后以一场力所能及的“普查”收场:凡是我能抓到手的,统统入册。我起头,是因为我盘算过:我近来写下的东西,有相当一部分,对一个对事物的阴阳双面尚不多少熟悉的读者来说,恐怕会整个“从头顶掠过”。费心给出若干醒目的这类成对的例子,以外加最近这些天里顺便溜进来的那些,或许不无用处。\dmark{}\origpage{490}接着,经不住我身上那位系统学的小魔鬼(或小天使,说不好……)的撺掇,我最终翻出了五年前在这个主题上的旧反思。当时我曾饶有兴味地“搜集”了一两百对十分有启发性的对子;它们按亲和性聚成约二十个组。由于那场反思是在我正在写的那部大名鼎鼎的“诗作”的边缘进行的,我忍不住把这些组按亲和与意义的承袭,鱼贯排开。昨晚,带着回顾重拾这场反思,脖子上不再套着诗的轭具,我找到了十八个组(而不是二十个),分组或许更严格了些。况且我疑心,应当还有许多别的组,甚至无数个组,对应于一些我在那项工作里(或许迄今为止都)不曾想过的把握实在的方式。至于我确实辑得的这十八个组,我尽力把它们装配成一幅图(或“图示”),依着联结它们的主要亲和纽带。其中有些纽带,还是在一遍遍画图草稿的过程中,才强加于我的注意。这份工作,与那十分熟悉的数学工作实在相近:当我们力求以尽可能醒目的方式,把一组或多或少复杂的关系(比如由“映射”给出的关系,以箭头表示)——存在于若干“集合”或“范畴”之间,它们充当所构“图示”的“顶点”——按图形方式把握起来。在那里,同样是一些本质上属于审美本性的要求——尤其对称与结构透明的要求——常常引导人们引入(因而在必要时发现乃至发明)一些起初不曾想到的“箭头”或纽带,有时甚至一些新的“顶点”。总之,经过五六次相继的草稿,我最终得到一幅图,大致呈圣诞树的形状;它暂时令我满意——更何况夜确实已深得犯规!

我心满意足地睡下;我感到自己没有虚掷光阴,即便那些注一个字也没有推进\fnote{42}{(*)}{作为补偿,我满可以为一项发明申请专利:一种新的诗体,即所谓“非线性”或“图表式”的诗。}。但我重新与一些分明多汁的事物接上了头\dmark{}\origpage{491}——这些组中的每一个都富于分量与神秘;而构成它们的(不如说,全体一道\textbf{指涉}着它,却丝毫无意穷尽它的)每一对阴阳,都有些精微而重要的东西要对我说:关于我所栖居的这个世界的本性,常常也关于我自己的本性。我以新的力量,重新找回五年前便已在场的那份感觉:阴阳——万物之中“女性”与“男性”——的精微游戏,是一条无可比拟的引线,通往对世界与自我的理解。它把我们径直引向那些本质的问题。常常,阴阳的“瑜伽”本身——我是说,单单留意事物与事件的那一个以阴阳平衡与失衡来表达的侧面——就为更好地理解这些问题、走向一个回答,供给了第一把钥匙。

我该向某些读者道个歉:过去一两页,我给人的印象也许是在谈“天使的性别”;他们可能连我说的那些阴阳“对子”究竟是什么都还不清楚,更不用说由其中若干对子组成的“组”,以及这些组最终又被装配成一幅“图”(数学到底还是有用的!)。这里我至少该举出一组。我想选昨天自发地从它开始的那一组:在反思中,它最终也显为“原始”组;其余各组似乎都通过一层层“衍生”从它生出(在我的那幅图上延续了八“代”……)。以下便是我列出的构成这一原始组的对子(可用其中第一对,称其为“行动——不行动”组):

\begin{itemize}
\item 行动——不行动
\item 活动——被动
\item 清醒——睡眠
\item 主体——客体
\item 生育——受孕\fnote{43}{(*)}{(11月6日)实际上,还有一个更原始的组,可以称为“\textbf{父——母}”组。关于这一“疏漏”,见“我们的母亲死亡——或行动与禁忌”一注(n\textsuperscript{o}~113)。“生育——受孕”与“执行——构想”这两对——我把它们列在了上面那个(号称“原始的”)“行动——不行动”组里——分明以更自然的方式,插入那个围绕“父——母”对子形成的“母组”之中。}
\item 执行——构想\textsuperscript{43(*)}
\item \dmark{}动力——平衡
\item 冲劲——立足\origpage{492}
\item 热切——坚忍
\item 猛进——耐心
\item 激情——宁静
\item 执着——超脱。
\end{itemize}

还有几组对子,是昨天反思的势头在今天早上又送来的十来个“迟到者”里的两个,我愿意把它们添在下面:

\begin{itemize}
\item 知——识
\item 解释——领会。
\end{itemize}

无须说破:在这些对子里,放在前头的,是“阳”的或“男性”的那个词——照我们父权社会的惯例,由男人给这对子起名。相反,中国传统社会比我们的社会父权得多,可当人们循中国的惯例谈论阴阳的关系时,却总把“阴”(“女性”)放在前头——比如说“阴阳平衡”(而非“阳阴平衡”)。这一惯例的意味,肯定在于那份原型直觉:阳生于阴;阴是两者中“最原始”的本原,而不是相反……

这里不是对这一对或那一对子加以评说的地方。对那看了它们却“毫无感觉”的读者,反正也是白费劲;而对那个感到被它们唤起的人——他感到(哪怕模糊地)它们每一个都有些话要对他说,关于世界、关于他自己,关于平衡与失衡、关于存在与事物的内在动力学——这个人无须周详的评注,尽可以把这份唤起当作他自己反思的起点。

\subsubsection*{(b) 敌对的配偶(阳葬阴(3))}
\tocsubsub{(b) 敌对的配偶(阳葬阴(3))}
\label{ch18:local:h13}

\notehead{111$'$}{}
\phantomsection\label{ch18:local:n12}

这里只有一点我想强调——它毫无例外地为一切“阴阳”对子所共有。我觉得,这也是为了理解阴与阳之间关系的本性、并由此理解这两个本原(或能量,或侧面,或力量……)在宇宙中之本性,最为要紧的一点。这就是:这些对子中任何一个的两端——比如行动——不行动,在\textbf{缺了另一端}的情况下\dmark{}\fnote{44}{(*)}{(10月16日)实际上,这份“缺席”从来不是全然的,依我看——在任何事物里,阴或阳都不曾以纯粹状态在场,而无其补充者的同时在场,无论多么微弱。所以我说的“失衡”,其特征不是两个补充端之一的全然缺席(那是从未实现的事),而是那一端的\textbf{过弱}状态。另一类失衡、或病态,则出现在\textbf{两}端都“缺席”之时——或更确切地说,都在场却都极弱。比如在“行动——不行动”这对子的情形,一种\textbf{躁动}状态:它并不严格地“行动”(除非为了自我延续、维持一场混乱),却四散着能量——想必可以算作这样一种“亏空型”(阴与阳两亏)的失衡。}\origpage{493},便构成一种严重的失衡状态;而(当所说的“缺席”近乎完全且持续时)在极限处,是一种导向毁灭的状态——毁灭发生这种失衡的事物(或存在),甚至连同他与他的周围。

比如,一种不间断的\textbf{行动}状态,不与充分的\textbf{不行动}期、休息期交替,便导向衰竭、疾病与(极限处的)死亡——这恰恰是近来于我格外切身的事\fnote{45}{(**)}{对此可参见第XI行列的前两条注(n\textsuperscript{o}~98、99),“亡故者(仍未死去……)”。}!但反过来,过度的无行动状态,又导向身体或心灵(视情形)的官能与机能的衰弱与硬化,极限处导向毁灭。在我那桩“患病事故”里,我恰好有一个\textbf{两}种失衡并存的例子:精神的行动过度,身体的无行动(外加对两者都没有充分的休息……)。

就这个个案而言,这份对阴阳之平衡-失衡“哲学”的“解释”,还停留在表面:因为它不曾触及一份积习已深的文化成见——它抬高阳的一端、行动,却\textbf{对立地}看待阴的一端、不行动。不行动被感受为一种“消极”的东西:不出产,从任何角度看都乏味;勉强被容许,作为一种无奈的权宜——它(唉)甚至强加于世上最善良的意愿,因为人总得时不时歇一歇,才好继续投入行动(否则,正如我方才解释的,就有操劳过度与天知道什么的风险……)。总而言之:不行动被看作行动的寒微婢女——(唉)不可或缺,除此之外不配注意、不配敬重。

\dmark{}\origpage{494}当然,这样一种对行动的“官方”抬举、以牺牲不行动为代价,随即有一份后果:在人身上发动抵抗机制(它们常常保持隐秘,或至少极度混乱),其表达是一种\textbf{相反}的价值化:行动,顿时显为生存的严酷必要所强加的东西,总之是那份“苦工”——在办公室、在工厂乃至在田间,烦得不能再烦,即便不那么烦也累得半死。行动的真正存在理由,是挣面包与住处(这是不可或缺的),而在此之上尤其,是要有美滋滋的闲暇(在有生之年),以及日后一份体面的退休、一份惬意而永久的休闲——那时人们将免除“苦工”这份可惜的义务。这一回,被(或多或少有意识地)抬举的是不行动(别名“闲暇”),而行动成了它的寒微婢女。于是有了一场角色的\textbf{颠倒},却始终是同一份失衡:即当事人(在文化制约的推动下)在自己生活的两个本质侧面或两极之间设立的那份\textbf{敌对};这敌对,经由其一面的专制支配与另一面的奴役状态,表达并延续着自己。

我觉得,最常见的是:这两种态度与价值化在同一个人身上叠加,一种在意识的层面占着街面,另一种在无意识的层面。而这两种相反失衡的叠加,分明并不产生平衡!平衡倒是自然出自对行动与不行动之真正本性的领会(即便这样一种领会停留于纯然“本能”,直接经由一种平衡的行为显形,而非经由一份形诸言语的“所知”)。\textbf{在完整意义上的行动之中,也有不行动}——我是说,它就在\textbf{行动的当下}在场,而不只是“事后”(因为事后总得歇一歇嘛!)。行动之中的这份“不行动”、这个“阳中之阴”,宛如一种深深的宁静,充当着发生于表层的一个运动的立足。比如,它显形于一只行进中的猫科动物所散发的那份完美松弛的印象——无论是随便哪只街猫,还是一头体格雄壮的母狮……而同样,\textbf{在真正的无行动之中},哪怕它全然彻底,\textbf{也有行动}。比如,\textbf{睡眠}富于它的梦:那些梦向我们诉说着我们自己;经由它们,我们过着另一种\textbf{生活}——一种更强烈、更精微的生活;对它,我们常常太沉睡、太怯懦,以至于无从在醒着的生活里去活。而且只需端详一个熟睡的婴儿,或仅仅从酣睡中被唤醒,就会感到:即便没有梦,真正的酣睡本身也是一种\textbf{工作}——以它自己的方式:某种全然吸收我们的事物,说到底,是把一份已然四散的能量“重新充满”——我们刚从其源头把它重新汲回……这又是那个“阴中之阳”\origpage{495};没有它,阴本身也会成为毁灭性的。

\dmark{}同样的方向的反思,想必也可以为醒着的不行动——睡眠时间之外的——发展出来。只消对人们知觉为“无行动”的这样或那样的状态,就实际情形作一番留心观察即可。人们会发觉:无行动之中有行动——哪怕是一个思想继续原地打转的贫瘠饶舌,尽管它已经停止了工作。但说实话,把那个纯然机械的运动称作“行动”是不恰当的:它仅凭惯性而延续——靠的是那台机器停不下来!而这份内在的躁动,当然不会给“无行动”带来一份使之有益的阴阳和谐。相反,倒可以是这样的:那些用来充实闲暇的种种活动(即便人们把闲暇体验为一种不行动)。甚至在完全休息的状态中——比如疗养期——也可以有行动;不然,这份休息或“不行动”就会沦为瘫软懒散,显然不利于疗养(也就是恢复被扰乱的平衡!)。比如,休息可以使人留心自己的身体与近旁的环境(它仿佛构成第二层皮肤……),由此生出认识,乃至一种相通;这本身便具有真实的“行动”性质。毫无疑问,学习确实是一种行动(因为它有一个无可争辩的效果:认识的出现……)。

逐一察看我纳入“行动——不行动”组的十四对对子(想必还可以找到更多自然归入其中的),可以发现:或许只有一个例外,其余每对都是第一端、即“男性”的一端,按我们文化传递并自童年灌输的反射性态度,被赋予声望与“价值”。这仍是我们文化中那种积习已深的失衡的标志:唯独抬高阳——我在前面已有机会提到它\fnote{46}{(*)}{见“阳葬阴(1)——或肌肉与肚肠”一注,n\textsuperscript{o}~106.}\dmark{}\origpage{496}。我遇见的几乎所有阴阳对子,也都可以作同样的观察。这一点实在触目;此前我从未花工夫如此细致地核对过。

在我方才写下的这些对子之中,唯一一个依我看堪称例外的,是\textbf{热情——宁静}这对子:因为在惯常用法里,“热情”一词常与一场爆发、一种暴力的形象相连;或不然,与一种\textbf{放纵}相连——那放纵,遗憾地贴近环绕“\textbf{堕落}”这样的词的那团联想之云。说来凑巧:放纵与堕落,指的正是以\textbf{阴}的、女性的过分支配为特征的心理失衡状态!而对称地,循着同一套按钮机制(它们揭示的是我们时下的制约,而非“宁静”这样一件事物的本性),“宁静”一词(与“热情”相对)便与\textbf{自制}的形象相连——即一种恰如其分地属于\textbf{男性}本质的品质。(实际上,“控制”的阴的对应者,绝不是“热情”,而是“交付”。)

于是,这里发生的是:由于人心对某些事物之本性的一份普遍混乱——它表达为人们用以指称它们的某些词的用法的同样混乱——“热情——宁静”这个阳-阴的对子,与另一对观念

\begin{center}松弛——控制\end{center}

混为一谈了;后者两端是阴-阳(却并不因此构成一个“对子”:在场的两端根本无意成婚!)。因此,依我看,所谓价值化规则的那个“例外”……

恰恰是一份格外有趣的确认!至于我辑得的其他几个例子——那些在一个阳-阴的对子里似乎被抬高的是“阴”端者——想必也是同样情形。

况且,我全然没有把握:我在这所谓“西方的”文明里所确认的这份世界观上的扭曲——它出自那份系统偏袒男性、对立于女性的成见——在中国传统里、乃至在今日的中国世界(或更普遍的“东方”世界)里,就真是轻微许多\dmark{}\origpage{497}。在日常生活的层面上,没有任何征兆能使我作此设想:无论经由我的东方朋友们的日常,还是经由中国或其他远东国家的传统与今日生活传到我耳中的回声——恰恰相反。依我看,倒不如说,一份对阴阳动力学的精微知觉,几乎排他地局限在某些\textbf{技艺的实践}之中——比如书法、诗、烹饪艺术,当然还有医学之艺\fnote{47}{(*)}{(10月21日)我在清单里忘了\textbf{占卜之艺}——《易经》或“变化之书”;它如今在欧洲与美洲的某些圈子里享有盛名。《易经》占卜语言的基本“词”,是那64个“卦”:它们不是别的,正是六位“爻”(阴爻与阳爻)的$2^6$种可能组合——从纯阴(阴的六次重复)到纯阳(阳的六次重复)。那里似乎有一种关于阴阳组合的极其精微的炼金术;据说它曾令Jung心醉神迷。这门炼金术的趣味(尤其作为“一部原型集”)依我看,先天而言,独立于它在占卜中的使用、以及人们对这类使用肯给的信用。}。

尤其后者,顶着“中医”之名、并借着针灸的几桩惊人成功,在过去二十年里终于在我们这里取得了公民权,并被赋予了威望。然而仍有许许多多的人不知道:在中医里,对身体、对气的循行及其扰动(它们构成我们称为“疾病”的病态)之把握的阿尔法与欧米伽,恰恰正在一份极其精微的阴阳辩证之中。这份辩证“行得通”这一事实——既然以它为根基的“中医”是有效的(包括在许多西方武库无能为力的个案里)——可以被视为阴阳这些“本原”或“侧面”或“模式”(把握的或存在的模式)之实在的一种“证明”:它们不是某些哲学家与其他诗人(不说是江湖术士)从帽子里掏出来的纯粹思辨。

人们当然会问:这样的证明有什么意味,乃至关于某一种世界观之效验的任何“证明”又有什么意味。且假定证明果真说服了人(也就是说,当事人肯被说服),甚至再外加一份厚礼:所论的视野深刻、因而是有益的——世上最好的证明,依然无力\textbf{传达\dmark{}\origpage{498}一份视野},更遑论一种世界观。死死“被说服”于一份依然陌生、依然未被理解的视野,对你毫无裨益。干脆说:那甚至没有意义——或者更确切地说,那份“信念”的真正意味,当事人并不比他所佯装修进其沉重视野的这份视野懂得更多。

当视野被理解、被消化之时,“证明”这个问题本身,便显得出奇地离奇——有点像当人们明明看见天是蓝的,却还要证明天是蓝的;或者像证明一朵所爱的花的香气是香的……

\subsubsection*{(c) 一半与全——或裂缝}
\tocsubsub{(c) 一半与全——或裂缝}
\label{ch18:local:h14}

\notehead{112}{}
\phantomsection\label{ch18:local:n13}

(10月17日)我关于“女性”与“男性”双面的最初反思,出自一场对我自己的反思。那是在1979年初;那时我还不知道“阴”“阳”这两个中国词,也不知道中国文化传统里有一套关于阴阳之不断游戏的精微“哲学”。这东西我是同年末才知道的,我想,是经由我的女儿、尤其我的女婿艾哈迈德——他那时开始对中医感兴趣,并在随后几年里深深咬住了它。他对我说的东西,大部分与我抵达的视野相互印证;这毫无使我惊讶之处。倘若说有惊讶,倒是在几个“对子”的个案里:其“自然的”阴阳角色,在中国传统里似乎是被颠倒的。我的反射(在这个案里,是强“阳”的反射!)曾是一种浮于表面的确信:这份“颠倒”必定出于一种文化扭曲——而且我不曾凑近细看\fnote{48}{(*)}{这份对着一个千年传统——它本可以教我更谨慎些——的武断自信的反应,与儿时的那个反射如出一辙:当年它使我驳回书本里教的那个(说实话够复杂的!)公式 $\pi=3.14\ldots$,而偏信我以自己的手段确证的 $\pi=3$。(见“化圆为方”一注,n\textsuperscript{o}~69.)诚然,关于这段阴阳的故事,我早已有充分的机会觉察:“女性”与“男性”之本性及其相互关系的知觉,被积习已深、力量可观的文化扭曲败坏着。但当时我还不曾觉察:对这份关系的精确而精微的知觉,在某些中国传统技艺的实践中是本质的东西,而且被推到高度精微的地步。}\dmark{}\origpage{499}——那是一个我的男女阴阳练习显得十分遥远的时刻,因为我那时正投入一场个人得多的禅观:关于我父母的生平、关于我的童年。想必是数月或数年之后,经由几路旁证,我才发觉:在某些个案里,我对这样那样“对子”中阴阳角色的把握,仍有点肤浅;我曾相当匆忙地,把一些本性不同的局面装进了同一个口袋——而中国的阴阳辩证,是小心区分它们的(注112$'$)。如今我明白:在我这里,对阴阳的把握,仍然相对粗糙、相对静态——尤其若与某些中国传统技艺(如医学:它与食疗及烹饪之艺血肉相连)的实践所要求的那种精微相比;在那里,这份把握最终成了第二天性。

我已有不止一次这样的印象:在这些技艺的修习者与从业者身上——无论东方人还是欧洲人——这份把握的精微停留在断片的形态:它在很大程度上,小心地只局限在该技艺的练习之内。在日常生活中,它更像一个寻常的“所知”那样起作用:干脆叠加在文化(及其他)制约的“所知”之上,而对后者多少形同虚设。换个说法:我有这样的印象——在这些“博识之至”的人们身上,世界观与自我观、知觉实在的压抑机制,与常人身上的毫无二致。

这份印象,又与我读过的两三份文本相互印证——它们的作者号称“在局中”,担着向人概述中国传统阴阳哲学的差事。(其中一位是法国的著名东方学家;名字我一时想不起来。)打动我的,是在这些文本里,阴与阳被呈献为“\textbf{对立}的”(或“相反”的)乃至\textbf{敌对}的本原(后一个词在其中一份文本里反复出现),而不是\textbf{互补}的。这份“对立”或“敌对”,据说有其典型表达:即在人类社会之内、在社会所设立的伴侣之内,女人与男人之间的那份敌对。

夫妇之间的敌对,确实是实在——东方与西方皆然。它深深扎根于文化,以至于有时俨然人类境况的一个(有时令人错愕的!)方面,乃至人身上或人类社会里冲突的根源。这份敌对的实在无可辩驳;它当然超出那些尽力驱邪的通行老生常谈。\origpage{500}这份“社会性的”实在,是一份不可追忆的制约的产物;它早早便在成形中的“自我”里扎根,并结构着它。然而,在这份实在之外,还有一份更深的实在——它来自更远得多的地方——它在性爱的冲动本身之中起决定作用。那是两性之\textbf{互补}的一份深刻而本质的实在:其中没有任何“敌对”的位置。这也是在一切物种之中分明显现的实在——唯独我们的物种例外;在我们这里,它在很大程度上被文化敌对所遮蔽,即被人与人类社会特有的那种\textbf{分裂}状态所遮蔽。

时下浪漫的老生常谈——“我俩”体——支配着大部分文学与媒体,把一份廉价赝品的“互补性”捧上云端,同时对男女之间那份令人不安的敌对侧面,罩上一块遮羞布;或(至多)把它当作一桩有点刺激的意外事件来打发——正好给一顿本来太过寡淡或甜腻的饭食加点辣。一旦越过这类安人心的套话,人们立刻直面这男女敌对的实在——一个显然普遍、且兼有一份无懈可击的顽强的实在:狗牙根般的顽强!但以这份无处不在、无可辩驳的实在为出发点,去设立一种阴阳、“女性”与“男性”的宇宙敌对,便是把人类社会与个人那份撕裂的、深刻分裂的状态——我们这个物种特有的一种疾病——投射到整个宇宙。这同时也是延续自己对自身之内\textbf{另一种}实在的无知(它接上那宇宙性的互补和谐之实在)——一份同样顽强(或不如说,不可摧毁)却更为隐蔽的实在。它逆着那些心照不宣地设立敌对的制约:无论在女人与男人之间、妻与夫之间,还是在我们身上作为“女人”的东西与作为“男人”的东西之间。

老实说,这种\textbf{二元}的或\textbf{好战}的宇宙视野——其中事物的一个侧面,与一个同样本质的“对称”侧面恒常交战——这份视野绝不是\textbf{反思}的果实:仿佛它会“出发”(如我方才所写)自人类伴侣与人类社会里冲突的实在,继而把它“推导”(或按我更准确的说法,“设立”)到整个宇宙。它不多不少,正是文化制约的忠实表达——简直可以说自动的表达——并且顺应该制约的一项本质功能:\textbf{维持人身上的冲突、分裂}。\origpage{501}分明,倘若我肯以那双我出生时领得的眼睛去凝望宇宙,并确认:处处——除了(据说……)在我自己身上、在我的同类之中——“女性”与“男性”都是彼此不可分解的互补者;正是它们的联姻与结合,诞生了造化中一切活物与“死物”里的和谐、创造之力与活的美——那么,维持我身上“女人”与“男人”之间那份既立的敌对,想必是无从谈起的事;或者不如说,这份敌对,\textbf{自那一刻起}便已化解。相反,倘若我自命在宇宙中处处“看见”并非在场的“对立”与“敌对”(即便如此做时,我循着一份可敬的、几度千年的传统),那绝非因为我用过了我的眼睛,而是因为我不过像人人一样,\textbf{重复}了自洪荒以来代代相传的东西;无论如何,是服从了文化共识那份无声而强制的指令——正是它,在我的身上牢牢设立了一份分裂、一场冲突;而我还要把它合理化(并由此延续它),称为一种“宇宙的必然”。

关于伴侣之中的敌对,更普遍地,关于女人与男人的敌对,当然大有可说——我信得过我的同类:关于它,包括中肯的东西,已经写下了许多。这里不是铺陈这个极为有趣的主题的地方,尤其不是铺陈这份敌对在我们父权社会里所采取的特殊形态的地方。我觉得,在那些分明见过其存在的人们中间,有许多人认为:社会的结构——它映照并落实男人对女人的支配——应为这份敌对负责。他们肯定有理;而且我疑心,在一个母权倾向鲜明的社会里,想必能找到一份类似的敌对,以或多或少对称的方式显形。我只想补一句:这份因果在我看来却是\textbf{间接的};它似乎经由一份更隐蔽的因果而行使——今天的反思已拂过其表面。伴侣中分裂的这份更隐蔽、更本质的缘由,是\textbf{个人内部}的分裂状态——无论女还是男——对着自己的冲动(尤其性的冲动)与自己的官能。我在其中看到男人与女人之间敌对的真正根源,一如看到他们在精神层面上那份\textbf{相互依附}的真正根源:我是说他们双方那种\textbf{内在自主的匮乏}\dmark{}\origpage{502}。

这一在自身之内的分裂,在于一种亲身的、隐秘的信念——无论男还是女——相信自己只是一个\textbf{一半}。这份信念的一个征兆,是那种弥散而阴险、从不被检视的\textbf{裂缝}感,或许还有\textbf{残缺}感:仿佛唯有另一性的伴侣,才能解救我们——至少临时地。在那些“大男子气”或“喀耳刻”的应景门面背后(还有许多别的门面),每一个人——男人一如女人——面对潜在的或现实的伴侣,都摆出一副\textbf{乞求者}的姿势:等着对方的(或多或少)好心,给他一场昙花一现的解救——他盼它彻底,它却总是跛脚——解救他那份可怜的破缸状态,甚至不如说碎缸状态:总之是\textbf{半只缸},正在寻找另一只,好与它多少重新粘合(而且粘得坏多于好,人们猜得到……)。

这份裂缝感,或者说,这份对我们真正本性的\textbf{无知}——对我们那超越与性别相连的生理特殊性的根本\textbf{统一}的无知——我身上的这份深层分裂,依我看,是纯然社会制约的产物。至少在婴儿的最初几天与几个月里,人们察觉不到它的踪迹。况且,这份制约绝不限于对“男性”的抬举、以牺牲“女性”为代价,或反过来。说到底,倘若我感到、接纳并且被接纳为\textbf{同时}是两者、“男人”\textbf{与}“女人”——带着一份可以随我人格的这一侧面或那一侧面而变动的主音;那份主音,绝不限于在生殖器官层面占支配地位的那个(诚然极其重要的)主音——那么,我周围抬举的是“男性”还是“女性”,就不再那么要紧了。在我的性冲动的层面上,我的“个人价值化”反正会往往投向与我相反(抱歉,我是想说:互补)的性别,而不会因此感到自己在这具\textbf{不同}的身体面前低人一等(也不高人一等)——一股强而深的冲动把我引向它。况且,无论与性别相连的价值化,还是任何别的价值化:社会共识所借贷的“价值”或威望(对自己或对他人)所占的分量,在一个不大(或全然不)被我所说的“裂缝”感触及的人身上,是相对次要的,甚至微不足道——即在一个身上活着那份自发的\textbf{安然确信}的人:那确信既非僭妄也非门面,而是一份对自身本性的完好无缺的认识的显形。\dmark{}\origpage{503}

“裂缝”或人身上的分裂\fnote{49}{(*)}{我在这里避免使用时下相当流行的“阉割”一词——一个暴力之至的词(这一回真是超级阳的!);它还有另一重不便:暗示一份无可补救、不可逆转的残缺的图像,从而刺激一份慌乱、反叛或认命的反应——这类反应,只会加强一种受阻状态,而无助于其朝渐进化解的方向演化。}并非仅仅是价值化的产物,还有一个征兆:这份分裂在男人身上与在女人身上同样肆虐——也就是说,肆虐于那份自命“抬举”他的共识的所谓“受益人”身上,而这份共识(在某种意义上)却把他连同其伴侣一起压断了腰。人们看到:这份分裂愈发尖锐、愈发暴烈,恰恰是当对一种性别(以“惠及”另一性别为名)的压制愈发强硬、愈发无情之时。\origpage{504}人们简直可以说:“社会”(压抑的源头与工具)在安置这些“压制”之时所遵循的原则,是让男人否认其阴的一面,让女人否认其阳的一面:“分而治之”!而这份由共识造出来、用以压断并奴役男人与女人的“分裂”,还同时在两块板上运作。最显见的一块板,是伴侣之内的分裂:它经由设立一种或多或少的专制的支配——一方性别支配另一方,男人支配女人,或反之——而获得\fnote{50}{(**)}{至少在表面上。不过,如前面所提示,往事情深处走,人们会发觉:伴侣中这份由男人支配女人所维系的分裂,有一份更深的“根”——几行之后我还要回到它。}。一方被认定统治另一方——而双方到头来都成了奴隶\fnote{51}{(***)}{况且,还是这样的奴隶:无论如何不肯与他们的锁链分手;那锁链,对他们比生命还珍贵……}。因为,当妻子或丈夫遭轻蔑时,被交付轻蔑的是双方——有时是旁人的轻蔑;但更深、尤其,是自我的轻蔑。

而在这里,我们接上了分裂游戏的“第二块板”——那更隐蔽的一块:即人本身的分裂。它是伴侣之分裂的隐秘弹簧;伴侣的分裂使它加剧,却不把它归结掉;而它绝不由对一种性别的抬举、以牺牲另一种为代价所单独造出。它宁可是一份无声而不间断的强制的产物:自我们最初的年月起,周围环境便对我们行使这份强制。它逼我们否认我们人格的整整一个“坡面”(“阴”的坡面,或“阳”的坡面\fnote{52}{(*)}{原则上,除意外情形,强制的方向,推着男人否认其阴的一面,女人否认其阳的一面。女人的处境更为微妙:她被认定要否认的,恰恰是那些被社会共识赋以威望、因而她想必感到受怂恿要去栽培的特质。她于是受着两股相反方向的压力;无意识要结构一份“可操作的”身份认同,其任务也随之倍加复杂。}),否则就有被拒斥之虞——被当作可笑,或不合体统;无论如何,当作不可容忍。

\subsubsection*{(d) 原型认识与制约}
\tocsubsub{(d) 原型认识与制约}
\label{ch18:local:h15}

\notehead{112$'$}{}
\phantomsection\label{ch18:local:n14}

\fnote{53}{(**)}{本注出自前一条注的一条脚注(见其首段的指涉)。}比如,在母体-胚胎与阴道-阴茎这两对里:阴阳角色的分配毫无疑义,阴的一端环绕并包含阳的一端。这曾使我匆忙断定:在“容者——所容”这对子里,“所容”是阳;而形-底、外-内、周-心这几对(在其中,正如我的确感到了的,前一端分明是阳,却同时是“容者”)却不曾使我警觉。实际上,在母体-胚胎与阴道-阴茎这两对里,我把重点错放在了在场的两端之关系的“几何的”或构形的侧面——一个次要的侧面;而在本案里决定角色分配的,是主要的侧面:\textbf{滋养者}相对于\textbf{被滋养者}(后者为阳)是阴;\textbf{进入者}相对于\textbf{被进入者}(后者为阴)是阳(同样,\textbf{给予者}相对于\textbf{领受者})。

我关于阴阳的反思,尽管如此有限,却在我身上立稳了一份亲身的信念:在人们对阴阳角色分配(或对一个给定的人身上“阴”或“阳”的“主音”)的种种个体领会差异之外——那些领会极易染上“文化扭曲”——这样一种分配(或“主音”)的“自然的”样子,确实存在。它有一份同样无可辩驳、“宇宙性”而不可改移的实在(就那些普遍本性的对子——如迄今所论者——中的角色分配而言),一如一条物理定律或一个数学关系;即便它既不能经由经验(按自然科学实践中该词的意义)来“确立”,也不能经由一份“证明”甚至一份“论证”来确立。\dmark{}阴阳的这份实在\origpage{505},是经由一种直接的知觉把握的;这份知觉可以(除别的途径外)经由一场足够深入的反思而发展、纯化。

我觉得,这样一场反思的一个主要效果,恰恰是使我们超越那份由周围文化在我们身上程序化的老生常谈式的反射,重新接上与实在本身的接触。实在,依我看,早已在场于心灵的深层——作为一种“原型认识”:它置身于文化制约的射程之外。反思的角色,是使人得以重新接上这份已在的认识,并细心把它从“所知”的表层——即文化制约——中澄清出来。

我在这方面着手的工作,对理解世界与我自己,并由此对我的日常“做”与我生活的航向,曾是重要的。这项工作(一如在许多别的场合)在我看来,像一次最初的突进,像一扇我刚刚推开的门:门外是一座宏阔的全景,尚待勘察。做成它的手段我样样在手——但我不知道有朝一日会不会去做\fnote{54}{(*)}{一如我不知道我在这里看到正在敞开的这类工作是否已经有人做过(说到底,是对宇宙中事物的诸品质及其把握方式,在阴阳互补的和谐之光照下,作一种“局部的与全局的”“地图”式的研究)。况且,这是一个纯然次要的问题:这里要做的,不是提交一篇这个或那个的博士论文,而是深化对世界与自己的领会——那领会,只能是个人工作的果实。}。姑且把数学放在一边,同样“多汁”、且更个人、更灼热的反思主题并不缺,它们无疑会先于对阴阳的一场更一般的反思之深化而占据优先……

\subsection{(4) 我们的母亲,死亡}
\label{ch18:local:h16}

\subsubsection*{(a) 行为}
\tocsubsub{(a) 行为}
\label{ch18:local:h17}

\notehead{113}{}
\phantomsection\label{ch18:local:n15}

(10月21日)三天没有写注了。我的日子被别的任务与事件占去。其中之一,是皮埃尔的来访——带着他的小女儿娜塔莉,昨晚到的。他打算住到明晚,并在此期间读一读葬礼已写成的部分。对我花了将近三个月写成的一份文本,这恐怕稍嫌短促……

我得以分给反思的时间,用来继续\dmark{}\origpage{506}摆弄“阴阳对子”与它们构成的组。这个题目自有其迷人的地方:它兼具探查一个数学“结构”的那份独特兴味——其本性在工作的进程中渐渐明确——与一场关于世界、关于存在的反思的兴味。每一个主要的阴阳对子,都代表一把“锁孔”(在无数别的锁孔之中),透出世界的一个侧面,或世界的一个角落。我迄今辑得的那些对子的“组”,似乎对应于把握宇宙事物的不同可能方式,宛如一扇扇向宇宙敞开的门,从各个不同的角度把它显示给我们。每一扇这样的“门”都有大量的锁孔——甚至数目无限的锁孔——可供窥视;且待哪天,或许干脆把门推开?眼下我只限于侦察出其中相当多的一批锁孔(我找到的足有两百来个),并逐一凑上眼睛看上片刻;同时每一次都意识到:这里有的是值得久久观看而不虚掷光阴的东西……

不过,我更急的,是先去瞧瞧这里或那里另一把还可窥视的锁孔,并且把这些门统统巡视一遍,多少弄清它们相互之间如何布列,或许还弄清:使人们得以侦察出这一扇或那一扇门之存在的那些锁孔,是循着怎样的“模式”布列的……

终于,一周多前侦察出的那十八扇“门”,又添了三扇,共二十一扇;它们布列成一幅图(我曾称它“大致呈圣诞树形”):如今带一根“主干”,由九个“顶点”(或“门”,或“组”,或“角”)构成,以垂直的“棱”或“纽带”相连;主干两侧各带六个别的顶点,它们与主干相连、彼此相连,构成“树枝”\dmark{}\origpage{507}。

\fnote{55}{(*)}{(10月24日)要预言会不会终于出现一些自然地插入我迄今辑得的任何一个组之外的阴阳对子——也就是说,是不是还有别的组或别的向世界敞开的“阴阳之门”,乃至数目无限?——我实在窘于断言。况且,我找不到别的组,绝不意味着不可能存在无限多个别的组——甚至无限多个逃离人类经验、逃离我们把握宇宙之手段的组。这使我记起,近几年里我不止一次被这样一份直觉打动:从蚂蚁或微小的蚜虫,直到已经与我们相当接近的哺乳动物,每一个动物物种都有一些逃离一切其他物种(当然包括我们)的把握宇宙的手段与知觉方式;以致就(不妨这样说的)感知把握我们周围事物的模式之丰富而言,我们的物种不“覆盖”也不“包含”任何别的物种,一如没有别的物种包含我们。而我顺势贸然说出的那句“一如没有别的物种包含我们”,在我看来又嫌仓促,乃至僭妄:鉴于在纯然感知的丰富与精微的层面上,我们物种的演化反倒往往走回头路、退化。只是在智力的层面、在心智图像的精微——尤其与语言相连的那些——上,我们胜过别的物种,依我看。自发来到我注意中的大部分阴阳对子属于这个特具“人性”的音区,而只有寥寥几把带有(此外)显见的感知意味——如影——光、冷——热、低——高,以及别的几把——这绝非偶然。}相当滑稽的是:最近几天出现的三个“新”组里,有一个恰恰是所有组中最显眼的、最原初或最原始的:即对应于对阴阳的那份头一个直觉——“女性”或“雌性”,与“男性”或“雄性”——的那个组。我觉得,它以“\textbf{父——母}”这个原型对子表达得最为醒目(胜过“男——女”——后者属于同一个组)。这个组带着强烈的性意味,显形于诸如“生育——受孕”或“阴茎——阴道”这样的对子;而这些对子本身,又属于环绕那个卓越地的行为——原型之\textbf{行为}——的联想之云:那创造性的拥抱,经由孩子的诞生——经由出自这行为的\textbf{作品}——(至少在潜能上)把女人变成母亲、把男人变成父亲。

与性爱冲动相连的这些意味,五年前的那场反思里始终占据前台。它们而且还在那部大名鼎鼎的“诗作”——当时那场反思凝聚而成的、约130页的东西——的通篇,享有一份近乎不间断的抒情强调;其效果,连最宽容的读者也会读得发腻。在近日反思所依凭的那份独一无二的参考文本里,我干脆“忘掉”了那组当然开道于游行队伍之前(而且不无道理)的对子,肯定正是对这份双重“蓄意”——诗的与Erotique的——的恼怒反应\fnote{56}{(*)}{形式上的这份蓄意,映出一种内在的态度:一个传讯使徒的角色之选定。对此可参见“非古鲁的古鲁——或三足马”一节(n\textsuperscript{o}~45)的结尾,及与之相关的注n\textsuperscript{o}~43.}。

那部著作的题目“乱伦颂”也有些挑衅的味道,\dmark{}\origpage{508}足以使人对其意图与其“讯息”产生误解。况且,这意图与讯息在写作进程中演变颇剧:诗的轭具并未妨碍一项深化的工作继续进行,一份澄清也确实发生了。一个最初的、主要的用意,是探测性爱冲动的某个侧面(我感到它深刻而本质)——就它经由我自己的经验所知的样子。所以说的,首先是男人身上的性爱冲动,或更确切地说:那份“阳”的冲动——它在爱的游戏与爱的行为中对应于“男性角色”,但它以可变的强度\fnote{57}{(*)}{(10月24日)这份冲动的在场,常被力量强大的压抑机制或多或少地彻底遮蔽。我有这样的印象:在男人身上,这份阳的冲动往往比与之互补的阴的冲动占优势;在女人身上则相反。但文化制约,以及对这些制约的种种“正面”或“负面”的内化方式,如此严厉地(而且常以复杂的方式)干扰着原初冲动的交互作用,以致有时很难从零星、短暂且往往变了质的表现背后辨认出它们。}在女人身上一如在男人身上在场。很久以来——或许一直如此——我知道这份冲动按其本性是“乱伦的”:它也是“回归母亲”的冲动、回到原初怀抱的冲动。

这场宏大的回归,在爱的游戏中被“排演”并重新经历,终于在一个寂灭、一个存在的熄灭、一场死亡之中达到顶点并完成。以全部的丰满活出爱的行为,也是活出自己的死——如一次“倒转的诞生”,使我们回到母亲的怀抱\fnote{58}{(**)}{此外,我深信:这份阳性的爱欲冲动的内容存在于一切生物物种,甚至存在于生物界之外;它对应着宇宙万物同一种深层的动力:每一个创造的过程(或“行为”),都是阴与阳、“母亲”与作为孩子的厄洛斯的一次拥抱;孩子回归母亲,消融于她之中。从这孩子回归母亲的“死亡”(或“倒转的诞生”)里,创造行为的果实——“作品”——仿佛从滋养的母胎中涌现。这是“孩子”、新事物的出现,由赋予它生命的“旧物”的死亡与更新之行为而来。在这个宇宙的维度上,原初的性冲动自古就在,远在人类出现之前,甚至早在我们的星球上出现(生物学意义的)生命之前。}。但这同时也是一并僭越两道威力可观的禁忌:乱伦的禁忌——它把“母亲”排除出爱欲对象;以及那份(至少在我们的文化里)把生与死分开、把它们对立起来如不共戴天之敌的禁忌——生与死\dmark{},诞生与死亡。然而我早已分明懂得:爱的行为既是一场死亡——它在高潮的痉挛中完成——又是一次诞生、一个存在的更新:出自这场死亡……宛如一株新苗从滋养的大地里婀娜升起——而那大地本身,由那些倾身于她的存在们的创造性分解构成……

正是在五年前那场关于爱之行为之意味的反思中,我终于懂得:“死亡”与“生命”是同一场紧紧相拥的伴侣中的妻与夫;生命永恒地诞生于死亡,又永恒地倾身于死亡。\fnote{59}{(*)}{(10月24日)奇怪的是:在其后几星期我辑得的一些阴阳对子中,竟没有“死——生命”这对子。也许是因为它与那对相近的“死——生”(或更好:“亡——生”)相混了——后者在册;于是前者便俨然与后者重复了。}或者不如说:生命永恒地倾身于\textbf{死亡},以便永恒地从\textbf{她}——那多产而滋养的母亲——重生;而她自己,又被她的无数儿女对她的永恒回归所不断滋养、不断更新。

而人类的伴侣——妻与夫、恋女与情郎——当其充分地活出把彼此相吸的那股冲动时,便宛如生命与死亡那无尽联姻的一则寓言:每夜爱之终了,情郎倾身于恋人之中、死于恋人之中,以便与她一同从这场死亡中重生——在他们共同的拥抱里……

在同一场反思之初,我曾把人格分裂的一个本质侧面,视觉化为一种“切分”、一道“水平”的切分:即由乱伦禁忌所设立的那道切分——它把孩童从母亲那里“切开”,一如它把生命从其母亲死亡那里切开,一如它也把一代人与生养它的那一代切开。

倘若我起先看到的是这道切分,那无疑恰恰因为我是免于此切分的人。然而,我的生活——一如每个人的生活——曾被另一道大切分深深打上印记;那是我在后来的反思中看到的、并称之为“竖直切分”的:它把每个存在之中“女性”与“男性”的两个“一半”分开来、互相对立,只容忍其一而排斥其二。这正是这一两周来我所投入的这场关于阴阳的长篇枝蔓所谈论的切分。\dmark{}\origpage{510}

我觉得,如今看来,这一(“竖直”的)分裂比另一道(“水平”的)更为关键;在某种意义上,后者为它所蕴含、为它所“包含”。说到底,把孩童与母亲分开、把生与死分开;把一份污秽感、一份排斥或羞耻感——加到死亡上、加到那把孩童与母亲连起来的冲动上——这分明也是在把两者彼此切开、使它们互相对立:把夫与妻——在那两场不可分解而原初的宇宙伴侣之中:母亲——孩子,死亡——生命——彼此切开\fnote{60}{(*)}{这里我按“自然的”阴阳次序书写这两对,以阴的一端——“原初的”一端——起头。关于“母亲——孩子”这对子,应当注意:“母亲”一词也出现在前面已谈过的另一个重要的原型对子里,即在一切之中最原始的“母——父”对子里;它把名字借给了它所描述的那个组。(此外,“母亲——孩子”这组的组名属于别的对子,即我以“因——果”对子之名所称的那一个。)而且,“母亲——孩子”这对子的阳端“孩子”,本身也属于另一个原型对子“老人——孩子”——它与十分有趣的“成熟——天真”这对子相邻。这两对插入我称为“高——低”的组里;那个组,是我迄今侦察出的所有组中(单就数目而言)最丰富的一个,含着许多别的可观的对子,如衰落——腾升,亡——生,毁灭——创造,忘却——学习,终——始……在列举这几对时,我几乎是强制了自己:违背积习,按阴阳的次序为它们命名。粗看之下,这个新次序显得有点滑稽,甚至离奇——总而言之,世界颠倒了!可凑近一看,人们发觉:这个不惯常的次序,向我们显出两端之关系的另一个侧面,一个对惯常侧面的补充:比如(照惯常)“生”先于“亡”——而我们方才分明看到,在更深的意义上,“亡”先于“生”。

我那场反思的总名字“收获与播种”,也是一样——它无疑是一个阴阳对子(这正是我此刻发现的!)。它的命名仍旧循着与惯常的阳——阴次序相反的次序:收获理应跟在播种之后,而非相反。然而这个名字是毫无歧义地自己强加于我的;任何时刻,甚至不曾闪过“这名字本可以是反的——‘播种与收获’”的念头。那是因为:一次次面对不合时宜的收获,总归把我的注意力引回到它们所自出的那些播种——仿佛收获的深意与功能,恰恰是执拗地把我带回到我那双早已被遗忘的亲手播种……}

有意思的是,这最后两对并不曾在《颂》里所辑的那些对子之列。而“死——生”这对\dmark{}\fnote{61}{(**)}{要注意:在“死——生”这对子里,“死”一词的意味,与“死——生命”对子里的不同:在前者它指一个行为(“辞世”的同义词),在后者指一个状态。德语里有两个不同的词:“Sterben”(不带“辞世”那份略显轻佻的意味)与“Todt”。法语里,依我看,最好以“亡——生”来称这对子;这就消除了“死”一词意味上的歧义。},则与我的恋爱经验有更直接的联系,是在册的。“母亲——孩子”与“死——生命”这两对,只是在最近这些天的反思里才出现——在同许多迄今逃脱我注意的别的对子之中;其中最有兴味的一个,是“恶——善”。那是可以称作“艰难的”对子中的一个(一如“死——生命”):那些威力极大的制约使我们把两端知觉为敌对的“相反者”,而不是不可分解的互补者。分明,五年前写《颂》时,这些制约在我身上比今天更强。不过,《颂》里已经有相当一批“艰难的对子”,其中有“混沌——秩序”与“毁灭——创造”……\origpage{511}

回顾起来,对各色阴阳对子之本性——作为不可分解的互补者构成的和谐整体——的一份多少深入的领会\fnote{62}{(*)}{我是指一种不停留于纯然智性的领会;它具体显形为一种改变了的、对他人、对世界或对自己的关系——因为为人之道变了。},如今在我眼前,俨然我们那通往发现世界与发现自我之旅中一级级有待越过的“门槛”。这样一个“门槛”,越是在所论对子“艰难”之时,便越是可观;也就是说,把它作为“对子”来把握所撞上的内在抵抗——文化制约的表达——越强。

\subsubsection*{(b) 挚爱者}
\tocsubsub{(b) 挚爱者}
\label{ch18:local:h18}

\notehead{114}{}
\phantomsection\label{ch18:local:n16}

(10月26日)昨天的反思\fnote{63}{(**)}{这里说的是前一天的反思,记在我排在今天这条注之后的注n\textsuperscript{o}~116里。}起步时有些费劲。这无疑是因为最近几天频频的打断。不过从前一天起,有一样东西在我身上还热乎着,我急于把它托付给纸面——哪怕只写几行。事后我大为赧然:它竟半路丢失了,被一些寻常杂事挤掉了!今天,我不忍心就这样过早地、仿佛出于一场误会地与它分手——在真正与它相识之前,干脆说。

我翻阅了《Zupfgeigenhansl》新近再版的版本\fnote{64}{(***)}{Wilhelm~Goldmann出版社(Wilhelm Goldmann Verlag),1981年。}——这部德国古老民歌的经典集子,约在本世纪初编定出版。据说它已无处可寻,但几位途经我家的德国朋友\dmark{}\origpage{512}带了一本来。前天动工之前,我匆匆翻了一页,有点像路过时与一位老朋友握握手。我撞上了那首歌《今天与明天》;我扫了一眼,没有真正停留——当时我也急着回到等我的工作。但这并不妨碍某样东西“咔嗒”了一声。我分明感到:这些如此单纯、外表如此天真的歌词,轻轻地触到了我身上某种深的东西——而且,这东西离我三天前勉力想唤起的那个,又如此之近。我正准备把这方面的笔记誊清。或许我模糊地感到:我方才扫过的那几节歌,比我那些以断然的简短、匆匆写就(又赶着去做别的事、仿佛顺带)、当日情绪尚且缺席的笔记,更是我想传达的东西的更忠实、更有说服力的信使。

今天早晨起身时,我试着把这几节歌译成法语——它们的曲调我并不知道,但两天来它们一直在我体内歌唱。这无疑是一种更好地找回它们的方式,更好地让它们的滋味与旋律渗入我体内。令我惊讶的是,我在另一门起初显得桀骜不驯的语言里,找回这部德语文本的节奏与音乐的一些部分,并没有太费劲;而字面的意思,始终贴得很近。以下便是这七节歌词,我尽力将它们译了出来\fnote{65}{(*)}{(10月29日)下面的版本,是随后三天里修订过的版本。那天晚上大家唱了歌,我得以学会这首歌的曲调。初版的改动,大部分是为了照顾歌唱文本的节奏与重音的要求。必要时不妨把音节在曲调的音符之间妥善分配;这支曲子可以用法文唱出,而不至于在任何时刻违背重音(一如近来某些新酿的法语歌不幸常有的情形)。}:

\begin{quote}\raggedright
“今天还在,明天\\
我将在你身旁;\\
可一到第三天,\\
我便要离开。”\par\medskip
“可你几时回来\\
我的爱,我温柔的意中人?”\\
“等红玫瑰落雪,\\
等新酒落雨!”\par\medskip
“玫瑰不落雪,\\
酒也不落雨;\\
就这样,我的爱,我温柔的意中人,\\
你也不再回来!”\par\medskip
在父亲的园子里\\
我躺下,睡着,\\
一个美丽的梦落到我身上,\\
白雪纷纷落在我身上。\par\medskip
\dmark{}而当我方才醒来,瞧——\\
纯然的空,纯然的无——\\
\origpage{513}却是那些娇美的红玫瑰\\
在我身上开花……\par\medskip
小伙子回来了,轻轻走过\\
那座美丽的园子,\\
头戴玫瑰花冠,\\
手捧一杯酒。\par
\newpage
他用脚轻轻一踢\\
那座小小的坟茔,\\
倒下——玫瑰落雪,\\
新酒也落雨……
\end{quote}

当我摸索着试图还原我所读到的东西时,我心中有一种喜悦、一种幸福:随着时刻流逝,这首歌仿佛变成了我的一部分。那里有这份剥落而温柔的美的东西,既宁静又切肤,一种由亲密相拥的喜悦与悲伤构成的重而深的美。我相信,很少有人不对这样一首歌多少有所触动——即便他们会抗拒这种触动;人们也常常抗拒一份不期而至的情绪,当我们身上某种我们不自知的深的东西,忽然进入共鸣,并向我们无声地诉说我们宁可不知道的事情。

首先正是梦,有能力在我们身上使那份必须保持隐藏、保持不被知晓、必须保持沉默的东西发出共鸣。也许唯有梦的语言,有能力触到我们身上这些隐秘的弦,并不由分说地使它们歌唱。而当你曾容它们歌唱片刻——哪怕是一首痛苦或沉重的悲歌——你忽然感到轻盈、仿佛焕然一新——像被大水洗过,仿佛一股丰沛的水流穿过你的存在,溶解并带走你身上一切纠结、坚硬而陈旧的东西……

当诗人准备拨动其中一根弦——它的歌会引发内心的洪水——他凭本能会借用梦的语言:既清澈又充满神秘——一种以图像与寓言言说的语言,以其表面的荒谬令理性错愕,又以其隐秘的分明径直抵达它要触动的地方!

这里无须“死”这个词出口,无须任何一个对醒着的理性暗示它的词。然而她在场;她那张迷雾的脸,正是挚爱者的脸。\origpage{514}那沉睡而遥远的挚爱者——你早已离她而去,却又近在咫尺——既是雪,又是落雪而降、又从雪中诞生的玫瑰……那把你吸向\textbf{她}的力量,宛如一股极深极强的波浪,一股来自那位召唤者的波浪、又把你带回她身边的波浪。而这召唤是切肤的悲伤,这回归是低声吟唱的喜悦;喜悦与悲伤本是一体,正是那股把你载向挚爱者的波浪——带着一场分娩的不容置辩之力的波浪。

而且无须提一个字,就已是那缠绵与渴望的冲动——你,孩子——被挚爱者唤入她体内的那个“小伙子”。只消一个梦讲到那位睡在父亲园子里、梦见落雪、醒来见玫瑰的她,你身上那股久被遗忘的波浪便也苏醒,应答着那位做梦又醒来、召唤又等待的她那缠绵的呼唤……

\subsubsection*{(c) 信使}
\tocsubsub{(c) 信使}
\label{ch18:local:h19}

\notehead{114$'$}{}
\phantomsection\label{ch18:local:n17}

这首古老的西里西亚歌谣,只是许多或老或新、歌咏挚爱者与死亡那神秘而切肤的交融的情歌中的一首。我刚刚抄录的这一首,或许多亏其含义饱满的意象之繁复、及其所激起的联想之丰富,而显得不同寻常。这里的主旨,不是把它们一一追索下去——在唤起其中最触动我的一两样之后。昨天与前天,当我的思绪回到那几节匆匆读过的歌时,并不是要深化那份起初还停留在表皮的情绪。倒是这份情绪,提请我注意到:“爱”与“死”的主题——或挚爱者与死亡的主题——显出何等紧密的联系,仿佛被某种神秘的咒语拴在一起!而且,越过那张挚爱者面孔的死亡主题,它们又接上诞生的主题——从睡—雪中觉醒—玫瑰;两者神秘地合一于那幅切肤的图景:玫瑰落雪,落在那位在父亲的园子里睡着、同时做梦又醒来的她的身上。

禁忌尽管把对死的排斥、死与生命及爱的不相容灌输给我们!但不得不相信:它逆着一份扎根极深的认识、或一股既强大又隐秘的冲动——以致那必须不惜一切代价分开的东西,竟以这样的顽强想要重聚,借符号与梦的迂回之路,穿过一代又一代、一个世纪又一个世纪传唱的歌与神话。

关于这些令人不安的交融,写就的渊博卷帙想必不少——为的是多少驱一驱邪。尽管有这些努力,肯定,在我们每个人身上的“某个地方”,这些顽强的联想的深意分明被知觉着——至少在那些我们没有刻意对身上的情绪关闭心门的时刻:它们接待那些信使,用梦那闪烁而有力的语言,向我们诉说着我们自己。

这份“深意”又一次向我们揭示——直接地、以一种基本元素般的力量——经由恋爱的经验:只要我们敢于充分地活出它、倾听它那分明的讯息。它便向我们诉说死亡与诞生的奥秘:二者在传递生命、更新恋人的那\textbf{行为}之中,密不可分。

无疑,我不是第一个:这份“扎根极深的认识”从它长期流放的幽暗深处上浮,变得全然有意识,并由此更加强烈地浸透我与死亡、与生命、与世界、与我自己的关系。然而我有这样的印象:见证这份有意识层面的认识的、已刊的书面文字,想必罕见。迄今我所知的仅有例子,是老子《道德经》里的三四节诗\fnote{66}{(*)}{(10月30日)我撞上《道德经》的这几段,是在1978年末。那是一份触目的、全然意外的确认:确认一些我感受强烈的事物(有些由来已久,有些新近才……),而且我似乎曾是唯一这样感受的人。这场“相遇”被经验为一种大的欢喜、一种无声的欢呼。这份欢喜、这份欢呼,承载了其后六七个月《乱伦颂》的孕育与写作。孕育,是在相遇后的几天或几星期里完成的。以一个更谦逊的音高,近日我又感到一份相似的欢喜:当我“认出”曾鼓舞一位匿名诗人(已死去几个世纪)的那份情绪——他歌唱那些落雪的玫瑰:荒诞地、奇迹般地从“纯然的空无”中诞生;或者不如说,经由我自己亲身的内在经验,重新找回那份\textbf{同一}的情绪——同一份认识的标志。人们也是在《道德经》里,跨越四千多年,重新找到这同一份认识——只是有一层分别:在中文文本里,这份认识以一个高度觉醒的意识的意象化、却绝非符号性的语言表达,而不是以梦的语言(那也是心灵深层的密码语言)表达。\hspace{2em}况且,我在这几节《道德经》诗里认出的内容,分明也逃过了我手头五六种不同译本(法、德、英)的译者们。我毫不惊讶。这类讯息——它们表达一种逆着千年制约的领会——只把其真正的意味(在用来表达它的词与图像之外)传递给这样的人:已经凭自己消化了的东西认识它的人;或那份消化的工作正在其中继续、且已然临近的人。}。

从另一方面(且有点悖论地),我同样有这样的印象:“爱——死”的交融,想必在某个时刻终于沦为一种浪漫的老生常谈、一块稳赚眼泪的“奶油馅饼”——专从最不情愿的眼里抠一滴讨好的泪。事实是:这套手法终于落得声名狼藉——(唉)到这种地步,以至于甚至在那些情感细腻的人们中间,有时也有一种倾向:把纯金与马口铁的粗劣赝品混为一谈。有人在这些地方看见过时的乃至可笑的调调,而那里分明有一份对一种隐藏实在的鲜活精微的知觉、一份远离一切“时尚”的精微表达。一个“好品味”的共识在此援助那些四面八方的内在抵抗:它们自动屏蔽一切鲜活而真实的情绪的闯入——无论是喜悦还是悲伤,是享受还是折磨——凡来搅扰熟悉的老一套者。也是同一机制,屡屡给爱的游戏的原初力量、及其高潮的完成设置壁垒。所幸,这些鼓舞性爱冲动的原型,单单保持被遮蔽、被放逐于意识领域之外,丝毫无碍于它们仍然在场——使该消失的消失、熄灭,以便爱之游戏的意味得以表达并完成,使最终的行为成为一个创造的行为、一次更新。\dmark{}但常常,也有一份隐秘的恐惧,给人们自以为在追求的那份“快乐”本身设置壁垒:人们被一股近在咫尺、既陌生又可畏的力量吓住了——它(倘不留神……)会像扫麦秸一样,扫走我们身上那位不惜一切代价要保住“控制权”的那一位。这样的恐惧,断不容许快乐靠近那个切肤的强度阈值——在那里,快乐同时又是折磨,二者在一场漫长而难忍的拥抱中结合,寻求解脱,以便最终消解、倾身于高潮的空无……\fnote{67}{(*)}{(10月28日)正是这份恐惧——显形为一种对快乐的拒斥——同时又推着人把快乐从恋爱经验的全体中孤立出来,把它化约为恋爱经验的全部,并奉为目的(有时心照不宣,有时明白宣示)。“爱”于是被化约为“寻欢”——说到底,是两个伴侣之间的一种礼尚往来,仿佛互邀去四星级的餐馆用餐,甚或去福利贝热歌舞厅(Folies Bergère)。这份惶恐地拴着皮带的“快乐”,与那股原初的冲动毫无共同之处,一如从大师亲手所绘的画上刮下的干漆屑之于那幅画;一如吹风机之于浩瀚的海风——那带着海与大地的芬芳的远风……}

(10月27日)我想我已经懂得了像《今天还在,明天……》这样的歌与梦的隐秘讯息——就它们所共有的本质而言。于是留下这个问题:究竟是怎样一股力量\dmark{}\origpage{517},如此执拗地推着人给这份“扎根极深的认识”——它想必比我们这个物种还古老——赋以声音;顶着那位乖戾而狭隘的检查官的警惕,逆着所有人把它表达出来,逃出家门、在梦那资源无穷的符号语言里自由驰骋?

倘若神话、歌与梦以无数的口吻不倦地向我们低语同一份讯息,那么同样真确的是:它们所诉说的那个囚徒,听它们也听不厌!诚然,那是一个自愿的囚徒;他小心不肯听。他被剥夺了空气、空间与光,却由四面墙安心地护着——那墙围起一种没有大意外、没有大奥秘的生存;若说有奥秘,或许只有尽头那场死亡,无限遥远……他的监狱保护他,避开墙外的那个未知——他佯装不知的未知。它既使他害怕,又使他着迷。正因为墙外的那个“彼岸”使他害怕,他的这座避难所式的监狱才对他比生命还珍贵。然而它又使他着迷、吸引他——违背他的本意;一如那些偶尔前来向他诉说彼岸的信使们吸引并迷住他。有时他向这份不寻常的吸引让步——只要能背着检查官——那位总监——行事:他一面不动声色地侧耳,一面却依旧处于“暂停”状态——他什么也没听见,尤其,什么也没听进去!

我方才向自己提出的问题,仿佛消失不见了——被一幅有说服力的图像藏掖了。一旦我记起那份讯息的效果——那种来到讯息之前的情绪、以及这份情绪的益处——问题便重新出现。

但老实说,任何触到一根深弦的情绪,都是四墙之外的彼岸的信使、是辽阔的信使。即便下一刻我们就努力抹去它的一切痕迹,它仍是有益的;它已经留下自己的痕迹,如一丝精微的香气——仿佛那些阴沉的墙挪开了哪怕一点点;或仿佛经由一道不曾料到的开口,有一缕哪怕极微弱的林野芬芳,透进了这消过毒的空气里。\origpage{518}

\dmark{}(10月28日)近十五天来,反思走进一个全然没有预料的方向——它与葬礼的主题没有明显的联系,甚至(似乎)与我自己的这个人也没有——这多少是违背我的本意的。我心底清楚:事情并非如此,我继续被卷入这些注中,不比以往少,甚至更多。这并不妨碍我在两种欲望之间摇摆:一边想“了结”,一边想开掘日复一日隐约瞥见的东西、追索那些最紧迫的联想——后者又与另一份关切合流:不放过任何能照亮我这场葬礼“侦查”的东西。显得最远的,有时恰恰也是最切身的……

总之,这十五天来——倘若不是自患病事故后重新提笔以来——我有一种(有时颇不好受的)印象:我在“赶工”、匆匆忙忙;仿佛每一条新注都是我又一道打开的括号(在一个会喊饶命的假想读者面前),而我理应尽快把它关上!肯定正是这些心境,或许更甚于最近几星期里相当多朋友在我这里的走动,造成了一份同样匆忙、有时略显杂乱的书写。我不得不随手把近来写的大部分注陆续誊清。这又拖慢了进度,吊着我那份眼看工作进展的急切胃口!

诚然,这些我有时佯装要趁热打铁、当作“众所周知”只为求个心安、并为一位刚刚“下船”的读者之益而费心显豁的主题——它们过于精微、分量过大,经不起如此漫不经心的态度。我忍不住在页页进程中觉察到这一点,并“拨正船头”——我是说,调整我的内在态度,顶着我本打算偷偷摸摸应付的东西的那份(不妨这样说的)重量!

这使我想起:这场我已投入将近四个星期、且远未终结的阴阳长篇反思,说到底,不过是在显豁一份瞬间的直觉——一份在我看来再单纯不过、甚至可以说明明白白的直觉;一份在5月12日的次日“闪”来的直觉——那时我刚写完关于某篇“葬礼悼词”的第一条注。一个月前\fnote{68}{(*)}{见“肌肉与肚肠(阳葬阴(1))”一注,n\textsuperscript{o}~106.}我接续那条注时,打算循着这一条联想(而非另几条在我看来趣味稍逊的联想)走下去;\dmark{}\origpage{519}我预计这顶多再花五六页。而那里,我已经越过了六十页的大关……

昨天,我在这个问题上停了笔:在爱与死之间、死与生之间、生命与死亡之间的联系的象征性唤起,其意味何在——以及这样一场唤起在我们身上激起的那份情绪,其意味又何在。在神话、歌与梦之中起作用的力量是什么,推着它们“以无数的口吻不倦地向我们低语同一份讯息”?而在我们身上——那些令人安心的监狱里的自愿囚徒——又是什么力量,屡屡以那份情绪回应它们,迎着那场唤起,表明它“命中了目标”、触到了它想触的地方?还有:梦的语言的那份奇特的力量从何而来——那种暗示而不点名的语言,传递着任何别的语言都不会传递的东西?

追索这些问题,也是更深地探测性爱冲动的作用——一如梦的作用——以及连接两者的深刻纽带:两者彼此滋养又被对方滋养,各自表达着自己,并经由一种它们共有的、逃避检查官的语言与对方交流。这同样是更深地探测原型与符号在性爱冲动中的作用、以及该冲动那些“符号性”满足的作用。

分明,这一切把我引向远远超出我所能合理指望“塞进”这场阴阳“枝蔓”的东西的限度——那场枝蔓(该提醒我自己了!),正延续在一场葬礼仪式的正中间!我觉得,是时候放下这根新“线”,回到三天前\fnote{69}{(*)}{见“失乐园”一注(n\textsuperscript{o}~116),它排在“本注”(注n\textsuperscript{o}~114)之后。}悬置的另一根“线”:当时它把我带回我自己的这个人。

\subsubsection*{(d) 安格拉——或告别与再会}
\tocsubsub{(d) 安格拉——或告别与再会}
\label{ch18:local:h20}

\notehead{115}{}
\phantomsection\label{ch18:local:n18}

\dmark{}\origpage{520}(10月30日)一两天以来,有几行诗在我脑际萦回,出自三年前写的一首诗。我先是用德文写成,次日又用法文重写。浮上心头的是头两节——第三节即末节仿佛从记忆里被抹去了,只除第一行“Ein~Kreis~schliesst~sich”——“一个圆环合拢”。(还有末行也一样,它重复了第一节的末行。)昨夜醒来时思绪又回到它;我终于起身,翻检纸堆。没费什么事就找回了那首诗——整齐归置毕竟有好处!它在此。

\begin{quote}\raggedright
果实沉沉\\
熟而重\\
我的生命俯身\\
为了归返\\
她之中\par\medskip
甜而稠的汁液\\
浸透了我\\
曾开花\\
脆弱的奶白的花\\
已成果实与酒\par\medskip
一个圆环合拢——\\
自我的怀抱\\
升起甘甜\\
画出它的环弧\\
又以弱音俯身\\
为了归返\\
她之中……
\end{quote}

我想,这是我写过的诗中唯一一首死亡之思\fnote{70}{(*)}{我更应当写:对我自己之死的思。1957年——我母亲去世那一年——写下的两首诗(各有几行),浸透着对这场死亡的预感。}分明在场的诗。它在这里以“她”之名现身。在前一晚的初稿里,它由德文词“Erde”——大地——唤起。顺便说,这三节的德文“译本”远非直译;第一节当初是这样来的:

\begin{quote}\raggedright
Voll und schwer\\
reife Frucht\\
neigt sich mein Leben\\
gen Ende\\
der Erde zu\par\medskip
Die süssen Säfte\\
die mich durchtränken\\
haben geblüht\\
weiche Blüten und wurden Frucht und Wein\par\medskip
Ein Kreis schliesst sich\\
aus meinem Schoss\\
steigt Süsse\\
kreist\\
und neigt sich\\
gen Ende\\
der Erde zu……
\end{quote}

\dmark{}\origpage{521}最终,刚刚把德文初稿重写一遍时,我竟不由自主地把它一气写到了底——那后面两节仿佛自发地从第一节里流出来!这三节,对我来说,是一首情诗(此外我写的诗,也几乎全是情诗)。倘若它致予的不是我自己,那便是致予她——致予那位在沉默中等待、准备迎接我的她……

同一天,我还写了两首别的诗,一在之前,一在之后。它们致予一位有血有肉的“挚爱者”:安格拉——“安琪儿”——一位高大、金发、苗条的姑娘,浑身洋溢着生命;一星期前,在那条夏天热浪蒸腾的路上遇见她,她正搭车。一两个钟头里,我们来得及彼此说了许多;然后就此别过。我本想把这些她所引发的诗送给她——包括在她被遇见的当天晚上写的另一首,以及还有一首(仍旧用德文,我们的共同语言):它诞生于那“三首(几乎)一挥而就”的次日。我也本愿我们相爱……但我失去了她的踪迹,一如她想必也失去了我的。

这场相遇所引发的诗的一个共同点:每一首都或者极强地“阳”,或者极强地“阴”。它们属于我写过的最炽烈的诗之列,每一首都一挥而就,几乎不加修改——仿佛它们早已齐备在那里,只等这场相遇的信号,好取得有形的言语之身。

乍看之下,在这些满载炽烈情欲张力的诗中间,竟有这另一首秋声阵阵、准备进入冬日长眠的诗,似乎奇怪。但只有感觉不到情欲的冲动与死亡之感之间那条深深纽带的人,才会诧异。在那些孤独的日子里,有一种对生命的炽烈知觉——由情欲、由支撑着它的原型意象的繁富所放大;同时,又有一种充分活过、正走向终点的生命所特有的宁静的超然:准备“归返她之中”。

与死亡——我们沉默的母亲、被感受为朋友而近在咫尺的她——契合的这样的心境,肯定被身体的一种大疲惫所促成:它把我们带回那些简单而本质的事物:我们的身体,爱,死……当时,我刚走出一场“数学的长期亢奋”——关于它,我已在《收获与播种》的导言里谈过\dmark{}\origpage{522}\fnote{71}{(*)}{见“梦与完成”,尤其第(iii)页。所论的“亢奋时期”,自1981年2月延续到6月;那也是“穿越Galois理论的长征”的时期(见“Galois的遗产”一节,n\textsuperscript{o}~7)。它通向一场关于我与数学之关系的漫长禅观(见“扫兴的老板——或高压锅”与“非古鲁的古鲁——或三足马”两节,n\textsuperscript{o}~43、45),自1981年7月19日延续到12月。献给安格拉的诗(以及献给“她”的那首),写于7月8日和9日(只有最早的一首作于7月1日)。}。我刚刚才开始从一场身体的衰竭中往上爬——那场颇有些疯狂的时期把我抛进的状态。它刚刚结束(来得与去时同样突然),借着一场碑文般有力的寓言梦的冲击;我当时肯倾听它的讯息\fnote{72}{(**)}{见注n\textsuperscript{o}~45的开头——前一条脚注已引用。}。那是可听可看的日子、倾听的日子——两波之间的一段“敏感时期”:身后,是一道长而宏阔的“数学”波;眼前,一道同样宏阔的“禅观”波已在酝酿……十来天后,它将开始它的冲刺——随着另一个梦:它的叙述开启《收获与播种》的导言,那幅“如我所是”的我自己的视野。

那是内在工作强烈的几个星期、无声的孕育、转变的几个星期。而这些情诗——音色不同于我此前写过的所有诗——正是这份强烈、这份丰满的果实与见证。

它们也是我写下的最后几首情诗。也许我身上有一种预感:这将是我最后一次恋爱,最后一次为挚爱者燃起歌唱的盛大焰火!一份预感:这些致予一位陌生姑娘的诗——我未曾相识却强烈地感到她的美——同时也是一场告别:向情歌、向我爱过的女人们告别;向我那正在这束光焰四射的花束中燃尽的爱的热情告别——它即将离我而去。而更隐秘、更深地:这也是一场告别(或许是一场再会):向着所有的女人——她们交相融合、在一个新的面孔下化为一个“她”。一张也许更遥远的面孔,淹没在雾中,在路的另一端——却同时非常近、非常温柔……

\subsection{(5) 拒受与接纳}
\label{ch18:local:h21}

\subsubsection*{(a) 失乐园}
\tocsubsub{(a) 失乐园}
\label{ch18:local:h22}

\notehead{116}{}
\phantomsection\label{ch18:local:n19}

\dmark{}\origpage{523}(10月25日)\fnote{73}{(*)}{(11月1日)本注早于前两条注;它们写于10月26至30日之间,并且是“行为”一注(n\textsuperscript{o}~113,10月21日)的直接延续与深化。本注倒是系于更早的10月17日那条注(n\textsuperscript{o}~112)的结尾,即“一半与全——或裂缝”。从这里起,反思遂分成两条平行的道路:一条关于死亡之感及其与爱欲冲动的联系,在依次列出的注113、114、115中继续;另一条则由本注n\textsuperscript{o}~116开启。}又过去了三天,我没有找到时间乘势写下去。第一天,星期一,主要被皮埃尔携女儿(两岁的)娜塔莉的来访占去;晚间很迟我才送他们去奥朗日赶夜车。再过几天,盘点这次来访带给我什么,还来得及——那是一次我已不再指望的来访……眼下,我更愿意断断续续地继续阴阳反思的线头。

这场反思,俨然一段哲学枝蔓,突然闯进某场它本无用武之地的侦查——只除了一点:它是从围绕某篇葬礼悼词的几缕联想中,不声不响地冒出来的……然而我分明感到:恰恰是借着这场“枝蔓”,我方才开始超越对构成葬礼的那些“原始事实”之总体的揭露阶段\fnote{74}{(**)}{(11月14日)这句顺口说出的话,并非深思熟虑,只有部分根据。至于一份更周详、更细化的概观,可参见“一场禅观的回顾——或一幅三联画”一注(n\textsuperscript{o}~127)。},终于稍稍逼近那些在种种似乎离奇出格的行为与举止背后起作用的力量……肯定并非偶然:正是经由这场“枝蔓”,我在自己毫无预料的情况下,被引至一种比《收获与播种》的任何其他时刻都更深的意义上牵连我自己的这个人的地步。这是最近的患病情节的一个意外果实——它恰好发生在我准备轻装把七周来的侦查带向其近在眼前的结论之时……

因此,这场“枝蔓”——有人会在其中看到一场亲身的忏悔,另一些人会看到一场玄思的形而上学——对我而言(比《收获与播种》的任何其他部分都更),处在葬礼的正中心、冲突的正中心。改变的只是眼光、看这件事的“视点”——而这一变如此苛烈,以致方才还在检视的那件事仿佛骤然消失!\origpage{524}我想,我们很快就会重新接上那看似半路丢失的联系——与“市井小事”即葬礼的联系。

不过,人们也可以把那桩市井小事忘了——它的主要功劳,毕竟是引发了这场“枝蔓”……昨天的一天,用来誊清前一条注的草稿——它写于四天前,我最终给它命名为“我们的母亲,死亡——或行为”。那份草稿有相当一部分涂改很重:一个征兆,表明表述还欠清晰;一些重要而精微的主题,当时是有点“从旁边”、在赶往别处的半路上混进反思的。老实说,起笔这条注时,我主要打算接上前一条注的线头——它名为“一半与全——或裂缝”,写于整整一周之前。可最终那根线头仍旧悬着;该是终于把它拾起来的时候了。

也是为那条注,我同样重誊了文本的相当一部分,缘由本质相同:沿途修正笨拙与晦涩。那是一场关于伴侣之内的分裂的反思的开端;它与个人内部的分裂血肉相连——更确切地说,与我(在四天前的“行为”一注里)所称的“竖直切分”血肉相连:它把我们身上原初之“全”的两“半”之一——阴或阳——“切开”,或削去。

在一个如今仍属于直觉领会、未形诸言语的层面上,我“懂得”、它对我“分明”:人本身内部的分裂(它俨然是全然由制约造出来的)是人类社会中无处不在的冲突的深层缘由。无论伴侣之内或家庭之内的冲突,还是更大群体内部的冲突、或把一些群体相互对立起来的冲突,直至各民族各国家之间的武装对抗。伴侣内的冲突——它把两位典型的、分明可辨的对手相互对立起来——满可以名正言顺地显为人类社会之内冲突的基本寓言、其不可再分的元素个案。“裂缝”那场反思的“要点”,主要是把伴侣内冲突的个案,归回到那个更根本、更“元素”的个案:每\dmark{}\origpage{525}个人自身内部的冲突——它使人自身的“一部分”与另一部分对立。

在七天前那场反思的眼光里,首先想到的,自然是我们身上“阴”与“阳”两“部分”之间的冲突——其中一方被接纳并妥善地抬举、吹胀,另一方被拒斥、被或多或少彻底地压抑。然而我心里明白,个人身上还有别的对抗,系于别的禁忌,而非性别只能有单一面向的禁忌。诚然,后一种禁忌与乱伦禁忌同样强,却更阴险,因为它带着一副明摆着的面貌——仿佛正因其明摆着,便免去了把它表述或命名的心思!尚未一步一步亲自检证,我便有这份印象(自《颂》的反思以来):这个禁忌是所有禁忌中最关键的一个;它在人身上设立的分裂或“切分”,是人身上积习已深的分裂之诸多侧面的终极根源。细心弄清事情在多大程度上确实如此,无疑会是一个最诱人的起点,足资开启一场“探索冲突之旅”。然而这里不是启程的地方——更不必说,在我面前摆着的、为我而设的那些旅程之中,我看到的起点有比这更灼热的……

誊清“一半与全——或裂缝”那条注的文本时,我还发觉:写它时,我不曾想到哪怕稍稍显豁一下,为什么我把个人内部的冲突视为伴侣冲突、社会冲突的深层缘由。这件事,我方才说过,属于那些我“懂得”过(却至今无须向自己“解释”过)的事情:它们由千百件日常小事那无声而雄辩的语言教给我、并向我证实,经年累月\fnote{75}{(*)}{这份“领悟”或信念,在我看来,并未被下面这桩我屡次得以作出的观察真正驳倒:由母亲与父亲结成的伴侣之内的分裂,以及表达这分裂的那些对抗姿态,会在孩子身上留下深深的印记,并常常支配着成年后的姿态与行为。至少在很大限度内,我们身上的分裂,是童年岁月里使我们母亲与父亲相互对立的那场分裂的印记与遗产——这样说无疑是站得住的。于是,要裁决个人内部的分裂究竟比伴侣之内的分裂更根本、更“元素”,抑或相反,这个问题有点像先有鸡还是先有蛋!
然而我确信:在一桩伴侣里,倘若夫妇中有一方是“一”、不与自己为敌,那么即便另一方对他怀着对抗的姿态,冲突也\textbf{不会}传递给这对夫妇的孩子。我相信这份确信的缘由是:在那种情形里,孩子会被他的父母之一\textbf{全然}接纳。幼童身上分裂的出现,在我看来不多不少,正是他的存在的一部分遭到周围环境\textbf{拒斥}的效应——而拒斥,首先是他的\textbf{双亲}。}。我并不是说:在这里花几页、或许卷帙浩繁的几册,来显豁或“解释”其“所以然”与“如何然”,毫无趣味。无疑,关于这个话题的几页,置于这些注中,并不比已在此安家的任何其他关于阴阳与冲突的篇幅更“不合适”。我肯定会在其中学到许多东西,一如继续那另一条反思的主题——我们身上被设立起来的阴阳冲突、作为我们内部分裂之终极缘由——也会教我许多。\dmark{}况且,这两个主题之一分明延续着\origpage{526}另一个,这使得它们各自愈发诱人!然而,眼下我毫无意愿朝那个方向走,哪怕一步。那不是我一周以来最想拾起、却始终悬着的“线”。

一周前,当我为那条注的反思收尾时\fnote{76}{(*)}{即“一半与全——或裂缝”一注(n\textsuperscript{o}~112)。},我忽然满心欢喜、精神一振:反思出人意料地重新接上了某样重要的东西——前几日我已有点失去视线的东西:接纳。这一接,是经由否定的路径恢复的,靠着那个为反思收尾、宛如一场出人意料的持续强音的词——“不可接受”。正是由于我们人格的整整一个“坡面”被周围环境当作“不可接受”而拒斥——首先由给我们定调的父母(或当父母缺位时,由代行其职者)——正是经由这份\textbf{不接纳},冲突才在我们身上安营扎寨。我们身上的冲突、分裂,不是别的,正是我们对自身一部分的退位——那一部分遭否认;是我们那未分裂的本性的退位。这份退位,是我们为被周围环境多少“接纳”而付出的、必须付出的代价。况且,那样的“接纳”并不是完整意义上的接纳——不是对我们真实所是的接纳。它宁可说是一份\textbf{报偿}\dmark{}\origpage{527}——报偿我们对某些\textbf{规范}的服从、报偿我们按规范自我塑形与浇铸——说到底,是一份\textbf{变形}、一份对我们存在的\textbf{戕害}的报偿;那变形与戕害,与我们的周围人自幼年所承受的如出一辙。

在前几条注的反思里,已有两次谈到接纳;两次,接纳都显为一件关键的事。第一次,在“天真(阴与阳的联姻)”一注(n\textsuperscript{o}~107)里:我在其中重提一份上溯到四年前一场禅观的确认:我身上一股未分裂之力得以绽放并充分舒展,竟发生在一个被冲突与潜伏的憎恨撕裂的家庭的语境里,\textbf{唯独因为我被我的父母、被我的环境全然接纳};冲突进驻我的存在,是在其后——五岁之后——在一个比我原生家庭“平静”得多的环境里。当然,亲人之间的冲突(至少在我那时候)远未达到我原生家庭那种加剧的(尽管蒙着面纱的)强度。然而,在那原生家庭里,我自己的这个人始终置身冲突之外。即便我偶尔表态站在哪一边,那也不是撕裂,而是一个未分裂的存在之自发的表达——一个从不曾尝过被自己人拒斥、以及对拒斥的恐惧之咬噬的存在。

如今,带着半个世纪的回顾,我明白:在我当时的新环境里,我身上这股天真的力量同样散发一种辐射,一种我愿称之为着魔的引力;宛如一座\textbf{失落的乐园}——无限遥远,令人一生怀恋——忽然借一个孩子的声音与目光向你发问。它当时为我赢得了强烈而持久的感情,追随我直到成年,直到那些如此爱过我的人一一谢世\fnote{77}{(*)}{我记得有七个人曾这样把他们的感情给了我,其中今日只有一位尚在人世。}。但\textbf{与此同时},理所当然:这股力量\textbf{不可能被容忍}——一如在一座修剪得笔直的观赏园子里,人们无法容忍某棵繁茂而恣肆的树或灌木:人们自以为爱它,却执拗地把它剪成方块、圆锥或球形……

按我对事件的复原\fnote{78}{(**)}{对我童年这些重要事件的复原,作于1980年3月。},这股力量坚持了或许两年、两年半,然后深潜下去,\dmark{}被贬入地下——那是在我终于决意“跟人人一样地做与是”之后:\origpage{528}十足十的肌肉与大脑(肚肠活该!),好落个清净!我终于步了后尘:凭着我周围所有成年人毫无破绽的共识,我\textbf{拒斥}并否认(以忽视的方式)了一切该被拒斥、该被忽视的东西。而且也凭着我父母自己的共识——他们到头来几乎不再给一点活着的音讯,在与他们的孩子尽可能远的地方过着他们的伟大的爱……

\subsubsection*{(b) 循环}
\tocsubsub{(b) 循环}
\label{ch18:local:h23}

\notehead{116$'$}{}
\phantomsection\label{ch18:local:n20}

(11月1日)我拾起恰好一周前中断的线头——10月26日,我曾出人意料地投身一场关于死亡之感在爱与情歌中的“诗的枝蔓”。

我刚重读了10月25日的前几页,并把其中最后一页誊清。我仿佛看见一个圆环正在合拢:它的轨迹两周前已经开始,随着“力量的绽放——或联姻”一注(n\textsuperscript{o}~107)。这轨迹在前面几页里收束:它们重拾并放大了10月17日那条注“一半与全——或裂缝”(n\textsuperscript{o}~112)的最后的“持续强音”。这个持续强音——或为当日反思收尾的“临终词”——归结为最后一个词的绝对律令:那个词,“不可接受”。

这个收尾的词,在形塑了我们生命的、形形色色令人错愕的大量制约之中,在我看来恰如其分地圈定了\textbf{那个}使我们身上发生分裂的决定性缘由:它就是我们生命最初年月里对我们的这个人的\textbf{不接纳}、\textbf{拒弃}\fnote{79}{(*)}{在这一点上,我本人的情况是个例外:直到六岁起,我才开始遭遇身边亲近之人这类姿态。}。它落实为对我们身上某些力量与冲动的失于接纳、予以拒斥——那些力量与冲动,是我们的存在、我们认识与创造之力的一个本质部分。对这些力量的压制,被一位忧心忡忡、毫不容情的内心检查官接了过来、据为己用,乃是对我们身上这种力量的一种戕害。其效应,往往无异于我们的创造官能一场真正的瘫痪\fnote{80}{(**)}{(11月2日)它往往还更显豁地表现为种种“卡壳”的效应:既无力在我们身陷的某种处境里“运转”,又无力从这毫无出路的处境中脱身……}。

\dmark{}这种不被接纳的力量,或这些“官能”,不是别的,正是那谦卑的能力:做我们自己。\origpage{529}这同时也是说:凭我们自身官能那份谦卑而完整的运用,活我们自己的生活,而不是过一种千篇一律、被编程的生活——那样一种生活,主要(且往往单单)靠\textbf{重复}、\textbf{模仿}的反射来推动。这些反射把我们封闭起来、隔离开来,宛如一副沉重的甲壳,僵硬而不透水,我们无论何时都无法从中脱出\fnote{81}{(*)}{睡眠与梦的时辰不计——那时甲壳会变轻:有时甚至会消失……}。

这副甲壳自我们幼年起便逐渐成形,年复一年,愈结愈厚。它最初的功能,无疑主要在于保护我们免遭亲近之人(往往出于好心)的侵凌,为我们从他们那里挣得一份多少还算善意的容忍。然而,这副甲壳不仅保护我们抵御外部世界——它还有另一层功能,或许更深、也更根本:使我们隔离,保护我们抵御\textbf{我们自己}——抵御我们身上那份知识与力量,它们被我们周围立法定规的那些无声共识宣布为“不可接受”、没有存在的理由。这在我们幼年便已如此,并年复一年愈发如此:一副\textbf{两面}的甲壳,一面朝“外”,一面朝“内”。它们保护着“我”、那位“东家”:一面抵御他惧怕的来自外部世界的侵凌(而且他一年比一年更胆怯!),另一面——也是\textbf{最要紧}的一面——抵御“工匠”那些扰人而不见容的奇想与乖张;说得直白些,抵御那个\textbf{野孩子}——难以逆料而又无所不能,纵然隔着三层厚厚的角质,耐火又防水、保证牢靠,他依然叫人心忧……

(11月2日)继“天真”一注(n\textsuperscript{o}~107)凸显出我最初几年被身边亲近之人接纳所起的作用之后,还有第二个时刻,“接纳”与“不接纳”处在反思的中心。那是在“接纳——阳中之阴”一注(n\textsuperscript{o}~110)里,我对自与孩童王“重逢”之日起在我身上发生的种种变化作了一次部分盘点。它们朝着一种逐步“回返”“童年状态”的方向。

\dmark{}这回返绝非“倒退”回某种先前状态——仿佛它有本事抹去我这旅人身上所留的我走过那条路的痕迹。\origpage{530}唯有凭\textbf{成熟}——一场内在工作的果实——我们才能重新接上那仿佛已经消失的天真,重新接上我们身上那个仿佛早已死去、早已下葬的孩子。而凡有成熟,无不同时也是些许的回返——回返孩子,回返质朴,回返孩子的天真。正因如此,一种活得完整的一生,宛如又一个正在“臻于圆满”的圆环;那是老境重新找回童年,是成熟重新找回天真——并或许终结于一场死亡,而这死亡准备着一次新生,一如冬天准备着一个新的春天……

在这份对一条尚未走完的回返之路的“盘点”里,显出:那“临终词”曾是\textbf{接纳},一如我的决裂之路、出发之路的临终词曾是\textbf{不接纳}、拒斥、拒绝。我的成熟不是别的,正是这样一个过程、这样一场内在的工作:藉着它,我逐步接纳了、迎纳了那些我长年以来拒绝、尽力清除、置之不理的东西。

这绝非一次“折返”,不是把一条已经走过的路再反着重走一遍;因而不是“倒退”——借用方才的说法。它毋宁像一个圆的上弧:延伸并接续那条已经画出的下弧,\textbf{诞生}于下弧,成了它滋养的根基、一次新跃的跳板……

(11月3日)昨天的注,收束在一幅意想不到的意象上——它不召自来,从反思中涌出。起初我带着几分迟疑接待它,生怕这意象旋即暗示出的那幅现实图景失于造作;生怕这意象“逼我就范”,逼我说出“牵强附会”的话。可一旦最后几行落笔,我在其上驻留片刻,便知道:我刚刚把手指按在了某一现实的意想不到而重要的侧面之上;这侧面也许我本已认识,却未充分消化——一个我素来倾向于忽略或遗忘的侧面。

\dmark{}多年来(118),我倾向于抬举朝着“接纳”方向去的东西,\origpage{531}而相反,主要以负面的眼光看待朝着“拒绝”方向去的东西。尽管这事情未必总是被清晰地表述出来,我却把这两类姿态——接纳与拒绝——感受为一对“对立”、一对“相反”:其中一个于我自己、于所有人都是“好的”,另一个则是“坏的”。

在这份未曾成形于言语的领会方式里,我始终是那套老生常谈的“二元论”世界观的囚徒(当然,我自己并无察觉)——那也就是我先前称之为“战争式”的世界观:它把一些事物当作对手对立起来,而一种更深的视域却向我们显明,它们是同一现实的\textbf{互补而不可分的}两面。当(10月25日,即十天之前)我开始眼下这场关于接纳与拒绝的反思时,我刚意识到:它们正是这一个月来——自这场关于阴阳的“枝蔓”开端以来——不断谈起的那些著名的阴阳“对子”、或“宇宙”对子之中一对里的妻与夫。因此我预计,反思将转向事物的这个侧面。两天来,它看上去似乎正在远离这个侧面。可昨日反思结尾的那几行——同一个循环的两道弧相互延续的那幅意象——出人意料地把我带回了那最初的直觉:它一直未曾形诸言语。

我倾向于以(即便不是全然)主要\textbf{负面}的眼光,看待从我八岁到四十八岁支配了我一生的那些\textbf{拒绝}:视之为一副有时令人不堪重负的\textbf{重担},我在我生命的四十年里一直拖着它,而在逝去的八年里,我终于甩掉了它(毋宁说,\textbf{开始}甩掉它)。这“重担”的那一面,是在禅观的发现之后、在与我身上的“孩子”“重逢”之后,才开始向我显露的。恰是在那一刻,我开始发现我生活中拒绝的进程——它表现为一种“超级阳”的循规蹈矩。事物的这个侧面绝非虚构。在先前只有一片“空白”、一片全然虚空之处知觉到它,是这八年里持续进行的成熟所结的果实之一。但与此同时,同一现实还有另一个侧面,同样真实、同样重要:它作为“\textbf{强大的行动原则}”的“正面”。这个侧面头一次出现(而且十分含蓄),是在10月5日那场禅观“阳葬阴——或超级父亲”(n\textsuperscript{o}~108)里,当时我写道:

\begin{quote}
\dmark{}“‘我将像他们一样’(而不是‘像我自己一样’)也意味着:我要把宝押在‘头脑’上——它在我身上并不比在任何人身上更差——\origpage{532}用他们自己的武器打败他们!”
\end{quote}

正是这股动力,宛如我1945至1969年间在数学中那场无度投入的活力——那股在四分之一世纪里滋养了一场发现之冲劲的力量\fnote{82}{(*)}{更确切地说,是这股冲劲中利己的成分,是这一“活力”的\textbf{利己}“因子”。}。无论人们选择以“正面”还是“负面”的眼光看待这样一场投入,清楚的是:那里确确实实有过\textbf{冲劲}、有过强烈的\textbf{行动}。在向生活学习的一侧,有着那副“有时令人不堪重负的重担”,从未被检视过——且不必说全然的停滞;而这同一副“重担”,恰恰在同一个时刻,却又滋养着一股认识的冲劲,赋予它活力。

自1970年我“出发”以来,我倾向于贬低、有时甚至否认这样一股冲劲——朝着对所谓外部世界作“科学”发现与理解的方向——所应得的“价值”。在《收获与播种》里,我已有好几次尝试去圈定这样一种发现与自我发现之间的共同侧面,以及反过来说,它们又在何处不同\fnote{83}{(**)}{尤其参见“欲望与禅观”“禁果”“孤独的历险”诸节(n\textsuperscript{o}~36、46、47)。}。这样说无疑是有根据的:朝着科学方向(哪怕是生物学,或“心理学”……)的发现冲劲,使我们远离我们自己、远离对自我的领会。当这样一份领会的角色被充分理解时,人们或许因此会倾向于在科学发现的冲劲中(以及在任何使我们“远离我们自己”的冲劲中)看出一种“恶”,或至少,看出成熟的一个“障碍”——并由此,看出我们自身充分绽放的一个障碍。(至少在我长时期所处的那种情形里:这股冲劲动员了绝大部分、乃至全部的心理能量。)然而,同样真实的是:我们所经历的\textbf{一切},都是我们学习生活、学习自我的原料。\textbf{材料}就在这里,任由我们让它转化为认识——只要容许一场成熟的工作在我们身上发端并继续。这也是为什么,我对自己经历过的一切并不懊悔,最终看到“其中一切都是好的,没有什么可抛弃”;也包括那些漫长的灵性停滞时期的荒漠——那是我为一场吞噬性的热情作无度投入所付的代价,我付出时不曾吝惜(而且闭着眼睛……)\dmark{}\origpage{533}。如今我看到,连这些荒漠也曾有东西要教给我,或许唯有它们才能教给我。这份代价我本无法省去——顶多,或许早在几年之后,我就本可以启动这循环的“第二道弧”——它的期限,我曾拖延了几十年。

也是在这一天,同样显出:对我自己、对他人的接纳——它诞生并成长于我成年的岁月——曾“滋养”于那些拒绝,它们标示了我一生最长的那段时光——即昨天唤起的那道循环的“下弧”,及其“滋养的根基”。诚然,在我生命的最初六年,我身上确有一种对自我的全然接纳,它无需任何先行的“拒绝”便能存在、舒展、确立。恰恰相反,它的绽放之所以可能,\textbf{正}因为它不曾被对抗、不曾被某一把拒绝的剪刀裁剪。但我童年身上的这份“接纳”,与我成年时的那份,不是“\textbf{同一份}”。它欠缺一个维度——那份仅凭我童年周围之人对我这个人的接纳所无法给予的维度。那是一份对拒绝的认识——对被他人、或被我自己拒斥(或我的某一部分被拒斥)的认识。这份认识,是经由拒绝的经验来到我这里的,也是经由轻蔑的经验——它是拒绝的众多面孔之一。

或许有些人生来就带着对拒绝的一种认识、一种领会,使他们能够在童年所暴露的种种拒绝面前依然保持\textbf{一}、依然天真而有知。我清楚地知道,我不是这样的人。他人的、和我自己的拒绝与轻蔑的经验,作为一片土壤,孕育出对拒绝与轻蔑的领会(无论多么不完善)——这份经验我无法省去。

\subsubsection*{(c) 伴侣——或“恶”之谜}
\tocsubsub{(c) 伴侣——或“恶”之谜}
\label{ch18:local:h24}

\notehead{117}{}
\phantomsection\label{ch18:local:n21}

我刚刚探测了我自己生活里拒绝与接纳之间关系的一个意想不到的侧面——它在昨天的反思中不期而至。不过,这里所说的“拒绝”并不是完整意义上的拒绝;我是说,一份被完全承担下来的拒绝——还差得远。这份拒绝也曾是一场对被拒之物的\textbf{长久逃遁}。它在于\textbf{不去看它}、\textbf{把它忽视},并由此在某种限度内,使它从我有意识的领会之场消失,也\dmark{}从他人可见的场中消失。\origpage{534}它曾是一种失谐、失衡状态的缘由与发条——在这里,是一种“超级阳”的失衡:它标示了我的成年,而它的若干关键机制至今仍在起作用。所以,这“拒绝”在此绝非以对称、甚或阴阳互补的角色,显于方才所谈的那种“接纳”(对我自己与对他人的)对面。相反,后者落在一场认识我自己的工作之中,朝着恢复一份被扰乱的和谐行进。所以,那里的接纳,是“充分知其所以然”的接纳,是完整意义上的接纳——而绝不是与方才名为“拒绝”的那场逃遁方向相反的又一场逃遁。

然而,“拒绝”与“接纳”之间,有一层比方才所探测的更明显的关系。它显现在两者都被“按完整意义”来取的时候。那时,它们是同一份和谐、同一种被完全承担下来的态度的\textbf{同时}而互补的两面。(而方才谈的,则是一段路程、一种进程的\textbf{先后}两面:途经一种失衡、失谐的状态,走向一份更新的平衡。)在这个视域里,没有那种把拒绝排除在外、对拒绝闭门的“真”接纳。也没有那种不诞生自接纳、不是接纳之一份可触的显现的“真”拒绝;没有那种不是同一件不可分之物的两“面”之一——阳的那一“面”——的拒绝:那物本有两面,其“阴”的或“母”的一面,便是接纳\fnote{84}{(*)}{有趣的是可注意到:接纳-拒绝这对偶子里,阴阳角色的这种“自然”分配(在法文里,由这对偶子两项一阴一阳的性属来表达)——在前一天反思结尾自发浮现在我心中的那幅意象里,却\textbf{颠倒}了。会有这样的互倒,没什么可奇怪的——正如在一对情人与情郎之中,其爱的关系并非凝定,爱的游戏里必不缺少角色翻转的时刻:好让活在情人身上的“阳”性爱欲冲动、与活在情郎身上的“阴”性爱欲冲动,畅行无碍。何况,我在“接纳(阳中之阴)”一注(n\textsuperscript{o}~110,该注第一部分末段)里,正谈到这种偶发的角色互换之重要。}。

一份把拒绝排除在外的“接纳”,不是接纳,而是一份讨好(对他人,或对自己,或兼而有之);或者——当涉及对他人的“接纳”时——一份串通、一份默契。全然接纳一个存在,无论这存在是自己还是他人,绝不意味着对其一言一行、种种习惯与癖好,作\dmark{}无条件的赞许。\origpage{535}这样一种无条件赞许,本身便是一场\textbf{逃遁},是对一份(往往雄辩的)现实拒绝去认识,而绝不是接纳。它非但造不出一个有利于更新、有利于与一份被遗忘的太一重新接上联系的“力场”,反而强化一种惰性,助长把人滞留在车辙里。

一个不同时是一次敞开的拒绝——它不同时也像一只向他人伸出的手(或“一根撑竿”),或像一个在自己与自己的关系中标出断裂与更新之点的跃起——这样的“拒绝”,真正是一场切断:它既“切断”又隔离拒绝者与被拒者。这仍是逃遁——逃开一份令人不快、甚或扰人的现实,一份沉沉压在我们安稳生活、我们种种便利之上的威胁——对这份现实,我们自信一刀便能躲开:“这没有存在的理由……”然而,\textbf{它有}!而且我们那道命令式的“拒绝”,丝毫不能拦阻事物如其所是,哪怕它们惹我们不快。恰恰相反,一如一场自动赞许的讨好,这样的拒绝强化着抵制创造之更新的种种惰性;它宛如一道\textbf{判决}:你是不可接受的,而且永远如此……

我并不自许在我这个人身上实现那份被完全承担下来的接纳与拒绝之和谐。恰恰相反,我知道远非如此——而我不敢肯定曾遇见过一位实现了这份和谐的存在。实现它,也就等于在你自己这个人身上,解开了“恶”的大谜:不义、谎言、恶毒、卑怯、轻蔑——以及那些被击中而喑哑无声的人们的苦难。这肯定也等于:充分领悟那份“善”——它就在那内在的跃起如此经常地给我们指为“恶”的东西之中。

拒绝战争,同时看到并接纳战争无处不在、在所有人身上;看到连我所爱的那些人也在他们身上携着它、传布着它,一如我自己也曾重拾它、携带它、传布它、传递它。拒绝战争,同时接纳它的存在,同时爱它那数不清的盲目的士兵。确定无疑地,唯有这个,别无其他,才也意味着:走出战争,走出冲突——停止传布战争。

\subsubsection*{(d) 阳扮阴——或主人的角色}
\tocsubsub{(d) 阳扮阴——或主人的角色}
\label{ch18:local:h25}

\notehead{118}{}
\phantomsection\label{ch18:local:n22}

\dmark{}(11月4日)\fnote{85}{(*)}{本注由“循环”一注(n\textsuperscript{o}~116)的一条脚注生出。参见11月3日诸注开头处的参见号。}这份“倾向”的出现\fnote{86}{(**)}{即抬举“接纳”、以与“拒绝”相对的那份倾向。},是在七十年代初——也就是说,在我从数学舞台“出发”之后的那些年头。\origpage{536}在一个与从前大不相同的氛围与朋友圈的影响下,我当时所信奉的那一整套“价值”,发生了一个严苛的转弯。带着回顾,我可以把这个转弯描述为:从一个“超级阳”的或“父权”的价值系统,过渡到一个几乎相反的、“阴”占强势主导的系统——一个“母权”系统。在这场颠倒中起过作用的影响里,还有对Krishnamurti的几次零星阅读——关于这一点,参见“Krishnamurti——或变成羁绊的解脱”一注(n\textsuperscript{o}~41)。

我当时之所以放任这些影响起作用——它们把我引向这样一场“意识形态”的转弯——无疑是因为(尽管我当时并未察觉)在我身上有一种深刻而紧迫的更新之需;而首先,是一种解脱的需要:从那份积习已深的“超级阳”姿态的重负下解脱出来。这同一份需要,无疑在1969年已经起过作用:那一年,正当数学活动强烈而多产之时,我忽然从数学上“脱钩”,转而关心生物学\fnote{87}{(***)}{我起初是从“分子生物学”这一头对它发生兴趣的,是受我的生物学家朋友Mircea~Dumitrescu的影响——是他把我领进这个迷人的世界。};继而在次年,离开(不存回头之念的)数学舞台,乃至科学研究。当时发生了一场突然而严苛的氛围与活动的改变,对此我在《愚妄与更新》(《收获与播种》的第一部分)里已屡有机会提及。

然而,把这些环境、活动乃至“价值”的引人瞩目的改变,看作一场“更新”、一场“解脱”,会是不确切的,或只部分为真。关于这一点,在“与Claude~Chevalley相遇——或自由与好心情”一节(n\textsuperscript{o}~11)里,我已表白得相当清楚。在眼下这场关于阴阳的反思的更透彻的光照下,我可以说:那个无疑显得在一切之中最有意义的改变——阳的价值被清除出去(甚至还没在我自己身上被辨认、更未被检视),让位给阴的价值——却丝毫没有改变“我”的(超级阳的)结构,至多稍稍缓和了由之而来的种种姿态与行为。诚然,我对外部世界的领会\dmark{}已发生相当大的转变,朝着一种突然拓宽的方向;但这一转变依然是零散的,几乎排他地限于智性的层面、“抉择”的层面。\origpage{537}事情不可能别样,只要这场转变恰恰只限于我的“外部世界”之景——在那景中,我自己的这个人并不出现,或只是偶然地、浮面地出现,主要经由我的“社会角色”及其种种含混与矛盾。与过去一样,我当时毫无些许的疑心:\textbf{我自己}身上也会有含混与矛盾!恰恰相反,我被一种不可动摇的确信驱动着:\textbf{我}这个人,是无矛盾的(尽管我当时已开始在四周处处辨认出他人身上的矛盾);尤其确信:在我有意识的欲望与有意识的事物之知、与我的无意识之间(倘若我的情形里竟有一个无意识的话——倘若它不是我意识的一件如实复制品的话……),有着完美的契合。

这份确信的第一道裂缝,直到1974年春才出现:我终于懂得,某种东西在我\textbf{自己}身上也必定出了毛病,而不只在别人身上——它是我与所有亲近之人的关系那场无情恶化的缘由(整个成年期间,我的生活似乎就缩小到了这关系之上)。当时,由于缺少一份对我自己的真正的\textbf{好奇}——它本会以钻进去为乐,看看背后有什么,并由此看一座沉重的建筑坍塌下来:一座由从未检视过的、天花乱坠的幻觉筑成的建筑——这道益人的裂缝的效应,便仍然有限。

这份对一种自然好奇的顽固阻塞,无疑主要缘于:我从未在他人身上遇见过这样的好奇——它本可让我起疑:在生活中,一如在数学里,凡有一个问题呈现,就有可看的东西,而看着,便能学到许多意想不到、极其有用的东西——换句话说:世上有\textbf{自我发现}这样一件事。我当时读过一些Krishnamurti,并得以确认:他说的某些东西是真的、深的、要紧的。于是我倾向于对他全盘照单收下。除枝节而外,我已在事实上默默采纳了Krishnamurti的世界观\fnote{88}{(*)}{(11月5日)这样一份“采纳”——它变成一种文化行囊——在我生活中的效应,始终极其有限。我的注意力被引向现实的某些此前完全逃过我知觉的侧面,但并未因此启动一场具有更新之力、深入的清理与消化的工作。倘若说,在1970与1976年之间(在我的数学舞台“出发”与禅观的发现之间),Krishnamurti在我的行程里曾是重要的,那与其说是因为我向他借来的那份“行囊”,远不如说是因为他已变成(当然,是在我不自知的情况下)一个默默的\textbf{模型}——我遵从着它,又不肯露出遵从的样子:总而言之,“非古鲁的古鲁”的模型,那位自称不是大师的大师。}。\dmark{}在我所说的那个时刻,这份行囊确确实实起了一场“羁绊”的作用——羁绊着一场真正的解脱、一场完整意义上的更新。\origpage{538}关于这一点,我恰好在那条已援引的注里作了说明(我刚刚重读过),我在其中力图圈定(Krishnamurti的)“教诲”在我自己的行程里所起的角色。

完整意义上的第一次“觉醒”,只在两年半以后才到来,随着禅观的发现。那也是对自我发现的发现:\textbf{有一个未知之物},它就是“我”,而我有力能穿透这物、认识它。这一关键的发现,是在一切教诲(无论有无大写)都已被遗忘的时刻完成的。那也正是“建筑”头一次坍塌的时刻——它由种种现成观念与形形色色的“教诲”筑成,靠一份巨大的惰性维持着;也是一份主动的好奇头一次出现的时刻——它往往顽皮,而始终满怀善意。

正是在这个转折之后,随着一份好奇在我身上绽放——首先朝向我自己的这个人,继而是“生活”,作为自然的果实——我才得以用一双新的眼睛,同时看Krishnamurti与他的讯息。带着回顾,我得以赏识那份讯息的丰富,同时辨认出它的限度与匮乏,以及那位\textbf{主人}(“the Teacher”,他的信徒与追随者如此称他)身上某些根本的矛盾。这些匮乏与矛盾中最重的一个,在我看来,正是我方才又一次擦过的那个:主人自身之中全然没有好奇。他的著述里没有任何东西让人起疑:在遥远的岁月里,这份景是否曾\textbf{诞生}于一个\textbf{人}——一个如你我一样陷在现成观念之网与从未被辨认的矛盾之中的人;这份景是否曾在一场强烈的、有时痛苦的、逆着巨大惰性之流的\textbf{工作}进程中,\textbf{从错误里澄清}出来;这项工作的诸阶段,或这些劳作中跨过的“门槛”,是否是一场场意想不到的\textbf{发现},每一场都撼动一整套积习已深的、\dmark{}由模仿与重复的\origpage{539}普遍机制所延续的观念\fnote{89}{(*)}{(11月5日)这些机制,分明是心灵的基本机制的一部分,在人与动物身上皆然。它们先于一切制约、一切学习(如幼童的语言学习,以及日常生活里几乎全部行为的学习)而在——没有它们,那些学习根本无从确立、无从展开。而在那位年轻的未来主人身上,它们并不比在任何人身上少在、少奏效。}。

所有这些事,那个孩子有一天都曾懂得,甚至曾认识——因为曾强烈地经历过。但主人把它们忘了,而且无意记起。他不愿做那个激情地\textbf{发现}、\textbf{学习}、并在发现中被转化的孩子;他要做一个一成不变的\textbf{主人}——凭天赋之知而不移地\textbf{知}着,并把一生献给传播他的\textbf{教诲},以造福凡俗大众。

他变成了他的追随者与信徒、那些信他的人要他成为的那一位:一份静态的、重复的、因而是抚慰人心的讯息之化身,一种新意识形态的使徒。总而言之,一位\textbf{非古鲁的古鲁}——一如我自己当年也曾是(或许正是仿效他的榜样\fnote{90}{(**)}{(11月5日)说真的,这个“或许”所带的存疑口气,根本不合宜!关于这一点,参见今天写下的倒数第二条脚注。})……

(11月15日)我把前一条注(11月4日)命名为“阳扮阴——或主人”。正如一场对我自己的禅观所应有的做法,这条注的主名关涉我自己的这个人:它指的是某一场“戏”——自1970年我离开科学世界之后\fnote{91}{(***)}{此外,1976年10月禅观发现的时刻,也标志着这场戏的陡然衰落;它好歹继续着,只是调门更低调,直到1981年——那一年,它终于被察觉、被拆除引信。关于这一点,参见已援引的“非古鲁的古鲁——或三足马”一节(n\textsuperscript{o}~45)。},我曾扮演了几年。至于第二个名字“主人”,它可以无差别地被解读为:或指涉我自己的这个人——指我在“阳扮阴”这场戏里所扮的角色、所摆的姿态;或指Krishnamurti——他曾默默充当我的模型。

实际上,从Krishnamurti的书里透出的价值,几乎是排他地阴的价值。在我头一次读Krishnamurti时(1970或1971年),我是头一次看到这样的价值被摆到前头来,并看到我自己那阳性的世界观(以及“人人”的世界观,出入无几)的限度与裂隙被透彻地圈定。这无疑就是那几章的阅读在我身上留下极深印象的缘由。六七年之后,我又有机会读到吕滕斯夫人写的那部Krishnamurti的漂亮传记。它印证了从他的书里已然透出的对他这个人的某种印象(尽管他在书里从不亲自出现)。今天我会这样表达:他气质里的基调,是\dmark{}强阴的\origpage{540}。

还须补上:透过他的全部著述,人们看到一条恒常的主导动机:把带阴色的品质、姿态与价值摆到前头,而(明言地或以沉默)贬低阳调的品质、姿态与价值。

所以,Krishnamurti的一生与教诲,实现了“\textbf{阴葬阳}”这一相当例外的心态——它与迄今为止最常见的那种心态、“阳葬阴”的心态正好逆行;而我自己的生活(至少直到我四十八岁)恰是后者的一份同样极端的例证。Krishnamurti的“超级阴”抉择\fnote{92}{(*)}{这些“抉择”无疑上溯到他的童年,更确切地说,上溯到他与他那些神智学监护人的最初接触。}有一个大功绩:逆着周边文化的基本价值而行。但这并不妨碍它们在我看来同样具有压制性(以他这个人的一部分压另一部分),一如我曾经的那些抉择。

然而,Krishnamurti的一生里,有一个非常触目而惹眼的“阳”的侧面——它无疑首先是被强加给他的:是他那些显赫的神智学监护人在他尚是孩童时,就为他定下的那号(未来的)“精神主人”、旗舰人物的角色。其后,在他生活的大转折——以那些翻天覆地、彻底撼动其世界之景的发现为标志(这些发现在其后成了“教诲”)——之后,这个“主人”或“引路人”的角色,似乎被他整个地内化了,被他自己接了过来,连同一份他个人(而非承自其神智学老师)的学说之传布。这份传布代表一种强烈乃至耗竭人的活动。它看上去很难朝着阴阳的一份\textbf{平衡}去,毋宁显为一份\textbf{强制}:由一个“我”强加给一种彻头彻尾禅观性的气质——这“我”在主人身上,与其在任何人身上一样强、一样具有侵夺性。在这个光照下看,眼下这条主要谈Krishnamurti的注“阳扮阴”,同样可以叫作“\textbf{阴扮阳}”。

这样,在我的一生里,有两次、以两种不同的方式,我玩过一些“戏”——它们恰似一种\textbf{颠倒}:颠倒那支配了某个人一生的种种姿态——那个人,在我行程的某段时期,本该成为我(同样默默的)形象招牌的默默模型;也颠倒我身上与之相应的某些姿态与架势。但透过这两种彼此相反的表达风格,我今天认出一种显然的亲缘。其一,在于一份\textbf{压制}的存在(无意识的,不言而喻),它\dmark{}生成阴阳自然平衡的\origpage{541}破裂\fnote{93}{(*)}{在这一份亲缘里,我们想必不乏人数众多的同道!}。其二,在于对一个\textbf{角色}的选择,在于\textbf{这角色之重}——它的刹车效应,乃至阻塞效应,落在一场绽放、一场成熟、一份领会或认识的进程之上。这个角色(或这份架势),在我身上与在我那位模型身上,是同一个——我也许不过是把它原封借来。那就是\textbf{主人}的角色。

\subsection{数学的阴与阳}
\label{ch18:local:h26}

\subsubsection*{(a) 最“大男子气”的艺术}
\tocsubsub{(a) 最“大男子气”的艺术}
\label{ch18:local:h27}

\notehead{119}{}
\phantomsection\label{ch18:local:n23}

(11月5日)想谈谈数学里的阴与阳,已有一阵子了。一项数学工作之中、或一种通向数学的进路之中的阴阳两面,只是在这几周关于阴阳的反思进程中,才向我显现。我预计:在这些注里对这个双面稍作探测,将是“言归正传”的最自然的方式——这些注,本应对“一个数学家的往昔”构成一场回顾。

从一开始(五年前)我关于阴阳的最初反思起,对我就已十分清楚的是:“做数学”,或许是已知的一切人类活动之中\textbf{最阳}、最“男性”的一种。老实说,一切全然智性的活动——尤其科学研究,以及更一般地,一切通常被称作“研究”的活动——都是一种阳占绝对优势的活动。我本要写:“带有强烈的阳失衡”,而事实上,当这种活动竟至吸收一个人的几乎全部能量时,情形正是如此。这种阳的优势(或失衡),经由列举许多阴阳对子而显现:对这些对子,分明是阳的那一项——若不说排他地——才“在”一项智性工作之中“在场”。我只限于举出其中几个;它们全都属于同一个“组”(或同一扇“通向世界的门”),我称之为“模糊——精确”组。(按:在这个对子以及接下来的那些对子里,是阴的一项放在前头。)

\begin{itemize}
\setlength{\itemsep}{0pt}
\setlength{\parsep}{0pt}
\item 感性——理性(或智性)
\item 本能——反思
\item 直觉——逻辑
\item 灵感——方法
\item \dmark{}视野——融贯\origpage{542}
\item 具体——抽象
\item 复杂——简单
\item 模糊——精确
\item 梦——现实
\item 未定——已定
\item 未言——已言
\item 无形——有形
\item 无限——有限
\item 无界——有界
\item 全(整体)——部分
\item 全局——局部(或细部)。
\end{itemize}

我刚把我的阴阳清单巡了一遍,又捡出一大批对子,它们都让人尝出纯粹智性活动那份超级阳的品格。只再举方才已想到的那些之中居首的一个: \textbf{身——心}对子。

有鉴于此,我觉得,在各种智性活动之中,代表着终极极阳的,是数学工作。这无疑首先缘于它极端抽象的品格:它在极大的限度内,不依赖任何借助感官经验、借助对外部世界——即我们生活于其中、我们身体活动于其中的那个世界——的理性观察的“支撑”。这种抽象中的极端品格,把数学同一切其他科学区分开来,把数学工作同一切其他智性工作区分开来,使它成为一门“纯理性”的科学、一份“纯理性”的工作。与实验科学和观测科学相对,它也是唯一一门其结果经由最严格意义上的\textbf{证明}来确立的科学:证明循着一条严格编码、原则上不出错的\textbf{方法}——所谓“\textbf{逻辑}”方法——行进,达到不容任何怀疑或保留的\textbf{确定性},不容许有逃过迄今所察案例的例外。这些都是在数学工作中、且唯独在数学工作中汇齐的极阳特征。

当然,这些特征自童年起就自然而然地吸引着我——我这个从头到尾选定了“头脑”、选定了极阳的人!\fnote{94}{(*)}{军事与战争的那个变体且不算:检阅、制服、挺着胸膛的立正,以及组织得无懈可击的屠杀与尸坑……}尤其在经历了战争与集中营的考验、直面种种似乎向最起码的理性也发起挑战的歧视与偏见之后,数学活动中(经由我中学年头所能认识的那么一点\dmark{})最让我着迷的,是它所给予的那份\textbf{权能}:凭一个简单的证明,赢取哪怕最不情愿的赞同;总而言之,\textbf{强使}他人首肯,无论他乐意与否——只须他同我一道接受数学的“游戏规则”。这些规则,早自我1940年在芒德中学最初接触学校数学之时(我得以去那里上学,尽管被拘禁在离那儿五六公里的里克罗营地),我就仿佛本来就认识它们、本能地感到它们,好像我从来就认识它们似的\fnote{95}{(*)}{这些最初的接触,就发生在我童年的“化圆为方”思考之后不久——见注~n\textsuperscript{o}~69。}。\origpage{543}当然,我比老师本人把它们感到更好:他毫无信念地背书给我们当时通行的那些老生常谈——“公设”(亦即欧几里得公设,唯一一条他同我们都有幸听说过……)与“公理”的分别,或三角形全等“三情形”的“\textbf{证明}”——一路照着教科书念,活像一个初领圣体的学生跟着念他的日课经。

五年之后,我被原子物理骤然显赫的名望所吸引,却先在蒙彼利埃大学注册学习物理,心里想着去入门物质结构与能量性质的奥秘。但我很快就明白:倘若我想入门奥秘,靠听课是不行的,只能靠自己的手段、孤独地去工作,有书没书都一样。既然我既没有那份天分、也没有那套装备,以这种方式学物理;我便把这事推到更合适的时日。于是我着手做数学,同时“远远地”听着几门课——没有一门能让我满足,或能给我超出通行教科书的东西。但考试反正总得考……

\subsubsection*{(b) 陌生的美人}
\tocsubsub{(b) 陌生的美人}
\label{ch18:local:h28}

\notehead{120}{}
\phantomsection\label{ch18:local:n24}

(11月6日)刚才重新翻阅昨天的注,我得以确认:我留了神,不曾重新落入某种混淆——把数学\textbf{工作}(一种阳占绝对优势的活动)同“\textbf{那}数学”混为一谈。法语如同德语,指称它的那个词是阴性的——这肯定不是偶然:一如涵盖着它的“\textbf{那}科学\dmark{}”,或更宽的词“\textbf{认识}”\origpage{544}\fnote{96}{(*)}{相反,“\textbf{知}”是阳性的,而在“认识——知”这对阴阳偶子里,“知”恰是“夫”。德语在此不那么分明,因为~Kennen~与~Wissen~两词(作为名词化的动词)都是\textbf{中性}的。},或“\textbf{实体}”。对本来意义上的数学家——我是说那个“做数学”的人(如同他“做爱”)——他与数学的关系之中角色的分配,毫无歧义可言:数学,即那未知的实体,他与之结交的实体,他穿透它而认识它。那时的数学,同他认识过或仅仅渴望过的任何女人一样是“女人”——他曾感到过那神秘的权能,那把他拉向她的吸引,带着那既极温柔、又无可回驳的力量。

我头一次察觉:把我拉向“女人”的冲动与把我拉向“数学”的冲动之间,有着深刻的同一——那是在《道德经》的几节诗章触发我写《乱伦颂》的几个月之前(顺带,也触发了我对“阴柔”与“阳刚”的第一场系统反思——它们的中国名字“阴”与“阳”,我当时还不知道)。那是六年之前,我在写一篇两页的短文,题目叫“以代纲领”,言外之意:为(C4的)“研究入门”课程而作——这篇短文构成那门“课程”的引言,或更确切地说,一场关于这门“课程”之精神的宣言。写完这篇笔底最自然不过流出的文字之后,我惊于其中意象之繁富——它们彼此相生,满载着爱欲的意味。我分明意识到:这既非偶然,也非单纯的文学故意的产物——那是支配了我成年生活的两种激情之间深刻亲缘的一个不含糊的信号。当时,我虽未想着以一场系统反思去深究此事(它只在几个月后,借着写《乱伦颂》才出现),我相信甚至也不曾把那骤然知觉的东西对自己清晰地表述;但我想我可以说:在那一刻,我悄无声息地学到了某样要紧的东西——我“发现”了某样东西\fnote{97}{(**)}{那是一场“阴”式的、“女性”式的“发现”——它经由我们对一份新知的迎纳而完成,怀着对来到我们心中之物默默敞开的姿态。这样的时刻,我相信,在我一生中是稀少的。无论如何,凡我留有记忆的发现时刻,几乎全带阳的、“男性”的调门。},某样此前整个逃过我知觉的东西。

当然,同大家一样,我听说过弗洛伊德、力比多的升华等等,但这毫不相干。纵然成吨的精神分析书籍、随你什么书,也替代不了这样的时刻:那时,一切理论、一切“行囊”都被遗忘,而忽然有什么\dmark{}“咔哒!”一声对上了\origpage{545}。正是在这样的时刻,我们对事物的认识得以更新。这与读书、听报告毫不相干——也就是说,与增益一份\textbf{知}\fnote{98}{(*)}{这份确认,并不被下面这件事驳倒:很可能,甚至大概,这次“意识到”(即一件在无意识中被知觉的东西上升到意识层面)曾得益于弗洛伊德共识的存在——我听说过它,但它不曾真的令我或热或冷。一份知可以助一份认识之绽放;但在我看来,更常见的情形,倒是它把一切绽放的企图扼杀在萌芽里——一如那些现成的“答案”,把一个(好)问题的绽放扼杀在萌芽里……
令人惊讶的是:尽管“人人都听说过”性爱冲动在(且说是艺术的或科学的)创造性中所起的作用,但在那些我一度置身其中的圈子的共识里,却看不到一丝痕迹。然而,惹眼的事实并不缺,它们本可早就让我有所警觉。比如,直到三年前,我一生中创造强烈的时期,尤其内在更新的时期,也都以一股强大的性能量之涌入为标志。不过,我的数学活动从来未曾伴随有意识的、性的意象或联想。但我记得,五十年代,在一次Bourbaki工作会议上,我曾有点错愕:一位同事兼朋友当着我的面,像谈世上最平常的事一样,谈起他数学工作中的一个特点:当他抵达一件艰难工作的尽头,他便感到一种强迫性的做爱的欲望(有无伴侣均可)——而且他对自己刚做的东西越满意,这欲望便越强。}。

当我想到“那数学”,我想到的肯定不是那可以被称作“数学”的全部\textbf{知}——从上古到今日,写进种种出版物、预印本、手稿与通信里的那份知。即便剔除重复,它想必也有几百万页排得密密麻麻的文字;也许十来吨的书,或几千册厚厚的卷帙,够装满一座宽敞的图书馆:这肯定不会让人血脉贲张——恰恰相反!说“那数学”,唯有在一幅\textbf{视野}、一份\textbf{领会}的语境里才多少有意思——而这些,在本质上是一些属于个人的事情,绝非集体的。有多少数学家,就有多少“数学”:其中每一位都对他那份数学有某种个人的经验,或广或狭;那经验的一个果实,就是他自己对“那数学”(即他所认识的那份)的领会、他自己的视野——它总归或多或少是零散的。这有点像“\textbf{那}女人”:在某些人眼里,她也许显得只是一个抽象、一个空洞的公式;然而她却有一种深刻的、强大的、无可辩驳的“实在”(至少对我而言),而所遇见或所认识的每一个女人\dmark{},都是她的一个化身、代表她的一个侧面\origpage{546};而同\textbf{一个}女人在另一个人的经验里,无疑又代表另一个化身、另一个侧面。

我在此的用意,绝不是要直面那份困难:把这一大堆互异的关于“那数学”的经验、领会、景“整合”为一个总体、一个统一体——何况,还正是在这样一个时代:依我看,数学成果的产出正经历着一场猛烈的“发散”,而恐怕没有一个数学家敢自夸,哪怕只在大轮廓上,认识我们这门科学里迄今完成的那份实质工作的全部或要害。我的用意倒是稍许检视阴阳在数学\textbf{工作}中的戏——也就是说,也在数学家(或某位数学家,从我自己开始)与“那数学”的关系中的戏。所以,被检视的东西是“数学家”或“某位数学家”(在其与数学的关系之中),而不是“那数学”本身。

\subsubsection*{(c) 欲与严格}
\tocsubsub{(c) 欲与严格}
\label{ch18:local:h29}

\notehead{121}{}
\phantomsection\label{ch18:local:n25}

(11月7日)在我们的智性官能、理性的层面上,“认识”一物,在一切之先,便是“\textbf{理解}”它。而在一场落在我们官能这个层面上的发现工作里,鼓舞着我们身上那个孩子的认识冲劲(独立于“我”、“东家”所特有的那些动机),便是\textbf{理解的欲望}。这或许正是智性的认识冲动区别于它姐姐——爱的冲动——的主要差别。这份理解的欲望,先于一切“方法”而在,科学的或其他。方法是一件工具,由欲望所塑,用来达成欲望的目的:穿透理性可及的未知,以达理解。认识诞生于认识的欲望——当想要认识的是理性时,便诞生于理解的欲望。\textbf{方法},欲望的器械,本身无力分娩任何认识——正如医生的产钳、甚至助产士那双娴熟的手,都不会分娩。但有时,当时机成熟、而它们懂得及时赶到,它们便能有益地助产那个新生儿……

中学生与大学生——即便不是全体,也必定是大半——应当都感受过数学里的那份\textbf{严格}:它由一些阴沉着脸的老师反复灌输给他们,显得像一种全然外在于他们谦卑个人的先天之物,不可理解、专断横行,由一位专断而\dmark{}不容情的上帝,颁给一位欧几里得——他被擢升为总大检查官\origpage{547},领的使命是叫无数世代的学生在课业面前面无人色,好歹把那大写C的文化整个吞下。我想必是少数几个例外之一:在与学校数学的关系中,我从未经历那一阶段——从第一次相遇起,便在一本六年级数学教科书那逼仄的框架里,本能地感到了严格的原始功能与意义:它是一件柔韧而效力惊人的工具,服务于对这些所谓“数学”之物的领会——那些理性凭一己便能全然认识的东西。这份“严格”,也像我前天反思中所谓的“数学游戏规则”、方才我所谓的“方法”的灵魂与神经。仅仅对它们惊鸿一瞥,便仿佛我从来就认识它们——仿佛是我\textbf{自己的}欲望,温柔地、爱恋地把它们塑成这样一把钥匙:它有力能使一个未知而神秘的世界向我开启;那世界被预感到的富饶,将显得永不枯竭……而恰恰是我自己的欲望,在我的中学与大学岁月里继续打磨这件工具——在遇见任何人能使我起疑:世上某处竟存在我的\textbf{同类}——那些如我一样,以探测未知为乐的人——之前;这把显然无人(包括我的老师们)知晓的钥匙,唯有它有力能把那扇门推开一道缝\fnote{99}{(*)}{然而,我在中学与大学学到的那么一点数学,本该也够让我明白:至少在过去,必定存在过像我一样的人——实际上,就是人们称作“数学家”的那些。而且,苏拉先生(我在大学的一位老师)还同我谈起勒贝格,说他解决了数学里最后一批未解决的问题,包括测度论里的(自1945年离开中学起,我就在做测度论)。但在那些年头(1945--48),我凭自己的手段澄清\textbf{我自己}向自己提出的那些问题的欲望,是那样的排他,以致排除了一切对过去或当下的数学家之存在、之著作、之其人的好奇。}。

\subsubsection*{(d) 上涨的海……}
\tocsubsub{(d) 上涨的海……}
\label{ch18:local:h30}

\notehead{122}{}
\phantomsection\label{ch18:local:n26}

(11月8日)三天以来,我的反思原则上一直“论数学中的阴与阳”,而我有种印象:它老也起不了步,因为我部分地被别的营生与事务吸住了。铺垫做了这么多,我一直没来到那个一开始就想来到的点上:在我自己的数学工作\dmark{}里,\origpage{548}占主导的,是\textbf{阴}的音符、“\textbf{女性}”的音符!

我是几周前察觉的,在这场阴阳反思的边边上,并联系着那场“由三部分葬礼悼词引发的联想”——它正是这条长枝蔓的出发点。(见“阳葬阴(1)——或肌肉与肚肠”一注的开头。)老实说,这场联想(我日后还有机会回到它)多少基于这样一个直觉:我的数学进路阳占强势。这个直觉相当自然,因为正是我的超级阳抉择,促成了我在数学中那份旷日持久的投入。但这并不妨碍这个直觉——或更确切地说,这个观念——是错的;只须容我自己稍事检视它,便会发觉真实情形恰恰相反。要说出乎意料,这就是:一场意外!我当时没有“趁热”在注里谈到它,为的是不打断当时反思的线索——那时我正试着圈定我感知阴阳的方式、以及由之透出的那份哲学。可如今,我们总算来到这里了!

这个关于我的数学进路之品格的错误观念,想必是在五六年前——我开始留意事物阴阳两面的那时——未经检视、仿佛天经地义地溜进我身上的。它想必是我那份阳性的、男性的形象招牌的一个残迹——残迹因纯粹的惰性而一直滞留在那儿,怪我只懒得起身,往那个角落里扫上一帚……

读者或许会觉得我在哄他:就在三天前,我还长篇大论地说,数学工作是一切超级阳活动中最超级阳的——在与数学的关系中,数学扮“女人”,数学家扮殷勤进取的情郎——可转眼之间,我竟提出这样的问题:就我这微末个人而言,我的工作或我的“进路”究竟是阴还是阳,并且(像世上最自然不过的事一样)断定:是阴!谁想得到呢!倘若这里有一副表面上的混淆,那缘于对下面这条普遍事实的误解:在一切事物之中——哪怕它是世上至阴或至阳的——都上演着阴与阳的动力学,借着两股原初力量的联姻。比如火,万物中最阳者、阳本身的象征,在其某些侧面却是阴的(此即“阳中之阴”)\dmark{};反过来,水,阴本身的象征,在其某些侧面与功能中却是阳的(此即“阴中之阳”)\origpage{549}。无须在此展开这两个格外有教益的例子——凡被这些论断(它们或许显得武断,或许像谜语)勾起好奇的读者,只须亲自追随系于火与水的那些联想,便能在这两个个案里,自己发现阳中之阴、阴中之阳的实在。而如果他是个数学家,或哪怕只是个对智性工作不陌生的人(纵使他既非数学家、也非科学家),他也将毫不费力地辨认出:面对任何一种智性工作——无论与其他较零碎的活动类型相比,它显得多么“阳”——都存在着互补的阴与阳的进路模式。

一个可能的出发点,是重拾三天前那场反思开头标出的那十五来个阴阳对子\fnote{100}{(*)}{见“最‘大男子气’的艺术”一注(n\textsuperscript{o}~119)。},当时我曾确认:拿这样的智性工作与其他活动类型相比——如做爱、唱歌、画画(画一幅画,或者,画一面墙也行)、侍弄园子等等——对其中每一个对子,都是阳项占优势(而这在数学工作那里尤其如此)。但这并不妨碍:倘若我们留在一种确定的活动之内——比如说做数学(其全然属阳,自不待言)——我们仍能分辨出一个(或有时是一种失)平衡:或阴或阳的特征,因数学家而异;有时在同一位数学家那里,也因工作而异。比如,在某些工作里,被摆到前头的是所展开理论的\textbf{逻辑}结构;在另一些里,则是\textbf{直觉}的侧面。有一种失衡,它显现在读者或听者身上,是那份大家熟得不能再熟的\textbf{不适}之感(有时作者身上也有):当两个必不可少的侧面之一被粗暴地忽视,以便让另一个独占“好处”。(而当两者都被粗暴忽视时,人们就把书扔进垃圾堆,或摔门离开教室!)当两个侧面都强有力地在场——或明言,或藏在字里行间——那便显现为另一种同样熟得不能再熟之感:和谐、美、平衡、满足之感。事情便是这样,与那支配着所循进路的“基调\dmark{}”无关——无论那基调朝“逻辑”的方向,还是朝“直觉”(或同样朝“结构”)或“实质”的方向\origpage{550}。无疑无须展开这个有教益的例子,去描述——比如说——当两个侧面之一被忽视时,症结何在(即圈定方才唤起的那种“不适”);读者凭自己的经验,对此早已了然!同样方向的论断,必将从三天前考虑过的大多数阴阳对子里透出。甚至或许对所有的对子都如此——尽管其中某些更为微妙,无疑需要比直觉——逻辑那一对更深入的检视,方能被充分把握。

现在,我该试着稍稍显豁这个事实,或者毋宁说“把它传递过去”:在我做数学的方式里,领舞的,是我那些阴的、“女性”的特征,而不是我那些“男性”的特征。倘若要做的是把这份印象追到底,在尽可能多的侧面下检验它,那个自然的想法(昨天它确实在我心头掠过)便是:在我所知的阴阳对子之中,检阅那些能代表一种智性工作的某一侧面或把握模式的对子(我猜想总有五十来个),并逐一查看,这两个“配偶”之中哪一个在我身上占优势。我预计:在所有情形里,经检视,两者之中确有一个将显出优势。

这样,在直觉——逻辑那一对里,我一眼便看出:两个侧面在我的数学工作里都强有力地在场。这便是平衡、和谐的标志之一,与其他同方向的标志并列。一如阴阳对子所应有的:对我而言(我是说在我的工作里),两个配偶真正不可分——一个理论的逻辑结构,与对它所处理之事的一份\textbf{领会}之深化,并步共进;也就是说,与对它一份越来越精细、越来越完整的\textbf{直觉}之展开,并步共进。或许在我已刊的著作里,按数学家行当的规矩,最显眼、对读者最醒目的,是阳的侧面——“结构”或“逻辑”或“方法”的侧面。然而,我分明知道:在我的工作里引领并支配着的、它的灵魂与存在的理由,是工作进程中形成的、用以把握数学事物之实在的\textbf{心灵意象}。

当然,为了借助数学语言,尽可能一丝不苟地圈定这些意象、以及它们所给予的把握,我从不吝惜气力\dmark{}\origpage{551}。或许正是在这持续不断的努力之中——把尚未言说的表述出来,把尚且模糊的弄精确——有着数学工作(或许也有一切创造性智性工作)特有的那种动力学:在多少还无形的\textbf{意象},与那个赋予它形式、并一路催生新的多少朦胧的意象的\textbf{语言}之间,一场持续的辩证——新的意象深化着先前那些,而它们同样呼唤着表述,好在轮到它们时得到形式……况且,正是这份永不停歇的工作——用语言去圈定那最初显现为不可名状而无形的“预感”、显现为未成形的“感觉”、显为一幅没入雾霭的意象的东西,尽可能精确、尽可能完美——正是这份工作,自我的童年起、直到今天,是数学发现工作里最让我着迷的东西。不过,倘若这里的“努力”似乎总是朝“语言”一侧——即朝表述、结构、逻辑那一侧,它们构成数学\textbf{方法}的关键成分;倘若(情势使然)一篇数学\textbf{文本}——它理应复现一项数学工作(或至少其果实)——的可见方面,主要也落在那里;这一切并不妨碍(至少在我身上):数学事物之领会的灵魂,并不在那个侧面,数学工作里起作用的活力或动机,也不在那里。我相信,在我的著作之中,这样的关系被颠倒了的工作——我曾仅凭其内在逻辑、凭融贯性的要求或形式主义自身的其他侧面来引导,而不是由一个内容、一种实质——一种以“几何”性的意象与直觉来显现的实质——所引导——必定极为罕见。

无论如何,我一生都无力阅读一篇数学文本——无论它多么平淡、多么简单——倘若我不能以我数学事物的经验为它给出一个“意义”;也就是说,倘若这文本不能在我心中唤起心灵意象、唤起将赋予它生命的直觉——一如鲜活的肌肉与器官的肉身赋予一具身体以生命,没有它们,身体便只剩一具骨架。这份无能,也把我和我的大多数数学同行区分开来;并且(如我有机会提过的)正是它,常常使我在Bourbaki小组的集体工作里难以容身——尤其在共同读书的场合:一连几个钟头,所有其他人都从容跟进,我却往往被晾在一边。

\asterisms

我刚刚追随了几场关于我数学工作的联想,系于“直觉——逻辑”这对子,以及紧随着它自行进场的那几个相邻的对子:无形——有形,未定——已定,未言——已言,模糊——精确,灵感——方法,视野——融贯……逐一检阅(如我曾盘算的)一切可以想见的、与智性工作相关的“对子”,并逐一探测两个配偶各自以何种方式、在多大程度上在我的数学工作里在场,以及两者之一是否显得“定着调子”、是哪一个——这无疑会很有教益。\origpage{552}它非但能让我更精微地把握\textbf{我的}数学工作的特质,也肯定能让我深化对一般数学工作之性质的理解,同样也深化对如此检阅过的每一个对子的把握。但这样一项系统的工作,分明会把我带得太远,超出眼下这场反思的合理边界。我觉得更自然的,是试着在这里重拾、并(若办得到)“传递”那些联想与意象——它们使我确信(无须再走更远):在我的数学工作里,确是我存在中那些“女性”的特征,倾向于悄悄定着调子;并由此为它们在我生活其他领域所受的压制,找到一场出乎意料的“报复”(在人们最想不到的地方!)。

拿证明一个尚属假设的定理这项任务来说吧(在某些人看来,数学工作似乎就缩小到这件事)。我看出两种极端的进法。其一是\textbf{锤与凿}的进法:摆在面前的问题被看作一颗大核桃,又硬又滑,要去够着它的内部,那被壳护住的滋养人的果肉。原理很简单:把凿刃抵住壳,用力敲。必要时,换几个不同的地方重来,直到壳裂——于是心满意足。当壳上有棱有凸起、可据以“下手”时,这条进法尤其诱人。有些情形,这种下手的“茬口”一眼便见;另一些情形,则须把核桃在各种方向细细翻转、细细勘探,才能找到攻击点。最难的个案,是那壳浑圆而均匀地坚硬:任你用力敲,凿刃打滑,只在表面上留下浅浅的划痕——人们终于厌倦而罢手。但有时,凭肌肉与耐力,终究也能成。

我仍可拿开核桃作比,来说明第二种进法。方才浮上心头的第一个寓言是:把核桃浸在一种软化液里——就用清水,何尝不可;不时搓一搓,好让它渗透得更透;其余的,交给时间。\origpage{553}几个星期、几个月过去,壳渐渐变软——时机成熟时,手上一压就够:壳便开了,像一只熟透的鳄梨!或者,把核桃留在太阳底下、雨水里,或许还有冬天的霜冻底下,让它熟。时机成熟时,一根纤嫩的幼芽,从那丰腴饱满的果肉里钻出,像玩儿似的穿透了壳——或者更确切地说,壳已自行开启,给它让路。

而几周前浮上我心头的意象,又自不同:那有待认识的未知之物,显现为一片泥土或紧实泥灰岩的地界,不肯让人穿透。你可以抡镐、使钎,甚或动用风镐:这是第一条进路,“凿”的进法(带锤或不带锤)。另一条,是\textbf{海}的进法。海悄然无声地向前推进,仿佛什么也不会破裂,什么也不移动;水还离得远,几乎听不见……然而它终归围住了那块倔强的实体;后者渐渐变成一座半岛,继而一座岛,继而一块礁石,终于也被淹没——仿佛它最终消融在那一片一望无际的海洋之中……

对我的某些工作多少不陌生的读者,认出这两种进法中哪一种是“我的”,将毫无困难——在《收获与播种》第一部分里,我已有机会就此表白,语境稍有不同\fnote{101}{(*)}{见“梦与证明”一节(n\textsuperscript{o}~8)。}。那是“海的进法”:凭淹没、吸收、消解——那种进法里,倘若你不是十二分留心,任何时刻都仿佛什么也没发生:每一事物在每一时刻都如此显然,尤其,如此自然,以致人们常常几乎不敢把它白纸黑字记下来,生怕显得自己在搞名堂,而不像大家那样抡着凿子……然而,这恰恰是我自幼凭本能实行的进法,从来不曾真正学过。

说到底,这也曾是Bourbaki的进法;而我与Bourbaki小组的相遇,在这一点上简直是天赐:它印证了我、鼓励了我\dmark{}\origpage{554}在这个自发属于我的“风格”里走下去——不然,我多半会发觉自己这一类的人,多少形单影只\fnote{102}{(*)}{在这份极阴的进路里,我的走法,甚至比Bourbaki的大多数朋友所愿意走的更远。这无疑是我终于在五十年代末离开那个小组的缘由之一。}。诚然,“形单影只”的处境,久已为我所熟悉,也并不怎么叫我难受。至于我的直觉性数学工作进法会不会“有效”——即首先(按当时的通行标准,尤其按评判一个初出茅庐的数学家的标准)我有没有能力解决那些尚无人能答的“未解决问题”——这我无法预知,也未曾过分操心。我自然的取向,更倾向于给自己提出我自己的问题,而不是想去解决别人提出的问题。而我的数学著作之所以证明为多产,确确实实主要凭新\textbf{问题}的发现、同样也凭新\textbf{观念}的发现,或凭新的\textbf{视点}、乃至新的“\textbf{世界}”,胜过凭我对已被提出的问题所能给出的“解决”。这股把我推向发现好问题、而非发现答案,推向发现好观念与好命题、远甚于推向证明的强劲冲动,也是我的数学进路里若干打着强烈“阴”印记的特征\fnote{103}{(**)}{况且我有这样一种印象:在我这里,别的任何研究工作,情形并无不同——尤其我所谓的“禅观”。}。这也是为什么,无疑地,当我看到我献给数学的最好的东西,被我从前的某些学生——即恰恰是它们的第一批受益者——漫不经心地或轻蔑地打发了的时候,我格外敏感。

无论如何,我只是事后才得以确认:我的自然数学进法,当我感到被别人提出的一个问题吸引、被它启发时——总之,当它“咔哒”一声对上了、而那问题旋即变成“我的”之时——也同样“行得通”。倘若我试着就这类情形开一份多少穷尽的清单,我猜想那清单会相当长。粗粗一望,便有四桩这样的情形,以其分量“出类拔萃”\fnote{104}{(***)}{我此时想到的问题,按其解决的时间先后,是下列这些:\par\noindent
1.~Riemann--Roch--Hirzebruch~公式在任意特征下的有效性。\par\noindent
2.~任意特征的代数闭域上一条代数曲线的“与特征互素”基本群的结构。\par\noindent
3.~有限域上有限型概形的~L~函数的有理性(它构成“Weil猜想”的一部分,也是通往这些猜想之证明——由Deligne完成——的重要一步)。\par\noindent
4.~定义在离散赋值环的分式域上的阿贝尔簇的半稳定约化。}。在这\dmark{}四桩个案里,那个假设性的定理终于大体得到了证明,凭的是“上涨之海”的进法——被某一套或大或小的理论淹没并消解,那理论远远超出当初要建立的那些结果\origpage{555}。况且我得以确认:我在这些情形(或同样在其他情形)里发展出的观念、想法、公式、方法,早已进入数学“众所周知”的领地——“人人”都知道、随意取用,而不问它们的来历\fnote{105}{(*)}{我自己也常常实行这份对所用的“众所周知”不问来历的轻心——例外是:那些我第一手知道来历的情形,或因为我多少亲历了它们的诞生,或因为我自己就是父亲。一如我在逝去的这些年里、尤其在我对葬礼的反思中所屡次确认的:这份起码的体贴,在我从前的某些学生、或数学世界里的近友们身上,却常常付之阙如——即便事关他们除我之外无从学得的东西、即便他们对那来历不容有任何怀疑。关于这一点,参见“掘墓人——或全体教众”一注(n\textsuperscript{o}~97)中的反思。}。

\subsubsection*{(e) 九个月与上涨的海}
\tocsubsub{(e) 九个月与上涨的海}
\label{ch18:local:h31}

\notehead{123}{}
\phantomsection\label{ch18:local:n27}

(11月9日)昨天唤起的那四桩个案——被“上涨之海的进法”所解决(或毋宁说“消解”)的未解决问题——还有一个共同点。那便是\textbf{J.-P.~Serre}在这四桩个案的每一桩里所扮演的角色。那首先是一个“雷管”的角色:促我“启动”这些问题——借用《导言》里一条提到这角色的脚注的说法(见“一个秘密的终结”,《导言》第8节)。实际上(如我当时所确认),情形显出:在我1955至1970年间——即从我离开泛函分析转向几何,直到我离开数学世界——所展开的那些主要的力量观念与大工程的发生学里,Serre都扮演过这样的角色。

我可以说,几乎没有夸张:从五十年代初直到1966年前后——即十五年间——我在“几何”(取极宽的意义,涵盖代数或解析几何、拓扑与算术)里学到的一切,要么是从Serre那里学的\dmark{},\origpage{556}要么是我在数学工作里自己学的。我想是1952年,Serre来南锡的时候(我在南锡待到1953年),他开始成为我一位得天独厚的对话者——而且在好些年里,对泛函分析之外的种种主题,他甚至是我\textbf{唯一}的对话者。我想,他跟我谈的头一件事,是Tor与Ext——我原把它们想得何等了不得,而他说:瞧,简单得跟请安问好一样……还有内射分解与投射分解、导函子与卫星函子的魔法——那时Cartan--Eilenberg那头“diplodocus”尚未出版。从那时起把我引向上同调的,是他刚同Cartan一道展开的、关于Stein解析空间的“定理A与定理B”。我想此前也听说过它们,但正是经过一两次与Serre的面谈,我才感到这些上同调式的命题——它们如此简单——所蕴藏的全部力量、全部几何的富饶。在他同我谈起它们之前,它们起初整个儿从我头顶掠过——那时,我还没有在一个空间的层上同调里“感到”几何的实质。我欣喜若狂,以致好些年里我一心打算:一旦把手头泛函分析里的工作(反正我分明不打算在那里久留!)收尾,便去研究解析空间。我之所以没有真正兑现这些打算,是因为Serre其间转向了代数几何,写下了他那篇著名的奠基之作“FAC”——它使先前在我看来乏味到极点的东西,变得可懂而极其迷人:它甚至迷人到这种地步,我没能抵住这份魅力,于是转向了代数几何,而不是解析空间。

倘若不加克制,我便会顺藤摸瓜,从这里讲起我同Serre关系的来历——那也几乎就是1952至1970年间我数学志趣的来历。此处不是地方。我只想补一句:理所当然地,把我“领进”上文唤起的那四个问题的“水”里的,正是Serre。这当然不是指把问题的确切表述给我指出来,如此而已。要紧的是:每一次,Serre都强烈地感到一个表述背后的丰富实质——冷不丁看到那表述,它多半不会叫我或热或冷——而且他能把这份对一份丰富、可触、神秘的实质的知觉“传递过去”;这份知觉,同时又是对这实质的\textbf{认识的欲望}、穿透它的欲望。这或许正是一切发现工作里最关键的时刻:那“咔哒”一声对上的时刻\dmark{}——尽管人们此刻对从哪儿下手接近未知、从哪儿进去,还没有任何想法,再模糊也没有。\origpage{557}这真正是“受孕”的时刻——从这一刻起,一场孕育的工作能够开始;倘若境况相宜,它便开始……

倘若说Serre在我的工作与我的数学著作里扮演了一个重要角色,那么依我看,他扮演得更重要的,是在那些关键的时刻的出现里——火花闪过、黑暗而不可见的劳作开始启动的时刻——而不是凭他那为我所不识的技术手段(他不时在恰当的时刻提供给我),或凭我在工作后继阶段向他借来的那些观念。

无疑,Serre之所以扮演这个特殊的角色,缘由之一是:我生性不爱靠读书去打听数学的时新,甚至也不爱靠读书——读那些讨论某套“众所周知”的理论的书或论文——去学那种理论的ABC。但凡可能,我喜欢经由“局内人”活生生的言语来打听。我有这份运气:从最初接触数学界(1948年)直到1970年我离开,从不缺一位内行而好意的对话者,给我讲那些可能使我感兴趣的事情。这也许造成了对这些对话者的依赖,但我从不曾这样感觉\fnote{106}{(*)}{头一个也是唯一的例外,出现在1981年——那已是我离开数学世界很久之后。那是在“穿越Galois理论的长征”之后,我求助于Deligne,把他当作我的阿那贝尔式(anabéliennes)反思的当然对话者。那时我分明感到,他有意利用这唯一对话者的处境,把我“折腾得团团转”——于是我便断绝了一切数学层面的关系,直到今天。关于这段情节,参见“两个转折”一注(n\textsuperscript{o}~66)。}。老实说,“依赖”的问题几乎无从谈起,只要我的对话者与我,就他所教给我的东西,怀着同一音高的兴趣。教一个渴望认识的人,对双方都有益,而且是“教者”自己学习的机会,与被他教的人一样。

方才给出的那个“缘由”,很好地解释了对话者在我数学家生涯中的重要性,却解释不了Serre所扮演的异常角色——依我看,它远远超过我所有其他“对话者”加在一起!确定无疑的是:Serre与我,彼此互补得妙极了。我们有强而多的共同志趣;我在他身上感到那同样的要求、那同样的严格——我在自己的工作里也是这样要求自己的。除此之外,我们的工作循着极不相同的“风格”进行。我有这样一种印象:我们对数学的进路\dmark{}与我们的工作相互\origpage{558}\textbf{补足},而几乎从不相互侵越。我所做的那类工作(以及我做它的方式),与Serre的那类工作大不相同。他会在一篇五十来页的文本里,为一门理论打下最初的地基,甚至肯花一年工夫,写一部中等篇幅的书,优雅而简洁地阐述某个令他心动的题目——但他肯定不会把五年生涯的大好时光、甚至十年或更多,用来经年累月、卷帙浩繁地展开一整套新的语言(此前人们没有它也过得极好),以便为比如说代数几何的某种新而多产的进法奠定基础。他引入过许许多多新的、多产的想法与观念,却不让自己被拖去把它们“怀”到足月、怀到底。而恰恰不止一次,这些想法与观念成了我的出发点,引向一项规模宏大的工作——它对我再合适不过,而要让Serre本人投身其中,则无从谈起。

在这里,一场联想不可抗拒地涌来。借着近日反思的光照,我看待我与数学工作、与我的“著作”的关系,更多是“\textbf{母性}的”,而非“父性的”。受孕的时刻,无论多么关键,在我这里只代表那场“工作”里微末的一小部分:在那场工作里,孕育中的东西、那个将要来的“孩子”,生长并发育。这份工作,正像一位孕妇的妊娠:它自受孕起启动,延续整整漫长的九个月……那把一个胎儿怀到足月、把它\textbf{分娩}出来所需的时间——也就是说,把一个\textbf{孩子}生到世上来,一个活的、\textbf{完整}的孩子,而不是一个婴儿的头、或一具躯干、一具婴儿的骨架、或天知道什么。母亲的角色,分明与父亲的角色(哪怕是世上最好的父亲……)大不相同:后者几乎只消播下一粒种子,便扬长而去,忙别的营生。显然,Serre的数学工作、他对数学的进路,是阳占强势主导的、“男性”的。他对付困难的进法,宁可说是凿与锤的进法,而极少是那上涨而淹没的海的进法、那浸润而消解的水的进法。而且他似乎乐于播下种子,并不过分操心它落到哪里,或它会不会触发受孕与分娩,甚至不操心由此可能生出的孩子会不会像他、会不会承他的名字。

\dmark{}一幅意象能帮我们把握某一现实的重要侧面,却穷尽不了那现实。\origpage{559}现实永远比任何想要表达它的意象更复杂、更丰富。我所得到的这些意象——它们不召自来,为表达对数学的两种不同进法:Serre的进法与我的进法——也是如此。Serre也曾完成过需要长劲的工作;我也曾播撒过观念的种子,其中有些发了芽,并由别人而非我怀到了足月。一如我的数学进法并不缺乏“男性气”(尽管底色是“女性”的),Serre的进法也不缺乏“女性味”,平衡着他那“男性”的底色。

在对一份未知实质的创造性进法里——无论数学的或其他——不可能不如此:没有阴与阳这两股原初的能量与冲动在同一存在者之中的联合而不可分的行动,便没有发现,没有认识,没有更新。恰在这两者的内在融合之中,存在着一个存在者、或一部著作的\textbf{美}——那份精微而难以把捉的品质,它借着那份和谐与满足的特殊感觉,向我们显现自己。这份品质,在我所认识的Serre的全部著作里都在场——无论当面聆听,还是经由他写下的文本。我认识的数学家之中,这份品质如此恒常、带着这份力量而在场者,极少。

\subsubsection*{(f) 阴的葬礼(阳葬阴(4))}
\tocsubsub{(f) 阴的葬礼(阳葬阴(4))}
\label{ch18:local:h32}

\notehead{124}{}
\phantomsection\label{ch18:local:n28}

(11月10日)昨天与前天的反思,远远没有穷尽我数学工作里那些打有强烈印记的、属阴品格的特征。顺着眼下这场关于数学中阴阳的反思,再往深处探测它们,对我也是一次绝好的机会,去深化对一般数学工作之性质的理解。数学中的阴阳这个题目,我原以为一天反思便可巡完,却已经连着花了五天,还觉得刚刚碰了碰边儿。它方才显出真面目:又是一个外表平淡的题目——愈是走近它、进入它,它便愈发广阔、愈发深邃。断断没有道理:我趁乱草草穷尽这个多汁的题目(甚至也不“巡完”它,哪怕是跑步通过),而葬礼仪式正进行到半中间——我可不愿意让它拖沓得漫无节制!

我只想再指出(不加评论,我保证!)我数学工作里这些朝着“阴”的、“女性”方向的“强烈印记”中的两个\dmark{}\origpage{560}。其一是对\textbf{一般}的偏爱,胜过对特殊者的偏爱(后者与前者“配对”或“成偶”)。另一个特征,我觉得还要更强,或者说得更准些,更要害、更要紧,也更宽广(就它\textbf{包含}第一个而言)。倘若有一场“追寻”贯穿了我整个数学家生涯——从十七岁(刚出中学校门)直到今天;一场不停歇的追寻,自我著作的发端起便打下了我全部著作(已刊或未刊)的印记——那便是对\textbf{统一}的追寻:穿过数学事物那无限的多样性、以及通向这些事物的种种可能的进法。在多样性之外辨认出、发现出这份统一——那份常常令人错愕的富饶(而丝毫不削减这份富饶);在差异与不似之外,认出共同的品质;一路走到类比与相似的根子跟前,去发现那份深刻的亲缘——这便是我的激情,我的一生。连那些差异本身——那无限而难以把捉的多样性的表达——也终于显为同一棵枝繁叶茂的大树的枝与杈:枝杈无限分岔,而每一根、每一枝、每一杈,都给我指明通往它们共同的树干的路。凭本能、凭天性,我的行程一直是\textbf{水}的行程——水总倾向于\textbf{往下}走:走向那树干、那些根须的行程。倘若我也爱在途中流连,那很少是在树梢,去探查叶片与纤弱的细枝;多半是在粗枝上、在树干与几条主根上,去认识它们的肌理,透过树皮,感到那滋养人的树液上涨的涌流\fnote{107}{(*)}{我自信能在这场穿过多样性而对统一的追寻里,辨认出打下了我一生印记的那三种激情的共同特征——爱欲的激情与禅观也在其内。甚至或许,在一切激情之外,这在我身上是把握实在的一个\textbf{模式}:我倾向于主要去看那些共同的品质与亲缘,把注意力与分量放在它们上面,而不是放在差异上(却也不曾起过把差异变没了的念头)。我注意到,远为常见的倾向是相反的倾向,即阳的倾向。它常常走到这一步:忽视或否认深刻的亲缘。(超级阳的倾向,我们文化的特征。它常常伴着这样一种反射:想把差异碾平,把一切对齐到同一个号称“完美”或“高等”的模子上,以成全一份虚假的“统一”——那是极度的贫瘠化,同时也是一场暴力。)对话者与我之间这些重音的差异,常常是聋子对话的缘由:两场平行的独白各自展开,却永不相交……}。

\asterisms

\dmark{}老实说,对这个不久前才发现的新事实,我还不知道该拿它怎么办,不知道怎么给它安一个位置:在我的数学进法里、在我“做数学”的方式里,那基本的调门,在我身上是强阴的、“女性”的。这与我已经暗示过的某一直觉同向:我存在深处的那个基本调门——我是说我身上“孩子”或“工匠”的调门,即那创造者的、在制约之外的调门(也就是说,在“我”、“东家”之外的)——它同样也是“女性”的,而非男性的。或许此刻我手中已样样在手,足以查清实情究竟如何:细细检视所有这些或朝这一方向、或朝那一方向的征象\fnote{108}{(*)}{我身上若干打有强烈印记的阳的特征,在我看来是\textbf{习得}的特征,来自制约,更确切地说,来自上溯到我童年的那份超级阳的形象招牌。这些特征之中,有对行动的一场无度的投入;朝向未来的极强的投射——即朝向完成我的任务;对一种首先是智性的发现工作的偏爱,以及思想的侵夺性角色;对一切不显得与我当下任务直接凝聚者关闭心门的姿态,尤其我对风景、季节等等的漫不经心。然而有一个阳的特征,我觉得是天生的而非习得的:那把我与\textbf{火}系在一起的极强的亲缘关系——不同于我与水的关系,水分明不是“我的元素”。况且据说,我的星盘也打有一个极强的阳失衡的印记:进入它的所有星座,都是“火象星座”,而没有任何一个水象星座。},去认出每一个的分量,以及从它们的总体里透出什么。倘若经由这样一份工作,我达不到一个“是”或“否”的可触的结果,它也肯定不是无用的:它将帮我更好地圈定我的无知——这份无知,此刻依然模糊、尚未定位,因为我不曾对它禅观。或许我会做这份工作,在《收获与播种》的工作完成之后,并顺着它的势头。但在这里,这同样不是地方。

不过,我之所以被引向这场关于阴阳的反思,是在一场反思的进程中——在那场反思里,我尤力图理解我自己与其他一些人(尤其那些曾是我学生的人)之间的某些关系。所以,在此最使我感兴趣的,是这个方才显出的“新事实”可能有的反响:对我与他人的关系、以及他人与我的关系。而这也正是我在“安放”、运用这个事实上的为难所在。这为难或许系于:这样一件事,除了我,恐怕无人曾有所察觉——至少不在意识的层面、不在表述出来的层面。无论如何,就记忆所及,我从不曾收到过任何可作此解的回声——除了一例不算——正如我也不记得收到过什么回声,把一幅“阴”的影像从我自身反射给我;尽管自我童年起(婴幼儿期不算)我所扮的那个角色,一直是强阳的\origpage{562};以致时至今日,这“男性”的品格似乎仍像一种第二(?)天性,仍以种种方式支配着我的生活。

诚然,某个特征在某个人身上(此处即我)不被意识层面所知觉这件事,并不必然拦阻它作用于与他人的关系。而这个特征确确实实被数学世界知觉到了——在我那些对我的著作多少熟悉的数学同行之中;并且这份知觉,在一个分明广阔得多的数学公众之中“油渍般扩散”开来——这在我毫无怀疑。当我在“葬礼悼词(1)——或恭维”里写道:“那支在这里细心料理我葬礼悼词的无名之笔,把今天已被交给轻蔑的东西,加倍再加倍地赏给了我”,我当时还无从一言道破:在那些我珍视的东西之中,被数学时尚“今天交给轻蔑”的,究竟是什么。但第二天,借着那场我日后必须回头的“联想”\fnote{109}{(*)}{见“肌肉与肚肠”一注(n\textsuperscript{o}~106)的开头,那场联想在那里头一次被唤起。},我已经感到(或许还不曾对自己表述,而它也还不像如今这般清晰):“这某样东西”不是别的,正是那一切被(往往在未成形的层面)承认为做数学的一种“阴”式的、“女性”方式的东西——那方式被默默等同于“瞎拼凑”(bombinage)、“胡话”(借用我的学生兼朋友Pierre~Deligne的恭维——针对那构成他全部著作之根基的文本)、“花架子”、“省力”等等。

当然,在那篇葬礼悼词(由同一位朋友皮埃尔所致)里——包括把我与他相提并论的那一段\fnote{110}{(**)}{见“葬礼悼词(2)——或力量与光环”一注(n\textsuperscript{o}~105)。}——恭维是题中之义!那里谈的绝不是什么胡话或瞎拼凑,而是“\textbf{泰坦般的}气象”、“二十卷著作”、“剥出了\textbf{本质问题}”、“最大的\textbf{自然普遍性}”(原文如此)、一个被他传达观念的那份\textbf{慷慨}所“\textbf{滋养}”的学派、“\textbf{传说般深邃}的理论”、“更新了\textbf{基础}”、“开辟了\textbf{新的应用}”、如此\textbf{自然}的观念,“以致我们难以想象它们曾花费的气力”(更不说它们“容易”了——不过这一点,我留神亲自说清了\fnote{111}{(***)}{见“陷阱——或省力与衰竭”一注(n\textsuperscript{o}~99)。})\dmark{}、“对术语的\textbf{高度关注}”(更不说“瞎拼凑”了)、代数K理论的“\textbf{鼻祖}”、在\textbf{一般}基域上引入的topos、“由Grothendieck\textbf{提示}的类比”、“那些\textbf{猜想}……依然同样无从下手”、“一如Grothendieck所\textbf{梦}想的那样”……\origpage{563}

在这些引文里,我给关键词划了着重号——它们全是一些指称对事物的阴式进法的词。这场由“配量得当的恭维”所行的葬礼里,那“完美的手法”就在于:对着这些品质——它们一面“被交给轻蔑”,一面却是真实的、在我这里价值连城——系统地使用夸张;而与此同时,对那些互补的侧面——今天独享荣耀的那些“男性”侧面,它们在我的著作里,也如在任何人的著作里一样强地在场,几乎无一例外——则拿橡皮整个儿、彻底地擦了个干净。

况且,在那篇谈Pierre~Deligne的文本里,被摆到聚光灯下的,恰恰是这些“男性”的侧面与价值——带“女性”味的一丝半点也被排除在外:既凭那几顶选定的形容语(“众所周知的\textbf{困难}”、“令人惊讶的\textbf{结果}”、“使$\ell$-adic上同调成为一种\textbf{强有力}的工具”、“\textbf{第一步}”、“\textbf{出奇地有用}”、“\textbf{敏捷}”、“\textbf{透彻}”、“对每个问题都作出\textbf{启人心智而富有建设性的回应}”、“\textbf{耀眼的发现}”),也凭对可触的成果的一份周详的罗列(而在给我画的那幅速写里,我的成果一项也没有被提及,甚至也不曾暗示:我的成果或许曾对Deligne的成果起过作用)。

我不后悔费神做了这份形容语的快速汇编——它的效应真正触目惊心!倘若在一份有结构的知识层面,对阴阳略有所知的人依然寥寥,那么只好相信:在我朋友皮埃尔的无意识里、如同在那位为他执笔者的无意识里,有着一份毫无破绽的知觉。它在这里被用于某一项事业:把该交给轻蔑的交给轻蔑,并给人群指定一位可供仰慕的英雄。

况且我很怀疑:我方才重读的这三篇短文,有过许许多多的读者。但读者或多或少,依我看是个次要问题。对我而言,这些文本所诉诸的,并非什么假想的潜在恩主(毕竟,替他的机构寻找恩主以筹措经费,不是我朋友皮埃尔操心的事),而是“\textbf{大会全体}”——它出现在与那篇同名的注(别名“掘墓人”,n\textsuperscript{o}~97)一同进行的那场反思里\dmark{}\origpage{564}。它们所传递的讯息,如同一份触目惊心而手法老到的速记,浓缩了无数同方向的讯息——它们出自我的朋友皮埃尔、出自那些曾是我朋友或我的学生的人中的另一些、或许还有其他一些人——并被这同一个大会全体接收、认可。倘若存在一个集体无意识(而此刻我相当倾向于相信它存在),那么毫无疑问:在这个大会全体(别名“数学共同体”)的无意识里,一如在那位主持我隆重葬礼的大司仪的无意识里,有着这份同样毫无破绽的知觉:什么是阴(受够了!),什么是阳(脱帽!)。

而这葬礼,忽然以一种新的、意想不到的光照显现在我眼前:在这光照里,我自己的这个人变得无关紧要,变成了那个“应当被交给轻蔑”者的\textbf{象征}。这不再是一个人的葬礼,也不是一部著作的葬礼,甚至不是一场不可容忍的异见的葬礼——而是“数学中的阴柔”的葬礼——而或许更深地,在每一位为葬礼悼词鼓掌的众多与会者身上:是\textbf{那被否认的、活在他自身之中的女人的葬礼}。

\subsubsection*{(g) 超级妈妈,还是超级爸爸}
\tocsubsub{(g) 超级妈妈,还是超级爸爸}
\label{ch18:local:h33}

\notehead{125}{}
\phantomsection\label{ch18:local:n29}

(11月11日)破例地(下不为例……)今天早上我醒得早,睡了不过四五个钟头。昨天反思那出人意料的收束,旋即启动了一场强烈的工作:去“安放”、消化这个方才显现的新事实——那是在我热一份丰盛的汤、睡前吃点东西的当儿,彼时已过凌晨三点。而一大早,这同一份工作又把我从睡梦里、从床上拽了起来……

倘若我说“出乎意料”的收束、“新”的事实,却还须补上:自这条关于阴阳的没完没了的“枝蔓”的最初几天起,我身上就有一种仿佛克制的等待:等一个“了结”,或至少等一场“会合”——它迟早要与某支队伍会合:那支在一场葬礼仪式中聚拢起来的队伍。看上去,我离葬礼的场地越来越远,甚至仿佛葬礼已被彻底遗忘——然而并非如此:它们一直在那儿,像弱音器里、像水印里。我从不曾真正离开它们。它们无声的在场,显于这份含蓄而恒常的等待、这份紧张与悬而未决之感——它把我推向那个尚且云雾迷蒙的点:那“会合”终须在那里完成\origpage{565}。我本能够预感到这个会合点的大致所在——它绕着某一场“联想”(我已不止一次唤起它,却始终未曾表述):它正是这场穿行阴阳、穿行我一生的意外旅程的出发点、最初的动因。这场旅程,总而言之,又将像一个大的循环,回到(或多或少……)它的出发点;或毋宁说,像沿着一条下行螺线转了一圈,把我带到那被探测之物更深一层的地方,带到这些葬礼的“正中心”(倘若我的预感不曾骗我)。

然而,正当我开始准备“降落”,并在一条还十足是“枝蔓”甚或“回锅”的“注”的最后一个段落的转弯处,我忽然发现自己已经登临葬礼仪式的正中,而且实实在在在它的正中心——有点像一个外星人,被弹射到恰好落在披着祭披的神甫面前、落在信众大会跟前;更糟的,像一个被认为已经死去、(几乎)已经下葬的亡人,忽然掀开棺盖(于是花圈与感人的墓志铭纷纷滚落!),以真身现身,裹着白色裹尸布、目光闪亮,活脱脱一只小魔鬼,在人们最不提防的时刻,从盒子里跳出来!

这样,昨天反思的收束,同时便是我方才所说的那份悬疑的了结——一份十分特别的悬疑,在“如铺展的海一般”的工作方式里,它是我熟得不能再熟的,无论数学工作或其他。但就在这场长悬疑之松弛的尾流里,旋即显出一份\textbf{困惑}。我相信,自那时起主要吸住我的,正是它;而且正是在不合时宜的钟点,是它把我从床上拽到打字机前。有困惑,本不足怪——差不多每次,当一个情势忽然显现在一种新的光照之下、而这光照初看似乎与旧的景相矛盾时,事情便是这样。于是头一件非做不可的工作,便是细细探测这些矛盾:检视它们在多大程度上是实在的,或仅仅是表面的——也就是说,是一种心智惰性的表达,那惰性不肯在不同的两道光照之下认出“同一个”东西。这份必不可少的工作便算完成:当一切不谐之音都消解于一份新的和谐(哪怕这和谐本身仍是暂时的),消解于一幅把先前诸幅局部景涵盖并统一起来的景——必要时修正、调整它们,并剔除那些终于显出根本为假的。在这样一幅焕新的景里,那个诞生了它的“旧”\dmark{}——即融入其中的那些更为零碎的景——本身也获得一份新的意义\fnote{112}{(*)}{可与“孩子与老天爷”“错误与发现”两节(n\textsuperscript{o}~1、2)中的反思对读。}\origpage{566}。

回到我的“困惑”。它在此:那个“了结”或“新的日子”,在于一幅忽然显出的意象——“数学中的阴柔”的“象征”的葬礼,大操大办,以我这个人而为化身;同时,它又是每一位葬礼与会者身上“被否认的女人”的投影;换句话说,这是一场象征性葬礼的意象,葬的是一位\textbf{超级母亲}——总而言之,一位赎罪的牺牲,代替那个女人——却极少是母亲——而她,正在每一位前来为葬礼鼓掌的与会者的黑暗地窖里,萎顿着。这幅意象似乎与\textbf{另一幅}相反的、尚且朦胧的意象相矛盾——后者在6月之前的反思进程中逐渐成形(在“掘墓人——或全体教众”一注达到顶点):那是一位\textbf{超级父亲}的意象,又受敬慕又遭惧怕,又有吸引力又遭憎恨,被他的孩子们“屠杀”,他被残害的躯体,就在这“同一场”葬礼中被交给嘲弄。两幅浓墨重彩的意象并排一放(倘若还需要的话),似乎要在滑稽与谵妄的边上打转;我很容易想象,这些按精神分析路数来的幻象,会激起怎样一场剥头皮的舞蹈——倘若竟有读者有那份气力,一路跟到这里!

我乐意把他们留给他们的舞蹈——它将给这场不同寻常的葬礼添一支最有效果的异国情调的插曲;而在此期间,我宁肯追随一场昨夜便已呈现在心中的联想——我相信,它能把这两幅所谓对立、乃至不可调和的意象或侧面调和起来,甚至使它们相爱、成婚。

\subsection{阴阳的颠倒}
\label{ch18:local:h34}

\subsubsection*{(a) 颠倒(1)——或激烈泼辣的妻子}
\tocsubsub{(a) 颠倒(1)——或激烈泼辣的妻子}
\label{ch18:local:h35}

\notehead{126}{}
\phantomsection\label{ch18:local:n30}

(11月12日)我原打算在注里继续这场联想——昨天注的结尾谈到过它,它能把从我葬礼中透出的那两幅表面敌对的意象“调和”起来、让它们“相爱”。正当我准备朝这个方向起笔,我感到一阵迟疑,而我不愿越过它。

这联想事关我母亲与我父亲的关系,事关1933年那个家庭的解体之意涵——那是凭我母亲的意志、带着我父亲的首肯(起初勉强而尴尬,继而急切而彻底)而发生的。这段关键的情节,在我父母结成的那对\dmark{}里打下一个颠倒的印记\origpage{567}:在这对偶子里,我父亲一直扮演男性价值的英雄化身——那份化身被张扬地崇拜着;而我母亲(倘若曾有一个专断而好支配的性格的话,便是她)则把自己打扮成被征服而乐在其中的女人的颜色,底下却是被持续不断的交锋打满了印记的日常生活。对孩子之牺牲的首肯,标志着\textbf{上帝与英雄崩塌}的时刻;继之而来的,是“胜利的轻蔑”的一场真正的放纵——在那位直到前夜还在扮演神魂颠倒的崇拜者的人身上;而她从此取代了那位跌落的、被阉割的而乐在其中的英雄,把他贬进那个受轻蔑的“女人”的角色——恰好就在同一时刻,她自己从这个角色里被解脱出来……

我说的这寥寥数语,如此梗概,我担心如此浓缩,以致它引起的误会恐怕数不胜数,而非帮助人理解某场葬礼的隐秘发条。然而我感到,此处不是把我方才以几句话勾勒的东西稍加展开的地方。要以最低限度的精微,复现一份复杂的实在——那实在被两位当事者着意搅浑——须要一场新的、长长的枝蔓,其规模并非语境所能辩解。我此刻不觉得自己有跳进去的兴致;尤其因为,这涉及一份牵涉他人的情势,而我自己的责任(作为共同行动者)在我看来并不真正牵涉其中。我自己、我的姐姐,在其中不是作为行动者出现,而是作为我母亲手里的\textbf{器械}:用来打倒那位被热烈敬慕与嫉妒的英雄,好取而代之,把他变成嘲弄的对象。

倘若这场五年前被耐心揭开的脚本\fnote{113}{(*)}{关于这一点,参见“表面与深度”“写作之颂”两注(n\textsuperscript{o}~101、102)。},是我所认识的这一类里最极端、最暴烈的,我此后却仍有的是机会,在其他伴侣之中察觉到全然相仿的脚本。对我父母生活所做的这份工作,大大帮助我睁开眼睛,看见那些此前整个逃过我知觉的事情。但在当时,我目瞪口呆——而也有得可呆!今天,我倾向于相信:除了那份色彩之暴烈与众不同,我在我父母结成的那对偶子里揭开的这种对抗关系,或多或少是伴侣关系的典型——或至少极其常见。因此,那位像我一样终于动用官能、去探测伴侣对抗(或男与女对抗)之隐秘发条的读者,对我在这里说的寥寥数语,将不至于特别惊讶(乃至震惊)。

倘若我试着撇开一件件个案各自特殊的东西,而从我有机会近距离看到、并有所领会的那些男与女的对抗里,提炼出共同点\dmark{}\origpage{568},那么得到的是下面这些。

\textbf{1.}~在女人身上,有一种朝向男人的钦慕与嫉妒的姿态,缘于罩在他身上的威望(常常言过其实)——那威望来自他的处境(尤其作为雄性),以及为之张目的种种品质(真实的或假定的)。

\textbf{2.}~常常,其中混有一份怨怼、甚或仇恨的成分,缘于男人(比如情郎或丈夫)与父亲之间的一份(理所当然是无意识的)混同。母亲对父亲的对抗关系,被女儿接了过来据为己有——她(或多或少充分地)与母亲认同。还常常加上一些(对父亲的)更直接的怨怼缘由(他的暴君姿态、缺爱、不关心、不体贴等等)。此后,这些“现成可用”的对抗之情(以及别的),便原封不动地投射到那位伴侣(现实的或潜在的)身上,无论他有没有“对得上号的长相”。

所以,方才(在第1条里)我写道,女人朝向男人的姿态(尤其钦慕与嫉妒)是“缘于威望等等”,这只是部分为真。我觉得,最常见的情形是:这些姿态里的\textbf{活力}\textbf{源自对父亲的关系}(哪怕父亲早已死去、早已下葬),而它进入行动,只在有限的程度上取决于伴侣个人的品性。

\textbf{3.}~作为对她那些劣等之感(不用说,完全出于主观)与蒙着面纱的对抗、乃至敌意或仇恨的补偿,有一种执念缠着她:要对伴侣行使权力(而按那份或多或少默契的普遍共识,被认定握有权威的,恰恰是他)。女人对权力的行使,凭她手头一切手段进行(其中最有力的,是她的身体,而尤其,是孩子\fnote{114}{(*)}{然而,那最主要的共同“手段”在此被略而不谈,因为它天性更为精微,不易用几句话说清。它在于某一套通行无阻的“战术”——那在关于阴阳的反思的后文“天鹅绒中的利爪”(注~n\textsuperscript{o}~137--140)部分里得到检视。}),而且它几乎总是隐秘的。伴随着它的那份犒赏,因此多半是无意识的,但它并不因此少了实在、少了分量。常常,这权力的游戏变得吞噬一切:它变成女人生活的主要内容,吸收她几乎全部的能量;其余的一切(包括爱的冲动与孩子)都从属于它,甚至被毫不犹豫地牺牲给它。

\textbf{4.}~\dmark{}最极端、被撕裂得最厉害的,是这样的个案:对雄性的钦慕与嫉妒——她要支配他,又显得臣服于他\origpage{569}——伴随着对一切属于阴柔性之事的轻蔑、乃至厌恶与仇恨:对她自己身为女人的处境。然而,恰恰只有玩弄她自己的“女性味”,她才能指望降服男人,或至少把他摆布到她的心意!这样,为满足她最强烈的利己冲动——“让男人听使唤”的冲动(甚至制服他、乃至毁掉他……)——她竟不得不全身投进一个她所憎恶的角色,一个她自觉可鄙、配不上她的角色。正是在这拒绝自身处境与天性的极端个案里——一份超级阳的、反阴的抉择——她将动用自己全部这些力量,去求一条摆脱她身上冲突的幻想的出路:一场\textbf{角色的颠倒}:她自己取代男人——那位英雄与主人,昔日的钦慕与嫉妒之对象,而今的跌落者——把他贬进她长久以来像一件卑贱号衣一样穿着的角色、那个她终得解脱的受轻蔑的角色……

我方才勾画的这幅速写,同样只是梗概,至多能在那种人身上“唤起”某一现实——他已在自己的那边、此地彼处知觉过它,却或许还不曾试图像这份粗略的描述那样,多少把它圈定;倘若我想给它一些立体感,我至少该试着把这整套相互敌对的感情与意愿所上演的各个\textbf{层面}(几乎全部无意识)弄清楚。此外,在这场无情利己机制的缠结里——爱的冲动在其中似乎被严格排除在外——也该试着给这冲动安一个位置:看它在多大程度上、以何种方式,助长这永世打转的转圈(也许就像风力,被一座精巧风车的翼板捕获,叫一副沉重的磨盘永远转下去……),以及在多大程度上,齿轮也会时而停下、归于沉寂,好让\textbf{别的东西}畅行无阻。

最后,我完全忘了谈在他身上、“伴侣”或主角身上上演着什么——仿佛他只是相对于她而存在,仿佛他只是迎面那位女人的吸引与厌拒、钦慕与嫉妒的\textbf{对象}。这份遗漏的缘由之一无疑是:在这对偶子的旋转木马里,扮演主动角色、全身投入、并往往在其中找到(退而求其次的)真正存在理由的,正是\textbf{她};而\textbf{他}却什么也没看见,他忙着别的事,而且\dmark{}天真得没人比得上\fnote{115}{(*)}{(11月23日)当然,倘若旋转木马转得起来,那是因为(他再“天真”也罢)\textbf{他}也在其中捞到他的好处,一如她——而把这些看牢,正是她的本职!依我看,她借此“拿住”他(而她自己也借此被拿住……)的两个主要“挂钩”,是虚荣,以及对一份情感与爱恋之安稳的需要——由一位稳定的伴侣作保。还有孩子……},\origpage{570}一次接一次地起反应,而不试着理解;而且(更要紧的)实在不曾理解——甚至(依我看)在无意识的层面也不曾。至少,这是自我开始留意这对偶子的旋转木马以来,我一贯得到的印象!不过同样真确的是:我对男人的角色知之甚少,因为我只能在我自己这微末个案里真正贴身观察它;而女人的那一侧的角色,我却不止一次有机会从第一排包厢里看到。

反正,即便我肯花上十页、乃至整整一卷的工夫,来充实我这份稍嫌太梗概的描述,对那位在这门事上还不曾“动用他的官能”、从未见过或感到过此类事情的读者,也终究是白费。至于那位多少“在局中”的读者,我说的这寥寥数语,尽管笨拙含混,肯定够把他重新浸回他凭自己早已知觉的那些事情里,并在他心中唤起不亚于我写下这份速描之时背后所存的那些意象与联想。

依我看,无须更多,便可看到那“缺失的环节”显形:它连接着对“超级父亲”的对抗——它在所说的那场象征性葬礼中找到表达——与对“阴柔”的轻蔑、拒斥,而更深地,与对自己身上“女人”的否认(后者或许将在一位“超级母亲”的“葬礼”的象征中找到表达,在一大堆滥情的、一语双关的颂词之下……)\fnote{116}{(**)}{(11月23日)这句“无须更多”,后来显出有点匆忙:一周之后,这条结论与这个“缺失的环节”已被忘得一干二净!至于要到达一个更有说服力的“缺失环节”所缺的那一步,参见昨天那条注“阴阳的颠倒(2)——或反叛”(n\textsuperscript{o}~132)。}。

\subsubsection*{回顾(1)——或一幅画的三联}
\tocsubsub{回顾(1)——或一幅画的三联}
\label{ch18:local:h36}

\notehead{127}{}
\phantomsection\label{ch18:local:n31}

\dmark{}(11月13日)现在时机似乎成熟了:试着以几笔大轮廓\origpage{571},勾出对葬礼的一幅更清晰、也更有层次的景——它(如我前天所写)“把先前诸幅局部景涵盖并统一起来,必要时修正、调整它们……”。粗粗一望,我看见三幅这样的先前的景:该把它们认作同一个\textbf{全}的三个局部侧面。

显出的第一个侧面,最明显也最简单化的,是“\textbf{对一场异见的报复}”的侧面——它是“掘墓人——或全体教众”一注(参见注~97)里主要摆在前头的侧面,即患病情节之前的最后一条注。在I至X各行列(事故之前的那些行列)的注里,依我看,它也是把那些\textbf{集体动机}圈定得最深的一条——“掘墓人”别名“(几乎)大会全体”的集体动机。

我刚刚重读了这条注。第二个侧面——我可以称之为“\textbf{超级父亲}的(远不止象征性的)\textbf{屠杀}与(象征性的)\textbf{葬礼}”——在那里头并不出现。这或许是因为:葬礼动机中的这一成分,并不真正涉及“大会全体”——它当时是我注意的焦点——而主要(若不说排他地)涉及“那些曾是我学生的人”。诚然,这些人——即便把他们那位无可争议的头领、我的朋友皮埃尔放在一边——在葬礼的操办中扮演了头等重要的角色:没有其中一些人的积极贡献、没有所有人的首肯,葬礼便无从举行。(关于这一点,参见“沉默”一注,参见注~n\textsuperscript{o}~84.)所以,主要是经由他们,“超级父亲”的侧面才对我显得对理解葬礼至关重要。

第一个侧面、“报复”的侧面,是自1976年Yves~Ladegaillerie的挫折以来,就强加给我的注意的\fnote{117}{(*)}{见“谁也拦不住进步!”“棺材2:锯成段的切割”两注(n\textsuperscript{o}~50、94)。};此后我倾向于遗忘这个侧面,但在随后几年里,它不时重新唤起我的记忆。它终于越过了那种只被“感到”的无形阶段,变成一份清晰而有层次的理解的实质——在那条谈“掘墓人”的引注里。第二个侧面、\dmark{}或“超级父亲”侧面,只是在《收获与播种》\origpage{572}的反思进程里\fnote{118}{(*)}{(11月29日)老实说,这个侧面早在许多年里,就已以一种皮肉直觉的形式,存在于我同Deligne的关系之中;只是在《收获与播种》的反思之前,我从不曾在此驻足。}才开始显出,而且起初\fnote{119}{(**)}{见“敌对的父亲(1)(2)”两节(n\textsuperscript{o}~29、30)。}与葬礼本身无关——葬礼本身,是我要在随后几个月里才发现的。在整个葬礼反思的进程中,这个侧面渐渐走出迷雾,终于在“屠杀”“躯体……”“……与身体”(87、88、89)诸注里,以惊人的形态现身。这些注写于5月12、16、17日;“掘墓人”那条注写于5月24日;患病情节6月10日登场,把诸注的写作中断了三个多月,9月22日才重新接上。至少很可能:倘若这场(实在不是时候的!)情节不曾在一个我正准备对整体作盘点、并画一条最后界线之时出现,我对葬礼的景,本会止步于5月12日至24日那两周里透出的景——一幅“两联”的景,各自停在各自的角落,而我想不到试着把它们拼起来。

然而,一直有一种弥散的感觉,像一层几乎察觉不到的薄雾,觉得那最终的答案始终没有被真正把握住;那是“在黑暗中摸索”的人的感觉(这个说法,肯定在我的葬礼诸注里出现过一两次)。“掘墓人”那条最后的注,多少起了点薄雾里一阵微风的作用:它能给人雾已散尽的幻觉,而其实它们只是稍稍挪了挪地方。或者说:这条注所重拾的那个侧面,在其中显得如此清晰、如此有立体感,以致对这一侧面的一份可触、透彻的理解之印象(绝非幻觉),连同伴随它的满足之感(在注的结尾,这感觉肯定显而易见)——这份印象与这份感觉,造出了一种亢奋,一种自觉马上就要够着目标的亢奋,使我或多或少忘了另一个同样有分量的联——“超级父亲”的侧面,它还“记在账上”呢!

第三个联,只在三天前才显出(患病那个不合时宜的情节出现后整整五个月)。那便是\dmark{}“\textbf{阴柔}”的(象征性)葬礼与(实实在在的)下葬这一侧面\origpage{573}——这“阴柔”被可视化成一位“\textbf{超级母亲}”,而她本人,正是由我这微末个人所化身的!这个侧面,是在一场全然未曾预料的、关于阴阳的长“枝蔓”的尽头显出的——那场枝蔓,终于落实了一份努力:把从某篇“葬礼悼词”里生出的一场“联想”以可懂的方式表达出来,而那篇悼词,本该为葬礼仪式收尾。这场著名的“联想”或“直觉”(我在“肌肉与肚肠(阳葬阴(1))”一注的开头处首次暗示过它)至今仍未被显豁——但万事俱备,而且我许诺要谈它,已有些时候了!

无论如何,一路走来,显出了大量的事实与直觉——其中一些对我而言是新的、意想不到的——而它们全都让我有效地重新接上了我生活的重要侧面,一如接上一般存在的重要侧面。这些事实之一——我数学工作的“基调”是“女性”的——看来与那场还在等它时辰的联想所依据的直觉之一相矛盾:那个直觉说,作为数学家(在其他方面亦然),我是一个十足\textbf{阳}的人物;一个系于葬礼“超级父亲”侧面的直觉。而这同一个事实——它看来与那场联想(整场阴阳反思都由它而生!)相矛盾——也一转眼便催出了那个此前一直逃过我的第三联:“超级母亲”的侧面。与此同时,(到头来)与那场似乎已被忘了近百页的“葬礼”的“会合”,也完成了。

要说“上涨的海”,这就是上涨的海——只须希望:最终的结果,我是说那幅我正准备从冥府里拽出来的“景”,配得上手段,配得上整整一海关于阴阳的哲学-弗洛伊德式的枝蔓……潮水已经启动(随着发令的那条注“肌肉与肚肠”),是10月2日;关键的“新事实”在随后几天便已显出\fnote{120}{(*)}{我记得,就在两天后,在“天真(阴与阳的联姻)”一注(n\textsuperscript{o}~107)里,所说的事实便已显出,并构成那条注里所谈的“种种迹象”之一(对它们未再多加说明),那些迹象“不止一次使我起疑:……在我身上占支配地位的,是那些‘女性’的品质……”。},\dmark{}正当我一天又一天地准备把那场著名的\origpage{574}“联想”终于写成白纸黑字——它出现于五个月前,5月12或13日,就在“葬礼悼词(1)——或恭维”一注的反思之后,与那条关键的注“屠杀”同日。但这个事实,直到五天前、11月8日,才在注里“被揭开”——在那之前,是连着三天关于数学中阴阳的三条预备的注。那就是“上涨的海……”那条注(122)。两天后,11月10日,随着“阴的葬礼(阳葬阴(4))”那条注(124),“超级母亲”登场了(但这个词,是在次日那条注“超级妈妈,还是超级爸爸?”(125)里才说出口的)。于是,这便是葬礼的“第三联”!

我并不是有意要投入这场回顾——借着当下的势头,从葬礼三个主要侧面相继显现的角度,对它作一番回顾(依我此刻看事情的方式)。在一部长程的禅观之中,这样偶一为之的回顾,每一次都被证明极其有益:它们给出反思进程的一幅总览,同时又给出其中某些主要“成果”的一幅新视景\fnote{121}{(*)}{这种回顾,在我看来,在一份数学工作里十分罕见;我自己也是从去年春天动笔写“Pursuing~Stacks”以来才实行它。倒是有一项惯常的工作实践,对一份进行中的数学工作的观念与成果,能起类似的“新视景”效应:把所展开理论的全部观念与命题,按当下此刻的理解看来最自然的次序,从头重述一遍。这样一份乍看纯乎例行的工作,常常导致理解的实质性深化——比如,借着新的编排对内在融贯的要求,显出一些同样“自然”的、此前未被见到的观念、性质、关系等等。有时,借着显出某些假设的偶然或人为品格,或整个出发点语境的逼仄品格,“重述”的工作会通向初旨的一场意想不到的拓宽,赋予最初展开的理论以新的维度与分量。}。或许最让这场回顾的假想读者吃惊的,是我绕了这么远的一场枝蔓,而不是马上来到那场著名的“联想”(它还欠着呢)跟前,把它谈完,好终于抵达葬礼之下那条著名的“最后界线\dmark{}”;\origpage{575}那条界线,我在9月29日“葬礼悼词(2)”那条注里,是那样急于画出——彼时我刚把6月悬着的反思的挽具重新套上!况且,三天后我开始下一条注“肌肉与肚肠”时,正是抱着这样的心态:那条注以对这场联想的暗示开头,却不给它任何具体说明。

我之所以当时没有给出说明,并且一天天、一星期又一星期地把它推后——已推了一个月零十天——绝非出于什么或早或晚会显形的故意。倘若我试着探测它的缘由,我会说:我当时想必凭本能感到——甚至无须对自己说破——在我当时所处的那个点上,劈头写出那场联想,将毫无意义;那会像一份单纯的“表述”,纯粹形式的、口头的,而被那些词句所覆盖的丰富实质——它们只是凭纯粹记忆的效应来到我笔下——将依旧被忽略、不被知觉。读者倘若是个数学家(或科学家,倘若不是数学家),肯定屡次认识过这样的情势,以及它所激起的不适:人们就这样面对一个表述——可以毫不费力地确认它完全精确,其中每个词的意味我们也多少知道——却分明感到它的“意义”与实质整个儿逃过我们。这样的情势,或许在那些\textbf{非}技术性、却表达着一份被作者强烈知觉的可触实质的文本里,还要常见得多;只是有一点不同:读者相当清楚地觉察到自己所读之物的意义正逃过自己,这种情形反倒少见得多。就眼下这个个案而言,还有更多——因为连\textbf{我自己}也一样:几个月来,我早已不在那篇葬礼悼词、以及系于它的那些联想的“水里”了;几年来,我也不曾真正“潜入”阴阳的实在(虽然一路走过时步步擦过它……)——即便是\textbf{我},我当时所能写下的、用来“说出”这场联想的东西,也会是口头的东西,并不真正被感到、被知觉。迫使自己屈从于此,或者更确切地说,强制自己去做,会是一种纯粹形式的做法,一种凭良心交差的做法:总算清偿了一份义务,总而言之“封了账”,同时留心“给足分量”、不在路上丢失那场“联想”——它(我记得清清楚楚!)曾经多汁而滚烫,却早已在记忆的一角里凉透、发霉了!

\dmark{}倘若我记得的东西,确确实实要用来深化一份依然零碎\origpage{576}的领会,那么对我而言清清楚楚:这上百页的“枝蔓”我当时省不得。它们构成那场贯穿整部《收获与播种》的反思里最深的部分。

我还无从预言:我正准备顺着它们的尾流析出的那幅葬礼之景,会留给我一份圆满满足之感,还是仍会剩下一些晦暗的角落或一些不谐之音——我或许会放弃去照亮、去解决它们,至少眼下如此,或在《收获与播种》之内如此。但无论如何,一如在我的数学著作里,我知道:这上百页中的每一页,一如至此写就的《收获与播种》六百页文本(出入无几)中的每一页,都有它独一无二的位置、它的讯息与它的功能,而我一页也省不得(有没有读者肯一路跟到那里,且由他去!)。纵然所追逐的目标已远(甚至全然被忘掉……),这些页中的每一页,都带给我它自己的收成——唯有它能带给我的收成。

\subsubsection*{回顾(2)——或死结}
\tocsubsub{回顾(2)——或死结}
\label{ch18:local:h37}

\notehead{127$'$}{}
\phantomsection\label{ch18:local:n32}

(11月17日)我刚度过了相当难受的四天,身边纷扰不断。谈不上顺势前进;我在诸注上的工作,仅限于一点后勤事务:重读那部分要送去打字誊清的文本,改正已誊清的部分。每条注的“初稿”——写下一注之前重读一遍——与送去复写的定稿誊清本之间,我至少要读三遍,三遍都用心,至少在前两遍里随时调整措辞。这样下去,《收获与播种》的文本我快要熟透了!但尤其要紧的是,我做了必要的一切,以确保送去复写的文本,确实是我真正拿得出的最好的东西——包括它的形式。除葬礼的一条注之外,《收获与播种》中凡我写就并重读过的各节各注,末次阅读都给我一份圆满满足之感。我感到,每一次我都已抵达:以我力所能及的清晰与有层次,说出了我要说的话——既不隐瞒写时于我已清晰、已懂、已知的任何东西,也不隐瞒那依旧晦暗、模糊、未懂、甚至全然神秘、未知的东西……

\dmark{}唯一的例外,是10月17日那条注“一半与全——或裂缝”——自那条注起,禅观的“线”分成了两股,沿着我所命名的两个主题(在“阴阳之钥”诸注的续篇里作副题):“我们的母亲,死亡”与“拒受与接纳”\fnote{122}{(*)}{把关于阴阳的“枝蔓”诸注用副题归组,这个需要几天前才感到。这也使我重新调整了当初给这些注起的名字;因此,这些注在有的地方被引用时,所用的名字与其定名略有出入(但号数照旧无误)。与此同时,也冒出了这组注的当然之名:“阴阳之钥”。}。所指的是那条注的最后部分,即那两三页:其中我谈个人内部的分裂,说它是伴侣之内、家庭之内、人类社会之内分裂与冲突的终极根源。这份直觉,最先显现在我“出发”离开科学世界之后的最初几年,并在岁月中发展、印证、深化,直到今天。它对我已变得如此“分明”(尽管我从不曾费心把它细细检视、面面检视),以致它有点像天经地义地溜进了反思,毫不费心以某一个“抓手”把它提出,好多少显出这份“分明”。可是,倘若这几页的阅读留给我一份模糊、不满足的印象,那肯定不单是“表述”笨拙的问题。毋宁说,我感到我当时是想一步跨过一个关于这个复杂主题的实质反思——对这场反思,我分明觉得要素样样在手,可以去做,但它终究没有做!在10月25日那条注(“失乐园”(116))里——它直接系于17日那条注(由它出发展开“拒受与接纳”的主题)——我先多少试着“补救”前一条注里我注意到的缺漏;但最终并没有说出多少更多的东西,只说了这么一句:至于一场可能的“探索冲突之旅”,“眼下我毫无意愿朝那个方向走”——罢了,那只好留待来日!

\asterisms

\dmark{}在四天前那条注里,我巡了一圈葬礼图景至此透出的三个侧面、或“三联”\origpage{577}。事后我想起:在葬礼反思的进程中,已有两个时刻,我感到、并写到:我触到了冲突的“死结”。那是在“死结”与“葬礼悼词(2)——或力量与光环”两条注里。这两条注,与《收获与播种》最早几节里(表面“十分一般”的)一些反思接上了头——“他人之无谬与对己之轻蔑”(第~n\textsuperscript{o}~4节)。正是\textbf{对己之轻蔑}、对我们身上那份给我们认识与创造之力的力量的不识,它同时也是\textbf{轻蔑他人}的源头:那永世不改的补偿反射——把自己抬到他人之上,好向自己“证明”自己的价值;比如,动用那贬低或压服的可怜权力,或干脆叫人受苦、叫人受害。

写那条注时,例子当然不缺。当时在我心头最在场的,是Pierre~Deligne:我曾多次见他动用他那使人灰心、甚至羞辱人的权力,其方式常常让我感到无法解释。只是在这条注写成两个月之后,我才开始发现“葬礼的全部辉煌”——4月19日的诸注(“一个梦的回忆——或Motif的诞生”与“下葬——或新父亲”(51)(52))可为佐证。同样逐步地,我发现了我的朋友皮埃尔作为我葬礼与出殡之大司仪的角色。六月之前关于葬礼的诸注(I至X行列),大部分以他这个人为中心。他也是我掌握着无可比拟地更丰富、更个人的材料的那一位——比其他众多与会者中的任何一位都丰富。所以,我感到“触到冲突心脏”的那两个时刻,处于我注意中心的还是他——他也是唯一一个与我保持经常联系、直到今天的人。

\subsubsection*{(d) 父母——或冲突的心脏}
\tocsubsub{(d) 父母——或冲突的心脏}
\label{ch18:local:h38}

\notehead{128}{}
\phantomsection\label{ch18:local:n33}

(11月18日)昨夜睡了十二个钟头——前几夜都睡得太短,正需要!我感到把一份已开始有点散架的能量重新打足了气——比昨天更有劲,可以从头拾起那条著名的“线”了。

\dmark{}在昨天谈到的那两个时刻里,我身上有过一种“闪光”,如此清晰、如此强烈,我压根想不到去怀疑它——我是说,怀疑它向我启示了某种真实的东西、在我这个个案里外在于我个人的东西;怀疑它不是某种纯粹主观之物,(比如说)一份单纯故意的产物:想看到某套我心心念念的“心理理论”得到应用——总之,不是那只被捕蝶人天助地网进网里的“蝴蝶”\fnote{123}{(*)}{这幅意象,见“孩子与海——或信与疑”一注(n\textsuperscript{o}~103)。}!\origpage{578}怀疑这样的征象——无论在禅观中、在数学中还是在别处——便干脆是放弃我的认识与发现之力。我有幸认识这份力量;倘若有一件事是我全然信赖的,那便是它。

我可以考虑在这个“闪光”里、在它所教给我的东西里,看出葬礼图景的第四个“联”——叠加到另外三个之上(它们已在11月13日那条注里巡过)。但我一上来便看到:它与“超级父亲”和“超级母亲”两个侧面亲密地系在一起——而这层显然的联系,远远超出我朋友这个人的范围。昨天我重新唤起的这份对我们身上“认识与创造之力”的不识,不是别的,正是对我们根本统一的不识——那统一,是我们存在之中“阴”与“阳”、“女性”与“男性”的品质、能量与力量的联姻之果。因为我们身上的“男性”,单凭它自己,并不能使我们能够认识、能够创造;我们身上的“女性”,单凭它自己,也同样不能。有认识与创造之力的,不是我们存在中那虚假而可怜的\textbf{一半},而是\textbf{全}、是我们存在的\textbf{总体}。它有这份力,不是作为一场追寻、一段长路、一个生成的终点——仿佛我们在一种暂时的无能状态里跋涉,一路慢慢积攒“权力”;而是:这份力,凭我们的天性便是我们的,自我们出生之日起,便作为无偿的赠礼授予了我们\fnote{124}{(**)}{而且无疑,早在我们出生之前很久……}。

而这份“对己之轻蔑”或“对己之不识”,不是别的,也正是那对着这份赠礼的\textbf{拒绝}:对这份根本统一的拒绝,对与之不可分的力量的拒绝。或者毋宁说,它像这份拒绝那不可分的影子:它是一份\textbf{无能之知}\dmark{}\fnote{125}{(*)}{正如我在一行之后要说明的:这份认识是“搅浑的”,就其本质内容而言,它依旧无意识。然而人们常常看到它露出小小的一角(像冰山的峰顶,其底座被小心翼翼地浸在水下……),借着某种\textbf{无能的信条}——它们不止一次让我目瞪口呆。它们以一份不容回驳、不许还嘴的\textbf{断言}的口气说出;在其背后,人们感到一种激烈的、桀骜的封闭——仿佛这份被当作不可触犯、神圣“事实”来申认的无能,是最贵重的财产,任何代价也不肯让渡……},由这份拒绝所设立的无能之知\origpage{580};一份固然怯懦、搅浑、未被承担的认识——它留神停在已知的东西上(那已知,又知得何其之坏……),生怕潜得更深,生怕见识那份被藏起、被这份蓄意培植的无能所阻塞的未知力量。

在我们这个超级阳的社会里,这份对我们统一的拒绝所取的最常见形式,便是日复一日、时时刻刻地埋葬我们身上的“阴”、“女性”。这恰恰便是那个“超级母亲的联”,别名“‘女性’的葬礼与下葬”——而尤其且\textbf{首要}地,是\textbf{自己身上}的女性的葬礼。

但我分明感到:对己之轻蔑与“超级父亲的联”——别名“父亲的屠杀与下葬”——之间,同样有一层直接而深刻的联系。这份被强烈预感到的联系,正是我现在想试着圈定的。换一种说法来表述这份“预感”、这份直觉:在我们身上的分裂,与对父亲的对抗之间,必定有一层直接而深刻的联系。

当然,这份“对抗”既有机会对着生身的父亲表达,也有机会对着童年里代行其职者表达,或对着任何其他这样的人表达:他在某一时刻、因某种缘由,充当了或多或少象征性的“替补父亲”,而那些原初的对抗冲动,便投射在他身上。所以,我的用意是圈定这些对抗冲动与姿态的深层缘由——它们如此常见,以致人们有时会倾向于把它们看作普遍的;这份缘由,须比一堆具体的怨怼更深——那些怨怼,固然常常全是可触的,是人们可以对生身者怀有的。我不止一次得以确认:这些怨怼,常常更像是一份说得通而又受欢迎的\textbf{合理化}\dmark{}\origpage{581}——为一桩对抗所作的合理化,而那对抗的真正的根、其暴烈与顽韧的缘由,却在别处。

我还可以换一种方式,以它自发呈现给我的形式,表述这份我正试着圈定的直觉:那就是,我怀有一种贴身的信念:在一个是“\textbf{一}”、未曾分裂的人身上——在一个以其存在的总体接纳他自己的人身上——对父亲的冲突、或对母亲的冲突,已经解决。他是\textbf{自主}的,对他的双亲都“\textbf{自由}”。那根继续把我们系在父母身上的脐带——它系得远比童年与青春期长久(而且最常见,贯穿整个成年,直到死亡)——在他身上,已经剪断。缆绳已经解开:不久之前,它们还拽着我们,不许我们真正出发,去\textbf{我们自己的旅程}——去发现我们的母亲,那世界\fnote{126}{(*)}{奇怪的是,在法文里,“世界”“宇宙”与“太空”三个词都是阳性的。德语里相应的词——die~Welt、das~All、der~Kosmos——则分属阴、中(中性在德语里往往是一种“超阴性”)、阳三性。依我看,这更契合这些词所指之物的本性。人们谈“太空”(Kosmos)时,其意味(撇开晚近发明的空间舱与外星人不论)是一份受法则支配的\textbf{秩序}——这观念恰与阳性相应(在这一点上,两种语言一致)。相反,“世界”与“宇宙”暗示的是一个\textbf{全}的意象:我们自己与一切别物,都是它的一个\textbf{部分};而且,这还是某种有待我们去\textbf{发现}、去\textbf{穿透}、去\textbf{认识}的东西。就这些在我看来是要害的方面而言,这两个词所指之物,是“阴”的、“女性”的——尤其相对于我们而言。我实在难以看出,法文为何仍然给它们派一个阳性。
顺带指出另一桩奇怪的“反常”(?)——这一回似乎是在德语里:“太阳”与“月亮”在德语里作~die~Sonne、der~Mond,其性属恰与法文相反,而法文的那一种看上去倒更“自然”。太阳直接让人联想到热与火,那是典型的阳性之物。或许这桩“反常”在北欧诸语里是共通的,因为在那些寒冷的国度,太阳的热从不被感到炙人灼人,而是被期待为一种恩惠、一种生命之源;在那里,太阳(连同大地)被感到像一位滋养的母亲,把热量布施给造物——他们凭这热量“滋养”自己,一如凭来自大地的食物……}。

这份贴身的信念,并非“一厢情愿”之想;它不是一个愿望的投影(为应景改名为“信念”)。它的源头确凿在我的亲身经历之中,而且首先\dmark{},在我与我自己父母的关系里所能确认的东西之中\origpage{582}。我这里指的是:八年前那个转折之后的岁月里,我与我父母的关系中所发生的深刻转变——那转折,以我身上那场“阴的苏醒”为标记,继之以其后几个月里禅观的发现,再继之以两天之后与我童年的“重逢”\fnote{127}{(*)}{我生活中的这些关键情节,我在“重逢(阴的苏醒(1))”与“接纳(阴的苏醒(2))”两注(n\textsuperscript{o}~109、110)里、以及在“欲望与禅观”一节(n\textsuperscript{o}~36)里谈到。}。我意识到,这个转折打上了一份\textbf{自主}的印记——它与先前对种种现成观念、尤其被采纳的观念的依赖,恰成对照。所有这些依赖中最深的,是对我父母的依赖:他们的价值与抉择,塑造了我的价值与抉择、塑造了我自己的世界观;而他们对自身、对彼此结成的伴侣关系、对他们与孩子的关系所抱持的那幅理想化图景,我也整个“成块地”、原封不动地接了过来——几乎分毫未改。自我童年起,我便在这套价值、抉择、意象上“运转”:它们全然不是我自己生活经验的果实、不是我对此经验消化工作的果实,而只是一份“行囊”。这份行囊,好大一部分由陈词滥调与自我陶醉的幻象构成,是我凭“信任”从我父母那里接来的;它们在我一生中,屡屡取代了对周围事物的一份直接的、活的知觉、一份创造性的知觉。

诚然,我所说的这份“自主”,是随着禅观之力的发现立刻显出的。凡我留神检视的东西,它(我相信)都是全面的。但这并不妨碍:许多现成观念——尤其、首要地是那些来自我父母的——起初仍凭纯粹的惰性留在原地,因为还不曾被检视。要看的东西太多,不可能一下子全看!更不必说,经过几个月强烈的工作之后,我让自己被“继续着的生活”分了心——尤其是几段恋情,自不待言\fnote{128}{(**)}{1976年禅观发现之后的那些年头里,我的爱情生活,比我一生其他任何时期都强烈,也更动荡。它肯定代表着一种分散、一场岔开——相对于禅观最初的那股冲劲而言;那冲劲,直到1979年8月,才(以它应有的规模)被重新拾起:那是一场关于我父母生活的长程禅观。(关于它,见“表面与深度”与“写作之颂”两注(n\textsuperscript{o}~101、102)。)然而,带着回顾,我意识到:这份分散,我当时还省不得——必须有某一份激情、某一份饥饿在我身上烧尽;而一路走来,我继续借着那些我做她们情郎的女人,学习那些我在已往的生活里不曾学全的东西。到了我当时的那个地步,我怀疑单凭对这段往昔的禅观,它本会教得会我。}。在其后将近两年里,我的禅观仅限于几场应景的、分量极有限的反思——当我面临某个尖锐冲突的情势、急切地需要看个分明之时。直到1979年8月之后(禅观发现后将近三年),才开始对种种现成观念作那场“大扫除”——尤其关于我父母、关于我自己的;它们继续堵塞着我,挡住我看这个我生活于其中的迷人世界。对我父母生活的工作,吸住了我七个月,直到次年3月。那时,我五十二岁在望。正是凭这份工作,那份我谈过的自主——它在某种意义上,三年间只停留在“潜在”——变成了全然现实的、完整的、不可逆转的。也正是凭这份工作、且唯凭它,我才能够\textbf{爱}我的父母,取这个词的完整意义;也就是说:也\textbf{接纳}他们所是的东西、或他们曾是的东西,连同这一切所牵涉的(而我那时开始隐约望见的)、尤其牵涉到我——他们的儿子——的东西。

倘若我感到需要做这份工作(注~128$_{1}$),并且能够做成它,那是因为三年之前,我已懂得接纳生命的这份赠礼——它在我出生之时授予我,却被我拒绝了四十年——我自身统一的赠礼。或者说:因为我已懂得\textbf{接纳我自己的天性}。正是凭对我自己的接纳、对我自己的爱,我才能够接纳、爱我的父母\fnote{129}{(*)}{这与“接纳(阴的苏醒(2))”一注(n\textsuperscript{o}~110)结尾处的反思相接。}。

我还可以说:正是唯凭这份工作,“对父母的冲突”才得到“解决”——这份冲突,几年前我甚至不曾疑心它的存在,而那时\dmark{},我的父母双双谢世已二十多年\origpage{584}。诚然,自我婴幼年起,我对父母姿态的底色,曾是一种敬佩的尊重、一种抬举、一种毫无保留的认同;而在他们死后,是对他们这个人、对他们音容的一种默默的尊崇。这不是人们惯常用“冲突”一词所指的关系——“冲突”暗示一种对抗、敌意的底色。在这份由我这个人带给他们的抬举里,我的父母当然自得其乐:他们觉得这样极好、合乎秩序——而天下的父母,不愿处在他们位置上的、或不为此自庆的,想必寥寥!只是在这份对我父母的工作之后、尤其在随之而来的对我童年的工作之后,我才能充分地、充分知其所以然地看清:我与父母之间那场田园诗般的关系,在多大的程度上曾是\textbf{假的}、造作的,不“\textbf{真}”。它之所以能维持,只凭着在一幅感人的图画上执拗地拿橡皮擦去一大堆“不入画”的东西:包括那些难受的时期(恰恰是尖锐的对抗,常常被感到是一场\textbf{撕裂}),或那些慢性的“越轨”,它们在我母亲与我之间的关系里,以同样毫不容情的规律性重演(即便频率低些),一如当年她与我父亲之间。更不必说那些在意识层面整个逃过我知觉的事情:比如那道我在八岁上、给父母画下的“大十字”——那时我在异乡已住了两年,母亲一年三四封匆匆的来信,便是他们二人全部的生的音讯……

而那深层缘由、那份\textbf{真正的}缘由——它使我把1933年夏(五岁)到1979/80年冬(五十一岁)之间我与父母的关系称作“冲突性的”——并不在于:这四十六年间有过一些冲突,把我和他们之一、或他们二人对立起来——无论这些冲突频繁或稀少、暴烈或潜伏、有意识或无意识。而毋宁在于:这份关系不曾被\textbf{承担},也不可能被承担(照它当时那样,我是说,不经深刻的转变)。它之所以能照我所活的那样被活、照我所看的那样被看,只凭对我的认识官能与领会官能的一份恒常而顽强的\textbf{压制};凭一份\textbf{拒绝}的执拗\origpage{585}:拒绝去认识这份关系的真正本性,或至少认识它的若干要害的侧面——那些侧面,本质上牵涉着我的双亲与我本人,牵涉着我所喂养的关于我们的那幅影像。换一种说法:这份关系当时所取的形态,是凭一份执拗的、不停歇的\textbf{逃遁}而自我延续的——逃开一份十足可触的实在;而那实在,同样执拗地一次又一次要让我认识它,可在我父母生前,我从不曾真正从中汲取种子。那清晰而无可否认的、把我和他们之一对立起来的冲突的那些情节——有时是撕裂性的——只是“我与父母关系之‘冲突性’”(即这份发生\textbf{在我自己身上}的压制与逃遁)的诸多或隐或显的征兆中的几桩。

换一种说法:与他人的“冲突性”关系,取这个词的深层意义,便是那份“分裂了的”关系——它凭一个压制、逃遁实在的进程而原样自我延续,又反过来助长这些进程在自己身上延续。关系中“冲突”、“分裂”的征兆,既可以是对抗品格的,也可以是效忠品格的;既可以是批评、乃至低估或轻蔑的故意,也可以是赞许或钦慕的故意。

就这样,未经寻索、也未可预料,我又回到了人们或许会称作我的哲学“老毛病”上:人与人之间的冲突,不过是各当事人身上冲突的“征兆”;或者说:社会中冲突的“源头”,是个人身上的冲突、分裂。(而父母,在这整个论说里,终究销声匿迹、不留痕迹了!)

这幅世界之景,看似全然忽略了那幅更简单化、也远为流行的景:两人之间的冲突,是双方各自“利益”或欲望冲突的结果——那些欲望“客观地”相互对立,即一方之满足,唯有损及另一方方能达成。这是举世公认的看待方式,无论涉及两个不同个人之间的冲突,还是同一个人内部的冲突。比如(在第一种情形里),这些不相容的“欲望”,在双方身上,可以是支配的欲望、定调子的欲望、掌舵的欲望——这当然是最常见的个案之一,包括在父母与孩子之间(也全然一样,在妻子与丈夫之间、或情郎与情人之间\dmark{}\origpage{586})。当然,我并不否认这种看待方式有任何实在性、任何效用——至少在某些个案里。但我看到:它只涉及一份浮面的实在,而一份更深的实在整个儿逃过它。为在这层意义上举一个例子,我指出:支配的欲望(或炫耀的欲望,或一般地,把自己抬到他人之上的欲望),它的根恰恰在那份“对己之轻蔑”、那份“对己之不识”里——方才已谈过;人们试图以种种足以\textbf{搅浑}、\textbf{补偿}这份对自己秘密的低估的姿态与行为,从它那里逃开。这样,在对立的欲望之“客观”冲突之外,人们在这个个案里看到:个人内部的冲突显出轮廓,作为如此这般欲望的缔造者——那些欲望,除了唤起并滋养对他人的对抗之外,别无他能。

当然,凭这几句评注,我并不打算穷尽冲突的两个侧面之间那份微妙而重要的关系的问题——我想把它们称作“浮面”侧面与“深层”侧面——而这无疑也不是地方。毋宁说,我感到需要回到对父亲的冲突、对父母的冲突这个主题上——我方才正要离开它。我一度可能给人(甚至让自己也一度信了几分钟!)这样一种印象:对某位父母的冲突,同对皮埃尔或对波莱的冲突,是一码事。然而我分明知道远非如此!我分明知道:\textbf{对父亲的冲突、对母亲的冲突,正是我们自己身上冲突的心脏}。

方才我已在这个意义上谈到我的“贴身信念”(我也乐意称之为我身上的一份\textbf{认识}、一件透彻懂得的事):在一个自身未曾分裂的人身上,对父母的冲突已经解决。这份认识,我说过,首先(我相信)来自我自己与父母关系中冲突之解决的经验\fnote{130}{(*)}{关于这一点,见紧随其后的那条脚注。}。换一种说法,便是:\textbf{对我们父母的接纳}(即对父母冲突的止息)\textbf{属于对我们自己的接纳的一部分}。他们(相对于我们而言),既是我们的\textbf{本源}\dmark{},也是我们的\textbf{制约}(或至少是其中好大一部分)\origpage{587}。前者(我们的本源)与我们的这个人不可分,无论我们的行程与命运如何;后者(我们的制约)深深扎根于我们,因而一如我们的本源,属于我们的这个人。否认我们母亲或父亲的真正实在——无论这拒绝以对抗表达,还是以效忠表达——也就是否认我们自己的一份要害的实质、否认我们的一生(就我们记忆所及的而言)……

还不止于此。正是经由我们的母亲与我们的父亲,先于一切他人,那曾在他们各自身上的冲突传递到了我们身上。(这正是片刻之前那个凝练的词“我们的制约”所表达的!)他们由此与我们身上的冲突相连,比世上任何别人都更贴近。而这冲突在我们身上的第一份外部的投影、一切投影中最古老也最关键的,便是对我们母亲与父亲的冲突。于是,在我看来,我们身上的冲突,与对我们双亲各自的冲突,不可分解地系在一起——它们宛如同一个冲突。方才我表达过那份“贴身信念”:当我们身上的冲突得到解决时(或至少,当它在根子上、在“阴对阳”的分裂里得到解决时),我们对父母的冲突也同样得到解决;或者换一种说法:我们身上冲突的解决,要经过对父母冲突的解决。但我确信,反过来说同样为真:一旦对父母的冲突得到解决,我们身上的冲突也随之解决\fnote{131}{(*)}{在这里,我或许给人造成一种印象:仿佛我在摆“那个已在自己身上解决了冲突的人”的架子。千真万确,我是毫无保留地说:对我父母的冲突已彻底解决。同样真确的是:我这个人身上的冲突,仍以种种方式让人感到,它没有消失。这件事,在《收获与播种》的每一页里,想必都显而易见;而我也不止一次有机会,在这个或那个具体个案里指出它。这看来便与这条脚注所评注的那句断言——“一旦对父母的冲突得到解决,我们身上的冲突也随之解决”——相矛盾。然而,在某种意义上(即我写这几行时所着意的那种),“冲突已在我身上解决”千真万确。至少,这份冲突里的某个本质的东西、在它根子上的东西,确确实实解决了——凭着这份对我自身统一的认识、这份对我自己的接纳。倘若把冲突比作一棵根深干壮的树,那么可以说:当根被切断、或已枯干时,树已经死了,只是凭习得的惰性,树干与几条主枝还立在原地,等待慢慢干枯、慢慢崩解。我能分明感到这些年来冲突这份渐进的“干枯”——它像一份曾经强而顽的掌控,在渐渐松脱。《收获与播种》的写作,在我看来,是这个进程中的一站,是过去八年里许多站中的一站。另一幅试图描述同一实在的意象,是一份徐徐铺开的深沉的宁静——像深海的宁静,不为搅动水面的涡流所动。关于这一点,我在“重逢(阴的苏醒(1))”与“接纳(阴的苏醒(2))”两注(n\textsuperscript{o}~109、110)里有更周详的表述。}。正是由此,我在与我们父母的关系里,看到一份在我们灵性历险里的\textbf{枢纽角色}——一份独一无二的角色:它不属于任何其他亲近之人,无论是配偶、孩子,还是朋友、老师或学生。

\asterisms

\notehead{128$_{1}$}{}
\phantomsection\label{ch18:local:n34}

\dmark{}(12月1日)\fnote{132}{(*)}{本注由前一条注“父母——或冲突的心脏”(n\textsuperscript{o}~128)的一条脚注生出。}“去认识我的父母”这件事之于我的重要性,是一个梦向我启示的——那梦来到1978年10月28日。那是关于我父亲弥留的一个梦。这场弥留,在痛苦挣扎的日日夜夜中迁延,周围是他身边人忙碌的冷漠——凭所有人的默契,他已被当作“已死之人”——“这宛如一道判决,它斩断一切怀疑,把他的死变成了事实。”我醒后把梦记述了下来;但在随后的三个月里,我回避了对它的一切反思,以致让它沉入半遗忘的昏影。总之,我当时把这个梦里对我谈到的、我父亲之死给“埋葬”了;而在这个梦里(它唤起我醒觉生活的一个关键侧面),我正“埋葬”着还活着的父亲。对这个梦的讯息——它那样清晰、那样透彻,美得撼人心魄——曾有过相当强大的抗拒。抗拒是在第一夜禅观结束时消解的,那是次年1月31日,以一场对梦之意义的执拗禅观为始;其后的三周里,又有四场禅观相继。

这个梦使我懂得:我与我父亲、我母亲的关系,是一份凝定的、“死去的”关系——它与一份活的实在隔断,对那份实在的知觉被压抑着——一如(在梦里)对那场被宣告为“子虚乌有”的弥留的知觉被压抑着,连同由之而来的那个自发的行动:去救助那个痛苦地、被所有人抛弃地、为活命而挣扎的人。

\dmark{}终结我内心这份隔绝的第一件事,便是去认识我的父母\origpage{589}。当时我丝毫没有料到这项任务的规模:我幻想“几个小时”便能“直捣题心”!经由我的童年、尤其经由它来认识我自己的想法,当时连我的边儿都没擦过。这份需要是后来才感到的:它将自发地从那场我正准备启程的旅程中流出。旅程只在六个月之后才开始,1979年8月——缘于“乱伦颂”情节那场长长的枝蔓(它在许多方面绝非无用)。(关于它,见“行为”一注(113)。)

与1976年10月18日的那个梦(它触发了“重逢”)一道,这个关于我父亲弥留的梦,是对我生活进程起过最强烈作用的两个梦之一。对它的讯息的抗拒,依我看,要强大得多。前一个梦的讯息,在醒来后的几个钟头里便被接收;而后一个的,被推拒了好几个月。它只是在九个月之后才开始兑现:凭我动身去作一场发现之旅——它延续到今天……

直到最近这几天,这个梦的意义与我正在眼下这场反思里力图穿透的葬礼之实在之间的关联,才来到我心头。这场葬礼——我在其中扮演“头号亡人”——不久前曾向我显为一场“事物的回转”(见同名注(73))。这一回,我又看到一场“事物的回转”,却是从一个全然意想不到的角度。因为在葬礼里,我相继扮作“父亲”与“母亲”。我从不曾想到:我自己也曾处于一个儿子的类似境地——“埋葬”活着的(哪怕是象征地、或凭默契)他的父亲或他的母亲——恰恰相反!而且我当时有十足的理由相信相反的情形;那些缘由,我在“屠杀”那条注的结尾头一次唤起(诚然,是在\textbf{屠杀}父亲的语境里,而非下葬的语境)。(在“天真(阴与阳的联姻)”一注(107)里,我更周详地回到它。)在“屠杀”那条注里写下关于我最初童年的这两段时,我肯定给了人(甚至我自己当时也处在)这样一种印象:我与父亲的关系,一生都没有冲突。对这份关系浮面一看,也会这样暗示。但已经在这条被评注的注里——“父母——或冲突的心脏”——在那里我并不止于这类皮肉的印象,事情分明显出\dmark{}:全然不是如此\origpage{590};那幅世界之景(它确确实实是我的,直到1979年1月31日)是我在大半成年生活里乐于喂养的诸幻觉之一。这份幻觉,从我终于费心检视那个关于父亲弥留之梦的意义的那一刻起,便清晰显形——那是迄今为止生活赐予我的所有梦之中最\textbf{美}的一个。这个梦以惊人的写实,呈现了冲突加于我与父亲关系之上的掌控——它也让我活出了这份冲突的\textbf{解决}。冲突之解决,凭的是我身上与那道宣告我父亲之死的共识之间的一场\textbf{决裂}——决裂旋即为\textbf{别的东西}敞开了门;也凭我父亲的一个爱的动作,向我示意:他已听见我那被扼住的喉咙无法向他迸出的呼喊……

这个梦的亲身经历——它是那份凝定的父母关系(它旋即重新活了过来……)的惊人寓言——与我已探测了将近九个月的葬礼之实在之间的深刻亲缘,此刻以一份明证之力呈现在我眼前。惹人注意的是:在这整场漫长的反思里,直到最近几天,这份亲缘的念头竟不曾擦过我的心。我是凭最大的偶然才“撞上”它的:当时我打算在一条脚注里,顺便指出某一个\textbf{梦}——在八年来宛如天助灯塔照着我路的众多梦之中——又一次(在对父母的反思的启动中)所扮演的角色。这个打算的效应,是把我与这个梦的亲身经历及其实质稍稍重新接上——我离穷尽它们还远得很。一旦这份联系重新接上,在此语境里,它与葬礼的亲缘便不可能不显形。

诚然,这份亲缘,眼下只涉及某个“死结”而已;而在这个梦里、在它所转录的实在里,既有死结,也有它的解决。何况那份解决——梦曾让我活出它,当夜我便知道了它的滋味与力量——要使它也成为一个在我醒觉生活里、在我与父亲及母亲的关系里被活出的实在,那是我、唯有我的本分。我可以去做,也可以不做——而有几个月,我选的恰是第二条路!今天——那场解决之后五年——在这份某种意义上对称的情势里,情形肯定依然如此:我在其中牵涉,而我扮的是被一道共识-判决所下葬的父亲——想当年,我可是那个恭恭敬敬活埋自己血肉之父的儿子\dmark{}\origpage{591}!而或许这一回,仍将是一场对我的亲身经历之意义的禅观——在此即对这场葬礼之意义的禅观——将解开我身陷的这个新的死结,并或许再消融我过去之重的一份余量。

至于这场禅观会不会对我之外的什么人有用处——对这场葬礼的某位当事人也许(在这场葬礼里,被下葬的不止我一个,而赶赴仪式的下葬者更是不计其数)——这轮不到我操心;我看到的别人身上的这个或那个死结解不解脱,也轮不到我操心。那是他的本分,我的本分已经够多了!但倘若趁我在世,它果真解脱了,我肯定是头一批得到消息的人之一,而且我会为此高兴……

\subsubsection*{(e) 敌对的父亲(3)——或阳葬阳}
\tocsubsub{(e) 敌对的父亲(3)——或阳葬阳}
\label{ch18:local:h39}

\notehead{129}{}
\phantomsection\label{ch18:local:n35}

显然,在前面的几页里\fnote{133}{(*)}{即注~n\textsuperscript{o}~128的几页,本注是它的直接延续。},我对\textbf{对父母的冲突}这个主题,不过刚刚擦了个边;甚至连对父亲的冲突——我的出发点——也谈不上。我从那里出发所追随的种种联想,似乎把我带离了它,而非深入它。在我方才关于对父母之冲突所说的话里,母亲与父亲的角色可以互换;同样无差别的是,这几页里所说的“我们”,指一个男人还是一个女人。然而,在我们与父母的关系里,母亲与父亲扮演的角色远非对称;而他们各自扮演的角色,关键地取决于“我们”是男孩还是女孩(后来长成男人,或女人)。

在这个个案里,对父亲的冲突(它以对父亲的象征性下葬、甚或以对他的屠杀来表达),首先令我感兴趣的,是那些我知道曾积极参与我下葬的人的个案——他们\textbf{全是男人}。于是,在我之结构的历程中,父亲是那个人们\textbf{与之认同}的人、那个人们\textbf{以他为模型}的人——在他的与他人的关系里(尤其与女人的关系里),也在他与自己的关系里。这份认同\dmark{}很少有不带一份相当大的“越轨”便能完成,而对父亲的对抗\origpage{592},便是它最顽的痕迹之一。此处不打算遍数这些偏差,以及哪怕最愿意以爸爸为榜样的小男孩也常会遇到的种种不顺;也不打算探究它们在与父亲的关系中往往如何表现。况且,我在这方面的亲身经验,实在太不合典型,以致做这样一份清单,我或许是最不称职的人:凭我自己的亲身经历,我并不贴身感到那些主要个案中任何一个的来龙去脉与特殊的“滋味”\fnote{134}{(*)}{可与“屠杀”一注(n\textsuperscript{o}~87)结尾处的反思对读。}。我这里的经验,主要是间接的:凭我对周围的观察,而首先,凭我的孩子们与我的关系。

在“越轨”的特殊品格之外、在由此汲取的对父亲的怨怼与积愤之外,却有一个共同的侧面,我在许多场合强烈地知觉过——那时,任何“解释性”的故意全然缺席。这便是:男孩或男人对父亲的对抗——父亲多少做过他的模型,而他正复现着这个模型,以“正”的或“负”的方式(凭模仿,或凭反着来),无论他愿不愿承认——这份对抗不是别的,正是一份对着\textbf{他自己}的对抗的格外雄辩、格外关键的侧面。更确切地说:经由对父亲的(或多或少明白说出的)拒斥,它是\textbf{他对自己一部分的拒斥}的外部征兆——是那样一些东西的征兆:凭着它们(在他不知不觉中,或违背他某些有意识或无意识的抉择),他像极了他所否认的那个模型——他的父亲。

这样一来,我稳稳落回了双脚——我看到那份被预感到的联系清晰起来:在“对己之轻蔑”(或“对己的拒斥(或不识)”)与“对父亲的对抗”之间——但我是从一个意想不到的一侧落回来的。我正准备在对父亲的对抗与自拒之间,找一层或多或少直接的联系——自拒取\textbf{他自己身上}女性之拒斥(或“下葬”)的形式。结果却似乎是:我落在了对\textbf{男性}的拒斥上——按“好逻辑”,这本该在我意料之中!然而我分明知道:这份拒斥,在男人身上不如他身上女性之拒斥那么显眼、藏得更深(我主要有机会谈的是后者),却几乎同样不罕见,压在他身上的分量,同样沉甸甸。它常常叠加在另一份之上,以致无论一个人的我以什么方式结构——以阴的颜色还是以阳的颜色——都注定不容于他自己!或者换一种说法:这份对父亲的拒斥、或对自身之中使我们像父亲的“男性”“阳刚”之物的拒斥,常常\textbf{与}毫无保留地采纳一套“阳”的、“大男子气”到底的价值系统\textbf{结伴而行}(既然“阴”的配重已被否认)\fnote{135}{(*)}{(11月29日)至少,就我所知的那些个案而言,这是远为常见的情形。}。

一个念头来到我心头:这份矛盾(一旦白纸黑字说出、写出,着实骇人!)无疑也是那场毫不容情的\textbf{竞争}里真正的\textbf{神经}——那竞争,是我们这个超级大男子气社会的特征之一(无论在科学的高层,还是在任何别的地方……)。因为,倘若“往上爬”与“超越”是十足十的超级阳价值,这些价值却不至于以这样的烈度被内化、其付诸实践也不至于以这样的暴烈进行(哪怕在“高层”里,那暴烈是压低了嗓门的……)——倘若在那位占了上风的对手身上——我们正要超越乃至排挤的对手——我们没有同时看见“父亲”那副可畏的影子显出轮廓:他既受钦慕又遭嫉妒,还被秘密地仇恨;他先我们而在,而他的存在本身,就我们记忆所及,一直是我们生活里的\textbf{大挑战}。

\subsubsection*{(f) 箭与浪}
\tocsubsub{(f) 箭与浪}
\label{ch18:local:h40}

\notehead{130}{}
\phantomsection\label{ch18:local:n36}

(11月19日)我又感到急不可耐,要接着把反思继续下去。老实说,已经一个星期了(自11月12日那条注“激烈泼辣的妻子(阴阳的颠倒)”(126)起),我天天感到自己就要进入“题心”了——就要来到我应承过自己的那幅葬礼总图跟前:它将把反思进程中析出的那些局部“联”统一起来\dmark{}——也一个星期了,那个“点”却一天天被往后推\origpage{594}。每天写完注(时候不早了,总得停下上床),我分明感到自己做成了一份非做不可的工作,“前进”了一步——却同时觉得,我想抵达的那个“点”也退了同样远!这里明显的诱惑,是一口气写下去,直到抵达那著名的“题心”。但经过这三年里的几次“健康事故”,我也分明知道:这正是该避开的失着。

况且,在心底我分明知道:我已经全身在那个“题心”里了。只是,我急不可耐,恨不得立刻把它整个摸透。这份抵达任务尽头的焦躁,这股朝着某个“点”、某个“题心”的冲劲——它在前方被我强烈地知觉着,或近或远,反正一样——这份“目标”对我的吸引,把我向前抛射,如一支直扑靶心的箭——这个在我看来最强烈地“\textbf{阳}”的侧面,刻画着我的\textbf{工作之外}的存在方式。它是“\textbf{东家}”的一个惹眼的侧面,是我身上受制约的、习得的东西的侧面。在我幼年所知的一切里,没有任何东西能预示这个品格——它在我的童年后期才显出,却如此强烈地打下了我整个成年生活直到今天的印记。

而在工作本身之中,这个侧面似乎几乎消失。我有这样一种印象:它残存的些许痕迹,此地彼处,不多不少,正是“东家”在工作中偶一为之的、不显眼的插手之标志(说实在的,他在这里无事可做!)。工作本身,顺着那匠人——他经我的双手,按他自己的节拍劳作——的心意,循着一股全然不同的气息进行。那急躁的烈性,退到一份宁静、平和而执拗的东西面前。不再有箭,急匆匆扑向靶心;只有一道浪,铺向极远,不知往哪里推进——那推动它的、活动的力量把它带向哪里,它便去哪里——一道浪跟着一道浪,再跟着又一道……这个运动里没有迟疑:在每一处、每一刻,都有一个它自己的方向托着它、把它往前拉。每一刻都有一场推进——说不上朝着什么;有一种“工作”在做着,在一种不知努力为何物的运行之中——却没有目标。“目标”这个念头本身,在这里显得出奇地乖谬——人们想把它放在哪里呢?!目标消失了,一如\dmark{}箭消失了。\origpage{595}倘若还有箭,那也不是一支颤动的箭,呼啸着扎进靶心、插在靶上折损自己——而是在这片相互追逐的浪的活动之体的\textbf{每一}处,都有一个毫不含糊的运动与力量,都有一场推进之中的方向,其精确与分明,不亚于一支箭——一支看不见却威严的箭,标示着这方向、这力量、这运动。

这样,依我看,在我的工作里,我同样“阴”,同样“海与运行”,达到了可能的极致。我相信,我一生中一切发现的工作、一切我满怀激情投身其中的工作,都是如此——首先,是我的数学工作与禅观的工作。而今,这幅意象不期而至,威严而骤然,描绘出我如何感受这份工作;我觉得,这幅意象同时也描绘了\textbf{我生活的运动}——自与我自己的重逢之日起,或许更早,或许早自我那场“有益的拔离”——拔离一个安乐的巢\fnote{136}{(*)}{见同名注(n\textsuperscript{o}~42)。}。至少,它描绘了我生活深层处的“如何”——即我昨天那条注的一条脚注里(就在几个小时之前)谈到的那个“宁静”:一份不为水面上的纷扰所动的宁静。在这份深沉的宁静里,有运动、有推进,却没有目标——目标消失了。

而今我也记起:3月份,正是这同一幅意象来到我心头,当时我谈到我那两种激情——禅观与数学——的显现,说它们像“一道道相互追逐的浪的高低起伏,像一份宽阔而平和的呼吸的气息……”\fnote{137}{(**)}{见“我的激情”一节(n\textsuperscript{o}~35)的结尾,这几行便摘自那里。}。八个月后的今天,我相信在这些意象里认出了我存在的自发运动——在它最自发的东西里、在我身上真正原初的东西里——在那个渴望认识的孩子身上,在“显得体面”的操心、与“求变成”的饥渴还未触及他之前……

\subsubsection*{(g) 冲突之谜}
\tocsubsub{(g) 冲突之谜}
\label{ch18:local:h41}

\notehead{131}{}
\phantomsection\label{ch18:local:n37}

\dmark{}(11月20日)昨天几乎整个晚上,都花在重读前天的注上\origpage{596}:一路修改,重打一页实在写得太挤的纸,写脚注(前天就已打算好的)——而已经午夜了!可我还急着当晚再往前走一点,哪怕一点点;于是又坐回打字机前,接上昨晚中断的“线”。接着,来的却是全然另一件事——箭与浪的意象。长久以来,我在箭的意象里认出自己,而浪的意象,在我看来说的是一种与我大相径庭的气质。这正是这场关于阴阳的反思进程中显出的意外之一:恰恰是浪的意象,以最触目、最贴切的方式,表达了我的存在里占上风的那个“基调”——当“东家”不在场、或至少退到别的东西后面时。意象升了起来,仿佛它早已在那儿备好,只等着那些终于使它成形的词。词不慌不忙、毫不犹豫地来了,而我所做的,不过是尽力\textbf{描述}——尽可能忠实,不省略、不歪曲——那尚处于弥散之感状态的东西。描述完成,已近凌晨两点。当夜我重读了这两页,几乎无须修改。最棘手的一段,是我试着描述那份对“箭”之连续无限的直觉——它们合拢起来,宛如一个力的“场”。这是一个强有力的观念,却不肯轻易让语言把它唤出。可我感到,这是整幅意象的一个重要侧面:“阴中之阳”的侧面。浪里有“箭”,有一股托着它向前的\textbf{冲劲}——循着一份它特有的运行,而不是\textbf{一支}箭的运行,毋宁说是整整一簇、一\textbf{连续}簇的运行,它柔韧地复现着浪的那份运行。我也分明知道:在我的工作里,我\textbf{同样}是“箭”;只是我之为箭,循着一种与我迄今想象不同的方式——因为我从不曾费心以某种注意去看这份工作、让自己浸入它,仿佛它是另一个人,以便知觉它特有的调门。八年来我一直在禅观,却至今不曾这样做;这无疑是因为,我在不知不觉中,始终是一个积习已深的故意的囚徒:认同于我身上的“\textbf{东家}”,而不是认同于那个匠人-孩子;也就是说,当我说“我”时,\origpage{597}首先(甚至常常排他地)想到的,是那个“东家”在前台时的那个我。而那些时刻,差不多也就是我工作之外的时刻。

教学的必要与波折(以及其他缘故),终于在禅观发现之后,把我的注意力引向我工作的\textbf{某些}特征——即那些我感到具有普遍品格的特征:它们应当存在于一切创造性的工作、一切发现的工作之中\fnote{138}{(*)}{我想,头一次写出、并唤起这些特征中某些特征的文本,是1978年10月的那份“以代纲领”(11月6日那条注“陌生的美人”(n\textsuperscript{o}~130[原书如此,当为120])里暗示过它)。那份文本之后,我没有费心把这方面的观察白纸黑字显豁、深化,直到今年《收获与播种》的反思。它的最初八节,基本上就献给这个主题,还不算反思进程中此地彼处的许多别的评注。}。但在这场阴阳反思之前,我还不曾想到去辨认我自己的工作里那些独特的特征——它们使我的工作区别于任何别人的。这些特征之一,依我看是所有特征中最要紧的,终于在11月8日那条注“上涨的海……”(122)里得到圈定。那条注里最初唤起的意象——在一桩有待证明的猜想这个典型语境里——在昨天那两条注里被重拾,却换了一道光照,脱开一切特殊语境。

我终于接上反思的线,从前天中断的地方续起。我起笔时\fnote{139}{(**)}{见“父母——或冲突的心脏”一注(n\textsuperscript{o}~128)。},本打算试着圈定对父亲之对抗的深层缘由——越过人们对父亲可能怀有的那些具体怨怼。顺着强劲涌现的种种联想,我离开了这个打算,主要被引去谈对父母的冲突——父亲或母亲,不分彼此。这份“冲突”,既可取效忠的形式(我的情形如此),也可取对抗的形式。自对我父母生活的工作以来,这份“\textbf{对父母的冲突}”向我显为真正\dmark{}“处在冲突的心脏”——我们身上的冲突的心脏。解决后者,我确信,不多不少,正是解决对父母的冲突,也就是说:\origpage{598}从他们那里自由,在灵性上充分自主,走\textbf{自己的}旅程……

回到男人身上对父亲的对抗,我重新接上了一份直觉——近年来它屡屡强加于我:这份对父亲之对抗的深层意义,是对我们身上使我们像父亲的那些东西的拒斥——是对我们这个人“\textbf{阳刚}”的方面与特征的拒斥。我把昨天反思(其实\fnote{140}{(*)}{实际上,这不是昨天那条注,而是前天那条注——我正准备接上的正是它。}是前天的)的最后这部分,单独抽出作一条注,命名为“敌对的父亲(3)——或阳葬阳”——这名字同时也暗示它与“敌对的父亲(1)(2)”两节(n\textsuperscript{o}~29、30)的联系:那“敌对的父亲”的主题,正是在那里头一次出现。

这样,前天反思开头所谈的葬礼的那个侧面——“对己之轻蔑”、或“对己之不识”、或“自拒”——显为一种联接的痕迹,或更好地说,一道“\textbf{铰链}”:铰接在前面的两联之间——“超级母亲——或‘女性’的下葬”之联,与“超级父亲——或父亲的屠杀与下葬”之联。这铰链的品格,一经清楚地知觉到下述两点,便显出来:在前一联里,“女性”在一切之先,是“\textbf{我们身上}的女性”(11月10日那条注“阴的葬礼(阳葬阴(4))”——“超级母亲”之联登场的那条——里,确实已被知觉);而且,“父亲”首先是“我们身上的男性”的象征性替补。于是,这两个相关侧面,显为完全对称的两联,对应着“自拒”的两个显然的“图式”:拒斥我们身上的“女人”(别名母亲),拒斥我们身上的“男人”(别名父亲)\fnote{141}{(**)}{我提醒一句:这两种“对称的”拒斥,叠在同一人身上的情形,绝非罕见。鉴于我们社会对阴的贬抑,无论如何,阴的拒斥以或多或少显著的形式在场的情形,想必相当常见。所以,我倾向于在对父亲的对抗里,看到一份对阴\textbf{与}阳的双重拒斥的(至少是盖然性的)征象。}。而对父母的冲突这个主题\dmark{}——它是对母亲的冲突与对父亲的冲突这两个不同主题的某种合取或叠加——也显为一道铰链。\origpage{599}或者更好地说,一如昨天反思里所看到的\fnote{142}{(*)}{见倒数第二条脚注。},这个主题显为与自拒的主题不可分:两者是同一份未分裂的实在——\textbf{我们身上的冲突}——的两个不同侧面。

在这一切之中,最初的那个打算——“圈定对父亲之对抗的深层\textbf{缘由}”——似乎始终悬着。我可以说:对父亲的对抗,是对抗自己、自拒所取的\textbf{形式}之一。于是,最初的问题仿佛一分为二。一方面,在哪些“缘由”之下,自拒在某些个案里取这份特殊的形式?探测它,也就或多或少要进入若干不同类型的典型情势——那些易于引发这份对抗的情势。

另一方面,我们回到那个更深、也更关键的问题:自拒的“\textbf{缘由}”,也就是我们身上冲突、分裂的缘由。我相信自己至少把握住了那份共同的\textbf{机制},冲突凭它代代相传:我们对自身的拒斥,不是别的,正是我们最初几年里周围环境对我们的拒斥之内化——至少,是我们身上某些侧面与某些冲动之被拒斥的内化;那些侧面与冲动,构成我们原初存在、我们创造官能的一个本质部分。我在“阴阳之钥”的“拒受与接纳”部分里触及事情的这方面(以及其他),尤其在最初两条注“失乐园”与“循环”(116)(116$'$)里。

把握住冲突传递的这份共同“机制”,却丝毫并不等于:已经懂得我们身上、并(经由我们)在人类社会里冲突的\textbf{缘由}。\textbf{为什么},古往今来、无远弗届(凭着世代传到我们手里的一致证词),“社会”竟不能容忍它的成员是一些\textbf{完整}的存在?也就是说,是一些充分\dmark{}握有\origpage{600}自己创造官能的存在——他们不必花费高昂的代价,压制自己之所是的一部分:那份被看得如此可耻(或如此可畏……)的部分,以致最好假装不知道它的存在,并默契地判定它\textbf{不存在}……

这在我,是存在的伟大奥秘之一,或许是最伟大的奥秘\fnote{143}{(*)}{这份提议纯属主观,它只反映这样一件事实:在“存在的伟大奥秘”之中,正是这一个,我以格外强烈的方式感到——那方式,超出单纯的智性好奇。它是唯一一个在我身上唤起一份\textbf{欲望}的:探测它、认识它、认识它的“最终答案”(在我这有限的官能所能认识的限度之内)。区别正如数学里的情形:一边是我“很有感觉”的未解决问题(我可以当场就投身其中),另一边是我按术语的意义“理解”的问题——我(浮面地)知觉其分量,它们却“不叫我或热或冷”。Riemann假设属于后者(无疑归因于我对解析数论的极端无知);“费马大定理”直到几年前也属于后者。是我的“阿那贝尔式”反思改变了我对后者的心态,尽管我对它所催生的那些工作,无知一如从前。}。

曾有一段时光——几年前还如此——我对压制与冲突这一普遍实在的姿态,是一种战斗的\textbf{反叛}:反叛那把“\textbf{利剑}”——它声称要把凭本性当为一、\textbf{本}为一的东西剖成两半。五年前我写《颂》时,我的心态犹是如此\fnote{144}{(**)}{这一情节,《收获与播种》里谈到过几次,最后一次是在“行为”一注(n\textsuperscript{o}~113)里。}。是其后那场关于我父母生活的长程禅观工作,改变了这份心态。凭这份工作——它日复一日让我贴身接触我父母身上冲突的种种显现,并让我耐心地由显现上溯其意义与缘由——我终于感到了冲突的\textbf{奥秘}。反叛的姿态消失了,仿佛从不曾有过。它曾是一种皮肉的反应,一份单纯的能量耗散。一场反叛——反谁?不是反某个人、或某群人,不是反那著名的“他们……”!我们全在同一条船上,而人们在这儿已经待了一百万年或两百万年了……反“上帝”反叛?就差这个了。

\dmark{}在心底,我分明早就知道(我甚至说不出从何时起,尽管有很长一段时间我装作不知道……):这世上的万事万物,都有它存在的充分理由;甚至,倘若人们懂得事物的底里,万事万物就其本身而言,肯定都是\textbf{好}的\origpage{601}。死亡与死亡的“彼岸”(倘若真有这样一个彼岸),也属于这些事物。这是一个奥秘;倘若我在这件事上有什么“\textbf{信}”,它绝不在关于彼岸之存在(或不存在)及其种种细节的“信条”里,而只在于这份单纯的确信:事物如其所是,已然圆满——包括一切涉及死亡的,也包括一切涉及诞生的,而诞生同样神秘。

许久以来,我却把“冲突”排除在这些事物之外——我把它当作一种“越轨”、一块不可容忍的污渍、创造的大合唱里一个顽固而乖谬(乃至令人义愤)的“走音”。只消我终于贴身去认识冲突,而不是浪费自己去装模作样地与它搏斗,我与它的关系便深刻地转变了。

死亡与“死后”的奥秘、诞生与“生前”的奥秘,并非我们这个物种所独有。它们所引出的问题,对一切活物都有意义——或许对一切事物都有意义,从电子到星云。冲突的奥秘,却依我看为人所独有,为人类这个物种所独有\fnote{145}{(*)}{(12月3日)人们或许会(有理地)反驳我:冲突——以个体之间或个体群之间的攻击性与交锋的形式——也存在于我们之外的其他物种之内。当我在这里谈“冲突”时,我想的是它在人类社会里所取的那种特殊形式,尤其是它与个人内部的\textbf{分裂}与\textbf{压制}之间的深刻联系——对其存在之大半的压制,尤其对其知觉实在之手段的压制、以及知觉本身的压制。种种形式的压制,在我看来,都扎根于那在我看来最要紧的压制之中:所谓“性”的压制——它灌输对自己身体的羞耻、对身体的功能与冲动的羞耻(或至少,对其中某些功能与冲动)。就我所知,这些机制是人类物种之外所不知道的。我或许不该把“冲突”“分裂”“压制”三个词当作近乎同义词来用,或至少当作指称同一实在之不同侧面的词。关于“冲突”一词对我所取的意义,我在“父母——或冲突的心脏”一注(n\textsuperscript{o}~128)里略有说明。}。它向我显为\textbf{那个}大奥秘:关乎\textbf{我们这个物种}的特殊意义、特殊命运的奥秘。民族学家与心理学家所给出的那些“解释”\dmark{}——至少我所听过的那些——分明不过是些\textbf{合理化}\origpage{602}:为那被承受、被内化的压制辩护,说它是社会良好运转乃至社会存在本身所不可少;一如在一群残臂者或独腿者的社会里,必不缺几位显赫的理论家,摆出一套看似无可辩驳的论证(无人想到反驳),证明:一个人们可以使用双臂(或双腿)的社会,无论如何也运转不了\fnote{146}{(*)}{一如奴隶制社会的时代,对“最优秀的头脑”(他们同样役使奴隶)一如对其他人,没有奴隶便没有社会,乃天经地义。据说,柏拉图非得“有幸”意外沦为奴隶,才开始换一副眼光看事情。}。这里说的是一些一望即穿的辩护,它们竭力以“科学”面孔的解释,把一个奥秘变没了。实际上,人类社会之内冲突(或压制)之起源与意义的问题,仍是纯粹的修辞——只要装模作样提出这问题的人,不曾经历一场强烈而深入的、对\textbf{他自己身上}冲突、对\textbf{他身上}冲突之起源的认识工作。没有这样一份自我认识,这问题(一如关于自由之本性、爱之本性、创造性之本性那些问题)便是中世纪那著名的“天使的性别”问题的现代翻版——不过是一场文体练习,好把反正得塞进去的东西塞进去。这个问题,严格说来,并不是一个“科学”问题——不是一个其检视不以一份\textbf{成熟}为前提、而只以某份预备的知识、某个智性力量或灵敏的程度为前提的问题\fnote{147}{(**)}{(12月3日)冲突之意义的问题不属于科学的管辖,这或许使人期待:人们能在神话与宗教里找到答案的要素。可依我看,全然不是如此。就我所知,神话与宗教的一个要害功能——若不说其主要功能——便是设立一份“法”:那法,就其实质而言,是一包禁令,压制便在一个个特殊的社会里凭它而落实。这法被宣称为神圣本质,便无须自我辩护,无须解释自己的“意义”,更无须解释它与其他社会的“法”所共有的意义。}。

在这个个案里,我要做的不是猜谜式地设想:压制在人类社会里是凭哪些机制设立的\dmark{}——也就是说,不是要为压制这个事实找一份\textbf{解释}\origpage{603}。即便人们达成了一个说得通的、乃至有说服力的脚本,我也不会因此觉得自己前进多少。它或许会照亮奥秘的某个有趣的侧面——“机械的”侧面——却穿透不了它。正如古生物学与分子生物学的周详成果、甚至达尔文那些深刻的思想,也真正穿透不了生命之出现、及其在三四十亿年里在地球上创造性绽放的那个奥秘。在冲突的奥秘里,使我感兴趣的,不是机械的、科学的侧面——一个\textbf{外在于我这个人}的侧面,一如那著名的“费马大定理”——而是冲突之\textbf{意义}的问题。这份意义\textbf{与我直接相关}、本质地相关,一如它与无数代以来相互撕扯、相互屠戮的无数男人与女人每个人直接相关,他们把从父母那里接来的冲突,又传给了自己的孩子。

冲突必定有一份\textbf{意义}、而我能够多少认识这份意义——这肯定属于我方才谈到的那份“信”。对我是显而易见的事——而那份熟得不能再熟的“奥秘之感”:此处有某种深的东西可探测——同时也告诉我:这“某种东西”恰恰\textbf{便是这意义}。所说的那份“信”,与我对自身官能的信彼此重合:当我的官能在此毫无一丝怀疑地向我启示,我面前有一个待发现的“意义”时。

或许有一天,这份意义将变得显然,仿佛我一直认识它!这个奥秘在我看来丝毫不遥远、不可及。它以切近之物的样子呈现给我:更贴身地认识它,只在我一念之间。而且肯定地,我此刻已望见一条接近它的路,或毋宁说一个侧面——它已经在向我友善地招手。毕竟,冲突有许多要教给我的东西,而且它已经教了我许多……

\subsubsection*{(h) 颠倒(2)——或暧昧的反叛}
\tocsubsub{(h) 颠倒(2)——或暧昧的反叛}
\label{ch18:local:h42}

\notehead{132}{}
\phantomsection\label{ch18:local:n38}

(11月22日)这已是连续两条注,我眼看自己登程于十足计划外的远足——这一回\dmark{},我留神头一件事就从那\textbf{预定}的东西开始,破例一回\origpage{604}。我想检视前一条注里唤起(未加说明)的“类型情势”之一:易于引发对父亲之对抗的情势,而更深地,引发对自身阳刚特征的(或多或少彻底的)拒斥的情势——这份拒斥在对父亲的拒斥中找到它的象征表达。那份情势,早在11月18日的反思里——以“敌对的父亲(3)——或阳葬阳”那条注收尾的反思——我就记起了。我当时的心愿,是在这个“类型情势”里(至少),把手指按在\textbf{对男性之拒斥与对女性之拒斥之间的一层直接联系}上。

离我最近的、我又曾长期下过工夫的个案,是我母亲的个案。她一生都怡然于一份几乎不加掩饰的、对一切女性之物的轻蔑;她以过度的男性价值为自己塑形;而与此同时,自她的青春期起,与男人的关系,对她一直是一份“彻头彻尾”对抗的关系\fnote{148}{(*)}{与对女性之轻蔑相反,这份透过一份激烈而动荡的感情生活透出来的、彻头彻尾的对抗,终生无意识。我只是在我1979年8月至1980年3月的工作中才察觉。}。我有这份大运气:我母亲曾极自由地同我谈起她自童年起的生活;我又掌握着极详尽的、直到她与我父亲共同生活最初年头的自传笔记,还有卷帙浩繁的通信。这是——加上与她接触时我自己的亲身经历所归还给我的——一份异常丰富的材料;况且我远未穷尽它。但我对它下的工夫已足够,使我无可置疑地感到:她身上那份我刚唤起的双重拒斥——对女性的拒斥、对男人的对抗——根在一份与父亲撕裂的关系里。她父亲,在许多方面是个讨人喜欢的人,慷慨、正直、重情,却在战后德国(我说的是1914--18那一战之后)一场漫长的社会坠落里变得刻毒——那样的坠落,当时何其多。老实说,这坠落开始得更早:从一位高车驷马的富人的地位,一路跌到走街串巷的擦鞋匠。在忧愁与失意的刺棒之下,他那易怒的脾性有时变成家庭暴政,而他那位体弱多病的妻子,首当其冲\dmark{}\origpage{605}。我母亲,对父亲一如对母亲依恋至深,被这些父权的暴政场面所激怒——那些场面,她母亲默默承受,有时实在撑不住了,却从不诉苦。这孩子激情地与母亲认同——那父权专横的牺牲品;而与此同时,她母亲所扮演的角色(牺牲品的角色、被动的角色——“女人的角色”……)在她看来难以容忍。一边是与母亲的认同,表达为一场反叛、一份对父亲彻头彻尾的对抗;\textbf{与此同时},又有那“我永远不做她那样的人(那个不反抗而忍受的人)”的跃起——这跃起必然同时意味着:“我永远不做女人们那样的人。”

但更深地,还有对那份权力的嫉妒——父亲的、男人的权力,它使他得以随心所欲地支配。我母亲的一生,被这份吞噬一切的支配欲所支配、所摧残;而首先,是支配并摧毁\textbf{男人}的欲望——恰恰是那在她身上唤起如此狂怒反叛之跃起的人、那凭其天性本该支配她的人——一如她父亲曾支配过她母亲:她默默承受着,脸色苍白,全然无力。

我本要在这里写:这份反思,如今“接上”了11月12日那条注“激烈泼辣的妻子(阴阳的颠倒)”(126)里进行的反思。由于对那条注的记忆已不很清晰,我刚刚重读了它。奇怪的是,我忘了那条注(一如今天这条)也是由我母亲的“个案”引发的。十天前,我对稍许展开这个个案还心存迟疑。今天我卷土重来、越过了这份迟疑(它后来我也忘了!);这无疑是因为:在被检视的情势里,有一个侧面当时还搅浑着。我也忘了:今天这条注的出发点——“把手指按在……对男性之拒斥与对女性之拒斥之间的一层直接联系上的心愿”——本就是十天前那场反思最初的动因:它自然接续前夜那条注“超级妈妈,还是超级爸爸?”(125)结尾处的那份追问。实际上,12日那场反思的最后一句:

“无须更多,便可看到那‘缺失的环节’显形……”,似乎是说,我当时自信当日之责(建立\dmark{}这样一层联系)已经完成。倘若我全然忘了自己早已揭示过\origpage{606}这层联系、甚至早在四天前那条注(今天的反思正接续它)之前就已提出过这个问题——那无疑是因为:我对刚才引用的那份辉煌结论,尚未全然信服——它写成于“敌对的父亲(3)——或阳葬阳”那条注之前不过六天。引全那句话,情形便更清楚:

“无须更多,便可看到那‘缺失的环节’显形:它连接着对‘超级父亲’的对抗(它在其下葬中找到象征表达)与对‘女性’的轻蔑、拒斥,而更深地,与对自己身上‘女人’的否认(它或许将在一位‘超级母亲’的‘葬礼’的象征中找到表达,在一大堆滥情的、一语双关的颂词之下……)。”

这份结论里缺了一步,使它显得匆忙:那便是“对超级父亲的对抗”与对“男性”的拒斥之间的联系——这层联系,直到11月18日那条引注“敌对的父亲(3)——或阳葬阳”里才在反思中显现。对父亲的对抗,那时向我显为一份远更要害的实在的象征表达:对自己身上阳一侧、“男性”一侧的拒斥。在对女性之拒斥的“对称”个案里,象征表达与其深层意义之间的联系,早在“超级母亲之联”显形时——11月10日那条注“阴的葬礼(阳葬阴(4))”(124)——便已被知觉。于是,11日那条注“超级妈妈,还是超级爸爸?”里显出的两“联”相对——父亲的下葬与母亲的下葬——前天才被看作自拒(或对己之轻蔑)的对称的显现,取\textbf{对自身男性之拒斥与对自身女性之拒斥}的双重面目。

况且,在18日那条注“敌对的父亲(3)——或阳葬阳”里,我把自己限于“主语”为\textbf{男人}的个案——然而,我所知道的最极端的个案,却是我母亲的!她在那份反思里被整个忘掉了,甚至已被忘了十天(若不算在11月17日那条注里,她藏在“我父母”这个词底下)。

正是我对我的孩子们、及其与我的关系的认识,四天前使我感到对父亲的对抗与自己身上对男性之拒斥之间的一层联系\dmark{}\origpage{607}。老实说,对我的五个孩子中我有机会相当贴身认识的那四个,近年来我不止一次感到:在他们对我——他们的父亲——积习已深的对抗姿态背后,有一份对他们自身阳刚一侧的拒斥,而尤其,是对他们身上那股冲劲的拒斥——它把他们掷向世界,却使他们像一个被否认的父亲!我从不曾问过自己这是不是一条普遍的事实;或者毋宁说,我身上有一种未说出口的预设:事情想必如此——但在四天前的反思之前,我从不曾感到需要把这件事对自己清晰表述,更不曾费心稍加检视。老实说,这类“普遍的”问题,根本不是我在禅观里向自己提出的那一类:禅观的用意要切近得多:理解我自己——首先经由我与他人的关系;并由此也多少理解“他人”——即那些与我发生关系的人。

当然,四天前的反思里,当我暗示应当有这层联系——对父亲的对抗是一份更深的冲突的表达,即对自身“男人”的拒斥——这还只是一个预设,由我极其有限的经验所暗示。这层联系在我看来至少说得通,尤其在男人身上;但我不自许“看见”这层联系在一般情形下成立。关于它,我没有那份我如此经常选作可靠向导的“贴身信念”。比如在我母亲的个案里,我看得分明:对父亲的对抗,是一份对\textbf{男人身上}阳刚特征的隐秘而恶毒的对抗之源,却全然不是对女人身上这类特征的——恰恰相反。诚然,单单充分推崇阳刚特征、在自己身上过度培植它们,或许并不必然意味着:人们全然接纳自己存在的阳一侧;那毕竟还意味着:也接纳那份“阳中之阴”——它自发地居于每一个以阳为“主导”的特征之中;而这当然不是我母亲的情形。

但反思正在这里取一种略嫌辩证的转向,它不叫我放心!我宁肯诉诸我对母亲这个人的直接知觉——它已由我对她、以及对我父亲生活的反思所磨细。我不记得曾有过这样的感觉:她身上有什么根本属“阳刚”的东西被拒斥\origpage{608}。相反,我在她身上强烈地知觉到这份矛盾、或毋宁说这份\textbf{撕裂}:她在自己身上培植(如一件件\textbf{武器})、并爱逾生命的,恰恰是那些在男人身上在她身上唤起如此烈性、如此狂暴的求战与求毁之饥渴的特征;而她的一生,正是被这份热病所磨蚀(并过早燃尽):不停地遭遇与迎战、并不遗余力地在他人身上压服那\textbf{同一份}力量——她押上全部赌注的力量,它摧残着她自己的生活,一如它摧残着所有她所疼爱之人的生活。

\subsection{主人们与仆人}
\label{ch18:local:h43}

\subsubsection*{(a) 天鹅绒的脚掌——或微笑}
\tocsubsub{(a) 天鹅绒的脚掌——或微笑}
\label{ch18:local:h44}

\notehead{133}{}
\phantomsection\label{ch18:local:n39}

(11月24日)前天那条注的反思里唤起的个案,并不是我所知的、印证这份预感的仅有的个案:父亲的超级阳失衡(无论这失衡取不取专制的形态),会在孩子们身上回响为一份对阳的拒斥,而这份拒斥又可以千姿百态的面目表达。男孩身上,就我所知、此刻在我想象中的个案而言,这份拒斥取对自身阳刚一侧(或多或少彻底的)压制之形——而这份拒斥,肯定将跟随他一生(除非一场深刻的更新,那当然极其罕见)。我母亲的个案使我确认:在女儿身上,情形不总是如此——除非我母亲身上也有过一份对自身阳刚一侧的拒斥,以更精微的方式表达,而至今逃过了我\fnote{149}{(*)}{一份邻近的情势,是一位禀性专断、侵夺性的母亲——超级阳失衡的征象。在我贴身知道的两个个案里,这在女儿身上表现为对她身上“阳刚”特征极深入的压制。}。而在她的个案里,相反,触目的是那极端相反的效应——她身上阳刚特征的过度发育(外加对一切女性之物的厌恶)。况且我还在别的方向知道一些个案,在\textbf{男人}身上(比如我母亲的父亲)——那是一场对父亲的\textbf{反叛},表达为一份强阳人格的发育,足以“势均力敌”地迎战父亲。既然我不曾有机会贴身认识这样的个案,我倾向于相信它应当更罕见。但这到底无关紧要。

\dmark{}倘若在或近或远我所知道的全部个案里有一个共同点,那便是:父亲的超级阳失衡,在孩子们身上回响为一份\textbf{失衡}——它可以朝阴的方向(或许最常见的情形),也可以朝阳的方向\fnote{150}{(*)}{当我在这里谈“朝阴方向的失衡”时,这并不意味着他的阴特征的(或许过度的、单侧的)发育,而毋宁是对阳特征的一份\textbf{压制}——这根本不是同一回事。在相反的、称“朝阳方向的失衡”的个案里,说的确是阳特征的“过度发育”,它常常伴着对某些阴特征或多或少深入的压制。}\origpage{609}。在我脑海在场的所有个案里(我并不打算在此对我所知道的个案作一份系统的清点),这份失衡伴着一份\textbf{对父亲的对抗关系}。我有这样一种印象:它同样伴着一份对另一些人的彻头彻尾对抗的姿态——那些阳特征打有强烈印记的\textbf{男人}——至少当那些特征不曾被互补的阴特征所平衡之时;也就是说,对着这样一些男人:在他们身上占上风的,是一种使人想起父亲的超级阳失衡。

这样一份超级阳的失衡(一如相反的失衡),固然易于在任何人身上唤起一份\textbf{不适},如我已有机会确认的\fnote{151}{(**)}{见“超级父亲(阳葬阴(2))”一注(n\textsuperscript{o}~108)。}。但这份不适并不必然表现为一份自动的对抗姿态——比如,它并不罕见地经由一份臣服、一份或多或少无条件的钦慕、或效忠的姿态而消解(或至少从意识的场域里消失)。

这里,一场联想来到心头:在数学世界之内、对我这个人(有着光环的)的关系里,最常见的,肯定正是这些调门——至少在那些同事(或学生)之中:他们(如我在别处所写)“并不觉得自己被一份相当的声名所保护”,或者(我要在此补上一句)那些人:某种内在的平衡、对自己力量的一份自发的认识,也未能使他们免于这类失衡。但无疑,这样一种“效忠”关系的本性,是它暗藏一份对抗;这份对抗,一俟合适的机会出现,便显形(公开地,或仍以隐蔽的方式)\dmark{}\origpage{610}……

我刚刚追随了几场联想,它们重拾并补足前天(在上一条注“阴阳的颠倒(2)——或反叛”里)的反思,并由此,也补足11月18日那条注“敌对的父亲(3)——或阳葬阳”的反思。它们使我明白:一位父母身上的某种阴阳失衡状态(在此即父亲的阳失衡),与它在孩子身上所取的那些回响之间的关系,绝无一义性——我先前说得匆忙。毫无疑问,父母的失衡(在此即父亲的)传递给孩子时所取的形式,必取决于许多别的因素:既有家庭环境的(尤其母亲这个人与她的姿态),也有孩子天生的禀性\fnote{152}{(*)}{这样,我确认:我母亲的三个弟弟(都比她小)各自走着一段与她(她多少像鸭群里的天鹅)全然不同的演进,也不同于其他兄弟。}。

但老实说,方才起笔反思时,我想走的方向并不是这一边。毋宁说,我想继续一场全然不同的联想——它自11月12日的反思起便在场,或头一次把阴阳角色的\textbf{颠倒}之动力学引入反思(在同名那条注里,“——或激烈泼辣的妻子”(126))。或许读者已在他那边做了这份联想——反正,当我在11月12日、随后前天22日唤起这个问题时,在我脑海的某处,确确实实像弱音器一样,有着另外两次“颠倒”已被谈起的记忆:在这场葬礼反思的进程中。第一次,是在V行列那条同名注里,“我的朋友皮埃尔”(4月28日的注(68$'$))。第二次,在9月30日那场反思里的一条脚注中,它属于“葬礼悼词(2)——或力量与光环”那条注。甚至还有第三次,但藏在字里行间:在两天后那场反思的开头,它开启了“阴阳之钥”的反思。(那是10月30日的注“肌肉与肚肠(阳葬阴(1))”(106)。)\dmark{}这里涉及的,是那场著名的“由三部分葬礼悼词引发的联想”的内容\origpage{611}——人们暗示过它;恰恰是它,当天便触发我踏上这条关于阴阳的枝蔓,我已沿着它走了将近两个月。或许,现在是时候揭开底牌了——自我谈起它以来;更不必说,自5月12日之后——“葬礼悼词(1)——或恭维”那条注之后——第二天起,我便想着它,已有六个多月。

这三桩情势的共同点,是它们涉及我的朋友兼昔日学生皮埃尔与我之间角色的一场“颠倒”。在那两桩被明白表述的情形里——方才重提过——我显为我昔日学生的“合作者”(甚至干脆是他的学生!)。第一次,是作为那个曾为“强有力的工具”之发展出过力的人——那$\ell$-adic上同调的“强有力的工具”,归在我那位才华横溢的前辈兼朋友名下(出力的方式当然乱糟糟,但人们也承认,有时有点意思)。第二次,当我们被相提并论时(因“以‘跨学科’手段联结了拓扑、代数几何与数论……”),同一份实在的颠倒,凭一个狡猾的手段被暗示出来:一场“排字”的遗忘——仿佛出于天大的巧合\fnote{153}{}{一如我在“屠杀”一注(n\textsuperscript{o}~87)里先前得以确认的:只要排字工人和搬家工人掺和进来,巧合常常把事情办得极好!}。这份颠倒的意味,当人们留意那些颂扬形容语的选择时,便变得比一个单纯的先后问题更别有用心(在此,是在一个我曾是唯一与Dieudonné一道、在科学层面上“启动”了它、却早已离开的机构之内):一边是“传说般深邃的理论”,另一边是“耀眼的发现”——后者还配得上着重号,人人都配,唯独我不配。这份意味,在“阴的葬礼(阳葬阴(4)”((124),11月10日)那条注的反思里,“触目惊心地”亮了出来——阴阳的反思凭它忽然“降落”在葬礼仪式的正中央:一边是(时而滥情到极点的)阴的、超级阴的形容语的堆积,另一边是阳的、超级阳的……

这正是早在5月12日“恭维”那条注的第二天,就已经打动过我的东西\dmark{}——早在两周前得以那样周详地显豁它之前\origpage{612}。按照我当时感受事情的方式(我日后还须在此重新检视),那里有过一场实在的真正\textbf{颠倒},或更确切地说,一场“颠倒”——被推向漫画式的极端——而其底下的实在,我感到是有层次的、平衡的。我把自己看作一个“阳”占强势主导的人,甚至超级阳——至少在我那些最显眼、最明白的特征上,尤其那些对他人彰显的特征上\fnote{154}{(*)}{而且,在我“出发”之前的那些年月,更甚于如今。}。相反,我在朋友皮埃尔身上感到一种基调属阴的气质——分明比我平衡得多,在我们常见面、他还是学生身份的那些年月,比当时的我要平衡。

况且我相信,这份对实在的把握在本质上是正确的。这些年里,有时——乃至最近\fnote{155}{(**)}{见“箭与浪”那条注(n\textsuperscript{o}~130,11月19日)。}——我偶然预感到我身上有一份原初的“阴”底色;而依我看,我是头一个也是唯一一个感到它的人——他人始终是经由我那些“阳刚”的特征(它们往往已经相当具有侵夺性)来把握我的\fnote{156}{(***)}{对我自己,亦然。},无论在意识的层面还是无意识的层面——至少就个人关系而言。而个人关系(恋情除外)所牵动的,主要乃至排他地,是我们身上的“东家”,那受制约的东西。阴阳反思进程中新显出的事实——\textbf{在我的工作里},我把握事物的方式阴占强势主导、“女性”——并不真正与我另外知道的东西相矛盾。它使后者更有层次,修正了其中一处我曾默默“一锅烩”的地方。细细权衡,我觉得:我心中那份骤然而强烈的印象——一场实在的漫画式“颠倒”,或更确切地说,一场蓄意\textbf{颠倒}的\textbf{意图}——这份“直觉”在本质上也是正确的,尽管粗糙。我现在想细细开掘的,正是这份直觉所未曾完全把握的实在。

\subsubsection*{(b) 兄弟与配偶——或双重签名}
\tocsubsub{(b) 兄弟与配偶——或双重签名}
\label{ch18:local:h45}

\notehead{134}{}
\phantomsection\label{ch18:local:n40}

\dmark{}(11月25日)我首先该试着更贴切地圈定这份印象——它于我显而易见:朋友皮埃尔这个人身上的“底色”,是一份\textbf{阴}的底色\origpage{613}。依我所知觉的,无论在“我”的层面上——我尤其在他与我的关系、与其他人的关系里看到它的表达——还是在他的工作里,也就是说,在他身上认识的冲动、创造官能的层面上,都是如此。

就第一个方面而言,分明地,他与我是气质\textbf{互补}的,还加一层细节:我身上那份过度的、“超级阳”的东西,有时似乎让他有些错愕。我相信,尤其是我那份朝着完成任务之实现的恒常的前投、那份对一切与此无关者的\textbf{隔绝},在他身上唤起一种难以置信的惊异——我在其中感到一缕温情的惋惜:那份惋惜,我曾在我母亲身上多次感到,当她看到我与周围事物的美如此隔绝之时\fnote{157}{(*)}{我母亲与我父亲一样,直到生命尽头都保持着一份与自然交融的能力,同时对周围的一切有敏锐的观察感——这两样,至今都是我所欠缺的。这或许是她存在中唯一未被她压制、得以自由绽放的“阴”的方面。另一方面,就“朝向目标的投射”而言——它是我“我”的主导特征之一——这或许也是我这个人唯一成功的方面:比我的母亲更阳!}。

在他身上,这不是本来意义的“不适”,不是拒斥某份实在的征兆。至少,我不记得有一次在他身上感到过对我的不适,也不记得有过拒斥、疏远、或哪怕我们之间的顶撞的印象。而我毫无怀疑:这绝不是他身上什么“外交式”的故意——决意什么都不表露。相反,他有时会表达我暗示的那份“惊异”,不带一丝局促,也不带一丝恼怒。分明地,我们关系里的基调——它至今\dmark{}从不曾变过\fnote{158}{(**)}{(11月26日)基调固然还是同情、吸引,但这并不妨碍:自我出发之后,年复一年,这份关系越来越凝定、僵化,失去了赋予它生命质量的东西。我感到自己面对一副“甲壳”,密封得如此完好,再没有什么能从任何一个方向通过。关于这一点,见“两个转折”与“墓”两注(n\textsuperscript{o}~66、71)。}——是一份温情的同情,没有一丝阴影从中穿过。

\origpage{614}这在我始终是一件奇怪的事,而且我相信,在我离开IHES的情节之前,任何人都不可能从中猜到什么(甚至那时,也只是在面对面的接触里“透”出来的层面)——即:自我与他相识的最初几年起,他与我的关系里就有一份深刻的、根本的含混:一份暗藏的对抗、至少一份想与我这人划清界限的愿望,以及排挤的愿望。这后一份,在我离开IHES的情节里(见“驱逐”一节(63)),以格外粗暴的方式显形(当场叫我目瞪口呆),尽管方式上无比地压低了嗓门。我的朋友,不久前刚被遴选为IHES第五位“终身教授”——主要凭我热心的奔走。在我们之间进行的那次“解释”(也许有过不止一次,我已说不清)里,他时时刻刻都保持着那份完美的、微笑的自然——带着善意亲切的一切表象,那是他如此讨人喜欢的地方。他向我解释,不带一丝犹豫或为难——更不带对抗或敌意、或暗藏的得意——说他早在青年时代就决定把一生与全部精力献给数学工作;说他这份对数学的奉献,无论好坏,对他而言必须置于一切之上;说我之所以等待同事们、尤其他本人的声援(以要求取消来自军部的经费),其缘由在他看来与数学全然无关;他当然惋惜:这对我是一票否决的情势,由于“公设”与他不同的生活观,我将离开IHES,为一个在他看来无关痛痒的缘由;但令他深感遗憾的是,他不能声援这份对他陌生的请求——其结局他毫不关心(注~134$_{1}$)。

我上面给出的,是我朋友那番话的“表层”内容、明白说出的内容,一如我的记忆所归还给我的,而我并不曾费心同时去复原一种表述的风格、或一场谈话的氛围——它的细节,除我此处所说的,我也不曾记住分毫。这段情节发生之时,我还没有丝毫疑心:在一番话那十分平淡(有时怪诞地荒谬)的表层内容背后,往往以弱音\dmark{}、而清清楚楚地,表达着一份全然不同的讯息\origpage{615}。这份讯息,肯定在无意识的层面被知觉着,却被惊惶地拒斥、逐出意识的场域。一如我在引注的“驱逐”那条注里所暗示的:要把一份其实足够刺眼的讯息清场出去,想必需要相当大的能量!而恰恰是在那条注里——写在十四年之后——我才头一次费心把这段情节交给有意识的注意,并把那份被否认了如此之久的意义清晰地表述出来。

我刚刚追随的,是呈现在我面前的联想线索中最强的一条。我追随它时,怀有某种迟疑——仿佛经由这场“枝蔓”,我正离开我的主要打算。可事后我明白:全然不是如此。无疑,一个人的影像与气质,当它自发地从他被牵涉其中的具体情势的描述里透出时,比一份号称能圈定它们的“特征”罗列,更活、更有说服力。但与其投入其中,我宁肯记下又一场联想,并踏上又一场枝蔓:把这里检视的关系,同Serre与我的关系相比。在我们两个人的关系层面,占上风的印象,绝不是与皮埃尔在一起的那份“\textbf{互补}”,而毋宁是两个气质之间的\textbf{亲缘}——双方都强“阳”。在十八年紧密的数学交往里,这份亲缘不止一次经由偶发的摩擦显形,表现为一时的冷淡,没有一次持续长久。依我的记忆,这些情节,缘于Serre那些不拘小节的焦躁举动,它们与我的敏感不“合拍”。Serre有时会恼火于我执着一个想法的顽固——当它在我看来重要时,我顶着逆风硬干;我一有机会就把它重新端出来,不管它“通不通过”,我深信(这份深信很少骗过我)自己握有“\textbf{对的}”观点。不知为什么,Serre对我的那些上同调“大家伙”发展出一种反感——也许他干脆过敏,一如André~Weil对一切“大家伙”过敏。另一方面,当我开始展开“我的”上同调yoga——在五十年代后半期——Serre实际上是我唯一不时的对话者;所以一开始就很难顺风顺水!我相信,他是在$\ell$-adic上同调自1963年发展起来之后,才谨慎地开始对这份工作产生兴趣的;同年,又有我对L函数有理性的证明草样(“三下五除二”)\dmark{}\origpage{616}\fnote{159}{(*)}{我记得的另一处摩擦点——无疑更带偶发性——是我坚持把代数群与形式概形中的取商理论(五十年代还被人误解)系于平坦等价关系之“有效性”的问题,乃至(后来)系于fpqc层的取商语境。这些观点,先由加布里埃尔与曼宁重拾,如今已是代数几何里几乎处处通行的通货,甚至在别的领域亦然。我觉得,Serre的迟疑,自我终于费心(因为似乎没有别人肯干这苦活)把第一条有效性定理——对平坦而有限的等价关系——白纸黑字证明出来那一刻起,便消散了。},才开始明白它们通向某个地方。

依我看,Serre与我的关系,是一份典型的阳-阳亲缘,与Deligne的关系正相反,那是一份阴-阳的互补。而在数学工作、数学进路的风格层面,情形恰恰颠倒过来。一如我在前面一条注里(“九个月与五分钟”(123))已有机会说过的:我感到Serre的进法与我的进法是\textbf{互补}的,取阳-阴互补之义。而恰恰是这份互补,成了那些偶发摩擦的缘由——摩擦缘于他与我双方都强阳的气质。

Deligne与我,在数学进路上的关系,则全然不同,这毫无疑问。我可以毫无保留地说:正是在Deligne身上,胜过任何别人,我体验过那份完美的\textbf{亲缘}——在我们彼此都感兴趣的数学问题之看法与进法上。每当我们的数学对话发生,这份体验便重新来一遍。我清清楚楚:这绝非一种偶然的情势——比如说,不是由于在那些关键的学习岁月里,我确确实实给过他影响。这份亲缘并非在长年的熟悉中发育——恰恰相反:它从我们最初接触起便在场,正是它,才是那份力量,几乎一夜之间造出那样一份系于我们共同激情的纽带\origpage{617}。这是两种数学进法之间的一份深刻亲缘\dmark{}——先于我们的相遇而在,而它(我深信)表达着双方原初气质的一个重要侧面:把握事物、发现事物时那份“阴”的“基调”\fnote{160}{(*)}{(11月26日)本注的反思,与“上涨的海”“九个月与五分钟”两注(n\textsuperscript{o}~122、123)的反思相续,似乎意味着:对任何人来说,都有一份“双重签名”、一份\textbf{双重}“基调”在:其一(无疑最显眼)关乎“东家”,即“我”的结构与支配它的机制;另一份关乎“匠人”,别名“孩子”,也就是说,也关乎认识的冲动、发现世界的冲动、创造的冲动(当然,也包括爱的冲动)。(诚然,把东家当作匠人、把匠人当作东家——也就是说,把鱼目当珍珠——是世上最平常的事,不过那是另一段故事……)
如此,在我身上,这双重基调是阳(东家)--阴(孩子);在Serre身上,是阳-阳;在Deligne身上,是阴-阴(对此,我心中没有一丝怀疑或犹豫)。在双方彼此好感的底子上,正是这份“记号”(或“调门”)的“分配”,使得在人与人的关系层面,我与Serre的关系是亲缘、我与Deligne的关系是互补;而就我们对数学的进路而言,恰相反。
四种可能的“分配”里,只剩双调阴-阳还“记在账上”。鉴于阴在我们大男子气社会里的失宠——这失宠倾向于主要作用在第一个调门(“东家调门”)上——我推测,双调阴-阳应当比阳-阳少见。可我至少认识一位知名的数学家,依我看正对应这份签名。当然,第二个调门、“原初调门”更难圈定,因为它常常被外部影响、被“跟人人一样”地做与是的心思搅浑。}。

这样一份贴身的信念,谈不上“证明”——正如我不曾想过要“证明”我自己数学工作里(姑且这么说)的基调是阴的、“女性”的。对这类事情,至多有时可能把一份感受从一个人“传递”给另一个人,在别人身上触发对某样东西的认识——那样东西,人们此前不曾对它留意;它逃过了他有意识的注意,却其实早已以弥散的形式被“登记”在某个地方。情形肯定被搅浑了——一如常常如此——被当事人为按时兴价值、按“阳”的“男性”价值自我塑形所作的那些努力。而我分明看到:他的数学著作、以及他所施加的那份(可观的)影响,都深深地打有他与我的关系那份暧昧的印记;我却怀疑:那些为抹去一份与我的底色相近、却被否认的底色而作的努力——是否终于成功\origpage{618}。诚然,严谨的那些要求——在我“出发”之前,它们尚未在他身上起作用——早就使他不再屈身(至少在打算刊布的文字里)理会那些远低于他的东西、或那些今日被视为该绝的东西。可是我觉得,在他刊布的东西里,他终究没能不追随那份自发属于他的进法风格。至少,这是我翻阅他那几份为数不多的抽印本时得到的印象——在我“亡故”十五年之后,他还肯把这些越过坟墓寄给我。

当然,我对Deligne数学进法的把握,首先汲取自我“亡故”之前的那些年,1965至1969之间。整整五年,我们强有力地都接在同样的事情上,数学交往不间断(除了他在比利时的那一年),而且比我与任何其他数学家——包括(依我看)Serre——的交往都强烈。我已有机会不止一次唤起这些年头\fnote{161}{(*)}{尤其见“孩子”“下葬”“驱逐”“授职”“死结”诸注(V行列,“我的朋友皮埃尔”内),以及“继承人”一注(IX行列,“我的学生”内)。},那是双方都创造强烈的岁月。在我朋友身上,它们以一个令人印象深刻的启程为标记——它却并不叫我惊讶,因为一切在我看来如此理所当然!那时,他对实质的那份极可靠的感觉——对最抽象的表象背后、或在最“一般的胡话”的表述之中的可触之物——还不曾被一种自满、也不曾被那后来才出现的葬礼综合征所遮蔽。他那时为这些主题(我可以说,极阴的主题)作过许多贡献——后来的共识(经他毫无保留的祝福)早已把它们开除出“正经”数学的行列\fnote{162}{(**)}{(11月26日)我顺带提醒:这些数学的一部分,曾在1981年的“反常层会议”上被大张旗鼓地开棺出土——而我连名字都不曾被提起;次年,又借着那“值得纪念的LN~900卷”。关于这一点,见“不义——或归来的意义”“赊账的学位论文与全险”“一个梦的回忆——或Motif的诞生”三注(n\textsuperscript{o}~75、81、51)。译按:原书此脚注角标作“18.5.4.4**”,当系排印之误。}:topos的形式主义、上同调的“大家伙”……我在SGA~4的引言里,怀着明显的快意,把它们逐一清点、郑重表彰\dmark{}\origpage{619}\fnote{163}{(*)}{(11月26日)这些评注,是在SGA~4彻底重编的第二版里补入的(重编主要关乎site与topos的部分)。它们会给人这样一种印象:仿佛Deligne曾参与平展上同调与$\ell$-adic上同调那份“强有力的工具”之主要观念与主要成果的绽放。于是,我在那里为Deligne、为我其他那些上同调家学生的水磨添了水——十年之后,他们瓜分一位亡故师父的躯体!}。还有其他这样的贡献(其中有些还更“带劲”,足以让他一上来就跻身“大明星”之列),在我1968/69年的那双份报告里——见“授职”那条注\fnote{164}{(**)}{我提醒:这双份报告已重印于《数学反思》第一卷。}。

\notehead{134$_{1}$}{}
\phantomsection\label{ch18:local:n41}

(11月26日)\fnote{165}{(***)}{本子注(系于前一条注“兄弟与配偶——或双重签名”~n\textsuperscript{o}~134)由该注的一条脚注生出。(参见该注第三段末尾处的参见号。)}典型的细节:那笔军费——当它关乎要我离开时,谁也不肯为它抬一抬小指——就在我离开的那一年,被悄无声息地取消了!谁也说不准,没准儿它会叫某位对这个章节有点计较的贵宾心里不痛快……况且,那笔经费在IHES的资源里只占一小部分(如果我没记错,是5\%)。不曾有人串通,可我那四位IHES的同事(不算所长)之间,却有一种漂亮的全体一致:抓住一个机会,把我甩掉(几乎同时,顺带把所长本人也甩掉了)。而我,还以为自己不可少、还被人爱着呢!

(12月6日)IHES的两位物理学家——Michel与吕埃勒——不满意:IHES的“物理”分部,与由Thom、Deligne和我(其中有两位“菲尔兹奖”!)代表的数学分部相比,显得有点像穷亲戚。Deligne的遴选,又加剧了这个失衡(那次遴选,本就是在Michel与吕埃勒毫无保留的同意下进行的——IHES科学委员会事实上全体一致,唯Thom例外)。物理学家与数学家之间有过协调,要向所长Léon~Motchane施压,以求在可能的限度内恢复两个分部之间的正当平衡\dmark{}\origpage{620}。不过我猜,我的物理学家同事们,看到这失衡被有效地补偿——而且比他们所敢望的早得多,借着我突然要走的前景——心里未必不乐意。

至于Thom,他怒不可遏:Deligne的遴选,竟是顶着他的正式反对进行的。他曾把我那份辉煌的“授职”报告里列举的Deligne的各项贡献——全未发表,分明全从他头顶掠过——称作不过是一些“练习”!Deligne登上IHES“终身教授”之位、与他平起平坐,叫他震惊的是:年轻的Deligne——当时25岁——竟尚未浑身挂满荣誉。按Thom的看法,登上这样的职位,只应作为“一项事业的加冕”而来。不到十年之后,便离那些英雄的年月那样远——当年,我把尚不知名的Heisuke~Hironaka迎进临时凑合的屋子……总之,Thom的怨气如此之深,以致他当时(据他亲口对我所说)琢磨着离开IHES,回到他在斯特拉斯堡的教授职位——那职位他留了一手(比我当年离开CNRS进IHES时更谨慎)。我热心地为Deligne作保,便成了他受挫的第一位的、主要的缘由;我猜Thom在心底里必定觉得:我不过咎由自取——我竟敢如此无礼,在把我那位辉煌的“被保护人”引进IHES不过几个月之后,自己便不得不离开!

至于所长,在他被终身教授们一致的意愿逼到墙角、催他下台的时刻,他便(凭一套他耍得炉火纯青的惯用战术)玩起了“分而治之”的把戏:拿军费问题当作一个方便的手段,声东击西,同时把最碍他的那位终身教授甩掉。(局势的颠倒,手法老到!而他对这笔经费之存在所守的那份秘密,在我看来,反倒成了催他下台的一条追加的、强制性的理由!)这并不妨碍:我走之后,事情到底没再拖多久,他的离任紧跟着我的离任——我,与他一样,自IHES最初那些寒酸而英雄的年头起便在其中,而正是我,与他一道、各凭自己的手段,保证了它的信誉与长存\dmark{}\origpage{621}。

\subsubsection*{(c) 阴,作为仆人——与新主人}
\tocsubsub{(c) 阴,作为仆人——与新主人}
\label{ch18:local:h46}

\notehead{135}{}
\phantomsection\label{ch18:local:n42}

(11月26日)在我出发前的那些年月里,Deligne与我之间的许多亲缘之中,有一份是这样的快意:他与我一样,喜欢在需要的时候,展开我所谓的“大家伙”。我作为数学家的能量,若不说全部,至少绝大部分,都献给了这样的任务。倘若说是盖一座房子,“做大家伙”便意味着:不止于画一张诱人的房子草图、甚至从不同角度画上两三张,也不止于画出带尺寸的详图;而是把用来盖房子的石头,一块一块地搬来、凿好;砌成墙,架上梁、椽,铺上瓦或石板瓦;安上门窗、盥洗盆、水槽、管道与檐槽;并且(倘若真打算自己住)一直安到窗帘与墙上的画。房子可以轩敞,也可以是只一间的小屋——做工里的精神却是一样的。而只要住进去,任凭你做得多么彻底、多么到头,你很快便发现:工作永远做不完,总会有新东西冒出来——至少,当“大家伙”——对不起,是房子——够大的时候。

1955至1970年间,我作为数学家的最大部分的精力,献给了启动并彻底展开四项\textbf{大}“大家伙”——当然,没有一项做到了头,见前文。按时间先后,它们是:上同调工具、概形、topos、Motif\fnote{166}{(*)}{“上同调工具”不曾等我到来便已存在。这里说的是某一种个人的进法,它尤其引向“平展上同调的掌握”(在我看来,它是Deligne所完成的Weil猜想之证明里主要的技术与概念的要素)。二十年之后,我在《追寻栈》(A~la~Poursuite~des~Champs)里重新沿这条进路走下去,朝着“非交换上同调”(或“同伦上同调”)的方向。至于“交换上同调”的方向,我在“我的孤儿”一注(n\textsuperscript{o}~46)的开头处对这份进法给出若干说明。这里所说的四项“大家伙”,大体对应引注里的五把“关键观念”——出入只在于:“上同调工具”对应其中\textbf{两}把观念(即导出范畴,与“六运算”的形式主义)。
耐人寻味的是:四项“大家伙”(或主要研究主题)之中,唯一在我的葬礼悼词(见注~n\textsuperscript{o}~104、105)里被点名的,是topos。像特意安排的一样,它也是被我的那些上同调家学生下葬的三项之中、在葬礼悼词的时刻还不曾被换别人爹妈“开棺出土”的那一项。(葬礼悼词在1983年;导出范畴1981年于反常层会议开棺,Motif则1982年,在那“值得纪念的LN~900卷”里。)}。这四项主导主题,彼此亲密地系在一起\dmark{}\origpage{622}——如同同一座农庄、同一个村落里的几座建筑,共同成就一份蓝图。而每一项“大家伙”,都强制地——我毫不曾寻索——把我引向展开另一些明显小一号的“大家伙”:有点像盖一座大房子、乃至整个村落,人们便被引去砌一座石灰窑、建一间木构与细木工场,等等。比如,每年都会重新感到需要给范畴的观念与构造的武库,添上两三件(小号的)“大家伙”。那些晚了十年、二十年才来的人——他们发现一切都做熟了,便舒舒服服住进现成的房子(还有另一些人,心底里知道该有个分寸)——对着那么多读不通的“胡话”(Deligne语)、对着头发里分出的四股(“Spitzfindigkeiten”,一位对我满怀善意的著名德国通信者这样称呼),耸起居高临下的肩膀\fnote{167}{(*)}{我的通信者和颜悦色地向我担保——为博我高兴——他自然明白,我的著作“在很大的限度内,没有这类毛病”(“weitgehend~frei~von~diesen~Übeln”)。他所说的“毛病”,指的是一旦有人胆敢在尚属猜想的地基上发展一门理论(如我就Motif所建议的)就免不了要掉进去的那些毛病(如形形色色范畴论者的“Spitzfindigkeiten”)。这里我们又见到那份对“数学之梦”的彻头彻尾的拒斥——见“被禁的梦”一节及随后三节(第5至8节)。这又是数学里对一切“阴”的、“女性”的进法或行径之自动压制的侧面之一。}。这些人对在光地上盖一座房子是怎么回事,毫无概念;而他们,恐怕永远不会去盖,只满足于在别人凭双手、凭整颗心为他们盖好的房子里,扮演房东。

\dmark{}我方才有点急躁,看上去把朋友皮埃尔也塞进了那些“对盖房子毫无概念”的人的口袋\origpage{623}。他非但见过我做工,而且他自己也满心欢喜地在盖他的房子,仿佛他来到世上二十年,不曾做过别的事。况且,这“大家伙”与盖房子的故事(倘若读者还没看出来……),又是另一个侧面、另一幅意象,用以圈定我此前曾试着以“上涨的海”的意象、继而以一道道相互追逐的浪的意象,多少把握住的那个东西\fnote{168}{(*)}{见“上涨的海”“箭与浪”两注(n\textsuperscript{o}~122、130)。}。说的是把握实在的“阴的样式”、“女性”的样式,以及与之相应的行径:浸润其中,从中析出一幅意象,柔韧而忠实地复现这份实在。于是,经由我自己的这个人绕了个弯,我又回到了最初的打算:把我心中那份强烈的知觉“传递过去”——Deligne的数学进法与我的之间,有一份亲缘、一份本质的亲缘。可在Deligne身上这个我方才试着以意象圈定的方面,自1970年我的“出发-亡故”之后,依我看,有过一场全面的“搅浑”:我相信,在他“之后”的出版物里,“大家伙”全然缺席。当然,他也无法合理地动用他那被申斥的师父身上的这个特征,来贬损师父,却又容许同一个特征照他自己的天性在他身上绽放。

诚然,倘若要做的不是追随一份内在的需要——一份原初冲动的表达——而只是靠积累“立得住”的\textbf{成果}来抬举声名,那么我的朋友,实在没有任何兴趣继续让(或大或小的)“大家伙”来缠身。就在我的时代,在Bourbaki小组之外(小组本身正投身于一项相当规模的“大家伙”!),这已经是一件颇不受待见的事。这毫不奇怪:毕竟,我们这个社会、科学世界的共识里,那份“超级阳”的马眼罩,不是昨天才戴上的。或许这正是主要的缘由:我所乐意盖的那些房子,在漫长的岁月里无人居住,除了那位匠人-泥瓦工自己(他同时又是建筑师、木匠等等)\origpage{624}。而直到今天,即便我工作中早已成为公共遗产的那部分(甚至在没有别的参考文献、唯我的文字可依的地方),对某些人而言(至少对那些不属于那个上流圈子、也不以居高临下为本分的人而言)仍环绕着一圈近乎令人畏惧的光晕,仿佛踏进去便需要近乎超人的官能。这是真的:那常常很费时间,也不可能不如此——既然一切确确实实是手工的、细部的,从头至尾,甚至在每一章的转弯处都有说明,告诉你走到哪儿了\fnote{169}{(*)}{我想,只是随着岁月的流逝,我才意识到纳入这类说明——它们往往纯乎启发式——的必要:尽可能给读者传达一份“方向”感与用意感;写的时候,它在我身上强而在场。今天,依我看,这比一丝不苟地写出关键的证明要紧得多:一旦读者感到去向何方、而那“何方”吸引着他,那些关键证明,他自会津津有味地复原、甚至凭空造出。}。我的学生们当年与我一道工作的时节,进入这水里,并不显得格外吃力。但那是在“可触成果”已经挣得数学建制背书的时节,我的学生们怀着稳操一张“保险牌”的自信工作。我有这样一种印象:从那以后,倒不止一个人乐得给“读不通”的版本续命\fnote{170}{(**)}{这件事,只有Deligne一人说得有鼻子有眼——他最近的来访中,还当面重提过。说的是SGA~4(其一大半以极度的细致展开topos的语言),被我的朋友判为“读不通”,以此为他那天才的“SGA~4$\frac{1}{2}$工程”作辩护。},顺应一股今日比我那时暴虐得多的时风。

可是,即便撇开时风的种种要求不谈,当人们计算收益与“回报”时,肯定会像躲瘟疫一样躲开“大家伙”。展开一项“大家伙”、并把它交到所有人手边,这是向科学共同体提供的一项\textbf{服务}——而这个共同体,往往不情愿地接受它。对这份完全可以理解的推拒,我向来不怎么为难;我分明知道,我手里有“好家伙”,人们或早或晚,总免不了要来的\dmark{}\origpage{625}。而即便他们来了,以“信誉”计的“收益”,也只能微薄。倘若我作一份数字的盘点——不是盘点我在1955至1970这十五年里引入并发展的那些观念、问题、想法(它们要么进入公共的、匿名的遗产,要么无声无息地下了葬,等候被大张旗鼓地开棺),而是盘点人们可称作“大定理”的东西——我怀疑连十个也凑不出。直接花在证明它们上面的时间总量,大约几个星期,至多几个月。1957年之前,一个也没有(Riemann--Roch--Grothendieck定理)——然而我知道,此前三年我并没有虚度。说不定,倘若这十五年里,我没有执拗地追随我身上那份理解的激情、信任它向我口授的那份进法——不管它“有利”与否(按这样那样的要求来算)、在上流世界里体面与否——眼下这些“大定理”竟一个也不会得到证明(而这本非我主要的操心)。这份进法,每一次都在于:从一份强烈的出发直觉、或一小把这样的直觉出发,把它们当作一根结实的、经得住一切的线,由它拽着我走向未知;而一边走——换个意象说——我总按捺不住,随着那未知渐渐为人所知,如一些粗粝的石头,人们一边凿它一边“认识”它,盖起一座座房子,很大的与不那么很大的,全都宜于居住——一座座房子:其中每一个角落与旮旯,都注定要成为不止一个人的亲切而熟悉的居所。门窗端正,开合顺滑,不夹不吱嘎;屋顶不漏,烟囱抽风。未必是巴黎圣母院,每一只面包箱里也不都藏着一个“大定理”——不过是些该盖的房子,我为居住而盖的。我盖它们、盖得漂亮而宽敞,从中找到我的快乐;我分明知道:我做的工作——独自或与人同做——是必须做的,而且每一刻都做到我当时所能做到的最好。

这份精神,我在五十年代的Bourbaki小组里也找到了,它使我在其中如鱼得水、“在家”一样——尽管有氛围与文化的差异、尽管有我在该提的地方提过的那些偶发的困难。至少在那时,我在其中找到的,仍是一份\textbf{服务}的精神\dmark{}\origpage{626}:为一项\textbf{任务}服务;而在任务之外,也是为其他人服务——那些像我们一样渴望懂得大大小小的事物、并要懂得个透彻到底的人们。这种“服务”并不呈现为严苛的义务或苦行。它自发而欢快地出自一份内在的需要;它表达着一种共同的东西,把这些如此不同的人系在一起。

而我在Cartan研讨班里也认出这同一份精神——那么多法国数学家在那里初试锋芒;后来(六十年代)又在我自己的研讨班里(它的代号为SGA,“布瓦玛丽代数几何研讨班”)。两个研讨班的一个差别是:我的那些,强有力地以方才唤起的“大家伙”的展开为中心(因此是“\textbf{我的}大家伙”),它们从来嫌人手不够;而Cartan逐年跟随的主题,更为驳杂。可依我看更要紧的,是两个研讨班共同的东西,尤其,是它们似乎曾有的那个要害功能、它们的\textbf{存在的理由}。老实说,我看到两个。这些研讨班的功能之一,与Bourbaki的用意相近,是准备并向所有人提供易上手的文本(我是说,实质上完整的),周详地展开那些重要而难入门的主题\fnote{171}{(*)}{“难入门”,或因为这些主题仍被理解得不完全,或因为它们只有寥寥几个入室者知道,而讨论它们的零散出版物只给出一幅不称职的图景。}。另一项功能,是构成一个\textbf{场所}:在那里,胸怀热望的年轻研究者——即便不是天才——也确保能学到数学家的手艺,就着正当时的种种问题,在高明而善意的长者身边。学手艺——也就是说,把手插进面里,并由此找到让人认识自己的机会。

看上去,我1970年的出发,标志着“大研讨班”的终结——至少在法国:那些\textbf{持久的}场所——年复一年,当代数学的某些大主题在其中开工;也是\textbf{善意的}、激励人的场所,为一切前来把手插进面里的人\dmark{}\origpage{627}。我不知道世上别处是否还有(或许在莫斯科,凭I.M.~Gelfand的推动?)。确定的是:这样的场所,分明不合时宜,一如那些白纸黑字、细致入微地写下来、为的是交到\textbf{所有人}手边的“大家伙”。

并非偶然:如今几乎再没有人写那种精心结构的、(暂时)穷尽的阐述——论那些熟透了十年、二十年的主题,分明至关重要,而在此期间却只有一撮“局内人”能够问津。属于数学“上流世界”的人,倘若不同时属于那一“撮”,一旦需要,尽可以毫不费力地让那撮人里的一位给他讲个明白——人家求之不得。至于其他人:见鬼去吧!六十年代,我看到一大摞书,声嘶力竭地要求有人来写。我自己倒很愿意写,可我不能事事同时做。据我所知,这些书至今一本也没有写出来\fnote{172}{(*)}{(11月28日)这里我该为我推动之下写出的那些论文作例外。鼓舞我的那份精神——我相信,至少在他们与我一道工作的时期,它也传给了我的学生——正是鼓舞我自己的工作的那份;也就是说,打个比方,“盖那些分明需要盖的房子”——尽管某一幢特定的“房子”需要不需要,常常只有我一个人感到。我有这样一种印象:一般而论(一个例外除外),这份感觉终于传到了学生身上,使他“挂上”某个题目,继而强有力地认同所选的题目。把Verdier除外——他不屑于把我们约定好的基础工作交到所有人手边,那工作至今等着有人来写——凡与我一道完成国家博士学位的学生的学位论文,都成了所谓的“标准参考文献”。都是些宜于居住的房子,而且没有任何两幢重复……}。然而,不止一个人(且不说别人,单在我昔日的学生们中间)曾足够“在局中”、有那份感觉与那份手艺,本可以毫不费力写出某一本来就该写(而且至今仍该写)的书。而从某些人后来的工作所传回我的那一点点来看,我并没有这样的印象:是他们那些更个人的工作之繁多与困难,使他们……不肯(“真遗憾,实在没有工夫!”)向那著名的“数学共同体”提供这项服务。而且对不止一个人,大可以打赌:写一本被阅读、被引用的书(即便书中所阐述的并非都出自他自己——但“怎么写”绝不是可以忽略的量……),比他那一叠或厚或薄的抽印本,更会使他知名\dmark{}\origpage{628}。分明地,叫这个那个以一种令人印象深刻的全体一致,拒绝把那依旧只属于少数人的东西交给所有人的,绝不是单纯的“没有工夫”——或者,不肯采取一种“服务”的姿态(哪怕此地彼处,花点时间写一本书也好)。这里,一场联想不可抗拒地涌来:1965/66年的SGA~5研讨班——被埋没了十一年,而埋没它的,恰恰是它最初也是唯一的受益人:我的朋友皮埃尔领头,还有我其他那些上同调家的学生!诚然,那个案子里有一具躯体要瓜分,动机因此有点特殊。但我也想到别的个案:已提供的服务填补了明显的空白,却被在位者们一挥手扫到一边\fnote{173}{(*)}{我这里当然想到Yves~Ladegaillerie的工作,以及Olivier~Leroy的工作——前面四条注与节里谈到过它们(“谁也拦不住进步!”“棺材2——或锯成段的切割”“科学院报告——或新伦理”“棺材4——或没有鲜花没有花圈的topos”,注~n\textsuperscript{o}~50、94、第33节、注~n\textsuperscript{o}~96)。}。人们会说,这些仍是一路特殊的个案:被瞄准的是我这个人,而明摆着,那些工作正是由我所启发的。然而,我在这一切之中分明感到一种“时代精神”,它超出一切个案。

我此刻正多少圈定着的“时代精神”的这个侧面,便是\textbf{笼罩服务姿态的那份污名}——我经由一大堆汇聚的征兆知觉到它;对我而言,它是明摆着的事实。谁都可以否认它,正如谁都可以亲自检视它、确认它。我这里的用意,不是向一位不情愿的读者“证明”它,而是试着把握它的意义。在眼下这场反思的光照里,有一个意义一望即知。服务的姿态,典型地是一份“阴”的、“女性”的姿态;它落在被贬抑者的份里,毫不足怪。我屡次仿佛知觉到的那层细节是:这样一种姿态,只配给那些没有本事摆“主人”姿态的人;照这份精神做成的工作,是\textbf{下等}的苦役,是给那些驾着伟大观念与“耀眼发现”之高车驷马的人打下手的小卒干的\dmark{}\origpage{629}。

然而我也知道,不止于此——因为不然的话,人们为什么要想方设法,拦住一伙(倘若居然找得到)有心的小卒,不许他们在自己的角落里安安静静地做那份合乎他们身份的低贱苦役:终于供给一些结实的参考文献,好叫人不必再满足于(当人们还肯说点什么的时候……)“人们知道……”或“人们可以证明……”,或更罕见也更老实的“我们将承认……”?!

八年前,Yves~Ladegaillerie为“安放”他的论文而屡遭挫折时,我头一次直面这份扰人的问题\fnote{174}{(*)}{关于这一点,见前一条脚注所引的两注(n\textsuperscript{o}~50、94)。}。老实说,那正是我对数学、对数学家世界的兴趣最边缘的时刻。我有点目瞪口呆,却不曾试着去澄清这个谜的意义。此后若干年,直到今年2月、随着《收获与播种》里继续的那场反思,我的态度几乎没变——出入无几。然而,由于不断地接收征兆——甚至并非有意——我终究渐渐不由自主地也接收到它或多或少的意义,或者毋宁说,它的诸般意义。我确实看到两个。一个关乎我这个人——那是针对我的葬礼综合征,我对它的巡礼还没有全然做完。另一个与任何具体的人无关。那是一份\textbf{对“信息”之占有与控制的独占姿态}——在“科学”的建制之内占上风的姿态,它使这建制变成所谓“科学”“共同体”之内一个君权神授的统治种姓\dmark{}\fnote{175}{(**)}{(12月6日)人们会注意到:支配欲是一份\textbf{超级阳}的失衡,而且是这类失衡远为最常见的形式。它对应于“主人——仆人”这对阴阳偶子(或“支配(或掌握)——服务”这对偶子)里,阴的、“女性”那一项的泯灭;它邻近“掌握——服务”这对偶子。}\origpage{630}。

这是我在“职业道德共识——与信息的控制”那条注里(只擦了一擦,只擦了一擦)轻轻触及的主题,在“年轻人的势利——或纯洁的卫道士”((25)(27))里也谈了一点。我怀疑:这是科学世界里的一桩\textbf{新事实},在逝去的二三十年里,蹑着狼步悄悄安顿下来的。我不相信我曾属于传播、迎纳这份“新伦理”——“两套砝码,两套量法”的伦理——的人\fnote{176}{(*)}{我不知道,在我上一辈或同辈的同事之中、乃至在更年轻的同事朋友之中,有多少人看见过它。我怀疑,在“那些曾以兄弟之情迎接我、走进那个变成我的世界的世界”的人们之中——《收获与播种》正是献给他们的——除Chevalley或可算一员,连一个人也没有。这当然属于那些我本愿与他谈谈的事情之列——可他已经不在,不能告诉我了……}。倘若我在这份新伦理的问世之中也有一份共同的责任,那毋宁是因为:我不曾看见它的到来\fnote{177}{(**)}{我不知道,在我上一辈或同辈的同事之中、乃至在更年轻的同事朋友之中,有多少人看见过它。我怀疑,在“那些曾以兄弟之情迎接我、走进那个变成我的世界的世界”的人们之中——《收获与播种》正是献给他们的——除Chevalley或可算一员,连一个人也没有。这当然属于那些我本愿与他谈谈的事情之列——可他已经不在,不能告诉我了……}。在最近的这几年之前,我从不曾疑心:那份全方位的信息——自1948年我最初接触科学世界起,我便自由地享有它——已在不知不觉的岁月里(我说不清从何时、如何),变成一份惊人的\textbf{特权}:我与一小撮伙伴共享的特权——一份\textbf{阶层的特权},借一个被用得有点滥的词说;而这词在此,依我看,恰恰表达着一份十足可触的实在。

但我的用意,不是给数学世界作一份“阶级分析”、罗列它的“力量关系”与“权力的手段”,也不是作一幅“风俗画”。该回到一个更有限的事情上了:在几位主要的当事人身上,就其要害的发条,理解我那场提前举行的葬礼的“市井小事”!

\subsubsection*{(d) 阴,作为仆人(2)——与慷慨}
\tocsubsub{(d) 阴,作为仆人(2)——与慷慨}
\label{ch18:local:h47}

\notehead{136}{}
\phantomsection\label{ch18:local:n43}

(11月28日)前面两条注,基本上是绕着Deligne与我之间那份阴-阴亲缘的主题、在数学工作与数学进路的层面上——所作的枝蔓。我不知道它们是否有助于把我对这份亲缘及其本性的知觉“传递过去”——对我而言,它无可怀疑。

我在别处写过:“在我的工作里,我同样‘阴’,同样‘海与运行’,达到了可能的极致。”细细想过之后,我要说:这话按字面并不真——“可能的极致”还可以更甚,因为(依我所知觉的)Deligne比我更甚。或者至少,那份“阴中之阳”,在我看来,在我身上比在他身上更显豁。我身上的那份烈性,在他身上取更沉着的姿态。在我大胆向前扑的地方,他不止一次停在一份审慎的观望上——而那观望,往往是有根据的。只要我有一星想法的苗头、一个可以进去的“抓手”,我便毫不犹豫地扎进一潭我感到有实质的数学泥沼,不操心先把出发的想法凑近端详端详(德语所谓“ihr~auf~den~Zahn~fühlen”……),也不预料这场混战的结局。有时,那想法站不住,缘于某个一眼可见的缘由——而我因一头扎进“汁水”里而把它放过了。我终究会明白过来——有时我觉得自己蠢透了;可我很少后悔自己扑了出去。正是这样、也只有这样,我才与一份未知的实质建立联系:蹭进去,不管“得当”与否。

我的朋友呢,先探测、先检视——然后才扑出去:当他感到稳当时才扑,未必稳在对岸——那要求过分——但至少稳在有着陆之处、稳在不至于空手而回。我从不曾在他的工作里有过任何\textbf{能量耗散}的印象——那份耗散,我自己常有——而毋宁是:在他那里,\textbf{招招中的}\origpage{632}。从这个观点看,他工作的风格带着一份\textbf{成熟}的印记,而我的,宁可带着一份\textbf{青春}的印记——烈性所至,有时乱糟糟。可我们初次相遇时,已近不惑的是我,他二十岁。而且不止一次,我在他身上感到对我的一份微笑的宽容——有点像一位善意的成年人对一个他喜爱的孩子:当他看到我又一次登上一项(小小的)“大家伙”、什么都不怀疑地……
此处唤起的这些侧面,想必难以在“誊清的”、已刊的工作里察觉——那些工作呈现的是一份反思的最终阶段、或至少相当靠后的阶段。我在工作中的要求,不比他低;这些注在手稿打字员或印刷者手中之前,我几乎从不交出,除非它们已达到满足我心中完全清晰之需的阶段。相反,在《数学反思》里(尤其《追寻栈》里)我所追随的那种书写风格中,工作里原初的行径,每页可见。读者可以在其中确认大量“失手”。它们全都幅度很小——多半在次日或隔日、甚至当天就被发现,并在其后几页里得到纠正。(事情如此,连我自己也惊讶——这就是我在别处谈过的、我数学工作里那份出奇的“轻捷”的一个征象\fnote{178}{(*)}{见“陷阱——或省力与衰竭”一注(n\textsuperscript{o}~99)。我觉得,这份“轻捷”,如今甚至比当年——我“出发”之前——更大。这在我看来,系于一场在逝去的十五年里在我身上完成的成熟;它在我的数学工作里、一如在别处,让人感到。}。)“小失手”在场的缘由之一,当然是我对一个七八年不曾触及的题目不再熟悉;而随着工作的推进、随着失去的联系渐渐接上,这些失手也愈发罕见。但这并不妨碍:这种一次次把我那份相当混沌的记忆所归还的东西、我当年多少知道的东西,不加迟疑当作“现钱”收下的做派,正极好地例解那份“闯将”的、有时乱糟糟的侧面——它(与其他侧面一道)构成我数学工作(或非数学工作)里“阴中之阳”的侧面。我深信:一篇同样自发的文字,倘若出自Deligne的笔,会更接近人们通常所认为“可刊”的东西——甚至,接近按他那套严格的标准的“可刊”。

\dmark{}我之所以在这里着力强调我朋友的工作风格与数学进法里那份“成熟”的品格、“阴之又阴”的品格,绝不是要借此暗示:他的工作里有任何失衡,也就是说,这份工作会打有“阳”的、“阳刚”的品质之缺乏或缺席的印记。倘若如此,他的著作便不会像Serre的著作、像我的著作一样,每页都带着那份精微而不会骗人的\textbf{美}的印记。但此处不是地方——一如在Serre的个案里、在我自己的个案里,我都不曾这样做——去逐项追索那阴阳、“女性”与“男性”的精微和谐:在他那份我所认识的已刊著作里,以及在我与他将近二十年的个人接触所让我认识的那部分工作里。

也不该以为:我对阴阳之均衡所作的这份确认,是一种老生常谈,会立刻适用于任何一个以这种那种身份扮“大数学家”的人。我方才唤起的那份对美的知觉,并不是在每一个在自己时代的数学上留下持久印记的数学家的著作面前,都同样在场、都同样程度地在场。这些人之中,我认识两位,如Deligne一样,在我看来无论工作还是人格都以阴为主导;而他们的工作,却从未在某一时刻给过我那份内在平衡的印象、那份从不叫人对美失望的印象。这两位同事之一,那份阴的失衡取如此极端的形态,以致他似乎全然无力——哪怕清晰地、正确地表述任何一个定义、任何一个命题(更不用说一个想法了……)——尽管他对许多事情有深刻的直觉,并且引入过许多重要而多产的想法。它们每一次都靠别人的工作才获得形体。分明地,他身上有一份罕见有效的压制,压着“阳”品格的力量——在他的工作里、在他的存在方式里都如此。这份压制取一种真正的无能之比例,包括在他的工作里:凭自己的手段,他绝无能力把最小的一件事做到底。他用一份自大狂的姿态补偿这份存在的无能,同时内化那些他乐于在自己身上培植的毛病,仿佛正是\textbf{凭着它们}\dmark{},他才得以设想\origpage{634}那些(在他眼里)使他成为\textbf{本}千年大智者的想法……\fnote{179}{(*)}{我在这里谈的,是我出发之前那些年月里得以确认的姿态与存在方式——那时我有机会与这位显赫的同事非正式地见面。自那以后,某种东西变了,并非不可能(尽管那会是罕见之至的事……)。}。

我在朋友皮埃尔身上感到一份反方向的压制:清除某些“阴”的特征,把他(或多或少成功地)引向一份超级阳的影像自我塑形。这份压制,当然离我方才唤起的那极端相反的个案还远。它不曾达到那样的地步:把那份美的感觉、那份不带一丝不适余味的满足之感——真正领会的征兆,时时刻刻给明晰与幽影、给奥秘各留其分——从读者或对话者身上抹去。这就是说:我朋友所选的那份“超级阳”的招牌,想必不曾侵越他的工作本身——我是说工作的时刻;在工作的时刻,“东家”的在场,想必也如(我相信)在Serre身上、在我身上一样地淡\fnote{180}{(**)}{关于这份匆忙的印象,我将在本注的子注~n\textsuperscript{o}~136$_{1}$(12月4日)的结尾处回头再谈。}。

可是,恰恰在选\textbf{题}的层面上,依我看,“东家”的角色才变得重要、乃至具有侵夺性。有那份固定不变的念头:与我这个人划清界限;由此,便拒斥追随他自己天性的某些癖好——它们在他身上与那被否认的主人影像系得太紧。于是,倘若他偶尔——如任何一个手段高强的人——证明了困难的定理(乃至“困难尽人皆知的定理”),甚至引入了好想法并把它们展开,他却决计不会“重新思想”——以他自己的方式、哪怕只在大轮廓上——一整门科学(比如拓扑,它分明需要这个……);更不会凭空创造一门新的科学,“把新的世界拖到光天化日之下”(如我在别处所写)(136$_{1}$)。然而,倘若说有谁,我对他的这份能力毫无怀疑,那正是他。倘若有一样东西,至今于他有所缺,那便是\textbf{慷慨}——真正的慷慨:它同时是一份沉静的确信,使我们追随我们自身天性的冲劲,任它把我们带向哪里,不操心有无勉励、有无“回报”\dmark{}\origpage{635}。

但也有那份单纯的快乐:盖大大小小的房子,让别人去住——不必非得是“一整门科学”或“一个新世界”的规模;搬石头、架梁,像随便哪个泥瓦工、木工,不怕因此被当作这个、像极了那个;或把(在某些人眼里)本该是极少数人禁脔的东西,交到所有人手边。这是一种服务的姿态、一份谦逊——它仍是方才唤起的那份慷慨的表达、那份对自己天性的忠实的表达。我的朋友把它换成了另一份自满的姿态(“我——做这样的工作!”)与一份种姓的姿态\fnote{181}{(*)}{这份“阶层”的姿态——在我的朋友身上、在那个“数学上流世界”里——首先出现在我3月份的反思里:在“职业道德共识——与信息的控制”与“年轻人的势利——或纯洁的卫道士”两条注(n\textsuperscript{o}~25、27)里;上周,它又在“阴,作为仆人——与新主人”那条注(n\textsuperscript{o}~135)里重现。},在那些号称“可接受”的工作主题的拣选上。

最后,还有第三份姿态或力量,“东家”凭它压在我朋友对工作主题的拣选上、压在他给自己圈定的探测对象上——一份给他划定强制性边界的力量。那便是“师父的下葬”的综合征,或\textbf{掘墓人的综合征}。这里,不只是止步于不提名那个必须被无视的人。还要把他的著作本身下葬,或更确切地说,把它干净利落地“锯断”——用\textbf{电锯},在他自己的工作里一如在别人的工作里,在从一具强健的树干上迸出的每一根主枝的水平上\fnote{182}{(**)}{我头一次直面“电锯”的实在,是5月19日,在那双注“继承人们……”“……与电锯”(n\textsuperscript{o}~91、92)的反思进程里;随后是四条棺材注(它们与“掘墓人”一道,构成葬礼的“灵车”,即X行列),5月21、22日(注~n\textsuperscript{o}~93--96)。}。一如我前天(在上一条注“阴,作为仆人——与新主人”里)还提醒的:在我“几何学家”时期——1955至1970年——所析出并展开的四个大主题之中,只有一个被我那辉煌的学生兼继任者公开“取用”;其余三个都被“锯断”了——当然,是以弱音\dmark{}\origpage{636}。其中一个主题1981年得到一份极局部的开棺,另一个次年——像几株孱弱的幼芽,从被截断的主枝那结痂的残桩上重新发出,人们还为这场合给它围上斑驳的花环与刺眼的霓虹,好掩人耳目……

\subsubsection*{注~136$_{1}$}
\notehead{136$_{1}$}{}
\phantomsection\label{ch18:local:n44}

(12月4日)\fnote{183}{(*)}{本子注(系于上一条注“阴,作为仆人(2)——与慷慨”~n\textsuperscript{o}~136)由该注的一条脚注生出。(参见该注结尾前第三段中的参见号。)}我自己的行径,不断地把我引向“重新思想”——彻底地——我在数学家路上遇到的东西:无论它是看似最微末的事物,还是有着“一整门科学”规模的事物。诚然,我和大家一样只有两条胳膊,我不可能每一次都把一项“彻底重做一整门科学”的工作纲领,实现到我在代数几何的情形里所达到的地步——那是从概形观念周围几个极简单的力量观念出发的。即便在这个个案里——我一口气投入了十二年数学家能量的很大一部分——我也远没有“收尾”预定的纲领:为此,我还得再要整整十二年!(而我出发之后,没有人操心去继续这项任务——它想必(错误地)显得吃力不讨好……)

作为我重思过一门科学、却肯定没有走得这么远的其他个案,我举出\textbf{同调代数}(交换的与非交换的——后者在我1955年最初的反思时还不存在),以及\textbf{拓扑}:随着\textbf{topos}观念的引入——它至今仍在等候它的时辰,好成为几何拓扑学家的日常面包,与人们今日惯于使用的种种“空间”“簇”观念平起平坐\fnote{184}{(**)}{可与3月底那条注“我的孤儿”(n\textsuperscript{o}~46)第二部分里的某些评注、及其子注~n\textsuperscript{o}~46$^5$至46$^7$对读。}。无疑,今日拓扑的某些重要部分,将不大被拓扑学中topos视点的系统展开所触及\dmark{}\origpage{637}。于是,这个视点依我看,宁可说是“凭空创造一门新科学”里的关键要素——这门科学实现了代数几何、拓扑与算术的综合(当我在五十年代乍到之时,它还全然出人意料)\fnote{185}{(*)}{见前一条脚注。(1985年3月11日)“全然出人意料”一语想必过分,因为这样一份综合的先见,已见于Weil猜想——它们曾如一股强大的灵感之源起作用。}。在新代数几何的奠基之外、并经由“平展上同调的掌握”(及随之而来的$\ell$-adic上同调的掌握),规划并施工这项尚在生成的新科学、为它发展坚实的技术基础——依我看,这才是我对我的时代的数学的主要贡献。那“Motif的yoga”——它仍然……尚属猜想,依我看像这份新科学的灵魂——或至少像它最要害的部分之一:它如此宏大,以致至今我还不曾想过给它一个名字。或许可以称它为\textbf{算术几何}:这名字暗示一幅“几何”的图景——人们将在“绝对基”$\operatorname{Spec}\mathbb{Z}$“之上”展开它;它容许向各种特征的“代数几何”之特化,也容许向“超越”的几何观念(在$\mathbb{R}$、$\mathbb{C}$或$\mathbb{Q}_\ell$等基域之上)的特化——经由“簇”(或更好地说,\textbf{多重簇})的解析或刚性解析的观念、及其种种变体。

我还看到又一门“新科学”——六十年代我就已隐约望见它,它汲取自1955年开始的同调代数反思。那是一场宏大的综合:综合来自同调代数的观念(它在与代数几何的需要——或更好地说,“算术几何”的需要——的接触中发展)、来自同伦代数、来自“一般拓扑”的topos版本,以及(自六十年代起尚在冥府的)(非严格的)$\infty$-范畴理论——或者,按我如今更愿意说的:$\infty$-叠($\infty$-champs)。我曾以为理所当然:这场综合将交由我那些上同调家的学生中的一些人接手——领头的是Verdier,他那著名的论文\fnote{186}{(*)}{关于这一点,见“赊账的学位论文与全险”一注(n\textsuperscript{o}~81)。}恰恰被认为朝这个方向走。依我看\origpage{638},发展一份令人满意的共同语言——具有一切可希求的普遍性与柔韧性——理应是几个年头的事:由一小撮心怀热望的研究者,做一场肯定引人入胜的工作。在我那些上同调家的学生作出一些极零碎的开端之后,我1970年的出发,敲响了放弃这项工作纲领的信号——连同许多别的我引以为心的纲领。这就是为什么,在1975年与Larry~Breen的书信往来里,我回到自己的某些想法,盼着能看到一幅世界之景重新活过来——我分明感到它“在路上”,而“人人”每每一与它照面,便小心翼翼地绕着走。在我致Larry~Breen的信里(重印于《追寻栈》第一章),我提议把这门尚在孕育的科学称作\textbf{拓扑代数}——一二十年以来,唯有我一个人隐约望见它\fnote{187}{(**)}{至多Deligne例外:我曾自信把一幅景传递给了他,而在我出发后的翌日,他便赶着把它与其余的一切一道埋葬。在《收获与播种》里,我屡次暗示我这整套“全方位几何”之奠基纲领中最古老的部分——尤其在“做梦人”(第~n\textsuperscript{o}~6节)里,以及在“我的孤儿”“本能与时尚——或强者的法则”“同伙”诸注(n\textsuperscript{o}~46、48、63$'$)里。}。终于,兴味索然、眼看无人肯接手这项二十年来燃着要做的工作,我便于1973年2月动起手来,写下《追寻栈》,至少以大轮廓画出我心目中该做之事的总设计。

分明地,方才所谈的“算术几何”,与拓扑代数,没有公度可言——而拓扑代数,依我看,它的一个主要角色,正是这门新几何发展中的“后勤支撑”。要使这门几何达到完全成熟的阶段——比如说,达到对Motif观念的掌握,堪比我们对平展上同调的掌握——想必须预期:好几代几何学家都得投身其中,比我目睹在场的那些更有活力、更大胆;更不必说,在\textbf{阿那贝尔式代数几何}的层面上,达到一份堪比的掌握——它(连同Motif)在我看来是算术几何当下便可辨认的两处“要害”之一\dmark{}\origpage{639}\fnote{188}{(*)}{(阿那贝尔式代数几何的若干主导观念,见《纲领概要》第2、3节。)凡言“要害”,我指这门“算术”几何的一个部分:它为这门几何带来直觉、引线与问题——相对于六十年代的所得,全然是新的。(这份“所得”,其实质在于一个框架、一门语言、一份同调与同伦的形式主义——为算术几何所涵盖的三门学科所共有。)或许,还该在上述两处“要害”之外,补上第三处:它与Motif亲密相连,即Langlands式的\textbf{自守形式}理论。我之所以对它略而不谈,是因为我对自守函数理论那份说来遗憾的、从来的无知。(我不知道机会会不会出现,逼我终于稍稍补上这份无知……)}。

还有第四个反思的方向,在我的数学家过去里进行过:朝着一门既有学科的“彻底”更新。那便是拓扑学里的“温和拓扑”进法——我在《纲领概要》(第5、6节)里对它略有铺陈。此处,一如自中学那些遥远的年头以来的许多次,似乎仍旧唯有我一个人感到那场该做的奠基工作的丰富与紧迫——而这份需要,在我看来,比以往任何时候都明显。我有一份极清晰的感受:温和拓扑视点的展开——照《纲领概要》所唤起的精神——将给拓扑学带来一场其分量堪比概形视点带给代数几何的更新,而这并不要求相当规模的能量投入。此外,我想,这样一门温和拓扑,终将在算术几何的发展中证明是一把珍贵的工具:比如,得以表述并证明“比较定理”——在复数域上有限型的分层概形所系的“投射有限”同伦结构(或更一般地,在该域上一个有限型的分层概形簇所系的同伦结构)与相应的“离散”同伦结构之间,后者以超越的途径定义,并以合宜的假设(尤其等奇异性假设)为模。这个问题,唯有在分层结构的一套精确“拆解理论”的术语里才有意义;而在“超越”拓扑的框架里,依我看,它需要引入“温和”的语境。

\asterisms

\dmark{}回到朋友Pierre~Deligne这个人:在1965至1970年与我数学上贴近接触的那些年里,他有的是机会,全面熟悉我方才以大轮廓巡过的这一整套观念与几何之景\origpage{640}。(温和拓扑的观念除外——如果我没记错,只是从七十年代初起,它们才开始发芽、并叫我着迷。)他对这份宏大纲领的角色是双重的、朝两个相反的方向。一方面,借着现成的$\ell$-adic上同调工具、借着(仍被遮蔽的)Motif理论的观念,他为算术几何纲领的发展作出了显著的贡献。最重要的,无疑是一门混合Hodge系数理论的启动,而尤其,是他关于Weil猜想、及其$\ell$-adic推广的工作。另一方面,撇开他工作直接需要的那些\textbf{工具}与观念(它们的来历,他曾系统地设法让人忘掉)不谈,他竭尽所能,叫其余的一切之自然发展落空:这便是“电锯效应”——我在葬礼反思的进程里已有的是机会谈它,包括(以暗示的方式)在上一条注(n\textsuperscript{o}~136)里。这电锯效应,曾因1981与1982年的几次局部开棺而部分被搅浑——“像几株孱弱的幼芽,重新发出……”——那是迫于眼前需要的骤然推力。(这些应景的开棺,上一条注的结尾处刚刚又提到过。)他还竭尽所能,不断地给人(从不曾明说……)这样一种印象:他所使用的那些观念、想法、技术、成果——他留神隐瞒其来历——应归功于他;而当他不曾慷慨地把它们派给昔日学生或合作者中的另一位时。

算总账——在这份对他那般顽韧地锯断与下葬之物的快速回顾之后——我回到上一条注里占上风的印象:我曾暗示,“东家”、利己的贪欲之插手他的工作,主要限于工作\textbf{主题}的拣选。毕竟,那掘墓人-电锯的姿态,在他的工作里分明可见——例外极少——\textbf{凡}有机会处\textbf{皆然};而我意识到,这些“机会”数不胜数!\textbf{这份掘墓人的综合征}(它无疑与超级阳价值的抬举亲密相连),依我看,\dmark{}对他的工作与著作\origpage{641}起过真正“具有侵夺性”的效应,与他的那些亲阳抉择的效应,毫无公度可言;而且这效应,绝不限于主题的拣选——“东家”把主题交到“匠人-孩子”手边,随后踮着脚尖退下。相反,依我看,整个工作进程中,“东家”几乎从不离开“匠人”:他太担心后者会忘记那些强制性的指令;换句话说,工作本身,常常被一些与发现工作之本性格格不入的\textbf{内在姿态}侵夺——那本性,是冲进未知。这件事,在葬礼反思的进程里曾被强烈地感到过多次;而我在这场关于阴阳的长反思中,曾倾向于把它忘了。

\subsection{天鹅绒中的利爪}
\label{ch18:local:h48}

\subsubsection*{(a) 天鹅绒的脚掌——或微笑}
\tocsubsub{(a) 天鹅绒的脚掌——或微笑}
\label{ch18:local:h49}

\notehead{137}{}
\phantomsection\label{ch18:local:n45}

(12月7日)一个多星期没有继续写注了,除了后勤的工作(包括给前面两条注写子注)。我不得不拔了三颗牙(快到六十的人,就是如此……)——必要而粗暴的一次入侵,害得我近来以略为减档的转速运转。我趁势转向积压的通信。那里,一切似乎已归于秩序……

在前面的四条注里(11月24至28日),我主要试着更贴切地圈定:Deligne与我之间,在气质与数学进路上的亲缘或互补关系——为的是给那“颠倒”定位:那份阴阳角色的颠倒,我曾自信在朋友竭力作出的自我呈现、与对我的呈现之中(至少在双方“数学”人格的层面上)知觉到。一路走来,关于他或我自己的实在的另一些侧面显了形;而在我们两个人之外,还有数学家世界的侧面,或者干脆说,人的世界的侧面。最终,依我看,最惹眼的新东西,是服务的姿态、以及这份姿态在科学世界里消失的种种征兆——我便以给这一阶段反思所起的“主人们与仆人”之名,试着暗示这一点。

\dmark{}回到“定位”某场颠倒这个最初的打算:此刻我有这样的印象\origpage{642},关于我与朋友的真实情势,我已圈定得足够贴近,可以给它下文了。一条不容回避的头一条确认是:那份关于阴阳角色颠倒的出发直觉——它在5月12日“葬礼悼词(1)——或恭维”那场反思的翌日来到我心头——确确实实是正确的。早在11月10日“阴的葬礼(阳葬阴(4))”(n\textsuperscript{o}~124)那条注的反思里,便已清楚:我的朋友竭力给他自己描画一份超阳刚的影像,给我描画一份超阴柔的影像。11月24日那条注“颠倒(3)——或阴葬阳”(n\textsuperscript{o}~133)里提出的问题是:这份呈现,是否确确实实是实在的一场“颠倒”。而“上涨的海……”那条注(n\textsuperscript{o}~122)里显出的“新事实”——即:一如我的朋友,我的数学进法的基调也是阴的、“女性”的——曾一度使人对此生疑。

但最后三条注的反思,驱散了这份疑团。一开始便已清楚:从来,我都被Deligne(一如被我的其他学生与昔日学生)知觉为——至少在意识的层面上——强阳的(也许太阳了……)\fnote{189}{(*)}{况且,时行的价值既然如此,我怀疑一份科学声名能否由一幅(被普遍承认、被普遍接收的)影像来承载,倘若那影像不是一幅“阳”的、乃至超级阳的影像。依我看,唯有在无意识的层面,我数学进法的“女性”本性,才曾被我的朋友兼昔日学生、以及数学公众(至少是与我工作的那类事情多少有接触的公众)所知觉。}。但事情显出:此外,在Deligne与我的数学层面的关系里,在一份强的阴-阴亲缘的底子上,还运作着一份阴-阳的\textbf{互补}(相对于那份扮演“原初”角色的亲缘,可以称作“次生的”)——在那份互补里,扮演“阳”的、阳刚角色的,确确实实是我:凭一份在我身上远比在他身上显豁的“阴中之阳”成分。

我在Deligne身上确认的那份故意——它依我看,从好些方向得到殷勤的回声\fnote{190}{(**)}{我这里想到《导言》里唤起的那“阵阵阴险的轻蔑与含蓄的嘲弄”(见《导言》10,“一份敬意的举动”)。当我看到昔日学生中最显赫的几位亲自定下这个调门时,我无须惊讶。这些年里传到我这里的大量“阵阵”之中,共同的,恰恰是一份居高临下的做派——对着我的数学进法与著作里打有强烈印记的“阴”特征。关于这一点,亦见6月23日那条脚注里的评注,注~n\textsuperscript{o}~96“棺材4——或没有鲜花没有花圈的topos”。}——因此向我显为确确实实一份\textbf{颠倒角色的故意}\dmark{}\origpage{643},更确切地说,\textbf{阴阳角色的颠倒}\fnote{191}{(*)}{这份颠倒角色的故意,头一次在我的反思里出现时,说的是师徒关系里的角色颠倒:我被呈现为我学生的“合作者”,而他本人则扮起平展上同调与$\ell$-adic上同调的\textbf{真正}奠基者与主人。(关于这一点,见“颠倒”与“葬礼悼词(1)——或恭维”两注(n\textsuperscript{o}~68$'$、104)。)耐人寻味的是:在“师——徒”这对“偶子”里,扮演阳的角色的(作为给予者、言说者)、“主动”的,正是师父;扮演阴的角色的(作为接受者、聆听者)、“被动”的,是徒弟。这里又是一次:我昔日学生漂亮地完成的颠倒,可以看作一份阴阳角色的颠倒——取同一方向(阴-阳变成阳-阴),一如构成我葬礼悼词之主要讯息的那场颠倒——那讯息,显形于“阴的葬礼(阳葬阴(4))”那条注里。}。依我看,这是葬礼的又一个重要侧面,叠加到前面已巡过的四个侧面之上(11月13与17日的“回顾(1)(2)”两条注,n\textsuperscript{o}~127、127$'$)。如今该做的,是把这五个亲密相连的侧面,拼成一幅连贯的葬礼总图。

这样一幅总图,若要令人信服,还应当在一份共同的视景里,汇集\textbf{三个相继的“层面”}。第一层面上,只有Deligne:我葬礼的大司仪,师父的非学生、非继承人——那师父被宣告死亡,没有存在的理由,也从未存在过……分明地,他是葬礼仪式的\textbf{中心}人物——除了亡人自己(而他,不过是个亡人,一个沉默的配角)。紧随其后,在第二层面上,是“我那些昔日学生的忙碌的一群,扛着铁锹与绳索”(凭记忆引用“葬礼的行列编成”里对各行列的罗列)。最后,在第三层面上,是(几乎)全体到场的会众:他们前来庆祝我的葬礼(以及四位共同亡人的葬礼——他们在钉得结结实实的“橡木棺材”里立正待命),并为下葬助一臂之力。

\dmark{}在这三个层面之间,似乎统治着一份完美的和谐、一份“\textbf{全体一致的协约}”——一如人们在任何一场照规矩举行的葬礼上\origpage{644}见到的:在满怀虔敬的哀恸的神甫、摆着合宜神色的亡人家属、与广大的会众之间:该起调的地方起调,该静默的地方静默,从不、从不走调。

顺着最后这幅意象说下去:我此刻看到自己被置于这样一个人的境地(不如那位亲爱的亡人自在——他分明置身事外……):面对如此感人的一幕,竟要无礼地猜测神甫、家属与普通信众在那些应景的庄严或悔罪神色背后,究竟怀着怎样的真实想法与动机。

反思已继续了一段时日,其暗中一贯的主导线索,便是那个打算:预备必要的条件,以便把握总图三个“层面”中最切近的那一个——披祭披的神甫的那一层,对不起,我是说我的朋友Pierre~Deligne的那一层。现在我愿把注意放到这一层面上。

我要一上来就说:总图诸侧面(或“联”)之中,在“掘墓人——或全体教众”那条注(n\textsuperscript{o}~97)里被摆到前头的那个——“对一场异见的报复”——在我看来,在我的朋友身上,即便它登场,也只扮演一份极不显眼的角色。我从不曾有过片刻的印象,觉得朋友皮埃尔哪怕丝毫感到被我的“异见”“牵连”。恰恰相反:这份异见是一场天赐的运气——他恐怕从来不敢梦想的运气——使他得以优雅地摆脱一位在场感过强的师父:就在那家机构里,他年方二十五,刚刚获得了数学世界里最令人艳羡(或至少最令人神往)的地位之一。这份异见在随后的月与年里愈演愈烈,依我看(或许不在意识的层面,但那无关紧要),被经验为一份更大的运气:它把一份可观的“遗产”拱手送到他嘴边,没有来自任何方向的抵抗之意向(如他随岁月推移而逐步确认的)\fnote{192}{(*)}{关于这份“遗产”,见“继承人”一注(n\textsuperscript{o}~90),以及“阴,作为仆人(2)——与慷慨”一注(n\textsuperscript{o}~136)的子注(n\textsuperscript{o}~136$_{1}$)。}。这份不期而至的\origpage{645}运气,他可不曾作态\dmark{}抱怨——哪怕在心底里、或在他自己不知不觉中!而且我觉得,等比例地,同样的确认,对我“出发”之前的大多数学生、无论如何对我那五位上同调家的学生中的每一位,想必都成立。他们之中若有人曾暗示过一份不满、一份因我的异见而来的受挫之感——无论是在心底里,还是以或多或少明白的方式\fnote{193}{(*)}{唯一一位曾让我听到这种调门(还带上一丝谴责的口气)的昔日学生,是Verdier,约莫一年之前。而在“生存与生活”的年代,他却似乎同情我的异见。甚至有过一段与他妻子伊冯娜友好合作的情节:记得是(如果我记不错)在Robert~Jaulin(伊冯娜曾是他的学生)倡议的一次巡回展览的组织上,我以“生存派”参与者的身份也参加了进去。}——我倾向于相信:那属于一种\textbf{合理化},为那对着天助地消失的师父的掘墓姿态所作的合理化,而远不是它的\textbf{缘由}(哪怕它是缘由之一)。使我在这份确信中愈发坚定——既就我那些上同调家的学生“一般”而言,就他们那位无可争议的头领Deligne而言——的是:那场即将到来的葬礼的先兆(只要合适的机会出现——而啊,意想不到的奇迹,它出现了!),早在我1970年出发之前便已可见,无论如何,早在那个著名的1965/66年SGA~5研讨班之后便已可见——那研讨班,注定要遭受我所知道的那场屠杀。这肯定不是偶然:如此完美地,他们五位全体\fnote{194}{(**)}{(12月12日)可我该把J.-P.~Jouanolou除外:他终于为研讨班接连编写了三讲,展开了一些观念与技术——那正是他自己的学位论文工作马上直接需要的。}都对那研讨班的命运漠不关心——他们学手艺的地方,连同那十二年间几乎唯有他们有特权认识、使用的漂亮数学一道。关于SGA~5的命运,我在反思中已说得够多,此处无须赘言。关于Deligne,我只提醒一句:在他1970年我出发之前所写的四篇文章里的三篇中,隐藏、或至少尽量遮掩并贬抑我的观念之影响的意图,分明可见——它并没有等我的“异见”。

\asterisms

\dmark{}那么,我的朋友身上那份对我对抗姿态——那份急欲取而代之\origpage{646}、抹而除之的贪心竞争者——的根与特殊本性,究竟是什么?这份姿态,自我与他相遇的最初几年起,便与一份温情的、信赖的同情、与数学层面的一份共融共存。我甚至确信:它自我相遇之日起便以弱音在场,无疑更早;而且,它与其说是被我身上的这个那个特点唤起的,不如说从一开始便是从我注定要在他身边扮演的那份角色中流出的——除非说,是那整个“特点”的总体:它们使我得以在他身边扮演那份角色。这也是他二十年来竭力抹去的角色;那份角色,无疑不出双方所求、纯凭情势之力,含有一个“父性”的侧面。而我心中毫无怀疑:冲突正是系于这一侧面打成结的——一份早已在他身上的冲突,远在他听说有人念出我的名字之前、甚至(无疑)在我们共同的师傅——数学——的名字被念出之前。

这份确信,老实说,不是一场反思的果实,我更不自许去“证明”它。毋宁说,它是在我出发之后的岁月里来到的,我自己也说不清何时如何;我相信,是渐渐的,凭着大大小小的征兆——它们没有一个曾让我驻足,哪怕一瞬——而它们全体,却终于留下一份认识的痕迹:弥散而欠完整,却终究是一份认识,有一天它就在那儿了……我当然可以凭一项费力的工作——把半埋的记忆挖出、逐个探测——深化并落实这份仍然几分难以称量的认识;而很可能,这样的工作会给我留下许多惊奇。可我不觉得自己有兴致去做。这无疑是因为(无论对错)我觉得:这真正说来不是我的工作,而是我朋友的工作——我那里将探测的东西,事关他的,远甚于事关我的。就我而言,我方才表述的这份直觉、或“认识”、或“信念”,已够满足我眼下理解的愿望;我对它毫无保留地信赖。

一如我一生中常常如此,我此刻面对的是一份\dmark{}\textbf{对父亲的对抗}\origpage{647}关系——我在其中扮演替代父亲、被他认作父亲的角色(依我看,与其说是收养他的父亲,不如说是被他认作父亲的人\fnote{195}{(*)}{(12月12日)写这几行时,我有意识:在这样一份角色“非对称”的断言里,该多么谨慎——何况这些角色是在无意识的层面上上演的。我猜,在这个层面上,在真正的数学交流之外,我想必曾在某个时刻,多少进入那被语境备好的“父性”角色。可这份角色,分明在我的生活里、在我与朋友的关系里,分量不能与我对数学的激情相比;它始终是偶发的,而且在我1970年从数学舞台“出发”之后,该不会再有痕迹。相反,我昔日学生对我这个人的眷恋——无论好也罢、(尤其)坏也罢——在其后的整整十五年里不断显形:既在他自己的工作里,也凭那份顶着逆风维持的、与我的持续的个人关系。}),连同我朋友那份颠倒阴阳角色的故意,旋即在我心中与“颠倒(2)——或暧昧的反叛”(n\textsuperscript{o}~132)那条注里唤起的情势联系在一起——我母亲与她父亲的关系,对我而言,是那份情势最极端的原型。然而,那份情势与朋友皮埃尔对我这份关系的差别,一望即知。在他与我的关系里,我从不曾知觉到一丝“反叛”的调门之影,或哪怕一份以些许恶毒、富于攻击性的方式、亮出爪或牙的对抗——哪怕是在一个微笑里。微笑,双方当然都不缺;但那是他一侧的同情之笑(如我所感到的),或有时是无邪的惊讶之笑,又有时几乎是忧伤的笑:当他得以确认(而我终于察觉到其中那一丝内心的得意)某些出其不意、以天鹅绒的脚掌落下的招式,在预定之处命中了。

换一种说法:这份对抗——无论它对着我,还是对着别人(当要经由他们打击那位亡故的师父,而他在他身上又始终好好地在场……)——从来、无一例外,取极阴的形式:那以最精微的温柔的全部表象为乐(并精于此道)、去打击去伤害、乃至去清除去压服的形式。他为数学家的招牌所作的刻意抉择是超级阳的(一如我的抉择也曾是——但没有比他更多成功);而依我看,在关系的层面上——至少对着我、对着那些他认为与我有牵连的人——那基调分明是、彻头彻尾是,超级阴的。(不过,在这件事上我要作唯一的、而且是重要的保留,日后还须回头谈它。)

另一处“一望即知”的差别——皮埃尔与我的关系、与那份“暧昧的反叛”之间:凭我对他家人的所知,我相信皮埃尔的父亲是一位性情温和而谦逊的人,因此绝非会激起一场反叛、再把反叛转嫁到替补父亲身上的那号“形象”\origpage{648}。

\subsubsection*{(b) 颠倒(4)——或夫妻的马戏}
\tocsubsub{(b) 颠倒(4)——或夫妻的马戏}
\label{ch18:local:h50}

\notehead{138}{}
\phantomsection\label{ch18:local:n46}

(12月8日)昨晚结束反思时,我有一种颇不好受的印象:觉得自己越来越看不懂了。睡前,我又逗留片刻,追随那场已写完的反思所引发的联想。我自信看到几点光亮浮现,我想它们将在今天的反思里充当我的灯具。

这些联想中最重要的一个,无疑系于我朋友身上那个“天鹅绒的脚掌”的侧面:以世上最天真的神态、“以最精微的温柔的全部表象”,津津有味地(而且有时入骨而不留情面地)抓挠。这幅意象,是在一场(与先前唤起的一份“反叛”情势所作的)已然触礁的比较的转弯处来到的,当场便向我显为意味深长,显为那份我打算探测的“对抗”的一个要害侧面。而且,回头看,“天真微笑与天鹅绒的脚掌”这幅意象的唤起——它复现着将近二十年贴身经历的精髓——依我看,正是昨天反思里的那个“痛点”、那个意想不到的“亮点”,当时我正在黑暗中摸索。倘若那份摸索与晦暗的印象在更远处仍占上风,那是因为:我太被片刻之前还在脑中的那些念头抓住——那些有待继续、有待安放的念头——以致不曾留心那精微的“咔哒”一声:它在意象显现的那一刻,便已在我身上响起。而其后的半个小时里,顺着系于这幅意象、与那场反思的另一个时刻的几场联想,注意又散了。直到现在,带着一天的回顾,接上被打断的反思之线,我才看到一份先前重读昨天的注时还不曾察觉的视景,正在就位。

倘若我留神追随所有联想中最强、与我亲身经历联系最紧密的那一条——暂且撇开其他更“有结构”、更“智性”的\dmark{}\origpage{649}——得到的是这个。我忽然看见自己回返到那一系列特殊的个案——或作为共同行动者经历、或作为近旁的证人——它们可以由一份单一的印象概括:\textbf{夫妻的马戏}——男与女这对偶子的马戏。这对偶子的马戏:成婚与否、有无孩子、年轻年老或少老颠倒、贫贱中拉着重负爬坡或富足中高车驷马——全然一样,夫妻的马戏并不因此改变。我忽然看见自己回返到它,经由这个马戏中打动我至深的一个侧面(老实说,过了很久,我才看穿其中的门道……):那非常特殊的、十足“天真相”的、“我什么也没说没做”的战术——“天鹅绒的脚掌”的战术——由女人在一个特定的游戏里使出:在那游戏里,始终是她以完美的手法、不动声色地领着;而始终是他跟着(而且常常挨着疼),什么都察觉不到。我所见过的夫妻之中,极少数不是循着这个调子运转的——变化无穷,自不待言:那留给双方即兴的天资,还有各自的脾气与种种情势。而且,就在今天,我还有机会看到一份格外眼花缭乱的演示——不过我谢绝在此为它枝蔓。

对这些马戏游戏的一份多少有色彩、有层次的描述——至少在大轮廓上——或者哪怕只是对其中调门的唤起(正是“她”那一侧的“天鹅绒的脚掌”),正是我刚重读过的11月12日那场反思——“颠倒(1)——或激烈泼辣的妻子”(n\textsuperscript{o}~126)——里缺席的那位要角。分明地,我逆着一阵迟疑推进那场反思,以致它终于摆出一副严峻的“力量与动机”分析的架势——那一天我分明不在状态!那也是在“阴阳之钥”里头一次谈到“阴阳的颠倒”。当时有点迷住我的那个极端个案——它直到昨天还在起作用——是我母亲的个案(重述于11月22日那条注“颠倒(2)——或暧昧的反叛”里,n\textsuperscript{o}~132)。可是,在我那份“四点分析尝试”里,我留神把这三“点”之前的几点圈定为:适用于我有机会多少贴身认识的大多数(即便不是全部)夫妻,而不必然以(暧昧的)“反叛”那份烈性的调门为主导(哪怕它以被遮蔽的形式在场)\dmark{}\origpage{650}。但这并不妨碍还有另一件共同的东西,它那天逃过了我。它只是昨夜才开始显出端倪——在那半小时里,我顺着“在状态”的那场反思的尾流,任思绪游荡。这件重要的共同东西——此前我只在“激烈泼辣的妻子”的极端个案里知觉过它——便是\textbf{阴阳角色颠倒}的精微游戏。

我犹豫:该写这份游戏是那份我方才暗示的权力游戏的“发条”,还是它\textbf{等同于}后者?肯定地,对她(而且常常也对他)构成男性角色之精髓、构成派给男人的角色的,是\textbf{权力的占有}——那占有固然常常是虚的,却无论如何从社会共识中汲取一份实在的元素。或许我曾倾向于低估这份实在元素的力道、男人的\textbf{象征}之力道:作为面对女人的一份\textbf{权威}的代表——尤其,它的力道,作为女人动机里的一个动力要素。我怀疑,对她而言,“做男人”、“做那个男人”,在一切之先,便是\textbf{行使权力}。在利己动机的层面上\fnote{196}{(*)}{阴阳角色的颠倒——在性爱冲动的层面上、在爱的游戏里——别处曾顺带谈到。(尤其见“接纳(阴的苏醒(2))”那条注。)性爱冲动,凭其本性,与自我的游戏无缘——尤其与权力的游戏无缘,即便自我渴望把它变成一件为自身目的服务的工具,并且(至少在某些窄窄的限度之内、以扭曲并残害那份原初冲动为代价)精于把它弄到手。倘若这两种“颠倒”之间——一边是阴与阳两股冲动那份\textbf{既}在情人身上、\textbf{又}在情郎身上的自由游戏,另一边是配偶一方对另一方那份无休止而阴险的权力炫耀的强迫游戏——果真有一层关系,那依我看,这关系多半只能是:两种形态,在任何时刻,都互相排除。},那“角色的颠倒”,想必不多不少,正是\textbf{女人对男人的权力行使}。

鉴于现存的共识,女人这份权力的行使,几乎只能以隐蔽的方式进行。它不在于发号施令,也不在于\dmark{}装模作样地做决定(期待决定被遵行),\origpage{651}而在于\textbf{让人运转}——尤其,把人折腾得团团转,而且始终不留痕迹。这便是那著名的夫妻旋转木马:它转动着,从不停歇!使它保持运转的战术……战术——由母亲传给女儿、由女人或少女传给小女孩、代代相传、无须言语的战术,便是昨天在路的转弯处唤起的那份战术:\textbf{“天鹅绒的脚掌”的战术}。只要留点心,人们便在无限的各式面目之下认出它:从“激烈泼辣的妻子”那极阳的个案——对我来说,以我母亲为化身——到“哀怨的(乃至愁苦的)妻子”那极阴的个案——我曾见它以另一位近亲为化身。

依我看,极少有女人不实行这份亘古的战术、不把它彻底掌握\fnote{197}{(*)}{同样真确的是:极少有男人,当“人家”对他们使出这份战术时,不立刻“入彀”。我自己便在我一生的大部分时间里,不战而降。它真正开始改变,只是随着禅观在我生活中的出现——那年我四十八岁(做好事,永远不嫌晚)。直到今天,我偶尔仍会上当。(说实话,不常,而且从来不会太久……)}。它主要是夫妻马戏里的日常实践,却并不限于此。依我看,女人对女人,它少有施展(或许仅仅因为:让一个女人“转起来”,比让一个男人难)。相反,在某些女人身上,这份战术变成一种第二天性,施于\textbf{一切}男人、或几乎一切——至少,施于那些被她知觉为具有鲜明阳刚品格的男人。

倘若我在这里说“战术”,这其实只表达了一个次要的侧面——恰恰是“战术”的侧面——的一份更要紧的实在:那份积习已深的内在姿态,对着一般“男人”的,或至少对着那个人的——父亲、情郎或丈夫,他在她的生活里扮演一份特权的\textbf{男人}角色,被(社会共识,或被她自己的抉择)授予一份\textbf{权威}。这份姿态,绝非总取支配欲的品格(如“激烈泼辣的妻子”的个案)——至少,不取人们通常理解“支配”一词的意义。它毋宁是一份饥渴——它有时变得吞噬一切:\textbf{不断地对另一方施加影响}\dmark{}、把他“维持在运转中”(言外之意:绕着她这个人运转……)\origpage{652}。为此,常常一切手段都是好的。其中一种行使行动——并由此行使权力的手段,便是\textbf{伤害}:而且有时,伤得越深越好,干脆把人放倒,直至毁掉——身体的或心理的——只要时机合适;而且始终,不着痕迹,“以最精微的温柔的全部表象”。我自己就不止一次被“放翻在地”!也常常,作为共同行动者或证人猝不及防,我被那伤人或毁人的举动之表面无端性惊得透不过气——它伴着一个天真的微笑、或一副心不在焉的神情,却始终不动声色,凭着一份不出错的直觉,抓住时机与地点,打击对方最能被击中的所在——无论那“对方”是父亲还是情郎、丈夫还是孩子,或仅仅一个相识、乃至一个陌生人(只要有机会出手、有机会命中……)。

\subsubsection*{(c) 天真的暴力——或传递}
\tocsubsub{(c) 天真的暴力——或传递}
\label{ch18:local:h51}

\notehead{139}{}
\phantomsection\label{ch18:local:n47}

(12月9日)我在这里触及那个极端、却绝非罕见的个案:\textbf{为暴力而暴力}、暴力与恶毒之中的\textbf{无端性}。这份暴力——无论它打击陌生人,还是最亲近、号称被爱的人——既非女人所专有,也非男人所专有;它既不“阴”,也不“阳”。可我此刻遇见它时所取的那份令人错愕而阴险的\textbf{形式}——戴着一副心不在焉、乃至天真温柔的假面——这份已然变得为我所熟稔的形式,依我看,尤其为女人所专有。这肯定系于那份“父权”的社会共识:它就女人而言,把权威与权力授予男人\dmark{}\fnote{198}{(*)}{况且,这份共识、以及男人在其与女人的关系中的权威,在最近几代人里已大遭侵蚀,如今愈发如此。对这份侵蚀,我将是最后一个抱怨的人!然而,法律与风俗里这份浮皮潦草的改变,似乎并不曾在两性关系的深层发条与“风格”里改变什么——尤其,不曾改变女人对着男人那份彻头彻尾、被精心遮蔽的对抗。这系于一个事实——本注反思的结尾处已强调过:这份对抗的姿态、及其借某份权力游戏(或权力颠倒的游戏)表达自身的手段,与其说是家庭之内种种“客观”条件的结果,远不如说是一份“\textbf{遗产}”之代代\textbf{传递}的结果。}\origpage{653}。这份形式,是她满足一份权力意志的手段——那份意志,纵然被迫(情势使然)走与男人不同的路,在她身上却并不因此少一分强制性、少一分吞噬性——恰恰相反!似乎不能在光天化日之下舒展、被预先注定过一份隐蔽的存在,只会加剧这份饥渴、并使它愈加蔓延,以致在许多个案里,真正“吞噬”她的生活、与她亲近之人的生活。

这份饥渴,远非总能(所幸如此!)达到全方位无端暴力的规模;而它得以施展的各个音区,也并不都落在暴力的调门里。虽然那份含蓄嘲弄的调门最常见——它们给一份蒙着面纱的对抗、或一份隐秘的敌意放风——那些单纯淘气的调门也并不被排除:其中带着一份宽容的亲昵,还透着几分顽皮。而如果说“天鹅绒的脚掌”这份经过考验的战术是女人的特权与首选武器,这份特权却绝非排他的。我不止一次、而且贴身地,见过男人挥舞这件武器\fnote{199}{(*)}{不过,在为我所知的个案里,我记下这份差别:倘若面对亲近或友好之人有一种表面“无端”(我是说,未被挑起)的暴力,那么每一次,那都是一个他对之怀有(哪怕不自觉)一份宿怨或长期敌意的人——那敌意落实为具体的怨怼(即便它们多半未曾成形)。这方面唯一的例外,是我的朋友Pierre~Deligne,在他对我的关系、以及在他归并到我这个人身上、当作属于我“势力范围”的那些人的关系里。于是,那是一份对抗与暴力的姿态(当然,压低了嗓门!),没有“个人”的缘由——我是说:没有缘于他对那些他竭力打击的人所怀的怨怼(真实的或想象的)。而反过来说,这是一种在许多女人身上遇见的行为——而且不仅(如这里)对着亲近的朋友、或单纯的相识乃至陌生人,也对着最亲近者中的某一个:比如情郎或丈夫(当然,而且优先),或兄弟、甚至自己的孩子。},其手法同样纯熟,毫不逊色\fnote{200}{(**)}{况且似乎,这份由无意识掌控的战术,总是从无意识那里继承那份“手法”与那份近乎不出错的稳当——它们在一份完全有意识的行动里,何其罕见。我不记得见过谁使这份战术,而不带着一份大师的手法。}。惹人注意的是:在这些个案的每一个里\dmark{},那据有了这件女人专属武器的男人,都是一个倾向于压抑自己存在中某些阳刚侧面、并(无疑由此)照着\textbf{母亲的模型}自我塑形的人\origpage{654}。

这同一份战术,也常见于孩子们——女孩男孩不分——对着父母、或对着代行父母之职的其他大人所玩的权力游戏之中,而且几乎是通则。这旋即引出另一份联想:在某些国家(过去的或当下的)里,作家或记者的情势——那里盛行着直接或间接的书报检查,使坦诚直白地表达真实思想感情成为不可能或有风险。后一种个案与前几种的主要差别在于:这里,诉诸间接的、蒙着面纱的、有时象征的表达,已不是无意识的作品,而是有意识的思想的作品。缘由无疑在于:当时存在着一份足够广布的共识——支持那些非正统的思想与感情(它们正要不动声色地“传递过去”)——以致当事人不再感到有义务把它们对自己也藏起来,生怕在自己眼中显得是个可憎的异类。唯有在极端酷烈的政治或宗教恐怖的个案里(如中世纪,或斯大林时代的苏联及其卫星国),非正统的意向才被迫(至少在某些人身上)再深潜一层:躲开内心检查官的目光,一如躲开那设立在风俗与警察机器里的检查。

所有这些例子似乎都暗示:“天鹅绒的脚掌”的风格(或“我什么也没说、没想、没要”的风格),或多或少自动地出现在任何一份多少持久的情势里:在那里,一份于我们不利的实力对比,使坦率直接地表达我们的感情、欲望、想法、意图成为不可能、或至少危险——尤其,那些对着被视为对我们施加强制之人(而尤其,那份恰恰自称要拦住我们表达真实感情的那个强制)的敌意或敌视之情\fnote{201}{(*)}{写这几行时,一个念头显现:我方才描述的情势,恰恰正是我们人人无一例外地在童年最初几年所面对的情势。我们无意识的很大一部分(那部分可以称作“垃圾滑槽”,在无意识的层面通常被知觉为一种“垃圾坑”)不是别的,正是我们童年的心灵对周围环境那份压力的回答:那压力强迫我们(这几乎是个生存问题)把身上一切触犯社会检查的东西,作为否弃的表示,深深地埋在自己眼前之外。这份检查,不久便内化为一位内心检查官,他阴沉的在场,保证这场早夭的下葬长治久安。然而,尽管有检查官,那些被规规矩矩埋葬的非正统冲动、认识与感情,却得以表达出来——有时以一种加剧而可畏的效力,以间接的、往往象征的、却依然充分具体的方式。“天鹅绒的脚掌”一栏,便提供一份格外“击中要害”的例子——而且常常,令人错愕……}。

\origpage{655}这并不是上述风格、及其所覆盖的内在姿态显形的\dmark{}唯一个案。常常,这份“实力对比”或多或少是虚的:它与其说对应一份“客观的”实在——顾及那个或那些被知觉为“压迫者”的人的真正姿态(或权力手段)——远不如说对应我们对它的\textbf{观念}(有意识的或无意识的)。这份观念,很少是对一份给定现实作细心的、明智的检视之果实;它几乎总是属于我们在幼年所接收的那一整包形形色色的制约——还须加上某些在我们人生极早的年代便已在我们身上完成的基本抉择。这样,无论在女孩还是在男孩身上,与\textbf{母亲}认同的那份抉择(当然是无意识的),便含着采纳一整套姿态与行为(尤其那些以“天鹅绒的脚掌”风格表达自己的),以及同时又采纳支撑它们的种种观念(多半无意识的,但这无关紧要)——比如关于某份实力对比的观念,及伴随这些观念的对抗反射。在相反的、与\textbf{父亲}认同的个案里——但当父亲本人在自己身上整合了某些典型“女性”的特征(或至少在我们社会中如此被视为的特征)时——可以想见,效应可能与第一种情形全然相仿。

我想在此抵达的点是:在我们当今的社会里——至少在我曾置身其中的那些圈子里——依我看,这份风格(“天鹅绒的脚掌”)、以及我此刻检视的这份“女性”内在姿态,只在很有限的限度内,才是对社会或我们童年周围的特殊情势(乃至我们成年某一时刻的情势)所设立之客观实力对比的一份自发的个人反应;它毋宁是一份“\textbf{遗产}”,得自我们父母中的一位(甚至有时同时得自两位?\dmark{}\origpage{656})——而那一位,又是从他自己的父母之一那里得来的。分明地,这份遗产优先循着\textbf{母系}的线路传递,首先由母亲传给女儿。但我不止一次贴身见过由母亲传给儿子的传递。没有任何东西让我以为,这份传递不能偶尔也由父亲传给儿子、甚至由父亲传给女儿。

\subsubsection*{(d) 奴隶与傀儡——或闸门}
\tocsubsub{(d) 奴隶与傀儡——或闸门}
\label{ch18:local:h52}

\notehead{140}{}
\phantomsection\label{ch18:local:n48}

(12月10日)我想回到几场系于\textbf{无端暴力}主题的联想。那正是昨天反思由之起笔的主题;随后我离开了它,转而检视权力游戏中的“女性”风格(或“天鹅绒的脚掌”),以及它作为对着他人的对抗姿态之表达手段(尤其,对着那些被感到强阳的、或以任何身份处于权威、声望或权力之位的男人)。

一如我昨天所提醒的:(表面)无端的暴力、“取乐”的暴力,并不比女人更属于男人。人人都曾有机会在路的转弯处猝然与它照面:或顶着“最精微的温柔”的面孔,或顶着皮靴踹进肚腹、一梭子机枪扫射的面孔。后一种风格,肯定是“阳”的风格,在眼下这些号称“和平”的时代、在我们这样的文明国度里,总归更罕见。对我们大多数人——教养良好的、在富足国度里处境或多或少优渥的人——这份直言其名的暴力,不属于日常经验,一如另一种——那份压低了嗓门的、天真相的暴力——属于日常。然而,只消翻一翻随便哪家大报的“社会新闻”栏,或听听新闻广播\fnote{202}{(*)}{这些事,说实话,我早已不做了,只满足于经旁人之手得到的偶尔的消息。},便会发觉:“硬”的无端暴力,即便在我们这里,依然随处可见。它还不总发展到\dmark{}顺手割断那位\origpage{657}匿名小老太太的喉咙——只因心血来潮要偷她的家。可当百无聊赖想冒险的年轻人“借走”人家大意地敞开在家门前的汽车时,把它丢在十公里二十公里外的沟里之前,却很少有不好好把它砸个稀烂的。即便在我有幸居住、不必过分担心的宁静乡间,一座最不起眼的农舍或棚屋,空不了多久,便被洗劫一空(这是实用性的),外加被痛痛快快地糟蹋(这是为了取乐)。在我方才唤起的这一切个案里,暴力的无端性显出得格外触目:因为被它打击的那个他(或她)是个陌生人,常常是个从未见过、也永远不会再见的人。

这是一种可以称作“\textbf{匿名}”的暴力。毫无疑问,战争从来便是这类暴力的集体狂欢——那是无偿杀人的机会称王的年代,是一个无名小卒的性命,面对扣动扳机的快感、面对体验自己叫一副无名的孱弱身影在眼前瘫倒的权力,分文不值……

倘若世上有一件事——就我记忆所及——每一次都叫我不知所措、哑口无言,那便是再度直面这份超出领会的暴力:那份只为打击与毁灭之乐而打击、毁灭的暴力。倘若世上有一件事在我们身上打下“恶”的那份不可磨灭之感,那不是死亡,也不是身体所能忍受的苦难,而是这件事。而当这样一份暴力——无论它取硬的面孔还是软的面孔,显得“大”还是“小”——由你所钟爱的一个存在冷不丁地朝你来,它注定击中得又狠又深,注定唤起(或重新唤起……)并向你倾覆一份无名的焦虑。这份焦虑的根扎得最深,当它找到童年、乃至婴幼期那片松软而新鲜的土壤扎下根来。这份焦虑——它是我童年生活、也是我成年生活里“世上守得最好的秘密”——是在我六岁那年,在母亲的手里,在我身上显现的。

直到五十一岁,1980年3月间,我才把这份焦虑在我生活中扎下根的那段情节挖出来。焦虑对我的掌控,早已在此前被大体拆除引信——至少在很大的限度内\dmark{}——随着禅观在我生活中的出现(1976年),它在其中逐步占据\origpage{658}越来越大的位置。我与焦虑的关系中的第三个决定性转折,发生在1982年7、8月间,在对日常生活里焦虑之机制的一份细心检视中。自我童年直到成年,那些制造焦虑的情势,一直是这样的情势:它们在我存在的未知深处,使我重新活过“那超出领会的东西”。也恰恰是那些情势:我又一次面对那份表面无法解释、无法把捉、无法还原的暴力的熟悉征象……这份暴力的骤然闯入,使一阵惶乱无措的焦虑之浪骤然复起、倾覆——随即被制住、被压下。这份本能反应,一如其旧,直到今天,出入无几\fnote{203}{(*)}{(12月14日)更确切的说法是:这份反应保持“一如其旧,出入无几”,\textbf{直到}我1982年7、8月禅观的时刻。从那以后,使我猝不及防的“挑衅”虽多,那份所说的“本能反应”却只现过一次,在一年前。它当时成为一场短短的、几个钟头的“应景”禅观的契机,把情势整个澄清。一旦一份内在的混乱情势被朴素地面对、被承担,那份伴随它来向我们传递我们混乱之讯的焦虑,便消失得无影无踪——只留下一份认识、一份焕新的宁静的痕迹。}。然而,这些年里若有什么改变了,那是由于一份\textbf{反思}的出现——它在焦虑的尾流里,使那曾顶着“超出领会的东西”、“谵妄”的吓人假面显形的东西,变得可懂、乃至显然;而尤其,两年以来,由于一份\textbf{对我自己的注视}的出现:一份对这份焦虑本身的关切而温情的注视——一份威严的反射运动,正想叫我把它对我自己也藏起来。或者换一种说法:我与焦虑的关系,尤其两年以来,已变成一份不再彻头彻尾拒斥、不再驯兽师、不再掘墓人的关系,而愈发是一份\textbf{迎纳}的关系:细心的、温情的,迎纳它带给我的、关于我自己的讯息——既关于我的现在,也关于我的过去、及其在我的现在里的行动。依我看,这是我至今跨出的最后一步,朝着一份对他人、越来越全面的内在\textbf{自主}——而首先,是对我的亲近之人、我的朋友\fnote{204}{(**)}{这“最后一步”,在“接纳”(n\textsuperscript{o}~110)的结尾处已谈过,取一份稍不同的光照:一份相对于\textbf{寻求赞许}或\textbf{寻求确认}之需要的解脱——“那份需要,真正构成了那个‘挂钩’,含蓄而坚不可摧,冲突凭它得以在我们身上‘挂上’,而我们由此……处于他人的依附之下……总而言之,凭它,他‘拿住’我们,并且(不动声色地)随心所欲摆布我们……”。(这一段,分明像是今天写的——可我发誓,一个字也没抄!)
我说不清,前面是否还有其他这样的“步”要跨——它们将给我那份回顾,把眼下这份自主看成仍是相对的、而非完整的(可我却倾向于——也许有点天真地——相信它完整……)。
一份对焦虑的松弛而专注的关系之绽放,确确实实代表一份在对他人的关系里的\textbf{解脱}。因为(如下一段所说),正是他人得以“随心所欲地摆布我们身上那份焦虑的闸门”的可能(尤其借着犒赏与拒斥的交替——由他配量、由他施用、凭他的手法),构成他对我们的主要的权力手段。}。

\dmark{}依我看,正是那份不直言其名的暴力、那份“女性”路数的暴力,最\origpage{659}有力地制造焦虑,远甚于那一拳打在脸上的更耀眼的暴力。那玩弄压低嗓门的暴力的人——并由此也在那些秘密的闸门上作文章:它们在别人身上释放无名的、无面孔的焦虑之波——手里握着一件比一份权威、或一份单纯的强制权力更可畏的武器。而且,带着一副天真的神态、随心所欲、凭兴致去摆弄那些焦虑的闸门,代表着一份权力——无疑更锋利、更可畏,即便它始终隐蔽:胜过一切事实的或法理的、由社会共识设立的权力。这便是女人对男人的“正当报复”,在一个男人自称(或曾自称)支配她的社会里;这也是“他”为那份虚幻的至尊(现在的,或过去的)所付的价钱。倘若她是\textbf{奴隶}(在我们的地界,她越来越不是了),那他在她手里便是\textbf{傀儡}、或相差无几(而他今天,一如既往地是)。

几年来,每当我直面一份无端暴力的情势——无论它施于我、还是施于他人,无论它以粗暴的方式、还是阴险的方式显形——便有一份联想以无可回驳的力量来到:\textbf{对己之轻蔑}——或者毋宁说,我\textbf{看见}这份对他自己的轻蔑,在那个公然或暗自以轻蔑待人的人身上。我毫无怀疑:这在我身上,不是\dmark{}一个简单的按钮机制、一个“哲学的”或“心理的”偏好,我\origpage{660}会乐意时不时拿出来遛遛——比如,作为凭一句有说服力的套话驱除我所谈的那份焦虑的手段,把一张万能的标签麻利地贴在一个吓人的陌生人身上。这只是一份\textbf{认识}:一份本质的、深刻的、(一旦看见便)显然的关系。

这份认识不“清除”任何东西;它只是让我得以给一个未知\textbf{定位}。它绝不是一名哨兵,守在那儿拦焦虑的路、或把焦虑逐出阵地。一份认识——取我所理解的意义——的本性不是这样。一份认识属于一份内在的\textbf{宁静},它有助于给宁静以根基。而那种使我们不停地想拦住“闯入者”、生怕它们搅乱一份拼凑出来的“宁静”的东西,是我们身上的一份躁动。我所谈的宁静不怕闯入者:它使我们迎纳闯入者。与焦虑的新相遇在水面制造的躁动,不打扰这份宁静,反而助成它。

\subsection{暴力——或游戏与刺棒}
\label{ch18:local:h53}

\subsubsection*{(a) 正义者的暴力}
\tocsubsub{(a) 正义者的暴力}
\label{ch18:local:h54}

\notehead{141}{}
\phantomsection\label{ch18:local:n49}

(12月13日)上一条注里,拿“奴隶”与“傀儡”开的那个玩笑,我肯定又找到了得罪所有人的办法,而且(倘若有人读的话……)要被所有人骂个遍!除非那位假想的读者(或女读者)满心欢喜地鼓掌——谁知道呢——深信那幅影像发射得漂亮、适用于全世界,唯独不适用于他自己(她自己);或许至多,还除开那位出言讥讽的作者。不过,凭这份假设,他给了我这个人一份根本配不上的信任。我至多敢承认:这几年来(尤其,自1982年7、8月那场关于焦虑的禅观以来),我开始走出、甚至已经走出了那著名的“马戏”——夫妻的马戏,固然;但也走出那些与它如手足般相似的其他马戏。《收获与播种》第一部分里,甚至有一节预告着这个色调,题为“旋转木马收摊了!”(第~n\textsuperscript{o}~41节,今年3月)。那里说的不是夫妻的马戏,而是一场数学的马戏:像所有人一样,我曾在其中转了半生。可同样真确的是:那节有着许诺性标题的文字几周之后——4月29日——出现了一条注“一脚踏进旋转木马”(n\textsuperscript{o}~72),它的名字似乎要唱另一种调门\dmark{}\origpage{661}!与从前的差别,或许在于:倘若此地彼处,我还会在某个旋转木马上转(而依我看,如今还继续吸引我的,唯有那数学的木马……),那么拽着那些让我原地打转的线的,是我自己(至少是我身上的某个人),而不是任何人;而且那些线,对我而言已不再隐形。

这些保留说完,我可以说:我成年生活的最大一部分(更确切地说,直到禅观发现的时刻),我“入彀”得飞快——一如所有人,仍然——无论在夫妻的旋转木马里(它足足欢快地转了二十年!),还是在别的木马里。我并不为此懊悔,因为我对一切这类木马的认识,首先归功于我自己转过的那些。我转了那么久,是因为学生学得慢——也肯定因为,在不止一个方面,我在其中找得到诱饵。到头来,它们的力量与魅力终于消散了——想必如此……

依我看,在所有这些木马里,我始终是“入彀”的那一个,而从来不是“让人入彀”的那一个。或者换一种说法:我相信,对那著名的“天鹅绒的脚掌”风格,我从不曾有过一丝一毫的癖向——我有时会狠狠地亮爪子;可我相信,从来不曾把爪子泡在一份天鹅绒的绒毛里。这是众多征象中的一个:它们证明,在我之结构的层面、“东家”的层面——即我身上受制约的东西——那基调是强“男性”的,这一次毫无歧义。阴的、“女性”的调门,反而支配着“孩子”的层面、我身上原初的东西的层面——也就是说,也在认识的冲动、创造官能的层面。

我还想就我生活中的“无端暴力”再说几句。三天前那条注里,我从被这份暴力瞄准的人、或至少在他人身上直面它的人的角度(哪怕只作为证人)谈到它,当时我写道:

“倘若世上有一件事——就我记忆所及——每一次都叫我不知所措、哑口无言,那便是再度直面这份超出领会的暴力:那份只为打击与毁灭之乐而打击、毁灭的暴力……”

\dmark{}这几行、及其后的几行,与实在相符——至少与我亲身经历的实在相符\origpage{662},肯定也与像我一样直面过这份暴力的无数男人与女人的实在相符。它们可能给人这样的印象:写下它们的人,自己与这份暴力整个无缘,一生免于这类谵妄。然而全然不是。我记得我一生中某些关系——共四段,其中三段落在我的童年或少年(八岁到十六岁之间):那些关系,打有敌意的印记,却没有任何具体的个人怨怼作根基,而以系统的、毫不容情的嘲弄表达自己,或以拳打脚踢及其他暴行表达。第一次,受害者是一个同学(还在德国):全班的受气包。情形拖了好几年,我记得。接着两个个案发生在战争期间:我(从一座法国集中营出来之后)住在瑞士救援会在利尼翁河畔勒尚邦的儿童之家——“盖斯皮”——1942至1944年间。这一次,“恶人”是我的一个同伴(他的父母,一如我的,该是被当作德国犹太人拘留了)和我们的两位监护员之一:两人都同我一样以德语为母语。他们俩,都多少是一群少年男女——有时毫不容情的一群,我就在其中——的受气包;可我相信,我给他们俩的生活,比全伙任何人给他们过的都更难受。同一个屋檐下的共处、难民营地那份朝不保夕的共同处境、盖世太保搜捕犹太人的恒常威胁,本该在我身上唤起团结与尊重之情——却全然没有。

在这三个个案里,我把那份恶意的箭头对准的人,天性温和、相当胆怯、绝不好斗——我于是把他归入“软”或“懦”一类,而这些品性也就被算作使他显得不体面的特征。在一个被暴力之风、被对个人之轻蔑吹得满目疮痍的时代,而我本人满怀对战争暴力、集中营暴力及其一切随行物的憎恶——我却感到:我对他人施加的那份轻蔑与暴力,完全名正言顺,理由单纯:我乐意把他归入“讨厌”一类(以及其他种种形容词……),此后,一切(或几乎一切)都变得被允许了,甚至变得高度可嘉。我这个自诩头脑“合乎逻辑”而公正的人,\dmark{}当时却不曾看到:我的行为、以及它那份由一份反感所作的辩护(那份反感的真正本性,我当然不曾\origpage{663}想过探测),与三十年代那些道貌岸然的德国人——对着“肮脏的犹太人”——的行为,一模一样(那些事,我在童年曾贴身见过);而且正是这些行为,使得那场席卷当时世界、空前绝后的暴力泛滥成为可能。我当然(循着我父母的脚步)装作远离那份暴力,视之为一场怪异的失序(有时甚至“超出领会”)。我满怀倨傲的居高临下,对着所有那些——士兵或平民——以这种那种方式同意在那些英雄的尸坑及其随行的种种可憎之事中充当主动或被动齿轮的人。而与此同时,以我微末的程度、在我自己有限的行动半径里,我做着所有人做的事……

倘若我试着辨明,这份服务于一份轻蔑与暴力之故意的、如此奇怪的盲目,其缘由何在,得到的是这个。我童年自五岁起自己曾承受的那些暴力——它们从不曾被当作暴力指给我注意——终于在我身上造出一份慢性的紧张状态:它始终无意识,又被一份百炼的意志小心翼翼地控制着。这份紧张、这份没有特定靶子的攻击性之积聚,造成一份攻击性放电的需要。可这份“需要”并非身体性的——在所有这些个案里,凭一项相宜的身体活动来泄愤的机会,从不曾缺——它分明是\textbf{心理}的。肯定有一份积聚起来的宿怨——当然主要无意识,且不曾落实为对某个特定的人(比如说,我的父母之一,或代行其职的某人)的可触怨怼:我本可以把宿怨之情转嫁到那人身上,并给它一份具体的、甚至暴力的表达。我身上想必有一份“空置的”暴力、一份弥散的、游荡的暴力,寻找一个可以卸下的靶子。似乎常常是动物(昆虫、蟾蜍、狗或猫、甚至牛或马……)为这类游荡的、寻找牺牲品的暴力付出代价。我的情形不是这样:我不记得我一生曾折磨过任何动物,无论大小。显然,我需要一个离我更近的替罪羊:一个\textbf{人}!而人找替罪羊,肯定从来不难找到。

\dmark{}我毫无怀疑:我方才写的,确实描述了实在的某一侧面。可我感到,这份描述\origpage{664}仍停在事情的表面:它只圈定某一个“机械论”的侧面,而不曾真正更深地进入那份无意识的亲身经历。眼下,在那份经历的位置上,有一大片“空白”、一片虚空。现在不是越过它的时刻与地方,去更深地探测这片“空白”覆盖着什么、什么东西溶解在这片“虚空”里。是不是那份著名的“对己之轻蔑”?——它三天前那条注里还以那样专断的姿态申张自己,而此刻,既然事关\textbf{我自己},它仿佛无影无踪地消失了?现在是时候了,机不可失,该把这件事弄个水落石出,澄清那份顽固而暧昧的“模糊”——它继续打着我对自己认识的印记,一如当年那份包围着焦虑在我生活中的角色与存在本身的“模糊”。焦虑,曾是我一生“守得最好的秘密”——我是这么觉得。会不会还有另一个秘密,守得更好,自我禅观以来,只在此地彼处、在两三次场合里被擦过?我有这样一种感觉:我仿佛已掌握了一切,足以弄清它的究竟——包括那份熟得不能再熟的兴致的骤然涌入:它告诉我,动身的时机已经成熟!可我感到,我不会在这里做这件事——在这场某种意义上“公开”的、或至少注定要拿去发表的禅观里。这场禅观,至少将在许多别的功效之中,有一功:意外地催熟一桩突然变得切近的问题——它终于被承认为理解我自己的关键,而此前,它不过在一张也许永远望不到头的长长的等候名单上,摆着“百问之一”的架势……

并非全无可能:我将还有机会遇到那三个男人中的一位(其中两位与我年纪相仿)——他们当年曾是我身上一份暴力与攻击性的无辜靶子;或者,至少我将有机会给他们中的某一位写信。能在充分明白事情原委的情况下郑重赔罪,这对我将是一件好事。或许,对他也将是一件好事。可奇怪的是:我没有这样的印象——觉得那三人中任何一个曾真正为此怨恨过我,觉得我的暴力曾在他身上唤起一份专对着我个人的敌意;毋宁说,依我看,他被卷入的整个语境,必定被他经验为一场劫难——一场他根本无从逃脱的劫难——而我这个人,更多地被知觉为这场劫难里众多配角中的一个,而不是那个毫不容情的折磨者\dmark{}\origpage{665}(而我是),被仇恨的。当然,我可能错了,而且可能永远无从知道——正如也可能,我将有幸有一天直面这份我在盲目中播下的业……

我相信,在“盖斯皮”情节之后的年月里,我身上有过一份成熟——尽管(就记忆所及)不曾有过任何对此的反思。总之,其后的岁月里,我身上形成了一些有效的反射:它们会拦住我,不再与全伙人对一个成员的集体暴力为伍。我不相信这类事在我成年生活里重演过,也不相信我曾受试探再去扮演那样的角色——我该感觉得到,那个角色多么虚伪,在那副欢快的、“讲体育精神”的外表底下,多么没有勇气。但这并不妨碍:战后,生活照样大量地在我面前堆起那些浸透蒙面暴力与焦虑的情势,在我身上延续那已然打下了我童年与少年印记的深层紧张。正是在这个语境里,才有第四段关系:它打有偶发的敌意与暴力的印记——我可以称之为“无端”的:不由具体的怨怼奠基或挑起,甚至(我相信)也不由任何可称作“挑衅”的行为挑起。那是与我一个儿子的关系。然而我知道:我对他的眷恋并不少,我“爱”他也并不少于我的其他孩子。可在无意识的某个层面上,我身上想必有一份对他这个人某些侧面的拒斥——恰恰是那些使他比他的兄弟姐妹更温和、更易受伤害、也更难把握的侧面。分明地,他全然不“入画”——比我的其他孩子更不——入我那些我本乐意在孩子身上看到实现的漂亮的超级阳影像;而尤其不入画的是:环绕着他人生最初两年、把他打得很深的那些非常严酷的情势,使他更难与父母结成信赖的关系。总之,在他还与我同住一个屋檐下的时节——直到他十岁上下——我有时会对他施以羞辱性的惩罚,以雷鸣般的嗓音颁布。这些东西,整个沉入了遗忘;连同某种终于浸透了家庭空气的氛围——是两三年前与他的姐妹和两个兄弟的几次对话,才恰逢其时地把这些事情多少\dmark{}重新捞回我的记忆\origpage{666}。或许有一天会来到:他也愿意同我谈这些——他,在我的孩子们之中,也许是那份浸透压低嗓门的焦虑与未被承担之紧张的家庭气氛的最大受害者;或至少,是在父亲手里“挨得最多”的一个——尽管他们每个人都分得了一份很足的“父母的包”。至少我知道——而且我为此高兴:拦着我的任何一个孩子与他父亲保持一份简单而信赖的关系、拦着我们一同谈论并探测那份沉重的过去的,不是一份他们还留着、还要设法藏起来的对我的畏惧。

但在这里,同样,在这些注里深究一份复杂的情势——它牵涉六七个人,一如牵涉我自己——不是地方。对我至关紧要的,首先是无遮无掩地确认:这份表面无端的暴力——它曾那么多次“叫我不知所措、哑口无言”,当我在他人身上遇见它时——曾偶尔在我生活中、在我自己的行为里此地彼处地显形。这份确认,不是怀着什么特殊的“意图”作出的;它不自许去“解释”或“原谅”任何人身上的无端暴力,一如那份暴力也不该被用来解释或原谅我的。并非不可能、甚至很可能:随着反思的深入,两份暴力——他人身上的与我身上的——终将相互照亮。这类事情,多半会自己到来,而且不招自来。我之所以作出这份确认,只因为它就在路上;而(除非甘愿不再真)我不能在这里不提它。

\subsubsection*{(b) 机械与自由}
\tocsubsub{(b) 机械与自由}
\label{ch18:local:h55}

\notehead{142}{}
\phantomsection\label{ch18:local:n50}

(12月14日)昨晚的反思,恰好提醒了我一件人们那么容易忘的事——而尤其(在此)是我,\textbf{我}那么容易忘的:我并不比任何人“更好”;我与所有人裁自同一块布;恰恰像我的某位朋友——我正准备把他架上板凳,置于一份毫不留情的注意的中心……

昨天,我给出了一份关于(表面)“无端”暴力之显形的描述:把它作为一份积聚起来\dmark{}的紧张与攻击性向某个替罪羊的放电——那人\origpage{667}因这样那样的缘由,恰好“对得上号”。这份“机械论”而浮面的描述——它肯定“众所周知”——可能助长一份同样“机械论”的\textbf{态度}:对着这份暴力,在自己身上或在他人身上。于是,暴力被看作一份不可规避的宿命,一份可叹地扎根于心灵结构本身的宿命——我们又能拿它怎么办!这样一种姿态,顶着一副“理性”或“科学”的面孔,依我看不过是\textbf{退位}的一份合理化:在自身与他人之中,对着一份创造性的\textbf{自由}的在场的退位——那份自由,向每个人敞开那份抉择:\textbf{承担}我们所处的情势,而不是被动地顺着那些划定的机制之坡道滑下去——它们时刻准备着接管我们。倘若人们难得动用这份“自由”的抉择为真,那么这份抉择的单纯\textbf{在场}、我们身上那些创造的可能之在场——无论我们用还是不用——已经把事情的本性整个改了。正是凭\textbf{这个},且唯凭这个,那些包含人与人关系、或人与自身、人与周围世界关系的情势,才有一份维度——当涉及的不再是人而是(比如说)计算机,无论多么完善,这份维度便缺席。也正是由此,对我们每个人显出那份\textbf{责任}的特权:对我们行为的、对我们行为之动机的责任。这份责任,绝不因我们常常求那份供我们使用的便利——把自己真实的动机对自己藏起来——而被解除。

回到眼下的个案作例解:倘若我曾扮过好人——同时动用我折磨某位不曾伤害过我的同学的权力——那是因为:在一份表面的“真诚”背后,我曾选择了一种粗劣而惊人的自欺态度,当时便与四十年后回顾时同样扎眼。那确确实实是一份\textbf{抉择}:没有任何东西强迫我做;它等于对我身上积聚的紧张与攻击性闭眼(同时当然还援引漂亮“非暴力”的念头),并把它们“悄悄地”(原文如此!)卸到随手可得的替罪羊身上。这样的暴力——也就是说,几乎全部在人类世界里肆虐的暴力与可憎之事——只有在一个\textbf{条件}下才能发生、其秘密的功能才能完成:这个条件本身严格保持秘密(哪怕它\dmark{}扎眼\origpage{668}之至);也就是说,得让自己给自己“把鱼目当珍珠”,凭一份信服玩一份粗劣的双重游戏,为了事态的需要,把我们最起码的认识官能遮起来。诚然,一直以来的空气在怂恿我们:我们始终看到我们的周围人赶着以他们的共识,批准那些为有他们附和的虚构效力的遁词——无论多么粗劣。而我自己的遁词,在我谈过的那些个案里,确确实实有着周围人的附和或默默的怂恿;不然,我无从把它维持下去、继续我的游戏。

而承担一份情势,不多不少,正是以\textbf{真诚}直面它——取这个词的完整意义;也就是说:不动用供我们使用的便利,凭粗劣的遁词,对它显而易见的来龙去脉闭眼。这也就是说,干脆地:动用我们健全的知觉与判断官能,而不为这样那样事态的需要留神把它们遮起来。事情或许显得奇怪,却同样简单而显然:当我们以这样的姿态——“无邪”的姿态——直面一份情势,它旋即深刻地转变,无论它曾显得多么混乱、多么打了死结。或者更好地说:倘若它果真“打了结”、许久纹丝不动,那是因为我们自己拦着它演化、拦着它按其本性“流淌”;是我们挡着它的自发运动,循着……一如我们自襁褓以来所有围绕我们的人的一致先例。只消\textbf{停止}绷紧,\textbf{停止}阻拦,那些仿佛凝定了的事物便重新动起来;卡住的东西松开了;积聚起来的生硬的紧张,终于得以释放,消解于一份新的、宽广的运动——它终于重现。

这份我们人人都有、又得众人怂恿的“便利”:把鱼目当珍珠,并由此堵住那本为流淌而造的东西——它其实毫不“舒服”!它为我们保有的那份怡然自得的内在不作为,我们付的是一笔高得吓人的价钱——一份内在的僵紧之价,加上为维持这份僵紧与“鱼目=珍珠”的虚构所需要的惊人能量投入。话虽如此,每一时刻,人人按自己的心意行事——这是我们的特权。而在每一\dmark{}时刻,凭我们所\origpage{669}做的,我们播种——为着我们自己,也为着他人。而我们之所\textbf{播}的收成,当下便开始。

\subsubsection*{(c) 贪婪——或蚀本买卖}
\tocsubsub{(c) 贪婪——或蚀本买卖}
\label{ch18:local:h56}

\notehead{143}{}
\phantomsection\label{ch18:local:n51}

或许该是时候回到葬礼的那个“第一层面”了——也就是回到我的大司仪、我的朋友皮埃尔在葬礼中所扮角色的来龙去脉。一周之前,我在“天鹅绒的脚掌——或微笑”那条注里(n\textsuperscript{o}~137,12月7日)已经回到过它,随后又经由那场(连着五条注的)关于“利爪”与“天鹅绒”的枝蔓离开了它。我感到,这场“枝蔓”——一如先前的许多场——不曾白费。

我之所以被引上这场枝蔓,恰恰是因为:我的朋友在扮演那份角色时,也许最惹眼的明显特征,便是那份最纯的“天鹅绒的脚掌”风格的执拗——没有一丝破裂的意向,为一份毫无破绽、从不直言其名的对抗效力\fnote{205}{(*)}{一如我别处已有机会强调的:对抗、或拒斥与嘲弄的故意“从不直言其名”,绝非朋友皮埃尔所专有;而(就我所知)适用于葬礼的全体与会者,无一例外。于是,在这场以嘲弄行事的“阴的葬礼”里,每一位与会者身上的底色(一如这样一场丧仪所应有)本身——是阴的!
关于葬礼这份“隐蔽”的品格,亦见“掘墓人——或全体教众”一注(n\textsuperscript{o}~97)。}。另一个惹眼的事实:在那副听得懂的微笑、和婉的神色之愉快而克制的表象背后,我的朋友身上,对着我本人、或对着他归入“我的”名下的人(在数学工作的层面上),不止一次表达出一份毫无歧义的、表面上无端的\textbf{害人}或\textbf{伤人}的意图。关于这一方向的诸多具体事实,我在葬礼第一部分里已说得够多,此处无须赘言。这确确实实是一份恶意的姿态(似乎严格圈定在科学活动的领域之内),是取一词之强义的“\textbf{暴力}”——即便它严格保持隐蔽:爪子永远泡在精美的、起绒的丝绸里。而这份暴力、这份恶意,带着那份最令人错愕的\textbf{无端性}的全部表象——仿佛它们施展出来,只为害人伤人之乐。

\dmark{}一如每次直面这样一份情势,它显得如此不可思议,以致人们常常\origpage{670}不敢相信自己健全官能的见证\fnote{206}{(*)}{关于这一点,见“中国皇帝的新衣”一注(n\textsuperscript{o}~77$'$)。}。否定这份见证——时下的惯常做法——是不承担一份情势的无数方式之一,并由此把它延续下去;更可取的,肯定是围着事情转一转,寻找或许逃过我们知觉的侧面,它们能提供一条接近的路径,使这份情势得以纳入我们的亲身经历。依我看,一生中从不曾经过这份无缘的恶意姿态的人,必定寥寥——而肯记得自己也曾如此,已是可能的第一步,好\textbf{靠近}那份情势的实况:时下的反射,更巴不得我们把它火速清场。再往深处探测一番也无妨:看看是否有什么暗藏的怨怼,是那份看似无端的暴力之缘由与发条——同样,视情形,也不妨认出某些“怨怼”的赝品,照它们的本相:比如我自己实行过的那一套,说某人是个可憎的家伙、不配任何客气等等。

可在眼下这个个案里,任凭我怎么探测,都看不出任何东西——近看远看——像一份\textbf{怨怼}:我的朋友可能(有理或无理地)怀有的、对着我的怨怼,或对着他择为恶意靶子的任何人之一。他自己,也从不曾有过片刻,暗示过哪怕一星半点朝这个方向的东西;更不必说,对我屡次就他那几桩叫我目瞪口呆的行为所作的探测,他从不曾有一刻承认:他身上可能怀着对任何人的、一丝敌意的姿态。我最终在他身上感到一份暗藏的得意——在我们偶尔见面时:当他端出他那套十足“客观”的好理由,带着那份他特有的、有点好玩的无邪惊讶的神色……总之,我进了一场他随心所欲、凭兴之所至引领的游戏,带着一份我久久未曾察觉的内心满足。(然而,把人折腾得团团转的,他可远不是第一个!)但我终究——迟走出来总比一直陷在里面好——走出了那个木马\fnote{207}{(**)}{那是1981年——即“两个转折”那条注(n\textsuperscript{o}~66)里所谈的“第二个转折”。}!

另一方面,倘若我探测我自己,检阅我与我这位\dmark{}朋友自将近二十年前(1965年)相遇以来的关系\origpage{671},我同样找不到任何痕迹:任何曾在某一时刻成为对我的怨怼之缘由的东西。按事物的常规的、浮面的意义,我可以说:在整个这段时间里,尤其在最初五年贴身接触里,我“只对他好”。可这份确认旋即使我记起另一份,不那么浮面的——我对他的那份\textbf{纵容}:它显形于“置身事外者”与“含混”两条注(n\textsuperscript{o}~67$'$、63$''$)的反思进程里。分明地,这份纵容对他绝不是“一件好事”;同样分明的是:我那年轻而辉煌的学生兼朋友对我的种种姿态,与我自己的种种姿态——尤其与这份纵容——曾在一个紧密的共生中发育。甚至可能:这份纵容,在无意识的某个层面上,曾被我的朋友(不仅被知觉——这反正显然——而且)感受为一份“怨怼”、一份也许太熟稔、被嚼烂了的老套脚本——在他那略带神童色彩的童年时代;而如今,它(哪怕不显眼地)又被端到他面前。他也许曾天真地以为:一脚踏进数学的“上流世界”,一切都将不同于他所认识的——结果不然,还是那套老把戏!(而凭他自己的刻意抉择,今天仍是那套老把戏,而且规模更大,更不必说了……)个中究竟,我恐怕永远无从知道。况且,把它澄清,也不是我的本分——假定我的天线细到足以单凭自己的手段做到。倘若有过“怨怼”,那至多是一份“贴补”的怨怼,它弹了一记指头,把“某样东西”推上路:某一场\textbf{游戏},由一份完全另一个量级的力所驱动;那份力的在场,我久已感到,其本性对我却仍是个谜。在离开葬礼图景的这个“第一层面”之前,我至少想试着掂量那份力的本性。

分明地,有一份\textbf{贪心}:取而代之、排挤、抹除的贪心;还有那份把他人与数学夫人辛劳与爱恋之果\textbf{据为己有}的贪心。可是,我清清楚楚:\textbf{不是}一份单纯的声名、钦慕、荣誉的“贪食症”,甚至不是权力的“贪食症”,才是他在葬礼中所扮角色的深层发条。在我对这份\dmark{}角色的反思进程中,我曾多少次骤然惊觉\origpage{672}:他那份下葬的\textbf{执念},使他把自己也埋了进去!凭他那份例外的天资、凭一份同样例外的时运,他分得了超越师父极远、在整个时代的数学上留下深深印记所需要的一切。只消他放任身上的孩子随意去玩,不用种种指令去烦它、不用这里一道栅栏那里一处禁行去烦它——只消简简单单照看好必要的事,严格地管好后勤——那么,他身上的“东家”(其贪心无疑与任何人半斤八两),既不必推,也不必拉,也不必动用手肘,本来就断然少不了各色可想而知的声名、钦慕、荣誉的标记,外加权力的标记,多到不知道拿它们怎么办——彼时,那位“郁郁寡欢者”却正玩得尽兴,也绝不给“东家”留下多少摆东家架子的闲暇……

分明地,按单纯的“功利”账算,陷进一场葬礼,是一桩很不划算的买卖:它粘在他腿上已十五年以上,若非那碍事的亡人忽然想起掀开棺盖、搅了仪式(恰在其时——正如常理所言——在人们最不提防的时刻!),它本来打算粘他一辈子!(关于那桩不合时宜的事件,对皮埃尔东家未来下注的影响,赌局已开……)或者换一种说法:我的朋友有那份材料(至少凭他的智性手段)、有那份显赫的资格,在数学里做一位\textbf{彼得大帝};可他偏偏选择了做一个小皮埃尔。这看上去千真万确是一桩蚀本买卖——至少,倘若所下的注,确确实实首先是那些虚荣的满足。

\subsubsection*{(d) 两种认识——或对认识的恐惧}
\tocsubsub{(d) 两种认识——或对认识的恐惧}
\label{ch18:local:h57}

\notehead{144}{}
\phantomsection\label{ch18:local:n52}

(12月15日)昨晚反思的收尾处,我身上有过一份轻微的不适,像那种人气十足地端出一份逻辑无懈可击的推理、却撇开一份弥散的感觉的人:总觉得有什么地方不对劲。这份“什么”在我一停笔时便显形了。表述它的一种含混方式是:无意识的“逻辑”——它主宰着我们最关键的抉择——绝不是那份寻常\dmark{}的有意识推理的逻辑,更不是那份“正统”推理的逻辑\origpage{673}。在此个案中,我对青年Deligne在(姑且说)六十年代后半期所握“王牌”的知觉、以及我赋予它们的分量(至少与任何一位消息灵通的理性数学家将赋予的分量同向)——这份知觉与这份分量(我乐意称之为“客观的”)与当事人自己的姿态与感情毫无关系;尤其,与那些关乎他自己能力的姿态与感情无关——而那些,分明是他手中所有王牌里的王牌。

可是我有这样一种印象:至少在意识的层面上、带着谦虚所要求的种种套语,我的朋友早已把那些久已传回他耳中的动听回声——关于他那不寻常的天资——纳为己有。但对我而言毫无怀疑:在更深的层面上——那个一生的主导性抉择被默默做出、无话可说的层面——那份“客观”的事物版本,当时便已(而今天依然)\textbf{失效}。占了它的位置的,是一份阴险的\textbf{怀疑}:任何一份价值的“证明”(或对他人之优越的证明……)都永远拔不了根——一份愈发顽韧的怀疑,因为它永远未被表述。我在朋友身上知觉到它,一如我在其他天资逊色者身上知觉到它:它是同一份。这份怀疑,是一份\textbf{贴身信念}的顽固信使:那份信念同样未被表达,埋得比这怀疑更深:一份对无能的贴身信念——根本的、无可救药的。它也就是那份“对己之轻蔑”——我在《收获与播种》的开头处、在一场仍取“一般”品格的反思的语境里谈到过它\fnote{208}{(*)}{见“他人之无谬与对己之轻蔑”一节(第~n\textsuperscript{o}~4节)。}。一两个月前,它又在一个仍然非个人的语境里、以一副不同的面孔重现:作为一份“裂缝感”\fnote{209}{(**)}{见“一半与全——或裂缝”那条注(n\textsuperscript{o}~112),10月17日。}——那份弥散的感觉:我头一次在我自己身上确认它,是在发现禅观后的第二天。而在葬礼反思的进程中,也有几次,对我朋友身上那份“无能的贴身信念”的骤然而敏锐的知觉,为某一仿佛公然违逆常识的情势投下新的光照……\fnote{210}{(***)}{关于这一点,见“颠倒(3)——或阴葬阳”那条注,其中(及其他)唤起了反思的某些这类“敏感时刻”。}。

\dmark{}我知道,这份贴身信念——在我朋友身上或任何人身上——本身宛如一份\textbf{认识的}\textbf{影子}:对那份确确实实存在的“裂缝”的认识、对一份被承受的“戕害”的认识——而这份戕害,直至今日,仍由当事人自己的默许所批准、所维持。然而,影子并不归还它所由出的那份认识——那份认识,一如一切认识,本身是有益的——它毋宁是它的一个巨大而畸形的漫画、一个稻草人版本。如此歪曲一份认识、使它面目全非的,是一份\textbf{恐惧}:恰恰是怕与这份认识本身接触的恐惧,怕让它从它一直被压抑的深处上浮的恐惧,怕承担它所忠实映照的那份谦卑实在的恐惧。

与这份可畏的认识接触、以一份完全有意识的目光认识那份在深层已知、却被逃避的实在——这才真正意味着:与我们身上那份(人们叫它“力量”也好、“孩子”也好)“自信已亡故一生之久”的东西,全面重新接上联系。因为确确实实是那份力量、唯有那份力量——童年的力量——使我们能够承担对自身那份打了裂缝、被残害、被瘫痪之物的认识。而承担它,也意味着:与那\textbf{另一份认识}重新接上——它先于对我们戕害的认识、且比它更要害:对我们身上那份力量之在场的原初认识——一份既非肌肉之力、亦非头脑之力,却把两者都包含在内的力量。

事情或许显得奇怪:这份失落的认识——关于那份“力量”、那份\textbf{创造权能}在我们身上的在场,作为我们真正本性的显而易见、不可摧毁的一部分——恰恰经由一份\textbf{无能状态}的发现与谦卑接纳而寻回;而那份无能状态,正是凭这份接纳本身得到解决。对一份无能状态的认识,覆盖并遮蔽着那份埋得更深的认识:对我们创造力量的认识。前者宛如一把钥匙,为我们打开后者;两者实际上不可分,如同一份\textbf{同一}认识的正面与反面\fnote{211}{(*)}{在这幅意象里,当然,“正面”是对无能状态的认识——即不真实的状态、“裂缝”的状态;而“反面”,藏得更深,是对我们未分裂的本性、我们的创造权能的认识。这些年里,我屡次得以确认:恰恰是“反面”——两者中埋得最深的那份认识——才是最强恐惧的对象、才招致最激烈的否认。让任何人不安的,并不是那份驯熟的、(或多或少)“有学问”的猴子的熟稔而平淡的状态,而是孩子的那份无邪:他按事物的本相感知它们、直呼其名,他照他所感去做去说,不因与“人家”期待他的样子不同而羞愧。},是同一份\textbf{恐惧}的对象。

\dmark{}当我谈到埋在我们每个人身上的那份“力量”时,这绝不是什么抽象而含混的东西\origpage{675}、什么“哲学家”的口头精致、或半吊子哲学家的心理学的口头精致。正是这份力量,使你能“做数学”(或“做爱”……),像一个孩子呼吸——也就是说,不必小心翼翼地约束自己不偏离前人留下的尾流、不必勤勤恳恳重复他们的手法与配方(或他们的陈词滥调……);也正是这份力量,给你勇气与谦逊——在你自己家里,也在别人家里——把猫叫作猫,不把鱼目当珍珠,哪怕这样做你会顶着最牢固的共识、或顶着你自己身上最积习已深、最运转纯熟的那些机制。\fnote{212}{(*)}{(12月16日)创造力量在每个人身上的作用——更新之力(或“孩子的力量”)——凭它的果实而辨认:无论凭手或心的作品,还是凭日常生活的种种事实,在与他人、与周围生灵与事物的关系中。我得以屡次确认:日常之中的创造性,远比经由“作品”(取常规之义——即凭手或凭心塑成的可触“产品”的创造性)的创造性罕见。
一份持续的创造性在某个人生活中的在场,是与那份在他身上的创造力量的一份持续“接触”的标志——哪怕何等零碎、何等不完美。这属于另一种本性,不同于“天资”的单纯在场、不同于为开发利用它们而作的持续能量投入——后者表达为一份或多或少的可观、或多或小时行的产品,却本身没有创造的德能、更新的德能。
在我的智性求索之中、尤其在我的数学工作里——天资平平(而投入可观)——我觉得,我与那份力量之间的这份“接触”、亦即我对它那份默默而深刻的认识,几乎原封未动。也就是说:出入无几,我在我生活的那个(诚然极零碎的)领域里,凭我(创造)手段的全部“运转”,几乎没有任何损耗、挪用或阻塞——没有惯常的那些“摩擦效应”。其中最常见的一种,是某种怯懦:它如此经常地使我们听不见内在的声音——那声音正叮嘱我们该做什么,当它所教的恰恰是“新”的东西,即把我们领上唯有我们自己踏足的小径。这类抑制,在我的数学关系里几乎缺席(而且依我看,随岁月推移愈发如此);可在我生活的其他领域,它一如在任何人身上,存在着——尤其恰恰在“日常生活”的领域。
回到数学活动:我在我那辉煌的昔日学生身上看到一份某种意义上颠倒的关系。他握有那份从来令我惊叹倾倒的天资——与我的不可同日而语。(诚然,我活得越久,越看清:要在科学或其他领域做出革新之作,这根本不是要害;关于这一点,见“阴,作为仆人(2)——与慷慨”那条注(n\textsuperscript{o}~136)里的反思。)他对数学的投入是可观的,一如我当年的投入;而自他最年幼起,他便享有格外有利的条件:为他的天资之绽放、为构想并打造一部与那些天资相称的作品。二十年之后,我还在等这部作品,而且意犹未尽!在他身上,肯定有某份与创造力量的“接触”——他所做的某些东西之美可为佐证;可这份接触是受扰的、受磨的。我的朋友对他的工作的关系、乃至在他的工作本身之中,是一份冲突的关系:工作随着岁月推移,愈发变成“东家”手里的一个\textbf{工具},用来满足\textbf{他的}饥渴——那些与孩子的认识渴、发现渴格格不入的饥渴。
我怀疑,这样一份冲突的关系,能不经先行承担而得到解决——也就是说,在一切之先:被认出。至少,我一生不曾一次见过这样的事不经此而完成。这正是使我写出下面这句话的缘由:对我们无能的认识,是寻回我们创造权能的全面认识的“钥匙”,并由此,也全面地寻回创造权能本身。在我的数学工作里,问题不曾提出,因为那份工作里不曾有深层的阻塞——一种相当于局部无能、会使我只凭我可能性的小部分“运转”的阻塞。可这个问题,一如对任何人,在日常生活的层面上、在我与他人的关系、与我自己的关系、与我身体及其冲动的层面上,提出过。正是在这个层面上,我一再经验到:对一份阻塞、一份“无能”的认识,确实是那把解放被囚的创造性的\textbf{钥匙}。}

顺带,头一个来到我笔下的例子,颇能说明问题——它肯定足以叫每一个(乃至不那么年轻的)热爱荣耀的年轻研究者心跳。谁不愿意做一门尚在孕育的科学的无畏先驱,并以此在所有教科书里占据显要地位,像开普勒——现代天文学之父!可当要做的(一如开普勒与其他人所做的),是在孤独中、在所有人的冷漠中(更不必说轻蔑或敌意)顽韧地纺自己的线,三十年、哪怕仅仅一年——那时,忽然就没了人影!人们愿意进教科书,总而言之愿意与名人并列;可人们也\textbf{怕}孤单——哪怕一年,甚至哪怕\dmark{}仅仅一天。而那个“认识”到他身上力量之在场的人(为认识它,他从不曾\origpage{677}必须谈起它——对别人、甚至对自己……)——那个人也分明知道:他是\textbf{孤单}的;而孤单,并不叫他不安。至于他会不会进教科书,那是他操心清单上的最后一项——尤其在他工作的时刻。

况且,恰好是这位开普勒,在他的工作本身之中,“顶着科学里最牢固的共识”——而且是竖立了数千年的共识,更不必说了。在他的时代(宗教裁判所还在),这比今天更不便:今天你很可能丢掉饭碗、或找不到饭碗,却不至于冒险上火刑柱。回到开普勒:我不知道在他的日常生活中,对着那些“最牢固的共识”,他如何自处;或许在那些事上,他跟所有人一样规规矩矩。确定的是:今天一如当年、一如从来,肯偏离这些共识一毫的人,也不多。想必还是那套老把戏:\textbf{怕孤单}——一份深层而几乎普世的需要之反面:对赞许的需要、对他人之确认的需要(哪怕世上只有一个肯赞许、肯确认的人)……\fnote{213}{(*)}{我由此经由另一条路径,接上一些已在“禁果”与“孤独的历险”两节(n\textsuperscript{o}~46、47)里、以及顺带在“接纳”那条注(n\textsuperscript{o}~110)里显现过的确认。}

\subsubsection*{(e) 秘密的神经}
\tocsubsub{(e) 秘密的神经}
\label{ch18:local:h58}

\notehead{145}{}
\phantomsection\label{ch18:local:n53}

可我又一次离开了我的话题!我起笔时,是针对这份确认:我昨晚的“推理”驴唇不对马嘴——当时我试图“传递”我身上的这份确信:我的朋友在葬礼中扮演我所知道的那份角色、以我所知的那种方式,其动机不是\textbf{贪婪}(对声名、钦慕、荣誉、权力的贪婪)。诚然,千真万确:拿一份孩子的冲劲换一份\textbf{角色},他做了一桩“蚀本买卖”——甚至按“回报”账、按声名那一侧的账,也是如此。可这什么也证明不了。依我看,这类“亏本账”,在我们主要\dmark{}投入与抉择的(无意识层面的)拣选里,几乎是绝对的通则,而绝不是例外。可即便那推理一文不值,我却毫不怀疑\origpage{678}:我想传递的,恰恰是一份对实在的知觉:\textbf{不是}这份贪心——它千真万确,而且在我朋友的生活里占据了越来越大的、真正吞噬性的份额——\textbf{不是它}构成我的朋友所扮角色的\textbf{神经}:作为我葬礼之操办中那个\textbf{关键}人物。

倘若我试着更贴切地圈定这份极清晰的感觉(至于“确立”它的正当性,则再无所提!),得到的是这个:是那份对抗或恶毒行为中的\textbf{无端性}——那份屡屡叫我目瞪口呆的无端性——与那份万金油的“解释”绝对不“对榫”:贪婪。至少,就声名、钦慕、荣誉而言,甚至就取常义的“权力”而言,我那辉煌的昔日学生兼朋友,在对付他昔日师父时使出那份他深得其妙的“含蓄而配量精巧的轻蔑”——无论在当下还是在更长的期限里,什么也没赢得;或在对付某位地位低于他的研究者、或其现在或过去的工作时,使出同一份轻蔑(或许配量不那么精巧),把那人本来便不如我牢固的自信吓散;或者对付另一个人——那人顶着朋友定调的普遍轻蔑英勇坚持,他便剥夺其坚持之果。倘若说在这最后一个个案里、一如在其他个案里,我的朋友曾作态把别人在孤独中(有时顶着长辈的轻蔑)成熟的果实据为己有,那么这份“收益”(取“\textbf{暂停}”的风格\fnote{214}{(*)}{关于我这位辉煌的朋友兼昔日学生这份据有风格,见“暂停!”(n\textsuperscript{o}~77)与“占有与轻蔑”(n\textsuperscript{o}~59$'$)两条注。}),一想到\textbf{据有者是谁},便寒酸到这种地步,以致那份端出的“解释”自己便化成青烟!

至于我,我分明知道——凭一目了然的认识:这类据有的“神经”,不是\textbf{那份}收益。我从中反而感到\textbf{某份权力的陶醉}——一份比常规意义的权力更精微、或许更醉人的权力:如同某位有分量的科学家惯常行使的那种——坐在各种委员会、理事会、评审团及同类机构里\dmark{},主管一所研究所\origpage{679},或指导一批年轻而辉煌的研究者,或对部长耳语。我所谈的那份“陶醉”,在“乖僻”那条注(n\textsuperscript{o}~76)里显现过(反思中头一次):我在那里骤然面对“一份\textbf{炫弄}之举,一份在如此全面的权力中的陶醉——它甚至容许自己(象征性地……)炫示它那‘乖僻地’剥夺他人的真正本性”。

那是一份炫弄之举:辉煌而张扬,却同时\textbf{隐蔽}、未曾言明,不动声色地滑进那里,甚至还为“反常层”这个奇怪的名称附上一段临时凑出的解释——多么自然啊,三言两语就能向你说明白——此外再列一张小清单,说说“哪些东西本该写进我们那篇谦逊而辉煌的文章”……\fnote{215}{(*)}{见“变戏法的人”那条注(n\textsuperscript{o}~75$'$)。}。

我在这里又认出了那份最纯的“天鹅绒的脚掌”风格——别名“暂停!”的风格;而且,在那份我已在不止一人身上变得熟稔的风格之齐一背后,我也感到那份\textbf{共同的神经}:那份强制性的、吞噬一切的\textbf{渴}:行使一份权力;一份\textbf{特定的权力},取一种特定的方式——猫玩老鼠的权力:当猫施展它的\textbf{大游戏},带着那份完美的优雅(唯独老鼠无从领教其妙处),并且肯定“以最精微的温柔”——或者,也是一位精明的妻子对她的傻大个丈夫的权力……

从我的朋友提出的个案出发,我已被引着谈过那份“风格”及其意义——在形形色色伴侣关系的一般语境里。那是在一周前的反思里,“颠倒(4)——或夫妻的马戏”那条注(n\textsuperscript{o}~138,12月8日)里。正是在那里,“天鹅绒的脚掌”(别名“暂停!”)游戏的“神经”头一次以它应有的清晰显形:作为一份\textbf{权力的游戏}。然而,作为一份性质十分特殊的权力游戏:游戏对行它之人的那份魔力、那份往往吞噬一切的魅力,恰恰在于经由它行使的\textbf{权力的隐蔽品格}——那份“看不见、摸不着”的品格:它使人得以玩弄对方(\textbf{玩弄}他,却从不\textbf{与他一起玩}……),随心所欲地叫他原地打转,始终领舞,而对方笨拙地一次接一次跟着转,对那些凭无形的线使出的、随兴所致的小小招式,报以笨拙的应答\dmark{}\origpage{680}……

终于白纸黑字写下那份多年来晦暗地感到的东西——我从不曾费心对自己明白表述——短短一份努力,把长久弥散的东西凝成言语,便已足够:昨天还向我显为“谜”的东西(即某位朋友身上一份“某种力量”的本性),骤然向我敞开其显然的意义!他身上这份“力量”、或(如我方才所写)这类看似“无法解释”(乃至“超出领会”)的行为的“神经”,我在12月8日那场反思里早已圈定。可尽管那场关键反思的出发点,确实是我那位辉煌朋友的某场“谜”样的游戏,滋养那场反思的,却是\textbf{另一份}亲身经历——比我与他这个人相联系的亲身经历更丰富、更强烈;一份已整个消化(或相差无几)的经历,它向我口授一份已然成形的认识——那份我与朋友皮埃尔零星关系的皮肉经历,当时本无力传给我。

当然,终究要理解的,正是那份经历,并由此把它全面承担;倘若我当时毫无内心保留,便投入那场关于“夫妻旋转木马”的枝蔓,那是因为我分明感到:那个木马,对我的朋友关系,有话要说。对他的念头,继续在背景里在场,像一份低调的底色。

然而,两者的全面“会合”,那天没有完成,随后几天也没有。无疑,时机还没有全然成熟。要使会合毫无保留、毫不费力、带着显然的从容完成,我须先“清理场地”:执拗而不慌不忙地,把最强制性地要求我注意的联想,一一追随过去。我不曾强迫事情;而且我知道,这恰恰是我该做的:照管那份执意呼唤我的东西,而不被一个“话题”、一根(反思的)“线”、乃至一份要收尾的纲领引开。

我这边锄草松土,天地之力在做它们的工。到了晚上,只须来收取熟透的果实:它落进那只张开迎候的手……

\subsubsection*{(f) 激情与饥渴——或逐步加码}
\tocsubsub{(f) 激情与饥渴——或逐步加码}
\label{ch18:local:h59}

\notehead{146}{}
\phantomsection\label{ch18:local:n54}

\dmark{}\origpage{681}(12月17日)依我看,随着前天的反思,发生了一份仿佛解冻的事情:一份一直犹豫不决、有点发懵的理解——面对堆积在我面前的一大堆事实与直觉,乱得像一座不成形的山——像一幅拼图,我不过勉强此地彼地拼上了几块。此刻我有这样的印象:我撞上了那幅有待复原的未知图景的\textbf{那}块“要害”拼块——围绕它,其余的拼块终于将毫不费力地各就各位。无论如何,我毫无怀疑:我确确实实触到了朋友皮埃尔在师父及其(或多或少)信徒们之下葬中所扮角色背后的“神经”,也同时触到了他与我这亡故师父之关系的“神经”。

这份玩弄某份权力的饥渴——不显眼地、带着一副天真无邪的神气,拉动那些无形的线——这份饥渴,必定早在我认识他之前很久便已存在,他本人与所有人都不知情。在我与他相识的最初几年——我出发的情节(1970年)之前——我不曾见它显形;这无疑是因为:在那些强烈学习、精微而有力的思想起飞的年月里,我朋友的能量被整个吸在别处。条件确实是理想的:足以作他那些例外手段的跳板。我出发的情节——先脱离我们共同隶属的那家机构,继而(次年内)脱离数学舞台——不仅是我自己灵性历险的关键转折,肯定也是他的转折。正是这份情节,骤然为他敞开一些权力的手段:前一天他还不敢梦想。首先是把一位前师父从那地方“排挤”出去的权力——那师父在那里占的位置太大,而他此前只求不显眼地与之保持距离\fnote{216}{(*)}{关于这份先求保持距离、继而排挤的心思,见“驱逐”(n\textsuperscript{o}~63)与“兄弟与配偶——或双重签名”(n\textsuperscript{o}~134)两条注,及后者的子注(n\textsuperscript{o}~134$_{1}$),以及“未竟的收成”一节(第~n\textsuperscript{o}~28节)。};继而,当那位师父显然要退出舞台时,是更醉人的权力:叫那所以亡故师父命名的“学派”无影无踪地消失;而这样做,最终是干脆利落地——在每一根主枝上(他自己栖身的那根除外)——砍断一份宏大纲领的绽放:它服务于一份宏大的愿景,而他本人曾长期以之为滋养\fnote{217}{(**)}{甚而,关于一个“学派”的清算与“电锯效应”,见“继承人”“共同继承人……”“……与电锯”(n\textsuperscript{o}~90、91、92)三条注,以及“灵车”行列的头四条棺材注(棺1至4,n\textsuperscript{o}~93--96)。关于那份被下葬的愿景,见“我的孤儿”(n\textsuperscript{o}~46)两条注里(以两道不同光照)给出的两份概观,以及“阴,作为仆人(2)——与慷慨”那条注的子注~n\textsuperscript{o}~136$_{1}$。
人们会注意到:正文里“而这样做……”一语(“……干脆利落地砍断一份宏大纲领的绽放……”)并不贴切。清算一个学派,是“干脆利落”地“锯断”一整批主枝的\textbf{第一}记“电锯”,但不是最后一记(所引的那些棺材注,n\textsuperscript{o}~93--96,即为明证)。}。

我朋友生活中这场大转折的意义,向我显为他这个人身上两股支配性力量的相互霸权之中的一场颠倒——依我看,它们凌驾于其余一切之上:数学的激情,与权力游戏的“饥渴”(“以天鹅绒的脚掌”)。前一份力量,在本质上取“冲动”的品格\fnote{218}{(*)}{数学的激情“取冲动的品格”、是“孩子”(别名“匠人”)的表达,这并不妨碍(同一段里有力地提醒)它也或多或少强有力地被“东家”的“饥渴”所占据——这属于“匠人”与“东家”关系的共同命运(我本人也毫不例外,一如任何人)。},后一份取利己的、习得的品格。转折之前,是认识的冲动支配着我朋友的生活(就我所知),而权力的饥渴或多或少沉睡,处于空置状态。在短短几年间一段令人眩晕的社会攀升之后\fnote{219}{(**)}{关于这一点,见“高升”那条注(n\textsuperscript{o}~63$'$)。},在一个骤然出现、逼人作严苛抉择的情势里,取胜的(我相信是断然取胜,且毫无交战之意)是权力及其秘密陶醉的诱惑,压过了认识的激情。后者并未退出舞台,却从此成了饥渴的封臣与谦卑的仆婢,成了后者手中的一个\textbf{工具}。激情(别名“匠人”)在饥渴——别名“东家”——的醋眼里操持它的活计;后者寸步不离。既然匠人有好工具(并非全被没收)、有一双好手,即便这样被辖制,他仍旧磕磕绊绊地维持着一份产量、维持着这家的名声。但已不复是当年:当年匠人(颇带几分顽童气)整日整日自得其乐,而东家远在天边,一季才来转一遭!

\dmark{}其后的演进,依我看,更多是数量性的,而非质量性的。那是东家某一份\textbf{战术}\origpage{683}的渐进演进——风格始终如一,而东家-匠人的关系纹丝不动。这位东家禀性谨慎,只在有赢的把握处才肯冒险。为此,须对地形有把握——或者,对“(几乎)全体大会”的默认附和有把握,从亡故师父那些昔日学生的更小圈子算起。顶风冒雨维持的那份个人关系之演进,是“地形知识”之演进的忠实映像。权力与轻蔑之游戏的胆量,有一份渐进的\textbf{加码},历时十二年,在1981年以反常层会议的壮举达于顶点:在普遍的亢奋中,一切克制(甚至一切谨慎)都欢天喜地地被抛进水里\fnote{220}{(*)}{关于“反常层会议”,见第七行列“会议——或Mebkhout的层与乖僻”,注~n\textsuperscript{o}~75--80。}。这样,竟须十二年,我的朋友才使自己确信:地形已如此相宜,任何谨慎都属多余:招招皆赢!时机分明成熟了:终于可以把那件秘密武器公之于众——\textbf{Motif}:次年便换了个爹被开棺出土\fnote{221}{(**)}{关于Motif的开棺,见“一个梦的回忆——或Motif的诞生”与“下葬——或新父亲”两条注(n\textsuperscript{o}~51、52)。}。

我没有兴致在此追溯这份十二年加码的逐步脚印——尽管我手段样样在手,完全可以做。那会是一份编年史家的工作——那类工作,我在葬礼第一部分那场未曾预料的“侦查”(或“中国皇帝的新衣”)里已经做得够多。这场加码的诸“步”,依我看,是“\textbf{试探}”——我的朋友朝一位沉默的大会掷出的一记记探针——每一次都得到同一回答:尽管放手干!将近十五年,大会是他的无声盟友与他的保人;而他,在不知不觉中、大概也不曾操心,是大会一件驯服的工具\fnote{222}{(***)}{见“掘墓人——或全体教众”一注(n\textsuperscript{o}~97)。}。

\subsubsection*{(g) 心软的爸爸}
\tocsubsub{(g) 心软的爸爸}
\label{ch18:local:h60}

\notehead{147}{}
\phantomsection\label{ch18:local:n55}

\dmark{}\origpage{684}我不知道我朋友这份饥渴,是否也施于我之外的其他人——施于那些嗅得出我的“气味”的更年轻数学家。这方面没有回声传到我这儿。可对我分明的是:正是经由他与我的关系、借着一份在科学世界里诚然罕见的情势,他身上那份在阴影里苟活的癖向,一夜之间变成一份吞噬性的饥渴。在我出发的情节里,当他一脸严肃地向我解释:他已把一生毫无保留地献给了数学\fnote{223}{(*)}{关于这段情节,见“兄弟与配偶——或双重签名”那条注(n\textsuperscript{o}~134)。}——他无疑“相信”他所说的话;而我本人,尽管有点发懵,也不曾想到怀疑他的话。可倘若我的耳朵更灵,或者更好地说,倘若我当时已有那份成熟,去倾听并信赖一份“更灵的耳朵”——它确确实实存在于我身上、一如存在于每个人——我本会知道:他关于他自己所说的话,前一天或许还是真的,可那一天已不是真的。那是一份高尚的缘由,为一桩可疑的行为而设:那桩行为的意义尽管昭然,他与我当时都没有那份质朴去直视它。那一天攫住他生活缰绳的、至今不曾松手的,是\textbf{别的东西},而不是那样的激情。

因此,是我这个人、或毋宁说他与我之关系中的某种东西,(借着合适的机会)成了触发器:触发那份支配他生活的力量之严苛的变质、触发他投入数学的意义与方向的严苛的变质。此刻该记起葬礼那些著名的“联”或“侧面”了——它们在11月13日的反思(“回顾(1)——或一幅画的三联”那条注,n\textsuperscript{o}~127)及紧随其后的注(“回顾(2)——或死结”,n\textsuperscript{o}~127$'$)里被摆到聚光灯下;此后,它们在路上多少有些被淡忘了。十天前,在“天鹅绒的脚掌,或微笑”那条注(n\textsuperscript{o}~137,12月7日)里,我曾多少作态记起它们。在那里,我尤其重新接上了一份直觉:那份“过继的父亲”的永世角色——我注定要在我年轻朋友身边扮演的;而依我看,它在他身上保存至今、至今仍在起作用。借着这场反思,我再次毫无保留地表达一份确信——它想必是在逝去的六七年里逐步\dmark{}成形、成体的(甚至可能更早\origpage{685}):“正是绕着这一侧面(他对我这个人的把握中的父性侧面),冲突打了结——一份早已在他身上的冲突,远在他听说有人念出我的名字之前……”(这便是那个著名的“超级父亲”之联;而“超级母亲”之联,至少眼下,仍在冥府。)

况且,只隔一页,那份著名的“微笑与天鹅绒的脚掌”风格便头一次匆匆亮相,成了注意的对象。系于它的联想,在随后几天里,似乎先把我引离我朋友这个人,也引离那份“父性”的隐蔽侧面——他在自己生活中派给我的角色。直到今天之前,这个侧面再没被谈过——人不可能同时想一切事,更不可能同时谈一切事!可是,就“想”而言,依我看,在某个地方、在那份模糊却依然在场、依然起作用的背景里,那个父性侧面的念头必定在场:它必定像一份有效而低调的兴奋剂,推动着那场关于“天鹅绒中的利爪”风格的长长枝蔓。毕竟(事后我才对自己说清楚;可它当时必已作为一份弥散却强有力的动机在场……),“父亲”的形象与那份著名的风格,绝非无缘——恰恰相反。甚至可以说:小女孩(或小男孩,也无妨)一生中见到的第一个被这份风格精微而利落地(虽不总是温柔地)引领的人,不是别人,正是爸爸!

而只要那无邪的小丫头(或小捣蛋)把这份风格、这份手艺收为己有——它几乎与此同时学会说话一般,变成一份第二天性,或相差无几——那么头一个试验品与受益人,毫无疑问,正是那位傻大个爸爸!通常,当我见过这份游戏上演时,其中还叠着一份宿怨的暗恨,外加一份嘲弄的故意。当然,在大多数家庭里,对着父亲的宿怨缘由并不缺——还不算那些由贤惠的妻子巧妙地暗示(乃至凭空炮制)出来的。可在我朋友身上,我从不曾有过片刻感到这样一份宿怨或暗恨的调门。当我见他“为取乐”而伤人害人时,那是\textbf{真的}(如我\dmark{}所感到的)\textbf{只为取乐}\origpage{686};我相信,不是为那份他施加的苦难或羞辱本身之乐,而毋宁是那份暗藏的陶醉:凭他的兴致、在他那手已臻化境的特殊风格里,行使一份\textbf{权力}——无疑更醉人、更够味,凭那份带“\textbf{乖僻}”“\textbf{禁忌}”意味的作料(为\textbf{取乐}而伤人、或叫人受苦)——而那,唯独他能做得起:精微地、不动声色地,此外一切照旧,想玩多久玩多久,管够……\fnote{224}{(*)}{作为周详的例解,尤其见“乖僻”那条注(n\textsuperscript{o}~76)。}

\subsubsection*{神经中的神经——或侏儒与巨人}
\tocsubsub{神经中的神经——或侏儒与巨人}
\label{ch18:local:h61}

\notehead{148}{}
\phantomsection\label{ch18:local:n56}

(12月18日)随着昨晚的反思,我感到葬礼图景的那个“第一层面”——以朋友皮埃尔与我的关系为中心——继续走出未被理解的迷雾、走出那份浑然未明的感受。好久以来,我便面对这样一项任务:在这个第一层面里,安插某个“超级父亲”之联;而我虽不曾真正对自己言明,那份联却似乎不很乐意体面地就位。倘若说我有过一个学生,始终让我感到与我“自在”至极、毫不紧张——就我记忆所及,任何时刻皆然——那正是他!诚然,我对我们最初几次见面的记忆已不多,也无从断言:当时他身上是否已有那份紧张——往往难以察觉、却实实在在——它在我们头一次接近一位被(以这样那样的名义)授予权威或声望的人、对其怀有特殊期待时出现。至少很可能:这样的紧张曾在场,而我对它的注意,不会多于对任何一位我偶然结识的年轻研究者。确定的是:倘若初接触时有过紧张,它很快便消失得无影无踪。借昨晚显出的意象说:他与我在一处,如同一个孩子(或曾经的孩子)与一位他从来不曾需要惧怕、也很少拒绝过他什么的心软爸爸在一处。

昨晚停笔后,我又把那份情势想了一遍。此刻向我显出:我朋友与我的关系,运作在两个截然分明的层面\dmark{}上,\origpage{687}而且(似乎)相互之间没有沟通。其一,大概在我们相遇之后的几个星期、几个月内便已确立,是个人关系的层面——也就是“心软爸爸”的层面:好得不能再好,毫不唬人,自己也有几分顽童气,包括在他的工作里;以致他对我的态度里有一层——我几乎想说\textbf{母性}的——色调:我已有一两次机会提到过:那是人们给一个毛手毛脚、有几分闹腾、尤其天真得没人比得上的孩子所给的那一种。而且,说得更进一步:在工作层面上、客观地说,他确实没有什么可被唬住的理由。当然,数学里我懂得许多他不懂的东西(他花几年便轻巧学会了),尤其,我有一份数学的经验,那还是他所缺的。可他有那份消化的敏捷、那份眼光的锐利——在纠缠混乱的情势里迅速辨明方向——常叫我叹为观止,而那是我所缺的。倘若我自己偶尔叫同事们吃惊,那主要凭我工作中那份不寻常的\textbf{快手},它主要归因(我相信)于我对数学工作的某种进法。可我的年轻而辉煌的朋友,断然没有理由为此唬住:他自己的快手,只要他肯动笔(那事他并不讨厌),分明比我的还要有效。

我朋友与我的关系的这个层面——“心软爸爸”的层面——依我看,涵盖了他对我的意识影像的全体,也无意识影像的很大一部分。依我看,正是这份影像,经由多半自童年起便已铺设的路径,唤起一份回应性的反射性欲望:那份著名的“天鹅绒中的利爪”的游戏——一份恰恰要求人对伙伴完全“自在”、完全“有把握”、并由此也自信的游戏\fnote{225}{(*)}{(12月29日)这份断言,只在表面上被某些个案(不包括我的朋友)驳倒:在那里,“领舞者”似乎(至少初看)被他所玩弄的人唬住、乃至慑服。可那是一份\textbf{架势},为事态的需要而设——而演员本人当然是头一个受骗的(我是说在意识的层面):这对赋予那份架势某种非即兴的“真实”外貌,必不可少!我所知道的这类游戏最极端的个案,是我母亲对我父亲。关于这一点,见“颠倒(1)——或激烈泼辣的妻子”与“颠倒(2)——或暧昧的反叛”两条注(n\textsuperscript{o}~126、132)。}。这是那份全然自信的层面:它安身于对一份情势的贴身认识,又进一步由经验印证——那份经验,被有意识与无意识的知觉与评判官能\dmark{}\origpage{688}以完全一致的方式解读。游戏本身是隐蔽的、对当事人自己也是无意识的(至少我这样推断);可那份自信之感、及作为其根基的对实在的知觉,属于意识的、理性的、“客观”的领域。

而另一个层面,则整个是无意识的(至少是我的印象)、不受控制的、也不可控的:一份非理性的本性,它似乎公然违逆、并嘲弄对“客观”实在(我方才唤起的)的一切有理有据的认识。在这个层面上,本来意义的个人关系——系于对他人的一份多少写实的知觉——消失了。我在其中显为一个巨人:强大而被暗暗嫉妒;而我的朋友自感侏儒,被他自己那份无可救药的微不足道之信念压垮,同时被一份疯狂的欲望吞噬:不是要他自己成为巨人——他以不可改变的境况生为侏儒——而是以这样那样的方式把自己\textbf{抬到}那个高度、至少\textbf{装出}巨人的样子;或者,更隐秘、更阴险地——那份疯狂的欲望:\textbf{做那个巨人本身},或至少,\textbf{冒充是他}。我在这份欲望里辨出又一层细节:它宛如那份欲望——在场于接近表面的层次的那份——在更深层的回声;表层那份欲望,恰恰在“以天鹅绒的脚掌”这份游戏里找到象征的满足,而它正是游戏的神经与发条:\textbf{角色颠倒}的欲望。在较浅的层次,涉及的是阴阳角色的颠倒:被支配-支配者、客体-主体。可这份关系在此并不合用:因为巨人绝无支配侏儒的操心——他满足于做巨人,并由此,在不知不觉、也不曾操心之中,做一份永恒而灼人的\textbf{挑战}:对着那个自感被一份无可救药的侏儒境况压垮的人……对方感到自己被圈在这份傲然的无知里,他把它感受为一份默契的轻蔑、一份冒犯。他心火烧灼,要颠倒的正是这份关系:他自己显为巨人,把巨人贬入微不足道——以\textbf{遗忘}致微不足道,不然便以\textbf{嘲弄}致微不足道,作为对那份他自感身处其中的无知与轻蔑的正当回报。

我方才说过,两个层面——“心软爸爸”与“巨人”——“似乎没有相互的沟通”。细细想过之后,眼下我倒觉得\origpage{689}:两者之间确有一份沟通——哪怕只凭那份颠倒的欲望:一个层面上的欲望,此刻显为对另一个层面所确认的相似欲望的一份“回声”。初看,我曾觉得:那个“侏儒-巨人”的更深层面上的角色颠倒,不属于阴阳型的角色颠倒。真确的是:这份颠倒确实不属于被支配-支配者的类型。可是,再想一层,毫无怀疑:巨人身上所体现的\textbf{价值},是阳的、超级阳的价值;而侏儒显为阴的非价值的化身——我是说,按我朋友的意识形态抉择的措辞;那些抉择,与我们关系最初年月里我自己的抉择,相去不远\fnote{226}{(*)}{阳的或超级阳价值之拣选上的这份一致,一直延续到我出发的1970年。其后的年月里,我有意识层面的价值系统,“翻转”向阴的、超级阴的抉择——见“阳扮阴——或主人的角色”那条注(n\textsuperscript{o}~118)。}。

这份断言将变得清楚,一旦我在“侏儒与巨人”的影像与实在之间架起一道桥——至少,显豁这份影像在我朋友与我关系的史前史与历史中的起源。关于“史前史”,几乎无须说明:这样一个有意识或无意识的影像,唯有借着那份埋藏极深的“对己之轻蔑”方能诞生——我在反思中已屡次唤起它;或者更好地说:这样一份影像,不是别的,正是这份轻蔑的一份可触的、多少具体的\textbf{落实}。甚至或许可以说:这份“秘密的信念”守候着一份能给它充当支撑的情势,并同时催生那幅表达它的稻草人影像。我相信:在心灵的万事万物之中——无论埋藏多深——都活着一份力,催促它表达自己,往往以象征的方式。这份表达,固然常常自身仍是无意识的,却并不因此少一分作用——恰恰相反——在日常生活的可见言行层面上。

回到这一次是我朋友与我这人关系的\textbf{历史}:它肯定也始于我们相遇之前。他大概是在1960年前后、在布鲁塞尔最初接触数学界的时刻,听说过我的——也就是说,比我们相遇\dmark{}\origpage{690}早四五年,当时他十六七岁\fnote{227}{(*)}{(12月29日)这条编年信息,我取自Pierre~Deligne1975年为(比利时)“国家科学研究基金会”“五年一度奖”撰写的那份两页的“传记小记”(Rue~d'Egmont~5,~1050~Bruxelles)。我打算在一条后续的注里回到这份小记:谈Deligne去年10月的来访。正是在那次来访中,我由他得知这份小记的存在;经我请求,他后来惠寄给我。也正是在这份小记里,我找到了我朋友身上某个影像的具体形态“侏儒与巨人”——一份弥散的观念,已在葬礼反思的进程中逐步析出。它开始显形于“下葬”那条注(n\textsuperscript{o}~61),并在“驱逐”“死结”“颠倒”“屠杀”“……与电锯”“葬礼悼词(2)——或力量与光环”各条注的反思中逐步明确。唯有借着本注,这份知觉才开始“落位”于葬礼“第一层面”的一份连贯总图。
(1985年3月)关于Deligne的传记小记,见“信条——或假中之真”那条注(n\textsuperscript{o}~166)。}。这肯定不是偶然:他请求教他数学家手艺的、正是我而非别人;或至少,请求教他那份即将成为他著作之中心主题与中心工具的东西(即代数几何)。在我们相遇之前,我以什么面貌向他显现(至少作为数学家),想必只能是我的招牌影像:它把我做成一份英雄式的、声名显赫的化身——数学世界里通行的那些主导价值的化身;而彼时,他自己还是个普通的学生,刚从中学毕业。他对我怀有的这份影像——它恰恰也是我乐于示人的那份——况且还不是一幅单纯的“埃皮纳勒”版画,专供渴望荣耀的中学生们做梦之用。它由可触的实在筑成;而他无疑颇有直觉,早在那些年头,便在与年富力强、十足在局内的数学家的接触中,嗅出它的气味。自1965年起,他更是比任何人都有资格,亲自来掂量我的斤两。我当时在他身上感到一份着迷:为一份向他敞开的景——它在我身上诞生、成熟于已逝的十年之中,并继续在他眼前铺展、发育。当时对我毫无怀疑:这些他“仿佛从来就认识”般纳为己有的景,将公开为他效力,充当灵感与工具,去发展更宏大的景、一部与他的手段相称的作品。事情全然不是如此——只是借着这份将近二十年之后、对一场葬礼的长程禅观的光照,我才隐约看清:那份对我所传递之物的精微而热烈的知觉,如何\textbf{同时}被用来充实、支撑一幅\textbf{稻草人影像}\origpage{691}——一份畸形的影像:它凭第一手的、无可辩驳的实在的要素而加固;一份足以\textbf{瘫痪}人的影像,一如它所表达的那份“贴身信念”。他对我在传递之物中的那份“伟大”与深度的知觉之敏锐——那份他独力(且毫不费力)整体纳为己有的知觉——正是这份敏锐与这份鲜活,曾构成他的力量,却旋即反过来指向他自己:把那幅畸形的影像,衬得愈发触目、愈发专断。

三天前,我曾自信触到了我的朋友将近十五年来所扮角色的“神经”——而当时确凿无疑:我刚触到一处要害:那份吞噬一切的\textbf{饥渴},渴着某一份\textbf{游戏}——一份精微的权力游戏;它同时,是一份特定角色颠倒之欲望的象征性的、昙花一现的满足……随着今天的反思,沉入更深的层,我觉得此刻触到了\textbf{神经中的神经}、那份更隐秘的\textbf{刺棒}:它不停地唤起并滋养这份饥渴。因为在“心软爸爸”的层面,固然有着机会与全副自由,安安全全地玩这份游戏:潇洒而精微地领舞,稳操每一招的胜券。可无疑,一旦没有刺棒,那轻易机会的魅力便钝了。而一如我昨天还确认的:面对心软的爸爸,没有那份咽下去的怨怼的刺棒、没有那份暗藏的宿怨——这正是人们叫它“心软”爸爸的缘由!这份缺席的刺棒,我方才忽然触到了:顺着种种联想,仿佛凭一份早已备好的认识的口授,我被引去描述那个“另一层面”——那份“不受控制也不可控”的层面:侏儒与巨人并肩居于其中。

最初那份直觉尚且含混的印象——两个层面之间没有相互沟通——顿时消散,让位于一份理解:它由“神经中的神经”与“刺棒”这双幅意象同时表达、同时引发。这次改用“层”的措辞——一些浮面,一些深层——我愿以第三幅意象再说一遍:深层滋养、维持着浮面的运动;深层是它们的根基,牢牢锚定在“我”的结构之中。没有这份根基,表面的躁动很快便会消散,终于让位给别的东西……

\subsection{另一个自己}
\label{ch18:local:h62}

\subsubsection*{(a) 缓期中的宿怨——或事物的回转(2)}
\tocsubsub{(a) 缓期中的宿怨——或事物的回转(2)}
\label{ch18:local:h63}

\notehead{149}{}
\phantomsection\label{ch18:local:n57}

\origpage{692}\dmark{}(12月20日)自五天前的反思以来——尤其自那天第二条注“秘密的神经”(n\textsuperscript{o}~145)里继续的反思以来——我感到:对葬礼图景那份著名“第一层面”的工作,骤然换了气象。在那场反思之前,我处于一个有点尴尬的位置:面对一幅拼图,却觉得对它不甚了了。自4月起,我便一直努力把拼块一一收拢、细心编目。倒不是缺拼块——不是的,我反倒觉得太多了!总之,拼块肯定足够拼出一幅图:也许只是局部,却是站得住的一幅。我掷上桌子的最后一块拼图,是那份“颠倒”(阴阳的)——它自“阴阳之钥”的发端起便被存而备用(作为一场我应许自己要回头谈的“联想”),终于在11月10日那条注“阴的葬礼(阳葬阴(4))”(n\textsuperscript{o}~124)里,以预料不到的强力闯了进来。其后的三十五天——直到五天前——大体用来把已出土的拼块翻来覆去,顺着最强制地要求我注意的种种联想\fnote{228}{(*)}{作为整场阴阳反思出发点的那个“拼块”(自10月初起),直到十四天后——11月24日,在“颠倒(3)——或阴葬阳”那条注(n\textsuperscript{o}~133)里——才重新登场、得到显豁。}。我本以为,这样做下去,那些拼块终会自行拼合,终于显出那幅未知的图景。事情全然不是如此。恰恰相反:它们继续互相顶牛,一如十张各不相同的报纸碎片被乱扔在那里,却要我来拼!我开始盘算:到头来,我是否会被迫作一份拼块的最终清点、再作一份关于其拼合的问号清单,就此收场……

五天前,情形变了:翻来覆去地摆弄那些著名的拼块、摩挲、嗅闻之后,终于有什么“咔哒”一声对上了——其中一块(那份藏在某个\textbf{风格}背后的\textbf{饥渴})骤然被承认为“要害”。我当即有一份\textbf{质变}的印象:一份此前一直缺席的\textbf{视景},已然\dmark{}开始就位\origpage{693}、要从那块拼块出发组织起来。仅仅两天后,我在下一条注(“激情与饥渴——或逐步加码”,n\textsuperscript{o}~146)里接续反思时,正是以这样的措辞表白自己的。而我的预感当天便已开始得到印证:随着“\textbf{心软爸爸}”那块拼块的现身——人们简直要说,它是应“要害拼块”之召而来的,为的恰恰是严丝合缝地嵌入!

那份“\textbf{超级父亲}”拼块——它一直在场(从《收获与播种》第一部分继承而来,并自“阴阳之钥”的发端起便被重拾\fnote{229}{(*)}{见“敌对的父亲(1)(2)”两节(第~n\textsuperscript{o}~29、30节),以及“超级父亲(阳葬阴(2))”那条注(n\textsuperscript{o}~108)。})——此刻似乎转眼成了盈亏账上的弃子,仿佛它只是不小心在那里迷了路。在“心软”新拼块\fnote{230}{(**)}{“新”拼块这个说法,或许不完全站得住。但至少,它是一块此前逃过编目的拼块——正因为它太显而易见了!}的新鲜印象之下,我倾向于忘掉:那份著名的超级父亲(这一回,一点儿也不“心软”)确确实实与朋友皮埃尔和我的关系有涉,即便它不曾领着那台戏(差得远了……)。可我终究在下一场工作里想起来了——必然如此:恰恰在那一刻,我正准备对自己解释:为什么拼图里那块永世的拼块,其实与此毫无干系!总之,它是“心软”拼块的“恰恰反面”——后者方才以那样的从容自行就位。可不然:细看之下,这块据称与游戏无关的拼块——它的轮廓一直最模糊——骤然显出了形状:“取得”了那幅力量意象的形状(道出这幅意象的,不是别人,正是朋友皮埃尔本人\fnote{231}{(***)}{个中细节,见上一条注“神经中的神经——或侏儒与巨人”(n\textsuperscript{o}~148)的最后一条脚注。}):\textbf{侏儒与巨人}。起初,见它以如此鲜明的轮廓重现,我料想它会与那对已就位的“要害双拼块”(由心软爸爸、与那份强制性的“让他转起来”的饥渴——一个电话这边、一个小动作那边……——构成)“不相沟通”。可它偏偏显形为“神经中的神经”、一份更要害的拼块,与已就位的拼图部分毫无摩擦、毫不脱胶地拼合!

这块拼块,顶着它的旧名“超级父亲”,其实已被擦过许多次\dmark{}\origpage{694},甚至被拿在手里、像其他拼块一样翻来覆去,甚至(我现在记起)被宣布为关键拼块、“图景的心脏”云云;可是——或许由于它不曾由一幅醒目的意象来化身(那意象由当事人自己提供),而尤其无疑由于它那荒谬的、有违常理的本性——按当下通行、举世公认的粗浅“常识”来看,它甚至全然滑稽——我对这块该死的拼块尴尬而近乎羞愧;它在我手里发烫:永远不会有人(包括那个还执拗地活在我身上的某位“我自己”……)肯拿它当回事!还不如悄悄打包收起,拿些更拿得出手的拼块来“玩”!

我方才谈“关键拼块”、“图景的心脏”云云——就那块变成了“侏儒与巨人”的拼块而言——我想的当然是“对己之轻蔑”的侧面,而非“超级父亲”的侧面。眼下,对这块刺棒拼块、或“神经中的神经”作“超级父亲”的指称,还为时过早,也无从论证。我是说:至少初看,那位著名的无面孔、双手硕大无朋的巨人,并不显得有多少父性。倘若非得给它一个名字,配得上它的,似乎是“超人”或“超级雄”,而不是“超级父亲”。所以,细细算来,后者眼下确确实实还“记在账上”——至少眼下如此,一如“超级母亲”那块拼块(或“联”):那份我日后同样要回头谈的。

眼下,最要紧的依我看,是试着以我的朋友这个人身上的阴阳动力学,给图景中已就位的部分定位——连同那份“秘密的神经”与更隐秘的“神经中的神经”。这方面,我手头有三件裸的事实。其中两件,由那份阴-阴的“双重签名”表达\fnote{232}{(*)}{“双重签名”的观念,是随着“兄弟与配偶——或双重签名”那条注(n\textsuperscript{o}~134)进入反思的。}:朋友皮埃尔基调属“阴”,无论在可称作“习得人格”的层面——它主要在与他人的关系的调门里表达——还是在“先天人格”或冲动的层面——它主要(至少对我这样一个外在观察者而言)经由自发的、免于“东家”插手的工作风格表达。第一件事实关乎习得人格、关乎“我”的结构(或以更形象的措辞,“东家”的\textbf{脑袋}),它似乎表明:这份结构是在童年、在人生最初几年里完成的,凭与一个“阴”性模型的认同。这在先验上并不排除:模型曾是父亲——倘若父亲本人(如我看来确系如此\dmark{}\origpage{695})便有一份基调属阴的“习得人格”。但另一方面,我朋友身上那份癖向:渴着一份权力游戏——它在我们这些地界(即便不是处处、事事)典型地(若不说排他地)是“女性”的,而更确切地说,它是妻子惯于与丈夫玩的那份游戏中的\textbf{头号}游戏——这份癖向使我推测:认同是与母亲这个人完成的;正是从她那里,他“继承”了这份饥渴(或一份对此类饥渴的癖向);也正是从她那里,他接过了那份合宜的“风格”(或“战术”):“天鹅绒中的利爪”。

有可能,父亲既是一位心软丈夫,又是一位心软父亲;而我的朋友,久已有的是机会拿他做自己第一个“试验品”,在他身上磨爪子(和绒毛!)。但也有可能:我朋友身上这份癖向或禀赋,一直闲置未用,直到与我相遇之后——因为头一个现成的靶子、即他的父亲,缺少足够鲜明的阳特征,不足以“\textbf{挑起}”这份饥渴、又不足以让那份对付倔家伙的考验过的战术\textbf{得手}。老实说,凡我记忆所及、落在我与朋友相识最初年头里的印象,没有一份能暗示:他对那份游戏已然熟稔,甚至已然实践过。无论如何,即便带着回顾,我在他与我的关系里、在他与别人的关系里,也辨不出一丝痕迹——比如稍微带点“被惯坏的孩子”味儿的做派。所以我反倒倾向于认为:他身上这份癖向当时还潜伏着;它得以发育、并取得我所知的那份对他生活与工作的掌控,是在我1970年“亡故”之后(他时年二十六),借着一份格外诱人的情势。

这里该回顾的“第三件事实”,是我的朋友对一份与通行价值相合的价值系统的抉择:对“阳刚”(或阳)价值的抉择。而这份抉择,在逝去的十五年里,依我看,在他身上愈发转向“超级阳”。在他的个案里,这份抉择有一道一望即知的矛盾:一面采纳\textbf{阳}的“官方”价值,一面却在大多数要害特征上,照一个\textbf{阴}的模型自我塑形\fnote{233}{(*)}{这是一类矛盾,尤其多见于女人;我自己的生活则与它无缘。}。而这并不是说:这份价值的抉择纯属“赝品”、只是为情势所需而竖起的假幌子、只在心灵的边缘层次\dmark{}流通\origpage{696}。侏儒与巨人的力量意象,从深层起作用;倘若对阳的抬举不曾也在这些层次里被内化,这份意象便会失去意义,连同它所引发的那份强制性的颠倒之饥渴。毫无怀疑:这份矛盾,必定给那份“裂缝”“阴险无能”的“贴身信念”追加一份活力——而与此同时(或许只缺童年里一个可供模仿的合宜“模型”),他自知(在心底)根本\textbf{不同于}他“\textbf{理应}是”的那个人!

倘若我的朋友——这在我看来是说得通的——不曾在他父亲身上找到那些照当时周围的共识\textbf{本该}在的、而他本可纳为己有的特征,这必定在他身上唤起一份弥散的\textbf{宿怨}:一份无法攀住任何具体怨怼的宿怨——对着一位心软爸爸,他唯一的缺点竟是太“软”!这份宿怨,苦无一个可以攀挂的“挂钩”,想必始终“空置”、\textbf{待命}一份相宜的靶子——恰恰是这样一份靶子:首先,(由语境)扮出父性面孔;而且,担纲此角的\textbf{资质}昭然,凭那份不容置疑、甚至耀眼乃至过分的特征的在场——正是他“原装”父亲身上所缺的那些特征。也恰恰是这些特征,使新来的“父亲”成为理想的\textbf{靶子}:在一场此地早已备好、只待相宜伙伴的“游戏”里——伙伴别名“替补父亲”,别名(我们终于来到这里!)“超级父亲”!

忽然,我觉得自己回到了一片非常熟悉的地界——直到此刻我才认出它。这是一片我做过二十年囚徒的地界:在我一生唯一的一场婚姻里(我的五个孩子中有三个出自这场婚姻)。在上面几行里,我完全不是存心(毋宁说像一个小心翼翼在黑暗中摸索、认识周围事物的人),也依次描述了那个曾是我妻子的人对其父亲、继而对我的关系中的要害力量。我无从说清,这份认识(或毋宁说那份无可辩驳的直觉)——关于这两股力量在她身上的无声而执拗的在场、及它们的相互关系——是何时、如何来到的。有一天我懂了,却从不曾费心稍加反思:那份无情地支配着我妻子与我的关系的力量——自我们成婚的最初日子起——正是被那份对我宿怨所驱动:怨我不曾在一段无依无靠的童年岁月里,守在她身边,做一个\textbf{别的、真正的}父亲……

\dmark{}诚然,我知道:我朋友的童年毫无“无依无靠”可言;而他发育成形、并为我所认识的那个人格——六十年代直至如今——与我前妻的人格相去甚远\origpage{697}。但这并不妨碍:在那些明显的差异之外,我看到——在正从阴影中浮现的图景那部分里——一份惊人的相似,与另一幅我熟之又熟的“图景”相似。这份相似,显于对父亲关系的本性(系于一份阳特征匮乏的父亲气质),以及它对一份成年关系的回响:在两人身上,这份成年关系都支配了其一生,都成为两人身上冲突力量之瞄准的焦点\fnote{234}{(*)}{(1985年2月19日)朋友皮埃尔与我的关系、与我妻子(自婚后最初日子起)与我的关系之间,确有一份惊人的亲缘。这份亲缘更超出对单一我的关系:两人都终于发育出一份癖向,把某些与我以情感相系的存在(一者是我的孩子,另一者是我的学生)变成\textbf{器械}——经由他们来打击我。}。

差一点,我便会略而不提第三份“相似”——它在我自己的生活里却不无分量:在这两份关系里,担纲的\textbf{主角},每一次都\textbf{不是别人,正是我}。而在这两个个案里,把我指派给那份我注定要扮演的“超级父亲”角色的,是(除一份不成熟之外)还有一样东西:它自我童年起,或许便比世上任何东西都更亲近于我——也是我在其中投入最无度的东西:一副天生超常的“阳刚”身量……

就这样,我以一道比八个月前更不同、更透彻的光照,重新寻回那份“事物的回转”之感\fnote{235}{(**)}{见“事物的回归(或一脚踩进菜盘)”那条注(n\textsuperscript{o}~73)。}——今天一如当年,带着一份难以置信的惊异的调门(事情似乎落得太“巧”,巧得不像是真的!)。而且,这一回也如此——只是调门比当年那阵骤然的大笑更收敛——有一份滑稽的知觉,为这些无情的“回转”添上幽默那份更柔和的音符。

\subsubsection*{(b) 天真与冲突——或绊脚石}
\tocsubsub{(b) 天真与冲突——或绊脚石}
\label{ch18:local:h64}

\notehead{150}{}
\phantomsection\label{ch18:local:n58}

(12月22日)昨天我又没找到时间在注上工作,只除了细心重读并修改前天的注。最近这些天,我的精力被通信\dmark{}\origpage{698}等任务岔开;我又有点按捺不住急切(这可不是新鲜事!):急着重新单独面对自己,推进已经开始的反思。在《收获与播种》的这第三部分——以眼下这场反思“阴阳之钥”为中心——写作分明更慢了:阴阳的动力学,是深入葬礼意义的那根恒常的引线。倘若不留神上好闹钟,以便在工作约三小时后安排一次休息(活动活动筋骨,或提醒自己夜已深、该停下了),一整夜会像一瞬般过去!那三个钟头,每一次都是这样过去的:我只觉得刚刚开始(或刚刚接上),只打出两三页可怜的纸——有时只有一两页——刚刚够巡完某场看似平淡、我本想一口气跨过去的联想……

按每小时或每天的页数计,有一种极慢的推进之感——对这份印象的自然反应,既然一份滚烫的实质正在我鼻子跟前拽着我向前,本该加倍再加倍地赶——一如直到最近几年我还惯做的那样。可我知道:这正是该避开的陷阱——发现工作中那份出奇的“轻捷”的陷阱\fnote{236}{(*)}{见“陷阱——或省力与衰竭”那条注(n\textsuperscript{o}~99)。}:只消轻轻“推”一下,便确信自己在切实前进——也许慢,但稳;像一个人双手扶住一辆钢犁的扶手,前面由一双强健而沉稳的公牛牵引:他慢慢地、稳稳地犁开一条路,一垄又一垄,穿过一片紧密的、有时倔强的、却又柔韧的土地——它顺从那把闪亮的犁铧:犁铧不慌不忙、精微地把它翻开、穿透、掀转,成大片大片冒着热气的黑土,把一份强烈而蠕动的地下生命带到光天化日之下。节拍或许慢,田地却广,每一垄犁沟似乎才动了动那片尚待开垦的荒地。可一天终了,一垄又一垄,田已犁遍;犁田人满意地收工:对他而言,这一天不曾虚度。他的辛劳与他的爱,是他的种子;他劳作中的快乐、每一垄尽头与长日终了的那份满足,便是他的收成与他的报偿。

\asterisms

\dmark{}\origpage{699}随着前天的反思——也许是在《收获与播种》的写作中头一次——我有这样的印象:我踏上了那片不确定的地界:尚未被直接知觉或感到、而依然(或许将始终)\textbf{属于假设}的地界。苦于没有能在看似半明半暗之夜里看东西的眼睛,我摸索着开出一条犹疑的路,毫无把握它是否“对”。当路分岔,我确实不曾掷硬币;我凭直觉与常识,为自己指出最说得通的继续方向——却毫无任何念头:它将把我引向何处。我这样循着、或这样画出的路,看上去处处“贴合”我所知的事实;这是个好兆头。但这并不排除:尤其在那些事实稀薄之处,另一条全然不同的路同样会“贴合”——前提也许是再把某条尚属生涩的事实挖掘一番……然后,在路的转弯处、在我自己的惊异之中,我忽然发现自己来到“一片非常熟悉的地界”:当年我曾长久而艰难地跋涉过、终于认识并离开的地界。一份几分钟前还显得晦暗的情势——被“想必”与“或许”的疑云包裹——骤然被另一份情势的光照亮:而后者,是被懂得的。当我追问那份朋友与我的关系中的冲突、在我与在他身上的遥远根源时,它们似乎经由一份骤然瞥见的深刻相似而显露——那份关系与另一份关系之间的相似:另一份曾以全然不同的分量,压了我漫长二十年的一生。

这份相似的显形,力量如此之大——我承认——那份迟疑、无把握、摸索之感旋即消散,让位于一份有把握、确信之感。当我在反思的结尾谈到那份感觉(“难以置信的惊异”):事情“落得太巧,巧得不像是真的”,那份感觉是对另一份的回答——另一份在底色里说:事情“落得太巧,巧得不\textbf{能}不是真的”!而后一份感觉——就我手头的事实现状而言,肯定仓促而毫无根据——其间并未调整过来:它始终作为底色在场,无论我愿不愿意。肯定地,若非借助某些我终于懂得、终于承担的经验——尤其,借我婚姻生活的漫长经验——那份“空置状态的宿怨”(总而言之,一份“缓期中的宿怨”)的念头,几乎无从来到;而恰恰是这同一个念头,又是那个“路的转弯处”:几瞬之间,它使我重新驶出那片婚姻经验的“非常熟悉的地界”。

\dmark{}\origpage{700}当然,可以说:一份无意识的故意,把我引到一个早已被指定的地点——那个地点,或许正说明一些关于我、关于这份故意的东西,而对他人身上的动机,则说明不了什么。同样可能:一份被承担的经验,使我得以把握他人身上的一份实在——否则,它会整个是个谜,怪我没有足够灵敏的“天线”(也怪我手头没有关于我朋友童年、他父母各自人格的可触事实)。

依我看,我离完成葬礼“第一层面的图景”的速写(信手写来!)已经很近。为拼上我手中剩下的最后几块拼图,必要时我会动用前一条注的反思中显出的领会要素(无论它们多么属于假设)。这倒也是一份办法:检验它们与我所知的其余事实的融贯。

前天的反思里,是拼图的“超级父亲”那块,显出了形状与轮廓。我起初把它认作“侏儒与巨人”那块——有点匆忙;在那块里,巨人毋宁显为一个尺寸骇人的“超人”,而不是“父亲”、或“超级父亲”。可后者那块,终于在同一场反思里重新显现——这一回,作为一份“缓期中的宿怨”的靶子:一份正在寻找靶子的宿怨——仿佛那位“超级父亲”恰是\textbf{被}这份宿怨所召、应这份召唤而来,兑现一份弥散的期待。倘若果真如此,可以说:倘若超级父亲(借这场合借得我的身量与我的面貌——它们分明是度身定做的)不曾出现在我朋友的生活里,那就得发明一个!至少,对那个我曾做其丈夫的人——而且我曾是“那支被盼了整整年轻一生的箭靶”——这份说法,对我已毫无假设可言。

于是,超级父亲显为“侏儒与巨人”那块拼块里那位“无面孔而双手硕大无朋的巨人”的“有面孔的一侧”。“侏儒”看巨人,想必主要看背影:巨人正忙着做他那著名的“力量展示”(10月5日“超级父亲”那条注(n\textsuperscript{o}~108)里谈到的)。于是,“超级父亲”那块拼块终于就位:与“侏儒与巨人”拼块的“巨人”一侧拼合。

至于“侏儒”那一侧,它的轮廓也在前天的反思里显得更清楚:它在此接上10月17日那条注“一半与全——或裂缝”(n\textsuperscript{o}~112)的反思。还是仿佛——常常如此——那份对“阴”、“女性”特征的永世拒斥(为了让位给“阳”、“男性”的特征),使我朋友“根本不同于他‘理应是’的那个人”,而他本人却照一个以“阴”为主导的模型自我塑形\dmark{}\origpage{701}。

这里须要强调:在整场已往的反思里,我从不曾想过、也不曾想过暗示:我朋友这个人打有以阴为主导的\textbf{失衡}的印记,即:在他习得人格的阳的、阳刚特征一侧,有一份匮乏、一个“空洞”。关于这一点,我提醒:至少在我认识他的最初几年,从他的这个人身上透出的那份印象,恰恰是一份\textbf{平衡}、一份和谐的印象——它使他对我、对我以为对所有当时认识他的人,都如此可爱。这份印象与另一份系得很紧——我在别处谈过\fnote{237}{(*)}{关于这一点,见“孩子”那条注(n\textsuperscript{o}~60),第五行列“我的朋友皮埃尔”内。}:他似乎保有了几分孩子的新鲜与天真——在他把握事物(尤其数学)的方式里;而且依我看,也在他对人的方式里。这份平衡、这份“新鲜”或“天真”,对我而言无可怀疑:它们是\textbf{事实},谈不上要把它们变没了。它们在我朋友身上表达为一份精微的感受性,而且时机一到,便以有层次而不含糊的方式,表达被知觉与被看见的东西。有一份坚定,正如有一份温柔。温柔在岁月中消隐了,只留下一副毡垫般柔化而空洞的甲壳——那消失了的温柔的甲壳;而坚定变成了封闭与坚硬,藏在一层借来的、矫饰的半音调式的门面背后。一份精微的阴阳平衡,在岁月中(无疑无人察觉)转变成了那永世的阳失衡——正是它,只是风格不同,曾支配我自己的生活,自我童年起。这是他的抉择;而抉择可以变——牌局从来不曾定死!无论如何,在我朋友的生活里,我从不曾知道有任何一段时期,打有阴失衡、疲软、放任或不一贯的印记;我也不认为曾有过。

这一切使得至少很可能:那个在他童年充当“模型”的人——那人肯定有强烈打上印记的阴特征——并不因此缺阳特征去与它们平衡。倘若(如我倾向于相信)那个人是他的母亲,那么我推想:她有着足够鲜明的阳特征(尤其对着父亲身上那些无疑较不鲜明的特征),以致显得是“最佳选择”——充当一个男孩的“男性”模型;而与此\dmark{}同时\origpage{702},凭这样的选择,助成一份和谐气质的绽放。

至此,一切似乎都尽善尽美,在最好的世界里,在一个(或许)没有任何不和搅扰的和睦家庭里。一切本都会尽善尽美——倘若没有那一小块绊脚石,取一份无声共识的形态、貌似再平淡不过:男孩理应像父亲,而不是像母亲……

\subsubsection*{(c) 天赐的情势——或封神}
\tocsubsub{(c) 天赐的情势——或封神}
\label{ch18:local:h65}

\notehead{151}{}
\phantomsection\label{ch18:local:n59}

(12月23日)依我看,要把葬礼第一层面的“拼图”拼完,只剩最后一块该就位。那是我在11月11日那条注“超级妈妈,还是超级爸爸?”(n\textsuperscript{o}~125)里称作“超级母亲”的那块。这份“超级”式的命名,起初由我葬礼悼词里那份以超级形容语一笔一笔画成的我的“肖像”所启发\fnote{238}{(*)}{见“葬礼悼词(1)(2)”(n\textsuperscript{o}~104、105)两条注,以及“阴的葬礼(阳葬阴(4))”(n\textsuperscript{o}~124)。}。肯定,一份对称的反射也起过作用:既然“超级父亲”已经在空气里,且不止一重!可是细细想过,我给那幅方才显现的影像起的名字,并不十分贴切。那幅超级阴的影像所唤起的,毫无“母性”的意味。倘若它与另一幅影像有对称关系,那是与“超人”的对称——钢的肌肉、IBM软件的头脑——而不是与“超级父亲”的对称。所以,在此情形,毋宁说该是“超女”或“超级妞”:垂到肚脐乃至更下的沉重乳脯(更不必说到膝盖……),臀部亦然,足以叫赫拉克勒斯做梦——至于头脑,且不必提……大抵是这样的调门。语言的不足,想必也多少逼了我的手:著名的“超人”(它本身也是晚近的发明——大力神赫拉克勒斯的一个现代版,分明已被时势超越)没有现成的“阴性”对等物。将就着,“超级妞”也罢,聊胜于无……

得承认:我拖着这块起错了名的拼块,拖了将近一个半月,不曾真正动它分毫,除了此地彼处以备忘的形式重提一句,作为一种许诺:总会去处理它,不过往后。最终,它实在不曾给我多少灵感;而这或许正系于那个贴不紧的名字\dmark{}\origpage{703}。毕竟,在我直到今天还在的数学世界里的所有朋友、(昔日)学生与其他同事之中,我会很难找出哪怕一位:对着他,我曾多少扮演过一个“母性”的角色;或者,我曾有印象,他觉得我担着这样一份角色。甚至那些对我而言我扮演的角色偏于“阴”、偏于接受的人——而不是“阳”的施教者、传达者、传递者的角色——也必定极少:粗粗一望,我几乎看不到谁(1952、53年我写论文的年头之后)——除了Serre,勉强……倘若我试着回忆我与别的数学家打交道时通常(甚至不妨说恒常)的心态,那主要是:我总有崭新的“地毯”要“铺出去”(借我当时流行的说法),更不必说那些(也出自我手)不那么新、却(依我看)几乎没用过、而在我看来为一家数学之宅的体面所不可或缺的“地毯”——在数学的某个我熟悉的街区里。换句话说,在我与我的数学“同行”的关系里——即便我们一同几乎只谈数学(在这点上,我该比我的任何同事朋友都糟!),我习得气质里的阳优势(或毋宁说超级阳的失衡),便重新主导一切,一如在任何其他关系里。甚至或许更强,鉴于我对数学的无度投入——那份投入,取利己的品格(无须说破),况且,恰恰由我积习已久的超级阳抉择所驱动!

正是这些显而易见的侧面——它们在我与其他数学家的关系的每一步显形——必定把它们那\textbf{另一份}事实遮蔽了,对同事们一如对我自己:我的数学工作风格、我对数学的进法,是阴占强势主导、“女性”的。依我看,正是这份特质——它在科学世界里显然颇为例外——也使这份风格如此\textbf{可辨认}、与任何别的数学家如此\textbf{不同}。这份风格确实“与众不同”,这一点,自我发表数学以来——至少自我的论文工作(1953年)以来——已由无数的回声传回我这里。况且,这份风格不曾少引起抵抗——我乐意称之为“发自脏腑”的抵抗:我是说,那些抵抗,依我看来当时(而且今天依我看依然)不能用可以称作“客观”或“理性”的“缘由”来辩解。这使我记起:我的论文(我在其中尤其引入了核型空间),投给《美国数学会回忆录》,曾被第一审稿人否决——一位在该领域工作过的、名声端正的数学家,他把我的工作看得多多少少有些“花架子”。多亏Dieudonné的有力斡旋,我的论文才不顾审稿人的否定意见得以发表\origpage{704}。几年前我得知:它属于过去二三十年数学文献里被引用最多的百篇文章之列\fnote{239}{(*)}{或许,此处我的记忆在骗我:说的可能是泛函分析里被引用最多的百篇(或二十篇?)文章。}。我推想,倘若我们前面还有二十年、三十年数学,同样的话对SGA~4也将成立——作为几何拓扑里topos视点的基本参考文献之一;而SGA~4,曾被我的辉煌朋友兼昔日学生Pierre~Deligne判为“读不通”(还有同一水质的其他考语\fnote{240}{(**)}{见“白板”那条注(n\textsuperscript{o}~67)。})。我知道(他自己也知道):它是我花时间最多、照料最极端的数学文本之一——尤其把所有关于site与topos、及范畴预备的部分,彻底重写、再彻底重写。这份格外照料的缘由,是:我分明感到,这是我正在奠基的“算术几何”的一块真正的奠基石——那份奠基,我已做了十年\fnote{241}{(***)}{这肯定也是那份缘由:Deligne如此执意要败坏这份文本的名声,以致他有时竟忘了自己偏爱的半音调式风格,毫无节制地来贬它!关于这一点,见前一条脚注已引的“白板”那条注。}。我也知道:当我做这份工作时,我久已(不想自夸)握有那份写数学的大师手法:既\textbf{清晰}——主导观念如一条无所不在的引线始终摆在前头——又\textbf{便览},便于查考\fnote{242}{(****)}{况且,正是借着熟悉SGA~4已誊清的部分(1965年,他刚到我的SGA~5研讨班),并亲手编写其中某些讲义(凭我的手稿笔记),这位Deligne在与我的接触中,学会了编写数学文本的手艺——尤其,把一份繁密而复杂的实质清晰地呈现出来。}。倘若我或许错了——提前四五十年的光阴,写(并让人写)一份周详的参考著作——那么,那些曾经成熟的年代(六十年代)骤然不再成熟,这笔账依我看记不到我头上!

\dmark{}这些与Deligne的联想,把我带回我出发之后的时期:同方向的回声,不止一次传回我这里——“如一阵阵阴险的轻蔑与含蓄的嘲弄”\origpage{705}。这层\textbf{嘲弄}的调门,是我出发之前那些“对我的工作风格的发自脏腑的抵抗”的征兆里所没有的——我方才唤起过。我在其中辨不出任何敌对的意图、或对我这个人多少恶毒的意图。甚至在Bourbaki之内,我也有机会唤起过这类征象\fnote{243}{(*)}{尤其见那条(无名的)注~n\textsuperscript{o}~5,《收获与播种》第一部分内。}——至少(如果我没记错)直到1957年前后:那一年,我关于Riemann--Roch--Hirzebruch--Grothendieck公式的工作,驱散了可能仍有的对我作为数学家之“扎实”的怀疑。1957至1970年(我“出发”之年)之间,我不记得知觉到对我工作风格的抵抗——除了偶尔在Serre身上\fnote{244}{(**)}{关于这一点,见“兄弟与配偶——或双重签名”那条注(n\textsuperscript{o}~134)。}——却从不带一丝敌意的调门:那宁可是一份皮肉的恼火反应。相反,我有这样的印象:我的朋友们有时感到不堪重负,因为我走得太快,而他们不愿把时间整个花在跟上我的全部著作——随着我把一部部砖头寄给他们,或(凭信件或当面)讲我正在炮制的东西。

我相信我已懂得那份对我的风格的“发自脏腑的抵抗”的本性——我方才唤起过。它的缘由,依我看,独立于那场后来才有的葬礼(那份抵抗,在其中终究扮演了一个重要角色)。这份抵抗,不是别的,正是对着一种“女性”\textbf{进法风格}——对着一门科学(在此即数学)——的(发自脏腑的)\textbf{反应}。这样的反应,在一个科学世界里,是常见的、“合乎事物本性的”——那个世界,与我们当今社会的任何其他局部小天地相比,不多让地被\textbf{阳刚价值}、及随之而来的感情、姿态、反应(尤其顾虑与拒斥的反应)所浸透。对那份特定的工作风格的抵抗反应——一份以“女性”为底色的创造性进法的化身——不过流自今日与最近几十年科学人的惯常制约——至少,流自我从来认识的那个科学世界。

一如任何出自制约的反应,这份反应实则毫无“理性”可言;而在它显形的那个人身上,有\dmark{}相当大的抵抗:仅只想到去检视它的意义,都有抵抗\origpage{706}。它被强烈地感到是\textbf{自身的辩白}——有点像道貌岸然的圈子里对“变态”的厌恶,或我们这里对“外来佬”的厌恶。可是,在我关切的个案里,我没有在这份反应本身之中感到一丝(有意识的或无意识的)对着我的敌意调门,而毋宁是一份\textbf{保留}、一份不利的成见——\textbf{单对着我的工作}。只是从那一时刻起——当我分明凭着我的风格(或不管我的风格,那都一样!)做出了此前人们不曾做出来的东西(而且事后也做不出真正别样的东西)——只是那时,这些保留才被收起,或许还带着几分不情愿……

无论如何,倘若在某些人身上,这些保留仍以默契、无意识的形式存留,那我太埋头于我的工作与我的任务,不曾知觉到它们。

老实说,依我看,至少很难想象:这样一份“发自脏腑的反应”,会如变戏法般消失,只因为某位先生证明了几样此前无人证出的定理。在有意识与无意识的接纳与拒斥的故意们做与散的那个层面上,两件事(“这样的工作方式不该被允许”与“某先生证明了某某定理”)真正毫无关联!

人们会说:既然如此,事情在我退出数学舞台之后起了变化,乃属正常——一旦我不再在场,去堵那些人的嘴:他们面对我的风格挑挑拣拣、撇嘴挑剔,却做不到用他们自己的风格做到同样的事。可这份“解释”瘸着腿走:因为它不顾那份嘲弄、那份压低嗓门的恶意的调门——它们从前不曾有。而且,凡我所知,没有任何东西能让我设想:1957至1970年间,我曾有时间把自己弄得叫我那些同行大会如此讨厌,以致一份宿怨或报复的动机会在我走后起作用。与我离开的那个世界的众多朋友,我保持过热忱的、有时深情的关系;而且(如我在别处所说)我不记得与任何一位数学家同事,在1970年之前有过一段敌意的关系。

然而,大会确实有过一份\textbf{后来的}对我的怨怼——一份集体“宿怨”的缘由,无论如何,一份集体的“报复”之举\dmark{}\origpage{707}:它虽保持默契,却“毫无破绽地有效”。这份“对一场异见的报复”的侧面,我在5月24日那条注“掘墓人——或全体教众”(n\textsuperscript{o}~97)里探测过。在那条注里,我把这些报复——对着我和那些不慎抬出我的人——的某个调门搁在了一边:恰恰是嘲弄的调门,它超出单纯的“不予受理”。而每当我感到那一“阵”,被点名的靶子,始终是某一份\textbf{风格}。换一种说法:使这份风格区别于一切其他的那份特质——它的“阴”或“女性”的本性——恰是天赐的情势:集体的无意识赶着抓住它,来洗刷一场异见所造成的冒犯,在以\textbf{排斥}施行的报复之外,再添上\textbf{嘲弄}的维度——一种嘲弄:它号称经由某一份风格,指认\textbf{无能}那份无可辩驳的征象。

而如今,随着“无能”这个词,某个不曾说出的东西终于被说出了:变得显然的是——这同一份“天赐的情势”,叠加到我“亡故”的情势之上,成了我的朋友兼昔日学生兼昔日继承人Pierre~Deligne的一次绝妙契机,使那份角色颠倒变得可触、可信、\textbf{被当真}——那个自感立于“\textbf{巨人}”跟前的“\textbf{侏儒}”那份疯狂而显然无望的欲望!(“骑在巨人的肩膀上”——借他简历里那句结语\fnote{245}{(*)}{关于这一点,见“神经中的神经——或侏儒与巨人”那条注(n\textsuperscript{o}~148)的最后一条脚注。})——从此,在众人眼里,当“巨人”的是他;而他将把那个“纯皮货”的巨人——千真万确!——指给全体大会嘲弄:一个吹牛大王、一个徒劳挥臂的“侏儒”——而那巨人毕竟曾(并且终究仍旧……)是“一份永恒而灼人的挑战,对着那个自感被无可救药的侏儒境况压垮的人……”

这场在“侏儒”与“巨人”角色分配中的惊人颠倒——在他与那个他者(那个被感到是一份\textbf{挑战}、必须不惜一切取而代之的人)之间——这场颠倒同时便是“\textbf{女性}”与“\textbf{男性}”角色的颠倒。分明地,正是作为\textbf{女性}的化身(臃肿、松弛、没有轮廓——从不曾被明白说出、却被热切否认),那个曾(并且终究仍旧……)是巨人的人,被交给\dmark{}人群(而首先交给变戏法的人本人……),当作可怜的侏儒、嘲弄的对象;也分明地,正是作为\textbf{男性}气概的英雄而典范的化身,那个曾是侏儒的人(而他,不管怎样、在他心底最深处“知道”自己是侏儒、且凭不可改变的境况仍是侏儒……)发现自己成了钢手的巨人,受着同一群赶来喝倒彩他者的人群的欢呼\origpage{708}。

这场颠倒,尽管是象征性的,分明与那个“私人”性质的“颠倒”毫无公度:后者凭一份考验过的战术(号称“天鹅绒的脚掌”)在一个狭小而无足轻重的“面对面”圈子里完成;一份可爱的小木马,他感到自己握着让“他者”运转、打转的线……侏儒让巨人转起来,行;可他终究还是、不可救药地还是侏儒!而侏儒的封神——他变成了巨人、还骑得更高,并把那个他骑在其上的人交给众人的嘲弄——这份封神,在大庭广众之下铺开:面对人山人海、欢天喜地的人群,他们来欢呼一位“侏儒”亡人的葬礼悼词,把它当作一场漂亮而令人陶醉的葬礼仪式分明的高潮。

\subsubsection*{(d) 否认(1)——或提醒}
\tocsubsub{(d) 否认(1)——或提醒}
\label{ch18:local:h66}

\notehead{152}{}
\phantomsection\label{ch18:local:n60}

(12月24日)随着昨天的反思,我有这样的印象:“拼装”葬礼图景的这个第一层面,已大体完成——至少,凭我手头现有的拼图“拼块”,我能做到的程度。当然,在葬礼反思的这第二部分(《收获与播种》的第三部分)里,我的用意已不再是收集事实材料(在“侦查”部分、I至X行列的进程中,我收集得尽够了),而是达到对葬礼\textbf{发条}的理解:经由众多当事人每一个身上那份秘密的动机(多半无意识)\fnote{246}{(*)}{(12月31日)这份“用意”,若按字面来取——鉴于“众多当事人”为数之众(哪怕只有十个!)——当然完全够不着。除朋友皮埃尔之外,我至多能形成一份总体的想法:在一个“集体无意识”里,多少圈定一些“动机”与“意图”——而那集体无意识,充其量不过近似地涵盖某个特定“当事人”的那些。}。这些动机,首要地,出自当事人与我这个人(作为“亡人”)关系的本性;或者更确切地说,出自我的所是之于他意味着什么——无论系于我从数学舞台的离开与否、系于环绕它的情况与否。

\dmark{}\origpage{709}这个“第一层面”,除我自己之外,便是那位在我葬礼上扮演“披祭披的神甫”、或“葬礼大司仪”角色的头号人物。他也是——在我出发之前的数学世界里曾为朋友或学生的人们之中——与我系得最紧的:凭一份格外强力的数学亲缘;也是唯一一个在我出发之后仍与我保持个人关系的人:那份关系延续至今。正是为这一切缘由,关于他,我掌握着一份材料:它的丰富,与葬礼与会者中其他任何人为我所知者,毫无公度可言。最后,在我认识的所有数学家之中\fnote{247}{(*)}{甚至在我认识的所有人之中——除两个例外。},他无疑也是远远最重的一个:那份他在自己生活中派给我的角色,压得最沉——分明远比通常派给师父的角色的分量更沉,哪怕那是人们倾注了全部身心的艺术(一如我自己曾倾注的那样)。这方面,我大约十年前便终于有所察觉:他派给我的那份角色,也漫溢到他的数学激情上(以及终于取而代之的东西上)。我身上这份知觉——多年来一直弥散——在葬礼反思的进程中,直到昨天,大大精确、充实起来。

依我看,随着昨天的反思,与那个以朋友皮埃尔和我的关系为中心的第一层面一道,终于就位、拼合的,还有“第三层面”:“(几乎)全体大会”——欢天喜地赶来,以它殷勤的首肯参与葬礼与下葬。如我昨天所写:5月24日那条注“掘墓人——或大会全体”的反思里析出的图景所缺的,是那份\textbf{嘲弄}的调门——它被放进对那个被当作亡人、“外人”、“圈外人”对待者的排斥之中。这份嘲弄的意义——它早在(11月10日)“阴的葬礼(阳葬阴(4))”那条注里便已清晰——昨天被重提并重新定位:这是对着被(在未成形的层面)感到为“女性”之物的嘲弄;而那份“女性”,经由(同样未曾成形的)一份对“女性”与“无能”的混同,成了“发自脏腑”的拒斥反应的对象:唯有男人——以他得胜的阳刚——才被认定是“力量”、创造力量的承载者。我也强调过:这类出自制约的出于本能的混同,全然不认常识与道理——当它所引发的观念与意象\dmark{}被带着那样强的信念与显然之力感到时,它们通常便被当作自身的辩白\origpage{710}。

然而有一个侧面,它随着“阴的葬礼”那条注的结语骤然以“闪光”显形,却还不曾被重拾。那条注里结束反思的几行如下:

“这不再是一个人的葬礼,也不是一部著作的葬礼,甚至不是一场不可容忍的异见的葬礼——而是‘数学中的阴柔’的葬礼——而或许更深地,在每一位为葬礼悼词鼓掌的众多与会者身上:是那被否认的、活在他自身之中的女人的葬礼。”

依我看,如今想来:这个侧面,在朋友皮埃尔本人的个案里,也曾或多或少被略而不提——尽管关于他,我第一手的事实并不缺!倘若这个侧面曾多少在场、或许被一位细心的读者感到,那必定是在字里行间:当时注意主要被“阴阳颠倒”侧面的种种角度吸住——(那个侧面,至少初看,似乎是我朋友在葬礼中的人格与特殊角色所专有的)。这份遗漏提醒我:日后(几天之内)我还须谈谈我朋友最后那次来访——10月10日至22日(10月21日那条注里提及,当时我许诺“几天之后”回头谈)。依我看,那将是最合宜的时刻,去检视“颠倒”的最后一个(?)角度——伴随着我朋友\textbf{这个人本身}之中原初阴阳平衡的颠倒。那又是一份\textbf{下葬}:下葬他身上某些原初的阴特征——在一份后起的、来占地盘的阳特征的鞭笞之下。我由此又以一道新的、更深的视景,站在那份惊人确认的跟前——它已不止一次强加于我\fnote{248}{(*)}{这份“确认”,头一次显形于“下葬”那条注(n\textsuperscript{o}~61)的反思里。}:以为自己下葬的是曾经的师父(而师父终究始终是一个朋友),可实际上,他亲手埋的,不是别人,正是\textbf{他自己}!

所以,倘若我再度回到“第三层面”或“底色层面”——那份“大会”别名“数学共同体”——方才引用的那几行便暗示:我在朋友皮埃尔的个案里感受得如此强烈的东西,或许对“每一位为葬礼悼词鼓掌的众多与会者”同样为真。依我看,正是这个侧面\dmark{}\origpage{711},还须稍加检视,我才会感到完全满意,才把葬礼图景的“底色层面”(连同前景)当作(暂时?)完成。

(12月25日)昨天,我借着平安夜,给自己放了个假,让自己痛快大干了一场:在注前一直待到凌晨三点过(下不为例!)。诚然,整个白天散给了别的任务;(重读完前一天的注)夜里只剩几个小时,倘若我想当天继续的话。一如常常:我终究甚至没能碰到我坐到白纸前时心里想着的任何东西!取而代之,我对葬礼“图景”的进度作了一点盘点,并突出了“第一层面”与“底色层面”中一个仍模糊的侧面:“\textbf{被否认的女人的下葬}”——她活在我的葬礼每一位与会者身上。

清楚的是:在这段引文里,“下葬”一词用作意象,指一份\textbf{否认}与\textbf{压制}的行为(或按通行的术语,“压抑”)。要谈否认、压制某样东西(在此,某样“活”在自己身上的东西),首先须确认:这份“东西”确确实实在场、“活着”(哪怕可怜地活着)。这里说的是每个存在者——无论女还是男——身上的“女人”:即其这个人身上由“女性”、“阴”品格的特征、品质、冲动或力量构成的那个“坡面”。奇怪的是,这个简单而要害的事实——在每个存在者身上,无论女还是男,同时活着“女人”\textbf{和}“男人”——至今普遍被忽视。我自己,也只在八年前——我四十七岁那年——才懂得它\fnote{249}{(*)}{关于这一点,见“接纳(阴的苏醒(2))”那条注(n\textsuperscript{o}~110)。}。

当然,“精神分析家”们早已“知道”它、谈论它。想必书很多,人人也多少听说过,一如我自己当年听说过。甚至,“人人”都乐于承认:这里头必有些真东西,只要说这话的是公认的内行、只要有为它写的书,等等。可是,“听说过”与“乐于承认……”,甚至读过一本、\dmark{}乃至十本\origpage{712}这方面的书,甚至(我斗胆说)自己写过一本、乃至几本——这本身并不意味着我们“知道”这件事;至少,并不意味着我们以一种比单纯死记现成套语更强、尤其更有用的方式知道它:诸如“弗洛伊德(或Jung,或老子……)说过……”。这类套语构成一份文化行囊、一张“有教养”者的名片——“懂”这“懂”那,甚至有时(配上文凭)是这或那的专家;就此,甚至可以承认它们有某种“效用”;确定的是:人人对自己这份在四面八方、在学校里、书里、“有趣的谈话”里攒起来的行囊,都宝贝得很,顶着逆风带着它,像一份亮闪闪而累赘的战利品,直到咽气。倘若我方才曾不敬地暗示:这份贵重的行囊“无用”,我的意思是:它对一件谁也不操心、甚至人人避之如瘟疫的事无用——即:自我的学习。或者换一种说法:这份行囊,对\textbf{承担自己的生活}无用——也就是说,对消化并同化自己亲身经历的实质无用,并由此,成熟、更新……

倘若我得以用几句话概括我这场关于阴阳的长反思的要害内容,那便是借着“提醒”这份“简单而要害的事实”——我方才恰恰提醒过。倘若有一位读者一路跟到这里,却还不曾以他自己的亲身经历感到这个事实——他身上有“女人”,哪怕他是男人;她身上有“男人”,哪怕她是女人——那么,他做了这番徒劳的努力来“跟着我”,却把时间花在给一份想必已然沉重的行囊,又压上一份新的分量:一份贴着“收获与播种”标签的分量。而倘若他是男人——即便他不属于那些葬礼与会者:在读我之前,他既不知情也不曾起疑——大可以打赌:他也在日复一日、在不知觉中,“下葬一个活在他自身之中的被否认的女人”(一如我自己当年也曾如此,且在我一生大部分时间里)。

一个男人“下葬”活在他身上的女人,一如一个女人“下葬”活在她身上的男人,有一千零一种方式\fnote{250}{(*)}{同样的话,也适用于“下葬活在他身上的男人”的男人、或“下葬活在她身上的女人”的女人:这类姿态,远不像人们想象的那么罕见。}——也就是说:去否认它、压制它。\dmark{}\origpage{713}“下葬”活在自己身上的东西,最常见的途径之一,是当那份东西在他人身上显形时,作出拒斥它的姿态或行为。这份拒斥,恰恰便是昨天在一个具体个案里我谈到的“发自脏腑的反应”。拒斥反应的发自脏腑的强力,其实\textbf{不是}(如我昨天似乎暗示的)因为被拒于他人的那份东西,顶撞了一套得到我们全盘而毫无保留之认同的“价值”总体。自知“强”的人,不会被“弱”的景象冒犯。反应的活力,恰恰相反,来自这样的事实:那份在他人身上被看见、“没有存在理由”的东西,\textbf{把我们自己牵连进去}。它像一份阴险的\textbf{提醒}——旋即被否认——提醒某件关乎我们的事:那件我们心底\textbf{知道}、却想对自己也对别人瞒住的事;一份提醒,由此便取一份无声而可畏的质问的调门。在这样的语境里,一份善意的宽容姿态——对着他人身上那份表面“毛病”——会显为一份危险的共谋自供,须不惜一切避开。而凭一份拒斥的姿态,相反,我们毫不含糊地与对方划清界限;总而言之,我们交出令人信服的保证(而首先,交给我们自己身上的内心检查官):我们自己白璧无瑕,我们始终合乎规矩、“道貌岸然”。拒斥反应,在是一份对某些\textbf{价值规范}的无条件的\textbf{效忠行为}(那些规范分判可敬与不可容忍)的同时,也是一份\textbf{象征性的下葬行为}:凭着它,我们身上那份“没有存在理由”的东西,被赶忙“归档”为一件“\textbf{不存在}”的东西。\textbf{至少不在我们身上}!

在这幅图景里,拒斥所取的形式——其变化无穷——依我看无关紧要。可以是那份被冒犯的拒斥,带着愤怒或厌恶的一切迹象;也可以是经由反讽、或“配量精巧”的轻蔑的拒斥。它可以以清晰不含糊的言辞表达,也可以仅仅被暗示:凭影射的或双关的话语,甚至无须言语,凭一个放在恰切之处的会意的微笑(或不笑……)。拒斥可以全然有意识,也可以只滞留在刚刚浮入眼帘的朦胧里,或躲进那份目光永不穿透的完全的阴影。

拒斥反应的强度,也变化无穷\dmark{}\origpage{714}:视那份“牵连”被感到是无足轻重,还是确实可畏。或许引发最强反应的“牵连”,是那些直接触及\textbf{性}的。这份极端的敏感,最近几代人里稍稍缓和了。可我确认:像爱欲冲动中所谓的“同性恋”与“自慰”(或说得好听些,“自恋”)的侧面这样普遍的东西,今天一如当年,仍引发强而有力的拒斥反应。至少,当你直面它的场合,不是一场关于罗马时代风俗、或深层心理学的“有趣谈话”,而是在你的日常生活里。甚至面对面,人们也很少谈起自己这个人身上这些性冲动侧面的显现(它们一般被感到是几处有点碍眼的“越轨”——至少这么说)。

在我这里关切的个案里,我出发脱离数学舞台之前所面对的拒斥反应,当然不曾有一份与我方才唤起的相当的力量。诚然,这份拒斥的对象——即“女性”的存在与做法,而人们本该“在男人之间”——确实有一份“性”的意味:取这个词更宽的意义,而不只是围着“屁股”及其余打转的言行的意味。我不怀疑:这份意味通常被感到——在无意识的层面\fnote{251}{(*)}{关于这一点,尤其见“阴的葬礼(阳葬阴(4))”那条注(n\textsuperscript{o}~124)。}。可它足够含蓄、足够间接,排除了任何粗暴的反应:超出对我作为数学家之“严肃”、“扎实”的一份单纯“保留”的反应。还要加上:我的“毛病”所处的领域——一份纯粹智性活动的领域——有助于给它一份相对平淡的外表:远(你上这儿来找什么呢……)远离男人-女人扭着腰肢跳肚皮舞那份令人不安、有伤风化的联想!但这并不妨碍:自我最初接触数学界(1948年)之后,还须将近十年,我的风格所引发的那些保留——在一个善好的小天地之内——才终于消失:至少,从我的视野里消失。可我出发之后,情形又变了:那份对我怀着善意、友情与尊重的氛围,骤然变了(在其后的六年里,我自己还不曾察觉)——被那个小天地所感到的一份“异见”、一份否认所变。

\asterisms

\dmark{}\origpage{715}老实说,我不敢肯定:这份氛围之变,是否真如我方才所说那般“骤然”。或者更好地说:我确认,我手头几乎没有事实,能使我对这份变化——1970年我出发之后如何发生、而1976年骤然(这一回,该用这个词了)摆到我面前——\textbf{如何}发生,形成任何想法\fnote{252}{(*)}{我提醒:那是借着我为使Yves~Ladegaillerie的论文得以发表所作的徒劳努力。这段情节,在“谁也拦不住进步!”与“棺材2——或锯成段的切割”两条注(n\textsuperscript{o}~50、94)里谈到。}。诚然,这整段时间里,我与那个我离开了的世界几乎毫无接触:本可让我感到某种“温度”及其演变的接触。对我清楚的是:在这份演变里,所有那些曾是我学生的人的小组、及其无可争议的头领Pierre~Deligne的姿态,扮演了一个决定性的角色。葬礼得以举行、引发葬礼的氛围得以确立,唯有凭一份“\textbf{全体一致的协约}”\fnote{253}{(**)}{关于“全体一致的协约”这份确认在反思中的头一次出现,见同名的(大写!)那条注(n\textsuperscript{o}~74)。}、一份毫无破绽的协约:它同时涵盖这场葬礼的“三个层面”:“继承人”(别名葬礼大司仪)、“共同继承人”或“近亲”的小组——由其他十一位“昔日的学生”构成——以及最后,“大会”(或许终究不是“全体”——日后还须回头谈……)。这份完美的协约如何设立、确立,对我仍属未知,或许将永远未知。眼下,我不觉得有兴致去探测它;我怀疑会有别人替我去探测(恰恰相反!)。

这提醒我:写上一条注“天赐的情势——或封神”时,有个问题擦过我心头:\textbf{二者之中},究竟是“大会”还是“披祭披的神甫”,代表了葬礼中运作的\textbf{主导}力量,而另一个多少算是它的“工具”\fnote{254}{(***)}{我提醒:在5月的反思里,在“掘墓人——或全体教众”那条注里,我曾意识到:我的朋友曾是一份毫无破绽的\textbf{集体意志}的“\textbf{工具}”。接下来的几行并不真正驳倒这份直觉,而是补充它:让“大会”与“披祭披的神甫”之间关系里的某种对称性保持开放。}。我当时没有驻足,甚至不确定这问题有没有意义——它分明像那道著名的先有鸡还是先有蛋!确定的是:两者之中任何一个——“神甫”也好、“大会”也好——都无法缺对方的助力来操办葬礼\dmark{}\origpage{716}。

可是另一个问题,依我看意义更清楚:两者之中,谁在这份工程里投入得更强。诚然,“大会”不是一个人,谈“它的”投入不很恰当。可同样真确的是:对我而言,这个人格化的实体,取得具体的形貌——凭十来个或二十个我熟识的\textbf{人}:与其中每一位,在一二十年间甚至更久,我保持着持续的、友好的关系。所以,当我谈大会的“投入”时,我具体想到的,是那些旧日朋友之中参与我葬礼的所有人的投入之“总和”。如此限定之后,依我看,这个问题便不再有什么修辞的成分。

对这个问题来到我面前的回答——不带一丝犹豫或怀疑的调门——是:在“继承人”的投入与大会的投入之间,\textbf{毫无公度可言};一如在一场普通的葬礼里毫无公度——何况遗产在继承人眼里越大(而大会里没有任何人能从中为自己捞到什么)、把他系于亡人的纽带(吸引的或冲突的)越强、在他的生活里扮演的角色越要害,便越甚。倘若这样的情势里有疑问,那多半只能来自亡故者近亲之中“共同继承人”的在场。(所以,这里说的是“第二层面”,而不是由大会的主体构成的“背景”。)在我关切的个案里,这些“近亲”与共同继承人之中,唯一一个其在我葬礼中所占的份额、可与主继承人Pierre~Deligne所占的相比者,依我看是Jean-Louis~Verdier:他扮演葬礼的\textbf{副司仪}。这份称呼并非凭空:在葬礼的进程中,我不止一次确确实实看到两位一同主持,配合得天衣无缝!可是,一如我在别处已写过的:撇开Jean-Louis~Verdier的某些公开行为不论,自我与他失去联系以来,关于他我知道得很少:太少,无疑,不足以对那份他与我的关系的来龙去脉、或他与那位显赫的“保护人”兼朋友的关系,形成一份多少周详的想法。

\subsubsection*{注~153}

\notehead{153}{}
\phantomsection\label{ch18:local:n61}

(12月26日)昨天的反思里,我试着精确化那份11月10日“闪光”般显出的直觉:在我葬礼的“每一位众多与会者”身上,我的葬礼代表着“活在他自身之中的被否认的女人”的象征性下葬\dmark{}\origpage{717}。当我在这里一次又一次说与会者的“每一位”时,这是一个稍嫌脱口而出的说法,或许最好不必完全按字面来取。至少,我深信:这份直觉,确确实实对每一位这样的人生效(而这样的人肯定很多):他们身上多多少少发生着那份对我特殊数学风格的“发自脏腑的拒斥反应”——那份反应,正是三天来我注意的中心。

另一方面,清楚的是:这样一份反应\textbf{不在}我的朋友皮埃尔身上,或至少,在我出发之前的五年里,那里非但没有它的痕迹,恰恰相反。正是我的数学进法风格与他自己的风格那份深刻的\textbf{亲缘},才造成了这些年里那份完美的交流;它也是我们在数学层面那份罕见的亲缘之缘由——那份亲缘,他与许多别人想必都曾感到,一如我自己曾感到。也正是这份亲缘,肯定是我这个数学家的人格、我的著作对他施加的那份\textbf{着迷}之缘由:不仅在那些年头(那时它“以正面”表达自己),也在其后的年月直至今天(那时它主要“以负面”表达,却同样雄辩\fnote{255}{(*)}{或者毋宁说:这份着迷,起初想必是“以正面”行的那份力(即\textbf{认同}于那个被感到相似的人之力)——在建立那份与我的暧昧而冲突的认同关系的那两股力之中。})。我不怀疑:倘若在那些最初的年头里,他对我把握数学事物的工作与进法,有过哪怕最轻微的保留、最轻微的不适,我不会觉察不到。

诚然,从那些年头起,我的朋友便竭尽所能,对外抹去那份在我身边为他所备的角色——哪怕只作为那个教过他、传给他一些有分量的东西的人,作为他从中取得工作所需的重要观念的人——更不必说,也抹去那份亲缘、乃至着迷的关系。我出发之后,便有对我这个人之否认的逐步加码:不仅凭沉默,也凭一份对我工作风格、对我引入的许多观念与想法的做作的轻蔑。我所知道的这类做作的\textbf{头一}份痕迹,定在1977年——借着“SGA~4$\frac{1}{2}$工程”\fnote{256}{(*)}{关于这一点,尤其见“两个转折”与“白板”两条注(n\textsuperscript{o}~66、67)。}\origpage{718}。

我没有试着亦步亦趋地追随这份加码的进程,也不觉得自己有兴致去追(一如昨天就一个切近的问题所说的)。

这份对一份与他自己近亲的进法风格、对一份他自己的著作由之而出的著作的否认,分明近于一份\textbf{对自己的否认}。方才,想着这份对我风格与著作的否认时(而我主要还沉浸在我1970年出发之前五年贴身数学接触的印象里),我本倾向于把它看轻,只给它一份某种意义上“战术”的意义:作为一份格外诱人的\textbf{手段},用来取而代之、满足对抗的冲动,借着抓住一份“天赐的情势”的运气。这确实便是三天前那条注“天赐的情势——或封神”(第~n\textsuperscript{o}~151)的调门。而我刚刚记起的这件事——我出发之前的年月里,没有对他自己风格或对我的拒斥姿态的痕迹——也正朝着这个方向,而不朝着昨天检视的情势的方向:即一份对“活在自己身上的女人”的否认(哪怕部分地经由对数学的某种进法)——一份\textbf{先于}葬礼之操办而有的否认。

但这并不妨碍:谁若拣了这样的手段,不管愿不愿意,\textbf{就得付它的价}。这份对某一份风格的“做作的轻蔑”,要奏效,便不但得对着别人演,也得、而且尤其\textbf{对着他自己}演。可是,人们无法当着别人、也当着自己,否认一份同时也是他自己骨子里的“风格”,\textbf{却照常实行着它},仿佛什么也没发生。这份对着别人的“战术性否认”,凭事物的逻辑,要经过一份对自己的否认、一份对自己一部分的\textbf{压制}——在此,即对他那份凭创造力量的原初本性而有的数学进法风格的压制。

这份确认,在此不是作为对一份事实的直接知觉之效而来到。它是一段简短反思的终点:动用已知的事实,从中得出常识的“结论”。我学会了提防这类结论(尤其在数学之外!),并只信那些事后由其他事实印证的\dmark{}\origpage{719}。可我在此恰好便利地记起:我曾被引着——以我对Deligne著作所知的措辞——确认:在他的著作里,找不到痕迹:他身上某些癖向(“阴”性的)——它们在我出发之前的年头里分明可见,我也在自己身上认得——在他的“之后”的著作里消失了。关于这一点,我在一个月前(11月26日与28日)的两条注“阴,作为仆人——与新主人”与“阴,作为仆人(2)——与慷慨”\fnote{257}{(*)}{即第~n\textsuperscript{o}~135、136两条注。还须加上第二条引注的子注(n\textsuperscript{o}~136$_{1}$)。}里表白得相当周详。这些事里最重要的,或许是一份谦逊:它让人看见(并描述,不怕显得傻气)那些极简单、极平淡的东西——此前无人肯对它们稍加注意。我自己献给数学的最好的东西\fnote{258}{(**)}{关于这一点,见前一条脚注所引的子注~n\textsuperscript{o}~136$_{1}$。},恰恰正是这一水质的。倘若我曾否认我天性的这份癖向——它可不是人人都喜欢的……——那么,我的著作、以及我那位最辉煌学生的著作的精要,都不会被写出来。这份癖向(或“倾向”)与另一份系得紧密,不然它的效应会极其有限。那又是一份谦逊、与“服务”的姿态:当要做的是认识并以精微的方式、面面俱到地描述那份被众人轻蔑的新东西时,不觉得自己的时间太贵,不肯为它花上十页,如果有必要(而不是满足于两行:东西在这儿——随你们怎么办!),甚至一万页;为它花上一整天(对一个并不缺别的苍蝇可拍的人而言……),甚至一生,如果有必要。

当我说到有待发现的“新的世界”时——调门或许有点倨傲——我说的不是别的,正是\textbf{这个}:看见并接收那份显得微末的东西,怀上它、滋养它九个月或九年,需要多久便多久,必要时在孤独中——好看见一件强健而活的东西发育、绽放:它本身便为生育与受孕而造。

倘若这份癖向——人们可以称它“母性”的——今日成了嘲弄的对象,那是“为了”那些被感到“阳刚”的姿态的利益:它们只容许一种可能的数学进法:“肌肉”的进法,把“肚肠”排除在外。“真正的数学”,又号称“硬数学”(或“硬邦邦的数学”),与那些(不怎么可口的)“软数学”(或“绵软的数学”——不说“瘫软的”,呕!)相对——\dmark{}那是十页五十页排得密密麻麻的证明,证明那些竞赛型定理(困难尽人皆知,不然便不算游戏!),万弹齐发——把一切“众所周知”的理论与观念、把左右可得的全部事实,统统打出去。\origpage{720}至于“木柴”,只消它在那儿;它在那儿就是为了这个!而那些耐心开荒的人——那些播种、栽种、施肥、剪枝的人,贯穿季季年年,好叫这些挺拔的大树之林生长铺展——它们如今待在那里如此合宜(那里曾是密不透风的莽丛),以致人们会以为它们自创世以来便在(大约充当背景布景,充当“万能木柴”的储备……)——这些人,只配下些大河文章(甚至大河书、大河书的大河丛书,倘若找得到疯到肯印的出版商),还外加读不通——他们是“软数学”的落伍者,不说“瘫软的”——不过,人再阳刚,礼貌总还是要讲的……

随着这段漂亮的咏叹,我忽然觉得回到了这场关于阴阳的长反思的出发点——10月初头一条注“肌肉与肚肠(阳葬阴(1))”(n\textsuperscript{o}~106)。确确实实是同一场下葬:踏着阅兵的步伐、吹着军号,下葬那份“女性”的东西——由铁臂,别名钢脑,别名超人的男性轻蔑所掩埋。这场下葬当然不只在数学的小天地里;它的分量超出一切个案——尽管个案恰好能多少贴身地让人嗅一嗅它的气味。而这份气味,确确实实是葬礼教给我的主要功课之一:在其中,我扮着未老先亡的亡人。

当我把注意的场子收得更紧,专注我朋友皮埃尔所扮演的特殊角色时,我在葬礼里看到又一层意义。我在其中辨出的,又一次是一场\textbf{颠倒}。一如我昨天预告过——没想到这么快便回头——这里,不再是一份系于他人的\textbf{关系}之中(真实的或虚拟的)的颠倒,而是一场发生在他\textbf{这个人本身}之中的颠倒。它不被为其自身的功德所求(或许作为一份“疯狂的欲望”的对象……),也不再限于纯然的象征(一如在一场漂亮的变戏法的收场处:那个自感“侏儒”的人,并不因此少觉得自己是侏儒,仿佛他并不曾刚说服自己已变成“巨人”……)。这是一场颠倒——我不说不可逆——但至少全然\textbf{实在}。它出发自一份和谐的平衡状态:“女性”与“男性”的创造冲动\dmark{}\origpage{721}的平衡,带一份女性的主导音符。它抵达一份战争与压制的状态:种种\textbf{姿态}与\textbf{架势}(利己的,一如一切姿态与架势)——打着“阳刚”的旗号——执拗地压制着那份\textbf{创造力量}:它被嘲弄、“象征性地”下葬,化作一具怪诞而瘫软的模拟像,长着“超级雌性”的面孔。

用不那么有层次的、但更形象、或许更触目的措辞:一个“\textbf{女性}”的存在——精微而强健、柔韧、\textbf{活着}——经一场永不停歇的变戏法,变形为一个“\textbf{阳刚}”的存在:不可摧毁、僵硬、而且\textbf{死}了。

\subsubsection*{(f) 布景——或“第二天性”}
\tocsubsub{(f) 布景——或“第二天性”}
\label{ch18:local:h67}

\notehead{154}{}
\phantomsection\label{ch18:local:n62}

(1985年1月1日)五天过去了,被种种营生占去。年尾正是写积压了几周几月的信的时机,更不必说几张贺年卡——回那些圣诞节前后收到的。还须做的是:堆肥——两个多月前收进的粪肥,加上园子里与开荒下来的植物废料、或从市政垃圾场运回的,要垒起几座肥堆,好叫开春时园子里有现成的好土。因地有坡度,还得再垒一级台地,靠着已有的那一座——那一座是给生活垃圾“日复一日”堆肥用的。这一切之下,我几乎没找到时间在注上工作,除了后勤。我极其细心地重读——此改彼修——了“主人们与仆人”部分以来(即自11月24日那条注“颠倒(3)——或阴葬阳”(n\textsuperscript{o}~133)起)的整场反思,补上为最近十五天的注预备好的脚注。主要是要有份可送去打字的手稿;但撇开实际问题不论,这次重读是有用的:重新取得近四五周反思的总览。一如在一场长程的数学反思中:当反思的“时刻”日复一日地被一份高度聚焦的注意的强光柱照着,反思的“线”、及其在数周乃至数月里走过的曲折轨迹,便倾向于在路上失落、没入半明半暗的混沌之中。我说不清:这在一切长程\dmark{}的研究工作\origpage{722}里是不是普遍现象,或者它系于我生活中那份“过去之下葬”的系统性机制——我已有机会暗示过它\fnote{259}{(*)}{这份机制,是在我童年发生的那场“翻转”的时刻启动的——我把它定在1936年夏(我九岁那年)。这一“我”结构化中的关键情节,在“超级父亲(阳葬阴(2))”那条注(n\textsuperscript{o}~108)里、及其子注~n\textsuperscript{o}~108$^1$里被暗示。}。无论如何,在一场长反思的日日夜夜、乃至数月之中,我身上有与它先前诸阶段的失联,表现为工作中一份渐增的不适。这份不适,终于由一场对刚完成的工作或多或少的深入回顾所解决:由它,那份逐渐松弛的联系重新接上。我观察到:这些“回顾性停留”在我的工作中扮演一个重要角色。每一次,我都重新启航,帆上新风,卸下了那份“不适”——它曾向我标示:对我进行中的工作在时间中的\textbf{连续}之整体知觉的一份渐进失落。在我的数学工作里,这样的回头并不罕见——甚至不妨说是通则:它引我把已做的工作彻底重想,并在一份新的视景里,既看已做的工作,也看将做的工作\fnote{260}{(**)}{关于“偶发的回顾”在长程工作中的角色,类似的反思,亦见“回顾(1)——或一幅画的三联”那条注(n\textsuperscript{o}~127)的第二部分,尤其其中所系的脚注。}。

但无论是一份数学工作,还是一场对我生活的禅观,我所谈的那份“不适”,始终是一份仍不完美的领会的征象:不仅(理所当然)是对尚待做的工作的领会,也同样是对已逝工作中所做之事的领会。这份不完美,实际上绝不归结为对反思各阶段及其时序的记性失灵(在一份数学反思里,这些是相对次要的方面:那里注意的对象是一份数学情势,它本身与研究者的心理特质、与这场检视的波折无缘)。它依我看,更是一份\textbf{太一}之缺的征象、一份不够充分的整合:把反思诸阶段之果实的那一整批部分领会,整合起来。这些部分领会,同样不完美、甚至属于假设,只要它们不曾被整合进一幅总景——在那幅总景里,它们相互照亮\dmark{}\origpage{723}。再用一次\textbf{拼图}的意象:对一份未知实质的探查,类似于拼一幅拼图的工作——但拼块不是预先给定的,须在工作进程中一一发现。更有甚者:每一块出土的拼块,起初只以一份含混而近似的形式显形,甚至相对于那个尚且未知的“正确”形状而言,被粗粗地歪曲着。反思的“局部”工作,在于把拼块一一辨出,并试着多少猜出每一块的轮廓:主要凭对所检拼块内部融贯的掂量、或凭对它与其他被预感到相邻的拼块的融贯掂量。可是,这些拼块的每一块,唯有当它们被拼进那幅它们由之而出的、尚且未知的总图,才显出其真正的本性、其精确而最终的形式。我所谈的那份“不适”,便是向我发出信号的:面对一堆被完美定位、却以一份或多或少的无形堆积着的拼块,该是把它们终于拼起来的时候了——或者,倘若(或多或少的局部)拼装已有,那是它仍太过零碎、或它歪着,须整个重来。要找到\textbf{那个}对的拼装,我撞见拼块的时序,无疑常常是次要的。可是,把拼块一一重新拿在手里(顺手便按那个顺序),怀着那份知道它们必定拼合、并等着它们各自落位的心境——无疑是工作中不可或缺的一步:为的是终于看见它们果真拼合。

上一条注(六天前)的“结语”,试着以言语圈定我心中一份强烈的印象:一场\textbf{变形}——它或许在朋友皮埃尔身上、在我离开数学舞台之后的十五年里,岁月流转,完成了。这些年里,我曾在此地彼处知觉过它的零散征象,它们有时叫我目瞪口呆;却(就我记忆所及)从不曾有任何时刻停下来,去形成一份对发生之事的\textbf{总}印象。得说:尽管我嗅出某阵“风”、以及我朋友在其中扮演的一份特殊角色(尤其借Motif的下葬——我对此只是朦朦胧胧地有所察觉\fnote{261}{(*)}{(2月30日)关于这份感觉——它长期保持着未成形、弥散的状态(直到去年4月19日发现“葬礼的全部辉煌”为止)——我尤其指出《收获与播种》第一部分里若干偶然的暗示(写于去年2、3月):关于\textbf{Motif}观念的命运,尤其见《导言》4(“一场追寻显然之物的旅程”)与“做梦人”一节(第~n\textsuperscript{o}~6节)。这份感觉的表述,在第一部分末节“一段过去的重量”(第~n\textsuperscript{o}~30节)的最后几页里大大清晰起来,自那个段落“我几乎可以把那封‘致……的信’(读作:Daniel~Quillen)视为……”起——它构成反思中的骤然转折。由这一天的反思最后阶段催生的头几条“注”——尤其月末写的双注“我的孤儿”与“对一份遗产的拒绝——或一个矛盾的价格”(n\textsuperscript{o}~50、51)——在相当程度上,给先前以弥散状态感到的东西“做了个盘点”:关于我的数学著作之待遇、关于对着它与我这个人吹来的那股“时风”。关于这份“弥散之感”就Motif所取的特定形态的一份描述,见“墓”(n\textsuperscript{o}~71)与紧随其后的“一脚踏进旋转木马”(n\textsuperscript{o}~72)两条注。}),对那场盛大的下葬——我的朋友\dmark{}\origpage{724}正手法老到地编排着的、对我这个人、对我全部著作的下葬——我远谈不上起疑。正是这场下葬在过去一年里的逐步发现,终于成了一记够强的\textbf{冲击}:它推动了我身上的一份惰性,并激励我终于“落定”在一份仿佛没入遥远过去之雾的情势上。所以,10月我朋友的来访,我也是以一份与当年会面时那份有点儿“例行公事”的心态全然不同的心态接待的:一份愕然惊诧的注意。正是在那次来访中,出现了这份印象、或毋宁说这份骤然知觉——它想必早已在场甚久,而我一直乐于无视:那份“变形”的知觉——恰恰是我在上一条注里经由另一条路径重新撞上的那个。倘若我又一次找回这份印象——这一次,经由对我朋友数学著作之所知——那肯定不是凭天大的巧合,而是被近两个月来与他这个人的直接接触所教给我的东西所引导。这份变形印象的确凿之力——它归结于一个“阳刚”的、不可摧毁的、僵硬的、而且\textbf{死}了的存在——当然不会作为一场比较、拼合事实(或其他性质的局部印象)的反思的终点而来,而只能凭一份直接的亲身经历而来:它一直未被说出。而这份经历,就在此刻,依然未说\fnote{262}{(*)}{(1985年2月30日)此刻它依然未说,尽管我终于写出了我朋友来访的记述——在“完成的义务——或真相的瞬间”那条注(n\textsuperscript{o}~163)里。}。

在上一条注里,我写道:这场“颠倒”(在我朋友这个人本身之中)、或这场“变形”(借用出现在“结语”里的说法\dmark{}\origpage{725}),并非“为其自身的功德所求”——还加了括注:“(或许作为一份‘疯狂的欲望’的对象……)”(即那份颠倒的欲望:“神经中的神经——或侏儒与巨人”那条注里谈到的)。可是,次日重读反思诸注时,我不再那么有把握了——也拿不准:我那份\textbf{把}我在葬礼中辨出的这两场“颠倒”\textbf{对立起来}的故意,是否真有根据。毕竟,在侏儒与巨人的意象里,“巨人”体现(如我不止一次强调的)着“阳刚”的价值,“侏儒”则被“阴柔”的非价值压垮。而即便这幅意象\textbf{坐落}在我朋友这个人之外——它被贴在他的关系上,对着另一个人(在此即我)——它在我朋友这个人之外,却毫无“客观”的存在:它反倒是——那场正发生在他\textbf{自己}身上的冲突实在——贴着外面的投影(贴在他的与某人的关系上)。换一种说法:侏儒与巨人的意象,显为那场真实\textbf{冲突}的象征性\textbf{布景}——那场冲突发生于比意象所居更深的层;而它不多不少,正是他这人身上“阴”“阳”两个“坡面”之间那场永世的\textbf{冲突}。

一份内在冲突的这样一份\textbf{外化}——它必须严格保持隐蔽——正属于无意识用来“清场”的少数尽人皆用的手法之列:尽可能“排去”那场原初的真实冲突,用一个显得更“可接受”、或至少不那么扰人的冲突取而代之。在此,所选的那幅避雷针意象本身仍然无意识(至少我推断如此);甚至,我倾向于相信:它被圈定在无意识中相当深、却比真实冲突之认识更接近浮面的层。(后者不多不少,正是“两面认识”的“正面”——“两种认识——或对认识的恐惧”那条注(n\textsuperscript{o}~144)里谈到的。)

这便暗示:上一条注括注里重提的那份“疯狂的欲望”——“\textbf{做那个巨人本身},或至少,\textbf{冒充是他}”——那份欲望,不多不少,是关于一场他自己身上“变形”之欲望的“外化的”\textbf{转写},就着侏儒与巨人的避雷针意象;一场变形——即使不真实,至少表面——在其中:他存在中被感到不可容忍的那份主导——那份“阴”调门的主导(被感到“软”、可鄙)——将被“颠倒”、变形为“阳”或“阳刚”调门的主导(被感到“英雄”、唯一值得艳羡)\dmark{}\origpage{726}。这两份欲望,远非在内在本质上有哪怕些微的对立,如今在我眼中反而是不可分割的:一份是另一份的影子、其\textbf{象征而可触的表达}。至于那场“变形”——我朋友来访时终于知觉到的(迟做总比不做好!)——它此刻显为那份“疯狂的”而强制性的欲望之\textbf{实现}或兑现;兑现,不是凭一份天赐的恩典的干预,而是作为“东家”那份执拗意志的长效效应:他要“校正瞄准”,照一些借来的特征把自己\textbf{重塑},并把这些特征强加给匠人-孩子(自然,此类“东家”式的操作,从不咨询匠人)。

上一条注里,我强调过这场“颠倒”(或这场“变形”)的\textbf{实在}品格。此刻我更清楚地辨出这份“实在”的本性与限度。它是一个\textbf{架势}的实在:竭力照一个被感到为理想目标的模型自我塑形。对模型的选择——即所采取的架势类型——无疑远早于我们的相遇。但依我看,在我们相遇时及其后的岁月里,投入这份架势的能量与耗散仍微乎其微。我朋友身上的这份投入之规模的骤然、严苛的变化,我相信,是由我的离去所创造的“非同寻常的”契机:先是从我的机构离去(我的朋友一夜之间必定在自己眼中显为偷偷地\textbf{取代了他的“对手”}),旋即又是我从数学舞台离去。第二个更重要的实在方面是:凭一份无度的投入,这份架势终于确确实实变成“\textbf{第二天性}”。正是在我们最近的会面中,我知觉到了这份“第二天性”。它压着一份巨大的惰性——一如我自己的情形也曾如此。这并不妨碍:在我这里,一场更新遂曾发生;而我身上发生了,并不削减我朋友身上那股反对他自身更新的惰性。

在他身上渐渐确立的这份“新”实在,并没有“解决”他身上的冲突——正如一个邻国占领一个国家,并不“解决”冲突。毋宁说,我朋友身上的冲突被“冻住”在某个“力量对比”之中;而很可能,它将如此直至生命的尽头\dmark{}\origpage{727}。无疑可以说:自我的结构——即行为机制——确确实实改变了,有时以触目的方式。然而,这类由“东家”的意志强加的变故,不改变原初的本性——匠人-孩子的创造力量的本性。它们只类似那些强加给匠人的刑具:匠人得想方设法,照样干活——在“东家”多疑的目光之下;当东家不干脆把工具从他手里夺走、好教匠人知道该怎么做的时候!

这并不妨碍:企业照转、照赚,而东家,大体上,满意。氛围确实够糟;可像大多数东家一样,他皮厚,不为所动,只要回报还好就行。

\subsubsection*{(g) 另一个自己——或认同与冲突}
\tocsubsub{(g) 另一个自己——或认同与冲突}
\label{ch18:local:h68}

\notehead{155}{}
\phantomsection\label{ch18:local:n63}

(1月2日)自12月24日那条注“否认(1)——或提醒”(n\textsuperscript{o}~152)以来,已有一周多了:我一直觉得,葬礼图景的前景层面,我差不多做完了。结果不然——已经连续三次,我得回到这个那个看着不太清楚的点上:大概只差三句话,好给最后一个ｉ点上最后一点。可每一次,这个“最后一点”都占了我整整一晚的工夫,而原来那个看着“不太清楚”的东西,其实一直相当晦暗;回头给它找它自己的光照,绝非奢侈。我预感今天也不会例外:我打算回到那个(最后一个?)点——它曾在下一条注“否认(2)——或变形”(n\textsuperscript{o}~153)里被一笔带过。那是我扮演“过继父亲”角色的关系独具的侧面之一:我朋友对我的“暧昧”\textbf{认同}。这一侧面,在引注的一条脚注里以三四行唤起。当夜它不再被谈;可次日,重读前一天的注,我感到必须回头谈它。昨晚接续反思时,我本打算从这里接下去;可最终,是前一场反思悬着的另一个“最后一点”占住了我,直到深夜。

在《收获与播种》中我屡次被引着记录:在某某朋友或学生的关系里,扮演收养的父亲或被收养的父亲\dmark{}\origpage{728}(的侧面)——每一次,都是借着这份关系中冲突性特征的显形。于是,无任何故意,我注意与强调的中心,总是这份“父性”意味的关系中的\textbf{冲突的}侧面。我分明感到:在这样的关系里,总有一份或多或少的\textbf{对父亲的认同}成分;唯有一点保留:这份认同有时取“负面”的形式——凭与一份被否认的父亲影像的“负面”(或反面)认同\fnote{263}{(*)}{尤其,这曾是我三个儿子与我的关系的情形——他们,这一回,全然不是“被收养的”,更不是“收养者”……}。这份认识留在背景里,不以可见的方式介入反思,却带来了它的一份贡献:对这样那样关系的一份弥散的把握、一幅尚模糊的轮廓之形成。我记得只有一次——我以为——以一般措辞、顺着认同的方向,在“敌对的父亲(1)”一节(n\textsuperscript{o}~29)的结尾表白过:

“……这便是那同一个对父亲的冲突原型之重现:那位父亲,既受钦慕又遭惧怕、既被爱又被恨——那个须要与其交锋、战胜、取代、乃至羞辱的\textbf{男人}……却也是那个人们暗中想成为的人:想夺他的力、据为己有——一个\textbf{另一个自己},被畏惧、被仇恨、被逃离……”

无须多言:在这些行间——写于一场“对我数学家往昔的回顾”——倘若有一个具体的个案在动笔时指引着我的笔,那便是与我的“隐蔽继承人”、那位不说自己名字的昔日学生的关系:Pierre~Deligne——尽管在写的那一刻,我对那场由他亲手编排的盛大葬礼,在意识的层面至少,毫无察觉!此刻复现这几行——写于九个多月前——我已惊于:它们仿佛预示、并(可谓)召唤着侏儒与巨人的意象——那意象仿佛正是为给这幅直觉以可触的形式而成形、落实的。然而,对我毫无怀疑:这幅意象绝非在我这个编年史作者兼探索者身上成形,而分明在我朋友自己身上;它也正是从非他莫属的某人那里取得\fnote{264}{(**)}{关于这一点,见“神经中的神经——或侏儒与巨人”那条注(n\textsuperscript{o}~148)的最后一条脚注。}!

冲突性的认同,清楚地显现在“那个人们暗中想成为的人”\dmark{}\origpage{729}这几个字里;更强、毫无含糊:“一个\textbf{另一个自己}”。在侏儒与巨人的意象里——它12月18日来到我的笔端(在那条注“神经中的神经——或侏儒与巨人”,n\textsuperscript{o}~148)——谈到的是那份“疯狂的欲望:\textbf{做那个巨人本身},或至少,\textbf{冒充是他}”——这几行似乎正是对方才引用的“那个人们暗中想成为的人”的回答。可这一次,我在这里停住(一天的本分一天尽!):离“另一个自己”还差一步——后者九个月前作为天经地义的事信手而来!诚然,这一次,既然是一份“看物件”的工作、在一个十足具体的个案里,便须迥然不同地小心、审慎——不同于那种语境:在那里,人们(不动声色!)作态抛出一份一般品性的断言,不涉及任何特定的人……

但想想:诚然,对那个渴望\textbf{象征性}满足的无意识而言,这确确实实只差一小步——它大可以凭自己制造的意象付这部账——从那份“疯狂的欲望”(而且显然有一股相当大的力)想成为这个那个,到与那所望之物的\textbf{认同行为}。要使这份认同——无论多么无意识——多少可信,并使它所带来的一切满足可以被带着最少的安全感来赏味,还须:它有某些(在此)被认同之人身上的“客观”相像特征的保票。我推想,在我关切的个案里——我朋友与我的关系——第一份“客观特征”、足以助成一份相似之感与一份认同行为的,是渗透在我们共同情人——数学——的进法之间的那股强亲缘。这想必便是那份“以正面”行之的力、那份“\textbf{认同于那个被感到相似的人}”的力——今天反思开头引的脚注里顺带谈到的。

可是,一如我在关于朋友与我的关系反思中已屡次提示的:自这份关系的最初几年起,他不曾漏看我自童年起所扮角色中的“超级阳”失衡侧面——那早已成为我的“第二天性”。我无从说:在意识的知觉层面,我朋友是否清楚地分辨了我这个人身上这两个截然不同的侧面。(我倾向于怀疑。)无论如何,我企业里“东家”的超级阳侧面,想必在他身上唤起了\dmark{}\origpage{730}两类截然不同的反应。一类——也是直到最近几个月我唯一知觉到的、也是他身上唯一有意识(我推断)的——偶尔以一份略带忧伤的惋惜姿态表达:我曾唤起过的;那份姿态从未离开友好或温情的调门。另一类反应,细看之下,自身也是“暧昧的”:由两份方向表面相反的成分组成。一份,“正面的”,朝着对我这个人的毫无保留的\textbf{抬举}:作为“英雄”价值、“超出常人的”价值的化身;那些普遍公认的价值固然,人们在少年时代像呼吸空气一样吸收它们,可他童年的直接近旁大概没有给他提供任何多少令人鼓舞的“模型”。这份成分——一如前面谈过的那份(全然另一性质的)\textbf{亲缘之感}——朝着对我的\textbf{认同},不带任何对抗因素。而那对抗因素,则进入另一份成分,或更好地说,进入这份我方才描述其“正面”的认同的另一面(或“\textbf{反面}”),它对我始终更费解。想必正是在这里,那份“父性”角色——我朋友凭我符合某个号称体现这类价值的理想“轮廓”而派给我的角色——扮演一个关键角色。当我试着摸黑探测——凭我手中极少、极细碎的材料——这份对一个“过继父亲”(特征十足的“超级父亲”! )的认同之强烈对抗内容的深层缘由时,两周前,我撞上了一份说得通的、但仍是假设的“脚本”:在12月30日那条注“缓期中的宿怨——或事物的回转(2)”里。

不是在这里回头谈那份脚本。依我看更有趣的,是把“侏儒与巨人”的意象(它刚在前天那条注里出现),放回这份我朋友对我的冲突性认同的光照下来看。于是一下子显出:意象中的两位主角——侏儒与巨人——\textbf{都不过是他自己};或者毋宁说,\textbf{他自己身上两个截然不同的侧面}。“侏儒”代表被我朋友感到为\textbf{原初}、“不变”的存在侧面:扎根于他的童年,远及他的记忆、无疑更远……它也是被感到为他这个人平淡、无足轻重、乃至可鄙的侧面。它是\textbf{被否认的}侧面,并由此,也是被感到“无可救药”、“不堪重负”的侧面:他存在的\textbf{耻辱的}、可鄙的\textbf{一极}。“巨人”相反,代表那个令人眩晕的\textbf{理想}:人们绝望于永远达不到\dmark{}\origpage{731};至多,有望些许相像——不惜自欺欺人、也欺他人,凭一切可用的手段。其中一个手段,便是取代那个显为这份理想之可羡化身的人,凭一切想得到的方式“证明”自己优于对手。至于巨人本身,此刻显为与对手和父亲有别:他是\textbf{被拔高、被吹大的一极:自我的理想、英雄的一极}。“东家”的至高犒赏,便是这一切:足以滋养那份错觉——人们\textbf{就是}这理想的一极、那份贪得无厌地自我膨胀的精神的投影。可正是这份对犒赏的饥渴,泄露出几分不安——“一份深深埋藏的怀疑”——它告诉我们:当事人“在心底最深处,并不上这些重要性与‘价值’的虚假征象的当”……\fnote{265}{(*)}{引号内的引文,出自“他人之无谬与对己之轻蔑”一节(第~n\textsuperscript{o}~4节)。}

在心灵的较浮面层次,这些“虚假征象”\fnote{266}{(**)}{这些征象尽管“虚假”,却常常终于铸成一份固若金汤的“第二天性”——“不可摧毁”(借用“否认(2)——或变形”那条注(n\textsuperscript{o}~153)里结语的说法)!}却属于方才所谈的“(或多或少)客观特征”之列:它们号称使一份对理想模型的认同行为“可信”(无论那模型保持“巨人”的非个人形式——活在他自身中的无面孔者——还是取敌对之父、对手的熟悉面容)。

\subsubsection*{(h) 敌对的兄弟——或传递(2)}
\tocsubsub{(h) 敌对的兄弟——或传递(2)}
\label{ch18:local:h69}

\notehead{156}{}
\phantomsection\label{ch18:local:n64}

(1月3日)昨天下午,趁等朋友登门的个把钟头空档,我翻了翻C.G.~Jung的自传——一位朋友凑巧刚带来。读到的寥寥几页,便强烈地钩住了我。这是我头一次手里拿着一篇Jung的文本;至今,我对他只有一份极模糊的概念:一个弗洛伊德的叛逆弟子——(凭零零星星传回我的回声)他懂得把奥秘那明暗微妙的运动,重新引回主人的笔直小径里。到此为止,相差无几。此刻,我感到一个像你像我一样活生生的人;而且,他不浪费时间绕弯子;尤其:一个直奔真正问题的人——那些他凭自己的光感到要害的问题,并且不满足于(当偶发的疑问古老如世界时)饱学之士那现成的答案\dmark{}\origpage{732}。

“传记”一面(那是要发表的)当然特别令我有兴味:因为我正写着的这些注,颇有几分像一部传记——而且与Jung的精神十分相近:外部事件始终从属于内在历险,它既是内在历险的揭示者,又是偶发的刺激者。令我震惊的是:Jung直到83岁才写一部自传(或更确切说,才为一部传记贡献他的一份),而且,尤其:在他一生的任何更早时刻,他从不曾费心深入地检视他自己的童年。在我看来,对弗洛伊德的弟子们,这本该是理所当然:要熟悉无意识之路,头几件事——若非头一件事——便是探一探他们自己身上的路!我甚至毫不怀疑:一份号称的“对无意识的认识”,若限于大学课程里学到的东西(哪怕由弗洛伊德本人这样一位显赫的大师讲授)、外加对若干“临床病例”的分析,则仍是一份未经整合的知识、一份零碎的、“死的”知识——一份本身不提供、也不促进对自己、对他人、对世界之理解的知识。

但同样真切:对自己这个人的勘探,凭其本性,不能列入一份制度化的“课程”之中——正如(且说一个“病人”)一份被扰动的心理平衡,就其根本而言,不能凭任何一位“学家”的干预为果实——他不过是施展万金油式的技术而已。那份“被扰动的平衡”,绝不限于神经崩溃或神经官能症那个社会不可接受的阶段;它几乎人人身上都查得出(程度多半\textbf{更深}而非更浅)。心理学家们自己(或民族学家、社会学家、其他“学家”),无论何种派别,也并不比旁人更豁免!而一份对被扰动平衡的真正恢复,绝非一份介入第三方“医疗行为”的本性。它是\textbf{当事人自己}的行动、非他人——\textbf{一份爱的行动}——他既可做、也可不做。它不是心理机制那无情运转的结果(无论有没有那位心理机制专家的干预),而是取这个词完整意义的一个行动:一份\textbf{创造},一场\textbf{重生}。

在写完上面那句断然武断的话之前——关于那份“号称的\dmark{}\origpage{733}‘对无意识的认识’”——我便意识到:语境会使人显得多么狂妄。对Jung的作品一无所知(前面刚谈过他),我竟仿佛把他、连他的“号称的”无意识认识,一起毫不客气地打发掉——就因为他显然不曾(在83岁之前)费心探一探他自己那无意识生长的土壤。可我倒猜:读他的传记,将会显出:即便不曾致力于这样一场“勘探”,Jung想必另有\textbf{别}的接触他自己无意识的路径(那些路想必自己也长期无意识)——由此,那被指摘断言的先决条件,想必不适用于他。

有另一桩全然不同的事,翻看词汇表时叫我错愕。在“四元性”条目下(按:这是法文版),Jung坚持数字四的“总体化”品格。直到十年前,我对数字的哲学或“神秘”运用还十分抵触——任何这种方向的思辨或谈论,在我看来都是胡话、孩子气、“变戏法”(德语叫Hokuspokus,指廉价的魔术)。我后来学到的关于《易经》(或“变化之书”)的那一点,使我没那么武断了。昨天,我把归给数字四的“宇宙”品格,与写“阴阳之钥”时自发形成的那份分组联系了起来:一般四或八条注合成一“包”,归在一个共同标题下。第一个“包”只有一条注;诚然——可(我在写完第六组“阴阳数学”时曾带着满意记下:那组有七条注,而不是八条)把它并入其后一个显得它最自然归属的组,又凑成八条(7+1=8),便又是四的倍数。这个“样式”一直延续到如今:最近完成的组,是第十组“暴力——或游戏与刺棒”(156$_{1}$)。得说:自第七组(“阴阳的颠倒”)起,我便任这份不召自来显出的“样式”引导——却从不曾给它寻、给它设一个“意义”,除了一份形式上的“数学”规整之感,被感到和谐。

这让我记起我写过的唯一另一篇堪称“宇宙”主题的文字——依旧以生活中的、创造行为中的阴阳动力学为中心\fnote{267}{(*)}{即《乱伦颂》——在注~n\textsuperscript{o}~43里谈到过(那里引的是“非古鲁的古鲁——或三足马”一节,第~n\textsuperscript{o}~45节),而主要在“行为”那条注(n\textsuperscript{o}~113)里,原书pp.507--509。亦见“事物的动力学(阴阳的和谐)”那条注(n\textsuperscript{o}~111)的开头。}\dmark{}\origpage{734}。这篇文字,表面上无任何初衷、也肯定从未费力,便按一份严格的数字编排归了组。我忘了它是怎样:方才一看(好奇者恒好奇!),显出:七节“判词”,每节四“诗章”。于是,这又是一份按四分组。诚然,判词的数目是七,不是四的倍数——于是按Jung的标准,“总体性”的品格对整部作品并不成立\fnote{268}{(**)}{这部拟议中的作品(顶着那个挑衅性的名字“乱伦颂”)本应含三部分(“天真”、“冲突(或坠落)”、“解脱(或寻回的童年)”),唯有第一部分做完了。这里说的就是它。},而只对构成它的七节“判词”成立。可这里我又有得脱身:那部著名的“诗集”,还配有一个天赐的“尾声”(更不必说一章没完没了的“序言”——我庆幸放下了),于是又有7+1=8——得救了!

该回到昨天的反思,从我中断的地方续起。我试过以我朋友对我的认同,理解他身上的侏儒与巨人意象。显出:“侏儒”与“巨人”代表(或“\textbf{布景}”着——借用昨天前那条注的说法)我朋友这个人身上的两个“\textbf{极端}”(我是说:“东家”设立为“极端”甚或更准确说是“两极端”的东西):一个“耻辱而可鄙的极端”,另一个“理想的、英雄的极端”。老实说,借一份重音与光照的差别,我在这里与前一天对同一力量意象找到的解释相接了:前天那条注“布景——或‘第二天性’”(n\textsuperscript{o}~154)。那里谈的是对那份由东家、“我”在存在的“阴”“阳”两个坡面之间设立的冲突的“布景”。这份对原初冲突的表述——以两个“坡面”的措辞——会对应一份未被歪曲的冲突认识;而我深信:这份认识,确确实实存在于(绝非不可触及的)心灵深层。昨天那份以两个“极端”的措辞的表述,代表的是一份\textbf{被歪曲的}冲突视景——被东家的一份故意歪曲:他抬举其中一个“坡面”\dmark{}\origpage{735},把它做成一个理想的、英雄的“极端”,贬抑另一个,把它做成另一个与前者相对立的极端:一个耻辱的、可鄙的极端。我猜,这份居间的意象,居于较浅的层、居间的层,或许与那幅外化的意象、那幅更接近意识表面、并与表层部分交叠的“侏儒与巨人的布景”部分共居\fnote{269}{(*)}{这份关于侏儒与巨人意象的推想,当然,系于那幅意象在其结语里如此明白的表达——在Pierre~Deligne亲自所写的那份传记小记里(暗示已于“神经中的神经——或侏儒与巨人”那条注(n\textsuperscript{o}~148)的最后一条脚注)。}。最后,在这些表层里,我提醒:统治着“心软爸爸”那份田园诗般的意象——边上有点瘫软的:一个恭敬而体贴入微的儿子,天鹅绒显露无遗,爪子隐在天鹅绒之面……

相较于前天的反思,昨天的反思更像是在给它增添细微层次,并由此稍稍磨锐其轮廓——却还不曾带来什么本质的新东西。诚然,当因时辰太晚而止笔时,我毫无已经抵达我所投入方向的尽头之感——“暧昧认同”的方向。事后回想,我意识到:由于一份积习——把自己“看作阳”——对我而言似乎天经地义:当涉及对我这个人的认同时,它只可能关乎我的阳特征。在此,在侏儒与巨人的舞台意象里,到目前为止,我是在“\textbf{巨人}”里认出自己的——虽然以一份歪曲的形式,却仍清晰可辨。然而,凭借我朋友身上的“颠倒”综合征之效,我被执拗地呈为“\textbf{侏儒}”\fnote{270}{(**)}{这个“侏儒”本身,不是别的,正是“超级妞”的隐喻——一个“假巨人”的面孔,臃肿、瘫软的形态……(1985年2月)}——这份混同(带着明显恶意的意图)被我当即拒绝:凭一份本能反射,它普遍而强力:直面一份嘲弄的意志——它拿我身上(在此即阴的)十分真实的特征当靶子,同时把同样十分真实的互补特征(后者享有抬举的共识)压回沉默——这样的情势,在我身上唤起那永世的反应:若非否认被指摘的特征,至少在默契中把它们缩小,并抬出那些恰恰被不公正地隐匿的特征,仿佛要去\textbf{对抗}它们。

凭这份“发自脏腑的”反应,我确确实实进入了冲突的轮舞——恰如我所该做的!它向我标示那个永世的“\textbf{挂钩}”:它在哪儿可以拿住我,把我拽进轮舞\dmark{}\origpage{736}。我自己的实在视景,也随之被歪曲了——作为对一份挑衅性歪曲的回答。所以,我昨天从牙缝里(或从打字机的键上)写道:

“‘第一份客观特征’——足以助成一份相似之感与一份认同行为——是渗透在我们共同情人——数学——的进法之间的那股强亲缘”。

这是白写的:我写它时,竟乐意忘了:这份“强亲缘”,由一份\textbf{阴}的、“女性”的\emph{进法}构成——在对事物的发现与认识中;而这恰恰是我作为与他“相似者”显形的一面:正因它,我与他一样,也显为“侏儒”:那是他秘而不宣、易受伤害、羞耻的一侧——待相宜时机一到,他便拉它上场,取而代之、“颠倒”。这份“天赐的情势”\fnote{271}{(*)}{见同名那条注,第~n\textsuperscript{o}~151。},即我认识冲动中的阴占主导,不仅是一份可疑朋友手中的\textbf{武器}——它首先也是、并首要地是,他对我认同的一份“客观根基”;这一回,不是对\textbf{父亲}的认同,而是对一个\textbf{兄长}、乃至一位“\textbf{姐姐}”的认同。

当我在这里用“客观”一词,是说要表达:这一回,涉及一份“认同”——它扎根的,不是“东家”想(或怕……)成为这个那个的虚构,而是一份\textbf{实在}:深刻、可触、无可置疑——双方原初本性之间那份\textbf{亲缘}的实在。无论如何,这份亲缘绝不会逃过他、一如逃过我;而我不怀疑:在某个深层面,这份亲缘的\textbf{意义}同样被知觉。而且我至少猜测——虽无十足的把握——这份知觉确确实实充当了他对我认同的材料。于是,这份认同是在\textbf{两个层面}上完成的:一方面是“理想”的层面,我在其中成为某类\textbf{价值}的化身——那些他但愿自己是其典范化身(哪怕仅在表面上,而模型显得遥不可及、理应着实实现理想)的价值;另一方面是“真实”的层面:认同凭一份被正确知觉的\textbf{亲缘之实}而确立——但那是些被公认\textbf{要命}的、可怜的共同特征上的亲缘\dmark{}\origpage{737}\fnote{272}{(*)}{这两个“层面”,于是对应着我这个人身上\emph{两种}彼此对立的“原型”:父亲(别名“巨人”)的原型,与兄弟、乃至姐妹(别名“侏儒”)的原型。后者也出现在“心软爸爸”的意象里——那意象由血肉之躯的父亲“如其所是”(唉!)所暗示,而非“如其所应是”……}。

此刻该记起:在我们相遇的当口、乃至其后十多年间,我身上肆虐着这份对我“女性”特征的压制——与我最终在我朋友身上确认的那份相同。带着回顾,依我看:我们相遇时,我朋友身上的这份压制已有几分程度,但主要尚属潜伏;无论如何,比我身上的弱得多。如我不止一次强调的:我这个人长期打有超级阳失衡的印记,而他透出的则是一个和谐平衡的印象。自那时起,他身上与我身上,有过\textbf{反向的演进}:在我朋友那边,从一个阴阳平衡的状态走向强的阳失衡;在我这边,从强的阳失衡走向一个(相对)的阴阳平衡。

随即浮现的念头是:我的朋友——或许凭这份对我的双重认同——(带着约莫三十年的时间差!)追随了我自己从八岁起所追随的那场演进:一份原初平衡的退化。有可能:一份对“阳刚”价值的适度高估——以“阴柔”价值为代价——在我的接触、或我所属圈子的接触之下,变成一份一竿子到底的高估。可是,一如我在别处强调过的:那场由他编排的葬礼的“神经”(或“活力”),以及在他自身变形(它同时是——在“东家”之手——他身上的孩子的下葬……)中的“神经”——这份神经,决计不能仅仅寓居在对这样那样、或极端或(甚至)疯狂的价值体系的单纯采纳里。同样的话,也适用于我这个人认同中的“神经”、以及这份认同在我朋友生活中扮演的过大角色。毫无疑问:是同一份“力”在运作,而它的根,深深扎入他的童年\fnote{273}{(**)}{关于朝着这个方向的更精确的直觉,尤见“缓期中的宿怨——或事物的回转(2)”那条注(n\textsuperscript{o}~149)。}。

这里又有另一个奇怪的念头浮现。仿佛是我一生四十年里拖过的那副最重的担子——那场“男性”对我身上“女性”的压制,它同时也近似于“大东家”对我身上孩子的压制——这份担子,被我的朋友“\textbf{接了过去}”\dmark{}\origpage{738}:恰在那他本可显得与此类担子无缘的时刻。那大约是在我的价值系统朝阴的方向\emph{翻转}的时刻——那场演进,预示了十五年后的“与孩子的重逢”;彼时,我骤然感到卸下了一副巨大的重负\fnote{274}{(*)}{关于价值系统这场“翻转”,在“阳扮阴——或主人的角色”那条注(n\textsuperscript{o}~118)里谈到;关于“重逢”,在同名那条注(n\textsuperscript{o}~109)里。}。随即便浮起与印度教\textbf{业}的念头的联想。对我清楚:在这八年里,我卸去了自童年起拖着的那份业的一个实质部分。我本会以为(而且现在我仍倾向以为):这份卸除,不是“以牺牲”任何人完成的;它不仅对我是善的,而且“对全世界都善”。我甚至可以说,我\textbf{分明知道}正是如此——即便竟显出:另一个人选择了(甚至,另一个人的命运注定要)把这份业接过去。诚然,这份我所卸下的业,我并不视之为一个“恶”。它曾是我一场\textbf{成熟}的滋养实质——那场成熟,正摆在我面前。我知道:我吃了、以它为滋养,是好的,对我对一切人都好;而一份认识,是在无知的滋养母体里成形的\fnote{275}{(**)}{朝着同一方向的反思,见“循环”那条注(n\textsuperscript{o}~116$'$)的结尾,尤其其中最后一段。}。我本觉得:这份实质或这份业,一旦转化为认识,不留残渣,它消失了。老实说,我不知道印度教或佛教传统在这个问题上教什么——是否有一份“业之守恒”的律法(类似于物质守恒),全然不受食入、消化、同化那些创造生命的过程之影响。

出于得体的顾虑,我刚刚在这些“生命过程”中,把\textbf{排泄}偷藏了过去。可它(一如整个机体的死亡)是一份回收所吸收之物的关键过程:让“死”的有机质,回到那永无穷尽的转化循环——活着有机质的循环——凭它,生命永恒地重生自死亡\fnote{276}{(***)}{关于生与死的循环,亦见“行为”那条注(n\textsuperscript{o}~113)。}。

\subsubsection*{注~156$_{1}$}
\notehead{156$_{1}$}{}
\phantomsection\label{ch18:local:n65}

(2月30日)这份“样式”,终于被最后的第十二组打破:它竟不幸地有\dmark{}\origpage{739}\textbf{六}条注,把构成“阴阳之钥”的注之总数推到62。我原预计这组“冲突与发现”会有\textbf{八}条注——那正与“总体性”的标准相符,并将把注的总数推到64=8×8=4×4×4——也是《易经》的卦数!我的预期没有实现,我很遗憾;可我也不肯为凑数而“作弊”,把关于Pierre~Deligne来访我处的那两条注塞进“阴阳之钥”——它们的天然位置,依我看,在“葬礼仪式”的续篇里,置于“阴阳之钥”\textbf{之后}。

尽管,我对第十二组仍留着一份不满之感:它是“阴阳之钥”十二部分中唯一不留下\textbf{统一}的灵感与用意之印象的。这份统一的缺失,依我看,并非系于“冲突与发现”这一主题本身,而系于反思进程中异己事件(有时扰人的)的闯入。

(3月7日)昨夜重读我归入一条注(n\textsuperscript{o}~162)、题为“信念与认识——或传递”的1月14日反思时,我感到了对这名字的不满\fnote{277}{(*)}{它也是“阴阳之钥”的最后一条注。}。一方面,这“主”题与副题乍看并不相合——事实上,它们分别对应反思中第一、第三“乐章”:彼此表面毫无关联:前者是对一份认识之绽放过程的描述(以一份骤然的\textbf{信念}的形式),后者是对业之无穷链条与“传递”、一代传一代、一个人传另一个人的唤起。何况,那份最贴身个人化的内容、那份对我自己“要害”的内容——它构成反思“第二乐章”的实质(而且它正是那道“桥”,从第一乐章通向第三乐章)——这份关键内容,竟没有出现在所选的名字里。(当然,对我毫无怀疑:这回偷梁换柱,绝非纯属偶然……)既然三个主题各有其重要,而我一时看不到任何把三者都唤起的“佳”名或双名,我终于\dmark{}明白\origpage{740}:最好把这条注一分为三,各取一个提示性的名字:“信念与认识”、“最烫的铁——或转折”、“无穷的链——或传递(2)”(n\textsuperscript{o}~162、162$'$、162$''$)。

我是事后才骤然明白:凭这场操作——可谓由反思的实质本身所口授——那份我拖了将近两个月的“美学”不满,也同时得到了解决:正当“阴阳之钥”这第十二个、也是最后一个部分(我先前唤它作“冲突与发现”)执拗着不肯(以自然的方式,自不待言)让自己补全为一组\textbf{八}条注,而只肯容纳已经写就的那六条。而我因没有向那份轻易的诱惑低头——“作弊”一把,把两条“信手拈来”、位置本在别处的注“贴”到“阴阳之钥”的末尾——得到了我的回报!这“阴阳之钥”的最后一部分(它终将取名“恶之谜——或冲突与发现”),同时也呈现出一副漂亮的对称结构:两个注包(各含三条注)围绕中心主题而聚拢,分居于两条“离题注”——一条关于Fujii~Guruji,一条关于我的僧友们——的两侧。

\subsection{冲突与发现——或恶之谜}
\label{ch18:local:h70}

\subsubsection*{(a) 无恨亦无容赦}
\tocsubsub{(a) 无恨亦无容赦}
\label{ch18:local:h71}

\notehead{157}{}
\phantomsection\label{ch18:local:n66}

(1月4日)在昨天与前天的反思里,我着力求取的,是与我的朋友认同于我这个人这一实在的接触,并由此辨明它的分量与含义。这是一件我做起来仍像在幽暗中摸索的事——不说在漆黑里的话。或者倒不如说:我的眼睛依旧闭着,我的眼睑对一束我依旧无力知觉的光,是不透明的。不管怎样,我不记得在我与我朋友的关系的任何时刻,曾“感到”或“看见”过这份认同;一如我不曾“感到”或“看见”过他对我怀有的敌对意向。然而我\textbf{知道},毫无怀疑的余地,凭一束丰富而相互印证的事实:这份对我这个人的认同,与这份如影随形的敌对,是\textbf{实在}——如同一个生而失明的人“会知道”太阳、日光、颜色、明与暗存在,纵然他从未见过它们。他知道它们,却不\textbf{认识}这些东西。或者说,倘若他毕竟有一种极为弥散的认识——经由一种或许更精微的触觉(或经由一个“回忆”——它并非只植根于他一个人的生命,而是植根于无数代有视觉的存在者\dmark{}\origpage{741}的生命,那些先他而在的人)——这份认识也仍是间接而微弱的,如同一副温暖而洪亮的人声,借一道遥远而不确定的回声,传到我们这里。

这两天的工夫,做得仍像一个权宜之计,像一份缺席的直接知觉的替身。凡是我所理解的“禅观”工夫,多多少少都是如此。这工夫时时在逆着一份\textbf{惰性}的水流\textbf{推进}——逆着眼睑那铅一般沉沉的惰性!诚然,在眼睛全然睁开、全然警觉的那些时刻,无须什么禅观、什么工夫:只消看,便见到。既然这样的时刻稀少,与其抱臂坐等它们到来,我宁可抢在前头,不计较这工夫是否笨拙而“缓慢”。它再慢,甚至有时比往常更慢——却从不曾原地踏步,也从不曾兜圈子。只要有工夫——我是说真正的工夫,由一个真正的愿望驱动——就有进展:某件事在被做成,在成形,在转变,此时不知不觉,彼时有目共睹……而有时,在一场笨拙而执拗的推进的尽头——那推进在一片无形无廓的幽暗里绵延数小时、数日,乃至数月、或许数年——奇迹发生了:盲人\textbf{看见}了!而被看见的,并非一道倏忽即逝、仿佛从不曾存在过的幻象,只留下一缕微弱的回忆痕迹。这是一份诞生于这些幽暗劳作的\textbf{认识},一份新的认识,亲切如在,恰似我们所爱之物的滋味。

我在前天的反思里写道:若说九个月前,在写“敌对的父亲(1)”那条注(我刚刚引过它)的末几行时,有一件个案,它的想法曾“引着我的笔”——那便是我朋友皮埃尔与我关系的个案。然而,当时与我更贴近的其他“个案”,想必也在这份反思的背景深处,呈于我的心头。当我在那里谈到一位“既受敬慕又受畏惧、既被爱又被恨的父亲”,继而又谈到“另一个自己,被畏惧、被恨、被躲避……”时,“畏惧”“恨”,乃至多半连“躲避”这个词,都\textbf{不}适用于朋友皮埃尔与我个人的关系。无论是凭直接知觉——哪怕何其倏忽轻微——还是凭对我所知的确凿事实的印证,我从来不曾得到过丝毫迹象,指向我的朋友对我怀有一份畏惧,或一份仇恨、哪怕仅仅一份敌意。反倒是真的\dmark{}\origpage{742},一如我不止一次有机会强调的。而恰恰是这个情况,使那份无缝无隙、表面上无缘无故的敌对显得如此令人错愕:它在已往十五年间一路渐强,顶着“暂停!”的风格——别名“天鹅绒的脚掌”\fnote{278}{(*)}{见“暂停!”与“天鹅绒的脚掌——或微笑”两条注(n\textsuperscript{o}s~77、137),以及紧随后者之后的那些注——它们构成“阴阳之钥”的“天鹅绒中的利爪”部分。}——最终抵达一种安然放肆的调门:(只要守得某些规矩)确信全然免于惩处……

这一份令人错愕、难解其故的演进,随即与另一份同样“令人错愕”、“难解其故”的演进(这一回,这可是委婉的说法!)联结在一起:同样在大约十五年间,那份与曾是我妻子的那位之间的伴侣关系里持续恶化的演进,连带波及我们所组建的家庭。由于没有任何迹象向我示警,说我的妻子对我怀有仇恨或长期敌意的意向,我花了十年,才在这段关系的无情恶化中(其间我最充沛的精力都被数学占去——它扮演着那著名的鸵鸟沙堆的角色……)终于确认:在我继续爱着的那个人身上,存在着一份执拗、神秘而毫不容情的毁灭意愿——它借着那些我疼爱的人,施于我身。那是1967年,早于我离开家宅五年,也早于这场冲突为我解开十年——我曾把它感受为我一生背负过的最重的担子。凭一份久已承担下来的关系所给出的纵深,我只能确认那至今于我仍是奥秘的东西:一份永不餍足的毁灭意愿,与此同时,却是\textbf{没有仇恨}、甚至没有敌意——面对那些成人或孩子,只要时机凑巧,便遭无情打击。

(在适当比例下)这是同一个奥秘,正如我现在在我朋友对我的关系中所面对的,带着这份区别:那份“执拗的毁灭意愿……借着那些我珍爱的人施于我身”,严格地局限在数学家世界的层面上,而且它的工具和人质,不是我“肉身所生”的孩子,而是象征性地充当这一角色的存在:那些多少“承我的名字”的学生及同类。在这两种情形中,我不仅察觉不出仇恨亦无敌意,更有甚者,对我还怀有同情、甚至常常是毋庸置疑的深情厚谊。

这些并不是我仅在他人身上直面一份伤人意愿——甚至一份摧毁的意愿\dmark{}\origpage{743}(在最强烈的意义上\fnote{279}{(*)}{所谓“最强烈的意义”,我指的不是为使痛苦而使痛苦,不是毁掉某样为对方所珍爱的有限之物,而是对他人施以心灵上(乃至身体上)毁灭的意愿;是(在可能之际)面对“超出理解的东西”种下一份洗刷不掉、摧垮一切的绝望的意愿。在“反常层会议”那副光鲜可亲的门面背后,我似乎在它最耀眼的两位活动家身上,重新发现了这个极端的维度……})——的仅有情境:在其中,我察觉不出丝毫仇恨或敌意的痕迹。对我一生打下最深烙印的那一例,发生在1933年,我生命的第六个年头,主角是我的母亲——就在那一年,我们构成的那个家——我的父母、我的姐姐和我——被永远地摧毁了\fnote{280}{(**)}{关于这段往事,见“超级父亲”那条注,n\textsuperscript{o}~108.}。

我近距离经历过的这类情境——无恨而欲毁灭,或欲伤人至深、深至所能,而我从中察觉不出丝毫敌意——彼此显得非常不同。我很怀疑能为它们找到一个共同的“解释”,或至少在主角们遥远的身世里找到一个共同的特征,足以暗示一条深的因果纽带\fnote{281}{(***)}{然而,一种恶毒的、深埋于内的对己之轻蔑,想必为所有这些情境所共有。或许,这样一种恶毒(当它不曾凭一个恩典的行动、凭一场深刻的内在转变而解消时——也就是说,只要它还不曾被“承担”)总要寻得出口,借种种毁灭性的行动、借一份毁灭的意愿表达出来;当它无法在他人身上找到目标时,便反身扑向自己。在不止一个男男女女身上,乃至在亲近的人们中间,我多次得以确认一份毁灭意愿的同时作用:既指向自身,又指向某个从亲近者中选出的外部目标(母亲、父亲、配偶或孩子……)。(1985年2月)亦见“无缘之暴力的缘由”那条注(n\textsuperscript{o}~159)里的反思——它比本注的反思晚三天,而本注显然为它作了铺垫。}。或许比一份解释更要紧、无论如何更在先的一件事,是先得\textbf{确认}这样一桩东西的存在:\textbf{无恨而毁灭的意愿}。我在这里接上了“无端的暴力”的主题——此前已从另一个侧面触及过它\fnote{282}{(****)}{见“天真的暴力”那条注,n\textsuperscript{o}~139.}。这里涉及的,是对一个亲近的存在、或一个被视作“朋友”的人施行的无端的暴力(有时是毁灭性的)。这样一种暴力——它难得自道其名\dmark{}\origpage{744}——在日常生活中单单\textbf{存在},就是每个人生活里的一个重要事实——人类生活中的重要事实之一。确认这一事实——顶着那些根深蒂固、时时撺掇我们把它藏掖过去的机制——是承担它的第一步。这一步,任何理论、任何推理、任何“进路”,都无法替我们省去。

我不知有朝一日我是否会\textbf{领会}这一事实;我觉得,领会它,也是“领会冲突”的一部分。对我清楚的是:这样一种领会,不可能出自一份“理论”,也不可能出自一份“经验”(单凭经验之力)。它不是某种累积(“知识”的或“经验”的)的“总和”,一如它不属于单纯智性的层面,甚至不属于单纯“智力”的层面\fnote{283}{(*)}{(3月5日)至少我知道:这样一种领会,唯有借对这份暴力\textbf{在我自身之中的}领会,才会来到我这里。}。我不敢肯定我认识什么人——哪怕只闻其名——在他身上活着这样一种领会。但我觉得:一个人,在一份无可否认、有千副面孔的实在面前,经过百次千次的闪避之后,终于抵达对这一事实的单单\textbf{确认}——谦卑地,无怨苦亦无反叛,无屈从亦无愤慨——如同确认一个可畏的\textbf{奥秘}:或许,它的意义逃过他的把握,他却预感到它的广度与深度;一个令他着迷、向他发问,却不再使他惊惧、不再使他惶惑的奥秘——这个人,就不曾白白活过。

\subsubsection*{(b) 领会与更新}
\tocsubsub{(b) 领会与更新}
\label{ch18:local:h72}

\notehead{158}{}
\phantomsection\label{ch18:local:n67}

(1月5日)并非预先谋定,昨天那份反思的结尾声调,依然全然是葬礼悼词的调门——只是这一回,由逝者本人亲口宣读(或咏唱)。人从来是自己侍候自己最好!

昨天,我再度直面“冲突之谜”最令人错愕的面向之一:一份无恨、亦无表面缘由的毁灭意愿,在暗处执拗而不懈地施于一个亲近的人,或这样的亲近者、朋友们。这样一份意愿,有时会终于失控加速,演成一场四面八方撒开的毁灭狂热:凡显出脆弱之相的,都成了受欢迎的靶子。这像一股按捺不住、专事倒行逆施的“行动”贪食症;当旁观的人(乃至刚吃了这份妨害的人)颇有幽默感\dmark{}\origpage{745},而那位“演员-提线人”支配他人的能耐又颇为有限时,它的重复性(像小丑把戏),连同它在牵线艺术上炉火纯青的绝技,效果可以滑稽之至。而当这些马戏游戏的妨害承受者中间有孩子在场——哪怕它们只在转喻的意义上“见血”——情形就更其严肃、更其后果攸关;或当那个被毁灭之渴攫住的人——无论男女——对某些同类握有可观乃至全凭己意的权力时。历史给我们传下了一些暴君的名字,他们被这样一种不加区分的毁灭疯狂攫住,把自家领地变成一座巨大的藏骸所。人们会想到伊凡雷帝,或斯大林,或中国某位皇帝(名字与朝代我都忘了)——他最终被走投无路、手持棍棒木桩的臣民们打倒了\fnote{284}{(*)}{这位皇帝惧怕民众暴动,曾禁止人民使用一切金属器物(如刀、叉之类)——凡可用作武器的——唯每村留一把刀,用坚固的铁链拴在一处公共场所。

三位人物有一个共同点:除这份毁灭之渴外,他们同样被\textbf{恐惧}攫住——怕遭暗杀,而在这恐惧之外,无疑还有对自身那无可避免的\textbf{死亡}的恐惧——可他们正四下播种着死亡。这个巧合决非偶然。我还注意到:斯大林(三人中唯一我有稍许翔实材料的一位)进入政治生涯时,恰恰是一位大行家——行在牵线的艺术,行在拿他人的虚荣与贪欲做杠杆来摆布人。他最早习得的风格,似乎是“天鹅绒的脚掌”,直到哪一天,连掩藏利爪都成了多余。

我未把我那位(前)同胞希特勒列入上述例子里,并非由于我对他怀有什么特别的同情,而是因为在他身上,我察觉不到刚才谈到的那种“四面八方”的毁灭癖。轻蔑、继而是毁灭的目标,是那些被指认为“他者”、“外人”的人:先是“犹太人”(还有纳粹行话津津乐道的共产党人与其他“犹太-布尔什维克”),继而是“亚洲人”和其他非雅利安的外来杂种。在希特勒治下,不沾犹太血统的正派德国人过得安逸至极——至少直到盟军头几轮大空袭、战局对他们真正急转直下之前。}。毫无疑问,在我们自己的地界也曾有过类似的情形,规模或许小些,而“历史”对它们则更其缄默……

昨天我写道——不带丝毫假谦虚——我不领会我刚确认过的那个“事实”:无恨的毁灭之渴。这绝不意味着我对此毫无想法,恰恰相反。我手里有的,明显不止一些单纯的“想法”,而正是几份强韧的直觉。它们诞生于、生长于我生命的沃土——那土壤富于冲突,那些冲突有时仿佛把它摧垮,如一场场无尽的风暴,在一片凝然不动的冬日大景里发作,毫不容情地拔走该拔走的\dmark{}\origpage{746}\fnote{285}{(*)}{这意象刚随笔势落纸,我便发觉它只是部分贴切——几乎带着一股“陈词滥调”的余味!在这余味上略一驻足,我找回了心里那个久已有之的蓄意图谋:“以阳观看我的一生”:运动、箭与风暴……

还来不及驻足,却分明觉出这意象有毛病(可它偏偏就是来到我笔下的那个意象,没办法!),我便在正文里“纠正弹道”,顺势接上“静候着的沉睡大地”——这可就是阴了!这正是把“假和音”(或“不协和音”)“解决”掉的那个和音。有一个意象,在许多方面都比风暴的意象——“拔走该拔走的”——更贴切,而且调门恰是更阴的:那虫子,啃噬“该啃噬的”——而它终究塌落——但一切都化作静候着的大地的腹中之物,而春天回来时……(接续不变!)。}。而一切,都为静候中的沉睡大地化作腹中之物。春天回来时,在僵卧于地、一动不动的巨树残干的凹处,浓烈的生命在攒动;到次年春天(若不是当年),那里已见青草与野花绽放。

这些“强韧的直觉”,我相信,全都关乎冲突的“成分”。其中某些成分,我稍稍谈过,又再谈过;头一个便是“对己之轻蔑”,以及它与我们原初存在的某些本质面向与力量的压制之间的关联——诸如“阴”“阳”这两个“坡面”,其中一个常遭否认。我也曾多次有机会谈到虚荣:它像一张名片,是我们之中冲突在场的所有征象里最普遍、最显眼的一个;在我看来,它是同一枚奖章的“正面”,其“反面”便是对己之轻蔑。还有对他人的轻蔑:对己之轻蔑向外的投射,同时又是它的掩体,或者更确切说,是它的宣泄口与驱魔仪。对他人的轻蔑,说到底无非是对其存在的刻意无视——把他当作一个有感受、与我们同等分享这个世界的存在。无端的暴力,唯有在这样的轻蔑的土壤上才能发芽、增殖。还有认识的恐惧,对实在的恐惧——这份恐惧的神经中枢,那个“黑点”,一份焦虑旋涡的震中——一遇风吹草动便待引爆,是对认识自己的恐惧:怕认出自己那些姿态与遁词,哪怕是最粗劣的;也怕认出我们之内那份创造的力——它日复一日被我们用同样的姿态与遁词拒斥、掩埋。

在我的生活里,恐惧现身于六岁,当时(照我看)还没有任何虚荣。虚荣想必是后来才出现的,在那场“翻转”发生的时刻(我猜想)——约在八岁\dmark{}\origpage{747}\fnote{286}{(*)}{关于这场“翻转”,见“超级父亲”那条注(n\textsuperscript{o}~108)。}。而恐惧也是第一个消失的,不留痕迹,随着一份既仁厚又不敬的好奇心的现身——它固然好奇,却绝不被那些装神弄鬼、阴森可怖的大场面机关布景震慑,诸如“黑点”公司一类。相反,自认识的恐惧消失以来的八年间,虚荣的机制一直原位未动,看不出变化。改变的,只是这些机制对我生活的把持:在一份警醒的好奇心在场的时刻,它们被拆除了引信——那份好奇心可不吃糊弄这一套!

我手里握着的,是冲突成分的一整个谱系——凭第一手、不带分毫怀疑,我知道:这些确确实实是成分,而且是本质性的成分。而多少年来,我也同样一切在手,只等我高兴的时刻,把这些成分“组装”起来:借我在自己身上、在他人身上得以观察的东西,细心阐明它们相邻与相依的关联。这是一件几天、几个星期——依我看甚至用不着几个月——的工夫,想必会非常有教益、非常有用。而我至今不曾费心去做,把优先让给了更直接切身的其他方向;这无疑是因为我心知肚明:对冲突的“领会”,不可能来自这样一种成分的“组装”——用一般性的措辞、而我这个人缺席其中的组装(除非我作为“例子”之一);也不可能凭把一定数目的单质并排放在一起、“组装”甚至混合——单质,即构成某一化合物的一些“成分”——就把这化合物复原出来。要让“复原”得以发生,首先得有一场“化学反应”——某种东西,以全然另一种方式把诸成分置于接触与作用之中,凭一种全然另类的力——那是单纯的“组装”或混合所不能为。

对生活之事的领会,亦复如是。单凭智力,勉强能辨认出“冲突”这样一桩东西的种种成分;而在成分已知的情况下,借助于有关的事实(第一手或第二手所知),它至少能把它们组装得貌似有理,甚至“正确”。这样的工夫自有用处:偶或在某一场冲突的情境里认得路径,从中得出一份多少精确的“病因学”——但这仍不是一份“对冲突的领会”。我倒是会说:朝这样一种领会迈进了一步,是在那一天——我的\textbf{与冲突的关系}得以转变的那一天\dmark{}\origpage{748}。我在这里说“我与冲突的关系”——这首先、自然,指的是我自己这个人身上的冲突,以及(由此出发)偶尔使我与这个或那个人对立的冲突;而最后,是我在日常生活中看到的、在或近或疏的人们身上运作的冲突——它常常借他们彼此之间的冲突表达出来。

在已往的八年里,确实有过这样一场朝着领会冲突的推进,也就是说:也是我与冲突的关系的一场转变——或不如说,接连的转变。我曾有机会提到其中两三段经历\fnote{287}{(*)}{关于这一点,尤见“接纳(阴的苏醒(2))”与“奴隶与傀儡——或闸门”两条注,n\textsuperscript{o}s~110、140.}。或许,对冲突的圆满领会,等值于对冲突之存在的圆满接纳——无论它在哪里,以何种方式显现\fnote{288}{(**)}{这样一份“圆满接纳”的意义,可能引起无数的误会。它与一份默许是全然不同的性质。它不排除\textbf{拒绝}——干脆而毫不含糊的拒绝——它把拒绝包含在内。关于这一点,见“伴侣——或‘恶’之谜”那条注(n\textsuperscript{o}~117)里的反思。}。显而易见,我离此尚远!而或许,对冲突的圆满领会,还意味着冲突在我自己这个人身上的全然消解。我离此更远!

然而关于那股力量的性质——它从一堆组装好的成分里,骤然激起一份更新此人的领会——我相信我还知道一件事。恰恰是这股力量,不“属于智力的层面”。我怀疑有任何智性工夫——譬如读书,且不管那些书多么博学、深刻或崇高——能丝毫促动它的显现。当它喷涌而出时,只在静默之中,只在接触我们这个人、我们经历里最切身之处;那是一样东西,任何书、任何人——哪怕是基督或佛陀——都永远无法向我们揭示。

当我说“最切身之处”,这并不意味着那是我们无法谈论的东西——对自己,或对他人——而且有时,谈一谈是有益的。但哪怕以天使之音、以先知之音来说,被说出的也不是那桩东西本身。这桩东西——早已认识,或许却深埋着;与它的接触,能骤然激起一份新的认识——这桩东西,天使不认识,先知不认识,乃至最亲近、最被深爱之人也不认识,唯独\textbf{你}认识。

回到冲突,回到“无恨的毁灭”\dmark{}\origpage{749}——在我看来,它是冲突最坚硬的“硬核”,最抗拒领会,也就是说:也最抗拒\textbf{接纳}。我还相信我知道,在摆在我面前、要更深入其中的下一步里,那桩“最切身”的东西是什么——我首先得重新找到与它的接触;在此处,它扮演的,正是那个出名的、被如此执拗地避开的“黑点”的角色!那是“无端的暴力”、对他人之轻蔑(或许还有“无恨的毁灭”)种种情境里的亲身经历——在其中,主角\textbf{是我}:行暴力的是我,从轻蔑中得计的是我。唯有接触这份实在——否则永无可能——我才有机会就那个出名的“对己之轻蔑”弄个水落石出,终于\textbf{看见}——在一切“毫无疑问”、一切“或许”之外——它是否确确实实是恶的深根,而不只是在“除我之外的所有人”身上!

\subsubsection*{(c) 无缘之暴力的缘由}
\tocsubsub{(c) 无缘之暴力的缘由}
\label{ch18:local:h73}

\notehead{159}{}
\phantomsection\label{ch18:local:n68}

(1月7日)前两条注里的反思,围绕着这桩怪东西的存在的奥秘打转:一份毁灭的意愿(或伤人的、羞辱人的、损害人的意愿),却在全然无恨、无敌意的情况下。促成这份反思的,是我朋友皮埃尔对我的关系,它随即唤起我的前妻对我的关系的联想。在整个关于葬礼的反思里,我不止一次被引着去确认,或去记起:在这两例以及别的情形里,在我这个人身上起作用的,正是某些特质——我自八岁起在自己身上培养的那些“超级阳刚”的特质——它们充当了这类敌对冲动的刺激物与“吸引子”。如果我没有弄错,这一点头一次谈到,是在10月5日那条注“超级父亲(阳葬阴(2))”(n\textsuperscript{o}~108)里。这层关联,在紧接着的10月9日那条注“重逢(阴的苏醒(1))”(n\textsuperscript{o}~109)里被重新接起。

在这条注里,我回到我一生中头一次知觉到这层关联的时刻。那是1976年10月18日,恰是与我心内孩子重逢的那一天,在见证这个日子——我成年生命里所有日子中最要紧的一个——的那几条注的末几行里。在那几行(已转录于所引的那条注)里,我谈到我曾爱过的三个女人的“隐秘的恨与怨怼”,其中一位,在那时还是我的妻子(尽管我已五年不曾与她同住)。凭回顾,我觉得:在我当时想及的三个个案里,每一个的这份“隐秘的恨”的印象\dmark{}\origpage{750},严格说来,并不对应实在——我是说,并不对应于任何时候我曾有的、对这样一份恨的直接知觉\fnote{289}{(*)}{(3月6日)写下这几行后,我记起:在我婚姻生活的岁月里,毕竟有过两段插曲——第一段几天,第二段几分钟——其间我觉得自己像被两束仇恨的强光袭中,它们从当时还是我妻子的人的眼中迸射出来。

第一次,我妻子正经受人们(委婉地)所谓的一场“神经抑郁”,时值我们婚姻的第五个年头(1962年)。这段插曲深深烙进了夫妻的生活与家庭的气氛。这也是我一生中——在所有我有清醒记忆的时刻里——活得最惨烈的时刻,给我留下的烙印也最深(而它也确实被指望如此)。

除非内里有一副异常稳固的根基(当年我远没有这份根基,那时还谈不上成熟),以我们为靶子的仇恨——当它来自所爱、所亲之人时更是如此——对我们的心灵有一种毁灭性的效果,只要它在我们心中唤起一种同样性质的、指向我们自身的毁灭性仇恨。我们身上似乎总有什么东西,非要不惜代价给“超出理解的东西”找出一个“意义”来,哪怕这个“意义”无非是一场自我谴责、一次无可挽回的自我拒斥:既然我们被恨(何况这恨的“缘由”我们全然无从把握……),那就是我们可恨……

我之所以被这段插曲创击得如此之深——在其后的六七年间,它始终像一把达摩克利斯之剑,悬在我的生命之上——想必是它与一场我童年的创伤经历发生了猛烈的共振。那经历已从清醒记忆中消失;可正因如此,每当我骤然直面一份无从解释的恶意或仇恨,它反而愈发作用——它们同样骤然、同样无从解释,一如那股在我五岁上袭中我的毁灭意愿,而它当年恰恰来自那样一个人:无论我在记忆里回溯多远,她都曾是宇宙那安宁而可靠的中心。

关于我有时成为其靶子的那份恶意或仇恨,我在生活里终于学到的重要一课是:我却根本不是它真正而直接的\textbf{缘由}(纵然我这人身上的某些方面——我既不否认也不推拒——有助于把它引到我身上)。可这份认识,此后多年仍嫌浮于表皮,不足以拆除那深植于我体内、每当我直面一份表面“无缘”的恶意或暴力便启动的机制。要拆除它,首先得让我溯到它的根,循着那些被遗忘的、浸满焦虑的日日夜夜寻踪而进——那时,我的母亲骤然间、神秘而不由分说地,成了一个陌生人,一个敌意的、可畏的陌生人……}。我曾知觉到、并有大把机会承受其效果的,是一份毁灭的意愿,或一份使人痛苦、伤人的意愿——既持久,又表面无从解释,无缘无故——我曾把它\textbf{解释}为一份恨的征象,一份“隐秘的”恨,因为它从不曾宣之于口。此外我还相信:对其中两位,正是在所引的那几行里——自我认识她们\dmark{}\origpage{751}以来头一次——我确认了这份在我看来是“隐秘的恨”的东西。在我当时所处的那一点上,我不可能不犯下我刚才指出的这层混淆。但这层混淆,丝毫不减损作出这份确认的重要性——在这份确认里,我把自己卷入得那样深,一如那些与我至近相连的女人们。

至于“怨怼”——它与“隐秘的恨”一并被谈到——从那一刻起我就感到:倘若我身上的一股“超级阳”的“某种力量”把这三个女人各自的怨怼引到了我这个人身上,那也是为着一些怨情——可我对它们毫无责任——为着一些创痛与损害,受于“早在她们认识我的存在之前、在一段缺爱童年的仓皇岁月里”。这份知觉,经年累月沉淀下来,如一段浓烈经历的果实;对我去年12月20日那条注“缓期中的宿怨——或事物的回转(2)”(n\textsuperscript{o}~149)里的反思,想必起过隐形向导的作用:在那条注里,出现了这样一份直觉:同一个怨怼移位的过程——一份原初的怨怼、一份“虚位以待的宿怨”的移位——想必也发生在我朋友皮埃尔身上,约在我们相遇的时刻,甚或更早。我所知的那些事实,至少使这份直觉显得可信。

然而与我前妻的情形相比,与重逢之后那场反思里谈到的另外两例相比,有一处重要的不同。的确,我全然不觉得我朋友的童年有多少“仓皇”或“缺爱”。这份不同,在我看来,显现在我朋友对我的那份敌对的调门里:它从不曾抵达那种激烈的音高——那样的音高,在那三段关系里,于我是何其熟悉。同样,在我朋友与我的关系里,敌对迹象的显现,起初极为隐微而零星;甚至在我1970年离开之后,又过了八年,这份敌对才直接而毫无疑义地表达出来,指向我这个人本身\fnote{290}{}{关于这一点,见“两个转折”那条注,n\textsuperscript{o}~66.}。这一切,恰与一份原初“怨怼”的存在相呼应:它始终弥散、无从称量,没有一个坚硬的“内核”与之相应——那本该是一份受辱或受害之感(哪怕藏于清醒目光之外),被感受为或许无可挽回……

在前一条注里,回想起那份毁伤他人的意愿\dmark{}\origpage{752}——它施于人身,却在\textbf{无恨亦无敌意}的情况下——一个表面的矛盾浮上心头(而且相当执拗),我本想立刻回头谈它。它是这样的。在我注意中心的两个个案里——涉及我从前的学生(和我推定的数学“继承人”)与曾是我妻子的那位——谈的确实是一份无意识的“宿怨”,他们把它转到了我这个人身上。可“宿怨”或“怨怼”这个观念本身,似乎与“敌意”或“不和”的观念连在一起:人们会想说:宿怨(或怨怼)是滋养一份敌意的可能方式之一——而且是最常见的方式之一。这份断言,在一种可称作“直接的”宿怨的情形里,肯定站得住:一份“真”宿怨,由针对当事人的一份怨情(真实的或想象的)所驱动,由当事人施加于我们的一份委屈或损害所驱动。但在我所关切的情形里,涉及的不是这样的宿怨,而是一份间接的宿怨——不妨说“凭代理”的宿怨:它从一副原初的潜在靶子上,挪往一副“认养的”或替代的靶子——前者出于这样那样的缘故而不配\fnote{291}{(*)}{这样的“缘故”有很多,它们常常使那个(有意或无意)造成损害、施加伤害的人,反倒“不配”充当一份怨怼或敌意——乃至一份仇恨或毁灭意愿——的靶子,尽管那意愿确确实实由他引发。最常见的或许(尤其当对方是母亲或父亲,或是一位凭其身份、社会地位而被视作不可触犯的人)是那道栏障:不敢违逆一份权威的禁忌——早已内化于心的禁忌。这类栏障力量极大。(对我而言,近十五年来,它们倾向于日渐消隐……)

反过来,也可能有这种情形:当事人“不配其分”,受不起一份与所受委屈相称的宿怨——显得太无足轻重、太躲闪或太怯懦,当不起那份本该名正言顺归于他的角色。

最后,我也可以设想:在某些情形里,所受的损害太过飘忽、太过精微(说到底,按当事人早已内化的通行共识,甚至“不存在”),除了生成一份弥散的宿怨之外,生成不出别的什么;那份宿怨无力“凝缩”,无法在一段本身调门温和、没有明显棱角的关系里成形、发力。这无疑只是前一情形的一个变体,出现在“缓期中的宿怨——或事物的回转(2)”那条注(n\textsuperscript{o}~149)的反思里。},后者则看来与事态的需要“对得上口径”。值得注意的是:这样一份“移位的宿怨”(这回可是实至名归的措辞!)——它作为那份执拗的力在运作,藏在种种态度、行为与行动背后,那些东西的性质,让人会说它们是被一份仇恨\dmark{}\origpage{753}、或一份“无缘的”敌意所驱动——这样一份“宿怨”,却\textbf{不带丝毫仇恨或敌意之感}!况且,正是“无端的暴力”——最强烈意义下的,即我此刻正在考察的这份——这两个\textbf{方面}的相互结合,使它如此令人错愕,如一桩真正“超出理解”的东西\fnote{292}{(*)}{关于这种“超出理解”的暴力(“德文作~unfassbar”),见“奴隶与傀儡——或闸门”那条注(n\textsuperscript{o}~140)。我在这里谈“最强烈意义”下的无端暴力时,当下只把它称作“超出理解的东西”而不另加限定;我那时心中所想的确切含义,在接下来的说明里被圈定:凭对那两个在它之中相互结合的“方面”的阐明。}:这份暴力全然没有理性的、可触的“缘由”,无论在吃这份妨害的人那一侧(他并未以伤人或损人的态度、行为或行动挑起它),还是在行使它的人那一侧(他并未被他对靶子怀有的——“有理或没理的”——仇恨或敌意之感所驱动)。

或许,在我所关切的这些情形里(人直面一份显得“无缘”、未被挑起的暴力),仇恨或敌意是否存在这个问题在这里是相对次要的;可以肯定的是,一如在我这里曾有的那样,在承受这份暴力的人的经历里,从被施的暴力变得有意识的时刻起,必会浮起一份“隐秘的恨”或“敌意”的印象——出自施加它的那一方。可这份印象,全然不是知觉的效果(那知觉仿佛被魔杖一点、骤然现身),而是生吞活剥的\textbf{等同}的效果:暴力=仇恨(或敌意)\fnote{293}{(**)}{(3月6日)不过在某些情形里,对一份确确实实在场的仇恨,完全可能有知觉,哪怕它根本不曾被挑起。(关于这一点,见本注上面另一条同日的脚注。)此时涉及的仇恨,除非在例外的情形,始终圈禁在无意识的深层,而且在那里还处于一种“虚位”状态,没有指定的靶子;可它恰恰是暗中驱动种种暴力行动的力量(多以阴险的形式)——而那些行动,却确确实实、毫不动摇地指向同一个选定的靶子……}。

但在我看来要紧得多的一件事,是确认:不仅“凭代理的宿怨”这样一桩东西存在着——它表面何其反常、何其悖理、何其违背那些最根深蒂固的“常识”反射:从它的“原初靶子”\dmark{}\origpage{754}(或诸原初靶子)挪往一副“替代靶子”(一副几乎纯为图方便的靶子!);而且要确认,更进一层:这是最常见的机制之一,在每个街角都会遇到,无论是在自己这个人身上(那恰是人们最不会想到去找它的地方……),还是在亲近者、朋友们身上。我甚至有这份印象:这一机制具有普遍的本性,它属于人心那些基本机制之列,属于构成“在实在面前逃避”综合征的那几种万能机制:拒绝去认识实在,害怕承担它。

说得更确切些,我有这份印象:今天,我把手指点在了所有“无端的暴力”情境——无一例外——共有的那个机关上。这份印象是带着一份骤然信念的力量浮现的,当时我正着手考察(三段之前)那个“表面的矛盾”。我那时感到:我一生积攒的、围绕这份“超出理解”的暴力那最疼痛点的、一大堆零碎而不相属的印象,骤然排起了秩序,霎时获得了一直欠缺的一份视角——这份视角出人意料地在那里现身——在一次反思收尾的转弯处,而我正准备给最后一个~i~点上最后一个点……

\subsubsection*{(d) Nichidatsu~Fujii~Guruji——或太阳与其行星}
\tocsubsub{(d) Nichidatsu~Fujii~Guruji——或太阳与其行星}
\label{ch18:local:h74}

\notehead{160}{}
\phantomsection\label{ch18:local:n69}

(1月8日)一周以来,一股少见的寒潮——零下十五度或更低,而当风从“旺图山”吹来(这名字名副其实——“旺图”(Ventoux)里就带着“风”),想必还要更冷。据说这股寒潮在世界各地肆虐(据一位听新闻的人说),而在南方,自1956年那场有名的寒冬与晚春以来,还不曾见过。我在德国的童年经历过这样的冷,但那时有雪,雪护着土地,给空气、给万物添上一层温柔。眼下这场无雪的冷,地表冻得像一整块冰。几天工夫,园子就被剃了个精光——春天来时,我们播下种下的还会剩些什么,我不知道。残存的韭葱、芹菜、莙荙菜、野苣、甜菜、刺菜蓟的叶子,都像冰做的叶子,冻菜一般。人们赶着一天天尽量多收,随收随吃,赶在解冻之前——一解冻,一切都得进堆肥。而昨天\dmark{}\origpage{755},厨房里的进水管冻住了,幸而楼下旧车房里还有自来水,那边不那么受寒。今天一位朋友带着便携燃气喷灯来了,把水弄通了。往后得让水细细地流着,免得它立刻又冻住。所幸饭厅里有一只好柴炉,我把活儿挪到了那里。坐在炉边,实在舒坦。我烧葡萄藤茬取暖,每天拿斧子劈,满满一葡萄筐还冒尖——天就是这么冷。当风整个下午刮个不停,光是顶着风劈上一刻钟、二十分钟的柴,就足以让人把手冻坏。更别提留在外面的车,再打不着火了——据说汽车这东西,顶不住大冷,防冻液也一样。还是那位好心的朋友,刚才把车发动了起来;可明天它还走得动吗——去校读那位我交付了活儿的打字员的打件?总之,冬天一股寒潮——有时是夏天一股热浪,或随便什么时候来一场好好的小病——就足以提醒我们几桩存在的实情;而当一切如意运转,人们容易把它们忘掉……

三个月来,不知不觉,我的工作节奏挪向了夜间的时辰。我工作到夜里两三点,睡到十一点、中午。就这天时,倘若我在床上由着自己一次,十二个小时睡下去毫不费事——反过来,一旦上了工,我就不再上床!目前我试着保持一份合理的平衡。作息的错位,我不太惊慌,只要觉还睡得好,不至于一连几个小时躺在床上睡不着、那台思考的机器还转个不停。即便眼下园里没什么活,每天总还有足够多样的事情可做,包括备柴火,还有这里那里一点小体操。我有一种满意的生计平衡之感:其中,发现的工作并不摆出吞没其余一切的架势,却也绝不受亏待。自9月22日重新上工以来,我平均每天得花上五六个钟头。不算多,可“收成”似乎不比从前差多少。“砍伐量”(每月约一百页)与写《收获与播种》前两部分时大致相当,相去无几。而从质的一面看,我毫无怀疑:最深的正是这第三部——它教我认识自己、认识他人,教得最多。

\asterisms

\dmark{}\origpage{756}Na mu myo ho ren ge kyo !

正当我写完这一小段回顾——关于冬天的严酷,关于我生计平衡的演变——接到了日本山妙法寺团体一位僧友的电话,告知他们尊崇的“preceptor”\fnote{294}{(*)}{“Preceptor”,英语词,与“teacher”约略相当,指“师父”、施教者。日本山妙法寺(Nihonzan Myohoji)是该团体日本名称的音译,意为“日本传道团”。这是一个佛教“传道”团体,主要宗旨是和平主义。详情见下文。}Nichidatsu~Fujii与世长辞——他更广为人知的名字是Fujii~Guruji,亲近的人则唤他“Osshosama”。我巴黎的朋友刚从东京的电话里得知消息,我猜想,Fujii~Guruji就在今天逝世\fnote{295}{(**)}{事后证实,他当时刚去世不过几个小时。消息传得真快!}。去年8月6日,他刚满一百岁;身体已衰弱,头脑却极佳。

奇巧的是,8月6日这一天,还是另两桩重要事件的纪念日,一桩具有历史的分量,一桩于我个人切身。它是广岛原子弹(1945年8月6日)的纪念日——日本人把它纪念作“广岛日”(Hiroshima day)。(正因如此,Fujii~Guruji的寿庆反倒多在7月末举行,好把8月6日前后的日子留给和平与反核的集会。)另一桩:我的父亲生于1890年8月6日,比Fujii~Guruji的诞生——一日不差——整整晚六年。

继Claude~Chevalley之后,Nichidatsu~Fujii是在《收获与播种》写作期间辞世的第二位、在我生命中起过不可忽视作用的人。面对他的离世(它其实算不得意外),我尤其庆幸:直到去年,我与他还有过一回满含暖意的通信。我曾受邀出席老师父的百岁寿庆典礼——它将在东京以罕见的排场举行。(人们甚至赶工印出一本关于他本人的证言小书,备作这一场合的献礼。)那给了我一个机会,像几乎每年一样,写几句预致的贺词,并为不能出席7月30日的典礼致歉——写这封信时,我自己还多多少少卧床。(诚然,我也不怎么喜欢大型公开典礼,可我当时觉得,信里不必提这个。无论如何,我一味固执地谢绝出席任何“大场合”\fnote{296}{(*)}{这类“大场合”中主要的,是“和平塔”(Shanti Stoupas)的落成典礼。这类塔——为世界和平而设的静修之所——的建造,溯自佛教世界一条极古老的传统(由印度的阿育王开启),也是Fujii~Guruji毕生的主要关切之一。他推动了世界各地大量和平塔的建造,其中三座在欧洲,一座在美国。},想必让不止一位僧友失望和伤心了\dmark{}\origpage{757},而他们邀请我,从不厌倦。)我还在信里添了几句话,谈病有益的一面:它不由分说,强使我们从诸务中“脱开”,把身体索要的还给身体。Fujii~Guruji自己,过去这一年也大半卧床;以他那倾于行动的性情、那份少有的精力,想来于他是沉重的。我已有七年多不曾收到Fujii~Guruji的私人音讯,却惊讶地收到他的一封来信——他卧床时亲自口授的。信(我刚刚重读了一遍)写于1984年7月13日。信中满纸体贴:他挂念我的健康,难过不能派人前来照料我。他也谈自己的健康,谈他以怎样的心绪承受这份被强加的无为。信以这样几句作结——那腔调非常“日本式”,须拿一颗(大大的!)盐粒来听——而它们向我显出的,或许比全信的其余部分更多:他的精气神一如既往\fnote{297}{(**)}{信是以日语口授的(那是古鲁吉唯一会说的语言),随即直接译成英文。所引几行的法译:“诚然,我是一个十分衰老、衰朽无用的人,即使我还能恢复正常的健康。可我仍愿意活下去,看看这世界怎样转下去。”}:

\begin{quote}
Indeed I am a very old decrepit man of no use even if I may get back to normal life. Yet still,
I would like to live and see how the world turns。
\end{quote}

随后,他果然还看着这世界转了近六个月……

我与日本山妙法寺团体的渊源,上溯到1974年。在这里,这些多插曲、方方面面都有的关系,连勾勒一个轮廓也谈不上——那得用一整卷书。它们是“生存与生活”那段插曲\fnote{298}{(***)}{《愚妄与更新》(《收获与播种》第一部分)里,多次提到这段插曲。“生存与生活”(起初就叫“生存”,再无其他)是一个团体的名字,起初以和平主义为宗旨,继而也兼生态关怀,1970年7月(在蒙特利尔大学一次“Summer School”的边缘)诞生于一群科学家(尤其是数学家)中间。它很快朝“文化革命”的方向演进,同时把听众扩展到科学界之外。它的主要行动手段,是同名的(半定期)公报;历任主持依次为Claude~Chevalley、我本人、Pierre~Samuel、Denis~Geygi(四人皆数学家)——还不算一个英文版,由Gordon~Edwards勉力维持(一位年轻的加拿大数学家,我在蒙特利尔结识他;他是团体与公报的少数发起人之一)。

第一期公报全部出自我的笔下(天真而满怀信念!),印了一千份,在尼斯国际数学家大会(1970年)上散发——那届大会(照例每四年一届)聚了数千名数学家。我本指望会有大规模的响应——记得响应一共两三份。我尤其感到同事们中间一片难堪!谈到科学家与军事机器的协作——它们已渗透到科学生活的各个角落——我几乎是一脚踩进了最肥腻的菜碟……正是在科学的“上流大世界”里,我感到最大的难堪——从那边传来的同情回应,最后只落到Chevalley与Samuel两人。倒是在我别处称作科学界“沼泽”的地方,我们的行动找到了某种共鸣。公报最后印到约一万五千份——顺带说,后勤是一桩疯狂的活儿,发行全靠手工。Didier~Savard那些多汁的插画,想必为我们这份小报的相对成功立了不少功。在我与Samuel先后离开之后,它终于滑向一个小小的左翼团体,行话尖锐、论断不容置辩,公报也终于寿终正寝。该理解的、该说的,在1968年那阵亢奋余温尚近的时候,已被理解、被说了。此后再把一张唱片无休无止地转下去、倒回来,已没有什么意思……}最丰厚的“落果”之列——那段插曲\dmark{}\origpage{758},随我1970年至1972年底的离开而来。1972或73年,某份(或某些?)日本报刊谈到过这个团体、同名的公报,也谈到我的“离开数学”与我的“轨迹”。那些“科学批判”的一面、对军事机器的揭发,或许还有“文明批判”的一面,想必多少“透”进了某篇文章,引起日本山妙法寺一位僧侣的注意。他把它讲给了别人,尤其是同城(鹿儿岛)一位更年轻的僧侣——后者在他的影响下出家,颇有几分“弟子”的气象。1974年春,这便是该团体头一位登上“西方”海岸的传道僧——确切地说,到了巴黎\fnote{299}{(*)}{他确确实实向我保证过:在佛教历史上,他是第一位远赴西方的佛教传道僧——这份材料是否可靠,我可不敢担保!再说,出家当传道士,对佛教是否真是一大“进步”,也未可定。打一开始,日本山妙法寺团体的这个面向便在我心中引起保留,而岁月只是让这份保留一再获得证实。}。几个星期后,他来找我,而且\dmark{}\origpage{759}不经通报,来到我当时住的那个偏僻村子——离蒙彼利埃五十来公里。自那个难忘的五月天起——正午的日头底下,我看见一个装束古怪的人,一路敲着鼓唱行,正朝着(没错儿……)我独自劳作的那个园子走来——自那一天起,我有了这份荣幸与喜悦:看到古鲁吉的众多信徒与支持者\fnote{300}{(*)}{恰恰是其中一位,有此“荣幸”:作为“身份不合规的外国人”,成了法国判例里某部“1949年法令”中某条相当不可思议的条款头一次被照字面适用的由头。而我也有幸进了轻罪法庭,罪名是“无偿收留并藏匿”了这么一个法外之人。关于这段插曲,见“我的告别——或异乡人们”那节(n\textsuperscript{o}~24)。},络绎不绝地来到我家中做客。与他们的接触,给了我很多。1976年11月初,我甚至有了这份殊荣与欢喜:在我朴拙的居所接待Fujii~Guruji本人——时年92岁——同行的是七或八位僧侣、尼众与弟子组成的一行人。前一年,我已见过他:在团体位于巴黎第十八区的殿堂的隆重落成典礼上。在例行的客套话之外,那时有过强烈的接触,有过即刻的相知。而这次数日逗留寒舍——一个更贴身、更个人的情境——自然给了我一份丰富得多的领会:既关于Fujii~Guruji其人,也关于他与团体的关系——他正是那团体的头颅与灵魂。

有意思的是,Fujii~Guruji的这次来访,与那场我生命中的关键转折——它完成于同年10月15日至18日之间,别处已经谈到\fnote{301}{(**)}{见“欲与禅观”那节(n\textsuperscript{o}~36)与“重逢(阴的苏醒(1))”那条注(n\textsuperscript{o}~109)。}——相隔不过两周。紧随那些危机与更新之日的几个星期,是我一生最激荡的几个星期:每一天都带来它意外的收成——内在的事件与发现。说实话,这次来访——几个星期前就谋划、筹备停当,由整整一个僧侣与尼众的团体簇拥着他们尊崇的师父——来到那里,像一段奇异的插剧,像当时吞没我全部存在的那场冒险中的一次岔神。正是对客人们的敬重——尤其对亲临寒舍的Fujii~Guruji的敬重——让我毕竟得以在这几天里,保有这场合所要求的从容。一如我常有的情形,唯有置身事件的中心,我才明白:它绝非什么“插剧”或“岔神”,而\dmark{}\origpage{760}正是我正在经历的那场冒险的一部分。顶着它那一副十足的“东方故事”的门面——极尽的体贴,不寻常的魅力——这个所谓的“插剧”,让我面对的是与我相似的男人和女人们,面对那些我向来认识的男人和女人们——只是处在没那么异国情调、外表没那么不寻常的情境里。正是感到了这份亲缘,我也才在我的客人们身上感到了朋友与兄弟,而不是直直走出《一千零一夜》的人物——对不止一个目瞪口呆的村民,想必正是后一种情形。而Fujii~Guruji本人——他的“亲近者们”还守着尊崇师父所要求的那份得体的距离,他却同我说话如此家常——我感到他非常、非常遥远(对我,一如对他的亲近者),却又同时切近,仿佛他是我的父亲,或一位宽厚仁慈的兄长。

而一如在一位父亲或兄长身上——哪怕最宽厚的——并不少见的那样,他对我怀有一份期待,对此他也并不掩饰;一份陪伴他的人们共有的期待——他们个个都是我的客人。而我也知道,我无法回应它。我的冒险与Fujii~Guruji的冒险相连——凭一些我难以辨明的纽带,或许比我能看到的更深;也与他的弟子们的冒险相连——他们闭着眼睛追随他。但这场冒险,既不属于我这位显赫而宽厚的客人,一如它不属于我的父亲——他于我同样显赫而宽厚,极其切近却又不同:另一个人,另一个命运。

要让他们“明白过来”:在一桩属于他们的事业里——我并不觉得它属于我——我不会是他们中的一员,这并非易事。照人们想必向Fujii~Guruji和他的信众们描画的那幅我的图景,这该是他们最想不到的事——个人层面的关系越是水乳交融,就越想不到:无论团体整体还是它的各个成员与我之间,都近乎一场真正的蜜月。也正是在这次来访期间,某些源于我所受教育的、积年的抗拒消散了:我加入了客人们,与他们一道唱他们的曼怛罗,以鼓伴奏:“Na mu myo ho ren ge kyo !”

这曼怛罗是他们宗教修行的根基,是他们的阿尔法与欧米伽。他们唱它,多伴着祈祷鼓,晨一小时,晚一小时。依日本先知日莲的教诲,这击鼓而歌,本身即是至善,给唱它的人、也给他周遭带来和平。所以,对我\dmark{}\origpage{761}的日本朋友们,这歌唱便是人们通常所说的“祈祷”。他们赋予它的意义——与日莲相合,也与他们直接的“preceptor”Fujii~Guruji相合——是一份\textbf{尊崇的行动}:对着向之咏唱的那个人,并透过这个人,对着宇宙间一切活着的存在——就其(依《法华经》)被许诺终将成佛、成为圆满智慧的化身而言。这七个音节,也用作对任何他人的问候,乃至对任何想要问候的存在,带着那层对他人身上本质神圣之物的尊崇之意。它们也用作饭前的谢恩。说实话,我觉得几乎没有哪一种场合——无论是惊诧、感怀还是凝神的时刻——对一位日莲的追随者是不宜说出那句神圣之言的。至于我,虽不分享僧友们的宗教信仰\fnote{302}{(*)}{我不觉得自己属于任何一种特定的宗教。因父母所受的教育,我是个无神论者(带一点反宗教的色彩),直到十四岁。我的博物学老师作过一场精彩的讲述,谈地球上生命演进的历史——它让我懂得,毫无分毫怀疑地:宇宙之中,有一种创造性的智慧在运作。这份理解,当时还停留在单纯智性的层面;它在我此后的成熟中扩展、精纯起来——这成熟,在我1970年离开数学舞台之后仍在继续。},却满心欢喜地加入他们——每当时机来临——行~Odaimoku,击鼓而歌他们所谓的“祈祷”。正是怀着对他们的纪念,怀着对他们师父Nichidatsu~Fujii~Guruji的一份深情尊崇的行动,我也把“祈祷”纳入了我的日常生活:每天两顿主餐之前各唱一遍——至少当我自家居家、或身在朋友家中、或与那些我知道不会因此不安的人同席时\fnote{303}{(**)}{我尤其避免在每周去学院吃的那顿饭上唱祈祷——与几个学生或同事同桌——我怕他们中会有人从中感到一种强制,觉得我借年长者或“东家”的地位把它强加于人。}。这是我从Fujii~Guruji和他那些我认识的弟子们那里受惠的贵重之物之一——他们给了我他们的情谊,对我或近或远参与他们传道事业的推拒,从不厌倦。

在日本,有数百万日莲宗的佛教徒,分属众多面貌迥异的宗派。日本山妙法寺团体按人数是最小的之一:几百位僧侣、尼众与活跃的同情者。然而它在日本内外都颇知名,在所有传统宗教团体之中,它以一份毫不含糊的政治担当自我区分\dmark{}\origpage{762}——其主旨是为和平而斗争,是反军国主义的行动,尤其反核。越战年代,它是(如果我没记错)唯一明确站在美国人对立面的佛教团体,并为反对美军基地留驻日本而斗争(那些基地,正是越战延续下去的后勤支柱)。近年,Fujii~Guruji还与美国印第安人解放运动的领袖们——“美国印第安人运动”(AIM)——有过密切接触。日本山妙法寺的僧侣们参加过美洲印第安人组织的进军,还不算世界各地的其他和平进军。印第安领袖们显然被Fujii~Guruji不寻常的人格吸引、折服。这位精力不屈不挠、年近百岁的老人,竟是另一种宗教信仰的大传道者——这个事实似乎丝毫不使他们为难。恰恰相反:尊者师父那一竿子到底的“反美”倾向里的宗教维度,无疑连同他的高龄,是他们把古鲁吉当作自家人迎接的原因之一——当作一位极受敬重的父亲或祖父,一个人们在他身上认出自己的人\fnote{304}{(*)}{为给出一个印象,说明联结印第安领袖们与古鲁吉其人的那份信任与尊崇,我在这里指出:在围绕“太阳舞”举行的那场盛大的年度入会节庆上,也有古鲁吉的弟子僧侣们参加,从日出到日落,擂着大祈祷鼓,合着那扣人心弦的节奏:\emph{Na mu myo ho ren ge kyo}!这些大鼓,由整段的树干凿成,绷着牛皮,声威不凡,而(我猜想)一连十二个钟头怕也难以忍受。(在巴黎殿堂的落成典礼上,我领教过两个钟头——经验是有结论的……)不管怎样,Robert~Jaulin(他与僧侣们一道,是受邀参加节庆的少数非印第安人之一)向我转述:印第安人以斯多葛式的坚忍,从头到尾承受了祖父古鲁吉的圣鼓——那祖父的鼓声,会是入会礼的诸多试炼之一……}。

无疑,这份宗教维度之于我也起了同样的作用——它使Fujii~Guruji于我更切近,尽管我并不自认任何一种界说分明的宗教信仰。若问在他身上最吸引我、最触动我的是什么,我看到好几种东西。最显眼的,是一份内在的\textbf{喜悦}。这份喜悦似乎自发地源于他这个人身上的一份\textbf{统一},或者或许不如说,源于一份\textbf{对他自己的忠实}。人们感到这个人是幸福的,因为他一生都毫不犹豫地做了他感到自己该做的事。他于我并非没有矛盾,却没有暧昧。他的某些行动或不行动的意义,我无从把握\dmark{}\origpage{763},但对这个人全然的正直,却不曾有半点疑问掠过我心。倘若确是如此,这并非凭对那些经他人转述所知的他的分析。只消见过他一次,便知道这是一个不知暧昧为何物的人,一个与自己深切相合的人。美国印第安人运动的领袖们必是感到了这一点,才在中间给了他他们所给的位置。他对那些自认是他一伙的人的超凡感召,无疑也正系于此:那些男男女女,意识形态与哲学倾向从彻底纯粹的马克思列宁主义,一直铺到一家连锁大百货公司总裁那种体面的随波逐流。把他们联起来的,不是对某部《法华经》的崇奉——他们中或许无人曾自大到去读它\fnote{305}{(*)}{古鲁吉不止一位弟子曾让我明白:摆出读《法华经》的样子,会被视作自大——纵然它已有日译本存在。唯有一个精神极为深邃的人——如他的师父Fujii~Guruji本人——才有能力、才配得上读这部神圣的文本;它远远超出俗人的智力。显然,这些男男女女的信仰,直接系于Fujii~Guruji其人,而非某个或多或少被神化的历史人物,如佛陀,或圆满的菩萨、先知日莲。};也不是某篇源出巴利语的祈祷文——经汉译辗转还原为日文,其教义正是崇奉这部《法华经》。把他们联起来的(或者该说:曾把他们联起来的?),是一个\textbf{人}——他对他们施加着一种感召;而他并不曾刻意寻求这份感召,一如太阳不曾寻求它的行星。

我还看到:这个人是\textbf{孤独的},而孤独并不压在他身上。那是他的自然处境,或许从来如此。这份孤独,与这份正直、这份与自己深切相合,在我看来,是同一桩东西的诸多不同面向。这同一桩东西还有一个面向:\textbf{力}——一种没有暴力的力,它也不在意自己是不是、或看起来“强”。这仍是太阳的力:太阳只消是它自己,它周遭便自会生出那个力的场,生出那些行星循行的轨道。

无疑,这也是我在《收获与播种》里不止一次谈到的那份“我们之中的\textbf{力}”——区别只在:在这个人身上,它全然显露,凡接近他的人都感觉得到;在另一个人身上,它埋得或深或浅,有时深到让人以为它并不存在。而如果我某些僧友看起来在自身之中否认它,他们所信奉崇敬的那部《法华经》,连同他们日复一日唱着的祈祷本身\dmark{}\origpage{764},却分明宣告:一种这样的力,活在造化中一切活着的东西里——如他们一样,也如他们尊崇的师父~Osshosama 本人一样,被许诺了佛陀的命运。

\subsubsection*{(e) 祈祷与冲突}
\tocsubsub{(e) 祈祷与冲突}
\label{ch18:local:h75}

\notehead{161}{}
\phantomsection\label{ch18:local:n70}

(1月13日)\fnote{306}{(*)}{(1月23日)本注的前半部分,是顶着强烈的抗拒写下的——抗拒在于要提及干扰我工作的种种风波。这些风波显出几分可笑的模样,单是提它们,就有点像自请领教鞭、任人责打!另一方面,这些风波——“真能把人锯断的”——在我的工作里已变得如此刺耳、如此侵扰,尤其这一两周,若装作无事一般讳而不提,反倒会是一种作假、一种见证里的不真实。此外,十天之后,我在“Jung——或‘恶’与‘善’的循环”那条注里,重新谈了我的这些倒霉事。
(3月7日)后面这条注——关于C.G.~Jung自传的一整套“读书注”的头一条——最终被划归《收获与播种》的末一部分,即由这部自传所引发的那部分反思。}又是整整四天,我不曾有工作的闲暇与平静——我是说,继续写注的闲暇。缘由主要在于:要把《收获与播种》这第三部分誊打成清稿,我遇到的困难简直不可思议。交人打字这类活,我做了三十多年,从未见过这样的事。显然,手头这份性质极为切身、甚至可谓私密的文章,在负责打字的人身上触发了(想必是无意识的)力量可观的反应,每一次都奔着对托付给他们的工作的真真切切的破坏而去。几个月间,同一个剧本一连上演了三次,只差些变奏,三位接续的秘书,她们彼此却不曾串通\fnote{307}{(**)}{那些对我怀有善意的人,在这里有得可说:该给我扣上一顶“受迫害妄想”的帽子了——搬家师傅公会之后,又轮到打字秘书们的公会总动员,一心跟我过不去!关于前科,见“屠杀”那条注(单看注的名字,关于我这个人,它已经说得够多了……)p.~538,谈的是我朋友Ionel~Bucur的搬家……}!而这第三回,又添上一段卑污的插曲:秘书J太太摆出要用那份交在她手里的、颇为异常的稿子作筹码,敲诈一笔赎金。这是一位资深的老秘书,行业经验十分老到\dmark{}\origpage{765}。头打的十一页无可挑剔,几乎没有一个错字——为的是显一显她的本事;而单单随后的十五页里,就跳行了十一行——把文本糟蹋到这般地步,我也少见!为取回我的稿子与打件,赎金的数目(除已打好文字议定的报酬之外)我没有去问——我毫无意愿鼓励这类行径。这意味着,我多半将不得不诉诸司法途径。

所幸我手头还留着稿子的一份草稿,需要时用得上。但这路数的马戏——尤其当它带上了重复性——真能把人“锯断”。当我设想我这块沉思而自传的拙砖会掀起怎样的困难与对抗时,我当然不曾想到:头一拨麻烦竟会从这一头袭来——从打字秘书们的公会(而非我可敬的数学家同行们的公会),而且是消耗战的路数!眼下,把这份文稿(一旦取回)再交到第四位秘书手里,我已不太热心:我没有任何依据,能预见她会比她将接替的那些人多一分怜惜。而自己来做秘书的活计,得投入整整一个月的时间——我绝对无意提供。

或许我将不得不放弃把这第三部分誊成清稿,直接以草稿的形式交予出版者。(排版付印的排字工们,我倒不预期会有同一路数的麻烦!)这尤其意味着:我将放弃把第三部分收入《收获与播种》的限印预版——它应由我的大学USTL经办,在同事与朋友间以私人名义散发。或者,倘若我终于找到一位把活做正派的秘书,也可以以后再印。这一部分(三部之中肯定最“难”的一部),我只寄给那些确实想收到它、并在已收到前两部的人之中明确提出请求的人。我确实巴不得把前两部印出来寄出去(第三部分,我倒不那么急)。这两部分的打字早已完工数月,是由USTL的秘书们经手(而且毫无问题)的。它们本可以早早就印出来,若不是我想在里头收入涵盖《收获与播种》全三部的总目录——而三个多月以来,我相信\dmark{}\origpage{766}自己已处在给这没完没了的第三部分收尾的关口。这里,我打算给自己宽限到本月底完成;不然,就先操办前两部的印行(《愚妄与更新》,与《葬礼I,或中国皇帝的新衣》),而不把第三部分(《葬礼II,或阴阳之钥》)完整而确定的目录收入其中。

而在这些不快的插曲之后,如今我得千方百计找回一份被打断的反思的线索。

1月9日,Fujii~Guruji在他第一百零一个年头上辞世,这给了我一个机会:借他这个人,唤起我生活的一个此前未曾触及的面向。既然无缘再见古鲁吉于灵床,无缘在他的亲近者们中间守一场灵,我便在他死后的那一夜,独自守了一夜的灵,直到清晨,记下这桩事件引发的一部分追忆与思绪。事后我想:借着这个时机,我也该试着说说:与Fujii~Guruji、与他的那些我曾过从甚密的弟子们的相遇,给了我什么。

五天前的那几条注里,我已谈到 Na mu myo ho ren ge kyo 之唱——它多年来已进入我的生活,是一桩善益。还有那份情谊:由Fujii~Guruji本人、也由他好几位弟子——年轻的与不那么年轻的——所给予的。正是这份情谊,无疑,给了我自他们那里领受的歌唱以价值与美——那歌唱本身,便是对造化中一切活物的一份尊崇与情谊的行动,连他们的人、连我的人,也在其内。

同样,我与日本山妙法寺僧侣与尼众的接触,是我头一遭、也是仅有的一遭,与这样一些男男女女的密切接触:他们主要的、乃至全部的倾注,都投向宗教动机的事务(一如长年间我自己的倾注,投向数学发现的劳作)。这让我有机会确认:一如别处,在超出某种因共同志业(所谓宗教的)而生的亲缘、超出对同一位强势而动人的人格的效忠之外,气质的差异、制约的差异、乃至深层\textbf{抉择}的差异,依旧同样醒目,在人与人的关系里依旧同样起作用。换个说法:某些人力图依某种宗教理想把自己\textbf{塑造}成什么(这里是“菩萨”的理想——佛陀教诲的不倦传扬者),这份努力所抵达的\dmark{}\origpage{767},是或多或少敏感应激的\textbf{姿态},而非一个\textbf{内在转变}的进程、一份成熟。另一方面,采纳一个“信经”(无论多么崇高)并全力投入一项所谓“宗教”的活动,对于惯常那些唯我机制的游戏,似乎并无实质影响。在修道院、女修院、寺院与各种教门的宗教团体里,冲突一如世界的任何角落一样在场。而宗教志业,又常常被当作诸多手段中的一件,用来把冲突打发掉——说服自己:凭信经之力,冲突已然消失。

同样真的还有:在不同的场合,我某位僧侣客人身上,有一种内在的平和与喜悦在放光——我感觉得到,凡接近他们的人都感觉得到;对他们自己、对一切人,都有益。显然,这样一种和谐、圆满、深切相合的状态,与一切“成为这个那个”的努力格格不入——那是一种“无须用力”的状态,一种全然自然的状态。

在我感到过这样一份光辉的四位僧侣身上,我有这份印象:那是他们惯常的状态,已历多年,乃至数十年。Fujii~Guruji本人,尤其如此。对我的另外两位朋友,我在别的场合也见过他们,与任何人一样纠结、一样撕裂。仿佛我认识他们时的那份和谐状态,连同作为其征象之一的那份对事物的自发领会,已化为一纸空文——仿佛它们不曾在他们身上留下任何痕迹。然而我深信:确有一份不可摧毁的“痕迹”,比记忆里记下的一道标记更深——一份其本性为\textbf{认识}的痕迹。一如任何人,这些朋友随时可以顾及那份在他们的存在的创造时刻存入他们之内的认识,让它起作用、让它结果;一如他们也可以无视它、埋葬它,总而言之“装傻”。说到底,这是世间最平常不过的事……

又有一个念头来到心头:这样一种全然自然的状态、这份与自己深切相合,连同伴随它的光辉,却\textbf{并非}那么寻常的东西。相当触目的是:在我得以接待到家里的那一小群僧侣中间——或数日,或数周——竟有这么多的人,我在他们身上找到了这种内在和谐的状态、这种十足意义上的\textbf{力}——谦卑与坚忍、温润与犀利在其中合一的那种力。这莫非到头来确确实实是一个信经、或表达信经的那份“祈祷”的作用?这祈祷——\dmark{}\origpage{768}显见得单凭它自己造不出一个恩典状态——却或许倾向于\textbf{促成}这样一个状态的显现,以及它日复一日的更新?毕竟,全副身心投入地唱一支美的歌,这本身就多少是一份“恩典状态”——而一支歌(或一篇祈祷)单是它的美,就已促使我们“全副身心投入”。

同样真的还有:最美的歌,当我们心不在焉地反复唱它时,依然不起作用——错在我们不曾向它打开自己。或者更确切说:我们这样反复唱的,\textbf{并不是}我们以为在唱的那支歌,我们的灵魂并不以它为滋养——纸做的玫瑰、塑料做的玫瑰不是玫瑰,蜜蜂也不会来采蜜。

\subsubsection*{(f) 信念与认识}
\tocsubsub{(f) 信念与认识}
\label{ch18:local:h76}

\notehead{162}{}
\phantomsection\label{ch18:local:n71}

(1月14日)一周前结束那份反思时,我有这份感觉:我把手指点在了一桩重要的东西上。当夜,我想以一句简峻如刻石的话,把这件“东西”表达在给一条注取的名字里:“无缘之暴力的缘由”(注~n\textsuperscript{o}~159)。我也知道:这道骤然领会的闪电,绝非一份反思的收梢,甚至不是它的终点——那份反思,一个多月以来\fnote{308}{(*)}{确切说,自12月7日那条注“天鹅绒的脚掌——或微笑”(n\textsuperscript{o}~137)起。}正围绕“无缘的暴力”或“无端的暴力”之谜打转。恰恰相反:这份骤然现身的新“视角”,宁可说更像一个新的起点。“移位”的机制——为久远岁月里受下的委屈与损害而生的宿怨或怨怼,从那些被感受为不可企及或“禁忌”的真正的责任者,移往一副\textbf{可以接受的}靶子——这个机制,我起初只是在生活中这样那样孤立的个案里零零星星地认出它,心里默默认它是无意识里一种古怪而飘忽的畸变;如今它终被认作“人心基本机制”之一。同时,它又显为“无缘的暴力”无数令人不安的显形的责任者:无论夫妻之间、情侣之间、父母与儿女之间肆虐的那一种,还是那“匿名”的暴力——它在战争或社会大动荡的年代抵达它的顶点。

我不知这些关联是否早已进入心理学或精神病学的入门常识\dmark{}\origpage{769}(且假定有这样一门“科学”),也不知我在这里就它们所说的,会不会显得像“精神分析的业余爱好者”的凭空幻构。既然我的主旨不是提交一篇心理学的博士论文,甚至不是为某个新旧理论挺枪出阵,而是借我这个人在场所处的种种情境去理解我的生活,那么,我偶尔把手指点在其上的这桩东西、或我看见骤然在此处彼处打开的那些“视角”,它们的“地位”如何,于我无关紧要。我很清楚:无论如何,只要我想理解哪怕最小的一桩东西,我就省不得一份个人的反思——无论在数学里,还是在生活里,在与我的生活这样那样相连的那些生活里。而越是当所要理解的东西乍看便像要挑战理性,当我看到此界内外人人像躲瘟疫一样,拿一堆让人安心的套话把它躲开——情形就更是如此。(而且我觉得:心理学业内人士,一旦自己的这个人直接成为问题,也并不比所有其他人更能例外。)

我很清楚:那道在“给最后一个~i~点上最后一个点”的转弯处现身的“骤然信念”——即“我刚把手指点在了所有无端暴力情境共有的那个机关上”——丝毫免不了我的这桩任务:把这份刚进入意识目光之域的新直觉、这份还远未脱开方才从雾里浮出的东西那片弥散光晕的新直觉,逐件地、针脚细密地加以检视。恰恰相反:这正是头一件该做的工夫;在其中,我已看到一大堆新问题冒头——有专属于这样那样个案的,也有一般的。如果说这份生吞活剥的“信念”里有任何确定性,或者更确切说,有一枚可靠认识的核,那么它告诉我的,绝不是:我刚给这份信念取的表述是“真的”、“对的”,或许没有重要的保留或修改;而倒是:我确确实实刚把手指点在了一桩\textbf{新的}(对我而言)、\textbf{本质性}的事实上,一份看待暴力的新\textbf{视角}确确实实刚刚确立\fnote{309}{(*)}{写下这几行时,一个比较强加于我:与我在1968年孟买会议上提出的关于代数闭链的“标准猜想”相比。它们当时(今天依然)在我眼中,与奇点消解一道,构成代数几何里最灼手的问题之一。提出这些猜想时,我分明感到:一份“新的视角……刚刚确立”——这一次,是在代数闭链之上,在它们与Hodge理论、与Weil猜想的关系之上。最触动我的是:我看到一道通往Weil猜想的进法正在显现——它将是“纯粹几何的”,我的意思是,无须(至少表面上)经由一道上同调理论的迂回。

一如我在别处强调过的(在“肌肉与肚肠”那条注的子注~n\textsuperscript{o}~106$^1$里),这份“新视角”的实在与它的分量,全然独立于这个问题(它仍悬于未来的雾中):这道猜想将被证实为真,还是为假。一道猜想,于我不是一场(或赢或输的)\textbf{赌注},而是\textbf{一记探针}——而无论答案如何,我们从中出来的,只能是“赢家”,我的意思是:带着一份更新的认识。(与“错误与发现”那节(n\textsuperscript{o}~2)里的反思比较。)假定这道猜想被证为假,凭直觉我已看到它的两三个“不那么乐观”的变式——它们随即把它精细化,其中最弱的一个,实质上等值于:一个域上半单Motif的“合理”理论的存在。

清理出这些变式,对一个多少在内行里的人,是一两个下午的练习(或许也是一场通往未知的漫长旅行的起点……)。清理出最初的表述(一如往常,受Serre一个想法的启发,发表于他的文章《Weil猜想的Kähler类似》)——那不是一道练习题,而确确实实是\textbf{一份发现};或者(借用Zoghman~Mebkhout信中的说法——它引于“一次教学的失败——或创造与自负”那条注,n\textsuperscript{o}~44$'$)\textbf{一份创造}。而当Zoghman怯生生地贸然说“我的学生们不太知道创造是什么”——那还是客气的说法;或者不如说,我会讲:他们曾知道,却早已忘了个干净——他们全副身心,都在给一辆灵车推轮子……}。至于这桩新事实、这份新视角的精确而细致的含义,它确切的分量,或许还有它未曾预料的延伸与回响——只消我向它投入必要的工夫,它们自会浮现。这份刚现身的“认识”特别告诉我的是:做这样一件工夫的时机成熟了——更深入一步,去领会暴力,无论如何去领会“无端的暴力”的时机成熟了;我献给这桩任务的每一小时、每一天,只要走到方才现身之物的尽头,都会使我更深地进入这份领会。我不记得:这样一种感觉——一桩新的、本质性的东西正在现身(纵然它还弥散而近似)的感觉,连同能够更深入领会这桩东西的那份贴身信念——曾经骗过我。如果说在我的研究里有一个可靠的向导,替我把投资“安放”到这个或那个方向,那便是这份\textbf{新}东西现身的感觉,与这份贴身的信念——它告诉我,什么时候时机成熟,可以更深入那个被瞥见的“新东西”,去认识它\fnote{310}{(*)}{与“孩子与海——或信与疑”那条注(n\textsuperscript{o}~103)比较。}。

\dmark{}\origpage{771}这并不意味着:每当时机成熟,该向这个方向纵身、该去认识这样一些东西了,我就确确实实纵身而去!早在把全部精力投入数学的岁月,这已不可能:渐渐地,我发觉自己先是十块铁、继而一百块铁同时在火里!\fnote{311}{(*)}{见“炉中百铁,或:干熬无济于事!”那条注(n\textsuperscript{o}~32)。}在禅观里——也就是说,在对我自己的发现里——亦复如是。在意识工夫的层面上,我们很遗憾,一次只能做一件事(而这已不算坏了,只要肯下工夫把它做好……)。对“一百块在火里的铁”之中的\textbf{一块}所下的工夫,的确可以循无意识的神秘途径,惠及所有其余,或至少其中若干——它可以让它们“热起来”,使我们一旦转向它们,它们便更乐于承接意识关注的铁砧上落下的锤击。但还得一开始就知道,从一百块里挑出“对的”那一块——对它的锻打,将同样推进对其他铁块的工夫——那些正与它一同热起来的铁块。

\subsubsection*{(g) 最烫的铁——或转折}
\tocsubsub{(g) 最烫的铁——或转折}
\label{ch18:local:h77}

\notehead{162$'$}{}
\phantomsection\label{ch18:local:n72}

在关于葬礼的反思里,我遇到了不少“铁”——它们要求我在它们上面下工夫,冷热不一,视情形而定。我觉得,它们都在工夫里热了起来,有的多些,有的少些。这些“铁”里的头一块,是对己之轻蔑的问题——在我自己这个人的情形里:它最初是顺带着、为良知起见而提出的,在《收获与播种》最初那个胚胎的边上\fnote{312}{(**)}{见那条注(n\textsuperscript{o}~2)——它指涉1983年6月的一节“他人之无谬与对己之轻蔑”(n\textsuperscript{o}~4)。}。它一直温而不热,直到12月13日(一个月零一天前)的反思——在“正义者的暴力——或发泄”那条注(n\textsuperscript{o}~141)里。我想,这是我一生中头一回,把一场反思——无论多么粗略——献给生活中的那几桩个案:我自己施行并让人承受“无缘的暴力”、“超出理解的暴力”的那些个案。近年来我不时想到它们,但总是顺带一想,不作停留,尤其:不曾为之写下一场反思。

然而,那份不自道其名的暴力,曾深深烙进我的生活——那是一桩关键性的东西,乃至是万般之中那桩最关键的东西:为了理解我的生活,以及一般而言的“生活”、人类的生活,我必须尽可能深地\dmark{}\origpage{772}理解它。可事情确是如此——这只要肯想一想便显而易见——却一直被遮蔽着。它最终浮现出来,像出于偶然,在12月13日之前那些日子的反思的边上——反思延续于合起来题为“天鹅绒中的利爪”(n\textsuperscript{o}s~133--136)的四条注。正是在这些注里,“暴力”头一次在《收获与播种》中被道出其名,并成为一份关注的对象。它一直处于关注的中心,直到如今,或至少直到1月7日(一周前)那条注“无缘之暴力的缘由”。

这个引人瞩目的题目,会让人以为最后这条注是那场持续了整整一月的暴力反思的一种登顶。而它确实也是主要果实之一。然而我心知肚明:倘若骤然有这份新视角现身,有这份贴身信念之感——关于某道骤然瞥见的关联——那是因为\textbf{我这个人自己}也直接卷入了方才现身的东西之中,在那“一生积攒的、一大堆零碎而不相属的印象”之中。所有这些印象里最后、也最新鲜的那一份——当时感到它确实很“零碎”、很不充分——恰恰要回溯到12月13日那场关于\textbf{我自身之中的暴力}的反思。这场反思,在浮泛的读者眼里,或许只是关于葬礼的追查里诸多插笔中的一笔;而于我,此刻,凭回顾,它却显为关于我自己的反思里一个神经中枢般的时刻、一个关键的转折(至少在潜能之中)。而且就在当天,我便感到自己终于起步,迈出了头一步,走进一个此前一直被我避开的方向——它将笔直地领我走向我这个人之内冲突的心脏。那块“温吞的铁”——它作为存照被搁在那里,已历十个月——骤然间烧得通红:只消我停步,鼓风吹打,它便烧成白热,向我显出一份形状、一道讯息。而今天,依然如此。

但显然,此处不是锻打这块铁的地方。在《收获与播种》进程中现身的所有铁块里,它于我确是最烫的一块;紧随其后,是那道与它密不可分、随“无缘之暴力的缘由”一同现身的铁。倘若孩子背上没有驮着一位太成年的东家——他执拗地铆在长线事务与它们强加的“优先事项”上——那么此刻\dmark{}\origpage{773},我必定早已纵身扑向这个方向——它领我走向我自身之内、他人之内冲突的心脏——无须自己掂量!可一如它的名字所示:下单的多半是东家,而不是孩子;决定投资的,也是东家。“恶之谜”因此要等一个更合宜的时刻——东家度假的时候(那是千载难逢的),或者他不再被尖端的“优先事项”压得过重的时候——比如终于写完《收获与播种》!

\subsubsection*{(h) 无穷的链——或传递(3)}
\tocsubsub{(h) 无穷的链——或传递(3)}
\label{ch18:local:h78}

\notehead{162$''$}{}
\phantomsection\label{ch18:local:n73}

但在回到葬礼之前,我想至少记下一周前那场反思所引发的联想之一——一条或许不如其他的显豁、因而倘若此刻不记下便有消散无痕之虞的联想。它与印度教的业之观念相连,与“敌对的兄弟——或传递”那条注(n\textsuperscript{o}~156)里出现的联想同向:朝着一份微弱的直觉——某种“业之守恒律”。

一个人身上那份原初的弥散宿怨——它此后转为表面上“无缘”的攻击性与暴力冲动——并非生于虚无。它是对真真切切受下的深层侵害的回应,尤其是童年最早期受下的那些侵害。诚然,可以认为:这些侵害中有许多,属压抑性质,算不上严格意义上的“暴力行动”——即出于伤人或损人的意图,尤其在父母对孩子。同样真的是:这样一种意图(几乎总是无意识的),其出现的情形,却远比通行共识所承认的多。但在一份业的创造或传递的视角里,当一份“暴力”确确实实发生——施加“一桩恶”,造成“一桩损害”——那么\textbf{意图}或\textbf{动机}(显明的或隐秘的)的问题,或许是次要的。这话我不敢说定。不管怎样,在大多数情形里,一双浮泛的眼睛会有这样的幻觉:所受的那桩“恶”化为一纸空文,它被承受下来,而一经承受,便“消失”了,不留痕迹。而事实是:那些把自身的焦虑、把自身“无法做自己”的无力撒进孩子心里的父母,最终从这同一群孩子手上,直接收到\dmark{}\origpage{774}他们当年播下的东西——这并不那么常见;或者至少,人们得到的印象是:他们从中收获的只是微不足道的一小部分!换个说法:在他们于孩子身上唤起的那份弥散宿怨里,只有一小撮会凝成一份“硬”宿怨、指向他们自己——对此他们叫苦连天,俨然它是天底下最黑的忘恩负义——仿佛这不证自明!可这份积攒的宿怨或“业”的其余部分,并不因此湮灭。它有效地找到用武之地,以一副或许显得无从解释的方式:凭那“移位”的机制,把宿怨挪往临时的靶子;有时是飘忽无定的靶子,有时也是特意搭配、专任其职、可以说备受娇宠的靶子——被孵育一生!

在寻常的时日里,业这桩强烈的工夫——如一个深植于人们生活中的脓疮——在幽暗中完成;而人人以无视它为己任,只肯把它看作此处彼处偶发的“越轨”,相对于那被视为正常与得体的东西。唯有在非常的时日——当战争或苦难肆虐(或在非常的处所,如监狱与疯人院)——这份地下的工夫才迸发出来,自由自在地铺展在大白天里:轻蔑与杀人癫狂的一场恣意的火光冲天,由大言煌煌的旗帜所激励——旗帜飘扬在“英雄的”藏骸所上空,飘扬在一座座赤裸而寒冷的城池上空……



% ================================================================
%  五 · 第四部 (18.3 起, 原书 p.633-929)
%  —— 部扉在 part4-ch18 片段之首, 承前文连续排印
% ================================================================
\markboth{收获与播种·第四部}{}

\partdivider{第四部}{葬礼(III)\\[0.35em]{\normalsize 或四种运算}}{葬礼(III)或四种运算}

\section{3.~最后的义务(或那次来访)}

\subsection{(1) 义务已尽——或真话时刻}

\notehead{163}{}

\dmark{}(2月16日)恰好整整一个月前,我开始了一场即兴的反思\origpage{775}——它由读C.G.~Jung的自传引发。我本打算花上几天,把读后最初那几份强烈印象写下来——而到今天,这些印象我还没有巡完一圈!它们在阅读的进程中丰富、转变——凭这份阅读所触发的工夫、凭我写下读书注的工夫。我恰好来得及巡完Jung青年时代头四章所引发的印象——那几章出自Jung本人手笔。我正准备拿这些印象,与后几章所引发的、乍看并不总是相合的其他印象相对勘。可就在今天要着手时,我意识到:这条岔路(已近百页……)实在不该待在另一条“岔路”里——那条我唤作“阴阳之钥”的岔路,本身就已够长。(一个月前,我还以为这条岔路已近尾声\fnote{313}{(*)}{(3月26日)写这一行时,我还抱着一个印象:我正在起头的这条注,会属于“阴阳之钥”。直到随后几天,我才明白:反思的另一个阶段早已开始。“阴阳之钥”于是以前一条注“无穷的链——或传递(3)”(n\textsuperscript{o}~162$''$)告终。}。)诚然,我关于Jung的读书注,确确实实嵌在阴阳的辩证法里;它们也的确不经意间引我把许多先前仅仅掠过的东西讲清楚了——无论关于我的生活,还是关于生活本身。可我觉得,这仍不足以让人在一道括弧之内,再开一道尺寸如此骇人的括弧——而后一道括弧,本身又位于我自己的下葬那场漫长反思的末章、“葬礼仪式”之内。终于该拾起这场反思,把它进行到底了!

到头来,这些读书注我既不打算收进“阴阳之钥”,也不打算收进葬礼——它与葬礼只有一线相当微弱的牵连。可以把这些注看作一个例证:印证我曾在“表面与深处”与“书写的颂歌”(n\textsuperscript{o}s~101、102)等注里,以一般措辞试图表达的东西。我还在犹豫:是把它们作为第四部分收进《收获与播种》,还是把它们单独成文,放进《数学反思》第2卷\fnote{314}{(**)}{(3月26日)最终,这些读书注将构成《收获与播种》的(不是第四部,而是)第五部、也是最后一部;它多半会收入《数学反思》第3卷(而非第2卷),与另一些性质更数学的文本放在一起。关于葬礼的整套注——它们构成《收获与播种》写作中的“第三口气”,始于去年9月22日;我本想把这一套做成《收获与播种》的第三部分——将被分成两个不同的部分,分别名为“阴阳之钥”与“四种运算”,构成《收获与播种》的第三部与第四部。}\dmark{}诚然,关于Jung生平的这场反思——它实际展开的这副样子——确确实实是我追寻已逾一年的那场漫长反思\origpage{776}的一个不可分割的部分;于我,它就叫《收获与播种》——而我直接卷入其中,一如我在这些注的任何其他地方。因此,仅仅因为这场反思不曾打招呼、就在一场葬礼的半中间冒了出来,又“溢出”葬礼的中心主题稍多了一点,就把它的这一部分从《收获与播种》里分出去,那是做作。

眼下,我要借Jung自传反思里的这道断口,言归正传,把这场葬礼仪式——倘若有此可能——终于进行到底!

如今该到时候了:就我朋友皮埃尔去年10月来我这里的来访,作一份小小的汇报。10月21日那条注(“行为”,n\textsuperscript{o}~113)里,我提过他的到来——其时他已于前一晚携女儿娜塔莉(两岁)抵达。在客人们离开之后(10月25日那条注“失乐园”,n\textsuperscript{o}~116)我写道:“过几天,该到时候盘点一下这次来访给了我什么——一次我已不再指望的来访……”这些“几天”,变成了将近四个月——可我总算到了这里!

我本想就这次相遇作一份“现场”的讲述——它于我,是在发现葬礼、发现其实在与意义的这场冒险中的一个重要插曲。可这一回,我感到自己被一份缄默之虑拽住了:不忍把我的朋友此行留给我的全部印象——繁多而鲜活的——原样端出。诚然,把其中一份印象收进反思(12月26日那条注“否认(2)——或变形”,n\textsuperscript{o}~153),我并无这样的踌躇。但提及某个时刻对某位朋友有过的一份印象,与对那个确切的“时刻”作一份现场的描绘——弥散的印象骤然变得显豁、无可否认的那个时刻——这是全然不同的两回事。后者有点像趁一位朋友不觉得自己被观看的时刻给他拍一张照片,进而不曾确得他的同意便\dmark{}让它传看\origpage{777}。所以,我只限于给出这次来访留给我的几份印象,并且(一如《收获与播种》别处\fnote{315}{(*)}{不过有一个例外——即“玩笑——或‘权复形’”那条注(n\textsuperscript{o}~83)里,我在一次电话交谈中给J.-L.~Verdier拍的那张“照片”。我还记得,为对那桩小场景作“现场”的描绘,我不得不压下心里的一份犹豫——我有点儿觉得,自己给我的昔日学生设了个圈套,这全然不是“我的风格”。当然,他扯满风帆一头扎进了这个圈套——那还是最大、最显眼的圈套之一——我又满心欢喜、对自己满意得很。活该!})绝不拍那些不谨慎的照片!

首先我得把这次来访\textbf{定位}一下。我原本的打算,是去皮埃尔家中看他\fnote{316}{(**)}{这一意向,我在“我的朋友们”那条注(n\textsuperscript{o}~79)的开头、及随后第一条脚注里表述过。},让他读《收获与播种》,包括葬礼。5月初,我曾写信给他,说我希望近期见他,给他读一篇文字——它首先是写给“数学世界里我昔日的朋友与昔日的学生们”的——在其中我“把自己整个放了进去”——“我不记得曾像照料这一篇那样照料过一篇文字”。我当时以为打字当月即可完工,便提议6月上半月去看他。最终,由于打字的延误,还不算给葬礼(如它当时所计划的,即大体上如今葬礼第一部分的内容)作最后润饰的工夫,我的来访一推再推;而7、8月间,皮埃尔不在法国。况且,对于那部我一意要亲手交给他、要他在任何他人之先阅读的作品,消息公布时他不曾显出分毫好奇。最终,6月间,我把《收获与播种》的第一部“愚妄与更新”寄给了他,想着:在把葬礼劈头砸给他之前,让他先有个了解——说不定我对我自己的反思会在他那里触响什么机关,触发点什么——谁知道呢!那时我已病倒十来天,反正近期去巴黎已无从谈起。

然而我急于让他读葬礼——在其中,皮埃尔是以触目惊心的方式被卷入的;我本想让他赶在度假出发之前,来我这里一读。怀着这样的心情,我在6月底把完整的《引言》连同葬礼的目录寄给了他——我以为这会给他一记冲击,他也会\dmark{}有心赶在出发之前来看我\origpage{778},细细了解我对那场出名的葬礼要说的话、以及在这场葬礼里分派给他的角色。可事实呢,直到8月底,我再没有得到过他的一点音讯——以致我疑心他是否收到了我的寄件。好一场悬念!他回来后的第二封信(署8月25日)里,终于对引言与目录说了几句——措辞在我看来闪躲之极。“我过去觉得,你不大知道你的‘孤儿们’曾被怎样的爱环绕着……”,他这样写道,并随信附上一份评注书目以为佐证——那分明是一份显见的善意,要驱散一桩他显然感受为令人痛心的误会。随后的下一封信(9月12日)里,他宣布10月7日将迁居普林斯顿,并说会设法在此之前抽空来我这里一趟。接着又音讯全无,我以为他已动身去了普林斯顿——并没有:往IHES打电话,我得知他的行期推迟了。而一周之后,正当我已不大指望短时间内见到他——他本人已活生生站在面前,还带着小娜塔莉!

(2月17日)相遇是在一种气氛里进行的——照一切表象,它是再平和、再友好不过的。一个凑巧在场的浮泛观察者,会指天誓日:皮埃尔正在啃一篇数学手稿,时不时把他的观察与批评——一个很“在行”的数学家的建设性意见——提交给我。对皮埃尔本人来说,事情本该不言自明:他是专程赶来的(出于对毕竟是他的“老师”的我的敬重),牺牲了一位大忙人宝贵的两天,为的是竭己所能,驱散一桩不幸溜进我心里去的、恼人的误会——天知道是怎样一场倒霉的机缘巧合促成了它。他的诚意与我的诚意,当然都在一切怀疑之上;连提也不必提起,事情就那么明摆着。而他的角色,是就一切物质性的细节之点给我以启发——那些在我的注里显得不甚分明之点,或我或许出错之点。他的阅读推进到哪里,他随手记下自己的观察,列成一张单子,在动身那天交给了我——我有这份好脑筋,当即用几个关键词把它好好记了下来。况且,他确确实实在两天里读完了葬礼I的大部,至少读完了(借目录、借正文内部的交叉指涉而定位的)一切直接关涉他本人的注。好一份\dmark{}业绩——要是想一想\origpage{779}:写这些注,我可是花了整整两个月的全时……

这两天里,小娜塔莉是最乖的孩子里最乖的。我几乎可以说不曾听见她的嗓音——无论是说话、喊叫还是啼哭。她在我家里似乎并不难受,却不怎么露面。至于她爸爸,他是十足模范的爸爸——要喂饭、要遛弯、要哄一位毫不苛求、绝不添乱的小姑娘睡觉,他总在需要的时刻随叫随到。他带着她来,他对我说,是因为筹备迁居普林斯顿的种种大事之后,妈妈忙着收拾家务,腾不出手再照管娜塔莉。但在这条实际而无可奈何的理由之外,我相信自己感到了另一条理由,它留在没说出口的地方:小姑娘的在场,给我朋友所畏惧的那场相遇的气氛,添了一笔柔和——纵然他或许连在心里也不愿承认这份畏惧。而这份在场,同时又是一个鲜活的、夺目的标志:表明他是在怎样的默契心境里,在迁居美国的手忙脚乱之中赶来的——一份昭然的诚意、一份同样显见的好意的心境。

至于我,我毫无搅扰我朋友的心思,要逼他触碰什么——他若感到愿意深入哪一个问题,我便随时奉陪。可偏偏,他最要紧的执着,是对我注里检视的众多情境中的任何一个,都决计不\textbf{进入实质}——在那些情境里,他作为数学家的操守(或干脆说他的操守)被明明白白打上了问号。一个听到我们谈话的旁观者——那谈话有时甚至岔进数学的讨论(我们之间,这已有三年多不曾发生了\fnote{317}{(*)}{关于Deligne与我之间一切数学交流的中断,见“两个转折”那条注(n\textsuperscript{o}~66)。}!)——决计猜不到:在我朋友评注的文本里,竟可能有什么东西,稍稍带上个人色彩地把他卷了进去。至于我,我感到我的朋友死死抓住这副虚构——那费劲维持着的“最佳世界里的昭然善意”的虚构。他小心翼翼避开一切可能让它爆裂的东西——凡是会显出:他想在我们之间、排除万难设立的那份默契“共识”根本不是实在,而恰恰是一场虚构、一根可供攀抓的“稻草”……

\dmark{}这两天里,我分明感到这情势有多假:那份平和温良、乐呵呵的门面\origpage{780}底下,压着多少焦虑。正应了那句话:悬梁者家里的绳子——人人心里想着,无人提起!可我到底还是就此说了一句——我想那是动身那天,午饭后。大意是说:我们这场相遇,竟走了茶馆客套的路数,我很有些透不过气;毕竟,在他正在读的这些注里,单在早他近四个月就该收到的引言里,我已经用相当分明、相当有力的措辞,谈到好几桩\textbf{出自他手}的\textbf{行为}。对这一点,他当真无话可答吗?他回答我,眼神模糊,挂着一丝惨淡的、近乎可怜的浅笑,说他在尽力“自保”——(就我记得的)却不说他究竟想从什么里头“自保”。无疑,我的追查,该被他感受为对一个生命的粗暴侵入——那生命至今于他想必是最安顺、最没有问题的,其中的一切想必甚至显得出奇地\textbf{驯服};驯服到这般地步,或许,以致他终于忘了:事情本可以是另一个样子。\textbf{承担}他自己置身其中的情势——就是说,仅仅与之直面,如其所是地检视它——那将是他对自己、对世界的视景的一场幅度巨大的颠覆,是自我那副僵硬结构的一场崩塌;大多数人宁愿死上千百回、(若做得到)把世界投入火与血,也不肯冒这样一个跃入未知的风险。正是为了躲开这一切,我的朋友执意(而且多半,今天依然执意)“自保”。

我本不该为此惊讶——这类剧本,我见过它重演成百上千次:那是面对事物实在的大恐惧,而尤在其上的,是面对内在更新之风险的大恐惧。我本不该惊讶,诚然;可每一次,我仍旧一次次惊讶——当我看见最刺目的显明遭到拒斥,当人们受着千般熬煎、也施人千般熬煎,只为躲开那个我凭可靠认识分明知道的东西:万般福祉中最大的福祉……

不管怎样,在我这番想“脱轨”的不成功尝试之后,谈话转了个弯便止住了。这两分钟,我想,是那两天里仅有的一次\textsuperscript{317(*)}——我们的谈话取了一份\textbf{个人的}转向\dmark{},或说出了什么,越过了那份“共识”的虚构——尽管相反的证据明摆着,那虚构仍被维持着\origpage{781}!我怕,一如常有的情形,这一回我身上那份深情却又直截了当的“圆融”——它本可以帮助我的朋友,给一种气氛祛魅——却不曾到场:那气氛,尽管表象相反,早已紧绷到极点,而且紧绷已有数月。我把自己限在家务、园子与写作的诸务之中,任我的朋友自去读书;共同进餐时也一样:在我之内,对我的年轻朋友,悬着一份无言的\textbf{期待}——一份\textbf{回答}的期待:对我借他手中那部文本所说话语的回答。这份期待,他不可能感不到——而且他心里清楚:他那几条可怜的物质性说明,可“回答”不了它!无疑,倘若我以这样那样的方式抢先把话挑明——哪怕从一场他活该受领的、精心的臭骂开头——反倒会让他松一口气:好歹在空无之处,立起一份\textbf{接触}。

诚然,在已往的十五年里,每当我试着同他提起一件切身的、于我紧要的事,我撞见的总是全然的沉默,或(若当面)那套规矩的惊诧语调——最纯粹的“天鹅绒的脚掌”风格。这个游戏,我自然不再想玩;自1981年的“转折”起,我早已离场,并无回头的意思\fnote{318}{(*)}{见已引过的“两个转折”那条注,n\textsuperscript{o}~66.}。但同样真的还有:这一回,在我们之间,分明有一个独一无二的“时刻”,它或许值得对一条规矩(或一种已成第二天性的习惯……)作一回通融——那条规矩说:不要去顶撞他人对谈论某事的推拒。有时候,“稍稍强按着头”或许是有益的(在一定的分寸之内)——就像带一个怕牙医(那怕是无来由的)的孩子去看牙医……

我说这些,可不是要怜惜我那位可怜的朋友皮埃尔——怪他没有从我这里得到他本可以指望的全部善意鼓励,诸如此类!毕竟,我有我的限度,一如人人都有,这很正常;况且,替那些把自己放进了某种情势(哪怕是身不由己)的人缓冲震荡——那情势迟早会以这样那样的方式落回他们头上——这未必是我的角色,更不是我的义务。

\dmark{}何况,10月22日晚,把皮埃尔与娜塔莉送到奥朗日车站之后,我毫无“白遇一场”、“错失良机”之感。我还不曾天真到指望有山呼海啸的收获——两个人就一桩把他们彼此深深牵连的问题谈到实质,这是何其稀罕!没有对话,这是明摆着的——可我仍感到自己学到了许多。已经有那些“物质性细节”了,其中不止一条相当有滋味,它们就某些发生过的“行动”的“脚本”问题、及其情境,给最后的~i~点上了最后的点。本注的续篇里我还会回到这里\fnote{319}{(*)}{见“把一个个~i~点清”那条注(n\textsuperscript{o}~164),它紧随本注之后。}。更要紧的是:在这两天里,我以新的眼睛观察了我的朋友——借我在葬礼反思里对他新学到的东西。可以说,我与他“重新认识”了一遍——在他与我的关系里、与物的关系里、与女儿的关系里……那一章仍是保留地——正是在这里,我今天这几条注开头所提到的那份自然的分寸,于我成为必须。

而在领会葬礼的视角里,还有另一条理由——比前两条更精微——使这场相遇的举行成为要紧。我相信,从我决意去巴黎会见我的朋友那一刻起,我就感到了这份要紧;可当时我多半说不出为什么,只知道:当有牵涉双方的大事之时,与当事人当面交谈,总是要紧的。可恰恰,我们不曾谈这些事——而我却有一种印象:关于葬礼的\textbf{实在},我尚该学的东西,我已经学到了。

我也可以这样说。在这场相遇之前,构成葬礼的那一整套情境、事实与举止,显得何其不可信、何其滑稽、何其谵妄:尽管几周几个月间积攒下来的、一切可触而无可否认的“证据”,尽管我已献给它的近三百页注——在我心底的某个地方,\textbf{我依旧无法相信}\fnote{320}{(**)}{面对我们健全官能的见证——当它们过于猛烈地撼动通行共识、或我们珍爱的看法时——这份\textbf{不相信},在“中国皇帝的新衣”那条注(n\textsuperscript{o}~77$'$)里已被唤起过。显然,写那条注,于我是抵达(至少部分)这份面对显明的不相信的手段:把手指点在这份积习的反应上。可这样做时,我把自己与这份不相信拉开距离——把它呈现为凡俗众生(成年人)的不相信,而把自己认同于“相信自己眼睛见证的小小孩”(“纵然他所见的相当闻所未闻、从未见过、被所有人无视与否认”)。这多半正是我写那条注时无意识的用意:与一份对我自己官能的不相信态度拉开距离,与一份“随大流”的群居本能拉开距离。这样的态度与本能,确确实实存在于我身上,一如在每个人身上;可(一如所有人)它们多半停留于无意识。于是,这像是把那使我自我疏离的东西从我心里驱除的一场尝试——而这尝试的主要结果,我相信,是把我想与之拉开距离的东西\textbf{压进了}无意识更深处。那份阴险的怀疑——它像一道暗藏的裂隙,在我对事物的认识里运作——并没有因此消除;那份倒霉的不相信,也谈不上“超越”(“至少部分”,原文如此)!\par 我再次发觉:在反思的那个时刻,反思还停留在低于我所称“禅观”的地方——禅观,是这样一场反思:其中那些幽暗而精微的内在运动(如这份隐秘的不相信,以及我写那条注时真正的动机——“驱除”这份碍事的不相信)始终处在一个警醒关注的对象之下。}!

而这样的事临到我,也\dmark{}并非头一遭——远非如此\origpage{783}:一份执拗的怀疑盘踞一时,那是抵抗旧物入档的残余——一副旧的事物的视景,它往往更舒适些,或更合于通行共识,比起接替它的那一副。有时,这份怀疑也不只是惰性的表达——对事物视景中创造性改变的惰性——它同样映出一个健康的成分、旧视景里\textbf{实在}的一面:那一面,或许曾被连同其余,过急地抛下了船!不管怎样,一如怀疑现身的每一次,该做的好事是意识到它(这并不总那么容易,鉴于那根深蒂固的、要“封住”不受欢迎的怀疑的反射),而做了这个,再细心检视它。我不记得有哪一次:我留心检视过一份怀疑,却不曾从中学到些有趣的东西(甚或,对我重要的东西),而且足以让一切怀疑烟消云散\fnote{321}{(*)}{更确切的说法是:怀疑转变成了\textbf{认识},认识取代了它的位置。这与人们驱赶(或“超越”!)一份怀疑时所发生的,毫无共同之处——后者使怀疑从目光里消失,而它其实已退隐(或遭流放……)于更深的、看不见的层。它离被解决(并转变为认识)比以往更远;它一如既往地运作着,像一道暗藏的\textbf{裂隙},一份不安,一桩仍被回避的工夫的信号。可与上一条脚注的评注对读。}。任何怀疑,都是一桩有待去做的工夫的无可置疑的信号。

\dmark{}就本个案而言——即我那份不曾言出的、全然无理性的怀疑:怀疑一场所谓的“葬礼”\origpage{784}的实在本身——我该承认:在这场与朋友的相遇之前,我甚至不曾抵达一切工夫的头一个先决条件:我不曾真正意识到它。它停留在一个单纯弥散的不安的状态,而且不自道其名——错在我没有审问它!不安及其意义,我是事后才发觉的——恰在它凭这场与朋友的相遇消散之时。此外我相信:无论他采取怎样的态度,这效果都会产生——无论是一种殷勤合作的姿态,急于向我提供一切缺失的“物质性细节”(事情正是如此),还是相反地——不妨说——一种激烈的、或许狂怒的否认,否认那些最昭然的事实。无论哪种情形,葬礼的\textbf{心灵}实在都不可能不向我显形——这一次,凭直接的知觉(而非凭文献的“归纳”、凭我所知其他事实的印证等等):眼看着我的对坐者干脆\textbf{无视}“最佳可能世界”版本里那些荒诞绝伦的荒诞——而恰恰是那些荒诞的巨大,当初使我在心底深处怀疑过那场葬礼的实在!

只举一个例子:竟需要Deligne本人亲口告诉我,他确确实实是从Zoghman~Mebkhout本人嘴里得知“好上帝定理”的——却不愿在自己与Beilinson、Bernstein合写的文章里指涉他\fnote{322}{(*)}{见“当值无名之辈与好上帝定理”那条注(n\textsuperscript{o}~48$'$)与“不义——或归来的意义”那条注(n\textsuperscript{o}~75),以及紧随后者的那些注——它们与它一同构成第VII行列“研讨会——或Mebkhout的层与反常性”。},出于对Kashiwara的(!)顾虑——(作为门外汉)他拿不准两人在这条定理里各占多少\fnote{323}{(**)}{当然,在Beilinson、Bernstein与Deligne发展那所谓“反常层”(为免称它们“Mebkhout的层”)的形式主义的文章里,指涉Kashiwara的,与指涉Zoghman~Mebkhout的一样,全无——而它整个是从那位始终不被提名的Mebkhout的哲学出发的。况且,Kashiwara在好上帝定理(别名:Mebkhout)里的角色,Deligne比我清楚:Kashiwara的可构性定理,使Mebkhout得以定义一个函子——它从一个由“连续”系数(微分算子的复形)构成的三角范畴,导向另一个由“离散”(可构)系数构成的范畴——这件事,世上在他之前无人想过去做,更无人猜到会有一个范畴等价。这恰恰是我花了十年左右(1956--66)所发展的对偶形式主义里那“缺失的环节”;而我的那些上同调家弟子们——Deligne领头——在我1970年离开之后,争先恐后把它埋掉了。}——竟需要我亲耳听\dmark{}Deligne说出这样的话,才亲眼看见这一奇怪的组合:细节上的诚意,与实质上、要害处异乎寻常而昭彰的自欺\origpage{785}。我没有觉得有必要让我的朋友留意他那别致的做法(它在“牵线人”那条注(n\textsuperscript{o}~75$'$)里已被揭示——他明明读过它!):为了那个“本该在其中找到位置”的结果,他怎样安排,好显出:那条定理光彩的作者,不是别人,正是他(或至少,那篇显赫文章的三位作者之一)!对这桩怪事——那个我唤作“反常层会议”的会议,主要是在Mebkhout前些年所发展的工作与哲学的余波里成形的(对此Deligne并不摆出争辩的样子\fnote{324}{(*)}{Deligne只限于摆出对我的视景稍加细化的样子,说依他看,MacPherson的想法在1981年6月吕米尼会议(即所谓“反常层会议”)上的影响,比Mebkhout的还要大。我身在局外,无从就事论事地讨论这个;而这显见得是个细节之点,几乎减轻不了所发生之事的巨大。况且Deligne并不否认:没有Mebkhout前些年的开拓性工作、没有他在全然孤独之中发展的那套哲学,无论那场会议,还是它所标志的代数簇上同调理论的那场大规模更新,都不会发生。

我得以理解到:MacPherson那个关于簇的“相交上同调”的想法——它独立于Mebkhout的想法而发展——一直有些束之高阁,直到Mebkhout的“哲学”以一道新的、意想不到的光把它照亮(这是Deligne发现的)。那是Mebkhout层理论(被错称作“反常层”,顶替了某个会议……)的全速起步。这个起步,是那场会议的主要\textbf{事件},而且(似乎)是我们理解代数簇上同调的历史上的一个转折。这份新理解的拱顶石,在我看正是那道“好上帝定理”——它自六十年代初便“悬在空中”,而我、以及后来的Deligne,都不曾把它提炼出来。}),而他的名字却严丝合缝地不见于登在《Astérisque》上的会议文集\fnote{325}{(**)}{“严丝合缝地缺席”这个说法,照字面为真——至少对文集的第1卷(由引言与Beilinson、Bernstein、Deligne的文章构成),它是会议的主干部分。第2卷里两篇文章的书目(一篇出自Brylinski,一篇出自马尔格朗日)里,有两处对Mebkhout的“暂停!”式提及——没有一处涉及“好上帝定理”的发明权。}——他也没有任何解释可提。他看那情形,像看一场不巧的\textbf{巧合}:他与任何人都不负责任。总之,我所称的葬礼,在我朋友皮埃尔那里,化约为二十或三十来桩这样的“巧合”。

\dmark{}我在这里重新遇到一个我在他身上熟知的把戏——而且不止在他身上:一个装傻的把戏,带着世上最无辜的神情,笃定自己永远不会被抓个正着\origpage{786}。而确实,有些时日了,我不再浪费时间,去说服什么人(比如)相信:某些所谓的“巧合”并非单纯的巧合。有时候指一指显眼的东西是有用的;可这做过之后,再去说服什么人相信:它们确确实实是\textbf{东西},而非胡思乱想——“你们到那儿找什么呀!”——那就是浪费时间。想说服自欺,是浪费时间:它是自觉的还是不自觉的,一样;它摆出愚蠢的脸孔还是精明的脸孔——还是一样。

然而在我们的相遇里变了的,给我的朋友添上一份他竭力控制、竭力遮掩的焦虑的,是:这一回,这把戏不再局限于两人之间、无人得见的那点无后果的小运动——而且对手还是一个\textbf{亡故者}!这一回,牌摊在桌面上,\textbf{这是一场公开的戏}。赌局已开,赌的是那位大名鼎鼎的教团将要认领、将要容忍什么。诚然,十或十五年来,它已经容忍了许多、认领了许多;它或许还会继续,谁知道呢?一如我的朋友皮埃尔,它或许并不在乎再多二三十桩“巧合”……

(2月18日)10月22日晚,当我终于把皮埃尔与娜塔莉送到奥朗日车站,我觉得自己蠢透了。皮埃尔一副样子:照着他给自己定的日程,一板一眼、巨细无遗地尽完了他的义务;而我感到一阵郁郁的挫败——这场终于举行的相遇(谈论它已有数月),竟什么也不曾说出、什么也不曾触及。

天黑了,小家伙(坐在后座)该睡了——紧赶着开车到车站,得四十来分钟。有很长一段时间,谁也没说话。是我先打破沉默的——那份在我心里寻找出口的不满推着我;那不满,无疑是对我自己,胜过对任何他人。可这不妨碍:我此行就是存心要撩一撩我的朋友。我对他说:我还没跟自己想明白,要不要对斯普林格出版社提起一桩诉讼,逼它把那部海盗卷册SGA~4$\frac{1}{2}$\dmark{}从流通中撤下来\origpage{787}——它登在《Lecture Notes》丛书里\fnote{326}{(*)}{关于这部卷册,尤见“同伙”“白板”“绿灯”“颠倒”四条注(n\textsuperscript{o}s~63$'''$、67、68、68$'$)。}。说实话,连这念头当初是何时掠过我的,我也说不大清了;我这会儿把它随口抛出来,也算顺带探一探我的朋友(德语所谓“ihm auf den Zahn fühlen”——探他的口风)。他没怎么反应,说真的;与其说我在与他对话,不如说我在独白——重拾一根我早在4月或5月就丢下的线头。我随即发觉,顺着它想下去:一桩单纯的法律角力,说到底成不了什么气候——把SGA~4$\frac{1}{2}$从它眼下的题名与装帧底下撤出流通,这事情的意义,除非创议出自别人而非我:或出自斯普林格,或者——谁知道呢——更好,出自Deligne本人。我只得接着说:由他Deligne来做一个这样的公开姿态——说到底,是对某些针对我的行径的一种悔罪——这在我看算不得奢侈。它会给一片亟需净化的气氛消消毒!

我的朋友以音节应答我的独白,此处彼处点缀一二。他让人听出:把手里那一大堆SGA~4$\frac{1}{2}$的存货贱价处理掉,斯普林格未必乐意——我回敬说:换一副封面就够了,它有先例,毫无问题\fnote{327}{(**)}{那是我与斯普林格出版社的第一次不幸遭遇:它出版了Hartshorne的讲义(据一门我在其中发展了局部上同调形式主义的课),标明的作者却是Hartshorne。那是《Lecture Notes》的n\textsuperscript{o}~41卷“局部上同调”,不得不更换封面。当时斯普林格还有这份礼貌:为这场错认致歉,并勤勉地补救错误。此后的家风,已大不相同……},那花不了它几个钱。而就算它贱卖存货——《Lecture Notes》一千多种里的一种,亏点损益,算得了什么!更别提Deligne:假如他真在乎,那几百万旧法郎——填补亏空所必需的——他总拿得出来……

有些话我没有说出口,可它藏在话里(而且肯定被听出来了):这里头攸关的东西,其价值或许超过我们中哪一个的一两个月的薪水。可我到底还是说了:在这类事情里,头一位要紧的,不是看\textbf{怎样}去做一件事(或者反过来,把去做的\textbf{障碍}一条条数出来),而是首先把自己\dmark{}\textbf{想做}什么弄明白\origpage{788}。这一步做了,其余便成了后勤问题,自然“跟上”(当它肯“跟上”的时候)。怪我这位寡言的对话者不肯把真实的感受说破,我把这当作不言自明:他分明看得出,把一桩亟需“清扫”的情势“清扫”一下会是好事——只是他对能为此做什么拿不定主意,无疑是为了“保全面子”,诸如此类。其实我压根“没搔着痒处”!我最终察觉这一点,是在站台上等火车的时候。那是当Deligne回到这个话头,带一点讪讪的神气对我说:最终他宁愿由\textbf{我}去联系斯普林格,谈SGA~4$\frac{1}{2}$的事。显见得,他无意掺和其中,甚至此刻都不肯就这部著作的命运表个态——而它明明是他著述的(自然,还有我的“协作”\fnote{328}{(*)}{关于这一点,见已引过的“颠倒”那条注,n\textsuperscript{o}~68$'$。}。)。

直到那时我才明白:我们在路上时,我的那番思量,的确是一场独白——而对我的朋友皮埃尔,那桩“SGA~4$\frac{1}{2}$—SGA~5”的“行动”里或许有什么不太“合规”的东西,仍旧不明不白。而我偏偏在这万般话题里的这一个上岔了出去,为自己的不满寻找出口,这决非偶然。正是这一行动——它连着一场对一件美好作品的按部就班的屠杀,在那作品里,我放进了自己最好的一部分\fnote{329}{(**)}{见“屠杀”那条注(n\textsuperscript{o}~87)与紧随其后的两条注。}——给我的触动最深:凭其中一口气——暴力的气(在屠杀里)与安然放肆的气(对着被屠杀的东西)。而我再次被触动,凭那份做作(我在朋友身上知之甚稔):仿佛这一切与他毫不相干,他与我对这个那个所怀的“想法”毫不相干。

火车很快就要到了;而这是头一回:借一份终于浮出水面的\textbf{情绪},我要用几句话,进入一桩于我紧要之事的\textbf{实质}。要把此刻的感受当面说出,用不了多久。那是真的情感——一个人的体面之感受了伤,出手的是他所亲近、而那亲近者戏耍了他——这不再是边上那点半科学化的文学,可以拿支铅笔一丝不苟地随读随注。

\dmark{}他一时愣住了,却仍竭力撑住他那副不动声色的架子\origpage{789}。我只得对他说了类似这样的话:“那么,你觉得那是个好题名了——SGA~4$\frac{1}{2}$,暗示这是排在SGA~5\textbf{之前}的东西——而在那里,早在十一年前,你学到了那些直到今天你天天在用的数学!”他回答我,语气像在背书:他把它题作SGA~4$\frac{1}{2}$,只是为了标明一种\textbf{逻辑}依赖关系,而非时间在先。

就这样,我有幸亲耳听到、并出自当事人自己之口:这出“闹剧”——它大到何等地步,以至于当初我白纸黑字读到它时——先是在他的笔下(在“SGA~4$\frac{1}{2}$”里),随后是在Illusie笔下(在那部名叫SGA~5的卷册里,它“合乎逻辑地”跟在我那位前任的卷册之后……)——竟难以相信自己的眼睛所见证的东西!

我只得对他说:他与我一样清楚,SGA~5整个“站得住”,不需要任何先决、任何猜想,它在逻辑上、在任何意义上都不依赖那些后来的添加。我对他说话时,直直盯着他的眼睛,他回答时也如此。他以同样平板的嗓音重复他的功课:SGA~5在逻辑上确乎依赖SGA~4$\frac{1}{2}$——可我从他闪躲的眼睛里看见:事情究竟如何,他与我一样清楚。他的眼睛,身不由己地,比他的嘴诚实。

于是,在我们之间,“真话时刻”到底还是来了——可没有任何仪器,相机或录音机,能把它侦测出来。当时,唯有他与我知道正在发生什么。

火车在随后几分钟里进了站,我记得。反正,对那一天来说,彼此再没有什么可说的了。

\subsection{(2) 把一个个~i~点清}

\notehead{164}{}

(2月20–21日)为给Deligne上一次(去年10月)来我家的访问作一份回顾的收尾,我想在这里把他乐意就若干点提供的说明清点一遍——那些点,在我关于葬礼I的反思注里或模糊,或有误。这也将给我一个机会,补充另外若干说明——由Deligne所提供的那些引发出来的\dmark{}\origpage{790}。

\textbf{I~Motif(《Lecture Notes~900》卷)。}

1. Deligne向我说明:LN~900卷\fnote{330}{(*)}{关于这部“值得纪念的卷册”的详情,见“一个梦的回忆——或Motif的诞生”与“下葬——或新父亲”两条注(n\textsuperscript{o}s~51、52)。}的主旨,是在一个数域~$K\subset\mathbb{C}$上——$\mathbb{Q}$的有限扩张——发展一套交换的、Motif的类域论。换言之,即确定“$\bar{K}$在$K$上的Motif Galois群,使之交换”。就此我记起:率先提出这个问题的正是我(事出有因!),在六十年代末。对一个选定的Motif概念,借“Motif范畴上的Betti自由函子”、借$K$在复数域$\mathbb{C}$中的给定嵌入,这个问题有一个确切的意义。其实,我当时给自己提出的问题还更一般些:确定“$\bar{K}/K$的\textbf{亚交换}Motif Galois群”——由完整的Motif Galois群导出:使之交换的,不是整个pro-代数群,而只是它的单位元分支。该当得到~Gal($\bar{K}$/$K$)~的一个全然典范的扩张:以诸乘法群~$L^*$——$L$取遍~$C/K$的有限子扩张——所对应的那些环面取投影极限而成的pro-环面为核。我记得Serre对这个问题极感兴趣,可无论他还是我(连同Deligne——我当然也把他拉了进来)都没能即席给出一个说得过去的“候选者”。这个问题随后坠入全然的遗忘,一如它所由出的Motif的yoga。这份沉默直到1979年才被Langlands的文章打破(Deligne在他1984年5月28日的信里,以一份Motif评注书目向我提示了它\fnote{331}{(**)}{即R.P.~Langlands的文章“Automorphic Representations, Shimura Varieties and Motives. Ein Märchen”(Corvallis会议),载Proc.~Symp.~Pure~Math.~33~(1979), AMS, vol.~II, pp.~205--246.}),在那篇文章里,我的Motif Galois群的想法头一次在文献里得到显式的表述。由于我无缘收到这篇文章的单行本,我不知道文中是否提到我这号人。Motif在文献里的下一次露面,似乎就是LN~900——在那里,凡指向我这个人、说我与该卷的主题及其主要问题有任何干系的提及,一概不见\fnote{332}{(***)}{(4月8日)我近来得知:Motif被用于Deligne1979年的一篇文章里(它登在与前一条脚注所引Langlands文章相同的那一卷里)。
(5月12日:这条“结尾”成了子注“先期起棺”(n\textsuperscript{o}~168~(iv)))。}。

2. Deligne向我说明:与我此前所认出的\dmark{}(凭某种“文风家风”……)相反\origpage{791},LN~900里Deligne--Milne的那篇文章——它从头重述了N.R.~Saavedra所发展的塔纳卡范畴的Galois理论——几乎全部出自Milne之手\fnote{333}{(*)}{关于Deligne--Milne的这篇文章,见“下葬——或新父亲”那条注(n\textsuperscript{o}~52),亦见后一条注“白板”(n\textsuperscript{o}~67)里的评注。}。Deligne还向我解释了Saavedra著作里的那个错误:它使人(若想拥有纤维函子的Galois--Poincaré理论的形式主义)不得不强化Saavedra对所谓“塔纳卡”范畴的定义。Deligne--Milne文章里的工作,确实仅限于做这个调整——错误一经发现,调整便是明摆着的。这倒顺带提出了一个极有趣的问题:对那些“真正的”塔纳卡范畴的$\otimes$-范畴,给出一个可操作的内部刻画(不妨更醒目地称作$\otimes$-\textbf{Galois--Poincaré范畴},因为正是为它们,才能发展一套Galois--Poincaré广群的理论\fnote{334}{(**)}{“(Galois--Poincaré)广群”这一称呼,好处在于暗示它与拓扑空间或topos的基本广群概念的近亲关系。不过按技术讲,‘(Galois--Poincaré)gerbe’的称呼会更贴切。所涉的是“纤维函子”的gerbe:它不只定义在所考虑的$\otimes$-范畴的基域~$k$上,而且定义在~$k$上概形的fpqc~site的任意对象上(对这site中形如~Spec($k'$)~的对象——$k'$为~$k$的扩张,乃至~$k$的\textbf{有限}扩张——予以特别的关注)。})。这个问题在该文中未曾触及,而且至今未得到令人满意的解决。显见得,那里要紧的不是提出或解决有趣的数学问题,而恰恰是给Saavedra的文章备下一份替身文献。(关于这一点,见“白板”那条注(n\textsuperscript{o}~67)的结尾。)\fnote{335}{(***)}{(5月12日)近来读到了所引的Saavedra那部书,如今看来:这部书,连同这个概念的名字(“塔纳卡范畴”)——那概念系我约1964年引入,并给了这部书它的书名——是一场\textbf{骗局}。我将在“第六颗钉子(钉进棺材)”那组注(n\textsuperscript{o}s~176$_{1}$至~176$_{7}$)里,对它作一番细密的拆解。}

3. 在葬礼I里,我不止一次强调:Deligne六十年代末所发展的Hodge-Deligne理论,只是迈向“$\mathbb{C}$上有限型概形上的‘Hodge-Deligne系数’”理论、及为此种系数的“六种运算的形式主义”的头一步。我曾(并且依然)深信:倘若不是Deligne存心作对——冲着我引入的某些强力观念(如六种运算的形式主义)——Hodge-Deligne理论今天早该“瓜熟蒂落”。Deligne强调:单是在~$\mathbb{C}$上有限型概形上定义一个Hodge-Deligne系数的范畴,就撞上了严峻的困难\dmark{}\origpage{792},而他自认无力克服。(他愈是如此,就愈该从理论一开始,就把这个问题明白地表述出来,连同与它休戚与共的、为此种系数的六种运算的形式主义——而Deligne始终避免这样做。)照他的说法,Mebkhout的视角与Mebkhout的层\fnote{336}{(*)}{这些层,即Deligne当年以“反常层”之名引入的那些。(关于这一点,见“不义——或归来的意义”与“反常性”两条注,n\textsuperscript{o}s~75、76.)他并不别扭,在我们的交谈里乐意把它们称作“Mebkhout的层”……}该当提供一条通往正确定义的进路。(而倘若没有这份存心作对,Deligne当然不会等Mebkhout去发展那套哲学——它逆着他的前辈们而行——并拿它去做那件显见得根本性的工作:十五年来它一直撂在那里,直到今天,除了我在《收获与播种》里亲手指出,文献中连提都无人提起!)

4. 我记错了:我曾以为,一个Motif的“按权过滤”是我引入的——它(对每个~$\ell$)反映为该Motif的$\ell$-adic实现上相应的过滤(以Frobenius特征值的绝对值为项定义的过滤)。其实,Deligne提醒我:我只以“虚拟”的权概念工作过(那等于以虚拟的Motif工作——某个合适的“Grothendieck群”的元素……)。这个重要的事实——我工作时所用的虚拟概念,该当对应一份典范的过滤,按“权递增”——是Deligne发现的\fnote{337}{(**)}{使Deligne相信一个Motif必有这样一份过滤(且必唯一)的启发式理由是:存在阿贝尔簇被环面作的非平凡扩张(其Motif的~$H^1$由此给出一个权~2~的Motif被权~1~的Motif作的非平凡扩张),但反方向不存在。这听来单薄——可我自己也差不多当场就被说服了:它太诱人,不可能是假的!更严肃的一条理由,在来自有限型域~$K$上Motif的$\ell$-adic表示的层面上,是去证明:一个权~i~的Galois模被另一个权~j~的Galois模作的扩张,当~i~<~j~时必为平凡。这个论断——它将证明与一个Motif相关联的$\ell$-adicGalois模(已是相当贴近Motif本身的对象……)存在一份典范的“按权递增”过滤——究竟当年是Deligne还是我证出来的,我已记不得了。}。这一发现——它一如“猜想性的Motif理论”,同样只是“猜想”——随即供出了那把钥匙:由此可在复数域上,给\textbf{Hodge-Deligne结构}(又称“混合Hodge结构”)作出成形的定义,把它当作Motif之上、及其Hodge实现之上那些“早已知道”的结构的“Hodge式”转写。

\dmark{}\origpage{793}按技术讲,我的想法在Hodge-Deligne结构的定义里的影响是双重的。一方面,经由一个Motif的权的概念——Deligne把它恰当地精确化为一副“\textbf{按权过滤}”的结构。另一方面,自五十年代起,我就一直强调:一张光滑代数簇~$X$的代数\textbf{De Rham上同调}——$X$不必真——之重要:它是一个比朴素的Hodge上同调(诸~$H^q(X,\Omega^p)$的直和)更丰富的不变量;后者借那条著名的谱序列与前者相连,谱序列关联于De Rham上同调的一份典范过滤(\textbf{De Rham过滤})。我是第一个定义代数De Rham上同调的人(在当年,谁也不曾想过去看一个微分算子复形——如De Rham复形——的整体超上同调),也是第一个坚持它的分次\textbf{过滤}结构的人——以与Hodge上同调的双分次结构相抗;自Hodge以来,双分次结构一直站在前台。在$X$真的情形(即拥有Hodge理论的情形,该理论蕴涵上述谱序列在特征零下退化),且在基域$\mathbb{C}$上,人们取这份过滤与复共轭过滤二者的“交”,便从De Rham上同调的过滤结构复原出双分次结构(借De Rham上同调的“实结构”——它同构于Betti上同调~$H^*(X,\mathbb{C})$)。随后我证明了(其时除我之外,还无人相信非真情形的De Rham上同调):对复数域上一张光滑概形~$X$,De Rham上同调(它有一个“纯粹代数”的意义)典范同构于复Betti上同调(由超越途径定义)。

话虽如此:一旦设定——$\mathbb{C}$上存在一个Motif概念(不必半单),以及一个$\mathbb{C}$-概形~$X$(当然不必真)的Motif上同调,以及一个$\mathbb{C}$上Motif的“Hodge实现”概念(合用的、有待找到的)——它(依我的想法)该当给光滑~$X$的Motif上同调配上一副“广义Hodge结构”(有待定义),以De Rham上同调~$H_{DR}(X)$为其底集合;那么,人们在这份De Rham上同调上读出的最初两副结构——即De Rham过滤(我早在五十年代引入)与按权过滤(Deligne引入,出自我关于虚拟权的想法,并细化了Serre的想法——那些想法本身出自Weil猜想)——恰恰就落在Deligne所引入的“混合Hodge结构”这个概念上。

\dmark{}\origpage{794}当然,这层想法的亲缘(164$_{1}$),Deligne知道得一清二楚。按行业的伦理(那份伦理,我不曾成功地传给他),他本该在他引入混合Hodge结构的那篇工作里把它明白指出\fnote{338}{(*)}{即文章“Théorie de Hodge II”(Pub.~Math.~IHES~40~(1971), pp.~5--58)。相反,在1970年尼斯大会上“Hodge~I”的摘要公告里,Serre与我倒被并在同一行里提及——一如我在“牺牲者”那条注(n\textsuperscript{o}~78$'$,第308页)里指出的。关于这方面的评注,见该注的子注~n\textsuperscript{o}~$78'^{1}$、$78'^{2}$。}。他却宁可在这篇工作里对它讳莫如深——那篇工作,也正是\textbf{他的博士论文};一如他判定,在这一特定场合,连那个曾是他师父的人的名字,也一并讳莫如深为宜。

5. 在那份关于Motif的评注书目(附于他去年8月25日的信)里,Deligne指明:“人们[!]之所以犹豫要不要在它上面构造——即在关于Motif的那几份‘经典文本’\fnote{339}{(**)}{即Kleiman、Manin、Demazure关于Motif的那几份零星的(‘经典’)文本,出版至1970年为止。它们几乎不出Motif的初始观念,对于我所发展的那份‘yoga’的精微——我曾试着把它传给愿意听的人——给不出丝毫印象。尤其,其中对Motif Galois群只字未提——而它恰恰是发展这份yoga的一条初始的根本动机。(见‘一个梦的回忆——或Motif的诞生’那条注,n\textsuperscript{o}~51.)}——之上构造——原因之一,是其中对代数闭链存在性猜想的使用:那些猜想并无真实的凭据;而Motif,它们,于我是无可怀疑的。”

对这份说明,我将回答:这些“经典文本”丝毫代表不了六十年代末的“技艺之现状”,差得还远;而Deligne学得这份“现状”,靠的也不是这些文本!他心里十分清楚:我的“标准猜想”只是诸多可能进法之中的一种——为在域上给一个(半单的)Motif概念作一份临时的“成形构造”——它丝毫不限制他得自我的那些想法的分量与内在动力。(关于这一点,见“一个梦的回忆——或Motif的诞生”那条注(n\textsuperscript{o}~51)的子注~n\textsuperscript{o}~51$^1$。)一手图二利:我离开之后,他一面竭力把标准猜想败坏成“无从下手”而乏味的东西,一面又竭力败坏某一种Motif的进法——仿佛那进法是我的,并且仿佛它代表一条死胡同,(照他的说法)与那些无望的猜想拆不开地绑在一起——以至于,在LN~900那部总算做了真正该做的工作的卷册里,把我的名字\dmark{}\origpage{795}知趣地讳莫如深,反倒算是对我仁慈……\fnote{340}{(*)}{还在他住在我家的头一天,Deligne就抢在任何我可能就此提出的问题之前,带着他最漂亮的微笑对我说:“你当真以为,并不是人人都已经知道,Motif是你引入的吗!”。而令人惊讶的恰恰是:不管我的朋友为让人忘掉它做了多少,我得以确认:这至今仍是人所共知的。可由于我的想法没有成文的文献,Deligne便大有余地去制造一种印象:仿佛我的贡献照例不过提出过一个模糊的一般观念(何况按它本来的样子根本无法使用,因为它依赖那些“今天与当年同样无从下手”的猜想……)——模糊到这般地步,以致它实在配不上一位做真功夫的严肃数学家,劳神为它哪怕作一次纯粹装点门面的引用……}。

6. 在同一份“评注书目”里,我读到:

\begin{quote}
从这个“经典”的观点看\fnote{341}{(**)}{见倒数第二条脚注。},文献里有一处令人惋惜的空缺:你对$\mathbb{F}_p$上Motif的塔纳卡$\otimes$-范畴的猜想式刻画——唯一等价,差非唯一的同构——连同那几个纤维函子(结晶的与$\ell$-adic的),参见Tate,~Classes d'isogénie des variétés abéliennes sur un corps fini,~Sém.~Bourbaki~352~(1968)。
\end{quote}

这是鳄鱼的眼泪:那处“令人惋惜的空缺”,除了我……之外,该由谁负责?正是我的朋友Pierre~Deligne本人——因为除我之外,世上知道那份“猜想式刻画”的数学家,恐怕唯有他一个……把它收进同一部LN~900,凑足分量,原只在他一念之间!况且,就我现在记得的,这份刻画毫无“猜想”可言——除了一点:得假定人们拥有一个所谓“$\mathbb{F}_p$上的Motif”的范畴,满足几条合理的条件——那是任何当得起这个名字的范畴都理应满足的。如果我没记错,所引Tate-本田的那条指涉蕴涵着:该范畴由Tate的Motif(及其逆)、与$\mathbb{F}_p$上定义的阿贝尔簇,乘法地生成。好东西是有一些的(我还略去许多),它们曾被我亲手交到那位辉煌的前学生手里,却被精心埋到了今天……

\textbf{II~平展上同调(“SGA~4$\frac{1}{2}$”、SGA~5、SGA~7、离散Riemann-Roch)。}

1. Deligne就葬礼I向我作的最初几条评论之一\dmark{},关乎一条猜想性定理的坎坷命运\origpage{796}——那条定理,是我在SGA~5里提炼出来的,名为“离散Riemann-Roch定理”。在“屠杀”那条注(n\textsuperscript{o}~87)的子注~n\textsuperscript{o}~87$^1$里,我对它作过相当细密的表述。Deligne向我指明:当他把我的猜想表述转达给MacPherson时,他自认扮演的是“邮差”、中介的角色。他没有给我的表述添任何新成分——把我的表述译成同调语言、从而使它对奇异空间也有意义的主意,出自MacPherson,而非Deligne。他对我说,当他收到MacPherson那篇文章的单行本——它以超越的论证,在解析-复的情形、在同调的语境里证明了我的猜想——他吃了一惊:猜想竟顶着“Deligne--Grothendieck猜想”的名字。他曾想过写信给MacPherson,纠正这个误会,可(连他自己也说不出为什么)最终没有写……

2. 与我所设想、所暗示的相反:在SGA~5口头讨论班进行之际,Deligne并未承担撰写该班一篇或几篇报告的承诺——例如关于结合于一个代数闭链的上同调类的报告(那一篇,他在讨论班十一年之后终于写成,收进由他编成的那部题为“SGA~4$\frac{1}{2}$”的卷册,不经任何手续\fnote{342}{(*)}{这一拆散SGA~5讨论班(还有许多其他)以肥了那部题为“SGA~4$\frac{1}{2}$”的卷册的行动,履行着两种职能,都朝着“颠倒”角色而去:把我弄成Deligne的“协作者”,并为“SGA~4$\frac{1}{2}$”对SGA~5的优先权主张(那主张,已由那个骗人的名字SGA~4$\frac{1}{2}$所暗示,又在引言里“字里行间”显豁出来——无论Deligne为SGA~4$\frac{1}{2}$写的引言,还是Illusie为SGA~5写的引言)提供支撑(在SGA~5里,经由那份被盗版的报告,指涉SGA~4$\frac{1}{2}$之处比比皆是)。关于这一点,亦见“颠倒”那条注(n\textsuperscript{o}~68$'$)里的评注——在那里,我终于识破那部海盗卷册所得的怪名字的意味,以及我的那篇代数闭链的报告何以在那卷册里现身。})。就此我提了一个问题:他是否不认为——能“现场”地在SGA~5里学到那些基本技术,是他此后全部著作的倚仗——这份特权,不给他强加一份\textbf{义务}或责任:尽力使这些技术经由SGA~5的迅速刊行,交到数学公众的手里。Deligne回答我:他\textbf{不这么认为}。关于Motif的哲学——它曾是他代数簇上同调(其著作的中心主题……)的主要灵感之源——同样的问题,我克制着没有问他。

\dmark{}\origpage{797}3. 正是Deligne主动向Verdier征询同意,要把Verdier关于导出范畴的工作里那份出名的“第0状态”收进“SGA~4$\frac{1}{2}$”。Verdier起初谢绝,认为那毫无意义(确切的措辞我记不得了)。最终是Illusie说服Verdier点了头。Verdier的第一反应,在我看来再自然不过,合乎单纯的数学常识。况且,导出范畴,Verdier多年前已决定埋葬——以一份大规模的逐件加工的形式,它有朝一日该构成他的博士论文——于是,发表一份大纲草图(它久已被文献大大覆盖)便显得滑稽突兀。我自信懂得Deligne与Illusie何以对发表这份“第0状态”如此执着——其中不提我的名字。至于Verdier何以收回他那最初合乎常识的反应,我自信有所感觉,并在“赊账的学位论文与全险”那条注(n\textsuperscript{o}~81)里谈到。

\par\noindent 4. 在“白板”那条注(n\textsuperscript{o}~67)里,我早已指出“这个讨论班”一语的歧义——它出现在SGA~4$\frac{1}{2}$的引言(第2页)里:“为着对L函数的应用,这个讨论班另含一个证明,它补足Frobenius态射的特殊情形。”这个歧义的说法,照上下文与文气,十有八九会被读作意指“SGA~4$\frac{1}{2}$”,从而暗示:母讨论班SGA~5并不含有L函数有理性的一份“完全”的证明。Deligne向我指明:在他心里,“这个讨论班”指的正是“SGA~5”。说实话,这份说明在我这里什么也没说明。我很清楚:SGA~5里有一个“完全”的证明——Deligne与我同样清楚——不过那是一条迹公式的证明,而且(与他所暗示的相反)它远远溢出“Frobenius态射的特殊情形”。但在Deligne的笔下,含混与歧义比比皆是,乃至昭然的假话,而它们全都指向同一个方向——这决非偶然:制造一种印象,关于我的著作、或Mebkhout及其他与我这个人相连者的著作,足以败坏它们的印象,同时抬举他自己的声望,或凭空捏造出这样的声望\fnote{343}{(*)}{尤其是暗示:Motif的枢纽想法、平展上同调、“好上帝定理”及与之相随的Mebkhout哲学,都是他的发明。(3月26日)就本个案与“这个讨论班”,亦见子注“双重含义——或行骗的艺术”(n\textsuperscript{o}~169$_7$)。}。

\dmark{}\origpage{798}5. 我借这个机会,就SGA~7~II(那部署为P.~Deligne与N.~Katz指导的讨论班)再添几条评注——关于它,我在那条(当时无名的\fnote{344}{(*)}{(3月26日)其间我已补上这个缺:把这条注收进目录,题为“一场屠杀的序曲”。})注~n\textsuperscript{o}~56里,已作过相当细密的表述。稍细一点的审视让我看到:在这一场里,N.~Katz没有客气,悄悄给那辆由Deligne利落驾着的灵车推轮子,而且方式很多。

Katz接受与Deligne一同署名为卷册与讨论班的合著者——这与四年前口头讨论班上实际发生的事,全然对不上。SGA~7讨论班的总体设计(它延续1967--69两个年头)出自于我;讨论班的招牌,是与Deligne共同指导。Katz在其中列名为一名合作者-报告人,另有若干位。但既然Katz接受了署名合著者(其中五篇报告出自他的手笔,但主要结果没有一条归于他),那么把卷册的整体操持、以及卷册里对我这个人的藏匿,也视他与Deligne同等的共同负责,便顺理成章。

我头一个想到的,是卷册引言(Deligne署名)里的藏匿:文中呈现的任何主题或结果,竟无一处让人想到与我有什么干系——而讨论班摆在突出位置的两条“关键结果”之一(即Lefschetz束的理论),早在我筹办SGA~7之前就由我发展,而且正是我设想办一个单值群主题讨论班的动机之一。在Katz呈现这一理论的报告(第XVIII讲)里——题名“Lefschetz束的上同调研究,作者N.~Katz”——我的名字不按惯例出现在题名里(“据A.~Grothendieck”),而是出现在N.~Katz名字之后一条简短的脚注里:“据Grothendieck的(简略)笔记”。那“简略”的考语,像是后加的,好冲淡这桩不巧的“Grothendieck笔记”在此扮演过角色的事实。可那些笔记纵然“简略”,它们代表的却是好几天工夫的结晶——面对一桩先验绝不显然的任务:把一些经由超越途径陈述并证明的结果,转写进一个技术上全然不同的语境。一如平展对偶、或一如Nielsen--Wecken的理论\dmark{}\origpage{799}\fnote{345}{(*)}{Deligne比他的朋友N.~Katz更少拘束,倒是压根不觉得值得提一句:在他所谓的“Nielsen--Wecken方法”里,我也有几分干系——关于这一点,见“白板”那条注(n\textsuperscript{o}~67)的子注~n\textsuperscript{o}~67$^1$。},经典的论证原样用不上,一切须从头重做:拿经典\textbf{结果}当引路的线头,而把它们传统的“证明”(如果那配得上这个称呼的话)整个忘掉。正常的倒是:即便有我那些细密的笔记相助,Katz也得费一番力气才能入水,一如在他之前我也曾费过——可这丝毫不意味着(至少,不按通行认可的游戏规则)他就是平展上同调里Lefschetz束理论的作者!

顺着这股劲头,在同一报告的引言(第225页)里,Katz摆出把Raynaud太太呈现为那定理作者的架势——特征p的代数曲线上、“与p互素”的驯顺基本群的构造定理。如果我记忆不差,正是这条定理(我1958年证明的,其时我尚未结识我未来的学生)——连同“蛮横的Lefschetz定理”——构成这份理论的深刻技术成分;而在不可约性定理的证明里,能以全副力量使用它,我曾深感快意。

在Katz第XXI讲的引言(第364--365页)里,描述完该报告的主要定理——关乎射影空间里的完全相交——之后,有这样一段:

\begin{quote}
存在一些归功于A.~Grothendieck的启发式论证,它们倚仗结晶上同调的yoga,以本质相同的方法,使一般的表述对一切光滑射影的X显得可信。
\end{quote}

这条评注暗示:仿佛我是从正文的方法(它出自一位未具名的作者——那多半只能是卷册的两位作者之一)获得启发,在它上面绣出那些“启发式论证”,以推广所证的结果。我记得的恰恰相反:是我的那些“启发式论证”(我远在讨论班之前,循着对Griffiths定理与Lefschetz束的反思\fnote{346}{(**)}{其实正是这些反思——一如我关于抽象代数几何里消没闭链理论的反思(我的另一场“超越理论的纯代数转写”!)——是SGA~7讨论班的起源。},独自发展出来)恰恰“走得通”\dmark{}\origpage{800}——何况不带任何猜想性成分——在X为完全相交的情形。况且,在前一报告(同样出自Katz)里——它专讲Griffiths定理——引言里说:“这里给出的证明(出自Grothendieck)是Griffiths原有证明——或多或少超越性的——在纯代数项下的转写。”这条评注给人的印象是:特征任意的Griffiths定理,证明不止一种,任凭挑选,而我有幸被选中了我的那种。其实,据我所知,此外并无别种证明。况且,凭我当时被迫下的工夫,我怀疑这证明是Griffiths证明的单纯“转写”——一如平展上同调里任何一条大关键定理的证明,都不是某个已有证明的“转写”;或者(干脆说到底)掌握概形的平展上同调,从来不是把寻常熟悉的普通上同调理论“译成纯代数项”的问题。

我清点了N.~Katz报告文本里指向我这个人的全部三条指涉(而在Deligne的八篇报告里,总共只有一条!)。在我看来,三条映出的是同一份蓄意。收尾之前我指出:卷册最后一篇报告——N.~Katz作,专讲特征p上一条L函数的“模p同余公式”——其文本里没有我的名字\fnote{347}{(*)}{这并不全对——它倒是出现了(于是这是指向我这个人第四条指涉):在第410页,与Deligne并提,感谢我们俩向作者解释过他所呈现的同余公式形状的几种等价改写。滑稽的是:他为这些辉煌变体开列的三条编号文献,在报告里一条也不存在——以致这份致谢,俨然一桩善意的恶作剧!(这已不是我在葬礼里遇到的第一桩……)}——连L函数那个寻常的上同调表达式也没提我。其实,正是结晶上同调项下的类似表达式(它还停留于猜想),多年前引我猜出了这条同余公式。我把这个猜想告诉过Deligne;他找到了一个惊人简单的证明,凭他那条对称Künneth公式(陈述于SGA~4~XVII~5.4.21)。我猜想:对这类事情全然在行的Katz,同样清楚这条猜想的来历,只是不觉得值得提。(他在文本里给出的,是与Deligne不同的证明,而且拙劣得多。)

滑稽的是,在SGA~7~II这最后一篇报告的引言末尾,人们读到:Deligne的证明“当收进SGA~5的重编本”(而SGA~5,它还无缘见识自己的头一“版”呢\dmark{}\origpage{801}!)。这可以让人猜想:早于“SGA~4$\frac{1}{2}$—SGA~5”那场行动五年,Deligne或许仍有意(那本属正常)在SGA~5将来的刊行版本里,收入他自1966年以来对平展上同调理论——发展于SGA~4、SGA~5——所作的种种补充\fnote{348}{(*)}{我猜想:正是SGA~7里那些藏匿未遭任何反应(来自任何一位在行的人),才鼓励了Deligne在节节攀爬中迈出下一步:“SGA~4$\frac{1}{2}$—SGA~5”那场行动的大骗局。}。

\textbf{III~Mebkhout的哲学(1981年6月吕米尼会议;Beilinson、Bernstein、Deligne论“反常层”的文章)。}

这里我把前一条注里就此所记的,再录一遍以备记忆。

1. Deligne对我说,他是在一次Bourbaki讨论班上与Mebkhout交谈时得知“好上帝定理”的\fnote{349}{(**)}{见“当值无名之辈与好上帝定理”那条注,n\textsuperscript{o}~48$'$。}——反正在1980年夏之前。这与我从Mebkhout那里听来的相印证:那条定理,是Deligne于1980年10月转告Bernstein与Beilinson的,随即被他们用进他们对Kazhdan-Lusztig猜想的证明里\fnote{350}{(***)}{见5月28日的脚注——附于“不义——或归来的意义”那条注(n\textsuperscript{o}~75);亦见“一种不义与无力的感觉”那条注(n\textsuperscript{o}~44$''$)。}。Deligne补充说:他在与Bernstein、Beilinson合写的文章里没有引Mebkhout,因为拿不准这条定理里Kashiwara该占几分\fnote{351}{(****)}{关于这方面的评注,见前一条注“义务已尽——或真话时刻”,尤其p.~784,及论“Kashiwara”的那条脚注。}。

2. Deligne并不否认:没有Mebkhout前些年的工作,1981年6月的吕米尼会议(他自己是会上最大的明星)不会举行。他只坚持补一句:MacPherson想法的角色,在他看“还要更根本”。至于Mebkhout的名字不见于会议文集,是否有几分怪异、几分反常——他从不曾暗示。

\textbf{IV~上同调的对偶形式主义、导出范畴(“好的文献”、“导出范畴的第0状态”)。}

\dmark{}\origpage{802}1. Deligne向我指明:Verdier那篇文章\fnote{352}{(*)}{即“好的文献们”那条注——这名字显然非它莫属!(n\textsuperscript{o}~82)所引的那篇文章。
(5月12日)关于Deligne这个叫人难以相信的版本,见“荣耀满天飞——或暧昧”那条注(n\textsuperscript{o}~170(ii)),第930、931页。}——它把(别的且不论、并且不提我名)结合于一个闭链的同调类与上同调类的形式主义(我1965/66年在SGA~5里发展的)认作己出——他是在1977年SGA~4$\frac{1}{2}$刊行之后才知道的,至少比那篇文章发表晚一年。这么说来,我曾经的印象——Verdier1976年那手漂亮的操作,是Deligne及其同伙次年那场大得多的行动的一只“试探气球”——倒像被它推翻了。

Deligne对我说:浏览Verdier那篇文章时,对他一目了然:那不过是把我SGA~5里的某些想法铺陈一遍。Verdier终于肯负责提供一份文献,他甚至为此满心欢喜。(“SGA~5的刊行本也许才是更相称的文献”——这个念头,想必从不曾掠过他的心头……)循着这层意思发问,Deligne回答我:他不曾留意Verdier的文章里没有我的名字——又补一句:他承认,他连想都不曾想到要问这个问题。我得到的印象是:他在无声地示意:这类事情在他操心的事里排最末,不值得停留……

2. 在Beilinson、Bernstein、Deligne那篇文章(葬礼I里多次引过;它由Deligne执笔整理,由他在吕米尼会议上宣读\fnote{353}{(**)}{关于这场“值得纪念的会议”与这篇文章,见“不义——或归来的意义”那条注,n\textsuperscript{o}~75.})里,平展上同调的对偶(我1963年发展的)被称作“Verdier对偶”\fnote{354}{(***)}{这场操作分几步完成。经我提议,Verdier在1963年之后,于寻常拓扑空间的语境里,发展了一套“六运算”的对偶理论,循着我在代数凝聚与平展语境里主持发展的那份总设计。这份对偶,不出所料,被我的上同调家弟子们洗礼作“Verdier对偶”或“Poincaré--Verdier对偶”,只字不提我这微不足道的人。另一方面,在1976年的那份“好文献”里,Verdier又在解析的语境里,不加指名地重述了我在五十年代于凝聚框架里发展的那份形式主义的一部分(对它不须改动分毫)。于是这份对偶,在解析框架里,又取名为“Verdier对偶”,或有时“Serre-Verdier对偶”——依旧不提我这个人——连Mebkhout也跟着大势走!可(妙的是一记回马枪)显然之极:代数的凝聚对偶,不过是那份超越的解析理论的“纯代数转写”,一如平展对偶是那份超越的拓扑理论的如此“转写”。于是势所必至:连它们也一并洗礼作“Verdier对偶”(Serre与Poincaré,情有可原,因为离得远)。据Deligne对我说,这正是“人人”争先恐后做了的。幕落……}。

\dmark{}\origpage{803}这个奇怪的称呼,我问过Deligne。他回答我(这一回,带着一丝窘迫):因为“人人”都这么叫。这个“人人”都是谁、以及这算什么理由——我没有细问;毕竟,这门理论该归功于谁,他Deligne一清二楚。

这让我想起一件久已触动我的事。至少在与我交谈、或给我写信时,Deligne从不使用“导出范畴”这个说法而不加上“Verdier的”。这每次都给我一份不快,可我(在发现葬礼之前)从不曾停下来探究其意,更谈不上把点重新点清。倘若当初肯费心对“SGA~4$\frac{1}{2}$”投去稍带好奇的一瞥,对Verdier那部在其中被起出的“论文”的“第0状态”——我无疑会停下来。(后一事的细节,见上II~3.)

\textbf{V~葬礼悼词}

1. IHES的那本周年小册子——我的那篇葬礼悼词\fnote{355}{(*)}{见“葬礼悼词(1)——或恭维”与“葬礼悼词(2)——或力量与光环”两条注,n\textsuperscript{o}s~104、105.}就载于其中——并非出自其创建人暨首任院长Léon~Motchane之手(我原以为是)。册子作者是谁——Deligne把身份告知了我——在此倒无关紧要。他向我确认:涉及我的那段,确确实实出自他的手笔;而那一段,连同涉及他Deligne本人的那一段(出自册子作者之手),在付印之前,都确确实实得到过他的“绿灯”。他写我的文字,起初更长,被册子作者删节了(经他同意)。涉及他自己的文字,Deligne也审阅并修改过。所以,这些文字确确实实代表着Deligne的立场——关于他的著作,与我的著作。

2. 我问过Deligne:我是否猜错了——我曾认定,在他任何一篇发表的东西里,他从不曾暗示:他或曾从我口中\dmark{}\origpage{804}学到过什么。他向我确认无误——只有一处保留。那涉及他写给(布鲁塞尔)国家科学研究基金会的一篇小传,为着“五年奖”的颁发。那项奖(我记得是1974年)授给他,奖励他对Weil猜想的证明。诚然(他补充说),这篇小传不算数学出版物,流布也极其有限。至于我,我根本不知道它存在。经我请求,随后几天里他给我送来一份复印件;这篇小传,我打算在下一注里再谈。

Deligne向我确认的这份对我的系统性否认,似乎不构成他的任何问题。他不觉得其中有什么怪异、值得停留的东西。鉴于他这份心态,我没有起意在这个方向问他任何问题——我不认为多问能榨出更多。为这份回顾收尾,我只添一句:凡属严格意义的“物质性事实”,我对Deligne的诚意并无怀疑——那诚意在我看是显见的\fnote{356}{(**)}{(5月12日)凭回顾,对这份印象毕竟浮现了某些保留,如前文一条脚注((*)第802页)所提到的那些。还显出来:有两条粗疏的事实性错误在我的注里,Deligne不曾向我指出——它们本很难逃过他的眼睛。(我曾漏看:他早在1970年Hodge~I里就披露了“权的yoga”的一部分,而且他1979年起就谈到Motif.)}。唯一的例外,是他的那个断言:SGA~5(1965/66年的)讨论班在逻辑上依赖SGA~4$\frac{1}{2}$\fnote{357}{(*)}{诚然,这一断言之来,并非Deligne自发地主动前来,给我送来“物质性说明”,以启发我、以显豁他那全然的诚意;而是在我刚刚当面道出我对“SGA~4$\frac{1}{2}$—SGA~5”那场令人难以相信的行动的感受之际,在“保全面子”这份未及预料的需要的压力之下。关于这一点,见前一条注“义务已尽——或真话时刻”的末一部分(2月18日)。}——那些结果自1973年起,在Deligne讲述其Weil猜想证明的各报告的边缘上发展而成。诚然,靠“俘获”母讨论班SGA~5的某些报告(尤其那篇关于结合于一个闭链的上同调类的),凭Illusie(他负责SGA~5的编辑)与许多其他人的默许,他赢得了这份辉煌的战果:SGA~5通篇塞满对SGA~4$\frac{1}{2}$的指涉,足以给(一位不甚留神、或不甚在行的)读者一种印象:SGA~5确确实实依赖SGA~4$\frac{1}{2}$,而后者在一切方面都摆出一份“更早”文本的架势。这手腾挪,在我们这门科学的编年史里想必独一无二;依我看,它足以把本世纪的七十年代,从数学所经历的一切其他时代里标出来。

\subsubsection*{注~164$_{1}$}
\notehead{164$_{1}$}{}

\dmark{}\origpage{805}关于那份出自Weil猜想的“权的哲学”,这份“亲缘”在我看来可总结如下。

a) 一如“我的孤儿”那条注的子注~n\textsuperscript{o}~46$^9$所说:作为Weil猜想背后那份“哲学”的一部分,Serre曾传给我某种“虚拟权的yoga”——在域上有限型概形的$\ell$-adic上同调的层面上。他不曾试图给出它一个精确的显式表述;对不同~$\ell$~所发生之事之间的关系,则全然神秘。

b) 自六十年代初起,引导我发展一份“Motif的yoga”的两条主要动机之一,恰恰是:把对不同~$\ell$~的那些“虚拟权结构”彼此连起来。(关于这一点,见“一个梦的回忆——或Motif的诞生”那条注(n\textsuperscript{o}~46),尤其第208页。)由此便清楚:这份结构该当在一个Motif一切可能的“实现”上重现——不只在$\ell$-adic实现上——尤其(在基域$\mathbb{C}$上)在De Rham-Hodge实现上。

c) 这份虚拟权的哲学——其终极的源头是Motif——经我告知Deligne之后,他为这份yoga添上一份重要的精确化:揣想一个Motif上的虚拟权结构,联系于一份(必然典范的)\textbf{按权递增}的\textbf{过滤}。于是这份过滤该当在Motif的一切实现上重现——无论$\ell$-adic实现,还是(在域$\mathbb{C}$上)De Rham-Hodge实现。Deligne的这份“揣想”,成了他“混合”Hodge结构理论(我称之为“Hodge-Deligne结构”)的起点,也是他为它们\textbf{成形}的定义的两条本质技术成分之一(另一条是De Rham过滤——我早在五十年代引入)。他为描述$\mathbb{C}$上任意有限型分离概形的“Hodge上同调”所作的尝试,其成功,可以被视作我们眼下就Motif上按权过滤存在性这份“揣想”之有效所拥有的主要的(甚或唯一的)“凭据”。

当然,阐明“$\mathbb{C}$上有限型概形上的Hodge系数”概念,本是我围绕Motif那份大工作纲领的一部分——Deligne对它是第一手的、日复一日地知情的:其方式,须使X上的一个Motif对应一份“Hodge实现”,并且对X上光滑而纯粹的Motif(例如那些取一张X上真而光滑的概形的“第i维X上Motif上同调”而得的Motif\dmark{}\origpage{806}),人们重新找回“Hodge结构族”的概念(六十年代尤其由Griffiths研究过)。而且,X变动时,这些“Hodge系数”的范畴须满足一份六运算的形式主义——在Motif的层面上映照同一份形式主义——Deligne的贡献,代表朝这份纲领之完成迈出的头一步:即(本质上)对X缩为一点的情形给出Hdg(X)范畴的描述\fnote{358}{(*)}{为求完善,须给Deligne的定义补上引入一个合适的\textbf{三角}范畴~Hdg$^*$(它莫非也是Hdg的导出范畴?)。他漏掉了这一步,在我看来是(此后还有别的)最早的信号之一:对导出范畴的yoga与六种运算的冷淡——它一直肆虐到1981年“反常层会议的转折”。},以及“实现”函子的描述——即本质上,在$\mathbb{C}$的有限型分离概形上构造一套上同调理论,取值于Hodge-Deligne结构这个范畴。

\section{4.~死之舞}

\subsection{(1) 为一副模糊的骸骨作安魂曲}

\notehead{165}{}

(2月22日)自他去年10月的来访算起——甚至更早,自他8月底的那几封信算起\fnote{359}{(**)}{见“义务已尽——或真话时刻”那条注(n\textsuperscript{o}~163)——我在那里给这次来访“定位”;亦见8月底的那两封信(它们是在近两个月的沉默之后到的——那沉默,继我寄出引言与葬礼目录之后)。}——我的朋友皮埃尔,成了昔日学生与好小子中的头一份:满载着一份显然动人的好意,要驱散那桩不巧溜进我们之间的误会,要让我感到他的善意与诚意。双方讲定:在《收获与播种》由我的大学(USTL)作预定预印之前,他对他读我注子的内容、乃至其存在,守口如瓶。他是否把话守全了,我不知道;反正好几路回音传到我这里\fnote{360}{(***)}{比如:我收到Illusie一份未署日期的预印本(我猜是最新的),出自某个未具名讨论班的一篇报告(文中特意说明:它不对应任何一次口头报告)。标题里——难以置信,却是真的——我的名字赫然在列:“据A.~Grothendieck的Barsotti-Tate群形变”,作者Luc~Illusie!引言里还有大段大段的“Grothendieck”——我疑心自己在做梦。分明是有什么事发生了……随信附有一封短信,他就若干“Grothendieck风格”的同伦代数之点向我讨教,并自问:“搞K理论的人们”(即Quillen等)为什么用层工作,而不用我二十多年前引入的装备库里的那些(伪凝聚或完美)复形。人们也的确要问个为什么……回信里,我不得不示意:这样的问题,轮不到他、也轮不到我任何一位昔日学生来问我。此后我再没有他的音讯。},我颇有印象:他准是两头都递过话——暗示:或许该是时候,给那位老师傅递几分体贴了(那位老师傅:人们偶尔在小圈子里谈起,却小心翼翼地避免在公开场合提名……)。

我另外有一种印象:骨子里,我的朋友并不相信(或至少不愿相信),我会\dmark{}\origpage{807}确确实实把葬礼付印,与《收获与播种》的第一部同时。第一部,倒是一副十足的“心软爸爸”的形象:唯恐点谁的名会让谁难受,对自己想到的种种自家的缺失,又随时准备当众认账。读那部“愚妄与更新”——我朋友度假动身之前、我给他寄出葬礼引言之前,我得过它的一点回音——不曾让他有分毫不安;恰恰相反——它多半还助长了一份自我满意的做派,那份做派,在他身上我早已熟稔——那丝居高临下、或至少是对那位显然已故的师父的护短的做派。葬礼,可就全然不是同一回事了:牌,骤然间干脆摊上了桌面!我猜想,读那篇引言,该给了他一记冲击——可惜当时我不在场,否则或许会有些什么发生。反正,他给自己留了回神的时间,才来见我——一阵风似的,在他迁居美国前五分钟。而他来得满载好意,相遇进行得何其合家欢、何其“心软”,以致——不妨说“以归谬法”——这反倒排除了:那位心软爸爸当真会把某一份不像他的文本当回事(那份文本,不说也罢,最好忘掉……),甚或把它散发给一群与他朋友皮埃尔同样体面、样样“正派”的人们,以及那位如他所一直认识的前亡故者……\fnote{361}{(*)}{不过,把我在葬礼上的全部各注公之于众——与《收获与播种》第一部一样——我的这份意图,从不曾有过分毫的迟疑;对此,我当然不曾留下任何暧昧。}

一如他所允诺的,回到比雷斯之后没几天,我的朋友就把我上面说的那篇小传\dmark{}\origpage{808}寄给了我——他1974年(或1975年)为(比利时)国家科学研究基金会而写的那篇\fnote{362}{(*)}{这篇小传头一次被提到,是在“神经中的神经——或侏儒与巨人”那条注(n\textsuperscript{o}~148)的最后一条脚注里。亦见前一条注n\textsuperscript{o}~164的结尾(V~2部分)。}。文字相当短,小小的两页;我那时饶有兴味地读过,刚才又重读了一遍(这已是第三遍,我想)。可乍一看,我并没有这份印象:这文本带来了什么新东西,值得我在葬礼里为它停留。诚然,藏匿的技巧——在我朋友身上我本已识得够多——在这里,在一篇百余行的紧凑文字里,得到了格外触目的演示。我的名字在其中出现四次(一如Serre,而Weil出现三次)——却没有任何东西让人想到:他与我打过交道的路数,或许不止是我那个讨论班(主题未说明)的一名无名听者(1965--66)。在提及我的四段文字里,有三段把我与另一位数学家并提(两次Serre,一次兰金),为的是避免给人印象:我曾在他身边扮演过稍稍特别的角色。这又是一门在别处久经考验的技巧\fnote{363}{(**)}{我这里想到那条一行字的简略指涉——在(尼斯大会上)Deligne宣布其Hodge理论结果的公告里——它把Serre(不提名)与“Grothendieck的猜想性Motif理论”并在一口气里。详情与评注,见“牺牲者”那条注(n\textsuperscript{o}~78$'$)的子注~n\textsuperscript{o}~$78'^{1}$。}。既然引起来不长,我就不揣冒昧,把我的微不足道的人现身的三个段落全文引出,以照亮那些不象我这样手头有这篇小传的读者。

第三个段落,接着(方才已作过的)对1965--66这一年的回顾——“在高等师范学校理想的氛围里,作为一名外籍寄宿生度过”——说下去\fnote{364}{(***)}{出于一个我无从把握的缘故,Henri~Cartan在这里没有具名。或许因为:凭我对他怀有的一份蓄意(见“格格不入者”那条注,n\textsuperscript{o}~67$'$)给他撑腰,Deligne执意避免任何“他或曾是谁的学生”的表象。“高师生”的身份,会即刻唤起“Cartan的学生”的联想;点名Cartan,只会坐实这份联想。}:

\begin{quote}
“在巴黎,我听了Grothendieck的讨论班与J.P.~Serre的课。每周三个小时的课\dmark{}\origpage{809};然而,尽管工作愉快而拼命,一周的其余时间也才勉强够我消化它们(165$_{1}$)。从Grothendieck那里,我学到了代数几何的现代技术;从Serre那里,学到了数论那令人着迷的美(165$_{2}$)。Serre的课专讲椭圆曲线理论,其中交织着……”,\end{quote}

随后继续铺陈Serre课程的魅力与花样。不在行的读者会想:作者所说那份“愉快而拼命的工作”,对象是这些课——每周三个钟头(言下之意:消化一个Grothendieck讨论班的“最大的自然一般性”,根本用不着下工夫……165$_{1}$)。

第五段,谈他对Weil猜想的证明,人们读到:

\begin{quote}
“我最可观的成功,是证明了‘Weil猜想’(……)。我无疑是由于既熟悉Grothendieck的著作,又——在一个全然不同的领域里——熟悉兰金关于模形式的著作,才做到这一点的。”
\end{quote}

人们将叹赏那个踌躇的“无疑”(放在那里,真是大师手笔!),叹赏那句“在一个全然不同的领域里”(暗示:我的著作与模形式毫无干系\fnote{365}{(*)}{诚然,“模形式”代表着我数学教养里一处令人惋惜的空洞(空洞还多着呢),一如解析数论——我对它至今不曾“挂上钩”。可我的见识毕竟足以知道:没有来自代数几何的想法——它给这门理论提供其“几何”的内容——对模形式的理解几乎无从谈起,而模形式理论最深刻的问题,又与\textbf{Motif}的(长期隐而不显的)在场密不可分。如下文将见,Motif同样不声不响地,出现在这篇小传的下一段里(别名:葬礼悼词(3)!)。}),尤其叹赏那个“既熟悉……又”——我就有这份荣幸,由它引荐登场——好把我那份宏大的奠基工作\fnote{366}{(**)}{关于概形的概念、与平展上同调形式主义的发展——Deligne小心避免提它们,除非在前一条引文里,借那个客气的、无我之称的婉辞“代数几何的现代技术”。},与一条从兰金那里借来的“一时一地”的技术想法,放在同一个天平上。

最后,在接下来那段谈DeligneHodge理论工作的文字里,有这样的句子:

\begin{quote}
“受到算术的启发\dmark{}\origpage{810}——尤其受到Grothendieck对Weil猜想之深刻意义的理解的启发——我把他的理论(以非平凡的方式)推广到任意簇的情形,并且(与Sullivan合作)推广到“形”的其他不变量,而不止上同调。这份理论的根已经颇老,有皮卡关于“两个独立变量的代数函数”的论著(约1890年);但如今人们对它的认识,大概不过剩一副模糊的骸骨。”\end{quote}

须要我费心把这段抄下来,我才恍然:“Grothendieck对Weil猜想之深刻意义的理解”云云,是我那位辉煌的前学生不点Motif之名的大师级“暂停!”手法——而你又抓不到他把它们讳莫如深的把柄!毫无疑问,“他的[亦即,\textbf{我的}]理论”——我此刻才反问起它来(整个这一段,前几遍阅读都逃过了我的注意)——只能是那份出名的Motif理论:四年以来,提它的名字都谈不上(而往后八年,更无人会提!)。那措辞模糊到这般地步——说穿了,除了一小撮在行的人(他们多半不曾象如今的我,有机缘读到这篇预备的葬礼悼词),谁也读不懂——以致连点破一句都嫌多余:他“推广”的那门“理论”,其实整个是猜想的!照上下文看,所云的“推广”,也只能是指Hodge-Deligne理论。这是我的朋友给自己的一点象征性的快意:他在这里断言(断言的地点,连同措辞那闪躲的模糊,保证他永远不会被驳)——Hodge-Deligne理论(它至今仍在童年)“推广”了我让他看的那幅宏大的Motif图景。可在这幅图景里,一门抵达全熟的“Hodge理论”,不过是全图诸多“层”中的一层\fnote{367}{(*)}{(2月27日)关于这方面的详情,尤见“墓上的旋律——或自负”那条注(n\textsuperscript{o}~167)。}。至于“形”的其他不变量,我早在六十年代就“了然于胸”(作为我的“Motif的yoga”的一部分):“任意”的代数簇(Deligne特意强调“任意”)有一个“Motif同伦型”,其高阶$\pi_i$($i\ge 2$)推广Motif的“几何”基本群,并且(对一个给定在数域$K$上的纤维函子)显形为$K$上的仿射pro-代数群。

\dmark{}\origpage{811}至于把皮卡抬出来作“这份理论的根”——在我看来,这整段是装点门面的赝品,插进来为着一石二鸟:既“显得有学问”,又顺势引出紧随其后的收尾一段\fnote{368}{(*)}{这收尾的一段,将是紧随本注的那条注(n\textsuperscript{o}~165)的对象。}。“模糊的骸骨”这个字眼,在我看来,同样是我朋友付给自己的另一份“象征性快意”的表达:他暗地里从那幅宏大的视景获得灵感,却把它始终埋着\fnote{369}{(**)}{Motif的视景,以两种方式保持“被埋”。一方面,对着\textbf{外界},对着数学公众:对Motif概念绝口不提(除了1970年Hodge~I里那半行的“暂停!”——参看注~$78'^{1}$)——直到1982年,这概念才在Deligne的(心照不宣的)父权之下“大张旗鼓地”被起出(见注~n\textsuperscript{o}~51及以下诸注)。另一方面,甚至在他自己的私人使用里,我看到:这份视景也被Deligne剥去了它真正的\textbf{元气}——那份使它\textbf{不止}是一种通行配方的汇编(为在代数簇的上同调里认路)的东西:一份足够宏大、足够深刻的\textbf{力之梦},足以充当灵感、地平线上的准线,予或许数代的算术几何学家。
Deligne(照旧不声不响)以“模糊的骸骨”这个字眼指称这份视景,令人惊心地显出他身上的\textbf{掘墓人心态}——在他与这个梦、与梦所由出的那位匠人的关系里,他始终守着这份心态。这已不是那种尚能感到一口气(如他当年曾感到)的心态,也不是体现一个梦的心态。体现一个梦,不是把它\textbf{拿来用}于自己的目的(同时还否认它……),而唯有\textbf{甘为其仆}。},却在心里把它当作——照总账——不过一副“模糊的骸骨”(照样是那副“暂停!”的做派,脸上却不露)。到头来,这些随手挥洒的藏匿,竟比我预想的更有滋味——我原本只打算顺带点一笔,算是了却良心。如今最触动我的,已不是(一如头几遍匆匆浮泛的阅读)那“暂停!”风格的完美——那风格早领教得倒胃。而是:这篇文字——它写于那篇葬礼悼词九年之前\fnote{370}{(***)}{见“葬礼悼词(1)——或恭维”与“葬礼悼词(2)——或力量与光环”两条注,n\textsuperscript{o}s~104、105.}——以触目的方式预示了后者;而且(依我看)在两个方面。一方面,凭那份照例的模糊:凡我这微不足道的人现身处,必得围着一团雾(与此相对,此处,Serre课程的回顾却伴着技术细节的铺陈)。另一方面,同向而行:对平展上同调或$\ell$-adic上同调——那份我从虚无中发展的、新的而根本的工具,没有它,Weil猜想再过一百年恐怕也证不出来——竟全然\dmark{}\origpage{812}无一词提及。其实,一如那篇葬礼悼词:“上同调”这个词,从未在与我的名字的关系里出口;同样无人暗示:Deligne对Weil猜想的证明,不过是一段长途的\textbf{最后一步}——那段路最长、也最新颖的部分,是另一个人在他之前走完的,那时我那位辉煌的学生尚未在数学的舞台上现身\fnote{371}{(*)}{这位“另一个人”的贡献,被Deligne用无我的字眼藏了起来,如“代数几何的现代技术”[或在他处,“有力的工具”]。}。

\subsubsection*{注~165$_{1}$}
\notehead{165$_{1}$}{}

一如我在下面几行里指出的:那份措辞不可抗拒地暗示,“每周三个小时的课”指的是刚谈到过的“J.P.~Serre的课”——而且两句之后还会再谈到它。其实,Serre每年只开一门课(在法兰西公学院),每周一个钟头。倘若试图消除歧义,把文本读成历年连年的Serre“课程”(与上下文所暗示的相反),又会撞上另一处不合:Serre每年换主题,绝不局限于椭圆曲线——可两句之后偏偏这么说。

我的朋友在这里挪用Serre这个人,想就那在他数学生涯的关键岁月里属于我的角色,混淆视听;有意思的是:据我所知,文献里唯一一处言明Deligne曾是我学生的——出自Serre的笔下,他借此(而不点破地)弥补了那位辉煌前学生亲手犯下的种种昭彰的省略。那是Serre1977年5月就Pierre~Deligne的工作所作的报告,为着颁发1978年菲尔兹奖章的国际委员会。该报告在1978年赫尔辛基大会颁发菲尔兹奖章之后公开。报告开头如此:

\begin{quote}
“Deligne最初的工作,直接受Grothendieck的启发——他是Grothendieck的学生——关乎代数几何的若干技术之点。我只限于提及:……”
\end{quote}

\dmark{}\origpage{813}再往下,Serre也提到我的想法与结果在Weil猜想之证明里的影响,以及(经由Motif)在Deligne模形式工作里的影响;但在Deligne--Mumford关于$(g,\nu)$型代数曲线之模多重簇的工作里,以及在Hodge-Deligne上同调的观念里——它与Motif的yoga及Weil猜想的关联,似乎逃过了他的眼睛——则没有提。(诚然,Deligne为此尽了遮掩之能事。)

菲尔兹奖章颁发之际关于Deligne的演说,照公认的惯例,原该是另一个场合:公开重提这份与我的关联——它此前一直被当事人捂着。出于一个我无从把握的缘故,受命介绍Deligne工作的数学家,不是J.P.~Serre,而是N.~Katz——SGA~7~II与Deligne“合著”的那位(关于这一点,见注~n\textsuperscript{o}~164(II~5))。不用说,N.~Katz对这层关联不着一字——尽管它于他一清二楚、而且是第一手的。(倒是,他顺手不露声色地弥补了那位桂冠得主关于我的若干处有些难堪的省略……)

\subsubsection*{注~165$_{2}$}
\notehead{165$_{2}$}{}

这里考语的拣选(给我的是“现代技术”,给Serre的是“令人着迷的美”)决非出于偶然。我在其中分明察觉我朋友的用意:把那份着迷——自我们相遇起(或许更早)便把他系于我这个人、系于我的著作的那份\textbf{着迷},他眼看它在他眼前日复一日地成形、铺展——(象征性地)打发掉。

在别的场合我也留意到我朋友的一份蓄意:把我的著作(尤其EGA即《代数几何原理》、SGA即“布瓦玛丽”代数几何讨论班)看成并呈现为某种“汇编”——搜罗一些或多或少技术性的结果,“人人”素来皆知,而我可嘉的努力在于把它们白纸黑字写下来,好终于补齐缺失的文献,从此不再谈起。可骨子里,他清楚该把账算在哪里:EGA与SGA的每一卷,呈现的都是我引入的想法——多少年间,我是它们唯一的持有者与辩护人;呈现的是无人梦到过的技术(除我之外),而我不得不以不倦的耐心把它们发展、检验、完善,直到它们磨合纯熟,随时可以进入“人所共知”的领地。这一点,没有人比他知道得更清楚;可与此同时,他彰示了十多年的那份蓄意\dmark{}\origpage{814},终于变成了一份“第二天性”,而他自己,成了它头一号(若非唯一的)受骗者。

就在几周前,我还为之一怔:我的朋友——自他10月造访寒舍以来,对我处处体贴——让人给我送来一份他与Heinze博士(在斯普林格负责“Ergebnisse der Mathematik”丛书)之间通信往来的抄件,谈EGA的重版计划(其中许多卷已售罄或濒临绝版)。在其回信中,Deligne毫无保留地推荐全部原样再版,“ne varietur”——几乎不作改动,理由是除了一处例外(EGA第三部第二部分,在那里用导出范畴叙述会更好(sic!)),这部论著“has aged very well”,它的巨大功绩正在于提供不可或缺的参考文献:“Thanks to it [EGA], in algebraic geometry (as opposed to analytic geometry, for instance) one can march securely on the ground without having to worry if this or that is indeed in the littérature”。(他接着提出若干建议,关于可能给某些卷册添附的附录,以及有能力提供这些附录的数学家……)

斯普林格出版社对我这个人的态度极具代表性:这场通信(关于我身为作者的著作的再版)是\textbf{同Deligne}进行的,而斯普林格甚至不曾先通报我这一计划。过了一个多月(在1月24日的信中),Heinze博士才顺带向我提起这事,算是了却良心——说Deligne教授先生“曾非常体贴地向我提供他1984年12月19日来信的抄件”(真是体贴……),并且“我们[斯普林格]愿听取您对此事[再版计划]的意见”(真是太荣幸了……)。我回复说:鉴于斯普林格出版社在重印方面的惯常做法(想想它在Lecture Notes里出版SGA~7与SGA~5,连知会我一声都不曾,更遑论征得我的同意),我认为向斯普林格出版社通报我那显然无足轻重的“意见”,纯属多余。事情便停在这里……

\subsection{(2) 信仰表白——或伪中之真}

\notehead{166}{}

(2月23日)归根结底,昨天在谈论我朋友皮埃尔的小传时,我还不曾抵达我的真正题旨。同“模糊的骸骨”(别名:Motif理论)的相遇是一段意外的插曲,当时我正准备顺势引出小传的最后一段——它紧接在方才所引的那最后一段之后\dmark{}\origpage{815}。现在终于迎来了“传记小记”的结语,也就是我打一开始就想谈到的地方:

\begin{quote}
“在结束时,我想强调与往昔数学家(从1800年至今)的著作保持接触对我何其宝贵,无论直接接触还是通过比我更博学的人(如A.~Weil与J.-P.~Serre)转介。我们‘是骑在巨人肩膀上的矮人’,最美丽的现代数学理论,都是由解决他们遗赠给我们的若干问题的希望所驱动的。

Pierre~Deligne”
\end{quote}

一如常见的情形,我对这几行——某种向往昔宣誓的信仰表白——的第一反应停留在表面,停留在字面含义——但我必定隐隐感到,在这之外另有玄机。这句引文(无疑出自某位著名数学家,而我理应“同所有人一样”知晓它)并不合我的意。我在其中感到一种刻意为之的谦逊、甚至谦卑,它十足是一副姿态,根本不符合事情的单纯实在。而且这一蓄意几乎近于荒谬:如果每一代人在格局上都比前一代“更小”,人类早就耗尽元气而灭绝了,沦为一群可悲的侏儒!我深知人身上的创造力在今天并不比一百年前、或一百个世纪前更少(无疑也不更多)。我也深知,单就数学而言,我所熟识的人们的这样那样的著作(包括我自己在内),完全配得上列在往昔最伟大的数学家名下。我也深知,\textbf{我}在做数学时的动机\fnote{372}{(*)}{顺带说一句,包括Pierre~Deligne本人!},以及我昔日在数学界的大多数朋友的动机,绝不在于“解决前人遗赠的若干问题的希望”!若非如此,我们的科学将无力自我更新——它将不再具有创造性。

在这篇借来的信仰表白中,更让我震惊、或者更确切地说是更让我\textbf{痛心}的,是我深知做出这一表白的人,比我所认识的任何其他人都更蒙受了令我赞叹的“禀赋”,而且我曾在他对待数学的进路中确认了一种“新鲜之感”,凭此他理应做出大事业,如极少有数学家曾有特权做出的那样\dmark{}\origpage{816}。我心中有一份痛惜,也有某种怨恨,因为在那个自诩于同往昔之人的交往中获得谦卑的人的姿态背后,我感到了一场\textbf{退却}。退却了那份在他体内的创造力——他似乎很久以前就已将它遗忘,而正是那份力量使他成为远\textbf{非其他}、远胜于骑在巨人肩膀上的矮人这一意象所暗示的存在\fnote{373}{(*)}{(2月25日)这种“退却”的印象,强烈地同葬礼悼词所引发的印象联系在一起。参见“葬礼悼词(2)——或力量与光环”那条注(n\textsuperscript{o}~105)末尾对此的唤起。}。

这是自初读小传以来,我第一次试图界定那次阅读最初在我心中唤起的情感。在随后的日子里,甚至并非出于我的刻意图谋,这层思考一直在发酵。尤其是最后这段话一直萦绕在我脑海中,像一件确乎不同寻常、难以释怀的事情。在这篇短小的传记文字结尾那表面荒谬的表白背后,我必定隐隐预感到一种\textbf{意义},它最初无疑停留在无意识层面,然后逐渐浮向表层,其间就我记忆所及并无严格意义上的反思。毕竟我很清楚,我的朋友皮埃尔并不比我更习惯于钻研往昔的著作。就算他读得更多,那也不是陈年古卷,而是那些在消息灵通的圈子里流传的最新抽印本与预印本,他总是最先拿到它们。我也同样清楚,我的朋友从我离开数学舞台以来(乃至在此之前)滋养其工作的灵感,绝非源自皮卡或上个世纪甚至本世纪其他受人尊崇的先驱。如果说他确实乐于“骑在某人的肩膀上”,那绝非在一份公开的、纯粹修辞性的信仰表白之中,而是秘密地、\textbf{切实地}骑在上面——毕竟在反思某场葬礼之际,我处于再合适不过的位置,去知晓\textbf{谁}是那个承受代价、吞下苦果的人!人们口头上用“往昔的伟人”替换了那位从不被提及的人\fnote{374}{(**)}{或者,如果实在无法回避,就假意以迂回的方式、“暂停!”的风格来提及他。}——而在上一段里,人们刚刚心照不宣地把Motif(别名“如今不过一副模糊的骸骨”)的父权归于那些伟人\dmark{}\origpage{817}——从而使这副替代面孔背后的\textbf{真正}身份变得更加昭彰。

我曾多次观察到,人身上有一种力量,似乎具有普遍的本性,驱使人在排除万难之中,往往以迂回和象征的方式,去表达那些无法公开显露的愿望与企图(无论清醒的还是潜意识的),从而给予它们一种出路和满足——这种满足(在“理性”的措辞与通行的共识看来)或许显得微不足道,却并不因此而不具实质意义。在某种意义上,这是一股推动我们(仿佛身不由己地)向愿意倾听的人宣告我们存在的真相的力量(而我们每个人心中都有一个听觉敏锐的孩子……),而且\textbf{哪怕}这样“宣告”出来的东西无论在他人面前还是在我们自己面前都是最严重的禁忌。表达这一力量的首选领地是梦,这也是为什么梦是通往我们自身认识的一把无比强有力的钥匙。然而,正是由于梦的亲密、个人性质——它只向我们自己显露我们自己——这一表达手段对我们来说绝不充分,因为它无法在\textbf{他人面前}、甚至在象征意义上在全世界面前证实我们存在的真相。正是因此,在每一个看似挑战理性的荒谬无稽背后,都隐藏着一种“意义”——或者更确切地说,荒谬无稽正是潜意识以无误的本能所选中的\textbf{特选表达手段},用以\textbf{宣告这一意义}:它既被隐藏,又被堂而皇之地昭示在所有人面前\fnote{375}{(*)}{关于借表面的荒谬无稽而大声宣扬的\textbf{意义}的另一个格外张扬的例子,见“玩笑——或‘权复形’”那条注(n\textsuperscript{o}~83)。亦见“表面与深处”那条注(n\textsuperscript{o}~101)末尾(第440页)以及紧随其后的“书写的颂歌”(n\textsuperscript{o}~102)中的评注。}!

这正是我在读到这个“荒谬之谈”之后的日子里朦胧感到的东西:一个“矮人”(天生本该是巨人)骑在一个“巨人”(其手段比那个自诩“矮人”的人远为质朴,而后者骑在他身上却又否认他……)的肩膀上。阻碍我清晰意识到这个荒谬所揭示的意义的原因之一\dmark{}\origpage{818}\fnote{376}{(**)}{另一个原因,而且在我看来是主要障碍的,是面对“相信自己眼睛的见证,哪怕所见之物闻所未闻且被所有人否认”时的一定程度的\textbf{怯懦}。最近在“义务已尽——或真话时刻”(n\textsuperscript{o}~163)那条注里我又遇到了它。尤见第782页脚注(**),我在那里探究了这种“不相信”。},无疑是我不愿在那张生硬的“巨人”画像中认出自己;或者更确切地说,是不愿在那张确实曾属于我的品牌形象中认出自己——而通过这个荒唐的荒谬无稽的意外转介,它突然在拷问着我!直到数周之后,在12月18日那条注“神经中的神经——或侏儒与巨人”(n\textsuperscript{o}~148)里,我才终于回到矮人与巨人的奇特意象上,这一回是在掌握实据的条件下,处在关于葬礼的反思语境完全准备好接纳它的时刻。这个意象(就在当天)随即显现为一个关键的“力量意象”,对于理解我朋友对我的关系具有决定性意义,更深层、更关键的是,它是理解我朋友同他自身关系(注定永远只能是局部的理解)的开端——也就是理解\textbf{他自身人格中分裂}所采取的特定形态。既然这场葬礼首先由我的前学生兼前继承人发起并实施\fnote{377}{(*)}{诚然,在这一“实施”中,他是与“教团”密切勾结而行动的,在某种程度上充当了实现集体意志的工具。(见“掘墓人——或全体教众”,n\textsuperscript{o}~97.)但这一力量意象完全可能同样存在于该教团的“集体无意识”层面,并在其众多成员的个体无意识中找到共鸣,特别是在曾是我学生的一些人身上(而不仅是Deligne一人)。(5月12日)自写下这些文字以来,这一直觉大有进展,如今已显而易见。关于这一点,见“信使(2)”那条注(n\textsuperscript{o}~181)。},那么正是这同一股力量,在我看来成了贯穿这场漫长调查中始终顽固运作的枢纽力量,成了它的真正\textbf{神经}。在它出现在笔记中之后的两周里,它一直处于反思的核心,贯穿12月18日(已引的“神经中的神经——或侏儒与巨人”)到1月3日“敌对的兄弟——或传递(2)”(n\textsuperscript{o}~156)之间接续的九条注。

这个看似平淡无奇的意象在我的反思中所扮演的枢纽力量意象角色的“有效性”,也就是这股力量在我朋友自身心灵中的\textbf{实在性}问题——作为深层冲突的表达与不可遏制之补偿行为的驱动力\dmark{}\origpage{819}\fnote{378}{(*)}{用“不可遏制”一词,我绝非暗示这股力量的存在具有不可抗拒的宿命性,从而免除我朋友的责任。这样一股力量在我们体内的运作,完全取决于我们乐于并执意回避对它的认识,以便坐享通过这种刻意的“无知”所“买来”的微小特权。代价是高昂的,诚然,但不去认识这一点正是“交易”本身的一部分。}——这一问题在我看来无法通过“证明”来解决,即无法通过一种旨在赢得任何善意且充分知情之对话者认同的“客观”进路来解决。对我而言,这种实在性无可怀疑,我的深切信念并非某种“论证性”步骤的结果。诚然,在刚才提到的两周反思中,它得到了深化(我在此不打算对其作“总结”或“盘点”)。但它从第一天起就已存在——从我在阅读之后第一次劳神将阅读带给我的灵感落于纸上那一刻起,仿佛听凭一个无声之音的口授\fnote{379}{(**)}{我感到这种由“无声之音”口授的意象在《收获与播种》中不止一次出现,每次都自然而然。这绝非某种文学套话的重复,而是(在我看来)发现过程的一个普遍方面,因情境不同而或显或隐。}。我随后“想起”了深处本已“知晓”的事情。我必定通过某种感知能力“知道”它——这种能力固然非同寻常,但比我们在对事物的\textbf{清醒}认识层面通常运用的能力精微得多。那些压制在我们内心“某处”被感知到、却与我们例行接受(或者说,\textbf{视而不见})周遭现实的逻辑不合之物的机制——无需多言,它们在我身上与在任何人身上一样运作。在这方面如果说我与他人有什么不同,那就是我终于意识到了它们在我体内的无声运作,尤其是自打我开始“禅观”以来:在不可遏制的好奇心驱动下,我有时会下功夫去\textbf{驻足}在那些我感受到的事情上,其效果便是将深层模糊隐匿之物\textbf{召唤上升}到意识表层,并使之成形。

而且初始的感知在给它赋予形状的\textbf{工作}过程中会发生转变,在光天化日之下成形\dmark{}\origpage{820}。这项工作同时是一场\textbf{沉淀},通过它,感知的转译(转译为明白晓畅的言语)逐渐摆脱最初玷污它的主观先验。在本例中,那些扭曲的先验之一(在方才所引的最后一条注中被察觉)便是我体内根深蒂固的倾向,促使我“以阳看待自己”,即便在显然是我存在的阴面——“我身上的女性”——提供理解钥匙的情形中亦然(或者至少提供了不可或缺的钥匙之一,或提供了一份细致理解所需的“视角”之一)。我在别处谈到过那些固然完全“主观”却无可置疑的迹象,它们向我昭示这样一项工作的\textbf{进展},并在我误入歧途或陷入暂时停滞时发出警告,一旦察觉,停滞便随即消解\fnote{380}{(*)}{关于这一点,见“孩子与海——或信与疑”那条注,n\textsuperscript{o}~103.}。

\subsection{(3) 墓上的旋律——或自负}

\notehead{167}{}

(2月25日)昨天的大部分时间都在给一位年轻同事Norman~Walter写长信,他似乎有志于投身Motif理论,不为决然谈不上鼓励人的大环境所吓倒。这一回是密密麻麻的八页纸(打字机打出),谈Motif范畴与最重要“系数范畴”的“六种运算”。这让我再次惊愕地意识到:自该问题被提出以来的二十年间(诚然不是在文献中……),除了相对于基概形~$X$互素的“$\ell$-adic系数”这唯一的例外之外,至今没有一种合用的概形上同调的系数“好范畴”被成形构造出来(sic!);而且那项工作,在三角范畴框架下(六运算形式主义不可或缺的框架)由Jouanolou完成,却从未发表。我自己甚至从未手持过这位学生博士论文的一份复本\fnote{381}{(**)}{Jouanolou的博士论文工作缺乏真正的信念(这使他区别于我离开前的所有其他学生),一拖再拖,答辩直到1970年之后才举行。无论是答辩委员会的组建还是答辩本身,我都不记得曾得到通知,更谈不上被联系过。答辩后Jouanolou也未觉得有必要给我寄一份论文复本。去年我写信向他索取,他(不带解释地)告知我深表遗憾已经一份不剩了……(5月12日)我的记忆在此出现了偏差——事实上Jouanolou的论文答辩早自1969年就举行了!关于此事的详情,见最后的注(在写下这些文字时尚未写就)“第六颗钉子(钉进棺材)”。}。这些都是普遍背离我六十年代所开创的基础纲领的鲜明迹象\dmark{}\origpage{821},我当年绝未料到这项工作不会借着已取得的势头继续推进,而是在我离开数学舞台后即刻被硬生生斩断(或被“电锯截断”……)。

当素数~$\ell$~在概形~$X$~上是\textbf{幂零的}时,比如所谓“$\ell$-adic系数”范畴~$\mathbb{Z}_\ell^*(X)$\fnote{382}{(*)}{在选定理论的基环说明之后的星号$*$意味着我们不仅处理朴素的可构层(此处为合适意义下的$\ell$-adic层),而且处理由层构成的“可构”复形,它们是合适三角范畴的对象(哪怕可构层范畴$\mathbb{Z}_\ell(X)$已知,三角范畴的描述也非易事)。当我们想处理Motif时(通常是指“iso-motif”,即差同源意义下的Motif,构成具有$\mathbb{Q}$-系数的范畴),用于“实现”这类(iso)motif的自然系数范畴自身必须具有$\mathbb{Q}$-系数,人们将采用$\mathbb{Q}_\ell(X)$、$\mathbb{Q}_\ell^*(X)$。当人们想同时处理所有的$\ell$时,最自然的是处理“adelic”层(或复形),其基环是adele环$\mathbb{A}$,它是由所有系数范畴$\mathbb{Z}_\ell^*(X)$的积张量$\mathbb{Q}$得到的。人们会注意到,当素数$\ell$不与概形$X$互素时,在$X$上“$\ell$-adic系数”的定义中,$\mathcal{O}(X)$的幂零元不能被忽略——它们干预在$X$于$\ell$处的奇异性中。对于$X$上的adelic系数更当如此,这使它们接近于下一段中将要谈到的De Rham-Mebkhout系数(目前尚属假想)。必须理解,这两种主要的系数类型,adelic系数与De Rham-Mebkhout系数(后者被赋予下文所提及的所有丰富结构),是同一个\textbf{Motif}的两种不同描述(弱化描述)或“实现”,由二者之一极其紧密地圈定。关于这种“忠实性”,我早在六十年代就提出了若干猜想,类似于所谓的“Tate猜想”(Deligne将它们连同其余一切一并埋葬了……)。我打算在致力于数学思想的《数学反思》卷中再回到这里。人们感受到这两种系数类型之间的强烈亲缘关系(adelic与De Rham-Mebkhout,后者“差同源”)。后者优于前者之处在于,使其在某些方面显得“更精细”的是,它们的自然基域是$\mathbb{Q}$,而对于前者则是大得多的adele环。},应当正是“结晶系数”范畴,带有Frobenius态射~$F$~与相应的\textbf{过滤}。这一三角范畴的成形构造,更不用说其上的六运算,仍在等待有人着手进行。至于通过我六十年代末隐约预见的一个“神秘函子”,将上述“普通的”$\ell$-adic情形与前述“结晶”情形进行“粘合”,以获得对~$\ell$~无限制的系数范畴~$\mathbb{Z}_\ell^*(X)$~的定义,即便在最简单的非平凡情形~$X=\mathrm{Spec}(\mathbb{Z}_\ell)$~下仍未完成\dmark{}\origpage{822}。至于一般概形上的De Rham-Hodge系数~$DRHdg^*(X)$\fnote{383}{(*)}{(5月12日)正如后文将看到的,这一临时起意的名称与记号被证明是不合宜的,应由记号$DRM^*(X)$或$\mathrm{Meb}^*(X)$代替,对偶于$DRD^*(X)$或$\mathrm{Del}^*(X)$,分别代表De Rham-Mebkhout与De Rham-Deligne系数。后者已被其生父遗弃,并在1985年由我在充分知情的情况下重新拾起,作为Grothendieck装备库中基础成分之一(通过Mebkhout)……},我并无清晰的描述思路,Deligne也未能给出真正令人满意的界定。这里的创新思想归功于Zoghman~Mebkhout——大家都知道他是在何等恶劣的逆境下工作的,以及在他的思想之深远影响(极为有限地)被承认之后,他本人遭遇了怎样的对待。无论如何,人们终于有了一条可靠的线索,借助于对绝对基~$\mathrm{Spec}(\mathbb{Z})$~的“晶体”施加有限性、完整性与正则性条件,并或许辅以“De Rham过滤”与“按权过滤”,来着手构建~$DRHdg^*(X)$~范畴——并有望构建出不仅局限于特征零、而且在给定的正特征下能够或多或少还原出“通常的”结晶系数(sic!)的东西。非同寻常的是,我似乎是世界上唯一感受到这项任务的人——Zoghman~Mebkhout本人在尝尽苦头之后,似乎也再无心思为\textbf{他的}哲学的奠基问题多花哪怕一天时间!我何苦多言,既然我看到Deligne在Hodge理论上以身作则,硬生生截断了“在我那个时代”曾激励他并催生出大有前途(却未兑现……)之进路的那股冲劲。我很清楚,复代数簇~$X$~之上的Hodge系数形式主义(甚至尚未成形)理应或多或少包含在刚才按我六十年代的语言习惯称之为“De Rham系数”或“De Rham-Hodge系数”的形式主义中——以唤起De Rham\textbf{过滤}对象与相应的分次对象(所谓的“Hodge对象”)之间的联系。鉴于Mebkhout哲学在理解这些系数范畴(诚然目前仍属假设)中的决定性作用,无疑最好将它们称为“\textbf{De Rham-Mebkhout系数}”(记号~$DRM^*(X)$),或勉强称为“De Rham-Hodge-Mebkhout系数”($DRHM^*(X)$)。当~$X$~在复数域~$\mathbb{C}$~上为有限型时,人们应当能够从De Rham-Mebkhout系数出发,或多或少“重言式地”复原出假设的Hodge系数范畴\dmark{}\origpage{823}$HDG^*(X)$(我绝不称之为Hodge-Deligne范畴,因为Deligne似乎竭尽全力隐藏这个问题,绝非揭示它!),以及其上的六种运算——只需在De Rham-Mebkhout系数上简单添加一层被称为“Betti结构”的额外结构(具有超越性质)。因此在我看来,目前摆在代数簇上同调“自然系数范畴”描述面前的主要问题\fnote{384}{(*)}{在某种意义上,这些问题是为Motif的yoga赋予其应有的精度与一般性(正如我自六十年代起所设想的那样)的预备性问题(或被默认解决的问题)。}如下:

1. 对于给定的素数~$\ell$~与\textbf{任意}概形~$X$(不必然“与~$\ell$~互素”),给出~$\ell$-adic系数范畴~$\mathbb{Z}_\ell^*(X)$~的描述,以及这些系数的六种运算形式主义。(这一问题或多或少等价于“神秘函子”问题。)

2. 对于任意概形~$X$,给出“De Rham-Mebkhout系数”范畴~$DRM^*(X)$~的描述,或对相对概形~$X/S$~给出类似范畴~$DRM^*(X/S)$~的描述(通过定义约定~$DRM^*(X)=DRM^*(X/\mathrm{Spec}(\mathbb{Z}))$),以及这些系数的六种运算形式主义。

对于问题2),根据人们希望在这些系数中引入的结构丰富程度,可能有几种变体。无论如何,“好上帝定理”(别名Mebkhout定理)先验地向我们表明(至少对于复数域上的有限型概形~$X$),必定存在Mebkhout式的结晶系数的六种变换形式主义,而无需在其上添加De Rham或/与按权的过滤。对于一般概形~$X$上的一个Motif(或一个“绝对系数”)所关联的De Rham-Mebkhout结晶复形~$K$,必然存在的第三类重要的额外结构,是对每个素数~$p$~给出一个“Frobenius态射”$K(p)^{(p)}\to K(p)$,其中~$K(p)$~表示对由模~$p$~约化得到的子概形~$X(p)$~的限制,上标$(p)$表示~$K(p)$~的“Frobenius化”,即它在Frobenius映射下的拉回。因此,根据人们打算在一个结晶复形中引入的额外结构(刚才提到的三种之中),先验地可以预见关于“De Rham-Mebkhout系数”概念的\textbf{八种}变体。唯有实地工作\dmark{}\origpage{824}才能向我们指明其中哪些变体确实能产生六种运算的形式主义。同样真实的是,出于Motif的需要,当我们旨在寻找能够尽可能贴近Motif并尽可能忠实丰富地描述其结构的简单“代数”对象时,先验地“最丰富的”系数显得“最好”。正是出于其巨大的丰富性,Hodge系数具有主要的魅力——甚至达到了这样的程度:如果Hodge猜想为真,人们可以希望在~$\mathbb{C}$~上完全重建Motif范畴,甚至对于~$\mathbb{C}$~上的任意有限型概形~$X$~亦然。

这提醒我注意到,某些结构可能是“多余的”,并从其他结构中推导出来(但确实是以一种极其隐蔽的方式,以至于很难用平实的语言明确表达)\fnote{385}{(*)}{作为朝同一方向的评注,我在此指出必须注意施加于某一类型“上同调系数”所关联的全部结构之上的或隐或显的相容性条件。我在此特别想到那些在“可Motif化”系数(即源自Motif的系数)上自动满足的相容性(具有或多或少的代数性质)。如果人们希望拥有一套“六种运算”的形式主义(哪怕独立于尽可能贴近Motif的愿望),很可能必须将它们强加于所考虑的系数之上。我特别想到Mebkhout理论中的完整性与正则性条件,以及(如果将De Rham过滤作为额外结构引入)联系De Rham过滤与Gauss-Manin联络的相容性条件。这些例子在我看来相当清楚地表明,描述上同调系数的“好范畴”这一根本任务,伴随着彻底探索和充分利用迄今为止在“代数簇上同调”上构想的所有结构以及联系这些结构的深层关系的约束。这从一开始就是Motif的yoga的主要宗旨:在一个差异性背后显现一份\textbf{统一},并同时提供一条可靠的线索,让人在不可穷尽的丰富性中认清方向。}。例如,在概形~$X$~关于另一个概形~$S$~的(相对)De Rham上同调上,我(在六十年代末)\fnote{386}{(**)}{(5月2日)实际上是在1966年。}揭示了一种无曲率的(绝对)典范联络的存在,我称之为\textbf{Gauss-Manin联络}。由此可知,Deligne赋予~$\mathbb{C}$~上光滑概形~$X$(乃至赋予~$\mathbb{C}$~上任意有限型概形~$X$)的Hodge-Deligne结构,典范地配有这样一个相对于素子域~$\mathbb{Q}$~的联络。只要Motif上同调本身已经能从其“Hodge实现”中复原出来\dmark{}\origpage{825},这意味着在任何可以称之为“Motif的”或“代数的”(即源自一个Motif)Hodge结构上,都将存在这样一个典范的Gauss-Manin联络。同样不难描述关联于一个Hodge-Deligne结构的其他更为精妙的典范结构,其存在性“源于Motif”:某些Galois群在~$\mathrm{Bet}(K)\otimes_\mathbb{Z}\mathbb{Z}_\ell$~上的作用(其中~$\mathrm{Bet}(K)$~是Hodge-Deligne结构~$K$~底层的“格”),以及在“模~$p$~约化”上的“Frobenius结构”(对几乎所有~$p$)。正是这种表面上毫无关联的结构的丰富多样性——其隐藏的纽带正是所有这些结构共有的\textbf{“motif”}——这种丰富性于我曾代表(且依然代表)代数簇上同调这一主题的特殊魅力,以及“Motif”的魅力,它们如同赋予这一具有无数变奏的主题以生命和意义的共同精致旋律\fnote{387}{(*)}{(3月26日)在对我昨天笔记开头(n\textsuperscript{o}~176)所涉及的各种类型上同调系数(用来“圈定Motif”)的(紧密相关的)问题以及一个上同调类(来自一个代数闭链)必须满足的“代数条件”进行了简短反思之后,我决定从《数学反思》第3卷(包含《收获与播种》之后的文字)开始,着手写一篇关于Motif、“系数”与标准猜想的实质性文本。我现在相信我已经掌握了对一个概形上的“该”Motif范畴进行成形描述的原则,至少在关键情形(通过极限过渡应当能够归结到的情形)下:在绝对基~$\mathbb{Z}$~上的有限型概形的情形。作为相对于我六十年代思想的唯一新成分,我引入了“Mebkhout范畴的理念”(由“好上帝定理”所表达)。此外,我假设“神秘函子”的问题已经解决(目前尚未解决!),它在我计划给出的完全描述中起着决定性作用。}。

如果说除我之外还有谁听到了并感受到了这首旋律,并在它喷涌展开之际长久地沉浸其中,那正是Pierre~Deligne。如果说我曾向谁托付过一样活着的东西——一样精致而强健、我倾注了自己最好心血、多年来以我的力量与爱所滋养的东西——那就是他。那是一样生来便要在光天化日之下绽放、生长与繁衍的东西——一颗种子,随时准备传递其内在的生命。昨天和今天的短暂接触,是一场与我久已失去联络之物的\textbf{重逢}——并非与词语或概念、亦非与僵死之物的重逢,而是与一样充盈着浓烈\textbf{生命}的东西的重逢\dmark{}\origpage{826}。它让我再次掂量出:我所留下的这件“东西”足够宏大、足够深邃,足以启发一位全副身心投入其中的数学家的整整一生,以及他身后的其他数学家——因为他的生命无疑不足以穷尽这项任务\fnote{388}{(*)}{(3月26日)我现在看来可能高估了这项任务的规模(但绝非其深度!)。关于这一点见同日的上一条脚注。}。

这是一种奇怪而令人欣慰的巧合:这次相遇发生在我刚刚经历了另一场同样不期而至的“相遇”之时——读到了我朋友的文字,在其中他避而不提其名,恰恰谈论着我托付于他的所有事物中最切中我心怀的那件东西:“如今人们对它的认识,大概不过剩一副模糊的骸骨……”

这几句话在过去的这三天里一直萦绕在我心头。我分明认出了那种自负——属于那个“任何东西都不够美丽以致他不屑为此欣喜”之人的自负。而且,并未刻意去想,关于“\textbf{坟墓}”的记忆又浮上心头\fnote{389}{(**)}{关于这种强烈且长久未曾言表的印象——它在我从巴黎“归来”并将《追寻栈》托付给Deligne之后一直纠缠着我——见“坟墓”那条注(n\textsuperscript{o}~71)。}。同一种印象在我心中复活了,带着同一个无声而执拗的意象。那件曾为我所珍爱的活着的事物,我当年曾以为托付给了充满爱意的双手——而在我失去它的这十五年间,它却在远离风、雨和阳光的恩泽的一座坟墓里衰败腐烂。今天我发现它已毫无血色,“一副骸骨……”,成了一个乐于\textbf{利用}它却绝不愿\textbf{付出}自己的人那居高临下的鄙夷对象。

\section{5.~四种运算(在遗骸上)}

\subsection{(0) 侦探——或玫瑰人生}

\notehead{167$'$}{}

\dmark{}\origpage{827}(4月22日)原定接在这里的注,很早以前就预定名为“四种运算”(这个名字将在下一条注的开头得到详尽的解释\fnote{390}{(*)}{(5月12日)在把原来那条旧注“沉默”(n\textsuperscript{o}~168)一分为四之后,“下一条注”指的是注~n\textsuperscript{o}~167$''$,即“四种运算——或一场调查的‘梳理’”。})。我原本打算把这一“梳理”工作——对一场当时在我看来已经结束的调查的梳理——写成一条注,顶多两条。从那时起已经过去了近两个月,鉴于意外波澜的涌现,直到此刻我甚至还没能完全巡览这个主题。在整整一年之后,发现葬礼的那一幕充满惊奇的剧本仿佛又以更高的调门重演了一遍!

最终,在目录中,出名的“四种运算”所指称的不再是一两注,而是一整套卷帙浩繁——恐怕略显芜杂——的\textbf{三十}条注与子注\fnote{391}{(**)}{(5月12日)自写下这些断言性的文字以来,这一数目又有所增加……}。它们分成八个部分(1)到(8),冠以(我希望是)富有提示性的名称——从(1)“藏宝”到(8)“第六颗钉子(钉进棺材)”。在此过程中,我不得不对四条注作了彻底的重构\fnote{392}{(***)}{(5月12日)这些注的篇幅已膨胀得难以容纳,最终每条又分成数条,成为注n\textsuperscript{o}s~168(i)--(iii)、169(i)--(v)、170(i)--(iii)、171(i)--(iv)。},它们曾构成“四种运算”在2月26日至3月1日期间写下的“初稿”。在3月22日(一个月前)那条注“门槛”(n\textsuperscript{o}~172)的开头,我就违背在《收获与播种》其余地方一贯秉持的精神作了解释。

所涉及的四条注是:\textbf{“沉默”,“操演”,“分赃”,“封神”}(n\textsuperscript{o}s~168, 169, 170, 171)\fnote{393}{(***)}{(5月12日)这些注的篇幅已膨胀得难以容纳,最终每条又分成数条,成为注n\textsuperscript{o}s~168(i)--(iii)、169(i)--(v)、170(i)--(iii)、171(i)--(iv)。},相继致力于勾勒一幅\dmark{}\origpage{828}全景,涵盖围绕我的著作、继而是Zoghman~Mebkhout的著作展开的四场“大操演”——隐匿与侵吞的操演。我建议读者最初只读这四条注,略过脚注(这里的脚注比《收获与播种》任何其他部分都更繁密),并略过正文中引用的子注(它们同样异乎寻常地繁多且厚重)。他随后可以沿着这一势头继续阅读接下来的四条主要注:\textbf{“门槛”“家庭相册”“节节加码(2)”“殡仪馆——为科学服务”}(n\textsuperscript{o}s~172--175),它们已经全无技术性可言。

好奇的读者若想更细致地了解这“四种运算”曲折迂回的迷宫,第二遍阅读时可以把脚注与子注也包括进去;甚至(若他尚未读过葬礼第一部,或觉得需要唤醒阅读的记忆),可以随时(一如我自己常做的)回头查阅文中大量指涉的葬礼I(或“中国皇帝的新衣”)的段落。

构成“四种运算”(或描述并评议它们)的三十条注,每一条的核心内容都没有技术性质。依我看,任何有兴趣、有头脑的读者,即便完全不是代数簇上同调的专家,甚至不是数学家、谈不上多少“科学素养”,也能读懂。对于仍然犹豫、怕一头栽进“诈骗艺术”种种奥秘的人,我尤其推荐以下子注:它们的内容在我看来最为丰厚,其意义显然超出“拆解”那些时而荒诞绝伦、却总是精心布置(专供巴不得受骗的人……)的“勾当”所能带来的兴趣。这些子注是“驱逐”(n\textsuperscript{o}~169$_1$),接着是“正经数学……”“……与无稽之谈”“遗产——或骗术与创造”(合起来是“公式”这一总题下五条子注中的头三条);最后是“封神”(n\textsuperscript{o}~171)那条注里关于Zoghman~Mebkhout奇异经历的四条子注:“眼光的绽放——或闯入者”“黑手党”“根”“掠夺的空白授权”(n\textsuperscript{o}s~171$_1$至171$_4$)。因此,我特别提请读者留意的是这八条子注(总数二十一条中的八条\fnote{394}{(*)}{(5月12日)其间已增至二十七条,还不算“第六颗钉子(钉进棺材)”(其中又有七条愉快可口的注)。})。

至于另外十三条子注,那些对其“文献价值”不感兴趣的读者,不妨在闲暇之余,权当阅读一部扣人心弦的侦探冒险传奇:在其中,一位半路出家的业余侦探(即区区在下)追踪并搜集种种“蛛丝马迹”——有的微妙难寻、若隐若现,有的则庞大显眼到无人再看得见的地步;而这些线索最终拼合为一幅色彩浓烈且不容置辩的(世态)风俗画卷:其中一位“面带微笑、和蔼可亲的当代凡尔杜先生(别名朗德吕)”\dmark{}\origpage{829},在四邻所有善良市民温柔(乃至赞叹)的注视下,对他那些纯洁无辜的受害者施以大卸八块与焚骨扬灰的极刑。他们久已习惯了这种稍显特殊的味道,显然不再让任何人感到不适。甚至不止一人效仿起他那位讨人喜欢又机敏过人的邻居,烟囱们争先恐后地嗡嗡作响、油烟四溢。

彻底开了眼的“侦探”,唯有蹑手蹑脚地退场:显见得,此地达成了一致的赞同,而一切都好极了,在这最佳的世界里……

\subsubsection*{四种运算——或一场调查的“梳理”}
\tocsubsub{四种运算——或一场调查的“梳理”}

\notehead{167$''$}{}

(2月26日)\fnote{395}{(*)}{这条注原本应当叫做“四种运算”,紧接在“墓上的旋律——或自负”(注n\textsuperscript{o}~167)之后,比前面那条(具有导论性质的)注“侦探——或玫瑰人生”(n\textsuperscript{o}~167$'$)早了近两个月。我建议先读后一条。}看来我对葬礼已经作了或多或少的巡视。当然是不完全的巡视,而且是暂时的——但就目前而言,我认为自己不会走得更远了。我感到如今需要的是抽身抽离的纵深,现在是收尾的时刻了。剩下的是盘点我在写作《收获与播种》这一即兴禅观过程中所学到的东西。

关于葬礼的反思构成了我绝大部分的工作。这一反思在两个截然不同的层面上有序推进。首先,在极其必要的“敬重之举”——即双重注“我的孤儿”与“对一份遗产的拒绝——或一个矛盾的价格”(n\textsuperscript{o}s~46, 47)之后,是对“全部辉煌中的”葬礼的渐进发现。七八年来我确实嗅到了这股风气——面对一部著作和某种风格的“轻微嘲讽之风”,以及面对那些表面上仍受其启发并在某种方式上\dmark{}\origpage{830}“顶着我的名字”的人的同样谨慎而无可挑剔的“闭门羹”。这是通过一种风尚和一个“无懈可击的共识”所体现的葬礼面向,在“掘墓人——或全体教众”以及先前的诸注(n\textsuperscript{o}s~93--97)中得到了审视,构成了第X行列即“灵车”。这一面向在过去岁月中由于未能花功夫进行周详的反思而一直处于弥散状态,在工作过程中得到了相当大的澄清,但对我而言并未增添真正崭新的事实。

相反,崭新的事实——我去年4月19日首次面对的事实,或者如果你愿意称之为“奇特事件”——是围绕我的著作、以及在离开数学舞台后唯一承担起“Grothendieck继承者”这一吃力不讨好且凶险角色的数学家Zoghman~Mebkhout的著作所展开的某场大规模行动。

去年4月19日所做的发现(1982年的Lecture Notes 900卷,在十二年的死寂沉寂之后\fnote{396}{(*)}{(4月19日)关于这“十二年”的更正,见子注“先期起棺”,n\textsuperscript{o}~168$^1$。}Motif在其中被起出,且完全未提及我这个人),成为了狭义上的调查的起点:一场关于我的著作所遭受命运的调查,首先是那些作为其首要和主要托管人的人——我的学生们——所给予的命运。这项调查揭露了大量事实,一个比一个出乎意料,随着时日推移,它们汇聚成一幅关于葬礼究竟为何物以及其主要行动者是谁的某种外部图景。这幅图景无疑并不完整,但它在极其精确和无可争辩的细节上足够丰富,足以满足我在这方面的全部好奇心。这是我刚才提到的反思的两个“层面”中的第一个。它基本上对应于关于葬礼反思的“第一口气”,从4月19日持续到6月10日左右,并以“疾病插曲”告终。

这也大体上对应于我笔记中的“葬礼I”(或“中国皇帝的新衣”)部分。此外还必须加上5月12日的“葬礼悼词(1)——或恭维”(n\textsuperscript{o}~104),它(无疑带有一点随意性)被划归到后续也是最后的行列“葬礼仪式”中,属于“葬礼II”的一部分。我还将紧接引述之注的下一条注——即“葬礼悼词(2)——或力量与光环”(n\textsuperscript{o}~105)\fnote{397}{(*)}{这条注早在5月12日写完前一条注“葬礼悼词(1)——或恭维”的次日就已预定。我那时意识到我刚刚略加审视的文本是一座真正的富矿,远未挖掘殆尽……(关于葬礼悼词的一些细节,见“封神”那条注的开头,n\textsuperscript{o}~171)。}\dmark{}\origpage{831}——归入这构成反思“第一层面”的“调查”之中,并且它还在下一条注“肌肉与肚肠(阳葬阴(1))”(n\textsuperscript{o}~106)的评注中延续。最后这两条注写于9月底至10月初。同样,在“葬礼悼词”谱系中,即在Deligne或多或少提及我的(极其罕见的)书面文献谱系中,人们可以将Deligne的小传最近所引发的两条注归入这项调查,即“为一副模糊的骸骨作安魂曲”与“信仰表白——或伪中之真”(n\textsuperscript{o}s~165, 166)。最后,还应加上“把一个个~i~点清”(n\textsuperscript{o}~164),它给出了一系列说明(主要是事实细节),大部分由Deligne本人在去年10月访问我处时提供\fnote{398}{(**)}{关于这次来访,见“义务已尽——或真话时刻”那条注(n\textsuperscript{o}~163)。}。

在结束了三个多月的一切智力活动的疾病插曲之后,反思的“第二口气”(或我刚才谈到的“第二层面”)是由理解4月和5月的调查所揭露的这一整套事实(其中某些确实非常重大甚至令人难以置信)的\textbf{意义}的努力所驱动的。这一反思的核心部分是“阴阳之钥”,在很大程度上独立于严格意义上的葬礼主题,尽管后者定期重现,每次都重新启动一场关于我个人、我的生活以及普遍存在的禅观。

显而易见,反思的这两个层面——“调查”与“禅观”——绝非相互独立或截然分开,而是彼此交融渗透。具体而言,这体现在贯穿葬礼第一部分的努力之中——去\textbf{理解}我日复一日所发现之事的意义;同时也体现在第二部分中依然出现的事实细节,它们补充了在初步“调查”中已获得的内容。

我目前的用意,是对在整个调查过程中逐日浮现的\textbf{事实}\dmark{}\origpage{832}作一个“盘点”或大纲式摘要——这些事实我至今从未劳神以稍具连贯性的方式加以整饬。因此这将是对我目前所知的针对我的著作\fnote{399}{(*)}{根据我所知的事实,这完全局限于我1955年至1970年间致力于发展关于概形上同调与(余)同调代数思想的著作部分。}以及Mebkhout著作的“大规模行动”的一场\textbf{梳理}。根据承受代价的是后者还是前者,以及根据被作为目标的是我著作的哪一部分,我实际上在其中区分出\textbf{四种}主要运算(总而言之即“四种运算”),我首先想对它们进行回顾。事实证明,它们在反思过程中引起我注意的顺序,也与它们在1970年我“离开”之后(甚至在此之前)启动的时间顺序相重合(最后两项仅有一点微小的倒置)。

\subsection{(1) 藏宝}

\subsubsection*{(a) 沉默(“Motifs”)}
\tocsubsub{(a) 沉默(“Motifs”)}

a1. \textbf{“Motifs”背景}

\notehead{168(i)}{}

\textbf{I~“Motifs”行动}

在从Serre的某些思想中获得启发、并出于为复代数簇的经典Betti上同调的各种已知(或预感到的)纯代数“化身”寻找一个共同“原则”(或“motif”)的愿望之下,我在六十年代初引入了“motif”的概念。在整个六十年代,尤其自1963年起\fnote{400}{(**)}{1963年是平展上同调(在1963/64年的SGA 4讨论班中发展)全速“启动”的一年,它终于为Motif反思的磨坊带来了充沛的水流,而此前这些反思多少带有一点思辨的意味。正是在接下来的这一年,我发展了“Motif Galois群”的形式主义,其详尽的概念基础(按照我提交给他的理论纲领)在N.~Saavedra的博士论文中得到发展,该论文直到1972年才发表(Springer Verlag, Lecture Notes n\textsuperscript{o}~265)。},在我整理基础理论的任务之余,我围绕这一主题发展出了一套既丰富又精确的“yoga”(或“哲学”)。这一宏大的理论当时仍属猜想性质,并且多半在未来的几代人中仍将保持为猜想性质\fnote{401}{(***)}{(8月8日)现在看来,这套理论并不像我曾以为的那样遥悬于“地平线上”——只要人们终于肯投身于它!关于这一点,见3月27日“守财奴与行将就木者”(n\textsuperscript{o}~177)中的评注。},然而它在当时(乃至直到今天)却提供了一个极可靠的向导,让人在涉及代数簇上同调的一切情境中认清方向,无论用于预见“人们有权期待什么”,还是用于启发该引入的“正确概念”,有时甚至是用于提供通向证明的进路。我在《收获与播种》引言\dmark{}\origpage{833}(“沉默的终结”,第XVIII页)中对此说道:

\begin{quote}
“在我有幸发现并公之于世的所有数学事物中,Motif的这一实在于我依然显得最为迷人,最富于奥秘——正处于‘几何’与‘算术’深层同一性的核心。而这一长久被忽视的实在所引向我的‘Motif的yoga’,也许是我在数学家生涯的这第一时期\fnote{402}{(*)}{如果我在此对‘我数学家生涯的这第一时期’作此限定,那是想到了‘阿那贝尔代数几何的yoga’,它在我看来具有相当的深度与广度。在《纲领概要》中对此有少许涉及,该文将作为《数学反思》的一部分附于《收获与播种》之后。}所提炼出的最强有力的发现工具。”
\end{quote}

除了为域上半单Motif范畴给出一种可能的显式构造(在众多其他构造之中)的临时草图之外,我在个人笔记中围绕这一主题所发展的思想一直停留在口头交流的阶段。我当时太专注于编写基础理论著作的大量其他任务\fnote{403}{(**)}{这首先是指EGA(代数几何原理,与Jean~Dieudonné合作)和SGA(布瓦玛丽代数几何讨论班),后者或由我单独撰写,或(特别是与学生们)合作撰写,依据我独创的指导思想与主控框架。在1959年至1969年间,这些无一例外成为标准参考文献的著作平均“产量”约为每年一千页。这一基础工作在我离开数学舞台之日便戛然而止。关于这一点见“阴,作为仆人,与新主人”(n\textsuperscript{o}~135)。},以至于抽不出所需的几个月闲暇来充实手稿笔记,以便为在我心中发展起来的内在视野编写一份足够“深入”以致在我看来可以发表的整体“主控大纲”。从1965年起直到1970年我离开数学舞台,我进行Motif(及其他)沉思的特选对话者,也是唯一完全领会了Motif的yoga并感受到其全部深远影响的人,就是Pierre~Deligne。

关于“Motif的yoga”的详细说明(比引文所属的引言部分更为详尽),读者可在“我的孤儿”(n\textsuperscript{o}~46)末尾找到,尤其是(关于该yoga的起源)在“一个梦的回忆——或Motif的诞生”(n\textsuperscript{o}~51)中。关于“Motif的yoga”融入六种运算形式主义(自我的离开以来直到今天依然被我的上同调弟子们\dmark{}\origpage{834}所忽视,尽管它是同调代数中的基础结构……),见“墓上的旋律——或自负”那条注(n\textsuperscript{o}~167)。关于围绕权的yoga(构成Motif的yoga的关键成分之一)与Hodge-Deligne理论(直接源于后一yoga)的思想亲缘(在文献中被彻底抹杀),见“把一个个~i~点清”那条注n\textsuperscript{o}~164(第II~4部分),以及紧随其后的子注(n\textsuperscript{o}~164$_{1}$)。

\subsubsection*{a2.~葬礼……}
\tocsubsub{a2.~葬礼……}

\notehead{168(ii)}{}

“Motifs”行动首先并在我离开数学舞台后即刻表现为:对Motif的yoga与“motif”一词本身的系统性\textbf{抹杀};继而在十二年的沉默之后\fnote{404}{(*)}{(4月8日)关于这“十二年”的更正,见本注“沉默”的子注“先期起棺”(n\textsuperscript{o}~168$^4$)。},随着该yoga的一个狭隘版本的起棺(于1982年),表现为对我这微不足道的亡故之人的\textbf{抹杀}——抹杀我与该yoga有任何瓜葛的事实。

对该yoga以“权的yoga”形式进行的首次公然抹杀,早在1968年就已发生,也就是早在我离开之前,出现在Deligne发表在《数学出版物》上关于谱序列退化的文章中。这首先在“罐头里的权与十二年秘密”那条注中谈及(写于发现起棺的“难忘卷册”之前),并在紧随其后的“驱逐”那条注的开头作了详尽阐述(注n\textsuperscript{o}s~49, 63)。

这场试探性的抹杀,在未遭任何反应之下\fnote{405}{(**)}{这种反应首先本该来自我。尽管那一篇文章呈现中的不诚实在我看来昭然若揭(参见引注,n\textsuperscript{o}~63),但我未能正视它,面对拿到文章读到有关段落时的“轻微不适”。关于我在反思过程中察觉到的自身某种顺从或暧昧所起的作用,见“暧昧”那条注,n\textsuperscript{o}~63$'$。至少在意识层面上,Deligne或我的任何其他学生可能存在不诚实念头从未掠过我的心头;或者说,我把它压制了下去,即便在不诚实显而易见并通过那从未被指认的“不适”向我发出警示之时。},在Deligne陈述其1968/69年发展的Hodge理论优美推广的Hodge~I、II、III文章中继续并加剧。虽然这一理论直接源于Motif的yoga(正如前文所回顾的),但在Hodge~II或Hodge~III中对此毫无提及——考虑到Hodge~II构成Deligne的博士论文,而他曾在我身边度过了其成长关键岁月\dmark{}\origpage{835},这一事实尤为触目\fnote{406}{(*)}{在Deligne与我之间存在某种默契,以抹杀其作为学生的身份——他是如此耀眼,以至于我无法自称曾是他的“导师”。我在“格格不入者”(n\textsuperscript{o}~67$'$)中揭示并检视了这一默契。}。

至于在1970年尼斯国际数学家大会上的简短“预告”Hodge~I,Deligne仅限于以半行晦涩的“暂停!”式提及指称“Grothendieck的猜想性Motif理论”(并同对Serre的一条假文献并列在一口气中,显然旨在掩人耳目\fnote{407}{(**)}{所涉及的是Serre关于Weil猜想的Kähler类似物的文章,它曾启发我关于“标准猜想”的工作。这是一篇出色的文章,在此绝非要加以贬低。但Deligne本人也很难解释这篇文章在何种程度上曾是其Hodge理论工作的“源泉”——无疑也从未有人想到去问他。近距离目睹了Deligne思想的萌发,我深知其源泉何在(参见前引注n\textsuperscript{o}~164$_{1}$),而且绝非Demazure关于Motif定义的初级ABC报告!Deligne引用这篇报告作为“Grothendieck的猜想性Motif理论”的凭据,以便给任何不深究的读者留下印象(而充分知情的人本就不多……),即所谓的“猜想性理论”仅限于Demazure报告中所写的内容,从而利用了该yoga在文献中缺乏更详尽发表痕迹的状况。})。这一做法在1974年温哥华国际数学家大会上关于“权的yoga”的报告中得以延续,会上Serre与我的名字都不再被提及。在这次交流中,如同在1970年尼斯大会的Hodge~I中一样,Deligne绝口不提他得自我的该yoga在Motif语境下(该语境被严格保密)的重要组成部分:权概念在“六种运算”下的行为,首先是通过$Rf_!$与$Rf_*$。这是众多例子中的一个,体现了一种如今司空见惯的做法,Deligne在我看来是其最早的推行者之一:即将某一数学领域中提出的“大问题”的知情权垄断给少数“圈内人”(甚至仅限其本人),以确保绝对霸权,而不是将其置于科学界共同支配之下以启发每个人\fnote{408}{(***)}{关于这种新心态(在写下这些文字时我从未遇到过痕迹),见“阴,作为仆人,与新主人”那条注,n\textsuperscript{o}~135,以及紧随其后的“操演”那条注末尾(4月28日)(x)。正是这种心态促使我给围绕我著作的前两场“大运算”的整个集合(n\textsuperscript{o}s~168至169$_8$)冠以“\textbf{藏宝}”之名。(x)这一末尾成为了“藏宝”那条注(n\textsuperscript{o}~169(v))。}。据我所知,\dmark{}\origpage{836}这一问题直到Deligne1980年在其“Weil~II”文章中(在$Rf_!$情形下)解决之前一直无人触碰,当然文中对我只字未提(正是向他传达了Motif语境下的相关猜想,而他所处理的$\ell$-adic语境不过是其反映,如同De Rham-Hodge系数语境一样……)。

在我了解(非常零碎地)Deligne工作或能够对其做出评价的范围内,我认为可以断言:他得自我的Motif的yoga构成了其全部工作的主要灵感源泉。他通过保持对Motif概念的绝对沉默直到1982年\fnote{409}{(*)}{(4月8日)关于更正,见前引子注“先期起棺”(n\textsuperscript{o}~168$^4$)。},将这一源泉隐匿了起来。唯一的例外(若无疏漏\fnote{410}{(**)}{如上一条脚注中所解释的,这一“暂停!”式提及的目的绝非为了承认或告知,而是为了(双重地)误导。至于Hodge-Deligne理论的思想谱系(在前引两条注中有所描述),我完全有理由相信除他之外无人知晓。})是那条“见证半行”,对于除他和我(以及或许Serre)之外的任何人来说都是无法理解的——正如稍早时他在关于谱序列退化的文章中晦涩地提及促使我猜想其退化结果“特例”的“权之考虑”(参见前注“驱逐”,n\textsuperscript{o}~63)。

\subsubsection*{a3.~……与起棺}
\tocsubsub{a3.~……与起棺}

\notehead{168(iii)}{}

随着Lecture Notes 900这一“难忘卷册”的出版,布景骤然改变\fnote{411}{(***)}{Springer Verlag, Lecture Notes in Mathematics, n\textsuperscript{o}~900, Hodge cycles, Motives, and Shimura Varieties, P. Deligne, J.S. Milne, A. Ogus, K.Y. Shih.}.Motif在其中大张旗鼓地被起出,初始yoga的一部分被公之于世。在这一卷册中,我的名字“顺带地”且仿佛出于极大偶然出现了两三次,绝不可能让读者怀疑我在其中发展的思想中起过任何作用。这些思想是以这样一种方式呈现的,即在读者心中无可怀疑:卷册的主要作者、辉煌的Pierre~Deligne刚刚发现了它们,并将它们新鲜端出。诚然,正如在尼斯或温哥华他未曾宣称是自己发现了首次在文献中阐明的权的yoga一样,这里也绝无任何一处\dmark{}\origpage{837}明确说是他发现了卷册中(表面上)首次阐述的所有这些优美思想——该卷册围绕一条确实出自他手笔的优美定理展开。这是他深谙的“暂停!”艺术,对此我首先在“暂停!”以及紧随其后的“中国皇帝的新衣”中作了评述(n\textsuperscript{o}s~77, 77$'$);亦见发现“难忘卷册”之前写下的更早的注:“下葬——或新父亲”、“新伦理——或抢夺的集市”以及“侵吞与轻蔑”(n\textsuperscript{o}s~52, 59, 59$'$)。

事实上,LN~900卷册中关于Motif的所有主要思想不仅早在六十年代就已为我所知(Deligne自1965年起完全有机会从我口中习得),而且该书的核心问题也早在六十年代末就由我提出(并且当然传达给了Deligne)。关于这方面的详情,见“把一个个~i~点清”那条注(n\textsuperscript{o}~164)(在该注的第一部分)。

正如我在《收获与播种》引言(“秘密的终结”,第XVII页)中所强调的,Deligne绝非我向其详述Motif的yoga的唯一之人,哪怕他是唯一深入将其内化之人。如果说在长达十年左右的时间里存在对这一yoga本身的彻底抹杀\fnote{412}{(*)}{根据Deligne好心提供给我的一份“Motif评注书目”,在我离开之后文献中关于Motif还出现过两次零星的工作,均直接源于六十年代的工作(在N.~Saavedra与我共同准备的论文中,以及在S.~Kleiman的一份报告中)。下一处指涉由R.P.~Langlands在1979年做出。之后便是1982年的LN~900。若无疏漏,在1982年之前Deligne发表的任何文本中都未出现“motif”一词,亦未在任何发表文本中(最多除注n\textsuperscript{o}s~165, 166中检视的传记小记之外)提及他曾从我口中习得任何东西……(4月8日)关于“若无疏漏”,见子注“先期起棺”(n\textsuperscript{o}~168$^4$)中的更正。},以及后来对我发现、发展和阐明它的角色的抹杀,这一抹杀若无我视为朋友的大批数学家的同谋是绝不可能完成的,特别是我的每一位(交换)“上同调学生”的同谋\fnote{413}{(**)}{我想可以说,我1970年之前的所有学生,唯一的例外是信女士(她自六十年代起就在越南工作),都知道(哪怕未必掌握)我关于Motif的想法,我曾在IHES就此做过一系列详尽的报告(1967年)。他们之中那些致力于代数簇上同调研究的人,在我看来对在主要“当事人”Deligne倡议下进行的下葬负有连带责任。这里主要涉及J.-L.~Verdier、L.~Illusie、P.~Berthelot,他们在所涉及的围绕我著作的其他三场运算中,表现得远比单纯的同谋更为主动。}。这一抹杀虽然服务于唯一一个人的可疑“利益”,却是由许多人共同的作为与不作为完成的。

\dmark{}\origpage{838}除Deligne和我的其他上同调学生之外,对这一“难忘卷册”LN~900负有最严重责任的,正是其\textbf{合著者},即\textbf{J.S.~Milne}、\textbf{A.~Ogus}与\textbf{K.Y.~Shih}。这些是我未曾谋面的数学家,我无意断定他们居心不良。对我而言,这丝毫不减免他们作为这部不同寻常卷册共同署名者的完全责任。

\subsubsection*{a4.~先期起棺}
\tocsubsub{a4.~先期起棺}

\notehead{168(iv)}{}

(4月8日)最近有人向我提示了Deligne1979年发表的文章“L函数的值与积分周期”(载于Proceedings of Symposia in Pure Mathematics, Vol.~33 (1979), part~2, pp.~313--346),与前文提及的R.P.~Langlands文章“自守表示、Shimura簇与Motif. 一个Corvallis童话”(pp.~205--246)刊于同一卷中。后一篇文章(而非Deligne的文章)出现在Deligne去年8月寄给我的Motif评注书目中,我曾一直以为正是Langlands这篇文章构成了在我离开后、1982年起棺之前文献中唯一谈及Motif之处(除了前文脚注中提及的Saavedra与Kleiman的报告)。

事实上,在Deligne引用的文章中,赫然有一章“第0章”,题为“Motif”,其开篇为:“我们回顾\textbf{归功于Grothendieck的Motif形式主义的一部分}”(着重号为我所加)。其阐述方式清晰表明,我为一个域上的(半单)Motif范畴给出的构造通则是多价的——文中接着指出“Grothendieck关于Motif的\textbf{定义之一}是通过……获得的”。在这方面该文是诚实的。确实,这里阐述的Motif“yoga”部分是最基础的部分,几乎在文献中早已存在(在Manin、Demazure、Kleiman、Saavedra的报告中),因此我的首创权\dmark{}\origpage{839}格外无可争议。(相反,在权的yoga以及后来的Motif Galois群中对我首创权的抹杀——以及对Serre首创权的抹杀——似乎未遭遇任何阻碍……)

正如我已经强调过的(在注“节节加码(2)”,n\textsuperscript{o}~174中),在1977年“下葬行动”达到暂时高潮(伴随“SGA~4$\frac{1}{2}$—SGA~5行动”)之后,出现了一段相对平静的时期,直到1981年反常层会议的“封神”,该会议标志着在对遗骸的瓜分中放弃了任何克制意图。(见注“封神”,n\textsuperscript{o}~171.)Deligne的这篇文章显然处于这一平静期的标志之下。我推测Langlands对Motif yoga的兴趣迫使Deligne最终泄露了Motif这一(早已公开的)秘密,而在当时的心理背景下,将死者的名字纯粹抹杀的时机尚未成熟。在随后的短短三年里,在这一对Motif的“怯懦起棺”与1982年LN~900卷册“大张旗鼓的起棺”之间,确实发生了一场令人惊叹的“节节加码”。

(4月22日)前一页所评述的(微小)发现在随后的日子里继续并大大扩充。我确实阅读了引述的R.P.~Langlands文章,并且最重要的是在次日读到了钉入我棺木的“第六颗钉子”\fnote{414}{(*)}{这是按发现顺序排列的第六颗“钉子”,但按制造年代算则是六颗中的第一颗,是由我朋友皮埃尔熟练地“打下”的,使用的是知名殡仪馆Springer Verlag GmbH(殡仪部“Lecture Notes in Mathematics”)提供的专利设备。},以我的(前学生)Neantro~Saavedra~Rivano所著、题为《塔纳卡范畴》的卷册形式出现。因此关于“Motifs行动”还有一段实质性的“后话”,我在题为“第六颗钉子(钉进棺材)”的一系列子注(n\textsuperscript{o}s~176$_{1}$至176$_{7}$)中予以展开。将这一后话移至调查“四种运算”的末尾在我看来更为妥当,因为在整场调查中浮现的新事实,尤其是在注“封神”(n\textsuperscript{o}~171)及其四条子注\fnote{415}{(**)}{(5月11日)自写下这些文字以来,引述的注已拆分为四条注,并增补了另外八条子注(n\textsuperscript{o}s~171(v)至(xii))。}中浮现的事实,对于准确安放这一“后话”并赋予其全部意义是必不可少的。

\subsubsection*{(b) 操演(“平展上同调”)}
\tocsubsub{(b) 操演(“平展上同调”)}

b1.~围绕“Weil猜想”的背景

\notehead{169(i)}{}

\dmark{}(2月27日)我现在来谈第二场“大运算”:\origpage{840}

\textbf{II~“平展上同调”行动。}

如同对待Motif一样,先用几句话说明背景将是有益的。

关于在任意域上的代数簇存在一种“上同调”理论——该理论将为这种簇(至少当它是射影且光滑时)赋予“上同调空间”,其系数域为特征零(例如$p$-adic域),其性质仿照复基域情形下(经由超越途径定义的)经典“Betti”上同调的已知性质——这一思想“在字里行间”存在于Weil著名的猜想陈述中(1949年)。无论如何,大约在1955年,Serre就是用上同调术语向我解释Weil猜想的——而且唯有在这些术语下,它们才真正具有吸引力。

当时无人知晓如何定义这样一种上同调,我甚至不确定除Serre和我之外是否还有其他人(如果这能算得上的话,连Weil本人也未必)具有其必然存在的深切信念。人们当时唯一能直接几何把握的只有$H^1$,通过阿贝尔簇理论及其有限阶点(由Weil发展),以及通过关联于非奇异射影代数簇的Albanese簇或Picard簇。对$H^1$的这一构造提示“自然”系数应当是$\ell$-adic域$\mathbb{Q}_\ell$,其中素数~$\ell$~\textbf{不同于}基域的特征。对于~$\ell$~等于特征的情形(当特征非零时),Serre在代数\textbf{曲线}情形下的初步结果提示人们应当可以取$k$(设为完美域)的Witt向量环的分式域作为系数。因此人们有望对\textbf{每一个}素数~$\ell$~拥有一套$\ell$-adic理论(对于$\ell=p$需打点折扣)——并且在合理意义下,它们应当“都给出相同的结果”。最后,当$k$为特征零时,鉴于人们拥有Hodge上同调空间(自Serre引入代数簇的“凝聚”上同调理论以来对任意$k$都有意义)与De Rham上同调空间(我受微分De Rham上同调启发引入),这些在当时提供了具有全部所期\dmark{}性质的上同调理论\origpage{841}\fnote{416}{(*)}{我早在五十年代就发展了结合于代数闭链的上同调类(De Rham或Hodge)形式主义。},并且它们应当同$\ell$-adic理论一样给出“相同的结果”。

这些问题处于我1955年至1970年(我离开数学舞台之年)间发表与未发表数学思考的核心。如果撇开我在凝聚上同调方面的工作不谈(“六运算”形式主义、Riemann-Roch-Grothendieck公式),大体可以说,我上同调工作的核心就在于给出这些问题的答案或大纲。至少在以Weil猜想为主要灵感源泉的视角下,我对上同调主题的思考具体化为四条大\textbf{潮流}或“\textbf{线索}”,紧密交织构成一张宏大织锦:

\textbf{线索~1.}——我(在合作者的协助下\fnote{417}{(**)}{发展平展上同调形式主义的主要合作者是M.~Artin与J.-L.~Verdier。$\ell$-adic系数则是在我的前学生P.~Jouanolou的论文中发展的(可惜他从未发表,我从未见过,如今已无从寻觅)。我打算在随附于《数学反思》的“历史”评注中给出关于平展上同调发展的全面描述。})发展了概形的\textbf{平展上同调}与\textbf{$\ell$-adic上同调}形式主义,对于与残余特征互素的~$\ell$,它具备拓扑空间熟悉的“离散”上同调的全部性质(乃至更多……)。除了三项技术性质的未决\fnote{418}{(***)}{这三项“未决问题”如下:
\textbf{a.}~对于特征与~$\ell$~互素的正则概形~$X$~的正则闭子概形~$Y$,“上同调纯性猜想”(平展版本)。当~$X$~与~$Y$~在正则基概形~$S$~上均光滑时(对于应用足够的情形),以及在特征零情形下(由Artin充分利用奇点消解证明),相关猜想已获证。
\textbf{b.}~更为严重的是对于诺特概形(需要时为优良概形)之态射~$f:X\to Y$,当~$f$~\textbf{非假定}为真时,关于更高直接像$R^i f_*$在$\ell$-adic“可构”框架下的\textbf{有限性定理}之有效性问题。我在基概形为代数闭域上的簇且在奇点消解与“上同调纯性”(参见a)假设下证明了该结果。然而我指出,在有限系数框架下(相对于$\ell$-adic系数而言),六运算的对偶形式主义(包括Poincaré对偶)在1963年由我在无需有限性条件下建立。
\textbf{c.}~优良正则概形上“双对偶定理”的有效性。类似情况。
Deligne(约在1973年?)对维数$\le 1$的正则概形~$S$~上的有限型概形态射给出的优雅证明大大改善了局面。在同一情形下,通过类似论证,Deligne也证明了双对偶定理。}问题之外,可以说自1963年起人们“原则上”就已拥有了\dmark{}\origpage{842}(伴随1963/64年SGA~4之后的SGA~5讨论班的展开)该上同调在平展上同调框架下的\textbf{完全构造}——以“六运算”对偶形式主义的形式。平展上同调的定义原则可追溯至1958年,我在1963年2月至3月证明了完整形式主义(包括“弱Lefschetz”型定理和平展语境下的上同调深度概念)所必需且充分的“关键结果”。

\textbf{线索~2.}——借助\textbf{Motif}的yoga,我发现了使人们能够将一个簇的各种$\ell$-adic(及其他)上同调视为该簇所共有的“Motif上同调”的不同“实现”的\textbf{哲学}。这一哲学在六十年代初萌生,伴随着直接受Weil猜想启发的“权的yoga”(以及Serre受其启发的关于簇的“虚Betti数”想法\fnote{419}{(*)}{关于这一点,见“我的孤儿”(n\textsuperscript{o}~46)的子注~n\textsuperscript{o}~46$^9$。})。它在1964年平展上同调启动的势头中丰富起来,引入了“Motif Galois群”的关键概念。

\textbf{线索~3.}——受Monsky-Washnitzer思想的启发(他们为特征~$p>0$~下的\textbf{光滑仿射}代数簇构建了常系数“$p$-adic”上同调理论),我在1968年给出了“$p$-adic上同调”的通用定义,我亦称之为\textbf{结晶上同调}\fnote{420}{(**)}{这一术语如今(且久已)由习惯确立,这一理论亦然。引导我走向这一理论的两项新思想(相对于Monsky与Washnitzer的思想而言)是:\textbf{晶体}(模等)的概念,与在“加厚”(初始无穷小加厚)之上的“生长”思想相联系;以及在所考虑的理想中引入\textbf{除幂}结构,以确保除幂“形式Poincaré引理”成立。奇怪的是,晶体的关键直觉被我的学生们弃之不顾,正如贯穿我全部上同调思考的线索——六运算形式主义一样。这在我看来是导致结晶上同调与Hodge-Deligne理论令人遗憾地陷入停滞的主要原因。此外在我看来显而易见的是,在这两种理论中,Zoghman~Mebkhout所发展的理论将扮演不可或缺的角色。我在这一方向上的建议(1978年致Berthelot)显然被置若罔闻……}。这一\dmark{}理论旨在涵盖非常数亦非局部常数的“系数”,并产生如同$\ell$-adic理论一样的“六运算”形式主义。从一开始就确立的是,对于\textbf{光滑}簇,这一上同调具有人们所期望的与De Rham上同调的联系,并且推广了Monsky-Washnitzer理论\fnote{421}{(*)}{P.~Berthelot以我的思想为起点的博士论文为光滑射影簇给出了其证明与对偶形式主义,至少足以给出有限域上此类簇的L函数的\textbf{结晶}上同调解释。然而距离掌握一套用于“结晶系数”的“六运算”形式主义,至今依然遥远。对于代数簇上的“Hodge系数”亦然!关于这一系数问题,见“墓上的旋律——或自负”(n\textsuperscript{o}~167)。该问题在六十年代对我而言十分明确,此后却被我的上同调弟子们埋葬至今。(4月23日)关于这一点见注“工地巡礼——或工具与视野”,n\textsuperscript{o}~170(i)。}\origpage{843}。

\dmark{}\textbf{线索~4.}——将\origpage{844}平展上同调及其近邻变体(涉及Zariski、fpqc、fppf等拓扑)、结晶上同调、超越语境下的“Betti”上同调以及更一般的任意拓扑空间之层上同调通过共同的“拓扑”直觉联系起来的统一几何概念,是“\textbf{site}”的概念,以及在其背后更内在、更隐秘的“\textbf{topos}”概念。自1964年起在随后的岁月里,这一概念逐渐走上前台。我在“我的孤儿”(n\textsuperscript{o}~46,第636--638页)中阐述了这一处于我工作核心、如今却被逐出几何学的概念的深远意义,在此仅引述如下段落:

\begin{quote}
“这一对概念[概形与topos]潜藏着代数几何与算术、以及拓扑学的大规模更新,通过将这些隔绝太久的‘世界’综合进一个共同的几何直觉之中。”\fnote{422}{(*)}{我在别处提议称这种尚在萌芽中的‘新科学’为\textbf{算术几何},它诞生于六十年代初Weil猜想的余波中,其中‘Motif的yoga’如同灵魂。我在此重申:‘在构建新代数几何之外,通过掌握平展拓扑,为一门仍在演进中的科学构建整体主控大纲,在我眼中是我对同时代数学的主要贡献。’}
\end{quote}

topos的语言与平展上同调形式主义在紧密相连的两个讨论班SGA~4(1963/64年)与SGA~5(1965/66年)中展开\fnote{423}{(**)}{SGA~4全面重修的第二版(三卷)于1972--73年发表在Lecture Notes n\textsuperscript{o}s~269, 270, 305中。关于SGA~5的波折,其遭受严重肢解的版本在口头讨论班结束后十一年的1977年发表于同套丛书第569卷。}。第一部\dmark{}与他人合作完成\fnote{424}{(*)}{site与topos语言的发展基于我1958年的最初想法,在M.~Artin、J.~Giraud、J.-L.~Verdier协助下完成。},除了topos语言外,发展了平展上同调的初步知识,包括对偶的关键起步陈述(六运算风格)。第二部中我几乎独力作战\fnote{425}{(**)}{唯一的例外是J.-P.~Serre作了若干关于有限群在代数簇上的作用以及与Artin导子相关联的Serre-Swan模的报告,这是我发展通用不动点公式所需的。对于所有其他报告,我是唯一的报告人。},极详尽地展开了\origpage{845}对偶形式主义,包括通向L函数上同调理论的不动点公式(构成Weil猜想的重要部分)。我在“遗骸……”那条注(n\textsuperscript{o}~88)中评价道:

\begin{quote}
“接续进行的两个讨论班SGA~4与SGA~5(于我犹如同一个讨论班)从虚无中发展了topos语言所代表的综合与发现的强有力工具,以及平展上同调这一极其精纯完备的工具——当时对其本质形式性质的理解甚至超过了普通空间的上同调理论。这构成我在数学中所做出的最深刻、最创新的贡献。这两部讨论班不可分割地联系在一起。它们在统一性中同时代表了\textbf{视野}与\textbf{工具}——topos,与上同调的完整形式主义。虽然视野在今天仍遭拒斥,但该工具二十多年来深刻更新了代数几何中最迷人的‘算术’面向。”
\end{quote}

\asterisms

\subsubsection*{b2.~四种操演}
\tocsubsub{b2.~四种操演}

\notehead{169(ii)}{}

\dmark{}\textbf{“平展上同调”行动}在于一方面将SGA~4与SGA~5中完成的奠基工作\textbf{贬损}\origpage{846}(贬为“无稽之谈”、瞎拼凑等),并顺势抹杀我在发现和发展这一工具中的角色;另一方面,\textbf{据为己有这一工具},即据为己有我在平展上同调中所发展的思想、技术与结果的发明权。这里的“受益者”又是Deligne\fnote{426}{(*)}{然而Verdier也获得了丰厚的回报,正如后文将看到的,他通过发表其“难忘文章”给肢解SGA~5发出了“开门红”,随后在1981年的反常层会议上大获其利。},正是其非凡的声望(既得益于其卓越的才能,也得益于其作为我默认“继承人”的地位)使得这样一场大规模的贬损与侵吞行动得以推行,且表面上不着一丝痕迹……

正是在1965/66年的SGA~5口头讨论班上,年轻的新来者Deligne第一次学到了概形语言、Grothendieck风格的同调代数以及平展上同调的新技术\fnote{427}{(**)}{这正是我在注“格格不入者”(n\textsuperscript{o}~67$'$)中回忆起的。Deligne正是在SGA~5中顺便学到了Poincaré对偶的‘六运算’形式主义。}——这些技术成为了他全部后续工作的基石。

在我这位前学生操纵的行动中(我称之为“SGA~4$\frac{1}{2}$—SGA~5行动”),我辨识出四种不可分割的“\textbf{操演}”:

\textbf{操演~1. 贬低}母讨论班\dmark{}SGA~4—SGA~5,斥之为“废话矿渣”及类似言辞:这在Deligne署名、名为“SGA~4$\frac{1}{2}$”(副题:平展上同调,载于Lecture Notes in Mathematics第569卷)的卷册前言中迂回且看似不经意地完成。关于Deligne学艺并获取其后续工作基础工具的双讨论班SGA~4—SGA~5所遭受的清算细节,见“白板”(n\textsuperscript{o}~67)。\origpage{847}

\textbf{操演~2. 破坏}我对SGA~5口头报告的整体撰写工作\fnote{428}{(*)}{正如我三条脚注前所指明的,对于SGA~5的全部报告,我都有详尽的手写笔记。只要分配下去,记录整理本该在两三年内完成。可现实是,这项工作被完全冷冻了十一年之久。}。通常这一整理本应在合理期限内(至多一两年)由我那些上同调学生完成,他们有幸学习了大部分技艺,并同讨论班其他听众一起成为了多年间唯一知晓这些思想技术的人。这也是他们熟悉这些实质内容与思想技术的最快方式,而在口头报告时这些内容往往让他们觉得有点“高深莫测”(当然总是敏捷过人的Deligne除外)。然而这一整理——或者更确切地说是\textbf{拒绝整理}——硬生生\textbf{拖延了十一年}——直到Deligne“亮绿灯”让Illusie最终处理这本不幸的SGA~5的撰写与出版的精确时刻(出于巧合),也就是确信它将在1977年发表在他自己手笔的一卷书\textbf{之后}的时刻!那卷书(写于1973年及随后几年)起初是为了普及平展上同调中证明Weil猜想(最后一环)所必需的“成分”,却在此机缘下被冠以奇特的名称“SGA~4$\frac{1}{2}$”。(这名字直到今天似乎除了我之外\dmark{}从未让任何人感到错愕、惊讶、乃至震惊……(169$_1$)\fnote{429}{(*)}{关于这一点,以及关于SGA这一缩写的最初真实含义,见“白板”(n\textsuperscript{o}~67)。})。\origpage{848}详情见注“绿灯”与“颠倒”(n\textsuperscript{o}s~68, 68$'$),在那里题为“SGA~4$\frac{1}{2}$”的卷册的含义开始向我显露,亦见注“沉默”与“团结”(n\textsuperscript{o}s~84, 85)。

\textbf{操演~3. 肢解原始讨论班}SGA~5,其出版版本(由我的前学生Luc~Illusie“操持”)已不过是一具被严重残害的“遗骸”。我在同名注“屠杀”(n\textsuperscript{o}~87)——关于葬礼反思中最长且最富启示性的注之一——中详尽检视了这场肆无忌惮的肢解,或者更确切地说,对托付给学生们的光辉讨论班的\textbf{屠杀}。

\textbf{操演~4. 击碎我关于平展上同调工作的统一性},这一工作由SGA~4与SGA~5两扇不可分割的门扉所体现,通过将其“切成两段”,“通过在两扇门扉之间强行插入一篇外来的、蔑视性的文本”\fnote{430}{(**)}{这句引文凭记忆引自注“遗骸……”(n\textsuperscript{o}~88)。},它顶着非同寻常的名字“SGA~4$\frac{1}{2}$”\fnote{431}{(***)}{副题:平展上同调——Pierre~Deligne著……这副题也同样言简意赅!}。\origpage{849}这个天才的名字言简意赅——关键在于能想得出来!单凭这一个名字,该卷册便自诩为平展上同调的\textbf{核心与根本文本},旨在\textbf{取代}“SGA~4与SGA~5冗长繁琐的报告”,后者“可被视为一系列离题赘述”,虽然“某些非常有趣”,但核心文本“应当允许使用者将其遗忘”。

此外,我这位天才的前学生兼朋友无需在此陷入冗长无用的辩解:单单这个简峻的名称“SGA~4$\frac{1}{2}$”\dmark{}便宣告并确立了这一文本相对于题为SGA~5的“离题赘述”而言具有无可辩驳的\textbf{在先性}(后者作为其必然结果,确实是在它之后出版的……),并同时确立了SGA~5相对于这一“在先”文本的(所谓的)\textbf{逻辑依赖性}。

关于SGA~5相对于这一伪作具有所谓的逻辑依赖性的离奇骗局,在该书的前言中被确凿宣称\fnote{432}{(*)}{我提醒读者,Deligne在上个月访问我处时,曾当面向我确认过这一点。},作者毫不脸红地宣布(且鉴于当今风气,在我之前显然无人觉得有何不妥……):“它的存在[‘SGA~4$\frac{1}{2}$’的存在]将使得不久后得以出版SGA~5\textbf{原貌}”(着重号为我所加)——请读作:处于被洗劫和大量盗窃后的\textbf{遗骸}状态……虽然我此前一周多已知晓我朋友的“Motifs”行动,但我仍花了整整两天时间(从4月26日的“白板”注到28日的“颠倒”注,n\textsuperscript{o}s~67, 68$'$)才领悟我这位高足这一显然荒唐的断言对我所代表的“谜团”的意义——并由此理解了看似平淡无奇的符号“SGA~4$\frac{1}{2}$”的含义,此前两天我甚至未曾在此停留。

相同的“逻辑依赖”骗局在Illusie为SGA~5所写的前言中被明确提示(169$_2$)\fnote{433}{(**)}{详情见本注的子注“好撒玛利亚人”(n\textsuperscript{o}~169$_2$)。}。对于不知情的读者而言,由于我的报告的迟到整理者们\fnote{434}{(***)}{(4月9日)经细致核实,所涉及的‘迟到整理者’主要是Illusie一人……}乐于在整理稿中塞满对“SGA~4$\frac{1}{2}$”的无数指涉,这一骗局显得更加可信\origpage{850}。这些指涉许多并非虚构,而是关联于原始讨论班的两篇报告(一篇由Illusie整理,另一篇——格外关键的——由Deligne整理\fnote{435}{(*)}{所涉及的是报告‘结合于闭链的上同调类,A.~Grothendieck著,P.~Deligne整理’。}),它们被不由分说地直接并入了题为“SGA~4$\frac{1}{2}$”的卷册中——既未通知我,亦未征求我的意见,而是仿佛在逝去的师父缺席下这本就理所当然属于他们一样……

\dmark{}这一强盗行径进一步使得我的前学生Deligne实现了\textbf{角色的绝妙倒错},得以在书本封面上(同样绝不与我商量……)将我呈现为他的\textbf{协作者}(为了平展上同调的发展!)\origpage{851}\fnote{436}{(*)}{这一幕安排(将我呈现为我学生Deligne的‘协作者’)格外滑稽,考虑到在整本书中Deligne独创的材料至多占三十到四十页……}——一个略显“糊涂”的协作者\fnote{437}{(**)}{在他寄给我的一份‘SGA~4$\frac{1}{2}$’提要中,我的名字甚至未被提及。},这倒是真的,但无论如何是个“协作者”……

至于那本被称为“SGA~4$\frac{1}{2}$”的盗版文本,除了两篇从原始SGA~5语境中剥离出来的报告,以及若干对算术应用格外重要的SGA~4—SGA~5结果的“摘要”,外加一章关于三角和应用的原创章节,以及最后Verdier“论文”的“第0状态”之外,它由对SGA~4—SGA~5中所发展的上同调形式主义的若干补充(固然十分有用\fnote{438}{(***)}{所涉及的是有限性结果(已在三条脚注前提及)以及关于具有有限系数的局部对偶定理。})组成。这本可以写成一篇约三十页(若算上‘三角和’一章则约五十页)内容略显杂糅的出色文章。如果我这位辉煌的前学生尚存一丝体面,将这几项补充各自放置在其所受启发并加以补充的两三篇SGA~5报告中\dmark{},本是不言而喻的\origpage{852}。然而相反,它们成了纯粹取消SGA~5第II讲的借口(得到了负责整理并以此‘弥补’的Illusie的祝福,将该讲转化为‘SGA~4$\frac{1}{2}$’中的一个附录),并顺理成章地将平展上同调的双对偶定理(由我在1963年根据五十年代发现的‘凝聚’类似物建立)重新命名为“Deligne定理”\fnote{439}{(*)}{\textbf{双对偶定理},或‘局部对偶定理’(两个名称均出于我手),在SGA~4$\frac{1}{2}$中被冠以Deligne之名,四年后Deligne又慷慨地将其转赠给他的朋友Verdier……}。

\subsubsection*{b3.~节节加码的各幕}
\tocsubsub{b3.~节节加码的各幕}

\notehead{169(iii)}{}

\dmark{}\origpage{853}(3月12日)“平展上同调”行动在从1966年到1977年的整整十一年间持续推进,介于SGA~5讨论班的结束与电锯卷册“SGA~4$\frac{1}{2}$”及随后的SGA~5屠杀版(所谓的“Illusie版”)\fnote{440}{(*)}{(3月12日)现在看来,将‘平展上同调’行动视为在1977年便告终止是不准确的。它在1981年反常层会议上达到了封神。}接踵出版之间。它的完成首先归功于我的五位“上同调”学生的共同参与,无论通过作为还是不作为:\textbf{P.~Deligne、L.~Illusie、J.-L.~Verdier、J.P.~Jouanolou、P.~Berthelot}\fnote{441}{(**)}{这种团结一致对这五位前学生中的每一个人来说,首先表现为在1970年我离开之后,在长达七年间对完成他们受托整理的SGA~5报告的彻底拖延。}。除Deligne之外,Illusie的责任在我看来最为沉重,因为正是他承担了屠杀版编辑的责任,从而充当了Deligne的顺从工具\dmark{}\origpage{854}\fnote{442}{(*)}{Illusie同样充当了\textbf{Verdier的同谋},为后者在导出范畴上的欺骗行为提供掩护。}。

Deligne侵吞平展上同调“真正”发明权的意图无可怀疑。整场“平展上同调”行动的精神为此提供了明证,这在我们科学的编年史中无疑是独一无二的。它同样在1975年Deligne的传记小记中被低调表达(其中对我交到他手中的上同调工具——它在其证明Weil猜想最后一环中起到了决定性作用\fnote{443}{(**)}{(3月12日)在这篇文本中,如同在Deligne的任何其他出版物中一样,绝无任何字句暗示他曾从我这里学到任何上同调工具……}——只字未提),并在八年后,在我称之为“葬礼悼词”(三篇)的这组简短而雄辩的文本(1983年)中被公然表达出来\fnote{444}{(***)}{在我关于葬礼的反思中,与葬礼悼词的相遇构成了最决定性的转折点之一。}。它们在“葬礼悼词(1)——或恭维”与“葬礼悼词(2)——或力量与光环”(n\textsuperscript{o}s~104, 105)这两条注中得到了其应得的审视(并在后来的注“阴的葬礼(阳葬阴(4))”,n\textsuperscript{o}~124中以更深刻的光照加以重温)。至于Deligne的自传性(且决非葬礼性的)“赞词”,我在“为一副模糊的骸骨作安魂曲”和“信仰表白——或伪中之真”(n\textsuperscript{o}s~165, 166)两条注中对其作了剖析\fnote{445}{(****)}{关于这一自传小记的详情,亦见前引注n\textsuperscript{o}~164末尾(第V~2部分)。}。

该行动在1977年达到顶点\fnote{446}{(*)}{(3月12日)这是一场纯粹暂时的‘登顶’!见本注‘操演’中今天所写的各条脚注(第853页)。},随着“SGA~4$\frac{1}{2}$(sic)—SGA~5”的相继出版\dmark{}\origpage{855}。这是一场长达十一年的针对我的著作与我个人的下葬中漫长\textbf{节节加码}的(暂时)终点,其每一步骤都因在学术界找到的默许鼓励而变得更加肆无忌惮。

——在前几个阶段里,面对其可疑本性,人们不是漠不关心就是无动于衷(甚或趋之若鹜……)。其中若干阶段我在上文已有提及,即方才回顾过的“Motifs”行动。我还注意到另外三段插曲,它们与“平展上同调”行动有更直接的关联,现在轮到我把它们检视一遍。

\textbf{插曲1.} 它关乎某一条“离散Riemann-Roch”型猜想的遭遇\dmark{}\origpage{856}。我于1966年在SGA~5口头讨论班上引入它,在最后一次报告里,我提炼并评述了若干未决问题与前所未有的猜想。这份报告在Illusie版里尸骨无存——文中对这条猜想绝口不提(并非无缘无故……),对当时提出的众多其他问题亦然。然而,讨论班七年后,这条猜想以解析语境的面目在MacPherson笔下重新现身,绝口不提任何SGA~5讨论班(或概形语境),并顶着“Deligne--Grothendieck猜想”这一不寻常的名字。所涉的是MacPherson那篇广为人知的文章\fnote{447}{(**)}{即R.~MacPherson《奇异代数簇的Chern类》,载Annals of Math.~(2)~100~(1974), pp.~423--432。他在该文中于解析语境证明了这条猜想。},他在其中于解析语境证明了这条猜想。在他去年10月来访时,Deligne向我指明:他在1972年仅限于把我的猜想\textbf{原封不动地}传达给MacPherson(那是他随SGA~5其他听众一起在口头讨论班上学到的)。他说他对MacPherson所起的名字感到惊讶,却并未因此劳神给他写信请他更正。关于这一点,见“把一个个~i~点清”那条注(n\textsuperscript{o}~164,第II~1部分);关于这条猜想本身更详细的说明,见“屠杀”那条注(n\textsuperscript{o}~87)的长子注~n\textsuperscript{o}~87$^1$\fnote{448}{(*)}{这条猜想将首次以其原初与完备的形式,出现在SGA~5的刊行本中(第VIII讲)——十一年后!而在题名之下,我的名字……}。

\textbf{插曲2.} 它关乎SGA~7讨论班的坎坷命运。该班致力于\textbf{平展上同调中的单值群}问题,1967至1969年间,在Deligne与我共同倡议和主持下进行。讨论班的起始想法与总体设计出自我手;Deligne在其中亦有数项贡献,最重要的一项是他在平展语境中对Picard-Lefschetz公式的证明。一如SGA~5,报告的整理工作也拖了数年——那恰是其不幸前任那场(非)整理之剧本的(开头部分的)重演!出版终究还是于1972与1973年在Lecture~Notes~n\textsuperscript{o}s~288, 340中实现了,由Deligne操持——其时我离开数学舞台已三年。经他倡议,讨论班被\textbf{硬切成两部分}:第一部署为我的名义指导,第二部署为他和N.~Katz(后者在讨论班第二年不过是一名普通的报告人)的共同指导\fnote{449}{(**)}{关于我在这场没有任何数学理由支持的\textbf{切断}中所看到的意义,见下文紧随其后的注(n\textsuperscript{o}~56),以及本注“操演”的子注“驱逐(2)”(n\textsuperscript{o}~169$_1$)。}。

在以我名义刊行的第一卷SGA~7~I中,我在讨论班开篇几场报告中系统阐述的消没闭链理论,被“砍”成Deligne一份二十页的摘要(其余报告由我本人及讨论班其他参加者在合理期限内完成整理)。至于以Deligne-Katz联合署名刊行的第二卷——我在其中对主要主题与结果之发展所起的作用,不减于第一卷——这一份贡献,却被系统性地抹杀了\dmark{}\origpage{857}。关于这方面详情,见“一场屠杀的序曲”那条注(我在其中试图厘清SGA~7这场微型行动的意义),尤其是“把一个个~i~点清”(第II~5部分),n\textsuperscript{o}s~56, 164。

这里我只限于提一提最大的一笔抹杀。它关乎我在平展上同调语境下,对“Lefschetz束”上同调理论及“不可约性定理”所作的转写。这批经典结果(凡是真被证明了的……)经由超越途径获得,这一转写(一如常有的情形)绝非机械照搬。我记得为此花了好几天乃至整整一星期。据我所知,至今没有别的已知证明可以处理其中的主要事实——除了我当年借助于谱序列、以及我于1958年确定的代数曲线“驯顺”基本群的“著名”结构所提炼出的那一个\fnote{450}{(*)}{在接下来将被引用的Katz报告的引言里,他似乎还给人以Raynaud夫人是那条定理作者的印象!}。这一理论被复录于SGA~7~II里Katz的一份报告(第XVIII讲)中——依据的是我提供给他的笔记。在卷册引言里,Lefschetz束理论(连同Deligne证明的Picard-Lefschetz公式)被呈为讨论班的“关键结果”之一,但对我在该卷所展开的任何主题中扮演的角色,却无一字提及。据我所知,文献中唯一一处哪怕暗示了我这层角色的,是Lefschetz束报告题名(以及作者名)之后一条简短而暧昧的脚注\fnote{451}{(**)}{这条脚注含混其辞,小心避免确认发明权——而当时的圈内人都心知肚明。}:“据Grothendieck的(简略)笔记”。

在Deligne同年(1973)发表于《数学出版物》的文章“Weil~I”中(169$_4$)\fnote{452}{(***)}{见子注“‘那条’猜想”(n\textsuperscript{o}~169$_4$),它源自此处的一条脚注。},Lefschetz束理论作为其Weil猜想证明中的一件重要技术成分出场。在这篇文章里,Deligne还没有胆量抹杀我在$\ell$-adic迹公式中的角色(那是他证明的另一件关键成分,其在知情圈内的发明权毕竟还太昭然\fnote{453}{(*)}{然而次年,在他那篇自传小记里(已在注n\textsuperscript{o}s~165, 166中检视),对上同调工具的任何提及都消失得无影无踪……});可是,当他郑重其事地表述他即将使用的Lefschetz理论的种种结果时,对我这个人却无一字提及。他仅限于指向SGA~7的相应报告——而一位不幸的读者,多半永远无从在那里觅得他朋友Katz那处难以捉摸的脚注……

\textbf{插曲3.} 我所知的“节节加码”最后一段插曲,发生在1976年——比“SGA~4$\frac{1}{2}$”行动的“高潮”早一年。它关乎Verdier那篇文章的命运\dmark{}\origpage{858}。

在这篇文章里——我曾把它称作“好的文献”——Verdier在解析空间上、对解析可构离散系数,复述了不差分毫地\textbf{我在1963年、尤其1965/66年在SGA~5的概形平展语境下}所发展出来的对偶形式主义。他处处只字不提它的来源——只除了以他特有的一类非常特别的玩笑方式……\fnote{454}{(**)}{将概形(或“空间”……)的\textbf{同调}定义为其在某“对偶化复形”中取值的超上同调——这一想法由我于1963年引入。}——并且对作者或许有所耳闻的任何SGA~5讨论班,也绝口不提。详情见“好的文献们”与“玩笑——或‘权复形’”两条注(永远是同样的“权”,错不了……)n\textsuperscript{o}s~82, 83。

正是从这份“难忘文章”开始,复解析空间上、对解析可构离散系数的对偶形式主义——它逐字逐句复录了我早在1963年、尤其1965/66年在SGA~5的概形平展语境下所发展的那套——便神不知鬼不觉地变成了“Verdier对偶”;人们尚须等待五年之后(在1981年6月吕米尼会议的亢奋中),同一套腾挪才也降临到平展对偶头上。但我在此提前说了(一如我方才对“难忘文章”那段插曲本身已提前的那样)那\textbf{第三}场大行动——这一回,Verdier是其主要的(若非唯一的)“受益者”——它将在下文谈到\fnote{455}{(*)}{见注“分赃”,n\textsuperscript{o}s~170(i)-(iii)。}。

\subsubsection*{b4.~放肆}
\tocsubsub{b4.~放肆}

\notehead{169(iv)}{}

\dmark{}\origpage{859}Verdier的这篇文章,于我如同投来一道意外的光照,照亮了SGA~5在我昔日某些学生手中所遭的命运。它让我看清:这些人从对我为他们最先发展于SGA~5的思想与技术的独占性知悉中,究竟能榨出何种“利益”。它也无可置疑地向我显明:我全体上同调学生与这类行动之间的同谋与团结。我把这篇文章称作“好的文献”时,不曾想到这名字取得如此贴切——它确实已经(如我从多方得到证实的那样)变成了一份标准的参考文献\fnote{456}{(*)}{至于\textbf{上同调}变体(在Verdier的文章中仅被一笔带过——他深知Deligne要把这一份留给自己……),见下一节。},他们当中谁也不可能不曾注意到它\dmark{}\origpage{860}。这正是我最终在注“沉默”与“团结”(n\textsuperscript{o}s~84, 85)中领悟到的东西。我知道我不必惊讶:在那部曾是SGA~5讨论班的Illusie版里,从头到尾对我的同调(及结合于闭链的同调类)形式主义毫无一字提及——确实也无须提,因为(十年之后)他的好伙计Verdier已经替代大家心满意足地提供了那缺失的文献……

Verdier提供的“好文献”,正如Deligne用以对Motif进行局部起棺的“难忘卷册”一样,在我看来纯粹就是彻头彻尾的剽窃。被称为“SGA~4$\frac{1}{2}$”的文本尚且不完全是这般模样\fnote{457}{(**)}{(3月21日)在合题为“公式”的那一系列子注(n\textsuperscript{o}s~169$_5$至169$_8$)里展开的反思向我表明,尽管某些“形式”仍被保留,但这种印象是错误的……}。在那里,某些形式仍被守着,遵循着那套必不可少的“暂停!”风格——它擅长于不断暗示虚假,却从不(或几乎从不……(169$_3$)\fnote{458}{(***)}{关于这一点见子注“特洛伊木马”(n\textsuperscript{o}~169$_3$),它源自此处本来用以阐明这“或几乎……”的一条脚注。})贸然把它明白\dmark{}\origpage{861}说出来。我与“SGA~4$\frac{1}{2}$”以及该风格在其中所取的特定形式(即轻蔑贬低的风格\fnote{459}{(*)}{这种“贬低”假装对一位“糊涂”(“虽说严密”……)且尽情杂乱的前辈所积累的“废话矿渣”一扫而空……})的首次交锋,是在“白板”那条注(n\textsuperscript{o}~67)里。

但这门行动尤其给我的触动,胜过一桩平庸的剽窃所能为的,在于它在放肆上所具有的某种维度。在我眼里,其余三场行动没有哪一场抵达了这一极端的维度\fnote{460}{(**)}{(3月11日)这种评价当然完全是主观的。顺带说,写下这行字时,我还有过一丝迟疑,想着那桩难以想象的反常层会议“行动”(即“行动IV”,下文还将谈到)。这场难忘的会议确确实实构成了借一个鲁莽的继承者(Zoghman~Mebkhout)的中介对我这个人下葬的一场真正的集体封神。正是在这一机缘下我意识到,这场封神同时是“平展上同调”行动的一个自然延续与最终结局,而“SGA~4$\frac{1}{2}$—SGA~5”那一幕,实际上不过是一场纯粹暂时的“登顶”。在后一场里,我的前学生Deligne仍免不了在此处彼处提及我这微不足道的人和我的著作,哪怕心不甘情不愿,并以此用傲慢的考语同其划清界限。而在1981年6月的吕米尼会议上相反,平展上同调虽处于公众注意的中心,我的名字(正如当值无名之辈Zoghman~Mebkhout的名字一样)却在任何时刻都不曾被吐露过……}。而且它对我的触动或许比其余三场中的任何一场都更强烈,因为它进一步像一桩暴力行为一样刺痛了我,像“为了好玩”而对一部我已完成且在其中倾注了整个自我的美好工作的一场屠杀——首要地,为了那些恰恰在事后乐于对它加以洗劫的人们而完成,好拿它充作他们自负的饲料,并且(顶着大雅之堂上名流君子的高尚门面)在其中张扬一种隐秘的狂妄以及那副自得轻蔑的神气\fnote{461}{(***)}{这种自负与轻蔑极其清晰地流露在这一卷名为“SGA~4$\frac{1}{2}$”(在我们这门科学的历史上在该类型中想必是独一无二的)的书本字里行间。它们也在这部卷册发表的当年(但调门更为隐蔽),出现在Pierre~Deligne对我的个人关系里。(见“两个转折”那条注,n\textsuperscript{o}~66.)我在那些曾是我学生的人之中这一个与那一个的怠慢轻浮中找到了它们——对我写给他谈及那些我挂心或使我伤心之事的信件置之不理。我轻盈而轻慢地,在对一项满怀爱意完成的工作的“Illusie版”(或屠杀版)的引言字里行间重新找到了它们,并在去年,在另一位学生的父式优越神气里(在“玩笑——或‘权复形’”那条注,n\textsuperscript{o}~83里谈到)重新找到了它们。}。

\subsubsection*{b5.~藏宝}
\tocsubsub{b5.~藏宝}

\notehead{169(v)}{}

\dmark{}\origpage{862}(2月28日)我方才检视的两场“行动”,连同第四场(所谓“反常层会议行动”)——下文还将谈及——皆有着共同的特征:在众多人的参与或纵容之下进行,而“受益者”(看来)仅此一人。这是这三场行动之间一个触目的共同点,它印证了在“掘墓人——或全体教众”(n\textsuperscript{o}~97)那条注里所继续的反思。

但我在这围绕Motif与平展上同调展开的前两场行动里,看到了一个更阴险的共同点,关乎激励着它们的某种精神。这里涉及的,是一种对待占有高水平、受限流通的科学信息的内心姿态;或者说到底,将信息局限在一个由利益同盟结合起来的少数人小圈子(乃至单一一个人)之内,只要把排他性的独占“收益”留给自己显得有利,便动用阻断流通的权力。

就这样,在我1970年“离开”之后,Deligne成了唯一一个(除我之外)深切同化了“Motif的yoga”并感受其全部深远影响的人——去作人们已知的那种使用。我的五位上同调学生(包括Deligne在内),或许还有两三位其他当年有毅力真正消化其精髓的SGA~5前听众,成了唯一独家支配我在这一讨论班中所发展的想法与技术的人。

在这两种情形里,无论在1965至1969年间与Deligne数不清的面对面长谈中,还是在1965/66年面对SGA~5听众的小圈子,若说确确实实是“首先为了他们、优于所有其他人”而在他们眼前长久地阐发并展开某种内在视景,我把这些于我珍贵的事物交到他们手中,却绝非作为某种“利益集团”的代表。于我,这是不言而喻的:我面对他们,就像面对同我一样的人——除了证明自己与为数学事物的共同认识作出贡献的自然愿望之外,还被一种服务精神所激励,面对着一个在空间与时间中皆无边界的“数学共同体”\fnote{462}{(*)}{关于这样一种“服务精神”,尤见下文同样引用的那条注“阴,作为仆人,与新主人”(n\textsuperscript{o}~135)。}。\dmark{}\origpage{863}而我交到他们手中的东西,我深知那绝非“稀罕物件”、绝非陈列品,而是活生生、灼热的事物,生来便要生长、便要播撒——而这也正是那些我所面对的人自始便感知到的\fnote{463}{(*)}{(4月10日)这并不妨碍他们之中的某些人在事后竭尽所能去贬损他们长久藏匿的东西——在他们起初(除Deligne外)费了九牛二虎之力才把握其意义与影响、并将其同化之后。我在这种贬损的调门里(它叠加在下文所谈的“藏宝”姿态之上),看到了一种双重补偿。一方面是在他们内心深处排遣一种不安(由于挪用了一样并不属于他们、而属于所有人的东西所引起的不安)——在他们自己的眼里假装贬低被挪用之物的价值。另一方面,则是面对被感受为一种超出他们之创造力化身的“父亲”的补偿(而他们无法承担那份同样安息在他们体内、一如安息在他们暗中引以为怨怼的那人身上的力量……)。我的“亡故”状态,以及直接继承人树立的榜样,创造了一种有利于“发泄”隐秘敌对的良机,“父亲”从此被感受为处于虚弱、低下的境地。}。

如果我面对他们,绝非像面对某种我托付了股票的股东(以不知什么共同“利益”的名义),而正是像面对以共同冒险与我相系的人们——因而那些人,会以充当我在他们耳边传递的“信息”的传递站为己任(哪怕随他们的心意在其中加入自己的东西,并在周遭引发回响……),正如我自己曾为了他们的利益而充当那份传递一样\fnote{464}{(**)}{因此,正是对这个“无边界的数学共同体”,我既直接面对它,又通过他们并穿越他们面对它。我在别处解释过(见第847页脚注(*)),为何在讨论班结束后的至少一年里,我不曾亲自负责将它誊清以供所有人支配。}。

相隔近二十年之后,我才意识到:在我们之间有着根本的误解——我们不曾“调谐在同一波长上”。我曾作为活生生的事物托付在自以为充满爱意之手里的东西,被像某种急于掩埋深藏的私藏宝藏一样囤积了起来。对私藏宝藏的占有代表着某种权力(鉴于其代价,这权力自然是可笑的……)——哪怕只是一种扣留的权力,一种阻拦(哪怕只在一时)一样生来便要绽放与播撒的活物绽放与播撒的权力。

我曾试图抓住在这场“误会”\fnote{465}{(***)}{在写下这些行与“误会”这个词时,一个联想不由自主地浮现在我心头:那是Zoghman~Mebkhout的信(引于注“教学的挫败——或创造与自负”,n\textsuperscript{o}~44$'$),其中谈到在我的学生与我之间存在“某种不理解”(然而把Deligne排除在外……)。我当时并不能肯定理解了他想说的是哪种“不理解”。这是否正是我在此谈到的这个“误会”——而他之所以将Deligne排除在外,是出于他蓄意(这位朋友的这种打算曾不止一次令我吃惊)只“从玫瑰色”看他?}中相互碰撞的两种本质迥异的姿态,在“阴,作为仆人,\dmark{}\origpage{864}与新主人”和“阴,作为仆人(2)——或慷慨”(n\textsuperscript{o}s~135, 136)这两条注中。我并不想在此摆出一副“服务姿态”模范化身的姿态,去同“等级姿态”相对立——后者把“知识”变成某种精英的独特标志,并(在风尚败坏的更深阶段)变成对他人施行任意权力的手段。正如在《愚妄与更新》(《收获与播种》第一部分)中的反思所表明的,现实要复杂得多。我曾在我本人身上、以及在我作为数学家的过去的某些行为中,察觉到今天所见到的普遍堕落的萌芽。而同样千真万确的是,我心中这股“服务冲动”曾是一大强大的行动引擎,尤其在推进我的数学文字工作、更具体地在不知疲倦地追求EGA与SGA两套奠基丛书的过程中\fnote{466}{(*)}{这项“不知疲倦的追求”往往直接违背了我心中同样强烈的另一种冲动:抛下所有牵制我的“任务”,奔向前方那不断召唤我(并且至今仍在召唤我……)的未知。}。

看来我未能把这股冲动的丝毫、或反映这股冲动的姿态,传递给我的学生。这项所开创的事业,在其体现为服务于一个共同体的姿态与性情的限度内,在我离开后便戛然而止\fnote{467}{(**)}{(4月10日)在重新誊写这几行时,我被这一局面的奇特讽刺深深震撼,其意义(正如整个下葬的意义一样)在此时此刻尚未被完全领会。正是那位全身心投入到服务于某个“数学共同体”的任务中的人,发现自己竟被从其自己的著作中驱逐了出去,而且得到了上述“共同体”默许且毫无保留的赞同——被恰恰把拒绝服务奉为等级信条与第二本能的人们所驱逐。
然而,一旦忆及我心中这股“服务冲动”所面向的“共同体”,绝非那个毫无保留地参与了我的下葬的社会学实体(带着其头面人物的“等级”等),而是前文所谈及的那个“在空间与时间中皆无边界的数学共同体”,这一表面的悖论在我看来便在很大程度上消解了。(关于这两个“共同体”之间的区分与混淆的评述,见后文注“尊重”(n\textsuperscript{o}~179)的第一条脚注。)}——仿佛突然被电\dmark{}\origpage{865}锯(或链锯……\fnote{468}{(*)}{关于“链锯”效应——它硬生生斩断了我几乎所有前学生的一项正在展开的著作的生机勃勃的势头——见“共同继承人……”与“……及链锯”两条注(n\textsuperscript{o}s~91, 92)。})截断了一样。

从我所离开的世界在此处彼处仍传达给我的回声里,我看到:这种自发的姿态(当年我与在我起步时接纳我的仁善前辈所共有的姿态),已经(正如那份仁善本身一样)在我曾所属的这个世界里成了一个陌生的异乡人。

\subsubsection*{b6.~驱逐}
\tocsubsub{b6.~驱逐}

\notehead{169$_1$}{}

(3月9日)\fnote{469}{(**)}{本子注源自主注“操演”(n\textsuperscript{o}~169)的一条脚注(见第848页脚注(*))。}“SGA”这一缩写,是“布瓦玛丽代数几何讨论班”(Séminaire de Géométrie Algébrique du Bois Marie)的代称。它指称(或至少,在六十年代曾指称)这样一些讨论班:1950至1969年间(自1962年起,与学生及其他合作者一道),我在其中展开了我那新的代数几何基础纲领,与EGA(《代数几何原理》)丛书那些(不那么“前沿”、文体也更规范的)文本并行\fnote{470}{(***)}{系与J.~Dieudonné合作撰写。}。这些讨论班自1962年起在“布瓦玛丽”举行——那是Bures-sur-Yvette的一处地名,IHES自1962年起即设于彼。说实话,头两期讨论班(1950至1962年间)是在巴黎一处因陋就简的场所(蒂耶尔研究所)进行的,听众从未超过十来个人,而我在他们面前严格地“单枪匹马”。“SGA”这个缩写就诞生于那些年头,当时还谈不上什么“布瓦玛丽”。我只是在后来,才给“代数几何讨论班”这个初始的名称添上这漂亮的别号,好让它少一点干巴巴的味道。

不言而喻,这些讨论班的后续,从SGA~1直到SGA~7,是按年代先后编号的。同样不言而喻:每一期讨论班的总体构想都出自于我。它受我那全局性与长远意图的启发:为代数几何奠定广阔的基础,并且——越来越同时地——为一种更广阔的“几何学”奠定基础;这种几何学,我从1963年前后起就强烈地感受到它,而它一直没有名字(今天我会以“算术几何”称之,即代数几何、拓扑\dmark{}与算术的综合\fnote{471}{(*)}{关于这一点,见第844页脚注(*)。})。这些讨论班中的最后一个\origpage{866}是SGA~7,它(与其前的各期不同)连续办了两个年度,1967--69,由我与Deligne共同主持。

那部名称玩着障眼法的卷册“SGA~4$\frac{1}{2}$”(如上文第847与851页所解释)由1973年之后的文本编成,因而也晚于SGA讨论班的最后一期——倘若不算那些从SGA~5里劫掠来的文本,以及Verdier一篇“博士论文”的那份著名的“第0状态”(第三场行动处将谈到它)。撇开日期问题不谈,构成这一卷册的文本之五花八门的性质,与我当年推进SGA丛书的精神毫无相干:SGA丛书的每一卷,都针对我的纲领中一个此前在任何地方都不曾展开过的部分,做出一项大规模的奠基工作——因此,它排绝一切“文摘”式的卷册,排绝对那些早已尽人皆知、打磨停当的结果的汇编,也排绝那些即使是新的、却属零星偶得的结果。退一步说,即便把Deligne那一卷命名为SGA~8(假定我会首肯)也不妥当:这样一个名称,会暗示(一种全无根据的想法)我在SGA~1至SGA~7各期讨论班中所推进的事业仍在延续。至于Deligne选定的“SGA~4$\frac{1}{2}$”这个标号,它就不只是“不妥当”了:它本身就是一场骗局与冒名顶替。在我看来,这一点对自1977年以来有机会读过此卷的众多数学家——他们同时又懂得SGA这个标号的含义,懂得它与我这个人、与我的著作、并由此与某种精神不可分离——本应一望即知。尽管如此,这一冒名顶替,在一个标准参考文献的名称里,已被“数学共同体”容忍了八年,表面上竟“不起一丝波澜”。我在其中——连同作为其自然延伸的1981年反常层会议——看到七、八十年代数学界的奇耻大辱;在我看来,这一耻辱在我们的科学史上没有先例。

作为这场把我从SGA丛书中驱逐出去的行动的先声,有一个插曲意在制造这样的印象:在SGA基础文本的发展中,我这个人只扮演一个偶然的、杂乱而次要的角色。这就是所谓的“SGA~7微型行动”。在注《操演》(n\textsuperscript{o}~169)的“第3幕”(一场节节加码)里谈到过它,而尤其(从我此处关心的视角看)是在\origpage{867}注《一场屠杀的序曲》(n\textsuperscript{o}~56)里。它说的是:把原初讨论班的一部分,以Deligne与Katz之名、将我排除在外,单独发表为SGA~7~II一卷(同时把我在其主体主题与若干关键结果的发展中所起的作用掉包隐去)。我在所引的那条注(n\textsuperscript{o}~56)中就此写道:“这一‘SGA~7’行动,绝非对SGA各期中所推进事业的延续;我把它感受为一种粗暴的‘锯击’(或电锯……),以一卷同我这个人公然划清界限的卷册,终结了SGA的序列——尽管它系于我的著作,并且同其他各卷一样带着它的印记。”SGA~7~I与SGA~7~II这两卷,尚未摆出那副对它们所从出的著作近乎掩饰的轻蔑与鄙夷的神气。倘若这一加码的步伐四年后终究迈了出去,那是因为先前的步子(其中包括这桩看似无奇的SGA~7微型行动)都“过去了”,从未(至少据我所知)在数学界激起半点反应。

我想以一段(想必是暂时的)富有教益的尾声,来结束这场由Deligne操办、得到“全体教众”默许的、把我从SGA中驱逐出去的行动。那是Byrne夫人的回答——她执掌Springer出版社的“Lecture Notes”丛书——口气十分“冷淡”。我曾写信给她,要求澄清一卷名为SGA~5的书:它于1977年以我的名义在“Lecture Notes”丛书出版,由Illusie编辑,而Springer出版社非但没有征求我的同意,甚至没有通知我它经手了这次出版。我从她的回信(一个月后才收到)中获悉:纠缠这类手续纯属多余,因为我自认为是该SGA~5卷的作者,这实属谬误——我不过是这期讨论班的“主持人”,仅此而已,它由Illusie编辑!(于是人们不免回过头来纳闷:那位已故的主持人,当年在那个讨论班里究竟干了些什么……)我抱着看看的念头,亲自写信给K.F.~Springer先生,述说我自1972年以来(那一年,SGA~7~I以同样的方式冠以我的名字出版——说真的,我对它的“作者”身份,并不比我对SGA~5的更多……)与Springer出版社打交道所经历的种种怪事。我至今仍在等他的回音……\fnote{472}{(*)}{(4月9日)关于这个故事的后续,见注“殡仪馆——为科学服务”(n\textsuperscript{o}~175)。}

\origpage{868}(3月16日)本子注得到了它当之无愧的名称:“驱逐(2)”。这个(2)提醒人们:先前已有一条同名注“驱逐”(n\textsuperscript{o}~63),我最近还有机会援引过它(连同“Motifs”那一行动)。那条注里(极其含蓄地……)提到的“驱逐”,是1970年我离开IHES那一幕中的驱逐;而那次离开,显然正合我那位年轻而杰出的朋友的心意——他彼时刚刚在取得的位子上安顿下来\fnote{473}{(*)}{关于我1970年离开IHES的插曲,见“未竟的收成”一节(n\textsuperscript{o}~28),以及注“有益的挣脱”“驱逐”“兄弟与配偶——或双重签名”(n\textsuperscript{o}s~42、63、134),最后见后引注的子注(n\textsuperscript{o}~134$_1$)。}。这两次“驱逐”——一次出于IHES,一次出于SGA丛书——之间的血脉渊源,在我看来一目了然。我在其中看到一种触目的递进,性质上仍是一次“节节加码”:第一次,只是把我这个人从一个机构中驱逐出去——我对那个机构固然依恋极深(我真的曾打算在那里终老,千真万确!),但我很快就脱身而出,不留一丝遗憾;第二次,则是把我从SGA中驱逐出去——而SGA(象征性地,甚至远不止象征性地)代表着我的数学家著作,一部直到今天我仍系念的著作。诚然,我被“驱逐”出IHES,十五年来早已尘埃落定;而我却始终怀疑,对一部我曾倾注十五年充实岁月的著作,事情是否同样如此。

我想到,当年我曾亲手为把我从SGA中驱逐出去提供过方便:我顺着自己自发的冲动,把那些曾在某些时刻全身心投入某期讨论班之发展的学生与合作者,呈现为与我同为该讨论班的“主持者”。这在我那个年代并非常例,在今天更绝无可能。我不知这样做对不对。一方面,这并不完全符合实情:我在其中所起的作用,同我的合作者们(纵然才华横溢、纵然如我一般全力以赴)所起的作用,绝无对称可言。因此,这种呈现方式,恰恰合乎我在注“上升”尤其是注“暧昧”(n\textsuperscript{o}s~63$'$、63$''$)中所考察的那种“暧昧”(或者说,对杰出的青年数学家的“纵容”)的方向。倘若由我设立的这种暧昧,竟鼓励了那些曾热烈参与SGA的人们中的一些人,在某个时刻来“驱逐”我(以或多或少局部的、或或多或少彻底的方式),那我怨他们便是我的不是!我不过是收获我自己播下的种子。

\origpage{869}但这丝毫不妨碍我在另一方面,对已然发生的事情作一份公开的陈述。另一方面也确实如此:我同某些合作者建立的那种关系,在他们看来可以是一种器重与信任的标记(它也确实是),并由此鼓励他们同我一样全身心投入这项任务。但我如今告诉自己,这样的器重与信任,完全可以同样清晰、同样鼓舞人心地表达出来,而不必因此染上暧昧。这就好比我在“购买”一份与任务相称的投入,所付出的代价是一份“好处”——而这份“好处”(回过头看)在我看来形迹可疑。因为,装出自己本不是的样子,是一种虚假的好处。而十分显然:在制造这样一种(即便不是彻头彻尾、至少也)边缘走样的假象这件事上,首当其冲负起责任的,是那个以长者身份出现的我,而不是任何别人。

渐渐地,这场反思在一道我落笔时全然不曾料到的“个案情形”的意外光照下,越来越趋近注“暧昧”的反思。我意识到:正如我同那位“年轻天才”(绝非无名之辈)Deligne的关系是虚假的——出于一种虚假的谦逊,我不肯承担那在我与他之间确实属于我的长者与“老师”的角色——同样,我同另一些才华横溢、在那项我当时视为“共同”的任务中全力以赴的年轻人之间的关系\fnote{474}{(*)}{直到注“绿灯”(n\textsuperscript{o}~68)的末尾,我才开始认识到:这绝非误会,而是一场彻底的幻灭。},也同样是虚假的。《下葬》中展开的反思,已经足够清楚地表明:倘若有所谓“共同任务”,那也只是一年顶两年的光阴——是让那个年轻人(姑且这么说)做出一篇博士论文的时间(这已经不算小了)。恰恰在我离开的1970年,响起了信号:这一整片明明灼手、前一天还实实在在“在我手中燃烧”的庞大“任务”,当即、几乎彻底地被弃置一旁\fnote{475}{(**)}{这一在离去的翌日便对一项纲领与灼手任务的即刻抛弃,在注“本能与时尚——或强者的法则”(n\textsuperscript{o}~48)中、尤其在与“共同继承人……”“……与电锯”两条注(n\textsuperscript{o}s~91、92)中曾被谈及;在那两条注里,我(依据传回我耳中的回声)试图逐一点检那些曾由我“1970年以前”的各色学生接手的主题,后来的下落。}。撇开Deligne关于Weil猜想的工作不谈,这同时也是每一个曾最深吸引我的重大主题长期停滞的开端\dmark{}\origpage{870}——这场停滞(除由不予点名的Mebkhout的哲学所触发的那场部分“重启”之外)至今仍在继续\fnote{476}{(*)}{(3月17日)这种“停滞”的印象,将在后面一条注里获得或许更具体的意义:我打算对那些曾列在我日程上、并在我离开之际被我的学生们以完美的齐一性弃置不顾的最“灼手”的主题,作一份简短的带评点的清单。(4月9日)关于这一点,见注“工地的巡礼——或工具与视野”(n\textsuperscript{o}~178)。}。

\subsubsection*{b7.~好撒玛利亚人}
\tocsubsub{b7.~好撒玛利亚人}

\notehead{169$_2$}{}

(3月13日)\fnote{477}{(**)}{本子注源自主注“操演”(n\textsuperscript{o}~169)的一条脚注(见第849页脚注(**))。至于把那套“暂停!”手法——专司把鱼鳔当灯笼卖给一位急不可耐、只求相信的“用户”——拆解得更周详些,见子注“特洛伊木马”与“公式”(n\textsuperscript{o}s~169$_3$与169$_5$--169$_8$)。}在SGA的这一篇序言里,Illusie热忱地感谢Deligne,其中尤其感谢后者“说服了他撰写……Lefschetz--Verdier公式的证明,从而为此讨论班的出版扫除了一道障碍”(着重号是我加的)。说白了:那障碍,便是Illusie十一年来不肯动笔撰写他承诺要撰写的文本的这一“缺乏信念”——而这缺乏,如上文所说,恰恰在好撒玛利亚人Deligne向好撒玛利亚人Illusie亮出“绿灯”、说他“可以放手去做”的那一刻,戛然而止……

这正是“假中之真”。至于这段文字显然力图暗示、却不必明说的那个“假”(循着那套久经考验、早已自成一派的手法……),则是:SGA~5讨论班须依赖那条公式(该公式在讨论班当时,仅凭奇点消解的假设方能成立;这一假设,后来已由Deligne在那部名为“SGA~4$\frac{1}{2}$”的“先行”卷册中提出的有限性结果,在最常见的情形下予以扫除)\fnote{478}{(***)}{关于这一点,见第841页脚注(***)与第850页脚注(*)。}。实际上,正如这两位朋友同我一样清楚无疑:Lefschetz--Verdier公式在SGA~5中所起的作用(一如它在我那$\ell$-进L函数上同调公式的证明中所起的作用)纯然是启发性的:它提供了\origpage{871}动机,促使人们去寻找并证明“显式”的不动点公式(即其中“局部项”可以被显式算出的那种)。于是,Illusie便同他的朋友一唱一和,制造出这样的印象:SGA~5确实(且是在一种他同他的朋友都不肯明确言明的意义上)从属于那部由此只能以“SGA~4$\frac{1}{2}$”为名的文本。

欲知详情,见注《屠杀》及其子注n\textsuperscript{o}~87$^2$。在这条注及其全部子注中,我终于得以弄清(迟悟总比不悟好):Illusie撰写的整篇序言,以及一般地说,Illusie版(或屠杀版)的整体呈现,乃是一份背信弃义的范本——它以一副漫不经心的姿态、伴着那份为其人格平添魅力的天真神气侍奉出来。

Illusie竭力营造的这份动人印象——倒霉的SGA~5讨论班之所以终得出版(晚了十一年,而且是付梓于我所知道的那个状态),全亏了好撒玛利亚人Deligne(自然还有第二位好撒玛利亚人Illusie)——看来“安然通过”,未遇任何阻碍。我在Serre1977年为菲尔兹奖国际评选委员会撰写的关于Deligne工作的报告中,重新找到了这一版本。我对Serre的十足诚意毫不怀疑;他本来也只是远远地追随那场口头讨论班的曲折周折——更何况此间早已逝水年华……他想必同所有人一样、不加质疑地,把Illusie序言里说出或暗示的东西当了真——那序言他总该有一天翻过,为的是看看(而他什么也没看出来!)……

有趣的是,据我所知,同一份Serre报告,也是全部文献中唯一一处(而且就在报告的第一句)提到Deligne曾是我学生的地方。而Deligne的任何一篇出版物,则决不会让哪怕一个读者猜想:这位作者或许从我这里学到过什么。

\subsubsection*{b8.~特洛伊木马}
\tocsubsub{b8.~特洛伊木马}

\notehead{169$_3$}{}

(3月10日)\fnote{479}{(*)}{本子注(属注《操演》)源自该注的一条脚注;见第860页脚注(***)。}在注“白板”的子注(n\textsuperscript{o}~67$_1$)里,我指出过两个例子:Deligne在其中越出了他素来的谨慎,当真“直言不讳地道出了假话”。为那位好奇、消息灵通、手头却没有那条注与其子注的读者着想,我在此提示:撇开他对SGA~4与SGA~5的那些“殷勤话”,\dmark{}以及对我这卑微\origpage{872}的人在各处或明或暗的大“遗忘”(这些在注《操演》及其脚注里已有这里那里的指涉),我所识破的那些明目张胆的骗局,集中在第2页的第3、4两段里(在“SGA~4、SGA~4$\frac{1}{2}$、SGA~5的阿里阿德涅之线”一节里——请欣赏这列壮美的队伍吧……)。这十七行文字,堪称“浑水摸鱼”之术的范本,大值得逐细剖析\fnote{480}{(*)}{关于所引两段中第二段的详细评注,见子注“双重含义——或行骗的艺术”(n\textsuperscript{o}~169$_7$)。}。

我在此只须指出:在所引第一段里,人们读到,为了“在étale上同调中建立一个与凝聚对偶相类似的对偶形式体系”,“Grothendieck曾使用奇点消解与纯性猜想”\fnote{481}{(**)}{原文在“纯性猜想”之后接着写道:“在SGA~4第XVI讲中于一个相对框架下[??]确立,并且——模去消解——在等特征情形下于SGA~4第XIX讲中确立。”这句“在一个相对框架下”(对任何不预先知情的读者都无从索解),正是掩盖下述事实的一种手法:该定理对一切特征的光滑代数簇早已成立。\par(3月17日)我直到此刻才注意到所引那段话结尾的妙处,它在头几遍阅读时竟“蒙混过关”:“若干发展见于SGA~5~I。在SGA~5~III中,人们展示这一形式体系[??]如何蕴含那条极为一般的Lefschetz--Verdier迹公式。”(着重号是我加的。)人们当会欣赏这“若干发展”之不着一字,而作者(他在别的场合是懂得精确的)随即接上“这一形式体系”(=若干发展?),说它“蕴含那条极为一般的迹公式”;好让紧接着的下一段一开篇便凸显:那条公式“在SGA~5的原始版本中”,“仅以猜想的方式得到确立”。我刚刚在SGA~5里核实了这些“若干发展”究竟是什么——第I讲的标题便告诉了我:“对偶化复形”,因而也包括双对偶定理。为什么说“若干发展”,而不说“对偶化复形的理论”或“双对偶定理”?其实并不更长,而且总不那么浑浊!这叫我想起:在那份著名的“有限性”报告里——亦即在“特洛伊木马”里——那位杰出的作者恰恰证明了一条“双对偶定理”,对我这卑微的人一字不提;而这条定理旋即(在由Illusie撰写的SGA~5第I讲引言里)被干脆命名为“Deligne定理”。的确,环环相扣……附记:关于这条双对偶定理(它竟被如此故作漫不经心地对待)的评注,见第852页那条长长的脚注(*)。}。紧跟着便补充说:在本卷中(多亏上苍与那位杰出的作者),这些“要点已由另一方法确立”(着重号是我加的),而该方法有效于“定义在0维或1维正则概形上的有限型概形”——也就是说,实际上有效于使用者会遇到的几乎所有情形。

如此,Deligne便竭力营造(甚至明白断言)这样一种印象:我所发展的étale对偶的整\origpage{873}个形式体系仍属猜想(至少在非零特征下如此),而这些“要点”终究只有靠他Deligne本人、靠本卷(也就是说,靠他的有限性结果——前面几条脚注已提及、他也随即指向的那些结果)方得确立。这可真是,瞧瞧!,正足以替那著名的“逻辑依赖”的虚构张目——SGA~5依赖于那部名为“SGA~4$\frac{1}{2}$”的文本(这层依赖由这名称本身、由“SGA~4—SGA~4$\frac{1}{2}$—SGA~5”这列壮美的队伍所设立),并由此替其引言中那句令人难以置信的断语(已在别处引用并评注过)开脱:“它的存在[‘SGA~4$\frac{1}{2}$’的存在]将使SGA~5不日得以按原貌出版。”

这便是Deligne的版本——它被悄悄塞进那部名为“SGA~4$\frac{1}{2}$”的电锯文本,并且毫无含糊地滑入我方才引用的“阿里阿德涅之线”那段文字之中。而实情是:早在1963年3月,我就已在étale框架中确立了六种运算的完整形式体系(远远越出通常的Poincaré对偶),唯一的限制性假设(显然不可少)是:须以“素于”所论概形之剩余特征的挠系数为系数\fnote{482}{(*)}{由此,“六个函子”以及关于它们的基本公式——其中最关键的是有限型态射的“对偶公式”(人们可把它看作迄今为止可以想见的、经典Poincaré对偶定理的最一般形态)——皆由我确立,自始至终不曾对系数强加任何有限性假设。况且,这一点Deligne比谁都清楚:恰恰是他,以我1963年的笔记为底本,把SGA~4中展开这一对偶形式体系(围绕上述对偶公式)的那份报告写成了详细的文本!}。唯独étale上同调的双对偶定理,我的证明才用到Deligne所指的那些假设。而这后一条定理,在我发现它之前,在上同调(各种“空间”的上同调)里属于前所未知的类型;况且,它在SGA~5讨论班中也只扮演过一个插曲性的角色,那是为了证明Lefschetz--Verdier公式\fnote{483}{(*)}{(3月17日)尽管如此,在所引第二段里,Deligne随即接着强调:该公式“仅以猜想的方式得到确立”,且“其中的局部项并未算出”(这一断语在数学上毫无意义,却有助于为SGA~5营造一片“废话矿渣”的印象,好让大家慈善地把它忘掉……)。我自己承认:一年前头几遍读这些段落时,我如遭雷击——这些奇怪地“不得要领”的评说,针对的又是一部人们本来就劝人遗忘的文本,我当时全然摸不着头脑。正是凭着事后的回顾,凭借一次细心的“按件”复核,那份掉包的意图才终于显形:它由一套精密而万全的“浑水摸鱼”式隐匿手法侍奉着,藏在我最初只当是表皮式恶意、沿着一支讨好的笔端信马由缰地挥洒的东西背后。欲知这套手法更周详的考察,见注《操演》的诸子注“公式”(n\textsuperscript{o}s~169$_5$--169$_9$)。},而该公式本身也只起纯启\origpage{874}发作用\fnote{484}{(**)}{正如我在下文子注“正经数学……”(n\textsuperscript{o}~169$_5$)中所强调的,这条公式曾在心理上十分重要:它为发展“显式”的不动点公式提供了动机。}。在Deligne那部伪作之中,上述双对偶定理的作用干脆等于零(若不算上它在其中于一些“帮衬性”的假设下得到证明,并且——在Illusie那副谄媚的笔端下、在其朋友的怂恿之中——摇身变成了“Deligne定理”……)。

此处无意贬低Deligne那些有限性结果的价值:它们确实填补了SGA~5的一项空白(如同其他许多空白之一),而这是事物之常。任何一门朝气蓬勃的数学理论都不曾“完成”!但人们不能不看到:Deligne利用了这份贡献——它再有用,也仍是有限度的(他还做过更深刻、更困难的事,而且并不费力气……)——把它吹胀到失度,使之成为一场规模空前的诈骗行动的“特洛伊木马”:这就是“étale上同调”行动。

这同一匹“特洛伊木马”,还重新出现在Deligne为名为“SGA~4$\frac{1}{2}$”的那部卷册写给Zentralblatt的“书评”里(前引之文;见第851页脚注(**))。在该书评的最后一段里,我读到:“人们证明,对于定义在一个1维正则概形S上的有限型概形,通常的那些上同调运算[更不必说那‘六种运算’,那是万万点不得名的!]把任何可构层变成可构层。”(着重号是我加的。)这话的表述方式,足以让人以为:在这位作者所呈献的这部杰作之前,人们手头\origpage{875}并没有任何一条针对étale上同调中那些著名的“通常运算”的有限性定理\fnote{485}{(*)}{这正合乎同一篇书评前文所说的“SGA~5的混乱状态”——而本卷本当“补其偏弊”。}。可我却有幸证明了第一条这样的有限性定理,而且是一切之中最关键的一条:关于函子$Rf_!$(带真支撑的上同调)的定理——而且这恰恰是在我发现这样一个函子于étale上同调中的定义之后的那几天里(如果我的记忆无误;当$f$为真时,它与“平凡的”$Rf_*$重合)。那是1963年2月,在我有幸结识我这位未来的学生之前,而在那时,除了我(勉强再算上Artin),还没有人十分有把握:étale上同调这东西究竟“存在”与否。它正是在那些日子里,开始真正存在的。

还剩下$Rf_*$的类似问题,它证明更为顽抗,而且事实上直到今天,仍未以它(想必)应有的全部一般性获得解决。其实,就在同一年(甚至同一个月),我已完成那些必要的“拆卸”(这在今天,随便谁来都能一挥而就……),表明从$Rf_*$的有限性出发,可以证明$Lf^!$与$R\mathcal{H}om$的有限性\fnote{486}{(**)}{至于六种运算中余下的两个,即$Lf^*$与$L\otimes$,它们把可构系数变成可构系数,这是平凡的。}。诚然,这从此成了étale上同调的“基本常识”,并且想必属于我那杰出的“前辈”“SGA~4$\frac{1}{2}$”注定要使人“遗忘”的那些“技术性离题”之列……

\subsubsection*{b9.~“那条”猜想}
\tocsubsub{b9.~“那条”猜想}

\notehead{169$_4$}{}


(3月12日)\fnote{487}{(***)}{本子注源自注“操演”(n\textsuperscript{o}~169)的一条脚注;见第857页脚注(***)。}自Deligne发表其著名文章“La Conjecture de Weil~I”(在其中他攻克了Weil猜想中我曾悬而未决的最后一环)以来,我不止一次注意到一件奇特的事,但直到最近几天才真正去探究:Deligne总是谈论\textbf{那条}Weil猜想,而此前人们习惯的说法始终是Weil\textbf{诸}猜想。

确实,在Weil1949年的经典论文《有限域上方程的解数》(载Bull. Amer. Math. Soc. 55 (1949), pp.~497--\origpage{876}508)中,这些猜想正是以一连串一个比一个更令人叹服的断言的形式呈现的;而我在五十年代中期从Serre口中习得它们时,也是作为这样一整套猜想群。诚然,在这个乍看五花八门的猜想集合里,有一种显然的灵感之\textbf{统一}:它首先来自与上同调形式主义相关联的直觉(经由Lefschetz公式),同时也(我至少如此揣想)来自Hodge理论。

通过创造并发展这样一套\textbf{上同调工具}——适用于基域任意的代数簇——我得以证明这些猜想中的大部分。我在Artin、Verdier及其他人的协助下完成了这项工作,把生命中紧凑的三年献给了一场实地精雕细琢的工夫——它凝结为两千页“不可卒读”的“废话矿渣”与“技术性离题”,正是这些,使得一位Deligne能在区区二十页密排文字里“砍下”最后一步……此外,受Serre发现的一个出色的Weil猜想“Kähler类似物”启发,我提炼出了(以我所谓关于代数闭链的“\textbf{标准猜想}”的形式)把Hodge理论\textbf{移植}到任意基域上去的原则——或更确切地说,把Hodge理论中从“代数”观点对复代数簇上代数闭链理论真正要紧的那部分加以移植。只要把这些猜想按其初始形式(或许过于乐观)稍加(且显然地)重新表述,它们至少在特征零中成立,并且在特征~$p>0$~中也“一定为真”(只要Weil诸猜想为真……)。

这决非巧合:同一位执意把Weil猜想“变为单数”的Deligne,也正是那个竭力抹杀其证明中他的老师所扮演角色的人;也正是他,还努力(鉴于普遍的冷漠,他成功了)把“标准猜想”败坏成一条死胡同、一条无法企及的路——说穿了,是一个\textbf{障碍}——好在通往Weil猜想的证明之路上,仰赖上帝(和他那微不足道的本人)把它绕了过去\fnote{488}{(*)}{(3月16日)关于这一双重“掉包-拆解”的若干细节,见《葬礼悼词》(注n\textsuperscript{o}~104、105),以及注171(x)开头关于这篇悼词的寥寥数语。关于掉包之术的更周详考察,见“公式”诸子注(n\textsuperscript{o}~169$_5$--169$_9$)之全体。(x)(5月11日)旧注“封神”的这一开头部分已与该注分离,成为一篇独立的注“珍宝”(n\textsuperscript{o}~170(iii))。}……

\subsubsection*{b10.~公式}
\tocsubsub{b10.~公式}

\subsubsection*{(a) 正经数学……}
\tocsubsub{(a) 正经数学……}

\notehead{169$_5$}{}

\dmark{}(3月17日)\origpage{877}对于定义在有限域~$k$~上的代数簇~$X$,出名的“Weil诸猜想”关乎与~$X$~相关联的“L函数”(所谓“Artin-Weil函数”)。它被定义为某一有理系数的形式级数,对该级数的了解等值于对~$X$~在~$k$~上及其全部有限扩张上的有理点个数的了解。诸猜想中的第一条断言是:该形式级数(常数项为1)是一个\textbf{有理函数}在$\mathbb{Q}$~上的级数展开。所有其余断言都关乎该有理函数的特殊形式与性质——在~$X$~为连通射影且非奇异的特殊情形下。这些猜想的核心是一条假定典范的公式,它把该有理函数呈为
\begin{quote}
$$L(t)=\frac{P_0(t)\,P_2(t)\cdots P_{2n}(t)}{P_1(t)\cdots P_{2n-1}(t)}$$
\end{quote}
其中诸~$P_i$($0\le i\le 2n$,$n=\dim X$)是常数项为1的整系数多项式。~$P_i$~的次数~$b_i$~应当充当~$X$~(更精确地说,对于域~$k$~的代数闭包~$\bar{k}$~上与之对应的簇~$\bar{X}$)的第~$i$~个“Betti数”。于是,当~$X$~由定义在特征零域~$K$~上的非奇异射影簇~$X_K$~经“特征~$p>0$~的约化”而来时,$b_i$~应等于该\textbf{复}代数簇的第~$i$~个Betti数(以超越途径定义)——后者经由~$K$~到~$\mathbb{C}$~的任一嵌入从~$X_K$~得来\fnote{489}{(*)}{Weil作这些猜想时,人们甚至还不知道如此定义的$b_i$与$K$到$\mathbb{C}$的选取无关。几年之后,这一点将由Serre的凝聚层上同调理论得出——该理论给了Hodge理论中更精细的不变量$h^{i,j}$一个“纯代数的”意义。}。该有理函数须满足一份\textbf{函数方程},等价于说:$P_{2n-i}$~的根恰是~$q^{n}/\xi_\alpha$,其中~$q=p^f$~是基域~$k$~的基数,$\xi_\alpha$~遍历~$P_i$~的根。(从义理上说,这该当“源于”簇~$\bar{X}$~那未具名、未定义的“上同调”之“Poincaré对偶”的存在。)

我相信Weil想必也猜到了:对~$i\le n$,$P_{2n-i}$~的零点恰是~$q^{n-i}\xi_\alpha$,其中~$\xi_\alpha$~遍历~$P_i$~的零点(或者——鉴于对偶条件——等价地说,$P_i$~的零点成对配组,每对的乘积等于~$q^i$)。这里的启发式“理由”是复非奇异射影簇\dmark{}上同调的另一条重要性质,\origpage{878}这一次由“Lefschetz定理”(所谓“蛮牛”版,即“困难”版)来表达。最后,Weil猜想中最末一条——Riemann猜想的“几何”类似物——说:$P_i$~的零点的倒数的绝对值都等于~$q^{-i/2}$(这条断言导致对~$X$~点数的高精度估计\fnote{490}{(*)}{由Weil猜想中的这最后一条,同时得出:$L$函数的写法$(L)$是\textbf{唯一的}。})。

对一般簇~$X$~的L函数的有理性,曾由Dwork于1950年以非上同调的“$p$-adic”方法确立。该方法的不便在于:L函数没有上同调解释,从而对其他猜想在非奇异射影~$X$~情形的攻克无从下手。在后者情形,拥有一个上同调形式主义(系数域~$R$~特征零)、含Poincaré对偶、含结合于闭链的上同调类形式主义(把相交变为杯积),就能本质上以“形式”的方式转录经典的“Lefschetz不动点公式”。把这一公式应用于~$\bar{X}$~的Frobenius自同态及其迭代,便将得到Weil所要求的形如(L)的表达式,其中诸~$P_i$~是~$R$~系数的多项式。这在Weil1949年陈述这些猜想时就该是清楚的,至少对五十年代的Serre和我而言是如此——这正是发展这套形式主义的初始动机。早在1963年3月,取~$R=\mathbb{Q}_\ell$($\ell\ne p$),这事就已做成。只剩两粒盐:

a) 先验地不清晰(尽管人们相信理当如此):多项式~$P_i(t)$——先验以$\ell$-adic整数环~$\mathbb{Z}_\ell$~为系数——实际上是\textbf{普通整数}系数的,而且不依赖于所取素数~$\ell$($\ell\ne p=\operatorname{car}k$).

b) 从非奇异射影~$X$~的L函数有理性,只有拥有奇点消解,才能推导出一般~$X$~的情形。

a) 提出的问题对Motif的yoga之绽放与发展、对与该yoga紧密相连的\textbf{标准猜想}的后来表述,扮演了关键角色。\origpage{879}它们也激励人们寻找一套\dmark{}\textbf{$p$-adic上同调理论}(后来由“结晶”理论实现),作为证明诸~$P_i$~系数之整性的可能进路——一旦(比如经由标准猜想的肯定解决)人们知道它们是有理的且不依赖于~$\ell$(\textbf{包括}$\ell=p$~的情形)。

不管怎样,到1963年人们已经拥有了L函数的表达式(L)(尽管它先验地依赖于~$\ell$~的选取)、函数方程、以及Betti数在特化下的良好行为。剩下的就是解决问题a)、证明关于~$P_i$~零点绝对值的断言、以及最后(为了凑足分量)关于~$P_i$~零点的“Lefschetz式”关系。这些正是十年后在Deligne的文章“La conjecture de Weil~I”(Pub.~Math.~de l'IHES~n\textsuperscript{o}~43~(1973), pp.~273--308)中所做的。

作为Deligne证明的成分,根本不需要比“普通”不动点公式更精致的东西——后者早在1963年初就已可用(不带任何“猜想性”成分)。如果我没记错,Deligne文章里仅有的另一项上同调成分,是我约1967或1968年发展的Lefschetz束的上同调理论(平展版),加上Picard-Lefschetz公式(由Deligne在平展框架下证明);两者都发表于SGA~7~II卷——对这部\textbf{前面已经谈到过}的卷册,我的名字不出所料地几乎消失了……

被称为“Lefschetz--Verdier”的“更精致”不动点公式,倒是扮演过一个重要的\textbf{心理}角色:它鼓励我去提炼L函数的上同调解释(L)——对任意簇(不必然射影非奇异)成立。Verdier的这条公式提醒我:不动点公式应当无须对~$X$~附加非奇异性条件(一如在普通Lefschetz公式情形中早已众所周知);而尤其,它把我的注意力引向\origpage{880}这样的事实:存在关于以\textbf{任意}(“可构”)\textbf{层}为系数的上同调的不动点公式——它把(以这种层为系数的上同调空间中的)交错和迹,解释为对应于自同态~$f:X\to X$~之不动点的“局部项”之和(当后者孤立时)。在这份启发式动机中,这条Lefschetz--Verdier公式在特征~$p>0$~下“仍属猜想”(因为还没有奇点消解,\dmark{}从而没有“双对偶定理”)这一事实,\textbf{完全无关紧要}\fnote{491}{(*)}{(3月20日)这一点是如此显著,以至于去年,我早已全然并且长久地忘记了这一事实,而在读到(出自Deligne的笔下)Lefschetz--Verdier公式“仅在SGA~5的原始版本中以猜想方式确立”时,惊得瞠目结舌。我在随后两天(3月18与19日)的反思中回到这一点。(在子注n\textsuperscript{o}~169$_6$与169$_7$中。)}。

一如常有的情形,关键的一步是找到“\textbf{正确的表述}”(在此处,是L函数的“上同调公式”的正确表述)。Verdier的公式启发我引入一个任意的(可构)$\ell$-adic层,取代此前一直隐而不显的常值层$\mathbb{Q}_\ell$。于是,仿照Weil对“普通”L函数的定义,须定义一个“以层~$\mathcal{F}$~为系数的”L函数。一旦想到这一点,定义便自动涌现:那就是我1964年12月Bourbaki报告(《Lefschetz公式与L函数的有理性》,Sém.~Bourbaki~279)中所给出的——在此无须重复。此外,Lefschetz--Verdier公式的合理“局部项”(用给定的系数层与Frobenius对应来表述)也同样自动涌现。最后(人要么大胆要么不大胆!),为什么索性在这里书写公式时,连“正统”Lefschetz--Verdier公式的\textbf{真性}假设也一并放弃,改用\textbf{真支撑上同调}?!

因此,关键的一步,这一回又是提炼出“\textbf{正确的}陈述”(此处即“\textbf{正确的}公式”),\textbf{足够一般}——从而足够\textbf{灵活},能经得起递归与“拆卸”,畅通无阻地通过。我(以及至今任何人都不)无法\textbf{直接}证明“普通”L函数的那条公式——对任意~$X$(甚至光滑但不真,或反之)——以常值$\ell$-adic层的真支撑上同调为项,而不经由层论推广。(一如在特征~$p>0$~下,我无法不经由把Riemann-Roch-Hirzebruch公式\textbf{首先}推广为对光滑代数簇\textbf{真态射}的层论公式,就\textbf{直接}证明它的\textbf{普通}版本——而且据我所知,至今也无人能够……)

\dmark{}\origpage{881}在那份Bourbaki报告里,我仅限于给出以普通$\ell$-adic层为“系数”的L函数公式的一般表述,并展示如何通过非常简单的拆卸,把问题归结为~$X$~为光滑射影曲线的情形。我心里清楚:一旦到达那一步,\textbf{就赢了}——因为对维数一的把握已经足够充分,使该公式的证明成为例行公事\fnote{492}{(*)}{我在此谈“例行工作”时,丝毫没有贬义。数学工作的十分之九,如果说不是更多的话,属于这一类——在我身上如此,在任何一位数学家的身上也如此,只要他有时机经历一些时刻,而那些时刻,恰恰是别的东西,是创造的时刻。在Verdier之后,我自己也曾花时间摇动那些可用技术的曲柄——精细而上了油的——去找到并证明一条令我满意(至少暂时如此)的一维不动点公式。那也像Verdier所做过的,是“例行工作”。}。我当时并未着手去提炼一维情形的好不动点公式并证明它;我觉得那更该轮到Verdier来干。他在次年给出了所谓“伍兹霍尔”不动点公式,足以覆盖Frobenius态射及L函数的应用。我读了他的表述,并不真正满意:因为他对其上同调对应所施加的条件(出于其证明的需要——该证明我未曾过目)在我看来有些人工——我倒希望有一条公式适用于代数曲线的\textbf{任何}自同态。SGA~5讨论班成了发展一条合我口味的这种公式的第一个好时机。(若无疏漏,那正是Illusie版第XII讲中奇迹般幸存下来的那条公式——尽管那场不幸的讨论班风波不断……)Weil诸猜想曾是初始动机与宝贵线索,引我“投身”于平展上同调(及其他)完整形式主义的发展。但我分明感到:这个上同调主题——它已在我工作的中心待了八九个年头,并将持续到1970年我离开——其意蕴比把我引到它跟前来的Weil诸猜想还要辽远。对我而言,Frobenius自同态不是上同调形式主义的“阿尔法与欧米伽”,而不过是众多自同态中的一员……

\origpage{882}依我看,Deligne“SGA~4$\frac{1}{2}$—SGA~5”行动的初始动机\dmark{}曾是\textbf{唯独}侵吞迹公式这一件,从而(作为“推论”)侵吞L函数。是路走到一半,这一图谋才扩大为对“平展上同调”整体的侵吞。我并且相信:这两个“份额”都太大;而且时至今日,尽管有“SGA~4$\frac{1}{2}$”、反常层会议以及诸如此类,“人们”(甚至那些消息并不怎么灵通的人)也都“知道”:创建$\ell$-adic上同调工具的并不是他,而且他也没有单凭一己之力证明“\textbf{那条}”Weil猜想。

尽管如此,为了给“平展上同调”这场行动收尾,我还想在这里稍稍追随我的朋友兼昔日学生Deligne的几番辗转——看他如何呈现那部名为“SGA~4$\frac{1}{2}$”卷册的“中心主题”\fnote{493}{(*)}{况且,在“SGA~4$\frac{1}{2}$”中没有任何一处说“报告”构成了“中心主题”,也没有一处说其主要宗旨是为“那条”Weil猜想提供平展上同调的主要成分。在为该卷撰写双重导言之际,一份规模及于整个平展与$\ell$-adic上同调的侵吞意图,想必已然在场。},即“那条”迹公式,由之导出L函数的上同调公式。它就是“迹公式报告”(在其书中被引作[报告],前引书,第76--109页)的主题。

Deligne在该卷中足足有\textbf{四个}地方,对迹公式作了多少带点“历史”味道的评注。该卷的读者,倘若事先并不知情——无论他读还是不读那四个段落(我们即将逐一检视)——从中得到的印象都将是:某位Grothendieck(一个有点马虎、且晚于“SGA~4$\frac{1}{2}$”一卷的讨论班的作者或主持者——人们尤其劝人千万别贸然去读它)关于L函数似乎也曾有过一些想法,当然难免有点纠缠不清,直到那部辉煌卷册的作者终于出面,给出可理解的陈述与站得住脚的证明。

在整部书中,对这位先生唯一准确的引证,是指向某一场Bourbaki报告(1964年),在一句“评注3.7。”(前引书,第88页)的转角处;而它出现在那里,是作为最后一棒,收在一串三条愈来愈技术的评注末尾\fnote{494}{(**)}{写下这些文字时,我正被一种强烈的“同一性”感觉所震撼:我在此探测的文风,与四年后为“以轻蔑”侵吞“好上帝定理”(别名Mebkhout)而展开的那一种。这些翻云覆雨,我是在“变戏法者”那条注里发现的(这名字当得起大写……),n\textsuperscript{o}~75。在那里,“敏感点”藏在一条评注4.1.9里(而不是3.7),更是乱得可以。进步是挡不住的……(3月22日)我不曾注意到,事实上在“SGA~4$\frac{1}{2}$”里还有对1974年那同一场Bourbaki报告的第二处指涉——一处以精湛的艺术奉上的指涉,就在“阿里阿德涅之线”里,一如将在子注“双重含义——或行骗的艺术”(n\textsuperscript{o}~169)中所见。}

其文如下:

\begin{quote}
“若承认$\mathbb{Q}_\ell$-层的形式主义……\textbf{很容易把证明归结}到$X_0$为光滑曲线、$F_0$光滑的情形。这在[2]§5中解释得很清楚(对3.1而言;3.2同样处理)。”
\end{quote}

\origpage{883}\dmark{}(着重号是我加的。)总之,这位除了以谀媚的代号[2]\fnote{495}{(*)}{轮到谁就是谁——1970年(在尼斯国际大会上),是Serre(在Deligne的报告“霍奇理论I”里)得到了代号[3]的待遇,而不被提名——就在那句神谕般的话里,人们头一回也是最后一回向所呈现的理论之若干“源头”示警……}之外再不获点名的先生,(诚然没有做出,但)解释了那点平凡的活儿——平凡到几乎不配在这条收尾评注里提一句,而且还好心让人以为:平凡归平凡,这里至少解释得很清楚。(我们从那位辉煌作者另一些评注中已经知道:清晰可不是这位含混先生的强项……)换言之:这一章“迹公式报告”的宗旨在于\textbf{做真正的工作},把那些平凡的边角留给干这个的人……

既然说到这里,不妨马上指出:正是在这一页上,有着我方才提及的四个段落之一——它们含有关于“那条”迹公式的历史性评注。那是第3.8节(紧随3.7的评注之后,合乎规矩)。在那里人们解释:证明3.2(即迹公式在该卷中唯一被显式讨论的情形——Frobenius对应的特殊情形)有“两种方法”。不用说,那位先生的名字在两者中都没有出现。人们区分方法A,称为“Lefschetz--Verdier”方法,和方法B,称为“Nielsen--Wecken”方法(这个名字对我而言也意味着什么……)。让我们看它怎么说:

“方法B:Nielsen--Wecken。一种受Nielsen--Wecken工作启发的方法,允许把3-2[即Frobenius的迹公式]归结为Weil证明过的一个特例;这将在以下诸段中解释。”

\origpage{884}事实上,第5节(pp.~100--106)的标题,恰如其分地,叫“\dmark{}\textbf{Nielsen--Wecken方法}”。前文说过该方法\textbf{受启发于}Nielsen--Wecken的工作——所以那位卷册作者把这方法冠名为“Nielsen--Wecken”的,想必纯粹出于谦虚。这就更清楚了,因为他们不是当下的人。如果读者碰巧翻阅了某份报告XII——文中从未引用它,而且是在一个人们建议他遗忘的讨论班中——的文献目录,他会知道:那是两位在四十年代初发表著作的人。如果他甚至读了他们的精美著作(那位辉煌作者,我敢打赌,从未亲手翻过),他会知道:他们的方法是三角剖分技术。那显然不是正文里用的方法。在没有相反声明的情形下,那卷册的谦逊作者,便同时也是该方法的作者。该方法没有标注日期——无疑又是出于谦虚,免得人们知道:果真是他第一个辛辛苦苦证明了这条出名的迹公式。

再看方法A,所谓“Lefschetz--Verdier”方法,人们怎么说它。这可不怎么鼓舞人心:

“若~$X_0$~是真的……Lefschetz--Verdier一般迹公式允许把3.2的右端表示为局部项之和,每个~$X^{p^n}$~的点一项。在SGA~5的\textbf{原始版本}中,这一公式仅在奇点消解的模意义下被证明[人们早就猜到会碰上麻烦!]。读者将在\textbf{最终版本}中找到无条件的证明[又谦虚得忘了提醒:正是多亏他,这一步才得救——反正人们是决计不会去翻那该死的SGA~5的]。对于曲线的情形(可以归结到该情形,见3.7),成分[???——我们放弃了……]反正都已可用。”

可如果它们已经可用(一位比其他人更清醒的读者,如果有的话,也许会问),为什么还要这样长篇大论地谈什么“仅在……模意义下被证明”的Lefschetz--Verdier公式等等等等?刚才不是说了:\textbf{真正}的工作在维数\textbf{一}中进行吗?答案:这就是所谓“乌贼战术”——喷墨汁,好在浑水里摸鱼!读者到了这一步,已经完全确信这方法肯定不是对的那一个。他以一双倦眼扫过下一段,那段给他补上了最后一击:

“为从Lefschetz--Verdier公式推导3.2,须能计算其局部项[天哪,这是哪门子苦役……!]。对一条曲线与Frobenius自同态[啊!他们泄气了!],这已由Artin与Verdier完成[还要两个人一起干!](见J.-L.\origpage{885}Verdier,\dmark{}平展上同调中的Lefschetz不动点定理,载于Local~Fields会议文集,Driebergen,Springer~Verlag~1967)以及SGA~5的\textbf{最终版本}[人们不免好奇:那原始版本到底长什么样,可怜的我们!]。”(此处及上文着重号皆为我所加,纯属恶意!)

显然是发善心,那位辉煌作者才免去指向那部注定被遗忘的讨论班的相应报告,或者哪怕暗示一声“那条”公式确确实实就在那里!一位不弃不馁而好奇的读者,若碰巧去那里面翻找,会发现一份报告XII,顶着不寻常的标题“代数几何中的Lefschetz与\textbf{Nielsen--Wecken}公式,A.~Grothendieck著[还是那位先生,我发誓!],I.~Bucur整理[不认识]”。那位先生和他的伙伴,八成是把他们辉煌前辈的报告抄了一遍,再兴致勃勃地堆上多余的细节……

在那份出名的“报告”里,没有任何东西能让读者想到:除了Lefschetz--Verdier公式(或者毋宁说该叫Lefschetz--Verdier-Deligne公式,反正也不怎么鼓舞人心——那位作者自己泄气的评注表明了这一点)之外,还存在一份显式的迹公式,适用于\textbf{任何}自同态,而非仅限于Frobenius。无论在援引Artin-Verdier的那一段,还是在(下面要引的)指向SGA~5的另一段(特别是不提名那位先生)里,都被暗示:这工作\textbf{仅}在Frobenius自同态的情形下被做过。人们跟Verdier是哥们儿(而且也向他证明了这一点),但迹公式嘛,这是明摆着的:给Verdier一个“暂停!”式引用可以(与Artin并提\fnote{496}{(*)}{我早已领教过Deligne这套久经考验的技艺:把水搅浑、把鱼淹死,好掉包某个某人(此处是Verdier,毕竟是哥们儿,而且人家将在别处获得实质性补偿)——办法是把他与另一个人一口气并列提名,这样就没法指责他不慷慨了!这便是所谓“以同化作稀释”的掉包法。这方法中的艺术,在于找到那位与待掉包者“配对”的先生。就我而言,我的朋友回回都找Serre……},淹没在一大段技术性的、索然无味的文字中间,读过即忘)——但毫无疑问而且没有弄错:迹公式,那是\textbf{他,Deligne}!

诚然,Deligne不止一根弦,把这些带着“历史”架势(sic!)的评注散布在\dmark{}\origpage{886}四个不同的地方,实有深意——好在这一处补上在另一处可能被人指摘遗漏(或夸大)的东西。那里他还可以退守到同一章的引言——一切都已经安排好了!那是一段七行的引言,值得一字不漏地引用\fnote{497}{(*)}{(3月20日)我在昨天的反思中回到这段导言。(参见子注“双重含义——或行骗的艺术”,n\textsuperscript{o}~169$_7$。)}。

\subsubsection*{(b) ……与无稽之谈}
\tocsubsub{(b) ……与无稽之谈}

\notehead{169$_6$}{}

(3月18日)昨天不得不在兴头上停笔,因为时候已实在太晚,而且显然当夜写不完“公式”了!在重新接上围绕那份公式的几番辗转之前,我想先借这个机会,就“Lefschetz--Verdier”这条优美的公式,稍许“把脚伸进菜碟里”。这条公式以一种完美的例证方式,阐明了一样在我看来至关紧要的东西——在《收获与播种》全程乃至《引言》中\fnote{498}{(**)}{见《引言》4,“一场追寻显然之物的旅程”。},我不止一次回过头来坚持谈过、但措辞或许始终有些过于“一般”的东西。

这条公式是一份\textbf{深刻}陈述的触目例证,而其证明是“平凡的”(169$_6$bis)。当Verdier告诉我,他提炼并证明了一条Lefschetz公式——针对任意(“真的”,当然)代数簇上的“上同调对应”(当时连定义都还没有!)与任意可构“系数”——我起初难以置信。或许“带有或多或少一般系数的Lefschetz公式”这一念头曾掠过我的脑海——至少我很久以前就为“局部常数”系数(即局部系统)写过这样一条。但\textbf{我不相信}对一般系数也行——那看上去太美好了,好得不像真的!Verdier说服我没花多久。当场写下公式并给我证明,多半花了十五分钟——还得算我慢,尤其涉及要确信某桩如此出乎意料的事情的时候!这正是人们可以称之为一“\textbf{平凡}证明”的东西——就“众所周知”的事物而言。而顺着当今刮的风(那阵风,J.H.C.~Whitehead早已嗅到了最初几缕\dmark{}\fnote{499}{(***)}{关于这一点,见“年轻人的势利——或纯洁的捍卫者”一注(n\textsuperscript{o}~27)。}),只需一步(大多数人轻快地迈出了这一步),就能把定理本身也归为“平凡”——十条或百条公式中的一条,从上同调形式主义里“自己掉下来”的\dmark{}\origpage{887}——此处的形式主义,是那个我在此前一年(1963)刚在平展框架中发展的\textbf{完整}体系:六种运算,\textbf{以及}双对偶定理。

我说Verdier(循Lefschetz开辟的路)发现的定理“深刻”,绝非因为(尽管这理由也成立)其证明所从出的那套形式主义本身“深刻”。那股时风,早就(而且恰恰有Verdier本人的无条件支持!)把这形式主义归档为“Grothendieck式的大厚片”,挥手即扫,同时又暗中使用那些“厚片”于每一步(而不点名)。至于该定理“仍属猜想”(某某先生以怜悯的口吻强调),或在一切特征中都已完全确立(它如今确实如此——多亏恰恰是顶着这位某某先生名字的“双对偶定理”)这个问题,在我看来,当我说它是一条深刻的定理时,同样无关宏旨——它实质性地丰富了我们对于各种“上同调主题”的理解(离散或连续系数,各种“簇”或“空间”……)。同样的,对普通Lefschetz公式——比如紧(或无论如何)微分流形及其有孤立不动点的自同态——也可这样说:“形式”证明,从上同调对偶形式主义出发,一页足矣,几行亦可。然而在两个情形里,都有过\textbf{创造}——某种崭新的、实质性的东西,它此前逃过了所有人的眼睛,它“还不存在”,而骤然间现身了……

在当前的情形里,“创造”确切地定位于何处?我相信,不止一位数学家——尤其不止一位曾是我的学生的人,他们当年本知道什么是创造,却早已把它忘到九霄云外——应当在这个案例上,或在与他们更近的某个类似案例上,好好沉思一番\fnote{500}{(*)}{当年“在圈内”的人并不多(而且如今也不多,鉴于事态的走向……)——但除了Verdier和我,总该还有三四个。Deligne嘛,他还没在这一带露面……}。我知道,如果当年我向我本人、或向某位当时在“上同调形式主义”行内的学生或同事,提议显式地陈述一条一般Lefschetz公式——\origpage{888}针对任意系数与任意“上同调对应”(定义由他们临时去下!)——在一张紧(抱歉,是真的)任意\dmark{}簇上,每一个人都会万无一失地到达:花上几个小时或几天,或至多几星期\fnote{501}{(*)}{当然,我在此假定所论的那个人确已对提出的问题“咬住了钩”,亦即我对此曾有过的“感觉”(若非如此,我也不会把它提出来!)确已“传了过去”,而那位学生或同事确确实实“发动”了。一件事要“传得过去”,可绝不是显而易见的——远非如此!}。一旦问题被提出(哪怕还相当模糊,而主要术语尚待定义……)并被\textbf{看见},去“解决”它(在此即:找到\textbf{正确的}表述,由已有的上同调形式主义所暗示)便成了一桩“\textbf{例行}”的事(Weil在同一意义上称之为“习题”)。这种“例行工作”确实需要嗅觉、最低限度的智力与想象力,但(一如我不止一次写过的)此时“是事物本身在向我们口授”如何着手——只要我们肯倾听它们。(而如果我们不会倾听数学事物,我们倒不如改行……)我所说的那道\textbf{火花}——它迸发出崭新之物——并不在这一类工作中\fnote{502}{(**)}{而火花就更不会迸发于那样一件打杂的工作里——也许是十年之后做的——它只不过证明:使某条证明“行得通”的诸项假设,确确实实在人们以为的地方得到了验证……}。

创造的时刻——触发发现进程的那道火花——在此处,是\textbf{问题被看见}并加以\textbf{“承担”}的那一刻:当意图从\textbf{真正去看}、从\textbf{追到底}、从弄个水落石出以“看清”这条人人自以为“理解”的Lefschetz公式之\textbf{确切}的“真实”有效域中诞生。迸出火花的,不是某种“精湛技艺”或“强大脑力”(通常意义上驾驭复杂技术或记忆嵌套局面的那种“力量”)。那是一份\textbf{天真}:人人都自以为理解Lefschetz公式,而我,可怜的人,尚无已理解之感,我想把这事弄个水落石出!在这一类情形中,一旦上路,\textbf{便已赢了}:事物向我们吹拂该做的事,我们去做。“追到底”在一种情形里,可以意味着证明“那个”正确的定理(在此处以已有的形式主义为项——该形式主义本身是“已确立的”还是“仍属猜想的”,无关宏旨)。在另一种情形,它可以\dmark{}\origpage{889}意味着:提炼出“\textbf{正确的}猜想”\fnote{503}{(*)}{两种情形——其一是“火花”(被“追到底”)把我们引向一条定理,另一是我们被引向一条猜想——并非天性不同。“追到底”意味着:使一份尚且弥散的直觉完全具体化,从它的所有侧面、用我们掌握的一切手段去探测它。一条定理在本性上并不比一条猜想更“完成”。有种种显然暂时的定理(甚至跛脚而歪斜的),也有种种猜想(例如Weil猜想全体)给人一种全然完成、完美无缺的印象。这并不妨碍这些Weil猜想曾是通向其他发展(起初和它们一样是猜想的)的起点,而那些发展更宏阔,并把它们涵盖在内。在这个意义上,可以说数学中没有任何东西,只要它还生气勃勃,是“完成了的”或“定局的”。};而这猜想常常自身是临时的,或许日后会被发现为假或不充分,须加以调整或拓宽——这同样无关宏旨。这条猜想是通往对事物(在此即数学事物)更深切的集体认识之路上的一个台阶,一个无法绕过的台阶\fnote{504}{(**)}{关于发现的动力学,以及“错误”在其中的关键角色,见(在《收获与播种》第一部分的)“错误与发现”一节(n\textsuperscript{o}~2)。}

深刻与丰饶是紧密相连的品质——后者在我看来是前者有形的标记。Verdier所发现的公式之丰饶的第一个标志,当年(如果不是在几天或几周之内的话——我说不准了)就到来了:这条公式曾是\textbf{主要}动机,引我写出以任意$\ell$-adic层“为系数”的L函数上同调公式。在\textbf{技术}上,我根本没有用到Lefschetz--Verdier公式——这一点无关紧要。确定无疑的是:没有这条公式作为引路线头,或者不如说:没有它向我吹拂“应当有什么东西”——一桩以层“为系数”的L函数——没有这声执拗的声音,我根本不会想到去提炼\textbf{正确的}概念及与之相伴的公式;诚然我日后无疑也会到达,但那须得\textbf{先}靠自己的力量发现另一条更具一般性、恰好“挡在路上”、\textbf{必须}被发现的公式。

在心理层面上,两种情境极为相似。正如Verdier须首先提炼“上同调对应”的概念,以精确化“Lefschetz公式问题”(超出“普通”公式之外),同样,我也须提炼“以层为系数的L函数”的概念,以\dmark{}精确化“L函数公式问题”(言下之意:\origpage{890}超出与真的、光滑的~$X$~相关联的“普通”L函数……)。“创造的时刻”——火花闪过的一刻——是我\textbf{看见这问题}的那一刻:定义这种广义L函数——而且我\textbf{承担了}它,把那个问题进行到底。一旦问题被清楚地看见,而且假定我能使“上同调形式主义”行内的随便哪个周围的人“接住”它,那么显然他必将情不自禁地解决它——以唯一自然而合理的方式,花上几天工夫(一如对我而言必曾如此),定义、陈述、证明,一气呵成\fnote{505}{(*)}{我在此把证明的最后一步单独挑出来——我当年把它搁置了(因为它想必不构成真正的问题),而它险些变得更长。}。诚然,那些归结到维数一的“拆卸”是“容易的”,甚至(如果你非要这么说)“平凡的”。发现并不在这类拆卸之中——那种随便谁来都做得跟我一样好(或不屑去做)的拆卸。\textbf{发现}在于一个\textbf{概念}:一个此前无人想到过的概念,尽管它是\textbf{显然的}:以层“为系数”的L函数。在这个概念及与之不可分的公式里,蕴含着(在素域$\mathbb{F}_p$~上有限型概形的语境中,或更一般地,在绝对基环$\mathbb{Z}$~上)把上同调中的“六种运算”——从函子$Rf_!$开始(即“\textbf{几何}”性质的操作)——解释为“L函数场”上的操作,即以“\textbf{算术}”为项,加以诠释的可能。这是1949年Weil猜想所开启的方向上的新一步:经由上同调主题,朝几何与算术之联姻迈进。

这两项发现在这份自诩为平展上同调与$\ell$-adic上同调\textbf{标准}参考书的文本——出自我曾有过的学生中最有天赋、最有名望者的手笔——里,变成了什么呢?

Lefschetz--Verdier公式——它启发了我,而我从未“使用”过它——变成了\textbf{稻草人}:被适时地高高举起,好让读者(他只求相信!)领会:某个讨论班悬在一根何其纤细而乏味(还是“猜想性的”,更不要说“局部项未被计算”)\origpage{891}的线上,而那个讨论班(“按照本卷的精神”)人们善意地绝口不以指涉(除非是为了贬损它……);同时不忘在此处彼处轻描淡写地提醒:如果那条不巧的(而且说到底毫无用处的)公式毕竟已不再“属于猜想”,那得归功于那部辉煌卷册的谦逊作者。

至于以层为系数的L函数的概念——它是构成全书核心的那份“报告”的中心概念——它在“报告”第1.6节(loc.~cit.~p.~80)悄无声息地登场,没有半句评注指示其动机或来历。定义嘛终究是定义,无须辩护。对这概念的来历——不得不说它有点装神弄鬼(尤其当你囫囵吞下的时候……)——心存疑问的读者,可以在Artin-Weil(但他们那个时代尚无$\ell$-adic层,它们显然是由作者在同一卷中\dmark{}引入的……)与(更可能的)这位辉煌作者之间做选择:后者正在利落地把他引向那条著名的“迹公式”。

后者在第3节(loc.~cit.~p.~86)被引入,该节以这样的字句开篇:“Grothendieck对L函数的\textbf{上同调解释}是如下定理:……”(接下来是公式3.1——注意:着重号是我加的)。除了那一章引言(下文还会回到它)之外,这是全章唯一一处\fnote{506}{(*)}{(4月9日)有一处例外(起初逃过了我的注意),带着一处“暂停!”式指涉(在第90页),指向“Grothendieck对导出范畴理论的一处本质运用”(用以在“非正统”情形中定义迹)。}提到某个名字的地方。因此还是那位先生,两页之后以缩写[2]被指涉(作为一个“清楚地解释了”若干“容易的约化”的人),他也给出了那条装神弄鬼的“解释”3.1——冷不防地抛在那里。这反正也没什么了不起,读者马上(而且毫不意外地)会看到:证明占不上半页(同在p.~86),而且本来就是“经典的”:它是那条给“报告”命名的出名的“\textbf{迹公式}”——即显然构成“真正定理”(3.2)的东西——的简单推论。后者的发明权,没有任何名字被提及——依旧是这份谦虚的癖好,偏偏在那些最辉煌的人身上!两页之后(如昨天所见\dmark{}\origpage{892}),Lefschetz、Verdier、Artin、Nielsen、Wecken的名字纷纷出笼——真是一场谦虚的豪宴——就是为了不说出是他!

我在这里想强调的、在我看来远远超出本案与这股骗术气息的事情,是这个。无论对那称之为(名副其实)“Lefschetz--Verdier”的公式,还是对L函数(“以层为系数”)的“上同调解释”,恰恰是\textbf{这一点}——即它们的发现是\textbf{创造之举}——在今天成为普遍低估的对象(即便不是轻慢的嘲弄),人们习惯以带贬义色彩的形容词来表达,如“\textbf{平凡}”“\textbf{幼稚}”“\textbf{显然}”“\textbf{容易}”“\textbf{猜想性}”,甚至“软数学”“梦想”“大话”“无稽之谈”,以及诸如此类的好听话,听凭各人即兴挥洒。而工作的另一部分——我始终知道(尤其,我觉得,从未\textbf{忘记})它是“额外的”、由事物的力量尾随而来的,如同“后勤”必然追随(只要人们肯躬身去做)——即那\textbf{技术}部分,它常被冠以“\textbf{困难}”之名,要“凭腕力”来做,我也曾称之为“例行工作”(不带任何贬义)——恰恰是工作的这一部分,被当今的共识所推崇,被捧到了排斥其余一切的地步。

对我而言,“困难”的概念全然相对:一件事在我尚未理解它时显得“困难”。我的工作便不是“凭腕力克服”困难,而是深入我的不理解之中,达到某种理解,把曾显得“困难”的东西变得“容易”\fnote{507}{(*)}{读者将注意到,这是对一项困难的“阴”的、“女性”进路——“上涨的海”那条路——的一番描述。我在此并不是说这是唯一可能的创造进路——还有那条所谓“锤子与凿子”的进路,“男性”的进路,唯一享尊荣(且不说今日是唯一被容许……)的一条。关于这两种可能的进路,见“上涨的海……”一注(n\textsuperscript{o}~122);关于时下对两种进路各自通行的态度,见“肌肉与肚肠(阳葬阴(1))”与“天赐的机缘——或封神”两注(n\textsuperscript{o}s~106、151),以及承后者而写的“否认(1)——或提醒”(n\textsuperscript{o}~152)。}。例如,我为“L函数公式”以及在其他场合所做的那些拆卸——它们今天被当作“平凡的”\dmark{}——对我而言并不比处理那些被称作“困难的”“不可约”情形更“容易”\origpage{893}。它们是工作的不同阶段,仅此而已\fnote{508}{(*)}{我所想到的情形,除“带系数的”函数L的一般公式之外,是我做过“拆卸”以回到维数(或相对维数)\textbf{一}之情形的地方,主要就是平展上同调中的两条基变换定理(对一真态射,以及经由一光滑态射)——它们构成那两条使该上同调“可活”(Deligne语)的关键陈述;还有$Rf_!$的“比较定理”,在平展上同调与超越上同调之间(对复域上有限型的概形)。(还有仿射态射的(所谓“弱”)Lefschetz定理。)从心理上说,一旦得以回到这类所谓“不可约”的情形,我就感到(多少)“赢了”,期待的定理必将“出来”,而经验在每一个这样的场合都确证了这感觉没有骗我。然而从技术上说,正是这些拆卸代表着那所谓“容易”的一步。凑巧的是一种当时令我甚为触动的“天意”:处理两条基变换定理各自的两个“不可约”情形所需的配料,都已由我发展出来(当时毫无自觉),为了第一条在SGA~1里,为了第二条在SGA~2里,早三年和早两年……}。某一阶段在另一阶段之后,或比它更长,并不意味着它更“困难”。在两个阶段里,都须\textbf{一份想法}:此处的“拆卸”的想法(这在以前从未有人想到要在这种情形下做——而且理由充分,因为那涉及针对任何非Frobenius对应的不动点公式!);彼处则是一个无疑更精微的想法,受一条比Lefschetz原始公式更精致的(出自Nielsen--Wecken\fnote{509}{(**)}{(4月10日)况且,Deligne正是和我、以及SGA~5的其他听众一样,从我口中得知这条“Nielsen--Wecken”公式及其在平展上同调中的转写的——这使他无须再去翻那三位作者的(德文)漂亮文章(发表于1941至1943年间),而后来的用处,其别致一如人们所知(见子注“真正的数学……”,n\textsuperscript{o}~169$_5$)。}的)不动点公式启发,通过对系数层的精心剖分来实施,以合适的导出范畴为项加以表达\fnote{510}{(***)}{在这项证明中,导出范畴的语言是不可或缺的。在我离开之后,直到名为SGA~4$\frac{1}{2}$的那一卷出版之年左右,我的上同调学生们对导出范畴建立了一种默契而有效的抵制——而它们曾是发展对偶形式主义(“六种运算”与双对偶)之概念关键工具的,在“凝聚”系数、继而“离散”系数的语境中。尽管它在Lefschetz--Verdier公式的证明中、同样也在平展语境中诸“经典”对偶公式的证明里扮演着关键角色,这一形式主义本身,作为一种数学结构、一个融贯的概念整体,却遭到了同样的抵制,而且抵制延续至今(且不说“六种运算”这个名字本身——它始终是个忌讳)。有可能是迹公式证明的种种需要,促使Deligne于1977年朝解除对导出范畴的抵制迈出了第一步:在海盗卷中出土了Verdier“论文”的一副骷髅般的“第0状态”文本(文中对我的名字只字未提)。关于这一点,见专论“行动III”的注“分赃”(n\textsuperscript{o}~170);至于“论文”这桩滑稽事件的详情,见“同伙”与“赊账的学位论文与全险”两注(n\textsuperscript{o}~63$'''$、81)。}\dmark{}。\origpage{894}第二个阶段确实更长:当须以它应得的全部一般性将其完善时(鉴于曲线上的自同态不止Frobenius一种),出现了整整一“地毯”的“斯坦伯格式”非交换迹,它们最终粘合在一起,须我细心发展\fnote{511}{(*)}{(4月23日)一般性被Illusie在其SGA~5屠杀版导言(第二段)中按“理”判为“多余”,同时毕恭毕敬地回声他那显赫的朋友Deligne——后者(不作进一步说明)指涉他会“剪除”的“无用细节”。这番贬损同时也一劳永逸地使他免于让读者疑心:在维数一存在一条比他为Frobenius所呈示的更一般的显式迹公式,其中他步步重演我的证明,却给人印象以为那是他的独创。见下一条子注“双重含义——或行骗的艺术”。}。那很漫长,也很“容易”——而它也是一桩\textbf{必须}做的事,这一点我十分清楚。但即便找到这类使工作“变容易”(或干脆使之成为可能……)的想法,在我看也属于“例行工作”的一部分。它给这工作平添一份魅力——那通常使它超越单纯摇动手柄的东西。

工作的\textbf{创造}性部分,却是那个\textbf{幼稚}的想法:那个人人早该看见的想法——若不是几年,就是几个世纪、几千年以来——而偏偏无人看见,虽然它一直在众目睽睽之下,人们每回都不得不绕一个大弯以免撞上它!

当你遭遇到这样一个想法——无论是你自己“撞上了它”(这可是实至名归的说法……),还是别人向你解释了它(如Verdier有一天向我解释了一个)——你觉得自己蠢透了:无法原谅自己此前竟没有看见它,当它恰恰是最自然、最显然、最“傻”的那一个——恕我直说……人们\textbf{早该}撞上它,这是肯定的;可偏偏没有……

看起来,当今(尤其当人们占着强势的时候……)在这种情形下,人们以圆滑来补偿\dmark{}:当\textbf{另一个人}(一位显赫的无名氏,或某个\origpage{895}早已入土的“亡者”……)不巧带着这样一个想法凑上前来(或某天曾凑上来……)。“可是我可怜的朋友,你跟我讲的这些\textbf{太平凡了}!为了向那不幸的人证明这有多平凡(顺便把他放回原位……),人们只消用说完一句话的工夫就把它重新做出来——你这就见识见识什么叫做数学!我们袖子里可比这些初出茅庐的(或被冷落一旁的……)多的是东西!只需稍加拉拽,呼——再一拉——阿布拉卡达布拉,好嘞!而且\textbf{瞧},一条像模像样的陈述就从帽子里出来了,甚至一整套理论,一点不掺水,这才是干活呐,没错!年轻人,回去穿好衣服,等你学会来这一套再回来吧!”

我在这里不假思索地,把我的“身后学生”Zoghman~Mebkhout的遭遇做了个速写:他当年不过是里尔(或天知道哪里)的一名微贱助教,落在我的“隐形学生”Pierre~Deligne——一所精选机构的头号台柱(其余且不论……)——的手中;那段遭遇发生在恩典之年1981年,而且延续至今……这便是“第四行动”,所谓“当值无名之辈”行动(或“反常层会议”,为免提名)——四场行动中最令人难以置信的一场。(关于它,见“封神”那条注,n\textsuperscript{o}~171.)

但同时,在写下前一段时,我隐约觉得在或多或少重写某段我已在另一场合写过的东西……

没过多久我就想起来了——那是在《收获与播种》第一部分,写于如今整整一年前,在“竞技数学”一节(名字不言自明)里,n\textsuperscript{o}~40(p.~105)。我在那里唤起的插曲与反常层会议插曲的区别在于:这一回,“当值无名之辈”的角色,由“那个踩着我花坛的小白脸”扮演,而那位高傲而“竞技”的大东家,不是我的某个可鄙前学生,而是我自己。诚然,我不认为自己走到了(象征性地,在此场合)把别人的想法据为己有的地步。但我无法凭良心起誓,而且须由当事人(二十年之后——但迟做总比不做好)告诉我\textbf{他}如何回忆那段插曲,它在我的记忆中有些模糊。他遭遇的不幸是:重新做了一些我早已烂熟的东西(其中包括经由“拆卸”从约化情形构造非约化概形的Picard概形……),而那“弄得不愉快”——这就是\dmark{}留给我的印象;但我不会起誓说他的进路(在一个不如我一般的框架里,这是肯定的)当真完全被我的进路所覆盖\fnote{512}{(*)}{我从不曾有机会把所说的皮卡概形“相对”构造(通过在幂零理想上“拆卸”)整理成稿并发表——那原定为《代数几何原理》后续一章的内容(而该章从未问世)。况且,当我谈对他人一个想法(无论大小)的“侵吞”时,未必是指通常意义上的抄袭,即把这个想法(哪怕是经过修改和完善的形式)奉上而不指明其来历——在我看来,这事正变得越来越平常。但侵吞也可以是那种漫不经心的轻蔑之侵吞,它的气息使发现的喜悦枯萎,仿佛只为挫败它而取乐,配上一副厌倦的“哦!只不过是这个……”的腔调。那腔调不言而喻地让人懂:我们刚刚被告知的东西,人家早就知道,不如说,从来就知道;而如果还不曾费心把它讲明白,那实在是因为它值不得费这个神……对于这类腔调(或其先祖),见(在《收获与播种》第一部分的)“挫败人的权力”一节(n\textsuperscript{o}~31)(在已引的“竞技数学”一注(n\textsuperscript{o}~40)中重提);以及(在70与80年代更严酷的氛围中)《葬礼I》中的“侵吞与轻蔑”(注n\textsuperscript{o}~59$'$)。}。

不管怎样,我必须在此再一次确认一份\textbf{亲缘}:它存在于六十年代至少在某些时刻属于我的那种态度,与我在我某些昔日学生身上所见到的那种态度之间。他们向我送回的那个人(从前的我)的一副确实被扭曲了的形象——一幅\origpage{896}多年来我一直想拒斥的形象。但如果《收获与播种》——它首先是对我数学过去的反思——曾有过\textbf{一个意义},那便是让我同样理解:即便我某些昔日学生乐于否认我,我也不该否认他们中的任何一个。经由他们回送到我这里的东西,属于我曾参与播种的那些收获,一如他们本身也参与了播种。而这份我以一支毫不留情的笔、在即将满三周之际做出的确认,不是对任何人的控诉,而恰恰是一份\textbf{确认}——它把我自己卷入得与任何人一样深。

\subsubsection*{(c) 遗产——或骗术与创造}
\tocsubsub{(c) 遗产——或骗术与创造}

\notehead{169$_6$bis}{}

(4月10日)\fnote{513}{(**)}{本子注源出前一条子注“……与无稽之谈”(n\textsuperscript{o}~169$_6$)的一条脚注;见第886页转引。}如众所周知,“平凡”一词在数学中的意义极为相对。在此,所谓“平凡”,我指的是:以人们理当“已知”的东西为项——即(在此情形)六种运算的形式主义,以及双对偶定理(后者在离散平展语境的特征~$p>0$~中,在Deligne找到证明之前仍属猜想……)。以这套形式主义为项,证明的原理可以在几分钟内(与陈述一起)以全然令人信服的方式\dmark{}\origpage{897}解释清楚。这当然不免于写一份正式的证明,那意味着验证若干烦冗的相容性。

在这种情形下的惯例,是定理的作者(尤其当它重要时)肯费心写下一份证明。在Verdier的情形,我毫不怀疑:这是他的全部成果中最深刻、就其意义而言最重要的一条——在他有幸(而且在此情形下理所当然)冠以自己名字的成果之中(借用Weil的惯用语)。然而他对待这条定理,一如对待导出范畴理论:既然荣誉横竖都能到手,他便不认为值得去做这项工作、把一份完整的证明交到所有人手中。

这是一个富有雄辩性的心态之标志——我在此彼处谈到过它,最近还在“操演”(n\textsuperscript{o}~169)那条注(2月28日)的末尾。我得以确认它已蔚然成风。诚然,那被称为(带上文保留)“Lefschetz--Verdier”的公式,确确实实是Verdier的一项创造之举——在他尚与我共事、尚对工作满怀热情的岁月。我看到两件事之间的直接关联:他从未尊重到去证明“他的”定理,以及\textbf{他的数学家生涯中再未有过类似的创造之举}。创造的时刻,唯有在我们“配得上它们”的时刻——亦即:处在接纳它们的状态——才降临我们……

这条被一位逃之夭夭的父亲弃置不顾的优美公式,此后还经历了奇特的坎坷。它先是成为我在SGA~5中最早报告之一(第III讲,1965年)的主题。Illusie承担了它的整理工作,却在十二年间不认为值得费这份工夫。而后,在他与Deligne的完美同谋中(我猜,还至少得到了Verdier的默认——Deligne将向他支付实质性的补偿),它变成了“特洛伊木马”(或如我下文所写的“稻草人”)的头颅,被熟练地操控,以使那场令人难以置信的冒名顶替——名为“SGA~4$\frac{1}{2}$”——显得可信。后者被凭空搭建,旨在埋葬他们三人共同的老师——也就是说,归根结底,该公式的“\textbf{祖父}”(那条公式,若非有我这微不足道的人与随我一同被埋葬的六种运算,恐怕再过一百年也写不出来……)。作为风俗画,这才是风俗画!

如果我那些亲爱的上同调前学生们,不是把自己挥霍在这类骗术里扮演矮子(他们不是),骑在一个巨人的\dmark{}\origpage{898}肩膀上(我也不是),而是——在这十五年里——让存在于他们身上的创造力自由发挥,一如让存在于我身上的创造力自由发挥,那么晶体系数、De Rham-Mebkhout 与 Hodge-Deligne诸理论,以及以“神秘函子”为钥匙的那一套,肯定早已抵达六种运算形式主义的“完全成熟阶段”了。而且(我有一两个星期以来就怀疑这一点……),那位曾是他们的师父的人的宏大梦想——那个“Motif”,本是用来成为旋律的,却在同样这几只手里变成了采邑、宝藏与“模糊的骨架”——它也早已化身为一首宏大的交响曲(绝非“猜想性的”,而同样是“完全成熟的”),并从今天起便成为\textbf{全人类的共同财富}。

\subsubsection*{(d) 双重含义——或行骗的艺术}

\tocsubsub{(d) 双重含义——或行骗的艺术}

\notehead{169$_7$}{}

(3月19日)但我必须回到我朋友Pierre~Deligne在介绍著名的“迹公式”时的“翻云覆雨”。咄咄怪事的是,他从不曾指明:就Weil猜想本身的应用而言(那想必是从实用角度被摆在第一位的,即便不是唯一的位置),根本不需要什么精致的公式与证明——“寻常的”Lefschetz公式(平展版)就足够了\fnote{514}{(*)}{(4月25日)我在此可能出错,因为我还不曾真正了解Deligne对Weil猜想最后一部分的证明,那关乎 Frobenius 特征值的绝对值。看来,Lefschetz束的使用引导他引入比函数$\zeta$(即“寻常”$L$函数)更一般的$L$函数。}。而以下这点当然并非偶然:他偏偏选中了关于“结合于一个闭链的上同调类”的那场报告,把它从 SGA~5 中“借”来,不加任何手续就并入他的摘要——而那场报告恰恰含有(除了“寻常”版的Poincaré对偶,平展版之外)建立“寻常”Lefschetz公式的关键配料——只消几汤匙的工夫。于是人们不禁要说:他本满可以免掉收入这份不三不四的“报告”,它不过为单独的 Frobenius 自同态建立了一条迹公式(同时还执拗地向读者隐瞒:他可以在别处(!)找到明显更一般、而且同样“显式”的公式)。可他既然费心写了这份“报告”,那想必出于两个相互牵连的理由。一方面,六十年代就已十分清楚:Weil猜想一经以“权”为项妥善重新表述,便对奇异簇与非常数“系数”仍保有意义。的确,人们从此可以把它们全然以几何的措辞来表述,无须明确提及\origpage{899}$L$\dmark{}函数的形式主义。我想,Deligne的文章“Weil猜想~II”里所做的正是如此(其中当然丝毫不曾暗示:对于他在其中所证明的那个主要陈述的提炼,我曾出过什么力)。然而,几何-上同调运算的算术诠释(以“带系数的”$L$函数为项)看来仍必定要扮演一个角色,而在那里,以我所发展的形态出现的一般$L$函数公式,将取得一个关键的位置。从长远的眼光看,那就必须在那部被命名为“SGA~4$\frac{1}{2}$”的书卷里提供一份引证。同时,鉴于一般迹公式(Lefschetz--Verdier式)已然构成上同调装备的一项重要配料,这也有助于造成那样的幻觉:这部书卷(如他所宣称的)确乎提供了一套基本完备的上同调装备,以满足“非专业的”$\ell$-adic 上同调使用者的需要。

我还得逐一检视那剩下的三段文字——在“SGA~4$\frac{1}{2}$”中装作就迹公式给出若干历史说明的段落共四段。我将按它们在书卷中出现的次序引用。前两段位于整部书卷的开头(《引言》第1页,以及“阿里阿德涅之线”第2页),其用意显然在于“摆明态度”。它们想必也是被读得最多的两段。第三段是“迹公式报告”一章的简短引言。(第四段前天已经谈到过,它属于同一份报告的正文,想必是所有四段中最少被读的一段。)

在那份置于“SGA~4、SGA~4$\frac{1}{2}$、SGA~5 之阿里阿德涅之线”之后的文献目录里,SGA 这个缩写被解释为“布瓦玛丽的代数几何讨论班”,而不消说丝毫不曾提及我个人。可我毕竟列名于 SGA~4、SGA~5 的主持人之列。况且,主持人的职务想必纯粹是柏拉图式的:把 SGA~4 与 SGA~5 的主要报告过一遍(就再也别谈它们了……),其中谈到Artin、Jouanolou、Houzel、Bucur的报告,却没有一篇是我的。在关于 SGA~4 与 SGA~5 的条目里,没有任何日期——我在整部书卷中不曾找到任何一处暗示,能让未事先知情的读者怀疑:SGA~5(“即将由 Lecture Notes 出版”)如其名义所示,确确实实是\textbf{晚于}那部名为“SGA~4$\frac{1}{2}$”的书卷的\fnote{515}{(*)}{竟没有只言片语能让读者猜到这个“不许阅读的讨论班”讲的是什么,连它的标题(“$\ell$-adic 上同调\textbf{与函数}$L$”)都无人知晓!}。万一对 SGA~5 中某场报告有所提及(通常并不指明是哪一场),反倒会被“明确”\dmark{}\origpage{900}说明那是“零状态”或“原始版本”(言下之意:臃肿而读不通,可想而知……)。这些对 SGA~5 的指涉(是给一位未获知情的读者的,而人们劝他既不要查 SGA~4,尤其不要查 SGA~5),在这位读者的心中便是指向一部\textbf{晚于}他手中正在读的这一部的文本。况且我怀疑,这类未获知情的读者占绝大多数,而其余的人(如我在别处所写)开始老去,并将体面地死去……

我引用《引言》第3段的第一页:

\begin{quote}
“‘迹公式报告’含有对 Frobenius 自同态之迹公式的一份完整证明。该证明是Grothendieck在 SGA~5 中所给出的那一条,剔除了一切无用细节。这份报告当使使用者得以忘掉 SGA~5;人们尽可以把后者看作一系列离题,其中某些很有意思。\textbf{它的存在将使人不久能以原样出版 SGA~5}。”(着重号是我加的。)
\end{quote}

这篇文字有\textbf{两个相反的意思},以精湛的艺术\textbf{同时}奉上。对于了解该公式(对 Frobenius 而言)历史的人,他也许会惊讶于这番满不在乎的呈献(而若他充分了解 SGA~5 讨论班的来龙去脉、以及它在培养那位辉煌而放浪的作者时所扮演的角色,他会加倍惊讶);但他会以为,作者至少指明了其证明的出处。而对于未获知情的读者,他得知:他手中这部书卷里的证明,也见于一部较晚的文本 SGA~5——那文本出自Grothendieck之手,塞满了无用细节,是那位先生在原证之外为取乐而添上的。所引的这段对原证本身语焉不详。如前天所见,在那份“报告”里亲自读那证明,几乎不容人怀疑“SGA~4$\frac{1}{2}$”这部书卷辉煌的作者就是它的生父。当然,人们到处都不屑于指明:写出迹公式这个念头是谁起的;说到底,只要不肯费心去证明,写点什么是不费分文的!也丝毫不提Verdier(是他头一个给出了我当年搁下的“关键情形”的证明)。无疑并非偶然:正当谈到迹公式、正当谈到“那条”猜想的核心时,作者才一味殷勤,满口“无用细节”“离题”(当然很有意思,人要么表现得有风度,要么不\dmark{}那样!)\origpage{901}劝人忘掉,最后又添上一句既含蓄又专断的提醒——“它的存在将使人不久能以原样出版 SGA~5”,仿佛 SGA~5 之所以能“立得住”、之所以可出版,全靠那部名为“SGA~4$\frac{1}{2}$”的文本的“存在”——而它想必正好为那位先生提供了所需的一切,好把在原初文本中以简单方式做出的东西,以复杂的方式呈献出来,如下文所示。

% ---------------- p.712 ----------------
在那份“阿里阿德涅之线”里,我早已(在“操演”那条注(n\textsuperscript{o}~169)的子注“特洛伊木马”(n\textsuperscript{o}~169$_3$)里)指出第2页那两个相连段落(第2、3段)共十七行文字,堪称“浑水摸鱼”之术的“典范”。第二段恰恰关乎那条著名的迹公式。这两段值得在此全文照录:

\begin{quote}
“平展上同调中存在着一种与凝聚对偶类似的对偶形式主义。为建立它,Grothendieck使用了奇点消解与纯粹性猜想(陈述见 [Cycle] 2.1.4),后者在 SGA~4 XVI 中于相对框架下、在 SGA~4 XIX 中于等特征并模去奇点消解的情形下得到确立。关键之点则在 [Th. finitude] 中以另一种方法确立,适用于正则零维或一维基概形上有限型的概形。若干发展见于 SGA~5 I.~在 SGA~5 III 中,人们展示这一形式主义如何蕴含那条极一般的Lefschetz--Verdier公式。

可见,在 SGA~5 的原始版本中,Lefschetz--Verdier公式只是以猜想方式确立的。此外,局部项在那里也不曾计算。为应用于$L$函数,该讨论班含有另一条证明,它在 Frobenius 态射的特殊情形下补足了前者。那就是见于 [Rapport] 的那一条。其他引证:关于陈述与拆卸之概览,见Grothendieck的Bourbaki报告 [5];关于(由Grothendieck做出的)把关键情形归结为Weil已处理过的一类情形的简要描述,见 [2] 第10节;关于对后一情形的$\ell$-adic 处理,见 [Cycle] 第3节。”
\end{quote}

第一段我已经在那条注里评注过了(另见该注第872页的脚注(**),关于那句千金难买的“若干发展见于 SGA I”)。如今且容我追随我的朋友在第二段里的翻云覆雨(或者说其中的某些——太多了,追不过来!)\origpage{902}。那\fnote{516}{(*)}{更确切地说,他让人清楚地以为:单单那份34页的“报告”就包含了(而且是更好地包含)SGA~5 中一切可能有用的东西(后者即便在屠杀版中也还有近500页)。好一堆“离题”,原来都是白费的!}\dmark{}前两句话,一遍又一遍地念叨那条永恒的Lefschetz--Verdier公式,仿佛整个 SGA~5(以及其中某条从未被明白道出姓名的证明,为了某条迹公式……)都死死依赖着它、终生依赖着它,显然属于“乌贼战术”之列:把清楚的事情搞乱,好浑水摸鱼\fnote{517}{(*)}{说Lefschetz--Verdier公式曾是“猜想性的”,这并不恰当。它是建立在“具备一套对偶形式主义(‘六种运算’与双对偶定理)”这一假设之上的,并且于1964年由Verdier在这一形态下确凿地获得了证明。该证明当然已在口头讨论班中给出,而且是完全的。在特征~$p>0$~中仍然“猜想性的”,是双对偶定理的有效性,而它(如我们所说)是在“SGA~4$\frac{1}{2}$”的“有限性”一章中得到确立的。

至于Lefschetz--Verdier公式的局部项,它们被“计算”的程度,与寻常Lefschetz公式(不动点未必“横截”的孤立不动点)不相上下,并且推广了后者中出现的经典“相交重数”。说这些项“不曾被计算”,与说一个向量空间(其定义未被\textbf{指明})的维数、或一个系数未定的多项式的根“不曾被计算”,含义同样多寡。“计算”,在这些场合一如在别处,意味着:在一个被指明的“具体情形”中(例如对Lefschetz--Verdier公式取维数一)确立两项之间的一项\textbf{等式},而两项中没有哪一项比另一项更“被计算”或更已知(例如,在Verdier定义的局部项、与某些结合于 Artin 导子的局部不变量之间……)}。

而那句一语双关的要紧话,却紧接在“淹鱼”之后:

“……该讨论班含有另一条证明,它在 Frobenius 态射的特殊情形下补足了前者。”

那位获知内情却又性急的读者(而哪位读者不性急呢……)有一瞬间被“该讨论班”这个说法的含糊所困扰——究竟是 SGA~5 呢,还是“SGA~4$\frac{1}{2}$”?——而由于他知道在 SGA~5 里确有一份完整的证明,便又一次这样裁定了:作者的确(以颇为含糊的方式)指涉了它。我去年4月头一遍读时(见“白板”一注,n\textsuperscript{o}~67)差点也这样上当,可到底对不上。我明明知道,我所给出的那条显式迹公式的证明,绝不限于“Frobenius 态射的特殊情形”。而且令我震惊的是:人们刚刚还在(用些莫名其妙的“论证”)强调某一种 SGA~5 的呈献(在它的“原始版本”里,拜托!)是“不完整”的——这边猜想性,那边局部项没算出来……有了这个被两个逗号夹住的“补足了”,这种断然的对照便不可抗拒地向未获知情的读者暗示:无须他自问,“该讨论班”显然\textbf{正是}\dmark{}他捧在\origpage{903}手里的“SGA~4$\frac{1}{2}$”这部书卷——而且人们紧接着就在下一句里告诉他去哪儿找:“那就是见于 [Rapport] 的那一条。”而事后去读所引那一章里的该证明,当然也绝不可能在这位读者心中引起丝毫怀疑\fnote{518}{(*)}{见前天子注“真正的数学……”(n\textsuperscript{o}~169$_5$)。}!

% ---------------- p.713 ----------------
要紧处中\textbf{那个}“另一”字被\textbf{着重}了,而这全然不是我朋友的惯常做法。这是两段引言文字中唯一被着重的词,而且如无差错,也是整部书卷中唯一的一个(标题、陈述与新引入的术语除外)。他如此刻意要使这个词跳出来,想必不会没有缘故。(只是到此刻,这件事才引起了我的注意。)这个“另一”的效力——尤其当它被这样推到前台时——在于强调:“那条”公式曾有\textbf{两种}证明:其一是\textbf{不完整}的,而人们刚刚才就这条“Lefschetz--Verdier”公式那绝不出色的处境说了几句(实在拿不出手!);(而在前天看过的著名的《报告》那段更技术的文字里,人们又恪尽职守地就这个凄惨的话题再度施压……)至于那条跛脚的办法究竟有没有多亏那位辉煌作者的有限性结果而终究走得通,谁也永远不会知道这事最后怎么收场。但在这个“反差”效果(归根结底与前天检视的是同一个)之后,驯顺的读者那心理反射便越发专断:与其要那种在某个含糊讨论班 SGA~5 里的\textbf{不完整}办法(它如此不完整,以至于连一份确切的指涉都压根谈不上\fnote{519}{(**)}{在名为“SGA~4$\frac{1}{2}$”的书卷中,我从未找到任何一处指及 SGA~5 中某场报告,其中或含有一条不动点公式的证明,或含有什么$L$函数的著名“上同调理论”。事实上(见下文)已经明确说过:“本着本书卷的精神,将不诉诸 SGA~5……”!}),那种办法人们当然永远乐得避之不及,倒不如在这个\textbf{讨论班}里得到那份好而扎实的证明,那份\textbf{完整}的证明,它已在专门为此而设的报告——“迹公式报告”——里向人们张开双臂,准没错,人们毫不费力就能找着……\fnote{520}{(***)}{最美妙的是:事实上,Deligne的证明正是他1966年在 SGA~5 讨论班上与其他听众一起学到的那一条的忠实复现。}。

那句“该讨论班”简直是神来之笔——我那处于“暂停”状态的朋友在这个词上是攻不破的。然而,无论是在所引那段文字里,还是\dmark{}在“报告”中那段更技术的上下文里(第88页,论及那种注定被遗忘的、所谓“Lefschetz--Verdier”的方法),他都还是又一次冒险\textbf{明明白白}地说出了(或者说,至少半明半暗地说出了)\fnote{521}{(*)}{“又一次”,因为他已在前一段里(这一回甚至说得更明白)“说了假话”,一如我们在子注“特洛伊木马”(n\textsuperscript{o}~169$_3$)中所见。}\textbf{假话}。在两处他都强调(这回倒用得着这个词了)有这么一种办法(人们猜得出,就是 SGA~5 里那条不幸遵循的办法,天知道是在它哪一场“臃肿”的报告里……)用来证明 Frobenius 的迹公式,无非是\textbf{使用Lefschetz--Verdier公式}。然而,除了阿利贝尔1982年的论文(它在维数一情形,对一条任意的、不动点孤立的上同调对应给出局部项的计算)之外,“关键情形”只有两种证明:Verdier的和我的,而其中没有一种(阿利贝尔的也一样)使用Lefschetz--Verdier公式!这原是一桩棘手的事,而且长期未获解决(如今看去也有点旁支):证明 SGA~5 中所提炼的那条显式公式里出现的局部项(对应于远比 Frobenius 对应一般的对应关系),确确实实就是Lefschetz--Verdier公式的那些局部项。Illusie终于把它验明了,照他在 SGA~5 屠杀版导言(第~VI~页)以及他的报告~IIIB“局部项计算”(第~139~页)的导言中所宣布的\fnote{522}{(**)}{至于Illusie这些突然努力的动机,见子注“恭贺——或新风格”(n\textsuperscript{o}~169$_9$),尤其是第916--918页。}。

% ---------------- p.714 ----------------
\origpage{904}如果说Deligne还是不辞辛劳地制造这虚假的印象,那并非没有缘由。因为由此他同时造成这样的印象:SGA~5(那个“技术性离题”的讨论班,“本着本书卷的精神将不指及它”,其用途就是让人“忘掉”它)依赖于那条“猜想性的”公式——况且按其原样也无法使用(局部项没算出来, sic……)——而它之最终得以确立,全赖Deligne在那部名字再明白不过的“SGA~4$\frac{1}{2}$”(读者正捧在手里)里做的事;而(即便只凭这一点)那个较晚的、“含糊的”讨论班 SGA~5 也就依赖于它……

至于所引那段最后一句(以“其他引证”(sic)开头),它同样是这一类把戏的典范,为的是不说那位含糊的Grothendieck先生早已给出过一份\origpage{905}完整的证明——早了整整十一\dmark{}年(在那个“较晚的”、注定被遗忘的讨论班里……)——并且该证明被忠实地复现于“报告”之中。必须造成的印象是:那位先生只做过一些初步的模糊约化,而困难的情形归功于Weil,并由作者以(一次“$\ell$-adic 处理”)辉煌地接手。引到一本读者或许听说过的Weil的权威著作,再加上一条内部引证,这很能增添分量——要么严谨、要么熟悉自己的经典,二者必居其一!偏偏巧得很,对Weil那本书的指涉没有任何日期,也没有章节或页码——那位辉煌作者似乎并不想鼓励读者到辉煌的书卷本身之外去翻检,而一涉及它,引证顿时变得再精确不过(章节、段落全有)。

那条有名的“Weil已处理过的结果”,说白了不过是代数曲线(代数闭域上连通光滑射影)情形下的\textbf{寻常}Lefschetz公式——Weil在四十年代以边缘的办法就能表述并证明它,那时还不曾掌握上同调工具(却用了雅可比簇来定义所缺的$\ell$-adic $H^1$)。在“抽象”代数几何的语境中把这条公式提炼出来,当时原是一个重要的新想法,而且想必正是它把Weil引向了他那著名的猜想。一旦掌握了上同调形式主义,所说的这条Lefschetz公式便大体上是平凡的了。但若是明白说出那位先生的约化就是约化到寻常Lefschetz公式(为此人们骄傲地、不提名字地指到辉煌书卷的“Cycle”一章——那一章恰好是从 SGA~5 抢来的……),那就会给人一个印象:连那条“约化”索性就是那神圣公式的一项\textbf{证明}了。您可别想这样!(*)

我可急着收尾了!剩下的就是“迹公式报告”那一章的引言,前引书第76页,照录如下(略去其最后两行,那指涉书卷作者的一篇概述性文章):

\begin{quote}
“在本文中,我试图尽可能直接地陈述Grothendieck的$L$函数上同调理论。我十分贴近地追随Grothendieck1966年春在 IHES 所给出的若干报告。本着本书卷的精神,将不诉诸 SGA~5——除了两处指及该讨论班报告~XV 中独立于其余部分的段落。”
\end{quote}

乍一看,人会觉得作者毫不隐讳地指出了他的来源,谈到\origpage{906}“Grothendieck\fnote{523}{(*)}{(5月11日)由此可见,这里全部的“暂停!”艺术,就在于在两个相隔遥远的地方(第2页与第88页)指到\textbf{两}处“约化”(!)(容易的,这不用说)都是那位先生做的(只被点名一次,第二次就没有了……),却不让天真的读者有丝毫可能疑心:同一位先生居然\textbf{发现了}并\textbf{证明了}一条迹公式;而且他的证明(注定被遗忘的)被忠实地复现于辉煌的“报告”之中……}的$L$\dmark{}函数上同调理论”,还补充说他“十分贴近地追随”我的若干报告。在\textbf{一部寻常的}书卷里,这没有什么可说的。但同样真切的是,\textbf{上下文}参与构成任何文本的意义。“SGA~4$\frac{1}{2}$”这部不寻常的书卷,深刻地改变了这段话的意义——对于一个天真的读者,他事先已被先前读到的内容提醒过,而事后在读《报告》本身的过程中还会更受教益。到那时他会得到这样的印象:那位慷慨的作者把那名叫Grothendieck的含糊先生记上一笔“$L$函数上同调理论”,实在是天大的善意——而那理论归根结底似乎不过是一个有点神神道道的“诠释”,而且说到底还是\textbf{平凡的}。它只需半页光景就得了证,作为一条“迹公式”的直接\textbf{推论},而那条迹公式可一点也不寒碜,而且除了辉煌书卷的那位太过谦虚的作者之外,谁也配不上说是它的父亲。

% ---------------- p.715 ----------------
诚然,文中说了:在他的“报告”里,作者“十分贴近地追随”那位先生1966年春在 IHES 给出的若干报告。关于这些必定臃肿不堪的报告,再没有多说一个字,除了书卷作者肯为其报告保留的部分之外,它们想必早已彻底失传。究竟是些关于 Frobenius 的呓语(为此人们还会慷慨地指到那个同一位先生“主持”的 SGA~5),还是些关于$\ell$-adic 层的一般议论,抑或是某些下文将谈到的“容易的约化”——全然一团模糊。无论如何,那些想必多半是些“无用细节”,而读《报告》正好为我们免此苦役,谢天谢地——人们要求的也不过多。于是,给那位先生盖上一层贞洁的薄纱,大家且干\textbf{真正的活}去!

而当涉及对那位先生的指涉时,我的朋友喜欢含糊其辞(只要他不是干脆不提),这一回却似乎不能怪他不精确了:在 IHES 给出的报告,1966年春。他若是再精确一丁点,就会添上一句:是 SGA~5 讨论班\textbf{上}的报告。

SGA~5?那不就是那个讨论班吗——它(不着日期)出现在“阿里阿德涅之线”的文献目录里,还注明“即将由 Lecture Notes 出版”?也就是说,人们(想必)已经明白,那个讨论班的所作所为,就是在它前面那个 SGA~4$\frac{1}{2}$(那才真是无可指摘)之上,再添些“离题”(其中某些很有意思,这我们同意)和些“无用细节”。别瞎扯了,SGA~5 哪是什么1966年春的事,您别逗了!而最好的证据就在您眼前,\dmark{}白纸黑字印在“迹公式报告”(Pierre~Deligne著)那篇\origpage{907}刚刚引过的引言里:

\begin{quote}
“本着本书卷的精神,\textbf{将不诉诸} SGA~5”。
\end{quote}

那么,这就清楚了,不是吗?!

\subsubsection*{(e) 变戏法者——或飞走的公式}

\tocsubsub{(e) 变戏法者——或飞走的公式}

\notehead{169$_8$}{}

(3月20日)我开始感到有些疲倦,若不说精疲力竭的话,为我已逾三个星期的这项工作——尤其是(逐细地)最近这几天——孜孜不倦地“拆解”——正是从那些构成一切的“小小无物”入手——我那位最有才华的学生的天才骗局拼凑,好让那些只求被唬弄的人(而他们想必人多势众……)在公共场合上受唬弄。我急着收尾,是的;可我并不后悔为此花去的时间,尽管我年近五十七,而更有意思(或至少更“可喜”)的事情并不缺少。这有点像数学里我称之(想必就在三天前)为“例行工作”的那种活儿——做着的时候牙咬得痒,明明知道这一切不过是粮草,可同时也清楚地知道它非做不可!并非出于什么严峻的“义务”或责任,而是因为一个人不能(至少只有我不能免俗)省掉这一步,倘若我想与所探测的东西建立亲密的接触、在其中“深入”。正是靠了这样的工作,日复一日、周复一周乃至年复一年地与被我们想认识的事物“摩擦”,我们才真正“认识”它们——而唯有出自这份认识(它是一场往往艰苦、其貌不扬的工作的果实),才有\textbf{别的东西}时而迸发,就是我前天谈到的“火花”,它骤然更新我们对事物的领会,以及这引领我们深入其中的工作本身。

正是凭着这般疲倦(它还不是厌倦,而是一份能量耗去的标志),我才得以同样充分地估量:我的朋友皮埃尔在这场微妙的、编排停当的行动(“SGA~4$\frac{1}{2}$”行动,或称“SGA~4$\frac{1}{2}$—SGA~5”行动)中必定耗费了多少惊人的能量。我说不出这艺术家的工作——它比数学家的活儿不知微妙多少,动用的官能全然另一层次——在多大程度上是有意识的,抑或全然是些无意识力量的作为。而这终究是个旁支的问题,只关乎他一个人。无论如何,这能量的分流,以及这投入一项与发现之冲动截然相反的任务——掘墓人-变戏法者的任务——的强度,想必一直庞大,\dmark{}而且(这于我毫无疑问)今天依然如此\fnote{524}{(*)}{这种落在“那条公式”上的侵吞执念,用简单明白的话说,真是疯了。一方面,这种侵吞由于事物的力量,势必大体上(若非完全地)停留在象征层面:一种自我给予的满足,仿佛自己\textbf{果然就是}“那位父亲”,或者仿佛自己\textbf{果然能够让}全世界都相信这一点。它的虚构性、象征性,只要回想一下Deligne本人在那部“SGA~4$\frac{1}{2}$—SGA~5”的骗局出现之前四年发表的文章《Weil猜想I》中所写的(第278页)“Grothendieck证明了Lefschetz公式”(对 Frobenius 对应而言),就已暴露无遗。诚然,仅仅几个月之后,在1974年2月的Bourbaki报告(n\textsuperscript{o}~446)中,Serre呈献Deligne这篇文章时,还(有理地)诧异于Lefschetz公式竟无一份已发表的证明(“人们自1966年以来就等着 SGA~5 的定稿版,它本该比现有的油印讲义更令人信服”),并趁机对 SGA~4 那1583页(它们“带着一切必要的细节、以及许许多多别的”)铺陈平展上同调形式主义的文字冷嘲热讽。Serre当然不会料到,这些针对一位缺席者的讥讽不会落到聋子耳朵里。我深信,它们必定在“让人忘掉”这一天才念头的萌芽中起了作用——忘掉 SGA~4 与 SGA~5 那一“废话矿渣”等等,正如公众借Serre本人的嘴似乎正在要求的……但即便抛开Weil猜想I不谈,仅就已出版的文本(包括 SGA~5 的屠杀版,它虽遭残损却仍是一份确凿的见证……)而言,父权的抹杀根本站不住,以最基本的数学常识也说不通。

此外,如我已强调过的,那条著名公式的制定是一份\textbf{纯粹例行}的工作,只要人知道自己想得到什么。我当年花了几天才把它的主要特征提炼出来——那把我引向一些与 Artin 导子相关的精确可除性问题,而Serre对这些问题早有现成答案,并且能用 Serre-Swan 模的语言优雅地表达。那段略长(但同样例行)的工作,是把受斯塔林斯工作启发(而他这篇东西恰好刚寄到我手里)的非交换迹形式主义仔细地理顺。这一切,都是那种有Deligne那份干劲(或者只有我这份更寒碜些的干劲)的人一年之内成打处理的事情!

诚然,在Deligne笔下,“迹公式”意味着\textbf{任意维数下} Frobenius 对应的迹公式,而他费心(在“SGA~4$\frac{1}{2}$”中)把它同他所谓的$L$函数的“上同调诠释”(“Grothendieck的”,谢谢!)区分开来。他把后者仅仅当作迹公式的一个\textbf{推论}呈献。(事实上,本着1964年我那场Bourbaki报告的精神,两条公式对我曾是\textbf{同义的},是“算术”与“几何”之间同一关系的等价表达式,其一是加性的、其二是乘性的。)

由此可见,在这场围绕“那条公式”的执念背后那真正的动机(固然还只是浮面的),全然不在于上同调装备,而在于尽可能缩小、乃至彻底抹去我个人在“\textbf{那条}”猜想的证明中扮演过角色这一事实。终究是她,在我看来(直到1981年6月的反常层会议为止),作为那场在我昔日的学生围绕被否认的师父身上打成的结的\textbf{重大}纠结点显现。}。这些侵吞-掉包的反射\origpage{908}——至少在与我的工作相关方面\dmark{}——以及对待任何公然打着我的印记的其他工作\origpage{909},终于在(那场漫长的“攀登”,即“亡故师父之葬礼”的过程中)获取了对他整个存在的如此支配,以至于它们成了仿佛第二天性,侵凌并覆盖了他的原初本性,那个“孩子”的本性,那奔向着发现世界的孩子……我不止一次在看似平常的处境中(其规模远不能与我刚刚稍稍细看过的“平展上同调”这样一场“行动”相比)就近看到这些反射无声的效能,它们在和颜悦色的天真外表下以完美的自如运作。还没等你明白过来发生了什么事(倘若你真能明白过来……),他已经占有了你在欢欣中创造的东西,而先是吹上一口含蓄而阴险的轻蔑的气息,使它枯萎。(同样真切的是,有这种气息的绝非他一人,远非如此;而这种气息今天看来很成了时令之气的一部分……)

% ---------------- p.716 ----------------
但是,这气息枯萎了另一个人所创造之物的美、枯萎了他的喜悦,也就枯萎了\textbf{一切}事物之美,以及那创造之力本身——它在他身上一如在我们每个人身上,使我们得以与事物相通、深刻地认识它。诚然,这并不妨碍做出些“困难的”事、受人羡慕、忌恨与畏惧。但他身上所承载的那部作品——我近年还见过它的若干先兆——至今还没有诞生。它将在某一天诞生(如果那天终将破晓),当某种东西已然崩塌,而当那位端坐于上的主-奴,一如他那被否认的师父曾经所做的那样,成了一个\textbf{仆人}。

如今这已是满满六十页的一厚摞了(还不算一大沓脚注!),将近三星期的劳作,我全数奉献给了“平展上同调”这一场行动。它无疑是所有行动中最浩大的一场,即便不是最“肥”的(那一场将在最后检视,就在那道名字当之无愧的注“封神”里……)。我意识到,饶是如此,我还没能把它的里里外外看全。一针引一针,这场预想中的“整理”——在一场“调查”中把“事实更新”——反倒让调查重新启动,引得我去稍稍细看那部不寻常的书卷“SGA~4$\frac{1}{2}$”,而从前我只在跑动中看过它一眼。这也是个机会,让我重看一遍那本阴郁的 SGA~5 Illusie版,而且目光比从前更清醒。如今我看出那两个盗匪之间有一丝不苟的默契:Illusie把自己整个交给Deligne调遣,好呈献一部全然\dmark{}\origpage{910}合乎那位显赫的庇护人与朋友心意的 SGA~5 版本。这番对 SGA~5 的呈献,是以低八度的音调,回应那份在电锯文本中铺展的贬损与轻蔑之气,并为其中设下的骗局提供含蓄而有效的支援。

屠杀版的导言通篇都是为了在未获提醒的读者心中造成这样的印象:这是一部“技术性离题”的书卷,叠在貌似\textbf{中心的}、而且\textbf{较早的}(!)文本“SGA~4$\frac{1}{2}$”之上。而在Illusie执笔的那些报告中,这一印象又因大量指及那部海盗文本而加强——每逢他用到他那位朋友认为宜于收入其摘要的成果,他都慷慨地指向它,哪怕在同一部 SGA~5 里、甚至早在 SGA~4 里就有“量身定制”的指涉\fnote{525}{(*)}{因此,紧支 Künneth 公式(在任意基概形之上)是(对导出范畴而言使用的)真态射基变换定理的直接推论,而该定理曾是平展上同调中第一个重大“突破”,时在1963年2月。它正以此身份列在 SGA~4 的那“废话矿渣”里——人们可不愿让Illusie指到它,既然有那部中心的文本(其用途恰恰是让人忘掉这些含糊的前身……)在向他张臂相迎……}。

我是在去年5月12日那道同名之注“屠杀”(n\textsuperscript{o}~87)以及它的子注中,逐渐发现一场不折不扣的屠戮的现实的。在那组注里,我终于着手对那场拆解做一番详尽的描述(虽说还不是穷尽的)——它在过去两周间渐渐显现在我眼前。由于我当时不曾像近三星期这样,把所谓“SGA~4$\frac{1}{2}$”里围绕“那条”公式精心设置的骗术一一拆开细看,去年我还不曾看清这个侧面:在 SGA~5 的Illusie版的整体呈献中,竟有这样一丝不苟的串通。为给“平展上同调”行动(别名“SGA~4$\frac{1}{2}$—SGA~5”)收尾,我还得说明这场串通在“那条”(不动点)公式的呈献中是何以表现的,给几个细节。

% ---------------- p.718 ----------------
我此前已(在子注“好撒玛利亚人”,n\textsuperscript{o}~169$_2$ 中)指出:Illusie在他的导言里帮着朋友唱和,好让人以为 SGA~5 的出版取决于那条永恒的Lefschetz--Verdier公式的证明。(该证明自1964年起就可获得,而我当然早已在口头讨论班里把它发展出来了;而Illusie虽在1965年承担了整理工作,十二年来却始终不认为值得守他的约……)\dmark{}我同样提醒:去年(在那条已引的注“屠杀”,n\textsuperscript{o}~87 里)我就已经发现了\origpage{911}原始讨论班报告~XI 的若干坎坷。那场报告与紧接其后的报告~XII 不可分割——后者发展了我那版(直到1981年为止所知最好的)一维Lefschetz公式——而它在Illusie版中完全消失了。据Illusie的导言所说,那场报告想必只是“Grothendieck的交换迹理论”(对“\textbf{非}交换迹”而言是天赐的笔误!)“推广了斯塔林斯的(非交换迹)理论”,而且是(同样天赐地)在一次搬家中消失的(!!)。事实上,那场报告发展的是必不可少的代数预备,为的是在下一场报告中描述局部项;在那里,我发展出局部项计算(或不如说\textbf{定义})的一般方法(经由一条“Nielsen--Wecken”式的公式\fnote{526}{(*)}{这条公式被Deligne侵吞了(丝毫不提我个人),连同那种把Nielsen--Wecken公式过渡到常系数(即“寻常的”)系数层的不动点公式的途径——而可构系数层则是任意的。关于这一点,见子注“真正的数学……”(n\textsuperscript{o}~169$_5$,第883--884页)。于是(身价使然……)同一个Deligne小心翼翼地丝毫不提那场“较晚的”讨论班 SGA~5 的报告~XII——那里“Nielsen--Wecken”这个名字正写在报告的标题里(“代数几何中的Nielsen--Wecken公式与Lefschetz公式”)。}),及其在一维情形下的显式应用(若我没记错,是靠了 Serre-Swan 模)。总之,Illusie用一场“新的”报告~IIIB 来“顶替”那场消失的原始报告~XI,并应景地称之为“局部项计算”(这想必也是,如无差错而且巧得出奇,正是被抹掉的那场报告的标题!),\textbf{而他自居其父}。就此他一石二鸟。一方面,这是一项\textbf{残毁}之举,初看仿佛无谓,却以这种\textbf{生硬的割裂}四处播撒混乱\fnote{527}{(**)}{这种残毁与混乱,在我的昔日学生Illusie奉我另一位昔日学生Deligne之命播撒的众多残毁之中,使后者得以居高临下地谈论 SGA~5 的“混乱状态”(“虽说严谨”,因为人到底是好牌手……),而“SGA~4$\frac{1}{2}$”(它再早不过!)据说应当对此“加以补救”……而这一切,都在那群忠诚信徒的慈祥目光下。恭喜!}——把一场报告从它的自然脉络中硬拽出来,在原地留下一个张着的窟窿,只为图个乐子把它塞到别处。这或许是绵密细致的Illusie强加于那部曾经辉煌的讨论班(而他自感骤然成了它的绝对主人……)的一切残毁之中,事后在我看来最暴力、最赤裸裸地招摇的一桩:我能白屠,我就屠——且带着与我的良好教养相称的细腻。恭喜了,Illusie,为这一类活儿,那可不是你跟我学的,而是跟另一个人学的,你把他奉为楷模与导师……

\origpage{912}\dmark{}这是一。而这第二下,同一块石头砸得漂亮,Illusie竟能\textbf{抹掉}我在1965年提炼出的那条不动点公式的\textbf{父权},同时又(尤其是)把\textbf{那条公式本身也抹掉}。自1965/66年起,它就是“\textbf{那条}”\textbf{一维情形下的好的不动点公式},远比Verdier前一年在伍兹霍尔发展的那条更一般(否则我也不必劳那个神),更遑论Deligne那部著名的“报告”里的那条(后者只限于 Frobenius 对应,却步步照搬我在一般情形下提炼出的证明)。直到几年前(晚了将近二十年)它才在阿利贝尔的论文\fnote{528}{(*)}{这篇论文是在Verdier的指导下准备的(没错,永远是那个Verdier),1981或1982年在蒙彼利埃通过(我手头没有确切引证)。它是十年艰苦劳作、显然颇为凄清的结局……}中得到改进,头一次处理了系数层上任意的上同调对应。

% ---------------- p.719 ----------------
Illusie设法把文本编排得使有关公式几乎无处可寻:在两场相互撕裂的报告(带官样文章的 IIIB(sic)与 XII)的技术岩浆中,没有任何东西(在这一篇或那一篇的引言里,也没有在别处)把读者的注意力引向这个构成两场报告之整体的核心成果、引向整个讨论班最重要的成果之一\fnote{529}{(**)}{技术上,正是\textbf{那条}关键公式(“不可约情形”)使得证明那条著名的“$L$函数公式”成为可能,后者等价于(任意维数下)Frobenius 对应的迹公式。这条公式的关键角色,早已由 SGA~5 这个讨论班的名字本身作证(那个名字在那部所谓“较早的”“SGA~4$\frac{1}{2}$”文本里从不曾被提及):“$\ell$-adic 上同调与$L$函数”。}!我甚至承认,我没能绝对确凿地弄清这条公式究竟在不在 SGA~5 里。鉴于文本故意造成的混乱,又鉴于我远离这个题材,我得花上几个小时甚至几天才能找着。真正让我困惑的,是全文丝毫不提 Serre-Swan 模——而(如果我的记忆没骗我)正是它们赋予了我提炼出的那条公式以优雅与概念上的简洁\fnote{530}{(***)}{很可能,在这里以及下一句里,我把Euler--Poincaré公式(见于报告~X)的结构与(报告~XII)那条Lefschetz公式的结构混为一谈了。在Bucur的报告(复述我的口头报告)中所呈现的那种形态的Euler--Poincaré公式里,Serre-Swan 模确实是显式登场的。}。而正是为了这条公式的需要,Serre才做了几场关于结合于 Artin 导子的Galois模的精彩报告——那些报告本来当然应该收入那部出版的讨论班,而它们最终却付诸流水了(连同原始讨论班的另外五六捆\origpage{913}报告——这对于Illusie、Deligne一干人倒不算什么……)。有可能,所说的这条不动点公式就是报告~XII(第431页)中的公式(6.3.1)。匆匆一瞥,无法把它与那几十条被密密麻麻编了号的其他公式区分开来,而它正被埋在其中。显然那位执笔者(Bucur)被那活儿压垮了——而那位才华出众的编辑(sic)Illusie,十五年来惯于做那些清晰完美的文书活儿,绝不会为他的朋友Bucur那些\textbf{正好合他心意}的笨拙\fnote{531}{(*)}{Illusie所作 SGA~5 屠杀版《引言》的最后几行,假装“向1976年死于癌症的 I.~Bucur致敬”——比屠杀版早一年。我不知道其中是否有因果关系——我对Bucur为人的诚实与忠贞毫无怀疑,他是绝不会让像这样一版东西轻易放过而不至少告知我的。无论如何,这场死后致敬所嵌入的那场行动的精神,赋予这致敬一种可疑的滋味。“那是只凭言辞自我交代,而本有一种更符合 Ionel Bucur的意愿与耿直、更能纪念他的方式,那就是缓和它的笨拙,而不是毫无愧色地利用它们。”}抬一抬小指头。恰恰相反,他还设法加剧混乱:使那条本已无处寻觅的关键公式,益发无法与Lefschetz--Verdier公式的那一条区分开来,也无法与它在“报告”中的特殊情形区分开来。在那篇著名的报告~IIIB(sic)的引言里,那位临时凑数的“父亲”Illusie写道:

\begin{quote}
“本报告~IIIB 的第二部分,\textbf{性质上技术性得多}[所以千万别去读它!],是\textbf{受启发}[!]于Grothendieck为建立曲线上某些上同调对应的Lefschetz公式所用的方法[我们可千万别去查是哪些!](见 XII[唯有聪明绝顶的人才知道去哪儿找“那条”公式!]以及(SGA~4$\frac{1}{2}$ 报告)[读者在那里不难找到那条公式,并得知它的\textbf{真正的父亲}是谁……])。”(着重号是我加的。)
\end{quote}

再往下,在同一篇引文中,人们说(即Illusie,不言自明)运用第5号\fnote{532}{(**)}{这回谈的是迹,而且是\textbf{非}交换迹——那些笔误-揶揄,是专为那位亡者保留的,至少只要他还不曾在场反驳……}\dmark{}“技术”:

\begin{quote}
“用以在第6节定义环不必交换的模复形之间上同调对\origpage{914}应的“Lefschetz--Verdier局部项”。”
\end{quote}

而偷偷赐给这些“局部项”的名字,正是我1965年为写那条\textbf{显式的}(“Lefschetz-Grothendieck”)公式而引入的——为的恰恰是不必指及一般Lefschetz--Verdier公式的局部项——这名字显然是为了维持那场由Deligne蓄意制造并维护的混淆,好像所说的那条显式公式在技术上\textbf{依赖于}Lefschetz--Verdier公式似的。再往下几行,我们更是惊喜地得知:“由Grothendieck在(XII 4.5)的Lefschetz公式中定义的局部项”\fnote{533}{(*)}{(5月12日)我被这个关于“我的”公式的异乎寻常的精确(XII 4-5)搞得莫名其妙,刚去查了所引的条目。我发现一条“猜想~4.5”(第415页),它似乎关乎定义局部项的可能性。大家早该料到,这位千金难买的先生又要给我们端出一条他的猜想来了,而不是一个真正的定义……}(偏偏不提这些正是刚刚被慷慨赐名为“Lefschetz--Verdier局部项”的那些)“确确实实就是Lefschetz--Verdier的局部项”(不过这回是\textbf{另一种}意思,当然了:那条\textbf{一般的}、“非显式的”、所谓Lefschetz--Verdier公式的那些)。

% ---------------- p.720 ----------------
要说在污泥里钓鱼的技艺,循着一种我太过熟悉的风格,这真是到家了!

卷首引言里也是同样的混淆手法,其中写道(第 VI 页第5行):

\begin{quote}
“对Lefschetz公式的应用见于报告 XII 与 IIIB。”(着重号是我加的),
\end{quote}

总而言之,为的是让读者毫无希望地迷失其中,连去找、连去试着去找都不必——去找\textbf{那条}一维情形下已知的唯一显式Lefschetz公式(至少直到1981年),那位(既不是Illusie、也不是他的主子Deligne,而是)亡故的前“主持人”(sic)——被他的那位“编辑”-掘墓人Illusie乐呵呵地屠杀了的讨论班的主持人,理所当然地不被提名——的公式\fnote{534}{(**)}{尽管 SGA~5 讨论班的一切本质成果,除了Lefschetz--Verdier公式与 Serre-Swan 模理论(后者不在屠杀版中)之外都出自我手,Illusie却把文本编排得使这些成果中\textbf{没有}一项(不止是那条所谓“Lefschetz”公式失落于某场报告 XII 之中的某处……)显出我卑微的本人与此有过任何干系。由此,他在那场把我这个人从 SGA 中排挤出去的行动中扮演了头等重要的角色——那行动由他的朋友Deligne长期筹备,并在注“葬礼队列——‘为科学尽职’”(n\textsuperscript{o}~175)中找到它的终局。(另见子注“驱逐(2)”,n\textsuperscript{o}~169$_1$。)}。

在原始讨论班中,那场被抹掉并改称为 IIIB(带着一位\dmark{}崭新的父亲)的报告~XI,本插在\origpage{915}六场报告\textbf{VIII至XIII}的序列里,那些报告围绕两个紧密相连的主题——Euler--Poincaré显式公式与Lefschetz公式——并以我在创办该讨论班之际提炼出的共同方法,本着同一精神展开。在那部分讨论班里,如同在其他部分,有一种论旨与视景的显而易见的统一。它被我那位昔日学生一丝不苟地屠戮了,而他正是利用了他作为一艘被他本人和我的其他上同调学生(以对师父的身后谢意)弄沉的讨论班的“编辑”(sic)这一角色。带着堪与缜密的Illusie相称的规律性,六场当中隔一场缺一场——即报告~IX、XI、XIII 三场——从屠杀版中消失了。报告~IX 出自Serre,呈献 Serre-Swan 模理论——眼看事态的走向,Serre宁可收回他的赌注,亲自操心把他那场漂亮的报告交到所有人手中。报告~XIII,据那位“编辑”在卷首引言中的解释,是多余的——显然那位不被提名的主持人不会数到十三——于是就给删掉了!报告~XI,如我们所见,凭一手漂亮的戏法,变成了报告~IIIB,几乎是报告~III 的附录(瞧瞧,瞧瞧——机缘真是巧合得很……),而报告~III 原称“Lefschetz--Verdier公式”,为着混淆的缘故被改称为干巴巴的“Lefschetz公式”。总之,这次“搬场”并非偶然——它始终朝着同一个方向,朝着Deligne-Illusie这对完美搭档从不松懈地维持的那个混淆:一边是Lefschetz--Verdier公式(那条“猜想性的”、“局部项没算出来”的,但终究还是靠Deligne与Illusie的合力得到证明了……),另一边是另一条公式,一条显式的,它必须严格保持隐秘,被小心翼翼地淹没在一大堆编了号、还带四位小数的公式、含沙射影的暗示与精心计算的含糊之中。再次恭喜,亲爱的昔日学生!于是,那场题为“平展上同调中的Euler--Poincaré公式”的报告~X,被掐去了前一场与后一场,悲惨地悬在虚空中\fnote{535}{(*)}{在没有反证的情况下,读者自会猜到:这条著名的所谓“Euler--Poincaré公式”当归于两位以他们的名字为它命名的杰出几何学家。试与上一条脚注比较。}。漂亮活儿,你的时间倒没白费……

\subsubsection*{(f) 恭贺——或新风格}

\tocsubsub{(f) 恭贺——或新风格}

\notehead{169$_9$}{}

(3月22日与4月29日)我想再回到Lefschetz--Verdier公式与那条\textbf{隐秘的}、\dmark{}\textbf{无处寻觅的}公式之间被蓄意维持的混淆。我刚在 SGA~5 里发现了一份相当详尽的\origpage{916}“术语索引”——人要么细心,要么不细心!出于好奇,我去翻“Lefschetz”词条,看我“那条”公式会不会在里面……唯一的指涉是“Lefschetz--Verdier公式(报告~III)”——而那场报告也已(如我们所见)被改称为“Lefschetz公式”了。于是读者得到确切的通知:在这个书卷里不存在(至少不在这里)别的“Lefschetz”公式,只有那条所谓“Lefschetz--Verdier公式”的(正是他刚从别处得知的那条:它是猜想性的,等等;SGA~5 死死依赖着它、终生依赖着它;而“SGA~4$\frac{1}{2}$”如其名义所示在此救急……)。是啊,漂亮活儿!

% ---------------- p.721 ----------------
我继续检视我昔日的学生Illusie——在我另一位昔日学生Deligne的麾下——的种种业绩。我接着引用那部屠杀版导言的引文\fnote{536}{(*)}{引文的开头见前一条子注“变戏法者——或飞走的公式”(n\textsuperscript{o}~169$_8$)。},就在“那条”Lefschetz--Verdier公式(永远是同一条)凭着数学的变戏法之术骤然倍增为“\textbf{若干}Lefschetz公式”的地方,可谁也没能说出是哪些。他接着写道(第 VI 页第6行):

\begin{quote}
“\textbf{报告 XII 的迹公式}[人们满心以为,绝不会有任何读者动念去把它找出来……]\textbf{是独立于报告 III 的一般公式而得到证明的},但人们表明((III B 6))其中出现的局部项确确实实就是一般公式的那些,\textbf{并且后者蕴含前者}。”(着重号是我加的。)
\end{quote}

两手空空,两兜干净——攻不破的Illusie,与那位变戏法的总领班同样攻不破!在先后追随了那一整群一味朝着同一方向的障眼术含混之后,我刚注意到:就在这里,在一个直到此刻都逃过我(一如它逃过这篇四页多导言的任何其他读者)的毫不起眼的迂回\fnote{537}{(**)}{Zoghman~Mebkhout是一位细心却来得稍晚的读者,他告诉我,他自己也曾受骗,深信那条不动点显式公式(对任意维数的 Frobenius,或对一维情形的一般对应关系)确确实实依赖于那条一般的(非显式的)Lefschetz--Verdier公式。由此可见,Illusie那句“暂停!”式的说法,逃过了他的注意一如逃过我的——而这正是所要的效果……

由于我1974年的那场Bourbaki报告,是在人们把一条一维情形下的\textbf{显式}公式明确写出之\textbf{前}写成的,呈献的是“带系数的”函数$L$的公式(可构$\ell$-adic 层中的系数,或与此等价的、Frobenius 对应在这类层中的不动点显式公式)。那时我以为,一维情形下 Frobenius 的显式公式的证明,会是一般Lefschetz--Verdier公式的一个推论——“只差把局部项写出来了”。于是,由于预支了一份尚待完成的工作(而且那活儿照例该归Verdier),我在这场Bourbaki报告里把这条\textbf{显式}公式命名为“Lefschetz--Verdier定理”。后来,无论是Verdier的“伍兹霍尔”证明,还是我的那条覆盖远为一般情形的证明,都不诉诸一般的Lefschetz--Verdier公式。情形对所有 SGA~5 的听众都是再清楚不过的,至少如此。但对于那些只知道我的Bourbaki报告而不知 SGA~5 的人(而 SGA~5 一直被扣到1977年),那里有一个误会,而Deligne与Illusie沆瀣一气,把它彻底利用来搭建那场“SGA~4$\frac{1}{2}$—SGA~5”的骗局(那针脚粗得一目了然)。

况且,就 SGA~5 对那部名字骗人的海盗文本的“逻辑依赖”这一骗术而言,它无论如何也站不住,即便那条显式公式当真依赖于那条“猜想性的”Lefschetz--Verdier公式。因为,正如Deligne本人在那段著名的“方法A”中顺便指出的(对于一个请求宽免的读者——见“真正的数学……”,n\textsuperscript{o}~169$_5$ 第884页),那位不被提名的先生的那些“容易的约化”,把事情归结到一维情形,而在那里“证明的各配料本来就已经齐备”。

所有这些骗术,只要端给一位或睡着、或性急、或正求之不得地要被唬弄的读者,就能成功。对于一位细心而挑剔的读者,整套巧妙的编排就显出它的本来面目:一场不知羞耻的诈骗。但看来,我是八年来这场骗局在数学市场上出现以来,第一位细心而挑剔的读者……}里,

\dmark{}人们半明半暗地说:报告~XII 的某条迹公式(读者只管尽力\origpage{917}去找是哪一条!)是\textbf{独立于}“报告~III 的一般公式”而获得证明的(而后者,为着合宜,同样不配有个名字,合乎所谓“刻意含糊”法……)——以便在同一次呼气、同一句话里接着(仿佛是为了“追回”一句合乎谨慎规则的“暂停!”式说法……)来一个“但人们表明……”。这个“但”指的是那份“柏拉图的”补篇:十二年来,谁也没有为它操过心,首先是Illusie与Verdier,也就是说:“我的”那些局部项——对不起,我是想说“其中出现的”那些(在这条报告~XII 的迹公式里,而那位作者永远不得被明白地提名\fnote{538}{(*)}{对于 SGA~5 的读者,看来该是Illusie——那场关于“局部项”的辉煌报告~IIIB——做那条从未被提名的公式的谦逊父亲。而对于那部自称“SGA~4$\frac{1}{2}$”的书卷的读者,他因此不曾听说过别的公式、只听说过“报告”里的那一条,父亲显然就是那部书卷辉煌的作者。而对于同时读这两者(倘若有人这么做)的读者,他就只好掷硬币猜正面,或者干脆把舌头交给猫……})

% ---------------- p.722 ----------------
\dmark{}——这些项就是那条永恒的“一般公式”的那些项——而关于给\origpage{918}各条公式的名字、以及去哪儿找它们的那种含糊,突然让位于一种堪与缜密的Illusie相称的范本式的精确:这一份“额外赏赐”的证明,却在 III B 6 里——倘若哪位读者想确证它确确实实在那里,他找起来毫不费力,那一条!

那么,对这恒等式为何突然有了兴致,既然整整十一年来,SGA~5 这个讨论班的整个命运,Illusie(一如我的其他上同调学生)都全然无动于衷?为的是能在同一句话里再接再厉地辉煌地(这事儿要么是天意,要么我就再也不知怎么说了)使“那条一般公式”(Lefschetz--Verdier的,就不点名了)\textbf{蕴含}“报告~XII 的那一条”(那位亡故者同样不被提名)。

这真是一手真正天才的戏法!我那位辉煌的昔日学生大汗淋漓(包括靠着数学上的活儿,一点不假),才挣得这句貌不惊人的收尾话的辉煌成果——而在Deligne眼里、以及在他那仆从的眼里,这一句却是性命攸关:Lefschetz--Verdier公式“蕴含”“报告~XII 的那条”(而人们刚刚才说过它是被独立证明的,可为了满足无意识那些全然象征性的快意,这又算得了什么!)。

而这个“\textbf{蕴含}”,就数学而言,性质十分特别——而我敢打赌,我是世上唯一的数学家,除了那个噱头的辉煌发明者(也许还有他的主子Deligne),能品出它的滋味。然而要理解它,并不需要是专家,甚至不必是数学家。两条公式,那条“一般的”(别名Lefschetz--Verdier)与“报告~XII 的那一条”(别名那位不被提名的亡者),分别写成这样的形式

\begin{quote}
$T = L$, \qquad $T = L'$,
\end{quote}

其中项$T$(交错迹和)在两条公式中是同一个,而项$L$、$L'$(局部项之和)则是临时定义的(其一由Verdier本着Lefschetz的精神,另一由那位亡故者本着Nielsen--Wecken-Grothendieck的精神)。

十一年后,Illusie(他的文书热忱一听见主子的招呼便骤然苏醒)做了一场突如其来的努力,本该用在更有价值的事业上,为的是去证明……
% ---------------- p.723 ----------------
$L = L'$\textbf{直接地}(?)(以及对各个局部项逐个地),为的是能说公式$T = L$“蕴含”$T = L'$(并由此隐含地,那部该屠杀的讨论班的公式$T = L'$——对$L$函数理论至关重要的公式——“依赖于”公式$T = L$,而后者在Deligne与他那救命的“SGA~4$\frac{1}{2}$”出现之前始终“属于猜想”——sic……)。

\dmark{}局势对于一个略知底细的人就变得更滑稽了,因为他意识到:若不是$L'$中那些局部项(那位不被提名的亡者的那些)那套稀奇古怪的定义直接被证明公式$T = L'$的运思本身“吹送”而来,\origpage{919}世上任何人都不会想出它来。说实话,我可以说:我在\textbf{甚至还没}定义右端$L'$及其各个局部项之\textbf{前},就找到了公式$T = L'$的一项“证明”——那些局部项是从证明里“冒”出来的,不多不少\fnote{539}{(*)}{我说明一下(这事当然不言而喻):在一切可以想见的应用中(不止于$L$函数公式,那仅关乎 Frobenius 对应),真正\textbf{贴切}的公式确是那条\textbf{显式}公式$T = L'$。从实用的角度看、并且就一维现象而言,Lefschetz--Verdier公式$T = L$只有一种历史(或启发)的兴趣,而Illusie的结果$L = L'$(或者更确切地说,出现在$L$中的那两类局部项与出现在$L'$中的那两类,是同样的)就更是如此了(至少在新命令到来之前)。这些原是再明白不过的事,而那两个同谋却设法(而且,照眼前的时势,还成功地)把它们搅浑。这让人不禁思索:我们眼前这场狂热的科学生产,究竟能有什么意义——而那一切都发生在一些与二十五年以来我们在几何与算术之关系上的认识所取得的关键进步息息相关的问题上,对普通数学常识如此粗暴的歪曲,竟丝毫未引起任何人的注意……}。

恭喜,这是第三次了,Illusie;也恭喜你,Deligne,你给他当了楷模。你们两位合力开创了数学中一种\textbf{新风格}的先河。而这风格已经收了门徒。

如今它已然成了“1980 风格”,显然前程似锦\fnote{540}{(**)}{关于这方面触目的例子,见“黑手党”一注(n\textsuperscript{o}~171$_{2}$)中若干“1980 风格”的样品,出自我们那几位大作者手笔:Brylinski、Kashiwara、Beilinson、Bernstein。显然,什么希望都是允许的!

(5月12日)至于这“新风格”另一些偶然的信徒——他们是在一位从未被提名的默默无闻的身后弟子的工作尾迹里崭露头角的——如今我可以补上 Malgrange、Laumon、Katz。(见“掠夺的空白授权”一注,n\textsuperscript{o}~171$_{4}$。)}。这便是变戏法-风格,别名“掘墓人的风格”,其全部艺术在于不断地\textbf{欺骗读者};不只在主要观念的\textbf{父权}上骗他,而且(顺带着)在它们的师承与相互关系上、在每一项的份量上、在什么是本质、什么是枝节上骗他——为的是让那该受抬举的人受抬举,让那该遭贬损的人遭贬损(或者以一个无精打采的手势、在一句若无其事的话的转折处把他埋葬……);而且\textbf{尤其是},为了一份刺激的\textbf{权力}之感:牵着读者的鼻子走,随心所欲地\dmark{}\origpage{920}炮制并拆解他的科学的历史,并决定那些声称要呈献的数学之物“是”什么、又不是什么。这便是那永远“\textbf{称王}”的艺术:轻轻牵动无形的(?)丝线,而永不、永不纡尊降贵去服侍。而这一切,为的是永远、彻底地享有“暂停!”的豁免:即便有哪位读者比别的更机灵,竟会亲自去看一看,竟会别出心裁地想动用(天知道……)他自己的眼光与官能(这少见,但说到底也不是不可能……),他也永远不能当场捉住你说了某句一经按字面理解、又无含糊或双关可逃的话,那确确实实、无可挽回地是\textbf{假的}。

这种句法上的艺术之艺术,在于那样一种风格条款,它看似一道难题,然而……随着那场记忆古怪的“反常层会议”(那是一场令人叹为观止的“SGA~4$\frac{1}{2}$—SGA~5”行动之后不过四年),人们得以看出这种新的、无害的技艺能走多远——在那场妙不可言的行动的变戏法绝技的种种施展之后,在对一部革新性作品的掉包中、在对那位曾长期孕育并在孤独中耐心成熟这部作品的人的肆意掠夺中……

向师父和学生、向Deligne和Illusie脱帽致敬!艺术家的活儿!你们二位当之无愧地赢得了教团全体的一致感谢。

% ---------------- p.724 ----------------
\subsection{(2) 分赃(“对偶——晶体”)}



\subsubsection*{a.~最后一份——或聋子的耳朵}

\tocsubsub{a.~最后一份——或聋子的耳朵}

\notehead{170(i)}{}

(2月28日)我这就说到围绕我数学著作的“四种运算”中的第三种(第四种且待下一条注,那一条将抹杀Zoghman~Mebkhout的作品)。

\begin{quote}
III “对偶——晶体”行动(或称:“美丽的残骸……”)。
\end{quote}

就我如今看事情的样子,这大体上是对我那部分著作中关于上同调、而尚未被\textbf{Pierre Deligne}以事实(或象征)方式侵吞的部分的一次\textbf{分赃}\fnote{541}{(*)}{(5月1日)尽管如此,仍该把\textbf{凝聚}语境中的对偶形式主义单独挑出来:它(与一个后来显得匆促的印象相反)显然还不曾被我的上同调学生中的任何一位、也不曾被我知识范围内的任何其他人所侵吞。诚然,唯一陈述我在这个主题上的大部分想法与结果的参考文献,是 R.~Hartshorne的《留数与对偶》(Residues and Duality),这使人可以指到它,而无须在任何一个时刻说出一个不受欢迎的名字……}。显然,他给自己留了最大的一份,带着 Motif 与平展上同调,更具体地说,是那套$\ell$-adic 的上同调工具。其余部分的分割(*),则在我的另两位上同调学生\origpage{921}\dmark{}\textbf{J.-L. Verdier}与\textbf{P. Berthelot}\fnote{542}{(*)}{(5月1日)自从写就以来,看来该添上“第四个盗匪”,即Neantro~Saavedra~Rivano,他经由那些应景地受洗为“坦纳克式”的范畴,侵吞 Motif Galois群的哲学。但他只是替Deligne当“稻草人父亲”,而后者在十年后“收回”了父权。事情的始末,见后续诸注“棺材上的第六颗钉子”,n\textsuperscript{o}~176$_1$至176$_7$。}之间进行。那默契是何时、怎样形成的,我说不上来;它似乎是这样:Berthelot得全部晶体上同调,而其余归Verdier——后者基本上并吞了一切围绕对偶瑜伽\fnote{543}{(**)}{见上一页的脚注。}的东西,以及作为其代数前提的导出范畴与三角范畴瑜伽。

关于Berthelot参与分我这具遗骸,我只有一个事实可依据,虽说那事实分量不轻。去年,我在“共同继承人……”(n\textsuperscript{o}~91)那条注的反思中偶然撞上它,并为它写了一篇短小的子注(n\textsuperscript{o}~911)。那就是我在那里引述的Berthelot的那篇综述文章\fnote{544}{(***)}{刚性几何与特征~$p$代数簇的上同调,Pierre Berthelot,载“$p$-adic 分析及其应用”CIRM 吕米尼学术会议(9月6--10日)。},它在一场题为“$p$-adic 分析及其应用”、于1982年9月举行的吕米尼学术会议上,为 Dwork-Monsky-Washnitzer 上同调与晶体上同调的一种“综合”(如他所说)呈献了主要的想法。在那篇引言(b)部分里,他给晶体上同调写了一小段历史,其气量狭隘,毕竟与我对晶体瑜伽所抱的远为宏阔的视野全然不合\fnote{545}{(****)}{关于这一点,见随后那篇子注“聋子的耳朵”(n\textsuperscript{o}~170(i)bis)。}。

我的名字在文章正文与文献目录中都缺席。至于若干评注与细节,读者可参看所引的那篇子注,这里无须重复。我只想补充一句:一旦把我这个人从图景中清除,那晶体上同调之父,不是别人,正是他Berthelot一个人,甚至不消他费心把这话明白说出来——某种侵吞的风格显然已经收了门徒……事实上,正是他那篇论文——他当年照着我的若干出发想法、与我一起筹备的——构成了这个晶体主题上第一部出版的作品(除了我自己\origpage{922}就若干出发想法\dmark{}所做的那份极其简略的勾勒之外\fnote{546}{(*)}{这些想法的唯一一份已发表的勾勒,出自我1966年11月与12月在 IHES 所做的五场报告,由 I.~Coates 与 0.~Jussila 执笔,即《晶体与概形的 De Rham 上同调》(Crystals and the De Rham Cohomology of Schemes),载《论概形上同调的十篇报告》(Dix exposés sur la Cohomologie des Schémas,North Holland, Amsterdam 1968)pp.~306--358。所有本质的出发想法都在那里有了勾勒,包括引入 Monsky-Washnitzer 式局部“加厚”的必要性(pp.~355--356)。})。他的论文为那个我曾向他提出的纲领的至少第一部分(170(i)bis)呈献了一项规模宏大的奠基工作。

这篇可纪念的“综述”写于1982年,即“反常层会议”(1982年6月吕米尼)之后一年;那场会议将在“第四行动”中谈到。我不曾费心去重读我手头Berthelot的抽印本,以便了解这场对我的葬礼的参与,究竟代表他与我和我的作品之关系上的一个晚期转折,还是一种更古老的态度之延续。若是前一种情形,那在某种意义上大可打赌:这个转折是对前一年随那场会议而来的、科学伦理总体败坏中那场骤然且疯狂的自我升级的一种回应。为此我提醒一句:同一年1982年,也因那部“值得纪念的书卷”LN~900 把 Motif 发掘出来\fnote{547}{(**)}{见“沉默”(n\textsuperscript{o}~168),尤其是“……与发掘”(n\textsuperscript{o}~168(iii))。}而出名——在那里,为这场行动买单的不再是一个无名的“当值无名之辈”(如在那场辉煌的会议上),而是一位“亡者”,他的名字毕竟还留在人们的记忆里(哪怕是出于遗憾……)。前一年的那场行动,已经足够清楚地表明:再也不需要任何克制——而“Motif 行动”事实上与“晶体行动”以及此前的所有行动一样,连一点涟漪也不曾激起……

\notehead{170(i)bis}{}

(170(i)bis)(2月28日与4月30日)\fnote{548}{(***)}{本子注源出上一条注“最后一份”的一条脚注。(****)(5月12日)事实上,那已是在1966年,见上文脚注(*)。}我在此所谓晶体理论(特征~$p>0$)的“第一部分”,指的是特征~$p$基概形上\textbf{真而光滑}的概形、以常系数(或“扭曲的常系数”)为项的晶体上同调。那时只需使用我在六十年代末\fnote{549}{(****)}{(5月12日)事实上,那已是在1966年,见上文脚注(*)。}(暂时地)引入的“寻常”或“无穷小”晶体 site 就够了。事实上,与Berthelot乐于赋予“晶体上同调”一词的狭隘含义相反,它对我来说从一开始就有一种远为\dmark{}宏大的含义,而我对谁都不曾\origpage{923}隐瞒,我的学生们显然已经把它忘了——为了在十或十五年之后又“重新发明”了它的一小块……

% ---------------- p.726 ----------------
一方面,我的晶体想法从一开始就绝不限于给定特征~$p>0$~的概形这一情形。我最早的晶体反思——在“带除幂的加厚”这个新念头来到我心中之前——处在\textbf{特征零}的概形上,那里除幂自动在场(正因此,它们往往不被人留意……)。这个研究方向由于Zoghman~Mebkhout的想法而重获生机,其自然的终点,将是在(先从)特征零概形上、针对“De Rham-Mebkhout 晶体系数”的六种运算形式主义——那个形式主义,我在“墓上的旋律——或自负”(n\textsuperscript{o}~167)那条注里已经提到过。早在六十年代,我就预见过一种不分特征的晶体上同调,其形态是在绝对基$\mathbb{Z}$上(例如)有限型概形的语境中的一套晶体“六种运算”形式主义。它本该涵盖$\mathbb{F}_p$($p$元素域)上有限型概形的“寻常”晶体理论(那个理论当年还在寻找——而今也还在寻找)。我深信:忘却并埋葬了亡故师父的这一视景(它却再简单、再富于启发性不过),正是晶体理论那可悲的停滞之因——而在它起步时那股蓬蓬勃勃的生气之后,已将近二十年。

\origpage{924}另一方面,且回到 Monsky-Washnitzer 的进路——它当年曾为把我“触发”到晶体上同调起了作用——我从一开始就清楚:有必要引入(为着一门不只适用于真而光滑的概形的理论)一个比“无穷小”site 更大的晶体 site,在其中,所设想的“加厚”将是适当的\textbf{拓扑}代数(带除幂理想)的谱,或许是 Monsky-Washnitzer 所用的那些(剥去了诸如光滑性这类无用假设的)(*)。提炼出“\textbf{对的}site”\dmark{}\fnote{550}{(*)}{如我在先前一条脚注中所说明的(见第922页),我在那篇关于晶体瑜伽的、唯一一篇已发表报告(1966年末)中谈到的正是这类 Monsky-Washnitzer 式加厚。从那时起我就已清楚:特征~$p>0$~的晶体上同调,将有很大一部分在特征零的刚性解析空间上进行。我当然不曾不把这告诉一切可能与此相关的人,首先肯定是我的学生Berthelot——一旦他选定投身晶体主题。在所引的那篇文章里,循着一种我认得的、而Berthelot未曾发明的风格,人们会以为他刚刚(晚了十五年)发现了与刚性解析几何之间那条未曾疑及的关联(x)。他在那里以一个“共同推广”(Monsky-Washnitzer 理论与晶体理论的)的天才发明者的姿态出现,并大言不惭地把那推广命名为“刚性上同调”(而它很快就会如愿地叫作“Berthelot上同调”)。我还得指出,Berthelot的这项工作,是“与 Ogus 共同进行的一场反思的延续”——而同一位 Ogus,正是在同一年(1982)以 LN~900 共同作者的身份,参与了那场“Motif”骗局而崭露头角。

系统性的埋葬在Berthelot后续的一篇文章(我手头有一份预印本)《刚性上同调与 Dwork 理论:指数和的情形》(未署日期)中继续进行。对$F$-晶体这一关键概念、或紧支上同调的概念(而后者是我1963年2月在代数几何中荣幸地引入的,早了二十年……),文中毫无对亡者的指涉。这些概念既然如此自然,实在不值得为那点区区小事费神……形式概形(在离散赋值环之上)的“泛纤维”的概念,作为一刚性解析空间,则被慷慨地记在我昔日的学生Raynaud的账上。而我早在Berthelot、Raynaud、乃至任何其他人还不曾听说过“刚性解析空间”这个词之前就已知道这一概念——因为能够定义这样一条泛纤维的需要,原是我预想存在一门“刚性解析几何”的两大动机之一,而后来也正是它,成为引导Tate把这样一门几何在形式上搭起来的导线之一:他的定义必须是使得“泛纤维”的概念成为同义反复……

(x)(1985年9月)事实上,头一个预见这样一门理论之存在的是 J.~Tate,时在1959年8月。关于这一点,见注~n\textsuperscript{o}~173$^d$)(“葬礼——或自然的下坡”),尤其是第1132页的脚注。}与“\textbf{对的}系数”,是我遗赠给我的上同调学生们的纲领之一部分(以Berthelot打头)——如今看来,那完全是白费。近来(趁写《收获与播种》之便)“顺带”想了想这件事,又记起一门囊括一切特征于一体的晶体理论之必要,我甚至自问:这些拓扑代数(Monsky-Washnitzer 式的,或任何合理变体)是不是也同样太“粗糙”了(如同受限形式幂级数一样),因为它们太“远离代数的东西”,是不是该用(在适当意义上)是“平展邻域”的“加厚”来代替它们。我想在《收获与播种》之后那部分《反思》(我揣测是第3卷)里回到这些问题,连同六种运算的瑜伽与“系数问题学”的阐述,特别是“De Rham-Mebkhout”类型的晶体系数。

Mebkhout还预见过:他的 D-模哲学该为晶体理论提供一个新的视点。但他这方面的建议——特别是1978年向Berthelot提的——出自一位含糊的无名者、一位不肯悔改的Grothendieck主义者,便落进了聋子的耳朵\fnote{551}{(*)}{耳朵聋并不妨碍同一个Berthelot,在我在上一条脚注里引述的那篇文章里,满不在乎地(在第3A节末尾)指到“复簇上$DX$-模理论的一个类似物”,而人们在刚性解析框架中“暂时”还不曾掌握它。当然谈不到在此提起某位含糊无名者的名字——他曾在四五年前来向他提过一些荒唐的建议,而这就更不必提了:因为前一年的某场学术会议(将在下一条注“封神”,n\textsuperscript{o}~171 中谈到)已经为这位无名者定下了清楚的调门。可以肯定,再过几年,借着那套众所周知的、所谓“Riemann--Hilbert-Deligne”哲学的真正之父的祝福,Berthelot就要以 D-模哲学在“刚性解析上同调”(也称作,尽管他自己避免这样称呼,“Berthelot上同调”)中的天才发明者身份亮相了。由此可见,凭眼下这年景,要远行,耳朵不必有多灵……}……

% ---------------- p.727 ----------------


\subsubsection*{b.~荣耀满天飞——或暧昧}

\tocsubsub{b.~荣耀满天飞——或暧昧}

\notehead{170(ii)}{}

\origpage{925}\dmark{}(2月28日)\fnote{552}{(*)}{这条注的文字曾于5月1日(铃兰日)重新整理,并在若干处作了修正。}为把“对偶行动”定位在 J.L.~Verdier那可疑的受益上,本该先说几句对偶瑜伽的话(所谓“六种运算”的——但那名字已沉没得无影无踪),那是我从五十年代后半期起发展出来的;还有导出范畴的瑜伽,它与前者实则不可分割。我在“我的孤儿”(n\textsuperscript{o}~46,尤其是第177--178页)那条注、以及它的子注 n\textsuperscript{o}~462(第186--187页)里,已就此说得相当周详;最后(作为对Verdier在我那同调代数观点之葬礼中所扮演角色的初步反思)在“本能与时尚——或强者的法则”(n\textsuperscript{o}~48)那条注里也谈过。我看再回头去讲是没用的,只建议读者需要时去参看,然后再继续读“对偶行动”的经过\fnote{553}{(**)}{(5月12日)另见“祖先”(n\textsuperscript{o}~171(i))与“工地的巡礼——或工具与视野”(n\textsuperscript{o}~178)两注,尤其是“六种运算”与“系数”两个工地(n\textsuperscript{o}~3、4)。}。

Verdier在这场分赃行动中的态度,比他那两位朋友的态度显得更暧昧,因为他\textbf{曾}有时同时\textbf{在两个台面}上玩,而那两个台面可能看似互相矛盾。我起先自己也难以在其中认出方向,那局面看去是那样混乱。一方面,在1967年他的论文答辩之后、尤其在我1970年离开之后,他努力(由于一些我所不解的理由)\textbf{埋葬并败坏}那套他得自我的同调代数与对偶瑜伽——而他原本把最充沛的精力、在整个六十年代直到论文答辩,都用于发展这些思想并以他自己的贡献使之丰富。另一方面,至少从1976年起(在他那篇论文(sic)答辩九年之后),并在Deligne的鼓励与有效支持下,他装作\textbf{侵吞}那些出发思想的父权(只要它们没有被抵制),也侵吞\dmark{}我围绕\origpage{926}平展对偶这一主题所发展的全部方法与成果\fnote{554}{(*)}{当然,在所说的那些“其他语境”中,平展语境出发时的困难、即需要一次“突破”以给平展上同调至少一点抓手(在缺乏以奇异单形、收缩法等手段做出的著名超越构造的情况下)——这一困难并不存在。我的学生们都找到了那样的情形:那场“突破”的重活已经由别人替他们干完了——他们只消搬自己的家具来,而往往那“别人”还奉送了家具。机会一到,他们就忙着把那人埋掉,好把自以为可以侵吞的东西据为己有,并嗤笑其余……}——那些方法只需略作变通(mutatis mutandis)便能套用到一切其他语境,诸如拓扑空间或复解析空间的语境。

关于Verdier对单单是导出范畴的态度,我曾在“赊账的学位论文与全险”(n\textsuperscript{o}~81)那注里设法指出这种暧昧的意义\fnote{555}{(**)}{写这条注时,我还不知道Verdier当年是怎样露了一手的——凭着1976年他提供的那“好的参考文献”——见下文“第2步”。}。读者还将在那里找到若干素材,尤其关于他的论文(至今仍未出版)与答辩周围那些古怪的情形。隔了一年回头看,这场反思中渐渐显出的那幅图景,在我看来或许是正确的(至多还需些修订),然而毕竟是皮毛的。我很清楚:Verdier的\textbf{真正}动机,不在任何一种近乎可怜的“回报计算”的层面上,而全然是另一种性质,并且本质地牵涉他对我这个人的矛盾关系。我琢磨,即便对一个皮毛的观察者而言,他的情形也格外触目:他以为埋葬了那位曾是他师父的人,实则日复一日、直到今天所埋葬的,不是别人,正是\textbf{他自己}与他自身之中的创造力。

% ---------------- p.728 ----------------
为把“对偶行动”兜一圈,我现在要对我所知道的这场行动、以及更一般地说Verdier对我的埋葬的参与,做一次简短的回顾。

\textbf{第1步}(1966--1976)。那是在我1970年离开之后——我说不准究竟何时了——Verdier通知我:他不再打算出版他的论文。我提醒一句:那篇论文本该呈献同调代数的新奠基,以导出范畴为着眼点。在我看来,他那篇论文工作的存在理由,就是交到所有人手中,提供一部参考文本,其意义可与 Cartan-Eilenberg 的书相比,并直接顺应\dmark{}五六十年代在我的工作和我的学生们的工作的尾迹中出现的那些新需要。\origpage{927}凭后见之明,我明白:那套新的上同调语言,当年(而我甚至可以说今天也依然……)只被我的上同调学生们完全消化,而Verdier的决定,于是就等于在这幅同调代数的新图景上划了一道大杠。由此,他那“论文”——只有二十五页,仅只呈献了一些想法的令人信服的勾勒,而那些想法他自己就说不该归他——也失去了意义,严格说来成了一份“冒牌论文”。可七十年代初,在(惊讶地)得知Verdier的决定时,我正如此强烈地埋头于与我当年那些数学兴趣截然相反的事务,这些问题对我其时是无穷地遥远。我不曾起念头去细看这件事——它是在匆匆间(我可以想见)夹在关于萨克莱原子废料桶裂缝丑闻的一场公开讨论、与为编写《生存与生活》简报的一次工作会议之间听来的。而当时我更不曾想到要去反应。我头一回终于“细看”Verdier这一举动之意义、并开始隐约看出它是一场刻意破坏的,是在那条已引的注“本能与时尚——或强者的法则”(n\textsuperscript{o}~48)里,几个星期之后,即在发现那场葬礼“以其全部辉煌”之后,又在那一注更为周详深入的“赊账的学位论文与全险”(n\textsuperscript{o}~81)里重提。

回想起来,有一点很清楚:Verdier在他自作自受的那份工作上的分心——那本是他与论文评审委员会订下的“诚意契约”(见所引注 n\textsuperscript{o}~81)的一部分——至少该上溯到1968或1969年;否则他那“论文”的撰写与出版,早在1970年我离开之前就已成事。我提醒一句:我还在1950年就把论文工作的纲领交给了他,而对一位像他当年那样有天份、有干劲的研究者,这个纲领加上一份新奠基的大规模撰写,满打满算也不过三四年工夫,连熟悉入门都算在内。同样真切的是:有某种心态,就是设法预先为一份“预定”的工作支取一笔信用,事后再也不必费神去把它做出来——这样一种心态,在我看来,早在1964年之后就已经显山露水:有那所谓“Lefschetz--Verdier”公式的坎坷,后来有\origpage{928}局部紧空间的对偶(顾名思义,就叫作“Verdier的”对偶),凭着六种运算的精神\fnote{556}{(*)}{关于这种相当特别的心态,见子注“遗产——或骗术与创造”(n\textsuperscript{o}~169$_6$bis),以及去年那两篇子注(n\textsuperscript{o}~812、813)——在那条已屡次引用的“赊账的论文与全保险”里。}\dmark{}(它们至今仍不被提名)。但整个六十年代,我身为囚徒,被自己的任务、被自己通过它们不懈追求的那幅图景所困,如同阿哈那难以捉摸而又无处不在的白鲸,我万万不曾疑心:在那位曾是我在自认为“共同”的任务中的亲密伙伴的人身上,有什么东西“不对劲”——正如我也毫不曾对我的上同调学生中任何别人起过疑。而凭二十年的后见之明,我现在才惊觉:我生命中有十年(若不是十五年、二十年),我活得与周围的现实全然\textbf{错位},而这不仅在我的家庭生活里(我在那里早已有所觉察),也同样在我全身心投入的职业生涯里……

% ---------------- p.729 ----------------
但且回到“第1步”。Verdier对我这个人和我的作品的暧昧关系,无论如何在 SGA~5 讨论班1966年结束之后就已显现:我的上同调学生中没有任何一位(他也不例外)觉得这个讨论班的整理工作与自己有关\fnote{557}{(**)}{回想起来,我不由要问:1964年(那时在我的接触下,他终于投身新的上同调技术的洪流)到1970年之间,Verdier的时间到底都用在什么上面了——因为他竟不肯着手并完成任何一项撰写工作,甚至对那些他即将自居为作者的理论也是如此。至于他那份清单:有分量却没有一项做到头,见那条被屡屡引用之注的子注~n\textsuperscript{o}~811.}——而那整理始终留在那些被任务压垮的、或不屑守信的自愿者(sic)手里。显然,就在那时,我的上同调学生们整体的情形已经腐烂,而我却丝毫不曾觉察,宁愿活在一个一切都只有秩序与美的世界里……直到十八年后,我才开始投去第一道怯懦的目光,看看当年究竟发生了什么——而那些岁月(仅在一年前)在我还显得田园牧歌\fnote{558}{(***)}{尤其见《愚妄与更新》中“一个没有冲突的世界?”一节(n\textsuperscript{o}~20),那里只有节名中的问号,能让人对那“牧歌”起一点疑心。}……

我1970年离开之后,甚至在他向我宣布他那“正式的”破坏基础工作的决定之前,Verdier在那些六十年代的暧昧,已由他与他的朋友Deligne当年那些各色小型骗局的串通得到确认——而他不可能不曾察觉:\origpage{929}\dmark{}我在 Hodge I、II、III 诸文中被抹掉\fnote{559}{(*)}{在那场“权复形”的玩笑(见同名之注,n\textsuperscript{o}~83)里,我几乎以为辨出一句暗讽,以挑衅的口吻,指向我的上同调学生中某一人的最早一桩昭然若揭的骗局——即Deligne1968年那篇论谱序列退化的文章。如果说当年我没看出其中名堂,那么我最出色的学生所树立的榜样,却并没有对所有人都白费!},然后是在我那单值群讨论班的出版版 SGA~7~II 中(它以Deligne与Katz的名义呈献,后者突如其来地占了一个亡者尚且温热的位子……)。同一年(1973年),他也不曾能不读到 Mac~Pherson 的那篇文章——那里解决了一条“Deligne--Grothendieck猜想”,而他心里清楚Deligne与此毫无关系。

直到1976年,在那些默许的并吞行动中,Verdier的角色看来主要还只是被动的。但相反,由于不出版那本该构成他论文的东西(那论文曾是“赊账”给他的\fnote{560}{(**)}{见已引注 n\textsuperscript{o}~81。}),他在我离开之前,就在那场对我的同调代数观点(他曾一度把它当作自己的)及其作为一种“日常”技术在代数几何、拓扑与代数中的使用的葬礼中,扮演了关键角色。正如他的朋友Illusie和Deligne一样,就这样破坏自己双手的工作,为着埋葬那位启发了他的人的乐趣,他当之无愧地赢得了教团毫无保留的拥戴……

这番蓄意的埋葬之意,在他1975年之后对Zoghman~Mebkhout那番令人沮丧的态度中也表现得明白——当时Mebkhout装作要取法我的对偶瑜伽和导出范畴的瑜伽。关于这一点,我仍请读者去参看那些已引的、更周详的注:“我的孤儿”“本能与时尚——或强者的法则”“赊账的学位论文与全险”(n\textsuperscript{o}~46、48、81),以及“当值无名之辈与好上帝定理”一注(n\textsuperscript{o}~48$'$)\fnote{561}{(***)}{(5月1日)另见“封神”一注(n\textsuperscript{o}~171)的子注“一个视景的萌发——或闯入者”(n\textsuperscript{o}~171$_{1}$)。}。

\origpage{930}\textbf{第2步}(1976).1976年出版了Verdier那篇“值得纪念的”发表在 Astérisque 上的文章\dmark{}\fnote{562}{(*)}{J.L.~Verdier,《结合于一个闭链的同调类》,Astérisque n\textsuperscript{o}~36 (SMF) p.~101--151 (1976)。},前面已作为“一场攀登的第3幕”谈过(见“操演”一注,n\textsuperscript{o}~169)。我提醒一句:那篇五十页的文章,除去几页他自己的东西,就是把一些我在 SGA~5 中早十年发展起来的观念与技巧原文照录,而丝毫不提我个人、也不提任何一个谈过这些事的讨论班。这次出版是我一年前在“反常层会议”的尾迹中(在“好的参考文献”一注,n\textsuperscript{o}~82 里)发现的,它当时以一道全新的光,照亮了Verdier本人和我的其他上同调学生们那份把 SGA~5 讨论班(以它这个名字、并带着属于它的父权)交给数学界的不情愿究竟意味着什么。

% ---------------- p.730 ----------------
这里无须回头再谈我昨天在那条已引的注(n\textsuperscript{o}~169)里对这篇文章所做的评论。作为一个可笑的细节,我只想补一句:正是Verdier这份“工作”(sic)的手稿——他上年(1975年)曾好心把它寄给Zoghman~Mebkhout——成了后者踏入簇之上同调的芝麻开门,并成了他对那位从此以“恩人”身份出现的人物的毫无保留之敬仰的根基。而且这份敬仰还很是耐久,我相信,它只是在Zoghman在反常层会议上的那些倒霉遭遇之后,才彻底瓦解。

Deligne告诉我\fnote{563}{(**)}{见“把一个个 i 点清”一注(n\textsuperscript{o}~164),第 IV 1 部分。},说他是在“SGA~4$\frac{1}{2}$”(sic)与 SGA~5 出版之后一年(1977)——才得知Verdier那篇文章的:这与我的信念正好相反,我认为那“好的参考文献”的出版,标志着那场终于导致那场规模空前的“SGA~4$\frac{1}{2}$—SGA~5”行动的骗局攀登中最后一个本质的步伐。细想之后,我难以相信Deligne的说法。他是我认识的消息最灵的数学家之一,而且自始至终与Verdier关系密切;他不可能不曾早已知道Verdier的计划、不可能不曾收到一份预印本(而这甚至早在Mebkhout之前)、也不可能不是在1976年头几个拿到抽印本的人之一。那篇文章填补了文献中一个张着的窟窿(鉴于1966年之后 SGA~5 讨论班一直不曾出版)——这是Deligne本人向我确认过的——\origpage{931}而他也不可能不曾至少费心把它浏览一遍:对一个像他这样“在圈内”的人,那不过是至多一刻钟的事\fnote{564}{(*)}{我甚至可以想见:比数学兴趣更强得多的(而那篇文章对作为 SGA~5 听众的Deligne本来毫无可学之处),必是能第一手、白纸黑字地得知对那位亡故师父的无瑕疵抹杀——循着他自己八年前就已开创的传统!}。无论如何,这一桩明火执仗的抄袭不曾引起 SGA~5 另外六七个明明“在圈内”的昔日听众中任何一个人的反应,这本身就是一切相关者之间一场无瑕疵串通的保证。时机已经成熟:对那母讨论班 SGA~5 的正式屠杀,以及把我的平展上同调作品炸成碎片……

\textbf{第3步}(1977)。在这场1977年发生的“SGA~4$\frac{1}{2}$—SGA~5”行动中,由Deligne发起、并由Illusie踊跃参与,Verdier这一次只扮演一个跑龙套的角色:他为那部名字骗人的“SGA~4$\frac{1}{2}$”这份寒酸的小册子添了一份砖瓦——即他那篇论文(sic)的某种“第0状态”(而论文本身,则彻底消失了……),为着合宜,在沉睡了十四年之后被专门挖了出来!在整部书卷里,无论在那篇把这份额外赏赐的文本(“已无从寻觅”——而且事出有因!)大大标举出来的引言中,还是在这份文本本身中,都没有一处暗示我曾在其间发展的思想中扮演过任何角色;也不曾(顺带地)提到这份文本当年注定要成为一篇论文。与Deligne一样,Verdier也不曾认为值得告知我这次出版(其故也一样),也不曾给我寄一部那部障眼的书卷。至于细节,见“同伙”一注(n\textsuperscript{o}~63$'''$,写在发现这场偷偷摸摸的发掘的激动之下),以及在那条已屡次引用的“赊账的学位论文与全险”(n\textsuperscript{o}~81)里更深入的反思。

% ---------------- p.731 ----------------
于是,在他那场不寻常的论文答辩十年之后,Verdier抓住了Deligne提供给他的机会,实质上是在同调代数中对“导出范畴”这一观点预先占下一份无可争辩且排他的父权,有他那位显赫的朋友的完全担保;而这正是他们两个还在继续对这一同样的观点维持着事实上的\textbf{抵制}\fnote{565}{(**)}{如我在先前一条脚注(注见第~页,底本页码空缺)中所解释的,在那部名为“SGA~4$\frac{1}{2}$”的棺材文本中,Deligne不得已曾在“那条”公式的证明中求助于导出范畴。想必正是这一点向他暗示了:该用一篇沉没的论文的“第0状态”来给他的书卷壮壮声势。事实上,直到1981年,这都没有改变对导出范畴的抵制局面。}的时候。那场抵制沉重地压在Zoghman~Mebkhout的工作上,判了他完全的孤独,并且一直维持到1981年的“反常层会议”。

\dmark{}\origpage{932}于是,在1977年,Verdier以一门上同调瑜伽的储备之父的身份出现——而那瑜伽,眼下还是一股合宜的、心照不宣的轻蔑的对象,可谁也说不准……此外,自上一年起,随着那“好的参考文献”的出版,他又以我发展出的那套对偶形式主义之一部分的父亲身份出现(关于结合于闭链的“离散”同调类与上同调类、双对偶性形式主义、以可构性表述的有限性定理等等)——这还不算局部紧空间的对偶,那也停留在一个暧昧的地位,一个等待的地位——一如赋予它以意义的导出范畴瑜伽。

\textbf{第4步}(\textbf{反常层会议},1981年6月)。那是迄今为止Verdier对我的葬礼之参与的登峰造极。那场会议把对Zoghman~Mebkhout——代数簇上同调中 D-模那套统一而丰产的观点的先驱——的肆意掠夺奉为神圣。作为那场会议的正式组织者(与 B.~Teissier一起),Verdier在其中扮演了头等重要的角色。我将在下一条注里,连同“第四行动”(所谓“反常层会议行动”或“当值无名之辈行动”)回到它。这里我只限于谈它对Verdier的直接后果,以“一份遗产之分赃”(而那遗赠的亡者则被小心翼翼地置之不理……)的名义。

那场会议奉导出范畴与三角范畴凯旋重返数学竞技场。以这些范畴的“父亲”身份(而他十五年来竭尽全力要埋葬它们),Verdier在Deligne之后,是那一事件的主要英雄。至少这是从那场会议的主要文章里透出的印象——那篇文章出自Deligne的笔下,一篇自成第一册、并构成会议《论文集》的扛鼎之作\fnote{566}{(*)}{《论文集》出版于 Astérisque n\textsuperscript{o}~100 (1982)——题为《奇异空间上的分析与拓扑》。事实上,那些标注1982年的《论文集》,直到1983年12月才印完,而Mebkhout是1984年1月才读到的。}。仿佛是巧合,一篇论文的(我当年绝不会接受它作为博士论文的)瘦骨伶仃而又救命的“第0状态”——它当年赶来给那部海盗文本“SGA~4$\frac{1}{2}$”救场——那书卷\dmark{}\origpage{933}一时未免单薄了些——瞧,它摇身一变,成了那件光彩夺目的呈堂证物:靠着对“SGA~4$\frac{1}{2}$”的一大阵引用,好让那位来路不明的父亲Verdier以谦逊的姿态,自诩为那场所谓“反常层”的大冲刺(而反常层与此毫不相干!)以及代数簇上同调一场新的、迟到的再出发的预见先驱(那是在一位没有人肯提名道姓的含糊无名者的足迹上……)。

同一篇文章(署名 Beilinson-Bernstein-Deligne)也奉平展语境中的六种运算形式主义(当然,那名字从来不曾出现过)强力回归,带着我在五十年代引入的、如今已被公认的那些记号。如我在别处所写\fnote{567}{(*)}{见“不义”一注(n\textsuperscript{o}~75),第288页。}:“在所引的那篇文章里……没有一页不是深深扎根于我的作品、并带着它的印记,甚至细致到我所引入的记号、以及每一步出现的那些概念所用的名字——那些名字,是我与那些概念初识、而它们还未被命名时,由我赋予它们的。”

我那平展对偶形式主义——十八年前发展的,而那时我的学生Verdier还在学上同调语言的 ABC——在普遍的狂喜中被改称为“Verdier对偶”\fnote{568}{(**)}{在记号索引里,那个我在1963年于平展语境中引入、并且(顺带说)构成Illusie版 SGA~5 报告~I 之主题的对偶函子,被叫作“Verdier对偶”。这名字在正文中到处重现(例如在第62、103页——只要运气好碰巧翻到……)。我发誓我没有半点编造!}。他那显赫的庇护人,在这样的喜庆日子里,可不会在区区小事上吝惜!亡者的名字在那篇文章\fnote{569}{(***)}{我的名字毕竟还出现在文献目录里,带着 EGA 这个缩写(而它总有一天该被一篇专门的文本取代……)。Mebkhout的名字也在正文和文献目录中缺席,整部书卷里都找不到他的痕迹。}里不出现,在由Teissier-Verdier署名的卷首引言里也不出现。那位含糊的无名者(Zoghman~Mebkhout,姑且不点名)的名字也不出现——而没有他,那篇文章、乃至那场辉煌的会议,本都不会问世……

论干活,那真是干了活!除了 Motif(它们不久就要跟上,就在第二年)之外,或许还有晶体瑜伽,一个从未被提名的亡者的上同调遗产之不成问题的分赃,就此告成,而且在一致同意与\textbf{皆大欢喜}中完成。

% ---------------- p.732 ----------------


\subsubsection*{c.~珍宝}

\tocsubsub{c.~珍宝}

\notehead{170(iii)}{}

\origpage{934}\dmark{}(3月1日)我在前几条注里检视的那三场“行动”,关乎我身后留下的“遗产”的分赃——那遗产,是我关于簇的上同调的、书面的与非书面的作品。这场分赃的直接受益者,是我的五位上同调学生中的三位,即Pierre~Deligne、Jean-Louis~Verdier和Pierre~Berthelot\fnote{570}{(*)}{(5月2日)事实上还该加上第四个“受益者”,那是我最近才发现的,即Neantro~Saavedra——在先前的脚注(第921页脚注(*))里已经谈过。}。但这三场行动(以及随后那第四场)中的每一场,都只能靠一大批或多少“搭在”簇之上同调这根线上的同事的串通(有时还有他们的积极支持)才办得成;而这批人中居首位的就是我的五位上同调学生——除我刚才点名的那三位,还有Luc~Illusie与Jean-Pierre~Jouanolou(*)。

这三场行动,以及即将谈到的第四场,在我看来是不可分解地缠在一起的——无论是在它们深刻的动机上,还是在它们最切实的变故上。最早的隐约征兆可追溯到1966至1968年,但更触目的表现,是在我1970年“离开”之后。这次离开,加上“大世界”数学界某种普遍的世风\fnote{571}{(**)}{(5月2日)那想必有一种双向的作用:某种心态败坏的状态(我在离开之前也曾参与其中),有利于我的一批昔日学生对我的作品进行变本加厉的掠夺与贬损,而他们日益增长的犬儒态度,无疑又反过来造成了今天我所目睹的那种多少带普遍性的腐败状态。},为这样一场规模宏大的行动创造了有利的外部条件——在我们这门科学的编年史上,这样的行动大概是独一无二的。

这场行动首先旨在\textbf{败坏}我引进数学的大部分重大\textbf{核心思想}\fnote{572}{(***)}{(5月2日)关于这方面的细节,见“我的孤儿”(n\textsuperscript{o}~45)一注、尤其是“工地的巡礼——或工具与视野”(n\textsuperscript{o}~178)一注。},并埋葬它们所嵌入的那幅统一的图景;然后,败坏或抹杀\textbf{匠人}在我所锻造的那些工具中的创造性角色——那些工具,是在这些思想的指点下、在那幅全局图景的启发下锻造的,而它们曾作为Deligne和我的其他上同调学生之作品的基础工具;最后,在最后一个阶段,侵吞那些有幸被我的学生们采用、或顶住他们加诸其上的\dmark{}抵制\fnote{573}{(*)}{(5月2日)我在那些我曾引进、后来被埋葬、却又顶住Deligne和我的其他上同调学生们设立的抵制而终于自立的思想与工具中,谨指出以下几项:导出范畴;Motif(诚然是狭隘版本);Galois--Poincaré-Grothendieck范畴的瑜伽(为着葬礼的需要被改称“坦纳克式的”);以及围绕场、gerbe 与 liens 诸概念的非交换上同调形式主义(由Giraud本着我在1955年之后引入的出发思想发展而成)。}而终于自立的\origpage{935}思想与工具的父权。

% ---------------- p.733 ----------------
这场行动在1982年以 Lecture Notes 900 那卷书的出版而告终——那部书卷奉 Motif 重返数学公共广场,形态狭隘(相对于我在六十年代渐渐看清的那幅图景),而父权则(心照不宣而显而易见)归于Deligne。它终于在第二年找到了尾声——在那份为庆祝 IHES 存在二十五周年而编印的小册子里端出的、分三个部分的《葬礼悼词》。

对这几篇文字所构成的那条“矿脉”之发现,是在去年5月12日\fnote{574}{(**)}{而正是在同一天,那场对 SGA~5 原始讨论班的无耻屠杀已经向我显露——它出自Illusie之手,并有我全体上同调学生的积极参与或踊跃串通,在“教团全体”的慈祥目光下……},在“葬礼悼词(1)——或恭维”(n\textsuperscript{o}~104)那条注里。近五个月之后,它在承续它的那条注(n\textsuperscript{o}~105)“葬礼悼词(2)——或力量与光环”\fnote{575}{(***)}{为了把对《葬礼悼词》的反思作一次意想不到的扩展,另见下一条注“肌肉与肚肠(阳葬阴(1))”(n\textsuperscript{o}~106)——它同时开启了“阴阳之钥”那场漫长的反思。}中继续。这里我只想用几句话召回这篇不寻常的“悼词”的精神与全部滋味。

那小册子(除别的之外)呈上一幅“肖像长廊”,由关于这座受庆机构历任与现任教授的一篇篇短文组成。在(出自Deligne笔下)献给我的那篇文字里——那文字本该唤起一部作品——却没有说出“上同调”或“Motif”这个词。“概形”一词也没有,也没有任何别的词能让人想到我曾发展过的一门理论、或我曾证明过的一条定理(哪怕它可能有过用处)。相反,我慷慨地得到一堆空洞的最高级赞辞与种种夸夸其谈的殷勤:“宏大的作品……”“二十卷……”\origpage{936}“最大的自然一般性……”\fnote{576}{(****)}{这法语-小黑人(petit-nègre)真是一项千金难买的妙笔,为的是以滑稽的方式(而且装作不经意地……)唤起一位庞大话痨的那份臃肿而多余的夸夸其谈……}\dmark{}“对术语的高度留意……”“他给自己的问题……沿他自己划定的路线……变得太困难了……”。这便是吹吹打打、在聚光灯下的葬礼,靠那篇送得恰到好处的“恭维”送来——它硕大而臃肿,正如那该被“缅怀”其记忆的亡者;同时又带着一股滑稽的暗示之精细,而那精细实在是那位笨拙的祖先所欠缺的……

在献给Deligne的那篇短文 里(并经他过目),没有任何东西能让人疑心我与“那”Weil诸猜想(“其难度有口皆碑”)的证明有任何关系——而那证明被恰当地大书特书。恰恰相反,文中强调:“这个结果显得越发惊人”,因为它不得不说是在与“一系列”出自我手的猜想作对(这位Grothendieck显然从来只会提这个!),而那些猜想(文中还补了一句,为免留下任何关于该如何看待它们的疑虑)“今天与当年一样高不可攀”(读作:当我竟有那个不幸的念头把它们陈述出来的时候……)。

这两幅微型肖像,再加上出色地补全它们的第三篇(仅用一句三行的凝练警句\fnote{577}{(*)}{这第三篇是我在那条已引的“葬礼悼词(2)——或力量与光环”的反思中发现的——而它立刻显得比另两篇合在一起还更有分量!正是它,启发了我给这条注取的名字“力量与光环”。})——都是些真正的珍宝,在献给一位“亡者”(而这一次他还并未过世!)的、以娴熟手法端出的《葬礼悼词》中,想必同样是同类中独一无二的。它们已在上面引过的那三条相继的注(n\textsuperscript{o}~104--106)里得到了应有的仔细检视,并在(晚几星期的)“阴的葬礼(阳葬阴(4))”(n\textsuperscript{o}~124)那条注里,在“阴阳倒转”这一更具穿透性的动态光照下重加检视。

% ---------------- p.734 ----------------
\subsection{(3) 封神(“De Rham 系数与 D-模”)}



\subsubsection*{a.~祖先}

\tocsubsub{a.~祖先}

\notehead{171(i)}{}

(3月1日与5月2--8日\fnote{578}{(**)}{(5月13日)本注与随后四条注,最初本是单独的一条注“封神”(n\textsuperscript{o}~171),写于3月1日。它还把上一条注“珍宝”(n\textsuperscript{o}~170(iii))包括在内。它曾在5月2日至8日间重新整理并大幅扩充——尤其是在其数学部分——并拆分为四条独立的注:“祖先”“作品……”“……与意外之喜”“光荣之日”(n\textsuperscript{o}~171(i)至(iv),再加上上面已提到的“珍宝”)。此外还有八篇子注(n\textsuperscript{o}~171(v)至(xii))关乎上述四条注,以及4月的四篇子注(n\textsuperscript{o}~171$_{1}$至171$_{4}$),讲述我的朋友Zoghman与“行内的法则”搏斗的种种奇遇——依照他本人向我讲的那样。正是这十六条注(n\textsuperscript{o}~171(i)至(xii)以及171$_{1}$至171$_{4}$)构成了眼下“四种运算”中的“封神”部分(而其中封神是第四场、且——在新命令到来之前——是最后一场……)。})在我那场预期中的葬礼里所区分的四场局部“运算”中的每一场里,明摆着是Deligne在充当指挥的角色(或者不如说,是葬礼上的\origpage{937}大司祭\dmark{}),我的另四位的上同调学生或多或少积极地参与,还有一大群数学家的串通——他们全都深知内情(而这内情显然不是要让他们不快的……)。这个“串通集团”在我正准备着手检视的第四场局部行动中,达到了惊人的、近乎难以置信的比例。

\begin{quote}
IV “当值无名之辈”行动(或称“反常层会议行动”)。
\end{quote}

那是一场对\textbf{Zoghman~Mebkhout作品}的行动——据我所知,他是唯一一位(在我离开数学舞台之后)敢于以“Grothendieck的延续者”身份冒险的数学家。

这场行动持续了十年上下,从1975年直到今天。冒着重复的危险,我先要唤起它的历史语境。

五十年代后半期,我在概形的语境中发展出了一套“凝聚对偶”的形式主义。这些反思——出于理解Serre对偶定理在解析几何中\fnote{579}{(*)}{我最早的对偶反思处在解析空间的框架中,且早于Serre的那些。凭着“平展式”的对偶技术与Poincaré-Grothendieck关于算子$\bar{\partial}$的引理(那是我刚证明的),我曾证明:若$X$是一 Stein 簇,则$H^i(X,\mathcal{O}_X)$(相应地$H^i(X,\omega_X)$)是 Fréchet 核空间,与$H^{n-i}(X,\omega_X)$(相应地$H^{n-i}(X,\mathcal{O}_X)$)处于完美对偶中。那时我不曾想到把同一方法套到向量丛的情形上(我没有看出那个极简单的代数事实:算子$\bar{\partial}$既是$\mathcal{O}_X$-线性的,便能延拓到取值于全纯向量丛的可微分形式),也不曾想到 Stein 之外的复簇(而它们是我当时唯一熟悉的)。Serre在一般情形下证明他的解析对偶定理,其证明与我在特殊情形下找到的几乎是同一个。}、尤其在代数几何中的确切意义与范围的愿望——是在一种近乎完全的孤独中进行的,除了我之外不曾有幸引起任何人的兴趣\origpage{938}\fnote{580}{(**)}{当然,在所有数学家中,我本该最指望Serre对我的凝聚对偶反思感兴趣。他曾感兴趣过,如果我记忆不差,是着迷于把他的对偶结果推广到域$k$上的光滑射影簇$X$上的一个凝聚层$\mathcal{F}$(未必局部自由),把$H^i(X,\mathcal{F})$的对偶等同于$\operatorname{Ext}^{n-i}_{\mathcal{O}_X}(X;\mathcal{F},\omega_X)$。这给 FAC 的一个“计算性”结果赋予了内蕴的几何意义(那个结果当年当然引过我的好奇与启发),而那是当$X$为射影空间时的情形。但除了这个成果——它是我在发现对偶的旅程中最早的成果之一,且与他所熟悉的东西还接近——之外,每当我想跟他谈对偶,Serre总是拒绝听。我想我几乎不曾试图跟任何别人谈过,除了(晚得多的)Hartshorne,他依我的想法在哈佛做了一个非常漂亮的讨论班,1966年出版(《Residues and Duality》,R.~Hartshorne,Lecture Notes in Mathematics,n\textsuperscript{o}~20,Springer Verlag)。}。正是这些反思\dmark{}使我逐步提炼出“导出范畴”这一观念 % ---------------- p.735 ----------------
——那些对象堪称各类空间与簇的同调与上同调形式主义中“天然的\textbf{系数}”,并嵌入一套关于带环空间的“六种运算”之形式主义的最初胚胎(带环的 topos 且留待以后)。这六种运算中有四种,早在1955年我那篇《论同调代数的若干问题》的工作里我就多少熟悉了\fnote{581}{(*)}{载 Tohoku Mathematical Journal, 9 (1957), p.~119--221。},只是还没有导出范畴的语言:凭着以后几年确立下来的记号(同时还有导出范畴的观点),那是些“内部”运算$L\otimes$与$\mathrm{RHom}$(层$\operatorname{Tor}_i$与$\operatorname{Ext}^i$的“总导出函子”形式,即我在《Tohoku》中引入的),以及“外部”运算$Lf^*$与$Rf_*$(逆像,与“Leray式”正像),构成两对(或两个双函子的)伴随函子。在$f$是一态射“浸入”$i:X\to Y$的情形,还要加上伴随函子$Ri_!$、$Ri^!$一对,分别体现“以零延拓”与“在$X$中(以其为支集)的局部上同调”两种运算。我的反思中的导线,是要达到一条\textbf{对偶定理}(整体的,那时还谈不上局部版本……),推广Serre为体上的光滑射影簇上的局部自由凝聚层所证明的那条。当时的问题,是要给出一个适用于任意凝聚层(或这类层的复形)、乃至准凝聚层的表述,对$X$既无光滑性也无射影性假设(只保留真态射,那在当时看来是本质的\fnote{582}{(**)}{关于这一点,见下文第940页的脚注(*)。})。此外,与我对 Riemann-Roch 定理的反思相仿,我感觉到正确的陈述不该关乎体上的一个簇,而该关乎一个随便怎样的概形的\origpage{939}真态射$f:X\to Y$。正是靠了历年的工作\fnote{583}{(*)}{不言自明,在这“几年的工作”中,我心头还压着不少别的炉火,远不止凝聚对偶的问题!那时我熟习着当时已知的代数几何奠基(以Serre的 FAC 观点为主要参考),熟习着Weil诸猜想的问题学,还熟习着在一门Serre的课上学到的相交重数形式主义——他在那里发展他的“Tor 的交错和”的想法。这一切,将要触发我在1957年投入 K-理论形式主义与 Riemann-Roch-Grothendieck 定理,而它们(就其精神而言)与我的对偶反思十分贴近。}、一次又一次的\dmark{}近似,整体对偶定理才逐渐从它那些多余的假设中澄清出来,同时导出范畴这一观念也从被预感的地带里走出,获得具体形态,并赋予那形式主义与那些陈述以一种\textbf{内蕴的意义}——缺了它,我会觉得自己根本无从工作!那正是为了达到一条令我完全满意的整体对偶陈述,我才引入了\textbf{对偶复形}的形式主义,并提炼出\textbf{双对偶定理},又发现(在适当的诺特假设下)存在一个内射的对偶复形、本质典范,我称之为“\textbf{残差复形}”,以及它的一套变差理论。整体对偶定理的一个最初陈述——有一阵我觉得那就是“对的”——是说:$Rf_*$与$X$上和$Y$上的对偶复形函子交换(对于两个互相“对应”的对偶复形)。只是后来我才发现:单就对偶复形而言(经由残差复形)的变差理论,能推广为一个性质全新的函子,即函子$Rf^!$,那个“异常的逆像”,在$X$上是局部性质的。由此,也就出现了对偶定理对真态射$f$的最终陈述:这个新函子是$Rf_*$的\textbf{右伴随},于是嵌进一串三个伴随函子

\begin{quote}
$Lf^*$, \quad $Rf_*$, \quad $Rf^!$。
\end{quote}

为求一套完全成熟的形式主义,只差一个函子$Rf_!$(对任意(分离的)有限型态射的“真支正像”)的描述——它推广浸入情形下已知的函子,对真态射化为$Rf_*$,并与$Rf^!$构成一对伴随函子$Rf_!$、$Rf^!$。我不记得五十年代曾为这一缺陷苦恼——那套形式主义超出凝聚(概形的或解析的)对偶之外的普遍范围,\origpage{940}\dmark{}当时仍为我所不及\fnote{584}{(*)}{当然,我已察觉到:即便在开浸入$f:X\hookrightarrow Y$的情形下(那里函子$Rf^!$与“限制到$X$”的函子$Lf^*$重合),后者在准凝聚层的语境中\textbf{不}容许左伴随。通常的左伴随$Rf_!$(“在$X$之外以零延拓”)不保持准凝聚性。另一方面,我也曾验证:即便抛开准凝聚性假设、甚至对一个以一个点为基的真态射,也不存在“对偶定理”。于是,在一般假设下无法定义一个$Rf_!$,这一不可能性在我看来已成定局、且出于事物之本性。

是Deligne在1965或66年(他才刚落脚!)察觉到:人们可以给$Rf_!$一个意义,并把凝聚对偶定理挽回给非真的、分离有限型态射——条件是改用一些作为准凝聚\textbf{pro-层}的(复形)系数。这个漂亮的想法却不曾得到人们本可期待的运气——正如它使之完善的凝聚对偶最初形式主义一样。

Deligne在特征零代数概形上构造“De Rham 系数”的尝试中,曾重拾这一想法并获成功;但那是个有前途的尝试,而他在我1970年离开之后不久便把它付诸流水了。六年之后,轮到Mebkhout去提炼出“那”一类好的“De Rham 系数”(晶体的)——而那是我已期待了十年之久的……}

% ---------------- p.736 ----------------
这一缺口直到1963年才充分显现在我眼前——那时我发现:在新生的平展上同调(以“离散”系数为项)语境中,存在一套在各方面都与凝聚形式主义类似的形式主义,而且还多一个函子$Rf_!$(“真支正像”),它对\textbf{任何}分离有限型态射皆有定义。况且,那正是因为我亦步亦趋地依循我前几年在凝聚情形下所做的工作(那工作除了我之外不曾引起任何人的兴趣),才得以(在一两个星期之内)从那两条基变换的关键定理出发,建立起所谓的“六种运算”完整形式主义。那是一个比此前在超越语境中所掌握的对偶形式主义(仅对拓扑簇及其上的局部系统)无比完善、无比强大的对偶形式主义,甚至比我在凝聚对偶中达到的那个更令人满意。

我的凝聚对偶工作,陈述在 R.~Hartshorne那场众所周知的讨论班\origpage{941}《Residues and Duality》(直到1966年才出版\fnote{585}{(**)}{所说的那场讨论班(出版于 Lecture Notes in Mathematics,n\textsuperscript{o}~20,Springer Verlag)陈述了我的凝聚对偶形式主义观念之主体,围绕六种运算形式主义、双对偶,以及一套“残差复形”理论(它们是对偶复形的典范内射代表)。这些观念在解析框架中被Verdier、而尤其是Ramis与吕盖重新采用。反之,Hartshorne的讨论班不曾包含若干更精细、与这形式主义密切相连的发展:一套残差理论(对任意基上有限型平坦的概形),以及一套差分的上同调理论——据我所知它们从未出版。

我还在五十年代发展了完美复形的“行列式模”形式主义,它本该最终被纳入 SGA~7,而其执笔者(照着 SGA~5 某些“执笔者”已确立的先例)两年后宣布弃权。

最后我指出:紧接我五十年代那些凝聚对偶反思,当时我也曾引入并略加发展纯代数版本的\textbf{Hodge 上同调}与\textbf{De Rham 上同调},尤其是结合于一个代数闭链(起初假定光滑)的上同调类形式主义,以及一套 Chern 类理论——仿照我在 Chow 理论中发展的那一套。}),而关于平展对偶的工作,则在 SGA~4 的一两章里、而尤其在专门献给它的 SGA~5 讨论班里。而直到写这几行的此刻,我才突然意识到:除了一些零星的先声文本(五十年代的 Cartan 与Bourbaki讨论班),不存在任何\textbf{已出版的}、出自我笔下的系统文本,陈述对偶的形式主义与瑜伽——无论是在凝聚语境中,还是在平展语境中。SGA~4 中专门论述这个主题的各场报告(围绕分离有限型态射唯一的“整体对偶定理”,证明$Rf_!$、$Rf^!$是伴随的),是由Deligne在讨论班之后两三年、依据我的手写笔记写成的\fnote{586}{(*)}{Deligne的这番整理处在 SGA~5 讨论班之\textbf{后}。况且Deligne不曾一味遵循我的笔记,而是遵循Verdier1965年在局部紧空间语境中引入的我的方法的一个变体(他大体上沿用了平展的模型)。那时在任何人心里,对偶的一切主要思想的父权、遑论平展对偶的父权,都毫无含糊;谁也不曾想到(想必连Deligne也不曾!)遵循我最初方法的一个变体,会在随后的二十年间被用来浑水摸鱼,把平展对偶归给Verdier(而Deligne则把“平展上同调……”那一“包”里的其余一口袋进自己的腰包……)}。至于 SGA~5 讨论班,它被我的上同调学生们实际禁闭了十一年,终于出版时(在Deligne1977年那部电锯文本之\textbf{后}),已被大肆掠夺、面目全非,被那位“编辑”(sic)Illusie奉其显赫朋友之命糟蹋得一塌糊涂\fnote{587}{(**)}{关于这一点,见“四种操演”一注(n\textsuperscript{o}~169$_2$),以及随后诸子注。}。

那正是在这片废墟中——它曾是我发展过的、也是\origpage{942}我全部几何学著作中最关键的讨论班之一(与 SGA~4 并列)——还留着我亲手写下、或至少依据我的手写笔记写下的唯一书面痕迹,稍稍唤起平展对偶的形式主义与瑜伽;而越过这套尚且片面的瑜伽,还有那套被它不可抗拒地暗示出的六种运算瑜伽。我的学生们细心抹去了后一套瑜伽的一切痕迹\fnote{588}{(*)}{(5月8日)我刚才匆匆浏览了我为 SGA~5 头三场报告做的手写笔记——去年应我的请求,Illusie好心把它们还给了我。(他是那些昔日执笔者中唯一肯费心把我托付给他们的笔记退还给我的人……)头一场报告是对此前的 SGA~4 讨论班就平展上同调形式主义及其与其他种种语境的关系所取得成就的一次广阔“巡礼”。第二场报告长篇宏论地发展了“六种变换”的“抽象”形式主义。那里有一份基本完备的清单,但还不曾着力去圈定那些典范同构之间的相容性。(那是一件性质更技术的工作,在一个我尤其想把“传递”出去、并清楚感受到其全部力量的对偶瑜伽之当口,它是无用的。)

不言而喻,在Illusie版中,这两场报告都不见踪影。我曾终于相信:(被讨论班那些更技术的方面缠住)我想必是疏于呈献那幅统一的图景了。回想起来,而在几乎恰好在发现 SGA~5 讨论班遭“屠杀”整一年之后,我今天似乎已把手指按在这场“屠杀”行动的要害上。那并不是这样或那样几场报告的消失——被一个Deligne并吞、被一个Verdier掠夺、被一个Serre从灾难中救出、或被一个Illusie为着纯粹的乐趣从一幅和谐的“整体”上撕走。而是这个讨论班的灵魂与要害,是这场由一个人独自完成的巨大工作中那根恒定而无处不在的导线——而Illusie着力把它从 SGA~5 中根除,不留(几乎)一丝痕迹。“六种运算”这个名字本身,在这场讨论班中缺席,正如它在我的学生们的著作中缺席——他们想必订有默契的约:只在极罕见的场合才说出这几个字,即其中某一位还与那位被宣告亡故的匠人(而他虽已亡故,却毕竟还得被蒙在鼓里……)狭路相逢之时。},它在整个六十年代以异常的暗示力量,给我的上同调工作以灵感。那才真正是“系数类型”\fnote{589}{(**)}{这个核心思想同样被我的上同调学生们根除并遗忘了。它也是在我对自己的作品及其坎坷作了“十五年后”的第一次回顾之后,最早重新强加于我的那些思想之一,就在“我的孤儿”(n\textsuperscript{o}~45)那条注里。那条注的名字,比我当年所能梦见的更贴切、更深刻,而它却写在发现“葬礼”(这个词的本义与强义)之前。六种运算与“上同调系数”这同一个核心思想,当《收获与播种》的反思重新让我与我的作品在那些曾是我学生们的人手中所受的命运接触时,几乎作为一个主旋律在处处回响。尤其见“墓上的旋律——或自负”(多少展开“系数类型”的“旋律”或变奏主题)与“工地的巡礼——或工具与视野”两注(n\textsuperscript{o}~167、178)。}这一核心思想的“要害”——而 Motif 的瑜伽,正是它的灵魂……

\origpage{943}这样一种畸形的局势——一门科学中的重要进展,体现于一幅新图景,却被那些曾是它第一批受益者与保管者的人亲手根除——若不是另一种局势,也同样堪称例外,即由我的骤然离开及其周遭条件所造成的那种局势,是不可能确立的。何况,事态的走向,早在我离开之前、在整个六十年代,就已由我所处的那种分裂局面预备好了——一方面,我被无穷无尽的基础工作占着,而那工作只有我能或才愿意承担\fnote{590}{(*)}{我提醒一句:这项维度宏大的基础工作,自我离开那一天起便戛然而止,直到今天依然如此。这是我在“藏宝”一注(n\textsuperscript{o}~169$_5$)里所谈的那个“误会”之一个雄辩的标志。人人都已准备好搬来自己的家具、永久住进我建造的那些房子里——但要再动手、拿起泥刀与铅垂线去继续建造和布置,哪怕只是在种种需要的紧迫压力之下,已经没有一个人了……},另一方面,又被一些主题往往离得很远的问题不断叩问……

% ---------------- p.738 ----------------
那些当时吸引着我的基础工作,以及那些因此往往更强烈、更直接地迷人的东西\fnote{591}{(**)}{若我当时肯听自己的,有多少次我会撂下那些我为众人操劳的没完没了的基础工作,去投入那场不断召唤我的未知冒险——那真正的冒险——而不把由我发现的新土地让别人去游逛的乐趣让给人。今天我看到:那些土地始终是处女地,或差不多,而那些我当年看作先驱的人,早在我离开之前就已选择做安逸的食利者……}即便是那些我曾让自己有空去深究和发展的主题(比如对偶),我也少有那份余暇,把研究结果整理成宜于出版的形式(合乎我那高标准的要求)。于是,我时常索性把整理的工作留给别人(那些我全然信任的人)——就像“对偶”这个主题,无论在凝聚框架还是平展离散框架中都是如此;或者把一些我知道是丰产的出发思想交给别人去发展(比如导出范畴的观念、或晶体上同调,这里只举这两个为例,同类例子还多得很)。在一种“正常”的局面下——即\origpage{944}一片诚意回应着我向一批有干劲的学生发出呼吁时的信任,他们在我的接触中学习他们的手艺、并为他们未来的作品学些有分量的根基——一切本是最好的\dmark{},对一切相关的人都是最大的好处,科学共同体也不例外。但实情是:这种不寻常的局面,把一份可观的\textbf{权力}放到了他们手中(而直到去年之前,这个念头从未在我脑中掠过……),尤其在我离开之后。我离开之后不久(甚至就在之前……),他们中某些人就忙不迭地滥用这份权力,去抹杀作品与图景,去贬损匠人,并把他们以为用得上的一些他亲手打造的工具据为己有。

我的凝聚对偶工作,似乎从未大受欢迎\fnote{592}{(*)}{如我在先前一条脚注中指出的,这些工作曾启发Verdier、Ramis、吕盖在解析空间的凝聚理论方面的工作。“一向以来”就清楚(至少对我如此):同一套形式主义必能在刚性解析语境中重现(而那个语境,据我得到的回声,也还停留在童年阶段)。

另一方面,Mebkhout告诉我,日本的分析学派相当大地从《Residues and Duality》中汲取了灵感——同时却绝口不提那位匠人的名字。凭眼前的时风,反过来的话倒要叫人惊讶了……}。相反,平展对偶的工作却立刻引来注意。但我认为更公允的说法是:引人注意的是,有人“想明白了”,不管用什么办法,在平展语境中证明了Poincaré对偶的类似物——而那对偶,我想人们在熟悉的定向拓扑簇的语境中已经认识了近百年。那对平展上同调来说曾是“一个好得分点”(人们几乎不再怀疑:对Weil诸猜想——“其难度有口皆碑”——来说那才是“对的”)。这就是说,那群在著名的猜想旁逡候着的数学公众,是以一个“消费者”的身份反应的——他不情愿承认并消化一幅新的、深刻的事物图景,只取一个面目熟悉的“结果”。二十多年后,我确认:六种运算与系数类型这个富有力量的图景(以一套简单得令人困惑的形式主义表达),仍然被所有人无视(唯一的例外是那位孤独的匠人);而每当有人敢提一句,它就成了冷嘲热讽或刻薄挖苦的对象\fnote{593}{(**)}{至于细节与评注,见子注“无用的细节”(n\textsuperscript{o}~171(v)),尤其是(a)部分“上千页的一捆……”。}。我那套装备中这样那样的零散成分,在这里那里被人使用,却丝毫不提我个人(而且早已找好的替代父亲),尤其是双对偶形式主义——自从\origpage{945}1981年那场难忘的会议之后、在相交上同调的大冲刺里。但那幅\dmark{}\textbf{图景}——带着孩子般的简洁与完美的优雅,而它曾给出过自身力量的有力见证\fnote{594}{(*)}{至于这些“有力的见证”的细节,见子注“无用的细节”(n\textsuperscript{o}~171(v)),(b)部分“无所事事的机器……”。}——却依然无人识得,只是那些宁可用轻蔑(和掠夺……)代替理解的人轻蔑的对象。

% ---------------- p.739 ----------------
如果说我亲手、并用心所做的东西,比它的时代超前了二十年、或许五十年,那不是因为我三十年前动手时遇到的那门\textbf{数学}还不成熟。那是因为人的不成熟\fnote{595}{(**)}{关于这个问题的一点初步反思,见子注“自由……”(n\textsuperscript{o}~171(vii))。}。而我的身后学生、我唯一的延续者Zoghman~Mebkhout,所面对的也正是同一种不成熟。我很有幸,在1970年离开之前,只以“不理解”的形式遭遇过它——而那不理解从未卸下友好的姿态。Zoghman~Mebkhout,来到一片与他所大胆延续的那份工作时代不同的数学场域,他在不理解与轻蔑之后、在他的某一项成果作为工具的价值终于被承认之后,还领教了前辈们的恶意与一个时代全部的不义——但我抢跑了……

我在数学中最重要的发现之一、而它至今仍几乎被所有人无视的,是那套我在五十年代开始发展的对偶形式主义的\textbf{无所不在}:“六种变换与双对偶的形式主义”同时适用于人们最初设想的“连续”系数(“凝聚”理论),与“离散”系数。这种无所不在,在1963年春显得像一次几乎难以置信的惊喜——正是靠了它、而非别的,我才得以发展出一套平展对偶形式主义,并达到我所谓的对平展上同调的“掌握”。而且早在那个时候,我就被一个问题困扰着(虽说不曾太停下来):一门“共同的”理论的可能性——无论是在概形框架、复解析框架、乃至拓扑框架中;一门能“笼住”两类系数的理论。De Rham 上同调(我的一位老朋友……)给出了这方面的第一个迹象,提示人们到“带可积联络的模”(或“分层模”,或许……)的方向去寻找一个“共同原则”。这些模生出一门“De Rham 上同调”(实质上以离散系数为项),于是被发现与凝聚上同调\origpage{946}处在关联之中。这个进路后来向我提示了“晶体”与“晶体上同调”的念头,然而它似乎还不足以(看来)提供那把钥匙,去描述一套完整的六种变换形式主义——针对那些在适当意义上同时涵盖离散(“可构”)系数与连续系数的“系数”类型。我的学生中,似乎\origpage{947}没有一个懂得感受这个问题\fnote{596}{(*)}{写这几行时,我关于这事的记忆还模糊。它后来变得清晰了,而我在子注“荒唐的问题”(n\textsuperscript{o}~171(vi))里更周详地回到这一点。}\fnote{597}{(**)}{我曾向Verdier谈过这个问题——那时他已(如我所建议的)发展了拓扑空间的(至少是其胚胎的)对偶理论,仿照我在平展语境中发展的那一套(关于这一点,见子注 n\textsuperscript{o}~81、81$'$)。那该是六十年代中期。显然当时没有“对上暗号”——那个问题的意义(或许有点含糊,实情如此)似乎从他嘴边溜过了。可我必定提到过 De Rham 上同调,既可微的、也复解析的,它把Serre对偶与Poincaré对偶联系起来,关乎这两种类型的系数。

(5月14日)况且,早在五十年代我就知道:可以把Serre的对偶定理推广到真光滑相对概形上局部自由层之间的微分算子复形的情形,以便也把 De Rham 上同调涵盖进来(因而,实质上,一门以离散系数为项的上同调)。那是一个与Mebkhout在解析框架中的对偶结果非常接近的对偶结果——而后者将在下一条注里谈到。当时我不曾沿着这条路走下去,尤其因为(我相信)我看不出该怎样用这些微分算子复形做出一个合适的“导出范畴”——因为缺乏一个好的“拟同构”观念。实情也是:我当时工作的孤独——做着一些(凝聚上同调)显然除我之外全世界无人感兴趣的问题——几乎不足以刺激我把又一项推广(用微分算子代替线性映射)堆到我在前几年于自己角落里已提炼出的那些推广之上。然而从Mebkhout的视点看,我本是万事俱备:那里通向相应的 D-模(一个微分算子复形的各个分量)的过渡,给出了一把具有完美简洁性的钥匙,去构造所需要的导出范畴。况且,早在1966年(尽管当时我没有清楚地意识到),我就有了一个对偶的视点,它本可以让我用“分层 pro-模”造出一个导出范畴(这个观念后来由Deligne在他那套 De Rham 系数的草图中发展出来,下文将谈到)。的确,给任意一个凝聚模结合其\textbf{无穷阶主部}所构成的 pro-模,后者带有典范的分层;或者给一个微分算子复形结合一个由这类分层 pro-模组成的复形,其晶体超上同调便与所考虑的微分算子复形的Zariski超上同调相同。(参见我的报告“Crystals and the De Rham Cohomology of schemes”,由 J.~Coates 与 O.~Jussila 记录,收入《Dix exposés sur la cohomologie des schémas》,North Holland,第~306--358~页,尤其是第~6~节。)于是,可以按通常的方法,借助相结合的分层 pro-模复形,为微分算子复形之间的(微分)态射定义“拟同构”的概念。}

% ---------------- p.740 ----------------
\dmark{}只有Deligne例外。他用了一整年讨论班(我想是在 IHES,1969/70年)去发展一套形式主义——那套形式主义至少使他能(对特征零体$k$上有限型的概形$X$)描述一些(所谓“De Rham 的”)上同调空间,它们在$k=\mathbb{C}$的情形下,重现通常的(以超越方式定义的)复“Betti 上同调”。他所用的系数是“分层 pro-模”以及这类 pro-模的复形。然而当时并不清楚这些系数是否能嵌进一套六种运算的形式主义\fnote{598}{(*)}{这里我的记忆仍然是模糊的,而且是个错误——先验地清楚,出于超越性质的启发式理由,这里\textbf{必定}存在一套六种运算形式主义。(细节见子注“……与阻碍”,n\textsuperscript{o}~171(viii)。)

我的错误显然来自一个刻意的(浮在意识表层的)说法:把一件或许看似无法解释的事情合理化、讲得通——即Deligne放弃一条“稳妥”而富于希望的研究方向。那理由绝非数学性质!},而Deligne放弃了继续走这条路。如果我没记错,当时最缺的(*)、令人能放心的一着是:用一种纯代数的术语(借助凝聚或前凝聚模与分层),对$X$上“代数可构”的$\mathbb{C}$-向量空间层之范畴做一个描述——该描述因而对任何特征零的基体皆有效,而这个范畴在基域为复体$\mathbb{C}$时本是以超越方式定义的\fnote{599}{(**)}{我提醒一句:这个可构性观念是我在五十年代引进的,带着众多变体(代数的、实解析的等等),而那正是我严格地独自一人对这些题目感兴趣的时刻。(见我去年的评注,在子注 n\textsuperscript{o}~463 中。)}。



\subsubsection*{b.~作品……}

\tocsubsub{b.~作品……}

\notehead{171(ii)}{}

Mebkhout的工作始于1972年,处在(技术上更艰难的)解析空间的超越语境中。在以后几年里,他在一种几乎完全的孤独中,熟悉我的上同调作品与\origpage{948}导出范畴形式主义\fnote{600}{(***)}{(5月14日)Mebkhout后来向我说明:1972年前后对数学文献的这第一批阅读,是日本Sato学派作者们的著作。他费了很大劲(如他所说)才理清头绪,那在他看来复杂得可怕。而他正是在那里找到了Hartshorne《Residues and Duality》一书的指涉——读那本书对他是真正的放松。诚然,那本书写得极其出色!我为那书写的几行引言,唤起其中所发展之形式主义的无所不在,大大启发了他。正是从那时起,他着手熟悉我的作品,而它后来成为他主要的灵感来源。在这一切工作与报告中,他留心清楚地指明这个来源。}——它们被那些曾经是我学生们的人弃置不顾。\dmark{}一条导线,在他的反思中渐渐取得第一重要的位置,那就是连续对偶与离散对偶之间惊人的平行。后者在这期间取了“Poincaré--Verdier对偶”的名字,而大世界里没有人(尤其是那位新“父亲”Verdier)装作对这个平行的深刻理由有所疑问。那是“实用”而短视的观点当道:使用我早就创造好的现成装备,而不提任何问题——尤其不提那种模糊至极、甚至荒唐的问题。这个问题在任何已出版的文本中都不曾提及,甚至(而我意识到这里是在自责……)在我自己笔下的文本中也不曾\fnote{601}{(*)}{(5月14日)然而我记得:在 SGA~5 讨论班期间,我心中始终挂着我那套形式主义的无所不在,而我不放过一个机会,指出我在离散平展上同调框架中所发展的观念与技术在其他语境中的可能变体。我很难相信,在口头讨论班上我竟没有提过综合这两类系数的问题,哪怕只是在最后那场关于未解决问题的报告中——它也从那部屠杀版中消失得无影无踪。不用说,这部精心删去了所有不符合“技术性离题文集”这一严整标签的内容的版本,绝不会提到这类问题……\par (5月19日)关于此事,另见子注“死去的书页”(n\textsuperscript{o}~171(xii))。}

% ---------------- p.741 ----------------
共同形式主义的表述本身,就离不开导出范畴。Mebkhout始终把它们作为工作工具,顶着时尚的风向与前辈们的轻蔑——首当其冲的,还是那位不知情愿还是勉强地充当这些范畴“父亲”的Verdier。相对于我引入的整套工具,Mebkhout新增的关键成分,是Sato及其学派的微局部分析。确切地说,他借用了他们在光滑复解析簇上的 $\mathcal{D}$-模概念(它等价于我在1965--1966年前后引入的“模的晶体”概念,而后者在更广泛的语境中仍有意义,尤其适用于奇异簇),更借用了 $\mathcal{D}$-凝聚性的概念以及凝聚 $\mathcal{D}$-模上精微的完整性条件。此外,他实质性地运用了Kashiwara1975年的一个定理:与一个\origpage{949}\dmark{}完整 $\mathcal{D}$-模相结合的微分算子复形,其上同调层是解析可构的。在Mebkhout两年前向我讲述之前,这些观点与结果我一无所知;Deligne在1969--1970年间思考德拉姆系数形式主义而没有继续下去时,想必同样不知道这些。正是\textbf{把这两股观念汇合起来},Mebkhout才得以统一把握光滑复解析簇 $X$ 上这两类系数,用微分算子复形来表达,或者更好、更精确地,用较为灵活的 $\mathcal{D}$-模语言,表述为具有凝聚上同调的 $\mathcal{D}$-模复形。\fnote{602}{(*)}{关于 D-模语言的细节、它与微分算子复形语言及晶体语言的关系,见子注“五张照片(D-模与晶体)”(n\textsuperscript{o}~171(ix)),(a)部分。}这正是他对当代数学的伟大贡献。

说得更确切些:若$X$是一光滑复解析空间,记$\mathrm{Cris}_{\mathrm{coh}}^*(X)$为导出范畴$D^*(X,D_X)$中由具有$D_X$-凝聚上同调的$D_X$-模复形组成的全子范畴,$\mathrm{Cons}^*(X,\mathbb{C})$为导出范畴$D^*(X,\mathbb{C}_X)$中由$X$上具有解析可构上同调的$\mathbb{C}$-向量空间层复形组成的全子范畴,而$\mathrm{Coh}^*(X)=D^*_{\mathrm{coh}}(X,\mathcal{O}_X)$为由具有凝聚上同调的$\mathcal{O}_X$-模复形组成的全子范畴。Mebkhout揭示出如下的基本函子图表:
% 原书图内的 coh 下标与正文的记号不完全一致,这里照录原图.
\begin{equation*}
\begin{tikzcd}[row sep=2.5em, column sep=1.7em]
\mathrm{Cons}^{*}_{\mathrm{coh}}(X,\mathbb{C}) \arrow[dr,"M"'] & &
\mathrm{Coh}^{*}(X) \arrow[dl,"N"] \\
& \mathrm{Cris}^{*}(X) &
\end{tikzcd}
\tag{Meb}
\end{equation*}

其中右边的函子$N$是“重言式的”函子——由显然的包含$\mathcal{O}_X\to D_X$所作的标量扩张函子的总导出。左边的函子$M$,或称“Mebkhout函子”,性质要深刻得多\fnote{603}{(**)}{至于一个紧密相关的函子$M_\infty$(在$D^\infty$-模语境中)的“显式”描述,见已引子注 n\textsuperscript{o}~171(ix)的(b)部分:“好上帝公式”。}。它是\textbf{全忠实的},而它的本质像,是由那样一些$D_X$-模复形组成的$\mathrm{Cris}_{\mathrm{coh}}^*(X)$的全子范畴——那些复形的上同调层不仅凝聚,而且还“完整的”且“正则的”。这些是精微的局部条件,前者由Sato学派引入,后者\origpage{950}由Mebkhout临时定义\fnote{604}{(*)}{“正则”这个名字当然是从单复变量函数的微分方程的“正则奇点”的经典术语借来的。若$i:U\hookrightarrow X$是补集$U=X-Y$(对$X$中的一个除子$Y$)的包含,则Mebkhout意义下的(对一个$X$上的 D-模复形$C$)“沿$Y$的”正则性,可以这样写:典范态射
\[
Ri^{\text{mér}}_{*}(C_U)\longrightarrow Ri_{*}(C_U)
\]
从$C$在$U$上的限制$C_U$的“亚纯正像”到通常的正像的那个典范态射,对所结合的 De Rham 复形诱导一个拟同构。

在$F_U$化为一个“局部系统”(即一个带可积联络的$\mathcal{O}_U$-凝聚层)的情形,这个观念等价于Deligne的那一个。后者显然也是受我的比较定理启发(其区别在于:Deligne绝口不提这一点,而Mebkhout则始终留心清楚地指明他的来源)。Mebkhout是在引入他自己的(超越性质的)定义之后,才得知Deligne的观念的。他先前不曾为他的条件寻找一个纯代数的描述。Deligne的工作表明:在所设想的特殊情形中,Deligne的代数条件蕴含Mebkhout的那个,而Mebkhout验证逆命题也成立。这就即刻提供了钥匙,去为Mebkhout的正则性条件作一个纯代数的描述——对任何一个具有凝聚上同调和完整的 D-模复形。

Mebkhout告诉我:日本人有一个“具正则奇点的微微分系统”的观念,他们把它用在全然不同的精神里(为的是分析、而非几何的需要)。在“好上帝定理”的大冲刺之后,那是一个现成的手段(在众多别的之中),为的是打乱牌面并抹杀Mebkhout的先驱工作。看起来两个观念是等价的——而鉴于这个题材中被蓄意造成的混乱状态,很可能谁也不曾费心去验证这一点。Mebkhout始终只用他1976年引入的那个正则性观念工作(它见于他的论文,两年后提交)。},而他尤其(如他所说)是从我的代数 De Rham 上同调与解析 De Rham 上同调之间的比较定理——为了我曾在离散平展上同调(即复 Betti 上同调)的框架中发展的那些观念与技术——得到启发。

% ---------------- p.742 ----------------
这些条件(而在我看来这正是它们的主要趣味所在)事实上是“纯代数的”,尤其在$X$被换成特征零的任意体上有限型概形(要光滑也无妨,但不必需)的情形下仍保持意义。

Mebkhout函子$M$(或称“好上帝函子”\fnote{605}{(**)}{至于“好上帝定理(或函子)”这个名字的来历与意义,见“当值无名之辈与好上帝定理”一注(n\textsuperscript{o}~68$'$)——那注写在我知道反常层会议那套神秘化之前、甚至在我知道“以其全部辉煌的葬礼”之前。})可描述为如下函子的拟逆:

$m : \mathrm{Cris}^{*}(X)_{\mathrm{hol.reg.}}\longrightarrow\mathrm{Cons}^{*}(X,\mathbb{C})$,

由以下式子定义

$m : F \longmapsto \mathrm{DR}(F)\overset{\mathrm{def}}{=}\mathrm{RHom}_{D_X}(\mathcal{O}_X, F)$,

\origpage{951}\dmark{}即(定义在整个$\mathrm{Cris}_{\mathrm{coh}}^*(X)$上的)那个把每个(具有凝聚上同调的)$D_X$-模复形结合到它所结合的微分算子复形(或“De Rham 复形”)的函子的限制\fnote{606}{(*)}{关于这一点,见已引的“五张照片(晶体与 D-模)”一注(n\textsuperscript{o}~171(ix)),(a)部分“‘De Rham 系数’相册”。}。Kashiwara的可构性定理蕴含:当$F$是完整的(从而,当它完整且正则时更是如此),$\mathrm{DR}(F)$确在$\mathrm{Cons}^*(X,\mathbb{C})$中,这使得函子$m$可以定义——那定义固然不言自明、孩子气,然而除了Mebkhout(而且直到五年后的“大冲刺”之际还是如此……)谁也不曾想到\fnote{607}{(**)}{(5月7日)人们不得不把那\textbf{两个}函子$m$、$M$——它们在一个方向和另一个方向建立起关键范畴的等价——称为\textbf{Mebkhout函子};$D^\infty$-模的函子$m_\infty$、$M_\infty$也同样。(关于后者,见所引的“五张照片”一注(n\textsuperscript{o}~171(ix))的(b)部分。)把这些函子与自然的对偶函子复合,人们得到另外两对互为拟逆的函子$(\delta,\Delta)$与$(\delta_\infty,\Delta_\infty)$,它们是反变的,在某些方面更方便(见所引注)。那四个就是“Mebkhout反函子”。}!那要先记得某一种瑜伽——导出范畴的瑜伽,而大家不约而同地决定把它连同引入它的那位亡故者以及他别的同类夸夸其谈一起埋葬\fnote{608}{(***)}{(5月7日)不止一次,Mebkhout被当作一个滑稽人——他以为在导出范畴之间写箭头(您倒是评评这个理!)再加上些$\mathrm{RHom}$,那就是做数学……可他并没有因此动摇,正如我当年(1955年)引入层的整体与局部$\operatorname{Ext}^i$(以及后来带$\mathrm{RHom}$或不带的)时一样——那些让所有的人都晕船,并为他们对我的最直言不讳的保留提供了理由(至少直到1957年,即 Riemann-Roch-Grothendieck 之年……)。

这一切都不曾妨碍Mebkhout相信自己的嗅觉,并跟着它走到它带他去的地方。他赤手空拳地投入工作,没有经验,没有任何人的帮助。他\textbf{确信}他预感到的定理必真——他手上的一切迹象都是相互印证的。稍许有点经验的话,就连他早已握有证明它的一切手段都会是显然的——用如今标准的手段,我的学生中随便哪一个三下两下就能证明。

然而独自一人、只能靠自己,那定理在他看来远得眩目、高不可攀——他几乎不敢指望有朝一日能证明它!

如果说他为证明它辛苦了近两年,那是因为他不曾具备我的学生们曾具备的优势:有一位慈善的长者搀扶,并在我的接触下学到某一种标准技术——可构层的拆卸技术,加上Hironaka的奇点消解。他所提炼出的陈述诚然深刻,而证明也同样深刻,但今天性质上已是标准的。回想起来,他当年要克服的困难主要是心理上的,更甚于技术上的:逆着潮流工作,而且完全只靠自己的光明……}

% ---------------- p.743 ----------------
\origpage{952}\dmark{}此外,那个\textbf{正则性}条件——超出完整性之上的——是由Mebkhout“量身”提炼的,其方式恰恰使得函子$m$受到这样的限制后,可望是全忠实的、甚至是范畴等价。他在1976年就达到了这个信念。他终于在1978年初、在他的论文中以一个非常相近的形式把它证明了出来\fnote{609}{(*)}{(5月5日)在他的论文里,Mebkhout对$D^\infty$-模陈述并证明了相应的等价定理,并给出了拟逆函子$M$的一个卓越的显式表达。关于这一点,见子注 171(ix)((b)部分),以及子注“一个视景的萌发——或闯入者”(n\textsuperscript{o}~171$_{1}$)。Mebkhout早在1976年就达到这样的信念:两个函子 $m$、$m_\infty$(因而还有标量扩张函子 $i$,即上面最后引的那条子注里谈到的那个)是等价,并得到了 $m_\infty$ 的拟逆函子的显式形式。他那论文里关于$m_\infty$的结果,是1978年的。从那时起,他就握有证明(类似的、但带着额外技术困难的)$m$情形的全部配料。

鉴于他的论文(1979年2月通过)所受到的一片冷淡,他那时不曾费力也为$m$的情形写一份像样的证明。那些配料与$m_\infty$的相同,并且受我的复代数簇的 De Rham 上同调比较定理之证明启发(他是1975年得知那个证明的),以及 SGA~5 的拆卸技术(他是从Verdier那“好的参考文献”里学到的,而 SGA~5 讨论班仍在我的好上同调学生们的小心禁闭中)。只是到1980年末,鉴于他的想法对 Kazhdan-Lusztig 猜想的证明日益重要,他才费力为$m$的情形(那里人们并不预先掌握一个拟逆函子)写了一份周详的证明。那证明出版于《另一个范畴等价》,Compositio Mathematica 51 (1984), p.~63--88(手稿收到日期 1981年6月10日)。
\par 我强调:1975年至1980年间,除了Kashiwara在1980年用几行文字提及它(下文“黑手党”子注,n\textsuperscript{o}~171$_2$,将谈到),在Zoghman~Mebkhout自己的著作之外,文献中没有任何地方谈到函子 $m$ 或 $m_\infty$,也没有谈到一种对偶“哲学”,把解析可构离散系数与完整正则 $\mathcal{D}$-模复形,或完整 $\mathcal{D}^{\infty}$-模复形精确地联系起来。正如我们将看到的,当这一联系的重要性终于因Kazhdan-Lusztig 及Deligne主持下的相交上同调热潮而得到承认时,Zoghman~Mebkhout的名字却在一份轻声细语、面带微笑、低调而又有效得冷酷无情的默契中,悄无声息地被抹去了……}。

\origpage{953}\dmark{}那才是Mebkhout带来的伟大的新定理,是八年顽强的、在完全孤独中持续的工作之加冕。它在一个简练的陈述里,包含着一整系列深刻的、一般性递增的结果,那些结果在1972到1980年间被耐心地提炼并逐个证明。至于这场孤独的旅程中发现“簇上同调中一种新哲学”的几座大路标,我请读者参看子注“三座路标——或天真”(n\textsuperscript{o}~171(x))。在这条注里,我的目的主要是用几句话描述那幅新的图景——它呈现在这位孤独匠人劳作的第一个漫长阶段之末——即Zoghman~Mebkhout。

关键的事实(Mebkhout早在1976年就清楚认识到的)是:范畴$\mathrm{Cons}^*(X,\mathbb{C})$(性质上“拓扑的”)可以借助于Mebkhout函子$M$,被解释为范畴$\mathrm{Cris}_{\mathrm{coh}}^*(X,\mathbb{C})$的一个全子范畴——而后者在“抽象”代数几何的框架中仍保持意义;它也可以“实质上”被解释为一种“导出范畴”,由通常意义上的微分算子复形构成\fnote{610}{(*)}{关于这两个视点之间的精确关系,我请读者参看被大量引用的子注“五张照片”(n\textsuperscript{o}~171(ix)),(a)部分。}。这个全子范畴由
% ---------------- p.744 ----------------
完整性(Sato式的)与正则性(Mebkhout式的)条件定义,显然就是我在六十年代就已预见的那种“De Rham 系数”的好范畴——它当年还缺失于我的装备中,在特征零的情形,为的是补全并相互联结(如一幅巨大的扇面)我在1963年提炼出的那些“$\ell$-adic 系数”;那也是Deligne在六十年代末曾试图把握而(似乎)不曾以一种令他满意的方式做到的范畴。这个范畴显然将在代数几何中(尤其是在基概形$X$上 Motif 范畴的描述中……)扮演一个本质的角色。在我看来,至少对我来说,这个范畴当得起的名字是“\textbf{De Rham-Mebkhout 系数范畴}”\fnote{611}{(**)}{对这个范畴的关键角色与意义的普遍不理解,早已显明在这样一个事实中:它至今还不曾得到一个名字、也不曾得到一个简练的记号。代替它的(在我看过的文本中),作者们只是含糊地指到“完整正则的微分系统”(愿有好运气的人在其中找得到门路!),指到这些系统与 $\mathcal{E}$-可构层之间的“构造”“对应”或“(据说众所周知的)关系”——而且当然,照例要严格地沉默掉那位孤独的匠人,是他发动了这场围绕数学俊髦世界的新热门点心“D-模”的大喧闹。},记作$\mathrm{DRM}^*(X)$(或$\mathrm{Meb}^*(X)$),或在概形框架中、当$X$是特征零的体$k$上有限型概形时,\origpage{954}记作$\mathrm{DRM}^*(X/k)$(或$\mathrm{Meb}^*(X/k)$)\dmark{}\fnote{612}{(*)}{在代数情形,除了局部的“正则性”条件之外,还必须加上一个“无穷远处”的正则性条件(在非真簇的情形),才能找到那些“好的”De Rham-Mebkhout 系数——它们在基域为复体时,将对应到$X_{\mathrm{an}}$上具有\textbf{代数}(而不只是解析)可构上同调层的$\mathbb{C}$-向量空间层复形。也正是为着这些系数,人们才有一条“比较定理”,推广我关于 De Rham 上同调的结果,即:晶体“总上同调”$R\Gamma_{\mathrm{cris}}$,无论取代数(Zariski)观点还是取超越意义,都是“同一个”。这个陈述本身又该被看作一个更完备的陈述的特例,即:代数观点的“六种运算”与超越观点的六种运算是“相容的”。

如果我的学生们不曾那么忙着埋葬师父的作品,那该是七十年代最初(若不及六十年代起……)他们就该提炼出那门势所必至的系数理论,以它全部的简洁与全部的力量……}。

正是通过上面那个函子图表(Meb)——它概括了Mebkhout的哲学(可追溯至1976年,并由他在往后几年中确立)——那些\textbf{凝聚晶体系数}(即$\mathrm{Cris}^{*}_{\mathrm{coh}}(X)$的对象)才能被看作是“离散”(可构)系数与“连续”(凝聚)系数的一种“共同推广”。由前者构成的范畴,无论如何通过Mebkhout函子$M$(一个性质深刻的函子),等同于由 De Rham-Mebkhout 系数构成的晶体凝聚范畴的\textbf{全子范畴}。情形对那个重言式的函子$N$就不那么好,它毫无全忠实之处。但为了安慰我们并把图景补全,可以补充说:在所涉及的每一个范畴中,人们都有一个自然的\textbf{对偶函子},它给出一条双对偶定理(对$\mathcal{O}_X$-模与$D_X$-模是“平凡的”,而在可构$\mathbb{C}$-向量空间层的情形要用到Hironaka奇点消解的全部力量)——依我当年在凝聚(交换)框架中、而后在(\origpage{955}1963年的)平展离散框架中所提炼的模型\fnote{613}{(**)}{(5月5日)把我的双对偶结果、以及可构性在$\operatorname{RHom}$之下的稳定性,从平展语境推广到解析语境,那其实是自动的、而且我早在1963年就已知道;那时Verdier已与我共事三年,正投入导出范畴(他已承担了为它做系统理论的工作)与凝聚对偶的瑜伽。正是从我口中,他学到了那些把凝聚对偶形式主义推广到离散系数情形的手法。如我们所见,十三年之后(1976年),他在“好的参考文献”中占有了那套解析复语境中的对偶与双对偶瑜伽,并与Deligne及其他上同调学生串通——他们全都深知内情。

在第二年(1977年)的 SGA~5 屠杀版中,Illusie保留了(在报告~I 中)双对偶定理,因此对同时读过两种文本的读者来说,Verdier的骗局显而易见——但显然所有人都视之为寻常(鉴于眼前这年景……)。相反,Illusie回避把可构性经由$\operatorname{RHom}$的稳定性这一结果收进去——而那当然是我在陈述并证明双对偶定理之\textbf{前}就给了的,而Verdier抄去的我的证明丝毫不依赖它。于是(这种事真做得出来!)Illusie只满足于在第二个论元是对偶复形时确立所说的稳定性!!那是替他的朋友Verdier遮掩的一种方式,使那部从头到尾(除了在他该提的地方提过的那三页)抄自我的 SGA~5 各场报告的文章,不那么显得出这一点。最美妙的是:所说的稳定性早已是对偶形式主义的一个直接推论(这并不妨碍:装作只对$G$是对偶复形时确立可构性经$\operatorname{RHom}$的稳定性,在数学上是荒唐的)。但那位殷勤的Illusie在他的报告里绝口不提这个推论,好保住那副样子:似乎出现在他那同伙的“好的参考文献”里的稳定性结果,确确实实是他的独创。

人们不禁要问:既然如此,Illusie为何还是保留了双对偶定理——要屠杀就屠杀到底,他为这一点不周又算得什么呢!但假如他把它掏空,他就势必也得同时掏空那条永恒的Lefschetz--Verdier公式(而它对那定理的运用是本质的)——那恰是“特洛伊木马的头颅”:那条公式,正是它在 SGA~5 中所谓的关键角色,该为他的另一位同伙那场无礼的“电锯”行动辩护,而那行动炸碎的是我的平展上同调作品的统一性。

再恭喜一次我的昔日学生Illusie,那位机灵的“编辑”-掘墓人……

};
% ---------------- p.745 ----------------
此外,这两个函子$M$、$N$是与自然的对偶函子相容的\fnote{614}{(*)}{对重言式的函子$N$而言,这个相容性本身是重言式的。相反,对Mebkhout函子$M$(或者说,对它的拟逆 $m=(G\mapsto\mathrm{DR}(G)=\mathrm{R\mathcal{H}om}_{\mathcal{D}}(\mathcal{O}_X,G))$)而言,那是一个深刻的结果,由Mebkhout在1976年证明(名为“局部对偶定理”),与 D-模的整体对偶定理同时——后者下文立刻就要谈到。可这并不妨碍“大家”现在把这个结果当作理所当然,而且尤其是(这就更理所当然了呢)丝毫从不提及某位含糊的无名者……}

此外,若$F$、$F'$是在$X$上处于对偶的晶体系数,Mebkhout证明:$F$与$F'$在$X$上的$\mathbb{C}$-向量空间“晶体上同调”复形\fnote{615}{(**)}{我提醒一句(参看“五张照片”,n\textsuperscript{o}~171(ix)):$F$在$X$上的(“绝对的”)晶体上同调定义为
\begin{align*}
R\Gamma_{\mathrm{cris}}(F)
 &\overset{\mathrm{def}}{=}\mathrm{RHom}_{\mathcal{D}}(\mathcal{O}_X,F)\\
 &\simeq R\Gamma\bigl(\mathrm{R\mathcal{H}om}_{\mathcal{D}}(\mathcal{O}_X,F)\bigr)
 =R\Gamma\bigl(\mathrm{DR}(F)\bigr).
\end{align*}
另一方面,下标 $!$ 表示真支撑上同调(这里是晶体上同调),即
\[
R\Gamma_{!}(F)\overset{\mathrm{def}}{=}
R\Gamma_{!}\mathrm{R\mathcal{H}om}_{\mathcal{D}}(\mathcal{O}_X,F).
\]}

\origpage{956}
\begin{quote}
$R\Gamma_{\mathrm{cris}}(F)_{!}$, \quad $R\Gamma_{\mathrm{cris}}(F')$
\end{quote}

作为拓扑向量空间的复形,由一个自然配对“处于对偶”,换句话说,人们有一个配对

\begin{quote}
$H^i_{\mathrm{cris}}(X,F)_{!} \times H^{-i}_{\mathrm{cris}}(X,F') \longrightarrow \mathbb{C}$
\end{quote}

(对一切整数$i$),而它\textbf{本身就是一条(拓扑向量空间的)对偶}。这条对偶定理“笼住”了离散系数情形下已知的(“绝对的”)对偶(Mebkhout称之为“Poincaré--Verdier对偶”),与凝聚系数情形下的对偶(Mebkhout称之为“Serre对偶”),合成为一种我愿称之为“Mebkhout对偶”的对偶——而他自己则称之为“Poincaré--Serre-Verdier对偶”\fnote{616}{(*)}{如我在别处(在“同伙”一注,n\textsuperscript{o}~63$'''$里)已经说到的:Mebkhout“不能不”向他的“恩人”Verdier脱帽致敬(自从后者把那条天赐的“好的参考文献”传给他之后)——凡有机会他都如此。然而两套对偶中任何一套(更遑论笼住它们的那一套)的本质观念,没有一件归功于Verdier。事实上,除了Poincaré对偶定理与Serre对偶定理的初始形态——它们当然曾是我的出发点——之外,一切本质观念都包含在我引入、并在凝聚与离散两个语境中长期发展起来的那套六种变换与双对偶形式主义里,而且是在孤独中发展的。

正是想着这一点,我去年在“受害者——或两种沉默”(n\textsuperscript{o}~78$'$)一注中写道:Mebkhout的那些“保护者”们“肯让他亲手抬着那口装着我的遗骸的棺材的一角”。当时我若还补一句才公道,那就是:Zoghman有那份勇气——尽管他清楚地感觉到俊髦世界里刮的是什么风——在每一篇文章里都明确重申他取法于我的思想,而不像大家那样:一边在(文字中)沉默地掠夺亡者,一边(在言谈中)摆出居高临下的姿态。

至于人们终于给那门我发展了多年、在完全孤独中发展的凝聚对偶理论起的名字“Serre对偶”,它格外有滋味(而Serre——他可没要求这么多——会比任何人都更欣赏!),因为Serre对我的对偶工作曾表现出全然的漠视,从而剥夺了我唯一可以指望的对话者来陪伴我的思索!我相信还可以说:那份漠视一直完好地保持到今天,连导出范畴的观念(以及其他无用细节……)也包括在内。}
% ---------------- p.746 ----------------
在我看来,这些是一个维度广阔的对偶纲领的最初几步,其中尤其包括(除别的之外(171(xi)))为特征零体上有限型概形上的 De Rham-Mebkhout 系数发展一套(双对偶与)六种运算的形式主义(且期待着更好)。鉴于Mebkhout不得不\origpage{957}工作的孤立条件与那种漠不关心的氛围,他当时不可能\textbf{为自己}发展一套完整的形式主义,像我在他所取法的两个语境中发展出的那种(171(xii))。在1972到1980这八年间他所提炼并证明的主要结果中(171(x)),在我看来、以我六十年代那个纲领的眼光看,最重要的当然是把那一类好的晶体系数——所谓“De Rham-Mebkhout”的——揭示出来的那个。巧的是,也正是这个结果,从1980年10月起,见识了最辉煌、甚至惊人的运气——而它却像从前的$\ell$-adic上同调、或特征$p$的晶体热门点心一样,只被当作一件\textbf{工具}侵占,从赋予它全部意义与力量的那幅图景上被撕走。

比Mebkhout的其他结果更甚,也一如我在发展双对偶与六种运算形式主义时那样,导出范畴的语言在这里是本质的——为的是揭示离散系数与凝聚系数之间那条简单而深刻的关系\fnote{617}{(*)}{(5月7日)确切地说,对一个完整 D-模(化归到零次的复形),好上帝函子一般结合出一个可构$\mathbb{C}$-向量空间层复形,它有不止一个非零的上同调层,反之亦然。最简单、最触目的例子是:在$X$上取一个除子$Y$,由此得包含$i:=X\setminus Y\hookrightarrow X$,再取 $i_{*}(\mathcal{O}_U)$ 中由沿 $Y$ 亚纯的函数构成的子层。Mebkhout早在1976年就得到的一个深刻结果(后来被吸入好上帝定理)是:那是一个完整正则的 D-模(在Mebkhout之前,谁也不曾想到把这个层看作一个 D-模,更不曾疑心它至少是凝聚的……)。它经好上帝函子的变换是$Ri_*(C_U)$,它至少在第0次与第1次有非零的上同调层。

那是Mebkhout哲学中为Deligne的进路所缺失的一面——后者得到的是一本词典,在可构 $\mathbb{C}$-向量空间层与$\operatorname{Coh}(\mathcal{O}_X)$(即 $\mathcal{O}_X$ 上凝聚模的范畴)的某些带分层的 pro-对象之间,而不必过渡到复形与导出范畴。(他终究还是留心让那些范畴出了场——那时我还在附近,而谁也想不到有一天那些范畴会被埋葬……)。那(至少乍看)是Deligne进路的一个优势,更贴近离散系数的直接几何直观——但那也是一个迹象,无疑说明他的进路较不深刻。我倾向于相信它还有自己的角色要扮演,但要与Mebkhout的视点“搭档”,而后者(我揣测)在某种意义上是对偶的。

(5月24日)关于这方面的细节,见子注“五张照片(晶体与 D-模)”(n\textsuperscript{o}~171(ix)),(c)部分,尤其是第1009页及以下。},它由好上帝定理所描述(别名:那位从不被点名的Mebkhout……)。于是,在那个\origpage{958}平展同调工具的创造近二十年之后(而今天所有人都理所当然地使用它,同时轻蔑那催生它的图景……),多亏这位隐身的身后弟子那(已成为“热门点心”的)结果,导出范畴的语言才突然获得平反(仿佛它从不曾被埋葬……),在聚光灯下、在人群的欢呼中——那人群是来向昨天的掘墓人欢呼的,而他们(谦逊地)扮演起新父亲的角色。但我又抢跑了……
\subsubsection*{c.~……与意外之喜}
\tocsubsub{c.~……与意外之喜}
\notehead{171(iii)}{}

Verdier多少充当着Mebkhout的“论文导师”;而Mebkhout此前七年的工作,完全是在孤独中进行的。他从来不曾对这个年轻人的工作产生丝毫兴趣——这年轻人显然既倔强又愚钝,不过是一个落伍的Grothendieck主义者,他以居高临下的姿态对待他。自1975年初次见面以来的四年里,他总共只给这个不知从哪里冒出来的无名小卒三次“会面”。我的其他上同调学生,没有一个
% ---------------- p.747 ----------------
肯屈尊对那位无名小卒的工作发生兴趣。他的工作对他们自己的研究的意义,他们全然看不出——而它却那么刺眼,连我这样一个十五年前就已从这一切“脱钩”的老朽都看得出来。他们太封闭在自己的掘墓-狂欢里、封闭在一种闷闷不乐的摇手柄日常里,无从领会一件新事物——它不带名片、不加修饰,只凭那些最简单、太显然不过的事物之力而来。他们早已埋葬自己的创造官能,只限于做名牌时兴货的消费者。然而后来,他们大大地向那位闯入者报复——他竟看到了他们以及所有人未曾看出的东西(而他们其实拥有一切,同他一样、甚至超出,足以看见、足以动手……)。但在这里我又抢跑了……

答辩在1979年2月15日举行,在一片普遍的冷淡中。Mebkhout把他的论文寄给所有他以为(对也罢、错也罢)对解析簇或代数簇的上同调感兴趣的数学家——当然从我的所有学生开始。在收到论文的全体人中,\textbf{没有一个}肯告知收到了寄件,或寄一个谢字。诚然,Mebkhout的论文比他某些文章更(在我看来)带着周围那逆境条件的印记——它在我看来卷帙浩繁、至少可以说不好进入,而那些不在圈内的人有理由不曾立刻咬钩。相\origpage{959}\dmark{}反,我发现Mebkhout对我口头解释他的哲学时说得完全清楚、且立刻令人信服;而没有任何理由说:他对Verdier(1976)、Berthelot(1978)、Illusie(1978)与Deligne(1979)所能作的解释,会比对我作的那些逊色。

那是在1979年6月的Bourbaki讨论班上,Deligne从Mebkhout口中得知那篇未被读过的论文中的“\textbf{Riemann--Hilbert对应}”。(那是Mebkhout给那个范畴等价(或那本“词典”)起的名字,前文刚谈到过。)显然,Verdier在那四年间从不曾想起就他的那位隐身学生的工作跟他多说一句——那工作的兴趣显然直到1981年反常层会议前后都全然逃过他(而在那次会议上,Deligne不得不负责给他解释那是怎么回事……)。而在Deligne那边,那只能立刻“对上暗号”——那正是他十年前自己撂下不管的那个问题的圆满而简练的解答!

在这种情势下,看似理所当然的反射(乃至我此刻都难以想象怎么能换一种做法……),是立刻向这位年轻的陌生人道贺——他终于找到了一个说实在深刻的、人们为之辛辛苦苦干了一整年、最后付诸流水的问题的谜底。世风已经大变……Deligne(当然依旧和和气气)只说了一句含糊的恭维(而那已足以叫天真的Zoghman心里发热——他确实不曾被惯坏,而且远未料到等待他的是什么):是的,他收到了他的论文,甚至读了它的引言,而他觉得那是些“美的数学”。对Zoghman那是黄道吉日!那当然是头一回(也同样是最后一回……)他有权得到一位大先生的恭维——那位人人都认识、人人都引用的先生……\fnote{618}{(*)}{(5月14日)况且那也是Mebkhout与Deligne之间唯一一次谈话的荣幸。

(6月7日)至于另一句恭维,在头一年(1978年6月)、且这次是出自Illusie之口,见“掠夺的空白授权——或高尚行当”一注(n\textsuperscript{o}~171$_{4}$),尤其是第1091页。}我说不出此刻以及随后一年里,Deligne脑子里关于这条他刚从一位无名者口中得知的卓越定理在想些什么。我揣度他必定会跟周围\origpage{960}\dmark{}的人谈起\fnote{619}{(*)}{(5月14日)细想之下,并且凭我从别处对Deligne的了解,我怀疑在他带着一个明确的主意和一个订好的计划去做之前,是否真的曾“跟周围的人谈起”。关于Deligne所扮演的十分特殊的角色、以及他让Beilinson与Bernstein这两个稻草人父亲所扮演的角色,见“父亲们的圆舞曲”一注(n\textsuperscript{o}~176$_{4}$)。(另见“冤大头买卖——或木偶戏”一注,n\textsuperscript{o}~172$^2$(e)。)}

% ---------------- p.748 ----------------
次年的10月,他把它私下通报给苏联数学家Beilinson与Bernstein\fnote{620}{(**)}{(5月14日)这是从Deligne给Mebkhout的一封信(1980年10月10日收到)中透露出来的。关于 Kazhdan-Lusztig 这一段的细节,见子注“黑手党”(n\textsuperscript{o}~171$_{2}$),(d)部分“总彩排”。},准是猜想到他们会用得上。同一年,事实上正是这个“对应”(当人们肯提名道姓时仍叫作“Riemann--Hilbert对应”,而Mebkhout的名字从不曾被说出)成了那场著名猜想\fnote{621}{(***)}{同一条猜想同时(相差几天)由 Brylinski-Kashiwara 独立证明——但其雷同程度却令人吃惊——带着同一主要配料、同一套掉包手法、以及对这个新事实的关键角色和它的作者的名字的同样抹杀。细节见已引的子注“黑手党”(n\textsuperscript{o}~171$_{2}$)(c)与(d)两部分。}——我几乎只知道它的名字,“Kazhdan-Lusztig 猜想”——的证明中那本质的配料、那此前一直缺失的\textbf{新事实}。那同时是代数簇上同调一场突然而壮观的新生的发令枪——它终于走出长达十余年的停滞(除了Deligne关于Weil诸猜想的工作)。这场意外的新生次年即见诸实行,那是1981年6月在主题“奇异空间上的分析与拓扑”下的吕米尼“事件”\fnote{622}{(****)}{会议的《论文集》出版于 Astérisque n\textsuperscript{o}~100 (1982)。而且那些《论文集》直到1983年12月才印出、1984年1月才问世,比卷上标注的日期晚了近两年。}。



\subsubsection*{d.~光荣之日}

\tocsubsub{d.~光荣之日}

\notehead{171(iv)}{}

关于这场“值得纪念的会议”,我请读者参看“不义——或归来的意义”(n\textsuperscript{o}~75)一注,以及随后那几注——它们写于发现之后还带着惊愕(这词不算过分)的热头上。这些注构成我所命名的葬礼第七行列:“会议——或Mebkhout层与反常性”。

我在这里只需提醒一句:在那场由 Bernard\textbf{Teissier}与 Jean-Louis\textbf{Verdier}签名的会议《论文集》引言中,那条著名的“Riemann--Hilbert对应”被作为会议的“Deus\origpage{961}\dmark{}ex machina”来介绍。在主要文章中也是如此——那篇文章(与所引引言一起)构成《论文集》的第一册,署名 A.A.~Beilinson、J.~Bernstein 与 P.~Deligne(而事实上由最后那位撰写、并在会议上报告,另两位共同作者缺席)。况且,头两位署名作者在上一年(1980年11月)就已由Mebkhout亲自(且独立于Deligne)告知了他那定理的来龙去脉——Mebkhout甚至为此特意去了一趟莫斯科\fnote{623}{(*)}{关于这段有教益的插曲,见已引的子注“黑手党”(n\textsuperscript{o}~171$_{2}$),(d)部分“总彩排(封神之前)”。}。Teissier也第一手地、而且早就知情——Verdier更不必说,他主持了Mebkhout的论文答辩……最后,我补充一句:临到末了才“急就章”地决定请Mebkhout作一个关于 D-模理论的报告(那门理论在场的人中除了他没有别人懂),Mebkhout就这样有机会把他谦逊地以 Riemann 与 Hilbert 之名命名的那个定理告知全场\fnote{624}{(**)}{(5月14日)关于这场古怪会议的与会者——那是一场盛大的“Grothendieck数学节”,却对亡故的祖先本人、以及对那位“有那份天赋让如此俊美的世界聚在一处”的隐身学生绝对沉默。作为“1970年前”的学生与会者,只有Deligne与Verdier,但已经足够牢牢占据舞台前排。蹊跷的是,Berthelot与Illusie(他们的工作可以说特别带着在此大张旗鼓地发掘出来的Mebkhout视点之缺席的印记)不在其列。相反,Contou-Carrère(“之后”的学生)却迷路闯了进去,庆幸自己被请来讲述他那套 Schubert 闭链的消解法。

我记得他回来后兴高采烈,完全认同于所有那些杰出而出名的人——他感到同他们平起平坐,而他们是来听他讲的,显然是感兴趣的,一点不假!他带着愧疚的神气向我谈起Mebkhout——后者曾对他打开心扉,却带着苦涩,而他也说不清为什么;可对孔图来说,无论如何日子显然是美的!

那是1981年6月。四个月后(为回应他向佩皮尼昂一个职位的单独应聘),那是结结实实的一记耳光,他当作一种屈辱和冒犯硬生生地受下。(关于这段插曲,见“棺材3——或那些稍稍太相对的雅可比簇”一注,n\textsuperscript{o}~95,尤其是第404--406页。写那注时我还不曾把它与Contou-Carrère参与那场辉煌会议的插曲联系起来——他大概被稍稍甩在了一边。)}
% ---------------- p.749 ----------------
而不对这件事的父权留下丝毫含糊(可想而知)。那件事有那份意想不到的天赋(无论对他还是对所有人)——让如此俊美的世界聚在一处。\origpage{962}\dmark{}而且,读者在会议的《论文集》中寻找Mebkhout那场报告的踪迹,也是白费。Verdier事后和颜悦色地向他解释:只有呈献\textbf{新}结果的文章才会收进《论文集》,而他博士论文里的结果已经两年以上了。读者在那部《论文集》中寻找关于这条著名定理之来历的任何文献引证、或哪怕最含糊一点的提示——同样白费:那条定理当然既不归功于 Riemann 也不归功于 Hilbert。读者也休想在其中找到Zoghman~Mebkhout这个名字。这个名字在第一册中不出现,无论在正文还是在文献目录中。在第二册里,它在文献目录中出现两次,是 Brylinski 与 Malgrange 的“暂停!”式引文(总不能说没人引过他!)——而那些引文与好上帝定理毫不相干——别名Riemann--Hilbert——别名Deligne(而尤其不别名Mebkhout)\fnote{625}{(*)}{(5月14日与26日)除了已提及的与会者,我具名得知的有 Brylinski、Malgrange 与 Laumon。三位都完全了解Mebkhout的工作,而他曾有机会对每一位详加说明,且远不止限于他在会议所作的那场报告。但这并不妨碍 Brylinski 与 Malgrange 在他们收进《论文集》的文章中——那文章本质地使用Mebkhout的思想与好上帝定理——既抹杀这些新思想与新工具之出现所扮演的关键角色,也抹杀它们的作者的名字。

至于 Laumon,他后来在一篇与 Katz 合写的文章中补了过来。而这位 N.~Katz 正是1973年以“SGA~7 行动”崭露头角的那一位——那件事在“一场攀登的若干插曲”一注(n\textsuperscript{o}~169$_3$,插曲2)中谈过。况且,他早在1979年就由Mebkhout亲自告知了他的结果(关于这一点,见“掠夺的空白授权”一注,n\textsuperscript{o}~171$_{4}$)。那是指《傅里叶变换与指数和的估计》一文(它同时也是 Laumon 的博士论文),那文章以预印本形式流通已有两年(我甚至承 Laumon 厚意得到一份)。这些作者们为$\ell$-adic 系数发展一种傅里叶变换,其模型是 Malgrange 1982年在 D-模情形引入的那一种(顺着那位无名者的工作的尾迹,且如所共料不提他的名字)。Mebkhout的工作构成 Malgrange 所发展的理论与 Laumon-Katz 二者的启发基础,一如它们曾是已引的 Beilinson-Bernstein-Deligne 那篇文章(关于他们误称的“反常层”)的启发基础。话虽如此,Laumon 与 Katz 也跟着大伙的调子走(文中与文献目录里都不提那位当值无名之辈——当然更不消提那位祖先……),循着 Deligne、Verdier、Berthelot、Illusie、Teissier、Malgrange、Brylinski、Kashiwara、Beilinson、Bernstein 定下的调门——我为这按字母排的次序道歉,总之这已有十二人直接而积极地卷入反常层会议那场辉煌的神秘化-骗局——还不算 Hotta 也在太平洋彼岸助了一臂之力,就成了十三!

Malgrange 在那篇文章中也并未被多引——显然有些作者小圈子互相狂引,却避开引隔壁的,哪怕他们一个比一个使劲地从人家那儿汲取。总之,一旦事关那位祖先或那位无名者,他们就全体一致。那往往确实是出色的数学——但我这老派人对这心态并非无动于衷,而它败掉了我阅读的胃口,乃至最终连做数学的胃口。至少不做他们那种数学。气味太难忍受……

我也曾对 J.L.~Verdier的文章《特化与温和单调层》投去一瞥,那文章发表在同一批《论文集》中。当然毫不意外,我在那里也看到了“Riemann--Hilbert对应”,却(在正文和文献目录里)丝毫不提那位他曾主持其论文答辩的无名者。他想必是不得已忘记了……那里还谈到一条平展 Riemann-Roch 定理(这名字我仿佛在哪儿听过……)——而我在 Laumon-Katz 的文章里也见过这个。既然双方都不曾提起某位亡者,我便对自己说:那条“定理”想必一定该归 Riemann 与 Roch 两位先生,正如那个在 SGA~5 的“技术性离题”与“废话”之中的特例一样(还不算那篇猜想的报告,被那位深谋远虑、机灵的“编辑”Illusie仿佛天意般地掏空了……)。

Mebkhout况且早在1977年就预感到他的哲学与傅里叶变换之间的联系——那时他正严格地独自一人对一门联结 D-模与离散系数的对偶瑜伽感兴趣(一如我自己当年对凝聚对偶、而后平展对偶形式主义的孤独)。那个“傅里叶变换”的直觉始终停留在含糊的状态——当时的语境对他继续走这条路并不比我1950年前后去拓展我的凝聚对偶理论、扩展到一门涵盖微分算子复形的理论(见第946页脚注(**))更容易。在Mebkhout为巴黎第七大学“奇异性讨论班”(1977--79)所作的报告“Poincaré对偶”之引言的第2页,有一条关于傅里叶变换的暗示。}

% ---------------- p.750 ----------------
\origpage{963}\dmark{}且回到那场有血有肉的会议,该相信,那些聚集在那里的杰出数学家,即便屈尊来听那位当值无名之辈给他们作的报告,也没有一个发觉:那位无名者作为自己的创造介绍给他们的“Riemann--Hilbert对应”,正是他们中最出色的一位早已辉煌地引入的那一个,作为他那辉煌报告的启发性的拱顶石——而那报告(据组织者Teissier与Verdier\fnote{626}{(*)}{那是从所引的会议《引言》中清楚透露出来的“隐含之见”,由 Teissier 与 Verdier 署名。}自己的估计)构成那场关于所谓(人们还真想知道为什么)“反常层”的辉煌会议的“钉子”。总之,该相信他们中没有一个曾惊讶:那位无名者的名字不曾在那场报告中被提起——那报告飞得那样高,固然无须为一个区区小辈费神;而两年半以后,随着《论文集》的问世(1984年初),那位无名者的名字也同样不在其中——既不在(已提及的)引言里,也不在Deligne及其合作者的相关文章里。而那篇文章\origpage{964}\dmark{}对这份对应的真正父权更不曾留下多少疑问,而主要作者兼报告者-变戏法者\fnote{627}{(*)}{关于我的朋友皮埃尔围绕那条从未被提名的定理之父权玩的花招,见去年那注“变戏法者”(n\textsuperscript{o}~75$'$)。},以他惯常的谦逊,也不曾提名道姓,连他那两位显赫先驱的名字也不曾。若说居然有人惊讶过,那他们直到今天也不曾让人知道——至少不曾让我知道,尤其不曾让那位提供笑料的主角知道:那位身后的、照理说完全不为人知的学生——一如既往——Zoghman~Mebkhout\fnote{628}{(**)}{(5月19日)关于我朋友Zoghman的那些坎坷遭遇——他天真地误入一个衣冠楚楚、满面和气的“硬汉”圈子——的详情,见以下几条子注:“一个愿景的萌发——或闯入者”“黑手党”“根”“掠夺的空白授权”(n\textsuperscript{o}~171$_{1}$至171$_{4}$)。}。

\subsubsection*{a1.~无用的细节}

\tocsubsub{a1.~无用的细节}

\notehead{171(v)}{}
\fnote{629}{(***)}{本注(分(a)(b)(c)三部分)源出“祖先”一注(n\textsuperscript{o}~171(i))的两条脚注——见第~944~页的脚注(**)与第~945~页的脚注(*)。}

\subsubsection*{(a) 上千页的一捆……}

\tocsubsub{(a) 上千页的一捆……}

\noindent (5月4日)连Serre也不曾免于这条规矩:他早已(如André~Weil一样)养成一种可憎的癖好,动辄宣判那些不合他心意的数学是“蠢话”。他与Weil那样的规格(人们或以为)本该使他们超越这类幼稚行径。而在此时此地(Deligne那“最后的二十页”姑且置而不论),恰恰是两三千页“Grothendieck式的蠢话”,最终证明了Weil猜想(还顺带证明了Weil与Serre做梦也不曾想到的不少别的东西)。这并没有引得Serre多一份谦逊:就在他陈述Deligne对这些猜想之最后一步的证明的那篇文字里(1974年2月的Bourbaki讨论班,第~446~讲),他恰恰借这个机会(措辞彬彬有礼,这自不待言),就SGA~4那“1583页”必定塞满的“无用细节”大加讥诮。在这种廉价的讥诮里,我辨认出的并非恶意,也并非不老实,而分明是一种无意识与轻率。他倒不辞劳苦,数出三大卷的页数(他自然留意不去读它们,其中的实质他也无从领会),又做了一道加法——好“优雅地”嘲弄一番。

% ---------------- p.751 ----------------
\origpage{965}\dmark{}但一切相互勾连:既有我当年对那些聪颖学生的迁就,也有Serre的“优雅”(而那是在葬礼已经发展了四年之后……)\fnote{630}{(*)}{(5月27日)对谈及Serre处加以接续的一段反思,见本注的(c)部分。}以及随后的一切。仅仅三年之后,人们几乎会在我的非学生Deligne笔下重新见到Serre那些说法本身、或它们的言外之意,还多一点恶意与无礼,连同那些被“剔除”的“无用细节”、那种“混乱状态”与那团“废话硬壳”——那正是这位Deligne学得手艺、觅得主要灵感的地方,而由他笔下的一篇寡淡摘要慈悲地“让人忘掉”。就这样,从迁就到随便、再到无礼,数学界在仅仅十年间,竟走到一种风习状态,连“体面”这简单的感情也仿佛消失了。

创造“那条猜想”所缺的新工具的,既不是Weil、也不是Serre,更不是Deligne;而是那个他们乐于讥讽的人——出于蓄意的无视或算计的恶意,效果并无太大差别。但我自己以万般的小心写了又写、让人写了又写,不倦地度过一月又一月、一年又一年,写出一部文本,它以其应有的全部宏阔,呈献那门语言与若干基础工具——为的是一幅崭新的、丰产的统一图景;我深知,而且明明白白地知道:在那被Serre、我的学生与一致的风尚置于脑后的1583页中,\textbf{没有一页}不曾被匠人称量又称量,没有一页不在其位上、不履行它的职能,而迄今写下的任何一页都不能取代它。这些页不是风尚的产物,也不是虚荣的产物——那种喜欢凌驾于他人之上的虚荣。它们是我的爱情的果实,是那些漫长而晦暗的劳动——它们准备着一场分娩。

至于我作品的这一部分,一如我那些已然进入共同财富的重大数学贡献,到今天为止,\textbf{没有谁}曾能重做我所做的事(仗着“蠢话”“无用细节”与“废话”),除非靠抄我(除了些不关紧要的变体)\fnote{631}{(**)}{(6月7日)我近来读到了 Fulton 的好书《Intersection Theory》(“Ergebnisse”丛书,Springer Verlag,1984年),发觉对 Riemann-Roch-Grothendieck 定理,应当算作例外。}。有些人抄(照抄、或在相近乃至崭新的语境里抄)并且说明(这种事如今已开始变得少之又少……);另一些人则扮演新父亲,并对他们无耻掠夺的作品、\origpage{966}\dmark{}对教会他们手艺的匠人摆出鄙夷的恩赐姿态。这种无礼之所以能滋生并铺展,只因为它找到了一份随时准备接待它的共识,而这首先来自那些(往往以其超常的身量)定调子的人。

% ---------------- p.752 ----------------
\subsubsection*{(b) 无所事事的机器……}

\tocsubsub{(b) 无所事事的机器……}

六种运算瑜伽是这幅在 SGA~4 与 SGA~5 两个讨论班中发展起来的“宏阔统一图景”的组成部分。我甚至要说,那场口头讨论班 SGA~5 的中心主题、或不如说它的“筋”与灵魂,就是这套瑜伽。因此Illusie才着意让它从屠杀版中消失(那部版子按其“照料”本该变成一册“技术性离题……”)。

在“祖先”一注(n\textsuperscript{o}~171(i),第945页)里我(未作进一步说明地)写道:六种运算这个富有力量的图景“给出过它自身力量的有力见证”。对我来说,这个力量最触目的具体标志——如果可以说的话——就在我们对平展上同调的掌握之中。为了在1963年达到这一掌握,那来自凝聚对偶的“六种运算”图景是我恒定的导线。况且我自认是世上唯一有资格就什么在这一工具的发展中具有决定性发言的人。

这里不言自明的是:在发现的进程里,所谓“启发性”的要素几乎总是决定性的。如果我谈一种观点或一幅图景的“力量”(那是与一条定理自身全然不同的另一种东西),那不能用严格技术的措辞衡量它。首要的是它那“暗示性的力量”——在发现的旅途中做谨慎而可靠的向导,在敏感的时刻向我们口授“那条”该引入的正确观念、“那个”该提炼并证明的正确陈述、“那门”仍待发展的理论。正是忘却了这样一幅向导图景(在把它埋葬之后),才使得代数簇上同调理论从六十年代那股强劲的冲劲,在我离开之后的年头里落得个混乱与萎靡。除了那桩“威望大问题”(即 Frobenius 特征值绝对值的问题)之外,一切本质的问题都被冥顽地回避了……

作为这幅图景(或在此情形,不如说这套形式主义)力量的另一标志,我举离散语境与凝聚语境中的Lefschetz--Verdier不动点公式。在这里,“六种运算”形式主义的角色既是\textbf{启发性}的(就这个意义而言:该公式被这\dmark{}形式主义\origpage{967}不可抗拒地暗示出来),也是\textbf{技术性}的(就这个意义而言:该形式主义也给出证明这条公式所必需且充分的手段)。诚然,鉴于葬礼,直到“相交上同调”与被命名为“反常”的层那场大冲刺之前(那时一部分形式主义被发掘出来,却不提匠人……),我所发展的那套上同调形式主义中只有极小一部分被使用过。但我自己清楚:凭着Weil诸猜想与 topos 直觉的无所不在,六种运算的图景是我1955--1970整个年间在上同调反思中的主要灵感来源\fnote{632}{(*)}{(5月15日)不言自明,这幅图景本身是在这一时期中渐渐取得形态的,出发点是我1955年那篇《论同调代数的若干问题》(载 Tohoku Math. Journal)中所含的最初萌芽。它在1963年达到完全成熟,随着平展上同调的突然启动。而这次启动(巧得出奇)几乎就发生在我引入那个“缺失的函子”$Rf_!$(真支正像)的那些日子里。但六种运算作为“富有力量的图景”与无所不在的导线,其角色(我相信)是在 SGA~5 讨论班里才变得充分自觉的。早在1966年,随着晶体上同调的启动,我就已清楚:第一个目标(超出那项“磨合”的有限纲领——它将在Berthelot的论文工作中完成)是达到一套针对“好的”晶体系数的六种运算(与双对偶)形式主义。只可惜须得由一个(被宣告亡故的)老朽从那为他预备的棺材里走出来,才使(在将近二十年之后,并借助一位当值无名之辈的想法,而他也被一同埋葬)这些“好系数”好歹终于\textbf{被定义}出来!至于它们的描述(尤其对$\mathbb{Z}$上有限型的概形),见《数学反思》第3卷(连同《收获与播种》第五亦即最后一部分)。}。这就是说,这幅图景的“力量”对我是显然的事,或不如说是我在数学家生涯的十五年间几乎日日经验过的现实。而且,这种经验在这些最近几周里又引人注目地得到再确认——一旦我重新接触那些“被遗弃的工地”:晶体系数与 De Rham 系数的工地、以及 Motif 的工地\fnote{633}{(**)}{(5月15日)关于“被弃置的工地”(或“荒废的”工地)这一意象,见葬礼仪式的第6幕(注~176$'$、177、178),尤其是三条所引注中的最后一条。在撰写《收获与播种》之余,我不过是在这里那里花上几个小时,处理 De Rham 晶体系数的问题与 Motif 的问题,便足以看到一个针对前者的令人信服的定义、以及至少一个构造后者的原理,在$\mathbb{Z}$上有限型概形这一关键语境中出现。(与前一脚注的评注相比较。)}

\origpage{968}\dmark{}这份关于某种富有力量的图景之力量的全部“主观”经验,也有一个“客观”的意义,那是不可一挥衣袖就抹掉的。只要愿意记住这样一件事,这个意义就显现出来:(除了少数例外)关于“抽象”代数簇与概形的上同调的那些主要观念与概念(今天所有人都使用它们,仿佛它们可以追溯到亚当与夏娃\fnote{634}{(*)}{关于这种数学成品“使用者”(或“消费者”)的心态——他忘记了(若他曾知道过……)创造是什么事——也关于亚当与夏娃以及上帝,我请读者参看去年那些注“一种不义与无力的感觉……”(n\textsuperscript{o}~44$''$)与“当值无名之辈与好上帝定理”(n\textsuperscript{o}~48$'$)。亦见“教学的失败(2)——或创造与浮夸”(注 n\textsuperscript{o}~44$'$)。}),正是由我——而不是别人——在1955--1970这同一时期里提炼出来的。(不言自明,我在此把我出发的 FAC 与Weil诸猜想都排除在外。)

这个伟大的时代,就数学而言(且据我迄今所能看到的),通向一种阴郁、呆滞与贫乏。其深层原因绝不位于技术层面。这种贫乏的标志之一,是一幅为激励并哺养宏大设计而生的有力图景,竟被埋葬、或交由嘲笑——而动手的正是它的保管人与第一批受益者。而另一个标志是:Deligne也罢、Verdier也罢、Berthelot也罢、Illusie也罢,尽管地位与威望带来一切便利、才华横溢、经验老到,没有一个做出一件在 De Rham 系数上非做不可的作品——那本是他们自己的研究(与被否决的图景……)的正统后裔;他们甚至没有本领在面对这件作品时,认出它的发明性与丰饶。而正是本着这同一种精神(因为一切又相互勾连……),一旦那件出自新作品之一的工具的意义终于被承认,他们就赶忙把它据为己有,甚至并不理解它,并把那位锻造了它的无名匠人——连同祖先——一道埋葬……
% ---------------- p.753 ----------------
\subsubsection*{(c) 像什么也不像的东西……——或枯萎}

\tocsubsub{(c) 像什么也不像的东西……——或枯萎}


(5月27日)\fnote{635}{(**)}{本注“无用的细节”的第三部分源出第一部分的一条脚注。见第965页脚注(*)中的转引。}我关于Serre所说的话是自然而然写出来的,并且出自对有关他的事物的一种知觉——那知觉想必是在最近几个星期或几个月里在我心中形成的。\dmark{}\origpage{969}然而在写这几行时,曾有一份残存的犹疑或困惑、或保留,针对我刚写下的东西。我那话里总而言之让人听出:Serre在这一场合上竟缺少“优雅”!

实情是:我认识Serre快三十年了,他一直就是我心目中“\textbf{优雅}”的化身。想必不是只有我这样看他。那是一种优雅,既在他的工作与作品中,也在他对他人的关系中;而那绝不是纯形式的。它包含工作中一丝不苟的诚实,以及对他人同等的诚实要求。不止一次,我注意到他在某个较不谨慎的同事面前判断的锐利——那位同事一心要把一个恼人的困难弄浑(为的是不必承认他不知该如何克服它),或把他自己的某个错误含混过去……因此这种优雅也包含一种\textbf{严谨},对他自己与对他人同样。

正是这一切,对我而言与Serre这个人不可分,它们必定介入了我刚提到的那份“保留”——在一份对事物的知觉不期然地抢在熟悉的知觉之前的时候。我无意为了其中一个的“好处”而排斥另一个。两者都有东西要教我,都是一个复杂现实的不同侧面,而那现实当然绝非静止。该由我把它们相互定位,以达致对一个与我有一段过去、有同情与尊重之情的人的细致入微的领会。

然而我刚谈到的这种“严谨”,并不曾延伸到Serre与数学和数学家的关系上的一切。我前面曾提到一种“不自觉”或“轻率”,我本可以同样称之为一种“\textbf{封闭}”。它与那种“谨慎与谦逊”的态度形成对照——我在大多数长辈身上遇到那种态度,他们同Serre一样在我起步时善待我,而有时(如他那样)还带热情。关于这一点我在下文(在“自由……”一注,n\textsuperscript{o}~171(vii)里)有更多表述,我在那里确认:那种态度曾是“熏陶过接纳我的那个环境的尊重气氛……”的一部分。

% ---------------- p.753（接上段） / p.754 ----------------
我在某些场合在Serre身上见到的“封闭”,并非始于昨日。我在五十年代下半期就感觉到它的最初征兆。我相信,它大大限制了他自六十年代起作品的深度\origpage{970}\dmark{}与广度。我感到这种“封闭”与以下各事之间有着联系:一些不同于他本人的数学进路;以及在他内心逐渐发展起来的一个蓄意之举——把他对数学事物与对数学本身的领会,封闭进一种(或如我真想写的,一副“眼罩”)纯技术的或技术主义的观点之中,从而对一切近似于一种\textbf{视野}的东西关闭门户;对某种超越了有形的、即时的、\textbf{可证明的}陈述(或陈述的集合)的东西关闭门户,或者(退一步说)超越了采取“硬邦邦”的猜想之形式的东西——那些猜想轮廓分明,总而言之是“封闭”的(只不过还剩下要把它们证明……)。回过头看,我觉得他终究是把他创造力的那一面推向了极致:那纯粹“阳”的一面、“超阳”的一面、“\textbf{雄性}”的一面。凭借他在同代数学家中、以及在随后两三代人身上所享有的非凡影响力,在我看来,Serre大大促成了我所见到的、在七十年代与八十年代肆虐的那种过度技术主义精神——那是如今唯一还被容忍的精神,而数学的其他任何进路都已成为众人嘲弄的对象。

借用 C.L.~Siegel的一个说法,我们今天目睹的是一种非同寻常的“扁平化”(Verflachung)\fnote{636}{(*)}{我从Serre刚刚收到的一封信中借来这个(德文)表达式。该表达式取自 C.L.~Siegel为 Hecke 著作集所写的序言。Serre在信的末尾引了 C-L.~Siegel的这一印象,随即补上一句:“那时这不公道,而在我看来,如今就更不公道了。”可是它在我心里“咔哒”作响,并且持续地起着作用。我那段关于Serre与我之间关系的简短反思,大概正出于此。况且我相信:Serre之所以引用Siegel,是因为从某种意义上说,这一出自我们时代一位大数学家之口的印象,必定早已在他心里起作用;那大概是数学的“玫瑰人生”中的一个不谐之音。这样的不谐之音当然不止一个,但这一个看来不那么容易排遣……“Flach”在德文中意为“平的”、“没有深度”;“Verflachung”指通向这样一种“平”的状态的过程,或指这样一个刚刚完成的过程的结局。在正文里,我着力跟随这个极富表现力、可惜无法原样翻译的词在我身上唤起的联想。当然,我完全不知道我感知此事的方式,与Siegel的感知是否有丝毫重合——Serre所引的Siegel原文,我不曾读过。},一种数学思想的“压平”、“收缩”:它被剥夺了一个维度——那个洞见的维度,那个梦境与神秘的维度,那个深度的维度——而在我看来,在此之前它与这个维度从未失去联系。我把这感受为一种\textbf{枯萎},一种思想的\textbf{硬化}:\origpage{971}它失去了活的柔韧、失去了滋养的品质——变成纯粹的\textbf{工具},\dmark{}僵直而冰冷,用来完美地执行那些“费力挣来的”任务、摆上公开拍卖的任务\fnote{637}{(*)}{“公开拍卖”这个意象,大概是从“招标”公告(sic)那里得来的暗示:“法国国家科学研究中心通讯”以及我作为这份可敬机构里一名初出茅庐的研究人员,定期收到的其他纸页上,满是这类公告。这一行话,连同许许多多别的征兆,表明这种发现工作的“压平”绝不限于我昔日熟知的那个圈子,也不限于数学这门科学。我至今还没见到纯数学的招标公告,但那不会太久——我很容易想象我当年的某些朋友或学生,一本正经地坐在包着软垫的门后,坐在某个名称令人望而生畏的委员会里,决定哪些“研究方向”应宣布为优先、哪些“进路策略”应予提倡、哪些“名列胜者”的课题组的“投标”宜予“采纳”进入“预选”,甚或授以头奖:由监护部颁发的官方津贴,经“主管委员会”出具肯定意见后每两年续发一次……},当目的感与方向感、当把这些任务本身视为一个浩大的整全之部分的感觉,都已被所有人遗忘。有一种深层的硬化,被一种狂热的肥大所掩盖。

思想的这种失衡,是另一种更为本质的失衡的征兆之一,也是更为深层的空虚与匮乏的征兆。思想的这种干枯在过去二十年里蔓延并扎根,与此同时,人与人之间交往中那些习见的细腻与敬重的形式也在遭受侵蚀——这不是偶然。同样不是偶然的是:这阵蔑视之风曾经刮起,而我终于感到了它的气息,它与一种或多或少普遍的腐败相伴而来;一年多以来,我一直在绕着这种腐败作巡视,未能穷尽。

% ---------------- p.755 ----------------
Serre直到今天,对这种四面八方环绕着他的腐败依然毫无知觉。我却记得他嗅觉灵敏。然而,光有灵敏的鼻子不够,还得用上它,听取它要告诉我们的东西,即使它向我们报出的气味足以让我们不安;甚至足以让我们忧虑,当它把我们自己牵连在内的时候。我很清楚,Serre同我一样,不会想到与狼同嚎,不会掠夺,不会搞阴谋,不会揭人短处——在“人人”都掠夺、都搞阴谋、都揭短的地方。他确实什么都不做——他只是捂住鼻子(这样一来少了一只手……),装作什么也没闻到。

而他绝非孤身一人——在我隐居中气味仍能传到的那个我们共同的世界里,那些曾是我朋友的人,一个也没有——一个也还没有跟我谈起过——\origpage{972}哪怕只是话里带一点——他\dmark{}闻到的、令他不适的气味。我那些同行中,无疑仍有许多人正直地从事着数学家这一值得敬重的行当。但在那些坐在前排的人中间,我\textbf{一个也不}认识——谁有这份质朴,去相信他自己健全官能(此处即嗅觉)的见证,而不是捂住鼻子,以免对自己说出:这里有什么东西发臭——也许该去看看吧……

但我回到Serre其人与我本人,回到我在他身上感到的那种“封闭”——它不知何时出现,并随着年月日益加重。我相信,他作品中最富多产性的那一部分,那最深刻地影响了其时数学的那一部分,位于早年,位于这种封闭出现之前,或者至少,位于它对Serre与数学、与数学家的关系取得决定性掌控之前。也正是那些年,那些五十年代,他与我的接触对我来说最为多产;也正是在那些年,Serre在我身上扮演了那种“引爆者”的角色,给我的作品某些最具决定性的推动。也是在那些年,我心中诞生并成长起一幅浩大的图景,它在那些年激励并滋养了我的作品,直到今天依然如此。我可以明明白白地、清醒地说:如果在“我”之外曾有人对这幅图景的萌发有所参与,那就是他,Serre,在那些年。而之所以可能如此,只是因为在这些多产而决定性的年头里,他心中有一种对数学事物本身(就其自身而言)的敞开,包括对那些尚未落入即刻之把握的事物,对那些起初似乎不愿被已经成形的语言之网眼所框定的事物——那些也许需要多年晦暗而耐心的劳作,甚至需要整整一生,才能凝聚为有形的实体,才能让一个\textbf{躯体}的四肢、形态与轮廓显现出来,活生生、强健有力,见证着在熟悉的已知者之语境中,一个\textbf{新的存在者}出其不意地出现。

我相信,在我认识Serre的头些年,直到五十年代末前后,他一直保持着一种感受力,去感受“创造”这种摸不着、微妙的事物,去感受那些为一次分娩做准备的、卑微的劳作。我相信,在某个时刻,他曾能感到一幅图景的萌发,以及赋予它形体的语言——如灵魂或精神,如躯体……那时有一种无需言说的温暖,有一种审慎而有效的随时可用性,在他能够协助一场并非他本人的、艰辛而紧张的劳作之处;而通过一份同情与一份期待,他毕竟参与了那场劳作。

\origpage{973}\dmark{}我无法说出,他在我们共同激情的那个层面上,这种鲜活究竟于何时、以何种方式钝化下去,让位于我方才试着描画的另一种东西。大约在六十年代初,甚至更早,他就不复感知森林,而只肯看见他合意的这一棵或那一棵树。其余的东西没有存在的理由。我想,看到我那样专注于不倦地开垦大片无形之地,耐心地在那里栽种这一切尚不像任何东西的事物,一副已经看见一片繁花森林的样子,他仅仅是恼火罢了\fnote{638}{(*)}{(6月17日)在我“巡访工地——或工具与视野”一注(n\textsuperscript{o}~178)里逐一巡视的六个“工地”之中,古时只有一个工地(“Motif”工地)曾使Serre多少发生兴趣——而且……当我最近在附言(P.S.)中不带评述地写信告诉他,自认为已掌握在 $\mathbb{Z}$ 上有限型概形之上构造 Motif 范畴之形式的原理时,他在回信中丝毫不提。看来这些“Grothendieck式的数学”再也无法使他或冷或热了……}。

% ---------------- p.756 ----------------
这并不曾阻止我继续开垦、继续栽种与补种、继续修剪、继续重新开垦与重新栽种——也不曾阻止我们一如既往地成为伙伴,依然花上几个小时又一个小时讨论数学(多半是在电话上)。当我有一个轮廓分明的问题,而且是一个并未被禁绝的问题,我通常尤其向他求教,指望他或许会有所烛照——事实上,他常常确实有所烛照。我继续从他那里学到许多东西,而他想必也从我这里学到那些能够令他感兴趣的东西。这比礼尚往来或互通有无更好——总有一份共同的激情把我们联结在一起,有火,有火花。

可他已经不再是我的灵感的源泉。这源泉如今唯有在我自身之中\fnote{639}{(**)}{(6月12日)关于这段对Serre与我之间关系之反思的续篇,见“家庭相册”一注(n\textsuperscript{o}~173)c.部(“众人之中的那一个——或认同”),6月11日所作,以及 d.、e. 部。}。

\subsubsection*{a2.~古怪的问题}

\tocsubsub{a2.~古怪的问题}

\notehead{171(vi)}{}

(5月5日)\fnote{640}{(***)}{本注源出“祖先”一注(n\textsuperscript{o}~171(i))的一条脚注——见第946页脚注(*)。}我在这里的记忆有些模糊,并在随后几个星期里重新清晰起来,那时我有机会就这些问题重新作些许接触。实际上,我心中有两个截然不同的问题,一个完全确切,另一个相当含糊。

第一个问题关乎这样一个需要:为那些尚待确切定义的“De Rham 系数”,廓清出一套完整的“六种变换”(variances)理论。我的晶体观念,无论在特征 $p>0$ 中\origpage{974}还是在零特征中,都提供一个极为确切开端——人们已经预先知道,\dmark{}什么将取代那些 $\ell$-adic“局部系统”(或“常值扭曲层”,即 Betti 情形,在超越语境中),而必须设法定义“带奇点的系数”,当然是在导出范畴的精神之中\fnote{641}{(*)}{同样清楚的是,当基域为 $\mathbb{C}$ 时,人们想要的是一类等价于如下范畴的东西:其上同调层为代数可构的 $\mathbb{C}$-向量空间层的复形之范畴。这一极为确切的提示暗示出:通过剥丝抽茧,关键的问题是把晶体局部系统(在一个不一定闭的子概形上)附着为一个周围环境概形上的晶体层。这基本上是 Deligne 在1969年所做的事,只是事实证明,人们找到的不是一个晶体层,而是一个\textbf{晶体 pro-层},这在当时便是一个重要的新观念(而且只要你肯一看,它就“显然”)。然而,系统地使用 pro-对象,将要求相当可观的基础性工作,Jouanolou 为其论文(关于 $\ell$-adic 系数)所做的工作给出了一个先声。那里本须再次捋起袖子……\par 从 Deligne 的观点与我的观点来看,Mebkhout 通过 D-模的新进路,从此就归结为把一个晶体 pro-层替换为一个晶体 ind-层(借助\textbf{通常的}凝聚对偶函子 $R\mathcal{H}\mathit{om}_X(-,\mathcal{O}_X)$),并\textbf{取归纳极限},以求得一个通常的晶体层,即(现在假设 $X$ 在零特征域上光滑)一个 $\mathcal{D}_X$-模。于是,Mebkhout 在1972至1976年间(从一个相反的“端点”出发,参见“三个路标”一注 n\textsuperscript{o}~171(x))所确立的那个出乎意料的“奇迹”,是这个 D-模是\textbf{凝聚的}(更确切地说,其同调层是凝聚的)。另一个同样出乎意料的奇迹是:人们能够以一些简单的、相对于Grothendieck的晶体观点而言性质全新的条件,来刻画这样得到的 D-模(或不如说 D-模的复形)——即那些“微观局部”的完整性条件,加上一个由 Mebkhout 引入并在此期间已为人所熟悉的“正则性”条件。\par (5月26日)关于 Mebkhout 的 De Rham 系数与 Deligne 的 De Rham 系数之间的对偶关系之细节,见“五张照片(晶体与 D-模)”一注(n\textsuperscript{o}~171(ix))c. 部。关于必须以“凝聚 pro-模形式的晶体”取代 Deligne 的“pro-凝聚模”观点,以及关于(尚未证明的)以不加修饰的晶体层(通过取射影极限)取代笨重的 pro-对象(晶体式或分层式)观点之可能性,见同一注的 c.、d. 部。}
% ---------------- p.757（本段承上） ----------------
,于是所缺的,便是晶体复形的一条恰当的“有限性”条件。在零特征中,给出答案的是“D-凝聚性”(不论我还是我的学生中谁都不曾想过——多么简单自然的主意!),再加上更微妙的完整性(condition d'holonomie)与正则性条件——正如“好上帝”即Mebkhout的哲学(在晶体瑜伽启动十二年之后)所教会我们的。我怀着好奇等待着我某个前学生终于动起来(不提那位当值无名之辈,也不提那位祖先——这是说好的……),以廓清特征 $p>0$ 中\origpage{975}对应的条件,或不如说,廓清在刚性解析语境的特征\dmark{}零中对应的条件。迟做总比不做好……\fnote{642}{(*)}{(5月26日)自这些行写就以来,作为我一番努力的一份意外成果——我本想为“登峰造极”写一份足以传世的事迹——我竟然(几乎不费吹灰之力)廓清了在我看来正是\textbf{恰当的}De Rham 系数定义,至少对 $\mathrm{Spec}\,\mathbb{Z}$ 上有限型概形如此(在我看来,这是最为至关紧要的情形)。当然,与我1966年的观念相比,本质性的新成分是那“含糊的无名氏”的哲学,我在这里(像所有人一样)将克制自己不去点名。\par 我所设想的针对 $\mathrm{Spec}(\mathbb{Z})$ 上有限型概形的方法,应当同样给出任意域(不论特征零与否)上有限型概形的好 De Rham 系数(Mebkhout式或Deligne式,任择其一)。我打算在《数学反思》第三卷“De Rham 系数”一部分里粗描这条进路,连同其他种种“技术性离题”,我的学生尽可前来抄录……}

我自己没有在六十年代继续追问这个问题,我另有足够多的事情要做,而且想着有Berthelot和Deligne在列,这事便落在可靠的人手中(可见人有时会看错……)。Deligne在1969/70年的工作原则上已提供一个零特征中的答案,倘若Deligne把这项工作做到底,大概会令我满意。

\noindent 可是在我心中,这样一套假想的 De Rham 系数理论,即使它能把“离散的”(以一种晶体上同调的形态)上同调与“凝聚的”上同调联系起来,也并\textbf{不}“统摄”凝聚对偶理论。于是我看不出一个Zariski凝聚层如何定义一个“包络晶体”(cristal enveloppant)\fnote{643}{(**)}{(5月26日)也许最好取那个包络“余晶体”(co-cristal)(见注171(ix)B部,其中对余晶体概念有所暗示)。我大概会在前一页边注所许诺的讲稿里回到这个问题。}(注意:在 D-模的语言里,这就是纯量环的扩张 $\mathcal{O}_X\to\mathcal{D}_X$,至少当 $X$ 光滑时如此)——而即使我看见了,所得的晶体(只取 $\mathcal{F}=\mathcal{O}_X$ 就已如此,它给出晶体 $\mathcal{D}_X$)也\textbf{并非} De Rham 型的。然而我暗自思忖:在一个复解析空间 $X$ 上,凝聚对偶(例如以Serre的形式,设 $X$ 光滑、系数局部自由)难道不能作为“离散”对偶的一个“特殊情形”得到吗——这“离散”对偶是Verdier按照 étale 理论的模型发展出来的。就这么着,它看起来有点古怪,并且立刻掀起一大簇问题:如何“以离散的语言”解释对偶模(最高次微分形式)$\omega_X$ 的作用,又如何处理那些怪异的病理现象,它们可在“离散的”对偶中没有对应物?

\origpage{976}\dmark{}是Mebkhout第一个(而且看来直到今天除我之外也是唯一一个人)领会到两种对偶之间确有一层深切的联系,但这联系并不\textbf{表述为}其中之一“统摄”另一个,而是表述为找出第三种对偶理论\fnote{644}{(*)}{关于这“第三种统摄另两种的对偶理论”的细节,见“……作品”一注(n\textsuperscript{o}~171(ii))。}——即 $\mathcal{D}_X$-模(或 $X$ 上“晶体”)的对偶理论,它“统摄”另两者,并且(除此之外)还在“离散”一边自我限于那些同调层为\textbf{解析可构}的 $\mathbb{C}$-向量空间复形。对我来说毫无疑义:这正是对那个“含糊的问题”(而且有点不着边际……)的“正确答案”,而我却从未有机会把它向我那位身后学生提出……

(5月15日)撰写“登峰造极”同时成了一个意外的机会,使我对Mebkhout的工作、对他引入簇的上同调研究的 D-模瑜伽多少熟悉起来。一路行来,这也打捞起一些沉没的记忆。我尤其注意到:早在五十年代末或六十年代初,我就比我十天以前动笔写本注开篇(“古怪的问题”)时自以为的程度,更接近“Mebkhout的哲学”。在任意基概形上的真光滑概形的框架里,我手中曾有一条对偶陈述(以一个相对微分算子复形与其“伴生”复形表述),它“统摄”凝聚对偶与 De Rham 上同调的对偶。技术上讲,这大略相当于Mebkhout对偶定理之代数版本的等价物(在复解析语境中,“……作品”一注(n\textsuperscript{o}~171(ii))已论及)。然而我那条对偶陈述并不使我满意,我不曾想过要发表它,甚至不曾想过为它做宣传,因为在所述形态下,它在我看来太靠近Serre对偶定理(在任意基概形上相对化——这是理所当然的),它便是后者一个多少直接的推论。要得到一条令我满意的陈述,我便须知道如何用微分算子复形作出一个“导出范畴”,以便仿照此前若干年间廓清的凝聚对偶理论之模型,以这些范畴的对象来表述一条内在的对偶陈述。

% ---------------- p.758 ----------------
\noindent 于是所缺的,便是\origpage{977}微分算子复形之间一个(微分的)态射的恰当“拟同构”概念,\dmark{}以便构成一个导出范畴(将这些拟同构形式地逆反)。显然,惯常的定义(经由相应的上同调层)在代数框架里不可用(在超越框架里恐怕同样不可用\fnote{645}{(*)}{我在这里弄错了。Mebkhout向我保证:对微分算子复形之间的一个(微分的)同态而言,它是拟同构(按相应 $\mathbb{C}$-向量空间层复形的朴素意义)当且仅当相应于关联 D-模复形的同态是一个拟同构。事实上(经由映射柱)这等价于说:一个微分算子复形按朴素意义拟零,当且仅当关联的 D-模复形拟零——这似乎是人人皆知的事(至少Mebkhout知道,他在那篇取之不竭的论文里给出了证明……)。})。转向相应的 D-模复形,如今为我昔日的困惑给出一个美妙简洁的答案!

既然看不出有什么现成的定义可给拟同构概念,我便没有试着去弄清楚这事存在与否,以及那里是否确然有一个非凡的导出范畴。那正是在我是唯一对导出范畴感兴趣的时候——那些由凝聚模及其之间的\textbf{线性}态射构成的导出范畴(远为不那么精微)……我并未清楚地感到:这“须廓清一个拟同构概念”的问题(同样有点含糊,若不说是荒诞)触及一个多产的奥秘,而这个奥秘容许一把孩子般简单的“钥匙”!并且有一类非凡的“系数”范畴,只等人来定义它。为此大概须得我的反思在这样一种气氛中继续推进:这反思在那里至少遇到一点兴趣与回响,哪怕只在\textbf{一个}既是局中人便算的对话者那里!

是 De Rham 上同调把我的注意力引向这个当然不言自明的事实:设 $X$ 为域 $k$ 上(比如说)一个代数簇,凝聚层的全局上同调空间不仅是 $\mathcal{O}_X$-线性同态的“函子”,甚至还是\textbf{所有} $k$-向量空间层同态的函子,特别地对微分算子也是如此。正是这个观察促使我萌芽了一份反思,关于一种“凝聚的”(或“拟凝聚的”)对偶理论,其中层之间的“态射”将是微分算子,而不必是线性的。这份反思夭折了,如上所述,而且到这样的地步:它甚至不曾留在\origpage{978}我记忆一角,作为(诸多事项之中的)一件终须\dmark{}有一天弄清楚的事——它沉没了(我想)在完全的遗忘之中,直到几天以前。甚至我约1966年时对晶体的零星反思,也不曾把它打捞回我的记忆,据我所知。然而,这晶体式的反思,在我当时毫无察觉之下(因为当时连这个问题本身都不记得!),将从1966年起为我提供\textbf{另一}把钥匙,某种意义上与Mebkhout那把“对偶”,为我昔日的困惑而设——经由与一个微分算子复形相关联的无穷阶主部复形。我在昨天写的一条页边注里提到了这一点(第946页注(**)),并且我打算在《数学反思》第三卷发展“系数类型”瑜伽的那一部分里详尽回到这里,尤其给出一个成形的定义,定义我所推测的“那些”
% ---------------- p.759 ----------------
好的德拉姆系数(Mebkhout式的或Deligne式的,任择其一),(比如说,在一个 $\mathbb{Z}$ 上有限型概形上)。

\noindent 技术上如此,甚至“心理上”也如此(就彼时早已提出的问题、以及赋予它们力量与生命的那幅总括图景而言),从六十年代后半期起,一切都已就绪,足以廓清这套德拉姆系数的定义。Deligne在我之后曾与正确的概念只差半步,而他不可能不廓清它——倘若那股他授予其对自己人生与作品之绝对权力的力量,不曾为他在这一道路上的反思早早地、斩钉截铁地画上句号……\fnote{646}{(*)}{关于此事,见“……与羁绊”子注(n\textsuperscript{o}~171(viii))。}

发现,并不是双臂倒挽,以锤或以槌猛砸一颗钉子、一柄凿子或一块钢楔。发现,首先是要懂得\textbf{倾听},以一份敬畏与一份全神贯注,倾听事物的声音。新事物并不如一道炫目的光柱,已然成形地从钻石中迸射而出;它也不出自一台机床,无论那机床如何精密、如何强有力。它并不轰然自报家门,披挂着它的贵族世系文书:我是这个,我是那个……它是一个卑微而脆弱的东西,一个娇嫩而活生生的东西,也许是一枚卑微的橡实,从中将长出一株橡树(倘若节令于它相宜……),或是一粒种子,将生出一茎,茎上生出一朵花。它不在聚光灯下诞生,甚至不在日光之下。它不是已知者的果实。它的母亲是夜与昏暗,是那捉不住、无轮廓的雾气——是那逃逸于想要\origpage{979}框定它的言辞之外的预感,是那仍在寻找自身的大古怪问题,或是\dmark{}那份如此含糊、如此飘忽却又如此真实的不自足,伴着那份说不出的(也赖不掉的……)感觉:有什么东西不对劲,有什么东西歪歪斜斜,石头底下藏着鳗鱼……

当我们懂得谦卑地倾听这些低声向我们诉说的话语,并执着地、狂热地追随它们那飘忽的讯息,那么——在历经晦暗而摸索的劳作,或许浑浊而无外观之后——雾气骤然具形、凝缩,成为\textbf{实体},坚固而可触,成为\textbf{形式},可见而清晰。在这孤独的一刻,在全神贯注与寂静之中,那新事物——夜与雾之女——显现了……

\subsubsection*{a3.~自由……}
\tocsubsub{a3.~自由……}
\notehead{171(vii)}{}

(5月4日)\fnote{647}{(*)}{本子注源出“祖先”一注(n\textsuperscript{o}~171(i))的一条页边注——见第945页注(**)。}我并不自称要摆出“成熟之人”或“智者”的姿态,周围尽是不成熟、不负责任的同类——我想,在《收获与播种》的书页里\fnote{648}{(**)}{(5月26日)我甚至可以说:若说撰写《收获与播种》在某一点上对我有所揭示,那揭示的恰是一种“未成熟”的状态,一份“智慧”的欠缺,而断不是它的反面。这也许是全部发现中最出人意料的一个,也是最关键的一个(就其即刻的意涵而言):我对某段过往、对我作为数学家的作品,依恋竟如此之强。这份依恋当初尚以较为克制的外形,在去年3月末向我显现,那是在《愚妄与更新》最后一条注“一段过往的重量”(n\textsuperscript{o}~50)的反思之中。正是直面“葬礼仪式”那粗暴的现实——尤其在其蓄意轻蔑与暴力的诸面相——在我身上启动了强大的自我中心的防卫反射。它们同时也向我揭示出那些将我系于一份过往的联结之强,而我从前竟曾相信那份过往已经离我而去。在过去一年之中,这些联结似乎获得了一份新的活力,而且常常(尤以近来为甚)我感到它们是一种\textbf{重量},说真的,一份累垮人的重量——正如其他曾经压在我身上、而后又化解了的重量……},从我这个人浮现出的形象大概并非如此。然而,至少在我与数学的关系中,我相信可以这样说:在我一生之中,始终保有一份货真价实的单纯\fnote{649}{(***)}{(5月16日)这里大概须对我“私人猎场”的一种占有姿态破个例外——我在《愚妄与更新》“好斗的数学”一节(n\textsuperscript{o}~40)里恰恰点到了它。这些“好斗的”心性本该促使我贬低他人的想法,每当那些想法我这边早已知晓之时。因此可以说(与我在正文中所断言的相反):在这些情形里,我的虚荣确确实实干扰了“我健全的判断”,并且往往倾向于在某一场合激起一种令人气馁的态度,而那时一句善意的鼓励本才是应有的姿态。不过我觉得,这样的处境在我数学家的一生中是例外的,而且并未曾对我的数学创造力构成羁绊。},\origpage{980}同时保有一份对我\dmark{}原初本性的忠诚。我的虚荣,侵入我生命之广不亚于侵入我任何一个同行的生命,然而(就我所记得的)它几乎不曾干扰我健全的判断,也不曾干扰我数学上的嗅觉\fnote{650}{(*)}{关于此事的若干保留意见,见前一条页边注。}。% ---------------- p.760 ----------------
\noindent 此外,也只是在我1970年离去之后,我才开始逐步地、且每一次都带着惊愕地发现:甚至在那些才具非凡的人身上,这种才具竟有时仿佛被“非理性”的成见所歼灭、所无望地封堵——而且越是……便越是顽固!我在这一方面最初的一些经验发生在1976年\fnote{651}{(**)}{(5月16日)这并不真是我在这方面的最初经验——前几年我另有若干次,尤其与Deligne一起,也在我的离去之前的往岁之中。但那些经验始终是零星的,而围绕拉德盖列叶博士论文的那一幕却令人震撼,因为五位数学家(个个水准极高)在行为与不作为上的完全一致,他们相互之间想必不曾串通一气。在那儿,在我与我的朋友皮埃尔一人之关系的种种变故之外,我第一次接触到了“葬礼仪式”。\par 然而,那些所谓“非理性”的因素在所谓“科学”思想中的惊人分量,远远超出“葬礼仪式”的语境,甚至超出一个时代。无须精通科学史(我全然不精通),便可看出:科学史每一步都铭刻着一种巨大惯性的效应,它抗拒一切新思想的萌发,也抗拒新思想在偶然萌发之后的茁壮成长。关于这一方向的反思,尤其见《愚妄与更新》头两部分(“劳作与发现”与“梦与梦者”),第1至8节。},在“进步是挡不住的!”一注(n\textsuperscript{o}~50)中亦有所追述;而关于此事的第一番书面反思,则在“掘墓人——或全体教众”(n\textsuperscript{o}~97)\fnote{652}{(***)}{这一反思在“阴阳之钥”中大有深化,尤其是在涉及同一“集会”的两条注:“天意攸关的境遇——或登峰造极”与“否决(一)——或召回”(n\textsuperscript{o}s~151、152)中。另可参看“肌肉与肚肠(阳葬阴(1))”一注(n\textsuperscript{o}~106),它开启了关于阴与阳的长程反思。}中继续展开,\origpage{981}其语境正是那\dmark{}“葬礼仪式”的特殊情形。同样,也只是渐次地、并且逆着巨大的惰性力量,我才意识到:这些“非理性的”原因并不因此便全不可解,只要人肯费神驻足、把它们探个明白。正因如此,我终于也“接受”了它们,尽管勉强……

回到我这个人,回到我与数学的关系。以我的工作方式而言,我倾向于靠一份份往往为时过早的臆测前行,而不讲求“审慎”\fnote{653}{(*)}{关于这种工作方式,尤其见“兄弟与配偶——或双重签名”一注(n\textsuperscript{o}~134),以及(《愚妄与更新》中的)“错误与发现”一节(n\textsuperscript{o}~2)。};但我把每一个出现的直觉(或“臆测”)追索到底,于是散布在工作最初阶段的大量错误终归自行剔除,让位给一种坚不可破的理解,而且(多半)确确实实触及事物的核心。当要对别人的工作下判断时,我自发的方式便全然不同,尤其当那工作立于我不熟悉的题目或不熟悉的调门上。我觉得,自己那时一向倾向于表现出审慎与谦逊。事实上,这恰恰是接纳我跻身其中的大多数前辈给我做出的榜样,如Cartan、Dieudonné、Chevalley、Schwartz、Leray——且只说这几位。我不记得听他们哪一位曾斩钉截铁地表过态,不论是为非还是为是,论一件他们摸不着其实质的工作。我现在才意识到:这份审慎原是我曾在他处谈及的那股\textbf{尊重}空气的一部分,它浸润着接纳我的那个圈子\fnote{654}{(**)}{另见“受欢迎的异乡人”一节,n\textsuperscript{o}~9。}。我觉得,在这圈子——我在二十多年的岁月里曾自居其中的圈子——最先退化的,正是这份审慎,而它是一份尊重的表征。也许我的记忆骗了我,也许我在自我欺骗;但我觉得,自己相对而言
并未怎样受那种尊重气氛退化的这一面影响。我想,我一直意识到自己在数学上无知有多广,也意识到自己无力领会他人的某项工作——每当那工作落在我的兴趣中心之外,而我的兴趣又往往是强烈聚焦于一处。

\origpage{982}\dmark{}至于那些我能够领会、因而也便能够赏识或评判的他人工作(只要我肯为此费点神),我同样记不起事后不得不承认有什么粗疏的判断失误,无论在贬在褒。我对自己那些观念与直觉的感觉也是这样,无论它关乎一个“好问题”之存在(或不存在),关乎有待探察的一份厚实材质,关乎某个观念的分量,还是关乎我对某一处境或某一事物那份或多或少的完整领会与深入领会。在所有这些场合,倘若确有差错,那差错每一次都是朝“少”的一边。是的——一个新题目或一个新观念之丰富,它在深度与广度上真正的分量,往往只是点点滴滴才充分显露,随着周,随着月,有时甚至是随着年。于是,一份起初正确的感觉(多半如此)——尽管它起初仍旧模糊而弥散——经由深浅粗细不等的“逐例功夫”而得到的这份渐次确认,于我们便如同一份意外、一抹惊叹,随着时辰与日子不断翻新。这无疑是研究工作(无论是数学的研究还是别的研究)那股非凡迷人之力的缘由:每一步,在我们眼前展现的实在,都在丰富、在精微、在深幽上越过我们最大胆的梦……

\noindent 不过我还是回到我对他人工作的领会,当那工作落在与我熟悉的题目上,甚至落在对我而言“火辣”的题目上。我想我可以说:我对一个观念真正分量的敏锐预感(这分量常常连作者本人都不自知),在我的作品中起过至为重要的作用。我在这里首先想到的,便是Serre所起的非凡作用,以及这么一件事:在1955至1970这十五年我作品异乎寻常地丰饶的岁月里,我的大多数观念、我的大多数重大投入,其出发点都在Serre的某个观念或某条进路上,尽管那观念有时貌似平淡。我打算在《主题纲要》\fnote{655}{(*)}{这些“历史评注”在《引论》之三“罗盘与行囊”中已有预告。}的“历史评注”里更详尽地谈这件事。但这么说,并不因此便是对Serre这个人有什么特别的敞开。(斟酌各比例而言)同样的事也在别的数学家身上发生过,无论在我作泛函分析家的过去,还是在我作几何学家\fnote{656}{(**)}{除却许多别的例子,我提一个为例:有一条原理,它把关于任一基座上“有限呈现”之相对概形情状的陈述,归约到该基座是一个\textbf{有限}局部环(甚或一个有限域)之谱的情形;这是一条意义重大的原理,而我是从 D. Lazard一个出色(且极特殊)的结果的一条引人注目的证法思路里把它提取出来的。关于此事,见“暂停!”一注(n\textsuperscript{o}~77)以及该注在第297页的那条页边注(***)。\par (5月16日)我不敢肯定:每逢我的进路受益于他人的一则观念,我都着意标明过。例如,我不记得在 EGA IV 相关的那一段里,我曾着意把Lazard列为那里所发展的一般归约方法之出处。那是当年看来无关宏旨的一桩疏忽。我想,Dieudonné(与我同为 EGA 撰述人)或Serre那些人,他们想必和我一样知道Lazard的这个结果乃是(无疑)同类中的头一个,也未必就认为非引不可(甚至只是适时)——何况这就不在Bourbaki的文风守则里!诚然,Bourbaki在历史评注里总会补上这笔账,而这类评注在 EGA 里、在我作品的别处都付阙如。今天,经受过数学界科学道德在七十、八十年代那可怕的败坏,我在标明出处上会比当年仔细得多——不仅在技术的意义上,也在启发的意义上标明,而后者常常远为关键。在已引用的那些“历史评注”里,我想我至少会补上我在这方面的一些疏漏。}的过去。

\origpage{983}\dmark{}我可以说,在我整个数学生涯中,我因着这份我刚刚试图稍稍描摹的、对待数学的单纯,得到了绰有余裕的“报偿”。这份单纯,
在我生命的别的圈子中常常令我失落,它本身却是一桩福祉。说实话,我作品的丰饶与力量,都归于这份单纯——它也正是\textbf{孩童}的单纯……

% ---------------- p.762 ----------------
\subsubsection*{a4.~……与羁绊}
\tocsubsub{a4.~……与羁绊}
\notehead{171(viii)}{}

\noindent (5月4日)\fnote{657}{(*)}{本子注源出“祖先”一注(n\textsuperscript{o}~171(i))的一条页边注,见第947页的注(*)。}我在这里说错了;在过去两个月里,我重新与本题稍稍接续,记忆也便弄清了些(纠正了些)。其实,Deligne的主要用意,恰恰就在于给出可构的离散(即 $\mathsf{C}$-向量空间)层的一个“纯粹代数的描述”,连同相宜的导出范畴\fnote{658}{(**)}{这里说的是那个范畴(在“作品……”一注(n\textsuperscript{o}~171(ii))中记作 $\mathrm{Cons}^{*}(X,\mathbb{C})$):它由 $X$ 上 $\mathsf{C}$-向量空间层的复形组成,其上同调层解析可构,并把它看作 $D^{*}(X,\mathsf{C}_{X})$ 的一个全子范畴。}。他所引入的系数(办法是在一个 pro-晶体层上施加一条临时的\origpage{984}“可构性”条件,该条件由存在一个适当的\dmark{}\textbf{“脱卸”}来定义,而那脱卸是照着我当初在平展情形或复解析情形引入的那个仿制的)——这些系数都是“量体裁衣”,专为回应这条希求。于是,六种运算的形式主义对这些系数\textbf{理应}存在(在特征零时),这事(在启发上)便“显而易见”;而且严格说来,它甚至能“粗暴而笨拙地”证出来:只要把“外尔原理”用得恰当,归约到基域是 $\mathsf{C}$ 的那个(已知的)情形。

\noindent 于是,只要人肯驻足一想,Deligne竟会撂下一条显然满含希望的进路,不再朝“系数范畴”的描述走去——而在六十年代中期,这一点已经再清楚不过:这些范畴就要在代数簇的上同调中起关键作用——那么这就成了个谜。于是他把最终提炼出一条某种意义上对偶、且更犀利的进路\fnote{659}{(*)}{我对此毫无疑义:倘若Deligne不曾撂下德拉姆系数这个题目(他是从我这里接过来的),他随即势必会(比那位当值无名之辈早八年)发现 D-模的“对偶”瑜伽,并因而熟悉Sato学派的思想。}的差事交给别人——八年后;而那进路随即\fnote{660}{(**)}{“随即”一词并不完全合乎实情的样子(它合乎的是“本该如此”的情形,假如……)。实则,从新的哲学与新的工具俱已就绪之时,到那些定调的人终于发觉此处有物可用(而且可顺手收入囊中……),中间过了三年。}刷新了几何中的上同调主题。这事先前并不怎样震着我,因为Deligne这一理论的发端就在我离去前不久,而那时还没有任何东西能够预示它将被赋予的命运。另一方面,自我离去之后,而且几乎直到最近这几个月,我与上同调主题全然断了联系。

\noindent 我近来曾想到——有些草草,未曾细究——Deligne对他自己投入了整整一年的理论竟如此疏离,其缘由或许在于:他对自己那条以“脱卸”定义的“可构性”判据并不满意。那判据也许显得过于简单,而确实:它远不如Mebkhout1976年在其“对偶”观点中提炼出的那条关于完整性与正则性的局部代数条件来得深。可深思之后,这份“解释”根本站不住脚!一条通往一个\origpage{985}枢机问题的进路“太过简单”,\dmark{}决不至于叫一位神志清明的数学家连那进路带那问题一并放手!他至多会在某天丢下他最初的进路——当他找到另一条,得以对同一问题得到一份更深、更全的视野时\fnote{661}{(*)}{其实,在这桩具体情形里,在我看来根本无所谓为着“老天爷”的进路(且不说出Mebkhout的名字)而“抛开”Deligne进路之事。两条进路互为补充:Deligne那条的长处在于更贴近几何直觉,Mebkhout那条则在技术上更简单(免得借助 pro-对象),并且在种种方面更深。}!
% ---------------- p.763 ----------------
\noindent 只要我对这桩古怪的情状稍稍驻足,事情便清楚起来:在这个情形里,一如在许多别的场合,我的朋友皮埃尔的动机全无数学的成分,甚至全无“理性”的成分。重新思量之下,我才察觉,围绕德拉姆系数的那整套问题意识——它唯有在六种运算与晶体瑜伽的视野中才获得它全部的意义(那瑜伽是我前几年引入晶体topos时提出的,而其精神恰恰就是六种运算……)\fnote{662}{(**)}{我还记得:Deligne讲他自己的理论时,一贯回避使用晶体语言;而那语言本可给他的理论一份更深的维度,把它嵌进一个业已存在的、topos式的上同调形式主义。同样,我如今也看出:无论是Berthelot,还是我别的上同调学生,都早已失去那种深层\textbf{统一性}的感觉——特征 $p$ 的晶体上同调与特征零的晶体现象(那正是他的讨论班的主题)之间的统一性。这些都是蓄意无视一种根本统一性的迹象;那统一性被任意切割,并因此被摧毁。这种蓄意,正是“封堵”的本性:自我中心的力量的介入,与求知冲动无涉。关于这种封堵在另一位上同调学生身上的例证(我虽曾认他天生具有细腻的直觉),参见“共同继承者们……”一注的子注 $n^{\circ}$912。}——我察觉那整套问题意识在我作品中和在我这个人中扎根之深,而且对人人都是明明白白地摆在眼前的。诚然,霍奇系数的问题意识也出自同一位师长,只是那学生早在其内心深处(也许连他自己也不自知)便已在疏离他。可那师承在外人眼里远不那样显豁(无人——连Serre在内——似乎有所察觉)\fnote{663}{(***)}{从Serre论Deligne工作的报告中,似乎确能看出这一点;该报告引在“为一副模糊的骸骨作安魂曲”一注的子注 $n^{\circ}$1651 中(尤其第813页)。关于这一师承的明白宣示,见“把一个个~i~点清”一注($n^{\circ}$164)第一部分第4节(尤其第793页)及其子注 $n^{\circ}$1641。},而最要紧的是:要做的那番意义重大的工作,其头一段并不落在一种表面上\textbf{大张旗鼓}Grothendieck化的视野(“六种运算”或别的……)之内——至少并不明明白白地摆在人人眼前。

\origpage{986}\dmark{}然而,我早已不止一次地着重指出:Hodge--Deligne的上同调理论,在六十年代末一次惊人的起步之后,至今仍停留在幼年阶段——那里被容许的系数唯有常数系数(或退一步说,“光滑”系数,即局部系统在“Hodge--Deligne”意义下的对应物),而连Leray的高次直接像 $R^{i}f_{*}$(且只说这些)那样要紧的运算,在那里竟也毫无通行的地位!“霍奇系数”的正确概念以及在其上建立恰当运算的问题,在Deligne的作品中(就我所知)甚至不曾被\textbf{提及}——而除非我记错,早在我有幸结识他之前,这问题于我便已熟稔。自我离去之后,多年来我有时曾问起:他究竟在等什么,好把他那非建立起来不可的理论最终建立起来(我终究也就厌倦了……)。他答得始终如一:“那太难了……”\fnote{664}{(*)}{后来,这一回答与Deligne手笔的《葬礼颂词》(或以赞词行葬礼)联在了一起,那篇颂词最近也曾提及(见“珍宝”一注 $n^{\circ}$170(iii))。这篇“颂词”以这样一问收场(此问当真值皮埃尔若干斤两……):“他在1970年离开 IHÉS,那时他对数学的热情正渐消渐隐。该相信他在那条为自己划定的路线上所提的问题,已经\textbf{变得太难了}吗?”(着重号是我加的)。这一殷勤的建议在《颂词》的第二段(专论Pierre~Deligne的那一段)里重述,我们从那里得知:逝者的某些猜想“今日仍如其时一般无从接近”,它们想必(至少这是明白的暗示)就是那位Deligne为证明某条“难点众所周知”的猜想所须克服的主要障碍。这些联想让我明白:在我朋友皮埃尔那句千篇一律的回答“那太难了……”里,含有一种嘲弄的弦外之音;这弦外之音必定给了他一种格外辛辣的满足,因为明摆着,那位高大的呆汉、懵懂的逝者万万想不到这弦外之音(更想不到自己已是逝者……)。}。这话并不能说服我,这是肯定——倘若我不曾抽身投入另一桩全然不同的冒险,我也会当即动手,把这“太难”的理论连同德拉姆系数的理论一并发展起来……

% ---------------- p.764 ----------------
\noindent 回过头来看,Hodge--Deligne理论中的停滞,与Deligne本人对待德拉姆系数主题的那副荒唐态度(\origpage{987}那态度在1981年6月吕米尼那场难忘的座谈会上达到顶峰\dmark{},那桩“反常”的不义之举将永远与它相连……),两者之间的平行使我吃惊。这两桩荒唐事,如今在我看来密切相关,而且关涉的层面全然不是数学的层面。诚然,显而易见,霍奇系数的形式主义之发展\textbf{从属于}德拉姆系数的形式主义之发展(这一点我早在1966年前后就已看得清清楚楚,而人们似乎是在那位从不被提起名字的身后学生的踪迹上,最近一两年才正在发现它……)。这个数学事实使两串事实之间的联系、以及它们各自的反常性质,都更加触目:因为这条“客观的”联系本是一份额外的强有力激励(至少对一个“充分掌握着自己能力”的人而言),促使他把两个理论都发展起来——它们本来只可能相互照亮、相互加强。

\noindent 这两个理论各自的停滞(德拉姆理论直到1981年的反常层会议为止,霍奇理论则直到今天),很大程度上在于上同调主题整体的萎靡不振,关于这种萎靡我不止一次有机会谈到过\fnote{665}{(*)}{关于这种萎靡不振,请特别参看“荒芜的工地”(《葬礼行列》,6.),尤其是“工地巡礼——或工具与视景”一注(n\textsuperscript{o}~178)。}。即便撇开人的精神维度不谈,即便只考虑通过“前沿”科学生产获取“收益”这一仅有的因素,这种停滞在我看来也以惊心动魄的方式同时说明了两件事:一是不为人知的隐秘的自我中心力量能在一个人身上取得怎样的统治——即便是在一门自称“无私”的科学的操持之中也是如此;二是这种统治那(显然)反常的性质,它在这里(至少初看之下)似乎总是在与所追求的目的背道而驰\fnote{666}{(**)}{至少,若把“目的”理解为摆在世人面前的那个(“科学的进步”,姑且说),或者甚至是那个丝毫不假的目的——通过积累迫使他人敬重与赞叹的作品来扩大声望,情形便是如此。然而在我看来,即便是这一种“收益”,比起最强大的那些隐秘力量所追逐的满足来,也只是次要的——而我的朋友正是选择让这些力量统治他的存在。}。

\subsubsection*{b1.~五张照片(晶体与 D-模)}

\tocsubsub{b1.~五张照片(晶体与 D-模)}

\notehead{171(ix)}{}

\fnote{667}{(*)}{本子注是“作品……”一注(n\textsuperscript{o}~171(ii))之下的一条子注,性质纯粹是数学的。读者若无意用数学的措辞,对Zoghman~Mebkhout的工作、以及作为簇的上同调理论中一门新的“系数理论”的“D-模瑜伽”,有哪怕些许的领会,尽可略去本子注。下面各页可以看作这套瑜伽、或“Mebkhout哲学”的一段简短导引,用的是一套晶体式的概念装备与整体视景。这一视景早在1966年就已在我心中成形,却被我的同调学生们Deligne、Berthelot、Illusie、Verdier系统地、几乎完全地遮蔽了,而他们本是它的主要保管人。唯一留下的文字痕迹,是我1966年在 IHÉS 所作报告“Crystals and the De Rham cohomology of schemes”的记录稿(I.~Coates 与 O.~Jussila 记录),载 \emph{Dix exposés sur la cohomologie étale des schémas},North-Holland Pub.~Cie(1968)。然而从技术观点看,这个报告包含着晶体上同调的全部启动观念。除了Mebkhout的工作之外,概念层面(或其他层面)似乎不曾有任何真正关键的进展——恰恰相反,我看到一种相对于我六十年代那些观念的惊人退步。不幸的是,那些观念在所引报告中只以极零碎的方式、或字里行间的方式出现;这里一如别处,最大的空缺,是完全没有明确提到德拉姆系数的问题意识,以及应当为这类系数建立的六种运算(与双对偶)形式主义(x)。我曾能够确认:Mebkhout虽然比任何别人都更熟悉我关于上同调的文字著作(以及我的学生们的那一部分),却全然不知这一原初的问题意识(直到两年前还是如此)——而在我看来,就数学“底层”而言(且不谈非理智性质的心理因素),这直到今天仍是他主要的障碍。今后我将把上述1966年的报告简称[Crystals]。\par(x)(6月16日)关于一条订正,见第990页脚注(**)。}

% ---------------- p.765 ----------------
\subsubsection*{(a) “德拉姆系数”图册}

\tocsubsub{(a) “德拉姆系数”图册}

\origpage{988}\dmark{}(5月4日与5月19--20日)我提醒一句:对一个光滑的复解析空间,记 $D_{X}$(或径直记 $D$)为 $X$ 上复解析微分算子的环层(更确切地说,是 $\mathsf{C}$-代数层)。Sato揭示的第一个关键事实,是这是一个\textbf{凝聚}的环层。第二个事实,性质是重言式的,却也同样关键,那就是:局部自由 $\mathcal{O}_{X}$-模的范畴——其中取作态射的,不只是 $\mathcal{O}_{X}$-线性态射,而是此类模之间的微分算子——经由一个先验上\textbf{反变}的函子,作为一个\textbf{全子范畴}嵌入局部自由 D-模的范畴,这个反变函子是\fnote{668}{(**)}{这里写下的同构是 $u\mapsto\varepsilon\circ u$,其中 $\varepsilon\colon D\to\mathcal{O}_{X}$ 即“增广”$\theta\mapsto\theta(1)$。}
\begin{equation*}
F\longmapsto \mathrm{Hom}_{\mathcal{O}_{X}}(F,D^{d})\xrightarrow{\ \sim\ }\mathrm{Op}_{\mathrm{diff}}(F,\mathcal{O}_{X}),
\tag{1}
\end{equation*}
其中 $D^{d}$ 表示 $D$ 配备其由典范\textbf{右} $D$-模结构诱导的 $D$-模结构,\origpage{989}该结构与 $D$ 在其自身上的左运算\dmark{}交换(这些运算使(1)的第二个成员成为一个 D-模)。这个全忠实函子进而在由自由模组成的全子范畴之间导出一个(反)\textbf{等价}。它不容许一个典范的拟逆函子,“与向开集的限制交换”——这想必就是为什么方才这第一个反变函子一般说来并不是一个等价。如果 $C$($C$ 犹如“cristal〔晶体〕”,见下文)表示一个局部自由的 D-模(甚至是自由的,也无妨),人们确实可以给它结合一个以函子方式依赖于 $C$ 的层:
\begin{equation*}
C\longmapsto \mathrm{Hom}_{D}(C,\mathcal{O}_{X}),
\tag{2}
\end{equation*}
这是一个反变函子,看上去可能为(1)提供“那个”拟逆函子的天然候选。麻烦在于,这个层(2)并非以自然的方式配备 $\mathcal{O}_{X}$-模结构,而只配备 $\mathsf{C}_{X}$-模结构(其中 $\mathsf{C}_{X}$ 是 $X$ 上由复数域 $\mathsf{C}$ 定义的常值层)。当 $C$ 由一个局部自由的 $\mathcal{O}_{X}$-模 $F$ 经由反变函子(1)得出时,(2)便典范同构于 $F$ 底下的那个 $\mathsf{C}$-向量空间层。

\noindent 函子(1)(如任何加法函子一样)延拓到复形的范畴:它把 $X$ 上的一个微分算子复形(通常意义下)变成一个局部自由的 $D_{X}$-模复形,而由此得到的(反)函子当然是全忠实的(对于第一个复形范畴中微分算子复形之间的微分态射)。正是在这个意义上,可以说 D-模复形(其分量为局部自由者)“\textbf{推广}”了 $X$ 上的微分算子复形。

\noindent D-模复形的视点相对于微分算子复形的视点有决定性的优势,就是直接嵌入环空间上模复形的瑜伽(最早发展于我1955年的文章《论同调代数的若干问题》\fnote{669}{(*)}{In Tohoku Mathematical Journal, 9 (1957) p.~121--138。}),并由此尤其嵌入\textbf{导出范畴}的瑜伽(那是我在所引文章之后的几年间提炼出来的)。关键观念“\textbf{拟同构}”,在采取微分复形之间的微分态射这一视点时并不跃然眼前,而一旦过渡到相结合的 D-模复形就变得显然。因此,\origpage{990}Mebkhout所引入的视点\fnote{670}{(**)}{(6月8日)此处应当读作:由Mebkhout为了一种新的系数理论的需要,引入Grothendieck的装备之中。不言而喻,“D-模的视点”应归功于Sato,但他是在一种截然不同的 optique〔取向〕中使用它的。},与其说是微分算子复形视点的一种\textbf{推广},不如说是一个\textbf{关键的放松}:\dmark{}正是借助这个视点、而且只有借助它,微分算子复形从今以后才能被用作一门新上同调理论的“系数”,并带着与之相联系的全部丰富直觉。

% ---------------- p.766 ----------------
\noindent 倘若我在德拉姆的系数理论与 $\ell$-进系数理论(后者事实上是Mebkhout发展其哲学的主要灵感来源之一)之间设立一个平行对照,那么我会说:这概念性质的第一步,这“孩童般的”一步,类似于我在1958年引入 étale 层概念时所迈出的那一步(其中已孕育着“topos”这一关键的统一概念)。在同一类比之中,“好上帝定理”(下文即将重述)类似于 étale 上同调中真态射的基变换定理——它(1963年)是 étale 上同调起步阶段的第一个大定理,短短几周之内就使人对 étale 上同调这件工具达到几乎完全的“驾驭”。而在 D-模的框架中(更一般地,在晶体的框架中),要达到对“晶体上同调”(或广义的“德拉姆上同调”,我早在六十年代就已预见这样一门理论)的同样驾驭,那相应的工作\textbf{仍有待完成}——自从Zoghman~Mebkhout终于完成第一次大突破以来,已经过去七年。

\noindent Mebkhout引入的这个新的系数范畴,“包含”(其确切含义见“作品……”一注,n\textsuperscript{o}~171(ii))既包含解析可构的“离散系数”,也包含Serre引入的凝聚系数(我曾把它们系统化为一套“凝聚系数”的上同调理论\fnote{671}{(*)}{这是我在五十年代后半叶在凝聚框架中发展的六种运算与双对偶的形式主义。}),它就是由上同调层(作为 D-模)\textbf{凝聚}的 D-模复形组成的范畴,视作通常的导出范畴 $D^{*}(X,D_{X})$ 的一个全子范畴
\begin{equation*}
D^{*}_{\mathrm{coh}}(X,D_{X}) \quad \text{或} \quad \mathrm{Cris}^{*}_{\mathrm{coh}}(X),
\tag{3}
\end{equation*}
若只限于上同调有界的复形(它们构成全子范畴 $\mathrm{Cris}^{b}_{\mathrm{coh}}(X)$),那么这样一个“系数”在\textbf{局部上}可由一个各次均为有限型自由 D-模、且次数有界的复形来表示;或者,本质上等价地,由一个次数有界的微分算子复形来表示\fnote{672}{(**)}{(6月16日——参看第988页注(*)之末尾。)Mebkhout刚刚向我指出,此说不尽确切——这一问题意识在 loc.~cit.~1.5~d)(第312页)处已略有提及。Mebkhout在他的工作“Poincaré 对偶”(巴黎第七大学“奇点”讨论班,1977--79)中明确提到了它,见于 §4.4(关于 D-模的相对对偶定理)的最后三行。}。

\dmark{}\origpage{991}使用导出范畴时,当然要把基本函子(1)与(2)替换为总体导出函子
\begin{equation*}
F\longmapsto \mathrm{RHom}_{\mathcal{O}_{X}}(F,D_{X})\ ,\quad C\longmapsto \mathrm{RHom}_{D}(C,\mathcal{O}_{X}).
\tag{4}
\end{equation*}
若要寻找与这两个函子同类型的\textbf{协变}函子,首先碰到的就是“标量扩张”函子(在所引注中记作 $N$):
\begin{equation*}
F\longmapsto D\otimes_{\mathcal{O}_{X}}F,
\tag{5}
\end{equation*}
(总张量积),其中张量积还要用到 $D$ 的 $\mathcal{O}_{X}$-右模结构,即 $D^{d}$\fnote{673}{(*)}{众所周知,$D$ 作为 $\mathcal{O}_{X}$-右模或左模都是平坦的(从 $D$ 的典范滤过以及相伴分次层的已知形状即可立见)。由此可知,(5) 中的“总”张量积其实就是一个通常的张量积。}。这个(就 $F$ 而言的)函子相对于(1)有一个缺点,就是不能延拓到那些变元 $F\to F'$ 仅仅是微分算子(而非线性映射)的态射。第二个函子(4)应当看作一个反函子
\begin{quote}
$\mathrm{Cris}^{*}_{\mathrm{coh}}(X)\longrightarrow D^{*}(X,\mathsf{C}_{X})$,
\end{quote}
它也有一个重要的协变“对应物”,由
\begin{equation*}
C\longmapsto \mathrm{RHom}_{D}(\mathcal{O}_{X},C)\overset{\mathrm{def}}{=}\mathrm{DR}(C) \quad \text{(与 $C$ 相结合的“德拉姆复形”),}
\tag{6}
\end{equation*}
给出,其中第二个成员确实借助 $\mathcal{O}_{X}$ 的、由局部自由的有限型 D-模构成的所谓“Spencer 分解”而明确地写成一个德\,拉姆型的复形。(这一分解取自通常的德\,拉姆复形,经由函子(1)取相结合的 D-模复形而得。)用晶体的话来说(下文将予阐明),函子 $\mathrm{DR}$ 明确写出为函子 $C\mapsto \mathrm{Hom}_{D}(\mathcal{O}_{X},C)$ 的总体导出函子,它给每个“D-模”(或“晶体”)结合由其在可变开集上的“水平”截面构成的 $\mathsf{C}$-向量空间层。这是一项\textbf{局部性质}的操作。对于一个“系数” $C$(即一个 D-模或此类复形)而言,恰当的(整体)“\textbf{积分}”(或整体上同调对象)概念,在这里并非通常的函子
\begin{quote}
$\mathrm{R}\Gamma_{X}(C)\simeq \mathrm{RHom}_{D}(X;D,C)$,
\end{quote}
而是(对我而言作为\textbf{晶体整体上同调}函子十分熟悉的)“水平(整体)截面”\origpage{992}函子 $C\mapsto \mathrm{Hom}_{D}(\mathcal{O}_{X},C)$ 的总体导出函子;我把这一总体导出记作 \dmark{}$\mathrm{R}\Gamma_{\mathrm{cris}}(C)$,使得有那些重言同构
\begin{equation*}
\mathrm{R}\Gamma_{\mathrm{cris}}(C)\overset{\mathrm{dfn}}{=}\mathrm{RHom}_{D}(\mathcal{O}_{X},C)\simeq \mathrm{R}\Gamma_{X}(\mathrm{DR}(C)),
\tag{7}
\end{equation*}
亦即 $C$ 在 $X$ 上的晶体上同调,是由取相结合的德\,拉姆复形的(整体)通常上同调而得到的。

\noindent 在 $\mathrm{Cris}^{*}_{\mathrm{coh}}(X)$ 中可以定义一个\textbf{对偶化函子},它给出一个双对偶定理,其模式就是我先后在(交换的)凝聚语境、随后在离散(étale)语境中提炼出的那些。我把它记作 $D$(如在所引诸语境中一样):
\begin{equation*}
D : \mathrm{Cris}^{*}_{\mathrm{coh}}(X)\overset{\sim}{\longrightarrow} \mathrm{Cris}^{*}_{\mathrm{coh}}(X).
\tag{8}
\end{equation*}
这是一个反等价,本质上是对合的,即有一个关于 $C$ 的函子性双对偶同构:
\begin{equation*}
C\simeq D(D(C)).
\tag{9}
\end{equation*}
\noindent {\small [译按:原书式(9)右端印作 $D(D(X))$;据本段的变元 $C$ 及双对偶的含义,此处改为 $D(D(C))$。]}

这个函子可以(通过复合)把反函子(1)与(2)变成协变函子。要记住的简单事实是:若 $C$ 与 $C'$ 互为“对偶”,则其一的德\,拉姆复形(6)等同于另一个的“余德\,拉姆”(2):
\begin{equation*}
\mathrm{RHom}_{D}(\mathcal{O}_{X},C)\simeq \mathrm{RHom}_{D}(C',\mathcal{O}_{X}).
\tag{10}
\end{equation*}
反之亦然。
在微分算子复形上,这个操作 $D$ 可表达出来(就次数而言相差一个 $n$ 的“平移”),办法是过渡到“相伴的”微分算子复形,其分量为 $\mathrm{Hom}_{\mathcal{O}_{X}}(F,\omega_{X})$,它是逐项取相伴算子而得的。于是,D-模的对偶化函子就与 Serre 对偶中那个熟悉的对偶化函子相容:
\begin{equation*}
F\mapsto \mathrm{Hom}_{\mathcal{O}_{X}}(F,\omega_{X})\simeq F\otimes_{\mathcal{O}_{X}}\omega_{X}.
\tag{11}
\end{equation*}
这里 $F$ 为有限型局部自由 $\mathcal{O}_{X}$-模。
\noindent {\small [译按:式(11)照录原书;右端的 $F$ 疑漏对偶记号。按左端的定义,应为 $F^{\vee}\otimes_{\mathcal{O}_{X}}\omega_{X}$,其中 $F^{\vee}=\mathrm{Hom}_{\mathcal{O}_{X}}(F,\mathcal{O}_{X})$。]}

其中 $\omega_{X}$ 表示 $X$ 上最高次微分形式的“对偶化模”。要注意的是,德\,拉姆函子
\begin{quote}
$\mathrm{DR} : D^{*}_{\mathrm{coh}}(X,D)\longrightarrow D^{*}(X,\mathsf{C})$,
\end{quote}
一般\textbf{不}与对偶化函子交换(在第二个范畴中取函子 $\mathrm{RHom}_{\mathsf{C}}(-,\mathsf{C}_{X})$)。但Mebkhout有一条深刻定理(人人都在用它,当然谁也不注明出处,仿佛那不过是一串简单的连环推理),就是说:对于\textbf{完整}的变元,因而对于诱导出的函子
\begin{quote}
$\mathrm{Cris}^{*}(X)_{\mathrm{hol}}\longrightarrow \mathrm{Cons}^{*}(X,\mathbb{C})\ (,\hookrightarrow D^{*}(X,\mathsf{C}))$
\end{quote}
\origpage{993}而言,与对偶化函子是交换的。我不在此“重述”完整性的条件\dmark{},只愿指出:一个 D-模复形是完整的,当且仅当它的上同调层都是完整 D-模;而这是一条在 $X$ 上局部性质的条件,并且还是“代数的”。另一方面,Kashiwara的可构性定理(Kashiwara陈述这定理时只针对一个完整\textbf{模},那时无论他自己还是任何人——除了Mebkhout——都不在导出范畴下工作……)蕴涵:德\,拉姆函子限制到完整复形上,确实取值于 $\mathrm{Cons}^{*}(X,\mathbb{C})$。引入Mebkhout的\textbf{正则性}概念——它同样是局部性质,并且是“代数的”\fnote{674}{(*)}{我提醒一句:Mebkhout对正则性的原始定义是超越性质的。关于一个“纯粹代数”的转译,请参见关于德拉姆系数的预定报告(“Mebkhout风格”或“Deligne风格”),载《数学反思》第三卷。}——就得到“好上帝函子”(又名Mebkhout函子)
\begin{equation*}
m : \mathrm{Cris}^{*}(X)_{\mathrm{hol.reg.}}\overset{\sim}{\longrightarrow} \mathrm{Cons}^{*}(X,\mathbb{C}),
\tag{12}
\end{equation*}
它这一次是一个\textbf{等价}(如我们在“作品……”一注(n\textsuperscript{o}~171(ii))中所见),因而与自然的对偶化函子相容。正是那个拟逆函子
\begin{equation*}
M : \mathrm{Cons}^{*}(X,\mathbb{C})\overset{\sim}{\longrightarrow} \mathrm{Cris}^{*}(X)_{\mathrm{hol.reg.}}\hookrightarrow \mathrm{Cris}^{*}_{\mathrm{coh}}(X),
\tag{13}
\end{equation*}
使得可以把 $X$ 上“可构离散系数”(即 $\mathsf{C}$-向量空间层)的范畴看作 $D^{*}(X,D)$ 的一个全子范畴、更确切地说是 $D^{*}_{\mathrm{cons}}(X,D)=\mathrm{Cris}^{*}_{\mathrm{coh}}(X)$ 的一个全子范畴,而人们将时而把它解释为一个“晶体”系数的范畴。

\noindent (5月19日)就目前而言,我们可以说,已经用三种不同的“语言”或“视点”,如同用三张不同的“照片”,描绘了同一个实在,或(本质上)同一类“系数”,即所谓“德拉姆系数”:有 $\mathsf{C}$-向量空间层及其复形这一视点(“拓扑”拍摄),配以一条“解析可构性”条件\fnote{675}{(**)}{我提醒一句:解析空间 $X$ 上的一个 $\mathsf{C}$-向量空间层称为“解析可构的”,如果在每一点的邻域内,它有一条合成列,其相继因子形如 $i_{!}(F)$,其中 $i : Y\to X$ 是 $X$ 的一个解析子空间 $Y=Z\backslash T$ 的包含映射(这里 $T\subset Z$ 是 $X$ 的两个解析闭子空间),而 $F$ 是 $Y$ 上一个局部自由的有限型 $\mathsf{C}$-层(或“$\mathsf{C}$-向量空间的局部系统”)。},起着有限性条件的作用(为了能写出 Riemann--Roch\origpage{994} 型的定理,这一点尤为不可或缺,这类定理涉及恰当的“Euler--Poincaré 示性数”与恰当的“\dmark{}Grothendieck 群”)。还有“微分算子复形”的视点,以完整性与正则性两条条件取代可构性条件。还有“D-模复形”的视点,随之而来的是凝聚性、完整性与正则性条件。

% ---------------- p.769 ----------------
\noindent 第二张“照片”(从“分析”的角度拍摄)很诱人,因为它用“经典”的语汇就能为我们所理解,而且它呈现的对象——即微分算子的复形——在我们看来“尺寸”适中,而 D-模,即便是凝聚的(且不说 $D$ 本身!),戴上“$\mathcal{O}_{X}$-模”这副眼镜去看,也显得庞大。然而就技术而言,反倒是它们提供的照片更为完整。事实上,虽然局部上“每个上同调凝聚、次数有界(姑且说)的 D-模复形都可以经由(1)用一个微分算子复形来表示”这件事看着“明白”,整体上还能如此可就难说了,除非对 $X$ 施加苛刻的假设(比如“Stein 簇”,或者——在代数框架中——一条拟射影性假设)\fnote{676}{(*)}{当然,没有什么能阻止人们从 $X$ 上微分算子复形的范畴、以及这些复形之间的微分态射出发,把那些“拟同构”(由过渡到相结合的 D-模复形来定义)形式地“反转”,从而构造出一个“导出范畴”。人们将会得到(我推测)$\mathrm{Cris}^{*}_{\mathrm{coh}}(X)$ 的一个\textbf{全}子范畴,但在缺乏“Stein”或“$X$ 射影”(在代数情形中至少拟射影)一类假设时,恐怕不会是这整个范畴。}。

\noindent “照片”1 的好处在于:当 $X$ 不再被假定为光滑、而是一个一般的复解析空间时,它仍然保有意义。相反,照目前这个样子,“照片”2 和 3 只有在光滑性假设之下才讲得通。人们当然仍可以在 $X$ 上不假定光滑就定义环层 $D_{X}$,也仍然有一部自明的“词典”,在微分算子复形(其分量为局部自由 $\mathcal{O}_{X}$-模)与 D-模复形(其分量局部自由)之间来回对译;只是(据说)$D_{X}$ 不再凝聚了,真可惜!在奇异情形中,恐怕很难照光滑情形中那个已知榜样,提炼出一条“好上帝定理”。而另一方面显而易见的是:在奇异情形中我们同样需要 2 或 3 那一类的照片,因为“照片”1 是\textbf{就其性质而言属于超越}的:天真地照搬它、用代数簇的 Zariski 拓扑或 étale 拓扑来表述,所得的“系数”将特殊得太离谱,不堪使用(因为比起超越拓扑来,\origpage{995}\dmark{}这些拓扑太过粗糙)。相反,“照片”2 和 3一开始把视野限制在“光滑”之内,却在“抽象”代数几何中仍保有意义(比如说,一开始在特征零的域上),这就构成(对我而言)它们主要的魅力。这就是说,非得给它们做\emph{放大},把奇异簇也纳入视野不可。

\noindent 这看来不曾让Mebkhout费神,他忙着别的要事——我向他提出这个问题时,他当即的念头是这样的。设 $X$ 作为一个闭解析子空间嵌入一个光滑簇 $X'$ 之中。于是范畴 $\mathrm{Cons}^{*}(X,\mathbb{C})$ 可解释为 $\mathrm{Cons}^{*}(X',\mathbb{C})$ 的全子范畴,其成员是在 $U=X'\setminus X$ 上限制为零的对象(即支撑含于 $X$ 的对象)。而由好上帝定理,后者也可用“照片”2 或 3 的语言,解释为 $X'$ 上那些在 $U$ 上限制为零的“德拉姆——Mebkhout系数”的范畴。停留在“德拉姆——Mebkhout系数”(即“照片”2、3)的框架内,先验地验证下面这点应当是容易的:这个范畴——在等价的意义下,而该等价本身除差一个唯一同构外即被确定——不依赖于为 $X$ 选定的“光滑化”$X'$。我自己就干过许许多多这样的事,我乐意相信它行得通。而倘若 $X$ 并不“可光滑化”,那也无妨(Mebkhout如是说),我们可以“做上同调下降”,从这些局部的碎片重建出一个整体范畴;或者引入 $X$ 的开集之“光滑化位”,在上面工作。看来确有办法应付,只是与其用一个“光滑化位”(Mebkhout为作出回应的需要而临时搭建的
回答,在一场始终停留在柏拉图式的对话中),这个位在我看来显得强烈冗余,为什么不索性用晶体site呢,它早已受过检验(即便看来它已被那些曾是我弟子的人,带着一种动人的浑然不觉,抛诸脑后……)?更何况,早在1966年我廓清晶体瑜伽的启动观念之时,我就已经看得很清楚:未来的“德拉姆系数”本来就应当用晶体的措辞来表达!

\noindent 这于是引我从抽屉深处取出一张照片,它可有时间积起灰尘了,可怜见的——然而一旦把它吹净,它在我看来恍若全新,且清晰得无以复加。况且,这是我去年写“我的孤儿们”一注(n\textsuperscript{o}~46)时(甚至还没遇见“葬礼”之前)最早想到的事情之一,当时我隐约感到,是时候有人怀着敬意谈一谈那些值得尊重的事物了……此外,自从Mebkhout\origpage{996}\dmark{}跟我谈起 D-模以来(1980年——天晓得,那时我可并不“在行”!),我就禁不住把它们想成是“晶体”,并且把“D-模”与“晶体”(即 $\mathcal{O}_{X}$-模的晶体)当同义词使用,当然(明显地)偏爱后者。

\noindent 我于是来到许诺过的第四张照片,即“晶体”照片。先设 $X$ 光滑。在 $X$ 上给定一个 D-模 $F$,等同于给定一个 $\mathcal{O}_{X}$-模,外加一道附加结构,而这道结构可以用若干等价的方式来表达。其中重言式的一种,是说把 $\mathcal{O}_{X}$ 在阿贝尔层 $F$ 上的运算“延拓”为环层 $D_{X}$(它包含 $\mathcal{O}_{X}$)的一次运算。由于 $D_{X}$ 由 $\mathcal{O}_{X}$ 与导子的加法子层生成,可见在 $F$ 上给出所谓“\textbf{可积联络}”,也就是说给出这样一个法则:对 $X$ 的每个开集 $U$ 上的每个导子 $\xi$,结合 $F$ 的一个“$\xi$-导子” $\theta_{\xi}$,关于 $\xi$ 是线性的,并且与导子的“括积”运算相容\fnote{677}{(*)}{当然,还须有一条与限制到开集相容的条件。}——二者是一回事。可以说,这是 $F$ 上一道“微分”性质的结构,阶为 1。

\noindent 由于我们处在特征零之中\fnote{678}{(**)}{在下面,我们可以免除一切特征假设(比如说,在光滑相对概形的框架中),办法是把沿对角线的 $X\times_{S}X$ 形式完备化,换成“S-除幂”的形式完备化。这同样导致,对 $X$ 上的 $\mathcal{O}_{X}$-模层 $F$,要把它的“无穷阶主部”的 pro-层 $P^{\infty}(F)$,换成“(无穷阶)S-除幂主部”。在对偶方面,这就等于把相对微分算子的环层 $D_{X/S}$(即使 $S$ 是诺特的,它也毫无凝聚性可言),换成 $\mathcal{O}_{X}$ 在 $\mathcal{O}_{S}$ 上的相对导子的“包络”环层(据Mebkhout向我保证,它会是凝聚的!)。事实上,这就是德拉姆系数的概念语境,它将推广Mebkhout的 D-模语境,尤其是为了发展一套 $\mathrm{Spec}\,\mathbb{Z}$ 上有限型概形的德拉姆系数理论。},这道结构也可以被诠释为一道更丰富的结构,即一道无穷阶的微分结构,我曾把它称为 $F$ 上的“\textbf{分层}”(此时 $F$ 便得名“\textbf{分层模}”)。表达一个分层的一种方式,是把它写成 $F$ 上(相对于态射 $X\to\mathrm{pt}$ 的)一道“\origpage{997}\textbf{无穷阶无穷小下降资料}”,或者更确切地说,写成\dmark{}沿对角线取 $X\times X$ 的形式完备化之上、$F$ 的两个拉回像(由两个典范投影 $\mathrm{pr}_{1}$ 与 $\mathrm{pr}_{2}$ 给出)之间的一个同构,它延拓对角线上的恒等,并且此外还满足一条适当的“传递性条件”。

\noindent 从一道可积联络过渡到一道“无穷小下降资料”(即分层结构),这是一个新观念——而且“平凡”,一如我有幸发现的所有新观念!然而,只有一旦用模的晶体这一概念对它重新诠释,它才取得全部的力量。事实上可以看出,$F$ 上所说的这道结构,也就等同于给出:对 $X$ 的每个开集 $U$ 的每个“无穷小邻域” $U'$,给出 $F|_{U}$ 到 $U'$ 的一个\textbf{延拓} $F_{U'}$(总之,$F$ 在无穷小邻域上方“生长”,犹如一颗“晶体”——晶体有种种品类,此处是模的晶体……)——这个延拓,就限制到 $U$ 的一个开子集 $V$ 的概念、以及同一个 $U$ 的无穷小邻域(或称“加厚”)$U'$、$U''$ 之间的态射(当然是诱导 $U$ 上恒等的态射)而言,表现得一如人们预料。

% ---------------- p.771 ----------------
\noindent 晶体视点的好处,在于所要研究的对象(即 D-模)可以诠释为一个适宜位\fnote{679}{(*)}{要注意,在晶体site上并不能找到所有的模层,而只能找到那些满足一条简单的附加条件的模层(在[Crystals]中称为“特殊的”层)。}上的“寻常”模层,这个位以交换局部环为环层,即由 $X$ 的各个开集 $U$ 的加厚 $U'$ 构成的“\textbf{晶体site}”(晶体结构层简单地就是 $U'\mapsto\Gamma(U',\mathcal{O}_{U'})$)。由此,人们就拥有一部与这种局面相联系的全部几何直觉武库。我在1966年发现并当时为之震惊的一条卓越关系,就是:晶体site(或它所对应的晶体topos)以结构层为系数的上同调(或者更一般地,以 $F$ 为系数,至少当 $F$ 在 $\mathcal{O}_{X}$ 上凝聚时),等同于 $X$ 的\textbf{德拉姆上同调}(此处是以 $F$ 为系数,即以 $\mathrm{DR}(F)$ 为系数的 $X$ 的寻常超上同调)。这就是晶体上同调的起点\fnote{680}{(**)}{在这里又是如此——那些启动观念“平凡”到这种地步,实在不必为这一点点东西而局促:一个人后来花了十五年的生命,去发展它的一个小角落(而把其余忘掉……)。}。

\origpage{998}\dmark{}于是我们有一部完美的词典,在我已引的1966年那些讲演\fnote{681}{(*)}{见[Crystals]报告,引于本子注的第一条脚注(第988页注(*))。}中有详尽的阐述,它在 $X$ 上四类对象、或一个 $\mathcal{O}_{X}$-模上四种类型的结构之间建立对应:
\begin{quote}
(Cr)\quad D-模\\
\qquad 带可积联络的 $\mathcal{O}_{X}$-模\\
\qquad 分层模(无穷阶无穷小下降资料)\\
\qquad $\mathcal{O}_{X}$-模的晶体
\end{quote}
这部词典有效,对 $F$ 不附加任何凝聚性或拟凝聚性类型的限制。然而要注意,若比较两个极端项
\begin{quote}
D-模 $\Longleftrightarrow$ $\mathcal{O}_{X}$-模的晶体,
\end{quote}
那么在这两个语境中“凝聚性”的自然概念\textbf{并不相对应}。晶体结构层是凝聚的,但晶体topos上的凝聚模,恰好对应于作为 $\mathcal{O}_{X}$-模凝聚的 D-模,此时它们甚至是有限型局部自由的。它们构成的范畴,经由相对于 $\mathsf{C}_{X}\to\mathcal{O}_{X}$ 的“标量扩张”函子,典范等价于局部自由的 $\mathsf{C}_{X}$-模层范畴,即 $X$ 上“$\mathsf{C}$-向量空间局部系统”的范畴。因此对这类对象,就有五种可能的描述(若算上表(Cr)先前的四项,就是五张“照片”)!但这只是些性质极其特殊的“系数”\fnote{682}{(**)}{事实上,在此是 D-凝聚性(我六十年代未曾察觉它)才是重要的有限性概念。},在(de~Rham--Mebkhout)那些我们感兴趣的系数之中。

\noindent 我们还是回到上面表(Cr)的四张照片,看看不再假定 $X$ 光滑时会怎样。所考虑的四类对象仍保有意义。另一方面,看来头两类并不构成重要的范畴——更确切地说,人们自然地遇到的、作为“具有几何意义”的所有 $D_{X}$-模与所有带可积联络的 $\mathcal{O}_{X}$-模,都(在一个显然的意义下)“来源”于分层模,而这些分层模仍可诠释为 $\mathcal{O}_{X}$-模的晶体,与光滑情形一样\fnote{683}{(***)}{这个论断下得仓促,而且照原样是错的。要使它为真,必须把由 $X$ 的开集的一切无穷小加厚构成的“晶体site”,换成由那些在 $X$ 上局部容许一个收缩的加厚构成的子位(称为“分层位”)(当 $X$ 光滑时,这个条件自动满足)。若给定 $X$ 上一个分层模 $F$,它经过这样一个收缩的拉回,除了差一个唯一同构之外,不依赖于所选定的收缩,由此得到 $F$ 到所考虑加厚上方的一个“典范延拓”。因此可见,当 $X$ 不光滑时,$F$ 上一道晶体结构比一个单纯的分层“更丰富”,因为它使我们能把 $F$(即把它“养大”)延拓到 $X$ 的开集的\emph{任意}无穷小邻域上方,尤其是(而这是一件特别重要的事)延拓到浸没在一个光滑环境空间中的 $X$ 的各阶无穷小邻域上方。事实上,在分层模概念与模的晶体概念之间,最新颖、最富生机的是后者。正是它注定要主宰德拉姆系数的理论。关于这一点,我“提醒”一句:对 $X$ 上真且光滑的相对概形 $Z$,$Z$ 在 $X$ 上的相对德拉姆上同调(无论在超越语境中还是在代数语境中……)都“不仅”配备一个分层,而且实实在在地配备一道晶体结构,使它在任何无穷小邻域上方“生长”。这是一个至关重要的数学事实,而Deligne早在1969年我离去之前就把它忘了,当时他用分层 pro-凝聚模、而不用更强的晶体版本、即 pro-凝聚模的晶体,来描述德拉姆类型的系数。应该说,我的名字与分层模概念的联系,不如我与模的晶体概念的联系那样显赫(这一概念如此自然,人们简直要赌咒说它必定上溯到上个世纪;而模的晶体概念则外观远不那么“传统”)。关于这一点,见“……与羁绊”(子注 n\textsuperscript{o}~171(viii))中的思考。}。

% ---------------- p.772 ----------------
\origpage{999}\dmark{}我承认,由于未曾就此深思,我还看不太清:对浸没在一个(姑且说)光滑的 $X'$ 中的 $X$,$X$ 上的晶体与 $X'$ 上的晶体之间的确切关系(即便 $X$ 本身光滑,也是如此)\fnote{684}{(*)}{(5月26日)随着余晶体(co-cristal)概念的引入,局面为我大为澄清,下文 D) 中将提及这一概念。}。

\noindent 可以肯定的是,晶体site、或不如说带环层结构的晶体topos $X_{\mathrm{cris}}$,以协变的方式依赖于解析空间 $X$,即:若 $f : X\to X'$ 是解析空间之间的一个态射,就由此导出
\begin{quote}
$f_{\mathrm{cris}} : X_{\mathrm{cris}}\longrightarrow X'_{\mathrm{cris}}$ ;
\end{quote}
由此特别有这些带环topos上模层的一个“直接像”函子。我们想要理解这个运算(尤其在闭浸入 $X\hookrightarrow X'$ 的情形),并理解在什么条件下一个晶体被变成晶体。在闭浸入的情形,我们还想要这个函子正合。这里的思想是这样的:若 $F$ 是导出范畴 $D^{*}(X_{\mathrm{cris}},\mathcal{O}_{X_{\mathrm{cris}}})$ 的一个对象,$F'$ 是它经由 $f_{\mathrm{cris}*}$ 的总体导出函子的像,并且再假定 $X'$ 光滑,那么 $F'$ 为\origpage{1000}完整正则这一条件\textbf{本不应}依赖于把 $X$ 浸没到一个\dmark{}(光滑)空间 $X'$ 中的所选浸入。若情形果真如此,那么就将把 $X$ 上de~Rham--Mebkhout的晶体系数的范畴,定义为(导出范畴的)由前述条件定义的全子范畴(显然在 $X$ 上是局部的)。

\noindent 因此,且不论一项二十年前就该做、而显然至今仍未做的基础工作(关于晶体的模上的基本运算),可以说在 $X$ 是任意解析空间(未必光滑)的情形,只剩下\textbf{两}张照片(而非四张)来描述我们所得的“德拉姆系数”:一是 $\mathrm{Cons}^{*}(X,\mathbb{C})$ 原封不动,二是(目前仍是假设性的、而且照原样我还看不真切的\fnote{685}{(*)}{我在下面提及一张“第五张照片”,它对我而言眼下就清晰得多,用以一种纯代数的晶体措辞捕捉“好的”德拉姆系数,且不加任何光滑性假设仍保意义。这张照片是从某种意义下与de~Rham--Mebkhout照片“对偶”的角度拍摄的。})de~Rham--Mebkhout系数 $\mathrm{DRM}^{*}(X)$ 范畴,我刚为它冒险提出一条定义原则。范畴 $\mathrm{Cons}^{*}(X,\mathbb{C})$,其描述从超越观点看毫无问题,然而一旦进入代数语境就\textbf{消失}了。这就使得廓清 $\mathrm{DRM}^{*}(X)$ 的一个好定义、使其在该语境中仍保意义的需要,变得显而易见。而对我来说同样清楚的是,这张照片的恰当“框架”——它于是看来(至少初看之下)是唯一剩下的一张——是由晶体的模形成的\fnote{686}{(**)}{我把 $X$ 上一个“\textbf{晶体的模}”称为晶体topos $X_{\mathrm{cris}}$ 上的一个模层。因此可以把模的晶体视为晶体的模的特例。}

\noindent 我况且承认,即便在 $X$ 光滑的情形,在Mebkhout对“德拉姆”系数的描述中,我也并不怎么认得过来:用好上帝函子来说,它不尊重自然的乘法结构;而据说,与之相容的是Mebkhout的反函子(下文(b)中将予讨论)。\fnote{687}{(***)}{这个“据说”是一种略显轻慢的说法(近乎“新风格”的腔调……),为的是一笔略过一条漂亮的定理,它照例出自同一位当值无名之辈(但我似乎得知,其年份比好上帝那一枚要新)。例如,对 $K$ 的两个解析闭子空间 $Y$ 与 $Z$,它蕴含局部上同调上下述公式,它漂亮得显然不可能是真的(然而……):
\begin{quote}
$R\Gamma^{\mathrm{alg}}_{Y\backslash Z}(\mathcal{O}_{X})\simeq R\Gamma^{\mathrm{alg}}_{Y}(\mathcal{O}_{X})\otimes^{L}R\Gamma^{\mathrm{alg}}_{Z}(\mathcal{O}_{X})$,
\end{quote}
我敢打赌,某些漂亮的先生们近日就会把它揣进兜里,俨然“仿佛他们一直就知道”——同时等着把它拍卖给其中最漂亮的那一位……}。更何况,这个函子不与“六种运算”交换。\origpage{1001}与Mebkhout系数相联系的直觉\dmark{},初看之下与离散系数相联系的直觉,性质似乎截然不同。从某种观点看,这是一个好处——我们拥有两张从截然不同角度拍摄的照片!这只是使习惯于从其中一个角度观看的人,更难在从另一角度拍摄的照片中辨认自己。

\noindent 事实上,除了已经浏览过的四张照片(我指的是就“德拉姆系数”而言)之外,还有\textbf{第五张}\fnote{688}{(*)}{这样,我就不只是兑现了本注“五张照片”这一标题的许诺:我实际上揭示了\textbf{两系}五张照片,第一系只描述“德拉姆系数”,第二系则描述一般的晶体系数。},我把它留作储备:那是Deligne的那一张,用分层 pro-模拍摄\fnote{689}{(**)}{如先前一条脚注(第998页注(***))所指出,这张德拉姆-Deligne照片是用一个略有扭曲的“镜头”拍摄的(且说其原因,也超出制造工匠的能力所及)。必须把它修版,还要把它放大,让它走出特征零的框框。这将在《数学反思》第三卷中完成,在那里我亲爱的前弟子们可以安心地来汲取一切“无用细节”和别的“技术性离题”,而那些是他们近二十年来一直未曾有闲情自己去找出来的——我放任他们凭一个华丽的主题自己摸爬滚打……}。它的优点,是极其贴近地“贴合”可构离散层的直觉:一个“零次”对象对应于一个同类型的对象,张量积与拉回像的概念通过Deligne的等价相互对应;因此六种运算的整体也将如此(它们确实可以用这两个概念来描述)。另一方面,从“德拉姆-Deligne”系数 $\mathrm{DRD}^{*}(X)$ 到“de~Rham--Mebkhout”系数 $\mathrm{DRM}^{*}(X)$ 的过渡运算,在我看来原则上特别好理解,用 $\mathcal{O}_{X}$-模上的运算(“$\mathcal{O}_{X}$-对偶”)来说(至少一开始,对光滑的 $X$ 是如此)——我在先前一条脚注\fnote{690}{(***)}{这条“先前的脚注”在此期间已变成本注“五张照片”的(c)部分。}中已经提及。因此我有这样一种感觉:在此我处在一片既坚实又熟悉的土地上,\origpage{1002}一旦有暇,它定能使我认明方向。我甚至本想在本注中\dmark{}勾勒Deligne的观点,并把它与Mebkhout的观点、与我在已引的1966年那些讲演中勾勒的形式主义联系起来。但这个子注开始变长,而且越来越变成一处离题支节!因此我宁愿把此事推迟到《数学反思》第三卷,在那里我还想给出 $\mathrm{Spec}\,\mathbb{Z}$ 上有限型概形上“好的”德拉姆系数(Deligne风格或Mebkhout风格,任选其一)的描述。

\subsubsection*{(b) 好上帝公式}

\tocsubsub{(b) 好上帝公式}

\noindent (5月5日与5月21日)我想在此回到对Mebkhout函子(亦称“好上帝函子”)的描述:
\begin{equation*}
M : \mathrm{Cons}^{*}(X,\mathbb{C})\longrightarrow \mathrm{Cris}^{*}_{\mathrm{coh}}(X)\overset{\mathrm{dfn}}{=}D^{*}_{\mathrm{coh}}(X,D_{X})
\tag{1}
\end{equation*}
,其中 $X$ 为一光滑复解析空间。

% ---------------- p.774 ----------------
如“作品……”一注(n\textsuperscript{o}~171(ii))中所说,这是一个性质深邃的函子,它定义为德拉姆函子 $\mathrm{DR}$ 限制到 $\mathrm{Cris}^{*}_{\mathrm{coh}}(X)$ 的全子范畴 $\mathrm{DRM}^{*}(X)$(即 $X$ 上“de~Rham--Mebkhout系数”的范畴)上的函子的拟逆:
\begin{equation*}
m\overset{\mathrm{dfn}}{=}\mathrm{DR}|_{\mathrm{DRM}^{*}(X)} : \mathrm{DRM}^{*}(X)=\mathrm{Cris}^{*}(X)_{\mathrm{hol\,reg}}\longrightarrow \mathrm{Cons}^{*}(X,\mathbb{C}),
\tag{2}
\end{equation*}
而这个函子结果是一个等价(“好上帝定理”)。事实上,Mebkhout还得到了函子 $M_{\infty}$ 的一个卓越的直接描述,它是把函子 $M$ 经过环同态
\begin{equation*}
D_{X}\longrightarrow D^{\infty}_{X}
\tag{3}
\end{equation*}
的“标量扩张”函子 $i$ 导出的,其中 $D^{\infty}_{X}$(或 $D_{X}^{\infty}$)表示“$X$ 上无穷阶微分算子”的环,即(由定义)把层 $\mathcal{O}_{X}$ 看作复拓扑向量空间层时它的($\mathsf{C}$-自同态)环。众所周知,$D^{\infty}_{X}$ 在 $D$ 左、右都忠实平坦,因此环扩张函子
\begin{equation*}
i : \mathrm{Cris}^{*}(X)=D(X,\mathcal{D})\longrightarrow D(X,\mathcal{D})\longrightarrow D(X,\mathcal{D}^{\infty})\overset{\mathrm{dfn}}{=}\mathrm{Cris}_{\infty}(X)
\tag{4}
\end{equation*}
\noindent {\small [译按:原书式(4)连续两次写出 $D(X,\mathcal{D})$,两者之间有箭头;此处照录,其重复疑为排印之误。]}

的总体导出函子由一个寻常的张量积来阐明。且提一句:人们不知道环 $D^{\infty}_{X}$ 是否凝聚,但显然没有它也过得去。人们把由 D-模复形中那些“完整”的组成的全子范畴
\begin{quote}
$\mathrm{Cris}^{*}_{\infty}(X)_{\mathrm{hol}}\hookrightarrow \mathrm{Cris}^{*}_{\infty}(X)$
\end{quote}
定义为:局部地(经由函子 $i$)可从一个完整的 D-模复形 $C$ 导出的条件。\origpage{1003}(由下文重述的双重好上帝定理将可得:那时甚至可以取 $C$ 同时完整且正则,即一个“de~Rham--Mebkhout”系数,并且这将在整个 $X$ 上除唯一同构外确定 $C$……)考虑函子 $M_{\infty}=iM$,它嵌入交换图表
\begin{equation*}
\begin{tikzcd}[row sep=2.5em, column sep=1.2em]
& \mathrm{Cons}^{*}(X,\mathbb{C}) \arrow[dl,"M"'] \arrow[dr,"M_{\infty}"] & \\
\begin{gathered}\mathrm{DRM}^{*}(X)\\[3pt]={}\mathrm{Cris}^{*}(X)_{\mathrm{hol\,reg}}\end{gathered}
  \arrow[rr,"i"] & & \mathrm{Cris}^{*}_{\infty}(X)_{\mathrm{hol}}
\end{tikzcd}
\tag{5}
\end{equation*}
结果发现(或者不如说,是那位无名工匠证明的……)函子 $M_{\infty}$ 本身也是一个范畴等价(因此 $i$ 也是)。人们也可以把它作为“德拉姆”型的函子 $m_{\infty}$(与 $m$ 类似,定义在 $\mathrm{Cris}^{*}_{\infty}(X)_{\mathrm{hol}}$ 上)的拟逆而得到。为了描述函子 $M_{\infty}$,更方便的是描述反函子
\begin{equation*}
\Delta_{\infty}\overset{\mathrm{dfn}}{=}M_{\infty}D=D_{\infty}M_{\infty}=i(MD)=i(DM),
\tag{6}
\end{equation*}
其中 $D$ 表示已提及的对偶化函子,在 $\mathrm{Cons}^{*}$ 或 $\mathrm{DRM}^{*}$ 中;而 $D_{\infty}$ 表示在 $\mathrm{Cris}^{*}_{\infty}(X)_{\mathrm{hol}}$(乃至 $\mathrm{Cris}^{*}_{\infty\mathrm{coh}}(X)$)中也存在的类似的对偶化函子。(注意:(5) 中出现的三个函子都与对偶化函子交换。)于是 $\Delta_{\infty}$ 的拟逆 $\delta_{\infty}$ 由公式

% ---------------- p.775 ----------------
\begin{equation*}
\delta_{\infty}\overset{\mathrm{dfn}}{=}Dm_{\infty}=m_{\infty}D_{\infty}
\tag{7}
\end{equation*}
给出,它与(6)类似。于是人们得到Mebkhout对 $\Delta_{\infty}$、$\delta_{\infty}$ 的表达式,由下面两个对称得惊人的公式给出:
\begin{equation*}
\left\{\begin{aligned}
\Delta_{\infty}(F)&=\mathrm{RHom}_{\mathsf{C}}(F,\mathcal{O}_{X}),\\
\delta_{\infty}(C)&=\mathrm{RHom}_{D}(C,\mathcal{O}_{X}).
\end{aligned}\right.
\tag{8}
\end{equation*}
要注意,在第一个公式中,第二个成员借助于 $D^{\infty}_{X}$ 在第二个参数 $\mathcal{O}_{X}$ 上的运算而继承了一个 $D^{\infty}_{X}$-结构;而在第二个公式中,第二个成员只是被诠释为一个 $\mathsf{C}$-向量空间层复形。这两个公式中的第二个,放在那里“以备查阅”,本质上重言式,它只是说:函子 $\delta_{\infty}$ 把一个 $D^{\infty}_{X}$-模复形 $C$ 结合到(由德拉姆函子 $\mathrm{DR}_{\infty}$ 与 $C$ 相结合的)那个微分算子(无穷阶)复形的“相伴”复形——这个复形被诠释为 $\mathsf{C}$-向量空间层复形。(人们由此确实得到一个上同调层可构的复形,这件事等价于Kashiwara的可构性定理。)

\origpage{1004}\dmark{}(相反,第一个函子 $\Delta_{\infty}$ 把可构层变成完整的(复形的)$D^{\infty}_{X}$-模,却是一条深刻的定理。仅由这一结果蕴含的有限性定理\fnote{691}{(*)}{例如,这个有限性结果蕴含:在 $X$ 上局部地,复形 $\mathrm{RHom}_{\infty}(F,\mathcal{O}_{X})$(在导出范畴中)同构于一个 $D^{\infty}_{X}$-模复形,它在每个次数上都是局部自由有限型的,并且它的上同调模(局部地)经标量扩张来源于凝聚的 D-模。事实上,甚至可以假定后者是完整且正则的。}(且不说完整性)本身就已经是一个引人注目的新结果。然而更加非凡的事情,\textbf{是这两个函子互为拟逆}。形式上,这件事使人想起双对偶关系,它既可以在范畴 $\mathrm{Cons}^{*}$ 中表达,也可以在范畴 $\mathrm{Cris}^{*}_{\infty}(X)_{\mathrm{hol}}$ 中表达——只不过(在两种情形中都表示为一个 $\mathrm{RHom}_{\infty}(-,\mathcal{O}_{X})$ 的)“对偶化”反函子把两个\textbf{不同的}范畴彼此相连。正是这种形式上的类比,使Mebkhout把断言同构
\begin{equation*}
\Delta_{\infty}\delta_{\infty}\simeq \mathrm{id}\qquad \text{在 }\mathrm{Cris}^{*}_{\infty}(X)_{\mathrm{coh}}\text{ 中}.
\tag{9}
\end{equation*}
的定理称为 $D^{\infty}_{X}$-模复形的“\textbf{双对偶定理}”(这个术语反正有引起混淆的危险)。这个关系,加上函子 $\delta_{\infty}$ 全忠实这一事实(或更确切地说,$\Delta_{\infty}$ 是它的一个伴随,这一点他把它纳入自己的双对偶定理的陈述中),Mebkhout早在1977年、在完整的好上帝定理之前就已得到。所谓的“双对偶”定理因此本质上意味着(正如它所启发的“我的”双对偶定理一样):一个\textbf{完整的} $D^{\infty}_{X}$-模复形,可以凭借对相结合的微分算子(无穷阶)复形的知识——即把该复形仅仅看成一个 $\mathsf{C}$-向量空间层复形(在适宜的导出范畴中)——把它\textbf{重构}为导出范畴的一个对象;更确切地说,可以凭借显式的反演公式(8)(第一个公式)来重构。更何况,完整 $D^{\infty}_{X}$-模复形之间的一个态射是拟同构,当且仅当无穷阶微分算子复形之间相应的态射在朴素意义下是拟同构(即在同调层上诱导一个同构)\fnote{692}{(**)}{(5月26日)事实上(如我在下文(c)开头所指出)Mebkhout证明这后一个结果时,甚至不借助任何完整性条件,其等价形式为:若与一个 $D^{\infty}_{X}$-模复形相结合的微分算子复形是拟零的,则该复形也是拟零的(D-模也是如此)。}。

% ---------------- p.776 ----------------
\dmark{}\origpage{1005}(Mebkhout的双对偶定理,在某种意义下构成好上帝定理的“一半”(就 $D^{\infty}_{X}$-模而言),当后者取其最强形式、即断言函子(8)互为拟逆时,这一点如此。)这是Mebkhout论文的核心结果,该论文于1980年1月提交。但这“一半”单独而论,已经是一个新结果,而且(据我所知)完全出乎意料。它是一个典型的结果,在Sato的观念与我的观念之间架起桥梁,但却是在我由来已久的纲领的视野中:以“连续”或“微分”的路子(并在导出范畴的视野中)来表述“离散系数”。就此而言,我觉得这个结果以其精神与灵感而论,完全逃出于日本分析学派的问题意识之外。Kashiwara的可构性定理似乎在这一点上只是一个“侧面”,而绝非一门新系数理论的出发点。正如1976至1980年间的出版物不容置疑地证明的那样,Mebkhout是当时唯一发展这样一种哲学的人。

\noindent 1978年1月,Mebkhout在巴黎过境时曾向Kashiwara讲述他的结果,当时他刚刚结束论文的撰写。应Kashiwara之请,天真的Mebkhout——他为终于找到一个似乎对他要说的东西感兴趣的人而满心欢喜——把还热乎乎的第三章寄到了普林斯顿,那一章里除别的之外就有所谓“双对偶”定理。那是1978年2月。三年之后,同一个结果(连同一个假装是证明的东西)出现在Kashiwara--Kawai的一篇著名文章中\fnote{693}{(*)}{M.~Kashiwara, T.~Kawai,\emph{On holonomic Systems of micro-differential équations}, III \emph{Systems with regular singularities},Publ. RIMS 17, 813--979 (1981). 从Mebkhout那里剽窃来的“重构定理”,见于这部长篇著作(1980年11月收到)的第4节。该著作的主要结果,是(5)中函子 $i$ 是范畴等价这一事实的一个弱化变体。因此这是Mebkhout(几何)理论的一个直接推论,是这些作者用解析的路子(独立于Mebkhout)得到的结论。详见“黑手党”子注 n\textsuperscript{o}~171(ii),(b)部分:“最初的麻烦——或太平洋对岸的卡伊德”。译按:此处171(ii)照录原书;所引标题“黑手党”实际对应171$_{2}$,疑为原书自引误。}。它为此番场合被重新命名为“reconstruction theorem”,而对某一位Zoghman~Mebkhout只字不提。那年也正是反常层会议值得纪念的一年——那光荣的一年,某种“新风格”\fnote{694}{(**)}{关于这种“新风格”(Kashiwara与Hotta是它在太平洋对岸的杰出信徒),见“祝贺——或:新风格”一注(n\textsuperscript{o}~169\textsubscript{9})。}强行征服了(且未遇丝毫抵抗……)数学的这一个部分,而在所有部分中,这里正是我当年惯常感到如鱼得水的地方……

\origpage{1006}
\subsubsection*{(c) 第五张照片(用“pro”)}

\tocsubsub{(c) 第五张照片(用“pro”)}

\dmark{}(5月21日)“双对偶定理”(9)属于1977年。为了证明 $D^{\infty}_{X}$-模“好上帝定理”的另一半——那时归结为证明函子 $\delta_{\infty}$ 本质上满射,第一个困难是证明:对 $\mathrm{Cons}^{*}$ 中的 $F$,用第一个公式(8)定义 $D^{\infty}_{X}$-模复形 $C=\Delta_{\infty}(F)$,则它至少局部地可以经由函子 $i$、借助于一个(完整、正则的)D-模复形得到。先验地,按照Mebkhout的观念(即遵循双重好上帝定理,它蕴涵(5)中的函子 $i$ 是一个等价),后者除唯一拟同构外必定唯一。

\noindent 我不曾试图理解Mebkhout最终在他的论文中是如何设法构造出这个 D-模的。在我看来,这里的情形应当用Deligne的观念来澄清,即与一个可构 $\mathsf{C}$-向量空间层 $F$ 相结合的 pro-凝聚层\fnote{695}{(*)}{这是他在1969--70年 IHÉS 讨论班上发展过、随后弃置不顾的观念。关于这一点,见“……与羁绊”子注(n\textsuperscript{o}~171(viii))。}。这个观念是他在 $X$ 上\textbf{代数簇}的语境中发展的,但必须能够 mutatis mutandis(作必要的变更)适用于解析情形,条件也许是只能在 $X$ 上“局部地”工作,或在 $X$ 的每个紧集上工作。与 $F$ 相结合的 pro-凝聚层,因此(至少在每个紧集 $K\subset X$ 上)是凝聚层的一个投射系统 $(F_{i})$(定义在 $K$ 的邻域上),它可以很简单地定义为 pro-表示以下函子的层
%JOIN:para
% ---------------- p.777 ----------------
\begin{quote}
$G\mapsto \mathrm{Hom}_{\mathsf{C}}(F,G)$
\end{quote}
(在 $X$ 上(在 $K$ 的邻域……)凝聚 $\mathcal{O}_{X}$-模 $G$ 的范畴上),这个函子正合于左,因此确实是 pro-可表示的。例如,若 $F$ 是 $X$ 的一个解析闭子空间 $Y$ 上的常值层 $\mathsf{C}_{Y}$,“以零延拓”到整个 $X$ 上,就得到由那些 $\mathcal{O}_{X_{n}}$ 组成的 pro-层,其中 $X_{n}$ 是 $Y$ 在 $X$ 中的无穷小邻域。(注意:这个投射系统的投射极限,是 $\mathcal{O}_{X}$ 沿 $Y$ 的形式完备化。)人们注意到(回到一般情形)pro-层 $(F_{i})$ 配备一个典范分层\fnote{696}{(**)}{pro-模的分层概念,与模的分层概念定义方式相同——即在前一日的注中(a)部分给出的描述,原则上每当有一个“相对”概念(如模、pro-模、相对概形等等)容许一个“拉回”概念、即给出一张以“簇”为基的“纤维范畴”时,就适用……要注意,若 $(F_{i})$ 是一个 pro-模,它的一个分层一般说来不能用一个 $F_{i}$ 的分层“相容”系统来描述——所考虑的对象,其性质远比分层模范畴的 pro-对象更为一般。}。\origpage{1007}Deligne的观念是:从 $X$ 上可构 $\mathsf{C}$-向量空间层的范畴、到分层 pro-凝聚层范畴的“\textbf{Deligne函子}”\dmark{},是\textbf{全忠实}的,从而使人能把前一个范畴(它性质上是超越的)诠释为分层 pro-凝聚层范畴的一个全子范畴。后者有一个纯代数的意义,而所说的全子范畴也能(以多少是重言式的方式\fnote{697}{(*)}{“重言式”,至少就这部已知的词典(最初由Deligne廓清)而言是如此:它在解析空间 $Y$ 中除子 $Z$ 的补集 $Y-Z$ 上局部常值的 $\mathsf{C}$-向量空间层(或“局部系统”),与 $Y-Z$ 上沿 $Z$ “正则”(Deligne意义下)的凝聚分层模之间。})用纯代数的术语来定义。这就是我将记作
\begin{quote}
(10)\quad $\mathrm{DRD}^{*}(X)$ 或 $\mathrm{Del}^{*}(X)$
\end{quote}
的范畴,它构成“\textbf{第五张照片}”,而我昨天不愿把它讲明白\fnote{698}{(**)}{最终,这一阐明(被称作“重言式”的!)在此也没有给出,至少没有当场给出。不过它将在更后面(第1011页)给出。要注意,记号(10)指的是“导出范畴”变体。}。况且我记得,Deligne当初费心把他的诠释(以及前述的全忠实性陈述)发展到能过渡到导出范畴(那时我的同调学生们、Deligne打头,还没有一致决定把这些范畴甩卖掉),而我用记号(10)所指的,当然正是“导出范畴”的版本。

\noindent 话虽如此,$\mathrm{RHom}_{\mathsf{C}}(F,\mathcal{O}_{X})$ 中的“代数部分”必须能够以一种十分自然的方式定义为一个(在适宜意义下的)归纳极限,即 $\mathrm{RHom}_{\mathcal{O}_{X}}(F_{i},\mathcal{O}_{X})$ 的归纳极限——并且特别地(过渡到上同调层),人们描述出典范箭头
\begin{quote}
(11)\quad $\varinjlim_{i}\mathrm{Ext}^{d}_{\mathcal{O}_{X}}(F_{i},\mathcal{O}_{X})\longrightarrow\mathrm{Ext}^{d}_{\mathsf{C}_{X}}(F,\mathcal{O}_{X})$ \quad ($\forall d\in\mathbb{Z}$)
\end{quote}
利用 pro-对象 $(F_{i})$ 上的分层与第二个参数 $\mathcal{O}_{X}$ 的重言式分层\dmark{},人们必须能够在(11)的第一个成员上定义一个分层、即一个 D-模结构,使得(11)与(相应的算子环的同态\origpage{1008}$D\to D^{\infty}_{X}$)相容。话虽如此,Mebkhout的好上帝定理必须能够精确化,即说(11)把第二个成员等同于由第一个成员经标量扩张而得到的 $D^{\infty}_{X}$-模\fnote{699}{(*)}{此外,当然,(11)的第一个成员(与Mebkhout的哲学一致)必须是一个\textbf{凝聚}、\textbf{完整}且正则的 D-模。}——这特别蕴含该箭头是一个\textbf{包含}。因此,左边一个成员应当被看成是右边一个成员中的某种\textbf{“代数部分”}(或“\textbf{亚纯部分}”),而后者性质上是
“超越的”。

\noindent 一般情形由前面这个特例 $F=i_{*}(\mathsf{C}_{Y})$(其中 $i : Y\to X$ 是 $X$ 的一个解析闭子空间的包含映射)而得到大量澄清。于是(11)的第二个成员是一个支撑在 $Y$ 中的局部上同调层——一个\textbf{超越的}不变量,而第一个成员
\begin{quote}
$\varinjlim_{n}\mathrm{Ext}^{d}_{\mathcal{O}_{X}}(\mathcal{O}_{X_{n}},\mathcal{O}_{X})$,
\end{quote}
却是我当年在概形框架中为局部上同调引入的那个众所周知的表达式。这个层在一点 $x\in Y$ 上的纤维,不是别的,正是结构层在 $Y$ 于 $X_{x}$ 上的“迹” $Y_{x}$ 中的支撑上、在 $\mathcal{O}_{X,x}$ 的谱 $X_{x}$ 上的局部上同调。

\noindent 从这个例子可以看出,Deligne的观念与我六十年代初就局部上同调主题所发展的那些观念何等接近\fnote{700}{(**)}{下文将显出,Deligne的观念也与我在1966年(在[Crystals]中)引入的观念密切相关:对任何微分算子复形,我把它的“形式化” $P^{\infty}(L^{\bullet})$ 看成一个分层 pro-模复形,或者更好地说,看成一个\textbf{晶体复形}的定义,其(整体)晶体上同调等同于 $L^{\bullet}$ 的(整体)上同调。}。无论如何,Mebkhout1972至1976年间工作的主题,正是在这个关键情形下研究箭头(11)
\begin{equation*}
\varinjlim_{n}\mathrm{Ext}^{d}_{\mathcal{O}_{X}}(\mathcal{O}_{X_{n}},\mathcal{O}_{X})=H^{d}_{Y}(\mathcal{O}_{X})^{\mathrm{alg}}\longrightarrow H^{d}_{Y}(\mathcal{O}_{X})
\tag{12}
\end{equation*}
他在这种情形下证明了上文宣告的关系,此外还证明了(这是我方才一度忘记纳入陈述的)(12)的第一个成员是一个\textbf{凝聚的} D-模,乃至完整且正则。\origpage{1009}由此起,(11)的类似陈述必须是经 dévissage(拆卸)而来的直接结果\fnote{701}{(***)}{(5月22日)我在这里有点急躁!可构 $\mathsf{C}$-层的“基本类型”比单纯的 $\mathsf{C}_{Y}$ 性质更为一般,(但一般定理的证明确实用了与1976年那个特例相同的技术。)},\dmark{}包括在 $F$ 不是可构 $\mathsf{C}$-向量空间层、而是 $\mathrm{Cons}^{*}(X,\mathbb{C})$ 中一个复形的情形下也是如此。唯一的一粒盐,除了Deligne函子形式的构造之外,是在一个分层 pro-模复形、取值于一个分层模复形(即取值于一个 D-模复形,此处是 $\mathcal{O}_{X}$)的 $\mathrm{RHom}_{\mathcal{O}_{X}}$ 的定义中,是作为一个 D-模复形(且作为导出范畴的一个对象)来取的。

\noindent 且不论这一粒盐,人们于是得到好上帝函子 $M$“代数的”(相对于好上帝函子 $M_{\infty}$“超越的”而言)、或者不如说相联系的反函子 $\Delta$ 及其拟逆 $\delta$
\begin{quote}
(13)\quad $\Delta=MD=DM$ ,\quad $\delta=mD=Dm$,
\end{quote}
的一个至为简单而概念化的描述,用的是一条双重公式,它把(8)复述一遍。但为了把它写出来,要利用Deligne的等价
\begin{equation*}
\mathrm{Del} : \mathrm{Cons}^{*}(X,\mathbb{C})\xrightarrow{\ \simeq\ }\mathrm{DRD}^{*}(X),
\tag{14}
\end{equation*}
人们宁愿看相应的函子 $\hat{\Delta}$、$\hat{\delta}$,它们在 $\mathrm{DRD}^{*}(X)$ 与 $\mathrm{DRM}^{*}(X)$ 之间,其中的帽子符号意在提醒:人们将在(“可构”那一侧)用 \textbf{pro-对象}来工作。于是人们得到那些卓越的公式(实质上包含在(8)中,但这一次联系的是“代数性质”的系数双方,并且用的也是“代数性质”的公式):
\begin{quote}
(15)\quad $\hat{\Delta}(C')=\mathrm{RHom}_{\mathcal{O}_{X}}(C',\mathcal{O}_{X})$ ,\quad $\hat{\delta}(C)=\mathrm{RHom}_{\mathcal{O}_{X}}(C,\mathcal{O}_{X})$。
\end{quote}
这里因此两次出现“同一个”公式,唯一的区别在于 $C'$ 此处是一个分层 pro-凝聚层复形(或者说\fnote{702}{(*)}{关于这一译法,见第1006页脚注(**)。},一个 pro-凝聚模的晶体复形),而 $C$ 是一个 D-模复形(本质上可以把它看成一个分层 ind-凝聚 $\mathcal{O}_{X}$-模复形,或者再看成一个 ind-凝聚模的晶体)。本质上正是“同一个”函子使人们从一类走到另一类,即“寻常的对偶化”(凝聚的)函子,我五十年代的老友……诚然,“显然”它必须交换 pro-对象与 ind-对象(不妨在后者中取归纳极限……)。

\origpage{1010}\dmark{}(当然,有一项基础工作要做,为的是给这些公式一个精确的意义——一项Deligne在他那个著名的沉没讨论班、或儒瓦诺卢在他那篇同样著名的沉没论文中所做的那类工作(自从反常层会议以来人人都引它,却没有一个人把它捧在手里……)。这是一项我确信或许会有点长、但本质上“连环推理”的工作。“硬”的部分包含在Mebkhout的好上帝定理中,由Mebkhout的公式(8)(也许不恰当地)称作“双对偶”公式者补全。而它们的代数翻译,即断言两个函子(15)互为拟逆,却实实在在是(实质上)凝聚 $\mathcal{O}_{X}$-模系数的“那个”寻常双对偶定理,只是佐以 ind-pro 的风味、并以分层为关键罢了(它们必须在对偶化函子中“通行”无阻)。

% ---------------- p.779 ----------------
\noindent 两类对偶对象之间的对应,用微分算子复形来观照则完美无遗(且不须任何基础工作作关键!)。(在这场对偶中,况且完整性条件(进而正则性条件)不扮演任何角色。)对这样一个复形 $L^{\bullet}$,把昨天(在(a)、(1)中)考虑的那个(反变的)函子
\begin{quote}
$F\mapsto \mathrm{Hom}_{\mathcal{O}_{X}}(F,\mathcal{D}_{d})$
\end{quote}
结合到一个分量局部自由有限型的 D-模复形,即 $C$。另一方面,把这个复形 $L^{\bullet}$“形式化”,过渡到无穷阶主部 $P^{\infty}(L^{i})$(被看作分层的 pro-模),就提供一个分层 pro-模复形 $C'=P^{\infty}(L^{\bullet})$。话虽如此,人们看到这两个复形由公式(15)相互对应,在其中此处显然 $\mathrm{RHom}$ 归结为 $\mathrm{Hom}$。(只须对各分量 $L^{i}$ 逐项验证这一对偶性,它就归结为下述多少是重言式的事实:从 $P^{\infty}(L^{i})$ 到 $\mathcal{O}_{X}$ 的“连续”线性同态,与从 $L^{i}$ 到 $D$ 的线性同态一样,恰好对应于从 $L^{i}$ 到 $\mathcal{O}_{X}$ 的微分算子,分别利用“万有”的(无穷阶)微分算子 $L^{i}\to P^{\infty}(L^{i})$,与由 $\theta\mapsto\theta(1)$ 给出的“增广” $D\to\mathcal{O}_{X}$。)由于至少局部地在 $X$ 上,$\mathrm{Cris}^{*}_{\mathrm{coh}}(X)$ 的每个对象(即每个上同调凝聚的 D-模复形)都可以借助一个微分算子复形 $L^{\bullet}$ 来描述,因此可以认为,只要对 $C$ 与 $C'$ 作适宜的 D-凝聚性与“D-pro-凝聚性”假设、彼此“对偶”,这个特例就在一切实际用途上完美地掌握了两类系数之间的对偶性(15)。\origpage{1011}于是只须发展我所提到的“连环推理”,在 $C'$ 或“pro”一侧局限于分层 pro-凝聚层复形\dmark{}(它们局部地可以(除拟同构外)描述为一个 $P^{\infty}(L^{\bullet})$)。

\noindent 相对于Deligne原本的进路,他所引入的 pro-凝聚模及其复形能够局部地用一个微分算子复形来实现,这件事本身是一个\textbf{全然出乎意料的现象},由Mebkhout的理论所带来。在我看来它本质上是
等同于\fnote{703}{(*)}{(5月26日)在这里我又“有点急躁”了,1976年的那个结果并不足够。请与第1008页脚注(***)的评注相比较。}上面提到的那个Mebkhout定理(可追溯至1976年,早在好上帝定理的证明之前),该定理关乎出现在上方(12)中的层 $H^{d}_{Y}(\mathcal{O}_{X})^{\mathrm{alg}}$ 的 D-凝聚性。那是一条深刻的定理,是四年工作的结晶,用尽了广中奇点消解的全部力量(更不消说那位廓清并证明了它的工匠的勇气——顶着普遍的冷漠)。我刚才指出的那个后果\textsuperscript{703(*)},是德拉姆系数(如我自1966年起所窥见的那样)与微分算子复形之间一层深刻的关系,这层关系我全然未曾预见(Deligne发展他通往德拉姆系数的第一条进路时,也同样未曾预见)。至于所考虑的微分算子复形上的完整性与正则性条件,它(事后而论,多亏那救场的好上帝定理)应当等同于Deligne的“有限性”(加“正则性”)条件(我在引入范畴 $\mathrm{DRD}^{*}(X)=\mathrm{Del}^{*}(X)$ 时一度未曾把它讲明)。条件如下:$P^{\infty}(L^{\bullet})$ 的上同调 pro-层在局部上由合成列“拆卸”,使得相继因子能(经由Deligne函子)由 $X$ 的子空间 $Y-Z$ 上的 $\mathsf{C}$-向量空间局部系统来描述(其中 $Z\subset Y\subset X$ 是 $X$ 的解析闭子空间)。为了最终给这条判据一副“代数”的面貌,只需把 $\mathsf{C}$-向量空间局部系统换成 $Y-Z$ 上一个\textbf{凝聚}分层,并附加条件:表达该分化的联络(注意可设 $Y-Z$ 光滑)在 $Z$ 的邻域上是“正则的”,取Deligne的意义\fnote{704}{(**)}{考虑到Deligne所廓清的范畴等价,即在 $Y-Z$ 上的 $\mathsf{C}$-向量空间局部系统、与 $Y-Z$ 上配备沿 $Z$ 的“亚纯结构”、且沿 $Z$ 联络正则的可积联络丛之间的等价,这条正则性条件在此以一种自然的方式引入。这种亚纯结构(蕴含把 $Y-Z$ 上的凝聚模延拓为 $Y$ 上的凝聚模的可能性,至少在 $Z$ 的每点邻域上如此)在方才给出的描述中是暗含的。除非我弄错,当在前一条条件中去掉正则性条件时(只需假定在 $Z$ 的邻域上给定 $E$ 的一个亚纯结构,以便能借Deligne的办法给它结合一个整个 $X$ 上的 pro-凝聚模),人们就得到完整性条件的一个“上同调式”描述。Sato的定义是借“微局部”的路子作出的——我承认,我还从未真正领会过它……}。(注意:相联系的 pro-层可以这样得到:\origpage{1012}把 $Y-Z=T$ 上拥有的\dmark{}晶体在 $T$ 的无穷小邻域上方养大,并沿 $Z$“压碎”,以便处处都有凝聚的层,而不只在 $Z$ 的补集上……)

\subsubsection*{(d) 晶体与余晶体——全忠实?}

\tocsubsub{(d) 晶体与余晶体——全忠实?}

\noindent 当不再假定 $X$ 光滑时,因此,除了性质上超越的“照片” $\mathrm{Cons}^{*}(X,\mathbb{C})$ 之外,还剩下两张“照片”(两者性质上都是晶体的)$\mathrm{DRM}^{*}(X)$ 与 $\mathrm{Del}^{*}(X)$,而它们却具有纯代数的意义。我昨天(在(a)中)为 $\mathrm{DRM}^{*}(X)$ 草拟了一条定义原则,今天又为范畴 $\mathrm{DRD}^{*}(X)$ 拟了一条。后者从眼前起对我来说就完全可理解。如我昨天所指出(见(a),第998页脚注(***)),这里应当把分层 pro-模的观点,用 pro-模(pro-凝聚的)\textbf{晶体}的观点来细化\fnote{705}{(*)}{(5月27日)细想起来,我甚至难以相信Deligne的定理 $\mathrm{Cons}^{*}(X,\mathbb{C})\cong\mathrm{Del}^{*}(X)$ 对非光滑的 $X$ 也成立,如果 $\mathrm{Del}^{*}(X)$ 像Deligne所做的那样、不借助晶体site来定义的话。也许正是为了意识到这一点,他最终宁愿把整套理论沉没掉,而不肯同意重新引入那个犯忌的位……(请与“……与羁绊”一注(n\textsuperscript{o}~171(viii))相比较。)}。这个观点眼下仅存的那点麻烦,就是它将要迫使人们发展的那套“pro”连环推理,而它(据我在此类事情上微薄的经验)有变成骇人规模的风险!这些 pro-模的晶体,给 $X$ 的每个开集 $U$ 的每个无穷小加厚 $U'$ 结合一个 $U'$ 上的 pro-凝聚模,“以与拉回像相容的方式”对应于加厚之间的态射 $U''\to U'$,它们甚至不能被诠释为晶体site(或同样地说,晶体
%JOIN:para
% ---------------- p.781 ----------------
topos $X_{\mathrm{cris}}$)上的 pro-层!因此人们先验地不能把诸如 $X_{\mathrm{cris}}$ 这样的(交换地)带环topos上模层的已知上同调形式主义应用到它们身上。

\noindent 这里的诱惑很大,就是取人们在每个加厚上拥有的那个 pro-层的投射极限。这样得到的是晶体的模(而非模的晶体),其“值”在每个 $U'$ 上既非凝聚也非拟凝聚。希望在于,至少对\origpage{1013}我们感兴趣的那类 pro-模晶体(尤其是那些由Deligne函子得到的)而言,这样一个 pro-模晶体能够从经过取极限而得到的晶体模 $C$\textbf{重构}出来,办法是在每个加厚 $U'$ 上取 Zariskian 层 $C_{U'}$($C$ 在 $U'$ 的 Zariski 开集上的限制)的“pro-凝聚包络”\fnote{706}{(*)}{我在这里谈到“Zariskian 层”(与“晶体的”相对)时,已悄悄滑回了概形的语境。偏爱解析语境的读者自会加以校正。}。在我看来这确实如此,至少对如上所说联系于 $Y-Z$ 上一个凝聚分层模的 pro-模晶体是这样,例如在取 $\mathcal{O}_{X}$ 沿 $Y-Z$ 的形式完备化、并在别处以零延拓(在加厚上亦然)的那个典型情形中就是这样。

\noindent 如果我的“希望”得到证明,那么 $X$ 上德拉姆-Deligne系数的范畴 $\mathrm{DRD}^{*}(X)$ 就可以诠释为寻常导出范畴 $D^{*}(X_{\mathrm{cris}},\mathcal{O}_{X_{\mathrm{cris}}})$ 的一个全子范畴,以上同调层上“有限性”与“正则性”类型的条件(它们本身如上所述以拆卸来描述)来定义。那将是一个令人困惑的简单描述,我早在1966年就可以把它给出,如果当时我有暇继续我的晶体思考的话……

\noindent 这个“基础的”问题(取极限是否合法)显然不依赖于 $X$ 是否光滑的问题——如果它不光滑,就把它浸没在一个光滑的 $X'$ 中而归结到光滑情形。如果这个观点(简直美得不像真的!)确实行得通,那么(现在在光滑情形下)顺便就该(我认为)把“双对偶”公式(代数版本)(15)诠释为 $\mathrm{RHom}_{\mathcal{O}_{X}}$ 的\textbf{寻常}公式,不必为 pro-问题费神(而只需小心搬运分层……)。这方面的第一个检验如下:若 $u : C_{1}\to C_{2}$ 是上同调凝聚的 D-模复形之间的一个态射,它经由朴素的对偶化函子 $\mathrm{RHom}_{\mathcal{O}_{X}}(-,\mathcal{O}_{X})$ 的像是一个拟同构,那么 $u$ 是否也是拟同构?但这一点(借一个映射柱的论证)归结为问:一个上同调凝聚的 D-模复形,若其“朴素对偶”为零(取导出范畴的意义,即上同调层为零),它本身是否也为零(取同样的意义)?或者再说一遍,如果有微分算子复形 $L^{\bullet}$,说相结合的 D-模复形上同调层为零,与说“形式化”的 $P^{\infty}(L^{\bullet})$(这一次不看作一个 pro-层复形,而看作一个寻常层复形(经取 $\varprojlim$))上同调层为零,是否一回事?Mebkhout肯定能告诉我……

\dmark{}\origpage{1014}(5月23日)昨晚我又给Mebkhout打了电话——况且数学的或历史的问题,我几乎每晚都给他打电话,已经有一两个星期了,加起来可要创下一笔天文数字的电话账单了!但《封神》,我这三个紧巴巴的星期以来孜孜不倦而又擦拭不休的那篇,值这个价……

\noindent 无论如何,Zoghman向我保证了一个结果,它看上去同我昨夜收尾时写到的那个“检验问题”相近:若 $C$ 属于 $\mathrm{Cris}^{*}_{\mathrm{coh}}$ 且使得相结合的算子复形 $L^{\bullet}=\mathrm{DR}(C)$ 是拟零的,则 $C$ 本身是拟零的(解析情形)。人们有一个层(的 $\mathsf{C}$-向量空间)复形之间的同态,由“无穷阶主部”给出
%JOIN:para
% ---------------- p.782 ----------------
\begin{quote}
$L^{\bullet}\longrightarrow P^{\infty}(L^{\bullet})$,
\end{quote}
由此有上同调层上的同态
\begin{quote}
(16)\quad $H^{i}(L^{\bullet})\longrightarrow H^{i}(P^{\infty}(L^{\bullet}))$\quad ($i\in\mathbb{Z}$)
\end{quote}
人们很想说这个同态(16)总是单的,并把第一个成员等同于第二个的“水平”截面子层(这将是某种“水平截面层”函子在适宜的分层 pro-模范畴上的正合性……)。单性便将蕴含:若第二个成员为零,则第一个成员也为零,因此若这一切对每个 $i$ 都成立(且据Mebkhout向我保证的那样),那么与 $L^{\bullet}$ 相结合的 D-模复形是拟零的——这正是我想要的。

\noindent (16)中的单性也意味着:对一个微分算子 $E\overset{d}{\longrightarrow}F$,以及 $F$ 的一个截面 $f$,若 $f$ 在每点 $x\in X$ 上“形式上”属于像中(过渡到该点的完备局部环),此外“形式解”(即方程 $d(g)=f$ 在 $g$ 中的解)可以取成对变动的 $x$ 解析地依赖,那么该方程局部地有解。Mebkhout告诉我,他不知道有这样的结果;然而这个问题如此自然,答案想必应该是已知的!

\noindent 为了结束“五张照片”,我想在此再回到那两张“晶体照片”上,一张对应于Mebkhout的 D-模观点,另一张对应于与之对偶的观点。不言而喻,应当在导出范畴的精神中工作——因此一种名副其实的“晶体”诠释必须顾及这一点。所以那两张晶体照片只有在\origpage{1015}那相应的函子\dmark{}(从范畴 $D^{b}_{\mathrm{coh}}(X,D)$(姑且说)出发、到一个适宜的晶体范畴、例如 $D^{b}(X_{\mathrm{cris}},\mathcal{O}_{X_{\mathrm{cris}}})$)本身是全忠实的时候,才是“全忠实”的。我希望情形正是如此,\textbf{甚至不必}给所考虑的 D-模复形\textbf{添上完整性与正则性条件}。

\noindent 最简单的情形无疑是第四张照片,它在于把 D-模的范畴诠释为模的晶体的范畴,由此有一个总体导出函子(不妨说是“Grothendieck的”——以抢在那些“无用细节”与“技术性离题”的爱好者之前……):
\begin{quote}
(17)\quad $G : D^{*}_{\mathrm{coh}}(X,D)\longrightarrow D^{*}(X_{\mathrm{cris}},\mathcal{O}_{X_{\mathrm{cris}}})$。
\end{quote}
这里关键的问题,是这个函子是否全忠实。只有在此时,记号 $\mathrm{Cris}^{*}_{\mathrm{coh}}(X)$ 用于第一个成员才是完全有道理的——并且同时,德拉姆上同调(至少在此情形中,在复解析框架或特征零的域上代数概形的框架中)的晶体观点也才有道理。为了证明全忠实性,比如说在代数几何中,经标准论证可归结到 $X$ 是仿射的情形(或在解析情形中,归结到一个多圆盘的情形),以及归结到第一个成员中所考虑的两个对象 $C$、$C'$(问题在于比较它们在一个方向与另一个方向上的 $\mathrm{Hom}$)都等于 $D$ 本身、只差一个次数平移的情形。(这个归结毫无问题地作出,至少假定 $C$、$C'$ 次数有界,从而局限于 $D^{b}_{\mathrm{coh}}(X,D)$,而这对应用而言看来是大为足够的。)于是最终被引导去验证公式
\begin{quote}
(18)\quad $\Gamma(X,D_{X})\xrightarrow{\ \sim\ }\mathrm{Hom}(G(D),G(D))$ ,\quad $\mathrm{Ext}^{i}_{\mathcal{O}_{X_{\mathrm{cris}}}}(X_{\mathrm{cris}};G(D),G(D))=0$\quad 对 $i>0$
\end{quote}
($X$ 仿射,相应地 Stein)。我不曾有暇去验证它\fnote{707}{(*)}{我有借口:一年多以来,我最清醒的时间都用来追踪我昔日学生中某几位的业绩了……},但毫不怀疑它是真的。某个非常接近的东西,我似乎在[Crystals](1966年)中已经证明过了\fnote{708}{(**)}{这是我已在别处提到过的结果:对概形 $X$ 上的一个微分算子复形 $L^{\bullet}$,(在 $X$ 对 $S$ 光滑相对的框架中,或在解析框架中无疑也同样) $L^{\bullet}$ 的“Zariskian”超上同调等同于它的形式化 $P^{\infty}(L^{\bullet})$ 的晶体超上同调。诚然,这个陈述更直接地关乎(17)的“对偶”箭头(19),也可以这样表达:对 $C$、$C'$ 上同调凝聚的 D-模复形,箭头 $\mathrm{Hom}(C,C')\to\mathrm{Hom}(G^{\circ}(C),G^{\circ}(C'))$ 在 $C=\mathcal{O}_{X}$ 的情形下是双射(这已经不错,且容许一切希望……)。}。

% ---------------- p.783 ----------------
\dmark{}\origpage{1016}(至于第五张照片,它有几种不同的印片。Deligne原初的印片是用分层 pro-凝聚模作出的。为了推广到 $X$ 非光滑的情形,第一次重要的修版,在于把所说的那些动物诠释为 pro-模的\textbf{晶体}。但那样就卷入了(毫无吸引力的!)无休止的 pro-上同调代数 pro-基础的齿轮——并且失去直接依附于 $X_{\mathrm{cris}}$ 的topos式直觉的好处。所以我宁愿(如果可能的话)干脆另取一张照片,视角几乎相同,经由一个\textbf{反变}函子(同样不妨说是“Grothendieck的”,识者自会领会……)
\begin{quote}
(19)\quad $G^{\circ} : D^{*}_{\mathrm{coh}}(X,D)^{\mathrm{opp}}\longrightarrow D^{*}(X_{\mathrm{cris}},\mathcal{O}_{X_{\mathrm{cris}}})$。
\end{quote}
可以说,这是从Deligne的照片出发、在 $X$ 的每个开集 $U$ 的每个无穷小加厚上粗暴地取投射极限层而得到的那一个。若第一个成员中的 $C$(以反变的方式,如(a)中公式(1)那样)结合到一个微分算子复形 $L^{\bullet}$,它经(19)的像可以这样得到:把 $P^{\infty}(L^{\bullet})$(复形 $L^{\bullet}$ 的“形式化”)看作一个分层 pro-模复形(在[Crystal]中引入的观念),或者再看成一个 pro-模的晶体复形,并在任何加厚上取投射极限。另一种说法是:对 $X$ 上任何局部自由的 $\mathcal{O}_{X}$-模(例如)$L$,结合一个晶体模(它\textbf{不是}一个模的晶体,除非我弄错),我把它记作 $P^{\infty}(L)_{\mathrm{cris}}$,其方式固然是“显然的”(而我的学生们早已把它忘掉了),该模以微分算子为中介函子性地依赖于 $L$,从而过渡到微分算子复形。

\noindent 前述对函子(19)的两种描述都还是不完整的,尤其是因为第一个成员的一个对象未必在整个 $X$ 上来自一个微分算子复形。我推测可以给这一启发性描述一个内在诠释,用公式
\begin{quote}
(20)\quad $G^{\circ}(C)\xrightarrow{\ \sim\ }\mathrm{RHom}_{\mathcal{O}_{X_{\mathrm{cris}}}}(G(C),\mathcal{O}_{X_{\mathrm{cris}}})$\quad (其中 $G$ 定义于(17))
\end{quote}
但不曾验证它正确。经标准论证,在这里人们又归结到(为了证明当 $C$ 如上所说结合到 $L^{\bullet}$ 时自然箭头(20)确实是一个同构)$C=D$ 的情形,而那时(20)归结为公式\dmark{}\origpage{1017}
\begin{quote}
(21)\quad $\mathrm{Ext}^{i}_{\mathcal{O}_{X_{\mathrm{cris}}}}(X_{\mathrm{cris}};G(D),\mathcal{O}_{X_{\mathrm{cris}}})=0$\quad 对 $i>0$,
\end{quote}
它们与(18)相当神似。

\noindent (19)的全忠实性之意,无论如何是相当清楚的,并且(如(17)那样)又经拆卸归结到 $C=D$、$C'=D[i]$(次数平移 $i$)的情形,于是归结为公式
\begin{quote}
(18.1)\quad $\Gamma(X,D)\simeq \mathrm{Hom}(P,P)$ ,\quad $\mathrm{Ext}^{i}_{\mathcal{O}_{X_{\mathrm{cris}}}}(X_{\mathrm{cris}};P,P)=0$\quad 对 $i>0$,
\end{quote}
其中记
\begin{quote}
$P=P^{\infty}(\mathcal{O}_{X})_{\mathrm{cris}}$,
\end{quote}
%JOIN:para
% ---------------- p.784 ----------------
它是 $X$ 上一个无比卓越的晶体代数。这里(就晶体 $\mathrm{Ext}^{i}$ 的零性而言)假定 $X$ 是仿射的(相应地 Stein)。

\noindent 最终,昨天在我看来还“简直美得不像真的”、而那时我还在透过Deligne的照片看事物的那个东西,一旦把事情写成不必添上完整性条件(更不必说正则性条件),就突然显出一切合情合理的模样。若上帝假我以寿命,且没有别人抢先替我做这工作,我希望赶在年底以前把那件事弄清楚(连同(21)与(18)的有效性),借《数学反思》第三卷中专门讨论德拉姆系数的那一部分。

\noindent 如我所说,是第五张照片——那张与离散系数相联系的拓扑学直觉“贴合”得最紧的——得我偏爱。若是有一天我心灰意冷地得知公式(22)是假的(而若是公式(18)如此,我还不至于那么烦恼,虽说它们技术上看着更没把握),那将说明必须回到(修过版的Deligne照片的)pro-观点——一个不怎么乐观的前景!无论如何,我毫不怀疑,除技术性调整之外,人们手里确实握着一张出色的照片,它在代数几何中尤其有效(甚至在异于特征零之域的基上),且不加任何光滑性假设。

\noindent 至于第四张照片,其忠实性取决于(18)的有效性,我再度承认,在光滑情形之外(甚至在光滑情形之中)“我还看不明白”它,且不能确定对非光滑的 $X$,我所提出的晶体诠释是否原样行得通。不过我觉得,我那些关于 D-模观点(尤其是我对这一观点的晶体诠\dmark{}\origpage{1018}释)的根深蒂固的变差困惑,正借助于引入一个与晶体概念相对偶的概念——我称之为\textbf{余晶体}(co-cristal)——而即将解决。就在不久前的昨天,那种(就闭浸入下 D-模的“变差”而言)一直存在的不适的弥散感觉,终于分娩出一个“好概念”(据我所知,还不曾真正写下什么)。它在“ind”一侧看着很合身,犹如晶体概念(为我所熟悉)在“pro”一侧那样。在光滑簇上,两个范畴(晶体与余晶体)典范等价(而这正是我必然倾向于把它们混同的原因——情有可原……),但对任意的 $X$ 就不再如此。局面与上同调环 $H^{\bullet}(X)$ 和同调群 $H_{\bullet}(X)$、Chow 环 $\mathrm{Ch}^{\bullet}(X)$ 和 Chow 群 $\mathrm{Ch}_{\bullet}(X)$、或Grothendieck环(请原谅我这处不规整……)$K^{\bullet}(X)$ 和Grothendieck群 $K_{\bullet}(X)$(再次恕罪)之间发生的事完全类似。在那里,当 $X$ 是(拓扑的、或代数的等等——视情形而定)光滑簇时,人们也曾长期混淆这两类对象。事后“可以解释”:第二项无论如何配备了一个第一项上的模结构(cap-积——在后两种情形中它由一位我不敢在此指名的祖先引入),而在光滑情形中,人们发现这个模是秩 1 自由的且配备一个典范基,这使它不幸与那个环(无疑美得多)混同。

\noindent 那么,对于 $X$ 上模的晶体的范畴 $\mathrm{Cris}^{\bullet}(X)$(以张量积配备“环”结构)与模的余晶体的范畴 $\mathrm{Cris}_{\bullet}(X)$(前者以一个 cap-积完美地在其上“作用”)之间,情形也是这样。

\noindent 但现在是停下这一长段数学离题的时候了,我承认它在一场漂亮的葬礼的编排中全然不合时宜。有兴趣知道下文(自然纷繁)的读者,只好去买《数学反思》第三卷(若他不吝惜钱财),在那里一位不知悔改的逝者打算继续他那些混乱的“技术性离题”\fnote{709}{(*)}{这一次,不言自明,是作为我另一位学生的“合作者”,而这位学生早已高升为“晶体”之父……}。

% ---------------- p.785 ----------------
\subsubsection*{(e) 好上帝的无处不在}

\tocsubsub{(e) 好上帝的无处不在}

\noindent (5月27日)一篇“最后的”脚注,昨天(在把《封神》头十二篇注交给打字之前)临急之际补进“五张照片”,结果又涨到“骇人的规模”,而我\dmark{}\origpage{1019}最终要把“这一长段数学离题”用最后一个(简短的)章节继续下去。这样一来,《五张照片》就将由(a)至(e)\textbf{五}个章节组成——可见一切都会圆全和完善……

\noindent 这是对Mebkhout“好上帝定理”(设想的)真正有效性域的一番评论,它(据我看来)远远超出复解析空间这一初始框架——不仅靠它带来的新\textbf{哲学}(而这哲学眼下已经刷新了上同调主题),而且也在一个技术意义上超出。

\noindent 一旦把 $X$(光滑)上可构的 $\mathsf{C}$-向量空间层,或者诠释为分层 pro-凝聚模(Deligne式),或者(通过在 $X$ 的开集的无穷小加厚上取投射极限)诠释为晶体层(Grothendieck式),那么“好上帝定理”(别名Mebkhout)就断言两个范畴的等价,而这两个范畴这一次\textbf{两者都是}“纯代数”性质的。换句话说,这定理现在有一个精确的意义,是在复解析语境以外的其他语境中:既在体上光滑概形的语境中(甚至无须假定特征零——关于这一点见上面第996页脚注(**);在特征 $p>0$ 中“除幂晶体”的观点在此是本质的),也在任意特征的刚性解析簇的语境中,又在 $\mathbb{Z}$ 上有限型光滑概形的语境中(此外还有许多……)。

\noindent 好上帝定理的“形式”部分关乎\textbf{所有}凝聚的 D-模复形,而不只是那些完整的,它说:好上帝函子经祖先修订校正之后(即本质上相对于结构层 $\mathcal{O}_{X}$ 的对偶)是范畴 $D_{\mathrm{coh}}(X,D_{X})=\mathrm{Cris}^{*}_{\mathrm{coh}}(X)$ 到所考虑的系数范畴 $\mathrm{Coeff}^{*}$(随感兴趣者的口味任选)的一个\textbf{全忠实}函子。把事情看透了,这应当或多或少是“连环推理”的。

\noindent 但在到达范畴中,人们“以拆卸”定义两个值得注意的全子范畴,即“完整系数”的与“完整正则系数”的(如(c)末尾以及第1011页脚注(**)中那样)。话虽如此,“Mebkhout广义定理”(就所考虑的语境而言)肯定毫无连环推理之嫌,但必定深刻,它将说两件事:
\begin{quote}
1.~\dmark{}\origpage{1020}完整“系数”的范畴 $\mathrm{Coeff}^{*}_{\mathrm{hol}}$ 在范畴 $\mathrm{Cris}^{*}_{\mathrm{coh}}(X)$ 经“Mebkhout-Grothendieck”(全忠实)函子的像之中。(注意:实质上,这个函子就是Mebkhout函子,但是着眼于 $\mathrm{Cris}^{*}_{\mathrm{coh}}(X)$ 全体看的,并且还经“祖先修订校正”,使目的落在 $\mathrm{Coeff}^{*}$ 中,而后者有纯代数的意义……)

2.~以“微局部”的“完整性”与“正则性”条件,用微分算子复形的术语,来刻画 $\mathrm{Coeff}^{*}_{\mathrm{hol}}$ 与 $\mathrm{Coeff}^{*}_{\mathrm{hol\,reg}}$ 的逆像。
\end{quote}
关于最后这一点(它对于我六十年代的纲领或许相对次要),在特征零中人们已经有一个现成的完整性条件。至于正则性条件,现在正好该看看日本人手里是否恰好有那个恰当的概念——但这可不是Mebkhout能告诉我的,因为他见得太多,不想再听人提起。

\noindent 至于我,我没有他那样的见识,我觉得正则性有\textbf{三个不同的侧面},它们相互补充:
\begin{quote}
1.~Deligne以拆卸在 $\mathrm{Coeff}^{*}_{\mathrm{hol}}$ 中廓清的“几何”侧面,归结为一个“局部系统”(例如带可积联络的丛)在一个除子邻域上的正则性条件。

2.~“微局部”或“日本式”的侧面,直接用微分算子复形的术语来表述(?)

3.~Mebkhout引入的“上同调式”侧面,这个侧面目前(据我看)只在复解析情形中被很好理解。我丝毫不知道它有没有机会推广到刚性解析情形。
\end{quote}
侧面 $3^{\circ}$ 当然将是关键的,每当需要为定义在一个完备赋值体上的代数簇、以完整系数,在“Zariskian”上同调与“刚性”上同调之间建立一个\textbf{比较定理}的时候。

\noindent 对于我六十年代宏大的“变差纲领”,最重要的恰恰是“几何”侧面。要紧的是为 $\mathrm{Coeff}^{*}_{\mathrm{hol\,reg}}$ 定义一套六种运算的形式主义。如果甚至能为 $\mathrm{Coeff}^{*}_{\mathrm{hol}}$ 找到一套,如Mebkhout似乎所信,那就更好。

\noindent 但是(如果我没弄错)Motif(我对此早有预感\dmark{}\origpage{1021}超过别的一切)只会生出同时完整且正则的系数来。

\noindent 我回到问题1,它有一个显而易见的变体,即一个“问题 $1'$”(更为谦虚),把 $\mathrm{Coeff}^{*}_{\mathrm{hol}}$ 换成 $\mathrm{Coeff}^{*}_{\mathrm{hol\,reg}}$。一旦证明了Mebkhout-Grothendieck函子的全忠实性,人们显然就归结到下面这件事:在 $X$ 的一个光滑(未必闭的)子簇 $Y$ 上,给定一个带可积联络的丛(或一个 $\mathsf{C}$-凝聚的 $F$-晶体 $C$,视所选的语境而定……),必要时给它附加一条Deligne的补充正则性条件(在 $\overline{Y}-Y$ 的各点上)。Deligne的办法(或由祖先修订以便过渡到晶体语境)使我们能给它们结合 $\mathrm{Coeff}^{*}$ 的一个对象(按定义它甚至会是“完整的”,乃至“完整正则的”)。这个对象在Mebkhout-Grothendieck函子的像中吗?或者,同样地说,在 $X$ 上局部地,所说的这个 $\mathrm{Coeff}^{*}$ 对象能用 $X$ 上的一个微分算子复形来描述吗,借祖先专利的办法,即过渡到该复形的“形式化”,把它或者诠释为Deligne式的复形,或者诠释为一个晶体复形?

\noindent 这个问题的答案,无论如何在复解析情形中、以及在特征零的体上光滑相对概形的情形中是肯定的(除非我弄错),甚至无须引入正则性条件。那就是我此前已经特意强调过的(在(c)中,第1011页)“由Mebkhout的理论带来的全然出乎意料的现象”\fnote{710}{(*)}{在当今,强调这样的事实(至少就在此所论的数学那一部分而言)已变成一桩真正的\textbf{公共卫生事业},在一个上同调主题的几乎全部出版物、而且(我怕)所有那些以今天享誉的署名发表的出版物,都被写成恰好\textbf{略去}使所有这些文本活生生的那些强有力的伟大观念、并\textbf{抹杀或根除}如此关键的工具(旧的或新出现的)、如此要害的概念、如此多产的观念的作用与出处。有一种智识上的\textbf{腐败}(一种更深\textbf{腐败}的标志……)当今在我们的科学中人人眼皮底下蔓延,对于任何别的科学、在任何别的历史时刻,我都未曾有所耳闻。}。在正则情形中(包括“在无穷远”处),这本质上就是好上帝定理。在一般情形中,如果我没弄错,这应当毫不费力地归结于我称之为“完整性的上同调判据”(或“Kashiwara可构性定理之逆”)者,它归功于Mebkhout,在下一篇注“三座路标——或天真”(n\textsuperscript{o}~171(x),见第1028页)中有所讨论。

% ---------------- p.786 ----------------
\subsubsection*{b2.~三座路标——或天真}

\tocsubsub{b2.~三座路标——或天真}

\origpage{1022}\notehead{171(x)}{}

\dmark{}\noindent (5月5日与5月23日)\fnote{711}{(*)}{本子注“三座路标”出自“作品……”一注(n\textsuperscript{o}~171(ii))的一条脚注。请见置于该注末尾附近的参照符号。}Mebkhout在1972至1980年间发展的哲学,可以归结为\textbf{三大定理},三者都紧密联系着我在五十与六十年代发展的那些观念,但其中没有一项是我(或任何人)当年预见到的\fnote{712}{(**)}{如我在“荒唐的问题”一注(n\textsuperscript{o}~171(vi))中所指出,然而我早已知道Mebkhout整体对偶定理的一个变体,是关于 $X/S$ 真光滑相对概形的,用相对微分算子复形表述。确切地说,若 $L^{\bullet}$ 与 $L'^{\bullet}$ 是这样的复形,彼此“相伴”,则 $Rf_{*}(L^{\bullet})$ 与 $Rf_{*}(L'^{\bullet})$,作为导出范畴 $D(S,\mathcal{O}_{S})$ 的对象,是“完美”复形(局部上可由次数有界的有限型自由模复形表示),并且按完美复形的通常意义互为对偶。在 $S=\mathrm{Spec}(\mathsf{C})$ 的情形,这个定理或多或少等价于Mebkhout的定理(限制在代数且真的解析簇情形),然而有一个重要区别:我当年缺少一个“导出范畴”的观点来处理微分算子复形。另一方面、而且尤其是,我丝毫没有想到过这些复形(附加以Mebkhout廓清的适宜条件)构成“离散系数”(即德拉姆系数)的完美替代物。另一方面,至少从1966年起我就清楚,必定存在这样一种代数可构 $\mathsf{C}$-向量空间系数的替代物,对任意特征的相对概形有意义,而我的晶体观念正是朝此方向的第一步。如将在[Crystals]中看到的(即前一篇注“五张照片(晶体与 D-模)”(n\textsuperscript{o}~171(ix))所引的那些讲演),我的晶体反思的内在逻辑却已然把我重新带到微分算子复形的门前。那时我已经离Mebkhout的哲学很近。需要我的同调学生们(尤其是Deligne、Berthelot、Illusie)被“葬礼”综合征堵住,才没有在此后几年间廓清这一哲学。(我本人当时全神贯注于别的基础工作,而把晶体主题交给我的学生们照料。)}

% ---------------- p.787 ----------------
\noindent 第一座大定理是Mebkhout1972至1976年间工作的主要成果。它关乎光滑复解析簇 $X$ 的结构层在解析闭子空间 $Y$ 中的支撑上的\textbf{局部上同调}层 $H^{i}_{Y}(\mathcal{O}_{X})$(此概念由Sato与我各自独立引入)。这里Mebkhout之前没有人想到去作的关键观察,是 $X$ 上无穷阶微分算子环 $D^{\infty}$ 的运算\fnote{713}{(***)}{关于这些算子的定义——其名字初看吓人,但它们引出一套在每一点上都与寻常微分算子形式主义平行的形式主义——见前一篇注“五张照片(晶体与 D-模)”(n\textsuperscript{o}~171(ix))的(b)部分。},由于它们\dmark{}\origpage{1023}作用于参数 $\mathcal{O}_{X}$,也作用于这些上同调层。另一方面,在代数几何的“Zariskian”框架中,我(在五十年代末?)曾把这些层描述为 $\mathrm{Ext}^{i}$ 层的归纳极限。这经类比引导Mebkhout引入局部上同调的一个“代数部分”,以及一个典范箭头
\begin{equation*}
H^{i}_{Y}(\mathcal{O}_{X})^{\mathrm{alg}}\overset{\mathrm{def}}{=}\varinjlim_{n}\mathrm{Ext}^{i}_{\mathcal{O}_{X}}(\mathcal{O}_{X_{n}},\mathcal{O}_{X})\longrightarrow H^{i}_{Y}(\mathcal{O}_{X})\overset{\mathrm{def}}{=}\mathrm{Ext}^{i}_{\mathsf{C}_{X}}(\mathsf{C}_{Y},\mathcal{O}_{X}),
\tag{1}
\end{equation*}
其中 $X_{n}$ 表示 $Y$ 在 $X$ 中的第 $n$ 个无穷小邻域,而 $\mathsf{C}_{X}$、$\mathsf{C}_{Y}$ 是 $X$ 相应地 $Y$ 上的常值层 $\mathsf{C}$(后者以零延拓到 $X-Y$ 上)。第二个本质观察,是这一次 $X$ 上寻常微分算子的环 $D$ 作用于第一个成员。人们早已知道,所得到的这类层,无论右边那个性质上超越的成员,还是左边那个性质上“代数”的成员,作为 $\mathcal{O}_{X}$-模都相当庞大——毫无凝聚性可言,这是肯定的。诚然,人们也曾有一种感觉(至少在代数那一侧):总还是有某类“有限性”或“余有限性”,其意义Mebkhout之前没有人想过去精确化。

\noindent Mebkhout那条卓越的定理,就是第一个成员是一个 D-模,而且此外第二个成员(它看着更加难以处理) 只是由第一个成员经环的更换 $D\longrightarrow D^{\infty}$ 得出。由于第二个环已知在第一个上平坦,这也顺带蕴含(1)是单的。\noindent 同时,鉴于凝聚性结果——关于第二个成员(Mebkhout之前没有人对它有任何理解)的这个结果可以被视为一条非常强的有限性定理——特别地,它作为一个 $D^{\infty}$-模是有限呈现的(但也许不凝聚,因为人们不知道 $D^{\infty}$ 本身是否凝聚)。

\noindent Mebkhout处理的头一种情形,即具正规交叉的除子的情形,是他1974年通过的第三阶段论文的课题。即便这种情形已非平凡,而且当然是全新的——Mebkhout所解决的那个问题本身从来没有人看见过。况且这种情形结果正是关键的情形,Mebkhout(以一般性递增的逐次近似)借奇点消解能够归结到它\fnote{714}{(*)}{关于Mebkhout的局部上同调定理,请特别见:\emph{La cohomologie locale d'une hypersurface},载 \emph{Fonctions de plusieurs variables complexes III},Lecture Notes in Mathematics n\textsuperscript{o}~670,p.~89--119,Springer-Verlag(1977),以及 \emph{Local Cohomology of analytic spaces},Publ. R.I.M.S. Kyoto Univers. 12,p.~247--256(1977)。}。

\origpage{1024}\dmark{}(我方才陈述的这个结果,单凭自身在我看来就有这样的分量:在稍稍正常一点的条件下,它本会为它的作者赢得国际声誉。同样,他所处理的头一个关键情形,已然显示了一种视觉上的独创性,“正常地”本会为他赢得他的长辈中那些有资格品鉴其妙处的人(例如我的每一位前学生,无一例外)的热烈鼓励。罢了……

\noindent 事实上,在这四年间,Mebkhout得出一个比我方才陈述的还要详备的结果。他证明他所研究的 D-模不仅凝聚,而且此外是\textbf{完整的}(这是他在日本学派中找到的概念),而且此外是\textbf{正则的}\fnote{715}{(*)}{Mebkhout对正则性的原始(超越的)定义,在“作品……”一注(n\textsuperscript{o}~171(ii))中有所回顾,见第950页脚注(*)。}(取他 ad hoc 定义的意义,其灵感来自我的代数-解析德拉姆上同调比较定理)。更进一步,他证明出发的那个可构 $\mathsf{C}$-向量空间层 $\mathsf{C}_{Y}$(它进入(1)第二个成员的定义)从 $D^{\infty}$-模复形 $\mathrm{RHom}_{D^{\infty}}(\mathsf{C}_{Y},\mathcal{O}_{X})=C$,凭那条非凡的反演公式
\begin{quote}
(2)\quad $\mathsf{C}_{Y}=\mathrm{RHom}_{D^{\infty}}(C,\mathcal{O}_{X})$。
\end{quote}
\textbf{重新构造出来}。从来没有人梦想过这样的公式——而在五年后那决定性的一天之前,也不会有人梦想它,那时这哲学的力量显露出来,同时给葬礼发出信号,在祖先身边、连同那位把它带来的人……为了梦想它,本不该已经埋葬了祖先的哲学(借导出范畴、带不带下划线的 $\mathrm{RHom}$ 以及别的“无用细节”之名);此外,还该懂得欣赏一个极平淡却充满神秘的几何局面(支撑在具正规交叉的除子中的局部上同调),并走到神秘\textbf{尽头}。这个“尽头”,它还不在我刚才描述的那座1976年的辉煌定理中——但早在那一刻,Mebkhout就对它有清晰的愿景:那就是双重的“好上帝定理”,一个用于完整正则 D-模,另一个用于完整 $D^{\infty}$-模,以及前面说过的双重反演(或“双对偶”)公式\fnote{716}{(**)}{在前一篇注“五张照片(晶体与 D-模)”(n\textsuperscript{o}~171(ix))的(b)部分。}。这同时也是离散(解析可构)系数与“连续”系数之间关系问题的一个美妙简洁的解。

\origpage{1025}\dmark{}(但我在预支。当他证明了构成他著作与哲学的第一座大路标的那个定理时,那个“尽头”,虽已被清晰洞察,在他看来仍远得令人眩晕。如果他在身边找到一位能干而仁慈的长辈,且有最低限度的经验与数学嗅觉,那人会纠正他:显然他已经离它很近,而所要克服的困难,一如发现工作中常见的那样(甚至不妨说,向来如此……),更多是心理上的,而非技术上的。但在着手追逐那无穷远之前,他先攻整体对偶定理——那个注定要“统罩”已知对偶定理的定理,无论对凝聚系数还是对离散系数。

% ---------------- p.789 ----------------
\noindent Mebkhout著作中无处不在的那种深刻动机,把局部上同调问题与整体对偶问题这两个问题联系起来,那就是对离散系数与连续系数之间一种本质统一的预感。那也正是我1966年晶体进路的指导线索,它努力以“连续”的术语把握“德拉姆系数”(性质上本质是离散的)……

\noindent 这里不是重温Mebkhout对偶定理陈述的地方\fnote{717}{(*)}{这个陈述在“作品……”一注(n\textsuperscript{o}~171(ii))中有所回顾。}。它的证明曾撞上严重的技术性困难,源于超越语境,而他以上同调下降与核 EVT(拓扑向量空间)的技术克服了它们(这些技术我本人也并不陌生,即便Mebkhout是至今唯一还坚持引用祖先的人……)。就他的对偶哲学而言,这个定理是一个本质的路标。如果同Mebkhout一起牢牢记住:把它应用到完整 D-模复形上,它除凝聚对偶之外还包含解析可构离散系数的整体对偶\fnote{718}{(**)}{在Mebkhout确立他的整体对偶定理(1976年)时,此外他还不曾证明每个解析可构的 $\mathsf{C}$-向量空间层都来自一个 D-模复形。但他对此深信不疑。},那么可以说,它本身已经孕育着Mebkhout式 D-模的全部哲学。它的分量,早在1980年(他的论文答辩后一年\fnote{719}{(***)}{见“墓外相会”一注(n\textsuperscript{o}~78)。})他第一次跟我谈起时,在我看来就是一件明摆的事。我不认为我曾有幸启发过任何一位\origpage{1026}与我接触的学生\dmark{}作出分量相当的工作\fnote{720}{(*)}{我在这里尤其想到那些跟我准备论文的学生。Deligne的情形是另一回事,因为他的论文在我离去之后进行,而且通篇不提我的名字,尽管他工作(关于Hodge--Deligne上同调)的初始灵感来自我对种种“系数”的问题意识,而后者也预见了“霍奇系数”的一套形式主义。Deligne的工作是朝此方向的第一步,远比Mebkhout所完成的那一步片面,而后者是在(与霍奇密切相关的)“德拉姆系数”方向上。诚然,Mebkhout因长辈的漠视与鄙薄而面对重重障碍,却不受那麻痹了我的学生们的葬礼综合征之苦。(关于这一点,见“……与羁绊”一注(n\textsuperscript{o}~171(viii))。)}。

\noindent Mebkhout当时还费了很大的劲才让这个定理得以发表,它浑身散发着“Grothendieck味儿”。(《数学年刊》把它退回了,让他明白这类东西够不上所需水平。结果它终究还是发表了,在 \emph{Mathematica Scandinavica} 上,1982年\fnote{721}{(**)}{整体对偶定理(用于凝聚 D-模),\emph{Mathematica Scandinavica} 50 (1982) p.~25--53。另见“Poincaré 对偶”,载巴黎第七大学奇点讨论班(Pub.~n\textsuperscript{o}~7),1977--1979,以及尤其“The Poincaré-Serre-Verdier duality”,载 \emph{Proceedings of the Conf. of Algebraic Geometry},Copenhague(1978),Lecture Notes in Mathematics n\textsuperscript{o}~732,p.~398--418,Springer Verlag(1979)。这两篇讲演的引言、尤其是第二篇的引言,代表了Mebkhout所带来的哲学的一幅草图,那是在他还是其唯一保管人与辩护人的时候。}。)我相信那正是他讲授 D-模哲学时最爱的主题,但其精神与日本人的截然不同。他对我说,这个定理确实有让听众或偶然的对话者惊叹的天赋,唯独那些属于建制圈子的人\fnote{722}{(***)}{(5月24日)这与我自己的观察若合符节。看来名人的处境容易使人养成这样一种自满,对“再美也不配她屈尊欣赏”的人。我不知道这些习性在当今整个科学界、乃至自来如此,是否通例。我起初受到一个不存在这种自满精神的圈子的接纳,曾是我极大的幸运。它想必是蹑手蹑脚地到来的,经年累月,在这些人那些人身上安家落户,渐渐无人(只有Chevalley除外……)察觉。一切看似如旧——然而一切已经不同。它已经像一层细薄的灰尘,覆盖着事物原初的新鲜。我也曾触及这灰尘,与别人一样。而今天,当我重新面对从前的学生或朋友时,我常有这种感觉:这灰尘已积成厚密的一层,形成一副密不透风的铠甲,透过他们向我发问……})例外,而且无一次例外。这是一件让我感到宽慰的事。它
%JOIN:para

表明这种餍足的自满精神(它使一切事物之美黯然失色,无论它多美)在数学界并未成为普遍。它主要(即便不说是排他地)肆虐于高层圈子,近十年来我确实有充分机会与它相识……

\dmark{}\origpage{1027}(应当用一个已提及的局部性质的结果来补全这个整体对偶定理,它同样深刻,说的是:上同调层凝聚的 D-模复形的自然对偶化函子(它把完整复形变成完整复形,完整正则复形亦然),在这些复形上此外还与德拉姆函子 $\mathrm{DR}$(“相结合的微分算子复形”,被看作上同调可构的 $\mathsf{C}$-向量空间层复形)相容,相容于我当年在这些复形上引入的那个自然对偶化函子\fnote{723}{(*)}{这就是那个对偶,经我的学生与旧友们的一致同意,它在此期间成了“Verdier对偶”(无论在复解析情形还是平展情形)……(关于这一点,例如见“好的参考文献”一注,n\textsuperscript{o}~82。)}。这种相容性显然是Mebkhout对偶形式主义的一个本质成分,为的是理解他的整体对偶定理的意义。出于一种我不明的原因,他把它称作“局部对偶定理”\fnote{724}{(**)}{这个结果就是以这个名称见于Mebkhout论文第三章的。他对我说,他这个名称(以及“双对偶定理”之名)的灵感来自我引入的术语——然而在我看来,“局部对偶定理”只是我所廓清的“双对偶定理”的另一个名字,它代表着后者的一个重要侧面,即“几何”侧面。这个相容性结果(Mebkhout向我解释)是他证明他在同一章中称作“双对偶定理”者的重要一步。(关于后者,见前一篇注“五张照片”的(b)部分。)证明的问题且搁在一边,从一种“哲学”或一套“瑜伽”的观点看,好上帝函子应当与对偶化函子交换,这当然是一件“显然”的事(因为既有好上帝!)。一则滑稽的细节:Kashiwara(Mebkhout1978年1月曾有机会与他面谈)\textbf{不相信}这个定理为真!这足以说明他多么不在行,而他恰缺少“六种运算”风格的几何视觉。这并不妨碍他后来在Mebkhout把自己的第三章寄给他(1978年2月)之后,把这结果据为己有(当然不提作者),收进他与Kawai那篇已引的大文章(见第1005页脚注(*))(loc.~cit.~第4节命题1.4.6)。正是在这部著作中,“双对偶定理”(loc.~cit.~第4节1.4--9)也被不经任何手续地据为己有(以“重构定理”之名)。这足以说明,生于巴黎的“新风格”诸位大师的太平洋对岸信徒们(他们出现在一个“Grothendieck学派”蒸发了、未留痕迹的地方),比起他们的法国同行来毫不逊色。我的(离散系数)双对偶定理也同样见于Kashiwara--Kawai这同一著作的同一不竭的第4节(命题1.4.2),然而当人们恬不知耻、不假思索地劫掠那位身后无名的学生(老板们公然弃之不顾)时,人们却按规矩向对面那位显赫的同行脱帽致意,恰如其分地引用Verdier提供的“好的参考文献”(而他本人也在劫掠一位从不被点名的逝者……)。这些骗局在消息灵通人士中固然臭名昭著,Mebkhout在这方面也有过若干反响。但显而易见,只要事关消灭那位不可引用的祖先及其倒霉的后继者,它们就被视为得体而受欢迎。}。这个深刻的定理,连同那著名的“对应”(人们肯赏脸称它“Riemann--Hilbert”时),被“所有人”(Verdier与Deligne打头)当作一件“众所周知”、不言自明的事,而且尤其从不提名某一位无名者(而“所有人”都清楚,尤其不要引用他)……

\dmark{}\origpage{1028}我终于来到Mebkhout著作中的第三座大路标。从技术上可以说它由三个(或至少两个)不同的定理构成,但它们联系得如此紧密,以致在Mebkhout的心中它们显得不可分离。早在1978年1月,他就证明了“$D^{\infty}$-模”那一侧面:即“相结合的德拉姆复形”函子限制到完整 $D^{\infty}$-模复形上的限制 $m_{\infty}$(即“Mebkhout函子”)是一个范畴等价(与上同调可构的 $\mathsf{C}$-向量空间层复形之间)。已知这个函子与对偶化函子交换,把这个定理改经相联系的反变函子 $\delta_{\infty}$ 来表述便是自然的,它由
\begin{equation*}
C\longmapsto \mathrm{RHom}_{D}(C,\mathcal{O}_{X})
\tag{3}
\end{equation*}
给出,
%JOIN:para
% ---------------- p.791 ----------------
而断言这个函子是一个(反)等价也是一回事。这个定理于是可以由Mebkhout那条绝妙的\textbf{反演公式}(或“重构”公式,或“双对偶”公式)来精确化,它把拟逆函子表达为
\begin{equation*}
F\longmapsto \mathrm{RHom}_{\mathsf{C}_{X}}(F,\mathcal{O}_{X})
\tag{4}
\end{equation*}
与此同时,Mebkhout还证明了Kashiwara可构性定理的一个\textbf{逆定理},即:若一个上同调凝聚的 $D^{\infty}$-模(或 D-模)复形使得相结合的德拉姆复形(作为 $\mathsf{C}$-向量空间层复形)上同调可构,则它是完整的(\textbf{完整性的上同调判据})。在 $D^{\infty}$-模复形的情形\origpage{1029}(其中不存在正则性问题),这因此蕴含:在那个\dmark{}导出范畴中(1978年直至1981年早就没人在里面工作了……),该复形(或不如说它的对偶)凭反演公式“重构”,除唯一同构外唯一。

\noindent 如我在别处所解释\fnote{725}{(*)}{见“作品……”一注(n\textsuperscript{o}~171(ii))第952页脚注(*)。},从那时起,Mebkhout手里就拥有证明好上帝定理(对 D-模亦然)所需的一切:即德拉姆函子限制到完整正则 D-模复形上的函子 $m$ 是一个范畴等价。这个结果给他的鼓舞较小,因为显然没有配套的反演公式\fnote{726}{(**)}{前面已经看到毕竟有一条——而我稍后还要回到这一点。}。无论如何,即便他那条绝妙的反演公式,对谁都不冷不热——从他论文的准导师Verdier开始(而Verdier后来还赏脸担任答辩委员会主席)。为了重新作技术努力去证明一件他反正确信、且感觉自己已握有全部证明手段的事,这可不是一种很令人鼓舞的氛围。他要到由一条据说高不可攀的猜想(这回不是Weil的,而是Kazhdan--Lusztig的那条)的证明所触发的“冲刺”启动之后,才去管它。

\noindent 这恰好像有意安排,正是人们突然急需的另一侧面。无论如何,当时“所有人”都那么急着使用这件崭新出现在市场上的“崭新工具”,而且所有人都心照不宣:尤其不要提出证明的问题——万一显出那工作早已由一个不可引用的人做完了呢。似乎除了当事人自己,没有人想到去把已经写就的 $D^{\infty}$-理论的各块抄录拼合起来,以证明 D-理论中该证明的那条定理。迄今唯一公开发表的证明\fnote{727}{(***)}{参考文献:\emph{Une autre équivalence de catégories},Compositio Mathematicae 51 (1984), 63--88。}看来确实是Mebkhout的那一篇,去年发表(而于1981年6月收到,正是那场难忘的反常层会议的同一个月……)。

\noindent 我在前一篇注(b)部分中解释过一个简单的原则,受Deligne通往德拉姆系数的进路启发,用以在(完整正则的)D-模框架中收复一条“反演公式”(或\origpage{1030}借Mebkhout的说法\dmark{}说,“双对偶”公式)。自从世界各地的人们开始就 D-模这门新的“时髦点心”办讨论班以来,我不知道这种十分自然的进路是否已被廓清——Mebkhout无论如何不曾有所耳闻。可以肯定的是,如果Deligne有那种“我那个时代”被视为理所当然的本能反应,那么正是他本人,在1979年6月得知一位无名者的美妙观念之后,就会鼓励他把其关键结果的 D-模那一侧面(更接近代数)的证明也写出来,并会向他建议他那条美妙反演公式的“pro”变体——说到底相当显然。同样,从那时起,对已经为此付过学费的Deligne来说,Mebkhout的观念将

% ---------------- p.792 ----------------
带来所缺的德拉姆系数(至少在特征零的体上代数几何中);理所当然的事,是鼓励他作出必需的调整,以陈述一个(不如说这一回是Mebkhout的)好上帝定理,用于复代数簇\fnote{728}{(*)}{如我已曾有机会指出的,在代数框架中,当需要以代数方式复述代数可构的离散系数时,除了局部的完整性与正则性条件之外,还应给所考虑的 D-模复形加上一条“Deligne-Mebkhout”的在无穷远处的正则性条件。}。

\noindent 但此一时彼一时。将不说一场代数簇上同调的新出发,是由一位无名的泛泛之辈孤独而顽强的努力完成的,他自承师法一位逝者,而风雅世界中好久以来没有人敢贸然说出他的名字\fnote{729}{(**)}{诚然,人们还没有找到办法为 EGA 与 SGA 找到替代的参考文献。但这些天降的缩略语中不含任何对一个必须保持沉默的名字的暗示。众所周知,SGA 这个缩略语指一个由布瓦玛丽主持、在一批十分体面且完全可具名的数学家推动下举办的代数几何讨论班,诸如 M.~Artin、J.L.~Verdier、P.~Deligne、L.~Illusie、P.~Berthelot、N.~Katz、P.~Jouanolou,甚至另一些较不出名但同样可引的人物。显然那里曾有过一个繁荣的代数几何“学派”,号称“布瓦玛丽学派”,其核心与灵魂是所引姓名中最光辉的那一位。关于这个“布瓦玛丽学派”以及作为其表现的 SGA 缩略语的更多细节,请特别见“驱逐(2)”与“殡仪馆——为科学服务”两注(n\textsuperscript{o}~169$_1$ 与 175)。(另见“双重含义——或骗术”一注(n\textsuperscript{o}~169$_7$)第899页第3段。)}。将不说复兴恰是通过十年来\origpage{1031}逝者的继承人们一边瓜分掠获物、一边埋葬的那种数学。天真的Mebkhout,若他想“活下来”\dmark{}并且“出人头地”,只须走“新风格”那条现成的路\fnote{730}{(*)}{关于这种风格(它取代了一个消失得无影无踪的“Grothendieck学派”……),见“恭贺——或新风格”一注(n\textsuperscript{o}~169$_9$)的结尾。},如其他杰出的年轻人(甚至有不太年轻的)急忙所做的那样。引用自己思想的(不可引用的)源头,这是怎样的怪癖,而把鱼搅浑、只引那些\textbf{该}引的人是多么简单。Mebkhout,我相信你的账到头了!

\noindent 你落进了一个不是为你而设的世界——而我仍然为此为你高兴,你不是为\textbf{那个世界}而造的。你做了你感觉自己该做的工作,不赶时髦,不盘算回报,只是信任你自己的本能——哪怕因此孤独地走你的路。你做了\textbf{你的}工作,而不是窥伺那些决定何为善、何为体面、何为不善的人发出的隐晦(和不那么隐晦的)信号。你没有为讨人喜欢而见风使舵,你没有看到黑却说“白”,或反过来——你是用\textbf{你的}眼睛看的。我不必为此向你道贺——你不曾寻求祝贺,不论是我们的还是任何人的。而为此我为你、为所有人感到高兴。

\subsubsection*{b3.~师父的角色(2)——或掘墓人们}

\tocsubsub{b3.~师父的角色(2)——或掘墓人们}

\notehead{171(xi)}{}

\noindent (5月5日)\fnote{731}{(**)}{本子注出自“作品……”一注(n\textsuperscript{o}~171(ii))的一条脚注。关于指向本子注的参照,见所引注末尾附近(第956页)。}这里自然的问题,当然是在代数几何中是否存在一套“六种运算”的形式主义,用于未必属于 DRM 类型的 D-模(或“晶体”),它将“统罩”我在凝聚与离散两种情形下引入的那两套——先假定,为明确思想,我们处在体 $\mathsf{C}$ 上。第一个困难源于 D-凝聚性概念对晶体的自然张量积概念、或类似的拉回像运算都不稳定\fnote{732}{(***)}{(5月22日)Mebkhout向我指出,他证明了完整性与正则性条件对(在 $\mathcal{O}_{X}$ 上的)总张量积运算以及对拉回像概念是稳定的,并且好上帝\textbf{反变}函子 $\delta$ 在其上交换。(相反,好上帝协变函子 $m$ 在其上并不交换,它把寻常拉回像变成非凡拉回像。)利用这个结果可以证明,对de~Rham--Mebkhout系数不存在一套“六种运算”的形式主义,能“延扩”张量积与拉回像这两个已知的基本运算。尤其,范畴 $\mathrm{DRM}^{b}(X)$ 不容许“内 $\mathrm{Hom}$”运算(扮演 $\mathrm{RHom}$ 的角色),且对 $f : X\to Y$,函子 $f^{*}$ 一般不容许右伴随 $Rf_{*}$。Mebkhout目下引入的函子 $Rf_{!}$(对光滑的 $X$、$Y$ 且 $f$ 固有)是 $f^{*}$ 的一个\textbf{左}伴随。(注意:de~Rham--Mebkhout系数上的 $Rf_{!}$ 运算的定义方式,使好上帝\textbf{协变}函子在其上交换,$Rf_{*}$ 亦然——无论对错……)因此可见,就用de~Rham--Mebkhout语境中现有的“自然”运算而论,它们本身\textbf{不}构成一套“六种运算”理论,而构成一种对偶理论。于是问题就在于看这种理论在多大程度上能延扩到不再假定完整且正则的(姑且说拟凝聚的)D-模上(例如仅仅完整——这个条件经张量积与拉回像保持)。似乎特别是,整体对偶公式可以为上同调凝聚(乃至仅仅拟凝聚)的 D-模复形、以及特征零的体 $K$ 上有限型分离概形的任意态射 $f : X\to Y$ 写出,以便同时统罩凝聚对偶定理与离散对偶定理,至少取如下形式:对偶化函子“交换”函子 $Rf_{*}$ 与 $Rf_{!}$。}。\origpage{1032}因此,为了希望有一套六种运算的形式主义,就必须\dmark{}用一个比 $\mathrm{Cris}^{*}_{\mathrm{coh}}(X)$ 更大的范畴来工作,

% ---------------- p.793 ----------------
也许是“拟凝聚”晶体的范畴(取一个显而易见的意义)——但那样就几乎无望收复一条双对偶定理!此外,经 $\mathcal{O}_{X}\to D_{X}$ 的自然标量扩张函子显然不与张量积交换——因此,即便对晶体有一套六种运算理论,它延扩那套(多亏Mebkhout,目下在实质上已为人知)de~Rham--Mebkhout晶体理论(借“结构搬运”从“离散”理论经好上帝函子得到),它也不会延扩凝聚 $\mathcal{O}_{X}$-模的那一套\fnote{733}{(*)}{这个陈述更应借一套“六种运算的对偶理论”来重新表述,见前一条脚注。}。这也许并不排除:仍然可能存在一条“整体对偶定理”(拟凝聚晶体版本),用于特征零的体上有限型概形的一个(姑且说)真态射,它(在一个显而易见的意义下)“统罩”de~Rham--Mebkhout晶体的“已知”对偶定理(实质上,仍是借结构搬运),以及凝聚情形中已知(不带引号)的类似对偶定理\fnote{734}{(**)}{可以设想这样一个对偶定理有三种不同形式。或说上与下的对偶化函子“交换”函子 $Rf_{!}$ 与 $Rf_{*}$,或说两个恰当定义的函子 $Rf_{!}$ 与 $Rf^{!}$ 互为伴随,或写一条“投影公式”(它统罩两种陈述): $Rf_{*}(\mathrm{RHom}(F,Rf^{!}(G)))\approx \mathrm{RHom}(Rf_{!}(F),G)$。}。

\dmark{}\origpage{1033}(令我相当震惊的是,Mebkhout本人连这最后一个问题都没有提出过,早在他达到自己的“绝对”对偶定理(对应于靶簇退化为一点的情形)的表述那一刻——就在最近他还显得不太“感受”到它\fnote{735}{(***)}{(6月8日)然而Mebkhout向我保证,他早就切实提出过这个问题。如果说我有过相反的印象,那必定是因为这个问题对他来说始终完全停留在柏拉图式的层面上。}。这使我深切地感到,某一种“哲学”——它在六十年代上半叶对我已成为第二天性,而且(我曾觉得……)对我的学生们也是——已被所有人忘却到何种程度,首先那些自命为其掘墓人而非传承者的人。而我看到,这也正是我离去后一门理论(概形的上同调)所经历的那场惊人停滞的主要原因,而我当年把它留在全盛之中。

\noindent 应该说,Mebkhout置身于复解析的超越语境中,而非概形语境。这就引入了可观的技术困难,在某种意义上说是“寄生的”,当需要达到对本质的变差现象的理解时更是如此。在这方面,他的长辈们又一次辜负了他们的任务,那就是把他们在我接触下获得的经验提供给新来者(正如我当年曾把我的经验提供给他们一样……),由此引导他(或至少给他一些启发),尤其是在选择投身的方向时。

% ---------------- p.794 ----------------
\noindent 但照亮与引导,这也是一种\textbf{服务},而他们早已毫无保留地选择了师父的角色。

\subsubsection*{b4.~死去的书页}

\tocsubsub{b4.~死去的书页}

\notehead{171(xii)}{}

\noindent (5月5日)\fnote{736}{(*)}{(5月22日)本子注与上一篇一样,出自“作品……”一注(n\textsuperscript{o}~171(ii))的一条脚注。见该注末尾的参照符号,第957页。}Mebkhout对我说,在我两年前我们相遇时跟他谈起之前\fnote{737}{(**)}{这次相遇在“墓外相会”一注(n\textsuperscript{o}~78)中有所讨论。},他还从未听人说出过“六种运算”这几个字——他很纳闷我想说的是哪些“运算”!显然,他从来没有想到(除我之外,似乎也没有任何别人想到)去回顾一套十分简单的上同调形式主义的主要成分,发现其中有六个基本函子或双函子,分成三对伴随函子,带着这样的箭头与相容性等等。这些事在我看来是如此显然,\origpage{1034}以致我原以为任何读者——无论\dmark{}是读《留数与对偶》阐述凝聚对偶要素的,还是读 SGA~4 或 SGA~5 阐述离散对偶要素的(况且基本用同一套公式)——都会(像我早在五十年代所做的那样,虽则我承认没有做到底)为自己立起一套多少系统、多少完备的主要同构与主要相容性的公式,因为只有这样、别无他法,才能领悟一门新语言的精神,把它同化为己有,使它成为“自己的”。毫无疑问,微积分的先驱们正是这样、也正是别样地做的,为的是在对无穷小获得一种精致而可靠的直觉,而那时他们尚缺概念工具去按后来(或重新)出现的严格性准则把握它们……

\noindent 隔了二十年回头来看,我意识到在所引的那些“参考文本”中——它们都以极大的细心乃至才华写成,而一切“真正的工作”(按时下的要求)都已做出,其顶峰是“那条”主对偶公式,即 $Rf_{!}$ 与 $Rf^{!}$ 之间的伴随公式(几乎是唯一被认为值得注意与努力的一条,哪怕第二天就忘掉它,如同没见过森林就忘了树……)——然而在所有这些文本中,\textbf{主要的}东西没有说出口,也没有从作者传给读者(假定它被作者本人很好地看到与感受到)。“主要的东西”是一套“瑜伽”,一种“哲学”,是一条穿越(在此情形下)代数几何中上同调丛林(以及别处)的万试万灵的指导线索。人们可以在“一切都已做完”(据说)之后,用五十页或一百页把它大加发挥;也可以只用几页把它唤起,而把发挥的工作留给读者,由他按自己的需要或自己的满足,以为己用的程度去发展它。

\noindent 正是那么几页——无论关于“六种运算”、关于Motif,还是关于许多别的东西\fnote{738}{(*)}{这些行写就之后,我曾能够确认:关于六种运算,我在此弄错了——事实上我是被 SGA~5 的屠杀版欺骗了,Illusie细心地把我在口头讨论班中大加发挥的“六种运算瑜伽”的一切痕迹根除了,连同一套被充分注释的完整公式。}——那些我强烈感受到、却不知道把它们写出来何其重要的书页——恰恰是它们,首先在我的文字著作中缺失了。我当时全神贯注于\origpage{1035}无休止的细心任务,\dmark{}为大家服务,做那沉重的“计件工作”,而那是唯一被认为该出版的东西——我不曾感觉到还有更本质的书页,而写那些书页我是\textbf{唯一}能做的人。\textbf{我要说的本质}东西没有进入写就的书页,只是口耳相传而已
%JOIN:para
% ---------------- p.795 ----------------
——在它传得过去的时候!至多,它也许是在无尽的奠基卷帙的字里行间——可当今有谁懂得读字里行间呢?

\noindent 因此,本质的东西,就是日复一日托付给那些我作为数学家的生活中堪称“亲近”的人、首先是我的学生们的东西。那是一件理所当然的事,毫无刻意。我不曾想到,以某种方式我也是在这样授予他们一份可观的\textbf{权力}。并不是我不曾感受到我所构思与传递之物的力量,但那种力量,它也是理所当然的。对我而言,肯定至少在数学中,“力量”与“美”过去是、现在仍然是同一回事。我不曾想到人们会滥用这些东西,这些对我充满平静而炽烈的生命、为生存与生育而造的东西。当我离去时,以着实再突然不过的方式,我对它们没有一丝忧虑。这些我从未想过要写的书页——我心中毫无疑问,它们的讯息早已被接纳并铭记,而这些“亲近”者将成为许多活生生的书页,说出那讯息,并以他们所能献上的最好东西丰富它。

\noindent 那些我曾以信任与尊重相待、视作弟弟并在他们身上认出我自己的人,选择了埋没与沉默。而当那一位到来时——忠于他自己,而他们在其中认出了我——这些饱有一切的人,却选择把他留在他们紧闭的门前——一个异乡人与一个闯入者。我不认识你!而这些未写的书页,这些白说的话,变成了死去的书页,死在这些门第高傲而紧闭的殷实人家中,而被拒认的弟弟只好尽力在自己心中把它们重新找回,经漫长而摸索的劳作。他独自不得不在那以千百卷计的无从解脱的丛林中开辟道路。经历过这一切的人,即便他曾有幸(如我当年)得到经验丰富而仁慈的向导们兄弟般的扶助,也深知我所说的是什么……

\noindent 他在日子与岁月中艰难地走出了自己的路——一条坎坷的路,事后看来有时没有罗盘,或至少没有别的罗盘,只有一种还在寻找自己的嗅觉,穿过\origpage{1036}一段\dmark{}经痛苦与艰辛获得的经验。他没有为自己的用途重写那些现成的书页、那些罗盘书页(它们在高门大户中变成了死去的书页)——除了零星碎片。他写了\textbf{另一些}书页,\textbf{他的}书页,痛苦地属于他自己。他坎坷地、顽固地写下它们,在所有人的漠然中。然而这些往往笨拙、配上个莽汉也相称的书页,我那些精明殷实的前学生(若他们肯屈尊读一读)定会以怜悯视之而一无所见——它们是\textbf{本该}写的书页,是那些我连想都不曾想过要写的书页的自然的、“显然的”续篇,因为它们在我看来太理所当然了……

\subsubsection*{c1.~一个愿景的萌发,或闯入者}

\tocsubsub{c1.~一个愿景的萌发,或闯入者}

\notehead{171$_{1}$}{}

\noindent (4月15日)\fnote{739}{(*)}{(5月30日)接下来这三篇注(n\textsuperscript{o}~171$_{1}$ 至 171$_{3}$)写于4月15日至18日(1985年),那时《封神》还只有十页光景。由于Zoghman~Mebkhout的来访引发的对“四种运算”反思的重新启动,这些注在5月份大为扩充。十页变成了一百多页,因而其中绝大部分比接下来这三篇注要晚。由此有若干局部的重复,某些事实或插曲在前面的注与接下来的注中都有提及或描述,只是角度不同。为保持文字的自发,我不愿作调整以消除这些重复。}我想借我的共同下葬者Zoghman~Mebkhout本人最近一次在我处的做客,提供一些关于他的奇特遭遇的最新鲜细节,如他在我们谈话的当口这儿那儿零星吝惜地告诉我的。

\noindent Zoghman曾有幸与他的“老板”\fnote{740}{(**)}{(5月24日)Mebkhout坚持认为“老板”一词(即便带引号)在此用得不妥。自1972年入门以来直到今天,他做工作都没有老板,只靠自己想办法。Verdier只不过是他论文答辩委员会的主席。除此之外,他的作用仅限于把“好的参考文献”传给Mebkhout,而这份文献大有用处,那时 SGA~5 仍在我的同调学生们联手扣押之下(而这正是为了诸如“好的参考文献”那种操作的需要……)。}J.L.~Verdier有三次“会谈”。第一次发生在1975年——他当时需要一条技术性的结果,而这条结果(事后看来)就包含在关于解析可构离散系数的双对偶定理之中——那时Zoghman连可构性这个概念都还不知道。(这是我很早——五十年代——就引入的一个概念,后来在平展拓扑的框架中又得到重提,见SGA~4。)在那个时刻,\origpage{1037}这个概念在分析中还根本谈不上“尽人皆知”,不像\dmark{}今天这样。事情恰恰是:这正是他的工作所需要的那个概念。Houzel(他与Verdier同期听过SGA~5,只是把我当时在那里讲的东西多少忘了些)建议他去找Verdier。这便是与那位大人物的第一次“会谈”。Verdier于是让他明白:他所要的东西(即两个具有同构“对偶”的离散复形必定同构)在某些技术性条件(恰恰就是“可构性”)之下成立,这些条件他会在Verdier正要交给他的那份手稿里找到现成的阐述。那份手稿正是“好的参考文献”\fnote{741}{(*)}{即 J.L.~Verdier的文章《结合于一个闭链的同调类》,Astérisque n\textsuperscript{o}~36 (SMF),第~101--151~页(1976年)。在相继的两条注“好的参考文献”与“玩笑——或‘权复形’”(n\textsuperscript{o}~82、83)中有详尽的谈及,在“攀登的若干插曲”一注(n\textsuperscript{o}~169~(iii))的第3插曲中也有较简略的谈及。},在那里面(在性质相同的种种壮举之间)他还装出发明了可构层、发现了双对偶定理(连同它的证明)的样子——而这些东西,他早在十二年之前(1963年)就从我口中得知了\fnote{742}{(**)}{早在五十年代下半期,我就对离散层的种种“可构性”概念(代数意义的、复解析意义的、实解析意义的、分片线性的——温和拓扑的语境则尚待到来……)发生了兴趣,与凝聚性概念并列,把它们当作在层的框架中表达有限性条件的自然概念;我还提出了这些概念在“六种运算”下的稳定性问题。平展上同调此后的进展(1963年及其后)促使我回到平展框架中的这些问题,并发展出使之能以统一方法处理的技术(拆卸与消解),它们同样适用于复代数簇与复解析簇的超越语境。双对偶定理——在平展框架中(以纯粹性与消解为前提)与在超越语境中同样成立(且证明相同)——我早在1963年就已提炼出来。它而且见于SGA~5的最早一讲(1965年),并在1977年Illusie版的大屠杀中幸存了下来。}。关于我的这个人,他在这个话题上一个字也没有提,无论在那次会谈里,还是在次年刊行的那份手稿里。反正Zoghman是心满意足地离去的,对那位大人物满怀感激——正是大人物,在那一刻、以及此后几年里(可构性概念将在他所有的工作中扮演关键角色),恰好供给了他所需的东西。

% ---------------- p.796 ----------------

\noindent 到了1976年初,他开始对对偶发生兴趣,并为我曾在凝聚情形与“平展”离散情形中发展、而后由Verdier在拓扑离散情形中接手的那两套对偶形式主义之间的相似所吸引。那正是这套形式主义多年以来\origpage{1038}已陷于废弃、而我的学生们对导出范畴实行着一场\dmark{}默契而严格的联合抵制的时刻——而导出范畴正是它的自然语言。“六种运算的形式主义”这一概念连同这个名字,自五十年代以来、贯穿整个六十年代,一直是我最重要的主导观念之一,却在我离去之后即刻变得(并且直到今天依然)是严格意义上的禁忌。(Zoghman两年前来见我的时候\fnote{743}{(*)}{这次来访在“墓外相会”一注(n\textsuperscript{o}~78)中有所谈及。关于针对“六种运算”实行的联合抵制,另见“死去的书页”一注(n\textsuperscript{o}~171(xii))。},还不曾听人说出过“六种运算”这几个字,起初也不知道我说的“运算”是些什么——而我还以为这是二十年来人人都熟悉的概念呢!)由此可见,投身这一方向的处境是何等不利——在那里,他注定要在全然的孤独中工作。但这并没有妨碍他早在1976年就提炼出一条对偶定理——关于非奇异复簇的——一举“统罩”Serre对偶定理与离散对偶定理(他称之为“Poincaré--Verdier对偶”),用的则是关于D-模复形的一条对偶陈述(其中还包含关于微分算子复形的一条整体对偶陈述)。他取的“系数”,其一般性也远远超出了Serre的情形(限于局部自由层)与Poincaré的情形(限于局部常值的离散层)——忠实于我以那套当时备受排斥的“六种运算”形式主义把这些主题引入时的精神。

% ---------------- p.797 ----------------
\noindent 两年前当Zoghman向我讲解这一定理时,我既感到它的趣味(这对我是一望而知的),也感到它的局限;因为以“六种运算”的精神来看,对我同样一望而知的是:“对的”陈述应当是关于解析空间的态射 $f : X\to Y$ 的一个陈述,取(比如说)两个函子 $Rf_{!}$ 与 $Rf^{!}$ 之间伴随关系的形式。诚然,置身超越语境带来了相当可观的额外困难,它们(在我看来)有力地遮住了Mebkhout的眼睛,使他看不清对偶之中那些本质代数机制的单纯性——而他身边竟无一人——尤其是昔日我的学生们中间——能够(或乐于……)让他体会到这一点。尽管如此,他毕竟把手指按在了一个重要的“原理”上——那就是:D-模理论(我自己更愿意称之为“晶体模”\fnote{744}{(**)}{这个“晶体”术语的缘由(显而易见)——它反映了一种更内在的D-模视野(我的学生们曾从我这里学到,却早已忘却)——见“我的孤儿”一注(n\textsuperscript{o}~46)中的评注(尤其第~179~页),以及子注~n\textsuperscript{o}~46$_{4}$(第~188~页)(x)。至于Mebkhout的哲学与我六十年代末提炼的晶体瑜伽之间那些显而易见的联系所遭遇的“健康官能的堵塞”,见“装神弄鬼”一注(n\textsuperscript{o}~85$'$,第~350--351~页)。\par(x)(5月24日)另见“五张照片(晶体与D-模)”一注(n\textsuperscript{o}~171(ix))。})提供着一个\origpage{1039}“公分母”,足以“统罩”起\dmark{}那些现象(尤其是对偶现象)——它们出现于离散上同调与凝聚上同调之中。沿这条轨道,倘若当时有一位“身在局中”的人给予鼓励,再加上最低限度的数学本能\fnote{745}{(*)}{这并不是说我的上同调前学生们缺乏“最低限度的数学本能”——否则他们中间当初谁也不可能随我做出那份好工作。但这种本能被埋葬师父的综合征带偏了道,或者堵死了。}与善意,那么毫无疑问,他会在随后三四年里,为(至少)特征零的代数几何发展出一套完整的六种运算形式主义,给出那同一形式主义(诚然备受排斥)在超越语境中那种形态的一个纯代数的、忠实的“范式”,用于代数可构的$\mathsf{C}$-向量空间层。

\noindent 他深深感到自己刚刚发现了某种重要的东西,于是满心欢喜地请求并获准与他的恩人作一次会谈,向他陈述自己的结果。这恰恰是\textbf{对}我十年或十二年前向Verdier提出、而他似乎毫无所动\fnote{746}{(**)}{(6月5日)关于这一点,见“祖先”一注(n\textsuperscript{o}~171(i)),尤其第~946~页的脚注(*)。}的那个问题的\textbf{回答}——他甚至很可能已把它忘得一干二净。无论如何,他对这个从天而降、正在做着他——Verdier——早已一笔勾销的事情的年轻人的善意,已经耗尽了。他连听也不想听Zoghman对定理的来龙去脉及其证明所作的解说。他(礼貌地)让他领会的要旨是:他Verdier早已不再相信圣诞老人,这个年轻人最好还是把东西收起来吧。

\noindent 奇怪得很,Zoghman周围竟\textbf{没有一个人}“接得住”这一结果\fnote{747}{(***)}{(6月3日)这里有过一场误会。如“三座路标——或天真”一注(n\textsuperscript{o}~171(x),第~1026~页)所说,这定理常常能让一位偶然的交谈者惊叹不已。但到目前为止,这惊叹似乎始终停留在柏拉图式的层面——定理并没有变成一件工具,一种人们熟知、随手使用乃至不加思索的东西。这肯定与下面的事实有关:为这结果显而易见的美而欣喜的人,从来不是那些“定调子”的人,不是那些决定什么“重要”、什么算“瞎拼凑”的人。(而在眼下这年头,“昨天”的瞎拼凑摇身变成“今天”的奶油甜点,可不是什么稀罕事……)——在他4月22日的评注里,Zoghman给我写道:“……面对这一定理,人们心里有一种别扭。有人暗地里嫉妒它。但鼓励过它的人实在寥寥,恰恰相反。”}——\origpage{1040}无疑这太像六十年代的“Grothendieck做派”了\dmark{},谢天谢地,如今人们早已超越了那一套!也许,两年前,我正是他所遇到的头一个体会到这一结果的重要性、以及它所孕育的那门新“哲学”的人——那就是在形形色色簇的上同调(先从代数的与解析的做起)之中,“离散”侧面与“微分”(或“解析”)侧面之间一场宏大综合的哲学。这个构成他论文其中一章的定理,最后于1982年刊于 Mathématica Scandinavica(第~50~卷,第~25--43~页)。同一篇文章曾投给 Annals of Mathematics,后者让这个年轻的妄自尊大者明白:他不够在这么一份有身份的刊物上发表的档次。

% ---------------- p.798 ----------------
\noindent 甚至直到今天,这定理在上流社会里通常仍遭无视或蔑视,尽管它已然孕育着那门新哲学——经由好上帝定理(别名Mebkhout),它提供了代数簇上同调一场壮观更新的手段。可是“人人”——包括我的上同调前学生们(尽管有一天我曾认为他们天赋一种健全的数学本能)——都一窝蜂扑向了那道新的“奶油甜点”,即某一件强大的工具(而“人人”却又装模作样,只肯以影射或迂回说法提名道姓\origpage{1041},比如“可构层与完整微分系统之间的关系”,或者“那本该正常地在这一套注里占有一席之地的东西”\fnote{748}{(*)}{这是凭记忆引自Beilinson--Bernstein-Deligne那篇“值得纪念的文章”(由Deligne执笔)——关于它,见“光荣之日”一注(n\textsuperscript{o}~171(iv))。关于这个不愧流传后世(作为一份提醒、也作为一纸警告……)的迂回说法及其来龙去脉,见“牵线人”一注(n\textsuperscript{o}~75$'$)。前面那处引文(“可构层与完整微分系统之间的关系”)摘自Beilinson--Bernstein(同在1981年)的一篇文章,下一子注(“黑手党”,n\textsuperscript{o}~171$_{2}$)将谈到它;在那里,读者还有幸结识Brylinski-Kashiwara为这门风格的开花结果所作的贡献——它服务于同一桩骗局。}……),扑向那“最新潮的玩意儿”(相交上同调),而\dmark{}那幅\textbf{视野}——它使这件工具有所由出——却依旧像先前一样无人理会,而这两者的父亲则被当作听差使唤。

\noindent 这里的情形,与我关于topos、导出范畴、六种运算、上同调系数、以及更进一步关于Motif的那幅宏大统摄视野的情形如出一辙。平展上同调与晶体上同调这类工具,正是从这幅视野中诞生的;而这幅视野本身——今天“人人”转动曲柄般使用的那些工具正是从中诞生的——在我离去那一天还强旺地活着,次日便被埋葬了。而我清楚地看到:我离去时正值蓬勃的这个辉煌主题\fnote{749}{(*)}{关于这种“阴郁的停滞”的印象,我第一次谈起是在“对一份遗产的拒绝——或一个矛盾的价格”一注(承“我的孤儿”之后,n\textsuperscript{o}~47,第~195~页)的末尾。自写那条注以来的一年里,这份印象只有加深的份儿——带着一条与我在所引注的子注~n\textsuperscript{o}~47$_{3}$中所表达的根本相同的保留:即Deligne关于Weil诸猜想(WeilI与II)的工作,以及“好上帝定理”(它把好上帝连同他的仆人Zoghman一并清洗出去)与相交上同调的“大冲刺”之后的新出发。但这些局部性的成功,在我看来,与那些我熟知的、此后“入驻”这个“辉煌主题”的人们的卓越——乃至超群——的手段,全然不相称;自离去至今十五年已经过去。也与我遗赠的那些主导观念的丰富与强劲全然不相称——而今天我重又见到它们时,它们已气息奄奄……},十五年后所呈现的惊人停滞,并非源于智力手段或天赋的匮乏(在我认识得那么好又那么坏的那些人中间,不止一人天赋出众),而是源于掘墓人的心性,或肆无忌惮的裙带之风,或两者兼有——那是一些与那种使人认出、并使人找到简单而本质之事物的天真无邪截然对立的心性。

\noindent 为发展他那门新哲学,Mebkhout汲取了导出范畴与六种运算的精神——那时导出范畴正被当作Grothendieck式的花架子,而他连听人说出“六种运算”这个名字的机会都不曾有过。今天,随着新工具的问世——它与导出范畴密不可分——人们大张旗鼓地把后者重新发掘出来,却对两个人的名字讳莫如深:一个是多年孤独劳作、把它们从虚无中提炼出来的人,另一个是从中汲取灵感、终于使一门把拓扑、复分析与代数几何连为一体的新系数理论开花结果的人——后者同样形单影只。

% ---------------- p.799 ----------------
\dmark{}Deligne们、Verdier及其同伙一窝蜂扑向这些簇新的时兴货,一面嚷嚷(那份\origpage{1042}照例必得、冠冕堂皇的谨慎,这还用说):“是我,是我!”。他们中间,还没有一个能在自己身上找到那份勇气与对自己的忠诚,把一份视野在孤独中养熟,沉重地怀上它几个月、几年,远离掌声——那时只有他们自己看得见,而他们所见之物在这世上无法与任何第二个人分享。

\noindent 可我扯远了,该回到我对\textbf{一个愿景之萌发}的讲述了。就在1976年——Mebkhout证明那条“统罩”Poincaré对偶与Serre对偶的对偶定理的同一年——他便得出了三个范畴等价的观念;这三个范畴分别体现同一实在、同一类型对象的“拓扑”侧面、“代数”侧面与(超越的)“解析”侧面。以一门“上同调系数”的普遍理论的眼光\fnote{750}{(*)}{“\textbf{系数类型}”各各不同的这一观念——其中每一个在我面前都显现为六种运算形式主义(与双对偶性)的一个特殊化身,或多或少切近于一切系数类型中最精细的那个、“绝对的”或“普遍的”类型,即“\textbf{Motif}”——或许正是引导我贯穿整个六十年代、尤其自1963年以降,在发展我对代数簇及其他对象的上同调视野时的那个最主要的主导观念。这一观念在我身上的力量,在我专门回顾我的著作及其在时风手中的坎坷遭遇的第一条注(“我的孤儿”,n\textsuperscript{o}~46)里,已清晰可见。在葬礼反思的各处,我以执着的口吻回到它,尤其在“墓上的旋律——或自负”与“工地的巡礼——或工具与视野”(n\textsuperscript{o}~167、178)两条注中。这也是一个首要的数学主题——在被我的上同调前学生们连同一场时风埋葬的那些主题之中——我打算在《收获与播种》之后加以发展,为它在数学思考中取得它应得的位置。\par 奇怪的是,我上同调著作的这一核心主导观念,以及表达它的那个(归根结底十分简单的)代数-范畴结构,在文献中从来不曾被明确陈述过,连我自己在六十年代也不曾(x)。它在我笔下的著作中见于字里行间,而它的传布主要通过口头的交流。在我心里,这原是理所当然:我的学生中间总会有一个,舍得花上所需的几天或几个星期,把这组观念系统地陈述出来——那时我正全力投入EGA与SGA的基础工作。\par 事后回望,我更清楚地意识到:一些非形式的文字(在此情形哪怕只有几页,也不必追求精确而系统的表述)——它们恰好把那些难得提名道姓、却隐藏在往往面目技术性的文字背后的“主导观念”传达出来——这样的文字对于给研究者指引方向、对于不时给一部容易被自己的技术性闷得透不过气的文献送进一口新鲜空气,是何等重要。说到这里,Zoghman对我说过:他在我的笔下找到的寥寥几段这样的文字,曾给他很大的帮助。其中,他近来还特别向我举出我附在Hartshorne那卷《留数与对偶》(原书拼作 \emph{Résignes and duality}——那卷书主要阐述我五十年代下半期在凝聚框架中发展的六种运算形式主义)卷首的几行引言。如今我掂量得出:倘若当年我肯费神在里面加进哪怕一两页非形式的文字,讲一讲“六种运算的瑜伽”、并强调它作为一条无处不在的引线在建造那些尚待诞生的上同调理论时的意义,这几行引言本会对他有用得多……\par(x)(5月24日与6月1日)这些行文写成之后,事情显明:早在SGA~5口头讨论班的最开头(在我的第二讲里),我就曾花大力气,长篇阔论地展开六种运算的“抽象”公式库——它后来主宰了讨论班的全部进程。(关于这一点,见“祖先”一注(n\textsuperscript{o}~171(i))第~942~页5月8日的脚注(*)。)此外,在整个口头讨论班期间,我不曾放过一个机会,不断指及我所发展的那套上同调形式主义的无所不在——它在原则上对“$\ell$-adic系数”之外的形形色色其他“系数”类型一概有效。而Illusie,却精心地把谈论六种运算形式主义的那一整讲、以及任何一丝超出讨论班主要对象这一特定语境的“上同调系数”视野的暗示,统统从屠杀版里剜了出去。\par 关于此事另见“死去的书页”一注(n\textsuperscript{o}~171(xii)),以及“无用的细节”一注(n\textsuperscript{o}~171(v)),(b)部分(“无所事事的机器……”)。},\origpage{1043}我把这些对象称作“德拉姆—\dmark{}Mebkhout系数”\fnote{751}{(*)}{(5月30日)在(后来写成的)“五张照片(晶体与D-模)”一注(n\textsuperscript{o}~171(ix))中,我采用了一套略有不同的术语,把“同一类型对象”径直称作“德拉姆系数”——在这里我们将给出它的三幅\textbf{描述}(或三张不同的“照片”)。其中两幅将分别冠以“de~Rham--Mebkhout系数”(或简称“Mebkhout系数”)“无穷阶”与“有限阶”之名。}。设$X$是一个光滑解析空间\fnote{752}{(**)}{(5月30日)在这几条注的初稿里,我因偏爱“代数几何”观点而忘乎所以,曾假定$X$是$\mathsf{C}$上的代数\textbf{簇}。这并不符合Mebkhout当初置身其中的框架;更何况,这还害我陈述了一个关于$D^{\infty}$-模复形的“好上帝定理”的变体——它只有当假定$X$为真(固有)时才如其所是地为真。于是我的心里出过一些误会,Mebkhout不得不温言提醒我注意。在把这些页重新誊清时,我作了应当作的订正。},那么一方面有 % ---------------- p.800 ----------------
由$X$上“可构”$\mathsf{C}$-向量空间复形组成的(“导出”)范畴 $\mathrm{Cons}^{*}(X,\mathbb{C})$,或径记作 $\mathrm{Cons}^{*}(X)$(\textbf{“拓扑”侧面});再是由那些具有凝聚上同调层的 $D^{\infty}_{X}$-模组成的复形构成的范畴\fnote{753}{(***)}{关于模理论与D-模理论的定义及最初那些连锁的小事实,读者可参看“五张照片(晶体与D-模)”一注(n\textsuperscript{o}~171(ix)),尤其是(a)与(b)两部分(“‘德拉姆系数’图册”与“好上帝公式”)。},\origpage{1044}\dmark{}它是无穷阶微分算子复形的推广,我记作 $\mathrm{DRM}^{*}_{\infty}(X)$(超越的\textbf{“解析”}侧面);最后是由那些具有凝聚上同调层的 $D_{X}$-模组成的复形构成的范畴——它是寻常微分算子复形(有限阶)的推广,我记作 $\mathrm{DRM}^{*}(X)$(\textbf{“代数”}侧面)。有一个重言式的标量扩张函子,从凝聚环 $D_{X}$ 到环 $D^{\infty}_{X}$:

\begin{quote}
$i\;:\;\mathrm{DRM}^{*}(X)\longrightarrow \mathrm{DRM}^{*}_{\infty}(X)$
\end{quote}

\noindent 它嵌入一个(本质上交换的)函子图表:

\begin{equation*}
\begin{tikzcd}[row sep=2.5em, column sep=2em]
\mathrm{DRM}^{*}(X) \arrow[rr,"i"] \arrow[dr,"m"'] & &
\mathrm{DRM}^{*}_{\infty}(X) \arrow[dl,"m_{\infty}"] \\
& \mathrm{Cons}^{*}(X) &
\end{tikzcd}
\tag{1}
\end{equation*}

\noindent 其中两条斜箭头是“相结合的 De Rham 复形”箭头\fnote{754}{(*)}{(5月24日)见已引的“五张照片……”一注(n\textsuperscript{o}~171(ix)),(a)部分。}——它不是别的,正是 $\mathrm{RHom}_{D}(\mathrm{Sp}_{*},\,\cdot\,)$,这里 $D=D_{X}$ 或 $D^{\infty}_{X}$,而 $\mathrm{Sp}_{*}$ 是 $\mathcal{O}_{X}$ 经由局部自由 D-模的“Spencer 消解”。

\noindent 那些竖直箭头的存在,来自“Kashiwara可构性定理”——它蕴含:与一个完整D-模复形相结合的 De Rham 复形,其上同调层是解析可构的。Kashiwara早在1975年就证明了这个重要的定理\fnote{755}{(**)}{Masaki Kashiwara,\emph{On the maximally overdetermined System of linear differential équations}, I,Publ. RIMS, Kyoto university 10 (1975),第~563--579~页。},但却是从一个全然不同的着眼点。他处理的是单独一个完整D-模,取它的 De Rham 复形,证明其上同调是可构的。直到1979年9月、直到好上帝定理引发的“大冲刺”为止,他同上流社会里任何别人一样,从不曾以导出范畴的精神工作;把那些竖直箭头写进图表(1)这个念头,压根不曾浮上任何人的心头!

\noindent 一旦(1)的三条箭头作为导出范畴之间的箭头写下\fnote{756}{(***)}{严格说来,更正确的说法或许是:这是寻常意义下导出范畴的(由“可构性”或凝聚性、完整性、正则性诸条件定义的)全子范畴。},问题便随之而来:它们果真是范畴等价吗?Mebkhout自1976年起就对这一点深信不疑。他的信念来自\dmark{}\origpage{1045}编列的一张约莫十个“典型的”例子的清单(复现于他与 Lê Dung Trang 合写的那篇综述文章里\fnote{757}{(*)}{Lê Dung Trang 与 Zoghman Mebkhout,\emph{Introduction to linear differential Systems},Proc. of Symposia in Pure Mathematics, Vol.~40 (1983), part 2,第~31--63~页。Zoghman向我推荐了这篇短文,说它是文献中对他自1976年以来所发展的那门哲学的最好导引。在该文书目里还可以找到一份(完备的?)Mebkhout关于这一主题的出版物清单,至少截至1983年。}),都是那些可以称作“初等的”可构$\mathsf{C}$-向量空间层——它们同时也是平展上同调理论中人人熟悉的、层的“拆卸”作业里不断出现的那类对象。就在这决定性的1976年,对其中每一个层,他都成功地在$D_{X}$(“\textbf{代数}”)上、也在$D^{\infty}_{X}$(“\textbf{解析}”)上构造出一个非同寻常的完整复形——以六种运算的眼光看,它具有非常简单的代数或解析的上同调意义,而它的 De Rham 复形恰是那个层。值得注意:他从出发时手里拿的是一个可构层、而非一个层复形,可在相当一批情形里,生出它的那个完整复形却绝非仅有一个上同调层。这向他清楚地表明:循着“六种运算”的精神(他并不知道这个名字……),倘若真有等价,它绝不可能从模层范畴(在$\mathsf{C}$上或在$D$上)彼此之间的等价推演出来;它唯有在过渡到导出范畴时才取得意义。

% ---------------- p.801 ----------------
\noindent 在我看来一清二楚的是:此时的\textbf{创造行为},就在于看出并写下那两条\textbf{显而易见的}箭头 $m$ 与 $m_{\infty}$——人人却都不屑一写,就在于提出那个“再幼稚不过”的问题:它们会不会碰巧正是范畴等价,从而给出可构$\mathsf{C}$-向量空间层(或层复形)这一拓扑观念的一个微分代数的、以及另一个微分解析的解释。\textbf{问题在那里,对这一问题之关键地位的清晰意识也在那里}——意识到它的分量;而由此,作为理所当然之事,一种把这问题\textbf{担在自己肩上}、要把它坚持到底的内在姿态。以那些“典型的”或“初等的”例子所作的预备性“实验”,便是朝这个方向迈出的第一步。

\noindent 那是婴孩般的一步,也是本质的一步——唯有知道如何独处的人才迈得出这一步。一旦迈出,我的上同调学生们随便哪一个,凭着在我身边、在SGA~4与SGA~5中学到的拆卸与消解技术\origpage{1046},都能在短短几天\dmark{}之内把它证明出来,或者几个星期——只消他“接得住”,当然,只要他能(像Mebkhout那样用整个内脏)感到这个问题的意义、它的\textbf{实质}。可他们中间没有一个人,连那个为提炼一套超越“六种运算”主导观念、\textbf{超越}其上——而它正是把连续系数与离散系数联结起来所仍欠缺的——的统摄视野而\textbf{宣告弃权}的Deligne也包括在内\fnote{758}{(*)}{(6月5日)重读时,这个提法在我看来显得仓促,而且有点“不着边际”。实际上,我的“六种运算这一主导观念”与一门“系数哲学”密不可分——后者(至少从1966年起就十分明确地)预见了一套“德拉姆系数理论”(与我的晶体观念血肉相连),它具有与$\ell$-adic系数理论相同的本质形式性质,并与后者(当$\ell$变动时)共同构成同一个终极对象类型——“Motif”——的种种不同的“实现”。Mebkhout完成于1972至1980年间的工作,在我看来是朝这一直觉的实现迈出的第一大步——迈进这一步的条件,实际上至少从1966年起、随着晶体瑜伽的启动就已完全成熟;那一年,一套德拉姆系数理论的问题(至少在我的心中)已被清晰地提出。假如我的上同调学生们中间——早自六十年代——无一人迈出这一步,在我看来这主要应归咎于种种堵塞自发创造力的机制;而那种创造力,他们中间没有一个人是缺乏的。关于这一点,见“……与羁绊”一注(n\textsuperscript{o}~171(viii))。},没有一个能看出Mebkhout的思想、这位依旧浑身透着我那股气息的无名之辈的思想的分量——尽管那分量一望即知……

\noindent 至于那位“无名之辈”,凭着他自己的手段与他自己的阅读,向自己提出范畴等价的问题,在他看来(而且也确实)是天底下最显而易见、最孩子气的事,或者由此达到确信:这些箭头货真价实地就是范畴等价。相反,由于缺乏经验、缺乏比他更有经验的长者的鼓励,那证明在他眼里成了一座大山,长久以来显得完全高不可攀。

\noindent 然而仅仅一年半之后,他便找到了一个证明——起初是对箭头$m_{\infty}$,在1978年3月。他对我说,在心理上,我的代数与超越 De Rham 上同调之间的比较定理曾给他很大帮助,把他引上了证明之路。出于一个我未曾很好理解的理由,他还把自己的定理(即:函子$m$——人们称它“好上帝”函子,恰恰是为了不提Mebkhout……——是一个等价)\origpage{1047}视为我的比较定理的一个“推广”。从这一刻起,他也知道\dmark{}自己手头已有必需的工具(连同Hironaka的消解技术),足以处理$m$的情形——对一个像我这样的代数几何学家来说,那才是远为有趣的情形。而他,身为分析家,起初倾注心力的本是函子$m_{\infty}$的情形——那原是他的偏好\fnote{759}{(*)}{(5月24日)另一个也许更强劲的缘由是:在$D^{\infty}$-模的情形,他手头有一条绝妙的反演公式——关于这一点,见“五张照片”一注(n\textsuperscript{o}~171(ix)),(b)部分“好上帝公式”。}。对那个在他看来多少有些次要的问题,他直到论文答辩之后才重新拾起,并在一个月之后(1979年3月)证明:函子$m$(就是那个今天人人以迂回说法使用、却从不肯写出来——免得提名道姓一位不可提名道姓的作者的……函子)确实是一个范畴等价\fnote{760}{(**)}{Mebkhout把“$m$是一个等价”这一事实的正式成文的证明(其原理与“解析的”好上帝函子$m_{\infty}$的那个证明相同),是两年之后、1980年末才写出的。这份证明陈述于两篇相继的文章中的第二篇(第一篇处理解析的好上帝函子$m_{\infty}$,重述其论文),即“一个范畴等价”与“另一个范畴等价”,载 Compositio Mathematica 51 (1984),第~51--62~页与第~63--88~页。(手稿1981年6月10日收到。)但早在1969年3月[原书如此,当为1979年]起、在其后的岁月里,凡有机会,他便传播这一结果(连同关于函子$m_{\infty}$的那个),尤其是同年6月就告知了Deligne。\par 我想,由于他极度的孤立,加上他作为分析家的“眼镜”,他不曾意识到:真正会引起Deligne之流兴趣的,恰恰是代数的好上帝函子——因为它在拓扑与代数几何之间搭起一座“桥”(算术尚待到来,而看来我是第一个、并且是唯一一个预见到它的人……),其分量堪比平展上同调工具所提供的那一座。否则,他定会立刻把它写成正式的文字并火速发表——尤其考虑到他误入的那个古怪圈子(他当时还一无所知……)的种种行规。不过,他在1980年3月的头一场噩运(与Kashiwara打交道),本该让他心里起疑(x)。\par 恰恰也是在同一个3月,Mebkhout在法国科学院《通报》上发表了关于“Riemann--Hilbert问题”的一篇短注(第~290~卷,1980年3月3日,Série A,第~415~页),其中他陈述了论文里(关于$m_{\infty}$)的等价定理,并谨慎地声明:“人们希望证明——以对偶定理[7]所用的上同调下降法为法——函子$S$[即我所谓的$m$]乃至$T$[即我所谓的$i$]也是范畴等价”。实际上,他的证明表明:它们是“在$X$上局部地”的等价;单是这一点,已经蕴含了著名的Kawai-Kashiwara定理(下一子注将谈到它),即:标量扩张函子$i$在完整正则$D_{X}$-模范畴与完整$D^{\infty}_{X}$-模范畴之间诱导一个等价。我顺带指出:Mebkhout的最终结果,即便应用于单个的模(而非模的复形),也远为有力,因为他同时断言:由“标量扩张”函子导出的那些典范箭头

$\mathrm{Ext}^{n}_{D_{X}}(M,N)\longrightarrow \mathrm{Ext}^{n}_{D^{\infty}_{X}}(M_{\infty},N_{\infty})$

也是同构(而不止于$n=0$)。\par(x)(5月25日)在4月24日的一封信里,Mebkhout还向我说明:“要知道,论文答辩之后我松了一口气。整整四年,我一直处在极大的紧张之中。”}。由此立刻得出:那个从“代数的”(他当时对它还不曾有多大兴趣)通往“解析的”(超越的)的“换环”函子$i$,同样是一个等价。

\asterisms

% ---------------- p.802 ----------------
\dmark{}1978年3月,Mebkhout与他的“恩人”Verdier有了第三次会谈——此前他们已两年未见。他向他讲解了(未来的)“好上帝定理”的来龙去脉;\origpage{1048}他谦逊地(这一谦逊害了他!)称之为“Riemann--Hilbert等价”。回想起来,Mebkhout确信:自己的解说多半从Verdier的头顶上飘了过去。可以肯定的是:Verdier压根没有意识到,他这位“受保护者”刚向他提交的是一些值得驻足的观念。他没有向身边任何人提起,连Deligne也没有——Deligne直到一年多以后,才从Mebkhout口中得知好上帝定理(连同那条号称“Poincaré--Serre-Verdier”的对偶定理——三年前这位Verdier本还死活不肯相信它……),那是1979年6月的Bourbaki讨论班(答辩之后四个月)。无论如何,Verdier开了绿灯,让Mebkhout把他的结果作为国家博士学位论文提出,并且同意组建并主持答辩委员会。答辩之所以推迟了一年才举行,是因为那个有名的“巴黎地区各大学论文委员会”办事的拖拉(对这个机构,Verdier可是珍爱如眼珠……)。

\noindent 正如我在一条先前的注里说过的\fnote{761}{(*)}{见“……与意外之喜”一注(n\textsuperscript{o}~171(iii))。},答辩在一种人人漠然的气氛中进行。Mebkhout固然可以把论文左一份右一份地寄出去,它依然默默无闻——连那块大部头的一份收执,都没有人屈尊回一张。

\noindent Mebkhout却没有被击倒。尽管种种迹象指向反面,他仍感到自己属于一个“家庭”——毕竟是一群做同一类数学的人:那类数学,他很大程度上是\origpage{1049}凭读我的著作学来的,而更多的是凭对这几部著作中的某种\textbf{精神}持开放、\dmark{}倾听的姿态学来的\fnote{762}{(*)}{人们不禁要问(或者要问我):我的著作里那种如此独特的著名“精神”究竟是什么,它竟启发了我这位“身后的学生”Zoghman~Mebkhout,却遭我其余所有学生——Deligne打头——以及一支亦步亦趋的时风的“排斥”。倘若我试着为这种精神找一份谱系(在我对数学史的那点支离破碎的知识允许的范围内),我会说:它属于\textbf{Galois、Riemann、Hilbert}的谱系。倘若我试着以心灵中种种作用力的动力学来界定它,我会说:这是一种通过“阴”“阳”两种创造力之和谐平衡而显自身的精神,其“基调”或“主导音”是\textbf{阴},是“女性的”。对数学中、乃至一般地对世界之发现中这一进路的更详尽的描述,可在“上涨的海”“九个月与五分钟”“阴的葬礼(阳葬阴(4))”(n\textsuperscript{o}~122、123、124)诸注的反思中逐步见出;这一反思又在“兄弟与配偶——或双重签名”“阴,作为仆人——与新主人”“阴,作为仆人(2)——与慷慨”(n\textsuperscript{o}~134、135、136)诸注中得到接续。至于当代世界对这种“精神”的某些“本能的”拒斥机制,见“天赐的情势——或封神”与“否认(1)——或提醒”两注(n\textsuperscript{o}~151、152)。}。他显然还不曾意识到——至少在意识的层面上——这种精神恰恰早已被那些构成他自信已然加入的“家庭”的人们所排斥;而对那些踩着厚绒地毯步入数学的翩翩君子来说,他不过是个拖拖拉拉的落魄鬼、一个闯入者。

% ---------------- p.803 ----------------
\subsubsection*{c2.~黑手党}

\tocsubsub{c2.~黑手党}

\notehead{171$_{2}$}{}

\noindent (4月15--17日)

\subsubsection*{(a) 家庭合影上的阴影}

\tocsubsub{(a) 家庭合影上的阴影}

\noindent 可我们这位朋友Zoghman,对这一切还蒙在鼓里;他虽则孤立,却并不不幸。自1973年起,他有幸在奥尔良谋得一个助教的职位,这让他得以安安静静做他感兴趣的数学——哪怕眼下感兴趣的只有他一个。他继续住在巴黎地区,继续出入讨论班,继续熟悉文献……

\noindent 倘若他稍稍琢磨过这件事,他本会发觉:在这个佯装不知有他、而他自信身属其中的“家庭”里,一切并不那么美满。凭着阅读我的著作,他终于多少觉察到:曾对他如天降甘露的那份“好的参考文献”,至少有相当一部分,根本不出自他“恩人”Verdier的自酿。可构性概念早在1963年就在SGA~4里得到长篇展开,比Verdier在那篇文章里装作发明它早了十二年。随着SGA~5于1977年问世——哪怕是以Illusie屠杀版的形式——事情变得清楚了:那著名的、关于解析或代数可构的$\mathsf{C}$-向量空间层复形的“Verdier双对偶”,被原封不动地\dmark{}\origpage{1050}抄自SGA~5的第一讲——那个名字古怪的书卷“SGA~4$\frac{1}{2}$”所指涉的正是这一讲,其说法是:“若干补充见 SGA~5 I”\fnote{763}{(*)}{关于这个叫人哭笑不得的委婉说法——它意在(这一回由他Deligne)把那条倒霉的双对偶定理据为己有——见“特洛伊木马”子注(n\textsuperscript{o}~169$_{3}$)第~872~页的脚注(**)。}!在同一部古怪的书卷里——它的作者乐于以十足的高傲俯视环绕着它的姊妹卷SGA~4与SGA~5——他还可能读到关于结合于一个闭链的上同调类的一讲;那一讲(不知何故)被人从那部“技术性离题”的书卷SGA~5(号称更晚……)中卸了下来;他由此同时可以看出:他那恩人的文章得名的那个主题的上同调侧面(同调侧面的对偶),同样是抄自SGA~5。然而,对“好的参考文献”里的这三个主题\fnote{764}{(**)}{这里说的是那“三个主题”:可构性,可构层的双对偶性,结合于一个闭链的上同调类(与同调类)。},却没有任何一处提到我的这个人或SGA~5……

% ---------------- p.804 ----------------
\noindent 他固然还无从知道:Verdier那篇文章剩下的部分(五十页里除去三页)是从我关于平展同调形式主义、关于结合于代数闭链的同调类的报告中“抄”来的——那些报告(仿佛有约在先)已从那部声名可哀的Illusie版中消失,连一丝提到它们存在过的影子都没有。但他手头的那点事实,对一个知情而警醒的人来说,确实已经绰绰有余。总而言之,这情形同我十年前翻阅Deligne关于谱序列退化那篇文章时的处境一模一样:在那篇文章里,他既把最初的动机、整个权的瑜伽(连同我这卑微之人的角色)掉包藏匿,也把Blanchard思想的贡献藏匿起来——那贡献恰恰是把“蛮横的”Lefschetz定理用于纤维之上\fnote{765}{(***)}{详情见“驱逐”一注(n\textsuperscript{o}~63)的开头,以及该注第~233~页的脚注(**)。}。像当年的我一样,Zoghman那时只得把对一件不快之事的清明感知压下去,对自己说(在此情形):这想必是师父与学生们之间惯常的一种“默契”——当学生们把他直接传授的观念、技术、结果当作自己的东西来呈现时,师父便睁一只眼闭一只眼\fnote{766}{(****)}{(5月30日)而且顺带还温言称他是个江湖骗子……}。正如在\dmark{}这类情形中常有的那样,这种解释(它正合Zoghman的心意)本就不乏真实的成分——而且还不仅如此。\origpage{1051}我不止一次确实曾是这类暧昧局面的当事人。(但也属实:在我离去之前,事情从不曾发展到这个地步——师父的著作成了遗骸,人们厚着脸皮瓜分它的碎片……)

\noindent 再说,在那个由所有对簇上同调感兴趣的人组成的扩大的家庭里——包括Sato学派的日本人——情形也不怎么美妙。那位Kashiwara,他1975年的可构造性定理曾为定义“好上帝函子”立下汗马功劳,却也装模作样,要把那些倒霉的可构造层——忽而间人人竞相争抢的东西!——的发明权据为己有。为着这份事业,他在所引文章的第2节里给它们重新施洗为“finitistic sheaves”,在该节里或多或少逐字照搬SGA~4中关于此主题的展开。据我从多方听到的,Sato学派熟悉我的上同调著作,哪怕他们引用我时吝啬得很\fnote{767}{(*)}{Mebkhout就此给我写道(1985年4月24日):“我在Sato的日本学派那里见到的唯一提到你的地方,涉及EGA~III第0章——而他们却在毫无顾忌地从你的著作里汲取灵感。”};很难相信Kashiwara不知道可构造性概念——至少在平展语境中,在那里它是全部理论的中心有限性概念。不言而喻,第二年Verdier对“finitiste”(原文如此)概念照样不提Kashiwara,正如他对某位亡故者、某个讨论班只字不提\fnote{768}{(**)}{仿佛有约在先:这个讨论班(SGA~5)恰恰就是(与SGA~4一道)经我的上同调学生们共同商定、并按他们头目Deligne的说法,注定要被“忘掉”的那一个(多亏了由他亲笔操刀的那份“电锯摘要”的出版……)。}。两人都是上流社会中人,说不定还同属一个“家庭”——可一到作者虚荣这牛排面前,各人只顾各人的……\fnote{769}{(***)}{(5月24日)Mebkhout提醒我,我在这里把画面涂得略黑了些。Verdier对Kashiwara的文章一无所知,对完整性(holonomie)概念同样一无所知——那是Mebkhout在1976年与Verdier的那次“会谈”中教给他的。(这发生在“好的参考文献”刊行之前——它似乎刊于1976年末;不过按情理,既然他知道自己与他的同行都在从同一处未提名的源头“抄”,也就不能指望他会去引Kashiwara……)反过来,Kashiwara不知道“好的参考文献”、也不知道我的双对偶定理(它在那里顶着Verdier的父名)——是Mebkhout在1978年1月把这两样东西连同他论文第三章的结果一并告知他的。而后,那些结果就在已引的Kashiwara--Kawai的文章里被毫无顾忌地(而且几乎不带证明地)据为己有——关于这一点,见“五张照片(晶体与D-模)”一注(n\textsuperscript{o}~171(ix)),尤其第~1005~页。Kashiwara不知道离散系数的双对偶定理这件事,连同这里或那里记录的许多别的迹象,表明他离Mebkhout那门直接源于我著作的对偶哲学有多远。}

% ---------------- p.805 ----------------
\dmark{}我看,对Zoghman来说\origpage{1052},宁可对自己说,某个他从未谋面的日本人\fnote{770}{(*)}{(5月24日)他终究见过他们一面,这些大名鼎鼎的日本人!Mebkhout就此给我写道(1985年4月22日):“1972年,Sato学派全员出动,来开一个关于超函数的会议。他们把自己的方法藏得很好。很长时间里,他们的成果都无法企及。围绕这个学派有某种神话,如今Kashiwara之所以能随心所欲,正赖于此。”\par(6月4日)应当说,倘若果真(如Mebkhout在此所暗示的)是Sato学派首创了那种为求支配而故示幽暗的手法,那么这手法在太平洋此岸也不乏徒众——如今他们比起师父来毫不逊色!而把反常层会议那场令人难以置信的蒙骗局操办起来的,恰恰是他们,而绝非Kashiwara之流:在这场蒙骗局里,Kashiwara被当作一枚方便的“小卒”用来铺路——随后便遭丢弃……}准是个“骗子”,也不愿不得不在下述事实面前认账:那是几位显赫的长辈——其中一位,于他犹如一位强有力而高不可攀的父亲、一位恩人;是他有机会在讨论班里晤对的长辈,而且他还荣幸地同他们以“你”相称(自Bourbaki时代以来,法国数学圈的惯常做派)。

\subsubsection*{(b) 最初的麻烦——或太平洋对岸的卡伊德}

\tocsubsub{(b) 最初的麻烦——或太平洋对岸的卡伊德}

\noindent 说来吊诡,Zoghman的麻烦,恰恰是从某个圈子开始察觉他随一整门哲学(诚然是一种看来分明已过时的哲学……)带来的某一件工具之威力那一天开始的。1979年6月,他已把这一点告知Deligne;Deligne聚精会神地听了他关于对偶定理的解说,而(可想而知)更加聚精会神地听了关于好上帝定理的解说。他甚至十分和蔼地对他说,他读过论文的引言,并且觉得这项工作里必有精彩的数学\fnote{771}{(**)}{(6月3日)前一年,在雷恩的$p$-adic分析讨论会上,Mebkhout已经从Illusie口中领受过一个同样无谓的恭维。关于这一点,见“掠夺的空白授权”一注(n\textsuperscript{o}~171$_{4}$),第~1091~页(尤其同页脚注(**))。}。那一天,Zoghman的日子真美——可惜好景不长。

\dmark{}同年,1979年9月,他参加了莱苏什讨论会\fnote{772}{(*)}{莱苏什讨论会(1979年9月1--13日)的《论文集》刊于 Lecture Notes in Physics n\textsuperscript{o}~126 (1980),Springer Verlag。这部《论文集》里,既有Mebkhout的报告《论Hilbert--Riemann问题》——它以完美清晰的方式陈述了他的全部哲学(我会称之为“德拉姆系数”的哲学),证明都有文献支撑——也有Kashiwara与Kawai提出的报告。任何善意的读者,把两篇文章一比即可验证:在这两位作者的文章里,既没有一丝这类哲学的萌芽,也没有半句提到好上帝定理之类的暗示。\par(6月4日)在其4月22日的评注信里,Mebkhout就前一年(1978年8月)举行的赫尔辛基国际数学家大会表达了同样的意思:“我必须说,我听了Kashiwara在赫尔辛基大会(1978年8月)上的报告,他是大会的主讲人之一。无论远远近近,那里没有任何一种哲学,可以同离散系数与连续系数之间的比较沾边。我特意把我一周前在哥本哈根所作的报告立即整理成文,交给那公认有裁判之责的数学共同体。而这位Kashiwara先生的报告,则刊于[赫尔辛基]大会《论文集》。”},在会上作了一个报告\origpage{1053} % ---------------- p.806 ----------------
《论Hilbert--Riemann问题》,陈述他的等价定理。他的报告看来是完全无人理睬地过去了。相反,讨论会的一个“压轴节目”,却是几天之前Kawai的一场报告,宣布了一个他与 M.~Kashiwara合作取得的、不寻常而出人意料的结果。以一种多少拐弯抹角、刻意叫人看不懂的形式(循Sato学派特有的风格\fnote{773}{(**)}{(6月4日)关于这一点,见一处先前的脚注(第~1052~页的注(*))。在我看来,恰恰是在反常层会议的尾迹里,这种蓄意晦暗的风格在太平洋此岸被精益求精,变成一门系统蒙骗与趁乱窃取的方法。}),这定理断言:在一个(光滑的)复解析簇上,从 $D$ 到 $D^{\infty}$ 的“标量扩张”函子在“具正则奇点的”完整D-模范畴与完整D-模范畴之间诱导一个\textbf{等价}。他们的证明后来写成一篇一百五十多页的长文,业已刊行\fnote{774}{(***)}{M.~Kashiwara, T.~Kawai,\emph{On holonomic Systems of microdifferential équations III}, System with regular singularities,Pub. RIMS 15,第~813--979~页(1981年)。}。

\noindent Mebkhout当场同所有听众一样,有些懵了。这个被炒得耸人听闻、却没有一个人太弄得清究竟是怎么回事的定理,对他却有一种说不上来的\dmark{}熟悉劲儿。随后几天里,他缓慢然而稳当地、按他惯常的方式\origpage{1054}反复琢磨这件事。我可以想象:在讨论会的喧嚣之中,单是把这定理整理成非日本人能看懂的形式,他少说也得花上一两天。从那一刻起,便稳操胜券了!

\noindent 我敢打赌,在场的西方人里,没有一个对“正则奇点”是什么有过半点概念。可Mebkhout,他几年前就为一门尚在寻找自身的“系数哲学”之需,定义过一个\textbf{正则}完整D-模的概念\fnote{775}{(*)}{关于Mebkhout对一个完整D-模复形(沿一个除子$Y$)的正则性的定义,见“作品……”一注(n\textsuperscript{o}~171(ii))第~950~页的脚注(*)。光说“正则”,意即:沿(每个开集上的)\textbf{一切}除子都正则。}。至少对他来说,这概念有一个十分精确的意义——而且,取合适的\textbf{导出范畴}、并且“走到镜子另一面去”,他知道如何以“可构造离散系数”的相应导出范畴来解释这个范畴。至少,他已在论文中长篇地证明了:就在“镜子另一面”的这同一离散系数范畴的术语之下,$D^{\infty}$-完整模范畴的类似解释——而且他清楚,自己手头已有证明“完整正则D-模”情形的类似物所需的一切。这他在论文里实际上已经做到,其形态是关于$X$的一个\textbf{局部}结果——单是这一条,已足以蕴含Kashiwara--Kawai那个“耸人听闻的结果”。由此,导出范畴的观点,连同连续系数与离散系数之间的游戏,给出了一个Kashiwara--Kawai型的结果,但原则上远为有力:它同时给出诸高阶 $\mathrm{Ext}^{i}$ 之间的同构,而不止于 $\mathrm{Hom}$ 的层面(那是单用D-模、而不用由这些模构造的导出范畴所能得到的全部)。有见于此,倘若这个日本人的“正则奇点”概念竟不等价于他自己的概念,那才是见了鬼——这样一来,那个声誉显赫的结果实际上是他的系数哲学的一个纯然推论;而对这门哲学,此前还没有一个人屈尊理睬过。

\noindent 当讨论会全员出动,来给一个无名之辈的报告捧场——那报告在日程上占着一行,谁也说不清为什么——而报告的末尾\fnote{776}{(**)}{(6月4日)实际上,Mebkhout早在报告开篇就特意暗示过这一点,天真地以为这准能抓住他的听众。},那人面不改色地一连串箭头加图表地摆开(那类玩意儿\dmark{}在六十年代还时兴,在严肃人士之间早已过时\origpage{1055}),宣布说:讨论会那个著名的“压轴节目”——没人能把它的大意复述出来,这只会叫它愈发唬人——那个“压轴节目”原来是一个范畴等价定理的直接推论(说得倒轻巧!),一个他在相应的\textbf{导出范畴}(这又是些什么动物?)之间得到的等价定理,外加一个看上去同它们不大相干的范畴;而这定理,他赌咒发誓说,早已写在一篇\textbf{论文}里(这可真是锦上添花!),而且早就寄给了Kashiwara先生以及满堂显赫同行中的许多人——这情形,怎么看怎么像一个拙劣的玩笑。一阵尴尬的沉默,几声会心的窃笑。那位Kashiwara先生亲自提出例行的那个问题,想必(正是)为了驱散这个不懂规矩的年轻人造成的难堪。不过说实在的,他看上去也有点被震住了;他准在纳闷自己是不是在做梦\fnote{777}{(*)}{(6月4日)Mebkhout就此给我写道(4月22日):“莱苏什讨论会之后,有人告诉我,这位Kashiwara先生觉得他与Kawai合写的文章空空如也。可他为了不老实地补上自己落下的功课,可没少下力气。打从[1975年证明其可构造性定理的]那篇文章算起,已有五年他没碰过离散系数了。[因那篇文章而得来的]骤然的名声——它全亏了另一个全然不同的问题——使他得以去忙一些更‘严肃’的事情——千万别是瞎拼凑!1975到1980年间,在一片普遍的敌意之中(这是我后来才明白的),我是\textbf{唯一}一个在发展我从你的著作里学来的这门孩子气的哲学的人。”}……而那位无名之辈呢,却没有被吓退——他差一点要在第一场报告之上再作第二场——真是无奇不有!

\noindent 一分钟之后,我们这位无名之辈Zoghman发现,自己独自一人立在黑板前,守着他那些漂亮的图表,面对一屋子走空了的座位……那一天,以及随后几天,没有一个人屈尊问一问那个不懂规矩的家伙那些所谓的“结果”的来龙去脉——那人错就错在被邀来这么一场身价显赫的讨论会。

\noindent 应当说,大场面的吹吹打打一旦过去,那位Kashiwara先生的脑袋里总归还是起了波澜的。事实是:不过短短几个月之后,在1979--80年度的古卢阿伊克-Schwartz讨论班上,在4月22日的一次口头报告\fnote{778}{(**)}{(6月4日)古卢阿伊克-Schwartz讨论班1979--80年度,M.~Kashiwara1980年4月22日的报告,“具正则奇点的线性偏微分方程的完整系统与可构造层”。关于这场值得纪念的报告会——\textbf{Mebkhout当时在场}——的详情,见“掠夺的空白授权”一注(n\textsuperscript{o}~171$_{4}$)。}里,他宣称\textbf{那是他自己的创造}——正是那个\dmark{}\origpage{1056}曾在某次讨论会上善于投下一阵冷场的定理!不过他倒还有那份“好意”,在第2页补了一句:

\begin{quote}
“请注意:该定理\textbf{也}由Mebkhout\textbf{以另一条途径}作了证明”(着重是我加的)\fnote{779}{(*)}{我在此引用的是那份书面讲稿的文字;它由Kashiwara在口头报告一年之后整理而成。详情见前一脚注所引的那条注。}
\end{quote}

\noindent 这个“也证明了”真可谓字字千金——要知道,这是一个他连同任何别人此前连想都不曾想到的定理,是他(不过几个月前)刚从当事人本人口中听来的,而他连人家差不多一年前寄给他的论文都懒得一读!倘若他先前就知道这个定理,那他肯定不会费神去写那167页密密麻麻的证明,来证明一个“蛮横的”分析结果——而那结果不过是它的一个直接推论,甚至是推论的推论。

% ---------------- p.807 ----------------
\noindent “以另一条途径”同样叫人绝倒。那份报告里没有一丝一毫证明的痕迹,在Kashiwara其后的著作、或其任何一位日本同行的著作里同样没有。Zoghman向我担保:文献里不存在他的定理的第二个证明,只有他自己的那一个;而我(鉴于那种证明的类型——我对它再熟悉不过,理由不言自明)极疑心人们永远也找不到第二个。那是一个对应于一种几何进路的证明,使用的是Hironaka式的奇点消解——这个工具对我(和我的学生们)已成了第二天性,而分析家们(尤其是Sato学派的那些)对它一无所知。以至于Kashiwara显然觉得自己连\textbf{照抄}Mebkhout的证明都力有不逮……

\noindent 这种白线缝得极粗的骗局,\textbf{只要有一种普遍的一致同意罩着它}就能行得通——而受损的(在此)是一个无名之辈。这满堂的高雅人物\fnote{780}{(**)}{(6月4日)若要看一看直接、积极参与围绕Zoghman~Mebkhout著作之蒙骗-骗局的那帮演员的“列队”(至少是我知情的那一些),见“黑手党”一注(n\textsuperscript{o}~171$_{2}$),(f)部分“演员列队——或黑手党”。不过这份列队并不齐全——要看一份更全的名单(一口气排出十三位国际知名数学家的名字),见“光荣之日”一注(n\textsuperscript{o}~171(iv))第~962~页的脚注(*)。其中还缺一个名字:R.~Remmert——他在此期间浮出水面(见已引的“黑手党”一注,(c1)部分“记忆的失灵——或新编历史”)——十三变成十四!(还没算那位隐名的审稿人——十四变成十五……)}大可不必客气:须知上述那位\dmark{}无名者,分明正被那样一些人弃之不顾\origpage{1057}——他们恰恰是最有资格第一手知情、且对当事人负有直接人身责任的人:J.L.~Verdier(答辩委员会主席)与 P.~Deligne(头一个感受到那结果之分量的人——他是前一年从Mebkhout口中得知它的)。

\noindent 既然说到了太平洋对岸的卡伊德Kashiwara,不妨就在这一章上收尾,讲一讲那位当值无名之辈被彻底清洗的尾声——顺着三年前1981年6月反常层会议上那道响亮的先例。这说的是 R.~Hotta与 M.~Kashiwara的一篇文章“半单李代数上的不变完整系统”(Inventiones Mathematicae 75,第~327--358~页),刊于1984年(1983年3月2日收到)。从引言第6行起便可看清:这篇文章是那道没完没了的、所谓好上帝(或当值无名之辈)的“Riemann--Hilbert对应”的又一应用。在这篇文章里,\textbf{那位无名者的名字不再被提起},书目里也没有他。Zoghman已知第二作者的心性,却无从预断第一作者是否也不老实,于是给他写信,告诉他自己是文中以举足轻重的方式用到的那条定理的作者,并就自己未曾以此被引用一事提出异议。结果呢?引文给了已引的Kawai-Kashiwara那篇文章(167页那篇)——那条定理在那里恰恰根本不见踪影\fnote{781}{(*)}{(5月25日)如他处已解释过的(在“五张照片(晶体与D-模)”一注,n\textsuperscript{o}~171(ix),尤其第~1005~页),那部著作只含好上帝定理的“一半”,而那一半是从Mebkhout论文第三章劫掠来的。}。Hotta回复他说:他觉得没有必要引用他,因为\textbf{反正人人都知道,那条对应是Kashiwara与Mebkhout的}。幕落……

% ---------------- p.808 ----------------
\subsubsection*{(c) 入门费——或有为青年}

\tocsubsub{(c) 入门费——或有为青年}

\noindent 但日本毕竟路远;我的朋友Zoghman经年累月地同遥远的日本人交锋,想必是因为:对他来说,直面一桩绝不限于天涯海角大陆的黑手党的现实,要痛苦得多——这黑手党无论在巴黎那些气派的讨论班里,还是在莫斯科或\origpage{1058}东京,都高踞\dmark{}首座。该回到这可爱的法兰西国土了,回到我亲爱的上同调前学生们组成的“小家庭”——以及(那稍大一点的)自我“亡故”的遥远时日以来围绕他们结成的那个圈子。

\noindent 消息有时跑得飞快。1979、1980年间,借着Deligne与莱苏什讨论会的东风,“人家”总该渐渐察觉:数学市场上刚刚冒出一个定理,说实在的颇有苗头,可惜出自一个落伍的Grothendieck派余孽;不过,这个叫人提不起劲的父权,有一个现成的替身——人所共知的日本分析家Kashiwara,他正求之不得,要出演那著名的“Riemann--Hilbert对应”的父亲。

\noindent 1980年1月,Mebkhout就他那条来得不是时候的定理,在巴黎第七大学 Lê Dung Trang 主持的“奇点讨论班”上作了一个报告。Jean-Louis~Brylinski没有到场听讲,但 Lê Dung Trang 向他谈起了它,并让他读了自己的笔记。据Brylinski自己后来向Mebkhout坦率讲出的:他一得知Mebkhout的定理,便脱口而出——有了这个,Kazhdan--Lusztig猜想就有法子证明了!(猜想嘛,照占卜师们的说法,是“无从下手的”。)

% ---------------- p.809 ----------------
\noindent 人们或许以为,Brylinski会去找当事人,请他把自己所需的那条定理所系的“完整性”与“正则性”条件的种种奥妙讲个明白。可照他自己向Mebkhout坦白讲出的:是“人家”劝他别去找当事人,而去找那位大人物Kashiwara。他没有点明那“人家”是谁。不过显然,他耳朵尖(脑子也快),而且那时候他也还同Mebkhout至今一样默默无闻。这种话无须说第二遍,他便去找Kashiwara打听去了——Kashiwara那时想必还在近旁。那是他最正当不过的权利。结果便是与Kashiwara合写的一篇文章,刊于 Inventiones Mathematicae 64(第~387--410~页),1981年(1980年12月19日收到),题为“Kazhdan-Lusztig 猜想与完整系统”。Brylinski一夜之间成了明星——这原是他应得的;而Kashiwara呢,又往那已够辉煌的战绩簿上添了一根花翎\fnote{782}{(*)}{把Kashiwara的名声同他刚刚找到的证明——一个Kashiwara毫无寸功的证明——撮合在一起,同时只字不提那位无名的年轻同行所扮演的关键角色:这就是Brylinski不加推辞地付出的“入门费”,为的是跻身某个“圈子”——名流们的圈子,本注“黑手党”之名即得自它……}。

\dmark{}一切本可以在这最好的世界里尽善尽美,但是……想必还是那个“人家”又出了主意\origpage{1059}:某个无名的泛泛之辈,提得越少越好。反正,寄给 Inventiones 的那部手稿里,\textbf{Mebkhout的名字不见踪影},正文里没有,书目里也没有。

\noindent Mebkhout看到了文章的预印本,就此向Brylinski提出了抗议,并致信 Inventiones 的编辑 R.~Remmert。Brylinski的反应很“柔顺”(那笔调如今对我已再熟悉不过……):在校样上于书目末尾(字母次序之外)添了三条“暂停!”式的Mebkhout引注(顺手为之!),但正文里对这位Mebkhout先生却始终一字不提\fnote{783}{(*)}{Brylinski-Kashiwara文章的引言,以对几位作者的鸣谢收尾,其中包括Jean-Louis~Verdier(而那位当值无名之辈,不必说,只字未提)。随后是第1节,概述“具正则奇点的完整微分系统”(这是日本人对完整正则D-模的叫法)。在该节开头几行,人们读到:“至于细节与证明,请读者参看[6, 15--17]。”文献[6]即Kashiwara1975年确立其可构造性定理的文章,而[15--17](校样上追加的)就是对Mebkhout的“暂停!”式引注。无论发生什么,“有为青年”Jean-Louis~Brylinski的荣誉总是无恙的……}。读这篇文章的人,倘若居然在书目末尾天晓得为什么多出一个“鼎鼎大名的无名之辈”的名字,只会想:准是为了照顾哪位朋友的情面,才把他放在那儿……

\noindent \textbf{Brylinski是靠一场骗局跨入名流的}。实情是:只要那件新工具不曾问世,他所证明的那个猜想便无从下手。且不说那件工具的\textbf{发明权},这篇文章里没有任何东西把这件新工具亮出来——它的角色从第6--8行起就被那句不咸不淡的“解释”(原文如此)掉了包:

\begin{quote}
“本文采用的方法,是把旗流形上具正则奇点(R.S.)的完整线性微分系统同 Verma 模联系起来,\textbf{并使用完整系统与可构造层之间的对应}。”
\end{quote}

\noindent (着重是我加的。)对这条未加限定的著名“对应”,没有片言只字的指涉或解释。“人家”想必让这位年轻后进心领神会:这条“对应”从此便属于人人皆知之事,根本无须为此\dmark{}\origpage{1060}援引一条专门的定理——那会掀起发明权这类枝节的、而且(尤其)不合时宜的问题。而Brylinski,这个有为青年,自然无须人说第二遍……

\noindent 至于Remmert,他把那位无名控告人的信转给了Brylinski-Kashiwara文章的审稿人。审稿人驳回了控诉,表示的意见是:“该结果\textbf{独立地为已知,而且大概更早,系由Kawai与Kashiwara所得}”,指涉的是他归于这两位作者的“重构定理”(指涉上引作者文章第~116~页,载 Guillemin 主编《微局部分析讨论班》,Annals of Math. Studies 丛书,n\textsuperscript{o}~93)。

% ---------------- p.810 ----------------
\noindent 审稿人的这一评价——他可是被认为知其所言——在两个层面上都骇人听闻,并且表明:他本人就是同一场骗局的当事人,与(此刻的)Kashiwara和Brylinski沆瀣一气。仅仅凭一个关于\dmark{}种种独立所得结果\origpage{1061}(据审稿人本人所表达的意见)之在先性的单纯\textbf{假定}\fnote{784}{(*)}{(6月4日)我在这里甚至撇开了一点:这个假定根本没有根据。而且,转达审稿人答复的Remmert的信(1981年1月26日),对(信中援引的)Guillemin 讨论班与Kashiwara报告的日期,也只字未提。我刚刚火速打电话追到意大利的Mebkhout,问他这条指涉及其日期的细节。我得知:Kashiwara的报告在1978年——就在Mebkhout把自己论文的第三章告知他(1978年1月)之后几个月——Kashiwara先生可没有浪费时间!由于论文答辩直到1979年2月才举行(这要归功于巴黎各大学论文委员会——那个 J.L.~Verdier珍爱有加的机构……——的拖拉),这本来可以为审稿人的在先性“假定”提供一个说得过去的依据——至少就“重构定理”而言。可倘若那位审稿人(姑且不说他老实不老实——而这一点,明摆着已经谈不上)尽心尽职地干了他的活,他本会注意到:在Kashiwara那份所引的报告里,没有任何东西像一个“重构定理”的\textbf{证明}。\par Mebkhout后来还再度出击,在1981年3月25日的一封信里,他强调:1\textsuperscript{o})审稿人所援引的定理是“他博士论文最重要的成果之一”,而且他已把这一结果连同其证明告知了Kashiwara(可他忘了说是\textbf{何时}——Zoghman从来不肯说这个!);2\textsuperscript{o})这条定理“对于确立所论的范畴等价,是远远不够的”。R.~Remmert没有屈尊回复这封信——它来自一个无名无援的控告人。\par Zoghman方才向我说明(我这样追问下去,终会查明一切的……):他是第二年、即1979年——他答辩的那一年——才在 Guillemin 讨论班上得悉Kashiwara那场骗局的。因此,那正是他头一回同“黑手党”圈子里通行的这类手法打照面。等到这年9月的莱苏什讨论会,他对大明星Kashiwara的底细,心里已经有数了。不过,既然他的哲学与结果连同证明都已白纸黑字地写出并刊行,他原不会想到:一旦其重要性得到承认,他的著作竟会被干脆利落地整个掉包。而他的进路之力量的第一个征兆,恰恰就在莱苏什讨论会上、关于Kashiwara--Kawai定理一事中显现。\par 当然,1978年1月,Mebkhout(那时他还没有任何理由起疑)向Kashiwara谈起的,不只是他所谓的“双对偶定理”(后来为一场骗局之需被重新施洗为“重构定理”),还有完整的好上帝定理——前者总而言之是它的“一半”(两半中较浅的那“一半”)。他对我说:对于双对偶定理,Kashiwara倒确实“接住了”,看来他此前已自己提出过这类问题;但显然,他全然不知如何证明它。(然而,Mebkhout的证明并不使用奇点消解。)至于好上帝定理,那可是彻底地从他头顶上飘过去了——以至于到莱苏什讨论会时,他把这事忘得干干净净。然而Mebkhout像对所有人一样,同年(1979年)年初就把论文全文寄给了他(也就是说,那时他还不曾察觉前一年 Guillemin 讨论班上的那场骗局)。还有一件事表明好上帝定理彻底溜出了这位卡伊德的掌握:就在同一份 Guillemin 讨论班报告里,他甚至没有想到——为了良心安稳,不妨这么说——把它也一并收进腰包(哪怕他根本不明白那是怎么回事……)。\par 我至今无缘把Kashiwara这份报告(*)拿到手里,于是不免自问:它会不会让一个不知情的读者得到这样的印象——Mebkhout发展的那门哲学,Kashiwara(而且靠他自己的手段,如他所说)至少从1978年起就已知晓。Zoghman答应给我寄一份这份报告的拷贝;他向我担保,它会让我打消这种误会。那里头(他说)堆砌着一批技术性的陈述,或多或少(不)可理解(Kashiwara不可能做得更少了……),没有证明,也没有一条显见的引线,更没有任何东西(同他同年赫尔辛基的报告、以及次年莱苏什讨论会的报告一样)像一个联结连续系数与离散系数的“系数哲学”。\par(x)(6月16日)Mebkhout向我说明:那份报告实际上是由\textbf{Kawai}宣读的,作为与Kashiwara的一项合作成果。},就认定那位假定的后作者(偏偏碰巧正是那位无名者……)可以全然不被引用,这已经够骇人听闻的了;这类做法,明摆着为最严重的滥用敞开大门(而且早已敞开多时……)\fnote{785}{(*)}{这与 R.~Hotta三年后以同样犬儒口吻表示的态度(见上文所引他对Mebkhout的答复)如出一辙:新的“规矩”,或者更恰当地说“圈子的法则”,就是引用身居权位者(哪怕名不副实),不引用无名之辈——哪怕后者的贡献是决定性的、并有无可抵赖的出版物为证。\par 我并不怀疑 R.~Remmert在此事上的善意。但我还是要指出:身为 Inventiones 的编辑,他在这场骗局中的责任无可推卸、是直接牵涉在内的——即便撇开下述事实不谈(这一点他当时本无从料想):他是被一位不诚实的审稿人引入了歧途。那位审稿人曾表示过一份“希望”(就当时的情形而言,这希望堪称犬儒):“出于礼貌,Brylinski与Kashiwara会提及Mebkhout的结果”。而身为编辑,R.~Remmert的职责恰恰在于:确保Mebkhout的结果在正文中得到应有的提及——不是作为“礼貌”,而是出于对数学家这一行当基本伦理规范的尊重。\par(5月30日)自写下这些文字以来,我又得知了一桩“新事实”:它以意想不到的方式照亮了 R.~Remmert在围绕Zoghman~Mebkhout著作那场骗局中所扮演的角色,表明他也积极参预了围绕我本人著作的那场骗局。于是我对他所抱的那份善意推定便就此消散(那是出于老习惯,也因为此前没有无可辩驳的反证)。感兴趣的读者可以在下文(注“黑手党”的)(c1)部分、在“记忆的失灵——或新编历史”这个标题下,找到关于这桩“新事实”的详情。}。可还有更甚的。他所引的那条“重构 % ---------------- p.811 ----------------
定理”(它\dmark{}同样是从Mebkhout的论文里劫掠来的\fnote{786}{(*)}{关于这一劫掠,见“五张照片(晶体与D-模)”一注(n\textsuperscript{o}~171(ix)),(b)部分(“好上帝公式”)末尾,第~1005~页。},在那里它顶着\origpage{1062}(不确切的)“双对偶定理”之名)离Brylinski-Kashiwara那篇问罪文章的证明中所用的范畴等价(号称“Riemann--Hilbert”等价)还远得很——那等价出自Mebkhout一人,而那条定理对它丝毫无所蕴涵\fnote{787}{(**)}{Mebkhout的“双对偶定理”与“好上帝”定理——前者构成后者的“一半”,是两半中较浅的那一半——两者之间的关系,见已引的那条注(同样在(b)部分)。前者无须诉诸消解,而完整的定理则动用Hironaka奇点消解的全部力量(那是一件典型的“几何”工具,日本学派至少直到八十年代初还对它一无所知)。}。在我这里,审稿人的不老实——他倚仗上同调建制的一致默契,去封杀一个无名之辈的名字与著作,以成全名流们的“好处”——不容丝毫怀疑。凡略具上同调-分析素养、且对这一迷人主题有起码兴趣的人,都能亲自核实事实,并且看出这是一场粗劣的戏法——而那位匿名审稿人竟趋奉其间\fnote{788}{(***)}{(5月30日)而 R.~Remmert,以 Inventiones 编辑的身份,对它作了毫无保留的襄助……}。这局面的暧昧余地就更小了:因为直到1981年,在Kashiwara那里、也在任何其他日本或其他微分系统的行家那里,“导出范畴”这个词还无人说出口\fnote{789}{(****)}{(5月25日)Mebkhout提醒我,这句斩钉截铁的话应当稍作限定。尽管导出范畴自我离去之后在法国实同禁忌,日本学派却一直对它作节制的使用。那是一件方便的技术手段(尤其为免于诉诸谱序列),却绝非一套为形形色色簇与空间的上同调中“系数”的内在几何视野“量身裁制”的语言。},更谈不上有\dmark{}丝毫反思朝着一种联结离散系数与连续系数的“哲学”的方向——\origpage{1063}而老实说,这门哲学在后来的那些含混其词的指涉里同样无影无踪——什么“(原文如此)完整系统与(原文如此)可构造层(再来一个原文如此)之间的对应”——这些翩翩君子,直到今天,也没有一个人拿出这份诚实:\textbf{把所涉的范畴、以及确立其等价的那些箭头,白纸黑字地(像我不久前所做的那样)明写出来}。相反,自1977年以来,Mebkhout的一系列讨论班报告、笔记与文章,为他自1972年起在全然孤独中完成的开拓工作作证\fnote{790}{(*)}{这些文章的清单,我在这里既不评述、也不开列,请看已引的Mebkhout与 Lê Dung Trang 的文章(载 Proceedings of Symposia in Pure Mathematics, 40 (1983), part 2)。(5月25日)另见“三座路标——或天真”一注(n\textsuperscript{o}~171(x))中随文给出的书目。}。

\noindent 我必须承认:在直面这件事、并久久地、从方方面面把它端详与检视之前\fnote{791}{(**)}{(6月1日)我先是在去年,5月2日至9日那一周(写“送葬行列VII”——题为“讨论会——或Mebkhout层与反常性”)做过这件事,而后近两个月以来,在写《登峰造极》时,又做了一遍,……},我连做梦也不曾想到:在学者的世界里,竟会发生这样一场恬不知耻的集体掠夺。而不得不对自己说:这场不义的蒙骗,主要是由我最亲近的两位昔日学生合手搬演上场的;更有甚者,那开场锣竟是由\textbf{我的著作的一个接棒人的出现}敲响的——那份著作,我曾满怀激情投入其中,把我所能奉献的最好的东西都放了进去\fnote{792}{(***)}{在把这一页(涂改得相当厉害)誊清时,我心中浮起一个念头:倘若我对这份著作的投入居然结出了(别的)这样的果实——出乎意料而不受欢迎的果实——那无疑是因为,在这份投入本身、在我当时的精神状态里,并不只有我在这里乐于强调的那个“我之最优”,同样还有“我之最劣”。这一点,《愚妄与更新》(《收获与播种》的第一部分)里固然已相当清楚地显明,可它也是这样一件事:种种强大的自我机制不断地逼着我把它忘掉!我开始意识到,这个“最劣”,去年的反思中只是匆匆一瞥;我并没有对它作过一次真正深入的检视,或作一次把它形形色色的面相真正细致地揭示出来的“巡礼”。这就是为什么我对它的认识仍然浮于表面,一如这份认识的作用(尤其在我同葬礼的关系中)仍然浮于表面。\par《收获与播种》这第四部分“四种运算”,主要是一份细致的清点工作——把与葬礼相连的种种原始事实重新清点一遍。可这份“后勤”式的活计,毕竟让我感到:对葬礼的一种更深入的理解,与其说将来自我近三个月来所做的这类工作,不如说将来自对《愚妄与更新》中那份工作的深化——也就是说:也来自对“我离去之前”那些遥远时日里的那个我的认识的深化……}。我离去之后,这份著作竟成了那些人\dmark{}觊觎的靶与猎物\origpage{1064}——而他们恰恰曾是与我最亲近的人;还有一种隐秘的暴力,越过了我本人与我的著作,击中的还有那些公开从中汲取灵感的人……

% ---------------- p.812 ----------------
\subsubsection*{(c1) 记忆的失灵——或新编历史}

\tocsubsub{(c1) 记忆的失灵——或新编历史}

\noindent (5月30日)写下前几页六周之后,我在讲述我朋友Zoghman坎坷遭遇的叙述里,在此打开一个括号,稍稍停留在一个“新事实”上——它在先前一处脚注(第~1061~页的注~793[原书如此,当为注~785])里已有暗示。接下去的几页,读者可以读作一份有趣的补遗:关于“新风格”的开花结果——那门风格别处已谈过\fnote{793}{(*)}{见“恭贺——或新风格”一注,n\textsuperscript{o}~169$_9$。},它擅长(让所有人满意地……)为当代数学的(在此情形下的某一)主题写出一部“新编历史”。急着知道我朋友Zoghman后续遭遇(他迷失在一场他无从预料的马戏里)的读者,可以径直跳到“总彩排(封神之前)”(下文(d)部分,日期为4月16日)。

\noindent 我读到了 S.~Bosch、U.~Güntzer 与 R.~Remmert 的书《Non Archimedian Analysis》的引言与书目\fnote{794}{(**)}{即 Grundlehren der Mathematik 丛书第~261~种(1984年)。}。这本书阐述刚性解析空间的理论,并(理所当然地)把 J.~Tate1962年的(“私人”)笔记《Rigid-analytic spaces》呈现为这门理论的出发点。引言里说明:R.~Remmert“曾得以获得一份”这份稀世文献的拷贝;而这份文献,在\dmark{}“簇”概念的名流集会中,多多少少算是一位新来者的出生证(解析的“簇”,在此情形)。

\noindent \origpage{1065}Remmert想必忘了:把这份文献交由 IHES(那时刚刚起步)油印、并给他以及另外几位复解析空间的专家寄送拷贝的,正是我——为的是让他们注意他们偏爱的主题这一出人意料的扩展。那是在他们中间还没有一个人装作对实数域或复数域之外的基域感兴趣的时节——不过世事难料……

\noindent Remmert想必还忘了:我当年之所以那么起劲地在朋友们中间传布这份宣告一个几何“宇宙”之诞生的文献,(除别的之外)是因为我同这次诞生有密切的干系。“刚性解析空间”这个名字本身,就是我想出来的——那时不仅Remmert,乃至Tate本人,都不曾听说过这个名字,连这名字所应表达的东西,也无人梦想到过。我是头一个看出:Tate的椭圆曲线“loxodromique”理论,应当成为一种当时尚不存在的“解析”簇的“取商”,并且应当引出Serre“GAGA”型的代数-解析比较定理。还有另一个动机给我指明了通往这同一类新对象的道路:需要能够对离散赋值环上有限型的形式概形定义“泛纤维”。

% ================================================================
%  18.5.4 注 (c1) 之尾 — (d) — (e) — (f) — c3 — c4 — Epilogue outre-tombe
%  对应原书 p.813–836（PDF 物理页与印刷页同）
%  源文: _work/chunks/chunk01–07.md
% ================================================================

% ---------------- p.813 ----------------
作为指向同一方向的第三个迹象:我曾听说,克拉斯纳(五六十年代的巴黎
数学圈里人尽皆知的一位怪才:他家里收容着一支猫大军,裹着一件厚重的俄式大衣,永远一副
乐呵呵的表情,在各讨论班之间走来走去……)——就是这样一位克拉斯纳,据说在非阿基米德
赋值域上“做解析延拓”。我所知仅此而已;我甚至不能确定可曾遇见过哪一位读过克拉斯纳这
方面工作的人——但这件事本身已够引人好奇。应当说,“解析延拓”这个词本身可没有让我心
跳加快的魔力(毋宁说相反,它叫我想起学生时代一些索然无味的记忆……);可一旦隐约察觉
到需要一类新的几何对象,它就难免让人心头一震……

回到Remmert——既然他的记性糟糕到这般地步,Tate的原文本(他还自夸握有一份)总该能替
他唤回记忆。在笔记里,Tate对我在这套理论的构想中\dmark\origpage{1066}所扮演的角色没有丝毫讳言\fnote{795}{(*)}{自Tate与我、他的家人与我的家人亲密往来的那些遥远日子以来,二十多年已经过去。多少年了,再没有他的只字音讯。这本书印出来,我的前学生与旧友当中没有人会看不到这篇引言里对我本人的这番抹除——可我也从未听说他们中哪一个为此动过声色。此一时,彼一时……},其中写道(凭记忆
引述):他“完全忠实地”遵循一位“总施工”(maître d'oeuvre——这是就通过“拼
合碎片”来构造该概念的方法而言),而这位总施工正是他从我这里得到的。此外我还为他提供
了某种“砌块”(或者说:在受限形式幂级数代数中进行“局部化”的“方法”),以应形式概形
纤维之需。他又以第二类“碎片”(或“方法”)补足这最初的一批,两者在某种意义上互为补充。

若不是我手里有“广义空间”这条主线(它后来成为\textbf{topos}),这个新概念恐怕根本不会问世
——平展上同调不会,晶体上同调不会,尾随其后涌现的许多东西也不会,甚至连那道最时兴的
“奶油甜点”、大名鼎鼎的 D-模也在内;这条主线的理论当时尚待建立,但四年来已被人隐隐预感。
正是这条直觉,为我指明了通往一类“簇”的道路——这类“簇”恰恰\textbf{越出了}寻常拓扑空间
(局部赋环空间)的框架。

自从John~Tate启动\textbf{刚性解析空间的局部理论}那一刻起,同样是我提出并推广
了关于这些新簇的首批关键“整体”定理的表述;这些表述早在第一批基础工作完成之前,就已
经在我心中成形:关于刚性解析空间上相对真概形的代数-解析比较定理;关于刚性解析空间的
真态射 $f$ 的 $R^i f_*$ 的有限性定理——这些问题由Kiehl在随后几年里解决\fnote{796}{(**)}{这里要指出,早在Tate打下刚性解析空间理论最初的地基时,我就清楚:他所处的那个框架仍是暂定的,丝毫不足以穷尽我试图用“刚性解析空间”这个名字表达的直觉内容——正如域上有限型概形穷尽不了“概形”一词所系的直觉。有一条线索指向Tate框架的实质扩充(凡愿意听的,我都讲过……),而提供它的是Tate自己:他在某个拓扑环上写出了一条“Tate泛椭圆曲线”(这个环是形式幂级数环 $\mathbb{Z}[[t]]$ 的一个子环,由那些在复平面开单位圆盘内收敛的级数组成——如果我记忆无误),而这个环显然应当被看作某一个刚性解析空间的“仿射坐标环”,其类型并不在Tate开出的清单之内。自我离去之后,一切基础问题都跌入人人不屑的境地;1962年Tate搭起的这台概念机器从那以后纹丝未动,也就不足为怪了。}\fnote{797}{(***)}{(6月4日)我还曾最先力主:对刚性解析空间,必须引入比Tate所设想的更一般的“点”(Tate的点只取基域的\textbf{有限}扩张为值)。这一必要性,既受代数几何之类比的提示,也出于一个愿望:为与所论刚性解析空间相伴的 topos 的“点”找到一个具体的解释。}。不过
话说回来,照眼下盛行的风气,\dmark\origpage{1067}预想新概念、物色总施工、提出供“真正的数学家”去解决
的问题,都被当作无关紧要之事——往极端里说,简直就是唬人的花招……

% ---------------- p.814 ----------------
不管怎样,这篇引言里没有一字提到我的名字与刚性解析空间有什么相干。克拉斯纳的名字同
样不见踪影——非但如此,Tate的理论还被说成引入了“一种足够丰富的结构,使不可能成为可
能:全不连通域上的解析延拓”——可早在1962年,这项(“不可能”的)解析延拓,十年也罢,
二十年三十年也罢(我说不准),早已是克拉斯纳“注册在案的正业”(若可以这样说)了。汗
牛充栋的参考文献里,克拉斯纳和我都无迹可寻。不过快到引言末尾,我的名字倒是顺带闪现了
一回,在“Grothendieck 拓扑”这个名字里;这个概念所引的是Artin1962年的笔记,
而对 SGA~4 里那份精益求精的打磨功夫(自1963年起,贯穿整个六十年代——却顶着一件显然
不受欢迎的父名……)则傲然无视,完全步了我那班前学生的后尘……;同样只字未提的,是当
年(1966年)我曾为刚性解析空间在晶体上同调的发展中指派的那个角色——那时Remmert(连同
他那些显赫的复分析家同事们)还全无迹象要理会这类所谓的“刚性解析”的怪“簇”(这名字
您给评评理……),它们是几位代数几何学家在自己角落里捣鼓出来的——仿佛复解析空间还
不够那些严肃的分析家与几何学家消遣似的……

\pmark\origpage{1068}关于书中所述这套理论的诞生史,只要是从第一手知情,就能看出这篇引言里招摇着同一种犬
儒做派——它同样流露在那份匿名审稿人对一位无名申诉者的答复里(并蒙这位Remmert先生祝福):显然,在作者们的心目中,这不过又是一桩“礼数”问题,说到底是一种“客
气”,施与不施全凭他们高兴:在新理论的孕育过程中起过关键作用的某某某之名,要不要写进
他们的“历史概述”(原文如此),悉听尊便。在他们看来(恐怕在数学界几乎所有建制中人看
来亦然——他们对此类歪曲照单全收,连眉头都不皱一下……),“历史”并不是\textbf{实际发生过
的事},而是一件可以由僭取书写之权者\textbf{乾纲独断}的东西,或者由一小撮人合议而定:什么有资格
存在,什么有资格算作发生过,全由他们说啦。

这班人就爱对苏联发生过、还在发生的事情津津乐道,而且一桩也不会放过(我知道自己在说
什么),忙着签署什么“捍卫自由”(思想自由之类……)的宣言——那是在\textbf{别人}家里;而在
\textbf{他们自己}掌权的地方,同样的谎言专制一天也没停过。

(6月3日)几天前在前几页里提到克拉斯纳那副独特而可亲的形象时,我心头冒出一个问题:
他还在世吗?他比我年长一两辈,我已有很久(少说十五年,也许二十年)没听人提起他的名字。
人物我还记得清清楚楚,他的名字却要过几秒钟才浮上来。(说实话,这一类事如今常在我身上
发生,年纪不饶人……)克拉斯纳以好客闻名,他的俄国血统本是另一个本可让我们接近的共同
点。可我那时一头扎在数学里,腾不出闲工夫只为“有趣”去交朋友。我们俩对待数学的路数,
想必正处在对跖点上。我们大概在Bourbaki讨论班的间隙闲聊过一两回,数学是肯定没有谈的。
那时真正让我着迷的,除了数学几乎没有别的……

% ---------------- p.815 ----------------
\pmark\origpage{1069}话虽如此,今天我收到Deligne的一张便条,不过寥寥数行,谈一桩无关紧要的实际琐事,
或许是想让自己还留在我的好记忆里(我们之间没有书信往来,算来已有几个月了);又或许是
想找个地方安插一条附言——我冒昧把它抄在这里(想必他会同意):

\begin{quote}
“又及:惊闻克拉斯纳已于半月前去世,不胜哀恸。我一直记得二十来年前他在布鲁塞尔所
做的一场滔滔不绝的长篇报告,我自然是云里雾里,却硬是撑到最后,成了所剩无几的听众之
一。令我吃惊的是,你那张五十年代人物表\fnote{798}{(*)}{这里对我在《收获与播种》第一部分(“Fatuité et Renouvellement”)中的说法有一个明显的误会。那一部分的主旨从来不是勾画一幅数学“五十年代群像”——哪怕只限巴黎圈子,或者Bourbaki周围那一圈。我的根本意图,是去发现我作为数学家的过去。正因如此,我才谈到我与这样那样同事或学生的关系——每当这层关系在我的生命中显得重要,或者能够照见我自己的时候。}里竟没有他——他在那个年代做了不少漂亮事,虽然与Bourbaki的精神格格不入,
而且有一种专把定义写得潦草马虎的天才。”
\end{quote}

这样一来,这里又多出一篇悼词,这一次是为我的一位同穴之人而作。这一篇里,我自信看出
一丝同情,或者毋宁说是某种情感的余光,那种情感往日曾经鲜活。可是,同在我的悼词里一样,
我的朋友皮埃尔决不会松一松牙关,说出——这一次是为了纪念一位一去不返的逝者——他津津
乐道却从不点名的那些“漂亮事”究竟\textbf{是什么}。然而他和我一样清楚:这些“事”为一个
如今正开得繁盛的理论的问世铺了路;而且,出于他或许心知肚明的原因,新主子们乐于抢先把
这位和善、潦草、“马虎”的先行者提前下葬(就葬在我的旁边)——他刚刚离世。他确凿无疑
是一位在超度量域上“做解析延拓”的人物——那时Tate、Remmert或者我还在“做”三角形全等
的种种情形和毕达哥拉斯定理,而我们的朋友皮埃尔还得让妈妈替他擤鼻涕(顺便擦洗……)呢!

% ---------------- (d) La Répétition Générale (avant Apothéose) — p.815–818 ----------------
\subsubsection*{(d) 总彩排(封神之前)}
\tocsubsub{(d) 总彩排(封神之前)}

(4月16日)可我必须回到我身后学生Zoghman~Mebkhout那一连串够味儿的“倒霉遭遇”上来。
\dmark\origpage{1070}1979年6月,当一个籍籍无名的陌生人——口口声声以Grothendieck的思想为宗——亲口告诉他,
那个让Deligne十年前整整一年绞尽脑汁、终未得到满意解答的关键问题\fnote{799}{(*)}{(5月25日)也可能,Deligne早就失去了对“关键”二字的分寸。关于这一点,见注“……以及羁绊”(n\textsuperscript{o}~171(viii)$'$)。}有了一个漂亮的解法时,Deligne的脑子里究竟发生了什么,我毫无概
念。以他一贯的心性,可想而知,他绝不会去祝贺这个年轻人在他Deligne栽过跟头的地方得了
手。可我深深感到,他那掘墓人的心性使他的嗅觉(我曾领教过的惊人嗅觉)失灵到这般地步,
以至于直到今天(六年之后),连他自己也没能把握住那个无名之辈的思想与眼光的真正分量。
他和所有人一样,最终只看见了那道“奶油甜点”:那件人人翘首以待的意外工具,那把砸开
“出了名的难题”的铁钎。可他也曾有一天,把别人传递给他的一整套广阔眼光据为己有——为
的是把那眼光连同生下它的人一并埋葬,再顺手抢走又一件工具,把它也改造成一柄“铁钎”……

我所知的Deligne对Mebkhout定理的最初一次反应,是一封未注明日期的短信(手写),写给Mebkhout,收到于1980年10月10日\fnote{800}{(**)}{这就是那份“以保密为条件提供、我在此不再多置一词……”的文件,“受害者”一注(打字稿第309页)里提到过它。事隔一年,Zoghman慨然允准我在此照录。}:

\begin{quote}
亲爱的Mebkhout:

我已把我那本你的论文寄给Bernstein和Beilinson:他们证明 Kazhdan--Lusztig 猜想用得
着你的结果(我有一份他们工作的俄文摘要,你要的话我可以转寄)。能不能再给我寄
一本?

谢谢。

P. Deligne
\end{quote}

从这封信我推想,Deligne必定已把“好上帝定理”的事告知了那两位苏联数学家,或许还暗示它
可以用来证明那个猜想;要么这是他自己悟出来的,要么当时已有风声,说Brylinski在这个
题目上有想法。Mebkhout那场“点着了”Brylinski的报告,早在1980年1月。Brylinski-Kashiwara
一边,Beilinson--Bernstein一边,两篇文章——用那个没有名字的定理(它的作者更没有名字)证
明了那条著名猜想——\dmark\origpage{1071}收到日期一为1980年12月19日,一为1980年12月8日,前后只差十一天。
巧合?

我甚至冒出过一个念头:Deligne1979年6月就比谁都先知道这件新工具(反正谁也不曾费神去
读那个无名之辈的厚砖头,Deligne也没有)——他为什么不自己想到把它用在那个猜想上,收
获一顶新桂冠,反倒去帮苏联同行摘桂冠?论头脑的敏捷,他不见得输给Brylinski吧?说不定
从那一刻起,他就已经瞥见一种可能:从旁门把“好上帝定理”本身的发明权捞回来——这发明权
(他心里必是这么感觉的)十年前就该是他的;说不定在他看来,正是由于一场不可饶恕的发牌
错误,那个不自量、欠打磨的狂妄后生才僭取了权利,去证明他Deligne早已久久钻研而终无决定
性成果的东西。他最终离成功只差那么一丝丝;换了别人在他流汗无功的地方收割,这可不公道
……但若想收回那件归根结底“理所应当”归于他的东西(依照某个人上人圈子里终于通行起来
的不成文法——他自认是那个圈子的中心和大佬……),那他就得以全然不同的细腻手法来操盘,
不能妄图一口吞下太多\fnote{801}{(*)}{这当然只是揣测:关于那条著名“对应”的据为己有之念,是否早在Deligne初闻其事时就已存在。就我自己而言,我深信不疑。诚然,上引那封信看上去会给人相反的揣测。而在我看来,这仍是那种挑衅姿态的又一个信号:他Deligne根本不必有半点顾忌——反正对方不过是个无名之辈,孤身一人对全世界,\textbf{横竖动弹不得};而他Deligne尽可以“自甘出丑”,正如他尽可以借“反常层”(\emph{faisceaux pervers})这个挑衅性的命名,以一种象征性却刺眼的方式,宣告自己真实的心性。关于这一点,见注“反常”(n\textsuperscript{o}~76),以及(在心理情境十分相近、却不那么极端的场合)注“玩笑——或‘权复形’”(n\textsuperscript{o}~83)。}。

话虽如此,Zoghman吃过与Kashiwara、与Brylinski打交道的怪插曲的亏,已有戒心,他审慎行事,决定亲自向Beilinson、Bernstein两位先生通报那个定理——Deligne说他们用得着它——以防Deligne这么大的
人物在谈起定理时,忘了提一句谁是它谦卑的作者。时机正好:下个月,1980年11月24日至28
日,莫斯科有一场“广义函数及其在数学物理中的应用会议”。Mebkhout在会上报告了他的定理,
已刊行的\dmark\origpage{1072}题目是“高维的Riemann--Hilbert问题”;他特意与Beilinson、Bernstein当面交谈,把自己
结果的来龙去脉原原本本讲给他们听。

他来的正是时候。会后不过十天,两位作者就寄出了他们关于 Kazhdan--Lusztig 的工作,以《法
国科学院院报》快报的形式(第292卷,1981年1月5日,第I辑--15):“群论——$g$-模的局部
化”,Alexander~Beilinson与Joseph~Bernstein署名,\textbf{由Pierre~Deligne转交}。理所当然,他们的
手稿上只字未提Mebkhout之名——显然,Deligne把那个无名之辈忘了个一干二净,尽管他明明把
此人的论文转寄给了他们,其目的正是……?谁懂,谁就懂吧!Mebkhout费了九牛二虎之力,才
让Beilinson(“两人中更老实的那一个”,Mebkhout一本正经地向我担保)相信:他们书目里引的
Kashiwara--Kawai那篇文章里,唯独没有那件“构造”(此处替代那句说了八百遍的“对应”)——而这件
“构造”,他们同所有人一样,也只敢用暗语提及。(多半是Deligne,一边把那篇无名之辈的论文
转寄给他们——想要的结果清清楚楚就在里面\fnote{802}{(*)}{(4月17日)论文里至少有一个非常接近的结果,尽管Beilinson--Bernstein(以及Brylinski-Kashiwara)所用的那套表述在其中并未逐字出现。细节见同日一条脚注(打字稿第1047页注(**))。}——一边必定向他们递了话:若非要给个参考文献,引Kashiwara的某一篇文章或许更稳妥些,哪一篇倒在
其次,反正没人会去查证到底。)对于那位亲自到场的无名之人,他们倒也当面许诺:会想着这
事的,会在Kashiwara那一处把引用改正过来。抱歉抱歉——我朋友Zoghman的倒霉史,确实老是那一套!

在这几位杰出作者的呈文里——由Deligne转交(他的信我刚才才照录,写于仅仅一个月之前)
——\textbf{Mebkhout的名字只字不见}。顺便说,Kashiwara的名字也不见(我已经瞥见那儿支起了一只耳朵
……)。不过,在呈文的最后一部分(证明 Kazhdan--Lusztig 的那一段)里,倒有一处偷带的两重
引注,引一件“\textbf{见于}[4]、[5]\textbf{的构造}”\fnote{803}{(**)}{读者不妨细品“\textbf{见于}……的构造”这一说法的含糊其辞:这件“构造”(或“对应”、或“关系”……)究竟出自谁手,问题被完全悬置;而仅仅六个月之后,在那场著名的讨论会上(见注“变戏法者”,n\textsuperscript{o}~75$'$),这个问题将以人们熟知的精湛手法得到“解决”:从Beilinson--Bernstein-Deligne的文章里,读者将获悉,那电报式的引注[4][5](在两处地方——那两处,想必那件构造恰好(托天之福)“有所陈述”)纯属客套,而“对应”的光辉父亲,果然是大家猜想的那一位……不过,撇开我刚提到的这套魔术手法不谈,引用一个新定理——又深又难——却只写成“见于……的构造”,仿佛那不过是件寻常“构造”,碰巧被丢在那里,作者们又碰巧选中它,拿来支撑他们辉煌的证明——这本身就已是一桩骗局。我在这里认出的,正是“SGA~4$\frac{1}{2}$--SGA~5”那场手术的同一副心肠:那场手术的要领,是(顺带)提一句“平展上同调形式‘有所陈述’于 SGA~4 与 SGA~5”(以及那位光辉作者不得不从中剔除的“废话矿渣”),随后挽起袖子,做作地开始“干\textbf{正经}数学……”。(5月25日)关于这种“新风格”,另见注“恭贺——或新风格”(n\textsuperscript{o}~169$_9$)。},\dmark\origpage{1073}这件“构造”(您准猜着了!)不是别
的,正是那位当值无名之辈那个从未被命名、其人更不曾被叫出名字的函子。参考文献[4]是Kashiwara的
一篇文章(临时的替身爸爸)。那篇文章里(Kawai-Kashiwara那篇也一样——它被扔进了盈亏账)自
然没有任何东西,哪怕远一点儿,像这几位作者所说的“构造”;况且那篇文章是1975年的\fnote{804}{(*)}{核实过了:就是前文引过的Kashiwara那篇文章,他在其中证明了他的可构造性定理——这个定理当然对定义那些“好上帝函子”起着关键作用(不过在1980年那阵抢购潮之前,除了Mebkhout,谁也不曾梦想过这些函子)。佯装把Kashiwara这个定理(无人打算跟他争)与“好上帝定理”混为一谈——后者无比深刻,层次完全不同——这是一桩粗劣的骗局。就证明而论,“好上帝定理”动用了Heisuke~Hironaka式奇点消解的全部威力。就“哲学”而论——这还要重要得多——它在同调形式主义内部架起了先前缺失的桥梁,沟通拓扑、代数与分析(算术尚待来日——前提是,我眼看着正在掘墓的某些先生,总有一天能找回他们健全的心智……)。},比莱苏什讨论会上某位无名之辈的报告早了差不多五年——正是那场报告让这位
Kashiwara先生开了窍:把“导出范畴”这个词说出口,毕竟不算傻事,而且可以依仗“强者即有理”这条再简单不过的法
则,把别人流汗挣来的功劳据为己有。至于参考文献[5],那是Mebkhout1979年9月在莱苏什讨
论会上的报告——恰恰就是在那场报告上,Kashiwara先生学到了:导出范畴原来是有用的——而且
除了用来坑害一位被他的大佬和前辈们晾在一边的无名之辈,还能派上别的用场……

\pmark\origpage{1074}同Brylinski-Kashiwara那篇文章一样,这里也没有任何东西,能让一个真正“不入局”的读
者起半点疑心:若不是那件新兴而天赐的工具现世——它被“见于……的构造”这句托词悄悄
藏起——这篇光辉的呈文根本不会问世。我还认出了那套屡试不爽的“溺毙鱼儿”术\fnote{805}{(*)}{关于这种号称“同化稀释法”的招数的其他例子,见子注“\textbf{正经}数学……”(n\textsuperscript{o}~169$_5$),即打字稿第885页脚注(*)。},号称“稀释法”:把要抹掉的那个人(却又想
给自己留“叫停”的余地,好在必要时说自己引过他……)\textbf{配对}到另一个与此事毫不相干、
或作用微不足道的人身上——仿佛在此处(在字里行间,却又明白无误地)说:我们把那个无名
小卒放在这儿(纯粹出于客套,也架不住他一再坚持),他和那件大名鼎鼎的“构造”(按新近确
立的共识,只能用暗语谈它,把它当作人人皆知的东西……)的干系,就跟1975年那篇文章的干
系一样多——那一年,“上流世界”还没有谁肯屈尊说出“导出范畴”这个词(顶多当笑话讲讲
……)。

\subsubsection*{(e) 冤大头买卖——或木偶戏}
\tocsubsub{(e) 冤大头买卖——或木偶戏}

我不后悔花这番功夫——既为我自己的备忘,也为哪位对这桩事感兴趣的数学读者——在这里
逐一清点围绕当值无名之辈那条定理的三桩前奏骗局。这三桩骗局,一桩是Kashiwara家的土产,一桩是
Brylinski-Kashiwara的土产(还劳了一位始终匿名的审稿人从旁襄助),一桩是Beilinson--Bernstein
的土产——而Deligne,则在幕后\fnote{806}{(**)}{(6月5日)Deligne“在幕后”扮演的角色,至少在第三段插曲里是明摆着的;第二段插曲,也有种种强烈的迹象指向同一方向。但看起来,Kashiwara早在1978年就已经在为他自己“打响第一枪”(就围绕Mebkhout著作的那些骗局而言)——当时(看来)Deligne还什么都蒙在鼓里。关于这一点,见本注 c) 部分(“入门费——或有为青年”),即打字稿第1050页脚注(*)。}。三桩在风
格上如出一辙,齐整得惊人,无须我再多着墨。这种风格,在我对\emph{葬礼}\fnote{807}{(***)}{关于这种风格,特别可看前已引过的注“恭贺——或新风格”(n\textsuperscript{o}~169$_9$)的结尾。}所
作的漫长调查一路之上,我早已领教到倒尽胃口;而在我那位天资最出众的学生——正是那位Pierre~Deligne——1968年的一篇文章里,它已经以刺眼的方式预先现出了形\fnote{808}{(****)}{见注“驱逐”(n\textsuperscript{o}~63)的开头。}。而这层情形也足以\dmark\origpage{1075}提醒我的记性:当年面对Deligne以及另一些在我看来天资出众的人,我抱着一种模棱两可、曲意逢迎的态度,
对今日铺展在我眼前的这股败坏,我也不是没有添上我自己的一份。

现在也清楚了:1981年6月那场“反常层会议”(Colloque~Pervers)的\emph{封神},离我们方
才清点的第三段插曲不过半年,绝非从天而降。说也奇怪,据我所知,这场讨论会是我离去之后头
一场、也是唯一一场专门(嘴上自然绝口不提,实际上却毫不含糊)用来把“Grothendieck数学”的
某个侧面重新发掘出土的讨论会——起因是一件猝然降临的新工具,谁料它竟无可替代。这件工
具,唯有在一种看待事物的路数里才活得下去——而时尚的种种“共识”,老早以前就把这种路数
归档为过时,还搭上几分隐隐的滑稽\fnote{809}{(*)}{至于这些“时尚共识”背后起作用的心理机制——它们与面对某种数学研究路数时某种“出自本能的”排斥反应彼此叠合——见前已引注“天赐的机缘——或封神”与“拒认——或召回”(n\textsuperscript{o}~151、152)。}。而仗着我那位天资夺
目的前学生的特殊天才,事情起了一场古怪的回旋:这条被他本人、被所有人弃绝的路数,竟在需
求的推动之下于事实中得到辉煌的证实;而恰恰又是借着同一场讨论会,这场证实成了机缘,把那
位已故而未被提名的大师——连同那位身后的学生(同样无人提名,他曾有幸(或不幸……)让这一大帮上流人物为他挪步前来)——彻底而最终地下葬。

% ---------------- p.819 ----------------
这场讨论会,可不是从天上掉下来的,决不是。我的朋友Pierre~Deligne有个特点(别的且不
论):他懂得等待,也懂得抓住有利的时机。围绕那道“奶油甜点”的三段插曲,加上对当值无名
之辈的提及早已被删削得几乎一干二净,分明向他表明:时机已经熟透,该神不知鬼不觉地、带着他那副
招牌式的笑容可掬与温文尔雅,把那件横竖都理应归他名下的东西收进兜里了。我猜他跟Verdier做过一番周密的串通:得让Verdier明白,如今时候已到,该锣鼓喧天把导出范畴和那份早已被
拒认的“父亲名分”从土里重新刨出来;与此同时,又把那个模糊的无名之辈、连同那位谢世多年的大
师,一并埋进舞台脚灯的光焰底下(就怕哪位不识相,偏偏记起:那些俨然“最新时尚”的美好
玩意儿,那位大师当年原也出过一份大力……)。

\pmark\origpage{1076}让Kashiwara给某条从未被点名的“好上帝定理”当个偷偷摸摸的爹,这一套一度还对付得过去
;而且恰恰也只在这样一个期间对付得过去:大家讲好了,那个定理既不会被点
名,也不会见诸文字。说到底,Kashiwara本人对这条定理也未必真有多么上心——这条定理,说不定他懂
的比Verdier本人还少;八成是机会凑手、积习相帮,他顺手就把它
揣进了腰包,仿佛出于无心。至于Deligne,那位懂得等待的人,他清楚得很:这条定理不会永远是
一条既无地址、也无姓名的定理。总而言之一句话:这是一条\textbf{在给自己寻找父亲的}定理,
要找还得找\textbf{配得上它的}父亲;而只有等到那“真正的”父名——按常理本就该属于他的那
一份(而且十二年前就该……)——成为举世公认、不容置疑的共识,它才够资格到光天化日之下
露面。那篇“反常”文章,同名讨论会的一颗明珠,正是朝这个方向立下的第一块界石,由头号当
事人以他惯常的娴熟手法亲手砌下。

我的感觉是:Beilinson和Bernstein多半还挺受用——想不到竟平白沾上了对那些层的父亲名分
(那些层也被称作“反常”——同样是叫错了的),而且绑在一起的还有一位至今这般显赫的黑帮
老大——可实际上,他们是被Deligne当了枪使,好在“万一……”之际给他充当“不在场证明”。
文章既然是这样写的,任何一位消息不够灵通的读者都只能认定:那件天赐“对应”的作者,不是
别人,当然是Deligne——尽管它从头到尾没被点过名,也没被明白说出过(反正人人都该早已知道
它……)。

剩下的,只有那一点暧昧的暗影(精心算计过的),藏在一句妙笔里:那条未被点名的“关系”,
“本该在这几篇笔记里找到自己的位置……”\fnote{810}{(*)}{关于这一点,见前已引注“变戏法者”(n\textsuperscript{o}~75$'$)。}。这正是“叫停”式的笔法:含蓄然而清楚地点拨一句——却又不
肯一字一句挑明——上述那条“关系”(既然通篇并无相反的说明),至少要归到那篇光辉文章的
\textbf{三位作者之一}名下,或者(退一万步讲)由三人共享。但同样明白的是:等到时机一到
(对那位懂得等待的人而言……),无论Beilinson还是Bernstein,都不会去向Deligne争一顶实际上已
然到手的父亲名分。想必是做成过一笔交易\fnote{811}{(**)}{会想到这样一桩“交易”,是联想到了两个类似的局面。其一是Deligne与Verdier之间的交易(或许是心照不宣的,但痕迹显然可见):后者“牺牲”掉 Lefschetz--Verdier 公式,把它划进盈亏账,以成全“SGA~4$\frac{1}{2}$--SGA~5”那场手术;作为交换,他却“收进”了过世者全部的“对偶”遗产,还外带导出范畴(一篇打折的文章)作添头。(详细的来龙去脉,见“公式”那组子注,n\textsuperscript{o}~169$_5$--169$_9$。)其二是Deligne跟一位被宣告已故的大师做成的那桩“交易”——反正此人早已从流通中销声匿迹,不愁他起而反应——事涉1967/69两个年头里合开的 SGA~7 讨论班:三年之后,这个讨论班被“对半分”掉,一半归那位已故者,另一半归Deligne和一位临时拉来的搭伙人。(详情例如见注“步步升级的插曲”(n\textsuperscript{o}~169~(iii)),插曲2.)这又让人联想到跟同一位已故者(他还蒙在鼓里)就那条(据MacPherson(Mac~Pherson)说)称作“Deligne--Grothendieck”的猜想做成的“交易”(见同一条前引注的插曲1):头一半归“邮差”Deligne——是他把一条猜想(此前一直由我那班上同调学生悉心保密)通风报信给了MacPherson;后一半归那位已故者,身份则是前者的“合作者”……}——\dmark\origpage{1077}即便不曾形诸明文,至少也彼此心照:Beilinson和Bernstein分得 Kazhdan--Lusztig 猜想,外加(为
了凑足分量——反正Brylinski-Kashiwara已经
% ---------------- p.820 ----------------
趴在上面了)对那些层的共同父亲名分——那些层,我猜想是出于众人的合议,才被叫作“反常”
的\fnote{812}{(*)}{见注“反常”(n\textsuperscript{o}~76)。};至于Deligne,则分得那条没有名字的著名“关
系”,只待那指日可待的一天到来——届时他的谦逊无须劳驾挪动一步——全世界都会管它叫“Deligne定理”。而这位未来的“父亲”鼻子也着实够灵,至少把这一点嗅了个分明:就这个孩子而论
(当初他宁可拒之门外,也不肯亲自动手把他生下来……),自己做成的是一笔“好买卖”
\fnote{813}{(**)}{这笔“好买卖”,我看同时又是一桩糟透的买卖;而且即便(甚至尤其……)在事事都遂了当事人心愿的情形下,也是如此:把宝贵的天赋与创造力,挥霍在扮黑帮的营生上。}。至于
Kashiwara,他的戏份儿已经演完:在那篇光辉文章里,就那件天赐的“关系”而论,再没有人提起他,正
如再没有人提起那位当值无名之辈。人多势众围攻一个,只要那是个无名小卒,好说;可地盘上一
旦清走了闯入者,那就各人顾各人……

% ---- 以下“(f)”节, 仍属原书 p.820 ----
\subsubsection*{(f) 演员列队——或黑手党}
\tocsubsub{(f) 演员列队——或黑手党}

那份才翻开三个星期的“家庭相册”\fnote{814}{(***)}{见3月22日那条同名注(n\textsuperscript{o}~173)。},想不到又添了几
张新面孔。看得出来,“这个家”丁口大为兴旺,我这副老朽之躯都快认不过来了,更何况时世也
已大不相同。这一回,按\dmark\origpage{1078}出场的先后,是\textbf{Kashiwara}、\textbf{Hotta}
\fnote{815}{(*)}{细心的读者也许会奇怪,在这份“演员列队”(在围绕Zoghman~Mebkhout著作的那场骗局兼蒙骗之中)里,怎么见不到Kawai的名字——他是Kashiwara那篇屡被引用的文章的合作者,其第4节肆无忌惮地劫掠了Mebkhout博士论文的第三章。(关于这一点,见注“五张照片(晶体与D-模)”(n\textsuperscript{o}~171(ix)),特别是第1005页。)Mebkhout一口咬定,不能把Kawai同Kashiwara装进“同一个口袋”(Kawai嘛,闭着眼睛跟着Kashiwara走也就是了……)。他向我形容Kawai是个有点儿跟不上趟的家伙,我有个印象,他几乎已经喜欢上这个Kawai了——总之一句话,这可是他的“好日本人”,在他那里,我去碰是万万使不得的!大概也正因如此,他才没有给Kawai写信(就像他给Kashiwara的另一位搭档Hotta写过信那样),去点破他与Kashiwara合写文章里的那些骗局,从而逼他公开表态,同他的搭档兼老板团结一致。}、
\textbf{Brylinski},还有Brylinski-Kashiwara投给《数学创造》的那篇文章的
\textbf{匿名审稿人}。一帮“硬茬儿”,这准保没错,反应招数训练有素;更何况只要事关坑害一
个无名路人,只消幕后大头目一个隐而不显的暗号,便个个心领神会、行动一致(甚至等不及暗号
……)。

我又一次认出了那副\textbf{黑手党}的做派\fnote{816}{(**)}{这种离奇的感觉,去年在注“讨论会”(n\textsuperscript{o}~75$'$)(猜得出是哪一场……)里就已攫住过我:那里一派收保护费的气氛,叫人恍然如梦,或者以为正“在看一部讲黑手党称霸某座遥远大都会阴暗底层的电影……”。这一回,在这趟穿行于当值无名之辈种种倒霉遭遇的漫游里,这种感觉又寸步不离地陪伴着我……}:在他们无可争
议的领地上称王称霸,而领地的核心,便是代数簇以及诸如此类的上同调理论。一帮才华耀眼而心
狠手辣的人物,脑瓜子个个无懈可击;在所谓“当值无名之辈”那场行动接连四段的插曲里,我一
路重新见识了他们的身手,直到“反常层会议”登上顶峰。除我方才点到的四位黑帮老大(其中一
位匿名)之外,还请诸君记起“硬核”的另外五名成员;这样一共九位,为埋葬\textbf{那位闯入
者,那个不是他们一伙的人}而全体出动。

有大头目\textbf{Pierre~Deligne}——那位永远懂得“下水”下得最少、进账却捞得最多的主儿。
有他的副手\textbf{Jean-Louis~Verdier},人称“大善人”——正是此公,主持过某位无名之辈某
篇论文的答辩委员会;又正是此公,当过那场令人难忘的讨论会的两位组织者之一,那场会肆无忌
惮地掠夺了同一位无名之辈。另一位主其事的是\textbf{Teissier},他同Verdier联名为那场令
人难忘的讨论会的令人难忘的文集,签下了那篇令人难忘的引言。跟旁人不同,他\dmark\origpage{1079}似乎不过
是跑了个龙套、当了个挂名替身,自己从中一无所得——唯一的乐趣,不过是讨得
他知道既有名望又毫无顾忌的那班人的欢心。最后\fnote{817}{(*)}{(5月25日)这个“最后”终究是言之过早——打那以后,又有别的帮派成员自行冒了头,进入了我的视线。关于这一点,见“光荣之日”注(n\textsuperscript{o}~171(iv))里打字稿第962页的脚注(*)。\par(5月30日)最新消息:又一名成员——R.~Remmert——刚刚被认了出来。关于这一点,见本注 c) 部分(“记忆的失灵——或新编历史”)。}还有\textbf{Beilinson}和\textbf{Bernstein}(对这两
位,就在这几页里
% ---------------- p.821 ----------------
我刚刚有过一番更充分的领教),被一根根看不见的提线轻轻地牵动着……

我还会等着,不急不躁,也不抱幻想,看未来还会给我们预备下哪些别的“反常层会议”,有整
个教团毫无保留的首肯作保,好为“科学”增最大的光荣,为“人类精神的荣誉”添彩。

\subsubsection*{c3.~根与孤独}
\tocsubsub{c3.~根与孤独}

\notehead{171$_{3}$}{}

(4月18日)一步一步跟完我朋友Zoghman那串倒霉遭遇的这第四天行将结束之际,我比去年更懂得了一些态度与心性——尤
其是冲着我来的那些——而直到去年,它们在我眼里还显得古怪。总而言之,凭着他那份自知分量的工作,他原以为进了“一个大家庭”,
有点像那位从不被人提起的已故宗师的家——的确,无人提起;可即便不提,他分明在场。谁承想,他
发现自己落进了一个鲨鱼的世界:个个彬彬有礼、甚至和蔼可亲,牙齿却毫不留情——转眼之间,他
带来的东西被剥夺得一干二净:八年孤独劳作的果实;然后有人让他明白:你,我们已经看够了——一个讨厌
的、不请自来的外人。换了他那个位置,没几个人能不留下创伤。我不知道他可曾对什么人敞开过他
的挫败,除了几声苦涩的暗示,又含混得反倒像在作证对他不利:一个怨气冲天的人,边缘上还有点不合群。

尽管没有人点我的名,我却俨然是这伙毫无顾忌、毫不留情的世界的“爸爸”;他实在没有什么理由
信任我。诚然,我们的初次见面——1980年,那时他还远在千里之外,猜不到等着他的是什么——曾
经打下过信任的底子;我分明感到,这个底子,顶着一切,一直保存到今天。说到底,他心里明白:我
纵然是一切鲨鱼的“爸爸”,也不会像他们那样行事。可\textbf{怨怼}是有的,这确定无疑;而且它乐于\dmark{}\origpage{1080}
摆出一副打心底不信任的架势——虽然(至少我感觉如此)那是“贴上去的”。

为自己认定的正当权益“斗争”,这容易——只要你属于一个团体,多小都行,同它声气相通。可那个
孤身对抗全世界的人、那个被排斥者、那个不请自来的外人,就像一棵被夺走了沃土的树。他身上的
力量帮不了他;那力量变成苦涩,反过来咬他自己,仿佛要同那个排斥他的全世界合唱一个调子。

当我手里捧起那本书——它为 Motif 掘墓复生,同时埋葬了让它们现形的工匠;那本书,署着一代
光辉人物(其中有我参与培养的)里最光彩的四位作者的大名——当我终于读到它(纯属偶然,因为
此前没有任何人觉得有什么值得特别提醒我的……),那一刻,在我结识数学家世界三十六年之后,
我第一次知道:\textbf{我孤身一人,对抗所有人}。过去八年里发生的许多事,忽然拼合起来,显出全部
意义。当一个人蓦地重新发现这份孤独,那滋味很是古怪。那一天我不得不缓一口气;此后一连几个
星期,一天一天地认清\emph{葬礼}的全部规模——一场与那部著作相称的\emph{葬礼}。

可这同Zoghman毫无共同之处——他的“自己人”在他还没真正扎下根之前,就把他“挂账”了。而命
运对我,是笑过的。多亏长者们接待了我(他们后来故去也罢、退休也罢、早就不再搞数学也罢,统
统不要紧)——多亏我年轻岁月里寻得的兄弟般的接待,我,得以在我自己选定的土壤里,“扎下根
来”。这些根,扎下去,长起来,随着岁月,变得深沉而有力。这些根,牢牢栽在一片不是“共识”、也
不是任何时尚的土壤里——扎得想必比任何以制造时尚、追随时尚为乐的人都深\fnote{818}{(*)}{如果说我在数学里也好、别处也好,从来不曾操心去追随或制造时尚,我知道,这正是我有幸在幼年便得以发展起来的那种深根的表现之一。既然我从一开始就在自身之中有了深根,投进我的大事业里的那份精力,就不会被补偿性的饥渴——比如领风气之先的饥渴,或保持并显得合乎必备之“调门”的饥渴——所分散。
关于我的童年和这些“根”(我想,我没有说出这个词),我在注“天真(阴阳的联姻)”(n\textsuperscript{o}~107)里说得具体。}。

\pmark\origpage{1081}说到底,我还“孤身对抗所有人”得起——想说什么就说什么,走我自己的路。

(5月25日)\fnote{819}{(*)}{下面两页,原本是作为注“……与意外之喜”(n\textsuperscript{o}~171(iii))的一条脚注写的。在哪里安插它们,我犹豫了一阵,最终决定收进眼下这条“根与孤独”里。因为这是“封神”全部注中唯一的例外:我试着从自己的经历出发,尽力去体会Zoghman本人如何经历了那些我为之作编年史的事件与处境。}不消多少想象力,就能体会Mebkhout的沮丧:他的核心结果之力一旦得到承认,他
就感到自己像一根麦秸似的被“扫掉”了\fnote{820}{(*)}{“扫掉”一语,借自Mebkhout的一封信(比正文所引那封早两天),相关段落我在这里照录:
“的确,[Kashiwara的]那个可构造性定理……让我爆发了。而且从那一刻起,像Deligne这样的人,一眨眼就
能找到我全部的结果,连‘好上帝定理’的各种形态在内,证明嘛,照你的说法,四个小菜的功夫。这就
说明了为什么这一切几天之内就被扫光了。”
我想,Mebkhout在这里把一个Deligne式的、默而不宣的“推理”,说得再确切不过:据他人\textbf{本来可
以}(而且\textbf{本应})凭他自己的手段、行囊与一切,“四个小菜的功夫”找到这些果实,所以就可以据为己有。这类推理——在类似的情境里,人们常常禁不住要做——唯一的别扭之处在于:\textbf{难的全是想
到}。而真正“想到”了的,是Mebkhout,绝不是Deligne,更不是别的什么人。创造不归“技术”管:技术,
在终于有人看见了无人曾看见的东西之后,顷刻间便可“扫光”一个局面,快过把它写下来。创造不在
“扫光”里,而在\textbf{看见}无人曾看见的东西之举里;用自己的眼睛看,不“追随”任何人。把这两者区分
开来——把创造之举,同摇一圈就转圆的那根曲柄区分开来——正是数学家操守的一部分。}。他给我写信(4月24日,在他新近来我家之后)说:“我花了八年,搭起
Kazhdan--Lusztig 证明所用的那些结果。他们用一个星期就把它证明了。”这一回,还是有一种羞
怯拉住了他,不肯把心里真真切切感到的那些话说到底;\dmark{}\origpage{1082}我在这里替他补上那句“没说出的”:而事情一旦做
成,“他们”便拿着那件别人在孤独里打造出来的崭新工具,彼此之间昂然招摇,同时让工匠明白:你,
我们已经看够了……

可这件事实在大得离谱,以至于当场,Zoghman还不大相信他自己那健全官能的见证——正如去年5
月2日,读到吕米尼讨论会文集时,我自己也迟迟难以相信自己官能的见证\fnote{821}{(*)}{关于这场讨论会(1981年6月),见注“不义——或归来的意义”或“光荣之日”(n\textsuperscript{o}s~75、171(iv))。说实话,去年5月头一个星期写下“送葬行列之七:讨论会——或Mebkhout层与反常”(n\textsuperscript{o}s~75--80),还不足以克服那近乎不可克服的惰性——在一个严格说来只有你一人使用这种见识的处境里,“相信自己健全官能的见证”何其难。直到五个月之后,终于“有血有肉地”面对了现实——我的朋友皮埃尔(Deligne)来到我隐居处看我——那股隐秘而顽强的怀疑,才终于消散。关于这一点,见注“义务已尽——或真话时刻”(n\textsuperscript{o}~163),尤其打字稿第782--784页。}。Zoghman是在去年1月读到这同一部
文集时——在 Kazhdan--Lusztig“总彩排”三年之后——才终于多少弄明白,实际发生的是什么事。

打击是可怕的,我自以为懂得——Zoghman当场以为他要把命搭上。幸好人这身子很结实——Zoghman
今天还活着,其间还结了婚,做了父亲……可我想,即便手捧“文集”的那一刻,他仍旧不能完全相
信。总有什么东西“卡着”。说不定,在我写这些行的时候,他仍旧没有全信。应当说,即便只按“理性”
或“客观”的字面讲,这件事也实在难以置信、实在巨大:直到今天,除了我(或许除了他,也未必
……),竟没有\textbf{一个人}敢相信自己眼睛、看见它——虽然它比一座主教堂还大!

可对一个被当头击中的人——被他的那些敬重的、什么都不缺的长辈,施以犬儒而\textbf{无缘无故}的
不义——这种事,只怕永远没法全信,只怕永远“\textbf{超出理解}”……而恰恰是这种事,能够毁掉一个
人的一生。它们之所以有这份\dmark{}\origpage{1083}毁灭之力,是因为一种幽暗的、绝望地被压下、却无可否认的知觉:
\textbf{毁掉}的\textbf{意图}——就这样无缘无故地,为了\textbf{取乐}:为了取乐于漫不经心的一挥手,碾碎那
对你珍贵的东西,碾碎(如果办得到)那构成你生命的实质与盐分的东西。正是这种“无缘无故”的
恶意里的反常快乐,才真正“超出理解”……

我想,Zoghman从来不曾真正对什么人谈过这些,大事之前也罢,之后也罢——至多几个单音节的
词,除了他自己,谁也解不开。单是 Kazhdan--Lusztig 那一幕,已经太大、太不像话,他不敢指望谁
会信。根深蒂固的共识,把最明显、最明白、最无可抵赖的事实,像麦秸一样扫掉。而这里说的,
是一件疼得那么切近、在他的存在里那么“一碰就疼”的事:只消想到,他若向谁袒露,那人可能拒斥这
份不速的信息,他面对“超出理解之物”的痛苦可能不被接纳——这个风险、这个概率,便膨胀成
\textbf{不堪忍受}之物;这种事,付任何代价也不去碰——哪怕为此当场憋死,要憋就憋……

两年前,他倒是对我“用单音节的词”说过。说不定,在他心底,他盼着我会听懂这些单音节——不只
按字面,而是同时听出他不敢高声说出的一切(也许连对他自己都不敢说……)。这盼望确实疯得
可以(在一个一切都疯得没边的处境里!);对我后来确凿得知的那些,我当年连想象的边都没沾
到。只能如此,谁叫我那时没有一份详尽的、原原本本的信息呢\fnote{822}{(*)}{(6月1日)更公平的说法或许是:以我这种有限的敞开与在场的状态,除非在极罕见的时刻,事情“不可能是别样”。然而我相信,我们人人都有“耳中之耳”,完全听得见没有说出的那部分——只不过,我们通常留神把这只耳朵捕到的讯息,排除在自觉注意的领域之外……}。而Zoghman那一边,也远不敢把这份信息
给我。这很疯;可这并没有妨碍他怨我。他总得怨一个人,怨一个足够亲近的、摸得着的人,好把“超
出理解之物”在他心里搅起的东西,至少卸下一部分在他身上,好让自己从那啃噬着他的东西里,解
脱那么一点点。

\subsubsection*{c4.~掠夺的空白授权——或“高尚行当”}
\tocsubsub{c4.~掠夺的空白授权——或“高尚行当”}

\notehead{171$_{4}$}{}
\dmark{}\origpage{1084}(6月2日)我给\emph{葬礼}画上“句号”、连同末尾那条注“自深处”(4月7日)一道的那份心满意足,到今天已满两个月——而这两个月来,我也一直在拼着老命,给\emph{葬礼}的最后一部分“做最后的润色”!这几乎是去年同一时节那出戏的原样重演——当时,我给那即将成为\emph{葬礼}第一部分的东西“收尾”,也是收个没完。
% ---------------- p.823 ----------------
和眼下一样,拖个没完的正是“最后一刻”——拖到我连吃喝都顾不上,最顾不上的还是睡觉。就这么一直拖下去,直到我的身体油尽灯枯,宣布弃权。那正是整整一年之前(差不了几天),我不得不撂下一切,足足三个多月,一门心思只顾从一场急性衰竭里爬出来\fnote{823}{(*)}{关于这一段,见注“事件——或身体与精神”(n\textsuperscript{o}~98)。}。可这一回我多了个心眼,格外当心,不再走同一条老路。这条命,我还想留着呢……

这一回,又是那场没完没了、一弹再弹的“调查”。我原先估计,撑死也就一条十来页的注,取名“四种运算”,把去年那阵旋风式调查的结果“归置归置”、做个总结。可如今,调查卷土重来、干得热火朝天,眼看要满四个月了;十页变成了三百页上下,而且这还不算(完全)完事!我再不敢作什么预言——从9月底重新动工算起,已经到第九个月了,我一直在“马上就要完工”!只有等到最后一批笔记全部打成清稿、通读修改完毕、交去复制的那一天,我才会知道它\textbf{真的}完成了。(到那之后,剩下的就不再是我的活了。)我唯一知道的,是我巴不得赶紧到那一步,就像巴不得看到一场漫长而磨人的疾病有个尽头;而我必须尽我所能走到终点,不被种种想象出来的期限推搡着走。只有等该\textbf{在此时}看见和说出的都已看见、都已说出,我才会停下来喘一口气。

害我吃了最多苦头的,正是那该死的“\emph{封神}”——我说不上为什么。这几样“四种运算”,是《收获与播种》全书里唯一一处写得磕磕绊绊、零敲碎打、一路费力挣扎的部分——按说本该是现成的活儿:不过是一次简单的“归置”,没错;没有任何东西会以触到\dmark{}\origpage{1085}要害的方式,把我这个人牵连进去、拿我这个人问罪,从而动员起抵抗的力量,一个“\textbf{摩擦}”。可天知道这摩擦有多少,而且在\emph{封神}这里,比其余一切加起来还多!它打哪儿来的?

早在“操演”那条注里,写起来就已经够费劲了。没完没了地拉长,就是从那里开始的。光这一道“运算”,最后就足足写了八十页——而如今,一个月之后,\emph{封神}已经写出了整整两倍的篇幅。话虽如此,除了“操演”里兴许有那么几页(边上带点儿太浓的“侦探”味儿……)(在那里,我兴许一头扎进一桩离谱透顶的“骗局”的絮叨细节里,扎得比必要的还深……)——除了这件写得眉目详尽、对“不在局中”的读者来说多半还有点磨人的“凭卷宗办案”的活儿之外,我并不觉得我最终码出来的那一摞摞上百页的东西是画蛇添足,是回锅的冷饭,是头发丝里切四瓣的功夫。让我一直提着气撑下来的,恰恰是\textbf{新的实质内容}大量地、出人意料地涌到我身上,它们我无论如何都得安置妥当,由不得我愿不愿意——其中甚至有,可不是嘛,数学的实质内容!有好些时刻我几乎招架不住:一时要落到白纸黑字上的东西实在太多,而且刻不容缓——这些事情一个个都还热乎着,甚至烫手,可你偏偏只能一件挨一件地去料理……

这样的丰富,本身就自带着一股强劲的工作推动力,丝毫不是那种会惹出“摩擦”的东西,恰恰相反。这摩擦,可以肯定,并非来自实质内容本身,而是来自我投进这份工作里的那股唯我的执念之力。说来像是悖论:恰恰是我那股“赶紧收场”、“一把甩到桌面上”的不耐——急着说出我该说的这样那样的事,而它们就发生在此时此刻,同我息息相关、触我至深——正是这种不耐(我想),制造了摩擦,制造了能量的涣散。摩擦是分裂的征兆,是几股力量朝相反的方向对拉,每一股都因对方的抵抗而愈激愈怒:
% ---------------- p.824 ----------------
一头是“赶紧收场”的急切——趁我还在打磨这块料,就想把它“甩手”交出去;另一头,是要求把眼下这一刻让我隐约望见的东西走到底,不满足于差不离,不让人推搡,也不把自己关进一个赶着“结项”的“计划”、一份预先定死的“日程”里。我很清楚:一旦把“意外”排除在外——这个防着我原地打转的搅局鬼——我的工作就失去了品质和意义。它\dmark{}\origpage{1086}就沦为“文牍差事”。这些年下来,我对这处看似不算什么、实则是一切的“小小差别”,已经变得非常敏感。这样的大转弯,如今偶尔还会启动,总在沉重已极的时刻——但从来没有维持多久。一旦走上这条路,那孩子把一切都扫到一边去——再想撑下去也是白搭。工作本身的兴致,那份\textbf{渴望}——它全然不同于攒页数的饿劲儿、不同于抢着画句号的急色——霎时间无影无踪,你便傻乎乎地守着涂黑纸张的活计。到那时,真就再没意思了——剩下的只有一件事:把航向拨正,而且马上就拨!

工作里总有一种\textbf{不耐}(我的一个老相识……),不停地把我往前拽。我觉得,这和自我跟这几样“四种运算”较上劲以来沉沉压在我身上的那种不耐,不是同一种。那另一种不耐,不是压下来的重量,而分明是一种往前拽的力量。它是胃口好的征兆,不是倦怠、疲乏或餍足的征兆。它不是累积的不耐,不是收尾的不耐,不是赶着把一份“计划”“结项”的不耐,而是要认识我面前那个即将向我把底交出的未知者的不耐。是一个赤身的孩子,独自站在无边的大海面前,急于纵身潜入其中去认识它……\fnote{824}{(*)}{这一形象,已见于注“孩子与海——或信与疑”(n\textsuperscript{o}~103)。}

不过,现在该回到我朋友Zoghman那串倒霉遭遇的叙述上来了——就在这条预定作为\emph{封神}最末一笔的注里。前面说过,这段叙述,Zoghman本人从来只是东一块西一块地零星递到我手里:随着信件、电话、会面,碰上什么算什么;可以肯定,我这边思考的推进和\emph{葬礼}的写作都因此受了影响,至少在专写我朋友那番坎坷的那一部分里是如此。如今我更懂得这种欲言又止的意味了,因为对“受害者”角色的一切执着(去年我还自以为从中看出几分)早已烟消云散(就算它当初确实在过)。我这边,想必也有过那么一些时刻,多少有些饱和了,态度上是那么一副“求您别再倒了,行行好!”的架势。这八成也没给他鼓劲儿。说实在的,那句“这儿是日本人、那儿是Kashiwara”的老调,Zoghman想必唱了总有四五年了,叫我不胜其烦;他跟他们之间确实领教过一些事,这不假。可我心里明白:他之所以领教过那些事,他的成果之所以就这样被交出去任人掠夺——近乎官方式地:“来吧,好人们,放量地拿,千万别客气……!”——那\textbf{不是}因为几个遥远\dmark{}\origpage{1087}的日本人。这\textbf{是因为“自己人”}:是“小家庭”\fnote{825}{(*)}{(6月16日)关于这一点,Mebkhout特意要我强调:他早已完全不再认同于这里所说的那个“小家庭”。}里的人——都是我们自家阵营里的体面人物,而他从不点名,除非引用他们的著作,并带着配得上他们崇高声望的全部敬意。

至于Kashiwara之流,我可再也不想听人提起了!眼看事情卡住了,Zoghman这时拿出他的智慧和耐心,索性撂开不谈,却并没有因此放下对我工作的关切,也没有停止在这里那里默默地、实实在在地帮我一把。

直到4月初他最后一次来我这儿,我才终于看了看那个“日本包裹”。起初还有点勉强。我原以为,自己准会陷在一堆解不开的超技术纠葛和读不懂的文章里腻烦得要死(没准儿还是日文的……),反正我是永远不会去读的——结果不然!事情简单得像打声招呼——有点像巴黎地铁(或者该说东京地铁)里“扒手的故事”。说真的,还挺有趣(至少,被摸走钱包的是别人的时候……)。
% ---------------- p.825 ----------------
这一来,Zoghman和我之间的僵局也松动了,他的倒霉史,这里那里、一道道闪电似的,也让我听到了些零头碎块。原先我多少按“技术情报卡片”的样式记下的几段插曲,如今被鲜活的临场追忆充实起来;恰恰是这一类东西,在科学文本那副纹丝不动的“立正”姿态里,甚至在同事之间的通信里,似乎永远被禁绝——您总不会想看吧!在“四种运算”里,我甚至不得不使劲摇晃自己,免得恰恰又落回那种腔调——“侦查结论”式的腔调(甚至,“诉苦清单”式的……)。Zoghman交出来的这些“碎块”,帮我从里面走了出来,帮我同一种活生生的实质保持接触。

Zoghman离我家的当天,我就重新拾起\emph{封神}的活儿,趁他讲给我听的东西还没凉,再写它一条两条子注。这就有了那几条注(或者说子注——写得多了,我自己也数不清了……):“眼光的绽放——或闯入者”、“黑手党”(后来我把它又分成七部分,每部分起了名字),还有“根与孤独”。我把整个儿火速寄给了他,请他在我送去打字之前提提意见。到这一步,我觉得自己多少有点在替他说话,我要确证:凡是我根据他讲的转述出来的东西,都得到他毫无保留的认可。他随即寄回了详细的意见(4月22日与24日的信)。

\pmark\origpage{1088}这些意见里,不乏这样的“碎块”,为我那些注里偶尔显得有点瘦骨嶙峋的事实骨架,敷上了鲜活的血肉。

也正是这样,我才知道Zoghman当时在场——就在那个令人难忘的1980年4月22日,Goulaouic--Schwartz讨论班上。就在那一天,Kashiwara把“好上帝定理”当成自家出产的定理当众宣布——而这条定理,是他几个月前在莱苏什讨论会上从Mebkhout亲口听来的\fnote{826}{(*)}{关于莱苏什讨论会与Goulaouic--Schwartz讨论班上的这段插曲,见注“黑手党”(n\textsuperscript{o}~171),部分(b)“最初的麻烦——或太平洋彼岸的黑帮老大”。}!这胆子实在大得没边,何况Mebkhout人就坐在底下,难怪听起来像天方夜谭。Mebkhout当场没有爆发(我直纳闷他是怎么忍住的……)。他客客气气地等报告讲完,“为的是公开抗议这种做派,向他提起莱苏什的那场会议、提起他的那个问题\fnote{827}{(**)}{指的是1979年9月莱苏什讨论会上,Mebkhout报告结束时Kashiwara所提的那个问题。关于此事,见上一条脚注中所引的那条注。}。Goulaouic请我私下解决我们之间的事。大厅几秒钟之内就空了”。

这便是那些“碎块”中的一块,由这段惜墨如金的叙述交了出来。后来我又在电话里听到些细节。这件事值得停下来细看。关于八十年代数学界的风气民情,它透露的东西可不少。这里涉及的,已不是某位心狠牙长的“黑帮老大”的心态——科学界传统价值败坏的一种极端症候——甚至也不是那班体面风光、处处周全的“建制中人”——在他们身上起作用的,是护着“自己人”的阶级本能。在这里,是整个大厅眨眼之间一空——转眼间一个人都不剩\fnote{828}{(***)}{这幅景象在我心头不可抗拒地勾起一连串联想,同三年前我亲历的一个场面如出一辙:那是一场Bourbaki讨论班的末尾,承人们好意赏我十分钟,让我谈谈某条坑害外国人的凶恶法令。关于此事,见“我的告别,或:外国人”一节,n\textsuperscript{o}~24。}!你们自己私了吧——我们嘛,什么也不想打听……

我真想知道,Goulaouic和那场讨论班上其他安安稳稳的听众,脑子里当时到底转过了什么念头——在那儿做报告的,可是一位声望斐然的外国讲演人(讲的题目,我相信他们当中没有谁真正熟悉)。这场插曲,说到底,本来够人好好想一想。但我\dmark{}\origpage{1089}恐怕他们谁也没费这份心,我更愿意猜想:全体一致,大家争先恐后把这段难堪的插曲忘了个干净。不过说到底,只要肯费点心思想一想,而不是拔腿就逃,这桩阴暗的事情里,毕竟还是有\textbf{一件}事是清楚的。Mebkhout的语气和措辞(何况他们认识此人,至少在各处讨论班里同他打过照面),使人几乎不容怀疑:
% ---------------- p.826 ----------------
这件事里必有\textbf{一个骗子}——不是Mebkhout,就是Kashiwara。

当然可能,他们在心底里早已提前断好了案:Mebkhout在编故事——那位尊贵的来访客人会去掠夺一个无名听众,这叫人怎么想象!这便等于说:面对一个无名之辈,名人无论做什么,都在一切嫌疑之上。这便是\textbf{掠夺的空白授权},交给有名望的人,用来对付求告无门的人。他要说的话,不会有人听:“你们自己私了吧!”

要不,他们就是一头埋进了怀疑里:谁知道谁说真话、谁说假话?(更何况,只要捂起耳朵,就更不知道了!)不错,一个Kashiwara,当着事主的面公然掠夺一个不知名的陌生人,这份蛮横的厚颜,看上去几乎难以置信。可说到底,还有更难置信的事:一个无名的陌生人(他们全都认识他,而他也从未靠骗子的把戏、靠那份胆量在他们眼前露过一手……)竟敢当众指控一位Kashiwara犯下粗暴的剽窃——假如他要说的纯属凭空编造……再退一步,就算他讲的或许有根有据,拿一句“你们自己私了吧!”把他打发到九霄云外,这依然是掠夺的空白授权。这就好比,一个人在大街上被一群穿燕尾服的恶棍洗劫,喊一声“抓贼呀!”——人们却冲他喊:“你们俩自己私了吧!”

况且据说,长期以来,纽约和其他美国大城市的下等街区,事情就是这样办的——在那里,谁也不想跟当地称王称霸的黑手党找麻烦。反正如今,数学界里事情就是这样办的(从什么时候起,我说不上来),在那些被当作“上等街区”的地方就是这样办的——比如Goulaouic--Schwartz讨论班\fnote{829}{(*)}{在这里我可以高兴地指明:那桩令人难忘的插曲发生的那一天,“他的”讨论班上,Laurent~Schwartz本人并不在场。后来他是否听人说起过,我就不得而知了。},或者在那班“搞”代数簇上同调的名流当中。

\pmark\origpage{1090}若按理性的说法、照字面来理解,“你们自己私了吧”这句话简直近于低能——何况无论怎样,两造之中必有一方是不老实的,这一点再清楚不过。在心理层面上,这句低能的话所表达的,是在一种被感到“碍事”的处境面前的一种\textbf{卸责}。这也是对下述明摆着事实的蓄意无视:遵守数学家这一行当基本伦理规则的问题,绝不是一件纯属“私下”的事,可以由僭取了蔑视这些规则之权的一方,同吃了亏的另一方,关起门来了结。这是一件\textbf{公共事务},一件关系到\textbf{每一个}数学家的事。

正是靠着普遍的冷漠,靠着每个人在承担个人责任时的惊慌失措,黑帮心态和“反常层会议”(Colloque~Pervers)那样无耻透顶的勾当,才能在科学界肆无忌惮地兴旺繁盛。一些人的惊慌,同另一些人的无耻,是\textbf{同一种腐败}的正反两面。1980年4月22日那一天拔腿就逃、捂着耳朵的那些人,同样促成了那场令人难忘的讨论会在次年的\emph{封神},其功劳丝毫不亚于那班凭空搭起这场辉煌大骗局、又神气活现跑去里面招摇过市的黑帮老大。

(6月3日)也正是在Mebkhout最后一次来我这里的时候,我从他那里听到了一些很有教益的细节,涉及同一场辉煌讨论会的某几位与会者,以及他们各自身上竞相盛开、一朵赛过一朵的“新风格”。我有机会翻了翻那次会议的报告集第二卷,里面有Verdier的文章、Brylinski--Malgrange的文章,又扫了一眼Laumon的博士论文(用一双比初次收到它那天更有经验、也更不分心的眼睛)。这篇论文其实是同N.~Katz合作的工作。关于这些工作里遵循的“新风格”,我在注“光荣之日”(上帝知道,这个名字它受之无愧……)的那条长长的脚注(打字稿第962页)里写了几句评论;在那条注里,为了更细的情节,我还指向了眼下这条注
% ---------------- p.827 ----------------
(当时它还没写出来)。说过的话,就要兑现!

Mebkhout给我讲过,他有怎样的荣幸和便宜,得以两次向N.~Katz谈起自己关于对偶的想法、关于连续系数与离散系数之间联系的想法。第一次是在雷恩的$p$-进分析讨论会上,1978年7月。当时他\dmark{}\origpage{1091}在“小圈子”里讲解了他的定理:复解析空间上D-模的整体对偶定理——那条为Serre对偶和Poincaré对偶加冕的定理\fnote{830}{(*)}{关于这条定理,注“作品……”与“三座路标——或天真”(n\textsuperscript{o}~171(ii)、(x))里都有谈到。}。在场的有Katz,有Illusie——正是在\emph{葬礼}里已经不止一次谈到过的那两位。Illusie呢,像他惯常那样和和气气、客客气气,觉得这真是漂亮极了——大致是这么句话\fnote{831}{(**)}{况且这本身还是一种分文不费的“客气”。两人的反应路数彼此不同(Illusie是“阴”,Katz是“阳”),实质却并无二致:只要是“无名先生”(Monsieur~Personne)说出来的,那就从这只耳朵进、打那只耳朵出!关于这一点,见注“蒙骗”(n\textsuperscript{o}~85$'$),尤其是打字稿第351页上我对Illusie的那几句观察。}。至于Katz——我猜,他在那儿是平生头一回听人谈起D-模(那阵子,这东西还远不是那场令人难忘的讨论会之后的那种大时髦)——他只干巴巴地撂下一句“这早就有人知道了!”,随即转身就走。既然跟他说话的,不过是个不知打哪儿冒出来的“无名先生”——而他N.~Katz,就在同一年,还要在几千位尊贵同行面前发表演说,向新科菲尔兹奖得主Pierre~Deligne致敬呢……——那当然只能“早就有人知道了”。

第二次,是1979年9月莱苏什讨论会之后不久\fnote{832}{(***)}{关于莱苏什讨论会,以及Kashiwara在Goulaouic--Schwartz讨论班上的那桩骗术,见注“黑手党”(n\textsuperscript{o}~171),部分(b),“最初的麻烦——或太平洋彼岸的黑帮老大”。}。Katz那时在法国高等科学研究所(IHÉS)。Katz在$p$-进微分系统方面的功力尽人皆知,而Mebkhout心里明白,这同他刚在莱苏什讲过的“好上帝定理”多少有些干系;于是他专程赶到法国高等科学研究所,把他在莱苏什的文章送上门去,同Katz谈自己的想法和结果。有了雷恩那场接待打底,可以说他这人是有点韧劲的,竟不知疲倦!反正,这一回的戏码跟上一回差不多。Katz又一次居高临下地接待了这个不知打哪儿来的无名之辈——此人竟敢第二次跑来催问,而且没准儿连招呼都没事先打。人一旦成了大人物,有时就不知道该怎么躲开这些缠磨人的家伙了……

一年之后,只消这同一批想法——由一个无名的陌生人,在孤独中久久揣着、慢慢养熟——跟着那场Katz先生遗憾未能赏光的辉煌讨论会的余威,被四处吹得震天响,说成是某位Deligne(或者某位Kashiwara,已经说不清了……)最新挖到的宝贝,它们便\dmark{}\origpage{1092}在这位大人物眼里陡然兼得了分量与斤两。替他把来龙去脉讲解清楚的,准是Laumon——Deligne最出众的弟子之一。而这位Laumon,对这些想法的来历也一清二楚,而且是第一手:因为正是那个无名的陌生人亲口告诉过他。可弟子以踏着老师的足迹为荣;而老师呢,早已明明白白、毫不含糊地示范过:对那个被判给沉默与湮没的人,该采取什么样的做派。

台前灯火的荣耀,归Deligne们,归Verdier们;Brylinski、Katz、Laumon诸公,也都赶在好时辰跑来,好从中分得自己的一份!乐声与吹吹打打归他们,感恩戴德的群众的欢呼也归他们——这群人欢天喜地蜂拥而至,庆祝这落入“新主子”们之手的“高尚行当”。

\subsubsection*{墓后尾声——或洗劫}
\tocsubsub{墓后尾声——或洗劫}

\notehead{171$'$}{}

\pmark\origpage{1093}(6月14日)直到一个月以前,我始终觉得,“葬礼”的精神所及,仅限于我间或称作数学“上流社会”或“大世界”的那个天地,尤其是那个天地里我素来出入、自己也身列其中的那些圈子。在朗格多克科学技术大学(USTL,蒙彼利埃)——它十二年来一直是我的挂靠机构——我不曾察觉到任何排斥的迹象,任何故示轻慢、失礼乃至粗暴的姿态,与那场十五年来方兴未艾的“葬礼”同一走向\fnote{833}{(*)}{在注93(第396页第3段)里,我尤其表达的就是这层意思。}。一桩新的事实刚刚闯入这幅安宁的图画,以一种严酷的方式改写了这幅图画,也改写了我同我的挂靠机构之间的关系。

% ---------------- p.828 ----------------
依照积习已深的机制,我起先并没有想到要把这桩新近的插曲收进我的见证《收获与播种》——乍一看,它出现在这里,仿佛“汤里落进了头发”,好不是地方。是顶住了几分切实的抗拒之后,我才终于承认:把这段插曲按下不表,便是有负于我的见证的精神。诚然,这段插曲余温犹在,何况我这一次“承受”得相当吃痛——这倒又给我方才影射的那些“积习已深的机制”平添了一分力气。可这一次,我领教这桩事故那些雄辩而不合时宜的教训时,感受如此鲜活,这本身也是一个征兆:它触及我至深——触及我的职业活动,触及我同我所隶属的职业圈子的纽带。所以,这恰恰属于《收获与播种》立志要作深入见证的那类事情:不允许有任何“保留的角落”,是我禁绝自己去触碰的——无论那“谨慎”是对我自己,还是对任何人,都用错了地方。

此外,在我对“葬礼”的思考这个更专门的框架之内,我感到一目了然:两者之间有着直接的关联。有可能,这类关联并不是单纯的因果关系:当地的某些同事,或许是终于把“葬礼”当真看在眼里,于是认定,从今往后他们也不妨“照样来上一场”。即便果真有这样一层因果,依我看,它触及的也只是事态一个附带的、偶然的侧面。然而另有一个更其本质的侧面,也是最令我震动的,无论在大写的“科学大世界”里,还是在一所外省的平凡大学里,皆无二致:那就是学术与大学环境中某种\pmark\origpage{1094}\textbf{退化}——也许是空前的退化:人际关系的品质在退化,礼数与尊重他人这些最起码的形式在退化,科学伦理亦然;而科学伦理,本就与尊重他人、尊重自己牢不可分。因此,下面这些页,读者不妨看作对“一个时代的风俗画”——在数学环境里,想必是一个时代的末日——的一份贡献(在围绕“葬礼”的思考一路提供的众多贡献之外)。

与其在这里把事件的经过或详或略地重述一遍,我更愿意转录四份\textbf{文件}——它们同样能把事件描述清楚。即:
\begin{enumerate}
\item 一封“致USTL数学任教同事书”,日期为5月28日,我在信中向他们通报了某种情况,并表示希望在全体大会上加以讨论;
\item 一份“答复”,出自Charles夫人——USTL数学楼的房舍负责人——形式为5月30日的一封通函,抬头具我的名,而实际上寄给全体数学教师;
\item 一项决议,由UER~5(第5教学与研究单位)全体大会于6月6日表决通过,大会议程为:“关于Grothendieck教授办公室搬迁的情况通报与讨论”;最后则是
\item 一封“致数学楼旧日同事书”,日期为次日6月7日。
\end{enumerate}

% ---------------- p.829 ----------------
我没有把两封信列入这些文件:一封是我5月21日致Charles夫人的信(第一份文件里提到它),另一封是我致USTL临时行政负责人Cano先生的信(这份文件里提到它,第四份文件——即“一场误会的尾声”——里也提到它);比起下面转录的文件所含的内容,这两封信在我看来没有提供任何新的信息。

对Charles夫人的信,我只作一条评注(“事实上很难联系上他”——这里的“他”,指的是区区在下,而这封信名义上正是写给我的):蒙彼利埃寄到我住处的信,一天即可到达;而多年以来,我离开住处,唯有去USTL的时候。

% ---------------- p.830 ----------------
\begin{quote}
\pmark\origpage{1095}朗格多克科学技术大学\\
数学研究所\\
数学楼里的一场洗劫\\
致USTL数学任教同事书\\
Alexandre~Grothendieck 撰

\hfill 蒙彼利埃,1985年5月28日

亲爱的同事:

上星期,我托UER的一位女秘书去我四楼的办公室取一份文稿,才得知那间办公室已被搬空,我的东西荡然无存——这一点我就在今天亲自核实过:只剩下一片光地板。从没有人通知我:我的办公室竟可以不经任何手续就被征用。因此我也就无从对这一举动表示同意,更谈不上准许任何人在我不在时进入我的办公室、触碰我的东西。当天我就给UER主任Lefranc先生(Lefranc)打了电话,把情况告诉他;此事(看来)出自Charles夫人的主动,这通电话似乎也证实了这一点。我向Lefranc先生言明:这种做法令我震惊;以如此粗暴的形式调换办公室,要我同意,绝无可能;我期待我的东西在最短期限内物归原位。他向我保证,会去办妥该办的事。

也就在同一天——5月21日,星期二——我写信给Charles夫人,告诉她:我把我办公室这场来得不是时候的“清空”视作滥用权力,并且把它感受为一种暴力;我期待她作出原原本本的解释,期待她毫无保留地道歉。否则,我将把问题提交大学理事会,由它裁定:对USTL的一位教师使出这类手段,是否应当被当作理所当然之事。

今天来到USTL,我得以证实:Charles夫人认为无须回复我的信(信的副本我也送交了Cano与Lefranc两位先生)。Lefranc先生同样认为无须就下述事实向我作任何解释:自从他向我保证会办妥该办的事、让我的东西回到办公室以来,一个星期已经过去,我的办公室依然空空如也。那些被席卷而去的东西如今安在,他也好,Charles夫人也好,都觉得无须告诉我。我辗转从几位秘书那里得知:这些东西存放在她们之中某一位的办公室里。此外,我曾在会议室里与Charles夫人打了个照面,她向我担保,她不过是执行UER主任Lefranc先生的指示,并请我就此事去找他,说这事与她不相干。在事态了结之前,Nguiffo~Boyom先生慨然愿意与我合用他的办公室。

% ---------------- p.831 ----------------
\pmark\origpage{1096}也许,只有我一个人觉得这里面有什么不对劲——一种暴力,一种蔑视;诚然,被人摆出这般架势、不经任何手续扫地出门的,也的确只有我一个。(倘若除我之外还有谁认为,这可不是他愿意在USTL工作的那种气氛,那就请他站出来让我知道,我会由衷高兴……\fnote{834}{(*)}{不言而喻,这样一个举动,对我来说只有一种情况下才有意义:大家明白,它使签署者本人担责,并准许我公开提及此事。})。至于我,我认为,在这场“误会”(姑且借用我一位同事那个可爱的委婉说法)之后,召开一次UER会议,让主任Lefranc先生和Charles夫人有机会就他们的意图与动机作出交代,也让UER的教师们说说,他们是否认为这类做法属正常(当它们落到别人头上的时候……)——这算不得什么奢侈。

我在USTL十二年,常常有机会领教Lefranc先生助人时的善意、热忱与高效——对此我心怀感谢。正因如此,眼见他甘当他人手中的工具,纵容一种专断与蔑视的气氛确立起来,我要收回对他的信任,遗憾也就更添一层。从现在起,我请求他要么切实承担UER主任之责,要么辞去职务。我也请求Charles夫人辞去她那UER“房舍负责人”的职务——这套职务,她滥用得津津有味。

敬候您的(或你的)答复

\hfill Alexandre~Grothendieck

又及:我天性乐于助人。去年,应Lefranc先生之请,我曾同意与Lapscher先生(Lapscher)对换办公室;而不久之后有人告诉我,他随后改变了自己的计划。不言而喻,我的同意,并不意味着我准许对我办公室来一场洗劫——无论在当时,还是在任何别的时候。
\end{quote}

% ---------------- p.832 ----------------
\begin{quote}
朗格多克科学技术大学

\hfill 1985年5月30日,星期四

\underline{数学}

\pmark\origpage{1097}J.~Charles夫人 “数学研究所房舍负责人”\\
致A.~Grothendieck先生,数学教授。

\begin{center}亲爱的同事:\end{center}

1. “数学研究所房舍负责人”的“工作”,从何处开始,到何处为止?这位“负责人”受理数学教师们提出的申请:\\
——或是为一位新任的教师(或研究人员)安排办公室\\
——或是为一位已有办公室的教师(或研究人员)另行安排办公室。

后一类申请,一般都举有工作上的理由:把同一研究组的成员归拢到一处。

这位“负责人”于是着手研究各种可能,首先是同UER~5的主任一道研究——数学研究楼的房舍,其正式委派的管理人,正是由USTL校长指定的那位主任。随后他同有关人士一起寻找可行的方案;调整须待各方谈妥之后方才实施。

2. 近年来照此办成的事情:\\
——几何研究组成员的归拢\\
——力学研究组成员的归拢

3. 这项“工作”中遇到的困难:\\
——实际接触到的每一个人,几乎都觉得自己的办公室归自己“所有”\\
——要强迫什么人“更换”办公室,看来绝无可能。

4. 我接到的最近一份申请,以及为所提问题寻求“方案”的经过评估:\\
——Lapscher(LAPSCHER)教授提出的申请:把Lapscher先生本人、他秘书的办公室、Micali先生归并到同一层。\\
——设想的第一种方案:在三楼与四楼之间交换办公室,好让“申请人”们归并到四楼。这次交换尤其涉及Grothendieck先生与Therond(THEROND)先生。Grothendieck先生经UER~5主任亲自接洽,\dmark{}\origpage{1098}向后者言明:\textbf{办公室在哪儿对他全然无所谓,只要有一间就行}。相反,Therond先生一度表示同意,随后却拒绝了任何搬迁。\\
——设想的第二种方案:我随后请Lapscher先生自己去联络他的同事,以提出另一种方案;这一点UER~5主任向他确认过。他的奔走情形,他不时向我们通报:五间办公室的“占用者”都同意作一次轮换;Grothendieck先生的同意,则出自他同UER~5主任的那次谈话。\\
——第二种方案的实施:UER~5主任得知这一致同意之后,为拟议中的办公室调整开了“绿灯”。

% ---------------- p.833 ----------------
Lapscher先生曾向我谈起钥匙的问题——搬迁启动而未完结的那段期间怎么办;我向他指出:\\
——新钥匙八成一把也没有,\\
——这场搬迁本可由全体有关者一齐动手,几个钟头做完,拖长它,在我看来并不可取。

Lapscher先生随后告知我:Grothendieck先生办公室里的物件,已被搬进他未来的办公室;这件事办成之前,未能先联系上Grothendieck先生。

应当指出:Grothendieck先生住处远在蒙彼利埃之外,目前又借调在法国国家科学研究中心(CNRS);事实上,要同他取得联系非常困难。

5. 以“负责人”的身份,我对这件或许可以称作“冲突”的事情的印象:\\
——我曾有机会向Grothendieck先生言明:我以UER~5的名义行事,不能亲自答复他的信;因此他应当向UER~5的主任索求答复。继这封发给大家的第二封信之后,我认为自己必须走出曾自我施加的“克制义务”。\\
——至少在搬动有关人士的东西之前先通报一声,这在我看来本是可取的\\
——把搬迁压缩到至多半个工作日之内,这在我看来本也是可取的。\\
——所设想的方案在我看来是妥当的:它丝毫改变不了有关各位的办公室占用率。

我不等待答复。

先生、亲爱的同事,请接受我最诚挚的问候。

附记:此信副本送下列各方知悉\\
——收讫Grothendieck先生1985年5月28日来信的全体数学教师;\\
——UER~5主任,他另外还收到Grothendieck先生1985年5月21日写给我的那封信的副本;\\
——USTL临时行政负责人,他已收到5月21日那封信的副本,我再随函附上5月28日那封信的副本。
\end{quote}

% ---------------- p.834 ----------------
\begin{quote}
朗格多克科学技术大学\\
数学研究所\\
\pmark\origpage{1099}数学研究所

1985年6月6日(星期四)18时会议纪要。

出席者:Auberson先生、Charles夫人、Ciulli先生、Contou-Carrère先生、Cuer先生、De~Lima先生、Delobel先生、De~Robert先生、Grothendieck先生、Hocquemiller先生、Escamilla先生、Hubert~Coulin小姐、Lefranc(LEFRANC)先生、Loupias先生、Meden夫人、Molino先生、Pierrot夫人、Pinchard先生、Saint~Pierre先生、Voisin小姐

经讨论后,到会者(19人)以16票赞成、3票弃权通过下列文本:

“数学教师们就Grothendieck先生的东西被搬动时所处的不可接受的情形,向Grothendieck先生致歉。他们共同承诺,警惕这类令人遗憾的事件不再发生。尤其必须明确:一间办公室的钥匙,未经占用者明确同意,任何人不得动用。”

\begin{center}
Lefranc先生\\
主任
\end{center}
\end{quote}

% ---------------- p.835 ----------------
\begin{quote}
\pmark\origpage{1100}朗格多克科学技术大学\\
数学研究所\\
一场“误会”的尾声\\
致数学楼旧日共事的同事们(教学人员与技术人员、第三阶段学生)书\\
Alexandre~Grothendieck

\hfill ……1985年6月7日

亲爱的同事:我在这里为我的办公室遭洗劫一事写下尾声——此事已在我5月28日的信里提起。那封信当时只寄给了数学教师,而它同样关乎、并以同一名义关乎在数学楼里占有一间办公室的每一个人,无论男女。我当初没有把信同时寄给技术人员和第三阶段的学生,是出于疏忽,也失之欠考虑——我(仓促地)以为,那样做等于给这桩插曲一个它并不应得的扩展。我在此向有关各位诚挚致歉——更何况,他们中有好几位(照理本该是“不知情者”……)曾向我表示同情,令我感动。大概也是因为这个疏忽,昨天专门讨论这桩插曲的UER全体大会,只限于“UER~5成员”参加。

在许许多多别的事情之外,这桩插曲还让我领教到:这类事在UER~5并非头一回发生——只是头一回,被瞄准的是一位“A级教师”。昨天表决通过的那份虔诚的决议,能否阻止这类事件重演——在普遍的冷漠之中(一如从前),尤其是对非在编教师、对第三阶段学生的冷漠——我不知道。我会留意向Mori夫人与Moure 夫人核实:她们是否确实收到UER主任的指示——今后无论以什么借口,都不得把任何一间办公室的钥匙交给任何人,也不得为任何人动用,除非得到其中一位占用者的明确许可。

我上一封信,结尾用的是这样几个字:“敬候您的(或你的)答复”。作为对这一期待的回应,我收到三份同情与声援的表示。它们来自Louis~Pinchard、Pierre~Molino与Christine~Voisin。同样,我还收到一份同一方向的表示,出自第三阶段学生Philippe~Delobel——他(与Christine~Voisin一样)曾随我读过DEA(深入研修文凭)。昨天有几位第三阶段学生出席全体大会,正是出于他的倡议。对他,也对我方才提到的所有那些人——他们(不含糊、不闪避)向我表示了声援——我乐于在此表达我的敬重与感谢。像这类的“严酷”经历,果实之一便是让人认出自己的朋友——倘若运气好,还有朋友可认……

我还收到另一封回应我的信,出自一位同事,他分明为发生的事情而喜不自胜,还借此机会温文尔雅地拿我开涮。这是我收到的这一方向的唯一回声。其余诸人那里,不少是这一位的全然冷漠、那一位的局促不安(其间不止一次,我触到一种没有说出口的恐惧:怕给人留下坏印象,从而断送晋升的机会,或危及本已朝不保夕的位置)。在那些人当中——凡是激动到肯为此劳神出席那场全体大会的人(大会是最后一刻偷偷摸摸才召集的,虽然一星期前就已排定……)——我感受最深的,是一种处心积虑、打定主意要把鱼搅浑的心计,口口声声“人人都和善,个个都可爱”。末了(在足足三刻钟的空谈之后),大家转而拿那个众望所归的“坏蛋”问罪,那位缺席者(多么凑巧),Lapscher先生——(据人们方才刚刚暗示出来的意思)\dmark{}\origpage{1101}发起这场偷袭的,正是他。至于走到点名道姓追究他那一步,那是谁也没有打算的,可怜的人——当然,对其他任何一个人,也同样如此,这不必说。

% ---------------- p.836 ----------------
在这桩洗劫插曲中以这种那种名分牵连在内的各“负责人”那里:一位Lapscher的无耻蛮横,一位Charles夫人“以取乐为务”的粗鄙(她在既成事实面前把这场偷袭遮掩过去,还添上她自酿的那一份放肆),以及USTL临时行政负责人Cano先生的失礼——我写信告知他事态原委,请他把问题提交大学理事会,他竟免作任何答复——这一切都令我震惊。然而最甚于一切的,是UER~5主任Lefranc先生那模棱两可的态度,令我错愕,令我痛心。从5月20日星期一(那一天,我把刚发现的情况和我对此事的感受告知了他)直到昨天,他始终不认为适宜之举,或者是向我说明究竟发生了什么,或者是毫不含糊地同一位Lapscher的强盗行径、一位Charles夫人的粗鄙划清界限。自始至终,他竭力维持那场倒霉“误会”的虚构,硬是给那些行径——在我看来无法容忍的行径——装出一副无关痛痒、甚至堪可敬重的模样。为的是不叫任何人难受,可以断言,他选择了“又护山羊又护白菜”:(多多地)照顾山羊,(稍稍地)照顾白菜。

在其他种种迹象之列,我也好好记下了一笔:我原以为可以列入朋友之数的许多人(其中有三位曾是我的学生)所守住的沉默;这一位招摇的漠然,那一位的窘迫不安,还有一位甜腻腻的欢欣雀跃。还有一位Micali的沉默(他是这场偷袭的共同受益人;而几年前,他就已有充分的机会确信:招得Charles先生与夫人不快,会带来多少麻烦……),以及Brun小姐的殷勤卖力——她领一位Lapscher的吩咐,客串了一回雇佣锁匠兼搬家工(等到这场行动的性质再无可疑之时,连一句遗憾的话也没有)。

在这一切的底子上——昨天我又见到那间十二年来一直是我的办公室,这一回它成了战场:我的东西(连家具一并)被七手八脚重新堆作一团(距那场“闪电”偷袭,足足过去了十五天……)——如今我已没有心绪再在里面重新安顿下来。人们向我担保,同样的插曲在我身上重演的可能性微乎其微;何况我尽可以先发制人,把那把迄今托付给Mori夫人与Moure 夫人的第二把钥匙收归己有。但只要物质上办得到——尤其在我借调法国国家科学研究中心的整个期间——我宁愿从今往后放弃在USTL使用办公室,不加争斗地把地方让给Lapscher们、Charles们以及各色同好。

如果能够避免,我不会再在USTL执教。在那儿,我自始至终——可以肯定——不过是个异乡人,一个故乡在别处的人:无论我对待数学的路数,还是我对待教学的路数,抑或我的生活方式,皆是如此。大学这个小天地所能教给我的,我想我已经学到了;而最后一“卷”,便是这桩插曲教给我的一切——它刚刚在一片皆大欢喜中落幕。我刚刚参与的这场UER~5会议,很可能是我的最后一场;这封信,也很可能是我有机会写给你们(或写给你)的最后一封。而这一回,我不等待答复。

\hfill Alexandre~Grothendieck
\end{quote}

% ================================================================
%  18.5.5 Le seuil + 18.5.6 L'album de famille (原书 p.836–864)
%  源文: _work/chunks/chunk08.md, chunk091.md, chunk092.md,
%        chunk093a.md, chunk093b.md, chunk094a1.md, chunk094a2.md, chunk094b.md
% ================================================================
\subsection{门槛}   % 18.5.5. Le seuil

\notehead{172}{}

\pmark\origpage{1102}(3月22日)我原以为顶多一两天、撑死十来页,就能把这些自从10月起我就打算重新巡视的
出名的“四种运算”走完一遍。可如今一鼓作气干了三个多星期,页数排到了一百上
下——居然还没有全然写完!初稿从2月26日写到3月1日,就花了四天。它给我的,不过是一块底布,
好(终究)在上头绣出一桩“事情”来,而不只是一纸侦查结论。3月1日的第二天重读初稿,那印象可
不妙:一张没完没了的“诉苦状”,而且,除了三四位真正在行的专家(还得假定他们有耐心读完
……),只怕谁也读不懂。我明白:至少得约略交代一下所为何事,也就是说,至少得铺一个语境——
否则就白搭\fnote{835}{(*)}{整部《收获与播种》的思考里,除这里之外,我曾对“即兴”写作方式作过这种(规模确实小些)让步的,只有“那张条子——或新伦理”一节(n\textsuperscript{o}~33)与“不义——或归来的意义”那条注(n\textsuperscript{o}~75)。}。

这就不可避免地带来一些重复,同\emph{葬礼}第一部分相比——可有些情形,重复不但有益,甚至
必不可少(在数学里,一如在别处)。而且在这些情形里,人们很快发现:那些所谓的“重复”其实并
非重复,因为被“重说”的东西,实际上是\textbf{重看了一遍、以新的眼光再看},并且照明已经不同。在
为“四种运算”铺设“语境”、交代我著作的某些侧面时,我觉得自己对这份著作学到了些东西,把它
摆得更正了。这么做,我也许并没有就我自己或他人学到什么真正新的东西,可我不后悔费了好几
天的功夫,把那份“诉苦状”式的初稿这样重写。这份著作,我放入过我所能给出的最好的东西;它
值得我怀着成熟给我的那份回顾,以不同的光线重新去认识它。正当我准备原原本本记录这份著作
自我离去后(交在了可靠的人手里,对此我从不怀疑……)遭受的一切之时,好事在于:我在它身上
\textbf{停一停},在它的位置上、在构成它之美的那个统一体上,哪怕只停几页——以此再一次,向那
些我亲眼看见被践踏的东西,表达我的敬意。

\pmark\origpage{1103}可这还不算完,差得远呢!放弃了“诉苦状”的写法——那种夹着编号、指向\emph{葬礼}第一部分
那些血肉更丰满的注的写法——我明白了:这些我重新拿起的笔记,同《收获与播种》里所有其他
章节与笔记一样,必须自身可读,把它们要说的实质,完整地传达出来,独立于那些指向思考\textbf{另一
时刻}之注的引用。这又带来许多“并非重复的重复”——也就是,以新的照明,重看将近一年前、在发
现的新鲜震动里逐日记下的东西。何况当时,那么多意外的、有时难以置信的事实劈头盖脸,压根谈
不上一场稍具方法的“侦查”。那时,我不过竭力去接住劈头砸下来的东西,勉力把它们“安放”进去,
细节顾不上。我的精力,大部分被\textbf{招架}这些我发现的原形毕露之物的那种\textbf{癫狂}、那种难以置
信(就像那篇皇帝新衣的故事里一样……\fnote{836}{(*)}{见同名之注,n\textsuperscript{o}~77$'$。})所吞掉;尤其是,要承受那股突然扑面而来的暴戾、犬儒与轻蔑之气——“顶着那副光鲜的面孔……”我太熟悉的面孔;那是另一个时代的气息,我经历过它,我没有忘……

而这最近三个星期,反倒成了补完去年那场风暴式侦查的机会:把某些文本再往细里翻一翻(SGA~5,
尤其是那卷所谓的“SGA~4½”)。于是生出一串(有一阵眼看不象话地收不了尾的!)多少原原本本
的脚注;其中一些成了子注,而其中一个(预定的名字叫“那个公式”)占用了我整整四天,又裂成
另外四个\fnote{837}{(*)}{(6月1日)如今已经变成六个了……}……有几次,我觉得永远写不完了——可到头来,它还是收敛了
\fnote{838}{(*)}{(6月1日)当然,只是一场临时的“收敛”,因为“封神”那条注,最后炸成了三十来条各自独立的注、子注等等,单是它们就有足足一百五十页!}。眼下撂下的,还有十来页涂改得太厉害、须得重写的页,以及“分赃”与“封
神”两条注的脚注,留待日后补上。眼下这样就可以了!哪怕日后还要回头管\dmark{}\origpage{1104}“后勤”,我急于收场,
急于不再耽搁,把“四种运算”这一章里我看到的还有实质可说的东西说出来。

我在\emph{葬礼}里区分出两个“面向”或“层次”:它们紧紧相系,却毕竟不同。它们之间(至少在
我眼中)相当分明地隔着一道\textbf{门槛}。

一方面,是“时尚之风”的面向(有时直到那股“嘲弄之气”——我在《收获与播种》里已不止一次
有机会说到它)。它主要表现在我在别处\fnote{839}{(*)}{见注“掘墓人——或教团全体”,n\textsuperscript{o}~97。}称作“自动排斥的态度”的东西上:这种态度常常拦腰截断数
学常识的简单反射,冲着某些人及其数学贡献而去——这里说的是我,以及另外几个人;他们被归
类(尽管当事人有时竭力想同我划清界限)为同我“串通一气”。在我的情形,要“排斥”(或“埋葬”)
我带来的\textbf{一切},当然不可能:好大一部分,早在我1970年告别数学舞台之前,就已进入日常通用
的公共领域\fnote{840}{(*)}{不过,确实,连某些早已进入“日常”通用(至少在我学生与近旁合作者的小圈子里)的观念与技术,自我离去后也被埋葬了。不妨说,尤其如此的是$\ell$-进上同调那件工具:我在 SGA~5 里(以 SGA~4 的关键结果为起点)极详尽地发展过它。它被我那班上同调学生压在斗底下,Deligne首当其冲,直到1977年,才以我所知的那个形状与精神被掘出土来。}。不过,
确实(而且我在一年前的注“我的孤儿们”(注 n\textsuperscript{o}~46)里第一次确认了这一点):我的著作——
写下的与未写下的——在\textbf{上同调主题}上,绝大部分自我离去的次日起就被埋葬了,头一份经手
的正是我的学生。(我引入的某些主题,四年后、七年后或十二年后被掘出土来,不提我这个人——
但那我们已经触到“第二层次”了……)

这样的自动排斥,人们尽可以惋惜:它有时违背对他人的最起码的体贴与敬意,而且在任何情形下
都与数学的常识与鉴别力无关。而当它打击的是那些手段有时十分出色的年轻数学家时,就更值得
惋惜:“轻蔑之齿”咬灭了一种欢乐,糟蹋了一份原本美好的激情,让它化为苦涩——投入眼看着被
糟蹋了(按那些说一不二的共识……)。还值得惋惜的是,\dmark{}\origpage{1105}当这种排斥打击的是那些简单而丰
饶的观念——它们早已充分证明自己,从无中生有地催生出如今“人人”看也不看就使用的强大工
具。在前一种情形(一个被毁掉的志向),伤害多半不可逆转;在后一种情形则不然——
因为迟早,那些简单而根本的观念,那些“在路上”的观念,终会显形或重新显形,终会进入公共遗
产。不管怎样,谁也不能合情合理地强求任何人,对一个人、一部著作、一个观念\textbf{往好里}想——
只要他(出于只有他自己知道的理由)一心想\textbf{往坏里}想,或干脆忘了它。这类问题,当然属于“个
人伦理”,而且属于其中精细而要紧的部分;但依我看,把它弄成一桩集体的“科学伦理”问题,恐怕
不妥;倘若硬要试,只怕药比病还凶……

我方才影射的第二个“面向”或“层次”,恰恰是这种集体伦理被破坏的层次。我说的那道\textbf{门槛},
是一条\textbf{共识}——据我所知,自有科学见诸文字记载以来,它在一切科学里都得到普遍接受。这条
共识规定:谁也不得把他从\dmark{}\origpage{1106}别人那里取来的观念\fnote{841}{(*)}{我在这里说“观念”,当然,在数学里绝不单指“结果”。常常,一个提得好的\textbf{问题},它触及一个此前无人看出的关键点,就比一个“结果”更重要,哪怕更难。即便这个问题还没有凝结成一个精确的\textbf{命题}——它本可以构成一个假设回答的胚胎,或者一个(尚属猜测的)多少完整的回答——情形也是如此。当然,从一个起初含混的问题里提炼出这样一个命题,是数学工作中本质而具创造性的一部分。把一个(也许深刻的)问题加工后的版本端出来,却隐去它的出处(即便加工是这位魔术师自己做的),同隐去一个形状完整的深刻命题的出处——借口是自己给出了证明——一样,都是剽窃,正如把别人给的证明当成自己的拿出来一样。
对丰饶\textbf{观念}的引入,道理相同——它们往往比好的命题更关键:因为“好命题”的问题,只有在已经提炼出好观念之后才谈得上。这里的道理相同:借口修改过、甚至改进过取自他人的观念,就隐去其出处,其不老实,同原封不动地“借用”这观念,毫无二致。最常见的情形是:提出一个问题(哪怕还含混),拿出一个命题或一个观念(哪怕不完善、临时),才是关键的一步,而不是人们后来加上的那些改进(更精确、更广、更深)。可即便情形不是这样,对那个在改进所受之物上做出原创工作的人,这仍旧不能当作隐匿其所受之物的“理由”(或者,同一回事:把它“拆掉”……)。
如我在别处已强调过的(在注“肌肉与肚肠(阳葬阴(1))”(n\textsuperscript{o}~106)的子注 n\textsuperscript{o}~106$^{1}$里),一条猜测性命题的“价值”,既不取决于其设想的难度,也不取决于它多“可信”,也不取决于它最终被证为真或为假。无论如何,人们愿意给一个数学观念(它表达于一个问题、一个命题、一个观念或一个证明之中)或一组观念的“价值”,很大程度上是主观的,难以成为科学伦理之共识的对象。所以,一个诚实的科学家,会标明他所使用的(显式地或默会地)、不属于“众所周知”的一切观念的出处,而不肯顺从那条斜坡:对他已在内心决断(也许是为了某桩可疑的案子……)反正“显而易见”、“不足道”、“不要紧”(或同类评语)的某个观念,隐去其出处。}

\pmark 这条共识责令我们:凡我们所陈述、使用或发展的观念,每次——至少当这些观念并非出自
我们自己的手笔、也还不是公共遗产、尚未被(不是三四个入室弟子,而是)“\textbf{所有人}”知晓的
时候——都要标明来历。

我不记得曾听人质疑过这条共识。在我身属数学圈子的岁月里——从1948年(二十岁的年轻新手,
去高等师范学院听Cartan的课)到1970年(我告别数学舞台)——我极少有机会目睹、甚至极少听说
有人公然违背这条共识或原则;只在一位同事兼朋友那里,他在这一章上稍欠检点\fnote{842}{(*)}{这位同事的情形,顺带谈过,在《收获与播种》第一部分,就在紧接着下一句所引的那一节里。经过一年多的回顾,这个“案子”如今所显出的分量,是我从前没有给它的。}。正如我在《收获与播种》第一部分里强调过的
(在“一个严守的波利钦内尔秘密”一节,n\textsuperscript{o}~21),遵守这条原则,对一个具备起码的诚实与自尊
的人,绝非自然而然之事。相反,它要求极大的警觉:童年便根深蒂固的反射,自然而然地推着我们
高估自己的功劳,把我们在自己之内消化他人观念的工作,同这些观念本身的构思混为一谈——两
者实属完全不同的层次。写上面所引那一节——一年多了——的时候,我显然对自己也没弄清:该给
这条共识多大的分量。那时我心里有一片含混(在思考的这个阶段,我自己也没有清楚意识到\dmark{}\origpage{1107}),连着一种弥漫的感受:对\textbf{他人}(比方说对我自己的学生)提出严格要求——在他们
同我的关系上遵守这条原则——是小气,是一种配不上我的卑微。所以那时,我之内有一道\textbf{暧昧};
我只是在6月1日那条同名之注(n\textsuperscript{o}~63$''$)的思考里,才把它看清。这场思考把这道暧昧驱散净尽;
而我当时意识到,这道暧昧,从最初(六十年代初)直到去年,一直沉沉地压着我同学生的关系。我
懂得了:数学家这一行(或更一般地,科学家这一行)的严谨,首先意味着对自己怀有极大的警觉,
遵守这条众共识中之首要者;同时也意味着对他人同样严格的要求,尤其对那些我们负有领入门墙
之责的人。

年复一年,我越来越懂得:这一行\textbf{绝不只是}一门技术上的窍门,甚至不只是发挥想象、解
决出了名的难题的能力。从某种意义上,这一点我从来就知道——可我低估了“伦理”的一面,或者
说\textbf{集体}的一面\fnote{843}{(*)}{我在这里无意说:一个处境的“伦理”面向,总是同时就是“集体”的面向,牵涉一个人同群体(这里即“同事”或“同侪”之群)的关系。不过,在我正在检视的这条“共识”上,情形确实如此。按照自童年起塑造了我的世界观的那些特殊条件作用,直到去年,我总倾向于低估(乃至无视)集体之物,而偏重个人之物。我的“数学历险”中“集体历险”的那一面,去年才清楚地向我显现,先是在“Galois的遗产”一节(n\textsuperscript{o}~7),但尤其《收获与播种》第一部分末尾的几节:“孤独的历险”与“一段过去的重量”(n\textsuperscript{o}s~47、50)。},以为它在善意的人们、体面的人们之间,“理所当然”。就这样,我为我刚
才说的那道“暧昧”备好了土壤;那道暧昧,也是(顶着一副虚假的“慷慨”的门面)对我的学生与
类似者的\textbf{纵容},而更隐蔽的,是对我自己的一种纵容。

我离开了这个“善意的人们、体面的人们”的圈子——它也曾是\textbf{我的}世界,我曾乐于自认同属
的世界。在去年4月19日之后的几个星期里,我稍微仔细地看了它一眼:离开不到十五年,我
在里面看到一种\textbf{腐败},做梦我也想象不出它的样子。

\pmark\origpage{1108}对我而言,这是一个谜:身为这个世界的一员,“搞数学”还能有什么\textbf{意义}——如果只
是把它当作一种\textbf{权力}的手段,或者(对身份低微的)当作一种谋一碗饭的手段,物质条件嘛,
说实在的还挺舒服(只要运气好,早早就“安顿”下来了……)。

\subsection{(5) 家庭相册}   % 18.5.6. (5) L'album de famille
\notehead{173}{}

\fnote{844}{(*)}{眼下这条“家庭相册”之注,最初本是前一道注“门槛”(同日,3月22日所写)的直接续篇。这一部分,如今构成其 a.部分(“一位备受簇拥的死者”);6月10日、11日又添了 b.(“新面孔——或吊嗓”)与 c.(“众人之上的那一位——或首肯”)两部分。下一道注“升级(2)”(n\textsuperscript{o}~174)——也是3月22日——紧接本注 a.部分(同日所写)而来。至于 b.、c. 两部分的脚注,日期是6月13日与14日。最后,6月18日又添了 d.(“最后一刻——或一个禁忌的终结”)这一部分。}

\subsubsection*{a.\ 一位备受簇拥的死者}
\tocsubsub{a.\ 一位备受簇拥的死者}

\pmark (3月22日)说得粗鲁些:\emph{葬礼}里有“时尚”的层次,也有“诈骗”的层次。也许不过是
我落伍了;我那个年代被看作诈骗的行径,如今只要干它的人属于上流社会,就早已成了完全正当、
体面的行当。也许那道“门槛”,早就荡然无存了呢?“第二层次”是\textbf{一场唯一的、规模浩大的诈骗
行动}:目标是我在上同调主题上的全部著作;继它之后,又是Zoghman~Mebkhout——那位鲁莽的继承
人、被埋葬的宗师的身后学生、固执而无名之辈——的著作。这场行动的总指挥,是另一位学
生:绝非身后的,倒是隐秘的,没错;他玩着一个心照不宣的角色——我著作的“继承人”——同时对
著作、对工匠,一并拒认、一并拆毁。他是我的朋友\textbf{Pierre~Deligne}。他的干练副官,不是别
人,正是随他一道选了“上同调”这条线的四位学生:\textbf{J.L.~Verdier、L.~Illusie、P.~Berthelot、
J.P.~Jouanolou}。死者想必被簇拥得妥妥帖帖:既有同享\emph{葬礼}哀荣的\textbf{共殓诸友}
\fnote{845}{(*)}{说实话,据我所知,“共殓者”不是一位,而是\textbf{四位}——即四道棺材注(棺材一至四,n\textsuperscript{o}s~93--96)所写的那四位。},又有生前曾是他的亲近者的那些人。至于临时凑数的殡仪工——由总指挥导演
的那场双重埋葬里,前来帮衬的——我看见另外七位“世界闻名”的数学家(借某张广告招贴的\dmark{}\origpage{1109}措辞
\fnote{846}{(*)}{指 IHÉS 1983年为建所二十五周年出的纪念小册子。关于它,见“悼词(一)(二)”两条注(n\textsuperscript{o}s~104、105),尤其第454页。}),在葬礼仪式进行之中不时露面;仪式已在家庭相册(又名“四种运算”)里
巡视过一遍。他们是(按其在仪式中的分量)\textbf{Teissier、Beilinson、Bernstein、Milne、Ogus、K.Y.~Shih、Katz}
\textbf{(N.~Katz)}。

\pmark 到此,我巡视了据我所知曾以这样那样名分\textbf{积极参与}“葬礼行动”的数学家。共十二位
\fnote{847}{(*)}{与(《收获与播种》第一部分)“耶稣与十二使徒”一节里相同的“十二”——那一节巡视了所有跟我做到国家博士论文层级的学生。确实,在我葬礼的积极参与者之中——这一回,是Springer殡仪股份有限公司那一边,而非信众教团那一边——还有 K.F.~Springer博士(那家可敬公司的共同老板)与Peters、Byrne两位博士;后文(注 n\textsuperscript{o}~175)还会谈到他们。于是是十五位!}。至于后四位,就我所知的事实,我无从断定他们存心不良。但我认为,他们的责任并不因此就轻了分毫。因为,他们乐于对他们所做的事视而不见,这本身是一种选择;选择,豁免不了他们对自己行为的责任。

至于那些直接串通的参与者,要我开一张哪怕不完全的名单,或估一估他们的人数——那人数肯定
是另一个量级——我实在无能为力。只需提醒一句:这里面有1981年6月吕米尼那场“难忘讨论会”
(人称反常层会议)的全体参加者;还有那卷顶着“SGA~4½”之名的文本的读者们之中,多少知道
SGA 这个符号意味着什么、却“听之任之”的所有人。

\pmark 我看见两份文字,它们见证着七十年代、八十年代数学里的一场\textbf{凋败},在我们这门科学
的历史上,恐怕还从来没有堪与比拟的先例。其中一份,凋败在它给自己起的名字里就绽开了——那
名字本身就是一次(天才的……)冒名:那份数字标着“SGA~4½”(作为常用的参考符号)、又称
“平展上同调”的文本——P.~Deligne署名,“协作”(除 L.~Illusie与 J.L.~Verdier等人之外,还)有
A.\textbf{Grothendieck}\fnote{848}{(*)}{关于这份“协作”的含义——它是Deligne所设骗局的一部分——见注“对调”(n\textsuperscript{o}~68$'$)。}。第二份文字,是1981年6月吕米尼\dmark{}\origpage{1110}讨论会的文集;尤其并且首要的,是其第一卷:由讨
论会引论(Teissier与Verdier署名)与讨论会的主报告(Beilinson、Bernstein、Deligne署名)构成。

倘若直接参与过这场凋败的、以这样那样名分参与的人之中,有一位能在自己身上找到那份朴素与
勇气,公开道个歉——或者更好,公开就\textbf{他自己}经历的那部分原委说个明白——这对所有
人都是一桩善事,也是那一代容忍了这等凋败的数学家的光荣。可这多半是奢望。

同样是奢望的,想必还有:让 J.L.~Verdier不再在高等师范学院占着Henri~Cartan曾占的那个位置。
那无疑是法国培养数学“接班人”的关键岗位。很久以前,当我得知Verdier——他曾是我的学生,
我曾钟爱于他——升到了这个职位,我自己觉得脸上有光(同时暗地里也有几分得意)。那时,我丝毫
没有怀疑过:Verdier会完美地承担曾属于Cartan的角色,让那些最有志于数学的年轻人在与他的接触中学好这门手艺。可今天我看见(多年前已在看,却从来没有以今天这般粗暴的清晰看见)我
错了;我在这里明白地说出来,不是要给他或任何人抹黑。我认为,他已丧失了指导研究的资格。这样
说,我并不推卸我自己的一份责任:是我把他教坏了(他同我所有其他学生一样)——那门我热爱并
且依然热爱的手艺。

\subsubsection*{b.\ 新面孔——或吊嗓}
\tocsubsub{b.\ 新面孔——或吊嗓}

(6月10日)“家庭相册”这条注动笔以来,两个半月过去了。我当然没料到,还须回到它上面来——
因为对\emph{葬礼}的调查又有了新的反复。主要是那场微薄的“封神”——我当时刚写了五六页——
炸开成一场大写的、足足一百五十页密密实实的\textbf{封神},让我随之发现了一些“新面孔”:它们理应
在家庭相册里占一席之地。也有一些早已熟识的面孔,如今显出他们同样属于那支大军——那些
在“诈骗”的层次上积极参与“葬礼运算”的人。我在这里“备而勿忘”地把他们再巡视一遍\dmark{}\origpage{1111},也好让
每一位当事人确知自己与名流为伴(不过,这事想必早就办妥了……)。新到的照片,按它们
引起我注意的先后,插入相册。

先说Springer出版公司那一边:\textbf{K.F.~Springer}(公司共同老板之一)、\textbf{Peters}、与
\textbf{Byrne}夫人(C.M.~Byrne)。细节见下文“殡仪馆——‘为科学服务’”那条注(n\textsuperscript{o}~175)。写本注开头、
3月22日的那几天,我刚收到 K.F.~Springer的信(日期3月15日)没几天;那封信,把我对那家可敬
的殡仪公司——忠于“为科学服务”的箴言——的精神的最后一点疑虑,一扫而空。

封神那一边(经由对当值无名之辈的埋葬),我知晓了\textbf{Kashiwara、Hotta、Brylinski、Malgrange(B.~Malgrange)、Laumon、Remmert}诸位的贡献;还不算一位\textbf{匿名审稿人}——其存心不良,
毋庸置疑;不过说真的,倘若要把与\emph{葬礼}沾亲带故的那些可疑文章与书籍的审稿人(reviewer) 也一一
点验,那准得另备一册相册。此外,我的老朋友Katz重新露面了,这一回的语境,让我对他本来(至少)
存着的那份善意推定,化为乌有。这样一来,据我所知曾以这样那样名分积极参与那场所谓“反常讨
论会”的骗局行动的、国际知名数学家,人数到了十四位(算上那位出名的匿名审稿人,十五位)。详
情备有文件,见“封神”,尤其“……以及横财”、“光荣之日”、“黑手党”、“掠夺的空白授权——或‘高
尚行当’”几条注(n\textsuperscript{o}s~171(iii)(iv)、171$_{2}$、171$_{4}$)。

最后,在“Motif 运算”那一边,出现了(迟做总比不做好)我的又一位学生,稍稍离大队人马有点距
离。事后,我几乎是被逼着把他算作我的“上同调”学生的第六位,尽管“我那个年代”,他对上同调是
什么,连半点概念都没有。他是Neantro~Saavedra~Rivano;显然,他被人当了“棋子”(当然,出
于他自己的自愿),而非为自己行动。他与Monsieur~Verdoux(化装成“献殷勤的骑士”)\dmark{}\origpage{1112}周旋的历险,已
在“第六根钉子(钉进棺材)”那组注(n\textsuperscript{o}s~176$_{1}$至176$_{7}$)里逐页复原——那是4月19、20日写的
(最后一道除外,尚待写出)。于是,在我“从前的”学生之中,积极参与宗师\emph{葬礼}
的,人数也到了六(十二之中)。Saavedra在\emph{葬礼}中的份额,其特别之处在于:他参与的那场
“塔纳卡范畴”(原文如此)运算,是第一场大规模行动,目标是变走我著作的一个重要侧面、以及我
发展的那门哲学(就其情形而言,是在 Motif 理论的带动之下、乘其便而发展的)的父亲身份。

把新到的面孔算进相册,并撇开Springer-殡仪馆一方的贡献、只算信众教团一方的,那么据我所
知曾积极参与\emph{葬礼}——按我那个年代人们管它叫“诈骗行动”的那种水平——的知名数学家,
人数到了十九位\fnote{849}{(*)}{算上那位出名的匿名审稿人,二十位。}。这些参与者中,只有三位——即与 P.~Deligne合署“难忘文集”《数学讲义》
第900种的那三位——在我看来,存心不良尚未坐实。

这份名单,当然远没有穷尽我的同事、旧学生或朋友——他们以这样那样名分、或主动或被动地参
与了我的葬仪,却没有走到参与一场铁板钉钉的诈骗的地步。我点出过三十来位;其中大多数,已在
对\emph{葬礼}的思考中提到;算上前面那些,凑足了结结实实的五十来位——而这还只是:最近八九
年,我在遥远的隐居处,身不由己地得知的那些;或是一场刻意保持得极有限的调查之中,自己撞上
我注意的那些。

这些数字本身已经雄辩,并且出人意料地印证了去年就已强加于我的那个印象:对我的著作、对我
这个人的\emph{葬礼},不是一个人的事业,也不是一个严格有限的集团(比如我离去前的学生,或我
那些“上同调学生”)的事业,而确实是一项集体事业——其水平,是“教团全体”的水平;或者至少,
是数学建制中曾经\dmark{}\origpage{1113}\textbf{见证}并参与了我1955至1970年间作为几何学家的著作之兴起与繁盛的那一
部分的水平。我1970年的离去,至少在这一片数学里,成了一个信号:对“Grothendieck式”
数学,一场即刻而严厉的\textbf{排斥反应}——人们把它感受为“女性气质的数学”的象征与化身
\fnote{850}{(*)}{关于对某种对待数学的进路风格的这类排斥反应,见“肌肉与肚肠(阳葬阴(一))”、“阴的葬礼(阳葬阴(四))”、“天赐的机缘——或封神”、“拒认(一)——或召回”、“拒认(二)——或变形”诸注(n\textsuperscript{o}s~106、124、151、152、153)。在“上涨的海……”、“九个月与五分钟”、“箭与浪”、“兄弟与配偶——或双重签名”、“阴,作为仆人——与新主人”、“阴,这仆人——或慷慨”诸注(n\textsuperscript{o}s~122、123、130、134、135、136)里,我试着勾勒“女性气质的数学”的某些显著特征,与互补的“男性”特征并列。}:在那
里,视野永远先于并启发着技术的层面;困难永远被化解,而不是被砍掉;与事物在表面参差之下的
深层统一的持续接触,让人每时每刻都能从偶然与次要的无定形团块里,认出什么是要紧的。与此同时,我的离去也是一个信号:一切概念性的工作骤然停摆——或者更确切地说,一切此类工作被
\textbf{褫夺公权},以“深化”为名,骤然打上嘲弄的烙印。

就这样,砍掉数学创造工作的一个本质“侧面”——“阴”的或“阴柔”的侧面——人们达到的是一种
惊人的“扁平化”、“削平”、“干涸”\fnote{851}{(*)}{关于这种“削平”的初步确认,见注“无用的细节”部分(c)“什么都不像的东西——或干涸”(n\textsuperscript{o}~171(v))。}。这件事(在我看来)是通过一个骤然而严厉的急转弯完成的,几乎一
夜之间。这件事如此奇怪、如此闻所未闻,以至于显得难以置信。我花了一年多对\emph{葬礼}的高强
度思考,才终于把握住发生了什么,才认了这份明摆着的事实。我不知道,近年或近几十年,或任何
别的时代,在科学的一支里、或任何别的动用我们创造之力的人类活动里,可曾有过可相比拟的转
向。

可我回到我的相册上来。除了方才点到的诸位,把另一些人对\emph{葬礼}的参与——在我看来毫无疑
义的——的名字收进来,我想是有用的。况且,我并不确信他们中哪一位对我怀有恶意;他们之中肯
定不止一位,\dmark{}\origpage{1114}甚至对我怀着同情乃至好感(回应着我自己心里相似的情感)。他们中也许
没有一个人,会不真心错愕地听说:竟然有一场“葬礼”在对我这个人、对我的著作举行;更不用说
听说自己据说还以这样那样的方式参与其中。在这里指名道姓,本身就会产生那个(对我而言受欢迎
的)效果:把这事告知他们;并且(如果他们自己关心)由此给我们之间一个说明的机会。当然,关于
我所感知到的——无论对错——他们对我的葬礼的直接参与、或经由“共殓者”的参与,我随时听候
诸位垂询,提供一切细节。我绝无意质疑他们中任何一位的善意与职业操守
\fnote{852}{(*)}{(6月16日)随着新近传到我这里的信息,这份善意推定,在Borel的情形里化为乌有了。照他去年与Zoghman~Mebkhout的通信——事由是Borel在苏黎世主持的一个 D-模理论讨论班——我早已知道:Mebkhout曾告知他,自己是那门理论(所谓“Riemann--Hilbert”)中心的范畴等价的作者,并给出了确切的参考文献,还寄去了他的全部著作,Borel本可轻易从中确证事实。这并没有妨碍Borel以照例的居高临下(甚至失礼)对待他。在刚于奥伯沃尔法赫结束的一场同题讨论会上(Algebraic~theory~of~Systems~of~partial~differential~equations,~Oberwolfach,~1985年6月9--15日),Borel做了头三个引言报告(题为“代数 D-模理论”),为“好上帝定理”铺路;\textbf{Mebkhout的名字,在这些报告里一次也没有被提起}——在随后的报告里也没有(Brylinski的报告里,除了一句顺带的“喊停式引用”之外)。照Mebkhout刚寄给我的会议纪要,这场讨论会——Borel在其中指挥乐队(顶替没有出席盛会的Deligne)——是四年前那场反常层会议的名副其实的\textbf{重演}。“黑手党”几乎全员到场:Verdier、Brylinski、Laumon、Malgrange,甚至(这一回)Kashiwara(此人在苏黎世那场讨论班里,早已扮演头等角色,尽管Mebkhout已就其为人向Borel作过周详的通报)。不用说(同苏黎世一样),没人觉得有必要请Mebkhout做个报告;而且(除了这位Mebkhout偶尔的插话——它们都掉进了冰窟)先祖的名字一次也没被提起(除了那个不合时宜的“Grothendieck群”里到底有他的在场)。双对偶定理在那里依旧叫“Verdier对偶”,连Borel的报告里也是。Mebkhout早在去年就再三提请注意:这份双对偶,是从 SGA~5 的第一讲抄来的——看来Borel对某种风格、对某位缺席者,发展出了一种过敏症,禁止他考虑这一类参考文献……他自己在《交叉同调》(Intersection~Homology,Birkhäuser,1984)里,也成了同一骗局的参与者——而那本书,是在Mebkhout向他点破Verdier的把戏\textbf{之后}出版的。
对Borel,我曾把善意推定保持到极限:五十年代我们同在Bourbaki,一起共事,我对他知之甚深。在真正可称为我在数学世界的“出身圈子”的成员里,他是第一位,我不得不在今日、以无可怀疑的方式,确认其直接参与——而且是在“诈骗”的层次上——葬礼。},其中不止一位,我甚至可以补一句:
他们的全然的善意与诚实,于我是不容置疑的。\dmark{}\origpage{1115}与其愚蠢地按字母序开一张清单(那种
事计算机比我强),我宁可按大致的时间顺序,列出为我的葬仪齐声合唱的信众们的名字:不是按他
们在葬礼仪式上露面的时刻(那多半我并不知道),而是按我清楚得知其参与的时刻。另一方面,我
要把我全体学生单独列出\fnote{853}{(*)}{我在这里说“我的学生”,指的是跟我做到博士论文层级、并且(除Deligne外)跟我做了博士论文的那些人。共十四位(其中两位在我“离去之后”),见“耶稣与十二使徒”一注(n\textsuperscript{o}~19)的巡视。}。除了Hoang~Xuan~Sinh夫人——她在越南工作,分明显然太远,
帮不上我的葬礼——我的学生之中,没有一位,不曾以这样那样的方式参与其中。此事我已在“沉默”
那条注(n\textsuperscript{o}~84)与“棺材1——或知恩的D-模”(n\textsuperscript{o}~93)的开头说明过;这里不是重谈的地方。对
于我的每一个学生,我认为最值得做的,是一份就发生了什么而作的深入的说明。

“为我的葬仪合唱”,分布在极不相同的调门上。我辨出了四种主要的调门;它们为我说成了一场第一
流的、气派十足的多声部葬礼!有对一切带着Grothendieck气味的数学之发展的“沉默而有效”的\textbf{抵制};
有\textbf{失礼}与不体贴——我\dmark{}\origpage{1116}离去之前,在数学世界里我从未遇到过的;在一两个极端的情形,它们采取
了几乎不加掩饰的\textbf{嘲弄}的形状。有\textbf{刻意无视}或贬低我的思想与观点在其个人著作中的影响、
或在当代数学某一部分中的影响——而在那些情形,这影响分明而关键;或把无可置疑地属于我的
结果与观念,安到一个第三者头上。最后,还有那种(“\textbf{鸵鸟}”)姿态:碰巧面对一场刺眼的骗局
的人,把头埋进沙子,装作什么也没看见、什么也没感觉到。

不用说,信众的唱诗班里,不止一位在同时吊着好几个调门的嗓子。

话说到这里,终于可以兑现那份答应过的清单\fnote{854}{(*)}{我没有把那八位“非上同调”学生的名字收进这份名单;连同上面已经巡视过的上同调学生的名字,它们可在已引之注(n\textsuperscript{o}~19)里找到。
同样,在我的“家庭相册”里,也理应列出我的同事或旧友之中那些我确知为“不予埋葬者”的名字——他们以绝无含糊的同情与敬重之表示证明了这一点。首先,关于我1983年所做的那份工作《追寻栈》(我还打算回头谈这份工作),我收到过热忱的鼓励,来自 J.~Benabou、N.J.~Baues、A.~Jouanolou,尤其来自Ronnie~Brown与Tim~Porter——他们(不止以一两种方式)在我做那份工作的整个期间给了我切实的帮助。
确实,这几位同事所属的环境,与我惯常自认同属的那个环境相当不同;而我那场气派十足的\emph{葬礼},也正是自然而然地发生在后一个环境里。作为属于或接近这一环境、并且最近(在过去一两年间)给了我同一方向之表示的数学家,我乐于在这里点出 B.~Lawvere、J.~Murre、D.~Mumford、I.M.~Gelfand,以及(并非最末!) J.P.~Serre。这最后一位,独享同时出现在两份“名单”上的殊荣——“埋葬者”的名单,与忠实朋友的名单!}了——它来充实我们的家庭相册:Eckmann、Dold、Campo(N.~A.~Campo)、Mazur、Poenaru、Epstein、Cartier、Quillen、Kuiper、MacPherson(R.~D.~Mac~Pherson)、Heisuke~Hironaka、Hirzebruch、Tits、陈省身(S.~S.~Chern)、Artin、Langlands、Rota、Goulaouic、Fulton、Borel、Tate、Serre。

\subsubsection*{c.\ 众人之上的那一位——或首肯}
\tocsubsub{c.\ 众人之上的那一位——或首肯}

(6月11日)昨夜敲出那份名单时,我觉得自己有点傻:那里老老实实排开的每一个名字,单是它本
身,就唤起一整团丰富的联想,而在这里全无踪影。可这里绝不打算一一停在那些名字上、停在它们
所唤起的东西上——那需要再写一卷,而我正急着结束\dmark{}\origpage{1117}这一卷!我向当事的人们致歉:如此草率地把
他们“贴”进一幅不太能启人深思的“(出席我葬礼的)合影”里。确实,他们中的大多数,已在《收获与
播种》的这里那里以这样那样的名分出现过,即便未必是作为我的葬仪的参与者。其中有四位,属于
我Bourbaki小组的朋友圈子;三十年乃至更久以前,我便与他们因工作而紧密相系(其中两位更有友
谊的纽带)。在这份石碑式的名单里,另有九位,我自觉以温暖的友谊之情与他们相连;那份情谊,到
我写这几行时还没有熄灭。可是,逝去的岁月里,不止一次,当我又面对这些昔日朋友中的某一位、或
我昔日学生中的某一位,一种奇怪的印象攫住了我:我心头那份依然完好的、向对方奔涌的同情之
流,所奔向的那个人,已经不在了——或者说,与那个人的联系已经断了,也许无可挽回地断了;另一
个人顶替了我从前认识的那个人——那曾是个充满热烈而颤动的生命的人——而新来的这位,似乎
已把他的一切痕迹抹得干干净净。那像是一场\textbf{干涸},一场枯竭;在一度有敏感而活的血肉之处,
生出了一层又硬又不怕水的壳……

在这本我才刚刚半开的家庭相册合上之前,我还想在匆匆插入相册的众人之中,单挑一位,稍稍多
停留片刻。他在这本相册里来得最晚。比之于我最终收进来的任何别人,在我这里都曾有(当然是不
自觉地)严肃的抗拒:不肯抛开一些关于我们关系的陈旧的现成图像,不肯认一份谦卑的明摆着的事
实。我说的是Jean-Pierre~Serre。

在《收获与播种》里,我不止一次有机会谈到Serre,通常指名道姓\fnote{855}{(*)}{在《愚妄与更新》里,有两三次我忍住没有点Serre的名;在那些时候,对我多少持批评态度的人,通常并不需要点名。《收获与播种》中我谈到Serre、谈到我们之间关系的最周详的段落,在“九个月与五分钟”、“兄弟与配偶——或双重签名”、“无用的细节”三条注里(n\textsuperscript{o}s~123、134、171(v))。}。这里那里说的那一点,想必已足以让人感到:他
在我的数学家过去里,扮演过一个无人可以相比的角色。这件事,在写\dmark{}\origpage{1118}《收获与播种》之前,我自己
也从未在上面停留过;它是我一路写下去才发现的。二十年间——从五十年代初,直到我告
别数学舞台的那一刻——他为我扮演了“得天独厚的对话者”的角色\fnote{856}{(*)}{1965到1969年间——Serre与我的关系依旧紧密——充当得天独厚的对话者的,更多是Deligne。缘由,想必在于性情上的强烈投合,尤其在于Deligne的一种开放态度(对我感到是我所当贡献者之精华的那种开放),而Serre常常缺这个。我下文还要回到这两种关系的截然不同的性质上——它们是我数学生涯中最紧密的两种关系。另见下一条脚注所引之注。};我的许多大的执念观念、大的投入,都是直接被Serre的思想(有时“貌似平淡”)激发
起来的。有某些时刻——我印象里尤其是五十年代下半叶,也许还有六十年代初——他与我之间有过
一种强烈的数学“共生”;我们的数学性情是互补的\fnote{857}{(*)}{关于这种互补性,以及Deligne与我之间的投合,见已引注“兄弟与配偶——或双重签名”(n\textsuperscript{o}~134)。};这种共生,每次都证明极其丰饶。Serre与我之间的关系,并不是“对称”的:比如,
Serre完全不像我这样,倾向于托付给一个或几个“得天独厚的对话者”,来了解他可能感兴趣、或他自
认为需要的东西。尽管如此(至少我揣想),我在他的数学家过去里,想必同样扮演过一个非同寻常
的角色;我不妨想象,我1970年那次不辞而别,在他的数学生活里,是一个断点(也许是某种平衡的断
点——在那平衡里,我代表着“阴”的一极),是一次骤然的转向,凭着一股突然出现的“空缺”。我不知
道……

不管怎样,Serre同我这个人、我的著作之间的这层紧密关系,数学界想必是知觉到的,哪怕它一直
在“不宣之列”。可以肯定:除Deligne之外,Serre被(有理由地)看作与我的著作“最亲近”的数学家。Deligne同我的著作、我的关系则全然不同——那是一种学生与“继承人”的关系。Deligne以我的思想、
我写下与未写下的著作为养分;而我的任何大的执念观念、大的投入,没有一样是由他唤起或激发
的。在某种意义上,他比Serre“更亲近”我:在他与我接触的那些年(1965--69)里,他身上没有Serre身
上的那种对我的著作、对我对待数学的进路的某些侧面的排斥反应;正是这一点,让他得以在仅仅\dmark{}\origpage{1119}
三四年之间(凭他出众的才能,也凭格外有利的种种机缘),把我此前数年诞生、发展的那套宏大的
统一视野,由内而外、整个地吸收进去。但他同我的关系,根子上是深深的暧昧——而且
他系统地玩弄着这层心照不宣的学生与继承人的关系:那对他是一种\textbf{权力}的手段;同时他拒认
这层关系,忙于把宗师、把宗师的视野一并埋葬……

Serre与我之间,没有这一类的暧昧——无论哪一方,任何时候,都不曾有过要从对方身上攫取“权
力”的念头、或把关系用作权力工具的念头。我甚至可以说,在我四十年代末加入的那个“Bourbaki
圈子”里,这类权力游戏并不存在;直到我1970年离去,我不曾目睹、更不曾(哪怕身不由己地)参演
过这类游戏\fnote{858}{(*)}{不过,我应当作一点保留,顾及某个完全在我不知不觉之中、在我某些学生之间、围绕我这个人我的著作进行的游戏。那游戏至少从1966年(SGA~5 讨论班结束之年)就开始了;第一个清晰可见的插曲,是Deligne1968年关于谱序列退化的文章(关于这一点,见“驱逐”一注,n\textsuperscript{o}~63)。我开始知悉这些把戏——它们确实是权力游戏——是去年,近二十年之后。诚然,积极的演员,并非当初接待我、我也自认同属的那个圈子的成员(在那个圈子里,即便凭着成熟给我的回顾,我依然看不出这类把戏)。他们组成了“接班人”。同样确实的是:我在这些接班人身上察觉的、相对母体圈子的品质败坏,与那个初始圈子的几乎每一位成员身上(几乎无一例外)悄悄发生的类似败坏,必然紧密相联。关于这一点,见“Bourbaki,或我的大机遇——及其反面”与“自深处”两节(n\textsuperscript{o}s~22、23)。}。另一种说法(关于Serre与我的关系,或关于我在Bourbaki圈子里
观察到的关系)也许是:任何时候,我在那里都没有察觉到丝毫\textbf{对抗}的成分——哪一方都没有
\fnote{859}{(*)}{不过,七十年代初“生存与生活”的插曲,应当算作例外。那一插曲让事情大白于天下:在许多我以为重要(今天依然如此)的问题上,我自己的伦理与意识形态抉择,同我的数学建制朋友们的抉择——Serre也在内——几乎全体处于对跖点。正是它,骤然终结了我同那个“建制”的认同感——我那时倾向于把它混同于一个理想(且田园诗般)的“数学共同体”。(关于这一点,见“‘数学共同体’:虚构与现实”一节,n\textsuperscript{o}~10.)这场意外的揭示,以及数月之内由此而来的“改换阵营”,当时使我采取了对某些旧友的敌对姿态;我从此把他们归入“反动派”等等。后来,我从这些斩钉截铁而肤浅的分类里退了出来。话虽如此,经过一个并不奇怪的翻转,Serre属于那样一些人:有一阵子,我把他们感知为“对手”,甚至“讨厌鬼”。我高兴地确认:这段插曲没有在他身上留下怨恨或敌意的痕迹——在我身上,也丝毫没有,这还用说么!}。偶有摩擦\dmark{}\origpage{1120},这确定无疑;前面谈到过,或许还要回头谈;但那是全然另一回事。Serre与我的关系之力
量,在我印象里,恰恰来自我们对同一位情人——数学——的共同激情,其间不掺任何唯我的“寄生”
成分:在其中,对方不是手段、不是工具、不是靶子。所以,最近当我同Serre恢复中断了十来年的通信,
我从他给我的两三封信的字里行间,重新找到完好无损的友谊与体贴的迹象,仿佛我们昨天才刚分
手。

再说,即便十多年没有机会通信,Serre偶一传来的音讯,全都朝着同一个方向:不变的友谊——而不
是葬礼的调门,后者是我许多旧友的情形。所以,直到最近的这几个星期,我都未曾想过:Serre会在
我的葬仪上扮演角色。凡从他那里传到我这里的一切、凡我知道的他的一切,似乎都朝着相反的方
向。而且确定无疑:单是他在数学舞台上的在场,就曾给\emph{葬礼}划定过某些界限(不得不承认,那
界限质朴之至……)。翻阅 J.S.~Milne那本《平展上同调》\fnote{860}{(*)}{普林斯顿大学出版社,普林斯顿,新泽西。正是这位 J.S.~Milne,两年之后,参与“难忘文集”《数学讲义》第900种的骗局(见“……以及掘墓”一注,n\textsuperscript{o}~168~(iii))。}(1980年出版,故而在那场难以置信的
“SGA~4½--SGA~5 运算”\textbf{之后}),我惊讶地看到:Milne“放心地”、几乎逐字逐句地沿用了Serre在
某次Bourbaki讨论班(1974年2月,n\textsuperscript{o}~446)上关于平展上同调之父的说法——理论“由Grothendieck发展,\dmark{}\origpage{1121}得 M.~Artin之助”\fnote{861}{(*)}{两年前,1978年赫尔辛基国际数学家大会上, N.~Katz(还是那个Katz)颂扬新科菲尔兹奖得主Pierre~Deligne的演说里,平展上同调理论被说成“由 M.~Artin与 A.~Grothendieck发展,方向为Grothendieck所预定”——可见字母序有时候还真管用……Milne在诸说之中选择沿袭Serre的说法而不是Katz,在我看来,是其善意的又一个征兆。}。有种种迹象表明,Milne只在 SGA~4 与 SGA~5 里零星翻过
\fnote{862}{(*)}{尤其令我惊讶的是:Milne(Mebkhout也一样——他可是我的著作的一位细心读者!)没有察觉 SGA~5 里有一条显式的Lefschetz公式,关于代数曲线上一般的上同调对应;那公式,被魔术师-骗子二人组Deligne与吕西耶妙手变走——说这是艺术家的活儿,再贴切不过!关于这一点,见“变戏法者——或飞走的公式”与“恭贺——或新风格”两条子注(n\textsuperscript{o}s~169$_8$、169$_9$)。};他一面沿袭Serre(在同一篇Bourbaki讲义里对 SGA~4 与 SGA~5 满不在乎地发议论),一面又沿袭Deligne(在其笔名“SGA~4½”的锯倒之作里,对这些讨论班肆无忌惮地拆台),为的是在他的引言里,把原
始文本 SGA~4 与 SGA~5 呈现为“难以进入”\fnote{863}{(*)}{至于 SGA~5 的出版版——托“编者”(原文如此)Illusie的“福”,它只是原讨论班一座毁容的废墟——Milne觉得它“难以进入”,情有可原。好撒玛利亚人Illusie竭尽全力(顺着好撒玛利亚人Deligne的心意),把它做成了一本难以消化的“技术性枝蔓”大杂烩……}——而这,恰恰是他的书(继Deligne三年前那本——说实话,边上
有点单薄——之后)理应来补救的情形;说得明白些:免得用户去费那无用而烦人的工夫,去读原始文
本。最高权威们的意见(这里首先是Serre,继之以Deligne,外加一位待在软壁棺材里一动不动、一声不
吭的死者……)——意见,一个Milne如每一个平常人一样闭着眼睛跟从(在殡葬的语境下,甚至抢着跟从……)——斩钉截铁地排除:这些文本会是别的什么,只会是“无用的细节”(或“废话矿
渣”……),而绝非一门新的“广义拓扑”(topos 版)的基础——这门新拓扑曾使(除其他理论外)Milne
书中所论理论的绽放成为可能;长远看,人们没法省掉这门新拓扑,正如没法省掉寻常的广义拓扑
——那种拓扑,Milne、Deligne、Serre(同我一样)有福气在学堂板凳上学过,所以他们顺从地(当作
理所当然)承认:那门拓扑的戏,值得那根蜡烛……

\pmark\origpage{1122}我想是去年,我才第一次匆匆看了Serre那篇Bourbaki讲义;我最近在“无用的细节”(n\textsuperscript{o}~171(v))
部分(a)“一摞摞千页的大书……”里谈到过它。Serre讥讽 SGA~4 共1583页的那一段,当时给我留下的
印象如此之淡,以至于一两个月前,为写《四种运算》重新拿起同一篇讲义时,我竟把它忘了个干
净。应当说,Serre对我那些出名的“千页大书”的这种拉开距离的态度,我早已深知——早在 SGA~4
系列出现之前——所以它毫无让我意外之处。我记得,这种“发自本能的排斥”第一次在Serre身上发
作、冲着我对待数学的某种进路,是围绕凝聚对偶理论——我在五十年代下半叶发展的。那才是潜
在的“千页大书”(至少),尤其若算上随附的一整套新的上同调代数——导出范畴版;但“大书”也罢,
潜在也罢,明摆着的一点是:Serre听都不愿意听,正如韦尔受不了看人白纸黑字写一个上同调群、
听人说出“拓扑向量空间”这几个字一样。

\pmark 然而这一回\fnote{864}{(*)}{其实,这已经只是我第三次手里捧起这份文本,它才“终于让我心头一动”。},当我重读Serre1974年的这份文字——底色则是一整年关于某场\emph{葬礼}的思考(那
葬礼,1974年已经“大行其道”四年了……)——那一段终于让我心头一动。它在我心中不动声色地起作
用,一天天、一星期星期。我意识到:Serre的那种态度——我最终习惯了它、它在我离去之前“并不
牵扯后果”——竟像是对那场葬礼的\textbf{一盏绿灯}。这一方向的头一件事,以明摆着的力量向我显形:
Serre的措辞本身(但“添上了恶意与无耻”),仅仅三年之后,就被一位Deligne急不可待地接了过去(或
更确切地说,怀着隐秘的快意),当作他那场难忘“操演”的“背景噪声”。

我在已引的那道注——5月4日——里第一次这样表态;这场思考在同注部分(c)(5月27日)“什么
都不像的东西——或干涸”里得到深化。那里,也是我第一次开启关于Serre与我之间关系的思考——\dmark{}\origpage{1123}
以\emph{葬礼}所提供的特殊照明\fnote{865}{(*)}{在前面一条脚注(打字稿第1117页注(*))里,我还点过另外两道注,其中我谈到过Serre与我的关系,但照明相当不同——那是“葬礼之前”的照明。}。写那几页时,我心中想必已经有一种弥漫的知觉:Serre在\emph{葬
礼}里扮演的关键角色。此后两星期,一场对整整一批事实与印象的整合与消化,想必在无声中进行;
抗拒对事物直接而细腻知觉的惯性之力,我相信,已不经战斗、不费力气地消解了。时机,在我看来
已经成熟,可以把知觉到的东西表述出来,尽我所能。

人们也许会想:Serre由来已久的、对我著作的某些侧面与部分拉开距离的倾向,不过起了一种阴差
阳错的偶然作用——可惜,恰好促成了一场同样阴差阳错的葬礼。然而那是一种肤浅的看法,碰不到
事情的深处。径直说到问题的核心吧:鉴于Serre同我这个人、我的著作之间独一无二的关系,鉴于他
对其同辈、以及其后各辈数学家的非凡影响,有一件事在我看来已经清楚:\textbf{若没有他对我葬礼的
心照不宣的首肯,这场葬礼原是举行不起来的}。

除了一位分明缺席的“死者”,这场\emph{葬礼}还有\textbf{两位主角}:他们的作为与不作为互相衔接、互
相补足,似乎没有丝毫摩擦与涂痕(但在我这里,并不打算谈什么串通——两位主角的调门实在太不
相同):他们是Pierre~Deligne,与Jean-Pierre~Serre。

关于前者,在这场关于\emph{葬礼}的长考的一开始,就已有长篇的交代:他代表着葬礼图景的“前景”,
作为葬仪大主祭,同时是那些由他倡导的运算的隐秘继承人与主要“受益人”(而且,早在死者的象征
性“死亡”之前就如此……)。Serre——这里第一次作为葬礼仪式的前排人物谈到他——则代表着图景
的“第三平面”,由“信众教团”构成的那一平面。

\pmark\origpage{1124}早在去年——或者更确切地说,早在我以它最粗粝、最悖谬的形状(以及这个名字)发现\emph{葬
礼}之前——我就清楚地知道:在一个我从不曾知道自己有仇敌的世界里,以这样的热忱来埋葬我的,
首先是\textbf{我昔日的}\textbf{朋友们};他们中的某些人,却并没有因此停止(哪怕只从唇边……)把自己算作
我的朋友——如今,我同样看清:在这些朋友之中,那些同时(而且尤其)是我的学生的人之外\fnote{866}{(*)}{在《收获与播种》的整个思考里,越来越清楚地显出:做过某个人的学生(比如做过我的学生)这个单纯的事实,就给一层关系打下烙印,赋予它一种特殊的质地,使它接近与父亲或母亲的关系。},真正
构成那场仪式的\textbf{台柱}的——作为教团的代表、作为全体信众首肯的担保人——正是他们中间
那位,在我们共同激情的层面上,曾与我最亲近的人。

Serre首肯的最耀眼的标记,当然不在我所熟知的那句随口的俏皮话里——那句话差点从我眼皮底
下溜走(虽说并没有逃过所有人的眼睛……)。真正令我驻足之后惊骇欲绝的标记,是\textbf{他所乐于
保持的蒙昧}:对那场就在他眼皮底下进行的葬礼——这样说毫不夸张——蒙昧到最后一刻
\fnote{867}{(*)}{可以说,在他那篇已引的1974年Bourbaki讲义里——他陈述Deligne对Weil猜想最后一角的证明——Serre的鼻子正埋在葬礼的正当中,却天真得连想都没想过记下一笔。我隐约感到他的不安:他发现自己面对一种表面荒谬的处境:我作报告(也在Bourbaki讨论班)给出某条 $\ell$-进上同调的 L 函数公式之证明的大纲,十年之后,我在其中承认(当作已知)的那条关键的“不动点公式”,居然仍未在文献中得到证明。Serre于是选择用一阵火气来\textbf{排遣}这份不安:讥讽 SGA~4 那出名的“1583页”(言下之意:而且\textbf{连}人们需要的那条公式\textbf{都没}提供)。这是一条省事的路:把一件不愉快的现实一笔带过(x)。然而他心里非常清楚(也许他乐于忘记……):在 SGA~5 讨论班里,我曾原原本本证明过一条不动点公式,其范围远远超出对 Frobenius 对应的那一条;他也知道,我的讲稿的成文,拖在那些所谓的义务“编者”手里,已经八年了。他即便乐于忘掉 SGA~5 的主题(“L 函数与 $\ell$-进上同调”——这标题好歹还是说了它该说的)及其内容,可他认识我二十多年,看我做数学看了二十多年,深知我把事情做半截不是我的习惯——恰恰相反(而且我做事情“不半截”到这样的程度,常令他不胜其烦,乃至忍无可忍……)。这本可以帮他刷新对 SGA~5 讨论班上所发生之事的记忆——他的脚,反正常踏进那扇门,足以大致知道我在那里做什么、进展如何。
显然,他既不想刷新记忆,也不想提问题。这只是众多例子中的一例:我的朋友宁可闭眼、塞住鼻子,也不肯去知悉一种他若接受就必然深深牵连他自己的真实。
(x)(6月22日)这几行写下之后,我有机会确认:这一类“不愉快的真实”,如今竟受到热烈的欢迎,如同天降的横财!关于这一点,见“家庭相册”的 d.、e. 两部分。}——那是一场对一部著作的葬礼;\dmark{}\origpage{1125}著作自其源头起就与他相系,比世上任何人都系得紧。而读《收获与播种》(假如他肯读这“一大摞”上千页的书的话……)会不会终于促使他动一动自己的鼻子(十五年来,馊味早就很大了……),以及其余的一切——对我而言,这是一个彻底的谜。可我深知:对他来说,同对葬仪的任何其他参与者一样(甚至更多),接纳我的讯息、动用他健全的官能,也意味着接受让他自己被深深牵连进去。

Serre的角色——率领那支前来观礼、为我葬仪合唱的信众教团——在我看来既典型,又例外。说它
例外,一是其极端性:他是与我最亲近的人,比教团任何成员都近;二是其非凡的身量(器局)
\fnote{868}{(*)}{还有第三种情况,赋予Serre在葬礼中的角色这种“例外”或“极端”的性质。他属于在我初次接触数学家世界时接待我的那群“善意的长者”。(关于这个群体,我有生以来第一次在“受欢迎的异乡人”(n\textsuperscript{o}~9)一节、继而在《收获与播种》引论(I~5,“一笔情愿的债”)里谈到过。)除了我们之间的友谊与好感之外,这也许是最主要的缘由,使得我花了一年多,才认了那份明摆着的事实,确认了Serre在我的数学葬礼中所扮演的关键角色。}:这
身量(器局)排除了深层动机中“补偿性的对抗”这类惯常的“寄生”成分
\fnote{869}{(*)}{这种“(表面上的)无因对抗”,我在这里那里已影射过两三次,尤其在(4月3日)“信使(2)”(n\textsuperscript{o}~182)一注(见下文)里。对我葬仪的绝大多数参与者而言——也许甚至是除Serre之外的全体——这种“原型式”的对抗肯定在起作用,这一点,我毫不怀疑。这股力量,在我看来,有别于那股通过“被自己否认的女人”之压抑(或“埋葬”)过程表达的力量。但这两种力量又紧密相联;在葬礼中它们结成一伙,以某种合金出现,常常难以分开。不过,我相信我已在其中辨认出葬礼中起作用的\textbf{两大力量}。可眼下,我很难说其中哪一方占先,是哪一方。我倾向于认为,是我先辨认出的那一方,即:对自己存在中阴柔侧面的压抑之力。
Serre的情形之所以在我眼前显得“典型”(同时又例外),想必是因为,在场两股力量中,正是后一种(我倾向于视为根本的那一种)在其中显出其全部力量,而丝毫没有另一方的痕迹(我称之为“寄生”的那一种——其义为:它会遮蔽对我所相信知觉到的\textbf{要害之物}的清晰把握)。我揣想(只要对原始事实与知觉的整合消化工作继续下去),未来几个月会带来更细腻的理解:在葬礼中、也在我以这样那样名分参与的其他冲突处境中,两股力量各占几分。}。如我方才强调的,我在\dmark{}\origpage{1126}Serre同我这个人、
我的著作的关系里,察觉不出丝毫对抗的痕迹;很清楚,在他的首肯之下起作用的深层力量里,也没
有对抗的痕迹。据我所知,除了那句出名的俏皮话,这份首肯纯然被动地表达:不过是一些\textbf{遗漏}
罢了。可是,这道给一场规模浩大的葬礼的心照不宣的“绿灯”——伴随的那些行动,有时庞大到仿佛
既定义了常识、又定义了体面——如今在我看来,恰是Deligne对同一场葬礼那积极而热切的参与的
不可或缺的\textbf{对应物}、某种意义上的“底片”\fnote{870}{(*)}{这里,Serre与Deligne在葬礼中的角色分配,有一处相当显著的\textbf{颠倒}:Serre的角色几乎全然被动,Deligne的角色则极度主动(虽说这位“操盘手”的角色,为事势所需、也合我朋友皮埃尔特有的风格,始终隐而不显)。然而,事实上,正是Serre其人,具有强烈得鲜明的“阳刚”主导;而Deligne其人,具有同样鲜明的“阴”(或“阴柔”)主导;并且(对两者皆然)无论在唯我机制的层面、“自我”及其条件作用的层面(亦即“\textbf{老板}”的层面),还是在发现之冲动的层面——那本源的、就其内在本性而言逸出条件作用的层面(即“\textbf{孩子}”的层面)——皆然。在Serre与Deligne这对截然相反的极端性情之间——葬礼的两大“台柱”——死者本人,则代表一种中间项:在“老板”一侧,阳刚主导强烈;在“工匠”(或“孩子”)一侧,阴柔主导同样强烈。(这一“基调”的分布,出现在注“兄弟与配偶——或双重签名”(n\textsuperscript{o}~134)里。)
阴阳角色之间的“对调”之力与机制,本是\textbf{主要}的思考主题:它催生了“阴阳之钥”这场长考,并作为底纹贯穿始终。它从《钥》的第一道注“肌肉与肚肠(阳葬阴(1))”(n\textsuperscript{o}~106)起就隐含在场,并在后续十一道注(n\textsuperscript{o}s~124、127、132、133、138、140、145、148、151、153、154)中或多或少走到注意的前台。而在这里,我意外地撞上了一个多少类似的“对调”处境——其动因,是葬礼中起作用的深层力量的内在逻辑。
近来,另一处似乎悖谬的阴阳角色“对调”,也打动了我——在这场不乏表面悖论的葬礼里!这一回,说的是“过早的死者”一方,与葬礼全体参与者一方,各自的角色。在集体无意识的意图层面上,这场对一个死者(按理,他只应安于与其身份相称的完全被动)的葬礼,首先是一场“女性气质的数学”的葬礼——一种带着强烈“阴柔”意涵的数学风格与进路的葬礼;而埋葬的教团,按理体现着“纯粹而坚硬”的阳刚价值,把那软塌塌的阴柔酥软之物交付给应得的蔑视。(比如见“阴的葬礼(阳葬阴(四))”与“天赐的机缘——或封神”两注,n\textsuperscript{o}s~124、151.)然而,处境的内在逻辑,却逼着这些“纯粹而坚硬”的参与者,在其中演一出典型的“阴”的或“阴柔”的戏:一场“天鹅绒爪子”的戏——半明半暗、沉默、遗漏、不动声色安插的暗示,或不断地暗示一事而佯言其反——总而言之,“喊停!”的那一套;我朋友皮埃尔是其中的至圣至师,而每一位埋葬者都不得不多少把它据为己有,事势使然。(关于这套风格,见“喊停!”一注,尤其“天鹅绒的脚掌——或微笑”与“对调(四)——或夫妻的马戏”两注,n\textsuperscript{o}s~77、137、138.)相反,是那位“死者”——那丰沛的阴柔酥软之物的化身——在人们最不提防的时刻,从舒适的棺材里坐起来,骤然接过一个他熟悉的“男子汉”角色:摊牌亮底,把好奇的鼻子和不逊的舌头、连同手电筒,伸进那些最精微暧昧的昏暗角落,粗声粗气地逐个点名,是猫说猫,是贼说贼——总而言之,一个真正的没规矩的家伙,一个不折不扣的搅局者,搅的正是那一团和气、细语呢喃的安稳局……}。而我朋友身上发生的这场变化,分量惊人。

\pmark\origpage{1127}我似乎确实深切地感到了在Serre身上起作用的那股力量。它所处的层次,比一种个人之间的对抗、或比谋求某种“好处”(按这个词的通常意义)更深。近来与他交换的那几封信,在这一点上很说明问题。说来,十五年来——自我离去以来的十五年里——我的朋友身上完成了一场\textbf{转变}
\fnote{871}{(*)}{“转变”(\textbf{transformation})这个词,立刻把我带到我朋友皮埃尔身上的“\textbf{变形}”(métamorphose)——那是去年10月他来访时,我第一次清楚知觉到的。(关于这一点,我在“拒认(二)——或变形”一注(n\textsuperscript{o}~153)里谈到。)“变形”一词更有力,对应着这一事实:在我朋友皮埃尔身上,发生了一次真正的翻转——从一种“阴”主导鲜明的本来性情,翻转到一副借来的、“彻头彻尾男子汉”的姿态。撇开这一点,我在两位朋友身上感到的转变,方向相同,动力相同:那股压抑“阴柔”特质的力量。},方向恰恰是那种对我对待数学的进路中的某些主导侧面的“发自本能的排斥”
\dmark{}\origpage{1128}。这些侧面,当年在Serre自己的进路里也曾存在,只是程度较轻——那是在他
作为数学家的过去里最丰饶的那些年:开放的、创造力澎湃的年月,尚未有一道\textbf{压抑}其创造个
性中这些侧面、压抑其内在“孩子”的工序立起来之前。这些,是创造力的“阴”的或“阴柔”的侧面与特
质。我在朋友身上以惊人的力度感到的转变,就是:从阴与阳两股创造之力和谐协作(而“阳”或“男
性”主导鲜明)的状态,转为一种“彻头彻尾男子汉”的失衡状态——“阴”的或“阴柔”的品质被毫不留
情地连根拔去。

说实话,正如两周前我已暗示过的(前引之注),这是一场演化的终点;五十年代,我已辨出它的最初
迹象,六十年代里它愈演愈烈。那时起,平衡就逐渐破裂,表现为视野的\textbf{窄化},表现为获准
上场的那份创造才能的清单的\textbf{缩水}。对我对待数学的某些主要侧面的排斥反应,以及渐渐地,
对一切真正构成我的著作的生命、深度与力量的东西的排斥——那不过是向外投射:是在他与我的
关系的层面上,一份远为深重的排斥的可触显形;被排斥的,是他自己存在的一个本质侧面,是他自
己的创造才能的一个本质侧面。

也许(如我方才暗示的),只要我还在跟前,同我的关系就对Serre的这一演化起着刹车的作用;在
他的生命里——五十年代,尤其六十年代——它代表着某种配重,由此成为一个相对平衡的因子。果真如
此,我的骤然离去,必定给那股压抑阴柔品质的力量放了行——那类力量,我并不陌生:它
也是作用于我自己生命的主导唯我之力之一;只是有一个显著的差别:在我这里,这份压抑之力只局
限于唯我机制的层面、我与他者的关系的层面,\dmark{}\origpage{1129}从不干涉我同数学夫人的恋情,也从不干涉(更一般
地)我在发现历险中的自发举步——无论数学的还是别的
\fnote{872}{(*)}{关于这股压抑之力在我自己生命中的角色,我在“超级爸爸(阳葬阴(二))”(n\textsuperscript{o}~108)一注里谈到过。我是在1976年开始察觉这股力量的;那一年,是我的心灵历险的一个关键转折。那转折,在“重逢(阴的苏醒(一))”与“接纳(阴的苏醒(二))”两条注(n\textsuperscript{o}s~109、110)里谈到。至于“阴柔”特质在我的数学工作中的主导地位(这些特质似乎在那里避了难,躲开了一切怀疑!),见“上涨的海……”一注,n\textsuperscript{o}~122。}。

回到\emph{葬礼}:眼下我做得最好的,莫过于引录11月10日那场思考的结尾几行——“阴的葬礼
(阳葬阴(四))”一注(n\textsuperscript{o}~124,打字稿第564页):
\begin{quote}
“……而这葬礼,骤然间以一种新的、意想不到的光呈现于我:连我这个人也变得次要了,我成
了那‘理应交付蔑视之物’的\textbf{象征}。这不再是某个人的葬礼,不再是一部著作的葬礼,甚至
不是一场不可容忍的异端的葬礼;而是‘数学中的阴柔’的葬礼——而或许更深一层,是在那众
多前来为悼词鼓掌的与礼者每一个人身上,\textbf{那活在他自己之内的、被自己拒认的女人的
葬礼}。”
\end{quote}
这最后一条直觉,是那一天、写下这最后两行的当口,以一道闪电突现的,像是一份意想不到的启
示,加在前几行所说的那份启示之上。此后几个星期,这道直觉一直作为底纹留在我的思考里,最终
在12月23日至26日连续的三道注里被重拾、深化:“天赐的机缘——或封神”、“拒认(一)——或召回”、
“拒认(二)——或变形”。

在这道直觉初次显现的那一天,以及在所引三道注的前两道——我在其中对它稍作探测——里,我
心里并没有一个具体的个案,除了(稍稍地)我的朋友皮埃尔(在第三道注里得到更周详的检视)。此
外我也深知:这一个案,对构成我\emph{葬礼}上出名的“第三平面”的信众教团之全体,绝无典型性。
所以,由于没有一个具体个案来作例证,我对那场突然瞥见的真实的把握,仍带着\dmark{}\origpage{1130}一层含混——被预感之物的那层含混:在某个层面上“知道”,却未曾充分而清晰地“看见”。我隐约
记得,我曾为这层含混稍感不安;曾有一阵冲动,想在我的朋友们——我知道他们参与了\emph{葬礼}
——找一个“有代表性”的,好把这份弥漫的认识“挂”上去,看它化身为一个可触的真实。

那时,Serre的念头,一次也没有掠过我的心头——他属于我旧友中那寥寥数人:我曾断然认定(至少
在意识的层面上),他至少不在参与我葬礼之列!而我那摸索着的思绪,当时(以及此前……)没有找
到那个在葬仪上须以一身体现“教团全体”的人,想必是因为,在我之内的某个地方,一清二楚:世上只
有\textbf{一个人}能演这个角色——而恰恰是这个人,我之内的一种沉重,让我以某种心照不宣而斩钉
截铁的禁忌,一开始就把他排除了……

如今,随着一场缓慢而幽暗的地下功课,那份沉重已经消散;现在,全光照亮了:恰恰是他——众人
之中恰是此人——那份寻找化身的直觉,在他身上以如此完美的方式落实,以至于人们几乎要相信:
正是他本人,让它在我之内浮现;正是他,从它显现的那一刻起,就赋予了它那种斩钉截铁、无可辩驳
的“已知之物”的力量\fnote{873}{(*)}{我甚至满心倾向于认为:这个“人们几乎要相信”,恰与事情的真实相合。这将再一次证明:我们的认识官能,远远超出那片可怜而微末的映像之域——我们只容许它们进入意识目光那片划得极窄的田地。(6月14日)前天的思考收尾的那道念头、那道骤然的直觉,也是在落笔之时以“闪电”方式显现的,没有明显的准备,也没有检视的打算。它带着某种“明摆着的力量”呈现在我面前。只是事后,我才记起:在紧挨着11月10日那段引文所出自之注的前一道注里,我恰好有机会相当周详地谈到Serre其人、谈到他与我的关系(而且,这还是《收获与播种》中的第一次)。}



\subsubsection*{d.\ 入殓——或自然的坡道}
\tocsubsub{d.\ 入殓——或自然的坡道}

\pmark\origpage{1131}(6月17日)现实,每一次都超出一切预感(无论多敏锐),甚至超出人们对它可能有的“知
识”;只有当我亲身与之摩挲——在路上的转弯处,而且往往全然出其不意——我才渐渐浸透它的滋味
与气味。即便这种接触似乎只是\textbf{印证}了已被预感或已被“知道”的东西,它却往往出人
意料,毫不客气地撼动一种几乎拔不掉的\textbf{怀疑}:对那确凿已知、已说、已写、已重述、又已重写的
东西的怀疑——尽管在某个层面上(那是巨大的沉重所在的层面),它依旧是一纸空文。我不止一次察
觉这份沉重\fnote{874}{(*)}{关于这份“沉重”、这份“不信自己健全官能之见证”,另见“义务已尽——或真话时刻”一注(n\textsuperscript{o}~163)第782至784页,尤其第782页脚注(**)。},我的不耐为之火起——一股顽强的沉重,执意要把我拽回熟
悉的观念与图像的车辙里,或拽回那些多少得到普遍首肯的观念与图像里;而我分明“知道”(或者说,
我之内的\textbf{别的}什么人、什么东西知道……)这些安顿妥帖的观念与图像是一场障眼戏,往往是明
摆着的障眼戏,根本站不住……思想,哪怕怀着要知道底细(那件既被“知道”又被否弃之事的底细)的
强烈渴望,单凭思想本身,无力抹去这份深深锚在自我结构之中的沉重。唯有与真实直接接触的斩钉
截铁之力,才有能力有时撼动这份沉重,侵蚀它、挪动它分毫——若谈不到真正抹去的话。

昨天我给Serre打了电话。只为一件求实的小事:关于Tate的“刚性解析空间”笔记——近日谈到的
\fnote{875}{(*)}{见“黑手党”一注(n\textsuperscript{o}~171)部分(c),“记忆的失灵——或新编历史”。}。我模模糊糊记得,那份文本原有一篇短引,交代这项工作的来历;我印象里,这篇引言在1971年
《数学创造》做的那一版里“丢”了。其实,Serre向我确认:Tate的笔记里,本来就没有这样的引言。那
不过是一些日志式的笔记,Tate把他关于刚性解析的思量寄给Serre,几乎像写信一样,而且(当然)
绝无发表的念头。我记得是我安排它们经法国高等科学研究所散发的(作者署名之下,加了一个小
标题“私人笔记,经(未经)本人许可刊布”\dmark{}\origpage{1132}),可我忘了Serre是中间人。反正,1962年刚性
解析空间诞生之时,除Tate与我,最“在局中”的就是Serre了。是他,或许早一两年,向我讲解过所谓
“Tate”椭圆曲线的理论——取一个完备离散赋值环的分式域 $K$。我被那——我记忆中是一场显式公
式的洪水(据说还是“经典的”)——砸得有点发懵;它们从我头顶掠过,没有“挂住”。但有一幅惊人
的几何图像留了下来,想必是Serre一句评点所致:归根到底,Tate椭圆曲线(至少它的“点”),是在
乘法群 $K^*$ 里“商化”、被一个同构于 $\mathbb{Z}$ 的离散子群相除而得。这正是复情形的类比:先把
$\mathbb{C}$ 除以一个因子 $\mathbb{Z}$ 得 $\mathbb{C}^*$,再除以一个因子 $\mathbb{Z}$,这次得到的就是一条椭
圆曲线。在复情形,商化在复解析的领域里有确切的意义,而Riemann-Serre型(即 GAGA 型)的定理保
证:最终的商(一条紧复曲线)具有典范的\textbf{代数}曲线结构。而在Tate的情形,唉,在 $K$ 这个完备赋
值域上那套让人不惯的解析空间里,得到的商是一个\textbf{全不连通}的紧解析空间,根本没有可能从
中取出一条椭圆曲线。然而(Serre当时必是这样对我说的)一切却仿佛……总之,Tate硬是靠着 $K^*$
和它的离散子群,凭着一手显式公式,造出了一条名副其实的椭圆曲线。

我记得清清楚楚:无论Serre还是Tate,都不相信Tate那套计算性的构造,会真有一个“解释”——以
$K$ 上一种“解析簇”的新概念来解释
\fnote{876}{(*)}{(1985年9月)正如去年7月我与Serre的通信所显现的,这里有一处我记忆的变形(Serre也有一处)。Tate(8月4日、10月16日)与我(8月18日、10月1日、10月19日)的几封信——都是写给Serre的——使事件的影片得以重排。是Tate(既非Serre,也非我)最先有了直觉与确信:应当存在一种“解析簇的新概念”,以便简单地解释“Tate椭圆曲线”的形式主义——那是在1959年8月前后。在我这里,它并没有“立刻让人心头一动”(照我记得的样子);我对Tate这一建议的第一反应相当怀疑,然后才开始思考这个问题。想必不久之后我就被说服了——一旦我察觉:既有的观念(尤其形式概形)无法解释与Tate椭圆曲线相联系的现象。此后两年,我相信,独自思考新概念的界定原理的,只有我一人;Tate与Serre都毫无头绪。这种情形持续到1961年10月:我向Tate提供了一门理论的总施工图。这立刻触发他去发展必需的基础,以便把握那些局部碎片(这项工作,若没有一个清楚的想法——日后如何把这些碎片拼装成整体的对象——本不会有太大意义)。更详细的评说与相关书信的引用,见《数学反思》第3卷预定的“历史评注”。}。

至于我,它当年是立刻让我心头一动的(*),我压根不曾想过,可以不把Tate曲线“看”作商化的结果
——为一个尚待提炼的合宜“簇”概念商化的结果——那正是让我着迷的那类工作!很可能也是Serre,
尽管满腹疑虑,向我提过:有一些人,至少克拉斯纳,在完备超度量赋值域——因而全不连通的域——
上“做解析延拓”。这似乎为我那个(有点怪诞的)希望添了水:不管怎样,应当会有一种“好的”解析簇
概念,比已知的更聪明,并且(凭“连通”一类的性质)接近实或复的、甚至代数的\dmark{}\origpage{1133}解析簇。可
再一次,在那个三人小组里,真正相信的只有我一个——至少我当时有这个印象。

这事在我脑子里转了好几个月,也许一年——说不清了。这处境让我想起一桩旧日的困惑:在当
时可用的概念框架里(靠赋环空间,概形与形式概形之类),人们无力赋一个形式概形——取所论的
离散赋值环 $A$——的\textbf{一般纤维}以意义。很快清楚:这本质上是同一桩困惑——我为给Tate的构造
以几何意义而寻找的那类“簇”,应当正是那种能给这出名的、尚不存在的“一般纤维”以意义的簇。
我终于有了第三条主线(克拉斯纳的风声之外):那是1968年出现的“广义拓扑空间”的直觉——当
时还没有得到“site”或“topos”这样的名字,因为我还没有开始边写边琢磨的概念工作——它将使我
得以定义那出名的、(隐含地)进入Weil猜想的“$\ell$-进Weil上同调”。这提示我:正如Weil上同调
的情形,我寻找的那种新“结构类”,不应到那老一套的寻常“赋环空间”里去找,或许要到那些“广义
空间”里去找——配以合宜的环层。

这些零零星星的直觉,何时才强大到足以推我开一个括号、从手头的任务里(主要是 EGA 与
SGA)抽出时间,动手做一点边写边琢磨的雏形工作——我说不清了。我知道的是:这项工作,一如既往
地,是在孤独中完成的——只有我一人“看见”那里有什么;也因此,只有我一人当时有能力做那把它
带到光天化日之下的最初的工作。我记得,我这里想几个钟头、那里想几个钟头,甚至有时一整天,
像逃学一般(虽然“日常”的工作并不缺!)。终于有一天,我咬紧缰绳,要弄个水落石出,真刀真枪地
干上了——我至少连续花了好几天,也许一两个星期。最难的是超越那些根深蒂固的思维习惯:它们
不停地要把我拽回已知的常规——“寻常”解析空间的车辙(我想,如今叫“松弛的”——德文作
“welk”)。我不得不重来三四次——一发觉自己又回到老辙里(马回马厩了!)就再爬出来\dmark{}\origpage{1134}。但分明,这\textbf{不是}那匹老马能顶用的时候……


这项工作做完,我心里有了底:再加一份我那时没有兴致去做的技术性工作,我已经立起了一个
“刚性解析空间”的观念(这是我给它的名字;“刚性”一词,为的是表达连通一类的性质——贴近代数
簇,而与所谓“松弛的”解析簇的性质恰成对跖)——它至少足以回应当时我心中的两个愿望:以这些
空间解释Tate的构造,以及解释形式概形的一般纤维。


我当时没有想得更远——我最急着回到那些被我暂时撂下的任务。倘若我再多玩一玩,我很快就会
发现:连闭环 $r\le r\le R$ 这样简单的空间(它们也配得一种“刚性解析”结构)都逃出了我的构造。是
Tate——我当然让他知道了我的思量——做了必要的调整,好把它们包括进来。除了我做的大半的
概念工作,还有一份更技术性的活儿要做:对所用的“砌块”(它们扮演仿射概形的角色)要有好的掌
握。这活儿,恰以其特有的优雅与细致,在Tate1962年的笔记里完成了\fnote{877}{(*)}{为求公允,我想可以说:我的工作与Tate的工作,代表刚性解析空间理论之绽放的两个阶段,同等不可或缺。我的那份,在于最初的视野(Tate与Serre都不曾有)与一份以概念为主的劳作——它也绝非没有必须迎面处理的技术侧面。Tate的那份,以技术为主,但其中同样有一份概念工作。我的工作,基调是“阴”的、“阴柔”的(所以,除了我已不在场,它还成了人人轻蔑的对象);Tate的工作,基调是“阳”,合乎好品味与好规矩的法度。\par 倘若我不曾把自己的工作推进到足以不再有任何怀疑的地步——怀疑一种好的“刚性解析”观念的\textbf{存在},以及对一门理论的总施工图的清晰视野——那么这门概念,恐怕今天还不会问世。因为,虽说人们“早晚”必然会发现并发展这个概念——它当然不是什么“发明”——其需要,在此后二十三年里,却从未以足够紧迫的方式出现,逼人“跨出那一步”。显然,是我第一个(1966年)预见到了刚性解析理论在最初的两个动机之外的应用领域:晶体上同调的发展。\par 此外,除我预见的三个用途——包括把Tate理论推广到一般的阿贝尔概形——我没有听闻别的几何应用。看来,后来“在这一行里工作”的人们,只从中看到关起门来发展理论的材料(既然理论已存在,而共识把它列入“严肃研究课题”),却不把它嵌进更宏大的几何视野。这是一个鲜明的案例:数学思想的原子化、碎块化——它与一切奠基工作的失宠相连,与一切不能归结为某种解“竞赛题”的技术绝活的失宠相连。一个特别雄辩的征兆是:无人尝试发展一个更一般的刚性解析空间观念——使Tate发展的那种,与“概形”之对“域上的代数簇”相仿——好把“变动的”完备赋值域上的诸刚性解析几何联系起来(尤其特征变动的,并包括实、复与“超度量”诸情形)。这一缺席,是近十五年数学在一切奠基工作层面上惊人停滞的众多征兆之一(这层面的关键,在此显然)。\par 回到Tate与我:同样可能的是,倘若我最初的“突破”没有在Tate那里“让人心头一动”、没有触发他的“第二回合”,刚性解析空间也不会存在。这事我原会在周围这里那里谈起的;可诱人的题目(包括看似更“紧迫”的题目)从来不缺,恐怕没有人会留意它——尤其今天,肯定不会:如今,引入如此怪诞之物的念头本身,就会让人想起某位,还是不点名的好……}

\pmark\origpage{1135}话说回来,还要过一阵,我才认了这份明摆着的事实:我所用的砌块,边边上短了点。它们对激
发我的那两个初始问题,已然够用——何必再往远里找!我很难改口。最终是Tate说服了我,以他那种
既平和又细致的方式:毕竟,不止这两个例子;就算我此生看来没遇到过圆环,也不构成排除它
们的理由。而显然,用我那套砌块,没有办法把它们“收编”(除非用无穷多块——那又差不多落回“松
弛”的老辙)。

我做好了我那份工作——在我自己的角落里,这很自然,反正我是唯一相信的人——不过当然,一旦走到(暂时的)终点,我也没有放过机会,把这事告诉两位主要当事人\dmark{}\origpage{1136}(实际上也是仅有的两位知情人):Serre与Tate。Tate那边分明心头一动;而Serre,我想,在我告诉他我的进展时,想必也被说服了。对此我没有确切的记忆,但倘若事情偏偏不是这样,我肯定会记得的。

\pmark 所以,昨天我给Serre打电话时,在我看来天经地义:他和我一样清楚,在这种新簇观念的诞生中,我扮演了什么角色。我没料到会有机会影射这件事;可当我跟他谈起Tate的笔记时,倒是他提醒我:那些笔记原封未动地刊在《数学创造》上,而且Remmert与另外两位作者刚出了一本专讲那些出名的刚性解析簇的书。这本书,就是我最近在“黑手党”一注部分(c1)(“记忆的失灵——或新编历史”)里谈到过的那本——我在其中痛斥Remmert“记忆的失灵”(可Tate的笔记本可以替他唤回记忆),服务于一种在我看来明摆着的蓄意刁难。我顺带向Serre提了一句——在给他写的上一封信里,我已有机会影射过某场葬礼\fnote{878}{(*)}{这是对那封信的复信里的话——那是我收到的Serre的最后一封信——其中Serre引了Siegel关于当代数学“扁平化”(Verflachung)的评语;我对它作了评说,并在“无用的细节”(n\textsuperscript{o}~171(v))部分(c)“什么都不像的东西——或干涸”里继续讨论。如我在该注所说,Serre曾把Siegel的这份印象斥为“不公正”——可我感到,Siegel竟这样想,多少让他心绪不宁。而这同一个词,他(肯定不是故意地)又拿来,打发我对一场葬礼的影射。不用说,他压根没有想到要问我:\textbf{什么}使我说有葬礼(我在信里一个字没提,宁可等他来问)。显然,案子早已审理完毕……},而这毫无疑问是一幅相当触目的写照。

头一件相当疯的事,是Serre(天知道他当年可是在最前排的!)——唉,连他也不记得了,而且是压根儿不记得了:我在那些出名的刚性解析簇里曾经有过一份!我简直瞠目结舌!这真叫人发疯——当我向他影射我以为自己有过的那份微薄份额、从激发我的那两个例子说起时,Serre记得的恰恰相反:差不多是我压根儿不想知道这些新簇,说(照他讲)有了形式概形,就已经有了所需的一切\dmark{}\origpage{1137}!当下我简直不敢相信自己的耳朵\fnote{879}{(*)}{事后回想,我明白了朋友的记忆(边上略有些失灵)发生了怎样的变形。既然我曾把形式概形当作主要的、几乎是唯一的向导,来提炼刚性解析空间的定义(以便给形式概形配一个刚性解析的一般纤维),他便(二十三年后)记成了:我曾斩钉截铁地坚称,不需要簇的新概念,因为“我的”形式概形已够用一切。(记忆的失灵,倒常常把事情“办”得挺好……)然而,连 $K^*$(我的第二条主线)也并非来自形式概形。反正,这里依旧是那句老话:案子早已审理完毕!}——然而不过几天前,我还刚刚以世上最平静的心情写下几页,谈的正是Serre在某场\emph{葬礼}里将要扮演的某个关键角色、某个“支柱”的角色。唉,这一回可好,我正正地站在\emph{葬礼}之中,就在我鼻子跟前、电话线的另一头,而且恰恰就是这位Serre本人,一如既往地从容自在,分明怀着世上最好的信念!(话说回来,我怎么也想象不出Serre会不怀好意,尤其在谈到数学的时候……)

我自然没有心思争辩,Serre更没有,不过到底还是有一场东拉西扯的闲谈,五分钟或十分钟。这十分钟真是用得其所:让我亲身摩挲到一场\emph{葬礼}的可触真实——色、香、味,一应俱全;那场葬礼,因为我一味只盯着纸看,已经变得有点遥远了。

我不得不说出的头一件事,就是“刚性解析空间”这个名字本身,是我给起的——还暗示,如果我未曾明说的话,那是在我还独自梦想着这些我如此命名的东西的时候。Serre有点愣住——显然,他同样不记得;但同样明显,我不是在编故事。不过不要紧,名字终究只是名字,而且毕竟“自然得很”……这句“自然得很”明白地暗示:正因为太自然了,它已不再意味着什么;任何人只要鼻子对着那件东西,都免不了就叫它这个名:“刚性解析”。总而言之,我的朋友就这个名字,给了我一句他并无此意的恭维——却配着“就这点小事……!”的调门。再说,关于这个,我什么也没发表过,不是么?那就无话可说了……

我越来越瞠目。发表也罢,不发表也罢,对我而言,现实分毫不改。一个女人怀胎九月,把孩子生到世\dmark{}\origpage{1138}上,孩子已经又蹦又跳、好好的,有人对她说:这不是你的孩子,因为什么也没发表,而且她也没本事
出示出生证——对说这种话的人,她必定当面嗤笑。说实话,我没有冲着Serre笑——那不是我的
做派;何况我还没回过神来。我也没有想争辩:连Tate本人在他的笔记里,都对我在理论起步中扮演
的角色没有丝毫讳言(此事,Serre看来同Remmert忘得一样干净
\fnote{880}{(*)}{我再一次感到:“无论如何,案子早已审理完毕”。既然Tate说他“完全忠实地”遵从我提供的一位总施工,那好——不过是个总施工罢了,一张潦草的图,随便哪个小把戏都能在沙子上画出来;一股淡淡的Grothendieck味儿,没错——Tate还算客气,真是够朋友,居然费神提了一笔……});而1972年,当我写“主题概要”\fnote{881}{(*)}{即那份日期为1972年的文本,对我当时数学贡献的一份略嫌干巴巴(也不太启人)的概览,为申请法兰西公学院教席而写(该席位后来给了 J.~Tits)。这份文本,加上更周详的历史评注,将刊于《数学反思》第3卷。《引论》之三(“罗盘与行囊”)里谈到过它。“主题概要”5 e)中,我写道:\textbf{刚性解析空间。}受“Tate椭圆曲线”之例、以及完备离散赋值环上“形式几何”之需的启发,我达成了完备赋值域上刚性解析簇观念的部分表述;它在其后 J.~Tate对这一观念的首次系统研究里发挥了作用。此外,我在特征 $>0$ 域上的代数簇上引入的“晶体”,有时可以解释为特征零域上某些类型的刚性解析空间上带可积联络的向量丛;这预示着特征 $p>0$ 的晶体上同调、与特征零的刚性解析簇上局部系统上同调之间,存在着深刻的关系。}——这还是“如假包换的Grothendieck”!——其中影射此事时,Serre也未曾有过一丝
皱眉(他的记性,想必从那时起就开工了)。反正都是白费劲——分明只要没有发表,我说什么都会被
当作刮过耳边的风……

可“没发表”三个字让我心头一动,我顺势接下去:恰恰,我的著作的很大一部分,就是没有发表的东
西,口耳相传的东西。我感到Serre又是一愣——在他看来,这想法似乎有点荒诞,像一种用词上的自相
矛盾:“著作——没发表……”,在他那里,这两者仿佛凑不到一块儿\dmark{}\origpage{1139}。我说出了“Motif”这个词;他
立刻扑了上去:这一回,他要纠正我那套葬礼观念,兴冲冲地向我宣布:两三年前,恰好出了一整本关
于 Motif 的书——真的,在“Motif”这一章上,我没有什么可抱怨的!

“那么,那本出名的书,你亲手捧过吗?”我问他(问得正巧——我琢磨着问他这个有趣的问题,已有
一阵了)。

捧过——但我也许是想笑吧,Serre回敬道;他当然熟悉那本书;他甚至像通读过全书的人那样谈起
它;想必他确实读过。倘若里面没有什么特别的东西,我本可以免了这一问——明摆着没有;可是(人
就是这样,我毫无办法!)我还是问了!见他似乎没听懂问题的意思,我告诉他:去年我捧起它时,我
几乎不敢相信自己的眼睛。

\pmark 我不得不说出了“诈骗”这个词,但我感到,那是一个纯粹委婉的说法。照我真正的感受——写
这几行时依然如此——那是一场\textbf{无耻之行};可我忍住没有把它说出口。说到底,我心里明白,用哪个
词都无所谓:十五年来,一直“馊得厉害”,而Serre选择什么也感觉不到(那正是我几天前刚写下的),
无论我说什么,都不会“通行”。

Serre这下打开了话匣子,仿佛他等的正是这个。诈骗?你做梦吧,可怜的朋友;可这本书正是Deligne
本人写的,还写得漂漂亮亮呢,没错——行,人人都很清楚是你引入的 Motif,可这不构成每次一说
到“Motif”这个词就要重念一遍经的理由,对吧?更不必说你从来没发表过一行字,而你的瑜伽靠的是
未证明的猜想(我觉得听见另有某人借Serre的嘴对我说话……);而那本书的妙处恰恰是:它不
用任何猜想,实际上,它不用你当年做过的\textbf{任何}东西……

语气斩钉截铁、不容回嘴,出自一位深知自己在谈什么、再无可学之人——还带一丝恼火:一个有
点赶时间的人,被一个死心眼的大胖子缠着,偏\dmark{}\origpage{1140}不肯明白最明摆着的事情。那气氛,容不
得人对任何东西稍作停留——一切早已定案、拍卖成交,分明如此。Serre的公理——关于这一行的伦
理、何为要紧、何为次要——分明变了;对此无可奈何。我只能连他的新公理一起,照单全收。

于是我无可奈何,只在“猜想”上接了话!我本可以对他说:Weil猜想也是猜想——然而他也好,谁也
好,从未想过把它们不当回事——但确实,那些猜想,Weil费心发表了!可我恰好在写“第六根钉子”
(钉向我的棺材)
\fnote{882}{(*)}{指的是那组注(n\textsuperscript{o}s~176$_{1}$至176$_{7}$)——我正在给它做最后润色;在那组注里,我拆解的恰恰是围绕 Motif Galois群观念、围绕Galois--Poincaré-Grothendieck范畴(应景地被封为“塔纳卡”的)的骗局——一场由Deligne策动、(最初)由“棋子”Saavedra经手的骗局……},于是我把话岔到“Motif Galois群”上:它可没有任何
“猜想”的成分;我发展过一整套高度精确的理论,关于Galois--Poincaré型的范畴——那正是那本出
名的书里用到的基础观念之一,而书里似乎认为无须对我这个人做任何影射。

Serre又扑向了这个影射;这一回,他又能纠正我的葬礼幻觉了——那套理论,白纸黑字发表在一本
书里,而且还是我的一个学生写的:Saavedra
\fnote{883}{(*)}{即同一位Neantro~Saavedra~Rivano的那本出名的“塔纳卡范畴”(原文如此),《数学讲义》第265种(1972),Springer出版公司。}
——难道不正是我让他做的那篇论文吗?这一回,分明又是一本他烂熟于心的书;他不止一次引用过
它
\fnote{884}{(*)}{我还知道:每当Serre有机会引用这本(几乎)不提我名字的书——而他本人并不觉得有什么不对——他总(出于我不知道的什么顾忌)小心地同时提一提我这个人。肯费这番神的,他该是最后一个了……}。“那么,那本书里,就再没有什么让你吃惊的了”——我又问他(而这一回,分明我早已知道
答案)。

没有,显然不曾让他吃惊:我的名字不在书里出现,无论是作为主题的作者,还是作为其中引入的附
属观念(诸如 Motif、晶体等等诸如此类)的作者——那些观念\dmark{}\origpage{1141}从头出现、作为例子展开,出自我
的又一位学生之笔,依旧没有我的名字。而这一回,Serre看来并没有记忆失灵——他还记
得(至少眼下……)那些观念归谁。倘若我这位朋友确有“失灵”,那总归不在“记性”这一层……

我们又就“塔纳卡范畴”这个名字聊了几分钟:我暗示,我认为它是一场骗局;而Serre有根有据,觉
得它再贴切不过。这一点,其实早在挑起这只新野兔之前,我心里就明白;我也知道\textbf{为什么}这名字,
对那位曾是我的朋友的人,如此贴切——而我这个孕育并分娩了这件东西的人,却觉得有话要说。

一如我们之间的常态,是Serre斩断了谈话——事实上,谈话确实拖得够久了。自始至终没有一刻“沟
通”;所以它留给我那份不满与不谐,想来正因为此。然而,正像最近收到的那两三封短信——以更强
的、无可辩驳的力量——这场简短的谈话,教了我许多。都是“知道”的事,没错,却被打了对折;知道,
却不信!而那份挫折感(至今没有消散),正是我在抗拒接纳那个讯息的征兆。

一个不速的讯息,当然。几个月前,我还毫不怀疑:Serre(一如我对他鲜活的记忆——犀利优雅、方正
不苟的化身),当他(迟做总比不做好……)凭借阅读“收获与播种”这份天赐的文本,得知某场\emph{葬
礼}的种种丑行(他当然万万料不到,可怜的人……)之时,必将拍案而起,投身这场战斗
\fnote{885}{(*)}{写“这一回”时,我想到了另外两次:我两次下本钱,试图向那个出名的“数学共同体”传递讯息——甚至动员它。第一次是1970年,我告别数学舞台之际,为了科学建制与军事机器的串通。第二次,层次低些,只在法国同事之间,为了一项打击在法外国人的不公正条文。(关于这一点,见“我的告别——或:外国人”一节,n\textsuperscript{o}~24.)两次,我的努力都遇到普遍的冷漠;Serre并不比我刚离开的那个圈子里的任何别的朋友(Chevalley与Samuel这两位除外)更例外。人们不妨就“收获与播种”这块大砖头,在那个建制里——Serre首当其冲——将产生(或不产生)何种效应,打个赌……}。这幅埃皮纳勒版画式的想象,在\dmark{}\origpage{1142}最近几个星期里消散了——多亏一次平淡的书信往
还。而昨天,我有幸亲眼看清,毫无半点疑义:Serre置身\emph{葬礼}之中,早已既久且深,而且他在其中
颇为受用。这一点——无须说(我也绝无半点讽刺)——是抱着世界上最好的信念的!

其实,我早已明白:“善意”绝非我大半生所觉得的那样,是一件如此简单、如此非黑即白的事。有一
种最常见的“善意”,不过是给自己递烟幕:一面体面的旗帜,遮着有时可疑的货物。我们的心灵由层
层叠叠的膜构成;目光越细,越看见某一层的“善意”,如何有时给底下一层的骗局充当掩护与托辞。

至于Serre的善意,我继续这样信任他:他绝不会写一本书、在实质上使用他人的思想而不明白说出
——哪怕那些思想从未发表、除了传达给他的人(倘若还在世)与他本人无人知晓。也就是说,我相
信我知道:Serre绝不会写昨天我们谈到的那一类书。我甚至可以说:单是\textbf{像Serre、像我}这样的人\fnote{886}{(*)}{在这里说“Serre或我”,我实际想的是:五十年代我们同属的那个圈子里的任何一位成员——那个圈子,我在《愚妄与更新》第三、四部分里试图略加勾勒,尤其在“Bourbaki,或我的大机遇——及其反面”一节里。不过,确实,我发觉即便在这个小圈子里,我也知道有两位成员“变坏了”(《收获与播种》里在适当的地方谈到过)。},
写一份面向公众的文本(这里是数学文本),就牵动根深蒂固的职业良知反射;它们(我相信)会消除、
或至少纠正某些“记忆的失灵”——这类失灵,在昨天那样一场随意的闲谈里,本算不得什么
\fnote{887}{(*)}{因此我毫不怀疑:倘若Serre是一本讲刚性解析空间的书的作者或合著者(Remmert正是),他不会由着自己滑上那条“自然的坡道”——把那位理当被略过的人略过;他会越过那些对那条坡道颇为迁就的“记忆失灵”——而他在一次私人谈话里,恰恰乐得滑上那条坡道。同样确实的是:若在十五年前,凭我当时认识于他的那份严谨,他不会任自己滑上这样一条坡道——我以为如此,哪怕只是在一次私人谈话里……}。而这一切,正合三星期前\dmark{}\origpage{1143}我在“什么都不像的东西——或干涸”一注(n\textsuperscript{o}~171(v),部分(c))里写下的话的方向:“我知道,Serre同我一样,绝不会想到跟着狼嚎,在‘人人’抢、骗、拆的地方去抢、去骗、去拆。”

话虽如此,我发现:这一切并不妨碍Serre——至少在某些情形里——实实在在受用于\textbf{别人的}抢、
骗、拆;而且是公开的、明摆着的,“在大庭广众之间”、“在聚光灯下”。他当然可以“怀着最好的信念”
如此——他不必脏自己的手:他只消对别人的掠夺、操纵与拆毁毫无保留地祝福;更何况他分文不取:
他不夸耀他人劳动的果实,却觉得别人(我几乎想说:那些领了专营执照的经销商)当众玩这种把戏
挺好。他收下的“收益”,比(边上有点不干不净的)出版物、比另一些人垂涎的银行账户,更精微。可
见它们毕竟举足轻重,足以造成我所认识的那个人令人瞠目的变形:今天(说不出从何时起),他闭
着眼睛、捂着鼻子,参与着普遍的败坏\fnote{888}{(*)}{这一对参与集体败坏的确认,与“掠夺的空白授权——或‘高尚行当’”一注(n\textsuperscript{o}~171$_{4}$,其名与其所指相符)——里所作的确认(针对1980年3月某讨论班的听众)相合;尤其打字稿第1090页第二段。}。

\subsubsection*{e.\ 最后一刻——或一个禁忌的终结}
\tocsubsub{e.\ 最后一刻——或一个禁忌的终结}

(6月18日)昨天我犹豫过:要不要给“家庭相册”那条注(n\textsuperscript{o}~173)再添第四部分,好“趁热”报
道前天与Serre的那通电话。诚然,那通电话留给我的是一种“不满之感、不谐之感”(如我昨天所写)
——而且这还是委婉语:那份刺痛逼近焦虑。这份不安催促我回到那段插曲,像对待一颗已经熟透的
脓疮:早该切开引流了\dmark{}\origpage{1144}。还有惯常的迁延。几个星期以来,朗格多克科学技术大学的复制室,就等着人
把《收获与播种》那册出名的分册IV的续篇送去——它迟迟不肯分娩;眼下时间已紧,就得抢印抢装,
赶在学院7月15日年度闭馆之前,何况不只我一家:学年将尽,各色论文蜂拥而至,人家有优先权。总
之,我对自己说:一本书总得知道收尾;若我继续往里塞“最后一刻”,明年这时候也完不了;够久了
……

可我终究还是动笔了——哪怕《收获与播种》的印刷要等到开学!它已经等了十五年(三十年也差
不多),再等两三个月无妨;只求容我从容看一看我该看的,说一说我该说的,不被种种“截止日期”推
搡……

那是艰苦的一天劳作——不如说一夜加大半个上午——我要让这份“加料”文本赶上今天的邮班。
事已成矣。

此刻,我有一种走到某项\textbf{必做}之工尽头的感觉。我忽然觉得轻了,仿佛卸下了一副我背了不知
多久、却从未察觉的重负。那想必是一种顽固幻觉的重量;它早在四十年代末,就在我之内扎下根来:
当时,一种领来的身份在我之内绽放——某个“共同体”(数学的)成员的身份、某个圈子的身份;那个
圈子,在我眼里满盛着温暖与生命。这种新身份的绽放,我在《愚妄与更新》里——“受欢迎的陌生
人”与“‘数学共同体’:虚构与现实”两节(n\textsuperscript{o}s~9、10),还有“Bourbaki,或我的大机遇——及其
反面”(n\textsuperscript{o}~22)一节——谈到过。诚然,这份认同,被1970年我离去前后的事件一扫而空,无影无
踪——那是在我投入战斗行动之后。可是,如今回顾,我意识到:对那个我已离弃、不再自认同属的圈
子,毕竟还系着一条\textbf{纽带}——一条也许无形、却力量巨大的纽带;它属于那份“过去的重量”(去年,
同名一节,n\textsuperscript{o}~50,我开始瞥见它)。我离开了那个圈子,没有回头的念头;可那个“家”\dmark{}\origpage{1145}——我为了
另一场历险而离弃的家——的一幅\textbf{图像},却在我之内活着,维系着那条纽带。这
幅图像,想必自我离去(以及更早,当然)直到《收获与播种》的思考,大体静止。《收获与播种》开始
修饰我对某段过去的图像,并勉力把现下的元素——常常令人错愕、常常不请自来——纳入其中。我
最终认了那份明摆着的事实:在我曾自认同属的圈子的后续、以及(据说)整个数学世界里,人心与
风俗发生了一场惊人的\textbf{败坏}。我意识到:这场败坏并非始于昨日;早在我离去之前,我就有足够的
时间,贡献过我的一份(至少,《愚妄与更新》的思考里瞥见过一份)。可是我有印象:我离去之
后,这场败坏有一场疯狂的升级;我的一些前学生,在其中扮演了头等的催化剂角色。

不管怎样——在对\emph{葬礼}的调查接踵而至的层层揭示之中,我心里始终维持着一种心照不宣
的“禁忌”,围着我旧日朋友中那些属于我青年时代那个圈子的人。我简直无法设想:他们中会有谁,
被这场我不断确认的深层败坏严重侵蚀、或“伤到”。当我说起“教团全体”对那些(至少在我眼里)超
出想象的行动的纵容时,我之内想必设有一道内在的“除外条款”,把那些在我看来应“立于一切嫌疑
之上”的人排除在外。他们分明毫不知情——他们想必在忙别的,肯定如此——可别怪他们!大致是这
样的调门。至于我的长者中最年长的几位,这种看法,我愿意相信,符合真实——或至少符合真实的
一个侧面。但对Serre、Cartier、Borel、Tate、Kuiper、Tits以及我熟知的其他人,却断然不然:他们与我
同代,正当盛年,完全融入我此刻检视的这个圈子,直到今天仍在其中行使着不小的权力、定着调
门,正如某些新来乍到者——他们在长者们的祝福之下,终于在那个圈子里结成了一伙毫无顾忌的
“黑手党”。

于是,在我据以看现实的那幅图像里——现实透过\emph{葬礼}这台头等的“显影剂”呈现——有一
道顽固而触目的矛盾。正是这道矛盾,无疑,\dmark{}\origpage{1146}在某个层面上被知觉、又在另一个层面上被否弃,造成了方才
说的那种逼近焦虑的“不安”——那焦虑,泄露了一场\textbf{分裂}。而那个人——比任何人都更能在我眼中体现这个圈子,体现这群被我内心的某个角色执意视为“亲近者”的人;而他本人,也曾是他们之中“最亲近”的一个——就是Jean-Pierre~Serre。正因此,被回避的矛盾之结,主要系在他身上。

\pmark 我开始怯怯地触及这道矛盾,不过六星期前,在“无用的细节”(n\textsuperscript{o}~171(v))一注的第一
部分(5月4日)。这场思考在同一注的第三部分(5月27日,即三周之后)——“什么都不像的东西——或
干涸”——里大为深化。一周之后(6月11日),我顶着顽强的内心抗拒,在本注的 c.部分(“众人之上的那
一位——或首肯”)里,重新回到Serre其人。这一回,Serre在\emph{葬礼}中的关键角色,终于大白于光
天之下。这是我对\emph{葬礼}的理解中的新的一大步——可矛盾的结,依旧没有触及!Serre其人,在我
这里(仿佛什么也不曾发生过)依然是无畏无瑕的“优雅”与“方正”的化身。“禁忌”完好无损!

是前天的那个电话,炸开了这道矛盾,把我的鼻子“直接摁进”(\emph{葬礼})里——不管我乐意与否。
随之,当然,是强大的抵抗之力(方才提到过的)的立即动员,以维持现状,而不是承认矛盾:以这样那
样的方式记录在案,由此解决它。我可以这样做,也可以不这样做。

\pmark 我跨出了那一步——而且我为此高兴。报偿是即时的:一场\textbf{解放},表现为轻快与如释重
负之感;一场内在张力的松解,诚然,但更是一场重负的卸除。

在《收获与播种》里,唯一可与之相比的解放时刻,同样标志了思考的第一次大转折:那是《愚妄与
更新》里的“竞技数学”,紧随其后的是“旋转木马收摊了!”(n\textsuperscript{o}s~40、41)。我还有印象:我方才“跨出”的
这新的一步,接续着去年那一步。何以如此\dmark{}\origpage{1147}、在何种意义上如此,我一时说不好。当时那声凯旋的欢呼
“旋转木马收摊了!”,肯定是过早了(次月我就察觉了)。可我刚迈的这新的一步,至少是又一步,把我带离那
台出名的旋转木马更远一步。将来会告诉我,实际情形是否如此。

经过昨天与6月11日的思考,我觉得已抵达一幅对\emph{葬礼}不那么含混的图景。此前,尤其
是这个“第三平面”模糊不清。6月11日的思考,让它“以可触的方式”化身于Serre其人;而这具化身,又
在昨天的思考里,取得了(这样说毫不夸张)一切具体的轮廓。

\pmark 说到底,在《收获与播种》这整个第四部分里,关于我与Serre关系的思考,在我看来,是对我理
解\emph{葬礼}最关键的部分——在“侦查补遗”与那幅数学大都会阴沟的浓墨重彩的图景之外。诚然,
倘若我不曾出于对我交给自己的题目的敬意,以我所能的全部细心,硬着头皮做这份“侦查的归置”,
并且尽力照亮沿途出现的每一个稍暗的角落,关于Serre的这场思考,恐怕也不会问世;而我对\emph{葬
礼}(以及我自身对它的牵连)的理解,将依旧如前含混。研究的工作里,一切互相维系!

\pmark \emph{葬礼}最后部分里那场思考的最有分量的部分,实际上是“最后一刻”才出现的。照理,这一
部分的“句号”,两个半月前(4月7日)就已画下。当时只剩十来页要重抄清稿、补几条脚注(一年前五
月末,情形也一样……)。意外之事,随即接踵而至:Zoghman来访——来读这(原则上已完成的)最后一
部分、向我提意见。它们凝结成了大约三百页的追加文字——其中,就有这几页:在葬礼(至今被回避)
的光照之下,我重新审视我与Serre的关系。


% ================================================================
%  18.5.7 -- 18.5.9 (原书 p.864–893)
%  源文: _work/chunks/chunk10.md, chunk11.md, chunk121.md, chunk122.md,
%        chunk123.md, chunk124.md
% ================================================================
\subsection{升级(2)}   % 18.5.7. L'escalade (2)

\notehead{174}{}

\dmark{}\origpage{1148}(3月22日)\fnote{889}{(*)}{(6月14日)眼下这道注,接续前一道注(同日所写)的 a.部分(“一位备受簇拥的死者”)而来。}我在别处已经强调过:其实并没有四种运算(对\textbf{一场}葬礼而言),只有
独一无二的一场“\textbf{葬礼运算}”。把它分成四大块,不过是为了叙述方便,其实是人为
的,而且(若太按字面当真)容易误导。确凿无疑,在那位导演——乐队指挥——葬礼主祭的身
体里,从来没有\textbf{四}只小魔鬼蹲在脑袋的四个角落里给他递台词,只有独一个!在关于阴与阳
的那场漫长冥想里\fnote{890}{(*)}{这就是构成《收获与播种》第三部分主体、Notes 104--162$''$的那场思考。},我曾试着与这只小魔鬼比从前更多地相识——从前,我不过时时察觉
它还在那儿折腾,转眼又去管别的。不过我并不声称已经充分认识了它;说到底,认识它或许根本
不是我的分内之事。确定无疑的是,它还像从前一样折腾个不停,而且未必会在我朋友咽气之前收
场。话虽如此,那场著名的“葬礼运算”至今仍在进行——就在我写这几行的此刻仍在进行。我不知
道,眼下这本“家庭相册”的散发,能不能至少终止全部局部运算中最庞大(也最不义)的那一项:
活埋一位年轻数学家——Zoghman~Mebkhout,“所有”在代数簇或复簇上同调里做工的人,四五年来
一直在用他的思想和结果……

既然分明只有一场运算,我就放下“四种”运算的虚构,不妨按时间先后,把我知道的主要插曲
与阶段勾一个轮廓。这里我不做了——我自认为已经做够了:在前四条主注(“沉默”、“操演”、
“分赃”、“封神”,n\textsuperscript{o}s~168、169、170、171)里,我把所知的全部插曲都收拢在一处,好奇
的读者自会把它在时间的标尺上排定次序。说来奇怪,若从“第二层次”或“运算”的观点(容我
用用这些委婉语)看,我1970年告别数学舞台的那一年,在这串插曲的接续之中竟不构成什么断
裂:在我印象里,它们自1966年 SGA~5 讨论班结束起,一直以相当均匀的节奏延续到1977年——
那一年,“SGA~4½--SGA~5”双双出版,SGA~5 又出了Illusie的\dmark{}\origpage{1149}编订版\fnote{891}{(*)}{(6月3日)这个印象应当修正,须把那场声势浩大的“Tannaka 范畴”(原文如此!)运算考虑进来:其第一插曲(与“稻草爸爸”N.~Saavedra搭档)在1972年(尾声则在1982年,由“真爸爸”Pierre~Deligne接力)。关于这一点,见那组注“第六根钉子(钉进棺材)”(n\textsuperscript{o}s~176$_{1}$--176$_{7}$)。}。这场运算,我看作一个骤然而惊人的\textbf{质变}。之前,是一次不动
声色的“收割”。从这里,我感到一阵暴烈与蔑视突然发作,死死缠住一个已被宣布“亡故”的缺席
者的作品。

在我那班上同调学生的这场上上下下齐心的集体发泄之后(“教团全体”在旁边看得心满意足),
似乎有过四年的平静。而从1966到1977的十一年间,我隔一两年就能辨认出一个类型分明的“插
曲”;1977到1981(“反常层会议”之年)之间,我却一无所知。相反,Deligne的长文《父亲猜想,
II》(La conjecture de père,~II),1980年发表于《数学公刊》——就在那场令人难以置信的讨论
会的前一年——照眼下的世道,几乎可以算作正常了……\fnote{892}{(*)}{当然,这项工作的主要结果没有一字提到我这个人;而那个结果的表述,本就是Deligne从我这里拿去的 Motif 瑜伽的一部分。反倒有一件事让我吃了一惊:我的名字,与Miller的名字一道,出现在文中一个段落里,说起那个除幂德拉姆复形——它大约在1976年由Miller和我各自独立地引入。1976年我曾在法国高等科学研究所就这个题目做过一次报告(顺便说,那是我一生中最后一次公开演讲),但明摆着我什么也不会发表。换了旁人,把这份纯属非正式的共同父亲身份略去不表,想必谁也不会察觉,更不会有人有半句闲话……}。也是在这一年,Deligne在一次Bourbaki讨论班上,从作者本
人嘴里得知了“好上帝定理”(又名Mebkhout定理)\fnote{893}{(*)}{(6月1日)其实,这一幕发生在前一年,即1979年6月,就在Bourbaki讨论班上。}。这是长期停滞的上同调主题冰雪骤消的开端。也是信号:第二年,葬礼
运算的第二场也是最后一场(?)高潮就此发动,这一回的调门是不义的,一切克制,连同最起码
的审慎,都被轻快地抛到了九霄云外。

翌年“难忘文集”LN~900那一幕(它为 Motif 掘墓复生,只字不提我这个人——去年某个4月19
日,这一幕曾让我百感交集……),还有同年Berthelot的报告(它把卑不足道的我从晶体上同调的
“历史”——原文如此——中一举清除),事后看来,都像是\dmark{}\origpage{1150}那场讨论会所发生之事的自然延伸,总
的说来也算得平淡无奇;那场讨论会的名字(或它的残余)也许会进历史,作为一次\textbf{警告}。
而翌年的“悼词”——无论在稍稍“摆一摆”它的人看来多么难以置信——同样像是这样一种延伸,
或者(如我先前所写\fnote{894}{(*)}{见注“宝石们”(n\textsuperscript{o}~170(iii))。})像一篇“尾声”。至于此后又过去的两年,它们不过是把一场光辉讨论会及其种种延伸
的“既得成果”,在文字里、在人心,一一盖棺论定……

这是一个显著的巧合——或者说,这分明\textbf{不是}“巧合”的效应——早在去年,在我还未见识
“SGA~4½--SGA~5”那场运算、也未见识反常层会议之前,我已经在我朋友皮埃尔与我的私人关系
中,辨认出恰好落在1977与1981这两个年头的两个“转折”。我第一次把它们纳入同一道目光、试
图探一探它们的意义,是在4月25日的注“两个转折”里——就在我发现\emph{葬礼}(读到那本难忘
的 LN~900)之后六天。两个转折发生的当时,早在多年之前,我远没有料到有一场\emph{葬礼}正在
暗中操办(至少在意识的层面上一无所知);要把其中任何一个同我所知的某件能照亮它的事联系
起来,我也会一筹莫展。

\subsection{殡仪馆——“为科学服务”}   % 18.5.8. Les pompes funèbres - “Im dienste der wissenschaft”

\notehead{175}{}

(3月23日)要把这场“葬礼运算”巡视完毕,还剩最后一位卖力的积极参与者的角色要看一看——
在这场关于葬礼的长考里,我已有机会屡屡“顺带”谈到他。这里说的是那家令人尊敬的Springer
出版公司(Springer~Verlag~GmbH,海德堡),科学书籍与期刊的出版者,大名鼎鼎——而且
还以“Im~Dienste~der~Wissenschaft”(“为科学服务”)的箴言而自豪\fnote{895}{(*)}{(6月1日)向Heinze博士打听之后得知,这其实算不上一句“箴言”,倒不如说是一条广告用语,英文作“Springer~for~Science”。}。这家出版社的数学书里,最负盛名的,大概要数“数学讲义
LECTURE NOTES IN MATHEMATICS”丛书。它恐怕也是世界上运气\dmark{}\origpage{1151}最惊人的科学丛书:二十来
年间出了一千多种。我想,我也为这份空前的成功出过一份力:这套丛书草创之初,正是靠我在六
十年代直至七十年代初,陆续放入的一大批我自己的或我的学生的著作给它撑的场子。我也曾以
“数学基础”(Grundlehren~der~Mathematik~und~ihrer~Grenzgebiete)丛书联合编辑的身份,与
Springer有过交往;三本书(包括 EGA~I 的重印本)经我的手在那里出版\fnote{896}{(*)}{另两本是Jean~Giraud与Monique~Hakim的学位论文(一是“场”与非交换 $1$-上同调的形式主义,一是任意赋环 topos 上的相对概形)。}。

1970年我告别数学舞台之后,我停止了一切编辑活动。仅凭惯性,我一直挂着那套丛书的编辑之
名,直到去年,才终于“正式”从Springer的一切编辑职务上退了下来。促使我这样做的是两个一
致的理由。其一,正当我重归“正统”的数学活动、重新发表数学著作之际,我要给这个“回归”划
出明确的界限:它对我绝不意味着重回某个“powerstructure”(权力与影响的架构),而只是回到
一份以发表为目的的\textbf{个人}数学工作。其二,自1976年起(经由Yves~Ladegaillerie
学位论文的那段插曲),我早有机会嗅到某种“葬礼的气味”——那时,对去年才发现的那场大规
模行动,我还毫无所觉。(关于这篇学位论文的插曲——它是我有幸启发过的最出色的论文之一——
见注“进步是挡不住的”(n\textsuperscript{o}~50),尤其是原委更详的注“二号棺材:或电锯式切块”,n\textsuperscript{o}~94.)这让我懂得:
“我所喜爱、我所愿意鼓励的那一类数学,在Springer出版公司里已经没有位置了”\fnote{897}{(*)}{这句引文(经翻译)出自去年2月18日致Peters博士的短信,信中我告知了他我退出“数学基础”编辑委员会的决定。Peters博士其实当时已经离开Springer(他如今在 Birkhäuser 出版社),通信便改由主管“数学基础”的Heinze博士接续。我曾要求把我那封信的副本寄给“数学基础”的每一位共同编辑(共十八位),并两次(1984年4月与1985年1月)向Heinze博士重申此请,他却始终认为无须告诉我副本究竟寄了没有(看来是\textbf{没有}寄)。我于是费神自己把信的副本一一寄给十八位编辑,并附了几句说明。他们之中我亲自结识七位,五位算得上我的朋友。只有一位(M.~Artin)费心回了信;其余诸位看来都没觉得有什么不对头——哪怕仅仅就\textbf{他自己}而言——:Springer公司居然没有费心把那封信(早在1984年2月)寄给他。};而且更重要的也许是,我在那里感到的气息\dmark{}\origpage{1152},并不鼓励我延续或重续与这家
公司哪怕稍微密切一点的往来。自去年2月我递交“数学基础”编辑委员会的辞呈以来,
过去的一年,只是把这份感受一再证实、一再加深。

不过这些都还是这场“葬礼运算”本体的边角——我昨天所说的那个“第二层次”,如今该回到它
了。据我所知,有\textbf{五本书}直接与这场运算相关\fnote{898}{(*)}{(6月1日)这几行写下之后,事情表明,还应在下面的清单上添入第六本书,它的书名本身就是一个骗局:“塔纳卡范畴”(Catégories~tannakiennes),Neantro~Saavedra~Rivano著。耐人寻味的是,这本书也出了Springer同一种“数学讲义”丛书。不过在那一桩运算里,Springer公司的责任在我看来并未牵连在内,不像其余五卷那样。至于“塔纳卡范畴”那场运算的细节,见那组注“第六根钉子(钉进棺材)”(n\textsuperscript{o}s~176$_{1}$--176$_{7}$)。}。按出版先后,它们是:SGA~7~I(1972年以我的名义出版)与 SGA~7~II(1973年
以Deligne-Katz的名义出版),整理的是1967/69年的单值群讨论班;题为“SGA~4½”(Deligne
编)的一卷,以及 SGA~5 的Illusie编订版(以我的名义出版,1977年);最后,是1982年以
Deligne--Milne--Ogus--Shih四人的共同署名问世、专为 Motif 掘墓复生立碑的那部“难
忘文集”。耐人寻味的是:这\textbf{五}卷皆出自\textbf{同}一家公司,皆在同一套“数学讲义”丛书\fnote{899}{(*)}{即第 288、340、569、589、900 卷。};前四卷出版时,主管“数学讲义”的是Peters博士\fnote{900}{(*)}{如我在倒数第二条脚注里指明的,Peters博士后来已离开Springer,转投 Birkhäuser 出版社。},最
后一卷出版时,主管丛书的是Byrne夫人。

这五次出版,其条件在我看来粗糙地不合规矩。我在别处已经指出:SGA~7~I 与 SGA~5 两卷
\textbf{以我的名义}于1972与1977年出版(LN~288与589),Springer却自始至终认为无须同我
打一声招呼,来征求我的同意,或仅仅把出版计划通知我。题为 SGA~7~II 与“SGA~4½”的两卷
的出版——它们打的也是 SGA 的旗号,而我以为这个旗号绝非人人可用,它与我的著作、我\dmark{}\origpage{1153}这
个人的干系,人所共知——事先同样没有征得我对使用这一旗号的同意,尽管(按理人们满可以期
待)我本该作为作者、或作为该卷(或它所整理的那个讨论班)的主编(或主编之一)列名其中
。最后,LN~900 一卷,在没有我的名分之下,呈现了那些概念、思想与构造;而凡消息
稍稍灵通的数学家都知道,这些是我引入的东西。在这一桩,明摆着(根本无须是 SGA~5 或 SGA~7
讨论班的那几个入室弟子之一)该卷就是人们通常所说的\textbf{剽窃}。我当然不指望Byrne夫人——出错
的话请指正——在该卷出版时主管“数学讲义”的那位——能凭自己的本事,看着稿子识破这场骗
局。但我想,一家严肃的出版社的本分,本该包括保证其出版物的严肃性,办法就是给自己配上胜
任的顾问。

这些顾问若诚实地做他们(想必)拿了报酬要做的事,也完全来得及向有关方面指出:SGA 这个
符号不是什么人都可以打的旗号,它自有\textbf{含义},而要尊重这含义,就该请教唯一有资格决定这个
旗号用法的那个人——也就是我自己。最后,就那卷顶着障眼法旗号“SGA~4½”的出版而言,作为
加重情节,只消翻一翻该卷的引论,或其后的那根“阿里阿德涅之线”,或第一章的引论,就能看
到:其中对待 SGA~4 与 SGA~5 两个讨论班,是何等漫不经心的轻蔑;而在消息稍稍灵通的人当中,
人尽皆知:那两个讨论班是六十年代中期的事,而顶着“SGA~4½”名目的一卷,装的却是七十年
代的伪作。依我看,对一个消息大体灵通、心智健全的人来说,这场骗局只能是一目了然的。因此,
更有十足的理由,在出版这样一卷、用上这样一个名头之前,先按正规手续征得我的同意。

依我看,在这五卷的每一卷的出版里,Springer出版公司的责任都牵连无疑——这五卷,正是围
绕我的著作、打着上同调主题所做的那场可歌可泣的诈骗大业中,一幕幕醒目的插曲。经由这些出
版,Springer公司成了这场非常行动的帮手与\textbf{贩运者}。我当然不能断言这是明知故犯。但我可以
说,这家公司在与我的关系上自\dmark{}\origpage{1154}1976年以来对我屡施的种种失礼(1970到1976年间,我想我没
有同它打过交道),也正朝着这场行动的方向,并且嵌入某种\textbf{精神},与它密不可分。

在注“操演”的子注“驱逐”(n\textsuperscript{o}~169$^{1}$)里,我曾提到我致Byrne夫人论 SGA~5 出版一事的信,和
她的复信——那复信,我必须说,真叫我目瞪口呆。(在这场“为科学服务”的光辉行动里,我“目
瞪口呆”既不是第一回,也不会是最后一回……)我从她的信里(日期是1985年2月15日)得
知:按照“一部著作包含多位作者的贡献时的惯常做法”(原文如此),无须专门知会我——我不
过是讨论班的\textbf{主编}罢了……SGA~5 的五位“作者”是Bucur、Houzel、伊
利耶、Jouanolou、Serre,独独不提我这卑不足道的人:我只以“主编”的身份列名——
想必纯属荣誉虚衔——为了这台光辉的讨论班,我说得未免太多\fnote{901}{(*)}{在那卷题为“SGA~4½”的书里那根著名的“阿里阿德涅之线”(贯穿 SGA~4 等等)之中,没有任何东西能让读者想到,我竟有幸在 SGA~4 与 SGA~5 里做过报告(相反,“有幸”为“SGA~4½”“合作”,倒是有的……)。关于这一点,见我在注“双重含义——或行骗的艺术”(n\textsuperscript{o}~169$_7$)里的观察,打字稿第899页。}。

早在收到这封富有教益的信之前,等得心里发闷(整整一个月杳无音信),我就拿出了我最漂亮
的笔(用德文),写信给 K.F.~Springer博士本人——公司主管负责人之一。那是一封漂亮的两页打
字信,我向他解释:公司与我的关系中绵延已久的一长串不快令我深为痛心;此外,还有若干起
对我不合规矩的粗暴做法,眼下我姑且奉上两桩,在我看来格外触目:“数学讲义”两卷(LN\textsuperscript{o}s~288、589)以我的名义出版,却认为无须征求我的意见;这两部书里,我在讨论班口头讲过的
思想、方法与结果,被删削,被残害,有时竟至于面目全非;这一事实,与环绕这两卷出版的种
种非常情形,凑在一起,对我而言绝不可能出于纯粹的偶然;我期待Springer公司公开地、毫无
保留地道歉,形式容当后议,一俟原则达成一致再行商定;我盼着他同我一样,把了结这一令人
不快、令人无法容忍的局面放在心上,找一个\dmark{}\origpage{1155}与事体相称的解决办法(“eine~dem~Fall~geziemende~Lösung~zu~finden”——这话说得还要更
雅些),落款则是规矩十足的“hochachtungsvoll”(谨上),出自我最漂亮的手笔。

开诚布公地说,我想我算是开诚布公了!这一回,Springer先生,总没法说,
没有人把情形亲自通报给他吧——而且是第一手的,通报者不是别人,正是当事的主要利害人
本人!

无巧不成书:回复(又过了一个来月)偏偏就在昨天到了我手上。它短得可怜,我禁不住要在此
全文照录(译文)。何况我还愣了一会儿,才明白这确实是上个月我那封漂亮信件的回复。回复如
下:

\begin{quote}
海德堡,1985年3月15日

亲爱的Grothendieck教授:

仍要感谢您2月9日的来信。Byrne博士夫人2月15日的来信想必已答复了您的问题。

此致,等等

K.F.~Springer
\end{quote}

这一回,至少我算彻底明白了!那班“消息灵通的人士”(前文已经说过的)想必跟他解释过:
何必为一个给他写信的、有点激动过头的先生费神——此人分明不属于上流社会。而且,倒也
不假……

还在等着收这封来自Springer殡仪股份有限公司管理层的、如此富有教益的回信的时候(承蒙
公司还肯屈尊,由大老板亲自署名作复,真是有礼),我有的是时间省察我自己的打算。这家可
敬的公司扮演的角色,在我看来实在太大;我一度设想过一场好看的大官司:以一位被冒犯的
“体面先生”的名义,作为无可名状的专横行径的受害者,索取天文数字的赔偿。可我又想:这样
一场官司,得吞掉多少精力。就算我赢了,并且拿到一笔叫人头晕的赔偿(往好处想!),那也是
若干年之后——我又落着什么好呢?我不缺钱,也不需要比我已有的更多;一场骗局,并不因为
某场官司打赢了或打输了,就多一分或少一分是骗局。我不会因此把世界变得更好,不会把我自己
变得更好,也不会把 K.F.~Springer先生和他麾下\dmark{}\origpage{1156}某些雇员的做派变好——当然更变不了他们对本
行的看法——即便我动员起我的律师,再让他们动员他们的\fnote{902}{(*)}{我还有过另一种念头:局面兴许会倒过来——这可敬的公司反过来告我一状,说我损害了它的名誉。这些“为科学服务”的人士,一碰到自己的名誉,想必是要较真的(只要涉及的是他们自己的名声……)。}。我也不会让某个上流社会里某种精神变得更好——那种精神,使得像Springer博士及其可敬公司(十三年来)甘为仆役的那类行动成为可能。我还剩下(我由衷希望)几年可
活——光阴似箭,而剩下的年月里,我满眼是叫人兴奋的事情。收集证据,好让法官们相信我与
SGA 有点干系——这类事想必不大叫人兴奋。我当年写它们,不是为了他们,也不是为了 K.F.~Springer先生,才费的力气……

至于那些(除我之外)我为之写作 SGA 的人:他们同这份东西——它(至少对我而言)仍是我
自身的一部分——所保持的关系,绝非与我无关。它是他们同我这个人的关系的一部分。说来奇怪,
这层关系我真正了解(或者说多少了解一点)的,只有我的五位上同调学生:正是托他们的福,
如今一位 K.F.~Springer博士才可以把我像块破布一样打发走,仿佛我对顶着 SGA 旗号的文字被如
何处置、或不如何处置,连发表意见的资格都没有——那位无名之辈的名字在不在封面上,都是
一样。

哪位数学读者若有耐心陪我读到此处,又若某一天也曾在 SGA(我说的是\textbf{真正的} SGA)里
出没,兴许愿意跟我说说他自己怎么看这件事。要是能收到一个字,来自某一位——他发觉,那部
我曾经独自一人整个儿扑在上面、花了十年生命、而\textbf{世上没有一个人}在工匠离场之后还肯费心
把它继续下去的著作——的确带着它的构思者与孕育者的印记,带着它成形之前那必需的全部岁
月的印记,并终于长成一间\textbf{人人可进的房子}\fnote{903}{(*)}{“盖房子”这个执念——而且是“对一切都好”的房子——在我的数学著作里扮演过可观的角色,早在五十年代初就是如此。这是我在劳动中的那股被我称作“服务的冲动”的具体表达;它属于那些赋予我的数学工作以生气的深层力量(尽管在“阴阳之钥”的思考之前,我自己都未曾察觉)。“房子”的原型第一次出现在我的思考里,是在11月26日的注“阴,作为仆人——与新主人”(n\textsuperscript{o}~135)里,来势之猛,出我意料。}——
收到这样一个字,我一定会高兴。况且,一间人人可进的房子,不是贫民区里一座公共小便池,谁
都觉得可以随意进去方便,在又破又黏的墙上涂写脏话……

\pmark\origpage{1157}而倘若读到这里的人,曾经是我的学生,或者曾经是我的朋友,却始终不觉得自己被
激发起写信的冲动、或当面谈一谈的冲动——哪怕只是为这件事,而不是为别的——那么让他知
道:他的沉默同样雄辩,而且,它是会被听见的。

\subsection{第六根钉子(钉进棺材)}   % 18.5.9. Le sixième clou (au cercueil)

\subsubsection{a.\ 预发掘}

\notehead{176$_{1}$}{}

\fnote{904}{(*)}{(6月16日)下面这一组注(n\textsuperscript{o}s~176$_{1}$至176$_{7}$),以“第六根钉子(钉进棺材)”为名,应当被看作“沉默”那组注(n\textsuperscript{o}s~168~(i)至(iv))的自然延续;那一组处理的是“Motif”运算,尤其是其中最后一道“预发掘”(n\textsuperscript{o}~168~(iv)),日期为4月8日。下面这些注——除最后一道(n\textsuperscript{o}~176$_{7}$)外——写于4月19日与20日。我之所以宁愿把它们撂在这里、“四种运算”的末尾,而不并入“Motif”运算,是因为此前几个星期里围绕其余三种运算——尤其是那场(所谓“反常层会议”或“当值无名之辈”的)运算——继续进行的思考,而那场运算正是“封神”那组注的题目——给刚刚出现的那个(同样出人意料的)“新事实”,投下了意想不到的光。这里提醒一句:写下面这些注的时候,原则上,我已经在\emph{葬礼}底下画过“句号”(其终极之注“朋友”(n\textsuperscript{o}~188)日期是4月7日),并且打算一两天内就把葬礼(III)的全部手稿送去打字。也就是说,这些笔记是在“最后一刻补遗”的心境里写下的……}

(4月19日)我终于有机会(4月10日)读到了“预发掘”一注(n\textsuperscript{o}~168$^{1}$)所引的Langlands的文章。照Deligne去年8月寄给我的那份 Motif“带评注书目”的说法,Langlands这篇文章,与Deligne刊在同一卷里的那篇(即前引之注的对象)一道,是自1970年我离去以来头一批用到 Motif 的文章\fnote{905}{(*)}{不过,1972年Kleiman与Saavedra的那些报告除外——它们承接了围绕 Motif 范畴描述的那几份微薄的“练手音阶”(可参照注“把一个个 i 点清”(n\textsuperscript{o}~164)里打字稿第794页脚注(**))。}。直到去年我还不知道Langlands的文章(Deligne的那篇也不知道),这情有可原:作者(同我的前学生一样)不认为有必要给我寄一份抽印本。人们也不禁要问:他何苦费这个神呢——翻一翻他的文章就清楚,我这卑不足道的人,同他文中所谓“自守表示、Shimura簇与 Motif”那个题目,严格说来毫无关系\dmark{}\origpage{1158}。我的名字(借用一句我的打字机恰好一年前就背熟的套语)在这篇文章里、在书目里,一处也不出现。不过,有几个我大约在1964年提炼的观念(或者是我梦见自己提炼过——看来我还在重蹈覆辙……),我自以为认出来了;我甚至把这个梦的追忆(或者也许是一个关于梦的追忆的梦……)白纸黑字写了出来,恰恰就在1984年4月19日\fnote{906}{(*)}{关于这一点,见注“一个梦的回忆——或Motif的诞生……”,n\textsuperscript{o}~51。}。我恍然觉得自己回到了整整一年前的同一天。

诚然,隔了这一年,我已有时间见怪不惊。纵有不快,也几乎算不上意外(人家会说,本来也没多
少东西……),更谈不上震动。何况,那本为三年后那部难忘文集 LN~900 鸣锣开道的先导文章,与
后者之间有一个不小的差别:我没有荣幸亲遇Langlands;他的“Motif Galois群”(或“基本群”)的瑜伽,
不是从我嘴里听来的(Deligne1965或1966年前后可是听过的)。不过,整个六十年代下半叶,我在
周围、在凡愿听者那里,把它讲了个够(Langlands总不是刚下船的外乡人……),所以我有一份揣测:
Langlands心里明镜似的,知道这门关于各种Galois群与基本群的新“几何”哲学——把它们看作合宜的
仿射 pro-代数群——是从哪儿来的。我猜他心知肚明:这门哲学,并非1972年从某位Neantro~Saavedra~Rivano的脑子里生出来的——此人此后从流通中销声匿迹,没有留下任何痕迹\fnote{907}{(*)}{据Deligne去年10月来访时向我暗示的,Saavedra实际上已经改了行(如今“搞经济”了),自1972年论文答辩之后,就再没碰过数学。}。我觉得,Langlands倘若觉得有必要,就此表个态,并不算什么奢侈。当然,照眼下的世
道,盼他费这份神,也许是我过分乐观了……

\subsubsection{b.\ 意外的惊喜}

\notehead{176$_{2}$}{}

\pmark\origpage{1159}好事从不单行。读到Langlands那篇文章的次日,我又有机会翻了一翻Neantro~Saavedra~Rivano的那本Langlands大量援引的书,书名叫“塔纳卡范畴”(《数学讲义》265卷,1972)
。

我离去之前有过九个(男)学生,Saavedra是其中唯一一个我后来再无音讯的;也因此,没有任何
风声向我提示,他多少沾上了某场\emph{葬礼}的“颜色”或“气味”。我于是草草下了结论——带着我惯
有的那份天真:他(哪怕只因为没了机会——听说他后来离开了数学的水域……)是所有学生中唯一
全然置身于“葬礼运算”精神之外的人。可是,同Jouanolou的情形一样,音讯稀少到这个地步,本该给
我提个醒。我当然知道:他跟我共事时预定要成为论文的那份东西,最终以“塔纳卡范畴”之名,于
1972年在“数学讲义”出了书;而那本书,上星期之前,我不记得自己曾费神看过一眼\fnote{908}{(*)}{(6月16日)Saavedra想必不认为有必要寄我这本书;我没有它的本子。不过,七十年代里,我也许亲手捧过它。我留着的印象,仅此而已:他做了一份细致周到、拿来就能用的活儿;至于这印象从何而来,我说不准。这印象尤其在场,是在写“清场”那注(n\textsuperscript{o}~67,尤其打字稿第252--253页)的时候;那里我评说了Deligne的这个“谜”:他几乎逐字“誊抄”了Saavedra跟我一起做的那篇论文。}。我满脑子别的任务,压根没想到:Saavedra从此杳无音信,连他的论文答辩也不
通知我一声、不请我进答辩委员会——按说我是最清楚底细的人——这有点奇怪。正是读了这本书,
才明白他为什么宁可不打扰我,而把论文“偷偷”过了关,交给一个我完全不知其组成的委员会
\fnote{909}{(*)}{这个答辩委员会组成的谜,在“第六根钉子”的第七道即最后一道注(n\textsuperscript{o}~176$_{7}$)里,以全然出人意料的方式解开了;此处不再多置一词……}。那时\emph{葬礼}已经在加速进行:因为委员会里没有一个人认为有必要通
知我答辩一事,更不用说请我参加委员会了(Jouanolou的论文答辩,大约也在同期,情形如出一辙)
\fnote{910}{(*)}{关于一次更正,见前一条脚注所引之注。}。

\pmark\origpage{1160}这本书,呈现的是那门“算术几何”的一个关键侧面——它的视野(总体眼光)六十年代里在我心中诞生
并成长(当时还没有名字),而 Motif 的瑜伽曾是(并且始终是\fnote{911}{(*)}{不过,其间,这份“灵魂”又添了“非阿贝尔瑜伽”的滋养;《纲领概要》里略微谈到过它。(关于这份文本,见《引论》之三“罗盘与行囊”。它反正将收入《数学反思》第3卷。)})它的灵魂。可以说,Saavedra这本书的实质,是对我的某些
思想的一份形状完整、细致周到的陈述——关于某些范畴的一门“Galois--Poincaré理论”(我做梦也
不会想到叫它们“塔纳卡的”……);那些思想,是我在Saavedra还拿不准肯不肯下功夫去熟悉、去
消化,以便能把它们收进其论文的“陈述性”部分的时候,长久的、不厌其烦地讲给他听的。我曾交给他
一批原原本本的手写笔记,命题钉是钉铆是铆,证明概略俱全;至今我还等着他肯把它们还给我
\fnote{912}{(*)}{把我的手写笔记这里那里散发给学生,按需取用,是我的老习惯——而他们要学的头一
课,就是辨认我的笔迹。向来讲好:他们用完,就把笔记还我——但据我所知,这个愿望很少被遵
守。这不过是众多征兆之一,表明我的学生并不怕我,而宁可把我看作一块“好脾气”的面团:干起活
来要求是严的,除此之外,随和得没边……}。当然,论文本身的题目,并不是去陈述一个其动机他完
全无从把握的人的思想。要做的,是显式地给出我在这里称为“Galois--Poincaré”范畴的那些“张量”范畴\fnote{913}{(*)}{为了不叫它们“Grothendieck范畴”!然而,在我有幸引入并命名(因而并不冠我之名)的众多范畴与新观念之中,若有一类非用这称呼不可——单凭体面,我就想写——正是这一类!(topos 除外,但它的名字我觉得本就好……)至于“塔纳卡范畴”这个名字——由一位天才的前学生悄悄塞进来的,又被全体一致的教团欣然采纳——这不多不少,正是一场骗局,如下文原原本本所述。(见“会等待的人……”之后的注,n\textsuperscript{o}~176$_{3}$。)}的一种“便利”的内蕴刻画,亦即:允许对它们作“Galois--Poincaré--Grothendieck式”的描述,把它们表达为基底环 $k=\mathrm{End}(1)$ 上一个“仿射(pro)代数 gerbe”的线性表示。当 $k$ 是域时,我曾以(我引入的术语)所谓“刚性”性质给出过这样一个条件;我记得自己在1964/65年、思考 Motif Galois群之始,就写出过一个完整的证明\dmark{}\origpage{1161}\fnote{914}{(*)}{我没有费时去我的笔记里核实这件事——关于 Motif Galois群的笔记(或者说,其中没交给Saavedra的残存部分)。反正第3卷里我要回去谈,多半在“Motif,吾爱”那一章里。}。我得把原理指给他;至于我在这题目上写下的笔记,我克制着没有交给他——因为
学他未来的行当,是他自己的事,不是我的事,活儿得他自己做。如果我记得不错,当时对我而言悬
而未决的,只有一件事:这样的Galois--Poincaré式理论,就基底环 $k$ 须满足的假设而论,其自然有
效域何在——我尤其关心 $k$ 取 $\mathbb{Z}$ 这样的环的情形(因为牵涉到 Motif 理论的应用)。

在我离去前的所有学生里,Saavedra是来得最晚的\fnote{915}{(*)}{如果我记得不错,Saavedra是1968或1969年要求跟我工作的,在我不期然告别数学舞台之前一两年。},也是准备最差的,而且(至少起初)最没
有干劲想“露一手”。所以我不抱希望:他会超出我给他出的那个非常有限的技术问题——它只要求最
微薄的知识(一点概形语言、线性代数、平坦下降、gerbe 的语言,如此而已)。而他书里第四
至第六章所处理的那些更精细的问题(纤维函子的滤套、实域上Galois--Poincaré范畴上的极化结构,
以及“可极化”的这类范畴的清单,及其对 Motif 范畴和许多变体的应用),要求的却是“四面八方”
的知识\fnote{916}{(*)}{这尤其要求对约化代数群结构理论、及其在实域上的分类的深入了解,还要熟悉整整一批观念:Motif、晶体、$F$-晶体、分层模、局部系统(而此人当初连拓扑空间的奇异基本群也只略知皮毛),再加上 Hodge 理论,以及那些从未在文献里显式陈述、只“藏在行间”的精微的“极化”性质。},从而要求可观的入门功夫;我不认为Saavedra供得起。我至多盼着他也许在论文后
面附一个摘要(多多少少由我口授),把那套理论中不会收进形状完整的陈述性工作里的要点列一
列。直到上星期我才被纠正过来,并且意识到:Saavedra交出了一份真正惊人、而且是以破纪录
的速度完成的活儿\dmark{}\origpage{1162}\fnote{917}{(*)}{关于这个“纪录”的更深的思考,以及它(明摆着的)解释,见注“Monsieur~Verdoux——或献殷勤的骑士”(n\textsuperscript{o}~176$_{5}$)。}。这活儿结晶为一本书:陈述细致周到、一丝不苟,分毫不差地
摹拟着、拿来就能用,以(在我看来)近乎穷尽的方式,呈现了我在六十年代发展的那套几何-代数
形式主义。从这个角度看,我认为他做了一份有用的、处处都无愧于人的工作;我方才说的那份“惊
讶”,确实是一份“意外的惊喜”。

这份工作,精确地说,在于把我提供的一整套观念、命题与证明,以“典范的”、可径直发表的形状
整理出来(按照我当年那些严格的标准)。这样的陈述性工作,当然属于数学家这一行的本分——无
论处理的是自己的观念与结果,还是他人的。与我许多同事不同,我不认为:在评估一篇论文或任
何出版的质量时——及至授予做这种工作的人数学“博士”的头衔、把他当作不折不扣的数学家
时——这样的工作可以被当作可忽略的量。不过,有一点在我看来至关重要:这一行的某种起码的
伦理必须得到尊重;凡工作在于陈述与发展他人观念之处,这一点必须清清楚楚地写明,不留丝毫
暧昧。

然而,在眼下这个案子里,全书除了一篇光辉引言末尾附着的、三行纯客套的、含混其辞的“致
谢”\fnote{918}{(*)}{这篇引言的实质,是逐字照抄四条主要命题——那是我向Saavedra指出的、有待发展的Galois--Poincaré瑜伽的“四根台柱”(至于纤维函子上的滤套所牵涉的问题,它们很难被概括成一条石碑式的命题,故未列入);但其中一条——本该构成其论文“最低纲领”的那条——被添加了一个显赫而明摆着的错误,使得它平凡地成为假的!下一道注(“会等待的人……”,n\textsuperscript{o}~176$_{3}$)就会谈到,尤其见已引的注“Monsieur~Verdoux——或献殷勤的骑士”(n\textsuperscript{o}~176$_{5}$)及其后一道注“肮脏的差事”(n\textsuperscript{o}~176$_{6}$)。}之外,没有任何东西能让读者猜到:书中展开的任何一个主题——从全书的主题本
身开始——跟我这卑不足道的人有什么相干。我恍然又回到读到那本纪念性的 Motif 掘墓文集的
那一天(到今天恰好整一年,一天不差)!我的名字在全书里几乎无处可见,只有两三处例外:那里
需要形状完整的参考文献,而除我之笔,别无可用者\dmark{}\origpage{1163}。

而且,这绝不仅是——怕被人看清作者“只是”在陈述别人的观念与结果而感到的那点尴尬(这——
尤其在这个案子里——若工作做得聪明,已经很不少了)。我从许多不会骗人的“小细节”里看清:这
里绝不止是稍微“收割”一把、给自己的纹章镀点金、然后退回幕后的意思。这真是\textbf{为葬礼而举
行的葬礼}。只举\textbf{一}个例子——上帝知道,我曾花了多少天、多少个星期,向Saavedra原原本本、
翻来覆去地讲解那些观念:他初来乍到,一无所知;晶体、$F$-晶体(在特征 $p>0$ 的情形顶替缺失
的 $p$-进“系数”,好定义 $L$-函数……)、分层模(及其与局部系统的关系),最后还有一点 Motif
的瑜伽(暂时以标准猜想作启发的底座);这一切,为的是借一大把例子让他明白,我拿这些Galois
--Poincaré范畴要到哪里去;也为的是——以防(谁说得准呢……)他有勇气有毅力,在预定的“最低纲
领”之外,至少收进一章典型的例子。他心里一清二楚,不必我多费唇舌:这些都是关键的几何观念,
而且并非开天辟地就有;把它们讲给他听、一遍遍讲给他听而不知疲倦的,不是别人,正是我——
它们是我在此前五到十年间引入的,用来伺候某一视野(总体眼光)的工具(即使那视野从他头顶掠过,
正如它掠过我所有学生的头顶,只除了一位\fnote{919}{(*)}{此人在师父一转身之后,就按捺不住把它埋葬了……})。可是,在他引入并稍稍发展这些观念的地方(在专述例子的第六章里),也如同在他佯装以作者自居的那部分理论正文里,我的名字照样不见踪
影。然而我很难设想,Saavedra会真的以为:读者无论多么孤陋寡闻,哪怕满心愿意相信他是这些
范畴(他慷慨地称之为“塔纳卡范畴”)的父亲,也不至于进一步以为:连 $F$-晶体、Motif 与“塔纳卡
装备”(原文如此)里的其他小玩意,也是这位Saavedra为应一时之需发明出来的。倘若这些观念被
当作信手拈来的寻常货色、或从最近的\dmark{}\origpage{1164}孤儿院里捡来的弃儿,那我在其中认出了一种\textbf{风格}——
一年来我为\emph{葬礼}兜圈子兜个没完,这种风格我再熟悉不过了……

是Mebkhout把那本书带给我的,他满心欢喜,总算能给我看一个我的学生的例子:这位至少还是
“诚实的”\fnote{920}{(*)}{(6月16日)再说,这事没成,他可是难过极了,还尽力来安抚我——这让我想起Kawai的情形(见打字稿第1078页脚注(*)),也想起了Beilinson:Mebkhout觉得Beilinson“比Bernstein更诚实”(见第1072页)——真像第欧根尼提着灯笼,只不过这一回,是在那伙一头扎进上同调这个名声实在不佳的题目的人当中,找一位诚实的数学家……}。他显然被引言末尾那三行致谢打动了——也难怪:1972年,感
谢某位已故者的风气,本来就已不大流行;从那以后,在我不止一位前学生那里,取而代之的调门,
不是冷嘲热讽,就是打趣,再不然就是彻底的沉默。话虽如此,这一回我到底得到一份“深深的感
激”:感谢我“把[作者]引入了这个题目”,感谢我的“劝告与鼓励……对完成这项工作必不可少”
\fnote{921}{(*)}{这些“致谢”,就情形而论,不过是一句玩笑话:人们会以为,我把作者“引入”了单复变函数这一“题目”,或任何同类的经典题目。其实,当我在一位正愁没论文的Saavedra面前谈起这桩事时,这个“题目”还并不存在——除非是在我之内、与 Motif 共生而发展出的那一份视野里,以及在我赋予它形状的手写笔记里。关于这份视野的诞生与发展,我在“一个梦的回忆——或Motif的诞生……”一注(n\textsuperscript{o}~51)里谈到;至于我的一位前学生(在众人纵容的目光下)对这根源所作的、漫不经心的彻底抹除,见其后一注“葬礼——或新爸爸”(n\textsuperscript{o}s~51、52)。(6月16日)Saavedra的这些致谢之所以更是一句“玩笑话”:其作者从来不曾费心,哪怕只寄我一本他的书、连同这些冒牌的致谢。既然我已经把“塔纳卡范畴”(原文如此)这场运算巡视完毕,如今我更明白:我这位前学生没有理由为他的“工作”(原文如此)感到骄傲,也不急于让我看到它。而照当时、乃至直到两年前事情的走势看,那位工匠似乎根本不大可能有机会看到它……}。这就是所谓的拿话头搪塞人;而在我看来,在\emph{葬礼}分明大行其道的时刻,在陈述
自己的工作时保有起码的诚实,才是表达“感激”的更有力的方式。
\subsubsection{c.\ 会等待的人}

\notehead{176$_{3}$}{}

\dmark{}\origpage{1165}说实在的,这本书一到手,我就看清:在那场难忘的“SGA~4½--SGA~5 运算”之前,整个
\emph{葬礼}里,还没有哪一幕的分量,赶得上这本名叫“塔纳卡范畴”——一个再平淡不过的书名——的
LN~265。此前的插曲\fnote{922}{(*)}{这里说的“插曲”,在“埋葬……”那条注(n\textsuperscript{o}~168~(ii))里匆匆勾勒过;它属于“Motif”运算那组注。}不过是或隐或显的“收割”,把某些重要思想的承续隐瞒起来。而这
里,我的“算术几何”视野的一个整个的关键侧面,被人不动声色地“挪用”了——经手的,竟是我的
所有学生当中看上去最“不起眼”的一位!

\pmark 的确,在这位学生背后,我凭一种不会骗人的笔法,认出了那个牵线的人;而且此公在Saavedra极尽殷勤的致谢名单里,赫然排在前列\fnote{923}{(*)}{就“数学”本身而论,这几位依次是(按出场顺序)我自己(不按字母序,算他有礼貌)、Berthelot与Deligne。}。Saavedra给这部书、给我引入的那个关键观念所起的名字,是一件精巧
的\textbf{褫夺之作}。论其斩钉截铁的效力,五年之后才被另一个名字超过——这回,那名字出自Deligne
本人之笔,落在另一部书上\fnote{924}{(*)}{正如下文(已引的注“Monsieur~Verdoux——或献殷勤的骑士”)将显现的,至少有强有力的揣测:与其读作“但这回出自Deligne本人之笔”,不如读作“并且同样出自Deligne本人之笔”。}。

\pmark 如果说,顶着一卷“锯倒之作”的“SGA~4½”之名,是一桩天才的冒名顶替,那么“塔纳卡范畴”
这个名字,就是一桩同样天才的\textbf{骗局}。哪怕对一个“平凡”或“中性”的Galois--Poincaré范畴——
等价于域 $k$ 上一个仿射群概形 $G$ 的有限维线性表示的范畴——我发展出的那套瑜伽,也十足是
“Grothendieck式”的:它由我早先在拓扑空间、概形或(更一般的)topos 的基本群情形发展的类似
瑜伽启发而来。把基本群定义为“覆盖层范畴上一个纤维函子”的自同构群,以及(同样古怪,因为
新,故而不惯)系统地同“不必\dmark{}\origpage{1166}\textbf{连通}”的平展覆盖的范畴打交道的想法,当年没少招来讥讽。我从
来不放在心上,深知那班爱开玩笑的家伙——自以为在学堂板凳上学过Galois理论或Poincaré理论就
算懂了——没有一个真正懂过;而且直到今天,他们之中谁也不能,不逐字照抄我做过的东西、不照
搬我以范畴等价给出的覆盖的Galois--Poincaré理论的表述,独自迈出(比方说)稍一般些的概形的覆
盖的Galois理论的\textbf{最初几步}\fnote{925}{(*)}{“稍一般些”在此可以精确理解为“一个非正规的概形”。在我之前,代数簇的基本群(由朗与Serre)只在正规簇的情形引入过,办法是把它描述为其函数域的“绝对”投射有限Galois群 $\mathrm{Gal}(\overline{K}/K)$ 的一个合宜的商。}\fnote{926}{(*)}{今天,这种表述基本群与覆盖之关系的方式——哪怕在普通拓扑空间(局部弧单连通)这个可以说是“教科书式”的特例里——也开始到处流通,而对我们共同的先祖只字不提,这还用说么……}。

同样,从一个“抽象”的有限维线性表示范畴出发——带着它自然的乘法结构与它自然的“忘却 $G$
之作用”的纤维函子——重建一个(为确定起见,取域上的)仿射群概形,把它作为\textbf{这个函子的自
同构的群概形}——这个想法,既不归功于塔纳卡(他可从未要求这么多),也不归我那微末的前
学生Saavedra,也不归我最出色的学生Deligne(令我遗憾——不过那时他还没在跟前),而是典型
“Grothendieck式”的。由此得到 $k$ 上仿射群概形与配有 $k$ 值纤维函子的刚 $k$-张量范畴之间的完
美对应,道理亦然。再者,倘若碰巧(特征非零域上的 Motif 范畴,就往往如此)一个刚张量范畴
(不幸,或者幸运加倍……)偏偏\textbf{没有}纤维函子可用,那么“代数群”就该换成一个代数
“gerbe”——这个想法,早在年轻的Deligne在数学里还没听人说过“gerbe”这个词、做梦也不曾想
到会有类似东西的年代,就已经被阐发得淋漓尽致。同样,当Giraud在六十年代自告奋勇,发展一整套
二维以内的非交换上同调代数的武器库\dmark{}\origpage{1167}——栈、gerbe 与 lien 的语言\fnote{927}{(*)}{这套形象化的术语是Giraud引入的,代替我从1955年起使用的临时术语(有点随手抓来,如“局部性质的纤维范畴”之类不成体统的名字;那些观念的根本性质,要求的是斩钉截铁、掷地有声的名字)。
(6月16日)在他那本书的引言第一页,Saavedra说到了Giraud“引入”的非交换同调代数“形式主义”。这是众多地方中的一处:在其中,我感觉到有一个比本书作者更机灵的人,在“搀着他的手”……正是同一位,一提“导出范畴”,就乐于紧接着补一句“由Verdier引入”(而这两处的底细,他心里一清二楚……)。}——讥笑声同样不缺。这类东西,如今、而且早就被Deligne之流叫作
“废话矿渣”。这些讥笑不曾让我却步\fnote{928}{(*)}{Giraud却不然——他在刚与我一同开辟这主题不久,就与之决裂,再无回头。关于这一点,见“共同继承人……”一注(尤其打字稿第386--387页)及其后一道注“……以及那把电锯”(n\textsuperscript{o}s~91、92)。};我知道自己在往哪里走——我看着这块“矿渣”以
完美精细的手法,抓住那些微妙而深刻的关系(我深知,任何别的“语言”都无力把握它们),心中是
“复苏”(如我在别处所写),却谈不上惊讶。

话虽如此,当这班讥笑者某一天发现一道被自己漏掉的“奶油甜点”——无论是那被人抢着命名为
“塔纳卡的”(且听下回分解……)范畴,还是某种“对应”或“关系”或“构造”(边上带点新Grothendieck味儿),被人用委婉语一笔带过、或被命名为“Riemann--Hilbert”(同样且听下回分解……)
\fnote{929}{(*)}{关于这后一个“且听下回分解”,见“反常层会议”的全套,尤其注“魔术师”与“冤大头买卖——或木偶戏”(n\textsuperscript{o}s~75$''$与 171$_{2}$~(e),后者属于长注“黑手党”(n\textsuperscript{o}~171$_{2}$)的一部分)。}——这时人人蜂拥而至,争当光辉的发明者。这就是本世纪七十年代、八十年代数学的
“时代精神”……可以确定的是:给这些范畴(我曾长久讲给他听的)起“塔纳卡范畴”这个真正光辉
的名字,这样的主意,不是一个Saavedra想得出来的。放任他自己,他绝不敢改动从我这里承袭的
术语,至少得先征求我的同意——这原是最起码的规矩!非得有上头的榜样与鼓励,他才敢\dmark{}\origpage{1168}这样把我
当作可忽略的量打发掉。再说,这可怜人为了熟悉必不可少的东西——若他想完成我交给他的那个雄心勃
勃的写作计划的哪怕一部分——手头的活儿已经够多了\fnote{930}{(*)}{他以仅仅两年的创纪录时间完成了这个计划——从我离去(那时计划几乎尚未动工,除了对概形基本技术的初步熟悉)算起。即便有Deligne在一旁搀扶(在我离去之前,他对这个学生毫无兴趣),这份成绩也简直堪称奇迹;这个“奇迹”,在注“Monsieur~Verdoux——或献殷勤的骑士”(n\textsuperscript{o}~176$_{5}$)里得到较细的检视。},哪有余暇再去翻文献、去读塔纳卡以及天知道什么的东西——他跟我共事的时候,肯定
从没听说过这些\fnote{931}{(*)}{提醒一句:Saavedra跟我工作,恰好在我离去前一两年(约1968、1969年),此后我几乎完全失去了他的踪迹。他那时的行囊,同任何一个来自第三世界(或我们外省院校)的三年级研究生相比,不厚也不薄。}。

这个名字之所以“光辉”,在于两种看似矛盾的品质的精巧结合。其一,在肤浅的观察者眼里,这
名字不显得全然荒唐。“人人”都约略记得存在一个“塔纳卡对偶”,乘法结构在其中扮演角色——
这看起来同那著名的 $\otimes$-范畴(某位Saavedra——这是谁呀?——称之为“塔纳卡的”)身上发生
的事,还真有几分像;那就“塔纳卡的”吧,为什么不呢!可是,对会等待的人,事情自己会成熟。十三
年过去;如今,顶着谁也没见过的无名作者之书的,换了一部三年前问世的、显赫得多的参考:光辉
的 LN~900,出自非别人而是Deligne之笔,外加一位名叫Milne的搭档。这两位大名鼎鼎的作者,从
头发展了他们同样称作“塔纳卡的”那类范畴的全部形式主义。显然,这是一个根本性的
观念,多年来已被Langlands、Deligne、Serre等人使用,前程似锦。当然,谁也不会相信:在这篇文章里被
顺带引了两三次的某位Saavedra,竟是这一关键观念、以及由它生出的那套精微形式主义的主人。
两位光辉作者的这篇文章,以其主笔人所熟知的全部娴熟,重拾这个题目,其笔调本身,就把这层意\dmark{}\origpage{1169}
思交代得没有半点含糊\fnote{932}{(*)}{关于这篇文章,尤其见“下葬——或新父亲”(n\textsuperscript{o}~52,尤其打字稿第214页)与“清场”(n\textsuperscript{o}~67,尤其打字稿第252--253页)两注。}。更不必说,他们还从Saavedra书里陈述的理论中,挑出一处粗糙已极的错误(以致
他们不得不从一个全然不同的定义出发——那个定义,终于像是正确的\fnote{933}{(*)}{关于Deligne的这桩绝技(Milne作陪衬),见屡次所引注“Monsieur~Verdoux——或献殷勤的骑士”(打字稿第1176页)的开头。}),叫人不免要问:这可怜的
Saavedra(当年必有某人——猜猜是谁……——费过心给他解释是怎么回事)到底懂不懂自己在
说什么。而Milne纵然光辉,且有幸与显赫的Deligne合署一篇文章、发展一个分明根本性的观念,
也决不会起念头认为自己能算它的父亲或哪怕共同父亲;Beilinson与Bernstein同样不会声称,他们发
明了(或哪怕共同发明了……)那句有名的“本该在这几篇笔记里找到自己位置的关系……”——他们
有幸与同一位显赫的Deligne合署了它,而在那之前,正是Deligne以他无上的好意,给他们指点了
Kazhdan--Lusztig 证明的门径……最后,又有谁会\textbf{认真}相信:那位把自己的名字(想必未经商
量)借给这一根本观念的、大名鼎鼎的塔纳卡,果真与此\textbf{有}几分干系?等到哪天世人都看清了这
个观念、以及随之而来的整套精微理论的\textbf{真正的父亲},他也不至于出面认领——倘若他还在世
的话。对此若还有半点疑虑,只消翻翻塔纳卡的著作——嫌多的话,就翻他论“塔纳卡对偶”的那
篇——就会发现:说到底,风马牛不相及……

这里,仍然是:几块界石一立,剩下交给时间。显然,这门理论——它将越来越多地显形为一种新的
\textbf{哲学}的技术手段,用以沟通几何与算术——在未来的年月里,注定要越来越走到数学舞台的中
央。五年、十年之内,再不会有人想到,就这个题目去援引某个无名作者的书;而那位想必搀过他手
的人,倒是不辞辛劳,写出了一份势在必行的陈述,由一位光辉的合作者协助,构成那部光辉的文集
的核心——Motif 的观念,终于在\dmark{}\origpage{1170}那片坚实的土地上得到发展。(至于那位含混而潦
草的先行者——他早已被遗忘——以及那道分明超出他能力的、惯常的猜测性“废话矿渣”,文集里
只字不提,倒也显得仁慈……)引用“P.~Deligne 与 J.~S.~Milne 的塔纳卡范畴”,将如同引用
FAC 或 GAGA(Serre的)或 SGA(法国高等科学研究所那部著名的匿名讨论班,人称“布瓦马
里”)一样,成为第二天性。而这么做的时候,任何人的心里,对这些新观念的父亲是谁,都不会有
半点含糊——那父亲断然不是合著者Milne,更不是塔纳卡,甚至也不是某位严格无名的作者(一个
叫Saavedra的):他在人家文章里被顺带提了两三次,为的是他(在他那卷书引言里)写过关于这
个题目的一份“出色的摘要”(带若干保留)。

不过,人们不会指望这门理论的父亲,让他去强暴自己那份出了名的谦虚,以至于把(按“消息灵
通”人士——他们在这类事情上说了算——的全体共识,理应如此称呼)本该如此称呼的东西,
叫作“Deligne范畴”(或者,在另一个天地里,“Deligne对应”)……

\subsubsection{d.\ 父亲们的圆舞曲}

\notehead{176$_{4}$}{}

(4月20日)昨天的思考,让我以新的眼睛看清一件去年——我刚一头栽进\emph{葬礼}的时候——曾让
我瞠目的事:“……这件看似荒唐的事:Deligne十年之后‘重做’Saavedra的论文*!”从去年4月19
日——我读到那本“难忘文集”LN~900 的日子——起,这件事就被谈到了:文集里(除其他美好东西之
外),Saavedra的论文被几乎逐字重现\fnote{934}{(*)}{见倒数第二条脚注所引诸注。}。一周之后,我在“清场”那条注里又回到它。当时我已得出“内
心确信”:这堆废话背后的意思,是光辉的Deligne(自任Saavedra的文书)的如下愿望:
\begin{quote}
“给自己一种幻觉中的解放感,解脱他肯定感到的一桩痛苦的义务:不得不时时援引那个他正要
取而代之、加以否定的本人,或哪怕任何一个援引了他的人。”
\end{quote}
可上星期,我第一次费心翻阅了这位“别人”的工作,惊讶地发现:他压根\dmark{}\origpage{1171}没想“援引我”(除那三行
冒牌的“深深感激”之外——明摆着是障眼的烟幕)。于是,一年前我的“内心确信”就站
不住脚了——里面肯定有对的成分,可仍剩下谜:总不至于就是那三行字(没有哪个读者会想到跑去
引言末尾把它们挖出来),竟能促使一个Deligne,去当一位早已作古的宗师门下最无名学生的抄写员
吧!更不必说,在这段引言的末尾,我几乎与他、与Berthelot并列出现:他们同我一样(人们会说
\fnote{935}{(*)}{只是有这样一点不同:是我“把他引入了这个题目”(原文如此),而且“他至今在数学家的养成上,还大大受惠于我”(这真是受之有愧)。}),有资格得到“他们在这一工作中慷慨给予的帮助与建议”的
致谢……

这个“谜”在昨天的思考里完全解开了;我不必去找,甚至不必刻意去想。写完之后重新回味,种种联
想浮现出来——写作之时它们必定已经在场,只是我毫无察觉,还在暗中引导我的笔。打动我的,不
仅是风格的相似,更是一整套\textbf{专利的据为己有之术},贯穿\emph{葬礼}的三大“运算”(四者之
中,Deligne本人是主要的——即便不是唯一的——“受益人”)。不妨称之为“临时替身爸爸”的手法:
悄悄把它摆上数学勒索的棋盘,好变走一份真实的父亲身份,而我朋友皮埃尔其人,则暂时隐在幕
后。一旦生身父亲被体体面面地彻底清出舞台,替身爸爸也就被变走,仿佛从来不曾存在;于是\textbf{真
正的父亲},谦和地微笑着登台,甚至无须亲口说破——因为那个不动声色牵线、又懂得等待的人,
诸事自会水到渠成,毫无阻拦:教团全体的同心一致,早已把他按他所应得的名分加了冕。

这套手法,直到前几天,我循着我朋友Zoghman在第四场运算(所谓“当值无名之辈”)各插曲中的
倒霉遭遇一路追索,才开始看清。这一场里的“替身爸爸”(为某件“对应”……)是\textbf{Kashiwara}——我
说不好,他是就这么从天上掉下来的,\dmark{}\origpage{1172}似有天助且纯属巧合,还是未来的真爸爸巧妙地让他明白:
一个无名之辈的那个没有像样父亲的结果,横在这里,弃之可惜\fnote{936}{(*)}{(6月16日)看来,对Mebkhout著作施行扒手手术的主意,正是企业家Kashiwara出的,而且早在1978年——其时Mebkhout刚把刚杀青的论文第三章拿给他看不过几个月。关于这一点,见注“黑手党”部分(b)(“最初的麻烦——或太平洋彼岸的黑帮老大”),打字稿第1060页脚注(*)。}\fnote{937}{(*)}{(6月16日)其实,两年前他就开始露头了——见前一条脚注。1980年3月的那一幕,是Goulaouic--Schwartz讨论班;已引之注以及“掠夺的空白授权——或‘高尚行当’”一注(n\textsuperscript{o}~171$_{4}$,尤其打字稿第1088--1090页)都谈到它。}。反正,皮埃尔朋友深谙
此道,在一个人为的、由“行家们”斩钉截铁的共识拼凑出来的父亲身份的暧昧上,把戏做得
滴水不漏——而这早在那件新东西的分量得到公认之前就动工了。替身爸爸Kashiwara,1980年3月即已
登场,若不是早在六个月前的莱苏什讨论会上就已露头;十五个月后,在1981年6月那场难忘的
讨论会上,他被变走,不留痕迹(而且看来他自己也不太计较)。在这里,变走的手法炉火纯青:再引
进两位——这次不妨叫作“名义上的共同爸爸”——Beilinson与Bernstein;他们登台,纯粹是一道讲究
格式的保险条款——“喊停!”——因为当然,谁也不会想象孩子是他们俩谁做的(虽说两位都沾了
光……)。

与“Motif 运算”的类比,着实惊人!当人们还能拿出来的、关于 Motif 的“外行废话”的父亲身份,
知名度还太高(尤其七十年代初),无从下手时,Motif 瑜伽有两个\textbf{关键侧面},还不曾有一行
字公开发表过,哪怕以影射的方式。其一,“权的瑜伽”,早在1970年就被\textbf{超级爸爸}据为己有,
纹丝不动——反正被变走的那部分,不过是“猜想”,跟一句纯形式的影射差不多。另一侧面却毫无猜
测成分,六十年代下半叶就已打磨停当。按安排,由一个有点跟不上趟的无名学生,来陈述这套瑜伽
的启动机制——技术上不难,但(至少直到那位不受欢迎的生身爸爸“死亡”前后)这件活儿看来超
出了那可怜人的能力。

\pmark\origpage{1173}正是这位学生,Saavedra,是现成的替身爸爸:有幕后人暂时的担保撑腰,足以让一个一心
只想忘掉该忘之人的教团点头称是;但与此同时(妙就妙在这里),这位“爸爸”分明“没有分量”。
时候一到,谁也不会——Saavedra想必比任何人都更不会——起念头,宣称他可能是某种新哲学的父亲;只要
肯稍微想一想,这种假设简直滑稽……这里,功成身退的替身爸爸被请出去,要晚到十年:1982年,
难忘的 LN~900 出版。应当说,从1972年(第一场所谓“Motif 运算”中抬出“替身爸爸”)到1980年
(第四场“当值无名之辈”运算中同样应运而生的替身爸爸登场),多少水从桥下流过,再用不着绕
弯子了!耐人寻味的是,这里又插进来一位“名义上、装装门面的共同爸爸”,好让替代父亲身份(说
到底,就是木偶戏里的父亲身份……)与\textbf{真正的}父亲身份之间“平稳”过渡(而且没人显得在
抢镜头)。我确信,Milne对那双按他人意志摆布他的无形之手,看得不会比Beilinson与Bernstein更清
楚。人人都分到了碎屑;人人(至少是有发言权的人……)都有一切理由心满意足
\fnote{938}{(*)}{(6月16日)在他与Saavedra、与一位(暂时)在幕后的Deligne之间必然达成的“交易”里——这位Deligne,只等时候成熟便重新露面……——Saavedra的“份额”是:一张国家博士学位文凭,以及作为“数学讲义”这一显赫丛书作者而获得的相对知名度——这足以为他启程回国、远离他只远远瞥见过一眼的枯燥数学追求,铺好道路……}。

\pmark 这一切,昨夜还让我想起那第三场大运算——那场直接惠及“未来无所不在的爸爸”的运算:
“平展上同调”。此前我已使自己相信:这场运算的初始动机\fnote{939}{(*)}{关于这一点,见“那个公式”那组注(n\textsuperscript{o}s~169$_5$--169$_9$)。这最初的用意,后来大大扩展开来——尤其见“悼词(一)——或恭维”(n\textsuperscript{o}~104)与“宝石们”(n\textsuperscript{o}~170(iii))两条注。},是据为己有某一\textbf{不动点公式}的
打算:人们可以把某个父亲身份不受欢迎的“L 函数公式”呈现为该公式的一个平凡推论。麻烦在于,
所论的迹\dmark{}\origpage{1174}公式偏偏带着\textbf{同样}不受欢迎的父亲身份。幸好,还有另一位可能的父亲——一位好朋友(不点名,Verdier),他做
出了两份公式:一份太一般(但启发式地至关重要),另一份略显局促,却足以“罩住”人们想要的东
西。可朋友归朋友,这里的“父亲”既不能是朋友,也不能是那个不受欢迎的死者——因为事关“那
个”著名猜想的那把钥匙公式\fnote{940}{(*)}{当然是“那个”Weil猜想。关于这一点,见注“那个猜想”(n\textsuperscript{o}~169$_4$)。}。糟糕的是,L 函数公式与它那不合时宜的父亲身份,偏偏尽人皆知;所以
这里的难点不在朋友(朋友之间,最后总能商量……),而在死者。更倒霉的是,他对“推论”的证明,
1964年就白纸黑字发表在一场Bourbaki讨论班上——只是(万幸!)该公式的那桩例行情形(呃,
对不起,我是说关键情形!——这公式或迹公式,反正一码事,但这话千万说不得……\fnote{941}{(*)}{这两条公式,实际上互为直接推论。既然我对其中之一(L 函数公式)的父亲身份尽人皆知,Deligne便设法(在那份名叫“SGA~4½”的难忘文本里)把它呈现为另一条的推论,并且使尽浑身解数,以远比我这微薄的证明(以及随附的命题)艰难得多的魔术-行骗手法,把自己装扮成后者的父亲。这套戏法见已引的那组公式注(n\textsuperscript{o}s~169$_5$--169$_9$)——在我们这门可敬科学的年鉴里,这大概是绝无仅有的一场大变活人。}),当时还来不及被验证。

这里的手法,是拿那位朋友充当他的超一般公式的父亲(这原是分毫不差的事实——只除了一点:
他从来没费神去证明它……),同时悄悄把它与那累赘的死者所\textbf{显式}证明的公式(对那一条,自
然讳莫如深,一字不提)混作一团,再把那条超一般公式\textbf{拆毁}(说它是猜想的、不完全的、说
白了,没用的)。这叫溺毙一条鱼:叫读者压根不去想,翻一翻某套 SGA~5 讨论班(人们本来就责成
他“忘掉”它)会就这个问题说些什么。至于朋友那条(边上略局促、却完全有效的)显式公式,大家
一致同意也不再提了,只在一个文本的末尾留下一处暧昧而装门面的引用——那文本拖沓冗长、令
人生畏,世上没有哪个读者有勇气读到底。总而言之,可以说:“替身爸爸”(这里是Verdier)确实出
了场,但与其说是凭他对一项结果的默许的“父亲身份”——这里要\textbf{彻底变走}的,恰恰是\dmark{}\origpage{1175}死者的那
项结果——不如说是凭他的默许,参与了一场搅浑水、拆两条“孩子”的把戏:那两个孩子
确凿是他的,搅浑水是为了在混战中变走第三个孩子——那孩子的父亲不可宣之于口;如今是一个
孤儿,谁也找不回他,谁也不上心去找他\fnote{942}{(*)}{关于这一点,见注“变戏法者——或飞走的公式”(n\textsuperscript{o}~169$_8$)——以及“家庭相册”一注打字稿第1121页的脚注(**):它表明Deligne与Illusie这对好撒玛利亚人的变走-蒸发之计,是如何大功告成的。}。在这套手法里,Illusie扮演一个辅助角色,同方才
的“名义共同爸爸”相仿——只是他的父亲身份,同Verdier的一样,任何时候都不许落在那神圣的
\textbf{对 Frobenius} 的迹公式头上——那是唯一算数的,理当(当然,以应有的细腻)留给Deligne一
人;它落在那个要被变走的、不可宣之于口的孩子身上——Illusie以他特有的、堪称楷模的虔诚,参
与了这件事。

\subsubsection{e.\ Monsieur~Verdoux——或献殷勤的骑士}

\notehead{176$_{5}$}{}

可我还想回到Saavedra的“论文”上来。大约在我告别数学舞台之际——1970年初(如果我记得
不错)——Saavedra才装作终于真正“上钩”了:此前一两年,他一直显得三心二意。他对
我说,他已提炼出我最初给他出的那个命题的一种表述与证明,使它适用于基底环 $k$ 为\textbf{任意}环
的情形。他甚至给我画了一张证明的草图,我只半心半意地听了听——我那时几乎全部精力,都被我
正经历的人生转折占着。没有细核Saavedra的话,我有一种印象:他总算起步了,往后可以自食其
力。也许,在一个我对真正的科研指导已经几乎无暇顾及的时刻,我有些性急,把愿望当成了现实
\fnote{943}{(*)}{至少,同我离去之前的投入相比;但并不比我后来在大多数同事——那些担任科研指导的同事——身上看到的更少。}。此后,据我记忆,再没有他的音讯\fnote{944}{(*)}{我的记忆在此稍稍背叛了我——见注 n\textsuperscript{o}~176$_{7}$,其中对此有出乎意料的揭示。}。直到上星期,我还一直以为:我给他出的那个最低纲领,他总该完
成了,也许还稍微\dmark{}\origpage{1176}多做了些——把 Motif 的情形也处理了(照Deligne去年8月随其 Motif 带评注书目
写给我的信所言)。

\textbf{我这才刚刚发现:压根没有这回事}。这可怜人花了三四年功夫,竟在他称之为“塔纳卡范畴”
的东西的\textbf{定义}本身里(我指的是那个内蕴性质的定义),犯了一个粗糙已极的错误\fnote{945}{(*)}{错误的来由,是Saavedra心里对我所说的张量范畴的“基底环”有混淆:它绝不是使该范畴“线性”、使张量积“双线性”的随便什么环,而是典范的环 $\mathrm{End}(1)$(1 为范畴的单位对象)。我给Saavedra讲理论入门常识的时候,他想必“不在局中”到这般地步,这些话全从他头顶掠过,沉入遗忘。Deligne——他似乎多少接了我的班,当上了Saavedra的指导(显然另有一本账……)——却绝不肯替他纠正航向。这让他得以(十年之后)悄悄拆掉Saavedra的纸牌楼,以救世天使、以及(这一回又是)人人翘首以待的真爸爸的面目出现……}:而这个定义,正是要用来证明
它蕴含“Galois式”的、以一个合宜 gerbe 的表示写出的描述。他在引言里陈述的定理三(那篇引言,
按理至少应当\textbf{陈述}我交给他的那理论的四条根本定理)因此\textbf{平凡地为假}。Deligne与Milne以指出这个显赫的错误为乐事,提议把所论范畴的“新”定义取作 gerbe 式的描述(先天地明摆
着,这才是对的,必要时改一改内蕴描述就是……),并一本正经地追问:Saavedra的“定义”(把那
愚蠢的错误剔除之后)是否蕴含“他们的”(原文如此)\fnote{946}{(*)}{前引书第160页(我可不编瞎话!)。}——而这恰恰原本就是预定要构成Saavedra博士论文的那桩题目!

这一处境,是十足的于布爸爸的翻版!而且一上来就是三十六个花样。譬如,交给Saavedra的那
桩题目——唯一需要一点原创贡献的部分,哪怕多么微薄(为一个尽可能一般的基底环上的Galois-Poincaré范畴,提炼好的内蕴条件)——即使在(我相信我在遇到Saavedra之前很久就已处理过
\fnote{947}{(*)}{那是1964或1965年,即Saavedra那篇出名的“论文”(原文如此)之前七八年;比Deligne--Milne二人组赶来救驾——也\textbf{不做}这份微薄的活儿——还早十七八年:那是我对最微薄的学生所期待的唯一“原创”工作……})基底环\dmark{}\origpage{1177} $k=\mathrm{End}(1)$ 为\textbf{域}的情形
——也根本没有被处理。Saavedra的“论文”因此精确地在于:虔诚地照抄理论中(Grothendieck式瑜伽的启动之外)早已由我全部做完的那一部分——基底域上的那部分——并且,以一个
骗人的定义和一个假定理的“证明”顶替那先于一切的必做之事;那“证明”(Deligne以指出这一点为
己任——前引书第160页)归结起来,不过是一个简单的循环论证!

还不止这些。论文站不住——而答辩委员会毫无察觉!想必没有一个委员真正弄懂了所为何事。可
这并没有促使他们中任何一位向我示意:至少他们中该有一位,有能力为那项他们一本正经摆出架
势来评审的工作的严肃性,提供一份有效的担保\fnote{948}{(*)}{这支可悲的委员会的组成,终究要(在坚持读到那里的读者面前)在“第六根钉子”的终极之注 176$_{7}$ 里揭晓——在我棺材上的第六根钉子……}。答辩居然举行了,而且没有让我参与——这只能拜Deligne的担保所赐;
(如Saavedra的致谢所暗示的)在我实际上从舞台上消失之后,Deligne必定稍稍跟进过他的工作
\fnote{949}{(*)}{Deligne对一个默默无闻、想做论文的学生的这份突如其来的关切,是在该理论的自然父亲(且不受欢迎……)死后才出现的——何苦来哉;而那位学生(显然被任务压垮了……)本该陈述的,正是这门理论。}。

再说,打那时起,Deligne居然没有察觉这个错误,我简直无法想象——我深知他的敏锐与犀利,直至
最小的细节——而这绝非什么“小细节”!当然,我曾把已达成的瑜伽的全部精妙之处讲给他听;我
给他解释的最初几桩里,若不包含这个反例——他同Milne如今佯装当作最新发现拿出来的那个反
例,它早在我思考瑜伽的最初就为我所知(我最终打算把这瑜伽就叫“Grothendieck式”,不必去麻烦
Galois与Poincaré的大名……)——那是根本不可能的。他若听任“门生”(原文如此)的“论文”(再如
此)里留着这样一个粗糙的错误——只等时机一到,便可以纯然干脆地搞垮那位(暂时的)“替身爸
爸”——那想必不是没有充分理由的。昨天的思考,把那些理由已经摆得再明白不过。

\pmark\origpage{1178}也许有人会说,我在编故事:Saavedra所说的“帮助与建议”,并不意味着Deligne当真费神
稍稍认真读过引言里概括理论之实质的那四条命题\fnote{950}{(*)}{纤维函子上的滤套的那些结果除外——它们更技术性,更难压缩进一条醒目的命题。}。这些命题,他早在认识当事人之前就熟悉了;那么,那至多是一时轻率:为一项工作作
保,却连一刻钟的功夫都不肯花,验一验引言里预告的主要命题是否成立。可在我心里毫无疑义:Deligne确实费了这个神。因为这项工作\textbf{不是随便什么工作}——它出自一个多少有点找不着北、正
愁没论文的学生之手。继我之后(而且还在Serre之前),Deligne是最有资格感受那份形式主义的全部
分量的人:那是已故宗师留下的、未写出(至少未发表)的遗产的一个关键侧面。他尽管对这一侧面
摆出他惯常的满不在乎的架势\fnote{951}{(*)}{关于这些架势,以及它们所伺候的据为己有之术,见注“据为己有与轻蔑”(n\textsuperscript{o}~59$'$)。},骨子里他比谁都清楚所为何事。他,光辉的Deligne,极端的精英主义
者,居然肯费神跟进一个分明天资平庸的人的工作——这断然不是为了当事人漂亮的眼睛、为了帮
他拿到(按时下的共识,更按他所自诩奉行的、推到极致的严格标准)一纸\textbf{冒牌文凭}。

这话一出口,人们立刻面对一个奇怪的矛盾。一方面,一个据称多年全职投入本题目的人,犯了如
此显赫的错误,很难不把它解读为根本上无能的征兆——那个摆在他面前的问题,哪怕单看其技术
面(其实并不深奥),在答辩之时、在书出版之时,似乎根本没有被把握。另一方面,同一个学生,跟
我一两年没做出什么名堂之后,不到两年,忽然习得一门堪称令人叹服的数学素养:代数群结构理论,
一般域与\dmark{}\origpage{1179}实域上的、硬邦邦的概形理论、Hodge 理论、Motif……不止如此——我不记得读过任何
出自他手笔的数学文字,哪怕几页;而且我深知(对手段微薄的学生尤其如此)学会给数学
“成文”,绝非易事——可翻阅那本以他之名问世的书,它的“成色”之精良令我震动。我曾想,单就技
术而论,这份显然以与 EGA、SGA 同等的标准参考书自任的文本,本可以出自我的手笔,或Deligne
的手笔,或我那四五个才具出众、娴熟于把一套交织复杂的思想与事实精确、完整而优雅地陈述出
来的学生中任何一位的手笔。我很清楚,这样的行文功力,比数学素养更甚,绝非可以临时抱佛脚(除
非是天纵之才,如那位Deligne,及寥寥数人);它(倘若终究能获得)只有经年累月的练笔方能获得。
我自己就花了十年,而那时我同所要表达的实质的接触是极强的。这份接触,同Saavedra对他论
文题目的接触,当然不可同日而语——他写了关于这个题目的文字,却始终没有懂它;而这文字,后
来(至少到1982年……)竟然\textbf{是}“那一份”好的参考文献,为那套精微而关键的形式主义。分明
是两件干脆“对不上榫”的事……

昨夜起就掠过我心头的一个念头,如今,当我把处境白纸黑字对自己讲述一遍之后,带着明摆着的
力量回来了:Saavedra——我深知他的可能,尤其他的局限——不可能是那本光辉著作的作者;这
书以纯然技术的面貌(诚然,在这方面却又是穷尽的、针脚细密的),陈述一门完全超出他的“哲学”
的基础:也许头三章——其中两章主要是人人早已熟知的函子式泛论,第三章呈现的是书中中心观
念的、Saavedra版的十足的骗局——这几章,本该构成他从未完成的“最低纲领”;也许这几章全
然出自Saavedra的手笔。第三章尽管是十足的骗局,却仍足以\dmark{}\origpage{1180}让人窥见要到哪里去——即某些 $\otimes$-范畴的“Grothendieck式”眼光(姑且不点名),或“gerbe 式”眼光;它给后面的第四至
六章赋予意义。一旦接受了 gerbe 式的描述(在Deligne与Milne那双料海盗式的文本里,它被明智地
径直取作那些所谓“塔纳卡”范畴的\textbf{定义}),后三章便是要被据为己有的那套形式主义的
心脏。我猜,这几章是Deligne统统亲笔写的,或部分由他、部分由Berthelot写的;而且比我交给
Saavedra的笔记详尽得多,以致他几乎只消逐字照抄——如果曾有人要他费这番形式主义的神
的话。他准觉得“赢了”:人家“送”他一篇论文,外带一枚文凭——虽然他心里分明感到,他自己做的
那点东西(甚至是在自欺欺人地以为它站得住的情形下),对一篇国家博士论文来说,恐怕太单薄
了。而Deligne(又化装成撒玛利亚人……)赢了:这便是所需要的参考文献——纵然不为眼下,也为
“日后”(对会等待的人……),而且在那份参考文献里,那不受欢迎的名字,至少在一切实际的用处
上,不再出现。

为让欢乐登峰造极,我再补一句:那位名叫Saavedra的,似乎已从流通中消失,再无踪迹。去年,预
备寄发(我眼看就要寄出)《收获与播种》的印装本之时,我曾翻检《国际数学家人名录》——那册
子厚得很,“人人”都在(名录就是干这个的);唯独当事人不在:Saavedra之下没有,里瓦诺之下没
有,连Neantro之名下也没有(我凭良心翻过)。于是,这段历史有了黑色侦探小说的模样。想到那位
满面春风、和蔼可亲的Deligne——俨然又一位Monsieur~Verdoux(别名朗吕),以那本称心如意的“恰好的参考文献”遂
了他那番折磨人的心愿之后(比他朋友Verdier的还早四年!\fnote{952}{(*)}{关于它,见那道(理应如此)取名的注“好的参考文献”,n\textsuperscript{o}~82。})——人们不寒而栗:眼看着他让那桩恶魔般的阴谋的“证物”消失:那可怜的Neantro~Saavedra~Rivano本人,被放进奥尔马伊(Les~Ormails)一座精致的小壁炉里,长时间地焚化——那壁炉,
本就是特地为这类用途设计的\fnote{953}{(*)}{“奥尔马伊”(Les~Ormails)是法国高等科学研究所(IHÉS)住宅区的名字;我朋友皮埃尔——别名Monsieur~Verdoux,别名朗吕(还化装成献殷勤的骑士)——在那里及时接班,顶替了某位已故者:那人被排挤出这块地盘,又被我朋友所擅长的那种不动声色的暗算,抛进了暗无天日的地牢。住宅区有十来幢家庭小楼,和一栋大些的房子,带舒适的单人工作室;那些工作室想必不久也会各自配上一座小巧的全能壁炉……}。

\pmark\origpage{1181}我宽慰自己:没听说Kashiwara或Verdier也从这个世界消失——说真的,就在前天,我还拨通了后
者的电话,问他(没什么信心,似乎也没成功)能不能给我带来另一位“失踪者”的消息:人人谈论、
却看来无一人见过的——Jouanolou的论文。那篇论文,我仍然一无所知;但至少,Verdier似乎
还活着,尽管他本人就是一份“证物”——而且我怀着良好的希望:Neantro~Saavedra~Rivano也
一样。

\subsubsection{f.\ 肮脏的差事}

\notehead{176$_{6}$}{}

话虽如此,Saavedra论文史的那些于布式面向,我甚至还没有巡视完——看来那些不同寻常的论文与论文人,我是
收集上了瘾!到这里,我已有了揣测(不说“内心确信”吧):Deligne(由一位殷勤而义务的助手协助)
之所以装模作样,在Saavedra答辩十年之后一本正经“誊抄”其论文,想必只是“取回”他当年乐意
“借”给Saavedra一段时间的东西(这段时间,恰够Saavedra把论文过了、然后消失);所以这
不过是天理昭昭——只除了一点:他“借”来一段时间的,是从那位从不被点名的死者那里“借”来的。
可是,借死者的东西不还(那还了得!),风纪倒也无亏:这一头,也是万事大吉。

最妙的还在后头:即便在第二位前学生也过了这道手之后(而且还是我有过的学生中最光
辉的一位),我交给Saavedra的那个谦卑(微末)的问题——那曾是二十多年前我的出发点,
也是我自认当时就解决了的第一桩事,在所论 $\otimes$-范畴的定义环是域的情形——这个谦卑(微末)的问题,
直到今天,即便在那个情形里,居然仍然没有“被解决”!Deligne满足于点破Saavedra那个粗糙的错
误(那错误肯定十多年前就被看破了,他只是等待时机……)。他一边在128页上誊抄先前那本参考文
献\dmark{}\origpage{1182},却无心去修复这个错误。何必费这个神呢——追求的目的分明已经达到!要做到那一步,他
心里就必须在这场行动之中,除了据为己有的饥渴之外,另有别的东西:一份醒着的关切,一份对他
所处理的数学实质的\textbf{敬意},一种超出眼前“收益”的视野。

如果我在六四、六五年前后,曾费神为那些可以用“代数 gerbe”表述的 $\otimes$-范畴提炼一套“Grothendieck式”的瑜伽,而不满足于可用群概形描述的那些,那是因为在“打动”我最深的那个例子——一个域
上的 Motif——里,众所周知(凭Serre一个非常简单的论证):当该域特征 $p>0$ 时,不存在“有理
于 $\mathbb{Q}$”(甚至有理于 $\mathbb{R}$)的纤维函子。这在当时\textbf{逼着我的手},把理论表达在像 gerbe
与 lien 的形式主义这样“不正经”的东西里;同时当然,也要找到简单代数性质的内蕴判据,保证这种
“Galois式”或“Grothendieck式”的视野几乎“总是”行得通,而且花费极小。我所提炼的判据(如果
我没弄错,还给出了证明)——由 $k$ 的一个扩域 $k'$ 上纤维函子的存在给出——二十年后,文献里仍
然没有建立!直到今天,照Saavedra、Deligne之流写下的文字,纵然人们对(比方说)有限域上
“Motif 上同调类”的形式主义承认我们想承认的一切,半单 Motif 范畴(比方说)是“Grothendieck式”
的(或照这些先生的叫法,“塔纳卡的”)这一点,仍然没有建立(至少文献里没有)。前后共 418+128
=546 页的文字,出自Saavedra(由一位Deligne与一位Berthelot协助)继而Deligne与Milne之笔——
到头来却连二十年前我的出发点都没有抵达:那出发点,是让我自己相信,“Motif Galois群”\textbf{确
实存在}。

是啊,一个Deligne何必费这个神呢——他早已忘了那视野;他要的声望反正已到手;他用来做
\textbf{他的} Motif 理论(它当然同某位死者的理论毫无干系……)的那些域,全是特征零的域——于
是他那些所谓的“塔纳卡”范畴,统统是“中性”的(或“平凡”的)。照这个账,自然犯不上为 gerbe 之类
大做文章\dmark{}\origpage{1183}——那从此不过是障眼的烟幕。犯不上,除非是为了\textbf{据为己有一个}
\textbf{连灵魂与精神都已被遗忘的东西的字面}。而这场令人瞠目、可悲的故事的收梢,我看到,是这
样的:正如那套被埋葬了十五年的 Motif 视野的入门常识,还得由那个老朽来做——他刚巡礼完那
场光辉的葬礼及其绝技,就得自己把这份小活儿干了:他的学生们在他“死后”,还没有谁把它放在
心上。因为他们早就忙于扮演主子,忙得没有几天工夫,哪怕片刻,再当一回\textbf{仆人}
\fnote{954}{(*)}{我在这里有点性急,装作把我所有学生同他们当中最光辉的一位装进同一个口袋。我先向他们中那些不乐意与如此光辉的人士为伍的诸位致歉!无论如何,我高兴地记得Giraud:他硬着头皮接下了(不期而至的)读Contou-Carrère论文的活儿;那心境,对Contou-Carrère、对我,无疑是一种“服务”的心境,也许对数学共同体也是如此;关于这一点,见“耶稣与十二使徒”一注(n\textsuperscript{o}~19)最末一段,打字稿第151页。}。

\subsubsection{g.\ 一场屠杀中的五篇论文——或孝道}

\notehead{176$_{7}$}{}

(6月19日)到今天恰好两个月,我乘着一阵旋风,动手写前面那些注(4月19、20日的),带着那
个现成的名字“第六根钉子(钉进棺材)”(n\textsuperscript{o}s~176$_{1}$至176$_{6}$,不计本注——它也在其列)。此前一星期,Zoghman~Mebkhout刚把Saavedra那本书带给我——只消看上一眼,我就已经明白是怎么回事了。

我必须承认,这一发现带来的震动,几乎不亚于一年前(整整一年前)那本 Motif 掘墓“难忘文集”
(《数学讲义》第900种)带来的震动。更确切地说:去年的那份震动,某种意义上重新出现了,被这
个与那场掘墓紧密相联的“运算”的发现猝然重新点着;这场运算(一开始就明摆着)曾为它作了准
备,规模也相仿。我再一次被那种安详的无耻之感攫住(简直要说被噎得透不过气)——\textbf{同样的}无耻(这一点,凭种种不会
骗人的迹象,一开始也明摆着)——冲着一件与我血肉相系的东西,一件世上除我之外无人长久怀抱
与滋养的东西……那感受如此强烈,近乎焦虑,连我自己都吃惊。

\pmark\origpage{1184}自发的反应、自然的出口,本该像去年那样:趁感受还新鲜就说出来,由此径直进入这新的侧
面——那场由我的亲近者对我施行的活埋的新侧面。可我忍住了\fnote{955}{(*)}{我到底还是乘着当时的激动写了四五页,但九天之后(4月19日)写下正文时,那些痕迹已所剩无几。};因为Mebkhout的来访,需要我至少有一些闲暇;何况他
有些话要对我说——即便不那么触动我的要害,我感觉得到,它们对他而言同样是“要害”,对\emph{葬
礼}同样意味深长。再说,趁新知还热,把它们记下来,在我看来也重要;而那本出名的“葬礼之书”的
来龙去脉,即便我日后再来清理,也不会溜掉。所以,朋友刚走,我(4月15日至18日)就去写他的倒
霉遭遇,成如今“封神”末尾的那组注(n\textsuperscript{o}s~171$_{1}$至171$_{4}$)。

也就是说,在来到那出名的“第六根钉子”之前,我已经有时间恢复镇定。说实话,刚才重读最初几
页,在我对那桩新揭发的丑闻所作的挖苦(而稍带距离)的描述里,我找不出最初袭我的心绪的痕迹
——那心绪曾让我失眠一夜,而我当时多么需要睡眠。这一回,我是真切地感到了那份“过去的重量”!

那是6月10日,在给\emph{葬礼}画上那个出名的“句号”之后三天——葬礼随之重新开动!当然,我远
没有料到它开动到了什么地步——还有(差不多)三百页要写!写完“第六根钉子”的第六道注(“肮
脏的差事”)时,我满以为连“四种运算”也一并走完了——只除十来页(第三、四场运算的)要重抄
清稿、补上预定的脚注。再过几天,我想就可以把葬礼(III)的全部手稿送去打字了。

然而,随后几天里(也许就在我以为钉完最后一根“钉子”的次日或隔日),发生了一场出人意料的
剧情突变,对此我还得回头再谈。这里,我自发的冲动本该是立刻扑上去。我又等了两个月才动手,
当然不是不情愿\dmark{}\origpage{1185}——是有更急的事要备好送去打字。从头重读《四种运算》,发觉这里那里大有充实的必
要——后事如何,大家都知道了!

不管怎样,今天(只要不再有新的意外——摸摸木头!),终于到了那个黄道吉日:我给\emph{葬礼}
画上真正的句号——照实说了:写下最后几页的日子;它们理应构成我在《收获与播种》之中关于
葬礼的思考的尾声。此后,只还剩写那封充当《收获与播种》前言的“信”——然后我想休息几天,
受之无愧,而且也实在需要……

写完前面六道注几天之后,我得知了Saavedra答辩委员会的组成——那支我在“Monsieur~Verdoux——
或献殷勤的骑士”倒数第二道注里痛加讥讽(罪有应得)的委员会。论文于1972年2月25日在奥赛的
理学院答辩,委员会由\textbf{Demazure}(J.~Demazure,报告人)、\textbf{卡斯泰勒}(Castelle)与
\textbf{Grothendieck}组成。

要说“剧情突变”,这才是剧情突变!于布戏的加冕礼!而且,面对这份官方来源的信息,我好不容
易才相信——我连出席过这样一场答辩的一丝记忆的影子都没留下。Verdoux-朗吕先生的故事,愈发精彩了!我姑且给
Demazure打了个电话,问他可记得同我一起,给一个叫Saavedra的人当过论文答辩委员会。Demazure也记不得多少了,但好歹足以向我保证:答辩确凿举行过(何时如何,他自然说不太清),我们
两位都在场,还有卡斯泰勒(连这名字我都不记得了……)。他知道的也不多,只记得自己是论文报
告人。还是我告诉他:论文官方说来,是一份25页的文本(想必这让他的报告人差事轻松了不少)。
这一来,吃惊的反倒是他。他答应想法给我寄一份论文来。我本来很想知道它长什么样,可至今没有
等到——看来(照Demazure几个星期后终于对我说的)这份论文已无处可寻;也许他也没怎么费神
去找。反正,同我一样,他的文件里也毫无踪迹。不过,这只是个细节……

\pmark\origpage{1186}这一来,我可够瞧的了!我对那支分明无能的委员会——“一本正经地摆出架势评审一项他
们‘未必弄懂了所为何事’的工作”——所发的那些冷嘲热讽,如今都泼回我自己身上!可以想见,我多么想收
回那些讥讽,保住体面,撑住场面——可是不,那是作假。整场\emph{葬礼}里的作假已经够多了,犯不着
再添上我的一份。再说一遍:那些讥讽完全是有道理的。如今知道了委员会的组成,我甚至可以指明:
正是我,先于任何人,最应当受这些讥讽。毕竟,Demazure与卡斯泰勒当时记住的,主要是:这篇论文,
Saavedra是跟我准备的,或至少是跟我开的头,题目是我给的。按说“在局中”的是我,而他们信任
我。说不定,那出名的25页——Demazure据称做过报告人的——本来是站得住的;即便那桩显
赫的蠢事也在其中,它也不过是一份理论概述,不在局中、信任我的Demazure,绝无机会察觉。

至于我——除了上课,两年来实际已脱离数学——这场答辩,我草草了结,想必不过是在奥赛的两堂
课之间、某场“生存与生活”的集会、或某场关于附近(萨克莱)核废料的公开讨论之间,纯然一道
行政手续。确定无疑的是:两年来,Saavedra的工作我早已没有跟进,任何人的也没有;而且我对
Saavedra的工作站得住,毫不怀疑。这份确信从何而来,如今说不清了。同此前我所有的学生不同,
对Saavedra,我没有任何直接的凭据——凭他跟我一起已完成的工作——来判定他的严肃。难道那
些年,我竟把我的学术职责轻率到这般地步,凭他的面色,就信任了他?倘若那本书(同年出版,那
25页的论文想必是它的摘要)当时已经杀青,并让我形成了印象,那“看一眼”确实印象极佳,说
不定我连想都没想过要去核实:那本该构成Saavedra个人贡献的部分。也可能,甚至更可能(此事
我毫无记忆了),我信赖了Deligne的判断——我离去之后,是他跟进了这项工作
\fnote{956}{(*)}{我连Deligne曾跟进Saavedra的工作这件事,都没有记忆。这是4月里,看Saavedra那本书的引言时才知道的。}。

\pmark\origpage{1187}无论哪种情形,我都必须承认:我负有同样的责任——我把科学博士的头衔,授予了一篇二十
三年后显出是\textbf{冒牌货}的论文(借已引之注的说法)。然而,我自己在不知不觉中成了这场骗局的
工具、并因轻率作保而负有责任,这一事实,并不改变骗局之为骗局。它反倒显得愈发天才。因为说到
底,牵线者的真正动机,当然不是让一个落魄的论文新手廉价地得一顶头衔,然后改行退场——而是
让一个绝不糊涂的人,细腻地、不动声色地,把某种在我之内诞生、并且在他还没有听到(在数学
里)说出“gerbe”或“Motif”这类字眼之前,就被我孕育至成熟的那份视野的父亲身份,据为己有。正是趁
我骤然全力投入物种生存与其他种种最紧迫的美好事业之时(那位前学生兼朋友曾对我说,他不得
不同这些保持距离,因为他全身心绝对奉献给唯一的数学
\fnote{957}{(*)}{关于这一点,见注“兄弟与配偶——或双重签名”(n\textsuperscript{o}~134),尤其打字稿第614--615页。})——趁我的精力完全被别处占尽的时刻,我那位光辉的学生兼朋友,完成了这桩
真正独一无二的障眼大戏:让我成了我自己被褫夺的工具!我当时的心境,完全切断了对昔日数学兴
趣的牵挂,对我那些学生——Deligne首当其冲,他们自 SGA~5 讨论班结束起,就已开始按自己的路
数玩起了小把戏——报以盲目的信任:随便什么名字,只要安到那些我只剩遥远记忆的出名范畴上,
我都会说:是,阿门!正如我对Verdier宣布不会再有新式同调代数的书时说了是、阿门;正如Deligne向我宣布,我们合搞的 SGA~7 讨论班的一半,将突然更换父亲身份时,我说了是、阿门……然而,
受骗者乖乖地点头、毫无所觉地同意,并不改变诈骗的性质——只给它添上一重背信。而当Serre之
流的其他占卜师也从中各得其所、毫无保留地祝福其事
\fnote{958}{(*)}{这份再明白不过的祝福,见“家庭相册”一注部分 d.(“入殓——或自然的坡道”)。},这件事便有了非同寻常的维度\dmark{}\origpage{1188}——它成了一个圈子、一个时代的集体败坏——但它并不因此就光彩起来,无论多么天才;它的无耻,一分也不少。

正如剧情突变式的好事从不单行:得知我前学生Saavedra答辩委员会的组成才几天,关于Jouanolou论文的相应信息也到了——那也是一篇有点特别的论文,我在思考中这里那里谈到过它
\fnote{959}{(*)}{这篇论文,在注“团结”(n\textsuperscript{o}~85$^{1}$,打字稿第349页)的子注,以及“共同继承人……”(n\textsuperscript{o}~91,第387--388页)一注里谈到过。另见“学生与纲领”一节(n\textsuperscript{o}~25)。}。同Saavedra一样,他也从未想过给我寄一份他那出名的论文(“人人都引(自反常层会议以来),却无人见过”)——于是我最终给他写了封略嫌干巴的信(4月25日),就这
篇论文的离奇遭遇问了他几个问题。他几乎当即回了信(5月1日):对实质性的问题闪烁其词(因为
“重提往事总归非常痛苦”),但行政细节却精确之至:论文于1969年7月3日在法国高等科学研究
所(巴黎)答辩,委员会主席为\textbf{P.~Samuel},考试官为\textbf{迪斯米耶}(J.~Dixmier)、
\textbf{Grothendieck}与\textbf{Verdier}。我的通信人还添了一句,不无恶意:“据我所知,委员会全
体成员都到了!”(此事 J.L.~Verdier不久后也在电话里向我证实。)

\pmark 又是如此:对这场答辩,我连一丝记忆的影子都没有——它分明也是“偷偷”过的(不得不这
样糟蹋我自己的形象,抱歉!)\fnote{960}{(*)}{答辩的日期,我记得,应在我已经“脱离”数学、转而关心生物学(尤其分子生物学)之后不久。}。我以为答辩是在斯特拉斯堡(因而把它放在七十年代初——知道Jouanolou那几年在斯特
拉斯堡有职位),多半是因为Verdier在这篇论文上的一条隐晦的引文(1975年2月Bourbaki讨论班
讲义 n\textsuperscript{o}~464),引作“J.P.~Jouanolou,~Thèse,~Fac.~Sc.~Strasbourg”(无日期,无题目)。可他是答辩委
员会的一员,同我一样——他的记性也同我一样失灵,或者说得更准,任性,把举行答辩的 IHP\dmark{}\origpage{1189}(Henri~Poincaré研究所)安到了斯特拉斯堡?谁懂,谁就懂吧!

还是这位Verdier,好意把自己持有的那份论文寄给了我。起先,看着这包208张散页\fnote{961}{(*)}{索邦大学论文管理处存有一份215页的存档本——Verdier的这份看来缺了六页。说不定,存放在该管理处的那份,是世上唯一完整的一份——而且,人家还向我保证,是装订成册的。那里想必有一个专为弃儿论文装订的部门,它们都是散着零件到的……},我以为那是一份草
稿的复印件——我记得当年Jouanolou跟我做这篇迟迟不完的论文时,我亲手捧过并细细评注过的正
是它。可Verdier向我确认:这确凿是那篇倒霉论文的定稿本;它看来从未荣幸地印过三四份以上
(我那本,连同我的批注,想必回到了Jouanolou手里,从此再没见过……),也从未装订过。

后来,Jouanolou惠赐的稍详细的说明(6月3日的信),加上给Verdier的电话,让我重新入了一点
水。显然,Jouanolou对他的论文到达了一个“饱和点”:他从一开始就毫无信念地做(而我也懒得把处
境看清楚\fnote{962}{(*)}{见已引的“学生与纲领”一节,n\textsuperscript{o}~25。})。到1969年,他想必到了一个堵点,以致无力哪怕稍作修改,来照顾我的
许多批注。我只好认了,随它去。反正,重读一遍之后,我觉得这份文本代表了一份认真的、可用的
成形工作,虽远非完美——分明好于“聊胜于无”,在没有更合我意的文本的情况下,可以充当一份
必不可少的参考文献\fnote{963}{(*)}{不管怎样,它至今仍是世上唯一的、以导出范畴的版本陈述$\ell$-进系数理论的文本——而且偏偏无处可购,喜上加喜。电锯从那里经过了……}。

\pmark 当然,我做梦也想不到:Jouanolou会以他自己的方式,报复他跟我做这项工作时的三心二意
\dmark{}\origpage{1190}——亲手凿沉它,几乎抹掉那份我如此珍视的出名的“参考文献”的一切痕迹!这又是
一次“往事回头”,我实在没什么可抱怨的(尽管怨气并不是没有!)。在同Jouanolou的关系里,我要
的是在他身上找到“臂膀”,来推一辆规模可观的大车的轮子。我预先认定,他Jouanolou是我事业的参
与者,却从来没有停下来想一想,那些执拗的迹象分明告诉我:根本不是这么回事。诚然,是他自己选
择来跟我工作的(同一位名声显赫的“老板”共事,想必他觉得划算,却没料到自己陷进了什么……);
也是他自己,从我不吝扶持的一大片题目里,自由选定了他的工作题目(这些题目,当然全连着那辆
“大车”——而那车,骨子里多半引不起他半点兴致)。换句话说:同每个人一样,Jouanolou也在他自己
之内、在他自己的愿望与抉择的层面上,同某些矛盾搏斗——这里,是在他的工作里。

我自己的矛盾,不在我同工作的关系里,而在一种对任务的如此的偏执,以致我在学生身上只看得
见受欢迎的臂膀,无法想象他们中任何一位会在我与之共事的工作中心怀分裂。有了对\emph{葬礼}
的长考所给的这一份额外的回顾,我如今发觉:在我的学生里,“分裂”的远不止Jouanolou一个。但他
是一个极端的个案:他是其中唯一没能认同自己所选任务的人;他的工作是在没有信念、没有欢乐
之中完成的。我在这一处境中的责任,在于我不肯真正去知悉它,宁可把本该次要的东西(完成我的
任务)\textbf{摆在前头},而把要紧的东西(学生所“选”的任务真正\textbf{属于}他、并且怀着欢乐去
做)撂在后头。

这想必就是为什么,Jouanolou也是我唯一曾在其身上察觉到怨怼的前学生(当然,那怨怼从不报姓
名)。怀怨是一个出口、一剂解药;它当然一无用处,只除了回避自己的问题(而人们很少往深里追
究)。可它毕竟是事出有因的;今天(二十年之后)我收割到它的某些果实,无权抱怨。

\pmark\origpage{1191}不到两个月里,接连面对Saavedra论文、继而Jouanolou论文的这些不寻常的插曲,让我切
身感到一件在《收获与播种》第一部分里才刚刚瞥见的事:早在我离去之前,以及紧接着的那些年
里,我的学生们和我之间,并不像我想当然的那样一切安好。譬如,在跟我做到国家博士论文层级的
学生们所通过的十二篇论文中,\textbf{四}篇赫然是针对宗师的“葬礼论文”!它们相继出现在1967至
1972年的五年之间,其中两场“葬礼论文”在我离去之前就已举行。第一场是Verdier1967年的论文,
缩减成一份28页的摘要,是我引入的新同调代数被埋葬的前奏;那套代数,Verdier曾受命去发展。
此事已有相当周详的交代\fnote{964}{(*)}{关于这一点,尤其见“赊账的学位论文与全险”与“随心所欲的荣耀——或暧昧”两条注(n\textsuperscript{o}s~81、170(ii))。},不必重复。第二场是Jouanolou1969年的论文,它为$\ell$-进上同调的形式主义——从导出范畴
与三角范畴的观点看(这对六种运算分明至关紧要;Verdier本该提供基本参考)——的埋葬盖了
章。第三场是Deligne1970年的(?)论文:一篇十足光辉的论文,也深深扎根于他从我这里得来的思
想\fnote{965}{(*)}{这是Deligne的“Hodge 理论 II”。关于这项工作在 Motif 瑜伽、在我的“系数理论”视野(包括一种“霍奇系数”的理论)之中的扎根,我在注“把一个个 i 点清”(n\textsuperscript{o}~164,尤其打字稿第739--740页)以及子注 n\textsuperscript{o}~164$_{1}$(第805--806页)里给出细节。同 M.~Raynaud与 C.~Contou-Carrère一样,Deligne自选了他的工作题目——包括论文题目——并没有等我给他出题,而且完全独立地做了这项工作,甚至在几乎完成之前都没跟我提起。尽管如此,他的工作(论混合 Hodge 结构)扎根于我的思想,比Raynaud与Contou-Carrère更深:后两位主要使用我带来的语言与技术,而他们各自追究的问题域,则是全然原创的。
的确(照今天刮的风),思想随风而去——尤其当它们没有发表的话(几天前,Serre又斩钉截铁地向
我解释过一遍……)。}——而我的名字在文中连提也不提!第四场是Saavedra的论文——方才长
篇谈过:不是那位推定的作者\fnote{966}{(*)}{这至少是我在“Monsieur~Verdoux——或献殷勤的骑士”(n\textsuperscript{o}~176$_{5}$)那一注里达到的确信。},而是别人,以他为人熟知的娴熟手法,陈述了第三个人关于 Motif Galois群的思想与\dmark{}\origpage{1192}结果(经由一套完整的所谓“塔纳卡”范畴理论——而且是整整四种!),却对我这卑不足道、亡故之人只字不提!

这四场葬礼式论文(它们正是“四种运算”——大写的那种!——的序曲)分明以许多方式互相勾连
\fnote{967}{(*)}{往深里探测这些联系,当然有趣——但如我在几行之下所说,现在不是时候。}。它们相续于不到五年之间,始于紧接 SGA~5 讨论班结束的那一年。那场讨论班,看来确是
我前学生们身上掘墓心性的起点与集结之地——而且早在我离去之前!这些心性先于我的离去而存
在,这是一个耐人寻味的情况,关乎\emph{葬礼}的那个“第二平面”——由“从前的”全体前学生组
成;这个情况,我还没有真正整合进一个总体的理解。眼下,这“第二平面”似乎是三个平面中最不被
理解的一个。但现在不是重新开启思考的时候;接下来的几个月,想必不会缺少新的材料,它们会由
我的前学生们亲手送来。到那时,再把它们拼成“第二平面”的一幅活的总体图景。

还有第五篇论文\fnote{968}{(*)}{在共十四篇论文之中——十四位(无论“从前的”还是“后来的”)跟我做到国家博士论文层级的学生。也就是说,按数目算,\textbf{三篇之中就有一篇以上}是葬礼论文——已经不算糟了!},对我而言,它也插进葬礼论文的序列,不过是一篇“之后”的论文——比前一序列晚了十
年。那是Contou-Carrère的论文,1982年12月通过;它也在不止一个方面很特别。它不同于前四篇之处在
于:Contou-Carrère为讨好要人们、为求得原谅自己曾多多少少当过我的学生而作出的英勇的掘墓努
力,并没有能避免Verdier(他以为选他做论文导师是聪明的
\fnote{969}{(*)}{而且那时候,我还一直(照Contou-Carrère本人向我保证的)以为自己是他的官方论文导师。我得知还有一位“平行”导师的存在——在这一对里,本该是我充当“备位”导师,以防万一……——是在Contou-Carrère不得不回头找我的时候;与此同时(局势已经变得有点太那个了),他向我揭示了Verdier扮演的角色。经年累月地玩着这种没谱的把戏,Contou-Carrère最后几乎不再做数学,难怪。可以说,他不是唯一的一个……})突然装模作样地要把他\textbf{搞沉}\dmark{}\origpage{1193},不宣而战——于是,不得已,他又回头找我。我本不必
出场充当论文导师:Contou-Carrère自己找到了题目、自己发展了方法,我也没有跟进他的工作,何况
他的工作所在的语境(约化群概形)我已有些生疏。尽管如此,他工作的初始想法——为 Schubert 闭
链的闭包搞一种“等变”的奇点消解方法——直接受启发于我(约1975或76年)向他原原本本解释过的
一个想法:关于约化群在自身上的伴随表示,其轨道闭包的典范的、同步的奇点消解
\fnote{970}{(*)}{六十年代末,我曾被Brieskorn关于所谓“有理”(曲面)奇点的漂亮工作、及其与某些单根系(根全等长的那些)的联系所吸引;我给自己出了一个问题(不用说,有点想入非非):用与根图相应的单代数群,直接描述一个有理奇点。就这样,我得到了所论奇点消解的一份非常简单(说白了,甚至是显然)的几何描述,靠的是 Killing 对;外加一整套精彩的猜想作为钥匙,后来我有些忘了,当年讲给过凡愿听的人。可我什么也没发表;而照Serre刚向我惠赐的新公理——归第一个捡到的人所有——我发觉,有人捡了不少,难免。有时候,换换公理,真方便……}。不用说,Contou-Carrère早就嗅准了上流社会里风往哪儿吹——那上流社会,他正合法地盼着跻
身;对这层渊源,他守口如瓶。倘若大家又都开始提这类不可捉摸的东西,那还了得:一个\textbf{观念}
(还没发表呢),据说是为了\textbf{激发}另一个(您给评评理……)——当然,除非被光荣引用的,正
是那些让所呈工作增光的名字之一(到那时,自然也无须解释为什么致谢滔滔,那致谢想必是名正
言顺的……)。

\asterisms

\centerline{\headfont “四种运算(对一具尸体)”终}


% ================================================================
%  18.6 Les chantiers désolés (原书 p.894–908)
%  源文: _work/chunks/chunk131.md, chunk132.md, chunk133a.md, chunk133b.md
% ================================================================
\section{荒废的工地}

\subsection{(1) 悬而未决之事}   % 18.6.1

\notehead{176$'$}{}

\pmark\origpage{1194}(3月25日)昨夜,我在床上辗转了几个钟头,让自己重新浸入“Motif 的瑜伽”里,而不是像
本该的那样安然入睡。稍后,我也没有接着写笔记,而是又花了一两个钟头,潦潦草
草画出一堆蕴涵图,涉及我所知的、一个德拉姆上同调类(比方说,取特征零域上一个非奇异射影簇的上同调类)
要成为“代数的”所需的那些内蕴条件。我一共找到\textbf{十二}种路数,统统是 Hodge 猜想和
Tate 猜想的变体\fnote{971}{(*)}{(3月27日)这十二种变体,应当每一种都能对任意一个特征零的基本概形 $X$ 给出 $X$ 上一类相应的“系数范畴”(“系数类型”一词,取注“墓上的旋律——或自负”(n\textsuperscript{o}~167)中所说的含义)。倘若所论猜想为真,这条系数范畴应当把 $X$ 上的 Motif 范畴作为一个全子范畴(三角的)包含进来(这个猜想,在 $X$ 取一个域的谱这一特殊情形下,说的恰好就是这件事……)。细节我留给《数学反思》第3卷中专论 Motif 理论的部分(“Motif,吾爱”)。这也等于说,这十二种人人皆知的猜想的变体,给出(至少先天地)同样多种特征零域上“Motif”的不同观念。这就让将来会有十一位我朋友皮埃尔的摹仿者,各自“发现”属于自己的那种 Motif 观念,同时装作不知道别人的——尤其是(十五年来照例如此……)装作不知道某位已故者的,他主要正是以偏爱无用细节而著称……}。与此同时,我也说服了自己:要在一个 $\mathbb{Z}$ 上有限型的概形上定义“那个”(三角的)Motif 范畴——或者至少给出它的一个极紧的近似(假如这仍不是“那一个”对的)——所需的材料,应当差不多已经在手,只消再有一套“神秘函子”的理论即可;这个函子,是我在六十年代末曾公设过的\fnote{972}{(*)}{这个“神秘函子”的问题——它在特征 $p$ 的晶体上同调(经由“滤套 $F$-晶体”观念,$F$ 即“Frobenius”)与特征零的 $p$-进上同调之间架起“缺失的一环”——对我们的代数簇上同调认识而言显然至关重要,可将近二十年过去,自我以再明白不过的方式把它提出以来,这个问题至今无人认真问津……}。

在这里就此铺开,当然不合时宜。可我现在已经看清:正是在此时此刻,眼看 Motif 这个主题在我
把它留在一伙靠不住的人手里十五年之后,落到这般可怜的荒废地步,把当年我已抵达的那些思想
的几条主线勾画出来,恰是其时。我没有心思再等下去,等哪天有了闲暇(等《追寻栈》写完之
后)再去写那部本该写的“\textbf{那}部”系统著作;那份对一个梦的原原本本的讲述,本该作为让这
个\textbf{梦}终于扎根于精心酝酿成熟(并且发表出来……)的表述沃土之中的第一个\dmark{}\origpage{1195}大步骤,好让它
依着自己的本性开花结果。除了已为这部“数学虚构”之书预告并宣布过的第一块界石——即所谓
“六种运算”的代数对偶形式主义的一个概要——我还要在《数学反思》第3卷里附上一篇短文,提出
几个与 Motif 和代数闭链相关联的关键问题\fnote{973}{(*)}{照我目前的出版计划,《收获与播种》的前四部分(以\emph{葬礼}的第三部分亦即最后一部分作结)应构成《数学反思》的第1、2卷。第3卷将由《收获与播种》的第五部分(读 C.G.~Jung自传的笔记)和若干较短的文本组成,其中多数已在《引论》中预告。这样,《追寻栈》的第一卷就预定为《数学反思》的第四卷。}。眼看它们在坟墓里凋敝,我于心不忍;我巴不得看见它们重见天日,重新加入季
节的轮转……

思考回到\emph{葬礼}上来,已经五个多星期,而且再不肯离开。想必正因如此,“孤儿们”的念头——
他们被抛在一个病了的世界里,无人认领——近日又执拗地回到我心头。上一次有根有据地谈到其
中一位孤儿的笔记,是“墓上的旋律——或自负”(n\textsuperscript{o}~167),主题与昨夜以及稍前的这段
Motif 思考正相毗邻(我方才提到过它)。那正好是一个月前,就在我动笔写一条预定(这早已预先
定好)称作“四种运算”的笔记的前一天,我还不知道等着我的会是什么!到头来,一条笔记变成了
十六\fnote{974}{(*)}{(5月9日)写下这几行才不过两个星期,临头冒出的新事实,又把“四种运算”这场调查重新推动起来;它已经添了足足二十来条新注和子注。},我还以为永远做不完了——可到头来,我终究还是把这几件拖泥带水的“运算”走完了一遍!

而在这里,我最想做的,首先是回到那些孤儿身边,至少逐个叫出他们的名字——这也许对他们
有点好处,对我则一定有。第一次谈到他们,是一年前,就在那道同名的注“我的孤儿们”里,去年
3月底,与紧随其后的那道注“拒受遗产——或矛盾的代价”(注 n\textsuperscript{o}s~46、47)一口气写
成。写下这些笔记、给它们起这些名字时,仿佛有一股冥冥中的先知引路,我还没有料到,我留下
的那些东西,究竟“孤”到了何等地步\dmark{}\origpage{1196}——其含义之沉重、之切肤,连做梦我也想象不到;也
没有料到,我当时初次怯怯地指出的那道“矛盾”,会延伸到多远。这段记忆立刻又勾起另一段,是
前一个月的:我看见自己写下——仿佛是另一个比我更透彻的人借我的手在写——“\textbf{谁也斗不
过腐败}”。那是在写“无爱的世界”一节(n\textsuperscript{o}~19)的时候。我还记得,看见“\textbf{腐
败}”这个词白纸黑字地立在眼前,我起初愣住了。我心里某个“讲道理”的声音在训斥我:真是的,
你下手未免太重——“腐败”可是个大字眼,可不能信口开河!你还是换一副腔调吧!

我着实自我盘问了好一阵子,也许有几分钟。随后我明白,我既不会改掉这个“大”字,也不会加
一条注,解释说这个词是笔下一篇酣畅之际溜出来的,不必太当真。从那个世界里,这里那里,陆
陆续续回到我心头的这些“呛人的气息”,我心底有个人——比那个负责挑选“讲道理”标签的“我”
更有眼光——早就一清二楚它们意味着什么,根本不等我费神替自己编个说法\fnote{975}{(*)}{这段经过,我先是在去年3月“那张条子——或新伦理”一节(n\textsuperscript{o}~33)里讲过;两个月后,在发现\emph{葬礼}之后,又在 Notes 93--97 那一组原原本本得多的笔记里讲了,它们构成送葬行列之十“灵车”(与掘墓人同行)。}……

我也清楚记得那一天,那天的思考突然变了质地、我心中那\textbf{另一个}接手执笔的确切时刻。
那紧接在一段追忆之后:我数学生涯最初几年,数学圈子里那团裹着我的温情——多亏长者们给予
的接待,一直深到他们家里:Schwartz家、Dieudonné家、戈德芒家……转折就发生在我接着写下“显
然,对今天许多年轻数学家来说,要命的正是被切断了……与一切情感的潮流、与温暖的联系
……正是这一点折断了工作的翅膀,又夺去了它比一份郁郁寡欢、朝不保夕的营生更深的意义……”——就在
此刻,这个\textbf{无爱的世界}陡然在我眼前浮现,活了起来,再一次向我发问……

去年,“我的孤儿们”这个名字,是不请自来地落到我离去时留下的一切头上的(它们被我托付给
的那些近亲宣告了“死亡”……)。想必是因为,这个名字\dmark{}\origpage{1197}道出了一种\textbf{实实在在}、触手
可及的真实:我所“留下”或“托付”的,既不是“物件”,也不是“财产”,而是一些\textbf{活
的东西}。如今想起它们,我依旧把它们当作活的东西来想:茁壮、丰饶、生来就要生长、
绽放、又要孕育和催生别的茁壮而丰饶的活的东西。我确实感到自己留下了一份“财富”,可那不
是银行家的财富,而是园丁的、砌墙匠的财富:他们用双手让这些繁花似锦的园子、这些宽敞好
客的房子从地里长出来。这种贵重(乃至脆弱)之物的感觉,最紧地系于那些\textbf{观念}、那些
\textbf{问题}、那些我知道生机勃勃的\textbf{大主题}——我把它们留在了更年轻的人手里;这些东
西还需要劳作与呵护——远甚于那些我打磨到位的工具,或者那些我一路盖到头、布置停当的“房
子”\fnote{976}{(*)}{关于我心里那股催我“盖房子”(数学的房子)的冲动,见注“阴,作为仆人——与新主人”(n\textsuperscript{o}~135)。}。自有别人照自己的意思去那里忙碌,做他们自己的饭,享他们自己的清福;哪间房子嫌小
了,他们会照自己的需要扩建——就像我自己也常常不得不一扩再扩,虽然当年看上去,我“眼看
得很远”了。可是,我的数学家过去之所以仍旧攥得住我,靠的恰恰是\textbf{那些悬而未决之事},
是那些刚刚在绝好的地基上、用那么美的石料开了工的工地(而工人们已经走了,拿走了他们中
意的,糟蹋了剩下的……)。正是这些被撂荒的\textbf{工地}——今天我重访时,只见它们横遭劫掠、
残破不堪——我现在要把它们逐一巡视。

\subsection{(2) 吝啬鬼与老朽}   % 18.6.2. L'avare et le croulant

\notehead{177}{}

(3月27日)昨天一整天都交给了庶务。为了送去打字,我不得不把\emph{葬礼}第三部分亦即最后
部分的开头五十页重读一遍——连这里那里的小小措辞修整,外加又添上的几条脚注,足足花了五
个钟头。“阴阳之钥”的打字眼看就要完成。自从为这部分打字吃了那些没完没了的苦头之后\fnote{977}{(*)}{关于这些“苦头”(这还是委婉的说法),见注“祈祷与冲突”(n\textsuperscript{o}~161)的开头,以及《收获与播种》第五部分即最后一部分的开卷之注“Jung——或‘恶’与‘善’的循环”的开头。},我最终退而求其次\dmark{}\origpage{1198},找
了一位系里的秘书,在正职之外替我做。苦头到头了,谢天谢地——她做得又认真又麻利,一
个星期三十来页,无可挑剔。总能到终点的。也该到终点了!

除此之外,在一个 $\mathbb{Z}$ 上有限型的概形上,以三角范畴的形状构造 Motif 范畴这个问题,
继续在我脑子里打转——我又在床上熬了大半夜想它,没有睡觉——当心啊!起初似乎我的想法只
对特征零的概形行得通(比方说 $\mathbb{Q}$ 上有限型的),在基底 $\mathrm{Spec}(\mathbb{Z})$
本身上就不大灵;随后我记起,六十年代我就原则上确定过有限域上 Motif 范畴的结构。假定把我
当年做的工作明确化,我终于看到,一般情形下至少有一个完整描述的原理显现出来——说实话,
够紧凑的,但在我看来绝非不可企及。比起我六十年代的想法,唯一的新配料,是Mebkhout的哲学,
它凝结在那条说来奇特的“好上帝定理”里。除此之外,我拿来作假设性配料的,是“神秘函子”的理
论。眼下它还拿不出来,这肯定不是因为它“无从下手”(借一个我早领教过的说法\fnote{978}{(*)}{这正是我那位杰出的前学生Deligne用以埋葬“标准猜想”的斩钉截铁的考语——于是,在将近二十年里,我那些大胆的同时代人竟没有一个敢向它们下手。完整引文见注“悼词(1)——或恭维”(n\textsuperscript{o}~104)。}),而
是因为,我所认识的那些在代数簇上同调里做工的人,连在数学里也丢掉了对本质之物的分寸——
当然,他们为一场需要悉心照料的葬礼忙得不可开交……

平心而论,Deligne围绕Weil猜想做的工作——“Weil~I”,尤其是“Weil~II”——终归会在紧要关
头派上用场,到要为那些被指望表达 Motif 的系数范畴搭建六种运算的时候。尽管如此,还是得
有一位“糊涂”的、如人所愿地老朽不堪的已故者,在十五年之后灵光一现,从他那些亲爱的学生
与继承人乐于把他安置其中的软壁棺材里坐起来——而他本人对这一切毫不知情,而且差不多把生
前懂得的那一点点也忘光了——这样,在基底概形 $S$ 上描述 Motif 范畴的问题才算\textbf{一字一句地提
出}来,而搭造的原理(把与一个 Motif 相关联的一切已知结构要素都考虑在内的那种)也才顺带\dmark{}\origpage{1199}——纯属巧合地——终于得到明白的表述\fnote{979}{(*)}{正如我昨天在思考里预告的,我打算把这一描述收进《数学反思》的下一卷,附上一幅“Motif 广景图”的(极简略的)总轮廓——暗中的 Motif 把戏,闹得够久了。这里就预先声明:所设想的构造原理,不依赖代数闭链上任何一种猜想,诸如“Hodge”或“Tate”(或者昨天谈到的那十二种变体之一)。}。

1982年那部关于 Motif 的“难忘文集”之后,这份“Motif 藏宝”——十来年间一直是一个人的专
属禁地、一个秘密——似乎变成了三四个人共有的宝藏;这几位彼此往来,一派合谋者的神情,或
者像某个秘密教派的入会大佬。可是,把几个简单的问题白纸黑字写下来、提交给公众注意,不过
需要几天;要是想稍加用心地把它们界定清楚,明白交代手头有哪些配料、还须发展哪些配料,也
不过几个星期。从1970年起十五年,从“难忘文集”起三年,先是一个人,后来是那么几位,竟没
有一个肯从自己固然宝贵的时间里拿出这么几天——更不必说几个星期——这必定有极好的理由,
尽管他们当中没有谁打算去探个究竟。可他们煞费经营的那种气氛、他们赖以自持的那种心气,本
身就已经是一场发现之旅的堕落:发现之旅沦落为凌驾于他人之上、乃至蔑视他人的手段。这样的
气氛天然传播腐败,与创造背道而驰——哪怕乐在其中者是天才里最耀眼的几位。守着这样的心境
——一头孵着金币的守财奴的心境——他们切断了自身之中的创造之力,正如他们乐于扼杀他人身
上的这种力量。

\subsection{(3) 巡视工地——或工具与视野}   % 18.6.3. Le tour des chantiers - ou outils et vision

\notehead{178}{}

(3月30日)前天我满五十七岁,稍稍歇了歇工。只为“阴阳之钥”的结尾做了一点打字校对,昨天
又继续了一点。这是一桩让人歇息而愉快的活儿——至少,当打字的人自己也肯上心、一份我全身
心投入的稿子不至于被弄得面目全非回到我手里的时候,是如此。这回,我给自己放了两天假,\dmark{}\origpage{1200}细
细重读了五十来页清稿,就为了这里那里揪出一个还没放对地方的逗号……

不过,工作的劲头并不在顶点。几个星期以来,我心里有一种悲哀在提醒我:比起把手头这些笔记
一路写到它们自然的终点,还有更要紧的事在等我。我写着,像在逆水行舟;然而我知道,只要不出
意外、不遇不可抗之力,我只有给\emph{葬礼}终于画上那个句号,才会停笔。可把这股悲哀压下去、
放逐出去——它于是重得像一块石头——不让它在这些笔记里发声(除了眼下这一刻,顺带地、影
影绰绰地),这已经相当清楚地表明:我的思考,早在不知多久之前,就不再有资格称作“冥想”了。
它落在一种分裂之中:一边是写的人(却提防着不肯整个人扑进去\fnote{980}{(*)}{可是,就在上一段里,我刚写过(心里并无半点保留):托付去打字的那些稿子,我是“全身心投入”的。可见同样的话(或差不多的话……),随语境不同,可以有不同含义,或指示不同的细微差别。}),一边是活着、感受着的人(却也不停下来
“摆弄”他所经历的、浸透其中的意义)。在那里我感到,到达那个“句号”(当然,也不是说就把尚
待看见、尚待说出的这一点点草草了结……)然后返回我自身,已是刻不容缓……

除了笔记的活儿,这几天还有另一件事在给我调剂。那是一场数学思考,不以我意志为转移地重新
开了头。几天来我已明白:一套 Motif 理论形状的构造,带着二十年前我所望见的那全部气魄,离
“地平线”其实并不像我想的那么远。甚至有可能,一套“完全成熟”的理论——连六种运算的完整
形式主义(外加双对偶)在内——对肯把自己整个儿扑上去的人(不以掘墓的心性糟蹋自己的创造
之力)来说,只是几年功夫的问题。我还看出,要显式描述“那个”Motif 范畴\dmark{}\origpage{1201}——比方说,取绝对基底 $\mathbb{Z}$ 上有限型的概形(人们总应当能归约到这种情形)——有\textbf{两把}“钥匙”\fnote{981}{(*)}{其实还有第三把“钥匙”,这里不提,因为在我看来(不管对错)它没那么棘手。这就是“de~Rham--Mebkhout系数”的恰切定义(先不管滤套与 $F$-结构),比方说,取绝对基底 $\mathbb{Z}$ 上光滑概形。这一定义应当同时提供特征 $p>0$ 的一般晶体系数“那一个”恰切定义的钥匙——而我那些亲爱的前学生们(这一回首当其冲是Berthelot),始终没能、或者始终不愿把它提炼出来。
1983年6月(快两年了),Mebkhout向我讲解他围绕“好上帝定理”的“哲学”时,我的印象是:他对复数域 $\mathbb{C}$ 上光滑概形的(可构造的)离散系数范畴给出的“纯代数”(“德拉姆型”)描述,恰是Deligne在(别处已提过的)他1969/70年间在法国高等科学研究所所给讨论班(若无记错)中采用、却从未发表的那条进路的对偶;那条进路,靠的是带联络的 pro-模。我猜,从一个观点到另一个观点的过渡,经由对偶化函子 $\mathrm{R}\mathcal{H}om(-,\mathcal{O}_X)$ 完成——相对于所论概形的结构层——它把有限型的 $\mathcal{D}_X$-模(可以看作配备了一个可积联络的“ind-凝聚的 $\mathcal{O}_X$-模”)变成“pro-凝聚”的模(同样配备可积联络)。Mebkhout观点的长处,在于它为“好系数”给出一个简单而深刻的代数表达($M$-凝聚、完整性、正则性),这正是Deligne所缺的。Deligne观点的长处,在于它给出的是(与所要表达的超越性质的系数之间的)一个等价,而非反等价,并且更宜于表达所论系数范畴的乘法结构(张量积)。我猜,实践中人们会常常愿意在两张互为对偶的图式上同时工作。Deligne的解读,经由可积联络的模(或 pro-模)这一直觉,在我看来更贴近直接的几何直观。这一点的表现,尤其是(若基域为 $\mathbb{C}$):一个可构造的 $\mathbb{C}$-向量空间层,对应的是唯一一个带联络的 pro-模,而不是这种 pro-模的一个复形。正因如此(读者想必猜得到,这让我多么遗憾……),我预料最终最宜于用来发展六种运算的形式主义、并充当 Motif 范畴构造中第三种配料——钥匙的,恰恰会是他的观点——尽管正是他,毫无惋惜地把它埋葬了,仿佛借此把那位被拒不认的宗师所遗留的系数问题也一并埋掉……(5月9日)关于这一点,另见子注“……以及羁绊”(n\textsuperscript{o}~171(viii)),以及“五张照片”(n\textsuperscript{o}~171(ix))。}:一把,是“神秘函子”的理论,要有足够的普遍性与灵活性,好过渡到相应的三角范畴,把de~Rham--Mebkhout系数与寻常的 $p$-进系数(特征零)连起来;另一把,是在一个\textbf{有限}域 $k$ 上显式构造 Motif 范畴的问题(最好经由一种“纯代数”的构造,不援引 $k$ 上的代数几何),此外还有从 $k$ 上分离有限型概形(先从射影光滑概形开始)到这一范畴的“Motif 上同调”函子。后者我在当年已经构造到了差一个等价的程度,用的是Weil猜想与 Tate 猜想作启发\fnote{982}{(*)}{如果我记得不错,我当时只限于描述半单的 Motif 范畴。按同一原理,这一构造稍作变通,就能给出一个颇可信的候选者,充当不必半单的 Motif 范畴。我在这里说“Motif”,实际上说的是“iso-Motif”,即差一个同源不计的 Motif。但若对每个素数 $\ell$ 使用“$\ell$-进实现”函子,便能从那里重建不经同源的 Motif 范畴(其中诸 Hom 将是 $\mathbb{Z}$ 上、而非 $\mathbb{Q}$ 上有限型的模)。
说到我的构造“以启发的方式”使用 Tate 猜想,不可按字面理解。倘若确实存在(比如说,在有限域上)射影光滑概形上的上同调类,它们是“Motif 的”(这含义正待提炼)而非“代数的”(即并非来自代数闭链),那么 Tate 猜想(同理还有 Hodge 猜想,这回是在 $\mathbb{C}$ 上)就应当改述,把“代数类”换成“Motif 类”。假定人们果真能(如我下文所提示)在有限域 $k$ 上射影光滑概形的范畴上,定义出典范的(并且在适当意义上应当是“普遍的”)上同调函子,取值于已经构造好的范畴(所谓“$k$ 上半单 Motif 范畴”),那么这就顺理成章地给出一种形状完整的定义,界定人们将要称作“Motif 类”的上同调类:即 $\mathrm{Hom}(T^i,H^{\bullet}_{mot}(X))$(2i 维情形)的元素,其中 $T$ 是 Tate 对象,而 $H^{\bullet}_{mot}$ 是所拟的假设函子。这正是为什么,这个函子的构造,在我看来眼下是形塑 Motif 理论(不再是六十年代那种假设性的,而是形状完整的)的\textbf{头号}关键问题。}。

我对这一\dmark{}\origpage{1202}构造的正确性没有半点怀疑。尚待完成的工作——无疑明显更精细——在于:用给定的
有限域 $k$ 把这条范畴“钉死”,尤其是定义“Motif 上同调”函子,哪怕先只在 $k$ 上阿贝尔概形
的范畴上定义(这应当就足以把所求的范畴“钉死”……)。这第二个问题,在我看来,比起神秘函
子,较少技术性,而更直接是“几何的”。而且,它向我显形为\textbf{那把}能解开标准猜想\fnote{983}{(*)}{“标准猜想”一词在这里不可按字面取义,正如前一条脚注里的“Tate 猜想”。毋宁说,这些猜想的表述里,应当把所论的闭链类扩大(起初只限于代数闭链)。在“重新调整过”的标准猜想的“最终”表述里(即便它们按原样已然成立),“代数的”上同调类还要再换成“Motif 的”类。关于标准猜想,我将在“Motif,吾爱”(《数学反思》第3卷)里更详细地回头来谈。}的钥匙,并由此也解开特征 $p>0$ 的上同调理论里那些恼人已极的整性问题。凡此种种,使得这个问题对我有一股强大的吸引力!今晚已是第三个晚上,我重新拿起笔记,打算快速巡视一番那些在我看来最灼热的主题——十五年前我告别数学舞台时,被我的学生们、也被所有人抛下的那些主题\fnote{984}{(*)}{这些主题的一个极简略的初步“巡视”,见去年那道注“我的孤儿们”(n\textsuperscript{o}~45)。}。这一回,我总算要成功了!

\asterisms

\pmark\origpage{1203}\textbf{工地一:Topos}。我在这里提到它,主要是备而勿忘;在注“我的孤儿们”(n\textsuperscript{o}~46)
里,我已经相当周详地说过它。由于我的某些前学生——Deligne首当其冲——乐于糟蹋这条
关键性的统一性观念,自我离去之后,它被打入冷宫,只落得在边缘上苟活。如我在所引注里提醒
的,topos 与形形色色的“多重叠”其实在几何里俯拾皆是——不过,当然,尽可以看不见它们,正如
人们几千年里看不见对称群、看不见集合、看不见数零一样。

一套贴着拓扑直觉、柔软而精微的关于 topos 的语言,已在 SGA~4 的前两卷里被精心发展完备
(就是那块著名的“废话矿渣”——光辉的“SGA~4½”一书第一讲的引言里,Deligne如此称呼它)。这是
Leray所创“层”的观念周遭的语言与直觉的自然结晶;层直觉与层工具发展中的这第二步(或这“第
二口气”),在我看来,分量不下于第一步(其临时表达即戈德芒那本名著)。正是这一视野,让
$\ell$-进与晶体的上同调工具得以问世——然后,恰恰是那班佯装把工具据为己有的人,把它无限期
地埋葬了。

SGA~4 中关于 topos 的种种展开,并不自命完备与终极,但我认为,对 topos 式视野的大多数
直接的几何应用而言,它们绰绰有余。正如一般拓扑或寻常的层论,“topos 式的一般拓扑”本身,
在我看来并不提出什么真正深刻的问题。它是一套精心打造的语言,服务于对形状的拓扑-几何直
觉的某种扩展;这种扩展,是事物本身口授给我们的。这一视野长期蒙受的污名,以及落在它头
上的\textbf{嘲弄},在我眼里属于七十年代、八十年代数学世界的几大凋败之列。

这里说的,不是一间需要有人来救活的“破败工地”,而是一所全然竣工、设备齐全的房子:那些
住过它、又本被唤来在其中工作与生活的人,选择离它而去,还顺手撵走了盖房子的工匠。房子宽
敞、健全,一切各安其位,一如工匠出门做别的营生的那一天。它若还缺什么,缺的\dmark{}\origpage{1204}不是他这双手
的劳作,也不是任何人的劳作。也许,工匠本人对这些东西的敬意——那双手怀着爱做
出来的、他知道是美的东西——会让这股暴戾与轻蔑的气息消散,让这间本为迎接而造的房子重新
学会迎接。

\textbf{工地二:上同调的语言}。这首先是指导出范畴的语言;再者,是我为非交换上同调引入
的诸观点——两者都在五十年代下半叶。

头一条路线,本该由Verdier那篇著名的“论文”来承载;而Verdier本人对自己论文的埋葬
\fnote{985}{(*)}{关于这一点,见注“赊账的学位论文与全险”(n\textsuperscript{o}~81),以及“随心所欲的荣耀——或暧昧”(n\textsuperscript{o}~170(ii))。},同时也就是代数同调中导出范畴这一观点的埋葬。这一观点,曾在六十年代代数几何中上同调主题的繁盛期——尤其在对偶形式主义、不动点公式(Lefschetz--Verdier型)的发展中——扮演过关键角色。实践的需要,早已显出Verdier六十年代初发展的三角范畴框架之不足;这一框架,至今没有像本该的那样得到更新。

“非交换”一路,有Giraud的论文打下的良好基础工作,但它止步于一种 1-叠的形式主义,只能把
上同调对象直接几何化到二维为止。以 $n$-叠与 $n$-gerbe 的语言发展非交换上同调形式主义——众
多例子急切地提示了这一点——一路撞上严重的概念困难。鉴于在某个上流社会里,基础
问题跌入了人人不屑——说得更直白些,人人轻蔑——的境地,这些困难从来无人问津,直到两年
多以前,我亲自去啃\fnote{986}{(*)}{这指的是1983年2月我致Daniel~Quillen的信中的思考;在那封信里,我发现可以“一跃而过”那道豁开的“炼狱”——随着 $n$ 增大,($n$ 不严格的)$n$-范畴(或按我如今的说法,$n$-叠)的形状化描述中,似乎总要冒出来的、越收越紧的相容性关系。$n=2$ 的情形已然不是一桩轻松活,而 $n=3$ 的全部相容性,我想还没人有勇气把它们一一写出。这封信,如我下文还要提到的,成了长途跋涉《追寻栈》的发令枪;次月,乘着已经启动的思考的余势,我动身了。
这封信没有被视为值得一读,也没有得到回复。一年多之后,我才收到当事人的一点评语;对此我在“一段过去的重量”一节(n\textsuperscript{o}~50)里有交代。(参见打字稿第136页第二段。)}。

\pmark\origpage{1205}如今我看见两条路线在一门新学科里汇合——我在别处\fnote{987}{(*)}{见注“阴,这仆人——或慷慨”的子注 n\textsuperscript{o}~136$_{1}$(尤其打字稿第638页)。}提议把它叫作\textbf{拓扑代数}(l'algèbre~topologique):它是传统同调代数(当然是导出范畴风格)、同伦代数、(尚在娘胎里的)$n$-范畴、$n$-广群及相应的叠与 gerbe 的形式主义,以及 topos 的视野——它如今提供了已知最广的“纯代数”框架,供安放拓扑直觉——这几家的合成。这样一种合成的初始观念,六十年代已经齐备,其中包括\textbf{dérivateur}的观念——它注定要取三角范畴这个不够用的观念而代之,并且同样适用于“非加性”的语境。同伦代数中的某些重要进展,比如Bousfield与Kan七十年代初发展的同伦极限与同伦余极限诸观念——他们并不知道我的想法(在我亲爱的学生们那里,这些想法被当作“Grothendieck式的老糊涂涂鸦”打发了)——正落在这些想法的延长线上。

两年前,我以致Quillen的那封信\fnote{988}{(*)}{关于这封信,尤其见“一段过去的重量”一节(n\textsuperscript{o}~50,打字稿第136页第二段)。}起手,为我看见该做的工作起草了一份总施工图。那封信是《追寻栈》一书的发令枪;其第一卷(“模型的历史”)业已基本完成,无疑将作为《数学反思》第4卷出版。我预计还需要一两卷、一两年功夫,才能完成这一前期勘探——勘探的是一片极其丰饶的矿脉;二十年后,我似乎依然是唯一看见它的人。所以,这确实是一处撂荒了十五年的工地——但近一年来,它已在我手中复活。《概要》的写作,继之《收获与播种》的写作,曾把它打断;然而一旦《收获与播种》和那些(篇幅都有限的)文本——它们将与《收获与播种》最后一部分一起构成《数学反思》第3卷——写完,我就打算回去,把它做到底。

\pmark\origpage{1206}\textbf{工地三:六种运算,双对偶}。这指的是我引入Poincaré式或Serre式对偶形式主义(取离散或连
续系数)的那个观点。我当年起的“六种运算”这个名字,被我那班上同调学生细心地铲除
干净了。他们不过这里那里挑几样合用的来使,却把这几样合成一体所成的那种结构(连同双对偶
形式主义)扔进了盈亏账,尤其是,把这一观点所提供的那条无可替代的主线扔了(这条主线,对提
炼好的“系数范畴”至关紧要,见下文)。二十多年来,这套形式主义早已存在并经受了检验;可在“局
中”的人里,没有谁(除非是在注定保密、我无从知晓的文字里)费过神,去提炼那“代数总目”——
凡有这种“万能”对偶、以六种运算形式主义表达的地方,通用的那一套代数公式\fnote{989}{(*)}{(5月9日)在 SGA~5 最早的一次报告里,我曾煞费苦心,把这套公式说得详详细细;它将像发动机的神经一样,驱动整个即将展开的讨论班。这一讲是 SGA~5 里最关键的一讲,却在“屠杀版”里不见了踪影。全卷之中,竟无一字暗示它曾存在!见注“先祖”(n\textsuperscript{o}~171(i))里打字稿第942页脚注(*)。}。

可见,这里严格说不是一处“撂荒的工地”(要补的形式化工作微不足道),而是一个被系统性回避
的丰饶观点(正如 topos 的遭遇)。这场抛弃,无疑在很大程度上,造成了我在代数簇上同调主题上
看到的可悲停滞(少数例外\fnote{990}{(*)}{这些“少数例外”,主要是(1981年以前)Deligne那两项重要工作 Weil~I、II,以及晶体上同调、Barsotti--Tate 群的 Dieudonné 理论在特征 $p>0$ 的一般基底上的若干零星结果(后者是我在1969年前后开创的)。如我在别处指出的,在“好上帝-Mebkhout定理”的余波里,有过一阵回潮(定理与作者照旧一个被无视……),尤其是Mebkhout层理论的发展(它们被错误地叫作“反常层”,顶了本该属于别人的名分……),由Deligne等人展开。}除外)——尤其与我1955到1970年间赋予这一主题的强劲起飞相比。

如我已在《引论》(I~8,“一个秘密的终结”)里预告的,继《收获与播种》之后\fnote{991}{(*)}{提醒一下:这指的是《数学反思》第3卷,原则上它也收《收获与播种》的最后一部分。},我打算收入一篇短
文,勾勒“六种运算”形式主义的要点。多亏我那些\dmark{}\origpage{1207}学生的照料,今天它的存在本身已无人知晓——
除了曾直接参与 SGA~4(1963/64)与 SGA~5(1965/66)两个讨论班的人,而他们显然忘了\fnote{992}{(*)}{巧了,这也正是那篇以“核心”自居、名为(何其讽刺!)“SGA~4½”的文本劝人千万别去读的两个讨论班……(5月29日)六种运算视野的分量,见注“无用的细节……”(n\textsuperscript{o}~170(v)),部分(b)(“一些什么也干不了的机器……”)。}。这样,凡我力所能及的,我都做了:为一件效力完美的工具、一个在
上同调主题里不断把我们径直引向关键问题的丰饶观点,恢复名誉(只要还找得到守着好工具的工
匠)。

我刚巡视过的这三处“撂荒的工地”(或房子,或工具……),牵涉到的,更是一种\textbf{共同的代数
语言}——用来表达最形形色色的几何情形——而不是某个个别的几何情形,比如代数簇的上同调。
第二处工地,我称之为“拓扑代数”的,有时会让我跟一些无疑深刻的问题擦肩(比如球面同伦群
的某些问题),那是偶然,不是本意。我在那里的主要动机,一如既往,是发展出普遍性、灵活性
都足够的代数工具,来服务于那门尚在襁褓中的\textbf{算术几何}:从Weil猜想那颗胚胎起,我花了
一生中十五年的好时光孕育它、接生它、喂养它。几何实质的精华正在这门几何里;这些年里,它
才是我同数学的那场恋爱真正的核心——今天依然如此。接下来还要巡视的三处“最灼热”的主题,
说的正是这精华。

\textbf{工地四:“系数问题”}。这个问题,在Weil猜想的表述本身里,已然萌动\fnote{993}{(*)}{关于这一点,见注“操演”(n\textsuperscript{o}~169)的开头,那里我评说了Weil猜想的初始问题域。(5月29日)这一开头已独立成注“‘Weil猜想’的语境”(n\textsuperscript{o}~169~(i))。}。整个六十年代,它是我在上同调里的兴趣中心。对当时已约略望见的几种主要系数类
型,它已被提得清清楚楚,带着全部必要的普遍性与精确性\fnote{994}{(*)}{似乎并没有出现什么新式的“系数类型”,超出我六十年代下半叶所预见的那些。}。这一问题域——对理解代数簇上同调显然至关重
要\dmark{}\origpage{1208}——我早在第一次回顾自己的作品、那道以敬意命名的注“我的孤儿们”(n\textsuperscript{o}~46)里就说
过,又在注“墓上的旋律——或自负”(n\textsuperscript{o}~167)里回到它。两条根本的主线:一条,是刚才说
过的六种运算与双对偶的形式主义;另一条,是需要在一个或普或泛的基底概形上,给出基底域上
已知的诸种“系数”的合宜推广——它们(哪怕只是默而不宣)进入了对该域上射影光滑概形范畴已
知的那些上同调函子的描述:$\ell$-进上同调、晶体上同调、德拉姆上同调,乃至($k=\mathbb{C}$ 时)
Betti 上同调或 Hodge 上同调。

我不认为这样说有什么过分:这一问题域已经以萌芽的形式\fnote{995}{(*)}{作这一说明,我丝毫无意贬低Deligne与Mebkhout上述贡献的原创性与重要性,正如我不会因为说算术几何“早已在Weil猜想中萌芽”,就去贬低我自己的贡献在这门几何的诞生与最初起飞中的原创性与重要性。}包含着那套至今仍在等待问世的“完全成熟”的“Hodge--Deligne理论”,也包含着同样仍在等待的“德拉姆--Mebkhout系数理论”\fnote{996}{(*)}{大体可以说,先是Deligne的贡献(1969年前后),后是Mebkhout的贡献(1975年以后),分别朝着两个很不一样的方向,回应了定义合宜的“德拉姆系数”的问题——这种系数将能把光滑概形的寻常德拉姆上同调嵌进一个六种运算的形式主义。Deligne只在概形 $\mathrm{Spec}(\mathbb{C})$ 上定义了一类“好”系数,并且对 $\mathbb{C}$ 上分离有限型概形的结构态射 $X\to\mathrm{Spec}(\mathbb{C})$、对 $X$ 上的常值系数(可惜!)定义了函子 $Rf_!$ 与 $Rf_*$。Mebkhout定义了一类“好”系数,原则上对特征零域 $K$ 上任意分离有限型的 $X$ 都成立——但他没有进而对这类概形之间的态射 $f:X\to Y$ 定义 $Rf_!$ 与 $Rf_*$,也没有对 $Rf_!$ 与 $Lf^!$ 发展对偶定理(除了 $Y=\mathrm{Spec}(K)$ 的情形——而且即便那时,也只在复解析簇的超越语境里,那无疑明显更难)。Mebkhout迄今发展的理论(应当说,是在一种令人丧气到极点的氛围里发展的)还有一个局限:它目前只对光滑的 $X$ 做成(我猜,是因为没有系统地使用晶体观点;晶体观点为微分算子层提供了一个满意的替代物,而后者在光滑情形是那样便利)。};而这两套理论之所以至今仍在襁褓之中,没有像 $\ell$-进系数理论(素数 $\ell$ 与特征互素)那样,一两年间便臻于成年,原因是同一个:我那班上同调学生,Deligne首当其冲,忙不迭地把宗师遗留的问题域连同宗师本人一道埋葬。

\pmark\origpage{1209}Deligne这边(他享尽了科学界宠儿的一切便利),与Mebkhout那边(他身陷的孤立再彻底不过——
而这孤立,恰恰是最该接纳他的那些人压在他身上的),所迈出的步子固然零碎,却仍不失为宝贵的
引线,让人得以提炼出某些关键的系数范畴。写那道已引的注“墓畔的旋律”时,这些重要的贡献就
在我心里。此后,我又更深地浸入了那套六十年代已然成形的“系数与 Motif 的瑜伽”,如今有了一幅
更精确、也更完整的图景。因此,我打算在《数学反思》第3卷里,接着六种运算形式主义的概要,回头
再谈系数问题(同时也谈 Motif 问题)。

现在只需说:我分明看见,在一个或普或泛的基底概形 $X$ 上,有本质上的三种“基本”系数类型\fnote{997}{(*)}{我说“基本”的系数类型,是想借这称呼提示:我眼下能窥见的其余各种重要的系数类型,都应当能借这几种来描述——或以合宜的方式把它们“组合”起来,或对它们添加合宜的结构增益,或两者并举。在为de~Rham--Mebkhout系数所设想的诸种结构增益之中,有(除了似乎“内在于”所论系数范畴的“权滤过”之外)一种“德拉姆滤套”,它在 Motif 式的应用中扮演头等角色。有可能,这层附加结构(从六种运算形式主义的眼光看),唯有当它与一种贝蒂式的“离散”结构相结合时,才有意义;这应当使我们得以表述这一滤套所须满足的好性质。关于这些问题,我打算在“Motif,吾爱”(《数学反思》第3卷)里更周详地回头来谈。}:$\ell$-进系数($\ell$ 为\textbf{任意}素数);\textbf{de~Rham--Mebkhout系数}\fnote{998}{(*)}{提醒一句:对这种de~Rham--Mebkhout型系数,我如今看见互为对偶的两种变体:Mebkhout的那一种,与我迟疑着不愿称作“Deligne的”那一种——尽管后者是一个被他拒认的孩子!(5月29日)关于这个被拒认的孩子,见注“……以及羁绊”(n\textsuperscript{o}~171(viii))。关于德拉姆系数的细节,见注“五张照片(晶体与D-模)”(n\textsuperscript{o}~171(ix))。}(主要对在基底概形 $S$ 上有限型的 $X$ 有意义,最\dmark{}\origpage{1210}重要的情形是 $S$ 为环 $\mathbb{Z}$、$\mathbb{Q}$或 $\mathbb{C}$ 的谱);最后,是 Betti 系数(对 $\mathbb{C}$ 上有限型的 $X$)。这三种之中,唯有第三种,在我看来眼下已经确定下来,不带任何假设成分。要定义第一种(哪怕只对绝对基底 $\mathbb{Z}$ 上有限型的 $X$),或描述它与第二种
的关系,关键配料在我看来是那套“神秘函子”的理论——我六十年代末就公设过它,这个问题,看
来也同其余的一切一道沉了下去……到时候,我自当在合适的地方更周详地回到它。

\textbf{工地五:Motif}。关于我朋友Pierre~Deligne对 Motif 的埋葬——蒙教团全体祝福——我已
经说得相当周详,此处不必再度铺陈。这里我倒想强调一件刚刚才向我显形的新事,它本该在十五
年、二十年前就显形的。一个月前,在一个或普或泛的基底概形(比方说 $\mathbb{Z}$ 上有限型,或
只取一个代数闭域的谱……)上“形状完整地”构造 Motif 范畴,在我看来还是一件分明“在地平
线上”的事,淹没在遥远将来的迷雾里。这种心境,无疑是那些已然遥远的日子的顽固遗产:当年
Motif 的思考起步时,基地再假设不过,那时连 $\ell$-进上同调的形式主义都还没有。对我而
言还有一条“从轻发落的情节”:1958到1970年间,那些近在眼前的奠基性写作任务,把我的精力吸
干了;我的 Motif 思考(以及关于其他在当时那些紧迫任务面前显得“奢侈”的主题的思考),一直
被压缩到一份可怜的口粮,那份口粮,我几乎怀着逃学孩子的负咎才拨给自己!不管怎样,我留在
脑子里的印象是:系数问题,是马上就该做的成熟的事(不过得靠别人,因为我已经忙在别处
……);而 Motif,眼下顶多配上一本“数学虚构”的书,要是哪天我抽得出闲,去写一写。说实在的,
倘若我真的一头扎进去写它,而不是耗在那些后来世上没有一个人肯接手的任务上——尽管人人都
乐于享用我做的东西……——事情的步调,想必很快就会改观……

\pmark\origpage{1211}话虽如此,我终于看清了这件事——它本身,一旦摆在眼前,原是明摆着的:只要你肯花功夫去描
述足够“细”的系数,也就是说,照顾到一个 Motif 相关联的全部已知结构,你最终描述的,就是\textbf{Motif
本身}。或者说得对些:你最终描述的是一个范畴,它将把(三角的)Motif 范畴作为一
个\textbf{全子范畴}包含在内(这已经很可观了)——正如复数域上的 Motif 范畴(若承认 Hodge 猜
想一个相当强的版本)作为 Hodge--Deligne结构范畴的全子范畴而出现。至于要确切刻画——用与手
头系数直接相配的“代数”语言——这个全子范畴\textbf{究竟}是什么,也就是说,\textbf{哪些}系数恰
好“是 Motif”,那就落进了恐怕要棘手得多的问题。它们关乎一个 Motif 相关联的种种几何-算术
结构之间的\textbf{相容性}(这些相容性,我想我在前面引过的注“墓畔的旋律”里已经影射过)。解决
这些问题的路数(它们在我看来,与“Motif 理论”的实际构造毫不相干),兴许真的“要等一百年”。
不管怎样,经验一次又一次让我们看到:这类关于一个问题多么“无从下手”的预言,没有多大意
思,顶多是在勇气本没有挂稳的地方,劝人泄气……

(4月1日)关于“\textbf{Motif 的Galois群}(或\textbf{基本群})”的形式主义,再添几句评注。这个
概念(1964年,在我有幸结识我未来的前学生Pierre~Deligne之前,我已经提炼出来并开始发展)带
来一些极其精确、极其精微的直觉和形式主义。它的存在及其要点,不依赖于人们为一个域上
的 Motif(或任意概形上的“光滑”Motif)的观念所选定的特定构造——只要那个构造满足几条合
理的条件。我曾把这样一项任务交给Neantro~Saavedra:把我大约1964年提炼的辞典整理成可
发表的形状,放进尽可能一般的语境——辞典的一边,是我当年称作“刚张量范畴”的那些范畴里
的几何($k$-线性范畴,带满足若干合宜条件的“张量积”运算,$k$ 此处是一个\textbf{域});另一
边,是 $k$ 上 pro-代数群的线性表示理论(更确切、\dmark{}\origpage{1212}更一般地说,是 $k$ 上的“pro-代数 gerbe”的表
示理论)。他把这项任务在他的论文里做到底,1972年由“数学讲义”出版(LN~265)
\fnote{999}{(*)}{(5月10日)这几行写下之后,我有机会读到了那本书——作者可没觉得有必要寄我一本。我看清了:书中,Saavedra俨然是那里所展现的新哲学的光辉发明人;他一丝不苟地遵照我交给他的那些笔记,而几乎不提我的名字(无论是书中引入的观念与关键结果,还是一些早已为人所知的观念,比如晶体、分层模或 Motif)。他为那个主要观念重新起的“塔纳卡范畴”这个名字,是一场天才绝伦的骗局——那名字,断不会比他所自命的这套理论,更是他自己的发明。何况这桩“父亲身份”本来就是暂时的:我的朋友皮埃尔已经在丛书问世十年之后,把必要的事情料理妥当,好叫它(如众人所期待的那样)回到那位早已内定的父亲名下。这场对一具尸身的辉煌手术的细节(在“SGA~4½--SGA~5”那场以同样不可模仿的手法进行的手术之前,它是第一场、也是唯一一场这种规模的),见那组注“第六根钉子(钉进棺材)”(n\textsuperscript{o}s~176$_{1}$至176$_{7}$)。}。这本辞典我当年还推进得更远(尤其涉及把滤套结构或分次结构等,翻译到某些纤维函子上,或涉及与塔纳卡范畴相关联的一个“极化”观念),这些在Saavedra的论文里都没有做\fnote{1000}{(*)}{(5月10日)这原来是一句错误的揣测。它出自我的一个信念:Saavedra断无能力“收尾”我给他指的那个计划——单是“pro-代数 gerbe 的线性表示”这一观点的驾驭,很长时间里似乎就超出了他;他的数学行囊又极其单薄。以Saavedra那点并不出众的本钱,从我离去(那时他对上同调、对代数群的构造一无所知)到书出版,不足两年,他竟有机会消化(而且消化得尽善尽美——书的成色即是明证)书中随手把玩的那一大堆四面八方的观念——这在我看来是不可设想的。关于这一点,见那组已引注“第六根钉子(钉进棺材)”里的注“Monsieur~Verdoux——或献殷勤的骑士”。},或者在“难忘文集”LN~900里都没有(在那里,Saavedra
的论文被重做了一遍,而 Motif Galois群的观念居于问题的中心——只是我的名字,在这个话题
上,同在所有其他关于 Motif 的话题上一样,一次也没有被提起)。

我还想指出:前面谈到过的那第一步——确定有限域上 Motif 范畴(差一个等价)——原本是对
该有限域的 Motif\dmark{}\origpage{1213} Galois群的确定\fnote{1001}{(*)}{(5月10日)我注意到,这一确定,也见于Saavedra那本取之不尽的书里(对我这卑不足道的人只字不提,这还用说么)。它用的是整体类域论的上同调理论(确定群 $H^2(\mathbb{Q},T)$,其中 $T$ 是 $\mathbb{Q}$ 上一个乘法型群)——所以,这也算我的前学生(其本钱俨然超乎常人)不到两年里就消化掉的东西之一……},该群应当是交换的(它被 Frobenius 元素拓扑地生成),实际
上是一个经某些 $\mathbb{Q}$ 上的 pro-代数环面扩张的 $\widehat{\mathbb{Z}}$(由 Frobenius 生成)\fnote{1002}{(*)}{这里说的是对\textbf{半单}Motif 分类的那一个 Motif Galois群。要得到一般的 Motif,还须用 $\mathbb{Q}$ 上的加法群 $\mathbb{G}_a$ 作它的乘积。}。第二步,是 $H^2(\mathbb{Q},T)$ 中那个元素的描述——按照Giraud的理论,它恰把纤维函
子的 $G$-gerbe 加以分类\fnote{1003}{(*)}{要点在于:这个类在所有 $\ell\neq p=\mathrm{car.}\,k$ 的位置上都变为零(多亏“$\ell$-进上同调”纤维函子的存在),而晶体纤维函子的存在,让我们对这一类在缺失的那个位置 $p$ 上的命运,有了足够的信息。}。

正如我在注“一个梦的回忆——或Motif的诞生……”(n\textsuperscript{o}~51)里所说,我是在寻找下述联系时撞
见 Motif Galois群的:对一个投射有限Galois群 $\mathrm{Gal}(\overline{K}/K)$,取变动的素数 $\ell$ 的各
个 $\ell$-进表示——例如取 $H^i(\overline{X},\mathbb{Q}_\ell)$ 得到的那些,其中 $X$ 是射影光滑概形,$i$ 为整
数(或它的某个合宜的子模)。Serre看的是Galois群在 $\mathrm{Aut}(V(\ell))$ 中的像——那是一个约化的
$\ell$-进李群——而这个像的结构(按李理论的意义)似乎与 $\ell$ 无关。正是在探寻这一现象(直到今
天,它本身仍属假设)的深层原因、把它与 Tate 猜想相联系的过程中,我发现了 Motif Galois群的
观念——紧接着“Motif”与“Motif 上同调”的观念之后。

如果说我曾把一件既简单又深刻的东西带到世上,如果说我的数学家生涯里有过创造之举,那正
是这条沟通几何与算术的关键观念的诞生。也正因为如此,去年那个难忘的4月19日,当看到这件东
西被人以那样一副气定神闲的做派据为己有——仿佛它是某人在一个技术段落的拐弯处,即兴哼出
的最末一句小玩意:瞧,简单得像剥白菜,只消把我们的那篇小文章里陈述塔纳卡范畴理论的命题
4.7.3 拿来用在这里就行——我给一股难以想象的\textbf{无耻}之气噎住了\fnote{1004}{(*)}{写这几行时,又逼着我想起一件十分相似的事:把一个 $\ell$-进层系数的 $L$-函数的定义引入,同样不提任何人,同样仿佛是信手拈来的最末一条老生常谈,由同一位光辉的作者即兴完成。关于这一点,见注“操演”的子注“……以及‘废话’”(n\textsuperscript{o}~169$_6$),打字稿第891页。}。七十年代里那些光辉的\dmark{}\origpage{1214}先例之后\fnote{1005}{(*)}{甚至早在六十年代——见“驱逐”一注(n\textsuperscript{o}~63)。},八十年代的数学,就是这样做的。

不过我扯远了,好吧——我本该是领人参观一处工地的,不是来抒情的。那就说明一下:同投射有限
基本群的情形一样,若 $X$ 是域 $k$ 上几何连通的概形,应当区分概形 $X$ 本身的 Motif 基本群,与
“\textbf{几何的}”Motif 基本群。两者并不重合,\textbf{哪怕} $k$ 代数闭——因为 $k$ 的 Motif 基本群
并不因此就是平凡的(它连通,仅此而已!)。所以要引入 $X$ 的“几何的”Motif 基本群;它的职责
(除了别的之外)是在与几何投射有限基本群 $\pi_1(\overline{X})$ 相关联的诸 $\ell$-进李群(作为其商)之间
建立联系。它的定义,是自然同态
$$\pi_1^{mot}(X)\to\pi_1^{mot}(\mathrm{Spec}(k))$$
的核(相对于 $X$ 上光滑 Motif 范畴上一个纤维函子的选取)。

我想要抵达的地方正在这里:这个核——不妨记作 $\pi_1^{mot}(X/k)$——应当是构造“$X$ 在 $k$ 上的
(几何的)Motif 同伦型”的第一步;这一同伦型,我前文已经顺带影射过\fnote{1006}{(*)}{见注“为一副模糊的骸骨作安魂曲”(n\textsuperscript{o}~165)。}。“这个同伦型”的形状完整的描述\fnote{1007}{(*)}{作为对象的类型,我预料它将是一个相对同伦型(按Illusie的意义),放在 $\mathrm{Spec}(K)$ 的 fpqc topos 的“扩张”topos(按Giraud的意义)里——该扩张对应着 $X$ 上光滑 Motif 范畴的纤维函子的(在这个 fpqc topos 上的)gerbe。这个同伦型在(刚才描述的)基底 topos 上的相对上同调,是拟凝聚的(甚至是“凝聚的”),并且可以等同于 $X$ 在 $K$ 上的 Motif 上同调。若用 $X$ 的一个复点(特征零的 $K$),取 Betti 纤维函子,则相应的纤维同伦型应当典范地同构于经超越途径与 $X\otimes_K\mathbb{C}$ 相关联的 $\mathbb{Q}$-同伦型(不计挠现象……)——至少当 $X\otimes_K\mathbb{C}$ 是 1-连通的时候。},它
的“上同调”应当不是别的,正是 $X$ 的 Motif 上同调;这属于“Motif”工地上一个有意思的概念性
工作,朝着一条分明不同的方向\dmark{}\origpage{1215}(而且很大程度上多半独立)——不同于那项中心任务:
Motif 范畴的实际构造,以及为它们建立六种运算的形式主义。

\textbf{工地六:标准猜想。}如我在前面一条脚注(打字稿第1202页脚注(*))里解释过的,这些猜
想可以在两种不同的意义上理解。首先是字面的意义,如我1967年在孟买讨论班上陈述它们的那样
\fnote{1008}{(*)}{Algebraic~Geometry,~Bombay~1968,~Oxford~University~Press~(1969)。}。照这个样子,在我看来,它们概括了代数闭链理论眼下最关键的问题——至少从这些闭链的所谓“同
调”等价的观点看是这样。

陈述这些猜想的时候,我的主要动机,其实并非朝向闭链本身,而是朝向它们(或许……)提供的
手段:在一域上建立一套半单 Motif 的理论,满足那些本该在十五年、二十年前就“人所共知”(至今
却仍讳莫如深……)的要求。我将在《数学反思》第3卷里指出这些猜想的种种减弱的变体——它们足
以建成这样一套理论(其中最弱的那个,几乎就是为了这个目的而必要且充分的)。如我在别处强
调过的:即便初始形态的猜想被证明在某个给定的域 $k$ 上成立(比如 $k$ 有限,甚至一切 $k$),这也
绝不意味着,应当称作“Motif 的”上同调类\fnote{1009}{(*)}{我想我可以提出一个合理的定义,界定射影光滑代数簇上的 Motif 上同调类——至少当基域特征零时如此。对一般情形,关键的(前文已经谈过的)是有限基域的情形。知道了后者之中 Motif 类的描述(且听下文分解),我想就能给出“那一个”Motif 类的恰切定义。可参照打字稿第1202页脚注(*)处的评注。}(人们可以指望这些类使 Hodge 猜想、Tate 猜想之类种种猜想
为真)就必然是代数的。倘若有一天发现存在非代数的 Motif 上同调类,那多半意味着:代数闭链在
Motif 理论中——也就是在代数簇上同调的算术-几何研究中——的分量,要比理论初创时我有理
由相信的更轻一些。话虽如此,我如今隐约望见的 Motif 理论的实际构造,先天地独立于关于
代数闭链的时下诸猜想(Hodge、Tate 或“标准”型)。

\pmark\origpage{1216}但这并不妨碍:一方面,标准猜想及其种种变体;另一方面,Hodge、Tate 猜想及其众多的变体
——这些猜想,尤其包含代数闭链的\textbf{存在性}命题(即上同调类的代数性),或者(在修改版
里)所谓“Motif 的”上同调类的存在性命题——彼此紧密相联,并且同主要“系数类型”的描述、乃
至同 Motif 范畴本身的描述,紧紧相连\fnote{1010}{(*)}{这并不与我刚才的断言相抵触——我说过,我隐约望见的 Motif 范畴(比方说一个域上的)构造,与所论的种种猜想“独立”(即“技术上”或“逻辑上”的独立)。我说的那些“紧密的联系”(比如,我在 Hodge、Tate 型猜想中看到的十二种变体,恰好提示了同样多的不同类型的上同调“系数”),属于启发而非技术——正如(所谓“猜想性的”)Lefschetz--Verdier公式与 Frobenius 对应的迹公式之间的联系一样。在后一情形,这个根本的启发式联系(并非逻辑依赖的联系),已在注“操演”的两条子注“\textbf{真正的}数学……”、“……以及‘废话’”(n\textsuperscript{o}s~169$_5$、169$_6$)里得到郑重的指明。}。

这里,又是一件近二十年来早该做的澄清、归置、通报的活儿,没有被做(尤其,没有公之于众),
做的人却宁愿到今天还在埋葬丰饶的观念(当它们还不曾发表),或者拆毁它们(当它们发表了),
好独占(眼前的)利益与(日后的)功劳,而不肯把这些迷人的问题通报出来、交到大家手里——
这些问题,对我们理解几何、拓扑与算术之间的纽带,至关紧要。我看出,这里欠缺的,绝不是能
力,甚至不是光辉的天分,而是一种朴素的诚实,以及在对一个派发声望与权力的“科学共同体”的
关系里,某种起码的\textbf{廉耻}——在那些并不因此觉得自己欠下丝毫义务、欠下丝毫“服务”式“回
报”的人身上。所以,虽然我同这个题目断了往来已十五年以上,虽然我已经说不上“局中”二字,
费力重新入水的却还是我:至少要在《数学反思》第3卷里,把我那些更年轻、更有天分的人漏下的,
尽力修补;到底,把他们没有慷慨去做的事做了。

到这里,我相信我已经把这些“工地”巡视了一遍——从建造那门“算术几何”的眼光看(这门几何
的基石,是我整个六十年代铺下的),它们如今(而且自我告别数学舞台那一刻起)正是“\dmark{}\origpage{1217}最灼热”
的。我绝无意说,我已经把那些实质问题——那些也许只有我看得见、而又放在我心
上的问题——粗粗走遍了。据我所知,它们仍停留在我离去时的原地,其中许多甚至还没来得及在
文献里得到显式的陈述。其中我要特别指出\textbf{离散Riemann--Roch猜想}(概形框架下的)\fnote{1011}{(*)}{这一猜想似乎第一次被显式陈述,是在那道名字意味深长的注“屠杀”的子注 n\textsuperscript{o}~87$_1$ 里——因为这条猜想,本就是 SGA~5 遭屠杀的诸事项之一,在Illusie的编订版里消失得连一个\textbf{名字}都没留下。}。同样,还有把\textbf{局部与几何整体类域论}推广为一条本质
上“几何”性质的\textbf{对偶}命题(经典“算术”命题作为推论给出)的工作。1976年致Larry~Breen的几封信里谈到了它,那几封信将作为附录收入《追寻栈》第一章(因此会出在
《数学反思》第4卷)。在这些命题里,主要的工作将在于对人们所处理的“系数”范畴的细心描述。其中
一个重要角色,由Serre发现的一种自对偶性扮演\fnote{1012}{(*)}{除了Serre这个漂亮的念头之外,影响我的,还有 Lang 在几何整体类域论中、Serre在局部类域论中引入的“几何”观点。}:特征 $p>0$ 的域 $k$ 上、幂幺代数群范畴里(差根同源不计)的一种自对偶(在我看来,除了当年听我讲过的屈指可数的几个人,至今无人知晓)。把这类命题推广到更高维的问题,(至少对我)是个彻底的谜(不过Milne在代数曲面的情形,或许有些眉目……)。

这些对偶问题,我记得可以上溯到五十年代末。那时我还顺手开创了一个与之相联的构造:与一个
域上有限型概形(先从这里开始……)相关联的、pro-代数群的“\textbf{雅可比}”(链)\textbf{复形}——
借与那些局部环相关联的合宜的“局部雅可比簇”,类比于我几年前在凝聚对偶里构造过的“残差”
或“对偶化”复形。所有这些对偶问题,在六十年代,都被推到了第二线,尤其是为了
\dmark{}\origpage{1218}发展平展与 $\ell$-进上同调的“废话”和 topos 的语言。我纲领里相对局部与整体雅可比的那一小段,约在1977年由Contou-Carrère完成(没有提我这卑不足道的人);而一见到Deligne与Raynaud对它摆出的脸色,他就赶忙收了起来\fnote{1013}{(*)}{见注“棺材3——或有点太相对的Jacobi簇”(n\textsuperscript{o}~95),及其子注 n\textsuperscript{o}~95$^{1}$,其中谈到Contou-Carrère在数学大世界里的几段倒霉遭遇。}。如今,要重新拾起、发展那些打着我太明显印记的观念(哪怕当事人力图把这印记
藏得最好),是需要一点勇气的。唯一死磕到底的,是Zoghman~Mebkhout;而他遭到的命运——在反常
层会议上登峰造极的那些“绝技”——足够清楚地显示了这里头的风险。

倘若我想要开列一张1955到1970年间我发现的漂亮问题的清单(它们当中的一些,当年我曾在周
围这里那里谈论过),我还得再花上几天;若想稍稍说得明白些,把来龙去脉都交代了,恐怕要几
个星期。这里不是做这件事的地方,我也怀疑我永远会去做。更何况,倘若我盼着有一天(谁知道
呢!)哪位年轻数学家肯在其中一个问题上扑进去,练练手,再让自己为人所知,那还是让他自己
重新发现的好——犯不上冒险,让人给贴上某种标签。

当心将来等着我们的那些反常层会议……

% ================================================================
%  18.7 Les fruits du soir (原书 p.908–916)
%  源文: _work/chunks/chunk14.md
% ================================================================
\section{(7) 暮熟的果实}

\subsection{(1) 尊重}   % 18.7.1

\notehead{179}{}

(4月2日)已经五个星期了(自2月26日那道“沉默”之注起——它揭开了集结在“四种运算”名目
之下的那组笔记),我一直在重新巡视有关\emph{葬礼}的种种“物质性”的、或者(多少)“技术性”
的事实。在“四种运算”里,我把自己限定在严格意义下的“诈骗”这一面——也就是越过了那道
“门槛”的一面;那道门槛,同名的注(n\textsuperscript{o}~172)里谈过,它分开\textbf{恶劣的心性}(表现为
“自动排斥”的反射,往往不顾一个数学家最起码的直觉)与公然的\textbf{蓄意刁难}、与铁板钉钉的剽
窃。在刚写完的“荒废的工地”里,我面对的主要是\emph{葬礼}的“第一层次”,门槛以下的那一面:
一场广阔视野\dmark{}\origpage{1219}与一些强有力执念观念的埋葬——当然,没有谁非得接着做它们不可,人人都有权
无视或遗忘它们——只是这样一来,也就“把自己埋葬了”,让自己的工作(至少让那部分直接蒙受
被拒视野的工作)或多或少归于不育。

在那里,我总算有了走完一遍的印象!就“工地巡视”(撂荒的那些)而言,它让我对我著作的\emph{葬
礼}有了一份原原本本的把握,同时又让我同十五年来失了音讯的一些主题,多少重新接上了头。这
尤其让我看清了:在接下来几卷《数学反思》里我打算白纸黑字写下的东西,孰急孰缓。我的宗旨,当
然不再是给一门临盆中的科学打细致的地基——这种事我做得够多了;倘若再没有旁人肯像我当年
那样,把自己交给这样一项任务,那就各人自便,大家都自便!我的宗旨毋宁是:提炼某些执念观
念——它们服务于一个1955到1970年间诞生的总体视野;而今天我重新找到它们(多亏我某些昔日学生的努力,又蒙全体默认)时,它们或被遗忘,或沦为笑柄,或被肆无忌惮地据为己有、
残害得失去大半元气。今天重拾它们,我终于松开了缰绳,放开心中那股求知的冲动——六十年代
里,为了没完没了的“服务”任务,我常常只给它一份可怜的口粮。那个时代一去不返了——然而我
知道,在我数学激情的这个新阶段里,服务的冲动,丝毫不减当年。我依旧要为那个由渴望知道的
头脑结成的理想“共同体”服务\fnote{1014}{(*)}{关于“数学共同体”,我第一次表明看法,是在《收获与播种》第一部分“‘数学共同体’:虚构与现实”一节(n\textsuperscript{o}~10)里。在这里我援引一个“由渴望知道的头脑结成的理想共同体”,看上去似乎我又退回到某种东西上去了——它的虚构性质,在所引那一节里曾显示得清清楚楚。但在《愚妄与更新》第八部分里,我有生以来第一次(迟做总比不做好……)确认了我自己的“求知历险”在数学层面上的一个集体维度。(关于这一点,见“孤独的历险”与“一段过去的重量”两节,n\textsuperscript{o}s~47、50,尤其打字稿第134、135页。)同样清楚的是:经历这场集体历险的那个“共同体”(或“集体”),与任何社会学意义上的实体——它体现于特定“时代”的特定“圈子”,带着特定的“心态”,或者(今天)带着特定的权力结构与阶级利益——完全是两回事。我所援引的这个“理想共同体”,“在空间与时间上都无边界”,对我而言丝毫不比社会学实体更不“真实”。它甚至更其根本,因为正是它(如我在同一句话的后文里所写)“继续给我的数学投入以一种比个人消遣、比自我扩张手段更深的意义”。它并不比我自己更“虚构”——我自觉是它的一分子,而且比当年更清醒。“虚构”不在于觉察到这样一个“共同体”的存在,而在于把它同一个我曾自认同属的圈子混为一谈。},它继
续给我的数学投入,以一种比个人消遣、比自我扩张的手段更深的意义。

\pmark\origpage{1220}在这些投入里,当然,“老板”也并没有比当年更缺席。面对那些在我眼里曾是数学世界里
“我的亲近者”的人的恶意与嘲弄,面对那些我爱过、也曾无条件信任的人对我最起码体面感的屡
屡伤害,我心中涌起一股按捺不住的冲动:在那些丢失了敬意的人面前,\textbf{以对这些活生生的、
茁壮而美丽的东西的敬意,来见证我对自己的敬意}——这些东西,是我亲手带到光天化日之下
的。也许,见证这份敬意的最好方式,就是在这段还归于我的宝贵岁月里,做这些东西的仆
人。因此,未来几年里,我打算在《数学反思》的后续各卷中发展的数学思考,将既是重新拾起一个\textbf{孩
童的游戏},又是一件\textbf{服务的礼物},又是一个\textbf{敬意的举动}。

在给\emph{葬礼}画上句号之前,我还想在“物质事实”之外,就这场思考教给我的东西,作一小结。
我先看它教给我的关于他人的事,最后再谈它教给我的关于我自己的事。

在整个思考过程中大白于天下的种种事实里,至今最触目的,仍然是七十年代、八十年代数学世界
里\textbf{风气与人心的败坏}。这种败坏,借成百上千桩“小节”表现出来,如同过去八九年间,一阵阵
回到我耳边的那些“小节”——它们却足以让人错愕,足以催生《收获与播种》第一部分的思考及
其根本追问:事情是怎么(又是何时)走到这一步的?而在这场我今天目睹的阴险而无情的败坏之
中,我的角色是什么,我的位置又在哪里?

\pmark\origpage{1221}这场败坏,在“SGA~4½--SGA~5”那样的行动里,或者在(更难以置信的)反常层会议那场
行动里,达到了顶点——其犬儒与轻蔑,远远超出了我所能想象的一切;哪怕就在我身不由己地发现
它们的前一天,也想象不到。

在这里,不适合回头去谈那些“小节”(其中不少,已在思考途中顺带指出),也不适合谈那些大行
动(它们自有小伎俩伺候)。显露在两者——小节与大骗局——之中的,是同一种精神。在可容忍与
下作之间有时不妨画一道“门槛”,这道门槛本身就相当脆弱、相当人为,是一根护栏;反正(看样
子)已经没有人在乎它了。经由这场把我本人牵连至深的\emph{葬礼},我得以也许比以往任何时候都
更近地看清这种精神——它当然既不是这场葬礼(为我这卑不足道的人发起)的专利,也不是数学
家世界的专利。对此我能说的只是:我不知道,在我们这个时代之外的任何时代,这种精神曾统治
过那个世界,或其他任何一门科学。这大概是众多征兆之一,表明一种文明——以及那无论如何还
在给它赋予意义的东西——分解到了终局。

\pmark 近来,我的思绪不止一次地流连于一个奇怪的巧合:十五年前我告别数学舞台,是受了科学
界某种腐败的冲击——那种腐败,我曾长期选择闭目不见(自以为离它远着呢)。而我是在那所我
本打算终老于此的机构里,猝然同它打了个照面\fnote{1015}{(*)}{关于这一点,见注“有益的挣脱”(n\textsuperscript{o}~42),以及“兄弟与夫妻——或双重署名”及其子注(n\textsuperscript{o}s~134、134$_{1}$)。}。在那里,事关科学家们与军事机器之间那几近普遍的、有利可图的
串通。军方对整个科学世界的这种阴险把持,同样是个新现象,只是(至少就我们今日所知的规模而
言)自第二次世界大战以来才出现。当然,那场“冲击”之所以搅乱了我预定的轨迹(我自己、人人
都预定的轨迹),以至于触发我无可挽回地离开那个我一直自认同属的世界(除了一项默而不宣的
保留……),是因为我心中有一种紧迫而急切的更新之需——这一点,我是事后回望才意识到的
\origpage{1222}。后来,我曾倾向于把触发那次不寻常离开的特殊诱因说得轻些。然而我也知道,那些惯性之力
有多么巨大(同时又是多么无形):它们要把我们永远拴在同一条“轨迹”上,敌视内在的更新——
这也让我掂量出,把我从一条划得那样牢固的轨迹上拔出来,需要何等猛烈的内心震撼。

我想说的是:触发我离开的那桩“特殊诱因”,肯定不是没有\textbf{意义}。至少在我离开后的头几个
月,多半甚至头一年里,这个意义强烈地在场。后来,在新印象的涌入之中、在这场头一波汹涌更新
的动势里,这个意义自然地退居幕后,终于从我眼前消失。可是,即使我不再感知我过去或当下的行
为及其果实的某个“意义”,这个意义并不因此消失。而我重返数学活动——它不可避免地带来同我
离弃的那个世界的更多接触——又冷不丁把我带回那段被遗忘的过去。因为这次“回归”(同当年那
次离开一样出人意料……)最初的收获之一,就是在我那个曾经属于自己的世界里,发现了另
一场腐败——一种据我所知它从未有过的腐败。若要给这件新东西起个名字,涌到我嘴边的是:\textbf{敬
意的丧失}。过去几年里,我不止一次痛切地感到它:当我看到“我昔日所爱的某人,不动声色地踩
压我如今所爱的另一个人——在那人身上,他本该认出我”。在对\emph{葬礼}的思考里,我又一次次重
新遇到它,而且语气更凶,这一回是冲着我亲手生下的某些东西,或者冲着某个胆敢从中汲取灵感的
继承者。在这些时刻,我真正领教了这种丢了敬意的精神的“气息”与“气味”。可我也明白,这种精
神“并不只在我家四邻吹拂”——尽管,恰恰是它吹到我身上、吹到我所爱的人身上,我才真正“认
识”它——正如人们只有亲口尝过,才知道苦果的味道。今天,这种精神已经成了时代精神……

我看得很清楚:这两场腐败——触发我离开的那场,与我“归来”时等着我的那场——并非没有关
联。若要用言词勾勒这层关联的隐隐感受,我要说:\dmark{}\origpage{1223}在科学家们面对军方的金钱(只提这一面)及
其便利时的那种顺手就沾的态度里,我察觉到一种对自己的不敬——个人层面与集体层面皆然
\fnote{1016}{(*)}{在这里,我冒着得罪某些昔日朋友的风险——他们把这种“顺手就沾的态度”当成自己的,却当然不觉得自己亏欠了对自己
的敬意!再说,也绝非确定:别的时代的科学家们,倘若集体地面对同量级的“诱惑”,反应就会不同。常言道,机会造就窃贼!}。而正是在对自己的敬意的丧失里,我认出了对他人的敬意、对出自
他人之手或造物之手的作品的敬意,所以丧失的根由。

我并不声称“理解”了这两场腐败中的任何一场。一方面,有“时代精神”——它的独特动势,在我
看来几乎完全脱出个人行动所能及。这股集体动势,对我始终是个彻底的谜,我从未想过要去探测
它。另一方面,是每一个具体的存在——秉着他的知觉与创造之能,驮着他那全部特殊条件作用的重
负——回应这时代精神的方式;他(自觉或不自觉地)把这种回应,做成了他个人历险的关键要素之
一。

在整个思考过程中,我花了很久,试图去把握\emph{葬礼}的两大主角各自的选择,以及这些选择背
后的力:死者,与葬礼主祭\fnote{1017}{(**)}{(6月22日)到了“最后一刻”,在“家庭相册”一注(n\textsuperscript{o}~173)的 c.(“众人之上的那一位——或首肯”)、d. 与 e. 三部分里,第三个“主角”终于向我显形。}。确定无疑的是,一路走来我学到了东
西;但绝不能说我完成了任务。我甚至可以说,就我的主角而言,我肯定\textbf{没有}完全成功。我收集
拼图的碎片,把它们拼了起来,我甚至相信碎片是对的、拼法大体不错——但对\textbf{整体}的认识,
始终缺着。它至今还是一堆碎片,而这些碎片,对我而言仍是\textbf{陌生的}——陌生于我这个人、我
的经历,因此也就是不可理解的。已做的功课,无疑会有助于我在别的场合多少认路,在该留神的地
方留神(年纪越大,我越发觉要留神的地方真多……)。但这一切,离真正的理解还差得远。于是我
心里冒出一个问题:往这个方向下的功夫,最终是不是一场障眼戏——或者至少,那个\textbf{目标}(在
某种冲突情境里“理解他人”)是不是\dmark{}\origpage{1224}一场障眼戏(虽然走过的\textbf{路}倒是教益丰富……)。我对自己说:真正理解\textbf{那个人}身上的冲突(或者任何一个曾与我亲密相系、而我看
到同样矛盾炸开的人身上的冲突),大概同时也就是\textbf{理解冲突本身}。而且我深知,这样的理解,
不可能来自对他人的冥想(他人永远逃脱我直接的认知),只能来自对我自身的冥想。“阴阳之钥”
这场长考若要显出它的丰饶,靠的不是对他人其人的偶然瞥视,而恰恰是回头看我自己的生命、我
自己的经历,以及我当年对它们的理解。

\subsection{馈赠}   % 18.7.2. Le don

\notehead{180}{}

(4月3日)说到底,我并不觉得有兴致,用几行或几页文字,为\emph{葬礼}中我的主角向我显形的一
切,作一番回顾。照现在的情形,这恐怕不过是一场文体的练习,而对那最残缺不全的理解的更新,
并无多少助益。眼下,我倒是急于抵达这场关于\emph{葬礼}的思考的句号!

当然,我也知道,这个句号绝不等于\emph{葬礼}本身的终结。接下来的几个月,随着这些孤独的果
实般的笔记招来形形色色的回声,想必会有孤独的思考给不了的惊喜与教益。也说不定,来到我这里
的惊喜并不都是苦味的;也许不远的将来还给我备着某种欢乐——正因为它稀罕,才愈发珍贵;正
如仅仅去年(丰收之年!),我就得到过欢乐:收到旧日同事与朋友中的三位——我素来特别敬重或
亲近的——写来的满纸热忱的信\fnote{1018}{(*)}{这里说的是 D.~Mumford、I.M.~Gelfand与 J.~Murre的来信。}。

至于《收获与播种》对数学世界的“时代精神”会有什么全局效应——无论多么微末——几乎不必
说,我对这一点没有半点幻想。至多,这套笔记的出版,或许能叫某桩空前的横逆收场,叫某桩太过刺
眼的不正常得以纠正——而且,我这么说也许还是乐观了。也可能,死者本人的不期复现——他被当
作死了多少年、都馊了——会给那团裹着他双手之作的、压低嗓门的嘲弄合唱,添一个\dmark{}\origpage{1225}终止,或至少
添一个更知趣的弱音器。而倘若这次复现不同时终结那场对一种视野、对一些强而丰
饶的观念的时髦抵制,那它或许至少能打动某一位比别人更大度的年轻数学家,让他毫无保留地从中
汲取灵感(不惜惹人嫌弃),怀着敬意把它们化为己有。

然而,我写《收获与播种》,不是为了这些事里的任何一件——它们当中有的兴许会附带着来,谁
知道呢!我写它,是“为我自己”的,当然,如同我写的一切——作为一场摸索着寻找自身的理解的
手段。但与此同时,对他人的念想——对那些我爱过、而我的历险把我带往\textbf{别处}的那天撇下的
人的念想——贯穿了《收获与播种》写作的始终\fnote{1019}{(*)}{这一念想,在《愚妄与更新》(《收获与播种》的第一部分)里不止一次得到表达。它兴许在后文各部分里不那么显眼,却同样在场。}。这些笔记,既是一场思考,有时也是一场冥
想,对我而言,曾经并且始终是一份\textbf{馈赠}——在我自己之外,送给那些我向他说话的人。当然
我也知道,这份馈赠,也许除了我自己,无人会接。但我并不因此后悔写了它。再说,今天若没有被
某个它属意的人接住,也许明天会被接住。这份既自发、又经长久酝酿的见证,其中每一页、每一个
字都来在其时、来在其所——它明天不会比今天更不真。而无论今天还是明天,若有出乎意料的、
被欣喜迎接的事,那将是得悉:我的馈赠被接住了,哪怕只被一个人——他在我之中认出了他自己……

\subsection{(3) 信使(2)}   % 18.7.3

\notehead{181}{}

比起\emph{葬礼}图景的“前景”,我更不觉得有兴致,就另两个“平面”——一个由“我那班生意兴
隆的学生组成的、扛着铁锹绳索的队伍”,一个由“教团全体”组成——对我在他们身上的所悟与所
惑,作一番原原本本的回顾。关于后者,以及它在\emph{葬礼}里的角色,我在注“掘墓人——或教团全
体”(n\textsuperscript{o}~97)里已经说得相当周详\fnote{1020}{(**)}{(6月22日)我对教团尚且模糊的知觉,最近在已引的注“家庭相册”(n\textsuperscript{o}~173)的 c.、d.、e. 三部分里,出人意料地变得具体了。}。至于我对我那些亲爱的前学生的角色与动机的困惑,它们最清楚地显
现在注“沉默”(n\textsuperscript{o}~84)里,却\dmark{}\origpage{1226}没有在思考的任何后续时刻被认真重审。所以,恰恰
是在这个层次——\emph{葬礼}图景的“第二平面”——我的功课最欠火候!\fnote{1021}{(*)}{(6月22日)不过,关于图景“第二平面”的思考,确有一份(微薄的)续篇:6月19日的注“一场屠杀中的五篇论文——或孝道”(n\textsuperscript{o}~176$_{7}$)。}。在那里,没有一桩功课可与我在
“掘墓人……”那条注里做的相比。图景的这一部分,后来在两道更晚的注里、在阴阳动势的光照下
得到了深化:“天赐的机缘——或封神”与“拒认(一)——或召回”(n\textsuperscript{o}s~151、152)。

“掘墓人——或教团全体”这条注,是思考“第一口气”里写下的最后一道,想必也是它的顶点。可
是,时隔将近一年,我已经不再相信:在我这卑不足道之人的\emph{葬礼}背后(它被看作“对异端的
报复”),那股显得相当明显的集体动机,在集体意志的层面上,真正触到了\emph{葬礼}的\textbf{神
经}。使我生疑的是:在我每一个学生的情形里,这股动机在我看来完全缺席,或者,同其他在场的
力量相比,无关痛痒\fnote{1022}{(**)}{这一事实,出现在注“天鹅绒的脚掌——或微笑”(n\textsuperscript{o}~137)的思考里,打字稿第644--645页。}。可是,\emph{葬礼}里最触目的事实之一,恰恰是它的三个前后相续的“平面”之间
存在的那种“全体一致”:它们的作为与不作为互相衔接、互相补足(仿佛由一个“天衣无缝”的共
同意志指挥着),同一场名副其实的葬礼上一般无二!在如此显著的全体一致里,在心境与作为的如
此整齐划一里,人们当能猜到一个共同的动机,一根同一的“神经”,在驱动着这一个与那一个。

我无意暗示:那股我这里那里察觉到的“弥漫的怨怼”——它起于我的“异端”,(浅浅地)被感受
为开小差,(更深地)被感受为不可容忍的挑衅——全然子虚乌有、全然不起作用。但我如今怀疑,
它的作用是否是决定性的;怀疑它是否就是那根共同的“神经”——难道它人人皆有,却偏偏不
在那些\textbf{在葬礼里戏份最重的人}身上!也就是说,不在我那些学生身上——他们恰恰是那
份遗产的第一批受托人……

这桩“缘由”(我那场“异端”——表面看还相当讲理)在我看来,同那股暴戾之气完全不成比
例——那股气,\dmark{}\origpage{1227}我是在“屠杀一个‘光辉讨论班’”那样的行动里、在教团默许的目光下感到的;同反常层会议上那场当众鼓掌的、同样凶暴的不义,也不成比例。这也不是说,我曾是一
个可憎的同事或老板,怕我怕到积攒的敌意不敢发作、非得等我死了葬了才一齐倾泻到我头上、倾泻
到人们多少“认出我”的那些人头上。从我这里那里收到的种种回声看,既没有\textbf{畏惧}的痕迹
——仿佛我这人曾叫人害怕,如今报了迟到的一箭之仇\fnote{1023}{(*)}{的确,在“愚妄与更新”里,我长篇谈过那种“知名人士”身边——从一个我也说不清何时开始的时刻——笼罩着的\textbf{畏惧},我自己也时或在我周围察觉它的迹象。但那是附着于“知名度”本身的弥散的畏惧,而不是对我这个人本身的——只要稍能建立一点个人接触,它就消失了。我印象里,在个人接触的层面上,人们与其说怕我,不如说把我看作“好脾气”。我确信,在那位“失手——或二十年之后”一节(n\textsuperscript{o}~27)里谈到的学生那里,情形并无不同:每次见面,某种“发怵”持续了很久。这种发怵,今天我看作一种弥漫性的内在不安(“Unsicherheit”)的征兆,它后来在对支配与轻蔑的姿态里找到了补偿与出口。在他众多的学生当中,我有机会认识的三个,个个都结结实实吃过他那副貌似“无缘无故”的恶意姿态的苦头。显然,那种安营扎寨、统治着数学圈子各处的精神,助长了这类悖谬的行径;而这些行径又回过头来塑造这种精神,给它打上那种叫人错愕的、闷声的野蛮印记……};也没有什么稍微具体的行迹可以
归\textbf{罪}于我、足以滋养某种敌意或暴戾(而那种敌意,却从来不报姓名)。

这正是我称作“无缘无故”或“无因”的暴力的典型情境。这种暴力,最终来到“阴阳之钥”这场长
考(它本身构成《收获与播种》的心脏)关注的中心,想来绝非偶然。这种暴力,我认识它远非昨
日始;我头一次同它打照面,并不在我的数学生涯里。即便我有时忘了它在人间世界的存在,那也
从来长不了:它自己会很快前来提醒我,好意不忘。至于当下——借着一个奇怪的“巧合”(我
承认,它常常不合时宜,或至少不受欢迎……),我不记得我此生曾在什么时候,同这样一种暴力的
\dmark{}\origpage{1228}熟悉的征兆打过这么执拗、这么重复、这么磨人的交道——自从我“重返数学”以来,尤其自写《收
获与播种》以来;而最近的这几个月、这几个星期,愈发如此。

\pmark 当然,这里有一个执拗的讯息,它一遍遍回来,而且想必会一直回来,直到被听见。在阴阳长考
的最后几个星期里,我开始侧耳倾听它——同时知道,它要对我说的,我还没有听完。此后两个月
里,一种地下的功课想必在无声中进行。我觉得,那\textbf{要紧而隐秘}\fnote{1024}{(*)}{写下这一行时,我心里清楚,“隐秘”一词在这里是个权宜之计,是对“共识”的让步。我常常有机会发现:当我得知一件我一生都不知道的事时,那件事绝非“隐秘”,恰恰相反,它明摆着、显而易见,有时甚至刺眼——只是我不肯看它而已。发现新事物的情形,大抵如此,无论那是数学的工作,还是认识自己的工作。这种失明的根由、这种常识或直觉官能的阻滞,绝不在这些官能的缺陷,而在于精神的一种近乎不可克服的惰性:它不肯离开共识已然铺好的车辙——无论那是全社会接纳的共识,还是自己所处的那个较小的圈子里的共识,甚至,它们只在我们内心缔结并加盖了印章,如同某位“老板”为他一己之便、同他自己缔结的条约条款……}的东西,开始从表面可见的(或
至少不那么难于承认的)杂事里沉淀出来。“侏儒与巨人”的形象(我朋友皮埃尔供认的)继续萦绕
着我。在这形象背后,我相信辨出一股力量可观的原型:它像童年所受压抑投下的影子,或者
影子之一。它的角色,该是为创造之力所受的压抑提供一个出口与补偿——那压抑,早已内化为那句
“说不出口的无能之信”……在这个预感到的原型里,我仿佛触到一股强大的引擎,驱动着
那些无端施暴的行径,打击那个被当作“巨人”、被当作完好无损之力的携带者的人——这些行径的
爆发,不需要别的“缘由”,只需要一个\textbf{合宜的时机},当下手的风险显得为零、或微不足道。

也许我已经说得太多,而这几行字,才刚刚拂过一道纤细而执拗的直觉——它向我示意:有一桩功
课要做,它还横在我的前头。对这桩功课而言,\emph{葬礼}只是材料之一,同许多别的材料一道,来自
我那所谓“私人”的生活\dmark{}\origpage{1229}。这里不是继续它、甚至不是接近它的地方。它的位置,不在这些
注定要发表的笔记里。

\subsection{(4) 失乐园(2)}   % 18.7.4

\notehead{182}{}

(4月4日)在这场已许诺的回顾——我的思考教给我关于他人的事——里,我的思绪不由自主地
执拗地回到我自己的身上。这对我是个好兆头——它标志着我心中那股回到要紧之事的强烈需要。
对我的认识,额外地带来对他人的理解,而不是相反。自我开始冥想以来,不止一次,“理解他人”
的心思,成了那桩要紧功课——认识我自己——的消遣与岔路。

在有意地回到我自己之前(逆着我那股要抵达那个有名“句号”的不耐!),我还想收入一份最近到
我手里的、关于我朋友皮埃尔的见证。这是我离开数学舞台以来,所知的此类见证中的孤例。它给我
朋友的照明,与我通过别的途径所知的,大异其趣。它也十分适时地再次提醒我:现实,永远比我步
履蹒跚地试图给它画的像更复杂、更丰富\fnote{1025}{(*)}{我绝无意借此暗示:人们(包括我自己不停地)为现实画一幅尽可能“忠实”的像、并随各种“信息”的到来不断调整这幅像的努力,是徒劳而无效的。恰恰相反,这里有一种效力极大的辩证法,让我们同现实接触、去“认识”现实。只有当这幅像(它按事物的本性,自带一份惰性)完全僵死、板结之时,它才成为把握现实的障碍,或者更确切地说,成为让我们的把握官能\textbf{失效}、把我们确实拥有的对现实的知识“蒸发”掉的\textbf{有效手段}。}。

这份见证不是直接的。它关乎一位外国数学家同Deligne的一次(或多或少偶然的)会面留下的印象;
那位同事把这次会面对我的通信人讲了(我猜还是热乎的),通信人又在一封信里把经过转述给我。
经我的通信人与那位讲述会面的同事(下文称他“Z”)同意,我在这里翻译信中涉及这次会面的部分。
我的通信人推测,这一幕应当发生在1981年。(注意:这也正是反常层会议之年——不过通信人与我
之间,倒不曾谈起过那场讨论会。)

\begin{quote}
\pmark\origpage{1230}“……有一天,Z.~去 Bures 参加一个会,发现自己置身一间上茶的屋子[看
来是法国高等科学研究所的‘茶室’],屋里有很多数学家。只见门开了,Deligne走进来。Z.~先生把
这一幕讲得相当生动:他看上去蔫蔫的,双臂低垂,人觉出他周身有一种孤立。别的人都定定地看
着他,有点像看一只珍禽,又没人知道能对他说点什么。Z.~靠窗独自坐着,Deligne有些迟疑,在他
身旁坐了下来。Z.~不知道跟他说什么好。这时他起了个念头,就简单地说:他觉得围绕‘平展拓扑’
等等的那一整套想法,还有您带来的那些新思想,非同寻常。[“您”,这里及下文,指的都是我,Grothendieck——我的通信人是转述给我的。]Deligne的眼睛顿时亮了,他对 Z.~说:是啊,那是数学里
最好的东西之一;听您的报告\fnote{1026}{(*)}{同前,这里的“您”指的就是我,Grothendieck。},那才叫美……他接着说:您想想这个,再想想那个……数出一大
串事情来——Z.~自己(据他向我说的)一窍不通,但他看见了谈话对方身上骤然涌起的热情。Deligne又说:多可惜,您\textsuperscript{(*)}退了!他确信,晶体上同调和许多别的东西,不至于落到现在这副半
生不熟的样子;要是您\textsuperscript{(*)}真还接着干,它们如今早就是站得稳稳的构造了,就像平展上同调
那样……”
\end{quote}

这个故事里有两件事击中了我。其一是那种孤立的印象——它似乎深深击中了 Z.~先生。这一印象究
竟出自Deligne生涯里某个特别的时刻,还是这样的孤立终于浸透了他同所有同侪的关系,我很难说。
我没有收到过任何指向后者的别的见证。

另一件同样触目的事,在回到我这里的种种回声之中同样独一无二,就是提到我的名字与某段过去
时,那份热情的骤然涌现。那段过去,他早已决意宣告无效作废。还有他在这段过去里的根。而在那段
过去里,也还有一种童年的清新——那种被他逐出自己的“成人”生活、\dmark{}\origpage{1231}逐出这个重要而受人仰慕的
男子的生活的清新。在他周围,识相的人想必都不去提那段过去,不去提他还只是众多学生
中的一个、还怀着一腔美好激情的那些年月——正如在阔人宅邸里,满屋老样式的家具,谁也不去
提寒素甚至窘迫的起家……

而偏偏是这个陌生人,鬼使神差地坐在他身旁,忽然热切地谈起这件从来无人谈起的事(至少不当着
他的面谈……)!当然,这就像那团精选的、拘谨的空气骤然消散了;一个陌生人的热情,唤醒了他
身上同样的热情,并且——在一瞬之间——把他重新接通到那眼遥远的泉源,那眼他被认为早已忘
却、早已失去的泉源……

% ================================================================
%  18.8 Découverte d'un passé (原书 p.916–926)
%  源文: _work/chunks/chunk15.md
% ================================================================
\section{发现一段过去}

\subsection{(1) 第一口气息——或确认}   % 18.8.1

\notehead{183}{}

我终于来到这场一个多月前开始的回顾与盘点里最切身私己的部分。还剩要快速巡视的,是这场思
考就\textbf{我自己}教给我的东西。

思考让我发现的第一件事,是一段\textbf{过去}——我的数学家过去;此前我从不曾费心在上面停留
过哪怕一瞬。在那光鲜无事、看似平淡的表面之下,我又一次看见了深处的种种展露出来:它们通常被忽略、被抹去(仿佛一把扫帚悄悄扫过),从我们习惯为自己、为周遭描画的那幅舒舒服服的自觉图像里扫了出去。在那些从未检视过(至少在我的数学生涯里如此)的“涂痕”(或扫出门外的尘屑……)之中,有一种阴险、有时咄咄逼
人的东西:我在同某些朋友的关系里的自负。打一开始,这种自负就取了一种数学精英主义的形状;
它秘而不宣,我毫无察觉,因为我的态度看上去太天经地义了。这种精英主义(或如Chevalley与盖杰所
称的“绩优主义”),想必逐年变硬。它结晶为那种“竞技”的姿态——我是在思考“第一口气”接近
尾声时才终于发现的。在一副规规矩矩的门面底下,这种姿态认可的,是一种占有式的、嫉妒的心性,
指向那些被当作我的“禁猎场”的东西,以及那些凭其出色的品质、我乐于放进禁猎场来的人。

\pmark\origpage{1232}所幸,这些非常“老板”的心性,并没有穷尽1948到1970年间,我同数学世界里的朋友、同事
与学生,同数学本身之间关系的内容——差得远呢。尽管如此,它们构成一条阴险的低音线,去年以
前我从未费神去注意它;直到《收获与播种》的第一部分(“第一口气”)。这场渐进的发现,在“竞技
数学”一节(n\textsuperscript{o}~40)里达到顶点。这一节,在我看来标志着思考中一次质变的时刻。当时我立刻
把它感受为一次\textbf{翻越垭口}——它蓦地为我打开一片新的全景……

再退后一年回望,如今我把1948到1970年间、我在数学家中度过的这段漫长的数学生涯的第一
期,看作一场\textbf{交易}:以我本来所属(人人本属)的“长子权”——去完完整整地(倘若这是我
的选择)经历一场独特而独一无二的历险——换来的,是一碗“扁豆汤”:与一个田园诗般、虚构的
“数学共同体”认同(我希望这认同毫无保留,却从未完全做到……),而那个共同体同时又发放
种种舒服的好处\fnote{1027}{(*)}{这就是前面一条脚注(打字稿第1219页注(*))里谈到的那层暧昧。}。借这幅图像,我并不声称已把那段岁月说尽——它太丰厚,套不进一句现成
的公式。但我觉得,它抓住了第一次思考阶段里首次显形的一个重要侧面。这个侧面,又出现在《收
获与播种》这一部分(事后)取得的名字“愚妄与更新”里。

这第一阶段最切身、最深沉的部分,由最后三个“章”\fnote{1028}{(**)}{当然(如我在《收获与播种》引论里所说明的),这些“章”——把由共同主题或特殊亲缘联系起来的相邻各节归拢在一处——是事后设立的,在写完(那将只是)《收获与播种》第一部分之后。在《愚妄与更新》里,我偶尔把它们当作《收获与播种》的“部”来援引(不要把它们同那五个“部”——“愚妄与更新”等等——混为一谈;1984年2月至今的整个思考,归拢为那五个部)。}:第六至第八:“收获”、“孩子自娱”与“孤独的历险”。在“收获”里,我先同我生命中的某些
时刻(不只是我的数学生涯,这一回不是)重新接上头——一些满蓄更新之力的时刻。仿佛受一股
未知之力、一道隐秘而急迫的声音驱使,我寻找着重拾那些\textbf{天真}的心境\dmark{}\origpage{1233},好跨过那道我
隐约还感到横在前面的门槛。当然,当时我无从预料:接下来还等着我的,是发现我对数学本身
的一种占有式的态度。我继续爬坡,不急不缓,仿佛双脚踩着一条只有它们自己才“看得见”的无形
之路。不必对自己说,我也知道它引我去该去的地方;而脚下的雾,一步一步,渐渐散了。

就这样,我抵达了旅途中这道新的门槛,或者说这道\textbf{垭口}:

\begin{quote}
“……一到那个地方,我就有了那种来到瞭望台的人的感觉:从那里,他望得见自己一路走来的景
色——行走之时,他每一步只能看见其中一角。而此刻,有了这种对广度与空间的知觉,这是一
种解放……”
\end{quote}

刚一越过思考的这个敏感点,它便深化为一场对我自身的冥想。第二天,我已感到需要引入“老板”
与“工匠”的形象——别名“孩子”;这形象,两三年来对我已相当熟悉。但我远没有料到,它在接下
来的思考里将证明何等有用——那时,近两个月来我一直自以为马上就要收尾,好马上回到我的数
学笔记,回到《追寻栈》去!

在构成“孩子自娱”这一“章”的四节里,我重新接上了我同数学关系的某些侧面与波折。这些,近
三年前(1981年7月至12月)我已细细探测过,可此后有的是时间,大半又忘了。这一回,我的宗
旨主要是让自己进入状态,去探测我那次不期然而重返长期数学投入的意义,并试着在如今支配我
生命的两种激情之间——表面看互不相容——“给自己”定位:数学,与冥想。

顺便说,这两股激情的“互斥”,如今在我看来已不像两年前那样严酷了。在《追寻栈》里,
数学思考有时让位于、甚至化作一场多少个人化的思考;在那里,我这个有感受、有情感、有(不只
是数学的)好奇、有命运的存在,不再完全缺席。反过来,在\dmark{}\origpage{1234}《收获与播种》这场对我自身的思考里,思考本身又把我同旧日的数学之爱重新接通,并且这里那里,成了数学思考的引信
\fnote{1029}{(*)}{(5月10日)这些“思考的引信”,至今已结了果:我对若干搁置了十五年的主题的理解,得到了更新。}。

这些共存的可能——乃至共生的可能——这两股我心中求知冲动的不同表达,按事物的本性,或许
注定要相当有限。但无论如何,去年的思考(甚至更早,三年前那场长考)已让我看清:这两股
激情绝不互相敌对,甚至本质并无不同。在思考的最后一部分“孤独的历险”里,我竭力贴身把
握:这两股激情究竟在何处不同,它们各自为我打开的“历险”又如何不同。正是在这个追问
中,我发现了一个明摆着、而我一辈子佯装不知的事实:数学是“\textbf{一场集体的历险}”;我个人的
数学历险,只有系于它所属的那场更宏大的集体历险,才获得意义。

说实话,在“孤独的历险”一节里,起初我只是顺带掠过这个事实;当时的宗旨,倒是要用言词把
另一件事勾勒出来——那件事我素来深知,却始终难以全盘接受:冥想,才是一场\textbf{孤独的}\textbf{历
险}。给一件“已知”的事作表述,这番功夫绝没有白费,远非如此!它深化了这份认识,并且顺带让
我发现(至少对我而言)一个明摆着的新事实:把我系于\textbf{另一场}历险——数学的历险,它倒是
集体的——的那条纽带(当时,我、或者说当时的我之中的某人或某物,本想同它保持距离……)。

地基由此打好了:第二天,在“分裂的确认”一节里,我径直进入我困惑的中心。首先是确认:“老板”
的下注,哪怕他想自欺(那本更合他的天性……),也只能是那场集体历险——唯一能给他“回报”
的历险。“孩子天性孤独”;\dmark{}\origpage{1235}唯有孩子,才会被一场世上无人想要的历险、一种切实可触且常常明摆着
的知识所吸引,而这知识,他无法同任何人分享。而此刻,在我这桩“事业”里,“小
家伙的偏好”,偏偏不合“老板”的意,正落在这一边。

这个确认,引出对我之内一场\textbf{分裂}的确认:\textbf{老板与孩子的分裂}。第一次,我在如此专注、
如此严格的心境里,作出这样一个确认。这不是一道\textbf{判词}——依照这样那样的“看法”或哲
学或随便什么拟出、自命多少具有普遍效力的判词。它只是一个简单的\textbf{确认},出自对一个极特
殊的个案的细心检视:我这卑不足道的人,在我发展的某个阶段。也许有一天,这道分裂会消失,
而老板照旧去做必要的事,任工匠孩子随意做工。这不是我今天操心的事,也不必是。一天的难处一
天当就够了……

(4月5日)确实,这道分裂,差不多九年前,已在一个梦里,借一场以震撼人心的力量搬演出来的
寓言,向我揭示过。那是在发现冥想——这种长期被我无视、随时待命的力量——的两天之后;正
是掘到这个梦的意义的深处,我才重新找到我之内那个不分裂的东西,\textbf{另一个}我——沉默而
不可见了那么久的,“一位至亲之人,被人当作死了一辈子……”(见《愚妄与更新》)。当时出现的
新东西、要紧的东西,并\textbf{不是}那道我熟得不能再熟的分裂,也不是梦以那样大的力量向我揭示
的、这道分裂的本性——它化身成两个亲爱而熟悉的存在,谁也没有名字,而他们本是\textbf{同一
个}——而是那场\textbf{重逢}:在四个钟头高强度冥想(如同分娩的剧烈阵痛)之后来到的重逢。

我当时就知道,此后几天、几星期里也一直知道:重逢并不等于分裂的终结。但托重逢的福,我用
新的眼睛看这道分裂了——它固然要紧,但在另一个更根本的真实面前,终究是“次要的”:那是我
重新找到的我之内那个\textbf{不分裂}的、不可摧毁的统一,我后来认出它就是“孩子”。这一双重的
\dmark{}\origpage{1236}认识,当时鲜活而敏锐地在场。随后几年里,它钝化了——那份对这道“次要”却真真切
切的分裂的认识,渐渐被收藏起来。而当“老板”由着自己一头扎进冥想(那匹出名的“三脚马”……),
他很想暗示(却始终没有胆量、或没有笨到明说……):有了冥想和这一切,分裂从此已成过往,几乎
可以说没有分裂了——顶多这里一小道涂痕、那里一道,好吧,咱们不否认,但总归几乎跟没有一样;
只消看看小家伙工匠多么兴高采烈地埋头干活,看看一团和气的老板怎样踮着脚尖走路,生怕惊扰
了他——真可谓琴瑟和鸣!我不免疑惑:去年那份思考——翻越垭口(“竞技数学”)之前的那份,
尤其当我出人意料地回顾“我的激情”时(同名的节,n\textsuperscript{o}~35)——是不是恰恰还带着这样
的调子:光照逼得人对玫瑰色起一丝疑心……

不管怎样,这道“分裂的确认”当时恰到好处地让我重新接上了一个我多年渐渐看不见的真实。同
时,它让我以新的照明、新的眼睛,重新看见八年前我曾知觉得清清楚楚的那道分裂。我可以毫无保
留、毫无怀疑地这样说,因为我记得清清楚楚:作那个“确认”的时刻,同重逢的插曲、同那插曲就
某道分裂及其本性教给我的东西,没有任何联想!这一联想,直到稍后——我重新拾起前一天的笔
记时——才姗姗而来。这再清楚不过地表明:那次插曲中显现的认识的“次要”(而且不受欢迎!)
的内容,被收藏得多么干净。而且当年做来想必格外容易:重逢那道关键转折之后,没有人就这份
内容作过思考;而最能表达它的那幅形象——“老板”与“工匠孩子”——要再过好几年才出现。

如今我似乎看清:这道重新作出的“分裂的确认”,才是《收获与播种》第一部分里我对自己学
到的最要紧的东西。这个确认,不过几行字,是思考这一部分里最短的几节之一。\dmark{}\origpage{1237}人们也许
会想:既然只为了走到这一步,或许犯不着再用一百五十页,去追索自负在我数学生涯中的种种曲
折显形。就流行的“常识”而论,这千真万确。可同样真确的是:这种用修枝刀裁出来的“常识”,根
本无力把握一场发现工作的幽深曲折——无论是认识自己,还是更粗朴的数学发现的工作\fnote{1030}{(*)}{科学发现的工作,在我看来之所以比认识自己的工作“更粗朴”,(我想)有两个缘故。其一,它几乎只调用我们的智性官能,也就是我们存在中极小的一部分。(而且,科学工作倾向于让这部分官能增生肥大,以其余官能、以人格的整体平衡为代价,及至把人变成一种电子计算机怪物……)其二,科学工作所调动的、对抗揭示真实的内在抵抗,通常与对抗自我认识的抵抗,不成比例。这也是为什么,“科学历险”很少——今日几乎从不——是一场“求真的历险”:一场调用我们谦卑的能力、调用我们承担不期而至之真实的勇气的历险——首先对我们自己,然后对外部世界。}。我有一个切身的信念:在这场《收获与播种》的长考
里,每件事都来在其所、来在其时,由先来的一切所准备、所酝酿。

\subsection{(2) 第二口气息——或调查}   % 18.8.2

\notehead{184}{}

(4月6日)随着这道关于分裂的短短确认——去年3月末(一年多前),我起初以为自己已经写完
了《收获与播种》。我远没有料到,它还要再来五倍之多(按页数算)!随后几天,我做些杂事,思绪
开始回到数学主题。可是,思考中还有一个悬着的“小尾巴”,继续在我脑子里打转。越过一层看似
纯然细节的困惑,我隐约感到:对于“老板”转向一项长期数学投入的背后那股力量,我还没有真正
走完一遍。或者说,即便我把它的主要机关都翻了出来,我的理解仍旧苍白而游移——只因为在这
件事上“摆”得还不够久,没能让它渗得更深。这个“最后的尾巴”,将成为我回过头去看那团模糊印
象的门径。思考的重新接续,发生在当时(并且此后三个星期里)被当作《收获与播种》终篇的那一
节\dmark{}\origpage{1238},它旋即得到名字“一段过去的重量”。这个名字,说尽了那场意外的发现:我的数学家过去的
这份\textbf{重量},以及把我继续系于那场集体历险的纽带之力。而且,那天我所瞥见的,还只是
冰山露出的那一点小小峰尖;那沉在水下的庞大部分,要在随后的几个月、整整一年里,才渐渐
显形……

为思考的第一口气收尾的这一节,同时像第二口气的引信与召唤。这份“过去的重量”,显然扎根于
我对一份著作的眷恋——与其说眷恋那部完成了的、做到终点的著作,不如说眷恋一些执念观念与
一些洞见;它们的丰饶与力量,我感觉得到、我“熟知”入骨,而我多少含混地、多年以来就意识
到:它们却在贫瘠而干旱的土地上萎顿着,那土地暗地里、阴险地怀着敌意……于是,“一段过去的
重量”这场思考——它唤回我的好记忆:那部著作,以及我同它的纽带——成为一道长注的机缘;自
“离去”以来,我第一次就这部著作、就它遭到的命运表明态度。十年来、十五年来含混感到的东西,
终于成形,化作言词——言词有时迟迟不肯来到;而一旦白纸黑字写下,它们清清楚楚地向我道出一
个讯息,直到那时我一直在回避去知悉的讯息。后来,鉴于这条一气写成的长注的篇幅,我把它分成
两道,取名“我的孤儿们”与“拒受遗产——或矛盾的代价”(n\textsuperscript{o}s~46、47)。

\enlargethispage{2pt}
可以认为,这双注,就是\emph{葬礼}思考的发令枪\fnote{1031}{(*)}{可惜,这一情节在“葬礼(I)(或中国皇帝的新衣)”的目录里看不出:那道双注构成送葬行列之二(“孤儿们”),而非行列之一(“身后的学生”)。缘由在于“愚妄与更新”最后一节“一段过去的重量”(n\textsuperscript{o}~50)——那几道注本应是对它的评说——之中,诸“注”(n\textsuperscript{o}s~44至47)引用的先后次序。}。三星期后,4月19日,思考继续推进:那本“难忘文集”LN~900激起的
震动,为 Motif 的掘墓复生、由“新爸爸”Deligne主持,作了证。思考的这“第二口气”,一直高密度
地持续到5月底、6月中——当时我再一次自以为马上要画上句号,货真价实的句号!——以一场病\dmark{}\origpage{1239}
中插曲告终\fnote{1032}{(*)}{关于这场病中插曲,见“事件——或身体与精神”与“陷阱——或便利与衰竭”两条注(n\textsuperscript{o}s~98、99)。}。

这第二口气,严格说不是一场对我自己、对我过去的思考,而毋宁是一场对刚发现的\emph{葬礼}的
“调查”,同时是一场努力:把那些明摆着却(大概鉴于我那拔不掉的天真!)令人瞠目、令人难以置
信的事实,边走边尽量“消化”。不过,它毕竟就我自己教给了我一些东西:主要是让我切身掂出,我
对我的过去、对我的著作的眷恋之力。我被触到了痛处:眼看着著作被一块块撕走——有的进了垃
圾堆,有的供人取笑,有的被肆无忌惮地据为己有,当作谁都可以拿的小玩意……

那时我知道,我并没有像我一度相信的那样(在翻过某个“垭口”、眼前展开广阔全景之后
的那阵欢腾里\fnote{1033}{(**)}{这阵欢腾,表达于“旋转木马收摊了!”一节(n\textsuperscript{o}~41);五六星期后,在注“一脚踏进旋转木马”(n\textsuperscript{o}~72)里被压低了调门。})“离开旋转木马”。换句话说,我得以掂量这份过去的全部\textbf{重量},以及把我系
于其上的那些唯我机制的全体之力。这是个大意外!

然而,在思考的第二阶段,我关于自己还有另一项发现,补足第一阶段的所获。第一阶段,我翻出的
主要是我身上自负态度的某个“反面”:以\textbf{排斥}的姿态,对待某些同事甚至朋友——因为这样
那样的缘故,我不把他们算进我自觉身属(当然,是默而不宣的!)的“精英”的世界。同一枚奖牌
的\textbf{正面},则是我在同更年轻的数学家(尤其是我的学生)的关系里的一种\textbf{纵容}与暧昧:我
可以说把他们“吸纳”进了“我的世界”;或者因为他们手段出色,或者仅仅因为我收他们做了学
生,他们从此便被我知觉为处于我的“庇护”之下。我是在5月10日的注“高升”(n\textsuperscript{o}~63$'$)里开始
点到这种态度的,继而是5月27日的注“格格不入者”(n\textsuperscript{o}~67$'$),两道都写我同我年轻而\dmark{}\origpage{1240}光辉的
朋友皮埃尔的关系。这场思考在6月1日的注“暧昧”(n\textsuperscript{o}~63$''$)里深化,转向我同
我的学生们的一般关系。正是在那里,我终于察觉某种暧昧:它从未被我察觉、检视过,一直跟着我
到近几年。而且,就在最近,我又在稍有不同的语境里同这种暧昧打了个照面:在子注“驱逐(2)”
(n\textsuperscript{o}~169$^{1}$)(其第二部分,日期为3月16日)里。我不得不在那里确认:把我本人从 SGA 讨论班
中排挤出去——那个讨论班,是我十年生命投入的总和\fnote{1034}{(*)}{关于这场排挤最新的一桩,见注“殡仪馆——‘为科学服务’”(n\textsuperscript{o}~175)。}——这一排挤,主要由我旧日学生中最亲近的某几
位操办,不过是我在他们面前乐于经营的那种暧昧态度的自然收成:关于他们在那部规模恢弘的
SGA 里、与我在其中的各自位置——他们中的某一位,曾在其中投入一两年光阴。

\subsection{(3) 第三口气息——或发现暴力}   % 18.8.3

\notehead{185}{}

还剩要巡视的,是思考的“第三口气”就我自己教给我的东西。第三口气始于去年9月22日(病中
插曲之后),眼下正接近尾声\fnote{1035}{(**)}{《收获与播种》的第五部分,我在这里暂且撇开;它最初本是葬礼仪式内部——甚至就在“阴阳之钥”内部——的一段“枝蔓”。写这篇回顾时,这一部分尚未完成,故未收入。(6月22日)随后几星期里事情显明:“阴阳之钥”之后的“四种运算”部分(即葬礼(3)),构成《收获与播种》的“第四口气”,同样不在这篇最终回顾之列。}。这里首先指的,是在“阴阳之钥”里继续的那场思考;在我看来,
它是《收获与播种》最切身、最深沉的部分。没有任何刻意的安排,我的这个人、我与世界的关系,
在那里最常处于注意的中心。当注意力似乎偶尔岔开,转向一些表面上更一般的主题,或在我朋友皮埃尔其人身上流连,它却总是回到中心、回到那位演员兼观察者、那个感受着、知觉着、追问着、探
测着的人身上,仿佛被一股无形之力吸回。所以,先于一切、尽管并不以之为目标,它是一场\textbf{对我
的生命、对我自己的冥想}——从一个意外的入口进入:那场\emph{葬礼}。

这也是思考中在我看来最丰饶的部分,是我从中获益最多的部分。许多“已知”的东西彼此各安其
位\dmark{}\origpage{1241};一些仅仅瞥见或预感到的东西,一些“知道”却被忽视、被淹没在日常杂沓的昏暗光晕里的东西,
开始从暗处浮出,既显出它们的重量,也显出它们的轮廓。像是一次新的敞开,一次向未知再次出发
的邀请——而在此刻,那场有名的“长期数学投入”眼看要让认识自己的工作一放就是好几年……

在这里,不适合原原本本地重述这场长考的各个阶段,也不适合为它教给我的一切开一张“清
单”。我宁可几句话说说:对认识我自己而言——作为一场仍在日复一日、月复一月、年复一年继续
的成熟的材料——什么在我看来最要紧。

这场思考开始时,我本意是开一个“括号”(顶多一两道注的工夫……),好让读者(顺便也让我自
己)“入水”,浸进一种阴阳(或“阴柔-阳刚”)的事物辩证洞见。开这个括号的缘由,是需要用阴
阳的直觉,去定位一个触目的印象:它来自对某一篇“悼词”的检视\fnote{1036}{(*)}{关于这篇“悼词”(其恭维之语,调配得法,施用有道……),见同名的两条注(n\textsuperscript{o}s~104、105),以及部分概括了它的注“宝石们”(n\textsuperscript{o}s~170(iii))。}:那是一种刻意的“角色对调”,发
生在一场原本的阴阳关系里。括号于10月2日打开。直到11月10日——在关于阴阳的把戏在我
的生命里、在一般的存在里、最后在数学发现的把戏里的整整一百页密不透风的思考之后——时
机才算成熟,可以至少把这个六个月前冒出的联想\textbf{说出来}\fnote{1037}{(**)}{在注“对调(3)——或阴葬阳”(n\textsuperscript{o}~137)里。};至于在知根知底的状态下去探测它,还要再过十四
天\fnote{1038}{(***)}{在注“天鹅绒的脚掌——或微笑”(n\textsuperscript{o}~137)的开头。}。(而著名的阴阳括号,要到将近两个月后、1月14日,才终于合拢——甚至过了一阵,我才察觉它早已合上了……)

\pmark\origpage{1242}很快,无须寻找也无须计划,“冲突”——在人的生命里、在人身上的冲突——占据了注意的
中心。\emph{葬礼}的发现骤然而强力地调动起来的唯我能量,在那里充当了一支援兵,让我再一次、
就在当下,直面多年以来一直向我发问的那个“冲突之谜”\fnote{1039}{(*)}{这一“发问”,主要开始被察觉,是在我1979年8月至1980年3月间持续的那场对我父母的漫长冥想之中。}。早在此前那些年里,这个谜已渐渐浮到我
想探测、想理解之事的前排,能走多远走多远——却从来没有“跨出那一步”、把自己整个儿投进去
……

渐渐地,思考之中显形出来的,是我生命里仿佛那枚“硬核”的东西,这个谜那可怕的震中,仿佛
“恶之谜”的心脏本身:一种可以称作“无缘无故”或“无因”的暴力——一种似乎单以伤人、害人、
毁人为快的暴力,一种从不报出姓名的暴力,常常闷声不响,顶着一副天真无辜、和蔼可亲的面孔,
因而在触及与摧残上愈发有效——“天鹅绒里的利爪”,纤细、敏捷、毫不留情……在“天鹅绒里的
利爪”那组注(n\textsuperscript{o}s~137--140)的思考中,注意最终落到的,正是这种暴力;直到《阴阳之钥》
终卷,它也一直处于注意的中心。它构成那把最高音:最后一道注里,那条“业”的“没有尽头的锁链”
——父母传给子女,子女传给孙辈,一代又一代,自太初以来——被唤起。

这是我一生中第一次,直面这个“无恨亦无情”的暴力之谜——一种深植于人类生活的暴力,自我
少年时代起,就在我生命里打下了洗不掉的印记。这也是我第一次,在我自己的存在里,确认这道印
记。这同时也是对一个简单事实的确认:这种暴力的\textbf{存在},它那可怕的无所不在——在我的生
命里,也在每一个人的生命里\fnote{1040}{(**)}{这一确认,构成注“无恨亦无容赦”(n\textsuperscript{o}~157)那场思考的高峰时刻。}。而这一个简简单单的确认,本身就在胚胎里含着对这个可怕事实的一
种\textbf{接受}。或许正是在这个确认里,藏着整部《收获与播种》的思考中我所\dmark{}\origpage{1243}学到的最重要的东西
(或者说,至少\textbf{开始}学到的东西)。

这并不是一个终局,不是一场思考的顶点。毋宁说,这仍是第一步,把我带过一道通向未知的门
槛。对我的行路与成熟而言,这个谦卑的步子,在我看来,分量比我所瞥见的那些“回答”的胚胎更
大——在随后几天里,我瞥见了“无因的暴力”之“因”这一问题的答案的胚胎\fnote{1041}{(*)}{见同名之注(n\textsuperscript{o}~159)。}。这个问题本身,只有
在充分看见并承认为其发问的那个事实的存在与分量之后,才获得它的全部意义——那分量,远比
一个“心理机制”的问题沉重。

有人会说,我这是跑题了:对一个(或我自以为)关乎人人的普遍心理事实的确认,属于客观知识的
领地,那是各门科学的专属(心理学、精神病学、社会学,或随便什么),不属于那出了名的“自我认
识”的领地(一种被感觉为含混而不可捉摸、倘若不算全然想入非非的领地)。可我看得很清楚——不
是含混不可捉摸地看,而是像看一个熟悉的、明摆着的数学事实一样清楚……:脱开认识自己,这样
的确认就失去了它活的意义——失去那使它不至于沦为一场哲学-心理学的文体练习、沦为一条
“论纲”的展开(当然十分有意思,如此等等……)的东西。这个确认,本身是一场\textbf{发现},一场贴
身的发现:世上没有任何人能替我作这个发现,我也无法替世上任何别人作它。这场发现,是那次发
现自己的旅程中,最近的一步——或者几乎是最新的。它把我相对于一件要紧而可怕的事定位:这件
事打过我烙印,而我直到那时一直执意忽视它,仿佛我是由于某种特殊的倒霉(也许系于我这卑不
足道的人的某些特点),才一生暴露于它;仿佛别人暴露于它、或施它于人——只要我肯费心睁眼看
看周围——也无非同样的倒霉。

而且这肯定不是偶然:对暴力的思考刚一起头,思考自身的内在逻辑,就引着我——也是平生第一
次——回头去看我记得的那几桩个案:我自己也曾\dmark{}\origpage{1244}\textbf{施}这种“超出理解的暴力”于人,当然,不
曾多想\fnote{1042}{(*)}{见《阴阳之钥》里“天鹅绒里的利爪”部分之后的那道注“正义者的暴力”(n\textsuperscript{o}~141)。}。这次回望的意义,不在于给我机会自责(还是在众目睽睽之下!)——这种事,我反倒完
全没有去做。而在于它为我打开了一扇门,通向对暴力更深的理解——一扇门,往后的日子里,何时
跨过它,只由我定。

\subsection{(4) 忠诚——或女性气质的数学}   % 18.8.4

\notehead{186}{}

这就是,以发现自己的旅程的眼光看,在我看来最要紧的东西。阴阳思考的最后阶段——以暴力为
中心——贯穿最后四部分:“天鹅绒里的利爪”、“暴力——或游戏与刺棒”、“另一个自己”、“冲突与
发现——或恶之谜”,从12月7日到1月14日(它们略多于《阴阳之钥》的三分之一)。

退后看去,我似乎觉得:《阴阳之钥》前八部分的主要角色,是最终把我带到这场关键的思考。我在
这准备性的部分里展开的许多东西,是多年以来我熟悉的事;却仍须“重提”,好让一位初来的读者
能够跟上,好给思考一种内在的融贯——否则它就有缺失之虞,或只有我一个人看得见。不时地,文
体显出这些内在心境的痕迹:写的人急着赶紧把重提打发完,好终于抵达“正题”——虽然,那些所谓
的重提,常常分量重得多,值得我在上面多停留片刻,比我那么急着要去的那个“正题”更值得(而
无论急不急,我到那“正题”,还是晚了一个多月……)。这些心境,我看在三节相连的部分里最为
可触:“夫妻”、“我们的母亲死神”、“拒与受”。即使在那里,真的,重新接上那些本该“已知”的事
情时,我仍旧不由自主同时重新认出了它们,而且有时是在新的光线下——哪怕是为了一些初看如
此非个人化的事情,比如清点那些由直接的亲缘联起来的阴阳“夫妻对”(或“锁孔”)——那些通往
世界的门。

但正是在接下来的三部分里(它们也正好先于以暴力为中心的最后四部分),我\dmark{}\origpage{1245}重新驶入此前未曾
勘探的岸域:“阴阳数学”、“阴阳的对调”、“主人们与仆人们”。

正是在这三部分的第一部分里,发生了那场“大意外”:它将为\emph{葬礼}的意义——或至少某种意
义——投下新的光。事情是这样的:在我对待数学的进路里,更一般地,在我走向发现世界的自发举
步里,我的存在的基调是\textbf{阴,“阴柔的”}。换个说法:尽管我的自我的条件化结构、我这桩事业
的“老板”是阳(不说“彻头彻尾的男子汉”吧),我本来的天性、我之中的“孩子”(他同时是那个
把孩子在玩耍中发现的东西做出形状的工匠……)却以“阴柔”为主。当然,单是这个特点,并不足
以把我个人对待数学的“风格”同任何别人区分开来。我甚至觉得,即使在数学家当中,这个基调(或
“主导”)本是阴的,也并不那么罕见。而我的情形里反常的(在我看来)倒是:在我的发现步履里、
尤其在我的数学工作里,我一生都对这本来天性完全忠诚,从不曾起过一丁点儿要修改它、纠正它的
念头——无论是出于内心检查员的意愿(反正他从来什么也没看见:在一门“男人之间”的事情里,
谁会怀疑一种“阴柔”的感受力与创造性的进路呢!),还是出于迁就外部世界、尤其是科学世界里流
行的好品味的法度。我毫不怀疑:主要正是靠着对我自己天性的这份忠诚——至少在我生命中这个
有限的领域里\fnote{1043}{(*)}{如我在《收获与播种》各处有机会一再说起的,自我八岁起(直到1976年,我四十八岁)支配我生命的两股最强的唯我之力之一,是对我身上“阴柔”特质的压抑,以利于那些被感受为“阳刚”的特质。只是到了“阴阳之钥”的思考里,我才意识到:这种压抑,并没有施行在我的数学工作上(后来,也没有施行在冥想、认识自己的工作上)。我存在中本来“阴柔”的主导,得以在一项通常被感受(也难怪)为最“阳刚”的活动里,尽情尽兴!(关于这一点,见注“最‘男子汉’的艺术”,n\textsuperscript{o}~119。)},我的数学创造力才
得以充分地、无阻碍地铺展,像一棵茁壮的树,稳稳扎根在大地里,依着昼夜、风与季节的节律自由
舒展。事情是这样成全的,尽管我的“天分”颇为平平,尽管开端根本算不上\dmark{}\origpage{1246}\textbf{吉兆}\fnote{1044}{(*)}{我说“天分平平”,绝非故作谦虚。这件事我一遍遍有机会确认:无论是同那些光辉的数学家打交道——他们把握要领、接受并消化新思想,比我敏锐得多——还是同这样一些没有名姓、没有像样数学底子的学生打交道,他们的好奇心与数学发明力一时被唤起。
我的“开端”(至少是我同数学家世界接触的开端,1948年),我在“受欢迎的异乡人”一节(n\textsuperscript{o}~9)里谈了一点。不过,还要早三年,1945年,我的“数学生涯”就开始了:我的精力的大半,投给了数学研究工作。直到1949、1950年前后,作为在法国的一个外国人,我以数学谋生的前景,还成问题得很。倘若这样的机会没有出现,我打算去学木匠——那是一门合我口味的营生。}。

当我在注“上涨的海……”(n\textsuperscript{o}~122)里,对我的数学进路作出这项出人意料的确认时
\fnote{1045}{(**)}{另见后文注“箭与浪”(n\textsuperscript{o}~130)。},它来得有点像一桩意外的奇事,有点“在边上”——因为在我的生活里,同他人的
关系,全都打着我阳的、超级阳的抉择的印记。只是在这场以冲突的动势为中心的思考的继续里,
借着对\emph{葬礼}的一次回望,我才意识到:我的数学家同侪同我的这个人、尤其同我的著作的关
系,被打上了这个不寻常的特点何等深的烙印——在他们心里唤起矜持的反射(若不是排斥),面对
一种被暗暗感受为“不合适”(不说“失体”)的进路风格。这类反应,早在我初入数学世界时就很有
些;只是在那宽厚的年月,它们被当时盛行的敬重他人的氛围调和着——至少在我落脚的那些数学
圈子里。后来,鉴于“Grothendieck结果之力”(借Borel致Mebkhout信中的话,那封信里提到了这些
“矜持”),它们也只好被压着。而自我告别数学舞台以来,它们成了常规,有时还在某副得体的低声
调门(那种低声调门,至今仍是规矩)背后舒舒服服地铺展:昔日的敬意早已风化、消失多时,而当事
人(据推定已死且已葬)也不在场,不能接话了……\emph{葬礼}的这个出人意料的面向——以我这卑
不足道的人为象征,埋葬“数学中的阴柔”——在“天赐的机缘——或封神”与“拒认——或召回”两
道注(n\textsuperscript{o}s~151、152)里得到探测;那是12月23日与24日,正在关于暴力的冥想的正当中。

\pmark\origpage{1247}还剩我这个人最后的一个面向想谈一谈——它是在写《阴阳之钥》时、在刚引的最后那一部
分“主人们与仆人们”里(那一部分紧挨着由“天鹅绒里的利爪”开启的思考转折之前)显形的。这
就是“服务的冲动”,以及它在我对数学投入的抉择里扮演的头等角色,以及它作为一股活的力,在
那些浩大而没有尽头的奠基任务里运转——我之后,还没有别人找到勇气(或谦卑……)去接手并把
它们继续下去。这个以非凡之力在我身上在场的面向,雄辩地见证了我本来天性中“阴柔”的主导:
它保存了下来(乃至避难……)在数学活动里(谁也不会想到去那里找它……)。

此刻我心中闪过一个念头:说不定,正是这股冲动,也以它非唯我的那一份,参与了那次“转向”——
转向高强度的数学活动,而把冥想的功课撂到了后台,一放就是不知多久。冥想,按其本性,是一项
孤独的功课;一项(在我看来)除非自欺、便无法纳入“为众人服务”、为某个“渴望知道的存在的理
想共同体”服务的投资眼光的功课。这么看来,似乎有一股深层的冲动,不同于唯我的求证之心、求
许之心,一股表达着个人与其所属的种属之间的深层纽带的冲动——它注定要在一场旷日持久的冥
想功课里受挫,按我对冥想的理解。而这也许是一条附加的缘由——除了那些来自唯我结构(“老
板”的心性)的缘由(它们本身已经足够强大)——使得这样的功课显得如此稀罕,以至于我拿不
准,可曾在别人身上见过它的痕迹。

% ================================================================
%  18.9 De Profundis (原书 p.926–929)
%  源文: _work/chunks/chunk16.md
% ================================================================
\section{自深处}

\subsection{(1) 感恩}

\notehead{187}{}

(4月7日)这场回顾与盘点——盘点全部《收获与播种》的思考教给了我什么——我相信已经走完
了一遍。唯有一处我从这份回顾里剔除出去,即《收获与播种》的第五部分\fnote{1046}{(*)}{(6月22日)还有第四部分(我眼下正在写的)也一样!见打字稿第1240页脚注(**)。},它至今没有写完。它最初本是“阴阳之钥”里的一段“枝
蔓”,后来延宕了整整一个月,落实为关于 C.G.~Jung自传的一百来页“读书笔记”。由于
这条枝蔓的收尾仍\dmark{}\origpage{1248}遥遥无期,我便把它搁置了。我最盼望的,是把\emph{葬礼}善始善终:写完、打
字、印出、四处寄发——然后,再没有人提它!

何况我预感到,这第五部分将就这同一个\emph{葬礼}给我带来意想不到的光亮——可不是嘛!——
通过我打算在其中对Jung与Sigmund~Freud之间关系的考察:有多年之久,弗洛伊德曾是尚
在寻找道路的年轻Jung的师父。初读自传里论及这层关系的那一章,我什么也没看出来——随后,
若干蹊跷之处引起了我的注意,我折回去看了几处,又把这一章重读了一遍。显然,这层关系绷
满了暧昧——弗洛伊德本人似乎强烈地感到了它,而Jung却津津有味地全然无视它(活脱脱像个
初来的讨论班新生……),把弗洛伊德的不安一概记在他那“神经症”的账上(他还饶有兴致地
把这条“神经症”描绘得有声有色,也许甚至绘声绘色得过了头,反而不大像真的……)。话虽如
此,种种联想随之而来,系于我的朋友、(同时也是)非学生Deligne与我的关系之上;这些联想我
打算追下去,兴许还要稍加开掘。我隐隐觉得,就所牵动的心理机制而言,\emph{葬礼}上发生的事,
绝非一桩绝无仅有、不合典型的巧合——恰恰相反!而我预感到,Jung与弗洛伊德的那层关系,
在这方面没准能提供更多的光亮。

不过对我而言,至少眼下,这第五部分(它也许会叫“Jung——或一场冒险的陷入泥沼”\fnote{1047}{(*)}{本想写“陷入泥沼”(enlisement),却看见自己写出了“葬礼”(enterrement)。这个口误提示的新名字——“Jung——或一场冒险的葬礼”——未必不贴切,甚至没准比原先拟的那个更切中要害。}),已经不再是
那场\emph{葬礼},哪怕它从中生出——我甚至要说:它已不再是《收获与播种》!它是“\textbf{此后}”
(l'Après)——同种种回声一样,那些回声,随着“愚妄与更新”、“葬礼(I)——或中国皇帝的新衣”
和“葬礼(III)——或四种运算”这三部分的寄出,势必要纷纷回到我这里,其中肯定少不了一些
尖酸刻薄的话\fnote{1048}{(**)}{不要把《收获与播种》的第四部分(副题为“四种运算”),与同在该部分之内、集结在这一名目之下
的那组笔记(n\textsuperscript{o}s~167$'$--176$_{7}$)混为一谈。}。到这里,等这部分打完清稿,就已经是一千来页,甚至更多——已经够多的
了!一天的难处一天当就够了……

\pmark\origpage{1249}这种赶紧收场、“把它寄出去”的急切,无疑首先是一匹嗅到火药味的战马的急切,急不可
耐地要冲进战阵\fnote{1049}{(*)}{类似的心绪,在“愚妄与更新”最后一节“一段过去的重量”(n\textsuperscript{o}~50)里已经谈过,只是视角稍有不
同:在那里,冲进战阵的不是战马,而是一头公牛,追赶着一块在它“鼻子跟前晃来晃去”的红布……}。但也许还有更深一层的,是一种愿望:巴望某一段过去从此脱开我而去。这一千来页,以惊
人的方式把这段过去的全部\textbf{重量}凝成了实体——把这项工作做到终了,直到最后一桩庶务
(最后一桩,想必就是把《收获与播种》寄给我寄发清单上预先开列的一百三十位收件人……\fnote{1050}{(**)}{到那时,那著名的“重量”将变得更加“惊人”:一上来就是二十万页(200×1000),而不是一千页!}),在我
看来,几乎出于本能,也正是我把它\textbf{抛掉}的时刻。幻觉?唯有将来知道……

就这样,我来到了那些“终止和弦”面前——在那个有名的“句号”之前;那个“句号”,一年多
来我一直以为它就在眼前,却一天天、一星期一星期、一个月一个月地被推远,总有料想不到的
东西涌来,要占一席之地。

还剩什么要说,在这些终止和弦里?有一份感激,它要借“感谢”说出。这场思考是孤独的果实,
然而我得到过许多方式的帮助。

最显而易见的帮助,来自Zoghman~Mebkhout,同样是以许多方式:凭他的耐心,他领我“入了门”,
走进“好上帝-Mebkhout定理”周身的哲学;凭他对我的信任,不管不顾地把他在我昔日学生们那里
遭遇到的困难与挫折向我和盘托出;凭他帮我厘清一片枝蔓丛生的数学文献——我早已同它断了
往来;最后,凭他自打得知这项工作那一刻起,就对它投以无私而友善的关注——他在其中看见
我倾注了身心,而首先(我相信),他认出并接纳了其中的\textbf{见证}。

我同样感激Pierre~Deligne,他曾专程前来看我,并(去年10月)读阅了当时已写出的\emph{葬礼}
部分,向我提出意见\dmark{}\origpage{1250}\fnote{1051}{(*)}{关于这次造访以及Deligne带给我的种种澄清,见“葬礼(III)”中构成“最后的义务(或造访)”一节的那两条注(n\textsuperscript{o}s~163、164)。}。这次造访对我多有助益。

最后,帮助过我的,还有朗格多克科学技术大学那些为我打字的秘书们的好意与友善氛围:Boulet
小姐(Mlle~Boulet)、Boucher夫人、Brun小姐、Cellier夫人、
Lacan小姐、Mori夫人。她们中有两位,不惜动用个人时间,为按时交
件承担了一部分打字工作,还不肯为此收取报酬——这份情义令我深深感动。此外,整部《收获
与播种》笔记的后一半,是Lacan小姐一个人打出来的,其用心与效率堪称典范。对她们每一位,
我都乐于在此表达我的感激。

我还想着所有那些人,所有那些女人和男人——在我工作的许多时刻,他们曾以常常是不受欢迎
的方式\fnote{1052}{(**)}{这几个月的笔记里,这里那里曾提到这类“打扰”。关于这一点,特别见注“信使(2)”(n\textsuperscript{o}~181)。}打断我的工作与安宁。当然,这些“打扰”本身——它们在某些时刻考验过我,有一些至今还在
我心里留着悲哀的残余——在我的工作里同样有它们的角色要扮演,同样带来一个讯息,听不听、
化不化,全在我自己。当悲哀或怨怼化为感激,我便会知道,这个讯息已被接纳……

\subsection{(2) 朋友}

\notehead{188}{}

\emph{葬礼}的这些终极和弦,近一年来早已名字现成:\emph{自深处}!(De~Profundis)在《引论》
(I~7,“葬仪的安排”)里,我甚至走得更远,竟(也许失之轻率……)预告说,死者那“完全的满
足”将构成“那场难忘葬礼的终曲与最后的和弦”。写下这句预言时(仿佛事已办讫),我是可以
原谅的——写那几行的时候(去年5月),这看上去确实是一个为期极近的预言,因为我自以为正
要抵达“自深处”的那些终止和弦。

诚然,比去年更切身地(当时,思考的“第二口气”行将告一段落),我意识到\dmark{}\origpage{1251}自己离真正“走
完”这葬礼还有多远——若把纯然物质性的事实除外的话(那些事实,我自觉得已经“尽在掌握”,
颇有几分自负\fnote{1053}{(*)}{(5月10日)然而,这几行写下之后,又有一个多月过去,耗在把新冒出来的事实勉力“塞进去”,临
时添了足足二十来条子注!})。倘若确实像我不时感到的那样,理解\emph{葬礼},也就意味着“理解冲
突”,那么很可能,我余生的岁月将不足以完成这番“走完”——至少走不到深处。

如此说来,我可以说,我此刻写下这条终极之注时的心境,已全然不同于写《\emph{葬礼}引论》时
的心境。这是否意味着,我是在“完全的满足”之情不在场的情况下,为这场思考收尾的?

我不以为然。愿景一旦深化,某项曾催生这愿景、为这深化作过铺垫的工作——它本可以被看
作“已臻完成”——便会在一个“彼岸”显现之时,显出自身的\textbf{未完成}。然而,工作的\textbf{意
义},以及它让我们尝到的满足或不满足,并不在它的完成之中,也不取决于这项工作注定完
成与否。工作的意义就在工作本身,就在\textbf{当下此刻}——在我们做事时的心境里,在我们倾
注其中的爱里(或爱的缺席里……)——而不在一个我们无从企及的假想的将来。

去年3月,甚至在我发现\emph{葬礼}之前,我在《引论》(I~1,“梦与实现”,打字稿第~iv~页)里写道:

\begin{quote}
“……我离开这项工作,带着一个知道已把工作做到终点的人的完全满足。没有任何一件事,
无论它多么‘微小’,被我避而不谈;没有任何一件事,是我念念于心要说、却终究没有说的,
以至于在此刻,会在我心里留下一星不满、一星遗憾的残余,无论它们多么‘微小’。”
\end{quote}

如今我知道,这份我原以为“做到终点”的工作,今天仍未到终点,也许永远到不了。但我也知
道,这毕竟是一件次要的事。那份“完全的满足”——在我写下这些力图贴它最近的话语的那一刻,
我曾强烈地感到过它——它一路伴随我写完了\dmark{}\origpage{1252}《收获与播种》。它是我的一位老朋友,早在
我数学生涯的全部岁月里,它就陪伴着我,低声告诉我:路走对了。后来,在冥想的功课里,我
又找到了它——正是同一位。

当我听不见它,工作便失去意义。正因如此,它的声音于我弥足珍贵,也正因如此,我在工作时
小心在意,从不让自己远离它。靠了这个,工作在我一生中始终是喜悦的源泉——那个全然以身
相许者的“完全的满足”。

在即将结束的这项工作里,情形并无二致——这项工作,即是“\textbf{收获}”,同时又是
“\textbf{播种}”。


\end{document}
